\documentclass[../../../main.tex]{subfiles}
\begin{document}
\subsection{Optimization For XXX Hamiltonian: No Treatment Whatsoever}
With known eigenvalue and eigenstates
\begin{gather*}
    \ket{\tilde{m }} = a_m \sum_{j=1} ^{N} \cos \left[ \frac{\pi }{2N }(m-1 ) (2j-1)\right]  \ket{\mathbf{j}}\\
    E_m = 2B + 4 J (1 - \cos \pi(m-1)/N)
\end{gather*}
the fidelity may now be computed
\begin{equation*}
    \braket{F }  = \frac{1 }{2 } + \frac{1 }{3} \mathrm{Re} (f_r) +\frac{1 }{6 }|f_r|^2
\end{equation*}
with transition amplitude
\begin{equation*}
    f_j(t) = \bra{\mathbf{j}} e ^{-iHt/\hbar }\ket{\mathbf{s}}
\end{equation*}

\begin{figure}
    \centering
    \resizebox{\textwidth}{!}{%% Creator: Matplotlib, PGF backend
%%
%% To include the figure in your LaTeX document, write
%%   \input{<filename>.pgf}
%%
%% Make sure the required packages are loaded in your preamble
%%   \usepackage{pgf}
%%
%% Also ensure that all the required font packages are loaded; for instance,
%% the lmodern package is sometimes necessary when using math font.
%%   \usepackage{lmodern}
%%
%% Figures using additional raster images can only be included by \input if
%% they are in the same directory as the main LaTeX file. For loading figures
%% from other directories you can use the `import` package
%%   \usepackage{import}
%%
%% and then include the figures with
%%   \import{<path to file>}{<filename>.pgf}
%%
%% Matplotlib used the following preamble
%%   
%%   \makeatletter\@ifpackageloaded{underscore}{}{\usepackage[strings]{underscore}}\makeatother
%%
\begingroup%
\makeatletter%
\begin{pgfpicture}%
\pgfpathrectangle{\pgfpointorigin}{\pgfqpoint{10.000000in}{6.000000in}}%
\pgfusepath{use as bounding box, clip}%
\begin{pgfscope}%
\pgfsetbuttcap%
\pgfsetmiterjoin%
\definecolor{currentfill}{rgb}{1.000000,1.000000,1.000000}%
\pgfsetfillcolor{currentfill}%
\pgfsetlinewidth{0.000000pt}%
\definecolor{currentstroke}{rgb}{1.000000,1.000000,1.000000}%
\pgfsetstrokecolor{currentstroke}%
\pgfsetdash{}{0pt}%
\pgfpathmoveto{\pgfqpoint{0.000000in}{0.000000in}}%
\pgfpathlineto{\pgfqpoint{10.000000in}{0.000000in}}%
\pgfpathlineto{\pgfqpoint{10.000000in}{6.000000in}}%
\pgfpathlineto{\pgfqpoint{0.000000in}{6.000000in}}%
\pgfpathlineto{\pgfqpoint{0.000000in}{0.000000in}}%
\pgfpathclose%
\pgfusepath{fill}%
\end{pgfscope}%
\begin{pgfscope}%
\pgfsetbuttcap%
\pgfsetmiterjoin%
\definecolor{currentfill}{rgb}{1.000000,1.000000,1.000000}%
\pgfsetfillcolor{currentfill}%
\pgfsetlinewidth{0.000000pt}%
\definecolor{currentstroke}{rgb}{0.000000,0.000000,0.000000}%
\pgfsetstrokecolor{currentstroke}%
\pgfsetstrokeopacity{0.000000}%
\pgfsetdash{}{0pt}%
\pgfpathmoveto{\pgfqpoint{1.250000in}{0.660000in}}%
\pgfpathlineto{\pgfqpoint{9.000000in}{0.660000in}}%
\pgfpathlineto{\pgfqpoint{9.000000in}{5.280000in}}%
\pgfpathlineto{\pgfqpoint{1.250000in}{5.280000in}}%
\pgfpathlineto{\pgfqpoint{1.250000in}{0.660000in}}%
\pgfpathclose%
\pgfusepath{fill}%
\end{pgfscope}%
\begin{pgfscope}%
\pgfpathrectangle{\pgfqpoint{1.250000in}{0.660000in}}{\pgfqpoint{7.750000in}{4.620000in}}%
\pgfusepath{clip}%
\pgfsetbuttcap%
\pgfsetmiterjoin%
\definecolor{currentfill}{rgb}{0.000000,0.000000,0.000000}%
\pgfsetfillcolor{currentfill}%
\pgfsetlinewidth{0.000000pt}%
\definecolor{currentstroke}{rgb}{0.000000,0.000000,0.000000}%
\pgfsetstrokecolor{currentstroke}%
\pgfsetstrokeopacity{0.000000}%
\pgfsetdash{}{0pt}%
\pgfpathmoveto{\pgfqpoint{1.262400in}{0.660000in}}%
\pgfpathlineto{\pgfqpoint{1.268600in}{0.660000in}}%
\pgfpathlineto{\pgfqpoint{1.268600in}{5.280000in}}%
\pgfpathlineto{\pgfqpoint{1.262400in}{5.280000in}}%
\pgfpathlineto{\pgfqpoint{1.262400in}{0.660000in}}%
\pgfpathclose%
\pgfusepath{fill}%
\end{pgfscope}%
\begin{pgfscope}%
\pgfpathrectangle{\pgfqpoint{1.250000in}{0.660000in}}{\pgfqpoint{7.750000in}{4.620000in}}%
\pgfusepath{clip}%
\pgfsetbuttcap%
\pgfsetmiterjoin%
\definecolor{currentfill}{rgb}{0.000000,0.000000,0.000000}%
\pgfsetfillcolor{currentfill}%
\pgfsetlinewidth{0.000000pt}%
\definecolor{currentstroke}{rgb}{0.000000,0.000000,0.000000}%
\pgfsetstrokecolor{currentstroke}%
\pgfsetstrokeopacity{0.000000}%
\pgfsetdash{}{0pt}%
\pgfpathmoveto{\pgfqpoint{1.270150in}{0.660000in}}%
\pgfpathlineto{\pgfqpoint{1.276350in}{0.660000in}}%
\pgfpathlineto{\pgfqpoint{1.276350in}{4.881179in}}%
\pgfpathlineto{\pgfqpoint{1.270150in}{4.881179in}}%
\pgfpathlineto{\pgfqpoint{1.270150in}{0.660000in}}%
\pgfpathclose%
\pgfusepath{fill}%
\end{pgfscope}%
\begin{pgfscope}%
\pgfpathrectangle{\pgfqpoint{1.250000in}{0.660000in}}{\pgfqpoint{7.750000in}{4.620000in}}%
\pgfusepath{clip}%
\pgfsetbuttcap%
\pgfsetmiterjoin%
\definecolor{currentfill}{rgb}{0.000000,0.000000,0.000000}%
\pgfsetfillcolor{currentfill}%
\pgfsetlinewidth{0.000000pt}%
\definecolor{currentstroke}{rgb}{0.000000,0.000000,0.000000}%
\pgfsetstrokecolor{currentstroke}%
\pgfsetstrokeopacity{0.000000}%
\pgfsetdash{}{0pt}%
\pgfpathmoveto{\pgfqpoint{1.277900in}{0.660000in}}%
\pgfpathlineto{\pgfqpoint{1.284100in}{0.660000in}}%
\pgfpathlineto{\pgfqpoint{1.284100in}{4.877283in}}%
\pgfpathlineto{\pgfqpoint{1.277900in}{4.877283in}}%
\pgfpathlineto{\pgfqpoint{1.277900in}{0.660000in}}%
\pgfpathclose%
\pgfusepath{fill}%
\end{pgfscope}%
\begin{pgfscope}%
\pgfpathrectangle{\pgfqpoint{1.250000in}{0.660000in}}{\pgfqpoint{7.750000in}{4.620000in}}%
\pgfusepath{clip}%
\pgfsetbuttcap%
\pgfsetmiterjoin%
\definecolor{currentfill}{rgb}{0.000000,0.000000,0.000000}%
\pgfsetfillcolor{currentfill}%
\pgfsetlinewidth{0.000000pt}%
\definecolor{currentstroke}{rgb}{0.000000,0.000000,0.000000}%
\pgfsetstrokecolor{currentstroke}%
\pgfsetstrokeopacity{0.000000}%
\pgfsetdash{}{0pt}%
\pgfpathmoveto{\pgfqpoint{1.285650in}{0.660000in}}%
\pgfpathlineto{\pgfqpoint{1.291850in}{0.660000in}}%
\pgfpathlineto{\pgfqpoint{1.291850in}{4.553155in}}%
\pgfpathlineto{\pgfqpoint{1.285650in}{4.553155in}}%
\pgfpathlineto{\pgfqpoint{1.285650in}{0.660000in}}%
\pgfpathclose%
\pgfusepath{fill}%
\end{pgfscope}%
\begin{pgfscope}%
\pgfpathrectangle{\pgfqpoint{1.250000in}{0.660000in}}{\pgfqpoint{7.750000in}{4.620000in}}%
\pgfusepath{clip}%
\pgfsetbuttcap%
\pgfsetmiterjoin%
\definecolor{currentfill}{rgb}{0.000000,0.000000,0.000000}%
\pgfsetfillcolor{currentfill}%
\pgfsetlinewidth{0.000000pt}%
\definecolor{currentstroke}{rgb}{0.000000,0.000000,0.000000}%
\pgfsetstrokecolor{currentstroke}%
\pgfsetstrokeopacity{0.000000}%
\pgfsetdash{}{0pt}%
\pgfpathmoveto{\pgfqpoint{1.293400in}{0.660000in}}%
\pgfpathlineto{\pgfqpoint{1.299600in}{0.660000in}}%
\pgfpathlineto{\pgfqpoint{1.299600in}{4.447041in}}%
\pgfpathlineto{\pgfqpoint{1.293400in}{4.447041in}}%
\pgfpathlineto{\pgfqpoint{1.293400in}{0.660000in}}%
\pgfpathclose%
\pgfusepath{fill}%
\end{pgfscope}%
\begin{pgfscope}%
\pgfpathrectangle{\pgfqpoint{1.250000in}{0.660000in}}{\pgfqpoint{7.750000in}{4.620000in}}%
\pgfusepath{clip}%
\pgfsetbuttcap%
\pgfsetmiterjoin%
\definecolor{currentfill}{rgb}{0.000000,0.000000,0.000000}%
\pgfsetfillcolor{currentfill}%
\pgfsetlinewidth{0.000000pt}%
\definecolor{currentstroke}{rgb}{0.000000,0.000000,0.000000}%
\pgfsetstrokecolor{currentstroke}%
\pgfsetstrokeopacity{0.000000}%
\pgfsetdash{}{0pt}%
\pgfpathmoveto{\pgfqpoint{1.301150in}{0.660000in}}%
\pgfpathlineto{\pgfqpoint{1.307350in}{0.660000in}}%
\pgfpathlineto{\pgfqpoint{1.307350in}{4.362558in}}%
\pgfpathlineto{\pgfqpoint{1.301150in}{4.362558in}}%
\pgfpathlineto{\pgfqpoint{1.301150in}{0.660000in}}%
\pgfpathclose%
\pgfusepath{fill}%
\end{pgfscope}%
\begin{pgfscope}%
\pgfpathrectangle{\pgfqpoint{1.250000in}{0.660000in}}{\pgfqpoint{7.750000in}{4.620000in}}%
\pgfusepath{clip}%
\pgfsetbuttcap%
\pgfsetmiterjoin%
\definecolor{currentfill}{rgb}{0.000000,0.000000,0.000000}%
\pgfsetfillcolor{currentfill}%
\pgfsetlinewidth{0.000000pt}%
\definecolor{currentstroke}{rgb}{0.000000,0.000000,0.000000}%
\pgfsetstrokecolor{currentstroke}%
\pgfsetstrokeopacity{0.000000}%
\pgfsetdash{}{0pt}%
\pgfpathmoveto{\pgfqpoint{1.308900in}{0.660000in}}%
\pgfpathlineto{\pgfqpoint{1.315100in}{0.660000in}}%
\pgfpathlineto{\pgfqpoint{1.315100in}{4.293179in}}%
\pgfpathlineto{\pgfqpoint{1.308900in}{4.293179in}}%
\pgfpathlineto{\pgfqpoint{1.308900in}{0.660000in}}%
\pgfpathclose%
\pgfusepath{fill}%
\end{pgfscope}%
\begin{pgfscope}%
\pgfpathrectangle{\pgfqpoint{1.250000in}{0.660000in}}{\pgfqpoint{7.750000in}{4.620000in}}%
\pgfusepath{clip}%
\pgfsetbuttcap%
\pgfsetmiterjoin%
\definecolor{currentfill}{rgb}{0.000000,0.000000,0.000000}%
\pgfsetfillcolor{currentfill}%
\pgfsetlinewidth{0.000000pt}%
\definecolor{currentstroke}{rgb}{0.000000,0.000000,0.000000}%
\pgfsetstrokecolor{currentstroke}%
\pgfsetstrokeopacity{0.000000}%
\pgfsetdash{}{0pt}%
\pgfpathmoveto{\pgfqpoint{1.316650in}{0.660000in}}%
\pgfpathlineto{\pgfqpoint{1.322850in}{0.660000in}}%
\pgfpathlineto{\pgfqpoint{1.322850in}{4.234800in}}%
\pgfpathlineto{\pgfqpoint{1.316650in}{4.234800in}}%
\pgfpathlineto{\pgfqpoint{1.316650in}{0.660000in}}%
\pgfpathclose%
\pgfusepath{fill}%
\end{pgfscope}%
\begin{pgfscope}%
\pgfpathrectangle{\pgfqpoint{1.250000in}{0.660000in}}{\pgfqpoint{7.750000in}{4.620000in}}%
\pgfusepath{clip}%
\pgfsetbuttcap%
\pgfsetmiterjoin%
\definecolor{currentfill}{rgb}{0.000000,0.000000,0.000000}%
\pgfsetfillcolor{currentfill}%
\pgfsetlinewidth{0.000000pt}%
\definecolor{currentstroke}{rgb}{0.000000,0.000000,0.000000}%
\pgfsetstrokecolor{currentstroke}%
\pgfsetstrokeopacity{0.000000}%
\pgfsetdash{}{0pt}%
\pgfpathmoveto{\pgfqpoint{1.324400in}{0.660000in}}%
\pgfpathlineto{\pgfqpoint{1.330600in}{0.660000in}}%
\pgfpathlineto{\pgfqpoint{1.330600in}{4.184742in}}%
\pgfpathlineto{\pgfqpoint{1.324400in}{4.184742in}}%
\pgfpathlineto{\pgfqpoint{1.324400in}{0.660000in}}%
\pgfpathclose%
\pgfusepath{fill}%
\end{pgfscope}%
\begin{pgfscope}%
\pgfpathrectangle{\pgfqpoint{1.250000in}{0.660000in}}{\pgfqpoint{7.750000in}{4.620000in}}%
\pgfusepath{clip}%
\pgfsetbuttcap%
\pgfsetmiterjoin%
\definecolor{currentfill}{rgb}{0.000000,0.000000,0.000000}%
\pgfsetfillcolor{currentfill}%
\pgfsetlinewidth{0.000000pt}%
\definecolor{currentstroke}{rgb}{0.000000,0.000000,0.000000}%
\pgfsetstrokecolor{currentstroke}%
\pgfsetstrokeopacity{0.000000}%
\pgfsetdash{}{0pt}%
\pgfpathmoveto{\pgfqpoint{1.332150in}{0.660000in}}%
\pgfpathlineto{\pgfqpoint{1.338350in}{0.660000in}}%
\pgfpathlineto{\pgfqpoint{1.338350in}{4.141163in}}%
\pgfpathlineto{\pgfqpoint{1.332150in}{4.141163in}}%
\pgfpathlineto{\pgfqpoint{1.332150in}{0.660000in}}%
\pgfpathclose%
\pgfusepath{fill}%
\end{pgfscope}%
\begin{pgfscope}%
\pgfpathrectangle{\pgfqpoint{1.250000in}{0.660000in}}{\pgfqpoint{7.750000in}{4.620000in}}%
\pgfusepath{clip}%
\pgfsetbuttcap%
\pgfsetmiterjoin%
\definecolor{currentfill}{rgb}{0.000000,0.000000,0.000000}%
\pgfsetfillcolor{currentfill}%
\pgfsetlinewidth{0.000000pt}%
\definecolor{currentstroke}{rgb}{0.000000,0.000000,0.000000}%
\pgfsetstrokecolor{currentstroke}%
\pgfsetstrokeopacity{0.000000}%
\pgfsetdash{}{0pt}%
\pgfpathmoveto{\pgfqpoint{1.339900in}{0.660000in}}%
\pgfpathlineto{\pgfqpoint{1.346100in}{0.660000in}}%
\pgfpathlineto{\pgfqpoint{1.346100in}{4.102751in}}%
\pgfpathlineto{\pgfqpoint{1.339900in}{4.102751in}}%
\pgfpathlineto{\pgfqpoint{1.339900in}{0.660000in}}%
\pgfpathclose%
\pgfusepath{fill}%
\end{pgfscope}%
\begin{pgfscope}%
\pgfpathrectangle{\pgfqpoint{1.250000in}{0.660000in}}{\pgfqpoint{7.750000in}{4.620000in}}%
\pgfusepath{clip}%
\pgfsetbuttcap%
\pgfsetmiterjoin%
\definecolor{currentfill}{rgb}{0.000000,0.000000,0.000000}%
\pgfsetfillcolor{currentfill}%
\pgfsetlinewidth{0.000000pt}%
\definecolor{currentstroke}{rgb}{0.000000,0.000000,0.000000}%
\pgfsetstrokecolor{currentstroke}%
\pgfsetstrokeopacity{0.000000}%
\pgfsetdash{}{0pt}%
\pgfpathmoveto{\pgfqpoint{1.347650in}{0.660000in}}%
\pgfpathlineto{\pgfqpoint{1.353850in}{0.660000in}}%
\pgfpathlineto{\pgfqpoint{1.353850in}{4.068539in}}%
\pgfpathlineto{\pgfqpoint{1.347650in}{4.068539in}}%
\pgfpathlineto{\pgfqpoint{1.347650in}{0.660000in}}%
\pgfpathclose%
\pgfusepath{fill}%
\end{pgfscope}%
\begin{pgfscope}%
\pgfpathrectangle{\pgfqpoint{1.250000in}{0.660000in}}{\pgfqpoint{7.750000in}{4.620000in}}%
\pgfusepath{clip}%
\pgfsetbuttcap%
\pgfsetmiterjoin%
\definecolor{currentfill}{rgb}{0.000000,0.000000,0.000000}%
\pgfsetfillcolor{currentfill}%
\pgfsetlinewidth{0.000000pt}%
\definecolor{currentstroke}{rgb}{0.000000,0.000000,0.000000}%
\pgfsetstrokecolor{currentstroke}%
\pgfsetstrokeopacity{0.000000}%
\pgfsetdash{}{0pt}%
\pgfpathmoveto{\pgfqpoint{1.355400in}{0.660000in}}%
\pgfpathlineto{\pgfqpoint{1.361600in}{0.660000in}}%
\pgfpathlineto{\pgfqpoint{1.361600in}{4.037797in}}%
\pgfpathlineto{\pgfqpoint{1.355400in}{4.037797in}}%
\pgfpathlineto{\pgfqpoint{1.355400in}{0.660000in}}%
\pgfpathclose%
\pgfusepath{fill}%
\end{pgfscope}%
\begin{pgfscope}%
\pgfpathrectangle{\pgfqpoint{1.250000in}{0.660000in}}{\pgfqpoint{7.750000in}{4.620000in}}%
\pgfusepath{clip}%
\pgfsetbuttcap%
\pgfsetmiterjoin%
\definecolor{currentfill}{rgb}{0.000000,0.000000,0.000000}%
\pgfsetfillcolor{currentfill}%
\pgfsetlinewidth{0.000000pt}%
\definecolor{currentstroke}{rgb}{0.000000,0.000000,0.000000}%
\pgfsetstrokecolor{currentstroke}%
\pgfsetstrokeopacity{0.000000}%
\pgfsetdash{}{0pt}%
\pgfpathmoveto{\pgfqpoint{1.363150in}{0.660000in}}%
\pgfpathlineto{\pgfqpoint{1.369350in}{0.660000in}}%
\pgfpathlineto{\pgfqpoint{1.369350in}{4.009965in}}%
\pgfpathlineto{\pgfqpoint{1.363150in}{4.009965in}}%
\pgfpathlineto{\pgfqpoint{1.363150in}{0.660000in}}%
\pgfpathclose%
\pgfusepath{fill}%
\end{pgfscope}%
\begin{pgfscope}%
\pgfpathrectangle{\pgfqpoint{1.250000in}{0.660000in}}{\pgfqpoint{7.750000in}{4.620000in}}%
\pgfusepath{clip}%
\pgfsetbuttcap%
\pgfsetmiterjoin%
\definecolor{currentfill}{rgb}{0.000000,0.000000,0.000000}%
\pgfsetfillcolor{currentfill}%
\pgfsetlinewidth{0.000000pt}%
\definecolor{currentstroke}{rgb}{0.000000,0.000000,0.000000}%
\pgfsetstrokecolor{currentstroke}%
\pgfsetstrokeopacity{0.000000}%
\pgfsetdash{}{0pt}%
\pgfpathmoveto{\pgfqpoint{1.370900in}{0.660000in}}%
\pgfpathlineto{\pgfqpoint{1.377100in}{0.660000in}}%
\pgfpathlineto{\pgfqpoint{1.377100in}{3.984601in}}%
\pgfpathlineto{\pgfqpoint{1.370900in}{3.984601in}}%
\pgfpathlineto{\pgfqpoint{1.370900in}{0.660000in}}%
\pgfpathclose%
\pgfusepath{fill}%
\end{pgfscope}%
\begin{pgfscope}%
\pgfpathrectangle{\pgfqpoint{1.250000in}{0.660000in}}{\pgfqpoint{7.750000in}{4.620000in}}%
\pgfusepath{clip}%
\pgfsetbuttcap%
\pgfsetmiterjoin%
\definecolor{currentfill}{rgb}{0.000000,0.000000,0.000000}%
\pgfsetfillcolor{currentfill}%
\pgfsetlinewidth{0.000000pt}%
\definecolor{currentstroke}{rgb}{0.000000,0.000000,0.000000}%
\pgfsetstrokecolor{currentstroke}%
\pgfsetstrokeopacity{0.000000}%
\pgfsetdash{}{0pt}%
\pgfpathmoveto{\pgfqpoint{1.378650in}{0.660000in}}%
\pgfpathlineto{\pgfqpoint{1.384850in}{0.660000in}}%
\pgfpathlineto{\pgfqpoint{1.384850in}{3.961352in}}%
\pgfpathlineto{\pgfqpoint{1.378650in}{3.961352in}}%
\pgfpathlineto{\pgfqpoint{1.378650in}{0.660000in}}%
\pgfpathclose%
\pgfusepath{fill}%
\end{pgfscope}%
\begin{pgfscope}%
\pgfpathrectangle{\pgfqpoint{1.250000in}{0.660000in}}{\pgfqpoint{7.750000in}{4.620000in}}%
\pgfusepath{clip}%
\pgfsetbuttcap%
\pgfsetmiterjoin%
\definecolor{currentfill}{rgb}{0.000000,0.000000,0.000000}%
\pgfsetfillcolor{currentfill}%
\pgfsetlinewidth{0.000000pt}%
\definecolor{currentstroke}{rgb}{0.000000,0.000000,0.000000}%
\pgfsetstrokecolor{currentstroke}%
\pgfsetstrokeopacity{0.000000}%
\pgfsetdash{}{0pt}%
\pgfpathmoveto{\pgfqpoint{1.386400in}{0.660000in}}%
\pgfpathlineto{\pgfqpoint{1.392600in}{0.660000in}}%
\pgfpathlineto{\pgfqpoint{1.392600in}{3.939933in}}%
\pgfpathlineto{\pgfqpoint{1.386400in}{3.939933in}}%
\pgfpathlineto{\pgfqpoint{1.386400in}{0.660000in}}%
\pgfpathclose%
\pgfusepath{fill}%
\end{pgfscope}%
\begin{pgfscope}%
\pgfpathrectangle{\pgfqpoint{1.250000in}{0.660000in}}{\pgfqpoint{7.750000in}{4.620000in}}%
\pgfusepath{clip}%
\pgfsetbuttcap%
\pgfsetmiterjoin%
\definecolor{currentfill}{rgb}{0.000000,0.000000,0.000000}%
\pgfsetfillcolor{currentfill}%
\pgfsetlinewidth{0.000000pt}%
\definecolor{currentstroke}{rgb}{0.000000,0.000000,0.000000}%
\pgfsetstrokecolor{currentstroke}%
\pgfsetstrokeopacity{0.000000}%
\pgfsetdash{}{0pt}%
\pgfpathmoveto{\pgfqpoint{1.394150in}{0.660000in}}%
\pgfpathlineto{\pgfqpoint{1.400350in}{0.660000in}}%
\pgfpathlineto{\pgfqpoint{1.400350in}{3.920112in}}%
\pgfpathlineto{\pgfqpoint{1.394150in}{3.920112in}}%
\pgfpathlineto{\pgfqpoint{1.394150in}{0.660000in}}%
\pgfpathclose%
\pgfusepath{fill}%
\end{pgfscope}%
\begin{pgfscope}%
\pgfpathrectangle{\pgfqpoint{1.250000in}{0.660000in}}{\pgfqpoint{7.750000in}{4.620000in}}%
\pgfusepath{clip}%
\pgfsetbuttcap%
\pgfsetmiterjoin%
\definecolor{currentfill}{rgb}{0.000000,0.000000,0.000000}%
\pgfsetfillcolor{currentfill}%
\pgfsetlinewidth{0.000000pt}%
\definecolor{currentstroke}{rgb}{0.000000,0.000000,0.000000}%
\pgfsetstrokecolor{currentstroke}%
\pgfsetstrokeopacity{0.000000}%
\pgfsetdash{}{0pt}%
\pgfpathmoveto{\pgfqpoint{1.401900in}{0.660000in}}%
\pgfpathlineto{\pgfqpoint{1.408100in}{0.660000in}}%
\pgfpathlineto{\pgfqpoint{1.408100in}{3.901694in}}%
\pgfpathlineto{\pgfqpoint{1.401900in}{3.901694in}}%
\pgfpathlineto{\pgfqpoint{1.401900in}{0.660000in}}%
\pgfpathclose%
\pgfusepath{fill}%
\end{pgfscope}%
\begin{pgfscope}%
\pgfpathrectangle{\pgfqpoint{1.250000in}{0.660000in}}{\pgfqpoint{7.750000in}{4.620000in}}%
\pgfusepath{clip}%
\pgfsetbuttcap%
\pgfsetmiterjoin%
\definecolor{currentfill}{rgb}{0.000000,0.000000,0.000000}%
\pgfsetfillcolor{currentfill}%
\pgfsetlinewidth{0.000000pt}%
\definecolor{currentstroke}{rgb}{0.000000,0.000000,0.000000}%
\pgfsetstrokecolor{currentstroke}%
\pgfsetstrokeopacity{0.000000}%
\pgfsetdash{}{0pt}%
\pgfpathmoveto{\pgfqpoint{1.409650in}{0.660000in}}%
\pgfpathlineto{\pgfqpoint{1.415850in}{0.660000in}}%
\pgfpathlineto{\pgfqpoint{1.415850in}{3.884517in}}%
\pgfpathlineto{\pgfqpoint{1.409650in}{3.884517in}}%
\pgfpathlineto{\pgfqpoint{1.409650in}{0.660000in}}%
\pgfpathclose%
\pgfusepath{fill}%
\end{pgfscope}%
\begin{pgfscope}%
\pgfpathrectangle{\pgfqpoint{1.250000in}{0.660000in}}{\pgfqpoint{7.750000in}{4.620000in}}%
\pgfusepath{clip}%
\pgfsetbuttcap%
\pgfsetmiterjoin%
\definecolor{currentfill}{rgb}{0.000000,0.000000,0.000000}%
\pgfsetfillcolor{currentfill}%
\pgfsetlinewidth{0.000000pt}%
\definecolor{currentstroke}{rgb}{0.000000,0.000000,0.000000}%
\pgfsetstrokecolor{currentstroke}%
\pgfsetstrokeopacity{0.000000}%
\pgfsetdash{}{0pt}%
\pgfpathmoveto{\pgfqpoint{1.417400in}{0.660000in}}%
\pgfpathlineto{\pgfqpoint{1.423600in}{0.660000in}}%
\pgfpathlineto{\pgfqpoint{1.423600in}{3.868444in}}%
\pgfpathlineto{\pgfqpoint{1.417400in}{3.868444in}}%
\pgfpathlineto{\pgfqpoint{1.417400in}{0.660000in}}%
\pgfpathclose%
\pgfusepath{fill}%
\end{pgfscope}%
\begin{pgfscope}%
\pgfpathrectangle{\pgfqpoint{1.250000in}{0.660000in}}{\pgfqpoint{7.750000in}{4.620000in}}%
\pgfusepath{clip}%
\pgfsetbuttcap%
\pgfsetmiterjoin%
\definecolor{currentfill}{rgb}{0.000000,0.000000,0.000000}%
\pgfsetfillcolor{currentfill}%
\pgfsetlinewidth{0.000000pt}%
\definecolor{currentstroke}{rgb}{0.000000,0.000000,0.000000}%
\pgfsetstrokecolor{currentstroke}%
\pgfsetstrokeopacity{0.000000}%
\pgfsetdash{}{0pt}%
\pgfpathmoveto{\pgfqpoint{1.425150in}{0.660000in}}%
\pgfpathlineto{\pgfqpoint{1.431350in}{0.660000in}}%
\pgfpathlineto{\pgfqpoint{1.431350in}{3.853360in}}%
\pgfpathlineto{\pgfqpoint{1.425150in}{3.853360in}}%
\pgfpathlineto{\pgfqpoint{1.425150in}{0.660000in}}%
\pgfpathclose%
\pgfusepath{fill}%
\end{pgfscope}%
\begin{pgfscope}%
\pgfpathrectangle{\pgfqpoint{1.250000in}{0.660000in}}{\pgfqpoint{7.750000in}{4.620000in}}%
\pgfusepath{clip}%
\pgfsetbuttcap%
\pgfsetmiterjoin%
\definecolor{currentfill}{rgb}{0.000000,0.000000,0.000000}%
\pgfsetfillcolor{currentfill}%
\pgfsetlinewidth{0.000000pt}%
\definecolor{currentstroke}{rgb}{0.000000,0.000000,0.000000}%
\pgfsetstrokecolor{currentstroke}%
\pgfsetstrokeopacity{0.000000}%
\pgfsetdash{}{0pt}%
\pgfpathmoveto{\pgfqpoint{1.432900in}{0.660000in}}%
\pgfpathlineto{\pgfqpoint{1.439100in}{0.660000in}}%
\pgfpathlineto{\pgfqpoint{1.439100in}{3.839164in}}%
\pgfpathlineto{\pgfqpoint{1.432900in}{3.839164in}}%
\pgfpathlineto{\pgfqpoint{1.432900in}{0.660000in}}%
\pgfpathclose%
\pgfusepath{fill}%
\end{pgfscope}%
\begin{pgfscope}%
\pgfpathrectangle{\pgfqpoint{1.250000in}{0.660000in}}{\pgfqpoint{7.750000in}{4.620000in}}%
\pgfusepath{clip}%
\pgfsetbuttcap%
\pgfsetmiterjoin%
\definecolor{currentfill}{rgb}{0.000000,0.000000,0.000000}%
\pgfsetfillcolor{currentfill}%
\pgfsetlinewidth{0.000000pt}%
\definecolor{currentstroke}{rgb}{0.000000,0.000000,0.000000}%
\pgfsetstrokecolor{currentstroke}%
\pgfsetstrokeopacity{0.000000}%
\pgfsetdash{}{0pt}%
\pgfpathmoveto{\pgfqpoint{1.440650in}{0.660000in}}%
\pgfpathlineto{\pgfqpoint{1.446850in}{0.660000in}}%
\pgfpathlineto{\pgfqpoint{1.446850in}{3.825771in}}%
\pgfpathlineto{\pgfqpoint{1.440650in}{3.825771in}}%
\pgfpathlineto{\pgfqpoint{1.440650in}{0.660000in}}%
\pgfpathclose%
\pgfusepath{fill}%
\end{pgfscope}%
\begin{pgfscope}%
\pgfpathrectangle{\pgfqpoint{1.250000in}{0.660000in}}{\pgfqpoint{7.750000in}{4.620000in}}%
\pgfusepath{clip}%
\pgfsetbuttcap%
\pgfsetmiterjoin%
\definecolor{currentfill}{rgb}{0.000000,0.000000,0.000000}%
\pgfsetfillcolor{currentfill}%
\pgfsetlinewidth{0.000000pt}%
\definecolor{currentstroke}{rgb}{0.000000,0.000000,0.000000}%
\pgfsetstrokecolor{currentstroke}%
\pgfsetstrokeopacity{0.000000}%
\pgfsetdash{}{0pt}%
\pgfpathmoveto{\pgfqpoint{1.448400in}{0.660000in}}%
\pgfpathlineto{\pgfqpoint{1.454600in}{0.660000in}}%
\pgfpathlineto{\pgfqpoint{1.454600in}{3.813106in}}%
\pgfpathlineto{\pgfqpoint{1.448400in}{3.813106in}}%
\pgfpathlineto{\pgfqpoint{1.448400in}{0.660000in}}%
\pgfpathclose%
\pgfusepath{fill}%
\end{pgfscope}%
\begin{pgfscope}%
\pgfpathrectangle{\pgfqpoint{1.250000in}{0.660000in}}{\pgfqpoint{7.750000in}{4.620000in}}%
\pgfusepath{clip}%
\pgfsetbuttcap%
\pgfsetmiterjoin%
\definecolor{currentfill}{rgb}{0.000000,0.000000,0.000000}%
\pgfsetfillcolor{currentfill}%
\pgfsetlinewidth{0.000000pt}%
\definecolor{currentstroke}{rgb}{0.000000,0.000000,0.000000}%
\pgfsetstrokecolor{currentstroke}%
\pgfsetstrokeopacity{0.000000}%
\pgfsetdash{}{0pt}%
\pgfpathmoveto{\pgfqpoint{1.456150in}{0.660000in}}%
\pgfpathlineto{\pgfqpoint{1.462350in}{0.660000in}}%
\pgfpathlineto{\pgfqpoint{1.462350in}{3.801103in}}%
\pgfpathlineto{\pgfqpoint{1.456150in}{3.801103in}}%
\pgfpathlineto{\pgfqpoint{1.456150in}{0.660000in}}%
\pgfpathclose%
\pgfusepath{fill}%
\end{pgfscope}%
\begin{pgfscope}%
\pgfpathrectangle{\pgfqpoint{1.250000in}{0.660000in}}{\pgfqpoint{7.750000in}{4.620000in}}%
\pgfusepath{clip}%
\pgfsetbuttcap%
\pgfsetmiterjoin%
\definecolor{currentfill}{rgb}{0.000000,0.000000,0.000000}%
\pgfsetfillcolor{currentfill}%
\pgfsetlinewidth{0.000000pt}%
\definecolor{currentstroke}{rgb}{0.000000,0.000000,0.000000}%
\pgfsetstrokecolor{currentstroke}%
\pgfsetstrokeopacity{0.000000}%
\pgfsetdash{}{0pt}%
\pgfpathmoveto{\pgfqpoint{1.463900in}{0.660000in}}%
\pgfpathlineto{\pgfqpoint{1.470100in}{0.660000in}}%
\pgfpathlineto{\pgfqpoint{1.470100in}{3.789705in}}%
\pgfpathlineto{\pgfqpoint{1.463900in}{3.789705in}}%
\pgfpathlineto{\pgfqpoint{1.463900in}{0.660000in}}%
\pgfpathclose%
\pgfusepath{fill}%
\end{pgfscope}%
\begin{pgfscope}%
\pgfpathrectangle{\pgfqpoint{1.250000in}{0.660000in}}{\pgfqpoint{7.750000in}{4.620000in}}%
\pgfusepath{clip}%
\pgfsetbuttcap%
\pgfsetmiterjoin%
\definecolor{currentfill}{rgb}{0.000000,0.000000,0.000000}%
\pgfsetfillcolor{currentfill}%
\pgfsetlinewidth{0.000000pt}%
\definecolor{currentstroke}{rgb}{0.000000,0.000000,0.000000}%
\pgfsetstrokecolor{currentstroke}%
\pgfsetstrokeopacity{0.000000}%
\pgfsetdash{}{0pt}%
\pgfpathmoveto{\pgfqpoint{1.471650in}{0.660000in}}%
\pgfpathlineto{\pgfqpoint{1.477850in}{0.660000in}}%
\pgfpathlineto{\pgfqpoint{1.477850in}{3.778862in}}%
\pgfpathlineto{\pgfqpoint{1.471650in}{3.778862in}}%
\pgfpathlineto{\pgfqpoint{1.471650in}{0.660000in}}%
\pgfpathclose%
\pgfusepath{fill}%
\end{pgfscope}%
\begin{pgfscope}%
\pgfpathrectangle{\pgfqpoint{1.250000in}{0.660000in}}{\pgfqpoint{7.750000in}{4.620000in}}%
\pgfusepath{clip}%
\pgfsetbuttcap%
\pgfsetmiterjoin%
\definecolor{currentfill}{rgb}{0.000000,0.000000,0.000000}%
\pgfsetfillcolor{currentfill}%
\pgfsetlinewidth{0.000000pt}%
\definecolor{currentstroke}{rgb}{0.000000,0.000000,0.000000}%
\pgfsetstrokecolor{currentstroke}%
\pgfsetstrokeopacity{0.000000}%
\pgfsetdash{}{0pt}%
\pgfpathmoveto{\pgfqpoint{1.479400in}{0.660000in}}%
\pgfpathlineto{\pgfqpoint{1.485600in}{0.660000in}}%
\pgfpathlineto{\pgfqpoint{1.485600in}{3.768528in}}%
\pgfpathlineto{\pgfqpoint{1.479400in}{3.768528in}}%
\pgfpathlineto{\pgfqpoint{1.479400in}{0.660000in}}%
\pgfpathclose%
\pgfusepath{fill}%
\end{pgfscope}%
\begin{pgfscope}%
\pgfpathrectangle{\pgfqpoint{1.250000in}{0.660000in}}{\pgfqpoint{7.750000in}{4.620000in}}%
\pgfusepath{clip}%
\pgfsetbuttcap%
\pgfsetmiterjoin%
\definecolor{currentfill}{rgb}{0.000000,0.000000,0.000000}%
\pgfsetfillcolor{currentfill}%
\pgfsetlinewidth{0.000000pt}%
\definecolor{currentstroke}{rgb}{0.000000,0.000000,0.000000}%
\pgfsetstrokecolor{currentstroke}%
\pgfsetstrokeopacity{0.000000}%
\pgfsetdash{}{0pt}%
\pgfpathmoveto{\pgfqpoint{1.487150in}{0.660000in}}%
\pgfpathlineto{\pgfqpoint{1.493350in}{0.660000in}}%
\pgfpathlineto{\pgfqpoint{1.493350in}{3.758666in}}%
\pgfpathlineto{\pgfqpoint{1.487150in}{3.758666in}}%
\pgfpathlineto{\pgfqpoint{1.487150in}{0.660000in}}%
\pgfpathclose%
\pgfusepath{fill}%
\end{pgfscope}%
\begin{pgfscope}%
\pgfpathrectangle{\pgfqpoint{1.250000in}{0.660000in}}{\pgfqpoint{7.750000in}{4.620000in}}%
\pgfusepath{clip}%
\pgfsetbuttcap%
\pgfsetmiterjoin%
\definecolor{currentfill}{rgb}{0.000000,0.000000,0.000000}%
\pgfsetfillcolor{currentfill}%
\pgfsetlinewidth{0.000000pt}%
\definecolor{currentstroke}{rgb}{0.000000,0.000000,0.000000}%
\pgfsetstrokecolor{currentstroke}%
\pgfsetstrokeopacity{0.000000}%
\pgfsetdash{}{0pt}%
\pgfpathmoveto{\pgfqpoint{1.494900in}{0.660000in}}%
\pgfpathlineto{\pgfqpoint{1.501100in}{0.660000in}}%
\pgfpathlineto{\pgfqpoint{1.501100in}{3.749237in}}%
\pgfpathlineto{\pgfqpoint{1.494900in}{3.749237in}}%
\pgfpathlineto{\pgfqpoint{1.494900in}{0.660000in}}%
\pgfpathclose%
\pgfusepath{fill}%
\end{pgfscope}%
\begin{pgfscope}%
\pgfpathrectangle{\pgfqpoint{1.250000in}{0.660000in}}{\pgfqpoint{7.750000in}{4.620000in}}%
\pgfusepath{clip}%
\pgfsetbuttcap%
\pgfsetmiterjoin%
\definecolor{currentfill}{rgb}{0.000000,0.000000,0.000000}%
\pgfsetfillcolor{currentfill}%
\pgfsetlinewidth{0.000000pt}%
\definecolor{currentstroke}{rgb}{0.000000,0.000000,0.000000}%
\pgfsetstrokecolor{currentstroke}%
\pgfsetstrokeopacity{0.000000}%
\pgfsetdash{}{0pt}%
\pgfpathmoveto{\pgfqpoint{1.502650in}{0.660000in}}%
\pgfpathlineto{\pgfqpoint{1.508850in}{0.660000in}}%
\pgfpathlineto{\pgfqpoint{1.508850in}{3.740212in}}%
\pgfpathlineto{\pgfqpoint{1.502650in}{3.740212in}}%
\pgfpathlineto{\pgfqpoint{1.502650in}{0.660000in}}%
\pgfpathclose%
\pgfusepath{fill}%
\end{pgfscope}%
\begin{pgfscope}%
\pgfpathrectangle{\pgfqpoint{1.250000in}{0.660000in}}{\pgfqpoint{7.750000in}{4.620000in}}%
\pgfusepath{clip}%
\pgfsetbuttcap%
\pgfsetmiterjoin%
\definecolor{currentfill}{rgb}{0.000000,0.000000,0.000000}%
\pgfsetfillcolor{currentfill}%
\pgfsetlinewidth{0.000000pt}%
\definecolor{currentstroke}{rgb}{0.000000,0.000000,0.000000}%
\pgfsetstrokecolor{currentstroke}%
\pgfsetstrokeopacity{0.000000}%
\pgfsetdash{}{0pt}%
\pgfpathmoveto{\pgfqpoint{1.510400in}{0.660000in}}%
\pgfpathlineto{\pgfqpoint{1.516600in}{0.660000in}}%
\pgfpathlineto{\pgfqpoint{1.516600in}{3.731561in}}%
\pgfpathlineto{\pgfqpoint{1.510400in}{3.731561in}}%
\pgfpathlineto{\pgfqpoint{1.510400in}{0.660000in}}%
\pgfpathclose%
\pgfusepath{fill}%
\end{pgfscope}%
\begin{pgfscope}%
\pgfpathrectangle{\pgfqpoint{1.250000in}{0.660000in}}{\pgfqpoint{7.750000in}{4.620000in}}%
\pgfusepath{clip}%
\pgfsetbuttcap%
\pgfsetmiterjoin%
\definecolor{currentfill}{rgb}{0.000000,0.000000,0.000000}%
\pgfsetfillcolor{currentfill}%
\pgfsetlinewidth{0.000000pt}%
\definecolor{currentstroke}{rgb}{0.000000,0.000000,0.000000}%
\pgfsetstrokecolor{currentstroke}%
\pgfsetstrokeopacity{0.000000}%
\pgfsetdash{}{0pt}%
\pgfpathmoveto{\pgfqpoint{1.518150in}{0.660000in}}%
\pgfpathlineto{\pgfqpoint{1.524350in}{0.660000in}}%
\pgfpathlineto{\pgfqpoint{1.524350in}{3.723258in}}%
\pgfpathlineto{\pgfqpoint{1.518150in}{3.723258in}}%
\pgfpathlineto{\pgfqpoint{1.518150in}{0.660000in}}%
\pgfpathclose%
\pgfusepath{fill}%
\end{pgfscope}%
\begin{pgfscope}%
\pgfpathrectangle{\pgfqpoint{1.250000in}{0.660000in}}{\pgfqpoint{7.750000in}{4.620000in}}%
\pgfusepath{clip}%
\pgfsetbuttcap%
\pgfsetmiterjoin%
\definecolor{currentfill}{rgb}{0.000000,0.000000,0.000000}%
\pgfsetfillcolor{currentfill}%
\pgfsetlinewidth{0.000000pt}%
\definecolor{currentstroke}{rgb}{0.000000,0.000000,0.000000}%
\pgfsetstrokecolor{currentstroke}%
\pgfsetstrokeopacity{0.000000}%
\pgfsetdash{}{0pt}%
\pgfpathmoveto{\pgfqpoint{1.525900in}{0.660000in}}%
\pgfpathlineto{\pgfqpoint{1.532100in}{0.660000in}}%
\pgfpathlineto{\pgfqpoint{1.532100in}{3.715280in}}%
\pgfpathlineto{\pgfqpoint{1.525900in}{3.715280in}}%
\pgfpathlineto{\pgfqpoint{1.525900in}{0.660000in}}%
\pgfpathclose%
\pgfusepath{fill}%
\end{pgfscope}%
\begin{pgfscope}%
\pgfpathrectangle{\pgfqpoint{1.250000in}{0.660000in}}{\pgfqpoint{7.750000in}{4.620000in}}%
\pgfusepath{clip}%
\pgfsetbuttcap%
\pgfsetmiterjoin%
\definecolor{currentfill}{rgb}{0.000000,0.000000,0.000000}%
\pgfsetfillcolor{currentfill}%
\pgfsetlinewidth{0.000000pt}%
\definecolor{currentstroke}{rgb}{0.000000,0.000000,0.000000}%
\pgfsetstrokecolor{currentstroke}%
\pgfsetstrokeopacity{0.000000}%
\pgfsetdash{}{0pt}%
\pgfpathmoveto{\pgfqpoint{1.533650in}{0.660000in}}%
\pgfpathlineto{\pgfqpoint{1.539850in}{0.660000in}}%
\pgfpathlineto{\pgfqpoint{1.539850in}{3.707605in}}%
\pgfpathlineto{\pgfqpoint{1.533650in}{3.707605in}}%
\pgfpathlineto{\pgfqpoint{1.533650in}{0.660000in}}%
\pgfpathclose%
\pgfusepath{fill}%
\end{pgfscope}%
\begin{pgfscope}%
\pgfpathrectangle{\pgfqpoint{1.250000in}{0.660000in}}{\pgfqpoint{7.750000in}{4.620000in}}%
\pgfusepath{clip}%
\pgfsetbuttcap%
\pgfsetmiterjoin%
\definecolor{currentfill}{rgb}{0.000000,0.000000,0.000000}%
\pgfsetfillcolor{currentfill}%
\pgfsetlinewidth{0.000000pt}%
\definecolor{currentstroke}{rgb}{0.000000,0.000000,0.000000}%
\pgfsetstrokecolor{currentstroke}%
\pgfsetstrokeopacity{0.000000}%
\pgfsetdash{}{0pt}%
\pgfpathmoveto{\pgfqpoint{1.541400in}{0.660000in}}%
\pgfpathlineto{\pgfqpoint{1.547600in}{0.660000in}}%
\pgfpathlineto{\pgfqpoint{1.547600in}{3.700216in}}%
\pgfpathlineto{\pgfqpoint{1.541400in}{3.700216in}}%
\pgfpathlineto{\pgfqpoint{1.541400in}{0.660000in}}%
\pgfpathclose%
\pgfusepath{fill}%
\end{pgfscope}%
\begin{pgfscope}%
\pgfpathrectangle{\pgfqpoint{1.250000in}{0.660000in}}{\pgfqpoint{7.750000in}{4.620000in}}%
\pgfusepath{clip}%
\pgfsetbuttcap%
\pgfsetmiterjoin%
\definecolor{currentfill}{rgb}{0.000000,0.000000,0.000000}%
\pgfsetfillcolor{currentfill}%
\pgfsetlinewidth{0.000000pt}%
\definecolor{currentstroke}{rgb}{0.000000,0.000000,0.000000}%
\pgfsetstrokecolor{currentstroke}%
\pgfsetstrokeopacity{0.000000}%
\pgfsetdash{}{0pt}%
\pgfpathmoveto{\pgfqpoint{1.549150in}{0.660000in}}%
\pgfpathlineto{\pgfqpoint{1.555350in}{0.660000in}}%
\pgfpathlineto{\pgfqpoint{1.555350in}{3.693092in}}%
\pgfpathlineto{\pgfqpoint{1.549150in}{3.693092in}}%
\pgfpathlineto{\pgfqpoint{1.549150in}{0.660000in}}%
\pgfpathclose%
\pgfusepath{fill}%
\end{pgfscope}%
\begin{pgfscope}%
\pgfpathrectangle{\pgfqpoint{1.250000in}{0.660000in}}{\pgfqpoint{7.750000in}{4.620000in}}%
\pgfusepath{clip}%
\pgfsetbuttcap%
\pgfsetmiterjoin%
\definecolor{currentfill}{rgb}{0.000000,0.000000,0.000000}%
\pgfsetfillcolor{currentfill}%
\pgfsetlinewidth{0.000000pt}%
\definecolor{currentstroke}{rgb}{0.000000,0.000000,0.000000}%
\pgfsetstrokecolor{currentstroke}%
\pgfsetstrokeopacity{0.000000}%
\pgfsetdash{}{0pt}%
\pgfpathmoveto{\pgfqpoint{1.556900in}{0.660000in}}%
\pgfpathlineto{\pgfqpoint{1.563100in}{0.660000in}}%
\pgfpathlineto{\pgfqpoint{1.563100in}{3.686220in}}%
\pgfpathlineto{\pgfqpoint{1.556900in}{3.686220in}}%
\pgfpathlineto{\pgfqpoint{1.556900in}{0.660000in}}%
\pgfpathclose%
\pgfusepath{fill}%
\end{pgfscope}%
\begin{pgfscope}%
\pgfpathrectangle{\pgfqpoint{1.250000in}{0.660000in}}{\pgfqpoint{7.750000in}{4.620000in}}%
\pgfusepath{clip}%
\pgfsetbuttcap%
\pgfsetmiterjoin%
\definecolor{currentfill}{rgb}{0.000000,0.000000,0.000000}%
\pgfsetfillcolor{currentfill}%
\pgfsetlinewidth{0.000000pt}%
\definecolor{currentstroke}{rgb}{0.000000,0.000000,0.000000}%
\pgfsetstrokecolor{currentstroke}%
\pgfsetstrokeopacity{0.000000}%
\pgfsetdash{}{0pt}%
\pgfpathmoveto{\pgfqpoint{1.564650in}{0.660000in}}%
\pgfpathlineto{\pgfqpoint{1.570850in}{0.660000in}}%
\pgfpathlineto{\pgfqpoint{1.570850in}{3.679583in}}%
\pgfpathlineto{\pgfqpoint{1.564650in}{3.679583in}}%
\pgfpathlineto{\pgfqpoint{1.564650in}{0.660000in}}%
\pgfpathclose%
\pgfusepath{fill}%
\end{pgfscope}%
\begin{pgfscope}%
\pgfpathrectangle{\pgfqpoint{1.250000in}{0.660000in}}{\pgfqpoint{7.750000in}{4.620000in}}%
\pgfusepath{clip}%
\pgfsetbuttcap%
\pgfsetmiterjoin%
\definecolor{currentfill}{rgb}{0.000000,0.000000,0.000000}%
\pgfsetfillcolor{currentfill}%
\pgfsetlinewidth{0.000000pt}%
\definecolor{currentstroke}{rgb}{0.000000,0.000000,0.000000}%
\pgfsetstrokecolor{currentstroke}%
\pgfsetstrokeopacity{0.000000}%
\pgfsetdash{}{0pt}%
\pgfpathmoveto{\pgfqpoint{1.572400in}{0.660000in}}%
\pgfpathlineto{\pgfqpoint{1.578600in}{0.660000in}}%
\pgfpathlineto{\pgfqpoint{1.578600in}{3.673169in}}%
\pgfpathlineto{\pgfqpoint{1.572400in}{3.673169in}}%
\pgfpathlineto{\pgfqpoint{1.572400in}{0.660000in}}%
\pgfpathclose%
\pgfusepath{fill}%
\end{pgfscope}%
\begin{pgfscope}%
\pgfpathrectangle{\pgfqpoint{1.250000in}{0.660000in}}{\pgfqpoint{7.750000in}{4.620000in}}%
\pgfusepath{clip}%
\pgfsetbuttcap%
\pgfsetmiterjoin%
\definecolor{currentfill}{rgb}{0.000000,0.000000,0.000000}%
\pgfsetfillcolor{currentfill}%
\pgfsetlinewidth{0.000000pt}%
\definecolor{currentstroke}{rgb}{0.000000,0.000000,0.000000}%
\pgfsetstrokecolor{currentstroke}%
\pgfsetstrokeopacity{0.000000}%
\pgfsetdash{}{0pt}%
\pgfpathmoveto{\pgfqpoint{1.580150in}{0.660000in}}%
\pgfpathlineto{\pgfqpoint{1.586350in}{0.660000in}}%
\pgfpathlineto{\pgfqpoint{1.586350in}{3.666964in}}%
\pgfpathlineto{\pgfqpoint{1.580150in}{3.666964in}}%
\pgfpathlineto{\pgfqpoint{1.580150in}{0.660000in}}%
\pgfpathclose%
\pgfusepath{fill}%
\end{pgfscope}%
\begin{pgfscope}%
\pgfpathrectangle{\pgfqpoint{1.250000in}{0.660000in}}{\pgfqpoint{7.750000in}{4.620000in}}%
\pgfusepath{clip}%
\pgfsetbuttcap%
\pgfsetmiterjoin%
\definecolor{currentfill}{rgb}{0.000000,0.000000,0.000000}%
\pgfsetfillcolor{currentfill}%
\pgfsetlinewidth{0.000000pt}%
\definecolor{currentstroke}{rgb}{0.000000,0.000000,0.000000}%
\pgfsetstrokecolor{currentstroke}%
\pgfsetstrokeopacity{0.000000}%
\pgfsetdash{}{0pt}%
\pgfpathmoveto{\pgfqpoint{1.587900in}{0.660000in}}%
\pgfpathlineto{\pgfqpoint{1.594100in}{0.660000in}}%
\pgfpathlineto{\pgfqpoint{1.594100in}{3.660958in}}%
\pgfpathlineto{\pgfqpoint{1.587900in}{3.660958in}}%
\pgfpathlineto{\pgfqpoint{1.587900in}{0.660000in}}%
\pgfpathclose%
\pgfusepath{fill}%
\end{pgfscope}%
\begin{pgfscope}%
\pgfpathrectangle{\pgfqpoint{1.250000in}{0.660000in}}{\pgfqpoint{7.750000in}{4.620000in}}%
\pgfusepath{clip}%
\pgfsetbuttcap%
\pgfsetmiterjoin%
\definecolor{currentfill}{rgb}{0.000000,0.000000,0.000000}%
\pgfsetfillcolor{currentfill}%
\pgfsetlinewidth{0.000000pt}%
\definecolor{currentstroke}{rgb}{0.000000,0.000000,0.000000}%
\pgfsetstrokecolor{currentstroke}%
\pgfsetstrokeopacity{0.000000}%
\pgfsetdash{}{0pt}%
\pgfpathmoveto{\pgfqpoint{1.595650in}{0.660000in}}%
\pgfpathlineto{\pgfqpoint{1.601850in}{0.660000in}}%
\pgfpathlineto{\pgfqpoint{1.601850in}{3.655140in}}%
\pgfpathlineto{\pgfqpoint{1.595650in}{3.655140in}}%
\pgfpathlineto{\pgfqpoint{1.595650in}{0.660000in}}%
\pgfpathclose%
\pgfusepath{fill}%
\end{pgfscope}%
\begin{pgfscope}%
\pgfpathrectangle{\pgfqpoint{1.250000in}{0.660000in}}{\pgfqpoint{7.750000in}{4.620000in}}%
\pgfusepath{clip}%
\pgfsetbuttcap%
\pgfsetmiterjoin%
\definecolor{currentfill}{rgb}{0.000000,0.000000,0.000000}%
\pgfsetfillcolor{currentfill}%
\pgfsetlinewidth{0.000000pt}%
\definecolor{currentstroke}{rgb}{0.000000,0.000000,0.000000}%
\pgfsetstrokecolor{currentstroke}%
\pgfsetstrokeopacity{0.000000}%
\pgfsetdash{}{0pt}%
\pgfpathmoveto{\pgfqpoint{1.603400in}{0.660000in}}%
\pgfpathlineto{\pgfqpoint{1.609600in}{0.660000in}}%
\pgfpathlineto{\pgfqpoint{1.609600in}{3.649499in}}%
\pgfpathlineto{\pgfqpoint{1.603400in}{3.649499in}}%
\pgfpathlineto{\pgfqpoint{1.603400in}{0.660000in}}%
\pgfpathclose%
\pgfusepath{fill}%
\end{pgfscope}%
\begin{pgfscope}%
\pgfpathrectangle{\pgfqpoint{1.250000in}{0.660000in}}{\pgfqpoint{7.750000in}{4.620000in}}%
\pgfusepath{clip}%
\pgfsetbuttcap%
\pgfsetmiterjoin%
\definecolor{currentfill}{rgb}{0.000000,0.000000,0.000000}%
\pgfsetfillcolor{currentfill}%
\pgfsetlinewidth{0.000000pt}%
\definecolor{currentstroke}{rgb}{0.000000,0.000000,0.000000}%
\pgfsetstrokecolor{currentstroke}%
\pgfsetstrokeopacity{0.000000}%
\pgfsetdash{}{0pt}%
\pgfpathmoveto{\pgfqpoint{1.611150in}{0.660000in}}%
\pgfpathlineto{\pgfqpoint{1.617350in}{0.660000in}}%
\pgfpathlineto{\pgfqpoint{1.617350in}{3.644027in}}%
\pgfpathlineto{\pgfqpoint{1.611150in}{3.644027in}}%
\pgfpathlineto{\pgfqpoint{1.611150in}{0.660000in}}%
\pgfpathclose%
\pgfusepath{fill}%
\end{pgfscope}%
\begin{pgfscope}%
\pgfpathrectangle{\pgfqpoint{1.250000in}{0.660000in}}{\pgfqpoint{7.750000in}{4.620000in}}%
\pgfusepath{clip}%
\pgfsetbuttcap%
\pgfsetmiterjoin%
\definecolor{currentfill}{rgb}{0.000000,0.000000,0.000000}%
\pgfsetfillcolor{currentfill}%
\pgfsetlinewidth{0.000000pt}%
\definecolor{currentstroke}{rgb}{0.000000,0.000000,0.000000}%
\pgfsetstrokecolor{currentstroke}%
\pgfsetstrokeopacity{0.000000}%
\pgfsetdash{}{0pt}%
\pgfpathmoveto{\pgfqpoint{1.618900in}{0.660000in}}%
\pgfpathlineto{\pgfqpoint{1.625100in}{0.660000in}}%
\pgfpathlineto{\pgfqpoint{1.625100in}{3.638716in}}%
\pgfpathlineto{\pgfqpoint{1.618900in}{3.638716in}}%
\pgfpathlineto{\pgfqpoint{1.618900in}{0.660000in}}%
\pgfpathclose%
\pgfusepath{fill}%
\end{pgfscope}%
\begin{pgfscope}%
\pgfpathrectangle{\pgfqpoint{1.250000in}{0.660000in}}{\pgfqpoint{7.750000in}{4.620000in}}%
\pgfusepath{clip}%
\pgfsetbuttcap%
\pgfsetmiterjoin%
\definecolor{currentfill}{rgb}{0.000000,0.000000,0.000000}%
\pgfsetfillcolor{currentfill}%
\pgfsetlinewidth{0.000000pt}%
\definecolor{currentstroke}{rgb}{0.000000,0.000000,0.000000}%
\pgfsetstrokecolor{currentstroke}%
\pgfsetstrokeopacity{0.000000}%
\pgfsetdash{}{0pt}%
\pgfpathmoveto{\pgfqpoint{1.626650in}{0.660000in}}%
\pgfpathlineto{\pgfqpoint{1.632850in}{0.660000in}}%
\pgfpathlineto{\pgfqpoint{1.632850in}{3.633556in}}%
\pgfpathlineto{\pgfqpoint{1.626650in}{3.633556in}}%
\pgfpathlineto{\pgfqpoint{1.626650in}{0.660000in}}%
\pgfpathclose%
\pgfusepath{fill}%
\end{pgfscope}%
\begin{pgfscope}%
\pgfpathrectangle{\pgfqpoint{1.250000in}{0.660000in}}{\pgfqpoint{7.750000in}{4.620000in}}%
\pgfusepath{clip}%
\pgfsetbuttcap%
\pgfsetmiterjoin%
\definecolor{currentfill}{rgb}{0.000000,0.000000,0.000000}%
\pgfsetfillcolor{currentfill}%
\pgfsetlinewidth{0.000000pt}%
\definecolor{currentstroke}{rgb}{0.000000,0.000000,0.000000}%
\pgfsetstrokecolor{currentstroke}%
\pgfsetstrokeopacity{0.000000}%
\pgfsetdash{}{0pt}%
\pgfpathmoveto{\pgfqpoint{1.634400in}{0.660000in}}%
\pgfpathlineto{\pgfqpoint{1.640600in}{0.660000in}}%
\pgfpathlineto{\pgfqpoint{1.640600in}{3.628542in}}%
\pgfpathlineto{\pgfqpoint{1.634400in}{3.628542in}}%
\pgfpathlineto{\pgfqpoint{1.634400in}{0.660000in}}%
\pgfpathclose%
\pgfusepath{fill}%
\end{pgfscope}%
\begin{pgfscope}%
\pgfpathrectangle{\pgfqpoint{1.250000in}{0.660000in}}{\pgfqpoint{7.750000in}{4.620000in}}%
\pgfusepath{clip}%
\pgfsetbuttcap%
\pgfsetmiterjoin%
\definecolor{currentfill}{rgb}{0.000000,0.000000,0.000000}%
\pgfsetfillcolor{currentfill}%
\pgfsetlinewidth{0.000000pt}%
\definecolor{currentstroke}{rgb}{0.000000,0.000000,0.000000}%
\pgfsetstrokecolor{currentstroke}%
\pgfsetstrokeopacity{0.000000}%
\pgfsetdash{}{0pt}%
\pgfpathmoveto{\pgfqpoint{1.642150in}{0.660000in}}%
\pgfpathlineto{\pgfqpoint{1.648350in}{0.660000in}}%
\pgfpathlineto{\pgfqpoint{1.648350in}{3.623665in}}%
\pgfpathlineto{\pgfqpoint{1.642150in}{3.623665in}}%
\pgfpathlineto{\pgfqpoint{1.642150in}{0.660000in}}%
\pgfpathclose%
\pgfusepath{fill}%
\end{pgfscope}%
\begin{pgfscope}%
\pgfpathrectangle{\pgfqpoint{1.250000in}{0.660000in}}{\pgfqpoint{7.750000in}{4.620000in}}%
\pgfusepath{clip}%
\pgfsetbuttcap%
\pgfsetmiterjoin%
\definecolor{currentfill}{rgb}{0.000000,0.000000,0.000000}%
\pgfsetfillcolor{currentfill}%
\pgfsetlinewidth{0.000000pt}%
\definecolor{currentstroke}{rgb}{0.000000,0.000000,0.000000}%
\pgfsetstrokecolor{currentstroke}%
\pgfsetstrokeopacity{0.000000}%
\pgfsetdash{}{0pt}%
\pgfpathmoveto{\pgfqpoint{1.649900in}{0.660000in}}%
\pgfpathlineto{\pgfqpoint{1.656100in}{0.660000in}}%
\pgfpathlineto{\pgfqpoint{1.656100in}{3.618920in}}%
\pgfpathlineto{\pgfqpoint{1.649900in}{3.618920in}}%
\pgfpathlineto{\pgfqpoint{1.649900in}{0.660000in}}%
\pgfpathclose%
\pgfusepath{fill}%
\end{pgfscope}%
\begin{pgfscope}%
\pgfpathrectangle{\pgfqpoint{1.250000in}{0.660000in}}{\pgfqpoint{7.750000in}{4.620000in}}%
\pgfusepath{clip}%
\pgfsetbuttcap%
\pgfsetmiterjoin%
\definecolor{currentfill}{rgb}{0.000000,0.000000,0.000000}%
\pgfsetfillcolor{currentfill}%
\pgfsetlinewidth{0.000000pt}%
\definecolor{currentstroke}{rgb}{0.000000,0.000000,0.000000}%
\pgfsetstrokecolor{currentstroke}%
\pgfsetstrokeopacity{0.000000}%
\pgfsetdash{}{0pt}%
\pgfpathmoveto{\pgfqpoint{1.657650in}{0.660000in}}%
\pgfpathlineto{\pgfqpoint{1.663850in}{0.660000in}}%
\pgfpathlineto{\pgfqpoint{1.663850in}{3.614302in}}%
\pgfpathlineto{\pgfqpoint{1.657650in}{3.614302in}}%
\pgfpathlineto{\pgfqpoint{1.657650in}{0.660000in}}%
\pgfpathclose%
\pgfusepath{fill}%
\end{pgfscope}%
\begin{pgfscope}%
\pgfpathrectangle{\pgfqpoint{1.250000in}{0.660000in}}{\pgfqpoint{7.750000in}{4.620000in}}%
\pgfusepath{clip}%
\pgfsetbuttcap%
\pgfsetmiterjoin%
\definecolor{currentfill}{rgb}{0.000000,0.000000,0.000000}%
\pgfsetfillcolor{currentfill}%
\pgfsetlinewidth{0.000000pt}%
\definecolor{currentstroke}{rgb}{0.000000,0.000000,0.000000}%
\pgfsetstrokecolor{currentstroke}%
\pgfsetstrokeopacity{0.000000}%
\pgfsetdash{}{0pt}%
\pgfpathmoveto{\pgfqpoint{1.665400in}{0.660000in}}%
\pgfpathlineto{\pgfqpoint{1.671600in}{0.660000in}}%
\pgfpathlineto{\pgfqpoint{1.671600in}{3.609803in}}%
\pgfpathlineto{\pgfqpoint{1.665400in}{3.609803in}}%
\pgfpathlineto{\pgfqpoint{1.665400in}{0.660000in}}%
\pgfpathclose%
\pgfusepath{fill}%
\end{pgfscope}%
\begin{pgfscope}%
\pgfpathrectangle{\pgfqpoint{1.250000in}{0.660000in}}{\pgfqpoint{7.750000in}{4.620000in}}%
\pgfusepath{clip}%
\pgfsetbuttcap%
\pgfsetmiterjoin%
\definecolor{currentfill}{rgb}{0.000000,0.000000,0.000000}%
\pgfsetfillcolor{currentfill}%
\pgfsetlinewidth{0.000000pt}%
\definecolor{currentstroke}{rgb}{0.000000,0.000000,0.000000}%
\pgfsetstrokecolor{currentstroke}%
\pgfsetstrokeopacity{0.000000}%
\pgfsetdash{}{0pt}%
\pgfpathmoveto{\pgfqpoint{1.673150in}{0.660000in}}%
\pgfpathlineto{\pgfqpoint{1.679350in}{0.660000in}}%
\pgfpathlineto{\pgfqpoint{1.679350in}{3.605419in}}%
\pgfpathlineto{\pgfqpoint{1.673150in}{3.605419in}}%
\pgfpathlineto{\pgfqpoint{1.673150in}{0.660000in}}%
\pgfpathclose%
\pgfusepath{fill}%
\end{pgfscope}%
\begin{pgfscope}%
\pgfpathrectangle{\pgfqpoint{1.250000in}{0.660000in}}{\pgfqpoint{7.750000in}{4.620000in}}%
\pgfusepath{clip}%
\pgfsetbuttcap%
\pgfsetmiterjoin%
\definecolor{currentfill}{rgb}{0.000000,0.000000,0.000000}%
\pgfsetfillcolor{currentfill}%
\pgfsetlinewidth{0.000000pt}%
\definecolor{currentstroke}{rgb}{0.000000,0.000000,0.000000}%
\pgfsetstrokecolor{currentstroke}%
\pgfsetstrokeopacity{0.000000}%
\pgfsetdash{}{0pt}%
\pgfpathmoveto{\pgfqpoint{1.680900in}{0.660000in}}%
\pgfpathlineto{\pgfqpoint{1.687100in}{0.660000in}}%
\pgfpathlineto{\pgfqpoint{1.687100in}{3.601144in}}%
\pgfpathlineto{\pgfqpoint{1.680900in}{3.601144in}}%
\pgfpathlineto{\pgfqpoint{1.680900in}{0.660000in}}%
\pgfpathclose%
\pgfusepath{fill}%
\end{pgfscope}%
\begin{pgfscope}%
\pgfpathrectangle{\pgfqpoint{1.250000in}{0.660000in}}{\pgfqpoint{7.750000in}{4.620000in}}%
\pgfusepath{clip}%
\pgfsetbuttcap%
\pgfsetmiterjoin%
\definecolor{currentfill}{rgb}{0.000000,0.000000,0.000000}%
\pgfsetfillcolor{currentfill}%
\pgfsetlinewidth{0.000000pt}%
\definecolor{currentstroke}{rgb}{0.000000,0.000000,0.000000}%
\pgfsetstrokecolor{currentstroke}%
\pgfsetstrokeopacity{0.000000}%
\pgfsetdash{}{0pt}%
\pgfpathmoveto{\pgfqpoint{1.688650in}{0.660000in}}%
\pgfpathlineto{\pgfqpoint{1.694850in}{0.660000in}}%
\pgfpathlineto{\pgfqpoint{1.694850in}{3.596975in}}%
\pgfpathlineto{\pgfqpoint{1.688650in}{3.596975in}}%
\pgfpathlineto{\pgfqpoint{1.688650in}{0.660000in}}%
\pgfpathclose%
\pgfusepath{fill}%
\end{pgfscope}%
\begin{pgfscope}%
\pgfpathrectangle{\pgfqpoint{1.250000in}{0.660000in}}{\pgfqpoint{7.750000in}{4.620000in}}%
\pgfusepath{clip}%
\pgfsetbuttcap%
\pgfsetmiterjoin%
\definecolor{currentfill}{rgb}{0.000000,0.000000,0.000000}%
\pgfsetfillcolor{currentfill}%
\pgfsetlinewidth{0.000000pt}%
\definecolor{currentstroke}{rgb}{0.000000,0.000000,0.000000}%
\pgfsetstrokecolor{currentstroke}%
\pgfsetstrokeopacity{0.000000}%
\pgfsetdash{}{0pt}%
\pgfpathmoveto{\pgfqpoint{1.696400in}{0.660000in}}%
\pgfpathlineto{\pgfqpoint{1.702600in}{0.660000in}}%
\pgfpathlineto{\pgfqpoint{1.702600in}{3.592907in}}%
\pgfpathlineto{\pgfqpoint{1.696400in}{3.592907in}}%
\pgfpathlineto{\pgfqpoint{1.696400in}{0.660000in}}%
\pgfpathclose%
\pgfusepath{fill}%
\end{pgfscope}%
\begin{pgfscope}%
\pgfpathrectangle{\pgfqpoint{1.250000in}{0.660000in}}{\pgfqpoint{7.750000in}{4.620000in}}%
\pgfusepath{clip}%
\pgfsetbuttcap%
\pgfsetmiterjoin%
\definecolor{currentfill}{rgb}{0.000000,0.000000,0.000000}%
\pgfsetfillcolor{currentfill}%
\pgfsetlinewidth{0.000000pt}%
\definecolor{currentstroke}{rgb}{0.000000,0.000000,0.000000}%
\pgfsetstrokecolor{currentstroke}%
\pgfsetstrokeopacity{0.000000}%
\pgfsetdash{}{0pt}%
\pgfpathmoveto{\pgfqpoint{1.704150in}{0.660000in}}%
\pgfpathlineto{\pgfqpoint{1.710350in}{0.660000in}}%
\pgfpathlineto{\pgfqpoint{1.710350in}{3.588936in}}%
\pgfpathlineto{\pgfqpoint{1.704150in}{3.588936in}}%
\pgfpathlineto{\pgfqpoint{1.704150in}{0.660000in}}%
\pgfpathclose%
\pgfusepath{fill}%
\end{pgfscope}%
\begin{pgfscope}%
\pgfpathrectangle{\pgfqpoint{1.250000in}{0.660000in}}{\pgfqpoint{7.750000in}{4.620000in}}%
\pgfusepath{clip}%
\pgfsetbuttcap%
\pgfsetmiterjoin%
\definecolor{currentfill}{rgb}{0.000000,0.000000,0.000000}%
\pgfsetfillcolor{currentfill}%
\pgfsetlinewidth{0.000000pt}%
\definecolor{currentstroke}{rgb}{0.000000,0.000000,0.000000}%
\pgfsetstrokecolor{currentstroke}%
\pgfsetstrokeopacity{0.000000}%
\pgfsetdash{}{0pt}%
\pgfpathmoveto{\pgfqpoint{1.711900in}{0.660000in}}%
\pgfpathlineto{\pgfqpoint{1.718100in}{0.660000in}}%
\pgfpathlineto{\pgfqpoint{1.718100in}{3.585058in}}%
\pgfpathlineto{\pgfqpoint{1.711900in}{3.585058in}}%
\pgfpathlineto{\pgfqpoint{1.711900in}{0.660000in}}%
\pgfpathclose%
\pgfusepath{fill}%
\end{pgfscope}%
\begin{pgfscope}%
\pgfpathrectangle{\pgfqpoint{1.250000in}{0.660000in}}{\pgfqpoint{7.750000in}{4.620000in}}%
\pgfusepath{clip}%
\pgfsetbuttcap%
\pgfsetmiterjoin%
\definecolor{currentfill}{rgb}{0.000000,0.000000,0.000000}%
\pgfsetfillcolor{currentfill}%
\pgfsetlinewidth{0.000000pt}%
\definecolor{currentstroke}{rgb}{0.000000,0.000000,0.000000}%
\pgfsetstrokecolor{currentstroke}%
\pgfsetstrokeopacity{0.000000}%
\pgfsetdash{}{0pt}%
\pgfpathmoveto{\pgfqpoint{1.719650in}{0.660000in}}%
\pgfpathlineto{\pgfqpoint{1.725850in}{0.660000in}}%
\pgfpathlineto{\pgfqpoint{1.725850in}{3.581268in}}%
\pgfpathlineto{\pgfqpoint{1.719650in}{3.581268in}}%
\pgfpathlineto{\pgfqpoint{1.719650in}{0.660000in}}%
\pgfpathclose%
\pgfusepath{fill}%
\end{pgfscope}%
\begin{pgfscope}%
\pgfpathrectangle{\pgfqpoint{1.250000in}{0.660000in}}{\pgfqpoint{7.750000in}{4.620000in}}%
\pgfusepath{clip}%
\pgfsetbuttcap%
\pgfsetmiterjoin%
\definecolor{currentfill}{rgb}{0.000000,0.000000,0.000000}%
\pgfsetfillcolor{currentfill}%
\pgfsetlinewidth{0.000000pt}%
\definecolor{currentstroke}{rgb}{0.000000,0.000000,0.000000}%
\pgfsetstrokecolor{currentstroke}%
\pgfsetstrokeopacity{0.000000}%
\pgfsetdash{}{0pt}%
\pgfpathmoveto{\pgfqpoint{1.727400in}{0.660000in}}%
\pgfpathlineto{\pgfqpoint{1.733600in}{0.660000in}}%
\pgfpathlineto{\pgfqpoint{1.733600in}{3.577564in}}%
\pgfpathlineto{\pgfqpoint{1.727400in}{3.577564in}}%
\pgfpathlineto{\pgfqpoint{1.727400in}{0.660000in}}%
\pgfpathclose%
\pgfusepath{fill}%
\end{pgfscope}%
\begin{pgfscope}%
\pgfpathrectangle{\pgfqpoint{1.250000in}{0.660000in}}{\pgfqpoint{7.750000in}{4.620000in}}%
\pgfusepath{clip}%
\pgfsetbuttcap%
\pgfsetmiterjoin%
\definecolor{currentfill}{rgb}{0.000000,0.000000,0.000000}%
\pgfsetfillcolor{currentfill}%
\pgfsetlinewidth{0.000000pt}%
\definecolor{currentstroke}{rgb}{0.000000,0.000000,0.000000}%
\pgfsetstrokecolor{currentstroke}%
\pgfsetstrokeopacity{0.000000}%
\pgfsetdash{}{0pt}%
\pgfpathmoveto{\pgfqpoint{1.735150in}{0.660000in}}%
\pgfpathlineto{\pgfqpoint{1.741350in}{0.660000in}}%
\pgfpathlineto{\pgfqpoint{1.741350in}{3.573943in}}%
\pgfpathlineto{\pgfqpoint{1.735150in}{3.573943in}}%
\pgfpathlineto{\pgfqpoint{1.735150in}{0.660000in}}%
\pgfpathclose%
\pgfusepath{fill}%
\end{pgfscope}%
\begin{pgfscope}%
\pgfpathrectangle{\pgfqpoint{1.250000in}{0.660000in}}{\pgfqpoint{7.750000in}{4.620000in}}%
\pgfusepath{clip}%
\pgfsetbuttcap%
\pgfsetmiterjoin%
\definecolor{currentfill}{rgb}{0.000000,0.000000,0.000000}%
\pgfsetfillcolor{currentfill}%
\pgfsetlinewidth{0.000000pt}%
\definecolor{currentstroke}{rgb}{0.000000,0.000000,0.000000}%
\pgfsetstrokecolor{currentstroke}%
\pgfsetstrokeopacity{0.000000}%
\pgfsetdash{}{0pt}%
\pgfpathmoveto{\pgfqpoint{1.742900in}{0.660000in}}%
\pgfpathlineto{\pgfqpoint{1.749100in}{0.660000in}}%
\pgfpathlineto{\pgfqpoint{1.749100in}{3.570401in}}%
\pgfpathlineto{\pgfqpoint{1.742900in}{3.570401in}}%
\pgfpathlineto{\pgfqpoint{1.742900in}{0.660000in}}%
\pgfpathclose%
\pgfusepath{fill}%
\end{pgfscope}%
\begin{pgfscope}%
\pgfpathrectangle{\pgfqpoint{1.250000in}{0.660000in}}{\pgfqpoint{7.750000in}{4.620000in}}%
\pgfusepath{clip}%
\pgfsetbuttcap%
\pgfsetmiterjoin%
\definecolor{currentfill}{rgb}{0.000000,0.000000,0.000000}%
\pgfsetfillcolor{currentfill}%
\pgfsetlinewidth{0.000000pt}%
\definecolor{currentstroke}{rgb}{0.000000,0.000000,0.000000}%
\pgfsetstrokecolor{currentstroke}%
\pgfsetstrokeopacity{0.000000}%
\pgfsetdash{}{0pt}%
\pgfpathmoveto{\pgfqpoint{1.750650in}{0.660000in}}%
\pgfpathlineto{\pgfqpoint{1.756850in}{0.660000in}}%
\pgfpathlineto{\pgfqpoint{1.756850in}{3.566935in}}%
\pgfpathlineto{\pgfqpoint{1.750650in}{3.566935in}}%
\pgfpathlineto{\pgfqpoint{1.750650in}{0.660000in}}%
\pgfpathclose%
\pgfusepath{fill}%
\end{pgfscope}%
\begin{pgfscope}%
\pgfpathrectangle{\pgfqpoint{1.250000in}{0.660000in}}{\pgfqpoint{7.750000in}{4.620000in}}%
\pgfusepath{clip}%
\pgfsetbuttcap%
\pgfsetmiterjoin%
\definecolor{currentfill}{rgb}{0.000000,0.000000,0.000000}%
\pgfsetfillcolor{currentfill}%
\pgfsetlinewidth{0.000000pt}%
\definecolor{currentstroke}{rgb}{0.000000,0.000000,0.000000}%
\pgfsetstrokecolor{currentstroke}%
\pgfsetstrokeopacity{0.000000}%
\pgfsetdash{}{0pt}%
\pgfpathmoveto{\pgfqpoint{1.758400in}{0.660000in}}%
\pgfpathlineto{\pgfqpoint{1.764600in}{0.660000in}}%
\pgfpathlineto{\pgfqpoint{1.764600in}{3.563544in}}%
\pgfpathlineto{\pgfqpoint{1.758400in}{3.563544in}}%
\pgfpathlineto{\pgfqpoint{1.758400in}{0.660000in}}%
\pgfpathclose%
\pgfusepath{fill}%
\end{pgfscope}%
\begin{pgfscope}%
\pgfpathrectangle{\pgfqpoint{1.250000in}{0.660000in}}{\pgfqpoint{7.750000in}{4.620000in}}%
\pgfusepath{clip}%
\pgfsetbuttcap%
\pgfsetmiterjoin%
\definecolor{currentfill}{rgb}{0.000000,0.000000,0.000000}%
\pgfsetfillcolor{currentfill}%
\pgfsetlinewidth{0.000000pt}%
\definecolor{currentstroke}{rgb}{0.000000,0.000000,0.000000}%
\pgfsetstrokecolor{currentstroke}%
\pgfsetstrokeopacity{0.000000}%
\pgfsetdash{}{0pt}%
\pgfpathmoveto{\pgfqpoint{1.766150in}{0.660000in}}%
\pgfpathlineto{\pgfqpoint{1.772350in}{0.660000in}}%
\pgfpathlineto{\pgfqpoint{1.772350in}{3.560223in}}%
\pgfpathlineto{\pgfqpoint{1.766150in}{3.560223in}}%
\pgfpathlineto{\pgfqpoint{1.766150in}{0.660000in}}%
\pgfpathclose%
\pgfusepath{fill}%
\end{pgfscope}%
\begin{pgfscope}%
\pgfpathrectangle{\pgfqpoint{1.250000in}{0.660000in}}{\pgfqpoint{7.750000in}{4.620000in}}%
\pgfusepath{clip}%
\pgfsetbuttcap%
\pgfsetmiterjoin%
\definecolor{currentfill}{rgb}{0.000000,0.000000,0.000000}%
\pgfsetfillcolor{currentfill}%
\pgfsetlinewidth{0.000000pt}%
\definecolor{currentstroke}{rgb}{0.000000,0.000000,0.000000}%
\pgfsetstrokecolor{currentstroke}%
\pgfsetstrokeopacity{0.000000}%
\pgfsetdash{}{0pt}%
\pgfpathmoveto{\pgfqpoint{1.773900in}{0.660000in}}%
\pgfpathlineto{\pgfqpoint{1.780100in}{0.660000in}}%
\pgfpathlineto{\pgfqpoint{1.780100in}{3.556970in}}%
\pgfpathlineto{\pgfqpoint{1.773900in}{3.556970in}}%
\pgfpathlineto{\pgfqpoint{1.773900in}{0.660000in}}%
\pgfpathclose%
\pgfusepath{fill}%
\end{pgfscope}%
\begin{pgfscope}%
\pgfpathrectangle{\pgfqpoint{1.250000in}{0.660000in}}{\pgfqpoint{7.750000in}{4.620000in}}%
\pgfusepath{clip}%
\pgfsetbuttcap%
\pgfsetmiterjoin%
\definecolor{currentfill}{rgb}{0.000000,0.000000,0.000000}%
\pgfsetfillcolor{currentfill}%
\pgfsetlinewidth{0.000000pt}%
\definecolor{currentstroke}{rgb}{0.000000,0.000000,0.000000}%
\pgfsetstrokecolor{currentstroke}%
\pgfsetstrokeopacity{0.000000}%
\pgfsetdash{}{0pt}%
\pgfpathmoveto{\pgfqpoint{1.781650in}{0.660000in}}%
\pgfpathlineto{\pgfqpoint{1.787850in}{0.660000in}}%
\pgfpathlineto{\pgfqpoint{1.787850in}{3.553783in}}%
\pgfpathlineto{\pgfqpoint{1.781650in}{3.553783in}}%
\pgfpathlineto{\pgfqpoint{1.781650in}{0.660000in}}%
\pgfpathclose%
\pgfusepath{fill}%
\end{pgfscope}%
\begin{pgfscope}%
\pgfpathrectangle{\pgfqpoint{1.250000in}{0.660000in}}{\pgfqpoint{7.750000in}{4.620000in}}%
\pgfusepath{clip}%
\pgfsetbuttcap%
\pgfsetmiterjoin%
\definecolor{currentfill}{rgb}{0.000000,0.000000,0.000000}%
\pgfsetfillcolor{currentfill}%
\pgfsetlinewidth{0.000000pt}%
\definecolor{currentstroke}{rgb}{0.000000,0.000000,0.000000}%
\pgfsetstrokecolor{currentstroke}%
\pgfsetstrokeopacity{0.000000}%
\pgfsetdash{}{0pt}%
\pgfpathmoveto{\pgfqpoint{1.789400in}{0.660000in}}%
\pgfpathlineto{\pgfqpoint{1.795600in}{0.660000in}}%
\pgfpathlineto{\pgfqpoint{1.795600in}{3.550660in}}%
\pgfpathlineto{\pgfqpoint{1.789400in}{3.550660in}}%
\pgfpathlineto{\pgfqpoint{1.789400in}{0.660000in}}%
\pgfpathclose%
\pgfusepath{fill}%
\end{pgfscope}%
\begin{pgfscope}%
\pgfpathrectangle{\pgfqpoint{1.250000in}{0.660000in}}{\pgfqpoint{7.750000in}{4.620000in}}%
\pgfusepath{clip}%
\pgfsetbuttcap%
\pgfsetmiterjoin%
\definecolor{currentfill}{rgb}{0.000000,0.000000,0.000000}%
\pgfsetfillcolor{currentfill}%
\pgfsetlinewidth{0.000000pt}%
\definecolor{currentstroke}{rgb}{0.000000,0.000000,0.000000}%
\pgfsetstrokecolor{currentstroke}%
\pgfsetstrokeopacity{0.000000}%
\pgfsetdash{}{0pt}%
\pgfpathmoveto{\pgfqpoint{1.797150in}{0.660000in}}%
\pgfpathlineto{\pgfqpoint{1.803350in}{0.660000in}}%
\pgfpathlineto{\pgfqpoint{1.803350in}{3.547599in}}%
\pgfpathlineto{\pgfqpoint{1.797150in}{3.547599in}}%
\pgfpathlineto{\pgfqpoint{1.797150in}{0.660000in}}%
\pgfpathclose%
\pgfusepath{fill}%
\end{pgfscope}%
\begin{pgfscope}%
\pgfpathrectangle{\pgfqpoint{1.250000in}{0.660000in}}{\pgfqpoint{7.750000in}{4.620000in}}%
\pgfusepath{clip}%
\pgfsetbuttcap%
\pgfsetmiterjoin%
\definecolor{currentfill}{rgb}{0.000000,0.000000,0.000000}%
\pgfsetfillcolor{currentfill}%
\pgfsetlinewidth{0.000000pt}%
\definecolor{currentstroke}{rgb}{0.000000,0.000000,0.000000}%
\pgfsetstrokecolor{currentstroke}%
\pgfsetstrokeopacity{0.000000}%
\pgfsetdash{}{0pt}%
\pgfpathmoveto{\pgfqpoint{1.804900in}{0.660000in}}%
\pgfpathlineto{\pgfqpoint{1.811100in}{0.660000in}}%
\pgfpathlineto{\pgfqpoint{1.811100in}{3.544597in}}%
\pgfpathlineto{\pgfqpoint{1.804900in}{3.544597in}}%
\pgfpathlineto{\pgfqpoint{1.804900in}{0.660000in}}%
\pgfpathclose%
\pgfusepath{fill}%
\end{pgfscope}%
\begin{pgfscope}%
\pgfpathrectangle{\pgfqpoint{1.250000in}{0.660000in}}{\pgfqpoint{7.750000in}{4.620000in}}%
\pgfusepath{clip}%
\pgfsetbuttcap%
\pgfsetmiterjoin%
\definecolor{currentfill}{rgb}{0.000000,0.000000,0.000000}%
\pgfsetfillcolor{currentfill}%
\pgfsetlinewidth{0.000000pt}%
\definecolor{currentstroke}{rgb}{0.000000,0.000000,0.000000}%
\pgfsetstrokecolor{currentstroke}%
\pgfsetstrokeopacity{0.000000}%
\pgfsetdash{}{0pt}%
\pgfpathmoveto{\pgfqpoint{1.812650in}{0.660000in}}%
\pgfpathlineto{\pgfqpoint{1.818850in}{0.660000in}}%
\pgfpathlineto{\pgfqpoint{1.818850in}{3.541653in}}%
\pgfpathlineto{\pgfqpoint{1.812650in}{3.541653in}}%
\pgfpathlineto{\pgfqpoint{1.812650in}{0.660000in}}%
\pgfpathclose%
\pgfusepath{fill}%
\end{pgfscope}%
\begin{pgfscope}%
\pgfpathrectangle{\pgfqpoint{1.250000in}{0.660000in}}{\pgfqpoint{7.750000in}{4.620000in}}%
\pgfusepath{clip}%
\pgfsetbuttcap%
\pgfsetmiterjoin%
\definecolor{currentfill}{rgb}{0.000000,0.000000,0.000000}%
\pgfsetfillcolor{currentfill}%
\pgfsetlinewidth{0.000000pt}%
\definecolor{currentstroke}{rgb}{0.000000,0.000000,0.000000}%
\pgfsetstrokecolor{currentstroke}%
\pgfsetstrokeopacity{0.000000}%
\pgfsetdash{}{0pt}%
\pgfpathmoveto{\pgfqpoint{1.820400in}{0.660000in}}%
\pgfpathlineto{\pgfqpoint{1.826600in}{0.660000in}}%
\pgfpathlineto{\pgfqpoint{1.826600in}{3.538765in}}%
\pgfpathlineto{\pgfqpoint{1.820400in}{3.538765in}}%
\pgfpathlineto{\pgfqpoint{1.820400in}{0.660000in}}%
\pgfpathclose%
\pgfusepath{fill}%
\end{pgfscope}%
\begin{pgfscope}%
\pgfpathrectangle{\pgfqpoint{1.250000in}{0.660000in}}{\pgfqpoint{7.750000in}{4.620000in}}%
\pgfusepath{clip}%
\pgfsetbuttcap%
\pgfsetmiterjoin%
\definecolor{currentfill}{rgb}{0.000000,0.000000,0.000000}%
\pgfsetfillcolor{currentfill}%
\pgfsetlinewidth{0.000000pt}%
\definecolor{currentstroke}{rgb}{0.000000,0.000000,0.000000}%
\pgfsetstrokecolor{currentstroke}%
\pgfsetstrokeopacity{0.000000}%
\pgfsetdash{}{0pt}%
\pgfpathmoveto{\pgfqpoint{1.828150in}{0.660000in}}%
\pgfpathlineto{\pgfqpoint{1.834350in}{0.660000in}}%
\pgfpathlineto{\pgfqpoint{1.834350in}{3.535931in}}%
\pgfpathlineto{\pgfqpoint{1.828150in}{3.535931in}}%
\pgfpathlineto{\pgfqpoint{1.828150in}{0.660000in}}%
\pgfpathclose%
\pgfusepath{fill}%
\end{pgfscope}%
\begin{pgfscope}%
\pgfpathrectangle{\pgfqpoint{1.250000in}{0.660000in}}{\pgfqpoint{7.750000in}{4.620000in}}%
\pgfusepath{clip}%
\pgfsetbuttcap%
\pgfsetmiterjoin%
\definecolor{currentfill}{rgb}{0.000000,0.000000,0.000000}%
\pgfsetfillcolor{currentfill}%
\pgfsetlinewidth{0.000000pt}%
\definecolor{currentstroke}{rgb}{0.000000,0.000000,0.000000}%
\pgfsetstrokecolor{currentstroke}%
\pgfsetstrokeopacity{0.000000}%
\pgfsetdash{}{0pt}%
\pgfpathmoveto{\pgfqpoint{1.835900in}{0.660000in}}%
\pgfpathlineto{\pgfqpoint{1.842100in}{0.660000in}}%
\pgfpathlineto{\pgfqpoint{1.842100in}{3.533148in}}%
\pgfpathlineto{\pgfqpoint{1.835900in}{3.533148in}}%
\pgfpathlineto{\pgfqpoint{1.835900in}{0.660000in}}%
\pgfpathclose%
\pgfusepath{fill}%
\end{pgfscope}%
\begin{pgfscope}%
\pgfpathrectangle{\pgfqpoint{1.250000in}{0.660000in}}{\pgfqpoint{7.750000in}{4.620000in}}%
\pgfusepath{clip}%
\pgfsetbuttcap%
\pgfsetmiterjoin%
\definecolor{currentfill}{rgb}{0.000000,0.000000,0.000000}%
\pgfsetfillcolor{currentfill}%
\pgfsetlinewidth{0.000000pt}%
\definecolor{currentstroke}{rgb}{0.000000,0.000000,0.000000}%
\pgfsetstrokecolor{currentstroke}%
\pgfsetstrokeopacity{0.000000}%
\pgfsetdash{}{0pt}%
\pgfpathmoveto{\pgfqpoint{1.843650in}{0.660000in}}%
\pgfpathlineto{\pgfqpoint{1.849850in}{0.660000in}}%
\pgfpathlineto{\pgfqpoint{1.849850in}{3.530416in}}%
\pgfpathlineto{\pgfqpoint{1.843650in}{3.530416in}}%
\pgfpathlineto{\pgfqpoint{1.843650in}{0.660000in}}%
\pgfpathclose%
\pgfusepath{fill}%
\end{pgfscope}%
\begin{pgfscope}%
\pgfpathrectangle{\pgfqpoint{1.250000in}{0.660000in}}{\pgfqpoint{7.750000in}{4.620000in}}%
\pgfusepath{clip}%
\pgfsetbuttcap%
\pgfsetmiterjoin%
\definecolor{currentfill}{rgb}{0.000000,0.000000,0.000000}%
\pgfsetfillcolor{currentfill}%
\pgfsetlinewidth{0.000000pt}%
\definecolor{currentstroke}{rgb}{0.000000,0.000000,0.000000}%
\pgfsetstrokecolor{currentstroke}%
\pgfsetstrokeopacity{0.000000}%
\pgfsetdash{}{0pt}%
\pgfpathmoveto{\pgfqpoint{1.851400in}{0.660000in}}%
\pgfpathlineto{\pgfqpoint{1.857600in}{0.660000in}}%
\pgfpathlineto{\pgfqpoint{1.857600in}{3.527734in}}%
\pgfpathlineto{\pgfqpoint{1.851400in}{3.527734in}}%
\pgfpathlineto{\pgfqpoint{1.851400in}{0.660000in}}%
\pgfpathclose%
\pgfusepath{fill}%
\end{pgfscope}%
\begin{pgfscope}%
\pgfpathrectangle{\pgfqpoint{1.250000in}{0.660000in}}{\pgfqpoint{7.750000in}{4.620000in}}%
\pgfusepath{clip}%
\pgfsetbuttcap%
\pgfsetmiterjoin%
\definecolor{currentfill}{rgb}{0.000000,0.000000,0.000000}%
\pgfsetfillcolor{currentfill}%
\pgfsetlinewidth{0.000000pt}%
\definecolor{currentstroke}{rgb}{0.000000,0.000000,0.000000}%
\pgfsetstrokecolor{currentstroke}%
\pgfsetstrokeopacity{0.000000}%
\pgfsetdash{}{0pt}%
\pgfpathmoveto{\pgfqpoint{1.859150in}{0.660000in}}%
\pgfpathlineto{\pgfqpoint{1.865350in}{0.660000in}}%
\pgfpathlineto{\pgfqpoint{1.865350in}{3.525098in}}%
\pgfpathlineto{\pgfqpoint{1.859150in}{3.525098in}}%
\pgfpathlineto{\pgfqpoint{1.859150in}{0.660000in}}%
\pgfpathclose%
\pgfusepath{fill}%
\end{pgfscope}%
\begin{pgfscope}%
\pgfpathrectangle{\pgfqpoint{1.250000in}{0.660000in}}{\pgfqpoint{7.750000in}{4.620000in}}%
\pgfusepath{clip}%
\pgfsetbuttcap%
\pgfsetmiterjoin%
\definecolor{currentfill}{rgb}{0.000000,0.000000,0.000000}%
\pgfsetfillcolor{currentfill}%
\pgfsetlinewidth{0.000000pt}%
\definecolor{currentstroke}{rgb}{0.000000,0.000000,0.000000}%
\pgfsetstrokecolor{currentstroke}%
\pgfsetstrokeopacity{0.000000}%
\pgfsetdash{}{0pt}%
\pgfpathmoveto{\pgfqpoint{1.866900in}{0.660000in}}%
\pgfpathlineto{\pgfqpoint{1.873100in}{0.660000in}}%
\pgfpathlineto{\pgfqpoint{1.873100in}{3.522510in}}%
\pgfpathlineto{\pgfqpoint{1.866900in}{3.522510in}}%
\pgfpathlineto{\pgfqpoint{1.866900in}{0.660000in}}%
\pgfpathclose%
\pgfusepath{fill}%
\end{pgfscope}%
\begin{pgfscope}%
\pgfpathrectangle{\pgfqpoint{1.250000in}{0.660000in}}{\pgfqpoint{7.750000in}{4.620000in}}%
\pgfusepath{clip}%
\pgfsetbuttcap%
\pgfsetmiterjoin%
\definecolor{currentfill}{rgb}{0.000000,0.000000,0.000000}%
\pgfsetfillcolor{currentfill}%
\pgfsetlinewidth{0.000000pt}%
\definecolor{currentstroke}{rgb}{0.000000,0.000000,0.000000}%
\pgfsetstrokecolor{currentstroke}%
\pgfsetstrokeopacity{0.000000}%
\pgfsetdash{}{0pt}%
\pgfpathmoveto{\pgfqpoint{1.874650in}{0.660000in}}%
\pgfpathlineto{\pgfqpoint{1.880850in}{0.660000in}}%
\pgfpathlineto{\pgfqpoint{1.880850in}{3.519965in}}%
\pgfpathlineto{\pgfqpoint{1.874650in}{3.519965in}}%
\pgfpathlineto{\pgfqpoint{1.874650in}{0.660000in}}%
\pgfpathclose%
\pgfusepath{fill}%
\end{pgfscope}%
\begin{pgfscope}%
\pgfpathrectangle{\pgfqpoint{1.250000in}{0.660000in}}{\pgfqpoint{7.750000in}{4.620000in}}%
\pgfusepath{clip}%
\pgfsetbuttcap%
\pgfsetmiterjoin%
\definecolor{currentfill}{rgb}{0.000000,0.000000,0.000000}%
\pgfsetfillcolor{currentfill}%
\pgfsetlinewidth{0.000000pt}%
\definecolor{currentstroke}{rgb}{0.000000,0.000000,0.000000}%
\pgfsetstrokecolor{currentstroke}%
\pgfsetstrokeopacity{0.000000}%
\pgfsetdash{}{0pt}%
\pgfpathmoveto{\pgfqpoint{1.882400in}{0.660000in}}%
\pgfpathlineto{\pgfqpoint{1.888600in}{0.660000in}}%
\pgfpathlineto{\pgfqpoint{1.888600in}{3.517464in}}%
\pgfpathlineto{\pgfqpoint{1.882400in}{3.517464in}}%
\pgfpathlineto{\pgfqpoint{1.882400in}{0.660000in}}%
\pgfpathclose%
\pgfusepath{fill}%
\end{pgfscope}%
\begin{pgfscope}%
\pgfpathrectangle{\pgfqpoint{1.250000in}{0.660000in}}{\pgfqpoint{7.750000in}{4.620000in}}%
\pgfusepath{clip}%
\pgfsetbuttcap%
\pgfsetmiterjoin%
\definecolor{currentfill}{rgb}{0.000000,0.000000,0.000000}%
\pgfsetfillcolor{currentfill}%
\pgfsetlinewidth{0.000000pt}%
\definecolor{currentstroke}{rgb}{0.000000,0.000000,0.000000}%
\pgfsetstrokecolor{currentstroke}%
\pgfsetstrokeopacity{0.000000}%
\pgfsetdash{}{0pt}%
\pgfpathmoveto{\pgfqpoint{1.890150in}{0.660000in}}%
\pgfpathlineto{\pgfqpoint{1.896350in}{0.660000in}}%
\pgfpathlineto{\pgfqpoint{1.896350in}{3.515006in}}%
\pgfpathlineto{\pgfqpoint{1.890150in}{3.515006in}}%
\pgfpathlineto{\pgfqpoint{1.890150in}{0.660000in}}%
\pgfpathclose%
\pgfusepath{fill}%
\end{pgfscope}%
\begin{pgfscope}%
\pgfpathrectangle{\pgfqpoint{1.250000in}{0.660000in}}{\pgfqpoint{7.750000in}{4.620000in}}%
\pgfusepath{clip}%
\pgfsetbuttcap%
\pgfsetmiterjoin%
\definecolor{currentfill}{rgb}{0.000000,0.000000,0.000000}%
\pgfsetfillcolor{currentfill}%
\pgfsetlinewidth{0.000000pt}%
\definecolor{currentstroke}{rgb}{0.000000,0.000000,0.000000}%
\pgfsetstrokecolor{currentstroke}%
\pgfsetstrokeopacity{0.000000}%
\pgfsetdash{}{0pt}%
\pgfpathmoveto{\pgfqpoint{1.897900in}{0.660000in}}%
\pgfpathlineto{\pgfqpoint{1.904100in}{0.660000in}}%
\pgfpathlineto{\pgfqpoint{1.904100in}{3.512587in}}%
\pgfpathlineto{\pgfqpoint{1.897900in}{3.512587in}}%
\pgfpathlineto{\pgfqpoint{1.897900in}{0.660000in}}%
\pgfpathclose%
\pgfusepath{fill}%
\end{pgfscope}%
\begin{pgfscope}%
\pgfpathrectangle{\pgfqpoint{1.250000in}{0.660000in}}{\pgfqpoint{7.750000in}{4.620000in}}%
\pgfusepath{clip}%
\pgfsetbuttcap%
\pgfsetmiterjoin%
\definecolor{currentfill}{rgb}{0.000000,0.000000,0.000000}%
\pgfsetfillcolor{currentfill}%
\pgfsetlinewidth{0.000000pt}%
\definecolor{currentstroke}{rgb}{0.000000,0.000000,0.000000}%
\pgfsetstrokecolor{currentstroke}%
\pgfsetstrokeopacity{0.000000}%
\pgfsetdash{}{0pt}%
\pgfpathmoveto{\pgfqpoint{1.905650in}{0.660000in}}%
\pgfpathlineto{\pgfqpoint{1.911850in}{0.660000in}}%
\pgfpathlineto{\pgfqpoint{1.911850in}{3.510209in}}%
\pgfpathlineto{\pgfqpoint{1.905650in}{3.510209in}}%
\pgfpathlineto{\pgfqpoint{1.905650in}{0.660000in}}%
\pgfpathclose%
\pgfusepath{fill}%
\end{pgfscope}%
\begin{pgfscope}%
\pgfpathrectangle{\pgfqpoint{1.250000in}{0.660000in}}{\pgfqpoint{7.750000in}{4.620000in}}%
\pgfusepath{clip}%
\pgfsetbuttcap%
\pgfsetmiterjoin%
\definecolor{currentfill}{rgb}{0.000000,0.000000,0.000000}%
\pgfsetfillcolor{currentfill}%
\pgfsetlinewidth{0.000000pt}%
\definecolor{currentstroke}{rgb}{0.000000,0.000000,0.000000}%
\pgfsetstrokecolor{currentstroke}%
\pgfsetstrokeopacity{0.000000}%
\pgfsetdash{}{0pt}%
\pgfpathmoveto{\pgfqpoint{1.913400in}{0.660000in}}%
\pgfpathlineto{\pgfqpoint{1.919600in}{0.660000in}}%
\pgfpathlineto{\pgfqpoint{1.919600in}{3.507870in}}%
\pgfpathlineto{\pgfqpoint{1.913400in}{3.507870in}}%
\pgfpathlineto{\pgfqpoint{1.913400in}{0.660000in}}%
\pgfpathclose%
\pgfusepath{fill}%
\end{pgfscope}%
\begin{pgfscope}%
\pgfpathrectangle{\pgfqpoint{1.250000in}{0.660000in}}{\pgfqpoint{7.750000in}{4.620000in}}%
\pgfusepath{clip}%
\pgfsetbuttcap%
\pgfsetmiterjoin%
\definecolor{currentfill}{rgb}{0.000000,0.000000,0.000000}%
\pgfsetfillcolor{currentfill}%
\pgfsetlinewidth{0.000000pt}%
\definecolor{currentstroke}{rgb}{0.000000,0.000000,0.000000}%
\pgfsetstrokecolor{currentstroke}%
\pgfsetstrokeopacity{0.000000}%
\pgfsetdash{}{0pt}%
\pgfpathmoveto{\pgfqpoint{1.921150in}{0.660000in}}%
\pgfpathlineto{\pgfqpoint{1.927350in}{0.660000in}}%
\pgfpathlineto{\pgfqpoint{1.927350in}{3.505568in}}%
\pgfpathlineto{\pgfqpoint{1.921150in}{3.505568in}}%
\pgfpathlineto{\pgfqpoint{1.921150in}{0.660000in}}%
\pgfpathclose%
\pgfusepath{fill}%
\end{pgfscope}%
\begin{pgfscope}%
\pgfpathrectangle{\pgfqpoint{1.250000in}{0.660000in}}{\pgfqpoint{7.750000in}{4.620000in}}%
\pgfusepath{clip}%
\pgfsetbuttcap%
\pgfsetmiterjoin%
\definecolor{currentfill}{rgb}{0.000000,0.000000,0.000000}%
\pgfsetfillcolor{currentfill}%
\pgfsetlinewidth{0.000000pt}%
\definecolor{currentstroke}{rgb}{0.000000,0.000000,0.000000}%
\pgfsetstrokecolor{currentstroke}%
\pgfsetstrokeopacity{0.000000}%
\pgfsetdash{}{0pt}%
\pgfpathmoveto{\pgfqpoint{1.928900in}{0.660000in}}%
\pgfpathlineto{\pgfqpoint{1.935100in}{0.660000in}}%
\pgfpathlineto{\pgfqpoint{1.935100in}{3.503303in}}%
\pgfpathlineto{\pgfqpoint{1.928900in}{3.503303in}}%
\pgfpathlineto{\pgfqpoint{1.928900in}{0.660000in}}%
\pgfpathclose%
\pgfusepath{fill}%
\end{pgfscope}%
\begin{pgfscope}%
\pgfpathrectangle{\pgfqpoint{1.250000in}{0.660000in}}{\pgfqpoint{7.750000in}{4.620000in}}%
\pgfusepath{clip}%
\pgfsetbuttcap%
\pgfsetmiterjoin%
\definecolor{currentfill}{rgb}{0.000000,0.000000,0.000000}%
\pgfsetfillcolor{currentfill}%
\pgfsetlinewidth{0.000000pt}%
\definecolor{currentstroke}{rgb}{0.000000,0.000000,0.000000}%
\pgfsetstrokecolor{currentstroke}%
\pgfsetstrokeopacity{0.000000}%
\pgfsetdash{}{0pt}%
\pgfpathmoveto{\pgfqpoint{1.936650in}{0.660000in}}%
\pgfpathlineto{\pgfqpoint{1.942850in}{0.660000in}}%
\pgfpathlineto{\pgfqpoint{1.942850in}{3.501073in}}%
\pgfpathlineto{\pgfqpoint{1.936650in}{3.501073in}}%
\pgfpathlineto{\pgfqpoint{1.936650in}{0.660000in}}%
\pgfpathclose%
\pgfusepath{fill}%
\end{pgfscope}%
\begin{pgfscope}%
\pgfpathrectangle{\pgfqpoint{1.250000in}{0.660000in}}{\pgfqpoint{7.750000in}{4.620000in}}%
\pgfusepath{clip}%
\pgfsetbuttcap%
\pgfsetmiterjoin%
\definecolor{currentfill}{rgb}{0.000000,0.000000,0.000000}%
\pgfsetfillcolor{currentfill}%
\pgfsetlinewidth{0.000000pt}%
\definecolor{currentstroke}{rgb}{0.000000,0.000000,0.000000}%
\pgfsetstrokecolor{currentstroke}%
\pgfsetstrokeopacity{0.000000}%
\pgfsetdash{}{0pt}%
\pgfpathmoveto{\pgfqpoint{1.944400in}{0.660000in}}%
\pgfpathlineto{\pgfqpoint{1.950600in}{0.660000in}}%
\pgfpathlineto{\pgfqpoint{1.950600in}{3.498878in}}%
\pgfpathlineto{\pgfqpoint{1.944400in}{3.498878in}}%
\pgfpathlineto{\pgfqpoint{1.944400in}{0.660000in}}%
\pgfpathclose%
\pgfusepath{fill}%
\end{pgfscope}%
\begin{pgfscope}%
\pgfpathrectangle{\pgfqpoint{1.250000in}{0.660000in}}{\pgfqpoint{7.750000in}{4.620000in}}%
\pgfusepath{clip}%
\pgfsetbuttcap%
\pgfsetmiterjoin%
\definecolor{currentfill}{rgb}{0.000000,0.000000,0.000000}%
\pgfsetfillcolor{currentfill}%
\pgfsetlinewidth{0.000000pt}%
\definecolor{currentstroke}{rgb}{0.000000,0.000000,0.000000}%
\pgfsetstrokecolor{currentstroke}%
\pgfsetstrokeopacity{0.000000}%
\pgfsetdash{}{0pt}%
\pgfpathmoveto{\pgfqpoint{1.952150in}{0.660000in}}%
\pgfpathlineto{\pgfqpoint{1.958350in}{0.660000in}}%
\pgfpathlineto{\pgfqpoint{1.958350in}{3.496717in}}%
\pgfpathlineto{\pgfqpoint{1.952150in}{3.496717in}}%
\pgfpathlineto{\pgfqpoint{1.952150in}{0.660000in}}%
\pgfpathclose%
\pgfusepath{fill}%
\end{pgfscope}%
\begin{pgfscope}%
\pgfpathrectangle{\pgfqpoint{1.250000in}{0.660000in}}{\pgfqpoint{7.750000in}{4.620000in}}%
\pgfusepath{clip}%
\pgfsetbuttcap%
\pgfsetmiterjoin%
\definecolor{currentfill}{rgb}{0.000000,0.000000,0.000000}%
\pgfsetfillcolor{currentfill}%
\pgfsetlinewidth{0.000000pt}%
\definecolor{currentstroke}{rgb}{0.000000,0.000000,0.000000}%
\pgfsetstrokecolor{currentstroke}%
\pgfsetstrokeopacity{0.000000}%
\pgfsetdash{}{0pt}%
\pgfpathmoveto{\pgfqpoint{1.959900in}{0.660000in}}%
\pgfpathlineto{\pgfqpoint{1.966100in}{0.660000in}}%
\pgfpathlineto{\pgfqpoint{1.966100in}{3.494588in}}%
\pgfpathlineto{\pgfqpoint{1.959900in}{3.494588in}}%
\pgfpathlineto{\pgfqpoint{1.959900in}{0.660000in}}%
\pgfpathclose%
\pgfusepath{fill}%
\end{pgfscope}%
\begin{pgfscope}%
\pgfpathrectangle{\pgfqpoint{1.250000in}{0.660000in}}{\pgfqpoint{7.750000in}{4.620000in}}%
\pgfusepath{clip}%
\pgfsetbuttcap%
\pgfsetmiterjoin%
\definecolor{currentfill}{rgb}{0.000000,0.000000,0.000000}%
\pgfsetfillcolor{currentfill}%
\pgfsetlinewidth{0.000000pt}%
\definecolor{currentstroke}{rgb}{0.000000,0.000000,0.000000}%
\pgfsetstrokecolor{currentstroke}%
\pgfsetstrokeopacity{0.000000}%
\pgfsetdash{}{0pt}%
\pgfpathmoveto{\pgfqpoint{1.967650in}{0.660000in}}%
\pgfpathlineto{\pgfqpoint{1.973850in}{0.660000in}}%
\pgfpathlineto{\pgfqpoint{1.973850in}{3.492491in}}%
\pgfpathlineto{\pgfqpoint{1.967650in}{3.492491in}}%
\pgfpathlineto{\pgfqpoint{1.967650in}{0.660000in}}%
\pgfpathclose%
\pgfusepath{fill}%
\end{pgfscope}%
\begin{pgfscope}%
\pgfpathrectangle{\pgfqpoint{1.250000in}{0.660000in}}{\pgfqpoint{7.750000in}{4.620000in}}%
\pgfusepath{clip}%
\pgfsetbuttcap%
\pgfsetmiterjoin%
\definecolor{currentfill}{rgb}{0.000000,0.000000,0.000000}%
\pgfsetfillcolor{currentfill}%
\pgfsetlinewidth{0.000000pt}%
\definecolor{currentstroke}{rgb}{0.000000,0.000000,0.000000}%
\pgfsetstrokecolor{currentstroke}%
\pgfsetstrokeopacity{0.000000}%
\pgfsetdash{}{0pt}%
\pgfpathmoveto{\pgfqpoint{1.975400in}{0.660000in}}%
\pgfpathlineto{\pgfqpoint{1.981600in}{0.660000in}}%
\pgfpathlineto{\pgfqpoint{1.981600in}{3.490426in}}%
\pgfpathlineto{\pgfqpoint{1.975400in}{3.490426in}}%
\pgfpathlineto{\pgfqpoint{1.975400in}{0.660000in}}%
\pgfpathclose%
\pgfusepath{fill}%
\end{pgfscope}%
\begin{pgfscope}%
\pgfpathrectangle{\pgfqpoint{1.250000in}{0.660000in}}{\pgfqpoint{7.750000in}{4.620000in}}%
\pgfusepath{clip}%
\pgfsetbuttcap%
\pgfsetmiterjoin%
\definecolor{currentfill}{rgb}{0.000000,0.000000,0.000000}%
\pgfsetfillcolor{currentfill}%
\pgfsetlinewidth{0.000000pt}%
\definecolor{currentstroke}{rgb}{0.000000,0.000000,0.000000}%
\pgfsetstrokecolor{currentstroke}%
\pgfsetstrokeopacity{0.000000}%
\pgfsetdash{}{0pt}%
\pgfpathmoveto{\pgfqpoint{1.983150in}{0.660000in}}%
\pgfpathlineto{\pgfqpoint{1.989350in}{0.660000in}}%
\pgfpathlineto{\pgfqpoint{1.989350in}{3.488391in}}%
\pgfpathlineto{\pgfqpoint{1.983150in}{3.488391in}}%
\pgfpathlineto{\pgfqpoint{1.983150in}{0.660000in}}%
\pgfpathclose%
\pgfusepath{fill}%
\end{pgfscope}%
\begin{pgfscope}%
\pgfpathrectangle{\pgfqpoint{1.250000in}{0.660000in}}{\pgfqpoint{7.750000in}{4.620000in}}%
\pgfusepath{clip}%
\pgfsetbuttcap%
\pgfsetmiterjoin%
\definecolor{currentfill}{rgb}{0.000000,0.000000,0.000000}%
\pgfsetfillcolor{currentfill}%
\pgfsetlinewidth{0.000000pt}%
\definecolor{currentstroke}{rgb}{0.000000,0.000000,0.000000}%
\pgfsetstrokecolor{currentstroke}%
\pgfsetstrokeopacity{0.000000}%
\pgfsetdash{}{0pt}%
\pgfpathmoveto{\pgfqpoint{1.990900in}{0.660000in}}%
\pgfpathlineto{\pgfqpoint{1.997100in}{0.660000in}}%
\pgfpathlineto{\pgfqpoint{1.997100in}{3.486384in}}%
\pgfpathlineto{\pgfqpoint{1.990900in}{3.486384in}}%
\pgfpathlineto{\pgfqpoint{1.990900in}{0.660000in}}%
\pgfpathclose%
\pgfusepath{fill}%
\end{pgfscope}%
\begin{pgfscope}%
\pgfpathrectangle{\pgfqpoint{1.250000in}{0.660000in}}{\pgfqpoint{7.750000in}{4.620000in}}%
\pgfusepath{clip}%
\pgfsetbuttcap%
\pgfsetmiterjoin%
\definecolor{currentfill}{rgb}{0.000000,0.000000,0.000000}%
\pgfsetfillcolor{currentfill}%
\pgfsetlinewidth{0.000000pt}%
\definecolor{currentstroke}{rgb}{0.000000,0.000000,0.000000}%
\pgfsetstrokecolor{currentstroke}%
\pgfsetstrokeopacity{0.000000}%
\pgfsetdash{}{0pt}%
\pgfpathmoveto{\pgfqpoint{1.998650in}{0.660000in}}%
\pgfpathlineto{\pgfqpoint{2.004850in}{0.660000in}}%
\pgfpathlineto{\pgfqpoint{2.004850in}{3.484409in}}%
\pgfpathlineto{\pgfqpoint{1.998650in}{3.484409in}}%
\pgfpathlineto{\pgfqpoint{1.998650in}{0.660000in}}%
\pgfpathclose%
\pgfusepath{fill}%
\end{pgfscope}%
\begin{pgfscope}%
\pgfpathrectangle{\pgfqpoint{1.250000in}{0.660000in}}{\pgfqpoint{7.750000in}{4.620000in}}%
\pgfusepath{clip}%
\pgfsetbuttcap%
\pgfsetmiterjoin%
\definecolor{currentfill}{rgb}{0.000000,0.000000,0.000000}%
\pgfsetfillcolor{currentfill}%
\pgfsetlinewidth{0.000000pt}%
\definecolor{currentstroke}{rgb}{0.000000,0.000000,0.000000}%
\pgfsetstrokecolor{currentstroke}%
\pgfsetstrokeopacity{0.000000}%
\pgfsetdash{}{0pt}%
\pgfpathmoveto{\pgfqpoint{2.006400in}{0.660000in}}%
\pgfpathlineto{\pgfqpoint{2.012600in}{0.660000in}}%
\pgfpathlineto{\pgfqpoint{2.012600in}{3.482460in}}%
\pgfpathlineto{\pgfqpoint{2.006400in}{3.482460in}}%
\pgfpathlineto{\pgfqpoint{2.006400in}{0.660000in}}%
\pgfpathclose%
\pgfusepath{fill}%
\end{pgfscope}%
\begin{pgfscope}%
\pgfpathrectangle{\pgfqpoint{1.250000in}{0.660000in}}{\pgfqpoint{7.750000in}{4.620000in}}%
\pgfusepath{clip}%
\pgfsetbuttcap%
\pgfsetmiterjoin%
\definecolor{currentfill}{rgb}{0.000000,0.000000,0.000000}%
\pgfsetfillcolor{currentfill}%
\pgfsetlinewidth{0.000000pt}%
\definecolor{currentstroke}{rgb}{0.000000,0.000000,0.000000}%
\pgfsetstrokecolor{currentstroke}%
\pgfsetstrokeopacity{0.000000}%
\pgfsetdash{}{0pt}%
\pgfpathmoveto{\pgfqpoint{2.014150in}{0.660000in}}%
\pgfpathlineto{\pgfqpoint{2.020350in}{0.660000in}}%
\pgfpathlineto{\pgfqpoint{2.020350in}{3.480539in}}%
\pgfpathlineto{\pgfqpoint{2.014150in}{3.480539in}}%
\pgfpathlineto{\pgfqpoint{2.014150in}{0.660000in}}%
\pgfpathclose%
\pgfusepath{fill}%
\end{pgfscope}%
\begin{pgfscope}%
\pgfpathrectangle{\pgfqpoint{1.250000in}{0.660000in}}{\pgfqpoint{7.750000in}{4.620000in}}%
\pgfusepath{clip}%
\pgfsetbuttcap%
\pgfsetmiterjoin%
\definecolor{currentfill}{rgb}{0.000000,0.000000,0.000000}%
\pgfsetfillcolor{currentfill}%
\pgfsetlinewidth{0.000000pt}%
\definecolor{currentstroke}{rgb}{0.000000,0.000000,0.000000}%
\pgfsetstrokecolor{currentstroke}%
\pgfsetstrokeopacity{0.000000}%
\pgfsetdash{}{0pt}%
\pgfpathmoveto{\pgfqpoint{2.021900in}{0.660000in}}%
\pgfpathlineto{\pgfqpoint{2.028100in}{0.660000in}}%
\pgfpathlineto{\pgfqpoint{2.028100in}{3.478645in}}%
\pgfpathlineto{\pgfqpoint{2.021900in}{3.478645in}}%
\pgfpathlineto{\pgfqpoint{2.021900in}{0.660000in}}%
\pgfpathclose%
\pgfusepath{fill}%
\end{pgfscope}%
\begin{pgfscope}%
\pgfpathrectangle{\pgfqpoint{1.250000in}{0.660000in}}{\pgfqpoint{7.750000in}{4.620000in}}%
\pgfusepath{clip}%
\pgfsetbuttcap%
\pgfsetmiterjoin%
\definecolor{currentfill}{rgb}{0.000000,0.000000,0.000000}%
\pgfsetfillcolor{currentfill}%
\pgfsetlinewidth{0.000000pt}%
\definecolor{currentstroke}{rgb}{0.000000,0.000000,0.000000}%
\pgfsetstrokecolor{currentstroke}%
\pgfsetstrokeopacity{0.000000}%
\pgfsetdash{}{0pt}%
\pgfpathmoveto{\pgfqpoint{2.029650in}{0.660000in}}%
\pgfpathlineto{\pgfqpoint{2.035850in}{0.660000in}}%
\pgfpathlineto{\pgfqpoint{2.035850in}{3.476777in}}%
\pgfpathlineto{\pgfqpoint{2.029650in}{3.476777in}}%
\pgfpathlineto{\pgfqpoint{2.029650in}{0.660000in}}%
\pgfpathclose%
\pgfusepath{fill}%
\end{pgfscope}%
\begin{pgfscope}%
\pgfpathrectangle{\pgfqpoint{1.250000in}{0.660000in}}{\pgfqpoint{7.750000in}{4.620000in}}%
\pgfusepath{clip}%
\pgfsetbuttcap%
\pgfsetmiterjoin%
\definecolor{currentfill}{rgb}{0.000000,0.000000,0.000000}%
\pgfsetfillcolor{currentfill}%
\pgfsetlinewidth{0.000000pt}%
\definecolor{currentstroke}{rgb}{0.000000,0.000000,0.000000}%
\pgfsetstrokecolor{currentstroke}%
\pgfsetstrokeopacity{0.000000}%
\pgfsetdash{}{0pt}%
\pgfpathmoveto{\pgfqpoint{2.037400in}{0.660000in}}%
\pgfpathlineto{\pgfqpoint{2.043600in}{0.660000in}}%
\pgfpathlineto{\pgfqpoint{2.043600in}{3.474935in}}%
\pgfpathlineto{\pgfqpoint{2.037400in}{3.474935in}}%
\pgfpathlineto{\pgfqpoint{2.037400in}{0.660000in}}%
\pgfpathclose%
\pgfusepath{fill}%
\end{pgfscope}%
\begin{pgfscope}%
\pgfpathrectangle{\pgfqpoint{1.250000in}{0.660000in}}{\pgfqpoint{7.750000in}{4.620000in}}%
\pgfusepath{clip}%
\pgfsetbuttcap%
\pgfsetmiterjoin%
\definecolor{currentfill}{rgb}{0.000000,0.000000,0.000000}%
\pgfsetfillcolor{currentfill}%
\pgfsetlinewidth{0.000000pt}%
\definecolor{currentstroke}{rgb}{0.000000,0.000000,0.000000}%
\pgfsetstrokecolor{currentstroke}%
\pgfsetstrokeopacity{0.000000}%
\pgfsetdash{}{0pt}%
\pgfpathmoveto{\pgfqpoint{2.045150in}{0.660000in}}%
\pgfpathlineto{\pgfqpoint{2.051350in}{0.660000in}}%
\pgfpathlineto{\pgfqpoint{2.051350in}{3.473117in}}%
\pgfpathlineto{\pgfqpoint{2.045150in}{3.473117in}}%
\pgfpathlineto{\pgfqpoint{2.045150in}{0.660000in}}%
\pgfpathclose%
\pgfusepath{fill}%
\end{pgfscope}%
\begin{pgfscope}%
\pgfpathrectangle{\pgfqpoint{1.250000in}{0.660000in}}{\pgfqpoint{7.750000in}{4.620000in}}%
\pgfusepath{clip}%
\pgfsetbuttcap%
\pgfsetmiterjoin%
\definecolor{currentfill}{rgb}{0.000000,0.000000,0.000000}%
\pgfsetfillcolor{currentfill}%
\pgfsetlinewidth{0.000000pt}%
\definecolor{currentstroke}{rgb}{0.000000,0.000000,0.000000}%
\pgfsetstrokecolor{currentstroke}%
\pgfsetstrokeopacity{0.000000}%
\pgfsetdash{}{0pt}%
\pgfpathmoveto{\pgfqpoint{2.052900in}{0.660000in}}%
\pgfpathlineto{\pgfqpoint{2.059100in}{0.660000in}}%
\pgfpathlineto{\pgfqpoint{2.059100in}{3.471325in}}%
\pgfpathlineto{\pgfqpoint{2.052900in}{3.471325in}}%
\pgfpathlineto{\pgfqpoint{2.052900in}{0.660000in}}%
\pgfpathclose%
\pgfusepath{fill}%
\end{pgfscope}%
\begin{pgfscope}%
\pgfpathrectangle{\pgfqpoint{1.250000in}{0.660000in}}{\pgfqpoint{7.750000in}{4.620000in}}%
\pgfusepath{clip}%
\pgfsetbuttcap%
\pgfsetmiterjoin%
\definecolor{currentfill}{rgb}{0.000000,0.000000,0.000000}%
\pgfsetfillcolor{currentfill}%
\pgfsetlinewidth{0.000000pt}%
\definecolor{currentstroke}{rgb}{0.000000,0.000000,0.000000}%
\pgfsetstrokecolor{currentstroke}%
\pgfsetstrokeopacity{0.000000}%
\pgfsetdash{}{0pt}%
\pgfpathmoveto{\pgfqpoint{2.060650in}{0.660000in}}%
\pgfpathlineto{\pgfqpoint{2.066850in}{0.660000in}}%
\pgfpathlineto{\pgfqpoint{2.066850in}{3.469555in}}%
\pgfpathlineto{\pgfqpoint{2.060650in}{3.469555in}}%
\pgfpathlineto{\pgfqpoint{2.060650in}{0.660000in}}%
\pgfpathclose%
\pgfusepath{fill}%
\end{pgfscope}%
\begin{pgfscope}%
\pgfpathrectangle{\pgfqpoint{1.250000in}{0.660000in}}{\pgfqpoint{7.750000in}{4.620000in}}%
\pgfusepath{clip}%
\pgfsetbuttcap%
\pgfsetmiterjoin%
\definecolor{currentfill}{rgb}{0.000000,0.000000,0.000000}%
\pgfsetfillcolor{currentfill}%
\pgfsetlinewidth{0.000000pt}%
\definecolor{currentstroke}{rgb}{0.000000,0.000000,0.000000}%
\pgfsetstrokecolor{currentstroke}%
\pgfsetstrokeopacity{0.000000}%
\pgfsetdash{}{0pt}%
\pgfpathmoveto{\pgfqpoint{2.068400in}{0.660000in}}%
\pgfpathlineto{\pgfqpoint{2.074600in}{0.660000in}}%
\pgfpathlineto{\pgfqpoint{2.074600in}{3.467810in}}%
\pgfpathlineto{\pgfqpoint{2.068400in}{3.467810in}}%
\pgfpathlineto{\pgfqpoint{2.068400in}{0.660000in}}%
\pgfpathclose%
\pgfusepath{fill}%
\end{pgfscope}%
\begin{pgfscope}%
\pgfpathrectangle{\pgfqpoint{1.250000in}{0.660000in}}{\pgfqpoint{7.750000in}{4.620000in}}%
\pgfusepath{clip}%
\pgfsetbuttcap%
\pgfsetmiterjoin%
\definecolor{currentfill}{rgb}{0.000000,0.000000,0.000000}%
\pgfsetfillcolor{currentfill}%
\pgfsetlinewidth{0.000000pt}%
\definecolor{currentstroke}{rgb}{0.000000,0.000000,0.000000}%
\pgfsetstrokecolor{currentstroke}%
\pgfsetstrokeopacity{0.000000}%
\pgfsetdash{}{0pt}%
\pgfpathmoveto{\pgfqpoint{2.076150in}{0.660000in}}%
\pgfpathlineto{\pgfqpoint{2.082350in}{0.660000in}}%
\pgfpathlineto{\pgfqpoint{2.082350in}{3.466086in}}%
\pgfpathlineto{\pgfqpoint{2.076150in}{3.466086in}}%
\pgfpathlineto{\pgfqpoint{2.076150in}{0.660000in}}%
\pgfpathclose%
\pgfusepath{fill}%
\end{pgfscope}%
\begin{pgfscope}%
\pgfpathrectangle{\pgfqpoint{1.250000in}{0.660000in}}{\pgfqpoint{7.750000in}{4.620000in}}%
\pgfusepath{clip}%
\pgfsetbuttcap%
\pgfsetmiterjoin%
\definecolor{currentfill}{rgb}{0.000000,0.000000,0.000000}%
\pgfsetfillcolor{currentfill}%
\pgfsetlinewidth{0.000000pt}%
\definecolor{currentstroke}{rgb}{0.000000,0.000000,0.000000}%
\pgfsetstrokecolor{currentstroke}%
\pgfsetstrokeopacity{0.000000}%
\pgfsetdash{}{0pt}%
\pgfpathmoveto{\pgfqpoint{2.083900in}{0.660000in}}%
\pgfpathlineto{\pgfqpoint{2.090100in}{0.660000in}}%
\pgfpathlineto{\pgfqpoint{2.090100in}{3.464386in}}%
\pgfpathlineto{\pgfqpoint{2.083900in}{3.464386in}}%
\pgfpathlineto{\pgfqpoint{2.083900in}{0.660000in}}%
\pgfpathclose%
\pgfusepath{fill}%
\end{pgfscope}%
\begin{pgfscope}%
\pgfpathrectangle{\pgfqpoint{1.250000in}{0.660000in}}{\pgfqpoint{7.750000in}{4.620000in}}%
\pgfusepath{clip}%
\pgfsetbuttcap%
\pgfsetmiterjoin%
\definecolor{currentfill}{rgb}{0.000000,0.000000,0.000000}%
\pgfsetfillcolor{currentfill}%
\pgfsetlinewidth{0.000000pt}%
\definecolor{currentstroke}{rgb}{0.000000,0.000000,0.000000}%
\pgfsetstrokecolor{currentstroke}%
\pgfsetstrokeopacity{0.000000}%
\pgfsetdash{}{0pt}%
\pgfpathmoveto{\pgfqpoint{2.091650in}{0.660000in}}%
\pgfpathlineto{\pgfqpoint{2.097850in}{0.660000in}}%
\pgfpathlineto{\pgfqpoint{2.097850in}{3.462706in}}%
\pgfpathlineto{\pgfqpoint{2.091650in}{3.462706in}}%
\pgfpathlineto{\pgfqpoint{2.091650in}{0.660000in}}%
\pgfpathclose%
\pgfusepath{fill}%
\end{pgfscope}%
\begin{pgfscope}%
\pgfpathrectangle{\pgfqpoint{1.250000in}{0.660000in}}{\pgfqpoint{7.750000in}{4.620000in}}%
\pgfusepath{clip}%
\pgfsetbuttcap%
\pgfsetmiterjoin%
\definecolor{currentfill}{rgb}{0.000000,0.000000,0.000000}%
\pgfsetfillcolor{currentfill}%
\pgfsetlinewidth{0.000000pt}%
\definecolor{currentstroke}{rgb}{0.000000,0.000000,0.000000}%
\pgfsetstrokecolor{currentstroke}%
\pgfsetstrokeopacity{0.000000}%
\pgfsetdash{}{0pt}%
\pgfpathmoveto{\pgfqpoint{2.099400in}{0.660000in}}%
\pgfpathlineto{\pgfqpoint{2.105600in}{0.660000in}}%
\pgfpathlineto{\pgfqpoint{2.105600in}{3.461049in}}%
\pgfpathlineto{\pgfqpoint{2.099400in}{3.461049in}}%
\pgfpathlineto{\pgfqpoint{2.099400in}{0.660000in}}%
\pgfpathclose%
\pgfusepath{fill}%
\end{pgfscope}%
\begin{pgfscope}%
\pgfpathrectangle{\pgfqpoint{1.250000in}{0.660000in}}{\pgfqpoint{7.750000in}{4.620000in}}%
\pgfusepath{clip}%
\pgfsetbuttcap%
\pgfsetmiterjoin%
\definecolor{currentfill}{rgb}{0.000000,0.000000,0.000000}%
\pgfsetfillcolor{currentfill}%
\pgfsetlinewidth{0.000000pt}%
\definecolor{currentstroke}{rgb}{0.000000,0.000000,0.000000}%
\pgfsetstrokecolor{currentstroke}%
\pgfsetstrokeopacity{0.000000}%
\pgfsetdash{}{0pt}%
\pgfpathmoveto{\pgfqpoint{2.107150in}{0.660000in}}%
\pgfpathlineto{\pgfqpoint{2.113350in}{0.660000in}}%
\pgfpathlineto{\pgfqpoint{2.113350in}{3.459411in}}%
\pgfpathlineto{\pgfqpoint{2.107150in}{3.459411in}}%
\pgfpathlineto{\pgfqpoint{2.107150in}{0.660000in}}%
\pgfpathclose%
\pgfusepath{fill}%
\end{pgfscope}%
\begin{pgfscope}%
\pgfpathrectangle{\pgfqpoint{1.250000in}{0.660000in}}{\pgfqpoint{7.750000in}{4.620000in}}%
\pgfusepath{clip}%
\pgfsetbuttcap%
\pgfsetmiterjoin%
\definecolor{currentfill}{rgb}{0.000000,0.000000,0.000000}%
\pgfsetfillcolor{currentfill}%
\pgfsetlinewidth{0.000000pt}%
\definecolor{currentstroke}{rgb}{0.000000,0.000000,0.000000}%
\pgfsetstrokecolor{currentstroke}%
\pgfsetstrokeopacity{0.000000}%
\pgfsetdash{}{0pt}%
\pgfpathmoveto{\pgfqpoint{2.114900in}{0.660000in}}%
\pgfpathlineto{\pgfqpoint{2.121100in}{0.660000in}}%
\pgfpathlineto{\pgfqpoint{2.121100in}{3.457795in}}%
\pgfpathlineto{\pgfqpoint{2.114900in}{3.457795in}}%
\pgfpathlineto{\pgfqpoint{2.114900in}{0.660000in}}%
\pgfpathclose%
\pgfusepath{fill}%
\end{pgfscope}%
\begin{pgfscope}%
\pgfpathrectangle{\pgfqpoint{1.250000in}{0.660000in}}{\pgfqpoint{7.750000in}{4.620000in}}%
\pgfusepath{clip}%
\pgfsetbuttcap%
\pgfsetmiterjoin%
\definecolor{currentfill}{rgb}{0.000000,0.000000,0.000000}%
\pgfsetfillcolor{currentfill}%
\pgfsetlinewidth{0.000000pt}%
\definecolor{currentstroke}{rgb}{0.000000,0.000000,0.000000}%
\pgfsetstrokecolor{currentstroke}%
\pgfsetstrokeopacity{0.000000}%
\pgfsetdash{}{0pt}%
\pgfpathmoveto{\pgfqpoint{2.122650in}{0.660000in}}%
\pgfpathlineto{\pgfqpoint{2.128850in}{0.660000in}}%
\pgfpathlineto{\pgfqpoint{2.128850in}{3.456199in}}%
\pgfpathlineto{\pgfqpoint{2.122650in}{3.456199in}}%
\pgfpathlineto{\pgfqpoint{2.122650in}{0.660000in}}%
\pgfpathclose%
\pgfusepath{fill}%
\end{pgfscope}%
\begin{pgfscope}%
\pgfpathrectangle{\pgfqpoint{1.250000in}{0.660000in}}{\pgfqpoint{7.750000in}{4.620000in}}%
\pgfusepath{clip}%
\pgfsetbuttcap%
\pgfsetmiterjoin%
\definecolor{currentfill}{rgb}{0.000000,0.000000,0.000000}%
\pgfsetfillcolor{currentfill}%
\pgfsetlinewidth{0.000000pt}%
\definecolor{currentstroke}{rgb}{0.000000,0.000000,0.000000}%
\pgfsetstrokecolor{currentstroke}%
\pgfsetstrokeopacity{0.000000}%
\pgfsetdash{}{0pt}%
\pgfpathmoveto{\pgfqpoint{2.130400in}{0.660000in}}%
\pgfpathlineto{\pgfqpoint{2.136600in}{0.660000in}}%
\pgfpathlineto{\pgfqpoint{2.136600in}{3.454621in}}%
\pgfpathlineto{\pgfqpoint{2.130400in}{3.454621in}}%
\pgfpathlineto{\pgfqpoint{2.130400in}{0.660000in}}%
\pgfpathclose%
\pgfusepath{fill}%
\end{pgfscope}%
\begin{pgfscope}%
\pgfpathrectangle{\pgfqpoint{1.250000in}{0.660000in}}{\pgfqpoint{7.750000in}{4.620000in}}%
\pgfusepath{clip}%
\pgfsetbuttcap%
\pgfsetmiterjoin%
\definecolor{currentfill}{rgb}{0.000000,0.000000,0.000000}%
\pgfsetfillcolor{currentfill}%
\pgfsetlinewidth{0.000000pt}%
\definecolor{currentstroke}{rgb}{0.000000,0.000000,0.000000}%
\pgfsetstrokecolor{currentstroke}%
\pgfsetstrokeopacity{0.000000}%
\pgfsetdash{}{0pt}%
\pgfpathmoveto{\pgfqpoint{2.138150in}{0.660000in}}%
\pgfpathlineto{\pgfqpoint{2.144350in}{0.660000in}}%
\pgfpathlineto{\pgfqpoint{2.144350in}{3.453064in}}%
\pgfpathlineto{\pgfqpoint{2.138150in}{3.453064in}}%
\pgfpathlineto{\pgfqpoint{2.138150in}{0.660000in}}%
\pgfpathclose%
\pgfusepath{fill}%
\end{pgfscope}%
\begin{pgfscope}%
\pgfpathrectangle{\pgfqpoint{1.250000in}{0.660000in}}{\pgfqpoint{7.750000in}{4.620000in}}%
\pgfusepath{clip}%
\pgfsetbuttcap%
\pgfsetmiterjoin%
\definecolor{currentfill}{rgb}{0.000000,0.000000,0.000000}%
\pgfsetfillcolor{currentfill}%
\pgfsetlinewidth{0.000000pt}%
\definecolor{currentstroke}{rgb}{0.000000,0.000000,0.000000}%
\pgfsetstrokecolor{currentstroke}%
\pgfsetstrokeopacity{0.000000}%
\pgfsetdash{}{0pt}%
\pgfpathmoveto{\pgfqpoint{2.145900in}{0.660000in}}%
\pgfpathlineto{\pgfqpoint{2.152100in}{0.660000in}}%
\pgfpathlineto{\pgfqpoint{2.152100in}{3.451524in}}%
\pgfpathlineto{\pgfqpoint{2.145900in}{3.451524in}}%
\pgfpathlineto{\pgfqpoint{2.145900in}{0.660000in}}%
\pgfpathclose%
\pgfusepath{fill}%
\end{pgfscope}%
\begin{pgfscope}%
\pgfpathrectangle{\pgfqpoint{1.250000in}{0.660000in}}{\pgfqpoint{7.750000in}{4.620000in}}%
\pgfusepath{clip}%
\pgfsetbuttcap%
\pgfsetmiterjoin%
\definecolor{currentfill}{rgb}{0.000000,0.000000,0.000000}%
\pgfsetfillcolor{currentfill}%
\pgfsetlinewidth{0.000000pt}%
\definecolor{currentstroke}{rgb}{0.000000,0.000000,0.000000}%
\pgfsetstrokecolor{currentstroke}%
\pgfsetstrokeopacity{0.000000}%
\pgfsetdash{}{0pt}%
\pgfpathmoveto{\pgfqpoint{2.153650in}{0.660000in}}%
\pgfpathlineto{\pgfqpoint{2.159850in}{0.660000in}}%
\pgfpathlineto{\pgfqpoint{2.159850in}{3.450003in}}%
\pgfpathlineto{\pgfqpoint{2.153650in}{3.450003in}}%
\pgfpathlineto{\pgfqpoint{2.153650in}{0.660000in}}%
\pgfpathclose%
\pgfusepath{fill}%
\end{pgfscope}%
\begin{pgfscope}%
\pgfpathrectangle{\pgfqpoint{1.250000in}{0.660000in}}{\pgfqpoint{7.750000in}{4.620000in}}%
\pgfusepath{clip}%
\pgfsetbuttcap%
\pgfsetmiterjoin%
\definecolor{currentfill}{rgb}{0.000000,0.000000,0.000000}%
\pgfsetfillcolor{currentfill}%
\pgfsetlinewidth{0.000000pt}%
\definecolor{currentstroke}{rgb}{0.000000,0.000000,0.000000}%
\pgfsetstrokecolor{currentstroke}%
\pgfsetstrokeopacity{0.000000}%
\pgfsetdash{}{0pt}%
\pgfpathmoveto{\pgfqpoint{2.161400in}{0.660000in}}%
\pgfpathlineto{\pgfqpoint{2.167600in}{0.660000in}}%
\pgfpathlineto{\pgfqpoint{2.167600in}{3.448501in}}%
\pgfpathlineto{\pgfqpoint{2.161400in}{3.448501in}}%
\pgfpathlineto{\pgfqpoint{2.161400in}{0.660000in}}%
\pgfpathclose%
\pgfusepath{fill}%
\end{pgfscope}%
\begin{pgfscope}%
\pgfpathrectangle{\pgfqpoint{1.250000in}{0.660000in}}{\pgfqpoint{7.750000in}{4.620000in}}%
\pgfusepath{clip}%
\pgfsetbuttcap%
\pgfsetmiterjoin%
\definecolor{currentfill}{rgb}{0.000000,0.000000,0.000000}%
\pgfsetfillcolor{currentfill}%
\pgfsetlinewidth{0.000000pt}%
\definecolor{currentstroke}{rgb}{0.000000,0.000000,0.000000}%
\pgfsetstrokecolor{currentstroke}%
\pgfsetstrokeopacity{0.000000}%
\pgfsetdash{}{0pt}%
\pgfpathmoveto{\pgfqpoint{2.169150in}{0.660000in}}%
\pgfpathlineto{\pgfqpoint{2.175350in}{0.660000in}}%
\pgfpathlineto{\pgfqpoint{2.175350in}{3.447016in}}%
\pgfpathlineto{\pgfqpoint{2.169150in}{3.447016in}}%
\pgfpathlineto{\pgfqpoint{2.169150in}{0.660000in}}%
\pgfpathclose%
\pgfusepath{fill}%
\end{pgfscope}%
\begin{pgfscope}%
\pgfpathrectangle{\pgfqpoint{1.250000in}{0.660000in}}{\pgfqpoint{7.750000in}{4.620000in}}%
\pgfusepath{clip}%
\pgfsetbuttcap%
\pgfsetmiterjoin%
\definecolor{currentfill}{rgb}{0.000000,0.000000,0.000000}%
\pgfsetfillcolor{currentfill}%
\pgfsetlinewidth{0.000000pt}%
\definecolor{currentstroke}{rgb}{0.000000,0.000000,0.000000}%
\pgfsetstrokecolor{currentstroke}%
\pgfsetstrokeopacity{0.000000}%
\pgfsetdash{}{0pt}%
\pgfpathmoveto{\pgfqpoint{2.176900in}{0.660000in}}%
\pgfpathlineto{\pgfqpoint{2.183100in}{0.660000in}}%
\pgfpathlineto{\pgfqpoint{2.183100in}{3.445547in}}%
\pgfpathlineto{\pgfqpoint{2.176900in}{3.445547in}}%
\pgfpathlineto{\pgfqpoint{2.176900in}{0.660000in}}%
\pgfpathclose%
\pgfusepath{fill}%
\end{pgfscope}%
\begin{pgfscope}%
\pgfpathrectangle{\pgfqpoint{1.250000in}{0.660000in}}{\pgfqpoint{7.750000in}{4.620000in}}%
\pgfusepath{clip}%
\pgfsetbuttcap%
\pgfsetmiterjoin%
\definecolor{currentfill}{rgb}{0.000000,0.000000,0.000000}%
\pgfsetfillcolor{currentfill}%
\pgfsetlinewidth{0.000000pt}%
\definecolor{currentstroke}{rgb}{0.000000,0.000000,0.000000}%
\pgfsetstrokecolor{currentstroke}%
\pgfsetstrokeopacity{0.000000}%
\pgfsetdash{}{0pt}%
\pgfpathmoveto{\pgfqpoint{2.184650in}{0.660000in}}%
\pgfpathlineto{\pgfqpoint{2.190850in}{0.660000in}}%
\pgfpathlineto{\pgfqpoint{2.190850in}{3.444095in}}%
\pgfpathlineto{\pgfqpoint{2.184650in}{3.444095in}}%
\pgfpathlineto{\pgfqpoint{2.184650in}{0.660000in}}%
\pgfpathclose%
\pgfusepath{fill}%
\end{pgfscope}%
\begin{pgfscope}%
\pgfpathrectangle{\pgfqpoint{1.250000in}{0.660000in}}{\pgfqpoint{7.750000in}{4.620000in}}%
\pgfusepath{clip}%
\pgfsetbuttcap%
\pgfsetmiterjoin%
\definecolor{currentfill}{rgb}{0.000000,0.000000,0.000000}%
\pgfsetfillcolor{currentfill}%
\pgfsetlinewidth{0.000000pt}%
\definecolor{currentstroke}{rgb}{0.000000,0.000000,0.000000}%
\pgfsetstrokecolor{currentstroke}%
\pgfsetstrokeopacity{0.000000}%
\pgfsetdash{}{0pt}%
\pgfpathmoveto{\pgfqpoint{2.192400in}{0.660000in}}%
\pgfpathlineto{\pgfqpoint{2.198600in}{0.660000in}}%
\pgfpathlineto{\pgfqpoint{2.198600in}{3.442662in}}%
\pgfpathlineto{\pgfqpoint{2.192400in}{3.442662in}}%
\pgfpathlineto{\pgfqpoint{2.192400in}{0.660000in}}%
\pgfpathclose%
\pgfusepath{fill}%
\end{pgfscope}%
\begin{pgfscope}%
\pgfpathrectangle{\pgfqpoint{1.250000in}{0.660000in}}{\pgfqpoint{7.750000in}{4.620000in}}%
\pgfusepath{clip}%
\pgfsetbuttcap%
\pgfsetmiterjoin%
\definecolor{currentfill}{rgb}{0.000000,0.000000,0.000000}%
\pgfsetfillcolor{currentfill}%
\pgfsetlinewidth{0.000000pt}%
\definecolor{currentstroke}{rgb}{0.000000,0.000000,0.000000}%
\pgfsetstrokecolor{currentstroke}%
\pgfsetstrokeopacity{0.000000}%
\pgfsetdash{}{0pt}%
\pgfpathmoveto{\pgfqpoint{2.200150in}{0.660000in}}%
\pgfpathlineto{\pgfqpoint{2.206350in}{0.660000in}}%
\pgfpathlineto{\pgfqpoint{2.206350in}{3.441244in}}%
\pgfpathlineto{\pgfqpoint{2.200150in}{3.441244in}}%
\pgfpathlineto{\pgfqpoint{2.200150in}{0.660000in}}%
\pgfpathclose%
\pgfusepath{fill}%
\end{pgfscope}%
\begin{pgfscope}%
\pgfpathrectangle{\pgfqpoint{1.250000in}{0.660000in}}{\pgfqpoint{7.750000in}{4.620000in}}%
\pgfusepath{clip}%
\pgfsetbuttcap%
\pgfsetmiterjoin%
\definecolor{currentfill}{rgb}{0.000000,0.000000,0.000000}%
\pgfsetfillcolor{currentfill}%
\pgfsetlinewidth{0.000000pt}%
\definecolor{currentstroke}{rgb}{0.000000,0.000000,0.000000}%
\pgfsetstrokecolor{currentstroke}%
\pgfsetstrokeopacity{0.000000}%
\pgfsetdash{}{0pt}%
\pgfpathmoveto{\pgfqpoint{2.207900in}{0.660000in}}%
\pgfpathlineto{\pgfqpoint{2.214100in}{0.660000in}}%
\pgfpathlineto{\pgfqpoint{2.214100in}{3.439841in}}%
\pgfpathlineto{\pgfqpoint{2.207900in}{3.439841in}}%
\pgfpathlineto{\pgfqpoint{2.207900in}{0.660000in}}%
\pgfpathclose%
\pgfusepath{fill}%
\end{pgfscope}%
\begin{pgfscope}%
\pgfpathrectangle{\pgfqpoint{1.250000in}{0.660000in}}{\pgfqpoint{7.750000in}{4.620000in}}%
\pgfusepath{clip}%
\pgfsetbuttcap%
\pgfsetmiterjoin%
\definecolor{currentfill}{rgb}{0.000000,0.000000,0.000000}%
\pgfsetfillcolor{currentfill}%
\pgfsetlinewidth{0.000000pt}%
\definecolor{currentstroke}{rgb}{0.000000,0.000000,0.000000}%
\pgfsetstrokecolor{currentstroke}%
\pgfsetstrokeopacity{0.000000}%
\pgfsetdash{}{0pt}%
\pgfpathmoveto{\pgfqpoint{2.215650in}{0.660000in}}%
\pgfpathlineto{\pgfqpoint{2.221850in}{0.660000in}}%
\pgfpathlineto{\pgfqpoint{2.221850in}{3.438454in}}%
\pgfpathlineto{\pgfqpoint{2.215650in}{3.438454in}}%
\pgfpathlineto{\pgfqpoint{2.215650in}{0.660000in}}%
\pgfpathclose%
\pgfusepath{fill}%
\end{pgfscope}%
\begin{pgfscope}%
\pgfpathrectangle{\pgfqpoint{1.250000in}{0.660000in}}{\pgfqpoint{7.750000in}{4.620000in}}%
\pgfusepath{clip}%
\pgfsetbuttcap%
\pgfsetmiterjoin%
\definecolor{currentfill}{rgb}{0.000000,0.000000,0.000000}%
\pgfsetfillcolor{currentfill}%
\pgfsetlinewidth{0.000000pt}%
\definecolor{currentstroke}{rgb}{0.000000,0.000000,0.000000}%
\pgfsetstrokecolor{currentstroke}%
\pgfsetstrokeopacity{0.000000}%
\pgfsetdash{}{0pt}%
\pgfpathmoveto{\pgfqpoint{2.223400in}{0.660000in}}%
\pgfpathlineto{\pgfqpoint{2.229600in}{0.660000in}}%
\pgfpathlineto{\pgfqpoint{2.229600in}{3.437082in}}%
\pgfpathlineto{\pgfqpoint{2.223400in}{3.437082in}}%
\pgfpathlineto{\pgfqpoint{2.223400in}{0.660000in}}%
\pgfpathclose%
\pgfusepath{fill}%
\end{pgfscope}%
\begin{pgfscope}%
\pgfpathrectangle{\pgfqpoint{1.250000in}{0.660000in}}{\pgfqpoint{7.750000in}{4.620000in}}%
\pgfusepath{clip}%
\pgfsetbuttcap%
\pgfsetmiterjoin%
\definecolor{currentfill}{rgb}{0.000000,0.000000,0.000000}%
\pgfsetfillcolor{currentfill}%
\pgfsetlinewidth{0.000000pt}%
\definecolor{currentstroke}{rgb}{0.000000,0.000000,0.000000}%
\pgfsetstrokecolor{currentstroke}%
\pgfsetstrokeopacity{0.000000}%
\pgfsetdash{}{0pt}%
\pgfpathmoveto{\pgfqpoint{2.231150in}{0.660000in}}%
\pgfpathlineto{\pgfqpoint{2.237350in}{0.660000in}}%
\pgfpathlineto{\pgfqpoint{2.237350in}{3.435726in}}%
\pgfpathlineto{\pgfqpoint{2.231150in}{3.435726in}}%
\pgfpathlineto{\pgfqpoint{2.231150in}{0.660000in}}%
\pgfpathclose%
\pgfusepath{fill}%
\end{pgfscope}%
\begin{pgfscope}%
\pgfpathrectangle{\pgfqpoint{1.250000in}{0.660000in}}{\pgfqpoint{7.750000in}{4.620000in}}%
\pgfusepath{clip}%
\pgfsetbuttcap%
\pgfsetmiterjoin%
\definecolor{currentfill}{rgb}{0.000000,0.000000,0.000000}%
\pgfsetfillcolor{currentfill}%
\pgfsetlinewidth{0.000000pt}%
\definecolor{currentstroke}{rgb}{0.000000,0.000000,0.000000}%
\pgfsetstrokecolor{currentstroke}%
\pgfsetstrokeopacity{0.000000}%
\pgfsetdash{}{0pt}%
\pgfpathmoveto{\pgfqpoint{2.238900in}{0.660000in}}%
\pgfpathlineto{\pgfqpoint{2.245100in}{0.660000in}}%
\pgfpathlineto{\pgfqpoint{2.245100in}{3.434385in}}%
\pgfpathlineto{\pgfqpoint{2.238900in}{3.434385in}}%
\pgfpathlineto{\pgfqpoint{2.238900in}{0.660000in}}%
\pgfpathclose%
\pgfusepath{fill}%
\end{pgfscope}%
\begin{pgfscope}%
\pgfpathrectangle{\pgfqpoint{1.250000in}{0.660000in}}{\pgfqpoint{7.750000in}{4.620000in}}%
\pgfusepath{clip}%
\pgfsetbuttcap%
\pgfsetmiterjoin%
\definecolor{currentfill}{rgb}{0.000000,0.000000,0.000000}%
\pgfsetfillcolor{currentfill}%
\pgfsetlinewidth{0.000000pt}%
\definecolor{currentstroke}{rgb}{0.000000,0.000000,0.000000}%
\pgfsetstrokecolor{currentstroke}%
\pgfsetstrokeopacity{0.000000}%
\pgfsetdash{}{0pt}%
\pgfpathmoveto{\pgfqpoint{2.246650in}{0.660000in}}%
\pgfpathlineto{\pgfqpoint{2.252850in}{0.660000in}}%
\pgfpathlineto{\pgfqpoint{2.252850in}{3.433058in}}%
\pgfpathlineto{\pgfqpoint{2.246650in}{3.433058in}}%
\pgfpathlineto{\pgfqpoint{2.246650in}{0.660000in}}%
\pgfpathclose%
\pgfusepath{fill}%
\end{pgfscope}%
\begin{pgfscope}%
\pgfpathrectangle{\pgfqpoint{1.250000in}{0.660000in}}{\pgfqpoint{7.750000in}{4.620000in}}%
\pgfusepath{clip}%
\pgfsetbuttcap%
\pgfsetmiterjoin%
\definecolor{currentfill}{rgb}{0.000000,0.000000,0.000000}%
\pgfsetfillcolor{currentfill}%
\pgfsetlinewidth{0.000000pt}%
\definecolor{currentstroke}{rgb}{0.000000,0.000000,0.000000}%
\pgfsetstrokecolor{currentstroke}%
\pgfsetstrokeopacity{0.000000}%
\pgfsetdash{}{0pt}%
\pgfpathmoveto{\pgfqpoint{2.254400in}{0.660000in}}%
\pgfpathlineto{\pgfqpoint{2.260600in}{0.660000in}}%
\pgfpathlineto{\pgfqpoint{2.260600in}{3.431745in}}%
\pgfpathlineto{\pgfqpoint{2.254400in}{3.431745in}}%
\pgfpathlineto{\pgfqpoint{2.254400in}{0.660000in}}%
\pgfpathclose%
\pgfusepath{fill}%
\end{pgfscope}%
\begin{pgfscope}%
\pgfpathrectangle{\pgfqpoint{1.250000in}{0.660000in}}{\pgfqpoint{7.750000in}{4.620000in}}%
\pgfusepath{clip}%
\pgfsetbuttcap%
\pgfsetmiterjoin%
\definecolor{currentfill}{rgb}{0.000000,0.000000,0.000000}%
\pgfsetfillcolor{currentfill}%
\pgfsetlinewidth{0.000000pt}%
\definecolor{currentstroke}{rgb}{0.000000,0.000000,0.000000}%
\pgfsetstrokecolor{currentstroke}%
\pgfsetstrokeopacity{0.000000}%
\pgfsetdash{}{0pt}%
\pgfpathmoveto{\pgfqpoint{2.262150in}{0.660000in}}%
\pgfpathlineto{\pgfqpoint{2.268350in}{0.660000in}}%
\pgfpathlineto{\pgfqpoint{2.268350in}{3.430446in}}%
\pgfpathlineto{\pgfqpoint{2.262150in}{3.430446in}}%
\pgfpathlineto{\pgfqpoint{2.262150in}{0.660000in}}%
\pgfpathclose%
\pgfusepath{fill}%
\end{pgfscope}%
\begin{pgfscope}%
\pgfpathrectangle{\pgfqpoint{1.250000in}{0.660000in}}{\pgfqpoint{7.750000in}{4.620000in}}%
\pgfusepath{clip}%
\pgfsetbuttcap%
\pgfsetmiterjoin%
\definecolor{currentfill}{rgb}{0.000000,0.000000,0.000000}%
\pgfsetfillcolor{currentfill}%
\pgfsetlinewidth{0.000000pt}%
\definecolor{currentstroke}{rgb}{0.000000,0.000000,0.000000}%
\pgfsetstrokecolor{currentstroke}%
\pgfsetstrokeopacity{0.000000}%
\pgfsetdash{}{0pt}%
\pgfpathmoveto{\pgfqpoint{2.269900in}{0.660000in}}%
\pgfpathlineto{\pgfqpoint{2.276100in}{0.660000in}}%
\pgfpathlineto{\pgfqpoint{2.276100in}{3.429161in}}%
\pgfpathlineto{\pgfqpoint{2.269900in}{3.429161in}}%
\pgfpathlineto{\pgfqpoint{2.269900in}{0.660000in}}%
\pgfpathclose%
\pgfusepath{fill}%
\end{pgfscope}%
\begin{pgfscope}%
\pgfpathrectangle{\pgfqpoint{1.250000in}{0.660000in}}{\pgfqpoint{7.750000in}{4.620000in}}%
\pgfusepath{clip}%
\pgfsetbuttcap%
\pgfsetmiterjoin%
\definecolor{currentfill}{rgb}{0.000000,0.000000,0.000000}%
\pgfsetfillcolor{currentfill}%
\pgfsetlinewidth{0.000000pt}%
\definecolor{currentstroke}{rgb}{0.000000,0.000000,0.000000}%
\pgfsetstrokecolor{currentstroke}%
\pgfsetstrokeopacity{0.000000}%
\pgfsetdash{}{0pt}%
\pgfpathmoveto{\pgfqpoint{2.277650in}{0.660000in}}%
\pgfpathlineto{\pgfqpoint{2.283850in}{0.660000in}}%
\pgfpathlineto{\pgfqpoint{2.283850in}{3.427889in}}%
\pgfpathlineto{\pgfqpoint{2.277650in}{3.427889in}}%
\pgfpathlineto{\pgfqpoint{2.277650in}{0.660000in}}%
\pgfpathclose%
\pgfusepath{fill}%
\end{pgfscope}%
\begin{pgfscope}%
\pgfpathrectangle{\pgfqpoint{1.250000in}{0.660000in}}{\pgfqpoint{7.750000in}{4.620000in}}%
\pgfusepath{clip}%
\pgfsetbuttcap%
\pgfsetmiterjoin%
\definecolor{currentfill}{rgb}{0.000000,0.000000,0.000000}%
\pgfsetfillcolor{currentfill}%
\pgfsetlinewidth{0.000000pt}%
\definecolor{currentstroke}{rgb}{0.000000,0.000000,0.000000}%
\pgfsetstrokecolor{currentstroke}%
\pgfsetstrokeopacity{0.000000}%
\pgfsetdash{}{0pt}%
\pgfpathmoveto{\pgfqpoint{2.285400in}{0.660000in}}%
\pgfpathlineto{\pgfqpoint{2.291600in}{0.660000in}}%
\pgfpathlineto{\pgfqpoint{2.291600in}{3.426631in}}%
\pgfpathlineto{\pgfqpoint{2.285400in}{3.426631in}}%
\pgfpathlineto{\pgfqpoint{2.285400in}{0.660000in}}%
\pgfpathclose%
\pgfusepath{fill}%
\end{pgfscope}%
\begin{pgfscope}%
\pgfpathrectangle{\pgfqpoint{1.250000in}{0.660000in}}{\pgfqpoint{7.750000in}{4.620000in}}%
\pgfusepath{clip}%
\pgfsetbuttcap%
\pgfsetmiterjoin%
\definecolor{currentfill}{rgb}{0.000000,0.000000,0.000000}%
\pgfsetfillcolor{currentfill}%
\pgfsetlinewidth{0.000000pt}%
\definecolor{currentstroke}{rgb}{0.000000,0.000000,0.000000}%
\pgfsetstrokecolor{currentstroke}%
\pgfsetstrokeopacity{0.000000}%
\pgfsetdash{}{0pt}%
\pgfpathmoveto{\pgfqpoint{2.293150in}{0.660000in}}%
\pgfpathlineto{\pgfqpoint{2.299350in}{0.660000in}}%
\pgfpathlineto{\pgfqpoint{2.299350in}{3.425385in}}%
\pgfpathlineto{\pgfqpoint{2.293150in}{3.425385in}}%
\pgfpathlineto{\pgfqpoint{2.293150in}{0.660000in}}%
\pgfpathclose%
\pgfusepath{fill}%
\end{pgfscope}%
\begin{pgfscope}%
\pgfpathrectangle{\pgfqpoint{1.250000in}{0.660000in}}{\pgfqpoint{7.750000in}{4.620000in}}%
\pgfusepath{clip}%
\pgfsetbuttcap%
\pgfsetmiterjoin%
\definecolor{currentfill}{rgb}{0.000000,0.000000,0.000000}%
\pgfsetfillcolor{currentfill}%
\pgfsetlinewidth{0.000000pt}%
\definecolor{currentstroke}{rgb}{0.000000,0.000000,0.000000}%
\pgfsetstrokecolor{currentstroke}%
\pgfsetstrokeopacity{0.000000}%
\pgfsetdash{}{0pt}%
\pgfpathmoveto{\pgfqpoint{2.300900in}{0.660000in}}%
\pgfpathlineto{\pgfqpoint{2.307100in}{0.660000in}}%
\pgfpathlineto{\pgfqpoint{2.307100in}{3.424152in}}%
\pgfpathlineto{\pgfqpoint{2.300900in}{3.424152in}}%
\pgfpathlineto{\pgfqpoint{2.300900in}{0.660000in}}%
\pgfpathclose%
\pgfusepath{fill}%
\end{pgfscope}%
\begin{pgfscope}%
\pgfpathrectangle{\pgfqpoint{1.250000in}{0.660000in}}{\pgfqpoint{7.750000in}{4.620000in}}%
\pgfusepath{clip}%
\pgfsetbuttcap%
\pgfsetmiterjoin%
\definecolor{currentfill}{rgb}{0.000000,0.000000,0.000000}%
\pgfsetfillcolor{currentfill}%
\pgfsetlinewidth{0.000000pt}%
\definecolor{currentstroke}{rgb}{0.000000,0.000000,0.000000}%
\pgfsetstrokecolor{currentstroke}%
\pgfsetstrokeopacity{0.000000}%
\pgfsetdash{}{0pt}%
\pgfpathmoveto{\pgfqpoint{2.308650in}{0.660000in}}%
\pgfpathlineto{\pgfqpoint{2.314850in}{0.660000in}}%
\pgfpathlineto{\pgfqpoint{2.314850in}{3.422932in}}%
\pgfpathlineto{\pgfqpoint{2.308650in}{3.422932in}}%
\pgfpathlineto{\pgfqpoint{2.308650in}{0.660000in}}%
\pgfpathclose%
\pgfusepath{fill}%
\end{pgfscope}%
\begin{pgfscope}%
\pgfpathrectangle{\pgfqpoint{1.250000in}{0.660000in}}{\pgfqpoint{7.750000in}{4.620000in}}%
\pgfusepath{clip}%
\pgfsetbuttcap%
\pgfsetmiterjoin%
\definecolor{currentfill}{rgb}{0.000000,0.000000,0.000000}%
\pgfsetfillcolor{currentfill}%
\pgfsetlinewidth{0.000000pt}%
\definecolor{currentstroke}{rgb}{0.000000,0.000000,0.000000}%
\pgfsetstrokecolor{currentstroke}%
\pgfsetstrokeopacity{0.000000}%
\pgfsetdash{}{0pt}%
\pgfpathmoveto{\pgfqpoint{2.316400in}{0.660000in}}%
\pgfpathlineto{\pgfqpoint{2.322600in}{0.660000in}}%
\pgfpathlineto{\pgfqpoint{2.322600in}{3.421724in}}%
\pgfpathlineto{\pgfqpoint{2.316400in}{3.421724in}}%
\pgfpathlineto{\pgfqpoint{2.316400in}{0.660000in}}%
\pgfpathclose%
\pgfusepath{fill}%
\end{pgfscope}%
\begin{pgfscope}%
\pgfpathrectangle{\pgfqpoint{1.250000in}{0.660000in}}{\pgfqpoint{7.750000in}{4.620000in}}%
\pgfusepath{clip}%
\pgfsetbuttcap%
\pgfsetmiterjoin%
\definecolor{currentfill}{rgb}{0.000000,0.000000,0.000000}%
\pgfsetfillcolor{currentfill}%
\pgfsetlinewidth{0.000000pt}%
\definecolor{currentstroke}{rgb}{0.000000,0.000000,0.000000}%
\pgfsetstrokecolor{currentstroke}%
\pgfsetstrokeopacity{0.000000}%
\pgfsetdash{}{0pt}%
\pgfpathmoveto{\pgfqpoint{2.324150in}{0.660000in}}%
\pgfpathlineto{\pgfqpoint{2.330350in}{0.660000in}}%
\pgfpathlineto{\pgfqpoint{2.330350in}{3.420528in}}%
\pgfpathlineto{\pgfqpoint{2.324150in}{3.420528in}}%
\pgfpathlineto{\pgfqpoint{2.324150in}{0.660000in}}%
\pgfpathclose%
\pgfusepath{fill}%
\end{pgfscope}%
\begin{pgfscope}%
\pgfpathrectangle{\pgfqpoint{1.250000in}{0.660000in}}{\pgfqpoint{7.750000in}{4.620000in}}%
\pgfusepath{clip}%
\pgfsetbuttcap%
\pgfsetmiterjoin%
\definecolor{currentfill}{rgb}{0.000000,0.000000,0.000000}%
\pgfsetfillcolor{currentfill}%
\pgfsetlinewidth{0.000000pt}%
\definecolor{currentstroke}{rgb}{0.000000,0.000000,0.000000}%
\pgfsetstrokecolor{currentstroke}%
\pgfsetstrokeopacity{0.000000}%
\pgfsetdash{}{0pt}%
\pgfpathmoveto{\pgfqpoint{2.331900in}{0.660000in}}%
\pgfpathlineto{\pgfqpoint{2.338100in}{0.660000in}}%
\pgfpathlineto{\pgfqpoint{2.338100in}{3.419345in}}%
\pgfpathlineto{\pgfqpoint{2.331900in}{3.419345in}}%
\pgfpathlineto{\pgfqpoint{2.331900in}{0.660000in}}%
\pgfpathclose%
\pgfusepath{fill}%
\end{pgfscope}%
\begin{pgfscope}%
\pgfpathrectangle{\pgfqpoint{1.250000in}{0.660000in}}{\pgfqpoint{7.750000in}{4.620000in}}%
\pgfusepath{clip}%
\pgfsetbuttcap%
\pgfsetmiterjoin%
\definecolor{currentfill}{rgb}{0.000000,0.000000,0.000000}%
\pgfsetfillcolor{currentfill}%
\pgfsetlinewidth{0.000000pt}%
\definecolor{currentstroke}{rgb}{0.000000,0.000000,0.000000}%
\pgfsetstrokecolor{currentstroke}%
\pgfsetstrokeopacity{0.000000}%
\pgfsetdash{}{0pt}%
\pgfpathmoveto{\pgfqpoint{2.339650in}{0.660000in}}%
\pgfpathlineto{\pgfqpoint{2.345850in}{0.660000in}}%
\pgfpathlineto{\pgfqpoint{2.345850in}{3.418172in}}%
\pgfpathlineto{\pgfqpoint{2.339650in}{3.418172in}}%
\pgfpathlineto{\pgfqpoint{2.339650in}{0.660000in}}%
\pgfpathclose%
\pgfusepath{fill}%
\end{pgfscope}%
\begin{pgfscope}%
\pgfpathrectangle{\pgfqpoint{1.250000in}{0.660000in}}{\pgfqpoint{7.750000in}{4.620000in}}%
\pgfusepath{clip}%
\pgfsetbuttcap%
\pgfsetmiterjoin%
\definecolor{currentfill}{rgb}{0.000000,0.000000,0.000000}%
\pgfsetfillcolor{currentfill}%
\pgfsetlinewidth{0.000000pt}%
\definecolor{currentstroke}{rgb}{0.000000,0.000000,0.000000}%
\pgfsetstrokecolor{currentstroke}%
\pgfsetstrokeopacity{0.000000}%
\pgfsetdash{}{0pt}%
\pgfpathmoveto{\pgfqpoint{2.347400in}{0.660000in}}%
\pgfpathlineto{\pgfqpoint{2.353600in}{0.660000in}}%
\pgfpathlineto{\pgfqpoint{2.353600in}{3.417012in}}%
\pgfpathlineto{\pgfqpoint{2.347400in}{3.417012in}}%
\pgfpathlineto{\pgfqpoint{2.347400in}{0.660000in}}%
\pgfpathclose%
\pgfusepath{fill}%
\end{pgfscope}%
\begin{pgfscope}%
\pgfpathrectangle{\pgfqpoint{1.250000in}{0.660000in}}{\pgfqpoint{7.750000in}{4.620000in}}%
\pgfusepath{clip}%
\pgfsetbuttcap%
\pgfsetmiterjoin%
\definecolor{currentfill}{rgb}{0.000000,0.000000,0.000000}%
\pgfsetfillcolor{currentfill}%
\pgfsetlinewidth{0.000000pt}%
\definecolor{currentstroke}{rgb}{0.000000,0.000000,0.000000}%
\pgfsetstrokecolor{currentstroke}%
\pgfsetstrokeopacity{0.000000}%
\pgfsetdash{}{0pt}%
\pgfpathmoveto{\pgfqpoint{2.355150in}{0.660000in}}%
\pgfpathlineto{\pgfqpoint{2.361350in}{0.660000in}}%
\pgfpathlineto{\pgfqpoint{2.361350in}{3.415862in}}%
\pgfpathlineto{\pgfqpoint{2.355150in}{3.415862in}}%
\pgfpathlineto{\pgfqpoint{2.355150in}{0.660000in}}%
\pgfpathclose%
\pgfusepath{fill}%
\end{pgfscope}%
\begin{pgfscope}%
\pgfpathrectangle{\pgfqpoint{1.250000in}{0.660000in}}{\pgfqpoint{7.750000in}{4.620000in}}%
\pgfusepath{clip}%
\pgfsetbuttcap%
\pgfsetmiterjoin%
\definecolor{currentfill}{rgb}{0.000000,0.000000,0.000000}%
\pgfsetfillcolor{currentfill}%
\pgfsetlinewidth{0.000000pt}%
\definecolor{currentstroke}{rgb}{0.000000,0.000000,0.000000}%
\pgfsetstrokecolor{currentstroke}%
\pgfsetstrokeopacity{0.000000}%
\pgfsetdash{}{0pt}%
\pgfpathmoveto{\pgfqpoint{2.362900in}{0.660000in}}%
\pgfpathlineto{\pgfqpoint{2.369100in}{0.660000in}}%
\pgfpathlineto{\pgfqpoint{2.369100in}{3.414724in}}%
\pgfpathlineto{\pgfqpoint{2.362900in}{3.414724in}}%
\pgfpathlineto{\pgfqpoint{2.362900in}{0.660000in}}%
\pgfpathclose%
\pgfusepath{fill}%
\end{pgfscope}%
\begin{pgfscope}%
\pgfpathrectangle{\pgfqpoint{1.250000in}{0.660000in}}{\pgfqpoint{7.750000in}{4.620000in}}%
\pgfusepath{clip}%
\pgfsetbuttcap%
\pgfsetmiterjoin%
\definecolor{currentfill}{rgb}{0.000000,0.000000,0.000000}%
\pgfsetfillcolor{currentfill}%
\pgfsetlinewidth{0.000000pt}%
\definecolor{currentstroke}{rgb}{0.000000,0.000000,0.000000}%
\pgfsetstrokecolor{currentstroke}%
\pgfsetstrokeopacity{0.000000}%
\pgfsetdash{}{0pt}%
\pgfpathmoveto{\pgfqpoint{2.370650in}{0.660000in}}%
\pgfpathlineto{\pgfqpoint{2.376850in}{0.660000in}}%
\pgfpathlineto{\pgfqpoint{2.376850in}{3.413596in}}%
\pgfpathlineto{\pgfqpoint{2.370650in}{3.413596in}}%
\pgfpathlineto{\pgfqpoint{2.370650in}{0.660000in}}%
\pgfpathclose%
\pgfusepath{fill}%
\end{pgfscope}%
\begin{pgfscope}%
\pgfpathrectangle{\pgfqpoint{1.250000in}{0.660000in}}{\pgfqpoint{7.750000in}{4.620000in}}%
\pgfusepath{clip}%
\pgfsetbuttcap%
\pgfsetmiterjoin%
\definecolor{currentfill}{rgb}{0.000000,0.000000,0.000000}%
\pgfsetfillcolor{currentfill}%
\pgfsetlinewidth{0.000000pt}%
\definecolor{currentstroke}{rgb}{0.000000,0.000000,0.000000}%
\pgfsetstrokecolor{currentstroke}%
\pgfsetstrokeopacity{0.000000}%
\pgfsetdash{}{0pt}%
\pgfpathmoveto{\pgfqpoint{2.378400in}{0.660000in}}%
\pgfpathlineto{\pgfqpoint{2.384600in}{0.660000in}}%
\pgfpathlineto{\pgfqpoint{2.384600in}{3.412479in}}%
\pgfpathlineto{\pgfqpoint{2.378400in}{3.412479in}}%
\pgfpathlineto{\pgfqpoint{2.378400in}{0.660000in}}%
\pgfpathclose%
\pgfusepath{fill}%
\end{pgfscope}%
\begin{pgfscope}%
\pgfpathrectangle{\pgfqpoint{1.250000in}{0.660000in}}{\pgfqpoint{7.750000in}{4.620000in}}%
\pgfusepath{clip}%
\pgfsetbuttcap%
\pgfsetmiterjoin%
\definecolor{currentfill}{rgb}{0.000000,0.000000,0.000000}%
\pgfsetfillcolor{currentfill}%
\pgfsetlinewidth{0.000000pt}%
\definecolor{currentstroke}{rgb}{0.000000,0.000000,0.000000}%
\pgfsetstrokecolor{currentstroke}%
\pgfsetstrokeopacity{0.000000}%
\pgfsetdash{}{0pt}%
\pgfpathmoveto{\pgfqpoint{2.386150in}{0.660000in}}%
\pgfpathlineto{\pgfqpoint{2.392350in}{0.660000in}}%
\pgfpathlineto{\pgfqpoint{2.392350in}{3.411372in}}%
\pgfpathlineto{\pgfqpoint{2.386150in}{3.411372in}}%
\pgfpathlineto{\pgfqpoint{2.386150in}{0.660000in}}%
\pgfpathclose%
\pgfusepath{fill}%
\end{pgfscope}%
\begin{pgfscope}%
\pgfpathrectangle{\pgfqpoint{1.250000in}{0.660000in}}{\pgfqpoint{7.750000in}{4.620000in}}%
\pgfusepath{clip}%
\pgfsetbuttcap%
\pgfsetmiterjoin%
\definecolor{currentfill}{rgb}{0.000000,0.000000,0.000000}%
\pgfsetfillcolor{currentfill}%
\pgfsetlinewidth{0.000000pt}%
\definecolor{currentstroke}{rgb}{0.000000,0.000000,0.000000}%
\pgfsetstrokecolor{currentstroke}%
\pgfsetstrokeopacity{0.000000}%
\pgfsetdash{}{0pt}%
\pgfpathmoveto{\pgfqpoint{2.393900in}{0.660000in}}%
\pgfpathlineto{\pgfqpoint{2.400100in}{0.660000in}}%
\pgfpathlineto{\pgfqpoint{2.400100in}{3.410277in}}%
\pgfpathlineto{\pgfqpoint{2.393900in}{3.410277in}}%
\pgfpathlineto{\pgfqpoint{2.393900in}{0.660000in}}%
\pgfpathclose%
\pgfusepath{fill}%
\end{pgfscope}%
\begin{pgfscope}%
\pgfpathrectangle{\pgfqpoint{1.250000in}{0.660000in}}{\pgfqpoint{7.750000in}{4.620000in}}%
\pgfusepath{clip}%
\pgfsetbuttcap%
\pgfsetmiterjoin%
\definecolor{currentfill}{rgb}{0.000000,0.000000,0.000000}%
\pgfsetfillcolor{currentfill}%
\pgfsetlinewidth{0.000000pt}%
\definecolor{currentstroke}{rgb}{0.000000,0.000000,0.000000}%
\pgfsetstrokecolor{currentstroke}%
\pgfsetstrokeopacity{0.000000}%
\pgfsetdash{}{0pt}%
\pgfpathmoveto{\pgfqpoint{2.401650in}{0.660000in}}%
\pgfpathlineto{\pgfqpoint{2.407850in}{0.660000in}}%
\pgfpathlineto{\pgfqpoint{2.407850in}{3.409192in}}%
\pgfpathlineto{\pgfqpoint{2.401650in}{3.409192in}}%
\pgfpathlineto{\pgfqpoint{2.401650in}{0.660000in}}%
\pgfpathclose%
\pgfusepath{fill}%
\end{pgfscope}%
\begin{pgfscope}%
\pgfpathrectangle{\pgfqpoint{1.250000in}{0.660000in}}{\pgfqpoint{7.750000in}{4.620000in}}%
\pgfusepath{clip}%
\pgfsetbuttcap%
\pgfsetmiterjoin%
\definecolor{currentfill}{rgb}{0.000000,0.000000,0.000000}%
\pgfsetfillcolor{currentfill}%
\pgfsetlinewidth{0.000000pt}%
\definecolor{currentstroke}{rgb}{0.000000,0.000000,0.000000}%
\pgfsetstrokecolor{currentstroke}%
\pgfsetstrokeopacity{0.000000}%
\pgfsetdash{}{0pt}%
\pgfpathmoveto{\pgfqpoint{2.409400in}{0.660000in}}%
\pgfpathlineto{\pgfqpoint{2.415600in}{0.660000in}}%
\pgfpathlineto{\pgfqpoint{2.415600in}{3.408117in}}%
\pgfpathlineto{\pgfqpoint{2.409400in}{3.408117in}}%
\pgfpathlineto{\pgfqpoint{2.409400in}{0.660000in}}%
\pgfpathclose%
\pgfusepath{fill}%
\end{pgfscope}%
\begin{pgfscope}%
\pgfpathrectangle{\pgfqpoint{1.250000in}{0.660000in}}{\pgfqpoint{7.750000in}{4.620000in}}%
\pgfusepath{clip}%
\pgfsetbuttcap%
\pgfsetmiterjoin%
\definecolor{currentfill}{rgb}{0.000000,0.000000,0.000000}%
\pgfsetfillcolor{currentfill}%
\pgfsetlinewidth{0.000000pt}%
\definecolor{currentstroke}{rgb}{0.000000,0.000000,0.000000}%
\pgfsetstrokecolor{currentstroke}%
\pgfsetstrokeopacity{0.000000}%
\pgfsetdash{}{0pt}%
\pgfpathmoveto{\pgfqpoint{2.417150in}{0.660000in}}%
\pgfpathlineto{\pgfqpoint{2.423350in}{0.660000in}}%
\pgfpathlineto{\pgfqpoint{2.423350in}{3.407051in}}%
\pgfpathlineto{\pgfqpoint{2.417150in}{3.407051in}}%
\pgfpathlineto{\pgfqpoint{2.417150in}{0.660000in}}%
\pgfpathclose%
\pgfusepath{fill}%
\end{pgfscope}%
\begin{pgfscope}%
\pgfpathrectangle{\pgfqpoint{1.250000in}{0.660000in}}{\pgfqpoint{7.750000in}{4.620000in}}%
\pgfusepath{clip}%
\pgfsetbuttcap%
\pgfsetmiterjoin%
\definecolor{currentfill}{rgb}{0.000000,0.000000,0.000000}%
\pgfsetfillcolor{currentfill}%
\pgfsetlinewidth{0.000000pt}%
\definecolor{currentstroke}{rgb}{0.000000,0.000000,0.000000}%
\pgfsetstrokecolor{currentstroke}%
\pgfsetstrokeopacity{0.000000}%
\pgfsetdash{}{0pt}%
\pgfpathmoveto{\pgfqpoint{2.424900in}{0.660000in}}%
\pgfpathlineto{\pgfqpoint{2.431100in}{0.660000in}}%
\pgfpathlineto{\pgfqpoint{2.431100in}{3.405995in}}%
\pgfpathlineto{\pgfqpoint{2.424900in}{3.405995in}}%
\pgfpathlineto{\pgfqpoint{2.424900in}{0.660000in}}%
\pgfpathclose%
\pgfusepath{fill}%
\end{pgfscope}%
\begin{pgfscope}%
\pgfpathrectangle{\pgfqpoint{1.250000in}{0.660000in}}{\pgfqpoint{7.750000in}{4.620000in}}%
\pgfusepath{clip}%
\pgfsetbuttcap%
\pgfsetmiterjoin%
\definecolor{currentfill}{rgb}{0.000000,0.000000,0.000000}%
\pgfsetfillcolor{currentfill}%
\pgfsetlinewidth{0.000000pt}%
\definecolor{currentstroke}{rgb}{0.000000,0.000000,0.000000}%
\pgfsetstrokecolor{currentstroke}%
\pgfsetstrokeopacity{0.000000}%
\pgfsetdash{}{0pt}%
\pgfpathmoveto{\pgfqpoint{2.432650in}{0.660000in}}%
\pgfpathlineto{\pgfqpoint{2.438850in}{0.660000in}}%
\pgfpathlineto{\pgfqpoint{2.438850in}{3.404949in}}%
\pgfpathlineto{\pgfqpoint{2.432650in}{3.404949in}}%
\pgfpathlineto{\pgfqpoint{2.432650in}{0.660000in}}%
\pgfpathclose%
\pgfusepath{fill}%
\end{pgfscope}%
\begin{pgfscope}%
\pgfpathrectangle{\pgfqpoint{1.250000in}{0.660000in}}{\pgfqpoint{7.750000in}{4.620000in}}%
\pgfusepath{clip}%
\pgfsetbuttcap%
\pgfsetmiterjoin%
\definecolor{currentfill}{rgb}{0.000000,0.000000,0.000000}%
\pgfsetfillcolor{currentfill}%
\pgfsetlinewidth{0.000000pt}%
\definecolor{currentstroke}{rgb}{0.000000,0.000000,0.000000}%
\pgfsetstrokecolor{currentstroke}%
\pgfsetstrokeopacity{0.000000}%
\pgfsetdash{}{0pt}%
\pgfpathmoveto{\pgfqpoint{2.440400in}{0.660000in}}%
\pgfpathlineto{\pgfqpoint{2.446600in}{0.660000in}}%
\pgfpathlineto{\pgfqpoint{2.446600in}{3.403911in}}%
\pgfpathlineto{\pgfqpoint{2.440400in}{3.403911in}}%
\pgfpathlineto{\pgfqpoint{2.440400in}{0.660000in}}%
\pgfpathclose%
\pgfusepath{fill}%
\end{pgfscope}%
\begin{pgfscope}%
\pgfpathrectangle{\pgfqpoint{1.250000in}{0.660000in}}{\pgfqpoint{7.750000in}{4.620000in}}%
\pgfusepath{clip}%
\pgfsetbuttcap%
\pgfsetmiterjoin%
\definecolor{currentfill}{rgb}{0.000000,0.000000,0.000000}%
\pgfsetfillcolor{currentfill}%
\pgfsetlinewidth{0.000000pt}%
\definecolor{currentstroke}{rgb}{0.000000,0.000000,0.000000}%
\pgfsetstrokecolor{currentstroke}%
\pgfsetstrokeopacity{0.000000}%
\pgfsetdash{}{0pt}%
\pgfpathmoveto{\pgfqpoint{2.448150in}{0.660000in}}%
\pgfpathlineto{\pgfqpoint{2.454350in}{0.660000in}}%
\pgfpathlineto{\pgfqpoint{2.454350in}{3.402885in}}%
\pgfpathlineto{\pgfqpoint{2.448150in}{3.402885in}}%
\pgfpathlineto{\pgfqpoint{2.448150in}{0.660000in}}%
\pgfpathclose%
\pgfusepath{fill}%
\end{pgfscope}%
\begin{pgfscope}%
\pgfpathrectangle{\pgfqpoint{1.250000in}{0.660000in}}{\pgfqpoint{7.750000in}{4.620000in}}%
\pgfusepath{clip}%
\pgfsetbuttcap%
\pgfsetmiterjoin%
\definecolor{currentfill}{rgb}{0.000000,0.000000,0.000000}%
\pgfsetfillcolor{currentfill}%
\pgfsetlinewidth{0.000000pt}%
\definecolor{currentstroke}{rgb}{0.000000,0.000000,0.000000}%
\pgfsetstrokecolor{currentstroke}%
\pgfsetstrokeopacity{0.000000}%
\pgfsetdash{}{0pt}%
\pgfpathmoveto{\pgfqpoint{2.455900in}{0.660000in}}%
\pgfpathlineto{\pgfqpoint{2.462100in}{0.660000in}}%
\pgfpathlineto{\pgfqpoint{2.462100in}{3.401867in}}%
\pgfpathlineto{\pgfqpoint{2.455900in}{3.401867in}}%
\pgfpathlineto{\pgfqpoint{2.455900in}{0.660000in}}%
\pgfpathclose%
\pgfusepath{fill}%
\end{pgfscope}%
\begin{pgfscope}%
\pgfpathrectangle{\pgfqpoint{1.250000in}{0.660000in}}{\pgfqpoint{7.750000in}{4.620000in}}%
\pgfusepath{clip}%
\pgfsetbuttcap%
\pgfsetmiterjoin%
\definecolor{currentfill}{rgb}{0.000000,0.000000,0.000000}%
\pgfsetfillcolor{currentfill}%
\pgfsetlinewidth{0.000000pt}%
\definecolor{currentstroke}{rgb}{0.000000,0.000000,0.000000}%
\pgfsetstrokecolor{currentstroke}%
\pgfsetstrokeopacity{0.000000}%
\pgfsetdash{}{0pt}%
\pgfpathmoveto{\pgfqpoint{2.463650in}{0.660000in}}%
\pgfpathlineto{\pgfqpoint{2.469850in}{0.660000in}}%
\pgfpathlineto{\pgfqpoint{2.469850in}{3.400858in}}%
\pgfpathlineto{\pgfqpoint{2.463650in}{3.400858in}}%
\pgfpathlineto{\pgfqpoint{2.463650in}{0.660000in}}%
\pgfpathclose%
\pgfusepath{fill}%
\end{pgfscope}%
\begin{pgfscope}%
\pgfpathrectangle{\pgfqpoint{1.250000in}{0.660000in}}{\pgfqpoint{7.750000in}{4.620000in}}%
\pgfusepath{clip}%
\pgfsetbuttcap%
\pgfsetmiterjoin%
\definecolor{currentfill}{rgb}{0.000000,0.000000,0.000000}%
\pgfsetfillcolor{currentfill}%
\pgfsetlinewidth{0.000000pt}%
\definecolor{currentstroke}{rgb}{0.000000,0.000000,0.000000}%
\pgfsetstrokecolor{currentstroke}%
\pgfsetstrokeopacity{0.000000}%
\pgfsetdash{}{0pt}%
\pgfpathmoveto{\pgfqpoint{2.471400in}{0.660000in}}%
\pgfpathlineto{\pgfqpoint{2.477600in}{0.660000in}}%
\pgfpathlineto{\pgfqpoint{2.477600in}{3.399857in}}%
\pgfpathlineto{\pgfqpoint{2.471400in}{3.399857in}}%
\pgfpathlineto{\pgfqpoint{2.471400in}{0.660000in}}%
\pgfpathclose%
\pgfusepath{fill}%
\end{pgfscope}%
\begin{pgfscope}%
\pgfpathrectangle{\pgfqpoint{1.250000in}{0.660000in}}{\pgfqpoint{7.750000in}{4.620000in}}%
\pgfusepath{clip}%
\pgfsetbuttcap%
\pgfsetmiterjoin%
\definecolor{currentfill}{rgb}{0.000000,0.000000,0.000000}%
\pgfsetfillcolor{currentfill}%
\pgfsetlinewidth{0.000000pt}%
\definecolor{currentstroke}{rgb}{0.000000,0.000000,0.000000}%
\pgfsetstrokecolor{currentstroke}%
\pgfsetstrokeopacity{0.000000}%
\pgfsetdash{}{0pt}%
\pgfpathmoveto{\pgfqpoint{2.479150in}{0.660000in}}%
\pgfpathlineto{\pgfqpoint{2.485350in}{0.660000in}}%
\pgfpathlineto{\pgfqpoint{2.485350in}{3.398864in}}%
\pgfpathlineto{\pgfqpoint{2.479150in}{3.398864in}}%
\pgfpathlineto{\pgfqpoint{2.479150in}{0.660000in}}%
\pgfpathclose%
\pgfusepath{fill}%
\end{pgfscope}%
\begin{pgfscope}%
\pgfpathrectangle{\pgfqpoint{1.250000in}{0.660000in}}{\pgfqpoint{7.750000in}{4.620000in}}%
\pgfusepath{clip}%
\pgfsetbuttcap%
\pgfsetmiterjoin%
\definecolor{currentfill}{rgb}{0.000000,0.000000,0.000000}%
\pgfsetfillcolor{currentfill}%
\pgfsetlinewidth{0.000000pt}%
\definecolor{currentstroke}{rgb}{0.000000,0.000000,0.000000}%
\pgfsetstrokecolor{currentstroke}%
\pgfsetstrokeopacity{0.000000}%
\pgfsetdash{}{0pt}%
\pgfpathmoveto{\pgfqpoint{2.486900in}{0.660000in}}%
\pgfpathlineto{\pgfqpoint{2.493100in}{0.660000in}}%
\pgfpathlineto{\pgfqpoint{2.493100in}{3.397882in}}%
\pgfpathlineto{\pgfqpoint{2.486900in}{3.397882in}}%
\pgfpathlineto{\pgfqpoint{2.486900in}{0.660000in}}%
\pgfpathclose%
\pgfusepath{fill}%
\end{pgfscope}%
\begin{pgfscope}%
\pgfpathrectangle{\pgfqpoint{1.250000in}{0.660000in}}{\pgfqpoint{7.750000in}{4.620000in}}%
\pgfusepath{clip}%
\pgfsetbuttcap%
\pgfsetmiterjoin%
\definecolor{currentfill}{rgb}{0.000000,0.000000,0.000000}%
\pgfsetfillcolor{currentfill}%
\pgfsetlinewidth{0.000000pt}%
\definecolor{currentstroke}{rgb}{0.000000,0.000000,0.000000}%
\pgfsetstrokecolor{currentstroke}%
\pgfsetstrokeopacity{0.000000}%
\pgfsetdash{}{0pt}%
\pgfpathmoveto{\pgfqpoint{2.494650in}{0.660000in}}%
\pgfpathlineto{\pgfqpoint{2.500850in}{0.660000in}}%
\pgfpathlineto{\pgfqpoint{2.500850in}{3.396908in}}%
\pgfpathlineto{\pgfqpoint{2.494650in}{3.396908in}}%
\pgfpathlineto{\pgfqpoint{2.494650in}{0.660000in}}%
\pgfpathclose%
\pgfusepath{fill}%
\end{pgfscope}%
\begin{pgfscope}%
\pgfpathrectangle{\pgfqpoint{1.250000in}{0.660000in}}{\pgfqpoint{7.750000in}{4.620000in}}%
\pgfusepath{clip}%
\pgfsetbuttcap%
\pgfsetmiterjoin%
\definecolor{currentfill}{rgb}{0.000000,0.000000,0.000000}%
\pgfsetfillcolor{currentfill}%
\pgfsetlinewidth{0.000000pt}%
\definecolor{currentstroke}{rgb}{0.000000,0.000000,0.000000}%
\pgfsetstrokecolor{currentstroke}%
\pgfsetstrokeopacity{0.000000}%
\pgfsetdash{}{0pt}%
\pgfpathmoveto{\pgfqpoint{2.502400in}{0.660000in}}%
\pgfpathlineto{\pgfqpoint{2.508600in}{0.660000in}}%
\pgfpathlineto{\pgfqpoint{2.508600in}{3.395942in}}%
\pgfpathlineto{\pgfqpoint{2.502400in}{3.395942in}}%
\pgfpathlineto{\pgfqpoint{2.502400in}{0.660000in}}%
\pgfpathclose%
\pgfusepath{fill}%
\end{pgfscope}%
\begin{pgfscope}%
\pgfpathrectangle{\pgfqpoint{1.250000in}{0.660000in}}{\pgfqpoint{7.750000in}{4.620000in}}%
\pgfusepath{clip}%
\pgfsetbuttcap%
\pgfsetmiterjoin%
\definecolor{currentfill}{rgb}{0.000000,0.000000,0.000000}%
\pgfsetfillcolor{currentfill}%
\pgfsetlinewidth{0.000000pt}%
\definecolor{currentstroke}{rgb}{0.000000,0.000000,0.000000}%
\pgfsetstrokecolor{currentstroke}%
\pgfsetstrokeopacity{0.000000}%
\pgfsetdash{}{0pt}%
\pgfpathmoveto{\pgfqpoint{2.510150in}{0.660000in}}%
\pgfpathlineto{\pgfqpoint{2.516350in}{0.660000in}}%
\pgfpathlineto{\pgfqpoint{2.516350in}{3.394983in}}%
\pgfpathlineto{\pgfqpoint{2.510150in}{3.394983in}}%
\pgfpathlineto{\pgfqpoint{2.510150in}{0.660000in}}%
\pgfpathclose%
\pgfusepath{fill}%
\end{pgfscope}%
\begin{pgfscope}%
\pgfpathrectangle{\pgfqpoint{1.250000in}{0.660000in}}{\pgfqpoint{7.750000in}{4.620000in}}%
\pgfusepath{clip}%
\pgfsetbuttcap%
\pgfsetmiterjoin%
\definecolor{currentfill}{rgb}{0.000000,0.000000,0.000000}%
\pgfsetfillcolor{currentfill}%
\pgfsetlinewidth{0.000000pt}%
\definecolor{currentstroke}{rgb}{0.000000,0.000000,0.000000}%
\pgfsetstrokecolor{currentstroke}%
\pgfsetstrokeopacity{0.000000}%
\pgfsetdash{}{0pt}%
\pgfpathmoveto{\pgfqpoint{2.517900in}{0.660000in}}%
\pgfpathlineto{\pgfqpoint{2.524100in}{0.660000in}}%
\pgfpathlineto{\pgfqpoint{2.524100in}{3.394034in}}%
\pgfpathlineto{\pgfqpoint{2.517900in}{3.394034in}}%
\pgfpathlineto{\pgfqpoint{2.517900in}{0.660000in}}%
\pgfpathclose%
\pgfusepath{fill}%
\end{pgfscope}%
\begin{pgfscope}%
\pgfpathrectangle{\pgfqpoint{1.250000in}{0.660000in}}{\pgfqpoint{7.750000in}{4.620000in}}%
\pgfusepath{clip}%
\pgfsetbuttcap%
\pgfsetmiterjoin%
\definecolor{currentfill}{rgb}{0.000000,0.000000,0.000000}%
\pgfsetfillcolor{currentfill}%
\pgfsetlinewidth{0.000000pt}%
\definecolor{currentstroke}{rgb}{0.000000,0.000000,0.000000}%
\pgfsetstrokecolor{currentstroke}%
\pgfsetstrokeopacity{0.000000}%
\pgfsetdash{}{0pt}%
\pgfpathmoveto{\pgfqpoint{2.525650in}{0.660000in}}%
\pgfpathlineto{\pgfqpoint{2.531850in}{0.660000in}}%
\pgfpathlineto{\pgfqpoint{2.531850in}{3.393093in}}%
\pgfpathlineto{\pgfqpoint{2.525650in}{3.393093in}}%
\pgfpathlineto{\pgfqpoint{2.525650in}{0.660000in}}%
\pgfpathclose%
\pgfusepath{fill}%
\end{pgfscope}%
\begin{pgfscope}%
\pgfpathrectangle{\pgfqpoint{1.250000in}{0.660000in}}{\pgfqpoint{7.750000in}{4.620000in}}%
\pgfusepath{clip}%
\pgfsetbuttcap%
\pgfsetmiterjoin%
\definecolor{currentfill}{rgb}{0.000000,0.000000,0.000000}%
\pgfsetfillcolor{currentfill}%
\pgfsetlinewidth{0.000000pt}%
\definecolor{currentstroke}{rgb}{0.000000,0.000000,0.000000}%
\pgfsetstrokecolor{currentstroke}%
\pgfsetstrokeopacity{0.000000}%
\pgfsetdash{}{0pt}%
\pgfpathmoveto{\pgfqpoint{2.533400in}{0.660000in}}%
\pgfpathlineto{\pgfqpoint{2.539600in}{0.660000in}}%
\pgfpathlineto{\pgfqpoint{2.539600in}{3.392158in}}%
\pgfpathlineto{\pgfqpoint{2.533400in}{3.392158in}}%
\pgfpathlineto{\pgfqpoint{2.533400in}{0.660000in}}%
\pgfpathclose%
\pgfusepath{fill}%
\end{pgfscope}%
\begin{pgfscope}%
\pgfpathrectangle{\pgfqpoint{1.250000in}{0.660000in}}{\pgfqpoint{7.750000in}{4.620000in}}%
\pgfusepath{clip}%
\pgfsetbuttcap%
\pgfsetmiterjoin%
\definecolor{currentfill}{rgb}{0.000000,0.000000,0.000000}%
\pgfsetfillcolor{currentfill}%
\pgfsetlinewidth{0.000000pt}%
\definecolor{currentstroke}{rgb}{0.000000,0.000000,0.000000}%
\pgfsetstrokecolor{currentstroke}%
\pgfsetstrokeopacity{0.000000}%
\pgfsetdash{}{0pt}%
\pgfpathmoveto{\pgfqpoint{2.541150in}{0.660000in}}%
\pgfpathlineto{\pgfqpoint{2.547350in}{0.660000in}}%
\pgfpathlineto{\pgfqpoint{2.547350in}{3.391232in}}%
\pgfpathlineto{\pgfqpoint{2.541150in}{3.391232in}}%
\pgfpathlineto{\pgfqpoint{2.541150in}{0.660000in}}%
\pgfpathclose%
\pgfusepath{fill}%
\end{pgfscope}%
\begin{pgfscope}%
\pgfpathrectangle{\pgfqpoint{1.250000in}{0.660000in}}{\pgfqpoint{7.750000in}{4.620000in}}%
\pgfusepath{clip}%
\pgfsetbuttcap%
\pgfsetmiterjoin%
\definecolor{currentfill}{rgb}{0.000000,0.000000,0.000000}%
\pgfsetfillcolor{currentfill}%
\pgfsetlinewidth{0.000000pt}%
\definecolor{currentstroke}{rgb}{0.000000,0.000000,0.000000}%
\pgfsetstrokecolor{currentstroke}%
\pgfsetstrokeopacity{0.000000}%
\pgfsetdash{}{0pt}%
\pgfpathmoveto{\pgfqpoint{2.548900in}{0.660000in}}%
\pgfpathlineto{\pgfqpoint{2.555100in}{0.660000in}}%
\pgfpathlineto{\pgfqpoint{2.555100in}{3.390314in}}%
\pgfpathlineto{\pgfqpoint{2.548900in}{3.390314in}}%
\pgfpathlineto{\pgfqpoint{2.548900in}{0.660000in}}%
\pgfpathclose%
\pgfusepath{fill}%
\end{pgfscope}%
\begin{pgfscope}%
\pgfpathrectangle{\pgfqpoint{1.250000in}{0.660000in}}{\pgfqpoint{7.750000in}{4.620000in}}%
\pgfusepath{clip}%
\pgfsetbuttcap%
\pgfsetmiterjoin%
\definecolor{currentfill}{rgb}{0.000000,0.000000,0.000000}%
\pgfsetfillcolor{currentfill}%
\pgfsetlinewidth{0.000000pt}%
\definecolor{currentstroke}{rgb}{0.000000,0.000000,0.000000}%
\pgfsetstrokecolor{currentstroke}%
\pgfsetstrokeopacity{0.000000}%
\pgfsetdash{}{0pt}%
\pgfpathmoveto{\pgfqpoint{2.556650in}{0.660000in}}%
\pgfpathlineto{\pgfqpoint{2.562850in}{0.660000in}}%
\pgfpathlineto{\pgfqpoint{2.562850in}{3.389403in}}%
\pgfpathlineto{\pgfqpoint{2.556650in}{3.389403in}}%
\pgfpathlineto{\pgfqpoint{2.556650in}{0.660000in}}%
\pgfpathclose%
\pgfusepath{fill}%
\end{pgfscope}%
\begin{pgfscope}%
\pgfpathrectangle{\pgfqpoint{1.250000in}{0.660000in}}{\pgfqpoint{7.750000in}{4.620000in}}%
\pgfusepath{clip}%
\pgfsetbuttcap%
\pgfsetmiterjoin%
\definecolor{currentfill}{rgb}{0.000000,0.000000,0.000000}%
\pgfsetfillcolor{currentfill}%
\pgfsetlinewidth{0.000000pt}%
\definecolor{currentstroke}{rgb}{0.000000,0.000000,0.000000}%
\pgfsetstrokecolor{currentstroke}%
\pgfsetstrokeopacity{0.000000}%
\pgfsetdash{}{0pt}%
\pgfpathmoveto{\pgfqpoint{2.564400in}{0.660000in}}%
\pgfpathlineto{\pgfqpoint{2.570600in}{0.660000in}}%
\pgfpathlineto{\pgfqpoint{2.570600in}{3.388498in}}%
\pgfpathlineto{\pgfqpoint{2.564400in}{3.388498in}}%
\pgfpathlineto{\pgfqpoint{2.564400in}{0.660000in}}%
\pgfpathclose%
\pgfusepath{fill}%
\end{pgfscope}%
\begin{pgfscope}%
\pgfpathrectangle{\pgfqpoint{1.250000in}{0.660000in}}{\pgfqpoint{7.750000in}{4.620000in}}%
\pgfusepath{clip}%
\pgfsetbuttcap%
\pgfsetmiterjoin%
\definecolor{currentfill}{rgb}{0.000000,0.000000,0.000000}%
\pgfsetfillcolor{currentfill}%
\pgfsetlinewidth{0.000000pt}%
\definecolor{currentstroke}{rgb}{0.000000,0.000000,0.000000}%
\pgfsetstrokecolor{currentstroke}%
\pgfsetstrokeopacity{0.000000}%
\pgfsetdash{}{0pt}%
\pgfpathmoveto{\pgfqpoint{2.572150in}{0.660000in}}%
\pgfpathlineto{\pgfqpoint{2.578350in}{0.660000in}}%
\pgfpathlineto{\pgfqpoint{2.578350in}{3.387603in}}%
\pgfpathlineto{\pgfqpoint{2.572150in}{3.387603in}}%
\pgfpathlineto{\pgfqpoint{2.572150in}{0.660000in}}%
\pgfpathclose%
\pgfusepath{fill}%
\end{pgfscope}%
\begin{pgfscope}%
\pgfpathrectangle{\pgfqpoint{1.250000in}{0.660000in}}{\pgfqpoint{7.750000in}{4.620000in}}%
\pgfusepath{clip}%
\pgfsetbuttcap%
\pgfsetmiterjoin%
\definecolor{currentfill}{rgb}{0.000000,0.000000,0.000000}%
\pgfsetfillcolor{currentfill}%
\pgfsetlinewidth{0.000000pt}%
\definecolor{currentstroke}{rgb}{0.000000,0.000000,0.000000}%
\pgfsetstrokecolor{currentstroke}%
\pgfsetstrokeopacity{0.000000}%
\pgfsetdash{}{0pt}%
\pgfpathmoveto{\pgfqpoint{2.579900in}{0.660000in}}%
\pgfpathlineto{\pgfqpoint{2.586100in}{0.660000in}}%
\pgfpathlineto{\pgfqpoint{2.586100in}{3.386715in}}%
\pgfpathlineto{\pgfqpoint{2.579900in}{3.386715in}}%
\pgfpathlineto{\pgfqpoint{2.579900in}{0.660000in}}%
\pgfpathclose%
\pgfusepath{fill}%
\end{pgfscope}%
\begin{pgfscope}%
\pgfpathrectangle{\pgfqpoint{1.250000in}{0.660000in}}{\pgfqpoint{7.750000in}{4.620000in}}%
\pgfusepath{clip}%
\pgfsetbuttcap%
\pgfsetmiterjoin%
\definecolor{currentfill}{rgb}{0.000000,0.000000,0.000000}%
\pgfsetfillcolor{currentfill}%
\pgfsetlinewidth{0.000000pt}%
\definecolor{currentstroke}{rgb}{0.000000,0.000000,0.000000}%
\pgfsetstrokecolor{currentstroke}%
\pgfsetstrokeopacity{0.000000}%
\pgfsetdash{}{0pt}%
\pgfpathmoveto{\pgfqpoint{2.587650in}{0.660000in}}%
\pgfpathlineto{\pgfqpoint{2.593850in}{0.660000in}}%
\pgfpathlineto{\pgfqpoint{2.593850in}{3.385832in}}%
\pgfpathlineto{\pgfqpoint{2.587650in}{3.385832in}}%
\pgfpathlineto{\pgfqpoint{2.587650in}{0.660000in}}%
\pgfpathclose%
\pgfusepath{fill}%
\end{pgfscope}%
\begin{pgfscope}%
\pgfpathrectangle{\pgfqpoint{1.250000in}{0.660000in}}{\pgfqpoint{7.750000in}{4.620000in}}%
\pgfusepath{clip}%
\pgfsetbuttcap%
\pgfsetmiterjoin%
\definecolor{currentfill}{rgb}{0.000000,0.000000,0.000000}%
\pgfsetfillcolor{currentfill}%
\pgfsetlinewidth{0.000000pt}%
\definecolor{currentstroke}{rgb}{0.000000,0.000000,0.000000}%
\pgfsetstrokecolor{currentstroke}%
\pgfsetstrokeopacity{0.000000}%
\pgfsetdash{}{0pt}%
\pgfpathmoveto{\pgfqpoint{2.595400in}{0.660000in}}%
\pgfpathlineto{\pgfqpoint{2.601600in}{0.660000in}}%
\pgfpathlineto{\pgfqpoint{2.601600in}{3.384958in}}%
\pgfpathlineto{\pgfqpoint{2.595400in}{3.384958in}}%
\pgfpathlineto{\pgfqpoint{2.595400in}{0.660000in}}%
\pgfpathclose%
\pgfusepath{fill}%
\end{pgfscope}%
\begin{pgfscope}%
\pgfpathrectangle{\pgfqpoint{1.250000in}{0.660000in}}{\pgfqpoint{7.750000in}{4.620000in}}%
\pgfusepath{clip}%
\pgfsetbuttcap%
\pgfsetmiterjoin%
\definecolor{currentfill}{rgb}{0.000000,0.000000,0.000000}%
\pgfsetfillcolor{currentfill}%
\pgfsetlinewidth{0.000000pt}%
\definecolor{currentstroke}{rgb}{0.000000,0.000000,0.000000}%
\pgfsetstrokecolor{currentstroke}%
\pgfsetstrokeopacity{0.000000}%
\pgfsetdash{}{0pt}%
\pgfpathmoveto{\pgfqpoint{2.603150in}{0.660000in}}%
\pgfpathlineto{\pgfqpoint{2.609350in}{0.660000in}}%
\pgfpathlineto{\pgfqpoint{2.609350in}{3.384091in}}%
\pgfpathlineto{\pgfqpoint{2.603150in}{3.384091in}}%
\pgfpathlineto{\pgfqpoint{2.603150in}{0.660000in}}%
\pgfpathclose%
\pgfusepath{fill}%
\end{pgfscope}%
\begin{pgfscope}%
\pgfpathrectangle{\pgfqpoint{1.250000in}{0.660000in}}{\pgfqpoint{7.750000in}{4.620000in}}%
\pgfusepath{clip}%
\pgfsetbuttcap%
\pgfsetmiterjoin%
\definecolor{currentfill}{rgb}{0.000000,0.000000,0.000000}%
\pgfsetfillcolor{currentfill}%
\pgfsetlinewidth{0.000000pt}%
\definecolor{currentstroke}{rgb}{0.000000,0.000000,0.000000}%
\pgfsetstrokecolor{currentstroke}%
\pgfsetstrokeopacity{0.000000}%
\pgfsetdash{}{0pt}%
\pgfpathmoveto{\pgfqpoint{2.610900in}{0.660000in}}%
\pgfpathlineto{\pgfqpoint{2.617100in}{0.660000in}}%
\pgfpathlineto{\pgfqpoint{2.617100in}{3.383228in}}%
\pgfpathlineto{\pgfqpoint{2.610900in}{3.383228in}}%
\pgfpathlineto{\pgfqpoint{2.610900in}{0.660000in}}%
\pgfpathclose%
\pgfusepath{fill}%
\end{pgfscope}%
\begin{pgfscope}%
\pgfpathrectangle{\pgfqpoint{1.250000in}{0.660000in}}{\pgfqpoint{7.750000in}{4.620000in}}%
\pgfusepath{clip}%
\pgfsetbuttcap%
\pgfsetmiterjoin%
\definecolor{currentfill}{rgb}{0.000000,0.000000,0.000000}%
\pgfsetfillcolor{currentfill}%
\pgfsetlinewidth{0.000000pt}%
\definecolor{currentstroke}{rgb}{0.000000,0.000000,0.000000}%
\pgfsetstrokecolor{currentstroke}%
\pgfsetstrokeopacity{0.000000}%
\pgfsetdash{}{0pt}%
\pgfpathmoveto{\pgfqpoint{2.618650in}{0.660000in}}%
\pgfpathlineto{\pgfqpoint{2.624850in}{0.660000in}}%
\pgfpathlineto{\pgfqpoint{2.624850in}{3.382375in}}%
\pgfpathlineto{\pgfqpoint{2.618650in}{3.382375in}}%
\pgfpathlineto{\pgfqpoint{2.618650in}{0.660000in}}%
\pgfpathclose%
\pgfusepath{fill}%
\end{pgfscope}%
\begin{pgfscope}%
\pgfpathrectangle{\pgfqpoint{1.250000in}{0.660000in}}{\pgfqpoint{7.750000in}{4.620000in}}%
\pgfusepath{clip}%
\pgfsetbuttcap%
\pgfsetmiterjoin%
\definecolor{currentfill}{rgb}{0.000000,0.000000,0.000000}%
\pgfsetfillcolor{currentfill}%
\pgfsetlinewidth{0.000000pt}%
\definecolor{currentstroke}{rgb}{0.000000,0.000000,0.000000}%
\pgfsetstrokecolor{currentstroke}%
\pgfsetstrokeopacity{0.000000}%
\pgfsetdash{}{0pt}%
\pgfpathmoveto{\pgfqpoint{2.626400in}{0.660000in}}%
\pgfpathlineto{\pgfqpoint{2.632600in}{0.660000in}}%
\pgfpathlineto{\pgfqpoint{2.632600in}{3.381528in}}%
\pgfpathlineto{\pgfqpoint{2.626400in}{3.381528in}}%
\pgfpathlineto{\pgfqpoint{2.626400in}{0.660000in}}%
\pgfpathclose%
\pgfusepath{fill}%
\end{pgfscope}%
\begin{pgfscope}%
\pgfpathrectangle{\pgfqpoint{1.250000in}{0.660000in}}{\pgfqpoint{7.750000in}{4.620000in}}%
\pgfusepath{clip}%
\pgfsetbuttcap%
\pgfsetmiterjoin%
\definecolor{currentfill}{rgb}{0.000000,0.000000,0.000000}%
\pgfsetfillcolor{currentfill}%
\pgfsetlinewidth{0.000000pt}%
\definecolor{currentstroke}{rgb}{0.000000,0.000000,0.000000}%
\pgfsetstrokecolor{currentstroke}%
\pgfsetstrokeopacity{0.000000}%
\pgfsetdash{}{0pt}%
\pgfpathmoveto{\pgfqpoint{2.634150in}{0.660000in}}%
\pgfpathlineto{\pgfqpoint{2.640350in}{0.660000in}}%
\pgfpathlineto{\pgfqpoint{2.640350in}{3.380686in}}%
\pgfpathlineto{\pgfqpoint{2.634150in}{3.380686in}}%
\pgfpathlineto{\pgfqpoint{2.634150in}{0.660000in}}%
\pgfpathclose%
\pgfusepath{fill}%
\end{pgfscope}%
\begin{pgfscope}%
\pgfpathrectangle{\pgfqpoint{1.250000in}{0.660000in}}{\pgfqpoint{7.750000in}{4.620000in}}%
\pgfusepath{clip}%
\pgfsetbuttcap%
\pgfsetmiterjoin%
\definecolor{currentfill}{rgb}{0.000000,0.000000,0.000000}%
\pgfsetfillcolor{currentfill}%
\pgfsetlinewidth{0.000000pt}%
\definecolor{currentstroke}{rgb}{0.000000,0.000000,0.000000}%
\pgfsetstrokecolor{currentstroke}%
\pgfsetstrokeopacity{0.000000}%
\pgfsetdash{}{0pt}%
\pgfpathmoveto{\pgfqpoint{2.641900in}{0.660000in}}%
\pgfpathlineto{\pgfqpoint{2.648100in}{0.660000in}}%
\pgfpathlineto{\pgfqpoint{2.648100in}{3.379853in}}%
\pgfpathlineto{\pgfqpoint{2.641900in}{3.379853in}}%
\pgfpathlineto{\pgfqpoint{2.641900in}{0.660000in}}%
\pgfpathclose%
\pgfusepath{fill}%
\end{pgfscope}%
\begin{pgfscope}%
\pgfpathrectangle{\pgfqpoint{1.250000in}{0.660000in}}{\pgfqpoint{7.750000in}{4.620000in}}%
\pgfusepath{clip}%
\pgfsetbuttcap%
\pgfsetmiterjoin%
\definecolor{currentfill}{rgb}{0.000000,0.000000,0.000000}%
\pgfsetfillcolor{currentfill}%
\pgfsetlinewidth{0.000000pt}%
\definecolor{currentstroke}{rgb}{0.000000,0.000000,0.000000}%
\pgfsetstrokecolor{currentstroke}%
\pgfsetstrokeopacity{0.000000}%
\pgfsetdash{}{0pt}%
\pgfpathmoveto{\pgfqpoint{2.649650in}{0.660000in}}%
\pgfpathlineto{\pgfqpoint{2.655850in}{0.660000in}}%
\pgfpathlineto{\pgfqpoint{2.655850in}{3.379024in}}%
\pgfpathlineto{\pgfqpoint{2.649650in}{3.379024in}}%
\pgfpathlineto{\pgfqpoint{2.649650in}{0.660000in}}%
\pgfpathclose%
\pgfusepath{fill}%
\end{pgfscope}%
\begin{pgfscope}%
\pgfpathrectangle{\pgfqpoint{1.250000in}{0.660000in}}{\pgfqpoint{7.750000in}{4.620000in}}%
\pgfusepath{clip}%
\pgfsetbuttcap%
\pgfsetmiterjoin%
\definecolor{currentfill}{rgb}{0.000000,0.000000,0.000000}%
\pgfsetfillcolor{currentfill}%
\pgfsetlinewidth{0.000000pt}%
\definecolor{currentstroke}{rgb}{0.000000,0.000000,0.000000}%
\pgfsetstrokecolor{currentstroke}%
\pgfsetstrokeopacity{0.000000}%
\pgfsetdash{}{0pt}%
\pgfpathmoveto{\pgfqpoint{2.657400in}{0.660000in}}%
\pgfpathlineto{\pgfqpoint{2.663600in}{0.660000in}}%
\pgfpathlineto{\pgfqpoint{2.663600in}{3.378202in}}%
\pgfpathlineto{\pgfqpoint{2.657400in}{3.378202in}}%
\pgfpathlineto{\pgfqpoint{2.657400in}{0.660000in}}%
\pgfpathclose%
\pgfusepath{fill}%
\end{pgfscope}%
\begin{pgfscope}%
\pgfpathrectangle{\pgfqpoint{1.250000in}{0.660000in}}{\pgfqpoint{7.750000in}{4.620000in}}%
\pgfusepath{clip}%
\pgfsetbuttcap%
\pgfsetmiterjoin%
\definecolor{currentfill}{rgb}{0.000000,0.000000,0.000000}%
\pgfsetfillcolor{currentfill}%
\pgfsetlinewidth{0.000000pt}%
\definecolor{currentstroke}{rgb}{0.000000,0.000000,0.000000}%
\pgfsetstrokecolor{currentstroke}%
\pgfsetstrokeopacity{0.000000}%
\pgfsetdash{}{0pt}%
\pgfpathmoveto{\pgfqpoint{2.665150in}{0.660000in}}%
\pgfpathlineto{\pgfqpoint{2.671350in}{0.660000in}}%
\pgfpathlineto{\pgfqpoint{2.671350in}{3.377388in}}%
\pgfpathlineto{\pgfqpoint{2.665150in}{3.377388in}}%
\pgfpathlineto{\pgfqpoint{2.665150in}{0.660000in}}%
\pgfpathclose%
\pgfusepath{fill}%
\end{pgfscope}%
\begin{pgfscope}%
\pgfpathrectangle{\pgfqpoint{1.250000in}{0.660000in}}{\pgfqpoint{7.750000in}{4.620000in}}%
\pgfusepath{clip}%
\pgfsetbuttcap%
\pgfsetmiterjoin%
\definecolor{currentfill}{rgb}{0.000000,0.000000,0.000000}%
\pgfsetfillcolor{currentfill}%
\pgfsetlinewidth{0.000000pt}%
\definecolor{currentstroke}{rgb}{0.000000,0.000000,0.000000}%
\pgfsetstrokecolor{currentstroke}%
\pgfsetstrokeopacity{0.000000}%
\pgfsetdash{}{0pt}%
\pgfpathmoveto{\pgfqpoint{2.672900in}{0.660000in}}%
\pgfpathlineto{\pgfqpoint{2.679100in}{0.660000in}}%
\pgfpathlineto{\pgfqpoint{2.679100in}{3.376577in}}%
\pgfpathlineto{\pgfqpoint{2.672900in}{3.376577in}}%
\pgfpathlineto{\pgfqpoint{2.672900in}{0.660000in}}%
\pgfpathclose%
\pgfusepath{fill}%
\end{pgfscope}%
\begin{pgfscope}%
\pgfpathrectangle{\pgfqpoint{1.250000in}{0.660000in}}{\pgfqpoint{7.750000in}{4.620000in}}%
\pgfusepath{clip}%
\pgfsetbuttcap%
\pgfsetmiterjoin%
\definecolor{currentfill}{rgb}{0.000000,0.000000,0.000000}%
\pgfsetfillcolor{currentfill}%
\pgfsetlinewidth{0.000000pt}%
\definecolor{currentstroke}{rgb}{0.000000,0.000000,0.000000}%
\pgfsetstrokecolor{currentstroke}%
\pgfsetstrokeopacity{0.000000}%
\pgfsetdash{}{0pt}%
\pgfpathmoveto{\pgfqpoint{2.680650in}{0.660000in}}%
\pgfpathlineto{\pgfqpoint{2.686850in}{0.660000in}}%
\pgfpathlineto{\pgfqpoint{2.686850in}{3.375775in}}%
\pgfpathlineto{\pgfqpoint{2.680650in}{3.375775in}}%
\pgfpathlineto{\pgfqpoint{2.680650in}{0.660000in}}%
\pgfpathclose%
\pgfusepath{fill}%
\end{pgfscope}%
\begin{pgfscope}%
\pgfpathrectangle{\pgfqpoint{1.250000in}{0.660000in}}{\pgfqpoint{7.750000in}{4.620000in}}%
\pgfusepath{clip}%
\pgfsetbuttcap%
\pgfsetmiterjoin%
\definecolor{currentfill}{rgb}{0.000000,0.000000,0.000000}%
\pgfsetfillcolor{currentfill}%
\pgfsetlinewidth{0.000000pt}%
\definecolor{currentstroke}{rgb}{0.000000,0.000000,0.000000}%
\pgfsetstrokecolor{currentstroke}%
\pgfsetstrokeopacity{0.000000}%
\pgfsetdash{}{0pt}%
\pgfpathmoveto{\pgfqpoint{2.688400in}{0.660000in}}%
\pgfpathlineto{\pgfqpoint{2.694600in}{0.660000in}}%
\pgfpathlineto{\pgfqpoint{2.694600in}{3.374978in}}%
\pgfpathlineto{\pgfqpoint{2.688400in}{3.374978in}}%
\pgfpathlineto{\pgfqpoint{2.688400in}{0.660000in}}%
\pgfpathclose%
\pgfusepath{fill}%
\end{pgfscope}%
\begin{pgfscope}%
\pgfpathrectangle{\pgfqpoint{1.250000in}{0.660000in}}{\pgfqpoint{7.750000in}{4.620000in}}%
\pgfusepath{clip}%
\pgfsetbuttcap%
\pgfsetmiterjoin%
\definecolor{currentfill}{rgb}{0.000000,0.000000,0.000000}%
\pgfsetfillcolor{currentfill}%
\pgfsetlinewidth{0.000000pt}%
\definecolor{currentstroke}{rgb}{0.000000,0.000000,0.000000}%
\pgfsetstrokecolor{currentstroke}%
\pgfsetstrokeopacity{0.000000}%
\pgfsetdash{}{0pt}%
\pgfpathmoveto{\pgfqpoint{2.696150in}{0.660000in}}%
\pgfpathlineto{\pgfqpoint{2.702350in}{0.660000in}}%
\pgfpathlineto{\pgfqpoint{2.702350in}{3.374186in}}%
\pgfpathlineto{\pgfqpoint{2.696150in}{3.374186in}}%
\pgfpathlineto{\pgfqpoint{2.696150in}{0.660000in}}%
\pgfpathclose%
\pgfusepath{fill}%
\end{pgfscope}%
\begin{pgfscope}%
\pgfpathrectangle{\pgfqpoint{1.250000in}{0.660000in}}{\pgfqpoint{7.750000in}{4.620000in}}%
\pgfusepath{clip}%
\pgfsetbuttcap%
\pgfsetmiterjoin%
\definecolor{currentfill}{rgb}{0.000000,0.000000,0.000000}%
\pgfsetfillcolor{currentfill}%
\pgfsetlinewidth{0.000000pt}%
\definecolor{currentstroke}{rgb}{0.000000,0.000000,0.000000}%
\pgfsetstrokecolor{currentstroke}%
\pgfsetstrokeopacity{0.000000}%
\pgfsetdash{}{0pt}%
\pgfpathmoveto{\pgfqpoint{2.703900in}{0.660000in}}%
\pgfpathlineto{\pgfqpoint{2.710100in}{0.660000in}}%
\pgfpathlineto{\pgfqpoint{2.710100in}{3.373401in}}%
\pgfpathlineto{\pgfqpoint{2.703900in}{3.373401in}}%
\pgfpathlineto{\pgfqpoint{2.703900in}{0.660000in}}%
\pgfpathclose%
\pgfusepath{fill}%
\end{pgfscope}%
\begin{pgfscope}%
\pgfpathrectangle{\pgfqpoint{1.250000in}{0.660000in}}{\pgfqpoint{7.750000in}{4.620000in}}%
\pgfusepath{clip}%
\pgfsetbuttcap%
\pgfsetmiterjoin%
\definecolor{currentfill}{rgb}{0.000000,0.000000,0.000000}%
\pgfsetfillcolor{currentfill}%
\pgfsetlinewidth{0.000000pt}%
\definecolor{currentstroke}{rgb}{0.000000,0.000000,0.000000}%
\pgfsetstrokecolor{currentstroke}%
\pgfsetstrokeopacity{0.000000}%
\pgfsetdash{}{0pt}%
\pgfpathmoveto{\pgfqpoint{2.711650in}{0.660000in}}%
\pgfpathlineto{\pgfqpoint{2.717850in}{0.660000in}}%
\pgfpathlineto{\pgfqpoint{2.717850in}{3.372620in}}%
\pgfpathlineto{\pgfqpoint{2.711650in}{3.372620in}}%
\pgfpathlineto{\pgfqpoint{2.711650in}{0.660000in}}%
\pgfpathclose%
\pgfusepath{fill}%
\end{pgfscope}%
\begin{pgfscope}%
\pgfpathrectangle{\pgfqpoint{1.250000in}{0.660000in}}{\pgfqpoint{7.750000in}{4.620000in}}%
\pgfusepath{clip}%
\pgfsetbuttcap%
\pgfsetmiterjoin%
\definecolor{currentfill}{rgb}{0.000000,0.000000,0.000000}%
\pgfsetfillcolor{currentfill}%
\pgfsetlinewidth{0.000000pt}%
\definecolor{currentstroke}{rgb}{0.000000,0.000000,0.000000}%
\pgfsetstrokecolor{currentstroke}%
\pgfsetstrokeopacity{0.000000}%
\pgfsetdash{}{0pt}%
\pgfpathmoveto{\pgfqpoint{2.719400in}{0.660000in}}%
\pgfpathlineto{\pgfqpoint{2.725600in}{0.660000in}}%
\pgfpathlineto{\pgfqpoint{2.725600in}{3.371847in}}%
\pgfpathlineto{\pgfqpoint{2.719400in}{3.371847in}}%
\pgfpathlineto{\pgfqpoint{2.719400in}{0.660000in}}%
\pgfpathclose%
\pgfusepath{fill}%
\end{pgfscope}%
\begin{pgfscope}%
\pgfpathrectangle{\pgfqpoint{1.250000in}{0.660000in}}{\pgfqpoint{7.750000in}{4.620000in}}%
\pgfusepath{clip}%
\pgfsetbuttcap%
\pgfsetmiterjoin%
\definecolor{currentfill}{rgb}{0.000000,0.000000,0.000000}%
\pgfsetfillcolor{currentfill}%
\pgfsetlinewidth{0.000000pt}%
\definecolor{currentstroke}{rgb}{0.000000,0.000000,0.000000}%
\pgfsetstrokecolor{currentstroke}%
\pgfsetstrokeopacity{0.000000}%
\pgfsetdash{}{0pt}%
\pgfpathmoveto{\pgfqpoint{2.727150in}{0.660000in}}%
\pgfpathlineto{\pgfqpoint{2.733350in}{0.660000in}}%
\pgfpathlineto{\pgfqpoint{2.733350in}{3.371079in}}%
\pgfpathlineto{\pgfqpoint{2.727150in}{3.371079in}}%
\pgfpathlineto{\pgfqpoint{2.727150in}{0.660000in}}%
\pgfpathclose%
\pgfusepath{fill}%
\end{pgfscope}%
\begin{pgfscope}%
\pgfpathrectangle{\pgfqpoint{1.250000in}{0.660000in}}{\pgfqpoint{7.750000in}{4.620000in}}%
\pgfusepath{clip}%
\pgfsetbuttcap%
\pgfsetmiterjoin%
\definecolor{currentfill}{rgb}{0.000000,0.000000,0.000000}%
\pgfsetfillcolor{currentfill}%
\pgfsetlinewidth{0.000000pt}%
\definecolor{currentstroke}{rgb}{0.000000,0.000000,0.000000}%
\pgfsetstrokecolor{currentstroke}%
\pgfsetstrokeopacity{0.000000}%
\pgfsetdash{}{0pt}%
\pgfpathmoveto{\pgfqpoint{2.734900in}{0.660000in}}%
\pgfpathlineto{\pgfqpoint{2.741100in}{0.660000in}}%
\pgfpathlineto{\pgfqpoint{2.741100in}{3.370316in}}%
\pgfpathlineto{\pgfqpoint{2.734900in}{3.370316in}}%
\pgfpathlineto{\pgfqpoint{2.734900in}{0.660000in}}%
\pgfpathclose%
\pgfusepath{fill}%
\end{pgfscope}%
\begin{pgfscope}%
\pgfpathrectangle{\pgfqpoint{1.250000in}{0.660000in}}{\pgfqpoint{7.750000in}{4.620000in}}%
\pgfusepath{clip}%
\pgfsetbuttcap%
\pgfsetmiterjoin%
\definecolor{currentfill}{rgb}{0.000000,0.000000,0.000000}%
\pgfsetfillcolor{currentfill}%
\pgfsetlinewidth{0.000000pt}%
\definecolor{currentstroke}{rgb}{0.000000,0.000000,0.000000}%
\pgfsetstrokecolor{currentstroke}%
\pgfsetstrokeopacity{0.000000}%
\pgfsetdash{}{0pt}%
\pgfpathmoveto{\pgfqpoint{2.742650in}{0.660000in}}%
\pgfpathlineto{\pgfqpoint{2.748850in}{0.660000in}}%
\pgfpathlineto{\pgfqpoint{2.748850in}{3.369559in}}%
\pgfpathlineto{\pgfqpoint{2.742650in}{3.369559in}}%
\pgfpathlineto{\pgfqpoint{2.742650in}{0.660000in}}%
\pgfpathclose%
\pgfusepath{fill}%
\end{pgfscope}%
\begin{pgfscope}%
\pgfpathrectangle{\pgfqpoint{1.250000in}{0.660000in}}{\pgfqpoint{7.750000in}{4.620000in}}%
\pgfusepath{clip}%
\pgfsetbuttcap%
\pgfsetmiterjoin%
\definecolor{currentfill}{rgb}{0.000000,0.000000,0.000000}%
\pgfsetfillcolor{currentfill}%
\pgfsetlinewidth{0.000000pt}%
\definecolor{currentstroke}{rgb}{0.000000,0.000000,0.000000}%
\pgfsetstrokecolor{currentstroke}%
\pgfsetstrokeopacity{0.000000}%
\pgfsetdash{}{0pt}%
\pgfpathmoveto{\pgfqpoint{2.750400in}{0.660000in}}%
\pgfpathlineto{\pgfqpoint{2.756600in}{0.660000in}}%
\pgfpathlineto{\pgfqpoint{2.756600in}{3.368806in}}%
\pgfpathlineto{\pgfqpoint{2.750400in}{3.368806in}}%
\pgfpathlineto{\pgfqpoint{2.750400in}{0.660000in}}%
\pgfpathclose%
\pgfusepath{fill}%
\end{pgfscope}%
\begin{pgfscope}%
\pgfpathrectangle{\pgfqpoint{1.250000in}{0.660000in}}{\pgfqpoint{7.750000in}{4.620000in}}%
\pgfusepath{clip}%
\pgfsetbuttcap%
\pgfsetmiterjoin%
\definecolor{currentfill}{rgb}{0.000000,0.000000,0.000000}%
\pgfsetfillcolor{currentfill}%
\pgfsetlinewidth{0.000000pt}%
\definecolor{currentstroke}{rgb}{0.000000,0.000000,0.000000}%
\pgfsetstrokecolor{currentstroke}%
\pgfsetstrokeopacity{0.000000}%
\pgfsetdash{}{0pt}%
\pgfpathmoveto{\pgfqpoint{2.758150in}{0.660000in}}%
\pgfpathlineto{\pgfqpoint{2.764350in}{0.660000in}}%
\pgfpathlineto{\pgfqpoint{2.764350in}{3.368061in}}%
\pgfpathlineto{\pgfqpoint{2.758150in}{3.368061in}}%
\pgfpathlineto{\pgfqpoint{2.758150in}{0.660000in}}%
\pgfpathclose%
\pgfusepath{fill}%
\end{pgfscope}%
\begin{pgfscope}%
\pgfpathrectangle{\pgfqpoint{1.250000in}{0.660000in}}{\pgfqpoint{7.750000in}{4.620000in}}%
\pgfusepath{clip}%
\pgfsetbuttcap%
\pgfsetmiterjoin%
\definecolor{currentfill}{rgb}{0.000000,0.000000,0.000000}%
\pgfsetfillcolor{currentfill}%
\pgfsetlinewidth{0.000000pt}%
\definecolor{currentstroke}{rgb}{0.000000,0.000000,0.000000}%
\pgfsetstrokecolor{currentstroke}%
\pgfsetstrokeopacity{0.000000}%
\pgfsetdash{}{0pt}%
\pgfpathmoveto{\pgfqpoint{2.765900in}{0.660000in}}%
\pgfpathlineto{\pgfqpoint{2.772100in}{0.660000in}}%
\pgfpathlineto{\pgfqpoint{2.772100in}{3.367317in}}%
\pgfpathlineto{\pgfqpoint{2.765900in}{3.367317in}}%
\pgfpathlineto{\pgfqpoint{2.765900in}{0.660000in}}%
\pgfpathclose%
\pgfusepath{fill}%
\end{pgfscope}%
\begin{pgfscope}%
\pgfpathrectangle{\pgfqpoint{1.250000in}{0.660000in}}{\pgfqpoint{7.750000in}{4.620000in}}%
\pgfusepath{clip}%
\pgfsetbuttcap%
\pgfsetmiterjoin%
\definecolor{currentfill}{rgb}{0.000000,0.000000,0.000000}%
\pgfsetfillcolor{currentfill}%
\pgfsetlinewidth{0.000000pt}%
\definecolor{currentstroke}{rgb}{0.000000,0.000000,0.000000}%
\pgfsetstrokecolor{currentstroke}%
\pgfsetstrokeopacity{0.000000}%
\pgfsetdash{}{0pt}%
\pgfpathmoveto{\pgfqpoint{2.773650in}{0.660000in}}%
\pgfpathlineto{\pgfqpoint{2.779850in}{0.660000in}}%
\pgfpathlineto{\pgfqpoint{2.779850in}{3.366583in}}%
\pgfpathlineto{\pgfqpoint{2.773650in}{3.366583in}}%
\pgfpathlineto{\pgfqpoint{2.773650in}{0.660000in}}%
\pgfpathclose%
\pgfusepath{fill}%
\end{pgfscope}%
\begin{pgfscope}%
\pgfpathrectangle{\pgfqpoint{1.250000in}{0.660000in}}{\pgfqpoint{7.750000in}{4.620000in}}%
\pgfusepath{clip}%
\pgfsetbuttcap%
\pgfsetmiterjoin%
\definecolor{currentfill}{rgb}{0.000000,0.000000,0.000000}%
\pgfsetfillcolor{currentfill}%
\pgfsetlinewidth{0.000000pt}%
\definecolor{currentstroke}{rgb}{0.000000,0.000000,0.000000}%
\pgfsetstrokecolor{currentstroke}%
\pgfsetstrokeopacity{0.000000}%
\pgfsetdash{}{0pt}%
\pgfpathmoveto{\pgfqpoint{2.781400in}{0.660000in}}%
\pgfpathlineto{\pgfqpoint{2.787600in}{0.660000in}}%
\pgfpathlineto{\pgfqpoint{2.787600in}{3.365851in}}%
\pgfpathlineto{\pgfqpoint{2.781400in}{3.365851in}}%
\pgfpathlineto{\pgfqpoint{2.781400in}{0.660000in}}%
\pgfpathclose%
\pgfusepath{fill}%
\end{pgfscope}%
\begin{pgfscope}%
\pgfpathrectangle{\pgfqpoint{1.250000in}{0.660000in}}{\pgfqpoint{7.750000in}{4.620000in}}%
\pgfusepath{clip}%
\pgfsetbuttcap%
\pgfsetmiterjoin%
\definecolor{currentfill}{rgb}{0.000000,0.000000,0.000000}%
\pgfsetfillcolor{currentfill}%
\pgfsetlinewidth{0.000000pt}%
\definecolor{currentstroke}{rgb}{0.000000,0.000000,0.000000}%
\pgfsetstrokecolor{currentstroke}%
\pgfsetstrokeopacity{0.000000}%
\pgfsetdash{}{0pt}%
\pgfpathmoveto{\pgfqpoint{2.789150in}{0.660000in}}%
\pgfpathlineto{\pgfqpoint{2.795350in}{0.660000in}}%
\pgfpathlineto{\pgfqpoint{2.795350in}{3.365127in}}%
\pgfpathlineto{\pgfqpoint{2.789150in}{3.365127in}}%
\pgfpathlineto{\pgfqpoint{2.789150in}{0.660000in}}%
\pgfpathclose%
\pgfusepath{fill}%
\end{pgfscope}%
\begin{pgfscope}%
\pgfpathrectangle{\pgfqpoint{1.250000in}{0.660000in}}{\pgfqpoint{7.750000in}{4.620000in}}%
\pgfusepath{clip}%
\pgfsetbuttcap%
\pgfsetmiterjoin%
\definecolor{currentfill}{rgb}{0.000000,0.000000,0.000000}%
\pgfsetfillcolor{currentfill}%
\pgfsetlinewidth{0.000000pt}%
\definecolor{currentstroke}{rgb}{0.000000,0.000000,0.000000}%
\pgfsetstrokecolor{currentstroke}%
\pgfsetstrokeopacity{0.000000}%
\pgfsetdash{}{0pt}%
\pgfpathmoveto{\pgfqpoint{2.796900in}{0.660000in}}%
\pgfpathlineto{\pgfqpoint{2.803100in}{0.660000in}}%
\pgfpathlineto{\pgfqpoint{2.803100in}{3.364405in}}%
\pgfpathlineto{\pgfqpoint{2.796900in}{3.364405in}}%
\pgfpathlineto{\pgfqpoint{2.796900in}{0.660000in}}%
\pgfpathclose%
\pgfusepath{fill}%
\end{pgfscope}%
\begin{pgfscope}%
\pgfpathrectangle{\pgfqpoint{1.250000in}{0.660000in}}{\pgfqpoint{7.750000in}{4.620000in}}%
\pgfusepath{clip}%
\pgfsetbuttcap%
\pgfsetmiterjoin%
\definecolor{currentfill}{rgb}{0.000000,0.000000,0.000000}%
\pgfsetfillcolor{currentfill}%
\pgfsetlinewidth{0.000000pt}%
\definecolor{currentstroke}{rgb}{0.000000,0.000000,0.000000}%
\pgfsetstrokecolor{currentstroke}%
\pgfsetstrokeopacity{0.000000}%
\pgfsetdash{}{0pt}%
\pgfpathmoveto{\pgfqpoint{2.804650in}{0.660000in}}%
\pgfpathlineto{\pgfqpoint{2.810850in}{0.660000in}}%
\pgfpathlineto{\pgfqpoint{2.810850in}{3.363690in}}%
\pgfpathlineto{\pgfqpoint{2.804650in}{3.363690in}}%
\pgfpathlineto{\pgfqpoint{2.804650in}{0.660000in}}%
\pgfpathclose%
\pgfusepath{fill}%
\end{pgfscope}%
\begin{pgfscope}%
\pgfpathrectangle{\pgfqpoint{1.250000in}{0.660000in}}{\pgfqpoint{7.750000in}{4.620000in}}%
\pgfusepath{clip}%
\pgfsetbuttcap%
\pgfsetmiterjoin%
\definecolor{currentfill}{rgb}{0.000000,0.000000,0.000000}%
\pgfsetfillcolor{currentfill}%
\pgfsetlinewidth{0.000000pt}%
\definecolor{currentstroke}{rgb}{0.000000,0.000000,0.000000}%
\pgfsetstrokecolor{currentstroke}%
\pgfsetstrokeopacity{0.000000}%
\pgfsetdash{}{0pt}%
\pgfpathmoveto{\pgfqpoint{2.812400in}{0.660000in}}%
\pgfpathlineto{\pgfqpoint{2.818600in}{0.660000in}}%
\pgfpathlineto{\pgfqpoint{2.818600in}{3.362979in}}%
\pgfpathlineto{\pgfqpoint{2.812400in}{3.362979in}}%
\pgfpathlineto{\pgfqpoint{2.812400in}{0.660000in}}%
\pgfpathclose%
\pgfusepath{fill}%
\end{pgfscope}%
\begin{pgfscope}%
\pgfpathrectangle{\pgfqpoint{1.250000in}{0.660000in}}{\pgfqpoint{7.750000in}{4.620000in}}%
\pgfusepath{clip}%
\pgfsetbuttcap%
\pgfsetmiterjoin%
\definecolor{currentfill}{rgb}{0.000000,0.000000,0.000000}%
\pgfsetfillcolor{currentfill}%
\pgfsetlinewidth{0.000000pt}%
\definecolor{currentstroke}{rgb}{0.000000,0.000000,0.000000}%
\pgfsetstrokecolor{currentstroke}%
\pgfsetstrokeopacity{0.000000}%
\pgfsetdash{}{0pt}%
\pgfpathmoveto{\pgfqpoint{2.820150in}{0.660000in}}%
\pgfpathlineto{\pgfqpoint{2.826350in}{0.660000in}}%
\pgfpathlineto{\pgfqpoint{2.826350in}{3.362273in}}%
\pgfpathlineto{\pgfqpoint{2.820150in}{3.362273in}}%
\pgfpathlineto{\pgfqpoint{2.820150in}{0.660000in}}%
\pgfpathclose%
\pgfusepath{fill}%
\end{pgfscope}%
\begin{pgfscope}%
\pgfpathrectangle{\pgfqpoint{1.250000in}{0.660000in}}{\pgfqpoint{7.750000in}{4.620000in}}%
\pgfusepath{clip}%
\pgfsetbuttcap%
\pgfsetmiterjoin%
\definecolor{currentfill}{rgb}{0.000000,0.000000,0.000000}%
\pgfsetfillcolor{currentfill}%
\pgfsetlinewidth{0.000000pt}%
\definecolor{currentstroke}{rgb}{0.000000,0.000000,0.000000}%
\pgfsetstrokecolor{currentstroke}%
\pgfsetstrokeopacity{0.000000}%
\pgfsetdash{}{0pt}%
\pgfpathmoveto{\pgfqpoint{2.827900in}{0.660000in}}%
\pgfpathlineto{\pgfqpoint{2.834100in}{0.660000in}}%
\pgfpathlineto{\pgfqpoint{2.834100in}{3.361571in}}%
\pgfpathlineto{\pgfqpoint{2.827900in}{3.361571in}}%
\pgfpathlineto{\pgfqpoint{2.827900in}{0.660000in}}%
\pgfpathclose%
\pgfusepath{fill}%
\end{pgfscope}%
\begin{pgfscope}%
\pgfpathrectangle{\pgfqpoint{1.250000in}{0.660000in}}{\pgfqpoint{7.750000in}{4.620000in}}%
\pgfusepath{clip}%
\pgfsetbuttcap%
\pgfsetmiterjoin%
\definecolor{currentfill}{rgb}{0.000000,0.000000,0.000000}%
\pgfsetfillcolor{currentfill}%
\pgfsetlinewidth{0.000000pt}%
\definecolor{currentstroke}{rgb}{0.000000,0.000000,0.000000}%
\pgfsetstrokecolor{currentstroke}%
\pgfsetstrokeopacity{0.000000}%
\pgfsetdash{}{0pt}%
\pgfpathmoveto{\pgfqpoint{2.835650in}{0.660000in}}%
\pgfpathlineto{\pgfqpoint{2.841850in}{0.660000in}}%
\pgfpathlineto{\pgfqpoint{2.841850in}{3.360875in}}%
\pgfpathlineto{\pgfqpoint{2.835650in}{3.360875in}}%
\pgfpathlineto{\pgfqpoint{2.835650in}{0.660000in}}%
\pgfpathclose%
\pgfusepath{fill}%
\end{pgfscope}%
\begin{pgfscope}%
\pgfpathrectangle{\pgfqpoint{1.250000in}{0.660000in}}{\pgfqpoint{7.750000in}{4.620000in}}%
\pgfusepath{clip}%
\pgfsetbuttcap%
\pgfsetmiterjoin%
\definecolor{currentfill}{rgb}{0.000000,0.000000,0.000000}%
\pgfsetfillcolor{currentfill}%
\pgfsetlinewidth{0.000000pt}%
\definecolor{currentstroke}{rgb}{0.000000,0.000000,0.000000}%
\pgfsetstrokecolor{currentstroke}%
\pgfsetstrokeopacity{0.000000}%
\pgfsetdash{}{0pt}%
\pgfpathmoveto{\pgfqpoint{2.843400in}{0.660000in}}%
\pgfpathlineto{\pgfqpoint{2.849600in}{0.660000in}}%
\pgfpathlineto{\pgfqpoint{2.849600in}{3.360183in}}%
\pgfpathlineto{\pgfqpoint{2.843400in}{3.360183in}}%
\pgfpathlineto{\pgfqpoint{2.843400in}{0.660000in}}%
\pgfpathclose%
\pgfusepath{fill}%
\end{pgfscope}%
\begin{pgfscope}%
\pgfpathrectangle{\pgfqpoint{1.250000in}{0.660000in}}{\pgfqpoint{7.750000in}{4.620000in}}%
\pgfusepath{clip}%
\pgfsetbuttcap%
\pgfsetmiterjoin%
\definecolor{currentfill}{rgb}{0.000000,0.000000,0.000000}%
\pgfsetfillcolor{currentfill}%
\pgfsetlinewidth{0.000000pt}%
\definecolor{currentstroke}{rgb}{0.000000,0.000000,0.000000}%
\pgfsetstrokecolor{currentstroke}%
\pgfsetstrokeopacity{0.000000}%
\pgfsetdash{}{0pt}%
\pgfpathmoveto{\pgfqpoint{2.851150in}{0.660000in}}%
\pgfpathlineto{\pgfqpoint{2.857350in}{0.660000in}}%
\pgfpathlineto{\pgfqpoint{2.857350in}{3.359496in}}%
\pgfpathlineto{\pgfqpoint{2.851150in}{3.359496in}}%
\pgfpathlineto{\pgfqpoint{2.851150in}{0.660000in}}%
\pgfpathclose%
\pgfusepath{fill}%
\end{pgfscope}%
\begin{pgfscope}%
\pgfpathrectangle{\pgfqpoint{1.250000in}{0.660000in}}{\pgfqpoint{7.750000in}{4.620000in}}%
\pgfusepath{clip}%
\pgfsetbuttcap%
\pgfsetmiterjoin%
\definecolor{currentfill}{rgb}{0.000000,0.000000,0.000000}%
\pgfsetfillcolor{currentfill}%
\pgfsetlinewidth{0.000000pt}%
\definecolor{currentstroke}{rgb}{0.000000,0.000000,0.000000}%
\pgfsetstrokecolor{currentstroke}%
\pgfsetstrokeopacity{0.000000}%
\pgfsetdash{}{0pt}%
\pgfpathmoveto{\pgfqpoint{2.858900in}{0.660000in}}%
\pgfpathlineto{\pgfqpoint{2.865100in}{0.660000in}}%
\pgfpathlineto{\pgfqpoint{2.865100in}{3.358814in}}%
\pgfpathlineto{\pgfqpoint{2.858900in}{3.358814in}}%
\pgfpathlineto{\pgfqpoint{2.858900in}{0.660000in}}%
\pgfpathclose%
\pgfusepath{fill}%
\end{pgfscope}%
\begin{pgfscope}%
\pgfpathrectangle{\pgfqpoint{1.250000in}{0.660000in}}{\pgfqpoint{7.750000in}{4.620000in}}%
\pgfusepath{clip}%
\pgfsetbuttcap%
\pgfsetmiterjoin%
\definecolor{currentfill}{rgb}{0.000000,0.000000,0.000000}%
\pgfsetfillcolor{currentfill}%
\pgfsetlinewidth{0.000000pt}%
\definecolor{currentstroke}{rgb}{0.000000,0.000000,0.000000}%
\pgfsetstrokecolor{currentstroke}%
\pgfsetstrokeopacity{0.000000}%
\pgfsetdash{}{0pt}%
\pgfpathmoveto{\pgfqpoint{2.866650in}{0.660000in}}%
\pgfpathlineto{\pgfqpoint{2.872850in}{0.660000in}}%
\pgfpathlineto{\pgfqpoint{2.872850in}{3.358136in}}%
\pgfpathlineto{\pgfqpoint{2.866650in}{3.358136in}}%
\pgfpathlineto{\pgfqpoint{2.866650in}{0.660000in}}%
\pgfpathclose%
\pgfusepath{fill}%
\end{pgfscope}%
\begin{pgfscope}%
\pgfpathrectangle{\pgfqpoint{1.250000in}{0.660000in}}{\pgfqpoint{7.750000in}{4.620000in}}%
\pgfusepath{clip}%
\pgfsetbuttcap%
\pgfsetmiterjoin%
\definecolor{currentfill}{rgb}{0.000000,0.000000,0.000000}%
\pgfsetfillcolor{currentfill}%
\pgfsetlinewidth{0.000000pt}%
\definecolor{currentstroke}{rgb}{0.000000,0.000000,0.000000}%
\pgfsetstrokecolor{currentstroke}%
\pgfsetstrokeopacity{0.000000}%
\pgfsetdash{}{0pt}%
\pgfpathmoveto{\pgfqpoint{2.874400in}{0.660000in}}%
\pgfpathlineto{\pgfqpoint{2.880600in}{0.660000in}}%
\pgfpathlineto{\pgfqpoint{2.880600in}{3.357462in}}%
\pgfpathlineto{\pgfqpoint{2.874400in}{3.357462in}}%
\pgfpathlineto{\pgfqpoint{2.874400in}{0.660000in}}%
\pgfpathclose%
\pgfusepath{fill}%
\end{pgfscope}%
\begin{pgfscope}%
\pgfpathrectangle{\pgfqpoint{1.250000in}{0.660000in}}{\pgfqpoint{7.750000in}{4.620000in}}%
\pgfusepath{clip}%
\pgfsetbuttcap%
\pgfsetmiterjoin%
\definecolor{currentfill}{rgb}{0.000000,0.000000,0.000000}%
\pgfsetfillcolor{currentfill}%
\pgfsetlinewidth{0.000000pt}%
\definecolor{currentstroke}{rgb}{0.000000,0.000000,0.000000}%
\pgfsetstrokecolor{currentstroke}%
\pgfsetstrokeopacity{0.000000}%
\pgfsetdash{}{0pt}%
\pgfpathmoveto{\pgfqpoint{2.882150in}{0.660000in}}%
\pgfpathlineto{\pgfqpoint{2.888350in}{0.660000in}}%
\pgfpathlineto{\pgfqpoint{2.888350in}{3.356793in}}%
\pgfpathlineto{\pgfqpoint{2.882150in}{3.356793in}}%
\pgfpathlineto{\pgfqpoint{2.882150in}{0.660000in}}%
\pgfpathclose%
\pgfusepath{fill}%
\end{pgfscope}%
\begin{pgfscope}%
\pgfpathrectangle{\pgfqpoint{1.250000in}{0.660000in}}{\pgfqpoint{7.750000in}{4.620000in}}%
\pgfusepath{clip}%
\pgfsetbuttcap%
\pgfsetmiterjoin%
\definecolor{currentfill}{rgb}{0.000000,0.000000,0.000000}%
\pgfsetfillcolor{currentfill}%
\pgfsetlinewidth{0.000000pt}%
\definecolor{currentstroke}{rgb}{0.000000,0.000000,0.000000}%
\pgfsetstrokecolor{currentstroke}%
\pgfsetstrokeopacity{0.000000}%
\pgfsetdash{}{0pt}%
\pgfpathmoveto{\pgfqpoint{2.889900in}{0.660000in}}%
\pgfpathlineto{\pgfqpoint{2.896100in}{0.660000in}}%
\pgfpathlineto{\pgfqpoint{2.896100in}{3.356128in}}%
\pgfpathlineto{\pgfqpoint{2.889900in}{3.356128in}}%
\pgfpathlineto{\pgfqpoint{2.889900in}{0.660000in}}%
\pgfpathclose%
\pgfusepath{fill}%
\end{pgfscope}%
\begin{pgfscope}%
\pgfpathrectangle{\pgfqpoint{1.250000in}{0.660000in}}{\pgfqpoint{7.750000in}{4.620000in}}%
\pgfusepath{clip}%
\pgfsetbuttcap%
\pgfsetmiterjoin%
\definecolor{currentfill}{rgb}{0.000000,0.000000,0.000000}%
\pgfsetfillcolor{currentfill}%
\pgfsetlinewidth{0.000000pt}%
\definecolor{currentstroke}{rgb}{0.000000,0.000000,0.000000}%
\pgfsetstrokecolor{currentstroke}%
\pgfsetstrokeopacity{0.000000}%
\pgfsetdash{}{0pt}%
\pgfpathmoveto{\pgfqpoint{2.897650in}{0.660000in}}%
\pgfpathlineto{\pgfqpoint{2.903850in}{0.660000in}}%
\pgfpathlineto{\pgfqpoint{2.903850in}{3.355468in}}%
\pgfpathlineto{\pgfqpoint{2.897650in}{3.355468in}}%
\pgfpathlineto{\pgfqpoint{2.897650in}{0.660000in}}%
\pgfpathclose%
\pgfusepath{fill}%
\end{pgfscope}%
\begin{pgfscope}%
\pgfpathrectangle{\pgfqpoint{1.250000in}{0.660000in}}{\pgfqpoint{7.750000in}{4.620000in}}%
\pgfusepath{clip}%
\pgfsetbuttcap%
\pgfsetmiterjoin%
\definecolor{currentfill}{rgb}{0.000000,0.000000,0.000000}%
\pgfsetfillcolor{currentfill}%
\pgfsetlinewidth{0.000000pt}%
\definecolor{currentstroke}{rgb}{0.000000,0.000000,0.000000}%
\pgfsetstrokecolor{currentstroke}%
\pgfsetstrokeopacity{0.000000}%
\pgfsetdash{}{0pt}%
\pgfpathmoveto{\pgfqpoint{2.905400in}{0.660000in}}%
\pgfpathlineto{\pgfqpoint{2.911600in}{0.660000in}}%
\pgfpathlineto{\pgfqpoint{2.911600in}{3.354811in}}%
\pgfpathlineto{\pgfqpoint{2.905400in}{3.354811in}}%
\pgfpathlineto{\pgfqpoint{2.905400in}{0.660000in}}%
\pgfpathclose%
\pgfusepath{fill}%
\end{pgfscope}%
\begin{pgfscope}%
\pgfpathrectangle{\pgfqpoint{1.250000in}{0.660000in}}{\pgfqpoint{7.750000in}{4.620000in}}%
\pgfusepath{clip}%
\pgfsetbuttcap%
\pgfsetmiterjoin%
\definecolor{currentfill}{rgb}{0.000000,0.000000,0.000000}%
\pgfsetfillcolor{currentfill}%
\pgfsetlinewidth{0.000000pt}%
\definecolor{currentstroke}{rgb}{0.000000,0.000000,0.000000}%
\pgfsetstrokecolor{currentstroke}%
\pgfsetstrokeopacity{0.000000}%
\pgfsetdash{}{0pt}%
\pgfpathmoveto{\pgfqpoint{2.913150in}{0.660000in}}%
\pgfpathlineto{\pgfqpoint{2.919350in}{0.660000in}}%
\pgfpathlineto{\pgfqpoint{2.919350in}{3.354160in}}%
\pgfpathlineto{\pgfqpoint{2.913150in}{3.354160in}}%
\pgfpathlineto{\pgfqpoint{2.913150in}{0.660000in}}%
\pgfpathclose%
\pgfusepath{fill}%
\end{pgfscope}%
\begin{pgfscope}%
\pgfpathrectangle{\pgfqpoint{1.250000in}{0.660000in}}{\pgfqpoint{7.750000in}{4.620000in}}%
\pgfusepath{clip}%
\pgfsetbuttcap%
\pgfsetmiterjoin%
\definecolor{currentfill}{rgb}{0.000000,0.000000,0.000000}%
\pgfsetfillcolor{currentfill}%
\pgfsetlinewidth{0.000000pt}%
\definecolor{currentstroke}{rgb}{0.000000,0.000000,0.000000}%
\pgfsetstrokecolor{currentstroke}%
\pgfsetstrokeopacity{0.000000}%
\pgfsetdash{}{0pt}%
\pgfpathmoveto{\pgfqpoint{2.920900in}{0.660000in}}%
\pgfpathlineto{\pgfqpoint{2.927100in}{0.660000in}}%
\pgfpathlineto{\pgfqpoint{2.927100in}{3.353511in}}%
\pgfpathlineto{\pgfqpoint{2.920900in}{3.353511in}}%
\pgfpathlineto{\pgfqpoint{2.920900in}{0.660000in}}%
\pgfpathclose%
\pgfusepath{fill}%
\end{pgfscope}%
\begin{pgfscope}%
\pgfpathrectangle{\pgfqpoint{1.250000in}{0.660000in}}{\pgfqpoint{7.750000in}{4.620000in}}%
\pgfusepath{clip}%
\pgfsetbuttcap%
\pgfsetmiterjoin%
\definecolor{currentfill}{rgb}{0.000000,0.000000,0.000000}%
\pgfsetfillcolor{currentfill}%
\pgfsetlinewidth{0.000000pt}%
\definecolor{currentstroke}{rgb}{0.000000,0.000000,0.000000}%
\pgfsetstrokecolor{currentstroke}%
\pgfsetstrokeopacity{0.000000}%
\pgfsetdash{}{0pt}%
\pgfpathmoveto{\pgfqpoint{2.928650in}{0.660000in}}%
\pgfpathlineto{\pgfqpoint{2.934850in}{0.660000in}}%
\pgfpathlineto{\pgfqpoint{2.934850in}{3.352869in}}%
\pgfpathlineto{\pgfqpoint{2.928650in}{3.352869in}}%
\pgfpathlineto{\pgfqpoint{2.928650in}{0.660000in}}%
\pgfpathclose%
\pgfusepath{fill}%
\end{pgfscope}%
\begin{pgfscope}%
\pgfpathrectangle{\pgfqpoint{1.250000in}{0.660000in}}{\pgfqpoint{7.750000in}{4.620000in}}%
\pgfusepath{clip}%
\pgfsetbuttcap%
\pgfsetmiterjoin%
\definecolor{currentfill}{rgb}{0.000000,0.000000,0.000000}%
\pgfsetfillcolor{currentfill}%
\pgfsetlinewidth{0.000000pt}%
\definecolor{currentstroke}{rgb}{0.000000,0.000000,0.000000}%
\pgfsetstrokecolor{currentstroke}%
\pgfsetstrokeopacity{0.000000}%
\pgfsetdash{}{0pt}%
\pgfpathmoveto{\pgfqpoint{2.936400in}{0.660000in}}%
\pgfpathlineto{\pgfqpoint{2.942600in}{0.660000in}}%
\pgfpathlineto{\pgfqpoint{2.942600in}{3.352228in}}%
\pgfpathlineto{\pgfqpoint{2.936400in}{3.352228in}}%
\pgfpathlineto{\pgfqpoint{2.936400in}{0.660000in}}%
\pgfpathclose%
\pgfusepath{fill}%
\end{pgfscope}%
\begin{pgfscope}%
\pgfpathrectangle{\pgfqpoint{1.250000in}{0.660000in}}{\pgfqpoint{7.750000in}{4.620000in}}%
\pgfusepath{clip}%
\pgfsetbuttcap%
\pgfsetmiterjoin%
\definecolor{currentfill}{rgb}{0.000000,0.000000,0.000000}%
\pgfsetfillcolor{currentfill}%
\pgfsetlinewidth{0.000000pt}%
\definecolor{currentstroke}{rgb}{0.000000,0.000000,0.000000}%
\pgfsetstrokecolor{currentstroke}%
\pgfsetstrokeopacity{0.000000}%
\pgfsetdash{}{0pt}%
\pgfpathmoveto{\pgfqpoint{2.944150in}{0.660000in}}%
\pgfpathlineto{\pgfqpoint{2.950350in}{0.660000in}}%
\pgfpathlineto{\pgfqpoint{2.950350in}{3.351594in}}%
\pgfpathlineto{\pgfqpoint{2.944150in}{3.351594in}}%
\pgfpathlineto{\pgfqpoint{2.944150in}{0.660000in}}%
\pgfpathclose%
\pgfusepath{fill}%
\end{pgfscope}%
\begin{pgfscope}%
\pgfpathrectangle{\pgfqpoint{1.250000in}{0.660000in}}{\pgfqpoint{7.750000in}{4.620000in}}%
\pgfusepath{clip}%
\pgfsetbuttcap%
\pgfsetmiterjoin%
\definecolor{currentfill}{rgb}{0.000000,0.000000,0.000000}%
\pgfsetfillcolor{currentfill}%
\pgfsetlinewidth{0.000000pt}%
\definecolor{currentstroke}{rgb}{0.000000,0.000000,0.000000}%
\pgfsetstrokecolor{currentstroke}%
\pgfsetstrokeopacity{0.000000}%
\pgfsetdash{}{0pt}%
\pgfpathmoveto{\pgfqpoint{2.951900in}{0.660000in}}%
\pgfpathlineto{\pgfqpoint{2.958100in}{0.660000in}}%
\pgfpathlineto{\pgfqpoint{2.958100in}{3.350962in}}%
\pgfpathlineto{\pgfqpoint{2.951900in}{3.350962in}}%
\pgfpathlineto{\pgfqpoint{2.951900in}{0.660000in}}%
\pgfpathclose%
\pgfusepath{fill}%
\end{pgfscope}%
\begin{pgfscope}%
\pgfpathrectangle{\pgfqpoint{1.250000in}{0.660000in}}{\pgfqpoint{7.750000in}{4.620000in}}%
\pgfusepath{clip}%
\pgfsetbuttcap%
\pgfsetmiterjoin%
\definecolor{currentfill}{rgb}{0.000000,0.000000,0.000000}%
\pgfsetfillcolor{currentfill}%
\pgfsetlinewidth{0.000000pt}%
\definecolor{currentstroke}{rgb}{0.000000,0.000000,0.000000}%
\pgfsetstrokecolor{currentstroke}%
\pgfsetstrokeopacity{0.000000}%
\pgfsetdash{}{0pt}%
\pgfpathmoveto{\pgfqpoint{2.959650in}{0.660000in}}%
\pgfpathlineto{\pgfqpoint{2.965850in}{0.660000in}}%
\pgfpathlineto{\pgfqpoint{2.965850in}{3.350335in}}%
\pgfpathlineto{\pgfqpoint{2.959650in}{3.350335in}}%
\pgfpathlineto{\pgfqpoint{2.959650in}{0.660000in}}%
\pgfpathclose%
\pgfusepath{fill}%
\end{pgfscope}%
\begin{pgfscope}%
\pgfpathrectangle{\pgfqpoint{1.250000in}{0.660000in}}{\pgfqpoint{7.750000in}{4.620000in}}%
\pgfusepath{clip}%
\pgfsetbuttcap%
\pgfsetmiterjoin%
\definecolor{currentfill}{rgb}{0.000000,0.000000,0.000000}%
\pgfsetfillcolor{currentfill}%
\pgfsetlinewidth{0.000000pt}%
\definecolor{currentstroke}{rgb}{0.000000,0.000000,0.000000}%
\pgfsetstrokecolor{currentstroke}%
\pgfsetstrokeopacity{0.000000}%
\pgfsetdash{}{0pt}%
\pgfpathmoveto{\pgfqpoint{2.967400in}{0.660000in}}%
\pgfpathlineto{\pgfqpoint{2.973600in}{0.660000in}}%
\pgfpathlineto{\pgfqpoint{2.973600in}{3.349711in}}%
\pgfpathlineto{\pgfqpoint{2.967400in}{3.349711in}}%
\pgfpathlineto{\pgfqpoint{2.967400in}{0.660000in}}%
\pgfpathclose%
\pgfusepath{fill}%
\end{pgfscope}%
\begin{pgfscope}%
\pgfpathrectangle{\pgfqpoint{1.250000in}{0.660000in}}{\pgfqpoint{7.750000in}{4.620000in}}%
\pgfusepath{clip}%
\pgfsetbuttcap%
\pgfsetmiterjoin%
\definecolor{currentfill}{rgb}{0.000000,0.000000,0.000000}%
\pgfsetfillcolor{currentfill}%
\pgfsetlinewidth{0.000000pt}%
\definecolor{currentstroke}{rgb}{0.000000,0.000000,0.000000}%
\pgfsetstrokecolor{currentstroke}%
\pgfsetstrokeopacity{0.000000}%
\pgfsetdash{}{0pt}%
\pgfpathmoveto{\pgfqpoint{2.975150in}{0.660000in}}%
\pgfpathlineto{\pgfqpoint{2.981350in}{0.660000in}}%
\pgfpathlineto{\pgfqpoint{2.981350in}{3.349092in}}%
\pgfpathlineto{\pgfqpoint{2.975150in}{3.349092in}}%
\pgfpathlineto{\pgfqpoint{2.975150in}{0.660000in}}%
\pgfpathclose%
\pgfusepath{fill}%
\end{pgfscope}%
\begin{pgfscope}%
\pgfpathrectangle{\pgfqpoint{1.250000in}{0.660000in}}{\pgfqpoint{7.750000in}{4.620000in}}%
\pgfusepath{clip}%
\pgfsetbuttcap%
\pgfsetmiterjoin%
\definecolor{currentfill}{rgb}{0.000000,0.000000,0.000000}%
\pgfsetfillcolor{currentfill}%
\pgfsetlinewidth{0.000000pt}%
\definecolor{currentstroke}{rgb}{0.000000,0.000000,0.000000}%
\pgfsetstrokecolor{currentstroke}%
\pgfsetstrokeopacity{0.000000}%
\pgfsetdash{}{0pt}%
\pgfpathmoveto{\pgfqpoint{2.982900in}{0.660000in}}%
\pgfpathlineto{\pgfqpoint{2.989100in}{0.660000in}}%
\pgfpathlineto{\pgfqpoint{2.989100in}{3.348476in}}%
\pgfpathlineto{\pgfqpoint{2.982900in}{3.348476in}}%
\pgfpathlineto{\pgfqpoint{2.982900in}{0.660000in}}%
\pgfpathclose%
\pgfusepath{fill}%
\end{pgfscope}%
\begin{pgfscope}%
\pgfpathrectangle{\pgfqpoint{1.250000in}{0.660000in}}{\pgfqpoint{7.750000in}{4.620000in}}%
\pgfusepath{clip}%
\pgfsetbuttcap%
\pgfsetmiterjoin%
\definecolor{currentfill}{rgb}{0.000000,0.000000,0.000000}%
\pgfsetfillcolor{currentfill}%
\pgfsetlinewidth{0.000000pt}%
\definecolor{currentstroke}{rgb}{0.000000,0.000000,0.000000}%
\pgfsetstrokecolor{currentstroke}%
\pgfsetstrokeopacity{0.000000}%
\pgfsetdash{}{0pt}%
\pgfpathmoveto{\pgfqpoint{2.990650in}{0.660000in}}%
\pgfpathlineto{\pgfqpoint{2.996850in}{0.660000in}}%
\pgfpathlineto{\pgfqpoint{2.996850in}{3.347863in}}%
\pgfpathlineto{\pgfqpoint{2.990650in}{3.347863in}}%
\pgfpathlineto{\pgfqpoint{2.990650in}{0.660000in}}%
\pgfpathclose%
\pgfusepath{fill}%
\end{pgfscope}%
\begin{pgfscope}%
\pgfpathrectangle{\pgfqpoint{1.250000in}{0.660000in}}{\pgfqpoint{7.750000in}{4.620000in}}%
\pgfusepath{clip}%
\pgfsetbuttcap%
\pgfsetmiterjoin%
\definecolor{currentfill}{rgb}{0.000000,0.000000,0.000000}%
\pgfsetfillcolor{currentfill}%
\pgfsetlinewidth{0.000000pt}%
\definecolor{currentstroke}{rgb}{0.000000,0.000000,0.000000}%
\pgfsetstrokecolor{currentstroke}%
\pgfsetstrokeopacity{0.000000}%
\pgfsetdash{}{0pt}%
\pgfpathmoveto{\pgfqpoint{2.998400in}{0.660000in}}%
\pgfpathlineto{\pgfqpoint{3.004600in}{0.660000in}}%
\pgfpathlineto{\pgfqpoint{3.004600in}{3.347256in}}%
\pgfpathlineto{\pgfqpoint{2.998400in}{3.347256in}}%
\pgfpathlineto{\pgfqpoint{2.998400in}{0.660000in}}%
\pgfpathclose%
\pgfusepath{fill}%
\end{pgfscope}%
\begin{pgfscope}%
\pgfpathrectangle{\pgfqpoint{1.250000in}{0.660000in}}{\pgfqpoint{7.750000in}{4.620000in}}%
\pgfusepath{clip}%
\pgfsetbuttcap%
\pgfsetmiterjoin%
\definecolor{currentfill}{rgb}{0.000000,0.000000,0.000000}%
\pgfsetfillcolor{currentfill}%
\pgfsetlinewidth{0.000000pt}%
\definecolor{currentstroke}{rgb}{0.000000,0.000000,0.000000}%
\pgfsetstrokecolor{currentstroke}%
\pgfsetstrokeopacity{0.000000}%
\pgfsetdash{}{0pt}%
\pgfpathmoveto{\pgfqpoint{3.006150in}{0.660000in}}%
\pgfpathlineto{\pgfqpoint{3.012350in}{0.660000in}}%
\pgfpathlineto{\pgfqpoint{3.012350in}{3.346650in}}%
\pgfpathlineto{\pgfqpoint{3.006150in}{3.346650in}}%
\pgfpathlineto{\pgfqpoint{3.006150in}{0.660000in}}%
\pgfpathclose%
\pgfusepath{fill}%
\end{pgfscope}%
\begin{pgfscope}%
\pgfpathrectangle{\pgfqpoint{1.250000in}{0.660000in}}{\pgfqpoint{7.750000in}{4.620000in}}%
\pgfusepath{clip}%
\pgfsetbuttcap%
\pgfsetmiterjoin%
\definecolor{currentfill}{rgb}{0.000000,0.000000,0.000000}%
\pgfsetfillcolor{currentfill}%
\pgfsetlinewidth{0.000000pt}%
\definecolor{currentstroke}{rgb}{0.000000,0.000000,0.000000}%
\pgfsetstrokecolor{currentstroke}%
\pgfsetstrokeopacity{0.000000}%
\pgfsetdash{}{0pt}%
\pgfpathmoveto{\pgfqpoint{3.013900in}{0.660000in}}%
\pgfpathlineto{\pgfqpoint{3.020100in}{0.660000in}}%
\pgfpathlineto{\pgfqpoint{3.020100in}{3.346051in}}%
\pgfpathlineto{\pgfqpoint{3.013900in}{3.346051in}}%
\pgfpathlineto{\pgfqpoint{3.013900in}{0.660000in}}%
\pgfpathclose%
\pgfusepath{fill}%
\end{pgfscope}%
\begin{pgfscope}%
\pgfpathrectangle{\pgfqpoint{1.250000in}{0.660000in}}{\pgfqpoint{7.750000in}{4.620000in}}%
\pgfusepath{clip}%
\pgfsetbuttcap%
\pgfsetmiterjoin%
\definecolor{currentfill}{rgb}{0.000000,0.000000,0.000000}%
\pgfsetfillcolor{currentfill}%
\pgfsetlinewidth{0.000000pt}%
\definecolor{currentstroke}{rgb}{0.000000,0.000000,0.000000}%
\pgfsetstrokecolor{currentstroke}%
\pgfsetstrokeopacity{0.000000}%
\pgfsetdash{}{0pt}%
\pgfpathmoveto{\pgfqpoint{3.021650in}{0.660000in}}%
\pgfpathlineto{\pgfqpoint{3.027850in}{0.660000in}}%
\pgfpathlineto{\pgfqpoint{3.027850in}{3.345453in}}%
\pgfpathlineto{\pgfqpoint{3.021650in}{3.345453in}}%
\pgfpathlineto{\pgfqpoint{3.021650in}{0.660000in}}%
\pgfpathclose%
\pgfusepath{fill}%
\end{pgfscope}%
\begin{pgfscope}%
\pgfpathrectangle{\pgfqpoint{1.250000in}{0.660000in}}{\pgfqpoint{7.750000in}{4.620000in}}%
\pgfusepath{clip}%
\pgfsetbuttcap%
\pgfsetmiterjoin%
\definecolor{currentfill}{rgb}{0.000000,0.000000,0.000000}%
\pgfsetfillcolor{currentfill}%
\pgfsetlinewidth{0.000000pt}%
\definecolor{currentstroke}{rgb}{0.000000,0.000000,0.000000}%
\pgfsetstrokecolor{currentstroke}%
\pgfsetstrokeopacity{0.000000}%
\pgfsetdash{}{0pt}%
\pgfpathmoveto{\pgfqpoint{3.029400in}{0.660000in}}%
\pgfpathlineto{\pgfqpoint{3.035600in}{0.660000in}}%
\pgfpathlineto{\pgfqpoint{3.035600in}{3.344859in}}%
\pgfpathlineto{\pgfqpoint{3.029400in}{3.344859in}}%
\pgfpathlineto{\pgfqpoint{3.029400in}{0.660000in}}%
\pgfpathclose%
\pgfusepath{fill}%
\end{pgfscope}%
\begin{pgfscope}%
\pgfpathrectangle{\pgfqpoint{1.250000in}{0.660000in}}{\pgfqpoint{7.750000in}{4.620000in}}%
\pgfusepath{clip}%
\pgfsetbuttcap%
\pgfsetmiterjoin%
\definecolor{currentfill}{rgb}{0.000000,0.000000,0.000000}%
\pgfsetfillcolor{currentfill}%
\pgfsetlinewidth{0.000000pt}%
\definecolor{currentstroke}{rgb}{0.000000,0.000000,0.000000}%
\pgfsetstrokecolor{currentstroke}%
\pgfsetstrokeopacity{0.000000}%
\pgfsetdash{}{0pt}%
\pgfpathmoveto{\pgfqpoint{3.037150in}{0.660000in}}%
\pgfpathlineto{\pgfqpoint{3.043350in}{0.660000in}}%
\pgfpathlineto{\pgfqpoint{3.043350in}{3.344270in}}%
\pgfpathlineto{\pgfqpoint{3.037150in}{3.344270in}}%
\pgfpathlineto{\pgfqpoint{3.037150in}{0.660000in}}%
\pgfpathclose%
\pgfusepath{fill}%
\end{pgfscope}%
\begin{pgfscope}%
\pgfpathrectangle{\pgfqpoint{1.250000in}{0.660000in}}{\pgfqpoint{7.750000in}{4.620000in}}%
\pgfusepath{clip}%
\pgfsetbuttcap%
\pgfsetmiterjoin%
\definecolor{currentfill}{rgb}{0.000000,0.000000,0.000000}%
\pgfsetfillcolor{currentfill}%
\pgfsetlinewidth{0.000000pt}%
\definecolor{currentstroke}{rgb}{0.000000,0.000000,0.000000}%
\pgfsetstrokecolor{currentstroke}%
\pgfsetstrokeopacity{0.000000}%
\pgfsetdash{}{0pt}%
\pgfpathmoveto{\pgfqpoint{3.044900in}{0.660000in}}%
\pgfpathlineto{\pgfqpoint{3.051100in}{0.660000in}}%
\pgfpathlineto{\pgfqpoint{3.051100in}{3.343681in}}%
\pgfpathlineto{\pgfqpoint{3.044900in}{3.343681in}}%
\pgfpathlineto{\pgfqpoint{3.044900in}{0.660000in}}%
\pgfpathclose%
\pgfusepath{fill}%
\end{pgfscope}%
\begin{pgfscope}%
\pgfpathrectangle{\pgfqpoint{1.250000in}{0.660000in}}{\pgfqpoint{7.750000in}{4.620000in}}%
\pgfusepath{clip}%
\pgfsetbuttcap%
\pgfsetmiterjoin%
\definecolor{currentfill}{rgb}{0.000000,0.000000,0.000000}%
\pgfsetfillcolor{currentfill}%
\pgfsetlinewidth{0.000000pt}%
\definecolor{currentstroke}{rgb}{0.000000,0.000000,0.000000}%
\pgfsetstrokecolor{currentstroke}%
\pgfsetstrokeopacity{0.000000}%
\pgfsetdash{}{0pt}%
\pgfpathmoveto{\pgfqpoint{3.052650in}{0.660000in}}%
\pgfpathlineto{\pgfqpoint{3.058850in}{0.660000in}}%
\pgfpathlineto{\pgfqpoint{3.058850in}{3.343100in}}%
\pgfpathlineto{\pgfqpoint{3.052650in}{3.343100in}}%
\pgfpathlineto{\pgfqpoint{3.052650in}{0.660000in}}%
\pgfpathclose%
\pgfusepath{fill}%
\end{pgfscope}%
\begin{pgfscope}%
\pgfpathrectangle{\pgfqpoint{1.250000in}{0.660000in}}{\pgfqpoint{7.750000in}{4.620000in}}%
\pgfusepath{clip}%
\pgfsetbuttcap%
\pgfsetmiterjoin%
\definecolor{currentfill}{rgb}{0.000000,0.000000,0.000000}%
\pgfsetfillcolor{currentfill}%
\pgfsetlinewidth{0.000000pt}%
\definecolor{currentstroke}{rgb}{0.000000,0.000000,0.000000}%
\pgfsetstrokecolor{currentstroke}%
\pgfsetstrokeopacity{0.000000}%
\pgfsetdash{}{0pt}%
\pgfpathmoveto{\pgfqpoint{3.060400in}{0.660000in}}%
\pgfpathlineto{\pgfqpoint{3.066600in}{0.660000in}}%
\pgfpathlineto{\pgfqpoint{3.066600in}{3.342519in}}%
\pgfpathlineto{\pgfqpoint{3.060400in}{3.342519in}}%
\pgfpathlineto{\pgfqpoint{3.060400in}{0.660000in}}%
\pgfpathclose%
\pgfusepath{fill}%
\end{pgfscope}%
\begin{pgfscope}%
\pgfpathrectangle{\pgfqpoint{1.250000in}{0.660000in}}{\pgfqpoint{7.750000in}{4.620000in}}%
\pgfusepath{clip}%
\pgfsetbuttcap%
\pgfsetmiterjoin%
\definecolor{currentfill}{rgb}{0.000000,0.000000,0.000000}%
\pgfsetfillcolor{currentfill}%
\pgfsetlinewidth{0.000000pt}%
\definecolor{currentstroke}{rgb}{0.000000,0.000000,0.000000}%
\pgfsetstrokecolor{currentstroke}%
\pgfsetstrokeopacity{0.000000}%
\pgfsetdash{}{0pt}%
\pgfpathmoveto{\pgfqpoint{3.068150in}{0.660000in}}%
\pgfpathlineto{\pgfqpoint{3.074350in}{0.660000in}}%
\pgfpathlineto{\pgfqpoint{3.074350in}{3.341944in}}%
\pgfpathlineto{\pgfqpoint{3.068150in}{3.341944in}}%
\pgfpathlineto{\pgfqpoint{3.068150in}{0.660000in}}%
\pgfpathclose%
\pgfusepath{fill}%
\end{pgfscope}%
\begin{pgfscope}%
\pgfpathrectangle{\pgfqpoint{1.250000in}{0.660000in}}{\pgfqpoint{7.750000in}{4.620000in}}%
\pgfusepath{clip}%
\pgfsetbuttcap%
\pgfsetmiterjoin%
\definecolor{currentfill}{rgb}{0.000000,0.000000,0.000000}%
\pgfsetfillcolor{currentfill}%
\pgfsetlinewidth{0.000000pt}%
\definecolor{currentstroke}{rgb}{0.000000,0.000000,0.000000}%
\pgfsetstrokecolor{currentstroke}%
\pgfsetstrokeopacity{0.000000}%
\pgfsetdash{}{0pt}%
\pgfpathmoveto{\pgfqpoint{3.075900in}{0.660000in}}%
\pgfpathlineto{\pgfqpoint{3.082100in}{0.660000in}}%
\pgfpathlineto{\pgfqpoint{3.082100in}{3.341371in}}%
\pgfpathlineto{\pgfqpoint{3.075900in}{3.341371in}}%
\pgfpathlineto{\pgfqpoint{3.075900in}{0.660000in}}%
\pgfpathclose%
\pgfusepath{fill}%
\end{pgfscope}%
\begin{pgfscope}%
\pgfpathrectangle{\pgfqpoint{1.250000in}{0.660000in}}{\pgfqpoint{7.750000in}{4.620000in}}%
\pgfusepath{clip}%
\pgfsetbuttcap%
\pgfsetmiterjoin%
\definecolor{currentfill}{rgb}{0.000000,0.000000,0.000000}%
\pgfsetfillcolor{currentfill}%
\pgfsetlinewidth{0.000000pt}%
\definecolor{currentstroke}{rgb}{0.000000,0.000000,0.000000}%
\pgfsetstrokecolor{currentstroke}%
\pgfsetstrokeopacity{0.000000}%
\pgfsetdash{}{0pt}%
\pgfpathmoveto{\pgfqpoint{3.083650in}{0.660000in}}%
\pgfpathlineto{\pgfqpoint{3.089850in}{0.660000in}}%
\pgfpathlineto{\pgfqpoint{3.089850in}{3.340800in}}%
\pgfpathlineto{\pgfqpoint{3.083650in}{3.340800in}}%
\pgfpathlineto{\pgfqpoint{3.083650in}{0.660000in}}%
\pgfpathclose%
\pgfusepath{fill}%
\end{pgfscope}%
\begin{pgfscope}%
\pgfpathrectangle{\pgfqpoint{1.250000in}{0.660000in}}{\pgfqpoint{7.750000in}{4.620000in}}%
\pgfusepath{clip}%
\pgfsetbuttcap%
\pgfsetmiterjoin%
\definecolor{currentfill}{rgb}{0.000000,0.000000,0.000000}%
\pgfsetfillcolor{currentfill}%
\pgfsetlinewidth{0.000000pt}%
\definecolor{currentstroke}{rgb}{0.000000,0.000000,0.000000}%
\pgfsetstrokecolor{currentstroke}%
\pgfsetstrokeopacity{0.000000}%
\pgfsetdash{}{0pt}%
\pgfpathmoveto{\pgfqpoint{3.091400in}{0.660000in}}%
\pgfpathlineto{\pgfqpoint{3.097600in}{0.660000in}}%
\pgfpathlineto{\pgfqpoint{3.097600in}{3.340236in}}%
\pgfpathlineto{\pgfqpoint{3.091400in}{3.340236in}}%
\pgfpathlineto{\pgfqpoint{3.091400in}{0.660000in}}%
\pgfpathclose%
\pgfusepath{fill}%
\end{pgfscope}%
\begin{pgfscope}%
\pgfpathrectangle{\pgfqpoint{1.250000in}{0.660000in}}{\pgfqpoint{7.750000in}{4.620000in}}%
\pgfusepath{clip}%
\pgfsetbuttcap%
\pgfsetmiterjoin%
\definecolor{currentfill}{rgb}{0.000000,0.000000,0.000000}%
\pgfsetfillcolor{currentfill}%
\pgfsetlinewidth{0.000000pt}%
\definecolor{currentstroke}{rgb}{0.000000,0.000000,0.000000}%
\pgfsetstrokecolor{currentstroke}%
\pgfsetstrokeopacity{0.000000}%
\pgfsetdash{}{0pt}%
\pgfpathmoveto{\pgfqpoint{3.099150in}{0.660000in}}%
\pgfpathlineto{\pgfqpoint{3.105350in}{0.660000in}}%
\pgfpathlineto{\pgfqpoint{3.105350in}{3.339672in}}%
\pgfpathlineto{\pgfqpoint{3.099150in}{3.339672in}}%
\pgfpathlineto{\pgfqpoint{3.099150in}{0.660000in}}%
\pgfpathclose%
\pgfusepath{fill}%
\end{pgfscope}%
\begin{pgfscope}%
\pgfpathrectangle{\pgfqpoint{1.250000in}{0.660000in}}{\pgfqpoint{7.750000in}{4.620000in}}%
\pgfusepath{clip}%
\pgfsetbuttcap%
\pgfsetmiterjoin%
\definecolor{currentfill}{rgb}{0.000000,0.000000,0.000000}%
\pgfsetfillcolor{currentfill}%
\pgfsetlinewidth{0.000000pt}%
\definecolor{currentstroke}{rgb}{0.000000,0.000000,0.000000}%
\pgfsetstrokecolor{currentstroke}%
\pgfsetstrokeopacity{0.000000}%
\pgfsetdash{}{0pt}%
\pgfpathmoveto{\pgfqpoint{3.106900in}{0.660000in}}%
\pgfpathlineto{\pgfqpoint{3.113100in}{0.660000in}}%
\pgfpathlineto{\pgfqpoint{3.113100in}{3.339113in}}%
\pgfpathlineto{\pgfqpoint{3.106900in}{3.339113in}}%
\pgfpathlineto{\pgfqpoint{3.106900in}{0.660000in}}%
\pgfpathclose%
\pgfusepath{fill}%
\end{pgfscope}%
\begin{pgfscope}%
\pgfpathrectangle{\pgfqpoint{1.250000in}{0.660000in}}{\pgfqpoint{7.750000in}{4.620000in}}%
\pgfusepath{clip}%
\pgfsetbuttcap%
\pgfsetmiterjoin%
\definecolor{currentfill}{rgb}{0.000000,0.000000,0.000000}%
\pgfsetfillcolor{currentfill}%
\pgfsetlinewidth{0.000000pt}%
\definecolor{currentstroke}{rgb}{0.000000,0.000000,0.000000}%
\pgfsetstrokecolor{currentstroke}%
\pgfsetstrokeopacity{0.000000}%
\pgfsetdash{}{0pt}%
\pgfpathmoveto{\pgfqpoint{3.114650in}{0.660000in}}%
\pgfpathlineto{\pgfqpoint{3.120850in}{0.660000in}}%
\pgfpathlineto{\pgfqpoint{3.120850in}{3.338557in}}%
\pgfpathlineto{\pgfqpoint{3.114650in}{3.338557in}}%
\pgfpathlineto{\pgfqpoint{3.114650in}{0.660000in}}%
\pgfpathclose%
\pgfusepath{fill}%
\end{pgfscope}%
\begin{pgfscope}%
\pgfpathrectangle{\pgfqpoint{1.250000in}{0.660000in}}{\pgfqpoint{7.750000in}{4.620000in}}%
\pgfusepath{clip}%
\pgfsetbuttcap%
\pgfsetmiterjoin%
\definecolor{currentfill}{rgb}{0.000000,0.000000,0.000000}%
\pgfsetfillcolor{currentfill}%
\pgfsetlinewidth{0.000000pt}%
\definecolor{currentstroke}{rgb}{0.000000,0.000000,0.000000}%
\pgfsetstrokecolor{currentstroke}%
\pgfsetstrokeopacity{0.000000}%
\pgfsetdash{}{0pt}%
\pgfpathmoveto{\pgfqpoint{3.122400in}{0.660000in}}%
\pgfpathlineto{\pgfqpoint{3.128600in}{0.660000in}}%
\pgfpathlineto{\pgfqpoint{3.128600in}{3.338002in}}%
\pgfpathlineto{\pgfqpoint{3.122400in}{3.338002in}}%
\pgfpathlineto{\pgfqpoint{3.122400in}{0.660000in}}%
\pgfpathclose%
\pgfusepath{fill}%
\end{pgfscope}%
\begin{pgfscope}%
\pgfpathrectangle{\pgfqpoint{1.250000in}{0.660000in}}{\pgfqpoint{7.750000in}{4.620000in}}%
\pgfusepath{clip}%
\pgfsetbuttcap%
\pgfsetmiterjoin%
\definecolor{currentfill}{rgb}{0.000000,0.000000,0.000000}%
\pgfsetfillcolor{currentfill}%
\pgfsetlinewidth{0.000000pt}%
\definecolor{currentstroke}{rgb}{0.000000,0.000000,0.000000}%
\pgfsetstrokecolor{currentstroke}%
\pgfsetstrokeopacity{0.000000}%
\pgfsetdash{}{0pt}%
\pgfpathmoveto{\pgfqpoint{3.130150in}{0.660000in}}%
\pgfpathlineto{\pgfqpoint{3.136350in}{0.660000in}}%
\pgfpathlineto{\pgfqpoint{3.136350in}{3.337454in}}%
\pgfpathlineto{\pgfqpoint{3.130150in}{3.337454in}}%
\pgfpathlineto{\pgfqpoint{3.130150in}{0.660000in}}%
\pgfpathclose%
\pgfusepath{fill}%
\end{pgfscope}%
\begin{pgfscope}%
\pgfpathrectangle{\pgfqpoint{1.250000in}{0.660000in}}{\pgfqpoint{7.750000in}{4.620000in}}%
\pgfusepath{clip}%
\pgfsetbuttcap%
\pgfsetmiterjoin%
\definecolor{currentfill}{rgb}{0.000000,0.000000,0.000000}%
\pgfsetfillcolor{currentfill}%
\pgfsetlinewidth{0.000000pt}%
\definecolor{currentstroke}{rgb}{0.000000,0.000000,0.000000}%
\pgfsetstrokecolor{currentstroke}%
\pgfsetstrokeopacity{0.000000}%
\pgfsetdash{}{0pt}%
\pgfpathmoveto{\pgfqpoint{3.137900in}{0.660000in}}%
\pgfpathlineto{\pgfqpoint{3.144100in}{0.660000in}}%
\pgfpathlineto{\pgfqpoint{3.144100in}{3.336906in}}%
\pgfpathlineto{\pgfqpoint{3.137900in}{3.336906in}}%
\pgfpathlineto{\pgfqpoint{3.137900in}{0.660000in}}%
\pgfpathclose%
\pgfusepath{fill}%
\end{pgfscope}%
\begin{pgfscope}%
\pgfpathrectangle{\pgfqpoint{1.250000in}{0.660000in}}{\pgfqpoint{7.750000in}{4.620000in}}%
\pgfusepath{clip}%
\pgfsetbuttcap%
\pgfsetmiterjoin%
\definecolor{currentfill}{rgb}{0.000000,0.000000,0.000000}%
\pgfsetfillcolor{currentfill}%
\pgfsetlinewidth{0.000000pt}%
\definecolor{currentstroke}{rgb}{0.000000,0.000000,0.000000}%
\pgfsetstrokecolor{currentstroke}%
\pgfsetstrokeopacity{0.000000}%
\pgfsetdash{}{0pt}%
\pgfpathmoveto{\pgfqpoint{3.145650in}{0.660000in}}%
\pgfpathlineto{\pgfqpoint{3.151850in}{0.660000in}}%
\pgfpathlineto{\pgfqpoint{3.151850in}{3.336362in}}%
\pgfpathlineto{\pgfqpoint{3.145650in}{3.336362in}}%
\pgfpathlineto{\pgfqpoint{3.145650in}{0.660000in}}%
\pgfpathclose%
\pgfusepath{fill}%
\end{pgfscope}%
\begin{pgfscope}%
\pgfpathrectangle{\pgfqpoint{1.250000in}{0.660000in}}{\pgfqpoint{7.750000in}{4.620000in}}%
\pgfusepath{clip}%
\pgfsetbuttcap%
\pgfsetmiterjoin%
\definecolor{currentfill}{rgb}{0.000000,0.000000,0.000000}%
\pgfsetfillcolor{currentfill}%
\pgfsetlinewidth{0.000000pt}%
\definecolor{currentstroke}{rgb}{0.000000,0.000000,0.000000}%
\pgfsetstrokecolor{currentstroke}%
\pgfsetstrokeopacity{0.000000}%
\pgfsetdash{}{0pt}%
\pgfpathmoveto{\pgfqpoint{3.153400in}{0.660000in}}%
\pgfpathlineto{\pgfqpoint{3.159600in}{0.660000in}}%
\pgfpathlineto{\pgfqpoint{3.159600in}{3.335822in}}%
\pgfpathlineto{\pgfqpoint{3.153400in}{3.335822in}}%
\pgfpathlineto{\pgfqpoint{3.153400in}{0.660000in}}%
\pgfpathclose%
\pgfusepath{fill}%
\end{pgfscope}%
\begin{pgfscope}%
\pgfpathrectangle{\pgfqpoint{1.250000in}{0.660000in}}{\pgfqpoint{7.750000in}{4.620000in}}%
\pgfusepath{clip}%
\pgfsetbuttcap%
\pgfsetmiterjoin%
\definecolor{currentfill}{rgb}{0.000000,0.000000,0.000000}%
\pgfsetfillcolor{currentfill}%
\pgfsetlinewidth{0.000000pt}%
\definecolor{currentstroke}{rgb}{0.000000,0.000000,0.000000}%
\pgfsetstrokecolor{currentstroke}%
\pgfsetstrokeopacity{0.000000}%
\pgfsetdash{}{0pt}%
\pgfpathmoveto{\pgfqpoint{3.161150in}{0.660000in}}%
\pgfpathlineto{\pgfqpoint{3.167350in}{0.660000in}}%
\pgfpathlineto{\pgfqpoint{3.167350in}{3.335283in}}%
\pgfpathlineto{\pgfqpoint{3.161150in}{3.335283in}}%
\pgfpathlineto{\pgfqpoint{3.161150in}{0.660000in}}%
\pgfpathclose%
\pgfusepath{fill}%
\end{pgfscope}%
\begin{pgfscope}%
\pgfpathrectangle{\pgfqpoint{1.250000in}{0.660000in}}{\pgfqpoint{7.750000in}{4.620000in}}%
\pgfusepath{clip}%
\pgfsetbuttcap%
\pgfsetmiterjoin%
\definecolor{currentfill}{rgb}{0.000000,0.000000,0.000000}%
\pgfsetfillcolor{currentfill}%
\pgfsetlinewidth{0.000000pt}%
\definecolor{currentstroke}{rgb}{0.000000,0.000000,0.000000}%
\pgfsetstrokecolor{currentstroke}%
\pgfsetstrokeopacity{0.000000}%
\pgfsetdash{}{0pt}%
\pgfpathmoveto{\pgfqpoint{3.168900in}{0.660000in}}%
\pgfpathlineto{\pgfqpoint{3.175100in}{0.660000in}}%
\pgfpathlineto{\pgfqpoint{3.175100in}{3.334749in}}%
\pgfpathlineto{\pgfqpoint{3.168900in}{3.334749in}}%
\pgfpathlineto{\pgfqpoint{3.168900in}{0.660000in}}%
\pgfpathclose%
\pgfusepath{fill}%
\end{pgfscope}%
\begin{pgfscope}%
\pgfpathrectangle{\pgfqpoint{1.250000in}{0.660000in}}{\pgfqpoint{7.750000in}{4.620000in}}%
\pgfusepath{clip}%
\pgfsetbuttcap%
\pgfsetmiterjoin%
\definecolor{currentfill}{rgb}{0.000000,0.000000,0.000000}%
\pgfsetfillcolor{currentfill}%
\pgfsetlinewidth{0.000000pt}%
\definecolor{currentstroke}{rgb}{0.000000,0.000000,0.000000}%
\pgfsetstrokecolor{currentstroke}%
\pgfsetstrokeopacity{0.000000}%
\pgfsetdash{}{0pt}%
\pgfpathmoveto{\pgfqpoint{3.176650in}{0.660000in}}%
\pgfpathlineto{\pgfqpoint{3.182850in}{0.660000in}}%
\pgfpathlineto{\pgfqpoint{3.182850in}{3.334218in}}%
\pgfpathlineto{\pgfqpoint{3.176650in}{3.334218in}}%
\pgfpathlineto{\pgfqpoint{3.176650in}{0.660000in}}%
\pgfpathclose%
\pgfusepath{fill}%
\end{pgfscope}%
\begin{pgfscope}%
\pgfpathrectangle{\pgfqpoint{1.250000in}{0.660000in}}{\pgfqpoint{7.750000in}{4.620000in}}%
\pgfusepath{clip}%
\pgfsetbuttcap%
\pgfsetmiterjoin%
\definecolor{currentfill}{rgb}{0.000000,0.000000,0.000000}%
\pgfsetfillcolor{currentfill}%
\pgfsetlinewidth{0.000000pt}%
\definecolor{currentstroke}{rgb}{0.000000,0.000000,0.000000}%
\pgfsetstrokecolor{currentstroke}%
\pgfsetstrokeopacity{0.000000}%
\pgfsetdash{}{0pt}%
\pgfpathmoveto{\pgfqpoint{3.184400in}{0.660000in}}%
\pgfpathlineto{\pgfqpoint{3.190600in}{0.660000in}}%
\pgfpathlineto{\pgfqpoint{3.190600in}{3.333687in}}%
\pgfpathlineto{\pgfqpoint{3.184400in}{3.333687in}}%
\pgfpathlineto{\pgfqpoint{3.184400in}{0.660000in}}%
\pgfpathclose%
\pgfusepath{fill}%
\end{pgfscope}%
\begin{pgfscope}%
\pgfpathrectangle{\pgfqpoint{1.250000in}{0.660000in}}{\pgfqpoint{7.750000in}{4.620000in}}%
\pgfusepath{clip}%
\pgfsetbuttcap%
\pgfsetmiterjoin%
\definecolor{currentfill}{rgb}{0.000000,0.000000,0.000000}%
\pgfsetfillcolor{currentfill}%
\pgfsetlinewidth{0.000000pt}%
\definecolor{currentstroke}{rgb}{0.000000,0.000000,0.000000}%
\pgfsetstrokecolor{currentstroke}%
\pgfsetstrokeopacity{0.000000}%
\pgfsetdash{}{0pt}%
\pgfpathmoveto{\pgfqpoint{3.192150in}{0.660000in}}%
\pgfpathlineto{\pgfqpoint{3.198350in}{0.660000in}}%
\pgfpathlineto{\pgfqpoint{3.198350in}{3.333163in}}%
\pgfpathlineto{\pgfqpoint{3.192150in}{3.333163in}}%
\pgfpathlineto{\pgfqpoint{3.192150in}{0.660000in}}%
\pgfpathclose%
\pgfusepath{fill}%
\end{pgfscope}%
\begin{pgfscope}%
\pgfpathrectangle{\pgfqpoint{1.250000in}{0.660000in}}{\pgfqpoint{7.750000in}{4.620000in}}%
\pgfusepath{clip}%
\pgfsetbuttcap%
\pgfsetmiterjoin%
\definecolor{currentfill}{rgb}{0.000000,0.000000,0.000000}%
\pgfsetfillcolor{currentfill}%
\pgfsetlinewidth{0.000000pt}%
\definecolor{currentstroke}{rgb}{0.000000,0.000000,0.000000}%
\pgfsetstrokecolor{currentstroke}%
\pgfsetstrokeopacity{0.000000}%
\pgfsetdash{}{0pt}%
\pgfpathmoveto{\pgfqpoint{3.199900in}{0.660000in}}%
\pgfpathlineto{\pgfqpoint{3.206100in}{0.660000in}}%
\pgfpathlineto{\pgfqpoint{3.206100in}{3.332640in}}%
\pgfpathlineto{\pgfqpoint{3.199900in}{3.332640in}}%
\pgfpathlineto{\pgfqpoint{3.199900in}{0.660000in}}%
\pgfpathclose%
\pgfusepath{fill}%
\end{pgfscope}%
\begin{pgfscope}%
\pgfpathrectangle{\pgfqpoint{1.250000in}{0.660000in}}{\pgfqpoint{7.750000in}{4.620000in}}%
\pgfusepath{clip}%
\pgfsetbuttcap%
\pgfsetmiterjoin%
\definecolor{currentfill}{rgb}{0.000000,0.000000,0.000000}%
\pgfsetfillcolor{currentfill}%
\pgfsetlinewidth{0.000000pt}%
\definecolor{currentstroke}{rgb}{0.000000,0.000000,0.000000}%
\pgfsetstrokecolor{currentstroke}%
\pgfsetstrokeopacity{0.000000}%
\pgfsetdash{}{0pt}%
\pgfpathmoveto{\pgfqpoint{3.207650in}{0.660000in}}%
\pgfpathlineto{\pgfqpoint{3.213850in}{0.660000in}}%
\pgfpathlineto{\pgfqpoint{3.213850in}{3.332119in}}%
\pgfpathlineto{\pgfqpoint{3.207650in}{3.332119in}}%
\pgfpathlineto{\pgfqpoint{3.207650in}{0.660000in}}%
\pgfpathclose%
\pgfusepath{fill}%
\end{pgfscope}%
\begin{pgfscope}%
\pgfpathrectangle{\pgfqpoint{1.250000in}{0.660000in}}{\pgfqpoint{7.750000in}{4.620000in}}%
\pgfusepath{clip}%
\pgfsetbuttcap%
\pgfsetmiterjoin%
\definecolor{currentfill}{rgb}{0.000000,0.000000,0.000000}%
\pgfsetfillcolor{currentfill}%
\pgfsetlinewidth{0.000000pt}%
\definecolor{currentstroke}{rgb}{0.000000,0.000000,0.000000}%
\pgfsetstrokecolor{currentstroke}%
\pgfsetstrokeopacity{0.000000}%
\pgfsetdash{}{0pt}%
\pgfpathmoveto{\pgfqpoint{3.215400in}{0.660000in}}%
\pgfpathlineto{\pgfqpoint{3.221600in}{0.660000in}}%
\pgfpathlineto{\pgfqpoint{3.221600in}{3.331603in}}%
\pgfpathlineto{\pgfqpoint{3.215400in}{3.331603in}}%
\pgfpathlineto{\pgfqpoint{3.215400in}{0.660000in}}%
\pgfpathclose%
\pgfusepath{fill}%
\end{pgfscope}%
\begin{pgfscope}%
\pgfpathrectangle{\pgfqpoint{1.250000in}{0.660000in}}{\pgfqpoint{7.750000in}{4.620000in}}%
\pgfusepath{clip}%
\pgfsetbuttcap%
\pgfsetmiterjoin%
\definecolor{currentfill}{rgb}{0.000000,0.000000,0.000000}%
\pgfsetfillcolor{currentfill}%
\pgfsetlinewidth{0.000000pt}%
\definecolor{currentstroke}{rgb}{0.000000,0.000000,0.000000}%
\pgfsetstrokecolor{currentstroke}%
\pgfsetstrokeopacity{0.000000}%
\pgfsetdash{}{0pt}%
\pgfpathmoveto{\pgfqpoint{3.223150in}{0.660000in}}%
\pgfpathlineto{\pgfqpoint{3.229350in}{0.660000in}}%
\pgfpathlineto{\pgfqpoint{3.229350in}{3.331087in}}%
\pgfpathlineto{\pgfqpoint{3.223150in}{3.331087in}}%
\pgfpathlineto{\pgfqpoint{3.223150in}{0.660000in}}%
\pgfpathclose%
\pgfusepath{fill}%
\end{pgfscope}%
\begin{pgfscope}%
\pgfpathrectangle{\pgfqpoint{1.250000in}{0.660000in}}{\pgfqpoint{7.750000in}{4.620000in}}%
\pgfusepath{clip}%
\pgfsetbuttcap%
\pgfsetmiterjoin%
\definecolor{currentfill}{rgb}{0.000000,0.000000,0.000000}%
\pgfsetfillcolor{currentfill}%
\pgfsetlinewidth{0.000000pt}%
\definecolor{currentstroke}{rgb}{0.000000,0.000000,0.000000}%
\pgfsetstrokecolor{currentstroke}%
\pgfsetstrokeopacity{0.000000}%
\pgfsetdash{}{0pt}%
\pgfpathmoveto{\pgfqpoint{3.230900in}{0.660000in}}%
\pgfpathlineto{\pgfqpoint{3.237100in}{0.660000in}}%
\pgfpathlineto{\pgfqpoint{3.237100in}{3.330575in}}%
\pgfpathlineto{\pgfqpoint{3.230900in}{3.330575in}}%
\pgfpathlineto{\pgfqpoint{3.230900in}{0.660000in}}%
\pgfpathclose%
\pgfusepath{fill}%
\end{pgfscope}%
\begin{pgfscope}%
\pgfpathrectangle{\pgfqpoint{1.250000in}{0.660000in}}{\pgfqpoint{7.750000in}{4.620000in}}%
\pgfusepath{clip}%
\pgfsetbuttcap%
\pgfsetmiterjoin%
\definecolor{currentfill}{rgb}{0.000000,0.000000,0.000000}%
\pgfsetfillcolor{currentfill}%
\pgfsetlinewidth{0.000000pt}%
\definecolor{currentstroke}{rgb}{0.000000,0.000000,0.000000}%
\pgfsetstrokecolor{currentstroke}%
\pgfsetstrokeopacity{0.000000}%
\pgfsetdash{}{0pt}%
\pgfpathmoveto{\pgfqpoint{3.238650in}{0.660000in}}%
\pgfpathlineto{\pgfqpoint{3.244850in}{0.660000in}}%
\pgfpathlineto{\pgfqpoint{3.244850in}{3.330068in}}%
\pgfpathlineto{\pgfqpoint{3.238650in}{3.330068in}}%
\pgfpathlineto{\pgfqpoint{3.238650in}{0.660000in}}%
\pgfpathclose%
\pgfusepath{fill}%
\end{pgfscope}%
\begin{pgfscope}%
\pgfpathrectangle{\pgfqpoint{1.250000in}{0.660000in}}{\pgfqpoint{7.750000in}{4.620000in}}%
\pgfusepath{clip}%
\pgfsetbuttcap%
\pgfsetmiterjoin%
\definecolor{currentfill}{rgb}{0.000000,0.000000,0.000000}%
\pgfsetfillcolor{currentfill}%
\pgfsetlinewidth{0.000000pt}%
\definecolor{currentstroke}{rgb}{0.000000,0.000000,0.000000}%
\pgfsetstrokecolor{currentstroke}%
\pgfsetstrokeopacity{0.000000}%
\pgfsetdash{}{0pt}%
\pgfpathmoveto{\pgfqpoint{3.246400in}{0.660000in}}%
\pgfpathlineto{\pgfqpoint{3.252600in}{0.660000in}}%
\pgfpathlineto{\pgfqpoint{3.252600in}{3.329560in}}%
\pgfpathlineto{\pgfqpoint{3.246400in}{3.329560in}}%
\pgfpathlineto{\pgfqpoint{3.246400in}{0.660000in}}%
\pgfpathclose%
\pgfusepath{fill}%
\end{pgfscope}%
\begin{pgfscope}%
\pgfpathrectangle{\pgfqpoint{1.250000in}{0.660000in}}{\pgfqpoint{7.750000in}{4.620000in}}%
\pgfusepath{clip}%
\pgfsetbuttcap%
\pgfsetmiterjoin%
\definecolor{currentfill}{rgb}{0.000000,0.000000,0.000000}%
\pgfsetfillcolor{currentfill}%
\pgfsetlinewidth{0.000000pt}%
\definecolor{currentstroke}{rgb}{0.000000,0.000000,0.000000}%
\pgfsetstrokecolor{currentstroke}%
\pgfsetstrokeopacity{0.000000}%
\pgfsetdash{}{0pt}%
\pgfpathmoveto{\pgfqpoint{3.254150in}{0.660000in}}%
\pgfpathlineto{\pgfqpoint{3.260350in}{0.660000in}}%
\pgfpathlineto{\pgfqpoint{3.260350in}{3.329057in}}%
\pgfpathlineto{\pgfqpoint{3.254150in}{3.329057in}}%
\pgfpathlineto{\pgfqpoint{3.254150in}{0.660000in}}%
\pgfpathclose%
\pgfusepath{fill}%
\end{pgfscope}%
\begin{pgfscope}%
\pgfpathrectangle{\pgfqpoint{1.250000in}{0.660000in}}{\pgfqpoint{7.750000in}{4.620000in}}%
\pgfusepath{clip}%
\pgfsetbuttcap%
\pgfsetmiterjoin%
\definecolor{currentfill}{rgb}{0.000000,0.000000,0.000000}%
\pgfsetfillcolor{currentfill}%
\pgfsetlinewidth{0.000000pt}%
\definecolor{currentstroke}{rgb}{0.000000,0.000000,0.000000}%
\pgfsetstrokecolor{currentstroke}%
\pgfsetstrokeopacity{0.000000}%
\pgfsetdash{}{0pt}%
\pgfpathmoveto{\pgfqpoint{3.261900in}{0.660000in}}%
\pgfpathlineto{\pgfqpoint{3.268100in}{0.660000in}}%
\pgfpathlineto{\pgfqpoint{3.268100in}{3.328557in}}%
\pgfpathlineto{\pgfqpoint{3.261900in}{3.328557in}}%
\pgfpathlineto{\pgfqpoint{3.261900in}{0.660000in}}%
\pgfpathclose%
\pgfusepath{fill}%
\end{pgfscope}%
\begin{pgfscope}%
\pgfpathrectangle{\pgfqpoint{1.250000in}{0.660000in}}{\pgfqpoint{7.750000in}{4.620000in}}%
\pgfusepath{clip}%
\pgfsetbuttcap%
\pgfsetmiterjoin%
\definecolor{currentfill}{rgb}{0.000000,0.000000,0.000000}%
\pgfsetfillcolor{currentfill}%
\pgfsetlinewidth{0.000000pt}%
\definecolor{currentstroke}{rgb}{0.000000,0.000000,0.000000}%
\pgfsetstrokecolor{currentstroke}%
\pgfsetstrokeopacity{0.000000}%
\pgfsetdash{}{0pt}%
\pgfpathmoveto{\pgfqpoint{3.269650in}{0.660000in}}%
\pgfpathlineto{\pgfqpoint{3.275850in}{0.660000in}}%
\pgfpathlineto{\pgfqpoint{3.275850in}{3.328057in}}%
\pgfpathlineto{\pgfqpoint{3.269650in}{3.328057in}}%
\pgfpathlineto{\pgfqpoint{3.269650in}{0.660000in}}%
\pgfpathclose%
\pgfusepath{fill}%
\end{pgfscope}%
\begin{pgfscope}%
\pgfpathrectangle{\pgfqpoint{1.250000in}{0.660000in}}{\pgfqpoint{7.750000in}{4.620000in}}%
\pgfusepath{clip}%
\pgfsetbuttcap%
\pgfsetmiterjoin%
\definecolor{currentfill}{rgb}{0.000000,0.000000,0.000000}%
\pgfsetfillcolor{currentfill}%
\pgfsetlinewidth{0.000000pt}%
\definecolor{currentstroke}{rgb}{0.000000,0.000000,0.000000}%
\pgfsetstrokecolor{currentstroke}%
\pgfsetstrokeopacity{0.000000}%
\pgfsetdash{}{0pt}%
\pgfpathmoveto{\pgfqpoint{3.277400in}{0.660000in}}%
\pgfpathlineto{\pgfqpoint{3.283600in}{0.660000in}}%
\pgfpathlineto{\pgfqpoint{3.283600in}{3.327562in}}%
\pgfpathlineto{\pgfqpoint{3.277400in}{3.327562in}}%
\pgfpathlineto{\pgfqpoint{3.277400in}{0.660000in}}%
\pgfpathclose%
\pgfusepath{fill}%
\end{pgfscope}%
\begin{pgfscope}%
\pgfpathrectangle{\pgfqpoint{1.250000in}{0.660000in}}{\pgfqpoint{7.750000in}{4.620000in}}%
\pgfusepath{clip}%
\pgfsetbuttcap%
\pgfsetmiterjoin%
\definecolor{currentfill}{rgb}{0.000000,0.000000,0.000000}%
\pgfsetfillcolor{currentfill}%
\pgfsetlinewidth{0.000000pt}%
\definecolor{currentstroke}{rgb}{0.000000,0.000000,0.000000}%
\pgfsetstrokecolor{currentstroke}%
\pgfsetstrokeopacity{0.000000}%
\pgfsetdash{}{0pt}%
\pgfpathmoveto{\pgfqpoint{3.285150in}{0.660000in}}%
\pgfpathlineto{\pgfqpoint{3.291350in}{0.660000in}}%
\pgfpathlineto{\pgfqpoint{3.291350in}{3.327070in}}%
\pgfpathlineto{\pgfqpoint{3.285150in}{3.327070in}}%
\pgfpathlineto{\pgfqpoint{3.285150in}{0.660000in}}%
\pgfpathclose%
\pgfusepath{fill}%
\end{pgfscope}%
\begin{pgfscope}%
\pgfpathrectangle{\pgfqpoint{1.250000in}{0.660000in}}{\pgfqpoint{7.750000in}{4.620000in}}%
\pgfusepath{clip}%
\pgfsetbuttcap%
\pgfsetmiterjoin%
\definecolor{currentfill}{rgb}{0.000000,0.000000,0.000000}%
\pgfsetfillcolor{currentfill}%
\pgfsetlinewidth{0.000000pt}%
\definecolor{currentstroke}{rgb}{0.000000,0.000000,0.000000}%
\pgfsetstrokecolor{currentstroke}%
\pgfsetstrokeopacity{0.000000}%
\pgfsetdash{}{0pt}%
\pgfpathmoveto{\pgfqpoint{3.292900in}{0.660000in}}%
\pgfpathlineto{\pgfqpoint{3.299100in}{0.660000in}}%
\pgfpathlineto{\pgfqpoint{3.299100in}{3.326578in}}%
\pgfpathlineto{\pgfqpoint{3.292900in}{3.326578in}}%
\pgfpathlineto{\pgfqpoint{3.292900in}{0.660000in}}%
\pgfpathclose%
\pgfusepath{fill}%
\end{pgfscope}%
\begin{pgfscope}%
\pgfpathrectangle{\pgfqpoint{1.250000in}{0.660000in}}{\pgfqpoint{7.750000in}{4.620000in}}%
\pgfusepath{clip}%
\pgfsetbuttcap%
\pgfsetmiterjoin%
\definecolor{currentfill}{rgb}{0.000000,0.000000,0.000000}%
\pgfsetfillcolor{currentfill}%
\pgfsetlinewidth{0.000000pt}%
\definecolor{currentstroke}{rgb}{0.000000,0.000000,0.000000}%
\pgfsetstrokecolor{currentstroke}%
\pgfsetstrokeopacity{0.000000}%
\pgfsetdash{}{0pt}%
\pgfpathmoveto{\pgfqpoint{3.300650in}{0.660000in}}%
\pgfpathlineto{\pgfqpoint{3.306850in}{0.660000in}}%
\pgfpathlineto{\pgfqpoint{3.306850in}{3.326091in}}%
\pgfpathlineto{\pgfqpoint{3.300650in}{3.326091in}}%
\pgfpathlineto{\pgfqpoint{3.300650in}{0.660000in}}%
\pgfpathclose%
\pgfusepath{fill}%
\end{pgfscope}%
\begin{pgfscope}%
\pgfpathrectangle{\pgfqpoint{1.250000in}{0.660000in}}{\pgfqpoint{7.750000in}{4.620000in}}%
\pgfusepath{clip}%
\pgfsetbuttcap%
\pgfsetmiterjoin%
\definecolor{currentfill}{rgb}{0.000000,0.000000,0.000000}%
\pgfsetfillcolor{currentfill}%
\pgfsetlinewidth{0.000000pt}%
\definecolor{currentstroke}{rgb}{0.000000,0.000000,0.000000}%
\pgfsetstrokecolor{currentstroke}%
\pgfsetstrokeopacity{0.000000}%
\pgfsetdash{}{0pt}%
\pgfpathmoveto{\pgfqpoint{3.308400in}{0.660000in}}%
\pgfpathlineto{\pgfqpoint{3.314600in}{0.660000in}}%
\pgfpathlineto{\pgfqpoint{3.314600in}{3.325607in}}%
\pgfpathlineto{\pgfqpoint{3.308400in}{3.325607in}}%
\pgfpathlineto{\pgfqpoint{3.308400in}{0.660000in}}%
\pgfpathclose%
\pgfusepath{fill}%
\end{pgfscope}%
\begin{pgfscope}%
\pgfpathrectangle{\pgfqpoint{1.250000in}{0.660000in}}{\pgfqpoint{7.750000in}{4.620000in}}%
\pgfusepath{clip}%
\pgfsetbuttcap%
\pgfsetmiterjoin%
\definecolor{currentfill}{rgb}{0.000000,0.000000,0.000000}%
\pgfsetfillcolor{currentfill}%
\pgfsetlinewidth{0.000000pt}%
\definecolor{currentstroke}{rgb}{0.000000,0.000000,0.000000}%
\pgfsetstrokecolor{currentstroke}%
\pgfsetstrokeopacity{0.000000}%
\pgfsetdash{}{0pt}%
\pgfpathmoveto{\pgfqpoint{3.316150in}{0.660000in}}%
\pgfpathlineto{\pgfqpoint{3.322350in}{0.660000in}}%
\pgfpathlineto{\pgfqpoint{3.322350in}{3.325122in}}%
\pgfpathlineto{\pgfqpoint{3.316150in}{3.325122in}}%
\pgfpathlineto{\pgfqpoint{3.316150in}{0.660000in}}%
\pgfpathclose%
\pgfusepath{fill}%
\end{pgfscope}%
\begin{pgfscope}%
\pgfpathrectangle{\pgfqpoint{1.250000in}{0.660000in}}{\pgfqpoint{7.750000in}{4.620000in}}%
\pgfusepath{clip}%
\pgfsetbuttcap%
\pgfsetmiterjoin%
\definecolor{currentfill}{rgb}{0.000000,0.000000,0.000000}%
\pgfsetfillcolor{currentfill}%
\pgfsetlinewidth{0.000000pt}%
\definecolor{currentstroke}{rgb}{0.000000,0.000000,0.000000}%
\pgfsetstrokecolor{currentstroke}%
\pgfsetstrokeopacity{0.000000}%
\pgfsetdash{}{0pt}%
\pgfpathmoveto{\pgfqpoint{3.323900in}{0.660000in}}%
\pgfpathlineto{\pgfqpoint{3.330100in}{0.660000in}}%
\pgfpathlineto{\pgfqpoint{3.330100in}{3.324643in}}%
\pgfpathlineto{\pgfqpoint{3.323900in}{3.324643in}}%
\pgfpathlineto{\pgfqpoint{3.323900in}{0.660000in}}%
\pgfpathclose%
\pgfusepath{fill}%
\end{pgfscope}%
\begin{pgfscope}%
\pgfpathrectangle{\pgfqpoint{1.250000in}{0.660000in}}{\pgfqpoint{7.750000in}{4.620000in}}%
\pgfusepath{clip}%
\pgfsetbuttcap%
\pgfsetmiterjoin%
\definecolor{currentfill}{rgb}{0.000000,0.000000,0.000000}%
\pgfsetfillcolor{currentfill}%
\pgfsetlinewidth{0.000000pt}%
\definecolor{currentstroke}{rgb}{0.000000,0.000000,0.000000}%
\pgfsetstrokecolor{currentstroke}%
\pgfsetstrokeopacity{0.000000}%
\pgfsetdash{}{0pt}%
\pgfpathmoveto{\pgfqpoint{3.331650in}{0.660000in}}%
\pgfpathlineto{\pgfqpoint{3.337850in}{0.660000in}}%
\pgfpathlineto{\pgfqpoint{3.337850in}{3.324166in}}%
\pgfpathlineto{\pgfqpoint{3.331650in}{3.324166in}}%
\pgfpathlineto{\pgfqpoint{3.331650in}{0.660000in}}%
\pgfpathclose%
\pgfusepath{fill}%
\end{pgfscope}%
\begin{pgfscope}%
\pgfpathrectangle{\pgfqpoint{1.250000in}{0.660000in}}{\pgfqpoint{7.750000in}{4.620000in}}%
\pgfusepath{clip}%
\pgfsetbuttcap%
\pgfsetmiterjoin%
\definecolor{currentfill}{rgb}{0.000000,0.000000,0.000000}%
\pgfsetfillcolor{currentfill}%
\pgfsetlinewidth{0.000000pt}%
\definecolor{currentstroke}{rgb}{0.000000,0.000000,0.000000}%
\pgfsetstrokecolor{currentstroke}%
\pgfsetstrokeopacity{0.000000}%
\pgfsetdash{}{0pt}%
\pgfpathmoveto{\pgfqpoint{3.339400in}{0.660000in}}%
\pgfpathlineto{\pgfqpoint{3.345600in}{0.660000in}}%
\pgfpathlineto{\pgfqpoint{3.345600in}{3.323689in}}%
\pgfpathlineto{\pgfqpoint{3.339400in}{3.323689in}}%
\pgfpathlineto{\pgfqpoint{3.339400in}{0.660000in}}%
\pgfpathclose%
\pgfusepath{fill}%
\end{pgfscope}%
\begin{pgfscope}%
\pgfpathrectangle{\pgfqpoint{1.250000in}{0.660000in}}{\pgfqpoint{7.750000in}{4.620000in}}%
\pgfusepath{clip}%
\pgfsetbuttcap%
\pgfsetmiterjoin%
\definecolor{currentfill}{rgb}{0.000000,0.000000,0.000000}%
\pgfsetfillcolor{currentfill}%
\pgfsetlinewidth{0.000000pt}%
\definecolor{currentstroke}{rgb}{0.000000,0.000000,0.000000}%
\pgfsetstrokecolor{currentstroke}%
\pgfsetstrokeopacity{0.000000}%
\pgfsetdash{}{0pt}%
\pgfpathmoveto{\pgfqpoint{3.347150in}{0.660000in}}%
\pgfpathlineto{\pgfqpoint{3.353350in}{0.660000in}}%
\pgfpathlineto{\pgfqpoint{3.353350in}{3.323217in}}%
\pgfpathlineto{\pgfqpoint{3.347150in}{3.323217in}}%
\pgfpathlineto{\pgfqpoint{3.347150in}{0.660000in}}%
\pgfpathclose%
\pgfusepath{fill}%
\end{pgfscope}%
\begin{pgfscope}%
\pgfpathrectangle{\pgfqpoint{1.250000in}{0.660000in}}{\pgfqpoint{7.750000in}{4.620000in}}%
\pgfusepath{clip}%
\pgfsetbuttcap%
\pgfsetmiterjoin%
\definecolor{currentfill}{rgb}{0.000000,0.000000,0.000000}%
\pgfsetfillcolor{currentfill}%
\pgfsetlinewidth{0.000000pt}%
\definecolor{currentstroke}{rgb}{0.000000,0.000000,0.000000}%
\pgfsetstrokecolor{currentstroke}%
\pgfsetstrokeopacity{0.000000}%
\pgfsetdash{}{0pt}%
\pgfpathmoveto{\pgfqpoint{3.354900in}{0.660000in}}%
\pgfpathlineto{\pgfqpoint{3.361100in}{0.660000in}}%
\pgfpathlineto{\pgfqpoint{3.361100in}{3.322747in}}%
\pgfpathlineto{\pgfqpoint{3.354900in}{3.322747in}}%
\pgfpathlineto{\pgfqpoint{3.354900in}{0.660000in}}%
\pgfpathclose%
\pgfusepath{fill}%
\end{pgfscope}%
\begin{pgfscope}%
\pgfpathrectangle{\pgfqpoint{1.250000in}{0.660000in}}{\pgfqpoint{7.750000in}{4.620000in}}%
\pgfusepath{clip}%
\pgfsetbuttcap%
\pgfsetmiterjoin%
\definecolor{currentfill}{rgb}{0.000000,0.000000,0.000000}%
\pgfsetfillcolor{currentfill}%
\pgfsetlinewidth{0.000000pt}%
\definecolor{currentstroke}{rgb}{0.000000,0.000000,0.000000}%
\pgfsetstrokecolor{currentstroke}%
\pgfsetstrokeopacity{0.000000}%
\pgfsetdash{}{0pt}%
\pgfpathmoveto{\pgfqpoint{3.362650in}{0.660000in}}%
\pgfpathlineto{\pgfqpoint{3.368850in}{0.660000in}}%
\pgfpathlineto{\pgfqpoint{3.368850in}{3.322278in}}%
\pgfpathlineto{\pgfqpoint{3.362650in}{3.322278in}}%
\pgfpathlineto{\pgfqpoint{3.362650in}{0.660000in}}%
\pgfpathclose%
\pgfusepath{fill}%
\end{pgfscope}%
\begin{pgfscope}%
\pgfpathrectangle{\pgfqpoint{1.250000in}{0.660000in}}{\pgfqpoint{7.750000in}{4.620000in}}%
\pgfusepath{clip}%
\pgfsetbuttcap%
\pgfsetmiterjoin%
\definecolor{currentfill}{rgb}{0.000000,0.000000,0.000000}%
\pgfsetfillcolor{currentfill}%
\pgfsetlinewidth{0.000000pt}%
\definecolor{currentstroke}{rgb}{0.000000,0.000000,0.000000}%
\pgfsetstrokecolor{currentstroke}%
\pgfsetstrokeopacity{0.000000}%
\pgfsetdash{}{0pt}%
\pgfpathmoveto{\pgfqpoint{3.370400in}{0.660000in}}%
\pgfpathlineto{\pgfqpoint{3.376600in}{0.660000in}}%
\pgfpathlineto{\pgfqpoint{3.376600in}{3.321812in}}%
\pgfpathlineto{\pgfqpoint{3.370400in}{3.321812in}}%
\pgfpathlineto{\pgfqpoint{3.370400in}{0.660000in}}%
\pgfpathclose%
\pgfusepath{fill}%
\end{pgfscope}%
\begin{pgfscope}%
\pgfpathrectangle{\pgfqpoint{1.250000in}{0.660000in}}{\pgfqpoint{7.750000in}{4.620000in}}%
\pgfusepath{clip}%
\pgfsetbuttcap%
\pgfsetmiterjoin%
\definecolor{currentfill}{rgb}{0.000000,0.000000,0.000000}%
\pgfsetfillcolor{currentfill}%
\pgfsetlinewidth{0.000000pt}%
\definecolor{currentstroke}{rgb}{0.000000,0.000000,0.000000}%
\pgfsetstrokecolor{currentstroke}%
\pgfsetstrokeopacity{0.000000}%
\pgfsetdash{}{0pt}%
\pgfpathmoveto{\pgfqpoint{3.378150in}{0.660000in}}%
\pgfpathlineto{\pgfqpoint{3.384350in}{0.660000in}}%
\pgfpathlineto{\pgfqpoint{3.384350in}{3.321350in}}%
\pgfpathlineto{\pgfqpoint{3.378150in}{3.321350in}}%
\pgfpathlineto{\pgfqpoint{3.378150in}{0.660000in}}%
\pgfpathclose%
\pgfusepath{fill}%
\end{pgfscope}%
\begin{pgfscope}%
\pgfpathrectangle{\pgfqpoint{1.250000in}{0.660000in}}{\pgfqpoint{7.750000in}{4.620000in}}%
\pgfusepath{clip}%
\pgfsetbuttcap%
\pgfsetmiterjoin%
\definecolor{currentfill}{rgb}{0.000000,0.000000,0.000000}%
\pgfsetfillcolor{currentfill}%
\pgfsetlinewidth{0.000000pt}%
\definecolor{currentstroke}{rgb}{0.000000,0.000000,0.000000}%
\pgfsetstrokecolor{currentstroke}%
\pgfsetstrokeopacity{0.000000}%
\pgfsetdash{}{0pt}%
\pgfpathmoveto{\pgfqpoint{3.385900in}{0.660000in}}%
\pgfpathlineto{\pgfqpoint{3.392100in}{0.660000in}}%
\pgfpathlineto{\pgfqpoint{3.392100in}{3.320888in}}%
\pgfpathlineto{\pgfqpoint{3.385900in}{3.320888in}}%
\pgfpathlineto{\pgfqpoint{3.385900in}{0.660000in}}%
\pgfpathclose%
\pgfusepath{fill}%
\end{pgfscope}%
\begin{pgfscope}%
\pgfpathrectangle{\pgfqpoint{1.250000in}{0.660000in}}{\pgfqpoint{7.750000in}{4.620000in}}%
\pgfusepath{clip}%
\pgfsetbuttcap%
\pgfsetmiterjoin%
\definecolor{currentfill}{rgb}{0.000000,0.000000,0.000000}%
\pgfsetfillcolor{currentfill}%
\pgfsetlinewidth{0.000000pt}%
\definecolor{currentstroke}{rgb}{0.000000,0.000000,0.000000}%
\pgfsetstrokecolor{currentstroke}%
\pgfsetstrokeopacity{0.000000}%
\pgfsetdash{}{0pt}%
\pgfpathmoveto{\pgfqpoint{3.393650in}{0.660000in}}%
\pgfpathlineto{\pgfqpoint{3.399850in}{0.660000in}}%
\pgfpathlineto{\pgfqpoint{3.399850in}{3.320428in}}%
\pgfpathlineto{\pgfqpoint{3.393650in}{3.320428in}}%
\pgfpathlineto{\pgfqpoint{3.393650in}{0.660000in}}%
\pgfpathclose%
\pgfusepath{fill}%
\end{pgfscope}%
\begin{pgfscope}%
\pgfpathrectangle{\pgfqpoint{1.250000in}{0.660000in}}{\pgfqpoint{7.750000in}{4.620000in}}%
\pgfusepath{clip}%
\pgfsetbuttcap%
\pgfsetmiterjoin%
\definecolor{currentfill}{rgb}{0.000000,0.000000,0.000000}%
\pgfsetfillcolor{currentfill}%
\pgfsetlinewidth{0.000000pt}%
\definecolor{currentstroke}{rgb}{0.000000,0.000000,0.000000}%
\pgfsetstrokecolor{currentstroke}%
\pgfsetstrokeopacity{0.000000}%
\pgfsetdash{}{0pt}%
\pgfpathmoveto{\pgfqpoint{3.401400in}{0.660000in}}%
\pgfpathlineto{\pgfqpoint{3.407600in}{0.660000in}}%
\pgfpathlineto{\pgfqpoint{3.407600in}{3.319974in}}%
\pgfpathlineto{\pgfqpoint{3.401400in}{3.319974in}}%
\pgfpathlineto{\pgfqpoint{3.401400in}{0.660000in}}%
\pgfpathclose%
\pgfusepath{fill}%
\end{pgfscope}%
\begin{pgfscope}%
\pgfpathrectangle{\pgfqpoint{1.250000in}{0.660000in}}{\pgfqpoint{7.750000in}{4.620000in}}%
\pgfusepath{clip}%
\pgfsetbuttcap%
\pgfsetmiterjoin%
\definecolor{currentfill}{rgb}{0.000000,0.000000,0.000000}%
\pgfsetfillcolor{currentfill}%
\pgfsetlinewidth{0.000000pt}%
\definecolor{currentstroke}{rgb}{0.000000,0.000000,0.000000}%
\pgfsetstrokecolor{currentstroke}%
\pgfsetstrokeopacity{0.000000}%
\pgfsetdash{}{0pt}%
\pgfpathmoveto{\pgfqpoint{3.409150in}{0.660000in}}%
\pgfpathlineto{\pgfqpoint{3.415350in}{0.660000in}}%
\pgfpathlineto{\pgfqpoint{3.415350in}{3.319519in}}%
\pgfpathlineto{\pgfqpoint{3.409150in}{3.319519in}}%
\pgfpathlineto{\pgfqpoint{3.409150in}{0.660000in}}%
\pgfpathclose%
\pgfusepath{fill}%
\end{pgfscope}%
\begin{pgfscope}%
\pgfpathrectangle{\pgfqpoint{1.250000in}{0.660000in}}{\pgfqpoint{7.750000in}{4.620000in}}%
\pgfusepath{clip}%
\pgfsetbuttcap%
\pgfsetmiterjoin%
\definecolor{currentfill}{rgb}{0.000000,0.000000,0.000000}%
\pgfsetfillcolor{currentfill}%
\pgfsetlinewidth{0.000000pt}%
\definecolor{currentstroke}{rgb}{0.000000,0.000000,0.000000}%
\pgfsetstrokecolor{currentstroke}%
\pgfsetstrokeopacity{0.000000}%
\pgfsetdash{}{0pt}%
\pgfpathmoveto{\pgfqpoint{3.416900in}{0.660000in}}%
\pgfpathlineto{\pgfqpoint{3.423100in}{0.660000in}}%
\pgfpathlineto{\pgfqpoint{3.423100in}{3.319065in}}%
\pgfpathlineto{\pgfqpoint{3.416900in}{3.319065in}}%
\pgfpathlineto{\pgfqpoint{3.416900in}{0.660000in}}%
\pgfpathclose%
\pgfusepath{fill}%
\end{pgfscope}%
\begin{pgfscope}%
\pgfpathrectangle{\pgfqpoint{1.250000in}{0.660000in}}{\pgfqpoint{7.750000in}{4.620000in}}%
\pgfusepath{clip}%
\pgfsetbuttcap%
\pgfsetmiterjoin%
\definecolor{currentfill}{rgb}{0.000000,0.000000,0.000000}%
\pgfsetfillcolor{currentfill}%
\pgfsetlinewidth{0.000000pt}%
\definecolor{currentstroke}{rgb}{0.000000,0.000000,0.000000}%
\pgfsetstrokecolor{currentstroke}%
\pgfsetstrokeopacity{0.000000}%
\pgfsetdash{}{0pt}%
\pgfpathmoveto{\pgfqpoint{3.424650in}{0.660000in}}%
\pgfpathlineto{\pgfqpoint{3.430850in}{0.660000in}}%
\pgfpathlineto{\pgfqpoint{3.430850in}{3.318617in}}%
\pgfpathlineto{\pgfqpoint{3.424650in}{3.318617in}}%
\pgfpathlineto{\pgfqpoint{3.424650in}{0.660000in}}%
\pgfpathclose%
\pgfusepath{fill}%
\end{pgfscope}%
\begin{pgfscope}%
\pgfpathrectangle{\pgfqpoint{1.250000in}{0.660000in}}{\pgfqpoint{7.750000in}{4.620000in}}%
\pgfusepath{clip}%
\pgfsetbuttcap%
\pgfsetmiterjoin%
\definecolor{currentfill}{rgb}{0.000000,0.000000,0.000000}%
\pgfsetfillcolor{currentfill}%
\pgfsetlinewidth{0.000000pt}%
\definecolor{currentstroke}{rgb}{0.000000,0.000000,0.000000}%
\pgfsetstrokecolor{currentstroke}%
\pgfsetstrokeopacity{0.000000}%
\pgfsetdash{}{0pt}%
\pgfpathmoveto{\pgfqpoint{3.432400in}{0.660000in}}%
\pgfpathlineto{\pgfqpoint{3.438600in}{0.660000in}}%
\pgfpathlineto{\pgfqpoint{3.438600in}{3.318170in}}%
\pgfpathlineto{\pgfqpoint{3.432400in}{3.318170in}}%
\pgfpathlineto{\pgfqpoint{3.432400in}{0.660000in}}%
\pgfpathclose%
\pgfusepath{fill}%
\end{pgfscope}%
\begin{pgfscope}%
\pgfpathrectangle{\pgfqpoint{1.250000in}{0.660000in}}{\pgfqpoint{7.750000in}{4.620000in}}%
\pgfusepath{clip}%
\pgfsetbuttcap%
\pgfsetmiterjoin%
\definecolor{currentfill}{rgb}{0.000000,0.000000,0.000000}%
\pgfsetfillcolor{currentfill}%
\pgfsetlinewidth{0.000000pt}%
\definecolor{currentstroke}{rgb}{0.000000,0.000000,0.000000}%
\pgfsetstrokecolor{currentstroke}%
\pgfsetstrokeopacity{0.000000}%
\pgfsetdash{}{0pt}%
\pgfpathmoveto{\pgfqpoint{3.440150in}{0.660000in}}%
\pgfpathlineto{\pgfqpoint{3.446350in}{0.660000in}}%
\pgfpathlineto{\pgfqpoint{3.446350in}{3.317723in}}%
\pgfpathlineto{\pgfqpoint{3.440150in}{3.317723in}}%
\pgfpathlineto{\pgfqpoint{3.440150in}{0.660000in}}%
\pgfpathclose%
\pgfusepath{fill}%
\end{pgfscope}%
\begin{pgfscope}%
\pgfpathrectangle{\pgfqpoint{1.250000in}{0.660000in}}{\pgfqpoint{7.750000in}{4.620000in}}%
\pgfusepath{clip}%
\pgfsetbuttcap%
\pgfsetmiterjoin%
\definecolor{currentfill}{rgb}{0.000000,0.000000,0.000000}%
\pgfsetfillcolor{currentfill}%
\pgfsetlinewidth{0.000000pt}%
\definecolor{currentstroke}{rgb}{0.000000,0.000000,0.000000}%
\pgfsetstrokecolor{currentstroke}%
\pgfsetstrokeopacity{0.000000}%
\pgfsetdash{}{0pt}%
\pgfpathmoveto{\pgfqpoint{3.447900in}{0.660000in}}%
\pgfpathlineto{\pgfqpoint{3.454100in}{0.660000in}}%
\pgfpathlineto{\pgfqpoint{3.454100in}{3.317280in}}%
\pgfpathlineto{\pgfqpoint{3.447900in}{3.317280in}}%
\pgfpathlineto{\pgfqpoint{3.447900in}{0.660000in}}%
\pgfpathclose%
\pgfusepath{fill}%
\end{pgfscope}%
\begin{pgfscope}%
\pgfpathrectangle{\pgfqpoint{1.250000in}{0.660000in}}{\pgfqpoint{7.750000in}{4.620000in}}%
\pgfusepath{clip}%
\pgfsetbuttcap%
\pgfsetmiterjoin%
\definecolor{currentfill}{rgb}{0.000000,0.000000,0.000000}%
\pgfsetfillcolor{currentfill}%
\pgfsetlinewidth{0.000000pt}%
\definecolor{currentstroke}{rgb}{0.000000,0.000000,0.000000}%
\pgfsetstrokecolor{currentstroke}%
\pgfsetstrokeopacity{0.000000}%
\pgfsetdash{}{0pt}%
\pgfpathmoveto{\pgfqpoint{3.455650in}{0.660000in}}%
\pgfpathlineto{\pgfqpoint{3.461850in}{0.660000in}}%
\pgfpathlineto{\pgfqpoint{3.461850in}{3.316841in}}%
\pgfpathlineto{\pgfqpoint{3.455650in}{3.316841in}}%
\pgfpathlineto{\pgfqpoint{3.455650in}{0.660000in}}%
\pgfpathclose%
\pgfusepath{fill}%
\end{pgfscope}%
\begin{pgfscope}%
\pgfpathrectangle{\pgfqpoint{1.250000in}{0.660000in}}{\pgfqpoint{7.750000in}{4.620000in}}%
\pgfusepath{clip}%
\pgfsetbuttcap%
\pgfsetmiterjoin%
\definecolor{currentfill}{rgb}{0.000000,0.000000,0.000000}%
\pgfsetfillcolor{currentfill}%
\pgfsetlinewidth{0.000000pt}%
\definecolor{currentstroke}{rgb}{0.000000,0.000000,0.000000}%
\pgfsetstrokecolor{currentstroke}%
\pgfsetstrokeopacity{0.000000}%
\pgfsetdash{}{0pt}%
\pgfpathmoveto{\pgfqpoint{3.463400in}{0.660000in}}%
\pgfpathlineto{\pgfqpoint{3.469600in}{0.660000in}}%
\pgfpathlineto{\pgfqpoint{3.469600in}{3.316401in}}%
\pgfpathlineto{\pgfqpoint{3.463400in}{3.316401in}}%
\pgfpathlineto{\pgfqpoint{3.463400in}{0.660000in}}%
\pgfpathclose%
\pgfusepath{fill}%
\end{pgfscope}%
\begin{pgfscope}%
\pgfpathrectangle{\pgfqpoint{1.250000in}{0.660000in}}{\pgfqpoint{7.750000in}{4.620000in}}%
\pgfusepath{clip}%
\pgfsetbuttcap%
\pgfsetmiterjoin%
\definecolor{currentfill}{rgb}{0.000000,0.000000,0.000000}%
\pgfsetfillcolor{currentfill}%
\pgfsetlinewidth{0.000000pt}%
\definecolor{currentstroke}{rgb}{0.000000,0.000000,0.000000}%
\pgfsetstrokecolor{currentstroke}%
\pgfsetstrokeopacity{0.000000}%
\pgfsetdash{}{0pt}%
\pgfpathmoveto{\pgfqpoint{3.471150in}{0.660000in}}%
\pgfpathlineto{\pgfqpoint{3.477350in}{0.660000in}}%
\pgfpathlineto{\pgfqpoint{3.477350in}{3.315962in}}%
\pgfpathlineto{\pgfqpoint{3.471150in}{3.315962in}}%
\pgfpathlineto{\pgfqpoint{3.471150in}{0.660000in}}%
\pgfpathclose%
\pgfusepath{fill}%
\end{pgfscope}%
\begin{pgfscope}%
\pgfpathrectangle{\pgfqpoint{1.250000in}{0.660000in}}{\pgfqpoint{7.750000in}{4.620000in}}%
\pgfusepath{clip}%
\pgfsetbuttcap%
\pgfsetmiterjoin%
\definecolor{currentfill}{rgb}{0.000000,0.000000,0.000000}%
\pgfsetfillcolor{currentfill}%
\pgfsetlinewidth{0.000000pt}%
\definecolor{currentstroke}{rgb}{0.000000,0.000000,0.000000}%
\pgfsetstrokecolor{currentstroke}%
\pgfsetstrokeopacity{0.000000}%
\pgfsetdash{}{0pt}%
\pgfpathmoveto{\pgfqpoint{3.478900in}{0.660000in}}%
\pgfpathlineto{\pgfqpoint{3.485100in}{0.660000in}}%
\pgfpathlineto{\pgfqpoint{3.485100in}{3.315530in}}%
\pgfpathlineto{\pgfqpoint{3.478900in}{3.315530in}}%
\pgfpathlineto{\pgfqpoint{3.478900in}{0.660000in}}%
\pgfpathclose%
\pgfusepath{fill}%
\end{pgfscope}%
\begin{pgfscope}%
\pgfpathrectangle{\pgfqpoint{1.250000in}{0.660000in}}{\pgfqpoint{7.750000in}{4.620000in}}%
\pgfusepath{clip}%
\pgfsetbuttcap%
\pgfsetmiterjoin%
\definecolor{currentfill}{rgb}{0.000000,0.000000,0.000000}%
\pgfsetfillcolor{currentfill}%
\pgfsetlinewidth{0.000000pt}%
\definecolor{currentstroke}{rgb}{0.000000,0.000000,0.000000}%
\pgfsetstrokecolor{currentstroke}%
\pgfsetstrokeopacity{0.000000}%
\pgfsetdash{}{0pt}%
\pgfpathmoveto{\pgfqpoint{3.486650in}{0.660000in}}%
\pgfpathlineto{\pgfqpoint{3.492850in}{0.660000in}}%
\pgfpathlineto{\pgfqpoint{3.492850in}{3.315098in}}%
\pgfpathlineto{\pgfqpoint{3.486650in}{3.315098in}}%
\pgfpathlineto{\pgfqpoint{3.486650in}{0.660000in}}%
\pgfpathclose%
\pgfusepath{fill}%
\end{pgfscope}%
\begin{pgfscope}%
\pgfpathrectangle{\pgfqpoint{1.250000in}{0.660000in}}{\pgfqpoint{7.750000in}{4.620000in}}%
\pgfusepath{clip}%
\pgfsetbuttcap%
\pgfsetmiterjoin%
\definecolor{currentfill}{rgb}{0.000000,0.000000,0.000000}%
\pgfsetfillcolor{currentfill}%
\pgfsetlinewidth{0.000000pt}%
\definecolor{currentstroke}{rgb}{0.000000,0.000000,0.000000}%
\pgfsetstrokecolor{currentstroke}%
\pgfsetstrokeopacity{0.000000}%
\pgfsetdash{}{0pt}%
\pgfpathmoveto{\pgfqpoint{3.494400in}{0.660000in}}%
\pgfpathlineto{\pgfqpoint{3.500600in}{0.660000in}}%
\pgfpathlineto{\pgfqpoint{3.500600in}{3.314665in}}%
\pgfpathlineto{\pgfqpoint{3.494400in}{3.314665in}}%
\pgfpathlineto{\pgfqpoint{3.494400in}{0.660000in}}%
\pgfpathclose%
\pgfusepath{fill}%
\end{pgfscope}%
\begin{pgfscope}%
\pgfpathrectangle{\pgfqpoint{1.250000in}{0.660000in}}{\pgfqpoint{7.750000in}{4.620000in}}%
\pgfusepath{clip}%
\pgfsetbuttcap%
\pgfsetmiterjoin%
\definecolor{currentfill}{rgb}{0.000000,0.000000,0.000000}%
\pgfsetfillcolor{currentfill}%
\pgfsetlinewidth{0.000000pt}%
\definecolor{currentstroke}{rgb}{0.000000,0.000000,0.000000}%
\pgfsetstrokecolor{currentstroke}%
\pgfsetstrokeopacity{0.000000}%
\pgfsetdash{}{0pt}%
\pgfpathmoveto{\pgfqpoint{3.502150in}{0.660000in}}%
\pgfpathlineto{\pgfqpoint{3.508350in}{0.660000in}}%
\pgfpathlineto{\pgfqpoint{3.508350in}{3.314237in}}%
\pgfpathlineto{\pgfqpoint{3.502150in}{3.314237in}}%
\pgfpathlineto{\pgfqpoint{3.502150in}{0.660000in}}%
\pgfpathclose%
\pgfusepath{fill}%
\end{pgfscope}%
\begin{pgfscope}%
\pgfpathrectangle{\pgfqpoint{1.250000in}{0.660000in}}{\pgfqpoint{7.750000in}{4.620000in}}%
\pgfusepath{clip}%
\pgfsetbuttcap%
\pgfsetmiterjoin%
\definecolor{currentfill}{rgb}{0.000000,0.000000,0.000000}%
\pgfsetfillcolor{currentfill}%
\pgfsetlinewidth{0.000000pt}%
\definecolor{currentstroke}{rgb}{0.000000,0.000000,0.000000}%
\pgfsetstrokecolor{currentstroke}%
\pgfsetstrokeopacity{0.000000}%
\pgfsetdash{}{0pt}%
\pgfpathmoveto{\pgfqpoint{3.509900in}{0.660000in}}%
\pgfpathlineto{\pgfqpoint{3.516100in}{0.660000in}}%
\pgfpathlineto{\pgfqpoint{3.516100in}{3.313812in}}%
\pgfpathlineto{\pgfqpoint{3.509900in}{3.313812in}}%
\pgfpathlineto{\pgfqpoint{3.509900in}{0.660000in}}%
\pgfpathclose%
\pgfusepath{fill}%
\end{pgfscope}%
\begin{pgfscope}%
\pgfpathrectangle{\pgfqpoint{1.250000in}{0.660000in}}{\pgfqpoint{7.750000in}{4.620000in}}%
\pgfusepath{clip}%
\pgfsetbuttcap%
\pgfsetmiterjoin%
\definecolor{currentfill}{rgb}{0.000000,0.000000,0.000000}%
\pgfsetfillcolor{currentfill}%
\pgfsetlinewidth{0.000000pt}%
\definecolor{currentstroke}{rgb}{0.000000,0.000000,0.000000}%
\pgfsetstrokecolor{currentstroke}%
\pgfsetstrokeopacity{0.000000}%
\pgfsetdash{}{0pt}%
\pgfpathmoveto{\pgfqpoint{3.517650in}{0.660000in}}%
\pgfpathlineto{\pgfqpoint{3.523850in}{0.660000in}}%
\pgfpathlineto{\pgfqpoint{3.523850in}{3.313387in}}%
\pgfpathlineto{\pgfqpoint{3.517650in}{3.313387in}}%
\pgfpathlineto{\pgfqpoint{3.517650in}{0.660000in}}%
\pgfpathclose%
\pgfusepath{fill}%
\end{pgfscope}%
\begin{pgfscope}%
\pgfpathrectangle{\pgfqpoint{1.250000in}{0.660000in}}{\pgfqpoint{7.750000in}{4.620000in}}%
\pgfusepath{clip}%
\pgfsetbuttcap%
\pgfsetmiterjoin%
\definecolor{currentfill}{rgb}{0.000000,0.000000,0.000000}%
\pgfsetfillcolor{currentfill}%
\pgfsetlinewidth{0.000000pt}%
\definecolor{currentstroke}{rgb}{0.000000,0.000000,0.000000}%
\pgfsetstrokecolor{currentstroke}%
\pgfsetstrokeopacity{0.000000}%
\pgfsetdash{}{0pt}%
\pgfpathmoveto{\pgfqpoint{3.525400in}{0.660000in}}%
\pgfpathlineto{\pgfqpoint{3.531600in}{0.660000in}}%
\pgfpathlineto{\pgfqpoint{3.531600in}{3.312962in}}%
\pgfpathlineto{\pgfqpoint{3.525400in}{3.312962in}}%
\pgfpathlineto{\pgfqpoint{3.525400in}{0.660000in}}%
\pgfpathclose%
\pgfusepath{fill}%
\end{pgfscope}%
\begin{pgfscope}%
\pgfpathrectangle{\pgfqpoint{1.250000in}{0.660000in}}{\pgfqpoint{7.750000in}{4.620000in}}%
\pgfusepath{clip}%
\pgfsetbuttcap%
\pgfsetmiterjoin%
\definecolor{currentfill}{rgb}{0.000000,0.000000,0.000000}%
\pgfsetfillcolor{currentfill}%
\pgfsetlinewidth{0.000000pt}%
\definecolor{currentstroke}{rgb}{0.000000,0.000000,0.000000}%
\pgfsetstrokecolor{currentstroke}%
\pgfsetstrokeopacity{0.000000}%
\pgfsetdash{}{0pt}%
\pgfpathmoveto{\pgfqpoint{3.533150in}{0.660000in}}%
\pgfpathlineto{\pgfqpoint{3.539350in}{0.660000in}}%
\pgfpathlineto{\pgfqpoint{3.539350in}{3.312544in}}%
\pgfpathlineto{\pgfqpoint{3.533150in}{3.312544in}}%
\pgfpathlineto{\pgfqpoint{3.533150in}{0.660000in}}%
\pgfpathclose%
\pgfusepath{fill}%
\end{pgfscope}%
\begin{pgfscope}%
\pgfpathrectangle{\pgfqpoint{1.250000in}{0.660000in}}{\pgfqpoint{7.750000in}{4.620000in}}%
\pgfusepath{clip}%
\pgfsetbuttcap%
\pgfsetmiterjoin%
\definecolor{currentfill}{rgb}{0.000000,0.000000,0.000000}%
\pgfsetfillcolor{currentfill}%
\pgfsetlinewidth{0.000000pt}%
\definecolor{currentstroke}{rgb}{0.000000,0.000000,0.000000}%
\pgfsetstrokecolor{currentstroke}%
\pgfsetstrokeopacity{0.000000}%
\pgfsetdash{}{0pt}%
\pgfpathmoveto{\pgfqpoint{3.540900in}{0.660000in}}%
\pgfpathlineto{\pgfqpoint{3.547100in}{0.660000in}}%
\pgfpathlineto{\pgfqpoint{3.547100in}{3.312126in}}%
\pgfpathlineto{\pgfqpoint{3.540900in}{3.312126in}}%
\pgfpathlineto{\pgfqpoint{3.540900in}{0.660000in}}%
\pgfpathclose%
\pgfusepath{fill}%
\end{pgfscope}%
\begin{pgfscope}%
\pgfpathrectangle{\pgfqpoint{1.250000in}{0.660000in}}{\pgfqpoint{7.750000in}{4.620000in}}%
\pgfusepath{clip}%
\pgfsetbuttcap%
\pgfsetmiterjoin%
\definecolor{currentfill}{rgb}{0.000000,0.000000,0.000000}%
\pgfsetfillcolor{currentfill}%
\pgfsetlinewidth{0.000000pt}%
\definecolor{currentstroke}{rgb}{0.000000,0.000000,0.000000}%
\pgfsetstrokecolor{currentstroke}%
\pgfsetstrokeopacity{0.000000}%
\pgfsetdash{}{0pt}%
\pgfpathmoveto{\pgfqpoint{3.548650in}{0.660000in}}%
\pgfpathlineto{\pgfqpoint{3.554850in}{0.660000in}}%
\pgfpathlineto{\pgfqpoint{3.554850in}{3.311708in}}%
\pgfpathlineto{\pgfqpoint{3.548650in}{3.311708in}}%
\pgfpathlineto{\pgfqpoint{3.548650in}{0.660000in}}%
\pgfpathclose%
\pgfusepath{fill}%
\end{pgfscope}%
\begin{pgfscope}%
\pgfpathrectangle{\pgfqpoint{1.250000in}{0.660000in}}{\pgfqpoint{7.750000in}{4.620000in}}%
\pgfusepath{clip}%
\pgfsetbuttcap%
\pgfsetmiterjoin%
\definecolor{currentfill}{rgb}{0.000000,0.000000,0.000000}%
\pgfsetfillcolor{currentfill}%
\pgfsetlinewidth{0.000000pt}%
\definecolor{currentstroke}{rgb}{0.000000,0.000000,0.000000}%
\pgfsetstrokecolor{currentstroke}%
\pgfsetstrokeopacity{0.000000}%
\pgfsetdash{}{0pt}%
\pgfpathmoveto{\pgfqpoint{3.556400in}{0.660000in}}%
\pgfpathlineto{\pgfqpoint{3.562600in}{0.660000in}}%
\pgfpathlineto{\pgfqpoint{3.562600in}{3.311293in}}%
\pgfpathlineto{\pgfqpoint{3.556400in}{3.311293in}}%
\pgfpathlineto{\pgfqpoint{3.556400in}{0.660000in}}%
\pgfpathclose%
\pgfusepath{fill}%
\end{pgfscope}%
\begin{pgfscope}%
\pgfpathrectangle{\pgfqpoint{1.250000in}{0.660000in}}{\pgfqpoint{7.750000in}{4.620000in}}%
\pgfusepath{clip}%
\pgfsetbuttcap%
\pgfsetmiterjoin%
\definecolor{currentfill}{rgb}{0.000000,0.000000,0.000000}%
\pgfsetfillcolor{currentfill}%
\pgfsetlinewidth{0.000000pt}%
\definecolor{currentstroke}{rgb}{0.000000,0.000000,0.000000}%
\pgfsetstrokecolor{currentstroke}%
\pgfsetstrokeopacity{0.000000}%
\pgfsetdash{}{0pt}%
\pgfpathmoveto{\pgfqpoint{3.564150in}{0.660000in}}%
\pgfpathlineto{\pgfqpoint{3.570350in}{0.660000in}}%
\pgfpathlineto{\pgfqpoint{3.570350in}{3.310882in}}%
\pgfpathlineto{\pgfqpoint{3.564150in}{3.310882in}}%
\pgfpathlineto{\pgfqpoint{3.564150in}{0.660000in}}%
\pgfpathclose%
\pgfusepath{fill}%
\end{pgfscope}%
\begin{pgfscope}%
\pgfpathrectangle{\pgfqpoint{1.250000in}{0.660000in}}{\pgfqpoint{7.750000in}{4.620000in}}%
\pgfusepath{clip}%
\pgfsetbuttcap%
\pgfsetmiterjoin%
\definecolor{currentfill}{rgb}{0.000000,0.000000,0.000000}%
\pgfsetfillcolor{currentfill}%
\pgfsetlinewidth{0.000000pt}%
\definecolor{currentstroke}{rgb}{0.000000,0.000000,0.000000}%
\pgfsetstrokecolor{currentstroke}%
\pgfsetstrokeopacity{0.000000}%
\pgfsetdash{}{0pt}%
\pgfpathmoveto{\pgfqpoint{3.571900in}{0.660000in}}%
\pgfpathlineto{\pgfqpoint{3.578100in}{0.660000in}}%
\pgfpathlineto{\pgfqpoint{3.578100in}{3.310471in}}%
\pgfpathlineto{\pgfqpoint{3.571900in}{3.310471in}}%
\pgfpathlineto{\pgfqpoint{3.571900in}{0.660000in}}%
\pgfpathclose%
\pgfusepath{fill}%
\end{pgfscope}%
\begin{pgfscope}%
\pgfpathrectangle{\pgfqpoint{1.250000in}{0.660000in}}{\pgfqpoint{7.750000in}{4.620000in}}%
\pgfusepath{clip}%
\pgfsetbuttcap%
\pgfsetmiterjoin%
\definecolor{currentfill}{rgb}{0.000000,0.000000,0.000000}%
\pgfsetfillcolor{currentfill}%
\pgfsetlinewidth{0.000000pt}%
\definecolor{currentstroke}{rgb}{0.000000,0.000000,0.000000}%
\pgfsetstrokecolor{currentstroke}%
\pgfsetstrokeopacity{0.000000}%
\pgfsetdash{}{0pt}%
\pgfpathmoveto{\pgfqpoint{3.579650in}{0.660000in}}%
\pgfpathlineto{\pgfqpoint{3.585850in}{0.660000in}}%
\pgfpathlineto{\pgfqpoint{3.585850in}{3.310060in}}%
\pgfpathlineto{\pgfqpoint{3.579650in}{3.310060in}}%
\pgfpathlineto{\pgfqpoint{3.579650in}{0.660000in}}%
\pgfpathclose%
\pgfusepath{fill}%
\end{pgfscope}%
\begin{pgfscope}%
\pgfpathrectangle{\pgfqpoint{1.250000in}{0.660000in}}{\pgfqpoint{7.750000in}{4.620000in}}%
\pgfusepath{clip}%
\pgfsetbuttcap%
\pgfsetmiterjoin%
\definecolor{currentfill}{rgb}{0.000000,0.000000,0.000000}%
\pgfsetfillcolor{currentfill}%
\pgfsetlinewidth{0.000000pt}%
\definecolor{currentstroke}{rgb}{0.000000,0.000000,0.000000}%
\pgfsetstrokecolor{currentstroke}%
\pgfsetstrokeopacity{0.000000}%
\pgfsetdash{}{0pt}%
\pgfpathmoveto{\pgfqpoint{3.587400in}{0.660000in}}%
\pgfpathlineto{\pgfqpoint{3.593600in}{0.660000in}}%
\pgfpathlineto{\pgfqpoint{3.593600in}{3.309654in}}%
\pgfpathlineto{\pgfqpoint{3.587400in}{3.309654in}}%
\pgfpathlineto{\pgfqpoint{3.587400in}{0.660000in}}%
\pgfpathclose%
\pgfusepath{fill}%
\end{pgfscope}%
\begin{pgfscope}%
\pgfpathrectangle{\pgfqpoint{1.250000in}{0.660000in}}{\pgfqpoint{7.750000in}{4.620000in}}%
\pgfusepath{clip}%
\pgfsetbuttcap%
\pgfsetmiterjoin%
\definecolor{currentfill}{rgb}{0.000000,0.000000,0.000000}%
\pgfsetfillcolor{currentfill}%
\pgfsetlinewidth{0.000000pt}%
\definecolor{currentstroke}{rgb}{0.000000,0.000000,0.000000}%
\pgfsetstrokecolor{currentstroke}%
\pgfsetstrokeopacity{0.000000}%
\pgfsetdash{}{0pt}%
\pgfpathmoveto{\pgfqpoint{3.595150in}{0.660000in}}%
\pgfpathlineto{\pgfqpoint{3.601350in}{0.660000in}}%
\pgfpathlineto{\pgfqpoint{3.601350in}{3.309250in}}%
\pgfpathlineto{\pgfqpoint{3.595150in}{3.309250in}}%
\pgfpathlineto{\pgfqpoint{3.595150in}{0.660000in}}%
\pgfpathclose%
\pgfusepath{fill}%
\end{pgfscope}%
\begin{pgfscope}%
\pgfpathrectangle{\pgfqpoint{1.250000in}{0.660000in}}{\pgfqpoint{7.750000in}{4.620000in}}%
\pgfusepath{clip}%
\pgfsetbuttcap%
\pgfsetmiterjoin%
\definecolor{currentfill}{rgb}{0.000000,0.000000,0.000000}%
\pgfsetfillcolor{currentfill}%
\pgfsetlinewidth{0.000000pt}%
\definecolor{currentstroke}{rgb}{0.000000,0.000000,0.000000}%
\pgfsetstrokecolor{currentstroke}%
\pgfsetstrokeopacity{0.000000}%
\pgfsetdash{}{0pt}%
\pgfpathmoveto{\pgfqpoint{3.602900in}{0.660000in}}%
\pgfpathlineto{\pgfqpoint{3.609100in}{0.660000in}}%
\pgfpathlineto{\pgfqpoint{3.609100in}{3.308846in}}%
\pgfpathlineto{\pgfqpoint{3.602900in}{3.308846in}}%
\pgfpathlineto{\pgfqpoint{3.602900in}{0.660000in}}%
\pgfpathclose%
\pgfusepath{fill}%
\end{pgfscope}%
\begin{pgfscope}%
\pgfpathrectangle{\pgfqpoint{1.250000in}{0.660000in}}{\pgfqpoint{7.750000in}{4.620000in}}%
\pgfusepath{clip}%
\pgfsetbuttcap%
\pgfsetmiterjoin%
\definecolor{currentfill}{rgb}{0.000000,0.000000,0.000000}%
\pgfsetfillcolor{currentfill}%
\pgfsetlinewidth{0.000000pt}%
\definecolor{currentstroke}{rgb}{0.000000,0.000000,0.000000}%
\pgfsetstrokecolor{currentstroke}%
\pgfsetstrokeopacity{0.000000}%
\pgfsetdash{}{0pt}%
\pgfpathmoveto{\pgfqpoint{3.610650in}{0.660000in}}%
\pgfpathlineto{\pgfqpoint{3.616850in}{0.660000in}}%
\pgfpathlineto{\pgfqpoint{3.616850in}{3.308443in}}%
\pgfpathlineto{\pgfqpoint{3.610650in}{3.308443in}}%
\pgfpathlineto{\pgfqpoint{3.610650in}{0.660000in}}%
\pgfpathclose%
\pgfusepath{fill}%
\end{pgfscope}%
\begin{pgfscope}%
\pgfpathrectangle{\pgfqpoint{1.250000in}{0.660000in}}{\pgfqpoint{7.750000in}{4.620000in}}%
\pgfusepath{clip}%
\pgfsetbuttcap%
\pgfsetmiterjoin%
\definecolor{currentfill}{rgb}{0.000000,0.000000,0.000000}%
\pgfsetfillcolor{currentfill}%
\pgfsetlinewidth{0.000000pt}%
\definecolor{currentstroke}{rgb}{0.000000,0.000000,0.000000}%
\pgfsetstrokecolor{currentstroke}%
\pgfsetstrokeopacity{0.000000}%
\pgfsetdash{}{0pt}%
\pgfpathmoveto{\pgfqpoint{3.618400in}{0.660000in}}%
\pgfpathlineto{\pgfqpoint{3.624600in}{0.660000in}}%
\pgfpathlineto{\pgfqpoint{3.624600in}{3.308045in}}%
\pgfpathlineto{\pgfqpoint{3.618400in}{3.308045in}}%
\pgfpathlineto{\pgfqpoint{3.618400in}{0.660000in}}%
\pgfpathclose%
\pgfusepath{fill}%
\end{pgfscope}%
\begin{pgfscope}%
\pgfpathrectangle{\pgfqpoint{1.250000in}{0.660000in}}{\pgfqpoint{7.750000in}{4.620000in}}%
\pgfusepath{clip}%
\pgfsetbuttcap%
\pgfsetmiterjoin%
\definecolor{currentfill}{rgb}{0.000000,0.000000,0.000000}%
\pgfsetfillcolor{currentfill}%
\pgfsetlinewidth{0.000000pt}%
\definecolor{currentstroke}{rgb}{0.000000,0.000000,0.000000}%
\pgfsetstrokecolor{currentstroke}%
\pgfsetstrokeopacity{0.000000}%
\pgfsetdash{}{0pt}%
\pgfpathmoveto{\pgfqpoint{3.626150in}{0.660000in}}%
\pgfpathlineto{\pgfqpoint{3.632350in}{0.660000in}}%
\pgfpathlineto{\pgfqpoint{3.632350in}{3.307648in}}%
\pgfpathlineto{\pgfqpoint{3.626150in}{3.307648in}}%
\pgfpathlineto{\pgfqpoint{3.626150in}{0.660000in}}%
\pgfpathclose%
\pgfusepath{fill}%
\end{pgfscope}%
\begin{pgfscope}%
\pgfpathrectangle{\pgfqpoint{1.250000in}{0.660000in}}{\pgfqpoint{7.750000in}{4.620000in}}%
\pgfusepath{clip}%
\pgfsetbuttcap%
\pgfsetmiterjoin%
\definecolor{currentfill}{rgb}{0.000000,0.000000,0.000000}%
\pgfsetfillcolor{currentfill}%
\pgfsetlinewidth{0.000000pt}%
\definecolor{currentstroke}{rgb}{0.000000,0.000000,0.000000}%
\pgfsetstrokecolor{currentstroke}%
\pgfsetstrokeopacity{0.000000}%
\pgfsetdash{}{0pt}%
\pgfpathmoveto{\pgfqpoint{3.633900in}{0.660000in}}%
\pgfpathlineto{\pgfqpoint{3.640100in}{0.660000in}}%
\pgfpathlineto{\pgfqpoint{3.640100in}{3.307251in}}%
\pgfpathlineto{\pgfqpoint{3.633900in}{3.307251in}}%
\pgfpathlineto{\pgfqpoint{3.633900in}{0.660000in}}%
\pgfpathclose%
\pgfusepath{fill}%
\end{pgfscope}%
\begin{pgfscope}%
\pgfpathrectangle{\pgfqpoint{1.250000in}{0.660000in}}{\pgfqpoint{7.750000in}{4.620000in}}%
\pgfusepath{clip}%
\pgfsetbuttcap%
\pgfsetmiterjoin%
\definecolor{currentfill}{rgb}{0.000000,0.000000,0.000000}%
\pgfsetfillcolor{currentfill}%
\pgfsetlinewidth{0.000000pt}%
\definecolor{currentstroke}{rgb}{0.000000,0.000000,0.000000}%
\pgfsetstrokecolor{currentstroke}%
\pgfsetstrokeopacity{0.000000}%
\pgfsetdash{}{0pt}%
\pgfpathmoveto{\pgfqpoint{3.641650in}{0.660000in}}%
\pgfpathlineto{\pgfqpoint{3.647850in}{0.660000in}}%
\pgfpathlineto{\pgfqpoint{3.647850in}{3.306854in}}%
\pgfpathlineto{\pgfqpoint{3.641650in}{3.306854in}}%
\pgfpathlineto{\pgfqpoint{3.641650in}{0.660000in}}%
\pgfpathclose%
\pgfusepath{fill}%
\end{pgfscope}%
\begin{pgfscope}%
\pgfpathrectangle{\pgfqpoint{1.250000in}{0.660000in}}{\pgfqpoint{7.750000in}{4.620000in}}%
\pgfusepath{clip}%
\pgfsetbuttcap%
\pgfsetmiterjoin%
\definecolor{currentfill}{rgb}{0.000000,0.000000,0.000000}%
\pgfsetfillcolor{currentfill}%
\pgfsetlinewidth{0.000000pt}%
\definecolor{currentstroke}{rgb}{0.000000,0.000000,0.000000}%
\pgfsetstrokecolor{currentstroke}%
\pgfsetstrokeopacity{0.000000}%
\pgfsetdash{}{0pt}%
\pgfpathmoveto{\pgfqpoint{3.649400in}{0.660000in}}%
\pgfpathlineto{\pgfqpoint{3.655600in}{0.660000in}}%
\pgfpathlineto{\pgfqpoint{3.655600in}{3.306464in}}%
\pgfpathlineto{\pgfqpoint{3.649400in}{3.306464in}}%
\pgfpathlineto{\pgfqpoint{3.649400in}{0.660000in}}%
\pgfpathclose%
\pgfusepath{fill}%
\end{pgfscope}%
\begin{pgfscope}%
\pgfpathrectangle{\pgfqpoint{1.250000in}{0.660000in}}{\pgfqpoint{7.750000in}{4.620000in}}%
\pgfusepath{clip}%
\pgfsetbuttcap%
\pgfsetmiterjoin%
\definecolor{currentfill}{rgb}{0.000000,0.000000,0.000000}%
\pgfsetfillcolor{currentfill}%
\pgfsetlinewidth{0.000000pt}%
\definecolor{currentstroke}{rgb}{0.000000,0.000000,0.000000}%
\pgfsetstrokecolor{currentstroke}%
\pgfsetstrokeopacity{0.000000}%
\pgfsetdash{}{0pt}%
\pgfpathmoveto{\pgfqpoint{3.657150in}{0.660000in}}%
\pgfpathlineto{\pgfqpoint{3.663350in}{0.660000in}}%
\pgfpathlineto{\pgfqpoint{3.663350in}{3.306074in}}%
\pgfpathlineto{\pgfqpoint{3.657150in}{3.306074in}}%
\pgfpathlineto{\pgfqpoint{3.657150in}{0.660000in}}%
\pgfpathclose%
\pgfusepath{fill}%
\end{pgfscope}%
\begin{pgfscope}%
\pgfpathrectangle{\pgfqpoint{1.250000in}{0.660000in}}{\pgfqpoint{7.750000in}{4.620000in}}%
\pgfusepath{clip}%
\pgfsetbuttcap%
\pgfsetmiterjoin%
\definecolor{currentfill}{rgb}{0.000000,0.000000,0.000000}%
\pgfsetfillcolor{currentfill}%
\pgfsetlinewidth{0.000000pt}%
\definecolor{currentstroke}{rgb}{0.000000,0.000000,0.000000}%
\pgfsetstrokecolor{currentstroke}%
\pgfsetstrokeopacity{0.000000}%
\pgfsetdash{}{0pt}%
\pgfpathmoveto{\pgfqpoint{3.664900in}{0.660000in}}%
\pgfpathlineto{\pgfqpoint{3.671100in}{0.660000in}}%
\pgfpathlineto{\pgfqpoint{3.671100in}{3.305684in}}%
\pgfpathlineto{\pgfqpoint{3.664900in}{3.305684in}}%
\pgfpathlineto{\pgfqpoint{3.664900in}{0.660000in}}%
\pgfpathclose%
\pgfusepath{fill}%
\end{pgfscope}%
\begin{pgfscope}%
\pgfpathrectangle{\pgfqpoint{1.250000in}{0.660000in}}{\pgfqpoint{7.750000in}{4.620000in}}%
\pgfusepath{clip}%
\pgfsetbuttcap%
\pgfsetmiterjoin%
\definecolor{currentfill}{rgb}{0.000000,0.000000,0.000000}%
\pgfsetfillcolor{currentfill}%
\pgfsetlinewidth{0.000000pt}%
\definecolor{currentstroke}{rgb}{0.000000,0.000000,0.000000}%
\pgfsetstrokecolor{currentstroke}%
\pgfsetstrokeopacity{0.000000}%
\pgfsetdash{}{0pt}%
\pgfpathmoveto{\pgfqpoint{3.672650in}{0.660000in}}%
\pgfpathlineto{\pgfqpoint{3.678850in}{0.660000in}}%
\pgfpathlineto{\pgfqpoint{3.678850in}{3.305294in}}%
\pgfpathlineto{\pgfqpoint{3.672650in}{3.305294in}}%
\pgfpathlineto{\pgfqpoint{3.672650in}{0.660000in}}%
\pgfpathclose%
\pgfusepath{fill}%
\end{pgfscope}%
\begin{pgfscope}%
\pgfpathrectangle{\pgfqpoint{1.250000in}{0.660000in}}{\pgfqpoint{7.750000in}{4.620000in}}%
\pgfusepath{clip}%
\pgfsetbuttcap%
\pgfsetmiterjoin%
\definecolor{currentfill}{rgb}{0.000000,0.000000,0.000000}%
\pgfsetfillcolor{currentfill}%
\pgfsetlinewidth{0.000000pt}%
\definecolor{currentstroke}{rgb}{0.000000,0.000000,0.000000}%
\pgfsetstrokecolor{currentstroke}%
\pgfsetstrokeopacity{0.000000}%
\pgfsetdash{}{0pt}%
\pgfpathmoveto{\pgfqpoint{3.680400in}{0.660000in}}%
\pgfpathlineto{\pgfqpoint{3.686600in}{0.660000in}}%
\pgfpathlineto{\pgfqpoint{3.686600in}{3.304911in}}%
\pgfpathlineto{\pgfqpoint{3.680400in}{3.304911in}}%
\pgfpathlineto{\pgfqpoint{3.680400in}{0.660000in}}%
\pgfpathclose%
\pgfusepath{fill}%
\end{pgfscope}%
\begin{pgfscope}%
\pgfpathrectangle{\pgfqpoint{1.250000in}{0.660000in}}{\pgfqpoint{7.750000in}{4.620000in}}%
\pgfusepath{clip}%
\pgfsetbuttcap%
\pgfsetmiterjoin%
\definecolor{currentfill}{rgb}{0.000000,0.000000,0.000000}%
\pgfsetfillcolor{currentfill}%
\pgfsetlinewidth{0.000000pt}%
\definecolor{currentstroke}{rgb}{0.000000,0.000000,0.000000}%
\pgfsetstrokecolor{currentstroke}%
\pgfsetstrokeopacity{0.000000}%
\pgfsetdash{}{0pt}%
\pgfpathmoveto{\pgfqpoint{3.688150in}{0.660000in}}%
\pgfpathlineto{\pgfqpoint{3.694350in}{0.660000in}}%
\pgfpathlineto{\pgfqpoint{3.694350in}{3.304528in}}%
\pgfpathlineto{\pgfqpoint{3.688150in}{3.304528in}}%
\pgfpathlineto{\pgfqpoint{3.688150in}{0.660000in}}%
\pgfpathclose%
\pgfusepath{fill}%
\end{pgfscope}%
\begin{pgfscope}%
\pgfpathrectangle{\pgfqpoint{1.250000in}{0.660000in}}{\pgfqpoint{7.750000in}{4.620000in}}%
\pgfusepath{clip}%
\pgfsetbuttcap%
\pgfsetmiterjoin%
\definecolor{currentfill}{rgb}{0.000000,0.000000,0.000000}%
\pgfsetfillcolor{currentfill}%
\pgfsetlinewidth{0.000000pt}%
\definecolor{currentstroke}{rgb}{0.000000,0.000000,0.000000}%
\pgfsetstrokecolor{currentstroke}%
\pgfsetstrokeopacity{0.000000}%
\pgfsetdash{}{0pt}%
\pgfpathmoveto{\pgfqpoint{3.695900in}{0.660000in}}%
\pgfpathlineto{\pgfqpoint{3.702100in}{0.660000in}}%
\pgfpathlineto{\pgfqpoint{3.702100in}{3.304145in}}%
\pgfpathlineto{\pgfqpoint{3.695900in}{3.304145in}}%
\pgfpathlineto{\pgfqpoint{3.695900in}{0.660000in}}%
\pgfpathclose%
\pgfusepath{fill}%
\end{pgfscope}%
\begin{pgfscope}%
\pgfpathrectangle{\pgfqpoint{1.250000in}{0.660000in}}{\pgfqpoint{7.750000in}{4.620000in}}%
\pgfusepath{clip}%
\pgfsetbuttcap%
\pgfsetmiterjoin%
\definecolor{currentfill}{rgb}{0.000000,0.000000,0.000000}%
\pgfsetfillcolor{currentfill}%
\pgfsetlinewidth{0.000000pt}%
\definecolor{currentstroke}{rgb}{0.000000,0.000000,0.000000}%
\pgfsetstrokecolor{currentstroke}%
\pgfsetstrokeopacity{0.000000}%
\pgfsetdash{}{0pt}%
\pgfpathmoveto{\pgfqpoint{3.703650in}{0.660000in}}%
\pgfpathlineto{\pgfqpoint{3.709850in}{0.660000in}}%
\pgfpathlineto{\pgfqpoint{3.709850in}{3.303761in}}%
\pgfpathlineto{\pgfqpoint{3.703650in}{3.303761in}}%
\pgfpathlineto{\pgfqpoint{3.703650in}{0.660000in}}%
\pgfpathclose%
\pgfusepath{fill}%
\end{pgfscope}%
\begin{pgfscope}%
\pgfpathrectangle{\pgfqpoint{1.250000in}{0.660000in}}{\pgfqpoint{7.750000in}{4.620000in}}%
\pgfusepath{clip}%
\pgfsetbuttcap%
\pgfsetmiterjoin%
\definecolor{currentfill}{rgb}{0.000000,0.000000,0.000000}%
\pgfsetfillcolor{currentfill}%
\pgfsetlinewidth{0.000000pt}%
\definecolor{currentstroke}{rgb}{0.000000,0.000000,0.000000}%
\pgfsetstrokecolor{currentstroke}%
\pgfsetstrokeopacity{0.000000}%
\pgfsetdash{}{0pt}%
\pgfpathmoveto{\pgfqpoint{3.711400in}{0.660000in}}%
\pgfpathlineto{\pgfqpoint{3.717600in}{0.660000in}}%
\pgfpathlineto{\pgfqpoint{3.717600in}{3.303385in}}%
\pgfpathlineto{\pgfqpoint{3.711400in}{3.303385in}}%
\pgfpathlineto{\pgfqpoint{3.711400in}{0.660000in}}%
\pgfpathclose%
\pgfusepath{fill}%
\end{pgfscope}%
\begin{pgfscope}%
\pgfpathrectangle{\pgfqpoint{1.250000in}{0.660000in}}{\pgfqpoint{7.750000in}{4.620000in}}%
\pgfusepath{clip}%
\pgfsetbuttcap%
\pgfsetmiterjoin%
\definecolor{currentfill}{rgb}{0.000000,0.000000,0.000000}%
\pgfsetfillcolor{currentfill}%
\pgfsetlinewidth{0.000000pt}%
\definecolor{currentstroke}{rgb}{0.000000,0.000000,0.000000}%
\pgfsetstrokecolor{currentstroke}%
\pgfsetstrokeopacity{0.000000}%
\pgfsetdash{}{0pt}%
\pgfpathmoveto{\pgfqpoint{3.719150in}{0.660000in}}%
\pgfpathlineto{\pgfqpoint{3.725350in}{0.660000in}}%
\pgfpathlineto{\pgfqpoint{3.725350in}{3.303008in}}%
\pgfpathlineto{\pgfqpoint{3.719150in}{3.303008in}}%
\pgfpathlineto{\pgfqpoint{3.719150in}{0.660000in}}%
\pgfpathclose%
\pgfusepath{fill}%
\end{pgfscope}%
\begin{pgfscope}%
\pgfpathrectangle{\pgfqpoint{1.250000in}{0.660000in}}{\pgfqpoint{7.750000in}{4.620000in}}%
\pgfusepath{clip}%
\pgfsetbuttcap%
\pgfsetmiterjoin%
\definecolor{currentfill}{rgb}{0.000000,0.000000,0.000000}%
\pgfsetfillcolor{currentfill}%
\pgfsetlinewidth{0.000000pt}%
\definecolor{currentstroke}{rgb}{0.000000,0.000000,0.000000}%
\pgfsetstrokecolor{currentstroke}%
\pgfsetstrokeopacity{0.000000}%
\pgfsetdash{}{0pt}%
\pgfpathmoveto{\pgfqpoint{3.726900in}{0.660000in}}%
\pgfpathlineto{\pgfqpoint{3.733100in}{0.660000in}}%
\pgfpathlineto{\pgfqpoint{3.733100in}{3.302632in}}%
\pgfpathlineto{\pgfqpoint{3.726900in}{3.302632in}}%
\pgfpathlineto{\pgfqpoint{3.726900in}{0.660000in}}%
\pgfpathclose%
\pgfusepath{fill}%
\end{pgfscope}%
\begin{pgfscope}%
\pgfpathrectangle{\pgfqpoint{1.250000in}{0.660000in}}{\pgfqpoint{7.750000in}{4.620000in}}%
\pgfusepath{clip}%
\pgfsetbuttcap%
\pgfsetmiterjoin%
\definecolor{currentfill}{rgb}{0.000000,0.000000,0.000000}%
\pgfsetfillcolor{currentfill}%
\pgfsetlinewidth{0.000000pt}%
\definecolor{currentstroke}{rgb}{0.000000,0.000000,0.000000}%
\pgfsetstrokecolor{currentstroke}%
\pgfsetstrokeopacity{0.000000}%
\pgfsetdash{}{0pt}%
\pgfpathmoveto{\pgfqpoint{3.734650in}{0.660000in}}%
\pgfpathlineto{\pgfqpoint{3.740850in}{0.660000in}}%
\pgfpathlineto{\pgfqpoint{3.740850in}{3.302256in}}%
\pgfpathlineto{\pgfqpoint{3.734650in}{3.302256in}}%
\pgfpathlineto{\pgfqpoint{3.734650in}{0.660000in}}%
\pgfpathclose%
\pgfusepath{fill}%
\end{pgfscope}%
\begin{pgfscope}%
\pgfpathrectangle{\pgfqpoint{1.250000in}{0.660000in}}{\pgfqpoint{7.750000in}{4.620000in}}%
\pgfusepath{clip}%
\pgfsetbuttcap%
\pgfsetmiterjoin%
\definecolor{currentfill}{rgb}{0.000000,0.000000,0.000000}%
\pgfsetfillcolor{currentfill}%
\pgfsetlinewidth{0.000000pt}%
\definecolor{currentstroke}{rgb}{0.000000,0.000000,0.000000}%
\pgfsetstrokecolor{currentstroke}%
\pgfsetstrokeopacity{0.000000}%
\pgfsetdash{}{0pt}%
\pgfpathmoveto{\pgfqpoint{3.742400in}{0.660000in}}%
\pgfpathlineto{\pgfqpoint{3.748600in}{0.660000in}}%
\pgfpathlineto{\pgfqpoint{3.748600in}{3.301885in}}%
\pgfpathlineto{\pgfqpoint{3.742400in}{3.301885in}}%
\pgfpathlineto{\pgfqpoint{3.742400in}{0.660000in}}%
\pgfpathclose%
\pgfusepath{fill}%
\end{pgfscope}%
\begin{pgfscope}%
\pgfpathrectangle{\pgfqpoint{1.250000in}{0.660000in}}{\pgfqpoint{7.750000in}{4.620000in}}%
\pgfusepath{clip}%
\pgfsetbuttcap%
\pgfsetmiterjoin%
\definecolor{currentfill}{rgb}{0.000000,0.000000,0.000000}%
\pgfsetfillcolor{currentfill}%
\pgfsetlinewidth{0.000000pt}%
\definecolor{currentstroke}{rgb}{0.000000,0.000000,0.000000}%
\pgfsetstrokecolor{currentstroke}%
\pgfsetstrokeopacity{0.000000}%
\pgfsetdash{}{0pt}%
\pgfpathmoveto{\pgfqpoint{3.750150in}{0.660000in}}%
\pgfpathlineto{\pgfqpoint{3.756350in}{0.660000in}}%
\pgfpathlineto{\pgfqpoint{3.756350in}{3.301515in}}%
\pgfpathlineto{\pgfqpoint{3.750150in}{3.301515in}}%
\pgfpathlineto{\pgfqpoint{3.750150in}{0.660000in}}%
\pgfpathclose%
\pgfusepath{fill}%
\end{pgfscope}%
\begin{pgfscope}%
\pgfpathrectangle{\pgfqpoint{1.250000in}{0.660000in}}{\pgfqpoint{7.750000in}{4.620000in}}%
\pgfusepath{clip}%
\pgfsetbuttcap%
\pgfsetmiterjoin%
\definecolor{currentfill}{rgb}{0.000000,0.000000,0.000000}%
\pgfsetfillcolor{currentfill}%
\pgfsetlinewidth{0.000000pt}%
\definecolor{currentstroke}{rgb}{0.000000,0.000000,0.000000}%
\pgfsetstrokecolor{currentstroke}%
\pgfsetstrokeopacity{0.000000}%
\pgfsetdash{}{0pt}%
\pgfpathmoveto{\pgfqpoint{3.757900in}{0.660000in}}%
\pgfpathlineto{\pgfqpoint{3.764100in}{0.660000in}}%
\pgfpathlineto{\pgfqpoint{3.764100in}{3.301145in}}%
\pgfpathlineto{\pgfqpoint{3.757900in}{3.301145in}}%
\pgfpathlineto{\pgfqpoint{3.757900in}{0.660000in}}%
\pgfpathclose%
\pgfusepath{fill}%
\end{pgfscope}%
\begin{pgfscope}%
\pgfpathrectangle{\pgfqpoint{1.250000in}{0.660000in}}{\pgfqpoint{7.750000in}{4.620000in}}%
\pgfusepath{clip}%
\pgfsetbuttcap%
\pgfsetmiterjoin%
\definecolor{currentfill}{rgb}{0.000000,0.000000,0.000000}%
\pgfsetfillcolor{currentfill}%
\pgfsetlinewidth{0.000000pt}%
\definecolor{currentstroke}{rgb}{0.000000,0.000000,0.000000}%
\pgfsetstrokecolor{currentstroke}%
\pgfsetstrokeopacity{0.000000}%
\pgfsetdash{}{0pt}%
\pgfpathmoveto{\pgfqpoint{3.765650in}{0.660000in}}%
\pgfpathlineto{\pgfqpoint{3.771850in}{0.660000in}}%
\pgfpathlineto{\pgfqpoint{3.771850in}{3.300776in}}%
\pgfpathlineto{\pgfqpoint{3.765650in}{3.300776in}}%
\pgfpathlineto{\pgfqpoint{3.765650in}{0.660000in}}%
\pgfpathclose%
\pgfusepath{fill}%
\end{pgfscope}%
\begin{pgfscope}%
\pgfpathrectangle{\pgfqpoint{1.250000in}{0.660000in}}{\pgfqpoint{7.750000in}{4.620000in}}%
\pgfusepath{clip}%
\pgfsetbuttcap%
\pgfsetmiterjoin%
\definecolor{currentfill}{rgb}{0.000000,0.000000,0.000000}%
\pgfsetfillcolor{currentfill}%
\pgfsetlinewidth{0.000000pt}%
\definecolor{currentstroke}{rgb}{0.000000,0.000000,0.000000}%
\pgfsetstrokecolor{currentstroke}%
\pgfsetstrokeopacity{0.000000}%
\pgfsetdash{}{0pt}%
\pgfpathmoveto{\pgfqpoint{3.773400in}{0.660000in}}%
\pgfpathlineto{\pgfqpoint{3.779600in}{0.660000in}}%
\pgfpathlineto{\pgfqpoint{3.779600in}{3.300410in}}%
\pgfpathlineto{\pgfqpoint{3.773400in}{3.300410in}}%
\pgfpathlineto{\pgfqpoint{3.773400in}{0.660000in}}%
\pgfpathclose%
\pgfusepath{fill}%
\end{pgfscope}%
\begin{pgfscope}%
\pgfpathrectangle{\pgfqpoint{1.250000in}{0.660000in}}{\pgfqpoint{7.750000in}{4.620000in}}%
\pgfusepath{clip}%
\pgfsetbuttcap%
\pgfsetmiterjoin%
\definecolor{currentfill}{rgb}{0.000000,0.000000,0.000000}%
\pgfsetfillcolor{currentfill}%
\pgfsetlinewidth{0.000000pt}%
\definecolor{currentstroke}{rgb}{0.000000,0.000000,0.000000}%
\pgfsetstrokecolor{currentstroke}%
\pgfsetstrokeopacity{0.000000}%
\pgfsetdash{}{0pt}%
\pgfpathmoveto{\pgfqpoint{3.781150in}{0.660000in}}%
\pgfpathlineto{\pgfqpoint{3.787350in}{0.660000in}}%
\pgfpathlineto{\pgfqpoint{3.787350in}{3.300047in}}%
\pgfpathlineto{\pgfqpoint{3.781150in}{3.300047in}}%
\pgfpathlineto{\pgfqpoint{3.781150in}{0.660000in}}%
\pgfpathclose%
\pgfusepath{fill}%
\end{pgfscope}%
\begin{pgfscope}%
\pgfpathrectangle{\pgfqpoint{1.250000in}{0.660000in}}{\pgfqpoint{7.750000in}{4.620000in}}%
\pgfusepath{clip}%
\pgfsetbuttcap%
\pgfsetmiterjoin%
\definecolor{currentfill}{rgb}{0.000000,0.000000,0.000000}%
\pgfsetfillcolor{currentfill}%
\pgfsetlinewidth{0.000000pt}%
\definecolor{currentstroke}{rgb}{0.000000,0.000000,0.000000}%
\pgfsetstrokecolor{currentstroke}%
\pgfsetstrokeopacity{0.000000}%
\pgfsetdash{}{0pt}%
\pgfpathmoveto{\pgfqpoint{3.788900in}{0.660000in}}%
\pgfpathlineto{\pgfqpoint{3.795100in}{0.660000in}}%
\pgfpathlineto{\pgfqpoint{3.795100in}{3.299684in}}%
\pgfpathlineto{\pgfqpoint{3.788900in}{3.299684in}}%
\pgfpathlineto{\pgfqpoint{3.788900in}{0.660000in}}%
\pgfpathclose%
\pgfusepath{fill}%
\end{pgfscope}%
\begin{pgfscope}%
\pgfpathrectangle{\pgfqpoint{1.250000in}{0.660000in}}{\pgfqpoint{7.750000in}{4.620000in}}%
\pgfusepath{clip}%
\pgfsetbuttcap%
\pgfsetmiterjoin%
\definecolor{currentfill}{rgb}{0.000000,0.000000,0.000000}%
\pgfsetfillcolor{currentfill}%
\pgfsetlinewidth{0.000000pt}%
\definecolor{currentstroke}{rgb}{0.000000,0.000000,0.000000}%
\pgfsetstrokecolor{currentstroke}%
\pgfsetstrokeopacity{0.000000}%
\pgfsetdash{}{0pt}%
\pgfpathmoveto{\pgfqpoint{3.796650in}{0.660000in}}%
\pgfpathlineto{\pgfqpoint{3.802850in}{0.660000in}}%
\pgfpathlineto{\pgfqpoint{3.802850in}{3.299321in}}%
\pgfpathlineto{\pgfqpoint{3.796650in}{3.299321in}}%
\pgfpathlineto{\pgfqpoint{3.796650in}{0.660000in}}%
\pgfpathclose%
\pgfusepath{fill}%
\end{pgfscope}%
\begin{pgfscope}%
\pgfpathrectangle{\pgfqpoint{1.250000in}{0.660000in}}{\pgfqpoint{7.750000in}{4.620000in}}%
\pgfusepath{clip}%
\pgfsetbuttcap%
\pgfsetmiterjoin%
\definecolor{currentfill}{rgb}{0.000000,0.000000,0.000000}%
\pgfsetfillcolor{currentfill}%
\pgfsetlinewidth{0.000000pt}%
\definecolor{currentstroke}{rgb}{0.000000,0.000000,0.000000}%
\pgfsetstrokecolor{currentstroke}%
\pgfsetstrokeopacity{0.000000}%
\pgfsetdash{}{0pt}%
\pgfpathmoveto{\pgfqpoint{3.804400in}{0.660000in}}%
\pgfpathlineto{\pgfqpoint{3.810600in}{0.660000in}}%
\pgfpathlineto{\pgfqpoint{3.810600in}{3.298960in}}%
\pgfpathlineto{\pgfqpoint{3.804400in}{3.298960in}}%
\pgfpathlineto{\pgfqpoint{3.804400in}{0.660000in}}%
\pgfpathclose%
\pgfusepath{fill}%
\end{pgfscope}%
\begin{pgfscope}%
\pgfpathrectangle{\pgfqpoint{1.250000in}{0.660000in}}{\pgfqpoint{7.750000in}{4.620000in}}%
\pgfusepath{clip}%
\pgfsetbuttcap%
\pgfsetmiterjoin%
\definecolor{currentfill}{rgb}{0.000000,0.000000,0.000000}%
\pgfsetfillcolor{currentfill}%
\pgfsetlinewidth{0.000000pt}%
\definecolor{currentstroke}{rgb}{0.000000,0.000000,0.000000}%
\pgfsetstrokecolor{currentstroke}%
\pgfsetstrokeopacity{0.000000}%
\pgfsetdash{}{0pt}%
\pgfpathmoveto{\pgfqpoint{3.812150in}{0.660000in}}%
\pgfpathlineto{\pgfqpoint{3.818350in}{0.660000in}}%
\pgfpathlineto{\pgfqpoint{3.818350in}{3.298603in}}%
\pgfpathlineto{\pgfqpoint{3.812150in}{3.298603in}}%
\pgfpathlineto{\pgfqpoint{3.812150in}{0.660000in}}%
\pgfpathclose%
\pgfusepath{fill}%
\end{pgfscope}%
\begin{pgfscope}%
\pgfpathrectangle{\pgfqpoint{1.250000in}{0.660000in}}{\pgfqpoint{7.750000in}{4.620000in}}%
\pgfusepath{clip}%
\pgfsetbuttcap%
\pgfsetmiterjoin%
\definecolor{currentfill}{rgb}{0.000000,0.000000,0.000000}%
\pgfsetfillcolor{currentfill}%
\pgfsetlinewidth{0.000000pt}%
\definecolor{currentstroke}{rgb}{0.000000,0.000000,0.000000}%
\pgfsetstrokecolor{currentstroke}%
\pgfsetstrokeopacity{0.000000}%
\pgfsetdash{}{0pt}%
\pgfpathmoveto{\pgfqpoint{3.819900in}{0.660000in}}%
\pgfpathlineto{\pgfqpoint{3.826100in}{0.660000in}}%
\pgfpathlineto{\pgfqpoint{3.826100in}{3.298247in}}%
\pgfpathlineto{\pgfqpoint{3.819900in}{3.298247in}}%
\pgfpathlineto{\pgfqpoint{3.819900in}{0.660000in}}%
\pgfpathclose%
\pgfusepath{fill}%
\end{pgfscope}%
\begin{pgfscope}%
\pgfpathrectangle{\pgfqpoint{1.250000in}{0.660000in}}{\pgfqpoint{7.750000in}{4.620000in}}%
\pgfusepath{clip}%
\pgfsetbuttcap%
\pgfsetmiterjoin%
\definecolor{currentfill}{rgb}{0.000000,0.000000,0.000000}%
\pgfsetfillcolor{currentfill}%
\pgfsetlinewidth{0.000000pt}%
\definecolor{currentstroke}{rgb}{0.000000,0.000000,0.000000}%
\pgfsetstrokecolor{currentstroke}%
\pgfsetstrokeopacity{0.000000}%
\pgfsetdash{}{0pt}%
\pgfpathmoveto{\pgfqpoint{3.827650in}{0.660000in}}%
\pgfpathlineto{\pgfqpoint{3.833850in}{0.660000in}}%
\pgfpathlineto{\pgfqpoint{3.833850in}{3.297890in}}%
\pgfpathlineto{\pgfqpoint{3.827650in}{3.297890in}}%
\pgfpathlineto{\pgfqpoint{3.827650in}{0.660000in}}%
\pgfpathclose%
\pgfusepath{fill}%
\end{pgfscope}%
\begin{pgfscope}%
\pgfpathrectangle{\pgfqpoint{1.250000in}{0.660000in}}{\pgfqpoint{7.750000in}{4.620000in}}%
\pgfusepath{clip}%
\pgfsetbuttcap%
\pgfsetmiterjoin%
\definecolor{currentfill}{rgb}{0.000000,0.000000,0.000000}%
\pgfsetfillcolor{currentfill}%
\pgfsetlinewidth{0.000000pt}%
\definecolor{currentstroke}{rgb}{0.000000,0.000000,0.000000}%
\pgfsetstrokecolor{currentstroke}%
\pgfsetstrokeopacity{0.000000}%
\pgfsetdash{}{0pt}%
\pgfpathmoveto{\pgfqpoint{3.835400in}{0.660000in}}%
\pgfpathlineto{\pgfqpoint{3.841600in}{0.660000in}}%
\pgfpathlineto{\pgfqpoint{3.841600in}{3.297534in}}%
\pgfpathlineto{\pgfqpoint{3.835400in}{3.297534in}}%
\pgfpathlineto{\pgfqpoint{3.835400in}{0.660000in}}%
\pgfpathclose%
\pgfusepath{fill}%
\end{pgfscope}%
\begin{pgfscope}%
\pgfpathrectangle{\pgfqpoint{1.250000in}{0.660000in}}{\pgfqpoint{7.750000in}{4.620000in}}%
\pgfusepath{clip}%
\pgfsetbuttcap%
\pgfsetmiterjoin%
\definecolor{currentfill}{rgb}{0.000000,0.000000,0.000000}%
\pgfsetfillcolor{currentfill}%
\pgfsetlinewidth{0.000000pt}%
\definecolor{currentstroke}{rgb}{0.000000,0.000000,0.000000}%
\pgfsetstrokecolor{currentstroke}%
\pgfsetstrokeopacity{0.000000}%
\pgfsetdash{}{0pt}%
\pgfpathmoveto{\pgfqpoint{3.843150in}{0.660000in}}%
\pgfpathlineto{\pgfqpoint{3.849350in}{0.660000in}}%
\pgfpathlineto{\pgfqpoint{3.849350in}{3.297183in}}%
\pgfpathlineto{\pgfqpoint{3.843150in}{3.297183in}}%
\pgfpathlineto{\pgfqpoint{3.843150in}{0.660000in}}%
\pgfpathclose%
\pgfusepath{fill}%
\end{pgfscope}%
\begin{pgfscope}%
\pgfpathrectangle{\pgfqpoint{1.250000in}{0.660000in}}{\pgfqpoint{7.750000in}{4.620000in}}%
\pgfusepath{clip}%
\pgfsetbuttcap%
\pgfsetmiterjoin%
\definecolor{currentfill}{rgb}{0.000000,0.000000,0.000000}%
\pgfsetfillcolor{currentfill}%
\pgfsetlinewidth{0.000000pt}%
\definecolor{currentstroke}{rgb}{0.000000,0.000000,0.000000}%
\pgfsetstrokecolor{currentstroke}%
\pgfsetstrokeopacity{0.000000}%
\pgfsetdash{}{0pt}%
\pgfpathmoveto{\pgfqpoint{3.850900in}{0.660000in}}%
\pgfpathlineto{\pgfqpoint{3.857100in}{0.660000in}}%
\pgfpathlineto{\pgfqpoint{3.857100in}{3.296833in}}%
\pgfpathlineto{\pgfqpoint{3.850900in}{3.296833in}}%
\pgfpathlineto{\pgfqpoint{3.850900in}{0.660000in}}%
\pgfpathclose%
\pgfusepath{fill}%
\end{pgfscope}%
\begin{pgfscope}%
\pgfpathrectangle{\pgfqpoint{1.250000in}{0.660000in}}{\pgfqpoint{7.750000in}{4.620000in}}%
\pgfusepath{clip}%
\pgfsetbuttcap%
\pgfsetmiterjoin%
\definecolor{currentfill}{rgb}{0.000000,0.000000,0.000000}%
\pgfsetfillcolor{currentfill}%
\pgfsetlinewidth{0.000000pt}%
\definecolor{currentstroke}{rgb}{0.000000,0.000000,0.000000}%
\pgfsetstrokecolor{currentstroke}%
\pgfsetstrokeopacity{0.000000}%
\pgfsetdash{}{0pt}%
\pgfpathmoveto{\pgfqpoint{3.858650in}{0.660000in}}%
\pgfpathlineto{\pgfqpoint{3.864850in}{0.660000in}}%
\pgfpathlineto{\pgfqpoint{3.864850in}{3.296483in}}%
\pgfpathlineto{\pgfqpoint{3.858650in}{3.296483in}}%
\pgfpathlineto{\pgfqpoint{3.858650in}{0.660000in}}%
\pgfpathclose%
\pgfusepath{fill}%
\end{pgfscope}%
\begin{pgfscope}%
\pgfpathrectangle{\pgfqpoint{1.250000in}{0.660000in}}{\pgfqpoint{7.750000in}{4.620000in}}%
\pgfusepath{clip}%
\pgfsetbuttcap%
\pgfsetmiterjoin%
\definecolor{currentfill}{rgb}{0.000000,0.000000,0.000000}%
\pgfsetfillcolor{currentfill}%
\pgfsetlinewidth{0.000000pt}%
\definecolor{currentstroke}{rgb}{0.000000,0.000000,0.000000}%
\pgfsetstrokecolor{currentstroke}%
\pgfsetstrokeopacity{0.000000}%
\pgfsetdash{}{0pt}%
\pgfpathmoveto{\pgfqpoint{3.866400in}{0.660000in}}%
\pgfpathlineto{\pgfqpoint{3.872600in}{0.660000in}}%
\pgfpathlineto{\pgfqpoint{3.872600in}{3.296133in}}%
\pgfpathlineto{\pgfqpoint{3.866400in}{3.296133in}}%
\pgfpathlineto{\pgfqpoint{3.866400in}{0.660000in}}%
\pgfpathclose%
\pgfusepath{fill}%
\end{pgfscope}%
\begin{pgfscope}%
\pgfpathrectangle{\pgfqpoint{1.250000in}{0.660000in}}{\pgfqpoint{7.750000in}{4.620000in}}%
\pgfusepath{clip}%
\pgfsetbuttcap%
\pgfsetmiterjoin%
\definecolor{currentfill}{rgb}{0.000000,0.000000,0.000000}%
\pgfsetfillcolor{currentfill}%
\pgfsetlinewidth{0.000000pt}%
\definecolor{currentstroke}{rgb}{0.000000,0.000000,0.000000}%
\pgfsetstrokecolor{currentstroke}%
\pgfsetstrokeopacity{0.000000}%
\pgfsetdash{}{0pt}%
\pgfpathmoveto{\pgfqpoint{3.874150in}{0.660000in}}%
\pgfpathlineto{\pgfqpoint{3.880350in}{0.660000in}}%
\pgfpathlineto{\pgfqpoint{3.880350in}{3.295786in}}%
\pgfpathlineto{\pgfqpoint{3.874150in}{3.295786in}}%
\pgfpathlineto{\pgfqpoint{3.874150in}{0.660000in}}%
\pgfpathclose%
\pgfusepath{fill}%
\end{pgfscope}%
\begin{pgfscope}%
\pgfpathrectangle{\pgfqpoint{1.250000in}{0.660000in}}{\pgfqpoint{7.750000in}{4.620000in}}%
\pgfusepath{clip}%
\pgfsetbuttcap%
\pgfsetmiterjoin%
\definecolor{currentfill}{rgb}{0.000000,0.000000,0.000000}%
\pgfsetfillcolor{currentfill}%
\pgfsetlinewidth{0.000000pt}%
\definecolor{currentstroke}{rgb}{0.000000,0.000000,0.000000}%
\pgfsetstrokecolor{currentstroke}%
\pgfsetstrokeopacity{0.000000}%
\pgfsetdash{}{0pt}%
\pgfpathmoveto{\pgfqpoint{3.881900in}{0.660000in}}%
\pgfpathlineto{\pgfqpoint{3.888100in}{0.660000in}}%
\pgfpathlineto{\pgfqpoint{3.888100in}{3.295442in}}%
\pgfpathlineto{\pgfqpoint{3.881900in}{3.295442in}}%
\pgfpathlineto{\pgfqpoint{3.881900in}{0.660000in}}%
\pgfpathclose%
\pgfusepath{fill}%
\end{pgfscope}%
\begin{pgfscope}%
\pgfpathrectangle{\pgfqpoint{1.250000in}{0.660000in}}{\pgfqpoint{7.750000in}{4.620000in}}%
\pgfusepath{clip}%
\pgfsetbuttcap%
\pgfsetmiterjoin%
\definecolor{currentfill}{rgb}{0.000000,0.000000,0.000000}%
\pgfsetfillcolor{currentfill}%
\pgfsetlinewidth{0.000000pt}%
\definecolor{currentstroke}{rgb}{0.000000,0.000000,0.000000}%
\pgfsetstrokecolor{currentstroke}%
\pgfsetstrokeopacity{0.000000}%
\pgfsetdash{}{0pt}%
\pgfpathmoveto{\pgfqpoint{3.889650in}{0.660000in}}%
\pgfpathlineto{\pgfqpoint{3.895850in}{0.660000in}}%
\pgfpathlineto{\pgfqpoint{3.895850in}{3.295099in}}%
\pgfpathlineto{\pgfqpoint{3.889650in}{3.295099in}}%
\pgfpathlineto{\pgfqpoint{3.889650in}{0.660000in}}%
\pgfpathclose%
\pgfusepath{fill}%
\end{pgfscope}%
\begin{pgfscope}%
\pgfpathrectangle{\pgfqpoint{1.250000in}{0.660000in}}{\pgfqpoint{7.750000in}{4.620000in}}%
\pgfusepath{clip}%
\pgfsetbuttcap%
\pgfsetmiterjoin%
\definecolor{currentfill}{rgb}{0.000000,0.000000,0.000000}%
\pgfsetfillcolor{currentfill}%
\pgfsetlinewidth{0.000000pt}%
\definecolor{currentstroke}{rgb}{0.000000,0.000000,0.000000}%
\pgfsetstrokecolor{currentstroke}%
\pgfsetstrokeopacity{0.000000}%
\pgfsetdash{}{0pt}%
\pgfpathmoveto{\pgfqpoint{3.897400in}{0.660000in}}%
\pgfpathlineto{\pgfqpoint{3.903600in}{0.660000in}}%
\pgfpathlineto{\pgfqpoint{3.903600in}{3.294755in}}%
\pgfpathlineto{\pgfqpoint{3.897400in}{3.294755in}}%
\pgfpathlineto{\pgfqpoint{3.897400in}{0.660000in}}%
\pgfpathclose%
\pgfusepath{fill}%
\end{pgfscope}%
\begin{pgfscope}%
\pgfpathrectangle{\pgfqpoint{1.250000in}{0.660000in}}{\pgfqpoint{7.750000in}{4.620000in}}%
\pgfusepath{clip}%
\pgfsetbuttcap%
\pgfsetmiterjoin%
\definecolor{currentfill}{rgb}{0.000000,0.000000,0.000000}%
\pgfsetfillcolor{currentfill}%
\pgfsetlinewidth{0.000000pt}%
\definecolor{currentstroke}{rgb}{0.000000,0.000000,0.000000}%
\pgfsetstrokecolor{currentstroke}%
\pgfsetstrokeopacity{0.000000}%
\pgfsetdash{}{0pt}%
\pgfpathmoveto{\pgfqpoint{3.905150in}{0.660000in}}%
\pgfpathlineto{\pgfqpoint{3.911350in}{0.660000in}}%
\pgfpathlineto{\pgfqpoint{3.911350in}{3.294411in}}%
\pgfpathlineto{\pgfqpoint{3.905150in}{3.294411in}}%
\pgfpathlineto{\pgfqpoint{3.905150in}{0.660000in}}%
\pgfpathclose%
\pgfusepath{fill}%
\end{pgfscope}%
\begin{pgfscope}%
\pgfpathrectangle{\pgfqpoint{1.250000in}{0.660000in}}{\pgfqpoint{7.750000in}{4.620000in}}%
\pgfusepath{clip}%
\pgfsetbuttcap%
\pgfsetmiterjoin%
\definecolor{currentfill}{rgb}{0.000000,0.000000,0.000000}%
\pgfsetfillcolor{currentfill}%
\pgfsetlinewidth{0.000000pt}%
\definecolor{currentstroke}{rgb}{0.000000,0.000000,0.000000}%
\pgfsetstrokecolor{currentstroke}%
\pgfsetstrokeopacity{0.000000}%
\pgfsetdash{}{0pt}%
\pgfpathmoveto{\pgfqpoint{3.912900in}{0.660000in}}%
\pgfpathlineto{\pgfqpoint{3.919100in}{0.660000in}}%
\pgfpathlineto{\pgfqpoint{3.919100in}{3.294073in}}%
\pgfpathlineto{\pgfqpoint{3.912900in}{3.294073in}}%
\pgfpathlineto{\pgfqpoint{3.912900in}{0.660000in}}%
\pgfpathclose%
\pgfusepath{fill}%
\end{pgfscope}%
\begin{pgfscope}%
\pgfpathrectangle{\pgfqpoint{1.250000in}{0.660000in}}{\pgfqpoint{7.750000in}{4.620000in}}%
\pgfusepath{clip}%
\pgfsetbuttcap%
\pgfsetmiterjoin%
\definecolor{currentfill}{rgb}{0.000000,0.000000,0.000000}%
\pgfsetfillcolor{currentfill}%
\pgfsetlinewidth{0.000000pt}%
\definecolor{currentstroke}{rgb}{0.000000,0.000000,0.000000}%
\pgfsetstrokecolor{currentstroke}%
\pgfsetstrokeopacity{0.000000}%
\pgfsetdash{}{0pt}%
\pgfpathmoveto{\pgfqpoint{3.920650in}{0.660000in}}%
\pgfpathlineto{\pgfqpoint{3.926850in}{0.660000in}}%
\pgfpathlineto{\pgfqpoint{3.926850in}{3.293736in}}%
\pgfpathlineto{\pgfqpoint{3.920650in}{3.293736in}}%
\pgfpathlineto{\pgfqpoint{3.920650in}{0.660000in}}%
\pgfpathclose%
\pgfusepath{fill}%
\end{pgfscope}%
\begin{pgfscope}%
\pgfpathrectangle{\pgfqpoint{1.250000in}{0.660000in}}{\pgfqpoint{7.750000in}{4.620000in}}%
\pgfusepath{clip}%
\pgfsetbuttcap%
\pgfsetmiterjoin%
\definecolor{currentfill}{rgb}{0.000000,0.000000,0.000000}%
\pgfsetfillcolor{currentfill}%
\pgfsetlinewidth{0.000000pt}%
\definecolor{currentstroke}{rgb}{0.000000,0.000000,0.000000}%
\pgfsetstrokecolor{currentstroke}%
\pgfsetstrokeopacity{0.000000}%
\pgfsetdash{}{0pt}%
\pgfpathmoveto{\pgfqpoint{3.928400in}{0.660000in}}%
\pgfpathlineto{\pgfqpoint{3.934600in}{0.660000in}}%
\pgfpathlineto{\pgfqpoint{3.934600in}{3.293399in}}%
\pgfpathlineto{\pgfqpoint{3.928400in}{3.293399in}}%
\pgfpathlineto{\pgfqpoint{3.928400in}{0.660000in}}%
\pgfpathclose%
\pgfusepath{fill}%
\end{pgfscope}%
\begin{pgfscope}%
\pgfpathrectangle{\pgfqpoint{1.250000in}{0.660000in}}{\pgfqpoint{7.750000in}{4.620000in}}%
\pgfusepath{clip}%
\pgfsetbuttcap%
\pgfsetmiterjoin%
\definecolor{currentfill}{rgb}{0.000000,0.000000,0.000000}%
\pgfsetfillcolor{currentfill}%
\pgfsetlinewidth{0.000000pt}%
\definecolor{currentstroke}{rgb}{0.000000,0.000000,0.000000}%
\pgfsetstrokecolor{currentstroke}%
\pgfsetstrokeopacity{0.000000}%
\pgfsetdash{}{0pt}%
\pgfpathmoveto{\pgfqpoint{3.936150in}{0.660000in}}%
\pgfpathlineto{\pgfqpoint{3.942350in}{0.660000in}}%
\pgfpathlineto{\pgfqpoint{3.942350in}{3.293061in}}%
\pgfpathlineto{\pgfqpoint{3.936150in}{3.293061in}}%
\pgfpathlineto{\pgfqpoint{3.936150in}{0.660000in}}%
\pgfpathclose%
\pgfusepath{fill}%
\end{pgfscope}%
\begin{pgfscope}%
\pgfpathrectangle{\pgfqpoint{1.250000in}{0.660000in}}{\pgfqpoint{7.750000in}{4.620000in}}%
\pgfusepath{clip}%
\pgfsetbuttcap%
\pgfsetmiterjoin%
\definecolor{currentfill}{rgb}{0.000000,0.000000,0.000000}%
\pgfsetfillcolor{currentfill}%
\pgfsetlinewidth{0.000000pt}%
\definecolor{currentstroke}{rgb}{0.000000,0.000000,0.000000}%
\pgfsetstrokecolor{currentstroke}%
\pgfsetstrokeopacity{0.000000}%
\pgfsetdash{}{0pt}%
\pgfpathmoveto{\pgfqpoint{3.943900in}{0.660000in}}%
\pgfpathlineto{\pgfqpoint{3.950100in}{0.660000in}}%
\pgfpathlineto{\pgfqpoint{3.950100in}{3.292725in}}%
\pgfpathlineto{\pgfqpoint{3.943900in}{3.292725in}}%
\pgfpathlineto{\pgfqpoint{3.943900in}{0.660000in}}%
\pgfpathclose%
\pgfusepath{fill}%
\end{pgfscope}%
\begin{pgfscope}%
\pgfpathrectangle{\pgfqpoint{1.250000in}{0.660000in}}{\pgfqpoint{7.750000in}{4.620000in}}%
\pgfusepath{clip}%
\pgfsetbuttcap%
\pgfsetmiterjoin%
\definecolor{currentfill}{rgb}{0.000000,0.000000,0.000000}%
\pgfsetfillcolor{currentfill}%
\pgfsetlinewidth{0.000000pt}%
\definecolor{currentstroke}{rgb}{0.000000,0.000000,0.000000}%
\pgfsetstrokecolor{currentstroke}%
\pgfsetstrokeopacity{0.000000}%
\pgfsetdash{}{0pt}%
\pgfpathmoveto{\pgfqpoint{3.951650in}{0.660000in}}%
\pgfpathlineto{\pgfqpoint{3.957850in}{0.660000in}}%
\pgfpathlineto{\pgfqpoint{3.957850in}{3.292394in}}%
\pgfpathlineto{\pgfqpoint{3.951650in}{3.292394in}}%
\pgfpathlineto{\pgfqpoint{3.951650in}{0.660000in}}%
\pgfpathclose%
\pgfusepath{fill}%
\end{pgfscope}%
\begin{pgfscope}%
\pgfpathrectangle{\pgfqpoint{1.250000in}{0.660000in}}{\pgfqpoint{7.750000in}{4.620000in}}%
\pgfusepath{clip}%
\pgfsetbuttcap%
\pgfsetmiterjoin%
\definecolor{currentfill}{rgb}{0.000000,0.000000,0.000000}%
\pgfsetfillcolor{currentfill}%
\pgfsetlinewidth{0.000000pt}%
\definecolor{currentstroke}{rgb}{0.000000,0.000000,0.000000}%
\pgfsetstrokecolor{currentstroke}%
\pgfsetstrokeopacity{0.000000}%
\pgfsetdash{}{0pt}%
\pgfpathmoveto{\pgfqpoint{3.959400in}{0.660000in}}%
\pgfpathlineto{\pgfqpoint{3.965600in}{0.660000in}}%
\pgfpathlineto{\pgfqpoint{3.965600in}{3.292063in}}%
\pgfpathlineto{\pgfqpoint{3.959400in}{3.292063in}}%
\pgfpathlineto{\pgfqpoint{3.959400in}{0.660000in}}%
\pgfpathclose%
\pgfusepath{fill}%
\end{pgfscope}%
\begin{pgfscope}%
\pgfpathrectangle{\pgfqpoint{1.250000in}{0.660000in}}{\pgfqpoint{7.750000in}{4.620000in}}%
\pgfusepath{clip}%
\pgfsetbuttcap%
\pgfsetmiterjoin%
\definecolor{currentfill}{rgb}{0.000000,0.000000,0.000000}%
\pgfsetfillcolor{currentfill}%
\pgfsetlinewidth{0.000000pt}%
\definecolor{currentstroke}{rgb}{0.000000,0.000000,0.000000}%
\pgfsetstrokecolor{currentstroke}%
\pgfsetstrokeopacity{0.000000}%
\pgfsetdash{}{0pt}%
\pgfpathmoveto{\pgfqpoint{3.967150in}{0.660000in}}%
\pgfpathlineto{\pgfqpoint{3.973350in}{0.660000in}}%
\pgfpathlineto{\pgfqpoint{3.973350in}{3.291732in}}%
\pgfpathlineto{\pgfqpoint{3.967150in}{3.291732in}}%
\pgfpathlineto{\pgfqpoint{3.967150in}{0.660000in}}%
\pgfpathclose%
\pgfusepath{fill}%
\end{pgfscope}%
\begin{pgfscope}%
\pgfpathrectangle{\pgfqpoint{1.250000in}{0.660000in}}{\pgfqpoint{7.750000in}{4.620000in}}%
\pgfusepath{clip}%
\pgfsetbuttcap%
\pgfsetmiterjoin%
\definecolor{currentfill}{rgb}{0.000000,0.000000,0.000000}%
\pgfsetfillcolor{currentfill}%
\pgfsetlinewidth{0.000000pt}%
\definecolor{currentstroke}{rgb}{0.000000,0.000000,0.000000}%
\pgfsetstrokecolor{currentstroke}%
\pgfsetstrokeopacity{0.000000}%
\pgfsetdash{}{0pt}%
\pgfpathmoveto{\pgfqpoint{3.974900in}{0.660000in}}%
\pgfpathlineto{\pgfqpoint{3.981100in}{0.660000in}}%
\pgfpathlineto{\pgfqpoint{3.981100in}{3.291401in}}%
\pgfpathlineto{\pgfqpoint{3.974900in}{3.291401in}}%
\pgfpathlineto{\pgfqpoint{3.974900in}{0.660000in}}%
\pgfpathclose%
\pgfusepath{fill}%
\end{pgfscope}%
\begin{pgfscope}%
\pgfpathrectangle{\pgfqpoint{1.250000in}{0.660000in}}{\pgfqpoint{7.750000in}{4.620000in}}%
\pgfusepath{clip}%
\pgfsetbuttcap%
\pgfsetmiterjoin%
\definecolor{currentfill}{rgb}{0.000000,0.000000,0.000000}%
\pgfsetfillcolor{currentfill}%
\pgfsetlinewidth{0.000000pt}%
\definecolor{currentstroke}{rgb}{0.000000,0.000000,0.000000}%
\pgfsetstrokecolor{currentstroke}%
\pgfsetstrokeopacity{0.000000}%
\pgfsetdash{}{0pt}%
\pgfpathmoveto{\pgfqpoint{3.982650in}{0.660000in}}%
\pgfpathlineto{\pgfqpoint{3.988850in}{0.660000in}}%
\pgfpathlineto{\pgfqpoint{3.988850in}{3.291072in}}%
\pgfpathlineto{\pgfqpoint{3.982650in}{3.291072in}}%
\pgfpathlineto{\pgfqpoint{3.982650in}{0.660000in}}%
\pgfpathclose%
\pgfusepath{fill}%
\end{pgfscope}%
\begin{pgfscope}%
\pgfpathrectangle{\pgfqpoint{1.250000in}{0.660000in}}{\pgfqpoint{7.750000in}{4.620000in}}%
\pgfusepath{clip}%
\pgfsetbuttcap%
\pgfsetmiterjoin%
\definecolor{currentfill}{rgb}{0.000000,0.000000,0.000000}%
\pgfsetfillcolor{currentfill}%
\pgfsetlinewidth{0.000000pt}%
\definecolor{currentstroke}{rgb}{0.000000,0.000000,0.000000}%
\pgfsetstrokecolor{currentstroke}%
\pgfsetstrokeopacity{0.000000}%
\pgfsetdash{}{0pt}%
\pgfpathmoveto{\pgfqpoint{3.990400in}{0.660000in}}%
\pgfpathlineto{\pgfqpoint{3.996600in}{0.660000in}}%
\pgfpathlineto{\pgfqpoint{3.996600in}{3.290747in}}%
\pgfpathlineto{\pgfqpoint{3.990400in}{3.290747in}}%
\pgfpathlineto{\pgfqpoint{3.990400in}{0.660000in}}%
\pgfpathclose%
\pgfusepath{fill}%
\end{pgfscope}%
\begin{pgfscope}%
\pgfpathrectangle{\pgfqpoint{1.250000in}{0.660000in}}{\pgfqpoint{7.750000in}{4.620000in}}%
\pgfusepath{clip}%
\pgfsetbuttcap%
\pgfsetmiterjoin%
\definecolor{currentfill}{rgb}{0.000000,0.000000,0.000000}%
\pgfsetfillcolor{currentfill}%
\pgfsetlinewidth{0.000000pt}%
\definecolor{currentstroke}{rgb}{0.000000,0.000000,0.000000}%
\pgfsetstrokecolor{currentstroke}%
\pgfsetstrokeopacity{0.000000}%
\pgfsetdash{}{0pt}%
\pgfpathmoveto{\pgfqpoint{3.998150in}{0.660000in}}%
\pgfpathlineto{\pgfqpoint{4.004350in}{0.660000in}}%
\pgfpathlineto{\pgfqpoint{4.004350in}{3.290423in}}%
\pgfpathlineto{\pgfqpoint{3.998150in}{3.290423in}}%
\pgfpathlineto{\pgfqpoint{3.998150in}{0.660000in}}%
\pgfpathclose%
\pgfusepath{fill}%
\end{pgfscope}%
\begin{pgfscope}%
\pgfpathrectangle{\pgfqpoint{1.250000in}{0.660000in}}{\pgfqpoint{7.750000in}{4.620000in}}%
\pgfusepath{clip}%
\pgfsetbuttcap%
\pgfsetmiterjoin%
\definecolor{currentfill}{rgb}{0.000000,0.000000,0.000000}%
\pgfsetfillcolor{currentfill}%
\pgfsetlinewidth{0.000000pt}%
\definecolor{currentstroke}{rgb}{0.000000,0.000000,0.000000}%
\pgfsetstrokecolor{currentstroke}%
\pgfsetstrokeopacity{0.000000}%
\pgfsetdash{}{0pt}%
\pgfpathmoveto{\pgfqpoint{4.005900in}{0.660000in}}%
\pgfpathlineto{\pgfqpoint{4.012100in}{0.660000in}}%
\pgfpathlineto{\pgfqpoint{4.012100in}{3.290098in}}%
\pgfpathlineto{\pgfqpoint{4.005900in}{3.290098in}}%
\pgfpathlineto{\pgfqpoint{4.005900in}{0.660000in}}%
\pgfpathclose%
\pgfusepath{fill}%
\end{pgfscope}%
\begin{pgfscope}%
\pgfpathrectangle{\pgfqpoint{1.250000in}{0.660000in}}{\pgfqpoint{7.750000in}{4.620000in}}%
\pgfusepath{clip}%
\pgfsetbuttcap%
\pgfsetmiterjoin%
\definecolor{currentfill}{rgb}{0.000000,0.000000,0.000000}%
\pgfsetfillcolor{currentfill}%
\pgfsetlinewidth{0.000000pt}%
\definecolor{currentstroke}{rgb}{0.000000,0.000000,0.000000}%
\pgfsetstrokecolor{currentstroke}%
\pgfsetstrokeopacity{0.000000}%
\pgfsetdash{}{0pt}%
\pgfpathmoveto{\pgfqpoint{4.013650in}{0.660000in}}%
\pgfpathlineto{\pgfqpoint{4.019850in}{0.660000in}}%
\pgfpathlineto{\pgfqpoint{4.019850in}{3.289773in}}%
\pgfpathlineto{\pgfqpoint{4.013650in}{3.289773in}}%
\pgfpathlineto{\pgfqpoint{4.013650in}{0.660000in}}%
\pgfpathclose%
\pgfusepath{fill}%
\end{pgfscope}%
\begin{pgfscope}%
\pgfpathrectangle{\pgfqpoint{1.250000in}{0.660000in}}{\pgfqpoint{7.750000in}{4.620000in}}%
\pgfusepath{clip}%
\pgfsetbuttcap%
\pgfsetmiterjoin%
\definecolor{currentfill}{rgb}{0.000000,0.000000,0.000000}%
\pgfsetfillcolor{currentfill}%
\pgfsetlinewidth{0.000000pt}%
\definecolor{currentstroke}{rgb}{0.000000,0.000000,0.000000}%
\pgfsetstrokecolor{currentstroke}%
\pgfsetstrokeopacity{0.000000}%
\pgfsetdash{}{0pt}%
\pgfpathmoveto{\pgfqpoint{4.021400in}{0.660000in}}%
\pgfpathlineto{\pgfqpoint{4.027600in}{0.660000in}}%
\pgfpathlineto{\pgfqpoint{4.027600in}{3.289451in}}%
\pgfpathlineto{\pgfqpoint{4.021400in}{3.289451in}}%
\pgfpathlineto{\pgfqpoint{4.021400in}{0.660000in}}%
\pgfpathclose%
\pgfusepath{fill}%
\end{pgfscope}%
\begin{pgfscope}%
\pgfpathrectangle{\pgfqpoint{1.250000in}{0.660000in}}{\pgfqpoint{7.750000in}{4.620000in}}%
\pgfusepath{clip}%
\pgfsetbuttcap%
\pgfsetmiterjoin%
\definecolor{currentfill}{rgb}{0.000000,0.000000,0.000000}%
\pgfsetfillcolor{currentfill}%
\pgfsetlinewidth{0.000000pt}%
\definecolor{currentstroke}{rgb}{0.000000,0.000000,0.000000}%
\pgfsetstrokecolor{currentstroke}%
\pgfsetstrokeopacity{0.000000}%
\pgfsetdash{}{0pt}%
\pgfpathmoveto{\pgfqpoint{4.029150in}{0.660000in}}%
\pgfpathlineto{\pgfqpoint{4.035350in}{0.660000in}}%
\pgfpathlineto{\pgfqpoint{4.035350in}{3.289132in}}%
\pgfpathlineto{\pgfqpoint{4.029150in}{3.289132in}}%
\pgfpathlineto{\pgfqpoint{4.029150in}{0.660000in}}%
\pgfpathclose%
\pgfusepath{fill}%
\end{pgfscope}%
\begin{pgfscope}%
\pgfpathrectangle{\pgfqpoint{1.250000in}{0.660000in}}{\pgfqpoint{7.750000in}{4.620000in}}%
\pgfusepath{clip}%
\pgfsetbuttcap%
\pgfsetmiterjoin%
\definecolor{currentfill}{rgb}{0.000000,0.000000,0.000000}%
\pgfsetfillcolor{currentfill}%
\pgfsetlinewidth{0.000000pt}%
\definecolor{currentstroke}{rgb}{0.000000,0.000000,0.000000}%
\pgfsetstrokecolor{currentstroke}%
\pgfsetstrokeopacity{0.000000}%
\pgfsetdash{}{0pt}%
\pgfpathmoveto{\pgfqpoint{4.036900in}{0.660000in}}%
\pgfpathlineto{\pgfqpoint{4.043100in}{0.660000in}}%
\pgfpathlineto{\pgfqpoint{4.043100in}{3.288814in}}%
\pgfpathlineto{\pgfqpoint{4.036900in}{3.288814in}}%
\pgfpathlineto{\pgfqpoint{4.036900in}{0.660000in}}%
\pgfpathclose%
\pgfusepath{fill}%
\end{pgfscope}%
\begin{pgfscope}%
\pgfpathrectangle{\pgfqpoint{1.250000in}{0.660000in}}{\pgfqpoint{7.750000in}{4.620000in}}%
\pgfusepath{clip}%
\pgfsetbuttcap%
\pgfsetmiterjoin%
\definecolor{currentfill}{rgb}{0.000000,0.000000,0.000000}%
\pgfsetfillcolor{currentfill}%
\pgfsetlinewidth{0.000000pt}%
\definecolor{currentstroke}{rgb}{0.000000,0.000000,0.000000}%
\pgfsetstrokecolor{currentstroke}%
\pgfsetstrokeopacity{0.000000}%
\pgfsetdash{}{0pt}%
\pgfpathmoveto{\pgfqpoint{4.044650in}{0.660000in}}%
\pgfpathlineto{\pgfqpoint{4.050850in}{0.660000in}}%
\pgfpathlineto{\pgfqpoint{4.050850in}{3.288495in}}%
\pgfpathlineto{\pgfqpoint{4.044650in}{3.288495in}}%
\pgfpathlineto{\pgfqpoint{4.044650in}{0.660000in}}%
\pgfpathclose%
\pgfusepath{fill}%
\end{pgfscope}%
\begin{pgfscope}%
\pgfpathrectangle{\pgfqpoint{1.250000in}{0.660000in}}{\pgfqpoint{7.750000in}{4.620000in}}%
\pgfusepath{clip}%
\pgfsetbuttcap%
\pgfsetmiterjoin%
\definecolor{currentfill}{rgb}{0.000000,0.000000,0.000000}%
\pgfsetfillcolor{currentfill}%
\pgfsetlinewidth{0.000000pt}%
\definecolor{currentstroke}{rgb}{0.000000,0.000000,0.000000}%
\pgfsetstrokecolor{currentstroke}%
\pgfsetstrokeopacity{0.000000}%
\pgfsetdash{}{0pt}%
\pgfpathmoveto{\pgfqpoint{4.052400in}{0.660000in}}%
\pgfpathlineto{\pgfqpoint{4.058600in}{0.660000in}}%
\pgfpathlineto{\pgfqpoint{4.058600in}{3.288177in}}%
\pgfpathlineto{\pgfqpoint{4.052400in}{3.288177in}}%
\pgfpathlineto{\pgfqpoint{4.052400in}{0.660000in}}%
\pgfpathclose%
\pgfusepath{fill}%
\end{pgfscope}%
\begin{pgfscope}%
\pgfpathrectangle{\pgfqpoint{1.250000in}{0.660000in}}{\pgfqpoint{7.750000in}{4.620000in}}%
\pgfusepath{clip}%
\pgfsetbuttcap%
\pgfsetmiterjoin%
\definecolor{currentfill}{rgb}{0.000000,0.000000,0.000000}%
\pgfsetfillcolor{currentfill}%
\pgfsetlinewidth{0.000000pt}%
\definecolor{currentstroke}{rgb}{0.000000,0.000000,0.000000}%
\pgfsetstrokecolor{currentstroke}%
\pgfsetstrokeopacity{0.000000}%
\pgfsetdash{}{0pt}%
\pgfpathmoveto{\pgfqpoint{4.060150in}{0.660000in}}%
\pgfpathlineto{\pgfqpoint{4.066350in}{0.660000in}}%
\pgfpathlineto{\pgfqpoint{4.066350in}{3.287860in}}%
\pgfpathlineto{\pgfqpoint{4.060150in}{3.287860in}}%
\pgfpathlineto{\pgfqpoint{4.060150in}{0.660000in}}%
\pgfpathclose%
\pgfusepath{fill}%
\end{pgfscope}%
\begin{pgfscope}%
\pgfpathrectangle{\pgfqpoint{1.250000in}{0.660000in}}{\pgfqpoint{7.750000in}{4.620000in}}%
\pgfusepath{clip}%
\pgfsetbuttcap%
\pgfsetmiterjoin%
\definecolor{currentfill}{rgb}{0.000000,0.000000,0.000000}%
\pgfsetfillcolor{currentfill}%
\pgfsetlinewidth{0.000000pt}%
\definecolor{currentstroke}{rgb}{0.000000,0.000000,0.000000}%
\pgfsetstrokecolor{currentstroke}%
\pgfsetstrokeopacity{0.000000}%
\pgfsetdash{}{0pt}%
\pgfpathmoveto{\pgfqpoint{4.067900in}{0.660000in}}%
\pgfpathlineto{\pgfqpoint{4.074100in}{0.660000in}}%
\pgfpathlineto{\pgfqpoint{4.074100in}{3.287548in}}%
\pgfpathlineto{\pgfqpoint{4.067900in}{3.287548in}}%
\pgfpathlineto{\pgfqpoint{4.067900in}{0.660000in}}%
\pgfpathclose%
\pgfusepath{fill}%
\end{pgfscope}%
\begin{pgfscope}%
\pgfpathrectangle{\pgfqpoint{1.250000in}{0.660000in}}{\pgfqpoint{7.750000in}{4.620000in}}%
\pgfusepath{clip}%
\pgfsetbuttcap%
\pgfsetmiterjoin%
\definecolor{currentfill}{rgb}{0.000000,0.000000,0.000000}%
\pgfsetfillcolor{currentfill}%
\pgfsetlinewidth{0.000000pt}%
\definecolor{currentstroke}{rgb}{0.000000,0.000000,0.000000}%
\pgfsetstrokecolor{currentstroke}%
\pgfsetstrokeopacity{0.000000}%
\pgfsetdash{}{0pt}%
\pgfpathmoveto{\pgfqpoint{4.075650in}{0.660000in}}%
\pgfpathlineto{\pgfqpoint{4.081850in}{0.660000in}}%
\pgfpathlineto{\pgfqpoint{4.081850in}{3.287236in}}%
\pgfpathlineto{\pgfqpoint{4.075650in}{3.287236in}}%
\pgfpathlineto{\pgfqpoint{4.075650in}{0.660000in}}%
\pgfpathclose%
\pgfusepath{fill}%
\end{pgfscope}%
\begin{pgfscope}%
\pgfpathrectangle{\pgfqpoint{1.250000in}{0.660000in}}{\pgfqpoint{7.750000in}{4.620000in}}%
\pgfusepath{clip}%
\pgfsetbuttcap%
\pgfsetmiterjoin%
\definecolor{currentfill}{rgb}{0.000000,0.000000,0.000000}%
\pgfsetfillcolor{currentfill}%
\pgfsetlinewidth{0.000000pt}%
\definecolor{currentstroke}{rgb}{0.000000,0.000000,0.000000}%
\pgfsetstrokecolor{currentstroke}%
\pgfsetstrokeopacity{0.000000}%
\pgfsetdash{}{0pt}%
\pgfpathmoveto{\pgfqpoint{4.083400in}{0.660000in}}%
\pgfpathlineto{\pgfqpoint{4.089600in}{0.660000in}}%
\pgfpathlineto{\pgfqpoint{4.089600in}{3.286923in}}%
\pgfpathlineto{\pgfqpoint{4.083400in}{3.286923in}}%
\pgfpathlineto{\pgfqpoint{4.083400in}{0.660000in}}%
\pgfpathclose%
\pgfusepath{fill}%
\end{pgfscope}%
\begin{pgfscope}%
\pgfpathrectangle{\pgfqpoint{1.250000in}{0.660000in}}{\pgfqpoint{7.750000in}{4.620000in}}%
\pgfusepath{clip}%
\pgfsetbuttcap%
\pgfsetmiterjoin%
\definecolor{currentfill}{rgb}{0.000000,0.000000,0.000000}%
\pgfsetfillcolor{currentfill}%
\pgfsetlinewidth{0.000000pt}%
\definecolor{currentstroke}{rgb}{0.000000,0.000000,0.000000}%
\pgfsetstrokecolor{currentstroke}%
\pgfsetstrokeopacity{0.000000}%
\pgfsetdash{}{0pt}%
\pgfpathmoveto{\pgfqpoint{4.091150in}{0.660000in}}%
\pgfpathlineto{\pgfqpoint{4.097350in}{0.660000in}}%
\pgfpathlineto{\pgfqpoint{4.097350in}{3.286611in}}%
\pgfpathlineto{\pgfqpoint{4.091150in}{3.286611in}}%
\pgfpathlineto{\pgfqpoint{4.091150in}{0.660000in}}%
\pgfpathclose%
\pgfusepath{fill}%
\end{pgfscope}%
\begin{pgfscope}%
\pgfpathrectangle{\pgfqpoint{1.250000in}{0.660000in}}{\pgfqpoint{7.750000in}{4.620000in}}%
\pgfusepath{clip}%
\pgfsetbuttcap%
\pgfsetmiterjoin%
\definecolor{currentfill}{rgb}{0.000000,0.000000,0.000000}%
\pgfsetfillcolor{currentfill}%
\pgfsetlinewidth{0.000000pt}%
\definecolor{currentstroke}{rgb}{0.000000,0.000000,0.000000}%
\pgfsetstrokecolor{currentstroke}%
\pgfsetstrokeopacity{0.000000}%
\pgfsetdash{}{0pt}%
\pgfpathmoveto{\pgfqpoint{4.098900in}{0.660000in}}%
\pgfpathlineto{\pgfqpoint{4.105100in}{0.660000in}}%
\pgfpathlineto{\pgfqpoint{4.105100in}{3.286299in}}%
\pgfpathlineto{\pgfqpoint{4.098900in}{3.286299in}}%
\pgfpathlineto{\pgfqpoint{4.098900in}{0.660000in}}%
\pgfpathclose%
\pgfusepath{fill}%
\end{pgfscope}%
\begin{pgfscope}%
\pgfpathrectangle{\pgfqpoint{1.250000in}{0.660000in}}{\pgfqpoint{7.750000in}{4.620000in}}%
\pgfusepath{clip}%
\pgfsetbuttcap%
\pgfsetmiterjoin%
\definecolor{currentfill}{rgb}{0.000000,0.000000,0.000000}%
\pgfsetfillcolor{currentfill}%
\pgfsetlinewidth{0.000000pt}%
\definecolor{currentstroke}{rgb}{0.000000,0.000000,0.000000}%
\pgfsetstrokecolor{currentstroke}%
\pgfsetstrokeopacity{0.000000}%
\pgfsetdash{}{0pt}%
\pgfpathmoveto{\pgfqpoint{4.106650in}{0.660000in}}%
\pgfpathlineto{\pgfqpoint{4.112850in}{0.660000in}}%
\pgfpathlineto{\pgfqpoint{4.112850in}{3.285993in}}%
\pgfpathlineto{\pgfqpoint{4.106650in}{3.285993in}}%
\pgfpathlineto{\pgfqpoint{4.106650in}{0.660000in}}%
\pgfpathclose%
\pgfusepath{fill}%
\end{pgfscope}%
\begin{pgfscope}%
\pgfpathrectangle{\pgfqpoint{1.250000in}{0.660000in}}{\pgfqpoint{7.750000in}{4.620000in}}%
\pgfusepath{clip}%
\pgfsetbuttcap%
\pgfsetmiterjoin%
\definecolor{currentfill}{rgb}{0.000000,0.000000,0.000000}%
\pgfsetfillcolor{currentfill}%
\pgfsetlinewidth{0.000000pt}%
\definecolor{currentstroke}{rgb}{0.000000,0.000000,0.000000}%
\pgfsetstrokecolor{currentstroke}%
\pgfsetstrokeopacity{0.000000}%
\pgfsetdash{}{0pt}%
\pgfpathmoveto{\pgfqpoint{4.114400in}{0.660000in}}%
\pgfpathlineto{\pgfqpoint{4.120600in}{0.660000in}}%
\pgfpathlineto{\pgfqpoint{4.120600in}{3.285687in}}%
\pgfpathlineto{\pgfqpoint{4.114400in}{3.285687in}}%
\pgfpathlineto{\pgfqpoint{4.114400in}{0.660000in}}%
\pgfpathclose%
\pgfusepath{fill}%
\end{pgfscope}%
\begin{pgfscope}%
\pgfpathrectangle{\pgfqpoint{1.250000in}{0.660000in}}{\pgfqpoint{7.750000in}{4.620000in}}%
\pgfusepath{clip}%
\pgfsetbuttcap%
\pgfsetmiterjoin%
\definecolor{currentfill}{rgb}{0.000000,0.000000,0.000000}%
\pgfsetfillcolor{currentfill}%
\pgfsetlinewidth{0.000000pt}%
\definecolor{currentstroke}{rgb}{0.000000,0.000000,0.000000}%
\pgfsetstrokecolor{currentstroke}%
\pgfsetstrokeopacity{0.000000}%
\pgfsetdash{}{0pt}%
\pgfpathmoveto{\pgfqpoint{4.122150in}{0.660000in}}%
\pgfpathlineto{\pgfqpoint{4.128350in}{0.660000in}}%
\pgfpathlineto{\pgfqpoint{4.128350in}{3.285381in}}%
\pgfpathlineto{\pgfqpoint{4.122150in}{3.285381in}}%
\pgfpathlineto{\pgfqpoint{4.122150in}{0.660000in}}%
\pgfpathclose%
\pgfusepath{fill}%
\end{pgfscope}%
\begin{pgfscope}%
\pgfpathrectangle{\pgfqpoint{1.250000in}{0.660000in}}{\pgfqpoint{7.750000in}{4.620000in}}%
\pgfusepath{clip}%
\pgfsetbuttcap%
\pgfsetmiterjoin%
\definecolor{currentfill}{rgb}{0.000000,0.000000,0.000000}%
\pgfsetfillcolor{currentfill}%
\pgfsetlinewidth{0.000000pt}%
\definecolor{currentstroke}{rgb}{0.000000,0.000000,0.000000}%
\pgfsetstrokecolor{currentstroke}%
\pgfsetstrokeopacity{0.000000}%
\pgfsetdash{}{0pt}%
\pgfpathmoveto{\pgfqpoint{4.129900in}{0.660000in}}%
\pgfpathlineto{\pgfqpoint{4.136100in}{0.660000in}}%
\pgfpathlineto{\pgfqpoint{4.136100in}{3.285074in}}%
\pgfpathlineto{\pgfqpoint{4.129900in}{3.285074in}}%
\pgfpathlineto{\pgfqpoint{4.129900in}{0.660000in}}%
\pgfpathclose%
\pgfusepath{fill}%
\end{pgfscope}%
\begin{pgfscope}%
\pgfpathrectangle{\pgfqpoint{1.250000in}{0.660000in}}{\pgfqpoint{7.750000in}{4.620000in}}%
\pgfusepath{clip}%
\pgfsetbuttcap%
\pgfsetmiterjoin%
\definecolor{currentfill}{rgb}{0.000000,0.000000,0.000000}%
\pgfsetfillcolor{currentfill}%
\pgfsetlinewidth{0.000000pt}%
\definecolor{currentstroke}{rgb}{0.000000,0.000000,0.000000}%
\pgfsetstrokecolor{currentstroke}%
\pgfsetstrokeopacity{0.000000}%
\pgfsetdash{}{0pt}%
\pgfpathmoveto{\pgfqpoint{4.137650in}{0.660000in}}%
\pgfpathlineto{\pgfqpoint{4.143850in}{0.660000in}}%
\pgfpathlineto{\pgfqpoint{4.143850in}{3.284768in}}%
\pgfpathlineto{\pgfqpoint{4.137650in}{3.284768in}}%
\pgfpathlineto{\pgfqpoint{4.137650in}{0.660000in}}%
\pgfpathclose%
\pgfusepath{fill}%
\end{pgfscope}%
\begin{pgfscope}%
\pgfpathrectangle{\pgfqpoint{1.250000in}{0.660000in}}{\pgfqpoint{7.750000in}{4.620000in}}%
\pgfusepath{clip}%
\pgfsetbuttcap%
\pgfsetmiterjoin%
\definecolor{currentfill}{rgb}{0.000000,0.000000,0.000000}%
\pgfsetfillcolor{currentfill}%
\pgfsetlinewidth{0.000000pt}%
\definecolor{currentstroke}{rgb}{0.000000,0.000000,0.000000}%
\pgfsetstrokecolor{currentstroke}%
\pgfsetstrokeopacity{0.000000}%
\pgfsetdash{}{0pt}%
\pgfpathmoveto{\pgfqpoint{4.145400in}{0.660000in}}%
\pgfpathlineto{\pgfqpoint{4.151600in}{0.660000in}}%
\pgfpathlineto{\pgfqpoint{4.151600in}{3.284467in}}%
\pgfpathlineto{\pgfqpoint{4.145400in}{3.284467in}}%
\pgfpathlineto{\pgfqpoint{4.145400in}{0.660000in}}%
\pgfpathclose%
\pgfusepath{fill}%
\end{pgfscope}%
\begin{pgfscope}%
\pgfpathrectangle{\pgfqpoint{1.250000in}{0.660000in}}{\pgfqpoint{7.750000in}{4.620000in}}%
\pgfusepath{clip}%
\pgfsetbuttcap%
\pgfsetmiterjoin%
\definecolor{currentfill}{rgb}{0.000000,0.000000,0.000000}%
\pgfsetfillcolor{currentfill}%
\pgfsetlinewidth{0.000000pt}%
\definecolor{currentstroke}{rgb}{0.000000,0.000000,0.000000}%
\pgfsetstrokecolor{currentstroke}%
\pgfsetstrokeopacity{0.000000}%
\pgfsetdash{}{0pt}%
\pgfpathmoveto{\pgfqpoint{4.153150in}{0.660000in}}%
\pgfpathlineto{\pgfqpoint{4.159350in}{0.660000in}}%
\pgfpathlineto{\pgfqpoint{4.159350in}{3.284167in}}%
\pgfpathlineto{\pgfqpoint{4.153150in}{3.284167in}}%
\pgfpathlineto{\pgfqpoint{4.153150in}{0.660000in}}%
\pgfpathclose%
\pgfusepath{fill}%
\end{pgfscope}%
\begin{pgfscope}%
\pgfpathrectangle{\pgfqpoint{1.250000in}{0.660000in}}{\pgfqpoint{7.750000in}{4.620000in}}%
\pgfusepath{clip}%
\pgfsetbuttcap%
\pgfsetmiterjoin%
\definecolor{currentfill}{rgb}{0.000000,0.000000,0.000000}%
\pgfsetfillcolor{currentfill}%
\pgfsetlinewidth{0.000000pt}%
\definecolor{currentstroke}{rgb}{0.000000,0.000000,0.000000}%
\pgfsetstrokecolor{currentstroke}%
\pgfsetstrokeopacity{0.000000}%
\pgfsetdash{}{0pt}%
\pgfpathmoveto{\pgfqpoint{4.160900in}{0.660000in}}%
\pgfpathlineto{\pgfqpoint{4.167100in}{0.660000in}}%
\pgfpathlineto{\pgfqpoint{4.167100in}{3.283866in}}%
\pgfpathlineto{\pgfqpoint{4.160900in}{3.283866in}}%
\pgfpathlineto{\pgfqpoint{4.160900in}{0.660000in}}%
\pgfpathclose%
\pgfusepath{fill}%
\end{pgfscope}%
\begin{pgfscope}%
\pgfpathrectangle{\pgfqpoint{1.250000in}{0.660000in}}{\pgfqpoint{7.750000in}{4.620000in}}%
\pgfusepath{clip}%
\pgfsetbuttcap%
\pgfsetmiterjoin%
\definecolor{currentfill}{rgb}{0.000000,0.000000,0.000000}%
\pgfsetfillcolor{currentfill}%
\pgfsetlinewidth{0.000000pt}%
\definecolor{currentstroke}{rgb}{0.000000,0.000000,0.000000}%
\pgfsetstrokecolor{currentstroke}%
\pgfsetstrokeopacity{0.000000}%
\pgfsetdash{}{0pt}%
\pgfpathmoveto{\pgfqpoint{4.168650in}{0.660000in}}%
\pgfpathlineto{\pgfqpoint{4.174850in}{0.660000in}}%
\pgfpathlineto{\pgfqpoint{4.174850in}{3.283566in}}%
\pgfpathlineto{\pgfqpoint{4.168650in}{3.283566in}}%
\pgfpathlineto{\pgfqpoint{4.168650in}{0.660000in}}%
\pgfpathclose%
\pgfusepath{fill}%
\end{pgfscope}%
\begin{pgfscope}%
\pgfpathrectangle{\pgfqpoint{1.250000in}{0.660000in}}{\pgfqpoint{7.750000in}{4.620000in}}%
\pgfusepath{clip}%
\pgfsetbuttcap%
\pgfsetmiterjoin%
\definecolor{currentfill}{rgb}{0.000000,0.000000,0.000000}%
\pgfsetfillcolor{currentfill}%
\pgfsetlinewidth{0.000000pt}%
\definecolor{currentstroke}{rgb}{0.000000,0.000000,0.000000}%
\pgfsetstrokecolor{currentstroke}%
\pgfsetstrokeopacity{0.000000}%
\pgfsetdash{}{0pt}%
\pgfpathmoveto{\pgfqpoint{4.176400in}{0.660000in}}%
\pgfpathlineto{\pgfqpoint{4.182600in}{0.660000in}}%
\pgfpathlineto{\pgfqpoint{4.182600in}{3.283266in}}%
\pgfpathlineto{\pgfqpoint{4.176400in}{3.283266in}}%
\pgfpathlineto{\pgfqpoint{4.176400in}{0.660000in}}%
\pgfpathclose%
\pgfusepath{fill}%
\end{pgfscope}%
\begin{pgfscope}%
\pgfpathrectangle{\pgfqpoint{1.250000in}{0.660000in}}{\pgfqpoint{7.750000in}{4.620000in}}%
\pgfusepath{clip}%
\pgfsetbuttcap%
\pgfsetmiterjoin%
\definecolor{currentfill}{rgb}{0.000000,0.000000,0.000000}%
\pgfsetfillcolor{currentfill}%
\pgfsetlinewidth{0.000000pt}%
\definecolor{currentstroke}{rgb}{0.000000,0.000000,0.000000}%
\pgfsetstrokecolor{currentstroke}%
\pgfsetstrokeopacity{0.000000}%
\pgfsetdash{}{0pt}%
\pgfpathmoveto{\pgfqpoint{4.184150in}{0.660000in}}%
\pgfpathlineto{\pgfqpoint{4.190350in}{0.660000in}}%
\pgfpathlineto{\pgfqpoint{4.190350in}{3.282968in}}%
\pgfpathlineto{\pgfqpoint{4.184150in}{3.282968in}}%
\pgfpathlineto{\pgfqpoint{4.184150in}{0.660000in}}%
\pgfpathclose%
\pgfusepath{fill}%
\end{pgfscope}%
\begin{pgfscope}%
\pgfpathrectangle{\pgfqpoint{1.250000in}{0.660000in}}{\pgfqpoint{7.750000in}{4.620000in}}%
\pgfusepath{clip}%
\pgfsetbuttcap%
\pgfsetmiterjoin%
\definecolor{currentfill}{rgb}{0.000000,0.000000,0.000000}%
\pgfsetfillcolor{currentfill}%
\pgfsetlinewidth{0.000000pt}%
\definecolor{currentstroke}{rgb}{0.000000,0.000000,0.000000}%
\pgfsetstrokecolor{currentstroke}%
\pgfsetstrokeopacity{0.000000}%
\pgfsetdash{}{0pt}%
\pgfpathmoveto{\pgfqpoint{4.191900in}{0.660000in}}%
\pgfpathlineto{\pgfqpoint{4.198100in}{0.660000in}}%
\pgfpathlineto{\pgfqpoint{4.198100in}{3.282674in}}%
\pgfpathlineto{\pgfqpoint{4.191900in}{3.282674in}}%
\pgfpathlineto{\pgfqpoint{4.191900in}{0.660000in}}%
\pgfpathclose%
\pgfusepath{fill}%
\end{pgfscope}%
\begin{pgfscope}%
\pgfpathrectangle{\pgfqpoint{1.250000in}{0.660000in}}{\pgfqpoint{7.750000in}{4.620000in}}%
\pgfusepath{clip}%
\pgfsetbuttcap%
\pgfsetmiterjoin%
\definecolor{currentfill}{rgb}{0.000000,0.000000,0.000000}%
\pgfsetfillcolor{currentfill}%
\pgfsetlinewidth{0.000000pt}%
\definecolor{currentstroke}{rgb}{0.000000,0.000000,0.000000}%
\pgfsetstrokecolor{currentstroke}%
\pgfsetstrokeopacity{0.000000}%
\pgfsetdash{}{0pt}%
\pgfpathmoveto{\pgfqpoint{4.199650in}{0.660000in}}%
\pgfpathlineto{\pgfqpoint{4.205850in}{0.660000in}}%
\pgfpathlineto{\pgfqpoint{4.205850in}{3.282379in}}%
\pgfpathlineto{\pgfqpoint{4.199650in}{3.282379in}}%
\pgfpathlineto{\pgfqpoint{4.199650in}{0.660000in}}%
\pgfpathclose%
\pgfusepath{fill}%
\end{pgfscope}%
\begin{pgfscope}%
\pgfpathrectangle{\pgfqpoint{1.250000in}{0.660000in}}{\pgfqpoint{7.750000in}{4.620000in}}%
\pgfusepath{clip}%
\pgfsetbuttcap%
\pgfsetmiterjoin%
\definecolor{currentfill}{rgb}{0.000000,0.000000,0.000000}%
\pgfsetfillcolor{currentfill}%
\pgfsetlinewidth{0.000000pt}%
\definecolor{currentstroke}{rgb}{0.000000,0.000000,0.000000}%
\pgfsetstrokecolor{currentstroke}%
\pgfsetstrokeopacity{0.000000}%
\pgfsetdash{}{0pt}%
\pgfpathmoveto{\pgfqpoint{4.207400in}{0.660000in}}%
\pgfpathlineto{\pgfqpoint{4.213600in}{0.660000in}}%
\pgfpathlineto{\pgfqpoint{4.213600in}{3.282085in}}%
\pgfpathlineto{\pgfqpoint{4.207400in}{3.282085in}}%
\pgfpathlineto{\pgfqpoint{4.207400in}{0.660000in}}%
\pgfpathclose%
\pgfusepath{fill}%
\end{pgfscope}%
\begin{pgfscope}%
\pgfpathrectangle{\pgfqpoint{1.250000in}{0.660000in}}{\pgfqpoint{7.750000in}{4.620000in}}%
\pgfusepath{clip}%
\pgfsetbuttcap%
\pgfsetmiterjoin%
\definecolor{currentfill}{rgb}{0.000000,0.000000,0.000000}%
\pgfsetfillcolor{currentfill}%
\pgfsetlinewidth{0.000000pt}%
\definecolor{currentstroke}{rgb}{0.000000,0.000000,0.000000}%
\pgfsetstrokecolor{currentstroke}%
\pgfsetstrokeopacity{0.000000}%
\pgfsetdash{}{0pt}%
\pgfpathmoveto{\pgfqpoint{4.215150in}{0.660000in}}%
\pgfpathlineto{\pgfqpoint{4.221350in}{0.660000in}}%
\pgfpathlineto{\pgfqpoint{4.221350in}{3.281791in}}%
\pgfpathlineto{\pgfqpoint{4.215150in}{3.281791in}}%
\pgfpathlineto{\pgfqpoint{4.215150in}{0.660000in}}%
\pgfpathclose%
\pgfusepath{fill}%
\end{pgfscope}%
\begin{pgfscope}%
\pgfpathrectangle{\pgfqpoint{1.250000in}{0.660000in}}{\pgfqpoint{7.750000in}{4.620000in}}%
\pgfusepath{clip}%
\pgfsetbuttcap%
\pgfsetmiterjoin%
\definecolor{currentfill}{rgb}{0.000000,0.000000,0.000000}%
\pgfsetfillcolor{currentfill}%
\pgfsetlinewidth{0.000000pt}%
\definecolor{currentstroke}{rgb}{0.000000,0.000000,0.000000}%
\pgfsetstrokecolor{currentstroke}%
\pgfsetstrokeopacity{0.000000}%
\pgfsetdash{}{0pt}%
\pgfpathmoveto{\pgfqpoint{4.222900in}{0.660000in}}%
\pgfpathlineto{\pgfqpoint{4.229100in}{0.660000in}}%
\pgfpathlineto{\pgfqpoint{4.229100in}{3.281497in}}%
\pgfpathlineto{\pgfqpoint{4.222900in}{3.281497in}}%
\pgfpathlineto{\pgfqpoint{4.222900in}{0.660000in}}%
\pgfpathclose%
\pgfusepath{fill}%
\end{pgfscope}%
\begin{pgfscope}%
\pgfpathrectangle{\pgfqpoint{1.250000in}{0.660000in}}{\pgfqpoint{7.750000in}{4.620000in}}%
\pgfusepath{clip}%
\pgfsetbuttcap%
\pgfsetmiterjoin%
\definecolor{currentfill}{rgb}{0.000000,0.000000,0.000000}%
\pgfsetfillcolor{currentfill}%
\pgfsetlinewidth{0.000000pt}%
\definecolor{currentstroke}{rgb}{0.000000,0.000000,0.000000}%
\pgfsetstrokecolor{currentstroke}%
\pgfsetstrokeopacity{0.000000}%
\pgfsetdash{}{0pt}%
\pgfpathmoveto{\pgfqpoint{4.230650in}{0.660000in}}%
\pgfpathlineto{\pgfqpoint{4.236850in}{0.660000in}}%
\pgfpathlineto{\pgfqpoint{4.236850in}{3.281207in}}%
\pgfpathlineto{\pgfqpoint{4.230650in}{3.281207in}}%
\pgfpathlineto{\pgfqpoint{4.230650in}{0.660000in}}%
\pgfpathclose%
\pgfusepath{fill}%
\end{pgfscope}%
\begin{pgfscope}%
\pgfpathrectangle{\pgfqpoint{1.250000in}{0.660000in}}{\pgfqpoint{7.750000in}{4.620000in}}%
\pgfusepath{clip}%
\pgfsetbuttcap%
\pgfsetmiterjoin%
\definecolor{currentfill}{rgb}{0.000000,0.000000,0.000000}%
\pgfsetfillcolor{currentfill}%
\pgfsetlinewidth{0.000000pt}%
\definecolor{currentstroke}{rgb}{0.000000,0.000000,0.000000}%
\pgfsetstrokecolor{currentstroke}%
\pgfsetstrokeopacity{0.000000}%
\pgfsetdash{}{0pt}%
\pgfpathmoveto{\pgfqpoint{4.238400in}{0.660000in}}%
\pgfpathlineto{\pgfqpoint{4.244600in}{0.660000in}}%
\pgfpathlineto{\pgfqpoint{4.244600in}{3.280919in}}%
\pgfpathlineto{\pgfqpoint{4.238400in}{3.280919in}}%
\pgfpathlineto{\pgfqpoint{4.238400in}{0.660000in}}%
\pgfpathclose%
\pgfusepath{fill}%
\end{pgfscope}%
\begin{pgfscope}%
\pgfpathrectangle{\pgfqpoint{1.250000in}{0.660000in}}{\pgfqpoint{7.750000in}{4.620000in}}%
\pgfusepath{clip}%
\pgfsetbuttcap%
\pgfsetmiterjoin%
\definecolor{currentfill}{rgb}{0.000000,0.000000,0.000000}%
\pgfsetfillcolor{currentfill}%
\pgfsetlinewidth{0.000000pt}%
\definecolor{currentstroke}{rgb}{0.000000,0.000000,0.000000}%
\pgfsetstrokecolor{currentstroke}%
\pgfsetstrokeopacity{0.000000}%
\pgfsetdash{}{0pt}%
\pgfpathmoveto{\pgfqpoint{4.246150in}{0.660000in}}%
\pgfpathlineto{\pgfqpoint{4.252350in}{0.660000in}}%
\pgfpathlineto{\pgfqpoint{4.252350in}{3.280630in}}%
\pgfpathlineto{\pgfqpoint{4.246150in}{3.280630in}}%
\pgfpathlineto{\pgfqpoint{4.246150in}{0.660000in}}%
\pgfpathclose%
\pgfusepath{fill}%
\end{pgfscope}%
\begin{pgfscope}%
\pgfpathrectangle{\pgfqpoint{1.250000in}{0.660000in}}{\pgfqpoint{7.750000in}{4.620000in}}%
\pgfusepath{clip}%
\pgfsetbuttcap%
\pgfsetmiterjoin%
\definecolor{currentfill}{rgb}{0.000000,0.000000,0.000000}%
\pgfsetfillcolor{currentfill}%
\pgfsetlinewidth{0.000000pt}%
\definecolor{currentstroke}{rgb}{0.000000,0.000000,0.000000}%
\pgfsetstrokecolor{currentstroke}%
\pgfsetstrokeopacity{0.000000}%
\pgfsetdash{}{0pt}%
\pgfpathmoveto{\pgfqpoint{4.253900in}{0.660000in}}%
\pgfpathlineto{\pgfqpoint{4.260100in}{0.660000in}}%
\pgfpathlineto{\pgfqpoint{4.260100in}{3.280342in}}%
\pgfpathlineto{\pgfqpoint{4.253900in}{3.280342in}}%
\pgfpathlineto{\pgfqpoint{4.253900in}{0.660000in}}%
\pgfpathclose%
\pgfusepath{fill}%
\end{pgfscope}%
\begin{pgfscope}%
\pgfpathrectangle{\pgfqpoint{1.250000in}{0.660000in}}{\pgfqpoint{7.750000in}{4.620000in}}%
\pgfusepath{clip}%
\pgfsetbuttcap%
\pgfsetmiterjoin%
\definecolor{currentfill}{rgb}{0.000000,0.000000,0.000000}%
\pgfsetfillcolor{currentfill}%
\pgfsetlinewidth{0.000000pt}%
\definecolor{currentstroke}{rgb}{0.000000,0.000000,0.000000}%
\pgfsetstrokecolor{currentstroke}%
\pgfsetstrokeopacity{0.000000}%
\pgfsetdash{}{0pt}%
\pgfpathmoveto{\pgfqpoint{4.261650in}{0.660000in}}%
\pgfpathlineto{\pgfqpoint{4.267850in}{0.660000in}}%
\pgfpathlineto{\pgfqpoint{4.267850in}{3.280054in}}%
\pgfpathlineto{\pgfqpoint{4.261650in}{3.280054in}}%
\pgfpathlineto{\pgfqpoint{4.261650in}{0.660000in}}%
\pgfpathclose%
\pgfusepath{fill}%
\end{pgfscope}%
\begin{pgfscope}%
\pgfpathrectangle{\pgfqpoint{1.250000in}{0.660000in}}{\pgfqpoint{7.750000in}{4.620000in}}%
\pgfusepath{clip}%
\pgfsetbuttcap%
\pgfsetmiterjoin%
\definecolor{currentfill}{rgb}{0.000000,0.000000,0.000000}%
\pgfsetfillcolor{currentfill}%
\pgfsetlinewidth{0.000000pt}%
\definecolor{currentstroke}{rgb}{0.000000,0.000000,0.000000}%
\pgfsetstrokecolor{currentstroke}%
\pgfsetstrokeopacity{0.000000}%
\pgfsetdash{}{0pt}%
\pgfpathmoveto{\pgfqpoint{4.269400in}{0.660000in}}%
\pgfpathlineto{\pgfqpoint{4.275600in}{0.660000in}}%
\pgfpathlineto{\pgfqpoint{4.275600in}{3.279766in}}%
\pgfpathlineto{\pgfqpoint{4.269400in}{3.279766in}}%
\pgfpathlineto{\pgfqpoint{4.269400in}{0.660000in}}%
\pgfpathclose%
\pgfusepath{fill}%
\end{pgfscope}%
\begin{pgfscope}%
\pgfpathrectangle{\pgfqpoint{1.250000in}{0.660000in}}{\pgfqpoint{7.750000in}{4.620000in}}%
\pgfusepath{clip}%
\pgfsetbuttcap%
\pgfsetmiterjoin%
\definecolor{currentfill}{rgb}{0.000000,0.000000,0.000000}%
\pgfsetfillcolor{currentfill}%
\pgfsetlinewidth{0.000000pt}%
\definecolor{currentstroke}{rgb}{0.000000,0.000000,0.000000}%
\pgfsetstrokecolor{currentstroke}%
\pgfsetstrokeopacity{0.000000}%
\pgfsetdash{}{0pt}%
\pgfpathmoveto{\pgfqpoint{4.277150in}{0.660000in}}%
\pgfpathlineto{\pgfqpoint{4.283350in}{0.660000in}}%
\pgfpathlineto{\pgfqpoint{4.283350in}{3.279483in}}%
\pgfpathlineto{\pgfqpoint{4.277150in}{3.279483in}}%
\pgfpathlineto{\pgfqpoint{4.277150in}{0.660000in}}%
\pgfpathclose%
\pgfusepath{fill}%
\end{pgfscope}%
\begin{pgfscope}%
\pgfpathrectangle{\pgfqpoint{1.250000in}{0.660000in}}{\pgfqpoint{7.750000in}{4.620000in}}%
\pgfusepath{clip}%
\pgfsetbuttcap%
\pgfsetmiterjoin%
\definecolor{currentfill}{rgb}{0.000000,0.000000,0.000000}%
\pgfsetfillcolor{currentfill}%
\pgfsetlinewidth{0.000000pt}%
\definecolor{currentstroke}{rgb}{0.000000,0.000000,0.000000}%
\pgfsetstrokecolor{currentstroke}%
\pgfsetstrokeopacity{0.000000}%
\pgfsetdash{}{0pt}%
\pgfpathmoveto{\pgfqpoint{4.284900in}{0.660000in}}%
\pgfpathlineto{\pgfqpoint{4.291100in}{0.660000in}}%
\pgfpathlineto{\pgfqpoint{4.291100in}{3.279201in}}%
\pgfpathlineto{\pgfqpoint{4.284900in}{3.279201in}}%
\pgfpathlineto{\pgfqpoint{4.284900in}{0.660000in}}%
\pgfpathclose%
\pgfusepath{fill}%
\end{pgfscope}%
\begin{pgfscope}%
\pgfpathrectangle{\pgfqpoint{1.250000in}{0.660000in}}{\pgfqpoint{7.750000in}{4.620000in}}%
\pgfusepath{clip}%
\pgfsetbuttcap%
\pgfsetmiterjoin%
\definecolor{currentfill}{rgb}{0.000000,0.000000,0.000000}%
\pgfsetfillcolor{currentfill}%
\pgfsetlinewidth{0.000000pt}%
\definecolor{currentstroke}{rgb}{0.000000,0.000000,0.000000}%
\pgfsetstrokecolor{currentstroke}%
\pgfsetstrokeopacity{0.000000}%
\pgfsetdash{}{0pt}%
\pgfpathmoveto{\pgfqpoint{4.292650in}{0.660000in}}%
\pgfpathlineto{\pgfqpoint{4.298850in}{0.660000in}}%
\pgfpathlineto{\pgfqpoint{4.298850in}{3.278918in}}%
\pgfpathlineto{\pgfqpoint{4.292650in}{3.278918in}}%
\pgfpathlineto{\pgfqpoint{4.292650in}{0.660000in}}%
\pgfpathclose%
\pgfusepath{fill}%
\end{pgfscope}%
\begin{pgfscope}%
\pgfpathrectangle{\pgfqpoint{1.250000in}{0.660000in}}{\pgfqpoint{7.750000in}{4.620000in}}%
\pgfusepath{clip}%
\pgfsetbuttcap%
\pgfsetmiterjoin%
\definecolor{currentfill}{rgb}{0.000000,0.000000,0.000000}%
\pgfsetfillcolor{currentfill}%
\pgfsetlinewidth{0.000000pt}%
\definecolor{currentstroke}{rgb}{0.000000,0.000000,0.000000}%
\pgfsetstrokecolor{currentstroke}%
\pgfsetstrokeopacity{0.000000}%
\pgfsetdash{}{0pt}%
\pgfpathmoveto{\pgfqpoint{4.300400in}{0.660000in}}%
\pgfpathlineto{\pgfqpoint{4.306600in}{0.660000in}}%
\pgfpathlineto{\pgfqpoint{4.306600in}{3.278636in}}%
\pgfpathlineto{\pgfqpoint{4.300400in}{3.278636in}}%
\pgfpathlineto{\pgfqpoint{4.300400in}{0.660000in}}%
\pgfpathclose%
\pgfusepath{fill}%
\end{pgfscope}%
\begin{pgfscope}%
\pgfpathrectangle{\pgfqpoint{1.250000in}{0.660000in}}{\pgfqpoint{7.750000in}{4.620000in}}%
\pgfusepath{clip}%
\pgfsetbuttcap%
\pgfsetmiterjoin%
\definecolor{currentfill}{rgb}{0.000000,0.000000,0.000000}%
\pgfsetfillcolor{currentfill}%
\pgfsetlinewidth{0.000000pt}%
\definecolor{currentstroke}{rgb}{0.000000,0.000000,0.000000}%
\pgfsetstrokecolor{currentstroke}%
\pgfsetstrokeopacity{0.000000}%
\pgfsetdash{}{0pt}%
\pgfpathmoveto{\pgfqpoint{4.308150in}{0.660000in}}%
\pgfpathlineto{\pgfqpoint{4.314350in}{0.660000in}}%
\pgfpathlineto{\pgfqpoint{4.314350in}{3.278353in}}%
\pgfpathlineto{\pgfqpoint{4.308150in}{3.278353in}}%
\pgfpathlineto{\pgfqpoint{4.308150in}{0.660000in}}%
\pgfpathclose%
\pgfusepath{fill}%
\end{pgfscope}%
\begin{pgfscope}%
\pgfpathrectangle{\pgfqpoint{1.250000in}{0.660000in}}{\pgfqpoint{7.750000in}{4.620000in}}%
\pgfusepath{clip}%
\pgfsetbuttcap%
\pgfsetmiterjoin%
\definecolor{currentfill}{rgb}{0.000000,0.000000,0.000000}%
\pgfsetfillcolor{currentfill}%
\pgfsetlinewidth{0.000000pt}%
\definecolor{currentstroke}{rgb}{0.000000,0.000000,0.000000}%
\pgfsetstrokecolor{currentstroke}%
\pgfsetstrokeopacity{0.000000}%
\pgfsetdash{}{0pt}%
\pgfpathmoveto{\pgfqpoint{4.315900in}{0.660000in}}%
\pgfpathlineto{\pgfqpoint{4.322100in}{0.660000in}}%
\pgfpathlineto{\pgfqpoint{4.322100in}{3.278072in}}%
\pgfpathlineto{\pgfqpoint{4.315900in}{3.278072in}}%
\pgfpathlineto{\pgfqpoint{4.315900in}{0.660000in}}%
\pgfpathclose%
\pgfusepath{fill}%
\end{pgfscope}%
\begin{pgfscope}%
\pgfpathrectangle{\pgfqpoint{1.250000in}{0.660000in}}{\pgfqpoint{7.750000in}{4.620000in}}%
\pgfusepath{clip}%
\pgfsetbuttcap%
\pgfsetmiterjoin%
\definecolor{currentfill}{rgb}{0.000000,0.000000,0.000000}%
\pgfsetfillcolor{currentfill}%
\pgfsetlinewidth{0.000000pt}%
\definecolor{currentstroke}{rgb}{0.000000,0.000000,0.000000}%
\pgfsetstrokecolor{currentstroke}%
\pgfsetstrokeopacity{0.000000}%
\pgfsetdash{}{0pt}%
\pgfpathmoveto{\pgfqpoint{4.323650in}{0.660000in}}%
\pgfpathlineto{\pgfqpoint{4.329850in}{0.660000in}}%
\pgfpathlineto{\pgfqpoint{4.329850in}{3.277795in}}%
\pgfpathlineto{\pgfqpoint{4.323650in}{3.277795in}}%
\pgfpathlineto{\pgfqpoint{4.323650in}{0.660000in}}%
\pgfpathclose%
\pgfusepath{fill}%
\end{pgfscope}%
\begin{pgfscope}%
\pgfpathrectangle{\pgfqpoint{1.250000in}{0.660000in}}{\pgfqpoint{7.750000in}{4.620000in}}%
\pgfusepath{clip}%
\pgfsetbuttcap%
\pgfsetmiterjoin%
\definecolor{currentfill}{rgb}{0.000000,0.000000,0.000000}%
\pgfsetfillcolor{currentfill}%
\pgfsetlinewidth{0.000000pt}%
\definecolor{currentstroke}{rgb}{0.000000,0.000000,0.000000}%
\pgfsetstrokecolor{currentstroke}%
\pgfsetstrokeopacity{0.000000}%
\pgfsetdash{}{0pt}%
\pgfpathmoveto{\pgfqpoint{4.331400in}{0.660000in}}%
\pgfpathlineto{\pgfqpoint{4.337600in}{0.660000in}}%
\pgfpathlineto{\pgfqpoint{4.337600in}{3.277518in}}%
\pgfpathlineto{\pgfqpoint{4.331400in}{3.277518in}}%
\pgfpathlineto{\pgfqpoint{4.331400in}{0.660000in}}%
\pgfpathclose%
\pgfusepath{fill}%
\end{pgfscope}%
\begin{pgfscope}%
\pgfpathrectangle{\pgfqpoint{1.250000in}{0.660000in}}{\pgfqpoint{7.750000in}{4.620000in}}%
\pgfusepath{clip}%
\pgfsetbuttcap%
\pgfsetmiterjoin%
\definecolor{currentfill}{rgb}{0.000000,0.000000,0.000000}%
\pgfsetfillcolor{currentfill}%
\pgfsetlinewidth{0.000000pt}%
\definecolor{currentstroke}{rgb}{0.000000,0.000000,0.000000}%
\pgfsetstrokecolor{currentstroke}%
\pgfsetstrokeopacity{0.000000}%
\pgfsetdash{}{0pt}%
\pgfpathmoveto{\pgfqpoint{4.339150in}{0.660000in}}%
\pgfpathlineto{\pgfqpoint{4.345350in}{0.660000in}}%
\pgfpathlineto{\pgfqpoint{4.345350in}{3.277242in}}%
\pgfpathlineto{\pgfqpoint{4.339150in}{3.277242in}}%
\pgfpathlineto{\pgfqpoint{4.339150in}{0.660000in}}%
\pgfpathclose%
\pgfusepath{fill}%
\end{pgfscope}%
\begin{pgfscope}%
\pgfpathrectangle{\pgfqpoint{1.250000in}{0.660000in}}{\pgfqpoint{7.750000in}{4.620000in}}%
\pgfusepath{clip}%
\pgfsetbuttcap%
\pgfsetmiterjoin%
\definecolor{currentfill}{rgb}{0.000000,0.000000,0.000000}%
\pgfsetfillcolor{currentfill}%
\pgfsetlinewidth{0.000000pt}%
\definecolor{currentstroke}{rgb}{0.000000,0.000000,0.000000}%
\pgfsetstrokecolor{currentstroke}%
\pgfsetstrokeopacity{0.000000}%
\pgfsetdash{}{0pt}%
\pgfpathmoveto{\pgfqpoint{4.346900in}{0.660000in}}%
\pgfpathlineto{\pgfqpoint{4.353100in}{0.660000in}}%
\pgfpathlineto{\pgfqpoint{4.353100in}{3.276965in}}%
\pgfpathlineto{\pgfqpoint{4.346900in}{3.276965in}}%
\pgfpathlineto{\pgfqpoint{4.346900in}{0.660000in}}%
\pgfpathclose%
\pgfusepath{fill}%
\end{pgfscope}%
\begin{pgfscope}%
\pgfpathrectangle{\pgfqpoint{1.250000in}{0.660000in}}{\pgfqpoint{7.750000in}{4.620000in}}%
\pgfusepath{clip}%
\pgfsetbuttcap%
\pgfsetmiterjoin%
\definecolor{currentfill}{rgb}{0.000000,0.000000,0.000000}%
\pgfsetfillcolor{currentfill}%
\pgfsetlinewidth{0.000000pt}%
\definecolor{currentstroke}{rgb}{0.000000,0.000000,0.000000}%
\pgfsetstrokecolor{currentstroke}%
\pgfsetstrokeopacity{0.000000}%
\pgfsetdash{}{0pt}%
\pgfpathmoveto{\pgfqpoint{4.354650in}{0.660000in}}%
\pgfpathlineto{\pgfqpoint{4.360850in}{0.660000in}}%
\pgfpathlineto{\pgfqpoint{4.360850in}{3.276688in}}%
\pgfpathlineto{\pgfqpoint{4.354650in}{3.276688in}}%
\pgfpathlineto{\pgfqpoint{4.354650in}{0.660000in}}%
\pgfpathclose%
\pgfusepath{fill}%
\end{pgfscope}%
\begin{pgfscope}%
\pgfpathrectangle{\pgfqpoint{1.250000in}{0.660000in}}{\pgfqpoint{7.750000in}{4.620000in}}%
\pgfusepath{clip}%
\pgfsetbuttcap%
\pgfsetmiterjoin%
\definecolor{currentfill}{rgb}{0.000000,0.000000,0.000000}%
\pgfsetfillcolor{currentfill}%
\pgfsetlinewidth{0.000000pt}%
\definecolor{currentstroke}{rgb}{0.000000,0.000000,0.000000}%
\pgfsetstrokecolor{currentstroke}%
\pgfsetstrokeopacity{0.000000}%
\pgfsetdash{}{0pt}%
\pgfpathmoveto{\pgfqpoint{4.362400in}{0.660000in}}%
\pgfpathlineto{\pgfqpoint{4.368600in}{0.660000in}}%
\pgfpathlineto{\pgfqpoint{4.368600in}{3.276413in}}%
\pgfpathlineto{\pgfqpoint{4.362400in}{3.276413in}}%
\pgfpathlineto{\pgfqpoint{4.362400in}{0.660000in}}%
\pgfpathclose%
\pgfusepath{fill}%
\end{pgfscope}%
\begin{pgfscope}%
\pgfpathrectangle{\pgfqpoint{1.250000in}{0.660000in}}{\pgfqpoint{7.750000in}{4.620000in}}%
\pgfusepath{clip}%
\pgfsetbuttcap%
\pgfsetmiterjoin%
\definecolor{currentfill}{rgb}{0.000000,0.000000,0.000000}%
\pgfsetfillcolor{currentfill}%
\pgfsetlinewidth{0.000000pt}%
\definecolor{currentstroke}{rgb}{0.000000,0.000000,0.000000}%
\pgfsetstrokecolor{currentstroke}%
\pgfsetstrokeopacity{0.000000}%
\pgfsetdash{}{0pt}%
\pgfpathmoveto{\pgfqpoint{4.370150in}{0.660000in}}%
\pgfpathlineto{\pgfqpoint{4.376350in}{0.660000in}}%
\pgfpathlineto{\pgfqpoint{4.376350in}{3.276142in}}%
\pgfpathlineto{\pgfqpoint{4.370150in}{3.276142in}}%
\pgfpathlineto{\pgfqpoint{4.370150in}{0.660000in}}%
\pgfpathclose%
\pgfusepath{fill}%
\end{pgfscope}%
\begin{pgfscope}%
\pgfpathrectangle{\pgfqpoint{1.250000in}{0.660000in}}{\pgfqpoint{7.750000in}{4.620000in}}%
\pgfusepath{clip}%
\pgfsetbuttcap%
\pgfsetmiterjoin%
\definecolor{currentfill}{rgb}{0.000000,0.000000,0.000000}%
\pgfsetfillcolor{currentfill}%
\pgfsetlinewidth{0.000000pt}%
\definecolor{currentstroke}{rgb}{0.000000,0.000000,0.000000}%
\pgfsetstrokecolor{currentstroke}%
\pgfsetstrokeopacity{0.000000}%
\pgfsetdash{}{0pt}%
\pgfpathmoveto{\pgfqpoint{4.377900in}{0.660000in}}%
\pgfpathlineto{\pgfqpoint{4.384100in}{0.660000in}}%
\pgfpathlineto{\pgfqpoint{4.384100in}{3.275870in}}%
\pgfpathlineto{\pgfqpoint{4.377900in}{3.275870in}}%
\pgfpathlineto{\pgfqpoint{4.377900in}{0.660000in}}%
\pgfpathclose%
\pgfusepath{fill}%
\end{pgfscope}%
\begin{pgfscope}%
\pgfpathrectangle{\pgfqpoint{1.250000in}{0.660000in}}{\pgfqpoint{7.750000in}{4.620000in}}%
\pgfusepath{clip}%
\pgfsetbuttcap%
\pgfsetmiterjoin%
\definecolor{currentfill}{rgb}{0.000000,0.000000,0.000000}%
\pgfsetfillcolor{currentfill}%
\pgfsetlinewidth{0.000000pt}%
\definecolor{currentstroke}{rgb}{0.000000,0.000000,0.000000}%
\pgfsetstrokecolor{currentstroke}%
\pgfsetstrokeopacity{0.000000}%
\pgfsetdash{}{0pt}%
\pgfpathmoveto{\pgfqpoint{4.385650in}{0.660000in}}%
\pgfpathlineto{\pgfqpoint{4.391850in}{0.660000in}}%
\pgfpathlineto{\pgfqpoint{4.391850in}{3.275599in}}%
\pgfpathlineto{\pgfqpoint{4.385650in}{3.275599in}}%
\pgfpathlineto{\pgfqpoint{4.385650in}{0.660000in}}%
\pgfpathclose%
\pgfusepath{fill}%
\end{pgfscope}%
\begin{pgfscope}%
\pgfpathrectangle{\pgfqpoint{1.250000in}{0.660000in}}{\pgfqpoint{7.750000in}{4.620000in}}%
\pgfusepath{clip}%
\pgfsetbuttcap%
\pgfsetmiterjoin%
\definecolor{currentfill}{rgb}{0.000000,0.000000,0.000000}%
\pgfsetfillcolor{currentfill}%
\pgfsetlinewidth{0.000000pt}%
\definecolor{currentstroke}{rgb}{0.000000,0.000000,0.000000}%
\pgfsetstrokecolor{currentstroke}%
\pgfsetstrokeopacity{0.000000}%
\pgfsetdash{}{0pt}%
\pgfpathmoveto{\pgfqpoint{4.393400in}{0.660000in}}%
\pgfpathlineto{\pgfqpoint{4.399600in}{0.660000in}}%
\pgfpathlineto{\pgfqpoint{4.399600in}{3.275328in}}%
\pgfpathlineto{\pgfqpoint{4.393400in}{3.275328in}}%
\pgfpathlineto{\pgfqpoint{4.393400in}{0.660000in}}%
\pgfpathclose%
\pgfusepath{fill}%
\end{pgfscope}%
\begin{pgfscope}%
\pgfpathrectangle{\pgfqpoint{1.250000in}{0.660000in}}{\pgfqpoint{7.750000in}{4.620000in}}%
\pgfusepath{clip}%
\pgfsetbuttcap%
\pgfsetmiterjoin%
\definecolor{currentfill}{rgb}{0.000000,0.000000,0.000000}%
\pgfsetfillcolor{currentfill}%
\pgfsetlinewidth{0.000000pt}%
\definecolor{currentstroke}{rgb}{0.000000,0.000000,0.000000}%
\pgfsetstrokecolor{currentstroke}%
\pgfsetstrokeopacity{0.000000}%
\pgfsetdash{}{0pt}%
\pgfpathmoveto{\pgfqpoint{4.401150in}{0.660000in}}%
\pgfpathlineto{\pgfqpoint{4.407350in}{0.660000in}}%
\pgfpathlineto{\pgfqpoint{4.407350in}{3.275057in}}%
\pgfpathlineto{\pgfqpoint{4.401150in}{3.275057in}}%
\pgfpathlineto{\pgfqpoint{4.401150in}{0.660000in}}%
\pgfpathclose%
\pgfusepath{fill}%
\end{pgfscope}%
\begin{pgfscope}%
\pgfpathrectangle{\pgfqpoint{1.250000in}{0.660000in}}{\pgfqpoint{7.750000in}{4.620000in}}%
\pgfusepath{clip}%
\pgfsetbuttcap%
\pgfsetmiterjoin%
\definecolor{currentfill}{rgb}{0.000000,0.000000,0.000000}%
\pgfsetfillcolor{currentfill}%
\pgfsetlinewidth{0.000000pt}%
\definecolor{currentstroke}{rgb}{0.000000,0.000000,0.000000}%
\pgfsetstrokecolor{currentstroke}%
\pgfsetstrokeopacity{0.000000}%
\pgfsetdash{}{0pt}%
\pgfpathmoveto{\pgfqpoint{4.408900in}{0.660000in}}%
\pgfpathlineto{\pgfqpoint{4.415100in}{0.660000in}}%
\pgfpathlineto{\pgfqpoint{4.415100in}{3.274787in}}%
\pgfpathlineto{\pgfqpoint{4.408900in}{3.274787in}}%
\pgfpathlineto{\pgfqpoint{4.408900in}{0.660000in}}%
\pgfpathclose%
\pgfusepath{fill}%
\end{pgfscope}%
\begin{pgfscope}%
\pgfpathrectangle{\pgfqpoint{1.250000in}{0.660000in}}{\pgfqpoint{7.750000in}{4.620000in}}%
\pgfusepath{clip}%
\pgfsetbuttcap%
\pgfsetmiterjoin%
\definecolor{currentfill}{rgb}{0.000000,0.000000,0.000000}%
\pgfsetfillcolor{currentfill}%
\pgfsetlinewidth{0.000000pt}%
\definecolor{currentstroke}{rgb}{0.000000,0.000000,0.000000}%
\pgfsetstrokecolor{currentstroke}%
\pgfsetstrokeopacity{0.000000}%
\pgfsetdash{}{0pt}%
\pgfpathmoveto{\pgfqpoint{4.416650in}{0.660000in}}%
\pgfpathlineto{\pgfqpoint{4.422850in}{0.660000in}}%
\pgfpathlineto{\pgfqpoint{4.422850in}{3.274521in}}%
\pgfpathlineto{\pgfqpoint{4.416650in}{3.274521in}}%
\pgfpathlineto{\pgfqpoint{4.416650in}{0.660000in}}%
\pgfpathclose%
\pgfusepath{fill}%
\end{pgfscope}%
\begin{pgfscope}%
\pgfpathrectangle{\pgfqpoint{1.250000in}{0.660000in}}{\pgfqpoint{7.750000in}{4.620000in}}%
\pgfusepath{clip}%
\pgfsetbuttcap%
\pgfsetmiterjoin%
\definecolor{currentfill}{rgb}{0.000000,0.000000,0.000000}%
\pgfsetfillcolor{currentfill}%
\pgfsetlinewidth{0.000000pt}%
\definecolor{currentstroke}{rgb}{0.000000,0.000000,0.000000}%
\pgfsetstrokecolor{currentstroke}%
\pgfsetstrokeopacity{0.000000}%
\pgfsetdash{}{0pt}%
\pgfpathmoveto{\pgfqpoint{4.424400in}{0.660000in}}%
\pgfpathlineto{\pgfqpoint{4.430600in}{0.660000in}}%
\pgfpathlineto{\pgfqpoint{4.430600in}{3.274256in}}%
\pgfpathlineto{\pgfqpoint{4.424400in}{3.274256in}}%
\pgfpathlineto{\pgfqpoint{4.424400in}{0.660000in}}%
\pgfpathclose%
\pgfusepath{fill}%
\end{pgfscope}%
\begin{pgfscope}%
\pgfpathrectangle{\pgfqpoint{1.250000in}{0.660000in}}{\pgfqpoint{7.750000in}{4.620000in}}%
\pgfusepath{clip}%
\pgfsetbuttcap%
\pgfsetmiterjoin%
\definecolor{currentfill}{rgb}{0.000000,0.000000,0.000000}%
\pgfsetfillcolor{currentfill}%
\pgfsetlinewidth{0.000000pt}%
\definecolor{currentstroke}{rgb}{0.000000,0.000000,0.000000}%
\pgfsetstrokecolor{currentstroke}%
\pgfsetstrokeopacity{0.000000}%
\pgfsetdash{}{0pt}%
\pgfpathmoveto{\pgfqpoint{4.432150in}{0.660000in}}%
\pgfpathlineto{\pgfqpoint{4.438350in}{0.660000in}}%
\pgfpathlineto{\pgfqpoint{4.438350in}{3.273991in}}%
\pgfpathlineto{\pgfqpoint{4.432150in}{3.273991in}}%
\pgfpathlineto{\pgfqpoint{4.432150in}{0.660000in}}%
\pgfpathclose%
\pgfusepath{fill}%
\end{pgfscope}%
\begin{pgfscope}%
\pgfpathrectangle{\pgfqpoint{1.250000in}{0.660000in}}{\pgfqpoint{7.750000in}{4.620000in}}%
\pgfusepath{clip}%
\pgfsetbuttcap%
\pgfsetmiterjoin%
\definecolor{currentfill}{rgb}{0.000000,0.000000,0.000000}%
\pgfsetfillcolor{currentfill}%
\pgfsetlinewidth{0.000000pt}%
\definecolor{currentstroke}{rgb}{0.000000,0.000000,0.000000}%
\pgfsetstrokecolor{currentstroke}%
\pgfsetstrokeopacity{0.000000}%
\pgfsetdash{}{0pt}%
\pgfpathmoveto{\pgfqpoint{4.439900in}{0.660000in}}%
\pgfpathlineto{\pgfqpoint{4.446100in}{0.660000in}}%
\pgfpathlineto{\pgfqpoint{4.446100in}{3.273725in}}%
\pgfpathlineto{\pgfqpoint{4.439900in}{3.273725in}}%
\pgfpathlineto{\pgfqpoint{4.439900in}{0.660000in}}%
\pgfpathclose%
\pgfusepath{fill}%
\end{pgfscope}%
\begin{pgfscope}%
\pgfpathrectangle{\pgfqpoint{1.250000in}{0.660000in}}{\pgfqpoint{7.750000in}{4.620000in}}%
\pgfusepath{clip}%
\pgfsetbuttcap%
\pgfsetmiterjoin%
\definecolor{currentfill}{rgb}{0.000000,0.000000,0.000000}%
\pgfsetfillcolor{currentfill}%
\pgfsetlinewidth{0.000000pt}%
\definecolor{currentstroke}{rgb}{0.000000,0.000000,0.000000}%
\pgfsetstrokecolor{currentstroke}%
\pgfsetstrokeopacity{0.000000}%
\pgfsetdash{}{0pt}%
\pgfpathmoveto{\pgfqpoint{4.447650in}{0.660000in}}%
\pgfpathlineto{\pgfqpoint{4.453850in}{0.660000in}}%
\pgfpathlineto{\pgfqpoint{4.453850in}{3.273460in}}%
\pgfpathlineto{\pgfqpoint{4.447650in}{3.273460in}}%
\pgfpathlineto{\pgfqpoint{4.447650in}{0.660000in}}%
\pgfpathclose%
\pgfusepath{fill}%
\end{pgfscope}%
\begin{pgfscope}%
\pgfpathrectangle{\pgfqpoint{1.250000in}{0.660000in}}{\pgfqpoint{7.750000in}{4.620000in}}%
\pgfusepath{clip}%
\pgfsetbuttcap%
\pgfsetmiterjoin%
\definecolor{currentfill}{rgb}{0.000000,0.000000,0.000000}%
\pgfsetfillcolor{currentfill}%
\pgfsetlinewidth{0.000000pt}%
\definecolor{currentstroke}{rgb}{0.000000,0.000000,0.000000}%
\pgfsetstrokecolor{currentstroke}%
\pgfsetstrokeopacity{0.000000}%
\pgfsetdash{}{0pt}%
\pgfpathmoveto{\pgfqpoint{4.455400in}{0.660000in}}%
\pgfpathlineto{\pgfqpoint{4.461600in}{0.660000in}}%
\pgfpathlineto{\pgfqpoint{4.461600in}{3.273194in}}%
\pgfpathlineto{\pgfqpoint{4.455400in}{3.273194in}}%
\pgfpathlineto{\pgfqpoint{4.455400in}{0.660000in}}%
\pgfpathclose%
\pgfusepath{fill}%
\end{pgfscope}%
\begin{pgfscope}%
\pgfpathrectangle{\pgfqpoint{1.250000in}{0.660000in}}{\pgfqpoint{7.750000in}{4.620000in}}%
\pgfusepath{clip}%
\pgfsetbuttcap%
\pgfsetmiterjoin%
\definecolor{currentfill}{rgb}{0.000000,0.000000,0.000000}%
\pgfsetfillcolor{currentfill}%
\pgfsetlinewidth{0.000000pt}%
\definecolor{currentstroke}{rgb}{0.000000,0.000000,0.000000}%
\pgfsetstrokecolor{currentstroke}%
\pgfsetstrokeopacity{0.000000}%
\pgfsetdash{}{0pt}%
\pgfpathmoveto{\pgfqpoint{4.463150in}{0.660000in}}%
\pgfpathlineto{\pgfqpoint{4.469350in}{0.660000in}}%
\pgfpathlineto{\pgfqpoint{4.469350in}{3.272933in}}%
\pgfpathlineto{\pgfqpoint{4.463150in}{3.272933in}}%
\pgfpathlineto{\pgfqpoint{4.463150in}{0.660000in}}%
\pgfpathclose%
\pgfusepath{fill}%
\end{pgfscope}%
\begin{pgfscope}%
\pgfpathrectangle{\pgfqpoint{1.250000in}{0.660000in}}{\pgfqpoint{7.750000in}{4.620000in}}%
\pgfusepath{clip}%
\pgfsetbuttcap%
\pgfsetmiterjoin%
\definecolor{currentfill}{rgb}{0.000000,0.000000,0.000000}%
\pgfsetfillcolor{currentfill}%
\pgfsetlinewidth{0.000000pt}%
\definecolor{currentstroke}{rgb}{0.000000,0.000000,0.000000}%
\pgfsetstrokecolor{currentstroke}%
\pgfsetstrokeopacity{0.000000}%
\pgfsetdash{}{0pt}%
\pgfpathmoveto{\pgfqpoint{4.470900in}{0.660000in}}%
\pgfpathlineto{\pgfqpoint{4.477100in}{0.660000in}}%
\pgfpathlineto{\pgfqpoint{4.477100in}{3.272673in}}%
\pgfpathlineto{\pgfqpoint{4.470900in}{3.272673in}}%
\pgfpathlineto{\pgfqpoint{4.470900in}{0.660000in}}%
\pgfpathclose%
\pgfusepath{fill}%
\end{pgfscope}%
\begin{pgfscope}%
\pgfpathrectangle{\pgfqpoint{1.250000in}{0.660000in}}{\pgfqpoint{7.750000in}{4.620000in}}%
\pgfusepath{clip}%
\pgfsetbuttcap%
\pgfsetmiterjoin%
\definecolor{currentfill}{rgb}{0.000000,0.000000,0.000000}%
\pgfsetfillcolor{currentfill}%
\pgfsetlinewidth{0.000000pt}%
\definecolor{currentstroke}{rgb}{0.000000,0.000000,0.000000}%
\pgfsetstrokecolor{currentstroke}%
\pgfsetstrokeopacity{0.000000}%
\pgfsetdash{}{0pt}%
\pgfpathmoveto{\pgfqpoint{4.478650in}{0.660000in}}%
\pgfpathlineto{\pgfqpoint{4.484850in}{0.660000in}}%
\pgfpathlineto{\pgfqpoint{4.484850in}{3.272414in}}%
\pgfpathlineto{\pgfqpoint{4.478650in}{3.272414in}}%
\pgfpathlineto{\pgfqpoint{4.478650in}{0.660000in}}%
\pgfpathclose%
\pgfusepath{fill}%
\end{pgfscope}%
\begin{pgfscope}%
\pgfpathrectangle{\pgfqpoint{1.250000in}{0.660000in}}{\pgfqpoint{7.750000in}{4.620000in}}%
\pgfusepath{clip}%
\pgfsetbuttcap%
\pgfsetmiterjoin%
\definecolor{currentfill}{rgb}{0.000000,0.000000,0.000000}%
\pgfsetfillcolor{currentfill}%
\pgfsetlinewidth{0.000000pt}%
\definecolor{currentstroke}{rgb}{0.000000,0.000000,0.000000}%
\pgfsetstrokecolor{currentstroke}%
\pgfsetstrokeopacity{0.000000}%
\pgfsetdash{}{0pt}%
\pgfpathmoveto{\pgfqpoint{4.486400in}{0.660000in}}%
\pgfpathlineto{\pgfqpoint{4.492600in}{0.660000in}}%
\pgfpathlineto{\pgfqpoint{4.492600in}{3.272154in}}%
\pgfpathlineto{\pgfqpoint{4.486400in}{3.272154in}}%
\pgfpathlineto{\pgfqpoint{4.486400in}{0.660000in}}%
\pgfpathclose%
\pgfusepath{fill}%
\end{pgfscope}%
\begin{pgfscope}%
\pgfpathrectangle{\pgfqpoint{1.250000in}{0.660000in}}{\pgfqpoint{7.750000in}{4.620000in}}%
\pgfusepath{clip}%
\pgfsetbuttcap%
\pgfsetmiterjoin%
\definecolor{currentfill}{rgb}{0.000000,0.000000,0.000000}%
\pgfsetfillcolor{currentfill}%
\pgfsetlinewidth{0.000000pt}%
\definecolor{currentstroke}{rgb}{0.000000,0.000000,0.000000}%
\pgfsetstrokecolor{currentstroke}%
\pgfsetstrokeopacity{0.000000}%
\pgfsetdash{}{0pt}%
\pgfpathmoveto{\pgfqpoint{4.494150in}{0.660000in}}%
\pgfpathlineto{\pgfqpoint{4.500350in}{0.660000in}}%
\pgfpathlineto{\pgfqpoint{4.500350in}{3.271894in}}%
\pgfpathlineto{\pgfqpoint{4.494150in}{3.271894in}}%
\pgfpathlineto{\pgfqpoint{4.494150in}{0.660000in}}%
\pgfpathclose%
\pgfusepath{fill}%
\end{pgfscope}%
\begin{pgfscope}%
\pgfpathrectangle{\pgfqpoint{1.250000in}{0.660000in}}{\pgfqpoint{7.750000in}{4.620000in}}%
\pgfusepath{clip}%
\pgfsetbuttcap%
\pgfsetmiterjoin%
\definecolor{currentfill}{rgb}{0.000000,0.000000,0.000000}%
\pgfsetfillcolor{currentfill}%
\pgfsetlinewidth{0.000000pt}%
\definecolor{currentstroke}{rgb}{0.000000,0.000000,0.000000}%
\pgfsetstrokecolor{currentstroke}%
\pgfsetstrokeopacity{0.000000}%
\pgfsetdash{}{0pt}%
\pgfpathmoveto{\pgfqpoint{4.501900in}{0.660000in}}%
\pgfpathlineto{\pgfqpoint{4.508100in}{0.660000in}}%
\pgfpathlineto{\pgfqpoint{4.508100in}{3.271634in}}%
\pgfpathlineto{\pgfqpoint{4.501900in}{3.271634in}}%
\pgfpathlineto{\pgfqpoint{4.501900in}{0.660000in}}%
\pgfpathclose%
\pgfusepath{fill}%
\end{pgfscope}%
\begin{pgfscope}%
\pgfpathrectangle{\pgfqpoint{1.250000in}{0.660000in}}{\pgfqpoint{7.750000in}{4.620000in}}%
\pgfusepath{clip}%
\pgfsetbuttcap%
\pgfsetmiterjoin%
\definecolor{currentfill}{rgb}{0.000000,0.000000,0.000000}%
\pgfsetfillcolor{currentfill}%
\pgfsetlinewidth{0.000000pt}%
\definecolor{currentstroke}{rgb}{0.000000,0.000000,0.000000}%
\pgfsetstrokecolor{currentstroke}%
\pgfsetstrokeopacity{0.000000}%
\pgfsetdash{}{0pt}%
\pgfpathmoveto{\pgfqpoint{4.509650in}{0.660000in}}%
\pgfpathlineto{\pgfqpoint{4.515850in}{0.660000in}}%
\pgfpathlineto{\pgfqpoint{4.515850in}{3.271376in}}%
\pgfpathlineto{\pgfqpoint{4.509650in}{3.271376in}}%
\pgfpathlineto{\pgfqpoint{4.509650in}{0.660000in}}%
\pgfpathclose%
\pgfusepath{fill}%
\end{pgfscope}%
\begin{pgfscope}%
\pgfpathrectangle{\pgfqpoint{1.250000in}{0.660000in}}{\pgfqpoint{7.750000in}{4.620000in}}%
\pgfusepath{clip}%
\pgfsetbuttcap%
\pgfsetmiterjoin%
\definecolor{currentfill}{rgb}{0.000000,0.000000,0.000000}%
\pgfsetfillcolor{currentfill}%
\pgfsetlinewidth{0.000000pt}%
\definecolor{currentstroke}{rgb}{0.000000,0.000000,0.000000}%
\pgfsetstrokecolor{currentstroke}%
\pgfsetstrokeopacity{0.000000}%
\pgfsetdash{}{0pt}%
\pgfpathmoveto{\pgfqpoint{4.517400in}{0.660000in}}%
\pgfpathlineto{\pgfqpoint{4.523600in}{0.660000in}}%
\pgfpathlineto{\pgfqpoint{4.523600in}{3.271122in}}%
\pgfpathlineto{\pgfqpoint{4.517400in}{3.271122in}}%
\pgfpathlineto{\pgfqpoint{4.517400in}{0.660000in}}%
\pgfpathclose%
\pgfusepath{fill}%
\end{pgfscope}%
\begin{pgfscope}%
\pgfpathrectangle{\pgfqpoint{1.250000in}{0.660000in}}{\pgfqpoint{7.750000in}{4.620000in}}%
\pgfusepath{clip}%
\pgfsetbuttcap%
\pgfsetmiterjoin%
\definecolor{currentfill}{rgb}{0.000000,0.000000,0.000000}%
\pgfsetfillcolor{currentfill}%
\pgfsetlinewidth{0.000000pt}%
\definecolor{currentstroke}{rgb}{0.000000,0.000000,0.000000}%
\pgfsetstrokecolor{currentstroke}%
\pgfsetstrokeopacity{0.000000}%
\pgfsetdash{}{0pt}%
\pgfpathmoveto{\pgfqpoint{4.525150in}{0.660000in}}%
\pgfpathlineto{\pgfqpoint{4.531350in}{0.660000in}}%
\pgfpathlineto{\pgfqpoint{4.531350in}{3.270867in}}%
\pgfpathlineto{\pgfqpoint{4.525150in}{3.270867in}}%
\pgfpathlineto{\pgfqpoint{4.525150in}{0.660000in}}%
\pgfpathclose%
\pgfusepath{fill}%
\end{pgfscope}%
\begin{pgfscope}%
\pgfpathrectangle{\pgfqpoint{1.250000in}{0.660000in}}{\pgfqpoint{7.750000in}{4.620000in}}%
\pgfusepath{clip}%
\pgfsetbuttcap%
\pgfsetmiterjoin%
\definecolor{currentfill}{rgb}{0.000000,0.000000,0.000000}%
\pgfsetfillcolor{currentfill}%
\pgfsetlinewidth{0.000000pt}%
\definecolor{currentstroke}{rgb}{0.000000,0.000000,0.000000}%
\pgfsetstrokecolor{currentstroke}%
\pgfsetstrokeopacity{0.000000}%
\pgfsetdash{}{0pt}%
\pgfpathmoveto{\pgfqpoint{4.532900in}{0.660000in}}%
\pgfpathlineto{\pgfqpoint{4.539100in}{0.660000in}}%
\pgfpathlineto{\pgfqpoint{4.539100in}{3.270613in}}%
\pgfpathlineto{\pgfqpoint{4.532900in}{3.270613in}}%
\pgfpathlineto{\pgfqpoint{4.532900in}{0.660000in}}%
\pgfpathclose%
\pgfusepath{fill}%
\end{pgfscope}%
\begin{pgfscope}%
\pgfpathrectangle{\pgfqpoint{1.250000in}{0.660000in}}{\pgfqpoint{7.750000in}{4.620000in}}%
\pgfusepath{clip}%
\pgfsetbuttcap%
\pgfsetmiterjoin%
\definecolor{currentfill}{rgb}{0.000000,0.000000,0.000000}%
\pgfsetfillcolor{currentfill}%
\pgfsetlinewidth{0.000000pt}%
\definecolor{currentstroke}{rgb}{0.000000,0.000000,0.000000}%
\pgfsetstrokecolor{currentstroke}%
\pgfsetstrokeopacity{0.000000}%
\pgfsetdash{}{0pt}%
\pgfpathmoveto{\pgfqpoint{4.540650in}{0.660000in}}%
\pgfpathlineto{\pgfqpoint{4.546850in}{0.660000in}}%
\pgfpathlineto{\pgfqpoint{4.546850in}{3.270359in}}%
\pgfpathlineto{\pgfqpoint{4.540650in}{3.270359in}}%
\pgfpathlineto{\pgfqpoint{4.540650in}{0.660000in}}%
\pgfpathclose%
\pgfusepath{fill}%
\end{pgfscope}%
\begin{pgfscope}%
\pgfpathrectangle{\pgfqpoint{1.250000in}{0.660000in}}{\pgfqpoint{7.750000in}{4.620000in}}%
\pgfusepath{clip}%
\pgfsetbuttcap%
\pgfsetmiterjoin%
\definecolor{currentfill}{rgb}{0.000000,0.000000,0.000000}%
\pgfsetfillcolor{currentfill}%
\pgfsetlinewidth{0.000000pt}%
\definecolor{currentstroke}{rgb}{0.000000,0.000000,0.000000}%
\pgfsetstrokecolor{currentstroke}%
\pgfsetstrokeopacity{0.000000}%
\pgfsetdash{}{0pt}%
\pgfpathmoveto{\pgfqpoint{4.548400in}{0.660000in}}%
\pgfpathlineto{\pgfqpoint{4.554600in}{0.660000in}}%
\pgfpathlineto{\pgfqpoint{4.554600in}{3.270104in}}%
\pgfpathlineto{\pgfqpoint{4.548400in}{3.270104in}}%
\pgfpathlineto{\pgfqpoint{4.548400in}{0.660000in}}%
\pgfpathclose%
\pgfusepath{fill}%
\end{pgfscope}%
\begin{pgfscope}%
\pgfpathrectangle{\pgfqpoint{1.250000in}{0.660000in}}{\pgfqpoint{7.750000in}{4.620000in}}%
\pgfusepath{clip}%
\pgfsetbuttcap%
\pgfsetmiterjoin%
\definecolor{currentfill}{rgb}{0.000000,0.000000,0.000000}%
\pgfsetfillcolor{currentfill}%
\pgfsetlinewidth{0.000000pt}%
\definecolor{currentstroke}{rgb}{0.000000,0.000000,0.000000}%
\pgfsetstrokecolor{currentstroke}%
\pgfsetstrokeopacity{0.000000}%
\pgfsetdash{}{0pt}%
\pgfpathmoveto{\pgfqpoint{4.556150in}{0.660000in}}%
\pgfpathlineto{\pgfqpoint{4.562350in}{0.660000in}}%
\pgfpathlineto{\pgfqpoint{4.562350in}{3.269850in}}%
\pgfpathlineto{\pgfqpoint{4.556150in}{3.269850in}}%
\pgfpathlineto{\pgfqpoint{4.556150in}{0.660000in}}%
\pgfpathclose%
\pgfusepath{fill}%
\end{pgfscope}%
\begin{pgfscope}%
\pgfpathrectangle{\pgfqpoint{1.250000in}{0.660000in}}{\pgfqpoint{7.750000in}{4.620000in}}%
\pgfusepath{clip}%
\pgfsetbuttcap%
\pgfsetmiterjoin%
\definecolor{currentfill}{rgb}{0.000000,0.000000,0.000000}%
\pgfsetfillcolor{currentfill}%
\pgfsetlinewidth{0.000000pt}%
\definecolor{currentstroke}{rgb}{0.000000,0.000000,0.000000}%
\pgfsetstrokecolor{currentstroke}%
\pgfsetstrokeopacity{0.000000}%
\pgfsetdash{}{0pt}%
\pgfpathmoveto{\pgfqpoint{4.563900in}{0.660000in}}%
\pgfpathlineto{\pgfqpoint{4.570100in}{0.660000in}}%
\pgfpathlineto{\pgfqpoint{4.570100in}{3.269600in}}%
\pgfpathlineto{\pgfqpoint{4.563900in}{3.269600in}}%
\pgfpathlineto{\pgfqpoint{4.563900in}{0.660000in}}%
\pgfpathclose%
\pgfusepath{fill}%
\end{pgfscope}%
\begin{pgfscope}%
\pgfpathrectangle{\pgfqpoint{1.250000in}{0.660000in}}{\pgfqpoint{7.750000in}{4.620000in}}%
\pgfusepath{clip}%
\pgfsetbuttcap%
\pgfsetmiterjoin%
\definecolor{currentfill}{rgb}{0.000000,0.000000,0.000000}%
\pgfsetfillcolor{currentfill}%
\pgfsetlinewidth{0.000000pt}%
\definecolor{currentstroke}{rgb}{0.000000,0.000000,0.000000}%
\pgfsetstrokecolor{currentstroke}%
\pgfsetstrokeopacity{0.000000}%
\pgfsetdash{}{0pt}%
\pgfpathmoveto{\pgfqpoint{4.571650in}{0.660000in}}%
\pgfpathlineto{\pgfqpoint{4.577850in}{0.660000in}}%
\pgfpathlineto{\pgfqpoint{4.577850in}{3.269351in}}%
\pgfpathlineto{\pgfqpoint{4.571650in}{3.269351in}}%
\pgfpathlineto{\pgfqpoint{4.571650in}{0.660000in}}%
\pgfpathclose%
\pgfusepath{fill}%
\end{pgfscope}%
\begin{pgfscope}%
\pgfpathrectangle{\pgfqpoint{1.250000in}{0.660000in}}{\pgfqpoint{7.750000in}{4.620000in}}%
\pgfusepath{clip}%
\pgfsetbuttcap%
\pgfsetmiterjoin%
\definecolor{currentfill}{rgb}{0.000000,0.000000,0.000000}%
\pgfsetfillcolor{currentfill}%
\pgfsetlinewidth{0.000000pt}%
\definecolor{currentstroke}{rgb}{0.000000,0.000000,0.000000}%
\pgfsetstrokecolor{currentstroke}%
\pgfsetstrokeopacity{0.000000}%
\pgfsetdash{}{0pt}%
\pgfpathmoveto{\pgfqpoint{4.579400in}{0.660000in}}%
\pgfpathlineto{\pgfqpoint{4.585600in}{0.660000in}}%
\pgfpathlineto{\pgfqpoint{4.585600in}{3.269102in}}%
\pgfpathlineto{\pgfqpoint{4.579400in}{3.269102in}}%
\pgfpathlineto{\pgfqpoint{4.579400in}{0.660000in}}%
\pgfpathclose%
\pgfusepath{fill}%
\end{pgfscope}%
\begin{pgfscope}%
\pgfpathrectangle{\pgfqpoint{1.250000in}{0.660000in}}{\pgfqpoint{7.750000in}{4.620000in}}%
\pgfusepath{clip}%
\pgfsetbuttcap%
\pgfsetmiterjoin%
\definecolor{currentfill}{rgb}{0.000000,0.000000,0.000000}%
\pgfsetfillcolor{currentfill}%
\pgfsetlinewidth{0.000000pt}%
\definecolor{currentstroke}{rgb}{0.000000,0.000000,0.000000}%
\pgfsetstrokecolor{currentstroke}%
\pgfsetstrokeopacity{0.000000}%
\pgfsetdash{}{0pt}%
\pgfpathmoveto{\pgfqpoint{4.587150in}{0.660000in}}%
\pgfpathlineto{\pgfqpoint{4.593350in}{0.660000in}}%
\pgfpathlineto{\pgfqpoint{4.593350in}{3.268853in}}%
\pgfpathlineto{\pgfqpoint{4.587150in}{3.268853in}}%
\pgfpathlineto{\pgfqpoint{4.587150in}{0.660000in}}%
\pgfpathclose%
\pgfusepath{fill}%
\end{pgfscope}%
\begin{pgfscope}%
\pgfpathrectangle{\pgfqpoint{1.250000in}{0.660000in}}{\pgfqpoint{7.750000in}{4.620000in}}%
\pgfusepath{clip}%
\pgfsetbuttcap%
\pgfsetmiterjoin%
\definecolor{currentfill}{rgb}{0.000000,0.000000,0.000000}%
\pgfsetfillcolor{currentfill}%
\pgfsetlinewidth{0.000000pt}%
\definecolor{currentstroke}{rgb}{0.000000,0.000000,0.000000}%
\pgfsetstrokecolor{currentstroke}%
\pgfsetstrokeopacity{0.000000}%
\pgfsetdash{}{0pt}%
\pgfpathmoveto{\pgfqpoint{4.594900in}{0.660000in}}%
\pgfpathlineto{\pgfqpoint{4.601100in}{0.660000in}}%
\pgfpathlineto{\pgfqpoint{4.601100in}{3.268604in}}%
\pgfpathlineto{\pgfqpoint{4.594900in}{3.268604in}}%
\pgfpathlineto{\pgfqpoint{4.594900in}{0.660000in}}%
\pgfpathclose%
\pgfusepath{fill}%
\end{pgfscope}%
\begin{pgfscope}%
\pgfpathrectangle{\pgfqpoint{1.250000in}{0.660000in}}{\pgfqpoint{7.750000in}{4.620000in}}%
\pgfusepath{clip}%
\pgfsetbuttcap%
\pgfsetmiterjoin%
\definecolor{currentfill}{rgb}{0.000000,0.000000,0.000000}%
\pgfsetfillcolor{currentfill}%
\pgfsetlinewidth{0.000000pt}%
\definecolor{currentstroke}{rgb}{0.000000,0.000000,0.000000}%
\pgfsetstrokecolor{currentstroke}%
\pgfsetstrokeopacity{0.000000}%
\pgfsetdash{}{0pt}%
\pgfpathmoveto{\pgfqpoint{4.602650in}{0.660000in}}%
\pgfpathlineto{\pgfqpoint{4.608850in}{0.660000in}}%
\pgfpathlineto{\pgfqpoint{4.608850in}{3.268355in}}%
\pgfpathlineto{\pgfqpoint{4.602650in}{3.268355in}}%
\pgfpathlineto{\pgfqpoint{4.602650in}{0.660000in}}%
\pgfpathclose%
\pgfusepath{fill}%
\end{pgfscope}%
\begin{pgfscope}%
\pgfpathrectangle{\pgfqpoint{1.250000in}{0.660000in}}{\pgfqpoint{7.750000in}{4.620000in}}%
\pgfusepath{clip}%
\pgfsetbuttcap%
\pgfsetmiterjoin%
\definecolor{currentfill}{rgb}{0.000000,0.000000,0.000000}%
\pgfsetfillcolor{currentfill}%
\pgfsetlinewidth{0.000000pt}%
\definecolor{currentstroke}{rgb}{0.000000,0.000000,0.000000}%
\pgfsetstrokecolor{currentstroke}%
\pgfsetstrokeopacity{0.000000}%
\pgfsetdash{}{0pt}%
\pgfpathmoveto{\pgfqpoint{4.610400in}{0.660000in}}%
\pgfpathlineto{\pgfqpoint{4.616600in}{0.660000in}}%
\pgfpathlineto{\pgfqpoint{4.616600in}{3.268106in}}%
\pgfpathlineto{\pgfqpoint{4.610400in}{3.268106in}}%
\pgfpathlineto{\pgfqpoint{4.610400in}{0.660000in}}%
\pgfpathclose%
\pgfusepath{fill}%
\end{pgfscope}%
\begin{pgfscope}%
\pgfpathrectangle{\pgfqpoint{1.250000in}{0.660000in}}{\pgfqpoint{7.750000in}{4.620000in}}%
\pgfusepath{clip}%
\pgfsetbuttcap%
\pgfsetmiterjoin%
\definecolor{currentfill}{rgb}{0.000000,0.000000,0.000000}%
\pgfsetfillcolor{currentfill}%
\pgfsetlinewidth{0.000000pt}%
\definecolor{currentstroke}{rgb}{0.000000,0.000000,0.000000}%
\pgfsetstrokecolor{currentstroke}%
\pgfsetstrokeopacity{0.000000}%
\pgfsetdash{}{0pt}%
\pgfpathmoveto{\pgfqpoint{4.618150in}{0.660000in}}%
\pgfpathlineto{\pgfqpoint{4.624350in}{0.660000in}}%
\pgfpathlineto{\pgfqpoint{4.624350in}{3.267862in}}%
\pgfpathlineto{\pgfqpoint{4.618150in}{3.267862in}}%
\pgfpathlineto{\pgfqpoint{4.618150in}{0.660000in}}%
\pgfpathclose%
\pgfusepath{fill}%
\end{pgfscope}%
\begin{pgfscope}%
\pgfpathrectangle{\pgfqpoint{1.250000in}{0.660000in}}{\pgfqpoint{7.750000in}{4.620000in}}%
\pgfusepath{clip}%
\pgfsetbuttcap%
\pgfsetmiterjoin%
\definecolor{currentfill}{rgb}{0.000000,0.000000,0.000000}%
\pgfsetfillcolor{currentfill}%
\pgfsetlinewidth{0.000000pt}%
\definecolor{currentstroke}{rgb}{0.000000,0.000000,0.000000}%
\pgfsetstrokecolor{currentstroke}%
\pgfsetstrokeopacity{0.000000}%
\pgfsetdash{}{0pt}%
\pgfpathmoveto{\pgfqpoint{4.625900in}{0.660000in}}%
\pgfpathlineto{\pgfqpoint{4.632100in}{0.660000in}}%
\pgfpathlineto{\pgfqpoint{4.632100in}{3.267619in}}%
\pgfpathlineto{\pgfqpoint{4.625900in}{3.267619in}}%
\pgfpathlineto{\pgfqpoint{4.625900in}{0.660000in}}%
\pgfpathclose%
\pgfusepath{fill}%
\end{pgfscope}%
\begin{pgfscope}%
\pgfpathrectangle{\pgfqpoint{1.250000in}{0.660000in}}{\pgfqpoint{7.750000in}{4.620000in}}%
\pgfusepath{clip}%
\pgfsetbuttcap%
\pgfsetmiterjoin%
\definecolor{currentfill}{rgb}{0.000000,0.000000,0.000000}%
\pgfsetfillcolor{currentfill}%
\pgfsetlinewidth{0.000000pt}%
\definecolor{currentstroke}{rgb}{0.000000,0.000000,0.000000}%
\pgfsetstrokecolor{currentstroke}%
\pgfsetstrokeopacity{0.000000}%
\pgfsetdash{}{0pt}%
\pgfpathmoveto{\pgfqpoint{4.633650in}{0.660000in}}%
\pgfpathlineto{\pgfqpoint{4.639850in}{0.660000in}}%
\pgfpathlineto{\pgfqpoint{4.639850in}{3.267375in}}%
\pgfpathlineto{\pgfqpoint{4.633650in}{3.267375in}}%
\pgfpathlineto{\pgfqpoint{4.633650in}{0.660000in}}%
\pgfpathclose%
\pgfusepath{fill}%
\end{pgfscope}%
\begin{pgfscope}%
\pgfpathrectangle{\pgfqpoint{1.250000in}{0.660000in}}{\pgfqpoint{7.750000in}{4.620000in}}%
\pgfusepath{clip}%
\pgfsetbuttcap%
\pgfsetmiterjoin%
\definecolor{currentfill}{rgb}{0.000000,0.000000,0.000000}%
\pgfsetfillcolor{currentfill}%
\pgfsetlinewidth{0.000000pt}%
\definecolor{currentstroke}{rgb}{0.000000,0.000000,0.000000}%
\pgfsetstrokecolor{currentstroke}%
\pgfsetstrokeopacity{0.000000}%
\pgfsetdash{}{0pt}%
\pgfpathmoveto{\pgfqpoint{4.641400in}{0.660000in}}%
\pgfpathlineto{\pgfqpoint{4.647600in}{0.660000in}}%
\pgfpathlineto{\pgfqpoint{4.647600in}{3.267132in}}%
\pgfpathlineto{\pgfqpoint{4.641400in}{3.267132in}}%
\pgfpathlineto{\pgfqpoint{4.641400in}{0.660000in}}%
\pgfpathclose%
\pgfusepath{fill}%
\end{pgfscope}%
\begin{pgfscope}%
\pgfpathrectangle{\pgfqpoint{1.250000in}{0.660000in}}{\pgfqpoint{7.750000in}{4.620000in}}%
\pgfusepath{clip}%
\pgfsetbuttcap%
\pgfsetmiterjoin%
\definecolor{currentfill}{rgb}{0.000000,0.000000,0.000000}%
\pgfsetfillcolor{currentfill}%
\pgfsetlinewidth{0.000000pt}%
\definecolor{currentstroke}{rgb}{0.000000,0.000000,0.000000}%
\pgfsetstrokecolor{currentstroke}%
\pgfsetstrokeopacity{0.000000}%
\pgfsetdash{}{0pt}%
\pgfpathmoveto{\pgfqpoint{4.649150in}{0.660000in}}%
\pgfpathlineto{\pgfqpoint{4.655350in}{0.660000in}}%
\pgfpathlineto{\pgfqpoint{4.655350in}{3.266888in}}%
\pgfpathlineto{\pgfqpoint{4.649150in}{3.266888in}}%
\pgfpathlineto{\pgfqpoint{4.649150in}{0.660000in}}%
\pgfpathclose%
\pgfusepath{fill}%
\end{pgfscope}%
\begin{pgfscope}%
\pgfpathrectangle{\pgfqpoint{1.250000in}{0.660000in}}{\pgfqpoint{7.750000in}{4.620000in}}%
\pgfusepath{clip}%
\pgfsetbuttcap%
\pgfsetmiterjoin%
\definecolor{currentfill}{rgb}{0.000000,0.000000,0.000000}%
\pgfsetfillcolor{currentfill}%
\pgfsetlinewidth{0.000000pt}%
\definecolor{currentstroke}{rgb}{0.000000,0.000000,0.000000}%
\pgfsetstrokecolor{currentstroke}%
\pgfsetstrokeopacity{0.000000}%
\pgfsetdash{}{0pt}%
\pgfpathmoveto{\pgfqpoint{4.656900in}{0.660000in}}%
\pgfpathlineto{\pgfqpoint{4.663100in}{0.660000in}}%
\pgfpathlineto{\pgfqpoint{4.663100in}{3.266645in}}%
\pgfpathlineto{\pgfqpoint{4.656900in}{3.266645in}}%
\pgfpathlineto{\pgfqpoint{4.656900in}{0.660000in}}%
\pgfpathclose%
\pgfusepath{fill}%
\end{pgfscope}%
\begin{pgfscope}%
\pgfpathrectangle{\pgfqpoint{1.250000in}{0.660000in}}{\pgfqpoint{7.750000in}{4.620000in}}%
\pgfusepath{clip}%
\pgfsetbuttcap%
\pgfsetmiterjoin%
\definecolor{currentfill}{rgb}{0.000000,0.000000,0.000000}%
\pgfsetfillcolor{currentfill}%
\pgfsetlinewidth{0.000000pt}%
\definecolor{currentstroke}{rgb}{0.000000,0.000000,0.000000}%
\pgfsetstrokecolor{currentstroke}%
\pgfsetstrokeopacity{0.000000}%
\pgfsetdash{}{0pt}%
\pgfpathmoveto{\pgfqpoint{4.664650in}{0.660000in}}%
\pgfpathlineto{\pgfqpoint{4.670850in}{0.660000in}}%
\pgfpathlineto{\pgfqpoint{4.670850in}{3.266401in}}%
\pgfpathlineto{\pgfqpoint{4.664650in}{3.266401in}}%
\pgfpathlineto{\pgfqpoint{4.664650in}{0.660000in}}%
\pgfpathclose%
\pgfusepath{fill}%
\end{pgfscope}%
\begin{pgfscope}%
\pgfpathrectangle{\pgfqpoint{1.250000in}{0.660000in}}{\pgfqpoint{7.750000in}{4.620000in}}%
\pgfusepath{clip}%
\pgfsetbuttcap%
\pgfsetmiterjoin%
\definecolor{currentfill}{rgb}{0.000000,0.000000,0.000000}%
\pgfsetfillcolor{currentfill}%
\pgfsetlinewidth{0.000000pt}%
\definecolor{currentstroke}{rgb}{0.000000,0.000000,0.000000}%
\pgfsetstrokecolor{currentstroke}%
\pgfsetstrokeopacity{0.000000}%
\pgfsetdash{}{0pt}%
\pgfpathmoveto{\pgfqpoint{4.672400in}{0.660000in}}%
\pgfpathlineto{\pgfqpoint{4.678600in}{0.660000in}}%
\pgfpathlineto{\pgfqpoint{4.678600in}{3.266163in}}%
\pgfpathlineto{\pgfqpoint{4.672400in}{3.266163in}}%
\pgfpathlineto{\pgfqpoint{4.672400in}{0.660000in}}%
\pgfpathclose%
\pgfusepath{fill}%
\end{pgfscope}%
\begin{pgfscope}%
\pgfpathrectangle{\pgfqpoint{1.250000in}{0.660000in}}{\pgfqpoint{7.750000in}{4.620000in}}%
\pgfusepath{clip}%
\pgfsetbuttcap%
\pgfsetmiterjoin%
\definecolor{currentfill}{rgb}{0.000000,0.000000,0.000000}%
\pgfsetfillcolor{currentfill}%
\pgfsetlinewidth{0.000000pt}%
\definecolor{currentstroke}{rgb}{0.000000,0.000000,0.000000}%
\pgfsetstrokecolor{currentstroke}%
\pgfsetstrokeopacity{0.000000}%
\pgfsetdash{}{0pt}%
\pgfpathmoveto{\pgfqpoint{4.680150in}{0.660000in}}%
\pgfpathlineto{\pgfqpoint{4.686350in}{0.660000in}}%
\pgfpathlineto{\pgfqpoint{4.686350in}{3.265924in}}%
\pgfpathlineto{\pgfqpoint{4.680150in}{3.265924in}}%
\pgfpathlineto{\pgfqpoint{4.680150in}{0.660000in}}%
\pgfpathclose%
\pgfusepath{fill}%
\end{pgfscope}%
\begin{pgfscope}%
\pgfpathrectangle{\pgfqpoint{1.250000in}{0.660000in}}{\pgfqpoint{7.750000in}{4.620000in}}%
\pgfusepath{clip}%
\pgfsetbuttcap%
\pgfsetmiterjoin%
\definecolor{currentfill}{rgb}{0.000000,0.000000,0.000000}%
\pgfsetfillcolor{currentfill}%
\pgfsetlinewidth{0.000000pt}%
\definecolor{currentstroke}{rgb}{0.000000,0.000000,0.000000}%
\pgfsetstrokecolor{currentstroke}%
\pgfsetstrokeopacity{0.000000}%
\pgfsetdash{}{0pt}%
\pgfpathmoveto{\pgfqpoint{4.687900in}{0.660000in}}%
\pgfpathlineto{\pgfqpoint{4.694100in}{0.660000in}}%
\pgfpathlineto{\pgfqpoint{4.694100in}{3.265686in}}%
\pgfpathlineto{\pgfqpoint{4.687900in}{3.265686in}}%
\pgfpathlineto{\pgfqpoint{4.687900in}{0.660000in}}%
\pgfpathclose%
\pgfusepath{fill}%
\end{pgfscope}%
\begin{pgfscope}%
\pgfpathrectangle{\pgfqpoint{1.250000in}{0.660000in}}{\pgfqpoint{7.750000in}{4.620000in}}%
\pgfusepath{clip}%
\pgfsetbuttcap%
\pgfsetmiterjoin%
\definecolor{currentfill}{rgb}{0.000000,0.000000,0.000000}%
\pgfsetfillcolor{currentfill}%
\pgfsetlinewidth{0.000000pt}%
\definecolor{currentstroke}{rgb}{0.000000,0.000000,0.000000}%
\pgfsetstrokecolor{currentstroke}%
\pgfsetstrokeopacity{0.000000}%
\pgfsetdash{}{0pt}%
\pgfpathmoveto{\pgfqpoint{4.695650in}{0.660000in}}%
\pgfpathlineto{\pgfqpoint{4.701850in}{0.660000in}}%
\pgfpathlineto{\pgfqpoint{4.701850in}{3.265448in}}%
\pgfpathlineto{\pgfqpoint{4.695650in}{3.265448in}}%
\pgfpathlineto{\pgfqpoint{4.695650in}{0.660000in}}%
\pgfpathclose%
\pgfusepath{fill}%
\end{pgfscope}%
\begin{pgfscope}%
\pgfpathrectangle{\pgfqpoint{1.250000in}{0.660000in}}{\pgfqpoint{7.750000in}{4.620000in}}%
\pgfusepath{clip}%
\pgfsetbuttcap%
\pgfsetmiterjoin%
\definecolor{currentfill}{rgb}{0.000000,0.000000,0.000000}%
\pgfsetfillcolor{currentfill}%
\pgfsetlinewidth{0.000000pt}%
\definecolor{currentstroke}{rgb}{0.000000,0.000000,0.000000}%
\pgfsetstrokecolor{currentstroke}%
\pgfsetstrokeopacity{0.000000}%
\pgfsetdash{}{0pt}%
\pgfpathmoveto{\pgfqpoint{4.703400in}{0.660000in}}%
\pgfpathlineto{\pgfqpoint{4.709600in}{0.660000in}}%
\pgfpathlineto{\pgfqpoint{4.709600in}{3.265210in}}%
\pgfpathlineto{\pgfqpoint{4.703400in}{3.265210in}}%
\pgfpathlineto{\pgfqpoint{4.703400in}{0.660000in}}%
\pgfpathclose%
\pgfusepath{fill}%
\end{pgfscope}%
\begin{pgfscope}%
\pgfpathrectangle{\pgfqpoint{1.250000in}{0.660000in}}{\pgfqpoint{7.750000in}{4.620000in}}%
\pgfusepath{clip}%
\pgfsetbuttcap%
\pgfsetmiterjoin%
\definecolor{currentfill}{rgb}{0.000000,0.000000,0.000000}%
\pgfsetfillcolor{currentfill}%
\pgfsetlinewidth{0.000000pt}%
\definecolor{currentstroke}{rgb}{0.000000,0.000000,0.000000}%
\pgfsetstrokecolor{currentstroke}%
\pgfsetstrokeopacity{0.000000}%
\pgfsetdash{}{0pt}%
\pgfpathmoveto{\pgfqpoint{4.711150in}{0.660000in}}%
\pgfpathlineto{\pgfqpoint{4.717350in}{0.660000in}}%
\pgfpathlineto{\pgfqpoint{4.717350in}{3.264972in}}%
\pgfpathlineto{\pgfqpoint{4.711150in}{3.264972in}}%
\pgfpathlineto{\pgfqpoint{4.711150in}{0.660000in}}%
\pgfpathclose%
\pgfusepath{fill}%
\end{pgfscope}%
\begin{pgfscope}%
\pgfpathrectangle{\pgfqpoint{1.250000in}{0.660000in}}{\pgfqpoint{7.750000in}{4.620000in}}%
\pgfusepath{clip}%
\pgfsetbuttcap%
\pgfsetmiterjoin%
\definecolor{currentfill}{rgb}{0.000000,0.000000,0.000000}%
\pgfsetfillcolor{currentfill}%
\pgfsetlinewidth{0.000000pt}%
\definecolor{currentstroke}{rgb}{0.000000,0.000000,0.000000}%
\pgfsetstrokecolor{currentstroke}%
\pgfsetstrokeopacity{0.000000}%
\pgfsetdash{}{0pt}%
\pgfpathmoveto{\pgfqpoint{4.718900in}{0.660000in}}%
\pgfpathlineto{\pgfqpoint{4.725100in}{0.660000in}}%
\pgfpathlineto{\pgfqpoint{4.725100in}{3.264734in}}%
\pgfpathlineto{\pgfqpoint{4.718900in}{3.264734in}}%
\pgfpathlineto{\pgfqpoint{4.718900in}{0.660000in}}%
\pgfpathclose%
\pgfusepath{fill}%
\end{pgfscope}%
\begin{pgfscope}%
\pgfpathrectangle{\pgfqpoint{1.250000in}{0.660000in}}{\pgfqpoint{7.750000in}{4.620000in}}%
\pgfusepath{clip}%
\pgfsetbuttcap%
\pgfsetmiterjoin%
\definecolor{currentfill}{rgb}{0.000000,0.000000,0.000000}%
\pgfsetfillcolor{currentfill}%
\pgfsetlinewidth{0.000000pt}%
\definecolor{currentstroke}{rgb}{0.000000,0.000000,0.000000}%
\pgfsetstrokecolor{currentstroke}%
\pgfsetstrokeopacity{0.000000}%
\pgfsetdash{}{0pt}%
\pgfpathmoveto{\pgfqpoint{4.726650in}{0.660000in}}%
\pgfpathlineto{\pgfqpoint{4.732850in}{0.660000in}}%
\pgfpathlineto{\pgfqpoint{4.732850in}{3.264499in}}%
\pgfpathlineto{\pgfqpoint{4.726650in}{3.264499in}}%
\pgfpathlineto{\pgfqpoint{4.726650in}{0.660000in}}%
\pgfpathclose%
\pgfusepath{fill}%
\end{pgfscope}%
\begin{pgfscope}%
\pgfpathrectangle{\pgfqpoint{1.250000in}{0.660000in}}{\pgfqpoint{7.750000in}{4.620000in}}%
\pgfusepath{clip}%
\pgfsetbuttcap%
\pgfsetmiterjoin%
\definecolor{currentfill}{rgb}{0.000000,0.000000,0.000000}%
\pgfsetfillcolor{currentfill}%
\pgfsetlinewidth{0.000000pt}%
\definecolor{currentstroke}{rgb}{0.000000,0.000000,0.000000}%
\pgfsetstrokecolor{currentstroke}%
\pgfsetstrokeopacity{0.000000}%
\pgfsetdash{}{0pt}%
\pgfpathmoveto{\pgfqpoint{4.734400in}{0.660000in}}%
\pgfpathlineto{\pgfqpoint{4.740600in}{0.660000in}}%
\pgfpathlineto{\pgfqpoint{4.740600in}{3.264266in}}%
\pgfpathlineto{\pgfqpoint{4.734400in}{3.264266in}}%
\pgfpathlineto{\pgfqpoint{4.734400in}{0.660000in}}%
\pgfpathclose%
\pgfusepath{fill}%
\end{pgfscope}%
\begin{pgfscope}%
\pgfpathrectangle{\pgfqpoint{1.250000in}{0.660000in}}{\pgfqpoint{7.750000in}{4.620000in}}%
\pgfusepath{clip}%
\pgfsetbuttcap%
\pgfsetmiterjoin%
\definecolor{currentfill}{rgb}{0.000000,0.000000,0.000000}%
\pgfsetfillcolor{currentfill}%
\pgfsetlinewidth{0.000000pt}%
\definecolor{currentstroke}{rgb}{0.000000,0.000000,0.000000}%
\pgfsetstrokecolor{currentstroke}%
\pgfsetstrokeopacity{0.000000}%
\pgfsetdash{}{0pt}%
\pgfpathmoveto{\pgfqpoint{4.742150in}{0.660000in}}%
\pgfpathlineto{\pgfqpoint{4.748350in}{0.660000in}}%
\pgfpathlineto{\pgfqpoint{4.748350in}{3.264033in}}%
\pgfpathlineto{\pgfqpoint{4.742150in}{3.264033in}}%
\pgfpathlineto{\pgfqpoint{4.742150in}{0.660000in}}%
\pgfpathclose%
\pgfusepath{fill}%
\end{pgfscope}%
\begin{pgfscope}%
\pgfpathrectangle{\pgfqpoint{1.250000in}{0.660000in}}{\pgfqpoint{7.750000in}{4.620000in}}%
\pgfusepath{clip}%
\pgfsetbuttcap%
\pgfsetmiterjoin%
\definecolor{currentfill}{rgb}{0.000000,0.000000,0.000000}%
\pgfsetfillcolor{currentfill}%
\pgfsetlinewidth{0.000000pt}%
\definecolor{currentstroke}{rgb}{0.000000,0.000000,0.000000}%
\pgfsetstrokecolor{currentstroke}%
\pgfsetstrokeopacity{0.000000}%
\pgfsetdash{}{0pt}%
\pgfpathmoveto{\pgfqpoint{4.749900in}{0.660000in}}%
\pgfpathlineto{\pgfqpoint{4.756100in}{0.660000in}}%
\pgfpathlineto{\pgfqpoint{4.756100in}{3.263800in}}%
\pgfpathlineto{\pgfqpoint{4.749900in}{3.263800in}}%
\pgfpathlineto{\pgfqpoint{4.749900in}{0.660000in}}%
\pgfpathclose%
\pgfusepath{fill}%
\end{pgfscope}%
\begin{pgfscope}%
\pgfpathrectangle{\pgfqpoint{1.250000in}{0.660000in}}{\pgfqpoint{7.750000in}{4.620000in}}%
\pgfusepath{clip}%
\pgfsetbuttcap%
\pgfsetmiterjoin%
\definecolor{currentfill}{rgb}{0.000000,0.000000,0.000000}%
\pgfsetfillcolor{currentfill}%
\pgfsetlinewidth{0.000000pt}%
\definecolor{currentstroke}{rgb}{0.000000,0.000000,0.000000}%
\pgfsetstrokecolor{currentstroke}%
\pgfsetstrokeopacity{0.000000}%
\pgfsetdash{}{0pt}%
\pgfpathmoveto{\pgfqpoint{4.757650in}{0.660000in}}%
\pgfpathlineto{\pgfqpoint{4.763850in}{0.660000in}}%
\pgfpathlineto{\pgfqpoint{4.763850in}{3.263568in}}%
\pgfpathlineto{\pgfqpoint{4.757650in}{3.263568in}}%
\pgfpathlineto{\pgfqpoint{4.757650in}{0.660000in}}%
\pgfpathclose%
\pgfusepath{fill}%
\end{pgfscope}%
\begin{pgfscope}%
\pgfpathrectangle{\pgfqpoint{1.250000in}{0.660000in}}{\pgfqpoint{7.750000in}{4.620000in}}%
\pgfusepath{clip}%
\pgfsetbuttcap%
\pgfsetmiterjoin%
\definecolor{currentfill}{rgb}{0.000000,0.000000,0.000000}%
\pgfsetfillcolor{currentfill}%
\pgfsetlinewidth{0.000000pt}%
\definecolor{currentstroke}{rgb}{0.000000,0.000000,0.000000}%
\pgfsetstrokecolor{currentstroke}%
\pgfsetstrokeopacity{0.000000}%
\pgfsetdash{}{0pt}%
\pgfpathmoveto{\pgfqpoint{4.765400in}{0.660000in}}%
\pgfpathlineto{\pgfqpoint{4.771600in}{0.660000in}}%
\pgfpathlineto{\pgfqpoint{4.771600in}{3.263335in}}%
\pgfpathlineto{\pgfqpoint{4.765400in}{3.263335in}}%
\pgfpathlineto{\pgfqpoint{4.765400in}{0.660000in}}%
\pgfpathclose%
\pgfusepath{fill}%
\end{pgfscope}%
\begin{pgfscope}%
\pgfpathrectangle{\pgfqpoint{1.250000in}{0.660000in}}{\pgfqpoint{7.750000in}{4.620000in}}%
\pgfusepath{clip}%
\pgfsetbuttcap%
\pgfsetmiterjoin%
\definecolor{currentfill}{rgb}{0.000000,0.000000,0.000000}%
\pgfsetfillcolor{currentfill}%
\pgfsetlinewidth{0.000000pt}%
\definecolor{currentstroke}{rgb}{0.000000,0.000000,0.000000}%
\pgfsetstrokecolor{currentstroke}%
\pgfsetstrokeopacity{0.000000}%
\pgfsetdash{}{0pt}%
\pgfpathmoveto{\pgfqpoint{4.773150in}{0.660000in}}%
\pgfpathlineto{\pgfqpoint{4.779350in}{0.660000in}}%
\pgfpathlineto{\pgfqpoint{4.779350in}{3.263102in}}%
\pgfpathlineto{\pgfqpoint{4.773150in}{3.263102in}}%
\pgfpathlineto{\pgfqpoint{4.773150in}{0.660000in}}%
\pgfpathclose%
\pgfusepath{fill}%
\end{pgfscope}%
\begin{pgfscope}%
\pgfpathrectangle{\pgfqpoint{1.250000in}{0.660000in}}{\pgfqpoint{7.750000in}{4.620000in}}%
\pgfusepath{clip}%
\pgfsetbuttcap%
\pgfsetmiterjoin%
\definecolor{currentfill}{rgb}{0.000000,0.000000,0.000000}%
\pgfsetfillcolor{currentfill}%
\pgfsetlinewidth{0.000000pt}%
\definecolor{currentstroke}{rgb}{0.000000,0.000000,0.000000}%
\pgfsetstrokecolor{currentstroke}%
\pgfsetstrokeopacity{0.000000}%
\pgfsetdash{}{0pt}%
\pgfpathmoveto{\pgfqpoint{4.780900in}{0.660000in}}%
\pgfpathlineto{\pgfqpoint{4.787100in}{0.660000in}}%
\pgfpathlineto{\pgfqpoint{4.787100in}{3.262871in}}%
\pgfpathlineto{\pgfqpoint{4.780900in}{3.262871in}}%
\pgfpathlineto{\pgfqpoint{4.780900in}{0.660000in}}%
\pgfpathclose%
\pgfusepath{fill}%
\end{pgfscope}%
\begin{pgfscope}%
\pgfpathrectangle{\pgfqpoint{1.250000in}{0.660000in}}{\pgfqpoint{7.750000in}{4.620000in}}%
\pgfusepath{clip}%
\pgfsetbuttcap%
\pgfsetmiterjoin%
\definecolor{currentfill}{rgb}{0.000000,0.000000,0.000000}%
\pgfsetfillcolor{currentfill}%
\pgfsetlinewidth{0.000000pt}%
\definecolor{currentstroke}{rgb}{0.000000,0.000000,0.000000}%
\pgfsetstrokecolor{currentstroke}%
\pgfsetstrokeopacity{0.000000}%
\pgfsetdash{}{0pt}%
\pgfpathmoveto{\pgfqpoint{4.788650in}{0.660000in}}%
\pgfpathlineto{\pgfqpoint{4.794850in}{0.660000in}}%
\pgfpathlineto{\pgfqpoint{4.794850in}{3.262643in}}%
\pgfpathlineto{\pgfqpoint{4.788650in}{3.262643in}}%
\pgfpathlineto{\pgfqpoint{4.788650in}{0.660000in}}%
\pgfpathclose%
\pgfusepath{fill}%
\end{pgfscope}%
\begin{pgfscope}%
\pgfpathrectangle{\pgfqpoint{1.250000in}{0.660000in}}{\pgfqpoint{7.750000in}{4.620000in}}%
\pgfusepath{clip}%
\pgfsetbuttcap%
\pgfsetmiterjoin%
\definecolor{currentfill}{rgb}{0.000000,0.000000,0.000000}%
\pgfsetfillcolor{currentfill}%
\pgfsetlinewidth{0.000000pt}%
\definecolor{currentstroke}{rgb}{0.000000,0.000000,0.000000}%
\pgfsetstrokecolor{currentstroke}%
\pgfsetstrokeopacity{0.000000}%
\pgfsetdash{}{0pt}%
\pgfpathmoveto{\pgfqpoint{4.796400in}{0.660000in}}%
\pgfpathlineto{\pgfqpoint{4.802600in}{0.660000in}}%
\pgfpathlineto{\pgfqpoint{4.802600in}{3.262415in}}%
\pgfpathlineto{\pgfqpoint{4.796400in}{3.262415in}}%
\pgfpathlineto{\pgfqpoint{4.796400in}{0.660000in}}%
\pgfpathclose%
\pgfusepath{fill}%
\end{pgfscope}%
\begin{pgfscope}%
\pgfpathrectangle{\pgfqpoint{1.250000in}{0.660000in}}{\pgfqpoint{7.750000in}{4.620000in}}%
\pgfusepath{clip}%
\pgfsetbuttcap%
\pgfsetmiterjoin%
\definecolor{currentfill}{rgb}{0.000000,0.000000,0.000000}%
\pgfsetfillcolor{currentfill}%
\pgfsetlinewidth{0.000000pt}%
\definecolor{currentstroke}{rgb}{0.000000,0.000000,0.000000}%
\pgfsetstrokecolor{currentstroke}%
\pgfsetstrokeopacity{0.000000}%
\pgfsetdash{}{0pt}%
\pgfpathmoveto{\pgfqpoint{4.804150in}{0.660000in}}%
\pgfpathlineto{\pgfqpoint{4.810350in}{0.660000in}}%
\pgfpathlineto{\pgfqpoint{4.810350in}{3.262188in}}%
\pgfpathlineto{\pgfqpoint{4.804150in}{3.262188in}}%
\pgfpathlineto{\pgfqpoint{4.804150in}{0.660000in}}%
\pgfpathclose%
\pgfusepath{fill}%
\end{pgfscope}%
\begin{pgfscope}%
\pgfpathrectangle{\pgfqpoint{1.250000in}{0.660000in}}{\pgfqpoint{7.750000in}{4.620000in}}%
\pgfusepath{clip}%
\pgfsetbuttcap%
\pgfsetmiterjoin%
\definecolor{currentfill}{rgb}{0.000000,0.000000,0.000000}%
\pgfsetfillcolor{currentfill}%
\pgfsetlinewidth{0.000000pt}%
\definecolor{currentstroke}{rgb}{0.000000,0.000000,0.000000}%
\pgfsetstrokecolor{currentstroke}%
\pgfsetstrokeopacity{0.000000}%
\pgfsetdash{}{0pt}%
\pgfpathmoveto{\pgfqpoint{4.811900in}{0.660000in}}%
\pgfpathlineto{\pgfqpoint{4.818100in}{0.660000in}}%
\pgfpathlineto{\pgfqpoint{4.818100in}{3.261960in}}%
\pgfpathlineto{\pgfqpoint{4.811900in}{3.261960in}}%
\pgfpathlineto{\pgfqpoint{4.811900in}{0.660000in}}%
\pgfpathclose%
\pgfusepath{fill}%
\end{pgfscope}%
\begin{pgfscope}%
\pgfpathrectangle{\pgfqpoint{1.250000in}{0.660000in}}{\pgfqpoint{7.750000in}{4.620000in}}%
\pgfusepath{clip}%
\pgfsetbuttcap%
\pgfsetmiterjoin%
\definecolor{currentfill}{rgb}{0.000000,0.000000,0.000000}%
\pgfsetfillcolor{currentfill}%
\pgfsetlinewidth{0.000000pt}%
\definecolor{currentstroke}{rgb}{0.000000,0.000000,0.000000}%
\pgfsetstrokecolor{currentstroke}%
\pgfsetstrokeopacity{0.000000}%
\pgfsetdash{}{0pt}%
\pgfpathmoveto{\pgfqpoint{4.819650in}{0.660000in}}%
\pgfpathlineto{\pgfqpoint{4.825850in}{0.660000in}}%
\pgfpathlineto{\pgfqpoint{4.825850in}{3.261732in}}%
\pgfpathlineto{\pgfqpoint{4.819650in}{3.261732in}}%
\pgfpathlineto{\pgfqpoint{4.819650in}{0.660000in}}%
\pgfpathclose%
\pgfusepath{fill}%
\end{pgfscope}%
\begin{pgfscope}%
\pgfpathrectangle{\pgfqpoint{1.250000in}{0.660000in}}{\pgfqpoint{7.750000in}{4.620000in}}%
\pgfusepath{clip}%
\pgfsetbuttcap%
\pgfsetmiterjoin%
\definecolor{currentfill}{rgb}{0.000000,0.000000,0.000000}%
\pgfsetfillcolor{currentfill}%
\pgfsetlinewidth{0.000000pt}%
\definecolor{currentstroke}{rgb}{0.000000,0.000000,0.000000}%
\pgfsetstrokecolor{currentstroke}%
\pgfsetstrokeopacity{0.000000}%
\pgfsetdash{}{0pt}%
\pgfpathmoveto{\pgfqpoint{4.827400in}{0.660000in}}%
\pgfpathlineto{\pgfqpoint{4.833600in}{0.660000in}}%
\pgfpathlineto{\pgfqpoint{4.833600in}{3.261505in}}%
\pgfpathlineto{\pgfqpoint{4.827400in}{3.261505in}}%
\pgfpathlineto{\pgfqpoint{4.827400in}{0.660000in}}%
\pgfpathclose%
\pgfusepath{fill}%
\end{pgfscope}%
\begin{pgfscope}%
\pgfpathrectangle{\pgfqpoint{1.250000in}{0.660000in}}{\pgfqpoint{7.750000in}{4.620000in}}%
\pgfusepath{clip}%
\pgfsetbuttcap%
\pgfsetmiterjoin%
\definecolor{currentfill}{rgb}{0.000000,0.000000,0.000000}%
\pgfsetfillcolor{currentfill}%
\pgfsetlinewidth{0.000000pt}%
\definecolor{currentstroke}{rgb}{0.000000,0.000000,0.000000}%
\pgfsetstrokecolor{currentstroke}%
\pgfsetstrokeopacity{0.000000}%
\pgfsetdash{}{0pt}%
\pgfpathmoveto{\pgfqpoint{4.835150in}{0.660000in}}%
\pgfpathlineto{\pgfqpoint{4.841350in}{0.660000in}}%
\pgfpathlineto{\pgfqpoint{4.841350in}{3.261277in}}%
\pgfpathlineto{\pgfqpoint{4.835150in}{3.261277in}}%
\pgfpathlineto{\pgfqpoint{4.835150in}{0.660000in}}%
\pgfpathclose%
\pgfusepath{fill}%
\end{pgfscope}%
\begin{pgfscope}%
\pgfpathrectangle{\pgfqpoint{1.250000in}{0.660000in}}{\pgfqpoint{7.750000in}{4.620000in}}%
\pgfusepath{clip}%
\pgfsetbuttcap%
\pgfsetmiterjoin%
\definecolor{currentfill}{rgb}{0.000000,0.000000,0.000000}%
\pgfsetfillcolor{currentfill}%
\pgfsetlinewidth{0.000000pt}%
\definecolor{currentstroke}{rgb}{0.000000,0.000000,0.000000}%
\pgfsetstrokecolor{currentstroke}%
\pgfsetstrokeopacity{0.000000}%
\pgfsetdash{}{0pt}%
\pgfpathmoveto{\pgfqpoint{4.842900in}{0.660000in}}%
\pgfpathlineto{\pgfqpoint{4.849100in}{0.660000in}}%
\pgfpathlineto{\pgfqpoint{4.849100in}{3.261053in}}%
\pgfpathlineto{\pgfqpoint{4.842900in}{3.261053in}}%
\pgfpathlineto{\pgfqpoint{4.842900in}{0.660000in}}%
\pgfpathclose%
\pgfusepath{fill}%
\end{pgfscope}%
\begin{pgfscope}%
\pgfpathrectangle{\pgfqpoint{1.250000in}{0.660000in}}{\pgfqpoint{7.750000in}{4.620000in}}%
\pgfusepath{clip}%
\pgfsetbuttcap%
\pgfsetmiterjoin%
\definecolor{currentfill}{rgb}{0.000000,0.000000,0.000000}%
\pgfsetfillcolor{currentfill}%
\pgfsetlinewidth{0.000000pt}%
\definecolor{currentstroke}{rgb}{0.000000,0.000000,0.000000}%
\pgfsetstrokecolor{currentstroke}%
\pgfsetstrokeopacity{0.000000}%
\pgfsetdash{}{0pt}%
\pgfpathmoveto{\pgfqpoint{4.850650in}{0.660000in}}%
\pgfpathlineto{\pgfqpoint{4.856850in}{0.660000in}}%
\pgfpathlineto{\pgfqpoint{4.856850in}{3.260830in}}%
\pgfpathlineto{\pgfqpoint{4.850650in}{3.260830in}}%
\pgfpathlineto{\pgfqpoint{4.850650in}{0.660000in}}%
\pgfpathclose%
\pgfusepath{fill}%
\end{pgfscope}%
\begin{pgfscope}%
\pgfpathrectangle{\pgfqpoint{1.250000in}{0.660000in}}{\pgfqpoint{7.750000in}{4.620000in}}%
\pgfusepath{clip}%
\pgfsetbuttcap%
\pgfsetmiterjoin%
\definecolor{currentfill}{rgb}{0.000000,0.000000,0.000000}%
\pgfsetfillcolor{currentfill}%
\pgfsetlinewidth{0.000000pt}%
\definecolor{currentstroke}{rgb}{0.000000,0.000000,0.000000}%
\pgfsetstrokecolor{currentstroke}%
\pgfsetstrokeopacity{0.000000}%
\pgfsetdash{}{0pt}%
\pgfpathmoveto{\pgfqpoint{4.858400in}{0.660000in}}%
\pgfpathlineto{\pgfqpoint{4.864600in}{0.660000in}}%
\pgfpathlineto{\pgfqpoint{4.864600in}{3.260608in}}%
\pgfpathlineto{\pgfqpoint{4.858400in}{3.260608in}}%
\pgfpathlineto{\pgfqpoint{4.858400in}{0.660000in}}%
\pgfpathclose%
\pgfusepath{fill}%
\end{pgfscope}%
\begin{pgfscope}%
\pgfpathrectangle{\pgfqpoint{1.250000in}{0.660000in}}{\pgfqpoint{7.750000in}{4.620000in}}%
\pgfusepath{clip}%
\pgfsetbuttcap%
\pgfsetmiterjoin%
\definecolor{currentfill}{rgb}{0.000000,0.000000,0.000000}%
\pgfsetfillcolor{currentfill}%
\pgfsetlinewidth{0.000000pt}%
\definecolor{currentstroke}{rgb}{0.000000,0.000000,0.000000}%
\pgfsetstrokecolor{currentstroke}%
\pgfsetstrokeopacity{0.000000}%
\pgfsetdash{}{0pt}%
\pgfpathmoveto{\pgfqpoint{4.866150in}{0.660000in}}%
\pgfpathlineto{\pgfqpoint{4.872350in}{0.660000in}}%
\pgfpathlineto{\pgfqpoint{4.872350in}{3.260385in}}%
\pgfpathlineto{\pgfqpoint{4.866150in}{3.260385in}}%
\pgfpathlineto{\pgfqpoint{4.866150in}{0.660000in}}%
\pgfpathclose%
\pgfusepath{fill}%
\end{pgfscope}%
\begin{pgfscope}%
\pgfpathrectangle{\pgfqpoint{1.250000in}{0.660000in}}{\pgfqpoint{7.750000in}{4.620000in}}%
\pgfusepath{clip}%
\pgfsetbuttcap%
\pgfsetmiterjoin%
\definecolor{currentfill}{rgb}{0.000000,0.000000,0.000000}%
\pgfsetfillcolor{currentfill}%
\pgfsetlinewidth{0.000000pt}%
\definecolor{currentstroke}{rgb}{0.000000,0.000000,0.000000}%
\pgfsetstrokecolor{currentstroke}%
\pgfsetstrokeopacity{0.000000}%
\pgfsetdash{}{0pt}%
\pgfpathmoveto{\pgfqpoint{4.873900in}{0.660000in}}%
\pgfpathlineto{\pgfqpoint{4.880100in}{0.660000in}}%
\pgfpathlineto{\pgfqpoint{4.880100in}{3.260163in}}%
\pgfpathlineto{\pgfqpoint{4.873900in}{3.260163in}}%
\pgfpathlineto{\pgfqpoint{4.873900in}{0.660000in}}%
\pgfpathclose%
\pgfusepath{fill}%
\end{pgfscope}%
\begin{pgfscope}%
\pgfpathrectangle{\pgfqpoint{1.250000in}{0.660000in}}{\pgfqpoint{7.750000in}{4.620000in}}%
\pgfusepath{clip}%
\pgfsetbuttcap%
\pgfsetmiterjoin%
\definecolor{currentfill}{rgb}{0.000000,0.000000,0.000000}%
\pgfsetfillcolor{currentfill}%
\pgfsetlinewidth{0.000000pt}%
\definecolor{currentstroke}{rgb}{0.000000,0.000000,0.000000}%
\pgfsetstrokecolor{currentstroke}%
\pgfsetstrokeopacity{0.000000}%
\pgfsetdash{}{0pt}%
\pgfpathmoveto{\pgfqpoint{4.881650in}{0.660000in}}%
\pgfpathlineto{\pgfqpoint{4.887850in}{0.660000in}}%
\pgfpathlineto{\pgfqpoint{4.887850in}{3.259940in}}%
\pgfpathlineto{\pgfqpoint{4.881650in}{3.259940in}}%
\pgfpathlineto{\pgfqpoint{4.881650in}{0.660000in}}%
\pgfpathclose%
\pgfusepath{fill}%
\end{pgfscope}%
\begin{pgfscope}%
\pgfpathrectangle{\pgfqpoint{1.250000in}{0.660000in}}{\pgfqpoint{7.750000in}{4.620000in}}%
\pgfusepath{clip}%
\pgfsetbuttcap%
\pgfsetmiterjoin%
\definecolor{currentfill}{rgb}{0.000000,0.000000,0.000000}%
\pgfsetfillcolor{currentfill}%
\pgfsetlinewidth{0.000000pt}%
\definecolor{currentstroke}{rgb}{0.000000,0.000000,0.000000}%
\pgfsetstrokecolor{currentstroke}%
\pgfsetstrokeopacity{0.000000}%
\pgfsetdash{}{0pt}%
\pgfpathmoveto{\pgfqpoint{4.889400in}{0.660000in}}%
\pgfpathlineto{\pgfqpoint{4.895600in}{0.660000in}}%
\pgfpathlineto{\pgfqpoint{4.895600in}{3.259718in}}%
\pgfpathlineto{\pgfqpoint{4.889400in}{3.259718in}}%
\pgfpathlineto{\pgfqpoint{4.889400in}{0.660000in}}%
\pgfpathclose%
\pgfusepath{fill}%
\end{pgfscope}%
\begin{pgfscope}%
\pgfpathrectangle{\pgfqpoint{1.250000in}{0.660000in}}{\pgfqpoint{7.750000in}{4.620000in}}%
\pgfusepath{clip}%
\pgfsetbuttcap%
\pgfsetmiterjoin%
\definecolor{currentfill}{rgb}{0.000000,0.000000,0.000000}%
\pgfsetfillcolor{currentfill}%
\pgfsetlinewidth{0.000000pt}%
\definecolor{currentstroke}{rgb}{0.000000,0.000000,0.000000}%
\pgfsetstrokecolor{currentstroke}%
\pgfsetstrokeopacity{0.000000}%
\pgfsetdash{}{0pt}%
\pgfpathmoveto{\pgfqpoint{4.897150in}{0.660000in}}%
\pgfpathlineto{\pgfqpoint{4.903350in}{0.660000in}}%
\pgfpathlineto{\pgfqpoint{4.903350in}{3.259495in}}%
\pgfpathlineto{\pgfqpoint{4.897150in}{3.259495in}}%
\pgfpathlineto{\pgfqpoint{4.897150in}{0.660000in}}%
\pgfpathclose%
\pgfusepath{fill}%
\end{pgfscope}%
\begin{pgfscope}%
\pgfpathrectangle{\pgfqpoint{1.250000in}{0.660000in}}{\pgfqpoint{7.750000in}{4.620000in}}%
\pgfusepath{clip}%
\pgfsetbuttcap%
\pgfsetmiterjoin%
\definecolor{currentfill}{rgb}{0.000000,0.000000,0.000000}%
\pgfsetfillcolor{currentfill}%
\pgfsetlinewidth{0.000000pt}%
\definecolor{currentstroke}{rgb}{0.000000,0.000000,0.000000}%
\pgfsetstrokecolor{currentstroke}%
\pgfsetstrokeopacity{0.000000}%
\pgfsetdash{}{0pt}%
\pgfpathmoveto{\pgfqpoint{4.904900in}{0.660000in}}%
\pgfpathlineto{\pgfqpoint{4.911100in}{0.660000in}}%
\pgfpathlineto{\pgfqpoint{4.911100in}{3.259278in}}%
\pgfpathlineto{\pgfqpoint{4.904900in}{3.259278in}}%
\pgfpathlineto{\pgfqpoint{4.904900in}{0.660000in}}%
\pgfpathclose%
\pgfusepath{fill}%
\end{pgfscope}%
\begin{pgfscope}%
\pgfpathrectangle{\pgfqpoint{1.250000in}{0.660000in}}{\pgfqpoint{7.750000in}{4.620000in}}%
\pgfusepath{clip}%
\pgfsetbuttcap%
\pgfsetmiterjoin%
\definecolor{currentfill}{rgb}{0.000000,0.000000,0.000000}%
\pgfsetfillcolor{currentfill}%
\pgfsetlinewidth{0.000000pt}%
\definecolor{currentstroke}{rgb}{0.000000,0.000000,0.000000}%
\pgfsetstrokecolor{currentstroke}%
\pgfsetstrokeopacity{0.000000}%
\pgfsetdash{}{0pt}%
\pgfpathmoveto{\pgfqpoint{4.912650in}{0.660000in}}%
\pgfpathlineto{\pgfqpoint{4.918850in}{0.660000in}}%
\pgfpathlineto{\pgfqpoint{4.918850in}{3.259060in}}%
\pgfpathlineto{\pgfqpoint{4.912650in}{3.259060in}}%
\pgfpathlineto{\pgfqpoint{4.912650in}{0.660000in}}%
\pgfpathclose%
\pgfusepath{fill}%
\end{pgfscope}%
\begin{pgfscope}%
\pgfpathrectangle{\pgfqpoint{1.250000in}{0.660000in}}{\pgfqpoint{7.750000in}{4.620000in}}%
\pgfusepath{clip}%
\pgfsetbuttcap%
\pgfsetmiterjoin%
\definecolor{currentfill}{rgb}{0.000000,0.000000,0.000000}%
\pgfsetfillcolor{currentfill}%
\pgfsetlinewidth{0.000000pt}%
\definecolor{currentstroke}{rgb}{0.000000,0.000000,0.000000}%
\pgfsetstrokecolor{currentstroke}%
\pgfsetstrokeopacity{0.000000}%
\pgfsetdash{}{0pt}%
\pgfpathmoveto{\pgfqpoint{4.920400in}{0.660000in}}%
\pgfpathlineto{\pgfqpoint{4.926600in}{0.660000in}}%
\pgfpathlineto{\pgfqpoint{4.926600in}{3.258843in}}%
\pgfpathlineto{\pgfqpoint{4.920400in}{3.258843in}}%
\pgfpathlineto{\pgfqpoint{4.920400in}{0.660000in}}%
\pgfpathclose%
\pgfusepath{fill}%
\end{pgfscope}%
\begin{pgfscope}%
\pgfpathrectangle{\pgfqpoint{1.250000in}{0.660000in}}{\pgfqpoint{7.750000in}{4.620000in}}%
\pgfusepath{clip}%
\pgfsetbuttcap%
\pgfsetmiterjoin%
\definecolor{currentfill}{rgb}{0.000000,0.000000,0.000000}%
\pgfsetfillcolor{currentfill}%
\pgfsetlinewidth{0.000000pt}%
\definecolor{currentstroke}{rgb}{0.000000,0.000000,0.000000}%
\pgfsetstrokecolor{currentstroke}%
\pgfsetstrokeopacity{0.000000}%
\pgfsetdash{}{0pt}%
\pgfpathmoveto{\pgfqpoint{4.928150in}{0.660000in}}%
\pgfpathlineto{\pgfqpoint{4.934350in}{0.660000in}}%
\pgfpathlineto{\pgfqpoint{4.934350in}{3.258625in}}%
\pgfpathlineto{\pgfqpoint{4.928150in}{3.258625in}}%
\pgfpathlineto{\pgfqpoint{4.928150in}{0.660000in}}%
\pgfpathclose%
\pgfusepath{fill}%
\end{pgfscope}%
\begin{pgfscope}%
\pgfpathrectangle{\pgfqpoint{1.250000in}{0.660000in}}{\pgfqpoint{7.750000in}{4.620000in}}%
\pgfusepath{clip}%
\pgfsetbuttcap%
\pgfsetmiterjoin%
\definecolor{currentfill}{rgb}{0.000000,0.000000,0.000000}%
\pgfsetfillcolor{currentfill}%
\pgfsetlinewidth{0.000000pt}%
\definecolor{currentstroke}{rgb}{0.000000,0.000000,0.000000}%
\pgfsetstrokecolor{currentstroke}%
\pgfsetstrokeopacity{0.000000}%
\pgfsetdash{}{0pt}%
\pgfpathmoveto{\pgfqpoint{4.935900in}{0.660000in}}%
\pgfpathlineto{\pgfqpoint{4.942100in}{0.660000in}}%
\pgfpathlineto{\pgfqpoint{4.942100in}{3.258408in}}%
\pgfpathlineto{\pgfqpoint{4.935900in}{3.258408in}}%
\pgfpathlineto{\pgfqpoint{4.935900in}{0.660000in}}%
\pgfpathclose%
\pgfusepath{fill}%
\end{pgfscope}%
\begin{pgfscope}%
\pgfpathrectangle{\pgfqpoint{1.250000in}{0.660000in}}{\pgfqpoint{7.750000in}{4.620000in}}%
\pgfusepath{clip}%
\pgfsetbuttcap%
\pgfsetmiterjoin%
\definecolor{currentfill}{rgb}{0.000000,0.000000,0.000000}%
\pgfsetfillcolor{currentfill}%
\pgfsetlinewidth{0.000000pt}%
\definecolor{currentstroke}{rgb}{0.000000,0.000000,0.000000}%
\pgfsetstrokecolor{currentstroke}%
\pgfsetstrokeopacity{0.000000}%
\pgfsetdash{}{0pt}%
\pgfpathmoveto{\pgfqpoint{4.943650in}{0.660000in}}%
\pgfpathlineto{\pgfqpoint{4.949850in}{0.660000in}}%
\pgfpathlineto{\pgfqpoint{4.949850in}{3.258190in}}%
\pgfpathlineto{\pgfqpoint{4.943650in}{3.258190in}}%
\pgfpathlineto{\pgfqpoint{4.943650in}{0.660000in}}%
\pgfpathclose%
\pgfusepath{fill}%
\end{pgfscope}%
\begin{pgfscope}%
\pgfpathrectangle{\pgfqpoint{1.250000in}{0.660000in}}{\pgfqpoint{7.750000in}{4.620000in}}%
\pgfusepath{clip}%
\pgfsetbuttcap%
\pgfsetmiterjoin%
\definecolor{currentfill}{rgb}{0.000000,0.000000,0.000000}%
\pgfsetfillcolor{currentfill}%
\pgfsetlinewidth{0.000000pt}%
\definecolor{currentstroke}{rgb}{0.000000,0.000000,0.000000}%
\pgfsetstrokecolor{currentstroke}%
\pgfsetstrokeopacity{0.000000}%
\pgfsetdash{}{0pt}%
\pgfpathmoveto{\pgfqpoint{4.951400in}{0.660000in}}%
\pgfpathlineto{\pgfqpoint{4.957600in}{0.660000in}}%
\pgfpathlineto{\pgfqpoint{4.957600in}{3.257973in}}%
\pgfpathlineto{\pgfqpoint{4.951400in}{3.257973in}}%
\pgfpathlineto{\pgfqpoint{4.951400in}{0.660000in}}%
\pgfpathclose%
\pgfusepath{fill}%
\end{pgfscope}%
\begin{pgfscope}%
\pgfpathrectangle{\pgfqpoint{1.250000in}{0.660000in}}{\pgfqpoint{7.750000in}{4.620000in}}%
\pgfusepath{clip}%
\pgfsetbuttcap%
\pgfsetmiterjoin%
\definecolor{currentfill}{rgb}{0.000000,0.000000,0.000000}%
\pgfsetfillcolor{currentfill}%
\pgfsetlinewidth{0.000000pt}%
\definecolor{currentstroke}{rgb}{0.000000,0.000000,0.000000}%
\pgfsetstrokecolor{currentstroke}%
\pgfsetstrokeopacity{0.000000}%
\pgfsetdash{}{0pt}%
\pgfpathmoveto{\pgfqpoint{4.959150in}{0.660000in}}%
\pgfpathlineto{\pgfqpoint{4.965350in}{0.660000in}}%
\pgfpathlineto{\pgfqpoint{4.965350in}{3.257756in}}%
\pgfpathlineto{\pgfqpoint{4.959150in}{3.257756in}}%
\pgfpathlineto{\pgfqpoint{4.959150in}{0.660000in}}%
\pgfpathclose%
\pgfusepath{fill}%
\end{pgfscope}%
\begin{pgfscope}%
\pgfpathrectangle{\pgfqpoint{1.250000in}{0.660000in}}{\pgfqpoint{7.750000in}{4.620000in}}%
\pgfusepath{clip}%
\pgfsetbuttcap%
\pgfsetmiterjoin%
\definecolor{currentfill}{rgb}{0.000000,0.000000,0.000000}%
\pgfsetfillcolor{currentfill}%
\pgfsetlinewidth{0.000000pt}%
\definecolor{currentstroke}{rgb}{0.000000,0.000000,0.000000}%
\pgfsetstrokecolor{currentstroke}%
\pgfsetstrokeopacity{0.000000}%
\pgfsetdash{}{0pt}%
\pgfpathmoveto{\pgfqpoint{4.966900in}{0.660000in}}%
\pgfpathlineto{\pgfqpoint{4.973100in}{0.660000in}}%
\pgfpathlineto{\pgfqpoint{4.973100in}{3.257543in}}%
\pgfpathlineto{\pgfqpoint{4.966900in}{3.257543in}}%
\pgfpathlineto{\pgfqpoint{4.966900in}{0.660000in}}%
\pgfpathclose%
\pgfusepath{fill}%
\end{pgfscope}%
\begin{pgfscope}%
\pgfpathrectangle{\pgfqpoint{1.250000in}{0.660000in}}{\pgfqpoint{7.750000in}{4.620000in}}%
\pgfusepath{clip}%
\pgfsetbuttcap%
\pgfsetmiterjoin%
\definecolor{currentfill}{rgb}{0.000000,0.000000,0.000000}%
\pgfsetfillcolor{currentfill}%
\pgfsetlinewidth{0.000000pt}%
\definecolor{currentstroke}{rgb}{0.000000,0.000000,0.000000}%
\pgfsetstrokecolor{currentstroke}%
\pgfsetstrokeopacity{0.000000}%
\pgfsetdash{}{0pt}%
\pgfpathmoveto{\pgfqpoint{4.974650in}{0.660000in}}%
\pgfpathlineto{\pgfqpoint{4.980850in}{0.660000in}}%
\pgfpathlineto{\pgfqpoint{4.980850in}{3.257330in}}%
\pgfpathlineto{\pgfqpoint{4.974650in}{3.257330in}}%
\pgfpathlineto{\pgfqpoint{4.974650in}{0.660000in}}%
\pgfpathclose%
\pgfusepath{fill}%
\end{pgfscope}%
\begin{pgfscope}%
\pgfpathrectangle{\pgfqpoint{1.250000in}{0.660000in}}{\pgfqpoint{7.750000in}{4.620000in}}%
\pgfusepath{clip}%
\pgfsetbuttcap%
\pgfsetmiterjoin%
\definecolor{currentfill}{rgb}{0.000000,0.000000,0.000000}%
\pgfsetfillcolor{currentfill}%
\pgfsetlinewidth{0.000000pt}%
\definecolor{currentstroke}{rgb}{0.000000,0.000000,0.000000}%
\pgfsetstrokecolor{currentstroke}%
\pgfsetstrokeopacity{0.000000}%
\pgfsetdash{}{0pt}%
\pgfpathmoveto{\pgfqpoint{4.982400in}{0.660000in}}%
\pgfpathlineto{\pgfqpoint{4.988600in}{0.660000in}}%
\pgfpathlineto{\pgfqpoint{4.988600in}{3.257118in}}%
\pgfpathlineto{\pgfqpoint{4.982400in}{3.257118in}}%
\pgfpathlineto{\pgfqpoint{4.982400in}{0.660000in}}%
\pgfpathclose%
\pgfusepath{fill}%
\end{pgfscope}%
\begin{pgfscope}%
\pgfpathrectangle{\pgfqpoint{1.250000in}{0.660000in}}{\pgfqpoint{7.750000in}{4.620000in}}%
\pgfusepath{clip}%
\pgfsetbuttcap%
\pgfsetmiterjoin%
\definecolor{currentfill}{rgb}{0.000000,0.000000,0.000000}%
\pgfsetfillcolor{currentfill}%
\pgfsetlinewidth{0.000000pt}%
\definecolor{currentstroke}{rgb}{0.000000,0.000000,0.000000}%
\pgfsetstrokecolor{currentstroke}%
\pgfsetstrokeopacity{0.000000}%
\pgfsetdash{}{0pt}%
\pgfpathmoveto{\pgfqpoint{4.990150in}{0.660000in}}%
\pgfpathlineto{\pgfqpoint{4.996350in}{0.660000in}}%
\pgfpathlineto{\pgfqpoint{4.996350in}{3.256905in}}%
\pgfpathlineto{\pgfqpoint{4.990150in}{3.256905in}}%
\pgfpathlineto{\pgfqpoint{4.990150in}{0.660000in}}%
\pgfpathclose%
\pgfusepath{fill}%
\end{pgfscope}%
\begin{pgfscope}%
\pgfpathrectangle{\pgfqpoint{1.250000in}{0.660000in}}{\pgfqpoint{7.750000in}{4.620000in}}%
\pgfusepath{clip}%
\pgfsetbuttcap%
\pgfsetmiterjoin%
\definecolor{currentfill}{rgb}{0.000000,0.000000,0.000000}%
\pgfsetfillcolor{currentfill}%
\pgfsetlinewidth{0.000000pt}%
\definecolor{currentstroke}{rgb}{0.000000,0.000000,0.000000}%
\pgfsetstrokecolor{currentstroke}%
\pgfsetstrokeopacity{0.000000}%
\pgfsetdash{}{0pt}%
\pgfpathmoveto{\pgfqpoint{4.997900in}{0.660000in}}%
\pgfpathlineto{\pgfqpoint{5.004100in}{0.660000in}}%
\pgfpathlineto{\pgfqpoint{5.004100in}{3.256693in}}%
\pgfpathlineto{\pgfqpoint{4.997900in}{3.256693in}}%
\pgfpathlineto{\pgfqpoint{4.997900in}{0.660000in}}%
\pgfpathclose%
\pgfusepath{fill}%
\end{pgfscope}%
\begin{pgfscope}%
\pgfpathrectangle{\pgfqpoint{1.250000in}{0.660000in}}{\pgfqpoint{7.750000in}{4.620000in}}%
\pgfusepath{clip}%
\pgfsetbuttcap%
\pgfsetmiterjoin%
\definecolor{currentfill}{rgb}{0.000000,0.000000,0.000000}%
\pgfsetfillcolor{currentfill}%
\pgfsetlinewidth{0.000000pt}%
\definecolor{currentstroke}{rgb}{0.000000,0.000000,0.000000}%
\pgfsetstrokecolor{currentstroke}%
\pgfsetstrokeopacity{0.000000}%
\pgfsetdash{}{0pt}%
\pgfpathmoveto{\pgfqpoint{5.005650in}{0.660000in}}%
\pgfpathlineto{\pgfqpoint{5.011850in}{0.660000in}}%
\pgfpathlineto{\pgfqpoint{5.011850in}{3.256481in}}%
\pgfpathlineto{\pgfqpoint{5.005650in}{3.256481in}}%
\pgfpathlineto{\pgfqpoint{5.005650in}{0.660000in}}%
\pgfpathclose%
\pgfusepath{fill}%
\end{pgfscope}%
\begin{pgfscope}%
\pgfpathrectangle{\pgfqpoint{1.250000in}{0.660000in}}{\pgfqpoint{7.750000in}{4.620000in}}%
\pgfusepath{clip}%
\pgfsetbuttcap%
\pgfsetmiterjoin%
\definecolor{currentfill}{rgb}{0.000000,0.000000,0.000000}%
\pgfsetfillcolor{currentfill}%
\pgfsetlinewidth{0.000000pt}%
\definecolor{currentstroke}{rgb}{0.000000,0.000000,0.000000}%
\pgfsetstrokecolor{currentstroke}%
\pgfsetstrokeopacity{0.000000}%
\pgfsetdash{}{0pt}%
\pgfpathmoveto{\pgfqpoint{5.013400in}{0.660000in}}%
\pgfpathlineto{\pgfqpoint{5.019600in}{0.660000in}}%
\pgfpathlineto{\pgfqpoint{5.019600in}{3.256268in}}%
\pgfpathlineto{\pgfqpoint{5.013400in}{3.256268in}}%
\pgfpathlineto{\pgfqpoint{5.013400in}{0.660000in}}%
\pgfpathclose%
\pgfusepath{fill}%
\end{pgfscope}%
\begin{pgfscope}%
\pgfpathrectangle{\pgfqpoint{1.250000in}{0.660000in}}{\pgfqpoint{7.750000in}{4.620000in}}%
\pgfusepath{clip}%
\pgfsetbuttcap%
\pgfsetmiterjoin%
\definecolor{currentfill}{rgb}{0.000000,0.000000,0.000000}%
\pgfsetfillcolor{currentfill}%
\pgfsetlinewidth{0.000000pt}%
\definecolor{currentstroke}{rgb}{0.000000,0.000000,0.000000}%
\pgfsetstrokecolor{currentstroke}%
\pgfsetstrokeopacity{0.000000}%
\pgfsetdash{}{0pt}%
\pgfpathmoveto{\pgfqpoint{5.021150in}{0.660000in}}%
\pgfpathlineto{\pgfqpoint{5.027350in}{0.660000in}}%
\pgfpathlineto{\pgfqpoint{5.027350in}{3.256056in}}%
\pgfpathlineto{\pgfqpoint{5.021150in}{3.256056in}}%
\pgfpathlineto{\pgfqpoint{5.021150in}{0.660000in}}%
\pgfpathclose%
\pgfusepath{fill}%
\end{pgfscope}%
\begin{pgfscope}%
\pgfpathrectangle{\pgfqpoint{1.250000in}{0.660000in}}{\pgfqpoint{7.750000in}{4.620000in}}%
\pgfusepath{clip}%
\pgfsetbuttcap%
\pgfsetmiterjoin%
\definecolor{currentfill}{rgb}{0.000000,0.000000,0.000000}%
\pgfsetfillcolor{currentfill}%
\pgfsetlinewidth{0.000000pt}%
\definecolor{currentstroke}{rgb}{0.000000,0.000000,0.000000}%
\pgfsetstrokecolor{currentstroke}%
\pgfsetstrokeopacity{0.000000}%
\pgfsetdash{}{0pt}%
\pgfpathmoveto{\pgfqpoint{5.028900in}{0.660000in}}%
\pgfpathlineto{\pgfqpoint{5.035100in}{0.660000in}}%
\pgfpathlineto{\pgfqpoint{5.035100in}{3.255847in}}%
\pgfpathlineto{\pgfqpoint{5.028900in}{3.255847in}}%
\pgfpathlineto{\pgfqpoint{5.028900in}{0.660000in}}%
\pgfpathclose%
\pgfusepath{fill}%
\end{pgfscope}%
\begin{pgfscope}%
\pgfpathrectangle{\pgfqpoint{1.250000in}{0.660000in}}{\pgfqpoint{7.750000in}{4.620000in}}%
\pgfusepath{clip}%
\pgfsetbuttcap%
\pgfsetmiterjoin%
\definecolor{currentfill}{rgb}{0.000000,0.000000,0.000000}%
\pgfsetfillcolor{currentfill}%
\pgfsetlinewidth{0.000000pt}%
\definecolor{currentstroke}{rgb}{0.000000,0.000000,0.000000}%
\pgfsetstrokecolor{currentstroke}%
\pgfsetstrokeopacity{0.000000}%
\pgfsetdash{}{0pt}%
\pgfpathmoveto{\pgfqpoint{5.036650in}{0.660000in}}%
\pgfpathlineto{\pgfqpoint{5.042850in}{0.660000in}}%
\pgfpathlineto{\pgfqpoint{5.042850in}{3.255639in}}%
\pgfpathlineto{\pgfqpoint{5.036650in}{3.255639in}}%
\pgfpathlineto{\pgfqpoint{5.036650in}{0.660000in}}%
\pgfpathclose%
\pgfusepath{fill}%
\end{pgfscope}%
\begin{pgfscope}%
\pgfpathrectangle{\pgfqpoint{1.250000in}{0.660000in}}{\pgfqpoint{7.750000in}{4.620000in}}%
\pgfusepath{clip}%
\pgfsetbuttcap%
\pgfsetmiterjoin%
\definecolor{currentfill}{rgb}{0.000000,0.000000,0.000000}%
\pgfsetfillcolor{currentfill}%
\pgfsetlinewidth{0.000000pt}%
\definecolor{currentstroke}{rgb}{0.000000,0.000000,0.000000}%
\pgfsetstrokecolor{currentstroke}%
\pgfsetstrokeopacity{0.000000}%
\pgfsetdash{}{0pt}%
\pgfpathmoveto{\pgfqpoint{5.044400in}{0.660000in}}%
\pgfpathlineto{\pgfqpoint{5.050600in}{0.660000in}}%
\pgfpathlineto{\pgfqpoint{5.050600in}{3.255432in}}%
\pgfpathlineto{\pgfqpoint{5.044400in}{3.255432in}}%
\pgfpathlineto{\pgfqpoint{5.044400in}{0.660000in}}%
\pgfpathclose%
\pgfusepath{fill}%
\end{pgfscope}%
\begin{pgfscope}%
\pgfpathrectangle{\pgfqpoint{1.250000in}{0.660000in}}{\pgfqpoint{7.750000in}{4.620000in}}%
\pgfusepath{clip}%
\pgfsetbuttcap%
\pgfsetmiterjoin%
\definecolor{currentfill}{rgb}{0.000000,0.000000,0.000000}%
\pgfsetfillcolor{currentfill}%
\pgfsetlinewidth{0.000000pt}%
\definecolor{currentstroke}{rgb}{0.000000,0.000000,0.000000}%
\pgfsetstrokecolor{currentstroke}%
\pgfsetstrokeopacity{0.000000}%
\pgfsetdash{}{0pt}%
\pgfpathmoveto{\pgfqpoint{5.052150in}{0.660000in}}%
\pgfpathlineto{\pgfqpoint{5.058350in}{0.660000in}}%
\pgfpathlineto{\pgfqpoint{5.058350in}{3.255224in}}%
\pgfpathlineto{\pgfqpoint{5.052150in}{3.255224in}}%
\pgfpathlineto{\pgfqpoint{5.052150in}{0.660000in}}%
\pgfpathclose%
\pgfusepath{fill}%
\end{pgfscope}%
\begin{pgfscope}%
\pgfpathrectangle{\pgfqpoint{1.250000in}{0.660000in}}{\pgfqpoint{7.750000in}{4.620000in}}%
\pgfusepath{clip}%
\pgfsetbuttcap%
\pgfsetmiterjoin%
\definecolor{currentfill}{rgb}{0.000000,0.000000,0.000000}%
\pgfsetfillcolor{currentfill}%
\pgfsetlinewidth{0.000000pt}%
\definecolor{currentstroke}{rgb}{0.000000,0.000000,0.000000}%
\pgfsetstrokecolor{currentstroke}%
\pgfsetstrokeopacity{0.000000}%
\pgfsetdash{}{0pt}%
\pgfpathmoveto{\pgfqpoint{5.059900in}{0.660000in}}%
\pgfpathlineto{\pgfqpoint{5.066100in}{0.660000in}}%
\pgfpathlineto{\pgfqpoint{5.066100in}{3.255017in}}%
\pgfpathlineto{\pgfqpoint{5.059900in}{3.255017in}}%
\pgfpathlineto{\pgfqpoint{5.059900in}{0.660000in}}%
\pgfpathclose%
\pgfusepath{fill}%
\end{pgfscope}%
\begin{pgfscope}%
\pgfpathrectangle{\pgfqpoint{1.250000in}{0.660000in}}{\pgfqpoint{7.750000in}{4.620000in}}%
\pgfusepath{clip}%
\pgfsetbuttcap%
\pgfsetmiterjoin%
\definecolor{currentfill}{rgb}{0.000000,0.000000,0.000000}%
\pgfsetfillcolor{currentfill}%
\pgfsetlinewidth{0.000000pt}%
\definecolor{currentstroke}{rgb}{0.000000,0.000000,0.000000}%
\pgfsetstrokecolor{currentstroke}%
\pgfsetstrokeopacity{0.000000}%
\pgfsetdash{}{0pt}%
\pgfpathmoveto{\pgfqpoint{5.067650in}{0.660000in}}%
\pgfpathlineto{\pgfqpoint{5.073850in}{0.660000in}}%
\pgfpathlineto{\pgfqpoint{5.073850in}{3.254809in}}%
\pgfpathlineto{\pgfqpoint{5.067650in}{3.254809in}}%
\pgfpathlineto{\pgfqpoint{5.067650in}{0.660000in}}%
\pgfpathclose%
\pgfusepath{fill}%
\end{pgfscope}%
\begin{pgfscope}%
\pgfpathrectangle{\pgfqpoint{1.250000in}{0.660000in}}{\pgfqpoint{7.750000in}{4.620000in}}%
\pgfusepath{clip}%
\pgfsetbuttcap%
\pgfsetmiterjoin%
\definecolor{currentfill}{rgb}{0.000000,0.000000,0.000000}%
\pgfsetfillcolor{currentfill}%
\pgfsetlinewidth{0.000000pt}%
\definecolor{currentstroke}{rgb}{0.000000,0.000000,0.000000}%
\pgfsetstrokecolor{currentstroke}%
\pgfsetstrokeopacity{0.000000}%
\pgfsetdash{}{0pt}%
\pgfpathmoveto{\pgfqpoint{5.075400in}{0.660000in}}%
\pgfpathlineto{\pgfqpoint{5.081600in}{0.660000in}}%
\pgfpathlineto{\pgfqpoint{5.081600in}{3.254602in}}%
\pgfpathlineto{\pgfqpoint{5.075400in}{3.254602in}}%
\pgfpathlineto{\pgfqpoint{5.075400in}{0.660000in}}%
\pgfpathclose%
\pgfusepath{fill}%
\end{pgfscope}%
\begin{pgfscope}%
\pgfpathrectangle{\pgfqpoint{1.250000in}{0.660000in}}{\pgfqpoint{7.750000in}{4.620000in}}%
\pgfusepath{clip}%
\pgfsetbuttcap%
\pgfsetmiterjoin%
\definecolor{currentfill}{rgb}{0.000000,0.000000,0.000000}%
\pgfsetfillcolor{currentfill}%
\pgfsetlinewidth{0.000000pt}%
\definecolor{currentstroke}{rgb}{0.000000,0.000000,0.000000}%
\pgfsetstrokecolor{currentstroke}%
\pgfsetstrokeopacity{0.000000}%
\pgfsetdash{}{0pt}%
\pgfpathmoveto{\pgfqpoint{5.083150in}{0.660000in}}%
\pgfpathlineto{\pgfqpoint{5.089350in}{0.660000in}}%
\pgfpathlineto{\pgfqpoint{5.089350in}{3.254394in}}%
\pgfpathlineto{\pgfqpoint{5.083150in}{3.254394in}}%
\pgfpathlineto{\pgfqpoint{5.083150in}{0.660000in}}%
\pgfpathclose%
\pgfusepath{fill}%
\end{pgfscope}%
\begin{pgfscope}%
\pgfpathrectangle{\pgfqpoint{1.250000in}{0.660000in}}{\pgfqpoint{7.750000in}{4.620000in}}%
\pgfusepath{clip}%
\pgfsetbuttcap%
\pgfsetmiterjoin%
\definecolor{currentfill}{rgb}{0.000000,0.000000,0.000000}%
\pgfsetfillcolor{currentfill}%
\pgfsetlinewidth{0.000000pt}%
\definecolor{currentstroke}{rgb}{0.000000,0.000000,0.000000}%
\pgfsetstrokecolor{currentstroke}%
\pgfsetstrokeopacity{0.000000}%
\pgfsetdash{}{0pt}%
\pgfpathmoveto{\pgfqpoint{5.090900in}{0.660000in}}%
\pgfpathlineto{\pgfqpoint{5.097100in}{0.660000in}}%
\pgfpathlineto{\pgfqpoint{5.097100in}{3.254188in}}%
\pgfpathlineto{\pgfqpoint{5.090900in}{3.254188in}}%
\pgfpathlineto{\pgfqpoint{5.090900in}{0.660000in}}%
\pgfpathclose%
\pgfusepath{fill}%
\end{pgfscope}%
\begin{pgfscope}%
\pgfpathrectangle{\pgfqpoint{1.250000in}{0.660000in}}{\pgfqpoint{7.750000in}{4.620000in}}%
\pgfusepath{clip}%
\pgfsetbuttcap%
\pgfsetmiterjoin%
\definecolor{currentfill}{rgb}{0.000000,0.000000,0.000000}%
\pgfsetfillcolor{currentfill}%
\pgfsetlinewidth{0.000000pt}%
\definecolor{currentstroke}{rgb}{0.000000,0.000000,0.000000}%
\pgfsetstrokecolor{currentstroke}%
\pgfsetstrokeopacity{0.000000}%
\pgfsetdash{}{0pt}%
\pgfpathmoveto{\pgfqpoint{5.098650in}{0.660000in}}%
\pgfpathlineto{\pgfqpoint{5.104850in}{0.660000in}}%
\pgfpathlineto{\pgfqpoint{5.104850in}{3.253985in}}%
\pgfpathlineto{\pgfqpoint{5.098650in}{3.253985in}}%
\pgfpathlineto{\pgfqpoint{5.098650in}{0.660000in}}%
\pgfpathclose%
\pgfusepath{fill}%
\end{pgfscope}%
\begin{pgfscope}%
\pgfpathrectangle{\pgfqpoint{1.250000in}{0.660000in}}{\pgfqpoint{7.750000in}{4.620000in}}%
\pgfusepath{clip}%
\pgfsetbuttcap%
\pgfsetmiterjoin%
\definecolor{currentfill}{rgb}{0.000000,0.000000,0.000000}%
\pgfsetfillcolor{currentfill}%
\pgfsetlinewidth{0.000000pt}%
\definecolor{currentstroke}{rgb}{0.000000,0.000000,0.000000}%
\pgfsetstrokecolor{currentstroke}%
\pgfsetstrokeopacity{0.000000}%
\pgfsetdash{}{0pt}%
\pgfpathmoveto{\pgfqpoint{5.106400in}{0.660000in}}%
\pgfpathlineto{\pgfqpoint{5.112600in}{0.660000in}}%
\pgfpathlineto{\pgfqpoint{5.112600in}{3.253783in}}%
\pgfpathlineto{\pgfqpoint{5.106400in}{3.253783in}}%
\pgfpathlineto{\pgfqpoint{5.106400in}{0.660000in}}%
\pgfpathclose%
\pgfusepath{fill}%
\end{pgfscope}%
\begin{pgfscope}%
\pgfpathrectangle{\pgfqpoint{1.250000in}{0.660000in}}{\pgfqpoint{7.750000in}{4.620000in}}%
\pgfusepath{clip}%
\pgfsetbuttcap%
\pgfsetmiterjoin%
\definecolor{currentfill}{rgb}{0.000000,0.000000,0.000000}%
\pgfsetfillcolor{currentfill}%
\pgfsetlinewidth{0.000000pt}%
\definecolor{currentstroke}{rgb}{0.000000,0.000000,0.000000}%
\pgfsetstrokecolor{currentstroke}%
\pgfsetstrokeopacity{0.000000}%
\pgfsetdash{}{0pt}%
\pgfpathmoveto{\pgfqpoint{5.114150in}{0.660000in}}%
\pgfpathlineto{\pgfqpoint{5.120350in}{0.660000in}}%
\pgfpathlineto{\pgfqpoint{5.120350in}{3.253580in}}%
\pgfpathlineto{\pgfqpoint{5.114150in}{3.253580in}}%
\pgfpathlineto{\pgfqpoint{5.114150in}{0.660000in}}%
\pgfpathclose%
\pgfusepath{fill}%
\end{pgfscope}%
\begin{pgfscope}%
\pgfpathrectangle{\pgfqpoint{1.250000in}{0.660000in}}{\pgfqpoint{7.750000in}{4.620000in}}%
\pgfusepath{clip}%
\pgfsetbuttcap%
\pgfsetmiterjoin%
\definecolor{currentfill}{rgb}{0.000000,0.000000,0.000000}%
\pgfsetfillcolor{currentfill}%
\pgfsetlinewidth{0.000000pt}%
\definecolor{currentstroke}{rgb}{0.000000,0.000000,0.000000}%
\pgfsetstrokecolor{currentstroke}%
\pgfsetstrokeopacity{0.000000}%
\pgfsetdash{}{0pt}%
\pgfpathmoveto{\pgfqpoint{5.121900in}{0.660000in}}%
\pgfpathlineto{\pgfqpoint{5.128100in}{0.660000in}}%
\pgfpathlineto{\pgfqpoint{5.128100in}{3.253378in}}%
\pgfpathlineto{\pgfqpoint{5.121900in}{3.253378in}}%
\pgfpathlineto{\pgfqpoint{5.121900in}{0.660000in}}%
\pgfpathclose%
\pgfusepath{fill}%
\end{pgfscope}%
\begin{pgfscope}%
\pgfpathrectangle{\pgfqpoint{1.250000in}{0.660000in}}{\pgfqpoint{7.750000in}{4.620000in}}%
\pgfusepath{clip}%
\pgfsetbuttcap%
\pgfsetmiterjoin%
\definecolor{currentfill}{rgb}{0.000000,0.000000,0.000000}%
\pgfsetfillcolor{currentfill}%
\pgfsetlinewidth{0.000000pt}%
\definecolor{currentstroke}{rgb}{0.000000,0.000000,0.000000}%
\pgfsetstrokecolor{currentstroke}%
\pgfsetstrokeopacity{0.000000}%
\pgfsetdash{}{0pt}%
\pgfpathmoveto{\pgfqpoint{5.129650in}{0.660000in}}%
\pgfpathlineto{\pgfqpoint{5.135850in}{0.660000in}}%
\pgfpathlineto{\pgfqpoint{5.135850in}{3.253175in}}%
\pgfpathlineto{\pgfqpoint{5.129650in}{3.253175in}}%
\pgfpathlineto{\pgfqpoint{5.129650in}{0.660000in}}%
\pgfpathclose%
\pgfusepath{fill}%
\end{pgfscope}%
\begin{pgfscope}%
\pgfpathrectangle{\pgfqpoint{1.250000in}{0.660000in}}{\pgfqpoint{7.750000in}{4.620000in}}%
\pgfusepath{clip}%
\pgfsetbuttcap%
\pgfsetmiterjoin%
\definecolor{currentfill}{rgb}{0.000000,0.000000,0.000000}%
\pgfsetfillcolor{currentfill}%
\pgfsetlinewidth{0.000000pt}%
\definecolor{currentstroke}{rgb}{0.000000,0.000000,0.000000}%
\pgfsetstrokecolor{currentstroke}%
\pgfsetstrokeopacity{0.000000}%
\pgfsetdash{}{0pt}%
\pgfpathmoveto{\pgfqpoint{5.137400in}{0.660000in}}%
\pgfpathlineto{\pgfqpoint{5.143600in}{0.660000in}}%
\pgfpathlineto{\pgfqpoint{5.143600in}{3.252972in}}%
\pgfpathlineto{\pgfqpoint{5.137400in}{3.252972in}}%
\pgfpathlineto{\pgfqpoint{5.137400in}{0.660000in}}%
\pgfpathclose%
\pgfusepath{fill}%
\end{pgfscope}%
\begin{pgfscope}%
\pgfpathrectangle{\pgfqpoint{1.250000in}{0.660000in}}{\pgfqpoint{7.750000in}{4.620000in}}%
\pgfusepath{clip}%
\pgfsetbuttcap%
\pgfsetmiterjoin%
\definecolor{currentfill}{rgb}{0.000000,0.000000,0.000000}%
\pgfsetfillcolor{currentfill}%
\pgfsetlinewidth{0.000000pt}%
\definecolor{currentstroke}{rgb}{0.000000,0.000000,0.000000}%
\pgfsetstrokecolor{currentstroke}%
\pgfsetstrokeopacity{0.000000}%
\pgfsetdash{}{0pt}%
\pgfpathmoveto{\pgfqpoint{5.145150in}{0.660000in}}%
\pgfpathlineto{\pgfqpoint{5.151350in}{0.660000in}}%
\pgfpathlineto{\pgfqpoint{5.151350in}{3.252770in}}%
\pgfpathlineto{\pgfqpoint{5.145150in}{3.252770in}}%
\pgfpathlineto{\pgfqpoint{5.145150in}{0.660000in}}%
\pgfpathclose%
\pgfusepath{fill}%
\end{pgfscope}%
\begin{pgfscope}%
\pgfpathrectangle{\pgfqpoint{1.250000in}{0.660000in}}{\pgfqpoint{7.750000in}{4.620000in}}%
\pgfusepath{clip}%
\pgfsetbuttcap%
\pgfsetmiterjoin%
\definecolor{currentfill}{rgb}{0.000000,0.000000,0.000000}%
\pgfsetfillcolor{currentfill}%
\pgfsetlinewidth{0.000000pt}%
\definecolor{currentstroke}{rgb}{0.000000,0.000000,0.000000}%
\pgfsetstrokecolor{currentstroke}%
\pgfsetstrokeopacity{0.000000}%
\pgfsetdash{}{0pt}%
\pgfpathmoveto{\pgfqpoint{5.152900in}{0.660000in}}%
\pgfpathlineto{\pgfqpoint{5.159100in}{0.660000in}}%
\pgfpathlineto{\pgfqpoint{5.159100in}{3.252567in}}%
\pgfpathlineto{\pgfqpoint{5.152900in}{3.252567in}}%
\pgfpathlineto{\pgfqpoint{5.152900in}{0.660000in}}%
\pgfpathclose%
\pgfusepath{fill}%
\end{pgfscope}%
\begin{pgfscope}%
\pgfpathrectangle{\pgfqpoint{1.250000in}{0.660000in}}{\pgfqpoint{7.750000in}{4.620000in}}%
\pgfusepath{clip}%
\pgfsetbuttcap%
\pgfsetmiterjoin%
\definecolor{currentfill}{rgb}{0.000000,0.000000,0.000000}%
\pgfsetfillcolor{currentfill}%
\pgfsetlinewidth{0.000000pt}%
\definecolor{currentstroke}{rgb}{0.000000,0.000000,0.000000}%
\pgfsetstrokecolor{currentstroke}%
\pgfsetstrokeopacity{0.000000}%
\pgfsetdash{}{0pt}%
\pgfpathmoveto{\pgfqpoint{5.160650in}{0.660000in}}%
\pgfpathlineto{\pgfqpoint{5.166850in}{0.660000in}}%
\pgfpathlineto{\pgfqpoint{5.166850in}{3.252367in}}%
\pgfpathlineto{\pgfqpoint{5.160650in}{3.252367in}}%
\pgfpathlineto{\pgfqpoint{5.160650in}{0.660000in}}%
\pgfpathclose%
\pgfusepath{fill}%
\end{pgfscope}%
\begin{pgfscope}%
\pgfpathrectangle{\pgfqpoint{1.250000in}{0.660000in}}{\pgfqpoint{7.750000in}{4.620000in}}%
\pgfusepath{clip}%
\pgfsetbuttcap%
\pgfsetmiterjoin%
\definecolor{currentfill}{rgb}{0.000000,0.000000,0.000000}%
\pgfsetfillcolor{currentfill}%
\pgfsetlinewidth{0.000000pt}%
\definecolor{currentstroke}{rgb}{0.000000,0.000000,0.000000}%
\pgfsetstrokecolor{currentstroke}%
\pgfsetstrokeopacity{0.000000}%
\pgfsetdash{}{0pt}%
\pgfpathmoveto{\pgfqpoint{5.168400in}{0.660000in}}%
\pgfpathlineto{\pgfqpoint{5.174600in}{0.660000in}}%
\pgfpathlineto{\pgfqpoint{5.174600in}{3.252169in}}%
\pgfpathlineto{\pgfqpoint{5.168400in}{3.252169in}}%
\pgfpathlineto{\pgfqpoint{5.168400in}{0.660000in}}%
\pgfpathclose%
\pgfusepath{fill}%
\end{pgfscope}%
\begin{pgfscope}%
\pgfpathrectangle{\pgfqpoint{1.250000in}{0.660000in}}{\pgfqpoint{7.750000in}{4.620000in}}%
\pgfusepath{clip}%
\pgfsetbuttcap%
\pgfsetmiterjoin%
\definecolor{currentfill}{rgb}{0.000000,0.000000,0.000000}%
\pgfsetfillcolor{currentfill}%
\pgfsetlinewidth{0.000000pt}%
\definecolor{currentstroke}{rgb}{0.000000,0.000000,0.000000}%
\pgfsetstrokecolor{currentstroke}%
\pgfsetstrokeopacity{0.000000}%
\pgfsetdash{}{0pt}%
\pgfpathmoveto{\pgfqpoint{5.176150in}{0.660000in}}%
\pgfpathlineto{\pgfqpoint{5.182350in}{0.660000in}}%
\pgfpathlineto{\pgfqpoint{5.182350in}{3.251971in}}%
\pgfpathlineto{\pgfqpoint{5.176150in}{3.251971in}}%
\pgfpathlineto{\pgfqpoint{5.176150in}{0.660000in}}%
\pgfpathclose%
\pgfusepath{fill}%
\end{pgfscope}%
\begin{pgfscope}%
\pgfpathrectangle{\pgfqpoint{1.250000in}{0.660000in}}{\pgfqpoint{7.750000in}{4.620000in}}%
\pgfusepath{clip}%
\pgfsetbuttcap%
\pgfsetmiterjoin%
\definecolor{currentfill}{rgb}{0.000000,0.000000,0.000000}%
\pgfsetfillcolor{currentfill}%
\pgfsetlinewidth{0.000000pt}%
\definecolor{currentstroke}{rgb}{0.000000,0.000000,0.000000}%
\pgfsetstrokecolor{currentstroke}%
\pgfsetstrokeopacity{0.000000}%
\pgfsetdash{}{0pt}%
\pgfpathmoveto{\pgfqpoint{5.183900in}{0.660000in}}%
\pgfpathlineto{\pgfqpoint{5.190100in}{0.660000in}}%
\pgfpathlineto{\pgfqpoint{5.190100in}{3.251774in}}%
\pgfpathlineto{\pgfqpoint{5.183900in}{3.251774in}}%
\pgfpathlineto{\pgfqpoint{5.183900in}{0.660000in}}%
\pgfpathclose%
\pgfusepath{fill}%
\end{pgfscope}%
\begin{pgfscope}%
\pgfpathrectangle{\pgfqpoint{1.250000in}{0.660000in}}{\pgfqpoint{7.750000in}{4.620000in}}%
\pgfusepath{clip}%
\pgfsetbuttcap%
\pgfsetmiterjoin%
\definecolor{currentfill}{rgb}{0.000000,0.000000,0.000000}%
\pgfsetfillcolor{currentfill}%
\pgfsetlinewidth{0.000000pt}%
\definecolor{currentstroke}{rgb}{0.000000,0.000000,0.000000}%
\pgfsetstrokecolor{currentstroke}%
\pgfsetstrokeopacity{0.000000}%
\pgfsetdash{}{0pt}%
\pgfpathmoveto{\pgfqpoint{5.191650in}{0.660000in}}%
\pgfpathlineto{\pgfqpoint{5.197850in}{0.660000in}}%
\pgfpathlineto{\pgfqpoint{5.197850in}{3.251576in}}%
\pgfpathlineto{\pgfqpoint{5.191650in}{3.251576in}}%
\pgfpathlineto{\pgfqpoint{5.191650in}{0.660000in}}%
\pgfpathclose%
\pgfusepath{fill}%
\end{pgfscope}%
\begin{pgfscope}%
\pgfpathrectangle{\pgfqpoint{1.250000in}{0.660000in}}{\pgfqpoint{7.750000in}{4.620000in}}%
\pgfusepath{clip}%
\pgfsetbuttcap%
\pgfsetmiterjoin%
\definecolor{currentfill}{rgb}{0.000000,0.000000,0.000000}%
\pgfsetfillcolor{currentfill}%
\pgfsetlinewidth{0.000000pt}%
\definecolor{currentstroke}{rgb}{0.000000,0.000000,0.000000}%
\pgfsetstrokecolor{currentstroke}%
\pgfsetstrokeopacity{0.000000}%
\pgfsetdash{}{0pt}%
\pgfpathmoveto{\pgfqpoint{5.199400in}{0.660000in}}%
\pgfpathlineto{\pgfqpoint{5.205600in}{0.660000in}}%
\pgfpathlineto{\pgfqpoint{5.205600in}{3.251378in}}%
\pgfpathlineto{\pgfqpoint{5.199400in}{3.251378in}}%
\pgfpathlineto{\pgfqpoint{5.199400in}{0.660000in}}%
\pgfpathclose%
\pgfusepath{fill}%
\end{pgfscope}%
\begin{pgfscope}%
\pgfpathrectangle{\pgfqpoint{1.250000in}{0.660000in}}{\pgfqpoint{7.750000in}{4.620000in}}%
\pgfusepath{clip}%
\pgfsetbuttcap%
\pgfsetmiterjoin%
\definecolor{currentfill}{rgb}{0.000000,0.000000,0.000000}%
\pgfsetfillcolor{currentfill}%
\pgfsetlinewidth{0.000000pt}%
\definecolor{currentstroke}{rgb}{0.000000,0.000000,0.000000}%
\pgfsetstrokecolor{currentstroke}%
\pgfsetstrokeopacity{0.000000}%
\pgfsetdash{}{0pt}%
\pgfpathmoveto{\pgfqpoint{5.207150in}{0.660000in}}%
\pgfpathlineto{\pgfqpoint{5.213350in}{0.660000in}}%
\pgfpathlineto{\pgfqpoint{5.213350in}{3.251181in}}%
\pgfpathlineto{\pgfqpoint{5.207150in}{3.251181in}}%
\pgfpathlineto{\pgfqpoint{5.207150in}{0.660000in}}%
\pgfpathclose%
\pgfusepath{fill}%
\end{pgfscope}%
\begin{pgfscope}%
\pgfpathrectangle{\pgfqpoint{1.250000in}{0.660000in}}{\pgfqpoint{7.750000in}{4.620000in}}%
\pgfusepath{clip}%
\pgfsetbuttcap%
\pgfsetmiterjoin%
\definecolor{currentfill}{rgb}{0.000000,0.000000,0.000000}%
\pgfsetfillcolor{currentfill}%
\pgfsetlinewidth{0.000000pt}%
\definecolor{currentstroke}{rgb}{0.000000,0.000000,0.000000}%
\pgfsetstrokecolor{currentstroke}%
\pgfsetstrokeopacity{0.000000}%
\pgfsetdash{}{0pt}%
\pgfpathmoveto{\pgfqpoint{5.214900in}{0.660000in}}%
\pgfpathlineto{\pgfqpoint{5.221100in}{0.660000in}}%
\pgfpathlineto{\pgfqpoint{5.221100in}{3.250983in}}%
\pgfpathlineto{\pgfqpoint{5.214900in}{3.250983in}}%
\pgfpathlineto{\pgfqpoint{5.214900in}{0.660000in}}%
\pgfpathclose%
\pgfusepath{fill}%
\end{pgfscope}%
\begin{pgfscope}%
\pgfpathrectangle{\pgfqpoint{1.250000in}{0.660000in}}{\pgfqpoint{7.750000in}{4.620000in}}%
\pgfusepath{clip}%
\pgfsetbuttcap%
\pgfsetmiterjoin%
\definecolor{currentfill}{rgb}{0.000000,0.000000,0.000000}%
\pgfsetfillcolor{currentfill}%
\pgfsetlinewidth{0.000000pt}%
\definecolor{currentstroke}{rgb}{0.000000,0.000000,0.000000}%
\pgfsetstrokecolor{currentstroke}%
\pgfsetstrokeopacity{0.000000}%
\pgfsetdash{}{0pt}%
\pgfpathmoveto{\pgfqpoint{5.222650in}{0.660000in}}%
\pgfpathlineto{\pgfqpoint{5.228850in}{0.660000in}}%
\pgfpathlineto{\pgfqpoint{5.228850in}{3.250785in}}%
\pgfpathlineto{\pgfqpoint{5.222650in}{3.250785in}}%
\pgfpathlineto{\pgfqpoint{5.222650in}{0.660000in}}%
\pgfpathclose%
\pgfusepath{fill}%
\end{pgfscope}%
\begin{pgfscope}%
\pgfpathrectangle{\pgfqpoint{1.250000in}{0.660000in}}{\pgfqpoint{7.750000in}{4.620000in}}%
\pgfusepath{clip}%
\pgfsetbuttcap%
\pgfsetmiterjoin%
\definecolor{currentfill}{rgb}{0.000000,0.000000,0.000000}%
\pgfsetfillcolor{currentfill}%
\pgfsetlinewidth{0.000000pt}%
\definecolor{currentstroke}{rgb}{0.000000,0.000000,0.000000}%
\pgfsetstrokecolor{currentstroke}%
\pgfsetstrokeopacity{0.000000}%
\pgfsetdash{}{0pt}%
\pgfpathmoveto{\pgfqpoint{5.230400in}{0.660000in}}%
\pgfpathlineto{\pgfqpoint{5.236600in}{0.660000in}}%
\pgfpathlineto{\pgfqpoint{5.236600in}{3.250589in}}%
\pgfpathlineto{\pgfqpoint{5.230400in}{3.250589in}}%
\pgfpathlineto{\pgfqpoint{5.230400in}{0.660000in}}%
\pgfpathclose%
\pgfusepath{fill}%
\end{pgfscope}%
\begin{pgfscope}%
\pgfpathrectangle{\pgfqpoint{1.250000in}{0.660000in}}{\pgfqpoint{7.750000in}{4.620000in}}%
\pgfusepath{clip}%
\pgfsetbuttcap%
\pgfsetmiterjoin%
\definecolor{currentfill}{rgb}{0.000000,0.000000,0.000000}%
\pgfsetfillcolor{currentfill}%
\pgfsetlinewidth{0.000000pt}%
\definecolor{currentstroke}{rgb}{0.000000,0.000000,0.000000}%
\pgfsetstrokecolor{currentstroke}%
\pgfsetstrokeopacity{0.000000}%
\pgfsetdash{}{0pt}%
\pgfpathmoveto{\pgfqpoint{5.238150in}{0.660000in}}%
\pgfpathlineto{\pgfqpoint{5.244350in}{0.660000in}}%
\pgfpathlineto{\pgfqpoint{5.244350in}{3.250397in}}%
\pgfpathlineto{\pgfqpoint{5.238150in}{3.250397in}}%
\pgfpathlineto{\pgfqpoint{5.238150in}{0.660000in}}%
\pgfpathclose%
\pgfusepath{fill}%
\end{pgfscope}%
\begin{pgfscope}%
\pgfpathrectangle{\pgfqpoint{1.250000in}{0.660000in}}{\pgfqpoint{7.750000in}{4.620000in}}%
\pgfusepath{clip}%
\pgfsetbuttcap%
\pgfsetmiterjoin%
\definecolor{currentfill}{rgb}{0.000000,0.000000,0.000000}%
\pgfsetfillcolor{currentfill}%
\pgfsetlinewidth{0.000000pt}%
\definecolor{currentstroke}{rgb}{0.000000,0.000000,0.000000}%
\pgfsetstrokecolor{currentstroke}%
\pgfsetstrokeopacity{0.000000}%
\pgfsetdash{}{0pt}%
\pgfpathmoveto{\pgfqpoint{5.245900in}{0.660000in}}%
\pgfpathlineto{\pgfqpoint{5.252100in}{0.660000in}}%
\pgfpathlineto{\pgfqpoint{5.252100in}{3.250204in}}%
\pgfpathlineto{\pgfqpoint{5.245900in}{3.250204in}}%
\pgfpathlineto{\pgfqpoint{5.245900in}{0.660000in}}%
\pgfpathclose%
\pgfusepath{fill}%
\end{pgfscope}%
\begin{pgfscope}%
\pgfpathrectangle{\pgfqpoint{1.250000in}{0.660000in}}{\pgfqpoint{7.750000in}{4.620000in}}%
\pgfusepath{clip}%
\pgfsetbuttcap%
\pgfsetmiterjoin%
\definecolor{currentfill}{rgb}{0.000000,0.000000,0.000000}%
\pgfsetfillcolor{currentfill}%
\pgfsetlinewidth{0.000000pt}%
\definecolor{currentstroke}{rgb}{0.000000,0.000000,0.000000}%
\pgfsetstrokecolor{currentstroke}%
\pgfsetstrokeopacity{0.000000}%
\pgfsetdash{}{0pt}%
\pgfpathmoveto{\pgfqpoint{5.253650in}{0.660000in}}%
\pgfpathlineto{\pgfqpoint{5.259850in}{0.660000in}}%
\pgfpathlineto{\pgfqpoint{5.259850in}{3.250011in}}%
\pgfpathlineto{\pgfqpoint{5.253650in}{3.250011in}}%
\pgfpathlineto{\pgfqpoint{5.253650in}{0.660000in}}%
\pgfpathclose%
\pgfusepath{fill}%
\end{pgfscope}%
\begin{pgfscope}%
\pgfpathrectangle{\pgfqpoint{1.250000in}{0.660000in}}{\pgfqpoint{7.750000in}{4.620000in}}%
\pgfusepath{clip}%
\pgfsetbuttcap%
\pgfsetmiterjoin%
\definecolor{currentfill}{rgb}{0.000000,0.000000,0.000000}%
\pgfsetfillcolor{currentfill}%
\pgfsetlinewidth{0.000000pt}%
\definecolor{currentstroke}{rgb}{0.000000,0.000000,0.000000}%
\pgfsetstrokecolor{currentstroke}%
\pgfsetstrokeopacity{0.000000}%
\pgfsetdash{}{0pt}%
\pgfpathmoveto{\pgfqpoint{5.261400in}{0.660000in}}%
\pgfpathlineto{\pgfqpoint{5.267600in}{0.660000in}}%
\pgfpathlineto{\pgfqpoint{5.267600in}{3.249818in}}%
\pgfpathlineto{\pgfqpoint{5.261400in}{3.249818in}}%
\pgfpathlineto{\pgfqpoint{5.261400in}{0.660000in}}%
\pgfpathclose%
\pgfusepath{fill}%
\end{pgfscope}%
\begin{pgfscope}%
\pgfpathrectangle{\pgfqpoint{1.250000in}{0.660000in}}{\pgfqpoint{7.750000in}{4.620000in}}%
\pgfusepath{clip}%
\pgfsetbuttcap%
\pgfsetmiterjoin%
\definecolor{currentfill}{rgb}{0.000000,0.000000,0.000000}%
\pgfsetfillcolor{currentfill}%
\pgfsetlinewidth{0.000000pt}%
\definecolor{currentstroke}{rgb}{0.000000,0.000000,0.000000}%
\pgfsetstrokecolor{currentstroke}%
\pgfsetstrokeopacity{0.000000}%
\pgfsetdash{}{0pt}%
\pgfpathmoveto{\pgfqpoint{5.269150in}{0.660000in}}%
\pgfpathlineto{\pgfqpoint{5.275350in}{0.660000in}}%
\pgfpathlineto{\pgfqpoint{5.275350in}{3.249625in}}%
\pgfpathlineto{\pgfqpoint{5.269150in}{3.249625in}}%
\pgfpathlineto{\pgfqpoint{5.269150in}{0.660000in}}%
\pgfpathclose%
\pgfusepath{fill}%
\end{pgfscope}%
\begin{pgfscope}%
\pgfpathrectangle{\pgfqpoint{1.250000in}{0.660000in}}{\pgfqpoint{7.750000in}{4.620000in}}%
\pgfusepath{clip}%
\pgfsetbuttcap%
\pgfsetmiterjoin%
\definecolor{currentfill}{rgb}{0.000000,0.000000,0.000000}%
\pgfsetfillcolor{currentfill}%
\pgfsetlinewidth{0.000000pt}%
\definecolor{currentstroke}{rgb}{0.000000,0.000000,0.000000}%
\pgfsetstrokecolor{currentstroke}%
\pgfsetstrokeopacity{0.000000}%
\pgfsetdash{}{0pt}%
\pgfpathmoveto{\pgfqpoint{5.276900in}{0.660000in}}%
\pgfpathlineto{\pgfqpoint{5.283100in}{0.660000in}}%
\pgfpathlineto{\pgfqpoint{5.283100in}{3.249432in}}%
\pgfpathlineto{\pgfqpoint{5.276900in}{3.249432in}}%
\pgfpathlineto{\pgfqpoint{5.276900in}{0.660000in}}%
\pgfpathclose%
\pgfusepath{fill}%
\end{pgfscope}%
\begin{pgfscope}%
\pgfpathrectangle{\pgfqpoint{1.250000in}{0.660000in}}{\pgfqpoint{7.750000in}{4.620000in}}%
\pgfusepath{clip}%
\pgfsetbuttcap%
\pgfsetmiterjoin%
\definecolor{currentfill}{rgb}{0.000000,0.000000,0.000000}%
\pgfsetfillcolor{currentfill}%
\pgfsetlinewidth{0.000000pt}%
\definecolor{currentstroke}{rgb}{0.000000,0.000000,0.000000}%
\pgfsetstrokecolor{currentstroke}%
\pgfsetstrokeopacity{0.000000}%
\pgfsetdash{}{0pt}%
\pgfpathmoveto{\pgfqpoint{5.284650in}{0.660000in}}%
\pgfpathlineto{\pgfqpoint{5.290850in}{0.660000in}}%
\pgfpathlineto{\pgfqpoint{5.290850in}{3.249239in}}%
\pgfpathlineto{\pgfqpoint{5.284650in}{3.249239in}}%
\pgfpathlineto{\pgfqpoint{5.284650in}{0.660000in}}%
\pgfpathclose%
\pgfusepath{fill}%
\end{pgfscope}%
\begin{pgfscope}%
\pgfpathrectangle{\pgfqpoint{1.250000in}{0.660000in}}{\pgfqpoint{7.750000in}{4.620000in}}%
\pgfusepath{clip}%
\pgfsetbuttcap%
\pgfsetmiterjoin%
\definecolor{currentfill}{rgb}{0.000000,0.000000,0.000000}%
\pgfsetfillcolor{currentfill}%
\pgfsetlinewidth{0.000000pt}%
\definecolor{currentstroke}{rgb}{0.000000,0.000000,0.000000}%
\pgfsetstrokecolor{currentstroke}%
\pgfsetstrokeopacity{0.000000}%
\pgfsetdash{}{0pt}%
\pgfpathmoveto{\pgfqpoint{5.292400in}{0.660000in}}%
\pgfpathlineto{\pgfqpoint{5.298600in}{0.660000in}}%
\pgfpathlineto{\pgfqpoint{5.298600in}{3.249046in}}%
\pgfpathlineto{\pgfqpoint{5.292400in}{3.249046in}}%
\pgfpathlineto{\pgfqpoint{5.292400in}{0.660000in}}%
\pgfpathclose%
\pgfusepath{fill}%
\end{pgfscope}%
\begin{pgfscope}%
\pgfpathrectangle{\pgfqpoint{1.250000in}{0.660000in}}{\pgfqpoint{7.750000in}{4.620000in}}%
\pgfusepath{clip}%
\pgfsetbuttcap%
\pgfsetmiterjoin%
\definecolor{currentfill}{rgb}{0.000000,0.000000,0.000000}%
\pgfsetfillcolor{currentfill}%
\pgfsetlinewidth{0.000000pt}%
\definecolor{currentstroke}{rgb}{0.000000,0.000000,0.000000}%
\pgfsetstrokecolor{currentstroke}%
\pgfsetstrokeopacity{0.000000}%
\pgfsetdash{}{0pt}%
\pgfpathmoveto{\pgfqpoint{5.300150in}{0.660000in}}%
\pgfpathlineto{\pgfqpoint{5.306350in}{0.660000in}}%
\pgfpathlineto{\pgfqpoint{5.306350in}{3.248854in}}%
\pgfpathlineto{\pgfqpoint{5.300150in}{3.248854in}}%
\pgfpathlineto{\pgfqpoint{5.300150in}{0.660000in}}%
\pgfpathclose%
\pgfusepath{fill}%
\end{pgfscope}%
\begin{pgfscope}%
\pgfpathrectangle{\pgfqpoint{1.250000in}{0.660000in}}{\pgfqpoint{7.750000in}{4.620000in}}%
\pgfusepath{clip}%
\pgfsetbuttcap%
\pgfsetmiterjoin%
\definecolor{currentfill}{rgb}{0.000000,0.000000,0.000000}%
\pgfsetfillcolor{currentfill}%
\pgfsetlinewidth{0.000000pt}%
\definecolor{currentstroke}{rgb}{0.000000,0.000000,0.000000}%
\pgfsetstrokecolor{currentstroke}%
\pgfsetstrokeopacity{0.000000}%
\pgfsetdash{}{0pt}%
\pgfpathmoveto{\pgfqpoint{5.307900in}{0.660000in}}%
\pgfpathlineto{\pgfqpoint{5.314100in}{0.660000in}}%
\pgfpathlineto{\pgfqpoint{5.314100in}{3.248666in}}%
\pgfpathlineto{\pgfqpoint{5.307900in}{3.248666in}}%
\pgfpathlineto{\pgfqpoint{5.307900in}{0.660000in}}%
\pgfpathclose%
\pgfusepath{fill}%
\end{pgfscope}%
\begin{pgfscope}%
\pgfpathrectangle{\pgfqpoint{1.250000in}{0.660000in}}{\pgfqpoint{7.750000in}{4.620000in}}%
\pgfusepath{clip}%
\pgfsetbuttcap%
\pgfsetmiterjoin%
\definecolor{currentfill}{rgb}{0.000000,0.000000,0.000000}%
\pgfsetfillcolor{currentfill}%
\pgfsetlinewidth{0.000000pt}%
\definecolor{currentstroke}{rgb}{0.000000,0.000000,0.000000}%
\pgfsetstrokecolor{currentstroke}%
\pgfsetstrokeopacity{0.000000}%
\pgfsetdash{}{0pt}%
\pgfpathmoveto{\pgfqpoint{5.315650in}{0.660000in}}%
\pgfpathlineto{\pgfqpoint{5.321850in}{0.660000in}}%
\pgfpathlineto{\pgfqpoint{5.321850in}{3.248477in}}%
\pgfpathlineto{\pgfqpoint{5.315650in}{3.248477in}}%
\pgfpathlineto{\pgfqpoint{5.315650in}{0.660000in}}%
\pgfpathclose%
\pgfusepath{fill}%
\end{pgfscope}%
\begin{pgfscope}%
\pgfpathrectangle{\pgfqpoint{1.250000in}{0.660000in}}{\pgfqpoint{7.750000in}{4.620000in}}%
\pgfusepath{clip}%
\pgfsetbuttcap%
\pgfsetmiterjoin%
\definecolor{currentfill}{rgb}{0.000000,0.000000,0.000000}%
\pgfsetfillcolor{currentfill}%
\pgfsetlinewidth{0.000000pt}%
\definecolor{currentstroke}{rgb}{0.000000,0.000000,0.000000}%
\pgfsetstrokecolor{currentstroke}%
\pgfsetstrokeopacity{0.000000}%
\pgfsetdash{}{0pt}%
\pgfpathmoveto{\pgfqpoint{5.323400in}{0.660000in}}%
\pgfpathlineto{\pgfqpoint{5.329600in}{0.660000in}}%
\pgfpathlineto{\pgfqpoint{5.329600in}{3.248289in}}%
\pgfpathlineto{\pgfqpoint{5.323400in}{3.248289in}}%
\pgfpathlineto{\pgfqpoint{5.323400in}{0.660000in}}%
\pgfpathclose%
\pgfusepath{fill}%
\end{pgfscope}%
\begin{pgfscope}%
\pgfpathrectangle{\pgfqpoint{1.250000in}{0.660000in}}{\pgfqpoint{7.750000in}{4.620000in}}%
\pgfusepath{clip}%
\pgfsetbuttcap%
\pgfsetmiterjoin%
\definecolor{currentfill}{rgb}{0.000000,0.000000,0.000000}%
\pgfsetfillcolor{currentfill}%
\pgfsetlinewidth{0.000000pt}%
\definecolor{currentstroke}{rgb}{0.000000,0.000000,0.000000}%
\pgfsetstrokecolor{currentstroke}%
\pgfsetstrokeopacity{0.000000}%
\pgfsetdash{}{0pt}%
\pgfpathmoveto{\pgfqpoint{5.331150in}{0.660000in}}%
\pgfpathlineto{\pgfqpoint{5.337350in}{0.660000in}}%
\pgfpathlineto{\pgfqpoint{5.337350in}{3.248101in}}%
\pgfpathlineto{\pgfqpoint{5.331150in}{3.248101in}}%
\pgfpathlineto{\pgfqpoint{5.331150in}{0.660000in}}%
\pgfpathclose%
\pgfusepath{fill}%
\end{pgfscope}%
\begin{pgfscope}%
\pgfpathrectangle{\pgfqpoint{1.250000in}{0.660000in}}{\pgfqpoint{7.750000in}{4.620000in}}%
\pgfusepath{clip}%
\pgfsetbuttcap%
\pgfsetmiterjoin%
\definecolor{currentfill}{rgb}{0.000000,0.000000,0.000000}%
\pgfsetfillcolor{currentfill}%
\pgfsetlinewidth{0.000000pt}%
\definecolor{currentstroke}{rgb}{0.000000,0.000000,0.000000}%
\pgfsetstrokecolor{currentstroke}%
\pgfsetstrokeopacity{0.000000}%
\pgfsetdash{}{0pt}%
\pgfpathmoveto{\pgfqpoint{5.338900in}{0.660000in}}%
\pgfpathlineto{\pgfqpoint{5.345100in}{0.660000in}}%
\pgfpathlineto{\pgfqpoint{5.345100in}{3.247913in}}%
\pgfpathlineto{\pgfqpoint{5.338900in}{3.247913in}}%
\pgfpathlineto{\pgfqpoint{5.338900in}{0.660000in}}%
\pgfpathclose%
\pgfusepath{fill}%
\end{pgfscope}%
\begin{pgfscope}%
\pgfpathrectangle{\pgfqpoint{1.250000in}{0.660000in}}{\pgfqpoint{7.750000in}{4.620000in}}%
\pgfusepath{clip}%
\pgfsetbuttcap%
\pgfsetmiterjoin%
\definecolor{currentfill}{rgb}{0.000000,0.000000,0.000000}%
\pgfsetfillcolor{currentfill}%
\pgfsetlinewidth{0.000000pt}%
\definecolor{currentstroke}{rgb}{0.000000,0.000000,0.000000}%
\pgfsetstrokecolor{currentstroke}%
\pgfsetstrokeopacity{0.000000}%
\pgfsetdash{}{0pt}%
\pgfpathmoveto{\pgfqpoint{5.346650in}{0.660000in}}%
\pgfpathlineto{\pgfqpoint{5.352850in}{0.660000in}}%
\pgfpathlineto{\pgfqpoint{5.352850in}{3.247725in}}%
\pgfpathlineto{\pgfqpoint{5.346650in}{3.247725in}}%
\pgfpathlineto{\pgfqpoint{5.346650in}{0.660000in}}%
\pgfpathclose%
\pgfusepath{fill}%
\end{pgfscope}%
\begin{pgfscope}%
\pgfpathrectangle{\pgfqpoint{1.250000in}{0.660000in}}{\pgfqpoint{7.750000in}{4.620000in}}%
\pgfusepath{clip}%
\pgfsetbuttcap%
\pgfsetmiterjoin%
\definecolor{currentfill}{rgb}{0.000000,0.000000,0.000000}%
\pgfsetfillcolor{currentfill}%
\pgfsetlinewidth{0.000000pt}%
\definecolor{currentstroke}{rgb}{0.000000,0.000000,0.000000}%
\pgfsetstrokecolor{currentstroke}%
\pgfsetstrokeopacity{0.000000}%
\pgfsetdash{}{0pt}%
\pgfpathmoveto{\pgfqpoint{5.354400in}{0.660000in}}%
\pgfpathlineto{\pgfqpoint{5.360600in}{0.660000in}}%
\pgfpathlineto{\pgfqpoint{5.360600in}{3.247536in}}%
\pgfpathlineto{\pgfqpoint{5.354400in}{3.247536in}}%
\pgfpathlineto{\pgfqpoint{5.354400in}{0.660000in}}%
\pgfpathclose%
\pgfusepath{fill}%
\end{pgfscope}%
\begin{pgfscope}%
\pgfpathrectangle{\pgfqpoint{1.250000in}{0.660000in}}{\pgfqpoint{7.750000in}{4.620000in}}%
\pgfusepath{clip}%
\pgfsetbuttcap%
\pgfsetmiterjoin%
\definecolor{currentfill}{rgb}{0.000000,0.000000,0.000000}%
\pgfsetfillcolor{currentfill}%
\pgfsetlinewidth{0.000000pt}%
\definecolor{currentstroke}{rgb}{0.000000,0.000000,0.000000}%
\pgfsetstrokecolor{currentstroke}%
\pgfsetstrokeopacity{0.000000}%
\pgfsetdash{}{0pt}%
\pgfpathmoveto{\pgfqpoint{5.362150in}{0.660000in}}%
\pgfpathlineto{\pgfqpoint{5.368350in}{0.660000in}}%
\pgfpathlineto{\pgfqpoint{5.368350in}{3.247348in}}%
\pgfpathlineto{\pgfqpoint{5.362150in}{3.247348in}}%
\pgfpathlineto{\pgfqpoint{5.362150in}{0.660000in}}%
\pgfpathclose%
\pgfusepath{fill}%
\end{pgfscope}%
\begin{pgfscope}%
\pgfpathrectangle{\pgfqpoint{1.250000in}{0.660000in}}{\pgfqpoint{7.750000in}{4.620000in}}%
\pgfusepath{clip}%
\pgfsetbuttcap%
\pgfsetmiterjoin%
\definecolor{currentfill}{rgb}{0.000000,0.000000,0.000000}%
\pgfsetfillcolor{currentfill}%
\pgfsetlinewidth{0.000000pt}%
\definecolor{currentstroke}{rgb}{0.000000,0.000000,0.000000}%
\pgfsetstrokecolor{currentstroke}%
\pgfsetstrokeopacity{0.000000}%
\pgfsetdash{}{0pt}%
\pgfpathmoveto{\pgfqpoint{5.369900in}{0.660000in}}%
\pgfpathlineto{\pgfqpoint{5.376100in}{0.660000in}}%
\pgfpathlineto{\pgfqpoint{5.376100in}{3.247160in}}%
\pgfpathlineto{\pgfqpoint{5.369900in}{3.247160in}}%
\pgfpathlineto{\pgfqpoint{5.369900in}{0.660000in}}%
\pgfpathclose%
\pgfusepath{fill}%
\end{pgfscope}%
\begin{pgfscope}%
\pgfpathrectangle{\pgfqpoint{1.250000in}{0.660000in}}{\pgfqpoint{7.750000in}{4.620000in}}%
\pgfusepath{clip}%
\pgfsetbuttcap%
\pgfsetmiterjoin%
\definecolor{currentfill}{rgb}{0.000000,0.000000,0.000000}%
\pgfsetfillcolor{currentfill}%
\pgfsetlinewidth{0.000000pt}%
\definecolor{currentstroke}{rgb}{0.000000,0.000000,0.000000}%
\pgfsetstrokecolor{currentstroke}%
\pgfsetstrokeopacity{0.000000}%
\pgfsetdash{}{0pt}%
\pgfpathmoveto{\pgfqpoint{5.377650in}{0.660000in}}%
\pgfpathlineto{\pgfqpoint{5.383850in}{0.660000in}}%
\pgfpathlineto{\pgfqpoint{5.383850in}{3.246974in}}%
\pgfpathlineto{\pgfqpoint{5.377650in}{3.246974in}}%
\pgfpathlineto{\pgfqpoint{5.377650in}{0.660000in}}%
\pgfpathclose%
\pgfusepath{fill}%
\end{pgfscope}%
\begin{pgfscope}%
\pgfpathrectangle{\pgfqpoint{1.250000in}{0.660000in}}{\pgfqpoint{7.750000in}{4.620000in}}%
\pgfusepath{clip}%
\pgfsetbuttcap%
\pgfsetmiterjoin%
\definecolor{currentfill}{rgb}{0.000000,0.000000,0.000000}%
\pgfsetfillcolor{currentfill}%
\pgfsetlinewidth{0.000000pt}%
\definecolor{currentstroke}{rgb}{0.000000,0.000000,0.000000}%
\pgfsetstrokecolor{currentstroke}%
\pgfsetstrokeopacity{0.000000}%
\pgfsetdash{}{0pt}%
\pgfpathmoveto{\pgfqpoint{5.385400in}{0.660000in}}%
\pgfpathlineto{\pgfqpoint{5.391600in}{0.660000in}}%
\pgfpathlineto{\pgfqpoint{5.391600in}{3.246791in}}%
\pgfpathlineto{\pgfqpoint{5.385400in}{3.246791in}}%
\pgfpathlineto{\pgfqpoint{5.385400in}{0.660000in}}%
\pgfpathclose%
\pgfusepath{fill}%
\end{pgfscope}%
\begin{pgfscope}%
\pgfpathrectangle{\pgfqpoint{1.250000in}{0.660000in}}{\pgfqpoint{7.750000in}{4.620000in}}%
\pgfusepath{clip}%
\pgfsetbuttcap%
\pgfsetmiterjoin%
\definecolor{currentfill}{rgb}{0.000000,0.000000,0.000000}%
\pgfsetfillcolor{currentfill}%
\pgfsetlinewidth{0.000000pt}%
\definecolor{currentstroke}{rgb}{0.000000,0.000000,0.000000}%
\pgfsetstrokecolor{currentstroke}%
\pgfsetstrokeopacity{0.000000}%
\pgfsetdash{}{0pt}%
\pgfpathmoveto{\pgfqpoint{5.393150in}{0.660000in}}%
\pgfpathlineto{\pgfqpoint{5.399350in}{0.660000in}}%
\pgfpathlineto{\pgfqpoint{5.399350in}{3.246607in}}%
\pgfpathlineto{\pgfqpoint{5.393150in}{3.246607in}}%
\pgfpathlineto{\pgfqpoint{5.393150in}{0.660000in}}%
\pgfpathclose%
\pgfusepath{fill}%
\end{pgfscope}%
\begin{pgfscope}%
\pgfpathrectangle{\pgfqpoint{1.250000in}{0.660000in}}{\pgfqpoint{7.750000in}{4.620000in}}%
\pgfusepath{clip}%
\pgfsetbuttcap%
\pgfsetmiterjoin%
\definecolor{currentfill}{rgb}{0.000000,0.000000,0.000000}%
\pgfsetfillcolor{currentfill}%
\pgfsetlinewidth{0.000000pt}%
\definecolor{currentstroke}{rgb}{0.000000,0.000000,0.000000}%
\pgfsetstrokecolor{currentstroke}%
\pgfsetstrokeopacity{0.000000}%
\pgfsetdash{}{0pt}%
\pgfpathmoveto{\pgfqpoint{5.400900in}{0.660000in}}%
\pgfpathlineto{\pgfqpoint{5.407100in}{0.660000in}}%
\pgfpathlineto{\pgfqpoint{5.407100in}{3.246424in}}%
\pgfpathlineto{\pgfqpoint{5.400900in}{3.246424in}}%
\pgfpathlineto{\pgfqpoint{5.400900in}{0.660000in}}%
\pgfpathclose%
\pgfusepath{fill}%
\end{pgfscope}%
\begin{pgfscope}%
\pgfpathrectangle{\pgfqpoint{1.250000in}{0.660000in}}{\pgfqpoint{7.750000in}{4.620000in}}%
\pgfusepath{clip}%
\pgfsetbuttcap%
\pgfsetmiterjoin%
\definecolor{currentfill}{rgb}{0.000000,0.000000,0.000000}%
\pgfsetfillcolor{currentfill}%
\pgfsetlinewidth{0.000000pt}%
\definecolor{currentstroke}{rgb}{0.000000,0.000000,0.000000}%
\pgfsetstrokecolor{currentstroke}%
\pgfsetstrokeopacity{0.000000}%
\pgfsetdash{}{0pt}%
\pgfpathmoveto{\pgfqpoint{5.408650in}{0.660000in}}%
\pgfpathlineto{\pgfqpoint{5.414850in}{0.660000in}}%
\pgfpathlineto{\pgfqpoint{5.414850in}{3.246240in}}%
\pgfpathlineto{\pgfqpoint{5.408650in}{3.246240in}}%
\pgfpathlineto{\pgfqpoint{5.408650in}{0.660000in}}%
\pgfpathclose%
\pgfusepath{fill}%
\end{pgfscope}%
\begin{pgfscope}%
\pgfpathrectangle{\pgfqpoint{1.250000in}{0.660000in}}{\pgfqpoint{7.750000in}{4.620000in}}%
\pgfusepath{clip}%
\pgfsetbuttcap%
\pgfsetmiterjoin%
\definecolor{currentfill}{rgb}{0.000000,0.000000,0.000000}%
\pgfsetfillcolor{currentfill}%
\pgfsetlinewidth{0.000000pt}%
\definecolor{currentstroke}{rgb}{0.000000,0.000000,0.000000}%
\pgfsetstrokecolor{currentstroke}%
\pgfsetstrokeopacity{0.000000}%
\pgfsetdash{}{0pt}%
\pgfpathmoveto{\pgfqpoint{5.416400in}{0.660000in}}%
\pgfpathlineto{\pgfqpoint{5.422600in}{0.660000in}}%
\pgfpathlineto{\pgfqpoint{5.422600in}{3.246057in}}%
\pgfpathlineto{\pgfqpoint{5.416400in}{3.246057in}}%
\pgfpathlineto{\pgfqpoint{5.416400in}{0.660000in}}%
\pgfpathclose%
\pgfusepath{fill}%
\end{pgfscope}%
\begin{pgfscope}%
\pgfpathrectangle{\pgfqpoint{1.250000in}{0.660000in}}{\pgfqpoint{7.750000in}{4.620000in}}%
\pgfusepath{clip}%
\pgfsetbuttcap%
\pgfsetmiterjoin%
\definecolor{currentfill}{rgb}{0.000000,0.000000,0.000000}%
\pgfsetfillcolor{currentfill}%
\pgfsetlinewidth{0.000000pt}%
\definecolor{currentstroke}{rgb}{0.000000,0.000000,0.000000}%
\pgfsetstrokecolor{currentstroke}%
\pgfsetstrokeopacity{0.000000}%
\pgfsetdash{}{0pt}%
\pgfpathmoveto{\pgfqpoint{5.424150in}{0.660000in}}%
\pgfpathlineto{\pgfqpoint{5.430350in}{0.660000in}}%
\pgfpathlineto{\pgfqpoint{5.430350in}{3.245873in}}%
\pgfpathlineto{\pgfqpoint{5.424150in}{3.245873in}}%
\pgfpathlineto{\pgfqpoint{5.424150in}{0.660000in}}%
\pgfpathclose%
\pgfusepath{fill}%
\end{pgfscope}%
\begin{pgfscope}%
\pgfpathrectangle{\pgfqpoint{1.250000in}{0.660000in}}{\pgfqpoint{7.750000in}{4.620000in}}%
\pgfusepath{clip}%
\pgfsetbuttcap%
\pgfsetmiterjoin%
\definecolor{currentfill}{rgb}{0.000000,0.000000,0.000000}%
\pgfsetfillcolor{currentfill}%
\pgfsetlinewidth{0.000000pt}%
\definecolor{currentstroke}{rgb}{0.000000,0.000000,0.000000}%
\pgfsetstrokecolor{currentstroke}%
\pgfsetstrokeopacity{0.000000}%
\pgfsetdash{}{0pt}%
\pgfpathmoveto{\pgfqpoint{5.431900in}{0.660000in}}%
\pgfpathlineto{\pgfqpoint{5.438100in}{0.660000in}}%
\pgfpathlineto{\pgfqpoint{5.438100in}{3.245690in}}%
\pgfpathlineto{\pgfqpoint{5.431900in}{3.245690in}}%
\pgfpathlineto{\pgfqpoint{5.431900in}{0.660000in}}%
\pgfpathclose%
\pgfusepath{fill}%
\end{pgfscope}%
\begin{pgfscope}%
\pgfpathrectangle{\pgfqpoint{1.250000in}{0.660000in}}{\pgfqpoint{7.750000in}{4.620000in}}%
\pgfusepath{clip}%
\pgfsetbuttcap%
\pgfsetmiterjoin%
\definecolor{currentfill}{rgb}{0.000000,0.000000,0.000000}%
\pgfsetfillcolor{currentfill}%
\pgfsetlinewidth{0.000000pt}%
\definecolor{currentstroke}{rgb}{0.000000,0.000000,0.000000}%
\pgfsetstrokecolor{currentstroke}%
\pgfsetstrokeopacity{0.000000}%
\pgfsetdash{}{0pt}%
\pgfpathmoveto{\pgfqpoint{5.439650in}{0.660000in}}%
\pgfpathlineto{\pgfqpoint{5.445850in}{0.660000in}}%
\pgfpathlineto{\pgfqpoint{5.445850in}{3.245506in}}%
\pgfpathlineto{\pgfqpoint{5.439650in}{3.245506in}}%
\pgfpathlineto{\pgfqpoint{5.439650in}{0.660000in}}%
\pgfpathclose%
\pgfusepath{fill}%
\end{pgfscope}%
\begin{pgfscope}%
\pgfpathrectangle{\pgfqpoint{1.250000in}{0.660000in}}{\pgfqpoint{7.750000in}{4.620000in}}%
\pgfusepath{clip}%
\pgfsetbuttcap%
\pgfsetmiterjoin%
\definecolor{currentfill}{rgb}{0.000000,0.000000,0.000000}%
\pgfsetfillcolor{currentfill}%
\pgfsetlinewidth{0.000000pt}%
\definecolor{currentstroke}{rgb}{0.000000,0.000000,0.000000}%
\pgfsetstrokecolor{currentstroke}%
\pgfsetstrokeopacity{0.000000}%
\pgfsetdash{}{0pt}%
\pgfpathmoveto{\pgfqpoint{5.447400in}{0.660000in}}%
\pgfpathlineto{\pgfqpoint{5.453600in}{0.660000in}}%
\pgfpathlineto{\pgfqpoint{5.453600in}{3.245323in}}%
\pgfpathlineto{\pgfqpoint{5.447400in}{3.245323in}}%
\pgfpathlineto{\pgfqpoint{5.447400in}{0.660000in}}%
\pgfpathclose%
\pgfusepath{fill}%
\end{pgfscope}%
\begin{pgfscope}%
\pgfpathrectangle{\pgfqpoint{1.250000in}{0.660000in}}{\pgfqpoint{7.750000in}{4.620000in}}%
\pgfusepath{clip}%
\pgfsetbuttcap%
\pgfsetmiterjoin%
\definecolor{currentfill}{rgb}{0.000000,0.000000,0.000000}%
\pgfsetfillcolor{currentfill}%
\pgfsetlinewidth{0.000000pt}%
\definecolor{currentstroke}{rgb}{0.000000,0.000000,0.000000}%
\pgfsetstrokecolor{currentstroke}%
\pgfsetstrokeopacity{0.000000}%
\pgfsetdash{}{0pt}%
\pgfpathmoveto{\pgfqpoint{5.455150in}{0.660000in}}%
\pgfpathlineto{\pgfqpoint{5.461350in}{0.660000in}}%
\pgfpathlineto{\pgfqpoint{5.461350in}{3.245142in}}%
\pgfpathlineto{\pgfqpoint{5.455150in}{3.245142in}}%
\pgfpathlineto{\pgfqpoint{5.455150in}{0.660000in}}%
\pgfpathclose%
\pgfusepath{fill}%
\end{pgfscope}%
\begin{pgfscope}%
\pgfpathrectangle{\pgfqpoint{1.250000in}{0.660000in}}{\pgfqpoint{7.750000in}{4.620000in}}%
\pgfusepath{clip}%
\pgfsetbuttcap%
\pgfsetmiterjoin%
\definecolor{currentfill}{rgb}{0.000000,0.000000,0.000000}%
\pgfsetfillcolor{currentfill}%
\pgfsetlinewidth{0.000000pt}%
\definecolor{currentstroke}{rgb}{0.000000,0.000000,0.000000}%
\pgfsetstrokecolor{currentstroke}%
\pgfsetstrokeopacity{0.000000}%
\pgfsetdash{}{0pt}%
\pgfpathmoveto{\pgfqpoint{5.462900in}{0.660000in}}%
\pgfpathlineto{\pgfqpoint{5.469100in}{0.660000in}}%
\pgfpathlineto{\pgfqpoint{5.469100in}{3.244963in}}%
\pgfpathlineto{\pgfqpoint{5.462900in}{3.244963in}}%
\pgfpathlineto{\pgfqpoint{5.462900in}{0.660000in}}%
\pgfpathclose%
\pgfusepath{fill}%
\end{pgfscope}%
\begin{pgfscope}%
\pgfpathrectangle{\pgfqpoint{1.250000in}{0.660000in}}{\pgfqpoint{7.750000in}{4.620000in}}%
\pgfusepath{clip}%
\pgfsetbuttcap%
\pgfsetmiterjoin%
\definecolor{currentfill}{rgb}{0.000000,0.000000,0.000000}%
\pgfsetfillcolor{currentfill}%
\pgfsetlinewidth{0.000000pt}%
\definecolor{currentstroke}{rgb}{0.000000,0.000000,0.000000}%
\pgfsetstrokecolor{currentstroke}%
\pgfsetstrokeopacity{0.000000}%
\pgfsetdash{}{0pt}%
\pgfpathmoveto{\pgfqpoint{5.470650in}{0.660000in}}%
\pgfpathlineto{\pgfqpoint{5.476850in}{0.660000in}}%
\pgfpathlineto{\pgfqpoint{5.476850in}{3.244784in}}%
\pgfpathlineto{\pgfqpoint{5.470650in}{3.244784in}}%
\pgfpathlineto{\pgfqpoint{5.470650in}{0.660000in}}%
\pgfpathclose%
\pgfusepath{fill}%
\end{pgfscope}%
\begin{pgfscope}%
\pgfpathrectangle{\pgfqpoint{1.250000in}{0.660000in}}{\pgfqpoint{7.750000in}{4.620000in}}%
\pgfusepath{clip}%
\pgfsetbuttcap%
\pgfsetmiterjoin%
\definecolor{currentfill}{rgb}{0.000000,0.000000,0.000000}%
\pgfsetfillcolor{currentfill}%
\pgfsetlinewidth{0.000000pt}%
\definecolor{currentstroke}{rgb}{0.000000,0.000000,0.000000}%
\pgfsetstrokecolor{currentstroke}%
\pgfsetstrokeopacity{0.000000}%
\pgfsetdash{}{0pt}%
\pgfpathmoveto{\pgfqpoint{5.478400in}{0.660000in}}%
\pgfpathlineto{\pgfqpoint{5.484600in}{0.660000in}}%
\pgfpathlineto{\pgfqpoint{5.484600in}{3.244606in}}%
\pgfpathlineto{\pgfqpoint{5.478400in}{3.244606in}}%
\pgfpathlineto{\pgfqpoint{5.478400in}{0.660000in}}%
\pgfpathclose%
\pgfusepath{fill}%
\end{pgfscope}%
\begin{pgfscope}%
\pgfpathrectangle{\pgfqpoint{1.250000in}{0.660000in}}{\pgfqpoint{7.750000in}{4.620000in}}%
\pgfusepath{clip}%
\pgfsetbuttcap%
\pgfsetmiterjoin%
\definecolor{currentfill}{rgb}{0.000000,0.000000,0.000000}%
\pgfsetfillcolor{currentfill}%
\pgfsetlinewidth{0.000000pt}%
\definecolor{currentstroke}{rgb}{0.000000,0.000000,0.000000}%
\pgfsetstrokecolor{currentstroke}%
\pgfsetstrokeopacity{0.000000}%
\pgfsetdash{}{0pt}%
\pgfpathmoveto{\pgfqpoint{5.486150in}{0.660000in}}%
\pgfpathlineto{\pgfqpoint{5.492350in}{0.660000in}}%
\pgfpathlineto{\pgfqpoint{5.492350in}{3.244427in}}%
\pgfpathlineto{\pgfqpoint{5.486150in}{3.244427in}}%
\pgfpathlineto{\pgfqpoint{5.486150in}{0.660000in}}%
\pgfpathclose%
\pgfusepath{fill}%
\end{pgfscope}%
\begin{pgfscope}%
\pgfpathrectangle{\pgfqpoint{1.250000in}{0.660000in}}{\pgfqpoint{7.750000in}{4.620000in}}%
\pgfusepath{clip}%
\pgfsetbuttcap%
\pgfsetmiterjoin%
\definecolor{currentfill}{rgb}{0.000000,0.000000,0.000000}%
\pgfsetfillcolor{currentfill}%
\pgfsetlinewidth{0.000000pt}%
\definecolor{currentstroke}{rgb}{0.000000,0.000000,0.000000}%
\pgfsetstrokecolor{currentstroke}%
\pgfsetstrokeopacity{0.000000}%
\pgfsetdash{}{0pt}%
\pgfpathmoveto{\pgfqpoint{5.493900in}{0.660000in}}%
\pgfpathlineto{\pgfqpoint{5.500100in}{0.660000in}}%
\pgfpathlineto{\pgfqpoint{5.500100in}{3.244248in}}%
\pgfpathlineto{\pgfqpoint{5.493900in}{3.244248in}}%
\pgfpathlineto{\pgfqpoint{5.493900in}{0.660000in}}%
\pgfpathclose%
\pgfusepath{fill}%
\end{pgfscope}%
\begin{pgfscope}%
\pgfpathrectangle{\pgfqpoint{1.250000in}{0.660000in}}{\pgfqpoint{7.750000in}{4.620000in}}%
\pgfusepath{clip}%
\pgfsetbuttcap%
\pgfsetmiterjoin%
\definecolor{currentfill}{rgb}{0.000000,0.000000,0.000000}%
\pgfsetfillcolor{currentfill}%
\pgfsetlinewidth{0.000000pt}%
\definecolor{currentstroke}{rgb}{0.000000,0.000000,0.000000}%
\pgfsetstrokecolor{currentstroke}%
\pgfsetstrokeopacity{0.000000}%
\pgfsetdash{}{0pt}%
\pgfpathmoveto{\pgfqpoint{5.501650in}{0.660000in}}%
\pgfpathlineto{\pgfqpoint{5.507850in}{0.660000in}}%
\pgfpathlineto{\pgfqpoint{5.507850in}{3.244069in}}%
\pgfpathlineto{\pgfqpoint{5.501650in}{3.244069in}}%
\pgfpathlineto{\pgfqpoint{5.501650in}{0.660000in}}%
\pgfpathclose%
\pgfusepath{fill}%
\end{pgfscope}%
\begin{pgfscope}%
\pgfpathrectangle{\pgfqpoint{1.250000in}{0.660000in}}{\pgfqpoint{7.750000in}{4.620000in}}%
\pgfusepath{clip}%
\pgfsetbuttcap%
\pgfsetmiterjoin%
\definecolor{currentfill}{rgb}{0.000000,0.000000,0.000000}%
\pgfsetfillcolor{currentfill}%
\pgfsetlinewidth{0.000000pt}%
\definecolor{currentstroke}{rgb}{0.000000,0.000000,0.000000}%
\pgfsetstrokecolor{currentstroke}%
\pgfsetstrokeopacity{0.000000}%
\pgfsetdash{}{0pt}%
\pgfpathmoveto{\pgfqpoint{5.509400in}{0.660000in}}%
\pgfpathlineto{\pgfqpoint{5.515600in}{0.660000in}}%
\pgfpathlineto{\pgfqpoint{5.515600in}{3.243890in}}%
\pgfpathlineto{\pgfqpoint{5.509400in}{3.243890in}}%
\pgfpathlineto{\pgfqpoint{5.509400in}{0.660000in}}%
\pgfpathclose%
\pgfusepath{fill}%
\end{pgfscope}%
\begin{pgfscope}%
\pgfpathrectangle{\pgfqpoint{1.250000in}{0.660000in}}{\pgfqpoint{7.750000in}{4.620000in}}%
\pgfusepath{clip}%
\pgfsetbuttcap%
\pgfsetmiterjoin%
\definecolor{currentfill}{rgb}{0.000000,0.000000,0.000000}%
\pgfsetfillcolor{currentfill}%
\pgfsetlinewidth{0.000000pt}%
\definecolor{currentstroke}{rgb}{0.000000,0.000000,0.000000}%
\pgfsetstrokecolor{currentstroke}%
\pgfsetstrokeopacity{0.000000}%
\pgfsetdash{}{0pt}%
\pgfpathmoveto{\pgfqpoint{5.517150in}{0.660000in}}%
\pgfpathlineto{\pgfqpoint{5.523350in}{0.660000in}}%
\pgfpathlineto{\pgfqpoint{5.523350in}{3.243711in}}%
\pgfpathlineto{\pgfqpoint{5.517150in}{3.243711in}}%
\pgfpathlineto{\pgfqpoint{5.517150in}{0.660000in}}%
\pgfpathclose%
\pgfusepath{fill}%
\end{pgfscope}%
\begin{pgfscope}%
\pgfpathrectangle{\pgfqpoint{1.250000in}{0.660000in}}{\pgfqpoint{7.750000in}{4.620000in}}%
\pgfusepath{clip}%
\pgfsetbuttcap%
\pgfsetmiterjoin%
\definecolor{currentfill}{rgb}{0.000000,0.000000,0.000000}%
\pgfsetfillcolor{currentfill}%
\pgfsetlinewidth{0.000000pt}%
\definecolor{currentstroke}{rgb}{0.000000,0.000000,0.000000}%
\pgfsetstrokecolor{currentstroke}%
\pgfsetstrokeopacity{0.000000}%
\pgfsetdash{}{0pt}%
\pgfpathmoveto{\pgfqpoint{5.524900in}{0.660000in}}%
\pgfpathlineto{\pgfqpoint{5.531100in}{0.660000in}}%
\pgfpathlineto{\pgfqpoint{5.531100in}{3.243532in}}%
\pgfpathlineto{\pgfqpoint{5.524900in}{3.243532in}}%
\pgfpathlineto{\pgfqpoint{5.524900in}{0.660000in}}%
\pgfpathclose%
\pgfusepath{fill}%
\end{pgfscope}%
\begin{pgfscope}%
\pgfpathrectangle{\pgfqpoint{1.250000in}{0.660000in}}{\pgfqpoint{7.750000in}{4.620000in}}%
\pgfusepath{clip}%
\pgfsetbuttcap%
\pgfsetmiterjoin%
\definecolor{currentfill}{rgb}{0.000000,0.000000,0.000000}%
\pgfsetfillcolor{currentfill}%
\pgfsetlinewidth{0.000000pt}%
\definecolor{currentstroke}{rgb}{0.000000,0.000000,0.000000}%
\pgfsetstrokecolor{currentstroke}%
\pgfsetstrokeopacity{0.000000}%
\pgfsetdash{}{0pt}%
\pgfpathmoveto{\pgfqpoint{5.532650in}{0.660000in}}%
\pgfpathlineto{\pgfqpoint{5.538850in}{0.660000in}}%
\pgfpathlineto{\pgfqpoint{5.538850in}{3.243355in}}%
\pgfpathlineto{\pgfqpoint{5.532650in}{3.243355in}}%
\pgfpathlineto{\pgfqpoint{5.532650in}{0.660000in}}%
\pgfpathclose%
\pgfusepath{fill}%
\end{pgfscope}%
\begin{pgfscope}%
\pgfpathrectangle{\pgfqpoint{1.250000in}{0.660000in}}{\pgfqpoint{7.750000in}{4.620000in}}%
\pgfusepath{clip}%
\pgfsetbuttcap%
\pgfsetmiterjoin%
\definecolor{currentfill}{rgb}{0.000000,0.000000,0.000000}%
\pgfsetfillcolor{currentfill}%
\pgfsetlinewidth{0.000000pt}%
\definecolor{currentstroke}{rgb}{0.000000,0.000000,0.000000}%
\pgfsetstrokecolor{currentstroke}%
\pgfsetstrokeopacity{0.000000}%
\pgfsetdash{}{0pt}%
\pgfpathmoveto{\pgfqpoint{5.540400in}{0.660000in}}%
\pgfpathlineto{\pgfqpoint{5.546600in}{0.660000in}}%
\pgfpathlineto{\pgfqpoint{5.546600in}{3.243181in}}%
\pgfpathlineto{\pgfqpoint{5.540400in}{3.243181in}}%
\pgfpathlineto{\pgfqpoint{5.540400in}{0.660000in}}%
\pgfpathclose%
\pgfusepath{fill}%
\end{pgfscope}%
\begin{pgfscope}%
\pgfpathrectangle{\pgfqpoint{1.250000in}{0.660000in}}{\pgfqpoint{7.750000in}{4.620000in}}%
\pgfusepath{clip}%
\pgfsetbuttcap%
\pgfsetmiterjoin%
\definecolor{currentfill}{rgb}{0.000000,0.000000,0.000000}%
\pgfsetfillcolor{currentfill}%
\pgfsetlinewidth{0.000000pt}%
\definecolor{currentstroke}{rgb}{0.000000,0.000000,0.000000}%
\pgfsetstrokecolor{currentstroke}%
\pgfsetstrokeopacity{0.000000}%
\pgfsetdash{}{0pt}%
\pgfpathmoveto{\pgfqpoint{5.548150in}{0.660000in}}%
\pgfpathlineto{\pgfqpoint{5.554350in}{0.660000in}}%
\pgfpathlineto{\pgfqpoint{5.554350in}{3.243007in}}%
\pgfpathlineto{\pgfqpoint{5.548150in}{3.243007in}}%
\pgfpathlineto{\pgfqpoint{5.548150in}{0.660000in}}%
\pgfpathclose%
\pgfusepath{fill}%
\end{pgfscope}%
\begin{pgfscope}%
\pgfpathrectangle{\pgfqpoint{1.250000in}{0.660000in}}{\pgfqpoint{7.750000in}{4.620000in}}%
\pgfusepath{clip}%
\pgfsetbuttcap%
\pgfsetmiterjoin%
\definecolor{currentfill}{rgb}{0.000000,0.000000,0.000000}%
\pgfsetfillcolor{currentfill}%
\pgfsetlinewidth{0.000000pt}%
\definecolor{currentstroke}{rgb}{0.000000,0.000000,0.000000}%
\pgfsetstrokecolor{currentstroke}%
\pgfsetstrokeopacity{0.000000}%
\pgfsetdash{}{0pt}%
\pgfpathmoveto{\pgfqpoint{5.555900in}{0.660000in}}%
\pgfpathlineto{\pgfqpoint{5.562100in}{0.660000in}}%
\pgfpathlineto{\pgfqpoint{5.562100in}{3.242832in}}%
\pgfpathlineto{\pgfqpoint{5.555900in}{3.242832in}}%
\pgfpathlineto{\pgfqpoint{5.555900in}{0.660000in}}%
\pgfpathclose%
\pgfusepath{fill}%
\end{pgfscope}%
\begin{pgfscope}%
\pgfpathrectangle{\pgfqpoint{1.250000in}{0.660000in}}{\pgfqpoint{7.750000in}{4.620000in}}%
\pgfusepath{clip}%
\pgfsetbuttcap%
\pgfsetmiterjoin%
\definecolor{currentfill}{rgb}{0.000000,0.000000,0.000000}%
\pgfsetfillcolor{currentfill}%
\pgfsetlinewidth{0.000000pt}%
\definecolor{currentstroke}{rgb}{0.000000,0.000000,0.000000}%
\pgfsetstrokecolor{currentstroke}%
\pgfsetstrokeopacity{0.000000}%
\pgfsetdash{}{0pt}%
\pgfpathmoveto{\pgfqpoint{5.563650in}{0.660000in}}%
\pgfpathlineto{\pgfqpoint{5.569850in}{0.660000in}}%
\pgfpathlineto{\pgfqpoint{5.569850in}{3.242658in}}%
\pgfpathlineto{\pgfqpoint{5.563650in}{3.242658in}}%
\pgfpathlineto{\pgfqpoint{5.563650in}{0.660000in}}%
\pgfpathclose%
\pgfusepath{fill}%
\end{pgfscope}%
\begin{pgfscope}%
\pgfpathrectangle{\pgfqpoint{1.250000in}{0.660000in}}{\pgfqpoint{7.750000in}{4.620000in}}%
\pgfusepath{clip}%
\pgfsetbuttcap%
\pgfsetmiterjoin%
\definecolor{currentfill}{rgb}{0.000000,0.000000,0.000000}%
\pgfsetfillcolor{currentfill}%
\pgfsetlinewidth{0.000000pt}%
\definecolor{currentstroke}{rgb}{0.000000,0.000000,0.000000}%
\pgfsetstrokecolor{currentstroke}%
\pgfsetstrokeopacity{0.000000}%
\pgfsetdash{}{0pt}%
\pgfpathmoveto{\pgfqpoint{5.571400in}{0.660000in}}%
\pgfpathlineto{\pgfqpoint{5.577600in}{0.660000in}}%
\pgfpathlineto{\pgfqpoint{5.577600in}{3.242484in}}%
\pgfpathlineto{\pgfqpoint{5.571400in}{3.242484in}}%
\pgfpathlineto{\pgfqpoint{5.571400in}{0.660000in}}%
\pgfpathclose%
\pgfusepath{fill}%
\end{pgfscope}%
\begin{pgfscope}%
\pgfpathrectangle{\pgfqpoint{1.250000in}{0.660000in}}{\pgfqpoint{7.750000in}{4.620000in}}%
\pgfusepath{clip}%
\pgfsetbuttcap%
\pgfsetmiterjoin%
\definecolor{currentfill}{rgb}{0.000000,0.000000,0.000000}%
\pgfsetfillcolor{currentfill}%
\pgfsetlinewidth{0.000000pt}%
\definecolor{currentstroke}{rgb}{0.000000,0.000000,0.000000}%
\pgfsetstrokecolor{currentstroke}%
\pgfsetstrokeopacity{0.000000}%
\pgfsetdash{}{0pt}%
\pgfpathmoveto{\pgfqpoint{5.579150in}{0.660000in}}%
\pgfpathlineto{\pgfqpoint{5.585350in}{0.660000in}}%
\pgfpathlineto{\pgfqpoint{5.585350in}{3.242309in}}%
\pgfpathlineto{\pgfqpoint{5.579150in}{3.242309in}}%
\pgfpathlineto{\pgfqpoint{5.579150in}{0.660000in}}%
\pgfpathclose%
\pgfusepath{fill}%
\end{pgfscope}%
\begin{pgfscope}%
\pgfpathrectangle{\pgfqpoint{1.250000in}{0.660000in}}{\pgfqpoint{7.750000in}{4.620000in}}%
\pgfusepath{clip}%
\pgfsetbuttcap%
\pgfsetmiterjoin%
\definecolor{currentfill}{rgb}{0.000000,0.000000,0.000000}%
\pgfsetfillcolor{currentfill}%
\pgfsetlinewidth{0.000000pt}%
\definecolor{currentstroke}{rgb}{0.000000,0.000000,0.000000}%
\pgfsetstrokecolor{currentstroke}%
\pgfsetstrokeopacity{0.000000}%
\pgfsetdash{}{0pt}%
\pgfpathmoveto{\pgfqpoint{5.586900in}{0.660000in}}%
\pgfpathlineto{\pgfqpoint{5.593100in}{0.660000in}}%
\pgfpathlineto{\pgfqpoint{5.593100in}{3.242135in}}%
\pgfpathlineto{\pgfqpoint{5.586900in}{3.242135in}}%
\pgfpathlineto{\pgfqpoint{5.586900in}{0.660000in}}%
\pgfpathclose%
\pgfusepath{fill}%
\end{pgfscope}%
\begin{pgfscope}%
\pgfpathrectangle{\pgfqpoint{1.250000in}{0.660000in}}{\pgfqpoint{7.750000in}{4.620000in}}%
\pgfusepath{clip}%
\pgfsetbuttcap%
\pgfsetmiterjoin%
\definecolor{currentfill}{rgb}{0.000000,0.000000,0.000000}%
\pgfsetfillcolor{currentfill}%
\pgfsetlinewidth{0.000000pt}%
\definecolor{currentstroke}{rgb}{0.000000,0.000000,0.000000}%
\pgfsetstrokecolor{currentstroke}%
\pgfsetstrokeopacity{0.000000}%
\pgfsetdash{}{0pt}%
\pgfpathmoveto{\pgfqpoint{5.594650in}{0.660000in}}%
\pgfpathlineto{\pgfqpoint{5.600850in}{0.660000in}}%
\pgfpathlineto{\pgfqpoint{5.600850in}{3.241960in}}%
\pgfpathlineto{\pgfqpoint{5.594650in}{3.241960in}}%
\pgfpathlineto{\pgfqpoint{5.594650in}{0.660000in}}%
\pgfpathclose%
\pgfusepath{fill}%
\end{pgfscope}%
\begin{pgfscope}%
\pgfpathrectangle{\pgfqpoint{1.250000in}{0.660000in}}{\pgfqpoint{7.750000in}{4.620000in}}%
\pgfusepath{clip}%
\pgfsetbuttcap%
\pgfsetmiterjoin%
\definecolor{currentfill}{rgb}{0.000000,0.000000,0.000000}%
\pgfsetfillcolor{currentfill}%
\pgfsetlinewidth{0.000000pt}%
\definecolor{currentstroke}{rgb}{0.000000,0.000000,0.000000}%
\pgfsetstrokecolor{currentstroke}%
\pgfsetstrokeopacity{0.000000}%
\pgfsetdash{}{0pt}%
\pgfpathmoveto{\pgfqpoint{5.602400in}{0.660000in}}%
\pgfpathlineto{\pgfqpoint{5.608600in}{0.660000in}}%
\pgfpathlineto{\pgfqpoint{5.608600in}{3.241786in}}%
\pgfpathlineto{\pgfqpoint{5.602400in}{3.241786in}}%
\pgfpathlineto{\pgfqpoint{5.602400in}{0.660000in}}%
\pgfpathclose%
\pgfusepath{fill}%
\end{pgfscope}%
\begin{pgfscope}%
\pgfpathrectangle{\pgfqpoint{1.250000in}{0.660000in}}{\pgfqpoint{7.750000in}{4.620000in}}%
\pgfusepath{clip}%
\pgfsetbuttcap%
\pgfsetmiterjoin%
\definecolor{currentfill}{rgb}{0.000000,0.000000,0.000000}%
\pgfsetfillcolor{currentfill}%
\pgfsetlinewidth{0.000000pt}%
\definecolor{currentstroke}{rgb}{0.000000,0.000000,0.000000}%
\pgfsetstrokecolor{currentstroke}%
\pgfsetstrokeopacity{0.000000}%
\pgfsetdash{}{0pt}%
\pgfpathmoveto{\pgfqpoint{5.610150in}{0.660000in}}%
\pgfpathlineto{\pgfqpoint{5.616350in}{0.660000in}}%
\pgfpathlineto{\pgfqpoint{5.616350in}{3.241612in}}%
\pgfpathlineto{\pgfqpoint{5.610150in}{3.241612in}}%
\pgfpathlineto{\pgfqpoint{5.610150in}{0.660000in}}%
\pgfpathclose%
\pgfusepath{fill}%
\end{pgfscope}%
\begin{pgfscope}%
\pgfpathrectangle{\pgfqpoint{1.250000in}{0.660000in}}{\pgfqpoint{7.750000in}{4.620000in}}%
\pgfusepath{clip}%
\pgfsetbuttcap%
\pgfsetmiterjoin%
\definecolor{currentfill}{rgb}{0.000000,0.000000,0.000000}%
\pgfsetfillcolor{currentfill}%
\pgfsetlinewidth{0.000000pt}%
\definecolor{currentstroke}{rgb}{0.000000,0.000000,0.000000}%
\pgfsetstrokecolor{currentstroke}%
\pgfsetstrokeopacity{0.000000}%
\pgfsetdash{}{0pt}%
\pgfpathmoveto{\pgfqpoint{5.617900in}{0.660000in}}%
\pgfpathlineto{\pgfqpoint{5.624100in}{0.660000in}}%
\pgfpathlineto{\pgfqpoint{5.624100in}{3.241442in}}%
\pgfpathlineto{\pgfqpoint{5.617900in}{3.241442in}}%
\pgfpathlineto{\pgfqpoint{5.617900in}{0.660000in}}%
\pgfpathclose%
\pgfusepath{fill}%
\end{pgfscope}%
\begin{pgfscope}%
\pgfpathrectangle{\pgfqpoint{1.250000in}{0.660000in}}{\pgfqpoint{7.750000in}{4.620000in}}%
\pgfusepath{clip}%
\pgfsetbuttcap%
\pgfsetmiterjoin%
\definecolor{currentfill}{rgb}{0.000000,0.000000,0.000000}%
\pgfsetfillcolor{currentfill}%
\pgfsetlinewidth{0.000000pt}%
\definecolor{currentstroke}{rgb}{0.000000,0.000000,0.000000}%
\pgfsetstrokecolor{currentstroke}%
\pgfsetstrokeopacity{0.000000}%
\pgfsetdash{}{0pt}%
\pgfpathmoveto{\pgfqpoint{5.625650in}{0.660000in}}%
\pgfpathlineto{\pgfqpoint{5.631850in}{0.660000in}}%
\pgfpathlineto{\pgfqpoint{5.631850in}{3.241272in}}%
\pgfpathlineto{\pgfqpoint{5.625650in}{3.241272in}}%
\pgfpathlineto{\pgfqpoint{5.625650in}{0.660000in}}%
\pgfpathclose%
\pgfusepath{fill}%
\end{pgfscope}%
\begin{pgfscope}%
\pgfpathrectangle{\pgfqpoint{1.250000in}{0.660000in}}{\pgfqpoint{7.750000in}{4.620000in}}%
\pgfusepath{clip}%
\pgfsetbuttcap%
\pgfsetmiterjoin%
\definecolor{currentfill}{rgb}{0.000000,0.000000,0.000000}%
\pgfsetfillcolor{currentfill}%
\pgfsetlinewidth{0.000000pt}%
\definecolor{currentstroke}{rgb}{0.000000,0.000000,0.000000}%
\pgfsetstrokecolor{currentstroke}%
\pgfsetstrokeopacity{0.000000}%
\pgfsetdash{}{0pt}%
\pgfpathmoveto{\pgfqpoint{5.633400in}{0.660000in}}%
\pgfpathlineto{\pgfqpoint{5.639600in}{0.660000in}}%
\pgfpathlineto{\pgfqpoint{5.639600in}{3.241102in}}%
\pgfpathlineto{\pgfqpoint{5.633400in}{3.241102in}}%
\pgfpathlineto{\pgfqpoint{5.633400in}{0.660000in}}%
\pgfpathclose%
\pgfusepath{fill}%
\end{pgfscope}%
\begin{pgfscope}%
\pgfpathrectangle{\pgfqpoint{1.250000in}{0.660000in}}{\pgfqpoint{7.750000in}{4.620000in}}%
\pgfusepath{clip}%
\pgfsetbuttcap%
\pgfsetmiterjoin%
\definecolor{currentfill}{rgb}{0.000000,0.000000,0.000000}%
\pgfsetfillcolor{currentfill}%
\pgfsetlinewidth{0.000000pt}%
\definecolor{currentstroke}{rgb}{0.000000,0.000000,0.000000}%
\pgfsetstrokecolor{currentstroke}%
\pgfsetstrokeopacity{0.000000}%
\pgfsetdash{}{0pt}%
\pgfpathmoveto{\pgfqpoint{5.641150in}{0.660000in}}%
\pgfpathlineto{\pgfqpoint{5.647350in}{0.660000in}}%
\pgfpathlineto{\pgfqpoint{5.647350in}{3.240932in}}%
\pgfpathlineto{\pgfqpoint{5.641150in}{3.240932in}}%
\pgfpathlineto{\pgfqpoint{5.641150in}{0.660000in}}%
\pgfpathclose%
\pgfusepath{fill}%
\end{pgfscope}%
\begin{pgfscope}%
\pgfpathrectangle{\pgfqpoint{1.250000in}{0.660000in}}{\pgfqpoint{7.750000in}{4.620000in}}%
\pgfusepath{clip}%
\pgfsetbuttcap%
\pgfsetmiterjoin%
\definecolor{currentfill}{rgb}{0.000000,0.000000,0.000000}%
\pgfsetfillcolor{currentfill}%
\pgfsetlinewidth{0.000000pt}%
\definecolor{currentstroke}{rgb}{0.000000,0.000000,0.000000}%
\pgfsetstrokecolor{currentstroke}%
\pgfsetstrokeopacity{0.000000}%
\pgfsetdash{}{0pt}%
\pgfpathmoveto{\pgfqpoint{5.648900in}{0.660000in}}%
\pgfpathlineto{\pgfqpoint{5.655100in}{0.660000in}}%
\pgfpathlineto{\pgfqpoint{5.655100in}{3.240762in}}%
\pgfpathlineto{\pgfqpoint{5.648900in}{3.240762in}}%
\pgfpathlineto{\pgfqpoint{5.648900in}{0.660000in}}%
\pgfpathclose%
\pgfusepath{fill}%
\end{pgfscope}%
\begin{pgfscope}%
\pgfpathrectangle{\pgfqpoint{1.250000in}{0.660000in}}{\pgfqpoint{7.750000in}{4.620000in}}%
\pgfusepath{clip}%
\pgfsetbuttcap%
\pgfsetmiterjoin%
\definecolor{currentfill}{rgb}{0.000000,0.000000,0.000000}%
\pgfsetfillcolor{currentfill}%
\pgfsetlinewidth{0.000000pt}%
\definecolor{currentstroke}{rgb}{0.000000,0.000000,0.000000}%
\pgfsetstrokecolor{currentstroke}%
\pgfsetstrokeopacity{0.000000}%
\pgfsetdash{}{0pt}%
\pgfpathmoveto{\pgfqpoint{5.656650in}{0.660000in}}%
\pgfpathlineto{\pgfqpoint{5.662850in}{0.660000in}}%
\pgfpathlineto{\pgfqpoint{5.662850in}{3.240592in}}%
\pgfpathlineto{\pgfqpoint{5.656650in}{3.240592in}}%
\pgfpathlineto{\pgfqpoint{5.656650in}{0.660000in}}%
\pgfpathclose%
\pgfusepath{fill}%
\end{pgfscope}%
\begin{pgfscope}%
\pgfpathrectangle{\pgfqpoint{1.250000in}{0.660000in}}{\pgfqpoint{7.750000in}{4.620000in}}%
\pgfusepath{clip}%
\pgfsetbuttcap%
\pgfsetmiterjoin%
\definecolor{currentfill}{rgb}{0.000000,0.000000,0.000000}%
\pgfsetfillcolor{currentfill}%
\pgfsetlinewidth{0.000000pt}%
\definecolor{currentstroke}{rgb}{0.000000,0.000000,0.000000}%
\pgfsetstrokecolor{currentstroke}%
\pgfsetstrokeopacity{0.000000}%
\pgfsetdash{}{0pt}%
\pgfpathmoveto{\pgfqpoint{5.664400in}{0.660000in}}%
\pgfpathlineto{\pgfqpoint{5.670600in}{0.660000in}}%
\pgfpathlineto{\pgfqpoint{5.670600in}{3.240422in}}%
\pgfpathlineto{\pgfqpoint{5.664400in}{3.240422in}}%
\pgfpathlineto{\pgfqpoint{5.664400in}{0.660000in}}%
\pgfpathclose%
\pgfusepath{fill}%
\end{pgfscope}%
\begin{pgfscope}%
\pgfpathrectangle{\pgfqpoint{1.250000in}{0.660000in}}{\pgfqpoint{7.750000in}{4.620000in}}%
\pgfusepath{clip}%
\pgfsetbuttcap%
\pgfsetmiterjoin%
\definecolor{currentfill}{rgb}{0.000000,0.000000,0.000000}%
\pgfsetfillcolor{currentfill}%
\pgfsetlinewidth{0.000000pt}%
\definecolor{currentstroke}{rgb}{0.000000,0.000000,0.000000}%
\pgfsetstrokecolor{currentstroke}%
\pgfsetstrokeopacity{0.000000}%
\pgfsetdash{}{0pt}%
\pgfpathmoveto{\pgfqpoint{5.672150in}{0.660000in}}%
\pgfpathlineto{\pgfqpoint{5.678350in}{0.660000in}}%
\pgfpathlineto{\pgfqpoint{5.678350in}{3.240253in}}%
\pgfpathlineto{\pgfqpoint{5.672150in}{3.240253in}}%
\pgfpathlineto{\pgfqpoint{5.672150in}{0.660000in}}%
\pgfpathclose%
\pgfusepath{fill}%
\end{pgfscope}%
\begin{pgfscope}%
\pgfpathrectangle{\pgfqpoint{1.250000in}{0.660000in}}{\pgfqpoint{7.750000in}{4.620000in}}%
\pgfusepath{clip}%
\pgfsetbuttcap%
\pgfsetmiterjoin%
\definecolor{currentfill}{rgb}{0.000000,0.000000,0.000000}%
\pgfsetfillcolor{currentfill}%
\pgfsetlinewidth{0.000000pt}%
\definecolor{currentstroke}{rgb}{0.000000,0.000000,0.000000}%
\pgfsetstrokecolor{currentstroke}%
\pgfsetstrokeopacity{0.000000}%
\pgfsetdash{}{0pt}%
\pgfpathmoveto{\pgfqpoint{5.679900in}{0.660000in}}%
\pgfpathlineto{\pgfqpoint{5.686100in}{0.660000in}}%
\pgfpathlineto{\pgfqpoint{5.686100in}{3.240083in}}%
\pgfpathlineto{\pgfqpoint{5.679900in}{3.240083in}}%
\pgfpathlineto{\pgfqpoint{5.679900in}{0.660000in}}%
\pgfpathclose%
\pgfusepath{fill}%
\end{pgfscope}%
\begin{pgfscope}%
\pgfpathrectangle{\pgfqpoint{1.250000in}{0.660000in}}{\pgfqpoint{7.750000in}{4.620000in}}%
\pgfusepath{clip}%
\pgfsetbuttcap%
\pgfsetmiterjoin%
\definecolor{currentfill}{rgb}{0.000000,0.000000,0.000000}%
\pgfsetfillcolor{currentfill}%
\pgfsetlinewidth{0.000000pt}%
\definecolor{currentstroke}{rgb}{0.000000,0.000000,0.000000}%
\pgfsetstrokecolor{currentstroke}%
\pgfsetstrokeopacity{0.000000}%
\pgfsetdash{}{0pt}%
\pgfpathmoveto{\pgfqpoint{5.687650in}{0.660000in}}%
\pgfpathlineto{\pgfqpoint{5.693850in}{0.660000in}}%
\pgfpathlineto{\pgfqpoint{5.693850in}{3.239913in}}%
\pgfpathlineto{\pgfqpoint{5.687650in}{3.239913in}}%
\pgfpathlineto{\pgfqpoint{5.687650in}{0.660000in}}%
\pgfpathclose%
\pgfusepath{fill}%
\end{pgfscope}%
\begin{pgfscope}%
\pgfpathrectangle{\pgfqpoint{1.250000in}{0.660000in}}{\pgfqpoint{7.750000in}{4.620000in}}%
\pgfusepath{clip}%
\pgfsetbuttcap%
\pgfsetmiterjoin%
\definecolor{currentfill}{rgb}{0.000000,0.000000,0.000000}%
\pgfsetfillcolor{currentfill}%
\pgfsetlinewidth{0.000000pt}%
\definecolor{currentstroke}{rgb}{0.000000,0.000000,0.000000}%
\pgfsetstrokecolor{currentstroke}%
\pgfsetstrokeopacity{0.000000}%
\pgfsetdash{}{0pt}%
\pgfpathmoveto{\pgfqpoint{5.695400in}{0.660000in}}%
\pgfpathlineto{\pgfqpoint{5.701600in}{0.660000in}}%
\pgfpathlineto{\pgfqpoint{5.701600in}{3.239744in}}%
\pgfpathlineto{\pgfqpoint{5.695400in}{3.239744in}}%
\pgfpathlineto{\pgfqpoint{5.695400in}{0.660000in}}%
\pgfpathclose%
\pgfusepath{fill}%
\end{pgfscope}%
\begin{pgfscope}%
\pgfpathrectangle{\pgfqpoint{1.250000in}{0.660000in}}{\pgfqpoint{7.750000in}{4.620000in}}%
\pgfusepath{clip}%
\pgfsetbuttcap%
\pgfsetmiterjoin%
\definecolor{currentfill}{rgb}{0.000000,0.000000,0.000000}%
\pgfsetfillcolor{currentfill}%
\pgfsetlinewidth{0.000000pt}%
\definecolor{currentstroke}{rgb}{0.000000,0.000000,0.000000}%
\pgfsetstrokecolor{currentstroke}%
\pgfsetstrokeopacity{0.000000}%
\pgfsetdash{}{0pt}%
\pgfpathmoveto{\pgfqpoint{5.703150in}{0.660000in}}%
\pgfpathlineto{\pgfqpoint{5.709350in}{0.660000in}}%
\pgfpathlineto{\pgfqpoint{5.709350in}{3.239578in}}%
\pgfpathlineto{\pgfqpoint{5.703150in}{3.239578in}}%
\pgfpathlineto{\pgfqpoint{5.703150in}{0.660000in}}%
\pgfpathclose%
\pgfusepath{fill}%
\end{pgfscope}%
\begin{pgfscope}%
\pgfpathrectangle{\pgfqpoint{1.250000in}{0.660000in}}{\pgfqpoint{7.750000in}{4.620000in}}%
\pgfusepath{clip}%
\pgfsetbuttcap%
\pgfsetmiterjoin%
\definecolor{currentfill}{rgb}{0.000000,0.000000,0.000000}%
\pgfsetfillcolor{currentfill}%
\pgfsetlinewidth{0.000000pt}%
\definecolor{currentstroke}{rgb}{0.000000,0.000000,0.000000}%
\pgfsetstrokecolor{currentstroke}%
\pgfsetstrokeopacity{0.000000}%
\pgfsetdash{}{0pt}%
\pgfpathmoveto{\pgfqpoint{5.710900in}{0.660000in}}%
\pgfpathlineto{\pgfqpoint{5.717100in}{0.660000in}}%
\pgfpathlineto{\pgfqpoint{5.717100in}{3.239413in}}%
\pgfpathlineto{\pgfqpoint{5.710900in}{3.239413in}}%
\pgfpathlineto{\pgfqpoint{5.710900in}{0.660000in}}%
\pgfpathclose%
\pgfusepath{fill}%
\end{pgfscope}%
\begin{pgfscope}%
\pgfpathrectangle{\pgfqpoint{1.250000in}{0.660000in}}{\pgfqpoint{7.750000in}{4.620000in}}%
\pgfusepath{clip}%
\pgfsetbuttcap%
\pgfsetmiterjoin%
\definecolor{currentfill}{rgb}{0.000000,0.000000,0.000000}%
\pgfsetfillcolor{currentfill}%
\pgfsetlinewidth{0.000000pt}%
\definecolor{currentstroke}{rgb}{0.000000,0.000000,0.000000}%
\pgfsetstrokecolor{currentstroke}%
\pgfsetstrokeopacity{0.000000}%
\pgfsetdash{}{0pt}%
\pgfpathmoveto{\pgfqpoint{5.718650in}{0.660000in}}%
\pgfpathlineto{\pgfqpoint{5.724850in}{0.660000in}}%
\pgfpathlineto{\pgfqpoint{5.724850in}{3.239247in}}%
\pgfpathlineto{\pgfqpoint{5.718650in}{3.239247in}}%
\pgfpathlineto{\pgfqpoint{5.718650in}{0.660000in}}%
\pgfpathclose%
\pgfusepath{fill}%
\end{pgfscope}%
\begin{pgfscope}%
\pgfpathrectangle{\pgfqpoint{1.250000in}{0.660000in}}{\pgfqpoint{7.750000in}{4.620000in}}%
\pgfusepath{clip}%
\pgfsetbuttcap%
\pgfsetmiterjoin%
\definecolor{currentfill}{rgb}{0.000000,0.000000,0.000000}%
\pgfsetfillcolor{currentfill}%
\pgfsetlinewidth{0.000000pt}%
\definecolor{currentstroke}{rgb}{0.000000,0.000000,0.000000}%
\pgfsetstrokecolor{currentstroke}%
\pgfsetstrokeopacity{0.000000}%
\pgfsetdash{}{0pt}%
\pgfpathmoveto{\pgfqpoint{5.726400in}{0.660000in}}%
\pgfpathlineto{\pgfqpoint{5.732600in}{0.660000in}}%
\pgfpathlineto{\pgfqpoint{5.732600in}{3.239082in}}%
\pgfpathlineto{\pgfqpoint{5.726400in}{3.239082in}}%
\pgfpathlineto{\pgfqpoint{5.726400in}{0.660000in}}%
\pgfpathclose%
\pgfusepath{fill}%
\end{pgfscope}%
\begin{pgfscope}%
\pgfpathrectangle{\pgfqpoint{1.250000in}{0.660000in}}{\pgfqpoint{7.750000in}{4.620000in}}%
\pgfusepath{clip}%
\pgfsetbuttcap%
\pgfsetmiterjoin%
\definecolor{currentfill}{rgb}{0.000000,0.000000,0.000000}%
\pgfsetfillcolor{currentfill}%
\pgfsetlinewidth{0.000000pt}%
\definecolor{currentstroke}{rgb}{0.000000,0.000000,0.000000}%
\pgfsetstrokecolor{currentstroke}%
\pgfsetstrokeopacity{0.000000}%
\pgfsetdash{}{0pt}%
\pgfpathmoveto{\pgfqpoint{5.734150in}{0.660000in}}%
\pgfpathlineto{\pgfqpoint{5.740350in}{0.660000in}}%
\pgfpathlineto{\pgfqpoint{5.740350in}{3.238916in}}%
\pgfpathlineto{\pgfqpoint{5.734150in}{3.238916in}}%
\pgfpathlineto{\pgfqpoint{5.734150in}{0.660000in}}%
\pgfpathclose%
\pgfusepath{fill}%
\end{pgfscope}%
\begin{pgfscope}%
\pgfpathrectangle{\pgfqpoint{1.250000in}{0.660000in}}{\pgfqpoint{7.750000in}{4.620000in}}%
\pgfusepath{clip}%
\pgfsetbuttcap%
\pgfsetmiterjoin%
\definecolor{currentfill}{rgb}{0.000000,0.000000,0.000000}%
\pgfsetfillcolor{currentfill}%
\pgfsetlinewidth{0.000000pt}%
\definecolor{currentstroke}{rgb}{0.000000,0.000000,0.000000}%
\pgfsetstrokecolor{currentstroke}%
\pgfsetstrokeopacity{0.000000}%
\pgfsetdash{}{0pt}%
\pgfpathmoveto{\pgfqpoint{5.741900in}{0.660000in}}%
\pgfpathlineto{\pgfqpoint{5.748100in}{0.660000in}}%
\pgfpathlineto{\pgfqpoint{5.748100in}{3.238751in}}%
\pgfpathlineto{\pgfqpoint{5.741900in}{3.238751in}}%
\pgfpathlineto{\pgfqpoint{5.741900in}{0.660000in}}%
\pgfpathclose%
\pgfusepath{fill}%
\end{pgfscope}%
\begin{pgfscope}%
\pgfpathrectangle{\pgfqpoint{1.250000in}{0.660000in}}{\pgfqpoint{7.750000in}{4.620000in}}%
\pgfusepath{clip}%
\pgfsetbuttcap%
\pgfsetmiterjoin%
\definecolor{currentfill}{rgb}{0.000000,0.000000,0.000000}%
\pgfsetfillcolor{currentfill}%
\pgfsetlinewidth{0.000000pt}%
\definecolor{currentstroke}{rgb}{0.000000,0.000000,0.000000}%
\pgfsetstrokecolor{currentstroke}%
\pgfsetstrokeopacity{0.000000}%
\pgfsetdash{}{0pt}%
\pgfpathmoveto{\pgfqpoint{5.749650in}{0.660000in}}%
\pgfpathlineto{\pgfqpoint{5.755850in}{0.660000in}}%
\pgfpathlineto{\pgfqpoint{5.755850in}{3.238585in}}%
\pgfpathlineto{\pgfqpoint{5.749650in}{3.238585in}}%
\pgfpathlineto{\pgfqpoint{5.749650in}{0.660000in}}%
\pgfpathclose%
\pgfusepath{fill}%
\end{pgfscope}%
\begin{pgfscope}%
\pgfpathrectangle{\pgfqpoint{1.250000in}{0.660000in}}{\pgfqpoint{7.750000in}{4.620000in}}%
\pgfusepath{clip}%
\pgfsetbuttcap%
\pgfsetmiterjoin%
\definecolor{currentfill}{rgb}{0.000000,0.000000,0.000000}%
\pgfsetfillcolor{currentfill}%
\pgfsetlinewidth{0.000000pt}%
\definecolor{currentstroke}{rgb}{0.000000,0.000000,0.000000}%
\pgfsetstrokecolor{currentstroke}%
\pgfsetstrokeopacity{0.000000}%
\pgfsetdash{}{0pt}%
\pgfpathmoveto{\pgfqpoint{5.757400in}{0.660000in}}%
\pgfpathlineto{\pgfqpoint{5.763600in}{0.660000in}}%
\pgfpathlineto{\pgfqpoint{5.763600in}{3.238420in}}%
\pgfpathlineto{\pgfqpoint{5.757400in}{3.238420in}}%
\pgfpathlineto{\pgfqpoint{5.757400in}{0.660000in}}%
\pgfpathclose%
\pgfusepath{fill}%
\end{pgfscope}%
\begin{pgfscope}%
\pgfpathrectangle{\pgfqpoint{1.250000in}{0.660000in}}{\pgfqpoint{7.750000in}{4.620000in}}%
\pgfusepath{clip}%
\pgfsetbuttcap%
\pgfsetmiterjoin%
\definecolor{currentfill}{rgb}{0.000000,0.000000,0.000000}%
\pgfsetfillcolor{currentfill}%
\pgfsetlinewidth{0.000000pt}%
\definecolor{currentstroke}{rgb}{0.000000,0.000000,0.000000}%
\pgfsetstrokecolor{currentstroke}%
\pgfsetstrokeopacity{0.000000}%
\pgfsetdash{}{0pt}%
\pgfpathmoveto{\pgfqpoint{5.765150in}{0.660000in}}%
\pgfpathlineto{\pgfqpoint{5.771350in}{0.660000in}}%
\pgfpathlineto{\pgfqpoint{5.771350in}{3.238254in}}%
\pgfpathlineto{\pgfqpoint{5.765150in}{3.238254in}}%
\pgfpathlineto{\pgfqpoint{5.765150in}{0.660000in}}%
\pgfpathclose%
\pgfusepath{fill}%
\end{pgfscope}%
\begin{pgfscope}%
\pgfpathrectangle{\pgfqpoint{1.250000in}{0.660000in}}{\pgfqpoint{7.750000in}{4.620000in}}%
\pgfusepath{clip}%
\pgfsetbuttcap%
\pgfsetmiterjoin%
\definecolor{currentfill}{rgb}{0.000000,0.000000,0.000000}%
\pgfsetfillcolor{currentfill}%
\pgfsetlinewidth{0.000000pt}%
\definecolor{currentstroke}{rgb}{0.000000,0.000000,0.000000}%
\pgfsetstrokecolor{currentstroke}%
\pgfsetstrokeopacity{0.000000}%
\pgfsetdash{}{0pt}%
\pgfpathmoveto{\pgfqpoint{5.772900in}{0.660000in}}%
\pgfpathlineto{\pgfqpoint{5.779100in}{0.660000in}}%
\pgfpathlineto{\pgfqpoint{5.779100in}{3.238089in}}%
\pgfpathlineto{\pgfqpoint{5.772900in}{3.238089in}}%
\pgfpathlineto{\pgfqpoint{5.772900in}{0.660000in}}%
\pgfpathclose%
\pgfusepath{fill}%
\end{pgfscope}%
\begin{pgfscope}%
\pgfpathrectangle{\pgfqpoint{1.250000in}{0.660000in}}{\pgfqpoint{7.750000in}{4.620000in}}%
\pgfusepath{clip}%
\pgfsetbuttcap%
\pgfsetmiterjoin%
\definecolor{currentfill}{rgb}{0.000000,0.000000,0.000000}%
\pgfsetfillcolor{currentfill}%
\pgfsetlinewidth{0.000000pt}%
\definecolor{currentstroke}{rgb}{0.000000,0.000000,0.000000}%
\pgfsetstrokecolor{currentstroke}%
\pgfsetstrokeopacity{0.000000}%
\pgfsetdash{}{0pt}%
\pgfpathmoveto{\pgfqpoint{5.780650in}{0.660000in}}%
\pgfpathlineto{\pgfqpoint{5.786850in}{0.660000in}}%
\pgfpathlineto{\pgfqpoint{5.786850in}{3.237924in}}%
\pgfpathlineto{\pgfqpoint{5.780650in}{3.237924in}}%
\pgfpathlineto{\pgfqpoint{5.780650in}{0.660000in}}%
\pgfpathclose%
\pgfusepath{fill}%
\end{pgfscope}%
\begin{pgfscope}%
\pgfpathrectangle{\pgfqpoint{1.250000in}{0.660000in}}{\pgfqpoint{7.750000in}{4.620000in}}%
\pgfusepath{clip}%
\pgfsetbuttcap%
\pgfsetmiterjoin%
\definecolor{currentfill}{rgb}{0.000000,0.000000,0.000000}%
\pgfsetfillcolor{currentfill}%
\pgfsetlinewidth{0.000000pt}%
\definecolor{currentstroke}{rgb}{0.000000,0.000000,0.000000}%
\pgfsetstrokecolor{currentstroke}%
\pgfsetstrokeopacity{0.000000}%
\pgfsetdash{}{0pt}%
\pgfpathmoveto{\pgfqpoint{5.788400in}{0.660000in}}%
\pgfpathlineto{\pgfqpoint{5.794600in}{0.660000in}}%
\pgfpathlineto{\pgfqpoint{5.794600in}{3.237762in}}%
\pgfpathlineto{\pgfqpoint{5.788400in}{3.237762in}}%
\pgfpathlineto{\pgfqpoint{5.788400in}{0.660000in}}%
\pgfpathclose%
\pgfusepath{fill}%
\end{pgfscope}%
\begin{pgfscope}%
\pgfpathrectangle{\pgfqpoint{1.250000in}{0.660000in}}{\pgfqpoint{7.750000in}{4.620000in}}%
\pgfusepath{clip}%
\pgfsetbuttcap%
\pgfsetmiterjoin%
\definecolor{currentfill}{rgb}{0.000000,0.000000,0.000000}%
\pgfsetfillcolor{currentfill}%
\pgfsetlinewidth{0.000000pt}%
\definecolor{currentstroke}{rgb}{0.000000,0.000000,0.000000}%
\pgfsetstrokecolor{currentstroke}%
\pgfsetstrokeopacity{0.000000}%
\pgfsetdash{}{0pt}%
\pgfpathmoveto{\pgfqpoint{5.796150in}{0.660000in}}%
\pgfpathlineto{\pgfqpoint{5.802350in}{0.660000in}}%
\pgfpathlineto{\pgfqpoint{5.802350in}{3.237601in}}%
\pgfpathlineto{\pgfqpoint{5.796150in}{3.237601in}}%
\pgfpathlineto{\pgfqpoint{5.796150in}{0.660000in}}%
\pgfpathclose%
\pgfusepath{fill}%
\end{pgfscope}%
\begin{pgfscope}%
\pgfpathrectangle{\pgfqpoint{1.250000in}{0.660000in}}{\pgfqpoint{7.750000in}{4.620000in}}%
\pgfusepath{clip}%
\pgfsetbuttcap%
\pgfsetmiterjoin%
\definecolor{currentfill}{rgb}{0.000000,0.000000,0.000000}%
\pgfsetfillcolor{currentfill}%
\pgfsetlinewidth{0.000000pt}%
\definecolor{currentstroke}{rgb}{0.000000,0.000000,0.000000}%
\pgfsetstrokecolor{currentstroke}%
\pgfsetstrokeopacity{0.000000}%
\pgfsetdash{}{0pt}%
\pgfpathmoveto{\pgfqpoint{5.803900in}{0.660000in}}%
\pgfpathlineto{\pgfqpoint{5.810100in}{0.660000in}}%
\pgfpathlineto{\pgfqpoint{5.810100in}{3.237440in}}%
\pgfpathlineto{\pgfqpoint{5.803900in}{3.237440in}}%
\pgfpathlineto{\pgfqpoint{5.803900in}{0.660000in}}%
\pgfpathclose%
\pgfusepath{fill}%
\end{pgfscope}%
\begin{pgfscope}%
\pgfpathrectangle{\pgfqpoint{1.250000in}{0.660000in}}{\pgfqpoint{7.750000in}{4.620000in}}%
\pgfusepath{clip}%
\pgfsetbuttcap%
\pgfsetmiterjoin%
\definecolor{currentfill}{rgb}{0.000000,0.000000,0.000000}%
\pgfsetfillcolor{currentfill}%
\pgfsetlinewidth{0.000000pt}%
\definecolor{currentstroke}{rgb}{0.000000,0.000000,0.000000}%
\pgfsetstrokecolor{currentstroke}%
\pgfsetstrokeopacity{0.000000}%
\pgfsetdash{}{0pt}%
\pgfpathmoveto{\pgfqpoint{5.811650in}{0.660000in}}%
\pgfpathlineto{\pgfqpoint{5.817850in}{0.660000in}}%
\pgfpathlineto{\pgfqpoint{5.817850in}{3.237279in}}%
\pgfpathlineto{\pgfqpoint{5.811650in}{3.237279in}}%
\pgfpathlineto{\pgfqpoint{5.811650in}{0.660000in}}%
\pgfpathclose%
\pgfusepath{fill}%
\end{pgfscope}%
\begin{pgfscope}%
\pgfpathrectangle{\pgfqpoint{1.250000in}{0.660000in}}{\pgfqpoint{7.750000in}{4.620000in}}%
\pgfusepath{clip}%
\pgfsetbuttcap%
\pgfsetmiterjoin%
\definecolor{currentfill}{rgb}{0.000000,0.000000,0.000000}%
\pgfsetfillcolor{currentfill}%
\pgfsetlinewidth{0.000000pt}%
\definecolor{currentstroke}{rgb}{0.000000,0.000000,0.000000}%
\pgfsetstrokecolor{currentstroke}%
\pgfsetstrokeopacity{0.000000}%
\pgfsetdash{}{0pt}%
\pgfpathmoveto{\pgfqpoint{5.819400in}{0.660000in}}%
\pgfpathlineto{\pgfqpoint{5.825600in}{0.660000in}}%
\pgfpathlineto{\pgfqpoint{5.825600in}{3.237118in}}%
\pgfpathlineto{\pgfqpoint{5.819400in}{3.237118in}}%
\pgfpathlineto{\pgfqpoint{5.819400in}{0.660000in}}%
\pgfpathclose%
\pgfusepath{fill}%
\end{pgfscope}%
\begin{pgfscope}%
\pgfpathrectangle{\pgfqpoint{1.250000in}{0.660000in}}{\pgfqpoint{7.750000in}{4.620000in}}%
\pgfusepath{clip}%
\pgfsetbuttcap%
\pgfsetmiterjoin%
\definecolor{currentfill}{rgb}{0.000000,0.000000,0.000000}%
\pgfsetfillcolor{currentfill}%
\pgfsetlinewidth{0.000000pt}%
\definecolor{currentstroke}{rgb}{0.000000,0.000000,0.000000}%
\pgfsetstrokecolor{currentstroke}%
\pgfsetstrokeopacity{0.000000}%
\pgfsetdash{}{0pt}%
\pgfpathmoveto{\pgfqpoint{5.827150in}{0.660000in}}%
\pgfpathlineto{\pgfqpoint{5.833350in}{0.660000in}}%
\pgfpathlineto{\pgfqpoint{5.833350in}{3.236957in}}%
\pgfpathlineto{\pgfqpoint{5.827150in}{3.236957in}}%
\pgfpathlineto{\pgfqpoint{5.827150in}{0.660000in}}%
\pgfpathclose%
\pgfusepath{fill}%
\end{pgfscope}%
\begin{pgfscope}%
\pgfpathrectangle{\pgfqpoint{1.250000in}{0.660000in}}{\pgfqpoint{7.750000in}{4.620000in}}%
\pgfusepath{clip}%
\pgfsetbuttcap%
\pgfsetmiterjoin%
\definecolor{currentfill}{rgb}{0.000000,0.000000,0.000000}%
\pgfsetfillcolor{currentfill}%
\pgfsetlinewidth{0.000000pt}%
\definecolor{currentstroke}{rgb}{0.000000,0.000000,0.000000}%
\pgfsetstrokecolor{currentstroke}%
\pgfsetstrokeopacity{0.000000}%
\pgfsetdash{}{0pt}%
\pgfpathmoveto{\pgfqpoint{5.834900in}{0.660000in}}%
\pgfpathlineto{\pgfqpoint{5.841100in}{0.660000in}}%
\pgfpathlineto{\pgfqpoint{5.841100in}{3.236796in}}%
\pgfpathlineto{\pgfqpoint{5.834900in}{3.236796in}}%
\pgfpathlineto{\pgfqpoint{5.834900in}{0.660000in}}%
\pgfpathclose%
\pgfusepath{fill}%
\end{pgfscope}%
\begin{pgfscope}%
\pgfpathrectangle{\pgfqpoint{1.250000in}{0.660000in}}{\pgfqpoint{7.750000in}{4.620000in}}%
\pgfusepath{clip}%
\pgfsetbuttcap%
\pgfsetmiterjoin%
\definecolor{currentfill}{rgb}{0.000000,0.000000,0.000000}%
\pgfsetfillcolor{currentfill}%
\pgfsetlinewidth{0.000000pt}%
\definecolor{currentstroke}{rgb}{0.000000,0.000000,0.000000}%
\pgfsetstrokecolor{currentstroke}%
\pgfsetstrokeopacity{0.000000}%
\pgfsetdash{}{0pt}%
\pgfpathmoveto{\pgfqpoint{5.842650in}{0.660000in}}%
\pgfpathlineto{\pgfqpoint{5.848850in}{0.660000in}}%
\pgfpathlineto{\pgfqpoint{5.848850in}{3.236635in}}%
\pgfpathlineto{\pgfqpoint{5.842650in}{3.236635in}}%
\pgfpathlineto{\pgfqpoint{5.842650in}{0.660000in}}%
\pgfpathclose%
\pgfusepath{fill}%
\end{pgfscope}%
\begin{pgfscope}%
\pgfpathrectangle{\pgfqpoint{1.250000in}{0.660000in}}{\pgfqpoint{7.750000in}{4.620000in}}%
\pgfusepath{clip}%
\pgfsetbuttcap%
\pgfsetmiterjoin%
\definecolor{currentfill}{rgb}{0.000000,0.000000,0.000000}%
\pgfsetfillcolor{currentfill}%
\pgfsetlinewidth{0.000000pt}%
\definecolor{currentstroke}{rgb}{0.000000,0.000000,0.000000}%
\pgfsetstrokecolor{currentstroke}%
\pgfsetstrokeopacity{0.000000}%
\pgfsetdash{}{0pt}%
\pgfpathmoveto{\pgfqpoint{5.850400in}{0.660000in}}%
\pgfpathlineto{\pgfqpoint{5.856600in}{0.660000in}}%
\pgfpathlineto{\pgfqpoint{5.856600in}{3.236473in}}%
\pgfpathlineto{\pgfqpoint{5.850400in}{3.236473in}}%
\pgfpathlineto{\pgfqpoint{5.850400in}{0.660000in}}%
\pgfpathclose%
\pgfusepath{fill}%
\end{pgfscope}%
\begin{pgfscope}%
\pgfpathrectangle{\pgfqpoint{1.250000in}{0.660000in}}{\pgfqpoint{7.750000in}{4.620000in}}%
\pgfusepath{clip}%
\pgfsetbuttcap%
\pgfsetmiterjoin%
\definecolor{currentfill}{rgb}{0.000000,0.000000,0.000000}%
\pgfsetfillcolor{currentfill}%
\pgfsetlinewidth{0.000000pt}%
\definecolor{currentstroke}{rgb}{0.000000,0.000000,0.000000}%
\pgfsetstrokecolor{currentstroke}%
\pgfsetstrokeopacity{0.000000}%
\pgfsetdash{}{0pt}%
\pgfpathmoveto{\pgfqpoint{5.858150in}{0.660000in}}%
\pgfpathlineto{\pgfqpoint{5.864350in}{0.660000in}}%
\pgfpathlineto{\pgfqpoint{5.864350in}{3.236312in}}%
\pgfpathlineto{\pgfqpoint{5.858150in}{3.236312in}}%
\pgfpathlineto{\pgfqpoint{5.858150in}{0.660000in}}%
\pgfpathclose%
\pgfusepath{fill}%
\end{pgfscope}%
\begin{pgfscope}%
\pgfpathrectangle{\pgfqpoint{1.250000in}{0.660000in}}{\pgfqpoint{7.750000in}{4.620000in}}%
\pgfusepath{clip}%
\pgfsetbuttcap%
\pgfsetmiterjoin%
\definecolor{currentfill}{rgb}{0.000000,0.000000,0.000000}%
\pgfsetfillcolor{currentfill}%
\pgfsetlinewidth{0.000000pt}%
\definecolor{currentstroke}{rgb}{0.000000,0.000000,0.000000}%
\pgfsetstrokecolor{currentstroke}%
\pgfsetstrokeopacity{0.000000}%
\pgfsetdash{}{0pt}%
\pgfpathmoveto{\pgfqpoint{5.865900in}{0.660000in}}%
\pgfpathlineto{\pgfqpoint{5.872100in}{0.660000in}}%
\pgfpathlineto{\pgfqpoint{5.872100in}{3.236151in}}%
\pgfpathlineto{\pgfqpoint{5.865900in}{3.236151in}}%
\pgfpathlineto{\pgfqpoint{5.865900in}{0.660000in}}%
\pgfpathclose%
\pgfusepath{fill}%
\end{pgfscope}%
\begin{pgfscope}%
\pgfpathrectangle{\pgfqpoint{1.250000in}{0.660000in}}{\pgfqpoint{7.750000in}{4.620000in}}%
\pgfusepath{clip}%
\pgfsetbuttcap%
\pgfsetmiterjoin%
\definecolor{currentfill}{rgb}{0.000000,0.000000,0.000000}%
\pgfsetfillcolor{currentfill}%
\pgfsetlinewidth{0.000000pt}%
\definecolor{currentstroke}{rgb}{0.000000,0.000000,0.000000}%
\pgfsetstrokecolor{currentstroke}%
\pgfsetstrokeopacity{0.000000}%
\pgfsetdash{}{0pt}%
\pgfpathmoveto{\pgfqpoint{5.873650in}{0.660000in}}%
\pgfpathlineto{\pgfqpoint{5.879850in}{0.660000in}}%
\pgfpathlineto{\pgfqpoint{5.879850in}{3.235992in}}%
\pgfpathlineto{\pgfqpoint{5.873650in}{3.235992in}}%
\pgfpathlineto{\pgfqpoint{5.873650in}{0.660000in}}%
\pgfpathclose%
\pgfusepath{fill}%
\end{pgfscope}%
\begin{pgfscope}%
\pgfpathrectangle{\pgfqpoint{1.250000in}{0.660000in}}{\pgfqpoint{7.750000in}{4.620000in}}%
\pgfusepath{clip}%
\pgfsetbuttcap%
\pgfsetmiterjoin%
\definecolor{currentfill}{rgb}{0.000000,0.000000,0.000000}%
\pgfsetfillcolor{currentfill}%
\pgfsetlinewidth{0.000000pt}%
\definecolor{currentstroke}{rgb}{0.000000,0.000000,0.000000}%
\pgfsetstrokecolor{currentstroke}%
\pgfsetstrokeopacity{0.000000}%
\pgfsetdash{}{0pt}%
\pgfpathmoveto{\pgfqpoint{5.881400in}{0.660000in}}%
\pgfpathlineto{\pgfqpoint{5.887600in}{0.660000in}}%
\pgfpathlineto{\pgfqpoint{5.887600in}{3.235835in}}%
\pgfpathlineto{\pgfqpoint{5.881400in}{3.235835in}}%
\pgfpathlineto{\pgfqpoint{5.881400in}{0.660000in}}%
\pgfpathclose%
\pgfusepath{fill}%
\end{pgfscope}%
\begin{pgfscope}%
\pgfpathrectangle{\pgfqpoint{1.250000in}{0.660000in}}{\pgfqpoint{7.750000in}{4.620000in}}%
\pgfusepath{clip}%
\pgfsetbuttcap%
\pgfsetmiterjoin%
\definecolor{currentfill}{rgb}{0.000000,0.000000,0.000000}%
\pgfsetfillcolor{currentfill}%
\pgfsetlinewidth{0.000000pt}%
\definecolor{currentstroke}{rgb}{0.000000,0.000000,0.000000}%
\pgfsetstrokecolor{currentstroke}%
\pgfsetstrokeopacity{0.000000}%
\pgfsetdash{}{0pt}%
\pgfpathmoveto{\pgfqpoint{5.889150in}{0.660000in}}%
\pgfpathlineto{\pgfqpoint{5.895350in}{0.660000in}}%
\pgfpathlineto{\pgfqpoint{5.895350in}{3.235678in}}%
\pgfpathlineto{\pgfqpoint{5.889150in}{3.235678in}}%
\pgfpathlineto{\pgfqpoint{5.889150in}{0.660000in}}%
\pgfpathclose%
\pgfusepath{fill}%
\end{pgfscope}%
\begin{pgfscope}%
\pgfpathrectangle{\pgfqpoint{1.250000in}{0.660000in}}{\pgfqpoint{7.750000in}{4.620000in}}%
\pgfusepath{clip}%
\pgfsetbuttcap%
\pgfsetmiterjoin%
\definecolor{currentfill}{rgb}{0.000000,0.000000,0.000000}%
\pgfsetfillcolor{currentfill}%
\pgfsetlinewidth{0.000000pt}%
\definecolor{currentstroke}{rgb}{0.000000,0.000000,0.000000}%
\pgfsetstrokecolor{currentstroke}%
\pgfsetstrokeopacity{0.000000}%
\pgfsetdash{}{0pt}%
\pgfpathmoveto{\pgfqpoint{5.896900in}{0.660000in}}%
\pgfpathlineto{\pgfqpoint{5.903100in}{0.660000in}}%
\pgfpathlineto{\pgfqpoint{5.903100in}{3.235522in}}%
\pgfpathlineto{\pgfqpoint{5.896900in}{3.235522in}}%
\pgfpathlineto{\pgfqpoint{5.896900in}{0.660000in}}%
\pgfpathclose%
\pgfusepath{fill}%
\end{pgfscope}%
\begin{pgfscope}%
\pgfpathrectangle{\pgfqpoint{1.250000in}{0.660000in}}{\pgfqpoint{7.750000in}{4.620000in}}%
\pgfusepath{clip}%
\pgfsetbuttcap%
\pgfsetmiterjoin%
\definecolor{currentfill}{rgb}{0.000000,0.000000,0.000000}%
\pgfsetfillcolor{currentfill}%
\pgfsetlinewidth{0.000000pt}%
\definecolor{currentstroke}{rgb}{0.000000,0.000000,0.000000}%
\pgfsetstrokecolor{currentstroke}%
\pgfsetstrokeopacity{0.000000}%
\pgfsetdash{}{0pt}%
\pgfpathmoveto{\pgfqpoint{5.904650in}{0.660000in}}%
\pgfpathlineto{\pgfqpoint{5.910850in}{0.660000in}}%
\pgfpathlineto{\pgfqpoint{5.910850in}{3.235365in}}%
\pgfpathlineto{\pgfqpoint{5.904650in}{3.235365in}}%
\pgfpathlineto{\pgfqpoint{5.904650in}{0.660000in}}%
\pgfpathclose%
\pgfusepath{fill}%
\end{pgfscope}%
\begin{pgfscope}%
\pgfpathrectangle{\pgfqpoint{1.250000in}{0.660000in}}{\pgfqpoint{7.750000in}{4.620000in}}%
\pgfusepath{clip}%
\pgfsetbuttcap%
\pgfsetmiterjoin%
\definecolor{currentfill}{rgb}{0.000000,0.000000,0.000000}%
\pgfsetfillcolor{currentfill}%
\pgfsetlinewidth{0.000000pt}%
\definecolor{currentstroke}{rgb}{0.000000,0.000000,0.000000}%
\pgfsetstrokecolor{currentstroke}%
\pgfsetstrokeopacity{0.000000}%
\pgfsetdash{}{0pt}%
\pgfpathmoveto{\pgfqpoint{5.912400in}{0.660000in}}%
\pgfpathlineto{\pgfqpoint{5.918600in}{0.660000in}}%
\pgfpathlineto{\pgfqpoint{5.918600in}{3.235208in}}%
\pgfpathlineto{\pgfqpoint{5.912400in}{3.235208in}}%
\pgfpathlineto{\pgfqpoint{5.912400in}{0.660000in}}%
\pgfpathclose%
\pgfusepath{fill}%
\end{pgfscope}%
\begin{pgfscope}%
\pgfpathrectangle{\pgfqpoint{1.250000in}{0.660000in}}{\pgfqpoint{7.750000in}{4.620000in}}%
\pgfusepath{clip}%
\pgfsetbuttcap%
\pgfsetmiterjoin%
\definecolor{currentfill}{rgb}{0.000000,0.000000,0.000000}%
\pgfsetfillcolor{currentfill}%
\pgfsetlinewidth{0.000000pt}%
\definecolor{currentstroke}{rgb}{0.000000,0.000000,0.000000}%
\pgfsetstrokecolor{currentstroke}%
\pgfsetstrokeopacity{0.000000}%
\pgfsetdash{}{0pt}%
\pgfpathmoveto{\pgfqpoint{5.920150in}{0.660000in}}%
\pgfpathlineto{\pgfqpoint{5.926350in}{0.660000in}}%
\pgfpathlineto{\pgfqpoint{5.926350in}{3.235051in}}%
\pgfpathlineto{\pgfqpoint{5.920150in}{3.235051in}}%
\pgfpathlineto{\pgfqpoint{5.920150in}{0.660000in}}%
\pgfpathclose%
\pgfusepath{fill}%
\end{pgfscope}%
\begin{pgfscope}%
\pgfpathrectangle{\pgfqpoint{1.250000in}{0.660000in}}{\pgfqpoint{7.750000in}{4.620000in}}%
\pgfusepath{clip}%
\pgfsetbuttcap%
\pgfsetmiterjoin%
\definecolor{currentfill}{rgb}{0.000000,0.000000,0.000000}%
\pgfsetfillcolor{currentfill}%
\pgfsetlinewidth{0.000000pt}%
\definecolor{currentstroke}{rgb}{0.000000,0.000000,0.000000}%
\pgfsetstrokecolor{currentstroke}%
\pgfsetstrokeopacity{0.000000}%
\pgfsetdash{}{0pt}%
\pgfpathmoveto{\pgfqpoint{5.927900in}{0.660000in}}%
\pgfpathlineto{\pgfqpoint{5.934100in}{0.660000in}}%
\pgfpathlineto{\pgfqpoint{5.934100in}{3.234894in}}%
\pgfpathlineto{\pgfqpoint{5.927900in}{3.234894in}}%
\pgfpathlineto{\pgfqpoint{5.927900in}{0.660000in}}%
\pgfpathclose%
\pgfusepath{fill}%
\end{pgfscope}%
\begin{pgfscope}%
\pgfpathrectangle{\pgfqpoint{1.250000in}{0.660000in}}{\pgfqpoint{7.750000in}{4.620000in}}%
\pgfusepath{clip}%
\pgfsetbuttcap%
\pgfsetmiterjoin%
\definecolor{currentfill}{rgb}{0.000000,0.000000,0.000000}%
\pgfsetfillcolor{currentfill}%
\pgfsetlinewidth{0.000000pt}%
\definecolor{currentstroke}{rgb}{0.000000,0.000000,0.000000}%
\pgfsetstrokecolor{currentstroke}%
\pgfsetstrokeopacity{0.000000}%
\pgfsetdash{}{0pt}%
\pgfpathmoveto{\pgfqpoint{5.935650in}{0.660000in}}%
\pgfpathlineto{\pgfqpoint{5.941850in}{0.660000in}}%
\pgfpathlineto{\pgfqpoint{5.941850in}{3.234737in}}%
\pgfpathlineto{\pgfqpoint{5.935650in}{3.234737in}}%
\pgfpathlineto{\pgfqpoint{5.935650in}{0.660000in}}%
\pgfpathclose%
\pgfusepath{fill}%
\end{pgfscope}%
\begin{pgfscope}%
\pgfpathrectangle{\pgfqpoint{1.250000in}{0.660000in}}{\pgfqpoint{7.750000in}{4.620000in}}%
\pgfusepath{clip}%
\pgfsetbuttcap%
\pgfsetmiterjoin%
\definecolor{currentfill}{rgb}{0.000000,0.000000,0.000000}%
\pgfsetfillcolor{currentfill}%
\pgfsetlinewidth{0.000000pt}%
\definecolor{currentstroke}{rgb}{0.000000,0.000000,0.000000}%
\pgfsetstrokecolor{currentstroke}%
\pgfsetstrokeopacity{0.000000}%
\pgfsetdash{}{0pt}%
\pgfpathmoveto{\pgfqpoint{5.943400in}{0.660000in}}%
\pgfpathlineto{\pgfqpoint{5.949600in}{0.660000in}}%
\pgfpathlineto{\pgfqpoint{5.949600in}{3.234581in}}%
\pgfpathlineto{\pgfqpoint{5.943400in}{3.234581in}}%
\pgfpathlineto{\pgfqpoint{5.943400in}{0.660000in}}%
\pgfpathclose%
\pgfusepath{fill}%
\end{pgfscope}%
\begin{pgfscope}%
\pgfpathrectangle{\pgfqpoint{1.250000in}{0.660000in}}{\pgfqpoint{7.750000in}{4.620000in}}%
\pgfusepath{clip}%
\pgfsetbuttcap%
\pgfsetmiterjoin%
\definecolor{currentfill}{rgb}{0.000000,0.000000,0.000000}%
\pgfsetfillcolor{currentfill}%
\pgfsetlinewidth{0.000000pt}%
\definecolor{currentstroke}{rgb}{0.000000,0.000000,0.000000}%
\pgfsetstrokecolor{currentstroke}%
\pgfsetstrokeopacity{0.000000}%
\pgfsetdash{}{0pt}%
\pgfpathmoveto{\pgfqpoint{5.951150in}{0.660000in}}%
\pgfpathlineto{\pgfqpoint{5.957350in}{0.660000in}}%
\pgfpathlineto{\pgfqpoint{5.957350in}{3.234424in}}%
\pgfpathlineto{\pgfqpoint{5.951150in}{3.234424in}}%
\pgfpathlineto{\pgfqpoint{5.951150in}{0.660000in}}%
\pgfpathclose%
\pgfusepath{fill}%
\end{pgfscope}%
\begin{pgfscope}%
\pgfpathrectangle{\pgfqpoint{1.250000in}{0.660000in}}{\pgfqpoint{7.750000in}{4.620000in}}%
\pgfusepath{clip}%
\pgfsetbuttcap%
\pgfsetmiterjoin%
\definecolor{currentfill}{rgb}{0.000000,0.000000,0.000000}%
\pgfsetfillcolor{currentfill}%
\pgfsetlinewidth{0.000000pt}%
\definecolor{currentstroke}{rgb}{0.000000,0.000000,0.000000}%
\pgfsetstrokecolor{currentstroke}%
\pgfsetstrokeopacity{0.000000}%
\pgfsetdash{}{0pt}%
\pgfpathmoveto{\pgfqpoint{5.958900in}{0.660000in}}%
\pgfpathlineto{\pgfqpoint{5.965100in}{0.660000in}}%
\pgfpathlineto{\pgfqpoint{5.965100in}{3.234267in}}%
\pgfpathlineto{\pgfqpoint{5.958900in}{3.234267in}}%
\pgfpathlineto{\pgfqpoint{5.958900in}{0.660000in}}%
\pgfpathclose%
\pgfusepath{fill}%
\end{pgfscope}%
\begin{pgfscope}%
\pgfpathrectangle{\pgfqpoint{1.250000in}{0.660000in}}{\pgfqpoint{7.750000in}{4.620000in}}%
\pgfusepath{clip}%
\pgfsetbuttcap%
\pgfsetmiterjoin%
\definecolor{currentfill}{rgb}{0.000000,0.000000,0.000000}%
\pgfsetfillcolor{currentfill}%
\pgfsetlinewidth{0.000000pt}%
\definecolor{currentstroke}{rgb}{0.000000,0.000000,0.000000}%
\pgfsetstrokecolor{currentstroke}%
\pgfsetstrokeopacity{0.000000}%
\pgfsetdash{}{0pt}%
\pgfpathmoveto{\pgfqpoint{5.966650in}{0.660000in}}%
\pgfpathlineto{\pgfqpoint{5.972850in}{0.660000in}}%
\pgfpathlineto{\pgfqpoint{5.972850in}{3.234113in}}%
\pgfpathlineto{\pgfqpoint{5.966650in}{3.234113in}}%
\pgfpathlineto{\pgfqpoint{5.966650in}{0.660000in}}%
\pgfpathclose%
\pgfusepath{fill}%
\end{pgfscope}%
\begin{pgfscope}%
\pgfpathrectangle{\pgfqpoint{1.250000in}{0.660000in}}{\pgfqpoint{7.750000in}{4.620000in}}%
\pgfusepath{clip}%
\pgfsetbuttcap%
\pgfsetmiterjoin%
\definecolor{currentfill}{rgb}{0.000000,0.000000,0.000000}%
\pgfsetfillcolor{currentfill}%
\pgfsetlinewidth{0.000000pt}%
\definecolor{currentstroke}{rgb}{0.000000,0.000000,0.000000}%
\pgfsetstrokecolor{currentstroke}%
\pgfsetstrokeopacity{0.000000}%
\pgfsetdash{}{0pt}%
\pgfpathmoveto{\pgfqpoint{5.974400in}{0.660000in}}%
\pgfpathlineto{\pgfqpoint{5.980600in}{0.660000in}}%
\pgfpathlineto{\pgfqpoint{5.980600in}{3.233960in}}%
\pgfpathlineto{\pgfqpoint{5.974400in}{3.233960in}}%
\pgfpathlineto{\pgfqpoint{5.974400in}{0.660000in}}%
\pgfpathclose%
\pgfusepath{fill}%
\end{pgfscope}%
\begin{pgfscope}%
\pgfpathrectangle{\pgfqpoint{1.250000in}{0.660000in}}{\pgfqpoint{7.750000in}{4.620000in}}%
\pgfusepath{clip}%
\pgfsetbuttcap%
\pgfsetmiterjoin%
\definecolor{currentfill}{rgb}{0.000000,0.000000,0.000000}%
\pgfsetfillcolor{currentfill}%
\pgfsetlinewidth{0.000000pt}%
\definecolor{currentstroke}{rgb}{0.000000,0.000000,0.000000}%
\pgfsetstrokecolor{currentstroke}%
\pgfsetstrokeopacity{0.000000}%
\pgfsetdash{}{0pt}%
\pgfpathmoveto{\pgfqpoint{5.982150in}{0.660000in}}%
\pgfpathlineto{\pgfqpoint{5.988350in}{0.660000in}}%
\pgfpathlineto{\pgfqpoint{5.988350in}{3.233807in}}%
\pgfpathlineto{\pgfqpoint{5.982150in}{3.233807in}}%
\pgfpathlineto{\pgfqpoint{5.982150in}{0.660000in}}%
\pgfpathclose%
\pgfusepath{fill}%
\end{pgfscope}%
\begin{pgfscope}%
\pgfpathrectangle{\pgfqpoint{1.250000in}{0.660000in}}{\pgfqpoint{7.750000in}{4.620000in}}%
\pgfusepath{clip}%
\pgfsetbuttcap%
\pgfsetmiterjoin%
\definecolor{currentfill}{rgb}{0.000000,0.000000,0.000000}%
\pgfsetfillcolor{currentfill}%
\pgfsetlinewidth{0.000000pt}%
\definecolor{currentstroke}{rgb}{0.000000,0.000000,0.000000}%
\pgfsetstrokecolor{currentstroke}%
\pgfsetstrokeopacity{0.000000}%
\pgfsetdash{}{0pt}%
\pgfpathmoveto{\pgfqpoint{5.989900in}{0.660000in}}%
\pgfpathlineto{\pgfqpoint{5.996100in}{0.660000in}}%
\pgfpathlineto{\pgfqpoint{5.996100in}{3.233655in}}%
\pgfpathlineto{\pgfqpoint{5.989900in}{3.233655in}}%
\pgfpathlineto{\pgfqpoint{5.989900in}{0.660000in}}%
\pgfpathclose%
\pgfusepath{fill}%
\end{pgfscope}%
\begin{pgfscope}%
\pgfpathrectangle{\pgfqpoint{1.250000in}{0.660000in}}{\pgfqpoint{7.750000in}{4.620000in}}%
\pgfusepath{clip}%
\pgfsetbuttcap%
\pgfsetmiterjoin%
\definecolor{currentfill}{rgb}{0.000000,0.000000,0.000000}%
\pgfsetfillcolor{currentfill}%
\pgfsetlinewidth{0.000000pt}%
\definecolor{currentstroke}{rgb}{0.000000,0.000000,0.000000}%
\pgfsetstrokecolor{currentstroke}%
\pgfsetstrokeopacity{0.000000}%
\pgfsetdash{}{0pt}%
\pgfpathmoveto{\pgfqpoint{5.997650in}{0.660000in}}%
\pgfpathlineto{\pgfqpoint{6.003850in}{0.660000in}}%
\pgfpathlineto{\pgfqpoint{6.003850in}{3.233502in}}%
\pgfpathlineto{\pgfqpoint{5.997650in}{3.233502in}}%
\pgfpathlineto{\pgfqpoint{5.997650in}{0.660000in}}%
\pgfpathclose%
\pgfusepath{fill}%
\end{pgfscope}%
\begin{pgfscope}%
\pgfpathrectangle{\pgfqpoint{1.250000in}{0.660000in}}{\pgfqpoint{7.750000in}{4.620000in}}%
\pgfusepath{clip}%
\pgfsetbuttcap%
\pgfsetmiterjoin%
\definecolor{currentfill}{rgb}{0.000000,0.000000,0.000000}%
\pgfsetfillcolor{currentfill}%
\pgfsetlinewidth{0.000000pt}%
\definecolor{currentstroke}{rgb}{0.000000,0.000000,0.000000}%
\pgfsetstrokecolor{currentstroke}%
\pgfsetstrokeopacity{0.000000}%
\pgfsetdash{}{0pt}%
\pgfpathmoveto{\pgfqpoint{6.005400in}{0.660000in}}%
\pgfpathlineto{\pgfqpoint{6.011600in}{0.660000in}}%
\pgfpathlineto{\pgfqpoint{6.011600in}{3.233350in}}%
\pgfpathlineto{\pgfqpoint{6.005400in}{3.233350in}}%
\pgfpathlineto{\pgfqpoint{6.005400in}{0.660000in}}%
\pgfpathclose%
\pgfusepath{fill}%
\end{pgfscope}%
\begin{pgfscope}%
\pgfpathrectangle{\pgfqpoint{1.250000in}{0.660000in}}{\pgfqpoint{7.750000in}{4.620000in}}%
\pgfusepath{clip}%
\pgfsetbuttcap%
\pgfsetmiterjoin%
\definecolor{currentfill}{rgb}{0.000000,0.000000,0.000000}%
\pgfsetfillcolor{currentfill}%
\pgfsetlinewidth{0.000000pt}%
\definecolor{currentstroke}{rgb}{0.000000,0.000000,0.000000}%
\pgfsetstrokecolor{currentstroke}%
\pgfsetstrokeopacity{0.000000}%
\pgfsetdash{}{0pt}%
\pgfpathmoveto{\pgfqpoint{6.013150in}{0.660000in}}%
\pgfpathlineto{\pgfqpoint{6.019350in}{0.660000in}}%
\pgfpathlineto{\pgfqpoint{6.019350in}{3.233197in}}%
\pgfpathlineto{\pgfqpoint{6.013150in}{3.233197in}}%
\pgfpathlineto{\pgfqpoint{6.013150in}{0.660000in}}%
\pgfpathclose%
\pgfusepath{fill}%
\end{pgfscope}%
\begin{pgfscope}%
\pgfpathrectangle{\pgfqpoint{1.250000in}{0.660000in}}{\pgfqpoint{7.750000in}{4.620000in}}%
\pgfusepath{clip}%
\pgfsetbuttcap%
\pgfsetmiterjoin%
\definecolor{currentfill}{rgb}{0.000000,0.000000,0.000000}%
\pgfsetfillcolor{currentfill}%
\pgfsetlinewidth{0.000000pt}%
\definecolor{currentstroke}{rgb}{0.000000,0.000000,0.000000}%
\pgfsetstrokecolor{currentstroke}%
\pgfsetstrokeopacity{0.000000}%
\pgfsetdash{}{0pt}%
\pgfpathmoveto{\pgfqpoint{6.020900in}{0.660000in}}%
\pgfpathlineto{\pgfqpoint{6.027100in}{0.660000in}}%
\pgfpathlineto{\pgfqpoint{6.027100in}{3.233044in}}%
\pgfpathlineto{\pgfqpoint{6.020900in}{3.233044in}}%
\pgfpathlineto{\pgfqpoint{6.020900in}{0.660000in}}%
\pgfpathclose%
\pgfusepath{fill}%
\end{pgfscope}%
\begin{pgfscope}%
\pgfpathrectangle{\pgfqpoint{1.250000in}{0.660000in}}{\pgfqpoint{7.750000in}{4.620000in}}%
\pgfusepath{clip}%
\pgfsetbuttcap%
\pgfsetmiterjoin%
\definecolor{currentfill}{rgb}{0.000000,0.000000,0.000000}%
\pgfsetfillcolor{currentfill}%
\pgfsetlinewidth{0.000000pt}%
\definecolor{currentstroke}{rgb}{0.000000,0.000000,0.000000}%
\pgfsetstrokecolor{currentstroke}%
\pgfsetstrokeopacity{0.000000}%
\pgfsetdash{}{0pt}%
\pgfpathmoveto{\pgfqpoint{6.028650in}{0.660000in}}%
\pgfpathlineto{\pgfqpoint{6.034850in}{0.660000in}}%
\pgfpathlineto{\pgfqpoint{6.034850in}{3.232892in}}%
\pgfpathlineto{\pgfqpoint{6.028650in}{3.232892in}}%
\pgfpathlineto{\pgfqpoint{6.028650in}{0.660000in}}%
\pgfpathclose%
\pgfusepath{fill}%
\end{pgfscope}%
\begin{pgfscope}%
\pgfpathrectangle{\pgfqpoint{1.250000in}{0.660000in}}{\pgfqpoint{7.750000in}{4.620000in}}%
\pgfusepath{clip}%
\pgfsetbuttcap%
\pgfsetmiterjoin%
\definecolor{currentfill}{rgb}{0.000000,0.000000,0.000000}%
\pgfsetfillcolor{currentfill}%
\pgfsetlinewidth{0.000000pt}%
\definecolor{currentstroke}{rgb}{0.000000,0.000000,0.000000}%
\pgfsetstrokecolor{currentstroke}%
\pgfsetstrokeopacity{0.000000}%
\pgfsetdash{}{0pt}%
\pgfpathmoveto{\pgfqpoint{6.036400in}{0.660000in}}%
\pgfpathlineto{\pgfqpoint{6.042600in}{0.660000in}}%
\pgfpathlineto{\pgfqpoint{6.042600in}{3.232739in}}%
\pgfpathlineto{\pgfqpoint{6.036400in}{3.232739in}}%
\pgfpathlineto{\pgfqpoint{6.036400in}{0.660000in}}%
\pgfpathclose%
\pgfusepath{fill}%
\end{pgfscope}%
\begin{pgfscope}%
\pgfpathrectangle{\pgfqpoint{1.250000in}{0.660000in}}{\pgfqpoint{7.750000in}{4.620000in}}%
\pgfusepath{clip}%
\pgfsetbuttcap%
\pgfsetmiterjoin%
\definecolor{currentfill}{rgb}{0.000000,0.000000,0.000000}%
\pgfsetfillcolor{currentfill}%
\pgfsetlinewidth{0.000000pt}%
\definecolor{currentstroke}{rgb}{0.000000,0.000000,0.000000}%
\pgfsetstrokecolor{currentstroke}%
\pgfsetstrokeopacity{0.000000}%
\pgfsetdash{}{0pt}%
\pgfpathmoveto{\pgfqpoint{6.044150in}{0.660000in}}%
\pgfpathlineto{\pgfqpoint{6.050350in}{0.660000in}}%
\pgfpathlineto{\pgfqpoint{6.050350in}{3.232587in}}%
\pgfpathlineto{\pgfqpoint{6.044150in}{3.232587in}}%
\pgfpathlineto{\pgfqpoint{6.044150in}{0.660000in}}%
\pgfpathclose%
\pgfusepath{fill}%
\end{pgfscope}%
\begin{pgfscope}%
\pgfpathrectangle{\pgfqpoint{1.250000in}{0.660000in}}{\pgfqpoint{7.750000in}{4.620000in}}%
\pgfusepath{clip}%
\pgfsetbuttcap%
\pgfsetmiterjoin%
\definecolor{currentfill}{rgb}{0.000000,0.000000,0.000000}%
\pgfsetfillcolor{currentfill}%
\pgfsetlinewidth{0.000000pt}%
\definecolor{currentstroke}{rgb}{0.000000,0.000000,0.000000}%
\pgfsetstrokecolor{currentstroke}%
\pgfsetstrokeopacity{0.000000}%
\pgfsetdash{}{0pt}%
\pgfpathmoveto{\pgfqpoint{6.051900in}{0.660000in}}%
\pgfpathlineto{\pgfqpoint{6.058100in}{0.660000in}}%
\pgfpathlineto{\pgfqpoint{6.058100in}{3.232434in}}%
\pgfpathlineto{\pgfqpoint{6.051900in}{3.232434in}}%
\pgfpathlineto{\pgfqpoint{6.051900in}{0.660000in}}%
\pgfpathclose%
\pgfusepath{fill}%
\end{pgfscope}%
\begin{pgfscope}%
\pgfpathrectangle{\pgfqpoint{1.250000in}{0.660000in}}{\pgfqpoint{7.750000in}{4.620000in}}%
\pgfusepath{clip}%
\pgfsetbuttcap%
\pgfsetmiterjoin%
\definecolor{currentfill}{rgb}{0.000000,0.000000,0.000000}%
\pgfsetfillcolor{currentfill}%
\pgfsetlinewidth{0.000000pt}%
\definecolor{currentstroke}{rgb}{0.000000,0.000000,0.000000}%
\pgfsetstrokecolor{currentstroke}%
\pgfsetstrokeopacity{0.000000}%
\pgfsetdash{}{0pt}%
\pgfpathmoveto{\pgfqpoint{6.059650in}{0.660000in}}%
\pgfpathlineto{\pgfqpoint{6.065850in}{0.660000in}}%
\pgfpathlineto{\pgfqpoint{6.065850in}{3.232282in}}%
\pgfpathlineto{\pgfqpoint{6.059650in}{3.232282in}}%
\pgfpathlineto{\pgfqpoint{6.059650in}{0.660000in}}%
\pgfpathclose%
\pgfusepath{fill}%
\end{pgfscope}%
\begin{pgfscope}%
\pgfpathrectangle{\pgfqpoint{1.250000in}{0.660000in}}{\pgfqpoint{7.750000in}{4.620000in}}%
\pgfusepath{clip}%
\pgfsetbuttcap%
\pgfsetmiterjoin%
\definecolor{currentfill}{rgb}{0.000000,0.000000,0.000000}%
\pgfsetfillcolor{currentfill}%
\pgfsetlinewidth{0.000000pt}%
\definecolor{currentstroke}{rgb}{0.000000,0.000000,0.000000}%
\pgfsetstrokecolor{currentstroke}%
\pgfsetstrokeopacity{0.000000}%
\pgfsetdash{}{0pt}%
\pgfpathmoveto{\pgfqpoint{6.067400in}{0.660000in}}%
\pgfpathlineto{\pgfqpoint{6.073600in}{0.660000in}}%
\pgfpathlineto{\pgfqpoint{6.073600in}{3.232134in}}%
\pgfpathlineto{\pgfqpoint{6.067400in}{3.232134in}}%
\pgfpathlineto{\pgfqpoint{6.067400in}{0.660000in}}%
\pgfpathclose%
\pgfusepath{fill}%
\end{pgfscope}%
\begin{pgfscope}%
\pgfpathrectangle{\pgfqpoint{1.250000in}{0.660000in}}{\pgfqpoint{7.750000in}{4.620000in}}%
\pgfusepath{clip}%
\pgfsetbuttcap%
\pgfsetmiterjoin%
\definecolor{currentfill}{rgb}{0.000000,0.000000,0.000000}%
\pgfsetfillcolor{currentfill}%
\pgfsetlinewidth{0.000000pt}%
\definecolor{currentstroke}{rgb}{0.000000,0.000000,0.000000}%
\pgfsetstrokecolor{currentstroke}%
\pgfsetstrokeopacity{0.000000}%
\pgfsetdash{}{0pt}%
\pgfpathmoveto{\pgfqpoint{6.075150in}{0.660000in}}%
\pgfpathlineto{\pgfqpoint{6.081350in}{0.660000in}}%
\pgfpathlineto{\pgfqpoint{6.081350in}{3.231986in}}%
\pgfpathlineto{\pgfqpoint{6.075150in}{3.231986in}}%
\pgfpathlineto{\pgfqpoint{6.075150in}{0.660000in}}%
\pgfpathclose%
\pgfusepath{fill}%
\end{pgfscope}%
\begin{pgfscope}%
\pgfpathrectangle{\pgfqpoint{1.250000in}{0.660000in}}{\pgfqpoint{7.750000in}{4.620000in}}%
\pgfusepath{clip}%
\pgfsetbuttcap%
\pgfsetmiterjoin%
\definecolor{currentfill}{rgb}{0.000000,0.000000,0.000000}%
\pgfsetfillcolor{currentfill}%
\pgfsetlinewidth{0.000000pt}%
\definecolor{currentstroke}{rgb}{0.000000,0.000000,0.000000}%
\pgfsetstrokecolor{currentstroke}%
\pgfsetstrokeopacity{0.000000}%
\pgfsetdash{}{0pt}%
\pgfpathmoveto{\pgfqpoint{6.082900in}{0.660000in}}%
\pgfpathlineto{\pgfqpoint{6.089100in}{0.660000in}}%
\pgfpathlineto{\pgfqpoint{6.089100in}{3.231837in}}%
\pgfpathlineto{\pgfqpoint{6.082900in}{3.231837in}}%
\pgfpathlineto{\pgfqpoint{6.082900in}{0.660000in}}%
\pgfpathclose%
\pgfusepath{fill}%
\end{pgfscope}%
\begin{pgfscope}%
\pgfpathrectangle{\pgfqpoint{1.250000in}{0.660000in}}{\pgfqpoint{7.750000in}{4.620000in}}%
\pgfusepath{clip}%
\pgfsetbuttcap%
\pgfsetmiterjoin%
\definecolor{currentfill}{rgb}{0.000000,0.000000,0.000000}%
\pgfsetfillcolor{currentfill}%
\pgfsetlinewidth{0.000000pt}%
\definecolor{currentstroke}{rgb}{0.000000,0.000000,0.000000}%
\pgfsetstrokecolor{currentstroke}%
\pgfsetstrokeopacity{0.000000}%
\pgfsetdash{}{0pt}%
\pgfpathmoveto{\pgfqpoint{6.090650in}{0.660000in}}%
\pgfpathlineto{\pgfqpoint{6.096850in}{0.660000in}}%
\pgfpathlineto{\pgfqpoint{6.096850in}{3.231689in}}%
\pgfpathlineto{\pgfqpoint{6.090650in}{3.231689in}}%
\pgfpathlineto{\pgfqpoint{6.090650in}{0.660000in}}%
\pgfpathclose%
\pgfusepath{fill}%
\end{pgfscope}%
\begin{pgfscope}%
\pgfpathrectangle{\pgfqpoint{1.250000in}{0.660000in}}{\pgfqpoint{7.750000in}{4.620000in}}%
\pgfusepath{clip}%
\pgfsetbuttcap%
\pgfsetmiterjoin%
\definecolor{currentfill}{rgb}{0.000000,0.000000,0.000000}%
\pgfsetfillcolor{currentfill}%
\pgfsetlinewidth{0.000000pt}%
\definecolor{currentstroke}{rgb}{0.000000,0.000000,0.000000}%
\pgfsetstrokecolor{currentstroke}%
\pgfsetstrokeopacity{0.000000}%
\pgfsetdash{}{0pt}%
\pgfpathmoveto{\pgfqpoint{6.098400in}{0.660000in}}%
\pgfpathlineto{\pgfqpoint{6.104600in}{0.660000in}}%
\pgfpathlineto{\pgfqpoint{6.104600in}{3.231540in}}%
\pgfpathlineto{\pgfqpoint{6.098400in}{3.231540in}}%
\pgfpathlineto{\pgfqpoint{6.098400in}{0.660000in}}%
\pgfpathclose%
\pgfusepath{fill}%
\end{pgfscope}%
\begin{pgfscope}%
\pgfpathrectangle{\pgfqpoint{1.250000in}{0.660000in}}{\pgfqpoint{7.750000in}{4.620000in}}%
\pgfusepath{clip}%
\pgfsetbuttcap%
\pgfsetmiterjoin%
\definecolor{currentfill}{rgb}{0.000000,0.000000,0.000000}%
\pgfsetfillcolor{currentfill}%
\pgfsetlinewidth{0.000000pt}%
\definecolor{currentstroke}{rgb}{0.000000,0.000000,0.000000}%
\pgfsetstrokecolor{currentstroke}%
\pgfsetstrokeopacity{0.000000}%
\pgfsetdash{}{0pt}%
\pgfpathmoveto{\pgfqpoint{6.106150in}{0.660000in}}%
\pgfpathlineto{\pgfqpoint{6.112350in}{0.660000in}}%
\pgfpathlineto{\pgfqpoint{6.112350in}{3.231392in}}%
\pgfpathlineto{\pgfqpoint{6.106150in}{3.231392in}}%
\pgfpathlineto{\pgfqpoint{6.106150in}{0.660000in}}%
\pgfpathclose%
\pgfusepath{fill}%
\end{pgfscope}%
\begin{pgfscope}%
\pgfpathrectangle{\pgfqpoint{1.250000in}{0.660000in}}{\pgfqpoint{7.750000in}{4.620000in}}%
\pgfusepath{clip}%
\pgfsetbuttcap%
\pgfsetmiterjoin%
\definecolor{currentfill}{rgb}{0.000000,0.000000,0.000000}%
\pgfsetfillcolor{currentfill}%
\pgfsetlinewidth{0.000000pt}%
\definecolor{currentstroke}{rgb}{0.000000,0.000000,0.000000}%
\pgfsetstrokecolor{currentstroke}%
\pgfsetstrokeopacity{0.000000}%
\pgfsetdash{}{0pt}%
\pgfpathmoveto{\pgfqpoint{6.113900in}{0.660000in}}%
\pgfpathlineto{\pgfqpoint{6.120100in}{0.660000in}}%
\pgfpathlineto{\pgfqpoint{6.120100in}{3.231243in}}%
\pgfpathlineto{\pgfqpoint{6.113900in}{3.231243in}}%
\pgfpathlineto{\pgfqpoint{6.113900in}{0.660000in}}%
\pgfpathclose%
\pgfusepath{fill}%
\end{pgfscope}%
\begin{pgfscope}%
\pgfpathrectangle{\pgfqpoint{1.250000in}{0.660000in}}{\pgfqpoint{7.750000in}{4.620000in}}%
\pgfusepath{clip}%
\pgfsetbuttcap%
\pgfsetmiterjoin%
\definecolor{currentfill}{rgb}{0.000000,0.000000,0.000000}%
\pgfsetfillcolor{currentfill}%
\pgfsetlinewidth{0.000000pt}%
\definecolor{currentstroke}{rgb}{0.000000,0.000000,0.000000}%
\pgfsetstrokecolor{currentstroke}%
\pgfsetstrokeopacity{0.000000}%
\pgfsetdash{}{0pt}%
\pgfpathmoveto{\pgfqpoint{6.121650in}{0.660000in}}%
\pgfpathlineto{\pgfqpoint{6.127850in}{0.660000in}}%
\pgfpathlineto{\pgfqpoint{6.127850in}{3.231095in}}%
\pgfpathlineto{\pgfqpoint{6.121650in}{3.231095in}}%
\pgfpathlineto{\pgfqpoint{6.121650in}{0.660000in}}%
\pgfpathclose%
\pgfusepath{fill}%
\end{pgfscope}%
\begin{pgfscope}%
\pgfpathrectangle{\pgfqpoint{1.250000in}{0.660000in}}{\pgfqpoint{7.750000in}{4.620000in}}%
\pgfusepath{clip}%
\pgfsetbuttcap%
\pgfsetmiterjoin%
\definecolor{currentfill}{rgb}{0.000000,0.000000,0.000000}%
\pgfsetfillcolor{currentfill}%
\pgfsetlinewidth{0.000000pt}%
\definecolor{currentstroke}{rgb}{0.000000,0.000000,0.000000}%
\pgfsetstrokecolor{currentstroke}%
\pgfsetstrokeopacity{0.000000}%
\pgfsetdash{}{0pt}%
\pgfpathmoveto{\pgfqpoint{6.129400in}{0.660000in}}%
\pgfpathlineto{\pgfqpoint{6.135600in}{0.660000in}}%
\pgfpathlineto{\pgfqpoint{6.135600in}{3.230947in}}%
\pgfpathlineto{\pgfqpoint{6.129400in}{3.230947in}}%
\pgfpathlineto{\pgfqpoint{6.129400in}{0.660000in}}%
\pgfpathclose%
\pgfusepath{fill}%
\end{pgfscope}%
\begin{pgfscope}%
\pgfpathrectangle{\pgfqpoint{1.250000in}{0.660000in}}{\pgfqpoint{7.750000in}{4.620000in}}%
\pgfusepath{clip}%
\pgfsetbuttcap%
\pgfsetmiterjoin%
\definecolor{currentfill}{rgb}{0.000000,0.000000,0.000000}%
\pgfsetfillcolor{currentfill}%
\pgfsetlinewidth{0.000000pt}%
\definecolor{currentstroke}{rgb}{0.000000,0.000000,0.000000}%
\pgfsetstrokecolor{currentstroke}%
\pgfsetstrokeopacity{0.000000}%
\pgfsetdash{}{0pt}%
\pgfpathmoveto{\pgfqpoint{6.137150in}{0.660000in}}%
\pgfpathlineto{\pgfqpoint{6.143350in}{0.660000in}}%
\pgfpathlineto{\pgfqpoint{6.143350in}{3.230798in}}%
\pgfpathlineto{\pgfqpoint{6.137150in}{3.230798in}}%
\pgfpathlineto{\pgfqpoint{6.137150in}{0.660000in}}%
\pgfpathclose%
\pgfusepath{fill}%
\end{pgfscope}%
\begin{pgfscope}%
\pgfpathrectangle{\pgfqpoint{1.250000in}{0.660000in}}{\pgfqpoint{7.750000in}{4.620000in}}%
\pgfusepath{clip}%
\pgfsetbuttcap%
\pgfsetmiterjoin%
\definecolor{currentfill}{rgb}{0.000000,0.000000,0.000000}%
\pgfsetfillcolor{currentfill}%
\pgfsetlinewidth{0.000000pt}%
\definecolor{currentstroke}{rgb}{0.000000,0.000000,0.000000}%
\pgfsetstrokecolor{currentstroke}%
\pgfsetstrokeopacity{0.000000}%
\pgfsetdash{}{0pt}%
\pgfpathmoveto{\pgfqpoint{6.144900in}{0.660000in}}%
\pgfpathlineto{\pgfqpoint{6.151100in}{0.660000in}}%
\pgfpathlineto{\pgfqpoint{6.151100in}{3.230650in}}%
\pgfpathlineto{\pgfqpoint{6.144900in}{3.230650in}}%
\pgfpathlineto{\pgfqpoint{6.144900in}{0.660000in}}%
\pgfpathclose%
\pgfusepath{fill}%
\end{pgfscope}%
\begin{pgfscope}%
\pgfpathrectangle{\pgfqpoint{1.250000in}{0.660000in}}{\pgfqpoint{7.750000in}{4.620000in}}%
\pgfusepath{clip}%
\pgfsetbuttcap%
\pgfsetmiterjoin%
\definecolor{currentfill}{rgb}{0.000000,0.000000,0.000000}%
\pgfsetfillcolor{currentfill}%
\pgfsetlinewidth{0.000000pt}%
\definecolor{currentstroke}{rgb}{0.000000,0.000000,0.000000}%
\pgfsetstrokecolor{currentstroke}%
\pgfsetstrokeopacity{0.000000}%
\pgfsetdash{}{0pt}%
\pgfpathmoveto{\pgfqpoint{6.152650in}{0.660000in}}%
\pgfpathlineto{\pgfqpoint{6.158850in}{0.660000in}}%
\pgfpathlineto{\pgfqpoint{6.158850in}{3.230501in}}%
\pgfpathlineto{\pgfqpoint{6.152650in}{3.230501in}}%
\pgfpathlineto{\pgfqpoint{6.152650in}{0.660000in}}%
\pgfpathclose%
\pgfusepath{fill}%
\end{pgfscope}%
\begin{pgfscope}%
\pgfpathrectangle{\pgfqpoint{1.250000in}{0.660000in}}{\pgfqpoint{7.750000in}{4.620000in}}%
\pgfusepath{clip}%
\pgfsetbuttcap%
\pgfsetmiterjoin%
\definecolor{currentfill}{rgb}{0.000000,0.000000,0.000000}%
\pgfsetfillcolor{currentfill}%
\pgfsetlinewidth{0.000000pt}%
\definecolor{currentstroke}{rgb}{0.000000,0.000000,0.000000}%
\pgfsetstrokecolor{currentstroke}%
\pgfsetstrokeopacity{0.000000}%
\pgfsetdash{}{0pt}%
\pgfpathmoveto{\pgfqpoint{6.160400in}{0.660000in}}%
\pgfpathlineto{\pgfqpoint{6.166600in}{0.660000in}}%
\pgfpathlineto{\pgfqpoint{6.166600in}{3.230355in}}%
\pgfpathlineto{\pgfqpoint{6.160400in}{3.230355in}}%
\pgfpathlineto{\pgfqpoint{6.160400in}{0.660000in}}%
\pgfpathclose%
\pgfusepath{fill}%
\end{pgfscope}%
\begin{pgfscope}%
\pgfpathrectangle{\pgfqpoint{1.250000in}{0.660000in}}{\pgfqpoint{7.750000in}{4.620000in}}%
\pgfusepath{clip}%
\pgfsetbuttcap%
\pgfsetmiterjoin%
\definecolor{currentfill}{rgb}{0.000000,0.000000,0.000000}%
\pgfsetfillcolor{currentfill}%
\pgfsetlinewidth{0.000000pt}%
\definecolor{currentstroke}{rgb}{0.000000,0.000000,0.000000}%
\pgfsetstrokecolor{currentstroke}%
\pgfsetstrokeopacity{0.000000}%
\pgfsetdash{}{0pt}%
\pgfpathmoveto{\pgfqpoint{6.168150in}{0.660000in}}%
\pgfpathlineto{\pgfqpoint{6.174350in}{0.660000in}}%
\pgfpathlineto{\pgfqpoint{6.174350in}{3.230211in}}%
\pgfpathlineto{\pgfqpoint{6.168150in}{3.230211in}}%
\pgfpathlineto{\pgfqpoint{6.168150in}{0.660000in}}%
\pgfpathclose%
\pgfusepath{fill}%
\end{pgfscope}%
\begin{pgfscope}%
\pgfpathrectangle{\pgfqpoint{1.250000in}{0.660000in}}{\pgfqpoint{7.750000in}{4.620000in}}%
\pgfusepath{clip}%
\pgfsetbuttcap%
\pgfsetmiterjoin%
\definecolor{currentfill}{rgb}{0.000000,0.000000,0.000000}%
\pgfsetfillcolor{currentfill}%
\pgfsetlinewidth{0.000000pt}%
\definecolor{currentstroke}{rgb}{0.000000,0.000000,0.000000}%
\pgfsetstrokecolor{currentstroke}%
\pgfsetstrokeopacity{0.000000}%
\pgfsetdash{}{0pt}%
\pgfpathmoveto{\pgfqpoint{6.175900in}{0.660000in}}%
\pgfpathlineto{\pgfqpoint{6.182100in}{0.660000in}}%
\pgfpathlineto{\pgfqpoint{6.182100in}{3.230066in}}%
\pgfpathlineto{\pgfqpoint{6.175900in}{3.230066in}}%
\pgfpathlineto{\pgfqpoint{6.175900in}{0.660000in}}%
\pgfpathclose%
\pgfusepath{fill}%
\end{pgfscope}%
\begin{pgfscope}%
\pgfpathrectangle{\pgfqpoint{1.250000in}{0.660000in}}{\pgfqpoint{7.750000in}{4.620000in}}%
\pgfusepath{clip}%
\pgfsetbuttcap%
\pgfsetmiterjoin%
\definecolor{currentfill}{rgb}{0.000000,0.000000,0.000000}%
\pgfsetfillcolor{currentfill}%
\pgfsetlinewidth{0.000000pt}%
\definecolor{currentstroke}{rgb}{0.000000,0.000000,0.000000}%
\pgfsetstrokecolor{currentstroke}%
\pgfsetstrokeopacity{0.000000}%
\pgfsetdash{}{0pt}%
\pgfpathmoveto{\pgfqpoint{6.183650in}{0.660000in}}%
\pgfpathlineto{\pgfqpoint{6.189850in}{0.660000in}}%
\pgfpathlineto{\pgfqpoint{6.189850in}{3.229922in}}%
\pgfpathlineto{\pgfqpoint{6.183650in}{3.229922in}}%
\pgfpathlineto{\pgfqpoint{6.183650in}{0.660000in}}%
\pgfpathclose%
\pgfusepath{fill}%
\end{pgfscope}%
\begin{pgfscope}%
\pgfpathrectangle{\pgfqpoint{1.250000in}{0.660000in}}{\pgfqpoint{7.750000in}{4.620000in}}%
\pgfusepath{clip}%
\pgfsetbuttcap%
\pgfsetmiterjoin%
\definecolor{currentfill}{rgb}{0.000000,0.000000,0.000000}%
\pgfsetfillcolor{currentfill}%
\pgfsetlinewidth{0.000000pt}%
\definecolor{currentstroke}{rgb}{0.000000,0.000000,0.000000}%
\pgfsetstrokecolor{currentstroke}%
\pgfsetstrokeopacity{0.000000}%
\pgfsetdash{}{0pt}%
\pgfpathmoveto{\pgfqpoint{6.191400in}{0.660000in}}%
\pgfpathlineto{\pgfqpoint{6.197600in}{0.660000in}}%
\pgfpathlineto{\pgfqpoint{6.197600in}{3.229778in}}%
\pgfpathlineto{\pgfqpoint{6.191400in}{3.229778in}}%
\pgfpathlineto{\pgfqpoint{6.191400in}{0.660000in}}%
\pgfpathclose%
\pgfusepath{fill}%
\end{pgfscope}%
\begin{pgfscope}%
\pgfpathrectangle{\pgfqpoint{1.250000in}{0.660000in}}{\pgfqpoint{7.750000in}{4.620000in}}%
\pgfusepath{clip}%
\pgfsetbuttcap%
\pgfsetmiterjoin%
\definecolor{currentfill}{rgb}{0.000000,0.000000,0.000000}%
\pgfsetfillcolor{currentfill}%
\pgfsetlinewidth{0.000000pt}%
\definecolor{currentstroke}{rgb}{0.000000,0.000000,0.000000}%
\pgfsetstrokecolor{currentstroke}%
\pgfsetstrokeopacity{0.000000}%
\pgfsetdash{}{0pt}%
\pgfpathmoveto{\pgfqpoint{6.199150in}{0.660000in}}%
\pgfpathlineto{\pgfqpoint{6.205350in}{0.660000in}}%
\pgfpathlineto{\pgfqpoint{6.205350in}{3.229633in}}%
\pgfpathlineto{\pgfqpoint{6.199150in}{3.229633in}}%
\pgfpathlineto{\pgfqpoint{6.199150in}{0.660000in}}%
\pgfpathclose%
\pgfusepath{fill}%
\end{pgfscope}%
\begin{pgfscope}%
\pgfpathrectangle{\pgfqpoint{1.250000in}{0.660000in}}{\pgfqpoint{7.750000in}{4.620000in}}%
\pgfusepath{clip}%
\pgfsetbuttcap%
\pgfsetmiterjoin%
\definecolor{currentfill}{rgb}{0.000000,0.000000,0.000000}%
\pgfsetfillcolor{currentfill}%
\pgfsetlinewidth{0.000000pt}%
\definecolor{currentstroke}{rgb}{0.000000,0.000000,0.000000}%
\pgfsetstrokecolor{currentstroke}%
\pgfsetstrokeopacity{0.000000}%
\pgfsetdash{}{0pt}%
\pgfpathmoveto{\pgfqpoint{6.206900in}{0.660000in}}%
\pgfpathlineto{\pgfqpoint{6.213100in}{0.660000in}}%
\pgfpathlineto{\pgfqpoint{6.213100in}{3.229489in}}%
\pgfpathlineto{\pgfqpoint{6.206900in}{3.229489in}}%
\pgfpathlineto{\pgfqpoint{6.206900in}{0.660000in}}%
\pgfpathclose%
\pgfusepath{fill}%
\end{pgfscope}%
\begin{pgfscope}%
\pgfpathrectangle{\pgfqpoint{1.250000in}{0.660000in}}{\pgfqpoint{7.750000in}{4.620000in}}%
\pgfusepath{clip}%
\pgfsetbuttcap%
\pgfsetmiterjoin%
\definecolor{currentfill}{rgb}{0.000000,0.000000,0.000000}%
\pgfsetfillcolor{currentfill}%
\pgfsetlinewidth{0.000000pt}%
\definecolor{currentstroke}{rgb}{0.000000,0.000000,0.000000}%
\pgfsetstrokecolor{currentstroke}%
\pgfsetstrokeopacity{0.000000}%
\pgfsetdash{}{0pt}%
\pgfpathmoveto{\pgfqpoint{6.214650in}{0.660000in}}%
\pgfpathlineto{\pgfqpoint{6.220850in}{0.660000in}}%
\pgfpathlineto{\pgfqpoint{6.220850in}{3.229345in}}%
\pgfpathlineto{\pgfqpoint{6.214650in}{3.229345in}}%
\pgfpathlineto{\pgfqpoint{6.214650in}{0.660000in}}%
\pgfpathclose%
\pgfusepath{fill}%
\end{pgfscope}%
\begin{pgfscope}%
\pgfpathrectangle{\pgfqpoint{1.250000in}{0.660000in}}{\pgfqpoint{7.750000in}{4.620000in}}%
\pgfusepath{clip}%
\pgfsetbuttcap%
\pgfsetmiterjoin%
\definecolor{currentfill}{rgb}{0.000000,0.000000,0.000000}%
\pgfsetfillcolor{currentfill}%
\pgfsetlinewidth{0.000000pt}%
\definecolor{currentstroke}{rgb}{0.000000,0.000000,0.000000}%
\pgfsetstrokecolor{currentstroke}%
\pgfsetstrokeopacity{0.000000}%
\pgfsetdash{}{0pt}%
\pgfpathmoveto{\pgfqpoint{6.222400in}{0.660000in}}%
\pgfpathlineto{\pgfqpoint{6.228600in}{0.660000in}}%
\pgfpathlineto{\pgfqpoint{6.228600in}{3.229201in}}%
\pgfpathlineto{\pgfqpoint{6.222400in}{3.229201in}}%
\pgfpathlineto{\pgfqpoint{6.222400in}{0.660000in}}%
\pgfpathclose%
\pgfusepath{fill}%
\end{pgfscope}%
\begin{pgfscope}%
\pgfpathrectangle{\pgfqpoint{1.250000in}{0.660000in}}{\pgfqpoint{7.750000in}{4.620000in}}%
\pgfusepath{clip}%
\pgfsetbuttcap%
\pgfsetmiterjoin%
\definecolor{currentfill}{rgb}{0.000000,0.000000,0.000000}%
\pgfsetfillcolor{currentfill}%
\pgfsetlinewidth{0.000000pt}%
\definecolor{currentstroke}{rgb}{0.000000,0.000000,0.000000}%
\pgfsetstrokecolor{currentstroke}%
\pgfsetstrokeopacity{0.000000}%
\pgfsetdash{}{0pt}%
\pgfpathmoveto{\pgfqpoint{6.230150in}{0.660000in}}%
\pgfpathlineto{\pgfqpoint{6.236350in}{0.660000in}}%
\pgfpathlineto{\pgfqpoint{6.236350in}{3.229056in}}%
\pgfpathlineto{\pgfqpoint{6.230150in}{3.229056in}}%
\pgfpathlineto{\pgfqpoint{6.230150in}{0.660000in}}%
\pgfpathclose%
\pgfusepath{fill}%
\end{pgfscope}%
\begin{pgfscope}%
\pgfpathrectangle{\pgfqpoint{1.250000in}{0.660000in}}{\pgfqpoint{7.750000in}{4.620000in}}%
\pgfusepath{clip}%
\pgfsetbuttcap%
\pgfsetmiterjoin%
\definecolor{currentfill}{rgb}{0.000000,0.000000,0.000000}%
\pgfsetfillcolor{currentfill}%
\pgfsetlinewidth{0.000000pt}%
\definecolor{currentstroke}{rgb}{0.000000,0.000000,0.000000}%
\pgfsetstrokecolor{currentstroke}%
\pgfsetstrokeopacity{0.000000}%
\pgfsetdash{}{0pt}%
\pgfpathmoveto{\pgfqpoint{6.237900in}{0.660000in}}%
\pgfpathlineto{\pgfqpoint{6.244100in}{0.660000in}}%
\pgfpathlineto{\pgfqpoint{6.244100in}{3.228912in}}%
\pgfpathlineto{\pgfqpoint{6.237900in}{3.228912in}}%
\pgfpathlineto{\pgfqpoint{6.237900in}{0.660000in}}%
\pgfpathclose%
\pgfusepath{fill}%
\end{pgfscope}%
\begin{pgfscope}%
\pgfpathrectangle{\pgfqpoint{1.250000in}{0.660000in}}{\pgfqpoint{7.750000in}{4.620000in}}%
\pgfusepath{clip}%
\pgfsetbuttcap%
\pgfsetmiterjoin%
\definecolor{currentfill}{rgb}{0.000000,0.000000,0.000000}%
\pgfsetfillcolor{currentfill}%
\pgfsetlinewidth{0.000000pt}%
\definecolor{currentstroke}{rgb}{0.000000,0.000000,0.000000}%
\pgfsetstrokecolor{currentstroke}%
\pgfsetstrokeopacity{0.000000}%
\pgfsetdash{}{0pt}%
\pgfpathmoveto{\pgfqpoint{6.245650in}{0.660000in}}%
\pgfpathlineto{\pgfqpoint{6.251850in}{0.660000in}}%
\pgfpathlineto{\pgfqpoint{6.251850in}{3.228768in}}%
\pgfpathlineto{\pgfqpoint{6.245650in}{3.228768in}}%
\pgfpathlineto{\pgfqpoint{6.245650in}{0.660000in}}%
\pgfpathclose%
\pgfusepath{fill}%
\end{pgfscope}%
\begin{pgfscope}%
\pgfpathrectangle{\pgfqpoint{1.250000in}{0.660000in}}{\pgfqpoint{7.750000in}{4.620000in}}%
\pgfusepath{clip}%
\pgfsetbuttcap%
\pgfsetmiterjoin%
\definecolor{currentfill}{rgb}{0.000000,0.000000,0.000000}%
\pgfsetfillcolor{currentfill}%
\pgfsetlinewidth{0.000000pt}%
\definecolor{currentstroke}{rgb}{0.000000,0.000000,0.000000}%
\pgfsetstrokecolor{currentstroke}%
\pgfsetstrokeopacity{0.000000}%
\pgfsetdash{}{0pt}%
\pgfpathmoveto{\pgfqpoint{6.253400in}{0.660000in}}%
\pgfpathlineto{\pgfqpoint{6.259600in}{0.660000in}}%
\pgfpathlineto{\pgfqpoint{6.259600in}{3.228623in}}%
\pgfpathlineto{\pgfqpoint{6.253400in}{3.228623in}}%
\pgfpathlineto{\pgfqpoint{6.253400in}{0.660000in}}%
\pgfpathclose%
\pgfusepath{fill}%
\end{pgfscope}%
\begin{pgfscope}%
\pgfpathrectangle{\pgfqpoint{1.250000in}{0.660000in}}{\pgfqpoint{7.750000in}{4.620000in}}%
\pgfusepath{clip}%
\pgfsetbuttcap%
\pgfsetmiterjoin%
\definecolor{currentfill}{rgb}{0.000000,0.000000,0.000000}%
\pgfsetfillcolor{currentfill}%
\pgfsetlinewidth{0.000000pt}%
\definecolor{currentstroke}{rgb}{0.000000,0.000000,0.000000}%
\pgfsetstrokecolor{currentstroke}%
\pgfsetstrokeopacity{0.000000}%
\pgfsetdash{}{0pt}%
\pgfpathmoveto{\pgfqpoint{6.261150in}{0.660000in}}%
\pgfpathlineto{\pgfqpoint{6.267350in}{0.660000in}}%
\pgfpathlineto{\pgfqpoint{6.267350in}{3.228480in}}%
\pgfpathlineto{\pgfqpoint{6.261150in}{3.228480in}}%
\pgfpathlineto{\pgfqpoint{6.261150in}{0.660000in}}%
\pgfpathclose%
\pgfusepath{fill}%
\end{pgfscope}%
\begin{pgfscope}%
\pgfpathrectangle{\pgfqpoint{1.250000in}{0.660000in}}{\pgfqpoint{7.750000in}{4.620000in}}%
\pgfusepath{clip}%
\pgfsetbuttcap%
\pgfsetmiterjoin%
\definecolor{currentfill}{rgb}{0.000000,0.000000,0.000000}%
\pgfsetfillcolor{currentfill}%
\pgfsetlinewidth{0.000000pt}%
\definecolor{currentstroke}{rgb}{0.000000,0.000000,0.000000}%
\pgfsetstrokecolor{currentstroke}%
\pgfsetstrokeopacity{0.000000}%
\pgfsetdash{}{0pt}%
\pgfpathmoveto{\pgfqpoint{6.268900in}{0.660000in}}%
\pgfpathlineto{\pgfqpoint{6.275100in}{0.660000in}}%
\pgfpathlineto{\pgfqpoint{6.275100in}{3.228340in}}%
\pgfpathlineto{\pgfqpoint{6.268900in}{3.228340in}}%
\pgfpathlineto{\pgfqpoint{6.268900in}{0.660000in}}%
\pgfpathclose%
\pgfusepath{fill}%
\end{pgfscope}%
\begin{pgfscope}%
\pgfpathrectangle{\pgfqpoint{1.250000in}{0.660000in}}{\pgfqpoint{7.750000in}{4.620000in}}%
\pgfusepath{clip}%
\pgfsetbuttcap%
\pgfsetmiterjoin%
\definecolor{currentfill}{rgb}{0.000000,0.000000,0.000000}%
\pgfsetfillcolor{currentfill}%
\pgfsetlinewidth{0.000000pt}%
\definecolor{currentstroke}{rgb}{0.000000,0.000000,0.000000}%
\pgfsetstrokecolor{currentstroke}%
\pgfsetstrokeopacity{0.000000}%
\pgfsetdash{}{0pt}%
\pgfpathmoveto{\pgfqpoint{6.276650in}{0.660000in}}%
\pgfpathlineto{\pgfqpoint{6.282850in}{0.660000in}}%
\pgfpathlineto{\pgfqpoint{6.282850in}{3.228200in}}%
\pgfpathlineto{\pgfqpoint{6.276650in}{3.228200in}}%
\pgfpathlineto{\pgfqpoint{6.276650in}{0.660000in}}%
\pgfpathclose%
\pgfusepath{fill}%
\end{pgfscope}%
\begin{pgfscope}%
\pgfpathrectangle{\pgfqpoint{1.250000in}{0.660000in}}{\pgfqpoint{7.750000in}{4.620000in}}%
\pgfusepath{clip}%
\pgfsetbuttcap%
\pgfsetmiterjoin%
\definecolor{currentfill}{rgb}{0.000000,0.000000,0.000000}%
\pgfsetfillcolor{currentfill}%
\pgfsetlinewidth{0.000000pt}%
\definecolor{currentstroke}{rgb}{0.000000,0.000000,0.000000}%
\pgfsetstrokecolor{currentstroke}%
\pgfsetstrokeopacity{0.000000}%
\pgfsetdash{}{0pt}%
\pgfpathmoveto{\pgfqpoint{6.284400in}{0.660000in}}%
\pgfpathlineto{\pgfqpoint{6.290600in}{0.660000in}}%
\pgfpathlineto{\pgfqpoint{6.290600in}{3.228060in}}%
\pgfpathlineto{\pgfqpoint{6.284400in}{3.228060in}}%
\pgfpathlineto{\pgfqpoint{6.284400in}{0.660000in}}%
\pgfpathclose%
\pgfusepath{fill}%
\end{pgfscope}%
\begin{pgfscope}%
\pgfpathrectangle{\pgfqpoint{1.250000in}{0.660000in}}{\pgfqpoint{7.750000in}{4.620000in}}%
\pgfusepath{clip}%
\pgfsetbuttcap%
\pgfsetmiterjoin%
\definecolor{currentfill}{rgb}{0.000000,0.000000,0.000000}%
\pgfsetfillcolor{currentfill}%
\pgfsetlinewidth{0.000000pt}%
\definecolor{currentstroke}{rgb}{0.000000,0.000000,0.000000}%
\pgfsetstrokecolor{currentstroke}%
\pgfsetstrokeopacity{0.000000}%
\pgfsetdash{}{0pt}%
\pgfpathmoveto{\pgfqpoint{6.292150in}{0.660000in}}%
\pgfpathlineto{\pgfqpoint{6.298350in}{0.660000in}}%
\pgfpathlineto{\pgfqpoint{6.298350in}{3.227919in}}%
\pgfpathlineto{\pgfqpoint{6.292150in}{3.227919in}}%
\pgfpathlineto{\pgfqpoint{6.292150in}{0.660000in}}%
\pgfpathclose%
\pgfusepath{fill}%
\end{pgfscope}%
\begin{pgfscope}%
\pgfpathrectangle{\pgfqpoint{1.250000in}{0.660000in}}{\pgfqpoint{7.750000in}{4.620000in}}%
\pgfusepath{clip}%
\pgfsetbuttcap%
\pgfsetmiterjoin%
\definecolor{currentfill}{rgb}{0.000000,0.000000,0.000000}%
\pgfsetfillcolor{currentfill}%
\pgfsetlinewidth{0.000000pt}%
\definecolor{currentstroke}{rgb}{0.000000,0.000000,0.000000}%
\pgfsetstrokecolor{currentstroke}%
\pgfsetstrokeopacity{0.000000}%
\pgfsetdash{}{0pt}%
\pgfpathmoveto{\pgfqpoint{6.299900in}{0.660000in}}%
\pgfpathlineto{\pgfqpoint{6.306100in}{0.660000in}}%
\pgfpathlineto{\pgfqpoint{6.306100in}{3.227779in}}%
\pgfpathlineto{\pgfqpoint{6.299900in}{3.227779in}}%
\pgfpathlineto{\pgfqpoint{6.299900in}{0.660000in}}%
\pgfpathclose%
\pgfusepath{fill}%
\end{pgfscope}%
\begin{pgfscope}%
\pgfpathrectangle{\pgfqpoint{1.250000in}{0.660000in}}{\pgfqpoint{7.750000in}{4.620000in}}%
\pgfusepath{clip}%
\pgfsetbuttcap%
\pgfsetmiterjoin%
\definecolor{currentfill}{rgb}{0.000000,0.000000,0.000000}%
\pgfsetfillcolor{currentfill}%
\pgfsetlinewidth{0.000000pt}%
\definecolor{currentstroke}{rgb}{0.000000,0.000000,0.000000}%
\pgfsetstrokecolor{currentstroke}%
\pgfsetstrokeopacity{0.000000}%
\pgfsetdash{}{0pt}%
\pgfpathmoveto{\pgfqpoint{6.307650in}{0.660000in}}%
\pgfpathlineto{\pgfqpoint{6.313850in}{0.660000in}}%
\pgfpathlineto{\pgfqpoint{6.313850in}{3.227639in}}%
\pgfpathlineto{\pgfqpoint{6.307650in}{3.227639in}}%
\pgfpathlineto{\pgfqpoint{6.307650in}{0.660000in}}%
\pgfpathclose%
\pgfusepath{fill}%
\end{pgfscope}%
\begin{pgfscope}%
\pgfpathrectangle{\pgfqpoint{1.250000in}{0.660000in}}{\pgfqpoint{7.750000in}{4.620000in}}%
\pgfusepath{clip}%
\pgfsetbuttcap%
\pgfsetmiterjoin%
\definecolor{currentfill}{rgb}{0.000000,0.000000,0.000000}%
\pgfsetfillcolor{currentfill}%
\pgfsetlinewidth{0.000000pt}%
\definecolor{currentstroke}{rgb}{0.000000,0.000000,0.000000}%
\pgfsetstrokecolor{currentstroke}%
\pgfsetstrokeopacity{0.000000}%
\pgfsetdash{}{0pt}%
\pgfpathmoveto{\pgfqpoint{6.315400in}{0.660000in}}%
\pgfpathlineto{\pgfqpoint{6.321600in}{0.660000in}}%
\pgfpathlineto{\pgfqpoint{6.321600in}{3.227499in}}%
\pgfpathlineto{\pgfqpoint{6.315400in}{3.227499in}}%
\pgfpathlineto{\pgfqpoint{6.315400in}{0.660000in}}%
\pgfpathclose%
\pgfusepath{fill}%
\end{pgfscope}%
\begin{pgfscope}%
\pgfpathrectangle{\pgfqpoint{1.250000in}{0.660000in}}{\pgfqpoint{7.750000in}{4.620000in}}%
\pgfusepath{clip}%
\pgfsetbuttcap%
\pgfsetmiterjoin%
\definecolor{currentfill}{rgb}{0.000000,0.000000,0.000000}%
\pgfsetfillcolor{currentfill}%
\pgfsetlinewidth{0.000000pt}%
\definecolor{currentstroke}{rgb}{0.000000,0.000000,0.000000}%
\pgfsetstrokecolor{currentstroke}%
\pgfsetstrokeopacity{0.000000}%
\pgfsetdash{}{0pt}%
\pgfpathmoveto{\pgfqpoint{6.323150in}{0.660000in}}%
\pgfpathlineto{\pgfqpoint{6.329350in}{0.660000in}}%
\pgfpathlineto{\pgfqpoint{6.329350in}{3.227358in}}%
\pgfpathlineto{\pgfqpoint{6.323150in}{3.227358in}}%
\pgfpathlineto{\pgfqpoint{6.323150in}{0.660000in}}%
\pgfpathclose%
\pgfusepath{fill}%
\end{pgfscope}%
\begin{pgfscope}%
\pgfpathrectangle{\pgfqpoint{1.250000in}{0.660000in}}{\pgfqpoint{7.750000in}{4.620000in}}%
\pgfusepath{clip}%
\pgfsetbuttcap%
\pgfsetmiterjoin%
\definecolor{currentfill}{rgb}{0.000000,0.000000,0.000000}%
\pgfsetfillcolor{currentfill}%
\pgfsetlinewidth{0.000000pt}%
\definecolor{currentstroke}{rgb}{0.000000,0.000000,0.000000}%
\pgfsetstrokecolor{currentstroke}%
\pgfsetstrokeopacity{0.000000}%
\pgfsetdash{}{0pt}%
\pgfpathmoveto{\pgfqpoint{6.330900in}{0.660000in}}%
\pgfpathlineto{\pgfqpoint{6.337100in}{0.660000in}}%
\pgfpathlineto{\pgfqpoint{6.337100in}{3.227218in}}%
\pgfpathlineto{\pgfqpoint{6.330900in}{3.227218in}}%
\pgfpathlineto{\pgfqpoint{6.330900in}{0.660000in}}%
\pgfpathclose%
\pgfusepath{fill}%
\end{pgfscope}%
\begin{pgfscope}%
\pgfpathrectangle{\pgfqpoint{1.250000in}{0.660000in}}{\pgfqpoint{7.750000in}{4.620000in}}%
\pgfusepath{clip}%
\pgfsetbuttcap%
\pgfsetmiterjoin%
\definecolor{currentfill}{rgb}{0.000000,0.000000,0.000000}%
\pgfsetfillcolor{currentfill}%
\pgfsetlinewidth{0.000000pt}%
\definecolor{currentstroke}{rgb}{0.000000,0.000000,0.000000}%
\pgfsetstrokecolor{currentstroke}%
\pgfsetstrokeopacity{0.000000}%
\pgfsetdash{}{0pt}%
\pgfpathmoveto{\pgfqpoint{6.338650in}{0.660000in}}%
\pgfpathlineto{\pgfqpoint{6.344850in}{0.660000in}}%
\pgfpathlineto{\pgfqpoint{6.344850in}{3.227078in}}%
\pgfpathlineto{\pgfqpoint{6.338650in}{3.227078in}}%
\pgfpathlineto{\pgfqpoint{6.338650in}{0.660000in}}%
\pgfpathclose%
\pgfusepath{fill}%
\end{pgfscope}%
\begin{pgfscope}%
\pgfpathrectangle{\pgfqpoint{1.250000in}{0.660000in}}{\pgfqpoint{7.750000in}{4.620000in}}%
\pgfusepath{clip}%
\pgfsetbuttcap%
\pgfsetmiterjoin%
\definecolor{currentfill}{rgb}{0.000000,0.000000,0.000000}%
\pgfsetfillcolor{currentfill}%
\pgfsetlinewidth{0.000000pt}%
\definecolor{currentstroke}{rgb}{0.000000,0.000000,0.000000}%
\pgfsetstrokecolor{currentstroke}%
\pgfsetstrokeopacity{0.000000}%
\pgfsetdash{}{0pt}%
\pgfpathmoveto{\pgfqpoint{6.346400in}{0.660000in}}%
\pgfpathlineto{\pgfqpoint{6.352600in}{0.660000in}}%
\pgfpathlineto{\pgfqpoint{6.352600in}{3.226938in}}%
\pgfpathlineto{\pgfqpoint{6.346400in}{3.226938in}}%
\pgfpathlineto{\pgfqpoint{6.346400in}{0.660000in}}%
\pgfpathclose%
\pgfusepath{fill}%
\end{pgfscope}%
\begin{pgfscope}%
\pgfpathrectangle{\pgfqpoint{1.250000in}{0.660000in}}{\pgfqpoint{7.750000in}{4.620000in}}%
\pgfusepath{clip}%
\pgfsetbuttcap%
\pgfsetmiterjoin%
\definecolor{currentfill}{rgb}{0.000000,0.000000,0.000000}%
\pgfsetfillcolor{currentfill}%
\pgfsetlinewidth{0.000000pt}%
\definecolor{currentstroke}{rgb}{0.000000,0.000000,0.000000}%
\pgfsetstrokecolor{currentstroke}%
\pgfsetstrokeopacity{0.000000}%
\pgfsetdash{}{0pt}%
\pgfpathmoveto{\pgfqpoint{6.354150in}{0.660000in}}%
\pgfpathlineto{\pgfqpoint{6.360350in}{0.660000in}}%
\pgfpathlineto{\pgfqpoint{6.360350in}{3.226797in}}%
\pgfpathlineto{\pgfqpoint{6.354150in}{3.226797in}}%
\pgfpathlineto{\pgfqpoint{6.354150in}{0.660000in}}%
\pgfpathclose%
\pgfusepath{fill}%
\end{pgfscope}%
\begin{pgfscope}%
\pgfpathrectangle{\pgfqpoint{1.250000in}{0.660000in}}{\pgfqpoint{7.750000in}{4.620000in}}%
\pgfusepath{clip}%
\pgfsetbuttcap%
\pgfsetmiterjoin%
\definecolor{currentfill}{rgb}{0.000000,0.000000,0.000000}%
\pgfsetfillcolor{currentfill}%
\pgfsetlinewidth{0.000000pt}%
\definecolor{currentstroke}{rgb}{0.000000,0.000000,0.000000}%
\pgfsetstrokecolor{currentstroke}%
\pgfsetstrokeopacity{0.000000}%
\pgfsetdash{}{0pt}%
\pgfpathmoveto{\pgfqpoint{6.361900in}{0.660000in}}%
\pgfpathlineto{\pgfqpoint{6.368100in}{0.660000in}}%
\pgfpathlineto{\pgfqpoint{6.368100in}{3.226657in}}%
\pgfpathlineto{\pgfqpoint{6.361900in}{3.226657in}}%
\pgfpathlineto{\pgfqpoint{6.361900in}{0.660000in}}%
\pgfpathclose%
\pgfusepath{fill}%
\end{pgfscope}%
\begin{pgfscope}%
\pgfpathrectangle{\pgfqpoint{1.250000in}{0.660000in}}{\pgfqpoint{7.750000in}{4.620000in}}%
\pgfusepath{clip}%
\pgfsetbuttcap%
\pgfsetmiterjoin%
\definecolor{currentfill}{rgb}{0.000000,0.000000,0.000000}%
\pgfsetfillcolor{currentfill}%
\pgfsetlinewidth{0.000000pt}%
\definecolor{currentstroke}{rgb}{0.000000,0.000000,0.000000}%
\pgfsetstrokecolor{currentstroke}%
\pgfsetstrokeopacity{0.000000}%
\pgfsetdash{}{0pt}%
\pgfpathmoveto{\pgfqpoint{6.369650in}{0.660000in}}%
\pgfpathlineto{\pgfqpoint{6.375850in}{0.660000in}}%
\pgfpathlineto{\pgfqpoint{6.375850in}{3.226519in}}%
\pgfpathlineto{\pgfqpoint{6.369650in}{3.226519in}}%
\pgfpathlineto{\pgfqpoint{6.369650in}{0.660000in}}%
\pgfpathclose%
\pgfusepath{fill}%
\end{pgfscope}%
\begin{pgfscope}%
\pgfpathrectangle{\pgfqpoint{1.250000in}{0.660000in}}{\pgfqpoint{7.750000in}{4.620000in}}%
\pgfusepath{clip}%
\pgfsetbuttcap%
\pgfsetmiterjoin%
\definecolor{currentfill}{rgb}{0.000000,0.000000,0.000000}%
\pgfsetfillcolor{currentfill}%
\pgfsetlinewidth{0.000000pt}%
\definecolor{currentstroke}{rgb}{0.000000,0.000000,0.000000}%
\pgfsetstrokecolor{currentstroke}%
\pgfsetstrokeopacity{0.000000}%
\pgfsetdash{}{0pt}%
\pgfpathmoveto{\pgfqpoint{6.377400in}{0.660000in}}%
\pgfpathlineto{\pgfqpoint{6.383600in}{0.660000in}}%
\pgfpathlineto{\pgfqpoint{6.383600in}{3.226383in}}%
\pgfpathlineto{\pgfqpoint{6.377400in}{3.226383in}}%
\pgfpathlineto{\pgfqpoint{6.377400in}{0.660000in}}%
\pgfpathclose%
\pgfusepath{fill}%
\end{pgfscope}%
\begin{pgfscope}%
\pgfpathrectangle{\pgfqpoint{1.250000in}{0.660000in}}{\pgfqpoint{7.750000in}{4.620000in}}%
\pgfusepath{clip}%
\pgfsetbuttcap%
\pgfsetmiterjoin%
\definecolor{currentfill}{rgb}{0.000000,0.000000,0.000000}%
\pgfsetfillcolor{currentfill}%
\pgfsetlinewidth{0.000000pt}%
\definecolor{currentstroke}{rgb}{0.000000,0.000000,0.000000}%
\pgfsetstrokecolor{currentstroke}%
\pgfsetstrokeopacity{0.000000}%
\pgfsetdash{}{0pt}%
\pgfpathmoveto{\pgfqpoint{6.385150in}{0.660000in}}%
\pgfpathlineto{\pgfqpoint{6.391350in}{0.660000in}}%
\pgfpathlineto{\pgfqpoint{6.391350in}{3.226247in}}%
\pgfpathlineto{\pgfqpoint{6.385150in}{3.226247in}}%
\pgfpathlineto{\pgfqpoint{6.385150in}{0.660000in}}%
\pgfpathclose%
\pgfusepath{fill}%
\end{pgfscope}%
\begin{pgfscope}%
\pgfpathrectangle{\pgfqpoint{1.250000in}{0.660000in}}{\pgfqpoint{7.750000in}{4.620000in}}%
\pgfusepath{clip}%
\pgfsetbuttcap%
\pgfsetmiterjoin%
\definecolor{currentfill}{rgb}{0.000000,0.000000,0.000000}%
\pgfsetfillcolor{currentfill}%
\pgfsetlinewidth{0.000000pt}%
\definecolor{currentstroke}{rgb}{0.000000,0.000000,0.000000}%
\pgfsetstrokecolor{currentstroke}%
\pgfsetstrokeopacity{0.000000}%
\pgfsetdash{}{0pt}%
\pgfpathmoveto{\pgfqpoint{6.392900in}{0.660000in}}%
\pgfpathlineto{\pgfqpoint{6.399100in}{0.660000in}}%
\pgfpathlineto{\pgfqpoint{6.399100in}{3.226111in}}%
\pgfpathlineto{\pgfqpoint{6.392900in}{3.226111in}}%
\pgfpathlineto{\pgfqpoint{6.392900in}{0.660000in}}%
\pgfpathclose%
\pgfusepath{fill}%
\end{pgfscope}%
\begin{pgfscope}%
\pgfpathrectangle{\pgfqpoint{1.250000in}{0.660000in}}{\pgfqpoint{7.750000in}{4.620000in}}%
\pgfusepath{clip}%
\pgfsetbuttcap%
\pgfsetmiterjoin%
\definecolor{currentfill}{rgb}{0.000000,0.000000,0.000000}%
\pgfsetfillcolor{currentfill}%
\pgfsetlinewidth{0.000000pt}%
\definecolor{currentstroke}{rgb}{0.000000,0.000000,0.000000}%
\pgfsetstrokecolor{currentstroke}%
\pgfsetstrokeopacity{0.000000}%
\pgfsetdash{}{0pt}%
\pgfpathmoveto{\pgfqpoint{6.400650in}{0.660000in}}%
\pgfpathlineto{\pgfqpoint{6.406850in}{0.660000in}}%
\pgfpathlineto{\pgfqpoint{6.406850in}{3.225974in}}%
\pgfpathlineto{\pgfqpoint{6.400650in}{3.225974in}}%
\pgfpathlineto{\pgfqpoint{6.400650in}{0.660000in}}%
\pgfpathclose%
\pgfusepath{fill}%
\end{pgfscope}%
\begin{pgfscope}%
\pgfpathrectangle{\pgfqpoint{1.250000in}{0.660000in}}{\pgfqpoint{7.750000in}{4.620000in}}%
\pgfusepath{clip}%
\pgfsetbuttcap%
\pgfsetmiterjoin%
\definecolor{currentfill}{rgb}{0.000000,0.000000,0.000000}%
\pgfsetfillcolor{currentfill}%
\pgfsetlinewidth{0.000000pt}%
\definecolor{currentstroke}{rgb}{0.000000,0.000000,0.000000}%
\pgfsetstrokecolor{currentstroke}%
\pgfsetstrokeopacity{0.000000}%
\pgfsetdash{}{0pt}%
\pgfpathmoveto{\pgfqpoint{6.408400in}{0.660000in}}%
\pgfpathlineto{\pgfqpoint{6.414600in}{0.660000in}}%
\pgfpathlineto{\pgfqpoint{6.414600in}{3.225838in}}%
\pgfpathlineto{\pgfqpoint{6.408400in}{3.225838in}}%
\pgfpathlineto{\pgfqpoint{6.408400in}{0.660000in}}%
\pgfpathclose%
\pgfusepath{fill}%
\end{pgfscope}%
\begin{pgfscope}%
\pgfpathrectangle{\pgfqpoint{1.250000in}{0.660000in}}{\pgfqpoint{7.750000in}{4.620000in}}%
\pgfusepath{clip}%
\pgfsetbuttcap%
\pgfsetmiterjoin%
\definecolor{currentfill}{rgb}{0.000000,0.000000,0.000000}%
\pgfsetfillcolor{currentfill}%
\pgfsetlinewidth{0.000000pt}%
\definecolor{currentstroke}{rgb}{0.000000,0.000000,0.000000}%
\pgfsetstrokecolor{currentstroke}%
\pgfsetstrokeopacity{0.000000}%
\pgfsetdash{}{0pt}%
\pgfpathmoveto{\pgfqpoint{6.416150in}{0.660000in}}%
\pgfpathlineto{\pgfqpoint{6.422350in}{0.660000in}}%
\pgfpathlineto{\pgfqpoint{6.422350in}{3.225702in}}%
\pgfpathlineto{\pgfqpoint{6.416150in}{3.225702in}}%
\pgfpathlineto{\pgfqpoint{6.416150in}{0.660000in}}%
\pgfpathclose%
\pgfusepath{fill}%
\end{pgfscope}%
\begin{pgfscope}%
\pgfpathrectangle{\pgfqpoint{1.250000in}{0.660000in}}{\pgfqpoint{7.750000in}{4.620000in}}%
\pgfusepath{clip}%
\pgfsetbuttcap%
\pgfsetmiterjoin%
\definecolor{currentfill}{rgb}{0.000000,0.000000,0.000000}%
\pgfsetfillcolor{currentfill}%
\pgfsetlinewidth{0.000000pt}%
\definecolor{currentstroke}{rgb}{0.000000,0.000000,0.000000}%
\pgfsetstrokecolor{currentstroke}%
\pgfsetstrokeopacity{0.000000}%
\pgfsetdash{}{0pt}%
\pgfpathmoveto{\pgfqpoint{6.423900in}{0.660000in}}%
\pgfpathlineto{\pgfqpoint{6.430100in}{0.660000in}}%
\pgfpathlineto{\pgfqpoint{6.430100in}{3.225566in}}%
\pgfpathlineto{\pgfqpoint{6.423900in}{3.225566in}}%
\pgfpathlineto{\pgfqpoint{6.423900in}{0.660000in}}%
\pgfpathclose%
\pgfusepath{fill}%
\end{pgfscope}%
\begin{pgfscope}%
\pgfpathrectangle{\pgfqpoint{1.250000in}{0.660000in}}{\pgfqpoint{7.750000in}{4.620000in}}%
\pgfusepath{clip}%
\pgfsetbuttcap%
\pgfsetmiterjoin%
\definecolor{currentfill}{rgb}{0.000000,0.000000,0.000000}%
\pgfsetfillcolor{currentfill}%
\pgfsetlinewidth{0.000000pt}%
\definecolor{currentstroke}{rgb}{0.000000,0.000000,0.000000}%
\pgfsetstrokecolor{currentstroke}%
\pgfsetstrokeopacity{0.000000}%
\pgfsetdash{}{0pt}%
\pgfpathmoveto{\pgfqpoint{6.431650in}{0.660000in}}%
\pgfpathlineto{\pgfqpoint{6.437850in}{0.660000in}}%
\pgfpathlineto{\pgfqpoint{6.437850in}{3.225429in}}%
\pgfpathlineto{\pgfqpoint{6.431650in}{3.225429in}}%
\pgfpathlineto{\pgfqpoint{6.431650in}{0.660000in}}%
\pgfpathclose%
\pgfusepath{fill}%
\end{pgfscope}%
\begin{pgfscope}%
\pgfpathrectangle{\pgfqpoint{1.250000in}{0.660000in}}{\pgfqpoint{7.750000in}{4.620000in}}%
\pgfusepath{clip}%
\pgfsetbuttcap%
\pgfsetmiterjoin%
\definecolor{currentfill}{rgb}{0.000000,0.000000,0.000000}%
\pgfsetfillcolor{currentfill}%
\pgfsetlinewidth{0.000000pt}%
\definecolor{currentstroke}{rgb}{0.000000,0.000000,0.000000}%
\pgfsetstrokecolor{currentstroke}%
\pgfsetstrokeopacity{0.000000}%
\pgfsetdash{}{0pt}%
\pgfpathmoveto{\pgfqpoint{6.439400in}{0.660000in}}%
\pgfpathlineto{\pgfqpoint{6.445600in}{0.660000in}}%
\pgfpathlineto{\pgfqpoint{6.445600in}{3.225293in}}%
\pgfpathlineto{\pgfqpoint{6.439400in}{3.225293in}}%
\pgfpathlineto{\pgfqpoint{6.439400in}{0.660000in}}%
\pgfpathclose%
\pgfusepath{fill}%
\end{pgfscope}%
\begin{pgfscope}%
\pgfpathrectangle{\pgfqpoint{1.250000in}{0.660000in}}{\pgfqpoint{7.750000in}{4.620000in}}%
\pgfusepath{clip}%
\pgfsetbuttcap%
\pgfsetmiterjoin%
\definecolor{currentfill}{rgb}{0.000000,0.000000,0.000000}%
\pgfsetfillcolor{currentfill}%
\pgfsetlinewidth{0.000000pt}%
\definecolor{currentstroke}{rgb}{0.000000,0.000000,0.000000}%
\pgfsetstrokecolor{currentstroke}%
\pgfsetstrokeopacity{0.000000}%
\pgfsetdash{}{0pt}%
\pgfpathmoveto{\pgfqpoint{6.447150in}{0.660000in}}%
\pgfpathlineto{\pgfqpoint{6.453350in}{0.660000in}}%
\pgfpathlineto{\pgfqpoint{6.453350in}{3.225157in}}%
\pgfpathlineto{\pgfqpoint{6.447150in}{3.225157in}}%
\pgfpathlineto{\pgfqpoint{6.447150in}{0.660000in}}%
\pgfpathclose%
\pgfusepath{fill}%
\end{pgfscope}%
\begin{pgfscope}%
\pgfpathrectangle{\pgfqpoint{1.250000in}{0.660000in}}{\pgfqpoint{7.750000in}{4.620000in}}%
\pgfusepath{clip}%
\pgfsetbuttcap%
\pgfsetmiterjoin%
\definecolor{currentfill}{rgb}{0.000000,0.000000,0.000000}%
\pgfsetfillcolor{currentfill}%
\pgfsetlinewidth{0.000000pt}%
\definecolor{currentstroke}{rgb}{0.000000,0.000000,0.000000}%
\pgfsetstrokecolor{currentstroke}%
\pgfsetstrokeopacity{0.000000}%
\pgfsetdash{}{0pt}%
\pgfpathmoveto{\pgfqpoint{6.454900in}{0.660000in}}%
\pgfpathlineto{\pgfqpoint{6.461100in}{0.660000in}}%
\pgfpathlineto{\pgfqpoint{6.461100in}{3.225021in}}%
\pgfpathlineto{\pgfqpoint{6.454900in}{3.225021in}}%
\pgfpathlineto{\pgfqpoint{6.454900in}{0.660000in}}%
\pgfpathclose%
\pgfusepath{fill}%
\end{pgfscope}%
\begin{pgfscope}%
\pgfpathrectangle{\pgfqpoint{1.250000in}{0.660000in}}{\pgfqpoint{7.750000in}{4.620000in}}%
\pgfusepath{clip}%
\pgfsetbuttcap%
\pgfsetmiterjoin%
\definecolor{currentfill}{rgb}{0.000000,0.000000,0.000000}%
\pgfsetfillcolor{currentfill}%
\pgfsetlinewidth{0.000000pt}%
\definecolor{currentstroke}{rgb}{0.000000,0.000000,0.000000}%
\pgfsetstrokecolor{currentstroke}%
\pgfsetstrokeopacity{0.000000}%
\pgfsetdash{}{0pt}%
\pgfpathmoveto{\pgfqpoint{6.462650in}{0.660000in}}%
\pgfpathlineto{\pgfqpoint{6.468850in}{0.660000in}}%
\pgfpathlineto{\pgfqpoint{6.468850in}{3.224884in}}%
\pgfpathlineto{\pgfqpoint{6.462650in}{3.224884in}}%
\pgfpathlineto{\pgfqpoint{6.462650in}{0.660000in}}%
\pgfpathclose%
\pgfusepath{fill}%
\end{pgfscope}%
\begin{pgfscope}%
\pgfpathrectangle{\pgfqpoint{1.250000in}{0.660000in}}{\pgfqpoint{7.750000in}{4.620000in}}%
\pgfusepath{clip}%
\pgfsetbuttcap%
\pgfsetmiterjoin%
\definecolor{currentfill}{rgb}{0.000000,0.000000,0.000000}%
\pgfsetfillcolor{currentfill}%
\pgfsetlinewidth{0.000000pt}%
\definecolor{currentstroke}{rgb}{0.000000,0.000000,0.000000}%
\pgfsetstrokecolor{currentstroke}%
\pgfsetstrokeopacity{0.000000}%
\pgfsetdash{}{0pt}%
\pgfpathmoveto{\pgfqpoint{6.470400in}{0.660000in}}%
\pgfpathlineto{\pgfqpoint{6.476600in}{0.660000in}}%
\pgfpathlineto{\pgfqpoint{6.476600in}{3.224748in}}%
\pgfpathlineto{\pgfqpoint{6.470400in}{3.224748in}}%
\pgfpathlineto{\pgfqpoint{6.470400in}{0.660000in}}%
\pgfpathclose%
\pgfusepath{fill}%
\end{pgfscope}%
\begin{pgfscope}%
\pgfpathrectangle{\pgfqpoint{1.250000in}{0.660000in}}{\pgfqpoint{7.750000in}{4.620000in}}%
\pgfusepath{clip}%
\pgfsetbuttcap%
\pgfsetmiterjoin%
\definecolor{currentfill}{rgb}{0.000000,0.000000,0.000000}%
\pgfsetfillcolor{currentfill}%
\pgfsetlinewidth{0.000000pt}%
\definecolor{currentstroke}{rgb}{0.000000,0.000000,0.000000}%
\pgfsetstrokecolor{currentstroke}%
\pgfsetstrokeopacity{0.000000}%
\pgfsetdash{}{0pt}%
\pgfpathmoveto{\pgfqpoint{6.478150in}{0.660000in}}%
\pgfpathlineto{\pgfqpoint{6.484350in}{0.660000in}}%
\pgfpathlineto{\pgfqpoint{6.484350in}{3.224614in}}%
\pgfpathlineto{\pgfqpoint{6.478150in}{3.224614in}}%
\pgfpathlineto{\pgfqpoint{6.478150in}{0.660000in}}%
\pgfpathclose%
\pgfusepath{fill}%
\end{pgfscope}%
\begin{pgfscope}%
\pgfpathrectangle{\pgfqpoint{1.250000in}{0.660000in}}{\pgfqpoint{7.750000in}{4.620000in}}%
\pgfusepath{clip}%
\pgfsetbuttcap%
\pgfsetmiterjoin%
\definecolor{currentfill}{rgb}{0.000000,0.000000,0.000000}%
\pgfsetfillcolor{currentfill}%
\pgfsetlinewidth{0.000000pt}%
\definecolor{currentstroke}{rgb}{0.000000,0.000000,0.000000}%
\pgfsetstrokecolor{currentstroke}%
\pgfsetstrokeopacity{0.000000}%
\pgfsetdash{}{0pt}%
\pgfpathmoveto{\pgfqpoint{6.485900in}{0.660000in}}%
\pgfpathlineto{\pgfqpoint{6.492100in}{0.660000in}}%
\pgfpathlineto{\pgfqpoint{6.492100in}{3.224481in}}%
\pgfpathlineto{\pgfqpoint{6.485900in}{3.224481in}}%
\pgfpathlineto{\pgfqpoint{6.485900in}{0.660000in}}%
\pgfpathclose%
\pgfusepath{fill}%
\end{pgfscope}%
\begin{pgfscope}%
\pgfpathrectangle{\pgfqpoint{1.250000in}{0.660000in}}{\pgfqpoint{7.750000in}{4.620000in}}%
\pgfusepath{clip}%
\pgfsetbuttcap%
\pgfsetmiterjoin%
\definecolor{currentfill}{rgb}{0.000000,0.000000,0.000000}%
\pgfsetfillcolor{currentfill}%
\pgfsetlinewidth{0.000000pt}%
\definecolor{currentstroke}{rgb}{0.000000,0.000000,0.000000}%
\pgfsetstrokecolor{currentstroke}%
\pgfsetstrokeopacity{0.000000}%
\pgfsetdash{}{0pt}%
\pgfpathmoveto{\pgfqpoint{6.493650in}{0.660000in}}%
\pgfpathlineto{\pgfqpoint{6.499850in}{0.660000in}}%
\pgfpathlineto{\pgfqpoint{6.499850in}{3.224349in}}%
\pgfpathlineto{\pgfqpoint{6.493650in}{3.224349in}}%
\pgfpathlineto{\pgfqpoint{6.493650in}{0.660000in}}%
\pgfpathclose%
\pgfusepath{fill}%
\end{pgfscope}%
\begin{pgfscope}%
\pgfpathrectangle{\pgfqpoint{1.250000in}{0.660000in}}{\pgfqpoint{7.750000in}{4.620000in}}%
\pgfusepath{clip}%
\pgfsetbuttcap%
\pgfsetmiterjoin%
\definecolor{currentfill}{rgb}{0.000000,0.000000,0.000000}%
\pgfsetfillcolor{currentfill}%
\pgfsetlinewidth{0.000000pt}%
\definecolor{currentstroke}{rgb}{0.000000,0.000000,0.000000}%
\pgfsetstrokecolor{currentstroke}%
\pgfsetstrokeopacity{0.000000}%
\pgfsetdash{}{0pt}%
\pgfpathmoveto{\pgfqpoint{6.501400in}{0.660000in}}%
\pgfpathlineto{\pgfqpoint{6.507600in}{0.660000in}}%
\pgfpathlineto{\pgfqpoint{6.507600in}{3.224217in}}%
\pgfpathlineto{\pgfqpoint{6.501400in}{3.224217in}}%
\pgfpathlineto{\pgfqpoint{6.501400in}{0.660000in}}%
\pgfpathclose%
\pgfusepath{fill}%
\end{pgfscope}%
\begin{pgfscope}%
\pgfpathrectangle{\pgfqpoint{1.250000in}{0.660000in}}{\pgfqpoint{7.750000in}{4.620000in}}%
\pgfusepath{clip}%
\pgfsetbuttcap%
\pgfsetmiterjoin%
\definecolor{currentfill}{rgb}{0.000000,0.000000,0.000000}%
\pgfsetfillcolor{currentfill}%
\pgfsetlinewidth{0.000000pt}%
\definecolor{currentstroke}{rgb}{0.000000,0.000000,0.000000}%
\pgfsetstrokecolor{currentstroke}%
\pgfsetstrokeopacity{0.000000}%
\pgfsetdash{}{0pt}%
\pgfpathmoveto{\pgfqpoint{6.509150in}{0.660000in}}%
\pgfpathlineto{\pgfqpoint{6.515350in}{0.660000in}}%
\pgfpathlineto{\pgfqpoint{6.515350in}{3.224085in}}%
\pgfpathlineto{\pgfqpoint{6.509150in}{3.224085in}}%
\pgfpathlineto{\pgfqpoint{6.509150in}{0.660000in}}%
\pgfpathclose%
\pgfusepath{fill}%
\end{pgfscope}%
\begin{pgfscope}%
\pgfpathrectangle{\pgfqpoint{1.250000in}{0.660000in}}{\pgfqpoint{7.750000in}{4.620000in}}%
\pgfusepath{clip}%
\pgfsetbuttcap%
\pgfsetmiterjoin%
\definecolor{currentfill}{rgb}{0.000000,0.000000,0.000000}%
\pgfsetfillcolor{currentfill}%
\pgfsetlinewidth{0.000000pt}%
\definecolor{currentstroke}{rgb}{0.000000,0.000000,0.000000}%
\pgfsetstrokecolor{currentstroke}%
\pgfsetstrokeopacity{0.000000}%
\pgfsetdash{}{0pt}%
\pgfpathmoveto{\pgfqpoint{6.516900in}{0.660000in}}%
\pgfpathlineto{\pgfqpoint{6.523100in}{0.660000in}}%
\pgfpathlineto{\pgfqpoint{6.523100in}{3.223952in}}%
\pgfpathlineto{\pgfqpoint{6.516900in}{3.223952in}}%
\pgfpathlineto{\pgfqpoint{6.516900in}{0.660000in}}%
\pgfpathclose%
\pgfusepath{fill}%
\end{pgfscope}%
\begin{pgfscope}%
\pgfpathrectangle{\pgfqpoint{1.250000in}{0.660000in}}{\pgfqpoint{7.750000in}{4.620000in}}%
\pgfusepath{clip}%
\pgfsetbuttcap%
\pgfsetmiterjoin%
\definecolor{currentfill}{rgb}{0.000000,0.000000,0.000000}%
\pgfsetfillcolor{currentfill}%
\pgfsetlinewidth{0.000000pt}%
\definecolor{currentstroke}{rgb}{0.000000,0.000000,0.000000}%
\pgfsetstrokecolor{currentstroke}%
\pgfsetstrokeopacity{0.000000}%
\pgfsetdash{}{0pt}%
\pgfpathmoveto{\pgfqpoint{6.524650in}{0.660000in}}%
\pgfpathlineto{\pgfqpoint{6.530850in}{0.660000in}}%
\pgfpathlineto{\pgfqpoint{6.530850in}{3.223820in}}%
\pgfpathlineto{\pgfqpoint{6.524650in}{3.223820in}}%
\pgfpathlineto{\pgfqpoint{6.524650in}{0.660000in}}%
\pgfpathclose%
\pgfusepath{fill}%
\end{pgfscope}%
\begin{pgfscope}%
\pgfpathrectangle{\pgfqpoint{1.250000in}{0.660000in}}{\pgfqpoint{7.750000in}{4.620000in}}%
\pgfusepath{clip}%
\pgfsetbuttcap%
\pgfsetmiterjoin%
\definecolor{currentfill}{rgb}{0.000000,0.000000,0.000000}%
\pgfsetfillcolor{currentfill}%
\pgfsetlinewidth{0.000000pt}%
\definecolor{currentstroke}{rgb}{0.000000,0.000000,0.000000}%
\pgfsetstrokecolor{currentstroke}%
\pgfsetstrokeopacity{0.000000}%
\pgfsetdash{}{0pt}%
\pgfpathmoveto{\pgfqpoint{6.532400in}{0.660000in}}%
\pgfpathlineto{\pgfqpoint{6.538600in}{0.660000in}}%
\pgfpathlineto{\pgfqpoint{6.538600in}{3.223688in}}%
\pgfpathlineto{\pgfqpoint{6.532400in}{3.223688in}}%
\pgfpathlineto{\pgfqpoint{6.532400in}{0.660000in}}%
\pgfpathclose%
\pgfusepath{fill}%
\end{pgfscope}%
\begin{pgfscope}%
\pgfpathrectangle{\pgfqpoint{1.250000in}{0.660000in}}{\pgfqpoint{7.750000in}{4.620000in}}%
\pgfusepath{clip}%
\pgfsetbuttcap%
\pgfsetmiterjoin%
\definecolor{currentfill}{rgb}{0.000000,0.000000,0.000000}%
\pgfsetfillcolor{currentfill}%
\pgfsetlinewidth{0.000000pt}%
\definecolor{currentstroke}{rgb}{0.000000,0.000000,0.000000}%
\pgfsetstrokecolor{currentstroke}%
\pgfsetstrokeopacity{0.000000}%
\pgfsetdash{}{0pt}%
\pgfpathmoveto{\pgfqpoint{6.540150in}{0.660000in}}%
\pgfpathlineto{\pgfqpoint{6.546350in}{0.660000in}}%
\pgfpathlineto{\pgfqpoint{6.546350in}{3.223555in}}%
\pgfpathlineto{\pgfqpoint{6.540150in}{3.223555in}}%
\pgfpathlineto{\pgfqpoint{6.540150in}{0.660000in}}%
\pgfpathclose%
\pgfusepath{fill}%
\end{pgfscope}%
\begin{pgfscope}%
\pgfpathrectangle{\pgfqpoint{1.250000in}{0.660000in}}{\pgfqpoint{7.750000in}{4.620000in}}%
\pgfusepath{clip}%
\pgfsetbuttcap%
\pgfsetmiterjoin%
\definecolor{currentfill}{rgb}{0.000000,0.000000,0.000000}%
\pgfsetfillcolor{currentfill}%
\pgfsetlinewidth{0.000000pt}%
\definecolor{currentstroke}{rgb}{0.000000,0.000000,0.000000}%
\pgfsetstrokecolor{currentstroke}%
\pgfsetstrokeopacity{0.000000}%
\pgfsetdash{}{0pt}%
\pgfpathmoveto{\pgfqpoint{6.547900in}{0.660000in}}%
\pgfpathlineto{\pgfqpoint{6.554100in}{0.660000in}}%
\pgfpathlineto{\pgfqpoint{6.554100in}{3.223423in}}%
\pgfpathlineto{\pgfqpoint{6.547900in}{3.223423in}}%
\pgfpathlineto{\pgfqpoint{6.547900in}{0.660000in}}%
\pgfpathclose%
\pgfusepath{fill}%
\end{pgfscope}%
\begin{pgfscope}%
\pgfpathrectangle{\pgfqpoint{1.250000in}{0.660000in}}{\pgfqpoint{7.750000in}{4.620000in}}%
\pgfusepath{clip}%
\pgfsetbuttcap%
\pgfsetmiterjoin%
\definecolor{currentfill}{rgb}{0.000000,0.000000,0.000000}%
\pgfsetfillcolor{currentfill}%
\pgfsetlinewidth{0.000000pt}%
\definecolor{currentstroke}{rgb}{0.000000,0.000000,0.000000}%
\pgfsetstrokecolor{currentstroke}%
\pgfsetstrokeopacity{0.000000}%
\pgfsetdash{}{0pt}%
\pgfpathmoveto{\pgfqpoint{6.555650in}{0.660000in}}%
\pgfpathlineto{\pgfqpoint{6.561850in}{0.660000in}}%
\pgfpathlineto{\pgfqpoint{6.561850in}{3.223291in}}%
\pgfpathlineto{\pgfqpoint{6.555650in}{3.223291in}}%
\pgfpathlineto{\pgfqpoint{6.555650in}{0.660000in}}%
\pgfpathclose%
\pgfusepath{fill}%
\end{pgfscope}%
\begin{pgfscope}%
\pgfpathrectangle{\pgfqpoint{1.250000in}{0.660000in}}{\pgfqpoint{7.750000in}{4.620000in}}%
\pgfusepath{clip}%
\pgfsetbuttcap%
\pgfsetmiterjoin%
\definecolor{currentfill}{rgb}{0.000000,0.000000,0.000000}%
\pgfsetfillcolor{currentfill}%
\pgfsetlinewidth{0.000000pt}%
\definecolor{currentstroke}{rgb}{0.000000,0.000000,0.000000}%
\pgfsetstrokecolor{currentstroke}%
\pgfsetstrokeopacity{0.000000}%
\pgfsetdash{}{0pt}%
\pgfpathmoveto{\pgfqpoint{6.563400in}{0.660000in}}%
\pgfpathlineto{\pgfqpoint{6.569600in}{0.660000in}}%
\pgfpathlineto{\pgfqpoint{6.569600in}{3.223158in}}%
\pgfpathlineto{\pgfqpoint{6.563400in}{3.223158in}}%
\pgfpathlineto{\pgfqpoint{6.563400in}{0.660000in}}%
\pgfpathclose%
\pgfusepath{fill}%
\end{pgfscope}%
\begin{pgfscope}%
\pgfpathrectangle{\pgfqpoint{1.250000in}{0.660000in}}{\pgfqpoint{7.750000in}{4.620000in}}%
\pgfusepath{clip}%
\pgfsetbuttcap%
\pgfsetmiterjoin%
\definecolor{currentfill}{rgb}{0.000000,0.000000,0.000000}%
\pgfsetfillcolor{currentfill}%
\pgfsetlinewidth{0.000000pt}%
\definecolor{currentstroke}{rgb}{0.000000,0.000000,0.000000}%
\pgfsetstrokecolor{currentstroke}%
\pgfsetstrokeopacity{0.000000}%
\pgfsetdash{}{0pt}%
\pgfpathmoveto{\pgfqpoint{6.571150in}{0.660000in}}%
\pgfpathlineto{\pgfqpoint{6.577350in}{0.660000in}}%
\pgfpathlineto{\pgfqpoint{6.577350in}{3.223026in}}%
\pgfpathlineto{\pgfqpoint{6.571150in}{3.223026in}}%
\pgfpathlineto{\pgfqpoint{6.571150in}{0.660000in}}%
\pgfpathclose%
\pgfusepath{fill}%
\end{pgfscope}%
\begin{pgfscope}%
\pgfpathrectangle{\pgfqpoint{1.250000in}{0.660000in}}{\pgfqpoint{7.750000in}{4.620000in}}%
\pgfusepath{clip}%
\pgfsetbuttcap%
\pgfsetmiterjoin%
\definecolor{currentfill}{rgb}{0.000000,0.000000,0.000000}%
\pgfsetfillcolor{currentfill}%
\pgfsetlinewidth{0.000000pt}%
\definecolor{currentstroke}{rgb}{0.000000,0.000000,0.000000}%
\pgfsetstrokecolor{currentstroke}%
\pgfsetstrokeopacity{0.000000}%
\pgfsetdash{}{0pt}%
\pgfpathmoveto{\pgfqpoint{6.578900in}{0.660000in}}%
\pgfpathlineto{\pgfqpoint{6.585100in}{0.660000in}}%
\pgfpathlineto{\pgfqpoint{6.585100in}{3.222894in}}%
\pgfpathlineto{\pgfqpoint{6.578900in}{3.222894in}}%
\pgfpathlineto{\pgfqpoint{6.578900in}{0.660000in}}%
\pgfpathclose%
\pgfusepath{fill}%
\end{pgfscope}%
\begin{pgfscope}%
\pgfpathrectangle{\pgfqpoint{1.250000in}{0.660000in}}{\pgfqpoint{7.750000in}{4.620000in}}%
\pgfusepath{clip}%
\pgfsetbuttcap%
\pgfsetmiterjoin%
\definecolor{currentfill}{rgb}{0.000000,0.000000,0.000000}%
\pgfsetfillcolor{currentfill}%
\pgfsetlinewidth{0.000000pt}%
\definecolor{currentstroke}{rgb}{0.000000,0.000000,0.000000}%
\pgfsetstrokecolor{currentstroke}%
\pgfsetstrokeopacity{0.000000}%
\pgfsetdash{}{0pt}%
\pgfpathmoveto{\pgfqpoint{6.586650in}{0.660000in}}%
\pgfpathlineto{\pgfqpoint{6.592850in}{0.660000in}}%
\pgfpathlineto{\pgfqpoint{6.592850in}{3.222761in}}%
\pgfpathlineto{\pgfqpoint{6.586650in}{3.222761in}}%
\pgfpathlineto{\pgfqpoint{6.586650in}{0.660000in}}%
\pgfpathclose%
\pgfusepath{fill}%
\end{pgfscope}%
\begin{pgfscope}%
\pgfpathrectangle{\pgfqpoint{1.250000in}{0.660000in}}{\pgfqpoint{7.750000in}{4.620000in}}%
\pgfusepath{clip}%
\pgfsetbuttcap%
\pgfsetmiterjoin%
\definecolor{currentfill}{rgb}{0.000000,0.000000,0.000000}%
\pgfsetfillcolor{currentfill}%
\pgfsetlinewidth{0.000000pt}%
\definecolor{currentstroke}{rgb}{0.000000,0.000000,0.000000}%
\pgfsetstrokecolor{currentstroke}%
\pgfsetstrokeopacity{0.000000}%
\pgfsetdash{}{0pt}%
\pgfpathmoveto{\pgfqpoint{6.594400in}{0.660000in}}%
\pgfpathlineto{\pgfqpoint{6.600600in}{0.660000in}}%
\pgfpathlineto{\pgfqpoint{6.600600in}{3.222632in}}%
\pgfpathlineto{\pgfqpoint{6.594400in}{3.222632in}}%
\pgfpathlineto{\pgfqpoint{6.594400in}{0.660000in}}%
\pgfpathclose%
\pgfusepath{fill}%
\end{pgfscope}%
\begin{pgfscope}%
\pgfpathrectangle{\pgfqpoint{1.250000in}{0.660000in}}{\pgfqpoint{7.750000in}{4.620000in}}%
\pgfusepath{clip}%
\pgfsetbuttcap%
\pgfsetmiterjoin%
\definecolor{currentfill}{rgb}{0.000000,0.000000,0.000000}%
\pgfsetfillcolor{currentfill}%
\pgfsetlinewidth{0.000000pt}%
\definecolor{currentstroke}{rgb}{0.000000,0.000000,0.000000}%
\pgfsetstrokecolor{currentstroke}%
\pgfsetstrokeopacity{0.000000}%
\pgfsetdash{}{0pt}%
\pgfpathmoveto{\pgfqpoint{6.602150in}{0.660000in}}%
\pgfpathlineto{\pgfqpoint{6.608350in}{0.660000in}}%
\pgfpathlineto{\pgfqpoint{6.608350in}{3.222504in}}%
\pgfpathlineto{\pgfqpoint{6.602150in}{3.222504in}}%
\pgfpathlineto{\pgfqpoint{6.602150in}{0.660000in}}%
\pgfpathclose%
\pgfusepath{fill}%
\end{pgfscope}%
\begin{pgfscope}%
\pgfpathrectangle{\pgfqpoint{1.250000in}{0.660000in}}{\pgfqpoint{7.750000in}{4.620000in}}%
\pgfusepath{clip}%
\pgfsetbuttcap%
\pgfsetmiterjoin%
\definecolor{currentfill}{rgb}{0.000000,0.000000,0.000000}%
\pgfsetfillcolor{currentfill}%
\pgfsetlinewidth{0.000000pt}%
\definecolor{currentstroke}{rgb}{0.000000,0.000000,0.000000}%
\pgfsetstrokecolor{currentstroke}%
\pgfsetstrokeopacity{0.000000}%
\pgfsetdash{}{0pt}%
\pgfpathmoveto{\pgfqpoint{6.609900in}{0.660000in}}%
\pgfpathlineto{\pgfqpoint{6.616100in}{0.660000in}}%
\pgfpathlineto{\pgfqpoint{6.616100in}{3.222375in}}%
\pgfpathlineto{\pgfqpoint{6.609900in}{3.222375in}}%
\pgfpathlineto{\pgfqpoint{6.609900in}{0.660000in}}%
\pgfpathclose%
\pgfusepath{fill}%
\end{pgfscope}%
\begin{pgfscope}%
\pgfpathrectangle{\pgfqpoint{1.250000in}{0.660000in}}{\pgfqpoint{7.750000in}{4.620000in}}%
\pgfusepath{clip}%
\pgfsetbuttcap%
\pgfsetmiterjoin%
\definecolor{currentfill}{rgb}{0.000000,0.000000,0.000000}%
\pgfsetfillcolor{currentfill}%
\pgfsetlinewidth{0.000000pt}%
\definecolor{currentstroke}{rgb}{0.000000,0.000000,0.000000}%
\pgfsetstrokecolor{currentstroke}%
\pgfsetstrokeopacity{0.000000}%
\pgfsetdash{}{0pt}%
\pgfpathmoveto{\pgfqpoint{6.617650in}{0.660000in}}%
\pgfpathlineto{\pgfqpoint{6.623850in}{0.660000in}}%
\pgfpathlineto{\pgfqpoint{6.623850in}{3.222247in}}%
\pgfpathlineto{\pgfqpoint{6.617650in}{3.222247in}}%
\pgfpathlineto{\pgfqpoint{6.617650in}{0.660000in}}%
\pgfpathclose%
\pgfusepath{fill}%
\end{pgfscope}%
\begin{pgfscope}%
\pgfpathrectangle{\pgfqpoint{1.250000in}{0.660000in}}{\pgfqpoint{7.750000in}{4.620000in}}%
\pgfusepath{clip}%
\pgfsetbuttcap%
\pgfsetmiterjoin%
\definecolor{currentfill}{rgb}{0.000000,0.000000,0.000000}%
\pgfsetfillcolor{currentfill}%
\pgfsetlinewidth{0.000000pt}%
\definecolor{currentstroke}{rgb}{0.000000,0.000000,0.000000}%
\pgfsetstrokecolor{currentstroke}%
\pgfsetstrokeopacity{0.000000}%
\pgfsetdash{}{0pt}%
\pgfpathmoveto{\pgfqpoint{6.625400in}{0.660000in}}%
\pgfpathlineto{\pgfqpoint{6.631600in}{0.660000in}}%
\pgfpathlineto{\pgfqpoint{6.631600in}{3.222119in}}%
\pgfpathlineto{\pgfqpoint{6.625400in}{3.222119in}}%
\pgfpathlineto{\pgfqpoint{6.625400in}{0.660000in}}%
\pgfpathclose%
\pgfusepath{fill}%
\end{pgfscope}%
\begin{pgfscope}%
\pgfpathrectangle{\pgfqpoint{1.250000in}{0.660000in}}{\pgfqpoint{7.750000in}{4.620000in}}%
\pgfusepath{clip}%
\pgfsetbuttcap%
\pgfsetmiterjoin%
\definecolor{currentfill}{rgb}{0.000000,0.000000,0.000000}%
\pgfsetfillcolor{currentfill}%
\pgfsetlinewidth{0.000000pt}%
\definecolor{currentstroke}{rgb}{0.000000,0.000000,0.000000}%
\pgfsetstrokecolor{currentstroke}%
\pgfsetstrokeopacity{0.000000}%
\pgfsetdash{}{0pt}%
\pgfpathmoveto{\pgfqpoint{6.633150in}{0.660000in}}%
\pgfpathlineto{\pgfqpoint{6.639350in}{0.660000in}}%
\pgfpathlineto{\pgfqpoint{6.639350in}{3.221990in}}%
\pgfpathlineto{\pgfqpoint{6.633150in}{3.221990in}}%
\pgfpathlineto{\pgfqpoint{6.633150in}{0.660000in}}%
\pgfpathclose%
\pgfusepath{fill}%
\end{pgfscope}%
\begin{pgfscope}%
\pgfpathrectangle{\pgfqpoint{1.250000in}{0.660000in}}{\pgfqpoint{7.750000in}{4.620000in}}%
\pgfusepath{clip}%
\pgfsetbuttcap%
\pgfsetmiterjoin%
\definecolor{currentfill}{rgb}{0.000000,0.000000,0.000000}%
\pgfsetfillcolor{currentfill}%
\pgfsetlinewidth{0.000000pt}%
\definecolor{currentstroke}{rgb}{0.000000,0.000000,0.000000}%
\pgfsetstrokecolor{currentstroke}%
\pgfsetstrokeopacity{0.000000}%
\pgfsetdash{}{0pt}%
\pgfpathmoveto{\pgfqpoint{6.640900in}{0.660000in}}%
\pgfpathlineto{\pgfqpoint{6.647100in}{0.660000in}}%
\pgfpathlineto{\pgfqpoint{6.647100in}{3.221862in}}%
\pgfpathlineto{\pgfqpoint{6.640900in}{3.221862in}}%
\pgfpathlineto{\pgfqpoint{6.640900in}{0.660000in}}%
\pgfpathclose%
\pgfusepath{fill}%
\end{pgfscope}%
\begin{pgfscope}%
\pgfpathrectangle{\pgfqpoint{1.250000in}{0.660000in}}{\pgfqpoint{7.750000in}{4.620000in}}%
\pgfusepath{clip}%
\pgfsetbuttcap%
\pgfsetmiterjoin%
\definecolor{currentfill}{rgb}{0.000000,0.000000,0.000000}%
\pgfsetfillcolor{currentfill}%
\pgfsetlinewidth{0.000000pt}%
\definecolor{currentstroke}{rgb}{0.000000,0.000000,0.000000}%
\pgfsetstrokecolor{currentstroke}%
\pgfsetstrokeopacity{0.000000}%
\pgfsetdash{}{0pt}%
\pgfpathmoveto{\pgfqpoint{6.648650in}{0.660000in}}%
\pgfpathlineto{\pgfqpoint{6.654850in}{0.660000in}}%
\pgfpathlineto{\pgfqpoint{6.654850in}{3.221733in}}%
\pgfpathlineto{\pgfqpoint{6.648650in}{3.221733in}}%
\pgfpathlineto{\pgfqpoint{6.648650in}{0.660000in}}%
\pgfpathclose%
\pgfusepath{fill}%
\end{pgfscope}%
\begin{pgfscope}%
\pgfpathrectangle{\pgfqpoint{1.250000in}{0.660000in}}{\pgfqpoint{7.750000in}{4.620000in}}%
\pgfusepath{clip}%
\pgfsetbuttcap%
\pgfsetmiterjoin%
\definecolor{currentfill}{rgb}{0.000000,0.000000,0.000000}%
\pgfsetfillcolor{currentfill}%
\pgfsetlinewidth{0.000000pt}%
\definecolor{currentstroke}{rgb}{0.000000,0.000000,0.000000}%
\pgfsetstrokecolor{currentstroke}%
\pgfsetstrokeopacity{0.000000}%
\pgfsetdash{}{0pt}%
\pgfpathmoveto{\pgfqpoint{6.656400in}{0.660000in}}%
\pgfpathlineto{\pgfqpoint{6.662600in}{0.660000in}}%
\pgfpathlineto{\pgfqpoint{6.662600in}{3.221605in}}%
\pgfpathlineto{\pgfqpoint{6.656400in}{3.221605in}}%
\pgfpathlineto{\pgfqpoint{6.656400in}{0.660000in}}%
\pgfpathclose%
\pgfusepath{fill}%
\end{pgfscope}%
\begin{pgfscope}%
\pgfpathrectangle{\pgfqpoint{1.250000in}{0.660000in}}{\pgfqpoint{7.750000in}{4.620000in}}%
\pgfusepath{clip}%
\pgfsetbuttcap%
\pgfsetmiterjoin%
\definecolor{currentfill}{rgb}{0.000000,0.000000,0.000000}%
\pgfsetfillcolor{currentfill}%
\pgfsetlinewidth{0.000000pt}%
\definecolor{currentstroke}{rgb}{0.000000,0.000000,0.000000}%
\pgfsetstrokecolor{currentstroke}%
\pgfsetstrokeopacity{0.000000}%
\pgfsetdash{}{0pt}%
\pgfpathmoveto{\pgfqpoint{6.664150in}{0.660000in}}%
\pgfpathlineto{\pgfqpoint{6.670350in}{0.660000in}}%
\pgfpathlineto{\pgfqpoint{6.670350in}{3.221476in}}%
\pgfpathlineto{\pgfqpoint{6.664150in}{3.221476in}}%
\pgfpathlineto{\pgfqpoint{6.664150in}{0.660000in}}%
\pgfpathclose%
\pgfusepath{fill}%
\end{pgfscope}%
\begin{pgfscope}%
\pgfpathrectangle{\pgfqpoint{1.250000in}{0.660000in}}{\pgfqpoint{7.750000in}{4.620000in}}%
\pgfusepath{clip}%
\pgfsetbuttcap%
\pgfsetmiterjoin%
\definecolor{currentfill}{rgb}{0.000000,0.000000,0.000000}%
\pgfsetfillcolor{currentfill}%
\pgfsetlinewidth{0.000000pt}%
\definecolor{currentstroke}{rgb}{0.000000,0.000000,0.000000}%
\pgfsetstrokecolor{currentstroke}%
\pgfsetstrokeopacity{0.000000}%
\pgfsetdash{}{0pt}%
\pgfpathmoveto{\pgfqpoint{6.671900in}{0.660000in}}%
\pgfpathlineto{\pgfqpoint{6.678100in}{0.660000in}}%
\pgfpathlineto{\pgfqpoint{6.678100in}{3.221348in}}%
\pgfpathlineto{\pgfqpoint{6.671900in}{3.221348in}}%
\pgfpathlineto{\pgfqpoint{6.671900in}{0.660000in}}%
\pgfpathclose%
\pgfusepath{fill}%
\end{pgfscope}%
\begin{pgfscope}%
\pgfpathrectangle{\pgfqpoint{1.250000in}{0.660000in}}{\pgfqpoint{7.750000in}{4.620000in}}%
\pgfusepath{clip}%
\pgfsetbuttcap%
\pgfsetmiterjoin%
\definecolor{currentfill}{rgb}{0.000000,0.000000,0.000000}%
\pgfsetfillcolor{currentfill}%
\pgfsetlinewidth{0.000000pt}%
\definecolor{currentstroke}{rgb}{0.000000,0.000000,0.000000}%
\pgfsetstrokecolor{currentstroke}%
\pgfsetstrokeopacity{0.000000}%
\pgfsetdash{}{0pt}%
\pgfpathmoveto{\pgfqpoint{6.679650in}{0.660000in}}%
\pgfpathlineto{\pgfqpoint{6.685850in}{0.660000in}}%
\pgfpathlineto{\pgfqpoint{6.685850in}{3.221220in}}%
\pgfpathlineto{\pgfqpoint{6.679650in}{3.221220in}}%
\pgfpathlineto{\pgfqpoint{6.679650in}{0.660000in}}%
\pgfpathclose%
\pgfusepath{fill}%
\end{pgfscope}%
\begin{pgfscope}%
\pgfpathrectangle{\pgfqpoint{1.250000in}{0.660000in}}{\pgfqpoint{7.750000in}{4.620000in}}%
\pgfusepath{clip}%
\pgfsetbuttcap%
\pgfsetmiterjoin%
\definecolor{currentfill}{rgb}{0.000000,0.000000,0.000000}%
\pgfsetfillcolor{currentfill}%
\pgfsetlinewidth{0.000000pt}%
\definecolor{currentstroke}{rgb}{0.000000,0.000000,0.000000}%
\pgfsetstrokecolor{currentstroke}%
\pgfsetstrokeopacity{0.000000}%
\pgfsetdash{}{0pt}%
\pgfpathmoveto{\pgfqpoint{6.687400in}{0.660000in}}%
\pgfpathlineto{\pgfqpoint{6.693600in}{0.660000in}}%
\pgfpathlineto{\pgfqpoint{6.693600in}{3.221091in}}%
\pgfpathlineto{\pgfqpoint{6.687400in}{3.221091in}}%
\pgfpathlineto{\pgfqpoint{6.687400in}{0.660000in}}%
\pgfpathclose%
\pgfusepath{fill}%
\end{pgfscope}%
\begin{pgfscope}%
\pgfpathrectangle{\pgfqpoint{1.250000in}{0.660000in}}{\pgfqpoint{7.750000in}{4.620000in}}%
\pgfusepath{clip}%
\pgfsetbuttcap%
\pgfsetmiterjoin%
\definecolor{currentfill}{rgb}{0.000000,0.000000,0.000000}%
\pgfsetfillcolor{currentfill}%
\pgfsetlinewidth{0.000000pt}%
\definecolor{currentstroke}{rgb}{0.000000,0.000000,0.000000}%
\pgfsetstrokecolor{currentstroke}%
\pgfsetstrokeopacity{0.000000}%
\pgfsetdash{}{0pt}%
\pgfpathmoveto{\pgfqpoint{6.695150in}{0.660000in}}%
\pgfpathlineto{\pgfqpoint{6.701350in}{0.660000in}}%
\pgfpathlineto{\pgfqpoint{6.701350in}{3.220963in}}%
\pgfpathlineto{\pgfqpoint{6.695150in}{3.220963in}}%
\pgfpathlineto{\pgfqpoint{6.695150in}{0.660000in}}%
\pgfpathclose%
\pgfusepath{fill}%
\end{pgfscope}%
\begin{pgfscope}%
\pgfpathrectangle{\pgfqpoint{1.250000in}{0.660000in}}{\pgfqpoint{7.750000in}{4.620000in}}%
\pgfusepath{clip}%
\pgfsetbuttcap%
\pgfsetmiterjoin%
\definecolor{currentfill}{rgb}{0.000000,0.000000,0.000000}%
\pgfsetfillcolor{currentfill}%
\pgfsetlinewidth{0.000000pt}%
\definecolor{currentstroke}{rgb}{0.000000,0.000000,0.000000}%
\pgfsetstrokecolor{currentstroke}%
\pgfsetstrokeopacity{0.000000}%
\pgfsetdash{}{0pt}%
\pgfpathmoveto{\pgfqpoint{6.702900in}{0.660000in}}%
\pgfpathlineto{\pgfqpoint{6.709100in}{0.660000in}}%
\pgfpathlineto{\pgfqpoint{6.709100in}{3.220834in}}%
\pgfpathlineto{\pgfqpoint{6.702900in}{3.220834in}}%
\pgfpathlineto{\pgfqpoint{6.702900in}{0.660000in}}%
\pgfpathclose%
\pgfusepath{fill}%
\end{pgfscope}%
\begin{pgfscope}%
\pgfpathrectangle{\pgfqpoint{1.250000in}{0.660000in}}{\pgfqpoint{7.750000in}{4.620000in}}%
\pgfusepath{clip}%
\pgfsetbuttcap%
\pgfsetmiterjoin%
\definecolor{currentfill}{rgb}{0.000000,0.000000,0.000000}%
\pgfsetfillcolor{currentfill}%
\pgfsetlinewidth{0.000000pt}%
\definecolor{currentstroke}{rgb}{0.000000,0.000000,0.000000}%
\pgfsetstrokecolor{currentstroke}%
\pgfsetstrokeopacity{0.000000}%
\pgfsetdash{}{0pt}%
\pgfpathmoveto{\pgfqpoint{6.710650in}{0.660000in}}%
\pgfpathlineto{\pgfqpoint{6.716850in}{0.660000in}}%
\pgfpathlineto{\pgfqpoint{6.716850in}{3.220708in}}%
\pgfpathlineto{\pgfqpoint{6.710650in}{3.220708in}}%
\pgfpathlineto{\pgfqpoint{6.710650in}{0.660000in}}%
\pgfpathclose%
\pgfusepath{fill}%
\end{pgfscope}%
\begin{pgfscope}%
\pgfpathrectangle{\pgfqpoint{1.250000in}{0.660000in}}{\pgfqpoint{7.750000in}{4.620000in}}%
\pgfusepath{clip}%
\pgfsetbuttcap%
\pgfsetmiterjoin%
\definecolor{currentfill}{rgb}{0.000000,0.000000,0.000000}%
\pgfsetfillcolor{currentfill}%
\pgfsetlinewidth{0.000000pt}%
\definecolor{currentstroke}{rgb}{0.000000,0.000000,0.000000}%
\pgfsetstrokecolor{currentstroke}%
\pgfsetstrokeopacity{0.000000}%
\pgfsetdash{}{0pt}%
\pgfpathmoveto{\pgfqpoint{6.718400in}{0.660000in}}%
\pgfpathlineto{\pgfqpoint{6.724600in}{0.660000in}}%
\pgfpathlineto{\pgfqpoint{6.724600in}{3.220584in}}%
\pgfpathlineto{\pgfqpoint{6.718400in}{3.220584in}}%
\pgfpathlineto{\pgfqpoint{6.718400in}{0.660000in}}%
\pgfpathclose%
\pgfusepath{fill}%
\end{pgfscope}%
\begin{pgfscope}%
\pgfpathrectangle{\pgfqpoint{1.250000in}{0.660000in}}{\pgfqpoint{7.750000in}{4.620000in}}%
\pgfusepath{clip}%
\pgfsetbuttcap%
\pgfsetmiterjoin%
\definecolor{currentfill}{rgb}{0.000000,0.000000,0.000000}%
\pgfsetfillcolor{currentfill}%
\pgfsetlinewidth{0.000000pt}%
\definecolor{currentstroke}{rgb}{0.000000,0.000000,0.000000}%
\pgfsetstrokecolor{currentstroke}%
\pgfsetstrokeopacity{0.000000}%
\pgfsetdash{}{0pt}%
\pgfpathmoveto{\pgfqpoint{6.726150in}{0.660000in}}%
\pgfpathlineto{\pgfqpoint{6.732350in}{0.660000in}}%
\pgfpathlineto{\pgfqpoint{6.732350in}{3.220459in}}%
\pgfpathlineto{\pgfqpoint{6.726150in}{3.220459in}}%
\pgfpathlineto{\pgfqpoint{6.726150in}{0.660000in}}%
\pgfpathclose%
\pgfusepath{fill}%
\end{pgfscope}%
\begin{pgfscope}%
\pgfpathrectangle{\pgfqpoint{1.250000in}{0.660000in}}{\pgfqpoint{7.750000in}{4.620000in}}%
\pgfusepath{clip}%
\pgfsetbuttcap%
\pgfsetmiterjoin%
\definecolor{currentfill}{rgb}{0.000000,0.000000,0.000000}%
\pgfsetfillcolor{currentfill}%
\pgfsetlinewidth{0.000000pt}%
\definecolor{currentstroke}{rgb}{0.000000,0.000000,0.000000}%
\pgfsetstrokecolor{currentstroke}%
\pgfsetstrokeopacity{0.000000}%
\pgfsetdash{}{0pt}%
\pgfpathmoveto{\pgfqpoint{6.733900in}{0.660000in}}%
\pgfpathlineto{\pgfqpoint{6.740100in}{0.660000in}}%
\pgfpathlineto{\pgfqpoint{6.740100in}{3.220334in}}%
\pgfpathlineto{\pgfqpoint{6.733900in}{3.220334in}}%
\pgfpathlineto{\pgfqpoint{6.733900in}{0.660000in}}%
\pgfpathclose%
\pgfusepath{fill}%
\end{pgfscope}%
\begin{pgfscope}%
\pgfpathrectangle{\pgfqpoint{1.250000in}{0.660000in}}{\pgfqpoint{7.750000in}{4.620000in}}%
\pgfusepath{clip}%
\pgfsetbuttcap%
\pgfsetmiterjoin%
\definecolor{currentfill}{rgb}{0.000000,0.000000,0.000000}%
\pgfsetfillcolor{currentfill}%
\pgfsetlinewidth{0.000000pt}%
\definecolor{currentstroke}{rgb}{0.000000,0.000000,0.000000}%
\pgfsetstrokecolor{currentstroke}%
\pgfsetstrokeopacity{0.000000}%
\pgfsetdash{}{0pt}%
\pgfpathmoveto{\pgfqpoint{6.741650in}{0.660000in}}%
\pgfpathlineto{\pgfqpoint{6.747850in}{0.660000in}}%
\pgfpathlineto{\pgfqpoint{6.747850in}{3.220210in}}%
\pgfpathlineto{\pgfqpoint{6.741650in}{3.220210in}}%
\pgfpathlineto{\pgfqpoint{6.741650in}{0.660000in}}%
\pgfpathclose%
\pgfusepath{fill}%
\end{pgfscope}%
\begin{pgfscope}%
\pgfpathrectangle{\pgfqpoint{1.250000in}{0.660000in}}{\pgfqpoint{7.750000in}{4.620000in}}%
\pgfusepath{clip}%
\pgfsetbuttcap%
\pgfsetmiterjoin%
\definecolor{currentfill}{rgb}{0.000000,0.000000,0.000000}%
\pgfsetfillcolor{currentfill}%
\pgfsetlinewidth{0.000000pt}%
\definecolor{currentstroke}{rgb}{0.000000,0.000000,0.000000}%
\pgfsetstrokecolor{currentstroke}%
\pgfsetstrokeopacity{0.000000}%
\pgfsetdash{}{0pt}%
\pgfpathmoveto{\pgfqpoint{6.749400in}{0.660000in}}%
\pgfpathlineto{\pgfqpoint{6.755600in}{0.660000in}}%
\pgfpathlineto{\pgfqpoint{6.755600in}{3.220085in}}%
\pgfpathlineto{\pgfqpoint{6.749400in}{3.220085in}}%
\pgfpathlineto{\pgfqpoint{6.749400in}{0.660000in}}%
\pgfpathclose%
\pgfusepath{fill}%
\end{pgfscope}%
\begin{pgfscope}%
\pgfpathrectangle{\pgfqpoint{1.250000in}{0.660000in}}{\pgfqpoint{7.750000in}{4.620000in}}%
\pgfusepath{clip}%
\pgfsetbuttcap%
\pgfsetmiterjoin%
\definecolor{currentfill}{rgb}{0.000000,0.000000,0.000000}%
\pgfsetfillcolor{currentfill}%
\pgfsetlinewidth{0.000000pt}%
\definecolor{currentstroke}{rgb}{0.000000,0.000000,0.000000}%
\pgfsetstrokecolor{currentstroke}%
\pgfsetstrokeopacity{0.000000}%
\pgfsetdash{}{0pt}%
\pgfpathmoveto{\pgfqpoint{6.757150in}{0.660000in}}%
\pgfpathlineto{\pgfqpoint{6.763350in}{0.660000in}}%
\pgfpathlineto{\pgfqpoint{6.763350in}{3.219960in}}%
\pgfpathlineto{\pgfqpoint{6.757150in}{3.219960in}}%
\pgfpathlineto{\pgfqpoint{6.757150in}{0.660000in}}%
\pgfpathclose%
\pgfusepath{fill}%
\end{pgfscope}%
\begin{pgfscope}%
\pgfpathrectangle{\pgfqpoint{1.250000in}{0.660000in}}{\pgfqpoint{7.750000in}{4.620000in}}%
\pgfusepath{clip}%
\pgfsetbuttcap%
\pgfsetmiterjoin%
\definecolor{currentfill}{rgb}{0.000000,0.000000,0.000000}%
\pgfsetfillcolor{currentfill}%
\pgfsetlinewidth{0.000000pt}%
\definecolor{currentstroke}{rgb}{0.000000,0.000000,0.000000}%
\pgfsetstrokecolor{currentstroke}%
\pgfsetstrokeopacity{0.000000}%
\pgfsetdash{}{0pt}%
\pgfpathmoveto{\pgfqpoint{6.764900in}{0.660000in}}%
\pgfpathlineto{\pgfqpoint{6.771100in}{0.660000in}}%
\pgfpathlineto{\pgfqpoint{6.771100in}{3.219836in}}%
\pgfpathlineto{\pgfqpoint{6.764900in}{3.219836in}}%
\pgfpathlineto{\pgfqpoint{6.764900in}{0.660000in}}%
\pgfpathclose%
\pgfusepath{fill}%
\end{pgfscope}%
\begin{pgfscope}%
\pgfpathrectangle{\pgfqpoint{1.250000in}{0.660000in}}{\pgfqpoint{7.750000in}{4.620000in}}%
\pgfusepath{clip}%
\pgfsetbuttcap%
\pgfsetmiterjoin%
\definecolor{currentfill}{rgb}{0.000000,0.000000,0.000000}%
\pgfsetfillcolor{currentfill}%
\pgfsetlinewidth{0.000000pt}%
\definecolor{currentstroke}{rgb}{0.000000,0.000000,0.000000}%
\pgfsetstrokecolor{currentstroke}%
\pgfsetstrokeopacity{0.000000}%
\pgfsetdash{}{0pt}%
\pgfpathmoveto{\pgfqpoint{6.772650in}{0.660000in}}%
\pgfpathlineto{\pgfqpoint{6.778850in}{0.660000in}}%
\pgfpathlineto{\pgfqpoint{6.778850in}{3.219711in}}%
\pgfpathlineto{\pgfqpoint{6.772650in}{3.219711in}}%
\pgfpathlineto{\pgfqpoint{6.772650in}{0.660000in}}%
\pgfpathclose%
\pgfusepath{fill}%
\end{pgfscope}%
\begin{pgfscope}%
\pgfpathrectangle{\pgfqpoint{1.250000in}{0.660000in}}{\pgfqpoint{7.750000in}{4.620000in}}%
\pgfusepath{clip}%
\pgfsetbuttcap%
\pgfsetmiterjoin%
\definecolor{currentfill}{rgb}{0.000000,0.000000,0.000000}%
\pgfsetfillcolor{currentfill}%
\pgfsetlinewidth{0.000000pt}%
\definecolor{currentstroke}{rgb}{0.000000,0.000000,0.000000}%
\pgfsetstrokecolor{currentstroke}%
\pgfsetstrokeopacity{0.000000}%
\pgfsetdash{}{0pt}%
\pgfpathmoveto{\pgfqpoint{6.780400in}{0.660000in}}%
\pgfpathlineto{\pgfqpoint{6.786600in}{0.660000in}}%
\pgfpathlineto{\pgfqpoint{6.786600in}{3.219587in}}%
\pgfpathlineto{\pgfqpoint{6.780400in}{3.219587in}}%
\pgfpathlineto{\pgfqpoint{6.780400in}{0.660000in}}%
\pgfpathclose%
\pgfusepath{fill}%
\end{pgfscope}%
\begin{pgfscope}%
\pgfpathrectangle{\pgfqpoint{1.250000in}{0.660000in}}{\pgfqpoint{7.750000in}{4.620000in}}%
\pgfusepath{clip}%
\pgfsetbuttcap%
\pgfsetmiterjoin%
\definecolor{currentfill}{rgb}{0.000000,0.000000,0.000000}%
\pgfsetfillcolor{currentfill}%
\pgfsetlinewidth{0.000000pt}%
\definecolor{currentstroke}{rgb}{0.000000,0.000000,0.000000}%
\pgfsetstrokecolor{currentstroke}%
\pgfsetstrokeopacity{0.000000}%
\pgfsetdash{}{0pt}%
\pgfpathmoveto{\pgfqpoint{6.788150in}{0.660000in}}%
\pgfpathlineto{\pgfqpoint{6.794350in}{0.660000in}}%
\pgfpathlineto{\pgfqpoint{6.794350in}{3.219462in}}%
\pgfpathlineto{\pgfqpoint{6.788150in}{3.219462in}}%
\pgfpathlineto{\pgfqpoint{6.788150in}{0.660000in}}%
\pgfpathclose%
\pgfusepath{fill}%
\end{pgfscope}%
\begin{pgfscope}%
\pgfpathrectangle{\pgfqpoint{1.250000in}{0.660000in}}{\pgfqpoint{7.750000in}{4.620000in}}%
\pgfusepath{clip}%
\pgfsetbuttcap%
\pgfsetmiterjoin%
\definecolor{currentfill}{rgb}{0.000000,0.000000,0.000000}%
\pgfsetfillcolor{currentfill}%
\pgfsetlinewidth{0.000000pt}%
\definecolor{currentstroke}{rgb}{0.000000,0.000000,0.000000}%
\pgfsetstrokecolor{currentstroke}%
\pgfsetstrokeopacity{0.000000}%
\pgfsetdash{}{0pt}%
\pgfpathmoveto{\pgfqpoint{6.795900in}{0.660000in}}%
\pgfpathlineto{\pgfqpoint{6.802100in}{0.660000in}}%
\pgfpathlineto{\pgfqpoint{6.802100in}{3.219337in}}%
\pgfpathlineto{\pgfqpoint{6.795900in}{3.219337in}}%
\pgfpathlineto{\pgfqpoint{6.795900in}{0.660000in}}%
\pgfpathclose%
\pgfusepath{fill}%
\end{pgfscope}%
\begin{pgfscope}%
\pgfpathrectangle{\pgfqpoint{1.250000in}{0.660000in}}{\pgfqpoint{7.750000in}{4.620000in}}%
\pgfusepath{clip}%
\pgfsetbuttcap%
\pgfsetmiterjoin%
\definecolor{currentfill}{rgb}{0.000000,0.000000,0.000000}%
\pgfsetfillcolor{currentfill}%
\pgfsetlinewidth{0.000000pt}%
\definecolor{currentstroke}{rgb}{0.000000,0.000000,0.000000}%
\pgfsetstrokecolor{currentstroke}%
\pgfsetstrokeopacity{0.000000}%
\pgfsetdash{}{0pt}%
\pgfpathmoveto{\pgfqpoint{6.803650in}{0.660000in}}%
\pgfpathlineto{\pgfqpoint{6.809850in}{0.660000in}}%
\pgfpathlineto{\pgfqpoint{6.809850in}{3.219213in}}%
\pgfpathlineto{\pgfqpoint{6.803650in}{3.219213in}}%
\pgfpathlineto{\pgfqpoint{6.803650in}{0.660000in}}%
\pgfpathclose%
\pgfusepath{fill}%
\end{pgfscope}%
\begin{pgfscope}%
\pgfpathrectangle{\pgfqpoint{1.250000in}{0.660000in}}{\pgfqpoint{7.750000in}{4.620000in}}%
\pgfusepath{clip}%
\pgfsetbuttcap%
\pgfsetmiterjoin%
\definecolor{currentfill}{rgb}{0.000000,0.000000,0.000000}%
\pgfsetfillcolor{currentfill}%
\pgfsetlinewidth{0.000000pt}%
\definecolor{currentstroke}{rgb}{0.000000,0.000000,0.000000}%
\pgfsetstrokecolor{currentstroke}%
\pgfsetstrokeopacity{0.000000}%
\pgfsetdash{}{0pt}%
\pgfpathmoveto{\pgfqpoint{6.811400in}{0.660000in}}%
\pgfpathlineto{\pgfqpoint{6.817600in}{0.660000in}}%
\pgfpathlineto{\pgfqpoint{6.817600in}{3.219088in}}%
\pgfpathlineto{\pgfqpoint{6.811400in}{3.219088in}}%
\pgfpathlineto{\pgfqpoint{6.811400in}{0.660000in}}%
\pgfpathclose%
\pgfusepath{fill}%
\end{pgfscope}%
\begin{pgfscope}%
\pgfpathrectangle{\pgfqpoint{1.250000in}{0.660000in}}{\pgfqpoint{7.750000in}{4.620000in}}%
\pgfusepath{clip}%
\pgfsetbuttcap%
\pgfsetmiterjoin%
\definecolor{currentfill}{rgb}{0.000000,0.000000,0.000000}%
\pgfsetfillcolor{currentfill}%
\pgfsetlinewidth{0.000000pt}%
\definecolor{currentstroke}{rgb}{0.000000,0.000000,0.000000}%
\pgfsetstrokecolor{currentstroke}%
\pgfsetstrokeopacity{0.000000}%
\pgfsetdash{}{0pt}%
\pgfpathmoveto{\pgfqpoint{6.819150in}{0.660000in}}%
\pgfpathlineto{\pgfqpoint{6.825350in}{0.660000in}}%
\pgfpathlineto{\pgfqpoint{6.825350in}{3.218964in}}%
\pgfpathlineto{\pgfqpoint{6.819150in}{3.218964in}}%
\pgfpathlineto{\pgfqpoint{6.819150in}{0.660000in}}%
\pgfpathclose%
\pgfusepath{fill}%
\end{pgfscope}%
\begin{pgfscope}%
\pgfpathrectangle{\pgfqpoint{1.250000in}{0.660000in}}{\pgfqpoint{7.750000in}{4.620000in}}%
\pgfusepath{clip}%
\pgfsetbuttcap%
\pgfsetmiterjoin%
\definecolor{currentfill}{rgb}{0.000000,0.000000,0.000000}%
\pgfsetfillcolor{currentfill}%
\pgfsetlinewidth{0.000000pt}%
\definecolor{currentstroke}{rgb}{0.000000,0.000000,0.000000}%
\pgfsetstrokecolor{currentstroke}%
\pgfsetstrokeopacity{0.000000}%
\pgfsetdash{}{0pt}%
\pgfpathmoveto{\pgfqpoint{6.826900in}{0.660000in}}%
\pgfpathlineto{\pgfqpoint{6.833100in}{0.660000in}}%
\pgfpathlineto{\pgfqpoint{6.833100in}{3.218839in}}%
\pgfpathlineto{\pgfqpoint{6.826900in}{3.218839in}}%
\pgfpathlineto{\pgfqpoint{6.826900in}{0.660000in}}%
\pgfpathclose%
\pgfusepath{fill}%
\end{pgfscope}%
\begin{pgfscope}%
\pgfpathrectangle{\pgfqpoint{1.250000in}{0.660000in}}{\pgfqpoint{7.750000in}{4.620000in}}%
\pgfusepath{clip}%
\pgfsetbuttcap%
\pgfsetmiterjoin%
\definecolor{currentfill}{rgb}{0.000000,0.000000,0.000000}%
\pgfsetfillcolor{currentfill}%
\pgfsetlinewidth{0.000000pt}%
\definecolor{currentstroke}{rgb}{0.000000,0.000000,0.000000}%
\pgfsetstrokecolor{currentstroke}%
\pgfsetstrokeopacity{0.000000}%
\pgfsetdash{}{0pt}%
\pgfpathmoveto{\pgfqpoint{6.834650in}{0.660000in}}%
\pgfpathlineto{\pgfqpoint{6.840850in}{0.660000in}}%
\pgfpathlineto{\pgfqpoint{6.840850in}{3.218718in}}%
\pgfpathlineto{\pgfqpoint{6.834650in}{3.218718in}}%
\pgfpathlineto{\pgfqpoint{6.834650in}{0.660000in}}%
\pgfpathclose%
\pgfusepath{fill}%
\end{pgfscope}%
\begin{pgfscope}%
\pgfpathrectangle{\pgfqpoint{1.250000in}{0.660000in}}{\pgfqpoint{7.750000in}{4.620000in}}%
\pgfusepath{clip}%
\pgfsetbuttcap%
\pgfsetmiterjoin%
\definecolor{currentfill}{rgb}{0.000000,0.000000,0.000000}%
\pgfsetfillcolor{currentfill}%
\pgfsetlinewidth{0.000000pt}%
\definecolor{currentstroke}{rgb}{0.000000,0.000000,0.000000}%
\pgfsetstrokecolor{currentstroke}%
\pgfsetstrokeopacity{0.000000}%
\pgfsetdash{}{0pt}%
\pgfpathmoveto{\pgfqpoint{6.842400in}{0.660000in}}%
\pgfpathlineto{\pgfqpoint{6.848600in}{0.660000in}}%
\pgfpathlineto{\pgfqpoint{6.848600in}{3.218597in}}%
\pgfpathlineto{\pgfqpoint{6.842400in}{3.218597in}}%
\pgfpathlineto{\pgfqpoint{6.842400in}{0.660000in}}%
\pgfpathclose%
\pgfusepath{fill}%
\end{pgfscope}%
\begin{pgfscope}%
\pgfpathrectangle{\pgfqpoint{1.250000in}{0.660000in}}{\pgfqpoint{7.750000in}{4.620000in}}%
\pgfusepath{clip}%
\pgfsetbuttcap%
\pgfsetmiterjoin%
\definecolor{currentfill}{rgb}{0.000000,0.000000,0.000000}%
\pgfsetfillcolor{currentfill}%
\pgfsetlinewidth{0.000000pt}%
\definecolor{currentstroke}{rgb}{0.000000,0.000000,0.000000}%
\pgfsetstrokecolor{currentstroke}%
\pgfsetstrokeopacity{0.000000}%
\pgfsetdash{}{0pt}%
\pgfpathmoveto{\pgfqpoint{6.850150in}{0.660000in}}%
\pgfpathlineto{\pgfqpoint{6.856350in}{0.660000in}}%
\pgfpathlineto{\pgfqpoint{6.856350in}{3.218476in}}%
\pgfpathlineto{\pgfqpoint{6.850150in}{3.218476in}}%
\pgfpathlineto{\pgfqpoint{6.850150in}{0.660000in}}%
\pgfpathclose%
\pgfusepath{fill}%
\end{pgfscope}%
\begin{pgfscope}%
\pgfpathrectangle{\pgfqpoint{1.250000in}{0.660000in}}{\pgfqpoint{7.750000in}{4.620000in}}%
\pgfusepath{clip}%
\pgfsetbuttcap%
\pgfsetmiterjoin%
\definecolor{currentfill}{rgb}{0.000000,0.000000,0.000000}%
\pgfsetfillcolor{currentfill}%
\pgfsetlinewidth{0.000000pt}%
\definecolor{currentstroke}{rgb}{0.000000,0.000000,0.000000}%
\pgfsetstrokecolor{currentstroke}%
\pgfsetstrokeopacity{0.000000}%
\pgfsetdash{}{0pt}%
\pgfpathmoveto{\pgfqpoint{6.857900in}{0.660000in}}%
\pgfpathlineto{\pgfqpoint{6.864100in}{0.660000in}}%
\pgfpathlineto{\pgfqpoint{6.864100in}{3.218355in}}%
\pgfpathlineto{\pgfqpoint{6.857900in}{3.218355in}}%
\pgfpathlineto{\pgfqpoint{6.857900in}{0.660000in}}%
\pgfpathclose%
\pgfusepath{fill}%
\end{pgfscope}%
\begin{pgfscope}%
\pgfpathrectangle{\pgfqpoint{1.250000in}{0.660000in}}{\pgfqpoint{7.750000in}{4.620000in}}%
\pgfusepath{clip}%
\pgfsetbuttcap%
\pgfsetmiterjoin%
\definecolor{currentfill}{rgb}{0.000000,0.000000,0.000000}%
\pgfsetfillcolor{currentfill}%
\pgfsetlinewidth{0.000000pt}%
\definecolor{currentstroke}{rgb}{0.000000,0.000000,0.000000}%
\pgfsetstrokecolor{currentstroke}%
\pgfsetstrokeopacity{0.000000}%
\pgfsetdash{}{0pt}%
\pgfpathmoveto{\pgfqpoint{6.865650in}{0.660000in}}%
\pgfpathlineto{\pgfqpoint{6.871850in}{0.660000in}}%
\pgfpathlineto{\pgfqpoint{6.871850in}{3.218234in}}%
\pgfpathlineto{\pgfqpoint{6.865650in}{3.218234in}}%
\pgfpathlineto{\pgfqpoint{6.865650in}{0.660000in}}%
\pgfpathclose%
\pgfusepath{fill}%
\end{pgfscope}%
\begin{pgfscope}%
\pgfpathrectangle{\pgfqpoint{1.250000in}{0.660000in}}{\pgfqpoint{7.750000in}{4.620000in}}%
\pgfusepath{clip}%
\pgfsetbuttcap%
\pgfsetmiterjoin%
\definecolor{currentfill}{rgb}{0.000000,0.000000,0.000000}%
\pgfsetfillcolor{currentfill}%
\pgfsetlinewidth{0.000000pt}%
\definecolor{currentstroke}{rgb}{0.000000,0.000000,0.000000}%
\pgfsetstrokecolor{currentstroke}%
\pgfsetstrokeopacity{0.000000}%
\pgfsetdash{}{0pt}%
\pgfpathmoveto{\pgfqpoint{6.873400in}{0.660000in}}%
\pgfpathlineto{\pgfqpoint{6.879600in}{0.660000in}}%
\pgfpathlineto{\pgfqpoint{6.879600in}{3.218113in}}%
\pgfpathlineto{\pgfqpoint{6.873400in}{3.218113in}}%
\pgfpathlineto{\pgfqpoint{6.873400in}{0.660000in}}%
\pgfpathclose%
\pgfusepath{fill}%
\end{pgfscope}%
\begin{pgfscope}%
\pgfpathrectangle{\pgfqpoint{1.250000in}{0.660000in}}{\pgfqpoint{7.750000in}{4.620000in}}%
\pgfusepath{clip}%
\pgfsetbuttcap%
\pgfsetmiterjoin%
\definecolor{currentfill}{rgb}{0.000000,0.000000,0.000000}%
\pgfsetfillcolor{currentfill}%
\pgfsetlinewidth{0.000000pt}%
\definecolor{currentstroke}{rgb}{0.000000,0.000000,0.000000}%
\pgfsetstrokecolor{currentstroke}%
\pgfsetstrokeopacity{0.000000}%
\pgfsetdash{}{0pt}%
\pgfpathmoveto{\pgfqpoint{6.881150in}{0.660000in}}%
\pgfpathlineto{\pgfqpoint{6.887350in}{0.660000in}}%
\pgfpathlineto{\pgfqpoint{6.887350in}{3.217992in}}%
\pgfpathlineto{\pgfqpoint{6.881150in}{3.217992in}}%
\pgfpathlineto{\pgfqpoint{6.881150in}{0.660000in}}%
\pgfpathclose%
\pgfusepath{fill}%
\end{pgfscope}%
\begin{pgfscope}%
\pgfpathrectangle{\pgfqpoint{1.250000in}{0.660000in}}{\pgfqpoint{7.750000in}{4.620000in}}%
\pgfusepath{clip}%
\pgfsetbuttcap%
\pgfsetmiterjoin%
\definecolor{currentfill}{rgb}{0.000000,0.000000,0.000000}%
\pgfsetfillcolor{currentfill}%
\pgfsetlinewidth{0.000000pt}%
\definecolor{currentstroke}{rgb}{0.000000,0.000000,0.000000}%
\pgfsetstrokecolor{currentstroke}%
\pgfsetstrokeopacity{0.000000}%
\pgfsetdash{}{0pt}%
\pgfpathmoveto{\pgfqpoint{6.888900in}{0.660000in}}%
\pgfpathlineto{\pgfqpoint{6.895100in}{0.660000in}}%
\pgfpathlineto{\pgfqpoint{6.895100in}{3.217872in}}%
\pgfpathlineto{\pgfqpoint{6.888900in}{3.217872in}}%
\pgfpathlineto{\pgfqpoint{6.888900in}{0.660000in}}%
\pgfpathclose%
\pgfusepath{fill}%
\end{pgfscope}%
\begin{pgfscope}%
\pgfpathrectangle{\pgfqpoint{1.250000in}{0.660000in}}{\pgfqpoint{7.750000in}{4.620000in}}%
\pgfusepath{clip}%
\pgfsetbuttcap%
\pgfsetmiterjoin%
\definecolor{currentfill}{rgb}{0.000000,0.000000,0.000000}%
\pgfsetfillcolor{currentfill}%
\pgfsetlinewidth{0.000000pt}%
\definecolor{currentstroke}{rgb}{0.000000,0.000000,0.000000}%
\pgfsetstrokecolor{currentstroke}%
\pgfsetstrokeopacity{0.000000}%
\pgfsetdash{}{0pt}%
\pgfpathmoveto{\pgfqpoint{6.896650in}{0.660000in}}%
\pgfpathlineto{\pgfqpoint{6.902850in}{0.660000in}}%
\pgfpathlineto{\pgfqpoint{6.902850in}{3.217751in}}%
\pgfpathlineto{\pgfqpoint{6.896650in}{3.217751in}}%
\pgfpathlineto{\pgfqpoint{6.896650in}{0.660000in}}%
\pgfpathclose%
\pgfusepath{fill}%
\end{pgfscope}%
\begin{pgfscope}%
\pgfpathrectangle{\pgfqpoint{1.250000in}{0.660000in}}{\pgfqpoint{7.750000in}{4.620000in}}%
\pgfusepath{clip}%
\pgfsetbuttcap%
\pgfsetmiterjoin%
\definecolor{currentfill}{rgb}{0.000000,0.000000,0.000000}%
\pgfsetfillcolor{currentfill}%
\pgfsetlinewidth{0.000000pt}%
\definecolor{currentstroke}{rgb}{0.000000,0.000000,0.000000}%
\pgfsetstrokecolor{currentstroke}%
\pgfsetstrokeopacity{0.000000}%
\pgfsetdash{}{0pt}%
\pgfpathmoveto{\pgfqpoint{6.904400in}{0.660000in}}%
\pgfpathlineto{\pgfqpoint{6.910600in}{0.660000in}}%
\pgfpathlineto{\pgfqpoint{6.910600in}{3.217630in}}%
\pgfpathlineto{\pgfqpoint{6.904400in}{3.217630in}}%
\pgfpathlineto{\pgfqpoint{6.904400in}{0.660000in}}%
\pgfpathclose%
\pgfusepath{fill}%
\end{pgfscope}%
\begin{pgfscope}%
\pgfpathrectangle{\pgfqpoint{1.250000in}{0.660000in}}{\pgfqpoint{7.750000in}{4.620000in}}%
\pgfusepath{clip}%
\pgfsetbuttcap%
\pgfsetmiterjoin%
\definecolor{currentfill}{rgb}{0.000000,0.000000,0.000000}%
\pgfsetfillcolor{currentfill}%
\pgfsetlinewidth{0.000000pt}%
\definecolor{currentstroke}{rgb}{0.000000,0.000000,0.000000}%
\pgfsetstrokecolor{currentstroke}%
\pgfsetstrokeopacity{0.000000}%
\pgfsetdash{}{0pt}%
\pgfpathmoveto{\pgfqpoint{6.912150in}{0.660000in}}%
\pgfpathlineto{\pgfqpoint{6.918350in}{0.660000in}}%
\pgfpathlineto{\pgfqpoint{6.918350in}{3.217509in}}%
\pgfpathlineto{\pgfqpoint{6.912150in}{3.217509in}}%
\pgfpathlineto{\pgfqpoint{6.912150in}{0.660000in}}%
\pgfpathclose%
\pgfusepath{fill}%
\end{pgfscope}%
\begin{pgfscope}%
\pgfpathrectangle{\pgfqpoint{1.250000in}{0.660000in}}{\pgfqpoint{7.750000in}{4.620000in}}%
\pgfusepath{clip}%
\pgfsetbuttcap%
\pgfsetmiterjoin%
\definecolor{currentfill}{rgb}{0.000000,0.000000,0.000000}%
\pgfsetfillcolor{currentfill}%
\pgfsetlinewidth{0.000000pt}%
\definecolor{currentstroke}{rgb}{0.000000,0.000000,0.000000}%
\pgfsetstrokecolor{currentstroke}%
\pgfsetstrokeopacity{0.000000}%
\pgfsetdash{}{0pt}%
\pgfpathmoveto{\pgfqpoint{6.919900in}{0.660000in}}%
\pgfpathlineto{\pgfqpoint{6.926100in}{0.660000in}}%
\pgfpathlineto{\pgfqpoint{6.926100in}{3.217388in}}%
\pgfpathlineto{\pgfqpoint{6.919900in}{3.217388in}}%
\pgfpathlineto{\pgfqpoint{6.919900in}{0.660000in}}%
\pgfpathclose%
\pgfusepath{fill}%
\end{pgfscope}%
\begin{pgfscope}%
\pgfpathrectangle{\pgfqpoint{1.250000in}{0.660000in}}{\pgfqpoint{7.750000in}{4.620000in}}%
\pgfusepath{clip}%
\pgfsetbuttcap%
\pgfsetmiterjoin%
\definecolor{currentfill}{rgb}{0.000000,0.000000,0.000000}%
\pgfsetfillcolor{currentfill}%
\pgfsetlinewidth{0.000000pt}%
\definecolor{currentstroke}{rgb}{0.000000,0.000000,0.000000}%
\pgfsetstrokecolor{currentstroke}%
\pgfsetstrokeopacity{0.000000}%
\pgfsetdash{}{0pt}%
\pgfpathmoveto{\pgfqpoint{6.927650in}{0.660000in}}%
\pgfpathlineto{\pgfqpoint{6.933850in}{0.660000in}}%
\pgfpathlineto{\pgfqpoint{6.933850in}{3.217267in}}%
\pgfpathlineto{\pgfqpoint{6.927650in}{3.217267in}}%
\pgfpathlineto{\pgfqpoint{6.927650in}{0.660000in}}%
\pgfpathclose%
\pgfusepath{fill}%
\end{pgfscope}%
\begin{pgfscope}%
\pgfpathrectangle{\pgfqpoint{1.250000in}{0.660000in}}{\pgfqpoint{7.750000in}{4.620000in}}%
\pgfusepath{clip}%
\pgfsetbuttcap%
\pgfsetmiterjoin%
\definecolor{currentfill}{rgb}{0.000000,0.000000,0.000000}%
\pgfsetfillcolor{currentfill}%
\pgfsetlinewidth{0.000000pt}%
\definecolor{currentstroke}{rgb}{0.000000,0.000000,0.000000}%
\pgfsetstrokecolor{currentstroke}%
\pgfsetstrokeopacity{0.000000}%
\pgfsetdash{}{0pt}%
\pgfpathmoveto{\pgfqpoint{6.935400in}{0.660000in}}%
\pgfpathlineto{\pgfqpoint{6.941600in}{0.660000in}}%
\pgfpathlineto{\pgfqpoint{6.941600in}{3.217146in}}%
\pgfpathlineto{\pgfqpoint{6.935400in}{3.217146in}}%
\pgfpathlineto{\pgfqpoint{6.935400in}{0.660000in}}%
\pgfpathclose%
\pgfusepath{fill}%
\end{pgfscope}%
\begin{pgfscope}%
\pgfpathrectangle{\pgfqpoint{1.250000in}{0.660000in}}{\pgfqpoint{7.750000in}{4.620000in}}%
\pgfusepath{clip}%
\pgfsetbuttcap%
\pgfsetmiterjoin%
\definecolor{currentfill}{rgb}{0.000000,0.000000,0.000000}%
\pgfsetfillcolor{currentfill}%
\pgfsetlinewidth{0.000000pt}%
\definecolor{currentstroke}{rgb}{0.000000,0.000000,0.000000}%
\pgfsetstrokecolor{currentstroke}%
\pgfsetstrokeopacity{0.000000}%
\pgfsetdash{}{0pt}%
\pgfpathmoveto{\pgfqpoint{6.943150in}{0.660000in}}%
\pgfpathlineto{\pgfqpoint{6.949350in}{0.660000in}}%
\pgfpathlineto{\pgfqpoint{6.949350in}{3.217026in}}%
\pgfpathlineto{\pgfqpoint{6.943150in}{3.217026in}}%
\pgfpathlineto{\pgfqpoint{6.943150in}{0.660000in}}%
\pgfpathclose%
\pgfusepath{fill}%
\end{pgfscope}%
\begin{pgfscope}%
\pgfpathrectangle{\pgfqpoint{1.250000in}{0.660000in}}{\pgfqpoint{7.750000in}{4.620000in}}%
\pgfusepath{clip}%
\pgfsetbuttcap%
\pgfsetmiterjoin%
\definecolor{currentfill}{rgb}{0.000000,0.000000,0.000000}%
\pgfsetfillcolor{currentfill}%
\pgfsetlinewidth{0.000000pt}%
\definecolor{currentstroke}{rgb}{0.000000,0.000000,0.000000}%
\pgfsetstrokecolor{currentstroke}%
\pgfsetstrokeopacity{0.000000}%
\pgfsetdash{}{0pt}%
\pgfpathmoveto{\pgfqpoint{6.950900in}{0.660000in}}%
\pgfpathlineto{\pgfqpoint{6.957100in}{0.660000in}}%
\pgfpathlineto{\pgfqpoint{6.957100in}{3.216905in}}%
\pgfpathlineto{\pgfqpoint{6.950900in}{3.216905in}}%
\pgfpathlineto{\pgfqpoint{6.950900in}{0.660000in}}%
\pgfpathclose%
\pgfusepath{fill}%
\end{pgfscope}%
\begin{pgfscope}%
\pgfpathrectangle{\pgfqpoint{1.250000in}{0.660000in}}{\pgfqpoint{7.750000in}{4.620000in}}%
\pgfusepath{clip}%
\pgfsetbuttcap%
\pgfsetmiterjoin%
\definecolor{currentfill}{rgb}{0.000000,0.000000,0.000000}%
\pgfsetfillcolor{currentfill}%
\pgfsetlinewidth{0.000000pt}%
\definecolor{currentstroke}{rgb}{0.000000,0.000000,0.000000}%
\pgfsetstrokecolor{currentstroke}%
\pgfsetstrokeopacity{0.000000}%
\pgfsetdash{}{0pt}%
\pgfpathmoveto{\pgfqpoint{6.958650in}{0.660000in}}%
\pgfpathlineto{\pgfqpoint{6.964850in}{0.660000in}}%
\pgfpathlineto{\pgfqpoint{6.964850in}{3.216786in}}%
\pgfpathlineto{\pgfqpoint{6.958650in}{3.216786in}}%
\pgfpathlineto{\pgfqpoint{6.958650in}{0.660000in}}%
\pgfpathclose%
\pgfusepath{fill}%
\end{pgfscope}%
\begin{pgfscope}%
\pgfpathrectangle{\pgfqpoint{1.250000in}{0.660000in}}{\pgfqpoint{7.750000in}{4.620000in}}%
\pgfusepath{clip}%
\pgfsetbuttcap%
\pgfsetmiterjoin%
\definecolor{currentfill}{rgb}{0.000000,0.000000,0.000000}%
\pgfsetfillcolor{currentfill}%
\pgfsetlinewidth{0.000000pt}%
\definecolor{currentstroke}{rgb}{0.000000,0.000000,0.000000}%
\pgfsetstrokecolor{currentstroke}%
\pgfsetstrokeopacity{0.000000}%
\pgfsetdash{}{0pt}%
\pgfpathmoveto{\pgfqpoint{6.966400in}{0.660000in}}%
\pgfpathlineto{\pgfqpoint{6.972600in}{0.660000in}}%
\pgfpathlineto{\pgfqpoint{6.972600in}{3.216668in}}%
\pgfpathlineto{\pgfqpoint{6.966400in}{3.216668in}}%
\pgfpathlineto{\pgfqpoint{6.966400in}{0.660000in}}%
\pgfpathclose%
\pgfusepath{fill}%
\end{pgfscope}%
\begin{pgfscope}%
\pgfpathrectangle{\pgfqpoint{1.250000in}{0.660000in}}{\pgfqpoint{7.750000in}{4.620000in}}%
\pgfusepath{clip}%
\pgfsetbuttcap%
\pgfsetmiterjoin%
\definecolor{currentfill}{rgb}{0.000000,0.000000,0.000000}%
\pgfsetfillcolor{currentfill}%
\pgfsetlinewidth{0.000000pt}%
\definecolor{currentstroke}{rgb}{0.000000,0.000000,0.000000}%
\pgfsetstrokecolor{currentstroke}%
\pgfsetstrokeopacity{0.000000}%
\pgfsetdash{}{0pt}%
\pgfpathmoveto{\pgfqpoint{6.974150in}{0.660000in}}%
\pgfpathlineto{\pgfqpoint{6.980350in}{0.660000in}}%
\pgfpathlineto{\pgfqpoint{6.980350in}{3.216551in}}%
\pgfpathlineto{\pgfqpoint{6.974150in}{3.216551in}}%
\pgfpathlineto{\pgfqpoint{6.974150in}{0.660000in}}%
\pgfpathclose%
\pgfusepath{fill}%
\end{pgfscope}%
\begin{pgfscope}%
\pgfpathrectangle{\pgfqpoint{1.250000in}{0.660000in}}{\pgfqpoint{7.750000in}{4.620000in}}%
\pgfusepath{clip}%
\pgfsetbuttcap%
\pgfsetmiterjoin%
\definecolor{currentfill}{rgb}{0.000000,0.000000,0.000000}%
\pgfsetfillcolor{currentfill}%
\pgfsetlinewidth{0.000000pt}%
\definecolor{currentstroke}{rgb}{0.000000,0.000000,0.000000}%
\pgfsetstrokecolor{currentstroke}%
\pgfsetstrokeopacity{0.000000}%
\pgfsetdash{}{0pt}%
\pgfpathmoveto{\pgfqpoint{6.981900in}{0.660000in}}%
\pgfpathlineto{\pgfqpoint{6.988100in}{0.660000in}}%
\pgfpathlineto{\pgfqpoint{6.988100in}{3.216434in}}%
\pgfpathlineto{\pgfqpoint{6.981900in}{3.216434in}}%
\pgfpathlineto{\pgfqpoint{6.981900in}{0.660000in}}%
\pgfpathclose%
\pgfusepath{fill}%
\end{pgfscope}%
\begin{pgfscope}%
\pgfpathrectangle{\pgfqpoint{1.250000in}{0.660000in}}{\pgfqpoint{7.750000in}{4.620000in}}%
\pgfusepath{clip}%
\pgfsetbuttcap%
\pgfsetmiterjoin%
\definecolor{currentfill}{rgb}{0.000000,0.000000,0.000000}%
\pgfsetfillcolor{currentfill}%
\pgfsetlinewidth{0.000000pt}%
\definecolor{currentstroke}{rgb}{0.000000,0.000000,0.000000}%
\pgfsetstrokecolor{currentstroke}%
\pgfsetstrokeopacity{0.000000}%
\pgfsetdash{}{0pt}%
\pgfpathmoveto{\pgfqpoint{6.989650in}{0.660000in}}%
\pgfpathlineto{\pgfqpoint{6.995850in}{0.660000in}}%
\pgfpathlineto{\pgfqpoint{6.995850in}{3.216317in}}%
\pgfpathlineto{\pgfqpoint{6.989650in}{3.216317in}}%
\pgfpathlineto{\pgfqpoint{6.989650in}{0.660000in}}%
\pgfpathclose%
\pgfusepath{fill}%
\end{pgfscope}%
\begin{pgfscope}%
\pgfpathrectangle{\pgfqpoint{1.250000in}{0.660000in}}{\pgfqpoint{7.750000in}{4.620000in}}%
\pgfusepath{clip}%
\pgfsetbuttcap%
\pgfsetmiterjoin%
\definecolor{currentfill}{rgb}{0.000000,0.000000,0.000000}%
\pgfsetfillcolor{currentfill}%
\pgfsetlinewidth{0.000000pt}%
\definecolor{currentstroke}{rgb}{0.000000,0.000000,0.000000}%
\pgfsetstrokecolor{currentstroke}%
\pgfsetstrokeopacity{0.000000}%
\pgfsetdash{}{0pt}%
\pgfpathmoveto{\pgfqpoint{6.997400in}{0.660000in}}%
\pgfpathlineto{\pgfqpoint{7.003600in}{0.660000in}}%
\pgfpathlineto{\pgfqpoint{7.003600in}{3.216200in}}%
\pgfpathlineto{\pgfqpoint{6.997400in}{3.216200in}}%
\pgfpathlineto{\pgfqpoint{6.997400in}{0.660000in}}%
\pgfpathclose%
\pgfusepath{fill}%
\end{pgfscope}%
\begin{pgfscope}%
\pgfpathrectangle{\pgfqpoint{1.250000in}{0.660000in}}{\pgfqpoint{7.750000in}{4.620000in}}%
\pgfusepath{clip}%
\pgfsetbuttcap%
\pgfsetmiterjoin%
\definecolor{currentfill}{rgb}{0.000000,0.000000,0.000000}%
\pgfsetfillcolor{currentfill}%
\pgfsetlinewidth{0.000000pt}%
\definecolor{currentstroke}{rgb}{0.000000,0.000000,0.000000}%
\pgfsetstrokecolor{currentstroke}%
\pgfsetstrokeopacity{0.000000}%
\pgfsetdash{}{0pt}%
\pgfpathmoveto{\pgfqpoint{7.005150in}{0.660000in}}%
\pgfpathlineto{\pgfqpoint{7.011350in}{0.660000in}}%
\pgfpathlineto{\pgfqpoint{7.011350in}{3.216083in}}%
\pgfpathlineto{\pgfqpoint{7.005150in}{3.216083in}}%
\pgfpathlineto{\pgfqpoint{7.005150in}{0.660000in}}%
\pgfpathclose%
\pgfusepath{fill}%
\end{pgfscope}%
\begin{pgfscope}%
\pgfpathrectangle{\pgfqpoint{1.250000in}{0.660000in}}{\pgfqpoint{7.750000in}{4.620000in}}%
\pgfusepath{clip}%
\pgfsetbuttcap%
\pgfsetmiterjoin%
\definecolor{currentfill}{rgb}{0.000000,0.000000,0.000000}%
\pgfsetfillcolor{currentfill}%
\pgfsetlinewidth{0.000000pt}%
\definecolor{currentstroke}{rgb}{0.000000,0.000000,0.000000}%
\pgfsetstrokecolor{currentstroke}%
\pgfsetstrokeopacity{0.000000}%
\pgfsetdash{}{0pt}%
\pgfpathmoveto{\pgfqpoint{7.012900in}{0.660000in}}%
\pgfpathlineto{\pgfqpoint{7.019100in}{0.660000in}}%
\pgfpathlineto{\pgfqpoint{7.019100in}{3.215965in}}%
\pgfpathlineto{\pgfqpoint{7.012900in}{3.215965in}}%
\pgfpathlineto{\pgfqpoint{7.012900in}{0.660000in}}%
\pgfpathclose%
\pgfusepath{fill}%
\end{pgfscope}%
\begin{pgfscope}%
\pgfpathrectangle{\pgfqpoint{1.250000in}{0.660000in}}{\pgfqpoint{7.750000in}{4.620000in}}%
\pgfusepath{clip}%
\pgfsetbuttcap%
\pgfsetmiterjoin%
\definecolor{currentfill}{rgb}{0.000000,0.000000,0.000000}%
\pgfsetfillcolor{currentfill}%
\pgfsetlinewidth{0.000000pt}%
\definecolor{currentstroke}{rgb}{0.000000,0.000000,0.000000}%
\pgfsetstrokecolor{currentstroke}%
\pgfsetstrokeopacity{0.000000}%
\pgfsetdash{}{0pt}%
\pgfpathmoveto{\pgfqpoint{7.020650in}{0.660000in}}%
\pgfpathlineto{\pgfqpoint{7.026850in}{0.660000in}}%
\pgfpathlineto{\pgfqpoint{7.026850in}{3.215848in}}%
\pgfpathlineto{\pgfqpoint{7.020650in}{3.215848in}}%
\pgfpathlineto{\pgfqpoint{7.020650in}{0.660000in}}%
\pgfpathclose%
\pgfusepath{fill}%
\end{pgfscope}%
\begin{pgfscope}%
\pgfpathrectangle{\pgfqpoint{1.250000in}{0.660000in}}{\pgfqpoint{7.750000in}{4.620000in}}%
\pgfusepath{clip}%
\pgfsetbuttcap%
\pgfsetmiterjoin%
\definecolor{currentfill}{rgb}{0.000000,0.000000,0.000000}%
\pgfsetfillcolor{currentfill}%
\pgfsetlinewidth{0.000000pt}%
\definecolor{currentstroke}{rgb}{0.000000,0.000000,0.000000}%
\pgfsetstrokecolor{currentstroke}%
\pgfsetstrokeopacity{0.000000}%
\pgfsetdash{}{0pt}%
\pgfpathmoveto{\pgfqpoint{7.028400in}{0.660000in}}%
\pgfpathlineto{\pgfqpoint{7.034600in}{0.660000in}}%
\pgfpathlineto{\pgfqpoint{7.034600in}{3.215731in}}%
\pgfpathlineto{\pgfqpoint{7.028400in}{3.215731in}}%
\pgfpathlineto{\pgfqpoint{7.028400in}{0.660000in}}%
\pgfpathclose%
\pgfusepath{fill}%
\end{pgfscope}%
\begin{pgfscope}%
\pgfpathrectangle{\pgfqpoint{1.250000in}{0.660000in}}{\pgfqpoint{7.750000in}{4.620000in}}%
\pgfusepath{clip}%
\pgfsetbuttcap%
\pgfsetmiterjoin%
\definecolor{currentfill}{rgb}{0.000000,0.000000,0.000000}%
\pgfsetfillcolor{currentfill}%
\pgfsetlinewidth{0.000000pt}%
\definecolor{currentstroke}{rgb}{0.000000,0.000000,0.000000}%
\pgfsetstrokecolor{currentstroke}%
\pgfsetstrokeopacity{0.000000}%
\pgfsetdash{}{0pt}%
\pgfpathmoveto{\pgfqpoint{7.036150in}{0.660000in}}%
\pgfpathlineto{\pgfqpoint{7.042350in}{0.660000in}}%
\pgfpathlineto{\pgfqpoint{7.042350in}{3.215614in}}%
\pgfpathlineto{\pgfqpoint{7.036150in}{3.215614in}}%
\pgfpathlineto{\pgfqpoint{7.036150in}{0.660000in}}%
\pgfpathclose%
\pgfusepath{fill}%
\end{pgfscope}%
\begin{pgfscope}%
\pgfpathrectangle{\pgfqpoint{1.250000in}{0.660000in}}{\pgfqpoint{7.750000in}{4.620000in}}%
\pgfusepath{clip}%
\pgfsetbuttcap%
\pgfsetmiterjoin%
\definecolor{currentfill}{rgb}{0.000000,0.000000,0.000000}%
\pgfsetfillcolor{currentfill}%
\pgfsetlinewidth{0.000000pt}%
\definecolor{currentstroke}{rgb}{0.000000,0.000000,0.000000}%
\pgfsetstrokecolor{currentstroke}%
\pgfsetstrokeopacity{0.000000}%
\pgfsetdash{}{0pt}%
\pgfpathmoveto{\pgfqpoint{7.043900in}{0.660000in}}%
\pgfpathlineto{\pgfqpoint{7.050100in}{0.660000in}}%
\pgfpathlineto{\pgfqpoint{7.050100in}{3.215497in}}%
\pgfpathlineto{\pgfqpoint{7.043900in}{3.215497in}}%
\pgfpathlineto{\pgfqpoint{7.043900in}{0.660000in}}%
\pgfpathclose%
\pgfusepath{fill}%
\end{pgfscope}%
\begin{pgfscope}%
\pgfpathrectangle{\pgfqpoint{1.250000in}{0.660000in}}{\pgfqpoint{7.750000in}{4.620000in}}%
\pgfusepath{clip}%
\pgfsetbuttcap%
\pgfsetmiterjoin%
\definecolor{currentfill}{rgb}{0.000000,0.000000,0.000000}%
\pgfsetfillcolor{currentfill}%
\pgfsetlinewidth{0.000000pt}%
\definecolor{currentstroke}{rgb}{0.000000,0.000000,0.000000}%
\pgfsetstrokecolor{currentstroke}%
\pgfsetstrokeopacity{0.000000}%
\pgfsetdash{}{0pt}%
\pgfpathmoveto{\pgfqpoint{7.051650in}{0.660000in}}%
\pgfpathlineto{\pgfqpoint{7.057850in}{0.660000in}}%
\pgfpathlineto{\pgfqpoint{7.057850in}{3.215380in}}%
\pgfpathlineto{\pgfqpoint{7.051650in}{3.215380in}}%
\pgfpathlineto{\pgfqpoint{7.051650in}{0.660000in}}%
\pgfpathclose%
\pgfusepath{fill}%
\end{pgfscope}%
\begin{pgfscope}%
\pgfpathrectangle{\pgfqpoint{1.250000in}{0.660000in}}{\pgfqpoint{7.750000in}{4.620000in}}%
\pgfusepath{clip}%
\pgfsetbuttcap%
\pgfsetmiterjoin%
\definecolor{currentfill}{rgb}{0.000000,0.000000,0.000000}%
\pgfsetfillcolor{currentfill}%
\pgfsetlinewidth{0.000000pt}%
\definecolor{currentstroke}{rgb}{0.000000,0.000000,0.000000}%
\pgfsetstrokecolor{currentstroke}%
\pgfsetstrokeopacity{0.000000}%
\pgfsetdash{}{0pt}%
\pgfpathmoveto{\pgfqpoint{7.059400in}{0.660000in}}%
\pgfpathlineto{\pgfqpoint{7.065600in}{0.660000in}}%
\pgfpathlineto{\pgfqpoint{7.065600in}{3.215263in}}%
\pgfpathlineto{\pgfqpoint{7.059400in}{3.215263in}}%
\pgfpathlineto{\pgfqpoint{7.059400in}{0.660000in}}%
\pgfpathclose%
\pgfusepath{fill}%
\end{pgfscope}%
\begin{pgfscope}%
\pgfpathrectangle{\pgfqpoint{1.250000in}{0.660000in}}{\pgfqpoint{7.750000in}{4.620000in}}%
\pgfusepath{clip}%
\pgfsetbuttcap%
\pgfsetmiterjoin%
\definecolor{currentfill}{rgb}{0.000000,0.000000,0.000000}%
\pgfsetfillcolor{currentfill}%
\pgfsetlinewidth{0.000000pt}%
\definecolor{currentstroke}{rgb}{0.000000,0.000000,0.000000}%
\pgfsetstrokecolor{currentstroke}%
\pgfsetstrokeopacity{0.000000}%
\pgfsetdash{}{0pt}%
\pgfpathmoveto{\pgfqpoint{7.067150in}{0.660000in}}%
\pgfpathlineto{\pgfqpoint{7.073350in}{0.660000in}}%
\pgfpathlineto{\pgfqpoint{7.073350in}{3.215145in}}%
\pgfpathlineto{\pgfqpoint{7.067150in}{3.215145in}}%
\pgfpathlineto{\pgfqpoint{7.067150in}{0.660000in}}%
\pgfpathclose%
\pgfusepath{fill}%
\end{pgfscope}%
\begin{pgfscope}%
\pgfpathrectangle{\pgfqpoint{1.250000in}{0.660000in}}{\pgfqpoint{7.750000in}{4.620000in}}%
\pgfusepath{clip}%
\pgfsetbuttcap%
\pgfsetmiterjoin%
\definecolor{currentfill}{rgb}{0.000000,0.000000,0.000000}%
\pgfsetfillcolor{currentfill}%
\pgfsetlinewidth{0.000000pt}%
\definecolor{currentstroke}{rgb}{0.000000,0.000000,0.000000}%
\pgfsetstrokecolor{currentstroke}%
\pgfsetstrokeopacity{0.000000}%
\pgfsetdash{}{0pt}%
\pgfpathmoveto{\pgfqpoint{7.074900in}{0.660000in}}%
\pgfpathlineto{\pgfqpoint{7.081100in}{0.660000in}}%
\pgfpathlineto{\pgfqpoint{7.081100in}{3.215028in}}%
\pgfpathlineto{\pgfqpoint{7.074900in}{3.215028in}}%
\pgfpathlineto{\pgfqpoint{7.074900in}{0.660000in}}%
\pgfpathclose%
\pgfusepath{fill}%
\end{pgfscope}%
\begin{pgfscope}%
\pgfpathrectangle{\pgfqpoint{1.250000in}{0.660000in}}{\pgfqpoint{7.750000in}{4.620000in}}%
\pgfusepath{clip}%
\pgfsetbuttcap%
\pgfsetmiterjoin%
\definecolor{currentfill}{rgb}{0.000000,0.000000,0.000000}%
\pgfsetfillcolor{currentfill}%
\pgfsetlinewidth{0.000000pt}%
\definecolor{currentstroke}{rgb}{0.000000,0.000000,0.000000}%
\pgfsetstrokecolor{currentstroke}%
\pgfsetstrokeopacity{0.000000}%
\pgfsetdash{}{0pt}%
\pgfpathmoveto{\pgfqpoint{7.082650in}{0.660000in}}%
\pgfpathlineto{\pgfqpoint{7.088850in}{0.660000in}}%
\pgfpathlineto{\pgfqpoint{7.088850in}{3.214911in}}%
\pgfpathlineto{\pgfqpoint{7.082650in}{3.214911in}}%
\pgfpathlineto{\pgfqpoint{7.082650in}{0.660000in}}%
\pgfpathclose%
\pgfusepath{fill}%
\end{pgfscope}%
\begin{pgfscope}%
\pgfpathrectangle{\pgfqpoint{1.250000in}{0.660000in}}{\pgfqpoint{7.750000in}{4.620000in}}%
\pgfusepath{clip}%
\pgfsetbuttcap%
\pgfsetmiterjoin%
\definecolor{currentfill}{rgb}{0.000000,0.000000,0.000000}%
\pgfsetfillcolor{currentfill}%
\pgfsetlinewidth{0.000000pt}%
\definecolor{currentstroke}{rgb}{0.000000,0.000000,0.000000}%
\pgfsetstrokecolor{currentstroke}%
\pgfsetstrokeopacity{0.000000}%
\pgfsetdash{}{0pt}%
\pgfpathmoveto{\pgfqpoint{7.090400in}{0.660000in}}%
\pgfpathlineto{\pgfqpoint{7.096600in}{0.660000in}}%
\pgfpathlineto{\pgfqpoint{7.096600in}{3.214796in}}%
\pgfpathlineto{\pgfqpoint{7.090400in}{3.214796in}}%
\pgfpathlineto{\pgfqpoint{7.090400in}{0.660000in}}%
\pgfpathclose%
\pgfusepath{fill}%
\end{pgfscope}%
\begin{pgfscope}%
\pgfpathrectangle{\pgfqpoint{1.250000in}{0.660000in}}{\pgfqpoint{7.750000in}{4.620000in}}%
\pgfusepath{clip}%
\pgfsetbuttcap%
\pgfsetmiterjoin%
\definecolor{currentfill}{rgb}{0.000000,0.000000,0.000000}%
\pgfsetfillcolor{currentfill}%
\pgfsetlinewidth{0.000000pt}%
\definecolor{currentstroke}{rgb}{0.000000,0.000000,0.000000}%
\pgfsetstrokecolor{currentstroke}%
\pgfsetstrokeopacity{0.000000}%
\pgfsetdash{}{0pt}%
\pgfpathmoveto{\pgfqpoint{7.098150in}{0.660000in}}%
\pgfpathlineto{\pgfqpoint{7.104350in}{0.660000in}}%
\pgfpathlineto{\pgfqpoint{7.104350in}{3.214682in}}%
\pgfpathlineto{\pgfqpoint{7.098150in}{3.214682in}}%
\pgfpathlineto{\pgfqpoint{7.098150in}{0.660000in}}%
\pgfpathclose%
\pgfusepath{fill}%
\end{pgfscope}%
\begin{pgfscope}%
\pgfpathrectangle{\pgfqpoint{1.250000in}{0.660000in}}{\pgfqpoint{7.750000in}{4.620000in}}%
\pgfusepath{clip}%
\pgfsetbuttcap%
\pgfsetmiterjoin%
\definecolor{currentfill}{rgb}{0.000000,0.000000,0.000000}%
\pgfsetfillcolor{currentfill}%
\pgfsetlinewidth{0.000000pt}%
\definecolor{currentstroke}{rgb}{0.000000,0.000000,0.000000}%
\pgfsetstrokecolor{currentstroke}%
\pgfsetstrokeopacity{0.000000}%
\pgfsetdash{}{0pt}%
\pgfpathmoveto{\pgfqpoint{7.105900in}{0.660000in}}%
\pgfpathlineto{\pgfqpoint{7.112100in}{0.660000in}}%
\pgfpathlineto{\pgfqpoint{7.112100in}{3.214569in}}%
\pgfpathlineto{\pgfqpoint{7.105900in}{3.214569in}}%
\pgfpathlineto{\pgfqpoint{7.105900in}{0.660000in}}%
\pgfpathclose%
\pgfusepath{fill}%
\end{pgfscope}%
\begin{pgfscope}%
\pgfpathrectangle{\pgfqpoint{1.250000in}{0.660000in}}{\pgfqpoint{7.750000in}{4.620000in}}%
\pgfusepath{clip}%
\pgfsetbuttcap%
\pgfsetmiterjoin%
\definecolor{currentfill}{rgb}{0.000000,0.000000,0.000000}%
\pgfsetfillcolor{currentfill}%
\pgfsetlinewidth{0.000000pt}%
\definecolor{currentstroke}{rgb}{0.000000,0.000000,0.000000}%
\pgfsetstrokecolor{currentstroke}%
\pgfsetstrokeopacity{0.000000}%
\pgfsetdash{}{0pt}%
\pgfpathmoveto{\pgfqpoint{7.113650in}{0.660000in}}%
\pgfpathlineto{\pgfqpoint{7.119850in}{0.660000in}}%
\pgfpathlineto{\pgfqpoint{7.119850in}{3.214455in}}%
\pgfpathlineto{\pgfqpoint{7.113650in}{3.214455in}}%
\pgfpathlineto{\pgfqpoint{7.113650in}{0.660000in}}%
\pgfpathclose%
\pgfusepath{fill}%
\end{pgfscope}%
\begin{pgfscope}%
\pgfpathrectangle{\pgfqpoint{1.250000in}{0.660000in}}{\pgfqpoint{7.750000in}{4.620000in}}%
\pgfusepath{clip}%
\pgfsetbuttcap%
\pgfsetmiterjoin%
\definecolor{currentfill}{rgb}{0.000000,0.000000,0.000000}%
\pgfsetfillcolor{currentfill}%
\pgfsetlinewidth{0.000000pt}%
\definecolor{currentstroke}{rgb}{0.000000,0.000000,0.000000}%
\pgfsetstrokecolor{currentstroke}%
\pgfsetstrokeopacity{0.000000}%
\pgfsetdash{}{0pt}%
\pgfpathmoveto{\pgfqpoint{7.121400in}{0.660000in}}%
\pgfpathlineto{\pgfqpoint{7.127600in}{0.660000in}}%
\pgfpathlineto{\pgfqpoint{7.127600in}{3.214342in}}%
\pgfpathlineto{\pgfqpoint{7.121400in}{3.214342in}}%
\pgfpathlineto{\pgfqpoint{7.121400in}{0.660000in}}%
\pgfpathclose%
\pgfusepath{fill}%
\end{pgfscope}%
\begin{pgfscope}%
\pgfpathrectangle{\pgfqpoint{1.250000in}{0.660000in}}{\pgfqpoint{7.750000in}{4.620000in}}%
\pgfusepath{clip}%
\pgfsetbuttcap%
\pgfsetmiterjoin%
\definecolor{currentfill}{rgb}{0.000000,0.000000,0.000000}%
\pgfsetfillcolor{currentfill}%
\pgfsetlinewidth{0.000000pt}%
\definecolor{currentstroke}{rgb}{0.000000,0.000000,0.000000}%
\pgfsetstrokecolor{currentstroke}%
\pgfsetstrokeopacity{0.000000}%
\pgfsetdash{}{0pt}%
\pgfpathmoveto{\pgfqpoint{7.129150in}{0.660000in}}%
\pgfpathlineto{\pgfqpoint{7.135350in}{0.660000in}}%
\pgfpathlineto{\pgfqpoint{7.135350in}{3.214228in}}%
\pgfpathlineto{\pgfqpoint{7.129150in}{3.214228in}}%
\pgfpathlineto{\pgfqpoint{7.129150in}{0.660000in}}%
\pgfpathclose%
\pgfusepath{fill}%
\end{pgfscope}%
\begin{pgfscope}%
\pgfpathrectangle{\pgfqpoint{1.250000in}{0.660000in}}{\pgfqpoint{7.750000in}{4.620000in}}%
\pgfusepath{clip}%
\pgfsetbuttcap%
\pgfsetmiterjoin%
\definecolor{currentfill}{rgb}{0.000000,0.000000,0.000000}%
\pgfsetfillcolor{currentfill}%
\pgfsetlinewidth{0.000000pt}%
\definecolor{currentstroke}{rgb}{0.000000,0.000000,0.000000}%
\pgfsetstrokecolor{currentstroke}%
\pgfsetstrokeopacity{0.000000}%
\pgfsetdash{}{0pt}%
\pgfpathmoveto{\pgfqpoint{7.136900in}{0.660000in}}%
\pgfpathlineto{\pgfqpoint{7.143100in}{0.660000in}}%
\pgfpathlineto{\pgfqpoint{7.143100in}{3.214114in}}%
\pgfpathlineto{\pgfqpoint{7.136900in}{3.214114in}}%
\pgfpathlineto{\pgfqpoint{7.136900in}{0.660000in}}%
\pgfpathclose%
\pgfusepath{fill}%
\end{pgfscope}%
\begin{pgfscope}%
\pgfpathrectangle{\pgfqpoint{1.250000in}{0.660000in}}{\pgfqpoint{7.750000in}{4.620000in}}%
\pgfusepath{clip}%
\pgfsetbuttcap%
\pgfsetmiterjoin%
\definecolor{currentfill}{rgb}{0.000000,0.000000,0.000000}%
\pgfsetfillcolor{currentfill}%
\pgfsetlinewidth{0.000000pt}%
\definecolor{currentstroke}{rgb}{0.000000,0.000000,0.000000}%
\pgfsetstrokecolor{currentstroke}%
\pgfsetstrokeopacity{0.000000}%
\pgfsetdash{}{0pt}%
\pgfpathmoveto{\pgfqpoint{7.144650in}{0.660000in}}%
\pgfpathlineto{\pgfqpoint{7.150850in}{0.660000in}}%
\pgfpathlineto{\pgfqpoint{7.150850in}{3.214001in}}%
\pgfpathlineto{\pgfqpoint{7.144650in}{3.214001in}}%
\pgfpathlineto{\pgfqpoint{7.144650in}{0.660000in}}%
\pgfpathclose%
\pgfusepath{fill}%
\end{pgfscope}%
\begin{pgfscope}%
\pgfpathrectangle{\pgfqpoint{1.250000in}{0.660000in}}{\pgfqpoint{7.750000in}{4.620000in}}%
\pgfusepath{clip}%
\pgfsetbuttcap%
\pgfsetmiterjoin%
\definecolor{currentfill}{rgb}{0.000000,0.000000,0.000000}%
\pgfsetfillcolor{currentfill}%
\pgfsetlinewidth{0.000000pt}%
\definecolor{currentstroke}{rgb}{0.000000,0.000000,0.000000}%
\pgfsetstrokecolor{currentstroke}%
\pgfsetstrokeopacity{0.000000}%
\pgfsetdash{}{0pt}%
\pgfpathmoveto{\pgfqpoint{7.152400in}{0.660000in}}%
\pgfpathlineto{\pgfqpoint{7.158600in}{0.660000in}}%
\pgfpathlineto{\pgfqpoint{7.158600in}{3.213887in}}%
\pgfpathlineto{\pgfqpoint{7.152400in}{3.213887in}}%
\pgfpathlineto{\pgfqpoint{7.152400in}{0.660000in}}%
\pgfpathclose%
\pgfusepath{fill}%
\end{pgfscope}%
\begin{pgfscope}%
\pgfpathrectangle{\pgfqpoint{1.250000in}{0.660000in}}{\pgfqpoint{7.750000in}{4.620000in}}%
\pgfusepath{clip}%
\pgfsetbuttcap%
\pgfsetmiterjoin%
\definecolor{currentfill}{rgb}{0.000000,0.000000,0.000000}%
\pgfsetfillcolor{currentfill}%
\pgfsetlinewidth{0.000000pt}%
\definecolor{currentstroke}{rgb}{0.000000,0.000000,0.000000}%
\pgfsetstrokecolor{currentstroke}%
\pgfsetstrokeopacity{0.000000}%
\pgfsetdash{}{0pt}%
\pgfpathmoveto{\pgfqpoint{7.160150in}{0.660000in}}%
\pgfpathlineto{\pgfqpoint{7.166350in}{0.660000in}}%
\pgfpathlineto{\pgfqpoint{7.166350in}{3.213774in}}%
\pgfpathlineto{\pgfqpoint{7.160150in}{3.213774in}}%
\pgfpathlineto{\pgfqpoint{7.160150in}{0.660000in}}%
\pgfpathclose%
\pgfusepath{fill}%
\end{pgfscope}%
\begin{pgfscope}%
\pgfpathrectangle{\pgfqpoint{1.250000in}{0.660000in}}{\pgfqpoint{7.750000in}{4.620000in}}%
\pgfusepath{clip}%
\pgfsetbuttcap%
\pgfsetmiterjoin%
\definecolor{currentfill}{rgb}{0.000000,0.000000,0.000000}%
\pgfsetfillcolor{currentfill}%
\pgfsetlinewidth{0.000000pt}%
\definecolor{currentstroke}{rgb}{0.000000,0.000000,0.000000}%
\pgfsetstrokecolor{currentstroke}%
\pgfsetstrokeopacity{0.000000}%
\pgfsetdash{}{0pt}%
\pgfpathmoveto{\pgfqpoint{7.167900in}{0.660000in}}%
\pgfpathlineto{\pgfqpoint{7.174100in}{0.660000in}}%
\pgfpathlineto{\pgfqpoint{7.174100in}{3.213660in}}%
\pgfpathlineto{\pgfqpoint{7.167900in}{3.213660in}}%
\pgfpathlineto{\pgfqpoint{7.167900in}{0.660000in}}%
\pgfpathclose%
\pgfusepath{fill}%
\end{pgfscope}%
\begin{pgfscope}%
\pgfpathrectangle{\pgfqpoint{1.250000in}{0.660000in}}{\pgfqpoint{7.750000in}{4.620000in}}%
\pgfusepath{clip}%
\pgfsetbuttcap%
\pgfsetmiterjoin%
\definecolor{currentfill}{rgb}{0.000000,0.000000,0.000000}%
\pgfsetfillcolor{currentfill}%
\pgfsetlinewidth{0.000000pt}%
\definecolor{currentstroke}{rgb}{0.000000,0.000000,0.000000}%
\pgfsetstrokecolor{currentstroke}%
\pgfsetstrokeopacity{0.000000}%
\pgfsetdash{}{0pt}%
\pgfpathmoveto{\pgfqpoint{7.175650in}{0.660000in}}%
\pgfpathlineto{\pgfqpoint{7.181850in}{0.660000in}}%
\pgfpathlineto{\pgfqpoint{7.181850in}{3.213547in}}%
\pgfpathlineto{\pgfqpoint{7.175650in}{3.213547in}}%
\pgfpathlineto{\pgfqpoint{7.175650in}{0.660000in}}%
\pgfpathclose%
\pgfusepath{fill}%
\end{pgfscope}%
\begin{pgfscope}%
\pgfpathrectangle{\pgfqpoint{1.250000in}{0.660000in}}{\pgfqpoint{7.750000in}{4.620000in}}%
\pgfusepath{clip}%
\pgfsetbuttcap%
\pgfsetmiterjoin%
\definecolor{currentfill}{rgb}{0.000000,0.000000,0.000000}%
\pgfsetfillcolor{currentfill}%
\pgfsetlinewidth{0.000000pt}%
\definecolor{currentstroke}{rgb}{0.000000,0.000000,0.000000}%
\pgfsetstrokecolor{currentstroke}%
\pgfsetstrokeopacity{0.000000}%
\pgfsetdash{}{0pt}%
\pgfpathmoveto{\pgfqpoint{7.183400in}{0.660000in}}%
\pgfpathlineto{\pgfqpoint{7.189600in}{0.660000in}}%
\pgfpathlineto{\pgfqpoint{7.189600in}{3.213433in}}%
\pgfpathlineto{\pgfqpoint{7.183400in}{3.213433in}}%
\pgfpathlineto{\pgfqpoint{7.183400in}{0.660000in}}%
\pgfpathclose%
\pgfusepath{fill}%
\end{pgfscope}%
\begin{pgfscope}%
\pgfpathrectangle{\pgfqpoint{1.250000in}{0.660000in}}{\pgfqpoint{7.750000in}{4.620000in}}%
\pgfusepath{clip}%
\pgfsetbuttcap%
\pgfsetmiterjoin%
\definecolor{currentfill}{rgb}{0.000000,0.000000,0.000000}%
\pgfsetfillcolor{currentfill}%
\pgfsetlinewidth{0.000000pt}%
\definecolor{currentstroke}{rgb}{0.000000,0.000000,0.000000}%
\pgfsetstrokecolor{currentstroke}%
\pgfsetstrokeopacity{0.000000}%
\pgfsetdash{}{0pt}%
\pgfpathmoveto{\pgfqpoint{7.191150in}{0.660000in}}%
\pgfpathlineto{\pgfqpoint{7.197350in}{0.660000in}}%
\pgfpathlineto{\pgfqpoint{7.197350in}{3.213320in}}%
\pgfpathlineto{\pgfqpoint{7.191150in}{3.213320in}}%
\pgfpathlineto{\pgfqpoint{7.191150in}{0.660000in}}%
\pgfpathclose%
\pgfusepath{fill}%
\end{pgfscope}%
\begin{pgfscope}%
\pgfpathrectangle{\pgfqpoint{1.250000in}{0.660000in}}{\pgfqpoint{7.750000in}{4.620000in}}%
\pgfusepath{clip}%
\pgfsetbuttcap%
\pgfsetmiterjoin%
\definecolor{currentfill}{rgb}{0.000000,0.000000,0.000000}%
\pgfsetfillcolor{currentfill}%
\pgfsetlinewidth{0.000000pt}%
\definecolor{currentstroke}{rgb}{0.000000,0.000000,0.000000}%
\pgfsetstrokecolor{currentstroke}%
\pgfsetstrokeopacity{0.000000}%
\pgfsetdash{}{0pt}%
\pgfpathmoveto{\pgfqpoint{7.198900in}{0.660000in}}%
\pgfpathlineto{\pgfqpoint{7.205100in}{0.660000in}}%
\pgfpathlineto{\pgfqpoint{7.205100in}{3.213206in}}%
\pgfpathlineto{\pgfqpoint{7.198900in}{3.213206in}}%
\pgfpathlineto{\pgfqpoint{7.198900in}{0.660000in}}%
\pgfpathclose%
\pgfusepath{fill}%
\end{pgfscope}%
\begin{pgfscope}%
\pgfpathrectangle{\pgfqpoint{1.250000in}{0.660000in}}{\pgfqpoint{7.750000in}{4.620000in}}%
\pgfusepath{clip}%
\pgfsetbuttcap%
\pgfsetmiterjoin%
\definecolor{currentfill}{rgb}{0.000000,0.000000,0.000000}%
\pgfsetfillcolor{currentfill}%
\pgfsetlinewidth{0.000000pt}%
\definecolor{currentstroke}{rgb}{0.000000,0.000000,0.000000}%
\pgfsetstrokecolor{currentstroke}%
\pgfsetstrokeopacity{0.000000}%
\pgfsetdash{}{0pt}%
\pgfpathmoveto{\pgfqpoint{7.206650in}{0.660000in}}%
\pgfpathlineto{\pgfqpoint{7.212850in}{0.660000in}}%
\pgfpathlineto{\pgfqpoint{7.212850in}{3.213093in}}%
\pgfpathlineto{\pgfqpoint{7.206650in}{3.213093in}}%
\pgfpathlineto{\pgfqpoint{7.206650in}{0.660000in}}%
\pgfpathclose%
\pgfusepath{fill}%
\end{pgfscope}%
\begin{pgfscope}%
\pgfpathrectangle{\pgfqpoint{1.250000in}{0.660000in}}{\pgfqpoint{7.750000in}{4.620000in}}%
\pgfusepath{clip}%
\pgfsetbuttcap%
\pgfsetmiterjoin%
\definecolor{currentfill}{rgb}{0.000000,0.000000,0.000000}%
\pgfsetfillcolor{currentfill}%
\pgfsetlinewidth{0.000000pt}%
\definecolor{currentstroke}{rgb}{0.000000,0.000000,0.000000}%
\pgfsetstrokecolor{currentstroke}%
\pgfsetstrokeopacity{0.000000}%
\pgfsetdash{}{0pt}%
\pgfpathmoveto{\pgfqpoint{7.214400in}{0.660000in}}%
\pgfpathlineto{\pgfqpoint{7.220600in}{0.660000in}}%
\pgfpathlineto{\pgfqpoint{7.220600in}{3.212979in}}%
\pgfpathlineto{\pgfqpoint{7.214400in}{3.212979in}}%
\pgfpathlineto{\pgfqpoint{7.214400in}{0.660000in}}%
\pgfpathclose%
\pgfusepath{fill}%
\end{pgfscope}%
\begin{pgfscope}%
\pgfpathrectangle{\pgfqpoint{1.250000in}{0.660000in}}{\pgfqpoint{7.750000in}{4.620000in}}%
\pgfusepath{clip}%
\pgfsetbuttcap%
\pgfsetmiterjoin%
\definecolor{currentfill}{rgb}{0.000000,0.000000,0.000000}%
\pgfsetfillcolor{currentfill}%
\pgfsetlinewidth{0.000000pt}%
\definecolor{currentstroke}{rgb}{0.000000,0.000000,0.000000}%
\pgfsetstrokecolor{currentstroke}%
\pgfsetstrokeopacity{0.000000}%
\pgfsetdash{}{0pt}%
\pgfpathmoveto{\pgfqpoint{7.222150in}{0.660000in}}%
\pgfpathlineto{\pgfqpoint{7.228350in}{0.660000in}}%
\pgfpathlineto{\pgfqpoint{7.228350in}{3.212866in}}%
\pgfpathlineto{\pgfqpoint{7.222150in}{3.212866in}}%
\pgfpathlineto{\pgfqpoint{7.222150in}{0.660000in}}%
\pgfpathclose%
\pgfusepath{fill}%
\end{pgfscope}%
\begin{pgfscope}%
\pgfpathrectangle{\pgfqpoint{1.250000in}{0.660000in}}{\pgfqpoint{7.750000in}{4.620000in}}%
\pgfusepath{clip}%
\pgfsetbuttcap%
\pgfsetmiterjoin%
\definecolor{currentfill}{rgb}{0.000000,0.000000,0.000000}%
\pgfsetfillcolor{currentfill}%
\pgfsetlinewidth{0.000000pt}%
\definecolor{currentstroke}{rgb}{0.000000,0.000000,0.000000}%
\pgfsetstrokecolor{currentstroke}%
\pgfsetstrokeopacity{0.000000}%
\pgfsetdash{}{0pt}%
\pgfpathmoveto{\pgfqpoint{7.229900in}{0.660000in}}%
\pgfpathlineto{\pgfqpoint{7.236100in}{0.660000in}}%
\pgfpathlineto{\pgfqpoint{7.236100in}{3.212755in}}%
\pgfpathlineto{\pgfqpoint{7.229900in}{3.212755in}}%
\pgfpathlineto{\pgfqpoint{7.229900in}{0.660000in}}%
\pgfpathclose%
\pgfusepath{fill}%
\end{pgfscope}%
\begin{pgfscope}%
\pgfpathrectangle{\pgfqpoint{1.250000in}{0.660000in}}{\pgfqpoint{7.750000in}{4.620000in}}%
\pgfusepath{clip}%
\pgfsetbuttcap%
\pgfsetmiterjoin%
\definecolor{currentfill}{rgb}{0.000000,0.000000,0.000000}%
\pgfsetfillcolor{currentfill}%
\pgfsetlinewidth{0.000000pt}%
\definecolor{currentstroke}{rgb}{0.000000,0.000000,0.000000}%
\pgfsetstrokecolor{currentstroke}%
\pgfsetstrokeopacity{0.000000}%
\pgfsetdash{}{0pt}%
\pgfpathmoveto{\pgfqpoint{7.237650in}{0.660000in}}%
\pgfpathlineto{\pgfqpoint{7.243850in}{0.660000in}}%
\pgfpathlineto{\pgfqpoint{7.243850in}{3.212645in}}%
\pgfpathlineto{\pgfqpoint{7.237650in}{3.212645in}}%
\pgfpathlineto{\pgfqpoint{7.237650in}{0.660000in}}%
\pgfpathclose%
\pgfusepath{fill}%
\end{pgfscope}%
\begin{pgfscope}%
\pgfpathrectangle{\pgfqpoint{1.250000in}{0.660000in}}{\pgfqpoint{7.750000in}{4.620000in}}%
\pgfusepath{clip}%
\pgfsetbuttcap%
\pgfsetmiterjoin%
\definecolor{currentfill}{rgb}{0.000000,0.000000,0.000000}%
\pgfsetfillcolor{currentfill}%
\pgfsetlinewidth{0.000000pt}%
\definecolor{currentstroke}{rgb}{0.000000,0.000000,0.000000}%
\pgfsetstrokecolor{currentstroke}%
\pgfsetstrokeopacity{0.000000}%
\pgfsetdash{}{0pt}%
\pgfpathmoveto{\pgfqpoint{7.245400in}{0.660000in}}%
\pgfpathlineto{\pgfqpoint{7.251600in}{0.660000in}}%
\pgfpathlineto{\pgfqpoint{7.251600in}{3.212535in}}%
\pgfpathlineto{\pgfqpoint{7.245400in}{3.212535in}}%
\pgfpathlineto{\pgfqpoint{7.245400in}{0.660000in}}%
\pgfpathclose%
\pgfusepath{fill}%
\end{pgfscope}%
\begin{pgfscope}%
\pgfpathrectangle{\pgfqpoint{1.250000in}{0.660000in}}{\pgfqpoint{7.750000in}{4.620000in}}%
\pgfusepath{clip}%
\pgfsetbuttcap%
\pgfsetmiterjoin%
\definecolor{currentfill}{rgb}{0.000000,0.000000,0.000000}%
\pgfsetfillcolor{currentfill}%
\pgfsetlinewidth{0.000000pt}%
\definecolor{currentstroke}{rgb}{0.000000,0.000000,0.000000}%
\pgfsetstrokecolor{currentstroke}%
\pgfsetstrokeopacity{0.000000}%
\pgfsetdash{}{0pt}%
\pgfpathmoveto{\pgfqpoint{7.253150in}{0.660000in}}%
\pgfpathlineto{\pgfqpoint{7.259350in}{0.660000in}}%
\pgfpathlineto{\pgfqpoint{7.259350in}{3.212425in}}%
\pgfpathlineto{\pgfqpoint{7.253150in}{3.212425in}}%
\pgfpathlineto{\pgfqpoint{7.253150in}{0.660000in}}%
\pgfpathclose%
\pgfusepath{fill}%
\end{pgfscope}%
\begin{pgfscope}%
\pgfpathrectangle{\pgfqpoint{1.250000in}{0.660000in}}{\pgfqpoint{7.750000in}{4.620000in}}%
\pgfusepath{clip}%
\pgfsetbuttcap%
\pgfsetmiterjoin%
\definecolor{currentfill}{rgb}{0.000000,0.000000,0.000000}%
\pgfsetfillcolor{currentfill}%
\pgfsetlinewidth{0.000000pt}%
\definecolor{currentstroke}{rgb}{0.000000,0.000000,0.000000}%
\pgfsetstrokecolor{currentstroke}%
\pgfsetstrokeopacity{0.000000}%
\pgfsetdash{}{0pt}%
\pgfpathmoveto{\pgfqpoint{7.260900in}{0.660000in}}%
\pgfpathlineto{\pgfqpoint{7.267100in}{0.660000in}}%
\pgfpathlineto{\pgfqpoint{7.267100in}{3.212315in}}%
\pgfpathlineto{\pgfqpoint{7.260900in}{3.212315in}}%
\pgfpathlineto{\pgfqpoint{7.260900in}{0.660000in}}%
\pgfpathclose%
\pgfusepath{fill}%
\end{pgfscope}%
\begin{pgfscope}%
\pgfpathrectangle{\pgfqpoint{1.250000in}{0.660000in}}{\pgfqpoint{7.750000in}{4.620000in}}%
\pgfusepath{clip}%
\pgfsetbuttcap%
\pgfsetmiterjoin%
\definecolor{currentfill}{rgb}{0.000000,0.000000,0.000000}%
\pgfsetfillcolor{currentfill}%
\pgfsetlinewidth{0.000000pt}%
\definecolor{currentstroke}{rgb}{0.000000,0.000000,0.000000}%
\pgfsetstrokecolor{currentstroke}%
\pgfsetstrokeopacity{0.000000}%
\pgfsetdash{}{0pt}%
\pgfpathmoveto{\pgfqpoint{7.268650in}{0.660000in}}%
\pgfpathlineto{\pgfqpoint{7.274850in}{0.660000in}}%
\pgfpathlineto{\pgfqpoint{7.274850in}{3.212205in}}%
\pgfpathlineto{\pgfqpoint{7.268650in}{3.212205in}}%
\pgfpathlineto{\pgfqpoint{7.268650in}{0.660000in}}%
\pgfpathclose%
\pgfusepath{fill}%
\end{pgfscope}%
\begin{pgfscope}%
\pgfpathrectangle{\pgfqpoint{1.250000in}{0.660000in}}{\pgfqpoint{7.750000in}{4.620000in}}%
\pgfusepath{clip}%
\pgfsetbuttcap%
\pgfsetmiterjoin%
\definecolor{currentfill}{rgb}{0.000000,0.000000,0.000000}%
\pgfsetfillcolor{currentfill}%
\pgfsetlinewidth{0.000000pt}%
\definecolor{currentstroke}{rgb}{0.000000,0.000000,0.000000}%
\pgfsetstrokecolor{currentstroke}%
\pgfsetstrokeopacity{0.000000}%
\pgfsetdash{}{0pt}%
\pgfpathmoveto{\pgfqpoint{7.276400in}{0.660000in}}%
\pgfpathlineto{\pgfqpoint{7.282600in}{0.660000in}}%
\pgfpathlineto{\pgfqpoint{7.282600in}{3.212096in}}%
\pgfpathlineto{\pgfqpoint{7.276400in}{3.212096in}}%
\pgfpathlineto{\pgfqpoint{7.276400in}{0.660000in}}%
\pgfpathclose%
\pgfusepath{fill}%
\end{pgfscope}%
\begin{pgfscope}%
\pgfpathrectangle{\pgfqpoint{1.250000in}{0.660000in}}{\pgfqpoint{7.750000in}{4.620000in}}%
\pgfusepath{clip}%
\pgfsetbuttcap%
\pgfsetmiterjoin%
\definecolor{currentfill}{rgb}{0.000000,0.000000,0.000000}%
\pgfsetfillcolor{currentfill}%
\pgfsetlinewidth{0.000000pt}%
\definecolor{currentstroke}{rgb}{0.000000,0.000000,0.000000}%
\pgfsetstrokecolor{currentstroke}%
\pgfsetstrokeopacity{0.000000}%
\pgfsetdash{}{0pt}%
\pgfpathmoveto{\pgfqpoint{7.284150in}{0.660000in}}%
\pgfpathlineto{\pgfqpoint{7.290350in}{0.660000in}}%
\pgfpathlineto{\pgfqpoint{7.290350in}{3.211986in}}%
\pgfpathlineto{\pgfqpoint{7.284150in}{3.211986in}}%
\pgfpathlineto{\pgfqpoint{7.284150in}{0.660000in}}%
\pgfpathclose%
\pgfusepath{fill}%
\end{pgfscope}%
\begin{pgfscope}%
\pgfpathrectangle{\pgfqpoint{1.250000in}{0.660000in}}{\pgfqpoint{7.750000in}{4.620000in}}%
\pgfusepath{clip}%
\pgfsetbuttcap%
\pgfsetmiterjoin%
\definecolor{currentfill}{rgb}{0.000000,0.000000,0.000000}%
\pgfsetfillcolor{currentfill}%
\pgfsetlinewidth{0.000000pt}%
\definecolor{currentstroke}{rgb}{0.000000,0.000000,0.000000}%
\pgfsetstrokecolor{currentstroke}%
\pgfsetstrokeopacity{0.000000}%
\pgfsetdash{}{0pt}%
\pgfpathmoveto{\pgfqpoint{7.291900in}{0.660000in}}%
\pgfpathlineto{\pgfqpoint{7.298100in}{0.660000in}}%
\pgfpathlineto{\pgfqpoint{7.298100in}{3.211876in}}%
\pgfpathlineto{\pgfqpoint{7.291900in}{3.211876in}}%
\pgfpathlineto{\pgfqpoint{7.291900in}{0.660000in}}%
\pgfpathclose%
\pgfusepath{fill}%
\end{pgfscope}%
\begin{pgfscope}%
\pgfpathrectangle{\pgfqpoint{1.250000in}{0.660000in}}{\pgfqpoint{7.750000in}{4.620000in}}%
\pgfusepath{clip}%
\pgfsetbuttcap%
\pgfsetmiterjoin%
\definecolor{currentfill}{rgb}{0.000000,0.000000,0.000000}%
\pgfsetfillcolor{currentfill}%
\pgfsetlinewidth{0.000000pt}%
\definecolor{currentstroke}{rgb}{0.000000,0.000000,0.000000}%
\pgfsetstrokecolor{currentstroke}%
\pgfsetstrokeopacity{0.000000}%
\pgfsetdash{}{0pt}%
\pgfpathmoveto{\pgfqpoint{7.299650in}{0.660000in}}%
\pgfpathlineto{\pgfqpoint{7.305850in}{0.660000in}}%
\pgfpathlineto{\pgfqpoint{7.305850in}{3.211766in}}%
\pgfpathlineto{\pgfqpoint{7.299650in}{3.211766in}}%
\pgfpathlineto{\pgfqpoint{7.299650in}{0.660000in}}%
\pgfpathclose%
\pgfusepath{fill}%
\end{pgfscope}%
\begin{pgfscope}%
\pgfpathrectangle{\pgfqpoint{1.250000in}{0.660000in}}{\pgfqpoint{7.750000in}{4.620000in}}%
\pgfusepath{clip}%
\pgfsetbuttcap%
\pgfsetmiterjoin%
\definecolor{currentfill}{rgb}{0.000000,0.000000,0.000000}%
\pgfsetfillcolor{currentfill}%
\pgfsetlinewidth{0.000000pt}%
\definecolor{currentstroke}{rgb}{0.000000,0.000000,0.000000}%
\pgfsetstrokecolor{currentstroke}%
\pgfsetstrokeopacity{0.000000}%
\pgfsetdash{}{0pt}%
\pgfpathmoveto{\pgfqpoint{7.307400in}{0.660000in}}%
\pgfpathlineto{\pgfqpoint{7.313600in}{0.660000in}}%
\pgfpathlineto{\pgfqpoint{7.313600in}{3.211656in}}%
\pgfpathlineto{\pgfqpoint{7.307400in}{3.211656in}}%
\pgfpathlineto{\pgfqpoint{7.307400in}{0.660000in}}%
\pgfpathclose%
\pgfusepath{fill}%
\end{pgfscope}%
\begin{pgfscope}%
\pgfpathrectangle{\pgfqpoint{1.250000in}{0.660000in}}{\pgfqpoint{7.750000in}{4.620000in}}%
\pgfusepath{clip}%
\pgfsetbuttcap%
\pgfsetmiterjoin%
\definecolor{currentfill}{rgb}{0.000000,0.000000,0.000000}%
\pgfsetfillcolor{currentfill}%
\pgfsetlinewidth{0.000000pt}%
\definecolor{currentstroke}{rgb}{0.000000,0.000000,0.000000}%
\pgfsetstrokecolor{currentstroke}%
\pgfsetstrokeopacity{0.000000}%
\pgfsetdash{}{0pt}%
\pgfpathmoveto{\pgfqpoint{7.315150in}{0.660000in}}%
\pgfpathlineto{\pgfqpoint{7.321350in}{0.660000in}}%
\pgfpathlineto{\pgfqpoint{7.321350in}{3.211546in}}%
\pgfpathlineto{\pgfqpoint{7.315150in}{3.211546in}}%
\pgfpathlineto{\pgfqpoint{7.315150in}{0.660000in}}%
\pgfpathclose%
\pgfusepath{fill}%
\end{pgfscope}%
\begin{pgfscope}%
\pgfpathrectangle{\pgfqpoint{1.250000in}{0.660000in}}{\pgfqpoint{7.750000in}{4.620000in}}%
\pgfusepath{clip}%
\pgfsetbuttcap%
\pgfsetmiterjoin%
\definecolor{currentfill}{rgb}{0.000000,0.000000,0.000000}%
\pgfsetfillcolor{currentfill}%
\pgfsetlinewidth{0.000000pt}%
\definecolor{currentstroke}{rgb}{0.000000,0.000000,0.000000}%
\pgfsetstrokecolor{currentstroke}%
\pgfsetstrokeopacity{0.000000}%
\pgfsetdash{}{0pt}%
\pgfpathmoveto{\pgfqpoint{7.322900in}{0.660000in}}%
\pgfpathlineto{\pgfqpoint{7.329100in}{0.660000in}}%
\pgfpathlineto{\pgfqpoint{7.329100in}{3.211436in}}%
\pgfpathlineto{\pgfqpoint{7.322900in}{3.211436in}}%
\pgfpathlineto{\pgfqpoint{7.322900in}{0.660000in}}%
\pgfpathclose%
\pgfusepath{fill}%
\end{pgfscope}%
\begin{pgfscope}%
\pgfpathrectangle{\pgfqpoint{1.250000in}{0.660000in}}{\pgfqpoint{7.750000in}{4.620000in}}%
\pgfusepath{clip}%
\pgfsetbuttcap%
\pgfsetmiterjoin%
\definecolor{currentfill}{rgb}{0.000000,0.000000,0.000000}%
\pgfsetfillcolor{currentfill}%
\pgfsetlinewidth{0.000000pt}%
\definecolor{currentstroke}{rgb}{0.000000,0.000000,0.000000}%
\pgfsetstrokecolor{currentstroke}%
\pgfsetstrokeopacity{0.000000}%
\pgfsetdash{}{0pt}%
\pgfpathmoveto{\pgfqpoint{7.330650in}{0.660000in}}%
\pgfpathlineto{\pgfqpoint{7.336850in}{0.660000in}}%
\pgfpathlineto{\pgfqpoint{7.336850in}{3.211326in}}%
\pgfpathlineto{\pgfqpoint{7.330650in}{3.211326in}}%
\pgfpathlineto{\pgfqpoint{7.330650in}{0.660000in}}%
\pgfpathclose%
\pgfusepath{fill}%
\end{pgfscope}%
\begin{pgfscope}%
\pgfpathrectangle{\pgfqpoint{1.250000in}{0.660000in}}{\pgfqpoint{7.750000in}{4.620000in}}%
\pgfusepath{clip}%
\pgfsetbuttcap%
\pgfsetmiterjoin%
\definecolor{currentfill}{rgb}{0.000000,0.000000,0.000000}%
\pgfsetfillcolor{currentfill}%
\pgfsetlinewidth{0.000000pt}%
\definecolor{currentstroke}{rgb}{0.000000,0.000000,0.000000}%
\pgfsetstrokecolor{currentstroke}%
\pgfsetstrokeopacity{0.000000}%
\pgfsetdash{}{0pt}%
\pgfpathmoveto{\pgfqpoint{7.338400in}{0.660000in}}%
\pgfpathlineto{\pgfqpoint{7.344600in}{0.660000in}}%
\pgfpathlineto{\pgfqpoint{7.344600in}{3.211216in}}%
\pgfpathlineto{\pgfqpoint{7.338400in}{3.211216in}}%
\pgfpathlineto{\pgfqpoint{7.338400in}{0.660000in}}%
\pgfpathclose%
\pgfusepath{fill}%
\end{pgfscope}%
\begin{pgfscope}%
\pgfpathrectangle{\pgfqpoint{1.250000in}{0.660000in}}{\pgfqpoint{7.750000in}{4.620000in}}%
\pgfusepath{clip}%
\pgfsetbuttcap%
\pgfsetmiterjoin%
\definecolor{currentfill}{rgb}{0.000000,0.000000,0.000000}%
\pgfsetfillcolor{currentfill}%
\pgfsetlinewidth{0.000000pt}%
\definecolor{currentstroke}{rgb}{0.000000,0.000000,0.000000}%
\pgfsetstrokecolor{currentstroke}%
\pgfsetstrokeopacity{0.000000}%
\pgfsetdash{}{0pt}%
\pgfpathmoveto{\pgfqpoint{7.346150in}{0.660000in}}%
\pgfpathlineto{\pgfqpoint{7.352350in}{0.660000in}}%
\pgfpathlineto{\pgfqpoint{7.352350in}{3.211106in}}%
\pgfpathlineto{\pgfqpoint{7.346150in}{3.211106in}}%
\pgfpathlineto{\pgfqpoint{7.346150in}{0.660000in}}%
\pgfpathclose%
\pgfusepath{fill}%
\end{pgfscope}%
\begin{pgfscope}%
\pgfpathrectangle{\pgfqpoint{1.250000in}{0.660000in}}{\pgfqpoint{7.750000in}{4.620000in}}%
\pgfusepath{clip}%
\pgfsetbuttcap%
\pgfsetmiterjoin%
\definecolor{currentfill}{rgb}{0.000000,0.000000,0.000000}%
\pgfsetfillcolor{currentfill}%
\pgfsetlinewidth{0.000000pt}%
\definecolor{currentstroke}{rgb}{0.000000,0.000000,0.000000}%
\pgfsetstrokecolor{currentstroke}%
\pgfsetstrokeopacity{0.000000}%
\pgfsetdash{}{0pt}%
\pgfpathmoveto{\pgfqpoint{7.353900in}{0.660000in}}%
\pgfpathlineto{\pgfqpoint{7.360100in}{0.660000in}}%
\pgfpathlineto{\pgfqpoint{7.360100in}{3.210996in}}%
\pgfpathlineto{\pgfqpoint{7.353900in}{3.210996in}}%
\pgfpathlineto{\pgfqpoint{7.353900in}{0.660000in}}%
\pgfpathclose%
\pgfusepath{fill}%
\end{pgfscope}%
\begin{pgfscope}%
\pgfpathrectangle{\pgfqpoint{1.250000in}{0.660000in}}{\pgfqpoint{7.750000in}{4.620000in}}%
\pgfusepath{clip}%
\pgfsetbuttcap%
\pgfsetmiterjoin%
\definecolor{currentfill}{rgb}{0.000000,0.000000,0.000000}%
\pgfsetfillcolor{currentfill}%
\pgfsetlinewidth{0.000000pt}%
\definecolor{currentstroke}{rgb}{0.000000,0.000000,0.000000}%
\pgfsetstrokecolor{currentstroke}%
\pgfsetstrokeopacity{0.000000}%
\pgfsetdash{}{0pt}%
\pgfpathmoveto{\pgfqpoint{7.361650in}{0.660000in}}%
\pgfpathlineto{\pgfqpoint{7.367850in}{0.660000in}}%
\pgfpathlineto{\pgfqpoint{7.367850in}{3.210886in}}%
\pgfpathlineto{\pgfqpoint{7.361650in}{3.210886in}}%
\pgfpathlineto{\pgfqpoint{7.361650in}{0.660000in}}%
\pgfpathclose%
\pgfusepath{fill}%
\end{pgfscope}%
\begin{pgfscope}%
\pgfpathrectangle{\pgfqpoint{1.250000in}{0.660000in}}{\pgfqpoint{7.750000in}{4.620000in}}%
\pgfusepath{clip}%
\pgfsetbuttcap%
\pgfsetmiterjoin%
\definecolor{currentfill}{rgb}{0.000000,0.000000,0.000000}%
\pgfsetfillcolor{currentfill}%
\pgfsetlinewidth{0.000000pt}%
\definecolor{currentstroke}{rgb}{0.000000,0.000000,0.000000}%
\pgfsetstrokecolor{currentstroke}%
\pgfsetstrokeopacity{0.000000}%
\pgfsetdash{}{0pt}%
\pgfpathmoveto{\pgfqpoint{7.369400in}{0.660000in}}%
\pgfpathlineto{\pgfqpoint{7.375600in}{0.660000in}}%
\pgfpathlineto{\pgfqpoint{7.375600in}{3.210778in}}%
\pgfpathlineto{\pgfqpoint{7.369400in}{3.210778in}}%
\pgfpathlineto{\pgfqpoint{7.369400in}{0.660000in}}%
\pgfpathclose%
\pgfusepath{fill}%
\end{pgfscope}%
\begin{pgfscope}%
\pgfpathrectangle{\pgfqpoint{1.250000in}{0.660000in}}{\pgfqpoint{7.750000in}{4.620000in}}%
\pgfusepath{clip}%
\pgfsetbuttcap%
\pgfsetmiterjoin%
\definecolor{currentfill}{rgb}{0.000000,0.000000,0.000000}%
\pgfsetfillcolor{currentfill}%
\pgfsetlinewidth{0.000000pt}%
\definecolor{currentstroke}{rgb}{0.000000,0.000000,0.000000}%
\pgfsetstrokecolor{currentstroke}%
\pgfsetstrokeopacity{0.000000}%
\pgfsetdash{}{0pt}%
\pgfpathmoveto{\pgfqpoint{7.377150in}{0.660000in}}%
\pgfpathlineto{\pgfqpoint{7.383350in}{0.660000in}}%
\pgfpathlineto{\pgfqpoint{7.383350in}{3.210671in}}%
\pgfpathlineto{\pgfqpoint{7.377150in}{3.210671in}}%
\pgfpathlineto{\pgfqpoint{7.377150in}{0.660000in}}%
\pgfpathclose%
\pgfusepath{fill}%
\end{pgfscope}%
\begin{pgfscope}%
\pgfpathrectangle{\pgfqpoint{1.250000in}{0.660000in}}{\pgfqpoint{7.750000in}{4.620000in}}%
\pgfusepath{clip}%
\pgfsetbuttcap%
\pgfsetmiterjoin%
\definecolor{currentfill}{rgb}{0.000000,0.000000,0.000000}%
\pgfsetfillcolor{currentfill}%
\pgfsetlinewidth{0.000000pt}%
\definecolor{currentstroke}{rgb}{0.000000,0.000000,0.000000}%
\pgfsetstrokecolor{currentstroke}%
\pgfsetstrokeopacity{0.000000}%
\pgfsetdash{}{0pt}%
\pgfpathmoveto{\pgfqpoint{7.384900in}{0.660000in}}%
\pgfpathlineto{\pgfqpoint{7.391100in}{0.660000in}}%
\pgfpathlineto{\pgfqpoint{7.391100in}{3.210565in}}%
\pgfpathlineto{\pgfqpoint{7.384900in}{3.210565in}}%
\pgfpathlineto{\pgfqpoint{7.384900in}{0.660000in}}%
\pgfpathclose%
\pgfusepath{fill}%
\end{pgfscope}%
\begin{pgfscope}%
\pgfpathrectangle{\pgfqpoint{1.250000in}{0.660000in}}{\pgfqpoint{7.750000in}{4.620000in}}%
\pgfusepath{clip}%
\pgfsetbuttcap%
\pgfsetmiterjoin%
\definecolor{currentfill}{rgb}{0.000000,0.000000,0.000000}%
\pgfsetfillcolor{currentfill}%
\pgfsetlinewidth{0.000000pt}%
\definecolor{currentstroke}{rgb}{0.000000,0.000000,0.000000}%
\pgfsetstrokecolor{currentstroke}%
\pgfsetstrokeopacity{0.000000}%
\pgfsetdash{}{0pt}%
\pgfpathmoveto{\pgfqpoint{7.392650in}{0.660000in}}%
\pgfpathlineto{\pgfqpoint{7.398850in}{0.660000in}}%
\pgfpathlineto{\pgfqpoint{7.398850in}{3.210459in}}%
\pgfpathlineto{\pgfqpoint{7.392650in}{3.210459in}}%
\pgfpathlineto{\pgfqpoint{7.392650in}{0.660000in}}%
\pgfpathclose%
\pgfusepath{fill}%
\end{pgfscope}%
\begin{pgfscope}%
\pgfpathrectangle{\pgfqpoint{1.250000in}{0.660000in}}{\pgfqpoint{7.750000in}{4.620000in}}%
\pgfusepath{clip}%
\pgfsetbuttcap%
\pgfsetmiterjoin%
\definecolor{currentfill}{rgb}{0.000000,0.000000,0.000000}%
\pgfsetfillcolor{currentfill}%
\pgfsetlinewidth{0.000000pt}%
\definecolor{currentstroke}{rgb}{0.000000,0.000000,0.000000}%
\pgfsetstrokecolor{currentstroke}%
\pgfsetstrokeopacity{0.000000}%
\pgfsetdash{}{0pt}%
\pgfpathmoveto{\pgfqpoint{7.400400in}{0.660000in}}%
\pgfpathlineto{\pgfqpoint{7.406600in}{0.660000in}}%
\pgfpathlineto{\pgfqpoint{7.406600in}{3.210352in}}%
\pgfpathlineto{\pgfqpoint{7.400400in}{3.210352in}}%
\pgfpathlineto{\pgfqpoint{7.400400in}{0.660000in}}%
\pgfpathclose%
\pgfusepath{fill}%
\end{pgfscope}%
\begin{pgfscope}%
\pgfpathrectangle{\pgfqpoint{1.250000in}{0.660000in}}{\pgfqpoint{7.750000in}{4.620000in}}%
\pgfusepath{clip}%
\pgfsetbuttcap%
\pgfsetmiterjoin%
\definecolor{currentfill}{rgb}{0.000000,0.000000,0.000000}%
\pgfsetfillcolor{currentfill}%
\pgfsetlinewidth{0.000000pt}%
\definecolor{currentstroke}{rgb}{0.000000,0.000000,0.000000}%
\pgfsetstrokecolor{currentstroke}%
\pgfsetstrokeopacity{0.000000}%
\pgfsetdash{}{0pt}%
\pgfpathmoveto{\pgfqpoint{7.408150in}{0.660000in}}%
\pgfpathlineto{\pgfqpoint{7.414350in}{0.660000in}}%
\pgfpathlineto{\pgfqpoint{7.414350in}{3.210246in}}%
\pgfpathlineto{\pgfqpoint{7.408150in}{3.210246in}}%
\pgfpathlineto{\pgfqpoint{7.408150in}{0.660000in}}%
\pgfpathclose%
\pgfusepath{fill}%
\end{pgfscope}%
\begin{pgfscope}%
\pgfpathrectangle{\pgfqpoint{1.250000in}{0.660000in}}{\pgfqpoint{7.750000in}{4.620000in}}%
\pgfusepath{clip}%
\pgfsetbuttcap%
\pgfsetmiterjoin%
\definecolor{currentfill}{rgb}{0.000000,0.000000,0.000000}%
\pgfsetfillcolor{currentfill}%
\pgfsetlinewidth{0.000000pt}%
\definecolor{currentstroke}{rgb}{0.000000,0.000000,0.000000}%
\pgfsetstrokecolor{currentstroke}%
\pgfsetstrokeopacity{0.000000}%
\pgfsetdash{}{0pt}%
\pgfpathmoveto{\pgfqpoint{7.415900in}{0.660000in}}%
\pgfpathlineto{\pgfqpoint{7.422100in}{0.660000in}}%
\pgfpathlineto{\pgfqpoint{7.422100in}{3.210139in}}%
\pgfpathlineto{\pgfqpoint{7.415900in}{3.210139in}}%
\pgfpathlineto{\pgfqpoint{7.415900in}{0.660000in}}%
\pgfpathclose%
\pgfusepath{fill}%
\end{pgfscope}%
\begin{pgfscope}%
\pgfpathrectangle{\pgfqpoint{1.250000in}{0.660000in}}{\pgfqpoint{7.750000in}{4.620000in}}%
\pgfusepath{clip}%
\pgfsetbuttcap%
\pgfsetmiterjoin%
\definecolor{currentfill}{rgb}{0.000000,0.000000,0.000000}%
\pgfsetfillcolor{currentfill}%
\pgfsetlinewidth{0.000000pt}%
\definecolor{currentstroke}{rgb}{0.000000,0.000000,0.000000}%
\pgfsetstrokecolor{currentstroke}%
\pgfsetstrokeopacity{0.000000}%
\pgfsetdash{}{0pt}%
\pgfpathmoveto{\pgfqpoint{7.423650in}{0.660000in}}%
\pgfpathlineto{\pgfqpoint{7.429850in}{0.660000in}}%
\pgfpathlineto{\pgfqpoint{7.429850in}{3.210033in}}%
\pgfpathlineto{\pgfqpoint{7.423650in}{3.210033in}}%
\pgfpathlineto{\pgfqpoint{7.423650in}{0.660000in}}%
\pgfpathclose%
\pgfusepath{fill}%
\end{pgfscope}%
\begin{pgfscope}%
\pgfpathrectangle{\pgfqpoint{1.250000in}{0.660000in}}{\pgfqpoint{7.750000in}{4.620000in}}%
\pgfusepath{clip}%
\pgfsetbuttcap%
\pgfsetmiterjoin%
\definecolor{currentfill}{rgb}{0.000000,0.000000,0.000000}%
\pgfsetfillcolor{currentfill}%
\pgfsetlinewidth{0.000000pt}%
\definecolor{currentstroke}{rgb}{0.000000,0.000000,0.000000}%
\pgfsetstrokecolor{currentstroke}%
\pgfsetstrokeopacity{0.000000}%
\pgfsetdash{}{0pt}%
\pgfpathmoveto{\pgfqpoint{7.431400in}{0.660000in}}%
\pgfpathlineto{\pgfqpoint{7.437600in}{0.660000in}}%
\pgfpathlineto{\pgfqpoint{7.437600in}{3.209926in}}%
\pgfpathlineto{\pgfqpoint{7.431400in}{3.209926in}}%
\pgfpathlineto{\pgfqpoint{7.431400in}{0.660000in}}%
\pgfpathclose%
\pgfusepath{fill}%
\end{pgfscope}%
\begin{pgfscope}%
\pgfpathrectangle{\pgfqpoint{1.250000in}{0.660000in}}{\pgfqpoint{7.750000in}{4.620000in}}%
\pgfusepath{clip}%
\pgfsetbuttcap%
\pgfsetmiterjoin%
\definecolor{currentfill}{rgb}{0.000000,0.000000,0.000000}%
\pgfsetfillcolor{currentfill}%
\pgfsetlinewidth{0.000000pt}%
\definecolor{currentstroke}{rgb}{0.000000,0.000000,0.000000}%
\pgfsetstrokecolor{currentstroke}%
\pgfsetstrokeopacity{0.000000}%
\pgfsetdash{}{0pt}%
\pgfpathmoveto{\pgfqpoint{7.439150in}{0.660000in}}%
\pgfpathlineto{\pgfqpoint{7.445350in}{0.660000in}}%
\pgfpathlineto{\pgfqpoint{7.445350in}{3.209820in}}%
\pgfpathlineto{\pgfqpoint{7.439150in}{3.209820in}}%
\pgfpathlineto{\pgfqpoint{7.439150in}{0.660000in}}%
\pgfpathclose%
\pgfusepath{fill}%
\end{pgfscope}%
\begin{pgfscope}%
\pgfpathrectangle{\pgfqpoint{1.250000in}{0.660000in}}{\pgfqpoint{7.750000in}{4.620000in}}%
\pgfusepath{clip}%
\pgfsetbuttcap%
\pgfsetmiterjoin%
\definecolor{currentfill}{rgb}{0.000000,0.000000,0.000000}%
\pgfsetfillcolor{currentfill}%
\pgfsetlinewidth{0.000000pt}%
\definecolor{currentstroke}{rgb}{0.000000,0.000000,0.000000}%
\pgfsetstrokecolor{currentstroke}%
\pgfsetstrokeopacity{0.000000}%
\pgfsetdash{}{0pt}%
\pgfpathmoveto{\pgfqpoint{7.446900in}{0.660000in}}%
\pgfpathlineto{\pgfqpoint{7.453100in}{0.660000in}}%
\pgfpathlineto{\pgfqpoint{7.453100in}{3.209714in}}%
\pgfpathlineto{\pgfqpoint{7.446900in}{3.209714in}}%
\pgfpathlineto{\pgfqpoint{7.446900in}{0.660000in}}%
\pgfpathclose%
\pgfusepath{fill}%
\end{pgfscope}%
\begin{pgfscope}%
\pgfpathrectangle{\pgfqpoint{1.250000in}{0.660000in}}{\pgfqpoint{7.750000in}{4.620000in}}%
\pgfusepath{clip}%
\pgfsetbuttcap%
\pgfsetmiterjoin%
\definecolor{currentfill}{rgb}{0.000000,0.000000,0.000000}%
\pgfsetfillcolor{currentfill}%
\pgfsetlinewidth{0.000000pt}%
\definecolor{currentstroke}{rgb}{0.000000,0.000000,0.000000}%
\pgfsetstrokecolor{currentstroke}%
\pgfsetstrokeopacity{0.000000}%
\pgfsetdash{}{0pt}%
\pgfpathmoveto{\pgfqpoint{7.454650in}{0.660000in}}%
\pgfpathlineto{\pgfqpoint{7.460850in}{0.660000in}}%
\pgfpathlineto{\pgfqpoint{7.460850in}{3.209607in}}%
\pgfpathlineto{\pgfqpoint{7.454650in}{3.209607in}}%
\pgfpathlineto{\pgfqpoint{7.454650in}{0.660000in}}%
\pgfpathclose%
\pgfusepath{fill}%
\end{pgfscope}%
\begin{pgfscope}%
\pgfpathrectangle{\pgfqpoint{1.250000in}{0.660000in}}{\pgfqpoint{7.750000in}{4.620000in}}%
\pgfusepath{clip}%
\pgfsetbuttcap%
\pgfsetmiterjoin%
\definecolor{currentfill}{rgb}{0.000000,0.000000,0.000000}%
\pgfsetfillcolor{currentfill}%
\pgfsetlinewidth{0.000000pt}%
\definecolor{currentstroke}{rgb}{0.000000,0.000000,0.000000}%
\pgfsetstrokecolor{currentstroke}%
\pgfsetstrokeopacity{0.000000}%
\pgfsetdash{}{0pt}%
\pgfpathmoveto{\pgfqpoint{7.462400in}{0.660000in}}%
\pgfpathlineto{\pgfqpoint{7.468600in}{0.660000in}}%
\pgfpathlineto{\pgfqpoint{7.468600in}{3.209501in}}%
\pgfpathlineto{\pgfqpoint{7.462400in}{3.209501in}}%
\pgfpathlineto{\pgfqpoint{7.462400in}{0.660000in}}%
\pgfpathclose%
\pgfusepath{fill}%
\end{pgfscope}%
\begin{pgfscope}%
\pgfpathrectangle{\pgfqpoint{1.250000in}{0.660000in}}{\pgfqpoint{7.750000in}{4.620000in}}%
\pgfusepath{clip}%
\pgfsetbuttcap%
\pgfsetmiterjoin%
\definecolor{currentfill}{rgb}{0.000000,0.000000,0.000000}%
\pgfsetfillcolor{currentfill}%
\pgfsetlinewidth{0.000000pt}%
\definecolor{currentstroke}{rgb}{0.000000,0.000000,0.000000}%
\pgfsetstrokecolor{currentstroke}%
\pgfsetstrokeopacity{0.000000}%
\pgfsetdash{}{0pt}%
\pgfpathmoveto{\pgfqpoint{7.470150in}{0.660000in}}%
\pgfpathlineto{\pgfqpoint{7.476350in}{0.660000in}}%
\pgfpathlineto{\pgfqpoint{7.476350in}{3.209394in}}%
\pgfpathlineto{\pgfqpoint{7.470150in}{3.209394in}}%
\pgfpathlineto{\pgfqpoint{7.470150in}{0.660000in}}%
\pgfpathclose%
\pgfusepath{fill}%
\end{pgfscope}%
\begin{pgfscope}%
\pgfpathrectangle{\pgfqpoint{1.250000in}{0.660000in}}{\pgfqpoint{7.750000in}{4.620000in}}%
\pgfusepath{clip}%
\pgfsetbuttcap%
\pgfsetmiterjoin%
\definecolor{currentfill}{rgb}{0.000000,0.000000,0.000000}%
\pgfsetfillcolor{currentfill}%
\pgfsetlinewidth{0.000000pt}%
\definecolor{currentstroke}{rgb}{0.000000,0.000000,0.000000}%
\pgfsetstrokecolor{currentstroke}%
\pgfsetstrokeopacity{0.000000}%
\pgfsetdash{}{0pt}%
\pgfpathmoveto{\pgfqpoint{7.477900in}{0.660000in}}%
\pgfpathlineto{\pgfqpoint{7.484100in}{0.660000in}}%
\pgfpathlineto{\pgfqpoint{7.484100in}{3.209288in}}%
\pgfpathlineto{\pgfqpoint{7.477900in}{3.209288in}}%
\pgfpathlineto{\pgfqpoint{7.477900in}{0.660000in}}%
\pgfpathclose%
\pgfusepath{fill}%
\end{pgfscope}%
\begin{pgfscope}%
\pgfpathrectangle{\pgfqpoint{1.250000in}{0.660000in}}{\pgfqpoint{7.750000in}{4.620000in}}%
\pgfusepath{clip}%
\pgfsetbuttcap%
\pgfsetmiterjoin%
\definecolor{currentfill}{rgb}{0.000000,0.000000,0.000000}%
\pgfsetfillcolor{currentfill}%
\pgfsetlinewidth{0.000000pt}%
\definecolor{currentstroke}{rgb}{0.000000,0.000000,0.000000}%
\pgfsetstrokecolor{currentstroke}%
\pgfsetstrokeopacity{0.000000}%
\pgfsetdash{}{0pt}%
\pgfpathmoveto{\pgfqpoint{7.485650in}{0.660000in}}%
\pgfpathlineto{\pgfqpoint{7.491850in}{0.660000in}}%
\pgfpathlineto{\pgfqpoint{7.491850in}{3.209182in}}%
\pgfpathlineto{\pgfqpoint{7.485650in}{3.209182in}}%
\pgfpathlineto{\pgfqpoint{7.485650in}{0.660000in}}%
\pgfpathclose%
\pgfusepath{fill}%
\end{pgfscope}%
\begin{pgfscope}%
\pgfpathrectangle{\pgfqpoint{1.250000in}{0.660000in}}{\pgfqpoint{7.750000in}{4.620000in}}%
\pgfusepath{clip}%
\pgfsetbuttcap%
\pgfsetmiterjoin%
\definecolor{currentfill}{rgb}{0.000000,0.000000,0.000000}%
\pgfsetfillcolor{currentfill}%
\pgfsetlinewidth{0.000000pt}%
\definecolor{currentstroke}{rgb}{0.000000,0.000000,0.000000}%
\pgfsetstrokecolor{currentstroke}%
\pgfsetstrokeopacity{0.000000}%
\pgfsetdash{}{0pt}%
\pgfpathmoveto{\pgfqpoint{7.493400in}{0.660000in}}%
\pgfpathlineto{\pgfqpoint{7.499600in}{0.660000in}}%
\pgfpathlineto{\pgfqpoint{7.499600in}{3.209075in}}%
\pgfpathlineto{\pgfqpoint{7.493400in}{3.209075in}}%
\pgfpathlineto{\pgfqpoint{7.493400in}{0.660000in}}%
\pgfpathclose%
\pgfusepath{fill}%
\end{pgfscope}%
\begin{pgfscope}%
\pgfpathrectangle{\pgfqpoint{1.250000in}{0.660000in}}{\pgfqpoint{7.750000in}{4.620000in}}%
\pgfusepath{clip}%
\pgfsetbuttcap%
\pgfsetmiterjoin%
\definecolor{currentfill}{rgb}{0.000000,0.000000,0.000000}%
\pgfsetfillcolor{currentfill}%
\pgfsetlinewidth{0.000000pt}%
\definecolor{currentstroke}{rgb}{0.000000,0.000000,0.000000}%
\pgfsetstrokecolor{currentstroke}%
\pgfsetstrokeopacity{0.000000}%
\pgfsetdash{}{0pt}%
\pgfpathmoveto{\pgfqpoint{7.501150in}{0.660000in}}%
\pgfpathlineto{\pgfqpoint{7.507350in}{0.660000in}}%
\pgfpathlineto{\pgfqpoint{7.507350in}{3.208969in}}%
\pgfpathlineto{\pgfqpoint{7.501150in}{3.208969in}}%
\pgfpathlineto{\pgfqpoint{7.501150in}{0.660000in}}%
\pgfpathclose%
\pgfusepath{fill}%
\end{pgfscope}%
\begin{pgfscope}%
\pgfpathrectangle{\pgfqpoint{1.250000in}{0.660000in}}{\pgfqpoint{7.750000in}{4.620000in}}%
\pgfusepath{clip}%
\pgfsetbuttcap%
\pgfsetmiterjoin%
\definecolor{currentfill}{rgb}{0.000000,0.000000,0.000000}%
\pgfsetfillcolor{currentfill}%
\pgfsetlinewidth{0.000000pt}%
\definecolor{currentstroke}{rgb}{0.000000,0.000000,0.000000}%
\pgfsetstrokecolor{currentstroke}%
\pgfsetstrokeopacity{0.000000}%
\pgfsetdash{}{0pt}%
\pgfpathmoveto{\pgfqpoint{7.508900in}{0.660000in}}%
\pgfpathlineto{\pgfqpoint{7.515100in}{0.660000in}}%
\pgfpathlineto{\pgfqpoint{7.515100in}{3.208862in}}%
\pgfpathlineto{\pgfqpoint{7.508900in}{3.208862in}}%
\pgfpathlineto{\pgfqpoint{7.508900in}{0.660000in}}%
\pgfpathclose%
\pgfusepath{fill}%
\end{pgfscope}%
\begin{pgfscope}%
\pgfpathrectangle{\pgfqpoint{1.250000in}{0.660000in}}{\pgfqpoint{7.750000in}{4.620000in}}%
\pgfusepath{clip}%
\pgfsetbuttcap%
\pgfsetmiterjoin%
\definecolor{currentfill}{rgb}{0.000000,0.000000,0.000000}%
\pgfsetfillcolor{currentfill}%
\pgfsetlinewidth{0.000000pt}%
\definecolor{currentstroke}{rgb}{0.000000,0.000000,0.000000}%
\pgfsetstrokecolor{currentstroke}%
\pgfsetstrokeopacity{0.000000}%
\pgfsetdash{}{0pt}%
\pgfpathmoveto{\pgfqpoint{7.516650in}{0.660000in}}%
\pgfpathlineto{\pgfqpoint{7.522850in}{0.660000in}}%
\pgfpathlineto{\pgfqpoint{7.522850in}{3.208757in}}%
\pgfpathlineto{\pgfqpoint{7.516650in}{3.208757in}}%
\pgfpathlineto{\pgfqpoint{7.516650in}{0.660000in}}%
\pgfpathclose%
\pgfusepath{fill}%
\end{pgfscope}%
\begin{pgfscope}%
\pgfpathrectangle{\pgfqpoint{1.250000in}{0.660000in}}{\pgfqpoint{7.750000in}{4.620000in}}%
\pgfusepath{clip}%
\pgfsetbuttcap%
\pgfsetmiterjoin%
\definecolor{currentfill}{rgb}{0.000000,0.000000,0.000000}%
\pgfsetfillcolor{currentfill}%
\pgfsetlinewidth{0.000000pt}%
\definecolor{currentstroke}{rgb}{0.000000,0.000000,0.000000}%
\pgfsetstrokecolor{currentstroke}%
\pgfsetstrokeopacity{0.000000}%
\pgfsetdash{}{0pt}%
\pgfpathmoveto{\pgfqpoint{7.524400in}{0.660000in}}%
\pgfpathlineto{\pgfqpoint{7.530600in}{0.660000in}}%
\pgfpathlineto{\pgfqpoint{7.530600in}{3.208654in}}%
\pgfpathlineto{\pgfqpoint{7.524400in}{3.208654in}}%
\pgfpathlineto{\pgfqpoint{7.524400in}{0.660000in}}%
\pgfpathclose%
\pgfusepath{fill}%
\end{pgfscope}%
\begin{pgfscope}%
\pgfpathrectangle{\pgfqpoint{1.250000in}{0.660000in}}{\pgfqpoint{7.750000in}{4.620000in}}%
\pgfusepath{clip}%
\pgfsetbuttcap%
\pgfsetmiterjoin%
\definecolor{currentfill}{rgb}{0.000000,0.000000,0.000000}%
\pgfsetfillcolor{currentfill}%
\pgfsetlinewidth{0.000000pt}%
\definecolor{currentstroke}{rgb}{0.000000,0.000000,0.000000}%
\pgfsetstrokecolor{currentstroke}%
\pgfsetstrokeopacity{0.000000}%
\pgfsetdash{}{0pt}%
\pgfpathmoveto{\pgfqpoint{7.532150in}{0.660000in}}%
\pgfpathlineto{\pgfqpoint{7.538350in}{0.660000in}}%
\pgfpathlineto{\pgfqpoint{7.538350in}{3.208551in}}%
\pgfpathlineto{\pgfqpoint{7.532150in}{3.208551in}}%
\pgfpathlineto{\pgfqpoint{7.532150in}{0.660000in}}%
\pgfpathclose%
\pgfusepath{fill}%
\end{pgfscope}%
\begin{pgfscope}%
\pgfpathrectangle{\pgfqpoint{1.250000in}{0.660000in}}{\pgfqpoint{7.750000in}{4.620000in}}%
\pgfusepath{clip}%
\pgfsetbuttcap%
\pgfsetmiterjoin%
\definecolor{currentfill}{rgb}{0.000000,0.000000,0.000000}%
\pgfsetfillcolor{currentfill}%
\pgfsetlinewidth{0.000000pt}%
\definecolor{currentstroke}{rgb}{0.000000,0.000000,0.000000}%
\pgfsetstrokecolor{currentstroke}%
\pgfsetstrokeopacity{0.000000}%
\pgfsetdash{}{0pt}%
\pgfpathmoveto{\pgfqpoint{7.539900in}{0.660000in}}%
\pgfpathlineto{\pgfqpoint{7.546100in}{0.660000in}}%
\pgfpathlineto{\pgfqpoint{7.546100in}{3.208448in}}%
\pgfpathlineto{\pgfqpoint{7.539900in}{3.208448in}}%
\pgfpathlineto{\pgfqpoint{7.539900in}{0.660000in}}%
\pgfpathclose%
\pgfusepath{fill}%
\end{pgfscope}%
\begin{pgfscope}%
\pgfpathrectangle{\pgfqpoint{1.250000in}{0.660000in}}{\pgfqpoint{7.750000in}{4.620000in}}%
\pgfusepath{clip}%
\pgfsetbuttcap%
\pgfsetmiterjoin%
\definecolor{currentfill}{rgb}{0.000000,0.000000,0.000000}%
\pgfsetfillcolor{currentfill}%
\pgfsetlinewidth{0.000000pt}%
\definecolor{currentstroke}{rgb}{0.000000,0.000000,0.000000}%
\pgfsetstrokecolor{currentstroke}%
\pgfsetstrokeopacity{0.000000}%
\pgfsetdash{}{0pt}%
\pgfpathmoveto{\pgfqpoint{7.547650in}{0.660000in}}%
\pgfpathlineto{\pgfqpoint{7.553850in}{0.660000in}}%
\pgfpathlineto{\pgfqpoint{7.553850in}{3.208345in}}%
\pgfpathlineto{\pgfqpoint{7.547650in}{3.208345in}}%
\pgfpathlineto{\pgfqpoint{7.547650in}{0.660000in}}%
\pgfpathclose%
\pgfusepath{fill}%
\end{pgfscope}%
\begin{pgfscope}%
\pgfpathrectangle{\pgfqpoint{1.250000in}{0.660000in}}{\pgfqpoint{7.750000in}{4.620000in}}%
\pgfusepath{clip}%
\pgfsetbuttcap%
\pgfsetmiterjoin%
\definecolor{currentfill}{rgb}{0.000000,0.000000,0.000000}%
\pgfsetfillcolor{currentfill}%
\pgfsetlinewidth{0.000000pt}%
\definecolor{currentstroke}{rgb}{0.000000,0.000000,0.000000}%
\pgfsetstrokecolor{currentstroke}%
\pgfsetstrokeopacity{0.000000}%
\pgfsetdash{}{0pt}%
\pgfpathmoveto{\pgfqpoint{7.555400in}{0.660000in}}%
\pgfpathlineto{\pgfqpoint{7.561600in}{0.660000in}}%
\pgfpathlineto{\pgfqpoint{7.561600in}{3.208242in}}%
\pgfpathlineto{\pgfqpoint{7.555400in}{3.208242in}}%
\pgfpathlineto{\pgfqpoint{7.555400in}{0.660000in}}%
\pgfpathclose%
\pgfusepath{fill}%
\end{pgfscope}%
\begin{pgfscope}%
\pgfpathrectangle{\pgfqpoint{1.250000in}{0.660000in}}{\pgfqpoint{7.750000in}{4.620000in}}%
\pgfusepath{clip}%
\pgfsetbuttcap%
\pgfsetmiterjoin%
\definecolor{currentfill}{rgb}{0.000000,0.000000,0.000000}%
\pgfsetfillcolor{currentfill}%
\pgfsetlinewidth{0.000000pt}%
\definecolor{currentstroke}{rgb}{0.000000,0.000000,0.000000}%
\pgfsetstrokecolor{currentstroke}%
\pgfsetstrokeopacity{0.000000}%
\pgfsetdash{}{0pt}%
\pgfpathmoveto{\pgfqpoint{7.563150in}{0.660000in}}%
\pgfpathlineto{\pgfqpoint{7.569350in}{0.660000in}}%
\pgfpathlineto{\pgfqpoint{7.569350in}{3.208139in}}%
\pgfpathlineto{\pgfqpoint{7.563150in}{3.208139in}}%
\pgfpathlineto{\pgfqpoint{7.563150in}{0.660000in}}%
\pgfpathclose%
\pgfusepath{fill}%
\end{pgfscope}%
\begin{pgfscope}%
\pgfpathrectangle{\pgfqpoint{1.250000in}{0.660000in}}{\pgfqpoint{7.750000in}{4.620000in}}%
\pgfusepath{clip}%
\pgfsetbuttcap%
\pgfsetmiterjoin%
\definecolor{currentfill}{rgb}{0.000000,0.000000,0.000000}%
\pgfsetfillcolor{currentfill}%
\pgfsetlinewidth{0.000000pt}%
\definecolor{currentstroke}{rgb}{0.000000,0.000000,0.000000}%
\pgfsetstrokecolor{currentstroke}%
\pgfsetstrokeopacity{0.000000}%
\pgfsetdash{}{0pt}%
\pgfpathmoveto{\pgfqpoint{7.570900in}{0.660000in}}%
\pgfpathlineto{\pgfqpoint{7.577100in}{0.660000in}}%
\pgfpathlineto{\pgfqpoint{7.577100in}{3.208036in}}%
\pgfpathlineto{\pgfqpoint{7.570900in}{3.208036in}}%
\pgfpathlineto{\pgfqpoint{7.570900in}{0.660000in}}%
\pgfpathclose%
\pgfusepath{fill}%
\end{pgfscope}%
\begin{pgfscope}%
\pgfpathrectangle{\pgfqpoint{1.250000in}{0.660000in}}{\pgfqpoint{7.750000in}{4.620000in}}%
\pgfusepath{clip}%
\pgfsetbuttcap%
\pgfsetmiterjoin%
\definecolor{currentfill}{rgb}{0.000000,0.000000,0.000000}%
\pgfsetfillcolor{currentfill}%
\pgfsetlinewidth{0.000000pt}%
\definecolor{currentstroke}{rgb}{0.000000,0.000000,0.000000}%
\pgfsetstrokecolor{currentstroke}%
\pgfsetstrokeopacity{0.000000}%
\pgfsetdash{}{0pt}%
\pgfpathmoveto{\pgfqpoint{7.578650in}{0.660000in}}%
\pgfpathlineto{\pgfqpoint{7.584850in}{0.660000in}}%
\pgfpathlineto{\pgfqpoint{7.584850in}{3.207933in}}%
\pgfpathlineto{\pgfqpoint{7.578650in}{3.207933in}}%
\pgfpathlineto{\pgfqpoint{7.578650in}{0.660000in}}%
\pgfpathclose%
\pgfusepath{fill}%
\end{pgfscope}%
\begin{pgfscope}%
\pgfpathrectangle{\pgfqpoint{1.250000in}{0.660000in}}{\pgfqpoint{7.750000in}{4.620000in}}%
\pgfusepath{clip}%
\pgfsetbuttcap%
\pgfsetmiterjoin%
\definecolor{currentfill}{rgb}{0.000000,0.000000,0.000000}%
\pgfsetfillcolor{currentfill}%
\pgfsetlinewidth{0.000000pt}%
\definecolor{currentstroke}{rgb}{0.000000,0.000000,0.000000}%
\pgfsetstrokecolor{currentstroke}%
\pgfsetstrokeopacity{0.000000}%
\pgfsetdash{}{0pt}%
\pgfpathmoveto{\pgfqpoint{7.586400in}{0.660000in}}%
\pgfpathlineto{\pgfqpoint{7.592600in}{0.660000in}}%
\pgfpathlineto{\pgfqpoint{7.592600in}{3.207830in}}%
\pgfpathlineto{\pgfqpoint{7.586400in}{3.207830in}}%
\pgfpathlineto{\pgfqpoint{7.586400in}{0.660000in}}%
\pgfpathclose%
\pgfusepath{fill}%
\end{pgfscope}%
\begin{pgfscope}%
\pgfpathrectangle{\pgfqpoint{1.250000in}{0.660000in}}{\pgfqpoint{7.750000in}{4.620000in}}%
\pgfusepath{clip}%
\pgfsetbuttcap%
\pgfsetmiterjoin%
\definecolor{currentfill}{rgb}{0.000000,0.000000,0.000000}%
\pgfsetfillcolor{currentfill}%
\pgfsetlinewidth{0.000000pt}%
\definecolor{currentstroke}{rgb}{0.000000,0.000000,0.000000}%
\pgfsetstrokecolor{currentstroke}%
\pgfsetstrokeopacity{0.000000}%
\pgfsetdash{}{0pt}%
\pgfpathmoveto{\pgfqpoint{7.594150in}{0.660000in}}%
\pgfpathlineto{\pgfqpoint{7.600350in}{0.660000in}}%
\pgfpathlineto{\pgfqpoint{7.600350in}{3.207727in}}%
\pgfpathlineto{\pgfqpoint{7.594150in}{3.207727in}}%
\pgfpathlineto{\pgfqpoint{7.594150in}{0.660000in}}%
\pgfpathclose%
\pgfusepath{fill}%
\end{pgfscope}%
\begin{pgfscope}%
\pgfpathrectangle{\pgfqpoint{1.250000in}{0.660000in}}{\pgfqpoint{7.750000in}{4.620000in}}%
\pgfusepath{clip}%
\pgfsetbuttcap%
\pgfsetmiterjoin%
\definecolor{currentfill}{rgb}{0.000000,0.000000,0.000000}%
\pgfsetfillcolor{currentfill}%
\pgfsetlinewidth{0.000000pt}%
\definecolor{currentstroke}{rgb}{0.000000,0.000000,0.000000}%
\pgfsetstrokecolor{currentstroke}%
\pgfsetstrokeopacity{0.000000}%
\pgfsetdash{}{0pt}%
\pgfpathmoveto{\pgfqpoint{7.601900in}{0.660000in}}%
\pgfpathlineto{\pgfqpoint{7.608100in}{0.660000in}}%
\pgfpathlineto{\pgfqpoint{7.608100in}{3.207624in}}%
\pgfpathlineto{\pgfqpoint{7.601900in}{3.207624in}}%
\pgfpathlineto{\pgfqpoint{7.601900in}{0.660000in}}%
\pgfpathclose%
\pgfusepath{fill}%
\end{pgfscope}%
\begin{pgfscope}%
\pgfpathrectangle{\pgfqpoint{1.250000in}{0.660000in}}{\pgfqpoint{7.750000in}{4.620000in}}%
\pgfusepath{clip}%
\pgfsetbuttcap%
\pgfsetmiterjoin%
\definecolor{currentfill}{rgb}{0.000000,0.000000,0.000000}%
\pgfsetfillcolor{currentfill}%
\pgfsetlinewidth{0.000000pt}%
\definecolor{currentstroke}{rgb}{0.000000,0.000000,0.000000}%
\pgfsetstrokecolor{currentstroke}%
\pgfsetstrokeopacity{0.000000}%
\pgfsetdash{}{0pt}%
\pgfpathmoveto{\pgfqpoint{7.609650in}{0.660000in}}%
\pgfpathlineto{\pgfqpoint{7.615850in}{0.660000in}}%
\pgfpathlineto{\pgfqpoint{7.615850in}{3.207522in}}%
\pgfpathlineto{\pgfqpoint{7.609650in}{3.207522in}}%
\pgfpathlineto{\pgfqpoint{7.609650in}{0.660000in}}%
\pgfpathclose%
\pgfusepath{fill}%
\end{pgfscope}%
\begin{pgfscope}%
\pgfpathrectangle{\pgfqpoint{1.250000in}{0.660000in}}{\pgfqpoint{7.750000in}{4.620000in}}%
\pgfusepath{clip}%
\pgfsetbuttcap%
\pgfsetmiterjoin%
\definecolor{currentfill}{rgb}{0.000000,0.000000,0.000000}%
\pgfsetfillcolor{currentfill}%
\pgfsetlinewidth{0.000000pt}%
\definecolor{currentstroke}{rgb}{0.000000,0.000000,0.000000}%
\pgfsetstrokecolor{currentstroke}%
\pgfsetstrokeopacity{0.000000}%
\pgfsetdash{}{0pt}%
\pgfpathmoveto{\pgfqpoint{7.617400in}{0.660000in}}%
\pgfpathlineto{\pgfqpoint{7.623600in}{0.660000in}}%
\pgfpathlineto{\pgfqpoint{7.623600in}{3.207419in}}%
\pgfpathlineto{\pgfqpoint{7.617400in}{3.207419in}}%
\pgfpathlineto{\pgfqpoint{7.617400in}{0.660000in}}%
\pgfpathclose%
\pgfusepath{fill}%
\end{pgfscope}%
\begin{pgfscope}%
\pgfpathrectangle{\pgfqpoint{1.250000in}{0.660000in}}{\pgfqpoint{7.750000in}{4.620000in}}%
\pgfusepath{clip}%
\pgfsetbuttcap%
\pgfsetmiterjoin%
\definecolor{currentfill}{rgb}{0.000000,0.000000,0.000000}%
\pgfsetfillcolor{currentfill}%
\pgfsetlinewidth{0.000000pt}%
\definecolor{currentstroke}{rgb}{0.000000,0.000000,0.000000}%
\pgfsetstrokecolor{currentstroke}%
\pgfsetstrokeopacity{0.000000}%
\pgfsetdash{}{0pt}%
\pgfpathmoveto{\pgfqpoint{7.625150in}{0.660000in}}%
\pgfpathlineto{\pgfqpoint{7.631350in}{0.660000in}}%
\pgfpathlineto{\pgfqpoint{7.631350in}{3.207316in}}%
\pgfpathlineto{\pgfqpoint{7.625150in}{3.207316in}}%
\pgfpathlineto{\pgfqpoint{7.625150in}{0.660000in}}%
\pgfpathclose%
\pgfusepath{fill}%
\end{pgfscope}%
\begin{pgfscope}%
\pgfpathrectangle{\pgfqpoint{1.250000in}{0.660000in}}{\pgfqpoint{7.750000in}{4.620000in}}%
\pgfusepath{clip}%
\pgfsetbuttcap%
\pgfsetmiterjoin%
\definecolor{currentfill}{rgb}{0.000000,0.000000,0.000000}%
\pgfsetfillcolor{currentfill}%
\pgfsetlinewidth{0.000000pt}%
\definecolor{currentstroke}{rgb}{0.000000,0.000000,0.000000}%
\pgfsetstrokecolor{currentstroke}%
\pgfsetstrokeopacity{0.000000}%
\pgfsetdash{}{0pt}%
\pgfpathmoveto{\pgfqpoint{7.632900in}{0.660000in}}%
\pgfpathlineto{\pgfqpoint{7.639100in}{0.660000in}}%
\pgfpathlineto{\pgfqpoint{7.639100in}{3.207213in}}%
\pgfpathlineto{\pgfqpoint{7.632900in}{3.207213in}}%
\pgfpathlineto{\pgfqpoint{7.632900in}{0.660000in}}%
\pgfpathclose%
\pgfusepath{fill}%
\end{pgfscope}%
\begin{pgfscope}%
\pgfpathrectangle{\pgfqpoint{1.250000in}{0.660000in}}{\pgfqpoint{7.750000in}{4.620000in}}%
\pgfusepath{clip}%
\pgfsetbuttcap%
\pgfsetmiterjoin%
\definecolor{currentfill}{rgb}{0.000000,0.000000,0.000000}%
\pgfsetfillcolor{currentfill}%
\pgfsetlinewidth{0.000000pt}%
\definecolor{currentstroke}{rgb}{0.000000,0.000000,0.000000}%
\pgfsetstrokecolor{currentstroke}%
\pgfsetstrokeopacity{0.000000}%
\pgfsetdash{}{0pt}%
\pgfpathmoveto{\pgfqpoint{7.640650in}{0.660000in}}%
\pgfpathlineto{\pgfqpoint{7.646850in}{0.660000in}}%
\pgfpathlineto{\pgfqpoint{7.646850in}{3.207110in}}%
\pgfpathlineto{\pgfqpoint{7.640650in}{3.207110in}}%
\pgfpathlineto{\pgfqpoint{7.640650in}{0.660000in}}%
\pgfpathclose%
\pgfusepath{fill}%
\end{pgfscope}%
\begin{pgfscope}%
\pgfpathrectangle{\pgfqpoint{1.250000in}{0.660000in}}{\pgfqpoint{7.750000in}{4.620000in}}%
\pgfusepath{clip}%
\pgfsetbuttcap%
\pgfsetmiterjoin%
\definecolor{currentfill}{rgb}{0.000000,0.000000,0.000000}%
\pgfsetfillcolor{currentfill}%
\pgfsetlinewidth{0.000000pt}%
\definecolor{currentstroke}{rgb}{0.000000,0.000000,0.000000}%
\pgfsetstrokecolor{currentstroke}%
\pgfsetstrokeopacity{0.000000}%
\pgfsetdash{}{0pt}%
\pgfpathmoveto{\pgfqpoint{7.648400in}{0.660000in}}%
\pgfpathlineto{\pgfqpoint{7.654600in}{0.660000in}}%
\pgfpathlineto{\pgfqpoint{7.654600in}{3.207007in}}%
\pgfpathlineto{\pgfqpoint{7.648400in}{3.207007in}}%
\pgfpathlineto{\pgfqpoint{7.648400in}{0.660000in}}%
\pgfpathclose%
\pgfusepath{fill}%
\end{pgfscope}%
\begin{pgfscope}%
\pgfpathrectangle{\pgfqpoint{1.250000in}{0.660000in}}{\pgfqpoint{7.750000in}{4.620000in}}%
\pgfusepath{clip}%
\pgfsetbuttcap%
\pgfsetmiterjoin%
\definecolor{currentfill}{rgb}{0.000000,0.000000,0.000000}%
\pgfsetfillcolor{currentfill}%
\pgfsetlinewidth{0.000000pt}%
\definecolor{currentstroke}{rgb}{0.000000,0.000000,0.000000}%
\pgfsetstrokecolor{currentstroke}%
\pgfsetstrokeopacity{0.000000}%
\pgfsetdash{}{0pt}%
\pgfpathmoveto{\pgfqpoint{7.656150in}{0.660000in}}%
\pgfpathlineto{\pgfqpoint{7.662350in}{0.660000in}}%
\pgfpathlineto{\pgfqpoint{7.662350in}{3.206904in}}%
\pgfpathlineto{\pgfqpoint{7.656150in}{3.206904in}}%
\pgfpathlineto{\pgfqpoint{7.656150in}{0.660000in}}%
\pgfpathclose%
\pgfusepath{fill}%
\end{pgfscope}%
\begin{pgfscope}%
\pgfpathrectangle{\pgfqpoint{1.250000in}{0.660000in}}{\pgfqpoint{7.750000in}{4.620000in}}%
\pgfusepath{clip}%
\pgfsetbuttcap%
\pgfsetmiterjoin%
\definecolor{currentfill}{rgb}{0.000000,0.000000,0.000000}%
\pgfsetfillcolor{currentfill}%
\pgfsetlinewidth{0.000000pt}%
\definecolor{currentstroke}{rgb}{0.000000,0.000000,0.000000}%
\pgfsetstrokecolor{currentstroke}%
\pgfsetstrokeopacity{0.000000}%
\pgfsetdash{}{0pt}%
\pgfpathmoveto{\pgfqpoint{7.663900in}{0.660000in}}%
\pgfpathlineto{\pgfqpoint{7.670100in}{0.660000in}}%
\pgfpathlineto{\pgfqpoint{7.670100in}{3.206801in}}%
\pgfpathlineto{\pgfqpoint{7.663900in}{3.206801in}}%
\pgfpathlineto{\pgfqpoint{7.663900in}{0.660000in}}%
\pgfpathclose%
\pgfusepath{fill}%
\end{pgfscope}%
\begin{pgfscope}%
\pgfpathrectangle{\pgfqpoint{1.250000in}{0.660000in}}{\pgfqpoint{7.750000in}{4.620000in}}%
\pgfusepath{clip}%
\pgfsetbuttcap%
\pgfsetmiterjoin%
\definecolor{currentfill}{rgb}{0.000000,0.000000,0.000000}%
\pgfsetfillcolor{currentfill}%
\pgfsetlinewidth{0.000000pt}%
\definecolor{currentstroke}{rgb}{0.000000,0.000000,0.000000}%
\pgfsetstrokecolor{currentstroke}%
\pgfsetstrokeopacity{0.000000}%
\pgfsetdash{}{0pt}%
\pgfpathmoveto{\pgfqpoint{7.671650in}{0.660000in}}%
\pgfpathlineto{\pgfqpoint{7.677850in}{0.660000in}}%
\pgfpathlineto{\pgfqpoint{7.677850in}{3.206699in}}%
\pgfpathlineto{\pgfqpoint{7.671650in}{3.206699in}}%
\pgfpathlineto{\pgfqpoint{7.671650in}{0.660000in}}%
\pgfpathclose%
\pgfusepath{fill}%
\end{pgfscope}%
\begin{pgfscope}%
\pgfpathrectangle{\pgfqpoint{1.250000in}{0.660000in}}{\pgfqpoint{7.750000in}{4.620000in}}%
\pgfusepath{clip}%
\pgfsetbuttcap%
\pgfsetmiterjoin%
\definecolor{currentfill}{rgb}{0.000000,0.000000,0.000000}%
\pgfsetfillcolor{currentfill}%
\pgfsetlinewidth{0.000000pt}%
\definecolor{currentstroke}{rgb}{0.000000,0.000000,0.000000}%
\pgfsetstrokecolor{currentstroke}%
\pgfsetstrokeopacity{0.000000}%
\pgfsetdash{}{0pt}%
\pgfpathmoveto{\pgfqpoint{7.679400in}{0.660000in}}%
\pgfpathlineto{\pgfqpoint{7.685600in}{0.660000in}}%
\pgfpathlineto{\pgfqpoint{7.685600in}{3.206600in}}%
\pgfpathlineto{\pgfqpoint{7.679400in}{3.206600in}}%
\pgfpathlineto{\pgfqpoint{7.679400in}{0.660000in}}%
\pgfpathclose%
\pgfusepath{fill}%
\end{pgfscope}%
\begin{pgfscope}%
\pgfpathrectangle{\pgfqpoint{1.250000in}{0.660000in}}{\pgfqpoint{7.750000in}{4.620000in}}%
\pgfusepath{clip}%
\pgfsetbuttcap%
\pgfsetmiterjoin%
\definecolor{currentfill}{rgb}{0.000000,0.000000,0.000000}%
\pgfsetfillcolor{currentfill}%
\pgfsetlinewidth{0.000000pt}%
\definecolor{currentstroke}{rgb}{0.000000,0.000000,0.000000}%
\pgfsetstrokecolor{currentstroke}%
\pgfsetstrokeopacity{0.000000}%
\pgfsetdash{}{0pt}%
\pgfpathmoveto{\pgfqpoint{7.687150in}{0.660000in}}%
\pgfpathlineto{\pgfqpoint{7.693350in}{0.660000in}}%
\pgfpathlineto{\pgfqpoint{7.693350in}{3.206500in}}%
\pgfpathlineto{\pgfqpoint{7.687150in}{3.206500in}}%
\pgfpathlineto{\pgfqpoint{7.687150in}{0.660000in}}%
\pgfpathclose%
\pgfusepath{fill}%
\end{pgfscope}%
\begin{pgfscope}%
\pgfpathrectangle{\pgfqpoint{1.250000in}{0.660000in}}{\pgfqpoint{7.750000in}{4.620000in}}%
\pgfusepath{clip}%
\pgfsetbuttcap%
\pgfsetmiterjoin%
\definecolor{currentfill}{rgb}{0.000000,0.000000,0.000000}%
\pgfsetfillcolor{currentfill}%
\pgfsetlinewidth{0.000000pt}%
\definecolor{currentstroke}{rgb}{0.000000,0.000000,0.000000}%
\pgfsetstrokecolor{currentstroke}%
\pgfsetstrokeopacity{0.000000}%
\pgfsetdash{}{0pt}%
\pgfpathmoveto{\pgfqpoint{7.694900in}{0.660000in}}%
\pgfpathlineto{\pgfqpoint{7.701100in}{0.660000in}}%
\pgfpathlineto{\pgfqpoint{7.701100in}{3.206401in}}%
\pgfpathlineto{\pgfqpoint{7.694900in}{3.206401in}}%
\pgfpathlineto{\pgfqpoint{7.694900in}{0.660000in}}%
\pgfpathclose%
\pgfusepath{fill}%
\end{pgfscope}%
\begin{pgfscope}%
\pgfpathrectangle{\pgfqpoint{1.250000in}{0.660000in}}{\pgfqpoint{7.750000in}{4.620000in}}%
\pgfusepath{clip}%
\pgfsetbuttcap%
\pgfsetmiterjoin%
\definecolor{currentfill}{rgb}{0.000000,0.000000,0.000000}%
\pgfsetfillcolor{currentfill}%
\pgfsetlinewidth{0.000000pt}%
\definecolor{currentstroke}{rgb}{0.000000,0.000000,0.000000}%
\pgfsetstrokecolor{currentstroke}%
\pgfsetstrokeopacity{0.000000}%
\pgfsetdash{}{0pt}%
\pgfpathmoveto{\pgfqpoint{7.702650in}{0.660000in}}%
\pgfpathlineto{\pgfqpoint{7.708850in}{0.660000in}}%
\pgfpathlineto{\pgfqpoint{7.708850in}{3.206301in}}%
\pgfpathlineto{\pgfqpoint{7.702650in}{3.206301in}}%
\pgfpathlineto{\pgfqpoint{7.702650in}{0.660000in}}%
\pgfpathclose%
\pgfusepath{fill}%
\end{pgfscope}%
\begin{pgfscope}%
\pgfpathrectangle{\pgfqpoint{1.250000in}{0.660000in}}{\pgfqpoint{7.750000in}{4.620000in}}%
\pgfusepath{clip}%
\pgfsetbuttcap%
\pgfsetmiterjoin%
\definecolor{currentfill}{rgb}{0.000000,0.000000,0.000000}%
\pgfsetfillcolor{currentfill}%
\pgfsetlinewidth{0.000000pt}%
\definecolor{currentstroke}{rgb}{0.000000,0.000000,0.000000}%
\pgfsetstrokecolor{currentstroke}%
\pgfsetstrokeopacity{0.000000}%
\pgfsetdash{}{0pt}%
\pgfpathmoveto{\pgfqpoint{7.710400in}{0.660000in}}%
\pgfpathlineto{\pgfqpoint{7.716600in}{0.660000in}}%
\pgfpathlineto{\pgfqpoint{7.716600in}{3.206202in}}%
\pgfpathlineto{\pgfqpoint{7.710400in}{3.206202in}}%
\pgfpathlineto{\pgfqpoint{7.710400in}{0.660000in}}%
\pgfpathclose%
\pgfusepath{fill}%
\end{pgfscope}%
\begin{pgfscope}%
\pgfpathrectangle{\pgfqpoint{1.250000in}{0.660000in}}{\pgfqpoint{7.750000in}{4.620000in}}%
\pgfusepath{clip}%
\pgfsetbuttcap%
\pgfsetmiterjoin%
\definecolor{currentfill}{rgb}{0.000000,0.000000,0.000000}%
\pgfsetfillcolor{currentfill}%
\pgfsetlinewidth{0.000000pt}%
\definecolor{currentstroke}{rgb}{0.000000,0.000000,0.000000}%
\pgfsetstrokecolor{currentstroke}%
\pgfsetstrokeopacity{0.000000}%
\pgfsetdash{}{0pt}%
\pgfpathmoveto{\pgfqpoint{7.718150in}{0.660000in}}%
\pgfpathlineto{\pgfqpoint{7.724350in}{0.660000in}}%
\pgfpathlineto{\pgfqpoint{7.724350in}{3.206102in}}%
\pgfpathlineto{\pgfqpoint{7.718150in}{3.206102in}}%
\pgfpathlineto{\pgfqpoint{7.718150in}{0.660000in}}%
\pgfpathclose%
\pgfusepath{fill}%
\end{pgfscope}%
\begin{pgfscope}%
\pgfpathrectangle{\pgfqpoint{1.250000in}{0.660000in}}{\pgfqpoint{7.750000in}{4.620000in}}%
\pgfusepath{clip}%
\pgfsetbuttcap%
\pgfsetmiterjoin%
\definecolor{currentfill}{rgb}{0.000000,0.000000,0.000000}%
\pgfsetfillcolor{currentfill}%
\pgfsetlinewidth{0.000000pt}%
\definecolor{currentstroke}{rgb}{0.000000,0.000000,0.000000}%
\pgfsetstrokecolor{currentstroke}%
\pgfsetstrokeopacity{0.000000}%
\pgfsetdash{}{0pt}%
\pgfpathmoveto{\pgfqpoint{7.725900in}{0.660000in}}%
\pgfpathlineto{\pgfqpoint{7.732100in}{0.660000in}}%
\pgfpathlineto{\pgfqpoint{7.732100in}{3.206002in}}%
\pgfpathlineto{\pgfqpoint{7.725900in}{3.206002in}}%
\pgfpathlineto{\pgfqpoint{7.725900in}{0.660000in}}%
\pgfpathclose%
\pgfusepath{fill}%
\end{pgfscope}%
\begin{pgfscope}%
\pgfpathrectangle{\pgfqpoint{1.250000in}{0.660000in}}{\pgfqpoint{7.750000in}{4.620000in}}%
\pgfusepath{clip}%
\pgfsetbuttcap%
\pgfsetmiterjoin%
\definecolor{currentfill}{rgb}{0.000000,0.000000,0.000000}%
\pgfsetfillcolor{currentfill}%
\pgfsetlinewidth{0.000000pt}%
\definecolor{currentstroke}{rgb}{0.000000,0.000000,0.000000}%
\pgfsetstrokecolor{currentstroke}%
\pgfsetstrokeopacity{0.000000}%
\pgfsetdash{}{0pt}%
\pgfpathmoveto{\pgfqpoint{7.733650in}{0.660000in}}%
\pgfpathlineto{\pgfqpoint{7.739850in}{0.660000in}}%
\pgfpathlineto{\pgfqpoint{7.739850in}{3.205903in}}%
\pgfpathlineto{\pgfqpoint{7.733650in}{3.205903in}}%
\pgfpathlineto{\pgfqpoint{7.733650in}{0.660000in}}%
\pgfpathclose%
\pgfusepath{fill}%
\end{pgfscope}%
\begin{pgfscope}%
\pgfpathrectangle{\pgfqpoint{1.250000in}{0.660000in}}{\pgfqpoint{7.750000in}{4.620000in}}%
\pgfusepath{clip}%
\pgfsetbuttcap%
\pgfsetmiterjoin%
\definecolor{currentfill}{rgb}{0.000000,0.000000,0.000000}%
\pgfsetfillcolor{currentfill}%
\pgfsetlinewidth{0.000000pt}%
\definecolor{currentstroke}{rgb}{0.000000,0.000000,0.000000}%
\pgfsetstrokecolor{currentstroke}%
\pgfsetstrokeopacity{0.000000}%
\pgfsetdash{}{0pt}%
\pgfpathmoveto{\pgfqpoint{7.741400in}{0.660000in}}%
\pgfpathlineto{\pgfqpoint{7.747600in}{0.660000in}}%
\pgfpathlineto{\pgfqpoint{7.747600in}{3.205803in}}%
\pgfpathlineto{\pgfqpoint{7.741400in}{3.205803in}}%
\pgfpathlineto{\pgfqpoint{7.741400in}{0.660000in}}%
\pgfpathclose%
\pgfusepath{fill}%
\end{pgfscope}%
\begin{pgfscope}%
\pgfpathrectangle{\pgfqpoint{1.250000in}{0.660000in}}{\pgfqpoint{7.750000in}{4.620000in}}%
\pgfusepath{clip}%
\pgfsetbuttcap%
\pgfsetmiterjoin%
\definecolor{currentfill}{rgb}{0.000000,0.000000,0.000000}%
\pgfsetfillcolor{currentfill}%
\pgfsetlinewidth{0.000000pt}%
\definecolor{currentstroke}{rgb}{0.000000,0.000000,0.000000}%
\pgfsetstrokecolor{currentstroke}%
\pgfsetstrokeopacity{0.000000}%
\pgfsetdash{}{0pt}%
\pgfpathmoveto{\pgfqpoint{7.749150in}{0.660000in}}%
\pgfpathlineto{\pgfqpoint{7.755350in}{0.660000in}}%
\pgfpathlineto{\pgfqpoint{7.755350in}{3.205704in}}%
\pgfpathlineto{\pgfqpoint{7.749150in}{3.205704in}}%
\pgfpathlineto{\pgfqpoint{7.749150in}{0.660000in}}%
\pgfpathclose%
\pgfusepath{fill}%
\end{pgfscope}%
\begin{pgfscope}%
\pgfpathrectangle{\pgfqpoint{1.250000in}{0.660000in}}{\pgfqpoint{7.750000in}{4.620000in}}%
\pgfusepath{clip}%
\pgfsetbuttcap%
\pgfsetmiterjoin%
\definecolor{currentfill}{rgb}{0.000000,0.000000,0.000000}%
\pgfsetfillcolor{currentfill}%
\pgfsetlinewidth{0.000000pt}%
\definecolor{currentstroke}{rgb}{0.000000,0.000000,0.000000}%
\pgfsetstrokecolor{currentstroke}%
\pgfsetstrokeopacity{0.000000}%
\pgfsetdash{}{0pt}%
\pgfpathmoveto{\pgfqpoint{7.756900in}{0.660000in}}%
\pgfpathlineto{\pgfqpoint{7.763100in}{0.660000in}}%
\pgfpathlineto{\pgfqpoint{7.763100in}{3.205604in}}%
\pgfpathlineto{\pgfqpoint{7.756900in}{3.205604in}}%
\pgfpathlineto{\pgfqpoint{7.756900in}{0.660000in}}%
\pgfpathclose%
\pgfusepath{fill}%
\end{pgfscope}%
\begin{pgfscope}%
\pgfpathrectangle{\pgfqpoint{1.250000in}{0.660000in}}{\pgfqpoint{7.750000in}{4.620000in}}%
\pgfusepath{clip}%
\pgfsetbuttcap%
\pgfsetmiterjoin%
\definecolor{currentfill}{rgb}{0.000000,0.000000,0.000000}%
\pgfsetfillcolor{currentfill}%
\pgfsetlinewidth{0.000000pt}%
\definecolor{currentstroke}{rgb}{0.000000,0.000000,0.000000}%
\pgfsetstrokecolor{currentstroke}%
\pgfsetstrokeopacity{0.000000}%
\pgfsetdash{}{0pt}%
\pgfpathmoveto{\pgfqpoint{7.764650in}{0.660000in}}%
\pgfpathlineto{\pgfqpoint{7.770850in}{0.660000in}}%
\pgfpathlineto{\pgfqpoint{7.770850in}{3.205505in}}%
\pgfpathlineto{\pgfqpoint{7.764650in}{3.205505in}}%
\pgfpathlineto{\pgfqpoint{7.764650in}{0.660000in}}%
\pgfpathclose%
\pgfusepath{fill}%
\end{pgfscope}%
\begin{pgfscope}%
\pgfpathrectangle{\pgfqpoint{1.250000in}{0.660000in}}{\pgfqpoint{7.750000in}{4.620000in}}%
\pgfusepath{clip}%
\pgfsetbuttcap%
\pgfsetmiterjoin%
\definecolor{currentfill}{rgb}{0.000000,0.000000,0.000000}%
\pgfsetfillcolor{currentfill}%
\pgfsetlinewidth{0.000000pt}%
\definecolor{currentstroke}{rgb}{0.000000,0.000000,0.000000}%
\pgfsetstrokecolor{currentstroke}%
\pgfsetstrokeopacity{0.000000}%
\pgfsetdash{}{0pt}%
\pgfpathmoveto{\pgfqpoint{7.772400in}{0.660000in}}%
\pgfpathlineto{\pgfqpoint{7.778600in}{0.660000in}}%
\pgfpathlineto{\pgfqpoint{7.778600in}{3.205405in}}%
\pgfpathlineto{\pgfqpoint{7.772400in}{3.205405in}}%
\pgfpathlineto{\pgfqpoint{7.772400in}{0.660000in}}%
\pgfpathclose%
\pgfusepath{fill}%
\end{pgfscope}%
\begin{pgfscope}%
\pgfpathrectangle{\pgfqpoint{1.250000in}{0.660000in}}{\pgfqpoint{7.750000in}{4.620000in}}%
\pgfusepath{clip}%
\pgfsetbuttcap%
\pgfsetmiterjoin%
\definecolor{currentfill}{rgb}{0.000000,0.000000,0.000000}%
\pgfsetfillcolor{currentfill}%
\pgfsetlinewidth{0.000000pt}%
\definecolor{currentstroke}{rgb}{0.000000,0.000000,0.000000}%
\pgfsetstrokecolor{currentstroke}%
\pgfsetstrokeopacity{0.000000}%
\pgfsetdash{}{0pt}%
\pgfpathmoveto{\pgfqpoint{7.780150in}{0.660000in}}%
\pgfpathlineto{\pgfqpoint{7.786350in}{0.660000in}}%
\pgfpathlineto{\pgfqpoint{7.786350in}{3.205306in}}%
\pgfpathlineto{\pgfqpoint{7.780150in}{3.205306in}}%
\pgfpathlineto{\pgfqpoint{7.780150in}{0.660000in}}%
\pgfpathclose%
\pgfusepath{fill}%
\end{pgfscope}%
\begin{pgfscope}%
\pgfpathrectangle{\pgfqpoint{1.250000in}{0.660000in}}{\pgfqpoint{7.750000in}{4.620000in}}%
\pgfusepath{clip}%
\pgfsetbuttcap%
\pgfsetmiterjoin%
\definecolor{currentfill}{rgb}{0.000000,0.000000,0.000000}%
\pgfsetfillcolor{currentfill}%
\pgfsetlinewidth{0.000000pt}%
\definecolor{currentstroke}{rgb}{0.000000,0.000000,0.000000}%
\pgfsetstrokecolor{currentstroke}%
\pgfsetstrokeopacity{0.000000}%
\pgfsetdash{}{0pt}%
\pgfpathmoveto{\pgfqpoint{7.787900in}{0.660000in}}%
\pgfpathlineto{\pgfqpoint{7.794100in}{0.660000in}}%
\pgfpathlineto{\pgfqpoint{7.794100in}{3.205206in}}%
\pgfpathlineto{\pgfqpoint{7.787900in}{3.205206in}}%
\pgfpathlineto{\pgfqpoint{7.787900in}{0.660000in}}%
\pgfpathclose%
\pgfusepath{fill}%
\end{pgfscope}%
\begin{pgfscope}%
\pgfpathrectangle{\pgfqpoint{1.250000in}{0.660000in}}{\pgfqpoint{7.750000in}{4.620000in}}%
\pgfusepath{clip}%
\pgfsetbuttcap%
\pgfsetmiterjoin%
\definecolor{currentfill}{rgb}{0.000000,0.000000,0.000000}%
\pgfsetfillcolor{currentfill}%
\pgfsetlinewidth{0.000000pt}%
\definecolor{currentstroke}{rgb}{0.000000,0.000000,0.000000}%
\pgfsetstrokecolor{currentstroke}%
\pgfsetstrokeopacity{0.000000}%
\pgfsetdash{}{0pt}%
\pgfpathmoveto{\pgfqpoint{7.795650in}{0.660000in}}%
\pgfpathlineto{\pgfqpoint{7.801850in}{0.660000in}}%
\pgfpathlineto{\pgfqpoint{7.801850in}{3.205106in}}%
\pgfpathlineto{\pgfqpoint{7.795650in}{3.205106in}}%
\pgfpathlineto{\pgfqpoint{7.795650in}{0.660000in}}%
\pgfpathclose%
\pgfusepath{fill}%
\end{pgfscope}%
\begin{pgfscope}%
\pgfpathrectangle{\pgfqpoint{1.250000in}{0.660000in}}{\pgfqpoint{7.750000in}{4.620000in}}%
\pgfusepath{clip}%
\pgfsetbuttcap%
\pgfsetmiterjoin%
\definecolor{currentfill}{rgb}{0.000000,0.000000,0.000000}%
\pgfsetfillcolor{currentfill}%
\pgfsetlinewidth{0.000000pt}%
\definecolor{currentstroke}{rgb}{0.000000,0.000000,0.000000}%
\pgfsetstrokecolor{currentstroke}%
\pgfsetstrokeopacity{0.000000}%
\pgfsetdash{}{0pt}%
\pgfpathmoveto{\pgfqpoint{7.803400in}{0.660000in}}%
\pgfpathlineto{\pgfqpoint{7.809600in}{0.660000in}}%
\pgfpathlineto{\pgfqpoint{7.809600in}{3.205007in}}%
\pgfpathlineto{\pgfqpoint{7.803400in}{3.205007in}}%
\pgfpathlineto{\pgfqpoint{7.803400in}{0.660000in}}%
\pgfpathclose%
\pgfusepath{fill}%
\end{pgfscope}%
\begin{pgfscope}%
\pgfpathrectangle{\pgfqpoint{1.250000in}{0.660000in}}{\pgfqpoint{7.750000in}{4.620000in}}%
\pgfusepath{clip}%
\pgfsetbuttcap%
\pgfsetmiterjoin%
\definecolor{currentfill}{rgb}{0.000000,0.000000,0.000000}%
\pgfsetfillcolor{currentfill}%
\pgfsetlinewidth{0.000000pt}%
\definecolor{currentstroke}{rgb}{0.000000,0.000000,0.000000}%
\pgfsetstrokecolor{currentstroke}%
\pgfsetstrokeopacity{0.000000}%
\pgfsetdash{}{0pt}%
\pgfpathmoveto{\pgfqpoint{7.811150in}{0.660000in}}%
\pgfpathlineto{\pgfqpoint{7.817350in}{0.660000in}}%
\pgfpathlineto{\pgfqpoint{7.817350in}{3.204907in}}%
\pgfpathlineto{\pgfqpoint{7.811150in}{3.204907in}}%
\pgfpathlineto{\pgfqpoint{7.811150in}{0.660000in}}%
\pgfpathclose%
\pgfusepath{fill}%
\end{pgfscope}%
\begin{pgfscope}%
\pgfpathrectangle{\pgfqpoint{1.250000in}{0.660000in}}{\pgfqpoint{7.750000in}{4.620000in}}%
\pgfusepath{clip}%
\pgfsetbuttcap%
\pgfsetmiterjoin%
\definecolor{currentfill}{rgb}{0.000000,0.000000,0.000000}%
\pgfsetfillcolor{currentfill}%
\pgfsetlinewidth{0.000000pt}%
\definecolor{currentstroke}{rgb}{0.000000,0.000000,0.000000}%
\pgfsetstrokecolor{currentstroke}%
\pgfsetstrokeopacity{0.000000}%
\pgfsetdash{}{0pt}%
\pgfpathmoveto{\pgfqpoint{7.818900in}{0.660000in}}%
\pgfpathlineto{\pgfqpoint{7.825100in}{0.660000in}}%
\pgfpathlineto{\pgfqpoint{7.825100in}{3.204808in}}%
\pgfpathlineto{\pgfqpoint{7.818900in}{3.204808in}}%
\pgfpathlineto{\pgfqpoint{7.818900in}{0.660000in}}%
\pgfpathclose%
\pgfusepath{fill}%
\end{pgfscope}%
\begin{pgfscope}%
\pgfpathrectangle{\pgfqpoint{1.250000in}{0.660000in}}{\pgfqpoint{7.750000in}{4.620000in}}%
\pgfusepath{clip}%
\pgfsetbuttcap%
\pgfsetmiterjoin%
\definecolor{currentfill}{rgb}{0.000000,0.000000,0.000000}%
\pgfsetfillcolor{currentfill}%
\pgfsetlinewidth{0.000000pt}%
\definecolor{currentstroke}{rgb}{0.000000,0.000000,0.000000}%
\pgfsetstrokecolor{currentstroke}%
\pgfsetstrokeopacity{0.000000}%
\pgfsetdash{}{0pt}%
\pgfpathmoveto{\pgfqpoint{7.826650in}{0.660000in}}%
\pgfpathlineto{\pgfqpoint{7.832850in}{0.660000in}}%
\pgfpathlineto{\pgfqpoint{7.832850in}{3.204708in}}%
\pgfpathlineto{\pgfqpoint{7.826650in}{3.204708in}}%
\pgfpathlineto{\pgfqpoint{7.826650in}{0.660000in}}%
\pgfpathclose%
\pgfusepath{fill}%
\end{pgfscope}%
\begin{pgfscope}%
\pgfpathrectangle{\pgfqpoint{1.250000in}{0.660000in}}{\pgfqpoint{7.750000in}{4.620000in}}%
\pgfusepath{clip}%
\pgfsetbuttcap%
\pgfsetmiterjoin%
\definecolor{currentfill}{rgb}{0.000000,0.000000,0.000000}%
\pgfsetfillcolor{currentfill}%
\pgfsetlinewidth{0.000000pt}%
\definecolor{currentstroke}{rgb}{0.000000,0.000000,0.000000}%
\pgfsetstrokecolor{currentstroke}%
\pgfsetstrokeopacity{0.000000}%
\pgfsetdash{}{0pt}%
\pgfpathmoveto{\pgfqpoint{7.834400in}{0.660000in}}%
\pgfpathlineto{\pgfqpoint{7.840600in}{0.660000in}}%
\pgfpathlineto{\pgfqpoint{7.840600in}{3.204611in}}%
\pgfpathlineto{\pgfqpoint{7.834400in}{3.204611in}}%
\pgfpathlineto{\pgfqpoint{7.834400in}{0.660000in}}%
\pgfpathclose%
\pgfusepath{fill}%
\end{pgfscope}%
\begin{pgfscope}%
\pgfpathrectangle{\pgfqpoint{1.250000in}{0.660000in}}{\pgfqpoint{7.750000in}{4.620000in}}%
\pgfusepath{clip}%
\pgfsetbuttcap%
\pgfsetmiterjoin%
\definecolor{currentfill}{rgb}{0.000000,0.000000,0.000000}%
\pgfsetfillcolor{currentfill}%
\pgfsetlinewidth{0.000000pt}%
\definecolor{currentstroke}{rgb}{0.000000,0.000000,0.000000}%
\pgfsetstrokecolor{currentstroke}%
\pgfsetstrokeopacity{0.000000}%
\pgfsetdash{}{0pt}%
\pgfpathmoveto{\pgfqpoint{7.842150in}{0.660000in}}%
\pgfpathlineto{\pgfqpoint{7.848350in}{0.660000in}}%
\pgfpathlineto{\pgfqpoint{7.848350in}{3.204515in}}%
\pgfpathlineto{\pgfqpoint{7.842150in}{3.204515in}}%
\pgfpathlineto{\pgfqpoint{7.842150in}{0.660000in}}%
\pgfpathclose%
\pgfusepath{fill}%
\end{pgfscope}%
\begin{pgfscope}%
\pgfpathrectangle{\pgfqpoint{1.250000in}{0.660000in}}{\pgfqpoint{7.750000in}{4.620000in}}%
\pgfusepath{clip}%
\pgfsetbuttcap%
\pgfsetmiterjoin%
\definecolor{currentfill}{rgb}{0.000000,0.000000,0.000000}%
\pgfsetfillcolor{currentfill}%
\pgfsetlinewidth{0.000000pt}%
\definecolor{currentstroke}{rgb}{0.000000,0.000000,0.000000}%
\pgfsetstrokecolor{currentstroke}%
\pgfsetstrokeopacity{0.000000}%
\pgfsetdash{}{0pt}%
\pgfpathmoveto{\pgfqpoint{7.849900in}{0.660000in}}%
\pgfpathlineto{\pgfqpoint{7.856100in}{0.660000in}}%
\pgfpathlineto{\pgfqpoint{7.856100in}{3.204419in}}%
\pgfpathlineto{\pgfqpoint{7.849900in}{3.204419in}}%
\pgfpathlineto{\pgfqpoint{7.849900in}{0.660000in}}%
\pgfpathclose%
\pgfusepath{fill}%
\end{pgfscope}%
\begin{pgfscope}%
\pgfpathrectangle{\pgfqpoint{1.250000in}{0.660000in}}{\pgfqpoint{7.750000in}{4.620000in}}%
\pgfusepath{clip}%
\pgfsetbuttcap%
\pgfsetmiterjoin%
\definecolor{currentfill}{rgb}{0.000000,0.000000,0.000000}%
\pgfsetfillcolor{currentfill}%
\pgfsetlinewidth{0.000000pt}%
\definecolor{currentstroke}{rgb}{0.000000,0.000000,0.000000}%
\pgfsetstrokecolor{currentstroke}%
\pgfsetstrokeopacity{0.000000}%
\pgfsetdash{}{0pt}%
\pgfpathmoveto{\pgfqpoint{7.857650in}{0.660000in}}%
\pgfpathlineto{\pgfqpoint{7.863850in}{0.660000in}}%
\pgfpathlineto{\pgfqpoint{7.863850in}{3.204322in}}%
\pgfpathlineto{\pgfqpoint{7.857650in}{3.204322in}}%
\pgfpathlineto{\pgfqpoint{7.857650in}{0.660000in}}%
\pgfpathclose%
\pgfusepath{fill}%
\end{pgfscope}%
\begin{pgfscope}%
\pgfpathrectangle{\pgfqpoint{1.250000in}{0.660000in}}{\pgfqpoint{7.750000in}{4.620000in}}%
\pgfusepath{clip}%
\pgfsetbuttcap%
\pgfsetmiterjoin%
\definecolor{currentfill}{rgb}{0.000000,0.000000,0.000000}%
\pgfsetfillcolor{currentfill}%
\pgfsetlinewidth{0.000000pt}%
\definecolor{currentstroke}{rgb}{0.000000,0.000000,0.000000}%
\pgfsetstrokecolor{currentstroke}%
\pgfsetstrokeopacity{0.000000}%
\pgfsetdash{}{0pt}%
\pgfpathmoveto{\pgfqpoint{7.865400in}{0.660000in}}%
\pgfpathlineto{\pgfqpoint{7.871600in}{0.660000in}}%
\pgfpathlineto{\pgfqpoint{7.871600in}{3.204226in}}%
\pgfpathlineto{\pgfqpoint{7.865400in}{3.204226in}}%
\pgfpathlineto{\pgfqpoint{7.865400in}{0.660000in}}%
\pgfpathclose%
\pgfusepath{fill}%
\end{pgfscope}%
\begin{pgfscope}%
\pgfpathrectangle{\pgfqpoint{1.250000in}{0.660000in}}{\pgfqpoint{7.750000in}{4.620000in}}%
\pgfusepath{clip}%
\pgfsetbuttcap%
\pgfsetmiterjoin%
\definecolor{currentfill}{rgb}{0.000000,0.000000,0.000000}%
\pgfsetfillcolor{currentfill}%
\pgfsetlinewidth{0.000000pt}%
\definecolor{currentstroke}{rgb}{0.000000,0.000000,0.000000}%
\pgfsetstrokecolor{currentstroke}%
\pgfsetstrokeopacity{0.000000}%
\pgfsetdash{}{0pt}%
\pgfpathmoveto{\pgfqpoint{7.873150in}{0.660000in}}%
\pgfpathlineto{\pgfqpoint{7.879350in}{0.660000in}}%
\pgfpathlineto{\pgfqpoint{7.879350in}{3.204130in}}%
\pgfpathlineto{\pgfqpoint{7.873150in}{3.204130in}}%
\pgfpathlineto{\pgfqpoint{7.873150in}{0.660000in}}%
\pgfpathclose%
\pgfusepath{fill}%
\end{pgfscope}%
\begin{pgfscope}%
\pgfpathrectangle{\pgfqpoint{1.250000in}{0.660000in}}{\pgfqpoint{7.750000in}{4.620000in}}%
\pgfusepath{clip}%
\pgfsetbuttcap%
\pgfsetmiterjoin%
\definecolor{currentfill}{rgb}{0.000000,0.000000,0.000000}%
\pgfsetfillcolor{currentfill}%
\pgfsetlinewidth{0.000000pt}%
\definecolor{currentstroke}{rgb}{0.000000,0.000000,0.000000}%
\pgfsetstrokecolor{currentstroke}%
\pgfsetstrokeopacity{0.000000}%
\pgfsetdash{}{0pt}%
\pgfpathmoveto{\pgfqpoint{7.880900in}{0.660000in}}%
\pgfpathlineto{\pgfqpoint{7.887100in}{0.660000in}}%
\pgfpathlineto{\pgfqpoint{7.887100in}{3.204034in}}%
\pgfpathlineto{\pgfqpoint{7.880900in}{3.204034in}}%
\pgfpathlineto{\pgfqpoint{7.880900in}{0.660000in}}%
\pgfpathclose%
\pgfusepath{fill}%
\end{pgfscope}%
\begin{pgfscope}%
\pgfpathrectangle{\pgfqpoint{1.250000in}{0.660000in}}{\pgfqpoint{7.750000in}{4.620000in}}%
\pgfusepath{clip}%
\pgfsetbuttcap%
\pgfsetmiterjoin%
\definecolor{currentfill}{rgb}{0.000000,0.000000,0.000000}%
\pgfsetfillcolor{currentfill}%
\pgfsetlinewidth{0.000000pt}%
\definecolor{currentstroke}{rgb}{0.000000,0.000000,0.000000}%
\pgfsetstrokecolor{currentstroke}%
\pgfsetstrokeopacity{0.000000}%
\pgfsetdash{}{0pt}%
\pgfpathmoveto{\pgfqpoint{7.888650in}{0.660000in}}%
\pgfpathlineto{\pgfqpoint{7.894850in}{0.660000in}}%
\pgfpathlineto{\pgfqpoint{7.894850in}{3.203938in}}%
\pgfpathlineto{\pgfqpoint{7.888650in}{3.203938in}}%
\pgfpathlineto{\pgfqpoint{7.888650in}{0.660000in}}%
\pgfpathclose%
\pgfusepath{fill}%
\end{pgfscope}%
\begin{pgfscope}%
\pgfpathrectangle{\pgfqpoint{1.250000in}{0.660000in}}{\pgfqpoint{7.750000in}{4.620000in}}%
\pgfusepath{clip}%
\pgfsetbuttcap%
\pgfsetmiterjoin%
\definecolor{currentfill}{rgb}{0.000000,0.000000,0.000000}%
\pgfsetfillcolor{currentfill}%
\pgfsetlinewidth{0.000000pt}%
\definecolor{currentstroke}{rgb}{0.000000,0.000000,0.000000}%
\pgfsetstrokecolor{currentstroke}%
\pgfsetstrokeopacity{0.000000}%
\pgfsetdash{}{0pt}%
\pgfpathmoveto{\pgfqpoint{7.896400in}{0.660000in}}%
\pgfpathlineto{\pgfqpoint{7.902600in}{0.660000in}}%
\pgfpathlineto{\pgfqpoint{7.902600in}{3.203841in}}%
\pgfpathlineto{\pgfqpoint{7.896400in}{3.203841in}}%
\pgfpathlineto{\pgfqpoint{7.896400in}{0.660000in}}%
\pgfpathclose%
\pgfusepath{fill}%
\end{pgfscope}%
\begin{pgfscope}%
\pgfpathrectangle{\pgfqpoint{1.250000in}{0.660000in}}{\pgfqpoint{7.750000in}{4.620000in}}%
\pgfusepath{clip}%
\pgfsetbuttcap%
\pgfsetmiterjoin%
\definecolor{currentfill}{rgb}{0.000000,0.000000,0.000000}%
\pgfsetfillcolor{currentfill}%
\pgfsetlinewidth{0.000000pt}%
\definecolor{currentstroke}{rgb}{0.000000,0.000000,0.000000}%
\pgfsetstrokecolor{currentstroke}%
\pgfsetstrokeopacity{0.000000}%
\pgfsetdash{}{0pt}%
\pgfpathmoveto{\pgfqpoint{7.904150in}{0.660000in}}%
\pgfpathlineto{\pgfqpoint{7.910350in}{0.660000in}}%
\pgfpathlineto{\pgfqpoint{7.910350in}{3.203745in}}%
\pgfpathlineto{\pgfqpoint{7.904150in}{3.203745in}}%
\pgfpathlineto{\pgfqpoint{7.904150in}{0.660000in}}%
\pgfpathclose%
\pgfusepath{fill}%
\end{pgfscope}%
\begin{pgfscope}%
\pgfpathrectangle{\pgfqpoint{1.250000in}{0.660000in}}{\pgfqpoint{7.750000in}{4.620000in}}%
\pgfusepath{clip}%
\pgfsetbuttcap%
\pgfsetmiterjoin%
\definecolor{currentfill}{rgb}{0.000000,0.000000,0.000000}%
\pgfsetfillcolor{currentfill}%
\pgfsetlinewidth{0.000000pt}%
\definecolor{currentstroke}{rgb}{0.000000,0.000000,0.000000}%
\pgfsetstrokecolor{currentstroke}%
\pgfsetstrokeopacity{0.000000}%
\pgfsetdash{}{0pt}%
\pgfpathmoveto{\pgfqpoint{7.911900in}{0.660000in}}%
\pgfpathlineto{\pgfqpoint{7.918100in}{0.660000in}}%
\pgfpathlineto{\pgfqpoint{7.918100in}{3.203649in}}%
\pgfpathlineto{\pgfqpoint{7.911900in}{3.203649in}}%
\pgfpathlineto{\pgfqpoint{7.911900in}{0.660000in}}%
\pgfpathclose%
\pgfusepath{fill}%
\end{pgfscope}%
\begin{pgfscope}%
\pgfpathrectangle{\pgfqpoint{1.250000in}{0.660000in}}{\pgfqpoint{7.750000in}{4.620000in}}%
\pgfusepath{clip}%
\pgfsetbuttcap%
\pgfsetmiterjoin%
\definecolor{currentfill}{rgb}{0.000000,0.000000,0.000000}%
\pgfsetfillcolor{currentfill}%
\pgfsetlinewidth{0.000000pt}%
\definecolor{currentstroke}{rgb}{0.000000,0.000000,0.000000}%
\pgfsetstrokecolor{currentstroke}%
\pgfsetstrokeopacity{0.000000}%
\pgfsetdash{}{0pt}%
\pgfpathmoveto{\pgfqpoint{7.919650in}{0.660000in}}%
\pgfpathlineto{\pgfqpoint{7.925850in}{0.660000in}}%
\pgfpathlineto{\pgfqpoint{7.925850in}{3.203553in}}%
\pgfpathlineto{\pgfqpoint{7.919650in}{3.203553in}}%
\pgfpathlineto{\pgfqpoint{7.919650in}{0.660000in}}%
\pgfpathclose%
\pgfusepath{fill}%
\end{pgfscope}%
\begin{pgfscope}%
\pgfpathrectangle{\pgfqpoint{1.250000in}{0.660000in}}{\pgfqpoint{7.750000in}{4.620000in}}%
\pgfusepath{clip}%
\pgfsetbuttcap%
\pgfsetmiterjoin%
\definecolor{currentfill}{rgb}{0.000000,0.000000,0.000000}%
\pgfsetfillcolor{currentfill}%
\pgfsetlinewidth{0.000000pt}%
\definecolor{currentstroke}{rgb}{0.000000,0.000000,0.000000}%
\pgfsetstrokecolor{currentstroke}%
\pgfsetstrokeopacity{0.000000}%
\pgfsetdash{}{0pt}%
\pgfpathmoveto{\pgfqpoint{7.927400in}{0.660000in}}%
\pgfpathlineto{\pgfqpoint{7.933600in}{0.660000in}}%
\pgfpathlineto{\pgfqpoint{7.933600in}{3.203457in}}%
\pgfpathlineto{\pgfqpoint{7.927400in}{3.203457in}}%
\pgfpathlineto{\pgfqpoint{7.927400in}{0.660000in}}%
\pgfpathclose%
\pgfusepath{fill}%
\end{pgfscope}%
\begin{pgfscope}%
\pgfpathrectangle{\pgfqpoint{1.250000in}{0.660000in}}{\pgfqpoint{7.750000in}{4.620000in}}%
\pgfusepath{clip}%
\pgfsetbuttcap%
\pgfsetmiterjoin%
\definecolor{currentfill}{rgb}{0.000000,0.000000,0.000000}%
\pgfsetfillcolor{currentfill}%
\pgfsetlinewidth{0.000000pt}%
\definecolor{currentstroke}{rgb}{0.000000,0.000000,0.000000}%
\pgfsetstrokecolor{currentstroke}%
\pgfsetstrokeopacity{0.000000}%
\pgfsetdash{}{0pt}%
\pgfpathmoveto{\pgfqpoint{7.935150in}{0.660000in}}%
\pgfpathlineto{\pgfqpoint{7.941350in}{0.660000in}}%
\pgfpathlineto{\pgfqpoint{7.941350in}{3.203360in}}%
\pgfpathlineto{\pgfqpoint{7.935150in}{3.203360in}}%
\pgfpathlineto{\pgfqpoint{7.935150in}{0.660000in}}%
\pgfpathclose%
\pgfusepath{fill}%
\end{pgfscope}%
\begin{pgfscope}%
\pgfpathrectangle{\pgfqpoint{1.250000in}{0.660000in}}{\pgfqpoint{7.750000in}{4.620000in}}%
\pgfusepath{clip}%
\pgfsetbuttcap%
\pgfsetmiterjoin%
\definecolor{currentfill}{rgb}{0.000000,0.000000,0.000000}%
\pgfsetfillcolor{currentfill}%
\pgfsetlinewidth{0.000000pt}%
\definecolor{currentstroke}{rgb}{0.000000,0.000000,0.000000}%
\pgfsetstrokecolor{currentstroke}%
\pgfsetstrokeopacity{0.000000}%
\pgfsetdash{}{0pt}%
\pgfpathmoveto{\pgfqpoint{7.942900in}{0.660000in}}%
\pgfpathlineto{\pgfqpoint{7.949100in}{0.660000in}}%
\pgfpathlineto{\pgfqpoint{7.949100in}{3.203264in}}%
\pgfpathlineto{\pgfqpoint{7.942900in}{3.203264in}}%
\pgfpathlineto{\pgfqpoint{7.942900in}{0.660000in}}%
\pgfpathclose%
\pgfusepath{fill}%
\end{pgfscope}%
\begin{pgfscope}%
\pgfpathrectangle{\pgfqpoint{1.250000in}{0.660000in}}{\pgfqpoint{7.750000in}{4.620000in}}%
\pgfusepath{clip}%
\pgfsetbuttcap%
\pgfsetmiterjoin%
\definecolor{currentfill}{rgb}{0.000000,0.000000,0.000000}%
\pgfsetfillcolor{currentfill}%
\pgfsetlinewidth{0.000000pt}%
\definecolor{currentstroke}{rgb}{0.000000,0.000000,0.000000}%
\pgfsetstrokecolor{currentstroke}%
\pgfsetstrokeopacity{0.000000}%
\pgfsetdash{}{0pt}%
\pgfpathmoveto{\pgfqpoint{7.950650in}{0.660000in}}%
\pgfpathlineto{\pgfqpoint{7.956850in}{0.660000in}}%
\pgfpathlineto{\pgfqpoint{7.956850in}{3.203168in}}%
\pgfpathlineto{\pgfqpoint{7.950650in}{3.203168in}}%
\pgfpathlineto{\pgfqpoint{7.950650in}{0.660000in}}%
\pgfpathclose%
\pgfusepath{fill}%
\end{pgfscope}%
\begin{pgfscope}%
\pgfpathrectangle{\pgfqpoint{1.250000in}{0.660000in}}{\pgfqpoint{7.750000in}{4.620000in}}%
\pgfusepath{clip}%
\pgfsetbuttcap%
\pgfsetmiterjoin%
\definecolor{currentfill}{rgb}{0.000000,0.000000,0.000000}%
\pgfsetfillcolor{currentfill}%
\pgfsetlinewidth{0.000000pt}%
\definecolor{currentstroke}{rgb}{0.000000,0.000000,0.000000}%
\pgfsetstrokecolor{currentstroke}%
\pgfsetstrokeopacity{0.000000}%
\pgfsetdash{}{0pt}%
\pgfpathmoveto{\pgfqpoint{7.958400in}{0.660000in}}%
\pgfpathlineto{\pgfqpoint{7.964600in}{0.660000in}}%
\pgfpathlineto{\pgfqpoint{7.964600in}{3.203072in}}%
\pgfpathlineto{\pgfqpoint{7.958400in}{3.203072in}}%
\pgfpathlineto{\pgfqpoint{7.958400in}{0.660000in}}%
\pgfpathclose%
\pgfusepath{fill}%
\end{pgfscope}%
\begin{pgfscope}%
\pgfpathrectangle{\pgfqpoint{1.250000in}{0.660000in}}{\pgfqpoint{7.750000in}{4.620000in}}%
\pgfusepath{clip}%
\pgfsetbuttcap%
\pgfsetmiterjoin%
\definecolor{currentfill}{rgb}{0.000000,0.000000,0.000000}%
\pgfsetfillcolor{currentfill}%
\pgfsetlinewidth{0.000000pt}%
\definecolor{currentstroke}{rgb}{0.000000,0.000000,0.000000}%
\pgfsetstrokecolor{currentstroke}%
\pgfsetstrokeopacity{0.000000}%
\pgfsetdash{}{0pt}%
\pgfpathmoveto{\pgfqpoint{7.966150in}{0.660000in}}%
\pgfpathlineto{\pgfqpoint{7.972350in}{0.660000in}}%
\pgfpathlineto{\pgfqpoint{7.972350in}{3.202975in}}%
\pgfpathlineto{\pgfqpoint{7.966150in}{3.202975in}}%
\pgfpathlineto{\pgfqpoint{7.966150in}{0.660000in}}%
\pgfpathclose%
\pgfusepath{fill}%
\end{pgfscope}%
\begin{pgfscope}%
\pgfpathrectangle{\pgfqpoint{1.250000in}{0.660000in}}{\pgfqpoint{7.750000in}{4.620000in}}%
\pgfusepath{clip}%
\pgfsetbuttcap%
\pgfsetmiterjoin%
\definecolor{currentfill}{rgb}{0.000000,0.000000,0.000000}%
\pgfsetfillcolor{currentfill}%
\pgfsetlinewidth{0.000000pt}%
\definecolor{currentstroke}{rgb}{0.000000,0.000000,0.000000}%
\pgfsetstrokecolor{currentstroke}%
\pgfsetstrokeopacity{0.000000}%
\pgfsetdash{}{0pt}%
\pgfpathmoveto{\pgfqpoint{7.973900in}{0.660000in}}%
\pgfpathlineto{\pgfqpoint{7.980100in}{0.660000in}}%
\pgfpathlineto{\pgfqpoint{7.980100in}{3.202879in}}%
\pgfpathlineto{\pgfqpoint{7.973900in}{3.202879in}}%
\pgfpathlineto{\pgfqpoint{7.973900in}{0.660000in}}%
\pgfpathclose%
\pgfusepath{fill}%
\end{pgfscope}%
\begin{pgfscope}%
\pgfpathrectangle{\pgfqpoint{1.250000in}{0.660000in}}{\pgfqpoint{7.750000in}{4.620000in}}%
\pgfusepath{clip}%
\pgfsetbuttcap%
\pgfsetmiterjoin%
\definecolor{currentfill}{rgb}{0.000000,0.000000,0.000000}%
\pgfsetfillcolor{currentfill}%
\pgfsetlinewidth{0.000000pt}%
\definecolor{currentstroke}{rgb}{0.000000,0.000000,0.000000}%
\pgfsetstrokecolor{currentstroke}%
\pgfsetstrokeopacity{0.000000}%
\pgfsetdash{}{0pt}%
\pgfpathmoveto{\pgfqpoint{7.981650in}{0.660000in}}%
\pgfpathlineto{\pgfqpoint{7.987850in}{0.660000in}}%
\pgfpathlineto{\pgfqpoint{7.987850in}{3.202783in}}%
\pgfpathlineto{\pgfqpoint{7.981650in}{3.202783in}}%
\pgfpathlineto{\pgfqpoint{7.981650in}{0.660000in}}%
\pgfpathclose%
\pgfusepath{fill}%
\end{pgfscope}%
\begin{pgfscope}%
\pgfpathrectangle{\pgfqpoint{1.250000in}{0.660000in}}{\pgfqpoint{7.750000in}{4.620000in}}%
\pgfusepath{clip}%
\pgfsetbuttcap%
\pgfsetmiterjoin%
\definecolor{currentfill}{rgb}{0.000000,0.000000,0.000000}%
\pgfsetfillcolor{currentfill}%
\pgfsetlinewidth{0.000000pt}%
\definecolor{currentstroke}{rgb}{0.000000,0.000000,0.000000}%
\pgfsetstrokecolor{currentstroke}%
\pgfsetstrokeopacity{0.000000}%
\pgfsetdash{}{0pt}%
\pgfpathmoveto{\pgfqpoint{7.989400in}{0.660000in}}%
\pgfpathlineto{\pgfqpoint{7.995600in}{0.660000in}}%
\pgfpathlineto{\pgfqpoint{7.995600in}{3.202687in}}%
\pgfpathlineto{\pgfqpoint{7.989400in}{3.202687in}}%
\pgfpathlineto{\pgfqpoint{7.989400in}{0.660000in}}%
\pgfpathclose%
\pgfusepath{fill}%
\end{pgfscope}%
\begin{pgfscope}%
\pgfpathrectangle{\pgfqpoint{1.250000in}{0.660000in}}{\pgfqpoint{7.750000in}{4.620000in}}%
\pgfusepath{clip}%
\pgfsetbuttcap%
\pgfsetmiterjoin%
\definecolor{currentfill}{rgb}{0.000000,0.000000,0.000000}%
\pgfsetfillcolor{currentfill}%
\pgfsetlinewidth{0.000000pt}%
\definecolor{currentstroke}{rgb}{0.000000,0.000000,0.000000}%
\pgfsetstrokecolor{currentstroke}%
\pgfsetstrokeopacity{0.000000}%
\pgfsetdash{}{0pt}%
\pgfpathmoveto{\pgfqpoint{7.997150in}{0.660000in}}%
\pgfpathlineto{\pgfqpoint{8.003350in}{0.660000in}}%
\pgfpathlineto{\pgfqpoint{8.003350in}{3.202591in}}%
\pgfpathlineto{\pgfqpoint{7.997150in}{3.202591in}}%
\pgfpathlineto{\pgfqpoint{7.997150in}{0.660000in}}%
\pgfpathclose%
\pgfusepath{fill}%
\end{pgfscope}%
\begin{pgfscope}%
\pgfpathrectangle{\pgfqpoint{1.250000in}{0.660000in}}{\pgfqpoint{7.750000in}{4.620000in}}%
\pgfusepath{clip}%
\pgfsetbuttcap%
\pgfsetmiterjoin%
\definecolor{currentfill}{rgb}{0.000000,0.000000,0.000000}%
\pgfsetfillcolor{currentfill}%
\pgfsetlinewidth{0.000000pt}%
\definecolor{currentstroke}{rgb}{0.000000,0.000000,0.000000}%
\pgfsetstrokecolor{currentstroke}%
\pgfsetstrokeopacity{0.000000}%
\pgfsetdash{}{0pt}%
\pgfpathmoveto{\pgfqpoint{8.004900in}{0.660000in}}%
\pgfpathlineto{\pgfqpoint{8.011100in}{0.660000in}}%
\pgfpathlineto{\pgfqpoint{8.011100in}{3.202498in}}%
\pgfpathlineto{\pgfqpoint{8.004900in}{3.202498in}}%
\pgfpathlineto{\pgfqpoint{8.004900in}{0.660000in}}%
\pgfpathclose%
\pgfusepath{fill}%
\end{pgfscope}%
\begin{pgfscope}%
\pgfpathrectangle{\pgfqpoint{1.250000in}{0.660000in}}{\pgfqpoint{7.750000in}{4.620000in}}%
\pgfusepath{clip}%
\pgfsetbuttcap%
\pgfsetmiterjoin%
\definecolor{currentfill}{rgb}{0.000000,0.000000,0.000000}%
\pgfsetfillcolor{currentfill}%
\pgfsetlinewidth{0.000000pt}%
\definecolor{currentstroke}{rgb}{0.000000,0.000000,0.000000}%
\pgfsetstrokecolor{currentstroke}%
\pgfsetstrokeopacity{0.000000}%
\pgfsetdash{}{0pt}%
\pgfpathmoveto{\pgfqpoint{8.012650in}{0.660000in}}%
\pgfpathlineto{\pgfqpoint{8.018850in}{0.660000in}}%
\pgfpathlineto{\pgfqpoint{8.018850in}{3.202405in}}%
\pgfpathlineto{\pgfqpoint{8.012650in}{3.202405in}}%
\pgfpathlineto{\pgfqpoint{8.012650in}{0.660000in}}%
\pgfpathclose%
\pgfusepath{fill}%
\end{pgfscope}%
\begin{pgfscope}%
\pgfpathrectangle{\pgfqpoint{1.250000in}{0.660000in}}{\pgfqpoint{7.750000in}{4.620000in}}%
\pgfusepath{clip}%
\pgfsetbuttcap%
\pgfsetmiterjoin%
\definecolor{currentfill}{rgb}{0.000000,0.000000,0.000000}%
\pgfsetfillcolor{currentfill}%
\pgfsetlinewidth{0.000000pt}%
\definecolor{currentstroke}{rgb}{0.000000,0.000000,0.000000}%
\pgfsetstrokecolor{currentstroke}%
\pgfsetstrokeopacity{0.000000}%
\pgfsetdash{}{0pt}%
\pgfpathmoveto{\pgfqpoint{8.020400in}{0.660000in}}%
\pgfpathlineto{\pgfqpoint{8.026600in}{0.660000in}}%
\pgfpathlineto{\pgfqpoint{8.026600in}{3.202312in}}%
\pgfpathlineto{\pgfqpoint{8.020400in}{3.202312in}}%
\pgfpathlineto{\pgfqpoint{8.020400in}{0.660000in}}%
\pgfpathclose%
\pgfusepath{fill}%
\end{pgfscope}%
\begin{pgfscope}%
\pgfpathrectangle{\pgfqpoint{1.250000in}{0.660000in}}{\pgfqpoint{7.750000in}{4.620000in}}%
\pgfusepath{clip}%
\pgfsetbuttcap%
\pgfsetmiterjoin%
\definecolor{currentfill}{rgb}{0.000000,0.000000,0.000000}%
\pgfsetfillcolor{currentfill}%
\pgfsetlinewidth{0.000000pt}%
\definecolor{currentstroke}{rgb}{0.000000,0.000000,0.000000}%
\pgfsetstrokecolor{currentstroke}%
\pgfsetstrokeopacity{0.000000}%
\pgfsetdash{}{0pt}%
\pgfpathmoveto{\pgfqpoint{8.028150in}{0.660000in}}%
\pgfpathlineto{\pgfqpoint{8.034350in}{0.660000in}}%
\pgfpathlineto{\pgfqpoint{8.034350in}{3.202219in}}%
\pgfpathlineto{\pgfqpoint{8.028150in}{3.202219in}}%
\pgfpathlineto{\pgfqpoint{8.028150in}{0.660000in}}%
\pgfpathclose%
\pgfusepath{fill}%
\end{pgfscope}%
\begin{pgfscope}%
\pgfpathrectangle{\pgfqpoint{1.250000in}{0.660000in}}{\pgfqpoint{7.750000in}{4.620000in}}%
\pgfusepath{clip}%
\pgfsetbuttcap%
\pgfsetmiterjoin%
\definecolor{currentfill}{rgb}{0.000000,0.000000,0.000000}%
\pgfsetfillcolor{currentfill}%
\pgfsetlinewidth{0.000000pt}%
\definecolor{currentstroke}{rgb}{0.000000,0.000000,0.000000}%
\pgfsetstrokecolor{currentstroke}%
\pgfsetstrokeopacity{0.000000}%
\pgfsetdash{}{0pt}%
\pgfpathmoveto{\pgfqpoint{8.035900in}{0.660000in}}%
\pgfpathlineto{\pgfqpoint{8.042100in}{0.660000in}}%
\pgfpathlineto{\pgfqpoint{8.042100in}{3.202126in}}%
\pgfpathlineto{\pgfqpoint{8.035900in}{3.202126in}}%
\pgfpathlineto{\pgfqpoint{8.035900in}{0.660000in}}%
\pgfpathclose%
\pgfusepath{fill}%
\end{pgfscope}%
\begin{pgfscope}%
\pgfpathrectangle{\pgfqpoint{1.250000in}{0.660000in}}{\pgfqpoint{7.750000in}{4.620000in}}%
\pgfusepath{clip}%
\pgfsetbuttcap%
\pgfsetmiterjoin%
\definecolor{currentfill}{rgb}{0.000000,0.000000,0.000000}%
\pgfsetfillcolor{currentfill}%
\pgfsetlinewidth{0.000000pt}%
\definecolor{currentstroke}{rgb}{0.000000,0.000000,0.000000}%
\pgfsetstrokecolor{currentstroke}%
\pgfsetstrokeopacity{0.000000}%
\pgfsetdash{}{0pt}%
\pgfpathmoveto{\pgfqpoint{8.043650in}{0.660000in}}%
\pgfpathlineto{\pgfqpoint{8.049850in}{0.660000in}}%
\pgfpathlineto{\pgfqpoint{8.049850in}{3.202033in}}%
\pgfpathlineto{\pgfqpoint{8.043650in}{3.202033in}}%
\pgfpathlineto{\pgfqpoint{8.043650in}{0.660000in}}%
\pgfpathclose%
\pgfusepath{fill}%
\end{pgfscope}%
\begin{pgfscope}%
\pgfpathrectangle{\pgfqpoint{1.250000in}{0.660000in}}{\pgfqpoint{7.750000in}{4.620000in}}%
\pgfusepath{clip}%
\pgfsetbuttcap%
\pgfsetmiterjoin%
\definecolor{currentfill}{rgb}{0.000000,0.000000,0.000000}%
\pgfsetfillcolor{currentfill}%
\pgfsetlinewidth{0.000000pt}%
\definecolor{currentstroke}{rgb}{0.000000,0.000000,0.000000}%
\pgfsetstrokecolor{currentstroke}%
\pgfsetstrokeopacity{0.000000}%
\pgfsetdash{}{0pt}%
\pgfpathmoveto{\pgfqpoint{8.051400in}{0.660000in}}%
\pgfpathlineto{\pgfqpoint{8.057600in}{0.660000in}}%
\pgfpathlineto{\pgfqpoint{8.057600in}{3.201940in}}%
\pgfpathlineto{\pgfqpoint{8.051400in}{3.201940in}}%
\pgfpathlineto{\pgfqpoint{8.051400in}{0.660000in}}%
\pgfpathclose%
\pgfusepath{fill}%
\end{pgfscope}%
\begin{pgfscope}%
\pgfpathrectangle{\pgfqpoint{1.250000in}{0.660000in}}{\pgfqpoint{7.750000in}{4.620000in}}%
\pgfusepath{clip}%
\pgfsetbuttcap%
\pgfsetmiterjoin%
\definecolor{currentfill}{rgb}{0.000000,0.000000,0.000000}%
\pgfsetfillcolor{currentfill}%
\pgfsetlinewidth{0.000000pt}%
\definecolor{currentstroke}{rgb}{0.000000,0.000000,0.000000}%
\pgfsetstrokecolor{currentstroke}%
\pgfsetstrokeopacity{0.000000}%
\pgfsetdash{}{0pt}%
\pgfpathmoveto{\pgfqpoint{8.059150in}{0.660000in}}%
\pgfpathlineto{\pgfqpoint{8.065350in}{0.660000in}}%
\pgfpathlineto{\pgfqpoint{8.065350in}{3.201848in}}%
\pgfpathlineto{\pgfqpoint{8.059150in}{3.201848in}}%
\pgfpathlineto{\pgfqpoint{8.059150in}{0.660000in}}%
\pgfpathclose%
\pgfusepath{fill}%
\end{pgfscope}%
\begin{pgfscope}%
\pgfpathrectangle{\pgfqpoint{1.250000in}{0.660000in}}{\pgfqpoint{7.750000in}{4.620000in}}%
\pgfusepath{clip}%
\pgfsetbuttcap%
\pgfsetmiterjoin%
\definecolor{currentfill}{rgb}{0.000000,0.000000,0.000000}%
\pgfsetfillcolor{currentfill}%
\pgfsetlinewidth{0.000000pt}%
\definecolor{currentstroke}{rgb}{0.000000,0.000000,0.000000}%
\pgfsetstrokecolor{currentstroke}%
\pgfsetstrokeopacity{0.000000}%
\pgfsetdash{}{0pt}%
\pgfpathmoveto{\pgfqpoint{8.066900in}{0.660000in}}%
\pgfpathlineto{\pgfqpoint{8.073100in}{0.660000in}}%
\pgfpathlineto{\pgfqpoint{8.073100in}{3.201755in}}%
\pgfpathlineto{\pgfqpoint{8.066900in}{3.201755in}}%
\pgfpathlineto{\pgfqpoint{8.066900in}{0.660000in}}%
\pgfpathclose%
\pgfusepath{fill}%
\end{pgfscope}%
\begin{pgfscope}%
\pgfpathrectangle{\pgfqpoint{1.250000in}{0.660000in}}{\pgfqpoint{7.750000in}{4.620000in}}%
\pgfusepath{clip}%
\pgfsetbuttcap%
\pgfsetmiterjoin%
\definecolor{currentfill}{rgb}{0.000000,0.000000,0.000000}%
\pgfsetfillcolor{currentfill}%
\pgfsetlinewidth{0.000000pt}%
\definecolor{currentstroke}{rgb}{0.000000,0.000000,0.000000}%
\pgfsetstrokecolor{currentstroke}%
\pgfsetstrokeopacity{0.000000}%
\pgfsetdash{}{0pt}%
\pgfpathmoveto{\pgfqpoint{8.074650in}{0.660000in}}%
\pgfpathlineto{\pgfqpoint{8.080850in}{0.660000in}}%
\pgfpathlineto{\pgfqpoint{8.080850in}{3.201662in}}%
\pgfpathlineto{\pgfqpoint{8.074650in}{3.201662in}}%
\pgfpathlineto{\pgfqpoint{8.074650in}{0.660000in}}%
\pgfpathclose%
\pgfusepath{fill}%
\end{pgfscope}%
\begin{pgfscope}%
\pgfpathrectangle{\pgfqpoint{1.250000in}{0.660000in}}{\pgfqpoint{7.750000in}{4.620000in}}%
\pgfusepath{clip}%
\pgfsetbuttcap%
\pgfsetmiterjoin%
\definecolor{currentfill}{rgb}{0.000000,0.000000,0.000000}%
\pgfsetfillcolor{currentfill}%
\pgfsetlinewidth{0.000000pt}%
\definecolor{currentstroke}{rgb}{0.000000,0.000000,0.000000}%
\pgfsetstrokecolor{currentstroke}%
\pgfsetstrokeopacity{0.000000}%
\pgfsetdash{}{0pt}%
\pgfpathmoveto{\pgfqpoint{8.082400in}{0.660000in}}%
\pgfpathlineto{\pgfqpoint{8.088600in}{0.660000in}}%
\pgfpathlineto{\pgfqpoint{8.088600in}{3.201569in}}%
\pgfpathlineto{\pgfqpoint{8.082400in}{3.201569in}}%
\pgfpathlineto{\pgfqpoint{8.082400in}{0.660000in}}%
\pgfpathclose%
\pgfusepath{fill}%
\end{pgfscope}%
\begin{pgfscope}%
\pgfpathrectangle{\pgfqpoint{1.250000in}{0.660000in}}{\pgfqpoint{7.750000in}{4.620000in}}%
\pgfusepath{clip}%
\pgfsetbuttcap%
\pgfsetmiterjoin%
\definecolor{currentfill}{rgb}{0.000000,0.000000,0.000000}%
\pgfsetfillcolor{currentfill}%
\pgfsetlinewidth{0.000000pt}%
\definecolor{currentstroke}{rgb}{0.000000,0.000000,0.000000}%
\pgfsetstrokecolor{currentstroke}%
\pgfsetstrokeopacity{0.000000}%
\pgfsetdash{}{0pt}%
\pgfpathmoveto{\pgfqpoint{8.090150in}{0.660000in}}%
\pgfpathlineto{\pgfqpoint{8.096350in}{0.660000in}}%
\pgfpathlineto{\pgfqpoint{8.096350in}{3.201476in}}%
\pgfpathlineto{\pgfqpoint{8.090150in}{3.201476in}}%
\pgfpathlineto{\pgfqpoint{8.090150in}{0.660000in}}%
\pgfpathclose%
\pgfusepath{fill}%
\end{pgfscope}%
\begin{pgfscope}%
\pgfpathrectangle{\pgfqpoint{1.250000in}{0.660000in}}{\pgfqpoint{7.750000in}{4.620000in}}%
\pgfusepath{clip}%
\pgfsetbuttcap%
\pgfsetmiterjoin%
\definecolor{currentfill}{rgb}{0.000000,0.000000,0.000000}%
\pgfsetfillcolor{currentfill}%
\pgfsetlinewidth{0.000000pt}%
\definecolor{currentstroke}{rgb}{0.000000,0.000000,0.000000}%
\pgfsetstrokecolor{currentstroke}%
\pgfsetstrokeopacity{0.000000}%
\pgfsetdash{}{0pt}%
\pgfpathmoveto{\pgfqpoint{8.097900in}{0.660000in}}%
\pgfpathlineto{\pgfqpoint{8.104100in}{0.660000in}}%
\pgfpathlineto{\pgfqpoint{8.104100in}{3.201383in}}%
\pgfpathlineto{\pgfqpoint{8.097900in}{3.201383in}}%
\pgfpathlineto{\pgfqpoint{8.097900in}{0.660000in}}%
\pgfpathclose%
\pgfusepath{fill}%
\end{pgfscope}%
\begin{pgfscope}%
\pgfpathrectangle{\pgfqpoint{1.250000in}{0.660000in}}{\pgfqpoint{7.750000in}{4.620000in}}%
\pgfusepath{clip}%
\pgfsetbuttcap%
\pgfsetmiterjoin%
\definecolor{currentfill}{rgb}{0.000000,0.000000,0.000000}%
\pgfsetfillcolor{currentfill}%
\pgfsetlinewidth{0.000000pt}%
\definecolor{currentstroke}{rgb}{0.000000,0.000000,0.000000}%
\pgfsetstrokecolor{currentstroke}%
\pgfsetstrokeopacity{0.000000}%
\pgfsetdash{}{0pt}%
\pgfpathmoveto{\pgfqpoint{8.105650in}{0.660000in}}%
\pgfpathlineto{\pgfqpoint{8.111850in}{0.660000in}}%
\pgfpathlineto{\pgfqpoint{8.111850in}{3.201290in}}%
\pgfpathlineto{\pgfqpoint{8.105650in}{3.201290in}}%
\pgfpathlineto{\pgfqpoint{8.105650in}{0.660000in}}%
\pgfpathclose%
\pgfusepath{fill}%
\end{pgfscope}%
\begin{pgfscope}%
\pgfpathrectangle{\pgfqpoint{1.250000in}{0.660000in}}{\pgfqpoint{7.750000in}{4.620000in}}%
\pgfusepath{clip}%
\pgfsetbuttcap%
\pgfsetmiterjoin%
\definecolor{currentfill}{rgb}{0.000000,0.000000,0.000000}%
\pgfsetfillcolor{currentfill}%
\pgfsetlinewidth{0.000000pt}%
\definecolor{currentstroke}{rgb}{0.000000,0.000000,0.000000}%
\pgfsetstrokecolor{currentstroke}%
\pgfsetstrokeopacity{0.000000}%
\pgfsetdash{}{0pt}%
\pgfpathmoveto{\pgfqpoint{8.113400in}{0.660000in}}%
\pgfpathlineto{\pgfqpoint{8.119600in}{0.660000in}}%
\pgfpathlineto{\pgfqpoint{8.119600in}{3.201197in}}%
\pgfpathlineto{\pgfqpoint{8.113400in}{3.201197in}}%
\pgfpathlineto{\pgfqpoint{8.113400in}{0.660000in}}%
\pgfpathclose%
\pgfusepath{fill}%
\end{pgfscope}%
\begin{pgfscope}%
\pgfpathrectangle{\pgfqpoint{1.250000in}{0.660000in}}{\pgfqpoint{7.750000in}{4.620000in}}%
\pgfusepath{clip}%
\pgfsetbuttcap%
\pgfsetmiterjoin%
\definecolor{currentfill}{rgb}{0.000000,0.000000,0.000000}%
\pgfsetfillcolor{currentfill}%
\pgfsetlinewidth{0.000000pt}%
\definecolor{currentstroke}{rgb}{0.000000,0.000000,0.000000}%
\pgfsetstrokecolor{currentstroke}%
\pgfsetstrokeopacity{0.000000}%
\pgfsetdash{}{0pt}%
\pgfpathmoveto{\pgfqpoint{8.121150in}{0.660000in}}%
\pgfpathlineto{\pgfqpoint{8.127350in}{0.660000in}}%
\pgfpathlineto{\pgfqpoint{8.127350in}{3.201104in}}%
\pgfpathlineto{\pgfqpoint{8.121150in}{3.201104in}}%
\pgfpathlineto{\pgfqpoint{8.121150in}{0.660000in}}%
\pgfpathclose%
\pgfusepath{fill}%
\end{pgfscope}%
\begin{pgfscope}%
\pgfpathrectangle{\pgfqpoint{1.250000in}{0.660000in}}{\pgfqpoint{7.750000in}{4.620000in}}%
\pgfusepath{clip}%
\pgfsetbuttcap%
\pgfsetmiterjoin%
\definecolor{currentfill}{rgb}{0.000000,0.000000,0.000000}%
\pgfsetfillcolor{currentfill}%
\pgfsetlinewidth{0.000000pt}%
\definecolor{currentstroke}{rgb}{0.000000,0.000000,0.000000}%
\pgfsetstrokecolor{currentstroke}%
\pgfsetstrokeopacity{0.000000}%
\pgfsetdash{}{0pt}%
\pgfpathmoveto{\pgfqpoint{8.128900in}{0.660000in}}%
\pgfpathlineto{\pgfqpoint{8.135100in}{0.660000in}}%
\pgfpathlineto{\pgfqpoint{8.135100in}{3.201011in}}%
\pgfpathlineto{\pgfqpoint{8.128900in}{3.201011in}}%
\pgfpathlineto{\pgfqpoint{8.128900in}{0.660000in}}%
\pgfpathclose%
\pgfusepath{fill}%
\end{pgfscope}%
\begin{pgfscope}%
\pgfpathrectangle{\pgfqpoint{1.250000in}{0.660000in}}{\pgfqpoint{7.750000in}{4.620000in}}%
\pgfusepath{clip}%
\pgfsetbuttcap%
\pgfsetmiterjoin%
\definecolor{currentfill}{rgb}{0.000000,0.000000,0.000000}%
\pgfsetfillcolor{currentfill}%
\pgfsetlinewidth{0.000000pt}%
\definecolor{currentstroke}{rgb}{0.000000,0.000000,0.000000}%
\pgfsetstrokecolor{currentstroke}%
\pgfsetstrokeopacity{0.000000}%
\pgfsetdash{}{0pt}%
\pgfpathmoveto{\pgfqpoint{8.136650in}{0.660000in}}%
\pgfpathlineto{\pgfqpoint{8.142850in}{0.660000in}}%
\pgfpathlineto{\pgfqpoint{8.142850in}{3.200918in}}%
\pgfpathlineto{\pgfqpoint{8.136650in}{3.200918in}}%
\pgfpathlineto{\pgfqpoint{8.136650in}{0.660000in}}%
\pgfpathclose%
\pgfusepath{fill}%
\end{pgfscope}%
\begin{pgfscope}%
\pgfpathrectangle{\pgfqpoint{1.250000in}{0.660000in}}{\pgfqpoint{7.750000in}{4.620000in}}%
\pgfusepath{clip}%
\pgfsetbuttcap%
\pgfsetmiterjoin%
\definecolor{currentfill}{rgb}{0.000000,0.000000,0.000000}%
\pgfsetfillcolor{currentfill}%
\pgfsetlinewidth{0.000000pt}%
\definecolor{currentstroke}{rgb}{0.000000,0.000000,0.000000}%
\pgfsetstrokecolor{currentstroke}%
\pgfsetstrokeopacity{0.000000}%
\pgfsetdash{}{0pt}%
\pgfpathmoveto{\pgfqpoint{8.144400in}{0.660000in}}%
\pgfpathlineto{\pgfqpoint{8.150600in}{0.660000in}}%
\pgfpathlineto{\pgfqpoint{8.150600in}{3.200825in}}%
\pgfpathlineto{\pgfqpoint{8.144400in}{3.200825in}}%
\pgfpathlineto{\pgfqpoint{8.144400in}{0.660000in}}%
\pgfpathclose%
\pgfusepath{fill}%
\end{pgfscope}%
\begin{pgfscope}%
\pgfpathrectangle{\pgfqpoint{1.250000in}{0.660000in}}{\pgfqpoint{7.750000in}{4.620000in}}%
\pgfusepath{clip}%
\pgfsetbuttcap%
\pgfsetmiterjoin%
\definecolor{currentfill}{rgb}{0.000000,0.000000,0.000000}%
\pgfsetfillcolor{currentfill}%
\pgfsetlinewidth{0.000000pt}%
\definecolor{currentstroke}{rgb}{0.000000,0.000000,0.000000}%
\pgfsetstrokecolor{currentstroke}%
\pgfsetstrokeopacity{0.000000}%
\pgfsetdash{}{0pt}%
\pgfpathmoveto{\pgfqpoint{8.152150in}{0.660000in}}%
\pgfpathlineto{\pgfqpoint{8.158350in}{0.660000in}}%
\pgfpathlineto{\pgfqpoint{8.158350in}{3.200732in}}%
\pgfpathlineto{\pgfqpoint{8.152150in}{3.200732in}}%
\pgfpathlineto{\pgfqpoint{8.152150in}{0.660000in}}%
\pgfpathclose%
\pgfusepath{fill}%
\end{pgfscope}%
\begin{pgfscope}%
\pgfpathrectangle{\pgfqpoint{1.250000in}{0.660000in}}{\pgfqpoint{7.750000in}{4.620000in}}%
\pgfusepath{clip}%
\pgfsetbuttcap%
\pgfsetmiterjoin%
\definecolor{currentfill}{rgb}{0.000000,0.000000,0.000000}%
\pgfsetfillcolor{currentfill}%
\pgfsetlinewidth{0.000000pt}%
\definecolor{currentstroke}{rgb}{0.000000,0.000000,0.000000}%
\pgfsetstrokecolor{currentstroke}%
\pgfsetstrokeopacity{0.000000}%
\pgfsetdash{}{0pt}%
\pgfpathmoveto{\pgfqpoint{8.159900in}{0.660000in}}%
\pgfpathlineto{\pgfqpoint{8.166100in}{0.660000in}}%
\pgfpathlineto{\pgfqpoint{8.166100in}{3.200639in}}%
\pgfpathlineto{\pgfqpoint{8.159900in}{3.200639in}}%
\pgfpathlineto{\pgfqpoint{8.159900in}{0.660000in}}%
\pgfpathclose%
\pgfusepath{fill}%
\end{pgfscope}%
\begin{pgfscope}%
\pgfpathrectangle{\pgfqpoint{1.250000in}{0.660000in}}{\pgfqpoint{7.750000in}{4.620000in}}%
\pgfusepath{clip}%
\pgfsetbuttcap%
\pgfsetmiterjoin%
\definecolor{currentfill}{rgb}{0.000000,0.000000,0.000000}%
\pgfsetfillcolor{currentfill}%
\pgfsetlinewidth{0.000000pt}%
\definecolor{currentstroke}{rgb}{0.000000,0.000000,0.000000}%
\pgfsetstrokecolor{currentstroke}%
\pgfsetstrokeopacity{0.000000}%
\pgfsetdash{}{0pt}%
\pgfpathmoveto{\pgfqpoint{8.167650in}{0.660000in}}%
\pgfpathlineto{\pgfqpoint{8.173850in}{0.660000in}}%
\pgfpathlineto{\pgfqpoint{8.173850in}{3.200546in}}%
\pgfpathlineto{\pgfqpoint{8.167650in}{3.200546in}}%
\pgfpathlineto{\pgfqpoint{8.167650in}{0.660000in}}%
\pgfpathclose%
\pgfusepath{fill}%
\end{pgfscope}%
\begin{pgfscope}%
\pgfpathrectangle{\pgfqpoint{1.250000in}{0.660000in}}{\pgfqpoint{7.750000in}{4.620000in}}%
\pgfusepath{clip}%
\pgfsetbuttcap%
\pgfsetmiterjoin%
\definecolor{currentfill}{rgb}{0.000000,0.000000,0.000000}%
\pgfsetfillcolor{currentfill}%
\pgfsetlinewidth{0.000000pt}%
\definecolor{currentstroke}{rgb}{0.000000,0.000000,0.000000}%
\pgfsetstrokecolor{currentstroke}%
\pgfsetstrokeopacity{0.000000}%
\pgfsetdash{}{0pt}%
\pgfpathmoveto{\pgfqpoint{8.175400in}{0.660000in}}%
\pgfpathlineto{\pgfqpoint{8.181600in}{0.660000in}}%
\pgfpathlineto{\pgfqpoint{8.181600in}{3.200455in}}%
\pgfpathlineto{\pgfqpoint{8.175400in}{3.200455in}}%
\pgfpathlineto{\pgfqpoint{8.175400in}{0.660000in}}%
\pgfpathclose%
\pgfusepath{fill}%
\end{pgfscope}%
\begin{pgfscope}%
\pgfpathrectangle{\pgfqpoint{1.250000in}{0.660000in}}{\pgfqpoint{7.750000in}{4.620000in}}%
\pgfusepath{clip}%
\pgfsetbuttcap%
\pgfsetmiterjoin%
\definecolor{currentfill}{rgb}{0.000000,0.000000,0.000000}%
\pgfsetfillcolor{currentfill}%
\pgfsetlinewidth{0.000000pt}%
\definecolor{currentstroke}{rgb}{0.000000,0.000000,0.000000}%
\pgfsetstrokecolor{currentstroke}%
\pgfsetstrokeopacity{0.000000}%
\pgfsetdash{}{0pt}%
\pgfpathmoveto{\pgfqpoint{8.183150in}{0.660000in}}%
\pgfpathlineto{\pgfqpoint{8.189350in}{0.660000in}}%
\pgfpathlineto{\pgfqpoint{8.189350in}{3.200366in}}%
\pgfpathlineto{\pgfqpoint{8.183150in}{3.200366in}}%
\pgfpathlineto{\pgfqpoint{8.183150in}{0.660000in}}%
\pgfpathclose%
\pgfusepath{fill}%
\end{pgfscope}%
\begin{pgfscope}%
\pgfpathrectangle{\pgfqpoint{1.250000in}{0.660000in}}{\pgfqpoint{7.750000in}{4.620000in}}%
\pgfusepath{clip}%
\pgfsetbuttcap%
\pgfsetmiterjoin%
\definecolor{currentfill}{rgb}{0.000000,0.000000,0.000000}%
\pgfsetfillcolor{currentfill}%
\pgfsetlinewidth{0.000000pt}%
\definecolor{currentstroke}{rgb}{0.000000,0.000000,0.000000}%
\pgfsetstrokecolor{currentstroke}%
\pgfsetstrokeopacity{0.000000}%
\pgfsetdash{}{0pt}%
\pgfpathmoveto{\pgfqpoint{8.190900in}{0.660000in}}%
\pgfpathlineto{\pgfqpoint{8.197100in}{0.660000in}}%
\pgfpathlineto{\pgfqpoint{8.197100in}{3.200276in}}%
\pgfpathlineto{\pgfqpoint{8.190900in}{3.200276in}}%
\pgfpathlineto{\pgfqpoint{8.190900in}{0.660000in}}%
\pgfpathclose%
\pgfusepath{fill}%
\end{pgfscope}%
\begin{pgfscope}%
\pgfpathrectangle{\pgfqpoint{1.250000in}{0.660000in}}{\pgfqpoint{7.750000in}{4.620000in}}%
\pgfusepath{clip}%
\pgfsetbuttcap%
\pgfsetmiterjoin%
\definecolor{currentfill}{rgb}{0.000000,0.000000,0.000000}%
\pgfsetfillcolor{currentfill}%
\pgfsetlinewidth{0.000000pt}%
\definecolor{currentstroke}{rgb}{0.000000,0.000000,0.000000}%
\pgfsetstrokecolor{currentstroke}%
\pgfsetstrokeopacity{0.000000}%
\pgfsetdash{}{0pt}%
\pgfpathmoveto{\pgfqpoint{8.198650in}{0.660000in}}%
\pgfpathlineto{\pgfqpoint{8.204850in}{0.660000in}}%
\pgfpathlineto{\pgfqpoint{8.204850in}{3.200186in}}%
\pgfpathlineto{\pgfqpoint{8.198650in}{3.200186in}}%
\pgfpathlineto{\pgfqpoint{8.198650in}{0.660000in}}%
\pgfpathclose%
\pgfusepath{fill}%
\end{pgfscope}%
\begin{pgfscope}%
\pgfpathrectangle{\pgfqpoint{1.250000in}{0.660000in}}{\pgfqpoint{7.750000in}{4.620000in}}%
\pgfusepath{clip}%
\pgfsetbuttcap%
\pgfsetmiterjoin%
\definecolor{currentfill}{rgb}{0.000000,0.000000,0.000000}%
\pgfsetfillcolor{currentfill}%
\pgfsetlinewidth{0.000000pt}%
\definecolor{currentstroke}{rgb}{0.000000,0.000000,0.000000}%
\pgfsetstrokecolor{currentstroke}%
\pgfsetstrokeopacity{0.000000}%
\pgfsetdash{}{0pt}%
\pgfpathmoveto{\pgfqpoint{8.206400in}{0.660000in}}%
\pgfpathlineto{\pgfqpoint{8.212600in}{0.660000in}}%
\pgfpathlineto{\pgfqpoint{8.212600in}{3.200097in}}%
\pgfpathlineto{\pgfqpoint{8.206400in}{3.200097in}}%
\pgfpathlineto{\pgfqpoint{8.206400in}{0.660000in}}%
\pgfpathclose%
\pgfusepath{fill}%
\end{pgfscope}%
\begin{pgfscope}%
\pgfpathrectangle{\pgfqpoint{1.250000in}{0.660000in}}{\pgfqpoint{7.750000in}{4.620000in}}%
\pgfusepath{clip}%
\pgfsetbuttcap%
\pgfsetmiterjoin%
\definecolor{currentfill}{rgb}{0.000000,0.000000,0.000000}%
\pgfsetfillcolor{currentfill}%
\pgfsetlinewidth{0.000000pt}%
\definecolor{currentstroke}{rgb}{0.000000,0.000000,0.000000}%
\pgfsetstrokecolor{currentstroke}%
\pgfsetstrokeopacity{0.000000}%
\pgfsetdash{}{0pt}%
\pgfpathmoveto{\pgfqpoint{8.214150in}{0.660000in}}%
\pgfpathlineto{\pgfqpoint{8.220350in}{0.660000in}}%
\pgfpathlineto{\pgfqpoint{8.220350in}{3.200007in}}%
\pgfpathlineto{\pgfqpoint{8.214150in}{3.200007in}}%
\pgfpathlineto{\pgfqpoint{8.214150in}{0.660000in}}%
\pgfpathclose%
\pgfusepath{fill}%
\end{pgfscope}%
\begin{pgfscope}%
\pgfpathrectangle{\pgfqpoint{1.250000in}{0.660000in}}{\pgfqpoint{7.750000in}{4.620000in}}%
\pgfusepath{clip}%
\pgfsetbuttcap%
\pgfsetmiterjoin%
\definecolor{currentfill}{rgb}{0.000000,0.000000,0.000000}%
\pgfsetfillcolor{currentfill}%
\pgfsetlinewidth{0.000000pt}%
\definecolor{currentstroke}{rgb}{0.000000,0.000000,0.000000}%
\pgfsetstrokecolor{currentstroke}%
\pgfsetstrokeopacity{0.000000}%
\pgfsetdash{}{0pt}%
\pgfpathmoveto{\pgfqpoint{8.221900in}{0.660000in}}%
\pgfpathlineto{\pgfqpoint{8.228100in}{0.660000in}}%
\pgfpathlineto{\pgfqpoint{8.228100in}{3.199917in}}%
\pgfpathlineto{\pgfqpoint{8.221900in}{3.199917in}}%
\pgfpathlineto{\pgfqpoint{8.221900in}{0.660000in}}%
\pgfpathclose%
\pgfusepath{fill}%
\end{pgfscope}%
\begin{pgfscope}%
\pgfpathrectangle{\pgfqpoint{1.250000in}{0.660000in}}{\pgfqpoint{7.750000in}{4.620000in}}%
\pgfusepath{clip}%
\pgfsetbuttcap%
\pgfsetmiterjoin%
\definecolor{currentfill}{rgb}{0.000000,0.000000,0.000000}%
\pgfsetfillcolor{currentfill}%
\pgfsetlinewidth{0.000000pt}%
\definecolor{currentstroke}{rgb}{0.000000,0.000000,0.000000}%
\pgfsetstrokecolor{currentstroke}%
\pgfsetstrokeopacity{0.000000}%
\pgfsetdash{}{0pt}%
\pgfpathmoveto{\pgfqpoint{8.229650in}{0.660000in}}%
\pgfpathlineto{\pgfqpoint{8.235850in}{0.660000in}}%
\pgfpathlineto{\pgfqpoint{8.235850in}{3.199828in}}%
\pgfpathlineto{\pgfqpoint{8.229650in}{3.199828in}}%
\pgfpathlineto{\pgfqpoint{8.229650in}{0.660000in}}%
\pgfpathclose%
\pgfusepath{fill}%
\end{pgfscope}%
\begin{pgfscope}%
\pgfpathrectangle{\pgfqpoint{1.250000in}{0.660000in}}{\pgfqpoint{7.750000in}{4.620000in}}%
\pgfusepath{clip}%
\pgfsetbuttcap%
\pgfsetmiterjoin%
\definecolor{currentfill}{rgb}{0.000000,0.000000,0.000000}%
\pgfsetfillcolor{currentfill}%
\pgfsetlinewidth{0.000000pt}%
\definecolor{currentstroke}{rgb}{0.000000,0.000000,0.000000}%
\pgfsetstrokecolor{currentstroke}%
\pgfsetstrokeopacity{0.000000}%
\pgfsetdash{}{0pt}%
\pgfpathmoveto{\pgfqpoint{8.237400in}{0.660000in}}%
\pgfpathlineto{\pgfqpoint{8.243600in}{0.660000in}}%
\pgfpathlineto{\pgfqpoint{8.243600in}{3.199738in}}%
\pgfpathlineto{\pgfqpoint{8.237400in}{3.199738in}}%
\pgfpathlineto{\pgfqpoint{8.237400in}{0.660000in}}%
\pgfpathclose%
\pgfusepath{fill}%
\end{pgfscope}%
\begin{pgfscope}%
\pgfpathrectangle{\pgfqpoint{1.250000in}{0.660000in}}{\pgfqpoint{7.750000in}{4.620000in}}%
\pgfusepath{clip}%
\pgfsetbuttcap%
\pgfsetmiterjoin%
\definecolor{currentfill}{rgb}{0.000000,0.000000,0.000000}%
\pgfsetfillcolor{currentfill}%
\pgfsetlinewidth{0.000000pt}%
\definecolor{currentstroke}{rgb}{0.000000,0.000000,0.000000}%
\pgfsetstrokecolor{currentstroke}%
\pgfsetstrokeopacity{0.000000}%
\pgfsetdash{}{0pt}%
\pgfpathmoveto{\pgfqpoint{8.245150in}{0.660000in}}%
\pgfpathlineto{\pgfqpoint{8.251350in}{0.660000in}}%
\pgfpathlineto{\pgfqpoint{8.251350in}{3.199648in}}%
\pgfpathlineto{\pgfqpoint{8.245150in}{3.199648in}}%
\pgfpathlineto{\pgfqpoint{8.245150in}{0.660000in}}%
\pgfpathclose%
\pgfusepath{fill}%
\end{pgfscope}%
\begin{pgfscope}%
\pgfpathrectangle{\pgfqpoint{1.250000in}{0.660000in}}{\pgfqpoint{7.750000in}{4.620000in}}%
\pgfusepath{clip}%
\pgfsetbuttcap%
\pgfsetmiterjoin%
\definecolor{currentfill}{rgb}{0.000000,0.000000,0.000000}%
\pgfsetfillcolor{currentfill}%
\pgfsetlinewidth{0.000000pt}%
\definecolor{currentstroke}{rgb}{0.000000,0.000000,0.000000}%
\pgfsetstrokecolor{currentstroke}%
\pgfsetstrokeopacity{0.000000}%
\pgfsetdash{}{0pt}%
\pgfpathmoveto{\pgfqpoint{8.252900in}{0.660000in}}%
\pgfpathlineto{\pgfqpoint{8.259100in}{0.660000in}}%
\pgfpathlineto{\pgfqpoint{8.259100in}{3.199558in}}%
\pgfpathlineto{\pgfqpoint{8.252900in}{3.199558in}}%
\pgfpathlineto{\pgfqpoint{8.252900in}{0.660000in}}%
\pgfpathclose%
\pgfusepath{fill}%
\end{pgfscope}%
\begin{pgfscope}%
\pgfpathrectangle{\pgfqpoint{1.250000in}{0.660000in}}{\pgfqpoint{7.750000in}{4.620000in}}%
\pgfusepath{clip}%
\pgfsetbuttcap%
\pgfsetmiterjoin%
\definecolor{currentfill}{rgb}{0.000000,0.000000,0.000000}%
\pgfsetfillcolor{currentfill}%
\pgfsetlinewidth{0.000000pt}%
\definecolor{currentstroke}{rgb}{0.000000,0.000000,0.000000}%
\pgfsetstrokecolor{currentstroke}%
\pgfsetstrokeopacity{0.000000}%
\pgfsetdash{}{0pt}%
\pgfpathmoveto{\pgfqpoint{8.260650in}{0.660000in}}%
\pgfpathlineto{\pgfqpoint{8.266850in}{0.660000in}}%
\pgfpathlineto{\pgfqpoint{8.266850in}{3.199469in}}%
\pgfpathlineto{\pgfqpoint{8.260650in}{3.199469in}}%
\pgfpathlineto{\pgfqpoint{8.260650in}{0.660000in}}%
\pgfpathclose%
\pgfusepath{fill}%
\end{pgfscope}%
\begin{pgfscope}%
\pgfpathrectangle{\pgfqpoint{1.250000in}{0.660000in}}{\pgfqpoint{7.750000in}{4.620000in}}%
\pgfusepath{clip}%
\pgfsetbuttcap%
\pgfsetmiterjoin%
\definecolor{currentfill}{rgb}{0.000000,0.000000,0.000000}%
\pgfsetfillcolor{currentfill}%
\pgfsetlinewidth{0.000000pt}%
\definecolor{currentstroke}{rgb}{0.000000,0.000000,0.000000}%
\pgfsetstrokecolor{currentstroke}%
\pgfsetstrokeopacity{0.000000}%
\pgfsetdash{}{0pt}%
\pgfpathmoveto{\pgfqpoint{8.268400in}{0.660000in}}%
\pgfpathlineto{\pgfqpoint{8.274600in}{0.660000in}}%
\pgfpathlineto{\pgfqpoint{8.274600in}{3.199379in}}%
\pgfpathlineto{\pgfqpoint{8.268400in}{3.199379in}}%
\pgfpathlineto{\pgfqpoint{8.268400in}{0.660000in}}%
\pgfpathclose%
\pgfusepath{fill}%
\end{pgfscope}%
\begin{pgfscope}%
\pgfpathrectangle{\pgfqpoint{1.250000in}{0.660000in}}{\pgfqpoint{7.750000in}{4.620000in}}%
\pgfusepath{clip}%
\pgfsetbuttcap%
\pgfsetmiterjoin%
\definecolor{currentfill}{rgb}{0.000000,0.000000,0.000000}%
\pgfsetfillcolor{currentfill}%
\pgfsetlinewidth{0.000000pt}%
\definecolor{currentstroke}{rgb}{0.000000,0.000000,0.000000}%
\pgfsetstrokecolor{currentstroke}%
\pgfsetstrokeopacity{0.000000}%
\pgfsetdash{}{0pt}%
\pgfpathmoveto{\pgfqpoint{8.276150in}{0.660000in}}%
\pgfpathlineto{\pgfqpoint{8.282350in}{0.660000in}}%
\pgfpathlineto{\pgfqpoint{8.282350in}{3.199289in}}%
\pgfpathlineto{\pgfqpoint{8.276150in}{3.199289in}}%
\pgfpathlineto{\pgfqpoint{8.276150in}{0.660000in}}%
\pgfpathclose%
\pgfusepath{fill}%
\end{pgfscope}%
\begin{pgfscope}%
\pgfpathrectangle{\pgfqpoint{1.250000in}{0.660000in}}{\pgfqpoint{7.750000in}{4.620000in}}%
\pgfusepath{clip}%
\pgfsetbuttcap%
\pgfsetmiterjoin%
\definecolor{currentfill}{rgb}{0.000000,0.000000,0.000000}%
\pgfsetfillcolor{currentfill}%
\pgfsetlinewidth{0.000000pt}%
\definecolor{currentstroke}{rgb}{0.000000,0.000000,0.000000}%
\pgfsetstrokecolor{currentstroke}%
\pgfsetstrokeopacity{0.000000}%
\pgfsetdash{}{0pt}%
\pgfpathmoveto{\pgfqpoint{8.283900in}{0.660000in}}%
\pgfpathlineto{\pgfqpoint{8.290100in}{0.660000in}}%
\pgfpathlineto{\pgfqpoint{8.290100in}{3.199200in}}%
\pgfpathlineto{\pgfqpoint{8.283900in}{3.199200in}}%
\pgfpathlineto{\pgfqpoint{8.283900in}{0.660000in}}%
\pgfpathclose%
\pgfusepath{fill}%
\end{pgfscope}%
\begin{pgfscope}%
\pgfpathrectangle{\pgfqpoint{1.250000in}{0.660000in}}{\pgfqpoint{7.750000in}{4.620000in}}%
\pgfusepath{clip}%
\pgfsetbuttcap%
\pgfsetmiterjoin%
\definecolor{currentfill}{rgb}{0.000000,0.000000,0.000000}%
\pgfsetfillcolor{currentfill}%
\pgfsetlinewidth{0.000000pt}%
\definecolor{currentstroke}{rgb}{0.000000,0.000000,0.000000}%
\pgfsetstrokecolor{currentstroke}%
\pgfsetstrokeopacity{0.000000}%
\pgfsetdash{}{0pt}%
\pgfpathmoveto{\pgfqpoint{8.291650in}{0.660000in}}%
\pgfpathlineto{\pgfqpoint{8.297850in}{0.660000in}}%
\pgfpathlineto{\pgfqpoint{8.297850in}{3.199110in}}%
\pgfpathlineto{\pgfqpoint{8.291650in}{3.199110in}}%
\pgfpathlineto{\pgfqpoint{8.291650in}{0.660000in}}%
\pgfpathclose%
\pgfusepath{fill}%
\end{pgfscope}%
\begin{pgfscope}%
\pgfpathrectangle{\pgfqpoint{1.250000in}{0.660000in}}{\pgfqpoint{7.750000in}{4.620000in}}%
\pgfusepath{clip}%
\pgfsetbuttcap%
\pgfsetmiterjoin%
\definecolor{currentfill}{rgb}{0.000000,0.000000,0.000000}%
\pgfsetfillcolor{currentfill}%
\pgfsetlinewidth{0.000000pt}%
\definecolor{currentstroke}{rgb}{0.000000,0.000000,0.000000}%
\pgfsetstrokecolor{currentstroke}%
\pgfsetstrokeopacity{0.000000}%
\pgfsetdash{}{0pt}%
\pgfpathmoveto{\pgfqpoint{8.299400in}{0.660000in}}%
\pgfpathlineto{\pgfqpoint{8.305600in}{0.660000in}}%
\pgfpathlineto{\pgfqpoint{8.305600in}{3.199020in}}%
\pgfpathlineto{\pgfqpoint{8.299400in}{3.199020in}}%
\pgfpathlineto{\pgfqpoint{8.299400in}{0.660000in}}%
\pgfpathclose%
\pgfusepath{fill}%
\end{pgfscope}%
\begin{pgfscope}%
\pgfpathrectangle{\pgfqpoint{1.250000in}{0.660000in}}{\pgfqpoint{7.750000in}{4.620000in}}%
\pgfusepath{clip}%
\pgfsetbuttcap%
\pgfsetmiterjoin%
\definecolor{currentfill}{rgb}{0.000000,0.000000,0.000000}%
\pgfsetfillcolor{currentfill}%
\pgfsetlinewidth{0.000000pt}%
\definecolor{currentstroke}{rgb}{0.000000,0.000000,0.000000}%
\pgfsetstrokecolor{currentstroke}%
\pgfsetstrokeopacity{0.000000}%
\pgfsetdash{}{0pt}%
\pgfpathmoveto{\pgfqpoint{8.307150in}{0.660000in}}%
\pgfpathlineto{\pgfqpoint{8.313350in}{0.660000in}}%
\pgfpathlineto{\pgfqpoint{8.313350in}{3.198930in}}%
\pgfpathlineto{\pgfqpoint{8.307150in}{3.198930in}}%
\pgfpathlineto{\pgfqpoint{8.307150in}{0.660000in}}%
\pgfpathclose%
\pgfusepath{fill}%
\end{pgfscope}%
\begin{pgfscope}%
\pgfpathrectangle{\pgfqpoint{1.250000in}{0.660000in}}{\pgfqpoint{7.750000in}{4.620000in}}%
\pgfusepath{clip}%
\pgfsetbuttcap%
\pgfsetmiterjoin%
\definecolor{currentfill}{rgb}{0.000000,0.000000,0.000000}%
\pgfsetfillcolor{currentfill}%
\pgfsetlinewidth{0.000000pt}%
\definecolor{currentstroke}{rgb}{0.000000,0.000000,0.000000}%
\pgfsetstrokecolor{currentstroke}%
\pgfsetstrokeopacity{0.000000}%
\pgfsetdash{}{0pt}%
\pgfpathmoveto{\pgfqpoint{8.314900in}{0.660000in}}%
\pgfpathlineto{\pgfqpoint{8.321100in}{0.660000in}}%
\pgfpathlineto{\pgfqpoint{8.321100in}{3.198841in}}%
\pgfpathlineto{\pgfqpoint{8.314900in}{3.198841in}}%
\pgfpathlineto{\pgfqpoint{8.314900in}{0.660000in}}%
\pgfpathclose%
\pgfusepath{fill}%
\end{pgfscope}%
\begin{pgfscope}%
\pgfpathrectangle{\pgfqpoint{1.250000in}{0.660000in}}{\pgfqpoint{7.750000in}{4.620000in}}%
\pgfusepath{clip}%
\pgfsetbuttcap%
\pgfsetmiterjoin%
\definecolor{currentfill}{rgb}{0.000000,0.000000,0.000000}%
\pgfsetfillcolor{currentfill}%
\pgfsetlinewidth{0.000000pt}%
\definecolor{currentstroke}{rgb}{0.000000,0.000000,0.000000}%
\pgfsetstrokecolor{currentstroke}%
\pgfsetstrokeopacity{0.000000}%
\pgfsetdash{}{0pt}%
\pgfpathmoveto{\pgfqpoint{8.322650in}{0.660000in}}%
\pgfpathlineto{\pgfqpoint{8.328850in}{0.660000in}}%
\pgfpathlineto{\pgfqpoint{8.328850in}{3.198751in}}%
\pgfpathlineto{\pgfqpoint{8.322650in}{3.198751in}}%
\pgfpathlineto{\pgfqpoint{8.322650in}{0.660000in}}%
\pgfpathclose%
\pgfusepath{fill}%
\end{pgfscope}%
\begin{pgfscope}%
\pgfpathrectangle{\pgfqpoint{1.250000in}{0.660000in}}{\pgfqpoint{7.750000in}{4.620000in}}%
\pgfusepath{clip}%
\pgfsetbuttcap%
\pgfsetmiterjoin%
\definecolor{currentfill}{rgb}{0.000000,0.000000,0.000000}%
\pgfsetfillcolor{currentfill}%
\pgfsetlinewidth{0.000000pt}%
\definecolor{currentstroke}{rgb}{0.000000,0.000000,0.000000}%
\pgfsetstrokecolor{currentstroke}%
\pgfsetstrokeopacity{0.000000}%
\pgfsetdash{}{0pt}%
\pgfpathmoveto{\pgfqpoint{8.330400in}{0.660000in}}%
\pgfpathlineto{\pgfqpoint{8.336600in}{0.660000in}}%
\pgfpathlineto{\pgfqpoint{8.336600in}{3.198661in}}%
\pgfpathlineto{\pgfqpoint{8.330400in}{3.198661in}}%
\pgfpathlineto{\pgfqpoint{8.330400in}{0.660000in}}%
\pgfpathclose%
\pgfusepath{fill}%
\end{pgfscope}%
\begin{pgfscope}%
\pgfpathrectangle{\pgfqpoint{1.250000in}{0.660000in}}{\pgfqpoint{7.750000in}{4.620000in}}%
\pgfusepath{clip}%
\pgfsetbuttcap%
\pgfsetmiterjoin%
\definecolor{currentfill}{rgb}{0.000000,0.000000,0.000000}%
\pgfsetfillcolor{currentfill}%
\pgfsetlinewidth{0.000000pt}%
\definecolor{currentstroke}{rgb}{0.000000,0.000000,0.000000}%
\pgfsetstrokecolor{currentstroke}%
\pgfsetstrokeopacity{0.000000}%
\pgfsetdash{}{0pt}%
\pgfpathmoveto{\pgfqpoint{8.338150in}{0.660000in}}%
\pgfpathlineto{\pgfqpoint{8.344350in}{0.660000in}}%
\pgfpathlineto{\pgfqpoint{8.344350in}{3.198572in}}%
\pgfpathlineto{\pgfqpoint{8.338150in}{3.198572in}}%
\pgfpathlineto{\pgfqpoint{8.338150in}{0.660000in}}%
\pgfpathclose%
\pgfusepath{fill}%
\end{pgfscope}%
\begin{pgfscope}%
\pgfpathrectangle{\pgfqpoint{1.250000in}{0.660000in}}{\pgfqpoint{7.750000in}{4.620000in}}%
\pgfusepath{clip}%
\pgfsetbuttcap%
\pgfsetmiterjoin%
\definecolor{currentfill}{rgb}{0.000000,0.000000,0.000000}%
\pgfsetfillcolor{currentfill}%
\pgfsetlinewidth{0.000000pt}%
\definecolor{currentstroke}{rgb}{0.000000,0.000000,0.000000}%
\pgfsetstrokecolor{currentstroke}%
\pgfsetstrokeopacity{0.000000}%
\pgfsetdash{}{0pt}%
\pgfpathmoveto{\pgfqpoint{8.345900in}{0.660000in}}%
\pgfpathlineto{\pgfqpoint{8.352100in}{0.660000in}}%
\pgfpathlineto{\pgfqpoint{8.352100in}{3.198482in}}%
\pgfpathlineto{\pgfqpoint{8.345900in}{3.198482in}}%
\pgfpathlineto{\pgfqpoint{8.345900in}{0.660000in}}%
\pgfpathclose%
\pgfusepath{fill}%
\end{pgfscope}%
\begin{pgfscope}%
\pgfpathrectangle{\pgfqpoint{1.250000in}{0.660000in}}{\pgfqpoint{7.750000in}{4.620000in}}%
\pgfusepath{clip}%
\pgfsetbuttcap%
\pgfsetmiterjoin%
\definecolor{currentfill}{rgb}{0.000000,0.000000,0.000000}%
\pgfsetfillcolor{currentfill}%
\pgfsetlinewidth{0.000000pt}%
\definecolor{currentstroke}{rgb}{0.000000,0.000000,0.000000}%
\pgfsetstrokecolor{currentstroke}%
\pgfsetstrokeopacity{0.000000}%
\pgfsetdash{}{0pt}%
\pgfpathmoveto{\pgfqpoint{8.353650in}{0.660000in}}%
\pgfpathlineto{\pgfqpoint{8.359850in}{0.660000in}}%
\pgfpathlineto{\pgfqpoint{8.359850in}{3.198392in}}%
\pgfpathlineto{\pgfqpoint{8.353650in}{3.198392in}}%
\pgfpathlineto{\pgfqpoint{8.353650in}{0.660000in}}%
\pgfpathclose%
\pgfusepath{fill}%
\end{pgfscope}%
\begin{pgfscope}%
\pgfpathrectangle{\pgfqpoint{1.250000in}{0.660000in}}{\pgfqpoint{7.750000in}{4.620000in}}%
\pgfusepath{clip}%
\pgfsetbuttcap%
\pgfsetmiterjoin%
\definecolor{currentfill}{rgb}{0.000000,0.000000,0.000000}%
\pgfsetfillcolor{currentfill}%
\pgfsetlinewidth{0.000000pt}%
\definecolor{currentstroke}{rgb}{0.000000,0.000000,0.000000}%
\pgfsetstrokecolor{currentstroke}%
\pgfsetstrokeopacity{0.000000}%
\pgfsetdash{}{0pt}%
\pgfpathmoveto{\pgfqpoint{8.361400in}{0.660000in}}%
\pgfpathlineto{\pgfqpoint{8.367600in}{0.660000in}}%
\pgfpathlineto{\pgfqpoint{8.367600in}{3.198306in}}%
\pgfpathlineto{\pgfqpoint{8.361400in}{3.198306in}}%
\pgfpathlineto{\pgfqpoint{8.361400in}{0.660000in}}%
\pgfpathclose%
\pgfusepath{fill}%
\end{pgfscope}%
\begin{pgfscope}%
\pgfpathrectangle{\pgfqpoint{1.250000in}{0.660000in}}{\pgfqpoint{7.750000in}{4.620000in}}%
\pgfusepath{clip}%
\pgfsetbuttcap%
\pgfsetmiterjoin%
\definecolor{currentfill}{rgb}{0.000000,0.000000,0.000000}%
\pgfsetfillcolor{currentfill}%
\pgfsetlinewidth{0.000000pt}%
\definecolor{currentstroke}{rgb}{0.000000,0.000000,0.000000}%
\pgfsetstrokecolor{currentstroke}%
\pgfsetstrokeopacity{0.000000}%
\pgfsetdash{}{0pt}%
\pgfpathmoveto{\pgfqpoint{8.369150in}{0.660000in}}%
\pgfpathlineto{\pgfqpoint{8.375350in}{0.660000in}}%
\pgfpathlineto{\pgfqpoint{8.375350in}{3.198219in}}%
\pgfpathlineto{\pgfqpoint{8.369150in}{3.198219in}}%
\pgfpathlineto{\pgfqpoint{8.369150in}{0.660000in}}%
\pgfpathclose%
\pgfusepath{fill}%
\end{pgfscope}%
\begin{pgfscope}%
\pgfpathrectangle{\pgfqpoint{1.250000in}{0.660000in}}{\pgfqpoint{7.750000in}{4.620000in}}%
\pgfusepath{clip}%
\pgfsetbuttcap%
\pgfsetmiterjoin%
\definecolor{currentfill}{rgb}{0.000000,0.000000,0.000000}%
\pgfsetfillcolor{currentfill}%
\pgfsetlinewidth{0.000000pt}%
\definecolor{currentstroke}{rgb}{0.000000,0.000000,0.000000}%
\pgfsetstrokecolor{currentstroke}%
\pgfsetstrokeopacity{0.000000}%
\pgfsetdash{}{0pt}%
\pgfpathmoveto{\pgfqpoint{8.376900in}{0.660000in}}%
\pgfpathlineto{\pgfqpoint{8.383100in}{0.660000in}}%
\pgfpathlineto{\pgfqpoint{8.383100in}{3.198133in}}%
\pgfpathlineto{\pgfqpoint{8.376900in}{3.198133in}}%
\pgfpathlineto{\pgfqpoint{8.376900in}{0.660000in}}%
\pgfpathclose%
\pgfusepath{fill}%
\end{pgfscope}%
\begin{pgfscope}%
\pgfpathrectangle{\pgfqpoint{1.250000in}{0.660000in}}{\pgfqpoint{7.750000in}{4.620000in}}%
\pgfusepath{clip}%
\pgfsetbuttcap%
\pgfsetmiterjoin%
\definecolor{currentfill}{rgb}{0.000000,0.000000,0.000000}%
\pgfsetfillcolor{currentfill}%
\pgfsetlinewidth{0.000000pt}%
\definecolor{currentstroke}{rgb}{0.000000,0.000000,0.000000}%
\pgfsetstrokecolor{currentstroke}%
\pgfsetstrokeopacity{0.000000}%
\pgfsetdash{}{0pt}%
\pgfpathmoveto{\pgfqpoint{8.384650in}{0.660000in}}%
\pgfpathlineto{\pgfqpoint{8.390850in}{0.660000in}}%
\pgfpathlineto{\pgfqpoint{8.390850in}{3.198046in}}%
\pgfpathlineto{\pgfqpoint{8.384650in}{3.198046in}}%
\pgfpathlineto{\pgfqpoint{8.384650in}{0.660000in}}%
\pgfpathclose%
\pgfusepath{fill}%
\end{pgfscope}%
\begin{pgfscope}%
\pgfpathrectangle{\pgfqpoint{1.250000in}{0.660000in}}{\pgfqpoint{7.750000in}{4.620000in}}%
\pgfusepath{clip}%
\pgfsetbuttcap%
\pgfsetmiterjoin%
\definecolor{currentfill}{rgb}{0.000000,0.000000,0.000000}%
\pgfsetfillcolor{currentfill}%
\pgfsetlinewidth{0.000000pt}%
\definecolor{currentstroke}{rgb}{0.000000,0.000000,0.000000}%
\pgfsetstrokecolor{currentstroke}%
\pgfsetstrokeopacity{0.000000}%
\pgfsetdash{}{0pt}%
\pgfpathmoveto{\pgfqpoint{8.392400in}{0.660000in}}%
\pgfpathlineto{\pgfqpoint{8.398600in}{0.660000in}}%
\pgfpathlineto{\pgfqpoint{8.398600in}{3.197960in}}%
\pgfpathlineto{\pgfqpoint{8.392400in}{3.197960in}}%
\pgfpathlineto{\pgfqpoint{8.392400in}{0.660000in}}%
\pgfpathclose%
\pgfusepath{fill}%
\end{pgfscope}%
\begin{pgfscope}%
\pgfpathrectangle{\pgfqpoint{1.250000in}{0.660000in}}{\pgfqpoint{7.750000in}{4.620000in}}%
\pgfusepath{clip}%
\pgfsetbuttcap%
\pgfsetmiterjoin%
\definecolor{currentfill}{rgb}{0.000000,0.000000,0.000000}%
\pgfsetfillcolor{currentfill}%
\pgfsetlinewidth{0.000000pt}%
\definecolor{currentstroke}{rgb}{0.000000,0.000000,0.000000}%
\pgfsetstrokecolor{currentstroke}%
\pgfsetstrokeopacity{0.000000}%
\pgfsetdash{}{0pt}%
\pgfpathmoveto{\pgfqpoint{8.400150in}{0.660000in}}%
\pgfpathlineto{\pgfqpoint{8.406350in}{0.660000in}}%
\pgfpathlineto{\pgfqpoint{8.406350in}{3.197873in}}%
\pgfpathlineto{\pgfqpoint{8.400150in}{3.197873in}}%
\pgfpathlineto{\pgfqpoint{8.400150in}{0.660000in}}%
\pgfpathclose%
\pgfusepath{fill}%
\end{pgfscope}%
\begin{pgfscope}%
\pgfpathrectangle{\pgfqpoint{1.250000in}{0.660000in}}{\pgfqpoint{7.750000in}{4.620000in}}%
\pgfusepath{clip}%
\pgfsetbuttcap%
\pgfsetmiterjoin%
\definecolor{currentfill}{rgb}{0.000000,0.000000,0.000000}%
\pgfsetfillcolor{currentfill}%
\pgfsetlinewidth{0.000000pt}%
\definecolor{currentstroke}{rgb}{0.000000,0.000000,0.000000}%
\pgfsetstrokecolor{currentstroke}%
\pgfsetstrokeopacity{0.000000}%
\pgfsetdash{}{0pt}%
\pgfpathmoveto{\pgfqpoint{8.407900in}{0.660000in}}%
\pgfpathlineto{\pgfqpoint{8.414100in}{0.660000in}}%
\pgfpathlineto{\pgfqpoint{8.414100in}{3.197787in}}%
\pgfpathlineto{\pgfqpoint{8.407900in}{3.197787in}}%
\pgfpathlineto{\pgfqpoint{8.407900in}{0.660000in}}%
\pgfpathclose%
\pgfusepath{fill}%
\end{pgfscope}%
\begin{pgfscope}%
\pgfpathrectangle{\pgfqpoint{1.250000in}{0.660000in}}{\pgfqpoint{7.750000in}{4.620000in}}%
\pgfusepath{clip}%
\pgfsetbuttcap%
\pgfsetmiterjoin%
\definecolor{currentfill}{rgb}{0.000000,0.000000,0.000000}%
\pgfsetfillcolor{currentfill}%
\pgfsetlinewidth{0.000000pt}%
\definecolor{currentstroke}{rgb}{0.000000,0.000000,0.000000}%
\pgfsetstrokecolor{currentstroke}%
\pgfsetstrokeopacity{0.000000}%
\pgfsetdash{}{0pt}%
\pgfpathmoveto{\pgfqpoint{8.415650in}{0.660000in}}%
\pgfpathlineto{\pgfqpoint{8.421850in}{0.660000in}}%
\pgfpathlineto{\pgfqpoint{8.421850in}{3.197700in}}%
\pgfpathlineto{\pgfqpoint{8.415650in}{3.197700in}}%
\pgfpathlineto{\pgfqpoint{8.415650in}{0.660000in}}%
\pgfpathclose%
\pgfusepath{fill}%
\end{pgfscope}%
\begin{pgfscope}%
\pgfpathrectangle{\pgfqpoint{1.250000in}{0.660000in}}{\pgfqpoint{7.750000in}{4.620000in}}%
\pgfusepath{clip}%
\pgfsetbuttcap%
\pgfsetmiterjoin%
\definecolor{currentfill}{rgb}{0.000000,0.000000,0.000000}%
\pgfsetfillcolor{currentfill}%
\pgfsetlinewidth{0.000000pt}%
\definecolor{currentstroke}{rgb}{0.000000,0.000000,0.000000}%
\pgfsetstrokecolor{currentstroke}%
\pgfsetstrokeopacity{0.000000}%
\pgfsetdash{}{0pt}%
\pgfpathmoveto{\pgfqpoint{8.423400in}{0.660000in}}%
\pgfpathlineto{\pgfqpoint{8.429600in}{0.660000in}}%
\pgfpathlineto{\pgfqpoint{8.429600in}{3.197614in}}%
\pgfpathlineto{\pgfqpoint{8.423400in}{3.197614in}}%
\pgfpathlineto{\pgfqpoint{8.423400in}{0.660000in}}%
\pgfpathclose%
\pgfusepath{fill}%
\end{pgfscope}%
\begin{pgfscope}%
\pgfpathrectangle{\pgfqpoint{1.250000in}{0.660000in}}{\pgfqpoint{7.750000in}{4.620000in}}%
\pgfusepath{clip}%
\pgfsetbuttcap%
\pgfsetmiterjoin%
\definecolor{currentfill}{rgb}{0.000000,0.000000,0.000000}%
\pgfsetfillcolor{currentfill}%
\pgfsetlinewidth{0.000000pt}%
\definecolor{currentstroke}{rgb}{0.000000,0.000000,0.000000}%
\pgfsetstrokecolor{currentstroke}%
\pgfsetstrokeopacity{0.000000}%
\pgfsetdash{}{0pt}%
\pgfpathmoveto{\pgfqpoint{8.431150in}{0.660000in}}%
\pgfpathlineto{\pgfqpoint{8.437350in}{0.660000in}}%
\pgfpathlineto{\pgfqpoint{8.437350in}{3.197527in}}%
\pgfpathlineto{\pgfqpoint{8.431150in}{3.197527in}}%
\pgfpathlineto{\pgfqpoint{8.431150in}{0.660000in}}%
\pgfpathclose%
\pgfusepath{fill}%
\end{pgfscope}%
\begin{pgfscope}%
\pgfpathrectangle{\pgfqpoint{1.250000in}{0.660000in}}{\pgfqpoint{7.750000in}{4.620000in}}%
\pgfusepath{clip}%
\pgfsetbuttcap%
\pgfsetmiterjoin%
\definecolor{currentfill}{rgb}{0.000000,0.000000,0.000000}%
\pgfsetfillcolor{currentfill}%
\pgfsetlinewidth{0.000000pt}%
\definecolor{currentstroke}{rgb}{0.000000,0.000000,0.000000}%
\pgfsetstrokecolor{currentstroke}%
\pgfsetstrokeopacity{0.000000}%
\pgfsetdash{}{0pt}%
\pgfpathmoveto{\pgfqpoint{8.438900in}{0.660000in}}%
\pgfpathlineto{\pgfqpoint{8.445100in}{0.660000in}}%
\pgfpathlineto{\pgfqpoint{8.445100in}{3.197440in}}%
\pgfpathlineto{\pgfqpoint{8.438900in}{3.197440in}}%
\pgfpathlineto{\pgfqpoint{8.438900in}{0.660000in}}%
\pgfpathclose%
\pgfusepath{fill}%
\end{pgfscope}%
\begin{pgfscope}%
\pgfpathrectangle{\pgfqpoint{1.250000in}{0.660000in}}{\pgfqpoint{7.750000in}{4.620000in}}%
\pgfusepath{clip}%
\pgfsetbuttcap%
\pgfsetmiterjoin%
\definecolor{currentfill}{rgb}{0.000000,0.000000,0.000000}%
\pgfsetfillcolor{currentfill}%
\pgfsetlinewidth{0.000000pt}%
\definecolor{currentstroke}{rgb}{0.000000,0.000000,0.000000}%
\pgfsetstrokecolor{currentstroke}%
\pgfsetstrokeopacity{0.000000}%
\pgfsetdash{}{0pt}%
\pgfpathmoveto{\pgfqpoint{8.446650in}{0.660000in}}%
\pgfpathlineto{\pgfqpoint{8.452850in}{0.660000in}}%
\pgfpathlineto{\pgfqpoint{8.452850in}{3.197354in}}%
\pgfpathlineto{\pgfqpoint{8.446650in}{3.197354in}}%
\pgfpathlineto{\pgfqpoint{8.446650in}{0.660000in}}%
\pgfpathclose%
\pgfusepath{fill}%
\end{pgfscope}%
\begin{pgfscope}%
\pgfpathrectangle{\pgfqpoint{1.250000in}{0.660000in}}{\pgfqpoint{7.750000in}{4.620000in}}%
\pgfusepath{clip}%
\pgfsetbuttcap%
\pgfsetmiterjoin%
\definecolor{currentfill}{rgb}{0.000000,0.000000,0.000000}%
\pgfsetfillcolor{currentfill}%
\pgfsetlinewidth{0.000000pt}%
\definecolor{currentstroke}{rgb}{0.000000,0.000000,0.000000}%
\pgfsetstrokecolor{currentstroke}%
\pgfsetstrokeopacity{0.000000}%
\pgfsetdash{}{0pt}%
\pgfpathmoveto{\pgfqpoint{8.454400in}{0.660000in}}%
\pgfpathlineto{\pgfqpoint{8.460600in}{0.660000in}}%
\pgfpathlineto{\pgfqpoint{8.460600in}{3.197267in}}%
\pgfpathlineto{\pgfqpoint{8.454400in}{3.197267in}}%
\pgfpathlineto{\pgfqpoint{8.454400in}{0.660000in}}%
\pgfpathclose%
\pgfusepath{fill}%
\end{pgfscope}%
\begin{pgfscope}%
\pgfpathrectangle{\pgfqpoint{1.250000in}{0.660000in}}{\pgfqpoint{7.750000in}{4.620000in}}%
\pgfusepath{clip}%
\pgfsetbuttcap%
\pgfsetmiterjoin%
\definecolor{currentfill}{rgb}{0.000000,0.000000,0.000000}%
\pgfsetfillcolor{currentfill}%
\pgfsetlinewidth{0.000000pt}%
\definecolor{currentstroke}{rgb}{0.000000,0.000000,0.000000}%
\pgfsetstrokecolor{currentstroke}%
\pgfsetstrokeopacity{0.000000}%
\pgfsetdash{}{0pt}%
\pgfpathmoveto{\pgfqpoint{8.462150in}{0.660000in}}%
\pgfpathlineto{\pgfqpoint{8.468350in}{0.660000in}}%
\pgfpathlineto{\pgfqpoint{8.468350in}{3.197181in}}%
\pgfpathlineto{\pgfqpoint{8.462150in}{3.197181in}}%
\pgfpathlineto{\pgfqpoint{8.462150in}{0.660000in}}%
\pgfpathclose%
\pgfusepath{fill}%
\end{pgfscope}%
\begin{pgfscope}%
\pgfpathrectangle{\pgfqpoint{1.250000in}{0.660000in}}{\pgfqpoint{7.750000in}{4.620000in}}%
\pgfusepath{clip}%
\pgfsetbuttcap%
\pgfsetmiterjoin%
\definecolor{currentfill}{rgb}{0.000000,0.000000,0.000000}%
\pgfsetfillcolor{currentfill}%
\pgfsetlinewidth{0.000000pt}%
\definecolor{currentstroke}{rgb}{0.000000,0.000000,0.000000}%
\pgfsetstrokecolor{currentstroke}%
\pgfsetstrokeopacity{0.000000}%
\pgfsetdash{}{0pt}%
\pgfpathmoveto{\pgfqpoint{8.469900in}{0.660000in}}%
\pgfpathlineto{\pgfqpoint{8.476100in}{0.660000in}}%
\pgfpathlineto{\pgfqpoint{8.476100in}{3.197094in}}%
\pgfpathlineto{\pgfqpoint{8.469900in}{3.197094in}}%
\pgfpathlineto{\pgfqpoint{8.469900in}{0.660000in}}%
\pgfpathclose%
\pgfusepath{fill}%
\end{pgfscope}%
\begin{pgfscope}%
\pgfpathrectangle{\pgfqpoint{1.250000in}{0.660000in}}{\pgfqpoint{7.750000in}{4.620000in}}%
\pgfusepath{clip}%
\pgfsetbuttcap%
\pgfsetmiterjoin%
\definecolor{currentfill}{rgb}{0.000000,0.000000,0.000000}%
\pgfsetfillcolor{currentfill}%
\pgfsetlinewidth{0.000000pt}%
\definecolor{currentstroke}{rgb}{0.000000,0.000000,0.000000}%
\pgfsetstrokecolor{currentstroke}%
\pgfsetstrokeopacity{0.000000}%
\pgfsetdash{}{0pt}%
\pgfpathmoveto{\pgfqpoint{8.477650in}{0.660000in}}%
\pgfpathlineto{\pgfqpoint{8.483850in}{0.660000in}}%
\pgfpathlineto{\pgfqpoint{8.483850in}{3.197008in}}%
\pgfpathlineto{\pgfqpoint{8.477650in}{3.197008in}}%
\pgfpathlineto{\pgfqpoint{8.477650in}{0.660000in}}%
\pgfpathclose%
\pgfusepath{fill}%
\end{pgfscope}%
\begin{pgfscope}%
\pgfpathrectangle{\pgfqpoint{1.250000in}{0.660000in}}{\pgfqpoint{7.750000in}{4.620000in}}%
\pgfusepath{clip}%
\pgfsetbuttcap%
\pgfsetmiterjoin%
\definecolor{currentfill}{rgb}{0.000000,0.000000,0.000000}%
\pgfsetfillcolor{currentfill}%
\pgfsetlinewidth{0.000000pt}%
\definecolor{currentstroke}{rgb}{0.000000,0.000000,0.000000}%
\pgfsetstrokecolor{currentstroke}%
\pgfsetstrokeopacity{0.000000}%
\pgfsetdash{}{0pt}%
\pgfpathmoveto{\pgfqpoint{8.485400in}{0.660000in}}%
\pgfpathlineto{\pgfqpoint{8.491600in}{0.660000in}}%
\pgfpathlineto{\pgfqpoint{8.491600in}{3.196921in}}%
\pgfpathlineto{\pgfqpoint{8.485400in}{3.196921in}}%
\pgfpathlineto{\pgfqpoint{8.485400in}{0.660000in}}%
\pgfpathclose%
\pgfusepath{fill}%
\end{pgfscope}%
\begin{pgfscope}%
\pgfpathrectangle{\pgfqpoint{1.250000in}{0.660000in}}{\pgfqpoint{7.750000in}{4.620000in}}%
\pgfusepath{clip}%
\pgfsetbuttcap%
\pgfsetmiterjoin%
\definecolor{currentfill}{rgb}{0.000000,0.000000,0.000000}%
\pgfsetfillcolor{currentfill}%
\pgfsetlinewidth{0.000000pt}%
\definecolor{currentstroke}{rgb}{0.000000,0.000000,0.000000}%
\pgfsetstrokecolor{currentstroke}%
\pgfsetstrokeopacity{0.000000}%
\pgfsetdash{}{0pt}%
\pgfpathmoveto{\pgfqpoint{8.493150in}{0.660000in}}%
\pgfpathlineto{\pgfqpoint{8.499350in}{0.660000in}}%
\pgfpathlineto{\pgfqpoint{8.499350in}{3.196835in}}%
\pgfpathlineto{\pgfqpoint{8.493150in}{3.196835in}}%
\pgfpathlineto{\pgfqpoint{8.493150in}{0.660000in}}%
\pgfpathclose%
\pgfusepath{fill}%
\end{pgfscope}%
\begin{pgfscope}%
\pgfpathrectangle{\pgfqpoint{1.250000in}{0.660000in}}{\pgfqpoint{7.750000in}{4.620000in}}%
\pgfusepath{clip}%
\pgfsetbuttcap%
\pgfsetmiterjoin%
\definecolor{currentfill}{rgb}{0.000000,0.000000,0.000000}%
\pgfsetfillcolor{currentfill}%
\pgfsetlinewidth{0.000000pt}%
\definecolor{currentstroke}{rgb}{0.000000,0.000000,0.000000}%
\pgfsetstrokecolor{currentstroke}%
\pgfsetstrokeopacity{0.000000}%
\pgfsetdash{}{0pt}%
\pgfpathmoveto{\pgfqpoint{8.500900in}{0.660000in}}%
\pgfpathlineto{\pgfqpoint{8.507100in}{0.660000in}}%
\pgfpathlineto{\pgfqpoint{8.507100in}{3.196748in}}%
\pgfpathlineto{\pgfqpoint{8.500900in}{3.196748in}}%
\pgfpathlineto{\pgfqpoint{8.500900in}{0.660000in}}%
\pgfpathclose%
\pgfusepath{fill}%
\end{pgfscope}%
\begin{pgfscope}%
\pgfpathrectangle{\pgfqpoint{1.250000in}{0.660000in}}{\pgfqpoint{7.750000in}{4.620000in}}%
\pgfusepath{clip}%
\pgfsetbuttcap%
\pgfsetmiterjoin%
\definecolor{currentfill}{rgb}{0.000000,0.000000,0.000000}%
\pgfsetfillcolor{currentfill}%
\pgfsetlinewidth{0.000000pt}%
\definecolor{currentstroke}{rgb}{0.000000,0.000000,0.000000}%
\pgfsetstrokecolor{currentstroke}%
\pgfsetstrokeopacity{0.000000}%
\pgfsetdash{}{0pt}%
\pgfpathmoveto{\pgfqpoint{8.508650in}{0.660000in}}%
\pgfpathlineto{\pgfqpoint{8.514850in}{0.660000in}}%
\pgfpathlineto{\pgfqpoint{8.514850in}{3.196662in}}%
\pgfpathlineto{\pgfqpoint{8.508650in}{3.196662in}}%
\pgfpathlineto{\pgfqpoint{8.508650in}{0.660000in}}%
\pgfpathclose%
\pgfusepath{fill}%
\end{pgfscope}%
\begin{pgfscope}%
\pgfpathrectangle{\pgfqpoint{1.250000in}{0.660000in}}{\pgfqpoint{7.750000in}{4.620000in}}%
\pgfusepath{clip}%
\pgfsetbuttcap%
\pgfsetmiterjoin%
\definecolor{currentfill}{rgb}{0.000000,0.000000,0.000000}%
\pgfsetfillcolor{currentfill}%
\pgfsetlinewidth{0.000000pt}%
\definecolor{currentstroke}{rgb}{0.000000,0.000000,0.000000}%
\pgfsetstrokecolor{currentstroke}%
\pgfsetstrokeopacity{0.000000}%
\pgfsetdash{}{0pt}%
\pgfpathmoveto{\pgfqpoint{8.516400in}{0.660000in}}%
\pgfpathlineto{\pgfqpoint{8.522600in}{0.660000in}}%
\pgfpathlineto{\pgfqpoint{8.522600in}{3.196575in}}%
\pgfpathlineto{\pgfqpoint{8.516400in}{3.196575in}}%
\pgfpathlineto{\pgfqpoint{8.516400in}{0.660000in}}%
\pgfpathclose%
\pgfusepath{fill}%
\end{pgfscope}%
\begin{pgfscope}%
\pgfpathrectangle{\pgfqpoint{1.250000in}{0.660000in}}{\pgfqpoint{7.750000in}{4.620000in}}%
\pgfusepath{clip}%
\pgfsetbuttcap%
\pgfsetmiterjoin%
\definecolor{currentfill}{rgb}{0.000000,0.000000,0.000000}%
\pgfsetfillcolor{currentfill}%
\pgfsetlinewidth{0.000000pt}%
\definecolor{currentstroke}{rgb}{0.000000,0.000000,0.000000}%
\pgfsetstrokecolor{currentstroke}%
\pgfsetstrokeopacity{0.000000}%
\pgfsetdash{}{0pt}%
\pgfpathmoveto{\pgfqpoint{8.524150in}{0.660000in}}%
\pgfpathlineto{\pgfqpoint{8.530350in}{0.660000in}}%
\pgfpathlineto{\pgfqpoint{8.530350in}{3.196488in}}%
\pgfpathlineto{\pgfqpoint{8.524150in}{3.196488in}}%
\pgfpathlineto{\pgfqpoint{8.524150in}{0.660000in}}%
\pgfpathclose%
\pgfusepath{fill}%
\end{pgfscope}%
\begin{pgfscope}%
\pgfpathrectangle{\pgfqpoint{1.250000in}{0.660000in}}{\pgfqpoint{7.750000in}{4.620000in}}%
\pgfusepath{clip}%
\pgfsetbuttcap%
\pgfsetmiterjoin%
\definecolor{currentfill}{rgb}{0.000000,0.000000,0.000000}%
\pgfsetfillcolor{currentfill}%
\pgfsetlinewidth{0.000000pt}%
\definecolor{currentstroke}{rgb}{0.000000,0.000000,0.000000}%
\pgfsetstrokecolor{currentstroke}%
\pgfsetstrokeopacity{0.000000}%
\pgfsetdash{}{0pt}%
\pgfpathmoveto{\pgfqpoint{8.531900in}{0.660000in}}%
\pgfpathlineto{\pgfqpoint{8.538100in}{0.660000in}}%
\pgfpathlineto{\pgfqpoint{8.538100in}{3.196402in}}%
\pgfpathlineto{\pgfqpoint{8.531900in}{3.196402in}}%
\pgfpathlineto{\pgfqpoint{8.531900in}{0.660000in}}%
\pgfpathclose%
\pgfusepath{fill}%
\end{pgfscope}%
\begin{pgfscope}%
\pgfpathrectangle{\pgfqpoint{1.250000in}{0.660000in}}{\pgfqpoint{7.750000in}{4.620000in}}%
\pgfusepath{clip}%
\pgfsetbuttcap%
\pgfsetmiterjoin%
\definecolor{currentfill}{rgb}{0.000000,0.000000,0.000000}%
\pgfsetfillcolor{currentfill}%
\pgfsetlinewidth{0.000000pt}%
\definecolor{currentstroke}{rgb}{0.000000,0.000000,0.000000}%
\pgfsetstrokecolor{currentstroke}%
\pgfsetstrokeopacity{0.000000}%
\pgfsetdash{}{0pt}%
\pgfpathmoveto{\pgfqpoint{8.539650in}{0.660000in}}%
\pgfpathlineto{\pgfqpoint{8.545850in}{0.660000in}}%
\pgfpathlineto{\pgfqpoint{8.545850in}{3.196315in}}%
\pgfpathlineto{\pgfqpoint{8.539650in}{3.196315in}}%
\pgfpathlineto{\pgfqpoint{8.539650in}{0.660000in}}%
\pgfpathclose%
\pgfusepath{fill}%
\end{pgfscope}%
\begin{pgfscope}%
\pgfpathrectangle{\pgfqpoint{1.250000in}{0.660000in}}{\pgfqpoint{7.750000in}{4.620000in}}%
\pgfusepath{clip}%
\pgfsetbuttcap%
\pgfsetmiterjoin%
\definecolor{currentfill}{rgb}{0.000000,0.000000,0.000000}%
\pgfsetfillcolor{currentfill}%
\pgfsetlinewidth{0.000000pt}%
\definecolor{currentstroke}{rgb}{0.000000,0.000000,0.000000}%
\pgfsetstrokecolor{currentstroke}%
\pgfsetstrokeopacity{0.000000}%
\pgfsetdash{}{0pt}%
\pgfpathmoveto{\pgfqpoint{8.547400in}{0.660000in}}%
\pgfpathlineto{\pgfqpoint{8.553600in}{0.660000in}}%
\pgfpathlineto{\pgfqpoint{8.553600in}{3.196230in}}%
\pgfpathlineto{\pgfqpoint{8.547400in}{3.196230in}}%
\pgfpathlineto{\pgfqpoint{8.547400in}{0.660000in}}%
\pgfpathclose%
\pgfusepath{fill}%
\end{pgfscope}%
\begin{pgfscope}%
\pgfpathrectangle{\pgfqpoint{1.250000in}{0.660000in}}{\pgfqpoint{7.750000in}{4.620000in}}%
\pgfusepath{clip}%
\pgfsetbuttcap%
\pgfsetmiterjoin%
\definecolor{currentfill}{rgb}{0.000000,0.000000,0.000000}%
\pgfsetfillcolor{currentfill}%
\pgfsetlinewidth{0.000000pt}%
\definecolor{currentstroke}{rgb}{0.000000,0.000000,0.000000}%
\pgfsetstrokecolor{currentstroke}%
\pgfsetstrokeopacity{0.000000}%
\pgfsetdash{}{0pt}%
\pgfpathmoveto{\pgfqpoint{8.555150in}{0.660000in}}%
\pgfpathlineto{\pgfqpoint{8.561350in}{0.660000in}}%
\pgfpathlineto{\pgfqpoint{8.561350in}{3.196147in}}%
\pgfpathlineto{\pgfqpoint{8.555150in}{3.196147in}}%
\pgfpathlineto{\pgfqpoint{8.555150in}{0.660000in}}%
\pgfpathclose%
\pgfusepath{fill}%
\end{pgfscope}%
\begin{pgfscope}%
\pgfpathrectangle{\pgfqpoint{1.250000in}{0.660000in}}{\pgfqpoint{7.750000in}{4.620000in}}%
\pgfusepath{clip}%
\pgfsetbuttcap%
\pgfsetmiterjoin%
\definecolor{currentfill}{rgb}{0.000000,0.000000,0.000000}%
\pgfsetfillcolor{currentfill}%
\pgfsetlinewidth{0.000000pt}%
\definecolor{currentstroke}{rgb}{0.000000,0.000000,0.000000}%
\pgfsetstrokecolor{currentstroke}%
\pgfsetstrokeopacity{0.000000}%
\pgfsetdash{}{0pt}%
\pgfpathmoveto{\pgfqpoint{8.562900in}{0.660000in}}%
\pgfpathlineto{\pgfqpoint{8.569100in}{0.660000in}}%
\pgfpathlineto{\pgfqpoint{8.569100in}{3.196064in}}%
\pgfpathlineto{\pgfqpoint{8.562900in}{3.196064in}}%
\pgfpathlineto{\pgfqpoint{8.562900in}{0.660000in}}%
\pgfpathclose%
\pgfusepath{fill}%
\end{pgfscope}%
\begin{pgfscope}%
\pgfpathrectangle{\pgfqpoint{1.250000in}{0.660000in}}{\pgfqpoint{7.750000in}{4.620000in}}%
\pgfusepath{clip}%
\pgfsetbuttcap%
\pgfsetmiterjoin%
\definecolor{currentfill}{rgb}{0.000000,0.000000,0.000000}%
\pgfsetfillcolor{currentfill}%
\pgfsetlinewidth{0.000000pt}%
\definecolor{currentstroke}{rgb}{0.000000,0.000000,0.000000}%
\pgfsetstrokecolor{currentstroke}%
\pgfsetstrokeopacity{0.000000}%
\pgfsetdash{}{0pt}%
\pgfpathmoveto{\pgfqpoint{8.570650in}{0.660000in}}%
\pgfpathlineto{\pgfqpoint{8.576850in}{0.660000in}}%
\pgfpathlineto{\pgfqpoint{8.576850in}{3.195980in}}%
\pgfpathlineto{\pgfqpoint{8.570650in}{3.195980in}}%
\pgfpathlineto{\pgfqpoint{8.570650in}{0.660000in}}%
\pgfpathclose%
\pgfusepath{fill}%
\end{pgfscope}%
\begin{pgfscope}%
\pgfpathrectangle{\pgfqpoint{1.250000in}{0.660000in}}{\pgfqpoint{7.750000in}{4.620000in}}%
\pgfusepath{clip}%
\pgfsetbuttcap%
\pgfsetmiterjoin%
\definecolor{currentfill}{rgb}{0.000000,0.000000,0.000000}%
\pgfsetfillcolor{currentfill}%
\pgfsetlinewidth{0.000000pt}%
\definecolor{currentstroke}{rgb}{0.000000,0.000000,0.000000}%
\pgfsetstrokecolor{currentstroke}%
\pgfsetstrokeopacity{0.000000}%
\pgfsetdash{}{0pt}%
\pgfpathmoveto{\pgfqpoint{8.578400in}{0.660000in}}%
\pgfpathlineto{\pgfqpoint{8.584600in}{0.660000in}}%
\pgfpathlineto{\pgfqpoint{8.584600in}{3.195897in}}%
\pgfpathlineto{\pgfqpoint{8.578400in}{3.195897in}}%
\pgfpathlineto{\pgfqpoint{8.578400in}{0.660000in}}%
\pgfpathclose%
\pgfusepath{fill}%
\end{pgfscope}%
\begin{pgfscope}%
\pgfpathrectangle{\pgfqpoint{1.250000in}{0.660000in}}{\pgfqpoint{7.750000in}{4.620000in}}%
\pgfusepath{clip}%
\pgfsetbuttcap%
\pgfsetmiterjoin%
\definecolor{currentfill}{rgb}{0.000000,0.000000,0.000000}%
\pgfsetfillcolor{currentfill}%
\pgfsetlinewidth{0.000000pt}%
\definecolor{currentstroke}{rgb}{0.000000,0.000000,0.000000}%
\pgfsetstrokecolor{currentstroke}%
\pgfsetstrokeopacity{0.000000}%
\pgfsetdash{}{0pt}%
\pgfpathmoveto{\pgfqpoint{8.586150in}{0.660000in}}%
\pgfpathlineto{\pgfqpoint{8.592350in}{0.660000in}}%
\pgfpathlineto{\pgfqpoint{8.592350in}{3.195813in}}%
\pgfpathlineto{\pgfqpoint{8.586150in}{3.195813in}}%
\pgfpathlineto{\pgfqpoint{8.586150in}{0.660000in}}%
\pgfpathclose%
\pgfusepath{fill}%
\end{pgfscope}%
\begin{pgfscope}%
\pgfpathrectangle{\pgfqpoint{1.250000in}{0.660000in}}{\pgfqpoint{7.750000in}{4.620000in}}%
\pgfusepath{clip}%
\pgfsetbuttcap%
\pgfsetmiterjoin%
\definecolor{currentfill}{rgb}{0.000000,0.000000,0.000000}%
\pgfsetfillcolor{currentfill}%
\pgfsetlinewidth{0.000000pt}%
\definecolor{currentstroke}{rgb}{0.000000,0.000000,0.000000}%
\pgfsetstrokecolor{currentstroke}%
\pgfsetstrokeopacity{0.000000}%
\pgfsetdash{}{0pt}%
\pgfpathmoveto{\pgfqpoint{8.593900in}{0.660000in}}%
\pgfpathlineto{\pgfqpoint{8.600100in}{0.660000in}}%
\pgfpathlineto{\pgfqpoint{8.600100in}{3.195730in}}%
\pgfpathlineto{\pgfqpoint{8.593900in}{3.195730in}}%
\pgfpathlineto{\pgfqpoint{8.593900in}{0.660000in}}%
\pgfpathclose%
\pgfusepath{fill}%
\end{pgfscope}%
\begin{pgfscope}%
\pgfpathrectangle{\pgfqpoint{1.250000in}{0.660000in}}{\pgfqpoint{7.750000in}{4.620000in}}%
\pgfusepath{clip}%
\pgfsetbuttcap%
\pgfsetmiterjoin%
\definecolor{currentfill}{rgb}{0.000000,0.000000,0.000000}%
\pgfsetfillcolor{currentfill}%
\pgfsetlinewidth{0.000000pt}%
\definecolor{currentstroke}{rgb}{0.000000,0.000000,0.000000}%
\pgfsetstrokecolor{currentstroke}%
\pgfsetstrokeopacity{0.000000}%
\pgfsetdash{}{0pt}%
\pgfpathmoveto{\pgfqpoint{8.601650in}{0.660000in}}%
\pgfpathlineto{\pgfqpoint{8.607850in}{0.660000in}}%
\pgfpathlineto{\pgfqpoint{8.607850in}{3.195646in}}%
\pgfpathlineto{\pgfqpoint{8.601650in}{3.195646in}}%
\pgfpathlineto{\pgfqpoint{8.601650in}{0.660000in}}%
\pgfpathclose%
\pgfusepath{fill}%
\end{pgfscope}%
\begin{pgfscope}%
\pgfpathrectangle{\pgfqpoint{1.250000in}{0.660000in}}{\pgfqpoint{7.750000in}{4.620000in}}%
\pgfusepath{clip}%
\pgfsetbuttcap%
\pgfsetmiterjoin%
\definecolor{currentfill}{rgb}{0.000000,0.000000,0.000000}%
\pgfsetfillcolor{currentfill}%
\pgfsetlinewidth{0.000000pt}%
\definecolor{currentstroke}{rgb}{0.000000,0.000000,0.000000}%
\pgfsetstrokecolor{currentstroke}%
\pgfsetstrokeopacity{0.000000}%
\pgfsetdash{}{0pt}%
\pgfpathmoveto{\pgfqpoint{8.609400in}{0.660000in}}%
\pgfpathlineto{\pgfqpoint{8.615600in}{0.660000in}}%
\pgfpathlineto{\pgfqpoint{8.615600in}{3.195563in}}%
\pgfpathlineto{\pgfqpoint{8.609400in}{3.195563in}}%
\pgfpathlineto{\pgfqpoint{8.609400in}{0.660000in}}%
\pgfpathclose%
\pgfusepath{fill}%
\end{pgfscope}%
\begin{pgfscope}%
\pgfpathrectangle{\pgfqpoint{1.250000in}{0.660000in}}{\pgfqpoint{7.750000in}{4.620000in}}%
\pgfusepath{clip}%
\pgfsetbuttcap%
\pgfsetmiterjoin%
\definecolor{currentfill}{rgb}{0.000000,0.000000,0.000000}%
\pgfsetfillcolor{currentfill}%
\pgfsetlinewidth{0.000000pt}%
\definecolor{currentstroke}{rgb}{0.000000,0.000000,0.000000}%
\pgfsetstrokecolor{currentstroke}%
\pgfsetstrokeopacity{0.000000}%
\pgfsetdash{}{0pt}%
\pgfpathmoveto{\pgfqpoint{8.617150in}{0.660000in}}%
\pgfpathlineto{\pgfqpoint{8.623350in}{0.660000in}}%
\pgfpathlineto{\pgfqpoint{8.623350in}{3.195480in}}%
\pgfpathlineto{\pgfqpoint{8.617150in}{3.195480in}}%
\pgfpathlineto{\pgfqpoint{8.617150in}{0.660000in}}%
\pgfpathclose%
\pgfusepath{fill}%
\end{pgfscope}%
\begin{pgfscope}%
\pgfpathrectangle{\pgfqpoint{1.250000in}{0.660000in}}{\pgfqpoint{7.750000in}{4.620000in}}%
\pgfusepath{clip}%
\pgfsetbuttcap%
\pgfsetmiterjoin%
\definecolor{currentfill}{rgb}{0.000000,0.000000,0.000000}%
\pgfsetfillcolor{currentfill}%
\pgfsetlinewidth{0.000000pt}%
\definecolor{currentstroke}{rgb}{0.000000,0.000000,0.000000}%
\pgfsetstrokecolor{currentstroke}%
\pgfsetstrokeopacity{0.000000}%
\pgfsetdash{}{0pt}%
\pgfpathmoveto{\pgfqpoint{8.624900in}{0.660000in}}%
\pgfpathlineto{\pgfqpoint{8.631100in}{0.660000in}}%
\pgfpathlineto{\pgfqpoint{8.631100in}{3.195396in}}%
\pgfpathlineto{\pgfqpoint{8.624900in}{3.195396in}}%
\pgfpathlineto{\pgfqpoint{8.624900in}{0.660000in}}%
\pgfpathclose%
\pgfusepath{fill}%
\end{pgfscope}%
\begin{pgfscope}%
\pgfpathrectangle{\pgfqpoint{1.250000in}{0.660000in}}{\pgfqpoint{7.750000in}{4.620000in}}%
\pgfusepath{clip}%
\pgfsetbuttcap%
\pgfsetmiterjoin%
\definecolor{currentfill}{rgb}{0.000000,0.000000,0.000000}%
\pgfsetfillcolor{currentfill}%
\pgfsetlinewidth{0.000000pt}%
\definecolor{currentstroke}{rgb}{0.000000,0.000000,0.000000}%
\pgfsetstrokecolor{currentstroke}%
\pgfsetstrokeopacity{0.000000}%
\pgfsetdash{}{0pt}%
\pgfpathmoveto{\pgfqpoint{8.632650in}{0.660000in}}%
\pgfpathlineto{\pgfqpoint{8.638850in}{0.660000in}}%
\pgfpathlineto{\pgfqpoint{8.638850in}{3.195313in}}%
\pgfpathlineto{\pgfqpoint{8.632650in}{3.195313in}}%
\pgfpathlineto{\pgfqpoint{8.632650in}{0.660000in}}%
\pgfpathclose%
\pgfusepath{fill}%
\end{pgfscope}%
\begin{pgfscope}%
\pgfpathrectangle{\pgfqpoint{1.250000in}{0.660000in}}{\pgfqpoint{7.750000in}{4.620000in}}%
\pgfusepath{clip}%
\pgfsetbuttcap%
\pgfsetmiterjoin%
\definecolor{currentfill}{rgb}{0.000000,0.000000,0.000000}%
\pgfsetfillcolor{currentfill}%
\pgfsetlinewidth{0.000000pt}%
\definecolor{currentstroke}{rgb}{0.000000,0.000000,0.000000}%
\pgfsetstrokecolor{currentstroke}%
\pgfsetstrokeopacity{0.000000}%
\pgfsetdash{}{0pt}%
\pgfpathmoveto{\pgfqpoint{8.640400in}{0.660000in}}%
\pgfpathlineto{\pgfqpoint{8.646600in}{0.660000in}}%
\pgfpathlineto{\pgfqpoint{8.646600in}{3.195229in}}%
\pgfpathlineto{\pgfqpoint{8.640400in}{3.195229in}}%
\pgfpathlineto{\pgfqpoint{8.640400in}{0.660000in}}%
\pgfpathclose%
\pgfusepath{fill}%
\end{pgfscope}%
\begin{pgfscope}%
\pgfpathrectangle{\pgfqpoint{1.250000in}{0.660000in}}{\pgfqpoint{7.750000in}{4.620000in}}%
\pgfusepath{clip}%
\pgfsetbuttcap%
\pgfsetmiterjoin%
\definecolor{currentfill}{rgb}{0.000000,0.000000,0.000000}%
\pgfsetfillcolor{currentfill}%
\pgfsetlinewidth{0.000000pt}%
\definecolor{currentstroke}{rgb}{0.000000,0.000000,0.000000}%
\pgfsetstrokecolor{currentstroke}%
\pgfsetstrokeopacity{0.000000}%
\pgfsetdash{}{0pt}%
\pgfpathmoveto{\pgfqpoint{8.648150in}{0.660000in}}%
\pgfpathlineto{\pgfqpoint{8.654350in}{0.660000in}}%
\pgfpathlineto{\pgfqpoint{8.654350in}{3.195146in}}%
\pgfpathlineto{\pgfqpoint{8.648150in}{3.195146in}}%
\pgfpathlineto{\pgfqpoint{8.648150in}{0.660000in}}%
\pgfpathclose%
\pgfusepath{fill}%
\end{pgfscope}%
\begin{pgfscope}%
\pgfpathrectangle{\pgfqpoint{1.250000in}{0.660000in}}{\pgfqpoint{7.750000in}{4.620000in}}%
\pgfusepath{clip}%
\pgfsetbuttcap%
\pgfsetmiterjoin%
\definecolor{currentfill}{rgb}{0.000000,0.000000,0.000000}%
\pgfsetfillcolor{currentfill}%
\pgfsetlinewidth{0.000000pt}%
\definecolor{currentstroke}{rgb}{0.000000,0.000000,0.000000}%
\pgfsetstrokecolor{currentstroke}%
\pgfsetstrokeopacity{0.000000}%
\pgfsetdash{}{0pt}%
\pgfpathmoveto{\pgfqpoint{8.655900in}{0.660000in}}%
\pgfpathlineto{\pgfqpoint{8.662100in}{0.660000in}}%
\pgfpathlineto{\pgfqpoint{8.662100in}{3.195062in}}%
\pgfpathlineto{\pgfqpoint{8.655900in}{3.195062in}}%
\pgfpathlineto{\pgfqpoint{8.655900in}{0.660000in}}%
\pgfpathclose%
\pgfusepath{fill}%
\end{pgfscope}%
\begin{pgfscope}%
\pgfpathrectangle{\pgfqpoint{1.250000in}{0.660000in}}{\pgfqpoint{7.750000in}{4.620000in}}%
\pgfusepath{clip}%
\pgfsetbuttcap%
\pgfsetmiterjoin%
\definecolor{currentfill}{rgb}{0.000000,0.000000,0.000000}%
\pgfsetfillcolor{currentfill}%
\pgfsetlinewidth{0.000000pt}%
\definecolor{currentstroke}{rgb}{0.000000,0.000000,0.000000}%
\pgfsetstrokecolor{currentstroke}%
\pgfsetstrokeopacity{0.000000}%
\pgfsetdash{}{0pt}%
\pgfpathmoveto{\pgfqpoint{8.663650in}{0.660000in}}%
\pgfpathlineto{\pgfqpoint{8.669850in}{0.660000in}}%
\pgfpathlineto{\pgfqpoint{8.669850in}{3.194979in}}%
\pgfpathlineto{\pgfqpoint{8.663650in}{3.194979in}}%
\pgfpathlineto{\pgfqpoint{8.663650in}{0.660000in}}%
\pgfpathclose%
\pgfusepath{fill}%
\end{pgfscope}%
\begin{pgfscope}%
\pgfpathrectangle{\pgfqpoint{1.250000in}{0.660000in}}{\pgfqpoint{7.750000in}{4.620000in}}%
\pgfusepath{clip}%
\pgfsetbuttcap%
\pgfsetmiterjoin%
\definecolor{currentfill}{rgb}{0.000000,0.000000,0.000000}%
\pgfsetfillcolor{currentfill}%
\pgfsetlinewidth{0.000000pt}%
\definecolor{currentstroke}{rgb}{0.000000,0.000000,0.000000}%
\pgfsetstrokecolor{currentstroke}%
\pgfsetstrokeopacity{0.000000}%
\pgfsetdash{}{0pt}%
\pgfpathmoveto{\pgfqpoint{8.671400in}{0.660000in}}%
\pgfpathlineto{\pgfqpoint{8.677600in}{0.660000in}}%
\pgfpathlineto{\pgfqpoint{8.677600in}{3.194895in}}%
\pgfpathlineto{\pgfqpoint{8.671400in}{3.194895in}}%
\pgfpathlineto{\pgfqpoint{8.671400in}{0.660000in}}%
\pgfpathclose%
\pgfusepath{fill}%
\end{pgfscope}%
\begin{pgfscope}%
\pgfpathrectangle{\pgfqpoint{1.250000in}{0.660000in}}{\pgfqpoint{7.750000in}{4.620000in}}%
\pgfusepath{clip}%
\pgfsetbuttcap%
\pgfsetmiterjoin%
\definecolor{currentfill}{rgb}{0.000000,0.000000,0.000000}%
\pgfsetfillcolor{currentfill}%
\pgfsetlinewidth{0.000000pt}%
\definecolor{currentstroke}{rgb}{0.000000,0.000000,0.000000}%
\pgfsetstrokecolor{currentstroke}%
\pgfsetstrokeopacity{0.000000}%
\pgfsetdash{}{0pt}%
\pgfpathmoveto{\pgfqpoint{8.679150in}{0.660000in}}%
\pgfpathlineto{\pgfqpoint{8.685350in}{0.660000in}}%
\pgfpathlineto{\pgfqpoint{8.685350in}{3.194812in}}%
\pgfpathlineto{\pgfqpoint{8.679150in}{3.194812in}}%
\pgfpathlineto{\pgfqpoint{8.679150in}{0.660000in}}%
\pgfpathclose%
\pgfusepath{fill}%
\end{pgfscope}%
\begin{pgfscope}%
\pgfpathrectangle{\pgfqpoint{1.250000in}{0.660000in}}{\pgfqpoint{7.750000in}{4.620000in}}%
\pgfusepath{clip}%
\pgfsetbuttcap%
\pgfsetmiterjoin%
\definecolor{currentfill}{rgb}{0.000000,0.000000,0.000000}%
\pgfsetfillcolor{currentfill}%
\pgfsetlinewidth{0.000000pt}%
\definecolor{currentstroke}{rgb}{0.000000,0.000000,0.000000}%
\pgfsetstrokecolor{currentstroke}%
\pgfsetstrokeopacity{0.000000}%
\pgfsetdash{}{0pt}%
\pgfpathmoveto{\pgfqpoint{8.686900in}{0.660000in}}%
\pgfpathlineto{\pgfqpoint{8.693100in}{0.660000in}}%
\pgfpathlineto{\pgfqpoint{8.693100in}{3.194729in}}%
\pgfpathlineto{\pgfqpoint{8.686900in}{3.194729in}}%
\pgfpathlineto{\pgfqpoint{8.686900in}{0.660000in}}%
\pgfpathclose%
\pgfusepath{fill}%
\end{pgfscope}%
\begin{pgfscope}%
\pgfpathrectangle{\pgfqpoint{1.250000in}{0.660000in}}{\pgfqpoint{7.750000in}{4.620000in}}%
\pgfusepath{clip}%
\pgfsetbuttcap%
\pgfsetmiterjoin%
\definecolor{currentfill}{rgb}{0.000000,0.000000,0.000000}%
\pgfsetfillcolor{currentfill}%
\pgfsetlinewidth{0.000000pt}%
\definecolor{currentstroke}{rgb}{0.000000,0.000000,0.000000}%
\pgfsetstrokecolor{currentstroke}%
\pgfsetstrokeopacity{0.000000}%
\pgfsetdash{}{0pt}%
\pgfpathmoveto{\pgfqpoint{8.694650in}{0.660000in}}%
\pgfpathlineto{\pgfqpoint{8.700850in}{0.660000in}}%
\pgfpathlineto{\pgfqpoint{8.700850in}{3.194645in}}%
\pgfpathlineto{\pgfqpoint{8.694650in}{3.194645in}}%
\pgfpathlineto{\pgfqpoint{8.694650in}{0.660000in}}%
\pgfpathclose%
\pgfusepath{fill}%
\end{pgfscope}%
\begin{pgfscope}%
\pgfpathrectangle{\pgfqpoint{1.250000in}{0.660000in}}{\pgfqpoint{7.750000in}{4.620000in}}%
\pgfusepath{clip}%
\pgfsetbuttcap%
\pgfsetmiterjoin%
\definecolor{currentfill}{rgb}{0.000000,0.000000,0.000000}%
\pgfsetfillcolor{currentfill}%
\pgfsetlinewidth{0.000000pt}%
\definecolor{currentstroke}{rgb}{0.000000,0.000000,0.000000}%
\pgfsetstrokecolor{currentstroke}%
\pgfsetstrokeopacity{0.000000}%
\pgfsetdash{}{0pt}%
\pgfpathmoveto{\pgfqpoint{8.702400in}{0.660000in}}%
\pgfpathlineto{\pgfqpoint{8.708600in}{0.660000in}}%
\pgfpathlineto{\pgfqpoint{8.708600in}{3.194562in}}%
\pgfpathlineto{\pgfqpoint{8.702400in}{3.194562in}}%
\pgfpathlineto{\pgfqpoint{8.702400in}{0.660000in}}%
\pgfpathclose%
\pgfusepath{fill}%
\end{pgfscope}%
\begin{pgfscope}%
\pgfpathrectangle{\pgfqpoint{1.250000in}{0.660000in}}{\pgfqpoint{7.750000in}{4.620000in}}%
\pgfusepath{clip}%
\pgfsetbuttcap%
\pgfsetmiterjoin%
\definecolor{currentfill}{rgb}{0.000000,0.000000,0.000000}%
\pgfsetfillcolor{currentfill}%
\pgfsetlinewidth{0.000000pt}%
\definecolor{currentstroke}{rgb}{0.000000,0.000000,0.000000}%
\pgfsetstrokecolor{currentstroke}%
\pgfsetstrokeopacity{0.000000}%
\pgfsetdash{}{0pt}%
\pgfpathmoveto{\pgfqpoint{8.710150in}{0.660000in}}%
\pgfpathlineto{\pgfqpoint{8.716350in}{0.660000in}}%
\pgfpathlineto{\pgfqpoint{8.716350in}{3.194478in}}%
\pgfpathlineto{\pgfqpoint{8.710150in}{3.194478in}}%
\pgfpathlineto{\pgfqpoint{8.710150in}{0.660000in}}%
\pgfpathclose%
\pgfusepath{fill}%
\end{pgfscope}%
\begin{pgfscope}%
\pgfpathrectangle{\pgfqpoint{1.250000in}{0.660000in}}{\pgfqpoint{7.750000in}{4.620000in}}%
\pgfusepath{clip}%
\pgfsetbuttcap%
\pgfsetmiterjoin%
\definecolor{currentfill}{rgb}{0.000000,0.000000,0.000000}%
\pgfsetfillcolor{currentfill}%
\pgfsetlinewidth{0.000000pt}%
\definecolor{currentstroke}{rgb}{0.000000,0.000000,0.000000}%
\pgfsetstrokecolor{currentstroke}%
\pgfsetstrokeopacity{0.000000}%
\pgfsetdash{}{0pt}%
\pgfpathmoveto{\pgfqpoint{8.717900in}{0.660000in}}%
\pgfpathlineto{\pgfqpoint{8.724100in}{0.660000in}}%
\pgfpathlineto{\pgfqpoint{8.724100in}{3.194395in}}%
\pgfpathlineto{\pgfqpoint{8.717900in}{3.194395in}}%
\pgfpathlineto{\pgfqpoint{8.717900in}{0.660000in}}%
\pgfpathclose%
\pgfusepath{fill}%
\end{pgfscope}%
\begin{pgfscope}%
\pgfpathrectangle{\pgfqpoint{1.250000in}{0.660000in}}{\pgfqpoint{7.750000in}{4.620000in}}%
\pgfusepath{clip}%
\pgfsetbuttcap%
\pgfsetmiterjoin%
\definecolor{currentfill}{rgb}{0.000000,0.000000,0.000000}%
\pgfsetfillcolor{currentfill}%
\pgfsetlinewidth{0.000000pt}%
\definecolor{currentstroke}{rgb}{0.000000,0.000000,0.000000}%
\pgfsetstrokecolor{currentstroke}%
\pgfsetstrokeopacity{0.000000}%
\pgfsetdash{}{0pt}%
\pgfpathmoveto{\pgfqpoint{8.725650in}{0.660000in}}%
\pgfpathlineto{\pgfqpoint{8.731850in}{0.660000in}}%
\pgfpathlineto{\pgfqpoint{8.731850in}{3.194311in}}%
\pgfpathlineto{\pgfqpoint{8.725650in}{3.194311in}}%
\pgfpathlineto{\pgfqpoint{8.725650in}{0.660000in}}%
\pgfpathclose%
\pgfusepath{fill}%
\end{pgfscope}%
\begin{pgfscope}%
\pgfpathrectangle{\pgfqpoint{1.250000in}{0.660000in}}{\pgfqpoint{7.750000in}{4.620000in}}%
\pgfusepath{clip}%
\pgfsetbuttcap%
\pgfsetmiterjoin%
\definecolor{currentfill}{rgb}{0.000000,0.000000,0.000000}%
\pgfsetfillcolor{currentfill}%
\pgfsetlinewidth{0.000000pt}%
\definecolor{currentstroke}{rgb}{0.000000,0.000000,0.000000}%
\pgfsetstrokecolor{currentstroke}%
\pgfsetstrokeopacity{0.000000}%
\pgfsetdash{}{0pt}%
\pgfpathmoveto{\pgfqpoint{8.733400in}{0.660000in}}%
\pgfpathlineto{\pgfqpoint{8.739600in}{0.660000in}}%
\pgfpathlineto{\pgfqpoint{8.739600in}{3.194228in}}%
\pgfpathlineto{\pgfqpoint{8.733400in}{3.194228in}}%
\pgfpathlineto{\pgfqpoint{8.733400in}{0.660000in}}%
\pgfpathclose%
\pgfusepath{fill}%
\end{pgfscope}%
\begin{pgfscope}%
\pgfpathrectangle{\pgfqpoint{1.250000in}{0.660000in}}{\pgfqpoint{7.750000in}{4.620000in}}%
\pgfusepath{clip}%
\pgfsetbuttcap%
\pgfsetmiterjoin%
\definecolor{currentfill}{rgb}{0.000000,0.000000,0.000000}%
\pgfsetfillcolor{currentfill}%
\pgfsetlinewidth{0.000000pt}%
\definecolor{currentstroke}{rgb}{0.000000,0.000000,0.000000}%
\pgfsetstrokecolor{currentstroke}%
\pgfsetstrokeopacity{0.000000}%
\pgfsetdash{}{0pt}%
\pgfpathmoveto{\pgfqpoint{8.741150in}{0.660000in}}%
\pgfpathlineto{\pgfqpoint{8.747350in}{0.660000in}}%
\pgfpathlineto{\pgfqpoint{8.747350in}{3.194144in}}%
\pgfpathlineto{\pgfqpoint{8.741150in}{3.194144in}}%
\pgfpathlineto{\pgfqpoint{8.741150in}{0.660000in}}%
\pgfpathclose%
\pgfusepath{fill}%
\end{pgfscope}%
\begin{pgfscope}%
\pgfpathrectangle{\pgfqpoint{1.250000in}{0.660000in}}{\pgfqpoint{7.750000in}{4.620000in}}%
\pgfusepath{clip}%
\pgfsetbuttcap%
\pgfsetmiterjoin%
\definecolor{currentfill}{rgb}{0.000000,0.000000,0.000000}%
\pgfsetfillcolor{currentfill}%
\pgfsetlinewidth{0.000000pt}%
\definecolor{currentstroke}{rgb}{0.000000,0.000000,0.000000}%
\pgfsetstrokecolor{currentstroke}%
\pgfsetstrokeopacity{0.000000}%
\pgfsetdash{}{0pt}%
\pgfpathmoveto{\pgfqpoint{8.748900in}{0.660000in}}%
\pgfpathlineto{\pgfqpoint{8.755100in}{0.660000in}}%
\pgfpathlineto{\pgfqpoint{8.755100in}{3.194064in}}%
\pgfpathlineto{\pgfqpoint{8.748900in}{3.194064in}}%
\pgfpathlineto{\pgfqpoint{8.748900in}{0.660000in}}%
\pgfpathclose%
\pgfusepath{fill}%
\end{pgfscope}%
\begin{pgfscope}%
\pgfpathrectangle{\pgfqpoint{1.250000in}{0.660000in}}{\pgfqpoint{7.750000in}{4.620000in}}%
\pgfusepath{clip}%
\pgfsetbuttcap%
\pgfsetmiterjoin%
\definecolor{currentfill}{rgb}{0.000000,0.000000,0.000000}%
\pgfsetfillcolor{currentfill}%
\pgfsetlinewidth{0.000000pt}%
\definecolor{currentstroke}{rgb}{0.000000,0.000000,0.000000}%
\pgfsetstrokecolor{currentstroke}%
\pgfsetstrokeopacity{0.000000}%
\pgfsetdash{}{0pt}%
\pgfpathmoveto{\pgfqpoint{8.756650in}{0.660000in}}%
\pgfpathlineto{\pgfqpoint{8.762850in}{0.660000in}}%
\pgfpathlineto{\pgfqpoint{8.762850in}{3.193983in}}%
\pgfpathlineto{\pgfqpoint{8.756650in}{3.193983in}}%
\pgfpathlineto{\pgfqpoint{8.756650in}{0.660000in}}%
\pgfpathclose%
\pgfusepath{fill}%
\end{pgfscope}%
\begin{pgfscope}%
\pgfpathrectangle{\pgfqpoint{1.250000in}{0.660000in}}{\pgfqpoint{7.750000in}{4.620000in}}%
\pgfusepath{clip}%
\pgfsetbuttcap%
\pgfsetmiterjoin%
\definecolor{currentfill}{rgb}{0.000000,0.000000,0.000000}%
\pgfsetfillcolor{currentfill}%
\pgfsetlinewidth{0.000000pt}%
\definecolor{currentstroke}{rgb}{0.000000,0.000000,0.000000}%
\pgfsetstrokecolor{currentstroke}%
\pgfsetstrokeopacity{0.000000}%
\pgfsetdash{}{0pt}%
\pgfpathmoveto{\pgfqpoint{8.764400in}{0.660000in}}%
\pgfpathlineto{\pgfqpoint{8.770600in}{0.660000in}}%
\pgfpathlineto{\pgfqpoint{8.770600in}{3.193903in}}%
\pgfpathlineto{\pgfqpoint{8.764400in}{3.193903in}}%
\pgfpathlineto{\pgfqpoint{8.764400in}{0.660000in}}%
\pgfpathclose%
\pgfusepath{fill}%
\end{pgfscope}%
\begin{pgfscope}%
\pgfpathrectangle{\pgfqpoint{1.250000in}{0.660000in}}{\pgfqpoint{7.750000in}{4.620000in}}%
\pgfusepath{clip}%
\pgfsetbuttcap%
\pgfsetmiterjoin%
\definecolor{currentfill}{rgb}{0.000000,0.000000,0.000000}%
\pgfsetfillcolor{currentfill}%
\pgfsetlinewidth{0.000000pt}%
\definecolor{currentstroke}{rgb}{0.000000,0.000000,0.000000}%
\pgfsetstrokecolor{currentstroke}%
\pgfsetstrokeopacity{0.000000}%
\pgfsetdash{}{0pt}%
\pgfpathmoveto{\pgfqpoint{8.772150in}{0.660000in}}%
\pgfpathlineto{\pgfqpoint{8.778350in}{0.660000in}}%
\pgfpathlineto{\pgfqpoint{8.778350in}{3.193822in}}%
\pgfpathlineto{\pgfqpoint{8.772150in}{3.193822in}}%
\pgfpathlineto{\pgfqpoint{8.772150in}{0.660000in}}%
\pgfpathclose%
\pgfusepath{fill}%
\end{pgfscope}%
\begin{pgfscope}%
\pgfpathrectangle{\pgfqpoint{1.250000in}{0.660000in}}{\pgfqpoint{7.750000in}{4.620000in}}%
\pgfusepath{clip}%
\pgfsetbuttcap%
\pgfsetmiterjoin%
\definecolor{currentfill}{rgb}{0.000000,0.000000,0.000000}%
\pgfsetfillcolor{currentfill}%
\pgfsetlinewidth{0.000000pt}%
\definecolor{currentstroke}{rgb}{0.000000,0.000000,0.000000}%
\pgfsetstrokecolor{currentstroke}%
\pgfsetstrokeopacity{0.000000}%
\pgfsetdash{}{0pt}%
\pgfpathmoveto{\pgfqpoint{8.779900in}{0.660000in}}%
\pgfpathlineto{\pgfqpoint{8.786100in}{0.660000in}}%
\pgfpathlineto{\pgfqpoint{8.786100in}{3.193742in}}%
\pgfpathlineto{\pgfqpoint{8.779900in}{3.193742in}}%
\pgfpathlineto{\pgfqpoint{8.779900in}{0.660000in}}%
\pgfpathclose%
\pgfusepath{fill}%
\end{pgfscope}%
\begin{pgfscope}%
\pgfpathrectangle{\pgfqpoint{1.250000in}{0.660000in}}{\pgfqpoint{7.750000in}{4.620000in}}%
\pgfusepath{clip}%
\pgfsetbuttcap%
\pgfsetmiterjoin%
\definecolor{currentfill}{rgb}{0.000000,0.000000,0.000000}%
\pgfsetfillcolor{currentfill}%
\pgfsetlinewidth{0.000000pt}%
\definecolor{currentstroke}{rgb}{0.000000,0.000000,0.000000}%
\pgfsetstrokecolor{currentstroke}%
\pgfsetstrokeopacity{0.000000}%
\pgfsetdash{}{0pt}%
\pgfpathmoveto{\pgfqpoint{8.787650in}{0.660000in}}%
\pgfpathlineto{\pgfqpoint{8.793850in}{0.660000in}}%
\pgfpathlineto{\pgfqpoint{8.793850in}{3.193662in}}%
\pgfpathlineto{\pgfqpoint{8.787650in}{3.193662in}}%
\pgfpathlineto{\pgfqpoint{8.787650in}{0.660000in}}%
\pgfpathclose%
\pgfusepath{fill}%
\end{pgfscope}%
\begin{pgfscope}%
\pgfpathrectangle{\pgfqpoint{1.250000in}{0.660000in}}{\pgfqpoint{7.750000in}{4.620000in}}%
\pgfusepath{clip}%
\pgfsetbuttcap%
\pgfsetmiterjoin%
\definecolor{currentfill}{rgb}{0.000000,0.000000,0.000000}%
\pgfsetfillcolor{currentfill}%
\pgfsetlinewidth{0.000000pt}%
\definecolor{currentstroke}{rgb}{0.000000,0.000000,0.000000}%
\pgfsetstrokecolor{currentstroke}%
\pgfsetstrokeopacity{0.000000}%
\pgfsetdash{}{0pt}%
\pgfpathmoveto{\pgfqpoint{8.795400in}{0.660000in}}%
\pgfpathlineto{\pgfqpoint{8.801600in}{0.660000in}}%
\pgfpathlineto{\pgfqpoint{8.801600in}{3.193581in}}%
\pgfpathlineto{\pgfqpoint{8.795400in}{3.193581in}}%
\pgfpathlineto{\pgfqpoint{8.795400in}{0.660000in}}%
\pgfpathclose%
\pgfusepath{fill}%
\end{pgfscope}%
\begin{pgfscope}%
\pgfpathrectangle{\pgfqpoint{1.250000in}{0.660000in}}{\pgfqpoint{7.750000in}{4.620000in}}%
\pgfusepath{clip}%
\pgfsetbuttcap%
\pgfsetmiterjoin%
\definecolor{currentfill}{rgb}{0.000000,0.000000,0.000000}%
\pgfsetfillcolor{currentfill}%
\pgfsetlinewidth{0.000000pt}%
\definecolor{currentstroke}{rgb}{0.000000,0.000000,0.000000}%
\pgfsetstrokecolor{currentstroke}%
\pgfsetstrokeopacity{0.000000}%
\pgfsetdash{}{0pt}%
\pgfpathmoveto{\pgfqpoint{8.803150in}{0.660000in}}%
\pgfpathlineto{\pgfqpoint{8.809350in}{0.660000in}}%
\pgfpathlineto{\pgfqpoint{8.809350in}{3.193501in}}%
\pgfpathlineto{\pgfqpoint{8.803150in}{3.193501in}}%
\pgfpathlineto{\pgfqpoint{8.803150in}{0.660000in}}%
\pgfpathclose%
\pgfusepath{fill}%
\end{pgfscope}%
\begin{pgfscope}%
\pgfpathrectangle{\pgfqpoint{1.250000in}{0.660000in}}{\pgfqpoint{7.750000in}{4.620000in}}%
\pgfusepath{clip}%
\pgfsetbuttcap%
\pgfsetmiterjoin%
\definecolor{currentfill}{rgb}{0.000000,0.000000,0.000000}%
\pgfsetfillcolor{currentfill}%
\pgfsetlinewidth{0.000000pt}%
\definecolor{currentstroke}{rgb}{0.000000,0.000000,0.000000}%
\pgfsetstrokecolor{currentstroke}%
\pgfsetstrokeopacity{0.000000}%
\pgfsetdash{}{0pt}%
\pgfpathmoveto{\pgfqpoint{8.810900in}{0.660000in}}%
\pgfpathlineto{\pgfqpoint{8.817100in}{0.660000in}}%
\pgfpathlineto{\pgfqpoint{8.817100in}{3.193420in}}%
\pgfpathlineto{\pgfqpoint{8.810900in}{3.193420in}}%
\pgfpathlineto{\pgfqpoint{8.810900in}{0.660000in}}%
\pgfpathclose%
\pgfusepath{fill}%
\end{pgfscope}%
\begin{pgfscope}%
\pgfpathrectangle{\pgfqpoint{1.250000in}{0.660000in}}{\pgfqpoint{7.750000in}{4.620000in}}%
\pgfusepath{clip}%
\pgfsetbuttcap%
\pgfsetmiterjoin%
\definecolor{currentfill}{rgb}{0.000000,0.000000,0.000000}%
\pgfsetfillcolor{currentfill}%
\pgfsetlinewidth{0.000000pt}%
\definecolor{currentstroke}{rgb}{0.000000,0.000000,0.000000}%
\pgfsetstrokecolor{currentstroke}%
\pgfsetstrokeopacity{0.000000}%
\pgfsetdash{}{0pt}%
\pgfpathmoveto{\pgfqpoint{8.818650in}{0.660000in}}%
\pgfpathlineto{\pgfqpoint{8.824850in}{0.660000in}}%
\pgfpathlineto{\pgfqpoint{8.824850in}{3.193340in}}%
\pgfpathlineto{\pgfqpoint{8.818650in}{3.193340in}}%
\pgfpathlineto{\pgfqpoint{8.818650in}{0.660000in}}%
\pgfpathclose%
\pgfusepath{fill}%
\end{pgfscope}%
\begin{pgfscope}%
\pgfpathrectangle{\pgfqpoint{1.250000in}{0.660000in}}{\pgfqpoint{7.750000in}{4.620000in}}%
\pgfusepath{clip}%
\pgfsetbuttcap%
\pgfsetmiterjoin%
\definecolor{currentfill}{rgb}{0.000000,0.000000,0.000000}%
\pgfsetfillcolor{currentfill}%
\pgfsetlinewidth{0.000000pt}%
\definecolor{currentstroke}{rgb}{0.000000,0.000000,0.000000}%
\pgfsetstrokecolor{currentstroke}%
\pgfsetstrokeopacity{0.000000}%
\pgfsetdash{}{0pt}%
\pgfpathmoveto{\pgfqpoint{8.826400in}{0.660000in}}%
\pgfpathlineto{\pgfqpoint{8.832600in}{0.660000in}}%
\pgfpathlineto{\pgfqpoint{8.832600in}{3.193260in}}%
\pgfpathlineto{\pgfqpoint{8.826400in}{3.193260in}}%
\pgfpathlineto{\pgfqpoint{8.826400in}{0.660000in}}%
\pgfpathclose%
\pgfusepath{fill}%
\end{pgfscope}%
\begin{pgfscope}%
\pgfpathrectangle{\pgfqpoint{1.250000in}{0.660000in}}{\pgfqpoint{7.750000in}{4.620000in}}%
\pgfusepath{clip}%
\pgfsetbuttcap%
\pgfsetmiterjoin%
\definecolor{currentfill}{rgb}{0.000000,0.000000,0.000000}%
\pgfsetfillcolor{currentfill}%
\pgfsetlinewidth{0.000000pt}%
\definecolor{currentstroke}{rgb}{0.000000,0.000000,0.000000}%
\pgfsetstrokecolor{currentstroke}%
\pgfsetstrokeopacity{0.000000}%
\pgfsetdash{}{0pt}%
\pgfpathmoveto{\pgfqpoint{8.834150in}{0.660000in}}%
\pgfpathlineto{\pgfqpoint{8.840350in}{0.660000in}}%
\pgfpathlineto{\pgfqpoint{8.840350in}{3.193179in}}%
\pgfpathlineto{\pgfqpoint{8.834150in}{3.193179in}}%
\pgfpathlineto{\pgfqpoint{8.834150in}{0.660000in}}%
\pgfpathclose%
\pgfusepath{fill}%
\end{pgfscope}%
\begin{pgfscope}%
\pgfpathrectangle{\pgfqpoint{1.250000in}{0.660000in}}{\pgfqpoint{7.750000in}{4.620000in}}%
\pgfusepath{clip}%
\pgfsetbuttcap%
\pgfsetmiterjoin%
\definecolor{currentfill}{rgb}{0.000000,0.000000,0.000000}%
\pgfsetfillcolor{currentfill}%
\pgfsetlinewidth{0.000000pt}%
\definecolor{currentstroke}{rgb}{0.000000,0.000000,0.000000}%
\pgfsetstrokecolor{currentstroke}%
\pgfsetstrokeopacity{0.000000}%
\pgfsetdash{}{0pt}%
\pgfpathmoveto{\pgfqpoint{8.841900in}{0.660000in}}%
\pgfpathlineto{\pgfqpoint{8.848100in}{0.660000in}}%
\pgfpathlineto{\pgfqpoint{8.848100in}{3.193099in}}%
\pgfpathlineto{\pgfqpoint{8.841900in}{3.193099in}}%
\pgfpathlineto{\pgfqpoint{8.841900in}{0.660000in}}%
\pgfpathclose%
\pgfusepath{fill}%
\end{pgfscope}%
\begin{pgfscope}%
\pgfpathrectangle{\pgfqpoint{1.250000in}{0.660000in}}{\pgfqpoint{7.750000in}{4.620000in}}%
\pgfusepath{clip}%
\pgfsetbuttcap%
\pgfsetmiterjoin%
\definecolor{currentfill}{rgb}{0.000000,0.000000,0.000000}%
\pgfsetfillcolor{currentfill}%
\pgfsetlinewidth{0.000000pt}%
\definecolor{currentstroke}{rgb}{0.000000,0.000000,0.000000}%
\pgfsetstrokecolor{currentstroke}%
\pgfsetstrokeopacity{0.000000}%
\pgfsetdash{}{0pt}%
\pgfpathmoveto{\pgfqpoint{8.849650in}{0.660000in}}%
\pgfpathlineto{\pgfqpoint{8.855850in}{0.660000in}}%
\pgfpathlineto{\pgfqpoint{8.855850in}{3.193019in}}%
\pgfpathlineto{\pgfqpoint{8.849650in}{3.193019in}}%
\pgfpathlineto{\pgfqpoint{8.849650in}{0.660000in}}%
\pgfpathclose%
\pgfusepath{fill}%
\end{pgfscope}%
\begin{pgfscope}%
\pgfpathrectangle{\pgfqpoint{1.250000in}{0.660000in}}{\pgfqpoint{7.750000in}{4.620000in}}%
\pgfusepath{clip}%
\pgfsetbuttcap%
\pgfsetmiterjoin%
\definecolor{currentfill}{rgb}{0.000000,0.000000,0.000000}%
\pgfsetfillcolor{currentfill}%
\pgfsetlinewidth{0.000000pt}%
\definecolor{currentstroke}{rgb}{0.000000,0.000000,0.000000}%
\pgfsetstrokecolor{currentstroke}%
\pgfsetstrokeopacity{0.000000}%
\pgfsetdash{}{0pt}%
\pgfpathmoveto{\pgfqpoint{8.857400in}{0.660000in}}%
\pgfpathlineto{\pgfqpoint{8.863600in}{0.660000in}}%
\pgfpathlineto{\pgfqpoint{8.863600in}{3.192938in}}%
\pgfpathlineto{\pgfqpoint{8.857400in}{3.192938in}}%
\pgfpathlineto{\pgfqpoint{8.857400in}{0.660000in}}%
\pgfpathclose%
\pgfusepath{fill}%
\end{pgfscope}%
\begin{pgfscope}%
\pgfpathrectangle{\pgfqpoint{1.250000in}{0.660000in}}{\pgfqpoint{7.750000in}{4.620000in}}%
\pgfusepath{clip}%
\pgfsetbuttcap%
\pgfsetmiterjoin%
\definecolor{currentfill}{rgb}{0.000000,0.000000,0.000000}%
\pgfsetfillcolor{currentfill}%
\pgfsetlinewidth{0.000000pt}%
\definecolor{currentstroke}{rgb}{0.000000,0.000000,0.000000}%
\pgfsetstrokecolor{currentstroke}%
\pgfsetstrokeopacity{0.000000}%
\pgfsetdash{}{0pt}%
\pgfpathmoveto{\pgfqpoint{8.865150in}{0.660000in}}%
\pgfpathlineto{\pgfqpoint{8.871350in}{0.660000in}}%
\pgfpathlineto{\pgfqpoint{8.871350in}{3.192858in}}%
\pgfpathlineto{\pgfqpoint{8.865150in}{3.192858in}}%
\pgfpathlineto{\pgfqpoint{8.865150in}{0.660000in}}%
\pgfpathclose%
\pgfusepath{fill}%
\end{pgfscope}%
\begin{pgfscope}%
\pgfpathrectangle{\pgfqpoint{1.250000in}{0.660000in}}{\pgfqpoint{7.750000in}{4.620000in}}%
\pgfusepath{clip}%
\pgfsetbuttcap%
\pgfsetmiterjoin%
\definecolor{currentfill}{rgb}{0.000000,0.000000,0.000000}%
\pgfsetfillcolor{currentfill}%
\pgfsetlinewidth{0.000000pt}%
\definecolor{currentstroke}{rgb}{0.000000,0.000000,0.000000}%
\pgfsetstrokecolor{currentstroke}%
\pgfsetstrokeopacity{0.000000}%
\pgfsetdash{}{0pt}%
\pgfpathmoveto{\pgfqpoint{8.872900in}{0.660000in}}%
\pgfpathlineto{\pgfqpoint{8.879100in}{0.660000in}}%
\pgfpathlineto{\pgfqpoint{8.879100in}{3.192777in}}%
\pgfpathlineto{\pgfqpoint{8.872900in}{3.192777in}}%
\pgfpathlineto{\pgfqpoint{8.872900in}{0.660000in}}%
\pgfpathclose%
\pgfusepath{fill}%
\end{pgfscope}%
\begin{pgfscope}%
\pgfpathrectangle{\pgfqpoint{1.250000in}{0.660000in}}{\pgfqpoint{7.750000in}{4.620000in}}%
\pgfusepath{clip}%
\pgfsetbuttcap%
\pgfsetmiterjoin%
\definecolor{currentfill}{rgb}{0.000000,0.000000,0.000000}%
\pgfsetfillcolor{currentfill}%
\pgfsetlinewidth{0.000000pt}%
\definecolor{currentstroke}{rgb}{0.000000,0.000000,0.000000}%
\pgfsetstrokecolor{currentstroke}%
\pgfsetstrokeopacity{0.000000}%
\pgfsetdash{}{0pt}%
\pgfpathmoveto{\pgfqpoint{8.880650in}{0.660000in}}%
\pgfpathlineto{\pgfqpoint{8.886850in}{0.660000in}}%
\pgfpathlineto{\pgfqpoint{8.886850in}{3.192697in}}%
\pgfpathlineto{\pgfqpoint{8.880650in}{3.192697in}}%
\pgfpathlineto{\pgfqpoint{8.880650in}{0.660000in}}%
\pgfpathclose%
\pgfusepath{fill}%
\end{pgfscope}%
\begin{pgfscope}%
\pgfpathrectangle{\pgfqpoint{1.250000in}{0.660000in}}{\pgfqpoint{7.750000in}{4.620000in}}%
\pgfusepath{clip}%
\pgfsetbuttcap%
\pgfsetmiterjoin%
\definecolor{currentfill}{rgb}{0.000000,0.000000,0.000000}%
\pgfsetfillcolor{currentfill}%
\pgfsetlinewidth{0.000000pt}%
\definecolor{currentstroke}{rgb}{0.000000,0.000000,0.000000}%
\pgfsetstrokecolor{currentstroke}%
\pgfsetstrokeopacity{0.000000}%
\pgfsetdash{}{0pt}%
\pgfpathmoveto{\pgfqpoint{8.888400in}{0.660000in}}%
\pgfpathlineto{\pgfqpoint{8.894600in}{0.660000in}}%
\pgfpathlineto{\pgfqpoint{8.894600in}{3.192617in}}%
\pgfpathlineto{\pgfqpoint{8.888400in}{3.192617in}}%
\pgfpathlineto{\pgfqpoint{8.888400in}{0.660000in}}%
\pgfpathclose%
\pgfusepath{fill}%
\end{pgfscope}%
\begin{pgfscope}%
\pgfpathrectangle{\pgfqpoint{1.250000in}{0.660000in}}{\pgfqpoint{7.750000in}{4.620000in}}%
\pgfusepath{clip}%
\pgfsetbuttcap%
\pgfsetmiterjoin%
\definecolor{currentfill}{rgb}{0.000000,0.000000,0.000000}%
\pgfsetfillcolor{currentfill}%
\pgfsetlinewidth{0.000000pt}%
\definecolor{currentstroke}{rgb}{0.000000,0.000000,0.000000}%
\pgfsetstrokecolor{currentstroke}%
\pgfsetstrokeopacity{0.000000}%
\pgfsetdash{}{0pt}%
\pgfpathmoveto{\pgfqpoint{8.896150in}{0.660000in}}%
\pgfpathlineto{\pgfqpoint{8.902350in}{0.660000in}}%
\pgfpathlineto{\pgfqpoint{8.902350in}{3.192536in}}%
\pgfpathlineto{\pgfqpoint{8.896150in}{3.192536in}}%
\pgfpathlineto{\pgfqpoint{8.896150in}{0.660000in}}%
\pgfpathclose%
\pgfusepath{fill}%
\end{pgfscope}%
\begin{pgfscope}%
\pgfpathrectangle{\pgfqpoint{1.250000in}{0.660000in}}{\pgfqpoint{7.750000in}{4.620000in}}%
\pgfusepath{clip}%
\pgfsetbuttcap%
\pgfsetmiterjoin%
\definecolor{currentfill}{rgb}{0.000000,0.000000,0.000000}%
\pgfsetfillcolor{currentfill}%
\pgfsetlinewidth{0.000000pt}%
\definecolor{currentstroke}{rgb}{0.000000,0.000000,0.000000}%
\pgfsetstrokecolor{currentstroke}%
\pgfsetstrokeopacity{0.000000}%
\pgfsetdash{}{0pt}%
\pgfpathmoveto{\pgfqpoint{8.903900in}{0.660000in}}%
\pgfpathlineto{\pgfqpoint{8.910100in}{0.660000in}}%
\pgfpathlineto{\pgfqpoint{8.910100in}{3.192456in}}%
\pgfpathlineto{\pgfqpoint{8.903900in}{3.192456in}}%
\pgfpathlineto{\pgfqpoint{8.903900in}{0.660000in}}%
\pgfpathclose%
\pgfusepath{fill}%
\end{pgfscope}%
\begin{pgfscope}%
\pgfpathrectangle{\pgfqpoint{1.250000in}{0.660000in}}{\pgfqpoint{7.750000in}{4.620000in}}%
\pgfusepath{clip}%
\pgfsetbuttcap%
\pgfsetmiterjoin%
\definecolor{currentfill}{rgb}{0.000000,0.000000,0.000000}%
\pgfsetfillcolor{currentfill}%
\pgfsetlinewidth{0.000000pt}%
\definecolor{currentstroke}{rgb}{0.000000,0.000000,0.000000}%
\pgfsetstrokecolor{currentstroke}%
\pgfsetstrokeopacity{0.000000}%
\pgfsetdash{}{0pt}%
\pgfpathmoveto{\pgfqpoint{8.911650in}{0.660000in}}%
\pgfpathlineto{\pgfqpoint{8.917850in}{0.660000in}}%
\pgfpathlineto{\pgfqpoint{8.917850in}{3.192375in}}%
\pgfpathlineto{\pgfqpoint{8.911650in}{3.192375in}}%
\pgfpathlineto{\pgfqpoint{8.911650in}{0.660000in}}%
\pgfpathclose%
\pgfusepath{fill}%
\end{pgfscope}%
\begin{pgfscope}%
\pgfpathrectangle{\pgfqpoint{1.250000in}{0.660000in}}{\pgfqpoint{7.750000in}{4.620000in}}%
\pgfusepath{clip}%
\pgfsetbuttcap%
\pgfsetmiterjoin%
\definecolor{currentfill}{rgb}{0.000000,0.000000,0.000000}%
\pgfsetfillcolor{currentfill}%
\pgfsetlinewidth{0.000000pt}%
\definecolor{currentstroke}{rgb}{0.000000,0.000000,0.000000}%
\pgfsetstrokecolor{currentstroke}%
\pgfsetstrokeopacity{0.000000}%
\pgfsetdash{}{0pt}%
\pgfpathmoveto{\pgfqpoint{8.919400in}{0.660000in}}%
\pgfpathlineto{\pgfqpoint{8.925600in}{0.660000in}}%
\pgfpathlineto{\pgfqpoint{8.925600in}{3.192295in}}%
\pgfpathlineto{\pgfqpoint{8.919400in}{3.192295in}}%
\pgfpathlineto{\pgfqpoint{8.919400in}{0.660000in}}%
\pgfpathclose%
\pgfusepath{fill}%
\end{pgfscope}%
\begin{pgfscope}%
\pgfpathrectangle{\pgfqpoint{1.250000in}{0.660000in}}{\pgfqpoint{7.750000in}{4.620000in}}%
\pgfusepath{clip}%
\pgfsetbuttcap%
\pgfsetmiterjoin%
\definecolor{currentfill}{rgb}{0.000000,0.000000,0.000000}%
\pgfsetfillcolor{currentfill}%
\pgfsetlinewidth{0.000000pt}%
\definecolor{currentstroke}{rgb}{0.000000,0.000000,0.000000}%
\pgfsetstrokecolor{currentstroke}%
\pgfsetstrokeopacity{0.000000}%
\pgfsetdash{}{0pt}%
\pgfpathmoveto{\pgfqpoint{8.927150in}{0.660000in}}%
\pgfpathlineto{\pgfqpoint{8.933350in}{0.660000in}}%
\pgfpathlineto{\pgfqpoint{8.933350in}{3.192215in}}%
\pgfpathlineto{\pgfqpoint{8.927150in}{3.192215in}}%
\pgfpathlineto{\pgfqpoint{8.927150in}{0.660000in}}%
\pgfpathclose%
\pgfusepath{fill}%
\end{pgfscope}%
\begin{pgfscope}%
\pgfpathrectangle{\pgfqpoint{1.250000in}{0.660000in}}{\pgfqpoint{7.750000in}{4.620000in}}%
\pgfusepath{clip}%
\pgfsetbuttcap%
\pgfsetmiterjoin%
\definecolor{currentfill}{rgb}{0.000000,0.000000,0.000000}%
\pgfsetfillcolor{currentfill}%
\pgfsetlinewidth{0.000000pt}%
\definecolor{currentstroke}{rgb}{0.000000,0.000000,0.000000}%
\pgfsetstrokecolor{currentstroke}%
\pgfsetstrokeopacity{0.000000}%
\pgfsetdash{}{0pt}%
\pgfpathmoveto{\pgfqpoint{8.934900in}{0.660000in}}%
\pgfpathlineto{\pgfqpoint{8.941100in}{0.660000in}}%
\pgfpathlineto{\pgfqpoint{8.941100in}{3.192134in}}%
\pgfpathlineto{\pgfqpoint{8.934900in}{3.192134in}}%
\pgfpathlineto{\pgfqpoint{8.934900in}{0.660000in}}%
\pgfpathclose%
\pgfusepath{fill}%
\end{pgfscope}%
\begin{pgfscope}%
\pgfpathrectangle{\pgfqpoint{1.250000in}{0.660000in}}{\pgfqpoint{7.750000in}{4.620000in}}%
\pgfusepath{clip}%
\pgfsetbuttcap%
\pgfsetmiterjoin%
\definecolor{currentfill}{rgb}{0.000000,0.000000,0.000000}%
\pgfsetfillcolor{currentfill}%
\pgfsetlinewidth{0.000000pt}%
\definecolor{currentstroke}{rgb}{0.000000,0.000000,0.000000}%
\pgfsetstrokecolor{currentstroke}%
\pgfsetstrokeopacity{0.000000}%
\pgfsetdash{}{0pt}%
\pgfpathmoveto{\pgfqpoint{8.942650in}{0.660000in}}%
\pgfpathlineto{\pgfqpoint{8.948850in}{0.660000in}}%
\pgfpathlineto{\pgfqpoint{8.948850in}{3.192054in}}%
\pgfpathlineto{\pgfqpoint{8.942650in}{3.192054in}}%
\pgfpathlineto{\pgfqpoint{8.942650in}{0.660000in}}%
\pgfpathclose%
\pgfusepath{fill}%
\end{pgfscope}%
\begin{pgfscope}%
\pgfpathrectangle{\pgfqpoint{1.250000in}{0.660000in}}{\pgfqpoint{7.750000in}{4.620000in}}%
\pgfusepath{clip}%
\pgfsetbuttcap%
\pgfsetmiterjoin%
\definecolor{currentfill}{rgb}{0.000000,0.000000,0.000000}%
\pgfsetfillcolor{currentfill}%
\pgfsetlinewidth{0.000000pt}%
\definecolor{currentstroke}{rgb}{0.000000,0.000000,0.000000}%
\pgfsetstrokecolor{currentstroke}%
\pgfsetstrokeopacity{0.000000}%
\pgfsetdash{}{0pt}%
\pgfpathmoveto{\pgfqpoint{8.950400in}{0.660000in}}%
\pgfpathlineto{\pgfqpoint{8.956600in}{0.660000in}}%
\pgfpathlineto{\pgfqpoint{8.956600in}{3.191973in}}%
\pgfpathlineto{\pgfqpoint{8.950400in}{3.191973in}}%
\pgfpathlineto{\pgfqpoint{8.950400in}{0.660000in}}%
\pgfpathclose%
\pgfusepath{fill}%
\end{pgfscope}%
\begin{pgfscope}%
\pgfpathrectangle{\pgfqpoint{1.250000in}{0.660000in}}{\pgfqpoint{7.750000in}{4.620000in}}%
\pgfusepath{clip}%
\pgfsetbuttcap%
\pgfsetmiterjoin%
\definecolor{currentfill}{rgb}{0.000000,0.000000,0.000000}%
\pgfsetfillcolor{currentfill}%
\pgfsetlinewidth{0.000000pt}%
\definecolor{currentstroke}{rgb}{0.000000,0.000000,0.000000}%
\pgfsetstrokecolor{currentstroke}%
\pgfsetstrokeopacity{0.000000}%
\pgfsetdash{}{0pt}%
\pgfpathmoveto{\pgfqpoint{8.958150in}{0.660000in}}%
\pgfpathlineto{\pgfqpoint{8.964350in}{0.660000in}}%
\pgfpathlineto{\pgfqpoint{8.964350in}{3.191896in}}%
\pgfpathlineto{\pgfqpoint{8.958150in}{3.191896in}}%
\pgfpathlineto{\pgfqpoint{8.958150in}{0.660000in}}%
\pgfpathclose%
\pgfusepath{fill}%
\end{pgfscope}%
\begin{pgfscope}%
\pgfpathrectangle{\pgfqpoint{1.250000in}{0.660000in}}{\pgfqpoint{7.750000in}{4.620000in}}%
\pgfusepath{clip}%
\pgfsetbuttcap%
\pgfsetmiterjoin%
\definecolor{currentfill}{rgb}{0.000000,0.000000,0.000000}%
\pgfsetfillcolor{currentfill}%
\pgfsetlinewidth{0.000000pt}%
\definecolor{currentstroke}{rgb}{0.000000,0.000000,0.000000}%
\pgfsetstrokecolor{currentstroke}%
\pgfsetstrokeopacity{0.000000}%
\pgfsetdash{}{0pt}%
\pgfpathmoveto{\pgfqpoint{8.965900in}{0.660000in}}%
\pgfpathlineto{\pgfqpoint{8.972100in}{0.660000in}}%
\pgfpathlineto{\pgfqpoint{8.972100in}{3.191818in}}%
\pgfpathlineto{\pgfqpoint{8.965900in}{3.191818in}}%
\pgfpathlineto{\pgfqpoint{8.965900in}{0.660000in}}%
\pgfpathclose%
\pgfusepath{fill}%
\end{pgfscope}%
\begin{pgfscope}%
\pgfpathrectangle{\pgfqpoint{1.250000in}{0.660000in}}{\pgfqpoint{7.750000in}{4.620000in}}%
\pgfusepath{clip}%
\pgfsetbuttcap%
\pgfsetmiterjoin%
\definecolor{currentfill}{rgb}{0.000000,0.000000,0.000000}%
\pgfsetfillcolor{currentfill}%
\pgfsetlinewidth{0.000000pt}%
\definecolor{currentstroke}{rgb}{0.000000,0.000000,0.000000}%
\pgfsetstrokecolor{currentstroke}%
\pgfsetstrokeopacity{0.000000}%
\pgfsetdash{}{0pt}%
\pgfpathmoveto{\pgfqpoint{8.973650in}{0.660000in}}%
\pgfpathlineto{\pgfqpoint{8.979850in}{0.660000in}}%
\pgfpathlineto{\pgfqpoint{8.979850in}{3.191741in}}%
\pgfpathlineto{\pgfqpoint{8.973650in}{3.191741in}}%
\pgfpathlineto{\pgfqpoint{8.973650in}{0.660000in}}%
\pgfpathclose%
\pgfusepath{fill}%
\end{pgfscope}%
\begin{pgfscope}%
\pgfpathrectangle{\pgfqpoint{1.250000in}{0.660000in}}{\pgfqpoint{7.750000in}{4.620000in}}%
\pgfusepath{clip}%
\pgfsetbuttcap%
\pgfsetmiterjoin%
\definecolor{currentfill}{rgb}{0.000000,0.000000,0.000000}%
\pgfsetfillcolor{currentfill}%
\pgfsetlinewidth{0.000000pt}%
\definecolor{currentstroke}{rgb}{0.000000,0.000000,0.000000}%
\pgfsetstrokecolor{currentstroke}%
\pgfsetstrokeopacity{0.000000}%
\pgfsetdash{}{0pt}%
\pgfpathmoveto{\pgfqpoint{8.981400in}{0.660000in}}%
\pgfpathlineto{\pgfqpoint{8.987600in}{0.660000in}}%
\pgfpathlineto{\pgfqpoint{8.987600in}{3.191664in}}%
\pgfpathlineto{\pgfqpoint{8.981400in}{3.191664in}}%
\pgfpathlineto{\pgfqpoint{8.981400in}{0.660000in}}%
\pgfpathclose%
\pgfusepath{fill}%
\end{pgfscope}%
\begin{pgfscope}%
\pgfpathrectangle{\pgfqpoint{1.250000in}{0.660000in}}{\pgfqpoint{7.750000in}{4.620000in}}%
\pgfusepath{clip}%
\pgfsetbuttcap%
\pgfsetmiterjoin%
\definecolor{currentfill}{rgb}{0.000000,0.000000,0.000000}%
\pgfsetfillcolor{currentfill}%
\pgfsetlinewidth{0.000000pt}%
\definecolor{currentstroke}{rgb}{0.000000,0.000000,0.000000}%
\pgfsetstrokecolor{currentstroke}%
\pgfsetstrokeopacity{0.000000}%
\pgfsetdash{}{0pt}%
\pgfpathmoveto{\pgfqpoint{8.989150in}{0.660000in}}%
\pgfpathlineto{\pgfqpoint{8.995350in}{0.660000in}}%
\pgfpathlineto{\pgfqpoint{8.995350in}{3.191586in}}%
\pgfpathlineto{\pgfqpoint{8.989150in}{3.191586in}}%
\pgfpathlineto{\pgfqpoint{8.989150in}{0.660000in}}%
\pgfpathclose%
\pgfusepath{fill}%
\end{pgfscope}%
\begin{pgfscope}%
\pgfsetbuttcap%
\pgfsetroundjoin%
\definecolor{currentfill}{rgb}{0.000000,0.000000,0.000000}%
\pgfsetfillcolor{currentfill}%
\pgfsetlinewidth{0.803000pt}%
\definecolor{currentstroke}{rgb}{0.000000,0.000000,0.000000}%
\pgfsetstrokecolor{currentstroke}%
\pgfsetdash{}{0pt}%
\pgfsys@defobject{currentmarker}{\pgfqpoint{0.000000in}{-0.048611in}}{\pgfqpoint{0.000000in}{0.000000in}}{%
\pgfpathmoveto{\pgfqpoint{0.000000in}{0.000000in}}%
\pgfpathlineto{\pgfqpoint{0.000000in}{-0.048611in}}%
\pgfusepath{stroke,fill}%
}%
\begin{pgfscope}%
\pgfsys@transformshift{1.250000in}{0.660000in}%
\pgfsys@useobject{currentmarker}{}%
\end{pgfscope}%
\end{pgfscope}%
\begin{pgfscope}%
\definecolor{textcolor}{rgb}{0.000000,0.000000,0.000000}%
\pgfsetstrokecolor{textcolor}%
\pgfsetfillcolor{textcolor}%
\pgftext[x=1.250000in,y=0.562778in,,top]{\color{textcolor}\rmfamily\fontsize{16.000000}{19.200000}\selectfont \(\displaystyle {0}\)}%
\end{pgfscope}%
\begin{pgfscope}%
\pgfsetbuttcap%
\pgfsetroundjoin%
\definecolor{currentfill}{rgb}{0.000000,0.000000,0.000000}%
\pgfsetfillcolor{currentfill}%
\pgfsetlinewidth{0.803000pt}%
\definecolor{currentstroke}{rgb}{0.000000,0.000000,0.000000}%
\pgfsetstrokecolor{currentstroke}%
\pgfsetdash{}{0pt}%
\pgfsys@defobject{currentmarker}{\pgfqpoint{0.000000in}{-0.048611in}}{\pgfqpoint{0.000000in}{0.000000in}}{%
\pgfpathmoveto{\pgfqpoint{0.000000in}{0.000000in}}%
\pgfpathlineto{\pgfqpoint{0.000000in}{-0.048611in}}%
\pgfusepath{stroke,fill}%
}%
\begin{pgfscope}%
\pgfsys@transformshift{2.800000in}{0.660000in}%
\pgfsys@useobject{currentmarker}{}%
\end{pgfscope}%
\end{pgfscope}%
\begin{pgfscope}%
\definecolor{textcolor}{rgb}{0.000000,0.000000,0.000000}%
\pgfsetstrokecolor{textcolor}%
\pgfsetfillcolor{textcolor}%
\pgftext[x=2.800000in,y=0.562778in,,top]{\color{textcolor}\rmfamily\fontsize{16.000000}{19.200000}\selectfont \(\displaystyle {200}\)}%
\end{pgfscope}%
\begin{pgfscope}%
\pgfsetbuttcap%
\pgfsetroundjoin%
\definecolor{currentfill}{rgb}{0.000000,0.000000,0.000000}%
\pgfsetfillcolor{currentfill}%
\pgfsetlinewidth{0.803000pt}%
\definecolor{currentstroke}{rgb}{0.000000,0.000000,0.000000}%
\pgfsetstrokecolor{currentstroke}%
\pgfsetdash{}{0pt}%
\pgfsys@defobject{currentmarker}{\pgfqpoint{0.000000in}{-0.048611in}}{\pgfqpoint{0.000000in}{0.000000in}}{%
\pgfpathmoveto{\pgfqpoint{0.000000in}{0.000000in}}%
\pgfpathlineto{\pgfqpoint{0.000000in}{-0.048611in}}%
\pgfusepath{stroke,fill}%
}%
\begin{pgfscope}%
\pgfsys@transformshift{4.350000in}{0.660000in}%
\pgfsys@useobject{currentmarker}{}%
\end{pgfscope}%
\end{pgfscope}%
\begin{pgfscope}%
\definecolor{textcolor}{rgb}{0.000000,0.000000,0.000000}%
\pgfsetstrokecolor{textcolor}%
\pgfsetfillcolor{textcolor}%
\pgftext[x=4.350000in,y=0.562778in,,top]{\color{textcolor}\rmfamily\fontsize{16.000000}{19.200000}\selectfont \(\displaystyle {400}\)}%
\end{pgfscope}%
\begin{pgfscope}%
\pgfsetbuttcap%
\pgfsetroundjoin%
\definecolor{currentfill}{rgb}{0.000000,0.000000,0.000000}%
\pgfsetfillcolor{currentfill}%
\pgfsetlinewidth{0.803000pt}%
\definecolor{currentstroke}{rgb}{0.000000,0.000000,0.000000}%
\pgfsetstrokecolor{currentstroke}%
\pgfsetdash{}{0pt}%
\pgfsys@defobject{currentmarker}{\pgfqpoint{0.000000in}{-0.048611in}}{\pgfqpoint{0.000000in}{0.000000in}}{%
\pgfpathmoveto{\pgfqpoint{0.000000in}{0.000000in}}%
\pgfpathlineto{\pgfqpoint{0.000000in}{-0.048611in}}%
\pgfusepath{stroke,fill}%
}%
\begin{pgfscope}%
\pgfsys@transformshift{5.900000in}{0.660000in}%
\pgfsys@useobject{currentmarker}{}%
\end{pgfscope}%
\end{pgfscope}%
\begin{pgfscope}%
\definecolor{textcolor}{rgb}{0.000000,0.000000,0.000000}%
\pgfsetstrokecolor{textcolor}%
\pgfsetfillcolor{textcolor}%
\pgftext[x=5.900000in,y=0.562778in,,top]{\color{textcolor}\rmfamily\fontsize{16.000000}{19.200000}\selectfont \(\displaystyle {600}\)}%
\end{pgfscope}%
\begin{pgfscope}%
\pgfsetbuttcap%
\pgfsetroundjoin%
\definecolor{currentfill}{rgb}{0.000000,0.000000,0.000000}%
\pgfsetfillcolor{currentfill}%
\pgfsetlinewidth{0.803000pt}%
\definecolor{currentstroke}{rgb}{0.000000,0.000000,0.000000}%
\pgfsetstrokecolor{currentstroke}%
\pgfsetdash{}{0pt}%
\pgfsys@defobject{currentmarker}{\pgfqpoint{0.000000in}{-0.048611in}}{\pgfqpoint{0.000000in}{0.000000in}}{%
\pgfpathmoveto{\pgfqpoint{0.000000in}{0.000000in}}%
\pgfpathlineto{\pgfqpoint{0.000000in}{-0.048611in}}%
\pgfusepath{stroke,fill}%
}%
\begin{pgfscope}%
\pgfsys@transformshift{7.450000in}{0.660000in}%
\pgfsys@useobject{currentmarker}{}%
\end{pgfscope}%
\end{pgfscope}%
\begin{pgfscope}%
\definecolor{textcolor}{rgb}{0.000000,0.000000,0.000000}%
\pgfsetstrokecolor{textcolor}%
\pgfsetfillcolor{textcolor}%
\pgftext[x=7.450000in,y=0.562778in,,top]{\color{textcolor}\rmfamily\fontsize{16.000000}{19.200000}\selectfont \(\displaystyle {800}\)}%
\end{pgfscope}%
\begin{pgfscope}%
\pgfsetbuttcap%
\pgfsetroundjoin%
\definecolor{currentfill}{rgb}{0.000000,0.000000,0.000000}%
\pgfsetfillcolor{currentfill}%
\pgfsetlinewidth{0.803000pt}%
\definecolor{currentstroke}{rgb}{0.000000,0.000000,0.000000}%
\pgfsetstrokecolor{currentstroke}%
\pgfsetdash{}{0pt}%
\pgfsys@defobject{currentmarker}{\pgfqpoint{0.000000in}{-0.048611in}}{\pgfqpoint{0.000000in}{0.000000in}}{%
\pgfpathmoveto{\pgfqpoint{0.000000in}{0.000000in}}%
\pgfpathlineto{\pgfqpoint{0.000000in}{-0.048611in}}%
\pgfusepath{stroke,fill}%
}%
\begin{pgfscope}%
\pgfsys@transformshift{9.000000in}{0.660000in}%
\pgfsys@useobject{currentmarker}{}%
\end{pgfscope}%
\end{pgfscope}%
\begin{pgfscope}%
\definecolor{textcolor}{rgb}{0.000000,0.000000,0.000000}%
\pgfsetstrokecolor{textcolor}%
\pgfsetfillcolor{textcolor}%
\pgftext[x=9.000000in,y=0.562778in,,top]{\color{textcolor}\rmfamily\fontsize{16.000000}{19.200000}\selectfont \(\displaystyle {1000}\)}%
\end{pgfscope}%
\begin{pgfscope}%
\pgfsetbuttcap%
\pgfsetroundjoin%
\definecolor{currentfill}{rgb}{0.000000,0.000000,0.000000}%
\pgfsetfillcolor{currentfill}%
\pgfsetlinewidth{0.803000pt}%
\definecolor{currentstroke}{rgb}{0.000000,0.000000,0.000000}%
\pgfsetstrokecolor{currentstroke}%
\pgfsetdash{}{0pt}%
\pgfsys@defobject{currentmarker}{\pgfqpoint{-0.048611in}{0.000000in}}{\pgfqpoint{-0.000000in}{0.000000in}}{%
\pgfpathmoveto{\pgfqpoint{-0.000000in}{0.000000in}}%
\pgfpathlineto{\pgfqpoint{-0.048611in}{0.000000in}}%
\pgfusepath{stroke,fill}%
}%
\begin{pgfscope}%
\pgfsys@transformshift{1.250000in}{0.660000in}%
\pgfsys@useobject{currentmarker}{}%
\end{pgfscope}%
\end{pgfscope}%
\begin{pgfscope}%
\definecolor{textcolor}{rgb}{0.000000,0.000000,0.000000}%
\pgfsetstrokecolor{textcolor}%
\pgfsetfillcolor{textcolor}%
\pgftext[x=0.867364in, y=0.576667in, left, base]{\color{textcolor}\rmfamily\fontsize{16.000000}{19.200000}\selectfont \(\displaystyle {0.0}\)}%
\end{pgfscope}%
\begin{pgfscope}%
\pgfsetbuttcap%
\pgfsetroundjoin%
\definecolor{currentfill}{rgb}{0.000000,0.000000,0.000000}%
\pgfsetfillcolor{currentfill}%
\pgfsetlinewidth{0.803000pt}%
\definecolor{currentstroke}{rgb}{0.000000,0.000000,0.000000}%
\pgfsetstrokecolor{currentstroke}%
\pgfsetdash{}{0pt}%
\pgfsys@defobject{currentmarker}{\pgfqpoint{-0.048611in}{0.000000in}}{\pgfqpoint{-0.000000in}{0.000000in}}{%
\pgfpathmoveto{\pgfqpoint{-0.000000in}{0.000000in}}%
\pgfpathlineto{\pgfqpoint{-0.048611in}{0.000000in}}%
\pgfusepath{stroke,fill}%
}%
\begin{pgfscope}%
\pgfsys@transformshift{1.250000in}{1.584000in}%
\pgfsys@useobject{currentmarker}{}%
\end{pgfscope}%
\end{pgfscope}%
\begin{pgfscope}%
\definecolor{textcolor}{rgb}{0.000000,0.000000,0.000000}%
\pgfsetstrokecolor{textcolor}%
\pgfsetfillcolor{textcolor}%
\pgftext[x=0.867364in, y=1.500667in, left, base]{\color{textcolor}\rmfamily\fontsize{16.000000}{19.200000}\selectfont \(\displaystyle {0.2}\)}%
\end{pgfscope}%
\begin{pgfscope}%
\pgfsetbuttcap%
\pgfsetroundjoin%
\definecolor{currentfill}{rgb}{0.000000,0.000000,0.000000}%
\pgfsetfillcolor{currentfill}%
\pgfsetlinewidth{0.803000pt}%
\definecolor{currentstroke}{rgb}{0.000000,0.000000,0.000000}%
\pgfsetstrokecolor{currentstroke}%
\pgfsetdash{}{0pt}%
\pgfsys@defobject{currentmarker}{\pgfqpoint{-0.048611in}{0.000000in}}{\pgfqpoint{-0.000000in}{0.000000in}}{%
\pgfpathmoveto{\pgfqpoint{-0.000000in}{0.000000in}}%
\pgfpathlineto{\pgfqpoint{-0.048611in}{0.000000in}}%
\pgfusepath{stroke,fill}%
}%
\begin{pgfscope}%
\pgfsys@transformshift{1.250000in}{2.508000in}%
\pgfsys@useobject{currentmarker}{}%
\end{pgfscope}%
\end{pgfscope}%
\begin{pgfscope}%
\definecolor{textcolor}{rgb}{0.000000,0.000000,0.000000}%
\pgfsetstrokecolor{textcolor}%
\pgfsetfillcolor{textcolor}%
\pgftext[x=0.867364in, y=2.424667in, left, base]{\color{textcolor}\rmfamily\fontsize{16.000000}{19.200000}\selectfont \(\displaystyle {0.4}\)}%
\end{pgfscope}%
\begin{pgfscope}%
\pgfsetbuttcap%
\pgfsetroundjoin%
\definecolor{currentfill}{rgb}{0.000000,0.000000,0.000000}%
\pgfsetfillcolor{currentfill}%
\pgfsetlinewidth{0.803000pt}%
\definecolor{currentstroke}{rgb}{0.000000,0.000000,0.000000}%
\pgfsetstrokecolor{currentstroke}%
\pgfsetdash{}{0pt}%
\pgfsys@defobject{currentmarker}{\pgfqpoint{-0.048611in}{0.000000in}}{\pgfqpoint{-0.000000in}{0.000000in}}{%
\pgfpathmoveto{\pgfqpoint{-0.000000in}{0.000000in}}%
\pgfpathlineto{\pgfqpoint{-0.048611in}{0.000000in}}%
\pgfusepath{stroke,fill}%
}%
\begin{pgfscope}%
\pgfsys@transformshift{1.250000in}{3.432000in}%
\pgfsys@useobject{currentmarker}{}%
\end{pgfscope}%
\end{pgfscope}%
\begin{pgfscope}%
\definecolor{textcolor}{rgb}{0.000000,0.000000,0.000000}%
\pgfsetstrokecolor{textcolor}%
\pgfsetfillcolor{textcolor}%
\pgftext[x=0.867364in, y=3.348667in, left, base]{\color{textcolor}\rmfamily\fontsize{16.000000}{19.200000}\selectfont \(\displaystyle {0.6}\)}%
\end{pgfscope}%
\begin{pgfscope}%
\pgfsetbuttcap%
\pgfsetroundjoin%
\definecolor{currentfill}{rgb}{0.000000,0.000000,0.000000}%
\pgfsetfillcolor{currentfill}%
\pgfsetlinewidth{0.803000pt}%
\definecolor{currentstroke}{rgb}{0.000000,0.000000,0.000000}%
\pgfsetstrokecolor{currentstroke}%
\pgfsetdash{}{0pt}%
\pgfsys@defobject{currentmarker}{\pgfqpoint{-0.048611in}{0.000000in}}{\pgfqpoint{-0.000000in}{0.000000in}}{%
\pgfpathmoveto{\pgfqpoint{-0.000000in}{0.000000in}}%
\pgfpathlineto{\pgfqpoint{-0.048611in}{0.000000in}}%
\pgfusepath{stroke,fill}%
}%
\begin{pgfscope}%
\pgfsys@transformshift{1.250000in}{4.356000in}%
\pgfsys@useobject{currentmarker}{}%
\end{pgfscope}%
\end{pgfscope}%
\begin{pgfscope}%
\definecolor{textcolor}{rgb}{0.000000,0.000000,0.000000}%
\pgfsetstrokecolor{textcolor}%
\pgfsetfillcolor{textcolor}%
\pgftext[x=0.867364in, y=4.272667in, left, base]{\color{textcolor}\rmfamily\fontsize{16.000000}{19.200000}\selectfont \(\displaystyle {0.8}\)}%
\end{pgfscope}%
\begin{pgfscope}%
\pgfsetbuttcap%
\pgfsetroundjoin%
\definecolor{currentfill}{rgb}{0.000000,0.000000,0.000000}%
\pgfsetfillcolor{currentfill}%
\pgfsetlinewidth{0.803000pt}%
\definecolor{currentstroke}{rgb}{0.000000,0.000000,0.000000}%
\pgfsetstrokecolor{currentstroke}%
\pgfsetdash{}{0pt}%
\pgfsys@defobject{currentmarker}{\pgfqpoint{-0.048611in}{0.000000in}}{\pgfqpoint{-0.000000in}{0.000000in}}{%
\pgfpathmoveto{\pgfqpoint{-0.000000in}{0.000000in}}%
\pgfpathlineto{\pgfqpoint{-0.048611in}{0.000000in}}%
\pgfusepath{stroke,fill}%
}%
\begin{pgfscope}%
\pgfsys@transformshift{1.250000in}{5.280000in}%
\pgfsys@useobject{currentmarker}{}%
\end{pgfscope}%
\end{pgfscope}%
\begin{pgfscope}%
\definecolor{textcolor}{rgb}{0.000000,0.000000,0.000000}%
\pgfsetstrokecolor{textcolor}%
\pgfsetfillcolor{textcolor}%
\pgftext[x=0.867364in, y=5.196667in, left, base]{\color{textcolor}\rmfamily\fontsize{16.000000}{19.200000}\selectfont \(\displaystyle {1.0}\)}%
\end{pgfscope}%
\begin{pgfscope}%
\definecolor{textcolor}{rgb}{0.000000,0.000000,0.000000}%
\pgfsetstrokecolor{textcolor}%
\pgfsetfillcolor{textcolor}%
\pgftext[x=0.811809in,y=2.970000in,,bottom,rotate=90.000000]{\color{textcolor}\rmfamily\fontsize{16.000000}{19.200000}\selectfont max \(\displaystyle F\)}%
\end{pgfscope}%
\begin{pgfscope}%
\pgfpathrectangle{\pgfqpoint{1.250000in}{0.660000in}}{\pgfqpoint{7.750000in}{4.620000in}}%
\pgfusepath{clip}%
\pgfsetrectcap%
\pgfsetroundjoin%
\pgfsetlinewidth{3.011250pt}%
\definecolor{currentstroke}{rgb}{0.121569,0.466667,0.705882}%
\pgfsetstrokecolor{currentstroke}%
\pgfsetdash{}{0pt}%
\pgfpathmoveto{\pgfqpoint{1.269461in}{5.290000in}}%
\pgfpathlineto{\pgfqpoint{1.273250in}{5.085671in}}%
\pgfpathlineto{\pgfqpoint{1.281000in}{4.836185in}}%
\pgfpathlineto{\pgfqpoint{1.288750in}{4.665367in}}%
\pgfpathlineto{\pgfqpoint{1.296500in}{4.538780in}}%
\pgfpathlineto{\pgfqpoint{1.304250in}{4.439993in}}%
\pgfpathlineto{\pgfqpoint{1.312000in}{4.360035in}}%
\pgfpathlineto{\pgfqpoint{1.319750in}{4.293535in}}%
\pgfpathlineto{\pgfqpoint{1.335250in}{4.188279in}}%
\pgfpathlineto{\pgfqpoint{1.350750in}{4.107764in}}%
\pgfpathlineto{\pgfqpoint{1.366250in}{4.043497in}}%
\pgfpathlineto{\pgfqpoint{1.381750in}{3.990592in}}%
\pgfpathlineto{\pgfqpoint{1.397250in}{3.946004in}}%
\pgfpathlineto{\pgfqpoint{1.412750in}{3.907728in}}%
\pgfpathlineto{\pgfqpoint{1.428250in}{3.874379in}}%
\pgfpathlineto{\pgfqpoint{1.443750in}{3.844965in}}%
\pgfpathlineto{\pgfqpoint{1.467000in}{3.806671in}}%
\pgfpathlineto{\pgfqpoint{1.490250in}{3.773855in}}%
\pgfpathlineto{\pgfqpoint{1.513500in}{3.745301in}}%
\pgfpathlineto{\pgfqpoint{1.536750in}{3.720142in}}%
\pgfpathlineto{\pgfqpoint{1.560000in}{3.697741in}}%
\pgfpathlineto{\pgfqpoint{1.591000in}{3.671351in}}%
\pgfpathlineto{\pgfqpoint{1.622000in}{3.648175in}}%
\pgfpathlineto{\pgfqpoint{1.653000in}{3.627598in}}%
\pgfpathlineto{\pgfqpoint{1.691750in}{3.604835in}}%
\pgfpathlineto{\pgfqpoint{1.730500in}{3.584743in}}%
\pgfpathlineto{\pgfqpoint{1.777000in}{3.563470in}}%
\pgfpathlineto{\pgfqpoint{1.823500in}{3.544701in}}%
\pgfpathlineto{\pgfqpoint{1.877750in}{3.525353in}}%
\pgfpathlineto{\pgfqpoint{1.939750in}{3.505943in}}%
\pgfpathlineto{\pgfqpoint{2.001750in}{3.488850in}}%
\pgfpathlineto{\pgfqpoint{2.071500in}{3.471850in}}%
\pgfpathlineto{\pgfqpoint{2.149000in}{3.455194in}}%
\pgfpathlineto{\pgfqpoint{2.234250in}{3.439058in}}%
\pgfpathlineto{\pgfqpoint{2.335000in}{3.422348in}}%
\pgfpathlineto{\pgfqpoint{2.443500in}{3.406627in}}%
\pgfpathlineto{\pgfqpoint{2.567500in}{3.390944in}}%
\pgfpathlineto{\pgfqpoint{2.707000in}{3.375596in}}%
\pgfpathlineto{\pgfqpoint{2.869750in}{3.360098in}}%
\pgfpathlineto{\pgfqpoint{3.048000in}{3.345441in}}%
\pgfpathlineto{\pgfqpoint{3.249500in}{3.331132in}}%
\pgfpathlineto{\pgfqpoint{3.482000in}{3.316930in}}%
\pgfpathlineto{\pgfqpoint{3.737750in}{3.303506in}}%
\pgfpathlineto{\pgfqpoint{4.047750in}{3.289597in}}%
\pgfpathlineto{\pgfqpoint{4.381000in}{3.276851in}}%
\pgfpathlineto{\pgfqpoint{4.768500in}{3.264200in}}%
\pgfpathlineto{\pgfqpoint{5.218000in}{3.251746in}}%
\pgfpathlineto{\pgfqpoint{5.737250in}{3.239589in}}%
\pgfpathlineto{\pgfqpoint{6.334000in}{3.227808in}}%
\pgfpathlineto{\pgfqpoint{7.008250in}{3.216602in}}%
\pgfpathlineto{\pgfqpoint{7.837500in}{3.205064in}}%
\pgfpathlineto{\pgfqpoint{8.883750in}{3.193083in}}%
\pgfpathlineto{\pgfqpoint{8.992250in}{3.191969in}}%
\pgfpathlineto{\pgfqpoint{8.992250in}{3.191969in}}%
\pgfusepath{stroke}%
\end{pgfscope}%
\begin{pgfscope}%
\pgfpathrectangle{\pgfqpoint{1.250000in}{0.660000in}}{\pgfqpoint{7.750000in}{4.620000in}}%
\pgfusepath{clip}%
\pgfsetrectcap%
\pgfsetroundjoin%
\pgfsetlinewidth{3.011250pt}%
\definecolor{currentstroke}{rgb}{1.000000,0.498039,0.054902}%
\pgfsetstrokecolor{currentstroke}%
\pgfsetdash{}{0pt}%
\pgfpathmoveto{\pgfqpoint{1.265500in}{3.901452in}}%
\pgfpathlineto{\pgfqpoint{1.273250in}{3.784265in}}%
\pgfpathlineto{\pgfqpoint{1.281000in}{3.708450in}}%
\pgfpathlineto{\pgfqpoint{1.288750in}{3.653861in}}%
\pgfpathlineto{\pgfqpoint{1.296500in}{3.611947in}}%
\pgfpathlineto{\pgfqpoint{1.304250in}{3.578347in}}%
\pgfpathlineto{\pgfqpoint{1.312000in}{3.550568in}}%
\pgfpathlineto{\pgfqpoint{1.327500in}{3.506798in}}%
\pgfpathlineto{\pgfqpoint{1.343000in}{3.473398in}}%
\pgfpathlineto{\pgfqpoint{1.358500in}{3.446740in}}%
\pgfpathlineto{\pgfqpoint{1.374000in}{3.424769in}}%
\pgfpathlineto{\pgfqpoint{1.389500in}{3.406221in}}%
\pgfpathlineto{\pgfqpoint{1.405000in}{3.390266in}}%
\pgfpathlineto{\pgfqpoint{1.428250in}{3.369996in}}%
\pgfpathlineto{\pgfqpoint{1.451500in}{3.353022in}}%
\pgfpathlineto{\pgfqpoint{1.474750in}{3.338515in}}%
\pgfpathlineto{\pgfqpoint{1.505750in}{3.322063in}}%
\pgfpathlineto{\pgfqpoint{1.536750in}{3.308116in}}%
\pgfpathlineto{\pgfqpoint{1.575500in}{3.293319in}}%
\pgfpathlineto{\pgfqpoint{1.622000in}{3.278442in}}%
\pgfpathlineto{\pgfqpoint{1.668500in}{3.265899in}}%
\pgfpathlineto{\pgfqpoint{1.722750in}{3.253466in}}%
\pgfpathlineto{\pgfqpoint{1.784750in}{3.241439in}}%
\pgfpathlineto{\pgfqpoint{1.862250in}{3.228830in}}%
\pgfpathlineto{\pgfqpoint{1.947500in}{3.217254in}}%
\pgfpathlineto{\pgfqpoint{2.048250in}{3.205830in}}%
\pgfpathlineto{\pgfqpoint{2.164500in}{3.194866in}}%
\pgfpathlineto{\pgfqpoint{2.304000in}{3.183975in}}%
\pgfpathlineto{\pgfqpoint{2.466750in}{3.173511in}}%
\pgfpathlineto{\pgfqpoint{2.660500in}{3.163290in}}%
\pgfpathlineto{\pgfqpoint{2.893000in}{3.153289in}}%
\pgfpathlineto{\pgfqpoint{3.172000in}{3.143562in}}%
\pgfpathlineto{\pgfqpoint{3.513000in}{3.133996in}}%
\pgfpathlineto{\pgfqpoint{3.923750in}{3.124785in}}%
\pgfpathlineto{\pgfqpoint{4.427500in}{3.115810in}}%
\pgfpathlineto{\pgfqpoint{5.047500in}{3.107102in}}%
\pgfpathlineto{\pgfqpoint{5.822500in}{3.098596in}}%
\pgfpathlineto{\pgfqpoint{6.775750in}{3.090479in}}%
\pgfpathlineto{\pgfqpoint{7.953750in}{3.082735in}}%
\pgfpathlineto{\pgfqpoint{8.992250in}{3.077295in}}%
\pgfpathlineto{\pgfqpoint{8.992250in}{3.077295in}}%
\pgfusepath{stroke}%
\end{pgfscope}%
\begin{pgfscope}%
\pgfpathrectangle{\pgfqpoint{1.250000in}{0.660000in}}{\pgfqpoint{7.750000in}{4.620000in}}%
\pgfusepath{clip}%
\pgfsetbuttcap%
\pgfsetroundjoin%
\pgfsetlinewidth{1.505625pt}%
\definecolor{currentstroke}{rgb}{0.121569,0.466667,0.705882}%
\pgfsetstrokecolor{currentstroke}%
\pgfsetdash{{5.550000pt}{2.400000pt}}{0.000000pt}%
\pgfpathmoveto{\pgfqpoint{1.250000in}{3.740000in}}%
\pgfpathlineto{\pgfqpoint{9.000000in}{3.740000in}}%
\pgfusepath{stroke}%
\end{pgfscope}%
\begin{pgfscope}%
\pgfpathrectangle{\pgfqpoint{1.250000in}{0.660000in}}{\pgfqpoint{7.750000in}{4.620000in}}%
\pgfusepath{clip}%
\pgfsetbuttcap%
\pgfsetroundjoin%
\definecolor{currentfill}{rgb}{0.172549,0.627451,0.172549}%
\pgfsetfillcolor{currentfill}%
\pgfsetlinewidth{1.003750pt}%
\definecolor{currentstroke}{rgb}{0.172549,0.627451,0.172549}%
\pgfsetstrokecolor{currentstroke}%
\pgfsetdash{}{0pt}%
\pgfsys@defobject{currentmarker}{\pgfqpoint{-0.041667in}{-0.041667in}}{\pgfqpoint{0.041667in}{0.041667in}}{%
\pgfpathmoveto{\pgfqpoint{0.000000in}{-0.041667in}}%
\pgfpathcurveto{\pgfqpoint{0.011050in}{-0.041667in}}{\pgfqpoint{0.021649in}{-0.037276in}}{\pgfqpoint{0.029463in}{-0.029463in}}%
\pgfpathcurveto{\pgfqpoint{0.037276in}{-0.021649in}}{\pgfqpoint{0.041667in}{-0.011050in}}{\pgfqpoint{0.041667in}{0.000000in}}%
\pgfpathcurveto{\pgfqpoint{0.041667in}{0.011050in}}{\pgfqpoint{0.037276in}{0.021649in}}{\pgfqpoint{0.029463in}{0.029463in}}%
\pgfpathcurveto{\pgfqpoint{0.021649in}{0.037276in}}{\pgfqpoint{0.011050in}{0.041667in}}{\pgfqpoint{0.000000in}{0.041667in}}%
\pgfpathcurveto{\pgfqpoint{-0.011050in}{0.041667in}}{\pgfqpoint{-0.021649in}{0.037276in}}{\pgfqpoint{-0.029463in}{0.029463in}}%
\pgfpathcurveto{\pgfqpoint{-0.037276in}{0.021649in}}{\pgfqpoint{-0.041667in}{0.011050in}}{\pgfqpoint{-0.041667in}{0.000000in}}%
\pgfpathcurveto{\pgfqpoint{-0.041667in}{-0.011050in}}{\pgfqpoint{-0.037276in}{-0.021649in}}{\pgfqpoint{-0.029463in}{-0.029463in}}%
\pgfpathcurveto{\pgfqpoint{-0.021649in}{-0.037276in}}{\pgfqpoint{-0.011050in}{-0.041667in}}{\pgfqpoint{0.000000in}{-0.041667in}}%
\pgfpathlineto{\pgfqpoint{0.000000in}{-0.041667in}}%
\pgfpathclose%
\pgfusepath{stroke,fill}%
}%
\begin{pgfscope}%
\pgfsys@transformshift{1.505750in}{3.740000in}%
\pgfsys@useobject{currentmarker}{}%
\end{pgfscope}%
\end{pgfscope}%
\begin{pgfscope}%
\pgfsetrectcap%
\pgfsetmiterjoin%
\pgfsetlinewidth{0.803000pt}%
\definecolor{currentstroke}{rgb}{0.000000,0.000000,0.000000}%
\pgfsetstrokecolor{currentstroke}%
\pgfsetdash{}{0pt}%
\pgfpathmoveto{\pgfqpoint{1.250000in}{0.660000in}}%
\pgfpathlineto{\pgfqpoint{1.250000in}{5.280000in}}%
\pgfusepath{stroke}%
\end{pgfscope}%
\begin{pgfscope}%
\pgfsetrectcap%
\pgfsetmiterjoin%
\pgfsetlinewidth{0.803000pt}%
\definecolor{currentstroke}{rgb}{0.000000,0.000000,0.000000}%
\pgfsetstrokecolor{currentstroke}%
\pgfsetdash{}{0pt}%
\pgfpathmoveto{\pgfqpoint{9.000000in}{0.660000in}}%
\pgfpathlineto{\pgfqpoint{9.000000in}{5.280000in}}%
\pgfusepath{stroke}%
\end{pgfscope}%
\begin{pgfscope}%
\pgfsetrectcap%
\pgfsetmiterjoin%
\pgfsetlinewidth{0.803000pt}%
\definecolor{currentstroke}{rgb}{0.000000,0.000000,0.000000}%
\pgfsetstrokecolor{currentstroke}%
\pgfsetdash{}{0pt}%
\pgfpathmoveto{\pgfqpoint{1.250000in}{0.660000in}}%
\pgfpathlineto{\pgfqpoint{9.000000in}{0.660000in}}%
\pgfusepath{stroke}%
\end{pgfscope}%
\begin{pgfscope}%
\pgfsetrectcap%
\pgfsetmiterjoin%
\pgfsetlinewidth{0.803000pt}%
\definecolor{currentstroke}{rgb}{0.000000,0.000000,0.000000}%
\pgfsetstrokecolor{currentstroke}%
\pgfsetdash{}{0pt}%
\pgfpathmoveto{\pgfqpoint{1.250000in}{5.280000in}}%
\pgfpathlineto{\pgfqpoint{9.000000in}{5.280000in}}%
\pgfusepath{stroke}%
\end{pgfscope}%
\begin{pgfscope}%
\definecolor{textcolor}{rgb}{0.000000,0.000000,1.000000}%
\pgfsetstrokecolor{textcolor}%
\pgfsetfillcolor{textcolor}%
\pgftext[x=1.544500in,y=3.832400in,left,base]{\color{textcolor}\rmfamily\fontsize{16.000000}{19.200000}\selectfont \(\displaystyle N = 33\)}%
\end{pgfscope}%
\begin{pgfscope}%
\pgfsetbuttcap%
\pgfsetmiterjoin%
\definecolor{currentfill}{rgb}{1.000000,1.000000,1.000000}%
\pgfsetfillcolor{currentfill}%
\pgfsetfillopacity{0.800000}%
\pgfsetlinewidth{1.003750pt}%
\definecolor{currentstroke}{rgb}{0.800000,0.800000,0.800000}%
\pgfsetstrokecolor{currentstroke}%
\pgfsetstrokeopacity{0.800000}%
\pgfsetdash{}{0pt}%
\pgfpathmoveto{\pgfqpoint{5.879782in}{3.777732in}}%
\pgfpathlineto{\pgfqpoint{8.844444in}{3.777732in}}%
\pgfpathquadraticcurveto{\pgfqpoint{8.888889in}{3.777732in}}{\pgfqpoint{8.888889in}{3.822177in}}%
\pgfpathlineto{\pgfqpoint{8.888889in}{5.124444in}}%
\pgfpathquadraticcurveto{\pgfqpoint{8.888889in}{5.168889in}}{\pgfqpoint{8.844444in}{5.168889in}}%
\pgfpathlineto{\pgfqpoint{5.879782in}{5.168889in}}%
\pgfpathquadraticcurveto{\pgfqpoint{5.835338in}{5.168889in}}{\pgfqpoint{5.835338in}{5.124444in}}%
\pgfpathlineto{\pgfqpoint{5.835338in}{3.822177in}}%
\pgfpathquadraticcurveto{\pgfqpoint{5.835338in}{3.777732in}}{\pgfqpoint{5.879782in}{3.777732in}}%
\pgfpathlineto{\pgfqpoint{5.879782in}{3.777732in}}%
\pgfpathclose%
\pgfusepath{stroke,fill}%
\end{pgfscope}%
\begin{pgfscope}%
\pgfsetrectcap%
\pgfsetroundjoin%
\pgfsetlinewidth{3.011250pt}%
\definecolor{currentstroke}{rgb}{0.121569,0.466667,0.705882}%
\pgfsetstrokecolor{currentstroke}%
\pgfsetdash{}{0pt}%
\pgfpathmoveto{\pgfqpoint{5.924227in}{4.991111in}}%
\pgfpathlineto{\pgfqpoint{6.146449in}{4.991111in}}%
\pgfpathlineto{\pgfqpoint{6.368671in}{4.991111in}}%
\pgfusepath{stroke}%
\end{pgfscope}%
\begin{pgfscope}%
\definecolor{textcolor}{rgb}{0.000000,0.000000,0.000000}%
\pgfsetstrokecolor{textcolor}%
\pgfsetfillcolor{textcolor}%
\pgftext[x=6.546449in,y=4.913333in,left,base]{\color{textcolor}\rmfamily\fontsize{16.000000}{19.200000}\selectfont Saddle}%
\end{pgfscope}%
\begin{pgfscope}%
\pgfsetrectcap%
\pgfsetroundjoin%
\pgfsetlinewidth{3.011250pt}%
\definecolor{currentstroke}{rgb}{1.000000,0.498039,0.054902}%
\pgfsetstrokecolor{currentstroke}%
\pgfsetdash{}{0pt}%
\pgfpathmoveto{\pgfqpoint{5.924227in}{4.666651in}}%
\pgfpathlineto{\pgfqpoint{6.146449in}{4.666651in}}%
\pgfpathlineto{\pgfqpoint{6.368671in}{4.666651in}}%
\pgfusepath{stroke}%
\end{pgfscope}%
\begin{pgfscope}%
\definecolor{textcolor}{rgb}{0.000000,0.000000,0.000000}%
\pgfsetstrokecolor{textcolor}%
\pgfsetfillcolor{textcolor}%
\pgftext[x=6.546449in,y=4.588874in,left,base]{\color{textcolor}\rmfamily\fontsize{16.000000}{19.200000}\selectfont Bessel}%
\end{pgfscope}%
\begin{pgfscope}%
\pgfsetbuttcap%
\pgfsetroundjoin%
\pgfsetlinewidth{1.505625pt}%
\definecolor{currentstroke}{rgb}{0.121569,0.466667,0.705882}%
\pgfsetstrokecolor{currentstroke}%
\pgfsetdash{{5.550000pt}{2.400000pt}}{0.000000pt}%
\pgfpathmoveto{\pgfqpoint{5.924227in}{4.328858in}}%
\pgfpathlineto{\pgfqpoint{6.146449in}{4.328858in}}%
\pgfpathlineto{\pgfqpoint{6.368671in}{4.328858in}}%
\pgfusepath{stroke}%
\end{pgfscope}%
\begin{pgfscope}%
\definecolor{textcolor}{rgb}{0.000000,0.000000,0.000000}%
\pgfsetstrokecolor{textcolor}%
\pgfsetfillcolor{textcolor}%
\pgftext[x=6.546449in,y=4.251081in,left,base]{\color{textcolor}\rmfamily\fontsize{16.000000}{19.200000}\selectfont Classical limit \(\displaystyle F=2/3\)}%
\end{pgfscope}%
\begin{pgfscope}%
\pgfsetbuttcap%
\pgfsetmiterjoin%
\definecolor{currentfill}{rgb}{0.000000,0.000000,0.000000}%
\pgfsetfillcolor{currentfill}%
\pgfsetlinewidth{0.000000pt}%
\definecolor{currentstroke}{rgb}{0.000000,0.000000,0.000000}%
\pgfsetstrokecolor{currentstroke}%
\pgfsetstrokeopacity{0.000000}%
\pgfsetdash{}{0pt}%
\pgfpathmoveto{\pgfqpoint{5.924227in}{3.913303in}}%
\pgfpathlineto{\pgfqpoint{6.368671in}{3.913303in}}%
\pgfpathlineto{\pgfqpoint{6.368671in}{4.068859in}}%
\pgfpathlineto{\pgfqpoint{5.924227in}{4.068859in}}%
\pgfpathlineto{\pgfqpoint{5.924227in}{3.913303in}}%
\pgfpathclose%
\pgfusepath{fill}%
\end{pgfscope}%
\begin{pgfscope}%
\definecolor{textcolor}{rgb}{0.000000,0.000000,0.000000}%
\pgfsetstrokecolor{textcolor}%
\pgfsetfillcolor{textcolor}%
\pgftext[x=6.546449in,y=3.913303in,left,base]{\color{textcolor}\rmfamily\fontsize{16.000000}{19.200000}\selectfont Numeric}%
\end{pgfscope}%
\end{pgfpicture}%
\makeatother%
\endgroup%
}
    \resizebox{\textwidth}{!}{%% Creator: Matplotlib, PGF backend
%%
%% To include the figure in your LaTeX document, write
%%   \input{<filename>.pgf}
%%
%% Make sure the required packages are loaded in your preamble
%%   \usepackage{pgf}
%%
%% Also ensure that all the required font packages are loaded; for instance,
%% the lmodern package is sometimes necessary when using math font.
%%   \usepackage{lmodern}
%%
%% Figures using additional raster images can only be included by \input if
%% they are in the same directory as the main LaTeX file. For loading figures
%% from other directories you can use the `import` package
%%   \usepackage{import}
%%
%% and then include the figures with
%%   \import{<path to file>}{<filename>.pgf}
%%
%% Matplotlib used the following preamble
%%   
%%   \makeatletter\@ifpackageloaded{underscore}{}{\usepackage[strings]{underscore}}\makeatother
%%
\begingroup%
\makeatletter%
\begin{pgfpicture}%
\pgfpathrectangle{\pgfpointorigin}{\pgfqpoint{10.000000in}{6.000000in}}%
\pgfusepath{use as bounding box, clip}%
\begin{pgfscope}%
\pgfsetbuttcap%
\pgfsetmiterjoin%
\definecolor{currentfill}{rgb}{1.000000,1.000000,1.000000}%
\pgfsetfillcolor{currentfill}%
\pgfsetlinewidth{0.000000pt}%
\definecolor{currentstroke}{rgb}{1.000000,1.000000,1.000000}%
\pgfsetstrokecolor{currentstroke}%
\pgfsetdash{}{0pt}%
\pgfpathmoveto{\pgfqpoint{0.000000in}{0.000000in}}%
\pgfpathlineto{\pgfqpoint{10.000000in}{0.000000in}}%
\pgfpathlineto{\pgfqpoint{10.000000in}{6.000000in}}%
\pgfpathlineto{\pgfqpoint{0.000000in}{6.000000in}}%
\pgfpathlineto{\pgfqpoint{0.000000in}{0.000000in}}%
\pgfpathclose%
\pgfusepath{fill}%
\end{pgfscope}%
\begin{pgfscope}%
\pgfsetbuttcap%
\pgfsetmiterjoin%
\definecolor{currentfill}{rgb}{1.000000,1.000000,1.000000}%
\pgfsetfillcolor{currentfill}%
\pgfsetlinewidth{0.000000pt}%
\definecolor{currentstroke}{rgb}{0.000000,0.000000,0.000000}%
\pgfsetstrokecolor{currentstroke}%
\pgfsetstrokeopacity{0.000000}%
\pgfsetdash{}{0pt}%
\pgfpathmoveto{\pgfqpoint{1.250000in}{0.660000in}}%
\pgfpathlineto{\pgfqpoint{9.000000in}{0.660000in}}%
\pgfpathlineto{\pgfqpoint{9.000000in}{5.280000in}}%
\pgfpathlineto{\pgfqpoint{1.250000in}{5.280000in}}%
\pgfpathlineto{\pgfqpoint{1.250000in}{0.660000in}}%
\pgfpathclose%
\pgfusepath{fill}%
\end{pgfscope}%
\begin{pgfscope}%
\pgfpathrectangle{\pgfqpoint{1.250000in}{0.660000in}}{\pgfqpoint{7.750000in}{4.620000in}}%
\pgfusepath{clip}%
\pgfsetbuttcap%
\pgfsetmiterjoin%
\definecolor{currentfill}{rgb}{0.000000,0.000000,0.000000}%
\pgfsetfillcolor{currentfill}%
\pgfsetlinewidth{0.000000pt}%
\definecolor{currentstroke}{rgb}{0.000000,0.000000,0.000000}%
\pgfsetstrokecolor{currentstroke}%
\pgfsetstrokeopacity{0.000000}%
\pgfsetdash{}{0pt}%
\pgfpathmoveto{\pgfqpoint{1.262400in}{0.660000in}}%
\pgfpathlineto{\pgfqpoint{1.268600in}{0.660000in}}%
\pgfpathlineto{\pgfqpoint{1.268600in}{5.280000in}}%
\pgfpathlineto{\pgfqpoint{1.262400in}{5.280000in}}%
\pgfpathlineto{\pgfqpoint{1.262400in}{0.660000in}}%
\pgfpathclose%
\pgfusepath{fill}%
\end{pgfscope}%
\begin{pgfscope}%
\pgfpathrectangle{\pgfqpoint{1.250000in}{0.660000in}}{\pgfqpoint{7.750000in}{4.620000in}}%
\pgfusepath{clip}%
\pgfsetbuttcap%
\pgfsetmiterjoin%
\definecolor{currentfill}{rgb}{0.000000,0.000000,0.000000}%
\pgfsetfillcolor{currentfill}%
\pgfsetlinewidth{0.000000pt}%
\definecolor{currentstroke}{rgb}{0.000000,0.000000,0.000000}%
\pgfsetstrokecolor{currentstroke}%
\pgfsetstrokeopacity{0.000000}%
\pgfsetdash{}{0pt}%
\pgfpathmoveto{\pgfqpoint{1.270150in}{0.660000in}}%
\pgfpathlineto{\pgfqpoint{1.276350in}{0.660000in}}%
\pgfpathlineto{\pgfqpoint{1.276350in}{4.881179in}}%
\pgfpathlineto{\pgfqpoint{1.270150in}{4.881179in}}%
\pgfpathlineto{\pgfqpoint{1.270150in}{0.660000in}}%
\pgfpathclose%
\pgfusepath{fill}%
\end{pgfscope}%
\begin{pgfscope}%
\pgfpathrectangle{\pgfqpoint{1.250000in}{0.660000in}}{\pgfqpoint{7.750000in}{4.620000in}}%
\pgfusepath{clip}%
\pgfsetbuttcap%
\pgfsetmiterjoin%
\definecolor{currentfill}{rgb}{0.000000,0.000000,0.000000}%
\pgfsetfillcolor{currentfill}%
\pgfsetlinewidth{0.000000pt}%
\definecolor{currentstroke}{rgb}{0.000000,0.000000,0.000000}%
\pgfsetstrokecolor{currentstroke}%
\pgfsetstrokeopacity{0.000000}%
\pgfsetdash{}{0pt}%
\pgfpathmoveto{\pgfqpoint{1.277900in}{0.660000in}}%
\pgfpathlineto{\pgfqpoint{1.284100in}{0.660000in}}%
\pgfpathlineto{\pgfqpoint{1.284100in}{5.279955in}}%
\pgfpathlineto{\pgfqpoint{1.277900in}{5.279955in}}%
\pgfpathlineto{\pgfqpoint{1.277900in}{0.660000in}}%
\pgfpathclose%
\pgfusepath{fill}%
\end{pgfscope}%
\begin{pgfscope}%
\pgfpathrectangle{\pgfqpoint{1.250000in}{0.660000in}}{\pgfqpoint{7.750000in}{4.620000in}}%
\pgfusepath{clip}%
\pgfsetbuttcap%
\pgfsetmiterjoin%
\definecolor{currentfill}{rgb}{0.000000,0.000000,0.000000}%
\pgfsetfillcolor{currentfill}%
\pgfsetlinewidth{0.000000pt}%
\definecolor{currentstroke}{rgb}{0.000000,0.000000,0.000000}%
\pgfsetstrokecolor{currentstroke}%
\pgfsetstrokeopacity{0.000000}%
\pgfsetdash{}{0pt}%
\pgfpathmoveto{\pgfqpoint{1.285650in}{0.660000in}}%
\pgfpathlineto{\pgfqpoint{1.291850in}{0.660000in}}%
\pgfpathlineto{\pgfqpoint{1.291850in}{5.130941in}}%
\pgfpathlineto{\pgfqpoint{1.285650in}{5.130941in}}%
\pgfpathlineto{\pgfqpoint{1.285650in}{0.660000in}}%
\pgfpathclose%
\pgfusepath{fill}%
\end{pgfscope}%
\begin{pgfscope}%
\pgfpathrectangle{\pgfqpoint{1.250000in}{0.660000in}}{\pgfqpoint{7.750000in}{4.620000in}}%
\pgfusepath{clip}%
\pgfsetbuttcap%
\pgfsetmiterjoin%
\definecolor{currentfill}{rgb}{0.000000,0.000000,0.000000}%
\pgfsetfillcolor{currentfill}%
\pgfsetlinewidth{0.000000pt}%
\definecolor{currentstroke}{rgb}{0.000000,0.000000,0.000000}%
\pgfsetstrokecolor{currentstroke}%
\pgfsetstrokeopacity{0.000000}%
\pgfsetdash{}{0pt}%
\pgfpathmoveto{\pgfqpoint{1.293400in}{0.660000in}}%
\pgfpathlineto{\pgfqpoint{1.299600in}{0.660000in}}%
\pgfpathlineto{\pgfqpoint{1.299600in}{4.558017in}}%
\pgfpathlineto{\pgfqpoint{1.293400in}{4.558017in}}%
\pgfpathlineto{\pgfqpoint{1.293400in}{0.660000in}}%
\pgfpathclose%
\pgfusepath{fill}%
\end{pgfscope}%
\begin{pgfscope}%
\pgfpathrectangle{\pgfqpoint{1.250000in}{0.660000in}}{\pgfqpoint{7.750000in}{4.620000in}}%
\pgfusepath{clip}%
\pgfsetbuttcap%
\pgfsetmiterjoin%
\definecolor{currentfill}{rgb}{0.000000,0.000000,0.000000}%
\pgfsetfillcolor{currentfill}%
\pgfsetlinewidth{0.000000pt}%
\definecolor{currentstroke}{rgb}{0.000000,0.000000,0.000000}%
\pgfsetstrokecolor{currentstroke}%
\pgfsetstrokeopacity{0.000000}%
\pgfsetdash{}{0pt}%
\pgfpathmoveto{\pgfqpoint{1.301150in}{0.660000in}}%
\pgfpathlineto{\pgfqpoint{1.307350in}{0.660000in}}%
\pgfpathlineto{\pgfqpoint{1.307350in}{5.201013in}}%
\pgfpathlineto{\pgfqpoint{1.301150in}{5.201013in}}%
\pgfpathlineto{\pgfqpoint{1.301150in}{0.660000in}}%
\pgfpathclose%
\pgfusepath{fill}%
\end{pgfscope}%
\begin{pgfscope}%
\pgfpathrectangle{\pgfqpoint{1.250000in}{0.660000in}}{\pgfqpoint{7.750000in}{4.620000in}}%
\pgfusepath{clip}%
\pgfsetbuttcap%
\pgfsetmiterjoin%
\definecolor{currentfill}{rgb}{0.000000,0.000000,0.000000}%
\pgfsetfillcolor{currentfill}%
\pgfsetlinewidth{0.000000pt}%
\definecolor{currentstroke}{rgb}{0.000000,0.000000,0.000000}%
\pgfsetstrokecolor{currentstroke}%
\pgfsetstrokeopacity{0.000000}%
\pgfsetdash{}{0pt}%
\pgfpathmoveto{\pgfqpoint{1.308900in}{0.660000in}}%
\pgfpathlineto{\pgfqpoint{1.315100in}{0.660000in}}%
\pgfpathlineto{\pgfqpoint{1.315100in}{5.252567in}}%
\pgfpathlineto{\pgfqpoint{1.308900in}{5.252567in}}%
\pgfpathlineto{\pgfqpoint{1.308900in}{0.660000in}}%
\pgfpathclose%
\pgfusepath{fill}%
\end{pgfscope}%
\begin{pgfscope}%
\pgfpathrectangle{\pgfqpoint{1.250000in}{0.660000in}}{\pgfqpoint{7.750000in}{4.620000in}}%
\pgfusepath{clip}%
\pgfsetbuttcap%
\pgfsetmiterjoin%
\definecolor{currentfill}{rgb}{0.000000,0.000000,0.000000}%
\pgfsetfillcolor{currentfill}%
\pgfsetlinewidth{0.000000pt}%
\definecolor{currentstroke}{rgb}{0.000000,0.000000,0.000000}%
\pgfsetstrokecolor{currentstroke}%
\pgfsetstrokeopacity{0.000000}%
\pgfsetdash{}{0pt}%
\pgfpathmoveto{\pgfqpoint{1.316650in}{0.660000in}}%
\pgfpathlineto{\pgfqpoint{1.322850in}{0.660000in}}%
\pgfpathlineto{\pgfqpoint{1.322850in}{4.872589in}}%
\pgfpathlineto{\pgfqpoint{1.316650in}{4.872589in}}%
\pgfpathlineto{\pgfqpoint{1.316650in}{0.660000in}}%
\pgfpathclose%
\pgfusepath{fill}%
\end{pgfscope}%
\begin{pgfscope}%
\pgfpathrectangle{\pgfqpoint{1.250000in}{0.660000in}}{\pgfqpoint{7.750000in}{4.620000in}}%
\pgfusepath{clip}%
\pgfsetbuttcap%
\pgfsetmiterjoin%
\definecolor{currentfill}{rgb}{0.000000,0.000000,0.000000}%
\pgfsetfillcolor{currentfill}%
\pgfsetlinewidth{0.000000pt}%
\definecolor{currentstroke}{rgb}{0.000000,0.000000,0.000000}%
\pgfsetstrokecolor{currentstroke}%
\pgfsetstrokeopacity{0.000000}%
\pgfsetdash{}{0pt}%
\pgfpathmoveto{\pgfqpoint{1.324400in}{0.660000in}}%
\pgfpathlineto{\pgfqpoint{1.330600in}{0.660000in}}%
\pgfpathlineto{\pgfqpoint{1.330600in}{4.954392in}}%
\pgfpathlineto{\pgfqpoint{1.324400in}{4.954392in}}%
\pgfpathlineto{\pgfqpoint{1.324400in}{0.660000in}}%
\pgfpathclose%
\pgfusepath{fill}%
\end{pgfscope}%
\begin{pgfscope}%
\pgfpathrectangle{\pgfqpoint{1.250000in}{0.660000in}}{\pgfqpoint{7.750000in}{4.620000in}}%
\pgfusepath{clip}%
\pgfsetbuttcap%
\pgfsetmiterjoin%
\definecolor{currentfill}{rgb}{0.000000,0.000000,0.000000}%
\pgfsetfillcolor{currentfill}%
\pgfsetlinewidth{0.000000pt}%
\definecolor{currentstroke}{rgb}{0.000000,0.000000,0.000000}%
\pgfsetstrokecolor{currentstroke}%
\pgfsetstrokeopacity{0.000000}%
\pgfsetdash{}{0pt}%
\pgfpathmoveto{\pgfqpoint{1.332150in}{0.660000in}}%
\pgfpathlineto{\pgfqpoint{1.338350in}{0.660000in}}%
\pgfpathlineto{\pgfqpoint{1.338350in}{5.162530in}}%
\pgfpathlineto{\pgfqpoint{1.332150in}{5.162530in}}%
\pgfpathlineto{\pgfqpoint{1.332150in}{0.660000in}}%
\pgfpathclose%
\pgfusepath{fill}%
\end{pgfscope}%
\begin{pgfscope}%
\pgfpathrectangle{\pgfqpoint{1.250000in}{0.660000in}}{\pgfqpoint{7.750000in}{4.620000in}}%
\pgfusepath{clip}%
\pgfsetbuttcap%
\pgfsetmiterjoin%
\definecolor{currentfill}{rgb}{0.000000,0.000000,0.000000}%
\pgfsetfillcolor{currentfill}%
\pgfsetlinewidth{0.000000pt}%
\definecolor{currentstroke}{rgb}{0.000000,0.000000,0.000000}%
\pgfsetstrokecolor{currentstroke}%
\pgfsetstrokeopacity{0.000000}%
\pgfsetdash{}{0pt}%
\pgfpathmoveto{\pgfqpoint{1.339900in}{0.660000in}}%
\pgfpathlineto{\pgfqpoint{1.346100in}{0.660000in}}%
\pgfpathlineto{\pgfqpoint{1.346100in}{4.542899in}}%
\pgfpathlineto{\pgfqpoint{1.339900in}{4.542899in}}%
\pgfpathlineto{\pgfqpoint{1.339900in}{0.660000in}}%
\pgfpathclose%
\pgfusepath{fill}%
\end{pgfscope}%
\begin{pgfscope}%
\pgfpathrectangle{\pgfqpoint{1.250000in}{0.660000in}}{\pgfqpoint{7.750000in}{4.620000in}}%
\pgfusepath{clip}%
\pgfsetbuttcap%
\pgfsetmiterjoin%
\definecolor{currentfill}{rgb}{0.000000,0.000000,0.000000}%
\pgfsetfillcolor{currentfill}%
\pgfsetlinewidth{0.000000pt}%
\definecolor{currentstroke}{rgb}{0.000000,0.000000,0.000000}%
\pgfsetstrokecolor{currentstroke}%
\pgfsetstrokeopacity{0.000000}%
\pgfsetdash{}{0pt}%
\pgfpathmoveto{\pgfqpoint{1.347650in}{0.660000in}}%
\pgfpathlineto{\pgfqpoint{1.353850in}{0.660000in}}%
\pgfpathlineto{\pgfqpoint{1.353850in}{5.036436in}}%
\pgfpathlineto{\pgfqpoint{1.347650in}{5.036436in}}%
\pgfpathlineto{\pgfqpoint{1.347650in}{0.660000in}}%
\pgfpathclose%
\pgfusepath{fill}%
\end{pgfscope}%
\begin{pgfscope}%
\pgfpathrectangle{\pgfqpoint{1.250000in}{0.660000in}}{\pgfqpoint{7.750000in}{4.620000in}}%
\pgfusepath{clip}%
\pgfsetbuttcap%
\pgfsetmiterjoin%
\definecolor{currentfill}{rgb}{0.000000,0.000000,0.000000}%
\pgfsetfillcolor{currentfill}%
\pgfsetlinewidth{0.000000pt}%
\definecolor{currentstroke}{rgb}{0.000000,0.000000,0.000000}%
\pgfsetstrokecolor{currentstroke}%
\pgfsetstrokeopacity{0.000000}%
\pgfsetdash{}{0pt}%
\pgfpathmoveto{\pgfqpoint{1.355400in}{0.660000in}}%
\pgfpathlineto{\pgfqpoint{1.361600in}{0.660000in}}%
\pgfpathlineto{\pgfqpoint{1.361600in}{4.993626in}}%
\pgfpathlineto{\pgfqpoint{1.355400in}{4.993626in}}%
\pgfpathlineto{\pgfqpoint{1.355400in}{0.660000in}}%
\pgfpathclose%
\pgfusepath{fill}%
\end{pgfscope}%
\begin{pgfscope}%
\pgfpathrectangle{\pgfqpoint{1.250000in}{0.660000in}}{\pgfqpoint{7.750000in}{4.620000in}}%
\pgfusepath{clip}%
\pgfsetbuttcap%
\pgfsetmiterjoin%
\definecolor{currentfill}{rgb}{0.000000,0.000000,0.000000}%
\pgfsetfillcolor{currentfill}%
\pgfsetlinewidth{0.000000pt}%
\definecolor{currentstroke}{rgb}{0.000000,0.000000,0.000000}%
\pgfsetstrokecolor{currentstroke}%
\pgfsetstrokeopacity{0.000000}%
\pgfsetdash{}{0pt}%
\pgfpathmoveto{\pgfqpoint{1.363150in}{0.660000in}}%
\pgfpathlineto{\pgfqpoint{1.369350in}{0.660000in}}%
\pgfpathlineto{\pgfqpoint{1.369350in}{4.427277in}}%
\pgfpathlineto{\pgfqpoint{1.363150in}{4.427277in}}%
\pgfpathlineto{\pgfqpoint{1.363150in}{0.660000in}}%
\pgfpathclose%
\pgfusepath{fill}%
\end{pgfscope}%
\begin{pgfscope}%
\pgfpathrectangle{\pgfqpoint{1.250000in}{0.660000in}}{\pgfqpoint{7.750000in}{4.620000in}}%
\pgfusepath{clip}%
\pgfsetbuttcap%
\pgfsetmiterjoin%
\definecolor{currentfill}{rgb}{0.000000,0.000000,0.000000}%
\pgfsetfillcolor{currentfill}%
\pgfsetlinewidth{0.000000pt}%
\definecolor{currentstroke}{rgb}{0.000000,0.000000,0.000000}%
\pgfsetstrokecolor{currentstroke}%
\pgfsetstrokeopacity{0.000000}%
\pgfsetdash{}{0pt}%
\pgfpathmoveto{\pgfqpoint{1.370900in}{0.660000in}}%
\pgfpathlineto{\pgfqpoint{1.377100in}{0.660000in}}%
\pgfpathlineto{\pgfqpoint{1.377100in}{4.932259in}}%
\pgfpathlineto{\pgfqpoint{1.370900in}{4.932259in}}%
\pgfpathlineto{\pgfqpoint{1.370900in}{0.660000in}}%
\pgfpathclose%
\pgfusepath{fill}%
\end{pgfscope}%
\begin{pgfscope}%
\pgfpathrectangle{\pgfqpoint{1.250000in}{0.660000in}}{\pgfqpoint{7.750000in}{4.620000in}}%
\pgfusepath{clip}%
\pgfsetbuttcap%
\pgfsetmiterjoin%
\definecolor{currentfill}{rgb}{0.000000,0.000000,0.000000}%
\pgfsetfillcolor{currentfill}%
\pgfsetlinewidth{0.000000pt}%
\definecolor{currentstroke}{rgb}{0.000000,0.000000,0.000000}%
\pgfsetstrokecolor{currentstroke}%
\pgfsetstrokeopacity{0.000000}%
\pgfsetdash{}{0pt}%
\pgfpathmoveto{\pgfqpoint{1.378650in}{0.660000in}}%
\pgfpathlineto{\pgfqpoint{1.384850in}{0.660000in}}%
\pgfpathlineto{\pgfqpoint{1.384850in}{4.888567in}}%
\pgfpathlineto{\pgfqpoint{1.378650in}{4.888567in}}%
\pgfpathlineto{\pgfqpoint{1.378650in}{0.660000in}}%
\pgfpathclose%
\pgfusepath{fill}%
\end{pgfscope}%
\begin{pgfscope}%
\pgfpathrectangle{\pgfqpoint{1.250000in}{0.660000in}}{\pgfqpoint{7.750000in}{4.620000in}}%
\pgfusepath{clip}%
\pgfsetbuttcap%
\pgfsetmiterjoin%
\definecolor{currentfill}{rgb}{0.000000,0.000000,0.000000}%
\pgfsetfillcolor{currentfill}%
\pgfsetlinewidth{0.000000pt}%
\definecolor{currentstroke}{rgb}{0.000000,0.000000,0.000000}%
\pgfsetstrokecolor{currentstroke}%
\pgfsetstrokeopacity{0.000000}%
\pgfsetdash{}{0pt}%
\pgfpathmoveto{\pgfqpoint{1.386400in}{0.660000in}}%
\pgfpathlineto{\pgfqpoint{1.392600in}{0.660000in}}%
\pgfpathlineto{\pgfqpoint{1.392600in}{4.423985in}}%
\pgfpathlineto{\pgfqpoint{1.386400in}{4.423985in}}%
\pgfpathlineto{\pgfqpoint{1.386400in}{0.660000in}}%
\pgfpathclose%
\pgfusepath{fill}%
\end{pgfscope}%
\begin{pgfscope}%
\pgfpathrectangle{\pgfqpoint{1.250000in}{0.660000in}}{\pgfqpoint{7.750000in}{4.620000in}}%
\pgfusepath{clip}%
\pgfsetbuttcap%
\pgfsetmiterjoin%
\definecolor{currentfill}{rgb}{0.000000,0.000000,0.000000}%
\pgfsetfillcolor{currentfill}%
\pgfsetlinewidth{0.000000pt}%
\definecolor{currentstroke}{rgb}{0.000000,0.000000,0.000000}%
\pgfsetstrokecolor{currentstroke}%
\pgfsetstrokeopacity{0.000000}%
\pgfsetdash{}{0pt}%
\pgfpathmoveto{\pgfqpoint{1.394150in}{0.660000in}}%
\pgfpathlineto{\pgfqpoint{1.400350in}{0.660000in}}%
\pgfpathlineto{\pgfqpoint{1.400350in}{4.689778in}}%
\pgfpathlineto{\pgfqpoint{1.394150in}{4.689778in}}%
\pgfpathlineto{\pgfqpoint{1.394150in}{0.660000in}}%
\pgfpathclose%
\pgfusepath{fill}%
\end{pgfscope}%
\begin{pgfscope}%
\pgfpathrectangle{\pgfqpoint{1.250000in}{0.660000in}}{\pgfqpoint{7.750000in}{4.620000in}}%
\pgfusepath{clip}%
\pgfsetbuttcap%
\pgfsetmiterjoin%
\definecolor{currentfill}{rgb}{0.000000,0.000000,0.000000}%
\pgfsetfillcolor{currentfill}%
\pgfsetlinewidth{0.000000pt}%
\definecolor{currentstroke}{rgb}{0.000000,0.000000,0.000000}%
\pgfsetstrokecolor{currentstroke}%
\pgfsetstrokeopacity{0.000000}%
\pgfsetdash{}{0pt}%
\pgfpathmoveto{\pgfqpoint{1.401900in}{0.660000in}}%
\pgfpathlineto{\pgfqpoint{1.408100in}{0.660000in}}%
\pgfpathlineto{\pgfqpoint{1.408100in}{4.637415in}}%
\pgfpathlineto{\pgfqpoint{1.401900in}{4.637415in}}%
\pgfpathlineto{\pgfqpoint{1.401900in}{0.660000in}}%
\pgfpathclose%
\pgfusepath{fill}%
\end{pgfscope}%
\begin{pgfscope}%
\pgfpathrectangle{\pgfqpoint{1.250000in}{0.660000in}}{\pgfqpoint{7.750000in}{4.620000in}}%
\pgfusepath{clip}%
\pgfsetbuttcap%
\pgfsetmiterjoin%
\definecolor{currentfill}{rgb}{0.000000,0.000000,0.000000}%
\pgfsetfillcolor{currentfill}%
\pgfsetlinewidth{0.000000pt}%
\definecolor{currentstroke}{rgb}{0.000000,0.000000,0.000000}%
\pgfsetstrokecolor{currentstroke}%
\pgfsetstrokeopacity{0.000000}%
\pgfsetdash{}{0pt}%
\pgfpathmoveto{\pgfqpoint{1.409650in}{0.660000in}}%
\pgfpathlineto{\pgfqpoint{1.415850in}{0.660000in}}%
\pgfpathlineto{\pgfqpoint{1.415850in}{4.520643in}}%
\pgfpathlineto{\pgfqpoint{1.409650in}{4.520643in}}%
\pgfpathlineto{\pgfqpoint{1.409650in}{0.660000in}}%
\pgfpathclose%
\pgfusepath{fill}%
\end{pgfscope}%
\begin{pgfscope}%
\pgfpathrectangle{\pgfqpoint{1.250000in}{0.660000in}}{\pgfqpoint{7.750000in}{4.620000in}}%
\pgfusepath{clip}%
\pgfsetbuttcap%
\pgfsetmiterjoin%
\definecolor{currentfill}{rgb}{0.000000,0.000000,0.000000}%
\pgfsetfillcolor{currentfill}%
\pgfsetlinewidth{0.000000pt}%
\definecolor{currentstroke}{rgb}{0.000000,0.000000,0.000000}%
\pgfsetstrokecolor{currentstroke}%
\pgfsetstrokeopacity{0.000000}%
\pgfsetdash{}{0pt}%
\pgfpathmoveto{\pgfqpoint{1.417400in}{0.660000in}}%
\pgfpathlineto{\pgfqpoint{1.423600in}{0.660000in}}%
\pgfpathlineto{\pgfqpoint{1.423600in}{4.850484in}}%
\pgfpathlineto{\pgfqpoint{1.417400in}{4.850484in}}%
\pgfpathlineto{\pgfqpoint{1.417400in}{0.660000in}}%
\pgfpathclose%
\pgfusepath{fill}%
\end{pgfscope}%
\begin{pgfscope}%
\pgfpathrectangle{\pgfqpoint{1.250000in}{0.660000in}}{\pgfqpoint{7.750000in}{4.620000in}}%
\pgfusepath{clip}%
\pgfsetbuttcap%
\pgfsetmiterjoin%
\definecolor{currentfill}{rgb}{0.000000,0.000000,0.000000}%
\pgfsetfillcolor{currentfill}%
\pgfsetlinewidth{0.000000pt}%
\definecolor{currentstroke}{rgb}{0.000000,0.000000,0.000000}%
\pgfsetstrokecolor{currentstroke}%
\pgfsetstrokeopacity{0.000000}%
\pgfsetdash{}{0pt}%
\pgfpathmoveto{\pgfqpoint{1.425150in}{0.660000in}}%
\pgfpathlineto{\pgfqpoint{1.431350in}{0.660000in}}%
\pgfpathlineto{\pgfqpoint{1.431350in}{4.564290in}}%
\pgfpathlineto{\pgfqpoint{1.425150in}{4.564290in}}%
\pgfpathlineto{\pgfqpoint{1.425150in}{0.660000in}}%
\pgfpathclose%
\pgfusepath{fill}%
\end{pgfscope}%
\begin{pgfscope}%
\pgfpathrectangle{\pgfqpoint{1.250000in}{0.660000in}}{\pgfqpoint{7.750000in}{4.620000in}}%
\pgfusepath{clip}%
\pgfsetbuttcap%
\pgfsetmiterjoin%
\definecolor{currentfill}{rgb}{0.000000,0.000000,0.000000}%
\pgfsetfillcolor{currentfill}%
\pgfsetlinewidth{0.000000pt}%
\definecolor{currentstroke}{rgb}{0.000000,0.000000,0.000000}%
\pgfsetstrokecolor{currentstroke}%
\pgfsetstrokeopacity{0.000000}%
\pgfsetdash{}{0pt}%
\pgfpathmoveto{\pgfqpoint{1.432900in}{0.660000in}}%
\pgfpathlineto{\pgfqpoint{1.439100in}{0.660000in}}%
\pgfpathlineto{\pgfqpoint{1.439100in}{4.338889in}}%
\pgfpathlineto{\pgfqpoint{1.432900in}{4.338889in}}%
\pgfpathlineto{\pgfqpoint{1.432900in}{0.660000in}}%
\pgfpathclose%
\pgfusepath{fill}%
\end{pgfscope}%
\begin{pgfscope}%
\pgfpathrectangle{\pgfqpoint{1.250000in}{0.660000in}}{\pgfqpoint{7.750000in}{4.620000in}}%
\pgfusepath{clip}%
\pgfsetbuttcap%
\pgfsetmiterjoin%
\definecolor{currentfill}{rgb}{0.000000,0.000000,0.000000}%
\pgfsetfillcolor{currentfill}%
\pgfsetlinewidth{0.000000pt}%
\definecolor{currentstroke}{rgb}{0.000000,0.000000,0.000000}%
\pgfsetstrokecolor{currentstroke}%
\pgfsetstrokeopacity{0.000000}%
\pgfsetdash{}{0pt}%
\pgfpathmoveto{\pgfqpoint{1.440650in}{0.660000in}}%
\pgfpathlineto{\pgfqpoint{1.446850in}{0.660000in}}%
\pgfpathlineto{\pgfqpoint{1.446850in}{4.501492in}}%
\pgfpathlineto{\pgfqpoint{1.440650in}{4.501492in}}%
\pgfpathlineto{\pgfqpoint{1.440650in}{0.660000in}}%
\pgfpathclose%
\pgfusepath{fill}%
\end{pgfscope}%
\begin{pgfscope}%
\pgfpathrectangle{\pgfqpoint{1.250000in}{0.660000in}}{\pgfqpoint{7.750000in}{4.620000in}}%
\pgfusepath{clip}%
\pgfsetbuttcap%
\pgfsetmiterjoin%
\definecolor{currentfill}{rgb}{0.000000,0.000000,0.000000}%
\pgfsetfillcolor{currentfill}%
\pgfsetlinewidth{0.000000pt}%
\definecolor{currentstroke}{rgb}{0.000000,0.000000,0.000000}%
\pgfsetstrokecolor{currentstroke}%
\pgfsetstrokeopacity{0.000000}%
\pgfsetdash{}{0pt}%
\pgfpathmoveto{\pgfqpoint{1.448400in}{0.660000in}}%
\pgfpathlineto{\pgfqpoint{1.454600in}{0.660000in}}%
\pgfpathlineto{\pgfqpoint{1.454600in}{4.472929in}}%
\pgfpathlineto{\pgfqpoint{1.448400in}{4.472929in}}%
\pgfpathlineto{\pgfqpoint{1.448400in}{0.660000in}}%
\pgfpathclose%
\pgfusepath{fill}%
\end{pgfscope}%
\begin{pgfscope}%
\pgfpathrectangle{\pgfqpoint{1.250000in}{0.660000in}}{\pgfqpoint{7.750000in}{4.620000in}}%
\pgfusepath{clip}%
\pgfsetbuttcap%
\pgfsetmiterjoin%
\definecolor{currentfill}{rgb}{0.000000,0.000000,0.000000}%
\pgfsetfillcolor{currentfill}%
\pgfsetlinewidth{0.000000pt}%
\definecolor{currentstroke}{rgb}{0.000000,0.000000,0.000000}%
\pgfsetstrokecolor{currentstroke}%
\pgfsetstrokeopacity{0.000000}%
\pgfsetdash{}{0pt}%
\pgfpathmoveto{\pgfqpoint{1.456150in}{0.660000in}}%
\pgfpathlineto{\pgfqpoint{1.462350in}{0.660000in}}%
\pgfpathlineto{\pgfqpoint{1.462350in}{4.442876in}}%
\pgfpathlineto{\pgfqpoint{1.456150in}{4.442876in}}%
\pgfpathlineto{\pgfqpoint{1.456150in}{0.660000in}}%
\pgfpathclose%
\pgfusepath{fill}%
\end{pgfscope}%
\begin{pgfscope}%
\pgfpathrectangle{\pgfqpoint{1.250000in}{0.660000in}}{\pgfqpoint{7.750000in}{4.620000in}}%
\pgfusepath{clip}%
\pgfsetbuttcap%
\pgfsetmiterjoin%
\definecolor{currentfill}{rgb}{0.000000,0.000000,0.000000}%
\pgfsetfillcolor{currentfill}%
\pgfsetlinewidth{0.000000pt}%
\definecolor{currentstroke}{rgb}{0.000000,0.000000,0.000000}%
\pgfsetstrokecolor{currentstroke}%
\pgfsetstrokeopacity{0.000000}%
\pgfsetdash{}{0pt}%
\pgfpathmoveto{\pgfqpoint{1.463900in}{0.660000in}}%
\pgfpathlineto{\pgfqpoint{1.470100in}{0.660000in}}%
\pgfpathlineto{\pgfqpoint{1.470100in}{4.474501in}}%
\pgfpathlineto{\pgfqpoint{1.463900in}{4.474501in}}%
\pgfpathlineto{\pgfqpoint{1.463900in}{0.660000in}}%
\pgfpathclose%
\pgfusepath{fill}%
\end{pgfscope}%
\begin{pgfscope}%
\pgfpathrectangle{\pgfqpoint{1.250000in}{0.660000in}}{\pgfqpoint{7.750000in}{4.620000in}}%
\pgfusepath{clip}%
\pgfsetbuttcap%
\pgfsetmiterjoin%
\definecolor{currentfill}{rgb}{0.000000,0.000000,0.000000}%
\pgfsetfillcolor{currentfill}%
\pgfsetlinewidth{0.000000pt}%
\definecolor{currentstroke}{rgb}{0.000000,0.000000,0.000000}%
\pgfsetstrokecolor{currentstroke}%
\pgfsetstrokeopacity{0.000000}%
\pgfsetdash{}{0pt}%
\pgfpathmoveto{\pgfqpoint{1.471650in}{0.660000in}}%
\pgfpathlineto{\pgfqpoint{1.477850in}{0.660000in}}%
\pgfpathlineto{\pgfqpoint{1.477850in}{4.418723in}}%
\pgfpathlineto{\pgfqpoint{1.471650in}{4.418723in}}%
\pgfpathlineto{\pgfqpoint{1.471650in}{0.660000in}}%
\pgfpathclose%
\pgfusepath{fill}%
\end{pgfscope}%
\begin{pgfscope}%
\pgfpathrectangle{\pgfqpoint{1.250000in}{0.660000in}}{\pgfqpoint{7.750000in}{4.620000in}}%
\pgfusepath{clip}%
\pgfsetbuttcap%
\pgfsetmiterjoin%
\definecolor{currentfill}{rgb}{0.000000,0.000000,0.000000}%
\pgfsetfillcolor{currentfill}%
\pgfsetlinewidth{0.000000pt}%
\definecolor{currentstroke}{rgb}{0.000000,0.000000,0.000000}%
\pgfsetstrokecolor{currentstroke}%
\pgfsetstrokeopacity{0.000000}%
\pgfsetdash{}{0pt}%
\pgfpathmoveto{\pgfqpoint{1.479400in}{0.660000in}}%
\pgfpathlineto{\pgfqpoint{1.485600in}{0.660000in}}%
\pgfpathlineto{\pgfqpoint{1.485600in}{4.315214in}}%
\pgfpathlineto{\pgfqpoint{1.479400in}{4.315214in}}%
\pgfpathlineto{\pgfqpoint{1.479400in}{0.660000in}}%
\pgfpathclose%
\pgfusepath{fill}%
\end{pgfscope}%
\begin{pgfscope}%
\pgfpathrectangle{\pgfqpoint{1.250000in}{0.660000in}}{\pgfqpoint{7.750000in}{4.620000in}}%
\pgfusepath{clip}%
\pgfsetbuttcap%
\pgfsetmiterjoin%
\definecolor{currentfill}{rgb}{0.000000,0.000000,0.000000}%
\pgfsetfillcolor{currentfill}%
\pgfsetlinewidth{0.000000pt}%
\definecolor{currentstroke}{rgb}{0.000000,0.000000,0.000000}%
\pgfsetstrokecolor{currentstroke}%
\pgfsetstrokeopacity{0.000000}%
\pgfsetdash{}{0pt}%
\pgfpathmoveto{\pgfqpoint{1.487150in}{0.660000in}}%
\pgfpathlineto{\pgfqpoint{1.493350in}{0.660000in}}%
\pgfpathlineto{\pgfqpoint{1.493350in}{4.525846in}}%
\pgfpathlineto{\pgfqpoint{1.487150in}{4.525846in}}%
\pgfpathlineto{\pgfqpoint{1.487150in}{0.660000in}}%
\pgfpathclose%
\pgfusepath{fill}%
\end{pgfscope}%
\begin{pgfscope}%
\pgfpathrectangle{\pgfqpoint{1.250000in}{0.660000in}}{\pgfqpoint{7.750000in}{4.620000in}}%
\pgfusepath{clip}%
\pgfsetbuttcap%
\pgfsetmiterjoin%
\definecolor{currentfill}{rgb}{0.000000,0.000000,0.000000}%
\pgfsetfillcolor{currentfill}%
\pgfsetlinewidth{0.000000pt}%
\definecolor{currentstroke}{rgb}{0.000000,0.000000,0.000000}%
\pgfsetstrokecolor{currentstroke}%
\pgfsetstrokeopacity{0.000000}%
\pgfsetdash{}{0pt}%
\pgfpathmoveto{\pgfqpoint{1.494900in}{0.660000in}}%
\pgfpathlineto{\pgfqpoint{1.501100in}{0.660000in}}%
\pgfpathlineto{\pgfqpoint{1.501100in}{4.422438in}}%
\pgfpathlineto{\pgfqpoint{1.494900in}{4.422438in}}%
\pgfpathlineto{\pgfqpoint{1.494900in}{0.660000in}}%
\pgfpathclose%
\pgfusepath{fill}%
\end{pgfscope}%
\begin{pgfscope}%
\pgfpathrectangle{\pgfqpoint{1.250000in}{0.660000in}}{\pgfqpoint{7.750000in}{4.620000in}}%
\pgfusepath{clip}%
\pgfsetbuttcap%
\pgfsetmiterjoin%
\definecolor{currentfill}{rgb}{0.000000,0.000000,0.000000}%
\pgfsetfillcolor{currentfill}%
\pgfsetlinewidth{0.000000pt}%
\definecolor{currentstroke}{rgb}{0.000000,0.000000,0.000000}%
\pgfsetstrokecolor{currentstroke}%
\pgfsetstrokeopacity{0.000000}%
\pgfsetdash{}{0pt}%
\pgfpathmoveto{\pgfqpoint{1.502650in}{0.660000in}}%
\pgfpathlineto{\pgfqpoint{1.508850in}{0.660000in}}%
\pgfpathlineto{\pgfqpoint{1.508850in}{4.332075in}}%
\pgfpathlineto{\pgfqpoint{1.502650in}{4.332075in}}%
\pgfpathlineto{\pgfqpoint{1.502650in}{0.660000in}}%
\pgfpathclose%
\pgfusepath{fill}%
\end{pgfscope}%
\begin{pgfscope}%
\pgfpathrectangle{\pgfqpoint{1.250000in}{0.660000in}}{\pgfqpoint{7.750000in}{4.620000in}}%
\pgfusepath{clip}%
\pgfsetbuttcap%
\pgfsetmiterjoin%
\definecolor{currentfill}{rgb}{0.000000,0.000000,0.000000}%
\pgfsetfillcolor{currentfill}%
\pgfsetlinewidth{0.000000pt}%
\definecolor{currentstroke}{rgb}{0.000000,0.000000,0.000000}%
\pgfsetstrokecolor{currentstroke}%
\pgfsetstrokeopacity{0.000000}%
\pgfsetdash{}{0pt}%
\pgfpathmoveto{\pgfqpoint{1.510400in}{0.660000in}}%
\pgfpathlineto{\pgfqpoint{1.516600in}{0.660000in}}%
\pgfpathlineto{\pgfqpoint{1.516600in}{4.373936in}}%
\pgfpathlineto{\pgfqpoint{1.510400in}{4.373936in}}%
\pgfpathlineto{\pgfqpoint{1.510400in}{0.660000in}}%
\pgfpathclose%
\pgfusepath{fill}%
\end{pgfscope}%
\begin{pgfscope}%
\pgfpathrectangle{\pgfqpoint{1.250000in}{0.660000in}}{\pgfqpoint{7.750000in}{4.620000in}}%
\pgfusepath{clip}%
\pgfsetbuttcap%
\pgfsetmiterjoin%
\definecolor{currentfill}{rgb}{0.000000,0.000000,0.000000}%
\pgfsetfillcolor{currentfill}%
\pgfsetlinewidth{0.000000pt}%
\definecolor{currentstroke}{rgb}{0.000000,0.000000,0.000000}%
\pgfsetstrokecolor{currentstroke}%
\pgfsetstrokeopacity{0.000000}%
\pgfsetdash{}{0pt}%
\pgfpathmoveto{\pgfqpoint{1.518150in}{0.660000in}}%
\pgfpathlineto{\pgfqpoint{1.524350in}{0.660000in}}%
\pgfpathlineto{\pgfqpoint{1.524350in}{4.232116in}}%
\pgfpathlineto{\pgfqpoint{1.518150in}{4.232116in}}%
\pgfpathlineto{\pgfqpoint{1.518150in}{0.660000in}}%
\pgfpathclose%
\pgfusepath{fill}%
\end{pgfscope}%
\begin{pgfscope}%
\pgfpathrectangle{\pgfqpoint{1.250000in}{0.660000in}}{\pgfqpoint{7.750000in}{4.620000in}}%
\pgfusepath{clip}%
\pgfsetbuttcap%
\pgfsetmiterjoin%
\definecolor{currentfill}{rgb}{0.000000,0.000000,0.000000}%
\pgfsetfillcolor{currentfill}%
\pgfsetlinewidth{0.000000pt}%
\definecolor{currentstroke}{rgb}{0.000000,0.000000,0.000000}%
\pgfsetstrokecolor{currentstroke}%
\pgfsetstrokeopacity{0.000000}%
\pgfsetdash{}{0pt}%
\pgfpathmoveto{\pgfqpoint{1.525900in}{0.660000in}}%
\pgfpathlineto{\pgfqpoint{1.532100in}{0.660000in}}%
\pgfpathlineto{\pgfqpoint{1.532100in}{4.238756in}}%
\pgfpathlineto{\pgfqpoint{1.525900in}{4.238756in}}%
\pgfpathlineto{\pgfqpoint{1.525900in}{0.660000in}}%
\pgfpathclose%
\pgfusepath{fill}%
\end{pgfscope}%
\begin{pgfscope}%
\pgfpathrectangle{\pgfqpoint{1.250000in}{0.660000in}}{\pgfqpoint{7.750000in}{4.620000in}}%
\pgfusepath{clip}%
\pgfsetbuttcap%
\pgfsetmiterjoin%
\definecolor{currentfill}{rgb}{0.000000,0.000000,0.000000}%
\pgfsetfillcolor{currentfill}%
\pgfsetlinewidth{0.000000pt}%
\definecolor{currentstroke}{rgb}{0.000000,0.000000,0.000000}%
\pgfsetstrokecolor{currentstroke}%
\pgfsetstrokeopacity{0.000000}%
\pgfsetdash{}{0pt}%
\pgfpathmoveto{\pgfqpoint{1.533650in}{0.660000in}}%
\pgfpathlineto{\pgfqpoint{1.539850in}{0.660000in}}%
\pgfpathlineto{\pgfqpoint{1.539850in}{4.161159in}}%
\pgfpathlineto{\pgfqpoint{1.533650in}{4.161159in}}%
\pgfpathlineto{\pgfqpoint{1.533650in}{0.660000in}}%
\pgfpathclose%
\pgfusepath{fill}%
\end{pgfscope}%
\begin{pgfscope}%
\pgfpathrectangle{\pgfqpoint{1.250000in}{0.660000in}}{\pgfqpoint{7.750000in}{4.620000in}}%
\pgfusepath{clip}%
\pgfsetbuttcap%
\pgfsetmiterjoin%
\definecolor{currentfill}{rgb}{0.000000,0.000000,0.000000}%
\pgfsetfillcolor{currentfill}%
\pgfsetlinewidth{0.000000pt}%
\definecolor{currentstroke}{rgb}{0.000000,0.000000,0.000000}%
\pgfsetstrokecolor{currentstroke}%
\pgfsetstrokeopacity{0.000000}%
\pgfsetdash{}{0pt}%
\pgfpathmoveto{\pgfqpoint{1.541400in}{0.660000in}}%
\pgfpathlineto{\pgfqpoint{1.547600in}{0.660000in}}%
\pgfpathlineto{\pgfqpoint{1.547600in}{4.096252in}}%
\pgfpathlineto{\pgfqpoint{1.541400in}{4.096252in}}%
\pgfpathlineto{\pgfqpoint{1.541400in}{0.660000in}}%
\pgfpathclose%
\pgfusepath{fill}%
\end{pgfscope}%
\begin{pgfscope}%
\pgfpathrectangle{\pgfqpoint{1.250000in}{0.660000in}}{\pgfqpoint{7.750000in}{4.620000in}}%
\pgfusepath{clip}%
\pgfsetbuttcap%
\pgfsetmiterjoin%
\definecolor{currentfill}{rgb}{0.000000,0.000000,0.000000}%
\pgfsetfillcolor{currentfill}%
\pgfsetlinewidth{0.000000pt}%
\definecolor{currentstroke}{rgb}{0.000000,0.000000,0.000000}%
\pgfsetstrokecolor{currentstroke}%
\pgfsetstrokeopacity{0.000000}%
\pgfsetdash{}{0pt}%
\pgfpathmoveto{\pgfqpoint{1.549150in}{0.660000in}}%
\pgfpathlineto{\pgfqpoint{1.555350in}{0.660000in}}%
\pgfpathlineto{\pgfqpoint{1.555350in}{4.106856in}}%
\pgfpathlineto{\pgfqpoint{1.549150in}{4.106856in}}%
\pgfpathlineto{\pgfqpoint{1.549150in}{0.660000in}}%
\pgfpathclose%
\pgfusepath{fill}%
\end{pgfscope}%
\begin{pgfscope}%
\pgfpathrectangle{\pgfqpoint{1.250000in}{0.660000in}}{\pgfqpoint{7.750000in}{4.620000in}}%
\pgfusepath{clip}%
\pgfsetbuttcap%
\pgfsetmiterjoin%
\definecolor{currentfill}{rgb}{0.000000,0.000000,0.000000}%
\pgfsetfillcolor{currentfill}%
\pgfsetlinewidth{0.000000pt}%
\definecolor{currentstroke}{rgb}{0.000000,0.000000,0.000000}%
\pgfsetstrokecolor{currentstroke}%
\pgfsetstrokeopacity{0.000000}%
\pgfsetdash{}{0pt}%
\pgfpathmoveto{\pgfqpoint{1.556900in}{0.660000in}}%
\pgfpathlineto{\pgfqpoint{1.563100in}{0.660000in}}%
\pgfpathlineto{\pgfqpoint{1.563100in}{4.457206in}}%
\pgfpathlineto{\pgfqpoint{1.556900in}{4.457206in}}%
\pgfpathlineto{\pgfqpoint{1.556900in}{0.660000in}}%
\pgfpathclose%
\pgfusepath{fill}%
\end{pgfscope}%
\begin{pgfscope}%
\pgfpathrectangle{\pgfqpoint{1.250000in}{0.660000in}}{\pgfqpoint{7.750000in}{4.620000in}}%
\pgfusepath{clip}%
\pgfsetbuttcap%
\pgfsetmiterjoin%
\definecolor{currentfill}{rgb}{0.000000,0.000000,0.000000}%
\pgfsetfillcolor{currentfill}%
\pgfsetlinewidth{0.000000pt}%
\definecolor{currentstroke}{rgb}{0.000000,0.000000,0.000000}%
\pgfsetstrokecolor{currentstroke}%
\pgfsetstrokeopacity{0.000000}%
\pgfsetdash{}{0pt}%
\pgfpathmoveto{\pgfqpoint{1.564650in}{0.660000in}}%
\pgfpathlineto{\pgfqpoint{1.570850in}{0.660000in}}%
\pgfpathlineto{\pgfqpoint{1.570850in}{4.137725in}}%
\pgfpathlineto{\pgfqpoint{1.564650in}{4.137725in}}%
\pgfpathlineto{\pgfqpoint{1.564650in}{0.660000in}}%
\pgfpathclose%
\pgfusepath{fill}%
\end{pgfscope}%
\begin{pgfscope}%
\pgfpathrectangle{\pgfqpoint{1.250000in}{0.660000in}}{\pgfqpoint{7.750000in}{4.620000in}}%
\pgfusepath{clip}%
\pgfsetbuttcap%
\pgfsetmiterjoin%
\definecolor{currentfill}{rgb}{0.000000,0.000000,0.000000}%
\pgfsetfillcolor{currentfill}%
\pgfsetlinewidth{0.000000pt}%
\definecolor{currentstroke}{rgb}{0.000000,0.000000,0.000000}%
\pgfsetstrokecolor{currentstroke}%
\pgfsetstrokeopacity{0.000000}%
\pgfsetdash{}{0pt}%
\pgfpathmoveto{\pgfqpoint{1.572400in}{0.660000in}}%
\pgfpathlineto{\pgfqpoint{1.578600in}{0.660000in}}%
\pgfpathlineto{\pgfqpoint{1.578600in}{4.215007in}}%
\pgfpathlineto{\pgfqpoint{1.572400in}{4.215007in}}%
\pgfpathlineto{\pgfqpoint{1.572400in}{0.660000in}}%
\pgfpathclose%
\pgfusepath{fill}%
\end{pgfscope}%
\begin{pgfscope}%
\pgfpathrectangle{\pgfqpoint{1.250000in}{0.660000in}}{\pgfqpoint{7.750000in}{4.620000in}}%
\pgfusepath{clip}%
\pgfsetbuttcap%
\pgfsetmiterjoin%
\definecolor{currentfill}{rgb}{0.000000,0.000000,0.000000}%
\pgfsetfillcolor{currentfill}%
\pgfsetlinewidth{0.000000pt}%
\definecolor{currentstroke}{rgb}{0.000000,0.000000,0.000000}%
\pgfsetstrokecolor{currentstroke}%
\pgfsetstrokeopacity{0.000000}%
\pgfsetdash{}{0pt}%
\pgfpathmoveto{\pgfqpoint{1.580150in}{0.660000in}}%
\pgfpathlineto{\pgfqpoint{1.586350in}{0.660000in}}%
\pgfpathlineto{\pgfqpoint{1.586350in}{4.066213in}}%
\pgfpathlineto{\pgfqpoint{1.580150in}{4.066213in}}%
\pgfpathlineto{\pgfqpoint{1.580150in}{0.660000in}}%
\pgfpathclose%
\pgfusepath{fill}%
\end{pgfscope}%
\begin{pgfscope}%
\pgfpathrectangle{\pgfqpoint{1.250000in}{0.660000in}}{\pgfqpoint{7.750000in}{4.620000in}}%
\pgfusepath{clip}%
\pgfsetbuttcap%
\pgfsetmiterjoin%
\definecolor{currentfill}{rgb}{0.000000,0.000000,0.000000}%
\pgfsetfillcolor{currentfill}%
\pgfsetlinewidth{0.000000pt}%
\definecolor{currentstroke}{rgb}{0.000000,0.000000,0.000000}%
\pgfsetstrokecolor{currentstroke}%
\pgfsetstrokeopacity{0.000000}%
\pgfsetdash{}{0pt}%
\pgfpathmoveto{\pgfqpoint{1.587900in}{0.660000in}}%
\pgfpathlineto{\pgfqpoint{1.594100in}{0.660000in}}%
\pgfpathlineto{\pgfqpoint{1.594100in}{4.094292in}}%
\pgfpathlineto{\pgfqpoint{1.587900in}{4.094292in}}%
\pgfpathlineto{\pgfqpoint{1.587900in}{0.660000in}}%
\pgfpathclose%
\pgfusepath{fill}%
\end{pgfscope}%
\begin{pgfscope}%
\pgfpathrectangle{\pgfqpoint{1.250000in}{0.660000in}}{\pgfqpoint{7.750000in}{4.620000in}}%
\pgfusepath{clip}%
\pgfsetbuttcap%
\pgfsetmiterjoin%
\definecolor{currentfill}{rgb}{0.000000,0.000000,0.000000}%
\pgfsetfillcolor{currentfill}%
\pgfsetlinewidth{0.000000pt}%
\definecolor{currentstroke}{rgb}{0.000000,0.000000,0.000000}%
\pgfsetstrokecolor{currentstroke}%
\pgfsetstrokeopacity{0.000000}%
\pgfsetdash{}{0pt}%
\pgfpathmoveto{\pgfqpoint{1.595650in}{0.660000in}}%
\pgfpathlineto{\pgfqpoint{1.601850in}{0.660000in}}%
\pgfpathlineto{\pgfqpoint{1.601850in}{4.013903in}}%
\pgfpathlineto{\pgfqpoint{1.595650in}{4.013903in}}%
\pgfpathlineto{\pgfqpoint{1.595650in}{0.660000in}}%
\pgfpathclose%
\pgfusepath{fill}%
\end{pgfscope}%
\begin{pgfscope}%
\pgfpathrectangle{\pgfqpoint{1.250000in}{0.660000in}}{\pgfqpoint{7.750000in}{4.620000in}}%
\pgfusepath{clip}%
\pgfsetbuttcap%
\pgfsetmiterjoin%
\definecolor{currentfill}{rgb}{0.000000,0.000000,0.000000}%
\pgfsetfillcolor{currentfill}%
\pgfsetlinewidth{0.000000pt}%
\definecolor{currentstroke}{rgb}{0.000000,0.000000,0.000000}%
\pgfsetstrokecolor{currentstroke}%
\pgfsetstrokeopacity{0.000000}%
\pgfsetdash{}{0pt}%
\pgfpathmoveto{\pgfqpoint{1.603400in}{0.660000in}}%
\pgfpathlineto{\pgfqpoint{1.609600in}{0.660000in}}%
\pgfpathlineto{\pgfqpoint{1.609600in}{4.252713in}}%
\pgfpathlineto{\pgfqpoint{1.603400in}{4.252713in}}%
\pgfpathlineto{\pgfqpoint{1.603400in}{0.660000in}}%
\pgfpathclose%
\pgfusepath{fill}%
\end{pgfscope}%
\begin{pgfscope}%
\pgfpathrectangle{\pgfqpoint{1.250000in}{0.660000in}}{\pgfqpoint{7.750000in}{4.620000in}}%
\pgfusepath{clip}%
\pgfsetbuttcap%
\pgfsetmiterjoin%
\definecolor{currentfill}{rgb}{0.000000,0.000000,0.000000}%
\pgfsetfillcolor{currentfill}%
\pgfsetlinewidth{0.000000pt}%
\definecolor{currentstroke}{rgb}{0.000000,0.000000,0.000000}%
\pgfsetstrokecolor{currentstroke}%
\pgfsetstrokeopacity{0.000000}%
\pgfsetdash{}{0pt}%
\pgfpathmoveto{\pgfqpoint{1.611150in}{0.660000in}}%
\pgfpathlineto{\pgfqpoint{1.617350in}{0.660000in}}%
\pgfpathlineto{\pgfqpoint{1.617350in}{4.236609in}}%
\pgfpathlineto{\pgfqpoint{1.611150in}{4.236609in}}%
\pgfpathlineto{\pgfqpoint{1.611150in}{0.660000in}}%
\pgfpathclose%
\pgfusepath{fill}%
\end{pgfscope}%
\begin{pgfscope}%
\pgfpathrectangle{\pgfqpoint{1.250000in}{0.660000in}}{\pgfqpoint{7.750000in}{4.620000in}}%
\pgfusepath{clip}%
\pgfsetbuttcap%
\pgfsetmiterjoin%
\definecolor{currentfill}{rgb}{0.000000,0.000000,0.000000}%
\pgfsetfillcolor{currentfill}%
\pgfsetlinewidth{0.000000pt}%
\definecolor{currentstroke}{rgb}{0.000000,0.000000,0.000000}%
\pgfsetstrokecolor{currentstroke}%
\pgfsetstrokeopacity{0.000000}%
\pgfsetdash{}{0pt}%
\pgfpathmoveto{\pgfqpoint{1.618900in}{0.660000in}}%
\pgfpathlineto{\pgfqpoint{1.625100in}{0.660000in}}%
\pgfpathlineto{\pgfqpoint{1.625100in}{4.029964in}}%
\pgfpathlineto{\pgfqpoint{1.618900in}{4.029964in}}%
\pgfpathlineto{\pgfqpoint{1.618900in}{0.660000in}}%
\pgfpathclose%
\pgfusepath{fill}%
\end{pgfscope}%
\begin{pgfscope}%
\pgfpathrectangle{\pgfqpoint{1.250000in}{0.660000in}}{\pgfqpoint{7.750000in}{4.620000in}}%
\pgfusepath{clip}%
\pgfsetbuttcap%
\pgfsetmiterjoin%
\definecolor{currentfill}{rgb}{0.000000,0.000000,0.000000}%
\pgfsetfillcolor{currentfill}%
\pgfsetlinewidth{0.000000pt}%
\definecolor{currentstroke}{rgb}{0.000000,0.000000,0.000000}%
\pgfsetstrokecolor{currentstroke}%
\pgfsetstrokeopacity{0.000000}%
\pgfsetdash{}{0pt}%
\pgfpathmoveto{\pgfqpoint{1.626650in}{0.660000in}}%
\pgfpathlineto{\pgfqpoint{1.632850in}{0.660000in}}%
\pgfpathlineto{\pgfqpoint{1.632850in}{4.004949in}}%
\pgfpathlineto{\pgfqpoint{1.626650in}{4.004949in}}%
\pgfpathlineto{\pgfqpoint{1.626650in}{0.660000in}}%
\pgfpathclose%
\pgfusepath{fill}%
\end{pgfscope}%
\begin{pgfscope}%
\pgfpathrectangle{\pgfqpoint{1.250000in}{0.660000in}}{\pgfqpoint{7.750000in}{4.620000in}}%
\pgfusepath{clip}%
\pgfsetbuttcap%
\pgfsetmiterjoin%
\definecolor{currentfill}{rgb}{0.000000,0.000000,0.000000}%
\pgfsetfillcolor{currentfill}%
\pgfsetlinewidth{0.000000pt}%
\definecolor{currentstroke}{rgb}{0.000000,0.000000,0.000000}%
\pgfsetstrokecolor{currentstroke}%
\pgfsetstrokeopacity{0.000000}%
\pgfsetdash{}{0pt}%
\pgfpathmoveto{\pgfqpoint{1.634400in}{0.660000in}}%
\pgfpathlineto{\pgfqpoint{1.640600in}{0.660000in}}%
\pgfpathlineto{\pgfqpoint{1.640600in}{4.064994in}}%
\pgfpathlineto{\pgfqpoint{1.634400in}{4.064994in}}%
\pgfpathlineto{\pgfqpoint{1.634400in}{0.660000in}}%
\pgfpathclose%
\pgfusepath{fill}%
\end{pgfscope}%
\begin{pgfscope}%
\pgfpathrectangle{\pgfqpoint{1.250000in}{0.660000in}}{\pgfqpoint{7.750000in}{4.620000in}}%
\pgfusepath{clip}%
\pgfsetbuttcap%
\pgfsetmiterjoin%
\definecolor{currentfill}{rgb}{0.000000,0.000000,0.000000}%
\pgfsetfillcolor{currentfill}%
\pgfsetlinewidth{0.000000pt}%
\definecolor{currentstroke}{rgb}{0.000000,0.000000,0.000000}%
\pgfsetstrokecolor{currentstroke}%
\pgfsetstrokeopacity{0.000000}%
\pgfsetdash{}{0pt}%
\pgfpathmoveto{\pgfqpoint{1.642150in}{0.660000in}}%
\pgfpathlineto{\pgfqpoint{1.648350in}{0.660000in}}%
\pgfpathlineto{\pgfqpoint{1.648350in}{4.017785in}}%
\pgfpathlineto{\pgfqpoint{1.642150in}{4.017785in}}%
\pgfpathlineto{\pgfqpoint{1.642150in}{0.660000in}}%
\pgfpathclose%
\pgfusepath{fill}%
\end{pgfscope}%
\begin{pgfscope}%
\pgfpathrectangle{\pgfqpoint{1.250000in}{0.660000in}}{\pgfqpoint{7.750000in}{4.620000in}}%
\pgfusepath{clip}%
\pgfsetbuttcap%
\pgfsetmiterjoin%
\definecolor{currentfill}{rgb}{0.000000,0.000000,0.000000}%
\pgfsetfillcolor{currentfill}%
\pgfsetlinewidth{0.000000pt}%
\definecolor{currentstroke}{rgb}{0.000000,0.000000,0.000000}%
\pgfsetstrokecolor{currentstroke}%
\pgfsetstrokeopacity{0.000000}%
\pgfsetdash{}{0pt}%
\pgfpathmoveto{\pgfqpoint{1.649900in}{0.660000in}}%
\pgfpathlineto{\pgfqpoint{1.656100in}{0.660000in}}%
\pgfpathlineto{\pgfqpoint{1.656100in}{4.013891in}}%
\pgfpathlineto{\pgfqpoint{1.649900in}{4.013891in}}%
\pgfpathlineto{\pgfqpoint{1.649900in}{0.660000in}}%
\pgfpathclose%
\pgfusepath{fill}%
\end{pgfscope}%
\begin{pgfscope}%
\pgfpathrectangle{\pgfqpoint{1.250000in}{0.660000in}}{\pgfqpoint{7.750000in}{4.620000in}}%
\pgfusepath{clip}%
\pgfsetbuttcap%
\pgfsetmiterjoin%
\definecolor{currentfill}{rgb}{0.000000,0.000000,0.000000}%
\pgfsetfillcolor{currentfill}%
\pgfsetlinewidth{0.000000pt}%
\definecolor{currentstroke}{rgb}{0.000000,0.000000,0.000000}%
\pgfsetstrokecolor{currentstroke}%
\pgfsetstrokeopacity{0.000000}%
\pgfsetdash{}{0pt}%
\pgfpathmoveto{\pgfqpoint{1.657650in}{0.660000in}}%
\pgfpathlineto{\pgfqpoint{1.663850in}{0.660000in}}%
\pgfpathlineto{\pgfqpoint{1.663850in}{4.053157in}}%
\pgfpathlineto{\pgfqpoint{1.657650in}{4.053157in}}%
\pgfpathlineto{\pgfqpoint{1.657650in}{0.660000in}}%
\pgfpathclose%
\pgfusepath{fill}%
\end{pgfscope}%
\begin{pgfscope}%
\pgfpathrectangle{\pgfqpoint{1.250000in}{0.660000in}}{\pgfqpoint{7.750000in}{4.620000in}}%
\pgfusepath{clip}%
\pgfsetbuttcap%
\pgfsetmiterjoin%
\definecolor{currentfill}{rgb}{0.000000,0.000000,0.000000}%
\pgfsetfillcolor{currentfill}%
\pgfsetlinewidth{0.000000pt}%
\definecolor{currentstroke}{rgb}{0.000000,0.000000,0.000000}%
\pgfsetstrokecolor{currentstroke}%
\pgfsetstrokeopacity{0.000000}%
\pgfsetdash{}{0pt}%
\pgfpathmoveto{\pgfqpoint{1.665400in}{0.660000in}}%
\pgfpathlineto{\pgfqpoint{1.671600in}{0.660000in}}%
\pgfpathlineto{\pgfqpoint{1.671600in}{3.937763in}}%
\pgfpathlineto{\pgfqpoint{1.665400in}{3.937763in}}%
\pgfpathlineto{\pgfqpoint{1.665400in}{0.660000in}}%
\pgfpathclose%
\pgfusepath{fill}%
\end{pgfscope}%
\begin{pgfscope}%
\pgfpathrectangle{\pgfqpoint{1.250000in}{0.660000in}}{\pgfqpoint{7.750000in}{4.620000in}}%
\pgfusepath{clip}%
\pgfsetbuttcap%
\pgfsetmiterjoin%
\definecolor{currentfill}{rgb}{0.000000,0.000000,0.000000}%
\pgfsetfillcolor{currentfill}%
\pgfsetlinewidth{0.000000pt}%
\definecolor{currentstroke}{rgb}{0.000000,0.000000,0.000000}%
\pgfsetstrokecolor{currentstroke}%
\pgfsetstrokeopacity{0.000000}%
\pgfsetdash{}{0pt}%
\pgfpathmoveto{\pgfqpoint{1.673150in}{0.660000in}}%
\pgfpathlineto{\pgfqpoint{1.679350in}{0.660000in}}%
\pgfpathlineto{\pgfqpoint{1.679350in}{3.939245in}}%
\pgfpathlineto{\pgfqpoint{1.673150in}{3.939245in}}%
\pgfpathlineto{\pgfqpoint{1.673150in}{0.660000in}}%
\pgfpathclose%
\pgfusepath{fill}%
\end{pgfscope}%
\begin{pgfscope}%
\pgfpathrectangle{\pgfqpoint{1.250000in}{0.660000in}}{\pgfqpoint{7.750000in}{4.620000in}}%
\pgfusepath{clip}%
\pgfsetbuttcap%
\pgfsetmiterjoin%
\definecolor{currentfill}{rgb}{0.000000,0.000000,0.000000}%
\pgfsetfillcolor{currentfill}%
\pgfsetlinewidth{0.000000pt}%
\definecolor{currentstroke}{rgb}{0.000000,0.000000,0.000000}%
\pgfsetstrokecolor{currentstroke}%
\pgfsetstrokeopacity{0.000000}%
\pgfsetdash{}{0pt}%
\pgfpathmoveto{\pgfqpoint{1.680900in}{0.660000in}}%
\pgfpathlineto{\pgfqpoint{1.687100in}{0.660000in}}%
\pgfpathlineto{\pgfqpoint{1.687100in}{3.972109in}}%
\pgfpathlineto{\pgfqpoint{1.680900in}{3.972109in}}%
\pgfpathlineto{\pgfqpoint{1.680900in}{0.660000in}}%
\pgfpathclose%
\pgfusepath{fill}%
\end{pgfscope}%
\begin{pgfscope}%
\pgfpathrectangle{\pgfqpoint{1.250000in}{0.660000in}}{\pgfqpoint{7.750000in}{4.620000in}}%
\pgfusepath{clip}%
\pgfsetbuttcap%
\pgfsetmiterjoin%
\definecolor{currentfill}{rgb}{0.000000,0.000000,0.000000}%
\pgfsetfillcolor{currentfill}%
\pgfsetlinewidth{0.000000pt}%
\definecolor{currentstroke}{rgb}{0.000000,0.000000,0.000000}%
\pgfsetstrokecolor{currentstroke}%
\pgfsetstrokeopacity{0.000000}%
\pgfsetdash{}{0pt}%
\pgfpathmoveto{\pgfqpoint{1.688650in}{0.660000in}}%
\pgfpathlineto{\pgfqpoint{1.694850in}{0.660000in}}%
\pgfpathlineto{\pgfqpoint{1.694850in}{3.983873in}}%
\pgfpathlineto{\pgfqpoint{1.688650in}{3.983873in}}%
\pgfpathlineto{\pgfqpoint{1.688650in}{0.660000in}}%
\pgfpathclose%
\pgfusepath{fill}%
\end{pgfscope}%
\begin{pgfscope}%
\pgfpathrectangle{\pgfqpoint{1.250000in}{0.660000in}}{\pgfqpoint{7.750000in}{4.620000in}}%
\pgfusepath{clip}%
\pgfsetbuttcap%
\pgfsetmiterjoin%
\definecolor{currentfill}{rgb}{0.000000,0.000000,0.000000}%
\pgfsetfillcolor{currentfill}%
\pgfsetlinewidth{0.000000pt}%
\definecolor{currentstroke}{rgb}{0.000000,0.000000,0.000000}%
\pgfsetstrokecolor{currentstroke}%
\pgfsetstrokeopacity{0.000000}%
\pgfsetdash{}{0pt}%
\pgfpathmoveto{\pgfqpoint{1.696400in}{0.660000in}}%
\pgfpathlineto{\pgfqpoint{1.702600in}{0.660000in}}%
\pgfpathlineto{\pgfqpoint{1.702600in}{3.942193in}}%
\pgfpathlineto{\pgfqpoint{1.696400in}{3.942193in}}%
\pgfpathlineto{\pgfqpoint{1.696400in}{0.660000in}}%
\pgfpathclose%
\pgfusepath{fill}%
\end{pgfscope}%
\begin{pgfscope}%
\pgfpathrectangle{\pgfqpoint{1.250000in}{0.660000in}}{\pgfqpoint{7.750000in}{4.620000in}}%
\pgfusepath{clip}%
\pgfsetbuttcap%
\pgfsetmiterjoin%
\definecolor{currentfill}{rgb}{0.000000,0.000000,0.000000}%
\pgfsetfillcolor{currentfill}%
\pgfsetlinewidth{0.000000pt}%
\definecolor{currentstroke}{rgb}{0.000000,0.000000,0.000000}%
\pgfsetstrokecolor{currentstroke}%
\pgfsetstrokeopacity{0.000000}%
\pgfsetdash{}{0pt}%
\pgfpathmoveto{\pgfqpoint{1.704150in}{0.660000in}}%
\pgfpathlineto{\pgfqpoint{1.710350in}{0.660000in}}%
\pgfpathlineto{\pgfqpoint{1.710350in}{3.976945in}}%
\pgfpathlineto{\pgfqpoint{1.704150in}{3.976945in}}%
\pgfpathlineto{\pgfqpoint{1.704150in}{0.660000in}}%
\pgfpathclose%
\pgfusepath{fill}%
\end{pgfscope}%
\begin{pgfscope}%
\pgfpathrectangle{\pgfqpoint{1.250000in}{0.660000in}}{\pgfqpoint{7.750000in}{4.620000in}}%
\pgfusepath{clip}%
\pgfsetbuttcap%
\pgfsetmiterjoin%
\definecolor{currentfill}{rgb}{0.000000,0.000000,0.000000}%
\pgfsetfillcolor{currentfill}%
\pgfsetlinewidth{0.000000pt}%
\definecolor{currentstroke}{rgb}{0.000000,0.000000,0.000000}%
\pgfsetstrokecolor{currentstroke}%
\pgfsetstrokeopacity{0.000000}%
\pgfsetdash{}{0pt}%
\pgfpathmoveto{\pgfqpoint{1.711900in}{0.660000in}}%
\pgfpathlineto{\pgfqpoint{1.718100in}{0.660000in}}%
\pgfpathlineto{\pgfqpoint{1.718100in}{3.893338in}}%
\pgfpathlineto{\pgfqpoint{1.711900in}{3.893338in}}%
\pgfpathlineto{\pgfqpoint{1.711900in}{0.660000in}}%
\pgfpathclose%
\pgfusepath{fill}%
\end{pgfscope}%
\begin{pgfscope}%
\pgfpathrectangle{\pgfqpoint{1.250000in}{0.660000in}}{\pgfqpoint{7.750000in}{4.620000in}}%
\pgfusepath{clip}%
\pgfsetbuttcap%
\pgfsetmiterjoin%
\definecolor{currentfill}{rgb}{0.000000,0.000000,0.000000}%
\pgfsetfillcolor{currentfill}%
\pgfsetlinewidth{0.000000pt}%
\definecolor{currentstroke}{rgb}{0.000000,0.000000,0.000000}%
\pgfsetstrokecolor{currentstroke}%
\pgfsetstrokeopacity{0.000000}%
\pgfsetdash{}{0pt}%
\pgfpathmoveto{\pgfqpoint{1.719650in}{0.660000in}}%
\pgfpathlineto{\pgfqpoint{1.725850in}{0.660000in}}%
\pgfpathlineto{\pgfqpoint{1.725850in}{4.084730in}}%
\pgfpathlineto{\pgfqpoint{1.719650in}{4.084730in}}%
\pgfpathlineto{\pgfqpoint{1.719650in}{0.660000in}}%
\pgfpathclose%
\pgfusepath{fill}%
\end{pgfscope}%
\begin{pgfscope}%
\pgfpathrectangle{\pgfqpoint{1.250000in}{0.660000in}}{\pgfqpoint{7.750000in}{4.620000in}}%
\pgfusepath{clip}%
\pgfsetbuttcap%
\pgfsetmiterjoin%
\definecolor{currentfill}{rgb}{0.000000,0.000000,0.000000}%
\pgfsetfillcolor{currentfill}%
\pgfsetlinewidth{0.000000pt}%
\definecolor{currentstroke}{rgb}{0.000000,0.000000,0.000000}%
\pgfsetstrokecolor{currentstroke}%
\pgfsetstrokeopacity{0.000000}%
\pgfsetdash{}{0pt}%
\pgfpathmoveto{\pgfqpoint{1.727400in}{0.660000in}}%
\pgfpathlineto{\pgfqpoint{1.733600in}{0.660000in}}%
\pgfpathlineto{\pgfqpoint{1.733600in}{3.947785in}}%
\pgfpathlineto{\pgfqpoint{1.727400in}{3.947785in}}%
\pgfpathlineto{\pgfqpoint{1.727400in}{0.660000in}}%
\pgfpathclose%
\pgfusepath{fill}%
\end{pgfscope}%
\begin{pgfscope}%
\pgfpathrectangle{\pgfqpoint{1.250000in}{0.660000in}}{\pgfqpoint{7.750000in}{4.620000in}}%
\pgfusepath{clip}%
\pgfsetbuttcap%
\pgfsetmiterjoin%
\definecolor{currentfill}{rgb}{0.000000,0.000000,0.000000}%
\pgfsetfillcolor{currentfill}%
\pgfsetlinewidth{0.000000pt}%
\definecolor{currentstroke}{rgb}{0.000000,0.000000,0.000000}%
\pgfsetstrokecolor{currentstroke}%
\pgfsetstrokeopacity{0.000000}%
\pgfsetdash{}{0pt}%
\pgfpathmoveto{\pgfqpoint{1.735150in}{0.660000in}}%
\pgfpathlineto{\pgfqpoint{1.741350in}{0.660000in}}%
\pgfpathlineto{\pgfqpoint{1.741350in}{3.949181in}}%
\pgfpathlineto{\pgfqpoint{1.735150in}{3.949181in}}%
\pgfpathlineto{\pgfqpoint{1.735150in}{0.660000in}}%
\pgfpathclose%
\pgfusepath{fill}%
\end{pgfscope}%
\begin{pgfscope}%
\pgfpathrectangle{\pgfqpoint{1.250000in}{0.660000in}}{\pgfqpoint{7.750000in}{4.620000in}}%
\pgfusepath{clip}%
\pgfsetbuttcap%
\pgfsetmiterjoin%
\definecolor{currentfill}{rgb}{0.000000,0.000000,0.000000}%
\pgfsetfillcolor{currentfill}%
\pgfsetlinewidth{0.000000pt}%
\definecolor{currentstroke}{rgb}{0.000000,0.000000,0.000000}%
\pgfsetstrokecolor{currentstroke}%
\pgfsetstrokeopacity{0.000000}%
\pgfsetdash{}{0pt}%
\pgfpathmoveto{\pgfqpoint{1.742900in}{0.660000in}}%
\pgfpathlineto{\pgfqpoint{1.749100in}{0.660000in}}%
\pgfpathlineto{\pgfqpoint{1.749100in}{3.859689in}}%
\pgfpathlineto{\pgfqpoint{1.742900in}{3.859689in}}%
\pgfpathlineto{\pgfqpoint{1.742900in}{0.660000in}}%
\pgfpathclose%
\pgfusepath{fill}%
\end{pgfscope}%
\begin{pgfscope}%
\pgfpathrectangle{\pgfqpoint{1.250000in}{0.660000in}}{\pgfqpoint{7.750000in}{4.620000in}}%
\pgfusepath{clip}%
\pgfsetbuttcap%
\pgfsetmiterjoin%
\definecolor{currentfill}{rgb}{0.000000,0.000000,0.000000}%
\pgfsetfillcolor{currentfill}%
\pgfsetlinewidth{0.000000pt}%
\definecolor{currentstroke}{rgb}{0.000000,0.000000,0.000000}%
\pgfsetstrokecolor{currentstroke}%
\pgfsetstrokeopacity{0.000000}%
\pgfsetdash{}{0pt}%
\pgfpathmoveto{\pgfqpoint{1.750650in}{0.660000in}}%
\pgfpathlineto{\pgfqpoint{1.756850in}{0.660000in}}%
\pgfpathlineto{\pgfqpoint{1.756850in}{4.064599in}}%
\pgfpathlineto{\pgfqpoint{1.750650in}{4.064599in}}%
\pgfpathlineto{\pgfqpoint{1.750650in}{0.660000in}}%
\pgfpathclose%
\pgfusepath{fill}%
\end{pgfscope}%
\begin{pgfscope}%
\pgfpathrectangle{\pgfqpoint{1.250000in}{0.660000in}}{\pgfqpoint{7.750000in}{4.620000in}}%
\pgfusepath{clip}%
\pgfsetbuttcap%
\pgfsetmiterjoin%
\definecolor{currentfill}{rgb}{0.000000,0.000000,0.000000}%
\pgfsetfillcolor{currentfill}%
\pgfsetlinewidth{0.000000pt}%
\definecolor{currentstroke}{rgb}{0.000000,0.000000,0.000000}%
\pgfsetstrokecolor{currentstroke}%
\pgfsetstrokeopacity{0.000000}%
\pgfsetdash{}{0pt}%
\pgfpathmoveto{\pgfqpoint{1.758400in}{0.660000in}}%
\pgfpathlineto{\pgfqpoint{1.764600in}{0.660000in}}%
\pgfpathlineto{\pgfqpoint{1.764600in}{3.886236in}}%
\pgfpathlineto{\pgfqpoint{1.758400in}{3.886236in}}%
\pgfpathlineto{\pgfqpoint{1.758400in}{0.660000in}}%
\pgfpathclose%
\pgfusepath{fill}%
\end{pgfscope}%
\begin{pgfscope}%
\pgfpathrectangle{\pgfqpoint{1.250000in}{0.660000in}}{\pgfqpoint{7.750000in}{4.620000in}}%
\pgfusepath{clip}%
\pgfsetbuttcap%
\pgfsetmiterjoin%
\definecolor{currentfill}{rgb}{0.000000,0.000000,0.000000}%
\pgfsetfillcolor{currentfill}%
\pgfsetlinewidth{0.000000pt}%
\definecolor{currentstroke}{rgb}{0.000000,0.000000,0.000000}%
\pgfsetstrokecolor{currentstroke}%
\pgfsetstrokeopacity{0.000000}%
\pgfsetdash{}{0pt}%
\pgfpathmoveto{\pgfqpoint{1.766150in}{0.660000in}}%
\pgfpathlineto{\pgfqpoint{1.772350in}{0.660000in}}%
\pgfpathlineto{\pgfqpoint{1.772350in}{3.963168in}}%
\pgfpathlineto{\pgfqpoint{1.766150in}{3.963168in}}%
\pgfpathlineto{\pgfqpoint{1.766150in}{0.660000in}}%
\pgfpathclose%
\pgfusepath{fill}%
\end{pgfscope}%
\begin{pgfscope}%
\pgfpathrectangle{\pgfqpoint{1.250000in}{0.660000in}}{\pgfqpoint{7.750000in}{4.620000in}}%
\pgfusepath{clip}%
\pgfsetbuttcap%
\pgfsetmiterjoin%
\definecolor{currentfill}{rgb}{0.000000,0.000000,0.000000}%
\pgfsetfillcolor{currentfill}%
\pgfsetlinewidth{0.000000pt}%
\definecolor{currentstroke}{rgb}{0.000000,0.000000,0.000000}%
\pgfsetstrokecolor{currentstroke}%
\pgfsetstrokeopacity{0.000000}%
\pgfsetdash{}{0pt}%
\pgfpathmoveto{\pgfqpoint{1.773900in}{0.660000in}}%
\pgfpathlineto{\pgfqpoint{1.780100in}{0.660000in}}%
\pgfpathlineto{\pgfqpoint{1.780100in}{3.833258in}}%
\pgfpathlineto{\pgfqpoint{1.773900in}{3.833258in}}%
\pgfpathlineto{\pgfqpoint{1.773900in}{0.660000in}}%
\pgfpathclose%
\pgfusepath{fill}%
\end{pgfscope}%
\begin{pgfscope}%
\pgfpathrectangle{\pgfqpoint{1.250000in}{0.660000in}}{\pgfqpoint{7.750000in}{4.620000in}}%
\pgfusepath{clip}%
\pgfsetbuttcap%
\pgfsetmiterjoin%
\definecolor{currentfill}{rgb}{0.000000,0.000000,0.000000}%
\pgfsetfillcolor{currentfill}%
\pgfsetlinewidth{0.000000pt}%
\definecolor{currentstroke}{rgb}{0.000000,0.000000,0.000000}%
\pgfsetstrokecolor{currentstroke}%
\pgfsetstrokeopacity{0.000000}%
\pgfsetdash{}{0pt}%
\pgfpathmoveto{\pgfqpoint{1.781650in}{0.660000in}}%
\pgfpathlineto{\pgfqpoint{1.787850in}{0.660000in}}%
\pgfpathlineto{\pgfqpoint{1.787850in}{3.954203in}}%
\pgfpathlineto{\pgfqpoint{1.781650in}{3.954203in}}%
\pgfpathlineto{\pgfqpoint{1.781650in}{0.660000in}}%
\pgfpathclose%
\pgfusepath{fill}%
\end{pgfscope}%
\begin{pgfscope}%
\pgfpathrectangle{\pgfqpoint{1.250000in}{0.660000in}}{\pgfqpoint{7.750000in}{4.620000in}}%
\pgfusepath{clip}%
\pgfsetbuttcap%
\pgfsetmiterjoin%
\definecolor{currentfill}{rgb}{0.000000,0.000000,0.000000}%
\pgfsetfillcolor{currentfill}%
\pgfsetlinewidth{0.000000pt}%
\definecolor{currentstroke}{rgb}{0.000000,0.000000,0.000000}%
\pgfsetstrokecolor{currentstroke}%
\pgfsetstrokeopacity{0.000000}%
\pgfsetdash{}{0pt}%
\pgfpathmoveto{\pgfqpoint{1.789400in}{0.660000in}}%
\pgfpathlineto{\pgfqpoint{1.795600in}{0.660000in}}%
\pgfpathlineto{\pgfqpoint{1.795600in}{3.868576in}}%
\pgfpathlineto{\pgfqpoint{1.789400in}{3.868576in}}%
\pgfpathlineto{\pgfqpoint{1.789400in}{0.660000in}}%
\pgfpathclose%
\pgfusepath{fill}%
\end{pgfscope}%
\begin{pgfscope}%
\pgfpathrectangle{\pgfqpoint{1.250000in}{0.660000in}}{\pgfqpoint{7.750000in}{4.620000in}}%
\pgfusepath{clip}%
\pgfsetbuttcap%
\pgfsetmiterjoin%
\definecolor{currentfill}{rgb}{0.000000,0.000000,0.000000}%
\pgfsetfillcolor{currentfill}%
\pgfsetlinewidth{0.000000pt}%
\definecolor{currentstroke}{rgb}{0.000000,0.000000,0.000000}%
\pgfsetstrokecolor{currentstroke}%
\pgfsetstrokeopacity{0.000000}%
\pgfsetdash{}{0pt}%
\pgfpathmoveto{\pgfqpoint{1.797150in}{0.660000in}}%
\pgfpathlineto{\pgfqpoint{1.803350in}{0.660000in}}%
\pgfpathlineto{\pgfqpoint{1.803350in}{4.094165in}}%
\pgfpathlineto{\pgfqpoint{1.797150in}{4.094165in}}%
\pgfpathlineto{\pgfqpoint{1.797150in}{0.660000in}}%
\pgfpathclose%
\pgfusepath{fill}%
\end{pgfscope}%
\begin{pgfscope}%
\pgfpathrectangle{\pgfqpoint{1.250000in}{0.660000in}}{\pgfqpoint{7.750000in}{4.620000in}}%
\pgfusepath{clip}%
\pgfsetbuttcap%
\pgfsetmiterjoin%
\definecolor{currentfill}{rgb}{0.000000,0.000000,0.000000}%
\pgfsetfillcolor{currentfill}%
\pgfsetlinewidth{0.000000pt}%
\definecolor{currentstroke}{rgb}{0.000000,0.000000,0.000000}%
\pgfsetstrokecolor{currentstroke}%
\pgfsetstrokeopacity{0.000000}%
\pgfsetdash{}{0pt}%
\pgfpathmoveto{\pgfqpoint{1.804900in}{0.660000in}}%
\pgfpathlineto{\pgfqpoint{1.811100in}{0.660000in}}%
\pgfpathlineto{\pgfqpoint{1.811100in}{3.804161in}}%
\pgfpathlineto{\pgfqpoint{1.804900in}{3.804161in}}%
\pgfpathlineto{\pgfqpoint{1.804900in}{0.660000in}}%
\pgfpathclose%
\pgfusepath{fill}%
\end{pgfscope}%
\begin{pgfscope}%
\pgfpathrectangle{\pgfqpoint{1.250000in}{0.660000in}}{\pgfqpoint{7.750000in}{4.620000in}}%
\pgfusepath{clip}%
\pgfsetbuttcap%
\pgfsetmiterjoin%
\definecolor{currentfill}{rgb}{0.000000,0.000000,0.000000}%
\pgfsetfillcolor{currentfill}%
\pgfsetlinewidth{0.000000pt}%
\definecolor{currentstroke}{rgb}{0.000000,0.000000,0.000000}%
\pgfsetstrokecolor{currentstroke}%
\pgfsetstrokeopacity{0.000000}%
\pgfsetdash{}{0pt}%
\pgfpathmoveto{\pgfqpoint{1.812650in}{0.660000in}}%
\pgfpathlineto{\pgfqpoint{1.818850in}{0.660000in}}%
\pgfpathlineto{\pgfqpoint{1.818850in}{3.799972in}}%
\pgfpathlineto{\pgfqpoint{1.812650in}{3.799972in}}%
\pgfpathlineto{\pgfqpoint{1.812650in}{0.660000in}}%
\pgfpathclose%
\pgfusepath{fill}%
\end{pgfscope}%
\begin{pgfscope}%
\pgfpathrectangle{\pgfqpoint{1.250000in}{0.660000in}}{\pgfqpoint{7.750000in}{4.620000in}}%
\pgfusepath{clip}%
\pgfsetbuttcap%
\pgfsetmiterjoin%
\definecolor{currentfill}{rgb}{0.000000,0.000000,0.000000}%
\pgfsetfillcolor{currentfill}%
\pgfsetlinewidth{0.000000pt}%
\definecolor{currentstroke}{rgb}{0.000000,0.000000,0.000000}%
\pgfsetstrokecolor{currentstroke}%
\pgfsetstrokeopacity{0.000000}%
\pgfsetdash{}{0pt}%
\pgfpathmoveto{\pgfqpoint{1.820400in}{0.660000in}}%
\pgfpathlineto{\pgfqpoint{1.826600in}{0.660000in}}%
\pgfpathlineto{\pgfqpoint{1.826600in}{3.815848in}}%
\pgfpathlineto{\pgfqpoint{1.820400in}{3.815848in}}%
\pgfpathlineto{\pgfqpoint{1.820400in}{0.660000in}}%
\pgfpathclose%
\pgfusepath{fill}%
\end{pgfscope}%
\begin{pgfscope}%
\pgfpathrectangle{\pgfqpoint{1.250000in}{0.660000in}}{\pgfqpoint{7.750000in}{4.620000in}}%
\pgfusepath{clip}%
\pgfsetbuttcap%
\pgfsetmiterjoin%
\definecolor{currentfill}{rgb}{0.000000,0.000000,0.000000}%
\pgfsetfillcolor{currentfill}%
\pgfsetlinewidth{0.000000pt}%
\definecolor{currentstroke}{rgb}{0.000000,0.000000,0.000000}%
\pgfsetstrokecolor{currentstroke}%
\pgfsetstrokeopacity{0.000000}%
\pgfsetdash{}{0pt}%
\pgfpathmoveto{\pgfqpoint{1.828150in}{0.660000in}}%
\pgfpathlineto{\pgfqpoint{1.834350in}{0.660000in}}%
\pgfpathlineto{\pgfqpoint{1.834350in}{3.831677in}}%
\pgfpathlineto{\pgfqpoint{1.828150in}{3.831677in}}%
\pgfpathlineto{\pgfqpoint{1.828150in}{0.660000in}}%
\pgfpathclose%
\pgfusepath{fill}%
\end{pgfscope}%
\begin{pgfscope}%
\pgfpathrectangle{\pgfqpoint{1.250000in}{0.660000in}}{\pgfqpoint{7.750000in}{4.620000in}}%
\pgfusepath{clip}%
\pgfsetbuttcap%
\pgfsetmiterjoin%
\definecolor{currentfill}{rgb}{0.000000,0.000000,0.000000}%
\pgfsetfillcolor{currentfill}%
\pgfsetlinewidth{0.000000pt}%
\definecolor{currentstroke}{rgb}{0.000000,0.000000,0.000000}%
\pgfsetstrokecolor{currentstroke}%
\pgfsetstrokeopacity{0.000000}%
\pgfsetdash{}{0pt}%
\pgfpathmoveto{\pgfqpoint{1.835900in}{0.660000in}}%
\pgfpathlineto{\pgfqpoint{1.842100in}{0.660000in}}%
\pgfpathlineto{\pgfqpoint{1.842100in}{3.783630in}}%
\pgfpathlineto{\pgfqpoint{1.835900in}{3.783630in}}%
\pgfpathlineto{\pgfqpoint{1.835900in}{0.660000in}}%
\pgfpathclose%
\pgfusepath{fill}%
\end{pgfscope}%
\begin{pgfscope}%
\pgfpathrectangle{\pgfqpoint{1.250000in}{0.660000in}}{\pgfqpoint{7.750000in}{4.620000in}}%
\pgfusepath{clip}%
\pgfsetbuttcap%
\pgfsetmiterjoin%
\definecolor{currentfill}{rgb}{0.000000,0.000000,0.000000}%
\pgfsetfillcolor{currentfill}%
\pgfsetlinewidth{0.000000pt}%
\definecolor{currentstroke}{rgb}{0.000000,0.000000,0.000000}%
\pgfsetstrokecolor{currentstroke}%
\pgfsetstrokeopacity{0.000000}%
\pgfsetdash{}{0pt}%
\pgfpathmoveto{\pgfqpoint{1.843650in}{0.660000in}}%
\pgfpathlineto{\pgfqpoint{1.849850in}{0.660000in}}%
\pgfpathlineto{\pgfqpoint{1.849850in}{3.856201in}}%
\pgfpathlineto{\pgfqpoint{1.843650in}{3.856201in}}%
\pgfpathlineto{\pgfqpoint{1.843650in}{0.660000in}}%
\pgfpathclose%
\pgfusepath{fill}%
\end{pgfscope}%
\begin{pgfscope}%
\pgfpathrectangle{\pgfqpoint{1.250000in}{0.660000in}}{\pgfqpoint{7.750000in}{4.620000in}}%
\pgfusepath{clip}%
\pgfsetbuttcap%
\pgfsetmiterjoin%
\definecolor{currentfill}{rgb}{0.000000,0.000000,0.000000}%
\pgfsetfillcolor{currentfill}%
\pgfsetlinewidth{0.000000pt}%
\definecolor{currentstroke}{rgb}{0.000000,0.000000,0.000000}%
\pgfsetstrokecolor{currentstroke}%
\pgfsetstrokeopacity{0.000000}%
\pgfsetdash{}{0pt}%
\pgfpathmoveto{\pgfqpoint{1.851400in}{0.660000in}}%
\pgfpathlineto{\pgfqpoint{1.857600in}{0.660000in}}%
\pgfpathlineto{\pgfqpoint{1.857600in}{3.789733in}}%
\pgfpathlineto{\pgfqpoint{1.851400in}{3.789733in}}%
\pgfpathlineto{\pgfqpoint{1.851400in}{0.660000in}}%
\pgfpathclose%
\pgfusepath{fill}%
\end{pgfscope}%
\begin{pgfscope}%
\pgfpathrectangle{\pgfqpoint{1.250000in}{0.660000in}}{\pgfqpoint{7.750000in}{4.620000in}}%
\pgfusepath{clip}%
\pgfsetbuttcap%
\pgfsetmiterjoin%
\definecolor{currentfill}{rgb}{0.000000,0.000000,0.000000}%
\pgfsetfillcolor{currentfill}%
\pgfsetlinewidth{0.000000pt}%
\definecolor{currentstroke}{rgb}{0.000000,0.000000,0.000000}%
\pgfsetstrokecolor{currentstroke}%
\pgfsetstrokeopacity{0.000000}%
\pgfsetdash{}{0pt}%
\pgfpathmoveto{\pgfqpoint{1.859150in}{0.660000in}}%
\pgfpathlineto{\pgfqpoint{1.865350in}{0.660000in}}%
\pgfpathlineto{\pgfqpoint{1.865350in}{3.743496in}}%
\pgfpathlineto{\pgfqpoint{1.859150in}{3.743496in}}%
\pgfpathlineto{\pgfqpoint{1.859150in}{0.660000in}}%
\pgfpathclose%
\pgfusepath{fill}%
\end{pgfscope}%
\begin{pgfscope}%
\pgfpathrectangle{\pgfqpoint{1.250000in}{0.660000in}}{\pgfqpoint{7.750000in}{4.620000in}}%
\pgfusepath{clip}%
\pgfsetbuttcap%
\pgfsetmiterjoin%
\definecolor{currentfill}{rgb}{0.000000,0.000000,0.000000}%
\pgfsetfillcolor{currentfill}%
\pgfsetlinewidth{0.000000pt}%
\definecolor{currentstroke}{rgb}{0.000000,0.000000,0.000000}%
\pgfsetstrokecolor{currentstroke}%
\pgfsetstrokeopacity{0.000000}%
\pgfsetdash{}{0pt}%
\pgfpathmoveto{\pgfqpoint{1.866900in}{0.660000in}}%
\pgfpathlineto{\pgfqpoint{1.873100in}{0.660000in}}%
\pgfpathlineto{\pgfqpoint{1.873100in}{3.833551in}}%
\pgfpathlineto{\pgfqpoint{1.866900in}{3.833551in}}%
\pgfpathlineto{\pgfqpoint{1.866900in}{0.660000in}}%
\pgfpathclose%
\pgfusepath{fill}%
\end{pgfscope}%
\begin{pgfscope}%
\pgfpathrectangle{\pgfqpoint{1.250000in}{0.660000in}}{\pgfqpoint{7.750000in}{4.620000in}}%
\pgfusepath{clip}%
\pgfsetbuttcap%
\pgfsetmiterjoin%
\definecolor{currentfill}{rgb}{0.000000,0.000000,0.000000}%
\pgfsetfillcolor{currentfill}%
\pgfsetlinewidth{0.000000pt}%
\definecolor{currentstroke}{rgb}{0.000000,0.000000,0.000000}%
\pgfsetstrokecolor{currentstroke}%
\pgfsetstrokeopacity{0.000000}%
\pgfsetdash{}{0pt}%
\pgfpathmoveto{\pgfqpoint{1.874650in}{0.660000in}}%
\pgfpathlineto{\pgfqpoint{1.880850in}{0.660000in}}%
\pgfpathlineto{\pgfqpoint{1.880850in}{3.772641in}}%
\pgfpathlineto{\pgfqpoint{1.874650in}{3.772641in}}%
\pgfpathlineto{\pgfqpoint{1.874650in}{0.660000in}}%
\pgfpathclose%
\pgfusepath{fill}%
\end{pgfscope}%
\begin{pgfscope}%
\pgfpathrectangle{\pgfqpoint{1.250000in}{0.660000in}}{\pgfqpoint{7.750000in}{4.620000in}}%
\pgfusepath{clip}%
\pgfsetbuttcap%
\pgfsetmiterjoin%
\definecolor{currentfill}{rgb}{0.000000,0.000000,0.000000}%
\pgfsetfillcolor{currentfill}%
\pgfsetlinewidth{0.000000pt}%
\definecolor{currentstroke}{rgb}{0.000000,0.000000,0.000000}%
\pgfsetstrokecolor{currentstroke}%
\pgfsetstrokeopacity{0.000000}%
\pgfsetdash{}{0pt}%
\pgfpathmoveto{\pgfqpoint{1.882400in}{0.660000in}}%
\pgfpathlineto{\pgfqpoint{1.888600in}{0.660000in}}%
\pgfpathlineto{\pgfqpoint{1.888600in}{3.860097in}}%
\pgfpathlineto{\pgfqpoint{1.882400in}{3.860097in}}%
\pgfpathlineto{\pgfqpoint{1.882400in}{0.660000in}}%
\pgfpathclose%
\pgfusepath{fill}%
\end{pgfscope}%
\begin{pgfscope}%
\pgfpathrectangle{\pgfqpoint{1.250000in}{0.660000in}}{\pgfqpoint{7.750000in}{4.620000in}}%
\pgfusepath{clip}%
\pgfsetbuttcap%
\pgfsetmiterjoin%
\definecolor{currentfill}{rgb}{0.000000,0.000000,0.000000}%
\pgfsetfillcolor{currentfill}%
\pgfsetlinewidth{0.000000pt}%
\definecolor{currentstroke}{rgb}{0.000000,0.000000,0.000000}%
\pgfsetstrokecolor{currentstroke}%
\pgfsetstrokeopacity{0.000000}%
\pgfsetdash{}{0pt}%
\pgfpathmoveto{\pgfqpoint{1.890150in}{0.660000in}}%
\pgfpathlineto{\pgfqpoint{1.896350in}{0.660000in}}%
\pgfpathlineto{\pgfqpoint{1.896350in}{3.786591in}}%
\pgfpathlineto{\pgfqpoint{1.890150in}{3.786591in}}%
\pgfpathlineto{\pgfqpoint{1.890150in}{0.660000in}}%
\pgfpathclose%
\pgfusepath{fill}%
\end{pgfscope}%
\begin{pgfscope}%
\pgfpathrectangle{\pgfqpoint{1.250000in}{0.660000in}}{\pgfqpoint{7.750000in}{4.620000in}}%
\pgfusepath{clip}%
\pgfsetbuttcap%
\pgfsetmiterjoin%
\definecolor{currentfill}{rgb}{0.000000,0.000000,0.000000}%
\pgfsetfillcolor{currentfill}%
\pgfsetlinewidth{0.000000pt}%
\definecolor{currentstroke}{rgb}{0.000000,0.000000,0.000000}%
\pgfsetstrokecolor{currentstroke}%
\pgfsetstrokeopacity{0.000000}%
\pgfsetdash{}{0pt}%
\pgfpathmoveto{\pgfqpoint{1.897900in}{0.660000in}}%
\pgfpathlineto{\pgfqpoint{1.904100in}{0.660000in}}%
\pgfpathlineto{\pgfqpoint{1.904100in}{3.797703in}}%
\pgfpathlineto{\pgfqpoint{1.897900in}{3.797703in}}%
\pgfpathlineto{\pgfqpoint{1.897900in}{0.660000in}}%
\pgfpathclose%
\pgfusepath{fill}%
\end{pgfscope}%
\begin{pgfscope}%
\pgfpathrectangle{\pgfqpoint{1.250000in}{0.660000in}}{\pgfqpoint{7.750000in}{4.620000in}}%
\pgfusepath{clip}%
\pgfsetbuttcap%
\pgfsetmiterjoin%
\definecolor{currentfill}{rgb}{0.000000,0.000000,0.000000}%
\pgfsetfillcolor{currentfill}%
\pgfsetlinewidth{0.000000pt}%
\definecolor{currentstroke}{rgb}{0.000000,0.000000,0.000000}%
\pgfsetstrokecolor{currentstroke}%
\pgfsetstrokeopacity{0.000000}%
\pgfsetdash{}{0pt}%
\pgfpathmoveto{\pgfqpoint{1.905650in}{0.660000in}}%
\pgfpathlineto{\pgfqpoint{1.911850in}{0.660000in}}%
\pgfpathlineto{\pgfqpoint{1.911850in}{3.892226in}}%
\pgfpathlineto{\pgfqpoint{1.905650in}{3.892226in}}%
\pgfpathlineto{\pgfqpoint{1.905650in}{0.660000in}}%
\pgfpathclose%
\pgfusepath{fill}%
\end{pgfscope}%
\begin{pgfscope}%
\pgfpathrectangle{\pgfqpoint{1.250000in}{0.660000in}}{\pgfqpoint{7.750000in}{4.620000in}}%
\pgfusepath{clip}%
\pgfsetbuttcap%
\pgfsetmiterjoin%
\definecolor{currentfill}{rgb}{0.000000,0.000000,0.000000}%
\pgfsetfillcolor{currentfill}%
\pgfsetlinewidth{0.000000pt}%
\definecolor{currentstroke}{rgb}{0.000000,0.000000,0.000000}%
\pgfsetstrokecolor{currentstroke}%
\pgfsetstrokeopacity{0.000000}%
\pgfsetdash{}{0pt}%
\pgfpathmoveto{\pgfqpoint{1.913400in}{0.660000in}}%
\pgfpathlineto{\pgfqpoint{1.919600in}{0.660000in}}%
\pgfpathlineto{\pgfqpoint{1.919600in}{3.787664in}}%
\pgfpathlineto{\pgfqpoint{1.913400in}{3.787664in}}%
\pgfpathlineto{\pgfqpoint{1.913400in}{0.660000in}}%
\pgfpathclose%
\pgfusepath{fill}%
\end{pgfscope}%
\begin{pgfscope}%
\pgfpathrectangle{\pgfqpoint{1.250000in}{0.660000in}}{\pgfqpoint{7.750000in}{4.620000in}}%
\pgfusepath{clip}%
\pgfsetbuttcap%
\pgfsetmiterjoin%
\definecolor{currentfill}{rgb}{0.000000,0.000000,0.000000}%
\pgfsetfillcolor{currentfill}%
\pgfsetlinewidth{0.000000pt}%
\definecolor{currentstroke}{rgb}{0.000000,0.000000,0.000000}%
\pgfsetstrokecolor{currentstroke}%
\pgfsetstrokeopacity{0.000000}%
\pgfsetdash{}{0pt}%
\pgfpathmoveto{\pgfqpoint{1.921150in}{0.660000in}}%
\pgfpathlineto{\pgfqpoint{1.927350in}{0.660000in}}%
\pgfpathlineto{\pgfqpoint{1.927350in}{3.695398in}}%
\pgfpathlineto{\pgfqpoint{1.921150in}{3.695398in}}%
\pgfpathlineto{\pgfqpoint{1.921150in}{0.660000in}}%
\pgfpathclose%
\pgfusepath{fill}%
\end{pgfscope}%
\begin{pgfscope}%
\pgfpathrectangle{\pgfqpoint{1.250000in}{0.660000in}}{\pgfqpoint{7.750000in}{4.620000in}}%
\pgfusepath{clip}%
\pgfsetbuttcap%
\pgfsetmiterjoin%
\definecolor{currentfill}{rgb}{0.000000,0.000000,0.000000}%
\pgfsetfillcolor{currentfill}%
\pgfsetlinewidth{0.000000pt}%
\definecolor{currentstroke}{rgb}{0.000000,0.000000,0.000000}%
\pgfsetstrokecolor{currentstroke}%
\pgfsetstrokeopacity{0.000000}%
\pgfsetdash{}{0pt}%
\pgfpathmoveto{\pgfqpoint{1.928900in}{0.660000in}}%
\pgfpathlineto{\pgfqpoint{1.935100in}{0.660000in}}%
\pgfpathlineto{\pgfqpoint{1.935100in}{3.785618in}}%
\pgfpathlineto{\pgfqpoint{1.928900in}{3.785618in}}%
\pgfpathlineto{\pgfqpoint{1.928900in}{0.660000in}}%
\pgfpathclose%
\pgfusepath{fill}%
\end{pgfscope}%
\begin{pgfscope}%
\pgfpathrectangle{\pgfqpoint{1.250000in}{0.660000in}}{\pgfqpoint{7.750000in}{4.620000in}}%
\pgfusepath{clip}%
\pgfsetbuttcap%
\pgfsetmiterjoin%
\definecolor{currentfill}{rgb}{0.000000,0.000000,0.000000}%
\pgfsetfillcolor{currentfill}%
\pgfsetlinewidth{0.000000pt}%
\definecolor{currentstroke}{rgb}{0.000000,0.000000,0.000000}%
\pgfsetstrokecolor{currentstroke}%
\pgfsetstrokeopacity{0.000000}%
\pgfsetdash{}{0pt}%
\pgfpathmoveto{\pgfqpoint{1.936650in}{0.660000in}}%
\pgfpathlineto{\pgfqpoint{1.942850in}{0.660000in}}%
\pgfpathlineto{\pgfqpoint{1.942850in}{3.738630in}}%
\pgfpathlineto{\pgfqpoint{1.936650in}{3.738630in}}%
\pgfpathlineto{\pgfqpoint{1.936650in}{0.660000in}}%
\pgfpathclose%
\pgfusepath{fill}%
\end{pgfscope}%
\begin{pgfscope}%
\pgfpathrectangle{\pgfqpoint{1.250000in}{0.660000in}}{\pgfqpoint{7.750000in}{4.620000in}}%
\pgfusepath{clip}%
\pgfsetbuttcap%
\pgfsetmiterjoin%
\definecolor{currentfill}{rgb}{0.000000,0.000000,0.000000}%
\pgfsetfillcolor{currentfill}%
\pgfsetlinewidth{0.000000pt}%
\definecolor{currentstroke}{rgb}{0.000000,0.000000,0.000000}%
\pgfsetstrokecolor{currentstroke}%
\pgfsetstrokeopacity{0.000000}%
\pgfsetdash{}{0pt}%
\pgfpathmoveto{\pgfqpoint{1.944400in}{0.660000in}}%
\pgfpathlineto{\pgfqpoint{1.950600in}{0.660000in}}%
\pgfpathlineto{\pgfqpoint{1.950600in}{3.705975in}}%
\pgfpathlineto{\pgfqpoint{1.944400in}{3.705975in}}%
\pgfpathlineto{\pgfqpoint{1.944400in}{0.660000in}}%
\pgfpathclose%
\pgfusepath{fill}%
\end{pgfscope}%
\begin{pgfscope}%
\pgfpathrectangle{\pgfqpoint{1.250000in}{0.660000in}}{\pgfqpoint{7.750000in}{4.620000in}}%
\pgfusepath{clip}%
\pgfsetbuttcap%
\pgfsetmiterjoin%
\definecolor{currentfill}{rgb}{0.000000,0.000000,0.000000}%
\pgfsetfillcolor{currentfill}%
\pgfsetlinewidth{0.000000pt}%
\definecolor{currentstroke}{rgb}{0.000000,0.000000,0.000000}%
\pgfsetstrokecolor{currentstroke}%
\pgfsetstrokeopacity{0.000000}%
\pgfsetdash{}{0pt}%
\pgfpathmoveto{\pgfqpoint{1.952150in}{0.660000in}}%
\pgfpathlineto{\pgfqpoint{1.958350in}{0.660000in}}%
\pgfpathlineto{\pgfqpoint{1.958350in}{3.766141in}}%
\pgfpathlineto{\pgfqpoint{1.952150in}{3.766141in}}%
\pgfpathlineto{\pgfqpoint{1.952150in}{0.660000in}}%
\pgfpathclose%
\pgfusepath{fill}%
\end{pgfscope}%
\begin{pgfscope}%
\pgfpathrectangle{\pgfqpoint{1.250000in}{0.660000in}}{\pgfqpoint{7.750000in}{4.620000in}}%
\pgfusepath{clip}%
\pgfsetbuttcap%
\pgfsetmiterjoin%
\definecolor{currentfill}{rgb}{0.000000,0.000000,0.000000}%
\pgfsetfillcolor{currentfill}%
\pgfsetlinewidth{0.000000pt}%
\definecolor{currentstroke}{rgb}{0.000000,0.000000,0.000000}%
\pgfsetstrokecolor{currentstroke}%
\pgfsetstrokeopacity{0.000000}%
\pgfsetdash{}{0pt}%
\pgfpathmoveto{\pgfqpoint{1.959900in}{0.660000in}}%
\pgfpathlineto{\pgfqpoint{1.966100in}{0.660000in}}%
\pgfpathlineto{\pgfqpoint{1.966100in}{3.796584in}}%
\pgfpathlineto{\pgfqpoint{1.959900in}{3.796584in}}%
\pgfpathlineto{\pgfqpoint{1.959900in}{0.660000in}}%
\pgfpathclose%
\pgfusepath{fill}%
\end{pgfscope}%
\begin{pgfscope}%
\pgfpathrectangle{\pgfqpoint{1.250000in}{0.660000in}}{\pgfqpoint{7.750000in}{4.620000in}}%
\pgfusepath{clip}%
\pgfsetbuttcap%
\pgfsetmiterjoin%
\definecolor{currentfill}{rgb}{0.000000,0.000000,0.000000}%
\pgfsetfillcolor{currentfill}%
\pgfsetlinewidth{0.000000pt}%
\definecolor{currentstroke}{rgb}{0.000000,0.000000,0.000000}%
\pgfsetstrokecolor{currentstroke}%
\pgfsetstrokeopacity{0.000000}%
\pgfsetdash{}{0pt}%
\pgfpathmoveto{\pgfqpoint{1.967650in}{0.660000in}}%
\pgfpathlineto{\pgfqpoint{1.973850in}{0.660000in}}%
\pgfpathlineto{\pgfqpoint{1.973850in}{3.730229in}}%
\pgfpathlineto{\pgfqpoint{1.967650in}{3.730229in}}%
\pgfpathlineto{\pgfqpoint{1.967650in}{0.660000in}}%
\pgfpathclose%
\pgfusepath{fill}%
\end{pgfscope}%
\begin{pgfscope}%
\pgfpathrectangle{\pgfqpoint{1.250000in}{0.660000in}}{\pgfqpoint{7.750000in}{4.620000in}}%
\pgfusepath{clip}%
\pgfsetbuttcap%
\pgfsetmiterjoin%
\definecolor{currentfill}{rgb}{0.000000,0.000000,0.000000}%
\pgfsetfillcolor{currentfill}%
\pgfsetlinewidth{0.000000pt}%
\definecolor{currentstroke}{rgb}{0.000000,0.000000,0.000000}%
\pgfsetstrokecolor{currentstroke}%
\pgfsetstrokeopacity{0.000000}%
\pgfsetdash{}{0pt}%
\pgfpathmoveto{\pgfqpoint{1.975400in}{0.660000in}}%
\pgfpathlineto{\pgfqpoint{1.981600in}{0.660000in}}%
\pgfpathlineto{\pgfqpoint{1.981600in}{3.679556in}}%
\pgfpathlineto{\pgfqpoint{1.975400in}{3.679556in}}%
\pgfpathlineto{\pgfqpoint{1.975400in}{0.660000in}}%
\pgfpathclose%
\pgfusepath{fill}%
\end{pgfscope}%
\begin{pgfscope}%
\pgfpathrectangle{\pgfqpoint{1.250000in}{0.660000in}}{\pgfqpoint{7.750000in}{4.620000in}}%
\pgfusepath{clip}%
\pgfsetbuttcap%
\pgfsetmiterjoin%
\definecolor{currentfill}{rgb}{0.000000,0.000000,0.000000}%
\pgfsetfillcolor{currentfill}%
\pgfsetlinewidth{0.000000pt}%
\definecolor{currentstroke}{rgb}{0.000000,0.000000,0.000000}%
\pgfsetstrokecolor{currentstroke}%
\pgfsetstrokeopacity{0.000000}%
\pgfsetdash{}{0pt}%
\pgfpathmoveto{\pgfqpoint{1.983150in}{0.660000in}}%
\pgfpathlineto{\pgfqpoint{1.989350in}{0.660000in}}%
\pgfpathlineto{\pgfqpoint{1.989350in}{3.744601in}}%
\pgfpathlineto{\pgfqpoint{1.983150in}{3.744601in}}%
\pgfpathlineto{\pgfqpoint{1.983150in}{0.660000in}}%
\pgfpathclose%
\pgfusepath{fill}%
\end{pgfscope}%
\begin{pgfscope}%
\pgfpathrectangle{\pgfqpoint{1.250000in}{0.660000in}}{\pgfqpoint{7.750000in}{4.620000in}}%
\pgfusepath{clip}%
\pgfsetbuttcap%
\pgfsetmiterjoin%
\definecolor{currentfill}{rgb}{0.000000,0.000000,0.000000}%
\pgfsetfillcolor{currentfill}%
\pgfsetlinewidth{0.000000pt}%
\definecolor{currentstroke}{rgb}{0.000000,0.000000,0.000000}%
\pgfsetstrokecolor{currentstroke}%
\pgfsetstrokeopacity{0.000000}%
\pgfsetdash{}{0pt}%
\pgfpathmoveto{\pgfqpoint{1.990900in}{0.660000in}}%
\pgfpathlineto{\pgfqpoint{1.997100in}{0.660000in}}%
\pgfpathlineto{\pgfqpoint{1.997100in}{3.692961in}}%
\pgfpathlineto{\pgfqpoint{1.990900in}{3.692961in}}%
\pgfpathlineto{\pgfqpoint{1.990900in}{0.660000in}}%
\pgfpathclose%
\pgfusepath{fill}%
\end{pgfscope}%
\begin{pgfscope}%
\pgfpathrectangle{\pgfqpoint{1.250000in}{0.660000in}}{\pgfqpoint{7.750000in}{4.620000in}}%
\pgfusepath{clip}%
\pgfsetbuttcap%
\pgfsetmiterjoin%
\definecolor{currentfill}{rgb}{0.000000,0.000000,0.000000}%
\pgfsetfillcolor{currentfill}%
\pgfsetlinewidth{0.000000pt}%
\definecolor{currentstroke}{rgb}{0.000000,0.000000,0.000000}%
\pgfsetstrokecolor{currentstroke}%
\pgfsetstrokeopacity{0.000000}%
\pgfsetdash{}{0pt}%
\pgfpathmoveto{\pgfqpoint{1.998650in}{0.660000in}}%
\pgfpathlineto{\pgfqpoint{2.004850in}{0.660000in}}%
\pgfpathlineto{\pgfqpoint{2.004850in}{3.737331in}}%
\pgfpathlineto{\pgfqpoint{1.998650in}{3.737331in}}%
\pgfpathlineto{\pgfqpoint{1.998650in}{0.660000in}}%
\pgfpathclose%
\pgfusepath{fill}%
\end{pgfscope}%
\begin{pgfscope}%
\pgfpathrectangle{\pgfqpoint{1.250000in}{0.660000in}}{\pgfqpoint{7.750000in}{4.620000in}}%
\pgfusepath{clip}%
\pgfsetbuttcap%
\pgfsetmiterjoin%
\definecolor{currentfill}{rgb}{0.000000,0.000000,0.000000}%
\pgfsetfillcolor{currentfill}%
\pgfsetlinewidth{0.000000pt}%
\definecolor{currentstroke}{rgb}{0.000000,0.000000,0.000000}%
\pgfsetstrokecolor{currentstroke}%
\pgfsetstrokeopacity{0.000000}%
\pgfsetdash{}{0pt}%
\pgfpathmoveto{\pgfqpoint{2.006400in}{0.660000in}}%
\pgfpathlineto{\pgfqpoint{2.012600in}{0.660000in}}%
\pgfpathlineto{\pgfqpoint{2.012600in}{3.687228in}}%
\pgfpathlineto{\pgfqpoint{2.006400in}{3.687228in}}%
\pgfpathlineto{\pgfqpoint{2.006400in}{0.660000in}}%
\pgfpathclose%
\pgfusepath{fill}%
\end{pgfscope}%
\begin{pgfscope}%
\pgfpathrectangle{\pgfqpoint{1.250000in}{0.660000in}}{\pgfqpoint{7.750000in}{4.620000in}}%
\pgfusepath{clip}%
\pgfsetbuttcap%
\pgfsetmiterjoin%
\definecolor{currentfill}{rgb}{0.000000,0.000000,0.000000}%
\pgfsetfillcolor{currentfill}%
\pgfsetlinewidth{0.000000pt}%
\definecolor{currentstroke}{rgb}{0.000000,0.000000,0.000000}%
\pgfsetstrokecolor{currentstroke}%
\pgfsetstrokeopacity{0.000000}%
\pgfsetdash{}{0pt}%
\pgfpathmoveto{\pgfqpoint{2.014150in}{0.660000in}}%
\pgfpathlineto{\pgfqpoint{2.020350in}{0.660000in}}%
\pgfpathlineto{\pgfqpoint{2.020350in}{3.740223in}}%
\pgfpathlineto{\pgfqpoint{2.014150in}{3.740223in}}%
\pgfpathlineto{\pgfqpoint{2.014150in}{0.660000in}}%
\pgfpathclose%
\pgfusepath{fill}%
\end{pgfscope}%
\begin{pgfscope}%
\pgfpathrectangle{\pgfqpoint{1.250000in}{0.660000in}}{\pgfqpoint{7.750000in}{4.620000in}}%
\pgfusepath{clip}%
\pgfsetbuttcap%
\pgfsetmiterjoin%
\definecolor{currentfill}{rgb}{0.000000,0.000000,0.000000}%
\pgfsetfillcolor{currentfill}%
\pgfsetlinewidth{0.000000pt}%
\definecolor{currentstroke}{rgb}{0.000000,0.000000,0.000000}%
\pgfsetstrokecolor{currentstroke}%
\pgfsetstrokeopacity{0.000000}%
\pgfsetdash{}{0pt}%
\pgfpathmoveto{\pgfqpoint{2.021900in}{0.660000in}}%
\pgfpathlineto{\pgfqpoint{2.028100in}{0.660000in}}%
\pgfpathlineto{\pgfqpoint{2.028100in}{3.675601in}}%
\pgfpathlineto{\pgfqpoint{2.021900in}{3.675601in}}%
\pgfpathlineto{\pgfqpoint{2.021900in}{0.660000in}}%
\pgfpathclose%
\pgfusepath{fill}%
\end{pgfscope}%
\begin{pgfscope}%
\pgfpathrectangle{\pgfqpoint{1.250000in}{0.660000in}}{\pgfqpoint{7.750000in}{4.620000in}}%
\pgfusepath{clip}%
\pgfsetbuttcap%
\pgfsetmiterjoin%
\definecolor{currentfill}{rgb}{0.000000,0.000000,0.000000}%
\pgfsetfillcolor{currentfill}%
\pgfsetlinewidth{0.000000pt}%
\definecolor{currentstroke}{rgb}{0.000000,0.000000,0.000000}%
\pgfsetstrokecolor{currentstroke}%
\pgfsetstrokeopacity{0.000000}%
\pgfsetdash{}{0pt}%
\pgfpathmoveto{\pgfqpoint{2.029650in}{0.660000in}}%
\pgfpathlineto{\pgfqpoint{2.035850in}{0.660000in}}%
\pgfpathlineto{\pgfqpoint{2.035850in}{3.824825in}}%
\pgfpathlineto{\pgfqpoint{2.029650in}{3.824825in}}%
\pgfpathlineto{\pgfqpoint{2.029650in}{0.660000in}}%
\pgfpathclose%
\pgfusepath{fill}%
\end{pgfscope}%
\begin{pgfscope}%
\pgfpathrectangle{\pgfqpoint{1.250000in}{0.660000in}}{\pgfqpoint{7.750000in}{4.620000in}}%
\pgfusepath{clip}%
\pgfsetbuttcap%
\pgfsetmiterjoin%
\definecolor{currentfill}{rgb}{0.000000,0.000000,0.000000}%
\pgfsetfillcolor{currentfill}%
\pgfsetlinewidth{0.000000pt}%
\definecolor{currentstroke}{rgb}{0.000000,0.000000,0.000000}%
\pgfsetstrokecolor{currentstroke}%
\pgfsetstrokeopacity{0.000000}%
\pgfsetdash{}{0pt}%
\pgfpathmoveto{\pgfqpoint{2.037400in}{0.660000in}}%
\pgfpathlineto{\pgfqpoint{2.043600in}{0.660000in}}%
\pgfpathlineto{\pgfqpoint{2.043600in}{3.658500in}}%
\pgfpathlineto{\pgfqpoint{2.037400in}{3.658500in}}%
\pgfpathlineto{\pgfqpoint{2.037400in}{0.660000in}}%
\pgfpathclose%
\pgfusepath{fill}%
\end{pgfscope}%
\begin{pgfscope}%
\pgfpathrectangle{\pgfqpoint{1.250000in}{0.660000in}}{\pgfqpoint{7.750000in}{4.620000in}}%
\pgfusepath{clip}%
\pgfsetbuttcap%
\pgfsetmiterjoin%
\definecolor{currentfill}{rgb}{0.000000,0.000000,0.000000}%
\pgfsetfillcolor{currentfill}%
\pgfsetlinewidth{0.000000pt}%
\definecolor{currentstroke}{rgb}{0.000000,0.000000,0.000000}%
\pgfsetstrokecolor{currentstroke}%
\pgfsetstrokeopacity{0.000000}%
\pgfsetdash{}{0pt}%
\pgfpathmoveto{\pgfqpoint{2.045150in}{0.660000in}}%
\pgfpathlineto{\pgfqpoint{2.051350in}{0.660000in}}%
\pgfpathlineto{\pgfqpoint{2.051350in}{3.713304in}}%
\pgfpathlineto{\pgfqpoint{2.045150in}{3.713304in}}%
\pgfpathlineto{\pgfqpoint{2.045150in}{0.660000in}}%
\pgfpathclose%
\pgfusepath{fill}%
\end{pgfscope}%
\begin{pgfscope}%
\pgfpathrectangle{\pgfqpoint{1.250000in}{0.660000in}}{\pgfqpoint{7.750000in}{4.620000in}}%
\pgfusepath{clip}%
\pgfsetbuttcap%
\pgfsetmiterjoin%
\definecolor{currentfill}{rgb}{0.000000,0.000000,0.000000}%
\pgfsetfillcolor{currentfill}%
\pgfsetlinewidth{0.000000pt}%
\definecolor{currentstroke}{rgb}{0.000000,0.000000,0.000000}%
\pgfsetstrokecolor{currentstroke}%
\pgfsetstrokeopacity{0.000000}%
\pgfsetdash{}{0pt}%
\pgfpathmoveto{\pgfqpoint{2.052900in}{0.660000in}}%
\pgfpathlineto{\pgfqpoint{2.059100in}{0.660000in}}%
\pgfpathlineto{\pgfqpoint{2.059100in}{3.630583in}}%
\pgfpathlineto{\pgfqpoint{2.052900in}{3.630583in}}%
\pgfpathlineto{\pgfqpoint{2.052900in}{0.660000in}}%
\pgfpathclose%
\pgfusepath{fill}%
\end{pgfscope}%
\begin{pgfscope}%
\pgfpathrectangle{\pgfqpoint{1.250000in}{0.660000in}}{\pgfqpoint{7.750000in}{4.620000in}}%
\pgfusepath{clip}%
\pgfsetbuttcap%
\pgfsetmiterjoin%
\definecolor{currentfill}{rgb}{0.000000,0.000000,0.000000}%
\pgfsetfillcolor{currentfill}%
\pgfsetlinewidth{0.000000pt}%
\definecolor{currentstroke}{rgb}{0.000000,0.000000,0.000000}%
\pgfsetstrokecolor{currentstroke}%
\pgfsetstrokeopacity{0.000000}%
\pgfsetdash{}{0pt}%
\pgfpathmoveto{\pgfqpoint{2.060650in}{0.660000in}}%
\pgfpathlineto{\pgfqpoint{2.066850in}{0.660000in}}%
\pgfpathlineto{\pgfqpoint{2.066850in}{3.689151in}}%
\pgfpathlineto{\pgfqpoint{2.060650in}{3.689151in}}%
\pgfpathlineto{\pgfqpoint{2.060650in}{0.660000in}}%
\pgfpathclose%
\pgfusepath{fill}%
\end{pgfscope}%
\begin{pgfscope}%
\pgfpathrectangle{\pgfqpoint{1.250000in}{0.660000in}}{\pgfqpoint{7.750000in}{4.620000in}}%
\pgfusepath{clip}%
\pgfsetbuttcap%
\pgfsetmiterjoin%
\definecolor{currentfill}{rgb}{0.000000,0.000000,0.000000}%
\pgfsetfillcolor{currentfill}%
\pgfsetlinewidth{0.000000pt}%
\definecolor{currentstroke}{rgb}{0.000000,0.000000,0.000000}%
\pgfsetstrokecolor{currentstroke}%
\pgfsetstrokeopacity{0.000000}%
\pgfsetdash{}{0pt}%
\pgfpathmoveto{\pgfqpoint{2.068400in}{0.660000in}}%
\pgfpathlineto{\pgfqpoint{2.074600in}{0.660000in}}%
\pgfpathlineto{\pgfqpoint{2.074600in}{3.679767in}}%
\pgfpathlineto{\pgfqpoint{2.068400in}{3.679767in}}%
\pgfpathlineto{\pgfqpoint{2.068400in}{0.660000in}}%
\pgfpathclose%
\pgfusepath{fill}%
\end{pgfscope}%
\begin{pgfscope}%
\pgfpathrectangle{\pgfqpoint{1.250000in}{0.660000in}}{\pgfqpoint{7.750000in}{4.620000in}}%
\pgfusepath{clip}%
\pgfsetbuttcap%
\pgfsetmiterjoin%
\definecolor{currentfill}{rgb}{0.000000,0.000000,0.000000}%
\pgfsetfillcolor{currentfill}%
\pgfsetlinewidth{0.000000pt}%
\definecolor{currentstroke}{rgb}{0.000000,0.000000,0.000000}%
\pgfsetstrokecolor{currentstroke}%
\pgfsetstrokeopacity{0.000000}%
\pgfsetdash{}{0pt}%
\pgfpathmoveto{\pgfqpoint{2.076150in}{0.660000in}}%
\pgfpathlineto{\pgfqpoint{2.082350in}{0.660000in}}%
\pgfpathlineto{\pgfqpoint{2.082350in}{3.635109in}}%
\pgfpathlineto{\pgfqpoint{2.076150in}{3.635109in}}%
\pgfpathlineto{\pgfqpoint{2.076150in}{0.660000in}}%
\pgfpathclose%
\pgfusepath{fill}%
\end{pgfscope}%
\begin{pgfscope}%
\pgfpathrectangle{\pgfqpoint{1.250000in}{0.660000in}}{\pgfqpoint{7.750000in}{4.620000in}}%
\pgfusepath{clip}%
\pgfsetbuttcap%
\pgfsetmiterjoin%
\definecolor{currentfill}{rgb}{0.000000,0.000000,0.000000}%
\pgfsetfillcolor{currentfill}%
\pgfsetlinewidth{0.000000pt}%
\definecolor{currentstroke}{rgb}{0.000000,0.000000,0.000000}%
\pgfsetstrokecolor{currentstroke}%
\pgfsetstrokeopacity{0.000000}%
\pgfsetdash{}{0pt}%
\pgfpathmoveto{\pgfqpoint{2.083900in}{0.660000in}}%
\pgfpathlineto{\pgfqpoint{2.090100in}{0.660000in}}%
\pgfpathlineto{\pgfqpoint{2.090100in}{3.659167in}}%
\pgfpathlineto{\pgfqpoint{2.083900in}{3.659167in}}%
\pgfpathlineto{\pgfqpoint{2.083900in}{0.660000in}}%
\pgfpathclose%
\pgfusepath{fill}%
\end{pgfscope}%
\begin{pgfscope}%
\pgfpathrectangle{\pgfqpoint{1.250000in}{0.660000in}}{\pgfqpoint{7.750000in}{4.620000in}}%
\pgfusepath{clip}%
\pgfsetbuttcap%
\pgfsetmiterjoin%
\definecolor{currentfill}{rgb}{0.000000,0.000000,0.000000}%
\pgfsetfillcolor{currentfill}%
\pgfsetlinewidth{0.000000pt}%
\definecolor{currentstroke}{rgb}{0.000000,0.000000,0.000000}%
\pgfsetstrokecolor{currentstroke}%
\pgfsetstrokeopacity{0.000000}%
\pgfsetdash{}{0pt}%
\pgfpathmoveto{\pgfqpoint{2.091650in}{0.660000in}}%
\pgfpathlineto{\pgfqpoint{2.097850in}{0.660000in}}%
\pgfpathlineto{\pgfqpoint{2.097850in}{3.667378in}}%
\pgfpathlineto{\pgfqpoint{2.091650in}{3.667378in}}%
\pgfpathlineto{\pgfqpoint{2.091650in}{0.660000in}}%
\pgfpathclose%
\pgfusepath{fill}%
\end{pgfscope}%
\begin{pgfscope}%
\pgfpathrectangle{\pgfqpoint{1.250000in}{0.660000in}}{\pgfqpoint{7.750000in}{4.620000in}}%
\pgfusepath{clip}%
\pgfsetbuttcap%
\pgfsetmiterjoin%
\definecolor{currentfill}{rgb}{0.000000,0.000000,0.000000}%
\pgfsetfillcolor{currentfill}%
\pgfsetlinewidth{0.000000pt}%
\definecolor{currentstroke}{rgb}{0.000000,0.000000,0.000000}%
\pgfsetstrokecolor{currentstroke}%
\pgfsetstrokeopacity{0.000000}%
\pgfsetdash{}{0pt}%
\pgfpathmoveto{\pgfqpoint{2.099400in}{0.660000in}}%
\pgfpathlineto{\pgfqpoint{2.105600in}{0.660000in}}%
\pgfpathlineto{\pgfqpoint{2.105600in}{3.647591in}}%
\pgfpathlineto{\pgfqpoint{2.099400in}{3.647591in}}%
\pgfpathlineto{\pgfqpoint{2.099400in}{0.660000in}}%
\pgfpathclose%
\pgfusepath{fill}%
\end{pgfscope}%
\begin{pgfscope}%
\pgfpathrectangle{\pgfqpoint{1.250000in}{0.660000in}}{\pgfqpoint{7.750000in}{4.620000in}}%
\pgfusepath{clip}%
\pgfsetbuttcap%
\pgfsetmiterjoin%
\definecolor{currentfill}{rgb}{0.000000,0.000000,0.000000}%
\pgfsetfillcolor{currentfill}%
\pgfsetlinewidth{0.000000pt}%
\definecolor{currentstroke}{rgb}{0.000000,0.000000,0.000000}%
\pgfsetstrokecolor{currentstroke}%
\pgfsetstrokeopacity{0.000000}%
\pgfsetdash{}{0pt}%
\pgfpathmoveto{\pgfqpoint{2.107150in}{0.660000in}}%
\pgfpathlineto{\pgfqpoint{2.113350in}{0.660000in}}%
\pgfpathlineto{\pgfqpoint{2.113350in}{3.621111in}}%
\pgfpathlineto{\pgfqpoint{2.107150in}{3.621111in}}%
\pgfpathlineto{\pgfqpoint{2.107150in}{0.660000in}}%
\pgfpathclose%
\pgfusepath{fill}%
\end{pgfscope}%
\begin{pgfscope}%
\pgfpathrectangle{\pgfqpoint{1.250000in}{0.660000in}}{\pgfqpoint{7.750000in}{4.620000in}}%
\pgfusepath{clip}%
\pgfsetbuttcap%
\pgfsetmiterjoin%
\definecolor{currentfill}{rgb}{0.000000,0.000000,0.000000}%
\pgfsetfillcolor{currentfill}%
\pgfsetlinewidth{0.000000pt}%
\definecolor{currentstroke}{rgb}{0.000000,0.000000,0.000000}%
\pgfsetstrokecolor{currentstroke}%
\pgfsetstrokeopacity{0.000000}%
\pgfsetdash{}{0pt}%
\pgfpathmoveto{\pgfqpoint{2.114900in}{0.660000in}}%
\pgfpathlineto{\pgfqpoint{2.121100in}{0.660000in}}%
\pgfpathlineto{\pgfqpoint{2.121100in}{3.762210in}}%
\pgfpathlineto{\pgfqpoint{2.114900in}{3.762210in}}%
\pgfpathlineto{\pgfqpoint{2.114900in}{0.660000in}}%
\pgfpathclose%
\pgfusepath{fill}%
\end{pgfscope}%
\begin{pgfscope}%
\pgfpathrectangle{\pgfqpoint{1.250000in}{0.660000in}}{\pgfqpoint{7.750000in}{4.620000in}}%
\pgfusepath{clip}%
\pgfsetbuttcap%
\pgfsetmiterjoin%
\definecolor{currentfill}{rgb}{0.000000,0.000000,0.000000}%
\pgfsetfillcolor{currentfill}%
\pgfsetlinewidth{0.000000pt}%
\definecolor{currentstroke}{rgb}{0.000000,0.000000,0.000000}%
\pgfsetstrokecolor{currentstroke}%
\pgfsetstrokeopacity{0.000000}%
\pgfsetdash{}{0pt}%
\pgfpathmoveto{\pgfqpoint{2.122650in}{0.660000in}}%
\pgfpathlineto{\pgfqpoint{2.128850in}{0.660000in}}%
\pgfpathlineto{\pgfqpoint{2.128850in}{3.702434in}}%
\pgfpathlineto{\pgfqpoint{2.122650in}{3.702434in}}%
\pgfpathlineto{\pgfqpoint{2.122650in}{0.660000in}}%
\pgfpathclose%
\pgfusepath{fill}%
\end{pgfscope}%
\begin{pgfscope}%
\pgfpathrectangle{\pgfqpoint{1.250000in}{0.660000in}}{\pgfqpoint{7.750000in}{4.620000in}}%
\pgfusepath{clip}%
\pgfsetbuttcap%
\pgfsetmiterjoin%
\definecolor{currentfill}{rgb}{0.000000,0.000000,0.000000}%
\pgfsetfillcolor{currentfill}%
\pgfsetlinewidth{0.000000pt}%
\definecolor{currentstroke}{rgb}{0.000000,0.000000,0.000000}%
\pgfsetstrokecolor{currentstroke}%
\pgfsetstrokeopacity{0.000000}%
\pgfsetdash{}{0pt}%
\pgfpathmoveto{\pgfqpoint{2.130400in}{0.660000in}}%
\pgfpathlineto{\pgfqpoint{2.136600in}{0.660000in}}%
\pgfpathlineto{\pgfqpoint{2.136600in}{3.630878in}}%
\pgfpathlineto{\pgfqpoint{2.130400in}{3.630878in}}%
\pgfpathlineto{\pgfqpoint{2.130400in}{0.660000in}}%
\pgfpathclose%
\pgfusepath{fill}%
\end{pgfscope}%
\begin{pgfscope}%
\pgfpathrectangle{\pgfqpoint{1.250000in}{0.660000in}}{\pgfqpoint{7.750000in}{4.620000in}}%
\pgfusepath{clip}%
\pgfsetbuttcap%
\pgfsetmiterjoin%
\definecolor{currentfill}{rgb}{0.000000,0.000000,0.000000}%
\pgfsetfillcolor{currentfill}%
\pgfsetlinewidth{0.000000pt}%
\definecolor{currentstroke}{rgb}{0.000000,0.000000,0.000000}%
\pgfsetstrokecolor{currentstroke}%
\pgfsetstrokeopacity{0.000000}%
\pgfsetdash{}{0pt}%
\pgfpathmoveto{\pgfqpoint{2.138150in}{0.660000in}}%
\pgfpathlineto{\pgfqpoint{2.144350in}{0.660000in}}%
\pgfpathlineto{\pgfqpoint{2.144350in}{3.665330in}}%
\pgfpathlineto{\pgfqpoint{2.138150in}{3.665330in}}%
\pgfpathlineto{\pgfqpoint{2.138150in}{0.660000in}}%
\pgfpathclose%
\pgfusepath{fill}%
\end{pgfscope}%
\begin{pgfscope}%
\pgfpathrectangle{\pgfqpoint{1.250000in}{0.660000in}}{\pgfqpoint{7.750000in}{4.620000in}}%
\pgfusepath{clip}%
\pgfsetbuttcap%
\pgfsetmiterjoin%
\definecolor{currentfill}{rgb}{0.000000,0.000000,0.000000}%
\pgfsetfillcolor{currentfill}%
\pgfsetlinewidth{0.000000pt}%
\definecolor{currentstroke}{rgb}{0.000000,0.000000,0.000000}%
\pgfsetstrokecolor{currentstroke}%
\pgfsetstrokeopacity{0.000000}%
\pgfsetdash{}{0pt}%
\pgfpathmoveto{\pgfqpoint{2.145900in}{0.660000in}}%
\pgfpathlineto{\pgfqpoint{2.152100in}{0.660000in}}%
\pgfpathlineto{\pgfqpoint{2.152100in}{3.566640in}}%
\pgfpathlineto{\pgfqpoint{2.145900in}{3.566640in}}%
\pgfpathlineto{\pgfqpoint{2.145900in}{0.660000in}}%
\pgfpathclose%
\pgfusepath{fill}%
\end{pgfscope}%
\begin{pgfscope}%
\pgfpathrectangle{\pgfqpoint{1.250000in}{0.660000in}}{\pgfqpoint{7.750000in}{4.620000in}}%
\pgfusepath{clip}%
\pgfsetbuttcap%
\pgfsetmiterjoin%
\definecolor{currentfill}{rgb}{0.000000,0.000000,0.000000}%
\pgfsetfillcolor{currentfill}%
\pgfsetlinewidth{0.000000pt}%
\definecolor{currentstroke}{rgb}{0.000000,0.000000,0.000000}%
\pgfsetstrokecolor{currentstroke}%
\pgfsetstrokeopacity{0.000000}%
\pgfsetdash{}{0pt}%
\pgfpathmoveto{\pgfqpoint{2.153650in}{0.660000in}}%
\pgfpathlineto{\pgfqpoint{2.159850in}{0.660000in}}%
\pgfpathlineto{\pgfqpoint{2.159850in}{3.630016in}}%
\pgfpathlineto{\pgfqpoint{2.153650in}{3.630016in}}%
\pgfpathlineto{\pgfqpoint{2.153650in}{0.660000in}}%
\pgfpathclose%
\pgfusepath{fill}%
\end{pgfscope}%
\begin{pgfscope}%
\pgfpathrectangle{\pgfqpoint{1.250000in}{0.660000in}}{\pgfqpoint{7.750000in}{4.620000in}}%
\pgfusepath{clip}%
\pgfsetbuttcap%
\pgfsetmiterjoin%
\definecolor{currentfill}{rgb}{0.000000,0.000000,0.000000}%
\pgfsetfillcolor{currentfill}%
\pgfsetlinewidth{0.000000pt}%
\definecolor{currentstroke}{rgb}{0.000000,0.000000,0.000000}%
\pgfsetstrokecolor{currentstroke}%
\pgfsetstrokeopacity{0.000000}%
\pgfsetdash{}{0pt}%
\pgfpathmoveto{\pgfqpoint{2.161400in}{0.660000in}}%
\pgfpathlineto{\pgfqpoint{2.167600in}{0.660000in}}%
\pgfpathlineto{\pgfqpoint{2.167600in}{3.613710in}}%
\pgfpathlineto{\pgfqpoint{2.161400in}{3.613710in}}%
\pgfpathlineto{\pgfqpoint{2.161400in}{0.660000in}}%
\pgfpathclose%
\pgfusepath{fill}%
\end{pgfscope}%
\begin{pgfscope}%
\pgfpathrectangle{\pgfqpoint{1.250000in}{0.660000in}}{\pgfqpoint{7.750000in}{4.620000in}}%
\pgfusepath{clip}%
\pgfsetbuttcap%
\pgfsetmiterjoin%
\definecolor{currentfill}{rgb}{0.000000,0.000000,0.000000}%
\pgfsetfillcolor{currentfill}%
\pgfsetlinewidth{0.000000pt}%
\definecolor{currentstroke}{rgb}{0.000000,0.000000,0.000000}%
\pgfsetstrokecolor{currentstroke}%
\pgfsetstrokeopacity{0.000000}%
\pgfsetdash{}{0pt}%
\pgfpathmoveto{\pgfqpoint{2.169150in}{0.660000in}}%
\pgfpathlineto{\pgfqpoint{2.175350in}{0.660000in}}%
\pgfpathlineto{\pgfqpoint{2.175350in}{3.648256in}}%
\pgfpathlineto{\pgfqpoint{2.169150in}{3.648256in}}%
\pgfpathlineto{\pgfqpoint{2.169150in}{0.660000in}}%
\pgfpathclose%
\pgfusepath{fill}%
\end{pgfscope}%
\begin{pgfscope}%
\pgfpathrectangle{\pgfqpoint{1.250000in}{0.660000in}}{\pgfqpoint{7.750000in}{4.620000in}}%
\pgfusepath{clip}%
\pgfsetbuttcap%
\pgfsetmiterjoin%
\definecolor{currentfill}{rgb}{0.000000,0.000000,0.000000}%
\pgfsetfillcolor{currentfill}%
\pgfsetlinewidth{0.000000pt}%
\definecolor{currentstroke}{rgb}{0.000000,0.000000,0.000000}%
\pgfsetstrokecolor{currentstroke}%
\pgfsetstrokeopacity{0.000000}%
\pgfsetdash{}{0pt}%
\pgfpathmoveto{\pgfqpoint{2.176900in}{0.660000in}}%
\pgfpathlineto{\pgfqpoint{2.183100in}{0.660000in}}%
\pgfpathlineto{\pgfqpoint{2.183100in}{3.648634in}}%
\pgfpathlineto{\pgfqpoint{2.176900in}{3.648634in}}%
\pgfpathlineto{\pgfqpoint{2.176900in}{0.660000in}}%
\pgfpathclose%
\pgfusepath{fill}%
\end{pgfscope}%
\begin{pgfscope}%
\pgfpathrectangle{\pgfqpoint{1.250000in}{0.660000in}}{\pgfqpoint{7.750000in}{4.620000in}}%
\pgfusepath{clip}%
\pgfsetbuttcap%
\pgfsetmiterjoin%
\definecolor{currentfill}{rgb}{0.000000,0.000000,0.000000}%
\pgfsetfillcolor{currentfill}%
\pgfsetlinewidth{0.000000pt}%
\definecolor{currentstroke}{rgb}{0.000000,0.000000,0.000000}%
\pgfsetstrokecolor{currentstroke}%
\pgfsetstrokeopacity{0.000000}%
\pgfsetdash{}{0pt}%
\pgfpathmoveto{\pgfqpoint{2.184650in}{0.660000in}}%
\pgfpathlineto{\pgfqpoint{2.190850in}{0.660000in}}%
\pgfpathlineto{\pgfqpoint{2.190850in}{3.611613in}}%
\pgfpathlineto{\pgfqpoint{2.184650in}{3.611613in}}%
\pgfpathlineto{\pgfqpoint{2.184650in}{0.660000in}}%
\pgfpathclose%
\pgfusepath{fill}%
\end{pgfscope}%
\begin{pgfscope}%
\pgfpathrectangle{\pgfqpoint{1.250000in}{0.660000in}}{\pgfqpoint{7.750000in}{4.620000in}}%
\pgfusepath{clip}%
\pgfsetbuttcap%
\pgfsetmiterjoin%
\definecolor{currentfill}{rgb}{0.000000,0.000000,0.000000}%
\pgfsetfillcolor{currentfill}%
\pgfsetlinewidth{0.000000pt}%
\definecolor{currentstroke}{rgb}{0.000000,0.000000,0.000000}%
\pgfsetstrokecolor{currentstroke}%
\pgfsetstrokeopacity{0.000000}%
\pgfsetdash{}{0pt}%
\pgfpathmoveto{\pgfqpoint{2.192400in}{0.660000in}}%
\pgfpathlineto{\pgfqpoint{2.198600in}{0.660000in}}%
\pgfpathlineto{\pgfqpoint{2.198600in}{3.718285in}}%
\pgfpathlineto{\pgfqpoint{2.192400in}{3.718285in}}%
\pgfpathlineto{\pgfqpoint{2.192400in}{0.660000in}}%
\pgfpathclose%
\pgfusepath{fill}%
\end{pgfscope}%
\begin{pgfscope}%
\pgfpathrectangle{\pgfqpoint{1.250000in}{0.660000in}}{\pgfqpoint{7.750000in}{4.620000in}}%
\pgfusepath{clip}%
\pgfsetbuttcap%
\pgfsetmiterjoin%
\definecolor{currentfill}{rgb}{0.000000,0.000000,0.000000}%
\pgfsetfillcolor{currentfill}%
\pgfsetlinewidth{0.000000pt}%
\definecolor{currentstroke}{rgb}{0.000000,0.000000,0.000000}%
\pgfsetstrokecolor{currentstroke}%
\pgfsetstrokeopacity{0.000000}%
\pgfsetdash{}{0pt}%
\pgfpathmoveto{\pgfqpoint{2.200150in}{0.660000in}}%
\pgfpathlineto{\pgfqpoint{2.206350in}{0.660000in}}%
\pgfpathlineto{\pgfqpoint{2.206350in}{3.775663in}}%
\pgfpathlineto{\pgfqpoint{2.200150in}{3.775663in}}%
\pgfpathlineto{\pgfqpoint{2.200150in}{0.660000in}}%
\pgfpathclose%
\pgfusepath{fill}%
\end{pgfscope}%
\begin{pgfscope}%
\pgfpathrectangle{\pgfqpoint{1.250000in}{0.660000in}}{\pgfqpoint{7.750000in}{4.620000in}}%
\pgfusepath{clip}%
\pgfsetbuttcap%
\pgfsetmiterjoin%
\definecolor{currentfill}{rgb}{0.000000,0.000000,0.000000}%
\pgfsetfillcolor{currentfill}%
\pgfsetlinewidth{0.000000pt}%
\definecolor{currentstroke}{rgb}{0.000000,0.000000,0.000000}%
\pgfsetstrokecolor{currentstroke}%
\pgfsetstrokeopacity{0.000000}%
\pgfsetdash{}{0pt}%
\pgfpathmoveto{\pgfqpoint{2.207900in}{0.660000in}}%
\pgfpathlineto{\pgfqpoint{2.214100in}{0.660000in}}%
\pgfpathlineto{\pgfqpoint{2.214100in}{3.671009in}}%
\pgfpathlineto{\pgfqpoint{2.207900in}{3.671009in}}%
\pgfpathlineto{\pgfqpoint{2.207900in}{0.660000in}}%
\pgfpathclose%
\pgfusepath{fill}%
\end{pgfscope}%
\begin{pgfscope}%
\pgfpathrectangle{\pgfqpoint{1.250000in}{0.660000in}}{\pgfqpoint{7.750000in}{4.620000in}}%
\pgfusepath{clip}%
\pgfsetbuttcap%
\pgfsetmiterjoin%
\definecolor{currentfill}{rgb}{0.000000,0.000000,0.000000}%
\pgfsetfillcolor{currentfill}%
\pgfsetlinewidth{0.000000pt}%
\definecolor{currentstroke}{rgb}{0.000000,0.000000,0.000000}%
\pgfsetstrokecolor{currentstroke}%
\pgfsetstrokeopacity{0.000000}%
\pgfsetdash{}{0pt}%
\pgfpathmoveto{\pgfqpoint{2.215650in}{0.660000in}}%
\pgfpathlineto{\pgfqpoint{2.221850in}{0.660000in}}%
\pgfpathlineto{\pgfqpoint{2.221850in}{3.599258in}}%
\pgfpathlineto{\pgfqpoint{2.215650in}{3.599258in}}%
\pgfpathlineto{\pgfqpoint{2.215650in}{0.660000in}}%
\pgfpathclose%
\pgfusepath{fill}%
\end{pgfscope}%
\begin{pgfscope}%
\pgfpathrectangle{\pgfqpoint{1.250000in}{0.660000in}}{\pgfqpoint{7.750000in}{4.620000in}}%
\pgfusepath{clip}%
\pgfsetbuttcap%
\pgfsetmiterjoin%
\definecolor{currentfill}{rgb}{0.000000,0.000000,0.000000}%
\pgfsetfillcolor{currentfill}%
\pgfsetlinewidth{0.000000pt}%
\definecolor{currentstroke}{rgb}{0.000000,0.000000,0.000000}%
\pgfsetstrokecolor{currentstroke}%
\pgfsetstrokeopacity{0.000000}%
\pgfsetdash{}{0pt}%
\pgfpathmoveto{\pgfqpoint{2.223400in}{0.660000in}}%
\pgfpathlineto{\pgfqpoint{2.229600in}{0.660000in}}%
\pgfpathlineto{\pgfqpoint{2.229600in}{3.734803in}}%
\pgfpathlineto{\pgfqpoint{2.223400in}{3.734803in}}%
\pgfpathlineto{\pgfqpoint{2.223400in}{0.660000in}}%
\pgfpathclose%
\pgfusepath{fill}%
\end{pgfscope}%
\begin{pgfscope}%
\pgfpathrectangle{\pgfqpoint{1.250000in}{0.660000in}}{\pgfqpoint{7.750000in}{4.620000in}}%
\pgfusepath{clip}%
\pgfsetbuttcap%
\pgfsetmiterjoin%
\definecolor{currentfill}{rgb}{0.000000,0.000000,0.000000}%
\pgfsetfillcolor{currentfill}%
\pgfsetlinewidth{0.000000pt}%
\definecolor{currentstroke}{rgb}{0.000000,0.000000,0.000000}%
\pgfsetstrokecolor{currentstroke}%
\pgfsetstrokeopacity{0.000000}%
\pgfsetdash{}{0pt}%
\pgfpathmoveto{\pgfqpoint{2.231150in}{0.660000in}}%
\pgfpathlineto{\pgfqpoint{2.237350in}{0.660000in}}%
\pgfpathlineto{\pgfqpoint{2.237350in}{3.553792in}}%
\pgfpathlineto{\pgfqpoint{2.231150in}{3.553792in}}%
\pgfpathlineto{\pgfqpoint{2.231150in}{0.660000in}}%
\pgfpathclose%
\pgfusepath{fill}%
\end{pgfscope}%
\begin{pgfscope}%
\pgfpathrectangle{\pgfqpoint{1.250000in}{0.660000in}}{\pgfqpoint{7.750000in}{4.620000in}}%
\pgfusepath{clip}%
\pgfsetbuttcap%
\pgfsetmiterjoin%
\definecolor{currentfill}{rgb}{0.000000,0.000000,0.000000}%
\pgfsetfillcolor{currentfill}%
\pgfsetlinewidth{0.000000pt}%
\definecolor{currentstroke}{rgb}{0.000000,0.000000,0.000000}%
\pgfsetstrokecolor{currentstroke}%
\pgfsetstrokeopacity{0.000000}%
\pgfsetdash{}{0pt}%
\pgfpathmoveto{\pgfqpoint{2.238900in}{0.660000in}}%
\pgfpathlineto{\pgfqpoint{2.245100in}{0.660000in}}%
\pgfpathlineto{\pgfqpoint{2.245100in}{3.565717in}}%
\pgfpathlineto{\pgfqpoint{2.238900in}{3.565717in}}%
\pgfpathlineto{\pgfqpoint{2.238900in}{0.660000in}}%
\pgfpathclose%
\pgfusepath{fill}%
\end{pgfscope}%
\begin{pgfscope}%
\pgfpathrectangle{\pgfqpoint{1.250000in}{0.660000in}}{\pgfqpoint{7.750000in}{4.620000in}}%
\pgfusepath{clip}%
\pgfsetbuttcap%
\pgfsetmiterjoin%
\definecolor{currentfill}{rgb}{0.000000,0.000000,0.000000}%
\pgfsetfillcolor{currentfill}%
\pgfsetlinewidth{0.000000pt}%
\definecolor{currentstroke}{rgb}{0.000000,0.000000,0.000000}%
\pgfsetstrokecolor{currentstroke}%
\pgfsetstrokeopacity{0.000000}%
\pgfsetdash{}{0pt}%
\pgfpathmoveto{\pgfqpoint{2.246650in}{0.660000in}}%
\pgfpathlineto{\pgfqpoint{2.252850in}{0.660000in}}%
\pgfpathlineto{\pgfqpoint{2.252850in}{3.693115in}}%
\pgfpathlineto{\pgfqpoint{2.246650in}{3.693115in}}%
\pgfpathlineto{\pgfqpoint{2.246650in}{0.660000in}}%
\pgfpathclose%
\pgfusepath{fill}%
\end{pgfscope}%
\begin{pgfscope}%
\pgfpathrectangle{\pgfqpoint{1.250000in}{0.660000in}}{\pgfqpoint{7.750000in}{4.620000in}}%
\pgfusepath{clip}%
\pgfsetbuttcap%
\pgfsetmiterjoin%
\definecolor{currentfill}{rgb}{0.000000,0.000000,0.000000}%
\pgfsetfillcolor{currentfill}%
\pgfsetlinewidth{0.000000pt}%
\definecolor{currentstroke}{rgb}{0.000000,0.000000,0.000000}%
\pgfsetstrokecolor{currentstroke}%
\pgfsetstrokeopacity{0.000000}%
\pgfsetdash{}{0pt}%
\pgfpathmoveto{\pgfqpoint{2.254400in}{0.660000in}}%
\pgfpathlineto{\pgfqpoint{2.260600in}{0.660000in}}%
\pgfpathlineto{\pgfqpoint{2.260600in}{3.589295in}}%
\pgfpathlineto{\pgfqpoint{2.254400in}{3.589295in}}%
\pgfpathlineto{\pgfqpoint{2.254400in}{0.660000in}}%
\pgfpathclose%
\pgfusepath{fill}%
\end{pgfscope}%
\begin{pgfscope}%
\pgfpathrectangle{\pgfqpoint{1.250000in}{0.660000in}}{\pgfqpoint{7.750000in}{4.620000in}}%
\pgfusepath{clip}%
\pgfsetbuttcap%
\pgfsetmiterjoin%
\definecolor{currentfill}{rgb}{0.000000,0.000000,0.000000}%
\pgfsetfillcolor{currentfill}%
\pgfsetlinewidth{0.000000pt}%
\definecolor{currentstroke}{rgb}{0.000000,0.000000,0.000000}%
\pgfsetstrokecolor{currentstroke}%
\pgfsetstrokeopacity{0.000000}%
\pgfsetdash{}{0pt}%
\pgfpathmoveto{\pgfqpoint{2.262150in}{0.660000in}}%
\pgfpathlineto{\pgfqpoint{2.268350in}{0.660000in}}%
\pgfpathlineto{\pgfqpoint{2.268350in}{3.680591in}}%
\pgfpathlineto{\pgfqpoint{2.262150in}{3.680591in}}%
\pgfpathlineto{\pgfqpoint{2.262150in}{0.660000in}}%
\pgfpathclose%
\pgfusepath{fill}%
\end{pgfscope}%
\begin{pgfscope}%
\pgfpathrectangle{\pgfqpoint{1.250000in}{0.660000in}}{\pgfqpoint{7.750000in}{4.620000in}}%
\pgfusepath{clip}%
\pgfsetbuttcap%
\pgfsetmiterjoin%
\definecolor{currentfill}{rgb}{0.000000,0.000000,0.000000}%
\pgfsetfillcolor{currentfill}%
\pgfsetlinewidth{0.000000pt}%
\definecolor{currentstroke}{rgb}{0.000000,0.000000,0.000000}%
\pgfsetstrokecolor{currentstroke}%
\pgfsetstrokeopacity{0.000000}%
\pgfsetdash{}{0pt}%
\pgfpathmoveto{\pgfqpoint{2.269900in}{0.660000in}}%
\pgfpathlineto{\pgfqpoint{2.276100in}{0.660000in}}%
\pgfpathlineto{\pgfqpoint{2.276100in}{3.539784in}}%
\pgfpathlineto{\pgfqpoint{2.269900in}{3.539784in}}%
\pgfpathlineto{\pgfqpoint{2.269900in}{0.660000in}}%
\pgfpathclose%
\pgfusepath{fill}%
\end{pgfscope}%
\begin{pgfscope}%
\pgfpathrectangle{\pgfqpoint{1.250000in}{0.660000in}}{\pgfqpoint{7.750000in}{4.620000in}}%
\pgfusepath{clip}%
\pgfsetbuttcap%
\pgfsetmiterjoin%
\definecolor{currentfill}{rgb}{0.000000,0.000000,0.000000}%
\pgfsetfillcolor{currentfill}%
\pgfsetlinewidth{0.000000pt}%
\definecolor{currentstroke}{rgb}{0.000000,0.000000,0.000000}%
\pgfsetstrokecolor{currentstroke}%
\pgfsetstrokeopacity{0.000000}%
\pgfsetdash{}{0pt}%
\pgfpathmoveto{\pgfqpoint{2.277650in}{0.660000in}}%
\pgfpathlineto{\pgfqpoint{2.283850in}{0.660000in}}%
\pgfpathlineto{\pgfqpoint{2.283850in}{3.541501in}}%
\pgfpathlineto{\pgfqpoint{2.277650in}{3.541501in}}%
\pgfpathlineto{\pgfqpoint{2.277650in}{0.660000in}}%
\pgfpathclose%
\pgfusepath{fill}%
\end{pgfscope}%
\begin{pgfscope}%
\pgfpathrectangle{\pgfqpoint{1.250000in}{0.660000in}}{\pgfqpoint{7.750000in}{4.620000in}}%
\pgfusepath{clip}%
\pgfsetbuttcap%
\pgfsetmiterjoin%
\definecolor{currentfill}{rgb}{0.000000,0.000000,0.000000}%
\pgfsetfillcolor{currentfill}%
\pgfsetlinewidth{0.000000pt}%
\definecolor{currentstroke}{rgb}{0.000000,0.000000,0.000000}%
\pgfsetstrokecolor{currentstroke}%
\pgfsetstrokeopacity{0.000000}%
\pgfsetdash{}{0pt}%
\pgfpathmoveto{\pgfqpoint{2.285400in}{0.660000in}}%
\pgfpathlineto{\pgfqpoint{2.291600in}{0.660000in}}%
\pgfpathlineto{\pgfqpoint{2.291600in}{3.584421in}}%
\pgfpathlineto{\pgfqpoint{2.285400in}{3.584421in}}%
\pgfpathlineto{\pgfqpoint{2.285400in}{0.660000in}}%
\pgfpathclose%
\pgfusepath{fill}%
\end{pgfscope}%
\begin{pgfscope}%
\pgfpathrectangle{\pgfqpoint{1.250000in}{0.660000in}}{\pgfqpoint{7.750000in}{4.620000in}}%
\pgfusepath{clip}%
\pgfsetbuttcap%
\pgfsetmiterjoin%
\definecolor{currentfill}{rgb}{0.000000,0.000000,0.000000}%
\pgfsetfillcolor{currentfill}%
\pgfsetlinewidth{0.000000pt}%
\definecolor{currentstroke}{rgb}{0.000000,0.000000,0.000000}%
\pgfsetstrokecolor{currentstroke}%
\pgfsetstrokeopacity{0.000000}%
\pgfsetdash{}{0pt}%
\pgfpathmoveto{\pgfqpoint{2.293150in}{0.660000in}}%
\pgfpathlineto{\pgfqpoint{2.299350in}{0.660000in}}%
\pgfpathlineto{\pgfqpoint{2.299350in}{3.631387in}}%
\pgfpathlineto{\pgfqpoint{2.293150in}{3.631387in}}%
\pgfpathlineto{\pgfqpoint{2.293150in}{0.660000in}}%
\pgfpathclose%
\pgfusepath{fill}%
\end{pgfscope}%
\begin{pgfscope}%
\pgfpathrectangle{\pgfqpoint{1.250000in}{0.660000in}}{\pgfqpoint{7.750000in}{4.620000in}}%
\pgfusepath{clip}%
\pgfsetbuttcap%
\pgfsetmiterjoin%
\definecolor{currentfill}{rgb}{0.000000,0.000000,0.000000}%
\pgfsetfillcolor{currentfill}%
\pgfsetlinewidth{0.000000pt}%
\definecolor{currentstroke}{rgb}{0.000000,0.000000,0.000000}%
\pgfsetstrokecolor{currentstroke}%
\pgfsetstrokeopacity{0.000000}%
\pgfsetdash{}{0pt}%
\pgfpathmoveto{\pgfqpoint{2.300900in}{0.660000in}}%
\pgfpathlineto{\pgfqpoint{2.307100in}{0.660000in}}%
\pgfpathlineto{\pgfqpoint{2.307100in}{3.618910in}}%
\pgfpathlineto{\pgfqpoint{2.300900in}{3.618910in}}%
\pgfpathlineto{\pgfqpoint{2.300900in}{0.660000in}}%
\pgfpathclose%
\pgfusepath{fill}%
\end{pgfscope}%
\begin{pgfscope}%
\pgfpathrectangle{\pgfqpoint{1.250000in}{0.660000in}}{\pgfqpoint{7.750000in}{4.620000in}}%
\pgfusepath{clip}%
\pgfsetbuttcap%
\pgfsetmiterjoin%
\definecolor{currentfill}{rgb}{0.000000,0.000000,0.000000}%
\pgfsetfillcolor{currentfill}%
\pgfsetlinewidth{0.000000pt}%
\definecolor{currentstroke}{rgb}{0.000000,0.000000,0.000000}%
\pgfsetstrokecolor{currentstroke}%
\pgfsetstrokeopacity{0.000000}%
\pgfsetdash{}{0pt}%
\pgfpathmoveto{\pgfqpoint{2.308650in}{0.660000in}}%
\pgfpathlineto{\pgfqpoint{2.314850in}{0.660000in}}%
\pgfpathlineto{\pgfqpoint{2.314850in}{3.563007in}}%
\pgfpathlineto{\pgfqpoint{2.308650in}{3.563007in}}%
\pgfpathlineto{\pgfqpoint{2.308650in}{0.660000in}}%
\pgfpathclose%
\pgfusepath{fill}%
\end{pgfscope}%
\begin{pgfscope}%
\pgfpathrectangle{\pgfqpoint{1.250000in}{0.660000in}}{\pgfqpoint{7.750000in}{4.620000in}}%
\pgfusepath{clip}%
\pgfsetbuttcap%
\pgfsetmiterjoin%
\definecolor{currentfill}{rgb}{0.000000,0.000000,0.000000}%
\pgfsetfillcolor{currentfill}%
\pgfsetlinewidth{0.000000pt}%
\definecolor{currentstroke}{rgb}{0.000000,0.000000,0.000000}%
\pgfsetstrokecolor{currentstroke}%
\pgfsetstrokeopacity{0.000000}%
\pgfsetdash{}{0pt}%
\pgfpathmoveto{\pgfqpoint{2.316400in}{0.660000in}}%
\pgfpathlineto{\pgfqpoint{2.322600in}{0.660000in}}%
\pgfpathlineto{\pgfqpoint{2.322600in}{3.568835in}}%
\pgfpathlineto{\pgfqpoint{2.316400in}{3.568835in}}%
\pgfpathlineto{\pgfqpoint{2.316400in}{0.660000in}}%
\pgfpathclose%
\pgfusepath{fill}%
\end{pgfscope}%
\begin{pgfscope}%
\pgfpathrectangle{\pgfqpoint{1.250000in}{0.660000in}}{\pgfqpoint{7.750000in}{4.620000in}}%
\pgfusepath{clip}%
\pgfsetbuttcap%
\pgfsetmiterjoin%
\definecolor{currentfill}{rgb}{0.000000,0.000000,0.000000}%
\pgfsetfillcolor{currentfill}%
\pgfsetlinewidth{0.000000pt}%
\definecolor{currentstroke}{rgb}{0.000000,0.000000,0.000000}%
\pgfsetstrokecolor{currentstroke}%
\pgfsetstrokeopacity{0.000000}%
\pgfsetdash{}{0pt}%
\pgfpathmoveto{\pgfqpoint{2.324150in}{0.660000in}}%
\pgfpathlineto{\pgfqpoint{2.330350in}{0.660000in}}%
\pgfpathlineto{\pgfqpoint{2.330350in}{3.546631in}}%
\pgfpathlineto{\pgfqpoint{2.324150in}{3.546631in}}%
\pgfpathlineto{\pgfqpoint{2.324150in}{0.660000in}}%
\pgfpathclose%
\pgfusepath{fill}%
\end{pgfscope}%
\begin{pgfscope}%
\pgfpathrectangle{\pgfqpoint{1.250000in}{0.660000in}}{\pgfqpoint{7.750000in}{4.620000in}}%
\pgfusepath{clip}%
\pgfsetbuttcap%
\pgfsetmiterjoin%
\definecolor{currentfill}{rgb}{0.000000,0.000000,0.000000}%
\pgfsetfillcolor{currentfill}%
\pgfsetlinewidth{0.000000pt}%
\definecolor{currentstroke}{rgb}{0.000000,0.000000,0.000000}%
\pgfsetstrokecolor{currentstroke}%
\pgfsetstrokeopacity{0.000000}%
\pgfsetdash{}{0pt}%
\pgfpathmoveto{\pgfqpoint{2.331900in}{0.660000in}}%
\pgfpathlineto{\pgfqpoint{2.338100in}{0.660000in}}%
\pgfpathlineto{\pgfqpoint{2.338100in}{3.623102in}}%
\pgfpathlineto{\pgfqpoint{2.331900in}{3.623102in}}%
\pgfpathlineto{\pgfqpoint{2.331900in}{0.660000in}}%
\pgfpathclose%
\pgfusepath{fill}%
\end{pgfscope}%
\begin{pgfscope}%
\pgfpathrectangle{\pgfqpoint{1.250000in}{0.660000in}}{\pgfqpoint{7.750000in}{4.620000in}}%
\pgfusepath{clip}%
\pgfsetbuttcap%
\pgfsetmiterjoin%
\definecolor{currentfill}{rgb}{0.000000,0.000000,0.000000}%
\pgfsetfillcolor{currentfill}%
\pgfsetlinewidth{0.000000pt}%
\definecolor{currentstroke}{rgb}{0.000000,0.000000,0.000000}%
\pgfsetstrokecolor{currentstroke}%
\pgfsetstrokeopacity{0.000000}%
\pgfsetdash{}{0pt}%
\pgfpathmoveto{\pgfqpoint{2.339650in}{0.660000in}}%
\pgfpathlineto{\pgfqpoint{2.345850in}{0.660000in}}%
\pgfpathlineto{\pgfqpoint{2.345850in}{3.522825in}}%
\pgfpathlineto{\pgfqpoint{2.339650in}{3.522825in}}%
\pgfpathlineto{\pgfqpoint{2.339650in}{0.660000in}}%
\pgfpathclose%
\pgfusepath{fill}%
\end{pgfscope}%
\begin{pgfscope}%
\pgfpathrectangle{\pgfqpoint{1.250000in}{0.660000in}}{\pgfqpoint{7.750000in}{4.620000in}}%
\pgfusepath{clip}%
\pgfsetbuttcap%
\pgfsetmiterjoin%
\definecolor{currentfill}{rgb}{0.000000,0.000000,0.000000}%
\pgfsetfillcolor{currentfill}%
\pgfsetlinewidth{0.000000pt}%
\definecolor{currentstroke}{rgb}{0.000000,0.000000,0.000000}%
\pgfsetstrokecolor{currentstroke}%
\pgfsetstrokeopacity{0.000000}%
\pgfsetdash{}{0pt}%
\pgfpathmoveto{\pgfqpoint{2.347400in}{0.660000in}}%
\pgfpathlineto{\pgfqpoint{2.353600in}{0.660000in}}%
\pgfpathlineto{\pgfqpoint{2.353600in}{3.524501in}}%
\pgfpathlineto{\pgfqpoint{2.347400in}{3.524501in}}%
\pgfpathlineto{\pgfqpoint{2.347400in}{0.660000in}}%
\pgfpathclose%
\pgfusepath{fill}%
\end{pgfscope}%
\begin{pgfscope}%
\pgfpathrectangle{\pgfqpoint{1.250000in}{0.660000in}}{\pgfqpoint{7.750000in}{4.620000in}}%
\pgfusepath{clip}%
\pgfsetbuttcap%
\pgfsetmiterjoin%
\definecolor{currentfill}{rgb}{0.000000,0.000000,0.000000}%
\pgfsetfillcolor{currentfill}%
\pgfsetlinewidth{0.000000pt}%
\definecolor{currentstroke}{rgb}{0.000000,0.000000,0.000000}%
\pgfsetstrokecolor{currentstroke}%
\pgfsetstrokeopacity{0.000000}%
\pgfsetdash{}{0pt}%
\pgfpathmoveto{\pgfqpoint{2.355150in}{0.660000in}}%
\pgfpathlineto{\pgfqpoint{2.361350in}{0.660000in}}%
\pgfpathlineto{\pgfqpoint{2.361350in}{3.580569in}}%
\pgfpathlineto{\pgfqpoint{2.355150in}{3.580569in}}%
\pgfpathlineto{\pgfqpoint{2.355150in}{0.660000in}}%
\pgfpathclose%
\pgfusepath{fill}%
\end{pgfscope}%
\begin{pgfscope}%
\pgfpathrectangle{\pgfqpoint{1.250000in}{0.660000in}}{\pgfqpoint{7.750000in}{4.620000in}}%
\pgfusepath{clip}%
\pgfsetbuttcap%
\pgfsetmiterjoin%
\definecolor{currentfill}{rgb}{0.000000,0.000000,0.000000}%
\pgfsetfillcolor{currentfill}%
\pgfsetlinewidth{0.000000pt}%
\definecolor{currentstroke}{rgb}{0.000000,0.000000,0.000000}%
\pgfsetstrokecolor{currentstroke}%
\pgfsetstrokeopacity{0.000000}%
\pgfsetdash{}{0pt}%
\pgfpathmoveto{\pgfqpoint{2.362900in}{0.660000in}}%
\pgfpathlineto{\pgfqpoint{2.369100in}{0.660000in}}%
\pgfpathlineto{\pgfqpoint{2.369100in}{3.548233in}}%
\pgfpathlineto{\pgfqpoint{2.362900in}{3.548233in}}%
\pgfpathlineto{\pgfqpoint{2.362900in}{0.660000in}}%
\pgfpathclose%
\pgfusepath{fill}%
\end{pgfscope}%
\begin{pgfscope}%
\pgfpathrectangle{\pgfqpoint{1.250000in}{0.660000in}}{\pgfqpoint{7.750000in}{4.620000in}}%
\pgfusepath{clip}%
\pgfsetbuttcap%
\pgfsetmiterjoin%
\definecolor{currentfill}{rgb}{0.000000,0.000000,0.000000}%
\pgfsetfillcolor{currentfill}%
\pgfsetlinewidth{0.000000pt}%
\definecolor{currentstroke}{rgb}{0.000000,0.000000,0.000000}%
\pgfsetstrokecolor{currentstroke}%
\pgfsetstrokeopacity{0.000000}%
\pgfsetdash{}{0pt}%
\pgfpathmoveto{\pgfqpoint{2.370650in}{0.660000in}}%
\pgfpathlineto{\pgfqpoint{2.376850in}{0.660000in}}%
\pgfpathlineto{\pgfqpoint{2.376850in}{3.562059in}}%
\pgfpathlineto{\pgfqpoint{2.370650in}{3.562059in}}%
\pgfpathlineto{\pgfqpoint{2.370650in}{0.660000in}}%
\pgfpathclose%
\pgfusepath{fill}%
\end{pgfscope}%
\begin{pgfscope}%
\pgfpathrectangle{\pgfqpoint{1.250000in}{0.660000in}}{\pgfqpoint{7.750000in}{4.620000in}}%
\pgfusepath{clip}%
\pgfsetbuttcap%
\pgfsetmiterjoin%
\definecolor{currentfill}{rgb}{0.000000,0.000000,0.000000}%
\pgfsetfillcolor{currentfill}%
\pgfsetlinewidth{0.000000pt}%
\definecolor{currentstroke}{rgb}{0.000000,0.000000,0.000000}%
\pgfsetstrokecolor{currentstroke}%
\pgfsetstrokeopacity{0.000000}%
\pgfsetdash{}{0pt}%
\pgfpathmoveto{\pgfqpoint{2.378400in}{0.660000in}}%
\pgfpathlineto{\pgfqpoint{2.384600in}{0.660000in}}%
\pgfpathlineto{\pgfqpoint{2.384600in}{3.517832in}}%
\pgfpathlineto{\pgfqpoint{2.378400in}{3.517832in}}%
\pgfpathlineto{\pgfqpoint{2.378400in}{0.660000in}}%
\pgfpathclose%
\pgfusepath{fill}%
\end{pgfscope}%
\begin{pgfscope}%
\pgfpathrectangle{\pgfqpoint{1.250000in}{0.660000in}}{\pgfqpoint{7.750000in}{4.620000in}}%
\pgfusepath{clip}%
\pgfsetbuttcap%
\pgfsetmiterjoin%
\definecolor{currentfill}{rgb}{0.000000,0.000000,0.000000}%
\pgfsetfillcolor{currentfill}%
\pgfsetlinewidth{0.000000pt}%
\definecolor{currentstroke}{rgb}{0.000000,0.000000,0.000000}%
\pgfsetstrokecolor{currentstroke}%
\pgfsetstrokeopacity{0.000000}%
\pgfsetdash{}{0pt}%
\pgfpathmoveto{\pgfqpoint{2.386150in}{0.660000in}}%
\pgfpathlineto{\pgfqpoint{2.392350in}{0.660000in}}%
\pgfpathlineto{\pgfqpoint{2.392350in}{3.554610in}}%
\pgfpathlineto{\pgfqpoint{2.386150in}{3.554610in}}%
\pgfpathlineto{\pgfqpoint{2.386150in}{0.660000in}}%
\pgfpathclose%
\pgfusepath{fill}%
\end{pgfscope}%
\begin{pgfscope}%
\pgfpathrectangle{\pgfqpoint{1.250000in}{0.660000in}}{\pgfqpoint{7.750000in}{4.620000in}}%
\pgfusepath{clip}%
\pgfsetbuttcap%
\pgfsetmiterjoin%
\definecolor{currentfill}{rgb}{0.000000,0.000000,0.000000}%
\pgfsetfillcolor{currentfill}%
\pgfsetlinewidth{0.000000pt}%
\definecolor{currentstroke}{rgb}{0.000000,0.000000,0.000000}%
\pgfsetstrokecolor{currentstroke}%
\pgfsetstrokeopacity{0.000000}%
\pgfsetdash{}{0pt}%
\pgfpathmoveto{\pgfqpoint{2.393900in}{0.660000in}}%
\pgfpathlineto{\pgfqpoint{2.400100in}{0.660000in}}%
\pgfpathlineto{\pgfqpoint{2.400100in}{3.587431in}}%
\pgfpathlineto{\pgfqpoint{2.393900in}{3.587431in}}%
\pgfpathlineto{\pgfqpoint{2.393900in}{0.660000in}}%
\pgfpathclose%
\pgfusepath{fill}%
\end{pgfscope}%
\begin{pgfscope}%
\pgfpathrectangle{\pgfqpoint{1.250000in}{0.660000in}}{\pgfqpoint{7.750000in}{4.620000in}}%
\pgfusepath{clip}%
\pgfsetbuttcap%
\pgfsetmiterjoin%
\definecolor{currentfill}{rgb}{0.000000,0.000000,0.000000}%
\pgfsetfillcolor{currentfill}%
\pgfsetlinewidth{0.000000pt}%
\definecolor{currentstroke}{rgb}{0.000000,0.000000,0.000000}%
\pgfsetstrokecolor{currentstroke}%
\pgfsetstrokeopacity{0.000000}%
\pgfsetdash{}{0pt}%
\pgfpathmoveto{\pgfqpoint{2.401650in}{0.660000in}}%
\pgfpathlineto{\pgfqpoint{2.407850in}{0.660000in}}%
\pgfpathlineto{\pgfqpoint{2.407850in}{3.577377in}}%
\pgfpathlineto{\pgfqpoint{2.401650in}{3.577377in}}%
\pgfpathlineto{\pgfqpoint{2.401650in}{0.660000in}}%
\pgfpathclose%
\pgfusepath{fill}%
\end{pgfscope}%
\begin{pgfscope}%
\pgfpathrectangle{\pgfqpoint{1.250000in}{0.660000in}}{\pgfqpoint{7.750000in}{4.620000in}}%
\pgfusepath{clip}%
\pgfsetbuttcap%
\pgfsetmiterjoin%
\definecolor{currentfill}{rgb}{0.000000,0.000000,0.000000}%
\pgfsetfillcolor{currentfill}%
\pgfsetlinewidth{0.000000pt}%
\definecolor{currentstroke}{rgb}{0.000000,0.000000,0.000000}%
\pgfsetstrokecolor{currentstroke}%
\pgfsetstrokeopacity{0.000000}%
\pgfsetdash{}{0pt}%
\pgfpathmoveto{\pgfqpoint{2.409400in}{0.660000in}}%
\pgfpathlineto{\pgfqpoint{2.415600in}{0.660000in}}%
\pgfpathlineto{\pgfqpoint{2.415600in}{3.551896in}}%
\pgfpathlineto{\pgfqpoint{2.409400in}{3.551896in}}%
\pgfpathlineto{\pgfqpoint{2.409400in}{0.660000in}}%
\pgfpathclose%
\pgfusepath{fill}%
\end{pgfscope}%
\begin{pgfscope}%
\pgfpathrectangle{\pgfqpoint{1.250000in}{0.660000in}}{\pgfqpoint{7.750000in}{4.620000in}}%
\pgfusepath{clip}%
\pgfsetbuttcap%
\pgfsetmiterjoin%
\definecolor{currentfill}{rgb}{0.000000,0.000000,0.000000}%
\pgfsetfillcolor{currentfill}%
\pgfsetlinewidth{0.000000pt}%
\definecolor{currentstroke}{rgb}{0.000000,0.000000,0.000000}%
\pgfsetstrokecolor{currentstroke}%
\pgfsetstrokeopacity{0.000000}%
\pgfsetdash{}{0pt}%
\pgfpathmoveto{\pgfqpoint{2.417150in}{0.660000in}}%
\pgfpathlineto{\pgfqpoint{2.423350in}{0.660000in}}%
\pgfpathlineto{\pgfqpoint{2.423350in}{3.589801in}}%
\pgfpathlineto{\pgfqpoint{2.417150in}{3.589801in}}%
\pgfpathlineto{\pgfqpoint{2.417150in}{0.660000in}}%
\pgfpathclose%
\pgfusepath{fill}%
\end{pgfscope}%
\begin{pgfscope}%
\pgfpathrectangle{\pgfqpoint{1.250000in}{0.660000in}}{\pgfqpoint{7.750000in}{4.620000in}}%
\pgfusepath{clip}%
\pgfsetbuttcap%
\pgfsetmiterjoin%
\definecolor{currentfill}{rgb}{0.000000,0.000000,0.000000}%
\pgfsetfillcolor{currentfill}%
\pgfsetlinewidth{0.000000pt}%
\definecolor{currentstroke}{rgb}{0.000000,0.000000,0.000000}%
\pgfsetstrokecolor{currentstroke}%
\pgfsetstrokeopacity{0.000000}%
\pgfsetdash{}{0pt}%
\pgfpathmoveto{\pgfqpoint{2.424900in}{0.660000in}}%
\pgfpathlineto{\pgfqpoint{2.431100in}{0.660000in}}%
\pgfpathlineto{\pgfqpoint{2.431100in}{3.563631in}}%
\pgfpathlineto{\pgfqpoint{2.424900in}{3.563631in}}%
\pgfpathlineto{\pgfqpoint{2.424900in}{0.660000in}}%
\pgfpathclose%
\pgfusepath{fill}%
\end{pgfscope}%
\begin{pgfscope}%
\pgfpathrectangle{\pgfqpoint{1.250000in}{0.660000in}}{\pgfqpoint{7.750000in}{4.620000in}}%
\pgfusepath{clip}%
\pgfsetbuttcap%
\pgfsetmiterjoin%
\definecolor{currentfill}{rgb}{0.000000,0.000000,0.000000}%
\pgfsetfillcolor{currentfill}%
\pgfsetlinewidth{0.000000pt}%
\definecolor{currentstroke}{rgb}{0.000000,0.000000,0.000000}%
\pgfsetstrokecolor{currentstroke}%
\pgfsetstrokeopacity{0.000000}%
\pgfsetdash{}{0pt}%
\pgfpathmoveto{\pgfqpoint{2.432650in}{0.660000in}}%
\pgfpathlineto{\pgfqpoint{2.438850in}{0.660000in}}%
\pgfpathlineto{\pgfqpoint{2.438850in}{3.565126in}}%
\pgfpathlineto{\pgfqpoint{2.432650in}{3.565126in}}%
\pgfpathlineto{\pgfqpoint{2.432650in}{0.660000in}}%
\pgfpathclose%
\pgfusepath{fill}%
\end{pgfscope}%
\begin{pgfscope}%
\pgfpathrectangle{\pgfqpoint{1.250000in}{0.660000in}}{\pgfqpoint{7.750000in}{4.620000in}}%
\pgfusepath{clip}%
\pgfsetbuttcap%
\pgfsetmiterjoin%
\definecolor{currentfill}{rgb}{0.000000,0.000000,0.000000}%
\pgfsetfillcolor{currentfill}%
\pgfsetlinewidth{0.000000pt}%
\definecolor{currentstroke}{rgb}{0.000000,0.000000,0.000000}%
\pgfsetstrokecolor{currentstroke}%
\pgfsetstrokeopacity{0.000000}%
\pgfsetdash{}{0pt}%
\pgfpathmoveto{\pgfqpoint{2.440400in}{0.660000in}}%
\pgfpathlineto{\pgfqpoint{2.446600in}{0.660000in}}%
\pgfpathlineto{\pgfqpoint{2.446600in}{3.587447in}}%
\pgfpathlineto{\pgfqpoint{2.440400in}{3.587447in}}%
\pgfpathlineto{\pgfqpoint{2.440400in}{0.660000in}}%
\pgfpathclose%
\pgfusepath{fill}%
\end{pgfscope}%
\begin{pgfscope}%
\pgfpathrectangle{\pgfqpoint{1.250000in}{0.660000in}}{\pgfqpoint{7.750000in}{4.620000in}}%
\pgfusepath{clip}%
\pgfsetbuttcap%
\pgfsetmiterjoin%
\definecolor{currentfill}{rgb}{0.000000,0.000000,0.000000}%
\pgfsetfillcolor{currentfill}%
\pgfsetlinewidth{0.000000pt}%
\definecolor{currentstroke}{rgb}{0.000000,0.000000,0.000000}%
\pgfsetstrokecolor{currentstroke}%
\pgfsetstrokeopacity{0.000000}%
\pgfsetdash{}{0pt}%
\pgfpathmoveto{\pgfqpoint{2.448150in}{0.660000in}}%
\pgfpathlineto{\pgfqpoint{2.454350in}{0.660000in}}%
\pgfpathlineto{\pgfqpoint{2.454350in}{3.539519in}}%
\pgfpathlineto{\pgfqpoint{2.448150in}{3.539519in}}%
\pgfpathlineto{\pgfqpoint{2.448150in}{0.660000in}}%
\pgfpathclose%
\pgfusepath{fill}%
\end{pgfscope}%
\begin{pgfscope}%
\pgfpathrectangle{\pgfqpoint{1.250000in}{0.660000in}}{\pgfqpoint{7.750000in}{4.620000in}}%
\pgfusepath{clip}%
\pgfsetbuttcap%
\pgfsetmiterjoin%
\definecolor{currentfill}{rgb}{0.000000,0.000000,0.000000}%
\pgfsetfillcolor{currentfill}%
\pgfsetlinewidth{0.000000pt}%
\definecolor{currentstroke}{rgb}{0.000000,0.000000,0.000000}%
\pgfsetstrokecolor{currentstroke}%
\pgfsetstrokeopacity{0.000000}%
\pgfsetdash{}{0pt}%
\pgfpathmoveto{\pgfqpoint{2.455900in}{0.660000in}}%
\pgfpathlineto{\pgfqpoint{2.462100in}{0.660000in}}%
\pgfpathlineto{\pgfqpoint{2.462100in}{3.509262in}}%
\pgfpathlineto{\pgfqpoint{2.455900in}{3.509262in}}%
\pgfpathlineto{\pgfqpoint{2.455900in}{0.660000in}}%
\pgfpathclose%
\pgfusepath{fill}%
\end{pgfscope}%
\begin{pgfscope}%
\pgfpathrectangle{\pgfqpoint{1.250000in}{0.660000in}}{\pgfqpoint{7.750000in}{4.620000in}}%
\pgfusepath{clip}%
\pgfsetbuttcap%
\pgfsetmiterjoin%
\definecolor{currentfill}{rgb}{0.000000,0.000000,0.000000}%
\pgfsetfillcolor{currentfill}%
\pgfsetlinewidth{0.000000pt}%
\definecolor{currentstroke}{rgb}{0.000000,0.000000,0.000000}%
\pgfsetstrokecolor{currentstroke}%
\pgfsetstrokeopacity{0.000000}%
\pgfsetdash{}{0pt}%
\pgfpathmoveto{\pgfqpoint{2.463650in}{0.660000in}}%
\pgfpathlineto{\pgfqpoint{2.469850in}{0.660000in}}%
\pgfpathlineto{\pgfqpoint{2.469850in}{3.522094in}}%
\pgfpathlineto{\pgfqpoint{2.463650in}{3.522094in}}%
\pgfpathlineto{\pgfqpoint{2.463650in}{0.660000in}}%
\pgfpathclose%
\pgfusepath{fill}%
\end{pgfscope}%
\begin{pgfscope}%
\pgfpathrectangle{\pgfqpoint{1.250000in}{0.660000in}}{\pgfqpoint{7.750000in}{4.620000in}}%
\pgfusepath{clip}%
\pgfsetbuttcap%
\pgfsetmiterjoin%
\definecolor{currentfill}{rgb}{0.000000,0.000000,0.000000}%
\pgfsetfillcolor{currentfill}%
\pgfsetlinewidth{0.000000pt}%
\definecolor{currentstroke}{rgb}{0.000000,0.000000,0.000000}%
\pgfsetstrokecolor{currentstroke}%
\pgfsetstrokeopacity{0.000000}%
\pgfsetdash{}{0pt}%
\pgfpathmoveto{\pgfqpoint{2.471400in}{0.660000in}}%
\pgfpathlineto{\pgfqpoint{2.477600in}{0.660000in}}%
\pgfpathlineto{\pgfqpoint{2.477600in}{3.669599in}}%
\pgfpathlineto{\pgfqpoint{2.471400in}{3.669599in}}%
\pgfpathlineto{\pgfqpoint{2.471400in}{0.660000in}}%
\pgfpathclose%
\pgfusepath{fill}%
\end{pgfscope}%
\begin{pgfscope}%
\pgfpathrectangle{\pgfqpoint{1.250000in}{0.660000in}}{\pgfqpoint{7.750000in}{4.620000in}}%
\pgfusepath{clip}%
\pgfsetbuttcap%
\pgfsetmiterjoin%
\definecolor{currentfill}{rgb}{0.000000,0.000000,0.000000}%
\pgfsetfillcolor{currentfill}%
\pgfsetlinewidth{0.000000pt}%
\definecolor{currentstroke}{rgb}{0.000000,0.000000,0.000000}%
\pgfsetstrokecolor{currentstroke}%
\pgfsetstrokeopacity{0.000000}%
\pgfsetdash{}{0pt}%
\pgfpathmoveto{\pgfqpoint{2.479150in}{0.660000in}}%
\pgfpathlineto{\pgfqpoint{2.485350in}{0.660000in}}%
\pgfpathlineto{\pgfqpoint{2.485350in}{3.508178in}}%
\pgfpathlineto{\pgfqpoint{2.479150in}{3.508178in}}%
\pgfpathlineto{\pgfqpoint{2.479150in}{0.660000in}}%
\pgfpathclose%
\pgfusepath{fill}%
\end{pgfscope}%
\begin{pgfscope}%
\pgfpathrectangle{\pgfqpoint{1.250000in}{0.660000in}}{\pgfqpoint{7.750000in}{4.620000in}}%
\pgfusepath{clip}%
\pgfsetbuttcap%
\pgfsetmiterjoin%
\definecolor{currentfill}{rgb}{0.000000,0.000000,0.000000}%
\pgfsetfillcolor{currentfill}%
\pgfsetlinewidth{0.000000pt}%
\definecolor{currentstroke}{rgb}{0.000000,0.000000,0.000000}%
\pgfsetstrokecolor{currentstroke}%
\pgfsetstrokeopacity{0.000000}%
\pgfsetdash{}{0pt}%
\pgfpathmoveto{\pgfqpoint{2.486900in}{0.660000in}}%
\pgfpathlineto{\pgfqpoint{2.493100in}{0.660000in}}%
\pgfpathlineto{\pgfqpoint{2.493100in}{3.478343in}}%
\pgfpathlineto{\pgfqpoint{2.486900in}{3.478343in}}%
\pgfpathlineto{\pgfqpoint{2.486900in}{0.660000in}}%
\pgfpathclose%
\pgfusepath{fill}%
\end{pgfscope}%
\begin{pgfscope}%
\pgfpathrectangle{\pgfqpoint{1.250000in}{0.660000in}}{\pgfqpoint{7.750000in}{4.620000in}}%
\pgfusepath{clip}%
\pgfsetbuttcap%
\pgfsetmiterjoin%
\definecolor{currentfill}{rgb}{0.000000,0.000000,0.000000}%
\pgfsetfillcolor{currentfill}%
\pgfsetlinewidth{0.000000pt}%
\definecolor{currentstroke}{rgb}{0.000000,0.000000,0.000000}%
\pgfsetstrokecolor{currentstroke}%
\pgfsetstrokeopacity{0.000000}%
\pgfsetdash{}{0pt}%
\pgfpathmoveto{\pgfqpoint{2.494650in}{0.660000in}}%
\pgfpathlineto{\pgfqpoint{2.500850in}{0.660000in}}%
\pgfpathlineto{\pgfqpoint{2.500850in}{3.495792in}}%
\pgfpathlineto{\pgfqpoint{2.494650in}{3.495792in}}%
\pgfpathlineto{\pgfqpoint{2.494650in}{0.660000in}}%
\pgfpathclose%
\pgfusepath{fill}%
\end{pgfscope}%
\begin{pgfscope}%
\pgfpathrectangle{\pgfqpoint{1.250000in}{0.660000in}}{\pgfqpoint{7.750000in}{4.620000in}}%
\pgfusepath{clip}%
\pgfsetbuttcap%
\pgfsetmiterjoin%
\definecolor{currentfill}{rgb}{0.000000,0.000000,0.000000}%
\pgfsetfillcolor{currentfill}%
\pgfsetlinewidth{0.000000pt}%
\definecolor{currentstroke}{rgb}{0.000000,0.000000,0.000000}%
\pgfsetstrokecolor{currentstroke}%
\pgfsetstrokeopacity{0.000000}%
\pgfsetdash{}{0pt}%
\pgfpathmoveto{\pgfqpoint{2.502400in}{0.660000in}}%
\pgfpathlineto{\pgfqpoint{2.508600in}{0.660000in}}%
\pgfpathlineto{\pgfqpoint{2.508600in}{3.475813in}}%
\pgfpathlineto{\pgfqpoint{2.502400in}{3.475813in}}%
\pgfpathlineto{\pgfqpoint{2.502400in}{0.660000in}}%
\pgfpathclose%
\pgfusepath{fill}%
\end{pgfscope}%
\begin{pgfscope}%
\pgfpathrectangle{\pgfqpoint{1.250000in}{0.660000in}}{\pgfqpoint{7.750000in}{4.620000in}}%
\pgfusepath{clip}%
\pgfsetbuttcap%
\pgfsetmiterjoin%
\definecolor{currentfill}{rgb}{0.000000,0.000000,0.000000}%
\pgfsetfillcolor{currentfill}%
\pgfsetlinewidth{0.000000pt}%
\definecolor{currentstroke}{rgb}{0.000000,0.000000,0.000000}%
\pgfsetstrokecolor{currentstroke}%
\pgfsetstrokeopacity{0.000000}%
\pgfsetdash{}{0pt}%
\pgfpathmoveto{\pgfqpoint{2.510150in}{0.660000in}}%
\pgfpathlineto{\pgfqpoint{2.516350in}{0.660000in}}%
\pgfpathlineto{\pgfqpoint{2.516350in}{3.469654in}}%
\pgfpathlineto{\pgfqpoint{2.510150in}{3.469654in}}%
\pgfpathlineto{\pgfqpoint{2.510150in}{0.660000in}}%
\pgfpathclose%
\pgfusepath{fill}%
\end{pgfscope}%
\begin{pgfscope}%
\pgfpathrectangle{\pgfqpoint{1.250000in}{0.660000in}}{\pgfqpoint{7.750000in}{4.620000in}}%
\pgfusepath{clip}%
\pgfsetbuttcap%
\pgfsetmiterjoin%
\definecolor{currentfill}{rgb}{0.000000,0.000000,0.000000}%
\pgfsetfillcolor{currentfill}%
\pgfsetlinewidth{0.000000pt}%
\definecolor{currentstroke}{rgb}{0.000000,0.000000,0.000000}%
\pgfsetstrokecolor{currentstroke}%
\pgfsetstrokeopacity{0.000000}%
\pgfsetdash{}{0pt}%
\pgfpathmoveto{\pgfqpoint{2.517900in}{0.660000in}}%
\pgfpathlineto{\pgfqpoint{2.524100in}{0.660000in}}%
\pgfpathlineto{\pgfqpoint{2.524100in}{3.469093in}}%
\pgfpathlineto{\pgfqpoint{2.517900in}{3.469093in}}%
\pgfpathlineto{\pgfqpoint{2.517900in}{0.660000in}}%
\pgfpathclose%
\pgfusepath{fill}%
\end{pgfscope}%
\begin{pgfscope}%
\pgfpathrectangle{\pgfqpoint{1.250000in}{0.660000in}}{\pgfqpoint{7.750000in}{4.620000in}}%
\pgfusepath{clip}%
\pgfsetbuttcap%
\pgfsetmiterjoin%
\definecolor{currentfill}{rgb}{0.000000,0.000000,0.000000}%
\pgfsetfillcolor{currentfill}%
\pgfsetlinewidth{0.000000pt}%
\definecolor{currentstroke}{rgb}{0.000000,0.000000,0.000000}%
\pgfsetstrokecolor{currentstroke}%
\pgfsetstrokeopacity{0.000000}%
\pgfsetdash{}{0pt}%
\pgfpathmoveto{\pgfqpoint{2.525650in}{0.660000in}}%
\pgfpathlineto{\pgfqpoint{2.531850in}{0.660000in}}%
\pgfpathlineto{\pgfqpoint{2.531850in}{3.513132in}}%
\pgfpathlineto{\pgfqpoint{2.525650in}{3.513132in}}%
\pgfpathlineto{\pgfqpoint{2.525650in}{0.660000in}}%
\pgfpathclose%
\pgfusepath{fill}%
\end{pgfscope}%
\begin{pgfscope}%
\pgfpathrectangle{\pgfqpoint{1.250000in}{0.660000in}}{\pgfqpoint{7.750000in}{4.620000in}}%
\pgfusepath{clip}%
\pgfsetbuttcap%
\pgfsetmiterjoin%
\definecolor{currentfill}{rgb}{0.000000,0.000000,0.000000}%
\pgfsetfillcolor{currentfill}%
\pgfsetlinewidth{0.000000pt}%
\definecolor{currentstroke}{rgb}{0.000000,0.000000,0.000000}%
\pgfsetstrokecolor{currentstroke}%
\pgfsetstrokeopacity{0.000000}%
\pgfsetdash{}{0pt}%
\pgfpathmoveto{\pgfqpoint{2.533400in}{0.660000in}}%
\pgfpathlineto{\pgfqpoint{2.539600in}{0.660000in}}%
\pgfpathlineto{\pgfqpoint{2.539600in}{3.507619in}}%
\pgfpathlineto{\pgfqpoint{2.533400in}{3.507619in}}%
\pgfpathlineto{\pgfqpoint{2.533400in}{0.660000in}}%
\pgfpathclose%
\pgfusepath{fill}%
\end{pgfscope}%
\begin{pgfscope}%
\pgfpathrectangle{\pgfqpoint{1.250000in}{0.660000in}}{\pgfqpoint{7.750000in}{4.620000in}}%
\pgfusepath{clip}%
\pgfsetbuttcap%
\pgfsetmiterjoin%
\definecolor{currentfill}{rgb}{0.000000,0.000000,0.000000}%
\pgfsetfillcolor{currentfill}%
\pgfsetlinewidth{0.000000pt}%
\definecolor{currentstroke}{rgb}{0.000000,0.000000,0.000000}%
\pgfsetstrokecolor{currentstroke}%
\pgfsetstrokeopacity{0.000000}%
\pgfsetdash{}{0pt}%
\pgfpathmoveto{\pgfqpoint{2.541150in}{0.660000in}}%
\pgfpathlineto{\pgfqpoint{2.547350in}{0.660000in}}%
\pgfpathlineto{\pgfqpoint{2.547350in}{3.490811in}}%
\pgfpathlineto{\pgfqpoint{2.541150in}{3.490811in}}%
\pgfpathlineto{\pgfqpoint{2.541150in}{0.660000in}}%
\pgfpathclose%
\pgfusepath{fill}%
\end{pgfscope}%
\begin{pgfscope}%
\pgfpathrectangle{\pgfqpoint{1.250000in}{0.660000in}}{\pgfqpoint{7.750000in}{4.620000in}}%
\pgfusepath{clip}%
\pgfsetbuttcap%
\pgfsetmiterjoin%
\definecolor{currentfill}{rgb}{0.000000,0.000000,0.000000}%
\pgfsetfillcolor{currentfill}%
\pgfsetlinewidth{0.000000pt}%
\definecolor{currentstroke}{rgb}{0.000000,0.000000,0.000000}%
\pgfsetstrokecolor{currentstroke}%
\pgfsetstrokeopacity{0.000000}%
\pgfsetdash{}{0pt}%
\pgfpathmoveto{\pgfqpoint{2.548900in}{0.660000in}}%
\pgfpathlineto{\pgfqpoint{2.555100in}{0.660000in}}%
\pgfpathlineto{\pgfqpoint{2.555100in}{3.517505in}}%
\pgfpathlineto{\pgfqpoint{2.548900in}{3.517505in}}%
\pgfpathlineto{\pgfqpoint{2.548900in}{0.660000in}}%
\pgfpathclose%
\pgfusepath{fill}%
\end{pgfscope}%
\begin{pgfscope}%
\pgfpathrectangle{\pgfqpoint{1.250000in}{0.660000in}}{\pgfqpoint{7.750000in}{4.620000in}}%
\pgfusepath{clip}%
\pgfsetbuttcap%
\pgfsetmiterjoin%
\definecolor{currentfill}{rgb}{0.000000,0.000000,0.000000}%
\pgfsetfillcolor{currentfill}%
\pgfsetlinewidth{0.000000pt}%
\definecolor{currentstroke}{rgb}{0.000000,0.000000,0.000000}%
\pgfsetstrokecolor{currentstroke}%
\pgfsetstrokeopacity{0.000000}%
\pgfsetdash{}{0pt}%
\pgfpathmoveto{\pgfqpoint{2.556650in}{0.660000in}}%
\pgfpathlineto{\pgfqpoint{2.562850in}{0.660000in}}%
\pgfpathlineto{\pgfqpoint{2.562850in}{3.474310in}}%
\pgfpathlineto{\pgfqpoint{2.556650in}{3.474310in}}%
\pgfpathlineto{\pgfqpoint{2.556650in}{0.660000in}}%
\pgfpathclose%
\pgfusepath{fill}%
\end{pgfscope}%
\begin{pgfscope}%
\pgfpathrectangle{\pgfqpoint{1.250000in}{0.660000in}}{\pgfqpoint{7.750000in}{4.620000in}}%
\pgfusepath{clip}%
\pgfsetbuttcap%
\pgfsetmiterjoin%
\definecolor{currentfill}{rgb}{0.000000,0.000000,0.000000}%
\pgfsetfillcolor{currentfill}%
\pgfsetlinewidth{0.000000pt}%
\definecolor{currentstroke}{rgb}{0.000000,0.000000,0.000000}%
\pgfsetstrokecolor{currentstroke}%
\pgfsetstrokeopacity{0.000000}%
\pgfsetdash{}{0pt}%
\pgfpathmoveto{\pgfqpoint{2.564400in}{0.660000in}}%
\pgfpathlineto{\pgfqpoint{2.570600in}{0.660000in}}%
\pgfpathlineto{\pgfqpoint{2.570600in}{3.478039in}}%
\pgfpathlineto{\pgfqpoint{2.564400in}{3.478039in}}%
\pgfpathlineto{\pgfqpoint{2.564400in}{0.660000in}}%
\pgfpathclose%
\pgfusepath{fill}%
\end{pgfscope}%
\begin{pgfscope}%
\pgfpathrectangle{\pgfqpoint{1.250000in}{0.660000in}}{\pgfqpoint{7.750000in}{4.620000in}}%
\pgfusepath{clip}%
\pgfsetbuttcap%
\pgfsetmiterjoin%
\definecolor{currentfill}{rgb}{0.000000,0.000000,0.000000}%
\pgfsetfillcolor{currentfill}%
\pgfsetlinewidth{0.000000pt}%
\definecolor{currentstroke}{rgb}{0.000000,0.000000,0.000000}%
\pgfsetstrokecolor{currentstroke}%
\pgfsetstrokeopacity{0.000000}%
\pgfsetdash{}{0pt}%
\pgfpathmoveto{\pgfqpoint{2.572150in}{0.660000in}}%
\pgfpathlineto{\pgfqpoint{2.578350in}{0.660000in}}%
\pgfpathlineto{\pgfqpoint{2.578350in}{3.477657in}}%
\pgfpathlineto{\pgfqpoint{2.572150in}{3.477657in}}%
\pgfpathlineto{\pgfqpoint{2.572150in}{0.660000in}}%
\pgfpathclose%
\pgfusepath{fill}%
\end{pgfscope}%
\begin{pgfscope}%
\pgfpathrectangle{\pgfqpoint{1.250000in}{0.660000in}}{\pgfqpoint{7.750000in}{4.620000in}}%
\pgfusepath{clip}%
\pgfsetbuttcap%
\pgfsetmiterjoin%
\definecolor{currentfill}{rgb}{0.000000,0.000000,0.000000}%
\pgfsetfillcolor{currentfill}%
\pgfsetlinewidth{0.000000pt}%
\definecolor{currentstroke}{rgb}{0.000000,0.000000,0.000000}%
\pgfsetstrokecolor{currentstroke}%
\pgfsetstrokeopacity{0.000000}%
\pgfsetdash{}{0pt}%
\pgfpathmoveto{\pgfqpoint{2.579900in}{0.660000in}}%
\pgfpathlineto{\pgfqpoint{2.586100in}{0.660000in}}%
\pgfpathlineto{\pgfqpoint{2.586100in}{3.511480in}}%
\pgfpathlineto{\pgfqpoint{2.579900in}{3.511480in}}%
\pgfpathlineto{\pgfqpoint{2.579900in}{0.660000in}}%
\pgfpathclose%
\pgfusepath{fill}%
\end{pgfscope}%
\begin{pgfscope}%
\pgfpathrectangle{\pgfqpoint{1.250000in}{0.660000in}}{\pgfqpoint{7.750000in}{4.620000in}}%
\pgfusepath{clip}%
\pgfsetbuttcap%
\pgfsetmiterjoin%
\definecolor{currentfill}{rgb}{0.000000,0.000000,0.000000}%
\pgfsetfillcolor{currentfill}%
\pgfsetlinewidth{0.000000pt}%
\definecolor{currentstroke}{rgb}{0.000000,0.000000,0.000000}%
\pgfsetstrokecolor{currentstroke}%
\pgfsetstrokeopacity{0.000000}%
\pgfsetdash{}{0pt}%
\pgfpathmoveto{\pgfqpoint{2.587650in}{0.660000in}}%
\pgfpathlineto{\pgfqpoint{2.593850in}{0.660000in}}%
\pgfpathlineto{\pgfqpoint{2.593850in}{3.567718in}}%
\pgfpathlineto{\pgfqpoint{2.587650in}{3.567718in}}%
\pgfpathlineto{\pgfqpoint{2.587650in}{0.660000in}}%
\pgfpathclose%
\pgfusepath{fill}%
\end{pgfscope}%
\begin{pgfscope}%
\pgfpathrectangle{\pgfqpoint{1.250000in}{0.660000in}}{\pgfqpoint{7.750000in}{4.620000in}}%
\pgfusepath{clip}%
\pgfsetbuttcap%
\pgfsetmiterjoin%
\definecolor{currentfill}{rgb}{0.000000,0.000000,0.000000}%
\pgfsetfillcolor{currentfill}%
\pgfsetlinewidth{0.000000pt}%
\definecolor{currentstroke}{rgb}{0.000000,0.000000,0.000000}%
\pgfsetstrokecolor{currentstroke}%
\pgfsetstrokeopacity{0.000000}%
\pgfsetdash{}{0pt}%
\pgfpathmoveto{\pgfqpoint{2.595400in}{0.660000in}}%
\pgfpathlineto{\pgfqpoint{2.601600in}{0.660000in}}%
\pgfpathlineto{\pgfqpoint{2.601600in}{3.502597in}}%
\pgfpathlineto{\pgfqpoint{2.595400in}{3.502597in}}%
\pgfpathlineto{\pgfqpoint{2.595400in}{0.660000in}}%
\pgfpathclose%
\pgfusepath{fill}%
\end{pgfscope}%
\begin{pgfscope}%
\pgfpathrectangle{\pgfqpoint{1.250000in}{0.660000in}}{\pgfqpoint{7.750000in}{4.620000in}}%
\pgfusepath{clip}%
\pgfsetbuttcap%
\pgfsetmiterjoin%
\definecolor{currentfill}{rgb}{0.000000,0.000000,0.000000}%
\pgfsetfillcolor{currentfill}%
\pgfsetlinewidth{0.000000pt}%
\definecolor{currentstroke}{rgb}{0.000000,0.000000,0.000000}%
\pgfsetstrokecolor{currentstroke}%
\pgfsetstrokeopacity{0.000000}%
\pgfsetdash{}{0pt}%
\pgfpathmoveto{\pgfqpoint{2.603150in}{0.660000in}}%
\pgfpathlineto{\pgfqpoint{2.609350in}{0.660000in}}%
\pgfpathlineto{\pgfqpoint{2.609350in}{3.481795in}}%
\pgfpathlineto{\pgfqpoint{2.603150in}{3.481795in}}%
\pgfpathlineto{\pgfqpoint{2.603150in}{0.660000in}}%
\pgfpathclose%
\pgfusepath{fill}%
\end{pgfscope}%
\begin{pgfscope}%
\pgfpathrectangle{\pgfqpoint{1.250000in}{0.660000in}}{\pgfqpoint{7.750000in}{4.620000in}}%
\pgfusepath{clip}%
\pgfsetbuttcap%
\pgfsetmiterjoin%
\definecolor{currentfill}{rgb}{0.000000,0.000000,0.000000}%
\pgfsetfillcolor{currentfill}%
\pgfsetlinewidth{0.000000pt}%
\definecolor{currentstroke}{rgb}{0.000000,0.000000,0.000000}%
\pgfsetstrokecolor{currentstroke}%
\pgfsetstrokeopacity{0.000000}%
\pgfsetdash{}{0pt}%
\pgfpathmoveto{\pgfqpoint{2.610900in}{0.660000in}}%
\pgfpathlineto{\pgfqpoint{2.617100in}{0.660000in}}%
\pgfpathlineto{\pgfqpoint{2.617100in}{3.502943in}}%
\pgfpathlineto{\pgfqpoint{2.610900in}{3.502943in}}%
\pgfpathlineto{\pgfqpoint{2.610900in}{0.660000in}}%
\pgfpathclose%
\pgfusepath{fill}%
\end{pgfscope}%
\begin{pgfscope}%
\pgfpathrectangle{\pgfqpoint{1.250000in}{0.660000in}}{\pgfqpoint{7.750000in}{4.620000in}}%
\pgfusepath{clip}%
\pgfsetbuttcap%
\pgfsetmiterjoin%
\definecolor{currentfill}{rgb}{0.000000,0.000000,0.000000}%
\pgfsetfillcolor{currentfill}%
\pgfsetlinewidth{0.000000pt}%
\definecolor{currentstroke}{rgb}{0.000000,0.000000,0.000000}%
\pgfsetstrokecolor{currentstroke}%
\pgfsetstrokeopacity{0.000000}%
\pgfsetdash{}{0pt}%
\pgfpathmoveto{\pgfqpoint{2.618650in}{0.660000in}}%
\pgfpathlineto{\pgfqpoint{2.624850in}{0.660000in}}%
\pgfpathlineto{\pgfqpoint{2.624850in}{3.477541in}}%
\pgfpathlineto{\pgfqpoint{2.618650in}{3.477541in}}%
\pgfpathlineto{\pgfqpoint{2.618650in}{0.660000in}}%
\pgfpathclose%
\pgfusepath{fill}%
\end{pgfscope}%
\begin{pgfscope}%
\pgfpathrectangle{\pgfqpoint{1.250000in}{0.660000in}}{\pgfqpoint{7.750000in}{4.620000in}}%
\pgfusepath{clip}%
\pgfsetbuttcap%
\pgfsetmiterjoin%
\definecolor{currentfill}{rgb}{0.000000,0.000000,0.000000}%
\pgfsetfillcolor{currentfill}%
\pgfsetlinewidth{0.000000pt}%
\definecolor{currentstroke}{rgb}{0.000000,0.000000,0.000000}%
\pgfsetstrokecolor{currentstroke}%
\pgfsetstrokeopacity{0.000000}%
\pgfsetdash{}{0pt}%
\pgfpathmoveto{\pgfqpoint{2.626400in}{0.660000in}}%
\pgfpathlineto{\pgfqpoint{2.632600in}{0.660000in}}%
\pgfpathlineto{\pgfqpoint{2.632600in}{3.485706in}}%
\pgfpathlineto{\pgfqpoint{2.626400in}{3.485706in}}%
\pgfpathlineto{\pgfqpoint{2.626400in}{0.660000in}}%
\pgfpathclose%
\pgfusepath{fill}%
\end{pgfscope}%
\begin{pgfscope}%
\pgfpathrectangle{\pgfqpoint{1.250000in}{0.660000in}}{\pgfqpoint{7.750000in}{4.620000in}}%
\pgfusepath{clip}%
\pgfsetbuttcap%
\pgfsetmiterjoin%
\definecolor{currentfill}{rgb}{0.000000,0.000000,0.000000}%
\pgfsetfillcolor{currentfill}%
\pgfsetlinewidth{0.000000pt}%
\definecolor{currentstroke}{rgb}{0.000000,0.000000,0.000000}%
\pgfsetstrokecolor{currentstroke}%
\pgfsetstrokeopacity{0.000000}%
\pgfsetdash{}{0pt}%
\pgfpathmoveto{\pgfqpoint{2.634150in}{0.660000in}}%
\pgfpathlineto{\pgfqpoint{2.640350in}{0.660000in}}%
\pgfpathlineto{\pgfqpoint{2.640350in}{3.548326in}}%
\pgfpathlineto{\pgfqpoint{2.634150in}{3.548326in}}%
\pgfpathlineto{\pgfqpoint{2.634150in}{0.660000in}}%
\pgfpathclose%
\pgfusepath{fill}%
\end{pgfscope}%
\begin{pgfscope}%
\pgfpathrectangle{\pgfqpoint{1.250000in}{0.660000in}}{\pgfqpoint{7.750000in}{4.620000in}}%
\pgfusepath{clip}%
\pgfsetbuttcap%
\pgfsetmiterjoin%
\definecolor{currentfill}{rgb}{0.000000,0.000000,0.000000}%
\pgfsetfillcolor{currentfill}%
\pgfsetlinewidth{0.000000pt}%
\definecolor{currentstroke}{rgb}{0.000000,0.000000,0.000000}%
\pgfsetstrokecolor{currentstroke}%
\pgfsetstrokeopacity{0.000000}%
\pgfsetdash{}{0pt}%
\pgfpathmoveto{\pgfqpoint{2.641900in}{0.660000in}}%
\pgfpathlineto{\pgfqpoint{2.648100in}{0.660000in}}%
\pgfpathlineto{\pgfqpoint{2.648100in}{3.552123in}}%
\pgfpathlineto{\pgfqpoint{2.641900in}{3.552123in}}%
\pgfpathlineto{\pgfqpoint{2.641900in}{0.660000in}}%
\pgfpathclose%
\pgfusepath{fill}%
\end{pgfscope}%
\begin{pgfscope}%
\pgfpathrectangle{\pgfqpoint{1.250000in}{0.660000in}}{\pgfqpoint{7.750000in}{4.620000in}}%
\pgfusepath{clip}%
\pgfsetbuttcap%
\pgfsetmiterjoin%
\definecolor{currentfill}{rgb}{0.000000,0.000000,0.000000}%
\pgfsetfillcolor{currentfill}%
\pgfsetlinewidth{0.000000pt}%
\definecolor{currentstroke}{rgb}{0.000000,0.000000,0.000000}%
\pgfsetstrokecolor{currentstroke}%
\pgfsetstrokeopacity{0.000000}%
\pgfsetdash{}{0pt}%
\pgfpathmoveto{\pgfqpoint{2.649650in}{0.660000in}}%
\pgfpathlineto{\pgfqpoint{2.655850in}{0.660000in}}%
\pgfpathlineto{\pgfqpoint{2.655850in}{3.482348in}}%
\pgfpathlineto{\pgfqpoint{2.649650in}{3.482348in}}%
\pgfpathlineto{\pgfqpoint{2.649650in}{0.660000in}}%
\pgfpathclose%
\pgfusepath{fill}%
\end{pgfscope}%
\begin{pgfscope}%
\pgfpathrectangle{\pgfqpoint{1.250000in}{0.660000in}}{\pgfqpoint{7.750000in}{4.620000in}}%
\pgfusepath{clip}%
\pgfsetbuttcap%
\pgfsetmiterjoin%
\definecolor{currentfill}{rgb}{0.000000,0.000000,0.000000}%
\pgfsetfillcolor{currentfill}%
\pgfsetlinewidth{0.000000pt}%
\definecolor{currentstroke}{rgb}{0.000000,0.000000,0.000000}%
\pgfsetstrokecolor{currentstroke}%
\pgfsetstrokeopacity{0.000000}%
\pgfsetdash{}{0pt}%
\pgfpathmoveto{\pgfqpoint{2.657400in}{0.660000in}}%
\pgfpathlineto{\pgfqpoint{2.663600in}{0.660000in}}%
\pgfpathlineto{\pgfqpoint{2.663600in}{3.489989in}}%
\pgfpathlineto{\pgfqpoint{2.657400in}{3.489989in}}%
\pgfpathlineto{\pgfqpoint{2.657400in}{0.660000in}}%
\pgfpathclose%
\pgfusepath{fill}%
\end{pgfscope}%
\begin{pgfscope}%
\pgfpathrectangle{\pgfqpoint{1.250000in}{0.660000in}}{\pgfqpoint{7.750000in}{4.620000in}}%
\pgfusepath{clip}%
\pgfsetbuttcap%
\pgfsetmiterjoin%
\definecolor{currentfill}{rgb}{0.000000,0.000000,0.000000}%
\pgfsetfillcolor{currentfill}%
\pgfsetlinewidth{0.000000pt}%
\definecolor{currentstroke}{rgb}{0.000000,0.000000,0.000000}%
\pgfsetstrokecolor{currentstroke}%
\pgfsetstrokeopacity{0.000000}%
\pgfsetdash{}{0pt}%
\pgfpathmoveto{\pgfqpoint{2.665150in}{0.660000in}}%
\pgfpathlineto{\pgfqpoint{2.671350in}{0.660000in}}%
\pgfpathlineto{\pgfqpoint{2.671350in}{3.445019in}}%
\pgfpathlineto{\pgfqpoint{2.665150in}{3.445019in}}%
\pgfpathlineto{\pgfqpoint{2.665150in}{0.660000in}}%
\pgfpathclose%
\pgfusepath{fill}%
\end{pgfscope}%
\begin{pgfscope}%
\pgfpathrectangle{\pgfqpoint{1.250000in}{0.660000in}}{\pgfqpoint{7.750000in}{4.620000in}}%
\pgfusepath{clip}%
\pgfsetbuttcap%
\pgfsetmiterjoin%
\definecolor{currentfill}{rgb}{0.000000,0.000000,0.000000}%
\pgfsetfillcolor{currentfill}%
\pgfsetlinewidth{0.000000pt}%
\definecolor{currentstroke}{rgb}{0.000000,0.000000,0.000000}%
\pgfsetstrokecolor{currentstroke}%
\pgfsetstrokeopacity{0.000000}%
\pgfsetdash{}{0pt}%
\pgfpathmoveto{\pgfqpoint{2.672900in}{0.660000in}}%
\pgfpathlineto{\pgfqpoint{2.679100in}{0.660000in}}%
\pgfpathlineto{\pgfqpoint{2.679100in}{3.440088in}}%
\pgfpathlineto{\pgfqpoint{2.672900in}{3.440088in}}%
\pgfpathlineto{\pgfqpoint{2.672900in}{0.660000in}}%
\pgfpathclose%
\pgfusepath{fill}%
\end{pgfscope}%
\begin{pgfscope}%
\pgfpathrectangle{\pgfqpoint{1.250000in}{0.660000in}}{\pgfqpoint{7.750000in}{4.620000in}}%
\pgfusepath{clip}%
\pgfsetbuttcap%
\pgfsetmiterjoin%
\definecolor{currentfill}{rgb}{0.000000,0.000000,0.000000}%
\pgfsetfillcolor{currentfill}%
\pgfsetlinewidth{0.000000pt}%
\definecolor{currentstroke}{rgb}{0.000000,0.000000,0.000000}%
\pgfsetstrokecolor{currentstroke}%
\pgfsetstrokeopacity{0.000000}%
\pgfsetdash{}{0pt}%
\pgfpathmoveto{\pgfqpoint{2.680650in}{0.660000in}}%
\pgfpathlineto{\pgfqpoint{2.686850in}{0.660000in}}%
\pgfpathlineto{\pgfqpoint{2.686850in}{3.539865in}}%
\pgfpathlineto{\pgfqpoint{2.680650in}{3.539865in}}%
\pgfpathlineto{\pgfqpoint{2.680650in}{0.660000in}}%
\pgfpathclose%
\pgfusepath{fill}%
\end{pgfscope}%
\begin{pgfscope}%
\pgfpathrectangle{\pgfqpoint{1.250000in}{0.660000in}}{\pgfqpoint{7.750000in}{4.620000in}}%
\pgfusepath{clip}%
\pgfsetbuttcap%
\pgfsetmiterjoin%
\definecolor{currentfill}{rgb}{0.000000,0.000000,0.000000}%
\pgfsetfillcolor{currentfill}%
\pgfsetlinewidth{0.000000pt}%
\definecolor{currentstroke}{rgb}{0.000000,0.000000,0.000000}%
\pgfsetstrokecolor{currentstroke}%
\pgfsetstrokeopacity{0.000000}%
\pgfsetdash{}{0pt}%
\pgfpathmoveto{\pgfqpoint{2.688400in}{0.660000in}}%
\pgfpathlineto{\pgfqpoint{2.694600in}{0.660000in}}%
\pgfpathlineto{\pgfqpoint{2.694600in}{3.486418in}}%
\pgfpathlineto{\pgfqpoint{2.688400in}{3.486418in}}%
\pgfpathlineto{\pgfqpoint{2.688400in}{0.660000in}}%
\pgfpathclose%
\pgfusepath{fill}%
\end{pgfscope}%
\begin{pgfscope}%
\pgfpathrectangle{\pgfqpoint{1.250000in}{0.660000in}}{\pgfqpoint{7.750000in}{4.620000in}}%
\pgfusepath{clip}%
\pgfsetbuttcap%
\pgfsetmiterjoin%
\definecolor{currentfill}{rgb}{0.000000,0.000000,0.000000}%
\pgfsetfillcolor{currentfill}%
\pgfsetlinewidth{0.000000pt}%
\definecolor{currentstroke}{rgb}{0.000000,0.000000,0.000000}%
\pgfsetstrokecolor{currentstroke}%
\pgfsetstrokeopacity{0.000000}%
\pgfsetdash{}{0pt}%
\pgfpathmoveto{\pgfqpoint{2.696150in}{0.660000in}}%
\pgfpathlineto{\pgfqpoint{2.702350in}{0.660000in}}%
\pgfpathlineto{\pgfqpoint{2.702350in}{3.466225in}}%
\pgfpathlineto{\pgfqpoint{2.696150in}{3.466225in}}%
\pgfpathlineto{\pgfqpoint{2.696150in}{0.660000in}}%
\pgfpathclose%
\pgfusepath{fill}%
\end{pgfscope}%
\begin{pgfscope}%
\pgfpathrectangle{\pgfqpoint{1.250000in}{0.660000in}}{\pgfqpoint{7.750000in}{4.620000in}}%
\pgfusepath{clip}%
\pgfsetbuttcap%
\pgfsetmiterjoin%
\definecolor{currentfill}{rgb}{0.000000,0.000000,0.000000}%
\pgfsetfillcolor{currentfill}%
\pgfsetlinewidth{0.000000pt}%
\definecolor{currentstroke}{rgb}{0.000000,0.000000,0.000000}%
\pgfsetstrokecolor{currentstroke}%
\pgfsetstrokeopacity{0.000000}%
\pgfsetdash{}{0pt}%
\pgfpathmoveto{\pgfqpoint{2.703900in}{0.660000in}}%
\pgfpathlineto{\pgfqpoint{2.710100in}{0.660000in}}%
\pgfpathlineto{\pgfqpoint{2.710100in}{3.432701in}}%
\pgfpathlineto{\pgfqpoint{2.703900in}{3.432701in}}%
\pgfpathlineto{\pgfqpoint{2.703900in}{0.660000in}}%
\pgfpathclose%
\pgfusepath{fill}%
\end{pgfscope}%
\begin{pgfscope}%
\pgfpathrectangle{\pgfqpoint{1.250000in}{0.660000in}}{\pgfqpoint{7.750000in}{4.620000in}}%
\pgfusepath{clip}%
\pgfsetbuttcap%
\pgfsetmiterjoin%
\definecolor{currentfill}{rgb}{0.000000,0.000000,0.000000}%
\pgfsetfillcolor{currentfill}%
\pgfsetlinewidth{0.000000pt}%
\definecolor{currentstroke}{rgb}{0.000000,0.000000,0.000000}%
\pgfsetstrokecolor{currentstroke}%
\pgfsetstrokeopacity{0.000000}%
\pgfsetdash{}{0pt}%
\pgfpathmoveto{\pgfqpoint{2.711650in}{0.660000in}}%
\pgfpathlineto{\pgfqpoint{2.717850in}{0.660000in}}%
\pgfpathlineto{\pgfqpoint{2.717850in}{3.425686in}}%
\pgfpathlineto{\pgfqpoint{2.711650in}{3.425686in}}%
\pgfpathlineto{\pgfqpoint{2.711650in}{0.660000in}}%
\pgfpathclose%
\pgfusepath{fill}%
\end{pgfscope}%
\begin{pgfscope}%
\pgfpathrectangle{\pgfqpoint{1.250000in}{0.660000in}}{\pgfqpoint{7.750000in}{4.620000in}}%
\pgfusepath{clip}%
\pgfsetbuttcap%
\pgfsetmiterjoin%
\definecolor{currentfill}{rgb}{0.000000,0.000000,0.000000}%
\pgfsetfillcolor{currentfill}%
\pgfsetlinewidth{0.000000pt}%
\definecolor{currentstroke}{rgb}{0.000000,0.000000,0.000000}%
\pgfsetstrokecolor{currentstroke}%
\pgfsetstrokeopacity{0.000000}%
\pgfsetdash{}{0pt}%
\pgfpathmoveto{\pgfqpoint{2.719400in}{0.660000in}}%
\pgfpathlineto{\pgfqpoint{2.725600in}{0.660000in}}%
\pgfpathlineto{\pgfqpoint{2.725600in}{3.506791in}}%
\pgfpathlineto{\pgfqpoint{2.719400in}{3.506791in}}%
\pgfpathlineto{\pgfqpoint{2.719400in}{0.660000in}}%
\pgfpathclose%
\pgfusepath{fill}%
\end{pgfscope}%
\begin{pgfscope}%
\pgfpathrectangle{\pgfqpoint{1.250000in}{0.660000in}}{\pgfqpoint{7.750000in}{4.620000in}}%
\pgfusepath{clip}%
\pgfsetbuttcap%
\pgfsetmiterjoin%
\definecolor{currentfill}{rgb}{0.000000,0.000000,0.000000}%
\pgfsetfillcolor{currentfill}%
\pgfsetlinewidth{0.000000pt}%
\definecolor{currentstroke}{rgb}{0.000000,0.000000,0.000000}%
\pgfsetstrokecolor{currentstroke}%
\pgfsetstrokeopacity{0.000000}%
\pgfsetdash{}{0pt}%
\pgfpathmoveto{\pgfqpoint{2.727150in}{0.660000in}}%
\pgfpathlineto{\pgfqpoint{2.733350in}{0.660000in}}%
\pgfpathlineto{\pgfqpoint{2.733350in}{3.566093in}}%
\pgfpathlineto{\pgfqpoint{2.727150in}{3.566093in}}%
\pgfpathlineto{\pgfqpoint{2.727150in}{0.660000in}}%
\pgfpathclose%
\pgfusepath{fill}%
\end{pgfscope}%
\begin{pgfscope}%
\pgfpathrectangle{\pgfqpoint{1.250000in}{0.660000in}}{\pgfqpoint{7.750000in}{4.620000in}}%
\pgfusepath{clip}%
\pgfsetbuttcap%
\pgfsetmiterjoin%
\definecolor{currentfill}{rgb}{0.000000,0.000000,0.000000}%
\pgfsetfillcolor{currentfill}%
\pgfsetlinewidth{0.000000pt}%
\definecolor{currentstroke}{rgb}{0.000000,0.000000,0.000000}%
\pgfsetstrokecolor{currentstroke}%
\pgfsetstrokeopacity{0.000000}%
\pgfsetdash{}{0pt}%
\pgfpathmoveto{\pgfqpoint{2.734900in}{0.660000in}}%
\pgfpathlineto{\pgfqpoint{2.741100in}{0.660000in}}%
\pgfpathlineto{\pgfqpoint{2.741100in}{3.439549in}}%
\pgfpathlineto{\pgfqpoint{2.734900in}{3.439549in}}%
\pgfpathlineto{\pgfqpoint{2.734900in}{0.660000in}}%
\pgfpathclose%
\pgfusepath{fill}%
\end{pgfscope}%
\begin{pgfscope}%
\pgfpathrectangle{\pgfqpoint{1.250000in}{0.660000in}}{\pgfqpoint{7.750000in}{4.620000in}}%
\pgfusepath{clip}%
\pgfsetbuttcap%
\pgfsetmiterjoin%
\definecolor{currentfill}{rgb}{0.000000,0.000000,0.000000}%
\pgfsetfillcolor{currentfill}%
\pgfsetlinewidth{0.000000pt}%
\definecolor{currentstroke}{rgb}{0.000000,0.000000,0.000000}%
\pgfsetstrokecolor{currentstroke}%
\pgfsetstrokeopacity{0.000000}%
\pgfsetdash{}{0pt}%
\pgfpathmoveto{\pgfqpoint{2.742650in}{0.660000in}}%
\pgfpathlineto{\pgfqpoint{2.748850in}{0.660000in}}%
\pgfpathlineto{\pgfqpoint{2.748850in}{3.443327in}}%
\pgfpathlineto{\pgfqpoint{2.742650in}{3.443327in}}%
\pgfpathlineto{\pgfqpoint{2.742650in}{0.660000in}}%
\pgfpathclose%
\pgfusepath{fill}%
\end{pgfscope}%
\begin{pgfscope}%
\pgfpathrectangle{\pgfqpoint{1.250000in}{0.660000in}}{\pgfqpoint{7.750000in}{4.620000in}}%
\pgfusepath{clip}%
\pgfsetbuttcap%
\pgfsetmiterjoin%
\definecolor{currentfill}{rgb}{0.000000,0.000000,0.000000}%
\pgfsetfillcolor{currentfill}%
\pgfsetlinewidth{0.000000pt}%
\definecolor{currentstroke}{rgb}{0.000000,0.000000,0.000000}%
\pgfsetstrokecolor{currentstroke}%
\pgfsetstrokeopacity{0.000000}%
\pgfsetdash{}{0pt}%
\pgfpathmoveto{\pgfqpoint{2.750400in}{0.660000in}}%
\pgfpathlineto{\pgfqpoint{2.756600in}{0.660000in}}%
\pgfpathlineto{\pgfqpoint{2.756600in}{3.480521in}}%
\pgfpathlineto{\pgfqpoint{2.750400in}{3.480521in}}%
\pgfpathlineto{\pgfqpoint{2.750400in}{0.660000in}}%
\pgfpathclose%
\pgfusepath{fill}%
\end{pgfscope}%
\begin{pgfscope}%
\pgfpathrectangle{\pgfqpoint{1.250000in}{0.660000in}}{\pgfqpoint{7.750000in}{4.620000in}}%
\pgfusepath{clip}%
\pgfsetbuttcap%
\pgfsetmiterjoin%
\definecolor{currentfill}{rgb}{0.000000,0.000000,0.000000}%
\pgfsetfillcolor{currentfill}%
\pgfsetlinewidth{0.000000pt}%
\definecolor{currentstroke}{rgb}{0.000000,0.000000,0.000000}%
\pgfsetstrokecolor{currentstroke}%
\pgfsetstrokeopacity{0.000000}%
\pgfsetdash{}{0pt}%
\pgfpathmoveto{\pgfqpoint{2.758150in}{0.660000in}}%
\pgfpathlineto{\pgfqpoint{2.764350in}{0.660000in}}%
\pgfpathlineto{\pgfqpoint{2.764350in}{3.519020in}}%
\pgfpathlineto{\pgfqpoint{2.758150in}{3.519020in}}%
\pgfpathlineto{\pgfqpoint{2.758150in}{0.660000in}}%
\pgfpathclose%
\pgfusepath{fill}%
\end{pgfscope}%
\begin{pgfscope}%
\pgfpathrectangle{\pgfqpoint{1.250000in}{0.660000in}}{\pgfqpoint{7.750000in}{4.620000in}}%
\pgfusepath{clip}%
\pgfsetbuttcap%
\pgfsetmiterjoin%
\definecolor{currentfill}{rgb}{0.000000,0.000000,0.000000}%
\pgfsetfillcolor{currentfill}%
\pgfsetlinewidth{0.000000pt}%
\definecolor{currentstroke}{rgb}{0.000000,0.000000,0.000000}%
\pgfsetstrokecolor{currentstroke}%
\pgfsetstrokeopacity{0.000000}%
\pgfsetdash{}{0pt}%
\pgfpathmoveto{\pgfqpoint{2.765900in}{0.660000in}}%
\pgfpathlineto{\pgfqpoint{2.772100in}{0.660000in}}%
\pgfpathlineto{\pgfqpoint{2.772100in}{3.463458in}}%
\pgfpathlineto{\pgfqpoint{2.765900in}{3.463458in}}%
\pgfpathlineto{\pgfqpoint{2.765900in}{0.660000in}}%
\pgfpathclose%
\pgfusepath{fill}%
\end{pgfscope}%
\begin{pgfscope}%
\pgfpathrectangle{\pgfqpoint{1.250000in}{0.660000in}}{\pgfqpoint{7.750000in}{4.620000in}}%
\pgfusepath{clip}%
\pgfsetbuttcap%
\pgfsetmiterjoin%
\definecolor{currentfill}{rgb}{0.000000,0.000000,0.000000}%
\pgfsetfillcolor{currentfill}%
\pgfsetlinewidth{0.000000pt}%
\definecolor{currentstroke}{rgb}{0.000000,0.000000,0.000000}%
\pgfsetstrokecolor{currentstroke}%
\pgfsetstrokeopacity{0.000000}%
\pgfsetdash{}{0pt}%
\pgfpathmoveto{\pgfqpoint{2.773650in}{0.660000in}}%
\pgfpathlineto{\pgfqpoint{2.779850in}{0.660000in}}%
\pgfpathlineto{\pgfqpoint{2.779850in}{3.511570in}}%
\pgfpathlineto{\pgfqpoint{2.773650in}{3.511570in}}%
\pgfpathlineto{\pgfqpoint{2.773650in}{0.660000in}}%
\pgfpathclose%
\pgfusepath{fill}%
\end{pgfscope}%
\begin{pgfscope}%
\pgfpathrectangle{\pgfqpoint{1.250000in}{0.660000in}}{\pgfqpoint{7.750000in}{4.620000in}}%
\pgfusepath{clip}%
\pgfsetbuttcap%
\pgfsetmiterjoin%
\definecolor{currentfill}{rgb}{0.000000,0.000000,0.000000}%
\pgfsetfillcolor{currentfill}%
\pgfsetlinewidth{0.000000pt}%
\definecolor{currentstroke}{rgb}{0.000000,0.000000,0.000000}%
\pgfsetstrokecolor{currentstroke}%
\pgfsetstrokeopacity{0.000000}%
\pgfsetdash{}{0pt}%
\pgfpathmoveto{\pgfqpoint{2.781400in}{0.660000in}}%
\pgfpathlineto{\pgfqpoint{2.787600in}{0.660000in}}%
\pgfpathlineto{\pgfqpoint{2.787600in}{3.473294in}}%
\pgfpathlineto{\pgfqpoint{2.781400in}{3.473294in}}%
\pgfpathlineto{\pgfqpoint{2.781400in}{0.660000in}}%
\pgfpathclose%
\pgfusepath{fill}%
\end{pgfscope}%
\begin{pgfscope}%
\pgfpathrectangle{\pgfqpoint{1.250000in}{0.660000in}}{\pgfqpoint{7.750000in}{4.620000in}}%
\pgfusepath{clip}%
\pgfsetbuttcap%
\pgfsetmiterjoin%
\definecolor{currentfill}{rgb}{0.000000,0.000000,0.000000}%
\pgfsetfillcolor{currentfill}%
\pgfsetlinewidth{0.000000pt}%
\definecolor{currentstroke}{rgb}{0.000000,0.000000,0.000000}%
\pgfsetstrokecolor{currentstroke}%
\pgfsetstrokeopacity{0.000000}%
\pgfsetdash{}{0pt}%
\pgfpathmoveto{\pgfqpoint{2.789150in}{0.660000in}}%
\pgfpathlineto{\pgfqpoint{2.795350in}{0.660000in}}%
\pgfpathlineto{\pgfqpoint{2.795350in}{3.438032in}}%
\pgfpathlineto{\pgfqpoint{2.789150in}{3.438032in}}%
\pgfpathlineto{\pgfqpoint{2.789150in}{0.660000in}}%
\pgfpathclose%
\pgfusepath{fill}%
\end{pgfscope}%
\begin{pgfscope}%
\pgfpathrectangle{\pgfqpoint{1.250000in}{0.660000in}}{\pgfqpoint{7.750000in}{4.620000in}}%
\pgfusepath{clip}%
\pgfsetbuttcap%
\pgfsetmiterjoin%
\definecolor{currentfill}{rgb}{0.000000,0.000000,0.000000}%
\pgfsetfillcolor{currentfill}%
\pgfsetlinewidth{0.000000pt}%
\definecolor{currentstroke}{rgb}{0.000000,0.000000,0.000000}%
\pgfsetstrokecolor{currentstroke}%
\pgfsetstrokeopacity{0.000000}%
\pgfsetdash{}{0pt}%
\pgfpathmoveto{\pgfqpoint{2.796900in}{0.660000in}}%
\pgfpathlineto{\pgfqpoint{2.803100in}{0.660000in}}%
\pgfpathlineto{\pgfqpoint{2.803100in}{3.444063in}}%
\pgfpathlineto{\pgfqpoint{2.796900in}{3.444063in}}%
\pgfpathlineto{\pgfqpoint{2.796900in}{0.660000in}}%
\pgfpathclose%
\pgfusepath{fill}%
\end{pgfscope}%
\begin{pgfscope}%
\pgfpathrectangle{\pgfqpoint{1.250000in}{0.660000in}}{\pgfqpoint{7.750000in}{4.620000in}}%
\pgfusepath{clip}%
\pgfsetbuttcap%
\pgfsetmiterjoin%
\definecolor{currentfill}{rgb}{0.000000,0.000000,0.000000}%
\pgfsetfillcolor{currentfill}%
\pgfsetlinewidth{0.000000pt}%
\definecolor{currentstroke}{rgb}{0.000000,0.000000,0.000000}%
\pgfsetstrokecolor{currentstroke}%
\pgfsetstrokeopacity{0.000000}%
\pgfsetdash{}{0pt}%
\pgfpathmoveto{\pgfqpoint{2.804650in}{0.660000in}}%
\pgfpathlineto{\pgfqpoint{2.810850in}{0.660000in}}%
\pgfpathlineto{\pgfqpoint{2.810850in}{3.472894in}}%
\pgfpathlineto{\pgfqpoint{2.804650in}{3.472894in}}%
\pgfpathlineto{\pgfqpoint{2.804650in}{0.660000in}}%
\pgfpathclose%
\pgfusepath{fill}%
\end{pgfscope}%
\begin{pgfscope}%
\pgfpathrectangle{\pgfqpoint{1.250000in}{0.660000in}}{\pgfqpoint{7.750000in}{4.620000in}}%
\pgfusepath{clip}%
\pgfsetbuttcap%
\pgfsetmiterjoin%
\definecolor{currentfill}{rgb}{0.000000,0.000000,0.000000}%
\pgfsetfillcolor{currentfill}%
\pgfsetlinewidth{0.000000pt}%
\definecolor{currentstroke}{rgb}{0.000000,0.000000,0.000000}%
\pgfsetstrokecolor{currentstroke}%
\pgfsetstrokeopacity{0.000000}%
\pgfsetdash{}{0pt}%
\pgfpathmoveto{\pgfqpoint{2.812400in}{0.660000in}}%
\pgfpathlineto{\pgfqpoint{2.818600in}{0.660000in}}%
\pgfpathlineto{\pgfqpoint{2.818600in}{3.451967in}}%
\pgfpathlineto{\pgfqpoint{2.812400in}{3.451967in}}%
\pgfpathlineto{\pgfqpoint{2.812400in}{0.660000in}}%
\pgfpathclose%
\pgfusepath{fill}%
\end{pgfscope}%
\begin{pgfscope}%
\pgfpathrectangle{\pgfqpoint{1.250000in}{0.660000in}}{\pgfqpoint{7.750000in}{4.620000in}}%
\pgfusepath{clip}%
\pgfsetbuttcap%
\pgfsetmiterjoin%
\definecolor{currentfill}{rgb}{0.000000,0.000000,0.000000}%
\pgfsetfillcolor{currentfill}%
\pgfsetlinewidth{0.000000pt}%
\definecolor{currentstroke}{rgb}{0.000000,0.000000,0.000000}%
\pgfsetstrokecolor{currentstroke}%
\pgfsetstrokeopacity{0.000000}%
\pgfsetdash{}{0pt}%
\pgfpathmoveto{\pgfqpoint{2.820150in}{0.660000in}}%
\pgfpathlineto{\pgfqpoint{2.826350in}{0.660000in}}%
\pgfpathlineto{\pgfqpoint{2.826350in}{3.456850in}}%
\pgfpathlineto{\pgfqpoint{2.820150in}{3.456850in}}%
\pgfpathlineto{\pgfqpoint{2.820150in}{0.660000in}}%
\pgfpathclose%
\pgfusepath{fill}%
\end{pgfscope}%
\begin{pgfscope}%
\pgfpathrectangle{\pgfqpoint{1.250000in}{0.660000in}}{\pgfqpoint{7.750000in}{4.620000in}}%
\pgfusepath{clip}%
\pgfsetbuttcap%
\pgfsetmiterjoin%
\definecolor{currentfill}{rgb}{0.000000,0.000000,0.000000}%
\pgfsetfillcolor{currentfill}%
\pgfsetlinewidth{0.000000pt}%
\definecolor{currentstroke}{rgb}{0.000000,0.000000,0.000000}%
\pgfsetstrokecolor{currentstroke}%
\pgfsetstrokeopacity{0.000000}%
\pgfsetdash{}{0pt}%
\pgfpathmoveto{\pgfqpoint{2.827900in}{0.660000in}}%
\pgfpathlineto{\pgfqpoint{2.834100in}{0.660000in}}%
\pgfpathlineto{\pgfqpoint{2.834100in}{3.457270in}}%
\pgfpathlineto{\pgfqpoint{2.827900in}{3.457270in}}%
\pgfpathlineto{\pgfqpoint{2.827900in}{0.660000in}}%
\pgfpathclose%
\pgfusepath{fill}%
\end{pgfscope}%
\begin{pgfscope}%
\pgfpathrectangle{\pgfqpoint{1.250000in}{0.660000in}}{\pgfqpoint{7.750000in}{4.620000in}}%
\pgfusepath{clip}%
\pgfsetbuttcap%
\pgfsetmiterjoin%
\definecolor{currentfill}{rgb}{0.000000,0.000000,0.000000}%
\pgfsetfillcolor{currentfill}%
\pgfsetlinewidth{0.000000pt}%
\definecolor{currentstroke}{rgb}{0.000000,0.000000,0.000000}%
\pgfsetstrokecolor{currentstroke}%
\pgfsetstrokeopacity{0.000000}%
\pgfsetdash{}{0pt}%
\pgfpathmoveto{\pgfqpoint{2.835650in}{0.660000in}}%
\pgfpathlineto{\pgfqpoint{2.841850in}{0.660000in}}%
\pgfpathlineto{\pgfqpoint{2.841850in}{3.458074in}}%
\pgfpathlineto{\pgfqpoint{2.835650in}{3.458074in}}%
\pgfpathlineto{\pgfqpoint{2.835650in}{0.660000in}}%
\pgfpathclose%
\pgfusepath{fill}%
\end{pgfscope}%
\begin{pgfscope}%
\pgfpathrectangle{\pgfqpoint{1.250000in}{0.660000in}}{\pgfqpoint{7.750000in}{4.620000in}}%
\pgfusepath{clip}%
\pgfsetbuttcap%
\pgfsetmiterjoin%
\definecolor{currentfill}{rgb}{0.000000,0.000000,0.000000}%
\pgfsetfillcolor{currentfill}%
\pgfsetlinewidth{0.000000pt}%
\definecolor{currentstroke}{rgb}{0.000000,0.000000,0.000000}%
\pgfsetstrokecolor{currentstroke}%
\pgfsetstrokeopacity{0.000000}%
\pgfsetdash{}{0pt}%
\pgfpathmoveto{\pgfqpoint{2.843400in}{0.660000in}}%
\pgfpathlineto{\pgfqpoint{2.849600in}{0.660000in}}%
\pgfpathlineto{\pgfqpoint{2.849600in}{3.422047in}}%
\pgfpathlineto{\pgfqpoint{2.843400in}{3.422047in}}%
\pgfpathlineto{\pgfqpoint{2.843400in}{0.660000in}}%
\pgfpathclose%
\pgfusepath{fill}%
\end{pgfscope}%
\begin{pgfscope}%
\pgfpathrectangle{\pgfqpoint{1.250000in}{0.660000in}}{\pgfqpoint{7.750000in}{4.620000in}}%
\pgfusepath{clip}%
\pgfsetbuttcap%
\pgfsetmiterjoin%
\definecolor{currentfill}{rgb}{0.000000,0.000000,0.000000}%
\pgfsetfillcolor{currentfill}%
\pgfsetlinewidth{0.000000pt}%
\definecolor{currentstroke}{rgb}{0.000000,0.000000,0.000000}%
\pgfsetstrokecolor{currentstroke}%
\pgfsetstrokeopacity{0.000000}%
\pgfsetdash{}{0pt}%
\pgfpathmoveto{\pgfqpoint{2.851150in}{0.660000in}}%
\pgfpathlineto{\pgfqpoint{2.857350in}{0.660000in}}%
\pgfpathlineto{\pgfqpoint{2.857350in}{3.456632in}}%
\pgfpathlineto{\pgfqpoint{2.851150in}{3.456632in}}%
\pgfpathlineto{\pgfqpoint{2.851150in}{0.660000in}}%
\pgfpathclose%
\pgfusepath{fill}%
\end{pgfscope}%
\begin{pgfscope}%
\pgfpathrectangle{\pgfqpoint{1.250000in}{0.660000in}}{\pgfqpoint{7.750000in}{4.620000in}}%
\pgfusepath{clip}%
\pgfsetbuttcap%
\pgfsetmiterjoin%
\definecolor{currentfill}{rgb}{0.000000,0.000000,0.000000}%
\pgfsetfillcolor{currentfill}%
\pgfsetlinewidth{0.000000pt}%
\definecolor{currentstroke}{rgb}{0.000000,0.000000,0.000000}%
\pgfsetstrokecolor{currentstroke}%
\pgfsetstrokeopacity{0.000000}%
\pgfsetdash{}{0pt}%
\pgfpathmoveto{\pgfqpoint{2.858900in}{0.660000in}}%
\pgfpathlineto{\pgfqpoint{2.865100in}{0.660000in}}%
\pgfpathlineto{\pgfqpoint{2.865100in}{3.421471in}}%
\pgfpathlineto{\pgfqpoint{2.858900in}{3.421471in}}%
\pgfpathlineto{\pgfqpoint{2.858900in}{0.660000in}}%
\pgfpathclose%
\pgfusepath{fill}%
\end{pgfscope}%
\begin{pgfscope}%
\pgfpathrectangle{\pgfqpoint{1.250000in}{0.660000in}}{\pgfqpoint{7.750000in}{4.620000in}}%
\pgfusepath{clip}%
\pgfsetbuttcap%
\pgfsetmiterjoin%
\definecolor{currentfill}{rgb}{0.000000,0.000000,0.000000}%
\pgfsetfillcolor{currentfill}%
\pgfsetlinewidth{0.000000pt}%
\definecolor{currentstroke}{rgb}{0.000000,0.000000,0.000000}%
\pgfsetstrokecolor{currentstroke}%
\pgfsetstrokeopacity{0.000000}%
\pgfsetdash{}{0pt}%
\pgfpathmoveto{\pgfqpoint{2.866650in}{0.660000in}}%
\pgfpathlineto{\pgfqpoint{2.872850in}{0.660000in}}%
\pgfpathlineto{\pgfqpoint{2.872850in}{3.421104in}}%
\pgfpathlineto{\pgfqpoint{2.866650in}{3.421104in}}%
\pgfpathlineto{\pgfqpoint{2.866650in}{0.660000in}}%
\pgfpathclose%
\pgfusepath{fill}%
\end{pgfscope}%
\begin{pgfscope}%
\pgfpathrectangle{\pgfqpoint{1.250000in}{0.660000in}}{\pgfqpoint{7.750000in}{4.620000in}}%
\pgfusepath{clip}%
\pgfsetbuttcap%
\pgfsetmiterjoin%
\definecolor{currentfill}{rgb}{0.000000,0.000000,0.000000}%
\pgfsetfillcolor{currentfill}%
\pgfsetlinewidth{0.000000pt}%
\definecolor{currentstroke}{rgb}{0.000000,0.000000,0.000000}%
\pgfsetstrokecolor{currentstroke}%
\pgfsetstrokeopacity{0.000000}%
\pgfsetdash{}{0pt}%
\pgfpathmoveto{\pgfqpoint{2.874400in}{0.660000in}}%
\pgfpathlineto{\pgfqpoint{2.880600in}{0.660000in}}%
\pgfpathlineto{\pgfqpoint{2.880600in}{3.474312in}}%
\pgfpathlineto{\pgfqpoint{2.874400in}{3.474312in}}%
\pgfpathlineto{\pgfqpoint{2.874400in}{0.660000in}}%
\pgfpathclose%
\pgfusepath{fill}%
\end{pgfscope}%
\begin{pgfscope}%
\pgfpathrectangle{\pgfqpoint{1.250000in}{0.660000in}}{\pgfqpoint{7.750000in}{4.620000in}}%
\pgfusepath{clip}%
\pgfsetbuttcap%
\pgfsetmiterjoin%
\definecolor{currentfill}{rgb}{0.000000,0.000000,0.000000}%
\pgfsetfillcolor{currentfill}%
\pgfsetlinewidth{0.000000pt}%
\definecolor{currentstroke}{rgb}{0.000000,0.000000,0.000000}%
\pgfsetstrokecolor{currentstroke}%
\pgfsetstrokeopacity{0.000000}%
\pgfsetdash{}{0pt}%
\pgfpathmoveto{\pgfqpoint{2.882150in}{0.660000in}}%
\pgfpathlineto{\pgfqpoint{2.888350in}{0.660000in}}%
\pgfpathlineto{\pgfqpoint{2.888350in}{3.431513in}}%
\pgfpathlineto{\pgfqpoint{2.882150in}{3.431513in}}%
\pgfpathlineto{\pgfqpoint{2.882150in}{0.660000in}}%
\pgfpathclose%
\pgfusepath{fill}%
\end{pgfscope}%
\begin{pgfscope}%
\pgfpathrectangle{\pgfqpoint{1.250000in}{0.660000in}}{\pgfqpoint{7.750000in}{4.620000in}}%
\pgfusepath{clip}%
\pgfsetbuttcap%
\pgfsetmiterjoin%
\definecolor{currentfill}{rgb}{0.000000,0.000000,0.000000}%
\pgfsetfillcolor{currentfill}%
\pgfsetlinewidth{0.000000pt}%
\definecolor{currentstroke}{rgb}{0.000000,0.000000,0.000000}%
\pgfsetstrokecolor{currentstroke}%
\pgfsetstrokeopacity{0.000000}%
\pgfsetdash{}{0pt}%
\pgfpathmoveto{\pgfqpoint{2.889900in}{0.660000in}}%
\pgfpathlineto{\pgfqpoint{2.896100in}{0.660000in}}%
\pgfpathlineto{\pgfqpoint{2.896100in}{3.466394in}}%
\pgfpathlineto{\pgfqpoint{2.889900in}{3.466394in}}%
\pgfpathlineto{\pgfqpoint{2.889900in}{0.660000in}}%
\pgfpathclose%
\pgfusepath{fill}%
\end{pgfscope}%
\begin{pgfscope}%
\pgfpathrectangle{\pgfqpoint{1.250000in}{0.660000in}}{\pgfqpoint{7.750000in}{4.620000in}}%
\pgfusepath{clip}%
\pgfsetbuttcap%
\pgfsetmiterjoin%
\definecolor{currentfill}{rgb}{0.000000,0.000000,0.000000}%
\pgfsetfillcolor{currentfill}%
\pgfsetlinewidth{0.000000pt}%
\definecolor{currentstroke}{rgb}{0.000000,0.000000,0.000000}%
\pgfsetstrokecolor{currentstroke}%
\pgfsetstrokeopacity{0.000000}%
\pgfsetdash{}{0pt}%
\pgfpathmoveto{\pgfqpoint{2.897650in}{0.660000in}}%
\pgfpathlineto{\pgfqpoint{2.903850in}{0.660000in}}%
\pgfpathlineto{\pgfqpoint{2.903850in}{3.478469in}}%
\pgfpathlineto{\pgfqpoint{2.897650in}{3.478469in}}%
\pgfpathlineto{\pgfqpoint{2.897650in}{0.660000in}}%
\pgfpathclose%
\pgfusepath{fill}%
\end{pgfscope}%
\begin{pgfscope}%
\pgfpathrectangle{\pgfqpoint{1.250000in}{0.660000in}}{\pgfqpoint{7.750000in}{4.620000in}}%
\pgfusepath{clip}%
\pgfsetbuttcap%
\pgfsetmiterjoin%
\definecolor{currentfill}{rgb}{0.000000,0.000000,0.000000}%
\pgfsetfillcolor{currentfill}%
\pgfsetlinewidth{0.000000pt}%
\definecolor{currentstroke}{rgb}{0.000000,0.000000,0.000000}%
\pgfsetstrokecolor{currentstroke}%
\pgfsetstrokeopacity{0.000000}%
\pgfsetdash{}{0pt}%
\pgfpathmoveto{\pgfqpoint{2.905400in}{0.660000in}}%
\pgfpathlineto{\pgfqpoint{2.911600in}{0.660000in}}%
\pgfpathlineto{\pgfqpoint{2.911600in}{3.403140in}}%
\pgfpathlineto{\pgfqpoint{2.905400in}{3.403140in}}%
\pgfpathlineto{\pgfqpoint{2.905400in}{0.660000in}}%
\pgfpathclose%
\pgfusepath{fill}%
\end{pgfscope}%
\begin{pgfscope}%
\pgfpathrectangle{\pgfqpoint{1.250000in}{0.660000in}}{\pgfqpoint{7.750000in}{4.620000in}}%
\pgfusepath{clip}%
\pgfsetbuttcap%
\pgfsetmiterjoin%
\definecolor{currentfill}{rgb}{0.000000,0.000000,0.000000}%
\pgfsetfillcolor{currentfill}%
\pgfsetlinewidth{0.000000pt}%
\definecolor{currentstroke}{rgb}{0.000000,0.000000,0.000000}%
\pgfsetstrokecolor{currentstroke}%
\pgfsetstrokeopacity{0.000000}%
\pgfsetdash{}{0pt}%
\pgfpathmoveto{\pgfqpoint{2.913150in}{0.660000in}}%
\pgfpathlineto{\pgfqpoint{2.919350in}{0.660000in}}%
\pgfpathlineto{\pgfqpoint{2.919350in}{3.422555in}}%
\pgfpathlineto{\pgfqpoint{2.913150in}{3.422555in}}%
\pgfpathlineto{\pgfqpoint{2.913150in}{0.660000in}}%
\pgfpathclose%
\pgfusepath{fill}%
\end{pgfscope}%
\begin{pgfscope}%
\pgfpathrectangle{\pgfqpoint{1.250000in}{0.660000in}}{\pgfqpoint{7.750000in}{4.620000in}}%
\pgfusepath{clip}%
\pgfsetbuttcap%
\pgfsetmiterjoin%
\definecolor{currentfill}{rgb}{0.000000,0.000000,0.000000}%
\pgfsetfillcolor{currentfill}%
\pgfsetlinewidth{0.000000pt}%
\definecolor{currentstroke}{rgb}{0.000000,0.000000,0.000000}%
\pgfsetstrokecolor{currentstroke}%
\pgfsetstrokeopacity{0.000000}%
\pgfsetdash{}{0pt}%
\pgfpathmoveto{\pgfqpoint{2.920900in}{0.660000in}}%
\pgfpathlineto{\pgfqpoint{2.927100in}{0.660000in}}%
\pgfpathlineto{\pgfqpoint{2.927100in}{3.414195in}}%
\pgfpathlineto{\pgfqpoint{2.920900in}{3.414195in}}%
\pgfpathlineto{\pgfqpoint{2.920900in}{0.660000in}}%
\pgfpathclose%
\pgfusepath{fill}%
\end{pgfscope}%
\begin{pgfscope}%
\pgfpathrectangle{\pgfqpoint{1.250000in}{0.660000in}}{\pgfqpoint{7.750000in}{4.620000in}}%
\pgfusepath{clip}%
\pgfsetbuttcap%
\pgfsetmiterjoin%
\definecolor{currentfill}{rgb}{0.000000,0.000000,0.000000}%
\pgfsetfillcolor{currentfill}%
\pgfsetlinewidth{0.000000pt}%
\definecolor{currentstroke}{rgb}{0.000000,0.000000,0.000000}%
\pgfsetstrokecolor{currentstroke}%
\pgfsetstrokeopacity{0.000000}%
\pgfsetdash{}{0pt}%
\pgfpathmoveto{\pgfqpoint{2.928650in}{0.660000in}}%
\pgfpathlineto{\pgfqpoint{2.934850in}{0.660000in}}%
\pgfpathlineto{\pgfqpoint{2.934850in}{3.495759in}}%
\pgfpathlineto{\pgfqpoint{2.928650in}{3.495759in}}%
\pgfpathlineto{\pgfqpoint{2.928650in}{0.660000in}}%
\pgfpathclose%
\pgfusepath{fill}%
\end{pgfscope}%
\begin{pgfscope}%
\pgfpathrectangle{\pgfqpoint{1.250000in}{0.660000in}}{\pgfqpoint{7.750000in}{4.620000in}}%
\pgfusepath{clip}%
\pgfsetbuttcap%
\pgfsetmiterjoin%
\definecolor{currentfill}{rgb}{0.000000,0.000000,0.000000}%
\pgfsetfillcolor{currentfill}%
\pgfsetlinewidth{0.000000pt}%
\definecolor{currentstroke}{rgb}{0.000000,0.000000,0.000000}%
\pgfsetstrokecolor{currentstroke}%
\pgfsetstrokeopacity{0.000000}%
\pgfsetdash{}{0pt}%
\pgfpathmoveto{\pgfqpoint{2.936400in}{0.660000in}}%
\pgfpathlineto{\pgfqpoint{2.942600in}{0.660000in}}%
\pgfpathlineto{\pgfqpoint{2.942600in}{3.411164in}}%
\pgfpathlineto{\pgfqpoint{2.936400in}{3.411164in}}%
\pgfpathlineto{\pgfqpoint{2.936400in}{0.660000in}}%
\pgfpathclose%
\pgfusepath{fill}%
\end{pgfscope}%
\begin{pgfscope}%
\pgfpathrectangle{\pgfqpoint{1.250000in}{0.660000in}}{\pgfqpoint{7.750000in}{4.620000in}}%
\pgfusepath{clip}%
\pgfsetbuttcap%
\pgfsetmiterjoin%
\definecolor{currentfill}{rgb}{0.000000,0.000000,0.000000}%
\pgfsetfillcolor{currentfill}%
\pgfsetlinewidth{0.000000pt}%
\definecolor{currentstroke}{rgb}{0.000000,0.000000,0.000000}%
\pgfsetstrokecolor{currentstroke}%
\pgfsetstrokeopacity{0.000000}%
\pgfsetdash{}{0pt}%
\pgfpathmoveto{\pgfqpoint{2.944150in}{0.660000in}}%
\pgfpathlineto{\pgfqpoint{2.950350in}{0.660000in}}%
\pgfpathlineto{\pgfqpoint{2.950350in}{3.475582in}}%
\pgfpathlineto{\pgfqpoint{2.944150in}{3.475582in}}%
\pgfpathlineto{\pgfqpoint{2.944150in}{0.660000in}}%
\pgfpathclose%
\pgfusepath{fill}%
\end{pgfscope}%
\begin{pgfscope}%
\pgfpathrectangle{\pgfqpoint{1.250000in}{0.660000in}}{\pgfqpoint{7.750000in}{4.620000in}}%
\pgfusepath{clip}%
\pgfsetbuttcap%
\pgfsetmiterjoin%
\definecolor{currentfill}{rgb}{0.000000,0.000000,0.000000}%
\pgfsetfillcolor{currentfill}%
\pgfsetlinewidth{0.000000pt}%
\definecolor{currentstroke}{rgb}{0.000000,0.000000,0.000000}%
\pgfsetstrokecolor{currentstroke}%
\pgfsetstrokeopacity{0.000000}%
\pgfsetdash{}{0pt}%
\pgfpathmoveto{\pgfqpoint{2.951900in}{0.660000in}}%
\pgfpathlineto{\pgfqpoint{2.958100in}{0.660000in}}%
\pgfpathlineto{\pgfqpoint{2.958100in}{3.396128in}}%
\pgfpathlineto{\pgfqpoint{2.951900in}{3.396128in}}%
\pgfpathlineto{\pgfqpoint{2.951900in}{0.660000in}}%
\pgfpathclose%
\pgfusepath{fill}%
\end{pgfscope}%
\begin{pgfscope}%
\pgfpathrectangle{\pgfqpoint{1.250000in}{0.660000in}}{\pgfqpoint{7.750000in}{4.620000in}}%
\pgfusepath{clip}%
\pgfsetbuttcap%
\pgfsetmiterjoin%
\definecolor{currentfill}{rgb}{0.000000,0.000000,0.000000}%
\pgfsetfillcolor{currentfill}%
\pgfsetlinewidth{0.000000pt}%
\definecolor{currentstroke}{rgb}{0.000000,0.000000,0.000000}%
\pgfsetstrokecolor{currentstroke}%
\pgfsetstrokeopacity{0.000000}%
\pgfsetdash{}{0pt}%
\pgfpathmoveto{\pgfqpoint{2.959650in}{0.660000in}}%
\pgfpathlineto{\pgfqpoint{2.965850in}{0.660000in}}%
\pgfpathlineto{\pgfqpoint{2.965850in}{3.429295in}}%
\pgfpathlineto{\pgfqpoint{2.959650in}{3.429295in}}%
\pgfpathlineto{\pgfqpoint{2.959650in}{0.660000in}}%
\pgfpathclose%
\pgfusepath{fill}%
\end{pgfscope}%
\begin{pgfscope}%
\pgfpathrectangle{\pgfqpoint{1.250000in}{0.660000in}}{\pgfqpoint{7.750000in}{4.620000in}}%
\pgfusepath{clip}%
\pgfsetbuttcap%
\pgfsetmiterjoin%
\definecolor{currentfill}{rgb}{0.000000,0.000000,0.000000}%
\pgfsetfillcolor{currentfill}%
\pgfsetlinewidth{0.000000pt}%
\definecolor{currentstroke}{rgb}{0.000000,0.000000,0.000000}%
\pgfsetstrokecolor{currentstroke}%
\pgfsetstrokeopacity{0.000000}%
\pgfsetdash{}{0pt}%
\pgfpathmoveto{\pgfqpoint{2.967400in}{0.660000in}}%
\pgfpathlineto{\pgfqpoint{2.973600in}{0.660000in}}%
\pgfpathlineto{\pgfqpoint{2.973600in}{3.423197in}}%
\pgfpathlineto{\pgfqpoint{2.967400in}{3.423197in}}%
\pgfpathlineto{\pgfqpoint{2.967400in}{0.660000in}}%
\pgfpathclose%
\pgfusepath{fill}%
\end{pgfscope}%
\begin{pgfscope}%
\pgfpathrectangle{\pgfqpoint{1.250000in}{0.660000in}}{\pgfqpoint{7.750000in}{4.620000in}}%
\pgfusepath{clip}%
\pgfsetbuttcap%
\pgfsetmiterjoin%
\definecolor{currentfill}{rgb}{0.000000,0.000000,0.000000}%
\pgfsetfillcolor{currentfill}%
\pgfsetlinewidth{0.000000pt}%
\definecolor{currentstroke}{rgb}{0.000000,0.000000,0.000000}%
\pgfsetstrokecolor{currentstroke}%
\pgfsetstrokeopacity{0.000000}%
\pgfsetdash{}{0pt}%
\pgfpathmoveto{\pgfqpoint{2.975150in}{0.660000in}}%
\pgfpathlineto{\pgfqpoint{2.981350in}{0.660000in}}%
\pgfpathlineto{\pgfqpoint{2.981350in}{3.444993in}}%
\pgfpathlineto{\pgfqpoint{2.975150in}{3.444993in}}%
\pgfpathlineto{\pgfqpoint{2.975150in}{0.660000in}}%
\pgfpathclose%
\pgfusepath{fill}%
\end{pgfscope}%
\begin{pgfscope}%
\pgfpathrectangle{\pgfqpoint{1.250000in}{0.660000in}}{\pgfqpoint{7.750000in}{4.620000in}}%
\pgfusepath{clip}%
\pgfsetbuttcap%
\pgfsetmiterjoin%
\definecolor{currentfill}{rgb}{0.000000,0.000000,0.000000}%
\pgfsetfillcolor{currentfill}%
\pgfsetlinewidth{0.000000pt}%
\definecolor{currentstroke}{rgb}{0.000000,0.000000,0.000000}%
\pgfsetstrokecolor{currentstroke}%
\pgfsetstrokeopacity{0.000000}%
\pgfsetdash{}{0pt}%
\pgfpathmoveto{\pgfqpoint{2.982900in}{0.660000in}}%
\pgfpathlineto{\pgfqpoint{2.989100in}{0.660000in}}%
\pgfpathlineto{\pgfqpoint{2.989100in}{3.446459in}}%
\pgfpathlineto{\pgfqpoint{2.982900in}{3.446459in}}%
\pgfpathlineto{\pgfqpoint{2.982900in}{0.660000in}}%
\pgfpathclose%
\pgfusepath{fill}%
\end{pgfscope}%
\begin{pgfscope}%
\pgfpathrectangle{\pgfqpoint{1.250000in}{0.660000in}}{\pgfqpoint{7.750000in}{4.620000in}}%
\pgfusepath{clip}%
\pgfsetbuttcap%
\pgfsetmiterjoin%
\definecolor{currentfill}{rgb}{0.000000,0.000000,0.000000}%
\pgfsetfillcolor{currentfill}%
\pgfsetlinewidth{0.000000pt}%
\definecolor{currentstroke}{rgb}{0.000000,0.000000,0.000000}%
\pgfsetstrokecolor{currentstroke}%
\pgfsetstrokeopacity{0.000000}%
\pgfsetdash{}{0pt}%
\pgfpathmoveto{\pgfqpoint{2.990650in}{0.660000in}}%
\pgfpathlineto{\pgfqpoint{2.996850in}{0.660000in}}%
\pgfpathlineto{\pgfqpoint{2.996850in}{3.419068in}}%
\pgfpathlineto{\pgfqpoint{2.990650in}{3.419068in}}%
\pgfpathlineto{\pgfqpoint{2.990650in}{0.660000in}}%
\pgfpathclose%
\pgfusepath{fill}%
\end{pgfscope}%
\begin{pgfscope}%
\pgfpathrectangle{\pgfqpoint{1.250000in}{0.660000in}}{\pgfqpoint{7.750000in}{4.620000in}}%
\pgfusepath{clip}%
\pgfsetbuttcap%
\pgfsetmiterjoin%
\definecolor{currentfill}{rgb}{0.000000,0.000000,0.000000}%
\pgfsetfillcolor{currentfill}%
\pgfsetlinewidth{0.000000pt}%
\definecolor{currentstroke}{rgb}{0.000000,0.000000,0.000000}%
\pgfsetstrokecolor{currentstroke}%
\pgfsetstrokeopacity{0.000000}%
\pgfsetdash{}{0pt}%
\pgfpathmoveto{\pgfqpoint{2.998400in}{0.660000in}}%
\pgfpathlineto{\pgfqpoint{3.004600in}{0.660000in}}%
\pgfpathlineto{\pgfqpoint{3.004600in}{3.449111in}}%
\pgfpathlineto{\pgfqpoint{2.998400in}{3.449111in}}%
\pgfpathlineto{\pgfqpoint{2.998400in}{0.660000in}}%
\pgfpathclose%
\pgfusepath{fill}%
\end{pgfscope}%
\begin{pgfscope}%
\pgfpathrectangle{\pgfqpoint{1.250000in}{0.660000in}}{\pgfqpoint{7.750000in}{4.620000in}}%
\pgfusepath{clip}%
\pgfsetbuttcap%
\pgfsetmiterjoin%
\definecolor{currentfill}{rgb}{0.000000,0.000000,0.000000}%
\pgfsetfillcolor{currentfill}%
\pgfsetlinewidth{0.000000pt}%
\definecolor{currentstroke}{rgb}{0.000000,0.000000,0.000000}%
\pgfsetstrokecolor{currentstroke}%
\pgfsetstrokeopacity{0.000000}%
\pgfsetdash{}{0pt}%
\pgfpathmoveto{\pgfqpoint{3.006150in}{0.660000in}}%
\pgfpathlineto{\pgfqpoint{3.012350in}{0.660000in}}%
\pgfpathlineto{\pgfqpoint{3.012350in}{3.426588in}}%
\pgfpathlineto{\pgfqpoint{3.006150in}{3.426588in}}%
\pgfpathlineto{\pgfqpoint{3.006150in}{0.660000in}}%
\pgfpathclose%
\pgfusepath{fill}%
\end{pgfscope}%
\begin{pgfscope}%
\pgfpathrectangle{\pgfqpoint{1.250000in}{0.660000in}}{\pgfqpoint{7.750000in}{4.620000in}}%
\pgfusepath{clip}%
\pgfsetbuttcap%
\pgfsetmiterjoin%
\definecolor{currentfill}{rgb}{0.000000,0.000000,0.000000}%
\pgfsetfillcolor{currentfill}%
\pgfsetlinewidth{0.000000pt}%
\definecolor{currentstroke}{rgb}{0.000000,0.000000,0.000000}%
\pgfsetstrokecolor{currentstroke}%
\pgfsetstrokeopacity{0.000000}%
\pgfsetdash{}{0pt}%
\pgfpathmoveto{\pgfqpoint{3.013900in}{0.660000in}}%
\pgfpathlineto{\pgfqpoint{3.020100in}{0.660000in}}%
\pgfpathlineto{\pgfqpoint{3.020100in}{3.439636in}}%
\pgfpathlineto{\pgfqpoint{3.013900in}{3.439636in}}%
\pgfpathlineto{\pgfqpoint{3.013900in}{0.660000in}}%
\pgfpathclose%
\pgfusepath{fill}%
\end{pgfscope}%
\begin{pgfscope}%
\pgfpathrectangle{\pgfqpoint{1.250000in}{0.660000in}}{\pgfqpoint{7.750000in}{4.620000in}}%
\pgfusepath{clip}%
\pgfsetbuttcap%
\pgfsetmiterjoin%
\definecolor{currentfill}{rgb}{0.000000,0.000000,0.000000}%
\pgfsetfillcolor{currentfill}%
\pgfsetlinewidth{0.000000pt}%
\definecolor{currentstroke}{rgb}{0.000000,0.000000,0.000000}%
\pgfsetstrokecolor{currentstroke}%
\pgfsetstrokeopacity{0.000000}%
\pgfsetdash{}{0pt}%
\pgfpathmoveto{\pgfqpoint{3.021650in}{0.660000in}}%
\pgfpathlineto{\pgfqpoint{3.027850in}{0.660000in}}%
\pgfpathlineto{\pgfqpoint{3.027850in}{3.455710in}}%
\pgfpathlineto{\pgfqpoint{3.021650in}{3.455710in}}%
\pgfpathlineto{\pgfqpoint{3.021650in}{0.660000in}}%
\pgfpathclose%
\pgfusepath{fill}%
\end{pgfscope}%
\begin{pgfscope}%
\pgfpathrectangle{\pgfqpoint{1.250000in}{0.660000in}}{\pgfqpoint{7.750000in}{4.620000in}}%
\pgfusepath{clip}%
\pgfsetbuttcap%
\pgfsetmiterjoin%
\definecolor{currentfill}{rgb}{0.000000,0.000000,0.000000}%
\pgfsetfillcolor{currentfill}%
\pgfsetlinewidth{0.000000pt}%
\definecolor{currentstroke}{rgb}{0.000000,0.000000,0.000000}%
\pgfsetstrokecolor{currentstroke}%
\pgfsetstrokeopacity{0.000000}%
\pgfsetdash{}{0pt}%
\pgfpathmoveto{\pgfqpoint{3.029400in}{0.660000in}}%
\pgfpathlineto{\pgfqpoint{3.035600in}{0.660000in}}%
\pgfpathlineto{\pgfqpoint{3.035600in}{3.434943in}}%
\pgfpathlineto{\pgfqpoint{3.029400in}{3.434943in}}%
\pgfpathlineto{\pgfqpoint{3.029400in}{0.660000in}}%
\pgfpathclose%
\pgfusepath{fill}%
\end{pgfscope}%
\begin{pgfscope}%
\pgfpathrectangle{\pgfqpoint{1.250000in}{0.660000in}}{\pgfqpoint{7.750000in}{4.620000in}}%
\pgfusepath{clip}%
\pgfsetbuttcap%
\pgfsetmiterjoin%
\definecolor{currentfill}{rgb}{0.000000,0.000000,0.000000}%
\pgfsetfillcolor{currentfill}%
\pgfsetlinewidth{0.000000pt}%
\definecolor{currentstroke}{rgb}{0.000000,0.000000,0.000000}%
\pgfsetstrokecolor{currentstroke}%
\pgfsetstrokeopacity{0.000000}%
\pgfsetdash{}{0pt}%
\pgfpathmoveto{\pgfqpoint{3.037150in}{0.660000in}}%
\pgfpathlineto{\pgfqpoint{3.043350in}{0.660000in}}%
\pgfpathlineto{\pgfqpoint{3.043350in}{3.403813in}}%
\pgfpathlineto{\pgfqpoint{3.037150in}{3.403813in}}%
\pgfpathlineto{\pgfqpoint{3.037150in}{0.660000in}}%
\pgfpathclose%
\pgfusepath{fill}%
\end{pgfscope}%
\begin{pgfscope}%
\pgfpathrectangle{\pgfqpoint{1.250000in}{0.660000in}}{\pgfqpoint{7.750000in}{4.620000in}}%
\pgfusepath{clip}%
\pgfsetbuttcap%
\pgfsetmiterjoin%
\definecolor{currentfill}{rgb}{0.000000,0.000000,0.000000}%
\pgfsetfillcolor{currentfill}%
\pgfsetlinewidth{0.000000pt}%
\definecolor{currentstroke}{rgb}{0.000000,0.000000,0.000000}%
\pgfsetstrokecolor{currentstroke}%
\pgfsetstrokeopacity{0.000000}%
\pgfsetdash{}{0pt}%
\pgfpathmoveto{\pgfqpoint{3.044900in}{0.660000in}}%
\pgfpathlineto{\pgfqpoint{3.051100in}{0.660000in}}%
\pgfpathlineto{\pgfqpoint{3.051100in}{3.381508in}}%
\pgfpathlineto{\pgfqpoint{3.044900in}{3.381508in}}%
\pgfpathlineto{\pgfqpoint{3.044900in}{0.660000in}}%
\pgfpathclose%
\pgfusepath{fill}%
\end{pgfscope}%
\begin{pgfscope}%
\pgfpathrectangle{\pgfqpoint{1.250000in}{0.660000in}}{\pgfqpoint{7.750000in}{4.620000in}}%
\pgfusepath{clip}%
\pgfsetbuttcap%
\pgfsetmiterjoin%
\definecolor{currentfill}{rgb}{0.000000,0.000000,0.000000}%
\pgfsetfillcolor{currentfill}%
\pgfsetlinewidth{0.000000pt}%
\definecolor{currentstroke}{rgb}{0.000000,0.000000,0.000000}%
\pgfsetstrokecolor{currentstroke}%
\pgfsetstrokeopacity{0.000000}%
\pgfsetdash{}{0pt}%
\pgfpathmoveto{\pgfqpoint{3.052650in}{0.660000in}}%
\pgfpathlineto{\pgfqpoint{3.058850in}{0.660000in}}%
\pgfpathlineto{\pgfqpoint{3.058850in}{3.409461in}}%
\pgfpathlineto{\pgfqpoint{3.052650in}{3.409461in}}%
\pgfpathlineto{\pgfqpoint{3.052650in}{0.660000in}}%
\pgfpathclose%
\pgfusepath{fill}%
\end{pgfscope}%
\begin{pgfscope}%
\pgfpathrectangle{\pgfqpoint{1.250000in}{0.660000in}}{\pgfqpoint{7.750000in}{4.620000in}}%
\pgfusepath{clip}%
\pgfsetbuttcap%
\pgfsetmiterjoin%
\definecolor{currentfill}{rgb}{0.000000,0.000000,0.000000}%
\pgfsetfillcolor{currentfill}%
\pgfsetlinewidth{0.000000pt}%
\definecolor{currentstroke}{rgb}{0.000000,0.000000,0.000000}%
\pgfsetstrokecolor{currentstroke}%
\pgfsetstrokeopacity{0.000000}%
\pgfsetdash{}{0pt}%
\pgfpathmoveto{\pgfqpoint{3.060400in}{0.660000in}}%
\pgfpathlineto{\pgfqpoint{3.066600in}{0.660000in}}%
\pgfpathlineto{\pgfqpoint{3.066600in}{3.476051in}}%
\pgfpathlineto{\pgfqpoint{3.060400in}{3.476051in}}%
\pgfpathlineto{\pgfqpoint{3.060400in}{0.660000in}}%
\pgfpathclose%
\pgfusepath{fill}%
\end{pgfscope}%
\begin{pgfscope}%
\pgfpathrectangle{\pgfqpoint{1.250000in}{0.660000in}}{\pgfqpoint{7.750000in}{4.620000in}}%
\pgfusepath{clip}%
\pgfsetbuttcap%
\pgfsetmiterjoin%
\definecolor{currentfill}{rgb}{0.000000,0.000000,0.000000}%
\pgfsetfillcolor{currentfill}%
\pgfsetlinewidth{0.000000pt}%
\definecolor{currentstroke}{rgb}{0.000000,0.000000,0.000000}%
\pgfsetstrokecolor{currentstroke}%
\pgfsetstrokeopacity{0.000000}%
\pgfsetdash{}{0pt}%
\pgfpathmoveto{\pgfqpoint{3.068150in}{0.660000in}}%
\pgfpathlineto{\pgfqpoint{3.074350in}{0.660000in}}%
\pgfpathlineto{\pgfqpoint{3.074350in}{3.421832in}}%
\pgfpathlineto{\pgfqpoint{3.068150in}{3.421832in}}%
\pgfpathlineto{\pgfqpoint{3.068150in}{0.660000in}}%
\pgfpathclose%
\pgfusepath{fill}%
\end{pgfscope}%
\begin{pgfscope}%
\pgfpathrectangle{\pgfqpoint{1.250000in}{0.660000in}}{\pgfqpoint{7.750000in}{4.620000in}}%
\pgfusepath{clip}%
\pgfsetbuttcap%
\pgfsetmiterjoin%
\definecolor{currentfill}{rgb}{0.000000,0.000000,0.000000}%
\pgfsetfillcolor{currentfill}%
\pgfsetlinewidth{0.000000pt}%
\definecolor{currentstroke}{rgb}{0.000000,0.000000,0.000000}%
\pgfsetstrokecolor{currentstroke}%
\pgfsetstrokeopacity{0.000000}%
\pgfsetdash{}{0pt}%
\pgfpathmoveto{\pgfqpoint{3.075900in}{0.660000in}}%
\pgfpathlineto{\pgfqpoint{3.082100in}{0.660000in}}%
\pgfpathlineto{\pgfqpoint{3.082100in}{3.427432in}}%
\pgfpathlineto{\pgfqpoint{3.075900in}{3.427432in}}%
\pgfpathlineto{\pgfqpoint{3.075900in}{0.660000in}}%
\pgfpathclose%
\pgfusepath{fill}%
\end{pgfscope}%
\begin{pgfscope}%
\pgfpathrectangle{\pgfqpoint{1.250000in}{0.660000in}}{\pgfqpoint{7.750000in}{4.620000in}}%
\pgfusepath{clip}%
\pgfsetbuttcap%
\pgfsetmiterjoin%
\definecolor{currentfill}{rgb}{0.000000,0.000000,0.000000}%
\pgfsetfillcolor{currentfill}%
\pgfsetlinewidth{0.000000pt}%
\definecolor{currentstroke}{rgb}{0.000000,0.000000,0.000000}%
\pgfsetstrokecolor{currentstroke}%
\pgfsetstrokeopacity{0.000000}%
\pgfsetdash{}{0pt}%
\pgfpathmoveto{\pgfqpoint{3.083650in}{0.660000in}}%
\pgfpathlineto{\pgfqpoint{3.089850in}{0.660000in}}%
\pgfpathlineto{\pgfqpoint{3.089850in}{3.422151in}}%
\pgfpathlineto{\pgfqpoint{3.083650in}{3.422151in}}%
\pgfpathlineto{\pgfqpoint{3.083650in}{0.660000in}}%
\pgfpathclose%
\pgfusepath{fill}%
\end{pgfscope}%
\begin{pgfscope}%
\pgfpathrectangle{\pgfqpoint{1.250000in}{0.660000in}}{\pgfqpoint{7.750000in}{4.620000in}}%
\pgfusepath{clip}%
\pgfsetbuttcap%
\pgfsetmiterjoin%
\definecolor{currentfill}{rgb}{0.000000,0.000000,0.000000}%
\pgfsetfillcolor{currentfill}%
\pgfsetlinewidth{0.000000pt}%
\definecolor{currentstroke}{rgb}{0.000000,0.000000,0.000000}%
\pgfsetstrokecolor{currentstroke}%
\pgfsetstrokeopacity{0.000000}%
\pgfsetdash{}{0pt}%
\pgfpathmoveto{\pgfqpoint{3.091400in}{0.660000in}}%
\pgfpathlineto{\pgfqpoint{3.097600in}{0.660000in}}%
\pgfpathlineto{\pgfqpoint{3.097600in}{3.397483in}}%
\pgfpathlineto{\pgfqpoint{3.091400in}{3.397483in}}%
\pgfpathlineto{\pgfqpoint{3.091400in}{0.660000in}}%
\pgfpathclose%
\pgfusepath{fill}%
\end{pgfscope}%
\begin{pgfscope}%
\pgfpathrectangle{\pgfqpoint{1.250000in}{0.660000in}}{\pgfqpoint{7.750000in}{4.620000in}}%
\pgfusepath{clip}%
\pgfsetbuttcap%
\pgfsetmiterjoin%
\definecolor{currentfill}{rgb}{0.000000,0.000000,0.000000}%
\pgfsetfillcolor{currentfill}%
\pgfsetlinewidth{0.000000pt}%
\definecolor{currentstroke}{rgb}{0.000000,0.000000,0.000000}%
\pgfsetstrokecolor{currentstroke}%
\pgfsetstrokeopacity{0.000000}%
\pgfsetdash{}{0pt}%
\pgfpathmoveto{\pgfqpoint{3.099150in}{0.660000in}}%
\pgfpathlineto{\pgfqpoint{3.105350in}{0.660000in}}%
\pgfpathlineto{\pgfqpoint{3.105350in}{3.464278in}}%
\pgfpathlineto{\pgfqpoint{3.099150in}{3.464278in}}%
\pgfpathlineto{\pgfqpoint{3.099150in}{0.660000in}}%
\pgfpathclose%
\pgfusepath{fill}%
\end{pgfscope}%
\begin{pgfscope}%
\pgfpathrectangle{\pgfqpoint{1.250000in}{0.660000in}}{\pgfqpoint{7.750000in}{4.620000in}}%
\pgfusepath{clip}%
\pgfsetbuttcap%
\pgfsetmiterjoin%
\definecolor{currentfill}{rgb}{0.000000,0.000000,0.000000}%
\pgfsetfillcolor{currentfill}%
\pgfsetlinewidth{0.000000pt}%
\definecolor{currentstroke}{rgb}{0.000000,0.000000,0.000000}%
\pgfsetstrokecolor{currentstroke}%
\pgfsetstrokeopacity{0.000000}%
\pgfsetdash{}{0pt}%
\pgfpathmoveto{\pgfqpoint{3.106900in}{0.660000in}}%
\pgfpathlineto{\pgfqpoint{3.113100in}{0.660000in}}%
\pgfpathlineto{\pgfqpoint{3.113100in}{3.371166in}}%
\pgfpathlineto{\pgfqpoint{3.106900in}{3.371166in}}%
\pgfpathlineto{\pgfqpoint{3.106900in}{0.660000in}}%
\pgfpathclose%
\pgfusepath{fill}%
\end{pgfscope}%
\begin{pgfscope}%
\pgfpathrectangle{\pgfqpoint{1.250000in}{0.660000in}}{\pgfqpoint{7.750000in}{4.620000in}}%
\pgfusepath{clip}%
\pgfsetbuttcap%
\pgfsetmiterjoin%
\definecolor{currentfill}{rgb}{0.000000,0.000000,0.000000}%
\pgfsetfillcolor{currentfill}%
\pgfsetlinewidth{0.000000pt}%
\definecolor{currentstroke}{rgb}{0.000000,0.000000,0.000000}%
\pgfsetstrokecolor{currentstroke}%
\pgfsetstrokeopacity{0.000000}%
\pgfsetdash{}{0pt}%
\pgfpathmoveto{\pgfqpoint{3.114650in}{0.660000in}}%
\pgfpathlineto{\pgfqpoint{3.120850in}{0.660000in}}%
\pgfpathlineto{\pgfqpoint{3.120850in}{3.373136in}}%
\pgfpathlineto{\pgfqpoint{3.114650in}{3.373136in}}%
\pgfpathlineto{\pgfqpoint{3.114650in}{0.660000in}}%
\pgfpathclose%
\pgfusepath{fill}%
\end{pgfscope}%
\begin{pgfscope}%
\pgfpathrectangle{\pgfqpoint{1.250000in}{0.660000in}}{\pgfqpoint{7.750000in}{4.620000in}}%
\pgfusepath{clip}%
\pgfsetbuttcap%
\pgfsetmiterjoin%
\definecolor{currentfill}{rgb}{0.000000,0.000000,0.000000}%
\pgfsetfillcolor{currentfill}%
\pgfsetlinewidth{0.000000pt}%
\definecolor{currentstroke}{rgb}{0.000000,0.000000,0.000000}%
\pgfsetstrokecolor{currentstroke}%
\pgfsetstrokeopacity{0.000000}%
\pgfsetdash{}{0pt}%
\pgfpathmoveto{\pgfqpoint{3.122400in}{0.660000in}}%
\pgfpathlineto{\pgfqpoint{3.128600in}{0.660000in}}%
\pgfpathlineto{\pgfqpoint{3.128600in}{3.394492in}}%
\pgfpathlineto{\pgfqpoint{3.122400in}{3.394492in}}%
\pgfpathlineto{\pgfqpoint{3.122400in}{0.660000in}}%
\pgfpathclose%
\pgfusepath{fill}%
\end{pgfscope}%
\begin{pgfscope}%
\pgfpathrectangle{\pgfqpoint{1.250000in}{0.660000in}}{\pgfqpoint{7.750000in}{4.620000in}}%
\pgfusepath{clip}%
\pgfsetbuttcap%
\pgfsetmiterjoin%
\definecolor{currentfill}{rgb}{0.000000,0.000000,0.000000}%
\pgfsetfillcolor{currentfill}%
\pgfsetlinewidth{0.000000pt}%
\definecolor{currentstroke}{rgb}{0.000000,0.000000,0.000000}%
\pgfsetstrokecolor{currentstroke}%
\pgfsetstrokeopacity{0.000000}%
\pgfsetdash{}{0pt}%
\pgfpathmoveto{\pgfqpoint{3.130150in}{0.660000in}}%
\pgfpathlineto{\pgfqpoint{3.136350in}{0.660000in}}%
\pgfpathlineto{\pgfqpoint{3.136350in}{3.456631in}}%
\pgfpathlineto{\pgfqpoint{3.130150in}{3.456631in}}%
\pgfpathlineto{\pgfqpoint{3.130150in}{0.660000in}}%
\pgfpathclose%
\pgfusepath{fill}%
\end{pgfscope}%
\begin{pgfscope}%
\pgfpathrectangle{\pgfqpoint{1.250000in}{0.660000in}}{\pgfqpoint{7.750000in}{4.620000in}}%
\pgfusepath{clip}%
\pgfsetbuttcap%
\pgfsetmiterjoin%
\definecolor{currentfill}{rgb}{0.000000,0.000000,0.000000}%
\pgfsetfillcolor{currentfill}%
\pgfsetlinewidth{0.000000pt}%
\definecolor{currentstroke}{rgb}{0.000000,0.000000,0.000000}%
\pgfsetstrokecolor{currentstroke}%
\pgfsetstrokeopacity{0.000000}%
\pgfsetdash{}{0pt}%
\pgfpathmoveto{\pgfqpoint{3.137900in}{0.660000in}}%
\pgfpathlineto{\pgfqpoint{3.144100in}{0.660000in}}%
\pgfpathlineto{\pgfqpoint{3.144100in}{3.375197in}}%
\pgfpathlineto{\pgfqpoint{3.137900in}{3.375197in}}%
\pgfpathlineto{\pgfqpoint{3.137900in}{0.660000in}}%
\pgfpathclose%
\pgfusepath{fill}%
\end{pgfscope}%
\begin{pgfscope}%
\pgfpathrectangle{\pgfqpoint{1.250000in}{0.660000in}}{\pgfqpoint{7.750000in}{4.620000in}}%
\pgfusepath{clip}%
\pgfsetbuttcap%
\pgfsetmiterjoin%
\definecolor{currentfill}{rgb}{0.000000,0.000000,0.000000}%
\pgfsetfillcolor{currentfill}%
\pgfsetlinewidth{0.000000pt}%
\definecolor{currentstroke}{rgb}{0.000000,0.000000,0.000000}%
\pgfsetstrokecolor{currentstroke}%
\pgfsetstrokeopacity{0.000000}%
\pgfsetdash{}{0pt}%
\pgfpathmoveto{\pgfqpoint{3.145650in}{0.660000in}}%
\pgfpathlineto{\pgfqpoint{3.151850in}{0.660000in}}%
\pgfpathlineto{\pgfqpoint{3.151850in}{3.391887in}}%
\pgfpathlineto{\pgfqpoint{3.145650in}{3.391887in}}%
\pgfpathlineto{\pgfqpoint{3.145650in}{0.660000in}}%
\pgfpathclose%
\pgfusepath{fill}%
\end{pgfscope}%
\begin{pgfscope}%
\pgfpathrectangle{\pgfqpoint{1.250000in}{0.660000in}}{\pgfqpoint{7.750000in}{4.620000in}}%
\pgfusepath{clip}%
\pgfsetbuttcap%
\pgfsetmiterjoin%
\definecolor{currentfill}{rgb}{0.000000,0.000000,0.000000}%
\pgfsetfillcolor{currentfill}%
\pgfsetlinewidth{0.000000pt}%
\definecolor{currentstroke}{rgb}{0.000000,0.000000,0.000000}%
\pgfsetstrokecolor{currentstroke}%
\pgfsetstrokeopacity{0.000000}%
\pgfsetdash{}{0pt}%
\pgfpathmoveto{\pgfqpoint{3.153400in}{0.660000in}}%
\pgfpathlineto{\pgfqpoint{3.159600in}{0.660000in}}%
\pgfpathlineto{\pgfqpoint{3.159600in}{3.369256in}}%
\pgfpathlineto{\pgfqpoint{3.153400in}{3.369256in}}%
\pgfpathlineto{\pgfqpoint{3.153400in}{0.660000in}}%
\pgfpathclose%
\pgfusepath{fill}%
\end{pgfscope}%
\begin{pgfscope}%
\pgfpathrectangle{\pgfqpoint{1.250000in}{0.660000in}}{\pgfqpoint{7.750000in}{4.620000in}}%
\pgfusepath{clip}%
\pgfsetbuttcap%
\pgfsetmiterjoin%
\definecolor{currentfill}{rgb}{0.000000,0.000000,0.000000}%
\pgfsetfillcolor{currentfill}%
\pgfsetlinewidth{0.000000pt}%
\definecolor{currentstroke}{rgb}{0.000000,0.000000,0.000000}%
\pgfsetstrokecolor{currentstroke}%
\pgfsetstrokeopacity{0.000000}%
\pgfsetdash{}{0pt}%
\pgfpathmoveto{\pgfqpoint{3.161150in}{0.660000in}}%
\pgfpathlineto{\pgfqpoint{3.167350in}{0.660000in}}%
\pgfpathlineto{\pgfqpoint{3.167350in}{3.386193in}}%
\pgfpathlineto{\pgfqpoint{3.161150in}{3.386193in}}%
\pgfpathlineto{\pgfqpoint{3.161150in}{0.660000in}}%
\pgfpathclose%
\pgfusepath{fill}%
\end{pgfscope}%
\begin{pgfscope}%
\pgfpathrectangle{\pgfqpoint{1.250000in}{0.660000in}}{\pgfqpoint{7.750000in}{4.620000in}}%
\pgfusepath{clip}%
\pgfsetbuttcap%
\pgfsetmiterjoin%
\definecolor{currentfill}{rgb}{0.000000,0.000000,0.000000}%
\pgfsetfillcolor{currentfill}%
\pgfsetlinewidth{0.000000pt}%
\definecolor{currentstroke}{rgb}{0.000000,0.000000,0.000000}%
\pgfsetstrokecolor{currentstroke}%
\pgfsetstrokeopacity{0.000000}%
\pgfsetdash{}{0pt}%
\pgfpathmoveto{\pgfqpoint{3.168900in}{0.660000in}}%
\pgfpathlineto{\pgfqpoint{3.175100in}{0.660000in}}%
\pgfpathlineto{\pgfqpoint{3.175100in}{3.418571in}}%
\pgfpathlineto{\pgfqpoint{3.168900in}{3.418571in}}%
\pgfpathlineto{\pgfqpoint{3.168900in}{0.660000in}}%
\pgfpathclose%
\pgfusepath{fill}%
\end{pgfscope}%
\begin{pgfscope}%
\pgfpathrectangle{\pgfqpoint{1.250000in}{0.660000in}}{\pgfqpoint{7.750000in}{4.620000in}}%
\pgfusepath{clip}%
\pgfsetbuttcap%
\pgfsetmiterjoin%
\definecolor{currentfill}{rgb}{0.000000,0.000000,0.000000}%
\pgfsetfillcolor{currentfill}%
\pgfsetlinewidth{0.000000pt}%
\definecolor{currentstroke}{rgb}{0.000000,0.000000,0.000000}%
\pgfsetstrokecolor{currentstroke}%
\pgfsetstrokeopacity{0.000000}%
\pgfsetdash{}{0pt}%
\pgfpathmoveto{\pgfqpoint{3.176650in}{0.660000in}}%
\pgfpathlineto{\pgfqpoint{3.182850in}{0.660000in}}%
\pgfpathlineto{\pgfqpoint{3.182850in}{3.371743in}}%
\pgfpathlineto{\pgfqpoint{3.176650in}{3.371743in}}%
\pgfpathlineto{\pgfqpoint{3.176650in}{0.660000in}}%
\pgfpathclose%
\pgfusepath{fill}%
\end{pgfscope}%
\begin{pgfscope}%
\pgfpathrectangle{\pgfqpoint{1.250000in}{0.660000in}}{\pgfqpoint{7.750000in}{4.620000in}}%
\pgfusepath{clip}%
\pgfsetbuttcap%
\pgfsetmiterjoin%
\definecolor{currentfill}{rgb}{0.000000,0.000000,0.000000}%
\pgfsetfillcolor{currentfill}%
\pgfsetlinewidth{0.000000pt}%
\definecolor{currentstroke}{rgb}{0.000000,0.000000,0.000000}%
\pgfsetstrokecolor{currentstroke}%
\pgfsetstrokeopacity{0.000000}%
\pgfsetdash{}{0pt}%
\pgfpathmoveto{\pgfqpoint{3.184400in}{0.660000in}}%
\pgfpathlineto{\pgfqpoint{3.190600in}{0.660000in}}%
\pgfpathlineto{\pgfqpoint{3.190600in}{3.358289in}}%
\pgfpathlineto{\pgfqpoint{3.184400in}{3.358289in}}%
\pgfpathlineto{\pgfqpoint{3.184400in}{0.660000in}}%
\pgfpathclose%
\pgfusepath{fill}%
\end{pgfscope}%
\begin{pgfscope}%
\pgfpathrectangle{\pgfqpoint{1.250000in}{0.660000in}}{\pgfqpoint{7.750000in}{4.620000in}}%
\pgfusepath{clip}%
\pgfsetbuttcap%
\pgfsetmiterjoin%
\definecolor{currentfill}{rgb}{0.000000,0.000000,0.000000}%
\pgfsetfillcolor{currentfill}%
\pgfsetlinewidth{0.000000pt}%
\definecolor{currentstroke}{rgb}{0.000000,0.000000,0.000000}%
\pgfsetstrokecolor{currentstroke}%
\pgfsetstrokeopacity{0.000000}%
\pgfsetdash{}{0pt}%
\pgfpathmoveto{\pgfqpoint{3.192150in}{0.660000in}}%
\pgfpathlineto{\pgfqpoint{3.198350in}{0.660000in}}%
\pgfpathlineto{\pgfqpoint{3.198350in}{3.439090in}}%
\pgfpathlineto{\pgfqpoint{3.192150in}{3.439090in}}%
\pgfpathlineto{\pgfqpoint{3.192150in}{0.660000in}}%
\pgfpathclose%
\pgfusepath{fill}%
\end{pgfscope}%
\begin{pgfscope}%
\pgfpathrectangle{\pgfqpoint{1.250000in}{0.660000in}}{\pgfqpoint{7.750000in}{4.620000in}}%
\pgfusepath{clip}%
\pgfsetbuttcap%
\pgfsetmiterjoin%
\definecolor{currentfill}{rgb}{0.000000,0.000000,0.000000}%
\pgfsetfillcolor{currentfill}%
\pgfsetlinewidth{0.000000pt}%
\definecolor{currentstroke}{rgb}{0.000000,0.000000,0.000000}%
\pgfsetstrokecolor{currentstroke}%
\pgfsetstrokeopacity{0.000000}%
\pgfsetdash{}{0pt}%
\pgfpathmoveto{\pgfqpoint{3.199900in}{0.660000in}}%
\pgfpathlineto{\pgfqpoint{3.206100in}{0.660000in}}%
\pgfpathlineto{\pgfqpoint{3.206100in}{3.341483in}}%
\pgfpathlineto{\pgfqpoint{3.199900in}{3.341483in}}%
\pgfpathlineto{\pgfqpoint{3.199900in}{0.660000in}}%
\pgfpathclose%
\pgfusepath{fill}%
\end{pgfscope}%
\begin{pgfscope}%
\pgfpathrectangle{\pgfqpoint{1.250000in}{0.660000in}}{\pgfqpoint{7.750000in}{4.620000in}}%
\pgfusepath{clip}%
\pgfsetbuttcap%
\pgfsetmiterjoin%
\definecolor{currentfill}{rgb}{0.000000,0.000000,0.000000}%
\pgfsetfillcolor{currentfill}%
\pgfsetlinewidth{0.000000pt}%
\definecolor{currentstroke}{rgb}{0.000000,0.000000,0.000000}%
\pgfsetstrokecolor{currentstroke}%
\pgfsetstrokeopacity{0.000000}%
\pgfsetdash{}{0pt}%
\pgfpathmoveto{\pgfqpoint{3.207650in}{0.660000in}}%
\pgfpathlineto{\pgfqpoint{3.213850in}{0.660000in}}%
\pgfpathlineto{\pgfqpoint{3.213850in}{3.383244in}}%
\pgfpathlineto{\pgfqpoint{3.207650in}{3.383244in}}%
\pgfpathlineto{\pgfqpoint{3.207650in}{0.660000in}}%
\pgfpathclose%
\pgfusepath{fill}%
\end{pgfscope}%
\begin{pgfscope}%
\pgfpathrectangle{\pgfqpoint{1.250000in}{0.660000in}}{\pgfqpoint{7.750000in}{4.620000in}}%
\pgfusepath{clip}%
\pgfsetbuttcap%
\pgfsetmiterjoin%
\definecolor{currentfill}{rgb}{0.000000,0.000000,0.000000}%
\pgfsetfillcolor{currentfill}%
\pgfsetlinewidth{0.000000pt}%
\definecolor{currentstroke}{rgb}{0.000000,0.000000,0.000000}%
\pgfsetstrokecolor{currentstroke}%
\pgfsetstrokeopacity{0.000000}%
\pgfsetdash{}{0pt}%
\pgfpathmoveto{\pgfqpoint{3.215400in}{0.660000in}}%
\pgfpathlineto{\pgfqpoint{3.221600in}{0.660000in}}%
\pgfpathlineto{\pgfqpoint{3.221600in}{3.471904in}}%
\pgfpathlineto{\pgfqpoint{3.215400in}{3.471904in}}%
\pgfpathlineto{\pgfqpoint{3.215400in}{0.660000in}}%
\pgfpathclose%
\pgfusepath{fill}%
\end{pgfscope}%
\begin{pgfscope}%
\pgfpathrectangle{\pgfqpoint{1.250000in}{0.660000in}}{\pgfqpoint{7.750000in}{4.620000in}}%
\pgfusepath{clip}%
\pgfsetbuttcap%
\pgfsetmiterjoin%
\definecolor{currentfill}{rgb}{0.000000,0.000000,0.000000}%
\pgfsetfillcolor{currentfill}%
\pgfsetlinewidth{0.000000pt}%
\definecolor{currentstroke}{rgb}{0.000000,0.000000,0.000000}%
\pgfsetstrokecolor{currentstroke}%
\pgfsetstrokeopacity{0.000000}%
\pgfsetdash{}{0pt}%
\pgfpathmoveto{\pgfqpoint{3.223150in}{0.660000in}}%
\pgfpathlineto{\pgfqpoint{3.229350in}{0.660000in}}%
\pgfpathlineto{\pgfqpoint{3.229350in}{3.371828in}}%
\pgfpathlineto{\pgfqpoint{3.223150in}{3.371828in}}%
\pgfpathlineto{\pgfqpoint{3.223150in}{0.660000in}}%
\pgfpathclose%
\pgfusepath{fill}%
\end{pgfscope}%
\begin{pgfscope}%
\pgfpathrectangle{\pgfqpoint{1.250000in}{0.660000in}}{\pgfqpoint{7.750000in}{4.620000in}}%
\pgfusepath{clip}%
\pgfsetbuttcap%
\pgfsetmiterjoin%
\definecolor{currentfill}{rgb}{0.000000,0.000000,0.000000}%
\pgfsetfillcolor{currentfill}%
\pgfsetlinewidth{0.000000pt}%
\definecolor{currentstroke}{rgb}{0.000000,0.000000,0.000000}%
\pgfsetstrokecolor{currentstroke}%
\pgfsetstrokeopacity{0.000000}%
\pgfsetdash{}{0pt}%
\pgfpathmoveto{\pgfqpoint{3.230900in}{0.660000in}}%
\pgfpathlineto{\pgfqpoint{3.237100in}{0.660000in}}%
\pgfpathlineto{\pgfqpoint{3.237100in}{3.403109in}}%
\pgfpathlineto{\pgfqpoint{3.230900in}{3.403109in}}%
\pgfpathlineto{\pgfqpoint{3.230900in}{0.660000in}}%
\pgfpathclose%
\pgfusepath{fill}%
\end{pgfscope}%
\begin{pgfscope}%
\pgfpathrectangle{\pgfqpoint{1.250000in}{0.660000in}}{\pgfqpoint{7.750000in}{4.620000in}}%
\pgfusepath{clip}%
\pgfsetbuttcap%
\pgfsetmiterjoin%
\definecolor{currentfill}{rgb}{0.000000,0.000000,0.000000}%
\pgfsetfillcolor{currentfill}%
\pgfsetlinewidth{0.000000pt}%
\definecolor{currentstroke}{rgb}{0.000000,0.000000,0.000000}%
\pgfsetstrokecolor{currentstroke}%
\pgfsetstrokeopacity{0.000000}%
\pgfsetdash{}{0pt}%
\pgfpathmoveto{\pgfqpoint{3.238650in}{0.660000in}}%
\pgfpathlineto{\pgfqpoint{3.244850in}{0.660000in}}%
\pgfpathlineto{\pgfqpoint{3.244850in}{3.366785in}}%
\pgfpathlineto{\pgfqpoint{3.238650in}{3.366785in}}%
\pgfpathlineto{\pgfqpoint{3.238650in}{0.660000in}}%
\pgfpathclose%
\pgfusepath{fill}%
\end{pgfscope}%
\begin{pgfscope}%
\pgfpathrectangle{\pgfqpoint{1.250000in}{0.660000in}}{\pgfqpoint{7.750000in}{4.620000in}}%
\pgfusepath{clip}%
\pgfsetbuttcap%
\pgfsetmiterjoin%
\definecolor{currentfill}{rgb}{0.000000,0.000000,0.000000}%
\pgfsetfillcolor{currentfill}%
\pgfsetlinewidth{0.000000pt}%
\definecolor{currentstroke}{rgb}{0.000000,0.000000,0.000000}%
\pgfsetstrokecolor{currentstroke}%
\pgfsetstrokeopacity{0.000000}%
\pgfsetdash{}{0pt}%
\pgfpathmoveto{\pgfqpoint{3.246400in}{0.660000in}}%
\pgfpathlineto{\pgfqpoint{3.252600in}{0.660000in}}%
\pgfpathlineto{\pgfqpoint{3.252600in}{3.414311in}}%
\pgfpathlineto{\pgfqpoint{3.246400in}{3.414311in}}%
\pgfpathlineto{\pgfqpoint{3.246400in}{0.660000in}}%
\pgfpathclose%
\pgfusepath{fill}%
\end{pgfscope}%
\begin{pgfscope}%
\pgfpathrectangle{\pgfqpoint{1.250000in}{0.660000in}}{\pgfqpoint{7.750000in}{4.620000in}}%
\pgfusepath{clip}%
\pgfsetbuttcap%
\pgfsetmiterjoin%
\definecolor{currentfill}{rgb}{0.000000,0.000000,0.000000}%
\pgfsetfillcolor{currentfill}%
\pgfsetlinewidth{0.000000pt}%
\definecolor{currentstroke}{rgb}{0.000000,0.000000,0.000000}%
\pgfsetstrokecolor{currentstroke}%
\pgfsetstrokeopacity{0.000000}%
\pgfsetdash{}{0pt}%
\pgfpathmoveto{\pgfqpoint{3.254150in}{0.660000in}}%
\pgfpathlineto{\pgfqpoint{3.260350in}{0.660000in}}%
\pgfpathlineto{\pgfqpoint{3.260350in}{3.395636in}}%
\pgfpathlineto{\pgfqpoint{3.254150in}{3.395636in}}%
\pgfpathlineto{\pgfqpoint{3.254150in}{0.660000in}}%
\pgfpathclose%
\pgfusepath{fill}%
\end{pgfscope}%
\begin{pgfscope}%
\pgfpathrectangle{\pgfqpoint{1.250000in}{0.660000in}}{\pgfqpoint{7.750000in}{4.620000in}}%
\pgfusepath{clip}%
\pgfsetbuttcap%
\pgfsetmiterjoin%
\definecolor{currentfill}{rgb}{0.000000,0.000000,0.000000}%
\pgfsetfillcolor{currentfill}%
\pgfsetlinewidth{0.000000pt}%
\definecolor{currentstroke}{rgb}{0.000000,0.000000,0.000000}%
\pgfsetstrokecolor{currentstroke}%
\pgfsetstrokeopacity{0.000000}%
\pgfsetdash{}{0pt}%
\pgfpathmoveto{\pgfqpoint{3.261900in}{0.660000in}}%
\pgfpathlineto{\pgfqpoint{3.268100in}{0.660000in}}%
\pgfpathlineto{\pgfqpoint{3.268100in}{3.366123in}}%
\pgfpathlineto{\pgfqpoint{3.261900in}{3.366123in}}%
\pgfpathlineto{\pgfqpoint{3.261900in}{0.660000in}}%
\pgfpathclose%
\pgfusepath{fill}%
\end{pgfscope}%
\begin{pgfscope}%
\pgfpathrectangle{\pgfqpoint{1.250000in}{0.660000in}}{\pgfqpoint{7.750000in}{4.620000in}}%
\pgfusepath{clip}%
\pgfsetbuttcap%
\pgfsetmiterjoin%
\definecolor{currentfill}{rgb}{0.000000,0.000000,0.000000}%
\pgfsetfillcolor{currentfill}%
\pgfsetlinewidth{0.000000pt}%
\definecolor{currentstroke}{rgb}{0.000000,0.000000,0.000000}%
\pgfsetstrokecolor{currentstroke}%
\pgfsetstrokeopacity{0.000000}%
\pgfsetdash{}{0pt}%
\pgfpathmoveto{\pgfqpoint{3.269650in}{0.660000in}}%
\pgfpathlineto{\pgfqpoint{3.275850in}{0.660000in}}%
\pgfpathlineto{\pgfqpoint{3.275850in}{3.378361in}}%
\pgfpathlineto{\pgfqpoint{3.269650in}{3.378361in}}%
\pgfpathlineto{\pgfqpoint{3.269650in}{0.660000in}}%
\pgfpathclose%
\pgfusepath{fill}%
\end{pgfscope}%
\begin{pgfscope}%
\pgfpathrectangle{\pgfqpoint{1.250000in}{0.660000in}}{\pgfqpoint{7.750000in}{4.620000in}}%
\pgfusepath{clip}%
\pgfsetbuttcap%
\pgfsetmiterjoin%
\definecolor{currentfill}{rgb}{0.000000,0.000000,0.000000}%
\pgfsetfillcolor{currentfill}%
\pgfsetlinewidth{0.000000pt}%
\definecolor{currentstroke}{rgb}{0.000000,0.000000,0.000000}%
\pgfsetstrokecolor{currentstroke}%
\pgfsetstrokeopacity{0.000000}%
\pgfsetdash{}{0pt}%
\pgfpathmoveto{\pgfqpoint{3.277400in}{0.660000in}}%
\pgfpathlineto{\pgfqpoint{3.283600in}{0.660000in}}%
\pgfpathlineto{\pgfqpoint{3.283600in}{3.399585in}}%
\pgfpathlineto{\pgfqpoint{3.277400in}{3.399585in}}%
\pgfpathlineto{\pgfqpoint{3.277400in}{0.660000in}}%
\pgfpathclose%
\pgfusepath{fill}%
\end{pgfscope}%
\begin{pgfscope}%
\pgfpathrectangle{\pgfqpoint{1.250000in}{0.660000in}}{\pgfqpoint{7.750000in}{4.620000in}}%
\pgfusepath{clip}%
\pgfsetbuttcap%
\pgfsetmiterjoin%
\definecolor{currentfill}{rgb}{0.000000,0.000000,0.000000}%
\pgfsetfillcolor{currentfill}%
\pgfsetlinewidth{0.000000pt}%
\definecolor{currentstroke}{rgb}{0.000000,0.000000,0.000000}%
\pgfsetstrokecolor{currentstroke}%
\pgfsetstrokeopacity{0.000000}%
\pgfsetdash{}{0pt}%
\pgfpathmoveto{\pgfqpoint{3.285150in}{0.660000in}}%
\pgfpathlineto{\pgfqpoint{3.291350in}{0.660000in}}%
\pgfpathlineto{\pgfqpoint{3.291350in}{3.389116in}}%
\pgfpathlineto{\pgfqpoint{3.285150in}{3.389116in}}%
\pgfpathlineto{\pgfqpoint{3.285150in}{0.660000in}}%
\pgfpathclose%
\pgfusepath{fill}%
\end{pgfscope}%
\begin{pgfscope}%
\pgfpathrectangle{\pgfqpoint{1.250000in}{0.660000in}}{\pgfqpoint{7.750000in}{4.620000in}}%
\pgfusepath{clip}%
\pgfsetbuttcap%
\pgfsetmiterjoin%
\definecolor{currentfill}{rgb}{0.000000,0.000000,0.000000}%
\pgfsetfillcolor{currentfill}%
\pgfsetlinewidth{0.000000pt}%
\definecolor{currentstroke}{rgb}{0.000000,0.000000,0.000000}%
\pgfsetstrokecolor{currentstroke}%
\pgfsetstrokeopacity{0.000000}%
\pgfsetdash{}{0pt}%
\pgfpathmoveto{\pgfqpoint{3.292900in}{0.660000in}}%
\pgfpathlineto{\pgfqpoint{3.299100in}{0.660000in}}%
\pgfpathlineto{\pgfqpoint{3.299100in}{3.364645in}}%
\pgfpathlineto{\pgfqpoint{3.292900in}{3.364645in}}%
\pgfpathlineto{\pgfqpoint{3.292900in}{0.660000in}}%
\pgfpathclose%
\pgfusepath{fill}%
\end{pgfscope}%
\begin{pgfscope}%
\pgfpathrectangle{\pgfqpoint{1.250000in}{0.660000in}}{\pgfqpoint{7.750000in}{4.620000in}}%
\pgfusepath{clip}%
\pgfsetbuttcap%
\pgfsetmiterjoin%
\definecolor{currentfill}{rgb}{0.000000,0.000000,0.000000}%
\pgfsetfillcolor{currentfill}%
\pgfsetlinewidth{0.000000pt}%
\definecolor{currentstroke}{rgb}{0.000000,0.000000,0.000000}%
\pgfsetstrokecolor{currentstroke}%
\pgfsetstrokeopacity{0.000000}%
\pgfsetdash{}{0pt}%
\pgfpathmoveto{\pgfqpoint{3.300650in}{0.660000in}}%
\pgfpathlineto{\pgfqpoint{3.306850in}{0.660000in}}%
\pgfpathlineto{\pgfqpoint{3.306850in}{3.399366in}}%
\pgfpathlineto{\pgfqpoint{3.300650in}{3.399366in}}%
\pgfpathlineto{\pgfqpoint{3.300650in}{0.660000in}}%
\pgfpathclose%
\pgfusepath{fill}%
\end{pgfscope}%
\begin{pgfscope}%
\pgfpathrectangle{\pgfqpoint{1.250000in}{0.660000in}}{\pgfqpoint{7.750000in}{4.620000in}}%
\pgfusepath{clip}%
\pgfsetbuttcap%
\pgfsetmiterjoin%
\definecolor{currentfill}{rgb}{0.000000,0.000000,0.000000}%
\pgfsetfillcolor{currentfill}%
\pgfsetlinewidth{0.000000pt}%
\definecolor{currentstroke}{rgb}{0.000000,0.000000,0.000000}%
\pgfsetstrokecolor{currentstroke}%
\pgfsetstrokeopacity{0.000000}%
\pgfsetdash{}{0pt}%
\pgfpathmoveto{\pgfqpoint{3.308400in}{0.660000in}}%
\pgfpathlineto{\pgfqpoint{3.314600in}{0.660000in}}%
\pgfpathlineto{\pgfqpoint{3.314600in}{3.435273in}}%
\pgfpathlineto{\pgfqpoint{3.308400in}{3.435273in}}%
\pgfpathlineto{\pgfqpoint{3.308400in}{0.660000in}}%
\pgfpathclose%
\pgfusepath{fill}%
\end{pgfscope}%
\begin{pgfscope}%
\pgfpathrectangle{\pgfqpoint{1.250000in}{0.660000in}}{\pgfqpoint{7.750000in}{4.620000in}}%
\pgfusepath{clip}%
\pgfsetbuttcap%
\pgfsetmiterjoin%
\definecolor{currentfill}{rgb}{0.000000,0.000000,0.000000}%
\pgfsetfillcolor{currentfill}%
\pgfsetlinewidth{0.000000pt}%
\definecolor{currentstroke}{rgb}{0.000000,0.000000,0.000000}%
\pgfsetstrokecolor{currentstroke}%
\pgfsetstrokeopacity{0.000000}%
\pgfsetdash{}{0pt}%
\pgfpathmoveto{\pgfqpoint{3.316150in}{0.660000in}}%
\pgfpathlineto{\pgfqpoint{3.322350in}{0.660000in}}%
\pgfpathlineto{\pgfqpoint{3.322350in}{3.432577in}}%
\pgfpathlineto{\pgfqpoint{3.316150in}{3.432577in}}%
\pgfpathlineto{\pgfqpoint{3.316150in}{0.660000in}}%
\pgfpathclose%
\pgfusepath{fill}%
\end{pgfscope}%
\begin{pgfscope}%
\pgfpathrectangle{\pgfqpoint{1.250000in}{0.660000in}}{\pgfqpoint{7.750000in}{4.620000in}}%
\pgfusepath{clip}%
\pgfsetbuttcap%
\pgfsetmiterjoin%
\definecolor{currentfill}{rgb}{0.000000,0.000000,0.000000}%
\pgfsetfillcolor{currentfill}%
\pgfsetlinewidth{0.000000pt}%
\definecolor{currentstroke}{rgb}{0.000000,0.000000,0.000000}%
\pgfsetstrokecolor{currentstroke}%
\pgfsetstrokeopacity{0.000000}%
\pgfsetdash{}{0pt}%
\pgfpathmoveto{\pgfqpoint{3.323900in}{0.660000in}}%
\pgfpathlineto{\pgfqpoint{3.330100in}{0.660000in}}%
\pgfpathlineto{\pgfqpoint{3.330100in}{3.386441in}}%
\pgfpathlineto{\pgfqpoint{3.323900in}{3.386441in}}%
\pgfpathlineto{\pgfqpoint{3.323900in}{0.660000in}}%
\pgfpathclose%
\pgfusepath{fill}%
\end{pgfscope}%
\begin{pgfscope}%
\pgfpathrectangle{\pgfqpoint{1.250000in}{0.660000in}}{\pgfqpoint{7.750000in}{4.620000in}}%
\pgfusepath{clip}%
\pgfsetbuttcap%
\pgfsetmiterjoin%
\definecolor{currentfill}{rgb}{0.000000,0.000000,0.000000}%
\pgfsetfillcolor{currentfill}%
\pgfsetlinewidth{0.000000pt}%
\definecolor{currentstroke}{rgb}{0.000000,0.000000,0.000000}%
\pgfsetstrokecolor{currentstroke}%
\pgfsetstrokeopacity{0.000000}%
\pgfsetdash{}{0pt}%
\pgfpathmoveto{\pgfqpoint{3.331650in}{0.660000in}}%
\pgfpathlineto{\pgfqpoint{3.337850in}{0.660000in}}%
\pgfpathlineto{\pgfqpoint{3.337850in}{3.345693in}}%
\pgfpathlineto{\pgfqpoint{3.331650in}{3.345693in}}%
\pgfpathlineto{\pgfqpoint{3.331650in}{0.660000in}}%
\pgfpathclose%
\pgfusepath{fill}%
\end{pgfscope}%
\begin{pgfscope}%
\pgfpathrectangle{\pgfqpoint{1.250000in}{0.660000in}}{\pgfqpoint{7.750000in}{4.620000in}}%
\pgfusepath{clip}%
\pgfsetbuttcap%
\pgfsetmiterjoin%
\definecolor{currentfill}{rgb}{0.000000,0.000000,0.000000}%
\pgfsetfillcolor{currentfill}%
\pgfsetlinewidth{0.000000pt}%
\definecolor{currentstroke}{rgb}{0.000000,0.000000,0.000000}%
\pgfsetstrokecolor{currentstroke}%
\pgfsetstrokeopacity{0.000000}%
\pgfsetdash{}{0pt}%
\pgfpathmoveto{\pgfqpoint{3.339400in}{0.660000in}}%
\pgfpathlineto{\pgfqpoint{3.345600in}{0.660000in}}%
\pgfpathlineto{\pgfqpoint{3.345600in}{3.369929in}}%
\pgfpathlineto{\pgfqpoint{3.339400in}{3.369929in}}%
\pgfpathlineto{\pgfqpoint{3.339400in}{0.660000in}}%
\pgfpathclose%
\pgfusepath{fill}%
\end{pgfscope}%
\begin{pgfscope}%
\pgfpathrectangle{\pgfqpoint{1.250000in}{0.660000in}}{\pgfqpoint{7.750000in}{4.620000in}}%
\pgfusepath{clip}%
\pgfsetbuttcap%
\pgfsetmiterjoin%
\definecolor{currentfill}{rgb}{0.000000,0.000000,0.000000}%
\pgfsetfillcolor{currentfill}%
\pgfsetlinewidth{0.000000pt}%
\definecolor{currentstroke}{rgb}{0.000000,0.000000,0.000000}%
\pgfsetstrokecolor{currentstroke}%
\pgfsetstrokeopacity{0.000000}%
\pgfsetdash{}{0pt}%
\pgfpathmoveto{\pgfqpoint{3.347150in}{0.660000in}}%
\pgfpathlineto{\pgfqpoint{3.353350in}{0.660000in}}%
\pgfpathlineto{\pgfqpoint{3.353350in}{3.357626in}}%
\pgfpathlineto{\pgfqpoint{3.347150in}{3.357626in}}%
\pgfpathlineto{\pgfqpoint{3.347150in}{0.660000in}}%
\pgfpathclose%
\pgfusepath{fill}%
\end{pgfscope}%
\begin{pgfscope}%
\pgfpathrectangle{\pgfqpoint{1.250000in}{0.660000in}}{\pgfqpoint{7.750000in}{4.620000in}}%
\pgfusepath{clip}%
\pgfsetbuttcap%
\pgfsetmiterjoin%
\definecolor{currentfill}{rgb}{0.000000,0.000000,0.000000}%
\pgfsetfillcolor{currentfill}%
\pgfsetlinewidth{0.000000pt}%
\definecolor{currentstroke}{rgb}{0.000000,0.000000,0.000000}%
\pgfsetstrokecolor{currentstroke}%
\pgfsetstrokeopacity{0.000000}%
\pgfsetdash{}{0pt}%
\pgfpathmoveto{\pgfqpoint{3.354900in}{0.660000in}}%
\pgfpathlineto{\pgfqpoint{3.361100in}{0.660000in}}%
\pgfpathlineto{\pgfqpoint{3.361100in}{3.380387in}}%
\pgfpathlineto{\pgfqpoint{3.354900in}{3.380387in}}%
\pgfpathlineto{\pgfqpoint{3.354900in}{0.660000in}}%
\pgfpathclose%
\pgfusepath{fill}%
\end{pgfscope}%
\begin{pgfscope}%
\pgfpathrectangle{\pgfqpoint{1.250000in}{0.660000in}}{\pgfqpoint{7.750000in}{4.620000in}}%
\pgfusepath{clip}%
\pgfsetbuttcap%
\pgfsetmiterjoin%
\definecolor{currentfill}{rgb}{0.000000,0.000000,0.000000}%
\pgfsetfillcolor{currentfill}%
\pgfsetlinewidth{0.000000pt}%
\definecolor{currentstroke}{rgb}{0.000000,0.000000,0.000000}%
\pgfsetstrokecolor{currentstroke}%
\pgfsetstrokeopacity{0.000000}%
\pgfsetdash{}{0pt}%
\pgfpathmoveto{\pgfqpoint{3.362650in}{0.660000in}}%
\pgfpathlineto{\pgfqpoint{3.368850in}{0.660000in}}%
\pgfpathlineto{\pgfqpoint{3.368850in}{3.340577in}}%
\pgfpathlineto{\pgfqpoint{3.362650in}{3.340577in}}%
\pgfpathlineto{\pgfqpoint{3.362650in}{0.660000in}}%
\pgfpathclose%
\pgfusepath{fill}%
\end{pgfscope}%
\begin{pgfscope}%
\pgfpathrectangle{\pgfqpoint{1.250000in}{0.660000in}}{\pgfqpoint{7.750000in}{4.620000in}}%
\pgfusepath{clip}%
\pgfsetbuttcap%
\pgfsetmiterjoin%
\definecolor{currentfill}{rgb}{0.000000,0.000000,0.000000}%
\pgfsetfillcolor{currentfill}%
\pgfsetlinewidth{0.000000pt}%
\definecolor{currentstroke}{rgb}{0.000000,0.000000,0.000000}%
\pgfsetstrokecolor{currentstroke}%
\pgfsetstrokeopacity{0.000000}%
\pgfsetdash{}{0pt}%
\pgfpathmoveto{\pgfqpoint{3.370400in}{0.660000in}}%
\pgfpathlineto{\pgfqpoint{3.376600in}{0.660000in}}%
\pgfpathlineto{\pgfqpoint{3.376600in}{3.411359in}}%
\pgfpathlineto{\pgfqpoint{3.370400in}{3.411359in}}%
\pgfpathlineto{\pgfqpoint{3.370400in}{0.660000in}}%
\pgfpathclose%
\pgfusepath{fill}%
\end{pgfscope}%
\begin{pgfscope}%
\pgfpathrectangle{\pgfqpoint{1.250000in}{0.660000in}}{\pgfqpoint{7.750000in}{4.620000in}}%
\pgfusepath{clip}%
\pgfsetbuttcap%
\pgfsetmiterjoin%
\definecolor{currentfill}{rgb}{0.000000,0.000000,0.000000}%
\pgfsetfillcolor{currentfill}%
\pgfsetlinewidth{0.000000pt}%
\definecolor{currentstroke}{rgb}{0.000000,0.000000,0.000000}%
\pgfsetstrokecolor{currentstroke}%
\pgfsetstrokeopacity{0.000000}%
\pgfsetdash{}{0pt}%
\pgfpathmoveto{\pgfqpoint{3.378150in}{0.660000in}}%
\pgfpathlineto{\pgfqpoint{3.384350in}{0.660000in}}%
\pgfpathlineto{\pgfqpoint{3.384350in}{3.364521in}}%
\pgfpathlineto{\pgfqpoint{3.378150in}{3.364521in}}%
\pgfpathlineto{\pgfqpoint{3.378150in}{0.660000in}}%
\pgfpathclose%
\pgfusepath{fill}%
\end{pgfscope}%
\begin{pgfscope}%
\pgfpathrectangle{\pgfqpoint{1.250000in}{0.660000in}}{\pgfqpoint{7.750000in}{4.620000in}}%
\pgfusepath{clip}%
\pgfsetbuttcap%
\pgfsetmiterjoin%
\definecolor{currentfill}{rgb}{0.000000,0.000000,0.000000}%
\pgfsetfillcolor{currentfill}%
\pgfsetlinewidth{0.000000pt}%
\definecolor{currentstroke}{rgb}{0.000000,0.000000,0.000000}%
\pgfsetstrokecolor{currentstroke}%
\pgfsetstrokeopacity{0.000000}%
\pgfsetdash{}{0pt}%
\pgfpathmoveto{\pgfqpoint{3.385900in}{0.660000in}}%
\pgfpathlineto{\pgfqpoint{3.392100in}{0.660000in}}%
\pgfpathlineto{\pgfqpoint{3.392100in}{3.360329in}}%
\pgfpathlineto{\pgfqpoint{3.385900in}{3.360329in}}%
\pgfpathlineto{\pgfqpoint{3.385900in}{0.660000in}}%
\pgfpathclose%
\pgfusepath{fill}%
\end{pgfscope}%
\begin{pgfscope}%
\pgfpathrectangle{\pgfqpoint{1.250000in}{0.660000in}}{\pgfqpoint{7.750000in}{4.620000in}}%
\pgfusepath{clip}%
\pgfsetbuttcap%
\pgfsetmiterjoin%
\definecolor{currentfill}{rgb}{0.000000,0.000000,0.000000}%
\pgfsetfillcolor{currentfill}%
\pgfsetlinewidth{0.000000pt}%
\definecolor{currentstroke}{rgb}{0.000000,0.000000,0.000000}%
\pgfsetstrokecolor{currentstroke}%
\pgfsetstrokeopacity{0.000000}%
\pgfsetdash{}{0pt}%
\pgfpathmoveto{\pgfqpoint{3.393650in}{0.660000in}}%
\pgfpathlineto{\pgfqpoint{3.399850in}{0.660000in}}%
\pgfpathlineto{\pgfqpoint{3.399850in}{3.383711in}}%
\pgfpathlineto{\pgfqpoint{3.393650in}{3.383711in}}%
\pgfpathlineto{\pgfqpoint{3.393650in}{0.660000in}}%
\pgfpathclose%
\pgfusepath{fill}%
\end{pgfscope}%
\begin{pgfscope}%
\pgfpathrectangle{\pgfqpoint{1.250000in}{0.660000in}}{\pgfqpoint{7.750000in}{4.620000in}}%
\pgfusepath{clip}%
\pgfsetbuttcap%
\pgfsetmiterjoin%
\definecolor{currentfill}{rgb}{0.000000,0.000000,0.000000}%
\pgfsetfillcolor{currentfill}%
\pgfsetlinewidth{0.000000pt}%
\definecolor{currentstroke}{rgb}{0.000000,0.000000,0.000000}%
\pgfsetstrokecolor{currentstroke}%
\pgfsetstrokeopacity{0.000000}%
\pgfsetdash{}{0pt}%
\pgfpathmoveto{\pgfqpoint{3.401400in}{0.660000in}}%
\pgfpathlineto{\pgfqpoint{3.407600in}{0.660000in}}%
\pgfpathlineto{\pgfqpoint{3.407600in}{3.375140in}}%
\pgfpathlineto{\pgfqpoint{3.401400in}{3.375140in}}%
\pgfpathlineto{\pgfqpoint{3.401400in}{0.660000in}}%
\pgfpathclose%
\pgfusepath{fill}%
\end{pgfscope}%
\begin{pgfscope}%
\pgfpathrectangle{\pgfqpoint{1.250000in}{0.660000in}}{\pgfqpoint{7.750000in}{4.620000in}}%
\pgfusepath{clip}%
\pgfsetbuttcap%
\pgfsetmiterjoin%
\definecolor{currentfill}{rgb}{0.000000,0.000000,0.000000}%
\pgfsetfillcolor{currentfill}%
\pgfsetlinewidth{0.000000pt}%
\definecolor{currentstroke}{rgb}{0.000000,0.000000,0.000000}%
\pgfsetstrokecolor{currentstroke}%
\pgfsetstrokeopacity{0.000000}%
\pgfsetdash{}{0pt}%
\pgfpathmoveto{\pgfqpoint{3.409150in}{0.660000in}}%
\pgfpathlineto{\pgfqpoint{3.415350in}{0.660000in}}%
\pgfpathlineto{\pgfqpoint{3.415350in}{3.386815in}}%
\pgfpathlineto{\pgfqpoint{3.409150in}{3.386815in}}%
\pgfpathlineto{\pgfqpoint{3.409150in}{0.660000in}}%
\pgfpathclose%
\pgfusepath{fill}%
\end{pgfscope}%
\begin{pgfscope}%
\pgfpathrectangle{\pgfqpoint{1.250000in}{0.660000in}}{\pgfqpoint{7.750000in}{4.620000in}}%
\pgfusepath{clip}%
\pgfsetbuttcap%
\pgfsetmiterjoin%
\definecolor{currentfill}{rgb}{0.000000,0.000000,0.000000}%
\pgfsetfillcolor{currentfill}%
\pgfsetlinewidth{0.000000pt}%
\definecolor{currentstroke}{rgb}{0.000000,0.000000,0.000000}%
\pgfsetstrokecolor{currentstroke}%
\pgfsetstrokeopacity{0.000000}%
\pgfsetdash{}{0pt}%
\pgfpathmoveto{\pgfqpoint{3.416900in}{0.660000in}}%
\pgfpathlineto{\pgfqpoint{3.423100in}{0.660000in}}%
\pgfpathlineto{\pgfqpoint{3.423100in}{3.413597in}}%
\pgfpathlineto{\pgfqpoint{3.416900in}{3.413597in}}%
\pgfpathlineto{\pgfqpoint{3.416900in}{0.660000in}}%
\pgfpathclose%
\pgfusepath{fill}%
\end{pgfscope}%
\begin{pgfscope}%
\pgfpathrectangle{\pgfqpoint{1.250000in}{0.660000in}}{\pgfqpoint{7.750000in}{4.620000in}}%
\pgfusepath{clip}%
\pgfsetbuttcap%
\pgfsetmiterjoin%
\definecolor{currentfill}{rgb}{0.000000,0.000000,0.000000}%
\pgfsetfillcolor{currentfill}%
\pgfsetlinewidth{0.000000pt}%
\definecolor{currentstroke}{rgb}{0.000000,0.000000,0.000000}%
\pgfsetstrokecolor{currentstroke}%
\pgfsetstrokeopacity{0.000000}%
\pgfsetdash{}{0pt}%
\pgfpathmoveto{\pgfqpoint{3.424650in}{0.660000in}}%
\pgfpathlineto{\pgfqpoint{3.430850in}{0.660000in}}%
\pgfpathlineto{\pgfqpoint{3.430850in}{3.373891in}}%
\pgfpathlineto{\pgfqpoint{3.424650in}{3.373891in}}%
\pgfpathlineto{\pgfqpoint{3.424650in}{0.660000in}}%
\pgfpathclose%
\pgfusepath{fill}%
\end{pgfscope}%
\begin{pgfscope}%
\pgfpathrectangle{\pgfqpoint{1.250000in}{0.660000in}}{\pgfqpoint{7.750000in}{4.620000in}}%
\pgfusepath{clip}%
\pgfsetbuttcap%
\pgfsetmiterjoin%
\definecolor{currentfill}{rgb}{0.000000,0.000000,0.000000}%
\pgfsetfillcolor{currentfill}%
\pgfsetlinewidth{0.000000pt}%
\definecolor{currentstroke}{rgb}{0.000000,0.000000,0.000000}%
\pgfsetstrokecolor{currentstroke}%
\pgfsetstrokeopacity{0.000000}%
\pgfsetdash{}{0pt}%
\pgfpathmoveto{\pgfqpoint{3.432400in}{0.660000in}}%
\pgfpathlineto{\pgfqpoint{3.438600in}{0.660000in}}%
\pgfpathlineto{\pgfqpoint{3.438600in}{3.385467in}}%
\pgfpathlineto{\pgfqpoint{3.432400in}{3.385467in}}%
\pgfpathlineto{\pgfqpoint{3.432400in}{0.660000in}}%
\pgfpathclose%
\pgfusepath{fill}%
\end{pgfscope}%
\begin{pgfscope}%
\pgfpathrectangle{\pgfqpoint{1.250000in}{0.660000in}}{\pgfqpoint{7.750000in}{4.620000in}}%
\pgfusepath{clip}%
\pgfsetbuttcap%
\pgfsetmiterjoin%
\definecolor{currentfill}{rgb}{0.000000,0.000000,0.000000}%
\pgfsetfillcolor{currentfill}%
\pgfsetlinewidth{0.000000pt}%
\definecolor{currentstroke}{rgb}{0.000000,0.000000,0.000000}%
\pgfsetstrokecolor{currentstroke}%
\pgfsetstrokeopacity{0.000000}%
\pgfsetdash{}{0pt}%
\pgfpathmoveto{\pgfqpoint{3.440150in}{0.660000in}}%
\pgfpathlineto{\pgfqpoint{3.446350in}{0.660000in}}%
\pgfpathlineto{\pgfqpoint{3.446350in}{3.351287in}}%
\pgfpathlineto{\pgfqpoint{3.440150in}{3.351287in}}%
\pgfpathlineto{\pgfqpoint{3.440150in}{0.660000in}}%
\pgfpathclose%
\pgfusepath{fill}%
\end{pgfscope}%
\begin{pgfscope}%
\pgfpathrectangle{\pgfqpoint{1.250000in}{0.660000in}}{\pgfqpoint{7.750000in}{4.620000in}}%
\pgfusepath{clip}%
\pgfsetbuttcap%
\pgfsetmiterjoin%
\definecolor{currentfill}{rgb}{0.000000,0.000000,0.000000}%
\pgfsetfillcolor{currentfill}%
\pgfsetlinewidth{0.000000pt}%
\definecolor{currentstroke}{rgb}{0.000000,0.000000,0.000000}%
\pgfsetstrokecolor{currentstroke}%
\pgfsetstrokeopacity{0.000000}%
\pgfsetdash{}{0pt}%
\pgfpathmoveto{\pgfqpoint{3.447900in}{0.660000in}}%
\pgfpathlineto{\pgfqpoint{3.454100in}{0.660000in}}%
\pgfpathlineto{\pgfqpoint{3.454100in}{3.340071in}}%
\pgfpathlineto{\pgfqpoint{3.447900in}{3.340071in}}%
\pgfpathlineto{\pgfqpoint{3.447900in}{0.660000in}}%
\pgfpathclose%
\pgfusepath{fill}%
\end{pgfscope}%
\begin{pgfscope}%
\pgfpathrectangle{\pgfqpoint{1.250000in}{0.660000in}}{\pgfqpoint{7.750000in}{4.620000in}}%
\pgfusepath{clip}%
\pgfsetbuttcap%
\pgfsetmiterjoin%
\definecolor{currentfill}{rgb}{0.000000,0.000000,0.000000}%
\pgfsetfillcolor{currentfill}%
\pgfsetlinewidth{0.000000pt}%
\definecolor{currentstroke}{rgb}{0.000000,0.000000,0.000000}%
\pgfsetstrokecolor{currentstroke}%
\pgfsetstrokeopacity{0.000000}%
\pgfsetdash{}{0pt}%
\pgfpathmoveto{\pgfqpoint{3.455650in}{0.660000in}}%
\pgfpathlineto{\pgfqpoint{3.461850in}{0.660000in}}%
\pgfpathlineto{\pgfqpoint{3.461850in}{3.353462in}}%
\pgfpathlineto{\pgfqpoint{3.455650in}{3.353462in}}%
\pgfpathlineto{\pgfqpoint{3.455650in}{0.660000in}}%
\pgfpathclose%
\pgfusepath{fill}%
\end{pgfscope}%
\begin{pgfscope}%
\pgfpathrectangle{\pgfqpoint{1.250000in}{0.660000in}}{\pgfqpoint{7.750000in}{4.620000in}}%
\pgfusepath{clip}%
\pgfsetbuttcap%
\pgfsetmiterjoin%
\definecolor{currentfill}{rgb}{0.000000,0.000000,0.000000}%
\pgfsetfillcolor{currentfill}%
\pgfsetlinewidth{0.000000pt}%
\definecolor{currentstroke}{rgb}{0.000000,0.000000,0.000000}%
\pgfsetstrokecolor{currentstroke}%
\pgfsetstrokeopacity{0.000000}%
\pgfsetdash{}{0pt}%
\pgfpathmoveto{\pgfqpoint{3.463400in}{0.660000in}}%
\pgfpathlineto{\pgfqpoint{3.469600in}{0.660000in}}%
\pgfpathlineto{\pgfqpoint{3.469600in}{3.363039in}}%
\pgfpathlineto{\pgfqpoint{3.463400in}{3.363039in}}%
\pgfpathlineto{\pgfqpoint{3.463400in}{0.660000in}}%
\pgfpathclose%
\pgfusepath{fill}%
\end{pgfscope}%
\begin{pgfscope}%
\pgfpathrectangle{\pgfqpoint{1.250000in}{0.660000in}}{\pgfqpoint{7.750000in}{4.620000in}}%
\pgfusepath{clip}%
\pgfsetbuttcap%
\pgfsetmiterjoin%
\definecolor{currentfill}{rgb}{0.000000,0.000000,0.000000}%
\pgfsetfillcolor{currentfill}%
\pgfsetlinewidth{0.000000pt}%
\definecolor{currentstroke}{rgb}{0.000000,0.000000,0.000000}%
\pgfsetstrokecolor{currentstroke}%
\pgfsetstrokeopacity{0.000000}%
\pgfsetdash{}{0pt}%
\pgfpathmoveto{\pgfqpoint{3.471150in}{0.660000in}}%
\pgfpathlineto{\pgfqpoint{3.477350in}{0.660000in}}%
\pgfpathlineto{\pgfqpoint{3.477350in}{3.347063in}}%
\pgfpathlineto{\pgfqpoint{3.471150in}{3.347063in}}%
\pgfpathlineto{\pgfqpoint{3.471150in}{0.660000in}}%
\pgfpathclose%
\pgfusepath{fill}%
\end{pgfscope}%
\begin{pgfscope}%
\pgfpathrectangle{\pgfqpoint{1.250000in}{0.660000in}}{\pgfqpoint{7.750000in}{4.620000in}}%
\pgfusepath{clip}%
\pgfsetbuttcap%
\pgfsetmiterjoin%
\definecolor{currentfill}{rgb}{0.000000,0.000000,0.000000}%
\pgfsetfillcolor{currentfill}%
\pgfsetlinewidth{0.000000pt}%
\definecolor{currentstroke}{rgb}{0.000000,0.000000,0.000000}%
\pgfsetstrokecolor{currentstroke}%
\pgfsetstrokeopacity{0.000000}%
\pgfsetdash{}{0pt}%
\pgfpathmoveto{\pgfqpoint{3.478900in}{0.660000in}}%
\pgfpathlineto{\pgfqpoint{3.485100in}{0.660000in}}%
\pgfpathlineto{\pgfqpoint{3.485100in}{3.331203in}}%
\pgfpathlineto{\pgfqpoint{3.478900in}{3.331203in}}%
\pgfpathlineto{\pgfqpoint{3.478900in}{0.660000in}}%
\pgfpathclose%
\pgfusepath{fill}%
\end{pgfscope}%
\begin{pgfscope}%
\pgfpathrectangle{\pgfqpoint{1.250000in}{0.660000in}}{\pgfqpoint{7.750000in}{4.620000in}}%
\pgfusepath{clip}%
\pgfsetbuttcap%
\pgfsetmiterjoin%
\definecolor{currentfill}{rgb}{0.000000,0.000000,0.000000}%
\pgfsetfillcolor{currentfill}%
\pgfsetlinewidth{0.000000pt}%
\definecolor{currentstroke}{rgb}{0.000000,0.000000,0.000000}%
\pgfsetstrokecolor{currentstroke}%
\pgfsetstrokeopacity{0.000000}%
\pgfsetdash{}{0pt}%
\pgfpathmoveto{\pgfqpoint{3.486650in}{0.660000in}}%
\pgfpathlineto{\pgfqpoint{3.492850in}{0.660000in}}%
\pgfpathlineto{\pgfqpoint{3.492850in}{3.367761in}}%
\pgfpathlineto{\pgfqpoint{3.486650in}{3.367761in}}%
\pgfpathlineto{\pgfqpoint{3.486650in}{0.660000in}}%
\pgfpathclose%
\pgfusepath{fill}%
\end{pgfscope}%
\begin{pgfscope}%
\pgfpathrectangle{\pgfqpoint{1.250000in}{0.660000in}}{\pgfqpoint{7.750000in}{4.620000in}}%
\pgfusepath{clip}%
\pgfsetbuttcap%
\pgfsetmiterjoin%
\definecolor{currentfill}{rgb}{0.000000,0.000000,0.000000}%
\pgfsetfillcolor{currentfill}%
\pgfsetlinewidth{0.000000pt}%
\definecolor{currentstroke}{rgb}{0.000000,0.000000,0.000000}%
\pgfsetstrokecolor{currentstroke}%
\pgfsetstrokeopacity{0.000000}%
\pgfsetdash{}{0pt}%
\pgfpathmoveto{\pgfqpoint{3.494400in}{0.660000in}}%
\pgfpathlineto{\pgfqpoint{3.500600in}{0.660000in}}%
\pgfpathlineto{\pgfqpoint{3.500600in}{3.386236in}}%
\pgfpathlineto{\pgfqpoint{3.494400in}{3.386236in}}%
\pgfpathlineto{\pgfqpoint{3.494400in}{0.660000in}}%
\pgfpathclose%
\pgfusepath{fill}%
\end{pgfscope}%
\begin{pgfscope}%
\pgfpathrectangle{\pgfqpoint{1.250000in}{0.660000in}}{\pgfqpoint{7.750000in}{4.620000in}}%
\pgfusepath{clip}%
\pgfsetbuttcap%
\pgfsetmiterjoin%
\definecolor{currentfill}{rgb}{0.000000,0.000000,0.000000}%
\pgfsetfillcolor{currentfill}%
\pgfsetlinewidth{0.000000pt}%
\definecolor{currentstroke}{rgb}{0.000000,0.000000,0.000000}%
\pgfsetstrokecolor{currentstroke}%
\pgfsetstrokeopacity{0.000000}%
\pgfsetdash{}{0pt}%
\pgfpathmoveto{\pgfqpoint{3.502150in}{0.660000in}}%
\pgfpathlineto{\pgfqpoint{3.508350in}{0.660000in}}%
\pgfpathlineto{\pgfqpoint{3.508350in}{3.334538in}}%
\pgfpathlineto{\pgfqpoint{3.502150in}{3.334538in}}%
\pgfpathlineto{\pgfqpoint{3.502150in}{0.660000in}}%
\pgfpathclose%
\pgfusepath{fill}%
\end{pgfscope}%
\begin{pgfscope}%
\pgfpathrectangle{\pgfqpoint{1.250000in}{0.660000in}}{\pgfqpoint{7.750000in}{4.620000in}}%
\pgfusepath{clip}%
\pgfsetbuttcap%
\pgfsetmiterjoin%
\definecolor{currentfill}{rgb}{0.000000,0.000000,0.000000}%
\pgfsetfillcolor{currentfill}%
\pgfsetlinewidth{0.000000pt}%
\definecolor{currentstroke}{rgb}{0.000000,0.000000,0.000000}%
\pgfsetstrokecolor{currentstroke}%
\pgfsetstrokeopacity{0.000000}%
\pgfsetdash{}{0pt}%
\pgfpathmoveto{\pgfqpoint{3.509900in}{0.660000in}}%
\pgfpathlineto{\pgfqpoint{3.516100in}{0.660000in}}%
\pgfpathlineto{\pgfqpoint{3.516100in}{3.333180in}}%
\pgfpathlineto{\pgfqpoint{3.509900in}{3.333180in}}%
\pgfpathlineto{\pgfqpoint{3.509900in}{0.660000in}}%
\pgfpathclose%
\pgfusepath{fill}%
\end{pgfscope}%
\begin{pgfscope}%
\pgfpathrectangle{\pgfqpoint{1.250000in}{0.660000in}}{\pgfqpoint{7.750000in}{4.620000in}}%
\pgfusepath{clip}%
\pgfsetbuttcap%
\pgfsetmiterjoin%
\definecolor{currentfill}{rgb}{0.000000,0.000000,0.000000}%
\pgfsetfillcolor{currentfill}%
\pgfsetlinewidth{0.000000pt}%
\definecolor{currentstroke}{rgb}{0.000000,0.000000,0.000000}%
\pgfsetstrokecolor{currentstroke}%
\pgfsetstrokeopacity{0.000000}%
\pgfsetdash{}{0pt}%
\pgfpathmoveto{\pgfqpoint{3.517650in}{0.660000in}}%
\pgfpathlineto{\pgfqpoint{3.523850in}{0.660000in}}%
\pgfpathlineto{\pgfqpoint{3.523850in}{3.347439in}}%
\pgfpathlineto{\pgfqpoint{3.517650in}{3.347439in}}%
\pgfpathlineto{\pgfqpoint{3.517650in}{0.660000in}}%
\pgfpathclose%
\pgfusepath{fill}%
\end{pgfscope}%
\begin{pgfscope}%
\pgfpathrectangle{\pgfqpoint{1.250000in}{0.660000in}}{\pgfqpoint{7.750000in}{4.620000in}}%
\pgfusepath{clip}%
\pgfsetbuttcap%
\pgfsetmiterjoin%
\definecolor{currentfill}{rgb}{0.000000,0.000000,0.000000}%
\pgfsetfillcolor{currentfill}%
\pgfsetlinewidth{0.000000pt}%
\definecolor{currentstroke}{rgb}{0.000000,0.000000,0.000000}%
\pgfsetstrokecolor{currentstroke}%
\pgfsetstrokeopacity{0.000000}%
\pgfsetdash{}{0pt}%
\pgfpathmoveto{\pgfqpoint{3.525400in}{0.660000in}}%
\pgfpathlineto{\pgfqpoint{3.531600in}{0.660000in}}%
\pgfpathlineto{\pgfqpoint{3.531600in}{3.351963in}}%
\pgfpathlineto{\pgfqpoint{3.525400in}{3.351963in}}%
\pgfpathlineto{\pgfqpoint{3.525400in}{0.660000in}}%
\pgfpathclose%
\pgfusepath{fill}%
\end{pgfscope}%
\begin{pgfscope}%
\pgfpathrectangle{\pgfqpoint{1.250000in}{0.660000in}}{\pgfqpoint{7.750000in}{4.620000in}}%
\pgfusepath{clip}%
\pgfsetbuttcap%
\pgfsetmiterjoin%
\definecolor{currentfill}{rgb}{0.000000,0.000000,0.000000}%
\pgfsetfillcolor{currentfill}%
\pgfsetlinewidth{0.000000pt}%
\definecolor{currentstroke}{rgb}{0.000000,0.000000,0.000000}%
\pgfsetstrokecolor{currentstroke}%
\pgfsetstrokeopacity{0.000000}%
\pgfsetdash{}{0pt}%
\pgfpathmoveto{\pgfqpoint{3.533150in}{0.660000in}}%
\pgfpathlineto{\pgfqpoint{3.539350in}{0.660000in}}%
\pgfpathlineto{\pgfqpoint{3.539350in}{3.341343in}}%
\pgfpathlineto{\pgfqpoint{3.533150in}{3.341343in}}%
\pgfpathlineto{\pgfqpoint{3.533150in}{0.660000in}}%
\pgfpathclose%
\pgfusepath{fill}%
\end{pgfscope}%
\begin{pgfscope}%
\pgfpathrectangle{\pgfqpoint{1.250000in}{0.660000in}}{\pgfqpoint{7.750000in}{4.620000in}}%
\pgfusepath{clip}%
\pgfsetbuttcap%
\pgfsetmiterjoin%
\definecolor{currentfill}{rgb}{0.000000,0.000000,0.000000}%
\pgfsetfillcolor{currentfill}%
\pgfsetlinewidth{0.000000pt}%
\definecolor{currentstroke}{rgb}{0.000000,0.000000,0.000000}%
\pgfsetstrokecolor{currentstroke}%
\pgfsetstrokeopacity{0.000000}%
\pgfsetdash{}{0pt}%
\pgfpathmoveto{\pgfqpoint{3.540900in}{0.660000in}}%
\pgfpathlineto{\pgfqpoint{3.547100in}{0.660000in}}%
\pgfpathlineto{\pgfqpoint{3.547100in}{3.400270in}}%
\pgfpathlineto{\pgfqpoint{3.540900in}{3.400270in}}%
\pgfpathlineto{\pgfqpoint{3.540900in}{0.660000in}}%
\pgfpathclose%
\pgfusepath{fill}%
\end{pgfscope}%
\begin{pgfscope}%
\pgfpathrectangle{\pgfqpoint{1.250000in}{0.660000in}}{\pgfqpoint{7.750000in}{4.620000in}}%
\pgfusepath{clip}%
\pgfsetbuttcap%
\pgfsetmiterjoin%
\definecolor{currentfill}{rgb}{0.000000,0.000000,0.000000}%
\pgfsetfillcolor{currentfill}%
\pgfsetlinewidth{0.000000pt}%
\definecolor{currentstroke}{rgb}{0.000000,0.000000,0.000000}%
\pgfsetstrokecolor{currentstroke}%
\pgfsetstrokeopacity{0.000000}%
\pgfsetdash{}{0pt}%
\pgfpathmoveto{\pgfqpoint{3.548650in}{0.660000in}}%
\pgfpathlineto{\pgfqpoint{3.554850in}{0.660000in}}%
\pgfpathlineto{\pgfqpoint{3.554850in}{3.334848in}}%
\pgfpathlineto{\pgfqpoint{3.548650in}{3.334848in}}%
\pgfpathlineto{\pgfqpoint{3.548650in}{0.660000in}}%
\pgfpathclose%
\pgfusepath{fill}%
\end{pgfscope}%
\begin{pgfscope}%
\pgfpathrectangle{\pgfqpoint{1.250000in}{0.660000in}}{\pgfqpoint{7.750000in}{4.620000in}}%
\pgfusepath{clip}%
\pgfsetbuttcap%
\pgfsetmiterjoin%
\definecolor{currentfill}{rgb}{0.000000,0.000000,0.000000}%
\pgfsetfillcolor{currentfill}%
\pgfsetlinewidth{0.000000pt}%
\definecolor{currentstroke}{rgb}{0.000000,0.000000,0.000000}%
\pgfsetstrokecolor{currentstroke}%
\pgfsetstrokeopacity{0.000000}%
\pgfsetdash{}{0pt}%
\pgfpathmoveto{\pgfqpoint{3.556400in}{0.660000in}}%
\pgfpathlineto{\pgfqpoint{3.562600in}{0.660000in}}%
\pgfpathlineto{\pgfqpoint{3.562600in}{3.376825in}}%
\pgfpathlineto{\pgfqpoint{3.556400in}{3.376825in}}%
\pgfpathlineto{\pgfqpoint{3.556400in}{0.660000in}}%
\pgfpathclose%
\pgfusepath{fill}%
\end{pgfscope}%
\begin{pgfscope}%
\pgfpathrectangle{\pgfqpoint{1.250000in}{0.660000in}}{\pgfqpoint{7.750000in}{4.620000in}}%
\pgfusepath{clip}%
\pgfsetbuttcap%
\pgfsetmiterjoin%
\definecolor{currentfill}{rgb}{0.000000,0.000000,0.000000}%
\pgfsetfillcolor{currentfill}%
\pgfsetlinewidth{0.000000pt}%
\definecolor{currentstroke}{rgb}{0.000000,0.000000,0.000000}%
\pgfsetstrokecolor{currentstroke}%
\pgfsetstrokeopacity{0.000000}%
\pgfsetdash{}{0pt}%
\pgfpathmoveto{\pgfqpoint{3.564150in}{0.660000in}}%
\pgfpathlineto{\pgfqpoint{3.570350in}{0.660000in}}%
\pgfpathlineto{\pgfqpoint{3.570350in}{3.355305in}}%
\pgfpathlineto{\pgfqpoint{3.564150in}{3.355305in}}%
\pgfpathlineto{\pgfqpoint{3.564150in}{0.660000in}}%
\pgfpathclose%
\pgfusepath{fill}%
\end{pgfscope}%
\begin{pgfscope}%
\pgfpathrectangle{\pgfqpoint{1.250000in}{0.660000in}}{\pgfqpoint{7.750000in}{4.620000in}}%
\pgfusepath{clip}%
\pgfsetbuttcap%
\pgfsetmiterjoin%
\definecolor{currentfill}{rgb}{0.000000,0.000000,0.000000}%
\pgfsetfillcolor{currentfill}%
\pgfsetlinewidth{0.000000pt}%
\definecolor{currentstroke}{rgb}{0.000000,0.000000,0.000000}%
\pgfsetstrokecolor{currentstroke}%
\pgfsetstrokeopacity{0.000000}%
\pgfsetdash{}{0pt}%
\pgfpathmoveto{\pgfqpoint{3.571900in}{0.660000in}}%
\pgfpathlineto{\pgfqpoint{3.578100in}{0.660000in}}%
\pgfpathlineto{\pgfqpoint{3.578100in}{3.309655in}}%
\pgfpathlineto{\pgfqpoint{3.571900in}{3.309655in}}%
\pgfpathlineto{\pgfqpoint{3.571900in}{0.660000in}}%
\pgfpathclose%
\pgfusepath{fill}%
\end{pgfscope}%
\begin{pgfscope}%
\pgfpathrectangle{\pgfqpoint{1.250000in}{0.660000in}}{\pgfqpoint{7.750000in}{4.620000in}}%
\pgfusepath{clip}%
\pgfsetbuttcap%
\pgfsetmiterjoin%
\definecolor{currentfill}{rgb}{0.000000,0.000000,0.000000}%
\pgfsetfillcolor{currentfill}%
\pgfsetlinewidth{0.000000pt}%
\definecolor{currentstroke}{rgb}{0.000000,0.000000,0.000000}%
\pgfsetstrokecolor{currentstroke}%
\pgfsetstrokeopacity{0.000000}%
\pgfsetdash{}{0pt}%
\pgfpathmoveto{\pgfqpoint{3.579650in}{0.660000in}}%
\pgfpathlineto{\pgfqpoint{3.585850in}{0.660000in}}%
\pgfpathlineto{\pgfqpoint{3.585850in}{3.339598in}}%
\pgfpathlineto{\pgfqpoint{3.579650in}{3.339598in}}%
\pgfpathlineto{\pgfqpoint{3.579650in}{0.660000in}}%
\pgfpathclose%
\pgfusepath{fill}%
\end{pgfscope}%
\begin{pgfscope}%
\pgfpathrectangle{\pgfqpoint{1.250000in}{0.660000in}}{\pgfqpoint{7.750000in}{4.620000in}}%
\pgfusepath{clip}%
\pgfsetbuttcap%
\pgfsetmiterjoin%
\definecolor{currentfill}{rgb}{0.000000,0.000000,0.000000}%
\pgfsetfillcolor{currentfill}%
\pgfsetlinewidth{0.000000pt}%
\definecolor{currentstroke}{rgb}{0.000000,0.000000,0.000000}%
\pgfsetstrokecolor{currentstroke}%
\pgfsetstrokeopacity{0.000000}%
\pgfsetdash{}{0pt}%
\pgfpathmoveto{\pgfqpoint{3.587400in}{0.660000in}}%
\pgfpathlineto{\pgfqpoint{3.593600in}{0.660000in}}%
\pgfpathlineto{\pgfqpoint{3.593600in}{3.353647in}}%
\pgfpathlineto{\pgfqpoint{3.587400in}{3.353647in}}%
\pgfpathlineto{\pgfqpoint{3.587400in}{0.660000in}}%
\pgfpathclose%
\pgfusepath{fill}%
\end{pgfscope}%
\begin{pgfscope}%
\pgfpathrectangle{\pgfqpoint{1.250000in}{0.660000in}}{\pgfqpoint{7.750000in}{4.620000in}}%
\pgfusepath{clip}%
\pgfsetbuttcap%
\pgfsetmiterjoin%
\definecolor{currentfill}{rgb}{0.000000,0.000000,0.000000}%
\pgfsetfillcolor{currentfill}%
\pgfsetlinewidth{0.000000pt}%
\definecolor{currentstroke}{rgb}{0.000000,0.000000,0.000000}%
\pgfsetstrokecolor{currentstroke}%
\pgfsetstrokeopacity{0.000000}%
\pgfsetdash{}{0pt}%
\pgfpathmoveto{\pgfqpoint{3.595150in}{0.660000in}}%
\pgfpathlineto{\pgfqpoint{3.601350in}{0.660000in}}%
\pgfpathlineto{\pgfqpoint{3.601350in}{3.309083in}}%
\pgfpathlineto{\pgfqpoint{3.595150in}{3.309083in}}%
\pgfpathlineto{\pgfqpoint{3.595150in}{0.660000in}}%
\pgfpathclose%
\pgfusepath{fill}%
\end{pgfscope}%
\begin{pgfscope}%
\pgfpathrectangle{\pgfqpoint{1.250000in}{0.660000in}}{\pgfqpoint{7.750000in}{4.620000in}}%
\pgfusepath{clip}%
\pgfsetbuttcap%
\pgfsetmiterjoin%
\definecolor{currentfill}{rgb}{0.000000,0.000000,0.000000}%
\pgfsetfillcolor{currentfill}%
\pgfsetlinewidth{0.000000pt}%
\definecolor{currentstroke}{rgb}{0.000000,0.000000,0.000000}%
\pgfsetstrokecolor{currentstroke}%
\pgfsetstrokeopacity{0.000000}%
\pgfsetdash{}{0pt}%
\pgfpathmoveto{\pgfqpoint{3.602900in}{0.660000in}}%
\pgfpathlineto{\pgfqpoint{3.609100in}{0.660000in}}%
\pgfpathlineto{\pgfqpoint{3.609100in}{3.358252in}}%
\pgfpathlineto{\pgfqpoint{3.602900in}{3.358252in}}%
\pgfpathlineto{\pgfqpoint{3.602900in}{0.660000in}}%
\pgfpathclose%
\pgfusepath{fill}%
\end{pgfscope}%
\begin{pgfscope}%
\pgfpathrectangle{\pgfqpoint{1.250000in}{0.660000in}}{\pgfqpoint{7.750000in}{4.620000in}}%
\pgfusepath{clip}%
\pgfsetbuttcap%
\pgfsetmiterjoin%
\definecolor{currentfill}{rgb}{0.000000,0.000000,0.000000}%
\pgfsetfillcolor{currentfill}%
\pgfsetlinewidth{0.000000pt}%
\definecolor{currentstroke}{rgb}{0.000000,0.000000,0.000000}%
\pgfsetstrokecolor{currentstroke}%
\pgfsetstrokeopacity{0.000000}%
\pgfsetdash{}{0pt}%
\pgfpathmoveto{\pgfqpoint{3.610650in}{0.660000in}}%
\pgfpathlineto{\pgfqpoint{3.616850in}{0.660000in}}%
\pgfpathlineto{\pgfqpoint{3.616850in}{3.377283in}}%
\pgfpathlineto{\pgfqpoint{3.610650in}{3.377283in}}%
\pgfpathlineto{\pgfqpoint{3.610650in}{0.660000in}}%
\pgfpathclose%
\pgfusepath{fill}%
\end{pgfscope}%
\begin{pgfscope}%
\pgfpathrectangle{\pgfqpoint{1.250000in}{0.660000in}}{\pgfqpoint{7.750000in}{4.620000in}}%
\pgfusepath{clip}%
\pgfsetbuttcap%
\pgfsetmiterjoin%
\definecolor{currentfill}{rgb}{0.000000,0.000000,0.000000}%
\pgfsetfillcolor{currentfill}%
\pgfsetlinewidth{0.000000pt}%
\definecolor{currentstroke}{rgb}{0.000000,0.000000,0.000000}%
\pgfsetstrokecolor{currentstroke}%
\pgfsetstrokeopacity{0.000000}%
\pgfsetdash{}{0pt}%
\pgfpathmoveto{\pgfqpoint{3.618400in}{0.660000in}}%
\pgfpathlineto{\pgfqpoint{3.624600in}{0.660000in}}%
\pgfpathlineto{\pgfqpoint{3.624600in}{3.383292in}}%
\pgfpathlineto{\pgfqpoint{3.618400in}{3.383292in}}%
\pgfpathlineto{\pgfqpoint{3.618400in}{0.660000in}}%
\pgfpathclose%
\pgfusepath{fill}%
\end{pgfscope}%
\begin{pgfscope}%
\pgfpathrectangle{\pgfqpoint{1.250000in}{0.660000in}}{\pgfqpoint{7.750000in}{4.620000in}}%
\pgfusepath{clip}%
\pgfsetbuttcap%
\pgfsetmiterjoin%
\definecolor{currentfill}{rgb}{0.000000,0.000000,0.000000}%
\pgfsetfillcolor{currentfill}%
\pgfsetlinewidth{0.000000pt}%
\definecolor{currentstroke}{rgb}{0.000000,0.000000,0.000000}%
\pgfsetstrokecolor{currentstroke}%
\pgfsetstrokeopacity{0.000000}%
\pgfsetdash{}{0pt}%
\pgfpathmoveto{\pgfqpoint{3.626150in}{0.660000in}}%
\pgfpathlineto{\pgfqpoint{3.632350in}{0.660000in}}%
\pgfpathlineto{\pgfqpoint{3.632350in}{3.363170in}}%
\pgfpathlineto{\pgfqpoint{3.626150in}{3.363170in}}%
\pgfpathlineto{\pgfqpoint{3.626150in}{0.660000in}}%
\pgfpathclose%
\pgfusepath{fill}%
\end{pgfscope}%
\begin{pgfscope}%
\pgfpathrectangle{\pgfqpoint{1.250000in}{0.660000in}}{\pgfqpoint{7.750000in}{4.620000in}}%
\pgfusepath{clip}%
\pgfsetbuttcap%
\pgfsetmiterjoin%
\definecolor{currentfill}{rgb}{0.000000,0.000000,0.000000}%
\pgfsetfillcolor{currentfill}%
\pgfsetlinewidth{0.000000pt}%
\definecolor{currentstroke}{rgb}{0.000000,0.000000,0.000000}%
\pgfsetstrokecolor{currentstroke}%
\pgfsetstrokeopacity{0.000000}%
\pgfsetdash{}{0pt}%
\pgfpathmoveto{\pgfqpoint{3.633900in}{0.660000in}}%
\pgfpathlineto{\pgfqpoint{3.640100in}{0.660000in}}%
\pgfpathlineto{\pgfqpoint{3.640100in}{3.345743in}}%
\pgfpathlineto{\pgfqpoint{3.633900in}{3.345743in}}%
\pgfpathlineto{\pgfqpoint{3.633900in}{0.660000in}}%
\pgfpathclose%
\pgfusepath{fill}%
\end{pgfscope}%
\begin{pgfscope}%
\pgfpathrectangle{\pgfqpoint{1.250000in}{0.660000in}}{\pgfqpoint{7.750000in}{4.620000in}}%
\pgfusepath{clip}%
\pgfsetbuttcap%
\pgfsetmiterjoin%
\definecolor{currentfill}{rgb}{0.000000,0.000000,0.000000}%
\pgfsetfillcolor{currentfill}%
\pgfsetlinewidth{0.000000pt}%
\definecolor{currentstroke}{rgb}{0.000000,0.000000,0.000000}%
\pgfsetstrokecolor{currentstroke}%
\pgfsetstrokeopacity{0.000000}%
\pgfsetdash{}{0pt}%
\pgfpathmoveto{\pgfqpoint{3.641650in}{0.660000in}}%
\pgfpathlineto{\pgfqpoint{3.647850in}{0.660000in}}%
\pgfpathlineto{\pgfqpoint{3.647850in}{3.386795in}}%
\pgfpathlineto{\pgfqpoint{3.641650in}{3.386795in}}%
\pgfpathlineto{\pgfqpoint{3.641650in}{0.660000in}}%
\pgfpathclose%
\pgfusepath{fill}%
\end{pgfscope}%
\begin{pgfscope}%
\pgfpathrectangle{\pgfqpoint{1.250000in}{0.660000in}}{\pgfqpoint{7.750000in}{4.620000in}}%
\pgfusepath{clip}%
\pgfsetbuttcap%
\pgfsetmiterjoin%
\definecolor{currentfill}{rgb}{0.000000,0.000000,0.000000}%
\pgfsetfillcolor{currentfill}%
\pgfsetlinewidth{0.000000pt}%
\definecolor{currentstroke}{rgb}{0.000000,0.000000,0.000000}%
\pgfsetstrokecolor{currentstroke}%
\pgfsetstrokeopacity{0.000000}%
\pgfsetdash{}{0pt}%
\pgfpathmoveto{\pgfqpoint{3.649400in}{0.660000in}}%
\pgfpathlineto{\pgfqpoint{3.655600in}{0.660000in}}%
\pgfpathlineto{\pgfqpoint{3.655600in}{3.364335in}}%
\pgfpathlineto{\pgfqpoint{3.649400in}{3.364335in}}%
\pgfpathlineto{\pgfqpoint{3.649400in}{0.660000in}}%
\pgfpathclose%
\pgfusepath{fill}%
\end{pgfscope}%
\begin{pgfscope}%
\pgfpathrectangle{\pgfqpoint{1.250000in}{0.660000in}}{\pgfqpoint{7.750000in}{4.620000in}}%
\pgfusepath{clip}%
\pgfsetbuttcap%
\pgfsetmiterjoin%
\definecolor{currentfill}{rgb}{0.000000,0.000000,0.000000}%
\pgfsetfillcolor{currentfill}%
\pgfsetlinewidth{0.000000pt}%
\definecolor{currentstroke}{rgb}{0.000000,0.000000,0.000000}%
\pgfsetstrokecolor{currentstroke}%
\pgfsetstrokeopacity{0.000000}%
\pgfsetdash{}{0pt}%
\pgfpathmoveto{\pgfqpoint{3.657150in}{0.660000in}}%
\pgfpathlineto{\pgfqpoint{3.663350in}{0.660000in}}%
\pgfpathlineto{\pgfqpoint{3.663350in}{3.397013in}}%
\pgfpathlineto{\pgfqpoint{3.657150in}{3.397013in}}%
\pgfpathlineto{\pgfqpoint{3.657150in}{0.660000in}}%
\pgfpathclose%
\pgfusepath{fill}%
\end{pgfscope}%
\begin{pgfscope}%
\pgfpathrectangle{\pgfqpoint{1.250000in}{0.660000in}}{\pgfqpoint{7.750000in}{4.620000in}}%
\pgfusepath{clip}%
\pgfsetbuttcap%
\pgfsetmiterjoin%
\definecolor{currentfill}{rgb}{0.000000,0.000000,0.000000}%
\pgfsetfillcolor{currentfill}%
\pgfsetlinewidth{0.000000pt}%
\definecolor{currentstroke}{rgb}{0.000000,0.000000,0.000000}%
\pgfsetstrokecolor{currentstroke}%
\pgfsetstrokeopacity{0.000000}%
\pgfsetdash{}{0pt}%
\pgfpathmoveto{\pgfqpoint{3.664900in}{0.660000in}}%
\pgfpathlineto{\pgfqpoint{3.671100in}{0.660000in}}%
\pgfpathlineto{\pgfqpoint{3.671100in}{3.434385in}}%
\pgfpathlineto{\pgfqpoint{3.664900in}{3.434385in}}%
\pgfpathlineto{\pgfqpoint{3.664900in}{0.660000in}}%
\pgfpathclose%
\pgfusepath{fill}%
\end{pgfscope}%
\begin{pgfscope}%
\pgfpathrectangle{\pgfqpoint{1.250000in}{0.660000in}}{\pgfqpoint{7.750000in}{4.620000in}}%
\pgfusepath{clip}%
\pgfsetbuttcap%
\pgfsetmiterjoin%
\definecolor{currentfill}{rgb}{0.000000,0.000000,0.000000}%
\pgfsetfillcolor{currentfill}%
\pgfsetlinewidth{0.000000pt}%
\definecolor{currentstroke}{rgb}{0.000000,0.000000,0.000000}%
\pgfsetstrokecolor{currentstroke}%
\pgfsetstrokeopacity{0.000000}%
\pgfsetdash{}{0pt}%
\pgfpathmoveto{\pgfqpoint{3.672650in}{0.660000in}}%
\pgfpathlineto{\pgfqpoint{3.678850in}{0.660000in}}%
\pgfpathlineto{\pgfqpoint{3.678850in}{3.342215in}}%
\pgfpathlineto{\pgfqpoint{3.672650in}{3.342215in}}%
\pgfpathlineto{\pgfqpoint{3.672650in}{0.660000in}}%
\pgfpathclose%
\pgfusepath{fill}%
\end{pgfscope}%
\begin{pgfscope}%
\pgfpathrectangle{\pgfqpoint{1.250000in}{0.660000in}}{\pgfqpoint{7.750000in}{4.620000in}}%
\pgfusepath{clip}%
\pgfsetbuttcap%
\pgfsetmiterjoin%
\definecolor{currentfill}{rgb}{0.000000,0.000000,0.000000}%
\pgfsetfillcolor{currentfill}%
\pgfsetlinewidth{0.000000pt}%
\definecolor{currentstroke}{rgb}{0.000000,0.000000,0.000000}%
\pgfsetstrokecolor{currentstroke}%
\pgfsetstrokeopacity{0.000000}%
\pgfsetdash{}{0pt}%
\pgfpathmoveto{\pgfqpoint{3.680400in}{0.660000in}}%
\pgfpathlineto{\pgfqpoint{3.686600in}{0.660000in}}%
\pgfpathlineto{\pgfqpoint{3.686600in}{3.312040in}}%
\pgfpathlineto{\pgfqpoint{3.680400in}{3.312040in}}%
\pgfpathlineto{\pgfqpoint{3.680400in}{0.660000in}}%
\pgfpathclose%
\pgfusepath{fill}%
\end{pgfscope}%
\begin{pgfscope}%
\pgfpathrectangle{\pgfqpoint{1.250000in}{0.660000in}}{\pgfqpoint{7.750000in}{4.620000in}}%
\pgfusepath{clip}%
\pgfsetbuttcap%
\pgfsetmiterjoin%
\definecolor{currentfill}{rgb}{0.000000,0.000000,0.000000}%
\pgfsetfillcolor{currentfill}%
\pgfsetlinewidth{0.000000pt}%
\definecolor{currentstroke}{rgb}{0.000000,0.000000,0.000000}%
\pgfsetstrokecolor{currentstroke}%
\pgfsetstrokeopacity{0.000000}%
\pgfsetdash{}{0pt}%
\pgfpathmoveto{\pgfqpoint{3.688150in}{0.660000in}}%
\pgfpathlineto{\pgfqpoint{3.694350in}{0.660000in}}%
\pgfpathlineto{\pgfqpoint{3.694350in}{3.328735in}}%
\pgfpathlineto{\pgfqpoint{3.688150in}{3.328735in}}%
\pgfpathlineto{\pgfqpoint{3.688150in}{0.660000in}}%
\pgfpathclose%
\pgfusepath{fill}%
\end{pgfscope}%
\begin{pgfscope}%
\pgfpathrectangle{\pgfqpoint{1.250000in}{0.660000in}}{\pgfqpoint{7.750000in}{4.620000in}}%
\pgfusepath{clip}%
\pgfsetbuttcap%
\pgfsetmiterjoin%
\definecolor{currentfill}{rgb}{0.000000,0.000000,0.000000}%
\pgfsetfillcolor{currentfill}%
\pgfsetlinewidth{0.000000pt}%
\definecolor{currentstroke}{rgb}{0.000000,0.000000,0.000000}%
\pgfsetstrokecolor{currentstroke}%
\pgfsetstrokeopacity{0.000000}%
\pgfsetdash{}{0pt}%
\pgfpathmoveto{\pgfqpoint{3.695900in}{0.660000in}}%
\pgfpathlineto{\pgfqpoint{3.702100in}{0.660000in}}%
\pgfpathlineto{\pgfqpoint{3.702100in}{3.334457in}}%
\pgfpathlineto{\pgfqpoint{3.695900in}{3.334457in}}%
\pgfpathlineto{\pgfqpoint{3.695900in}{0.660000in}}%
\pgfpathclose%
\pgfusepath{fill}%
\end{pgfscope}%
\begin{pgfscope}%
\pgfpathrectangle{\pgfqpoint{1.250000in}{0.660000in}}{\pgfqpoint{7.750000in}{4.620000in}}%
\pgfusepath{clip}%
\pgfsetbuttcap%
\pgfsetmiterjoin%
\definecolor{currentfill}{rgb}{0.000000,0.000000,0.000000}%
\pgfsetfillcolor{currentfill}%
\pgfsetlinewidth{0.000000pt}%
\definecolor{currentstroke}{rgb}{0.000000,0.000000,0.000000}%
\pgfsetstrokecolor{currentstroke}%
\pgfsetstrokeopacity{0.000000}%
\pgfsetdash{}{0pt}%
\pgfpathmoveto{\pgfqpoint{3.703650in}{0.660000in}}%
\pgfpathlineto{\pgfqpoint{3.709850in}{0.660000in}}%
\pgfpathlineto{\pgfqpoint{3.709850in}{3.345063in}}%
\pgfpathlineto{\pgfqpoint{3.703650in}{3.345063in}}%
\pgfpathlineto{\pgfqpoint{3.703650in}{0.660000in}}%
\pgfpathclose%
\pgfusepath{fill}%
\end{pgfscope}%
\begin{pgfscope}%
\pgfpathrectangle{\pgfqpoint{1.250000in}{0.660000in}}{\pgfqpoint{7.750000in}{4.620000in}}%
\pgfusepath{clip}%
\pgfsetbuttcap%
\pgfsetmiterjoin%
\definecolor{currentfill}{rgb}{0.000000,0.000000,0.000000}%
\pgfsetfillcolor{currentfill}%
\pgfsetlinewidth{0.000000pt}%
\definecolor{currentstroke}{rgb}{0.000000,0.000000,0.000000}%
\pgfsetstrokecolor{currentstroke}%
\pgfsetstrokeopacity{0.000000}%
\pgfsetdash{}{0pt}%
\pgfpathmoveto{\pgfqpoint{3.711400in}{0.660000in}}%
\pgfpathlineto{\pgfqpoint{3.717600in}{0.660000in}}%
\pgfpathlineto{\pgfqpoint{3.717600in}{3.307795in}}%
\pgfpathlineto{\pgfqpoint{3.711400in}{3.307795in}}%
\pgfpathlineto{\pgfqpoint{3.711400in}{0.660000in}}%
\pgfpathclose%
\pgfusepath{fill}%
\end{pgfscope}%
\begin{pgfscope}%
\pgfpathrectangle{\pgfqpoint{1.250000in}{0.660000in}}{\pgfqpoint{7.750000in}{4.620000in}}%
\pgfusepath{clip}%
\pgfsetbuttcap%
\pgfsetmiterjoin%
\definecolor{currentfill}{rgb}{0.000000,0.000000,0.000000}%
\pgfsetfillcolor{currentfill}%
\pgfsetlinewidth{0.000000pt}%
\definecolor{currentstroke}{rgb}{0.000000,0.000000,0.000000}%
\pgfsetstrokecolor{currentstroke}%
\pgfsetstrokeopacity{0.000000}%
\pgfsetdash{}{0pt}%
\pgfpathmoveto{\pgfqpoint{3.719150in}{0.660000in}}%
\pgfpathlineto{\pgfqpoint{3.725350in}{0.660000in}}%
\pgfpathlineto{\pgfqpoint{3.725350in}{3.307352in}}%
\pgfpathlineto{\pgfqpoint{3.719150in}{3.307352in}}%
\pgfpathlineto{\pgfqpoint{3.719150in}{0.660000in}}%
\pgfpathclose%
\pgfusepath{fill}%
\end{pgfscope}%
\begin{pgfscope}%
\pgfpathrectangle{\pgfqpoint{1.250000in}{0.660000in}}{\pgfqpoint{7.750000in}{4.620000in}}%
\pgfusepath{clip}%
\pgfsetbuttcap%
\pgfsetmiterjoin%
\definecolor{currentfill}{rgb}{0.000000,0.000000,0.000000}%
\pgfsetfillcolor{currentfill}%
\pgfsetlinewidth{0.000000pt}%
\definecolor{currentstroke}{rgb}{0.000000,0.000000,0.000000}%
\pgfsetstrokecolor{currentstroke}%
\pgfsetstrokeopacity{0.000000}%
\pgfsetdash{}{0pt}%
\pgfpathmoveto{\pgfqpoint{3.726900in}{0.660000in}}%
\pgfpathlineto{\pgfqpoint{3.733100in}{0.660000in}}%
\pgfpathlineto{\pgfqpoint{3.733100in}{3.342923in}}%
\pgfpathlineto{\pgfqpoint{3.726900in}{3.342923in}}%
\pgfpathlineto{\pgfqpoint{3.726900in}{0.660000in}}%
\pgfpathclose%
\pgfusepath{fill}%
\end{pgfscope}%
\begin{pgfscope}%
\pgfpathrectangle{\pgfqpoint{1.250000in}{0.660000in}}{\pgfqpoint{7.750000in}{4.620000in}}%
\pgfusepath{clip}%
\pgfsetbuttcap%
\pgfsetmiterjoin%
\definecolor{currentfill}{rgb}{0.000000,0.000000,0.000000}%
\pgfsetfillcolor{currentfill}%
\pgfsetlinewidth{0.000000pt}%
\definecolor{currentstroke}{rgb}{0.000000,0.000000,0.000000}%
\pgfsetstrokecolor{currentstroke}%
\pgfsetstrokeopacity{0.000000}%
\pgfsetdash{}{0pt}%
\pgfpathmoveto{\pgfqpoint{3.734650in}{0.660000in}}%
\pgfpathlineto{\pgfqpoint{3.740850in}{0.660000in}}%
\pgfpathlineto{\pgfqpoint{3.740850in}{3.337175in}}%
\pgfpathlineto{\pgfqpoint{3.734650in}{3.337175in}}%
\pgfpathlineto{\pgfqpoint{3.734650in}{0.660000in}}%
\pgfpathclose%
\pgfusepath{fill}%
\end{pgfscope}%
\begin{pgfscope}%
\pgfpathrectangle{\pgfqpoint{1.250000in}{0.660000in}}{\pgfqpoint{7.750000in}{4.620000in}}%
\pgfusepath{clip}%
\pgfsetbuttcap%
\pgfsetmiterjoin%
\definecolor{currentfill}{rgb}{0.000000,0.000000,0.000000}%
\pgfsetfillcolor{currentfill}%
\pgfsetlinewidth{0.000000pt}%
\definecolor{currentstroke}{rgb}{0.000000,0.000000,0.000000}%
\pgfsetstrokecolor{currentstroke}%
\pgfsetstrokeopacity{0.000000}%
\pgfsetdash{}{0pt}%
\pgfpathmoveto{\pgfqpoint{3.742400in}{0.660000in}}%
\pgfpathlineto{\pgfqpoint{3.748600in}{0.660000in}}%
\pgfpathlineto{\pgfqpoint{3.748600in}{3.359053in}}%
\pgfpathlineto{\pgfqpoint{3.742400in}{3.359053in}}%
\pgfpathlineto{\pgfqpoint{3.742400in}{0.660000in}}%
\pgfpathclose%
\pgfusepath{fill}%
\end{pgfscope}%
\begin{pgfscope}%
\pgfpathrectangle{\pgfqpoint{1.250000in}{0.660000in}}{\pgfqpoint{7.750000in}{4.620000in}}%
\pgfusepath{clip}%
\pgfsetbuttcap%
\pgfsetmiterjoin%
\definecolor{currentfill}{rgb}{0.000000,0.000000,0.000000}%
\pgfsetfillcolor{currentfill}%
\pgfsetlinewidth{0.000000pt}%
\definecolor{currentstroke}{rgb}{0.000000,0.000000,0.000000}%
\pgfsetstrokecolor{currentstroke}%
\pgfsetstrokeopacity{0.000000}%
\pgfsetdash{}{0pt}%
\pgfpathmoveto{\pgfqpoint{3.750150in}{0.660000in}}%
\pgfpathlineto{\pgfqpoint{3.756350in}{0.660000in}}%
\pgfpathlineto{\pgfqpoint{3.756350in}{3.339515in}}%
\pgfpathlineto{\pgfqpoint{3.750150in}{3.339515in}}%
\pgfpathlineto{\pgfqpoint{3.750150in}{0.660000in}}%
\pgfpathclose%
\pgfusepath{fill}%
\end{pgfscope}%
\begin{pgfscope}%
\pgfpathrectangle{\pgfqpoint{1.250000in}{0.660000in}}{\pgfqpoint{7.750000in}{4.620000in}}%
\pgfusepath{clip}%
\pgfsetbuttcap%
\pgfsetmiterjoin%
\definecolor{currentfill}{rgb}{0.000000,0.000000,0.000000}%
\pgfsetfillcolor{currentfill}%
\pgfsetlinewidth{0.000000pt}%
\definecolor{currentstroke}{rgb}{0.000000,0.000000,0.000000}%
\pgfsetstrokecolor{currentstroke}%
\pgfsetstrokeopacity{0.000000}%
\pgfsetdash{}{0pt}%
\pgfpathmoveto{\pgfqpoint{3.757900in}{0.660000in}}%
\pgfpathlineto{\pgfqpoint{3.764100in}{0.660000in}}%
\pgfpathlineto{\pgfqpoint{3.764100in}{3.334759in}}%
\pgfpathlineto{\pgfqpoint{3.757900in}{3.334759in}}%
\pgfpathlineto{\pgfqpoint{3.757900in}{0.660000in}}%
\pgfpathclose%
\pgfusepath{fill}%
\end{pgfscope}%
\begin{pgfscope}%
\pgfpathrectangle{\pgfqpoint{1.250000in}{0.660000in}}{\pgfqpoint{7.750000in}{4.620000in}}%
\pgfusepath{clip}%
\pgfsetbuttcap%
\pgfsetmiterjoin%
\definecolor{currentfill}{rgb}{0.000000,0.000000,0.000000}%
\pgfsetfillcolor{currentfill}%
\pgfsetlinewidth{0.000000pt}%
\definecolor{currentstroke}{rgb}{0.000000,0.000000,0.000000}%
\pgfsetstrokecolor{currentstroke}%
\pgfsetstrokeopacity{0.000000}%
\pgfsetdash{}{0pt}%
\pgfpathmoveto{\pgfqpoint{3.765650in}{0.660000in}}%
\pgfpathlineto{\pgfqpoint{3.771850in}{0.660000in}}%
\pgfpathlineto{\pgfqpoint{3.771850in}{3.322483in}}%
\pgfpathlineto{\pgfqpoint{3.765650in}{3.322483in}}%
\pgfpathlineto{\pgfqpoint{3.765650in}{0.660000in}}%
\pgfpathclose%
\pgfusepath{fill}%
\end{pgfscope}%
\begin{pgfscope}%
\pgfpathrectangle{\pgfqpoint{1.250000in}{0.660000in}}{\pgfqpoint{7.750000in}{4.620000in}}%
\pgfusepath{clip}%
\pgfsetbuttcap%
\pgfsetmiterjoin%
\definecolor{currentfill}{rgb}{0.000000,0.000000,0.000000}%
\pgfsetfillcolor{currentfill}%
\pgfsetlinewidth{0.000000pt}%
\definecolor{currentstroke}{rgb}{0.000000,0.000000,0.000000}%
\pgfsetstrokecolor{currentstroke}%
\pgfsetstrokeopacity{0.000000}%
\pgfsetdash{}{0pt}%
\pgfpathmoveto{\pgfqpoint{3.773400in}{0.660000in}}%
\pgfpathlineto{\pgfqpoint{3.779600in}{0.660000in}}%
\pgfpathlineto{\pgfqpoint{3.779600in}{3.340881in}}%
\pgfpathlineto{\pgfqpoint{3.773400in}{3.340881in}}%
\pgfpathlineto{\pgfqpoint{3.773400in}{0.660000in}}%
\pgfpathclose%
\pgfusepath{fill}%
\end{pgfscope}%
\begin{pgfscope}%
\pgfpathrectangle{\pgfqpoint{1.250000in}{0.660000in}}{\pgfqpoint{7.750000in}{4.620000in}}%
\pgfusepath{clip}%
\pgfsetbuttcap%
\pgfsetmiterjoin%
\definecolor{currentfill}{rgb}{0.000000,0.000000,0.000000}%
\pgfsetfillcolor{currentfill}%
\pgfsetlinewidth{0.000000pt}%
\definecolor{currentstroke}{rgb}{0.000000,0.000000,0.000000}%
\pgfsetstrokecolor{currentstroke}%
\pgfsetstrokeopacity{0.000000}%
\pgfsetdash{}{0pt}%
\pgfpathmoveto{\pgfqpoint{3.781150in}{0.660000in}}%
\pgfpathlineto{\pgfqpoint{3.787350in}{0.660000in}}%
\pgfpathlineto{\pgfqpoint{3.787350in}{3.345138in}}%
\pgfpathlineto{\pgfqpoint{3.781150in}{3.345138in}}%
\pgfpathlineto{\pgfqpoint{3.781150in}{0.660000in}}%
\pgfpathclose%
\pgfusepath{fill}%
\end{pgfscope}%
\begin{pgfscope}%
\pgfpathrectangle{\pgfqpoint{1.250000in}{0.660000in}}{\pgfqpoint{7.750000in}{4.620000in}}%
\pgfusepath{clip}%
\pgfsetbuttcap%
\pgfsetmiterjoin%
\definecolor{currentfill}{rgb}{0.000000,0.000000,0.000000}%
\pgfsetfillcolor{currentfill}%
\pgfsetlinewidth{0.000000pt}%
\definecolor{currentstroke}{rgb}{0.000000,0.000000,0.000000}%
\pgfsetstrokecolor{currentstroke}%
\pgfsetstrokeopacity{0.000000}%
\pgfsetdash{}{0pt}%
\pgfpathmoveto{\pgfqpoint{3.788900in}{0.660000in}}%
\pgfpathlineto{\pgfqpoint{3.795100in}{0.660000in}}%
\pgfpathlineto{\pgfqpoint{3.795100in}{3.324509in}}%
\pgfpathlineto{\pgfqpoint{3.788900in}{3.324509in}}%
\pgfpathlineto{\pgfqpoint{3.788900in}{0.660000in}}%
\pgfpathclose%
\pgfusepath{fill}%
\end{pgfscope}%
\begin{pgfscope}%
\pgfpathrectangle{\pgfqpoint{1.250000in}{0.660000in}}{\pgfqpoint{7.750000in}{4.620000in}}%
\pgfusepath{clip}%
\pgfsetbuttcap%
\pgfsetmiterjoin%
\definecolor{currentfill}{rgb}{0.000000,0.000000,0.000000}%
\pgfsetfillcolor{currentfill}%
\pgfsetlinewidth{0.000000pt}%
\definecolor{currentstroke}{rgb}{0.000000,0.000000,0.000000}%
\pgfsetstrokecolor{currentstroke}%
\pgfsetstrokeopacity{0.000000}%
\pgfsetdash{}{0pt}%
\pgfpathmoveto{\pgfqpoint{3.796650in}{0.660000in}}%
\pgfpathlineto{\pgfqpoint{3.802850in}{0.660000in}}%
\pgfpathlineto{\pgfqpoint{3.802850in}{3.355890in}}%
\pgfpathlineto{\pgfqpoint{3.796650in}{3.355890in}}%
\pgfpathlineto{\pgfqpoint{3.796650in}{0.660000in}}%
\pgfpathclose%
\pgfusepath{fill}%
\end{pgfscope}%
\begin{pgfscope}%
\pgfpathrectangle{\pgfqpoint{1.250000in}{0.660000in}}{\pgfqpoint{7.750000in}{4.620000in}}%
\pgfusepath{clip}%
\pgfsetbuttcap%
\pgfsetmiterjoin%
\definecolor{currentfill}{rgb}{0.000000,0.000000,0.000000}%
\pgfsetfillcolor{currentfill}%
\pgfsetlinewidth{0.000000pt}%
\definecolor{currentstroke}{rgb}{0.000000,0.000000,0.000000}%
\pgfsetstrokecolor{currentstroke}%
\pgfsetstrokeopacity{0.000000}%
\pgfsetdash{}{0pt}%
\pgfpathmoveto{\pgfqpoint{3.804400in}{0.660000in}}%
\pgfpathlineto{\pgfqpoint{3.810600in}{0.660000in}}%
\pgfpathlineto{\pgfqpoint{3.810600in}{3.300232in}}%
\pgfpathlineto{\pgfqpoint{3.804400in}{3.300232in}}%
\pgfpathlineto{\pgfqpoint{3.804400in}{0.660000in}}%
\pgfpathclose%
\pgfusepath{fill}%
\end{pgfscope}%
\begin{pgfscope}%
\pgfpathrectangle{\pgfqpoint{1.250000in}{0.660000in}}{\pgfqpoint{7.750000in}{4.620000in}}%
\pgfusepath{clip}%
\pgfsetbuttcap%
\pgfsetmiterjoin%
\definecolor{currentfill}{rgb}{0.000000,0.000000,0.000000}%
\pgfsetfillcolor{currentfill}%
\pgfsetlinewidth{0.000000pt}%
\definecolor{currentstroke}{rgb}{0.000000,0.000000,0.000000}%
\pgfsetstrokecolor{currentstroke}%
\pgfsetstrokeopacity{0.000000}%
\pgfsetdash{}{0pt}%
\pgfpathmoveto{\pgfqpoint{3.812150in}{0.660000in}}%
\pgfpathlineto{\pgfqpoint{3.818350in}{0.660000in}}%
\pgfpathlineto{\pgfqpoint{3.818350in}{3.304544in}}%
\pgfpathlineto{\pgfqpoint{3.812150in}{3.304544in}}%
\pgfpathlineto{\pgfqpoint{3.812150in}{0.660000in}}%
\pgfpathclose%
\pgfusepath{fill}%
\end{pgfscope}%
\begin{pgfscope}%
\pgfpathrectangle{\pgfqpoint{1.250000in}{0.660000in}}{\pgfqpoint{7.750000in}{4.620000in}}%
\pgfusepath{clip}%
\pgfsetbuttcap%
\pgfsetmiterjoin%
\definecolor{currentfill}{rgb}{0.000000,0.000000,0.000000}%
\pgfsetfillcolor{currentfill}%
\pgfsetlinewidth{0.000000pt}%
\definecolor{currentstroke}{rgb}{0.000000,0.000000,0.000000}%
\pgfsetstrokecolor{currentstroke}%
\pgfsetstrokeopacity{0.000000}%
\pgfsetdash{}{0pt}%
\pgfpathmoveto{\pgfqpoint{3.819900in}{0.660000in}}%
\pgfpathlineto{\pgfqpoint{3.826100in}{0.660000in}}%
\pgfpathlineto{\pgfqpoint{3.826100in}{3.297970in}}%
\pgfpathlineto{\pgfqpoint{3.819900in}{3.297970in}}%
\pgfpathlineto{\pgfqpoint{3.819900in}{0.660000in}}%
\pgfpathclose%
\pgfusepath{fill}%
\end{pgfscope}%
\begin{pgfscope}%
\pgfpathrectangle{\pgfqpoint{1.250000in}{0.660000in}}{\pgfqpoint{7.750000in}{4.620000in}}%
\pgfusepath{clip}%
\pgfsetbuttcap%
\pgfsetmiterjoin%
\definecolor{currentfill}{rgb}{0.000000,0.000000,0.000000}%
\pgfsetfillcolor{currentfill}%
\pgfsetlinewidth{0.000000pt}%
\definecolor{currentstroke}{rgb}{0.000000,0.000000,0.000000}%
\pgfsetstrokecolor{currentstroke}%
\pgfsetstrokeopacity{0.000000}%
\pgfsetdash{}{0pt}%
\pgfpathmoveto{\pgfqpoint{3.827650in}{0.660000in}}%
\pgfpathlineto{\pgfqpoint{3.833850in}{0.660000in}}%
\pgfpathlineto{\pgfqpoint{3.833850in}{3.316279in}}%
\pgfpathlineto{\pgfqpoint{3.827650in}{3.316279in}}%
\pgfpathlineto{\pgfqpoint{3.827650in}{0.660000in}}%
\pgfpathclose%
\pgfusepath{fill}%
\end{pgfscope}%
\begin{pgfscope}%
\pgfpathrectangle{\pgfqpoint{1.250000in}{0.660000in}}{\pgfqpoint{7.750000in}{4.620000in}}%
\pgfusepath{clip}%
\pgfsetbuttcap%
\pgfsetmiterjoin%
\definecolor{currentfill}{rgb}{0.000000,0.000000,0.000000}%
\pgfsetfillcolor{currentfill}%
\pgfsetlinewidth{0.000000pt}%
\definecolor{currentstroke}{rgb}{0.000000,0.000000,0.000000}%
\pgfsetstrokecolor{currentstroke}%
\pgfsetstrokeopacity{0.000000}%
\pgfsetdash{}{0pt}%
\pgfpathmoveto{\pgfqpoint{3.835400in}{0.660000in}}%
\pgfpathlineto{\pgfqpoint{3.841600in}{0.660000in}}%
\pgfpathlineto{\pgfqpoint{3.841600in}{3.298748in}}%
\pgfpathlineto{\pgfqpoint{3.835400in}{3.298748in}}%
\pgfpathlineto{\pgfqpoint{3.835400in}{0.660000in}}%
\pgfpathclose%
\pgfusepath{fill}%
\end{pgfscope}%
\begin{pgfscope}%
\pgfpathrectangle{\pgfqpoint{1.250000in}{0.660000in}}{\pgfqpoint{7.750000in}{4.620000in}}%
\pgfusepath{clip}%
\pgfsetbuttcap%
\pgfsetmiterjoin%
\definecolor{currentfill}{rgb}{0.000000,0.000000,0.000000}%
\pgfsetfillcolor{currentfill}%
\pgfsetlinewidth{0.000000pt}%
\definecolor{currentstroke}{rgb}{0.000000,0.000000,0.000000}%
\pgfsetstrokecolor{currentstroke}%
\pgfsetstrokeopacity{0.000000}%
\pgfsetdash{}{0pt}%
\pgfpathmoveto{\pgfqpoint{3.843150in}{0.660000in}}%
\pgfpathlineto{\pgfqpoint{3.849350in}{0.660000in}}%
\pgfpathlineto{\pgfqpoint{3.849350in}{3.327185in}}%
\pgfpathlineto{\pgfqpoint{3.843150in}{3.327185in}}%
\pgfpathlineto{\pgfqpoint{3.843150in}{0.660000in}}%
\pgfpathclose%
\pgfusepath{fill}%
\end{pgfscope}%
\begin{pgfscope}%
\pgfpathrectangle{\pgfqpoint{1.250000in}{0.660000in}}{\pgfqpoint{7.750000in}{4.620000in}}%
\pgfusepath{clip}%
\pgfsetbuttcap%
\pgfsetmiterjoin%
\definecolor{currentfill}{rgb}{0.000000,0.000000,0.000000}%
\pgfsetfillcolor{currentfill}%
\pgfsetlinewidth{0.000000pt}%
\definecolor{currentstroke}{rgb}{0.000000,0.000000,0.000000}%
\pgfsetstrokecolor{currentstroke}%
\pgfsetstrokeopacity{0.000000}%
\pgfsetdash{}{0pt}%
\pgfpathmoveto{\pgfqpoint{3.850900in}{0.660000in}}%
\pgfpathlineto{\pgfqpoint{3.857100in}{0.660000in}}%
\pgfpathlineto{\pgfqpoint{3.857100in}{3.367562in}}%
\pgfpathlineto{\pgfqpoint{3.850900in}{3.367562in}}%
\pgfpathlineto{\pgfqpoint{3.850900in}{0.660000in}}%
\pgfpathclose%
\pgfusepath{fill}%
\end{pgfscope}%
\begin{pgfscope}%
\pgfpathrectangle{\pgfqpoint{1.250000in}{0.660000in}}{\pgfqpoint{7.750000in}{4.620000in}}%
\pgfusepath{clip}%
\pgfsetbuttcap%
\pgfsetmiterjoin%
\definecolor{currentfill}{rgb}{0.000000,0.000000,0.000000}%
\pgfsetfillcolor{currentfill}%
\pgfsetlinewidth{0.000000pt}%
\definecolor{currentstroke}{rgb}{0.000000,0.000000,0.000000}%
\pgfsetstrokecolor{currentstroke}%
\pgfsetstrokeopacity{0.000000}%
\pgfsetdash{}{0pt}%
\pgfpathmoveto{\pgfqpoint{3.858650in}{0.660000in}}%
\pgfpathlineto{\pgfqpoint{3.864850in}{0.660000in}}%
\pgfpathlineto{\pgfqpoint{3.864850in}{3.333216in}}%
\pgfpathlineto{\pgfqpoint{3.858650in}{3.333216in}}%
\pgfpathlineto{\pgfqpoint{3.858650in}{0.660000in}}%
\pgfpathclose%
\pgfusepath{fill}%
\end{pgfscope}%
\begin{pgfscope}%
\pgfpathrectangle{\pgfqpoint{1.250000in}{0.660000in}}{\pgfqpoint{7.750000in}{4.620000in}}%
\pgfusepath{clip}%
\pgfsetbuttcap%
\pgfsetmiterjoin%
\definecolor{currentfill}{rgb}{0.000000,0.000000,0.000000}%
\pgfsetfillcolor{currentfill}%
\pgfsetlinewidth{0.000000pt}%
\definecolor{currentstroke}{rgb}{0.000000,0.000000,0.000000}%
\pgfsetstrokecolor{currentstroke}%
\pgfsetstrokeopacity{0.000000}%
\pgfsetdash{}{0pt}%
\pgfpathmoveto{\pgfqpoint{3.866400in}{0.660000in}}%
\pgfpathlineto{\pgfqpoint{3.872600in}{0.660000in}}%
\pgfpathlineto{\pgfqpoint{3.872600in}{3.346964in}}%
\pgfpathlineto{\pgfqpoint{3.866400in}{3.346964in}}%
\pgfpathlineto{\pgfqpoint{3.866400in}{0.660000in}}%
\pgfpathclose%
\pgfusepath{fill}%
\end{pgfscope}%
\begin{pgfscope}%
\pgfpathrectangle{\pgfqpoint{1.250000in}{0.660000in}}{\pgfqpoint{7.750000in}{4.620000in}}%
\pgfusepath{clip}%
\pgfsetbuttcap%
\pgfsetmiterjoin%
\definecolor{currentfill}{rgb}{0.000000,0.000000,0.000000}%
\pgfsetfillcolor{currentfill}%
\pgfsetlinewidth{0.000000pt}%
\definecolor{currentstroke}{rgb}{0.000000,0.000000,0.000000}%
\pgfsetstrokecolor{currentstroke}%
\pgfsetstrokeopacity{0.000000}%
\pgfsetdash{}{0pt}%
\pgfpathmoveto{\pgfqpoint{3.874150in}{0.660000in}}%
\pgfpathlineto{\pgfqpoint{3.880350in}{0.660000in}}%
\pgfpathlineto{\pgfqpoint{3.880350in}{3.373648in}}%
\pgfpathlineto{\pgfqpoint{3.874150in}{3.373648in}}%
\pgfpathlineto{\pgfqpoint{3.874150in}{0.660000in}}%
\pgfpathclose%
\pgfusepath{fill}%
\end{pgfscope}%
\begin{pgfscope}%
\pgfpathrectangle{\pgfqpoint{1.250000in}{0.660000in}}{\pgfqpoint{7.750000in}{4.620000in}}%
\pgfusepath{clip}%
\pgfsetbuttcap%
\pgfsetmiterjoin%
\definecolor{currentfill}{rgb}{0.000000,0.000000,0.000000}%
\pgfsetfillcolor{currentfill}%
\pgfsetlinewidth{0.000000pt}%
\definecolor{currentstroke}{rgb}{0.000000,0.000000,0.000000}%
\pgfsetstrokecolor{currentstroke}%
\pgfsetstrokeopacity{0.000000}%
\pgfsetdash{}{0pt}%
\pgfpathmoveto{\pgfqpoint{3.881900in}{0.660000in}}%
\pgfpathlineto{\pgfqpoint{3.888100in}{0.660000in}}%
\pgfpathlineto{\pgfqpoint{3.888100in}{3.299288in}}%
\pgfpathlineto{\pgfqpoint{3.881900in}{3.299288in}}%
\pgfpathlineto{\pgfqpoint{3.881900in}{0.660000in}}%
\pgfpathclose%
\pgfusepath{fill}%
\end{pgfscope}%
\begin{pgfscope}%
\pgfpathrectangle{\pgfqpoint{1.250000in}{0.660000in}}{\pgfqpoint{7.750000in}{4.620000in}}%
\pgfusepath{clip}%
\pgfsetbuttcap%
\pgfsetmiterjoin%
\definecolor{currentfill}{rgb}{0.000000,0.000000,0.000000}%
\pgfsetfillcolor{currentfill}%
\pgfsetlinewidth{0.000000pt}%
\definecolor{currentstroke}{rgb}{0.000000,0.000000,0.000000}%
\pgfsetstrokecolor{currentstroke}%
\pgfsetstrokeopacity{0.000000}%
\pgfsetdash{}{0pt}%
\pgfpathmoveto{\pgfqpoint{3.889650in}{0.660000in}}%
\pgfpathlineto{\pgfqpoint{3.895850in}{0.660000in}}%
\pgfpathlineto{\pgfqpoint{3.895850in}{3.339859in}}%
\pgfpathlineto{\pgfqpoint{3.889650in}{3.339859in}}%
\pgfpathlineto{\pgfqpoint{3.889650in}{0.660000in}}%
\pgfpathclose%
\pgfusepath{fill}%
\end{pgfscope}%
\begin{pgfscope}%
\pgfpathrectangle{\pgfqpoint{1.250000in}{0.660000in}}{\pgfqpoint{7.750000in}{4.620000in}}%
\pgfusepath{clip}%
\pgfsetbuttcap%
\pgfsetmiterjoin%
\definecolor{currentfill}{rgb}{0.000000,0.000000,0.000000}%
\pgfsetfillcolor{currentfill}%
\pgfsetlinewidth{0.000000pt}%
\definecolor{currentstroke}{rgb}{0.000000,0.000000,0.000000}%
\pgfsetstrokecolor{currentstroke}%
\pgfsetstrokeopacity{0.000000}%
\pgfsetdash{}{0pt}%
\pgfpathmoveto{\pgfqpoint{3.897400in}{0.660000in}}%
\pgfpathlineto{\pgfqpoint{3.903600in}{0.660000in}}%
\pgfpathlineto{\pgfqpoint{3.903600in}{3.295851in}}%
\pgfpathlineto{\pgfqpoint{3.897400in}{3.295851in}}%
\pgfpathlineto{\pgfqpoint{3.897400in}{0.660000in}}%
\pgfpathclose%
\pgfusepath{fill}%
\end{pgfscope}%
\begin{pgfscope}%
\pgfpathrectangle{\pgfqpoint{1.250000in}{0.660000in}}{\pgfqpoint{7.750000in}{4.620000in}}%
\pgfusepath{clip}%
\pgfsetbuttcap%
\pgfsetmiterjoin%
\definecolor{currentfill}{rgb}{0.000000,0.000000,0.000000}%
\pgfsetfillcolor{currentfill}%
\pgfsetlinewidth{0.000000pt}%
\definecolor{currentstroke}{rgb}{0.000000,0.000000,0.000000}%
\pgfsetstrokecolor{currentstroke}%
\pgfsetstrokeopacity{0.000000}%
\pgfsetdash{}{0pt}%
\pgfpathmoveto{\pgfqpoint{3.905150in}{0.660000in}}%
\pgfpathlineto{\pgfqpoint{3.911350in}{0.660000in}}%
\pgfpathlineto{\pgfqpoint{3.911350in}{3.325779in}}%
\pgfpathlineto{\pgfqpoint{3.905150in}{3.325779in}}%
\pgfpathlineto{\pgfqpoint{3.905150in}{0.660000in}}%
\pgfpathclose%
\pgfusepath{fill}%
\end{pgfscope}%
\begin{pgfscope}%
\pgfpathrectangle{\pgfqpoint{1.250000in}{0.660000in}}{\pgfqpoint{7.750000in}{4.620000in}}%
\pgfusepath{clip}%
\pgfsetbuttcap%
\pgfsetmiterjoin%
\definecolor{currentfill}{rgb}{0.000000,0.000000,0.000000}%
\pgfsetfillcolor{currentfill}%
\pgfsetlinewidth{0.000000pt}%
\definecolor{currentstroke}{rgb}{0.000000,0.000000,0.000000}%
\pgfsetstrokecolor{currentstroke}%
\pgfsetstrokeopacity{0.000000}%
\pgfsetdash{}{0pt}%
\pgfpathmoveto{\pgfqpoint{3.912900in}{0.660000in}}%
\pgfpathlineto{\pgfqpoint{3.919100in}{0.660000in}}%
\pgfpathlineto{\pgfqpoint{3.919100in}{3.295365in}}%
\pgfpathlineto{\pgfqpoint{3.912900in}{3.295365in}}%
\pgfpathlineto{\pgfqpoint{3.912900in}{0.660000in}}%
\pgfpathclose%
\pgfusepath{fill}%
\end{pgfscope}%
\begin{pgfscope}%
\pgfpathrectangle{\pgfqpoint{1.250000in}{0.660000in}}{\pgfqpoint{7.750000in}{4.620000in}}%
\pgfusepath{clip}%
\pgfsetbuttcap%
\pgfsetmiterjoin%
\definecolor{currentfill}{rgb}{0.000000,0.000000,0.000000}%
\pgfsetfillcolor{currentfill}%
\pgfsetlinewidth{0.000000pt}%
\definecolor{currentstroke}{rgb}{0.000000,0.000000,0.000000}%
\pgfsetstrokecolor{currentstroke}%
\pgfsetstrokeopacity{0.000000}%
\pgfsetdash{}{0pt}%
\pgfpathmoveto{\pgfqpoint{3.920650in}{0.660000in}}%
\pgfpathlineto{\pgfqpoint{3.926850in}{0.660000in}}%
\pgfpathlineto{\pgfqpoint{3.926850in}{3.298045in}}%
\pgfpathlineto{\pgfqpoint{3.920650in}{3.298045in}}%
\pgfpathlineto{\pgfqpoint{3.920650in}{0.660000in}}%
\pgfpathclose%
\pgfusepath{fill}%
\end{pgfscope}%
\begin{pgfscope}%
\pgfpathrectangle{\pgfqpoint{1.250000in}{0.660000in}}{\pgfqpoint{7.750000in}{4.620000in}}%
\pgfusepath{clip}%
\pgfsetbuttcap%
\pgfsetmiterjoin%
\definecolor{currentfill}{rgb}{0.000000,0.000000,0.000000}%
\pgfsetfillcolor{currentfill}%
\pgfsetlinewidth{0.000000pt}%
\definecolor{currentstroke}{rgb}{0.000000,0.000000,0.000000}%
\pgfsetstrokecolor{currentstroke}%
\pgfsetstrokeopacity{0.000000}%
\pgfsetdash{}{0pt}%
\pgfpathmoveto{\pgfqpoint{3.928400in}{0.660000in}}%
\pgfpathlineto{\pgfqpoint{3.934600in}{0.660000in}}%
\pgfpathlineto{\pgfqpoint{3.934600in}{3.328712in}}%
\pgfpathlineto{\pgfqpoint{3.928400in}{3.328712in}}%
\pgfpathlineto{\pgfqpoint{3.928400in}{0.660000in}}%
\pgfpathclose%
\pgfusepath{fill}%
\end{pgfscope}%
\begin{pgfscope}%
\pgfpathrectangle{\pgfqpoint{1.250000in}{0.660000in}}{\pgfqpoint{7.750000in}{4.620000in}}%
\pgfusepath{clip}%
\pgfsetbuttcap%
\pgfsetmiterjoin%
\definecolor{currentfill}{rgb}{0.000000,0.000000,0.000000}%
\pgfsetfillcolor{currentfill}%
\pgfsetlinewidth{0.000000pt}%
\definecolor{currentstroke}{rgb}{0.000000,0.000000,0.000000}%
\pgfsetstrokecolor{currentstroke}%
\pgfsetstrokeopacity{0.000000}%
\pgfsetdash{}{0pt}%
\pgfpathmoveto{\pgfqpoint{3.936150in}{0.660000in}}%
\pgfpathlineto{\pgfqpoint{3.942350in}{0.660000in}}%
\pgfpathlineto{\pgfqpoint{3.942350in}{3.309740in}}%
\pgfpathlineto{\pgfqpoint{3.936150in}{3.309740in}}%
\pgfpathlineto{\pgfqpoint{3.936150in}{0.660000in}}%
\pgfpathclose%
\pgfusepath{fill}%
\end{pgfscope}%
\begin{pgfscope}%
\pgfpathrectangle{\pgfqpoint{1.250000in}{0.660000in}}{\pgfqpoint{7.750000in}{4.620000in}}%
\pgfusepath{clip}%
\pgfsetbuttcap%
\pgfsetmiterjoin%
\definecolor{currentfill}{rgb}{0.000000,0.000000,0.000000}%
\pgfsetfillcolor{currentfill}%
\pgfsetlinewidth{0.000000pt}%
\definecolor{currentstroke}{rgb}{0.000000,0.000000,0.000000}%
\pgfsetstrokecolor{currentstroke}%
\pgfsetstrokeopacity{0.000000}%
\pgfsetdash{}{0pt}%
\pgfpathmoveto{\pgfqpoint{3.943900in}{0.660000in}}%
\pgfpathlineto{\pgfqpoint{3.950100in}{0.660000in}}%
\pgfpathlineto{\pgfqpoint{3.950100in}{3.302763in}}%
\pgfpathlineto{\pgfqpoint{3.943900in}{3.302763in}}%
\pgfpathlineto{\pgfqpoint{3.943900in}{0.660000in}}%
\pgfpathclose%
\pgfusepath{fill}%
\end{pgfscope}%
\begin{pgfscope}%
\pgfpathrectangle{\pgfqpoint{1.250000in}{0.660000in}}{\pgfqpoint{7.750000in}{4.620000in}}%
\pgfusepath{clip}%
\pgfsetbuttcap%
\pgfsetmiterjoin%
\definecolor{currentfill}{rgb}{0.000000,0.000000,0.000000}%
\pgfsetfillcolor{currentfill}%
\pgfsetlinewidth{0.000000pt}%
\definecolor{currentstroke}{rgb}{0.000000,0.000000,0.000000}%
\pgfsetstrokecolor{currentstroke}%
\pgfsetstrokeopacity{0.000000}%
\pgfsetdash{}{0pt}%
\pgfpathmoveto{\pgfqpoint{3.951650in}{0.660000in}}%
\pgfpathlineto{\pgfqpoint{3.957850in}{0.660000in}}%
\pgfpathlineto{\pgfqpoint{3.957850in}{3.301889in}}%
\pgfpathlineto{\pgfqpoint{3.951650in}{3.301889in}}%
\pgfpathlineto{\pgfqpoint{3.951650in}{0.660000in}}%
\pgfpathclose%
\pgfusepath{fill}%
\end{pgfscope}%
\begin{pgfscope}%
\pgfpathrectangle{\pgfqpoint{1.250000in}{0.660000in}}{\pgfqpoint{7.750000in}{4.620000in}}%
\pgfusepath{clip}%
\pgfsetbuttcap%
\pgfsetmiterjoin%
\definecolor{currentfill}{rgb}{0.000000,0.000000,0.000000}%
\pgfsetfillcolor{currentfill}%
\pgfsetlinewidth{0.000000pt}%
\definecolor{currentstroke}{rgb}{0.000000,0.000000,0.000000}%
\pgfsetstrokecolor{currentstroke}%
\pgfsetstrokeopacity{0.000000}%
\pgfsetdash{}{0pt}%
\pgfpathmoveto{\pgfqpoint{3.959400in}{0.660000in}}%
\pgfpathlineto{\pgfqpoint{3.965600in}{0.660000in}}%
\pgfpathlineto{\pgfqpoint{3.965600in}{3.318105in}}%
\pgfpathlineto{\pgfqpoint{3.959400in}{3.318105in}}%
\pgfpathlineto{\pgfqpoint{3.959400in}{0.660000in}}%
\pgfpathclose%
\pgfusepath{fill}%
\end{pgfscope}%
\begin{pgfscope}%
\pgfpathrectangle{\pgfqpoint{1.250000in}{0.660000in}}{\pgfqpoint{7.750000in}{4.620000in}}%
\pgfusepath{clip}%
\pgfsetbuttcap%
\pgfsetmiterjoin%
\definecolor{currentfill}{rgb}{0.000000,0.000000,0.000000}%
\pgfsetfillcolor{currentfill}%
\pgfsetlinewidth{0.000000pt}%
\definecolor{currentstroke}{rgb}{0.000000,0.000000,0.000000}%
\pgfsetstrokecolor{currentstroke}%
\pgfsetstrokeopacity{0.000000}%
\pgfsetdash{}{0pt}%
\pgfpathmoveto{\pgfqpoint{3.967150in}{0.660000in}}%
\pgfpathlineto{\pgfqpoint{3.973350in}{0.660000in}}%
\pgfpathlineto{\pgfqpoint{3.973350in}{3.313942in}}%
\pgfpathlineto{\pgfqpoint{3.967150in}{3.313942in}}%
\pgfpathlineto{\pgfqpoint{3.967150in}{0.660000in}}%
\pgfpathclose%
\pgfusepath{fill}%
\end{pgfscope}%
\begin{pgfscope}%
\pgfpathrectangle{\pgfqpoint{1.250000in}{0.660000in}}{\pgfqpoint{7.750000in}{4.620000in}}%
\pgfusepath{clip}%
\pgfsetbuttcap%
\pgfsetmiterjoin%
\definecolor{currentfill}{rgb}{0.000000,0.000000,0.000000}%
\pgfsetfillcolor{currentfill}%
\pgfsetlinewidth{0.000000pt}%
\definecolor{currentstroke}{rgb}{0.000000,0.000000,0.000000}%
\pgfsetstrokecolor{currentstroke}%
\pgfsetstrokeopacity{0.000000}%
\pgfsetdash{}{0pt}%
\pgfpathmoveto{\pgfqpoint{3.974900in}{0.660000in}}%
\pgfpathlineto{\pgfqpoint{3.981100in}{0.660000in}}%
\pgfpathlineto{\pgfqpoint{3.981100in}{3.309345in}}%
\pgfpathlineto{\pgfqpoint{3.974900in}{3.309345in}}%
\pgfpathlineto{\pgfqpoint{3.974900in}{0.660000in}}%
\pgfpathclose%
\pgfusepath{fill}%
\end{pgfscope}%
\begin{pgfscope}%
\pgfpathrectangle{\pgfqpoint{1.250000in}{0.660000in}}{\pgfqpoint{7.750000in}{4.620000in}}%
\pgfusepath{clip}%
\pgfsetbuttcap%
\pgfsetmiterjoin%
\definecolor{currentfill}{rgb}{0.000000,0.000000,0.000000}%
\pgfsetfillcolor{currentfill}%
\pgfsetlinewidth{0.000000pt}%
\definecolor{currentstroke}{rgb}{0.000000,0.000000,0.000000}%
\pgfsetstrokecolor{currentstroke}%
\pgfsetstrokeopacity{0.000000}%
\pgfsetdash{}{0pt}%
\pgfpathmoveto{\pgfqpoint{3.982650in}{0.660000in}}%
\pgfpathlineto{\pgfqpoint{3.988850in}{0.660000in}}%
\pgfpathlineto{\pgfqpoint{3.988850in}{3.340511in}}%
\pgfpathlineto{\pgfqpoint{3.982650in}{3.340511in}}%
\pgfpathlineto{\pgfqpoint{3.982650in}{0.660000in}}%
\pgfpathclose%
\pgfusepath{fill}%
\end{pgfscope}%
\begin{pgfscope}%
\pgfpathrectangle{\pgfqpoint{1.250000in}{0.660000in}}{\pgfqpoint{7.750000in}{4.620000in}}%
\pgfusepath{clip}%
\pgfsetbuttcap%
\pgfsetmiterjoin%
\definecolor{currentfill}{rgb}{0.000000,0.000000,0.000000}%
\pgfsetfillcolor{currentfill}%
\pgfsetlinewidth{0.000000pt}%
\definecolor{currentstroke}{rgb}{0.000000,0.000000,0.000000}%
\pgfsetstrokecolor{currentstroke}%
\pgfsetstrokeopacity{0.000000}%
\pgfsetdash{}{0pt}%
\pgfpathmoveto{\pgfqpoint{3.990400in}{0.660000in}}%
\pgfpathlineto{\pgfqpoint{3.996600in}{0.660000in}}%
\pgfpathlineto{\pgfqpoint{3.996600in}{3.301701in}}%
\pgfpathlineto{\pgfqpoint{3.990400in}{3.301701in}}%
\pgfpathlineto{\pgfqpoint{3.990400in}{0.660000in}}%
\pgfpathclose%
\pgfusepath{fill}%
\end{pgfscope}%
\begin{pgfscope}%
\pgfpathrectangle{\pgfqpoint{1.250000in}{0.660000in}}{\pgfqpoint{7.750000in}{4.620000in}}%
\pgfusepath{clip}%
\pgfsetbuttcap%
\pgfsetmiterjoin%
\definecolor{currentfill}{rgb}{0.000000,0.000000,0.000000}%
\pgfsetfillcolor{currentfill}%
\pgfsetlinewidth{0.000000pt}%
\definecolor{currentstroke}{rgb}{0.000000,0.000000,0.000000}%
\pgfsetstrokecolor{currentstroke}%
\pgfsetstrokeopacity{0.000000}%
\pgfsetdash{}{0pt}%
\pgfpathmoveto{\pgfqpoint{3.998150in}{0.660000in}}%
\pgfpathlineto{\pgfqpoint{4.004350in}{0.660000in}}%
\pgfpathlineto{\pgfqpoint{4.004350in}{3.321340in}}%
\pgfpathlineto{\pgfqpoint{3.998150in}{3.321340in}}%
\pgfpathlineto{\pgfqpoint{3.998150in}{0.660000in}}%
\pgfpathclose%
\pgfusepath{fill}%
\end{pgfscope}%
\begin{pgfscope}%
\pgfpathrectangle{\pgfqpoint{1.250000in}{0.660000in}}{\pgfqpoint{7.750000in}{4.620000in}}%
\pgfusepath{clip}%
\pgfsetbuttcap%
\pgfsetmiterjoin%
\definecolor{currentfill}{rgb}{0.000000,0.000000,0.000000}%
\pgfsetfillcolor{currentfill}%
\pgfsetlinewidth{0.000000pt}%
\definecolor{currentstroke}{rgb}{0.000000,0.000000,0.000000}%
\pgfsetstrokecolor{currentstroke}%
\pgfsetstrokeopacity{0.000000}%
\pgfsetdash{}{0pt}%
\pgfpathmoveto{\pgfqpoint{4.005900in}{0.660000in}}%
\pgfpathlineto{\pgfqpoint{4.012100in}{0.660000in}}%
\pgfpathlineto{\pgfqpoint{4.012100in}{3.302560in}}%
\pgfpathlineto{\pgfqpoint{4.005900in}{3.302560in}}%
\pgfpathlineto{\pgfqpoint{4.005900in}{0.660000in}}%
\pgfpathclose%
\pgfusepath{fill}%
\end{pgfscope}%
\begin{pgfscope}%
\pgfpathrectangle{\pgfqpoint{1.250000in}{0.660000in}}{\pgfqpoint{7.750000in}{4.620000in}}%
\pgfusepath{clip}%
\pgfsetbuttcap%
\pgfsetmiterjoin%
\definecolor{currentfill}{rgb}{0.000000,0.000000,0.000000}%
\pgfsetfillcolor{currentfill}%
\pgfsetlinewidth{0.000000pt}%
\definecolor{currentstroke}{rgb}{0.000000,0.000000,0.000000}%
\pgfsetstrokecolor{currentstroke}%
\pgfsetstrokeopacity{0.000000}%
\pgfsetdash{}{0pt}%
\pgfpathmoveto{\pgfqpoint{4.013650in}{0.660000in}}%
\pgfpathlineto{\pgfqpoint{4.019850in}{0.660000in}}%
\pgfpathlineto{\pgfqpoint{4.019850in}{3.290296in}}%
\pgfpathlineto{\pgfqpoint{4.013650in}{3.290296in}}%
\pgfpathlineto{\pgfqpoint{4.013650in}{0.660000in}}%
\pgfpathclose%
\pgfusepath{fill}%
\end{pgfscope}%
\begin{pgfscope}%
\pgfpathrectangle{\pgfqpoint{1.250000in}{0.660000in}}{\pgfqpoint{7.750000in}{4.620000in}}%
\pgfusepath{clip}%
\pgfsetbuttcap%
\pgfsetmiterjoin%
\definecolor{currentfill}{rgb}{0.000000,0.000000,0.000000}%
\pgfsetfillcolor{currentfill}%
\pgfsetlinewidth{0.000000pt}%
\definecolor{currentstroke}{rgb}{0.000000,0.000000,0.000000}%
\pgfsetstrokecolor{currentstroke}%
\pgfsetstrokeopacity{0.000000}%
\pgfsetdash{}{0pt}%
\pgfpathmoveto{\pgfqpoint{4.021400in}{0.660000in}}%
\pgfpathlineto{\pgfqpoint{4.027600in}{0.660000in}}%
\pgfpathlineto{\pgfqpoint{4.027600in}{3.325987in}}%
\pgfpathlineto{\pgfqpoint{4.021400in}{3.325987in}}%
\pgfpathlineto{\pgfqpoint{4.021400in}{0.660000in}}%
\pgfpathclose%
\pgfusepath{fill}%
\end{pgfscope}%
\begin{pgfscope}%
\pgfpathrectangle{\pgfqpoint{1.250000in}{0.660000in}}{\pgfqpoint{7.750000in}{4.620000in}}%
\pgfusepath{clip}%
\pgfsetbuttcap%
\pgfsetmiterjoin%
\definecolor{currentfill}{rgb}{0.000000,0.000000,0.000000}%
\pgfsetfillcolor{currentfill}%
\pgfsetlinewidth{0.000000pt}%
\definecolor{currentstroke}{rgb}{0.000000,0.000000,0.000000}%
\pgfsetstrokecolor{currentstroke}%
\pgfsetstrokeopacity{0.000000}%
\pgfsetdash{}{0pt}%
\pgfpathmoveto{\pgfqpoint{4.029150in}{0.660000in}}%
\pgfpathlineto{\pgfqpoint{4.035350in}{0.660000in}}%
\pgfpathlineto{\pgfqpoint{4.035350in}{3.337393in}}%
\pgfpathlineto{\pgfqpoint{4.029150in}{3.337393in}}%
\pgfpathlineto{\pgfqpoint{4.029150in}{0.660000in}}%
\pgfpathclose%
\pgfusepath{fill}%
\end{pgfscope}%
\begin{pgfscope}%
\pgfpathrectangle{\pgfqpoint{1.250000in}{0.660000in}}{\pgfqpoint{7.750000in}{4.620000in}}%
\pgfusepath{clip}%
\pgfsetbuttcap%
\pgfsetmiterjoin%
\definecolor{currentfill}{rgb}{0.000000,0.000000,0.000000}%
\pgfsetfillcolor{currentfill}%
\pgfsetlinewidth{0.000000pt}%
\definecolor{currentstroke}{rgb}{0.000000,0.000000,0.000000}%
\pgfsetstrokecolor{currentstroke}%
\pgfsetstrokeopacity{0.000000}%
\pgfsetdash{}{0pt}%
\pgfpathmoveto{\pgfqpoint{4.036900in}{0.660000in}}%
\pgfpathlineto{\pgfqpoint{4.043100in}{0.660000in}}%
\pgfpathlineto{\pgfqpoint{4.043100in}{3.315508in}}%
\pgfpathlineto{\pgfqpoint{4.036900in}{3.315508in}}%
\pgfpathlineto{\pgfqpoint{4.036900in}{0.660000in}}%
\pgfpathclose%
\pgfusepath{fill}%
\end{pgfscope}%
\begin{pgfscope}%
\pgfpathrectangle{\pgfqpoint{1.250000in}{0.660000in}}{\pgfqpoint{7.750000in}{4.620000in}}%
\pgfusepath{clip}%
\pgfsetbuttcap%
\pgfsetmiterjoin%
\definecolor{currentfill}{rgb}{0.000000,0.000000,0.000000}%
\pgfsetfillcolor{currentfill}%
\pgfsetlinewidth{0.000000pt}%
\definecolor{currentstroke}{rgb}{0.000000,0.000000,0.000000}%
\pgfsetstrokecolor{currentstroke}%
\pgfsetstrokeopacity{0.000000}%
\pgfsetdash{}{0pt}%
\pgfpathmoveto{\pgfqpoint{4.044650in}{0.660000in}}%
\pgfpathlineto{\pgfqpoint{4.050850in}{0.660000in}}%
\pgfpathlineto{\pgfqpoint{4.050850in}{3.306119in}}%
\pgfpathlineto{\pgfqpoint{4.044650in}{3.306119in}}%
\pgfpathlineto{\pgfqpoint{4.044650in}{0.660000in}}%
\pgfpathclose%
\pgfusepath{fill}%
\end{pgfscope}%
\begin{pgfscope}%
\pgfpathrectangle{\pgfqpoint{1.250000in}{0.660000in}}{\pgfqpoint{7.750000in}{4.620000in}}%
\pgfusepath{clip}%
\pgfsetbuttcap%
\pgfsetmiterjoin%
\definecolor{currentfill}{rgb}{0.000000,0.000000,0.000000}%
\pgfsetfillcolor{currentfill}%
\pgfsetlinewidth{0.000000pt}%
\definecolor{currentstroke}{rgb}{0.000000,0.000000,0.000000}%
\pgfsetstrokecolor{currentstroke}%
\pgfsetstrokeopacity{0.000000}%
\pgfsetdash{}{0pt}%
\pgfpathmoveto{\pgfqpoint{4.052400in}{0.660000in}}%
\pgfpathlineto{\pgfqpoint{4.058600in}{0.660000in}}%
\pgfpathlineto{\pgfqpoint{4.058600in}{3.294201in}}%
\pgfpathlineto{\pgfqpoint{4.052400in}{3.294201in}}%
\pgfpathlineto{\pgfqpoint{4.052400in}{0.660000in}}%
\pgfpathclose%
\pgfusepath{fill}%
\end{pgfscope}%
\begin{pgfscope}%
\pgfpathrectangle{\pgfqpoint{1.250000in}{0.660000in}}{\pgfqpoint{7.750000in}{4.620000in}}%
\pgfusepath{clip}%
\pgfsetbuttcap%
\pgfsetmiterjoin%
\definecolor{currentfill}{rgb}{0.000000,0.000000,0.000000}%
\pgfsetfillcolor{currentfill}%
\pgfsetlinewidth{0.000000pt}%
\definecolor{currentstroke}{rgb}{0.000000,0.000000,0.000000}%
\pgfsetstrokecolor{currentstroke}%
\pgfsetstrokeopacity{0.000000}%
\pgfsetdash{}{0pt}%
\pgfpathmoveto{\pgfqpoint{4.060150in}{0.660000in}}%
\pgfpathlineto{\pgfqpoint{4.066350in}{0.660000in}}%
\pgfpathlineto{\pgfqpoint{4.066350in}{3.324198in}}%
\pgfpathlineto{\pgfqpoint{4.060150in}{3.324198in}}%
\pgfpathlineto{\pgfqpoint{4.060150in}{0.660000in}}%
\pgfpathclose%
\pgfusepath{fill}%
\end{pgfscope}%
\begin{pgfscope}%
\pgfpathrectangle{\pgfqpoint{1.250000in}{0.660000in}}{\pgfqpoint{7.750000in}{4.620000in}}%
\pgfusepath{clip}%
\pgfsetbuttcap%
\pgfsetmiterjoin%
\definecolor{currentfill}{rgb}{0.000000,0.000000,0.000000}%
\pgfsetfillcolor{currentfill}%
\pgfsetlinewidth{0.000000pt}%
\definecolor{currentstroke}{rgb}{0.000000,0.000000,0.000000}%
\pgfsetstrokecolor{currentstroke}%
\pgfsetstrokeopacity{0.000000}%
\pgfsetdash{}{0pt}%
\pgfpathmoveto{\pgfqpoint{4.067900in}{0.660000in}}%
\pgfpathlineto{\pgfqpoint{4.074100in}{0.660000in}}%
\pgfpathlineto{\pgfqpoint{4.074100in}{3.303634in}}%
\pgfpathlineto{\pgfqpoint{4.067900in}{3.303634in}}%
\pgfpathlineto{\pgfqpoint{4.067900in}{0.660000in}}%
\pgfpathclose%
\pgfusepath{fill}%
\end{pgfscope}%
\begin{pgfscope}%
\pgfpathrectangle{\pgfqpoint{1.250000in}{0.660000in}}{\pgfqpoint{7.750000in}{4.620000in}}%
\pgfusepath{clip}%
\pgfsetbuttcap%
\pgfsetmiterjoin%
\definecolor{currentfill}{rgb}{0.000000,0.000000,0.000000}%
\pgfsetfillcolor{currentfill}%
\pgfsetlinewidth{0.000000pt}%
\definecolor{currentstroke}{rgb}{0.000000,0.000000,0.000000}%
\pgfsetstrokecolor{currentstroke}%
\pgfsetstrokeopacity{0.000000}%
\pgfsetdash{}{0pt}%
\pgfpathmoveto{\pgfqpoint{4.075650in}{0.660000in}}%
\pgfpathlineto{\pgfqpoint{4.081850in}{0.660000in}}%
\pgfpathlineto{\pgfqpoint{4.081850in}{3.318684in}}%
\pgfpathlineto{\pgfqpoint{4.075650in}{3.318684in}}%
\pgfpathlineto{\pgfqpoint{4.075650in}{0.660000in}}%
\pgfpathclose%
\pgfusepath{fill}%
\end{pgfscope}%
\begin{pgfscope}%
\pgfpathrectangle{\pgfqpoint{1.250000in}{0.660000in}}{\pgfqpoint{7.750000in}{4.620000in}}%
\pgfusepath{clip}%
\pgfsetbuttcap%
\pgfsetmiterjoin%
\definecolor{currentfill}{rgb}{0.000000,0.000000,0.000000}%
\pgfsetfillcolor{currentfill}%
\pgfsetlinewidth{0.000000pt}%
\definecolor{currentstroke}{rgb}{0.000000,0.000000,0.000000}%
\pgfsetstrokecolor{currentstroke}%
\pgfsetstrokeopacity{0.000000}%
\pgfsetdash{}{0pt}%
\pgfpathmoveto{\pgfqpoint{4.083400in}{0.660000in}}%
\pgfpathlineto{\pgfqpoint{4.089600in}{0.660000in}}%
\pgfpathlineto{\pgfqpoint{4.089600in}{3.306261in}}%
\pgfpathlineto{\pgfqpoint{4.083400in}{3.306261in}}%
\pgfpathlineto{\pgfqpoint{4.083400in}{0.660000in}}%
\pgfpathclose%
\pgfusepath{fill}%
\end{pgfscope}%
\begin{pgfscope}%
\pgfpathrectangle{\pgfqpoint{1.250000in}{0.660000in}}{\pgfqpoint{7.750000in}{4.620000in}}%
\pgfusepath{clip}%
\pgfsetbuttcap%
\pgfsetmiterjoin%
\definecolor{currentfill}{rgb}{0.000000,0.000000,0.000000}%
\pgfsetfillcolor{currentfill}%
\pgfsetlinewidth{0.000000pt}%
\definecolor{currentstroke}{rgb}{0.000000,0.000000,0.000000}%
\pgfsetstrokecolor{currentstroke}%
\pgfsetstrokeopacity{0.000000}%
\pgfsetdash{}{0pt}%
\pgfpathmoveto{\pgfqpoint{4.091150in}{0.660000in}}%
\pgfpathlineto{\pgfqpoint{4.097350in}{0.660000in}}%
\pgfpathlineto{\pgfqpoint{4.097350in}{3.306157in}}%
\pgfpathlineto{\pgfqpoint{4.091150in}{3.306157in}}%
\pgfpathlineto{\pgfqpoint{4.091150in}{0.660000in}}%
\pgfpathclose%
\pgfusepath{fill}%
\end{pgfscope}%
\begin{pgfscope}%
\pgfpathrectangle{\pgfqpoint{1.250000in}{0.660000in}}{\pgfqpoint{7.750000in}{4.620000in}}%
\pgfusepath{clip}%
\pgfsetbuttcap%
\pgfsetmiterjoin%
\definecolor{currentfill}{rgb}{0.000000,0.000000,0.000000}%
\pgfsetfillcolor{currentfill}%
\pgfsetlinewidth{0.000000pt}%
\definecolor{currentstroke}{rgb}{0.000000,0.000000,0.000000}%
\pgfsetstrokecolor{currentstroke}%
\pgfsetstrokeopacity{0.000000}%
\pgfsetdash{}{0pt}%
\pgfpathmoveto{\pgfqpoint{4.098900in}{0.660000in}}%
\pgfpathlineto{\pgfqpoint{4.105100in}{0.660000in}}%
\pgfpathlineto{\pgfqpoint{4.105100in}{3.286282in}}%
\pgfpathlineto{\pgfqpoint{4.098900in}{3.286282in}}%
\pgfpathlineto{\pgfqpoint{4.098900in}{0.660000in}}%
\pgfpathclose%
\pgfusepath{fill}%
\end{pgfscope}%
\begin{pgfscope}%
\pgfpathrectangle{\pgfqpoint{1.250000in}{0.660000in}}{\pgfqpoint{7.750000in}{4.620000in}}%
\pgfusepath{clip}%
\pgfsetbuttcap%
\pgfsetmiterjoin%
\definecolor{currentfill}{rgb}{0.000000,0.000000,0.000000}%
\pgfsetfillcolor{currentfill}%
\pgfsetlinewidth{0.000000pt}%
\definecolor{currentstroke}{rgb}{0.000000,0.000000,0.000000}%
\pgfsetstrokecolor{currentstroke}%
\pgfsetstrokeopacity{0.000000}%
\pgfsetdash{}{0pt}%
\pgfpathmoveto{\pgfqpoint{4.106650in}{0.660000in}}%
\pgfpathlineto{\pgfqpoint{4.112850in}{0.660000in}}%
\pgfpathlineto{\pgfqpoint{4.112850in}{3.310094in}}%
\pgfpathlineto{\pgfqpoint{4.106650in}{3.310094in}}%
\pgfpathlineto{\pgfqpoint{4.106650in}{0.660000in}}%
\pgfpathclose%
\pgfusepath{fill}%
\end{pgfscope}%
\begin{pgfscope}%
\pgfpathrectangle{\pgfqpoint{1.250000in}{0.660000in}}{\pgfqpoint{7.750000in}{4.620000in}}%
\pgfusepath{clip}%
\pgfsetbuttcap%
\pgfsetmiterjoin%
\definecolor{currentfill}{rgb}{0.000000,0.000000,0.000000}%
\pgfsetfillcolor{currentfill}%
\pgfsetlinewidth{0.000000pt}%
\definecolor{currentstroke}{rgb}{0.000000,0.000000,0.000000}%
\pgfsetstrokecolor{currentstroke}%
\pgfsetstrokeopacity{0.000000}%
\pgfsetdash{}{0pt}%
\pgfpathmoveto{\pgfqpoint{4.114400in}{0.660000in}}%
\pgfpathlineto{\pgfqpoint{4.120600in}{0.660000in}}%
\pgfpathlineto{\pgfqpoint{4.120600in}{3.306417in}}%
\pgfpathlineto{\pgfqpoint{4.114400in}{3.306417in}}%
\pgfpathlineto{\pgfqpoint{4.114400in}{0.660000in}}%
\pgfpathclose%
\pgfusepath{fill}%
\end{pgfscope}%
\begin{pgfscope}%
\pgfpathrectangle{\pgfqpoint{1.250000in}{0.660000in}}{\pgfqpoint{7.750000in}{4.620000in}}%
\pgfusepath{clip}%
\pgfsetbuttcap%
\pgfsetmiterjoin%
\definecolor{currentfill}{rgb}{0.000000,0.000000,0.000000}%
\pgfsetfillcolor{currentfill}%
\pgfsetlinewidth{0.000000pt}%
\definecolor{currentstroke}{rgb}{0.000000,0.000000,0.000000}%
\pgfsetstrokecolor{currentstroke}%
\pgfsetstrokeopacity{0.000000}%
\pgfsetdash{}{0pt}%
\pgfpathmoveto{\pgfqpoint{4.122150in}{0.660000in}}%
\pgfpathlineto{\pgfqpoint{4.128350in}{0.660000in}}%
\pgfpathlineto{\pgfqpoint{4.128350in}{3.309014in}}%
\pgfpathlineto{\pgfqpoint{4.122150in}{3.309014in}}%
\pgfpathlineto{\pgfqpoint{4.122150in}{0.660000in}}%
\pgfpathclose%
\pgfusepath{fill}%
\end{pgfscope}%
\begin{pgfscope}%
\pgfpathrectangle{\pgfqpoint{1.250000in}{0.660000in}}{\pgfqpoint{7.750000in}{4.620000in}}%
\pgfusepath{clip}%
\pgfsetbuttcap%
\pgfsetmiterjoin%
\definecolor{currentfill}{rgb}{0.000000,0.000000,0.000000}%
\pgfsetfillcolor{currentfill}%
\pgfsetlinewidth{0.000000pt}%
\definecolor{currentstroke}{rgb}{0.000000,0.000000,0.000000}%
\pgfsetstrokecolor{currentstroke}%
\pgfsetstrokeopacity{0.000000}%
\pgfsetdash{}{0pt}%
\pgfpathmoveto{\pgfqpoint{4.129900in}{0.660000in}}%
\pgfpathlineto{\pgfqpoint{4.136100in}{0.660000in}}%
\pgfpathlineto{\pgfqpoint{4.136100in}{3.298594in}}%
\pgfpathlineto{\pgfqpoint{4.129900in}{3.298594in}}%
\pgfpathlineto{\pgfqpoint{4.129900in}{0.660000in}}%
\pgfpathclose%
\pgfusepath{fill}%
\end{pgfscope}%
\begin{pgfscope}%
\pgfpathrectangle{\pgfqpoint{1.250000in}{0.660000in}}{\pgfqpoint{7.750000in}{4.620000in}}%
\pgfusepath{clip}%
\pgfsetbuttcap%
\pgfsetmiterjoin%
\definecolor{currentfill}{rgb}{0.000000,0.000000,0.000000}%
\pgfsetfillcolor{currentfill}%
\pgfsetlinewidth{0.000000pt}%
\definecolor{currentstroke}{rgb}{0.000000,0.000000,0.000000}%
\pgfsetstrokecolor{currentstroke}%
\pgfsetstrokeopacity{0.000000}%
\pgfsetdash{}{0pt}%
\pgfpathmoveto{\pgfqpoint{4.137650in}{0.660000in}}%
\pgfpathlineto{\pgfqpoint{4.143850in}{0.660000in}}%
\pgfpathlineto{\pgfqpoint{4.143850in}{3.302974in}}%
\pgfpathlineto{\pgfqpoint{4.137650in}{3.302974in}}%
\pgfpathlineto{\pgfqpoint{4.137650in}{0.660000in}}%
\pgfpathclose%
\pgfusepath{fill}%
\end{pgfscope}%
\begin{pgfscope}%
\pgfpathrectangle{\pgfqpoint{1.250000in}{0.660000in}}{\pgfqpoint{7.750000in}{4.620000in}}%
\pgfusepath{clip}%
\pgfsetbuttcap%
\pgfsetmiterjoin%
\definecolor{currentfill}{rgb}{0.000000,0.000000,0.000000}%
\pgfsetfillcolor{currentfill}%
\pgfsetlinewidth{0.000000pt}%
\definecolor{currentstroke}{rgb}{0.000000,0.000000,0.000000}%
\pgfsetstrokecolor{currentstroke}%
\pgfsetstrokeopacity{0.000000}%
\pgfsetdash{}{0pt}%
\pgfpathmoveto{\pgfqpoint{4.145400in}{0.660000in}}%
\pgfpathlineto{\pgfqpoint{4.151600in}{0.660000in}}%
\pgfpathlineto{\pgfqpoint{4.151600in}{3.339305in}}%
\pgfpathlineto{\pgfqpoint{4.145400in}{3.339305in}}%
\pgfpathlineto{\pgfqpoint{4.145400in}{0.660000in}}%
\pgfpathclose%
\pgfusepath{fill}%
\end{pgfscope}%
\begin{pgfscope}%
\pgfpathrectangle{\pgfqpoint{1.250000in}{0.660000in}}{\pgfqpoint{7.750000in}{4.620000in}}%
\pgfusepath{clip}%
\pgfsetbuttcap%
\pgfsetmiterjoin%
\definecolor{currentfill}{rgb}{0.000000,0.000000,0.000000}%
\pgfsetfillcolor{currentfill}%
\pgfsetlinewidth{0.000000pt}%
\definecolor{currentstroke}{rgb}{0.000000,0.000000,0.000000}%
\pgfsetstrokecolor{currentstroke}%
\pgfsetstrokeopacity{0.000000}%
\pgfsetdash{}{0pt}%
\pgfpathmoveto{\pgfqpoint{4.153150in}{0.660000in}}%
\pgfpathlineto{\pgfqpoint{4.159350in}{0.660000in}}%
\pgfpathlineto{\pgfqpoint{4.159350in}{3.329497in}}%
\pgfpathlineto{\pgfqpoint{4.153150in}{3.329497in}}%
\pgfpathlineto{\pgfqpoint{4.153150in}{0.660000in}}%
\pgfpathclose%
\pgfusepath{fill}%
\end{pgfscope}%
\begin{pgfscope}%
\pgfpathrectangle{\pgfqpoint{1.250000in}{0.660000in}}{\pgfqpoint{7.750000in}{4.620000in}}%
\pgfusepath{clip}%
\pgfsetbuttcap%
\pgfsetmiterjoin%
\definecolor{currentfill}{rgb}{0.000000,0.000000,0.000000}%
\pgfsetfillcolor{currentfill}%
\pgfsetlinewidth{0.000000pt}%
\definecolor{currentstroke}{rgb}{0.000000,0.000000,0.000000}%
\pgfsetstrokecolor{currentstroke}%
\pgfsetstrokeopacity{0.000000}%
\pgfsetdash{}{0pt}%
\pgfpathmoveto{\pgfqpoint{4.160900in}{0.660000in}}%
\pgfpathlineto{\pgfqpoint{4.167100in}{0.660000in}}%
\pgfpathlineto{\pgfqpoint{4.167100in}{3.342495in}}%
\pgfpathlineto{\pgfqpoint{4.160900in}{3.342495in}}%
\pgfpathlineto{\pgfqpoint{4.160900in}{0.660000in}}%
\pgfpathclose%
\pgfusepath{fill}%
\end{pgfscope}%
\begin{pgfscope}%
\pgfpathrectangle{\pgfqpoint{1.250000in}{0.660000in}}{\pgfqpoint{7.750000in}{4.620000in}}%
\pgfusepath{clip}%
\pgfsetbuttcap%
\pgfsetmiterjoin%
\definecolor{currentfill}{rgb}{0.000000,0.000000,0.000000}%
\pgfsetfillcolor{currentfill}%
\pgfsetlinewidth{0.000000pt}%
\definecolor{currentstroke}{rgb}{0.000000,0.000000,0.000000}%
\pgfsetstrokecolor{currentstroke}%
\pgfsetstrokeopacity{0.000000}%
\pgfsetdash{}{0pt}%
\pgfpathmoveto{\pgfqpoint{4.168650in}{0.660000in}}%
\pgfpathlineto{\pgfqpoint{4.174850in}{0.660000in}}%
\pgfpathlineto{\pgfqpoint{4.174850in}{3.289742in}}%
\pgfpathlineto{\pgfqpoint{4.168650in}{3.289742in}}%
\pgfpathlineto{\pgfqpoint{4.168650in}{0.660000in}}%
\pgfpathclose%
\pgfusepath{fill}%
\end{pgfscope}%
\begin{pgfscope}%
\pgfpathrectangle{\pgfqpoint{1.250000in}{0.660000in}}{\pgfqpoint{7.750000in}{4.620000in}}%
\pgfusepath{clip}%
\pgfsetbuttcap%
\pgfsetmiterjoin%
\definecolor{currentfill}{rgb}{0.000000,0.000000,0.000000}%
\pgfsetfillcolor{currentfill}%
\pgfsetlinewidth{0.000000pt}%
\definecolor{currentstroke}{rgb}{0.000000,0.000000,0.000000}%
\pgfsetstrokecolor{currentstroke}%
\pgfsetstrokeopacity{0.000000}%
\pgfsetdash{}{0pt}%
\pgfpathmoveto{\pgfqpoint{4.176400in}{0.660000in}}%
\pgfpathlineto{\pgfqpoint{4.182600in}{0.660000in}}%
\pgfpathlineto{\pgfqpoint{4.182600in}{3.305226in}}%
\pgfpathlineto{\pgfqpoint{4.176400in}{3.305226in}}%
\pgfpathlineto{\pgfqpoint{4.176400in}{0.660000in}}%
\pgfpathclose%
\pgfusepath{fill}%
\end{pgfscope}%
\begin{pgfscope}%
\pgfpathrectangle{\pgfqpoint{1.250000in}{0.660000in}}{\pgfqpoint{7.750000in}{4.620000in}}%
\pgfusepath{clip}%
\pgfsetbuttcap%
\pgfsetmiterjoin%
\definecolor{currentfill}{rgb}{0.000000,0.000000,0.000000}%
\pgfsetfillcolor{currentfill}%
\pgfsetlinewidth{0.000000pt}%
\definecolor{currentstroke}{rgb}{0.000000,0.000000,0.000000}%
\pgfsetstrokecolor{currentstroke}%
\pgfsetstrokeopacity{0.000000}%
\pgfsetdash{}{0pt}%
\pgfpathmoveto{\pgfqpoint{4.184150in}{0.660000in}}%
\pgfpathlineto{\pgfqpoint{4.190350in}{0.660000in}}%
\pgfpathlineto{\pgfqpoint{4.190350in}{3.308575in}}%
\pgfpathlineto{\pgfqpoint{4.184150in}{3.308575in}}%
\pgfpathlineto{\pgfqpoint{4.184150in}{0.660000in}}%
\pgfpathclose%
\pgfusepath{fill}%
\end{pgfscope}%
\begin{pgfscope}%
\pgfpathrectangle{\pgfqpoint{1.250000in}{0.660000in}}{\pgfqpoint{7.750000in}{4.620000in}}%
\pgfusepath{clip}%
\pgfsetbuttcap%
\pgfsetmiterjoin%
\definecolor{currentfill}{rgb}{0.000000,0.000000,0.000000}%
\pgfsetfillcolor{currentfill}%
\pgfsetlinewidth{0.000000pt}%
\definecolor{currentstroke}{rgb}{0.000000,0.000000,0.000000}%
\pgfsetstrokecolor{currentstroke}%
\pgfsetstrokeopacity{0.000000}%
\pgfsetdash{}{0pt}%
\pgfpathmoveto{\pgfqpoint{4.191900in}{0.660000in}}%
\pgfpathlineto{\pgfqpoint{4.198100in}{0.660000in}}%
\pgfpathlineto{\pgfqpoint{4.198100in}{3.311292in}}%
\pgfpathlineto{\pgfqpoint{4.191900in}{3.311292in}}%
\pgfpathlineto{\pgfqpoint{4.191900in}{0.660000in}}%
\pgfpathclose%
\pgfusepath{fill}%
\end{pgfscope}%
\begin{pgfscope}%
\pgfpathrectangle{\pgfqpoint{1.250000in}{0.660000in}}{\pgfqpoint{7.750000in}{4.620000in}}%
\pgfusepath{clip}%
\pgfsetbuttcap%
\pgfsetmiterjoin%
\definecolor{currentfill}{rgb}{0.000000,0.000000,0.000000}%
\pgfsetfillcolor{currentfill}%
\pgfsetlinewidth{0.000000pt}%
\definecolor{currentstroke}{rgb}{0.000000,0.000000,0.000000}%
\pgfsetstrokecolor{currentstroke}%
\pgfsetstrokeopacity{0.000000}%
\pgfsetdash{}{0pt}%
\pgfpathmoveto{\pgfqpoint{4.199650in}{0.660000in}}%
\pgfpathlineto{\pgfqpoint{4.205850in}{0.660000in}}%
\pgfpathlineto{\pgfqpoint{4.205850in}{3.286811in}}%
\pgfpathlineto{\pgfqpoint{4.199650in}{3.286811in}}%
\pgfpathlineto{\pgfqpoint{4.199650in}{0.660000in}}%
\pgfpathclose%
\pgfusepath{fill}%
\end{pgfscope}%
\begin{pgfscope}%
\pgfpathrectangle{\pgfqpoint{1.250000in}{0.660000in}}{\pgfqpoint{7.750000in}{4.620000in}}%
\pgfusepath{clip}%
\pgfsetbuttcap%
\pgfsetmiterjoin%
\definecolor{currentfill}{rgb}{0.000000,0.000000,0.000000}%
\pgfsetfillcolor{currentfill}%
\pgfsetlinewidth{0.000000pt}%
\definecolor{currentstroke}{rgb}{0.000000,0.000000,0.000000}%
\pgfsetstrokecolor{currentstroke}%
\pgfsetstrokeopacity{0.000000}%
\pgfsetdash{}{0pt}%
\pgfpathmoveto{\pgfqpoint{4.207400in}{0.660000in}}%
\pgfpathlineto{\pgfqpoint{4.213600in}{0.660000in}}%
\pgfpathlineto{\pgfqpoint{4.213600in}{3.334920in}}%
\pgfpathlineto{\pgfqpoint{4.207400in}{3.334920in}}%
\pgfpathlineto{\pgfqpoint{4.207400in}{0.660000in}}%
\pgfpathclose%
\pgfusepath{fill}%
\end{pgfscope}%
\begin{pgfscope}%
\pgfpathrectangle{\pgfqpoint{1.250000in}{0.660000in}}{\pgfqpoint{7.750000in}{4.620000in}}%
\pgfusepath{clip}%
\pgfsetbuttcap%
\pgfsetmiterjoin%
\definecolor{currentfill}{rgb}{0.000000,0.000000,0.000000}%
\pgfsetfillcolor{currentfill}%
\pgfsetlinewidth{0.000000pt}%
\definecolor{currentstroke}{rgb}{0.000000,0.000000,0.000000}%
\pgfsetstrokecolor{currentstroke}%
\pgfsetstrokeopacity{0.000000}%
\pgfsetdash{}{0pt}%
\pgfpathmoveto{\pgfqpoint{4.215150in}{0.660000in}}%
\pgfpathlineto{\pgfqpoint{4.221350in}{0.660000in}}%
\pgfpathlineto{\pgfqpoint{4.221350in}{3.307746in}}%
\pgfpathlineto{\pgfqpoint{4.215150in}{3.307746in}}%
\pgfpathlineto{\pgfqpoint{4.215150in}{0.660000in}}%
\pgfpathclose%
\pgfusepath{fill}%
\end{pgfscope}%
\begin{pgfscope}%
\pgfpathrectangle{\pgfqpoint{1.250000in}{0.660000in}}{\pgfqpoint{7.750000in}{4.620000in}}%
\pgfusepath{clip}%
\pgfsetbuttcap%
\pgfsetmiterjoin%
\definecolor{currentfill}{rgb}{0.000000,0.000000,0.000000}%
\pgfsetfillcolor{currentfill}%
\pgfsetlinewidth{0.000000pt}%
\definecolor{currentstroke}{rgb}{0.000000,0.000000,0.000000}%
\pgfsetstrokecolor{currentstroke}%
\pgfsetstrokeopacity{0.000000}%
\pgfsetdash{}{0pt}%
\pgfpathmoveto{\pgfqpoint{4.222900in}{0.660000in}}%
\pgfpathlineto{\pgfqpoint{4.229100in}{0.660000in}}%
\pgfpathlineto{\pgfqpoint{4.229100in}{3.301680in}}%
\pgfpathlineto{\pgfqpoint{4.222900in}{3.301680in}}%
\pgfpathlineto{\pgfqpoint{4.222900in}{0.660000in}}%
\pgfpathclose%
\pgfusepath{fill}%
\end{pgfscope}%
\begin{pgfscope}%
\pgfpathrectangle{\pgfqpoint{1.250000in}{0.660000in}}{\pgfqpoint{7.750000in}{4.620000in}}%
\pgfusepath{clip}%
\pgfsetbuttcap%
\pgfsetmiterjoin%
\definecolor{currentfill}{rgb}{0.000000,0.000000,0.000000}%
\pgfsetfillcolor{currentfill}%
\pgfsetlinewidth{0.000000pt}%
\definecolor{currentstroke}{rgb}{0.000000,0.000000,0.000000}%
\pgfsetstrokecolor{currentstroke}%
\pgfsetstrokeopacity{0.000000}%
\pgfsetdash{}{0pt}%
\pgfpathmoveto{\pgfqpoint{4.230650in}{0.660000in}}%
\pgfpathlineto{\pgfqpoint{4.236850in}{0.660000in}}%
\pgfpathlineto{\pgfqpoint{4.236850in}{3.311977in}}%
\pgfpathlineto{\pgfqpoint{4.230650in}{3.311977in}}%
\pgfpathlineto{\pgfqpoint{4.230650in}{0.660000in}}%
\pgfpathclose%
\pgfusepath{fill}%
\end{pgfscope}%
\begin{pgfscope}%
\pgfpathrectangle{\pgfqpoint{1.250000in}{0.660000in}}{\pgfqpoint{7.750000in}{4.620000in}}%
\pgfusepath{clip}%
\pgfsetbuttcap%
\pgfsetmiterjoin%
\definecolor{currentfill}{rgb}{0.000000,0.000000,0.000000}%
\pgfsetfillcolor{currentfill}%
\pgfsetlinewidth{0.000000pt}%
\definecolor{currentstroke}{rgb}{0.000000,0.000000,0.000000}%
\pgfsetstrokecolor{currentstroke}%
\pgfsetstrokeopacity{0.000000}%
\pgfsetdash{}{0pt}%
\pgfpathmoveto{\pgfqpoint{4.238400in}{0.660000in}}%
\pgfpathlineto{\pgfqpoint{4.244600in}{0.660000in}}%
\pgfpathlineto{\pgfqpoint{4.244600in}{3.334906in}}%
\pgfpathlineto{\pgfqpoint{4.238400in}{3.334906in}}%
\pgfpathlineto{\pgfqpoint{4.238400in}{0.660000in}}%
\pgfpathclose%
\pgfusepath{fill}%
\end{pgfscope}%
\begin{pgfscope}%
\pgfpathrectangle{\pgfqpoint{1.250000in}{0.660000in}}{\pgfqpoint{7.750000in}{4.620000in}}%
\pgfusepath{clip}%
\pgfsetbuttcap%
\pgfsetmiterjoin%
\definecolor{currentfill}{rgb}{0.000000,0.000000,0.000000}%
\pgfsetfillcolor{currentfill}%
\pgfsetlinewidth{0.000000pt}%
\definecolor{currentstroke}{rgb}{0.000000,0.000000,0.000000}%
\pgfsetstrokecolor{currentstroke}%
\pgfsetstrokeopacity{0.000000}%
\pgfsetdash{}{0pt}%
\pgfpathmoveto{\pgfqpoint{4.246150in}{0.660000in}}%
\pgfpathlineto{\pgfqpoint{4.252350in}{0.660000in}}%
\pgfpathlineto{\pgfqpoint{4.252350in}{3.280577in}}%
\pgfpathlineto{\pgfqpoint{4.246150in}{3.280577in}}%
\pgfpathlineto{\pgfqpoint{4.246150in}{0.660000in}}%
\pgfpathclose%
\pgfusepath{fill}%
\end{pgfscope}%
\begin{pgfscope}%
\pgfpathrectangle{\pgfqpoint{1.250000in}{0.660000in}}{\pgfqpoint{7.750000in}{4.620000in}}%
\pgfusepath{clip}%
\pgfsetbuttcap%
\pgfsetmiterjoin%
\definecolor{currentfill}{rgb}{0.000000,0.000000,0.000000}%
\pgfsetfillcolor{currentfill}%
\pgfsetlinewidth{0.000000pt}%
\definecolor{currentstroke}{rgb}{0.000000,0.000000,0.000000}%
\pgfsetstrokecolor{currentstroke}%
\pgfsetstrokeopacity{0.000000}%
\pgfsetdash{}{0pt}%
\pgfpathmoveto{\pgfqpoint{4.253900in}{0.660000in}}%
\pgfpathlineto{\pgfqpoint{4.260100in}{0.660000in}}%
\pgfpathlineto{\pgfqpoint{4.260100in}{3.280886in}}%
\pgfpathlineto{\pgfqpoint{4.253900in}{3.280886in}}%
\pgfpathlineto{\pgfqpoint{4.253900in}{0.660000in}}%
\pgfpathclose%
\pgfusepath{fill}%
\end{pgfscope}%
\begin{pgfscope}%
\pgfpathrectangle{\pgfqpoint{1.250000in}{0.660000in}}{\pgfqpoint{7.750000in}{4.620000in}}%
\pgfusepath{clip}%
\pgfsetbuttcap%
\pgfsetmiterjoin%
\definecolor{currentfill}{rgb}{0.000000,0.000000,0.000000}%
\pgfsetfillcolor{currentfill}%
\pgfsetlinewidth{0.000000pt}%
\definecolor{currentstroke}{rgb}{0.000000,0.000000,0.000000}%
\pgfsetstrokecolor{currentstroke}%
\pgfsetstrokeopacity{0.000000}%
\pgfsetdash{}{0pt}%
\pgfpathmoveto{\pgfqpoint{4.261650in}{0.660000in}}%
\pgfpathlineto{\pgfqpoint{4.267850in}{0.660000in}}%
\pgfpathlineto{\pgfqpoint{4.267850in}{3.290913in}}%
\pgfpathlineto{\pgfqpoint{4.261650in}{3.290913in}}%
\pgfpathlineto{\pgfqpoint{4.261650in}{0.660000in}}%
\pgfpathclose%
\pgfusepath{fill}%
\end{pgfscope}%
\begin{pgfscope}%
\pgfpathrectangle{\pgfqpoint{1.250000in}{0.660000in}}{\pgfqpoint{7.750000in}{4.620000in}}%
\pgfusepath{clip}%
\pgfsetbuttcap%
\pgfsetmiterjoin%
\definecolor{currentfill}{rgb}{0.000000,0.000000,0.000000}%
\pgfsetfillcolor{currentfill}%
\pgfsetlinewidth{0.000000pt}%
\definecolor{currentstroke}{rgb}{0.000000,0.000000,0.000000}%
\pgfsetstrokecolor{currentstroke}%
\pgfsetstrokeopacity{0.000000}%
\pgfsetdash{}{0pt}%
\pgfpathmoveto{\pgfqpoint{4.269400in}{0.660000in}}%
\pgfpathlineto{\pgfqpoint{4.275600in}{0.660000in}}%
\pgfpathlineto{\pgfqpoint{4.275600in}{3.309323in}}%
\pgfpathlineto{\pgfqpoint{4.269400in}{3.309323in}}%
\pgfpathlineto{\pgfqpoint{4.269400in}{0.660000in}}%
\pgfpathclose%
\pgfusepath{fill}%
\end{pgfscope}%
\begin{pgfscope}%
\pgfpathrectangle{\pgfqpoint{1.250000in}{0.660000in}}{\pgfqpoint{7.750000in}{4.620000in}}%
\pgfusepath{clip}%
\pgfsetbuttcap%
\pgfsetmiterjoin%
\definecolor{currentfill}{rgb}{0.000000,0.000000,0.000000}%
\pgfsetfillcolor{currentfill}%
\pgfsetlinewidth{0.000000pt}%
\definecolor{currentstroke}{rgb}{0.000000,0.000000,0.000000}%
\pgfsetstrokecolor{currentstroke}%
\pgfsetstrokeopacity{0.000000}%
\pgfsetdash{}{0pt}%
\pgfpathmoveto{\pgfqpoint{4.277150in}{0.660000in}}%
\pgfpathlineto{\pgfqpoint{4.283350in}{0.660000in}}%
\pgfpathlineto{\pgfqpoint{4.283350in}{3.289876in}}%
\pgfpathlineto{\pgfqpoint{4.277150in}{3.289876in}}%
\pgfpathlineto{\pgfqpoint{4.277150in}{0.660000in}}%
\pgfpathclose%
\pgfusepath{fill}%
\end{pgfscope}%
\begin{pgfscope}%
\pgfpathrectangle{\pgfqpoint{1.250000in}{0.660000in}}{\pgfqpoint{7.750000in}{4.620000in}}%
\pgfusepath{clip}%
\pgfsetbuttcap%
\pgfsetmiterjoin%
\definecolor{currentfill}{rgb}{0.000000,0.000000,0.000000}%
\pgfsetfillcolor{currentfill}%
\pgfsetlinewidth{0.000000pt}%
\definecolor{currentstroke}{rgb}{0.000000,0.000000,0.000000}%
\pgfsetstrokecolor{currentstroke}%
\pgfsetstrokeopacity{0.000000}%
\pgfsetdash{}{0pt}%
\pgfpathmoveto{\pgfqpoint{4.284900in}{0.660000in}}%
\pgfpathlineto{\pgfqpoint{4.291100in}{0.660000in}}%
\pgfpathlineto{\pgfqpoint{4.291100in}{3.287668in}}%
\pgfpathlineto{\pgfqpoint{4.284900in}{3.287668in}}%
\pgfpathlineto{\pgfqpoint{4.284900in}{0.660000in}}%
\pgfpathclose%
\pgfusepath{fill}%
\end{pgfscope}%
\begin{pgfscope}%
\pgfpathrectangle{\pgfqpoint{1.250000in}{0.660000in}}{\pgfqpoint{7.750000in}{4.620000in}}%
\pgfusepath{clip}%
\pgfsetbuttcap%
\pgfsetmiterjoin%
\definecolor{currentfill}{rgb}{0.000000,0.000000,0.000000}%
\pgfsetfillcolor{currentfill}%
\pgfsetlinewidth{0.000000pt}%
\definecolor{currentstroke}{rgb}{0.000000,0.000000,0.000000}%
\pgfsetstrokecolor{currentstroke}%
\pgfsetstrokeopacity{0.000000}%
\pgfsetdash{}{0pt}%
\pgfpathmoveto{\pgfqpoint{4.292650in}{0.660000in}}%
\pgfpathlineto{\pgfqpoint{4.298850in}{0.660000in}}%
\pgfpathlineto{\pgfqpoint{4.298850in}{3.345586in}}%
\pgfpathlineto{\pgfqpoint{4.292650in}{3.345586in}}%
\pgfpathlineto{\pgfqpoint{4.292650in}{0.660000in}}%
\pgfpathclose%
\pgfusepath{fill}%
\end{pgfscope}%
\begin{pgfscope}%
\pgfpathrectangle{\pgfqpoint{1.250000in}{0.660000in}}{\pgfqpoint{7.750000in}{4.620000in}}%
\pgfusepath{clip}%
\pgfsetbuttcap%
\pgfsetmiterjoin%
\definecolor{currentfill}{rgb}{0.000000,0.000000,0.000000}%
\pgfsetfillcolor{currentfill}%
\pgfsetlinewidth{0.000000pt}%
\definecolor{currentstroke}{rgb}{0.000000,0.000000,0.000000}%
\pgfsetstrokecolor{currentstroke}%
\pgfsetstrokeopacity{0.000000}%
\pgfsetdash{}{0pt}%
\pgfpathmoveto{\pgfqpoint{4.300400in}{0.660000in}}%
\pgfpathlineto{\pgfqpoint{4.306600in}{0.660000in}}%
\pgfpathlineto{\pgfqpoint{4.306600in}{3.300431in}}%
\pgfpathlineto{\pgfqpoint{4.300400in}{3.300431in}}%
\pgfpathlineto{\pgfqpoint{4.300400in}{0.660000in}}%
\pgfpathclose%
\pgfusepath{fill}%
\end{pgfscope}%
\begin{pgfscope}%
\pgfpathrectangle{\pgfqpoint{1.250000in}{0.660000in}}{\pgfqpoint{7.750000in}{4.620000in}}%
\pgfusepath{clip}%
\pgfsetbuttcap%
\pgfsetmiterjoin%
\definecolor{currentfill}{rgb}{0.000000,0.000000,0.000000}%
\pgfsetfillcolor{currentfill}%
\pgfsetlinewidth{0.000000pt}%
\definecolor{currentstroke}{rgb}{0.000000,0.000000,0.000000}%
\pgfsetstrokecolor{currentstroke}%
\pgfsetstrokeopacity{0.000000}%
\pgfsetdash{}{0pt}%
\pgfpathmoveto{\pgfqpoint{4.308150in}{0.660000in}}%
\pgfpathlineto{\pgfqpoint{4.314350in}{0.660000in}}%
\pgfpathlineto{\pgfqpoint{4.314350in}{3.291417in}}%
\pgfpathlineto{\pgfqpoint{4.308150in}{3.291417in}}%
\pgfpathlineto{\pgfqpoint{4.308150in}{0.660000in}}%
\pgfpathclose%
\pgfusepath{fill}%
\end{pgfscope}%
\begin{pgfscope}%
\pgfpathrectangle{\pgfqpoint{1.250000in}{0.660000in}}{\pgfqpoint{7.750000in}{4.620000in}}%
\pgfusepath{clip}%
\pgfsetbuttcap%
\pgfsetmiterjoin%
\definecolor{currentfill}{rgb}{0.000000,0.000000,0.000000}%
\pgfsetfillcolor{currentfill}%
\pgfsetlinewidth{0.000000pt}%
\definecolor{currentstroke}{rgb}{0.000000,0.000000,0.000000}%
\pgfsetstrokecolor{currentstroke}%
\pgfsetstrokeopacity{0.000000}%
\pgfsetdash{}{0pt}%
\pgfpathmoveto{\pgfqpoint{4.315900in}{0.660000in}}%
\pgfpathlineto{\pgfqpoint{4.322100in}{0.660000in}}%
\pgfpathlineto{\pgfqpoint{4.322100in}{3.327761in}}%
\pgfpathlineto{\pgfqpoint{4.315900in}{3.327761in}}%
\pgfpathlineto{\pgfqpoint{4.315900in}{0.660000in}}%
\pgfpathclose%
\pgfusepath{fill}%
\end{pgfscope}%
\begin{pgfscope}%
\pgfpathrectangle{\pgfqpoint{1.250000in}{0.660000in}}{\pgfqpoint{7.750000in}{4.620000in}}%
\pgfusepath{clip}%
\pgfsetbuttcap%
\pgfsetmiterjoin%
\definecolor{currentfill}{rgb}{0.000000,0.000000,0.000000}%
\pgfsetfillcolor{currentfill}%
\pgfsetlinewidth{0.000000pt}%
\definecolor{currentstroke}{rgb}{0.000000,0.000000,0.000000}%
\pgfsetstrokecolor{currentstroke}%
\pgfsetstrokeopacity{0.000000}%
\pgfsetdash{}{0pt}%
\pgfpathmoveto{\pgfqpoint{4.323650in}{0.660000in}}%
\pgfpathlineto{\pgfqpoint{4.329850in}{0.660000in}}%
\pgfpathlineto{\pgfqpoint{4.329850in}{3.364454in}}%
\pgfpathlineto{\pgfqpoint{4.323650in}{3.364454in}}%
\pgfpathlineto{\pgfqpoint{4.323650in}{0.660000in}}%
\pgfpathclose%
\pgfusepath{fill}%
\end{pgfscope}%
\begin{pgfscope}%
\pgfpathrectangle{\pgfqpoint{1.250000in}{0.660000in}}{\pgfqpoint{7.750000in}{4.620000in}}%
\pgfusepath{clip}%
\pgfsetbuttcap%
\pgfsetmiterjoin%
\definecolor{currentfill}{rgb}{0.000000,0.000000,0.000000}%
\pgfsetfillcolor{currentfill}%
\pgfsetlinewidth{0.000000pt}%
\definecolor{currentstroke}{rgb}{0.000000,0.000000,0.000000}%
\pgfsetstrokecolor{currentstroke}%
\pgfsetstrokeopacity{0.000000}%
\pgfsetdash{}{0pt}%
\pgfpathmoveto{\pgfqpoint{4.331400in}{0.660000in}}%
\pgfpathlineto{\pgfqpoint{4.337600in}{0.660000in}}%
\pgfpathlineto{\pgfqpoint{4.337600in}{3.306794in}}%
\pgfpathlineto{\pgfqpoint{4.331400in}{3.306794in}}%
\pgfpathlineto{\pgfqpoint{4.331400in}{0.660000in}}%
\pgfpathclose%
\pgfusepath{fill}%
\end{pgfscope}%
\begin{pgfscope}%
\pgfpathrectangle{\pgfqpoint{1.250000in}{0.660000in}}{\pgfqpoint{7.750000in}{4.620000in}}%
\pgfusepath{clip}%
\pgfsetbuttcap%
\pgfsetmiterjoin%
\definecolor{currentfill}{rgb}{0.000000,0.000000,0.000000}%
\pgfsetfillcolor{currentfill}%
\pgfsetlinewidth{0.000000pt}%
\definecolor{currentstroke}{rgb}{0.000000,0.000000,0.000000}%
\pgfsetstrokecolor{currentstroke}%
\pgfsetstrokeopacity{0.000000}%
\pgfsetdash{}{0pt}%
\pgfpathmoveto{\pgfqpoint{4.339150in}{0.660000in}}%
\pgfpathlineto{\pgfqpoint{4.345350in}{0.660000in}}%
\pgfpathlineto{\pgfqpoint{4.345350in}{3.294313in}}%
\pgfpathlineto{\pgfqpoint{4.339150in}{3.294313in}}%
\pgfpathlineto{\pgfqpoint{4.339150in}{0.660000in}}%
\pgfpathclose%
\pgfusepath{fill}%
\end{pgfscope}%
\begin{pgfscope}%
\pgfpathrectangle{\pgfqpoint{1.250000in}{0.660000in}}{\pgfqpoint{7.750000in}{4.620000in}}%
\pgfusepath{clip}%
\pgfsetbuttcap%
\pgfsetmiterjoin%
\definecolor{currentfill}{rgb}{0.000000,0.000000,0.000000}%
\pgfsetfillcolor{currentfill}%
\pgfsetlinewidth{0.000000pt}%
\definecolor{currentstroke}{rgb}{0.000000,0.000000,0.000000}%
\pgfsetstrokecolor{currentstroke}%
\pgfsetstrokeopacity{0.000000}%
\pgfsetdash{}{0pt}%
\pgfpathmoveto{\pgfqpoint{4.346900in}{0.660000in}}%
\pgfpathlineto{\pgfqpoint{4.353100in}{0.660000in}}%
\pgfpathlineto{\pgfqpoint{4.353100in}{3.302203in}}%
\pgfpathlineto{\pgfqpoint{4.346900in}{3.302203in}}%
\pgfpathlineto{\pgfqpoint{4.346900in}{0.660000in}}%
\pgfpathclose%
\pgfusepath{fill}%
\end{pgfscope}%
\begin{pgfscope}%
\pgfpathrectangle{\pgfqpoint{1.250000in}{0.660000in}}{\pgfqpoint{7.750000in}{4.620000in}}%
\pgfusepath{clip}%
\pgfsetbuttcap%
\pgfsetmiterjoin%
\definecolor{currentfill}{rgb}{0.000000,0.000000,0.000000}%
\pgfsetfillcolor{currentfill}%
\pgfsetlinewidth{0.000000pt}%
\definecolor{currentstroke}{rgb}{0.000000,0.000000,0.000000}%
\pgfsetstrokecolor{currentstroke}%
\pgfsetstrokeopacity{0.000000}%
\pgfsetdash{}{0pt}%
\pgfpathmoveto{\pgfqpoint{4.354650in}{0.660000in}}%
\pgfpathlineto{\pgfqpoint{4.360850in}{0.660000in}}%
\pgfpathlineto{\pgfqpoint{4.360850in}{3.333188in}}%
\pgfpathlineto{\pgfqpoint{4.354650in}{3.333188in}}%
\pgfpathlineto{\pgfqpoint{4.354650in}{0.660000in}}%
\pgfpathclose%
\pgfusepath{fill}%
\end{pgfscope}%
\begin{pgfscope}%
\pgfpathrectangle{\pgfqpoint{1.250000in}{0.660000in}}{\pgfqpoint{7.750000in}{4.620000in}}%
\pgfusepath{clip}%
\pgfsetbuttcap%
\pgfsetmiterjoin%
\definecolor{currentfill}{rgb}{0.000000,0.000000,0.000000}%
\pgfsetfillcolor{currentfill}%
\pgfsetlinewidth{0.000000pt}%
\definecolor{currentstroke}{rgb}{0.000000,0.000000,0.000000}%
\pgfsetstrokecolor{currentstroke}%
\pgfsetstrokeopacity{0.000000}%
\pgfsetdash{}{0pt}%
\pgfpathmoveto{\pgfqpoint{4.362400in}{0.660000in}}%
\pgfpathlineto{\pgfqpoint{4.368600in}{0.660000in}}%
\pgfpathlineto{\pgfqpoint{4.368600in}{3.294800in}}%
\pgfpathlineto{\pgfqpoint{4.362400in}{3.294800in}}%
\pgfpathlineto{\pgfqpoint{4.362400in}{0.660000in}}%
\pgfpathclose%
\pgfusepath{fill}%
\end{pgfscope}%
\begin{pgfscope}%
\pgfpathrectangle{\pgfqpoint{1.250000in}{0.660000in}}{\pgfqpoint{7.750000in}{4.620000in}}%
\pgfusepath{clip}%
\pgfsetbuttcap%
\pgfsetmiterjoin%
\definecolor{currentfill}{rgb}{0.000000,0.000000,0.000000}%
\pgfsetfillcolor{currentfill}%
\pgfsetlinewidth{0.000000pt}%
\definecolor{currentstroke}{rgb}{0.000000,0.000000,0.000000}%
\pgfsetstrokecolor{currentstroke}%
\pgfsetstrokeopacity{0.000000}%
\pgfsetdash{}{0pt}%
\pgfpathmoveto{\pgfqpoint{4.370150in}{0.660000in}}%
\pgfpathlineto{\pgfqpoint{4.376350in}{0.660000in}}%
\pgfpathlineto{\pgfqpoint{4.376350in}{3.290431in}}%
\pgfpathlineto{\pgfqpoint{4.370150in}{3.290431in}}%
\pgfpathlineto{\pgfqpoint{4.370150in}{0.660000in}}%
\pgfpathclose%
\pgfusepath{fill}%
\end{pgfscope}%
\begin{pgfscope}%
\pgfpathrectangle{\pgfqpoint{1.250000in}{0.660000in}}{\pgfqpoint{7.750000in}{4.620000in}}%
\pgfusepath{clip}%
\pgfsetbuttcap%
\pgfsetmiterjoin%
\definecolor{currentfill}{rgb}{0.000000,0.000000,0.000000}%
\pgfsetfillcolor{currentfill}%
\pgfsetlinewidth{0.000000pt}%
\definecolor{currentstroke}{rgb}{0.000000,0.000000,0.000000}%
\pgfsetstrokecolor{currentstroke}%
\pgfsetstrokeopacity{0.000000}%
\pgfsetdash{}{0pt}%
\pgfpathmoveto{\pgfqpoint{4.377900in}{0.660000in}}%
\pgfpathlineto{\pgfqpoint{4.384100in}{0.660000in}}%
\pgfpathlineto{\pgfqpoint{4.384100in}{3.289147in}}%
\pgfpathlineto{\pgfqpoint{4.377900in}{3.289147in}}%
\pgfpathlineto{\pgfqpoint{4.377900in}{0.660000in}}%
\pgfpathclose%
\pgfusepath{fill}%
\end{pgfscope}%
\begin{pgfscope}%
\pgfpathrectangle{\pgfqpoint{1.250000in}{0.660000in}}{\pgfqpoint{7.750000in}{4.620000in}}%
\pgfusepath{clip}%
\pgfsetbuttcap%
\pgfsetmiterjoin%
\definecolor{currentfill}{rgb}{0.000000,0.000000,0.000000}%
\pgfsetfillcolor{currentfill}%
\pgfsetlinewidth{0.000000pt}%
\definecolor{currentstroke}{rgb}{0.000000,0.000000,0.000000}%
\pgfsetstrokecolor{currentstroke}%
\pgfsetstrokeopacity{0.000000}%
\pgfsetdash{}{0pt}%
\pgfpathmoveto{\pgfqpoint{4.385650in}{0.660000in}}%
\pgfpathlineto{\pgfqpoint{4.391850in}{0.660000in}}%
\pgfpathlineto{\pgfqpoint{4.391850in}{3.288532in}}%
\pgfpathlineto{\pgfqpoint{4.385650in}{3.288532in}}%
\pgfpathlineto{\pgfqpoint{4.385650in}{0.660000in}}%
\pgfpathclose%
\pgfusepath{fill}%
\end{pgfscope}%
\begin{pgfscope}%
\pgfpathrectangle{\pgfqpoint{1.250000in}{0.660000in}}{\pgfqpoint{7.750000in}{4.620000in}}%
\pgfusepath{clip}%
\pgfsetbuttcap%
\pgfsetmiterjoin%
\definecolor{currentfill}{rgb}{0.000000,0.000000,0.000000}%
\pgfsetfillcolor{currentfill}%
\pgfsetlinewidth{0.000000pt}%
\definecolor{currentstroke}{rgb}{0.000000,0.000000,0.000000}%
\pgfsetstrokecolor{currentstroke}%
\pgfsetstrokeopacity{0.000000}%
\pgfsetdash{}{0pt}%
\pgfpathmoveto{\pgfqpoint{4.393400in}{0.660000in}}%
\pgfpathlineto{\pgfqpoint{4.399600in}{0.660000in}}%
\pgfpathlineto{\pgfqpoint{4.399600in}{3.282799in}}%
\pgfpathlineto{\pgfqpoint{4.393400in}{3.282799in}}%
\pgfpathlineto{\pgfqpoint{4.393400in}{0.660000in}}%
\pgfpathclose%
\pgfusepath{fill}%
\end{pgfscope}%
\begin{pgfscope}%
\pgfpathrectangle{\pgfqpoint{1.250000in}{0.660000in}}{\pgfqpoint{7.750000in}{4.620000in}}%
\pgfusepath{clip}%
\pgfsetbuttcap%
\pgfsetmiterjoin%
\definecolor{currentfill}{rgb}{0.000000,0.000000,0.000000}%
\pgfsetfillcolor{currentfill}%
\pgfsetlinewidth{0.000000pt}%
\definecolor{currentstroke}{rgb}{0.000000,0.000000,0.000000}%
\pgfsetstrokecolor{currentstroke}%
\pgfsetstrokeopacity{0.000000}%
\pgfsetdash{}{0pt}%
\pgfpathmoveto{\pgfqpoint{4.401150in}{0.660000in}}%
\pgfpathlineto{\pgfqpoint{4.407350in}{0.660000in}}%
\pgfpathlineto{\pgfqpoint{4.407350in}{3.284203in}}%
\pgfpathlineto{\pgfqpoint{4.401150in}{3.284203in}}%
\pgfpathlineto{\pgfqpoint{4.401150in}{0.660000in}}%
\pgfpathclose%
\pgfusepath{fill}%
\end{pgfscope}%
\begin{pgfscope}%
\pgfpathrectangle{\pgfqpoint{1.250000in}{0.660000in}}{\pgfqpoint{7.750000in}{4.620000in}}%
\pgfusepath{clip}%
\pgfsetbuttcap%
\pgfsetmiterjoin%
\definecolor{currentfill}{rgb}{0.000000,0.000000,0.000000}%
\pgfsetfillcolor{currentfill}%
\pgfsetlinewidth{0.000000pt}%
\definecolor{currentstroke}{rgb}{0.000000,0.000000,0.000000}%
\pgfsetstrokecolor{currentstroke}%
\pgfsetstrokeopacity{0.000000}%
\pgfsetdash{}{0pt}%
\pgfpathmoveto{\pgfqpoint{4.408900in}{0.660000in}}%
\pgfpathlineto{\pgfqpoint{4.415100in}{0.660000in}}%
\pgfpathlineto{\pgfqpoint{4.415100in}{3.280253in}}%
\pgfpathlineto{\pgfqpoint{4.408900in}{3.280253in}}%
\pgfpathlineto{\pgfqpoint{4.408900in}{0.660000in}}%
\pgfpathclose%
\pgfusepath{fill}%
\end{pgfscope}%
\begin{pgfscope}%
\pgfpathrectangle{\pgfqpoint{1.250000in}{0.660000in}}{\pgfqpoint{7.750000in}{4.620000in}}%
\pgfusepath{clip}%
\pgfsetbuttcap%
\pgfsetmiterjoin%
\definecolor{currentfill}{rgb}{0.000000,0.000000,0.000000}%
\pgfsetfillcolor{currentfill}%
\pgfsetlinewidth{0.000000pt}%
\definecolor{currentstroke}{rgb}{0.000000,0.000000,0.000000}%
\pgfsetstrokecolor{currentstroke}%
\pgfsetstrokeopacity{0.000000}%
\pgfsetdash{}{0pt}%
\pgfpathmoveto{\pgfqpoint{4.416650in}{0.660000in}}%
\pgfpathlineto{\pgfqpoint{4.422850in}{0.660000in}}%
\pgfpathlineto{\pgfqpoint{4.422850in}{3.294800in}}%
\pgfpathlineto{\pgfqpoint{4.416650in}{3.294800in}}%
\pgfpathlineto{\pgfqpoint{4.416650in}{0.660000in}}%
\pgfpathclose%
\pgfusepath{fill}%
\end{pgfscope}%
\begin{pgfscope}%
\pgfpathrectangle{\pgfqpoint{1.250000in}{0.660000in}}{\pgfqpoint{7.750000in}{4.620000in}}%
\pgfusepath{clip}%
\pgfsetbuttcap%
\pgfsetmiterjoin%
\definecolor{currentfill}{rgb}{0.000000,0.000000,0.000000}%
\pgfsetfillcolor{currentfill}%
\pgfsetlinewidth{0.000000pt}%
\definecolor{currentstroke}{rgb}{0.000000,0.000000,0.000000}%
\pgfsetstrokecolor{currentstroke}%
\pgfsetstrokeopacity{0.000000}%
\pgfsetdash{}{0pt}%
\pgfpathmoveto{\pgfqpoint{4.424400in}{0.660000in}}%
\pgfpathlineto{\pgfqpoint{4.430600in}{0.660000in}}%
\pgfpathlineto{\pgfqpoint{4.430600in}{3.274202in}}%
\pgfpathlineto{\pgfqpoint{4.424400in}{3.274202in}}%
\pgfpathlineto{\pgfqpoint{4.424400in}{0.660000in}}%
\pgfpathclose%
\pgfusepath{fill}%
\end{pgfscope}%
\begin{pgfscope}%
\pgfpathrectangle{\pgfqpoint{1.250000in}{0.660000in}}{\pgfqpoint{7.750000in}{4.620000in}}%
\pgfusepath{clip}%
\pgfsetbuttcap%
\pgfsetmiterjoin%
\definecolor{currentfill}{rgb}{0.000000,0.000000,0.000000}%
\pgfsetfillcolor{currentfill}%
\pgfsetlinewidth{0.000000pt}%
\definecolor{currentstroke}{rgb}{0.000000,0.000000,0.000000}%
\pgfsetstrokecolor{currentstroke}%
\pgfsetstrokeopacity{0.000000}%
\pgfsetdash{}{0pt}%
\pgfpathmoveto{\pgfqpoint{4.432150in}{0.660000in}}%
\pgfpathlineto{\pgfqpoint{4.438350in}{0.660000in}}%
\pgfpathlineto{\pgfqpoint{4.438350in}{3.273984in}}%
\pgfpathlineto{\pgfqpoint{4.432150in}{3.273984in}}%
\pgfpathlineto{\pgfqpoint{4.432150in}{0.660000in}}%
\pgfpathclose%
\pgfusepath{fill}%
\end{pgfscope}%
\begin{pgfscope}%
\pgfpathrectangle{\pgfqpoint{1.250000in}{0.660000in}}{\pgfqpoint{7.750000in}{4.620000in}}%
\pgfusepath{clip}%
\pgfsetbuttcap%
\pgfsetmiterjoin%
\definecolor{currentfill}{rgb}{0.000000,0.000000,0.000000}%
\pgfsetfillcolor{currentfill}%
\pgfsetlinewidth{0.000000pt}%
\definecolor{currentstroke}{rgb}{0.000000,0.000000,0.000000}%
\pgfsetstrokecolor{currentstroke}%
\pgfsetstrokeopacity{0.000000}%
\pgfsetdash{}{0pt}%
\pgfpathmoveto{\pgfqpoint{4.439900in}{0.660000in}}%
\pgfpathlineto{\pgfqpoint{4.446100in}{0.660000in}}%
\pgfpathlineto{\pgfqpoint{4.446100in}{3.302447in}}%
\pgfpathlineto{\pgfqpoint{4.439900in}{3.302447in}}%
\pgfpathlineto{\pgfqpoint{4.439900in}{0.660000in}}%
\pgfpathclose%
\pgfusepath{fill}%
\end{pgfscope}%
\begin{pgfscope}%
\pgfpathrectangle{\pgfqpoint{1.250000in}{0.660000in}}{\pgfqpoint{7.750000in}{4.620000in}}%
\pgfusepath{clip}%
\pgfsetbuttcap%
\pgfsetmiterjoin%
\definecolor{currentfill}{rgb}{0.000000,0.000000,0.000000}%
\pgfsetfillcolor{currentfill}%
\pgfsetlinewidth{0.000000pt}%
\definecolor{currentstroke}{rgb}{0.000000,0.000000,0.000000}%
\pgfsetstrokecolor{currentstroke}%
\pgfsetstrokeopacity{0.000000}%
\pgfsetdash{}{0pt}%
\pgfpathmoveto{\pgfqpoint{4.447650in}{0.660000in}}%
\pgfpathlineto{\pgfqpoint{4.453850in}{0.660000in}}%
\pgfpathlineto{\pgfqpoint{4.453850in}{3.294520in}}%
\pgfpathlineto{\pgfqpoint{4.447650in}{3.294520in}}%
\pgfpathlineto{\pgfqpoint{4.447650in}{0.660000in}}%
\pgfpathclose%
\pgfusepath{fill}%
\end{pgfscope}%
\begin{pgfscope}%
\pgfpathrectangle{\pgfqpoint{1.250000in}{0.660000in}}{\pgfqpoint{7.750000in}{4.620000in}}%
\pgfusepath{clip}%
\pgfsetbuttcap%
\pgfsetmiterjoin%
\definecolor{currentfill}{rgb}{0.000000,0.000000,0.000000}%
\pgfsetfillcolor{currentfill}%
\pgfsetlinewidth{0.000000pt}%
\definecolor{currentstroke}{rgb}{0.000000,0.000000,0.000000}%
\pgfsetstrokecolor{currentstroke}%
\pgfsetstrokeopacity{0.000000}%
\pgfsetdash{}{0pt}%
\pgfpathmoveto{\pgfqpoint{4.455400in}{0.660000in}}%
\pgfpathlineto{\pgfqpoint{4.461600in}{0.660000in}}%
\pgfpathlineto{\pgfqpoint{4.461600in}{3.291642in}}%
\pgfpathlineto{\pgfqpoint{4.455400in}{3.291642in}}%
\pgfpathlineto{\pgfqpoint{4.455400in}{0.660000in}}%
\pgfpathclose%
\pgfusepath{fill}%
\end{pgfscope}%
\begin{pgfscope}%
\pgfpathrectangle{\pgfqpoint{1.250000in}{0.660000in}}{\pgfqpoint{7.750000in}{4.620000in}}%
\pgfusepath{clip}%
\pgfsetbuttcap%
\pgfsetmiterjoin%
\definecolor{currentfill}{rgb}{0.000000,0.000000,0.000000}%
\pgfsetfillcolor{currentfill}%
\pgfsetlinewidth{0.000000pt}%
\definecolor{currentstroke}{rgb}{0.000000,0.000000,0.000000}%
\pgfsetstrokecolor{currentstroke}%
\pgfsetstrokeopacity{0.000000}%
\pgfsetdash{}{0pt}%
\pgfpathmoveto{\pgfqpoint{4.463150in}{0.660000in}}%
\pgfpathlineto{\pgfqpoint{4.469350in}{0.660000in}}%
\pgfpathlineto{\pgfqpoint{4.469350in}{3.283010in}}%
\pgfpathlineto{\pgfqpoint{4.463150in}{3.283010in}}%
\pgfpathlineto{\pgfqpoint{4.463150in}{0.660000in}}%
\pgfpathclose%
\pgfusepath{fill}%
\end{pgfscope}%
\begin{pgfscope}%
\pgfpathrectangle{\pgfqpoint{1.250000in}{0.660000in}}{\pgfqpoint{7.750000in}{4.620000in}}%
\pgfusepath{clip}%
\pgfsetbuttcap%
\pgfsetmiterjoin%
\definecolor{currentfill}{rgb}{0.000000,0.000000,0.000000}%
\pgfsetfillcolor{currentfill}%
\pgfsetlinewidth{0.000000pt}%
\definecolor{currentstroke}{rgb}{0.000000,0.000000,0.000000}%
\pgfsetstrokecolor{currentstroke}%
\pgfsetstrokeopacity{0.000000}%
\pgfsetdash{}{0pt}%
\pgfpathmoveto{\pgfqpoint{4.470900in}{0.660000in}}%
\pgfpathlineto{\pgfqpoint{4.477100in}{0.660000in}}%
\pgfpathlineto{\pgfqpoint{4.477100in}{3.280746in}}%
\pgfpathlineto{\pgfqpoint{4.470900in}{3.280746in}}%
\pgfpathlineto{\pgfqpoint{4.470900in}{0.660000in}}%
\pgfpathclose%
\pgfusepath{fill}%
\end{pgfscope}%
\begin{pgfscope}%
\pgfpathrectangle{\pgfqpoint{1.250000in}{0.660000in}}{\pgfqpoint{7.750000in}{4.620000in}}%
\pgfusepath{clip}%
\pgfsetbuttcap%
\pgfsetmiterjoin%
\definecolor{currentfill}{rgb}{0.000000,0.000000,0.000000}%
\pgfsetfillcolor{currentfill}%
\pgfsetlinewidth{0.000000pt}%
\definecolor{currentstroke}{rgb}{0.000000,0.000000,0.000000}%
\pgfsetstrokecolor{currentstroke}%
\pgfsetstrokeopacity{0.000000}%
\pgfsetdash{}{0pt}%
\pgfpathmoveto{\pgfqpoint{4.478650in}{0.660000in}}%
\pgfpathlineto{\pgfqpoint{4.484850in}{0.660000in}}%
\pgfpathlineto{\pgfqpoint{4.484850in}{3.273760in}}%
\pgfpathlineto{\pgfqpoint{4.478650in}{3.273760in}}%
\pgfpathlineto{\pgfqpoint{4.478650in}{0.660000in}}%
\pgfpathclose%
\pgfusepath{fill}%
\end{pgfscope}%
\begin{pgfscope}%
\pgfpathrectangle{\pgfqpoint{1.250000in}{0.660000in}}{\pgfqpoint{7.750000in}{4.620000in}}%
\pgfusepath{clip}%
\pgfsetbuttcap%
\pgfsetmiterjoin%
\definecolor{currentfill}{rgb}{0.000000,0.000000,0.000000}%
\pgfsetfillcolor{currentfill}%
\pgfsetlinewidth{0.000000pt}%
\definecolor{currentstroke}{rgb}{0.000000,0.000000,0.000000}%
\pgfsetstrokecolor{currentstroke}%
\pgfsetstrokeopacity{0.000000}%
\pgfsetdash{}{0pt}%
\pgfpathmoveto{\pgfqpoint{4.486400in}{0.660000in}}%
\pgfpathlineto{\pgfqpoint{4.492600in}{0.660000in}}%
\pgfpathlineto{\pgfqpoint{4.492600in}{3.284249in}}%
\pgfpathlineto{\pgfqpoint{4.486400in}{3.284249in}}%
\pgfpathlineto{\pgfqpoint{4.486400in}{0.660000in}}%
\pgfpathclose%
\pgfusepath{fill}%
\end{pgfscope}%
\begin{pgfscope}%
\pgfpathrectangle{\pgfqpoint{1.250000in}{0.660000in}}{\pgfqpoint{7.750000in}{4.620000in}}%
\pgfusepath{clip}%
\pgfsetbuttcap%
\pgfsetmiterjoin%
\definecolor{currentfill}{rgb}{0.000000,0.000000,0.000000}%
\pgfsetfillcolor{currentfill}%
\pgfsetlinewidth{0.000000pt}%
\definecolor{currentstroke}{rgb}{0.000000,0.000000,0.000000}%
\pgfsetstrokecolor{currentstroke}%
\pgfsetstrokeopacity{0.000000}%
\pgfsetdash{}{0pt}%
\pgfpathmoveto{\pgfqpoint{4.494150in}{0.660000in}}%
\pgfpathlineto{\pgfqpoint{4.500350in}{0.660000in}}%
\pgfpathlineto{\pgfqpoint{4.500350in}{3.276099in}}%
\pgfpathlineto{\pgfqpoint{4.494150in}{3.276099in}}%
\pgfpathlineto{\pgfqpoint{4.494150in}{0.660000in}}%
\pgfpathclose%
\pgfusepath{fill}%
\end{pgfscope}%
\begin{pgfscope}%
\pgfpathrectangle{\pgfqpoint{1.250000in}{0.660000in}}{\pgfqpoint{7.750000in}{4.620000in}}%
\pgfusepath{clip}%
\pgfsetbuttcap%
\pgfsetmiterjoin%
\definecolor{currentfill}{rgb}{0.000000,0.000000,0.000000}%
\pgfsetfillcolor{currentfill}%
\pgfsetlinewidth{0.000000pt}%
\definecolor{currentstroke}{rgb}{0.000000,0.000000,0.000000}%
\pgfsetstrokecolor{currentstroke}%
\pgfsetstrokeopacity{0.000000}%
\pgfsetdash{}{0pt}%
\pgfpathmoveto{\pgfqpoint{4.501900in}{0.660000in}}%
\pgfpathlineto{\pgfqpoint{4.508100in}{0.660000in}}%
\pgfpathlineto{\pgfqpoint{4.508100in}{3.284125in}}%
\pgfpathlineto{\pgfqpoint{4.501900in}{3.284125in}}%
\pgfpathlineto{\pgfqpoint{4.501900in}{0.660000in}}%
\pgfpathclose%
\pgfusepath{fill}%
\end{pgfscope}%
\begin{pgfscope}%
\pgfpathrectangle{\pgfqpoint{1.250000in}{0.660000in}}{\pgfqpoint{7.750000in}{4.620000in}}%
\pgfusepath{clip}%
\pgfsetbuttcap%
\pgfsetmiterjoin%
\definecolor{currentfill}{rgb}{0.000000,0.000000,0.000000}%
\pgfsetfillcolor{currentfill}%
\pgfsetlinewidth{0.000000pt}%
\definecolor{currentstroke}{rgb}{0.000000,0.000000,0.000000}%
\pgfsetstrokecolor{currentstroke}%
\pgfsetstrokeopacity{0.000000}%
\pgfsetdash{}{0pt}%
\pgfpathmoveto{\pgfqpoint{4.509650in}{0.660000in}}%
\pgfpathlineto{\pgfqpoint{4.515850in}{0.660000in}}%
\pgfpathlineto{\pgfqpoint{4.515850in}{3.278241in}}%
\pgfpathlineto{\pgfqpoint{4.509650in}{3.278241in}}%
\pgfpathlineto{\pgfqpoint{4.509650in}{0.660000in}}%
\pgfpathclose%
\pgfusepath{fill}%
\end{pgfscope}%
\begin{pgfscope}%
\pgfpathrectangle{\pgfqpoint{1.250000in}{0.660000in}}{\pgfqpoint{7.750000in}{4.620000in}}%
\pgfusepath{clip}%
\pgfsetbuttcap%
\pgfsetmiterjoin%
\definecolor{currentfill}{rgb}{0.000000,0.000000,0.000000}%
\pgfsetfillcolor{currentfill}%
\pgfsetlinewidth{0.000000pt}%
\definecolor{currentstroke}{rgb}{0.000000,0.000000,0.000000}%
\pgfsetstrokecolor{currentstroke}%
\pgfsetstrokeopacity{0.000000}%
\pgfsetdash{}{0pt}%
\pgfpathmoveto{\pgfqpoint{4.517400in}{0.660000in}}%
\pgfpathlineto{\pgfqpoint{4.523600in}{0.660000in}}%
\pgfpathlineto{\pgfqpoint{4.523600in}{3.272338in}}%
\pgfpathlineto{\pgfqpoint{4.517400in}{3.272338in}}%
\pgfpathlineto{\pgfqpoint{4.517400in}{0.660000in}}%
\pgfpathclose%
\pgfusepath{fill}%
\end{pgfscope}%
\begin{pgfscope}%
\pgfpathrectangle{\pgfqpoint{1.250000in}{0.660000in}}{\pgfqpoint{7.750000in}{4.620000in}}%
\pgfusepath{clip}%
\pgfsetbuttcap%
\pgfsetmiterjoin%
\definecolor{currentfill}{rgb}{0.000000,0.000000,0.000000}%
\pgfsetfillcolor{currentfill}%
\pgfsetlinewidth{0.000000pt}%
\definecolor{currentstroke}{rgb}{0.000000,0.000000,0.000000}%
\pgfsetstrokecolor{currentstroke}%
\pgfsetstrokeopacity{0.000000}%
\pgfsetdash{}{0pt}%
\pgfpathmoveto{\pgfqpoint{4.525150in}{0.660000in}}%
\pgfpathlineto{\pgfqpoint{4.531350in}{0.660000in}}%
\pgfpathlineto{\pgfqpoint{4.531350in}{3.270842in}}%
\pgfpathlineto{\pgfqpoint{4.525150in}{3.270842in}}%
\pgfpathlineto{\pgfqpoint{4.525150in}{0.660000in}}%
\pgfpathclose%
\pgfusepath{fill}%
\end{pgfscope}%
\begin{pgfscope}%
\pgfpathrectangle{\pgfqpoint{1.250000in}{0.660000in}}{\pgfqpoint{7.750000in}{4.620000in}}%
\pgfusepath{clip}%
\pgfsetbuttcap%
\pgfsetmiterjoin%
\definecolor{currentfill}{rgb}{0.000000,0.000000,0.000000}%
\pgfsetfillcolor{currentfill}%
\pgfsetlinewidth{0.000000pt}%
\definecolor{currentstroke}{rgb}{0.000000,0.000000,0.000000}%
\pgfsetstrokecolor{currentstroke}%
\pgfsetstrokeopacity{0.000000}%
\pgfsetdash{}{0pt}%
\pgfpathmoveto{\pgfqpoint{4.532900in}{0.660000in}}%
\pgfpathlineto{\pgfqpoint{4.539100in}{0.660000in}}%
\pgfpathlineto{\pgfqpoint{4.539100in}{3.276638in}}%
\pgfpathlineto{\pgfqpoint{4.532900in}{3.276638in}}%
\pgfpathlineto{\pgfqpoint{4.532900in}{0.660000in}}%
\pgfpathclose%
\pgfusepath{fill}%
\end{pgfscope}%
\begin{pgfscope}%
\pgfpathrectangle{\pgfqpoint{1.250000in}{0.660000in}}{\pgfqpoint{7.750000in}{4.620000in}}%
\pgfusepath{clip}%
\pgfsetbuttcap%
\pgfsetmiterjoin%
\definecolor{currentfill}{rgb}{0.000000,0.000000,0.000000}%
\pgfsetfillcolor{currentfill}%
\pgfsetlinewidth{0.000000pt}%
\definecolor{currentstroke}{rgb}{0.000000,0.000000,0.000000}%
\pgfsetstrokecolor{currentstroke}%
\pgfsetstrokeopacity{0.000000}%
\pgfsetdash{}{0pt}%
\pgfpathmoveto{\pgfqpoint{4.540650in}{0.660000in}}%
\pgfpathlineto{\pgfqpoint{4.546850in}{0.660000in}}%
\pgfpathlineto{\pgfqpoint{4.546850in}{3.323364in}}%
\pgfpathlineto{\pgfqpoint{4.540650in}{3.323364in}}%
\pgfpathlineto{\pgfqpoint{4.540650in}{0.660000in}}%
\pgfpathclose%
\pgfusepath{fill}%
\end{pgfscope}%
\begin{pgfscope}%
\pgfpathrectangle{\pgfqpoint{1.250000in}{0.660000in}}{\pgfqpoint{7.750000in}{4.620000in}}%
\pgfusepath{clip}%
\pgfsetbuttcap%
\pgfsetmiterjoin%
\definecolor{currentfill}{rgb}{0.000000,0.000000,0.000000}%
\pgfsetfillcolor{currentfill}%
\pgfsetlinewidth{0.000000pt}%
\definecolor{currentstroke}{rgb}{0.000000,0.000000,0.000000}%
\pgfsetstrokecolor{currentstroke}%
\pgfsetstrokeopacity{0.000000}%
\pgfsetdash{}{0pt}%
\pgfpathmoveto{\pgfqpoint{4.548400in}{0.660000in}}%
\pgfpathlineto{\pgfqpoint{4.554600in}{0.660000in}}%
\pgfpathlineto{\pgfqpoint{4.554600in}{3.306318in}}%
\pgfpathlineto{\pgfqpoint{4.548400in}{3.306318in}}%
\pgfpathlineto{\pgfqpoint{4.548400in}{0.660000in}}%
\pgfpathclose%
\pgfusepath{fill}%
\end{pgfscope}%
\begin{pgfscope}%
\pgfpathrectangle{\pgfqpoint{1.250000in}{0.660000in}}{\pgfqpoint{7.750000in}{4.620000in}}%
\pgfusepath{clip}%
\pgfsetbuttcap%
\pgfsetmiterjoin%
\definecolor{currentfill}{rgb}{0.000000,0.000000,0.000000}%
\pgfsetfillcolor{currentfill}%
\pgfsetlinewidth{0.000000pt}%
\definecolor{currentstroke}{rgb}{0.000000,0.000000,0.000000}%
\pgfsetstrokecolor{currentstroke}%
\pgfsetstrokeopacity{0.000000}%
\pgfsetdash{}{0pt}%
\pgfpathmoveto{\pgfqpoint{4.556150in}{0.660000in}}%
\pgfpathlineto{\pgfqpoint{4.562350in}{0.660000in}}%
\pgfpathlineto{\pgfqpoint{4.562350in}{3.307457in}}%
\pgfpathlineto{\pgfqpoint{4.556150in}{3.307457in}}%
\pgfpathlineto{\pgfqpoint{4.556150in}{0.660000in}}%
\pgfpathclose%
\pgfusepath{fill}%
\end{pgfscope}%
\begin{pgfscope}%
\pgfpathrectangle{\pgfqpoint{1.250000in}{0.660000in}}{\pgfqpoint{7.750000in}{4.620000in}}%
\pgfusepath{clip}%
\pgfsetbuttcap%
\pgfsetmiterjoin%
\definecolor{currentfill}{rgb}{0.000000,0.000000,0.000000}%
\pgfsetfillcolor{currentfill}%
\pgfsetlinewidth{0.000000pt}%
\definecolor{currentstroke}{rgb}{0.000000,0.000000,0.000000}%
\pgfsetstrokecolor{currentstroke}%
\pgfsetstrokeopacity{0.000000}%
\pgfsetdash{}{0pt}%
\pgfpathmoveto{\pgfqpoint{4.563900in}{0.660000in}}%
\pgfpathlineto{\pgfqpoint{4.570100in}{0.660000in}}%
\pgfpathlineto{\pgfqpoint{4.570100in}{3.301191in}}%
\pgfpathlineto{\pgfqpoint{4.563900in}{3.301191in}}%
\pgfpathlineto{\pgfqpoint{4.563900in}{0.660000in}}%
\pgfpathclose%
\pgfusepath{fill}%
\end{pgfscope}%
\begin{pgfscope}%
\pgfpathrectangle{\pgfqpoint{1.250000in}{0.660000in}}{\pgfqpoint{7.750000in}{4.620000in}}%
\pgfusepath{clip}%
\pgfsetbuttcap%
\pgfsetmiterjoin%
\definecolor{currentfill}{rgb}{0.000000,0.000000,0.000000}%
\pgfsetfillcolor{currentfill}%
\pgfsetlinewidth{0.000000pt}%
\definecolor{currentstroke}{rgb}{0.000000,0.000000,0.000000}%
\pgfsetstrokecolor{currentstroke}%
\pgfsetstrokeopacity{0.000000}%
\pgfsetdash{}{0pt}%
\pgfpathmoveto{\pgfqpoint{4.571650in}{0.660000in}}%
\pgfpathlineto{\pgfqpoint{4.577850in}{0.660000in}}%
\pgfpathlineto{\pgfqpoint{4.577850in}{3.349735in}}%
\pgfpathlineto{\pgfqpoint{4.571650in}{3.349735in}}%
\pgfpathlineto{\pgfqpoint{4.571650in}{0.660000in}}%
\pgfpathclose%
\pgfusepath{fill}%
\end{pgfscope}%
\begin{pgfscope}%
\pgfpathrectangle{\pgfqpoint{1.250000in}{0.660000in}}{\pgfqpoint{7.750000in}{4.620000in}}%
\pgfusepath{clip}%
\pgfsetbuttcap%
\pgfsetmiterjoin%
\definecolor{currentfill}{rgb}{0.000000,0.000000,0.000000}%
\pgfsetfillcolor{currentfill}%
\pgfsetlinewidth{0.000000pt}%
\definecolor{currentstroke}{rgb}{0.000000,0.000000,0.000000}%
\pgfsetstrokecolor{currentstroke}%
\pgfsetstrokeopacity{0.000000}%
\pgfsetdash{}{0pt}%
\pgfpathmoveto{\pgfqpoint{4.579400in}{0.660000in}}%
\pgfpathlineto{\pgfqpoint{4.585600in}{0.660000in}}%
\pgfpathlineto{\pgfqpoint{4.585600in}{3.279290in}}%
\pgfpathlineto{\pgfqpoint{4.579400in}{3.279290in}}%
\pgfpathlineto{\pgfqpoint{4.579400in}{0.660000in}}%
\pgfpathclose%
\pgfusepath{fill}%
\end{pgfscope}%
\begin{pgfscope}%
\pgfpathrectangle{\pgfqpoint{1.250000in}{0.660000in}}{\pgfqpoint{7.750000in}{4.620000in}}%
\pgfusepath{clip}%
\pgfsetbuttcap%
\pgfsetmiterjoin%
\definecolor{currentfill}{rgb}{0.000000,0.000000,0.000000}%
\pgfsetfillcolor{currentfill}%
\pgfsetlinewidth{0.000000pt}%
\definecolor{currentstroke}{rgb}{0.000000,0.000000,0.000000}%
\pgfsetstrokecolor{currentstroke}%
\pgfsetstrokeopacity{0.000000}%
\pgfsetdash{}{0pt}%
\pgfpathmoveto{\pgfqpoint{4.587150in}{0.660000in}}%
\pgfpathlineto{\pgfqpoint{4.593350in}{0.660000in}}%
\pgfpathlineto{\pgfqpoint{4.593350in}{3.291760in}}%
\pgfpathlineto{\pgfqpoint{4.587150in}{3.291760in}}%
\pgfpathlineto{\pgfqpoint{4.587150in}{0.660000in}}%
\pgfpathclose%
\pgfusepath{fill}%
\end{pgfscope}%
\begin{pgfscope}%
\pgfpathrectangle{\pgfqpoint{1.250000in}{0.660000in}}{\pgfqpoint{7.750000in}{4.620000in}}%
\pgfusepath{clip}%
\pgfsetbuttcap%
\pgfsetmiterjoin%
\definecolor{currentfill}{rgb}{0.000000,0.000000,0.000000}%
\pgfsetfillcolor{currentfill}%
\pgfsetlinewidth{0.000000pt}%
\definecolor{currentstroke}{rgb}{0.000000,0.000000,0.000000}%
\pgfsetstrokecolor{currentstroke}%
\pgfsetstrokeopacity{0.000000}%
\pgfsetdash{}{0pt}%
\pgfpathmoveto{\pgfqpoint{4.594900in}{0.660000in}}%
\pgfpathlineto{\pgfqpoint{4.601100in}{0.660000in}}%
\pgfpathlineto{\pgfqpoint{4.601100in}{3.289127in}}%
\pgfpathlineto{\pgfqpoint{4.594900in}{3.289127in}}%
\pgfpathlineto{\pgfqpoint{4.594900in}{0.660000in}}%
\pgfpathclose%
\pgfusepath{fill}%
\end{pgfscope}%
\begin{pgfscope}%
\pgfpathrectangle{\pgfqpoint{1.250000in}{0.660000in}}{\pgfqpoint{7.750000in}{4.620000in}}%
\pgfusepath{clip}%
\pgfsetbuttcap%
\pgfsetmiterjoin%
\definecolor{currentfill}{rgb}{0.000000,0.000000,0.000000}%
\pgfsetfillcolor{currentfill}%
\pgfsetlinewidth{0.000000pt}%
\definecolor{currentstroke}{rgb}{0.000000,0.000000,0.000000}%
\pgfsetstrokecolor{currentstroke}%
\pgfsetstrokeopacity{0.000000}%
\pgfsetdash{}{0pt}%
\pgfpathmoveto{\pgfqpoint{4.602650in}{0.660000in}}%
\pgfpathlineto{\pgfqpoint{4.608850in}{0.660000in}}%
\pgfpathlineto{\pgfqpoint{4.608850in}{3.278261in}}%
\pgfpathlineto{\pgfqpoint{4.602650in}{3.278261in}}%
\pgfpathlineto{\pgfqpoint{4.602650in}{0.660000in}}%
\pgfpathclose%
\pgfusepath{fill}%
\end{pgfscope}%
\begin{pgfscope}%
\pgfpathrectangle{\pgfqpoint{1.250000in}{0.660000in}}{\pgfqpoint{7.750000in}{4.620000in}}%
\pgfusepath{clip}%
\pgfsetbuttcap%
\pgfsetmiterjoin%
\definecolor{currentfill}{rgb}{0.000000,0.000000,0.000000}%
\pgfsetfillcolor{currentfill}%
\pgfsetlinewidth{0.000000pt}%
\definecolor{currentstroke}{rgb}{0.000000,0.000000,0.000000}%
\pgfsetstrokecolor{currentstroke}%
\pgfsetstrokeopacity{0.000000}%
\pgfsetdash{}{0pt}%
\pgfpathmoveto{\pgfqpoint{4.610400in}{0.660000in}}%
\pgfpathlineto{\pgfqpoint{4.616600in}{0.660000in}}%
\pgfpathlineto{\pgfqpoint{4.616600in}{3.274383in}}%
\pgfpathlineto{\pgfqpoint{4.610400in}{3.274383in}}%
\pgfpathlineto{\pgfqpoint{4.610400in}{0.660000in}}%
\pgfpathclose%
\pgfusepath{fill}%
\end{pgfscope}%
\begin{pgfscope}%
\pgfpathrectangle{\pgfqpoint{1.250000in}{0.660000in}}{\pgfqpoint{7.750000in}{4.620000in}}%
\pgfusepath{clip}%
\pgfsetbuttcap%
\pgfsetmiterjoin%
\definecolor{currentfill}{rgb}{0.000000,0.000000,0.000000}%
\pgfsetfillcolor{currentfill}%
\pgfsetlinewidth{0.000000pt}%
\definecolor{currentstroke}{rgb}{0.000000,0.000000,0.000000}%
\pgfsetstrokecolor{currentstroke}%
\pgfsetstrokeopacity{0.000000}%
\pgfsetdash{}{0pt}%
\pgfpathmoveto{\pgfqpoint{4.618150in}{0.660000in}}%
\pgfpathlineto{\pgfqpoint{4.624350in}{0.660000in}}%
\pgfpathlineto{\pgfqpoint{4.624350in}{3.270331in}}%
\pgfpathlineto{\pgfqpoint{4.618150in}{3.270331in}}%
\pgfpathlineto{\pgfqpoint{4.618150in}{0.660000in}}%
\pgfpathclose%
\pgfusepath{fill}%
\end{pgfscope}%
\begin{pgfscope}%
\pgfpathrectangle{\pgfqpoint{1.250000in}{0.660000in}}{\pgfqpoint{7.750000in}{4.620000in}}%
\pgfusepath{clip}%
\pgfsetbuttcap%
\pgfsetmiterjoin%
\definecolor{currentfill}{rgb}{0.000000,0.000000,0.000000}%
\pgfsetfillcolor{currentfill}%
\pgfsetlinewidth{0.000000pt}%
\definecolor{currentstroke}{rgb}{0.000000,0.000000,0.000000}%
\pgfsetstrokecolor{currentstroke}%
\pgfsetstrokeopacity{0.000000}%
\pgfsetdash{}{0pt}%
\pgfpathmoveto{\pgfqpoint{4.625900in}{0.660000in}}%
\pgfpathlineto{\pgfqpoint{4.632100in}{0.660000in}}%
\pgfpathlineto{\pgfqpoint{4.632100in}{3.278943in}}%
\pgfpathlineto{\pgfqpoint{4.625900in}{3.278943in}}%
\pgfpathlineto{\pgfqpoint{4.625900in}{0.660000in}}%
\pgfpathclose%
\pgfusepath{fill}%
\end{pgfscope}%
\begin{pgfscope}%
\pgfpathrectangle{\pgfqpoint{1.250000in}{0.660000in}}{\pgfqpoint{7.750000in}{4.620000in}}%
\pgfusepath{clip}%
\pgfsetbuttcap%
\pgfsetmiterjoin%
\definecolor{currentfill}{rgb}{0.000000,0.000000,0.000000}%
\pgfsetfillcolor{currentfill}%
\pgfsetlinewidth{0.000000pt}%
\definecolor{currentstroke}{rgb}{0.000000,0.000000,0.000000}%
\pgfsetstrokecolor{currentstroke}%
\pgfsetstrokeopacity{0.000000}%
\pgfsetdash{}{0pt}%
\pgfpathmoveto{\pgfqpoint{4.633650in}{0.660000in}}%
\pgfpathlineto{\pgfqpoint{4.639850in}{0.660000in}}%
\pgfpathlineto{\pgfqpoint{4.639850in}{3.311610in}}%
\pgfpathlineto{\pgfqpoint{4.633650in}{3.311610in}}%
\pgfpathlineto{\pgfqpoint{4.633650in}{0.660000in}}%
\pgfpathclose%
\pgfusepath{fill}%
\end{pgfscope}%
\begin{pgfscope}%
\pgfpathrectangle{\pgfqpoint{1.250000in}{0.660000in}}{\pgfqpoint{7.750000in}{4.620000in}}%
\pgfusepath{clip}%
\pgfsetbuttcap%
\pgfsetmiterjoin%
\definecolor{currentfill}{rgb}{0.000000,0.000000,0.000000}%
\pgfsetfillcolor{currentfill}%
\pgfsetlinewidth{0.000000pt}%
\definecolor{currentstroke}{rgb}{0.000000,0.000000,0.000000}%
\pgfsetstrokecolor{currentstroke}%
\pgfsetstrokeopacity{0.000000}%
\pgfsetdash{}{0pt}%
\pgfpathmoveto{\pgfqpoint{4.641400in}{0.660000in}}%
\pgfpathlineto{\pgfqpoint{4.647600in}{0.660000in}}%
\pgfpathlineto{\pgfqpoint{4.647600in}{3.267129in}}%
\pgfpathlineto{\pgfqpoint{4.641400in}{3.267129in}}%
\pgfpathlineto{\pgfqpoint{4.641400in}{0.660000in}}%
\pgfpathclose%
\pgfusepath{fill}%
\end{pgfscope}%
\begin{pgfscope}%
\pgfpathrectangle{\pgfqpoint{1.250000in}{0.660000in}}{\pgfqpoint{7.750000in}{4.620000in}}%
\pgfusepath{clip}%
\pgfsetbuttcap%
\pgfsetmiterjoin%
\definecolor{currentfill}{rgb}{0.000000,0.000000,0.000000}%
\pgfsetfillcolor{currentfill}%
\pgfsetlinewidth{0.000000pt}%
\definecolor{currentstroke}{rgb}{0.000000,0.000000,0.000000}%
\pgfsetstrokecolor{currentstroke}%
\pgfsetstrokeopacity{0.000000}%
\pgfsetdash{}{0pt}%
\pgfpathmoveto{\pgfqpoint{4.649150in}{0.660000in}}%
\pgfpathlineto{\pgfqpoint{4.655350in}{0.660000in}}%
\pgfpathlineto{\pgfqpoint{4.655350in}{3.289097in}}%
\pgfpathlineto{\pgfqpoint{4.649150in}{3.289097in}}%
\pgfpathlineto{\pgfqpoint{4.649150in}{0.660000in}}%
\pgfpathclose%
\pgfusepath{fill}%
\end{pgfscope}%
\begin{pgfscope}%
\pgfpathrectangle{\pgfqpoint{1.250000in}{0.660000in}}{\pgfqpoint{7.750000in}{4.620000in}}%
\pgfusepath{clip}%
\pgfsetbuttcap%
\pgfsetmiterjoin%
\definecolor{currentfill}{rgb}{0.000000,0.000000,0.000000}%
\pgfsetfillcolor{currentfill}%
\pgfsetlinewidth{0.000000pt}%
\definecolor{currentstroke}{rgb}{0.000000,0.000000,0.000000}%
\pgfsetstrokecolor{currentstroke}%
\pgfsetstrokeopacity{0.000000}%
\pgfsetdash{}{0pt}%
\pgfpathmoveto{\pgfqpoint{4.656900in}{0.660000in}}%
\pgfpathlineto{\pgfqpoint{4.663100in}{0.660000in}}%
\pgfpathlineto{\pgfqpoint{4.663100in}{3.296117in}}%
\pgfpathlineto{\pgfqpoint{4.656900in}{3.296117in}}%
\pgfpathlineto{\pgfqpoint{4.656900in}{0.660000in}}%
\pgfpathclose%
\pgfusepath{fill}%
\end{pgfscope}%
\begin{pgfscope}%
\pgfpathrectangle{\pgfqpoint{1.250000in}{0.660000in}}{\pgfqpoint{7.750000in}{4.620000in}}%
\pgfusepath{clip}%
\pgfsetbuttcap%
\pgfsetmiterjoin%
\definecolor{currentfill}{rgb}{0.000000,0.000000,0.000000}%
\pgfsetfillcolor{currentfill}%
\pgfsetlinewidth{0.000000pt}%
\definecolor{currentstroke}{rgb}{0.000000,0.000000,0.000000}%
\pgfsetstrokecolor{currentstroke}%
\pgfsetstrokeopacity{0.000000}%
\pgfsetdash{}{0pt}%
\pgfpathmoveto{\pgfqpoint{4.664650in}{0.660000in}}%
\pgfpathlineto{\pgfqpoint{4.670850in}{0.660000in}}%
\pgfpathlineto{\pgfqpoint{4.670850in}{3.289927in}}%
\pgfpathlineto{\pgfqpoint{4.664650in}{3.289927in}}%
\pgfpathlineto{\pgfqpoint{4.664650in}{0.660000in}}%
\pgfpathclose%
\pgfusepath{fill}%
\end{pgfscope}%
\begin{pgfscope}%
\pgfpathrectangle{\pgfqpoint{1.250000in}{0.660000in}}{\pgfqpoint{7.750000in}{4.620000in}}%
\pgfusepath{clip}%
\pgfsetbuttcap%
\pgfsetmiterjoin%
\definecolor{currentfill}{rgb}{0.000000,0.000000,0.000000}%
\pgfsetfillcolor{currentfill}%
\pgfsetlinewidth{0.000000pt}%
\definecolor{currentstroke}{rgb}{0.000000,0.000000,0.000000}%
\pgfsetstrokecolor{currentstroke}%
\pgfsetstrokeopacity{0.000000}%
\pgfsetdash{}{0pt}%
\pgfpathmoveto{\pgfqpoint{4.672400in}{0.660000in}}%
\pgfpathlineto{\pgfqpoint{4.678600in}{0.660000in}}%
\pgfpathlineto{\pgfqpoint{4.678600in}{3.276258in}}%
\pgfpathlineto{\pgfqpoint{4.672400in}{3.276258in}}%
\pgfpathlineto{\pgfqpoint{4.672400in}{0.660000in}}%
\pgfpathclose%
\pgfusepath{fill}%
\end{pgfscope}%
\begin{pgfscope}%
\pgfpathrectangle{\pgfqpoint{1.250000in}{0.660000in}}{\pgfqpoint{7.750000in}{4.620000in}}%
\pgfusepath{clip}%
\pgfsetbuttcap%
\pgfsetmiterjoin%
\definecolor{currentfill}{rgb}{0.000000,0.000000,0.000000}%
\pgfsetfillcolor{currentfill}%
\pgfsetlinewidth{0.000000pt}%
\definecolor{currentstroke}{rgb}{0.000000,0.000000,0.000000}%
\pgfsetstrokecolor{currentstroke}%
\pgfsetstrokeopacity{0.000000}%
\pgfsetdash{}{0pt}%
\pgfpathmoveto{\pgfqpoint{4.680150in}{0.660000in}}%
\pgfpathlineto{\pgfqpoint{4.686350in}{0.660000in}}%
\pgfpathlineto{\pgfqpoint{4.686350in}{3.282253in}}%
\pgfpathlineto{\pgfqpoint{4.680150in}{3.282253in}}%
\pgfpathlineto{\pgfqpoint{4.680150in}{0.660000in}}%
\pgfpathclose%
\pgfusepath{fill}%
\end{pgfscope}%
\begin{pgfscope}%
\pgfpathrectangle{\pgfqpoint{1.250000in}{0.660000in}}{\pgfqpoint{7.750000in}{4.620000in}}%
\pgfusepath{clip}%
\pgfsetbuttcap%
\pgfsetmiterjoin%
\definecolor{currentfill}{rgb}{0.000000,0.000000,0.000000}%
\pgfsetfillcolor{currentfill}%
\pgfsetlinewidth{0.000000pt}%
\definecolor{currentstroke}{rgb}{0.000000,0.000000,0.000000}%
\pgfsetstrokecolor{currentstroke}%
\pgfsetstrokeopacity{0.000000}%
\pgfsetdash{}{0pt}%
\pgfpathmoveto{\pgfqpoint{4.687900in}{0.660000in}}%
\pgfpathlineto{\pgfqpoint{4.694100in}{0.660000in}}%
\pgfpathlineto{\pgfqpoint{4.694100in}{3.286882in}}%
\pgfpathlineto{\pgfqpoint{4.687900in}{3.286882in}}%
\pgfpathlineto{\pgfqpoint{4.687900in}{0.660000in}}%
\pgfpathclose%
\pgfusepath{fill}%
\end{pgfscope}%
\begin{pgfscope}%
\pgfpathrectangle{\pgfqpoint{1.250000in}{0.660000in}}{\pgfqpoint{7.750000in}{4.620000in}}%
\pgfusepath{clip}%
\pgfsetbuttcap%
\pgfsetmiterjoin%
\definecolor{currentfill}{rgb}{0.000000,0.000000,0.000000}%
\pgfsetfillcolor{currentfill}%
\pgfsetlinewidth{0.000000pt}%
\definecolor{currentstroke}{rgb}{0.000000,0.000000,0.000000}%
\pgfsetstrokecolor{currentstroke}%
\pgfsetstrokeopacity{0.000000}%
\pgfsetdash{}{0pt}%
\pgfpathmoveto{\pgfqpoint{4.695650in}{0.660000in}}%
\pgfpathlineto{\pgfqpoint{4.701850in}{0.660000in}}%
\pgfpathlineto{\pgfqpoint{4.701850in}{3.283070in}}%
\pgfpathlineto{\pgfqpoint{4.695650in}{3.283070in}}%
\pgfpathlineto{\pgfqpoint{4.695650in}{0.660000in}}%
\pgfpathclose%
\pgfusepath{fill}%
\end{pgfscope}%
\begin{pgfscope}%
\pgfpathrectangle{\pgfqpoint{1.250000in}{0.660000in}}{\pgfqpoint{7.750000in}{4.620000in}}%
\pgfusepath{clip}%
\pgfsetbuttcap%
\pgfsetmiterjoin%
\definecolor{currentfill}{rgb}{0.000000,0.000000,0.000000}%
\pgfsetfillcolor{currentfill}%
\pgfsetlinewidth{0.000000pt}%
\definecolor{currentstroke}{rgb}{0.000000,0.000000,0.000000}%
\pgfsetstrokecolor{currentstroke}%
\pgfsetstrokeopacity{0.000000}%
\pgfsetdash{}{0pt}%
\pgfpathmoveto{\pgfqpoint{4.703400in}{0.660000in}}%
\pgfpathlineto{\pgfqpoint{4.709600in}{0.660000in}}%
\pgfpathlineto{\pgfqpoint{4.709600in}{3.314567in}}%
\pgfpathlineto{\pgfqpoint{4.703400in}{3.314567in}}%
\pgfpathlineto{\pgfqpoint{4.703400in}{0.660000in}}%
\pgfpathclose%
\pgfusepath{fill}%
\end{pgfscope}%
\begin{pgfscope}%
\pgfpathrectangle{\pgfqpoint{1.250000in}{0.660000in}}{\pgfqpoint{7.750000in}{4.620000in}}%
\pgfusepath{clip}%
\pgfsetbuttcap%
\pgfsetmiterjoin%
\definecolor{currentfill}{rgb}{0.000000,0.000000,0.000000}%
\pgfsetfillcolor{currentfill}%
\pgfsetlinewidth{0.000000pt}%
\definecolor{currentstroke}{rgb}{0.000000,0.000000,0.000000}%
\pgfsetstrokecolor{currentstroke}%
\pgfsetstrokeopacity{0.000000}%
\pgfsetdash{}{0pt}%
\pgfpathmoveto{\pgfqpoint{4.711150in}{0.660000in}}%
\pgfpathlineto{\pgfqpoint{4.717350in}{0.660000in}}%
\pgfpathlineto{\pgfqpoint{4.717350in}{3.306829in}}%
\pgfpathlineto{\pgfqpoint{4.711150in}{3.306829in}}%
\pgfpathlineto{\pgfqpoint{4.711150in}{0.660000in}}%
\pgfpathclose%
\pgfusepath{fill}%
\end{pgfscope}%
\begin{pgfscope}%
\pgfpathrectangle{\pgfqpoint{1.250000in}{0.660000in}}{\pgfqpoint{7.750000in}{4.620000in}}%
\pgfusepath{clip}%
\pgfsetbuttcap%
\pgfsetmiterjoin%
\definecolor{currentfill}{rgb}{0.000000,0.000000,0.000000}%
\pgfsetfillcolor{currentfill}%
\pgfsetlinewidth{0.000000pt}%
\definecolor{currentstroke}{rgb}{0.000000,0.000000,0.000000}%
\pgfsetstrokecolor{currentstroke}%
\pgfsetstrokeopacity{0.000000}%
\pgfsetdash{}{0pt}%
\pgfpathmoveto{\pgfqpoint{4.718900in}{0.660000in}}%
\pgfpathlineto{\pgfqpoint{4.725100in}{0.660000in}}%
\pgfpathlineto{\pgfqpoint{4.725100in}{3.273647in}}%
\pgfpathlineto{\pgfqpoint{4.718900in}{3.273647in}}%
\pgfpathlineto{\pgfqpoint{4.718900in}{0.660000in}}%
\pgfpathclose%
\pgfusepath{fill}%
\end{pgfscope}%
\begin{pgfscope}%
\pgfpathrectangle{\pgfqpoint{1.250000in}{0.660000in}}{\pgfqpoint{7.750000in}{4.620000in}}%
\pgfusepath{clip}%
\pgfsetbuttcap%
\pgfsetmiterjoin%
\definecolor{currentfill}{rgb}{0.000000,0.000000,0.000000}%
\pgfsetfillcolor{currentfill}%
\pgfsetlinewidth{0.000000pt}%
\definecolor{currentstroke}{rgb}{0.000000,0.000000,0.000000}%
\pgfsetstrokecolor{currentstroke}%
\pgfsetstrokeopacity{0.000000}%
\pgfsetdash{}{0pt}%
\pgfpathmoveto{\pgfqpoint{4.726650in}{0.660000in}}%
\pgfpathlineto{\pgfqpoint{4.732850in}{0.660000in}}%
\pgfpathlineto{\pgfqpoint{4.732850in}{3.285843in}}%
\pgfpathlineto{\pgfqpoint{4.726650in}{3.285843in}}%
\pgfpathlineto{\pgfqpoint{4.726650in}{0.660000in}}%
\pgfpathclose%
\pgfusepath{fill}%
\end{pgfscope}%
\begin{pgfscope}%
\pgfpathrectangle{\pgfqpoint{1.250000in}{0.660000in}}{\pgfqpoint{7.750000in}{4.620000in}}%
\pgfusepath{clip}%
\pgfsetbuttcap%
\pgfsetmiterjoin%
\definecolor{currentfill}{rgb}{0.000000,0.000000,0.000000}%
\pgfsetfillcolor{currentfill}%
\pgfsetlinewidth{0.000000pt}%
\definecolor{currentstroke}{rgb}{0.000000,0.000000,0.000000}%
\pgfsetstrokecolor{currentstroke}%
\pgfsetstrokeopacity{0.000000}%
\pgfsetdash{}{0pt}%
\pgfpathmoveto{\pgfqpoint{4.734400in}{0.660000in}}%
\pgfpathlineto{\pgfqpoint{4.740600in}{0.660000in}}%
\pgfpathlineto{\pgfqpoint{4.740600in}{3.264205in}}%
\pgfpathlineto{\pgfqpoint{4.734400in}{3.264205in}}%
\pgfpathlineto{\pgfqpoint{4.734400in}{0.660000in}}%
\pgfpathclose%
\pgfusepath{fill}%
\end{pgfscope}%
\begin{pgfscope}%
\pgfpathrectangle{\pgfqpoint{1.250000in}{0.660000in}}{\pgfqpoint{7.750000in}{4.620000in}}%
\pgfusepath{clip}%
\pgfsetbuttcap%
\pgfsetmiterjoin%
\definecolor{currentfill}{rgb}{0.000000,0.000000,0.000000}%
\pgfsetfillcolor{currentfill}%
\pgfsetlinewidth{0.000000pt}%
\definecolor{currentstroke}{rgb}{0.000000,0.000000,0.000000}%
\pgfsetstrokecolor{currentstroke}%
\pgfsetstrokeopacity{0.000000}%
\pgfsetdash{}{0pt}%
\pgfpathmoveto{\pgfqpoint{4.742150in}{0.660000in}}%
\pgfpathlineto{\pgfqpoint{4.748350in}{0.660000in}}%
\pgfpathlineto{\pgfqpoint{4.748350in}{3.279320in}}%
\pgfpathlineto{\pgfqpoint{4.742150in}{3.279320in}}%
\pgfpathlineto{\pgfqpoint{4.742150in}{0.660000in}}%
\pgfpathclose%
\pgfusepath{fill}%
\end{pgfscope}%
\begin{pgfscope}%
\pgfpathrectangle{\pgfqpoint{1.250000in}{0.660000in}}{\pgfqpoint{7.750000in}{4.620000in}}%
\pgfusepath{clip}%
\pgfsetbuttcap%
\pgfsetmiterjoin%
\definecolor{currentfill}{rgb}{0.000000,0.000000,0.000000}%
\pgfsetfillcolor{currentfill}%
\pgfsetlinewidth{0.000000pt}%
\definecolor{currentstroke}{rgb}{0.000000,0.000000,0.000000}%
\pgfsetstrokecolor{currentstroke}%
\pgfsetstrokeopacity{0.000000}%
\pgfsetdash{}{0pt}%
\pgfpathmoveto{\pgfqpoint{4.749900in}{0.660000in}}%
\pgfpathlineto{\pgfqpoint{4.756100in}{0.660000in}}%
\pgfpathlineto{\pgfqpoint{4.756100in}{3.274123in}}%
\pgfpathlineto{\pgfqpoint{4.749900in}{3.274123in}}%
\pgfpathlineto{\pgfqpoint{4.749900in}{0.660000in}}%
\pgfpathclose%
\pgfusepath{fill}%
\end{pgfscope}%
\begin{pgfscope}%
\pgfpathrectangle{\pgfqpoint{1.250000in}{0.660000in}}{\pgfqpoint{7.750000in}{4.620000in}}%
\pgfusepath{clip}%
\pgfsetbuttcap%
\pgfsetmiterjoin%
\definecolor{currentfill}{rgb}{0.000000,0.000000,0.000000}%
\pgfsetfillcolor{currentfill}%
\pgfsetlinewidth{0.000000pt}%
\definecolor{currentstroke}{rgb}{0.000000,0.000000,0.000000}%
\pgfsetstrokecolor{currentstroke}%
\pgfsetstrokeopacity{0.000000}%
\pgfsetdash{}{0pt}%
\pgfpathmoveto{\pgfqpoint{4.757650in}{0.660000in}}%
\pgfpathlineto{\pgfqpoint{4.763850in}{0.660000in}}%
\pgfpathlineto{\pgfqpoint{4.763850in}{3.279841in}}%
\pgfpathlineto{\pgfqpoint{4.757650in}{3.279841in}}%
\pgfpathlineto{\pgfqpoint{4.757650in}{0.660000in}}%
\pgfpathclose%
\pgfusepath{fill}%
\end{pgfscope}%
\begin{pgfscope}%
\pgfpathrectangle{\pgfqpoint{1.250000in}{0.660000in}}{\pgfqpoint{7.750000in}{4.620000in}}%
\pgfusepath{clip}%
\pgfsetbuttcap%
\pgfsetmiterjoin%
\definecolor{currentfill}{rgb}{0.000000,0.000000,0.000000}%
\pgfsetfillcolor{currentfill}%
\pgfsetlinewidth{0.000000pt}%
\definecolor{currentstroke}{rgb}{0.000000,0.000000,0.000000}%
\pgfsetstrokecolor{currentstroke}%
\pgfsetstrokeopacity{0.000000}%
\pgfsetdash{}{0pt}%
\pgfpathmoveto{\pgfqpoint{4.765400in}{0.660000in}}%
\pgfpathlineto{\pgfqpoint{4.771600in}{0.660000in}}%
\pgfpathlineto{\pgfqpoint{4.771600in}{3.271856in}}%
\pgfpathlineto{\pgfqpoint{4.765400in}{3.271856in}}%
\pgfpathlineto{\pgfqpoint{4.765400in}{0.660000in}}%
\pgfpathclose%
\pgfusepath{fill}%
\end{pgfscope}%
\begin{pgfscope}%
\pgfpathrectangle{\pgfqpoint{1.250000in}{0.660000in}}{\pgfqpoint{7.750000in}{4.620000in}}%
\pgfusepath{clip}%
\pgfsetbuttcap%
\pgfsetmiterjoin%
\definecolor{currentfill}{rgb}{0.000000,0.000000,0.000000}%
\pgfsetfillcolor{currentfill}%
\pgfsetlinewidth{0.000000pt}%
\definecolor{currentstroke}{rgb}{0.000000,0.000000,0.000000}%
\pgfsetstrokecolor{currentstroke}%
\pgfsetstrokeopacity{0.000000}%
\pgfsetdash{}{0pt}%
\pgfpathmoveto{\pgfqpoint{4.773150in}{0.660000in}}%
\pgfpathlineto{\pgfqpoint{4.779350in}{0.660000in}}%
\pgfpathlineto{\pgfqpoint{4.779350in}{3.269465in}}%
\pgfpathlineto{\pgfqpoint{4.773150in}{3.269465in}}%
\pgfpathlineto{\pgfqpoint{4.773150in}{0.660000in}}%
\pgfpathclose%
\pgfusepath{fill}%
\end{pgfscope}%
\begin{pgfscope}%
\pgfpathrectangle{\pgfqpoint{1.250000in}{0.660000in}}{\pgfqpoint{7.750000in}{4.620000in}}%
\pgfusepath{clip}%
\pgfsetbuttcap%
\pgfsetmiterjoin%
\definecolor{currentfill}{rgb}{0.000000,0.000000,0.000000}%
\pgfsetfillcolor{currentfill}%
\pgfsetlinewidth{0.000000pt}%
\definecolor{currentstroke}{rgb}{0.000000,0.000000,0.000000}%
\pgfsetstrokecolor{currentstroke}%
\pgfsetstrokeopacity{0.000000}%
\pgfsetdash{}{0pt}%
\pgfpathmoveto{\pgfqpoint{4.780900in}{0.660000in}}%
\pgfpathlineto{\pgfqpoint{4.787100in}{0.660000in}}%
\pgfpathlineto{\pgfqpoint{4.787100in}{3.262675in}}%
\pgfpathlineto{\pgfqpoint{4.780900in}{3.262675in}}%
\pgfpathlineto{\pgfqpoint{4.780900in}{0.660000in}}%
\pgfpathclose%
\pgfusepath{fill}%
\end{pgfscope}%
\begin{pgfscope}%
\pgfpathrectangle{\pgfqpoint{1.250000in}{0.660000in}}{\pgfqpoint{7.750000in}{4.620000in}}%
\pgfusepath{clip}%
\pgfsetbuttcap%
\pgfsetmiterjoin%
\definecolor{currentfill}{rgb}{0.000000,0.000000,0.000000}%
\pgfsetfillcolor{currentfill}%
\pgfsetlinewidth{0.000000pt}%
\definecolor{currentstroke}{rgb}{0.000000,0.000000,0.000000}%
\pgfsetstrokecolor{currentstroke}%
\pgfsetstrokeopacity{0.000000}%
\pgfsetdash{}{0pt}%
\pgfpathmoveto{\pgfqpoint{4.788650in}{0.660000in}}%
\pgfpathlineto{\pgfqpoint{4.794850in}{0.660000in}}%
\pgfpathlineto{\pgfqpoint{4.794850in}{3.333848in}}%
\pgfpathlineto{\pgfqpoint{4.788650in}{3.333848in}}%
\pgfpathlineto{\pgfqpoint{4.788650in}{0.660000in}}%
\pgfpathclose%
\pgfusepath{fill}%
\end{pgfscope}%
\begin{pgfscope}%
\pgfpathrectangle{\pgfqpoint{1.250000in}{0.660000in}}{\pgfqpoint{7.750000in}{4.620000in}}%
\pgfusepath{clip}%
\pgfsetbuttcap%
\pgfsetmiterjoin%
\definecolor{currentfill}{rgb}{0.000000,0.000000,0.000000}%
\pgfsetfillcolor{currentfill}%
\pgfsetlinewidth{0.000000pt}%
\definecolor{currentstroke}{rgb}{0.000000,0.000000,0.000000}%
\pgfsetstrokecolor{currentstroke}%
\pgfsetstrokeopacity{0.000000}%
\pgfsetdash{}{0pt}%
\pgfpathmoveto{\pgfqpoint{4.796400in}{0.660000in}}%
\pgfpathlineto{\pgfqpoint{4.802600in}{0.660000in}}%
\pgfpathlineto{\pgfqpoint{4.802600in}{3.302927in}}%
\pgfpathlineto{\pgfqpoint{4.796400in}{3.302927in}}%
\pgfpathlineto{\pgfqpoint{4.796400in}{0.660000in}}%
\pgfpathclose%
\pgfusepath{fill}%
\end{pgfscope}%
\begin{pgfscope}%
\pgfpathrectangle{\pgfqpoint{1.250000in}{0.660000in}}{\pgfqpoint{7.750000in}{4.620000in}}%
\pgfusepath{clip}%
\pgfsetbuttcap%
\pgfsetmiterjoin%
\definecolor{currentfill}{rgb}{0.000000,0.000000,0.000000}%
\pgfsetfillcolor{currentfill}%
\pgfsetlinewidth{0.000000pt}%
\definecolor{currentstroke}{rgb}{0.000000,0.000000,0.000000}%
\pgfsetstrokecolor{currentstroke}%
\pgfsetstrokeopacity{0.000000}%
\pgfsetdash{}{0pt}%
\pgfpathmoveto{\pgfqpoint{4.804150in}{0.660000in}}%
\pgfpathlineto{\pgfqpoint{4.810350in}{0.660000in}}%
\pgfpathlineto{\pgfqpoint{4.810350in}{3.291121in}}%
\pgfpathlineto{\pgfqpoint{4.804150in}{3.291121in}}%
\pgfpathlineto{\pgfqpoint{4.804150in}{0.660000in}}%
\pgfpathclose%
\pgfusepath{fill}%
\end{pgfscope}%
\begin{pgfscope}%
\pgfpathrectangle{\pgfqpoint{1.250000in}{0.660000in}}{\pgfqpoint{7.750000in}{4.620000in}}%
\pgfusepath{clip}%
\pgfsetbuttcap%
\pgfsetmiterjoin%
\definecolor{currentfill}{rgb}{0.000000,0.000000,0.000000}%
\pgfsetfillcolor{currentfill}%
\pgfsetlinewidth{0.000000pt}%
\definecolor{currentstroke}{rgb}{0.000000,0.000000,0.000000}%
\pgfsetstrokecolor{currentstroke}%
\pgfsetstrokeopacity{0.000000}%
\pgfsetdash{}{0pt}%
\pgfpathmoveto{\pgfqpoint{4.811900in}{0.660000in}}%
\pgfpathlineto{\pgfqpoint{4.818100in}{0.660000in}}%
\pgfpathlineto{\pgfqpoint{4.818100in}{3.316126in}}%
\pgfpathlineto{\pgfqpoint{4.811900in}{3.316126in}}%
\pgfpathlineto{\pgfqpoint{4.811900in}{0.660000in}}%
\pgfpathclose%
\pgfusepath{fill}%
\end{pgfscope}%
\begin{pgfscope}%
\pgfpathrectangle{\pgfqpoint{1.250000in}{0.660000in}}{\pgfqpoint{7.750000in}{4.620000in}}%
\pgfusepath{clip}%
\pgfsetbuttcap%
\pgfsetmiterjoin%
\definecolor{currentfill}{rgb}{0.000000,0.000000,0.000000}%
\pgfsetfillcolor{currentfill}%
\pgfsetlinewidth{0.000000pt}%
\definecolor{currentstroke}{rgb}{0.000000,0.000000,0.000000}%
\pgfsetstrokecolor{currentstroke}%
\pgfsetstrokeopacity{0.000000}%
\pgfsetdash{}{0pt}%
\pgfpathmoveto{\pgfqpoint{4.819650in}{0.660000in}}%
\pgfpathlineto{\pgfqpoint{4.825850in}{0.660000in}}%
\pgfpathlineto{\pgfqpoint{4.825850in}{3.292626in}}%
\pgfpathlineto{\pgfqpoint{4.819650in}{3.292626in}}%
\pgfpathlineto{\pgfqpoint{4.819650in}{0.660000in}}%
\pgfpathclose%
\pgfusepath{fill}%
\end{pgfscope}%
\begin{pgfscope}%
\pgfpathrectangle{\pgfqpoint{1.250000in}{0.660000in}}{\pgfqpoint{7.750000in}{4.620000in}}%
\pgfusepath{clip}%
\pgfsetbuttcap%
\pgfsetmiterjoin%
\definecolor{currentfill}{rgb}{0.000000,0.000000,0.000000}%
\pgfsetfillcolor{currentfill}%
\pgfsetlinewidth{0.000000pt}%
\definecolor{currentstroke}{rgb}{0.000000,0.000000,0.000000}%
\pgfsetstrokecolor{currentstroke}%
\pgfsetstrokeopacity{0.000000}%
\pgfsetdash{}{0pt}%
\pgfpathmoveto{\pgfqpoint{4.827400in}{0.660000in}}%
\pgfpathlineto{\pgfqpoint{4.833600in}{0.660000in}}%
\pgfpathlineto{\pgfqpoint{4.833600in}{3.261095in}}%
\pgfpathlineto{\pgfqpoint{4.827400in}{3.261095in}}%
\pgfpathlineto{\pgfqpoint{4.827400in}{0.660000in}}%
\pgfpathclose%
\pgfusepath{fill}%
\end{pgfscope}%
\begin{pgfscope}%
\pgfpathrectangle{\pgfqpoint{1.250000in}{0.660000in}}{\pgfqpoint{7.750000in}{4.620000in}}%
\pgfusepath{clip}%
\pgfsetbuttcap%
\pgfsetmiterjoin%
\definecolor{currentfill}{rgb}{0.000000,0.000000,0.000000}%
\pgfsetfillcolor{currentfill}%
\pgfsetlinewidth{0.000000pt}%
\definecolor{currentstroke}{rgb}{0.000000,0.000000,0.000000}%
\pgfsetstrokecolor{currentstroke}%
\pgfsetstrokeopacity{0.000000}%
\pgfsetdash{}{0pt}%
\pgfpathmoveto{\pgfqpoint{4.835150in}{0.660000in}}%
\pgfpathlineto{\pgfqpoint{4.841350in}{0.660000in}}%
\pgfpathlineto{\pgfqpoint{4.841350in}{3.293124in}}%
\pgfpathlineto{\pgfqpoint{4.835150in}{3.293124in}}%
\pgfpathlineto{\pgfqpoint{4.835150in}{0.660000in}}%
\pgfpathclose%
\pgfusepath{fill}%
\end{pgfscope}%
\begin{pgfscope}%
\pgfpathrectangle{\pgfqpoint{1.250000in}{0.660000in}}{\pgfqpoint{7.750000in}{4.620000in}}%
\pgfusepath{clip}%
\pgfsetbuttcap%
\pgfsetmiterjoin%
\definecolor{currentfill}{rgb}{0.000000,0.000000,0.000000}%
\pgfsetfillcolor{currentfill}%
\pgfsetlinewidth{0.000000pt}%
\definecolor{currentstroke}{rgb}{0.000000,0.000000,0.000000}%
\pgfsetstrokecolor{currentstroke}%
\pgfsetstrokeopacity{0.000000}%
\pgfsetdash{}{0pt}%
\pgfpathmoveto{\pgfqpoint{4.842900in}{0.660000in}}%
\pgfpathlineto{\pgfqpoint{4.849100in}{0.660000in}}%
\pgfpathlineto{\pgfqpoint{4.849100in}{3.262545in}}%
\pgfpathlineto{\pgfqpoint{4.842900in}{3.262545in}}%
\pgfpathlineto{\pgfqpoint{4.842900in}{0.660000in}}%
\pgfpathclose%
\pgfusepath{fill}%
\end{pgfscope}%
\begin{pgfscope}%
\pgfpathrectangle{\pgfqpoint{1.250000in}{0.660000in}}{\pgfqpoint{7.750000in}{4.620000in}}%
\pgfusepath{clip}%
\pgfsetbuttcap%
\pgfsetmiterjoin%
\definecolor{currentfill}{rgb}{0.000000,0.000000,0.000000}%
\pgfsetfillcolor{currentfill}%
\pgfsetlinewidth{0.000000pt}%
\definecolor{currentstroke}{rgb}{0.000000,0.000000,0.000000}%
\pgfsetstrokecolor{currentstroke}%
\pgfsetstrokeopacity{0.000000}%
\pgfsetdash{}{0pt}%
\pgfpathmoveto{\pgfqpoint{4.850650in}{0.660000in}}%
\pgfpathlineto{\pgfqpoint{4.856850in}{0.660000in}}%
\pgfpathlineto{\pgfqpoint{4.856850in}{3.294268in}}%
\pgfpathlineto{\pgfqpoint{4.850650in}{3.294268in}}%
\pgfpathlineto{\pgfqpoint{4.850650in}{0.660000in}}%
\pgfpathclose%
\pgfusepath{fill}%
\end{pgfscope}%
\begin{pgfscope}%
\pgfpathrectangle{\pgfqpoint{1.250000in}{0.660000in}}{\pgfqpoint{7.750000in}{4.620000in}}%
\pgfusepath{clip}%
\pgfsetbuttcap%
\pgfsetmiterjoin%
\definecolor{currentfill}{rgb}{0.000000,0.000000,0.000000}%
\pgfsetfillcolor{currentfill}%
\pgfsetlinewidth{0.000000pt}%
\definecolor{currentstroke}{rgb}{0.000000,0.000000,0.000000}%
\pgfsetstrokecolor{currentstroke}%
\pgfsetstrokeopacity{0.000000}%
\pgfsetdash{}{0pt}%
\pgfpathmoveto{\pgfqpoint{4.858400in}{0.660000in}}%
\pgfpathlineto{\pgfqpoint{4.864600in}{0.660000in}}%
\pgfpathlineto{\pgfqpoint{4.864600in}{3.277232in}}%
\pgfpathlineto{\pgfqpoint{4.858400in}{3.277232in}}%
\pgfpathlineto{\pgfqpoint{4.858400in}{0.660000in}}%
\pgfpathclose%
\pgfusepath{fill}%
\end{pgfscope}%
\begin{pgfscope}%
\pgfpathrectangle{\pgfqpoint{1.250000in}{0.660000in}}{\pgfqpoint{7.750000in}{4.620000in}}%
\pgfusepath{clip}%
\pgfsetbuttcap%
\pgfsetmiterjoin%
\definecolor{currentfill}{rgb}{0.000000,0.000000,0.000000}%
\pgfsetfillcolor{currentfill}%
\pgfsetlinewidth{0.000000pt}%
\definecolor{currentstroke}{rgb}{0.000000,0.000000,0.000000}%
\pgfsetstrokecolor{currentstroke}%
\pgfsetstrokeopacity{0.000000}%
\pgfsetdash{}{0pt}%
\pgfpathmoveto{\pgfqpoint{4.866150in}{0.660000in}}%
\pgfpathlineto{\pgfqpoint{4.872350in}{0.660000in}}%
\pgfpathlineto{\pgfqpoint{4.872350in}{3.261312in}}%
\pgfpathlineto{\pgfqpoint{4.866150in}{3.261312in}}%
\pgfpathlineto{\pgfqpoint{4.866150in}{0.660000in}}%
\pgfpathclose%
\pgfusepath{fill}%
\end{pgfscope}%
\begin{pgfscope}%
\pgfpathrectangle{\pgfqpoint{1.250000in}{0.660000in}}{\pgfqpoint{7.750000in}{4.620000in}}%
\pgfusepath{clip}%
\pgfsetbuttcap%
\pgfsetmiterjoin%
\definecolor{currentfill}{rgb}{0.000000,0.000000,0.000000}%
\pgfsetfillcolor{currentfill}%
\pgfsetlinewidth{0.000000pt}%
\definecolor{currentstroke}{rgb}{0.000000,0.000000,0.000000}%
\pgfsetstrokecolor{currentstroke}%
\pgfsetstrokeopacity{0.000000}%
\pgfsetdash{}{0pt}%
\pgfpathmoveto{\pgfqpoint{4.873900in}{0.660000in}}%
\pgfpathlineto{\pgfqpoint{4.880100in}{0.660000in}}%
\pgfpathlineto{\pgfqpoint{4.880100in}{3.260163in}}%
\pgfpathlineto{\pgfqpoint{4.873900in}{3.260163in}}%
\pgfpathlineto{\pgfqpoint{4.873900in}{0.660000in}}%
\pgfpathclose%
\pgfusepath{fill}%
\end{pgfscope}%
\begin{pgfscope}%
\pgfpathrectangle{\pgfqpoint{1.250000in}{0.660000in}}{\pgfqpoint{7.750000in}{4.620000in}}%
\pgfusepath{clip}%
\pgfsetbuttcap%
\pgfsetmiterjoin%
\definecolor{currentfill}{rgb}{0.000000,0.000000,0.000000}%
\pgfsetfillcolor{currentfill}%
\pgfsetlinewidth{0.000000pt}%
\definecolor{currentstroke}{rgb}{0.000000,0.000000,0.000000}%
\pgfsetstrokecolor{currentstroke}%
\pgfsetstrokeopacity{0.000000}%
\pgfsetdash{}{0pt}%
\pgfpathmoveto{\pgfqpoint{4.881650in}{0.660000in}}%
\pgfpathlineto{\pgfqpoint{4.887850in}{0.660000in}}%
\pgfpathlineto{\pgfqpoint{4.887850in}{3.292628in}}%
\pgfpathlineto{\pgfqpoint{4.881650in}{3.292628in}}%
\pgfpathlineto{\pgfqpoint{4.881650in}{0.660000in}}%
\pgfpathclose%
\pgfusepath{fill}%
\end{pgfscope}%
\begin{pgfscope}%
\pgfpathrectangle{\pgfqpoint{1.250000in}{0.660000in}}{\pgfqpoint{7.750000in}{4.620000in}}%
\pgfusepath{clip}%
\pgfsetbuttcap%
\pgfsetmiterjoin%
\definecolor{currentfill}{rgb}{0.000000,0.000000,0.000000}%
\pgfsetfillcolor{currentfill}%
\pgfsetlinewidth{0.000000pt}%
\definecolor{currentstroke}{rgb}{0.000000,0.000000,0.000000}%
\pgfsetstrokecolor{currentstroke}%
\pgfsetstrokeopacity{0.000000}%
\pgfsetdash{}{0pt}%
\pgfpathmoveto{\pgfqpoint{4.889400in}{0.660000in}}%
\pgfpathlineto{\pgfqpoint{4.895600in}{0.660000in}}%
\pgfpathlineto{\pgfqpoint{4.895600in}{3.278201in}}%
\pgfpathlineto{\pgfqpoint{4.889400in}{3.278201in}}%
\pgfpathlineto{\pgfqpoint{4.889400in}{0.660000in}}%
\pgfpathclose%
\pgfusepath{fill}%
\end{pgfscope}%
\begin{pgfscope}%
\pgfpathrectangle{\pgfqpoint{1.250000in}{0.660000in}}{\pgfqpoint{7.750000in}{4.620000in}}%
\pgfusepath{clip}%
\pgfsetbuttcap%
\pgfsetmiterjoin%
\definecolor{currentfill}{rgb}{0.000000,0.000000,0.000000}%
\pgfsetfillcolor{currentfill}%
\pgfsetlinewidth{0.000000pt}%
\definecolor{currentstroke}{rgb}{0.000000,0.000000,0.000000}%
\pgfsetstrokecolor{currentstroke}%
\pgfsetstrokeopacity{0.000000}%
\pgfsetdash{}{0pt}%
\pgfpathmoveto{\pgfqpoint{4.897150in}{0.660000in}}%
\pgfpathlineto{\pgfqpoint{4.903350in}{0.660000in}}%
\pgfpathlineto{\pgfqpoint{4.903350in}{3.267474in}}%
\pgfpathlineto{\pgfqpoint{4.897150in}{3.267474in}}%
\pgfpathlineto{\pgfqpoint{4.897150in}{0.660000in}}%
\pgfpathclose%
\pgfusepath{fill}%
\end{pgfscope}%
\begin{pgfscope}%
\pgfpathrectangle{\pgfqpoint{1.250000in}{0.660000in}}{\pgfqpoint{7.750000in}{4.620000in}}%
\pgfusepath{clip}%
\pgfsetbuttcap%
\pgfsetmiterjoin%
\definecolor{currentfill}{rgb}{0.000000,0.000000,0.000000}%
\pgfsetfillcolor{currentfill}%
\pgfsetlinewidth{0.000000pt}%
\definecolor{currentstroke}{rgb}{0.000000,0.000000,0.000000}%
\pgfsetstrokecolor{currentstroke}%
\pgfsetstrokeopacity{0.000000}%
\pgfsetdash{}{0pt}%
\pgfpathmoveto{\pgfqpoint{4.904900in}{0.660000in}}%
\pgfpathlineto{\pgfqpoint{4.911100in}{0.660000in}}%
\pgfpathlineto{\pgfqpoint{4.911100in}{3.259103in}}%
\pgfpathlineto{\pgfqpoint{4.904900in}{3.259103in}}%
\pgfpathlineto{\pgfqpoint{4.904900in}{0.660000in}}%
\pgfpathclose%
\pgfusepath{fill}%
\end{pgfscope}%
\begin{pgfscope}%
\pgfpathrectangle{\pgfqpoint{1.250000in}{0.660000in}}{\pgfqpoint{7.750000in}{4.620000in}}%
\pgfusepath{clip}%
\pgfsetbuttcap%
\pgfsetmiterjoin%
\definecolor{currentfill}{rgb}{0.000000,0.000000,0.000000}%
\pgfsetfillcolor{currentfill}%
\pgfsetlinewidth{0.000000pt}%
\definecolor{currentstroke}{rgb}{0.000000,0.000000,0.000000}%
\pgfsetstrokecolor{currentstroke}%
\pgfsetstrokeopacity{0.000000}%
\pgfsetdash{}{0pt}%
\pgfpathmoveto{\pgfqpoint{4.912650in}{0.660000in}}%
\pgfpathlineto{\pgfqpoint{4.918850in}{0.660000in}}%
\pgfpathlineto{\pgfqpoint{4.918850in}{3.259379in}}%
\pgfpathlineto{\pgfqpoint{4.912650in}{3.259379in}}%
\pgfpathlineto{\pgfqpoint{4.912650in}{0.660000in}}%
\pgfpathclose%
\pgfusepath{fill}%
\end{pgfscope}%
\begin{pgfscope}%
\pgfpathrectangle{\pgfqpoint{1.250000in}{0.660000in}}{\pgfqpoint{7.750000in}{4.620000in}}%
\pgfusepath{clip}%
\pgfsetbuttcap%
\pgfsetmiterjoin%
\definecolor{currentfill}{rgb}{0.000000,0.000000,0.000000}%
\pgfsetfillcolor{currentfill}%
\pgfsetlinewidth{0.000000pt}%
\definecolor{currentstroke}{rgb}{0.000000,0.000000,0.000000}%
\pgfsetstrokecolor{currentstroke}%
\pgfsetstrokeopacity{0.000000}%
\pgfsetdash{}{0pt}%
\pgfpathmoveto{\pgfqpoint{4.920400in}{0.660000in}}%
\pgfpathlineto{\pgfqpoint{4.926600in}{0.660000in}}%
\pgfpathlineto{\pgfqpoint{4.926600in}{3.269708in}}%
\pgfpathlineto{\pgfqpoint{4.920400in}{3.269708in}}%
\pgfpathlineto{\pgfqpoint{4.920400in}{0.660000in}}%
\pgfpathclose%
\pgfusepath{fill}%
\end{pgfscope}%
\begin{pgfscope}%
\pgfpathrectangle{\pgfqpoint{1.250000in}{0.660000in}}{\pgfqpoint{7.750000in}{4.620000in}}%
\pgfusepath{clip}%
\pgfsetbuttcap%
\pgfsetmiterjoin%
\definecolor{currentfill}{rgb}{0.000000,0.000000,0.000000}%
\pgfsetfillcolor{currentfill}%
\pgfsetlinewidth{0.000000pt}%
\definecolor{currentstroke}{rgb}{0.000000,0.000000,0.000000}%
\pgfsetstrokecolor{currentstroke}%
\pgfsetstrokeopacity{0.000000}%
\pgfsetdash{}{0pt}%
\pgfpathmoveto{\pgfqpoint{4.928150in}{0.660000in}}%
\pgfpathlineto{\pgfqpoint{4.934350in}{0.660000in}}%
\pgfpathlineto{\pgfqpoint{4.934350in}{3.270999in}}%
\pgfpathlineto{\pgfqpoint{4.928150in}{3.270999in}}%
\pgfpathlineto{\pgfqpoint{4.928150in}{0.660000in}}%
\pgfpathclose%
\pgfusepath{fill}%
\end{pgfscope}%
\begin{pgfscope}%
\pgfpathrectangle{\pgfqpoint{1.250000in}{0.660000in}}{\pgfqpoint{7.750000in}{4.620000in}}%
\pgfusepath{clip}%
\pgfsetbuttcap%
\pgfsetmiterjoin%
\definecolor{currentfill}{rgb}{0.000000,0.000000,0.000000}%
\pgfsetfillcolor{currentfill}%
\pgfsetlinewidth{0.000000pt}%
\definecolor{currentstroke}{rgb}{0.000000,0.000000,0.000000}%
\pgfsetstrokecolor{currentstroke}%
\pgfsetstrokeopacity{0.000000}%
\pgfsetdash{}{0pt}%
\pgfpathmoveto{\pgfqpoint{4.935900in}{0.660000in}}%
\pgfpathlineto{\pgfqpoint{4.942100in}{0.660000in}}%
\pgfpathlineto{\pgfqpoint{4.942100in}{3.290887in}}%
\pgfpathlineto{\pgfqpoint{4.935900in}{3.290887in}}%
\pgfpathlineto{\pgfqpoint{4.935900in}{0.660000in}}%
\pgfpathclose%
\pgfusepath{fill}%
\end{pgfscope}%
\begin{pgfscope}%
\pgfpathrectangle{\pgfqpoint{1.250000in}{0.660000in}}{\pgfqpoint{7.750000in}{4.620000in}}%
\pgfusepath{clip}%
\pgfsetbuttcap%
\pgfsetmiterjoin%
\definecolor{currentfill}{rgb}{0.000000,0.000000,0.000000}%
\pgfsetfillcolor{currentfill}%
\pgfsetlinewidth{0.000000pt}%
\definecolor{currentstroke}{rgb}{0.000000,0.000000,0.000000}%
\pgfsetstrokecolor{currentstroke}%
\pgfsetstrokeopacity{0.000000}%
\pgfsetdash{}{0pt}%
\pgfpathmoveto{\pgfqpoint{4.943650in}{0.660000in}}%
\pgfpathlineto{\pgfqpoint{4.949850in}{0.660000in}}%
\pgfpathlineto{\pgfqpoint{4.949850in}{3.263773in}}%
\pgfpathlineto{\pgfqpoint{4.943650in}{3.263773in}}%
\pgfpathlineto{\pgfqpoint{4.943650in}{0.660000in}}%
\pgfpathclose%
\pgfusepath{fill}%
\end{pgfscope}%
\begin{pgfscope}%
\pgfpathrectangle{\pgfqpoint{1.250000in}{0.660000in}}{\pgfqpoint{7.750000in}{4.620000in}}%
\pgfusepath{clip}%
\pgfsetbuttcap%
\pgfsetmiterjoin%
\definecolor{currentfill}{rgb}{0.000000,0.000000,0.000000}%
\pgfsetfillcolor{currentfill}%
\pgfsetlinewidth{0.000000pt}%
\definecolor{currentstroke}{rgb}{0.000000,0.000000,0.000000}%
\pgfsetstrokecolor{currentstroke}%
\pgfsetstrokeopacity{0.000000}%
\pgfsetdash{}{0pt}%
\pgfpathmoveto{\pgfqpoint{4.951400in}{0.660000in}}%
\pgfpathlineto{\pgfqpoint{4.957600in}{0.660000in}}%
\pgfpathlineto{\pgfqpoint{4.957600in}{3.307896in}}%
\pgfpathlineto{\pgfqpoint{4.951400in}{3.307896in}}%
\pgfpathlineto{\pgfqpoint{4.951400in}{0.660000in}}%
\pgfpathclose%
\pgfusepath{fill}%
\end{pgfscope}%
\begin{pgfscope}%
\pgfpathrectangle{\pgfqpoint{1.250000in}{0.660000in}}{\pgfqpoint{7.750000in}{4.620000in}}%
\pgfusepath{clip}%
\pgfsetbuttcap%
\pgfsetmiterjoin%
\definecolor{currentfill}{rgb}{0.000000,0.000000,0.000000}%
\pgfsetfillcolor{currentfill}%
\pgfsetlinewidth{0.000000pt}%
\definecolor{currentstroke}{rgb}{0.000000,0.000000,0.000000}%
\pgfsetstrokecolor{currentstroke}%
\pgfsetstrokeopacity{0.000000}%
\pgfsetdash{}{0pt}%
\pgfpathmoveto{\pgfqpoint{4.959150in}{0.660000in}}%
\pgfpathlineto{\pgfqpoint{4.965350in}{0.660000in}}%
\pgfpathlineto{\pgfqpoint{4.965350in}{3.257381in}}%
\pgfpathlineto{\pgfqpoint{4.959150in}{3.257381in}}%
\pgfpathlineto{\pgfqpoint{4.959150in}{0.660000in}}%
\pgfpathclose%
\pgfusepath{fill}%
\end{pgfscope}%
\begin{pgfscope}%
\pgfpathrectangle{\pgfqpoint{1.250000in}{0.660000in}}{\pgfqpoint{7.750000in}{4.620000in}}%
\pgfusepath{clip}%
\pgfsetbuttcap%
\pgfsetmiterjoin%
\definecolor{currentfill}{rgb}{0.000000,0.000000,0.000000}%
\pgfsetfillcolor{currentfill}%
\pgfsetlinewidth{0.000000pt}%
\definecolor{currentstroke}{rgb}{0.000000,0.000000,0.000000}%
\pgfsetstrokecolor{currentstroke}%
\pgfsetstrokeopacity{0.000000}%
\pgfsetdash{}{0pt}%
\pgfpathmoveto{\pgfqpoint{4.966900in}{0.660000in}}%
\pgfpathlineto{\pgfqpoint{4.973100in}{0.660000in}}%
\pgfpathlineto{\pgfqpoint{4.973100in}{3.259234in}}%
\pgfpathlineto{\pgfqpoint{4.966900in}{3.259234in}}%
\pgfpathlineto{\pgfqpoint{4.966900in}{0.660000in}}%
\pgfpathclose%
\pgfusepath{fill}%
\end{pgfscope}%
\begin{pgfscope}%
\pgfpathrectangle{\pgfqpoint{1.250000in}{0.660000in}}{\pgfqpoint{7.750000in}{4.620000in}}%
\pgfusepath{clip}%
\pgfsetbuttcap%
\pgfsetmiterjoin%
\definecolor{currentfill}{rgb}{0.000000,0.000000,0.000000}%
\pgfsetfillcolor{currentfill}%
\pgfsetlinewidth{0.000000pt}%
\definecolor{currentstroke}{rgb}{0.000000,0.000000,0.000000}%
\pgfsetstrokecolor{currentstroke}%
\pgfsetstrokeopacity{0.000000}%
\pgfsetdash{}{0pt}%
\pgfpathmoveto{\pgfqpoint{4.974650in}{0.660000in}}%
\pgfpathlineto{\pgfqpoint{4.980850in}{0.660000in}}%
\pgfpathlineto{\pgfqpoint{4.980850in}{3.265226in}}%
\pgfpathlineto{\pgfqpoint{4.974650in}{3.265226in}}%
\pgfpathlineto{\pgfqpoint{4.974650in}{0.660000in}}%
\pgfpathclose%
\pgfusepath{fill}%
\end{pgfscope}%
\begin{pgfscope}%
\pgfpathrectangle{\pgfqpoint{1.250000in}{0.660000in}}{\pgfqpoint{7.750000in}{4.620000in}}%
\pgfusepath{clip}%
\pgfsetbuttcap%
\pgfsetmiterjoin%
\definecolor{currentfill}{rgb}{0.000000,0.000000,0.000000}%
\pgfsetfillcolor{currentfill}%
\pgfsetlinewidth{0.000000pt}%
\definecolor{currentstroke}{rgb}{0.000000,0.000000,0.000000}%
\pgfsetstrokecolor{currentstroke}%
\pgfsetstrokeopacity{0.000000}%
\pgfsetdash{}{0pt}%
\pgfpathmoveto{\pgfqpoint{4.982400in}{0.660000in}}%
\pgfpathlineto{\pgfqpoint{4.988600in}{0.660000in}}%
\pgfpathlineto{\pgfqpoint{4.988600in}{3.257013in}}%
\pgfpathlineto{\pgfqpoint{4.982400in}{3.257013in}}%
\pgfpathlineto{\pgfqpoint{4.982400in}{0.660000in}}%
\pgfpathclose%
\pgfusepath{fill}%
\end{pgfscope}%
\begin{pgfscope}%
\pgfpathrectangle{\pgfqpoint{1.250000in}{0.660000in}}{\pgfqpoint{7.750000in}{4.620000in}}%
\pgfusepath{clip}%
\pgfsetbuttcap%
\pgfsetmiterjoin%
\definecolor{currentfill}{rgb}{0.000000,0.000000,0.000000}%
\pgfsetfillcolor{currentfill}%
\pgfsetlinewidth{0.000000pt}%
\definecolor{currentstroke}{rgb}{0.000000,0.000000,0.000000}%
\pgfsetstrokecolor{currentstroke}%
\pgfsetstrokeopacity{0.000000}%
\pgfsetdash{}{0pt}%
\pgfpathmoveto{\pgfqpoint{4.990150in}{0.660000in}}%
\pgfpathlineto{\pgfqpoint{4.996350in}{0.660000in}}%
\pgfpathlineto{\pgfqpoint{4.996350in}{3.275653in}}%
\pgfpathlineto{\pgfqpoint{4.990150in}{3.275653in}}%
\pgfpathlineto{\pgfqpoint{4.990150in}{0.660000in}}%
\pgfpathclose%
\pgfusepath{fill}%
\end{pgfscope}%
\begin{pgfscope}%
\pgfpathrectangle{\pgfqpoint{1.250000in}{0.660000in}}{\pgfqpoint{7.750000in}{4.620000in}}%
\pgfusepath{clip}%
\pgfsetbuttcap%
\pgfsetmiterjoin%
\definecolor{currentfill}{rgb}{0.000000,0.000000,0.000000}%
\pgfsetfillcolor{currentfill}%
\pgfsetlinewidth{0.000000pt}%
\definecolor{currentstroke}{rgb}{0.000000,0.000000,0.000000}%
\pgfsetstrokecolor{currentstroke}%
\pgfsetstrokeopacity{0.000000}%
\pgfsetdash{}{0pt}%
\pgfpathmoveto{\pgfqpoint{4.997900in}{0.660000in}}%
\pgfpathlineto{\pgfqpoint{5.004100in}{0.660000in}}%
\pgfpathlineto{\pgfqpoint{5.004100in}{3.279186in}}%
\pgfpathlineto{\pgfqpoint{4.997900in}{3.279186in}}%
\pgfpathlineto{\pgfqpoint{4.997900in}{0.660000in}}%
\pgfpathclose%
\pgfusepath{fill}%
\end{pgfscope}%
\begin{pgfscope}%
\pgfpathrectangle{\pgfqpoint{1.250000in}{0.660000in}}{\pgfqpoint{7.750000in}{4.620000in}}%
\pgfusepath{clip}%
\pgfsetbuttcap%
\pgfsetmiterjoin%
\definecolor{currentfill}{rgb}{0.000000,0.000000,0.000000}%
\pgfsetfillcolor{currentfill}%
\pgfsetlinewidth{0.000000pt}%
\definecolor{currentstroke}{rgb}{0.000000,0.000000,0.000000}%
\pgfsetstrokecolor{currentstroke}%
\pgfsetstrokeopacity{0.000000}%
\pgfsetdash{}{0pt}%
\pgfpathmoveto{\pgfqpoint{5.005650in}{0.660000in}}%
\pgfpathlineto{\pgfqpoint{5.011850in}{0.660000in}}%
\pgfpathlineto{\pgfqpoint{5.011850in}{3.256480in}}%
\pgfpathlineto{\pgfqpoint{5.005650in}{3.256480in}}%
\pgfpathlineto{\pgfqpoint{5.005650in}{0.660000in}}%
\pgfpathclose%
\pgfusepath{fill}%
\end{pgfscope}%
\begin{pgfscope}%
\pgfpathrectangle{\pgfqpoint{1.250000in}{0.660000in}}{\pgfqpoint{7.750000in}{4.620000in}}%
\pgfusepath{clip}%
\pgfsetbuttcap%
\pgfsetmiterjoin%
\definecolor{currentfill}{rgb}{0.000000,0.000000,0.000000}%
\pgfsetfillcolor{currentfill}%
\pgfsetlinewidth{0.000000pt}%
\definecolor{currentstroke}{rgb}{0.000000,0.000000,0.000000}%
\pgfsetstrokecolor{currentstroke}%
\pgfsetstrokeopacity{0.000000}%
\pgfsetdash{}{0pt}%
\pgfpathmoveto{\pgfqpoint{5.013400in}{0.660000in}}%
\pgfpathlineto{\pgfqpoint{5.019600in}{0.660000in}}%
\pgfpathlineto{\pgfqpoint{5.019600in}{3.256268in}}%
\pgfpathlineto{\pgfqpoint{5.013400in}{3.256268in}}%
\pgfpathlineto{\pgfqpoint{5.013400in}{0.660000in}}%
\pgfpathclose%
\pgfusepath{fill}%
\end{pgfscope}%
\begin{pgfscope}%
\pgfpathrectangle{\pgfqpoint{1.250000in}{0.660000in}}{\pgfqpoint{7.750000in}{4.620000in}}%
\pgfusepath{clip}%
\pgfsetbuttcap%
\pgfsetmiterjoin%
\definecolor{currentfill}{rgb}{0.000000,0.000000,0.000000}%
\pgfsetfillcolor{currentfill}%
\pgfsetlinewidth{0.000000pt}%
\definecolor{currentstroke}{rgb}{0.000000,0.000000,0.000000}%
\pgfsetstrokecolor{currentstroke}%
\pgfsetstrokeopacity{0.000000}%
\pgfsetdash{}{0pt}%
\pgfpathmoveto{\pgfqpoint{5.021150in}{0.660000in}}%
\pgfpathlineto{\pgfqpoint{5.027350in}{0.660000in}}%
\pgfpathlineto{\pgfqpoint{5.027350in}{3.256040in}}%
\pgfpathlineto{\pgfqpoint{5.021150in}{3.256040in}}%
\pgfpathlineto{\pgfqpoint{5.021150in}{0.660000in}}%
\pgfpathclose%
\pgfusepath{fill}%
\end{pgfscope}%
\begin{pgfscope}%
\pgfpathrectangle{\pgfqpoint{1.250000in}{0.660000in}}{\pgfqpoint{7.750000in}{4.620000in}}%
\pgfusepath{clip}%
\pgfsetbuttcap%
\pgfsetmiterjoin%
\definecolor{currentfill}{rgb}{0.000000,0.000000,0.000000}%
\pgfsetfillcolor{currentfill}%
\pgfsetlinewidth{0.000000pt}%
\definecolor{currentstroke}{rgb}{0.000000,0.000000,0.000000}%
\pgfsetstrokecolor{currentstroke}%
\pgfsetstrokeopacity{0.000000}%
\pgfsetdash{}{0pt}%
\pgfpathmoveto{\pgfqpoint{5.028900in}{0.660000in}}%
\pgfpathlineto{\pgfqpoint{5.035100in}{0.660000in}}%
\pgfpathlineto{\pgfqpoint{5.035100in}{3.255796in}}%
\pgfpathlineto{\pgfqpoint{5.028900in}{3.255796in}}%
\pgfpathlineto{\pgfqpoint{5.028900in}{0.660000in}}%
\pgfpathclose%
\pgfusepath{fill}%
\end{pgfscope}%
\begin{pgfscope}%
\pgfpathrectangle{\pgfqpoint{1.250000in}{0.660000in}}{\pgfqpoint{7.750000in}{4.620000in}}%
\pgfusepath{clip}%
\pgfsetbuttcap%
\pgfsetmiterjoin%
\definecolor{currentfill}{rgb}{0.000000,0.000000,0.000000}%
\pgfsetfillcolor{currentfill}%
\pgfsetlinewidth{0.000000pt}%
\definecolor{currentstroke}{rgb}{0.000000,0.000000,0.000000}%
\pgfsetstrokecolor{currentstroke}%
\pgfsetstrokeopacity{0.000000}%
\pgfsetdash{}{0pt}%
\pgfpathmoveto{\pgfqpoint{5.036650in}{0.660000in}}%
\pgfpathlineto{\pgfqpoint{5.042850in}{0.660000in}}%
\pgfpathlineto{\pgfqpoint{5.042850in}{3.255537in}}%
\pgfpathlineto{\pgfqpoint{5.036650in}{3.255537in}}%
\pgfpathlineto{\pgfqpoint{5.036650in}{0.660000in}}%
\pgfpathclose%
\pgfusepath{fill}%
\end{pgfscope}%
\begin{pgfscope}%
\pgfpathrectangle{\pgfqpoint{1.250000in}{0.660000in}}{\pgfqpoint{7.750000in}{4.620000in}}%
\pgfusepath{clip}%
\pgfsetbuttcap%
\pgfsetmiterjoin%
\definecolor{currentfill}{rgb}{0.000000,0.000000,0.000000}%
\pgfsetfillcolor{currentfill}%
\pgfsetlinewidth{0.000000pt}%
\definecolor{currentstroke}{rgb}{0.000000,0.000000,0.000000}%
\pgfsetstrokecolor{currentstroke}%
\pgfsetstrokeopacity{0.000000}%
\pgfsetdash{}{0pt}%
\pgfpathmoveto{\pgfqpoint{5.044400in}{0.660000in}}%
\pgfpathlineto{\pgfqpoint{5.050600in}{0.660000in}}%
\pgfpathlineto{\pgfqpoint{5.050600in}{3.302886in}}%
\pgfpathlineto{\pgfqpoint{5.044400in}{3.302886in}}%
\pgfpathlineto{\pgfqpoint{5.044400in}{0.660000in}}%
\pgfpathclose%
\pgfusepath{fill}%
\end{pgfscope}%
\begin{pgfscope}%
\pgfpathrectangle{\pgfqpoint{1.250000in}{0.660000in}}{\pgfqpoint{7.750000in}{4.620000in}}%
\pgfusepath{clip}%
\pgfsetbuttcap%
\pgfsetmiterjoin%
\definecolor{currentfill}{rgb}{0.000000,0.000000,0.000000}%
\pgfsetfillcolor{currentfill}%
\pgfsetlinewidth{0.000000pt}%
\definecolor{currentstroke}{rgb}{0.000000,0.000000,0.000000}%
\pgfsetstrokecolor{currentstroke}%
\pgfsetstrokeopacity{0.000000}%
\pgfsetdash{}{0pt}%
\pgfpathmoveto{\pgfqpoint{5.052150in}{0.660000in}}%
\pgfpathlineto{\pgfqpoint{5.058350in}{0.660000in}}%
\pgfpathlineto{\pgfqpoint{5.058350in}{3.302533in}}%
\pgfpathlineto{\pgfqpoint{5.052150in}{3.302533in}}%
\pgfpathlineto{\pgfqpoint{5.052150in}{0.660000in}}%
\pgfpathclose%
\pgfusepath{fill}%
\end{pgfscope}%
\begin{pgfscope}%
\pgfpathrectangle{\pgfqpoint{1.250000in}{0.660000in}}{\pgfqpoint{7.750000in}{4.620000in}}%
\pgfusepath{clip}%
\pgfsetbuttcap%
\pgfsetmiterjoin%
\definecolor{currentfill}{rgb}{0.000000,0.000000,0.000000}%
\pgfsetfillcolor{currentfill}%
\pgfsetlinewidth{0.000000pt}%
\definecolor{currentstroke}{rgb}{0.000000,0.000000,0.000000}%
\pgfsetstrokecolor{currentstroke}%
\pgfsetstrokeopacity{0.000000}%
\pgfsetdash{}{0pt}%
\pgfpathmoveto{\pgfqpoint{5.059900in}{0.660000in}}%
\pgfpathlineto{\pgfqpoint{5.066100in}{0.660000in}}%
\pgfpathlineto{\pgfqpoint{5.066100in}{3.274690in}}%
\pgfpathlineto{\pgfqpoint{5.059900in}{3.274690in}}%
\pgfpathlineto{\pgfqpoint{5.059900in}{0.660000in}}%
\pgfpathclose%
\pgfusepath{fill}%
\end{pgfscope}%
\begin{pgfscope}%
\pgfpathrectangle{\pgfqpoint{1.250000in}{0.660000in}}{\pgfqpoint{7.750000in}{4.620000in}}%
\pgfusepath{clip}%
\pgfsetbuttcap%
\pgfsetmiterjoin%
\definecolor{currentfill}{rgb}{0.000000,0.000000,0.000000}%
\pgfsetfillcolor{currentfill}%
\pgfsetlinewidth{0.000000pt}%
\definecolor{currentstroke}{rgb}{0.000000,0.000000,0.000000}%
\pgfsetstrokecolor{currentstroke}%
\pgfsetstrokeopacity{0.000000}%
\pgfsetdash{}{0pt}%
\pgfpathmoveto{\pgfqpoint{5.067650in}{0.660000in}}%
\pgfpathlineto{\pgfqpoint{5.073850in}{0.660000in}}%
\pgfpathlineto{\pgfqpoint{5.073850in}{3.254360in}}%
\pgfpathlineto{\pgfqpoint{5.067650in}{3.254360in}}%
\pgfpathlineto{\pgfqpoint{5.067650in}{0.660000in}}%
\pgfpathclose%
\pgfusepath{fill}%
\end{pgfscope}%
\begin{pgfscope}%
\pgfpathrectangle{\pgfqpoint{1.250000in}{0.660000in}}{\pgfqpoint{7.750000in}{4.620000in}}%
\pgfusepath{clip}%
\pgfsetbuttcap%
\pgfsetmiterjoin%
\definecolor{currentfill}{rgb}{0.000000,0.000000,0.000000}%
\pgfsetfillcolor{currentfill}%
\pgfsetlinewidth{0.000000pt}%
\definecolor{currentstroke}{rgb}{0.000000,0.000000,0.000000}%
\pgfsetstrokecolor{currentstroke}%
\pgfsetstrokeopacity{0.000000}%
\pgfsetdash{}{0pt}%
\pgfpathmoveto{\pgfqpoint{5.075400in}{0.660000in}}%
\pgfpathlineto{\pgfqpoint{5.081600in}{0.660000in}}%
\pgfpathlineto{\pgfqpoint{5.081600in}{3.254033in}}%
\pgfpathlineto{\pgfqpoint{5.075400in}{3.254033in}}%
\pgfpathlineto{\pgfqpoint{5.075400in}{0.660000in}}%
\pgfpathclose%
\pgfusepath{fill}%
\end{pgfscope}%
\begin{pgfscope}%
\pgfpathrectangle{\pgfqpoint{1.250000in}{0.660000in}}{\pgfqpoint{7.750000in}{4.620000in}}%
\pgfusepath{clip}%
\pgfsetbuttcap%
\pgfsetmiterjoin%
\definecolor{currentfill}{rgb}{0.000000,0.000000,0.000000}%
\pgfsetfillcolor{currentfill}%
\pgfsetlinewidth{0.000000pt}%
\definecolor{currentstroke}{rgb}{0.000000,0.000000,0.000000}%
\pgfsetstrokecolor{currentstroke}%
\pgfsetstrokeopacity{0.000000}%
\pgfsetdash{}{0pt}%
\pgfpathmoveto{\pgfqpoint{5.083150in}{0.660000in}}%
\pgfpathlineto{\pgfqpoint{5.089350in}{0.660000in}}%
\pgfpathlineto{\pgfqpoint{5.089350in}{3.253694in}}%
\pgfpathlineto{\pgfqpoint{5.083150in}{3.253694in}}%
\pgfpathlineto{\pgfqpoint{5.083150in}{0.660000in}}%
\pgfpathclose%
\pgfusepath{fill}%
\end{pgfscope}%
\begin{pgfscope}%
\pgfpathrectangle{\pgfqpoint{1.250000in}{0.660000in}}{\pgfqpoint{7.750000in}{4.620000in}}%
\pgfusepath{clip}%
\pgfsetbuttcap%
\pgfsetmiterjoin%
\definecolor{currentfill}{rgb}{0.000000,0.000000,0.000000}%
\pgfsetfillcolor{currentfill}%
\pgfsetlinewidth{0.000000pt}%
\definecolor{currentstroke}{rgb}{0.000000,0.000000,0.000000}%
\pgfsetstrokecolor{currentstroke}%
\pgfsetstrokeopacity{0.000000}%
\pgfsetdash{}{0pt}%
\pgfpathmoveto{\pgfqpoint{5.090900in}{0.660000in}}%
\pgfpathlineto{\pgfqpoint{5.097100in}{0.660000in}}%
\pgfpathlineto{\pgfqpoint{5.097100in}{3.256740in}}%
\pgfpathlineto{\pgfqpoint{5.090900in}{3.256740in}}%
\pgfpathlineto{\pgfqpoint{5.090900in}{0.660000in}}%
\pgfpathclose%
\pgfusepath{fill}%
\end{pgfscope}%
\begin{pgfscope}%
\pgfpathrectangle{\pgfqpoint{1.250000in}{0.660000in}}{\pgfqpoint{7.750000in}{4.620000in}}%
\pgfusepath{clip}%
\pgfsetbuttcap%
\pgfsetmiterjoin%
\definecolor{currentfill}{rgb}{0.000000,0.000000,0.000000}%
\pgfsetfillcolor{currentfill}%
\pgfsetlinewidth{0.000000pt}%
\definecolor{currentstroke}{rgb}{0.000000,0.000000,0.000000}%
\pgfsetstrokecolor{currentstroke}%
\pgfsetstrokeopacity{0.000000}%
\pgfsetdash{}{0pt}%
\pgfpathmoveto{\pgfqpoint{5.098650in}{0.660000in}}%
\pgfpathlineto{\pgfqpoint{5.104850in}{0.660000in}}%
\pgfpathlineto{\pgfqpoint{5.104850in}{3.277757in}}%
\pgfpathlineto{\pgfqpoint{5.098650in}{3.277757in}}%
\pgfpathlineto{\pgfqpoint{5.098650in}{0.660000in}}%
\pgfpathclose%
\pgfusepath{fill}%
\end{pgfscope}%
\begin{pgfscope}%
\pgfpathrectangle{\pgfqpoint{1.250000in}{0.660000in}}{\pgfqpoint{7.750000in}{4.620000in}}%
\pgfusepath{clip}%
\pgfsetbuttcap%
\pgfsetmiterjoin%
\definecolor{currentfill}{rgb}{0.000000,0.000000,0.000000}%
\pgfsetfillcolor{currentfill}%
\pgfsetlinewidth{0.000000pt}%
\definecolor{currentstroke}{rgb}{0.000000,0.000000,0.000000}%
\pgfsetstrokecolor{currentstroke}%
\pgfsetstrokeopacity{0.000000}%
\pgfsetdash{}{0pt}%
\pgfpathmoveto{\pgfqpoint{5.106400in}{0.660000in}}%
\pgfpathlineto{\pgfqpoint{5.112600in}{0.660000in}}%
\pgfpathlineto{\pgfqpoint{5.112600in}{3.255545in}}%
\pgfpathlineto{\pgfqpoint{5.106400in}{3.255545in}}%
\pgfpathlineto{\pgfqpoint{5.106400in}{0.660000in}}%
\pgfpathclose%
\pgfusepath{fill}%
\end{pgfscope}%
\begin{pgfscope}%
\pgfpathrectangle{\pgfqpoint{1.250000in}{0.660000in}}{\pgfqpoint{7.750000in}{4.620000in}}%
\pgfusepath{clip}%
\pgfsetbuttcap%
\pgfsetmiterjoin%
\definecolor{currentfill}{rgb}{0.000000,0.000000,0.000000}%
\pgfsetfillcolor{currentfill}%
\pgfsetlinewidth{0.000000pt}%
\definecolor{currentstroke}{rgb}{0.000000,0.000000,0.000000}%
\pgfsetstrokecolor{currentstroke}%
\pgfsetstrokeopacity{0.000000}%
\pgfsetdash{}{0pt}%
\pgfpathmoveto{\pgfqpoint{5.114150in}{0.660000in}}%
\pgfpathlineto{\pgfqpoint{5.120350in}{0.660000in}}%
\pgfpathlineto{\pgfqpoint{5.120350in}{3.253311in}}%
\pgfpathlineto{\pgfqpoint{5.114150in}{3.253311in}}%
\pgfpathlineto{\pgfqpoint{5.114150in}{0.660000in}}%
\pgfpathclose%
\pgfusepath{fill}%
\end{pgfscope}%
\begin{pgfscope}%
\pgfpathrectangle{\pgfqpoint{1.250000in}{0.660000in}}{\pgfqpoint{7.750000in}{4.620000in}}%
\pgfusepath{clip}%
\pgfsetbuttcap%
\pgfsetmiterjoin%
\definecolor{currentfill}{rgb}{0.000000,0.000000,0.000000}%
\pgfsetfillcolor{currentfill}%
\pgfsetlinewidth{0.000000pt}%
\definecolor{currentstroke}{rgb}{0.000000,0.000000,0.000000}%
\pgfsetstrokecolor{currentstroke}%
\pgfsetstrokeopacity{0.000000}%
\pgfsetdash{}{0pt}%
\pgfpathmoveto{\pgfqpoint{5.121900in}{0.660000in}}%
\pgfpathlineto{\pgfqpoint{5.128100in}{0.660000in}}%
\pgfpathlineto{\pgfqpoint{5.128100in}{3.273788in}}%
\pgfpathlineto{\pgfqpoint{5.121900in}{3.273788in}}%
\pgfpathlineto{\pgfqpoint{5.121900in}{0.660000in}}%
\pgfpathclose%
\pgfusepath{fill}%
\end{pgfscope}%
\begin{pgfscope}%
\pgfpathrectangle{\pgfqpoint{1.250000in}{0.660000in}}{\pgfqpoint{7.750000in}{4.620000in}}%
\pgfusepath{clip}%
\pgfsetbuttcap%
\pgfsetmiterjoin%
\definecolor{currentfill}{rgb}{0.000000,0.000000,0.000000}%
\pgfsetfillcolor{currentfill}%
\pgfsetlinewidth{0.000000pt}%
\definecolor{currentstroke}{rgb}{0.000000,0.000000,0.000000}%
\pgfsetstrokecolor{currentstroke}%
\pgfsetstrokeopacity{0.000000}%
\pgfsetdash{}{0pt}%
\pgfpathmoveto{\pgfqpoint{5.129650in}{0.660000in}}%
\pgfpathlineto{\pgfqpoint{5.135850in}{0.660000in}}%
\pgfpathlineto{\pgfqpoint{5.135850in}{3.253051in}}%
\pgfpathlineto{\pgfqpoint{5.129650in}{3.253051in}}%
\pgfpathlineto{\pgfqpoint{5.129650in}{0.660000in}}%
\pgfpathclose%
\pgfusepath{fill}%
\end{pgfscope}%
\begin{pgfscope}%
\pgfpathrectangle{\pgfqpoint{1.250000in}{0.660000in}}{\pgfqpoint{7.750000in}{4.620000in}}%
\pgfusepath{clip}%
\pgfsetbuttcap%
\pgfsetmiterjoin%
\definecolor{currentfill}{rgb}{0.000000,0.000000,0.000000}%
\pgfsetfillcolor{currentfill}%
\pgfsetlinewidth{0.000000pt}%
\definecolor{currentstroke}{rgb}{0.000000,0.000000,0.000000}%
\pgfsetstrokecolor{currentstroke}%
\pgfsetstrokeopacity{0.000000}%
\pgfsetdash{}{0pt}%
\pgfpathmoveto{\pgfqpoint{5.137400in}{0.660000in}}%
\pgfpathlineto{\pgfqpoint{5.143600in}{0.660000in}}%
\pgfpathlineto{\pgfqpoint{5.143600in}{3.252900in}}%
\pgfpathlineto{\pgfqpoint{5.137400in}{3.252900in}}%
\pgfpathlineto{\pgfqpoint{5.137400in}{0.660000in}}%
\pgfpathclose%
\pgfusepath{fill}%
\end{pgfscope}%
\begin{pgfscope}%
\pgfpathrectangle{\pgfqpoint{1.250000in}{0.660000in}}{\pgfqpoint{7.750000in}{4.620000in}}%
\pgfusepath{clip}%
\pgfsetbuttcap%
\pgfsetmiterjoin%
\definecolor{currentfill}{rgb}{0.000000,0.000000,0.000000}%
\pgfsetfillcolor{currentfill}%
\pgfsetlinewidth{0.000000pt}%
\definecolor{currentstroke}{rgb}{0.000000,0.000000,0.000000}%
\pgfsetstrokecolor{currentstroke}%
\pgfsetstrokeopacity{0.000000}%
\pgfsetdash{}{0pt}%
\pgfpathmoveto{\pgfqpoint{5.145150in}{0.660000in}}%
\pgfpathlineto{\pgfqpoint{5.151350in}{0.660000in}}%
\pgfpathlineto{\pgfqpoint{5.151350in}{3.263917in}}%
\pgfpathlineto{\pgfqpoint{5.145150in}{3.263917in}}%
\pgfpathlineto{\pgfqpoint{5.145150in}{0.660000in}}%
\pgfpathclose%
\pgfusepath{fill}%
\end{pgfscope}%
\begin{pgfscope}%
\pgfpathrectangle{\pgfqpoint{1.250000in}{0.660000in}}{\pgfqpoint{7.750000in}{4.620000in}}%
\pgfusepath{clip}%
\pgfsetbuttcap%
\pgfsetmiterjoin%
\definecolor{currentfill}{rgb}{0.000000,0.000000,0.000000}%
\pgfsetfillcolor{currentfill}%
\pgfsetlinewidth{0.000000pt}%
\definecolor{currentstroke}{rgb}{0.000000,0.000000,0.000000}%
\pgfsetstrokecolor{currentstroke}%
\pgfsetstrokeopacity{0.000000}%
\pgfsetdash{}{0pt}%
\pgfpathmoveto{\pgfqpoint{5.152900in}{0.660000in}}%
\pgfpathlineto{\pgfqpoint{5.159100in}{0.660000in}}%
\pgfpathlineto{\pgfqpoint{5.159100in}{3.319151in}}%
\pgfpathlineto{\pgfqpoint{5.152900in}{3.319151in}}%
\pgfpathlineto{\pgfqpoint{5.152900in}{0.660000in}}%
\pgfpathclose%
\pgfusepath{fill}%
\end{pgfscope}%
\begin{pgfscope}%
\pgfpathrectangle{\pgfqpoint{1.250000in}{0.660000in}}{\pgfqpoint{7.750000in}{4.620000in}}%
\pgfusepath{clip}%
\pgfsetbuttcap%
\pgfsetmiterjoin%
\definecolor{currentfill}{rgb}{0.000000,0.000000,0.000000}%
\pgfsetfillcolor{currentfill}%
\pgfsetlinewidth{0.000000pt}%
\definecolor{currentstroke}{rgb}{0.000000,0.000000,0.000000}%
\pgfsetstrokecolor{currentstroke}%
\pgfsetstrokeopacity{0.000000}%
\pgfsetdash{}{0pt}%
\pgfpathmoveto{\pgfqpoint{5.160650in}{0.660000in}}%
\pgfpathlineto{\pgfqpoint{5.166850in}{0.660000in}}%
\pgfpathlineto{\pgfqpoint{5.166850in}{3.283754in}}%
\pgfpathlineto{\pgfqpoint{5.160650in}{3.283754in}}%
\pgfpathlineto{\pgfqpoint{5.160650in}{0.660000in}}%
\pgfpathclose%
\pgfusepath{fill}%
\end{pgfscope}%
\begin{pgfscope}%
\pgfpathrectangle{\pgfqpoint{1.250000in}{0.660000in}}{\pgfqpoint{7.750000in}{4.620000in}}%
\pgfusepath{clip}%
\pgfsetbuttcap%
\pgfsetmiterjoin%
\definecolor{currentfill}{rgb}{0.000000,0.000000,0.000000}%
\pgfsetfillcolor{currentfill}%
\pgfsetlinewidth{0.000000pt}%
\definecolor{currentstroke}{rgb}{0.000000,0.000000,0.000000}%
\pgfsetstrokecolor{currentstroke}%
\pgfsetstrokeopacity{0.000000}%
\pgfsetdash{}{0pt}%
\pgfpathmoveto{\pgfqpoint{5.168400in}{0.660000in}}%
\pgfpathlineto{\pgfqpoint{5.174600in}{0.660000in}}%
\pgfpathlineto{\pgfqpoint{5.174600in}{3.274963in}}%
\pgfpathlineto{\pgfqpoint{5.168400in}{3.274963in}}%
\pgfpathlineto{\pgfqpoint{5.168400in}{0.660000in}}%
\pgfpathclose%
\pgfusepath{fill}%
\end{pgfscope}%
\begin{pgfscope}%
\pgfpathrectangle{\pgfqpoint{1.250000in}{0.660000in}}{\pgfqpoint{7.750000in}{4.620000in}}%
\pgfusepath{clip}%
\pgfsetbuttcap%
\pgfsetmiterjoin%
\definecolor{currentfill}{rgb}{0.000000,0.000000,0.000000}%
\pgfsetfillcolor{currentfill}%
\pgfsetlinewidth{0.000000pt}%
\definecolor{currentstroke}{rgb}{0.000000,0.000000,0.000000}%
\pgfsetstrokecolor{currentstroke}%
\pgfsetstrokeopacity{0.000000}%
\pgfsetdash{}{0pt}%
\pgfpathmoveto{\pgfqpoint{5.176150in}{0.660000in}}%
\pgfpathlineto{\pgfqpoint{5.182350in}{0.660000in}}%
\pgfpathlineto{\pgfqpoint{5.182350in}{3.251957in}}%
\pgfpathlineto{\pgfqpoint{5.176150in}{3.251957in}}%
\pgfpathlineto{\pgfqpoint{5.176150in}{0.660000in}}%
\pgfpathclose%
\pgfusepath{fill}%
\end{pgfscope}%
\begin{pgfscope}%
\pgfpathrectangle{\pgfqpoint{1.250000in}{0.660000in}}{\pgfqpoint{7.750000in}{4.620000in}}%
\pgfusepath{clip}%
\pgfsetbuttcap%
\pgfsetmiterjoin%
\definecolor{currentfill}{rgb}{0.000000,0.000000,0.000000}%
\pgfsetfillcolor{currentfill}%
\pgfsetlinewidth{0.000000pt}%
\definecolor{currentstroke}{rgb}{0.000000,0.000000,0.000000}%
\pgfsetstrokecolor{currentstroke}%
\pgfsetstrokeopacity{0.000000}%
\pgfsetdash{}{0pt}%
\pgfpathmoveto{\pgfqpoint{5.183900in}{0.660000in}}%
\pgfpathlineto{\pgfqpoint{5.190100in}{0.660000in}}%
\pgfpathlineto{\pgfqpoint{5.190100in}{3.273706in}}%
\pgfpathlineto{\pgfqpoint{5.183900in}{3.273706in}}%
\pgfpathlineto{\pgfqpoint{5.183900in}{0.660000in}}%
\pgfpathclose%
\pgfusepath{fill}%
\end{pgfscope}%
\begin{pgfscope}%
\pgfpathrectangle{\pgfqpoint{1.250000in}{0.660000in}}{\pgfqpoint{7.750000in}{4.620000in}}%
\pgfusepath{clip}%
\pgfsetbuttcap%
\pgfsetmiterjoin%
\definecolor{currentfill}{rgb}{0.000000,0.000000,0.000000}%
\pgfsetfillcolor{currentfill}%
\pgfsetlinewidth{0.000000pt}%
\definecolor{currentstroke}{rgb}{0.000000,0.000000,0.000000}%
\pgfsetstrokecolor{currentstroke}%
\pgfsetstrokeopacity{0.000000}%
\pgfsetdash{}{0pt}%
\pgfpathmoveto{\pgfqpoint{5.191650in}{0.660000in}}%
\pgfpathlineto{\pgfqpoint{5.197850in}{0.660000in}}%
\pgfpathlineto{\pgfqpoint{5.197850in}{3.251498in}}%
\pgfpathlineto{\pgfqpoint{5.191650in}{3.251498in}}%
\pgfpathlineto{\pgfqpoint{5.191650in}{0.660000in}}%
\pgfpathclose%
\pgfusepath{fill}%
\end{pgfscope}%
\begin{pgfscope}%
\pgfpathrectangle{\pgfqpoint{1.250000in}{0.660000in}}{\pgfqpoint{7.750000in}{4.620000in}}%
\pgfusepath{clip}%
\pgfsetbuttcap%
\pgfsetmiterjoin%
\definecolor{currentfill}{rgb}{0.000000,0.000000,0.000000}%
\pgfsetfillcolor{currentfill}%
\pgfsetlinewidth{0.000000pt}%
\definecolor{currentstroke}{rgb}{0.000000,0.000000,0.000000}%
\pgfsetstrokecolor{currentstroke}%
\pgfsetstrokeopacity{0.000000}%
\pgfsetdash{}{0pt}%
\pgfpathmoveto{\pgfqpoint{5.199400in}{0.660000in}}%
\pgfpathlineto{\pgfqpoint{5.205600in}{0.660000in}}%
\pgfpathlineto{\pgfqpoint{5.205600in}{3.253408in}}%
\pgfpathlineto{\pgfqpoint{5.199400in}{3.253408in}}%
\pgfpathlineto{\pgfqpoint{5.199400in}{0.660000in}}%
\pgfpathclose%
\pgfusepath{fill}%
\end{pgfscope}%
\begin{pgfscope}%
\pgfpathrectangle{\pgfqpoint{1.250000in}{0.660000in}}{\pgfqpoint{7.750000in}{4.620000in}}%
\pgfusepath{clip}%
\pgfsetbuttcap%
\pgfsetmiterjoin%
\definecolor{currentfill}{rgb}{0.000000,0.000000,0.000000}%
\pgfsetfillcolor{currentfill}%
\pgfsetlinewidth{0.000000pt}%
\definecolor{currentstroke}{rgb}{0.000000,0.000000,0.000000}%
\pgfsetstrokecolor{currentstroke}%
\pgfsetstrokeopacity{0.000000}%
\pgfsetdash{}{0pt}%
\pgfpathmoveto{\pgfqpoint{5.207150in}{0.660000in}}%
\pgfpathlineto{\pgfqpoint{5.213350in}{0.660000in}}%
\pgfpathlineto{\pgfqpoint{5.213350in}{3.250999in}}%
\pgfpathlineto{\pgfqpoint{5.207150in}{3.250999in}}%
\pgfpathlineto{\pgfqpoint{5.207150in}{0.660000in}}%
\pgfpathclose%
\pgfusepath{fill}%
\end{pgfscope}%
\begin{pgfscope}%
\pgfpathrectangle{\pgfqpoint{1.250000in}{0.660000in}}{\pgfqpoint{7.750000in}{4.620000in}}%
\pgfusepath{clip}%
\pgfsetbuttcap%
\pgfsetmiterjoin%
\definecolor{currentfill}{rgb}{0.000000,0.000000,0.000000}%
\pgfsetfillcolor{currentfill}%
\pgfsetlinewidth{0.000000pt}%
\definecolor{currentstroke}{rgb}{0.000000,0.000000,0.000000}%
\pgfsetstrokecolor{currentstroke}%
\pgfsetstrokeopacity{0.000000}%
\pgfsetdash{}{0pt}%
\pgfpathmoveto{\pgfqpoint{5.214900in}{0.660000in}}%
\pgfpathlineto{\pgfqpoint{5.221100in}{0.660000in}}%
\pgfpathlineto{\pgfqpoint{5.221100in}{3.296244in}}%
\pgfpathlineto{\pgfqpoint{5.214900in}{3.296244in}}%
\pgfpathlineto{\pgfqpoint{5.214900in}{0.660000in}}%
\pgfpathclose%
\pgfusepath{fill}%
\end{pgfscope}%
\begin{pgfscope}%
\pgfpathrectangle{\pgfqpoint{1.250000in}{0.660000in}}{\pgfqpoint{7.750000in}{4.620000in}}%
\pgfusepath{clip}%
\pgfsetbuttcap%
\pgfsetmiterjoin%
\definecolor{currentfill}{rgb}{0.000000,0.000000,0.000000}%
\pgfsetfillcolor{currentfill}%
\pgfsetlinewidth{0.000000pt}%
\definecolor{currentstroke}{rgb}{0.000000,0.000000,0.000000}%
\pgfsetstrokecolor{currentstroke}%
\pgfsetstrokeopacity{0.000000}%
\pgfsetdash{}{0pt}%
\pgfpathmoveto{\pgfqpoint{5.222650in}{0.660000in}}%
\pgfpathlineto{\pgfqpoint{5.228850in}{0.660000in}}%
\pgfpathlineto{\pgfqpoint{5.228850in}{3.250461in}}%
\pgfpathlineto{\pgfqpoint{5.222650in}{3.250461in}}%
\pgfpathlineto{\pgfqpoint{5.222650in}{0.660000in}}%
\pgfpathclose%
\pgfusepath{fill}%
\end{pgfscope}%
\begin{pgfscope}%
\pgfpathrectangle{\pgfqpoint{1.250000in}{0.660000in}}{\pgfqpoint{7.750000in}{4.620000in}}%
\pgfusepath{clip}%
\pgfsetbuttcap%
\pgfsetmiterjoin%
\definecolor{currentfill}{rgb}{0.000000,0.000000,0.000000}%
\pgfsetfillcolor{currentfill}%
\pgfsetlinewidth{0.000000pt}%
\definecolor{currentstroke}{rgb}{0.000000,0.000000,0.000000}%
\pgfsetstrokecolor{currentstroke}%
\pgfsetstrokeopacity{0.000000}%
\pgfsetdash{}{0pt}%
\pgfpathmoveto{\pgfqpoint{5.230400in}{0.660000in}}%
\pgfpathlineto{\pgfqpoint{5.236600in}{0.660000in}}%
\pgfpathlineto{\pgfqpoint{5.236600in}{3.250179in}}%
\pgfpathlineto{\pgfqpoint{5.230400in}{3.250179in}}%
\pgfpathlineto{\pgfqpoint{5.230400in}{0.660000in}}%
\pgfpathclose%
\pgfusepath{fill}%
\end{pgfscope}%
\begin{pgfscope}%
\pgfpathrectangle{\pgfqpoint{1.250000in}{0.660000in}}{\pgfqpoint{7.750000in}{4.620000in}}%
\pgfusepath{clip}%
\pgfsetbuttcap%
\pgfsetmiterjoin%
\definecolor{currentfill}{rgb}{0.000000,0.000000,0.000000}%
\pgfsetfillcolor{currentfill}%
\pgfsetlinewidth{0.000000pt}%
\definecolor{currentstroke}{rgb}{0.000000,0.000000,0.000000}%
\pgfsetstrokecolor{currentstroke}%
\pgfsetstrokeopacity{0.000000}%
\pgfsetdash{}{0pt}%
\pgfpathmoveto{\pgfqpoint{5.238150in}{0.660000in}}%
\pgfpathlineto{\pgfqpoint{5.244350in}{0.660000in}}%
\pgfpathlineto{\pgfqpoint{5.244350in}{3.266652in}}%
\pgfpathlineto{\pgfqpoint{5.238150in}{3.266652in}}%
\pgfpathlineto{\pgfqpoint{5.238150in}{0.660000in}}%
\pgfpathclose%
\pgfusepath{fill}%
\end{pgfscope}%
\begin{pgfscope}%
\pgfpathrectangle{\pgfqpoint{1.250000in}{0.660000in}}{\pgfqpoint{7.750000in}{4.620000in}}%
\pgfusepath{clip}%
\pgfsetbuttcap%
\pgfsetmiterjoin%
\definecolor{currentfill}{rgb}{0.000000,0.000000,0.000000}%
\pgfsetfillcolor{currentfill}%
\pgfsetlinewidth{0.000000pt}%
\definecolor{currentstroke}{rgb}{0.000000,0.000000,0.000000}%
\pgfsetstrokecolor{currentstroke}%
\pgfsetstrokeopacity{0.000000}%
\pgfsetdash{}{0pt}%
\pgfpathmoveto{\pgfqpoint{5.245900in}{0.660000in}}%
\pgfpathlineto{\pgfqpoint{5.252100in}{0.660000in}}%
\pgfpathlineto{\pgfqpoint{5.252100in}{3.263207in}}%
\pgfpathlineto{\pgfqpoint{5.245900in}{3.263207in}}%
\pgfpathlineto{\pgfqpoint{5.245900in}{0.660000in}}%
\pgfpathclose%
\pgfusepath{fill}%
\end{pgfscope}%
\begin{pgfscope}%
\pgfpathrectangle{\pgfqpoint{1.250000in}{0.660000in}}{\pgfqpoint{7.750000in}{4.620000in}}%
\pgfusepath{clip}%
\pgfsetbuttcap%
\pgfsetmiterjoin%
\definecolor{currentfill}{rgb}{0.000000,0.000000,0.000000}%
\pgfsetfillcolor{currentfill}%
\pgfsetlinewidth{0.000000pt}%
\definecolor{currentstroke}{rgb}{0.000000,0.000000,0.000000}%
\pgfsetstrokecolor{currentstroke}%
\pgfsetstrokeopacity{0.000000}%
\pgfsetdash{}{0pt}%
\pgfpathmoveto{\pgfqpoint{5.253650in}{0.660000in}}%
\pgfpathlineto{\pgfqpoint{5.259850in}{0.660000in}}%
\pgfpathlineto{\pgfqpoint{5.259850in}{3.306733in}}%
\pgfpathlineto{\pgfqpoint{5.253650in}{3.306733in}}%
\pgfpathlineto{\pgfqpoint{5.253650in}{0.660000in}}%
\pgfpathclose%
\pgfusepath{fill}%
\end{pgfscope}%
\begin{pgfscope}%
\pgfpathrectangle{\pgfqpoint{1.250000in}{0.660000in}}{\pgfqpoint{7.750000in}{4.620000in}}%
\pgfusepath{clip}%
\pgfsetbuttcap%
\pgfsetmiterjoin%
\definecolor{currentfill}{rgb}{0.000000,0.000000,0.000000}%
\pgfsetfillcolor{currentfill}%
\pgfsetlinewidth{0.000000pt}%
\definecolor{currentstroke}{rgb}{0.000000,0.000000,0.000000}%
\pgfsetstrokecolor{currentstroke}%
\pgfsetstrokeopacity{0.000000}%
\pgfsetdash{}{0pt}%
\pgfpathmoveto{\pgfqpoint{5.261400in}{0.660000in}}%
\pgfpathlineto{\pgfqpoint{5.267600in}{0.660000in}}%
\pgfpathlineto{\pgfqpoint{5.267600in}{3.249263in}}%
\pgfpathlineto{\pgfqpoint{5.261400in}{3.249263in}}%
\pgfpathlineto{\pgfqpoint{5.261400in}{0.660000in}}%
\pgfpathclose%
\pgfusepath{fill}%
\end{pgfscope}%
\begin{pgfscope}%
\pgfpathrectangle{\pgfqpoint{1.250000in}{0.660000in}}{\pgfqpoint{7.750000in}{4.620000in}}%
\pgfusepath{clip}%
\pgfsetbuttcap%
\pgfsetmiterjoin%
\definecolor{currentfill}{rgb}{0.000000,0.000000,0.000000}%
\pgfsetfillcolor{currentfill}%
\pgfsetlinewidth{0.000000pt}%
\definecolor{currentstroke}{rgb}{0.000000,0.000000,0.000000}%
\pgfsetstrokecolor{currentstroke}%
\pgfsetstrokeopacity{0.000000}%
\pgfsetdash{}{0pt}%
\pgfpathmoveto{\pgfqpoint{5.269150in}{0.660000in}}%
\pgfpathlineto{\pgfqpoint{5.275350in}{0.660000in}}%
\pgfpathlineto{\pgfqpoint{5.275350in}{3.249532in}}%
\pgfpathlineto{\pgfqpoint{5.269150in}{3.249532in}}%
\pgfpathlineto{\pgfqpoint{5.269150in}{0.660000in}}%
\pgfpathclose%
\pgfusepath{fill}%
\end{pgfscope}%
\begin{pgfscope}%
\pgfpathrectangle{\pgfqpoint{1.250000in}{0.660000in}}{\pgfqpoint{7.750000in}{4.620000in}}%
\pgfusepath{clip}%
\pgfsetbuttcap%
\pgfsetmiterjoin%
\definecolor{currentfill}{rgb}{0.000000,0.000000,0.000000}%
\pgfsetfillcolor{currentfill}%
\pgfsetlinewidth{0.000000pt}%
\definecolor{currentstroke}{rgb}{0.000000,0.000000,0.000000}%
\pgfsetstrokecolor{currentstroke}%
\pgfsetstrokeopacity{0.000000}%
\pgfsetdash{}{0pt}%
\pgfpathmoveto{\pgfqpoint{5.276900in}{0.660000in}}%
\pgfpathlineto{\pgfqpoint{5.283100in}{0.660000in}}%
\pgfpathlineto{\pgfqpoint{5.283100in}{3.252113in}}%
\pgfpathlineto{\pgfqpoint{5.276900in}{3.252113in}}%
\pgfpathlineto{\pgfqpoint{5.276900in}{0.660000in}}%
\pgfpathclose%
\pgfusepath{fill}%
\end{pgfscope}%
\begin{pgfscope}%
\pgfpathrectangle{\pgfqpoint{1.250000in}{0.660000in}}{\pgfqpoint{7.750000in}{4.620000in}}%
\pgfusepath{clip}%
\pgfsetbuttcap%
\pgfsetmiterjoin%
\definecolor{currentfill}{rgb}{0.000000,0.000000,0.000000}%
\pgfsetfillcolor{currentfill}%
\pgfsetlinewidth{0.000000pt}%
\definecolor{currentstroke}{rgb}{0.000000,0.000000,0.000000}%
\pgfsetstrokecolor{currentstroke}%
\pgfsetstrokeopacity{0.000000}%
\pgfsetdash{}{0pt}%
\pgfpathmoveto{\pgfqpoint{5.284650in}{0.660000in}}%
\pgfpathlineto{\pgfqpoint{5.290850in}{0.660000in}}%
\pgfpathlineto{\pgfqpoint{5.290850in}{3.256608in}}%
\pgfpathlineto{\pgfqpoint{5.284650in}{3.256608in}}%
\pgfpathlineto{\pgfqpoint{5.284650in}{0.660000in}}%
\pgfpathclose%
\pgfusepath{fill}%
\end{pgfscope}%
\begin{pgfscope}%
\pgfpathrectangle{\pgfqpoint{1.250000in}{0.660000in}}{\pgfqpoint{7.750000in}{4.620000in}}%
\pgfusepath{clip}%
\pgfsetbuttcap%
\pgfsetmiterjoin%
\definecolor{currentfill}{rgb}{0.000000,0.000000,0.000000}%
\pgfsetfillcolor{currentfill}%
\pgfsetlinewidth{0.000000pt}%
\definecolor{currentstroke}{rgb}{0.000000,0.000000,0.000000}%
\pgfsetstrokecolor{currentstroke}%
\pgfsetstrokeopacity{0.000000}%
\pgfsetdash{}{0pt}%
\pgfpathmoveto{\pgfqpoint{5.292400in}{0.660000in}}%
\pgfpathlineto{\pgfqpoint{5.298600in}{0.660000in}}%
\pgfpathlineto{\pgfqpoint{5.298600in}{3.248834in}}%
\pgfpathlineto{\pgfqpoint{5.292400in}{3.248834in}}%
\pgfpathlineto{\pgfqpoint{5.292400in}{0.660000in}}%
\pgfpathclose%
\pgfusepath{fill}%
\end{pgfscope}%
\begin{pgfscope}%
\pgfpathrectangle{\pgfqpoint{1.250000in}{0.660000in}}{\pgfqpoint{7.750000in}{4.620000in}}%
\pgfusepath{clip}%
\pgfsetbuttcap%
\pgfsetmiterjoin%
\definecolor{currentfill}{rgb}{0.000000,0.000000,0.000000}%
\pgfsetfillcolor{currentfill}%
\pgfsetlinewidth{0.000000pt}%
\definecolor{currentstroke}{rgb}{0.000000,0.000000,0.000000}%
\pgfsetstrokecolor{currentstroke}%
\pgfsetstrokeopacity{0.000000}%
\pgfsetdash{}{0pt}%
\pgfpathmoveto{\pgfqpoint{5.300150in}{0.660000in}}%
\pgfpathlineto{\pgfqpoint{5.306350in}{0.660000in}}%
\pgfpathlineto{\pgfqpoint{5.306350in}{3.248702in}}%
\pgfpathlineto{\pgfqpoint{5.300150in}{3.248702in}}%
\pgfpathlineto{\pgfqpoint{5.300150in}{0.660000in}}%
\pgfpathclose%
\pgfusepath{fill}%
\end{pgfscope}%
\begin{pgfscope}%
\pgfpathrectangle{\pgfqpoint{1.250000in}{0.660000in}}{\pgfqpoint{7.750000in}{4.620000in}}%
\pgfusepath{clip}%
\pgfsetbuttcap%
\pgfsetmiterjoin%
\definecolor{currentfill}{rgb}{0.000000,0.000000,0.000000}%
\pgfsetfillcolor{currentfill}%
\pgfsetlinewidth{0.000000pt}%
\definecolor{currentstroke}{rgb}{0.000000,0.000000,0.000000}%
\pgfsetstrokecolor{currentstroke}%
\pgfsetstrokeopacity{0.000000}%
\pgfsetdash{}{0pt}%
\pgfpathmoveto{\pgfqpoint{5.307900in}{0.660000in}}%
\pgfpathlineto{\pgfqpoint{5.314100in}{0.660000in}}%
\pgfpathlineto{\pgfqpoint{5.314100in}{3.248561in}}%
\pgfpathlineto{\pgfqpoint{5.307900in}{3.248561in}}%
\pgfpathlineto{\pgfqpoint{5.307900in}{0.660000in}}%
\pgfpathclose%
\pgfusepath{fill}%
\end{pgfscope}%
\begin{pgfscope}%
\pgfpathrectangle{\pgfqpoint{1.250000in}{0.660000in}}{\pgfqpoint{7.750000in}{4.620000in}}%
\pgfusepath{clip}%
\pgfsetbuttcap%
\pgfsetmiterjoin%
\definecolor{currentfill}{rgb}{0.000000,0.000000,0.000000}%
\pgfsetfillcolor{currentfill}%
\pgfsetlinewidth{0.000000pt}%
\definecolor{currentstroke}{rgb}{0.000000,0.000000,0.000000}%
\pgfsetstrokecolor{currentstroke}%
\pgfsetstrokeopacity{0.000000}%
\pgfsetdash{}{0pt}%
\pgfpathmoveto{\pgfqpoint{5.315650in}{0.660000in}}%
\pgfpathlineto{\pgfqpoint{5.321850in}{0.660000in}}%
\pgfpathlineto{\pgfqpoint{5.321850in}{3.248410in}}%
\pgfpathlineto{\pgfqpoint{5.315650in}{3.248410in}}%
\pgfpathlineto{\pgfqpoint{5.315650in}{0.660000in}}%
\pgfpathclose%
\pgfusepath{fill}%
\end{pgfscope}%
\begin{pgfscope}%
\pgfpathrectangle{\pgfqpoint{1.250000in}{0.660000in}}{\pgfqpoint{7.750000in}{4.620000in}}%
\pgfusepath{clip}%
\pgfsetbuttcap%
\pgfsetmiterjoin%
\definecolor{currentfill}{rgb}{0.000000,0.000000,0.000000}%
\pgfsetfillcolor{currentfill}%
\pgfsetlinewidth{0.000000pt}%
\definecolor{currentstroke}{rgb}{0.000000,0.000000,0.000000}%
\pgfsetstrokecolor{currentstroke}%
\pgfsetstrokeopacity{0.000000}%
\pgfsetdash{}{0pt}%
\pgfpathmoveto{\pgfqpoint{5.323400in}{0.660000in}}%
\pgfpathlineto{\pgfqpoint{5.329600in}{0.660000in}}%
\pgfpathlineto{\pgfqpoint{5.329600in}{3.248251in}}%
\pgfpathlineto{\pgfqpoint{5.323400in}{3.248251in}}%
\pgfpathlineto{\pgfqpoint{5.323400in}{0.660000in}}%
\pgfpathclose%
\pgfusepath{fill}%
\end{pgfscope}%
\begin{pgfscope}%
\pgfpathrectangle{\pgfqpoint{1.250000in}{0.660000in}}{\pgfqpoint{7.750000in}{4.620000in}}%
\pgfusepath{clip}%
\pgfsetbuttcap%
\pgfsetmiterjoin%
\definecolor{currentfill}{rgb}{0.000000,0.000000,0.000000}%
\pgfsetfillcolor{currentfill}%
\pgfsetlinewidth{0.000000pt}%
\definecolor{currentstroke}{rgb}{0.000000,0.000000,0.000000}%
\pgfsetstrokecolor{currentstroke}%
\pgfsetstrokeopacity{0.000000}%
\pgfsetdash{}{0pt}%
\pgfpathmoveto{\pgfqpoint{5.331150in}{0.660000in}}%
\pgfpathlineto{\pgfqpoint{5.337350in}{0.660000in}}%
\pgfpathlineto{\pgfqpoint{5.337350in}{3.248084in}}%
\pgfpathlineto{\pgfqpoint{5.331150in}{3.248084in}}%
\pgfpathlineto{\pgfqpoint{5.331150in}{0.660000in}}%
\pgfpathclose%
\pgfusepath{fill}%
\end{pgfscope}%
\begin{pgfscope}%
\pgfpathrectangle{\pgfqpoint{1.250000in}{0.660000in}}{\pgfqpoint{7.750000in}{4.620000in}}%
\pgfusepath{clip}%
\pgfsetbuttcap%
\pgfsetmiterjoin%
\definecolor{currentfill}{rgb}{0.000000,0.000000,0.000000}%
\pgfsetfillcolor{currentfill}%
\pgfsetlinewidth{0.000000pt}%
\definecolor{currentstroke}{rgb}{0.000000,0.000000,0.000000}%
\pgfsetstrokecolor{currentstroke}%
\pgfsetstrokeopacity{0.000000}%
\pgfsetdash{}{0pt}%
\pgfpathmoveto{\pgfqpoint{5.338900in}{0.660000in}}%
\pgfpathlineto{\pgfqpoint{5.345100in}{0.660000in}}%
\pgfpathlineto{\pgfqpoint{5.345100in}{3.247908in}}%
\pgfpathlineto{\pgfqpoint{5.338900in}{3.247908in}}%
\pgfpathlineto{\pgfqpoint{5.338900in}{0.660000in}}%
\pgfpathclose%
\pgfusepath{fill}%
\end{pgfscope}%
\begin{pgfscope}%
\pgfpathrectangle{\pgfqpoint{1.250000in}{0.660000in}}{\pgfqpoint{7.750000in}{4.620000in}}%
\pgfusepath{clip}%
\pgfsetbuttcap%
\pgfsetmiterjoin%
\definecolor{currentfill}{rgb}{0.000000,0.000000,0.000000}%
\pgfsetfillcolor{currentfill}%
\pgfsetlinewidth{0.000000pt}%
\definecolor{currentstroke}{rgb}{0.000000,0.000000,0.000000}%
\pgfsetstrokecolor{currentstroke}%
\pgfsetstrokeopacity{0.000000}%
\pgfsetdash{}{0pt}%
\pgfpathmoveto{\pgfqpoint{5.346650in}{0.660000in}}%
\pgfpathlineto{\pgfqpoint{5.352850in}{0.660000in}}%
\pgfpathlineto{\pgfqpoint{5.352850in}{3.256113in}}%
\pgfpathlineto{\pgfqpoint{5.346650in}{3.256113in}}%
\pgfpathlineto{\pgfqpoint{5.346650in}{0.660000in}}%
\pgfpathclose%
\pgfusepath{fill}%
\end{pgfscope}%
\begin{pgfscope}%
\pgfpathrectangle{\pgfqpoint{1.250000in}{0.660000in}}{\pgfqpoint{7.750000in}{4.620000in}}%
\pgfusepath{clip}%
\pgfsetbuttcap%
\pgfsetmiterjoin%
\definecolor{currentfill}{rgb}{0.000000,0.000000,0.000000}%
\pgfsetfillcolor{currentfill}%
\pgfsetlinewidth{0.000000pt}%
\definecolor{currentstroke}{rgb}{0.000000,0.000000,0.000000}%
\pgfsetstrokecolor{currentstroke}%
\pgfsetstrokeopacity{0.000000}%
\pgfsetdash{}{0pt}%
\pgfpathmoveto{\pgfqpoint{5.354400in}{0.660000in}}%
\pgfpathlineto{\pgfqpoint{5.360600in}{0.660000in}}%
\pgfpathlineto{\pgfqpoint{5.360600in}{3.247534in}}%
\pgfpathlineto{\pgfqpoint{5.354400in}{3.247534in}}%
\pgfpathlineto{\pgfqpoint{5.354400in}{0.660000in}}%
\pgfpathclose%
\pgfusepath{fill}%
\end{pgfscope}%
\begin{pgfscope}%
\pgfpathrectangle{\pgfqpoint{1.250000in}{0.660000in}}{\pgfqpoint{7.750000in}{4.620000in}}%
\pgfusepath{clip}%
\pgfsetbuttcap%
\pgfsetmiterjoin%
\definecolor{currentfill}{rgb}{0.000000,0.000000,0.000000}%
\pgfsetfillcolor{currentfill}%
\pgfsetlinewidth{0.000000pt}%
\definecolor{currentstroke}{rgb}{0.000000,0.000000,0.000000}%
\pgfsetstrokecolor{currentstroke}%
\pgfsetstrokeopacity{0.000000}%
\pgfsetdash{}{0pt}%
\pgfpathmoveto{\pgfqpoint{5.362150in}{0.660000in}}%
\pgfpathlineto{\pgfqpoint{5.368350in}{0.660000in}}%
\pgfpathlineto{\pgfqpoint{5.368350in}{3.251283in}}%
\pgfpathlineto{\pgfqpoint{5.362150in}{3.251283in}}%
\pgfpathlineto{\pgfqpoint{5.362150in}{0.660000in}}%
\pgfpathclose%
\pgfusepath{fill}%
\end{pgfscope}%
\begin{pgfscope}%
\pgfpathrectangle{\pgfqpoint{1.250000in}{0.660000in}}{\pgfqpoint{7.750000in}{4.620000in}}%
\pgfusepath{clip}%
\pgfsetbuttcap%
\pgfsetmiterjoin%
\definecolor{currentfill}{rgb}{0.000000,0.000000,0.000000}%
\pgfsetfillcolor{currentfill}%
\pgfsetlinewidth{0.000000pt}%
\definecolor{currentstroke}{rgb}{0.000000,0.000000,0.000000}%
\pgfsetstrokecolor{currentstroke}%
\pgfsetstrokeopacity{0.000000}%
\pgfsetdash{}{0pt}%
\pgfpathmoveto{\pgfqpoint{5.369900in}{0.660000in}}%
\pgfpathlineto{\pgfqpoint{5.376100in}{0.660000in}}%
\pgfpathlineto{\pgfqpoint{5.376100in}{3.253922in}}%
\pgfpathlineto{\pgfqpoint{5.369900in}{3.253922in}}%
\pgfpathlineto{\pgfqpoint{5.369900in}{0.660000in}}%
\pgfpathclose%
\pgfusepath{fill}%
\end{pgfscope}%
\begin{pgfscope}%
\pgfpathrectangle{\pgfqpoint{1.250000in}{0.660000in}}{\pgfqpoint{7.750000in}{4.620000in}}%
\pgfusepath{clip}%
\pgfsetbuttcap%
\pgfsetmiterjoin%
\definecolor{currentfill}{rgb}{0.000000,0.000000,0.000000}%
\pgfsetfillcolor{currentfill}%
\pgfsetlinewidth{0.000000pt}%
\definecolor{currentstroke}{rgb}{0.000000,0.000000,0.000000}%
\pgfsetstrokecolor{currentstroke}%
\pgfsetstrokeopacity{0.000000}%
\pgfsetdash{}{0pt}%
\pgfpathmoveto{\pgfqpoint{5.377650in}{0.660000in}}%
\pgfpathlineto{\pgfqpoint{5.383850in}{0.660000in}}%
\pgfpathlineto{\pgfqpoint{5.383850in}{3.263022in}}%
\pgfpathlineto{\pgfqpoint{5.377650in}{3.263022in}}%
\pgfpathlineto{\pgfqpoint{5.377650in}{0.660000in}}%
\pgfpathclose%
\pgfusepath{fill}%
\end{pgfscope}%
\begin{pgfscope}%
\pgfpathrectangle{\pgfqpoint{1.250000in}{0.660000in}}{\pgfqpoint{7.750000in}{4.620000in}}%
\pgfusepath{clip}%
\pgfsetbuttcap%
\pgfsetmiterjoin%
\definecolor{currentfill}{rgb}{0.000000,0.000000,0.000000}%
\pgfsetfillcolor{currentfill}%
\pgfsetlinewidth{0.000000pt}%
\definecolor{currentstroke}{rgb}{0.000000,0.000000,0.000000}%
\pgfsetstrokecolor{currentstroke}%
\pgfsetstrokeopacity{0.000000}%
\pgfsetdash{}{0pt}%
\pgfpathmoveto{\pgfqpoint{5.385400in}{0.660000in}}%
\pgfpathlineto{\pgfqpoint{5.391600in}{0.660000in}}%
\pgfpathlineto{\pgfqpoint{5.391600in}{3.253138in}}%
\pgfpathlineto{\pgfqpoint{5.385400in}{3.253138in}}%
\pgfpathlineto{\pgfqpoint{5.385400in}{0.660000in}}%
\pgfpathclose%
\pgfusepath{fill}%
\end{pgfscope}%
\begin{pgfscope}%
\pgfpathrectangle{\pgfqpoint{1.250000in}{0.660000in}}{\pgfqpoint{7.750000in}{4.620000in}}%
\pgfusepath{clip}%
\pgfsetbuttcap%
\pgfsetmiterjoin%
\definecolor{currentfill}{rgb}{0.000000,0.000000,0.000000}%
\pgfsetfillcolor{currentfill}%
\pgfsetlinewidth{0.000000pt}%
\definecolor{currentstroke}{rgb}{0.000000,0.000000,0.000000}%
\pgfsetstrokecolor{currentstroke}%
\pgfsetstrokeopacity{0.000000}%
\pgfsetdash{}{0pt}%
\pgfpathmoveto{\pgfqpoint{5.393150in}{0.660000in}}%
\pgfpathlineto{\pgfqpoint{5.399350in}{0.660000in}}%
\pgfpathlineto{\pgfqpoint{5.399350in}{3.248477in}}%
\pgfpathlineto{\pgfqpoint{5.393150in}{3.248477in}}%
\pgfpathlineto{\pgfqpoint{5.393150in}{0.660000in}}%
\pgfpathclose%
\pgfusepath{fill}%
\end{pgfscope}%
\begin{pgfscope}%
\pgfpathrectangle{\pgfqpoint{1.250000in}{0.660000in}}{\pgfqpoint{7.750000in}{4.620000in}}%
\pgfusepath{clip}%
\pgfsetbuttcap%
\pgfsetmiterjoin%
\definecolor{currentfill}{rgb}{0.000000,0.000000,0.000000}%
\pgfsetfillcolor{currentfill}%
\pgfsetlinewidth{0.000000pt}%
\definecolor{currentstroke}{rgb}{0.000000,0.000000,0.000000}%
\pgfsetstrokecolor{currentstroke}%
\pgfsetstrokeopacity{0.000000}%
\pgfsetdash{}{0pt}%
\pgfpathmoveto{\pgfqpoint{5.400900in}{0.660000in}}%
\pgfpathlineto{\pgfqpoint{5.407100in}{0.660000in}}%
\pgfpathlineto{\pgfqpoint{5.407100in}{3.246251in}}%
\pgfpathlineto{\pgfqpoint{5.400900in}{3.246251in}}%
\pgfpathlineto{\pgfqpoint{5.400900in}{0.660000in}}%
\pgfpathclose%
\pgfusepath{fill}%
\end{pgfscope}%
\begin{pgfscope}%
\pgfpathrectangle{\pgfqpoint{1.250000in}{0.660000in}}{\pgfqpoint{7.750000in}{4.620000in}}%
\pgfusepath{clip}%
\pgfsetbuttcap%
\pgfsetmiterjoin%
\definecolor{currentfill}{rgb}{0.000000,0.000000,0.000000}%
\pgfsetfillcolor{currentfill}%
\pgfsetlinewidth{0.000000pt}%
\definecolor{currentstroke}{rgb}{0.000000,0.000000,0.000000}%
\pgfsetstrokecolor{currentstroke}%
\pgfsetstrokeopacity{0.000000}%
\pgfsetdash{}{0pt}%
\pgfpathmoveto{\pgfqpoint{5.408650in}{0.660000in}}%
\pgfpathlineto{\pgfqpoint{5.414850in}{0.660000in}}%
\pgfpathlineto{\pgfqpoint{5.414850in}{3.276231in}}%
\pgfpathlineto{\pgfqpoint{5.408650in}{3.276231in}}%
\pgfpathlineto{\pgfqpoint{5.408650in}{0.660000in}}%
\pgfpathclose%
\pgfusepath{fill}%
\end{pgfscope}%
\begin{pgfscope}%
\pgfpathrectangle{\pgfqpoint{1.250000in}{0.660000in}}{\pgfqpoint{7.750000in}{4.620000in}}%
\pgfusepath{clip}%
\pgfsetbuttcap%
\pgfsetmiterjoin%
\definecolor{currentfill}{rgb}{0.000000,0.000000,0.000000}%
\pgfsetfillcolor{currentfill}%
\pgfsetlinewidth{0.000000pt}%
\definecolor{currentstroke}{rgb}{0.000000,0.000000,0.000000}%
\pgfsetstrokecolor{currentstroke}%
\pgfsetstrokeopacity{0.000000}%
\pgfsetdash{}{0pt}%
\pgfpathmoveto{\pgfqpoint{5.416400in}{0.660000in}}%
\pgfpathlineto{\pgfqpoint{5.422600in}{0.660000in}}%
\pgfpathlineto{\pgfqpoint{5.422600in}{3.308333in}}%
\pgfpathlineto{\pgfqpoint{5.416400in}{3.308333in}}%
\pgfpathlineto{\pgfqpoint{5.416400in}{0.660000in}}%
\pgfpathclose%
\pgfusepath{fill}%
\end{pgfscope}%
\begin{pgfscope}%
\pgfpathrectangle{\pgfqpoint{1.250000in}{0.660000in}}{\pgfqpoint{7.750000in}{4.620000in}}%
\pgfusepath{clip}%
\pgfsetbuttcap%
\pgfsetmiterjoin%
\definecolor{currentfill}{rgb}{0.000000,0.000000,0.000000}%
\pgfsetfillcolor{currentfill}%
\pgfsetlinewidth{0.000000pt}%
\definecolor{currentstroke}{rgb}{0.000000,0.000000,0.000000}%
\pgfsetstrokecolor{currentstroke}%
\pgfsetstrokeopacity{0.000000}%
\pgfsetdash{}{0pt}%
\pgfpathmoveto{\pgfqpoint{5.424150in}{0.660000in}}%
\pgfpathlineto{\pgfqpoint{5.430350in}{0.660000in}}%
\pgfpathlineto{\pgfqpoint{5.430350in}{3.292961in}}%
\pgfpathlineto{\pgfqpoint{5.424150in}{3.292961in}}%
\pgfpathlineto{\pgfqpoint{5.424150in}{0.660000in}}%
\pgfpathclose%
\pgfusepath{fill}%
\end{pgfscope}%
\begin{pgfscope}%
\pgfpathrectangle{\pgfqpoint{1.250000in}{0.660000in}}{\pgfqpoint{7.750000in}{4.620000in}}%
\pgfusepath{clip}%
\pgfsetbuttcap%
\pgfsetmiterjoin%
\definecolor{currentfill}{rgb}{0.000000,0.000000,0.000000}%
\pgfsetfillcolor{currentfill}%
\pgfsetlinewidth{0.000000pt}%
\definecolor{currentstroke}{rgb}{0.000000,0.000000,0.000000}%
\pgfsetstrokecolor{currentstroke}%
\pgfsetstrokeopacity{0.000000}%
\pgfsetdash{}{0pt}%
\pgfpathmoveto{\pgfqpoint{5.431900in}{0.660000in}}%
\pgfpathlineto{\pgfqpoint{5.438100in}{0.660000in}}%
\pgfpathlineto{\pgfqpoint{5.438100in}{3.245282in}}%
\pgfpathlineto{\pgfqpoint{5.431900in}{3.245282in}}%
\pgfpathlineto{\pgfqpoint{5.431900in}{0.660000in}}%
\pgfpathclose%
\pgfusepath{fill}%
\end{pgfscope}%
\begin{pgfscope}%
\pgfpathrectangle{\pgfqpoint{1.250000in}{0.660000in}}{\pgfqpoint{7.750000in}{4.620000in}}%
\pgfusepath{clip}%
\pgfsetbuttcap%
\pgfsetmiterjoin%
\definecolor{currentfill}{rgb}{0.000000,0.000000,0.000000}%
\pgfsetfillcolor{currentfill}%
\pgfsetlinewidth{0.000000pt}%
\definecolor{currentstroke}{rgb}{0.000000,0.000000,0.000000}%
\pgfsetstrokecolor{currentstroke}%
\pgfsetstrokeopacity{0.000000}%
\pgfsetdash{}{0pt}%
\pgfpathmoveto{\pgfqpoint{5.439650in}{0.660000in}}%
\pgfpathlineto{\pgfqpoint{5.445850in}{0.660000in}}%
\pgfpathlineto{\pgfqpoint{5.445850in}{3.252995in}}%
\pgfpathlineto{\pgfqpoint{5.439650in}{3.252995in}}%
\pgfpathlineto{\pgfqpoint{5.439650in}{0.660000in}}%
\pgfpathclose%
\pgfusepath{fill}%
\end{pgfscope}%
\begin{pgfscope}%
\pgfpathrectangle{\pgfqpoint{1.250000in}{0.660000in}}{\pgfqpoint{7.750000in}{4.620000in}}%
\pgfusepath{clip}%
\pgfsetbuttcap%
\pgfsetmiterjoin%
\definecolor{currentfill}{rgb}{0.000000,0.000000,0.000000}%
\pgfsetfillcolor{currentfill}%
\pgfsetlinewidth{0.000000pt}%
\definecolor{currentstroke}{rgb}{0.000000,0.000000,0.000000}%
\pgfsetstrokecolor{currentstroke}%
\pgfsetstrokeopacity{0.000000}%
\pgfsetdash{}{0pt}%
\pgfpathmoveto{\pgfqpoint{5.447400in}{0.660000in}}%
\pgfpathlineto{\pgfqpoint{5.453600in}{0.660000in}}%
\pgfpathlineto{\pgfqpoint{5.453600in}{3.244769in}}%
\pgfpathlineto{\pgfqpoint{5.447400in}{3.244769in}}%
\pgfpathlineto{\pgfqpoint{5.447400in}{0.660000in}}%
\pgfpathclose%
\pgfusepath{fill}%
\end{pgfscope}%
\begin{pgfscope}%
\pgfpathrectangle{\pgfqpoint{1.250000in}{0.660000in}}{\pgfqpoint{7.750000in}{4.620000in}}%
\pgfusepath{clip}%
\pgfsetbuttcap%
\pgfsetmiterjoin%
\definecolor{currentfill}{rgb}{0.000000,0.000000,0.000000}%
\pgfsetfillcolor{currentfill}%
\pgfsetlinewidth{0.000000pt}%
\definecolor{currentstroke}{rgb}{0.000000,0.000000,0.000000}%
\pgfsetstrokecolor{currentstroke}%
\pgfsetstrokeopacity{0.000000}%
\pgfsetdash{}{0pt}%
\pgfpathmoveto{\pgfqpoint{5.455150in}{0.660000in}}%
\pgfpathlineto{\pgfqpoint{5.461350in}{0.660000in}}%
\pgfpathlineto{\pgfqpoint{5.461350in}{3.244507in}}%
\pgfpathlineto{\pgfqpoint{5.455150in}{3.244507in}}%
\pgfpathlineto{\pgfqpoint{5.455150in}{0.660000in}}%
\pgfpathclose%
\pgfusepath{fill}%
\end{pgfscope}%
\begin{pgfscope}%
\pgfpathrectangle{\pgfqpoint{1.250000in}{0.660000in}}{\pgfqpoint{7.750000in}{4.620000in}}%
\pgfusepath{clip}%
\pgfsetbuttcap%
\pgfsetmiterjoin%
\definecolor{currentfill}{rgb}{0.000000,0.000000,0.000000}%
\pgfsetfillcolor{currentfill}%
\pgfsetlinewidth{0.000000pt}%
\definecolor{currentstroke}{rgb}{0.000000,0.000000,0.000000}%
\pgfsetstrokecolor{currentstroke}%
\pgfsetstrokeopacity{0.000000}%
\pgfsetdash{}{0pt}%
\pgfpathmoveto{\pgfqpoint{5.462900in}{0.660000in}}%
\pgfpathlineto{\pgfqpoint{5.469100in}{0.660000in}}%
\pgfpathlineto{\pgfqpoint{5.469100in}{3.244353in}}%
\pgfpathlineto{\pgfqpoint{5.462900in}{3.244353in}}%
\pgfpathlineto{\pgfqpoint{5.462900in}{0.660000in}}%
\pgfpathclose%
\pgfusepath{fill}%
\end{pgfscope}%
\begin{pgfscope}%
\pgfpathrectangle{\pgfqpoint{1.250000in}{0.660000in}}{\pgfqpoint{7.750000in}{4.620000in}}%
\pgfusepath{clip}%
\pgfsetbuttcap%
\pgfsetmiterjoin%
\definecolor{currentfill}{rgb}{0.000000,0.000000,0.000000}%
\pgfsetfillcolor{currentfill}%
\pgfsetlinewidth{0.000000pt}%
\definecolor{currentstroke}{rgb}{0.000000,0.000000,0.000000}%
\pgfsetstrokecolor{currentstroke}%
\pgfsetstrokeopacity{0.000000}%
\pgfsetdash{}{0pt}%
\pgfpathmoveto{\pgfqpoint{5.470650in}{0.660000in}}%
\pgfpathlineto{\pgfqpoint{5.476850in}{0.660000in}}%
\pgfpathlineto{\pgfqpoint{5.476850in}{3.244255in}}%
\pgfpathlineto{\pgfqpoint{5.470650in}{3.244255in}}%
\pgfpathlineto{\pgfqpoint{5.470650in}{0.660000in}}%
\pgfpathclose%
\pgfusepath{fill}%
\end{pgfscope}%
\begin{pgfscope}%
\pgfpathrectangle{\pgfqpoint{1.250000in}{0.660000in}}{\pgfqpoint{7.750000in}{4.620000in}}%
\pgfusepath{clip}%
\pgfsetbuttcap%
\pgfsetmiterjoin%
\definecolor{currentfill}{rgb}{0.000000,0.000000,0.000000}%
\pgfsetfillcolor{currentfill}%
\pgfsetlinewidth{0.000000pt}%
\definecolor{currentstroke}{rgb}{0.000000,0.000000,0.000000}%
\pgfsetstrokecolor{currentstroke}%
\pgfsetstrokeopacity{0.000000}%
\pgfsetdash{}{0pt}%
\pgfpathmoveto{\pgfqpoint{5.478400in}{0.660000in}}%
\pgfpathlineto{\pgfqpoint{5.484600in}{0.660000in}}%
\pgfpathlineto{\pgfqpoint{5.484600in}{3.256674in}}%
\pgfpathlineto{\pgfqpoint{5.478400in}{3.256674in}}%
\pgfpathlineto{\pgfqpoint{5.478400in}{0.660000in}}%
\pgfpathclose%
\pgfusepath{fill}%
\end{pgfscope}%
\begin{pgfscope}%
\pgfpathrectangle{\pgfqpoint{1.250000in}{0.660000in}}{\pgfqpoint{7.750000in}{4.620000in}}%
\pgfusepath{clip}%
\pgfsetbuttcap%
\pgfsetmiterjoin%
\definecolor{currentfill}{rgb}{0.000000,0.000000,0.000000}%
\pgfsetfillcolor{currentfill}%
\pgfsetlinewidth{0.000000pt}%
\definecolor{currentstroke}{rgb}{0.000000,0.000000,0.000000}%
\pgfsetstrokecolor{currentstroke}%
\pgfsetstrokeopacity{0.000000}%
\pgfsetdash{}{0pt}%
\pgfpathmoveto{\pgfqpoint{5.486150in}{0.660000in}}%
\pgfpathlineto{\pgfqpoint{5.492350in}{0.660000in}}%
\pgfpathlineto{\pgfqpoint{5.492350in}{3.250088in}}%
\pgfpathlineto{\pgfqpoint{5.486150in}{3.250088in}}%
\pgfpathlineto{\pgfqpoint{5.486150in}{0.660000in}}%
\pgfpathclose%
\pgfusepath{fill}%
\end{pgfscope}%
\begin{pgfscope}%
\pgfpathrectangle{\pgfqpoint{1.250000in}{0.660000in}}{\pgfqpoint{7.750000in}{4.620000in}}%
\pgfusepath{clip}%
\pgfsetbuttcap%
\pgfsetmiterjoin%
\definecolor{currentfill}{rgb}{0.000000,0.000000,0.000000}%
\pgfsetfillcolor{currentfill}%
\pgfsetlinewidth{0.000000pt}%
\definecolor{currentstroke}{rgb}{0.000000,0.000000,0.000000}%
\pgfsetstrokecolor{currentstroke}%
\pgfsetstrokeopacity{0.000000}%
\pgfsetdash{}{0pt}%
\pgfpathmoveto{\pgfqpoint{5.493900in}{0.660000in}}%
\pgfpathlineto{\pgfqpoint{5.500100in}{0.660000in}}%
\pgfpathlineto{\pgfqpoint{5.500100in}{3.272553in}}%
\pgfpathlineto{\pgfqpoint{5.493900in}{3.272553in}}%
\pgfpathlineto{\pgfqpoint{5.493900in}{0.660000in}}%
\pgfpathclose%
\pgfusepath{fill}%
\end{pgfscope}%
\begin{pgfscope}%
\pgfpathrectangle{\pgfqpoint{1.250000in}{0.660000in}}{\pgfqpoint{7.750000in}{4.620000in}}%
\pgfusepath{clip}%
\pgfsetbuttcap%
\pgfsetmiterjoin%
\definecolor{currentfill}{rgb}{0.000000,0.000000,0.000000}%
\pgfsetfillcolor{currentfill}%
\pgfsetlinewidth{0.000000pt}%
\definecolor{currentstroke}{rgb}{0.000000,0.000000,0.000000}%
\pgfsetstrokecolor{currentstroke}%
\pgfsetstrokeopacity{0.000000}%
\pgfsetdash{}{0pt}%
\pgfpathmoveto{\pgfqpoint{5.501650in}{0.660000in}}%
\pgfpathlineto{\pgfqpoint{5.507850in}{0.660000in}}%
\pgfpathlineto{\pgfqpoint{5.507850in}{3.253668in}}%
\pgfpathlineto{\pgfqpoint{5.501650in}{3.253668in}}%
\pgfpathlineto{\pgfqpoint{5.501650in}{0.660000in}}%
\pgfpathclose%
\pgfusepath{fill}%
\end{pgfscope}%
\begin{pgfscope}%
\pgfpathrectangle{\pgfqpoint{1.250000in}{0.660000in}}{\pgfqpoint{7.750000in}{4.620000in}}%
\pgfusepath{clip}%
\pgfsetbuttcap%
\pgfsetmiterjoin%
\definecolor{currentfill}{rgb}{0.000000,0.000000,0.000000}%
\pgfsetfillcolor{currentfill}%
\pgfsetlinewidth{0.000000pt}%
\definecolor{currentstroke}{rgb}{0.000000,0.000000,0.000000}%
\pgfsetstrokecolor{currentstroke}%
\pgfsetstrokeopacity{0.000000}%
\pgfsetdash{}{0pt}%
\pgfpathmoveto{\pgfqpoint{5.509400in}{0.660000in}}%
\pgfpathlineto{\pgfqpoint{5.515600in}{0.660000in}}%
\pgfpathlineto{\pgfqpoint{5.515600in}{3.252477in}}%
\pgfpathlineto{\pgfqpoint{5.509400in}{3.252477in}}%
\pgfpathlineto{\pgfqpoint{5.509400in}{0.660000in}}%
\pgfpathclose%
\pgfusepath{fill}%
\end{pgfscope}%
\begin{pgfscope}%
\pgfpathrectangle{\pgfqpoint{1.250000in}{0.660000in}}{\pgfqpoint{7.750000in}{4.620000in}}%
\pgfusepath{clip}%
\pgfsetbuttcap%
\pgfsetmiterjoin%
\definecolor{currentfill}{rgb}{0.000000,0.000000,0.000000}%
\pgfsetfillcolor{currentfill}%
\pgfsetlinewidth{0.000000pt}%
\definecolor{currentstroke}{rgb}{0.000000,0.000000,0.000000}%
\pgfsetstrokecolor{currentstroke}%
\pgfsetstrokeopacity{0.000000}%
\pgfsetdash{}{0pt}%
\pgfpathmoveto{\pgfqpoint{5.517150in}{0.660000in}}%
\pgfpathlineto{\pgfqpoint{5.523350in}{0.660000in}}%
\pgfpathlineto{\pgfqpoint{5.523350in}{3.263863in}}%
\pgfpathlineto{\pgfqpoint{5.517150in}{3.263863in}}%
\pgfpathlineto{\pgfqpoint{5.517150in}{0.660000in}}%
\pgfpathclose%
\pgfusepath{fill}%
\end{pgfscope}%
\begin{pgfscope}%
\pgfpathrectangle{\pgfqpoint{1.250000in}{0.660000in}}{\pgfqpoint{7.750000in}{4.620000in}}%
\pgfusepath{clip}%
\pgfsetbuttcap%
\pgfsetmiterjoin%
\definecolor{currentfill}{rgb}{0.000000,0.000000,0.000000}%
\pgfsetfillcolor{currentfill}%
\pgfsetlinewidth{0.000000pt}%
\definecolor{currentstroke}{rgb}{0.000000,0.000000,0.000000}%
\pgfsetstrokecolor{currentstroke}%
\pgfsetstrokeopacity{0.000000}%
\pgfsetdash{}{0pt}%
\pgfpathmoveto{\pgfqpoint{5.524900in}{0.660000in}}%
\pgfpathlineto{\pgfqpoint{5.531100in}{0.660000in}}%
\pgfpathlineto{\pgfqpoint{5.531100in}{3.253108in}}%
\pgfpathlineto{\pgfqpoint{5.524900in}{3.253108in}}%
\pgfpathlineto{\pgfqpoint{5.524900in}{0.660000in}}%
\pgfpathclose%
\pgfusepath{fill}%
\end{pgfscope}%
\begin{pgfscope}%
\pgfpathrectangle{\pgfqpoint{1.250000in}{0.660000in}}{\pgfqpoint{7.750000in}{4.620000in}}%
\pgfusepath{clip}%
\pgfsetbuttcap%
\pgfsetmiterjoin%
\definecolor{currentfill}{rgb}{0.000000,0.000000,0.000000}%
\pgfsetfillcolor{currentfill}%
\pgfsetlinewidth{0.000000pt}%
\definecolor{currentstroke}{rgb}{0.000000,0.000000,0.000000}%
\pgfsetstrokecolor{currentstroke}%
\pgfsetstrokeopacity{0.000000}%
\pgfsetdash{}{0pt}%
\pgfpathmoveto{\pgfqpoint{5.532650in}{0.660000in}}%
\pgfpathlineto{\pgfqpoint{5.538850in}{0.660000in}}%
\pgfpathlineto{\pgfqpoint{5.538850in}{3.243259in}}%
\pgfpathlineto{\pgfqpoint{5.532650in}{3.243259in}}%
\pgfpathlineto{\pgfqpoint{5.532650in}{0.660000in}}%
\pgfpathclose%
\pgfusepath{fill}%
\end{pgfscope}%
\begin{pgfscope}%
\pgfpathrectangle{\pgfqpoint{1.250000in}{0.660000in}}{\pgfqpoint{7.750000in}{4.620000in}}%
\pgfusepath{clip}%
\pgfsetbuttcap%
\pgfsetmiterjoin%
\definecolor{currentfill}{rgb}{0.000000,0.000000,0.000000}%
\pgfsetfillcolor{currentfill}%
\pgfsetlinewidth{0.000000pt}%
\definecolor{currentstroke}{rgb}{0.000000,0.000000,0.000000}%
\pgfsetstrokecolor{currentstroke}%
\pgfsetstrokeopacity{0.000000}%
\pgfsetdash{}{0pt}%
\pgfpathmoveto{\pgfqpoint{5.540400in}{0.660000in}}%
\pgfpathlineto{\pgfqpoint{5.546600in}{0.660000in}}%
\pgfpathlineto{\pgfqpoint{5.546600in}{3.243565in}}%
\pgfpathlineto{\pgfqpoint{5.540400in}{3.243565in}}%
\pgfpathlineto{\pgfqpoint{5.540400in}{0.660000in}}%
\pgfpathclose%
\pgfusepath{fill}%
\end{pgfscope}%
\begin{pgfscope}%
\pgfpathrectangle{\pgfqpoint{1.250000in}{0.660000in}}{\pgfqpoint{7.750000in}{4.620000in}}%
\pgfusepath{clip}%
\pgfsetbuttcap%
\pgfsetmiterjoin%
\definecolor{currentfill}{rgb}{0.000000,0.000000,0.000000}%
\pgfsetfillcolor{currentfill}%
\pgfsetlinewidth{0.000000pt}%
\definecolor{currentstroke}{rgb}{0.000000,0.000000,0.000000}%
\pgfsetstrokecolor{currentstroke}%
\pgfsetstrokeopacity{0.000000}%
\pgfsetdash{}{0pt}%
\pgfpathmoveto{\pgfqpoint{5.548150in}{0.660000in}}%
\pgfpathlineto{\pgfqpoint{5.554350in}{0.660000in}}%
\pgfpathlineto{\pgfqpoint{5.554350in}{3.250381in}}%
\pgfpathlineto{\pgfqpoint{5.548150in}{3.250381in}}%
\pgfpathlineto{\pgfqpoint{5.548150in}{0.660000in}}%
\pgfpathclose%
\pgfusepath{fill}%
\end{pgfscope}%
\begin{pgfscope}%
\pgfpathrectangle{\pgfqpoint{1.250000in}{0.660000in}}{\pgfqpoint{7.750000in}{4.620000in}}%
\pgfusepath{clip}%
\pgfsetbuttcap%
\pgfsetmiterjoin%
\definecolor{currentfill}{rgb}{0.000000,0.000000,0.000000}%
\pgfsetfillcolor{currentfill}%
\pgfsetlinewidth{0.000000pt}%
\definecolor{currentstroke}{rgb}{0.000000,0.000000,0.000000}%
\pgfsetstrokecolor{currentstroke}%
\pgfsetstrokeopacity{0.000000}%
\pgfsetdash{}{0pt}%
\pgfpathmoveto{\pgfqpoint{5.555900in}{0.660000in}}%
\pgfpathlineto{\pgfqpoint{5.562100in}{0.660000in}}%
\pgfpathlineto{\pgfqpoint{5.562100in}{3.242802in}}%
\pgfpathlineto{\pgfqpoint{5.555900in}{3.242802in}}%
\pgfpathlineto{\pgfqpoint{5.555900in}{0.660000in}}%
\pgfpathclose%
\pgfusepath{fill}%
\end{pgfscope}%
\begin{pgfscope}%
\pgfpathrectangle{\pgfqpoint{1.250000in}{0.660000in}}{\pgfqpoint{7.750000in}{4.620000in}}%
\pgfusepath{clip}%
\pgfsetbuttcap%
\pgfsetmiterjoin%
\definecolor{currentfill}{rgb}{0.000000,0.000000,0.000000}%
\pgfsetfillcolor{currentfill}%
\pgfsetlinewidth{0.000000pt}%
\definecolor{currentstroke}{rgb}{0.000000,0.000000,0.000000}%
\pgfsetstrokecolor{currentstroke}%
\pgfsetstrokeopacity{0.000000}%
\pgfsetdash{}{0pt}%
\pgfpathmoveto{\pgfqpoint{5.563650in}{0.660000in}}%
\pgfpathlineto{\pgfqpoint{5.569850in}{0.660000in}}%
\pgfpathlineto{\pgfqpoint{5.569850in}{3.242641in}}%
\pgfpathlineto{\pgfqpoint{5.563650in}{3.242641in}}%
\pgfpathlineto{\pgfqpoint{5.563650in}{0.660000in}}%
\pgfpathclose%
\pgfusepath{fill}%
\end{pgfscope}%
\begin{pgfscope}%
\pgfpathrectangle{\pgfqpoint{1.250000in}{0.660000in}}{\pgfqpoint{7.750000in}{4.620000in}}%
\pgfusepath{clip}%
\pgfsetbuttcap%
\pgfsetmiterjoin%
\definecolor{currentfill}{rgb}{0.000000,0.000000,0.000000}%
\pgfsetfillcolor{currentfill}%
\pgfsetlinewidth{0.000000pt}%
\definecolor{currentstroke}{rgb}{0.000000,0.000000,0.000000}%
\pgfsetstrokecolor{currentstroke}%
\pgfsetstrokeopacity{0.000000}%
\pgfsetdash{}{0pt}%
\pgfpathmoveto{\pgfqpoint{5.571400in}{0.660000in}}%
\pgfpathlineto{\pgfqpoint{5.577600in}{0.660000in}}%
\pgfpathlineto{\pgfqpoint{5.577600in}{3.242476in}}%
\pgfpathlineto{\pgfqpoint{5.571400in}{3.242476in}}%
\pgfpathlineto{\pgfqpoint{5.571400in}{0.660000in}}%
\pgfpathclose%
\pgfusepath{fill}%
\end{pgfscope}%
\begin{pgfscope}%
\pgfpathrectangle{\pgfqpoint{1.250000in}{0.660000in}}{\pgfqpoint{7.750000in}{4.620000in}}%
\pgfusepath{clip}%
\pgfsetbuttcap%
\pgfsetmiterjoin%
\definecolor{currentfill}{rgb}{0.000000,0.000000,0.000000}%
\pgfsetfillcolor{currentfill}%
\pgfsetlinewidth{0.000000pt}%
\definecolor{currentstroke}{rgb}{0.000000,0.000000,0.000000}%
\pgfsetstrokecolor{currentstroke}%
\pgfsetstrokeopacity{0.000000}%
\pgfsetdash{}{0pt}%
\pgfpathmoveto{\pgfqpoint{5.579150in}{0.660000in}}%
\pgfpathlineto{\pgfqpoint{5.585350in}{0.660000in}}%
\pgfpathlineto{\pgfqpoint{5.585350in}{3.268441in}}%
\pgfpathlineto{\pgfqpoint{5.579150in}{3.268441in}}%
\pgfpathlineto{\pgfqpoint{5.579150in}{0.660000in}}%
\pgfpathclose%
\pgfusepath{fill}%
\end{pgfscope}%
\begin{pgfscope}%
\pgfpathrectangle{\pgfqpoint{1.250000in}{0.660000in}}{\pgfqpoint{7.750000in}{4.620000in}}%
\pgfusepath{clip}%
\pgfsetbuttcap%
\pgfsetmiterjoin%
\definecolor{currentfill}{rgb}{0.000000,0.000000,0.000000}%
\pgfsetfillcolor{currentfill}%
\pgfsetlinewidth{0.000000pt}%
\definecolor{currentstroke}{rgb}{0.000000,0.000000,0.000000}%
\pgfsetstrokecolor{currentstroke}%
\pgfsetstrokeopacity{0.000000}%
\pgfsetdash{}{0pt}%
\pgfpathmoveto{\pgfqpoint{5.586900in}{0.660000in}}%
\pgfpathlineto{\pgfqpoint{5.593100in}{0.660000in}}%
\pgfpathlineto{\pgfqpoint{5.593100in}{3.242136in}}%
\pgfpathlineto{\pgfqpoint{5.586900in}{3.242136in}}%
\pgfpathlineto{\pgfqpoint{5.586900in}{0.660000in}}%
\pgfpathclose%
\pgfusepath{fill}%
\end{pgfscope}%
\begin{pgfscope}%
\pgfpathrectangle{\pgfqpoint{1.250000in}{0.660000in}}{\pgfqpoint{7.750000in}{4.620000in}}%
\pgfusepath{clip}%
\pgfsetbuttcap%
\pgfsetmiterjoin%
\definecolor{currentfill}{rgb}{0.000000,0.000000,0.000000}%
\pgfsetfillcolor{currentfill}%
\pgfsetlinewidth{0.000000pt}%
\definecolor{currentstroke}{rgb}{0.000000,0.000000,0.000000}%
\pgfsetstrokecolor{currentstroke}%
\pgfsetstrokeopacity{0.000000}%
\pgfsetdash{}{0pt}%
\pgfpathmoveto{\pgfqpoint{5.594650in}{0.660000in}}%
\pgfpathlineto{\pgfqpoint{5.600850in}{0.660000in}}%
\pgfpathlineto{\pgfqpoint{5.600850in}{3.259155in}}%
\pgfpathlineto{\pgfqpoint{5.594650in}{3.259155in}}%
\pgfpathlineto{\pgfqpoint{5.594650in}{0.660000in}}%
\pgfpathclose%
\pgfusepath{fill}%
\end{pgfscope}%
\begin{pgfscope}%
\pgfpathrectangle{\pgfqpoint{1.250000in}{0.660000in}}{\pgfqpoint{7.750000in}{4.620000in}}%
\pgfusepath{clip}%
\pgfsetbuttcap%
\pgfsetmiterjoin%
\definecolor{currentfill}{rgb}{0.000000,0.000000,0.000000}%
\pgfsetfillcolor{currentfill}%
\pgfsetlinewidth{0.000000pt}%
\definecolor{currentstroke}{rgb}{0.000000,0.000000,0.000000}%
\pgfsetstrokecolor{currentstroke}%
\pgfsetstrokeopacity{0.000000}%
\pgfsetdash{}{0pt}%
\pgfpathmoveto{\pgfqpoint{5.602400in}{0.660000in}}%
\pgfpathlineto{\pgfqpoint{5.608600in}{0.660000in}}%
\pgfpathlineto{\pgfqpoint{5.608600in}{3.249099in}}%
\pgfpathlineto{\pgfqpoint{5.602400in}{3.249099in}}%
\pgfpathlineto{\pgfqpoint{5.602400in}{0.660000in}}%
\pgfpathclose%
\pgfusepath{fill}%
\end{pgfscope}%
\begin{pgfscope}%
\pgfpathrectangle{\pgfqpoint{1.250000in}{0.660000in}}{\pgfqpoint{7.750000in}{4.620000in}}%
\pgfusepath{clip}%
\pgfsetbuttcap%
\pgfsetmiterjoin%
\definecolor{currentfill}{rgb}{0.000000,0.000000,0.000000}%
\pgfsetfillcolor{currentfill}%
\pgfsetlinewidth{0.000000pt}%
\definecolor{currentstroke}{rgb}{0.000000,0.000000,0.000000}%
\pgfsetstrokecolor{currentstroke}%
\pgfsetstrokeopacity{0.000000}%
\pgfsetdash{}{0pt}%
\pgfpathmoveto{\pgfqpoint{5.610150in}{0.660000in}}%
\pgfpathlineto{\pgfqpoint{5.616350in}{0.660000in}}%
\pgfpathlineto{\pgfqpoint{5.616350in}{3.241601in}}%
\pgfpathlineto{\pgfqpoint{5.610150in}{3.241601in}}%
\pgfpathlineto{\pgfqpoint{5.610150in}{0.660000in}}%
\pgfpathclose%
\pgfusepath{fill}%
\end{pgfscope}%
\begin{pgfscope}%
\pgfpathrectangle{\pgfqpoint{1.250000in}{0.660000in}}{\pgfqpoint{7.750000in}{4.620000in}}%
\pgfusepath{clip}%
\pgfsetbuttcap%
\pgfsetmiterjoin%
\definecolor{currentfill}{rgb}{0.000000,0.000000,0.000000}%
\pgfsetfillcolor{currentfill}%
\pgfsetlinewidth{0.000000pt}%
\definecolor{currentstroke}{rgb}{0.000000,0.000000,0.000000}%
\pgfsetstrokecolor{currentstroke}%
\pgfsetstrokeopacity{0.000000}%
\pgfsetdash{}{0pt}%
\pgfpathmoveto{\pgfqpoint{5.617900in}{0.660000in}}%
\pgfpathlineto{\pgfqpoint{5.624100in}{0.660000in}}%
\pgfpathlineto{\pgfqpoint{5.624100in}{3.244909in}}%
\pgfpathlineto{\pgfqpoint{5.617900in}{3.244909in}}%
\pgfpathlineto{\pgfqpoint{5.617900in}{0.660000in}}%
\pgfpathclose%
\pgfusepath{fill}%
\end{pgfscope}%
\begin{pgfscope}%
\pgfpathrectangle{\pgfqpoint{1.250000in}{0.660000in}}{\pgfqpoint{7.750000in}{4.620000in}}%
\pgfusepath{clip}%
\pgfsetbuttcap%
\pgfsetmiterjoin%
\definecolor{currentfill}{rgb}{0.000000,0.000000,0.000000}%
\pgfsetfillcolor{currentfill}%
\pgfsetlinewidth{0.000000pt}%
\definecolor{currentstroke}{rgb}{0.000000,0.000000,0.000000}%
\pgfsetstrokecolor{currentstroke}%
\pgfsetstrokeopacity{0.000000}%
\pgfsetdash{}{0pt}%
\pgfpathmoveto{\pgfqpoint{5.625650in}{0.660000in}}%
\pgfpathlineto{\pgfqpoint{5.631850in}{0.660000in}}%
\pgfpathlineto{\pgfqpoint{5.631850in}{3.285900in}}%
\pgfpathlineto{\pgfqpoint{5.625650in}{3.285900in}}%
\pgfpathlineto{\pgfqpoint{5.625650in}{0.660000in}}%
\pgfpathclose%
\pgfusepath{fill}%
\end{pgfscope}%
\begin{pgfscope}%
\pgfpathrectangle{\pgfqpoint{1.250000in}{0.660000in}}{\pgfqpoint{7.750000in}{4.620000in}}%
\pgfusepath{clip}%
\pgfsetbuttcap%
\pgfsetmiterjoin%
\definecolor{currentfill}{rgb}{0.000000,0.000000,0.000000}%
\pgfsetfillcolor{currentfill}%
\pgfsetlinewidth{0.000000pt}%
\definecolor{currentstroke}{rgb}{0.000000,0.000000,0.000000}%
\pgfsetstrokecolor{currentstroke}%
\pgfsetstrokeopacity{0.000000}%
\pgfsetdash{}{0pt}%
\pgfpathmoveto{\pgfqpoint{5.633400in}{0.660000in}}%
\pgfpathlineto{\pgfqpoint{5.639600in}{0.660000in}}%
\pgfpathlineto{\pgfqpoint{5.639600in}{3.265978in}}%
\pgfpathlineto{\pgfqpoint{5.633400in}{3.265978in}}%
\pgfpathlineto{\pgfqpoint{5.633400in}{0.660000in}}%
\pgfpathclose%
\pgfusepath{fill}%
\end{pgfscope}%
\begin{pgfscope}%
\pgfpathrectangle{\pgfqpoint{1.250000in}{0.660000in}}{\pgfqpoint{7.750000in}{4.620000in}}%
\pgfusepath{clip}%
\pgfsetbuttcap%
\pgfsetmiterjoin%
\definecolor{currentfill}{rgb}{0.000000,0.000000,0.000000}%
\pgfsetfillcolor{currentfill}%
\pgfsetlinewidth{0.000000pt}%
\definecolor{currentstroke}{rgb}{0.000000,0.000000,0.000000}%
\pgfsetstrokecolor{currentstroke}%
\pgfsetstrokeopacity{0.000000}%
\pgfsetdash{}{0pt}%
\pgfpathmoveto{\pgfqpoint{5.641150in}{0.660000in}}%
\pgfpathlineto{\pgfqpoint{5.647350in}{0.660000in}}%
\pgfpathlineto{\pgfqpoint{5.647350in}{3.260948in}}%
\pgfpathlineto{\pgfqpoint{5.641150in}{3.260948in}}%
\pgfpathlineto{\pgfqpoint{5.641150in}{0.660000in}}%
\pgfpathclose%
\pgfusepath{fill}%
\end{pgfscope}%
\begin{pgfscope}%
\pgfpathrectangle{\pgfqpoint{1.250000in}{0.660000in}}{\pgfqpoint{7.750000in}{4.620000in}}%
\pgfusepath{clip}%
\pgfsetbuttcap%
\pgfsetmiterjoin%
\definecolor{currentfill}{rgb}{0.000000,0.000000,0.000000}%
\pgfsetfillcolor{currentfill}%
\pgfsetlinewidth{0.000000pt}%
\definecolor{currentstroke}{rgb}{0.000000,0.000000,0.000000}%
\pgfsetstrokecolor{currentstroke}%
\pgfsetstrokeopacity{0.000000}%
\pgfsetdash{}{0pt}%
\pgfpathmoveto{\pgfqpoint{5.648900in}{0.660000in}}%
\pgfpathlineto{\pgfqpoint{5.655100in}{0.660000in}}%
\pgfpathlineto{\pgfqpoint{5.655100in}{3.240654in}}%
\pgfpathlineto{\pgfqpoint{5.648900in}{3.240654in}}%
\pgfpathlineto{\pgfqpoint{5.648900in}{0.660000in}}%
\pgfpathclose%
\pgfusepath{fill}%
\end{pgfscope}%
\begin{pgfscope}%
\pgfpathrectangle{\pgfqpoint{1.250000in}{0.660000in}}{\pgfqpoint{7.750000in}{4.620000in}}%
\pgfusepath{clip}%
\pgfsetbuttcap%
\pgfsetmiterjoin%
\definecolor{currentfill}{rgb}{0.000000,0.000000,0.000000}%
\pgfsetfillcolor{currentfill}%
\pgfsetlinewidth{0.000000pt}%
\definecolor{currentstroke}{rgb}{0.000000,0.000000,0.000000}%
\pgfsetstrokecolor{currentstroke}%
\pgfsetstrokeopacity{0.000000}%
\pgfsetdash{}{0pt}%
\pgfpathmoveto{\pgfqpoint{5.656650in}{0.660000in}}%
\pgfpathlineto{\pgfqpoint{5.662850in}{0.660000in}}%
\pgfpathlineto{\pgfqpoint{5.662850in}{3.240458in}}%
\pgfpathlineto{\pgfqpoint{5.656650in}{3.240458in}}%
\pgfpathlineto{\pgfqpoint{5.656650in}{0.660000in}}%
\pgfpathclose%
\pgfusepath{fill}%
\end{pgfscope}%
\begin{pgfscope}%
\pgfpathrectangle{\pgfqpoint{1.250000in}{0.660000in}}{\pgfqpoint{7.750000in}{4.620000in}}%
\pgfusepath{clip}%
\pgfsetbuttcap%
\pgfsetmiterjoin%
\definecolor{currentfill}{rgb}{0.000000,0.000000,0.000000}%
\pgfsetfillcolor{currentfill}%
\pgfsetlinewidth{0.000000pt}%
\definecolor{currentstroke}{rgb}{0.000000,0.000000,0.000000}%
\pgfsetstrokecolor{currentstroke}%
\pgfsetstrokeopacity{0.000000}%
\pgfsetdash{}{0pt}%
\pgfpathmoveto{\pgfqpoint{5.664400in}{0.660000in}}%
\pgfpathlineto{\pgfqpoint{5.670600in}{0.660000in}}%
\pgfpathlineto{\pgfqpoint{5.670600in}{3.240260in}}%
\pgfpathlineto{\pgfqpoint{5.664400in}{3.240260in}}%
\pgfpathlineto{\pgfqpoint{5.664400in}{0.660000in}}%
\pgfpathclose%
\pgfusepath{fill}%
\end{pgfscope}%
\begin{pgfscope}%
\pgfpathrectangle{\pgfqpoint{1.250000in}{0.660000in}}{\pgfqpoint{7.750000in}{4.620000in}}%
\pgfusepath{clip}%
\pgfsetbuttcap%
\pgfsetmiterjoin%
\definecolor{currentfill}{rgb}{0.000000,0.000000,0.000000}%
\pgfsetfillcolor{currentfill}%
\pgfsetlinewidth{0.000000pt}%
\definecolor{currentstroke}{rgb}{0.000000,0.000000,0.000000}%
\pgfsetstrokecolor{currentstroke}%
\pgfsetstrokeopacity{0.000000}%
\pgfsetdash{}{0pt}%
\pgfpathmoveto{\pgfqpoint{5.672150in}{0.660000in}}%
\pgfpathlineto{\pgfqpoint{5.678350in}{0.660000in}}%
\pgfpathlineto{\pgfqpoint{5.678350in}{3.240060in}}%
\pgfpathlineto{\pgfqpoint{5.672150in}{3.240060in}}%
\pgfpathlineto{\pgfqpoint{5.672150in}{0.660000in}}%
\pgfpathclose%
\pgfusepath{fill}%
\end{pgfscope}%
\begin{pgfscope}%
\pgfpathrectangle{\pgfqpoint{1.250000in}{0.660000in}}{\pgfqpoint{7.750000in}{4.620000in}}%
\pgfusepath{clip}%
\pgfsetbuttcap%
\pgfsetmiterjoin%
\definecolor{currentfill}{rgb}{0.000000,0.000000,0.000000}%
\pgfsetfillcolor{currentfill}%
\pgfsetlinewidth{0.000000pt}%
\definecolor{currentstroke}{rgb}{0.000000,0.000000,0.000000}%
\pgfsetstrokecolor{currentstroke}%
\pgfsetstrokeopacity{0.000000}%
\pgfsetdash{}{0pt}%
\pgfpathmoveto{\pgfqpoint{5.679900in}{0.660000in}}%
\pgfpathlineto{\pgfqpoint{5.686100in}{0.660000in}}%
\pgfpathlineto{\pgfqpoint{5.686100in}{3.239859in}}%
\pgfpathlineto{\pgfqpoint{5.679900in}{3.239859in}}%
\pgfpathlineto{\pgfqpoint{5.679900in}{0.660000in}}%
\pgfpathclose%
\pgfusepath{fill}%
\end{pgfscope}%
\begin{pgfscope}%
\pgfpathrectangle{\pgfqpoint{1.250000in}{0.660000in}}{\pgfqpoint{7.750000in}{4.620000in}}%
\pgfusepath{clip}%
\pgfsetbuttcap%
\pgfsetmiterjoin%
\definecolor{currentfill}{rgb}{0.000000,0.000000,0.000000}%
\pgfsetfillcolor{currentfill}%
\pgfsetlinewidth{0.000000pt}%
\definecolor{currentstroke}{rgb}{0.000000,0.000000,0.000000}%
\pgfsetstrokecolor{currentstroke}%
\pgfsetstrokeopacity{0.000000}%
\pgfsetdash{}{0pt}%
\pgfpathmoveto{\pgfqpoint{5.687650in}{0.660000in}}%
\pgfpathlineto{\pgfqpoint{5.693850in}{0.660000in}}%
\pgfpathlineto{\pgfqpoint{5.693850in}{3.239656in}}%
\pgfpathlineto{\pgfqpoint{5.687650in}{3.239656in}}%
\pgfpathlineto{\pgfqpoint{5.687650in}{0.660000in}}%
\pgfpathclose%
\pgfusepath{fill}%
\end{pgfscope}%
\begin{pgfscope}%
\pgfpathrectangle{\pgfqpoint{1.250000in}{0.660000in}}{\pgfqpoint{7.750000in}{4.620000in}}%
\pgfusepath{clip}%
\pgfsetbuttcap%
\pgfsetmiterjoin%
\definecolor{currentfill}{rgb}{0.000000,0.000000,0.000000}%
\pgfsetfillcolor{currentfill}%
\pgfsetlinewidth{0.000000pt}%
\definecolor{currentstroke}{rgb}{0.000000,0.000000,0.000000}%
\pgfsetstrokecolor{currentstroke}%
\pgfsetstrokeopacity{0.000000}%
\pgfsetdash{}{0pt}%
\pgfpathmoveto{\pgfqpoint{5.695400in}{0.660000in}}%
\pgfpathlineto{\pgfqpoint{5.701600in}{0.660000in}}%
\pgfpathlineto{\pgfqpoint{5.701600in}{3.239452in}}%
\pgfpathlineto{\pgfqpoint{5.695400in}{3.239452in}}%
\pgfpathlineto{\pgfqpoint{5.695400in}{0.660000in}}%
\pgfpathclose%
\pgfusepath{fill}%
\end{pgfscope}%
\begin{pgfscope}%
\pgfpathrectangle{\pgfqpoint{1.250000in}{0.660000in}}{\pgfqpoint{7.750000in}{4.620000in}}%
\pgfusepath{clip}%
\pgfsetbuttcap%
\pgfsetmiterjoin%
\definecolor{currentfill}{rgb}{0.000000,0.000000,0.000000}%
\pgfsetfillcolor{currentfill}%
\pgfsetlinewidth{0.000000pt}%
\definecolor{currentstroke}{rgb}{0.000000,0.000000,0.000000}%
\pgfsetstrokecolor{currentstroke}%
\pgfsetstrokeopacity{0.000000}%
\pgfsetdash{}{0pt}%
\pgfpathmoveto{\pgfqpoint{5.703150in}{0.660000in}}%
\pgfpathlineto{\pgfqpoint{5.709350in}{0.660000in}}%
\pgfpathlineto{\pgfqpoint{5.709350in}{3.239246in}}%
\pgfpathlineto{\pgfqpoint{5.703150in}{3.239246in}}%
\pgfpathlineto{\pgfqpoint{5.703150in}{0.660000in}}%
\pgfpathclose%
\pgfusepath{fill}%
\end{pgfscope}%
\begin{pgfscope}%
\pgfpathrectangle{\pgfqpoint{1.250000in}{0.660000in}}{\pgfqpoint{7.750000in}{4.620000in}}%
\pgfusepath{clip}%
\pgfsetbuttcap%
\pgfsetmiterjoin%
\definecolor{currentfill}{rgb}{0.000000,0.000000,0.000000}%
\pgfsetfillcolor{currentfill}%
\pgfsetlinewidth{0.000000pt}%
\definecolor{currentstroke}{rgb}{0.000000,0.000000,0.000000}%
\pgfsetstrokecolor{currentstroke}%
\pgfsetstrokeopacity{0.000000}%
\pgfsetdash{}{0pt}%
\pgfpathmoveto{\pgfqpoint{5.710900in}{0.660000in}}%
\pgfpathlineto{\pgfqpoint{5.717100in}{0.660000in}}%
\pgfpathlineto{\pgfqpoint{5.717100in}{3.239039in}}%
\pgfpathlineto{\pgfqpoint{5.710900in}{3.239039in}}%
\pgfpathlineto{\pgfqpoint{5.710900in}{0.660000in}}%
\pgfpathclose%
\pgfusepath{fill}%
\end{pgfscope}%
\begin{pgfscope}%
\pgfpathrectangle{\pgfqpoint{1.250000in}{0.660000in}}{\pgfqpoint{7.750000in}{4.620000in}}%
\pgfusepath{clip}%
\pgfsetbuttcap%
\pgfsetmiterjoin%
\definecolor{currentfill}{rgb}{0.000000,0.000000,0.000000}%
\pgfsetfillcolor{currentfill}%
\pgfsetlinewidth{0.000000pt}%
\definecolor{currentstroke}{rgb}{0.000000,0.000000,0.000000}%
\pgfsetstrokecolor{currentstroke}%
\pgfsetstrokeopacity{0.000000}%
\pgfsetdash{}{0pt}%
\pgfpathmoveto{\pgfqpoint{5.718650in}{0.660000in}}%
\pgfpathlineto{\pgfqpoint{5.724850in}{0.660000in}}%
\pgfpathlineto{\pgfqpoint{5.724850in}{3.238832in}}%
\pgfpathlineto{\pgfqpoint{5.718650in}{3.238832in}}%
\pgfpathlineto{\pgfqpoint{5.718650in}{0.660000in}}%
\pgfpathclose%
\pgfusepath{fill}%
\end{pgfscope}%
\begin{pgfscope}%
\pgfpathrectangle{\pgfqpoint{1.250000in}{0.660000in}}{\pgfqpoint{7.750000in}{4.620000in}}%
\pgfusepath{clip}%
\pgfsetbuttcap%
\pgfsetmiterjoin%
\definecolor{currentfill}{rgb}{0.000000,0.000000,0.000000}%
\pgfsetfillcolor{currentfill}%
\pgfsetlinewidth{0.000000pt}%
\definecolor{currentstroke}{rgb}{0.000000,0.000000,0.000000}%
\pgfsetstrokecolor{currentstroke}%
\pgfsetstrokeopacity{0.000000}%
\pgfsetdash{}{0pt}%
\pgfpathmoveto{\pgfqpoint{5.726400in}{0.660000in}}%
\pgfpathlineto{\pgfqpoint{5.732600in}{0.660000in}}%
\pgfpathlineto{\pgfqpoint{5.732600in}{3.238623in}}%
\pgfpathlineto{\pgfqpoint{5.726400in}{3.238623in}}%
\pgfpathlineto{\pgfqpoint{5.726400in}{0.660000in}}%
\pgfpathclose%
\pgfusepath{fill}%
\end{pgfscope}%
\begin{pgfscope}%
\pgfpathrectangle{\pgfqpoint{1.250000in}{0.660000in}}{\pgfqpoint{7.750000in}{4.620000in}}%
\pgfusepath{clip}%
\pgfsetbuttcap%
\pgfsetmiterjoin%
\definecolor{currentfill}{rgb}{0.000000,0.000000,0.000000}%
\pgfsetfillcolor{currentfill}%
\pgfsetlinewidth{0.000000pt}%
\definecolor{currentstroke}{rgb}{0.000000,0.000000,0.000000}%
\pgfsetstrokecolor{currentstroke}%
\pgfsetstrokeopacity{0.000000}%
\pgfsetdash{}{0pt}%
\pgfpathmoveto{\pgfqpoint{5.734150in}{0.660000in}}%
\pgfpathlineto{\pgfqpoint{5.740350in}{0.660000in}}%
\pgfpathlineto{\pgfqpoint{5.740350in}{3.238414in}}%
\pgfpathlineto{\pgfqpoint{5.734150in}{3.238414in}}%
\pgfpathlineto{\pgfqpoint{5.734150in}{0.660000in}}%
\pgfpathclose%
\pgfusepath{fill}%
\end{pgfscope}%
\begin{pgfscope}%
\pgfpathrectangle{\pgfqpoint{1.250000in}{0.660000in}}{\pgfqpoint{7.750000in}{4.620000in}}%
\pgfusepath{clip}%
\pgfsetbuttcap%
\pgfsetmiterjoin%
\definecolor{currentfill}{rgb}{0.000000,0.000000,0.000000}%
\pgfsetfillcolor{currentfill}%
\pgfsetlinewidth{0.000000pt}%
\definecolor{currentstroke}{rgb}{0.000000,0.000000,0.000000}%
\pgfsetstrokecolor{currentstroke}%
\pgfsetstrokeopacity{0.000000}%
\pgfsetdash{}{0pt}%
\pgfpathmoveto{\pgfqpoint{5.741900in}{0.660000in}}%
\pgfpathlineto{\pgfqpoint{5.748100in}{0.660000in}}%
\pgfpathlineto{\pgfqpoint{5.748100in}{3.238204in}}%
\pgfpathlineto{\pgfqpoint{5.741900in}{3.238204in}}%
\pgfpathlineto{\pgfqpoint{5.741900in}{0.660000in}}%
\pgfpathclose%
\pgfusepath{fill}%
\end{pgfscope}%
\begin{pgfscope}%
\pgfpathrectangle{\pgfqpoint{1.250000in}{0.660000in}}{\pgfqpoint{7.750000in}{4.620000in}}%
\pgfusepath{clip}%
\pgfsetbuttcap%
\pgfsetmiterjoin%
\definecolor{currentfill}{rgb}{0.000000,0.000000,0.000000}%
\pgfsetfillcolor{currentfill}%
\pgfsetlinewidth{0.000000pt}%
\definecolor{currentstroke}{rgb}{0.000000,0.000000,0.000000}%
\pgfsetstrokecolor{currentstroke}%
\pgfsetstrokeopacity{0.000000}%
\pgfsetdash{}{0pt}%
\pgfpathmoveto{\pgfqpoint{5.749650in}{0.660000in}}%
\pgfpathlineto{\pgfqpoint{5.755850in}{0.660000in}}%
\pgfpathlineto{\pgfqpoint{5.755850in}{3.238251in}}%
\pgfpathlineto{\pgfqpoint{5.749650in}{3.238251in}}%
\pgfpathlineto{\pgfqpoint{5.749650in}{0.660000in}}%
\pgfpathclose%
\pgfusepath{fill}%
\end{pgfscope}%
\begin{pgfscope}%
\pgfpathrectangle{\pgfqpoint{1.250000in}{0.660000in}}{\pgfqpoint{7.750000in}{4.620000in}}%
\pgfusepath{clip}%
\pgfsetbuttcap%
\pgfsetmiterjoin%
\definecolor{currentfill}{rgb}{0.000000,0.000000,0.000000}%
\pgfsetfillcolor{currentfill}%
\pgfsetlinewidth{0.000000pt}%
\definecolor{currentstroke}{rgb}{0.000000,0.000000,0.000000}%
\pgfsetstrokecolor{currentstroke}%
\pgfsetstrokeopacity{0.000000}%
\pgfsetdash{}{0pt}%
\pgfpathmoveto{\pgfqpoint{5.757400in}{0.660000in}}%
\pgfpathlineto{\pgfqpoint{5.763600in}{0.660000in}}%
\pgfpathlineto{\pgfqpoint{5.763600in}{3.238112in}}%
\pgfpathlineto{\pgfqpoint{5.757400in}{3.238112in}}%
\pgfpathlineto{\pgfqpoint{5.757400in}{0.660000in}}%
\pgfpathclose%
\pgfusepath{fill}%
\end{pgfscope}%
\begin{pgfscope}%
\pgfpathrectangle{\pgfqpoint{1.250000in}{0.660000in}}{\pgfqpoint{7.750000in}{4.620000in}}%
\pgfusepath{clip}%
\pgfsetbuttcap%
\pgfsetmiterjoin%
\definecolor{currentfill}{rgb}{0.000000,0.000000,0.000000}%
\pgfsetfillcolor{currentfill}%
\pgfsetlinewidth{0.000000pt}%
\definecolor{currentstroke}{rgb}{0.000000,0.000000,0.000000}%
\pgfsetstrokecolor{currentstroke}%
\pgfsetstrokeopacity{0.000000}%
\pgfsetdash{}{0pt}%
\pgfpathmoveto{\pgfqpoint{5.765150in}{0.660000in}}%
\pgfpathlineto{\pgfqpoint{5.771350in}{0.660000in}}%
\pgfpathlineto{\pgfqpoint{5.771350in}{3.247660in}}%
\pgfpathlineto{\pgfqpoint{5.765150in}{3.247660in}}%
\pgfpathlineto{\pgfqpoint{5.765150in}{0.660000in}}%
\pgfpathclose%
\pgfusepath{fill}%
\end{pgfscope}%
\begin{pgfscope}%
\pgfpathrectangle{\pgfqpoint{1.250000in}{0.660000in}}{\pgfqpoint{7.750000in}{4.620000in}}%
\pgfusepath{clip}%
\pgfsetbuttcap%
\pgfsetmiterjoin%
\definecolor{currentfill}{rgb}{0.000000,0.000000,0.000000}%
\pgfsetfillcolor{currentfill}%
\pgfsetlinewidth{0.000000pt}%
\definecolor{currentstroke}{rgb}{0.000000,0.000000,0.000000}%
\pgfsetstrokecolor{currentstroke}%
\pgfsetstrokeopacity{0.000000}%
\pgfsetdash{}{0pt}%
\pgfpathmoveto{\pgfqpoint{5.772900in}{0.660000in}}%
\pgfpathlineto{\pgfqpoint{5.779100in}{0.660000in}}%
\pgfpathlineto{\pgfqpoint{5.779100in}{3.252116in}}%
\pgfpathlineto{\pgfqpoint{5.772900in}{3.252116in}}%
\pgfpathlineto{\pgfqpoint{5.772900in}{0.660000in}}%
\pgfpathclose%
\pgfusepath{fill}%
\end{pgfscope}%
\begin{pgfscope}%
\pgfpathrectangle{\pgfqpoint{1.250000in}{0.660000in}}{\pgfqpoint{7.750000in}{4.620000in}}%
\pgfusepath{clip}%
\pgfsetbuttcap%
\pgfsetmiterjoin%
\definecolor{currentfill}{rgb}{0.000000,0.000000,0.000000}%
\pgfsetfillcolor{currentfill}%
\pgfsetlinewidth{0.000000pt}%
\definecolor{currentstroke}{rgb}{0.000000,0.000000,0.000000}%
\pgfsetstrokecolor{currentstroke}%
\pgfsetstrokeopacity{0.000000}%
\pgfsetdash{}{0pt}%
\pgfpathmoveto{\pgfqpoint{5.780650in}{0.660000in}}%
\pgfpathlineto{\pgfqpoint{5.786850in}{0.660000in}}%
\pgfpathlineto{\pgfqpoint{5.786850in}{3.237441in}}%
\pgfpathlineto{\pgfqpoint{5.780650in}{3.237441in}}%
\pgfpathlineto{\pgfqpoint{5.780650in}{0.660000in}}%
\pgfpathclose%
\pgfusepath{fill}%
\end{pgfscope}%
\begin{pgfscope}%
\pgfpathrectangle{\pgfqpoint{1.250000in}{0.660000in}}{\pgfqpoint{7.750000in}{4.620000in}}%
\pgfusepath{clip}%
\pgfsetbuttcap%
\pgfsetmiterjoin%
\definecolor{currentfill}{rgb}{0.000000,0.000000,0.000000}%
\pgfsetfillcolor{currentfill}%
\pgfsetlinewidth{0.000000pt}%
\definecolor{currentstroke}{rgb}{0.000000,0.000000,0.000000}%
\pgfsetstrokecolor{currentstroke}%
\pgfsetstrokeopacity{0.000000}%
\pgfsetdash{}{0pt}%
\pgfpathmoveto{\pgfqpoint{5.788400in}{0.660000in}}%
\pgfpathlineto{\pgfqpoint{5.794600in}{0.660000in}}%
\pgfpathlineto{\pgfqpoint{5.794600in}{3.237318in}}%
\pgfpathlineto{\pgfqpoint{5.788400in}{3.237318in}}%
\pgfpathlineto{\pgfqpoint{5.788400in}{0.660000in}}%
\pgfpathclose%
\pgfusepath{fill}%
\end{pgfscope}%
\begin{pgfscope}%
\pgfpathrectangle{\pgfqpoint{1.250000in}{0.660000in}}{\pgfqpoint{7.750000in}{4.620000in}}%
\pgfusepath{clip}%
\pgfsetbuttcap%
\pgfsetmiterjoin%
\definecolor{currentfill}{rgb}{0.000000,0.000000,0.000000}%
\pgfsetfillcolor{currentfill}%
\pgfsetlinewidth{0.000000pt}%
\definecolor{currentstroke}{rgb}{0.000000,0.000000,0.000000}%
\pgfsetstrokecolor{currentstroke}%
\pgfsetstrokeopacity{0.000000}%
\pgfsetdash{}{0pt}%
\pgfpathmoveto{\pgfqpoint{5.796150in}{0.660000in}}%
\pgfpathlineto{\pgfqpoint{5.802350in}{0.660000in}}%
\pgfpathlineto{\pgfqpoint{5.802350in}{3.251001in}}%
\pgfpathlineto{\pgfqpoint{5.796150in}{3.251001in}}%
\pgfpathlineto{\pgfqpoint{5.796150in}{0.660000in}}%
\pgfpathclose%
\pgfusepath{fill}%
\end{pgfscope}%
\begin{pgfscope}%
\pgfpathrectangle{\pgfqpoint{1.250000in}{0.660000in}}{\pgfqpoint{7.750000in}{4.620000in}}%
\pgfusepath{clip}%
\pgfsetbuttcap%
\pgfsetmiterjoin%
\definecolor{currentfill}{rgb}{0.000000,0.000000,0.000000}%
\pgfsetfillcolor{currentfill}%
\pgfsetlinewidth{0.000000pt}%
\definecolor{currentstroke}{rgb}{0.000000,0.000000,0.000000}%
\pgfsetstrokecolor{currentstroke}%
\pgfsetstrokeopacity{0.000000}%
\pgfsetdash{}{0pt}%
\pgfpathmoveto{\pgfqpoint{5.803900in}{0.660000in}}%
\pgfpathlineto{\pgfqpoint{5.810100in}{0.660000in}}%
\pgfpathlineto{\pgfqpoint{5.810100in}{3.237237in}}%
\pgfpathlineto{\pgfqpoint{5.803900in}{3.237237in}}%
\pgfpathlineto{\pgfqpoint{5.803900in}{0.660000in}}%
\pgfpathclose%
\pgfusepath{fill}%
\end{pgfscope}%
\begin{pgfscope}%
\pgfpathrectangle{\pgfqpoint{1.250000in}{0.660000in}}{\pgfqpoint{7.750000in}{4.620000in}}%
\pgfusepath{clip}%
\pgfsetbuttcap%
\pgfsetmiterjoin%
\definecolor{currentfill}{rgb}{0.000000,0.000000,0.000000}%
\pgfsetfillcolor{currentfill}%
\pgfsetlinewidth{0.000000pt}%
\definecolor{currentstroke}{rgb}{0.000000,0.000000,0.000000}%
\pgfsetstrokecolor{currentstroke}%
\pgfsetstrokeopacity{0.000000}%
\pgfsetdash{}{0pt}%
\pgfpathmoveto{\pgfqpoint{5.811650in}{0.660000in}}%
\pgfpathlineto{\pgfqpoint{5.817850in}{0.660000in}}%
\pgfpathlineto{\pgfqpoint{5.817850in}{3.238683in}}%
\pgfpathlineto{\pgfqpoint{5.811650in}{3.238683in}}%
\pgfpathlineto{\pgfqpoint{5.811650in}{0.660000in}}%
\pgfpathclose%
\pgfusepath{fill}%
\end{pgfscope}%
\begin{pgfscope}%
\pgfpathrectangle{\pgfqpoint{1.250000in}{0.660000in}}{\pgfqpoint{7.750000in}{4.620000in}}%
\pgfusepath{clip}%
\pgfsetbuttcap%
\pgfsetmiterjoin%
\definecolor{currentfill}{rgb}{0.000000,0.000000,0.000000}%
\pgfsetfillcolor{currentfill}%
\pgfsetlinewidth{0.000000pt}%
\definecolor{currentstroke}{rgb}{0.000000,0.000000,0.000000}%
\pgfsetstrokecolor{currentstroke}%
\pgfsetstrokeopacity{0.000000}%
\pgfsetdash{}{0pt}%
\pgfpathmoveto{\pgfqpoint{5.819400in}{0.660000in}}%
\pgfpathlineto{\pgfqpoint{5.825600in}{0.660000in}}%
\pgfpathlineto{\pgfqpoint{5.825600in}{3.236806in}}%
\pgfpathlineto{\pgfqpoint{5.819400in}{3.236806in}}%
\pgfpathlineto{\pgfqpoint{5.819400in}{0.660000in}}%
\pgfpathclose%
\pgfusepath{fill}%
\end{pgfscope}%
\begin{pgfscope}%
\pgfpathrectangle{\pgfqpoint{1.250000in}{0.660000in}}{\pgfqpoint{7.750000in}{4.620000in}}%
\pgfusepath{clip}%
\pgfsetbuttcap%
\pgfsetmiterjoin%
\definecolor{currentfill}{rgb}{0.000000,0.000000,0.000000}%
\pgfsetfillcolor{currentfill}%
\pgfsetlinewidth{0.000000pt}%
\definecolor{currentstroke}{rgb}{0.000000,0.000000,0.000000}%
\pgfsetstrokecolor{currentstroke}%
\pgfsetstrokeopacity{0.000000}%
\pgfsetdash{}{0pt}%
\pgfpathmoveto{\pgfqpoint{5.827150in}{0.660000in}}%
\pgfpathlineto{\pgfqpoint{5.833350in}{0.660000in}}%
\pgfpathlineto{\pgfqpoint{5.833350in}{3.245140in}}%
\pgfpathlineto{\pgfqpoint{5.827150in}{3.245140in}}%
\pgfpathlineto{\pgfqpoint{5.827150in}{0.660000in}}%
\pgfpathclose%
\pgfusepath{fill}%
\end{pgfscope}%
\begin{pgfscope}%
\pgfpathrectangle{\pgfqpoint{1.250000in}{0.660000in}}{\pgfqpoint{7.750000in}{4.620000in}}%
\pgfusepath{clip}%
\pgfsetbuttcap%
\pgfsetmiterjoin%
\definecolor{currentfill}{rgb}{0.000000,0.000000,0.000000}%
\pgfsetfillcolor{currentfill}%
\pgfsetlinewidth{0.000000pt}%
\definecolor{currentstroke}{rgb}{0.000000,0.000000,0.000000}%
\pgfsetstrokecolor{currentstroke}%
\pgfsetstrokeopacity{0.000000}%
\pgfsetdash{}{0pt}%
\pgfpathmoveto{\pgfqpoint{5.834900in}{0.660000in}}%
\pgfpathlineto{\pgfqpoint{5.841100in}{0.660000in}}%
\pgfpathlineto{\pgfqpoint{5.841100in}{3.236540in}}%
\pgfpathlineto{\pgfqpoint{5.834900in}{3.236540in}}%
\pgfpathlineto{\pgfqpoint{5.834900in}{0.660000in}}%
\pgfpathclose%
\pgfusepath{fill}%
\end{pgfscope}%
\begin{pgfscope}%
\pgfpathrectangle{\pgfqpoint{1.250000in}{0.660000in}}{\pgfqpoint{7.750000in}{4.620000in}}%
\pgfusepath{clip}%
\pgfsetbuttcap%
\pgfsetmiterjoin%
\definecolor{currentfill}{rgb}{0.000000,0.000000,0.000000}%
\pgfsetfillcolor{currentfill}%
\pgfsetlinewidth{0.000000pt}%
\definecolor{currentstroke}{rgb}{0.000000,0.000000,0.000000}%
\pgfsetstrokecolor{currentstroke}%
\pgfsetstrokeopacity{0.000000}%
\pgfsetdash{}{0pt}%
\pgfpathmoveto{\pgfqpoint{5.842650in}{0.660000in}}%
\pgfpathlineto{\pgfqpoint{5.848850in}{0.660000in}}%
\pgfpathlineto{\pgfqpoint{5.848850in}{3.242226in}}%
\pgfpathlineto{\pgfqpoint{5.842650in}{3.242226in}}%
\pgfpathlineto{\pgfqpoint{5.842650in}{0.660000in}}%
\pgfpathclose%
\pgfusepath{fill}%
\end{pgfscope}%
\begin{pgfscope}%
\pgfpathrectangle{\pgfqpoint{1.250000in}{0.660000in}}{\pgfqpoint{7.750000in}{4.620000in}}%
\pgfusepath{clip}%
\pgfsetbuttcap%
\pgfsetmiterjoin%
\definecolor{currentfill}{rgb}{0.000000,0.000000,0.000000}%
\pgfsetfillcolor{currentfill}%
\pgfsetlinewidth{0.000000pt}%
\definecolor{currentstroke}{rgb}{0.000000,0.000000,0.000000}%
\pgfsetstrokecolor{currentstroke}%
\pgfsetstrokeopacity{0.000000}%
\pgfsetdash{}{0pt}%
\pgfpathmoveto{\pgfqpoint{5.850400in}{0.660000in}}%
\pgfpathlineto{\pgfqpoint{5.856600in}{0.660000in}}%
\pgfpathlineto{\pgfqpoint{5.856600in}{3.249496in}}%
\pgfpathlineto{\pgfqpoint{5.850400in}{3.249496in}}%
\pgfpathlineto{\pgfqpoint{5.850400in}{0.660000in}}%
\pgfpathclose%
\pgfusepath{fill}%
\end{pgfscope}%
\begin{pgfscope}%
\pgfpathrectangle{\pgfqpoint{1.250000in}{0.660000in}}{\pgfqpoint{7.750000in}{4.620000in}}%
\pgfusepath{clip}%
\pgfsetbuttcap%
\pgfsetmiterjoin%
\definecolor{currentfill}{rgb}{0.000000,0.000000,0.000000}%
\pgfsetfillcolor{currentfill}%
\pgfsetlinewidth{0.000000pt}%
\definecolor{currentstroke}{rgb}{0.000000,0.000000,0.000000}%
\pgfsetstrokecolor{currentstroke}%
\pgfsetstrokeopacity{0.000000}%
\pgfsetdash{}{0pt}%
\pgfpathmoveto{\pgfqpoint{5.858150in}{0.660000in}}%
\pgfpathlineto{\pgfqpoint{5.864350in}{0.660000in}}%
\pgfpathlineto{\pgfqpoint{5.864350in}{3.308930in}}%
\pgfpathlineto{\pgfqpoint{5.858150in}{3.308930in}}%
\pgfpathlineto{\pgfqpoint{5.858150in}{0.660000in}}%
\pgfpathclose%
\pgfusepath{fill}%
\end{pgfscope}%
\begin{pgfscope}%
\pgfpathrectangle{\pgfqpoint{1.250000in}{0.660000in}}{\pgfqpoint{7.750000in}{4.620000in}}%
\pgfusepath{clip}%
\pgfsetbuttcap%
\pgfsetmiterjoin%
\definecolor{currentfill}{rgb}{0.000000,0.000000,0.000000}%
\pgfsetfillcolor{currentfill}%
\pgfsetlinewidth{0.000000pt}%
\definecolor{currentstroke}{rgb}{0.000000,0.000000,0.000000}%
\pgfsetstrokecolor{currentstroke}%
\pgfsetstrokeopacity{0.000000}%
\pgfsetdash{}{0pt}%
\pgfpathmoveto{\pgfqpoint{5.865900in}{0.660000in}}%
\pgfpathlineto{\pgfqpoint{5.872100in}{0.660000in}}%
\pgfpathlineto{\pgfqpoint{5.872100in}{3.235989in}}%
\pgfpathlineto{\pgfqpoint{5.865900in}{3.235989in}}%
\pgfpathlineto{\pgfqpoint{5.865900in}{0.660000in}}%
\pgfpathclose%
\pgfusepath{fill}%
\end{pgfscope}%
\begin{pgfscope}%
\pgfpathrectangle{\pgfqpoint{1.250000in}{0.660000in}}{\pgfqpoint{7.750000in}{4.620000in}}%
\pgfusepath{clip}%
\pgfsetbuttcap%
\pgfsetmiterjoin%
\definecolor{currentfill}{rgb}{0.000000,0.000000,0.000000}%
\pgfsetfillcolor{currentfill}%
\pgfsetlinewidth{0.000000pt}%
\definecolor{currentstroke}{rgb}{0.000000,0.000000,0.000000}%
\pgfsetstrokecolor{currentstroke}%
\pgfsetstrokeopacity{0.000000}%
\pgfsetdash{}{0pt}%
\pgfpathmoveto{\pgfqpoint{5.873650in}{0.660000in}}%
\pgfpathlineto{\pgfqpoint{5.879850in}{0.660000in}}%
\pgfpathlineto{\pgfqpoint{5.879850in}{3.235848in}}%
\pgfpathlineto{\pgfqpoint{5.873650in}{3.235848in}}%
\pgfpathlineto{\pgfqpoint{5.873650in}{0.660000in}}%
\pgfpathclose%
\pgfusepath{fill}%
\end{pgfscope}%
\begin{pgfscope}%
\pgfpathrectangle{\pgfqpoint{1.250000in}{0.660000in}}{\pgfqpoint{7.750000in}{4.620000in}}%
\pgfusepath{clip}%
\pgfsetbuttcap%
\pgfsetmiterjoin%
\definecolor{currentfill}{rgb}{0.000000,0.000000,0.000000}%
\pgfsetfillcolor{currentfill}%
\pgfsetlinewidth{0.000000pt}%
\definecolor{currentstroke}{rgb}{0.000000,0.000000,0.000000}%
\pgfsetstrokecolor{currentstroke}%
\pgfsetstrokeopacity{0.000000}%
\pgfsetdash{}{0pt}%
\pgfpathmoveto{\pgfqpoint{5.881400in}{0.660000in}}%
\pgfpathlineto{\pgfqpoint{5.887600in}{0.660000in}}%
\pgfpathlineto{\pgfqpoint{5.887600in}{3.235706in}}%
\pgfpathlineto{\pgfqpoint{5.881400in}{3.235706in}}%
\pgfpathlineto{\pgfqpoint{5.881400in}{0.660000in}}%
\pgfpathclose%
\pgfusepath{fill}%
\end{pgfscope}%
\begin{pgfscope}%
\pgfpathrectangle{\pgfqpoint{1.250000in}{0.660000in}}{\pgfqpoint{7.750000in}{4.620000in}}%
\pgfusepath{clip}%
\pgfsetbuttcap%
\pgfsetmiterjoin%
\definecolor{currentfill}{rgb}{0.000000,0.000000,0.000000}%
\pgfsetfillcolor{currentfill}%
\pgfsetlinewidth{0.000000pt}%
\definecolor{currentstroke}{rgb}{0.000000,0.000000,0.000000}%
\pgfsetstrokecolor{currentstroke}%
\pgfsetstrokeopacity{0.000000}%
\pgfsetdash{}{0pt}%
\pgfpathmoveto{\pgfqpoint{5.889150in}{0.660000in}}%
\pgfpathlineto{\pgfqpoint{5.895350in}{0.660000in}}%
\pgfpathlineto{\pgfqpoint{5.895350in}{3.235563in}}%
\pgfpathlineto{\pgfqpoint{5.889150in}{3.235563in}}%
\pgfpathlineto{\pgfqpoint{5.889150in}{0.660000in}}%
\pgfpathclose%
\pgfusepath{fill}%
\end{pgfscope}%
\begin{pgfscope}%
\pgfpathrectangle{\pgfqpoint{1.250000in}{0.660000in}}{\pgfqpoint{7.750000in}{4.620000in}}%
\pgfusepath{clip}%
\pgfsetbuttcap%
\pgfsetmiterjoin%
\definecolor{currentfill}{rgb}{0.000000,0.000000,0.000000}%
\pgfsetfillcolor{currentfill}%
\pgfsetlinewidth{0.000000pt}%
\definecolor{currentstroke}{rgb}{0.000000,0.000000,0.000000}%
\pgfsetstrokecolor{currentstroke}%
\pgfsetstrokeopacity{0.000000}%
\pgfsetdash{}{0pt}%
\pgfpathmoveto{\pgfqpoint{5.896900in}{0.660000in}}%
\pgfpathlineto{\pgfqpoint{5.903100in}{0.660000in}}%
\pgfpathlineto{\pgfqpoint{5.903100in}{3.235419in}}%
\pgfpathlineto{\pgfqpoint{5.896900in}{3.235419in}}%
\pgfpathlineto{\pgfqpoint{5.896900in}{0.660000in}}%
\pgfpathclose%
\pgfusepath{fill}%
\end{pgfscope}%
\begin{pgfscope}%
\pgfpathrectangle{\pgfqpoint{1.250000in}{0.660000in}}{\pgfqpoint{7.750000in}{4.620000in}}%
\pgfusepath{clip}%
\pgfsetbuttcap%
\pgfsetmiterjoin%
\definecolor{currentfill}{rgb}{0.000000,0.000000,0.000000}%
\pgfsetfillcolor{currentfill}%
\pgfsetlinewidth{0.000000pt}%
\definecolor{currentstroke}{rgb}{0.000000,0.000000,0.000000}%
\pgfsetstrokecolor{currentstroke}%
\pgfsetstrokeopacity{0.000000}%
\pgfsetdash{}{0pt}%
\pgfpathmoveto{\pgfqpoint{5.904650in}{0.660000in}}%
\pgfpathlineto{\pgfqpoint{5.910850in}{0.660000in}}%
\pgfpathlineto{\pgfqpoint{5.910850in}{3.235274in}}%
\pgfpathlineto{\pgfqpoint{5.904650in}{3.235274in}}%
\pgfpathlineto{\pgfqpoint{5.904650in}{0.660000in}}%
\pgfpathclose%
\pgfusepath{fill}%
\end{pgfscope}%
\begin{pgfscope}%
\pgfpathrectangle{\pgfqpoint{1.250000in}{0.660000in}}{\pgfqpoint{7.750000in}{4.620000in}}%
\pgfusepath{clip}%
\pgfsetbuttcap%
\pgfsetmiterjoin%
\definecolor{currentfill}{rgb}{0.000000,0.000000,0.000000}%
\pgfsetfillcolor{currentfill}%
\pgfsetlinewidth{0.000000pt}%
\definecolor{currentstroke}{rgb}{0.000000,0.000000,0.000000}%
\pgfsetstrokecolor{currentstroke}%
\pgfsetstrokeopacity{0.000000}%
\pgfsetdash{}{0pt}%
\pgfpathmoveto{\pgfqpoint{5.912400in}{0.660000in}}%
\pgfpathlineto{\pgfqpoint{5.918600in}{0.660000in}}%
\pgfpathlineto{\pgfqpoint{5.918600in}{3.235128in}}%
\pgfpathlineto{\pgfqpoint{5.912400in}{3.235128in}}%
\pgfpathlineto{\pgfqpoint{5.912400in}{0.660000in}}%
\pgfpathclose%
\pgfusepath{fill}%
\end{pgfscope}%
\begin{pgfscope}%
\pgfpathrectangle{\pgfqpoint{1.250000in}{0.660000in}}{\pgfqpoint{7.750000in}{4.620000in}}%
\pgfusepath{clip}%
\pgfsetbuttcap%
\pgfsetmiterjoin%
\definecolor{currentfill}{rgb}{0.000000,0.000000,0.000000}%
\pgfsetfillcolor{currentfill}%
\pgfsetlinewidth{0.000000pt}%
\definecolor{currentstroke}{rgb}{0.000000,0.000000,0.000000}%
\pgfsetstrokecolor{currentstroke}%
\pgfsetstrokeopacity{0.000000}%
\pgfsetdash{}{0pt}%
\pgfpathmoveto{\pgfqpoint{5.920150in}{0.660000in}}%
\pgfpathlineto{\pgfqpoint{5.926350in}{0.660000in}}%
\pgfpathlineto{\pgfqpoint{5.926350in}{3.234982in}}%
\pgfpathlineto{\pgfqpoint{5.920150in}{3.234982in}}%
\pgfpathlineto{\pgfqpoint{5.920150in}{0.660000in}}%
\pgfpathclose%
\pgfusepath{fill}%
\end{pgfscope}%
\begin{pgfscope}%
\pgfpathrectangle{\pgfqpoint{1.250000in}{0.660000in}}{\pgfqpoint{7.750000in}{4.620000in}}%
\pgfusepath{clip}%
\pgfsetbuttcap%
\pgfsetmiterjoin%
\definecolor{currentfill}{rgb}{0.000000,0.000000,0.000000}%
\pgfsetfillcolor{currentfill}%
\pgfsetlinewidth{0.000000pt}%
\definecolor{currentstroke}{rgb}{0.000000,0.000000,0.000000}%
\pgfsetstrokecolor{currentstroke}%
\pgfsetstrokeopacity{0.000000}%
\pgfsetdash{}{0pt}%
\pgfpathmoveto{\pgfqpoint{5.927900in}{0.660000in}}%
\pgfpathlineto{\pgfqpoint{5.934100in}{0.660000in}}%
\pgfpathlineto{\pgfqpoint{5.934100in}{3.234835in}}%
\pgfpathlineto{\pgfqpoint{5.927900in}{3.234835in}}%
\pgfpathlineto{\pgfqpoint{5.927900in}{0.660000in}}%
\pgfpathclose%
\pgfusepath{fill}%
\end{pgfscope}%
\begin{pgfscope}%
\pgfpathrectangle{\pgfqpoint{1.250000in}{0.660000in}}{\pgfqpoint{7.750000in}{4.620000in}}%
\pgfusepath{clip}%
\pgfsetbuttcap%
\pgfsetmiterjoin%
\definecolor{currentfill}{rgb}{0.000000,0.000000,0.000000}%
\pgfsetfillcolor{currentfill}%
\pgfsetlinewidth{0.000000pt}%
\definecolor{currentstroke}{rgb}{0.000000,0.000000,0.000000}%
\pgfsetstrokecolor{currentstroke}%
\pgfsetstrokeopacity{0.000000}%
\pgfsetdash{}{0pt}%
\pgfpathmoveto{\pgfqpoint{5.935650in}{0.660000in}}%
\pgfpathlineto{\pgfqpoint{5.941850in}{0.660000in}}%
\pgfpathlineto{\pgfqpoint{5.941850in}{3.234687in}}%
\pgfpathlineto{\pgfqpoint{5.935650in}{3.234687in}}%
\pgfpathlineto{\pgfqpoint{5.935650in}{0.660000in}}%
\pgfpathclose%
\pgfusepath{fill}%
\end{pgfscope}%
\begin{pgfscope}%
\pgfpathrectangle{\pgfqpoint{1.250000in}{0.660000in}}{\pgfqpoint{7.750000in}{4.620000in}}%
\pgfusepath{clip}%
\pgfsetbuttcap%
\pgfsetmiterjoin%
\definecolor{currentfill}{rgb}{0.000000,0.000000,0.000000}%
\pgfsetfillcolor{currentfill}%
\pgfsetlinewidth{0.000000pt}%
\definecolor{currentstroke}{rgb}{0.000000,0.000000,0.000000}%
\pgfsetstrokecolor{currentstroke}%
\pgfsetstrokeopacity{0.000000}%
\pgfsetdash{}{0pt}%
\pgfpathmoveto{\pgfqpoint{5.943400in}{0.660000in}}%
\pgfpathlineto{\pgfqpoint{5.949600in}{0.660000in}}%
\pgfpathlineto{\pgfqpoint{5.949600in}{3.236161in}}%
\pgfpathlineto{\pgfqpoint{5.943400in}{3.236161in}}%
\pgfpathlineto{\pgfqpoint{5.943400in}{0.660000in}}%
\pgfpathclose%
\pgfusepath{fill}%
\end{pgfscope}%
\begin{pgfscope}%
\pgfpathrectangle{\pgfqpoint{1.250000in}{0.660000in}}{\pgfqpoint{7.750000in}{4.620000in}}%
\pgfusepath{clip}%
\pgfsetbuttcap%
\pgfsetmiterjoin%
\definecolor{currentfill}{rgb}{0.000000,0.000000,0.000000}%
\pgfsetfillcolor{currentfill}%
\pgfsetlinewidth{0.000000pt}%
\definecolor{currentstroke}{rgb}{0.000000,0.000000,0.000000}%
\pgfsetstrokecolor{currentstroke}%
\pgfsetstrokeopacity{0.000000}%
\pgfsetdash{}{0pt}%
\pgfpathmoveto{\pgfqpoint{5.951150in}{0.660000in}}%
\pgfpathlineto{\pgfqpoint{5.957350in}{0.660000in}}%
\pgfpathlineto{\pgfqpoint{5.957350in}{3.234391in}}%
\pgfpathlineto{\pgfqpoint{5.951150in}{3.234391in}}%
\pgfpathlineto{\pgfqpoint{5.951150in}{0.660000in}}%
\pgfpathclose%
\pgfusepath{fill}%
\end{pgfscope}%
\begin{pgfscope}%
\pgfpathrectangle{\pgfqpoint{1.250000in}{0.660000in}}{\pgfqpoint{7.750000in}{4.620000in}}%
\pgfusepath{clip}%
\pgfsetbuttcap%
\pgfsetmiterjoin%
\definecolor{currentfill}{rgb}{0.000000,0.000000,0.000000}%
\pgfsetfillcolor{currentfill}%
\pgfsetlinewidth{0.000000pt}%
\definecolor{currentstroke}{rgb}{0.000000,0.000000,0.000000}%
\pgfsetstrokecolor{currentstroke}%
\pgfsetstrokeopacity{0.000000}%
\pgfsetdash{}{0pt}%
\pgfpathmoveto{\pgfqpoint{5.958900in}{0.660000in}}%
\pgfpathlineto{\pgfqpoint{5.965100in}{0.660000in}}%
\pgfpathlineto{\pgfqpoint{5.965100in}{3.234242in}}%
\pgfpathlineto{\pgfqpoint{5.958900in}{3.234242in}}%
\pgfpathlineto{\pgfqpoint{5.958900in}{0.660000in}}%
\pgfpathclose%
\pgfusepath{fill}%
\end{pgfscope}%
\begin{pgfscope}%
\pgfpathrectangle{\pgfqpoint{1.250000in}{0.660000in}}{\pgfqpoint{7.750000in}{4.620000in}}%
\pgfusepath{clip}%
\pgfsetbuttcap%
\pgfsetmiterjoin%
\definecolor{currentfill}{rgb}{0.000000,0.000000,0.000000}%
\pgfsetfillcolor{currentfill}%
\pgfsetlinewidth{0.000000pt}%
\definecolor{currentstroke}{rgb}{0.000000,0.000000,0.000000}%
\pgfsetstrokecolor{currentstroke}%
\pgfsetstrokeopacity{0.000000}%
\pgfsetdash{}{0pt}%
\pgfpathmoveto{\pgfqpoint{5.966650in}{0.660000in}}%
\pgfpathlineto{\pgfqpoint{5.972850in}{0.660000in}}%
\pgfpathlineto{\pgfqpoint{5.972850in}{3.265904in}}%
\pgfpathlineto{\pgfqpoint{5.966650in}{3.265904in}}%
\pgfpathlineto{\pgfqpoint{5.966650in}{0.660000in}}%
\pgfpathclose%
\pgfusepath{fill}%
\end{pgfscope}%
\begin{pgfscope}%
\pgfpathrectangle{\pgfqpoint{1.250000in}{0.660000in}}{\pgfqpoint{7.750000in}{4.620000in}}%
\pgfusepath{clip}%
\pgfsetbuttcap%
\pgfsetmiterjoin%
\definecolor{currentfill}{rgb}{0.000000,0.000000,0.000000}%
\pgfsetfillcolor{currentfill}%
\pgfsetlinewidth{0.000000pt}%
\definecolor{currentstroke}{rgb}{0.000000,0.000000,0.000000}%
\pgfsetstrokecolor{currentstroke}%
\pgfsetstrokeopacity{0.000000}%
\pgfsetdash{}{0pt}%
\pgfpathmoveto{\pgfqpoint{5.974400in}{0.660000in}}%
\pgfpathlineto{\pgfqpoint{5.980600in}{0.660000in}}%
\pgfpathlineto{\pgfqpoint{5.980600in}{3.269079in}}%
\pgfpathlineto{\pgfqpoint{5.974400in}{3.269079in}}%
\pgfpathlineto{\pgfqpoint{5.974400in}{0.660000in}}%
\pgfpathclose%
\pgfusepath{fill}%
\end{pgfscope}%
\begin{pgfscope}%
\pgfpathrectangle{\pgfqpoint{1.250000in}{0.660000in}}{\pgfqpoint{7.750000in}{4.620000in}}%
\pgfusepath{clip}%
\pgfsetbuttcap%
\pgfsetmiterjoin%
\definecolor{currentfill}{rgb}{0.000000,0.000000,0.000000}%
\pgfsetfillcolor{currentfill}%
\pgfsetlinewidth{0.000000pt}%
\definecolor{currentstroke}{rgb}{0.000000,0.000000,0.000000}%
\pgfsetstrokecolor{currentstroke}%
\pgfsetstrokeopacity{0.000000}%
\pgfsetdash{}{0pt}%
\pgfpathmoveto{\pgfqpoint{5.982150in}{0.660000in}}%
\pgfpathlineto{\pgfqpoint{5.988350in}{0.660000in}}%
\pgfpathlineto{\pgfqpoint{5.988350in}{3.257895in}}%
\pgfpathlineto{\pgfqpoint{5.982150in}{3.257895in}}%
\pgfpathlineto{\pgfqpoint{5.982150in}{0.660000in}}%
\pgfpathclose%
\pgfusepath{fill}%
\end{pgfscope}%
\begin{pgfscope}%
\pgfpathrectangle{\pgfqpoint{1.250000in}{0.660000in}}{\pgfqpoint{7.750000in}{4.620000in}}%
\pgfusepath{clip}%
\pgfsetbuttcap%
\pgfsetmiterjoin%
\definecolor{currentfill}{rgb}{0.000000,0.000000,0.000000}%
\pgfsetfillcolor{currentfill}%
\pgfsetlinewidth{0.000000pt}%
\definecolor{currentstroke}{rgb}{0.000000,0.000000,0.000000}%
\pgfsetstrokecolor{currentstroke}%
\pgfsetstrokeopacity{0.000000}%
\pgfsetdash{}{0pt}%
\pgfpathmoveto{\pgfqpoint{5.989900in}{0.660000in}}%
\pgfpathlineto{\pgfqpoint{5.996100in}{0.660000in}}%
\pgfpathlineto{\pgfqpoint{5.996100in}{3.233642in}}%
\pgfpathlineto{\pgfqpoint{5.989900in}{3.233642in}}%
\pgfpathlineto{\pgfqpoint{5.989900in}{0.660000in}}%
\pgfpathclose%
\pgfusepath{fill}%
\end{pgfscope}%
\begin{pgfscope}%
\pgfpathrectangle{\pgfqpoint{1.250000in}{0.660000in}}{\pgfqpoint{7.750000in}{4.620000in}}%
\pgfusepath{clip}%
\pgfsetbuttcap%
\pgfsetmiterjoin%
\definecolor{currentfill}{rgb}{0.000000,0.000000,0.000000}%
\pgfsetfillcolor{currentfill}%
\pgfsetlinewidth{0.000000pt}%
\definecolor{currentstroke}{rgb}{0.000000,0.000000,0.000000}%
\pgfsetstrokecolor{currentstroke}%
\pgfsetstrokeopacity{0.000000}%
\pgfsetdash{}{0pt}%
\pgfpathmoveto{\pgfqpoint{5.997650in}{0.660000in}}%
\pgfpathlineto{\pgfqpoint{6.003850in}{0.660000in}}%
\pgfpathlineto{\pgfqpoint{6.003850in}{3.233492in}}%
\pgfpathlineto{\pgfqpoint{5.997650in}{3.233492in}}%
\pgfpathlineto{\pgfqpoint{5.997650in}{0.660000in}}%
\pgfpathclose%
\pgfusepath{fill}%
\end{pgfscope}%
\begin{pgfscope}%
\pgfpathrectangle{\pgfqpoint{1.250000in}{0.660000in}}{\pgfqpoint{7.750000in}{4.620000in}}%
\pgfusepath{clip}%
\pgfsetbuttcap%
\pgfsetmiterjoin%
\definecolor{currentfill}{rgb}{0.000000,0.000000,0.000000}%
\pgfsetfillcolor{currentfill}%
\pgfsetlinewidth{0.000000pt}%
\definecolor{currentstroke}{rgb}{0.000000,0.000000,0.000000}%
\pgfsetstrokecolor{currentstroke}%
\pgfsetstrokeopacity{0.000000}%
\pgfsetdash{}{0pt}%
\pgfpathmoveto{\pgfqpoint{6.005400in}{0.660000in}}%
\pgfpathlineto{\pgfqpoint{6.011600in}{0.660000in}}%
\pgfpathlineto{\pgfqpoint{6.011600in}{3.233342in}}%
\pgfpathlineto{\pgfqpoint{6.005400in}{3.233342in}}%
\pgfpathlineto{\pgfqpoint{6.005400in}{0.660000in}}%
\pgfpathclose%
\pgfusepath{fill}%
\end{pgfscope}%
\begin{pgfscope}%
\pgfpathrectangle{\pgfqpoint{1.250000in}{0.660000in}}{\pgfqpoint{7.750000in}{4.620000in}}%
\pgfusepath{clip}%
\pgfsetbuttcap%
\pgfsetmiterjoin%
\definecolor{currentfill}{rgb}{0.000000,0.000000,0.000000}%
\pgfsetfillcolor{currentfill}%
\pgfsetlinewidth{0.000000pt}%
\definecolor{currentstroke}{rgb}{0.000000,0.000000,0.000000}%
\pgfsetstrokecolor{currentstroke}%
\pgfsetstrokeopacity{0.000000}%
\pgfsetdash{}{0pt}%
\pgfpathmoveto{\pgfqpoint{6.013150in}{0.660000in}}%
\pgfpathlineto{\pgfqpoint{6.019350in}{0.660000in}}%
\pgfpathlineto{\pgfqpoint{6.019350in}{3.238719in}}%
\pgfpathlineto{\pgfqpoint{6.013150in}{3.238719in}}%
\pgfpathlineto{\pgfqpoint{6.013150in}{0.660000in}}%
\pgfpathclose%
\pgfusepath{fill}%
\end{pgfscope}%
\begin{pgfscope}%
\pgfpathrectangle{\pgfqpoint{1.250000in}{0.660000in}}{\pgfqpoint{7.750000in}{4.620000in}}%
\pgfusepath{clip}%
\pgfsetbuttcap%
\pgfsetmiterjoin%
\definecolor{currentfill}{rgb}{0.000000,0.000000,0.000000}%
\pgfsetfillcolor{currentfill}%
\pgfsetlinewidth{0.000000pt}%
\definecolor{currentstroke}{rgb}{0.000000,0.000000,0.000000}%
\pgfsetstrokecolor{currentstroke}%
\pgfsetstrokeopacity{0.000000}%
\pgfsetdash{}{0pt}%
\pgfpathmoveto{\pgfqpoint{6.020900in}{0.660000in}}%
\pgfpathlineto{\pgfqpoint{6.027100in}{0.660000in}}%
\pgfpathlineto{\pgfqpoint{6.027100in}{3.233040in}}%
\pgfpathlineto{\pgfqpoint{6.020900in}{3.233040in}}%
\pgfpathlineto{\pgfqpoint{6.020900in}{0.660000in}}%
\pgfpathclose%
\pgfusepath{fill}%
\end{pgfscope}%
\begin{pgfscope}%
\pgfpathrectangle{\pgfqpoint{1.250000in}{0.660000in}}{\pgfqpoint{7.750000in}{4.620000in}}%
\pgfusepath{clip}%
\pgfsetbuttcap%
\pgfsetmiterjoin%
\definecolor{currentfill}{rgb}{0.000000,0.000000,0.000000}%
\pgfsetfillcolor{currentfill}%
\pgfsetlinewidth{0.000000pt}%
\definecolor{currentstroke}{rgb}{0.000000,0.000000,0.000000}%
\pgfsetstrokecolor{currentstroke}%
\pgfsetstrokeopacity{0.000000}%
\pgfsetdash{}{0pt}%
\pgfpathmoveto{\pgfqpoint{6.028650in}{0.660000in}}%
\pgfpathlineto{\pgfqpoint{6.034850in}{0.660000in}}%
\pgfpathlineto{\pgfqpoint{6.034850in}{3.240848in}}%
\pgfpathlineto{\pgfqpoint{6.028650in}{3.240848in}}%
\pgfpathlineto{\pgfqpoint{6.028650in}{0.660000in}}%
\pgfpathclose%
\pgfusepath{fill}%
\end{pgfscope}%
\begin{pgfscope}%
\pgfpathrectangle{\pgfqpoint{1.250000in}{0.660000in}}{\pgfqpoint{7.750000in}{4.620000in}}%
\pgfusepath{clip}%
\pgfsetbuttcap%
\pgfsetmiterjoin%
\definecolor{currentfill}{rgb}{0.000000,0.000000,0.000000}%
\pgfsetfillcolor{currentfill}%
\pgfsetlinewidth{0.000000pt}%
\definecolor{currentstroke}{rgb}{0.000000,0.000000,0.000000}%
\pgfsetstrokecolor{currentstroke}%
\pgfsetstrokeopacity{0.000000}%
\pgfsetdash{}{0pt}%
\pgfpathmoveto{\pgfqpoint{6.036400in}{0.660000in}}%
\pgfpathlineto{\pgfqpoint{6.042600in}{0.660000in}}%
\pgfpathlineto{\pgfqpoint{6.042600in}{3.232739in}}%
\pgfpathlineto{\pgfqpoint{6.036400in}{3.232739in}}%
\pgfpathlineto{\pgfqpoint{6.036400in}{0.660000in}}%
\pgfpathclose%
\pgfusepath{fill}%
\end{pgfscope}%
\begin{pgfscope}%
\pgfpathrectangle{\pgfqpoint{1.250000in}{0.660000in}}{\pgfqpoint{7.750000in}{4.620000in}}%
\pgfusepath{clip}%
\pgfsetbuttcap%
\pgfsetmiterjoin%
\definecolor{currentfill}{rgb}{0.000000,0.000000,0.000000}%
\pgfsetfillcolor{currentfill}%
\pgfsetlinewidth{0.000000pt}%
\definecolor{currentstroke}{rgb}{0.000000,0.000000,0.000000}%
\pgfsetstrokecolor{currentstroke}%
\pgfsetstrokeopacity{0.000000}%
\pgfsetdash{}{0pt}%
\pgfpathmoveto{\pgfqpoint{6.044150in}{0.660000in}}%
\pgfpathlineto{\pgfqpoint{6.050350in}{0.660000in}}%
\pgfpathlineto{\pgfqpoint{6.050350in}{3.232589in}}%
\pgfpathlineto{\pgfqpoint{6.044150in}{3.232589in}}%
\pgfpathlineto{\pgfqpoint{6.044150in}{0.660000in}}%
\pgfpathclose%
\pgfusepath{fill}%
\end{pgfscope}%
\begin{pgfscope}%
\pgfpathrectangle{\pgfqpoint{1.250000in}{0.660000in}}{\pgfqpoint{7.750000in}{4.620000in}}%
\pgfusepath{clip}%
\pgfsetbuttcap%
\pgfsetmiterjoin%
\definecolor{currentfill}{rgb}{0.000000,0.000000,0.000000}%
\pgfsetfillcolor{currentfill}%
\pgfsetlinewidth{0.000000pt}%
\definecolor{currentstroke}{rgb}{0.000000,0.000000,0.000000}%
\pgfsetstrokecolor{currentstroke}%
\pgfsetstrokeopacity{0.000000}%
\pgfsetdash{}{0pt}%
\pgfpathmoveto{\pgfqpoint{6.051900in}{0.660000in}}%
\pgfpathlineto{\pgfqpoint{6.058100in}{0.660000in}}%
\pgfpathlineto{\pgfqpoint{6.058100in}{3.232438in}}%
\pgfpathlineto{\pgfqpoint{6.051900in}{3.232438in}}%
\pgfpathlineto{\pgfqpoint{6.051900in}{0.660000in}}%
\pgfpathclose%
\pgfusepath{fill}%
\end{pgfscope}%
\begin{pgfscope}%
\pgfpathrectangle{\pgfqpoint{1.250000in}{0.660000in}}{\pgfqpoint{7.750000in}{4.620000in}}%
\pgfusepath{clip}%
\pgfsetbuttcap%
\pgfsetmiterjoin%
\definecolor{currentfill}{rgb}{0.000000,0.000000,0.000000}%
\pgfsetfillcolor{currentfill}%
\pgfsetlinewidth{0.000000pt}%
\definecolor{currentstroke}{rgb}{0.000000,0.000000,0.000000}%
\pgfsetstrokecolor{currentstroke}%
\pgfsetstrokeopacity{0.000000}%
\pgfsetdash{}{0pt}%
\pgfpathmoveto{\pgfqpoint{6.059650in}{0.660000in}}%
\pgfpathlineto{\pgfqpoint{6.065850in}{0.660000in}}%
\pgfpathlineto{\pgfqpoint{6.065850in}{3.257060in}}%
\pgfpathlineto{\pgfqpoint{6.059650in}{3.257060in}}%
\pgfpathlineto{\pgfqpoint{6.059650in}{0.660000in}}%
\pgfpathclose%
\pgfusepath{fill}%
\end{pgfscope}%
\begin{pgfscope}%
\pgfpathrectangle{\pgfqpoint{1.250000in}{0.660000in}}{\pgfqpoint{7.750000in}{4.620000in}}%
\pgfusepath{clip}%
\pgfsetbuttcap%
\pgfsetmiterjoin%
\definecolor{currentfill}{rgb}{0.000000,0.000000,0.000000}%
\pgfsetfillcolor{currentfill}%
\pgfsetlinewidth{0.000000pt}%
\definecolor{currentstroke}{rgb}{0.000000,0.000000,0.000000}%
\pgfsetstrokecolor{currentstroke}%
\pgfsetstrokeopacity{0.000000}%
\pgfsetdash{}{0pt}%
\pgfpathmoveto{\pgfqpoint{6.067400in}{0.660000in}}%
\pgfpathlineto{\pgfqpoint{6.073600in}{0.660000in}}%
\pgfpathlineto{\pgfqpoint{6.073600in}{3.232138in}}%
\pgfpathlineto{\pgfqpoint{6.067400in}{3.232138in}}%
\pgfpathlineto{\pgfqpoint{6.067400in}{0.660000in}}%
\pgfpathclose%
\pgfusepath{fill}%
\end{pgfscope}%
\begin{pgfscope}%
\pgfpathrectangle{\pgfqpoint{1.250000in}{0.660000in}}{\pgfqpoint{7.750000in}{4.620000in}}%
\pgfusepath{clip}%
\pgfsetbuttcap%
\pgfsetmiterjoin%
\definecolor{currentfill}{rgb}{0.000000,0.000000,0.000000}%
\pgfsetfillcolor{currentfill}%
\pgfsetlinewidth{0.000000pt}%
\definecolor{currentstroke}{rgb}{0.000000,0.000000,0.000000}%
\pgfsetstrokecolor{currentstroke}%
\pgfsetstrokeopacity{0.000000}%
\pgfsetdash{}{0pt}%
\pgfpathmoveto{\pgfqpoint{6.075150in}{0.660000in}}%
\pgfpathlineto{\pgfqpoint{6.081350in}{0.660000in}}%
\pgfpathlineto{\pgfqpoint{6.081350in}{3.248141in}}%
\pgfpathlineto{\pgfqpoint{6.075150in}{3.248141in}}%
\pgfpathlineto{\pgfqpoint{6.075150in}{0.660000in}}%
\pgfpathclose%
\pgfusepath{fill}%
\end{pgfscope}%
\begin{pgfscope}%
\pgfpathrectangle{\pgfqpoint{1.250000in}{0.660000in}}{\pgfqpoint{7.750000in}{4.620000in}}%
\pgfusepath{clip}%
\pgfsetbuttcap%
\pgfsetmiterjoin%
\definecolor{currentfill}{rgb}{0.000000,0.000000,0.000000}%
\pgfsetfillcolor{currentfill}%
\pgfsetlinewidth{0.000000pt}%
\definecolor{currentstroke}{rgb}{0.000000,0.000000,0.000000}%
\pgfsetstrokecolor{currentstroke}%
\pgfsetstrokeopacity{0.000000}%
\pgfsetdash{}{0pt}%
\pgfpathmoveto{\pgfqpoint{6.082900in}{0.660000in}}%
\pgfpathlineto{\pgfqpoint{6.089100in}{0.660000in}}%
\pgfpathlineto{\pgfqpoint{6.089100in}{3.231839in}}%
\pgfpathlineto{\pgfqpoint{6.082900in}{3.231839in}}%
\pgfpathlineto{\pgfqpoint{6.082900in}{0.660000in}}%
\pgfpathclose%
\pgfusepath{fill}%
\end{pgfscope}%
\begin{pgfscope}%
\pgfpathrectangle{\pgfqpoint{1.250000in}{0.660000in}}{\pgfqpoint{7.750000in}{4.620000in}}%
\pgfusepath{clip}%
\pgfsetbuttcap%
\pgfsetmiterjoin%
\definecolor{currentfill}{rgb}{0.000000,0.000000,0.000000}%
\pgfsetfillcolor{currentfill}%
\pgfsetlinewidth{0.000000pt}%
\definecolor{currentstroke}{rgb}{0.000000,0.000000,0.000000}%
\pgfsetstrokecolor{currentstroke}%
\pgfsetstrokeopacity{0.000000}%
\pgfsetdash{}{0pt}%
\pgfpathmoveto{\pgfqpoint{6.090650in}{0.660000in}}%
\pgfpathlineto{\pgfqpoint{6.096850in}{0.660000in}}%
\pgfpathlineto{\pgfqpoint{6.096850in}{3.238555in}}%
\pgfpathlineto{\pgfqpoint{6.090650in}{3.238555in}}%
\pgfpathlineto{\pgfqpoint{6.090650in}{0.660000in}}%
\pgfpathclose%
\pgfusepath{fill}%
\end{pgfscope}%
\begin{pgfscope}%
\pgfpathrectangle{\pgfqpoint{1.250000in}{0.660000in}}{\pgfqpoint{7.750000in}{4.620000in}}%
\pgfusepath{clip}%
\pgfsetbuttcap%
\pgfsetmiterjoin%
\definecolor{currentfill}{rgb}{0.000000,0.000000,0.000000}%
\pgfsetfillcolor{currentfill}%
\pgfsetlinewidth{0.000000pt}%
\definecolor{currentstroke}{rgb}{0.000000,0.000000,0.000000}%
\pgfsetstrokecolor{currentstroke}%
\pgfsetstrokeopacity{0.000000}%
\pgfsetdash{}{0pt}%
\pgfpathmoveto{\pgfqpoint{6.098400in}{0.660000in}}%
\pgfpathlineto{\pgfqpoint{6.104600in}{0.660000in}}%
\pgfpathlineto{\pgfqpoint{6.104600in}{3.231540in}}%
\pgfpathlineto{\pgfqpoint{6.098400in}{3.231540in}}%
\pgfpathlineto{\pgfqpoint{6.098400in}{0.660000in}}%
\pgfpathclose%
\pgfusepath{fill}%
\end{pgfscope}%
\begin{pgfscope}%
\pgfpathrectangle{\pgfqpoint{1.250000in}{0.660000in}}{\pgfqpoint{7.750000in}{4.620000in}}%
\pgfusepath{clip}%
\pgfsetbuttcap%
\pgfsetmiterjoin%
\definecolor{currentfill}{rgb}{0.000000,0.000000,0.000000}%
\pgfsetfillcolor{currentfill}%
\pgfsetlinewidth{0.000000pt}%
\definecolor{currentstroke}{rgb}{0.000000,0.000000,0.000000}%
\pgfsetstrokecolor{currentstroke}%
\pgfsetstrokeopacity{0.000000}%
\pgfsetdash{}{0pt}%
\pgfpathmoveto{\pgfqpoint{6.106150in}{0.660000in}}%
\pgfpathlineto{\pgfqpoint{6.112350in}{0.660000in}}%
\pgfpathlineto{\pgfqpoint{6.112350in}{3.231391in}}%
\pgfpathlineto{\pgfqpoint{6.106150in}{3.231391in}}%
\pgfpathlineto{\pgfqpoint{6.106150in}{0.660000in}}%
\pgfpathclose%
\pgfusepath{fill}%
\end{pgfscope}%
\begin{pgfscope}%
\pgfpathrectangle{\pgfqpoint{1.250000in}{0.660000in}}{\pgfqpoint{7.750000in}{4.620000in}}%
\pgfusepath{clip}%
\pgfsetbuttcap%
\pgfsetmiterjoin%
\definecolor{currentfill}{rgb}{0.000000,0.000000,0.000000}%
\pgfsetfillcolor{currentfill}%
\pgfsetlinewidth{0.000000pt}%
\definecolor{currentstroke}{rgb}{0.000000,0.000000,0.000000}%
\pgfsetstrokecolor{currentstroke}%
\pgfsetstrokeopacity{0.000000}%
\pgfsetdash{}{0pt}%
\pgfpathmoveto{\pgfqpoint{6.113900in}{0.660000in}}%
\pgfpathlineto{\pgfqpoint{6.120100in}{0.660000in}}%
\pgfpathlineto{\pgfqpoint{6.120100in}{3.231243in}}%
\pgfpathlineto{\pgfqpoint{6.113900in}{3.231243in}}%
\pgfpathlineto{\pgfqpoint{6.113900in}{0.660000in}}%
\pgfpathclose%
\pgfusepath{fill}%
\end{pgfscope}%
\begin{pgfscope}%
\pgfpathrectangle{\pgfqpoint{1.250000in}{0.660000in}}{\pgfqpoint{7.750000in}{4.620000in}}%
\pgfusepath{clip}%
\pgfsetbuttcap%
\pgfsetmiterjoin%
\definecolor{currentfill}{rgb}{0.000000,0.000000,0.000000}%
\pgfsetfillcolor{currentfill}%
\pgfsetlinewidth{0.000000pt}%
\definecolor{currentstroke}{rgb}{0.000000,0.000000,0.000000}%
\pgfsetstrokecolor{currentstroke}%
\pgfsetstrokeopacity{0.000000}%
\pgfsetdash{}{0pt}%
\pgfpathmoveto{\pgfqpoint{6.121650in}{0.660000in}}%
\pgfpathlineto{\pgfqpoint{6.127850in}{0.660000in}}%
\pgfpathlineto{\pgfqpoint{6.127850in}{3.248937in}}%
\pgfpathlineto{\pgfqpoint{6.121650in}{3.248937in}}%
\pgfpathlineto{\pgfqpoint{6.121650in}{0.660000in}}%
\pgfpathclose%
\pgfusepath{fill}%
\end{pgfscope}%
\begin{pgfscope}%
\pgfpathrectangle{\pgfqpoint{1.250000in}{0.660000in}}{\pgfqpoint{7.750000in}{4.620000in}}%
\pgfusepath{clip}%
\pgfsetbuttcap%
\pgfsetmiterjoin%
\definecolor{currentfill}{rgb}{0.000000,0.000000,0.000000}%
\pgfsetfillcolor{currentfill}%
\pgfsetlinewidth{0.000000pt}%
\definecolor{currentstroke}{rgb}{0.000000,0.000000,0.000000}%
\pgfsetstrokecolor{currentstroke}%
\pgfsetstrokeopacity{0.000000}%
\pgfsetdash{}{0pt}%
\pgfpathmoveto{\pgfqpoint{6.129400in}{0.660000in}}%
\pgfpathlineto{\pgfqpoint{6.135600in}{0.660000in}}%
\pgfpathlineto{\pgfqpoint{6.135600in}{3.238677in}}%
\pgfpathlineto{\pgfqpoint{6.129400in}{3.238677in}}%
\pgfpathlineto{\pgfqpoint{6.129400in}{0.660000in}}%
\pgfpathclose%
\pgfusepath{fill}%
\end{pgfscope}%
\begin{pgfscope}%
\pgfpathrectangle{\pgfqpoint{1.250000in}{0.660000in}}{\pgfqpoint{7.750000in}{4.620000in}}%
\pgfusepath{clip}%
\pgfsetbuttcap%
\pgfsetmiterjoin%
\definecolor{currentfill}{rgb}{0.000000,0.000000,0.000000}%
\pgfsetfillcolor{currentfill}%
\pgfsetlinewidth{0.000000pt}%
\definecolor{currentstroke}{rgb}{0.000000,0.000000,0.000000}%
\pgfsetstrokecolor{currentstroke}%
\pgfsetstrokeopacity{0.000000}%
\pgfsetdash{}{0pt}%
\pgfpathmoveto{\pgfqpoint{6.137150in}{0.660000in}}%
\pgfpathlineto{\pgfqpoint{6.143350in}{0.660000in}}%
\pgfpathlineto{\pgfqpoint{6.143350in}{3.230799in}}%
\pgfpathlineto{\pgfqpoint{6.137150in}{3.230799in}}%
\pgfpathlineto{\pgfqpoint{6.137150in}{0.660000in}}%
\pgfpathclose%
\pgfusepath{fill}%
\end{pgfscope}%
\begin{pgfscope}%
\pgfpathrectangle{\pgfqpoint{1.250000in}{0.660000in}}{\pgfqpoint{7.750000in}{4.620000in}}%
\pgfusepath{clip}%
\pgfsetbuttcap%
\pgfsetmiterjoin%
\definecolor{currentfill}{rgb}{0.000000,0.000000,0.000000}%
\pgfsetfillcolor{currentfill}%
\pgfsetlinewidth{0.000000pt}%
\definecolor{currentstroke}{rgb}{0.000000,0.000000,0.000000}%
\pgfsetstrokecolor{currentstroke}%
\pgfsetstrokeopacity{0.000000}%
\pgfsetdash{}{0pt}%
\pgfpathmoveto{\pgfqpoint{6.144900in}{0.660000in}}%
\pgfpathlineto{\pgfqpoint{6.151100in}{0.660000in}}%
\pgfpathlineto{\pgfqpoint{6.151100in}{3.244482in}}%
\pgfpathlineto{\pgfqpoint{6.144900in}{3.244482in}}%
\pgfpathlineto{\pgfqpoint{6.144900in}{0.660000in}}%
\pgfpathclose%
\pgfusepath{fill}%
\end{pgfscope}%
\begin{pgfscope}%
\pgfpathrectangle{\pgfqpoint{1.250000in}{0.660000in}}{\pgfqpoint{7.750000in}{4.620000in}}%
\pgfusepath{clip}%
\pgfsetbuttcap%
\pgfsetmiterjoin%
\definecolor{currentfill}{rgb}{0.000000,0.000000,0.000000}%
\pgfsetfillcolor{currentfill}%
\pgfsetlinewidth{0.000000pt}%
\definecolor{currentstroke}{rgb}{0.000000,0.000000,0.000000}%
\pgfsetstrokecolor{currentstroke}%
\pgfsetstrokeopacity{0.000000}%
\pgfsetdash{}{0pt}%
\pgfpathmoveto{\pgfqpoint{6.152650in}{0.660000in}}%
\pgfpathlineto{\pgfqpoint{6.158850in}{0.660000in}}%
\pgfpathlineto{\pgfqpoint{6.158850in}{3.230504in}}%
\pgfpathlineto{\pgfqpoint{6.152650in}{3.230504in}}%
\pgfpathlineto{\pgfqpoint{6.152650in}{0.660000in}}%
\pgfpathclose%
\pgfusepath{fill}%
\end{pgfscope}%
\begin{pgfscope}%
\pgfpathrectangle{\pgfqpoint{1.250000in}{0.660000in}}{\pgfqpoint{7.750000in}{4.620000in}}%
\pgfusepath{clip}%
\pgfsetbuttcap%
\pgfsetmiterjoin%
\definecolor{currentfill}{rgb}{0.000000,0.000000,0.000000}%
\pgfsetfillcolor{currentfill}%
\pgfsetlinewidth{0.000000pt}%
\definecolor{currentstroke}{rgb}{0.000000,0.000000,0.000000}%
\pgfsetstrokecolor{currentstroke}%
\pgfsetstrokeopacity{0.000000}%
\pgfsetdash{}{0pt}%
\pgfpathmoveto{\pgfqpoint{6.160400in}{0.660000in}}%
\pgfpathlineto{\pgfqpoint{6.166600in}{0.660000in}}%
\pgfpathlineto{\pgfqpoint{6.166600in}{3.254317in}}%
\pgfpathlineto{\pgfqpoint{6.160400in}{3.254317in}}%
\pgfpathlineto{\pgfqpoint{6.160400in}{0.660000in}}%
\pgfpathclose%
\pgfusepath{fill}%
\end{pgfscope}%
\begin{pgfscope}%
\pgfpathrectangle{\pgfqpoint{1.250000in}{0.660000in}}{\pgfqpoint{7.750000in}{4.620000in}}%
\pgfusepath{clip}%
\pgfsetbuttcap%
\pgfsetmiterjoin%
\definecolor{currentfill}{rgb}{0.000000,0.000000,0.000000}%
\pgfsetfillcolor{currentfill}%
\pgfsetlinewidth{0.000000pt}%
\definecolor{currentstroke}{rgb}{0.000000,0.000000,0.000000}%
\pgfsetstrokecolor{currentstroke}%
\pgfsetstrokeopacity{0.000000}%
\pgfsetdash{}{0pt}%
\pgfpathmoveto{\pgfqpoint{6.168150in}{0.660000in}}%
\pgfpathlineto{\pgfqpoint{6.174350in}{0.660000in}}%
\pgfpathlineto{\pgfqpoint{6.174350in}{3.230212in}}%
\pgfpathlineto{\pgfqpoint{6.168150in}{3.230212in}}%
\pgfpathlineto{\pgfqpoint{6.168150in}{0.660000in}}%
\pgfpathclose%
\pgfusepath{fill}%
\end{pgfscope}%
\begin{pgfscope}%
\pgfpathrectangle{\pgfqpoint{1.250000in}{0.660000in}}{\pgfqpoint{7.750000in}{4.620000in}}%
\pgfusepath{clip}%
\pgfsetbuttcap%
\pgfsetmiterjoin%
\definecolor{currentfill}{rgb}{0.000000,0.000000,0.000000}%
\pgfsetfillcolor{currentfill}%
\pgfsetlinewidth{0.000000pt}%
\definecolor{currentstroke}{rgb}{0.000000,0.000000,0.000000}%
\pgfsetstrokecolor{currentstroke}%
\pgfsetstrokeopacity{0.000000}%
\pgfsetdash{}{0pt}%
\pgfpathmoveto{\pgfqpoint{6.175900in}{0.660000in}}%
\pgfpathlineto{\pgfqpoint{6.182100in}{0.660000in}}%
\pgfpathlineto{\pgfqpoint{6.182100in}{3.235064in}}%
\pgfpathlineto{\pgfqpoint{6.175900in}{3.235064in}}%
\pgfpathlineto{\pgfqpoint{6.175900in}{0.660000in}}%
\pgfpathclose%
\pgfusepath{fill}%
\end{pgfscope}%
\begin{pgfscope}%
\pgfpathrectangle{\pgfqpoint{1.250000in}{0.660000in}}{\pgfqpoint{7.750000in}{4.620000in}}%
\pgfusepath{clip}%
\pgfsetbuttcap%
\pgfsetmiterjoin%
\definecolor{currentfill}{rgb}{0.000000,0.000000,0.000000}%
\pgfsetfillcolor{currentfill}%
\pgfsetlinewidth{0.000000pt}%
\definecolor{currentstroke}{rgb}{0.000000,0.000000,0.000000}%
\pgfsetstrokecolor{currentstroke}%
\pgfsetstrokeopacity{0.000000}%
\pgfsetdash{}{0pt}%
\pgfpathmoveto{\pgfqpoint{6.183650in}{0.660000in}}%
\pgfpathlineto{\pgfqpoint{6.189850in}{0.660000in}}%
\pgfpathlineto{\pgfqpoint{6.189850in}{3.245152in}}%
\pgfpathlineto{\pgfqpoint{6.183650in}{3.245152in}}%
\pgfpathlineto{\pgfqpoint{6.183650in}{0.660000in}}%
\pgfpathclose%
\pgfusepath{fill}%
\end{pgfscope}%
\begin{pgfscope}%
\pgfpathrectangle{\pgfqpoint{1.250000in}{0.660000in}}{\pgfqpoint{7.750000in}{4.620000in}}%
\pgfusepath{clip}%
\pgfsetbuttcap%
\pgfsetmiterjoin%
\definecolor{currentfill}{rgb}{0.000000,0.000000,0.000000}%
\pgfsetfillcolor{currentfill}%
\pgfsetlinewidth{0.000000pt}%
\definecolor{currentstroke}{rgb}{0.000000,0.000000,0.000000}%
\pgfsetstrokecolor{currentstroke}%
\pgfsetstrokeopacity{0.000000}%
\pgfsetdash{}{0pt}%
\pgfpathmoveto{\pgfqpoint{6.191400in}{0.660000in}}%
\pgfpathlineto{\pgfqpoint{6.197600in}{0.660000in}}%
\pgfpathlineto{\pgfqpoint{6.197600in}{3.229775in}}%
\pgfpathlineto{\pgfqpoint{6.191400in}{3.229775in}}%
\pgfpathlineto{\pgfqpoint{6.191400in}{0.660000in}}%
\pgfpathclose%
\pgfusepath{fill}%
\end{pgfscope}%
\begin{pgfscope}%
\pgfpathrectangle{\pgfqpoint{1.250000in}{0.660000in}}{\pgfqpoint{7.750000in}{4.620000in}}%
\pgfusepath{clip}%
\pgfsetbuttcap%
\pgfsetmiterjoin%
\definecolor{currentfill}{rgb}{0.000000,0.000000,0.000000}%
\pgfsetfillcolor{currentfill}%
\pgfsetlinewidth{0.000000pt}%
\definecolor{currentstroke}{rgb}{0.000000,0.000000,0.000000}%
\pgfsetstrokecolor{currentstroke}%
\pgfsetstrokeopacity{0.000000}%
\pgfsetdash{}{0pt}%
\pgfpathmoveto{\pgfqpoint{6.199150in}{0.660000in}}%
\pgfpathlineto{\pgfqpoint{6.205350in}{0.660000in}}%
\pgfpathlineto{\pgfqpoint{6.205350in}{3.229630in}}%
\pgfpathlineto{\pgfqpoint{6.199150in}{3.229630in}}%
\pgfpathlineto{\pgfqpoint{6.199150in}{0.660000in}}%
\pgfpathclose%
\pgfusepath{fill}%
\end{pgfscope}%
\begin{pgfscope}%
\pgfpathrectangle{\pgfqpoint{1.250000in}{0.660000in}}{\pgfqpoint{7.750000in}{4.620000in}}%
\pgfusepath{clip}%
\pgfsetbuttcap%
\pgfsetmiterjoin%
\definecolor{currentfill}{rgb}{0.000000,0.000000,0.000000}%
\pgfsetfillcolor{currentfill}%
\pgfsetlinewidth{0.000000pt}%
\definecolor{currentstroke}{rgb}{0.000000,0.000000,0.000000}%
\pgfsetstrokecolor{currentstroke}%
\pgfsetstrokeopacity{0.000000}%
\pgfsetdash{}{0pt}%
\pgfpathmoveto{\pgfqpoint{6.206900in}{0.660000in}}%
\pgfpathlineto{\pgfqpoint{6.213100in}{0.660000in}}%
\pgfpathlineto{\pgfqpoint{6.213100in}{3.229486in}}%
\pgfpathlineto{\pgfqpoint{6.206900in}{3.229486in}}%
\pgfpathlineto{\pgfqpoint{6.206900in}{0.660000in}}%
\pgfpathclose%
\pgfusepath{fill}%
\end{pgfscope}%
\begin{pgfscope}%
\pgfpathrectangle{\pgfqpoint{1.250000in}{0.660000in}}{\pgfqpoint{7.750000in}{4.620000in}}%
\pgfusepath{clip}%
\pgfsetbuttcap%
\pgfsetmiterjoin%
\definecolor{currentfill}{rgb}{0.000000,0.000000,0.000000}%
\pgfsetfillcolor{currentfill}%
\pgfsetlinewidth{0.000000pt}%
\definecolor{currentstroke}{rgb}{0.000000,0.000000,0.000000}%
\pgfsetstrokecolor{currentstroke}%
\pgfsetstrokeopacity{0.000000}%
\pgfsetdash{}{0pt}%
\pgfpathmoveto{\pgfqpoint{6.214650in}{0.660000in}}%
\pgfpathlineto{\pgfqpoint{6.220850in}{0.660000in}}%
\pgfpathlineto{\pgfqpoint{6.220850in}{3.229342in}}%
\pgfpathlineto{\pgfqpoint{6.214650in}{3.229342in}}%
\pgfpathlineto{\pgfqpoint{6.214650in}{0.660000in}}%
\pgfpathclose%
\pgfusepath{fill}%
\end{pgfscope}%
\begin{pgfscope}%
\pgfpathrectangle{\pgfqpoint{1.250000in}{0.660000in}}{\pgfqpoint{7.750000in}{4.620000in}}%
\pgfusepath{clip}%
\pgfsetbuttcap%
\pgfsetmiterjoin%
\definecolor{currentfill}{rgb}{0.000000,0.000000,0.000000}%
\pgfsetfillcolor{currentfill}%
\pgfsetlinewidth{0.000000pt}%
\definecolor{currentstroke}{rgb}{0.000000,0.000000,0.000000}%
\pgfsetstrokecolor{currentstroke}%
\pgfsetstrokeopacity{0.000000}%
\pgfsetdash{}{0pt}%
\pgfpathmoveto{\pgfqpoint{6.222400in}{0.660000in}}%
\pgfpathlineto{\pgfqpoint{6.228600in}{0.660000in}}%
\pgfpathlineto{\pgfqpoint{6.228600in}{3.229198in}}%
\pgfpathlineto{\pgfqpoint{6.222400in}{3.229198in}}%
\pgfpathlineto{\pgfqpoint{6.222400in}{0.660000in}}%
\pgfpathclose%
\pgfusepath{fill}%
\end{pgfscope}%
\begin{pgfscope}%
\pgfpathrectangle{\pgfqpoint{1.250000in}{0.660000in}}{\pgfqpoint{7.750000in}{4.620000in}}%
\pgfusepath{clip}%
\pgfsetbuttcap%
\pgfsetmiterjoin%
\definecolor{currentfill}{rgb}{0.000000,0.000000,0.000000}%
\pgfsetfillcolor{currentfill}%
\pgfsetlinewidth{0.000000pt}%
\definecolor{currentstroke}{rgb}{0.000000,0.000000,0.000000}%
\pgfsetstrokecolor{currentstroke}%
\pgfsetstrokeopacity{0.000000}%
\pgfsetdash{}{0pt}%
\pgfpathmoveto{\pgfqpoint{6.230150in}{0.660000in}}%
\pgfpathlineto{\pgfqpoint{6.236350in}{0.660000in}}%
\pgfpathlineto{\pgfqpoint{6.236350in}{3.229055in}}%
\pgfpathlineto{\pgfqpoint{6.230150in}{3.229055in}}%
\pgfpathlineto{\pgfqpoint{6.230150in}{0.660000in}}%
\pgfpathclose%
\pgfusepath{fill}%
\end{pgfscope}%
\begin{pgfscope}%
\pgfpathrectangle{\pgfqpoint{1.250000in}{0.660000in}}{\pgfqpoint{7.750000in}{4.620000in}}%
\pgfusepath{clip}%
\pgfsetbuttcap%
\pgfsetmiterjoin%
\definecolor{currentfill}{rgb}{0.000000,0.000000,0.000000}%
\pgfsetfillcolor{currentfill}%
\pgfsetlinewidth{0.000000pt}%
\definecolor{currentstroke}{rgb}{0.000000,0.000000,0.000000}%
\pgfsetstrokecolor{currentstroke}%
\pgfsetstrokeopacity{0.000000}%
\pgfsetdash{}{0pt}%
\pgfpathmoveto{\pgfqpoint{6.237900in}{0.660000in}}%
\pgfpathlineto{\pgfqpoint{6.244100in}{0.660000in}}%
\pgfpathlineto{\pgfqpoint{6.244100in}{3.228912in}}%
\pgfpathlineto{\pgfqpoint{6.237900in}{3.228912in}}%
\pgfpathlineto{\pgfqpoint{6.237900in}{0.660000in}}%
\pgfpathclose%
\pgfusepath{fill}%
\end{pgfscope}%
\begin{pgfscope}%
\pgfpathrectangle{\pgfqpoint{1.250000in}{0.660000in}}{\pgfqpoint{7.750000in}{4.620000in}}%
\pgfusepath{clip}%
\pgfsetbuttcap%
\pgfsetmiterjoin%
\definecolor{currentfill}{rgb}{0.000000,0.000000,0.000000}%
\pgfsetfillcolor{currentfill}%
\pgfsetlinewidth{0.000000pt}%
\definecolor{currentstroke}{rgb}{0.000000,0.000000,0.000000}%
\pgfsetstrokecolor{currentstroke}%
\pgfsetstrokeopacity{0.000000}%
\pgfsetdash{}{0pt}%
\pgfpathmoveto{\pgfqpoint{6.245650in}{0.660000in}}%
\pgfpathlineto{\pgfqpoint{6.251850in}{0.660000in}}%
\pgfpathlineto{\pgfqpoint{6.251850in}{3.228770in}}%
\pgfpathlineto{\pgfqpoint{6.245650in}{3.228770in}}%
\pgfpathlineto{\pgfqpoint{6.245650in}{0.660000in}}%
\pgfpathclose%
\pgfusepath{fill}%
\end{pgfscope}%
\begin{pgfscope}%
\pgfpathrectangle{\pgfqpoint{1.250000in}{0.660000in}}{\pgfqpoint{7.750000in}{4.620000in}}%
\pgfusepath{clip}%
\pgfsetbuttcap%
\pgfsetmiterjoin%
\definecolor{currentfill}{rgb}{0.000000,0.000000,0.000000}%
\pgfsetfillcolor{currentfill}%
\pgfsetlinewidth{0.000000pt}%
\definecolor{currentstroke}{rgb}{0.000000,0.000000,0.000000}%
\pgfsetstrokecolor{currentstroke}%
\pgfsetstrokeopacity{0.000000}%
\pgfsetdash{}{0pt}%
\pgfpathmoveto{\pgfqpoint{6.253400in}{0.660000in}}%
\pgfpathlineto{\pgfqpoint{6.259600in}{0.660000in}}%
\pgfpathlineto{\pgfqpoint{6.259600in}{3.246540in}}%
\pgfpathlineto{\pgfqpoint{6.253400in}{3.246540in}}%
\pgfpathlineto{\pgfqpoint{6.253400in}{0.660000in}}%
\pgfpathclose%
\pgfusepath{fill}%
\end{pgfscope}%
\begin{pgfscope}%
\pgfpathrectangle{\pgfqpoint{1.250000in}{0.660000in}}{\pgfqpoint{7.750000in}{4.620000in}}%
\pgfusepath{clip}%
\pgfsetbuttcap%
\pgfsetmiterjoin%
\definecolor{currentfill}{rgb}{0.000000,0.000000,0.000000}%
\pgfsetfillcolor{currentfill}%
\pgfsetlinewidth{0.000000pt}%
\definecolor{currentstroke}{rgb}{0.000000,0.000000,0.000000}%
\pgfsetstrokecolor{currentstroke}%
\pgfsetstrokeopacity{0.000000}%
\pgfsetdash{}{0pt}%
\pgfpathmoveto{\pgfqpoint{6.261150in}{0.660000in}}%
\pgfpathlineto{\pgfqpoint{6.267350in}{0.660000in}}%
\pgfpathlineto{\pgfqpoint{6.267350in}{3.252817in}}%
\pgfpathlineto{\pgfqpoint{6.261150in}{3.252817in}}%
\pgfpathlineto{\pgfqpoint{6.261150in}{0.660000in}}%
\pgfpathclose%
\pgfusepath{fill}%
\end{pgfscope}%
\begin{pgfscope}%
\pgfpathrectangle{\pgfqpoint{1.250000in}{0.660000in}}{\pgfqpoint{7.750000in}{4.620000in}}%
\pgfusepath{clip}%
\pgfsetbuttcap%
\pgfsetmiterjoin%
\definecolor{currentfill}{rgb}{0.000000,0.000000,0.000000}%
\pgfsetfillcolor{currentfill}%
\pgfsetlinewidth{0.000000pt}%
\definecolor{currentstroke}{rgb}{0.000000,0.000000,0.000000}%
\pgfsetstrokecolor{currentstroke}%
\pgfsetstrokeopacity{0.000000}%
\pgfsetdash{}{0pt}%
\pgfpathmoveto{\pgfqpoint{6.268900in}{0.660000in}}%
\pgfpathlineto{\pgfqpoint{6.275100in}{0.660000in}}%
\pgfpathlineto{\pgfqpoint{6.275100in}{3.233641in}}%
\pgfpathlineto{\pgfqpoint{6.268900in}{3.233641in}}%
\pgfpathlineto{\pgfqpoint{6.268900in}{0.660000in}}%
\pgfpathclose%
\pgfusepath{fill}%
\end{pgfscope}%
\begin{pgfscope}%
\pgfpathrectangle{\pgfqpoint{1.250000in}{0.660000in}}{\pgfqpoint{7.750000in}{4.620000in}}%
\pgfusepath{clip}%
\pgfsetbuttcap%
\pgfsetmiterjoin%
\definecolor{currentfill}{rgb}{0.000000,0.000000,0.000000}%
\pgfsetfillcolor{currentfill}%
\pgfsetlinewidth{0.000000pt}%
\definecolor{currentstroke}{rgb}{0.000000,0.000000,0.000000}%
\pgfsetstrokecolor{currentstroke}%
\pgfsetstrokeopacity{0.000000}%
\pgfsetdash{}{0pt}%
\pgfpathmoveto{\pgfqpoint{6.276650in}{0.660000in}}%
\pgfpathlineto{\pgfqpoint{6.282850in}{0.660000in}}%
\pgfpathlineto{\pgfqpoint{6.282850in}{3.228202in}}%
\pgfpathlineto{\pgfqpoint{6.276650in}{3.228202in}}%
\pgfpathlineto{\pgfqpoint{6.276650in}{0.660000in}}%
\pgfpathclose%
\pgfusepath{fill}%
\end{pgfscope}%
\begin{pgfscope}%
\pgfpathrectangle{\pgfqpoint{1.250000in}{0.660000in}}{\pgfqpoint{7.750000in}{4.620000in}}%
\pgfusepath{clip}%
\pgfsetbuttcap%
\pgfsetmiterjoin%
\definecolor{currentfill}{rgb}{0.000000,0.000000,0.000000}%
\pgfsetfillcolor{currentfill}%
\pgfsetlinewidth{0.000000pt}%
\definecolor{currentstroke}{rgb}{0.000000,0.000000,0.000000}%
\pgfsetstrokecolor{currentstroke}%
\pgfsetstrokeopacity{0.000000}%
\pgfsetdash{}{0pt}%
\pgfpathmoveto{\pgfqpoint{6.284400in}{0.660000in}}%
\pgfpathlineto{\pgfqpoint{6.290600in}{0.660000in}}%
\pgfpathlineto{\pgfqpoint{6.290600in}{3.228061in}}%
\pgfpathlineto{\pgfqpoint{6.284400in}{3.228061in}}%
\pgfpathlineto{\pgfqpoint{6.284400in}{0.660000in}}%
\pgfpathclose%
\pgfusepath{fill}%
\end{pgfscope}%
\begin{pgfscope}%
\pgfpathrectangle{\pgfqpoint{1.250000in}{0.660000in}}{\pgfqpoint{7.750000in}{4.620000in}}%
\pgfusepath{clip}%
\pgfsetbuttcap%
\pgfsetmiterjoin%
\definecolor{currentfill}{rgb}{0.000000,0.000000,0.000000}%
\pgfsetfillcolor{currentfill}%
\pgfsetlinewidth{0.000000pt}%
\definecolor{currentstroke}{rgb}{0.000000,0.000000,0.000000}%
\pgfsetstrokecolor{currentstroke}%
\pgfsetstrokeopacity{0.000000}%
\pgfsetdash{}{0pt}%
\pgfpathmoveto{\pgfqpoint{6.292150in}{0.660000in}}%
\pgfpathlineto{\pgfqpoint{6.298350in}{0.660000in}}%
\pgfpathlineto{\pgfqpoint{6.298350in}{3.227920in}}%
\pgfpathlineto{\pgfqpoint{6.292150in}{3.227920in}}%
\pgfpathlineto{\pgfqpoint{6.292150in}{0.660000in}}%
\pgfpathclose%
\pgfusepath{fill}%
\end{pgfscope}%
\begin{pgfscope}%
\pgfpathrectangle{\pgfqpoint{1.250000in}{0.660000in}}{\pgfqpoint{7.750000in}{4.620000in}}%
\pgfusepath{clip}%
\pgfsetbuttcap%
\pgfsetmiterjoin%
\definecolor{currentfill}{rgb}{0.000000,0.000000,0.000000}%
\pgfsetfillcolor{currentfill}%
\pgfsetlinewidth{0.000000pt}%
\definecolor{currentstroke}{rgb}{0.000000,0.000000,0.000000}%
\pgfsetstrokecolor{currentstroke}%
\pgfsetstrokeopacity{0.000000}%
\pgfsetdash{}{0pt}%
\pgfpathmoveto{\pgfqpoint{6.299900in}{0.660000in}}%
\pgfpathlineto{\pgfqpoint{6.306100in}{0.660000in}}%
\pgfpathlineto{\pgfqpoint{6.306100in}{3.238693in}}%
\pgfpathlineto{\pgfqpoint{6.299900in}{3.238693in}}%
\pgfpathlineto{\pgfqpoint{6.299900in}{0.660000in}}%
\pgfpathclose%
\pgfusepath{fill}%
\end{pgfscope}%
\begin{pgfscope}%
\pgfpathrectangle{\pgfqpoint{1.250000in}{0.660000in}}{\pgfqpoint{7.750000in}{4.620000in}}%
\pgfusepath{clip}%
\pgfsetbuttcap%
\pgfsetmiterjoin%
\definecolor{currentfill}{rgb}{0.000000,0.000000,0.000000}%
\pgfsetfillcolor{currentfill}%
\pgfsetlinewidth{0.000000pt}%
\definecolor{currentstroke}{rgb}{0.000000,0.000000,0.000000}%
\pgfsetstrokecolor{currentstroke}%
\pgfsetstrokeopacity{0.000000}%
\pgfsetdash{}{0pt}%
\pgfpathmoveto{\pgfqpoint{6.307650in}{0.660000in}}%
\pgfpathlineto{\pgfqpoint{6.313850in}{0.660000in}}%
\pgfpathlineto{\pgfqpoint{6.313850in}{3.261782in}}%
\pgfpathlineto{\pgfqpoint{6.307650in}{3.261782in}}%
\pgfpathlineto{\pgfqpoint{6.307650in}{0.660000in}}%
\pgfpathclose%
\pgfusepath{fill}%
\end{pgfscope}%
\begin{pgfscope}%
\pgfpathrectangle{\pgfqpoint{1.250000in}{0.660000in}}{\pgfqpoint{7.750000in}{4.620000in}}%
\pgfusepath{clip}%
\pgfsetbuttcap%
\pgfsetmiterjoin%
\definecolor{currentfill}{rgb}{0.000000,0.000000,0.000000}%
\pgfsetfillcolor{currentfill}%
\pgfsetlinewidth{0.000000pt}%
\definecolor{currentstroke}{rgb}{0.000000,0.000000,0.000000}%
\pgfsetstrokecolor{currentstroke}%
\pgfsetstrokeopacity{0.000000}%
\pgfsetdash{}{0pt}%
\pgfpathmoveto{\pgfqpoint{6.315400in}{0.660000in}}%
\pgfpathlineto{\pgfqpoint{6.321600in}{0.660000in}}%
\pgfpathlineto{\pgfqpoint{6.321600in}{3.227498in}}%
\pgfpathlineto{\pgfqpoint{6.315400in}{3.227498in}}%
\pgfpathlineto{\pgfqpoint{6.315400in}{0.660000in}}%
\pgfpathclose%
\pgfusepath{fill}%
\end{pgfscope}%
\begin{pgfscope}%
\pgfpathrectangle{\pgfqpoint{1.250000in}{0.660000in}}{\pgfqpoint{7.750000in}{4.620000in}}%
\pgfusepath{clip}%
\pgfsetbuttcap%
\pgfsetmiterjoin%
\definecolor{currentfill}{rgb}{0.000000,0.000000,0.000000}%
\pgfsetfillcolor{currentfill}%
\pgfsetlinewidth{0.000000pt}%
\definecolor{currentstroke}{rgb}{0.000000,0.000000,0.000000}%
\pgfsetstrokecolor{currentstroke}%
\pgfsetstrokeopacity{0.000000}%
\pgfsetdash{}{0pt}%
\pgfpathmoveto{\pgfqpoint{6.323150in}{0.660000in}}%
\pgfpathlineto{\pgfqpoint{6.329350in}{0.660000in}}%
\pgfpathlineto{\pgfqpoint{6.329350in}{3.243212in}}%
\pgfpathlineto{\pgfqpoint{6.323150in}{3.243212in}}%
\pgfpathlineto{\pgfqpoint{6.323150in}{0.660000in}}%
\pgfpathclose%
\pgfusepath{fill}%
\end{pgfscope}%
\begin{pgfscope}%
\pgfpathrectangle{\pgfqpoint{1.250000in}{0.660000in}}{\pgfqpoint{7.750000in}{4.620000in}}%
\pgfusepath{clip}%
\pgfsetbuttcap%
\pgfsetmiterjoin%
\definecolor{currentfill}{rgb}{0.000000,0.000000,0.000000}%
\pgfsetfillcolor{currentfill}%
\pgfsetlinewidth{0.000000pt}%
\definecolor{currentstroke}{rgb}{0.000000,0.000000,0.000000}%
\pgfsetstrokecolor{currentstroke}%
\pgfsetstrokeopacity{0.000000}%
\pgfsetdash{}{0pt}%
\pgfpathmoveto{\pgfqpoint{6.330900in}{0.660000in}}%
\pgfpathlineto{\pgfqpoint{6.337100in}{0.660000in}}%
\pgfpathlineto{\pgfqpoint{6.337100in}{3.227218in}}%
\pgfpathlineto{\pgfqpoint{6.330900in}{3.227218in}}%
\pgfpathlineto{\pgfqpoint{6.330900in}{0.660000in}}%
\pgfpathclose%
\pgfusepath{fill}%
\end{pgfscope}%
\begin{pgfscope}%
\pgfpathrectangle{\pgfqpoint{1.250000in}{0.660000in}}{\pgfqpoint{7.750000in}{4.620000in}}%
\pgfusepath{clip}%
\pgfsetbuttcap%
\pgfsetmiterjoin%
\definecolor{currentfill}{rgb}{0.000000,0.000000,0.000000}%
\pgfsetfillcolor{currentfill}%
\pgfsetlinewidth{0.000000pt}%
\definecolor{currentstroke}{rgb}{0.000000,0.000000,0.000000}%
\pgfsetstrokecolor{currentstroke}%
\pgfsetstrokeopacity{0.000000}%
\pgfsetdash{}{0pt}%
\pgfpathmoveto{\pgfqpoint{6.338650in}{0.660000in}}%
\pgfpathlineto{\pgfqpoint{6.344850in}{0.660000in}}%
\pgfpathlineto{\pgfqpoint{6.344850in}{3.227078in}}%
\pgfpathlineto{\pgfqpoint{6.338650in}{3.227078in}}%
\pgfpathlineto{\pgfqpoint{6.338650in}{0.660000in}}%
\pgfpathclose%
\pgfusepath{fill}%
\end{pgfscope}%
\begin{pgfscope}%
\pgfpathrectangle{\pgfqpoint{1.250000in}{0.660000in}}{\pgfqpoint{7.750000in}{4.620000in}}%
\pgfusepath{clip}%
\pgfsetbuttcap%
\pgfsetmiterjoin%
\definecolor{currentfill}{rgb}{0.000000,0.000000,0.000000}%
\pgfsetfillcolor{currentfill}%
\pgfsetlinewidth{0.000000pt}%
\definecolor{currentstroke}{rgb}{0.000000,0.000000,0.000000}%
\pgfsetstrokecolor{currentstroke}%
\pgfsetstrokeopacity{0.000000}%
\pgfsetdash{}{0pt}%
\pgfpathmoveto{\pgfqpoint{6.346400in}{0.660000in}}%
\pgfpathlineto{\pgfqpoint{6.352600in}{0.660000in}}%
\pgfpathlineto{\pgfqpoint{6.352600in}{3.226938in}}%
\pgfpathlineto{\pgfqpoint{6.346400in}{3.226938in}}%
\pgfpathlineto{\pgfqpoint{6.346400in}{0.660000in}}%
\pgfpathclose%
\pgfusepath{fill}%
\end{pgfscope}%
\begin{pgfscope}%
\pgfpathrectangle{\pgfqpoint{1.250000in}{0.660000in}}{\pgfqpoint{7.750000in}{4.620000in}}%
\pgfusepath{clip}%
\pgfsetbuttcap%
\pgfsetmiterjoin%
\definecolor{currentfill}{rgb}{0.000000,0.000000,0.000000}%
\pgfsetfillcolor{currentfill}%
\pgfsetlinewidth{0.000000pt}%
\definecolor{currentstroke}{rgb}{0.000000,0.000000,0.000000}%
\pgfsetstrokecolor{currentstroke}%
\pgfsetstrokeopacity{0.000000}%
\pgfsetdash{}{0pt}%
\pgfpathmoveto{\pgfqpoint{6.354150in}{0.660000in}}%
\pgfpathlineto{\pgfqpoint{6.360350in}{0.660000in}}%
\pgfpathlineto{\pgfqpoint{6.360350in}{3.226798in}}%
\pgfpathlineto{\pgfqpoint{6.354150in}{3.226798in}}%
\pgfpathlineto{\pgfqpoint{6.354150in}{0.660000in}}%
\pgfpathclose%
\pgfusepath{fill}%
\end{pgfscope}%
\begin{pgfscope}%
\pgfpathrectangle{\pgfqpoint{1.250000in}{0.660000in}}{\pgfqpoint{7.750000in}{4.620000in}}%
\pgfusepath{clip}%
\pgfsetbuttcap%
\pgfsetmiterjoin%
\definecolor{currentfill}{rgb}{0.000000,0.000000,0.000000}%
\pgfsetfillcolor{currentfill}%
\pgfsetlinewidth{0.000000pt}%
\definecolor{currentstroke}{rgb}{0.000000,0.000000,0.000000}%
\pgfsetstrokecolor{currentstroke}%
\pgfsetstrokeopacity{0.000000}%
\pgfsetdash{}{0pt}%
\pgfpathmoveto{\pgfqpoint{6.361900in}{0.660000in}}%
\pgfpathlineto{\pgfqpoint{6.368100in}{0.660000in}}%
\pgfpathlineto{\pgfqpoint{6.368100in}{3.226658in}}%
\pgfpathlineto{\pgfqpoint{6.361900in}{3.226658in}}%
\pgfpathlineto{\pgfqpoint{6.361900in}{0.660000in}}%
\pgfpathclose%
\pgfusepath{fill}%
\end{pgfscope}%
\begin{pgfscope}%
\pgfpathrectangle{\pgfqpoint{1.250000in}{0.660000in}}{\pgfqpoint{7.750000in}{4.620000in}}%
\pgfusepath{clip}%
\pgfsetbuttcap%
\pgfsetmiterjoin%
\definecolor{currentfill}{rgb}{0.000000,0.000000,0.000000}%
\pgfsetfillcolor{currentfill}%
\pgfsetlinewidth{0.000000pt}%
\definecolor{currentstroke}{rgb}{0.000000,0.000000,0.000000}%
\pgfsetstrokecolor{currentstroke}%
\pgfsetstrokeopacity{0.000000}%
\pgfsetdash{}{0pt}%
\pgfpathmoveto{\pgfqpoint{6.369650in}{0.660000in}}%
\pgfpathlineto{\pgfqpoint{6.375850in}{0.660000in}}%
\pgfpathlineto{\pgfqpoint{6.375850in}{3.226518in}}%
\pgfpathlineto{\pgfqpoint{6.369650in}{3.226518in}}%
\pgfpathlineto{\pgfqpoint{6.369650in}{0.660000in}}%
\pgfpathclose%
\pgfusepath{fill}%
\end{pgfscope}%
\begin{pgfscope}%
\pgfpathrectangle{\pgfqpoint{1.250000in}{0.660000in}}{\pgfqpoint{7.750000in}{4.620000in}}%
\pgfusepath{clip}%
\pgfsetbuttcap%
\pgfsetmiterjoin%
\definecolor{currentfill}{rgb}{0.000000,0.000000,0.000000}%
\pgfsetfillcolor{currentfill}%
\pgfsetlinewidth{0.000000pt}%
\definecolor{currentstroke}{rgb}{0.000000,0.000000,0.000000}%
\pgfsetstrokecolor{currentstroke}%
\pgfsetstrokeopacity{0.000000}%
\pgfsetdash{}{0pt}%
\pgfpathmoveto{\pgfqpoint{6.377400in}{0.660000in}}%
\pgfpathlineto{\pgfqpoint{6.383600in}{0.660000in}}%
\pgfpathlineto{\pgfqpoint{6.383600in}{3.226378in}}%
\pgfpathlineto{\pgfqpoint{6.377400in}{3.226378in}}%
\pgfpathlineto{\pgfqpoint{6.377400in}{0.660000in}}%
\pgfpathclose%
\pgfusepath{fill}%
\end{pgfscope}%
\begin{pgfscope}%
\pgfpathrectangle{\pgfqpoint{1.250000in}{0.660000in}}{\pgfqpoint{7.750000in}{4.620000in}}%
\pgfusepath{clip}%
\pgfsetbuttcap%
\pgfsetmiterjoin%
\definecolor{currentfill}{rgb}{0.000000,0.000000,0.000000}%
\pgfsetfillcolor{currentfill}%
\pgfsetlinewidth{0.000000pt}%
\definecolor{currentstroke}{rgb}{0.000000,0.000000,0.000000}%
\pgfsetstrokecolor{currentstroke}%
\pgfsetstrokeopacity{0.000000}%
\pgfsetdash{}{0pt}%
\pgfpathmoveto{\pgfqpoint{6.385150in}{0.660000in}}%
\pgfpathlineto{\pgfqpoint{6.391350in}{0.660000in}}%
\pgfpathlineto{\pgfqpoint{6.391350in}{3.226238in}}%
\pgfpathlineto{\pgfqpoint{6.385150in}{3.226238in}}%
\pgfpathlineto{\pgfqpoint{6.385150in}{0.660000in}}%
\pgfpathclose%
\pgfusepath{fill}%
\end{pgfscope}%
\begin{pgfscope}%
\pgfpathrectangle{\pgfqpoint{1.250000in}{0.660000in}}{\pgfqpoint{7.750000in}{4.620000in}}%
\pgfusepath{clip}%
\pgfsetbuttcap%
\pgfsetmiterjoin%
\definecolor{currentfill}{rgb}{0.000000,0.000000,0.000000}%
\pgfsetfillcolor{currentfill}%
\pgfsetlinewidth{0.000000pt}%
\definecolor{currentstroke}{rgb}{0.000000,0.000000,0.000000}%
\pgfsetstrokecolor{currentstroke}%
\pgfsetstrokeopacity{0.000000}%
\pgfsetdash{}{0pt}%
\pgfpathmoveto{\pgfqpoint{6.392900in}{0.660000in}}%
\pgfpathlineto{\pgfqpoint{6.399100in}{0.660000in}}%
\pgfpathlineto{\pgfqpoint{6.399100in}{3.235880in}}%
\pgfpathlineto{\pgfqpoint{6.392900in}{3.235880in}}%
\pgfpathlineto{\pgfqpoint{6.392900in}{0.660000in}}%
\pgfpathclose%
\pgfusepath{fill}%
\end{pgfscope}%
\begin{pgfscope}%
\pgfpathrectangle{\pgfqpoint{1.250000in}{0.660000in}}{\pgfqpoint{7.750000in}{4.620000in}}%
\pgfusepath{clip}%
\pgfsetbuttcap%
\pgfsetmiterjoin%
\definecolor{currentfill}{rgb}{0.000000,0.000000,0.000000}%
\pgfsetfillcolor{currentfill}%
\pgfsetlinewidth{0.000000pt}%
\definecolor{currentstroke}{rgb}{0.000000,0.000000,0.000000}%
\pgfsetstrokecolor{currentstroke}%
\pgfsetstrokeopacity{0.000000}%
\pgfsetdash{}{0pt}%
\pgfpathmoveto{\pgfqpoint{6.400650in}{0.660000in}}%
\pgfpathlineto{\pgfqpoint{6.406850in}{0.660000in}}%
\pgfpathlineto{\pgfqpoint{6.406850in}{3.242335in}}%
\pgfpathlineto{\pgfqpoint{6.400650in}{3.242335in}}%
\pgfpathlineto{\pgfqpoint{6.400650in}{0.660000in}}%
\pgfpathclose%
\pgfusepath{fill}%
\end{pgfscope}%
\begin{pgfscope}%
\pgfpathrectangle{\pgfqpoint{1.250000in}{0.660000in}}{\pgfqpoint{7.750000in}{4.620000in}}%
\pgfusepath{clip}%
\pgfsetbuttcap%
\pgfsetmiterjoin%
\definecolor{currentfill}{rgb}{0.000000,0.000000,0.000000}%
\pgfsetfillcolor{currentfill}%
\pgfsetlinewidth{0.000000pt}%
\definecolor{currentstroke}{rgb}{0.000000,0.000000,0.000000}%
\pgfsetstrokecolor{currentstroke}%
\pgfsetstrokeopacity{0.000000}%
\pgfsetdash{}{0pt}%
\pgfpathmoveto{\pgfqpoint{6.408400in}{0.660000in}}%
\pgfpathlineto{\pgfqpoint{6.414600in}{0.660000in}}%
\pgfpathlineto{\pgfqpoint{6.414600in}{3.225816in}}%
\pgfpathlineto{\pgfqpoint{6.408400in}{3.225816in}}%
\pgfpathlineto{\pgfqpoint{6.408400in}{0.660000in}}%
\pgfpathclose%
\pgfusepath{fill}%
\end{pgfscope}%
\begin{pgfscope}%
\pgfpathrectangle{\pgfqpoint{1.250000in}{0.660000in}}{\pgfqpoint{7.750000in}{4.620000in}}%
\pgfusepath{clip}%
\pgfsetbuttcap%
\pgfsetmiterjoin%
\definecolor{currentfill}{rgb}{0.000000,0.000000,0.000000}%
\pgfsetfillcolor{currentfill}%
\pgfsetlinewidth{0.000000pt}%
\definecolor{currentstroke}{rgb}{0.000000,0.000000,0.000000}%
\pgfsetstrokecolor{currentstroke}%
\pgfsetstrokeopacity{0.000000}%
\pgfsetdash{}{0pt}%
\pgfpathmoveto{\pgfqpoint{6.416150in}{0.660000in}}%
\pgfpathlineto{\pgfqpoint{6.422350in}{0.660000in}}%
\pgfpathlineto{\pgfqpoint{6.422350in}{3.225675in}}%
\pgfpathlineto{\pgfqpoint{6.416150in}{3.225675in}}%
\pgfpathlineto{\pgfqpoint{6.416150in}{0.660000in}}%
\pgfpathclose%
\pgfusepath{fill}%
\end{pgfscope}%
\begin{pgfscope}%
\pgfpathrectangle{\pgfqpoint{1.250000in}{0.660000in}}{\pgfqpoint{7.750000in}{4.620000in}}%
\pgfusepath{clip}%
\pgfsetbuttcap%
\pgfsetmiterjoin%
\definecolor{currentfill}{rgb}{0.000000,0.000000,0.000000}%
\pgfsetfillcolor{currentfill}%
\pgfsetlinewidth{0.000000pt}%
\definecolor{currentstroke}{rgb}{0.000000,0.000000,0.000000}%
\pgfsetstrokecolor{currentstroke}%
\pgfsetstrokeopacity{0.000000}%
\pgfsetdash{}{0pt}%
\pgfpathmoveto{\pgfqpoint{6.423900in}{0.660000in}}%
\pgfpathlineto{\pgfqpoint{6.430100in}{0.660000in}}%
\pgfpathlineto{\pgfqpoint{6.430100in}{3.225534in}}%
\pgfpathlineto{\pgfqpoint{6.423900in}{3.225534in}}%
\pgfpathlineto{\pgfqpoint{6.423900in}{0.660000in}}%
\pgfpathclose%
\pgfusepath{fill}%
\end{pgfscope}%
\begin{pgfscope}%
\pgfpathrectangle{\pgfqpoint{1.250000in}{0.660000in}}{\pgfqpoint{7.750000in}{4.620000in}}%
\pgfusepath{clip}%
\pgfsetbuttcap%
\pgfsetmiterjoin%
\definecolor{currentfill}{rgb}{0.000000,0.000000,0.000000}%
\pgfsetfillcolor{currentfill}%
\pgfsetlinewidth{0.000000pt}%
\definecolor{currentstroke}{rgb}{0.000000,0.000000,0.000000}%
\pgfsetstrokecolor{currentstroke}%
\pgfsetstrokeopacity{0.000000}%
\pgfsetdash{}{0pt}%
\pgfpathmoveto{\pgfqpoint{6.431650in}{0.660000in}}%
\pgfpathlineto{\pgfqpoint{6.437850in}{0.660000in}}%
\pgfpathlineto{\pgfqpoint{6.437850in}{3.225392in}}%
\pgfpathlineto{\pgfqpoint{6.431650in}{3.225392in}}%
\pgfpathlineto{\pgfqpoint{6.431650in}{0.660000in}}%
\pgfpathclose%
\pgfusepath{fill}%
\end{pgfscope}%
\begin{pgfscope}%
\pgfpathrectangle{\pgfqpoint{1.250000in}{0.660000in}}{\pgfqpoint{7.750000in}{4.620000in}}%
\pgfusepath{clip}%
\pgfsetbuttcap%
\pgfsetmiterjoin%
\definecolor{currentfill}{rgb}{0.000000,0.000000,0.000000}%
\pgfsetfillcolor{currentfill}%
\pgfsetlinewidth{0.000000pt}%
\definecolor{currentstroke}{rgb}{0.000000,0.000000,0.000000}%
\pgfsetstrokecolor{currentstroke}%
\pgfsetstrokeopacity{0.000000}%
\pgfsetdash{}{0pt}%
\pgfpathmoveto{\pgfqpoint{6.439400in}{0.660000in}}%
\pgfpathlineto{\pgfqpoint{6.445600in}{0.660000in}}%
\pgfpathlineto{\pgfqpoint{6.445600in}{3.225250in}}%
\pgfpathlineto{\pgfqpoint{6.439400in}{3.225250in}}%
\pgfpathlineto{\pgfqpoint{6.439400in}{0.660000in}}%
\pgfpathclose%
\pgfusepath{fill}%
\end{pgfscope}%
\begin{pgfscope}%
\pgfpathrectangle{\pgfqpoint{1.250000in}{0.660000in}}{\pgfqpoint{7.750000in}{4.620000in}}%
\pgfusepath{clip}%
\pgfsetbuttcap%
\pgfsetmiterjoin%
\definecolor{currentfill}{rgb}{0.000000,0.000000,0.000000}%
\pgfsetfillcolor{currentfill}%
\pgfsetlinewidth{0.000000pt}%
\definecolor{currentstroke}{rgb}{0.000000,0.000000,0.000000}%
\pgfsetstrokecolor{currentstroke}%
\pgfsetstrokeopacity{0.000000}%
\pgfsetdash{}{0pt}%
\pgfpathmoveto{\pgfqpoint{6.447150in}{0.660000in}}%
\pgfpathlineto{\pgfqpoint{6.453350in}{0.660000in}}%
\pgfpathlineto{\pgfqpoint{6.453350in}{3.225107in}}%
\pgfpathlineto{\pgfqpoint{6.447150in}{3.225107in}}%
\pgfpathlineto{\pgfqpoint{6.447150in}{0.660000in}}%
\pgfpathclose%
\pgfusepath{fill}%
\end{pgfscope}%
\begin{pgfscope}%
\pgfpathrectangle{\pgfqpoint{1.250000in}{0.660000in}}{\pgfqpoint{7.750000in}{4.620000in}}%
\pgfusepath{clip}%
\pgfsetbuttcap%
\pgfsetmiterjoin%
\definecolor{currentfill}{rgb}{0.000000,0.000000,0.000000}%
\pgfsetfillcolor{currentfill}%
\pgfsetlinewidth{0.000000pt}%
\definecolor{currentstroke}{rgb}{0.000000,0.000000,0.000000}%
\pgfsetstrokecolor{currentstroke}%
\pgfsetstrokeopacity{0.000000}%
\pgfsetdash{}{0pt}%
\pgfpathmoveto{\pgfqpoint{6.454900in}{0.660000in}}%
\pgfpathlineto{\pgfqpoint{6.461100in}{0.660000in}}%
\pgfpathlineto{\pgfqpoint{6.461100in}{3.224964in}}%
\pgfpathlineto{\pgfqpoint{6.454900in}{3.224964in}}%
\pgfpathlineto{\pgfqpoint{6.454900in}{0.660000in}}%
\pgfpathclose%
\pgfusepath{fill}%
\end{pgfscope}%
\begin{pgfscope}%
\pgfpathrectangle{\pgfqpoint{1.250000in}{0.660000in}}{\pgfqpoint{7.750000in}{4.620000in}}%
\pgfusepath{clip}%
\pgfsetbuttcap%
\pgfsetmiterjoin%
\definecolor{currentfill}{rgb}{0.000000,0.000000,0.000000}%
\pgfsetfillcolor{currentfill}%
\pgfsetlinewidth{0.000000pt}%
\definecolor{currentstroke}{rgb}{0.000000,0.000000,0.000000}%
\pgfsetstrokecolor{currentstroke}%
\pgfsetstrokeopacity{0.000000}%
\pgfsetdash{}{0pt}%
\pgfpathmoveto{\pgfqpoint{6.462650in}{0.660000in}}%
\pgfpathlineto{\pgfqpoint{6.468850in}{0.660000in}}%
\pgfpathlineto{\pgfqpoint{6.468850in}{3.224820in}}%
\pgfpathlineto{\pgfqpoint{6.462650in}{3.224820in}}%
\pgfpathlineto{\pgfqpoint{6.462650in}{0.660000in}}%
\pgfpathclose%
\pgfusepath{fill}%
\end{pgfscope}%
\begin{pgfscope}%
\pgfpathrectangle{\pgfqpoint{1.250000in}{0.660000in}}{\pgfqpoint{7.750000in}{4.620000in}}%
\pgfusepath{clip}%
\pgfsetbuttcap%
\pgfsetmiterjoin%
\definecolor{currentfill}{rgb}{0.000000,0.000000,0.000000}%
\pgfsetfillcolor{currentfill}%
\pgfsetlinewidth{0.000000pt}%
\definecolor{currentstroke}{rgb}{0.000000,0.000000,0.000000}%
\pgfsetstrokecolor{currentstroke}%
\pgfsetstrokeopacity{0.000000}%
\pgfsetdash{}{0pt}%
\pgfpathmoveto{\pgfqpoint{6.470400in}{0.660000in}}%
\pgfpathlineto{\pgfqpoint{6.476600in}{0.660000in}}%
\pgfpathlineto{\pgfqpoint{6.476600in}{3.224676in}}%
\pgfpathlineto{\pgfqpoint{6.470400in}{3.224676in}}%
\pgfpathlineto{\pgfqpoint{6.470400in}{0.660000in}}%
\pgfpathclose%
\pgfusepath{fill}%
\end{pgfscope}%
\begin{pgfscope}%
\pgfpathrectangle{\pgfqpoint{1.250000in}{0.660000in}}{\pgfqpoint{7.750000in}{4.620000in}}%
\pgfusepath{clip}%
\pgfsetbuttcap%
\pgfsetmiterjoin%
\definecolor{currentfill}{rgb}{0.000000,0.000000,0.000000}%
\pgfsetfillcolor{currentfill}%
\pgfsetlinewidth{0.000000pt}%
\definecolor{currentstroke}{rgb}{0.000000,0.000000,0.000000}%
\pgfsetstrokecolor{currentstroke}%
\pgfsetstrokeopacity{0.000000}%
\pgfsetdash{}{0pt}%
\pgfpathmoveto{\pgfqpoint{6.478150in}{0.660000in}}%
\pgfpathlineto{\pgfqpoint{6.484350in}{0.660000in}}%
\pgfpathlineto{\pgfqpoint{6.484350in}{3.224531in}}%
\pgfpathlineto{\pgfqpoint{6.478150in}{3.224531in}}%
\pgfpathlineto{\pgfqpoint{6.478150in}{0.660000in}}%
\pgfpathclose%
\pgfusepath{fill}%
\end{pgfscope}%
\begin{pgfscope}%
\pgfpathrectangle{\pgfqpoint{1.250000in}{0.660000in}}{\pgfqpoint{7.750000in}{4.620000in}}%
\pgfusepath{clip}%
\pgfsetbuttcap%
\pgfsetmiterjoin%
\definecolor{currentfill}{rgb}{0.000000,0.000000,0.000000}%
\pgfsetfillcolor{currentfill}%
\pgfsetlinewidth{0.000000pt}%
\definecolor{currentstroke}{rgb}{0.000000,0.000000,0.000000}%
\pgfsetstrokecolor{currentstroke}%
\pgfsetstrokeopacity{0.000000}%
\pgfsetdash{}{0pt}%
\pgfpathmoveto{\pgfqpoint{6.485900in}{0.660000in}}%
\pgfpathlineto{\pgfqpoint{6.492100in}{0.660000in}}%
\pgfpathlineto{\pgfqpoint{6.492100in}{3.224385in}}%
\pgfpathlineto{\pgfqpoint{6.485900in}{3.224385in}}%
\pgfpathlineto{\pgfqpoint{6.485900in}{0.660000in}}%
\pgfpathclose%
\pgfusepath{fill}%
\end{pgfscope}%
\begin{pgfscope}%
\pgfpathrectangle{\pgfqpoint{1.250000in}{0.660000in}}{\pgfqpoint{7.750000in}{4.620000in}}%
\pgfusepath{clip}%
\pgfsetbuttcap%
\pgfsetmiterjoin%
\definecolor{currentfill}{rgb}{0.000000,0.000000,0.000000}%
\pgfsetfillcolor{currentfill}%
\pgfsetlinewidth{0.000000pt}%
\definecolor{currentstroke}{rgb}{0.000000,0.000000,0.000000}%
\pgfsetstrokecolor{currentstroke}%
\pgfsetstrokeopacity{0.000000}%
\pgfsetdash{}{0pt}%
\pgfpathmoveto{\pgfqpoint{6.493650in}{0.660000in}}%
\pgfpathlineto{\pgfqpoint{6.499850in}{0.660000in}}%
\pgfpathlineto{\pgfqpoint{6.499850in}{3.224238in}}%
\pgfpathlineto{\pgfqpoint{6.493650in}{3.224238in}}%
\pgfpathlineto{\pgfqpoint{6.493650in}{0.660000in}}%
\pgfpathclose%
\pgfusepath{fill}%
\end{pgfscope}%
\begin{pgfscope}%
\pgfpathrectangle{\pgfqpoint{1.250000in}{0.660000in}}{\pgfqpoint{7.750000in}{4.620000in}}%
\pgfusepath{clip}%
\pgfsetbuttcap%
\pgfsetmiterjoin%
\definecolor{currentfill}{rgb}{0.000000,0.000000,0.000000}%
\pgfsetfillcolor{currentfill}%
\pgfsetlinewidth{0.000000pt}%
\definecolor{currentstroke}{rgb}{0.000000,0.000000,0.000000}%
\pgfsetstrokecolor{currentstroke}%
\pgfsetstrokeopacity{0.000000}%
\pgfsetdash{}{0pt}%
\pgfpathmoveto{\pgfqpoint{6.501400in}{0.660000in}}%
\pgfpathlineto{\pgfqpoint{6.507600in}{0.660000in}}%
\pgfpathlineto{\pgfqpoint{6.507600in}{3.243741in}}%
\pgfpathlineto{\pgfqpoint{6.501400in}{3.243741in}}%
\pgfpathlineto{\pgfqpoint{6.501400in}{0.660000in}}%
\pgfpathclose%
\pgfusepath{fill}%
\end{pgfscope}%
\begin{pgfscope}%
\pgfpathrectangle{\pgfqpoint{1.250000in}{0.660000in}}{\pgfqpoint{7.750000in}{4.620000in}}%
\pgfusepath{clip}%
\pgfsetbuttcap%
\pgfsetmiterjoin%
\definecolor{currentfill}{rgb}{0.000000,0.000000,0.000000}%
\pgfsetfillcolor{currentfill}%
\pgfsetlinewidth{0.000000pt}%
\definecolor{currentstroke}{rgb}{0.000000,0.000000,0.000000}%
\pgfsetstrokecolor{currentstroke}%
\pgfsetstrokeopacity{0.000000}%
\pgfsetdash{}{0pt}%
\pgfpathmoveto{\pgfqpoint{6.509150in}{0.660000in}}%
\pgfpathlineto{\pgfqpoint{6.515350in}{0.660000in}}%
\pgfpathlineto{\pgfqpoint{6.515350in}{3.223943in}}%
\pgfpathlineto{\pgfqpoint{6.509150in}{3.223943in}}%
\pgfpathlineto{\pgfqpoint{6.509150in}{0.660000in}}%
\pgfpathclose%
\pgfusepath{fill}%
\end{pgfscope}%
\begin{pgfscope}%
\pgfpathrectangle{\pgfqpoint{1.250000in}{0.660000in}}{\pgfqpoint{7.750000in}{4.620000in}}%
\pgfusepath{clip}%
\pgfsetbuttcap%
\pgfsetmiterjoin%
\definecolor{currentfill}{rgb}{0.000000,0.000000,0.000000}%
\pgfsetfillcolor{currentfill}%
\pgfsetlinewidth{0.000000pt}%
\definecolor{currentstroke}{rgb}{0.000000,0.000000,0.000000}%
\pgfsetstrokecolor{currentstroke}%
\pgfsetstrokeopacity{0.000000}%
\pgfsetdash{}{0pt}%
\pgfpathmoveto{\pgfqpoint{6.516900in}{0.660000in}}%
\pgfpathlineto{\pgfqpoint{6.523100in}{0.660000in}}%
\pgfpathlineto{\pgfqpoint{6.523100in}{3.223794in}}%
\pgfpathlineto{\pgfqpoint{6.516900in}{3.223794in}}%
\pgfpathlineto{\pgfqpoint{6.516900in}{0.660000in}}%
\pgfpathclose%
\pgfusepath{fill}%
\end{pgfscope}%
\begin{pgfscope}%
\pgfpathrectangle{\pgfqpoint{1.250000in}{0.660000in}}{\pgfqpoint{7.750000in}{4.620000in}}%
\pgfusepath{clip}%
\pgfsetbuttcap%
\pgfsetmiterjoin%
\definecolor{currentfill}{rgb}{0.000000,0.000000,0.000000}%
\pgfsetfillcolor{currentfill}%
\pgfsetlinewidth{0.000000pt}%
\definecolor{currentstroke}{rgb}{0.000000,0.000000,0.000000}%
\pgfsetstrokecolor{currentstroke}%
\pgfsetstrokeopacity{0.000000}%
\pgfsetdash{}{0pt}%
\pgfpathmoveto{\pgfqpoint{6.524650in}{0.660000in}}%
\pgfpathlineto{\pgfqpoint{6.530850in}{0.660000in}}%
\pgfpathlineto{\pgfqpoint{6.530850in}{3.223643in}}%
\pgfpathlineto{\pgfqpoint{6.524650in}{3.223643in}}%
\pgfpathlineto{\pgfqpoint{6.524650in}{0.660000in}}%
\pgfpathclose%
\pgfusepath{fill}%
\end{pgfscope}%
\begin{pgfscope}%
\pgfpathrectangle{\pgfqpoint{1.250000in}{0.660000in}}{\pgfqpoint{7.750000in}{4.620000in}}%
\pgfusepath{clip}%
\pgfsetbuttcap%
\pgfsetmiterjoin%
\definecolor{currentfill}{rgb}{0.000000,0.000000,0.000000}%
\pgfsetfillcolor{currentfill}%
\pgfsetlinewidth{0.000000pt}%
\definecolor{currentstroke}{rgb}{0.000000,0.000000,0.000000}%
\pgfsetstrokecolor{currentstroke}%
\pgfsetstrokeopacity{0.000000}%
\pgfsetdash{}{0pt}%
\pgfpathmoveto{\pgfqpoint{6.532400in}{0.660000in}}%
\pgfpathlineto{\pgfqpoint{6.538600in}{0.660000in}}%
\pgfpathlineto{\pgfqpoint{6.538600in}{3.231112in}}%
\pgfpathlineto{\pgfqpoint{6.532400in}{3.231112in}}%
\pgfpathlineto{\pgfqpoint{6.532400in}{0.660000in}}%
\pgfpathclose%
\pgfusepath{fill}%
\end{pgfscope}%
\begin{pgfscope}%
\pgfpathrectangle{\pgfqpoint{1.250000in}{0.660000in}}{\pgfqpoint{7.750000in}{4.620000in}}%
\pgfusepath{clip}%
\pgfsetbuttcap%
\pgfsetmiterjoin%
\definecolor{currentfill}{rgb}{0.000000,0.000000,0.000000}%
\pgfsetfillcolor{currentfill}%
\pgfsetlinewidth{0.000000pt}%
\definecolor{currentstroke}{rgb}{0.000000,0.000000,0.000000}%
\pgfsetstrokecolor{currentstroke}%
\pgfsetstrokeopacity{0.000000}%
\pgfsetdash{}{0pt}%
\pgfpathmoveto{\pgfqpoint{6.540150in}{0.660000in}}%
\pgfpathlineto{\pgfqpoint{6.546350in}{0.660000in}}%
\pgfpathlineto{\pgfqpoint{6.546350in}{3.253929in}}%
\pgfpathlineto{\pgfqpoint{6.540150in}{3.253929in}}%
\pgfpathlineto{\pgfqpoint{6.540150in}{0.660000in}}%
\pgfpathclose%
\pgfusepath{fill}%
\end{pgfscope}%
\begin{pgfscope}%
\pgfpathrectangle{\pgfqpoint{1.250000in}{0.660000in}}{\pgfqpoint{7.750000in}{4.620000in}}%
\pgfusepath{clip}%
\pgfsetbuttcap%
\pgfsetmiterjoin%
\definecolor{currentfill}{rgb}{0.000000,0.000000,0.000000}%
\pgfsetfillcolor{currentfill}%
\pgfsetlinewidth{0.000000pt}%
\definecolor{currentstroke}{rgb}{0.000000,0.000000,0.000000}%
\pgfsetstrokecolor{currentstroke}%
\pgfsetstrokeopacity{0.000000}%
\pgfsetdash{}{0pt}%
\pgfpathmoveto{\pgfqpoint{6.547900in}{0.660000in}}%
\pgfpathlineto{\pgfqpoint{6.554100in}{0.660000in}}%
\pgfpathlineto{\pgfqpoint{6.554100in}{3.223186in}}%
\pgfpathlineto{\pgfqpoint{6.547900in}{3.223186in}}%
\pgfpathlineto{\pgfqpoint{6.547900in}{0.660000in}}%
\pgfpathclose%
\pgfusepath{fill}%
\end{pgfscope}%
\begin{pgfscope}%
\pgfpathrectangle{\pgfqpoint{1.250000in}{0.660000in}}{\pgfqpoint{7.750000in}{4.620000in}}%
\pgfusepath{clip}%
\pgfsetbuttcap%
\pgfsetmiterjoin%
\definecolor{currentfill}{rgb}{0.000000,0.000000,0.000000}%
\pgfsetfillcolor{currentfill}%
\pgfsetlinewidth{0.000000pt}%
\definecolor{currentstroke}{rgb}{0.000000,0.000000,0.000000}%
\pgfsetstrokecolor{currentstroke}%
\pgfsetstrokeopacity{0.000000}%
\pgfsetdash{}{0pt}%
\pgfpathmoveto{\pgfqpoint{6.555650in}{0.660000in}}%
\pgfpathlineto{\pgfqpoint{6.561850in}{0.660000in}}%
\pgfpathlineto{\pgfqpoint{6.561850in}{3.223031in}}%
\pgfpathlineto{\pgfqpoint{6.555650in}{3.223031in}}%
\pgfpathlineto{\pgfqpoint{6.555650in}{0.660000in}}%
\pgfpathclose%
\pgfusepath{fill}%
\end{pgfscope}%
\begin{pgfscope}%
\pgfpathrectangle{\pgfqpoint{1.250000in}{0.660000in}}{\pgfqpoint{7.750000in}{4.620000in}}%
\pgfusepath{clip}%
\pgfsetbuttcap%
\pgfsetmiterjoin%
\definecolor{currentfill}{rgb}{0.000000,0.000000,0.000000}%
\pgfsetfillcolor{currentfill}%
\pgfsetlinewidth{0.000000pt}%
\definecolor{currentstroke}{rgb}{0.000000,0.000000,0.000000}%
\pgfsetstrokecolor{currentstroke}%
\pgfsetstrokeopacity{0.000000}%
\pgfsetdash{}{0pt}%
\pgfpathmoveto{\pgfqpoint{6.563400in}{0.660000in}}%
\pgfpathlineto{\pgfqpoint{6.569600in}{0.660000in}}%
\pgfpathlineto{\pgfqpoint{6.569600in}{3.235411in}}%
\pgfpathlineto{\pgfqpoint{6.563400in}{3.235411in}}%
\pgfpathlineto{\pgfqpoint{6.563400in}{0.660000in}}%
\pgfpathclose%
\pgfusepath{fill}%
\end{pgfscope}%
\begin{pgfscope}%
\pgfpathrectangle{\pgfqpoint{1.250000in}{0.660000in}}{\pgfqpoint{7.750000in}{4.620000in}}%
\pgfusepath{clip}%
\pgfsetbuttcap%
\pgfsetmiterjoin%
\definecolor{currentfill}{rgb}{0.000000,0.000000,0.000000}%
\pgfsetfillcolor{currentfill}%
\pgfsetlinewidth{0.000000pt}%
\definecolor{currentstroke}{rgb}{0.000000,0.000000,0.000000}%
\pgfsetstrokecolor{currentstroke}%
\pgfsetstrokeopacity{0.000000}%
\pgfsetdash{}{0pt}%
\pgfpathmoveto{\pgfqpoint{6.571150in}{0.660000in}}%
\pgfpathlineto{\pgfqpoint{6.577350in}{0.660000in}}%
\pgfpathlineto{\pgfqpoint{6.577350in}{3.222717in}}%
\pgfpathlineto{\pgfqpoint{6.571150in}{3.222717in}}%
\pgfpathlineto{\pgfqpoint{6.571150in}{0.660000in}}%
\pgfpathclose%
\pgfusepath{fill}%
\end{pgfscope}%
\begin{pgfscope}%
\pgfpathrectangle{\pgfqpoint{1.250000in}{0.660000in}}{\pgfqpoint{7.750000in}{4.620000in}}%
\pgfusepath{clip}%
\pgfsetbuttcap%
\pgfsetmiterjoin%
\definecolor{currentfill}{rgb}{0.000000,0.000000,0.000000}%
\pgfsetfillcolor{currentfill}%
\pgfsetlinewidth{0.000000pt}%
\definecolor{currentstroke}{rgb}{0.000000,0.000000,0.000000}%
\pgfsetstrokecolor{currentstroke}%
\pgfsetstrokeopacity{0.000000}%
\pgfsetdash{}{0pt}%
\pgfpathmoveto{\pgfqpoint{6.578900in}{0.660000in}}%
\pgfpathlineto{\pgfqpoint{6.585100in}{0.660000in}}%
\pgfpathlineto{\pgfqpoint{6.585100in}{3.227885in}}%
\pgfpathlineto{\pgfqpoint{6.578900in}{3.227885in}}%
\pgfpathlineto{\pgfqpoint{6.578900in}{0.660000in}}%
\pgfpathclose%
\pgfusepath{fill}%
\end{pgfscope}%
\begin{pgfscope}%
\pgfpathrectangle{\pgfqpoint{1.250000in}{0.660000in}}{\pgfqpoint{7.750000in}{4.620000in}}%
\pgfusepath{clip}%
\pgfsetbuttcap%
\pgfsetmiterjoin%
\definecolor{currentfill}{rgb}{0.000000,0.000000,0.000000}%
\pgfsetfillcolor{currentfill}%
\pgfsetlinewidth{0.000000pt}%
\definecolor{currentstroke}{rgb}{0.000000,0.000000,0.000000}%
\pgfsetstrokecolor{currentstroke}%
\pgfsetstrokeopacity{0.000000}%
\pgfsetdash{}{0pt}%
\pgfpathmoveto{\pgfqpoint{6.586650in}{0.660000in}}%
\pgfpathlineto{\pgfqpoint{6.592850in}{0.660000in}}%
\pgfpathlineto{\pgfqpoint{6.592850in}{3.222397in}}%
\pgfpathlineto{\pgfqpoint{6.586650in}{3.222397in}}%
\pgfpathlineto{\pgfqpoint{6.586650in}{0.660000in}}%
\pgfpathclose%
\pgfusepath{fill}%
\end{pgfscope}%
\begin{pgfscope}%
\pgfpathrectangle{\pgfqpoint{1.250000in}{0.660000in}}{\pgfqpoint{7.750000in}{4.620000in}}%
\pgfusepath{clip}%
\pgfsetbuttcap%
\pgfsetmiterjoin%
\definecolor{currentfill}{rgb}{0.000000,0.000000,0.000000}%
\pgfsetfillcolor{currentfill}%
\pgfsetlinewidth{0.000000pt}%
\definecolor{currentstroke}{rgb}{0.000000,0.000000,0.000000}%
\pgfsetstrokecolor{currentstroke}%
\pgfsetstrokeopacity{0.000000}%
\pgfsetdash{}{0pt}%
\pgfpathmoveto{\pgfqpoint{6.594400in}{0.660000in}}%
\pgfpathlineto{\pgfqpoint{6.600600in}{0.660000in}}%
\pgfpathlineto{\pgfqpoint{6.600600in}{3.222234in}}%
\pgfpathlineto{\pgfqpoint{6.594400in}{3.222234in}}%
\pgfpathlineto{\pgfqpoint{6.594400in}{0.660000in}}%
\pgfpathclose%
\pgfusepath{fill}%
\end{pgfscope}%
\begin{pgfscope}%
\pgfpathrectangle{\pgfqpoint{1.250000in}{0.660000in}}{\pgfqpoint{7.750000in}{4.620000in}}%
\pgfusepath{clip}%
\pgfsetbuttcap%
\pgfsetmiterjoin%
\definecolor{currentfill}{rgb}{0.000000,0.000000,0.000000}%
\pgfsetfillcolor{currentfill}%
\pgfsetlinewidth{0.000000pt}%
\definecolor{currentstroke}{rgb}{0.000000,0.000000,0.000000}%
\pgfsetstrokecolor{currentstroke}%
\pgfsetstrokeopacity{0.000000}%
\pgfsetdash{}{0pt}%
\pgfpathmoveto{\pgfqpoint{6.602150in}{0.660000in}}%
\pgfpathlineto{\pgfqpoint{6.608350in}{0.660000in}}%
\pgfpathlineto{\pgfqpoint{6.608350in}{3.222070in}}%
\pgfpathlineto{\pgfqpoint{6.602150in}{3.222070in}}%
\pgfpathlineto{\pgfqpoint{6.602150in}{0.660000in}}%
\pgfpathclose%
\pgfusepath{fill}%
\end{pgfscope}%
\begin{pgfscope}%
\pgfpathrectangle{\pgfqpoint{1.250000in}{0.660000in}}{\pgfqpoint{7.750000in}{4.620000in}}%
\pgfusepath{clip}%
\pgfsetbuttcap%
\pgfsetmiterjoin%
\definecolor{currentfill}{rgb}{0.000000,0.000000,0.000000}%
\pgfsetfillcolor{currentfill}%
\pgfsetlinewidth{0.000000pt}%
\definecolor{currentstroke}{rgb}{0.000000,0.000000,0.000000}%
\pgfsetstrokecolor{currentstroke}%
\pgfsetstrokeopacity{0.000000}%
\pgfsetdash{}{0pt}%
\pgfpathmoveto{\pgfqpoint{6.609900in}{0.660000in}}%
\pgfpathlineto{\pgfqpoint{6.616100in}{0.660000in}}%
\pgfpathlineto{\pgfqpoint{6.616100in}{3.221904in}}%
\pgfpathlineto{\pgfqpoint{6.609900in}{3.221904in}}%
\pgfpathlineto{\pgfqpoint{6.609900in}{0.660000in}}%
\pgfpathclose%
\pgfusepath{fill}%
\end{pgfscope}%
\begin{pgfscope}%
\pgfpathrectangle{\pgfqpoint{1.250000in}{0.660000in}}{\pgfqpoint{7.750000in}{4.620000in}}%
\pgfusepath{clip}%
\pgfsetbuttcap%
\pgfsetmiterjoin%
\definecolor{currentfill}{rgb}{0.000000,0.000000,0.000000}%
\pgfsetfillcolor{currentfill}%
\pgfsetlinewidth{0.000000pt}%
\definecolor{currentstroke}{rgb}{0.000000,0.000000,0.000000}%
\pgfsetstrokecolor{currentstroke}%
\pgfsetstrokeopacity{0.000000}%
\pgfsetdash{}{0pt}%
\pgfpathmoveto{\pgfqpoint{6.617650in}{0.660000in}}%
\pgfpathlineto{\pgfqpoint{6.623850in}{0.660000in}}%
\pgfpathlineto{\pgfqpoint{6.623850in}{3.221953in}}%
\pgfpathlineto{\pgfqpoint{6.617650in}{3.221953in}}%
\pgfpathlineto{\pgfqpoint{6.617650in}{0.660000in}}%
\pgfpathclose%
\pgfusepath{fill}%
\end{pgfscope}%
\begin{pgfscope}%
\pgfpathrectangle{\pgfqpoint{1.250000in}{0.660000in}}{\pgfqpoint{7.750000in}{4.620000in}}%
\pgfusepath{clip}%
\pgfsetbuttcap%
\pgfsetmiterjoin%
\definecolor{currentfill}{rgb}{0.000000,0.000000,0.000000}%
\pgfsetfillcolor{currentfill}%
\pgfsetlinewidth{0.000000pt}%
\definecolor{currentstroke}{rgb}{0.000000,0.000000,0.000000}%
\pgfsetstrokecolor{currentstroke}%
\pgfsetstrokeopacity{0.000000}%
\pgfsetdash{}{0pt}%
\pgfpathmoveto{\pgfqpoint{6.625400in}{0.660000in}}%
\pgfpathlineto{\pgfqpoint{6.631600in}{0.660000in}}%
\pgfpathlineto{\pgfqpoint{6.631600in}{3.233395in}}%
\pgfpathlineto{\pgfqpoint{6.625400in}{3.233395in}}%
\pgfpathlineto{\pgfqpoint{6.625400in}{0.660000in}}%
\pgfpathclose%
\pgfusepath{fill}%
\end{pgfscope}%
\begin{pgfscope}%
\pgfpathrectangle{\pgfqpoint{1.250000in}{0.660000in}}{\pgfqpoint{7.750000in}{4.620000in}}%
\pgfusepath{clip}%
\pgfsetbuttcap%
\pgfsetmiterjoin%
\definecolor{currentfill}{rgb}{0.000000,0.000000,0.000000}%
\pgfsetfillcolor{currentfill}%
\pgfsetlinewidth{0.000000pt}%
\definecolor{currentstroke}{rgb}{0.000000,0.000000,0.000000}%
\pgfsetstrokecolor{currentstroke}%
\pgfsetstrokeopacity{0.000000}%
\pgfsetdash{}{0pt}%
\pgfpathmoveto{\pgfqpoint{6.633150in}{0.660000in}}%
\pgfpathlineto{\pgfqpoint{6.639350in}{0.660000in}}%
\pgfpathlineto{\pgfqpoint{6.639350in}{3.221486in}}%
\pgfpathlineto{\pgfqpoint{6.633150in}{3.221486in}}%
\pgfpathlineto{\pgfqpoint{6.633150in}{0.660000in}}%
\pgfpathclose%
\pgfusepath{fill}%
\end{pgfscope}%
\begin{pgfscope}%
\pgfpathrectangle{\pgfqpoint{1.250000in}{0.660000in}}{\pgfqpoint{7.750000in}{4.620000in}}%
\pgfusepath{clip}%
\pgfsetbuttcap%
\pgfsetmiterjoin%
\definecolor{currentfill}{rgb}{0.000000,0.000000,0.000000}%
\pgfsetfillcolor{currentfill}%
\pgfsetlinewidth{0.000000pt}%
\definecolor{currentstroke}{rgb}{0.000000,0.000000,0.000000}%
\pgfsetstrokecolor{currentstroke}%
\pgfsetstrokeopacity{0.000000}%
\pgfsetdash{}{0pt}%
\pgfpathmoveto{\pgfqpoint{6.640900in}{0.660000in}}%
\pgfpathlineto{\pgfqpoint{6.647100in}{0.660000in}}%
\pgfpathlineto{\pgfqpoint{6.647100in}{3.221397in}}%
\pgfpathlineto{\pgfqpoint{6.640900in}{3.221397in}}%
\pgfpathlineto{\pgfqpoint{6.640900in}{0.660000in}}%
\pgfpathclose%
\pgfusepath{fill}%
\end{pgfscope}%
\begin{pgfscope}%
\pgfpathrectangle{\pgfqpoint{1.250000in}{0.660000in}}{\pgfqpoint{7.750000in}{4.620000in}}%
\pgfusepath{clip}%
\pgfsetbuttcap%
\pgfsetmiterjoin%
\definecolor{currentfill}{rgb}{0.000000,0.000000,0.000000}%
\pgfsetfillcolor{currentfill}%
\pgfsetlinewidth{0.000000pt}%
\definecolor{currentstroke}{rgb}{0.000000,0.000000,0.000000}%
\pgfsetstrokecolor{currentstroke}%
\pgfsetstrokeopacity{0.000000}%
\pgfsetdash{}{0pt}%
\pgfpathmoveto{\pgfqpoint{6.648650in}{0.660000in}}%
\pgfpathlineto{\pgfqpoint{6.654850in}{0.660000in}}%
\pgfpathlineto{\pgfqpoint{6.654850in}{3.221307in}}%
\pgfpathlineto{\pgfqpoint{6.648650in}{3.221307in}}%
\pgfpathlineto{\pgfqpoint{6.648650in}{0.660000in}}%
\pgfpathclose%
\pgfusepath{fill}%
\end{pgfscope}%
\begin{pgfscope}%
\pgfpathrectangle{\pgfqpoint{1.250000in}{0.660000in}}{\pgfqpoint{7.750000in}{4.620000in}}%
\pgfusepath{clip}%
\pgfsetbuttcap%
\pgfsetmiterjoin%
\definecolor{currentfill}{rgb}{0.000000,0.000000,0.000000}%
\pgfsetfillcolor{currentfill}%
\pgfsetlinewidth{0.000000pt}%
\definecolor{currentstroke}{rgb}{0.000000,0.000000,0.000000}%
\pgfsetstrokecolor{currentstroke}%
\pgfsetstrokeopacity{0.000000}%
\pgfsetdash{}{0pt}%
\pgfpathmoveto{\pgfqpoint{6.656400in}{0.660000in}}%
\pgfpathlineto{\pgfqpoint{6.662600in}{0.660000in}}%
\pgfpathlineto{\pgfqpoint{6.662600in}{3.221217in}}%
\pgfpathlineto{\pgfqpoint{6.656400in}{3.221217in}}%
\pgfpathlineto{\pgfqpoint{6.656400in}{0.660000in}}%
\pgfpathclose%
\pgfusepath{fill}%
\end{pgfscope}%
\begin{pgfscope}%
\pgfpathrectangle{\pgfqpoint{1.250000in}{0.660000in}}{\pgfqpoint{7.750000in}{4.620000in}}%
\pgfusepath{clip}%
\pgfsetbuttcap%
\pgfsetmiterjoin%
\definecolor{currentfill}{rgb}{0.000000,0.000000,0.000000}%
\pgfsetfillcolor{currentfill}%
\pgfsetlinewidth{0.000000pt}%
\definecolor{currentstroke}{rgb}{0.000000,0.000000,0.000000}%
\pgfsetstrokecolor{currentstroke}%
\pgfsetstrokeopacity{0.000000}%
\pgfsetdash{}{0pt}%
\pgfpathmoveto{\pgfqpoint{6.664150in}{0.660000in}}%
\pgfpathlineto{\pgfqpoint{6.670350in}{0.660000in}}%
\pgfpathlineto{\pgfqpoint{6.670350in}{3.221125in}}%
\pgfpathlineto{\pgfqpoint{6.664150in}{3.221125in}}%
\pgfpathlineto{\pgfqpoint{6.664150in}{0.660000in}}%
\pgfpathclose%
\pgfusepath{fill}%
\end{pgfscope}%
\begin{pgfscope}%
\pgfpathrectangle{\pgfqpoint{1.250000in}{0.660000in}}{\pgfqpoint{7.750000in}{4.620000in}}%
\pgfusepath{clip}%
\pgfsetbuttcap%
\pgfsetmiterjoin%
\definecolor{currentfill}{rgb}{0.000000,0.000000,0.000000}%
\pgfsetfillcolor{currentfill}%
\pgfsetlinewidth{0.000000pt}%
\definecolor{currentstroke}{rgb}{0.000000,0.000000,0.000000}%
\pgfsetstrokecolor{currentstroke}%
\pgfsetstrokeopacity{0.000000}%
\pgfsetdash{}{0pt}%
\pgfpathmoveto{\pgfqpoint{6.671900in}{0.660000in}}%
\pgfpathlineto{\pgfqpoint{6.678100in}{0.660000in}}%
\pgfpathlineto{\pgfqpoint{6.678100in}{3.221032in}}%
\pgfpathlineto{\pgfqpoint{6.671900in}{3.221032in}}%
\pgfpathlineto{\pgfqpoint{6.671900in}{0.660000in}}%
\pgfpathclose%
\pgfusepath{fill}%
\end{pgfscope}%
\begin{pgfscope}%
\pgfpathrectangle{\pgfqpoint{1.250000in}{0.660000in}}{\pgfqpoint{7.750000in}{4.620000in}}%
\pgfusepath{clip}%
\pgfsetbuttcap%
\pgfsetmiterjoin%
\definecolor{currentfill}{rgb}{0.000000,0.000000,0.000000}%
\pgfsetfillcolor{currentfill}%
\pgfsetlinewidth{0.000000pt}%
\definecolor{currentstroke}{rgb}{0.000000,0.000000,0.000000}%
\pgfsetstrokecolor{currentstroke}%
\pgfsetstrokeopacity{0.000000}%
\pgfsetdash{}{0pt}%
\pgfpathmoveto{\pgfqpoint{6.679650in}{0.660000in}}%
\pgfpathlineto{\pgfqpoint{6.685850in}{0.660000in}}%
\pgfpathlineto{\pgfqpoint{6.685850in}{3.220939in}}%
\pgfpathlineto{\pgfqpoint{6.679650in}{3.220939in}}%
\pgfpathlineto{\pgfqpoint{6.679650in}{0.660000in}}%
\pgfpathclose%
\pgfusepath{fill}%
\end{pgfscope}%
\begin{pgfscope}%
\pgfpathrectangle{\pgfqpoint{1.250000in}{0.660000in}}{\pgfqpoint{7.750000in}{4.620000in}}%
\pgfusepath{clip}%
\pgfsetbuttcap%
\pgfsetmiterjoin%
\definecolor{currentfill}{rgb}{0.000000,0.000000,0.000000}%
\pgfsetfillcolor{currentfill}%
\pgfsetlinewidth{0.000000pt}%
\definecolor{currentstroke}{rgb}{0.000000,0.000000,0.000000}%
\pgfsetstrokecolor{currentstroke}%
\pgfsetstrokeopacity{0.000000}%
\pgfsetdash{}{0pt}%
\pgfpathmoveto{\pgfqpoint{6.687400in}{0.660000in}}%
\pgfpathlineto{\pgfqpoint{6.693600in}{0.660000in}}%
\pgfpathlineto{\pgfqpoint{6.693600in}{3.229278in}}%
\pgfpathlineto{\pgfqpoint{6.687400in}{3.229278in}}%
\pgfpathlineto{\pgfqpoint{6.687400in}{0.660000in}}%
\pgfpathclose%
\pgfusepath{fill}%
\end{pgfscope}%
\begin{pgfscope}%
\pgfpathrectangle{\pgfqpoint{1.250000in}{0.660000in}}{\pgfqpoint{7.750000in}{4.620000in}}%
\pgfusepath{clip}%
\pgfsetbuttcap%
\pgfsetmiterjoin%
\definecolor{currentfill}{rgb}{0.000000,0.000000,0.000000}%
\pgfsetfillcolor{currentfill}%
\pgfsetlinewidth{0.000000pt}%
\definecolor{currentstroke}{rgb}{0.000000,0.000000,0.000000}%
\pgfsetstrokecolor{currentstroke}%
\pgfsetstrokeopacity{0.000000}%
\pgfsetdash{}{0pt}%
\pgfpathmoveto{\pgfqpoint{6.695150in}{0.660000in}}%
\pgfpathlineto{\pgfqpoint{6.701350in}{0.660000in}}%
\pgfpathlineto{\pgfqpoint{6.701350in}{3.220748in}}%
\pgfpathlineto{\pgfqpoint{6.695150in}{3.220748in}}%
\pgfpathlineto{\pgfqpoint{6.695150in}{0.660000in}}%
\pgfpathclose%
\pgfusepath{fill}%
\end{pgfscope}%
\begin{pgfscope}%
\pgfpathrectangle{\pgfqpoint{1.250000in}{0.660000in}}{\pgfqpoint{7.750000in}{4.620000in}}%
\pgfusepath{clip}%
\pgfsetbuttcap%
\pgfsetmiterjoin%
\definecolor{currentfill}{rgb}{0.000000,0.000000,0.000000}%
\pgfsetfillcolor{currentfill}%
\pgfsetlinewidth{0.000000pt}%
\definecolor{currentstroke}{rgb}{0.000000,0.000000,0.000000}%
\pgfsetstrokecolor{currentstroke}%
\pgfsetstrokeopacity{0.000000}%
\pgfsetdash{}{0pt}%
\pgfpathmoveto{\pgfqpoint{6.702900in}{0.660000in}}%
\pgfpathlineto{\pgfqpoint{6.709100in}{0.660000in}}%
\pgfpathlineto{\pgfqpoint{6.709100in}{3.220650in}}%
\pgfpathlineto{\pgfqpoint{6.702900in}{3.220650in}}%
\pgfpathlineto{\pgfqpoint{6.702900in}{0.660000in}}%
\pgfpathclose%
\pgfusepath{fill}%
\end{pgfscope}%
\begin{pgfscope}%
\pgfpathrectangle{\pgfqpoint{1.250000in}{0.660000in}}{\pgfqpoint{7.750000in}{4.620000in}}%
\pgfusepath{clip}%
\pgfsetbuttcap%
\pgfsetmiterjoin%
\definecolor{currentfill}{rgb}{0.000000,0.000000,0.000000}%
\pgfsetfillcolor{currentfill}%
\pgfsetlinewidth{0.000000pt}%
\definecolor{currentstroke}{rgb}{0.000000,0.000000,0.000000}%
\pgfsetstrokecolor{currentstroke}%
\pgfsetstrokeopacity{0.000000}%
\pgfsetdash{}{0pt}%
\pgfpathmoveto{\pgfqpoint{6.710650in}{0.660000in}}%
\pgfpathlineto{\pgfqpoint{6.716850in}{0.660000in}}%
\pgfpathlineto{\pgfqpoint{6.716850in}{3.220551in}}%
\pgfpathlineto{\pgfqpoint{6.710650in}{3.220551in}}%
\pgfpathlineto{\pgfqpoint{6.710650in}{0.660000in}}%
\pgfpathclose%
\pgfusepath{fill}%
\end{pgfscope}%
\begin{pgfscope}%
\pgfpathrectangle{\pgfqpoint{1.250000in}{0.660000in}}{\pgfqpoint{7.750000in}{4.620000in}}%
\pgfusepath{clip}%
\pgfsetbuttcap%
\pgfsetmiterjoin%
\definecolor{currentfill}{rgb}{0.000000,0.000000,0.000000}%
\pgfsetfillcolor{currentfill}%
\pgfsetlinewidth{0.000000pt}%
\definecolor{currentstroke}{rgb}{0.000000,0.000000,0.000000}%
\pgfsetstrokecolor{currentstroke}%
\pgfsetstrokeopacity{0.000000}%
\pgfsetdash{}{0pt}%
\pgfpathmoveto{\pgfqpoint{6.718400in}{0.660000in}}%
\pgfpathlineto{\pgfqpoint{6.724600in}{0.660000in}}%
\pgfpathlineto{\pgfqpoint{6.724600in}{3.220451in}}%
\pgfpathlineto{\pgfqpoint{6.718400in}{3.220451in}}%
\pgfpathlineto{\pgfqpoint{6.718400in}{0.660000in}}%
\pgfpathclose%
\pgfusepath{fill}%
\end{pgfscope}%
\begin{pgfscope}%
\pgfpathrectangle{\pgfqpoint{1.250000in}{0.660000in}}{\pgfqpoint{7.750000in}{4.620000in}}%
\pgfusepath{clip}%
\pgfsetbuttcap%
\pgfsetmiterjoin%
\definecolor{currentfill}{rgb}{0.000000,0.000000,0.000000}%
\pgfsetfillcolor{currentfill}%
\pgfsetlinewidth{0.000000pt}%
\definecolor{currentstroke}{rgb}{0.000000,0.000000,0.000000}%
\pgfsetstrokecolor{currentstroke}%
\pgfsetstrokeopacity{0.000000}%
\pgfsetdash{}{0pt}%
\pgfpathmoveto{\pgfqpoint{6.726150in}{0.660000in}}%
\pgfpathlineto{\pgfqpoint{6.732350in}{0.660000in}}%
\pgfpathlineto{\pgfqpoint{6.732350in}{3.220348in}}%
\pgfpathlineto{\pgfqpoint{6.726150in}{3.220348in}}%
\pgfpathlineto{\pgfqpoint{6.726150in}{0.660000in}}%
\pgfpathclose%
\pgfusepath{fill}%
\end{pgfscope}%
\begin{pgfscope}%
\pgfpathrectangle{\pgfqpoint{1.250000in}{0.660000in}}{\pgfqpoint{7.750000in}{4.620000in}}%
\pgfusepath{clip}%
\pgfsetbuttcap%
\pgfsetmiterjoin%
\definecolor{currentfill}{rgb}{0.000000,0.000000,0.000000}%
\pgfsetfillcolor{currentfill}%
\pgfsetlinewidth{0.000000pt}%
\definecolor{currentstroke}{rgb}{0.000000,0.000000,0.000000}%
\pgfsetstrokecolor{currentstroke}%
\pgfsetstrokeopacity{0.000000}%
\pgfsetdash{}{0pt}%
\pgfpathmoveto{\pgfqpoint{6.733900in}{0.660000in}}%
\pgfpathlineto{\pgfqpoint{6.740100in}{0.660000in}}%
\pgfpathlineto{\pgfqpoint{6.740100in}{3.226217in}}%
\pgfpathlineto{\pgfqpoint{6.733900in}{3.226217in}}%
\pgfpathlineto{\pgfqpoint{6.733900in}{0.660000in}}%
\pgfpathclose%
\pgfusepath{fill}%
\end{pgfscope}%
\begin{pgfscope}%
\pgfpathrectangle{\pgfqpoint{1.250000in}{0.660000in}}{\pgfqpoint{7.750000in}{4.620000in}}%
\pgfusepath{clip}%
\pgfsetbuttcap%
\pgfsetmiterjoin%
\definecolor{currentfill}{rgb}{0.000000,0.000000,0.000000}%
\pgfsetfillcolor{currentfill}%
\pgfsetlinewidth{0.000000pt}%
\definecolor{currentstroke}{rgb}{0.000000,0.000000,0.000000}%
\pgfsetstrokecolor{currentstroke}%
\pgfsetstrokeopacity{0.000000}%
\pgfsetdash{}{0pt}%
\pgfpathmoveto{\pgfqpoint{6.741650in}{0.660000in}}%
\pgfpathlineto{\pgfqpoint{6.747850in}{0.660000in}}%
\pgfpathlineto{\pgfqpoint{6.747850in}{3.223533in}}%
\pgfpathlineto{\pgfqpoint{6.741650in}{3.223533in}}%
\pgfpathlineto{\pgfqpoint{6.741650in}{0.660000in}}%
\pgfpathclose%
\pgfusepath{fill}%
\end{pgfscope}%
\begin{pgfscope}%
\pgfpathrectangle{\pgfqpoint{1.250000in}{0.660000in}}{\pgfqpoint{7.750000in}{4.620000in}}%
\pgfusepath{clip}%
\pgfsetbuttcap%
\pgfsetmiterjoin%
\definecolor{currentfill}{rgb}{0.000000,0.000000,0.000000}%
\pgfsetfillcolor{currentfill}%
\pgfsetlinewidth{0.000000pt}%
\definecolor{currentstroke}{rgb}{0.000000,0.000000,0.000000}%
\pgfsetstrokecolor{currentstroke}%
\pgfsetstrokeopacity{0.000000}%
\pgfsetdash{}{0pt}%
\pgfpathmoveto{\pgfqpoint{6.749400in}{0.660000in}}%
\pgfpathlineto{\pgfqpoint{6.755600in}{0.660000in}}%
\pgfpathlineto{\pgfqpoint{6.755600in}{3.224298in}}%
\pgfpathlineto{\pgfqpoint{6.749400in}{3.224298in}}%
\pgfpathlineto{\pgfqpoint{6.749400in}{0.660000in}}%
\pgfpathclose%
\pgfusepath{fill}%
\end{pgfscope}%
\begin{pgfscope}%
\pgfpathrectangle{\pgfqpoint{1.250000in}{0.660000in}}{\pgfqpoint{7.750000in}{4.620000in}}%
\pgfusepath{clip}%
\pgfsetbuttcap%
\pgfsetmiterjoin%
\definecolor{currentfill}{rgb}{0.000000,0.000000,0.000000}%
\pgfsetfillcolor{currentfill}%
\pgfsetlinewidth{0.000000pt}%
\definecolor{currentstroke}{rgb}{0.000000,0.000000,0.000000}%
\pgfsetstrokecolor{currentstroke}%
\pgfsetstrokeopacity{0.000000}%
\pgfsetdash{}{0pt}%
\pgfpathmoveto{\pgfqpoint{6.757150in}{0.660000in}}%
\pgfpathlineto{\pgfqpoint{6.763350in}{0.660000in}}%
\pgfpathlineto{\pgfqpoint{6.763350in}{3.219922in}}%
\pgfpathlineto{\pgfqpoint{6.757150in}{3.219922in}}%
\pgfpathlineto{\pgfqpoint{6.757150in}{0.660000in}}%
\pgfpathclose%
\pgfusepath{fill}%
\end{pgfscope}%
\begin{pgfscope}%
\pgfpathrectangle{\pgfqpoint{1.250000in}{0.660000in}}{\pgfqpoint{7.750000in}{4.620000in}}%
\pgfusepath{clip}%
\pgfsetbuttcap%
\pgfsetmiterjoin%
\definecolor{currentfill}{rgb}{0.000000,0.000000,0.000000}%
\pgfsetfillcolor{currentfill}%
\pgfsetlinewidth{0.000000pt}%
\definecolor{currentstroke}{rgb}{0.000000,0.000000,0.000000}%
\pgfsetstrokecolor{currentstroke}%
\pgfsetstrokeopacity{0.000000}%
\pgfsetdash{}{0pt}%
\pgfpathmoveto{\pgfqpoint{6.764900in}{0.660000in}}%
\pgfpathlineto{\pgfqpoint{6.771100in}{0.660000in}}%
\pgfpathlineto{\pgfqpoint{6.771100in}{3.219810in}}%
\pgfpathlineto{\pgfqpoint{6.764900in}{3.219810in}}%
\pgfpathlineto{\pgfqpoint{6.764900in}{0.660000in}}%
\pgfpathclose%
\pgfusepath{fill}%
\end{pgfscope}%
\begin{pgfscope}%
\pgfpathrectangle{\pgfqpoint{1.250000in}{0.660000in}}{\pgfqpoint{7.750000in}{4.620000in}}%
\pgfusepath{clip}%
\pgfsetbuttcap%
\pgfsetmiterjoin%
\definecolor{currentfill}{rgb}{0.000000,0.000000,0.000000}%
\pgfsetfillcolor{currentfill}%
\pgfsetlinewidth{0.000000pt}%
\definecolor{currentstroke}{rgb}{0.000000,0.000000,0.000000}%
\pgfsetstrokecolor{currentstroke}%
\pgfsetstrokeopacity{0.000000}%
\pgfsetdash{}{0pt}%
\pgfpathmoveto{\pgfqpoint{6.772650in}{0.660000in}}%
\pgfpathlineto{\pgfqpoint{6.778850in}{0.660000in}}%
\pgfpathlineto{\pgfqpoint{6.778850in}{3.219696in}}%
\pgfpathlineto{\pgfqpoint{6.772650in}{3.219696in}}%
\pgfpathlineto{\pgfqpoint{6.772650in}{0.660000in}}%
\pgfpathclose%
\pgfusepath{fill}%
\end{pgfscope}%
\begin{pgfscope}%
\pgfpathrectangle{\pgfqpoint{1.250000in}{0.660000in}}{\pgfqpoint{7.750000in}{4.620000in}}%
\pgfusepath{clip}%
\pgfsetbuttcap%
\pgfsetmiterjoin%
\definecolor{currentfill}{rgb}{0.000000,0.000000,0.000000}%
\pgfsetfillcolor{currentfill}%
\pgfsetlinewidth{0.000000pt}%
\definecolor{currentstroke}{rgb}{0.000000,0.000000,0.000000}%
\pgfsetstrokecolor{currentstroke}%
\pgfsetstrokeopacity{0.000000}%
\pgfsetdash{}{0pt}%
\pgfpathmoveto{\pgfqpoint{6.780400in}{0.660000in}}%
\pgfpathlineto{\pgfqpoint{6.786600in}{0.660000in}}%
\pgfpathlineto{\pgfqpoint{6.786600in}{3.240988in}}%
\pgfpathlineto{\pgfqpoint{6.780400in}{3.240988in}}%
\pgfpathlineto{\pgfqpoint{6.780400in}{0.660000in}}%
\pgfpathclose%
\pgfusepath{fill}%
\end{pgfscope}%
\begin{pgfscope}%
\pgfpathrectangle{\pgfqpoint{1.250000in}{0.660000in}}{\pgfqpoint{7.750000in}{4.620000in}}%
\pgfusepath{clip}%
\pgfsetbuttcap%
\pgfsetmiterjoin%
\definecolor{currentfill}{rgb}{0.000000,0.000000,0.000000}%
\pgfsetfillcolor{currentfill}%
\pgfsetlinewidth{0.000000pt}%
\definecolor{currentstroke}{rgb}{0.000000,0.000000,0.000000}%
\pgfsetstrokecolor{currentstroke}%
\pgfsetstrokeopacity{0.000000}%
\pgfsetdash{}{0pt}%
\pgfpathmoveto{\pgfqpoint{6.788150in}{0.660000in}}%
\pgfpathlineto{\pgfqpoint{6.794350in}{0.660000in}}%
\pgfpathlineto{\pgfqpoint{6.794350in}{3.219460in}}%
\pgfpathlineto{\pgfqpoint{6.788150in}{3.219460in}}%
\pgfpathlineto{\pgfqpoint{6.788150in}{0.660000in}}%
\pgfpathclose%
\pgfusepath{fill}%
\end{pgfscope}%
\begin{pgfscope}%
\pgfpathrectangle{\pgfqpoint{1.250000in}{0.660000in}}{\pgfqpoint{7.750000in}{4.620000in}}%
\pgfusepath{clip}%
\pgfsetbuttcap%
\pgfsetmiterjoin%
\definecolor{currentfill}{rgb}{0.000000,0.000000,0.000000}%
\pgfsetfillcolor{currentfill}%
\pgfsetlinewidth{0.000000pt}%
\definecolor{currentstroke}{rgb}{0.000000,0.000000,0.000000}%
\pgfsetstrokecolor{currentstroke}%
\pgfsetstrokeopacity{0.000000}%
\pgfsetdash{}{0pt}%
\pgfpathmoveto{\pgfqpoint{6.795900in}{0.660000in}}%
\pgfpathlineto{\pgfqpoint{6.802100in}{0.660000in}}%
\pgfpathlineto{\pgfqpoint{6.802100in}{3.219339in}}%
\pgfpathlineto{\pgfqpoint{6.795900in}{3.219339in}}%
\pgfpathlineto{\pgfqpoint{6.795900in}{0.660000in}}%
\pgfpathclose%
\pgfusepath{fill}%
\end{pgfscope}%
\begin{pgfscope}%
\pgfpathrectangle{\pgfqpoint{1.250000in}{0.660000in}}{\pgfqpoint{7.750000in}{4.620000in}}%
\pgfusepath{clip}%
\pgfsetbuttcap%
\pgfsetmiterjoin%
\definecolor{currentfill}{rgb}{0.000000,0.000000,0.000000}%
\pgfsetfillcolor{currentfill}%
\pgfsetlinewidth{0.000000pt}%
\definecolor{currentstroke}{rgb}{0.000000,0.000000,0.000000}%
\pgfsetstrokecolor{currentstroke}%
\pgfsetstrokeopacity{0.000000}%
\pgfsetdash{}{0pt}%
\pgfpathmoveto{\pgfqpoint{6.803650in}{0.660000in}}%
\pgfpathlineto{\pgfqpoint{6.809850in}{0.660000in}}%
\pgfpathlineto{\pgfqpoint{6.809850in}{3.219215in}}%
\pgfpathlineto{\pgfqpoint{6.803650in}{3.219215in}}%
\pgfpathlineto{\pgfqpoint{6.803650in}{0.660000in}}%
\pgfpathclose%
\pgfusepath{fill}%
\end{pgfscope}%
\begin{pgfscope}%
\pgfpathrectangle{\pgfqpoint{1.250000in}{0.660000in}}{\pgfqpoint{7.750000in}{4.620000in}}%
\pgfusepath{clip}%
\pgfsetbuttcap%
\pgfsetmiterjoin%
\definecolor{currentfill}{rgb}{0.000000,0.000000,0.000000}%
\pgfsetfillcolor{currentfill}%
\pgfsetlinewidth{0.000000pt}%
\definecolor{currentstroke}{rgb}{0.000000,0.000000,0.000000}%
\pgfsetstrokecolor{currentstroke}%
\pgfsetstrokeopacity{0.000000}%
\pgfsetdash{}{0pt}%
\pgfpathmoveto{\pgfqpoint{6.811400in}{0.660000in}}%
\pgfpathlineto{\pgfqpoint{6.817600in}{0.660000in}}%
\pgfpathlineto{\pgfqpoint{6.817600in}{3.219088in}}%
\pgfpathlineto{\pgfqpoint{6.811400in}{3.219088in}}%
\pgfpathlineto{\pgfqpoint{6.811400in}{0.660000in}}%
\pgfpathclose%
\pgfusepath{fill}%
\end{pgfscope}%
\begin{pgfscope}%
\pgfpathrectangle{\pgfqpoint{1.250000in}{0.660000in}}{\pgfqpoint{7.750000in}{4.620000in}}%
\pgfusepath{clip}%
\pgfsetbuttcap%
\pgfsetmiterjoin%
\definecolor{currentfill}{rgb}{0.000000,0.000000,0.000000}%
\pgfsetfillcolor{currentfill}%
\pgfsetlinewidth{0.000000pt}%
\definecolor{currentstroke}{rgb}{0.000000,0.000000,0.000000}%
\pgfsetstrokecolor{currentstroke}%
\pgfsetstrokeopacity{0.000000}%
\pgfsetdash{}{0pt}%
\pgfpathmoveto{\pgfqpoint{6.819150in}{0.660000in}}%
\pgfpathlineto{\pgfqpoint{6.825350in}{0.660000in}}%
\pgfpathlineto{\pgfqpoint{6.825350in}{3.218958in}}%
\pgfpathlineto{\pgfqpoint{6.819150in}{3.218958in}}%
\pgfpathlineto{\pgfqpoint{6.819150in}{0.660000in}}%
\pgfpathclose%
\pgfusepath{fill}%
\end{pgfscope}%
\begin{pgfscope}%
\pgfpathrectangle{\pgfqpoint{1.250000in}{0.660000in}}{\pgfqpoint{7.750000in}{4.620000in}}%
\pgfusepath{clip}%
\pgfsetbuttcap%
\pgfsetmiterjoin%
\definecolor{currentfill}{rgb}{0.000000,0.000000,0.000000}%
\pgfsetfillcolor{currentfill}%
\pgfsetlinewidth{0.000000pt}%
\definecolor{currentstroke}{rgb}{0.000000,0.000000,0.000000}%
\pgfsetstrokecolor{currentstroke}%
\pgfsetstrokeopacity{0.000000}%
\pgfsetdash{}{0pt}%
\pgfpathmoveto{\pgfqpoint{6.826900in}{0.660000in}}%
\pgfpathlineto{\pgfqpoint{6.833100in}{0.660000in}}%
\pgfpathlineto{\pgfqpoint{6.833100in}{3.218826in}}%
\pgfpathlineto{\pgfqpoint{6.826900in}{3.218826in}}%
\pgfpathlineto{\pgfqpoint{6.826900in}{0.660000in}}%
\pgfpathclose%
\pgfusepath{fill}%
\end{pgfscope}%
\begin{pgfscope}%
\pgfpathrectangle{\pgfqpoint{1.250000in}{0.660000in}}{\pgfqpoint{7.750000in}{4.620000in}}%
\pgfusepath{clip}%
\pgfsetbuttcap%
\pgfsetmiterjoin%
\definecolor{currentfill}{rgb}{0.000000,0.000000,0.000000}%
\pgfsetfillcolor{currentfill}%
\pgfsetlinewidth{0.000000pt}%
\definecolor{currentstroke}{rgb}{0.000000,0.000000,0.000000}%
\pgfsetstrokecolor{currentstroke}%
\pgfsetstrokeopacity{0.000000}%
\pgfsetdash{}{0pt}%
\pgfpathmoveto{\pgfqpoint{6.834650in}{0.660000in}}%
\pgfpathlineto{\pgfqpoint{6.840850in}{0.660000in}}%
\pgfpathlineto{\pgfqpoint{6.840850in}{3.218690in}}%
\pgfpathlineto{\pgfqpoint{6.834650in}{3.218690in}}%
\pgfpathlineto{\pgfqpoint{6.834650in}{0.660000in}}%
\pgfpathclose%
\pgfusepath{fill}%
\end{pgfscope}%
\begin{pgfscope}%
\pgfpathrectangle{\pgfqpoint{1.250000in}{0.660000in}}{\pgfqpoint{7.750000in}{4.620000in}}%
\pgfusepath{clip}%
\pgfsetbuttcap%
\pgfsetmiterjoin%
\definecolor{currentfill}{rgb}{0.000000,0.000000,0.000000}%
\pgfsetfillcolor{currentfill}%
\pgfsetlinewidth{0.000000pt}%
\definecolor{currentstroke}{rgb}{0.000000,0.000000,0.000000}%
\pgfsetstrokecolor{currentstroke}%
\pgfsetstrokeopacity{0.000000}%
\pgfsetdash{}{0pt}%
\pgfpathmoveto{\pgfqpoint{6.842400in}{0.660000in}}%
\pgfpathlineto{\pgfqpoint{6.848600in}{0.660000in}}%
\pgfpathlineto{\pgfqpoint{6.848600in}{3.218551in}}%
\pgfpathlineto{\pgfqpoint{6.842400in}{3.218551in}}%
\pgfpathlineto{\pgfqpoint{6.842400in}{0.660000in}}%
\pgfpathclose%
\pgfusepath{fill}%
\end{pgfscope}%
\begin{pgfscope}%
\pgfpathrectangle{\pgfqpoint{1.250000in}{0.660000in}}{\pgfqpoint{7.750000in}{4.620000in}}%
\pgfusepath{clip}%
\pgfsetbuttcap%
\pgfsetmiterjoin%
\definecolor{currentfill}{rgb}{0.000000,0.000000,0.000000}%
\pgfsetfillcolor{currentfill}%
\pgfsetlinewidth{0.000000pt}%
\definecolor{currentstroke}{rgb}{0.000000,0.000000,0.000000}%
\pgfsetstrokecolor{currentstroke}%
\pgfsetstrokeopacity{0.000000}%
\pgfsetdash{}{0pt}%
\pgfpathmoveto{\pgfqpoint{6.850150in}{0.660000in}}%
\pgfpathlineto{\pgfqpoint{6.856350in}{0.660000in}}%
\pgfpathlineto{\pgfqpoint{6.856350in}{3.218408in}}%
\pgfpathlineto{\pgfqpoint{6.850150in}{3.218408in}}%
\pgfpathlineto{\pgfqpoint{6.850150in}{0.660000in}}%
\pgfpathclose%
\pgfusepath{fill}%
\end{pgfscope}%
\begin{pgfscope}%
\pgfpathrectangle{\pgfqpoint{1.250000in}{0.660000in}}{\pgfqpoint{7.750000in}{4.620000in}}%
\pgfusepath{clip}%
\pgfsetbuttcap%
\pgfsetmiterjoin%
\definecolor{currentfill}{rgb}{0.000000,0.000000,0.000000}%
\pgfsetfillcolor{currentfill}%
\pgfsetlinewidth{0.000000pt}%
\definecolor{currentstroke}{rgb}{0.000000,0.000000,0.000000}%
\pgfsetstrokecolor{currentstroke}%
\pgfsetstrokeopacity{0.000000}%
\pgfsetdash{}{0pt}%
\pgfpathmoveto{\pgfqpoint{6.857900in}{0.660000in}}%
\pgfpathlineto{\pgfqpoint{6.864100in}{0.660000in}}%
\pgfpathlineto{\pgfqpoint{6.864100in}{3.221742in}}%
\pgfpathlineto{\pgfqpoint{6.857900in}{3.221742in}}%
\pgfpathlineto{\pgfqpoint{6.857900in}{0.660000in}}%
\pgfpathclose%
\pgfusepath{fill}%
\end{pgfscope}%
\begin{pgfscope}%
\pgfpathrectangle{\pgfqpoint{1.250000in}{0.660000in}}{\pgfqpoint{7.750000in}{4.620000in}}%
\pgfusepath{clip}%
\pgfsetbuttcap%
\pgfsetmiterjoin%
\definecolor{currentfill}{rgb}{0.000000,0.000000,0.000000}%
\pgfsetfillcolor{currentfill}%
\pgfsetlinewidth{0.000000pt}%
\definecolor{currentstroke}{rgb}{0.000000,0.000000,0.000000}%
\pgfsetstrokecolor{currentstroke}%
\pgfsetstrokeopacity{0.000000}%
\pgfsetdash{}{0pt}%
\pgfpathmoveto{\pgfqpoint{6.865650in}{0.660000in}}%
\pgfpathlineto{\pgfqpoint{6.871850in}{0.660000in}}%
\pgfpathlineto{\pgfqpoint{6.871850in}{3.224672in}}%
\pgfpathlineto{\pgfqpoint{6.865650in}{3.224672in}}%
\pgfpathlineto{\pgfqpoint{6.865650in}{0.660000in}}%
\pgfpathclose%
\pgfusepath{fill}%
\end{pgfscope}%
\begin{pgfscope}%
\pgfpathrectangle{\pgfqpoint{1.250000in}{0.660000in}}{\pgfqpoint{7.750000in}{4.620000in}}%
\pgfusepath{clip}%
\pgfsetbuttcap%
\pgfsetmiterjoin%
\definecolor{currentfill}{rgb}{0.000000,0.000000,0.000000}%
\pgfsetfillcolor{currentfill}%
\pgfsetlinewidth{0.000000pt}%
\definecolor{currentstroke}{rgb}{0.000000,0.000000,0.000000}%
\pgfsetstrokecolor{currentstroke}%
\pgfsetstrokeopacity{0.000000}%
\pgfsetdash{}{0pt}%
\pgfpathmoveto{\pgfqpoint{6.873400in}{0.660000in}}%
\pgfpathlineto{\pgfqpoint{6.879600in}{0.660000in}}%
\pgfpathlineto{\pgfqpoint{6.879600in}{3.217960in}}%
\pgfpathlineto{\pgfqpoint{6.873400in}{3.217960in}}%
\pgfpathlineto{\pgfqpoint{6.873400in}{0.660000in}}%
\pgfpathclose%
\pgfusepath{fill}%
\end{pgfscope}%
\begin{pgfscope}%
\pgfpathrectangle{\pgfqpoint{1.250000in}{0.660000in}}{\pgfqpoint{7.750000in}{4.620000in}}%
\pgfusepath{clip}%
\pgfsetbuttcap%
\pgfsetmiterjoin%
\definecolor{currentfill}{rgb}{0.000000,0.000000,0.000000}%
\pgfsetfillcolor{currentfill}%
\pgfsetlinewidth{0.000000pt}%
\definecolor{currentstroke}{rgb}{0.000000,0.000000,0.000000}%
\pgfsetstrokecolor{currentstroke}%
\pgfsetstrokeopacity{0.000000}%
\pgfsetdash{}{0pt}%
\pgfpathmoveto{\pgfqpoint{6.881150in}{0.660000in}}%
\pgfpathlineto{\pgfqpoint{6.887350in}{0.660000in}}%
\pgfpathlineto{\pgfqpoint{6.887350in}{3.217803in}}%
\pgfpathlineto{\pgfqpoint{6.881150in}{3.217803in}}%
\pgfpathlineto{\pgfqpoint{6.881150in}{0.660000in}}%
\pgfpathclose%
\pgfusepath{fill}%
\end{pgfscope}%
\begin{pgfscope}%
\pgfpathrectangle{\pgfqpoint{1.250000in}{0.660000in}}{\pgfqpoint{7.750000in}{4.620000in}}%
\pgfusepath{clip}%
\pgfsetbuttcap%
\pgfsetmiterjoin%
\definecolor{currentfill}{rgb}{0.000000,0.000000,0.000000}%
\pgfsetfillcolor{currentfill}%
\pgfsetlinewidth{0.000000pt}%
\definecolor{currentstroke}{rgb}{0.000000,0.000000,0.000000}%
\pgfsetstrokecolor{currentstroke}%
\pgfsetstrokeopacity{0.000000}%
\pgfsetdash{}{0pt}%
\pgfpathmoveto{\pgfqpoint{6.888900in}{0.660000in}}%
\pgfpathlineto{\pgfqpoint{6.895100in}{0.660000in}}%
\pgfpathlineto{\pgfqpoint{6.895100in}{3.217642in}}%
\pgfpathlineto{\pgfqpoint{6.888900in}{3.217642in}}%
\pgfpathlineto{\pgfqpoint{6.888900in}{0.660000in}}%
\pgfpathclose%
\pgfusepath{fill}%
\end{pgfscope}%
\begin{pgfscope}%
\pgfpathrectangle{\pgfqpoint{1.250000in}{0.660000in}}{\pgfqpoint{7.750000in}{4.620000in}}%
\pgfusepath{clip}%
\pgfsetbuttcap%
\pgfsetmiterjoin%
\definecolor{currentfill}{rgb}{0.000000,0.000000,0.000000}%
\pgfsetfillcolor{currentfill}%
\pgfsetlinewidth{0.000000pt}%
\definecolor{currentstroke}{rgb}{0.000000,0.000000,0.000000}%
\pgfsetstrokecolor{currentstroke}%
\pgfsetstrokeopacity{0.000000}%
\pgfsetdash{}{0pt}%
\pgfpathmoveto{\pgfqpoint{6.896650in}{0.660000in}}%
\pgfpathlineto{\pgfqpoint{6.902850in}{0.660000in}}%
\pgfpathlineto{\pgfqpoint{6.902850in}{3.217477in}}%
\pgfpathlineto{\pgfqpoint{6.896650in}{3.217477in}}%
\pgfpathlineto{\pgfqpoint{6.896650in}{0.660000in}}%
\pgfpathclose%
\pgfusepath{fill}%
\end{pgfscope}%
\begin{pgfscope}%
\pgfpathrectangle{\pgfqpoint{1.250000in}{0.660000in}}{\pgfqpoint{7.750000in}{4.620000in}}%
\pgfusepath{clip}%
\pgfsetbuttcap%
\pgfsetmiterjoin%
\definecolor{currentfill}{rgb}{0.000000,0.000000,0.000000}%
\pgfsetfillcolor{currentfill}%
\pgfsetlinewidth{0.000000pt}%
\definecolor{currentstroke}{rgb}{0.000000,0.000000,0.000000}%
\pgfsetstrokecolor{currentstroke}%
\pgfsetstrokeopacity{0.000000}%
\pgfsetdash{}{0pt}%
\pgfpathmoveto{\pgfqpoint{6.904400in}{0.660000in}}%
\pgfpathlineto{\pgfqpoint{6.910600in}{0.660000in}}%
\pgfpathlineto{\pgfqpoint{6.910600in}{3.219357in}}%
\pgfpathlineto{\pgfqpoint{6.904400in}{3.219357in}}%
\pgfpathlineto{\pgfqpoint{6.904400in}{0.660000in}}%
\pgfpathclose%
\pgfusepath{fill}%
\end{pgfscope}%
\begin{pgfscope}%
\pgfpathrectangle{\pgfqpoint{1.250000in}{0.660000in}}{\pgfqpoint{7.750000in}{4.620000in}}%
\pgfusepath{clip}%
\pgfsetbuttcap%
\pgfsetmiterjoin%
\definecolor{currentfill}{rgb}{0.000000,0.000000,0.000000}%
\pgfsetfillcolor{currentfill}%
\pgfsetlinewidth{0.000000pt}%
\definecolor{currentstroke}{rgb}{0.000000,0.000000,0.000000}%
\pgfsetstrokecolor{currentstroke}%
\pgfsetstrokeopacity{0.000000}%
\pgfsetdash{}{0pt}%
\pgfpathmoveto{\pgfqpoint{6.912150in}{0.660000in}}%
\pgfpathlineto{\pgfqpoint{6.918350in}{0.660000in}}%
\pgfpathlineto{\pgfqpoint{6.918350in}{3.217135in}}%
\pgfpathlineto{\pgfqpoint{6.912150in}{3.217135in}}%
\pgfpathlineto{\pgfqpoint{6.912150in}{0.660000in}}%
\pgfpathclose%
\pgfusepath{fill}%
\end{pgfscope}%
\begin{pgfscope}%
\pgfpathrectangle{\pgfqpoint{1.250000in}{0.660000in}}{\pgfqpoint{7.750000in}{4.620000in}}%
\pgfusepath{clip}%
\pgfsetbuttcap%
\pgfsetmiterjoin%
\definecolor{currentfill}{rgb}{0.000000,0.000000,0.000000}%
\pgfsetfillcolor{currentfill}%
\pgfsetlinewidth{0.000000pt}%
\definecolor{currentstroke}{rgb}{0.000000,0.000000,0.000000}%
\pgfsetstrokecolor{currentstroke}%
\pgfsetstrokeopacity{0.000000}%
\pgfsetdash{}{0pt}%
\pgfpathmoveto{\pgfqpoint{6.919900in}{0.660000in}}%
\pgfpathlineto{\pgfqpoint{6.926100in}{0.660000in}}%
\pgfpathlineto{\pgfqpoint{6.926100in}{3.216957in}}%
\pgfpathlineto{\pgfqpoint{6.919900in}{3.216957in}}%
\pgfpathlineto{\pgfqpoint{6.919900in}{0.660000in}}%
\pgfpathclose%
\pgfusepath{fill}%
\end{pgfscope}%
\begin{pgfscope}%
\pgfpathrectangle{\pgfqpoint{1.250000in}{0.660000in}}{\pgfqpoint{7.750000in}{4.620000in}}%
\pgfusepath{clip}%
\pgfsetbuttcap%
\pgfsetmiterjoin%
\definecolor{currentfill}{rgb}{0.000000,0.000000,0.000000}%
\pgfsetfillcolor{currentfill}%
\pgfsetlinewidth{0.000000pt}%
\definecolor{currentstroke}{rgb}{0.000000,0.000000,0.000000}%
\pgfsetstrokecolor{currentstroke}%
\pgfsetstrokeopacity{0.000000}%
\pgfsetdash{}{0pt}%
\pgfpathmoveto{\pgfqpoint{6.927650in}{0.660000in}}%
\pgfpathlineto{\pgfqpoint{6.933850in}{0.660000in}}%
\pgfpathlineto{\pgfqpoint{6.933850in}{3.216774in}}%
\pgfpathlineto{\pgfqpoint{6.927650in}{3.216774in}}%
\pgfpathlineto{\pgfqpoint{6.927650in}{0.660000in}}%
\pgfpathclose%
\pgfusepath{fill}%
\end{pgfscope}%
\begin{pgfscope}%
\pgfpathrectangle{\pgfqpoint{1.250000in}{0.660000in}}{\pgfqpoint{7.750000in}{4.620000in}}%
\pgfusepath{clip}%
\pgfsetbuttcap%
\pgfsetmiterjoin%
\definecolor{currentfill}{rgb}{0.000000,0.000000,0.000000}%
\pgfsetfillcolor{currentfill}%
\pgfsetlinewidth{0.000000pt}%
\definecolor{currentstroke}{rgb}{0.000000,0.000000,0.000000}%
\pgfsetstrokecolor{currentstroke}%
\pgfsetstrokeopacity{0.000000}%
\pgfsetdash{}{0pt}%
\pgfpathmoveto{\pgfqpoint{6.935400in}{0.660000in}}%
\pgfpathlineto{\pgfqpoint{6.941600in}{0.660000in}}%
\pgfpathlineto{\pgfqpoint{6.941600in}{3.216660in}}%
\pgfpathlineto{\pgfqpoint{6.935400in}{3.216660in}}%
\pgfpathlineto{\pgfqpoint{6.935400in}{0.660000in}}%
\pgfpathclose%
\pgfusepath{fill}%
\end{pgfscope}%
\begin{pgfscope}%
\pgfpathrectangle{\pgfqpoint{1.250000in}{0.660000in}}{\pgfqpoint{7.750000in}{4.620000in}}%
\pgfusepath{clip}%
\pgfsetbuttcap%
\pgfsetmiterjoin%
\definecolor{currentfill}{rgb}{0.000000,0.000000,0.000000}%
\pgfsetfillcolor{currentfill}%
\pgfsetlinewidth{0.000000pt}%
\definecolor{currentstroke}{rgb}{0.000000,0.000000,0.000000}%
\pgfsetstrokecolor{currentstroke}%
\pgfsetstrokeopacity{0.000000}%
\pgfsetdash{}{0pt}%
\pgfpathmoveto{\pgfqpoint{6.943150in}{0.660000in}}%
\pgfpathlineto{\pgfqpoint{6.949350in}{0.660000in}}%
\pgfpathlineto{\pgfqpoint{6.949350in}{3.216603in}}%
\pgfpathlineto{\pgfqpoint{6.943150in}{3.216603in}}%
\pgfpathlineto{\pgfqpoint{6.943150in}{0.660000in}}%
\pgfpathclose%
\pgfusepath{fill}%
\end{pgfscope}%
\begin{pgfscope}%
\pgfpathrectangle{\pgfqpoint{1.250000in}{0.660000in}}{\pgfqpoint{7.750000in}{4.620000in}}%
\pgfusepath{clip}%
\pgfsetbuttcap%
\pgfsetmiterjoin%
\definecolor{currentfill}{rgb}{0.000000,0.000000,0.000000}%
\pgfsetfillcolor{currentfill}%
\pgfsetlinewidth{0.000000pt}%
\definecolor{currentstroke}{rgb}{0.000000,0.000000,0.000000}%
\pgfsetstrokecolor{currentstroke}%
\pgfsetstrokeopacity{0.000000}%
\pgfsetdash{}{0pt}%
\pgfpathmoveto{\pgfqpoint{6.950900in}{0.660000in}}%
\pgfpathlineto{\pgfqpoint{6.957100in}{0.660000in}}%
\pgfpathlineto{\pgfqpoint{6.957100in}{3.216542in}}%
\pgfpathlineto{\pgfqpoint{6.950900in}{3.216542in}}%
\pgfpathlineto{\pgfqpoint{6.950900in}{0.660000in}}%
\pgfpathclose%
\pgfusepath{fill}%
\end{pgfscope}%
\begin{pgfscope}%
\pgfpathrectangle{\pgfqpoint{1.250000in}{0.660000in}}{\pgfqpoint{7.750000in}{4.620000in}}%
\pgfusepath{clip}%
\pgfsetbuttcap%
\pgfsetmiterjoin%
\definecolor{currentfill}{rgb}{0.000000,0.000000,0.000000}%
\pgfsetfillcolor{currentfill}%
\pgfsetlinewidth{0.000000pt}%
\definecolor{currentstroke}{rgb}{0.000000,0.000000,0.000000}%
\pgfsetstrokecolor{currentstroke}%
\pgfsetstrokeopacity{0.000000}%
\pgfsetdash{}{0pt}%
\pgfpathmoveto{\pgfqpoint{6.958650in}{0.660000in}}%
\pgfpathlineto{\pgfqpoint{6.964850in}{0.660000in}}%
\pgfpathlineto{\pgfqpoint{6.964850in}{3.216478in}}%
\pgfpathlineto{\pgfqpoint{6.958650in}{3.216478in}}%
\pgfpathlineto{\pgfqpoint{6.958650in}{0.660000in}}%
\pgfpathclose%
\pgfusepath{fill}%
\end{pgfscope}%
\begin{pgfscope}%
\pgfpathrectangle{\pgfqpoint{1.250000in}{0.660000in}}{\pgfqpoint{7.750000in}{4.620000in}}%
\pgfusepath{clip}%
\pgfsetbuttcap%
\pgfsetmiterjoin%
\definecolor{currentfill}{rgb}{0.000000,0.000000,0.000000}%
\pgfsetfillcolor{currentfill}%
\pgfsetlinewidth{0.000000pt}%
\definecolor{currentstroke}{rgb}{0.000000,0.000000,0.000000}%
\pgfsetstrokecolor{currentstroke}%
\pgfsetstrokeopacity{0.000000}%
\pgfsetdash{}{0pt}%
\pgfpathmoveto{\pgfqpoint{6.966400in}{0.660000in}}%
\pgfpathlineto{\pgfqpoint{6.972600in}{0.660000in}}%
\pgfpathlineto{\pgfqpoint{6.972600in}{3.216410in}}%
\pgfpathlineto{\pgfqpoint{6.966400in}{3.216410in}}%
\pgfpathlineto{\pgfqpoint{6.966400in}{0.660000in}}%
\pgfpathclose%
\pgfusepath{fill}%
\end{pgfscope}%
\begin{pgfscope}%
\pgfpathrectangle{\pgfqpoint{1.250000in}{0.660000in}}{\pgfqpoint{7.750000in}{4.620000in}}%
\pgfusepath{clip}%
\pgfsetbuttcap%
\pgfsetmiterjoin%
\definecolor{currentfill}{rgb}{0.000000,0.000000,0.000000}%
\pgfsetfillcolor{currentfill}%
\pgfsetlinewidth{0.000000pt}%
\definecolor{currentstroke}{rgb}{0.000000,0.000000,0.000000}%
\pgfsetstrokecolor{currentstroke}%
\pgfsetstrokeopacity{0.000000}%
\pgfsetdash{}{0pt}%
\pgfpathmoveto{\pgfqpoint{6.974150in}{0.660000in}}%
\pgfpathlineto{\pgfqpoint{6.980350in}{0.660000in}}%
\pgfpathlineto{\pgfqpoint{6.980350in}{3.216338in}}%
\pgfpathlineto{\pgfqpoint{6.974150in}{3.216338in}}%
\pgfpathlineto{\pgfqpoint{6.974150in}{0.660000in}}%
\pgfpathclose%
\pgfusepath{fill}%
\end{pgfscope}%
\begin{pgfscope}%
\pgfpathrectangle{\pgfqpoint{1.250000in}{0.660000in}}{\pgfqpoint{7.750000in}{4.620000in}}%
\pgfusepath{clip}%
\pgfsetbuttcap%
\pgfsetmiterjoin%
\definecolor{currentfill}{rgb}{0.000000,0.000000,0.000000}%
\pgfsetfillcolor{currentfill}%
\pgfsetlinewidth{0.000000pt}%
\definecolor{currentstroke}{rgb}{0.000000,0.000000,0.000000}%
\pgfsetstrokecolor{currentstroke}%
\pgfsetstrokeopacity{0.000000}%
\pgfsetdash{}{0pt}%
\pgfpathmoveto{\pgfqpoint{6.981900in}{0.660000in}}%
\pgfpathlineto{\pgfqpoint{6.988100in}{0.660000in}}%
\pgfpathlineto{\pgfqpoint{6.988100in}{3.216262in}}%
\pgfpathlineto{\pgfqpoint{6.981900in}{3.216262in}}%
\pgfpathlineto{\pgfqpoint{6.981900in}{0.660000in}}%
\pgfpathclose%
\pgfusepath{fill}%
\end{pgfscope}%
\begin{pgfscope}%
\pgfpathrectangle{\pgfqpoint{1.250000in}{0.660000in}}{\pgfqpoint{7.750000in}{4.620000in}}%
\pgfusepath{clip}%
\pgfsetbuttcap%
\pgfsetmiterjoin%
\definecolor{currentfill}{rgb}{0.000000,0.000000,0.000000}%
\pgfsetfillcolor{currentfill}%
\pgfsetlinewidth{0.000000pt}%
\definecolor{currentstroke}{rgb}{0.000000,0.000000,0.000000}%
\pgfsetstrokecolor{currentstroke}%
\pgfsetstrokeopacity{0.000000}%
\pgfsetdash{}{0pt}%
\pgfpathmoveto{\pgfqpoint{6.989650in}{0.660000in}}%
\pgfpathlineto{\pgfqpoint{6.995850in}{0.660000in}}%
\pgfpathlineto{\pgfqpoint{6.995850in}{3.216183in}}%
\pgfpathlineto{\pgfqpoint{6.989650in}{3.216183in}}%
\pgfpathlineto{\pgfqpoint{6.989650in}{0.660000in}}%
\pgfpathclose%
\pgfusepath{fill}%
\end{pgfscope}%
\begin{pgfscope}%
\pgfpathrectangle{\pgfqpoint{1.250000in}{0.660000in}}{\pgfqpoint{7.750000in}{4.620000in}}%
\pgfusepath{clip}%
\pgfsetbuttcap%
\pgfsetmiterjoin%
\definecolor{currentfill}{rgb}{0.000000,0.000000,0.000000}%
\pgfsetfillcolor{currentfill}%
\pgfsetlinewidth{0.000000pt}%
\definecolor{currentstroke}{rgb}{0.000000,0.000000,0.000000}%
\pgfsetstrokecolor{currentstroke}%
\pgfsetstrokeopacity{0.000000}%
\pgfsetdash{}{0pt}%
\pgfpathmoveto{\pgfqpoint{6.997400in}{0.660000in}}%
\pgfpathlineto{\pgfqpoint{7.003600in}{0.660000in}}%
\pgfpathlineto{\pgfqpoint{7.003600in}{3.216099in}}%
\pgfpathlineto{\pgfqpoint{6.997400in}{3.216099in}}%
\pgfpathlineto{\pgfqpoint{6.997400in}{0.660000in}}%
\pgfpathclose%
\pgfusepath{fill}%
\end{pgfscope}%
\begin{pgfscope}%
\pgfpathrectangle{\pgfqpoint{1.250000in}{0.660000in}}{\pgfqpoint{7.750000in}{4.620000in}}%
\pgfusepath{clip}%
\pgfsetbuttcap%
\pgfsetmiterjoin%
\definecolor{currentfill}{rgb}{0.000000,0.000000,0.000000}%
\pgfsetfillcolor{currentfill}%
\pgfsetlinewidth{0.000000pt}%
\definecolor{currentstroke}{rgb}{0.000000,0.000000,0.000000}%
\pgfsetstrokecolor{currentstroke}%
\pgfsetstrokeopacity{0.000000}%
\pgfsetdash{}{0pt}%
\pgfpathmoveto{\pgfqpoint{7.005150in}{0.660000in}}%
\pgfpathlineto{\pgfqpoint{7.011350in}{0.660000in}}%
\pgfpathlineto{\pgfqpoint{7.011350in}{3.225540in}}%
\pgfpathlineto{\pgfqpoint{7.005150in}{3.225540in}}%
\pgfpathlineto{\pgfqpoint{7.005150in}{0.660000in}}%
\pgfpathclose%
\pgfusepath{fill}%
\end{pgfscope}%
\begin{pgfscope}%
\pgfpathrectangle{\pgfqpoint{1.250000in}{0.660000in}}{\pgfqpoint{7.750000in}{4.620000in}}%
\pgfusepath{clip}%
\pgfsetbuttcap%
\pgfsetmiterjoin%
\definecolor{currentfill}{rgb}{0.000000,0.000000,0.000000}%
\pgfsetfillcolor{currentfill}%
\pgfsetlinewidth{0.000000pt}%
\definecolor{currentstroke}{rgb}{0.000000,0.000000,0.000000}%
\pgfsetstrokecolor{currentstroke}%
\pgfsetstrokeopacity{0.000000}%
\pgfsetdash{}{0pt}%
\pgfpathmoveto{\pgfqpoint{7.012900in}{0.660000in}}%
\pgfpathlineto{\pgfqpoint{7.019100in}{0.660000in}}%
\pgfpathlineto{\pgfqpoint{7.019100in}{3.249664in}}%
\pgfpathlineto{\pgfqpoint{7.012900in}{3.249664in}}%
\pgfpathlineto{\pgfqpoint{7.012900in}{0.660000in}}%
\pgfpathclose%
\pgfusepath{fill}%
\end{pgfscope}%
\begin{pgfscope}%
\pgfpathrectangle{\pgfqpoint{1.250000in}{0.660000in}}{\pgfqpoint{7.750000in}{4.620000in}}%
\pgfusepath{clip}%
\pgfsetbuttcap%
\pgfsetmiterjoin%
\definecolor{currentfill}{rgb}{0.000000,0.000000,0.000000}%
\pgfsetfillcolor{currentfill}%
\pgfsetlinewidth{0.000000pt}%
\definecolor{currentstroke}{rgb}{0.000000,0.000000,0.000000}%
\pgfsetstrokecolor{currentstroke}%
\pgfsetstrokeopacity{0.000000}%
\pgfsetdash{}{0pt}%
\pgfpathmoveto{\pgfqpoint{7.020650in}{0.660000in}}%
\pgfpathlineto{\pgfqpoint{7.026850in}{0.660000in}}%
\pgfpathlineto{\pgfqpoint{7.026850in}{3.215820in}}%
\pgfpathlineto{\pgfqpoint{7.020650in}{3.215820in}}%
\pgfpathlineto{\pgfqpoint{7.020650in}{0.660000in}}%
\pgfpathclose%
\pgfusepath{fill}%
\end{pgfscope}%
\begin{pgfscope}%
\pgfpathrectangle{\pgfqpoint{1.250000in}{0.660000in}}{\pgfqpoint{7.750000in}{4.620000in}}%
\pgfusepath{clip}%
\pgfsetbuttcap%
\pgfsetmiterjoin%
\definecolor{currentfill}{rgb}{0.000000,0.000000,0.000000}%
\pgfsetfillcolor{currentfill}%
\pgfsetlinewidth{0.000000pt}%
\definecolor{currentstroke}{rgb}{0.000000,0.000000,0.000000}%
\pgfsetstrokecolor{currentstroke}%
\pgfsetstrokeopacity{0.000000}%
\pgfsetdash{}{0pt}%
\pgfpathmoveto{\pgfqpoint{7.028400in}{0.660000in}}%
\pgfpathlineto{\pgfqpoint{7.034600in}{0.660000in}}%
\pgfpathlineto{\pgfqpoint{7.034600in}{3.215717in}}%
\pgfpathlineto{\pgfqpoint{7.028400in}{3.215717in}}%
\pgfpathlineto{\pgfqpoint{7.028400in}{0.660000in}}%
\pgfpathclose%
\pgfusepath{fill}%
\end{pgfscope}%
\begin{pgfscope}%
\pgfpathrectangle{\pgfqpoint{1.250000in}{0.660000in}}{\pgfqpoint{7.750000in}{4.620000in}}%
\pgfusepath{clip}%
\pgfsetbuttcap%
\pgfsetmiterjoin%
\definecolor{currentfill}{rgb}{0.000000,0.000000,0.000000}%
\pgfsetfillcolor{currentfill}%
\pgfsetlinewidth{0.000000pt}%
\definecolor{currentstroke}{rgb}{0.000000,0.000000,0.000000}%
\pgfsetstrokecolor{currentstroke}%
\pgfsetstrokeopacity{0.000000}%
\pgfsetdash{}{0pt}%
\pgfpathmoveto{\pgfqpoint{7.036150in}{0.660000in}}%
\pgfpathlineto{\pgfqpoint{7.042350in}{0.660000in}}%
\pgfpathlineto{\pgfqpoint{7.042350in}{3.215610in}}%
\pgfpathlineto{\pgfqpoint{7.036150in}{3.215610in}}%
\pgfpathlineto{\pgfqpoint{7.036150in}{0.660000in}}%
\pgfpathclose%
\pgfusepath{fill}%
\end{pgfscope}%
\begin{pgfscope}%
\pgfpathrectangle{\pgfqpoint{1.250000in}{0.660000in}}{\pgfqpoint{7.750000in}{4.620000in}}%
\pgfusepath{clip}%
\pgfsetbuttcap%
\pgfsetmiterjoin%
\definecolor{currentfill}{rgb}{0.000000,0.000000,0.000000}%
\pgfsetfillcolor{currentfill}%
\pgfsetlinewidth{0.000000pt}%
\definecolor{currentstroke}{rgb}{0.000000,0.000000,0.000000}%
\pgfsetstrokecolor{currentstroke}%
\pgfsetstrokeopacity{0.000000}%
\pgfsetdash{}{0pt}%
\pgfpathmoveto{\pgfqpoint{7.043900in}{0.660000in}}%
\pgfpathlineto{\pgfqpoint{7.050100in}{0.660000in}}%
\pgfpathlineto{\pgfqpoint{7.050100in}{3.218474in}}%
\pgfpathlineto{\pgfqpoint{7.043900in}{3.218474in}}%
\pgfpathlineto{\pgfqpoint{7.043900in}{0.660000in}}%
\pgfpathclose%
\pgfusepath{fill}%
\end{pgfscope}%
\begin{pgfscope}%
\pgfpathrectangle{\pgfqpoint{1.250000in}{0.660000in}}{\pgfqpoint{7.750000in}{4.620000in}}%
\pgfusepath{clip}%
\pgfsetbuttcap%
\pgfsetmiterjoin%
\definecolor{currentfill}{rgb}{0.000000,0.000000,0.000000}%
\pgfsetfillcolor{currentfill}%
\pgfsetlinewidth{0.000000pt}%
\definecolor{currentstroke}{rgb}{0.000000,0.000000,0.000000}%
\pgfsetstrokecolor{currentstroke}%
\pgfsetstrokeopacity{0.000000}%
\pgfsetdash{}{0pt}%
\pgfpathmoveto{\pgfqpoint{7.051650in}{0.660000in}}%
\pgfpathlineto{\pgfqpoint{7.057850in}{0.660000in}}%
\pgfpathlineto{\pgfqpoint{7.057850in}{3.220184in}}%
\pgfpathlineto{\pgfqpoint{7.051650in}{3.220184in}}%
\pgfpathlineto{\pgfqpoint{7.051650in}{0.660000in}}%
\pgfpathclose%
\pgfusepath{fill}%
\end{pgfscope}%
\begin{pgfscope}%
\pgfpathrectangle{\pgfqpoint{1.250000in}{0.660000in}}{\pgfqpoint{7.750000in}{4.620000in}}%
\pgfusepath{clip}%
\pgfsetbuttcap%
\pgfsetmiterjoin%
\definecolor{currentfill}{rgb}{0.000000,0.000000,0.000000}%
\pgfsetfillcolor{currentfill}%
\pgfsetlinewidth{0.000000pt}%
\definecolor{currentstroke}{rgb}{0.000000,0.000000,0.000000}%
\pgfsetstrokecolor{currentstroke}%
\pgfsetstrokeopacity{0.000000}%
\pgfsetdash{}{0pt}%
\pgfpathmoveto{\pgfqpoint{7.059400in}{0.660000in}}%
\pgfpathlineto{\pgfqpoint{7.065600in}{0.660000in}}%
\pgfpathlineto{\pgfqpoint{7.065600in}{3.215257in}}%
\pgfpathlineto{\pgfqpoint{7.059400in}{3.215257in}}%
\pgfpathlineto{\pgfqpoint{7.059400in}{0.660000in}}%
\pgfpathclose%
\pgfusepath{fill}%
\end{pgfscope}%
\begin{pgfscope}%
\pgfpathrectangle{\pgfqpoint{1.250000in}{0.660000in}}{\pgfqpoint{7.750000in}{4.620000in}}%
\pgfusepath{clip}%
\pgfsetbuttcap%
\pgfsetmiterjoin%
\definecolor{currentfill}{rgb}{0.000000,0.000000,0.000000}%
\pgfsetfillcolor{currentfill}%
\pgfsetlinewidth{0.000000pt}%
\definecolor{currentstroke}{rgb}{0.000000,0.000000,0.000000}%
\pgfsetstrokecolor{currentstroke}%
\pgfsetstrokeopacity{0.000000}%
\pgfsetdash{}{0pt}%
\pgfpathmoveto{\pgfqpoint{7.067150in}{0.660000in}}%
\pgfpathlineto{\pgfqpoint{7.073350in}{0.660000in}}%
\pgfpathlineto{\pgfqpoint{7.073350in}{3.215128in}}%
\pgfpathlineto{\pgfqpoint{7.067150in}{3.215128in}}%
\pgfpathlineto{\pgfqpoint{7.067150in}{0.660000in}}%
\pgfpathclose%
\pgfusepath{fill}%
\end{pgfscope}%
\begin{pgfscope}%
\pgfpathrectangle{\pgfqpoint{1.250000in}{0.660000in}}{\pgfqpoint{7.750000in}{4.620000in}}%
\pgfusepath{clip}%
\pgfsetbuttcap%
\pgfsetmiterjoin%
\definecolor{currentfill}{rgb}{0.000000,0.000000,0.000000}%
\pgfsetfillcolor{currentfill}%
\pgfsetlinewidth{0.000000pt}%
\definecolor{currentstroke}{rgb}{0.000000,0.000000,0.000000}%
\pgfsetstrokecolor{currentstroke}%
\pgfsetstrokeopacity{0.000000}%
\pgfsetdash{}{0pt}%
\pgfpathmoveto{\pgfqpoint{7.074900in}{0.660000in}}%
\pgfpathlineto{\pgfqpoint{7.081100in}{0.660000in}}%
\pgfpathlineto{\pgfqpoint{7.081100in}{3.214994in}}%
\pgfpathlineto{\pgfqpoint{7.074900in}{3.214994in}}%
\pgfpathlineto{\pgfqpoint{7.074900in}{0.660000in}}%
\pgfpathclose%
\pgfusepath{fill}%
\end{pgfscope}%
\begin{pgfscope}%
\pgfpathrectangle{\pgfqpoint{1.250000in}{0.660000in}}{\pgfqpoint{7.750000in}{4.620000in}}%
\pgfusepath{clip}%
\pgfsetbuttcap%
\pgfsetmiterjoin%
\definecolor{currentfill}{rgb}{0.000000,0.000000,0.000000}%
\pgfsetfillcolor{currentfill}%
\pgfsetlinewidth{0.000000pt}%
\definecolor{currentstroke}{rgb}{0.000000,0.000000,0.000000}%
\pgfsetstrokecolor{currentstroke}%
\pgfsetstrokeopacity{0.000000}%
\pgfsetdash{}{0pt}%
\pgfpathmoveto{\pgfqpoint{7.082650in}{0.660000in}}%
\pgfpathlineto{\pgfqpoint{7.088850in}{0.660000in}}%
\pgfpathlineto{\pgfqpoint{7.088850in}{3.214853in}}%
\pgfpathlineto{\pgfqpoint{7.082650in}{3.214853in}}%
\pgfpathlineto{\pgfqpoint{7.082650in}{0.660000in}}%
\pgfpathclose%
\pgfusepath{fill}%
\end{pgfscope}%
\begin{pgfscope}%
\pgfpathrectangle{\pgfqpoint{1.250000in}{0.660000in}}{\pgfqpoint{7.750000in}{4.620000in}}%
\pgfusepath{clip}%
\pgfsetbuttcap%
\pgfsetmiterjoin%
\definecolor{currentfill}{rgb}{0.000000,0.000000,0.000000}%
\pgfsetfillcolor{currentfill}%
\pgfsetlinewidth{0.000000pt}%
\definecolor{currentstroke}{rgb}{0.000000,0.000000,0.000000}%
\pgfsetstrokecolor{currentstroke}%
\pgfsetstrokeopacity{0.000000}%
\pgfsetdash{}{0pt}%
\pgfpathmoveto{\pgfqpoint{7.090400in}{0.660000in}}%
\pgfpathlineto{\pgfqpoint{7.096600in}{0.660000in}}%
\pgfpathlineto{\pgfqpoint{7.096600in}{3.230581in}}%
\pgfpathlineto{\pgfqpoint{7.090400in}{3.230581in}}%
\pgfpathlineto{\pgfqpoint{7.090400in}{0.660000in}}%
\pgfpathclose%
\pgfusepath{fill}%
\end{pgfscope}%
\begin{pgfscope}%
\pgfpathrectangle{\pgfqpoint{1.250000in}{0.660000in}}{\pgfqpoint{7.750000in}{4.620000in}}%
\pgfusepath{clip}%
\pgfsetbuttcap%
\pgfsetmiterjoin%
\definecolor{currentfill}{rgb}{0.000000,0.000000,0.000000}%
\pgfsetfillcolor{currentfill}%
\pgfsetlinewidth{0.000000pt}%
\definecolor{currentstroke}{rgb}{0.000000,0.000000,0.000000}%
\pgfsetstrokecolor{currentstroke}%
\pgfsetstrokeopacity{0.000000}%
\pgfsetdash{}{0pt}%
\pgfpathmoveto{\pgfqpoint{7.098150in}{0.660000in}}%
\pgfpathlineto{\pgfqpoint{7.104350in}{0.660000in}}%
\pgfpathlineto{\pgfqpoint{7.104350in}{3.214555in}}%
\pgfpathlineto{\pgfqpoint{7.098150in}{3.214555in}}%
\pgfpathlineto{\pgfqpoint{7.098150in}{0.660000in}}%
\pgfpathclose%
\pgfusepath{fill}%
\end{pgfscope}%
\begin{pgfscope}%
\pgfpathrectangle{\pgfqpoint{1.250000in}{0.660000in}}{\pgfqpoint{7.750000in}{4.620000in}}%
\pgfusepath{clip}%
\pgfsetbuttcap%
\pgfsetmiterjoin%
\definecolor{currentfill}{rgb}{0.000000,0.000000,0.000000}%
\pgfsetfillcolor{currentfill}%
\pgfsetlinewidth{0.000000pt}%
\definecolor{currentstroke}{rgb}{0.000000,0.000000,0.000000}%
\pgfsetstrokecolor{currentstroke}%
\pgfsetstrokeopacity{0.000000}%
\pgfsetdash{}{0pt}%
\pgfpathmoveto{\pgfqpoint{7.105900in}{0.660000in}}%
\pgfpathlineto{\pgfqpoint{7.112100in}{0.660000in}}%
\pgfpathlineto{\pgfqpoint{7.112100in}{3.214396in}}%
\pgfpathlineto{\pgfqpoint{7.105900in}{3.214396in}}%
\pgfpathlineto{\pgfqpoint{7.105900in}{0.660000in}}%
\pgfpathclose%
\pgfusepath{fill}%
\end{pgfscope}%
\begin{pgfscope}%
\pgfpathrectangle{\pgfqpoint{1.250000in}{0.660000in}}{\pgfqpoint{7.750000in}{4.620000in}}%
\pgfusepath{clip}%
\pgfsetbuttcap%
\pgfsetmiterjoin%
\definecolor{currentfill}{rgb}{0.000000,0.000000,0.000000}%
\pgfsetfillcolor{currentfill}%
\pgfsetlinewidth{0.000000pt}%
\definecolor{currentstroke}{rgb}{0.000000,0.000000,0.000000}%
\pgfsetstrokecolor{currentstroke}%
\pgfsetstrokeopacity{0.000000}%
\pgfsetdash{}{0pt}%
\pgfpathmoveto{\pgfqpoint{7.113650in}{0.660000in}}%
\pgfpathlineto{\pgfqpoint{7.119850in}{0.660000in}}%
\pgfpathlineto{\pgfqpoint{7.119850in}{3.214230in}}%
\pgfpathlineto{\pgfqpoint{7.113650in}{3.214230in}}%
\pgfpathlineto{\pgfqpoint{7.113650in}{0.660000in}}%
\pgfpathclose%
\pgfusepath{fill}%
\end{pgfscope}%
\begin{pgfscope}%
\pgfpathrectangle{\pgfqpoint{1.250000in}{0.660000in}}{\pgfqpoint{7.750000in}{4.620000in}}%
\pgfusepath{clip}%
\pgfsetbuttcap%
\pgfsetmiterjoin%
\definecolor{currentfill}{rgb}{0.000000,0.000000,0.000000}%
\pgfsetfillcolor{currentfill}%
\pgfsetlinewidth{0.000000pt}%
\definecolor{currentstroke}{rgb}{0.000000,0.000000,0.000000}%
\pgfsetstrokecolor{currentstroke}%
\pgfsetstrokeopacity{0.000000}%
\pgfsetdash{}{0pt}%
\pgfpathmoveto{\pgfqpoint{7.121400in}{0.660000in}}%
\pgfpathlineto{\pgfqpoint{7.127600in}{0.660000in}}%
\pgfpathlineto{\pgfqpoint{7.127600in}{3.214058in}}%
\pgfpathlineto{\pgfqpoint{7.121400in}{3.214058in}}%
\pgfpathlineto{\pgfqpoint{7.121400in}{0.660000in}}%
\pgfpathclose%
\pgfusepath{fill}%
\end{pgfscope}%
\begin{pgfscope}%
\pgfpathrectangle{\pgfqpoint{1.250000in}{0.660000in}}{\pgfqpoint{7.750000in}{4.620000in}}%
\pgfusepath{clip}%
\pgfsetbuttcap%
\pgfsetmiterjoin%
\definecolor{currentfill}{rgb}{0.000000,0.000000,0.000000}%
\pgfsetfillcolor{currentfill}%
\pgfsetlinewidth{0.000000pt}%
\definecolor{currentstroke}{rgb}{0.000000,0.000000,0.000000}%
\pgfsetstrokecolor{currentstroke}%
\pgfsetstrokeopacity{0.000000}%
\pgfsetdash{}{0pt}%
\pgfpathmoveto{\pgfqpoint{7.129150in}{0.660000in}}%
\pgfpathlineto{\pgfqpoint{7.135350in}{0.660000in}}%
\pgfpathlineto{\pgfqpoint{7.135350in}{3.213879in}}%
\pgfpathlineto{\pgfqpoint{7.129150in}{3.213879in}}%
\pgfpathlineto{\pgfqpoint{7.129150in}{0.660000in}}%
\pgfpathclose%
\pgfusepath{fill}%
\end{pgfscope}%
\begin{pgfscope}%
\pgfpathrectangle{\pgfqpoint{1.250000in}{0.660000in}}{\pgfqpoint{7.750000in}{4.620000in}}%
\pgfusepath{clip}%
\pgfsetbuttcap%
\pgfsetmiterjoin%
\definecolor{currentfill}{rgb}{0.000000,0.000000,0.000000}%
\pgfsetfillcolor{currentfill}%
\pgfsetlinewidth{0.000000pt}%
\definecolor{currentstroke}{rgb}{0.000000,0.000000,0.000000}%
\pgfsetstrokecolor{currentstroke}%
\pgfsetstrokeopacity{0.000000}%
\pgfsetdash{}{0pt}%
\pgfpathmoveto{\pgfqpoint{7.136900in}{0.660000in}}%
\pgfpathlineto{\pgfqpoint{7.143100in}{0.660000in}}%
\pgfpathlineto{\pgfqpoint{7.143100in}{3.213693in}}%
\pgfpathlineto{\pgfqpoint{7.136900in}{3.213693in}}%
\pgfpathlineto{\pgfqpoint{7.136900in}{0.660000in}}%
\pgfpathclose%
\pgfusepath{fill}%
\end{pgfscope}%
\begin{pgfscope}%
\pgfpathrectangle{\pgfqpoint{1.250000in}{0.660000in}}{\pgfqpoint{7.750000in}{4.620000in}}%
\pgfusepath{clip}%
\pgfsetbuttcap%
\pgfsetmiterjoin%
\definecolor{currentfill}{rgb}{0.000000,0.000000,0.000000}%
\pgfsetfillcolor{currentfill}%
\pgfsetlinewidth{0.000000pt}%
\definecolor{currentstroke}{rgb}{0.000000,0.000000,0.000000}%
\pgfsetstrokecolor{currentstroke}%
\pgfsetstrokeopacity{0.000000}%
\pgfsetdash{}{0pt}%
\pgfpathmoveto{\pgfqpoint{7.144650in}{0.660000in}}%
\pgfpathlineto{\pgfqpoint{7.150850in}{0.660000in}}%
\pgfpathlineto{\pgfqpoint{7.150850in}{3.213500in}}%
\pgfpathlineto{\pgfqpoint{7.144650in}{3.213500in}}%
\pgfpathlineto{\pgfqpoint{7.144650in}{0.660000in}}%
\pgfpathclose%
\pgfusepath{fill}%
\end{pgfscope}%
\begin{pgfscope}%
\pgfpathrectangle{\pgfqpoint{1.250000in}{0.660000in}}{\pgfqpoint{7.750000in}{4.620000in}}%
\pgfusepath{clip}%
\pgfsetbuttcap%
\pgfsetmiterjoin%
\definecolor{currentfill}{rgb}{0.000000,0.000000,0.000000}%
\pgfsetfillcolor{currentfill}%
\pgfsetlinewidth{0.000000pt}%
\definecolor{currentstroke}{rgb}{0.000000,0.000000,0.000000}%
\pgfsetstrokecolor{currentstroke}%
\pgfsetstrokeopacity{0.000000}%
\pgfsetdash{}{0pt}%
\pgfpathmoveto{\pgfqpoint{7.152400in}{0.660000in}}%
\pgfpathlineto{\pgfqpoint{7.158600in}{0.660000in}}%
\pgfpathlineto{\pgfqpoint{7.158600in}{3.213446in}}%
\pgfpathlineto{\pgfqpoint{7.152400in}{3.213446in}}%
\pgfpathlineto{\pgfqpoint{7.152400in}{0.660000in}}%
\pgfpathclose%
\pgfusepath{fill}%
\end{pgfscope}%
\begin{pgfscope}%
\pgfpathrectangle{\pgfqpoint{1.250000in}{0.660000in}}{\pgfqpoint{7.750000in}{4.620000in}}%
\pgfusepath{clip}%
\pgfsetbuttcap%
\pgfsetmiterjoin%
\definecolor{currentfill}{rgb}{0.000000,0.000000,0.000000}%
\pgfsetfillcolor{currentfill}%
\pgfsetlinewidth{0.000000pt}%
\definecolor{currentstroke}{rgb}{0.000000,0.000000,0.000000}%
\pgfsetstrokecolor{currentstroke}%
\pgfsetstrokeopacity{0.000000}%
\pgfsetdash{}{0pt}%
\pgfpathmoveto{\pgfqpoint{7.160150in}{0.660000in}}%
\pgfpathlineto{\pgfqpoint{7.166350in}{0.660000in}}%
\pgfpathlineto{\pgfqpoint{7.166350in}{3.213406in}}%
\pgfpathlineto{\pgfqpoint{7.160150in}{3.213406in}}%
\pgfpathlineto{\pgfqpoint{7.160150in}{0.660000in}}%
\pgfpathclose%
\pgfusepath{fill}%
\end{pgfscope}%
\begin{pgfscope}%
\pgfpathrectangle{\pgfqpoint{1.250000in}{0.660000in}}{\pgfqpoint{7.750000in}{4.620000in}}%
\pgfusepath{clip}%
\pgfsetbuttcap%
\pgfsetmiterjoin%
\definecolor{currentfill}{rgb}{0.000000,0.000000,0.000000}%
\pgfsetfillcolor{currentfill}%
\pgfsetlinewidth{0.000000pt}%
\definecolor{currentstroke}{rgb}{0.000000,0.000000,0.000000}%
\pgfsetstrokecolor{currentstroke}%
\pgfsetstrokeopacity{0.000000}%
\pgfsetdash{}{0pt}%
\pgfpathmoveto{\pgfqpoint{7.167900in}{0.660000in}}%
\pgfpathlineto{\pgfqpoint{7.174100in}{0.660000in}}%
\pgfpathlineto{\pgfqpoint{7.174100in}{3.213360in}}%
\pgfpathlineto{\pgfqpoint{7.167900in}{3.213360in}}%
\pgfpathlineto{\pgfqpoint{7.167900in}{0.660000in}}%
\pgfpathclose%
\pgfusepath{fill}%
\end{pgfscope}%
\begin{pgfscope}%
\pgfpathrectangle{\pgfqpoint{1.250000in}{0.660000in}}{\pgfqpoint{7.750000in}{4.620000in}}%
\pgfusepath{clip}%
\pgfsetbuttcap%
\pgfsetmiterjoin%
\definecolor{currentfill}{rgb}{0.000000,0.000000,0.000000}%
\pgfsetfillcolor{currentfill}%
\pgfsetlinewidth{0.000000pt}%
\definecolor{currentstroke}{rgb}{0.000000,0.000000,0.000000}%
\pgfsetstrokecolor{currentstroke}%
\pgfsetstrokeopacity{0.000000}%
\pgfsetdash{}{0pt}%
\pgfpathmoveto{\pgfqpoint{7.175650in}{0.660000in}}%
\pgfpathlineto{\pgfqpoint{7.181850in}{0.660000in}}%
\pgfpathlineto{\pgfqpoint{7.181850in}{3.213308in}}%
\pgfpathlineto{\pgfqpoint{7.175650in}{3.213308in}}%
\pgfpathlineto{\pgfqpoint{7.175650in}{0.660000in}}%
\pgfpathclose%
\pgfusepath{fill}%
\end{pgfscope}%
\begin{pgfscope}%
\pgfpathrectangle{\pgfqpoint{1.250000in}{0.660000in}}{\pgfqpoint{7.750000in}{4.620000in}}%
\pgfusepath{clip}%
\pgfsetbuttcap%
\pgfsetmiterjoin%
\definecolor{currentfill}{rgb}{0.000000,0.000000,0.000000}%
\pgfsetfillcolor{currentfill}%
\pgfsetlinewidth{0.000000pt}%
\definecolor{currentstroke}{rgb}{0.000000,0.000000,0.000000}%
\pgfsetstrokecolor{currentstroke}%
\pgfsetstrokeopacity{0.000000}%
\pgfsetdash{}{0pt}%
\pgfpathmoveto{\pgfqpoint{7.183400in}{0.660000in}}%
\pgfpathlineto{\pgfqpoint{7.189600in}{0.660000in}}%
\pgfpathlineto{\pgfqpoint{7.189600in}{3.213250in}}%
\pgfpathlineto{\pgfqpoint{7.183400in}{3.213250in}}%
\pgfpathlineto{\pgfqpoint{7.183400in}{0.660000in}}%
\pgfpathclose%
\pgfusepath{fill}%
\end{pgfscope}%
\begin{pgfscope}%
\pgfpathrectangle{\pgfqpoint{1.250000in}{0.660000in}}{\pgfqpoint{7.750000in}{4.620000in}}%
\pgfusepath{clip}%
\pgfsetbuttcap%
\pgfsetmiterjoin%
\definecolor{currentfill}{rgb}{0.000000,0.000000,0.000000}%
\pgfsetfillcolor{currentfill}%
\pgfsetlinewidth{0.000000pt}%
\definecolor{currentstroke}{rgb}{0.000000,0.000000,0.000000}%
\pgfsetstrokecolor{currentstroke}%
\pgfsetstrokeopacity{0.000000}%
\pgfsetdash{}{0pt}%
\pgfpathmoveto{\pgfqpoint{7.191150in}{0.660000in}}%
\pgfpathlineto{\pgfqpoint{7.197350in}{0.660000in}}%
\pgfpathlineto{\pgfqpoint{7.197350in}{3.213185in}}%
\pgfpathlineto{\pgfqpoint{7.191150in}{3.213185in}}%
\pgfpathlineto{\pgfqpoint{7.191150in}{0.660000in}}%
\pgfpathclose%
\pgfusepath{fill}%
\end{pgfscope}%
\begin{pgfscope}%
\pgfpathrectangle{\pgfqpoint{1.250000in}{0.660000in}}{\pgfqpoint{7.750000in}{4.620000in}}%
\pgfusepath{clip}%
\pgfsetbuttcap%
\pgfsetmiterjoin%
\definecolor{currentfill}{rgb}{0.000000,0.000000,0.000000}%
\pgfsetfillcolor{currentfill}%
\pgfsetlinewidth{0.000000pt}%
\definecolor{currentstroke}{rgb}{0.000000,0.000000,0.000000}%
\pgfsetstrokecolor{currentstroke}%
\pgfsetstrokeopacity{0.000000}%
\pgfsetdash{}{0pt}%
\pgfpathmoveto{\pgfqpoint{7.198900in}{0.660000in}}%
\pgfpathlineto{\pgfqpoint{7.205100in}{0.660000in}}%
\pgfpathlineto{\pgfqpoint{7.205100in}{3.213114in}}%
\pgfpathlineto{\pgfqpoint{7.198900in}{3.213114in}}%
\pgfpathlineto{\pgfqpoint{7.198900in}{0.660000in}}%
\pgfpathclose%
\pgfusepath{fill}%
\end{pgfscope}%
\begin{pgfscope}%
\pgfpathrectangle{\pgfqpoint{1.250000in}{0.660000in}}{\pgfqpoint{7.750000in}{4.620000in}}%
\pgfusepath{clip}%
\pgfsetbuttcap%
\pgfsetmiterjoin%
\definecolor{currentfill}{rgb}{0.000000,0.000000,0.000000}%
\pgfsetfillcolor{currentfill}%
\pgfsetlinewidth{0.000000pt}%
\definecolor{currentstroke}{rgb}{0.000000,0.000000,0.000000}%
\pgfsetstrokecolor{currentstroke}%
\pgfsetstrokeopacity{0.000000}%
\pgfsetdash{}{0pt}%
\pgfpathmoveto{\pgfqpoint{7.206650in}{0.660000in}}%
\pgfpathlineto{\pgfqpoint{7.212850in}{0.660000in}}%
\pgfpathlineto{\pgfqpoint{7.212850in}{3.224055in}}%
\pgfpathlineto{\pgfqpoint{7.206650in}{3.224055in}}%
\pgfpathlineto{\pgfqpoint{7.206650in}{0.660000in}}%
\pgfpathclose%
\pgfusepath{fill}%
\end{pgfscope}%
\begin{pgfscope}%
\pgfpathrectangle{\pgfqpoint{1.250000in}{0.660000in}}{\pgfqpoint{7.750000in}{4.620000in}}%
\pgfusepath{clip}%
\pgfsetbuttcap%
\pgfsetmiterjoin%
\definecolor{currentfill}{rgb}{0.000000,0.000000,0.000000}%
\pgfsetfillcolor{currentfill}%
\pgfsetlinewidth{0.000000pt}%
\definecolor{currentstroke}{rgb}{0.000000,0.000000,0.000000}%
\pgfsetstrokecolor{currentstroke}%
\pgfsetstrokeopacity{0.000000}%
\pgfsetdash{}{0pt}%
\pgfpathmoveto{\pgfqpoint{7.214400in}{0.660000in}}%
\pgfpathlineto{\pgfqpoint{7.220600in}{0.660000in}}%
\pgfpathlineto{\pgfqpoint{7.220600in}{3.214276in}}%
\pgfpathlineto{\pgfqpoint{7.214400in}{3.214276in}}%
\pgfpathlineto{\pgfqpoint{7.214400in}{0.660000in}}%
\pgfpathclose%
\pgfusepath{fill}%
\end{pgfscope}%
\begin{pgfscope}%
\pgfpathrectangle{\pgfqpoint{1.250000in}{0.660000in}}{\pgfqpoint{7.750000in}{4.620000in}}%
\pgfusepath{clip}%
\pgfsetbuttcap%
\pgfsetmiterjoin%
\definecolor{currentfill}{rgb}{0.000000,0.000000,0.000000}%
\pgfsetfillcolor{currentfill}%
\pgfsetlinewidth{0.000000pt}%
\definecolor{currentstroke}{rgb}{0.000000,0.000000,0.000000}%
\pgfsetstrokecolor{currentstroke}%
\pgfsetstrokeopacity{0.000000}%
\pgfsetdash{}{0pt}%
\pgfpathmoveto{\pgfqpoint{7.222150in}{0.660000in}}%
\pgfpathlineto{\pgfqpoint{7.228350in}{0.660000in}}%
\pgfpathlineto{\pgfqpoint{7.228350in}{3.221968in}}%
\pgfpathlineto{\pgfqpoint{7.222150in}{3.221968in}}%
\pgfpathlineto{\pgfqpoint{7.222150in}{0.660000in}}%
\pgfpathclose%
\pgfusepath{fill}%
\end{pgfscope}%
\begin{pgfscope}%
\pgfpathrectangle{\pgfqpoint{1.250000in}{0.660000in}}{\pgfqpoint{7.750000in}{4.620000in}}%
\pgfusepath{clip}%
\pgfsetbuttcap%
\pgfsetmiterjoin%
\definecolor{currentfill}{rgb}{0.000000,0.000000,0.000000}%
\pgfsetfillcolor{currentfill}%
\pgfsetlinewidth{0.000000pt}%
\definecolor{currentstroke}{rgb}{0.000000,0.000000,0.000000}%
\pgfsetstrokecolor{currentstroke}%
\pgfsetstrokeopacity{0.000000}%
\pgfsetdash{}{0pt}%
\pgfpathmoveto{\pgfqpoint{7.229900in}{0.660000in}}%
\pgfpathlineto{\pgfqpoint{7.236100in}{0.660000in}}%
\pgfpathlineto{\pgfqpoint{7.236100in}{3.212758in}}%
\pgfpathlineto{\pgfqpoint{7.229900in}{3.212758in}}%
\pgfpathlineto{\pgfqpoint{7.229900in}{0.660000in}}%
\pgfpathclose%
\pgfusepath{fill}%
\end{pgfscope}%
\begin{pgfscope}%
\pgfpathrectangle{\pgfqpoint{1.250000in}{0.660000in}}{\pgfqpoint{7.750000in}{4.620000in}}%
\pgfusepath{clip}%
\pgfsetbuttcap%
\pgfsetmiterjoin%
\definecolor{currentfill}{rgb}{0.000000,0.000000,0.000000}%
\pgfsetfillcolor{currentfill}%
\pgfsetlinewidth{0.000000pt}%
\definecolor{currentstroke}{rgb}{0.000000,0.000000,0.000000}%
\pgfsetstrokecolor{currentstroke}%
\pgfsetstrokeopacity{0.000000}%
\pgfsetdash{}{0pt}%
\pgfpathmoveto{\pgfqpoint{7.237650in}{0.660000in}}%
\pgfpathlineto{\pgfqpoint{7.243850in}{0.660000in}}%
\pgfpathlineto{\pgfqpoint{7.243850in}{3.212882in}}%
\pgfpathlineto{\pgfqpoint{7.237650in}{3.212882in}}%
\pgfpathlineto{\pgfqpoint{7.237650in}{0.660000in}}%
\pgfpathclose%
\pgfusepath{fill}%
\end{pgfscope}%
\begin{pgfscope}%
\pgfpathrectangle{\pgfqpoint{1.250000in}{0.660000in}}{\pgfqpoint{7.750000in}{4.620000in}}%
\pgfusepath{clip}%
\pgfsetbuttcap%
\pgfsetmiterjoin%
\definecolor{currentfill}{rgb}{0.000000,0.000000,0.000000}%
\pgfsetfillcolor{currentfill}%
\pgfsetlinewidth{0.000000pt}%
\definecolor{currentstroke}{rgb}{0.000000,0.000000,0.000000}%
\pgfsetstrokecolor{currentstroke}%
\pgfsetstrokeopacity{0.000000}%
\pgfsetdash{}{0pt}%
\pgfpathmoveto{\pgfqpoint{7.245400in}{0.660000in}}%
\pgfpathlineto{\pgfqpoint{7.251600in}{0.660000in}}%
\pgfpathlineto{\pgfqpoint{7.251600in}{3.212535in}}%
\pgfpathlineto{\pgfqpoint{7.245400in}{3.212535in}}%
\pgfpathlineto{\pgfqpoint{7.245400in}{0.660000in}}%
\pgfpathclose%
\pgfusepath{fill}%
\end{pgfscope}%
\begin{pgfscope}%
\pgfpathrectangle{\pgfqpoint{1.250000in}{0.660000in}}{\pgfqpoint{7.750000in}{4.620000in}}%
\pgfusepath{clip}%
\pgfsetbuttcap%
\pgfsetmiterjoin%
\definecolor{currentfill}{rgb}{0.000000,0.000000,0.000000}%
\pgfsetfillcolor{currentfill}%
\pgfsetlinewidth{0.000000pt}%
\definecolor{currentstroke}{rgb}{0.000000,0.000000,0.000000}%
\pgfsetstrokecolor{currentstroke}%
\pgfsetstrokeopacity{0.000000}%
\pgfsetdash{}{0pt}%
\pgfpathmoveto{\pgfqpoint{7.253150in}{0.660000in}}%
\pgfpathlineto{\pgfqpoint{7.259350in}{0.660000in}}%
\pgfpathlineto{\pgfqpoint{7.259350in}{3.212412in}}%
\pgfpathlineto{\pgfqpoint{7.253150in}{3.212412in}}%
\pgfpathlineto{\pgfqpoint{7.253150in}{0.660000in}}%
\pgfpathclose%
\pgfusepath{fill}%
\end{pgfscope}%
\begin{pgfscope}%
\pgfpathrectangle{\pgfqpoint{1.250000in}{0.660000in}}{\pgfqpoint{7.750000in}{4.620000in}}%
\pgfusepath{clip}%
\pgfsetbuttcap%
\pgfsetmiterjoin%
\definecolor{currentfill}{rgb}{0.000000,0.000000,0.000000}%
\pgfsetfillcolor{currentfill}%
\pgfsetlinewidth{0.000000pt}%
\definecolor{currentstroke}{rgb}{0.000000,0.000000,0.000000}%
\pgfsetstrokecolor{currentstroke}%
\pgfsetstrokeopacity{0.000000}%
\pgfsetdash{}{0pt}%
\pgfpathmoveto{\pgfqpoint{7.260900in}{0.660000in}}%
\pgfpathlineto{\pgfqpoint{7.267100in}{0.660000in}}%
\pgfpathlineto{\pgfqpoint{7.267100in}{3.212406in}}%
\pgfpathlineto{\pgfqpoint{7.260900in}{3.212406in}}%
\pgfpathlineto{\pgfqpoint{7.260900in}{0.660000in}}%
\pgfpathclose%
\pgfusepath{fill}%
\end{pgfscope}%
\begin{pgfscope}%
\pgfpathrectangle{\pgfqpoint{1.250000in}{0.660000in}}{\pgfqpoint{7.750000in}{4.620000in}}%
\pgfusepath{clip}%
\pgfsetbuttcap%
\pgfsetmiterjoin%
\definecolor{currentfill}{rgb}{0.000000,0.000000,0.000000}%
\pgfsetfillcolor{currentfill}%
\pgfsetlinewidth{0.000000pt}%
\definecolor{currentstroke}{rgb}{0.000000,0.000000,0.000000}%
\pgfsetstrokecolor{currentstroke}%
\pgfsetstrokeopacity{0.000000}%
\pgfsetdash{}{0pt}%
\pgfpathmoveto{\pgfqpoint{7.268650in}{0.660000in}}%
\pgfpathlineto{\pgfqpoint{7.274850in}{0.660000in}}%
\pgfpathlineto{\pgfqpoint{7.274850in}{3.212141in}}%
\pgfpathlineto{\pgfqpoint{7.268650in}{3.212141in}}%
\pgfpathlineto{\pgfqpoint{7.268650in}{0.660000in}}%
\pgfpathclose%
\pgfusepath{fill}%
\end{pgfscope}%
\begin{pgfscope}%
\pgfpathrectangle{\pgfqpoint{1.250000in}{0.660000in}}{\pgfqpoint{7.750000in}{4.620000in}}%
\pgfusepath{clip}%
\pgfsetbuttcap%
\pgfsetmiterjoin%
\definecolor{currentfill}{rgb}{0.000000,0.000000,0.000000}%
\pgfsetfillcolor{currentfill}%
\pgfsetlinewidth{0.000000pt}%
\definecolor{currentstroke}{rgb}{0.000000,0.000000,0.000000}%
\pgfsetstrokecolor{currentstroke}%
\pgfsetstrokeopacity{0.000000}%
\pgfsetdash{}{0pt}%
\pgfpathmoveto{\pgfqpoint{7.276400in}{0.660000in}}%
\pgfpathlineto{\pgfqpoint{7.282600in}{0.660000in}}%
\pgfpathlineto{\pgfqpoint{7.282600in}{3.211993in}}%
\pgfpathlineto{\pgfqpoint{7.276400in}{3.211993in}}%
\pgfpathlineto{\pgfqpoint{7.276400in}{0.660000in}}%
\pgfpathclose%
\pgfusepath{fill}%
\end{pgfscope}%
\begin{pgfscope}%
\pgfpathrectangle{\pgfqpoint{1.250000in}{0.660000in}}{\pgfqpoint{7.750000in}{4.620000in}}%
\pgfusepath{clip}%
\pgfsetbuttcap%
\pgfsetmiterjoin%
\definecolor{currentfill}{rgb}{0.000000,0.000000,0.000000}%
\pgfsetfillcolor{currentfill}%
\pgfsetlinewidth{0.000000pt}%
\definecolor{currentstroke}{rgb}{0.000000,0.000000,0.000000}%
\pgfsetstrokecolor{currentstroke}%
\pgfsetstrokeopacity{0.000000}%
\pgfsetdash{}{0pt}%
\pgfpathmoveto{\pgfqpoint{7.284150in}{0.660000in}}%
\pgfpathlineto{\pgfqpoint{7.290350in}{0.660000in}}%
\pgfpathlineto{\pgfqpoint{7.290350in}{3.215356in}}%
\pgfpathlineto{\pgfqpoint{7.284150in}{3.215356in}}%
\pgfpathlineto{\pgfqpoint{7.284150in}{0.660000in}}%
\pgfpathclose%
\pgfusepath{fill}%
\end{pgfscope}%
\begin{pgfscope}%
\pgfpathrectangle{\pgfqpoint{1.250000in}{0.660000in}}{\pgfqpoint{7.750000in}{4.620000in}}%
\pgfusepath{clip}%
\pgfsetbuttcap%
\pgfsetmiterjoin%
\definecolor{currentfill}{rgb}{0.000000,0.000000,0.000000}%
\pgfsetfillcolor{currentfill}%
\pgfsetlinewidth{0.000000pt}%
\definecolor{currentstroke}{rgb}{0.000000,0.000000,0.000000}%
\pgfsetstrokecolor{currentstroke}%
\pgfsetstrokeopacity{0.000000}%
\pgfsetdash{}{0pt}%
\pgfpathmoveto{\pgfqpoint{7.291900in}{0.660000in}}%
\pgfpathlineto{\pgfqpoint{7.298100in}{0.660000in}}%
\pgfpathlineto{\pgfqpoint{7.298100in}{3.214317in}}%
\pgfpathlineto{\pgfqpoint{7.291900in}{3.214317in}}%
\pgfpathlineto{\pgfqpoint{7.291900in}{0.660000in}}%
\pgfpathclose%
\pgfusepath{fill}%
\end{pgfscope}%
\begin{pgfscope}%
\pgfpathrectangle{\pgfqpoint{1.250000in}{0.660000in}}{\pgfqpoint{7.750000in}{4.620000in}}%
\pgfusepath{clip}%
\pgfsetbuttcap%
\pgfsetmiterjoin%
\definecolor{currentfill}{rgb}{0.000000,0.000000,0.000000}%
\pgfsetfillcolor{currentfill}%
\pgfsetlinewidth{0.000000pt}%
\definecolor{currentstroke}{rgb}{0.000000,0.000000,0.000000}%
\pgfsetstrokecolor{currentstroke}%
\pgfsetstrokeopacity{0.000000}%
\pgfsetdash{}{0pt}%
\pgfpathmoveto{\pgfqpoint{7.299650in}{0.660000in}}%
\pgfpathlineto{\pgfqpoint{7.305850in}{0.660000in}}%
\pgfpathlineto{\pgfqpoint{7.305850in}{3.211496in}}%
\pgfpathlineto{\pgfqpoint{7.299650in}{3.211496in}}%
\pgfpathlineto{\pgfqpoint{7.299650in}{0.660000in}}%
\pgfpathclose%
\pgfusepath{fill}%
\end{pgfscope}%
\begin{pgfscope}%
\pgfpathrectangle{\pgfqpoint{1.250000in}{0.660000in}}{\pgfqpoint{7.750000in}{4.620000in}}%
\pgfusepath{clip}%
\pgfsetbuttcap%
\pgfsetmiterjoin%
\definecolor{currentfill}{rgb}{0.000000,0.000000,0.000000}%
\pgfsetfillcolor{currentfill}%
\pgfsetlinewidth{0.000000pt}%
\definecolor{currentstroke}{rgb}{0.000000,0.000000,0.000000}%
\pgfsetstrokecolor{currentstroke}%
\pgfsetstrokeopacity{0.000000}%
\pgfsetdash{}{0pt}%
\pgfpathmoveto{\pgfqpoint{7.307400in}{0.660000in}}%
\pgfpathlineto{\pgfqpoint{7.313600in}{0.660000in}}%
\pgfpathlineto{\pgfqpoint{7.313600in}{3.218134in}}%
\pgfpathlineto{\pgfqpoint{7.307400in}{3.218134in}}%
\pgfpathlineto{\pgfqpoint{7.307400in}{0.660000in}}%
\pgfpathclose%
\pgfusepath{fill}%
\end{pgfscope}%
\begin{pgfscope}%
\pgfpathrectangle{\pgfqpoint{1.250000in}{0.660000in}}{\pgfqpoint{7.750000in}{4.620000in}}%
\pgfusepath{clip}%
\pgfsetbuttcap%
\pgfsetmiterjoin%
\definecolor{currentfill}{rgb}{0.000000,0.000000,0.000000}%
\pgfsetfillcolor{currentfill}%
\pgfsetlinewidth{0.000000pt}%
\definecolor{currentstroke}{rgb}{0.000000,0.000000,0.000000}%
\pgfsetstrokecolor{currentstroke}%
\pgfsetstrokeopacity{0.000000}%
\pgfsetdash{}{0pt}%
\pgfpathmoveto{\pgfqpoint{7.315150in}{0.660000in}}%
\pgfpathlineto{\pgfqpoint{7.321350in}{0.660000in}}%
\pgfpathlineto{\pgfqpoint{7.321350in}{3.211118in}}%
\pgfpathlineto{\pgfqpoint{7.315150in}{3.211118in}}%
\pgfpathlineto{\pgfqpoint{7.315150in}{0.660000in}}%
\pgfpathclose%
\pgfusepath{fill}%
\end{pgfscope}%
\begin{pgfscope}%
\pgfpathrectangle{\pgfqpoint{1.250000in}{0.660000in}}{\pgfqpoint{7.750000in}{4.620000in}}%
\pgfusepath{clip}%
\pgfsetbuttcap%
\pgfsetmiterjoin%
\definecolor{currentfill}{rgb}{0.000000,0.000000,0.000000}%
\pgfsetfillcolor{currentfill}%
\pgfsetlinewidth{0.000000pt}%
\definecolor{currentstroke}{rgb}{0.000000,0.000000,0.000000}%
\pgfsetstrokecolor{currentstroke}%
\pgfsetstrokeopacity{0.000000}%
\pgfsetdash{}{0pt}%
\pgfpathmoveto{\pgfqpoint{7.322900in}{0.660000in}}%
\pgfpathlineto{\pgfqpoint{7.329100in}{0.660000in}}%
\pgfpathlineto{\pgfqpoint{7.329100in}{3.210959in}}%
\pgfpathlineto{\pgfqpoint{7.322900in}{3.210959in}}%
\pgfpathlineto{\pgfqpoint{7.322900in}{0.660000in}}%
\pgfpathclose%
\pgfusepath{fill}%
\end{pgfscope}%
\begin{pgfscope}%
\pgfpathrectangle{\pgfqpoint{1.250000in}{0.660000in}}{\pgfqpoint{7.750000in}{4.620000in}}%
\pgfusepath{clip}%
\pgfsetbuttcap%
\pgfsetmiterjoin%
\definecolor{currentfill}{rgb}{0.000000,0.000000,0.000000}%
\pgfsetfillcolor{currentfill}%
\pgfsetlinewidth{0.000000pt}%
\definecolor{currentstroke}{rgb}{0.000000,0.000000,0.000000}%
\pgfsetstrokecolor{currentstroke}%
\pgfsetstrokeopacity{0.000000}%
\pgfsetdash{}{0pt}%
\pgfpathmoveto{\pgfqpoint{7.330650in}{0.660000in}}%
\pgfpathlineto{\pgfqpoint{7.336850in}{0.660000in}}%
\pgfpathlineto{\pgfqpoint{7.336850in}{3.210938in}}%
\pgfpathlineto{\pgfqpoint{7.330650in}{3.210938in}}%
\pgfpathlineto{\pgfqpoint{7.330650in}{0.660000in}}%
\pgfpathclose%
\pgfusepath{fill}%
\end{pgfscope}%
\begin{pgfscope}%
\pgfpathrectangle{\pgfqpoint{1.250000in}{0.660000in}}{\pgfqpoint{7.750000in}{4.620000in}}%
\pgfusepath{clip}%
\pgfsetbuttcap%
\pgfsetmiterjoin%
\definecolor{currentfill}{rgb}{0.000000,0.000000,0.000000}%
\pgfsetfillcolor{currentfill}%
\pgfsetlinewidth{0.000000pt}%
\definecolor{currentstroke}{rgb}{0.000000,0.000000,0.000000}%
\pgfsetstrokecolor{currentstroke}%
\pgfsetstrokeopacity{0.000000}%
\pgfsetdash{}{0pt}%
\pgfpathmoveto{\pgfqpoint{7.338400in}{0.660000in}}%
\pgfpathlineto{\pgfqpoint{7.344600in}{0.660000in}}%
\pgfpathlineto{\pgfqpoint{7.344600in}{3.210908in}}%
\pgfpathlineto{\pgfqpoint{7.338400in}{3.210908in}}%
\pgfpathlineto{\pgfqpoint{7.338400in}{0.660000in}}%
\pgfpathclose%
\pgfusepath{fill}%
\end{pgfscope}%
\begin{pgfscope}%
\pgfpathrectangle{\pgfqpoint{1.250000in}{0.660000in}}{\pgfqpoint{7.750000in}{4.620000in}}%
\pgfusepath{clip}%
\pgfsetbuttcap%
\pgfsetmiterjoin%
\definecolor{currentfill}{rgb}{0.000000,0.000000,0.000000}%
\pgfsetfillcolor{currentfill}%
\pgfsetlinewidth{0.000000pt}%
\definecolor{currentstroke}{rgb}{0.000000,0.000000,0.000000}%
\pgfsetstrokecolor{currentstroke}%
\pgfsetstrokeopacity{0.000000}%
\pgfsetdash{}{0pt}%
\pgfpathmoveto{\pgfqpoint{7.346150in}{0.660000in}}%
\pgfpathlineto{\pgfqpoint{7.352350in}{0.660000in}}%
\pgfpathlineto{\pgfqpoint{7.352350in}{3.210871in}}%
\pgfpathlineto{\pgfqpoint{7.346150in}{3.210871in}}%
\pgfpathlineto{\pgfqpoint{7.346150in}{0.660000in}}%
\pgfpathclose%
\pgfusepath{fill}%
\end{pgfscope}%
\begin{pgfscope}%
\pgfpathrectangle{\pgfqpoint{1.250000in}{0.660000in}}{\pgfqpoint{7.750000in}{4.620000in}}%
\pgfusepath{clip}%
\pgfsetbuttcap%
\pgfsetmiterjoin%
\definecolor{currentfill}{rgb}{0.000000,0.000000,0.000000}%
\pgfsetfillcolor{currentfill}%
\pgfsetlinewidth{0.000000pt}%
\definecolor{currentstroke}{rgb}{0.000000,0.000000,0.000000}%
\pgfsetstrokecolor{currentstroke}%
\pgfsetstrokeopacity{0.000000}%
\pgfsetdash{}{0pt}%
\pgfpathmoveto{\pgfqpoint{7.353900in}{0.660000in}}%
\pgfpathlineto{\pgfqpoint{7.360100in}{0.660000in}}%
\pgfpathlineto{\pgfqpoint{7.360100in}{3.210824in}}%
\pgfpathlineto{\pgfqpoint{7.353900in}{3.210824in}}%
\pgfpathlineto{\pgfqpoint{7.353900in}{0.660000in}}%
\pgfpathclose%
\pgfusepath{fill}%
\end{pgfscope}%
\begin{pgfscope}%
\pgfpathrectangle{\pgfqpoint{1.250000in}{0.660000in}}{\pgfqpoint{7.750000in}{4.620000in}}%
\pgfusepath{clip}%
\pgfsetbuttcap%
\pgfsetmiterjoin%
\definecolor{currentfill}{rgb}{0.000000,0.000000,0.000000}%
\pgfsetfillcolor{currentfill}%
\pgfsetlinewidth{0.000000pt}%
\definecolor{currentstroke}{rgb}{0.000000,0.000000,0.000000}%
\pgfsetstrokecolor{currentstroke}%
\pgfsetstrokeopacity{0.000000}%
\pgfsetdash{}{0pt}%
\pgfpathmoveto{\pgfqpoint{7.361650in}{0.660000in}}%
\pgfpathlineto{\pgfqpoint{7.367850in}{0.660000in}}%
\pgfpathlineto{\pgfqpoint{7.367850in}{3.210769in}}%
\pgfpathlineto{\pgfqpoint{7.361650in}{3.210769in}}%
\pgfpathlineto{\pgfqpoint{7.361650in}{0.660000in}}%
\pgfpathclose%
\pgfusepath{fill}%
\end{pgfscope}%
\begin{pgfscope}%
\pgfpathrectangle{\pgfqpoint{1.250000in}{0.660000in}}{\pgfqpoint{7.750000in}{4.620000in}}%
\pgfusepath{clip}%
\pgfsetbuttcap%
\pgfsetmiterjoin%
\definecolor{currentfill}{rgb}{0.000000,0.000000,0.000000}%
\pgfsetfillcolor{currentfill}%
\pgfsetlinewidth{0.000000pt}%
\definecolor{currentstroke}{rgb}{0.000000,0.000000,0.000000}%
\pgfsetstrokecolor{currentstroke}%
\pgfsetstrokeopacity{0.000000}%
\pgfsetdash{}{0pt}%
\pgfpathmoveto{\pgfqpoint{7.369400in}{0.660000in}}%
\pgfpathlineto{\pgfqpoint{7.375600in}{0.660000in}}%
\pgfpathlineto{\pgfqpoint{7.375600in}{3.210705in}}%
\pgfpathlineto{\pgfqpoint{7.369400in}{3.210705in}}%
\pgfpathlineto{\pgfqpoint{7.369400in}{0.660000in}}%
\pgfpathclose%
\pgfusepath{fill}%
\end{pgfscope}%
\begin{pgfscope}%
\pgfpathrectangle{\pgfqpoint{1.250000in}{0.660000in}}{\pgfqpoint{7.750000in}{4.620000in}}%
\pgfusepath{clip}%
\pgfsetbuttcap%
\pgfsetmiterjoin%
\definecolor{currentfill}{rgb}{0.000000,0.000000,0.000000}%
\pgfsetfillcolor{currentfill}%
\pgfsetlinewidth{0.000000pt}%
\definecolor{currentstroke}{rgb}{0.000000,0.000000,0.000000}%
\pgfsetstrokecolor{currentstroke}%
\pgfsetstrokeopacity{0.000000}%
\pgfsetdash{}{0pt}%
\pgfpathmoveto{\pgfqpoint{7.377150in}{0.660000in}}%
\pgfpathlineto{\pgfqpoint{7.383350in}{0.660000in}}%
\pgfpathlineto{\pgfqpoint{7.383350in}{3.210632in}}%
\pgfpathlineto{\pgfqpoint{7.377150in}{3.210632in}}%
\pgfpathlineto{\pgfqpoint{7.377150in}{0.660000in}}%
\pgfpathclose%
\pgfusepath{fill}%
\end{pgfscope}%
\begin{pgfscope}%
\pgfpathrectangle{\pgfqpoint{1.250000in}{0.660000in}}{\pgfqpoint{7.750000in}{4.620000in}}%
\pgfusepath{clip}%
\pgfsetbuttcap%
\pgfsetmiterjoin%
\definecolor{currentfill}{rgb}{0.000000,0.000000,0.000000}%
\pgfsetfillcolor{currentfill}%
\pgfsetlinewidth{0.000000pt}%
\definecolor{currentstroke}{rgb}{0.000000,0.000000,0.000000}%
\pgfsetstrokecolor{currentstroke}%
\pgfsetstrokeopacity{0.000000}%
\pgfsetdash{}{0pt}%
\pgfpathmoveto{\pgfqpoint{7.384900in}{0.660000in}}%
\pgfpathlineto{\pgfqpoint{7.391100in}{0.660000in}}%
\pgfpathlineto{\pgfqpoint{7.391100in}{3.210549in}}%
\pgfpathlineto{\pgfqpoint{7.384900in}{3.210549in}}%
\pgfpathlineto{\pgfqpoint{7.384900in}{0.660000in}}%
\pgfpathclose%
\pgfusepath{fill}%
\end{pgfscope}%
\begin{pgfscope}%
\pgfpathrectangle{\pgfqpoint{1.250000in}{0.660000in}}{\pgfqpoint{7.750000in}{4.620000in}}%
\pgfusepath{clip}%
\pgfsetbuttcap%
\pgfsetmiterjoin%
\definecolor{currentfill}{rgb}{0.000000,0.000000,0.000000}%
\pgfsetfillcolor{currentfill}%
\pgfsetlinewidth{0.000000pt}%
\definecolor{currentstroke}{rgb}{0.000000,0.000000,0.000000}%
\pgfsetstrokecolor{currentstroke}%
\pgfsetstrokeopacity{0.000000}%
\pgfsetdash{}{0pt}%
\pgfpathmoveto{\pgfqpoint{7.392650in}{0.660000in}}%
\pgfpathlineto{\pgfqpoint{7.398850in}{0.660000in}}%
\pgfpathlineto{\pgfqpoint{7.398850in}{3.210744in}}%
\pgfpathlineto{\pgfqpoint{7.392650in}{3.210744in}}%
\pgfpathlineto{\pgfqpoint{7.392650in}{0.660000in}}%
\pgfpathclose%
\pgfusepath{fill}%
\end{pgfscope}%
\begin{pgfscope}%
\pgfpathrectangle{\pgfqpoint{1.250000in}{0.660000in}}{\pgfqpoint{7.750000in}{4.620000in}}%
\pgfusepath{clip}%
\pgfsetbuttcap%
\pgfsetmiterjoin%
\definecolor{currentfill}{rgb}{0.000000,0.000000,0.000000}%
\pgfsetfillcolor{currentfill}%
\pgfsetlinewidth{0.000000pt}%
\definecolor{currentstroke}{rgb}{0.000000,0.000000,0.000000}%
\pgfsetstrokecolor{currentstroke}%
\pgfsetstrokeopacity{0.000000}%
\pgfsetdash{}{0pt}%
\pgfpathmoveto{\pgfqpoint{7.400400in}{0.660000in}}%
\pgfpathlineto{\pgfqpoint{7.406600in}{0.660000in}}%
\pgfpathlineto{\pgfqpoint{7.406600in}{3.229997in}}%
\pgfpathlineto{\pgfqpoint{7.400400in}{3.229997in}}%
\pgfpathlineto{\pgfqpoint{7.400400in}{0.660000in}}%
\pgfpathclose%
\pgfusepath{fill}%
\end{pgfscope}%
\begin{pgfscope}%
\pgfpathrectangle{\pgfqpoint{1.250000in}{0.660000in}}{\pgfqpoint{7.750000in}{4.620000in}}%
\pgfusepath{clip}%
\pgfsetbuttcap%
\pgfsetmiterjoin%
\definecolor{currentfill}{rgb}{0.000000,0.000000,0.000000}%
\pgfsetfillcolor{currentfill}%
\pgfsetlinewidth{0.000000pt}%
\definecolor{currentstroke}{rgb}{0.000000,0.000000,0.000000}%
\pgfsetstrokecolor{currentstroke}%
\pgfsetstrokeopacity{0.000000}%
\pgfsetdash{}{0pt}%
\pgfpathmoveto{\pgfqpoint{7.408150in}{0.660000in}}%
\pgfpathlineto{\pgfqpoint{7.414350in}{0.660000in}}%
\pgfpathlineto{\pgfqpoint{7.414350in}{3.210242in}}%
\pgfpathlineto{\pgfqpoint{7.408150in}{3.210242in}}%
\pgfpathlineto{\pgfqpoint{7.408150in}{0.660000in}}%
\pgfpathclose%
\pgfusepath{fill}%
\end{pgfscope}%
\begin{pgfscope}%
\pgfpathrectangle{\pgfqpoint{1.250000in}{0.660000in}}{\pgfqpoint{7.750000in}{4.620000in}}%
\pgfusepath{clip}%
\pgfsetbuttcap%
\pgfsetmiterjoin%
\definecolor{currentfill}{rgb}{0.000000,0.000000,0.000000}%
\pgfsetfillcolor{currentfill}%
\pgfsetlinewidth{0.000000pt}%
\definecolor{currentstroke}{rgb}{0.000000,0.000000,0.000000}%
\pgfsetstrokecolor{currentstroke}%
\pgfsetstrokeopacity{0.000000}%
\pgfsetdash{}{0pt}%
\pgfpathmoveto{\pgfqpoint{7.415900in}{0.660000in}}%
\pgfpathlineto{\pgfqpoint{7.422100in}{0.660000in}}%
\pgfpathlineto{\pgfqpoint{7.422100in}{3.210119in}}%
\pgfpathlineto{\pgfqpoint{7.415900in}{3.210119in}}%
\pgfpathlineto{\pgfqpoint{7.415900in}{0.660000in}}%
\pgfpathclose%
\pgfusepath{fill}%
\end{pgfscope}%
\begin{pgfscope}%
\pgfpathrectangle{\pgfqpoint{1.250000in}{0.660000in}}{\pgfqpoint{7.750000in}{4.620000in}}%
\pgfusepath{clip}%
\pgfsetbuttcap%
\pgfsetmiterjoin%
\definecolor{currentfill}{rgb}{0.000000,0.000000,0.000000}%
\pgfsetfillcolor{currentfill}%
\pgfsetlinewidth{0.000000pt}%
\definecolor{currentstroke}{rgb}{0.000000,0.000000,0.000000}%
\pgfsetstrokecolor{currentstroke}%
\pgfsetstrokeopacity{0.000000}%
\pgfsetdash{}{0pt}%
\pgfpathmoveto{\pgfqpoint{7.423650in}{0.660000in}}%
\pgfpathlineto{\pgfqpoint{7.429850in}{0.660000in}}%
\pgfpathlineto{\pgfqpoint{7.429850in}{3.209986in}}%
\pgfpathlineto{\pgfqpoint{7.423650in}{3.209986in}}%
\pgfpathlineto{\pgfqpoint{7.423650in}{0.660000in}}%
\pgfpathclose%
\pgfusepath{fill}%
\end{pgfscope}%
\begin{pgfscope}%
\pgfpathrectangle{\pgfqpoint{1.250000in}{0.660000in}}{\pgfqpoint{7.750000in}{4.620000in}}%
\pgfusepath{clip}%
\pgfsetbuttcap%
\pgfsetmiterjoin%
\definecolor{currentfill}{rgb}{0.000000,0.000000,0.000000}%
\pgfsetfillcolor{currentfill}%
\pgfsetlinewidth{0.000000pt}%
\definecolor{currentstroke}{rgb}{0.000000,0.000000,0.000000}%
\pgfsetstrokecolor{currentstroke}%
\pgfsetstrokeopacity{0.000000}%
\pgfsetdash{}{0pt}%
\pgfpathmoveto{\pgfqpoint{7.431400in}{0.660000in}}%
\pgfpathlineto{\pgfqpoint{7.437600in}{0.660000in}}%
\pgfpathlineto{\pgfqpoint{7.437600in}{3.209842in}}%
\pgfpathlineto{\pgfqpoint{7.431400in}{3.209842in}}%
\pgfpathlineto{\pgfqpoint{7.431400in}{0.660000in}}%
\pgfpathclose%
\pgfusepath{fill}%
\end{pgfscope}%
\begin{pgfscope}%
\pgfpathrectangle{\pgfqpoint{1.250000in}{0.660000in}}{\pgfqpoint{7.750000in}{4.620000in}}%
\pgfusepath{clip}%
\pgfsetbuttcap%
\pgfsetmiterjoin%
\definecolor{currentfill}{rgb}{0.000000,0.000000,0.000000}%
\pgfsetfillcolor{currentfill}%
\pgfsetlinewidth{0.000000pt}%
\definecolor{currentstroke}{rgb}{0.000000,0.000000,0.000000}%
\pgfsetstrokecolor{currentstroke}%
\pgfsetstrokeopacity{0.000000}%
\pgfsetdash{}{0pt}%
\pgfpathmoveto{\pgfqpoint{7.439150in}{0.660000in}}%
\pgfpathlineto{\pgfqpoint{7.445350in}{0.660000in}}%
\pgfpathlineto{\pgfqpoint{7.445350in}{3.209687in}}%
\pgfpathlineto{\pgfqpoint{7.439150in}{3.209687in}}%
\pgfpathlineto{\pgfqpoint{7.439150in}{0.660000in}}%
\pgfpathclose%
\pgfusepath{fill}%
\end{pgfscope}%
\begin{pgfscope}%
\pgfpathrectangle{\pgfqpoint{1.250000in}{0.660000in}}{\pgfqpoint{7.750000in}{4.620000in}}%
\pgfusepath{clip}%
\pgfsetbuttcap%
\pgfsetmiterjoin%
\definecolor{currentfill}{rgb}{0.000000,0.000000,0.000000}%
\pgfsetfillcolor{currentfill}%
\pgfsetlinewidth{0.000000pt}%
\definecolor{currentstroke}{rgb}{0.000000,0.000000,0.000000}%
\pgfsetstrokecolor{currentstroke}%
\pgfsetstrokeopacity{0.000000}%
\pgfsetdash{}{0pt}%
\pgfpathmoveto{\pgfqpoint{7.446900in}{0.660000in}}%
\pgfpathlineto{\pgfqpoint{7.453100in}{0.660000in}}%
\pgfpathlineto{\pgfqpoint{7.453100in}{3.209521in}}%
\pgfpathlineto{\pgfqpoint{7.446900in}{3.209521in}}%
\pgfpathlineto{\pgfqpoint{7.446900in}{0.660000in}}%
\pgfpathclose%
\pgfusepath{fill}%
\end{pgfscope}%
\begin{pgfscope}%
\pgfpathrectangle{\pgfqpoint{1.250000in}{0.660000in}}{\pgfqpoint{7.750000in}{4.620000in}}%
\pgfusepath{clip}%
\pgfsetbuttcap%
\pgfsetmiterjoin%
\definecolor{currentfill}{rgb}{0.000000,0.000000,0.000000}%
\pgfsetfillcolor{currentfill}%
\pgfsetlinewidth{0.000000pt}%
\definecolor{currentstroke}{rgb}{0.000000,0.000000,0.000000}%
\pgfsetstrokecolor{currentstroke}%
\pgfsetstrokeopacity{0.000000}%
\pgfsetdash{}{0pt}%
\pgfpathmoveto{\pgfqpoint{7.454650in}{0.660000in}}%
\pgfpathlineto{\pgfqpoint{7.460850in}{0.660000in}}%
\pgfpathlineto{\pgfqpoint{7.460850in}{3.209343in}}%
\pgfpathlineto{\pgfqpoint{7.454650in}{3.209343in}}%
\pgfpathlineto{\pgfqpoint{7.454650in}{0.660000in}}%
\pgfpathclose%
\pgfusepath{fill}%
\end{pgfscope}%
\begin{pgfscope}%
\pgfpathrectangle{\pgfqpoint{1.250000in}{0.660000in}}{\pgfqpoint{7.750000in}{4.620000in}}%
\pgfusepath{clip}%
\pgfsetbuttcap%
\pgfsetmiterjoin%
\definecolor{currentfill}{rgb}{0.000000,0.000000,0.000000}%
\pgfsetfillcolor{currentfill}%
\pgfsetlinewidth{0.000000pt}%
\definecolor{currentstroke}{rgb}{0.000000,0.000000,0.000000}%
\pgfsetstrokecolor{currentstroke}%
\pgfsetstrokeopacity{0.000000}%
\pgfsetdash{}{0pt}%
\pgfpathmoveto{\pgfqpoint{7.462400in}{0.660000in}}%
\pgfpathlineto{\pgfqpoint{7.468600in}{0.660000in}}%
\pgfpathlineto{\pgfqpoint{7.468600in}{3.209154in}}%
\pgfpathlineto{\pgfqpoint{7.462400in}{3.209154in}}%
\pgfpathlineto{\pgfqpoint{7.462400in}{0.660000in}}%
\pgfpathclose%
\pgfusepath{fill}%
\end{pgfscope}%
\begin{pgfscope}%
\pgfpathrectangle{\pgfqpoint{1.250000in}{0.660000in}}{\pgfqpoint{7.750000in}{4.620000in}}%
\pgfusepath{clip}%
\pgfsetbuttcap%
\pgfsetmiterjoin%
\definecolor{currentfill}{rgb}{0.000000,0.000000,0.000000}%
\pgfsetfillcolor{currentfill}%
\pgfsetlinewidth{0.000000pt}%
\definecolor{currentstroke}{rgb}{0.000000,0.000000,0.000000}%
\pgfsetstrokecolor{currentstroke}%
\pgfsetstrokeopacity{0.000000}%
\pgfsetdash{}{0pt}%
\pgfpathmoveto{\pgfqpoint{7.470150in}{0.660000in}}%
\pgfpathlineto{\pgfqpoint{7.476350in}{0.660000in}}%
\pgfpathlineto{\pgfqpoint{7.476350in}{3.208953in}}%
\pgfpathlineto{\pgfqpoint{7.470150in}{3.208953in}}%
\pgfpathlineto{\pgfqpoint{7.470150in}{0.660000in}}%
\pgfpathclose%
\pgfusepath{fill}%
\end{pgfscope}%
\begin{pgfscope}%
\pgfpathrectangle{\pgfqpoint{1.250000in}{0.660000in}}{\pgfqpoint{7.750000in}{4.620000in}}%
\pgfusepath{clip}%
\pgfsetbuttcap%
\pgfsetmiterjoin%
\definecolor{currentfill}{rgb}{0.000000,0.000000,0.000000}%
\pgfsetfillcolor{currentfill}%
\pgfsetlinewidth{0.000000pt}%
\definecolor{currentstroke}{rgb}{0.000000,0.000000,0.000000}%
\pgfsetstrokecolor{currentstroke}%
\pgfsetstrokeopacity{0.000000}%
\pgfsetdash{}{0pt}%
\pgfpathmoveto{\pgfqpoint{7.477900in}{0.660000in}}%
\pgfpathlineto{\pgfqpoint{7.484100in}{0.660000in}}%
\pgfpathlineto{\pgfqpoint{7.484100in}{3.208855in}}%
\pgfpathlineto{\pgfqpoint{7.477900in}{3.208855in}}%
\pgfpathlineto{\pgfqpoint{7.477900in}{0.660000in}}%
\pgfpathclose%
\pgfusepath{fill}%
\end{pgfscope}%
\begin{pgfscope}%
\pgfpathrectangle{\pgfqpoint{1.250000in}{0.660000in}}{\pgfqpoint{7.750000in}{4.620000in}}%
\pgfusepath{clip}%
\pgfsetbuttcap%
\pgfsetmiterjoin%
\definecolor{currentfill}{rgb}{0.000000,0.000000,0.000000}%
\pgfsetfillcolor{currentfill}%
\pgfsetlinewidth{0.000000pt}%
\definecolor{currentstroke}{rgb}{0.000000,0.000000,0.000000}%
\pgfsetstrokecolor{currentstroke}%
\pgfsetstrokeopacity{0.000000}%
\pgfsetdash{}{0pt}%
\pgfpathmoveto{\pgfqpoint{7.485650in}{0.660000in}}%
\pgfpathlineto{\pgfqpoint{7.491850in}{0.660000in}}%
\pgfpathlineto{\pgfqpoint{7.491850in}{3.208843in}}%
\pgfpathlineto{\pgfqpoint{7.485650in}{3.208843in}}%
\pgfpathlineto{\pgfqpoint{7.485650in}{0.660000in}}%
\pgfpathclose%
\pgfusepath{fill}%
\end{pgfscope}%
\begin{pgfscope}%
\pgfpathrectangle{\pgfqpoint{1.250000in}{0.660000in}}{\pgfqpoint{7.750000in}{4.620000in}}%
\pgfusepath{clip}%
\pgfsetbuttcap%
\pgfsetmiterjoin%
\definecolor{currentfill}{rgb}{0.000000,0.000000,0.000000}%
\pgfsetfillcolor{currentfill}%
\pgfsetlinewidth{0.000000pt}%
\definecolor{currentstroke}{rgb}{0.000000,0.000000,0.000000}%
\pgfsetstrokecolor{currentstroke}%
\pgfsetstrokeopacity{0.000000}%
\pgfsetdash{}{0pt}%
\pgfpathmoveto{\pgfqpoint{7.493400in}{0.660000in}}%
\pgfpathlineto{\pgfqpoint{7.499600in}{0.660000in}}%
\pgfpathlineto{\pgfqpoint{7.499600in}{3.208820in}}%
\pgfpathlineto{\pgfqpoint{7.493400in}{3.208820in}}%
\pgfpathlineto{\pgfqpoint{7.493400in}{0.660000in}}%
\pgfpathclose%
\pgfusepath{fill}%
\end{pgfscope}%
\begin{pgfscope}%
\pgfpathrectangle{\pgfqpoint{1.250000in}{0.660000in}}{\pgfqpoint{7.750000in}{4.620000in}}%
\pgfusepath{clip}%
\pgfsetbuttcap%
\pgfsetmiterjoin%
\definecolor{currentfill}{rgb}{0.000000,0.000000,0.000000}%
\pgfsetfillcolor{currentfill}%
\pgfsetlinewidth{0.000000pt}%
\definecolor{currentstroke}{rgb}{0.000000,0.000000,0.000000}%
\pgfsetstrokecolor{currentstroke}%
\pgfsetstrokeopacity{0.000000}%
\pgfsetdash{}{0pt}%
\pgfpathmoveto{\pgfqpoint{7.501150in}{0.660000in}}%
\pgfpathlineto{\pgfqpoint{7.507350in}{0.660000in}}%
\pgfpathlineto{\pgfqpoint{7.507350in}{3.208787in}}%
\pgfpathlineto{\pgfqpoint{7.501150in}{3.208787in}}%
\pgfpathlineto{\pgfqpoint{7.501150in}{0.660000in}}%
\pgfpathclose%
\pgfusepath{fill}%
\end{pgfscope}%
\begin{pgfscope}%
\pgfpathrectangle{\pgfqpoint{1.250000in}{0.660000in}}{\pgfqpoint{7.750000in}{4.620000in}}%
\pgfusepath{clip}%
\pgfsetbuttcap%
\pgfsetmiterjoin%
\definecolor{currentfill}{rgb}{0.000000,0.000000,0.000000}%
\pgfsetfillcolor{currentfill}%
\pgfsetlinewidth{0.000000pt}%
\definecolor{currentstroke}{rgb}{0.000000,0.000000,0.000000}%
\pgfsetstrokecolor{currentstroke}%
\pgfsetstrokeopacity{0.000000}%
\pgfsetdash{}{0pt}%
\pgfpathmoveto{\pgfqpoint{7.508900in}{0.660000in}}%
\pgfpathlineto{\pgfqpoint{7.515100in}{0.660000in}}%
\pgfpathlineto{\pgfqpoint{7.515100in}{3.208742in}}%
\pgfpathlineto{\pgfqpoint{7.508900in}{3.208742in}}%
\pgfpathlineto{\pgfqpoint{7.508900in}{0.660000in}}%
\pgfpathclose%
\pgfusepath{fill}%
\end{pgfscope}%
\begin{pgfscope}%
\pgfpathrectangle{\pgfqpoint{1.250000in}{0.660000in}}{\pgfqpoint{7.750000in}{4.620000in}}%
\pgfusepath{clip}%
\pgfsetbuttcap%
\pgfsetmiterjoin%
\definecolor{currentfill}{rgb}{0.000000,0.000000,0.000000}%
\pgfsetfillcolor{currentfill}%
\pgfsetlinewidth{0.000000pt}%
\definecolor{currentstroke}{rgb}{0.000000,0.000000,0.000000}%
\pgfsetstrokecolor{currentstroke}%
\pgfsetstrokeopacity{0.000000}%
\pgfsetdash{}{0pt}%
\pgfpathmoveto{\pgfqpoint{7.516650in}{0.660000in}}%
\pgfpathlineto{\pgfqpoint{7.522850in}{0.660000in}}%
\pgfpathlineto{\pgfqpoint{7.522850in}{3.208686in}}%
\pgfpathlineto{\pgfqpoint{7.516650in}{3.208686in}}%
\pgfpathlineto{\pgfqpoint{7.516650in}{0.660000in}}%
\pgfpathclose%
\pgfusepath{fill}%
\end{pgfscope}%
\begin{pgfscope}%
\pgfpathrectangle{\pgfqpoint{1.250000in}{0.660000in}}{\pgfqpoint{7.750000in}{4.620000in}}%
\pgfusepath{clip}%
\pgfsetbuttcap%
\pgfsetmiterjoin%
\definecolor{currentfill}{rgb}{0.000000,0.000000,0.000000}%
\pgfsetfillcolor{currentfill}%
\pgfsetlinewidth{0.000000pt}%
\definecolor{currentstroke}{rgb}{0.000000,0.000000,0.000000}%
\pgfsetstrokecolor{currentstroke}%
\pgfsetstrokeopacity{0.000000}%
\pgfsetdash{}{0pt}%
\pgfpathmoveto{\pgfqpoint{7.524400in}{0.660000in}}%
\pgfpathlineto{\pgfqpoint{7.530600in}{0.660000in}}%
\pgfpathlineto{\pgfqpoint{7.530600in}{3.208619in}}%
\pgfpathlineto{\pgfqpoint{7.524400in}{3.208619in}}%
\pgfpathlineto{\pgfqpoint{7.524400in}{0.660000in}}%
\pgfpathclose%
\pgfusepath{fill}%
\end{pgfscope}%
\begin{pgfscope}%
\pgfpathrectangle{\pgfqpoint{1.250000in}{0.660000in}}{\pgfqpoint{7.750000in}{4.620000in}}%
\pgfusepath{clip}%
\pgfsetbuttcap%
\pgfsetmiterjoin%
\definecolor{currentfill}{rgb}{0.000000,0.000000,0.000000}%
\pgfsetfillcolor{currentfill}%
\pgfsetlinewidth{0.000000pt}%
\definecolor{currentstroke}{rgb}{0.000000,0.000000,0.000000}%
\pgfsetstrokecolor{currentstroke}%
\pgfsetstrokeopacity{0.000000}%
\pgfsetdash{}{0pt}%
\pgfpathmoveto{\pgfqpoint{7.532150in}{0.660000in}}%
\pgfpathlineto{\pgfqpoint{7.538350in}{0.660000in}}%
\pgfpathlineto{\pgfqpoint{7.538350in}{3.208540in}}%
\pgfpathlineto{\pgfqpoint{7.532150in}{3.208540in}}%
\pgfpathlineto{\pgfqpoint{7.532150in}{0.660000in}}%
\pgfpathclose%
\pgfusepath{fill}%
\end{pgfscope}%
\begin{pgfscope}%
\pgfpathrectangle{\pgfqpoint{1.250000in}{0.660000in}}{\pgfqpoint{7.750000in}{4.620000in}}%
\pgfusepath{clip}%
\pgfsetbuttcap%
\pgfsetmiterjoin%
\definecolor{currentfill}{rgb}{0.000000,0.000000,0.000000}%
\pgfsetfillcolor{currentfill}%
\pgfsetlinewidth{0.000000pt}%
\definecolor{currentstroke}{rgb}{0.000000,0.000000,0.000000}%
\pgfsetstrokecolor{currentstroke}%
\pgfsetstrokeopacity{0.000000}%
\pgfsetdash{}{0pt}%
\pgfpathmoveto{\pgfqpoint{7.539900in}{0.660000in}}%
\pgfpathlineto{\pgfqpoint{7.546100in}{0.660000in}}%
\pgfpathlineto{\pgfqpoint{7.546100in}{3.208450in}}%
\pgfpathlineto{\pgfqpoint{7.539900in}{3.208450in}}%
\pgfpathlineto{\pgfqpoint{7.539900in}{0.660000in}}%
\pgfpathclose%
\pgfusepath{fill}%
\end{pgfscope}%
\begin{pgfscope}%
\pgfpathrectangle{\pgfqpoint{1.250000in}{0.660000in}}{\pgfqpoint{7.750000in}{4.620000in}}%
\pgfusepath{clip}%
\pgfsetbuttcap%
\pgfsetmiterjoin%
\definecolor{currentfill}{rgb}{0.000000,0.000000,0.000000}%
\pgfsetfillcolor{currentfill}%
\pgfsetlinewidth{0.000000pt}%
\definecolor{currentstroke}{rgb}{0.000000,0.000000,0.000000}%
\pgfsetstrokecolor{currentstroke}%
\pgfsetstrokeopacity{0.000000}%
\pgfsetdash{}{0pt}%
\pgfpathmoveto{\pgfqpoint{7.547650in}{0.660000in}}%
\pgfpathlineto{\pgfqpoint{7.553850in}{0.660000in}}%
\pgfpathlineto{\pgfqpoint{7.553850in}{3.208347in}}%
\pgfpathlineto{\pgfqpoint{7.547650in}{3.208347in}}%
\pgfpathlineto{\pgfqpoint{7.547650in}{0.660000in}}%
\pgfpathclose%
\pgfusepath{fill}%
\end{pgfscope}%
\begin{pgfscope}%
\pgfpathrectangle{\pgfqpoint{1.250000in}{0.660000in}}{\pgfqpoint{7.750000in}{4.620000in}}%
\pgfusepath{clip}%
\pgfsetbuttcap%
\pgfsetmiterjoin%
\definecolor{currentfill}{rgb}{0.000000,0.000000,0.000000}%
\pgfsetfillcolor{currentfill}%
\pgfsetlinewidth{0.000000pt}%
\definecolor{currentstroke}{rgb}{0.000000,0.000000,0.000000}%
\pgfsetstrokecolor{currentstroke}%
\pgfsetstrokeopacity{0.000000}%
\pgfsetdash{}{0pt}%
\pgfpathmoveto{\pgfqpoint{7.555400in}{0.660000in}}%
\pgfpathlineto{\pgfqpoint{7.561600in}{0.660000in}}%
\pgfpathlineto{\pgfqpoint{7.561600in}{3.208231in}}%
\pgfpathlineto{\pgfqpoint{7.555400in}{3.208231in}}%
\pgfpathlineto{\pgfqpoint{7.555400in}{0.660000in}}%
\pgfpathclose%
\pgfusepath{fill}%
\end{pgfscope}%
\begin{pgfscope}%
\pgfpathrectangle{\pgfqpoint{1.250000in}{0.660000in}}{\pgfqpoint{7.750000in}{4.620000in}}%
\pgfusepath{clip}%
\pgfsetbuttcap%
\pgfsetmiterjoin%
\definecolor{currentfill}{rgb}{0.000000,0.000000,0.000000}%
\pgfsetfillcolor{currentfill}%
\pgfsetlinewidth{0.000000pt}%
\definecolor{currentstroke}{rgb}{0.000000,0.000000,0.000000}%
\pgfsetstrokecolor{currentstroke}%
\pgfsetstrokeopacity{0.000000}%
\pgfsetdash{}{0pt}%
\pgfpathmoveto{\pgfqpoint{7.563150in}{0.660000in}}%
\pgfpathlineto{\pgfqpoint{7.569350in}{0.660000in}}%
\pgfpathlineto{\pgfqpoint{7.569350in}{3.208103in}}%
\pgfpathlineto{\pgfqpoint{7.563150in}{3.208103in}}%
\pgfpathlineto{\pgfqpoint{7.563150in}{0.660000in}}%
\pgfpathclose%
\pgfusepath{fill}%
\end{pgfscope}%
\begin{pgfscope}%
\pgfpathrectangle{\pgfqpoint{1.250000in}{0.660000in}}{\pgfqpoint{7.750000in}{4.620000in}}%
\pgfusepath{clip}%
\pgfsetbuttcap%
\pgfsetmiterjoin%
\definecolor{currentfill}{rgb}{0.000000,0.000000,0.000000}%
\pgfsetfillcolor{currentfill}%
\pgfsetlinewidth{0.000000pt}%
\definecolor{currentstroke}{rgb}{0.000000,0.000000,0.000000}%
\pgfsetstrokecolor{currentstroke}%
\pgfsetstrokeopacity{0.000000}%
\pgfsetdash{}{0pt}%
\pgfpathmoveto{\pgfqpoint{7.570900in}{0.660000in}}%
\pgfpathlineto{\pgfqpoint{7.577100in}{0.660000in}}%
\pgfpathlineto{\pgfqpoint{7.577100in}{3.207962in}}%
\pgfpathlineto{\pgfqpoint{7.570900in}{3.207962in}}%
\pgfpathlineto{\pgfqpoint{7.570900in}{0.660000in}}%
\pgfpathclose%
\pgfusepath{fill}%
\end{pgfscope}%
\begin{pgfscope}%
\pgfpathrectangle{\pgfqpoint{1.250000in}{0.660000in}}{\pgfqpoint{7.750000in}{4.620000in}}%
\pgfusepath{clip}%
\pgfsetbuttcap%
\pgfsetmiterjoin%
\definecolor{currentfill}{rgb}{0.000000,0.000000,0.000000}%
\pgfsetfillcolor{currentfill}%
\pgfsetlinewidth{0.000000pt}%
\definecolor{currentstroke}{rgb}{0.000000,0.000000,0.000000}%
\pgfsetstrokecolor{currentstroke}%
\pgfsetstrokeopacity{0.000000}%
\pgfsetdash{}{0pt}%
\pgfpathmoveto{\pgfqpoint{7.578650in}{0.660000in}}%
\pgfpathlineto{\pgfqpoint{7.584850in}{0.660000in}}%
\pgfpathlineto{\pgfqpoint{7.584850in}{3.207808in}}%
\pgfpathlineto{\pgfqpoint{7.578650in}{3.207808in}}%
\pgfpathlineto{\pgfqpoint{7.578650in}{0.660000in}}%
\pgfpathclose%
\pgfusepath{fill}%
\end{pgfscope}%
\begin{pgfscope}%
\pgfpathrectangle{\pgfqpoint{1.250000in}{0.660000in}}{\pgfqpoint{7.750000in}{4.620000in}}%
\pgfusepath{clip}%
\pgfsetbuttcap%
\pgfsetmiterjoin%
\definecolor{currentfill}{rgb}{0.000000,0.000000,0.000000}%
\pgfsetfillcolor{currentfill}%
\pgfsetlinewidth{0.000000pt}%
\definecolor{currentstroke}{rgb}{0.000000,0.000000,0.000000}%
\pgfsetstrokecolor{currentstroke}%
\pgfsetstrokeopacity{0.000000}%
\pgfsetdash{}{0pt}%
\pgfpathmoveto{\pgfqpoint{7.586400in}{0.660000in}}%
\pgfpathlineto{\pgfqpoint{7.592600in}{0.660000in}}%
\pgfpathlineto{\pgfqpoint{7.592600in}{3.207640in}}%
\pgfpathlineto{\pgfqpoint{7.586400in}{3.207640in}}%
\pgfpathlineto{\pgfqpoint{7.586400in}{0.660000in}}%
\pgfpathclose%
\pgfusepath{fill}%
\end{pgfscope}%
\begin{pgfscope}%
\pgfpathrectangle{\pgfqpoint{1.250000in}{0.660000in}}{\pgfqpoint{7.750000in}{4.620000in}}%
\pgfusepath{clip}%
\pgfsetbuttcap%
\pgfsetmiterjoin%
\definecolor{currentfill}{rgb}{0.000000,0.000000,0.000000}%
\pgfsetfillcolor{currentfill}%
\pgfsetlinewidth{0.000000pt}%
\definecolor{currentstroke}{rgb}{0.000000,0.000000,0.000000}%
\pgfsetstrokecolor{currentstroke}%
\pgfsetstrokeopacity{0.000000}%
\pgfsetdash{}{0pt}%
\pgfpathmoveto{\pgfqpoint{7.594150in}{0.660000in}}%
\pgfpathlineto{\pgfqpoint{7.600350in}{0.660000in}}%
\pgfpathlineto{\pgfqpoint{7.600350in}{3.207458in}}%
\pgfpathlineto{\pgfqpoint{7.594150in}{3.207458in}}%
\pgfpathlineto{\pgfqpoint{7.594150in}{0.660000in}}%
\pgfpathclose%
\pgfusepath{fill}%
\end{pgfscope}%
\begin{pgfscope}%
\pgfpathrectangle{\pgfqpoint{1.250000in}{0.660000in}}{\pgfqpoint{7.750000in}{4.620000in}}%
\pgfusepath{clip}%
\pgfsetbuttcap%
\pgfsetmiterjoin%
\definecolor{currentfill}{rgb}{0.000000,0.000000,0.000000}%
\pgfsetfillcolor{currentfill}%
\pgfsetlinewidth{0.000000pt}%
\definecolor{currentstroke}{rgb}{0.000000,0.000000,0.000000}%
\pgfsetstrokecolor{currentstroke}%
\pgfsetstrokeopacity{0.000000}%
\pgfsetdash{}{0pt}%
\pgfpathmoveto{\pgfqpoint{7.601900in}{0.660000in}}%
\pgfpathlineto{\pgfqpoint{7.608100in}{0.660000in}}%
\pgfpathlineto{\pgfqpoint{7.608100in}{3.207263in}}%
\pgfpathlineto{\pgfqpoint{7.601900in}{3.207263in}}%
\pgfpathlineto{\pgfqpoint{7.601900in}{0.660000in}}%
\pgfpathclose%
\pgfusepath{fill}%
\end{pgfscope}%
\begin{pgfscope}%
\pgfpathrectangle{\pgfqpoint{1.250000in}{0.660000in}}{\pgfqpoint{7.750000in}{4.620000in}}%
\pgfusepath{clip}%
\pgfsetbuttcap%
\pgfsetmiterjoin%
\definecolor{currentfill}{rgb}{0.000000,0.000000,0.000000}%
\pgfsetfillcolor{currentfill}%
\pgfsetlinewidth{0.000000pt}%
\definecolor{currentstroke}{rgb}{0.000000,0.000000,0.000000}%
\pgfsetstrokecolor{currentstroke}%
\pgfsetstrokeopacity{0.000000}%
\pgfsetdash{}{0pt}%
\pgfpathmoveto{\pgfqpoint{7.609650in}{0.660000in}}%
\pgfpathlineto{\pgfqpoint{7.615850in}{0.660000in}}%
\pgfpathlineto{\pgfqpoint{7.615850in}{3.212696in}}%
\pgfpathlineto{\pgfqpoint{7.609650in}{3.212696in}}%
\pgfpathlineto{\pgfqpoint{7.609650in}{0.660000in}}%
\pgfpathclose%
\pgfusepath{fill}%
\end{pgfscope}%
\begin{pgfscope}%
\pgfpathrectangle{\pgfqpoint{1.250000in}{0.660000in}}{\pgfqpoint{7.750000in}{4.620000in}}%
\pgfusepath{clip}%
\pgfsetbuttcap%
\pgfsetmiterjoin%
\definecolor{currentfill}{rgb}{0.000000,0.000000,0.000000}%
\pgfsetfillcolor{currentfill}%
\pgfsetlinewidth{0.000000pt}%
\definecolor{currentstroke}{rgb}{0.000000,0.000000,0.000000}%
\pgfsetstrokecolor{currentstroke}%
\pgfsetstrokeopacity{0.000000}%
\pgfsetdash{}{0pt}%
\pgfpathmoveto{\pgfqpoint{7.617400in}{0.660000in}}%
\pgfpathlineto{\pgfqpoint{7.623600in}{0.660000in}}%
\pgfpathlineto{\pgfqpoint{7.623600in}{3.209780in}}%
\pgfpathlineto{\pgfqpoint{7.617400in}{3.209780in}}%
\pgfpathlineto{\pgfqpoint{7.617400in}{0.660000in}}%
\pgfpathclose%
\pgfusepath{fill}%
\end{pgfscope}%
\begin{pgfscope}%
\pgfpathrectangle{\pgfqpoint{1.250000in}{0.660000in}}{\pgfqpoint{7.750000in}{4.620000in}}%
\pgfusepath{clip}%
\pgfsetbuttcap%
\pgfsetmiterjoin%
\definecolor{currentfill}{rgb}{0.000000,0.000000,0.000000}%
\pgfsetfillcolor{currentfill}%
\pgfsetlinewidth{0.000000pt}%
\definecolor{currentstroke}{rgb}{0.000000,0.000000,0.000000}%
\pgfsetstrokecolor{currentstroke}%
\pgfsetstrokeopacity{0.000000}%
\pgfsetdash{}{0pt}%
\pgfpathmoveto{\pgfqpoint{7.625150in}{0.660000in}}%
\pgfpathlineto{\pgfqpoint{7.631350in}{0.660000in}}%
\pgfpathlineto{\pgfqpoint{7.631350in}{3.207025in}}%
\pgfpathlineto{\pgfqpoint{7.625150in}{3.207025in}}%
\pgfpathlineto{\pgfqpoint{7.625150in}{0.660000in}}%
\pgfpathclose%
\pgfusepath{fill}%
\end{pgfscope}%
\begin{pgfscope}%
\pgfpathrectangle{\pgfqpoint{1.250000in}{0.660000in}}{\pgfqpoint{7.750000in}{4.620000in}}%
\pgfusepath{clip}%
\pgfsetbuttcap%
\pgfsetmiterjoin%
\definecolor{currentfill}{rgb}{0.000000,0.000000,0.000000}%
\pgfsetfillcolor{currentfill}%
\pgfsetlinewidth{0.000000pt}%
\definecolor{currentstroke}{rgb}{0.000000,0.000000,0.000000}%
\pgfsetstrokecolor{currentstroke}%
\pgfsetstrokeopacity{0.000000}%
\pgfsetdash{}{0pt}%
\pgfpathmoveto{\pgfqpoint{7.632900in}{0.660000in}}%
\pgfpathlineto{\pgfqpoint{7.639100in}{0.660000in}}%
\pgfpathlineto{\pgfqpoint{7.639100in}{3.207005in}}%
\pgfpathlineto{\pgfqpoint{7.632900in}{3.207005in}}%
\pgfpathlineto{\pgfqpoint{7.632900in}{0.660000in}}%
\pgfpathclose%
\pgfusepath{fill}%
\end{pgfscope}%
\begin{pgfscope}%
\pgfpathrectangle{\pgfqpoint{1.250000in}{0.660000in}}{\pgfqpoint{7.750000in}{4.620000in}}%
\pgfusepath{clip}%
\pgfsetbuttcap%
\pgfsetmiterjoin%
\definecolor{currentfill}{rgb}{0.000000,0.000000,0.000000}%
\pgfsetfillcolor{currentfill}%
\pgfsetlinewidth{0.000000pt}%
\definecolor{currentstroke}{rgb}{0.000000,0.000000,0.000000}%
\pgfsetstrokecolor{currentstroke}%
\pgfsetstrokeopacity{0.000000}%
\pgfsetdash{}{0pt}%
\pgfpathmoveto{\pgfqpoint{7.640650in}{0.660000in}}%
\pgfpathlineto{\pgfqpoint{7.646850in}{0.660000in}}%
\pgfpathlineto{\pgfqpoint{7.646850in}{3.208523in}}%
\pgfpathlineto{\pgfqpoint{7.640650in}{3.208523in}}%
\pgfpathlineto{\pgfqpoint{7.640650in}{0.660000in}}%
\pgfpathclose%
\pgfusepath{fill}%
\end{pgfscope}%
\begin{pgfscope}%
\pgfpathrectangle{\pgfqpoint{1.250000in}{0.660000in}}{\pgfqpoint{7.750000in}{4.620000in}}%
\pgfusepath{clip}%
\pgfsetbuttcap%
\pgfsetmiterjoin%
\definecolor{currentfill}{rgb}{0.000000,0.000000,0.000000}%
\pgfsetfillcolor{currentfill}%
\pgfsetlinewidth{0.000000pt}%
\definecolor{currentstroke}{rgb}{0.000000,0.000000,0.000000}%
\pgfsetstrokecolor{currentstroke}%
\pgfsetstrokeopacity{0.000000}%
\pgfsetdash{}{0pt}%
\pgfpathmoveto{\pgfqpoint{7.648400in}{0.660000in}}%
\pgfpathlineto{\pgfqpoint{7.654600in}{0.660000in}}%
\pgfpathlineto{\pgfqpoint{7.654600in}{3.221619in}}%
\pgfpathlineto{\pgfqpoint{7.648400in}{3.221619in}}%
\pgfpathlineto{\pgfqpoint{7.648400in}{0.660000in}}%
\pgfpathclose%
\pgfusepath{fill}%
\end{pgfscope}%
\begin{pgfscope}%
\pgfpathrectangle{\pgfqpoint{1.250000in}{0.660000in}}{\pgfqpoint{7.750000in}{4.620000in}}%
\pgfusepath{clip}%
\pgfsetbuttcap%
\pgfsetmiterjoin%
\definecolor{currentfill}{rgb}{0.000000,0.000000,0.000000}%
\pgfsetfillcolor{currentfill}%
\pgfsetlinewidth{0.000000pt}%
\definecolor{currentstroke}{rgb}{0.000000,0.000000,0.000000}%
\pgfsetstrokecolor{currentstroke}%
\pgfsetstrokeopacity{0.000000}%
\pgfsetdash{}{0pt}%
\pgfpathmoveto{\pgfqpoint{7.656150in}{0.660000in}}%
\pgfpathlineto{\pgfqpoint{7.662350in}{0.660000in}}%
\pgfpathlineto{\pgfqpoint{7.662350in}{3.206867in}}%
\pgfpathlineto{\pgfqpoint{7.656150in}{3.206867in}}%
\pgfpathlineto{\pgfqpoint{7.656150in}{0.660000in}}%
\pgfpathclose%
\pgfusepath{fill}%
\end{pgfscope}%
\begin{pgfscope}%
\pgfpathrectangle{\pgfqpoint{1.250000in}{0.660000in}}{\pgfqpoint{7.750000in}{4.620000in}}%
\pgfusepath{clip}%
\pgfsetbuttcap%
\pgfsetmiterjoin%
\definecolor{currentfill}{rgb}{0.000000,0.000000,0.000000}%
\pgfsetfillcolor{currentfill}%
\pgfsetlinewidth{0.000000pt}%
\definecolor{currentstroke}{rgb}{0.000000,0.000000,0.000000}%
\pgfsetstrokecolor{currentstroke}%
\pgfsetstrokeopacity{0.000000}%
\pgfsetdash{}{0pt}%
\pgfpathmoveto{\pgfqpoint{7.663900in}{0.660000in}}%
\pgfpathlineto{\pgfqpoint{7.670100in}{0.660000in}}%
\pgfpathlineto{\pgfqpoint{7.670100in}{3.206793in}}%
\pgfpathlineto{\pgfqpoint{7.663900in}{3.206793in}}%
\pgfpathlineto{\pgfqpoint{7.663900in}{0.660000in}}%
\pgfpathclose%
\pgfusepath{fill}%
\end{pgfscope}%
\begin{pgfscope}%
\pgfpathrectangle{\pgfqpoint{1.250000in}{0.660000in}}{\pgfqpoint{7.750000in}{4.620000in}}%
\pgfusepath{clip}%
\pgfsetbuttcap%
\pgfsetmiterjoin%
\definecolor{currentfill}{rgb}{0.000000,0.000000,0.000000}%
\pgfsetfillcolor{currentfill}%
\pgfsetlinewidth{0.000000pt}%
\definecolor{currentstroke}{rgb}{0.000000,0.000000,0.000000}%
\pgfsetstrokecolor{currentstroke}%
\pgfsetstrokeopacity{0.000000}%
\pgfsetdash{}{0pt}%
\pgfpathmoveto{\pgfqpoint{7.671650in}{0.660000in}}%
\pgfpathlineto{\pgfqpoint{7.677850in}{0.660000in}}%
\pgfpathlineto{\pgfqpoint{7.677850in}{3.206706in}}%
\pgfpathlineto{\pgfqpoint{7.671650in}{3.206706in}}%
\pgfpathlineto{\pgfqpoint{7.671650in}{0.660000in}}%
\pgfpathclose%
\pgfusepath{fill}%
\end{pgfscope}%
\begin{pgfscope}%
\pgfpathrectangle{\pgfqpoint{1.250000in}{0.660000in}}{\pgfqpoint{7.750000in}{4.620000in}}%
\pgfusepath{clip}%
\pgfsetbuttcap%
\pgfsetmiterjoin%
\definecolor{currentfill}{rgb}{0.000000,0.000000,0.000000}%
\pgfsetfillcolor{currentfill}%
\pgfsetlinewidth{0.000000pt}%
\definecolor{currentstroke}{rgb}{0.000000,0.000000,0.000000}%
\pgfsetstrokecolor{currentstroke}%
\pgfsetstrokeopacity{0.000000}%
\pgfsetdash{}{0pt}%
\pgfpathmoveto{\pgfqpoint{7.679400in}{0.660000in}}%
\pgfpathlineto{\pgfqpoint{7.685600in}{0.660000in}}%
\pgfpathlineto{\pgfqpoint{7.685600in}{3.206604in}}%
\pgfpathlineto{\pgfqpoint{7.679400in}{3.206604in}}%
\pgfpathlineto{\pgfqpoint{7.679400in}{0.660000in}}%
\pgfpathclose%
\pgfusepath{fill}%
\end{pgfscope}%
\begin{pgfscope}%
\pgfpathrectangle{\pgfqpoint{1.250000in}{0.660000in}}{\pgfqpoint{7.750000in}{4.620000in}}%
\pgfusepath{clip}%
\pgfsetbuttcap%
\pgfsetmiterjoin%
\definecolor{currentfill}{rgb}{0.000000,0.000000,0.000000}%
\pgfsetfillcolor{currentfill}%
\pgfsetlinewidth{0.000000pt}%
\definecolor{currentstroke}{rgb}{0.000000,0.000000,0.000000}%
\pgfsetstrokecolor{currentstroke}%
\pgfsetstrokeopacity{0.000000}%
\pgfsetdash{}{0pt}%
\pgfpathmoveto{\pgfqpoint{7.687150in}{0.660000in}}%
\pgfpathlineto{\pgfqpoint{7.693350in}{0.660000in}}%
\pgfpathlineto{\pgfqpoint{7.693350in}{3.206487in}}%
\pgfpathlineto{\pgfqpoint{7.687150in}{3.206487in}}%
\pgfpathlineto{\pgfqpoint{7.687150in}{0.660000in}}%
\pgfpathclose%
\pgfusepath{fill}%
\end{pgfscope}%
\begin{pgfscope}%
\pgfpathrectangle{\pgfqpoint{1.250000in}{0.660000in}}{\pgfqpoint{7.750000in}{4.620000in}}%
\pgfusepath{clip}%
\pgfsetbuttcap%
\pgfsetmiterjoin%
\definecolor{currentfill}{rgb}{0.000000,0.000000,0.000000}%
\pgfsetfillcolor{currentfill}%
\pgfsetlinewidth{0.000000pt}%
\definecolor{currentstroke}{rgb}{0.000000,0.000000,0.000000}%
\pgfsetstrokecolor{currentstroke}%
\pgfsetstrokeopacity{0.000000}%
\pgfsetdash{}{0pt}%
\pgfpathmoveto{\pgfqpoint{7.694900in}{0.660000in}}%
\pgfpathlineto{\pgfqpoint{7.701100in}{0.660000in}}%
\pgfpathlineto{\pgfqpoint{7.701100in}{3.206356in}}%
\pgfpathlineto{\pgfqpoint{7.694900in}{3.206356in}}%
\pgfpathlineto{\pgfqpoint{7.694900in}{0.660000in}}%
\pgfpathclose%
\pgfusepath{fill}%
\end{pgfscope}%
\begin{pgfscope}%
\pgfpathrectangle{\pgfqpoint{1.250000in}{0.660000in}}{\pgfqpoint{7.750000in}{4.620000in}}%
\pgfusepath{clip}%
\pgfsetbuttcap%
\pgfsetmiterjoin%
\definecolor{currentfill}{rgb}{0.000000,0.000000,0.000000}%
\pgfsetfillcolor{currentfill}%
\pgfsetlinewidth{0.000000pt}%
\definecolor{currentstroke}{rgb}{0.000000,0.000000,0.000000}%
\pgfsetstrokecolor{currentstroke}%
\pgfsetstrokeopacity{0.000000}%
\pgfsetdash{}{0pt}%
\pgfpathmoveto{\pgfqpoint{7.702650in}{0.660000in}}%
\pgfpathlineto{\pgfqpoint{7.708850in}{0.660000in}}%
\pgfpathlineto{\pgfqpoint{7.708850in}{3.206209in}}%
\pgfpathlineto{\pgfqpoint{7.702650in}{3.206209in}}%
\pgfpathlineto{\pgfqpoint{7.702650in}{0.660000in}}%
\pgfpathclose%
\pgfusepath{fill}%
\end{pgfscope}%
\begin{pgfscope}%
\pgfpathrectangle{\pgfqpoint{1.250000in}{0.660000in}}{\pgfqpoint{7.750000in}{4.620000in}}%
\pgfusepath{clip}%
\pgfsetbuttcap%
\pgfsetmiterjoin%
\definecolor{currentfill}{rgb}{0.000000,0.000000,0.000000}%
\pgfsetfillcolor{currentfill}%
\pgfsetlinewidth{0.000000pt}%
\definecolor{currentstroke}{rgb}{0.000000,0.000000,0.000000}%
\pgfsetstrokecolor{currentstroke}%
\pgfsetstrokeopacity{0.000000}%
\pgfsetdash{}{0pt}%
\pgfpathmoveto{\pgfqpoint{7.710400in}{0.660000in}}%
\pgfpathlineto{\pgfqpoint{7.716600in}{0.660000in}}%
\pgfpathlineto{\pgfqpoint{7.716600in}{3.209131in}}%
\pgfpathlineto{\pgfqpoint{7.710400in}{3.209131in}}%
\pgfpathlineto{\pgfqpoint{7.710400in}{0.660000in}}%
\pgfpathclose%
\pgfusepath{fill}%
\end{pgfscope}%
\begin{pgfscope}%
\pgfpathrectangle{\pgfqpoint{1.250000in}{0.660000in}}{\pgfqpoint{7.750000in}{4.620000in}}%
\pgfusepath{clip}%
\pgfsetbuttcap%
\pgfsetmiterjoin%
\definecolor{currentfill}{rgb}{0.000000,0.000000,0.000000}%
\pgfsetfillcolor{currentfill}%
\pgfsetlinewidth{0.000000pt}%
\definecolor{currentstroke}{rgb}{0.000000,0.000000,0.000000}%
\pgfsetstrokecolor{currentstroke}%
\pgfsetstrokeopacity{0.000000}%
\pgfsetdash{}{0pt}%
\pgfpathmoveto{\pgfqpoint{7.718150in}{0.660000in}}%
\pgfpathlineto{\pgfqpoint{7.724350in}{0.660000in}}%
\pgfpathlineto{\pgfqpoint{7.724350in}{3.205868in}}%
\pgfpathlineto{\pgfqpoint{7.718150in}{3.205868in}}%
\pgfpathlineto{\pgfqpoint{7.718150in}{0.660000in}}%
\pgfpathclose%
\pgfusepath{fill}%
\end{pgfscope}%
\begin{pgfscope}%
\pgfpathrectangle{\pgfqpoint{1.250000in}{0.660000in}}{\pgfqpoint{7.750000in}{4.620000in}}%
\pgfusepath{clip}%
\pgfsetbuttcap%
\pgfsetmiterjoin%
\definecolor{currentfill}{rgb}{0.000000,0.000000,0.000000}%
\pgfsetfillcolor{currentfill}%
\pgfsetlinewidth{0.000000pt}%
\definecolor{currentstroke}{rgb}{0.000000,0.000000,0.000000}%
\pgfsetstrokecolor{currentstroke}%
\pgfsetstrokeopacity{0.000000}%
\pgfsetdash{}{0pt}%
\pgfpathmoveto{\pgfqpoint{7.725900in}{0.660000in}}%
\pgfpathlineto{\pgfqpoint{7.732100in}{0.660000in}}%
\pgfpathlineto{\pgfqpoint{7.732100in}{3.205674in}}%
\pgfpathlineto{\pgfqpoint{7.725900in}{3.205674in}}%
\pgfpathlineto{\pgfqpoint{7.725900in}{0.660000in}}%
\pgfpathclose%
\pgfusepath{fill}%
\end{pgfscope}%
\begin{pgfscope}%
\pgfpathrectangle{\pgfqpoint{1.250000in}{0.660000in}}{\pgfqpoint{7.750000in}{4.620000in}}%
\pgfusepath{clip}%
\pgfsetbuttcap%
\pgfsetmiterjoin%
\definecolor{currentfill}{rgb}{0.000000,0.000000,0.000000}%
\pgfsetfillcolor{currentfill}%
\pgfsetlinewidth{0.000000pt}%
\definecolor{currentstroke}{rgb}{0.000000,0.000000,0.000000}%
\pgfsetstrokecolor{currentstroke}%
\pgfsetstrokeopacity{0.000000}%
\pgfsetdash{}{0pt}%
\pgfpathmoveto{\pgfqpoint{7.733650in}{0.660000in}}%
\pgfpathlineto{\pgfqpoint{7.739850in}{0.660000in}}%
\pgfpathlineto{\pgfqpoint{7.739850in}{3.205463in}}%
\pgfpathlineto{\pgfqpoint{7.733650in}{3.205463in}}%
\pgfpathlineto{\pgfqpoint{7.733650in}{0.660000in}}%
\pgfpathclose%
\pgfusepath{fill}%
\end{pgfscope}%
\begin{pgfscope}%
\pgfpathrectangle{\pgfqpoint{1.250000in}{0.660000in}}{\pgfqpoint{7.750000in}{4.620000in}}%
\pgfusepath{clip}%
\pgfsetbuttcap%
\pgfsetmiterjoin%
\definecolor{currentfill}{rgb}{0.000000,0.000000,0.000000}%
\pgfsetfillcolor{currentfill}%
\pgfsetlinewidth{0.000000pt}%
\definecolor{currentstroke}{rgb}{0.000000,0.000000,0.000000}%
\pgfsetstrokecolor{currentstroke}%
\pgfsetstrokeopacity{0.000000}%
\pgfsetdash{}{0pt}%
\pgfpathmoveto{\pgfqpoint{7.741400in}{0.660000in}}%
\pgfpathlineto{\pgfqpoint{7.747600in}{0.660000in}}%
\pgfpathlineto{\pgfqpoint{7.747600in}{3.205410in}}%
\pgfpathlineto{\pgfqpoint{7.741400in}{3.205410in}}%
\pgfpathlineto{\pgfqpoint{7.741400in}{0.660000in}}%
\pgfpathclose%
\pgfusepath{fill}%
\end{pgfscope}%
\begin{pgfscope}%
\pgfpathrectangle{\pgfqpoint{1.250000in}{0.660000in}}{\pgfqpoint{7.750000in}{4.620000in}}%
\pgfusepath{clip}%
\pgfsetbuttcap%
\pgfsetmiterjoin%
\definecolor{currentfill}{rgb}{0.000000,0.000000,0.000000}%
\pgfsetfillcolor{currentfill}%
\pgfsetlinewidth{0.000000pt}%
\definecolor{currentstroke}{rgb}{0.000000,0.000000,0.000000}%
\pgfsetstrokecolor{currentstroke}%
\pgfsetstrokeopacity{0.000000}%
\pgfsetdash{}{0pt}%
\pgfpathmoveto{\pgfqpoint{7.749150in}{0.660000in}}%
\pgfpathlineto{\pgfqpoint{7.755350in}{0.660000in}}%
\pgfpathlineto{\pgfqpoint{7.755350in}{3.211903in}}%
\pgfpathlineto{\pgfqpoint{7.749150in}{3.211903in}}%
\pgfpathlineto{\pgfqpoint{7.749150in}{0.660000in}}%
\pgfpathclose%
\pgfusepath{fill}%
\end{pgfscope}%
\begin{pgfscope}%
\pgfpathrectangle{\pgfqpoint{1.250000in}{0.660000in}}{\pgfqpoint{7.750000in}{4.620000in}}%
\pgfusepath{clip}%
\pgfsetbuttcap%
\pgfsetmiterjoin%
\definecolor{currentfill}{rgb}{0.000000,0.000000,0.000000}%
\pgfsetfillcolor{currentfill}%
\pgfsetlinewidth{0.000000pt}%
\definecolor{currentstroke}{rgb}{0.000000,0.000000,0.000000}%
\pgfsetstrokecolor{currentstroke}%
\pgfsetstrokeopacity{0.000000}%
\pgfsetdash{}{0pt}%
\pgfpathmoveto{\pgfqpoint{7.756900in}{0.660000in}}%
\pgfpathlineto{\pgfqpoint{7.763100in}{0.660000in}}%
\pgfpathlineto{\pgfqpoint{7.763100in}{3.205401in}}%
\pgfpathlineto{\pgfqpoint{7.756900in}{3.205401in}}%
\pgfpathlineto{\pgfqpoint{7.756900in}{0.660000in}}%
\pgfpathclose%
\pgfusepath{fill}%
\end{pgfscope}%
\begin{pgfscope}%
\pgfpathrectangle{\pgfqpoint{1.250000in}{0.660000in}}{\pgfqpoint{7.750000in}{4.620000in}}%
\pgfusepath{clip}%
\pgfsetbuttcap%
\pgfsetmiterjoin%
\definecolor{currentfill}{rgb}{0.000000,0.000000,0.000000}%
\pgfsetfillcolor{currentfill}%
\pgfsetlinewidth{0.000000pt}%
\definecolor{currentstroke}{rgb}{0.000000,0.000000,0.000000}%
\pgfsetstrokecolor{currentstroke}%
\pgfsetstrokeopacity{0.000000}%
\pgfsetdash{}{0pt}%
\pgfpathmoveto{\pgfqpoint{7.764650in}{0.660000in}}%
\pgfpathlineto{\pgfqpoint{7.770850in}{0.660000in}}%
\pgfpathlineto{\pgfqpoint{7.770850in}{3.205374in}}%
\pgfpathlineto{\pgfqpoint{7.764650in}{3.205374in}}%
\pgfpathlineto{\pgfqpoint{7.764650in}{0.660000in}}%
\pgfpathclose%
\pgfusepath{fill}%
\end{pgfscope}%
\begin{pgfscope}%
\pgfpathrectangle{\pgfqpoint{1.250000in}{0.660000in}}{\pgfqpoint{7.750000in}{4.620000in}}%
\pgfusepath{clip}%
\pgfsetbuttcap%
\pgfsetmiterjoin%
\definecolor{currentfill}{rgb}{0.000000,0.000000,0.000000}%
\pgfsetfillcolor{currentfill}%
\pgfsetlinewidth{0.000000pt}%
\definecolor{currentstroke}{rgb}{0.000000,0.000000,0.000000}%
\pgfsetstrokecolor{currentstroke}%
\pgfsetstrokeopacity{0.000000}%
\pgfsetdash{}{0pt}%
\pgfpathmoveto{\pgfqpoint{7.772400in}{0.660000in}}%
\pgfpathlineto{\pgfqpoint{7.778600in}{0.660000in}}%
\pgfpathlineto{\pgfqpoint{7.778600in}{3.205332in}}%
\pgfpathlineto{\pgfqpoint{7.772400in}{3.205332in}}%
\pgfpathlineto{\pgfqpoint{7.772400in}{0.660000in}}%
\pgfpathclose%
\pgfusepath{fill}%
\end{pgfscope}%
\begin{pgfscope}%
\pgfpathrectangle{\pgfqpoint{1.250000in}{0.660000in}}{\pgfqpoint{7.750000in}{4.620000in}}%
\pgfusepath{clip}%
\pgfsetbuttcap%
\pgfsetmiterjoin%
\definecolor{currentfill}{rgb}{0.000000,0.000000,0.000000}%
\pgfsetfillcolor{currentfill}%
\pgfsetlinewidth{0.000000pt}%
\definecolor{currentstroke}{rgb}{0.000000,0.000000,0.000000}%
\pgfsetstrokecolor{currentstroke}%
\pgfsetstrokeopacity{0.000000}%
\pgfsetdash{}{0pt}%
\pgfpathmoveto{\pgfqpoint{7.780150in}{0.660000in}}%
\pgfpathlineto{\pgfqpoint{7.786350in}{0.660000in}}%
\pgfpathlineto{\pgfqpoint{7.786350in}{3.210649in}}%
\pgfpathlineto{\pgfqpoint{7.780150in}{3.210649in}}%
\pgfpathlineto{\pgfqpoint{7.780150in}{0.660000in}}%
\pgfpathclose%
\pgfusepath{fill}%
\end{pgfscope}%
\begin{pgfscope}%
\pgfpathrectangle{\pgfqpoint{1.250000in}{0.660000in}}{\pgfqpoint{7.750000in}{4.620000in}}%
\pgfusepath{clip}%
\pgfsetbuttcap%
\pgfsetmiterjoin%
\definecolor{currentfill}{rgb}{0.000000,0.000000,0.000000}%
\pgfsetfillcolor{currentfill}%
\pgfsetlinewidth{0.000000pt}%
\definecolor{currentstroke}{rgb}{0.000000,0.000000,0.000000}%
\pgfsetstrokecolor{currentstroke}%
\pgfsetstrokeopacity{0.000000}%
\pgfsetdash{}{0pt}%
\pgfpathmoveto{\pgfqpoint{7.787900in}{0.660000in}}%
\pgfpathlineto{\pgfqpoint{7.794100in}{0.660000in}}%
\pgfpathlineto{\pgfqpoint{7.794100in}{3.205200in}}%
\pgfpathlineto{\pgfqpoint{7.787900in}{3.205200in}}%
\pgfpathlineto{\pgfqpoint{7.787900in}{0.660000in}}%
\pgfpathclose%
\pgfusepath{fill}%
\end{pgfscope}%
\begin{pgfscope}%
\pgfpathrectangle{\pgfqpoint{1.250000in}{0.660000in}}{\pgfqpoint{7.750000in}{4.620000in}}%
\pgfusepath{clip}%
\pgfsetbuttcap%
\pgfsetmiterjoin%
\definecolor{currentfill}{rgb}{0.000000,0.000000,0.000000}%
\pgfsetfillcolor{currentfill}%
\pgfsetlinewidth{0.000000pt}%
\definecolor{currentstroke}{rgb}{0.000000,0.000000,0.000000}%
\pgfsetstrokecolor{currentstroke}%
\pgfsetstrokeopacity{0.000000}%
\pgfsetdash{}{0pt}%
\pgfpathmoveto{\pgfqpoint{7.795650in}{0.660000in}}%
\pgfpathlineto{\pgfqpoint{7.801850in}{0.660000in}}%
\pgfpathlineto{\pgfqpoint{7.801850in}{3.205110in}}%
\pgfpathlineto{\pgfqpoint{7.795650in}{3.205110in}}%
\pgfpathlineto{\pgfqpoint{7.795650in}{0.660000in}}%
\pgfpathclose%
\pgfusepath{fill}%
\end{pgfscope}%
\begin{pgfscope}%
\pgfpathrectangle{\pgfqpoint{1.250000in}{0.660000in}}{\pgfqpoint{7.750000in}{4.620000in}}%
\pgfusepath{clip}%
\pgfsetbuttcap%
\pgfsetmiterjoin%
\definecolor{currentfill}{rgb}{0.000000,0.000000,0.000000}%
\pgfsetfillcolor{currentfill}%
\pgfsetlinewidth{0.000000pt}%
\definecolor{currentstroke}{rgb}{0.000000,0.000000,0.000000}%
\pgfsetstrokecolor{currentstroke}%
\pgfsetstrokeopacity{0.000000}%
\pgfsetdash{}{0pt}%
\pgfpathmoveto{\pgfqpoint{7.803400in}{0.660000in}}%
\pgfpathlineto{\pgfqpoint{7.809600in}{0.660000in}}%
\pgfpathlineto{\pgfqpoint{7.809600in}{3.205003in}}%
\pgfpathlineto{\pgfqpoint{7.803400in}{3.205003in}}%
\pgfpathlineto{\pgfqpoint{7.803400in}{0.660000in}}%
\pgfpathclose%
\pgfusepath{fill}%
\end{pgfscope}%
\begin{pgfscope}%
\pgfpathrectangle{\pgfqpoint{1.250000in}{0.660000in}}{\pgfqpoint{7.750000in}{4.620000in}}%
\pgfusepath{clip}%
\pgfsetbuttcap%
\pgfsetmiterjoin%
\definecolor{currentfill}{rgb}{0.000000,0.000000,0.000000}%
\pgfsetfillcolor{currentfill}%
\pgfsetlinewidth{0.000000pt}%
\definecolor{currentstroke}{rgb}{0.000000,0.000000,0.000000}%
\pgfsetstrokecolor{currentstroke}%
\pgfsetstrokeopacity{0.000000}%
\pgfsetdash{}{0pt}%
\pgfpathmoveto{\pgfqpoint{7.811150in}{0.660000in}}%
\pgfpathlineto{\pgfqpoint{7.817350in}{0.660000in}}%
\pgfpathlineto{\pgfqpoint{7.817350in}{3.204879in}}%
\pgfpathlineto{\pgfqpoint{7.811150in}{3.204879in}}%
\pgfpathlineto{\pgfqpoint{7.811150in}{0.660000in}}%
\pgfpathclose%
\pgfusepath{fill}%
\end{pgfscope}%
\begin{pgfscope}%
\pgfpathrectangle{\pgfqpoint{1.250000in}{0.660000in}}{\pgfqpoint{7.750000in}{4.620000in}}%
\pgfusepath{clip}%
\pgfsetbuttcap%
\pgfsetmiterjoin%
\definecolor{currentfill}{rgb}{0.000000,0.000000,0.000000}%
\pgfsetfillcolor{currentfill}%
\pgfsetlinewidth{0.000000pt}%
\definecolor{currentstroke}{rgb}{0.000000,0.000000,0.000000}%
\pgfsetstrokecolor{currentstroke}%
\pgfsetstrokeopacity{0.000000}%
\pgfsetdash{}{0pt}%
\pgfpathmoveto{\pgfqpoint{7.818900in}{0.660000in}}%
\pgfpathlineto{\pgfqpoint{7.825100in}{0.660000in}}%
\pgfpathlineto{\pgfqpoint{7.825100in}{3.213208in}}%
\pgfpathlineto{\pgfqpoint{7.818900in}{3.213208in}}%
\pgfpathlineto{\pgfqpoint{7.818900in}{0.660000in}}%
\pgfpathclose%
\pgfusepath{fill}%
\end{pgfscope}%
\begin{pgfscope}%
\pgfpathrectangle{\pgfqpoint{1.250000in}{0.660000in}}{\pgfqpoint{7.750000in}{4.620000in}}%
\pgfusepath{clip}%
\pgfsetbuttcap%
\pgfsetmiterjoin%
\definecolor{currentfill}{rgb}{0.000000,0.000000,0.000000}%
\pgfsetfillcolor{currentfill}%
\pgfsetlinewidth{0.000000pt}%
\definecolor{currentstroke}{rgb}{0.000000,0.000000,0.000000}%
\pgfsetstrokecolor{currentstroke}%
\pgfsetstrokeopacity{0.000000}%
\pgfsetdash{}{0pt}%
\pgfpathmoveto{\pgfqpoint{7.826650in}{0.660000in}}%
\pgfpathlineto{\pgfqpoint{7.832850in}{0.660000in}}%
\pgfpathlineto{\pgfqpoint{7.832850in}{3.213545in}}%
\pgfpathlineto{\pgfqpoint{7.826650in}{3.213545in}}%
\pgfpathlineto{\pgfqpoint{7.826650in}{0.660000in}}%
\pgfpathclose%
\pgfusepath{fill}%
\end{pgfscope}%
\begin{pgfscope}%
\pgfpathrectangle{\pgfqpoint{1.250000in}{0.660000in}}{\pgfqpoint{7.750000in}{4.620000in}}%
\pgfusepath{clip}%
\pgfsetbuttcap%
\pgfsetmiterjoin%
\definecolor{currentfill}{rgb}{0.000000,0.000000,0.000000}%
\pgfsetfillcolor{currentfill}%
\pgfsetlinewidth{0.000000pt}%
\definecolor{currentstroke}{rgb}{0.000000,0.000000,0.000000}%
\pgfsetstrokecolor{currentstroke}%
\pgfsetstrokeopacity{0.000000}%
\pgfsetdash{}{0pt}%
\pgfpathmoveto{\pgfqpoint{7.834400in}{0.660000in}}%
\pgfpathlineto{\pgfqpoint{7.840600in}{0.660000in}}%
\pgfpathlineto{\pgfqpoint{7.840600in}{3.204402in}}%
\pgfpathlineto{\pgfqpoint{7.834400in}{3.204402in}}%
\pgfpathlineto{\pgfqpoint{7.834400in}{0.660000in}}%
\pgfpathclose%
\pgfusepath{fill}%
\end{pgfscope}%
\begin{pgfscope}%
\pgfpathrectangle{\pgfqpoint{1.250000in}{0.660000in}}{\pgfqpoint{7.750000in}{4.620000in}}%
\pgfusepath{clip}%
\pgfsetbuttcap%
\pgfsetmiterjoin%
\definecolor{currentfill}{rgb}{0.000000,0.000000,0.000000}%
\pgfsetfillcolor{currentfill}%
\pgfsetlinewidth{0.000000pt}%
\definecolor{currentstroke}{rgb}{0.000000,0.000000,0.000000}%
\pgfsetstrokecolor{currentstroke}%
\pgfsetstrokeopacity{0.000000}%
\pgfsetdash{}{0pt}%
\pgfpathmoveto{\pgfqpoint{7.842150in}{0.660000in}}%
\pgfpathlineto{\pgfqpoint{7.848350in}{0.660000in}}%
\pgfpathlineto{\pgfqpoint{7.848350in}{3.204207in}}%
\pgfpathlineto{\pgfqpoint{7.842150in}{3.204207in}}%
\pgfpathlineto{\pgfqpoint{7.842150in}{0.660000in}}%
\pgfpathclose%
\pgfusepath{fill}%
\end{pgfscope}%
\begin{pgfscope}%
\pgfpathrectangle{\pgfqpoint{1.250000in}{0.660000in}}{\pgfqpoint{7.750000in}{4.620000in}}%
\pgfusepath{clip}%
\pgfsetbuttcap%
\pgfsetmiterjoin%
\definecolor{currentfill}{rgb}{0.000000,0.000000,0.000000}%
\pgfsetfillcolor{currentfill}%
\pgfsetlinewidth{0.000000pt}%
\definecolor{currentstroke}{rgb}{0.000000,0.000000,0.000000}%
\pgfsetstrokecolor{currentstroke}%
\pgfsetstrokeopacity{0.000000}%
\pgfsetdash{}{0pt}%
\pgfpathmoveto{\pgfqpoint{7.849900in}{0.660000in}}%
\pgfpathlineto{\pgfqpoint{7.856100in}{0.660000in}}%
\pgfpathlineto{\pgfqpoint{7.856100in}{3.203993in}}%
\pgfpathlineto{\pgfqpoint{7.849900in}{3.203993in}}%
\pgfpathlineto{\pgfqpoint{7.849900in}{0.660000in}}%
\pgfpathclose%
\pgfusepath{fill}%
\end{pgfscope}%
\begin{pgfscope}%
\pgfpathrectangle{\pgfqpoint{1.250000in}{0.660000in}}{\pgfqpoint{7.750000in}{4.620000in}}%
\pgfusepath{clip}%
\pgfsetbuttcap%
\pgfsetmiterjoin%
\definecolor{currentfill}{rgb}{0.000000,0.000000,0.000000}%
\pgfsetfillcolor{currentfill}%
\pgfsetlinewidth{0.000000pt}%
\definecolor{currentstroke}{rgb}{0.000000,0.000000,0.000000}%
\pgfsetstrokecolor{currentstroke}%
\pgfsetstrokeopacity{0.000000}%
\pgfsetdash{}{0pt}%
\pgfpathmoveto{\pgfqpoint{7.857650in}{0.660000in}}%
\pgfpathlineto{\pgfqpoint{7.863850in}{0.660000in}}%
\pgfpathlineto{\pgfqpoint{7.863850in}{3.203942in}}%
\pgfpathlineto{\pgfqpoint{7.857650in}{3.203942in}}%
\pgfpathlineto{\pgfqpoint{7.857650in}{0.660000in}}%
\pgfpathclose%
\pgfusepath{fill}%
\end{pgfscope}%
\begin{pgfscope}%
\pgfpathrectangle{\pgfqpoint{1.250000in}{0.660000in}}{\pgfqpoint{7.750000in}{4.620000in}}%
\pgfusepath{clip}%
\pgfsetbuttcap%
\pgfsetmiterjoin%
\definecolor{currentfill}{rgb}{0.000000,0.000000,0.000000}%
\pgfsetfillcolor{currentfill}%
\pgfsetlinewidth{0.000000pt}%
\definecolor{currentstroke}{rgb}{0.000000,0.000000,0.000000}%
\pgfsetstrokecolor{currentstroke}%
\pgfsetstrokeopacity{0.000000}%
\pgfsetdash{}{0pt}%
\pgfpathmoveto{\pgfqpoint{7.865400in}{0.660000in}}%
\pgfpathlineto{\pgfqpoint{7.871600in}{0.660000in}}%
\pgfpathlineto{\pgfqpoint{7.871600in}{3.203952in}}%
\pgfpathlineto{\pgfqpoint{7.865400in}{3.203952in}}%
\pgfpathlineto{\pgfqpoint{7.865400in}{0.660000in}}%
\pgfpathclose%
\pgfusepath{fill}%
\end{pgfscope}%
\begin{pgfscope}%
\pgfpathrectangle{\pgfqpoint{1.250000in}{0.660000in}}{\pgfqpoint{7.750000in}{4.620000in}}%
\pgfusepath{clip}%
\pgfsetbuttcap%
\pgfsetmiterjoin%
\definecolor{currentfill}{rgb}{0.000000,0.000000,0.000000}%
\pgfsetfillcolor{currentfill}%
\pgfsetlinewidth{0.000000pt}%
\definecolor{currentstroke}{rgb}{0.000000,0.000000,0.000000}%
\pgfsetstrokecolor{currentstroke}%
\pgfsetstrokeopacity{0.000000}%
\pgfsetdash{}{0pt}%
\pgfpathmoveto{\pgfqpoint{7.873150in}{0.660000in}}%
\pgfpathlineto{\pgfqpoint{7.879350in}{0.660000in}}%
\pgfpathlineto{\pgfqpoint{7.879350in}{3.203945in}}%
\pgfpathlineto{\pgfqpoint{7.873150in}{3.203945in}}%
\pgfpathlineto{\pgfqpoint{7.873150in}{0.660000in}}%
\pgfpathclose%
\pgfusepath{fill}%
\end{pgfscope}%
\begin{pgfscope}%
\pgfpathrectangle{\pgfqpoint{1.250000in}{0.660000in}}{\pgfqpoint{7.750000in}{4.620000in}}%
\pgfusepath{clip}%
\pgfsetbuttcap%
\pgfsetmiterjoin%
\definecolor{currentfill}{rgb}{0.000000,0.000000,0.000000}%
\pgfsetfillcolor{currentfill}%
\pgfsetlinewidth{0.000000pt}%
\definecolor{currentstroke}{rgb}{0.000000,0.000000,0.000000}%
\pgfsetstrokecolor{currentstroke}%
\pgfsetstrokeopacity{0.000000}%
\pgfsetdash{}{0pt}%
\pgfpathmoveto{\pgfqpoint{7.880900in}{0.660000in}}%
\pgfpathlineto{\pgfqpoint{7.887100in}{0.660000in}}%
\pgfpathlineto{\pgfqpoint{7.887100in}{3.203921in}}%
\pgfpathlineto{\pgfqpoint{7.880900in}{3.203921in}}%
\pgfpathlineto{\pgfqpoint{7.880900in}{0.660000in}}%
\pgfpathclose%
\pgfusepath{fill}%
\end{pgfscope}%
\begin{pgfscope}%
\pgfpathrectangle{\pgfqpoint{1.250000in}{0.660000in}}{\pgfqpoint{7.750000in}{4.620000in}}%
\pgfusepath{clip}%
\pgfsetbuttcap%
\pgfsetmiterjoin%
\definecolor{currentfill}{rgb}{0.000000,0.000000,0.000000}%
\pgfsetfillcolor{currentfill}%
\pgfsetlinewidth{0.000000pt}%
\definecolor{currentstroke}{rgb}{0.000000,0.000000,0.000000}%
\pgfsetstrokecolor{currentstroke}%
\pgfsetstrokeopacity{0.000000}%
\pgfsetdash{}{0pt}%
\pgfpathmoveto{\pgfqpoint{7.888650in}{0.660000in}}%
\pgfpathlineto{\pgfqpoint{7.894850in}{0.660000in}}%
\pgfpathlineto{\pgfqpoint{7.894850in}{3.203880in}}%
\pgfpathlineto{\pgfqpoint{7.888650in}{3.203880in}}%
\pgfpathlineto{\pgfqpoint{7.888650in}{0.660000in}}%
\pgfpathclose%
\pgfusepath{fill}%
\end{pgfscope}%
\begin{pgfscope}%
\pgfpathrectangle{\pgfqpoint{1.250000in}{0.660000in}}{\pgfqpoint{7.750000in}{4.620000in}}%
\pgfusepath{clip}%
\pgfsetbuttcap%
\pgfsetmiterjoin%
\definecolor{currentfill}{rgb}{0.000000,0.000000,0.000000}%
\pgfsetfillcolor{currentfill}%
\pgfsetlinewidth{0.000000pt}%
\definecolor{currentstroke}{rgb}{0.000000,0.000000,0.000000}%
\pgfsetstrokecolor{currentstroke}%
\pgfsetstrokeopacity{0.000000}%
\pgfsetdash{}{0pt}%
\pgfpathmoveto{\pgfqpoint{7.896400in}{0.660000in}}%
\pgfpathlineto{\pgfqpoint{7.902600in}{0.660000in}}%
\pgfpathlineto{\pgfqpoint{7.902600in}{3.203820in}}%
\pgfpathlineto{\pgfqpoint{7.896400in}{3.203820in}}%
\pgfpathlineto{\pgfqpoint{7.896400in}{0.660000in}}%
\pgfpathclose%
\pgfusepath{fill}%
\end{pgfscope}%
\begin{pgfscope}%
\pgfpathrectangle{\pgfqpoint{1.250000in}{0.660000in}}{\pgfqpoint{7.750000in}{4.620000in}}%
\pgfusepath{clip}%
\pgfsetbuttcap%
\pgfsetmiterjoin%
\definecolor{currentfill}{rgb}{0.000000,0.000000,0.000000}%
\pgfsetfillcolor{currentfill}%
\pgfsetlinewidth{0.000000pt}%
\definecolor{currentstroke}{rgb}{0.000000,0.000000,0.000000}%
\pgfsetstrokecolor{currentstroke}%
\pgfsetstrokeopacity{0.000000}%
\pgfsetdash{}{0pt}%
\pgfpathmoveto{\pgfqpoint{7.904150in}{0.660000in}}%
\pgfpathlineto{\pgfqpoint{7.910350in}{0.660000in}}%
\pgfpathlineto{\pgfqpoint{7.910350in}{3.203743in}}%
\pgfpathlineto{\pgfqpoint{7.904150in}{3.203743in}}%
\pgfpathlineto{\pgfqpoint{7.904150in}{0.660000in}}%
\pgfpathclose%
\pgfusepath{fill}%
\end{pgfscope}%
\begin{pgfscope}%
\pgfpathrectangle{\pgfqpoint{1.250000in}{0.660000in}}{\pgfqpoint{7.750000in}{4.620000in}}%
\pgfusepath{clip}%
\pgfsetbuttcap%
\pgfsetmiterjoin%
\definecolor{currentfill}{rgb}{0.000000,0.000000,0.000000}%
\pgfsetfillcolor{currentfill}%
\pgfsetlinewidth{0.000000pt}%
\definecolor{currentstroke}{rgb}{0.000000,0.000000,0.000000}%
\pgfsetstrokecolor{currentstroke}%
\pgfsetstrokeopacity{0.000000}%
\pgfsetdash{}{0pt}%
\pgfpathmoveto{\pgfqpoint{7.911900in}{0.660000in}}%
\pgfpathlineto{\pgfqpoint{7.918100in}{0.660000in}}%
\pgfpathlineto{\pgfqpoint{7.918100in}{3.203647in}}%
\pgfpathlineto{\pgfqpoint{7.911900in}{3.203647in}}%
\pgfpathlineto{\pgfqpoint{7.911900in}{0.660000in}}%
\pgfpathclose%
\pgfusepath{fill}%
\end{pgfscope}%
\begin{pgfscope}%
\pgfpathrectangle{\pgfqpoint{1.250000in}{0.660000in}}{\pgfqpoint{7.750000in}{4.620000in}}%
\pgfusepath{clip}%
\pgfsetbuttcap%
\pgfsetmiterjoin%
\definecolor{currentfill}{rgb}{0.000000,0.000000,0.000000}%
\pgfsetfillcolor{currentfill}%
\pgfsetlinewidth{0.000000pt}%
\definecolor{currentstroke}{rgb}{0.000000,0.000000,0.000000}%
\pgfsetstrokecolor{currentstroke}%
\pgfsetstrokeopacity{0.000000}%
\pgfsetdash{}{0pt}%
\pgfpathmoveto{\pgfqpoint{7.919650in}{0.660000in}}%
\pgfpathlineto{\pgfqpoint{7.925850in}{0.660000in}}%
\pgfpathlineto{\pgfqpoint{7.925850in}{3.203532in}}%
\pgfpathlineto{\pgfqpoint{7.919650in}{3.203532in}}%
\pgfpathlineto{\pgfqpoint{7.919650in}{0.660000in}}%
\pgfpathclose%
\pgfusepath{fill}%
\end{pgfscope}%
\begin{pgfscope}%
\pgfpathrectangle{\pgfqpoint{1.250000in}{0.660000in}}{\pgfqpoint{7.750000in}{4.620000in}}%
\pgfusepath{clip}%
\pgfsetbuttcap%
\pgfsetmiterjoin%
\definecolor{currentfill}{rgb}{0.000000,0.000000,0.000000}%
\pgfsetfillcolor{currentfill}%
\pgfsetlinewidth{0.000000pt}%
\definecolor{currentstroke}{rgb}{0.000000,0.000000,0.000000}%
\pgfsetstrokecolor{currentstroke}%
\pgfsetstrokeopacity{0.000000}%
\pgfsetdash{}{0pt}%
\pgfpathmoveto{\pgfqpoint{7.927400in}{0.660000in}}%
\pgfpathlineto{\pgfqpoint{7.933600in}{0.660000in}}%
\pgfpathlineto{\pgfqpoint{7.933600in}{3.203398in}}%
\pgfpathlineto{\pgfqpoint{7.927400in}{3.203398in}}%
\pgfpathlineto{\pgfqpoint{7.927400in}{0.660000in}}%
\pgfpathclose%
\pgfusepath{fill}%
\end{pgfscope}%
\begin{pgfscope}%
\pgfpathrectangle{\pgfqpoint{1.250000in}{0.660000in}}{\pgfqpoint{7.750000in}{4.620000in}}%
\pgfusepath{clip}%
\pgfsetbuttcap%
\pgfsetmiterjoin%
\definecolor{currentfill}{rgb}{0.000000,0.000000,0.000000}%
\pgfsetfillcolor{currentfill}%
\pgfsetlinewidth{0.000000pt}%
\definecolor{currentstroke}{rgb}{0.000000,0.000000,0.000000}%
\pgfsetstrokecolor{currentstroke}%
\pgfsetstrokeopacity{0.000000}%
\pgfsetdash{}{0pt}%
\pgfpathmoveto{\pgfqpoint{7.935150in}{0.660000in}}%
\pgfpathlineto{\pgfqpoint{7.941350in}{0.660000in}}%
\pgfpathlineto{\pgfqpoint{7.941350in}{3.203244in}}%
\pgfpathlineto{\pgfqpoint{7.935150in}{3.203244in}}%
\pgfpathlineto{\pgfqpoint{7.935150in}{0.660000in}}%
\pgfpathclose%
\pgfusepath{fill}%
\end{pgfscope}%
\begin{pgfscope}%
\pgfpathrectangle{\pgfqpoint{1.250000in}{0.660000in}}{\pgfqpoint{7.750000in}{4.620000in}}%
\pgfusepath{clip}%
\pgfsetbuttcap%
\pgfsetmiterjoin%
\definecolor{currentfill}{rgb}{0.000000,0.000000,0.000000}%
\pgfsetfillcolor{currentfill}%
\pgfsetlinewidth{0.000000pt}%
\definecolor{currentstroke}{rgb}{0.000000,0.000000,0.000000}%
\pgfsetstrokecolor{currentstroke}%
\pgfsetstrokeopacity{0.000000}%
\pgfsetdash{}{0pt}%
\pgfpathmoveto{\pgfqpoint{7.942900in}{0.660000in}}%
\pgfpathlineto{\pgfqpoint{7.949100in}{0.660000in}}%
\pgfpathlineto{\pgfqpoint{7.949100in}{3.215927in}}%
\pgfpathlineto{\pgfqpoint{7.942900in}{3.215927in}}%
\pgfpathlineto{\pgfqpoint{7.942900in}{0.660000in}}%
\pgfpathclose%
\pgfusepath{fill}%
\end{pgfscope}%
\begin{pgfscope}%
\pgfpathrectangle{\pgfqpoint{1.250000in}{0.660000in}}{\pgfqpoint{7.750000in}{4.620000in}}%
\pgfusepath{clip}%
\pgfsetbuttcap%
\pgfsetmiterjoin%
\definecolor{currentfill}{rgb}{0.000000,0.000000,0.000000}%
\pgfsetfillcolor{currentfill}%
\pgfsetlinewidth{0.000000pt}%
\definecolor{currentstroke}{rgb}{0.000000,0.000000,0.000000}%
\pgfsetstrokecolor{currentstroke}%
\pgfsetstrokeopacity{0.000000}%
\pgfsetdash{}{0pt}%
\pgfpathmoveto{\pgfqpoint{7.950650in}{0.660000in}}%
\pgfpathlineto{\pgfqpoint{7.956850in}{0.660000in}}%
\pgfpathlineto{\pgfqpoint{7.956850in}{3.202877in}}%
\pgfpathlineto{\pgfqpoint{7.950650in}{3.202877in}}%
\pgfpathlineto{\pgfqpoint{7.950650in}{0.660000in}}%
\pgfpathclose%
\pgfusepath{fill}%
\end{pgfscope}%
\begin{pgfscope}%
\pgfpathrectangle{\pgfqpoint{1.250000in}{0.660000in}}{\pgfqpoint{7.750000in}{4.620000in}}%
\pgfusepath{clip}%
\pgfsetbuttcap%
\pgfsetmiterjoin%
\definecolor{currentfill}{rgb}{0.000000,0.000000,0.000000}%
\pgfsetfillcolor{currentfill}%
\pgfsetlinewidth{0.000000pt}%
\definecolor{currentstroke}{rgb}{0.000000,0.000000,0.000000}%
\pgfsetstrokecolor{currentstroke}%
\pgfsetstrokeopacity{0.000000}%
\pgfsetdash{}{0pt}%
\pgfpathmoveto{\pgfqpoint{7.958400in}{0.660000in}}%
\pgfpathlineto{\pgfqpoint{7.964600in}{0.660000in}}%
\pgfpathlineto{\pgfqpoint{7.964600in}{3.202662in}}%
\pgfpathlineto{\pgfqpoint{7.958400in}{3.202662in}}%
\pgfpathlineto{\pgfqpoint{7.958400in}{0.660000in}}%
\pgfpathclose%
\pgfusepath{fill}%
\end{pgfscope}%
\begin{pgfscope}%
\pgfpathrectangle{\pgfqpoint{1.250000in}{0.660000in}}{\pgfqpoint{7.750000in}{4.620000in}}%
\pgfusepath{clip}%
\pgfsetbuttcap%
\pgfsetmiterjoin%
\definecolor{currentfill}{rgb}{0.000000,0.000000,0.000000}%
\pgfsetfillcolor{currentfill}%
\pgfsetlinewidth{0.000000pt}%
\definecolor{currentstroke}{rgb}{0.000000,0.000000,0.000000}%
\pgfsetstrokecolor{currentstroke}%
\pgfsetstrokeopacity{0.000000}%
\pgfsetdash{}{0pt}%
\pgfpathmoveto{\pgfqpoint{7.966150in}{0.660000in}}%
\pgfpathlineto{\pgfqpoint{7.972350in}{0.660000in}}%
\pgfpathlineto{\pgfqpoint{7.972350in}{3.202595in}}%
\pgfpathlineto{\pgfqpoint{7.966150in}{3.202595in}}%
\pgfpathlineto{\pgfqpoint{7.966150in}{0.660000in}}%
\pgfpathclose%
\pgfusepath{fill}%
\end{pgfscope}%
\begin{pgfscope}%
\pgfpathrectangle{\pgfqpoint{1.250000in}{0.660000in}}{\pgfqpoint{7.750000in}{4.620000in}}%
\pgfusepath{clip}%
\pgfsetbuttcap%
\pgfsetmiterjoin%
\definecolor{currentfill}{rgb}{0.000000,0.000000,0.000000}%
\pgfsetfillcolor{currentfill}%
\pgfsetlinewidth{0.000000pt}%
\definecolor{currentstroke}{rgb}{0.000000,0.000000,0.000000}%
\pgfsetstrokecolor{currentstroke}%
\pgfsetstrokeopacity{0.000000}%
\pgfsetdash{}{0pt}%
\pgfpathmoveto{\pgfqpoint{7.973900in}{0.660000in}}%
\pgfpathlineto{\pgfqpoint{7.980100in}{0.660000in}}%
\pgfpathlineto{\pgfqpoint{7.980100in}{3.202612in}}%
\pgfpathlineto{\pgfqpoint{7.973900in}{3.202612in}}%
\pgfpathlineto{\pgfqpoint{7.973900in}{0.660000in}}%
\pgfpathclose%
\pgfusepath{fill}%
\end{pgfscope}%
\begin{pgfscope}%
\pgfpathrectangle{\pgfqpoint{1.250000in}{0.660000in}}{\pgfqpoint{7.750000in}{4.620000in}}%
\pgfusepath{clip}%
\pgfsetbuttcap%
\pgfsetmiterjoin%
\definecolor{currentfill}{rgb}{0.000000,0.000000,0.000000}%
\pgfsetfillcolor{currentfill}%
\pgfsetlinewidth{0.000000pt}%
\definecolor{currentstroke}{rgb}{0.000000,0.000000,0.000000}%
\pgfsetstrokecolor{currentstroke}%
\pgfsetstrokeopacity{0.000000}%
\pgfsetdash{}{0pt}%
\pgfpathmoveto{\pgfqpoint{7.981650in}{0.660000in}}%
\pgfpathlineto{\pgfqpoint{7.987850in}{0.660000in}}%
\pgfpathlineto{\pgfqpoint{7.987850in}{3.202611in}}%
\pgfpathlineto{\pgfqpoint{7.981650in}{3.202611in}}%
\pgfpathlineto{\pgfqpoint{7.981650in}{0.660000in}}%
\pgfpathclose%
\pgfusepath{fill}%
\end{pgfscope}%
\begin{pgfscope}%
\pgfpathrectangle{\pgfqpoint{1.250000in}{0.660000in}}{\pgfqpoint{7.750000in}{4.620000in}}%
\pgfusepath{clip}%
\pgfsetbuttcap%
\pgfsetmiterjoin%
\definecolor{currentfill}{rgb}{0.000000,0.000000,0.000000}%
\pgfsetfillcolor{currentfill}%
\pgfsetlinewidth{0.000000pt}%
\definecolor{currentstroke}{rgb}{0.000000,0.000000,0.000000}%
\pgfsetstrokecolor{currentstroke}%
\pgfsetstrokeopacity{0.000000}%
\pgfsetdash{}{0pt}%
\pgfpathmoveto{\pgfqpoint{7.989400in}{0.660000in}}%
\pgfpathlineto{\pgfqpoint{7.995600in}{0.660000in}}%
\pgfpathlineto{\pgfqpoint{7.995600in}{3.202591in}}%
\pgfpathlineto{\pgfqpoint{7.989400in}{3.202591in}}%
\pgfpathlineto{\pgfqpoint{7.989400in}{0.660000in}}%
\pgfpathclose%
\pgfusepath{fill}%
\end{pgfscope}%
\begin{pgfscope}%
\pgfpathrectangle{\pgfqpoint{1.250000in}{0.660000in}}{\pgfqpoint{7.750000in}{4.620000in}}%
\pgfusepath{clip}%
\pgfsetbuttcap%
\pgfsetmiterjoin%
\definecolor{currentfill}{rgb}{0.000000,0.000000,0.000000}%
\pgfsetfillcolor{currentfill}%
\pgfsetlinewidth{0.000000pt}%
\definecolor{currentstroke}{rgb}{0.000000,0.000000,0.000000}%
\pgfsetstrokecolor{currentstroke}%
\pgfsetstrokeopacity{0.000000}%
\pgfsetdash{}{0pt}%
\pgfpathmoveto{\pgfqpoint{7.997150in}{0.660000in}}%
\pgfpathlineto{\pgfqpoint{8.003350in}{0.660000in}}%
\pgfpathlineto{\pgfqpoint{8.003350in}{3.202551in}}%
\pgfpathlineto{\pgfqpoint{7.997150in}{3.202551in}}%
\pgfpathlineto{\pgfqpoint{7.997150in}{0.660000in}}%
\pgfpathclose%
\pgfusepath{fill}%
\end{pgfscope}%
\begin{pgfscope}%
\pgfpathrectangle{\pgfqpoint{1.250000in}{0.660000in}}{\pgfqpoint{7.750000in}{4.620000in}}%
\pgfusepath{clip}%
\pgfsetbuttcap%
\pgfsetmiterjoin%
\definecolor{currentfill}{rgb}{0.000000,0.000000,0.000000}%
\pgfsetfillcolor{currentfill}%
\pgfsetlinewidth{0.000000pt}%
\definecolor{currentstroke}{rgb}{0.000000,0.000000,0.000000}%
\pgfsetstrokecolor{currentstroke}%
\pgfsetstrokeopacity{0.000000}%
\pgfsetdash{}{0pt}%
\pgfpathmoveto{\pgfqpoint{8.004900in}{0.660000in}}%
\pgfpathlineto{\pgfqpoint{8.011100in}{0.660000in}}%
\pgfpathlineto{\pgfqpoint{8.011100in}{3.202491in}}%
\pgfpathlineto{\pgfqpoint{8.004900in}{3.202491in}}%
\pgfpathlineto{\pgfqpoint{8.004900in}{0.660000in}}%
\pgfpathclose%
\pgfusepath{fill}%
\end{pgfscope}%
\begin{pgfscope}%
\pgfpathrectangle{\pgfqpoint{1.250000in}{0.660000in}}{\pgfqpoint{7.750000in}{4.620000in}}%
\pgfusepath{clip}%
\pgfsetbuttcap%
\pgfsetmiterjoin%
\definecolor{currentfill}{rgb}{0.000000,0.000000,0.000000}%
\pgfsetfillcolor{currentfill}%
\pgfsetlinewidth{0.000000pt}%
\definecolor{currentstroke}{rgb}{0.000000,0.000000,0.000000}%
\pgfsetstrokecolor{currentstroke}%
\pgfsetstrokeopacity{0.000000}%
\pgfsetdash{}{0pt}%
\pgfpathmoveto{\pgfqpoint{8.012650in}{0.660000in}}%
\pgfpathlineto{\pgfqpoint{8.018850in}{0.660000in}}%
\pgfpathlineto{\pgfqpoint{8.018850in}{3.202411in}}%
\pgfpathlineto{\pgfqpoint{8.012650in}{3.202411in}}%
\pgfpathlineto{\pgfqpoint{8.012650in}{0.660000in}}%
\pgfpathclose%
\pgfusepath{fill}%
\end{pgfscope}%
\begin{pgfscope}%
\pgfpathrectangle{\pgfqpoint{1.250000in}{0.660000in}}{\pgfqpoint{7.750000in}{4.620000in}}%
\pgfusepath{clip}%
\pgfsetbuttcap%
\pgfsetmiterjoin%
\definecolor{currentfill}{rgb}{0.000000,0.000000,0.000000}%
\pgfsetfillcolor{currentfill}%
\pgfsetlinewidth{0.000000pt}%
\definecolor{currentstroke}{rgb}{0.000000,0.000000,0.000000}%
\pgfsetstrokecolor{currentstroke}%
\pgfsetstrokeopacity{0.000000}%
\pgfsetdash{}{0pt}%
\pgfpathmoveto{\pgfqpoint{8.020400in}{0.660000in}}%
\pgfpathlineto{\pgfqpoint{8.026600in}{0.660000in}}%
\pgfpathlineto{\pgfqpoint{8.026600in}{3.202310in}}%
\pgfpathlineto{\pgfqpoint{8.020400in}{3.202310in}}%
\pgfpathlineto{\pgfqpoint{8.020400in}{0.660000in}}%
\pgfpathclose%
\pgfusepath{fill}%
\end{pgfscope}%
\begin{pgfscope}%
\pgfpathrectangle{\pgfqpoint{1.250000in}{0.660000in}}{\pgfqpoint{7.750000in}{4.620000in}}%
\pgfusepath{clip}%
\pgfsetbuttcap%
\pgfsetmiterjoin%
\definecolor{currentfill}{rgb}{0.000000,0.000000,0.000000}%
\pgfsetfillcolor{currentfill}%
\pgfsetlinewidth{0.000000pt}%
\definecolor{currentstroke}{rgb}{0.000000,0.000000,0.000000}%
\pgfsetstrokecolor{currentstroke}%
\pgfsetstrokeopacity{0.000000}%
\pgfsetdash{}{0pt}%
\pgfpathmoveto{\pgfqpoint{8.028150in}{0.660000in}}%
\pgfpathlineto{\pgfqpoint{8.034350in}{0.660000in}}%
\pgfpathlineto{\pgfqpoint{8.034350in}{3.202188in}}%
\pgfpathlineto{\pgfqpoint{8.028150in}{3.202188in}}%
\pgfpathlineto{\pgfqpoint{8.028150in}{0.660000in}}%
\pgfpathclose%
\pgfusepath{fill}%
\end{pgfscope}%
\begin{pgfscope}%
\pgfpathrectangle{\pgfqpoint{1.250000in}{0.660000in}}{\pgfqpoint{7.750000in}{4.620000in}}%
\pgfusepath{clip}%
\pgfsetbuttcap%
\pgfsetmiterjoin%
\definecolor{currentfill}{rgb}{0.000000,0.000000,0.000000}%
\pgfsetfillcolor{currentfill}%
\pgfsetlinewidth{0.000000pt}%
\definecolor{currentstroke}{rgb}{0.000000,0.000000,0.000000}%
\pgfsetstrokecolor{currentstroke}%
\pgfsetstrokeopacity{0.000000}%
\pgfsetdash{}{0pt}%
\pgfpathmoveto{\pgfqpoint{8.035900in}{0.660000in}}%
\pgfpathlineto{\pgfqpoint{8.042100in}{0.660000in}}%
\pgfpathlineto{\pgfqpoint{8.042100in}{3.202045in}}%
\pgfpathlineto{\pgfqpoint{8.035900in}{3.202045in}}%
\pgfpathlineto{\pgfqpoint{8.035900in}{0.660000in}}%
\pgfpathclose%
\pgfusepath{fill}%
\end{pgfscope}%
\begin{pgfscope}%
\pgfpathrectangle{\pgfqpoint{1.250000in}{0.660000in}}{\pgfqpoint{7.750000in}{4.620000in}}%
\pgfusepath{clip}%
\pgfsetbuttcap%
\pgfsetmiterjoin%
\definecolor{currentfill}{rgb}{0.000000,0.000000,0.000000}%
\pgfsetfillcolor{currentfill}%
\pgfsetlinewidth{0.000000pt}%
\definecolor{currentstroke}{rgb}{0.000000,0.000000,0.000000}%
\pgfsetstrokecolor{currentstroke}%
\pgfsetstrokeopacity{0.000000}%
\pgfsetdash{}{0pt}%
\pgfpathmoveto{\pgfqpoint{8.043650in}{0.660000in}}%
\pgfpathlineto{\pgfqpoint{8.049850in}{0.660000in}}%
\pgfpathlineto{\pgfqpoint{8.049850in}{3.204703in}}%
\pgfpathlineto{\pgfqpoint{8.043650in}{3.204703in}}%
\pgfpathlineto{\pgfqpoint{8.043650in}{0.660000in}}%
\pgfpathclose%
\pgfusepath{fill}%
\end{pgfscope}%
\begin{pgfscope}%
\pgfpathrectangle{\pgfqpoint{1.250000in}{0.660000in}}{\pgfqpoint{7.750000in}{4.620000in}}%
\pgfusepath{clip}%
\pgfsetbuttcap%
\pgfsetmiterjoin%
\definecolor{currentfill}{rgb}{0.000000,0.000000,0.000000}%
\pgfsetfillcolor{currentfill}%
\pgfsetlinewidth{0.000000pt}%
\definecolor{currentstroke}{rgb}{0.000000,0.000000,0.000000}%
\pgfsetstrokecolor{currentstroke}%
\pgfsetstrokeopacity{0.000000}%
\pgfsetdash{}{0pt}%
\pgfpathmoveto{\pgfqpoint{8.051400in}{0.660000in}}%
\pgfpathlineto{\pgfqpoint{8.057600in}{0.660000in}}%
\pgfpathlineto{\pgfqpoint{8.057600in}{3.201693in}}%
\pgfpathlineto{\pgfqpoint{8.051400in}{3.201693in}}%
\pgfpathlineto{\pgfqpoint{8.051400in}{0.660000in}}%
\pgfpathclose%
\pgfusepath{fill}%
\end{pgfscope}%
\begin{pgfscope}%
\pgfpathrectangle{\pgfqpoint{1.250000in}{0.660000in}}{\pgfqpoint{7.750000in}{4.620000in}}%
\pgfusepath{clip}%
\pgfsetbuttcap%
\pgfsetmiterjoin%
\definecolor{currentfill}{rgb}{0.000000,0.000000,0.000000}%
\pgfsetfillcolor{currentfill}%
\pgfsetlinewidth{0.000000pt}%
\definecolor{currentstroke}{rgb}{0.000000,0.000000,0.000000}%
\pgfsetstrokecolor{currentstroke}%
\pgfsetstrokeopacity{0.000000}%
\pgfsetdash{}{0pt}%
\pgfpathmoveto{\pgfqpoint{8.059150in}{0.660000in}}%
\pgfpathlineto{\pgfqpoint{8.065350in}{0.660000in}}%
\pgfpathlineto{\pgfqpoint{8.065350in}{3.201483in}}%
\pgfpathlineto{\pgfqpoint{8.059150in}{3.201483in}}%
\pgfpathlineto{\pgfqpoint{8.059150in}{0.660000in}}%
\pgfpathclose%
\pgfusepath{fill}%
\end{pgfscope}%
\begin{pgfscope}%
\pgfpathrectangle{\pgfqpoint{1.250000in}{0.660000in}}{\pgfqpoint{7.750000in}{4.620000in}}%
\pgfusepath{clip}%
\pgfsetbuttcap%
\pgfsetmiterjoin%
\definecolor{currentfill}{rgb}{0.000000,0.000000,0.000000}%
\pgfsetfillcolor{currentfill}%
\pgfsetlinewidth{0.000000pt}%
\definecolor{currentstroke}{rgb}{0.000000,0.000000,0.000000}%
\pgfsetstrokecolor{currentstroke}%
\pgfsetstrokeopacity{0.000000}%
\pgfsetdash{}{0pt}%
\pgfpathmoveto{\pgfqpoint{8.066900in}{0.660000in}}%
\pgfpathlineto{\pgfqpoint{8.073100in}{0.660000in}}%
\pgfpathlineto{\pgfqpoint{8.073100in}{3.201339in}}%
\pgfpathlineto{\pgfqpoint{8.066900in}{3.201339in}}%
\pgfpathlineto{\pgfqpoint{8.066900in}{0.660000in}}%
\pgfpathclose%
\pgfusepath{fill}%
\end{pgfscope}%
\begin{pgfscope}%
\pgfpathrectangle{\pgfqpoint{1.250000in}{0.660000in}}{\pgfqpoint{7.750000in}{4.620000in}}%
\pgfusepath{clip}%
\pgfsetbuttcap%
\pgfsetmiterjoin%
\definecolor{currentfill}{rgb}{0.000000,0.000000,0.000000}%
\pgfsetfillcolor{currentfill}%
\pgfsetlinewidth{0.000000pt}%
\definecolor{currentstroke}{rgb}{0.000000,0.000000,0.000000}%
\pgfsetstrokecolor{currentstroke}%
\pgfsetstrokeopacity{0.000000}%
\pgfsetdash{}{0pt}%
\pgfpathmoveto{\pgfqpoint{8.074650in}{0.660000in}}%
\pgfpathlineto{\pgfqpoint{8.080850in}{0.660000in}}%
\pgfpathlineto{\pgfqpoint{8.080850in}{3.201368in}}%
\pgfpathlineto{\pgfqpoint{8.074650in}{3.201368in}}%
\pgfpathlineto{\pgfqpoint{8.074650in}{0.660000in}}%
\pgfpathclose%
\pgfusepath{fill}%
\end{pgfscope}%
\begin{pgfscope}%
\pgfpathrectangle{\pgfqpoint{1.250000in}{0.660000in}}{\pgfqpoint{7.750000in}{4.620000in}}%
\pgfusepath{clip}%
\pgfsetbuttcap%
\pgfsetmiterjoin%
\definecolor{currentfill}{rgb}{0.000000,0.000000,0.000000}%
\pgfsetfillcolor{currentfill}%
\pgfsetlinewidth{0.000000pt}%
\definecolor{currentstroke}{rgb}{0.000000,0.000000,0.000000}%
\pgfsetstrokecolor{currentstroke}%
\pgfsetstrokeopacity{0.000000}%
\pgfsetdash{}{0pt}%
\pgfpathmoveto{\pgfqpoint{8.082400in}{0.660000in}}%
\pgfpathlineto{\pgfqpoint{8.088600in}{0.660000in}}%
\pgfpathlineto{\pgfqpoint{8.088600in}{3.201377in}}%
\pgfpathlineto{\pgfqpoint{8.082400in}{3.201377in}}%
\pgfpathlineto{\pgfqpoint{8.082400in}{0.660000in}}%
\pgfpathclose%
\pgfusepath{fill}%
\end{pgfscope}%
\begin{pgfscope}%
\pgfpathrectangle{\pgfqpoint{1.250000in}{0.660000in}}{\pgfqpoint{7.750000in}{4.620000in}}%
\pgfusepath{clip}%
\pgfsetbuttcap%
\pgfsetmiterjoin%
\definecolor{currentfill}{rgb}{0.000000,0.000000,0.000000}%
\pgfsetfillcolor{currentfill}%
\pgfsetlinewidth{0.000000pt}%
\definecolor{currentstroke}{rgb}{0.000000,0.000000,0.000000}%
\pgfsetstrokecolor{currentstroke}%
\pgfsetstrokeopacity{0.000000}%
\pgfsetdash{}{0pt}%
\pgfpathmoveto{\pgfqpoint{8.090150in}{0.660000in}}%
\pgfpathlineto{\pgfqpoint{8.096350in}{0.660000in}}%
\pgfpathlineto{\pgfqpoint{8.096350in}{3.201365in}}%
\pgfpathlineto{\pgfqpoint{8.090150in}{3.201365in}}%
\pgfpathlineto{\pgfqpoint{8.090150in}{0.660000in}}%
\pgfpathclose%
\pgfusepath{fill}%
\end{pgfscope}%
\begin{pgfscope}%
\pgfpathrectangle{\pgfqpoint{1.250000in}{0.660000in}}{\pgfqpoint{7.750000in}{4.620000in}}%
\pgfusepath{clip}%
\pgfsetbuttcap%
\pgfsetmiterjoin%
\definecolor{currentfill}{rgb}{0.000000,0.000000,0.000000}%
\pgfsetfillcolor{currentfill}%
\pgfsetlinewidth{0.000000pt}%
\definecolor{currentstroke}{rgb}{0.000000,0.000000,0.000000}%
\pgfsetstrokecolor{currentstroke}%
\pgfsetstrokeopacity{0.000000}%
\pgfsetdash{}{0pt}%
\pgfpathmoveto{\pgfqpoint{8.097900in}{0.660000in}}%
\pgfpathlineto{\pgfqpoint{8.104100in}{0.660000in}}%
\pgfpathlineto{\pgfqpoint{8.104100in}{3.201331in}}%
\pgfpathlineto{\pgfqpoint{8.097900in}{3.201331in}}%
\pgfpathlineto{\pgfqpoint{8.097900in}{0.660000in}}%
\pgfpathclose%
\pgfusepath{fill}%
\end{pgfscope}%
\begin{pgfscope}%
\pgfpathrectangle{\pgfqpoint{1.250000in}{0.660000in}}{\pgfqpoint{7.750000in}{4.620000in}}%
\pgfusepath{clip}%
\pgfsetbuttcap%
\pgfsetmiterjoin%
\definecolor{currentfill}{rgb}{0.000000,0.000000,0.000000}%
\pgfsetfillcolor{currentfill}%
\pgfsetlinewidth{0.000000pt}%
\definecolor{currentstroke}{rgb}{0.000000,0.000000,0.000000}%
\pgfsetstrokecolor{currentstroke}%
\pgfsetstrokeopacity{0.000000}%
\pgfsetdash{}{0pt}%
\pgfpathmoveto{\pgfqpoint{8.105650in}{0.660000in}}%
\pgfpathlineto{\pgfqpoint{8.111850in}{0.660000in}}%
\pgfpathlineto{\pgfqpoint{8.111850in}{3.201276in}}%
\pgfpathlineto{\pgfqpoint{8.105650in}{3.201276in}}%
\pgfpathlineto{\pgfqpoint{8.105650in}{0.660000in}}%
\pgfpathclose%
\pgfusepath{fill}%
\end{pgfscope}%
\begin{pgfscope}%
\pgfpathrectangle{\pgfqpoint{1.250000in}{0.660000in}}{\pgfqpoint{7.750000in}{4.620000in}}%
\pgfusepath{clip}%
\pgfsetbuttcap%
\pgfsetmiterjoin%
\definecolor{currentfill}{rgb}{0.000000,0.000000,0.000000}%
\pgfsetfillcolor{currentfill}%
\pgfsetlinewidth{0.000000pt}%
\definecolor{currentstroke}{rgb}{0.000000,0.000000,0.000000}%
\pgfsetstrokecolor{currentstroke}%
\pgfsetstrokeopacity{0.000000}%
\pgfsetdash{}{0pt}%
\pgfpathmoveto{\pgfqpoint{8.113400in}{0.660000in}}%
\pgfpathlineto{\pgfqpoint{8.119600in}{0.660000in}}%
\pgfpathlineto{\pgfqpoint{8.119600in}{3.201198in}}%
\pgfpathlineto{\pgfqpoint{8.113400in}{3.201198in}}%
\pgfpathlineto{\pgfqpoint{8.113400in}{0.660000in}}%
\pgfpathclose%
\pgfusepath{fill}%
\end{pgfscope}%
\begin{pgfscope}%
\pgfpathrectangle{\pgfqpoint{1.250000in}{0.660000in}}{\pgfqpoint{7.750000in}{4.620000in}}%
\pgfusepath{clip}%
\pgfsetbuttcap%
\pgfsetmiterjoin%
\definecolor{currentfill}{rgb}{0.000000,0.000000,0.000000}%
\pgfsetfillcolor{currentfill}%
\pgfsetlinewidth{0.000000pt}%
\definecolor{currentstroke}{rgb}{0.000000,0.000000,0.000000}%
\pgfsetstrokecolor{currentstroke}%
\pgfsetstrokeopacity{0.000000}%
\pgfsetdash{}{0pt}%
\pgfpathmoveto{\pgfqpoint{8.121150in}{0.660000in}}%
\pgfpathlineto{\pgfqpoint{8.127350in}{0.660000in}}%
\pgfpathlineto{\pgfqpoint{8.127350in}{3.201098in}}%
\pgfpathlineto{\pgfqpoint{8.121150in}{3.201098in}}%
\pgfpathlineto{\pgfqpoint{8.121150in}{0.660000in}}%
\pgfpathclose%
\pgfusepath{fill}%
\end{pgfscope}%
\begin{pgfscope}%
\pgfpathrectangle{\pgfqpoint{1.250000in}{0.660000in}}{\pgfqpoint{7.750000in}{4.620000in}}%
\pgfusepath{clip}%
\pgfsetbuttcap%
\pgfsetmiterjoin%
\definecolor{currentfill}{rgb}{0.000000,0.000000,0.000000}%
\pgfsetfillcolor{currentfill}%
\pgfsetlinewidth{0.000000pt}%
\definecolor{currentstroke}{rgb}{0.000000,0.000000,0.000000}%
\pgfsetstrokecolor{currentstroke}%
\pgfsetstrokeopacity{0.000000}%
\pgfsetdash{}{0pt}%
\pgfpathmoveto{\pgfqpoint{8.128900in}{0.660000in}}%
\pgfpathlineto{\pgfqpoint{8.135100in}{0.660000in}}%
\pgfpathlineto{\pgfqpoint{8.135100in}{3.200974in}}%
\pgfpathlineto{\pgfqpoint{8.128900in}{3.200974in}}%
\pgfpathlineto{\pgfqpoint{8.128900in}{0.660000in}}%
\pgfpathclose%
\pgfusepath{fill}%
\end{pgfscope}%
\begin{pgfscope}%
\pgfpathrectangle{\pgfqpoint{1.250000in}{0.660000in}}{\pgfqpoint{7.750000in}{4.620000in}}%
\pgfusepath{clip}%
\pgfsetbuttcap%
\pgfsetmiterjoin%
\definecolor{currentfill}{rgb}{0.000000,0.000000,0.000000}%
\pgfsetfillcolor{currentfill}%
\pgfsetlinewidth{0.000000pt}%
\definecolor{currentstroke}{rgb}{0.000000,0.000000,0.000000}%
\pgfsetstrokecolor{currentstroke}%
\pgfsetstrokeopacity{0.000000}%
\pgfsetdash{}{0pt}%
\pgfpathmoveto{\pgfqpoint{8.136650in}{0.660000in}}%
\pgfpathlineto{\pgfqpoint{8.142850in}{0.660000in}}%
\pgfpathlineto{\pgfqpoint{8.142850in}{3.200827in}}%
\pgfpathlineto{\pgfqpoint{8.136650in}{3.200827in}}%
\pgfpathlineto{\pgfqpoint{8.136650in}{0.660000in}}%
\pgfpathclose%
\pgfusepath{fill}%
\end{pgfscope}%
\begin{pgfscope}%
\pgfpathrectangle{\pgfqpoint{1.250000in}{0.660000in}}{\pgfqpoint{7.750000in}{4.620000in}}%
\pgfusepath{clip}%
\pgfsetbuttcap%
\pgfsetmiterjoin%
\definecolor{currentfill}{rgb}{0.000000,0.000000,0.000000}%
\pgfsetfillcolor{currentfill}%
\pgfsetlinewidth{0.000000pt}%
\definecolor{currentstroke}{rgb}{0.000000,0.000000,0.000000}%
\pgfsetstrokecolor{currentstroke}%
\pgfsetstrokeopacity{0.000000}%
\pgfsetdash{}{0pt}%
\pgfpathmoveto{\pgfqpoint{8.144400in}{0.660000in}}%
\pgfpathlineto{\pgfqpoint{8.150600in}{0.660000in}}%
\pgfpathlineto{\pgfqpoint{8.150600in}{3.200657in}}%
\pgfpathlineto{\pgfqpoint{8.144400in}{3.200657in}}%
\pgfpathlineto{\pgfqpoint{8.144400in}{0.660000in}}%
\pgfpathclose%
\pgfusepath{fill}%
\end{pgfscope}%
\begin{pgfscope}%
\pgfpathrectangle{\pgfqpoint{1.250000in}{0.660000in}}{\pgfqpoint{7.750000in}{4.620000in}}%
\pgfusepath{clip}%
\pgfsetbuttcap%
\pgfsetmiterjoin%
\definecolor{currentfill}{rgb}{0.000000,0.000000,0.000000}%
\pgfsetfillcolor{currentfill}%
\pgfsetlinewidth{0.000000pt}%
\definecolor{currentstroke}{rgb}{0.000000,0.000000,0.000000}%
\pgfsetstrokecolor{currentstroke}%
\pgfsetstrokeopacity{0.000000}%
\pgfsetdash{}{0pt}%
\pgfpathmoveto{\pgfqpoint{8.152150in}{0.660000in}}%
\pgfpathlineto{\pgfqpoint{8.158350in}{0.660000in}}%
\pgfpathlineto{\pgfqpoint{8.158350in}{3.207162in}}%
\pgfpathlineto{\pgfqpoint{8.152150in}{3.207162in}}%
\pgfpathlineto{\pgfqpoint{8.152150in}{0.660000in}}%
\pgfpathclose%
\pgfusepath{fill}%
\end{pgfscope}%
\begin{pgfscope}%
\pgfpathrectangle{\pgfqpoint{1.250000in}{0.660000in}}{\pgfqpoint{7.750000in}{4.620000in}}%
\pgfusepath{clip}%
\pgfsetbuttcap%
\pgfsetmiterjoin%
\definecolor{currentfill}{rgb}{0.000000,0.000000,0.000000}%
\pgfsetfillcolor{currentfill}%
\pgfsetlinewidth{0.000000pt}%
\definecolor{currentstroke}{rgb}{0.000000,0.000000,0.000000}%
\pgfsetstrokecolor{currentstroke}%
\pgfsetstrokeopacity{0.000000}%
\pgfsetdash{}{0pt}%
\pgfpathmoveto{\pgfqpoint{8.159900in}{0.660000in}}%
\pgfpathlineto{\pgfqpoint{8.166100in}{0.660000in}}%
\pgfpathlineto{\pgfqpoint{8.166100in}{3.200242in}}%
\pgfpathlineto{\pgfqpoint{8.159900in}{3.200242in}}%
\pgfpathlineto{\pgfqpoint{8.159900in}{0.660000in}}%
\pgfpathclose%
\pgfusepath{fill}%
\end{pgfscope}%
\begin{pgfscope}%
\pgfpathrectangle{\pgfqpoint{1.250000in}{0.660000in}}{\pgfqpoint{7.750000in}{4.620000in}}%
\pgfusepath{clip}%
\pgfsetbuttcap%
\pgfsetmiterjoin%
\definecolor{currentfill}{rgb}{0.000000,0.000000,0.000000}%
\pgfsetfillcolor{currentfill}%
\pgfsetlinewidth{0.000000pt}%
\definecolor{currentstroke}{rgb}{0.000000,0.000000,0.000000}%
\pgfsetstrokecolor{currentstroke}%
\pgfsetstrokeopacity{0.000000}%
\pgfsetdash{}{0pt}%
\pgfpathmoveto{\pgfqpoint{8.167650in}{0.660000in}}%
\pgfpathlineto{\pgfqpoint{8.173850in}{0.660000in}}%
\pgfpathlineto{\pgfqpoint{8.173850in}{3.200192in}}%
\pgfpathlineto{\pgfqpoint{8.167650in}{3.200192in}}%
\pgfpathlineto{\pgfqpoint{8.167650in}{0.660000in}}%
\pgfpathclose%
\pgfusepath{fill}%
\end{pgfscope}%
\begin{pgfscope}%
\pgfpathrectangle{\pgfqpoint{1.250000in}{0.660000in}}{\pgfqpoint{7.750000in}{4.620000in}}%
\pgfusepath{clip}%
\pgfsetbuttcap%
\pgfsetmiterjoin%
\definecolor{currentfill}{rgb}{0.000000,0.000000,0.000000}%
\pgfsetfillcolor{currentfill}%
\pgfsetlinewidth{0.000000pt}%
\definecolor{currentstroke}{rgb}{0.000000,0.000000,0.000000}%
\pgfsetstrokecolor{currentstroke}%
\pgfsetstrokeopacity{0.000000}%
\pgfsetdash{}{0pt}%
\pgfpathmoveto{\pgfqpoint{8.175400in}{0.660000in}}%
\pgfpathlineto{\pgfqpoint{8.181600in}{0.660000in}}%
\pgfpathlineto{\pgfqpoint{8.181600in}{3.200251in}}%
\pgfpathlineto{\pgfqpoint{8.175400in}{3.200251in}}%
\pgfpathlineto{\pgfqpoint{8.175400in}{0.660000in}}%
\pgfpathclose%
\pgfusepath{fill}%
\end{pgfscope}%
\begin{pgfscope}%
\pgfpathrectangle{\pgfqpoint{1.250000in}{0.660000in}}{\pgfqpoint{7.750000in}{4.620000in}}%
\pgfusepath{clip}%
\pgfsetbuttcap%
\pgfsetmiterjoin%
\definecolor{currentfill}{rgb}{0.000000,0.000000,0.000000}%
\pgfsetfillcolor{currentfill}%
\pgfsetlinewidth{0.000000pt}%
\definecolor{currentstroke}{rgb}{0.000000,0.000000,0.000000}%
\pgfsetstrokecolor{currentstroke}%
\pgfsetstrokeopacity{0.000000}%
\pgfsetdash{}{0pt}%
\pgfpathmoveto{\pgfqpoint{8.183150in}{0.660000in}}%
\pgfpathlineto{\pgfqpoint{8.189350in}{0.660000in}}%
\pgfpathlineto{\pgfqpoint{8.189350in}{3.204942in}}%
\pgfpathlineto{\pgfqpoint{8.183150in}{3.204942in}}%
\pgfpathlineto{\pgfqpoint{8.183150in}{0.660000in}}%
\pgfpathclose%
\pgfusepath{fill}%
\end{pgfscope}%
\begin{pgfscope}%
\pgfpathrectangle{\pgfqpoint{1.250000in}{0.660000in}}{\pgfqpoint{7.750000in}{4.620000in}}%
\pgfusepath{clip}%
\pgfsetbuttcap%
\pgfsetmiterjoin%
\definecolor{currentfill}{rgb}{0.000000,0.000000,0.000000}%
\pgfsetfillcolor{currentfill}%
\pgfsetlinewidth{0.000000pt}%
\definecolor{currentstroke}{rgb}{0.000000,0.000000,0.000000}%
\pgfsetstrokecolor{currentstroke}%
\pgfsetstrokeopacity{0.000000}%
\pgfsetdash{}{0pt}%
\pgfpathmoveto{\pgfqpoint{8.190900in}{0.660000in}}%
\pgfpathlineto{\pgfqpoint{8.197100in}{0.660000in}}%
\pgfpathlineto{\pgfqpoint{8.197100in}{3.200205in}}%
\pgfpathlineto{\pgfqpoint{8.190900in}{3.200205in}}%
\pgfpathlineto{\pgfqpoint{8.190900in}{0.660000in}}%
\pgfpathclose%
\pgfusepath{fill}%
\end{pgfscope}%
\begin{pgfscope}%
\pgfpathrectangle{\pgfqpoint{1.250000in}{0.660000in}}{\pgfqpoint{7.750000in}{4.620000in}}%
\pgfusepath{clip}%
\pgfsetbuttcap%
\pgfsetmiterjoin%
\definecolor{currentfill}{rgb}{0.000000,0.000000,0.000000}%
\pgfsetfillcolor{currentfill}%
\pgfsetlinewidth{0.000000pt}%
\definecolor{currentstroke}{rgb}{0.000000,0.000000,0.000000}%
\pgfsetstrokecolor{currentstroke}%
\pgfsetstrokeopacity{0.000000}%
\pgfsetdash{}{0pt}%
\pgfpathmoveto{\pgfqpoint{8.198650in}{0.660000in}}%
\pgfpathlineto{\pgfqpoint{8.204850in}{0.660000in}}%
\pgfpathlineto{\pgfqpoint{8.204850in}{3.200163in}}%
\pgfpathlineto{\pgfqpoint{8.198650in}{3.200163in}}%
\pgfpathlineto{\pgfqpoint{8.198650in}{0.660000in}}%
\pgfpathclose%
\pgfusepath{fill}%
\end{pgfscope}%
\begin{pgfscope}%
\pgfpathrectangle{\pgfqpoint{1.250000in}{0.660000in}}{\pgfqpoint{7.750000in}{4.620000in}}%
\pgfusepath{clip}%
\pgfsetbuttcap%
\pgfsetmiterjoin%
\definecolor{currentfill}{rgb}{0.000000,0.000000,0.000000}%
\pgfsetfillcolor{currentfill}%
\pgfsetlinewidth{0.000000pt}%
\definecolor{currentstroke}{rgb}{0.000000,0.000000,0.000000}%
\pgfsetstrokecolor{currentstroke}%
\pgfsetstrokeopacity{0.000000}%
\pgfsetdash{}{0pt}%
\pgfpathmoveto{\pgfqpoint{8.206400in}{0.660000in}}%
\pgfpathlineto{\pgfqpoint{8.212600in}{0.660000in}}%
\pgfpathlineto{\pgfqpoint{8.212600in}{3.200097in}}%
\pgfpathlineto{\pgfqpoint{8.206400in}{3.200097in}}%
\pgfpathlineto{\pgfqpoint{8.206400in}{0.660000in}}%
\pgfpathclose%
\pgfusepath{fill}%
\end{pgfscope}%
\begin{pgfscope}%
\pgfpathrectangle{\pgfqpoint{1.250000in}{0.660000in}}{\pgfqpoint{7.750000in}{4.620000in}}%
\pgfusepath{clip}%
\pgfsetbuttcap%
\pgfsetmiterjoin%
\definecolor{currentfill}{rgb}{0.000000,0.000000,0.000000}%
\pgfsetfillcolor{currentfill}%
\pgfsetlinewidth{0.000000pt}%
\definecolor{currentstroke}{rgb}{0.000000,0.000000,0.000000}%
\pgfsetstrokecolor{currentstroke}%
\pgfsetstrokeopacity{0.000000}%
\pgfsetdash{}{0pt}%
\pgfpathmoveto{\pgfqpoint{8.214150in}{0.660000in}}%
\pgfpathlineto{\pgfqpoint{8.220350in}{0.660000in}}%
\pgfpathlineto{\pgfqpoint{8.220350in}{3.201918in}}%
\pgfpathlineto{\pgfqpoint{8.214150in}{3.201918in}}%
\pgfpathlineto{\pgfqpoint{8.214150in}{0.660000in}}%
\pgfpathclose%
\pgfusepath{fill}%
\end{pgfscope}%
\begin{pgfscope}%
\pgfpathrectangle{\pgfqpoint{1.250000in}{0.660000in}}{\pgfqpoint{7.750000in}{4.620000in}}%
\pgfusepath{clip}%
\pgfsetbuttcap%
\pgfsetmiterjoin%
\definecolor{currentfill}{rgb}{0.000000,0.000000,0.000000}%
\pgfsetfillcolor{currentfill}%
\pgfsetlinewidth{0.000000pt}%
\definecolor{currentstroke}{rgb}{0.000000,0.000000,0.000000}%
\pgfsetstrokecolor{currentstroke}%
\pgfsetstrokeopacity{0.000000}%
\pgfsetdash{}{0pt}%
\pgfpathmoveto{\pgfqpoint{8.221900in}{0.660000in}}%
\pgfpathlineto{\pgfqpoint{8.228100in}{0.660000in}}%
\pgfpathlineto{\pgfqpoint{8.228100in}{3.199891in}}%
\pgfpathlineto{\pgfqpoint{8.221900in}{3.199891in}}%
\pgfpathlineto{\pgfqpoint{8.221900in}{0.660000in}}%
\pgfpathclose%
\pgfusepath{fill}%
\end{pgfscope}%
\begin{pgfscope}%
\pgfpathrectangle{\pgfqpoint{1.250000in}{0.660000in}}{\pgfqpoint{7.750000in}{4.620000in}}%
\pgfusepath{clip}%
\pgfsetbuttcap%
\pgfsetmiterjoin%
\definecolor{currentfill}{rgb}{0.000000,0.000000,0.000000}%
\pgfsetfillcolor{currentfill}%
\pgfsetlinewidth{0.000000pt}%
\definecolor{currentstroke}{rgb}{0.000000,0.000000,0.000000}%
\pgfsetstrokecolor{currentstroke}%
\pgfsetstrokeopacity{0.000000}%
\pgfsetdash{}{0pt}%
\pgfpathmoveto{\pgfqpoint{8.229650in}{0.660000in}}%
\pgfpathlineto{\pgfqpoint{8.235850in}{0.660000in}}%
\pgfpathlineto{\pgfqpoint{8.235850in}{3.199751in}}%
\pgfpathlineto{\pgfqpoint{8.229650in}{3.199751in}}%
\pgfpathlineto{\pgfqpoint{8.229650in}{0.660000in}}%
\pgfpathclose%
\pgfusepath{fill}%
\end{pgfscope}%
\begin{pgfscope}%
\pgfpathrectangle{\pgfqpoint{1.250000in}{0.660000in}}{\pgfqpoint{7.750000in}{4.620000in}}%
\pgfusepath{clip}%
\pgfsetbuttcap%
\pgfsetmiterjoin%
\definecolor{currentfill}{rgb}{0.000000,0.000000,0.000000}%
\pgfsetfillcolor{currentfill}%
\pgfsetlinewidth{0.000000pt}%
\definecolor{currentstroke}{rgb}{0.000000,0.000000,0.000000}%
\pgfsetstrokecolor{currentstroke}%
\pgfsetstrokeopacity{0.000000}%
\pgfsetdash{}{0pt}%
\pgfpathmoveto{\pgfqpoint{8.237400in}{0.660000in}}%
\pgfpathlineto{\pgfqpoint{8.243600in}{0.660000in}}%
\pgfpathlineto{\pgfqpoint{8.243600in}{3.199584in}}%
\pgfpathlineto{\pgfqpoint{8.237400in}{3.199584in}}%
\pgfpathlineto{\pgfqpoint{8.237400in}{0.660000in}}%
\pgfpathclose%
\pgfusepath{fill}%
\end{pgfscope}%
\begin{pgfscope}%
\pgfpathrectangle{\pgfqpoint{1.250000in}{0.660000in}}{\pgfqpoint{7.750000in}{4.620000in}}%
\pgfusepath{clip}%
\pgfsetbuttcap%
\pgfsetmiterjoin%
\definecolor{currentfill}{rgb}{0.000000,0.000000,0.000000}%
\pgfsetfillcolor{currentfill}%
\pgfsetlinewidth{0.000000pt}%
\definecolor{currentstroke}{rgb}{0.000000,0.000000,0.000000}%
\pgfsetstrokecolor{currentstroke}%
\pgfsetstrokeopacity{0.000000}%
\pgfsetdash{}{0pt}%
\pgfpathmoveto{\pgfqpoint{8.245150in}{0.660000in}}%
\pgfpathlineto{\pgfqpoint{8.251350in}{0.660000in}}%
\pgfpathlineto{\pgfqpoint{8.251350in}{3.199392in}}%
\pgfpathlineto{\pgfqpoint{8.245150in}{3.199392in}}%
\pgfpathlineto{\pgfqpoint{8.245150in}{0.660000in}}%
\pgfpathclose%
\pgfusepath{fill}%
\end{pgfscope}%
\begin{pgfscope}%
\pgfpathrectangle{\pgfqpoint{1.250000in}{0.660000in}}{\pgfqpoint{7.750000in}{4.620000in}}%
\pgfusepath{clip}%
\pgfsetbuttcap%
\pgfsetmiterjoin%
\definecolor{currentfill}{rgb}{0.000000,0.000000,0.000000}%
\pgfsetfillcolor{currentfill}%
\pgfsetlinewidth{0.000000pt}%
\definecolor{currentstroke}{rgb}{0.000000,0.000000,0.000000}%
\pgfsetstrokecolor{currentstroke}%
\pgfsetstrokeopacity{0.000000}%
\pgfsetdash{}{0pt}%
\pgfpathmoveto{\pgfqpoint{8.252900in}{0.660000in}}%
\pgfpathlineto{\pgfqpoint{8.259100in}{0.660000in}}%
\pgfpathlineto{\pgfqpoint{8.259100in}{3.199173in}}%
\pgfpathlineto{\pgfqpoint{8.252900in}{3.199173in}}%
\pgfpathlineto{\pgfqpoint{8.252900in}{0.660000in}}%
\pgfpathclose%
\pgfusepath{fill}%
\end{pgfscope}%
\begin{pgfscope}%
\pgfpathrectangle{\pgfqpoint{1.250000in}{0.660000in}}{\pgfqpoint{7.750000in}{4.620000in}}%
\pgfusepath{clip}%
\pgfsetbuttcap%
\pgfsetmiterjoin%
\definecolor{currentfill}{rgb}{0.000000,0.000000,0.000000}%
\pgfsetfillcolor{currentfill}%
\pgfsetlinewidth{0.000000pt}%
\definecolor{currentstroke}{rgb}{0.000000,0.000000,0.000000}%
\pgfsetstrokecolor{currentstroke}%
\pgfsetstrokeopacity{0.000000}%
\pgfsetdash{}{0pt}%
\pgfpathmoveto{\pgfqpoint{8.260650in}{0.660000in}}%
\pgfpathlineto{\pgfqpoint{8.266850in}{0.660000in}}%
\pgfpathlineto{\pgfqpoint{8.266850in}{3.199101in}}%
\pgfpathlineto{\pgfqpoint{8.260650in}{3.199101in}}%
\pgfpathlineto{\pgfqpoint{8.260650in}{0.660000in}}%
\pgfpathclose%
\pgfusepath{fill}%
\end{pgfscope}%
\begin{pgfscope}%
\pgfpathrectangle{\pgfqpoint{1.250000in}{0.660000in}}{\pgfqpoint{7.750000in}{4.620000in}}%
\pgfusepath{clip}%
\pgfsetbuttcap%
\pgfsetmiterjoin%
\definecolor{currentfill}{rgb}{0.000000,0.000000,0.000000}%
\pgfsetfillcolor{currentfill}%
\pgfsetlinewidth{0.000000pt}%
\definecolor{currentstroke}{rgb}{0.000000,0.000000,0.000000}%
\pgfsetstrokecolor{currentstroke}%
\pgfsetstrokeopacity{0.000000}%
\pgfsetdash{}{0pt}%
\pgfpathmoveto{\pgfqpoint{8.268400in}{0.660000in}}%
\pgfpathlineto{\pgfqpoint{8.274600in}{0.660000in}}%
\pgfpathlineto{\pgfqpoint{8.274600in}{3.199135in}}%
\pgfpathlineto{\pgfqpoint{8.268400in}{3.199135in}}%
\pgfpathlineto{\pgfqpoint{8.268400in}{0.660000in}}%
\pgfpathclose%
\pgfusepath{fill}%
\end{pgfscope}%
\begin{pgfscope}%
\pgfpathrectangle{\pgfqpoint{1.250000in}{0.660000in}}{\pgfqpoint{7.750000in}{4.620000in}}%
\pgfusepath{clip}%
\pgfsetbuttcap%
\pgfsetmiterjoin%
\definecolor{currentfill}{rgb}{0.000000,0.000000,0.000000}%
\pgfsetfillcolor{currentfill}%
\pgfsetlinewidth{0.000000pt}%
\definecolor{currentstroke}{rgb}{0.000000,0.000000,0.000000}%
\pgfsetstrokecolor{currentstroke}%
\pgfsetstrokeopacity{0.000000}%
\pgfsetdash{}{0pt}%
\pgfpathmoveto{\pgfqpoint{8.276150in}{0.660000in}}%
\pgfpathlineto{\pgfqpoint{8.282350in}{0.660000in}}%
\pgfpathlineto{\pgfqpoint{8.282350in}{3.199144in}}%
\pgfpathlineto{\pgfqpoint{8.276150in}{3.199144in}}%
\pgfpathlineto{\pgfqpoint{8.276150in}{0.660000in}}%
\pgfpathclose%
\pgfusepath{fill}%
\end{pgfscope}%
\begin{pgfscope}%
\pgfpathrectangle{\pgfqpoint{1.250000in}{0.660000in}}{\pgfqpoint{7.750000in}{4.620000in}}%
\pgfusepath{clip}%
\pgfsetbuttcap%
\pgfsetmiterjoin%
\definecolor{currentfill}{rgb}{0.000000,0.000000,0.000000}%
\pgfsetfillcolor{currentfill}%
\pgfsetlinewidth{0.000000pt}%
\definecolor{currentstroke}{rgb}{0.000000,0.000000,0.000000}%
\pgfsetstrokecolor{currentstroke}%
\pgfsetstrokeopacity{0.000000}%
\pgfsetdash{}{0pt}%
\pgfpathmoveto{\pgfqpoint{8.283900in}{0.660000in}}%
\pgfpathlineto{\pgfqpoint{8.290100in}{0.660000in}}%
\pgfpathlineto{\pgfqpoint{8.290100in}{3.199128in}}%
\pgfpathlineto{\pgfqpoint{8.283900in}{3.199128in}}%
\pgfpathlineto{\pgfqpoint{8.283900in}{0.660000in}}%
\pgfpathclose%
\pgfusepath{fill}%
\end{pgfscope}%
\begin{pgfscope}%
\pgfpathrectangle{\pgfqpoint{1.250000in}{0.660000in}}{\pgfqpoint{7.750000in}{4.620000in}}%
\pgfusepath{clip}%
\pgfsetbuttcap%
\pgfsetmiterjoin%
\definecolor{currentfill}{rgb}{0.000000,0.000000,0.000000}%
\pgfsetfillcolor{currentfill}%
\pgfsetlinewidth{0.000000pt}%
\definecolor{currentstroke}{rgb}{0.000000,0.000000,0.000000}%
\pgfsetstrokecolor{currentstroke}%
\pgfsetstrokeopacity{0.000000}%
\pgfsetdash{}{0pt}%
\pgfpathmoveto{\pgfqpoint{8.291650in}{0.660000in}}%
\pgfpathlineto{\pgfqpoint{8.297850in}{0.660000in}}%
\pgfpathlineto{\pgfqpoint{8.297850in}{3.199087in}}%
\pgfpathlineto{\pgfqpoint{8.291650in}{3.199087in}}%
\pgfpathlineto{\pgfqpoint{8.291650in}{0.660000in}}%
\pgfpathclose%
\pgfusepath{fill}%
\end{pgfscope}%
\begin{pgfscope}%
\pgfpathrectangle{\pgfqpoint{1.250000in}{0.660000in}}{\pgfqpoint{7.750000in}{4.620000in}}%
\pgfusepath{clip}%
\pgfsetbuttcap%
\pgfsetmiterjoin%
\definecolor{currentfill}{rgb}{0.000000,0.000000,0.000000}%
\pgfsetfillcolor{currentfill}%
\pgfsetlinewidth{0.000000pt}%
\definecolor{currentstroke}{rgb}{0.000000,0.000000,0.000000}%
\pgfsetstrokecolor{currentstroke}%
\pgfsetstrokeopacity{0.000000}%
\pgfsetdash{}{0pt}%
\pgfpathmoveto{\pgfqpoint{8.299400in}{0.660000in}}%
\pgfpathlineto{\pgfqpoint{8.305600in}{0.660000in}}%
\pgfpathlineto{\pgfqpoint{8.305600in}{3.199020in}}%
\pgfpathlineto{\pgfqpoint{8.299400in}{3.199020in}}%
\pgfpathlineto{\pgfqpoint{8.299400in}{0.660000in}}%
\pgfpathclose%
\pgfusepath{fill}%
\end{pgfscope}%
\begin{pgfscope}%
\pgfpathrectangle{\pgfqpoint{1.250000in}{0.660000in}}{\pgfqpoint{7.750000in}{4.620000in}}%
\pgfusepath{clip}%
\pgfsetbuttcap%
\pgfsetmiterjoin%
\definecolor{currentfill}{rgb}{0.000000,0.000000,0.000000}%
\pgfsetfillcolor{currentfill}%
\pgfsetlinewidth{0.000000pt}%
\definecolor{currentstroke}{rgb}{0.000000,0.000000,0.000000}%
\pgfsetstrokecolor{currentstroke}%
\pgfsetstrokeopacity{0.000000}%
\pgfsetdash{}{0pt}%
\pgfpathmoveto{\pgfqpoint{8.307150in}{0.660000in}}%
\pgfpathlineto{\pgfqpoint{8.313350in}{0.660000in}}%
\pgfpathlineto{\pgfqpoint{8.313350in}{3.198927in}}%
\pgfpathlineto{\pgfqpoint{8.307150in}{3.198927in}}%
\pgfpathlineto{\pgfqpoint{8.307150in}{0.660000in}}%
\pgfpathclose%
\pgfusepath{fill}%
\end{pgfscope}%
\begin{pgfscope}%
\pgfpathrectangle{\pgfqpoint{1.250000in}{0.660000in}}{\pgfqpoint{7.750000in}{4.620000in}}%
\pgfusepath{clip}%
\pgfsetbuttcap%
\pgfsetmiterjoin%
\definecolor{currentfill}{rgb}{0.000000,0.000000,0.000000}%
\pgfsetfillcolor{currentfill}%
\pgfsetlinewidth{0.000000pt}%
\definecolor{currentstroke}{rgb}{0.000000,0.000000,0.000000}%
\pgfsetstrokecolor{currentstroke}%
\pgfsetstrokeopacity{0.000000}%
\pgfsetdash{}{0pt}%
\pgfpathmoveto{\pgfqpoint{8.314900in}{0.660000in}}%
\pgfpathlineto{\pgfqpoint{8.321100in}{0.660000in}}%
\pgfpathlineto{\pgfqpoint{8.321100in}{3.198807in}}%
\pgfpathlineto{\pgfqpoint{8.314900in}{3.198807in}}%
\pgfpathlineto{\pgfqpoint{8.314900in}{0.660000in}}%
\pgfpathclose%
\pgfusepath{fill}%
\end{pgfscope}%
\begin{pgfscope}%
\pgfpathrectangle{\pgfqpoint{1.250000in}{0.660000in}}{\pgfqpoint{7.750000in}{4.620000in}}%
\pgfusepath{clip}%
\pgfsetbuttcap%
\pgfsetmiterjoin%
\definecolor{currentfill}{rgb}{0.000000,0.000000,0.000000}%
\pgfsetfillcolor{currentfill}%
\pgfsetlinewidth{0.000000pt}%
\definecolor{currentstroke}{rgb}{0.000000,0.000000,0.000000}%
\pgfsetstrokecolor{currentstroke}%
\pgfsetstrokeopacity{0.000000}%
\pgfsetdash{}{0pt}%
\pgfpathmoveto{\pgfqpoint{8.322650in}{0.660000in}}%
\pgfpathlineto{\pgfqpoint{8.328850in}{0.660000in}}%
\pgfpathlineto{\pgfqpoint{8.328850in}{3.198660in}}%
\pgfpathlineto{\pgfqpoint{8.322650in}{3.198660in}}%
\pgfpathlineto{\pgfqpoint{8.322650in}{0.660000in}}%
\pgfpathclose%
\pgfusepath{fill}%
\end{pgfscope}%
\begin{pgfscope}%
\pgfpathrectangle{\pgfqpoint{1.250000in}{0.660000in}}{\pgfqpoint{7.750000in}{4.620000in}}%
\pgfusepath{clip}%
\pgfsetbuttcap%
\pgfsetmiterjoin%
\definecolor{currentfill}{rgb}{0.000000,0.000000,0.000000}%
\pgfsetfillcolor{currentfill}%
\pgfsetlinewidth{0.000000pt}%
\definecolor{currentstroke}{rgb}{0.000000,0.000000,0.000000}%
\pgfsetstrokecolor{currentstroke}%
\pgfsetstrokeopacity{0.000000}%
\pgfsetdash{}{0pt}%
\pgfpathmoveto{\pgfqpoint{8.330400in}{0.660000in}}%
\pgfpathlineto{\pgfqpoint{8.336600in}{0.660000in}}%
\pgfpathlineto{\pgfqpoint{8.336600in}{3.198485in}}%
\pgfpathlineto{\pgfqpoint{8.330400in}{3.198485in}}%
\pgfpathlineto{\pgfqpoint{8.330400in}{0.660000in}}%
\pgfpathclose%
\pgfusepath{fill}%
\end{pgfscope}%
\begin{pgfscope}%
\pgfpathrectangle{\pgfqpoint{1.250000in}{0.660000in}}{\pgfqpoint{7.750000in}{4.620000in}}%
\pgfusepath{clip}%
\pgfsetbuttcap%
\pgfsetmiterjoin%
\definecolor{currentfill}{rgb}{0.000000,0.000000,0.000000}%
\pgfsetfillcolor{currentfill}%
\pgfsetlinewidth{0.000000pt}%
\definecolor{currentstroke}{rgb}{0.000000,0.000000,0.000000}%
\pgfsetstrokecolor{currentstroke}%
\pgfsetstrokeopacity{0.000000}%
\pgfsetdash{}{0pt}%
\pgfpathmoveto{\pgfqpoint{8.338150in}{0.660000in}}%
\pgfpathlineto{\pgfqpoint{8.344350in}{0.660000in}}%
\pgfpathlineto{\pgfqpoint{8.344350in}{3.203821in}}%
\pgfpathlineto{\pgfqpoint{8.338150in}{3.203821in}}%
\pgfpathlineto{\pgfqpoint{8.338150in}{0.660000in}}%
\pgfpathclose%
\pgfusepath{fill}%
\end{pgfscope}%
\begin{pgfscope}%
\pgfpathrectangle{\pgfqpoint{1.250000in}{0.660000in}}{\pgfqpoint{7.750000in}{4.620000in}}%
\pgfusepath{clip}%
\pgfsetbuttcap%
\pgfsetmiterjoin%
\definecolor{currentfill}{rgb}{0.000000,0.000000,0.000000}%
\pgfsetfillcolor{currentfill}%
\pgfsetlinewidth{0.000000pt}%
\definecolor{currentstroke}{rgb}{0.000000,0.000000,0.000000}%
\pgfsetstrokecolor{currentstroke}%
\pgfsetstrokeopacity{0.000000}%
\pgfsetdash{}{0pt}%
\pgfpathmoveto{\pgfqpoint{8.345900in}{0.660000in}}%
\pgfpathlineto{\pgfqpoint{8.352100in}{0.660000in}}%
\pgfpathlineto{\pgfqpoint{8.352100in}{3.198049in}}%
\pgfpathlineto{\pgfqpoint{8.345900in}{3.198049in}}%
\pgfpathlineto{\pgfqpoint{8.345900in}{0.660000in}}%
\pgfpathclose%
\pgfusepath{fill}%
\end{pgfscope}%
\begin{pgfscope}%
\pgfpathrectangle{\pgfqpoint{1.250000in}{0.660000in}}{\pgfqpoint{7.750000in}{4.620000in}}%
\pgfusepath{clip}%
\pgfsetbuttcap%
\pgfsetmiterjoin%
\definecolor{currentfill}{rgb}{0.000000,0.000000,0.000000}%
\pgfsetfillcolor{currentfill}%
\pgfsetlinewidth{0.000000pt}%
\definecolor{currentstroke}{rgb}{0.000000,0.000000,0.000000}%
\pgfsetstrokecolor{currentstroke}%
\pgfsetstrokeopacity{0.000000}%
\pgfsetdash{}{0pt}%
\pgfpathmoveto{\pgfqpoint{8.353650in}{0.660000in}}%
\pgfpathlineto{\pgfqpoint{8.359850in}{0.660000in}}%
\pgfpathlineto{\pgfqpoint{8.359850in}{3.215639in}}%
\pgfpathlineto{\pgfqpoint{8.353650in}{3.215639in}}%
\pgfpathlineto{\pgfqpoint{8.353650in}{0.660000in}}%
\pgfpathclose%
\pgfusepath{fill}%
\end{pgfscope}%
\begin{pgfscope}%
\pgfpathrectangle{\pgfqpoint{1.250000in}{0.660000in}}{\pgfqpoint{7.750000in}{4.620000in}}%
\pgfusepath{clip}%
\pgfsetbuttcap%
\pgfsetmiterjoin%
\definecolor{currentfill}{rgb}{0.000000,0.000000,0.000000}%
\pgfsetfillcolor{currentfill}%
\pgfsetlinewidth{0.000000pt}%
\definecolor{currentstroke}{rgb}{0.000000,0.000000,0.000000}%
\pgfsetstrokecolor{currentstroke}%
\pgfsetstrokeopacity{0.000000}%
\pgfsetdash{}{0pt}%
\pgfpathmoveto{\pgfqpoint{8.361400in}{0.660000in}}%
\pgfpathlineto{\pgfqpoint{8.367600in}{0.660000in}}%
\pgfpathlineto{\pgfqpoint{8.367600in}{3.198122in}}%
\pgfpathlineto{\pgfqpoint{8.361400in}{3.198122in}}%
\pgfpathlineto{\pgfqpoint{8.361400in}{0.660000in}}%
\pgfpathclose%
\pgfusepath{fill}%
\end{pgfscope}%
\begin{pgfscope}%
\pgfpathrectangle{\pgfqpoint{1.250000in}{0.660000in}}{\pgfqpoint{7.750000in}{4.620000in}}%
\pgfusepath{clip}%
\pgfsetbuttcap%
\pgfsetmiterjoin%
\definecolor{currentfill}{rgb}{0.000000,0.000000,0.000000}%
\pgfsetfillcolor{currentfill}%
\pgfsetlinewidth{0.000000pt}%
\definecolor{currentstroke}{rgb}{0.000000,0.000000,0.000000}%
\pgfsetstrokecolor{currentstroke}%
\pgfsetstrokeopacity{0.000000}%
\pgfsetdash{}{0pt}%
\pgfpathmoveto{\pgfqpoint{8.369150in}{0.660000in}}%
\pgfpathlineto{\pgfqpoint{8.375350in}{0.660000in}}%
\pgfpathlineto{\pgfqpoint{8.375350in}{3.222011in}}%
\pgfpathlineto{\pgfqpoint{8.369150in}{3.222011in}}%
\pgfpathlineto{\pgfqpoint{8.369150in}{0.660000in}}%
\pgfpathclose%
\pgfusepath{fill}%
\end{pgfscope}%
\begin{pgfscope}%
\pgfpathrectangle{\pgfqpoint{1.250000in}{0.660000in}}{\pgfqpoint{7.750000in}{4.620000in}}%
\pgfusepath{clip}%
\pgfsetbuttcap%
\pgfsetmiterjoin%
\definecolor{currentfill}{rgb}{0.000000,0.000000,0.000000}%
\pgfsetfillcolor{currentfill}%
\pgfsetlinewidth{0.000000pt}%
\definecolor{currentstroke}{rgb}{0.000000,0.000000,0.000000}%
\pgfsetstrokecolor{currentstroke}%
\pgfsetstrokeopacity{0.000000}%
\pgfsetdash{}{0pt}%
\pgfpathmoveto{\pgfqpoint{8.376900in}{0.660000in}}%
\pgfpathlineto{\pgfqpoint{8.383100in}{0.660000in}}%
\pgfpathlineto{\pgfqpoint{8.383100in}{3.198097in}}%
\pgfpathlineto{\pgfqpoint{8.376900in}{3.198097in}}%
\pgfpathlineto{\pgfqpoint{8.376900in}{0.660000in}}%
\pgfpathclose%
\pgfusepath{fill}%
\end{pgfscope}%
\begin{pgfscope}%
\pgfpathrectangle{\pgfqpoint{1.250000in}{0.660000in}}{\pgfqpoint{7.750000in}{4.620000in}}%
\pgfusepath{clip}%
\pgfsetbuttcap%
\pgfsetmiterjoin%
\definecolor{currentfill}{rgb}{0.000000,0.000000,0.000000}%
\pgfsetfillcolor{currentfill}%
\pgfsetlinewidth{0.000000pt}%
\definecolor{currentstroke}{rgb}{0.000000,0.000000,0.000000}%
\pgfsetstrokecolor{currentstroke}%
\pgfsetstrokeopacity{0.000000}%
\pgfsetdash{}{0pt}%
\pgfpathmoveto{\pgfqpoint{8.384650in}{0.660000in}}%
\pgfpathlineto{\pgfqpoint{8.390850in}{0.660000in}}%
\pgfpathlineto{\pgfqpoint{8.390850in}{3.198043in}}%
\pgfpathlineto{\pgfqpoint{8.384650in}{3.198043in}}%
\pgfpathlineto{\pgfqpoint{8.384650in}{0.660000in}}%
\pgfpathclose%
\pgfusepath{fill}%
\end{pgfscope}%
\begin{pgfscope}%
\pgfpathrectangle{\pgfqpoint{1.250000in}{0.660000in}}{\pgfqpoint{7.750000in}{4.620000in}}%
\pgfusepath{clip}%
\pgfsetbuttcap%
\pgfsetmiterjoin%
\definecolor{currentfill}{rgb}{0.000000,0.000000,0.000000}%
\pgfsetfillcolor{currentfill}%
\pgfsetlinewidth{0.000000pt}%
\definecolor{currentstroke}{rgb}{0.000000,0.000000,0.000000}%
\pgfsetstrokecolor{currentstroke}%
\pgfsetstrokeopacity{0.000000}%
\pgfsetdash{}{0pt}%
\pgfpathmoveto{\pgfqpoint{8.392400in}{0.660000in}}%
\pgfpathlineto{\pgfqpoint{8.398600in}{0.660000in}}%
\pgfpathlineto{\pgfqpoint{8.398600in}{3.197962in}}%
\pgfpathlineto{\pgfqpoint{8.392400in}{3.197962in}}%
\pgfpathlineto{\pgfqpoint{8.392400in}{0.660000in}}%
\pgfpathclose%
\pgfusepath{fill}%
\end{pgfscope}%
\begin{pgfscope}%
\pgfpathrectangle{\pgfqpoint{1.250000in}{0.660000in}}{\pgfqpoint{7.750000in}{4.620000in}}%
\pgfusepath{clip}%
\pgfsetbuttcap%
\pgfsetmiterjoin%
\definecolor{currentfill}{rgb}{0.000000,0.000000,0.000000}%
\pgfsetfillcolor{currentfill}%
\pgfsetlinewidth{0.000000pt}%
\definecolor{currentstroke}{rgb}{0.000000,0.000000,0.000000}%
\pgfsetstrokecolor{currentstroke}%
\pgfsetstrokeopacity{0.000000}%
\pgfsetdash{}{0pt}%
\pgfpathmoveto{\pgfqpoint{8.400150in}{0.660000in}}%
\pgfpathlineto{\pgfqpoint{8.406350in}{0.660000in}}%
\pgfpathlineto{\pgfqpoint{8.406350in}{3.197852in}}%
\pgfpathlineto{\pgfqpoint{8.400150in}{3.197852in}}%
\pgfpathlineto{\pgfqpoint{8.400150in}{0.660000in}}%
\pgfpathclose%
\pgfusepath{fill}%
\end{pgfscope}%
\begin{pgfscope}%
\pgfpathrectangle{\pgfqpoint{1.250000in}{0.660000in}}{\pgfqpoint{7.750000in}{4.620000in}}%
\pgfusepath{clip}%
\pgfsetbuttcap%
\pgfsetmiterjoin%
\definecolor{currentfill}{rgb}{0.000000,0.000000,0.000000}%
\pgfsetfillcolor{currentfill}%
\pgfsetlinewidth{0.000000pt}%
\definecolor{currentstroke}{rgb}{0.000000,0.000000,0.000000}%
\pgfsetstrokecolor{currentstroke}%
\pgfsetstrokeopacity{0.000000}%
\pgfsetdash{}{0pt}%
\pgfpathmoveto{\pgfqpoint{8.407900in}{0.660000in}}%
\pgfpathlineto{\pgfqpoint{8.414100in}{0.660000in}}%
\pgfpathlineto{\pgfqpoint{8.414100in}{3.197713in}}%
\pgfpathlineto{\pgfqpoint{8.407900in}{3.197713in}}%
\pgfpathlineto{\pgfqpoint{8.407900in}{0.660000in}}%
\pgfpathclose%
\pgfusepath{fill}%
\end{pgfscope}%
\begin{pgfscope}%
\pgfpathrectangle{\pgfqpoint{1.250000in}{0.660000in}}{\pgfqpoint{7.750000in}{4.620000in}}%
\pgfusepath{clip}%
\pgfsetbuttcap%
\pgfsetmiterjoin%
\definecolor{currentfill}{rgb}{0.000000,0.000000,0.000000}%
\pgfsetfillcolor{currentfill}%
\pgfsetlinewidth{0.000000pt}%
\definecolor{currentstroke}{rgb}{0.000000,0.000000,0.000000}%
\pgfsetstrokecolor{currentstroke}%
\pgfsetstrokeopacity{0.000000}%
\pgfsetdash{}{0pt}%
\pgfpathmoveto{\pgfqpoint{8.415650in}{0.660000in}}%
\pgfpathlineto{\pgfqpoint{8.421850in}{0.660000in}}%
\pgfpathlineto{\pgfqpoint{8.421850in}{3.197544in}}%
\pgfpathlineto{\pgfqpoint{8.415650in}{3.197544in}}%
\pgfpathlineto{\pgfqpoint{8.415650in}{0.660000in}}%
\pgfpathclose%
\pgfusepath{fill}%
\end{pgfscope}%
\begin{pgfscope}%
\pgfpathrectangle{\pgfqpoint{1.250000in}{0.660000in}}{\pgfqpoint{7.750000in}{4.620000in}}%
\pgfusepath{clip}%
\pgfsetbuttcap%
\pgfsetmiterjoin%
\definecolor{currentfill}{rgb}{0.000000,0.000000,0.000000}%
\pgfsetfillcolor{currentfill}%
\pgfsetlinewidth{0.000000pt}%
\definecolor{currentstroke}{rgb}{0.000000,0.000000,0.000000}%
\pgfsetstrokecolor{currentstroke}%
\pgfsetstrokeopacity{0.000000}%
\pgfsetdash{}{0pt}%
\pgfpathmoveto{\pgfqpoint{8.423400in}{0.660000in}}%
\pgfpathlineto{\pgfqpoint{8.429600in}{0.660000in}}%
\pgfpathlineto{\pgfqpoint{8.429600in}{3.197346in}}%
\pgfpathlineto{\pgfqpoint{8.423400in}{3.197346in}}%
\pgfpathlineto{\pgfqpoint{8.423400in}{0.660000in}}%
\pgfpathclose%
\pgfusepath{fill}%
\end{pgfscope}%
\begin{pgfscope}%
\pgfpathrectangle{\pgfqpoint{1.250000in}{0.660000in}}{\pgfqpoint{7.750000in}{4.620000in}}%
\pgfusepath{clip}%
\pgfsetbuttcap%
\pgfsetmiterjoin%
\definecolor{currentfill}{rgb}{0.000000,0.000000,0.000000}%
\pgfsetfillcolor{currentfill}%
\pgfsetlinewidth{0.000000pt}%
\definecolor{currentstroke}{rgb}{0.000000,0.000000,0.000000}%
\pgfsetstrokecolor{currentstroke}%
\pgfsetstrokeopacity{0.000000}%
\pgfsetdash{}{0pt}%
\pgfpathmoveto{\pgfqpoint{8.431150in}{0.660000in}}%
\pgfpathlineto{\pgfqpoint{8.437350in}{0.660000in}}%
\pgfpathlineto{\pgfqpoint{8.437350in}{3.197117in}}%
\pgfpathlineto{\pgfqpoint{8.431150in}{3.197117in}}%
\pgfpathlineto{\pgfqpoint{8.431150in}{0.660000in}}%
\pgfpathclose%
\pgfusepath{fill}%
\end{pgfscope}%
\begin{pgfscope}%
\pgfpathrectangle{\pgfqpoint{1.250000in}{0.660000in}}{\pgfqpoint{7.750000in}{4.620000in}}%
\pgfusepath{clip}%
\pgfsetbuttcap%
\pgfsetmiterjoin%
\definecolor{currentfill}{rgb}{0.000000,0.000000,0.000000}%
\pgfsetfillcolor{currentfill}%
\pgfsetlinewidth{0.000000pt}%
\definecolor{currentstroke}{rgb}{0.000000,0.000000,0.000000}%
\pgfsetstrokecolor{currentstroke}%
\pgfsetstrokeopacity{0.000000}%
\pgfsetdash{}{0pt}%
\pgfpathmoveto{\pgfqpoint{8.438900in}{0.660000in}}%
\pgfpathlineto{\pgfqpoint{8.445100in}{0.660000in}}%
\pgfpathlineto{\pgfqpoint{8.445100in}{3.197119in}}%
\pgfpathlineto{\pgfqpoint{8.438900in}{3.197119in}}%
\pgfpathlineto{\pgfqpoint{8.438900in}{0.660000in}}%
\pgfpathclose%
\pgfusepath{fill}%
\end{pgfscope}%
\begin{pgfscope}%
\pgfpathrectangle{\pgfqpoint{1.250000in}{0.660000in}}{\pgfqpoint{7.750000in}{4.620000in}}%
\pgfusepath{clip}%
\pgfsetbuttcap%
\pgfsetmiterjoin%
\definecolor{currentfill}{rgb}{0.000000,0.000000,0.000000}%
\pgfsetfillcolor{currentfill}%
\pgfsetlinewidth{0.000000pt}%
\definecolor{currentstroke}{rgb}{0.000000,0.000000,0.000000}%
\pgfsetstrokecolor{currentstroke}%
\pgfsetstrokeopacity{0.000000}%
\pgfsetdash{}{0pt}%
\pgfpathmoveto{\pgfqpoint{8.446650in}{0.660000in}}%
\pgfpathlineto{\pgfqpoint{8.452850in}{0.660000in}}%
\pgfpathlineto{\pgfqpoint{8.452850in}{3.197154in}}%
\pgfpathlineto{\pgfqpoint{8.446650in}{3.197154in}}%
\pgfpathlineto{\pgfqpoint{8.446650in}{0.660000in}}%
\pgfpathclose%
\pgfusepath{fill}%
\end{pgfscope}%
\begin{pgfscope}%
\pgfpathrectangle{\pgfqpoint{1.250000in}{0.660000in}}{\pgfqpoint{7.750000in}{4.620000in}}%
\pgfusepath{clip}%
\pgfsetbuttcap%
\pgfsetmiterjoin%
\definecolor{currentfill}{rgb}{0.000000,0.000000,0.000000}%
\pgfsetfillcolor{currentfill}%
\pgfsetlinewidth{0.000000pt}%
\definecolor{currentstroke}{rgb}{0.000000,0.000000,0.000000}%
\pgfsetstrokecolor{currentstroke}%
\pgfsetstrokeopacity{0.000000}%
\pgfsetdash{}{0pt}%
\pgfpathmoveto{\pgfqpoint{8.454400in}{0.660000in}}%
\pgfpathlineto{\pgfqpoint{8.460600in}{0.660000in}}%
\pgfpathlineto{\pgfqpoint{8.460600in}{3.197161in}}%
\pgfpathlineto{\pgfqpoint{8.454400in}{3.197161in}}%
\pgfpathlineto{\pgfqpoint{8.454400in}{0.660000in}}%
\pgfpathclose%
\pgfusepath{fill}%
\end{pgfscope}%
\begin{pgfscope}%
\pgfpathrectangle{\pgfqpoint{1.250000in}{0.660000in}}{\pgfqpoint{7.750000in}{4.620000in}}%
\pgfusepath{clip}%
\pgfsetbuttcap%
\pgfsetmiterjoin%
\definecolor{currentfill}{rgb}{0.000000,0.000000,0.000000}%
\pgfsetfillcolor{currentfill}%
\pgfsetlinewidth{0.000000pt}%
\definecolor{currentstroke}{rgb}{0.000000,0.000000,0.000000}%
\pgfsetstrokecolor{currentstroke}%
\pgfsetstrokeopacity{0.000000}%
\pgfsetdash{}{0pt}%
\pgfpathmoveto{\pgfqpoint{8.462150in}{0.660000in}}%
\pgfpathlineto{\pgfqpoint{8.468350in}{0.660000in}}%
\pgfpathlineto{\pgfqpoint{8.468350in}{3.197139in}}%
\pgfpathlineto{\pgfqpoint{8.462150in}{3.197139in}}%
\pgfpathlineto{\pgfqpoint{8.462150in}{0.660000in}}%
\pgfpathclose%
\pgfusepath{fill}%
\end{pgfscope}%
\begin{pgfscope}%
\pgfpathrectangle{\pgfqpoint{1.250000in}{0.660000in}}{\pgfqpoint{7.750000in}{4.620000in}}%
\pgfusepath{clip}%
\pgfsetbuttcap%
\pgfsetmiterjoin%
\definecolor{currentfill}{rgb}{0.000000,0.000000,0.000000}%
\pgfsetfillcolor{currentfill}%
\pgfsetlinewidth{0.000000pt}%
\definecolor{currentstroke}{rgb}{0.000000,0.000000,0.000000}%
\pgfsetstrokecolor{currentstroke}%
\pgfsetstrokeopacity{0.000000}%
\pgfsetdash{}{0pt}%
\pgfpathmoveto{\pgfqpoint{8.469900in}{0.660000in}}%
\pgfpathlineto{\pgfqpoint{8.476100in}{0.660000in}}%
\pgfpathlineto{\pgfqpoint{8.476100in}{3.197088in}}%
\pgfpathlineto{\pgfqpoint{8.469900in}{3.197088in}}%
\pgfpathlineto{\pgfqpoint{8.469900in}{0.660000in}}%
\pgfpathclose%
\pgfusepath{fill}%
\end{pgfscope}%
\begin{pgfscope}%
\pgfpathrectangle{\pgfqpoint{1.250000in}{0.660000in}}{\pgfqpoint{7.750000in}{4.620000in}}%
\pgfusepath{clip}%
\pgfsetbuttcap%
\pgfsetmiterjoin%
\definecolor{currentfill}{rgb}{0.000000,0.000000,0.000000}%
\pgfsetfillcolor{currentfill}%
\pgfsetlinewidth{0.000000pt}%
\definecolor{currentstroke}{rgb}{0.000000,0.000000,0.000000}%
\pgfsetstrokecolor{currentstroke}%
\pgfsetstrokeopacity{0.000000}%
\pgfsetdash{}{0pt}%
\pgfpathmoveto{\pgfqpoint{8.477650in}{0.660000in}}%
\pgfpathlineto{\pgfqpoint{8.483850in}{0.660000in}}%
\pgfpathlineto{\pgfqpoint{8.483850in}{3.197007in}}%
\pgfpathlineto{\pgfqpoint{8.477650in}{3.197007in}}%
\pgfpathlineto{\pgfqpoint{8.477650in}{0.660000in}}%
\pgfpathclose%
\pgfusepath{fill}%
\end{pgfscope}%
\begin{pgfscope}%
\pgfpathrectangle{\pgfqpoint{1.250000in}{0.660000in}}{\pgfqpoint{7.750000in}{4.620000in}}%
\pgfusepath{clip}%
\pgfsetbuttcap%
\pgfsetmiterjoin%
\definecolor{currentfill}{rgb}{0.000000,0.000000,0.000000}%
\pgfsetfillcolor{currentfill}%
\pgfsetlinewidth{0.000000pt}%
\definecolor{currentstroke}{rgb}{0.000000,0.000000,0.000000}%
\pgfsetstrokecolor{currentstroke}%
\pgfsetstrokeopacity{0.000000}%
\pgfsetdash{}{0pt}%
\pgfpathmoveto{\pgfqpoint{8.485400in}{0.660000in}}%
\pgfpathlineto{\pgfqpoint{8.491600in}{0.660000in}}%
\pgfpathlineto{\pgfqpoint{8.491600in}{3.196896in}}%
\pgfpathlineto{\pgfqpoint{8.485400in}{3.196896in}}%
\pgfpathlineto{\pgfqpoint{8.485400in}{0.660000in}}%
\pgfpathclose%
\pgfusepath{fill}%
\end{pgfscope}%
\begin{pgfscope}%
\pgfpathrectangle{\pgfqpoint{1.250000in}{0.660000in}}{\pgfqpoint{7.750000in}{4.620000in}}%
\pgfusepath{clip}%
\pgfsetbuttcap%
\pgfsetmiterjoin%
\definecolor{currentfill}{rgb}{0.000000,0.000000,0.000000}%
\pgfsetfillcolor{currentfill}%
\pgfsetlinewidth{0.000000pt}%
\definecolor{currentstroke}{rgb}{0.000000,0.000000,0.000000}%
\pgfsetstrokecolor{currentstroke}%
\pgfsetstrokeopacity{0.000000}%
\pgfsetdash{}{0pt}%
\pgfpathmoveto{\pgfqpoint{8.493150in}{0.660000in}}%
\pgfpathlineto{\pgfqpoint{8.499350in}{0.660000in}}%
\pgfpathlineto{\pgfqpoint{8.499350in}{3.196754in}}%
\pgfpathlineto{\pgfqpoint{8.493150in}{3.196754in}}%
\pgfpathlineto{\pgfqpoint{8.493150in}{0.660000in}}%
\pgfpathclose%
\pgfusepath{fill}%
\end{pgfscope}%
\begin{pgfscope}%
\pgfpathrectangle{\pgfqpoint{1.250000in}{0.660000in}}{\pgfqpoint{7.750000in}{4.620000in}}%
\pgfusepath{clip}%
\pgfsetbuttcap%
\pgfsetmiterjoin%
\definecolor{currentfill}{rgb}{0.000000,0.000000,0.000000}%
\pgfsetfillcolor{currentfill}%
\pgfsetlinewidth{0.000000pt}%
\definecolor{currentstroke}{rgb}{0.000000,0.000000,0.000000}%
\pgfsetstrokecolor{currentstroke}%
\pgfsetstrokeopacity{0.000000}%
\pgfsetdash{}{0pt}%
\pgfpathmoveto{\pgfqpoint{8.500900in}{0.660000in}}%
\pgfpathlineto{\pgfqpoint{8.507100in}{0.660000in}}%
\pgfpathlineto{\pgfqpoint{8.507100in}{3.196581in}}%
\pgfpathlineto{\pgfqpoint{8.500900in}{3.196581in}}%
\pgfpathlineto{\pgfqpoint{8.500900in}{0.660000in}}%
\pgfpathclose%
\pgfusepath{fill}%
\end{pgfscope}%
\begin{pgfscope}%
\pgfpathrectangle{\pgfqpoint{1.250000in}{0.660000in}}{\pgfqpoint{7.750000in}{4.620000in}}%
\pgfusepath{clip}%
\pgfsetbuttcap%
\pgfsetmiterjoin%
\definecolor{currentfill}{rgb}{0.000000,0.000000,0.000000}%
\pgfsetfillcolor{currentfill}%
\pgfsetlinewidth{0.000000pt}%
\definecolor{currentstroke}{rgb}{0.000000,0.000000,0.000000}%
\pgfsetstrokecolor{currentstroke}%
\pgfsetstrokeopacity{0.000000}%
\pgfsetdash{}{0pt}%
\pgfpathmoveto{\pgfqpoint{8.508650in}{0.660000in}}%
\pgfpathlineto{\pgfqpoint{8.514850in}{0.660000in}}%
\pgfpathlineto{\pgfqpoint{8.514850in}{3.196375in}}%
\pgfpathlineto{\pgfqpoint{8.508650in}{3.196375in}}%
\pgfpathlineto{\pgfqpoint{8.508650in}{0.660000in}}%
\pgfpathclose%
\pgfusepath{fill}%
\end{pgfscope}%
\begin{pgfscope}%
\pgfpathrectangle{\pgfqpoint{1.250000in}{0.660000in}}{\pgfqpoint{7.750000in}{4.620000in}}%
\pgfusepath{clip}%
\pgfsetbuttcap%
\pgfsetmiterjoin%
\definecolor{currentfill}{rgb}{0.000000,0.000000,0.000000}%
\pgfsetfillcolor{currentfill}%
\pgfsetlinewidth{0.000000pt}%
\definecolor{currentstroke}{rgb}{0.000000,0.000000,0.000000}%
\pgfsetstrokecolor{currentstroke}%
\pgfsetstrokeopacity{0.000000}%
\pgfsetdash{}{0pt}%
\pgfpathmoveto{\pgfqpoint{8.516400in}{0.660000in}}%
\pgfpathlineto{\pgfqpoint{8.522600in}{0.660000in}}%
\pgfpathlineto{\pgfqpoint{8.522600in}{3.196148in}}%
\pgfpathlineto{\pgfqpoint{8.516400in}{3.196148in}}%
\pgfpathlineto{\pgfqpoint{8.516400in}{0.660000in}}%
\pgfpathclose%
\pgfusepath{fill}%
\end{pgfscope}%
\begin{pgfscope}%
\pgfpathrectangle{\pgfqpoint{1.250000in}{0.660000in}}{\pgfqpoint{7.750000in}{4.620000in}}%
\pgfusepath{clip}%
\pgfsetbuttcap%
\pgfsetmiterjoin%
\definecolor{currentfill}{rgb}{0.000000,0.000000,0.000000}%
\pgfsetfillcolor{currentfill}%
\pgfsetlinewidth{0.000000pt}%
\definecolor{currentstroke}{rgb}{0.000000,0.000000,0.000000}%
\pgfsetstrokecolor{currentstroke}%
\pgfsetstrokeopacity{0.000000}%
\pgfsetdash{}{0pt}%
\pgfpathmoveto{\pgfqpoint{8.524150in}{0.660000in}}%
\pgfpathlineto{\pgfqpoint{8.530350in}{0.660000in}}%
\pgfpathlineto{\pgfqpoint{8.530350in}{3.196209in}}%
\pgfpathlineto{\pgfqpoint{8.524150in}{3.196209in}}%
\pgfpathlineto{\pgfqpoint{8.524150in}{0.660000in}}%
\pgfpathclose%
\pgfusepath{fill}%
\end{pgfscope}%
\begin{pgfscope}%
\pgfpathrectangle{\pgfqpoint{1.250000in}{0.660000in}}{\pgfqpoint{7.750000in}{4.620000in}}%
\pgfusepath{clip}%
\pgfsetbuttcap%
\pgfsetmiterjoin%
\definecolor{currentfill}{rgb}{0.000000,0.000000,0.000000}%
\pgfsetfillcolor{currentfill}%
\pgfsetlinewidth{0.000000pt}%
\definecolor{currentstroke}{rgb}{0.000000,0.000000,0.000000}%
\pgfsetstrokecolor{currentstroke}%
\pgfsetstrokeopacity{0.000000}%
\pgfsetdash{}{0pt}%
\pgfpathmoveto{\pgfqpoint{8.531900in}{0.660000in}}%
\pgfpathlineto{\pgfqpoint{8.538100in}{0.660000in}}%
\pgfpathlineto{\pgfqpoint{8.538100in}{3.196241in}}%
\pgfpathlineto{\pgfqpoint{8.531900in}{3.196241in}}%
\pgfpathlineto{\pgfqpoint{8.531900in}{0.660000in}}%
\pgfpathclose%
\pgfusepath{fill}%
\end{pgfscope}%
\begin{pgfscope}%
\pgfpathrectangle{\pgfqpoint{1.250000in}{0.660000in}}{\pgfqpoint{7.750000in}{4.620000in}}%
\pgfusepath{clip}%
\pgfsetbuttcap%
\pgfsetmiterjoin%
\definecolor{currentfill}{rgb}{0.000000,0.000000,0.000000}%
\pgfsetfillcolor{currentfill}%
\pgfsetlinewidth{0.000000pt}%
\definecolor{currentstroke}{rgb}{0.000000,0.000000,0.000000}%
\pgfsetstrokecolor{currentstroke}%
\pgfsetstrokeopacity{0.000000}%
\pgfsetdash{}{0pt}%
\pgfpathmoveto{\pgfqpoint{8.539650in}{0.660000in}}%
\pgfpathlineto{\pgfqpoint{8.545850in}{0.660000in}}%
\pgfpathlineto{\pgfqpoint{8.545850in}{3.197221in}}%
\pgfpathlineto{\pgfqpoint{8.539650in}{3.197221in}}%
\pgfpathlineto{\pgfqpoint{8.539650in}{0.660000in}}%
\pgfpathclose%
\pgfusepath{fill}%
\end{pgfscope}%
\begin{pgfscope}%
\pgfpathrectangle{\pgfqpoint{1.250000in}{0.660000in}}{\pgfqpoint{7.750000in}{4.620000in}}%
\pgfusepath{clip}%
\pgfsetbuttcap%
\pgfsetmiterjoin%
\definecolor{currentfill}{rgb}{0.000000,0.000000,0.000000}%
\pgfsetfillcolor{currentfill}%
\pgfsetlinewidth{0.000000pt}%
\definecolor{currentstroke}{rgb}{0.000000,0.000000,0.000000}%
\pgfsetstrokecolor{currentstroke}%
\pgfsetstrokeopacity{0.000000}%
\pgfsetdash{}{0pt}%
\pgfpathmoveto{\pgfqpoint{8.547400in}{0.660000in}}%
\pgfpathlineto{\pgfqpoint{8.553600in}{0.660000in}}%
\pgfpathlineto{\pgfqpoint{8.553600in}{3.196214in}}%
\pgfpathlineto{\pgfqpoint{8.547400in}{3.196214in}}%
\pgfpathlineto{\pgfqpoint{8.547400in}{0.660000in}}%
\pgfpathclose%
\pgfusepath{fill}%
\end{pgfscope}%
\begin{pgfscope}%
\pgfpathrectangle{\pgfqpoint{1.250000in}{0.660000in}}{\pgfqpoint{7.750000in}{4.620000in}}%
\pgfusepath{clip}%
\pgfsetbuttcap%
\pgfsetmiterjoin%
\definecolor{currentfill}{rgb}{0.000000,0.000000,0.000000}%
\pgfsetfillcolor{currentfill}%
\pgfsetlinewidth{0.000000pt}%
\definecolor{currentstroke}{rgb}{0.000000,0.000000,0.000000}%
\pgfsetstrokecolor{currentstroke}%
\pgfsetstrokeopacity{0.000000}%
\pgfsetdash{}{0pt}%
\pgfpathmoveto{\pgfqpoint{8.555150in}{0.660000in}}%
\pgfpathlineto{\pgfqpoint{8.561350in}{0.660000in}}%
\pgfpathlineto{\pgfqpoint{8.561350in}{3.196153in}}%
\pgfpathlineto{\pgfqpoint{8.555150in}{3.196153in}}%
\pgfpathlineto{\pgfqpoint{8.555150in}{0.660000in}}%
\pgfpathclose%
\pgfusepath{fill}%
\end{pgfscope}%
\begin{pgfscope}%
\pgfpathrectangle{\pgfqpoint{1.250000in}{0.660000in}}{\pgfqpoint{7.750000in}{4.620000in}}%
\pgfusepath{clip}%
\pgfsetbuttcap%
\pgfsetmiterjoin%
\definecolor{currentfill}{rgb}{0.000000,0.000000,0.000000}%
\pgfsetfillcolor{currentfill}%
\pgfsetlinewidth{0.000000pt}%
\definecolor{currentstroke}{rgb}{0.000000,0.000000,0.000000}%
\pgfsetstrokecolor{currentstroke}%
\pgfsetstrokeopacity{0.000000}%
\pgfsetdash{}{0pt}%
\pgfpathmoveto{\pgfqpoint{8.562900in}{0.660000in}}%
\pgfpathlineto{\pgfqpoint{8.569100in}{0.660000in}}%
\pgfpathlineto{\pgfqpoint{8.569100in}{3.196061in}}%
\pgfpathlineto{\pgfqpoint{8.562900in}{3.196061in}}%
\pgfpathlineto{\pgfqpoint{8.562900in}{0.660000in}}%
\pgfpathclose%
\pgfusepath{fill}%
\end{pgfscope}%
\begin{pgfscope}%
\pgfpathrectangle{\pgfqpoint{1.250000in}{0.660000in}}{\pgfqpoint{7.750000in}{4.620000in}}%
\pgfusepath{clip}%
\pgfsetbuttcap%
\pgfsetmiterjoin%
\definecolor{currentfill}{rgb}{0.000000,0.000000,0.000000}%
\pgfsetfillcolor{currentfill}%
\pgfsetlinewidth{0.000000pt}%
\definecolor{currentstroke}{rgb}{0.000000,0.000000,0.000000}%
\pgfsetstrokecolor{currentstroke}%
\pgfsetstrokeopacity{0.000000}%
\pgfsetdash{}{0pt}%
\pgfpathmoveto{\pgfqpoint{8.570650in}{0.660000in}}%
\pgfpathlineto{\pgfqpoint{8.576850in}{0.660000in}}%
\pgfpathlineto{\pgfqpoint{8.576850in}{3.195937in}}%
\pgfpathlineto{\pgfqpoint{8.570650in}{3.195937in}}%
\pgfpathlineto{\pgfqpoint{8.570650in}{0.660000in}}%
\pgfpathclose%
\pgfusepath{fill}%
\end{pgfscope}%
\begin{pgfscope}%
\pgfpathrectangle{\pgfqpoint{1.250000in}{0.660000in}}{\pgfqpoint{7.750000in}{4.620000in}}%
\pgfusepath{clip}%
\pgfsetbuttcap%
\pgfsetmiterjoin%
\definecolor{currentfill}{rgb}{0.000000,0.000000,0.000000}%
\pgfsetfillcolor{currentfill}%
\pgfsetlinewidth{0.000000pt}%
\definecolor{currentstroke}{rgb}{0.000000,0.000000,0.000000}%
\pgfsetstrokecolor{currentstroke}%
\pgfsetstrokeopacity{0.000000}%
\pgfsetdash{}{0pt}%
\pgfpathmoveto{\pgfqpoint{8.578400in}{0.660000in}}%
\pgfpathlineto{\pgfqpoint{8.584600in}{0.660000in}}%
\pgfpathlineto{\pgfqpoint{8.584600in}{3.195780in}}%
\pgfpathlineto{\pgfqpoint{8.578400in}{3.195780in}}%
\pgfpathlineto{\pgfqpoint{8.578400in}{0.660000in}}%
\pgfpathclose%
\pgfusepath{fill}%
\end{pgfscope}%
\begin{pgfscope}%
\pgfpathrectangle{\pgfqpoint{1.250000in}{0.660000in}}{\pgfqpoint{7.750000in}{4.620000in}}%
\pgfusepath{clip}%
\pgfsetbuttcap%
\pgfsetmiterjoin%
\definecolor{currentfill}{rgb}{0.000000,0.000000,0.000000}%
\pgfsetfillcolor{currentfill}%
\pgfsetlinewidth{0.000000pt}%
\definecolor{currentstroke}{rgb}{0.000000,0.000000,0.000000}%
\pgfsetstrokecolor{currentstroke}%
\pgfsetstrokeopacity{0.000000}%
\pgfsetdash{}{0pt}%
\pgfpathmoveto{\pgfqpoint{8.586150in}{0.660000in}}%
\pgfpathlineto{\pgfqpoint{8.592350in}{0.660000in}}%
\pgfpathlineto{\pgfqpoint{8.592350in}{3.195589in}}%
\pgfpathlineto{\pgfqpoint{8.586150in}{3.195589in}}%
\pgfpathlineto{\pgfqpoint{8.586150in}{0.660000in}}%
\pgfpathclose%
\pgfusepath{fill}%
\end{pgfscope}%
\begin{pgfscope}%
\pgfpathrectangle{\pgfqpoint{1.250000in}{0.660000in}}{\pgfqpoint{7.750000in}{4.620000in}}%
\pgfusepath{clip}%
\pgfsetbuttcap%
\pgfsetmiterjoin%
\definecolor{currentfill}{rgb}{0.000000,0.000000,0.000000}%
\pgfsetfillcolor{currentfill}%
\pgfsetlinewidth{0.000000pt}%
\definecolor{currentstroke}{rgb}{0.000000,0.000000,0.000000}%
\pgfsetstrokecolor{currentstroke}%
\pgfsetstrokeopacity{0.000000}%
\pgfsetdash{}{0pt}%
\pgfpathmoveto{\pgfqpoint{8.593900in}{0.660000in}}%
\pgfpathlineto{\pgfqpoint{8.600100in}{0.660000in}}%
\pgfpathlineto{\pgfqpoint{8.600100in}{3.195364in}}%
\pgfpathlineto{\pgfqpoint{8.593900in}{3.195364in}}%
\pgfpathlineto{\pgfqpoint{8.593900in}{0.660000in}}%
\pgfpathclose%
\pgfusepath{fill}%
\end{pgfscope}%
\begin{pgfscope}%
\pgfpathrectangle{\pgfqpoint{1.250000in}{0.660000in}}{\pgfqpoint{7.750000in}{4.620000in}}%
\pgfusepath{clip}%
\pgfsetbuttcap%
\pgfsetmiterjoin%
\definecolor{currentfill}{rgb}{0.000000,0.000000,0.000000}%
\pgfsetfillcolor{currentfill}%
\pgfsetlinewidth{0.000000pt}%
\definecolor{currentstroke}{rgb}{0.000000,0.000000,0.000000}%
\pgfsetstrokecolor{currentstroke}%
\pgfsetstrokeopacity{0.000000}%
\pgfsetdash{}{0pt}%
\pgfpathmoveto{\pgfqpoint{8.601650in}{0.660000in}}%
\pgfpathlineto{\pgfqpoint{8.607850in}{0.660000in}}%
\pgfpathlineto{\pgfqpoint{8.607850in}{3.195308in}}%
\pgfpathlineto{\pgfqpoint{8.601650in}{3.195308in}}%
\pgfpathlineto{\pgfqpoint{8.601650in}{0.660000in}}%
\pgfpathclose%
\pgfusepath{fill}%
\end{pgfscope}%
\begin{pgfscope}%
\pgfpathrectangle{\pgfqpoint{1.250000in}{0.660000in}}{\pgfqpoint{7.750000in}{4.620000in}}%
\pgfusepath{clip}%
\pgfsetbuttcap%
\pgfsetmiterjoin%
\definecolor{currentfill}{rgb}{0.000000,0.000000,0.000000}%
\pgfsetfillcolor{currentfill}%
\pgfsetlinewidth{0.000000pt}%
\definecolor{currentstroke}{rgb}{0.000000,0.000000,0.000000}%
\pgfsetstrokecolor{currentstroke}%
\pgfsetstrokeopacity{0.000000}%
\pgfsetdash{}{0pt}%
\pgfpathmoveto{\pgfqpoint{8.609400in}{0.660000in}}%
\pgfpathlineto{\pgfqpoint{8.615600in}{0.660000in}}%
\pgfpathlineto{\pgfqpoint{8.615600in}{3.195356in}}%
\pgfpathlineto{\pgfqpoint{8.609400in}{3.195356in}}%
\pgfpathlineto{\pgfqpoint{8.609400in}{0.660000in}}%
\pgfpathclose%
\pgfusepath{fill}%
\end{pgfscope}%
\begin{pgfscope}%
\pgfpathrectangle{\pgfqpoint{1.250000in}{0.660000in}}{\pgfqpoint{7.750000in}{4.620000in}}%
\pgfusepath{clip}%
\pgfsetbuttcap%
\pgfsetmiterjoin%
\definecolor{currentfill}{rgb}{0.000000,0.000000,0.000000}%
\pgfsetfillcolor{currentfill}%
\pgfsetlinewidth{0.000000pt}%
\definecolor{currentstroke}{rgb}{0.000000,0.000000,0.000000}%
\pgfsetstrokecolor{currentstroke}%
\pgfsetstrokeopacity{0.000000}%
\pgfsetdash{}{0pt}%
\pgfpathmoveto{\pgfqpoint{8.617150in}{0.660000in}}%
\pgfpathlineto{\pgfqpoint{8.623350in}{0.660000in}}%
\pgfpathlineto{\pgfqpoint{8.623350in}{3.195372in}}%
\pgfpathlineto{\pgfqpoint{8.617150in}{3.195372in}}%
\pgfpathlineto{\pgfqpoint{8.617150in}{0.660000in}}%
\pgfpathclose%
\pgfusepath{fill}%
\end{pgfscope}%
\begin{pgfscope}%
\pgfpathrectangle{\pgfqpoint{1.250000in}{0.660000in}}{\pgfqpoint{7.750000in}{4.620000in}}%
\pgfusepath{clip}%
\pgfsetbuttcap%
\pgfsetmiterjoin%
\definecolor{currentfill}{rgb}{0.000000,0.000000,0.000000}%
\pgfsetfillcolor{currentfill}%
\pgfsetlinewidth{0.000000pt}%
\definecolor{currentstroke}{rgb}{0.000000,0.000000,0.000000}%
\pgfsetstrokecolor{currentstroke}%
\pgfsetstrokeopacity{0.000000}%
\pgfsetdash{}{0pt}%
\pgfpathmoveto{\pgfqpoint{8.624900in}{0.660000in}}%
\pgfpathlineto{\pgfqpoint{8.631100in}{0.660000in}}%
\pgfpathlineto{\pgfqpoint{8.631100in}{3.195356in}}%
\pgfpathlineto{\pgfqpoint{8.624900in}{3.195356in}}%
\pgfpathlineto{\pgfqpoint{8.624900in}{0.660000in}}%
\pgfpathclose%
\pgfusepath{fill}%
\end{pgfscope}%
\begin{pgfscope}%
\pgfpathrectangle{\pgfqpoint{1.250000in}{0.660000in}}{\pgfqpoint{7.750000in}{4.620000in}}%
\pgfusepath{clip}%
\pgfsetbuttcap%
\pgfsetmiterjoin%
\definecolor{currentfill}{rgb}{0.000000,0.000000,0.000000}%
\pgfsetfillcolor{currentfill}%
\pgfsetlinewidth{0.000000pt}%
\definecolor{currentstroke}{rgb}{0.000000,0.000000,0.000000}%
\pgfsetstrokecolor{currentstroke}%
\pgfsetstrokeopacity{0.000000}%
\pgfsetdash{}{0pt}%
\pgfpathmoveto{\pgfqpoint{8.632650in}{0.660000in}}%
\pgfpathlineto{\pgfqpoint{8.638850in}{0.660000in}}%
\pgfpathlineto{\pgfqpoint{8.638850in}{3.195308in}}%
\pgfpathlineto{\pgfqpoint{8.632650in}{3.195308in}}%
\pgfpathlineto{\pgfqpoint{8.632650in}{0.660000in}}%
\pgfpathclose%
\pgfusepath{fill}%
\end{pgfscope}%
\begin{pgfscope}%
\pgfpathrectangle{\pgfqpoint{1.250000in}{0.660000in}}{\pgfqpoint{7.750000in}{4.620000in}}%
\pgfusepath{clip}%
\pgfsetbuttcap%
\pgfsetmiterjoin%
\definecolor{currentfill}{rgb}{0.000000,0.000000,0.000000}%
\pgfsetfillcolor{currentfill}%
\pgfsetlinewidth{0.000000pt}%
\definecolor{currentstroke}{rgb}{0.000000,0.000000,0.000000}%
\pgfsetstrokecolor{currentstroke}%
\pgfsetstrokeopacity{0.000000}%
\pgfsetdash{}{0pt}%
\pgfpathmoveto{\pgfqpoint{8.640400in}{0.660000in}}%
\pgfpathlineto{\pgfqpoint{8.646600in}{0.660000in}}%
\pgfpathlineto{\pgfqpoint{8.646600in}{3.197265in}}%
\pgfpathlineto{\pgfqpoint{8.640400in}{3.197265in}}%
\pgfpathlineto{\pgfqpoint{8.640400in}{0.660000in}}%
\pgfpathclose%
\pgfusepath{fill}%
\end{pgfscope}%
\begin{pgfscope}%
\pgfpathrectangle{\pgfqpoint{1.250000in}{0.660000in}}{\pgfqpoint{7.750000in}{4.620000in}}%
\pgfusepath{clip}%
\pgfsetbuttcap%
\pgfsetmiterjoin%
\definecolor{currentfill}{rgb}{0.000000,0.000000,0.000000}%
\pgfsetfillcolor{currentfill}%
\pgfsetlinewidth{0.000000pt}%
\definecolor{currentstroke}{rgb}{0.000000,0.000000,0.000000}%
\pgfsetstrokecolor{currentstroke}%
\pgfsetstrokeopacity{0.000000}%
\pgfsetdash{}{0pt}%
\pgfpathmoveto{\pgfqpoint{8.648150in}{0.660000in}}%
\pgfpathlineto{\pgfqpoint{8.654350in}{0.660000in}}%
\pgfpathlineto{\pgfqpoint{8.654350in}{3.195110in}}%
\pgfpathlineto{\pgfqpoint{8.648150in}{3.195110in}}%
\pgfpathlineto{\pgfqpoint{8.648150in}{0.660000in}}%
\pgfpathclose%
\pgfusepath{fill}%
\end{pgfscope}%
\begin{pgfscope}%
\pgfpathrectangle{\pgfqpoint{1.250000in}{0.660000in}}{\pgfqpoint{7.750000in}{4.620000in}}%
\pgfusepath{clip}%
\pgfsetbuttcap%
\pgfsetmiterjoin%
\definecolor{currentfill}{rgb}{0.000000,0.000000,0.000000}%
\pgfsetfillcolor{currentfill}%
\pgfsetlinewidth{0.000000pt}%
\definecolor{currentstroke}{rgb}{0.000000,0.000000,0.000000}%
\pgfsetstrokecolor{currentstroke}%
\pgfsetstrokeopacity{0.000000}%
\pgfsetdash{}{0pt}%
\pgfpathmoveto{\pgfqpoint{8.655900in}{0.660000in}}%
\pgfpathlineto{\pgfqpoint{8.662100in}{0.660000in}}%
\pgfpathlineto{\pgfqpoint{8.662100in}{3.194959in}}%
\pgfpathlineto{\pgfqpoint{8.655900in}{3.194959in}}%
\pgfpathlineto{\pgfqpoint{8.655900in}{0.660000in}}%
\pgfpathclose%
\pgfusepath{fill}%
\end{pgfscope}%
\begin{pgfscope}%
\pgfpathrectangle{\pgfqpoint{1.250000in}{0.660000in}}{\pgfqpoint{7.750000in}{4.620000in}}%
\pgfusepath{clip}%
\pgfsetbuttcap%
\pgfsetmiterjoin%
\definecolor{currentfill}{rgb}{0.000000,0.000000,0.000000}%
\pgfsetfillcolor{currentfill}%
\pgfsetlinewidth{0.000000pt}%
\definecolor{currentstroke}{rgb}{0.000000,0.000000,0.000000}%
\pgfsetstrokecolor{currentstroke}%
\pgfsetstrokeopacity{0.000000}%
\pgfsetdash{}{0pt}%
\pgfpathmoveto{\pgfqpoint{8.663650in}{0.660000in}}%
\pgfpathlineto{\pgfqpoint{8.669850in}{0.660000in}}%
\pgfpathlineto{\pgfqpoint{8.669850in}{3.194773in}}%
\pgfpathlineto{\pgfqpoint{8.663650in}{3.194773in}}%
\pgfpathlineto{\pgfqpoint{8.663650in}{0.660000in}}%
\pgfpathclose%
\pgfusepath{fill}%
\end{pgfscope}%
\begin{pgfscope}%
\pgfpathrectangle{\pgfqpoint{1.250000in}{0.660000in}}{\pgfqpoint{7.750000in}{4.620000in}}%
\pgfusepath{clip}%
\pgfsetbuttcap%
\pgfsetmiterjoin%
\definecolor{currentfill}{rgb}{0.000000,0.000000,0.000000}%
\pgfsetfillcolor{currentfill}%
\pgfsetlinewidth{0.000000pt}%
\definecolor{currentstroke}{rgb}{0.000000,0.000000,0.000000}%
\pgfsetstrokecolor{currentstroke}%
\pgfsetstrokeopacity{0.000000}%
\pgfsetdash{}{0pt}%
\pgfpathmoveto{\pgfqpoint{8.671400in}{0.660000in}}%
\pgfpathlineto{\pgfqpoint{8.677600in}{0.660000in}}%
\pgfpathlineto{\pgfqpoint{8.677600in}{3.194552in}}%
\pgfpathlineto{\pgfqpoint{8.671400in}{3.194552in}}%
\pgfpathlineto{\pgfqpoint{8.671400in}{0.660000in}}%
\pgfpathclose%
\pgfusepath{fill}%
\end{pgfscope}%
\begin{pgfscope}%
\pgfpathrectangle{\pgfqpoint{1.250000in}{0.660000in}}{\pgfqpoint{7.750000in}{4.620000in}}%
\pgfusepath{clip}%
\pgfsetbuttcap%
\pgfsetmiterjoin%
\definecolor{currentfill}{rgb}{0.000000,0.000000,0.000000}%
\pgfsetfillcolor{currentfill}%
\pgfsetlinewidth{0.000000pt}%
\definecolor{currentstroke}{rgb}{0.000000,0.000000,0.000000}%
\pgfsetstrokecolor{currentstroke}%
\pgfsetstrokeopacity{0.000000}%
\pgfsetdash{}{0pt}%
\pgfpathmoveto{\pgfqpoint{8.679150in}{0.660000in}}%
\pgfpathlineto{\pgfqpoint{8.685350in}{0.660000in}}%
\pgfpathlineto{\pgfqpoint{8.685350in}{3.194461in}}%
\pgfpathlineto{\pgfqpoint{8.679150in}{3.194461in}}%
\pgfpathlineto{\pgfqpoint{8.679150in}{0.660000in}}%
\pgfpathclose%
\pgfusepath{fill}%
\end{pgfscope}%
\begin{pgfscope}%
\pgfpathrectangle{\pgfqpoint{1.250000in}{0.660000in}}{\pgfqpoint{7.750000in}{4.620000in}}%
\pgfusepath{clip}%
\pgfsetbuttcap%
\pgfsetmiterjoin%
\definecolor{currentfill}{rgb}{0.000000,0.000000,0.000000}%
\pgfsetfillcolor{currentfill}%
\pgfsetlinewidth{0.000000pt}%
\definecolor{currentstroke}{rgb}{0.000000,0.000000,0.000000}%
\pgfsetstrokecolor{currentstroke}%
\pgfsetstrokeopacity{0.000000}%
\pgfsetdash{}{0pt}%
\pgfpathmoveto{\pgfqpoint{8.686900in}{0.660000in}}%
\pgfpathlineto{\pgfqpoint{8.693100in}{0.660000in}}%
\pgfpathlineto{\pgfqpoint{8.693100in}{3.194516in}}%
\pgfpathlineto{\pgfqpoint{8.686900in}{3.194516in}}%
\pgfpathlineto{\pgfqpoint{8.686900in}{0.660000in}}%
\pgfpathclose%
\pgfusepath{fill}%
\end{pgfscope}%
\begin{pgfscope}%
\pgfpathrectangle{\pgfqpoint{1.250000in}{0.660000in}}{\pgfqpoint{7.750000in}{4.620000in}}%
\pgfusepath{clip}%
\pgfsetbuttcap%
\pgfsetmiterjoin%
\definecolor{currentfill}{rgb}{0.000000,0.000000,0.000000}%
\pgfsetfillcolor{currentfill}%
\pgfsetlinewidth{0.000000pt}%
\definecolor{currentstroke}{rgb}{0.000000,0.000000,0.000000}%
\pgfsetstrokecolor{currentstroke}%
\pgfsetstrokeopacity{0.000000}%
\pgfsetdash{}{0pt}%
\pgfpathmoveto{\pgfqpoint{8.694650in}{0.660000in}}%
\pgfpathlineto{\pgfqpoint{8.700850in}{0.660000in}}%
\pgfpathlineto{\pgfqpoint{8.700850in}{3.194537in}}%
\pgfpathlineto{\pgfqpoint{8.694650in}{3.194537in}}%
\pgfpathlineto{\pgfqpoint{8.694650in}{0.660000in}}%
\pgfpathclose%
\pgfusepath{fill}%
\end{pgfscope}%
\begin{pgfscope}%
\pgfpathrectangle{\pgfqpoint{1.250000in}{0.660000in}}{\pgfqpoint{7.750000in}{4.620000in}}%
\pgfusepath{clip}%
\pgfsetbuttcap%
\pgfsetmiterjoin%
\definecolor{currentfill}{rgb}{0.000000,0.000000,0.000000}%
\pgfsetfillcolor{currentfill}%
\pgfsetlinewidth{0.000000pt}%
\definecolor{currentstroke}{rgb}{0.000000,0.000000,0.000000}%
\pgfsetstrokecolor{currentstroke}%
\pgfsetstrokeopacity{0.000000}%
\pgfsetdash{}{0pt}%
\pgfpathmoveto{\pgfqpoint{8.702400in}{0.660000in}}%
\pgfpathlineto{\pgfqpoint{8.708600in}{0.660000in}}%
\pgfpathlineto{\pgfqpoint{8.708600in}{3.194525in}}%
\pgfpathlineto{\pgfqpoint{8.702400in}{3.194525in}}%
\pgfpathlineto{\pgfqpoint{8.702400in}{0.660000in}}%
\pgfpathclose%
\pgfusepath{fill}%
\end{pgfscope}%
\begin{pgfscope}%
\pgfpathrectangle{\pgfqpoint{1.250000in}{0.660000in}}{\pgfqpoint{7.750000in}{4.620000in}}%
\pgfusepath{clip}%
\pgfsetbuttcap%
\pgfsetmiterjoin%
\definecolor{currentfill}{rgb}{0.000000,0.000000,0.000000}%
\pgfsetfillcolor{currentfill}%
\pgfsetlinewidth{0.000000pt}%
\definecolor{currentstroke}{rgb}{0.000000,0.000000,0.000000}%
\pgfsetstrokecolor{currentstroke}%
\pgfsetstrokeopacity{0.000000}%
\pgfsetdash{}{0pt}%
\pgfpathmoveto{\pgfqpoint{8.710150in}{0.660000in}}%
\pgfpathlineto{\pgfqpoint{8.716350in}{0.660000in}}%
\pgfpathlineto{\pgfqpoint{8.716350in}{3.194478in}}%
\pgfpathlineto{\pgfqpoint{8.710150in}{3.194478in}}%
\pgfpathlineto{\pgfqpoint{8.710150in}{0.660000in}}%
\pgfpathclose%
\pgfusepath{fill}%
\end{pgfscope}%
\begin{pgfscope}%
\pgfpathrectangle{\pgfqpoint{1.250000in}{0.660000in}}{\pgfqpoint{7.750000in}{4.620000in}}%
\pgfusepath{clip}%
\pgfsetbuttcap%
\pgfsetmiterjoin%
\definecolor{currentfill}{rgb}{0.000000,0.000000,0.000000}%
\pgfsetfillcolor{currentfill}%
\pgfsetlinewidth{0.000000pt}%
\definecolor{currentstroke}{rgb}{0.000000,0.000000,0.000000}%
\pgfsetstrokecolor{currentstroke}%
\pgfsetstrokeopacity{0.000000}%
\pgfsetdash{}{0pt}%
\pgfpathmoveto{\pgfqpoint{8.717900in}{0.660000in}}%
\pgfpathlineto{\pgfqpoint{8.724100in}{0.660000in}}%
\pgfpathlineto{\pgfqpoint{8.724100in}{3.194396in}}%
\pgfpathlineto{\pgfqpoint{8.717900in}{3.194396in}}%
\pgfpathlineto{\pgfqpoint{8.717900in}{0.660000in}}%
\pgfpathclose%
\pgfusepath{fill}%
\end{pgfscope}%
\begin{pgfscope}%
\pgfpathrectangle{\pgfqpoint{1.250000in}{0.660000in}}{\pgfqpoint{7.750000in}{4.620000in}}%
\pgfusepath{clip}%
\pgfsetbuttcap%
\pgfsetmiterjoin%
\definecolor{currentfill}{rgb}{0.000000,0.000000,0.000000}%
\pgfsetfillcolor{currentfill}%
\pgfsetlinewidth{0.000000pt}%
\definecolor{currentstroke}{rgb}{0.000000,0.000000,0.000000}%
\pgfsetstrokecolor{currentstroke}%
\pgfsetstrokeopacity{0.000000}%
\pgfsetdash{}{0pt}%
\pgfpathmoveto{\pgfqpoint{8.725650in}{0.660000in}}%
\pgfpathlineto{\pgfqpoint{8.731850in}{0.660000in}}%
\pgfpathlineto{\pgfqpoint{8.731850in}{3.194279in}}%
\pgfpathlineto{\pgfqpoint{8.725650in}{3.194279in}}%
\pgfpathlineto{\pgfqpoint{8.725650in}{0.660000in}}%
\pgfpathclose%
\pgfusepath{fill}%
\end{pgfscope}%
\begin{pgfscope}%
\pgfpathrectangle{\pgfqpoint{1.250000in}{0.660000in}}{\pgfqpoint{7.750000in}{4.620000in}}%
\pgfusepath{clip}%
\pgfsetbuttcap%
\pgfsetmiterjoin%
\definecolor{currentfill}{rgb}{0.000000,0.000000,0.000000}%
\pgfsetfillcolor{currentfill}%
\pgfsetlinewidth{0.000000pt}%
\definecolor{currentstroke}{rgb}{0.000000,0.000000,0.000000}%
\pgfsetstrokecolor{currentstroke}%
\pgfsetstrokeopacity{0.000000}%
\pgfsetdash{}{0pt}%
\pgfpathmoveto{\pgfqpoint{8.733400in}{0.660000in}}%
\pgfpathlineto{\pgfqpoint{8.739600in}{0.660000in}}%
\pgfpathlineto{\pgfqpoint{8.739600in}{3.194125in}}%
\pgfpathlineto{\pgfqpoint{8.733400in}{3.194125in}}%
\pgfpathlineto{\pgfqpoint{8.733400in}{0.660000in}}%
\pgfpathclose%
\pgfusepath{fill}%
\end{pgfscope}%
\begin{pgfscope}%
\pgfpathrectangle{\pgfqpoint{1.250000in}{0.660000in}}{\pgfqpoint{7.750000in}{4.620000in}}%
\pgfusepath{clip}%
\pgfsetbuttcap%
\pgfsetmiterjoin%
\definecolor{currentfill}{rgb}{0.000000,0.000000,0.000000}%
\pgfsetfillcolor{currentfill}%
\pgfsetlinewidth{0.000000pt}%
\definecolor{currentstroke}{rgb}{0.000000,0.000000,0.000000}%
\pgfsetstrokecolor{currentstroke}%
\pgfsetstrokeopacity{0.000000}%
\pgfsetdash{}{0pt}%
\pgfpathmoveto{\pgfqpoint{8.741150in}{0.660000in}}%
\pgfpathlineto{\pgfqpoint{8.747350in}{0.660000in}}%
\pgfpathlineto{\pgfqpoint{8.747350in}{3.193934in}}%
\pgfpathlineto{\pgfqpoint{8.741150in}{3.193934in}}%
\pgfpathlineto{\pgfqpoint{8.741150in}{0.660000in}}%
\pgfpathclose%
\pgfusepath{fill}%
\end{pgfscope}%
\begin{pgfscope}%
\pgfpathrectangle{\pgfqpoint{1.250000in}{0.660000in}}{\pgfqpoint{7.750000in}{4.620000in}}%
\pgfusepath{clip}%
\pgfsetbuttcap%
\pgfsetmiterjoin%
\definecolor{currentfill}{rgb}{0.000000,0.000000,0.000000}%
\pgfsetfillcolor{currentfill}%
\pgfsetlinewidth{0.000000pt}%
\definecolor{currentstroke}{rgb}{0.000000,0.000000,0.000000}%
\pgfsetstrokecolor{currentstroke}%
\pgfsetstrokeopacity{0.000000}%
\pgfsetdash{}{0pt}%
\pgfpathmoveto{\pgfqpoint{8.748900in}{0.660000in}}%
\pgfpathlineto{\pgfqpoint{8.755100in}{0.660000in}}%
\pgfpathlineto{\pgfqpoint{8.755100in}{3.193705in}}%
\pgfpathlineto{\pgfqpoint{8.748900in}{3.193705in}}%
\pgfpathlineto{\pgfqpoint{8.748900in}{0.660000in}}%
\pgfpathclose%
\pgfusepath{fill}%
\end{pgfscope}%
\begin{pgfscope}%
\pgfpathrectangle{\pgfqpoint{1.250000in}{0.660000in}}{\pgfqpoint{7.750000in}{4.620000in}}%
\pgfusepath{clip}%
\pgfsetbuttcap%
\pgfsetmiterjoin%
\definecolor{currentfill}{rgb}{0.000000,0.000000,0.000000}%
\pgfsetfillcolor{currentfill}%
\pgfsetlinewidth{0.000000pt}%
\definecolor{currentstroke}{rgb}{0.000000,0.000000,0.000000}%
\pgfsetstrokecolor{currentstroke}%
\pgfsetstrokeopacity{0.000000}%
\pgfsetdash{}{0pt}%
\pgfpathmoveto{\pgfqpoint{8.756650in}{0.660000in}}%
\pgfpathlineto{\pgfqpoint{8.762850in}{0.660000in}}%
\pgfpathlineto{\pgfqpoint{8.762850in}{3.193669in}}%
\pgfpathlineto{\pgfqpoint{8.756650in}{3.193669in}}%
\pgfpathlineto{\pgfqpoint{8.756650in}{0.660000in}}%
\pgfpathclose%
\pgfusepath{fill}%
\end{pgfscope}%
\begin{pgfscope}%
\pgfpathrectangle{\pgfqpoint{1.250000in}{0.660000in}}{\pgfqpoint{7.750000in}{4.620000in}}%
\pgfusepath{clip}%
\pgfsetbuttcap%
\pgfsetmiterjoin%
\definecolor{currentfill}{rgb}{0.000000,0.000000,0.000000}%
\pgfsetfillcolor{currentfill}%
\pgfsetlinewidth{0.000000pt}%
\definecolor{currentstroke}{rgb}{0.000000,0.000000,0.000000}%
\pgfsetstrokecolor{currentstroke}%
\pgfsetstrokeopacity{0.000000}%
\pgfsetdash{}{0pt}%
\pgfpathmoveto{\pgfqpoint{8.764400in}{0.660000in}}%
\pgfpathlineto{\pgfqpoint{8.770600in}{0.660000in}}%
\pgfpathlineto{\pgfqpoint{8.770600in}{3.193721in}}%
\pgfpathlineto{\pgfqpoint{8.764400in}{3.193721in}}%
\pgfpathlineto{\pgfqpoint{8.764400in}{0.660000in}}%
\pgfpathclose%
\pgfusepath{fill}%
\end{pgfscope}%
\begin{pgfscope}%
\pgfpathrectangle{\pgfqpoint{1.250000in}{0.660000in}}{\pgfqpoint{7.750000in}{4.620000in}}%
\pgfusepath{clip}%
\pgfsetbuttcap%
\pgfsetmiterjoin%
\definecolor{currentfill}{rgb}{0.000000,0.000000,0.000000}%
\pgfsetfillcolor{currentfill}%
\pgfsetlinewidth{0.000000pt}%
\definecolor{currentstroke}{rgb}{0.000000,0.000000,0.000000}%
\pgfsetstrokecolor{currentstroke}%
\pgfsetstrokeopacity{0.000000}%
\pgfsetdash{}{0pt}%
\pgfpathmoveto{\pgfqpoint{8.772150in}{0.660000in}}%
\pgfpathlineto{\pgfqpoint{8.778350in}{0.660000in}}%
\pgfpathlineto{\pgfqpoint{8.778350in}{3.193738in}}%
\pgfpathlineto{\pgfqpoint{8.772150in}{3.193738in}}%
\pgfpathlineto{\pgfqpoint{8.772150in}{0.660000in}}%
\pgfpathclose%
\pgfusepath{fill}%
\end{pgfscope}%
\begin{pgfscope}%
\pgfpathrectangle{\pgfqpoint{1.250000in}{0.660000in}}{\pgfqpoint{7.750000in}{4.620000in}}%
\pgfusepath{clip}%
\pgfsetbuttcap%
\pgfsetmiterjoin%
\definecolor{currentfill}{rgb}{0.000000,0.000000,0.000000}%
\pgfsetfillcolor{currentfill}%
\pgfsetlinewidth{0.000000pt}%
\definecolor{currentstroke}{rgb}{0.000000,0.000000,0.000000}%
\pgfsetstrokecolor{currentstroke}%
\pgfsetstrokeopacity{0.000000}%
\pgfsetdash{}{0pt}%
\pgfpathmoveto{\pgfqpoint{8.779900in}{0.660000in}}%
\pgfpathlineto{\pgfqpoint{8.786100in}{0.660000in}}%
\pgfpathlineto{\pgfqpoint{8.786100in}{3.193719in}}%
\pgfpathlineto{\pgfqpoint{8.779900in}{3.193719in}}%
\pgfpathlineto{\pgfqpoint{8.779900in}{0.660000in}}%
\pgfpathclose%
\pgfusepath{fill}%
\end{pgfscope}%
\begin{pgfscope}%
\pgfpathrectangle{\pgfqpoint{1.250000in}{0.660000in}}{\pgfqpoint{7.750000in}{4.620000in}}%
\pgfusepath{clip}%
\pgfsetbuttcap%
\pgfsetmiterjoin%
\definecolor{currentfill}{rgb}{0.000000,0.000000,0.000000}%
\pgfsetfillcolor{currentfill}%
\pgfsetlinewidth{0.000000pt}%
\definecolor{currentstroke}{rgb}{0.000000,0.000000,0.000000}%
\pgfsetstrokecolor{currentstroke}%
\pgfsetstrokeopacity{0.000000}%
\pgfsetdash{}{0pt}%
\pgfpathmoveto{\pgfqpoint{8.787650in}{0.660000in}}%
\pgfpathlineto{\pgfqpoint{8.793850in}{0.660000in}}%
\pgfpathlineto{\pgfqpoint{8.793850in}{3.193664in}}%
\pgfpathlineto{\pgfqpoint{8.787650in}{3.193664in}}%
\pgfpathlineto{\pgfqpoint{8.787650in}{0.660000in}}%
\pgfpathclose%
\pgfusepath{fill}%
\end{pgfscope}%
\begin{pgfscope}%
\pgfpathrectangle{\pgfqpoint{1.250000in}{0.660000in}}{\pgfqpoint{7.750000in}{4.620000in}}%
\pgfusepath{clip}%
\pgfsetbuttcap%
\pgfsetmiterjoin%
\definecolor{currentfill}{rgb}{0.000000,0.000000,0.000000}%
\pgfsetfillcolor{currentfill}%
\pgfsetlinewidth{0.000000pt}%
\definecolor{currentstroke}{rgb}{0.000000,0.000000,0.000000}%
\pgfsetstrokecolor{currentstroke}%
\pgfsetstrokeopacity{0.000000}%
\pgfsetdash{}{0pt}%
\pgfpathmoveto{\pgfqpoint{8.795400in}{0.660000in}}%
\pgfpathlineto{\pgfqpoint{8.801600in}{0.660000in}}%
\pgfpathlineto{\pgfqpoint{8.801600in}{3.193572in}}%
\pgfpathlineto{\pgfqpoint{8.795400in}{3.193572in}}%
\pgfpathlineto{\pgfqpoint{8.795400in}{0.660000in}}%
\pgfpathclose%
\pgfusepath{fill}%
\end{pgfscope}%
\begin{pgfscope}%
\pgfpathrectangle{\pgfqpoint{1.250000in}{0.660000in}}{\pgfqpoint{7.750000in}{4.620000in}}%
\pgfusepath{clip}%
\pgfsetbuttcap%
\pgfsetmiterjoin%
\definecolor{currentfill}{rgb}{0.000000,0.000000,0.000000}%
\pgfsetfillcolor{currentfill}%
\pgfsetlinewidth{0.000000pt}%
\definecolor{currentstroke}{rgb}{0.000000,0.000000,0.000000}%
\pgfsetstrokecolor{currentstroke}%
\pgfsetstrokeopacity{0.000000}%
\pgfsetdash{}{0pt}%
\pgfpathmoveto{\pgfqpoint{8.803150in}{0.660000in}}%
\pgfpathlineto{\pgfqpoint{8.809350in}{0.660000in}}%
\pgfpathlineto{\pgfqpoint{8.809350in}{3.202493in}}%
\pgfpathlineto{\pgfqpoint{8.803150in}{3.202493in}}%
\pgfpathlineto{\pgfqpoint{8.803150in}{0.660000in}}%
\pgfpathclose%
\pgfusepath{fill}%
\end{pgfscope}%
\begin{pgfscope}%
\pgfpathrectangle{\pgfqpoint{1.250000in}{0.660000in}}{\pgfqpoint{7.750000in}{4.620000in}}%
\pgfusepath{clip}%
\pgfsetbuttcap%
\pgfsetmiterjoin%
\definecolor{currentfill}{rgb}{0.000000,0.000000,0.000000}%
\pgfsetfillcolor{currentfill}%
\pgfsetlinewidth{0.000000pt}%
\definecolor{currentstroke}{rgb}{0.000000,0.000000,0.000000}%
\pgfsetstrokecolor{currentstroke}%
\pgfsetstrokeopacity{0.000000}%
\pgfsetdash{}{0pt}%
\pgfpathmoveto{\pgfqpoint{8.810900in}{0.660000in}}%
\pgfpathlineto{\pgfqpoint{8.817100in}{0.660000in}}%
\pgfpathlineto{\pgfqpoint{8.817100in}{3.200886in}}%
\pgfpathlineto{\pgfqpoint{8.810900in}{3.200886in}}%
\pgfpathlineto{\pgfqpoint{8.810900in}{0.660000in}}%
\pgfpathclose%
\pgfusepath{fill}%
\end{pgfscope}%
\begin{pgfscope}%
\pgfpathrectangle{\pgfqpoint{1.250000in}{0.660000in}}{\pgfqpoint{7.750000in}{4.620000in}}%
\pgfusepath{clip}%
\pgfsetbuttcap%
\pgfsetmiterjoin%
\definecolor{currentfill}{rgb}{0.000000,0.000000,0.000000}%
\pgfsetfillcolor{currentfill}%
\pgfsetlinewidth{0.000000pt}%
\definecolor{currentstroke}{rgb}{0.000000,0.000000,0.000000}%
\pgfsetstrokecolor{currentstroke}%
\pgfsetstrokeopacity{0.000000}%
\pgfsetdash{}{0pt}%
\pgfpathmoveto{\pgfqpoint{8.818650in}{0.660000in}}%
\pgfpathlineto{\pgfqpoint{8.824850in}{0.660000in}}%
\pgfpathlineto{\pgfqpoint{8.824850in}{3.193067in}}%
\pgfpathlineto{\pgfqpoint{8.818650in}{3.193067in}}%
\pgfpathlineto{\pgfqpoint{8.818650in}{0.660000in}}%
\pgfpathclose%
\pgfusepath{fill}%
\end{pgfscope}%
\begin{pgfscope}%
\pgfpathrectangle{\pgfqpoint{1.250000in}{0.660000in}}{\pgfqpoint{7.750000in}{4.620000in}}%
\pgfusepath{clip}%
\pgfsetbuttcap%
\pgfsetmiterjoin%
\definecolor{currentfill}{rgb}{0.000000,0.000000,0.000000}%
\pgfsetfillcolor{currentfill}%
\pgfsetlinewidth{0.000000pt}%
\definecolor{currentstroke}{rgb}{0.000000,0.000000,0.000000}%
\pgfsetstrokecolor{currentstroke}%
\pgfsetstrokeopacity{0.000000}%
\pgfsetdash{}{0pt}%
\pgfpathmoveto{\pgfqpoint{8.826400in}{0.660000in}}%
\pgfpathlineto{\pgfqpoint{8.832600in}{0.660000in}}%
\pgfpathlineto{\pgfqpoint{8.832600in}{3.192850in}}%
\pgfpathlineto{\pgfqpoint{8.826400in}{3.192850in}}%
\pgfpathlineto{\pgfqpoint{8.826400in}{0.660000in}}%
\pgfpathclose%
\pgfusepath{fill}%
\end{pgfscope}%
\begin{pgfscope}%
\pgfpathrectangle{\pgfqpoint{1.250000in}{0.660000in}}{\pgfqpoint{7.750000in}{4.620000in}}%
\pgfusepath{clip}%
\pgfsetbuttcap%
\pgfsetmiterjoin%
\definecolor{currentfill}{rgb}{0.000000,0.000000,0.000000}%
\pgfsetfillcolor{currentfill}%
\pgfsetlinewidth{0.000000pt}%
\definecolor{currentstroke}{rgb}{0.000000,0.000000,0.000000}%
\pgfsetstrokecolor{currentstroke}%
\pgfsetstrokeopacity{0.000000}%
\pgfsetdash{}{0pt}%
\pgfpathmoveto{\pgfqpoint{8.834150in}{0.660000in}}%
\pgfpathlineto{\pgfqpoint{8.840350in}{0.660000in}}%
\pgfpathlineto{\pgfqpoint{8.840350in}{3.192925in}}%
\pgfpathlineto{\pgfqpoint{8.834150in}{3.192925in}}%
\pgfpathlineto{\pgfqpoint{8.834150in}{0.660000in}}%
\pgfpathclose%
\pgfusepath{fill}%
\end{pgfscope}%
\begin{pgfscope}%
\pgfpathrectangle{\pgfqpoint{1.250000in}{0.660000in}}{\pgfqpoint{7.750000in}{4.620000in}}%
\pgfusepath{clip}%
\pgfsetbuttcap%
\pgfsetmiterjoin%
\definecolor{currentfill}{rgb}{0.000000,0.000000,0.000000}%
\pgfsetfillcolor{currentfill}%
\pgfsetlinewidth{0.000000pt}%
\definecolor{currentstroke}{rgb}{0.000000,0.000000,0.000000}%
\pgfsetstrokecolor{currentstroke}%
\pgfsetstrokeopacity{0.000000}%
\pgfsetdash{}{0pt}%
\pgfpathmoveto{\pgfqpoint{8.841900in}{0.660000in}}%
\pgfpathlineto{\pgfqpoint{8.848100in}{0.660000in}}%
\pgfpathlineto{\pgfqpoint{8.848100in}{3.192964in}}%
\pgfpathlineto{\pgfqpoint{8.841900in}{3.192964in}}%
\pgfpathlineto{\pgfqpoint{8.841900in}{0.660000in}}%
\pgfpathclose%
\pgfusepath{fill}%
\end{pgfscope}%
\begin{pgfscope}%
\pgfpathrectangle{\pgfqpoint{1.250000in}{0.660000in}}{\pgfqpoint{7.750000in}{4.620000in}}%
\pgfusepath{clip}%
\pgfsetbuttcap%
\pgfsetmiterjoin%
\definecolor{currentfill}{rgb}{0.000000,0.000000,0.000000}%
\pgfsetfillcolor{currentfill}%
\pgfsetlinewidth{0.000000pt}%
\definecolor{currentstroke}{rgb}{0.000000,0.000000,0.000000}%
\pgfsetstrokecolor{currentstroke}%
\pgfsetstrokeopacity{0.000000}%
\pgfsetdash{}{0pt}%
\pgfpathmoveto{\pgfqpoint{8.849650in}{0.660000in}}%
\pgfpathlineto{\pgfqpoint{8.855850in}{0.660000in}}%
\pgfpathlineto{\pgfqpoint{8.855850in}{3.192966in}}%
\pgfpathlineto{\pgfqpoint{8.849650in}{3.192966in}}%
\pgfpathlineto{\pgfqpoint{8.849650in}{0.660000in}}%
\pgfpathclose%
\pgfusepath{fill}%
\end{pgfscope}%
\begin{pgfscope}%
\pgfpathrectangle{\pgfqpoint{1.250000in}{0.660000in}}{\pgfqpoint{7.750000in}{4.620000in}}%
\pgfusepath{clip}%
\pgfsetbuttcap%
\pgfsetmiterjoin%
\definecolor{currentfill}{rgb}{0.000000,0.000000,0.000000}%
\pgfsetfillcolor{currentfill}%
\pgfsetlinewidth{0.000000pt}%
\definecolor{currentstroke}{rgb}{0.000000,0.000000,0.000000}%
\pgfsetstrokecolor{currentstroke}%
\pgfsetstrokeopacity{0.000000}%
\pgfsetdash{}{0pt}%
\pgfpathmoveto{\pgfqpoint{8.857400in}{0.660000in}}%
\pgfpathlineto{\pgfqpoint{8.863600in}{0.660000in}}%
\pgfpathlineto{\pgfqpoint{8.863600in}{3.192930in}}%
\pgfpathlineto{\pgfqpoint{8.857400in}{3.192930in}}%
\pgfpathlineto{\pgfqpoint{8.857400in}{0.660000in}}%
\pgfpathclose%
\pgfusepath{fill}%
\end{pgfscope}%
\begin{pgfscope}%
\pgfpathrectangle{\pgfqpoint{1.250000in}{0.660000in}}{\pgfqpoint{7.750000in}{4.620000in}}%
\pgfusepath{clip}%
\pgfsetbuttcap%
\pgfsetmiterjoin%
\definecolor{currentfill}{rgb}{0.000000,0.000000,0.000000}%
\pgfsetfillcolor{currentfill}%
\pgfsetlinewidth{0.000000pt}%
\definecolor{currentstroke}{rgb}{0.000000,0.000000,0.000000}%
\pgfsetstrokecolor{currentstroke}%
\pgfsetstrokeopacity{0.000000}%
\pgfsetdash{}{0pt}%
\pgfpathmoveto{\pgfqpoint{8.865150in}{0.660000in}}%
\pgfpathlineto{\pgfqpoint{8.871350in}{0.660000in}}%
\pgfpathlineto{\pgfqpoint{8.871350in}{3.192856in}}%
\pgfpathlineto{\pgfqpoint{8.865150in}{3.192856in}}%
\pgfpathlineto{\pgfqpoint{8.865150in}{0.660000in}}%
\pgfpathclose%
\pgfusepath{fill}%
\end{pgfscope}%
\begin{pgfscope}%
\pgfpathrectangle{\pgfqpoint{1.250000in}{0.660000in}}{\pgfqpoint{7.750000in}{4.620000in}}%
\pgfusepath{clip}%
\pgfsetbuttcap%
\pgfsetmiterjoin%
\definecolor{currentfill}{rgb}{0.000000,0.000000,0.000000}%
\pgfsetfillcolor{currentfill}%
\pgfsetlinewidth{0.000000pt}%
\definecolor{currentstroke}{rgb}{0.000000,0.000000,0.000000}%
\pgfsetstrokecolor{currentstroke}%
\pgfsetstrokeopacity{0.000000}%
\pgfsetdash{}{0pt}%
\pgfpathmoveto{\pgfqpoint{8.872900in}{0.660000in}}%
\pgfpathlineto{\pgfqpoint{8.879100in}{0.660000in}}%
\pgfpathlineto{\pgfqpoint{8.879100in}{3.192743in}}%
\pgfpathlineto{\pgfqpoint{8.872900in}{3.192743in}}%
\pgfpathlineto{\pgfqpoint{8.872900in}{0.660000in}}%
\pgfpathclose%
\pgfusepath{fill}%
\end{pgfscope}%
\begin{pgfscope}%
\pgfpathrectangle{\pgfqpoint{1.250000in}{0.660000in}}{\pgfqpoint{7.750000in}{4.620000in}}%
\pgfusepath{clip}%
\pgfsetbuttcap%
\pgfsetmiterjoin%
\definecolor{currentfill}{rgb}{0.000000,0.000000,0.000000}%
\pgfsetfillcolor{currentfill}%
\pgfsetlinewidth{0.000000pt}%
\definecolor{currentstroke}{rgb}{0.000000,0.000000,0.000000}%
\pgfsetstrokecolor{currentstroke}%
\pgfsetstrokeopacity{0.000000}%
\pgfsetdash{}{0pt}%
\pgfpathmoveto{\pgfqpoint{8.880650in}{0.660000in}}%
\pgfpathlineto{\pgfqpoint{8.886850in}{0.660000in}}%
\pgfpathlineto{\pgfqpoint{8.886850in}{3.192590in}}%
\pgfpathlineto{\pgfqpoint{8.880650in}{3.192590in}}%
\pgfpathlineto{\pgfqpoint{8.880650in}{0.660000in}}%
\pgfpathclose%
\pgfusepath{fill}%
\end{pgfscope}%
\begin{pgfscope}%
\pgfpathrectangle{\pgfqpoint{1.250000in}{0.660000in}}{\pgfqpoint{7.750000in}{4.620000in}}%
\pgfusepath{clip}%
\pgfsetbuttcap%
\pgfsetmiterjoin%
\definecolor{currentfill}{rgb}{0.000000,0.000000,0.000000}%
\pgfsetfillcolor{currentfill}%
\pgfsetlinewidth{0.000000pt}%
\definecolor{currentstroke}{rgb}{0.000000,0.000000,0.000000}%
\pgfsetstrokecolor{currentstroke}%
\pgfsetstrokeopacity{0.000000}%
\pgfsetdash{}{0pt}%
\pgfpathmoveto{\pgfqpoint{8.888400in}{0.660000in}}%
\pgfpathlineto{\pgfqpoint{8.894600in}{0.660000in}}%
\pgfpathlineto{\pgfqpoint{8.894600in}{3.192396in}}%
\pgfpathlineto{\pgfqpoint{8.888400in}{3.192396in}}%
\pgfpathlineto{\pgfqpoint{8.888400in}{0.660000in}}%
\pgfpathclose%
\pgfusepath{fill}%
\end{pgfscope}%
\begin{pgfscope}%
\pgfpathrectangle{\pgfqpoint{1.250000in}{0.660000in}}{\pgfqpoint{7.750000in}{4.620000in}}%
\pgfusepath{clip}%
\pgfsetbuttcap%
\pgfsetmiterjoin%
\definecolor{currentfill}{rgb}{0.000000,0.000000,0.000000}%
\pgfsetfillcolor{currentfill}%
\pgfsetlinewidth{0.000000pt}%
\definecolor{currentstroke}{rgb}{0.000000,0.000000,0.000000}%
\pgfsetstrokecolor{currentstroke}%
\pgfsetstrokeopacity{0.000000}%
\pgfsetdash{}{0pt}%
\pgfpathmoveto{\pgfqpoint{8.896150in}{0.660000in}}%
\pgfpathlineto{\pgfqpoint{8.902350in}{0.660000in}}%
\pgfpathlineto{\pgfqpoint{8.902350in}{3.192162in}}%
\pgfpathlineto{\pgfqpoint{8.896150in}{3.192162in}}%
\pgfpathlineto{\pgfqpoint{8.896150in}{0.660000in}}%
\pgfpathclose%
\pgfusepath{fill}%
\end{pgfscope}%
\begin{pgfscope}%
\pgfpathrectangle{\pgfqpoint{1.250000in}{0.660000in}}{\pgfqpoint{7.750000in}{4.620000in}}%
\pgfusepath{clip}%
\pgfsetbuttcap%
\pgfsetmiterjoin%
\definecolor{currentfill}{rgb}{0.000000,0.000000,0.000000}%
\pgfsetfillcolor{currentfill}%
\pgfsetlinewidth{0.000000pt}%
\definecolor{currentstroke}{rgb}{0.000000,0.000000,0.000000}%
\pgfsetstrokecolor{currentstroke}%
\pgfsetstrokeopacity{0.000000}%
\pgfsetdash{}{0pt}%
\pgfpathmoveto{\pgfqpoint{8.903900in}{0.660000in}}%
\pgfpathlineto{\pgfqpoint{8.910100in}{0.660000in}}%
\pgfpathlineto{\pgfqpoint{8.910100in}{3.192162in}}%
\pgfpathlineto{\pgfqpoint{8.903900in}{3.192162in}}%
\pgfpathlineto{\pgfqpoint{8.903900in}{0.660000in}}%
\pgfpathclose%
\pgfusepath{fill}%
\end{pgfscope}%
\begin{pgfscope}%
\pgfpathrectangle{\pgfqpoint{1.250000in}{0.660000in}}{\pgfqpoint{7.750000in}{4.620000in}}%
\pgfusepath{clip}%
\pgfsetbuttcap%
\pgfsetmiterjoin%
\definecolor{currentfill}{rgb}{0.000000,0.000000,0.000000}%
\pgfsetfillcolor{currentfill}%
\pgfsetlinewidth{0.000000pt}%
\definecolor{currentstroke}{rgb}{0.000000,0.000000,0.000000}%
\pgfsetstrokecolor{currentstroke}%
\pgfsetstrokeopacity{0.000000}%
\pgfsetdash{}{0pt}%
\pgfpathmoveto{\pgfqpoint{8.911650in}{0.660000in}}%
\pgfpathlineto{\pgfqpoint{8.917850in}{0.660000in}}%
\pgfpathlineto{\pgfqpoint{8.917850in}{3.192215in}}%
\pgfpathlineto{\pgfqpoint{8.911650in}{3.192215in}}%
\pgfpathlineto{\pgfqpoint{8.911650in}{0.660000in}}%
\pgfpathclose%
\pgfusepath{fill}%
\end{pgfscope}%
\begin{pgfscope}%
\pgfpathrectangle{\pgfqpoint{1.250000in}{0.660000in}}{\pgfqpoint{7.750000in}{4.620000in}}%
\pgfusepath{clip}%
\pgfsetbuttcap%
\pgfsetmiterjoin%
\definecolor{currentfill}{rgb}{0.000000,0.000000,0.000000}%
\pgfsetfillcolor{currentfill}%
\pgfsetlinewidth{0.000000pt}%
\definecolor{currentstroke}{rgb}{0.000000,0.000000,0.000000}%
\pgfsetstrokecolor{currentstroke}%
\pgfsetstrokeopacity{0.000000}%
\pgfsetdash{}{0pt}%
\pgfpathmoveto{\pgfqpoint{8.919400in}{0.660000in}}%
\pgfpathlineto{\pgfqpoint{8.925600in}{0.660000in}}%
\pgfpathlineto{\pgfqpoint{8.925600in}{3.192230in}}%
\pgfpathlineto{\pgfqpoint{8.919400in}{3.192230in}}%
\pgfpathlineto{\pgfqpoint{8.919400in}{0.660000in}}%
\pgfpathclose%
\pgfusepath{fill}%
\end{pgfscope}%
\begin{pgfscope}%
\pgfpathrectangle{\pgfqpoint{1.250000in}{0.660000in}}{\pgfqpoint{7.750000in}{4.620000in}}%
\pgfusepath{clip}%
\pgfsetbuttcap%
\pgfsetmiterjoin%
\definecolor{currentfill}{rgb}{0.000000,0.000000,0.000000}%
\pgfsetfillcolor{currentfill}%
\pgfsetlinewidth{0.000000pt}%
\definecolor{currentstroke}{rgb}{0.000000,0.000000,0.000000}%
\pgfsetstrokecolor{currentstroke}%
\pgfsetstrokeopacity{0.000000}%
\pgfsetdash{}{0pt}%
\pgfpathmoveto{\pgfqpoint{8.927150in}{0.660000in}}%
\pgfpathlineto{\pgfqpoint{8.933350in}{0.660000in}}%
\pgfpathlineto{\pgfqpoint{8.933350in}{3.192206in}}%
\pgfpathlineto{\pgfqpoint{8.927150in}{3.192206in}}%
\pgfpathlineto{\pgfqpoint{8.927150in}{0.660000in}}%
\pgfpathclose%
\pgfusepath{fill}%
\end{pgfscope}%
\begin{pgfscope}%
\pgfpathrectangle{\pgfqpoint{1.250000in}{0.660000in}}{\pgfqpoint{7.750000in}{4.620000in}}%
\pgfusepath{clip}%
\pgfsetbuttcap%
\pgfsetmiterjoin%
\definecolor{currentfill}{rgb}{0.000000,0.000000,0.000000}%
\pgfsetfillcolor{currentfill}%
\pgfsetlinewidth{0.000000pt}%
\definecolor{currentstroke}{rgb}{0.000000,0.000000,0.000000}%
\pgfsetstrokecolor{currentstroke}%
\pgfsetstrokeopacity{0.000000}%
\pgfsetdash{}{0pt}%
\pgfpathmoveto{\pgfqpoint{8.934900in}{0.660000in}}%
\pgfpathlineto{\pgfqpoint{8.941100in}{0.660000in}}%
\pgfpathlineto{\pgfqpoint{8.941100in}{3.192141in}}%
\pgfpathlineto{\pgfqpoint{8.934900in}{3.192141in}}%
\pgfpathlineto{\pgfqpoint{8.934900in}{0.660000in}}%
\pgfpathclose%
\pgfusepath{fill}%
\end{pgfscope}%
\begin{pgfscope}%
\pgfpathrectangle{\pgfqpoint{1.250000in}{0.660000in}}{\pgfqpoint{7.750000in}{4.620000in}}%
\pgfusepath{clip}%
\pgfsetbuttcap%
\pgfsetmiterjoin%
\definecolor{currentfill}{rgb}{0.000000,0.000000,0.000000}%
\pgfsetfillcolor{currentfill}%
\pgfsetlinewidth{0.000000pt}%
\definecolor{currentstroke}{rgb}{0.000000,0.000000,0.000000}%
\pgfsetstrokecolor{currentstroke}%
\pgfsetstrokeopacity{0.000000}%
\pgfsetdash{}{0pt}%
\pgfpathmoveto{\pgfqpoint{8.942650in}{0.660000in}}%
\pgfpathlineto{\pgfqpoint{8.948850in}{0.660000in}}%
\pgfpathlineto{\pgfqpoint{8.948850in}{3.192037in}}%
\pgfpathlineto{\pgfqpoint{8.942650in}{3.192037in}}%
\pgfpathlineto{\pgfqpoint{8.942650in}{0.660000in}}%
\pgfpathclose%
\pgfusepath{fill}%
\end{pgfscope}%
\begin{pgfscope}%
\pgfpathrectangle{\pgfqpoint{1.250000in}{0.660000in}}{\pgfqpoint{7.750000in}{4.620000in}}%
\pgfusepath{clip}%
\pgfsetbuttcap%
\pgfsetmiterjoin%
\definecolor{currentfill}{rgb}{0.000000,0.000000,0.000000}%
\pgfsetfillcolor{currentfill}%
\pgfsetlinewidth{0.000000pt}%
\definecolor{currentstroke}{rgb}{0.000000,0.000000,0.000000}%
\pgfsetstrokecolor{currentstroke}%
\pgfsetstrokeopacity{0.000000}%
\pgfsetdash{}{0pt}%
\pgfpathmoveto{\pgfqpoint{8.950400in}{0.660000in}}%
\pgfpathlineto{\pgfqpoint{8.956600in}{0.660000in}}%
\pgfpathlineto{\pgfqpoint{8.956600in}{3.191891in}}%
\pgfpathlineto{\pgfqpoint{8.950400in}{3.191891in}}%
\pgfpathlineto{\pgfqpoint{8.950400in}{0.660000in}}%
\pgfpathclose%
\pgfusepath{fill}%
\end{pgfscope}%
\begin{pgfscope}%
\pgfpathrectangle{\pgfqpoint{1.250000in}{0.660000in}}{\pgfqpoint{7.750000in}{4.620000in}}%
\pgfusepath{clip}%
\pgfsetbuttcap%
\pgfsetmiterjoin%
\definecolor{currentfill}{rgb}{0.000000,0.000000,0.000000}%
\pgfsetfillcolor{currentfill}%
\pgfsetlinewidth{0.000000pt}%
\definecolor{currentstroke}{rgb}{0.000000,0.000000,0.000000}%
\pgfsetstrokecolor{currentstroke}%
\pgfsetstrokeopacity{0.000000}%
\pgfsetdash{}{0pt}%
\pgfpathmoveto{\pgfqpoint{8.958150in}{0.660000in}}%
\pgfpathlineto{\pgfqpoint{8.964350in}{0.660000in}}%
\pgfpathlineto{\pgfqpoint{8.964350in}{3.191703in}}%
\pgfpathlineto{\pgfqpoint{8.958150in}{3.191703in}}%
\pgfpathlineto{\pgfqpoint{8.958150in}{0.660000in}}%
\pgfpathclose%
\pgfusepath{fill}%
\end{pgfscope}%
\begin{pgfscope}%
\pgfpathrectangle{\pgfqpoint{1.250000in}{0.660000in}}{\pgfqpoint{7.750000in}{4.620000in}}%
\pgfusepath{clip}%
\pgfsetbuttcap%
\pgfsetmiterjoin%
\definecolor{currentfill}{rgb}{0.000000,0.000000,0.000000}%
\pgfsetfillcolor{currentfill}%
\pgfsetlinewidth{0.000000pt}%
\definecolor{currentstroke}{rgb}{0.000000,0.000000,0.000000}%
\pgfsetstrokecolor{currentstroke}%
\pgfsetstrokeopacity{0.000000}%
\pgfsetdash{}{0pt}%
\pgfpathmoveto{\pgfqpoint{8.965900in}{0.660000in}}%
\pgfpathlineto{\pgfqpoint{8.972100in}{0.660000in}}%
\pgfpathlineto{\pgfqpoint{8.972100in}{3.191472in}}%
\pgfpathlineto{\pgfqpoint{8.965900in}{3.191472in}}%
\pgfpathlineto{\pgfqpoint{8.965900in}{0.660000in}}%
\pgfpathclose%
\pgfusepath{fill}%
\end{pgfscope}%
\begin{pgfscope}%
\pgfpathrectangle{\pgfqpoint{1.250000in}{0.660000in}}{\pgfqpoint{7.750000in}{4.620000in}}%
\pgfusepath{clip}%
\pgfsetbuttcap%
\pgfsetmiterjoin%
\definecolor{currentfill}{rgb}{0.000000,0.000000,0.000000}%
\pgfsetfillcolor{currentfill}%
\pgfsetlinewidth{0.000000pt}%
\definecolor{currentstroke}{rgb}{0.000000,0.000000,0.000000}%
\pgfsetstrokecolor{currentstroke}%
\pgfsetstrokeopacity{0.000000}%
\pgfsetdash{}{0pt}%
\pgfpathmoveto{\pgfqpoint{8.973650in}{0.660000in}}%
\pgfpathlineto{\pgfqpoint{8.979850in}{0.660000in}}%
\pgfpathlineto{\pgfqpoint{8.979850in}{3.192002in}}%
\pgfpathlineto{\pgfqpoint{8.973650in}{3.192002in}}%
\pgfpathlineto{\pgfqpoint{8.973650in}{0.660000in}}%
\pgfpathclose%
\pgfusepath{fill}%
\end{pgfscope}%
\begin{pgfscope}%
\pgfpathrectangle{\pgfqpoint{1.250000in}{0.660000in}}{\pgfqpoint{7.750000in}{4.620000in}}%
\pgfusepath{clip}%
\pgfsetbuttcap%
\pgfsetmiterjoin%
\definecolor{currentfill}{rgb}{0.000000,0.000000,0.000000}%
\pgfsetfillcolor{currentfill}%
\pgfsetlinewidth{0.000000pt}%
\definecolor{currentstroke}{rgb}{0.000000,0.000000,0.000000}%
\pgfsetstrokecolor{currentstroke}%
\pgfsetstrokeopacity{0.000000}%
\pgfsetdash{}{0pt}%
\pgfpathmoveto{\pgfqpoint{8.981400in}{0.660000in}}%
\pgfpathlineto{\pgfqpoint{8.987600in}{0.660000in}}%
\pgfpathlineto{\pgfqpoint{8.987600in}{3.191498in}}%
\pgfpathlineto{\pgfqpoint{8.981400in}{3.191498in}}%
\pgfpathlineto{\pgfqpoint{8.981400in}{0.660000in}}%
\pgfpathclose%
\pgfusepath{fill}%
\end{pgfscope}%
\begin{pgfscope}%
\pgfpathrectangle{\pgfqpoint{1.250000in}{0.660000in}}{\pgfqpoint{7.750000in}{4.620000in}}%
\pgfusepath{clip}%
\pgfsetbuttcap%
\pgfsetmiterjoin%
\definecolor{currentfill}{rgb}{0.000000,0.000000,0.000000}%
\pgfsetfillcolor{currentfill}%
\pgfsetlinewidth{0.000000pt}%
\definecolor{currentstroke}{rgb}{0.000000,0.000000,0.000000}%
\pgfsetstrokecolor{currentstroke}%
\pgfsetstrokeopacity{0.000000}%
\pgfsetdash{}{0pt}%
\pgfpathmoveto{\pgfqpoint{8.989150in}{0.660000in}}%
\pgfpathlineto{\pgfqpoint{8.995350in}{0.660000in}}%
\pgfpathlineto{\pgfqpoint{8.995350in}{3.191517in}}%
\pgfpathlineto{\pgfqpoint{8.989150in}{3.191517in}}%
\pgfpathlineto{\pgfqpoint{8.989150in}{0.660000in}}%
\pgfpathclose%
\pgfusepath{fill}%
\end{pgfscope}%
\begin{pgfscope}%
\pgfpathrectangle{\pgfqpoint{1.250000in}{0.660000in}}{\pgfqpoint{7.750000in}{4.620000in}}%
\pgfusepath{clip}%
\pgfsetbuttcap%
\pgfsetmiterjoin%
\definecolor{currentfill}{rgb}{0.000000,0.000000,0.000000}%
\pgfsetfillcolor{currentfill}%
\pgfsetlinewidth{0.000000pt}%
\definecolor{currentstroke}{rgb}{0.000000,0.000000,0.000000}%
\pgfsetstrokecolor{currentstroke}%
\pgfsetstrokeopacity{0.000000}%
\pgfsetdash{}{0pt}%
\pgfpathmoveto{\pgfqpoint{8.996900in}{0.660000in}}%
\pgfpathlineto{\pgfqpoint{9.003100in}{0.660000in}}%
\pgfpathlineto{\pgfqpoint{9.003100in}{3.191496in}}%
\pgfpathlineto{\pgfqpoint{8.996900in}{3.191496in}}%
\pgfpathlineto{\pgfqpoint{8.996900in}{0.660000in}}%
\pgfpathclose%
\pgfusepath{fill}%
\end{pgfscope}%
\begin{pgfscope}%
\pgfsetbuttcap%
\pgfsetroundjoin%
\definecolor{currentfill}{rgb}{0.000000,0.000000,0.000000}%
\pgfsetfillcolor{currentfill}%
\pgfsetlinewidth{0.803000pt}%
\definecolor{currentstroke}{rgb}{0.000000,0.000000,0.000000}%
\pgfsetstrokecolor{currentstroke}%
\pgfsetdash{}{0pt}%
\pgfsys@defobject{currentmarker}{\pgfqpoint{0.000000in}{-0.048611in}}{\pgfqpoint{0.000000in}{0.000000in}}{%
\pgfpathmoveto{\pgfqpoint{0.000000in}{0.000000in}}%
\pgfpathlineto{\pgfqpoint{0.000000in}{-0.048611in}}%
\pgfusepath{stroke,fill}%
}%
\begin{pgfscope}%
\pgfsys@transformshift{1.250000in}{0.660000in}%
\pgfsys@useobject{currentmarker}{}%
\end{pgfscope}%
\end{pgfscope}%
\begin{pgfscope}%
\definecolor{textcolor}{rgb}{0.000000,0.000000,0.000000}%
\pgfsetstrokecolor{textcolor}%
\pgfsetfillcolor{textcolor}%
\pgftext[x=1.250000in,y=0.562778in,,top]{\color{textcolor}\rmfamily\fontsize{16.000000}{19.200000}\selectfont \(\displaystyle {0}\)}%
\end{pgfscope}%
\begin{pgfscope}%
\pgfsetbuttcap%
\pgfsetroundjoin%
\definecolor{currentfill}{rgb}{0.000000,0.000000,0.000000}%
\pgfsetfillcolor{currentfill}%
\pgfsetlinewidth{0.803000pt}%
\definecolor{currentstroke}{rgb}{0.000000,0.000000,0.000000}%
\pgfsetstrokecolor{currentstroke}%
\pgfsetdash{}{0pt}%
\pgfsys@defobject{currentmarker}{\pgfqpoint{0.000000in}{-0.048611in}}{\pgfqpoint{0.000000in}{0.000000in}}{%
\pgfpathmoveto{\pgfqpoint{0.000000in}{0.000000in}}%
\pgfpathlineto{\pgfqpoint{0.000000in}{-0.048611in}}%
\pgfusepath{stroke,fill}%
}%
\begin{pgfscope}%
\pgfsys@transformshift{2.800000in}{0.660000in}%
\pgfsys@useobject{currentmarker}{}%
\end{pgfscope}%
\end{pgfscope}%
\begin{pgfscope}%
\definecolor{textcolor}{rgb}{0.000000,0.000000,0.000000}%
\pgfsetstrokecolor{textcolor}%
\pgfsetfillcolor{textcolor}%
\pgftext[x=2.800000in,y=0.562778in,,top]{\color{textcolor}\rmfamily\fontsize{16.000000}{19.200000}\selectfont \(\displaystyle {200}\)}%
\end{pgfscope}%
\begin{pgfscope}%
\pgfsetbuttcap%
\pgfsetroundjoin%
\definecolor{currentfill}{rgb}{0.000000,0.000000,0.000000}%
\pgfsetfillcolor{currentfill}%
\pgfsetlinewidth{0.803000pt}%
\definecolor{currentstroke}{rgb}{0.000000,0.000000,0.000000}%
\pgfsetstrokecolor{currentstroke}%
\pgfsetdash{}{0pt}%
\pgfsys@defobject{currentmarker}{\pgfqpoint{0.000000in}{-0.048611in}}{\pgfqpoint{0.000000in}{0.000000in}}{%
\pgfpathmoveto{\pgfqpoint{0.000000in}{0.000000in}}%
\pgfpathlineto{\pgfqpoint{0.000000in}{-0.048611in}}%
\pgfusepath{stroke,fill}%
}%
\begin{pgfscope}%
\pgfsys@transformshift{4.350000in}{0.660000in}%
\pgfsys@useobject{currentmarker}{}%
\end{pgfscope}%
\end{pgfscope}%
\begin{pgfscope}%
\definecolor{textcolor}{rgb}{0.000000,0.000000,0.000000}%
\pgfsetstrokecolor{textcolor}%
\pgfsetfillcolor{textcolor}%
\pgftext[x=4.350000in,y=0.562778in,,top]{\color{textcolor}\rmfamily\fontsize{16.000000}{19.200000}\selectfont \(\displaystyle {400}\)}%
\end{pgfscope}%
\begin{pgfscope}%
\pgfsetbuttcap%
\pgfsetroundjoin%
\definecolor{currentfill}{rgb}{0.000000,0.000000,0.000000}%
\pgfsetfillcolor{currentfill}%
\pgfsetlinewidth{0.803000pt}%
\definecolor{currentstroke}{rgb}{0.000000,0.000000,0.000000}%
\pgfsetstrokecolor{currentstroke}%
\pgfsetdash{}{0pt}%
\pgfsys@defobject{currentmarker}{\pgfqpoint{0.000000in}{-0.048611in}}{\pgfqpoint{0.000000in}{0.000000in}}{%
\pgfpathmoveto{\pgfqpoint{0.000000in}{0.000000in}}%
\pgfpathlineto{\pgfqpoint{0.000000in}{-0.048611in}}%
\pgfusepath{stroke,fill}%
}%
\begin{pgfscope}%
\pgfsys@transformshift{5.900000in}{0.660000in}%
\pgfsys@useobject{currentmarker}{}%
\end{pgfscope}%
\end{pgfscope}%
\begin{pgfscope}%
\definecolor{textcolor}{rgb}{0.000000,0.000000,0.000000}%
\pgfsetstrokecolor{textcolor}%
\pgfsetfillcolor{textcolor}%
\pgftext[x=5.900000in,y=0.562778in,,top]{\color{textcolor}\rmfamily\fontsize{16.000000}{19.200000}\selectfont \(\displaystyle {600}\)}%
\end{pgfscope}%
\begin{pgfscope}%
\pgfsetbuttcap%
\pgfsetroundjoin%
\definecolor{currentfill}{rgb}{0.000000,0.000000,0.000000}%
\pgfsetfillcolor{currentfill}%
\pgfsetlinewidth{0.803000pt}%
\definecolor{currentstroke}{rgb}{0.000000,0.000000,0.000000}%
\pgfsetstrokecolor{currentstroke}%
\pgfsetdash{}{0pt}%
\pgfsys@defobject{currentmarker}{\pgfqpoint{0.000000in}{-0.048611in}}{\pgfqpoint{0.000000in}{0.000000in}}{%
\pgfpathmoveto{\pgfqpoint{0.000000in}{0.000000in}}%
\pgfpathlineto{\pgfqpoint{0.000000in}{-0.048611in}}%
\pgfusepath{stroke,fill}%
}%
\begin{pgfscope}%
\pgfsys@transformshift{7.450000in}{0.660000in}%
\pgfsys@useobject{currentmarker}{}%
\end{pgfscope}%
\end{pgfscope}%
\begin{pgfscope}%
\definecolor{textcolor}{rgb}{0.000000,0.000000,0.000000}%
\pgfsetstrokecolor{textcolor}%
\pgfsetfillcolor{textcolor}%
\pgftext[x=7.450000in,y=0.562778in,,top]{\color{textcolor}\rmfamily\fontsize{16.000000}{19.200000}\selectfont \(\displaystyle {800}\)}%
\end{pgfscope}%
\begin{pgfscope}%
\pgfsetbuttcap%
\pgfsetroundjoin%
\definecolor{currentfill}{rgb}{0.000000,0.000000,0.000000}%
\pgfsetfillcolor{currentfill}%
\pgfsetlinewidth{0.803000pt}%
\definecolor{currentstroke}{rgb}{0.000000,0.000000,0.000000}%
\pgfsetstrokecolor{currentstroke}%
\pgfsetdash{}{0pt}%
\pgfsys@defobject{currentmarker}{\pgfqpoint{0.000000in}{-0.048611in}}{\pgfqpoint{0.000000in}{0.000000in}}{%
\pgfpathmoveto{\pgfqpoint{0.000000in}{0.000000in}}%
\pgfpathlineto{\pgfqpoint{0.000000in}{-0.048611in}}%
\pgfusepath{stroke,fill}%
}%
\begin{pgfscope}%
\pgfsys@transformshift{9.000000in}{0.660000in}%
\pgfsys@useobject{currentmarker}{}%
\end{pgfscope}%
\end{pgfscope}%
\begin{pgfscope}%
\definecolor{textcolor}{rgb}{0.000000,0.000000,0.000000}%
\pgfsetstrokecolor{textcolor}%
\pgfsetfillcolor{textcolor}%
\pgftext[x=9.000000in,y=0.562778in,,top]{\color{textcolor}\rmfamily\fontsize{16.000000}{19.200000}\selectfont \(\displaystyle {1000}\)}%
\end{pgfscope}%
\begin{pgfscope}%
\pgfsetbuttcap%
\pgfsetroundjoin%
\definecolor{currentfill}{rgb}{0.000000,0.000000,0.000000}%
\pgfsetfillcolor{currentfill}%
\pgfsetlinewidth{0.803000pt}%
\definecolor{currentstroke}{rgb}{0.000000,0.000000,0.000000}%
\pgfsetstrokecolor{currentstroke}%
\pgfsetdash{}{0pt}%
\pgfsys@defobject{currentmarker}{\pgfqpoint{-0.048611in}{0.000000in}}{\pgfqpoint{-0.000000in}{0.000000in}}{%
\pgfpathmoveto{\pgfqpoint{-0.000000in}{0.000000in}}%
\pgfpathlineto{\pgfqpoint{-0.048611in}{0.000000in}}%
\pgfusepath{stroke,fill}%
}%
\begin{pgfscope}%
\pgfsys@transformshift{1.250000in}{0.660000in}%
\pgfsys@useobject{currentmarker}{}%
\end{pgfscope}%
\end{pgfscope}%
\begin{pgfscope}%
\definecolor{textcolor}{rgb}{0.000000,0.000000,0.000000}%
\pgfsetstrokecolor{textcolor}%
\pgfsetfillcolor{textcolor}%
\pgftext[x=0.867364in, y=0.576667in, left, base]{\color{textcolor}\rmfamily\fontsize{16.000000}{19.200000}\selectfont \(\displaystyle {0.0}\)}%
\end{pgfscope}%
\begin{pgfscope}%
\pgfsetbuttcap%
\pgfsetroundjoin%
\definecolor{currentfill}{rgb}{0.000000,0.000000,0.000000}%
\pgfsetfillcolor{currentfill}%
\pgfsetlinewidth{0.803000pt}%
\definecolor{currentstroke}{rgb}{0.000000,0.000000,0.000000}%
\pgfsetstrokecolor{currentstroke}%
\pgfsetdash{}{0pt}%
\pgfsys@defobject{currentmarker}{\pgfqpoint{-0.048611in}{0.000000in}}{\pgfqpoint{-0.000000in}{0.000000in}}{%
\pgfpathmoveto{\pgfqpoint{-0.000000in}{0.000000in}}%
\pgfpathlineto{\pgfqpoint{-0.048611in}{0.000000in}}%
\pgfusepath{stroke,fill}%
}%
\begin{pgfscope}%
\pgfsys@transformshift{1.250000in}{1.584000in}%
\pgfsys@useobject{currentmarker}{}%
\end{pgfscope}%
\end{pgfscope}%
\begin{pgfscope}%
\definecolor{textcolor}{rgb}{0.000000,0.000000,0.000000}%
\pgfsetstrokecolor{textcolor}%
\pgfsetfillcolor{textcolor}%
\pgftext[x=0.867364in, y=1.500667in, left, base]{\color{textcolor}\rmfamily\fontsize{16.000000}{19.200000}\selectfont \(\displaystyle {0.2}\)}%
\end{pgfscope}%
\begin{pgfscope}%
\pgfsetbuttcap%
\pgfsetroundjoin%
\definecolor{currentfill}{rgb}{0.000000,0.000000,0.000000}%
\pgfsetfillcolor{currentfill}%
\pgfsetlinewidth{0.803000pt}%
\definecolor{currentstroke}{rgb}{0.000000,0.000000,0.000000}%
\pgfsetstrokecolor{currentstroke}%
\pgfsetdash{}{0pt}%
\pgfsys@defobject{currentmarker}{\pgfqpoint{-0.048611in}{0.000000in}}{\pgfqpoint{-0.000000in}{0.000000in}}{%
\pgfpathmoveto{\pgfqpoint{-0.000000in}{0.000000in}}%
\pgfpathlineto{\pgfqpoint{-0.048611in}{0.000000in}}%
\pgfusepath{stroke,fill}%
}%
\begin{pgfscope}%
\pgfsys@transformshift{1.250000in}{2.508000in}%
\pgfsys@useobject{currentmarker}{}%
\end{pgfscope}%
\end{pgfscope}%
\begin{pgfscope}%
\definecolor{textcolor}{rgb}{0.000000,0.000000,0.000000}%
\pgfsetstrokecolor{textcolor}%
\pgfsetfillcolor{textcolor}%
\pgftext[x=0.867364in, y=2.424667in, left, base]{\color{textcolor}\rmfamily\fontsize{16.000000}{19.200000}\selectfont \(\displaystyle {0.4}\)}%
\end{pgfscope}%
\begin{pgfscope}%
\pgfsetbuttcap%
\pgfsetroundjoin%
\definecolor{currentfill}{rgb}{0.000000,0.000000,0.000000}%
\pgfsetfillcolor{currentfill}%
\pgfsetlinewidth{0.803000pt}%
\definecolor{currentstroke}{rgb}{0.000000,0.000000,0.000000}%
\pgfsetstrokecolor{currentstroke}%
\pgfsetdash{}{0pt}%
\pgfsys@defobject{currentmarker}{\pgfqpoint{-0.048611in}{0.000000in}}{\pgfqpoint{-0.000000in}{0.000000in}}{%
\pgfpathmoveto{\pgfqpoint{-0.000000in}{0.000000in}}%
\pgfpathlineto{\pgfqpoint{-0.048611in}{0.000000in}}%
\pgfusepath{stroke,fill}%
}%
\begin{pgfscope}%
\pgfsys@transformshift{1.250000in}{3.432000in}%
\pgfsys@useobject{currentmarker}{}%
\end{pgfscope}%
\end{pgfscope}%
\begin{pgfscope}%
\definecolor{textcolor}{rgb}{0.000000,0.000000,0.000000}%
\pgfsetstrokecolor{textcolor}%
\pgfsetfillcolor{textcolor}%
\pgftext[x=0.867364in, y=3.348667in, left, base]{\color{textcolor}\rmfamily\fontsize{16.000000}{19.200000}\selectfont \(\displaystyle {0.6}\)}%
\end{pgfscope}%
\begin{pgfscope}%
\pgfsetbuttcap%
\pgfsetroundjoin%
\definecolor{currentfill}{rgb}{0.000000,0.000000,0.000000}%
\pgfsetfillcolor{currentfill}%
\pgfsetlinewidth{0.803000pt}%
\definecolor{currentstroke}{rgb}{0.000000,0.000000,0.000000}%
\pgfsetstrokecolor{currentstroke}%
\pgfsetdash{}{0pt}%
\pgfsys@defobject{currentmarker}{\pgfqpoint{-0.048611in}{0.000000in}}{\pgfqpoint{-0.000000in}{0.000000in}}{%
\pgfpathmoveto{\pgfqpoint{-0.000000in}{0.000000in}}%
\pgfpathlineto{\pgfqpoint{-0.048611in}{0.000000in}}%
\pgfusepath{stroke,fill}%
}%
\begin{pgfscope}%
\pgfsys@transformshift{1.250000in}{4.356000in}%
\pgfsys@useobject{currentmarker}{}%
\end{pgfscope}%
\end{pgfscope}%
\begin{pgfscope}%
\definecolor{textcolor}{rgb}{0.000000,0.000000,0.000000}%
\pgfsetstrokecolor{textcolor}%
\pgfsetfillcolor{textcolor}%
\pgftext[x=0.867364in, y=4.272667in, left, base]{\color{textcolor}\rmfamily\fontsize{16.000000}{19.200000}\selectfont \(\displaystyle {0.8}\)}%
\end{pgfscope}%
\begin{pgfscope}%
\pgfsetbuttcap%
\pgfsetroundjoin%
\definecolor{currentfill}{rgb}{0.000000,0.000000,0.000000}%
\pgfsetfillcolor{currentfill}%
\pgfsetlinewidth{0.803000pt}%
\definecolor{currentstroke}{rgb}{0.000000,0.000000,0.000000}%
\pgfsetstrokecolor{currentstroke}%
\pgfsetdash{}{0pt}%
\pgfsys@defobject{currentmarker}{\pgfqpoint{-0.048611in}{0.000000in}}{\pgfqpoint{-0.000000in}{0.000000in}}{%
\pgfpathmoveto{\pgfqpoint{-0.000000in}{0.000000in}}%
\pgfpathlineto{\pgfqpoint{-0.048611in}{0.000000in}}%
\pgfusepath{stroke,fill}%
}%
\begin{pgfscope}%
\pgfsys@transformshift{1.250000in}{5.280000in}%
\pgfsys@useobject{currentmarker}{}%
\end{pgfscope}%
\end{pgfscope}%
\begin{pgfscope}%
\definecolor{textcolor}{rgb}{0.000000,0.000000,0.000000}%
\pgfsetstrokecolor{textcolor}%
\pgfsetfillcolor{textcolor}%
\pgftext[x=0.867364in, y=5.196667in, left, base]{\color{textcolor}\rmfamily\fontsize{16.000000}{19.200000}\selectfont \(\displaystyle {1.0}\)}%
\end{pgfscope}%
\begin{pgfscope}%
\definecolor{textcolor}{rgb}{0.000000,0.000000,0.000000}%
\pgfsetstrokecolor{textcolor}%
\pgfsetfillcolor{textcolor}%
\pgftext[x=0.811809in,y=2.970000in,,bottom,rotate=90.000000]{\color{textcolor}\rmfamily\fontsize{16.000000}{19.200000}\selectfont max \(\displaystyle F\)}%
\end{pgfscope}%
\begin{pgfscope}%
\pgfpathrectangle{\pgfqpoint{1.250000in}{0.660000in}}{\pgfqpoint{7.750000in}{4.620000in}}%
\pgfusepath{clip}%
\pgfsetrectcap%
\pgfsetroundjoin%
\pgfsetlinewidth{3.011250pt}%
\definecolor{currentstroke}{rgb}{0.121569,0.466667,0.705882}%
\pgfsetstrokecolor{currentstroke}%
\pgfsetdash{}{0pt}%
\pgfpathmoveto{\pgfqpoint{1.268575in}{5.290000in}}%
\pgfpathlineto{\pgfqpoint{1.273250in}{5.058052in}}%
\pgfpathlineto{\pgfqpoint{1.281000in}{4.750694in}}%
\pgfpathlineto{\pgfqpoint{1.288750in}{4.665030in}}%
\pgfpathlineto{\pgfqpoint{1.296500in}{4.507699in}}%
\pgfpathlineto{\pgfqpoint{1.304250in}{4.409740in}}%
\pgfpathlineto{\pgfqpoint{1.312000in}{4.360033in}}%
\pgfpathlineto{\pgfqpoint{1.319750in}{4.273298in}}%
\pgfpathlineto{\pgfqpoint{1.327500in}{4.214430in}}%
\pgfpathlineto{\pgfqpoint{1.335250in}{4.188076in}}%
\pgfpathlineto{\pgfqpoint{1.350750in}{4.083513in}}%
\pgfpathlineto{\pgfqpoint{1.358500in}{4.072284in}}%
\pgfpathlineto{\pgfqpoint{1.366250in}{4.039889in}}%
\pgfpathlineto{\pgfqpoint{1.374000in}{3.993694in}}%
\pgfpathlineto{\pgfqpoint{1.381750in}{3.985403in}}%
\pgfpathlineto{\pgfqpoint{1.389500in}{3.967042in}}%
\pgfpathlineto{\pgfqpoint{1.397250in}{3.935861in}}%
\pgfpathlineto{\pgfqpoint{1.405000in}{3.915147in}}%
\pgfpathlineto{\pgfqpoint{1.412750in}{3.907272in}}%
\pgfpathlineto{\pgfqpoint{1.420500in}{3.887483in}}%
\pgfpathlineto{\pgfqpoint{1.428250in}{3.858880in}}%
\pgfpathlineto{\pgfqpoint{1.436000in}{3.855517in}}%
\pgfpathlineto{\pgfqpoint{1.443750in}{3.844841in}}%
\pgfpathlineto{\pgfqpoint{1.459250in}{3.808796in}}%
\pgfpathlineto{\pgfqpoint{1.467000in}{3.805689in}}%
\pgfpathlineto{\pgfqpoint{1.474750in}{3.794143in}}%
\pgfpathlineto{\pgfqpoint{1.482500in}{3.775877in}}%
\pgfpathlineto{\pgfqpoint{1.490250in}{3.768580in}}%
\pgfpathlineto{\pgfqpoint{1.498000in}{3.763786in}}%
\pgfpathlineto{\pgfqpoint{1.505750in}{3.752262in}}%
\pgfpathlineto{\pgfqpoint{1.513500in}{3.735298in}}%
\pgfpathlineto{\pgfqpoint{1.521250in}{3.733692in}}%
\pgfpathlineto{\pgfqpoint{1.529000in}{3.728189in}}%
\pgfpathlineto{\pgfqpoint{1.536750in}{3.717125in}}%
\pgfpathlineto{\pgfqpoint{1.544500in}{3.703346in}}%
\pgfpathlineto{\pgfqpoint{1.552250in}{3.703280in}}%
\pgfpathlineto{\pgfqpoint{1.560000in}{3.697622in}}%
\pgfpathlineto{\pgfqpoint{1.575500in}{3.677261in}}%
\pgfpathlineto{\pgfqpoint{1.583250in}{3.676586in}}%
\pgfpathlineto{\pgfqpoint{1.591000in}{3.671098in}}%
\pgfpathlineto{\pgfqpoint{1.598750in}{3.661502in}}%
\pgfpathlineto{\pgfqpoint{1.606500in}{3.653868in}}%
\pgfpathlineto{\pgfqpoint{1.614250in}{3.652971in}}%
\pgfpathlineto{\pgfqpoint{1.622000in}{3.647852in}}%
\pgfpathlineto{\pgfqpoint{1.629750in}{3.639097in}}%
\pgfpathlineto{\pgfqpoint{1.637500in}{3.632772in}}%
\pgfpathlineto{\pgfqpoint{1.645250in}{3.631908in}}%
\pgfpathlineto{\pgfqpoint{1.653000in}{3.627283in}}%
\pgfpathlineto{\pgfqpoint{1.660750in}{3.619390in}}%
\pgfpathlineto{\pgfqpoint{1.668500in}{3.613620in}}%
\pgfpathlineto{\pgfqpoint{1.676250in}{3.612965in}}%
\pgfpathlineto{\pgfqpoint{1.684000in}{3.608914in}}%
\pgfpathlineto{\pgfqpoint{1.699500in}{3.596108in}}%
\pgfpathlineto{\pgfqpoint{1.707250in}{3.595783in}}%
\pgfpathlineto{\pgfqpoint{1.715000in}{3.592358in}}%
\pgfpathlineto{\pgfqpoint{1.730500in}{3.579971in}}%
\pgfpathlineto{\pgfqpoint{1.738250in}{3.580062in}}%
\pgfpathlineto{\pgfqpoint{1.746000in}{3.577298in}}%
\pgfpathlineto{\pgfqpoint{1.753750in}{3.571994in}}%
\pgfpathlineto{\pgfqpoint{1.761500in}{3.564977in}}%
\pgfpathlineto{\pgfqpoint{1.769250in}{3.565550in}}%
\pgfpathlineto{\pgfqpoint{1.777000in}{3.563470in}}%
\pgfpathlineto{\pgfqpoint{1.784750in}{3.559015in}}%
\pgfpathlineto{\pgfqpoint{1.792500in}{3.552453in}}%
\pgfpathlineto{\pgfqpoint{1.808000in}{3.550650in}}%
\pgfpathlineto{\pgfqpoint{1.815750in}{3.547034in}}%
\pgfpathlineto{\pgfqpoint{1.823500in}{3.541420in}}%
\pgfpathlineto{\pgfqpoint{1.831250in}{3.539313in}}%
\pgfpathlineto{\pgfqpoint{1.839000in}{3.538646in}}%
\pgfpathlineto{\pgfqpoint{1.846750in}{3.535860in}}%
\pgfpathlineto{\pgfqpoint{1.862250in}{3.527235in}}%
\pgfpathlineto{\pgfqpoint{1.870000in}{3.527289in}}%
\pgfpathlineto{\pgfqpoint{1.877750in}{3.525327in}}%
\pgfpathlineto{\pgfqpoint{1.885500in}{3.521540in}}%
\pgfpathlineto{\pgfqpoint{1.893250in}{3.516114in}}%
\pgfpathlineto{\pgfqpoint{1.901000in}{3.516431in}}%
\pgfpathlineto{\pgfqpoint{1.908750in}{3.515287in}}%
\pgfpathlineto{\pgfqpoint{1.916500in}{3.512388in}}%
\pgfpathlineto{\pgfqpoint{1.924250in}{3.507902in}}%
\pgfpathlineto{\pgfqpoint{1.932000in}{3.505940in}}%
\pgfpathlineto{\pgfqpoint{1.939750in}{3.505608in}}%
\pgfpathlineto{\pgfqpoint{1.947500in}{3.503584in}}%
\pgfpathlineto{\pgfqpoint{1.963000in}{3.495695in}}%
\pgfpathlineto{\pgfqpoint{1.970750in}{3.496170in}}%
\pgfpathlineto{\pgfqpoint{1.978500in}{3.495010in}}%
\pgfpathlineto{\pgfqpoint{1.986250in}{3.492351in}}%
\pgfpathlineto{\pgfqpoint{1.994000in}{3.488328in}}%
\pgfpathlineto{\pgfqpoint{2.001750in}{3.486865in}}%
\pgfpathlineto{\pgfqpoint{2.009500in}{3.486560in}}%
\pgfpathlineto{\pgfqpoint{2.017250in}{3.484795in}}%
\pgfpathlineto{\pgfqpoint{2.032750in}{3.477591in}}%
\pgfpathlineto{\pgfqpoint{2.040500in}{3.478134in}}%
\pgfpathlineto{\pgfqpoint{2.048250in}{3.477252in}}%
\pgfpathlineto{\pgfqpoint{2.056000in}{3.475056in}}%
\pgfpathlineto{\pgfqpoint{2.063750in}{3.471657in}}%
\pgfpathlineto{\pgfqpoint{2.071500in}{3.469637in}}%
\pgfpathlineto{\pgfqpoint{2.079250in}{3.469630in}}%
\pgfpathlineto{\pgfqpoint{2.087000in}{3.468333in}}%
\pgfpathlineto{\pgfqpoint{2.110250in}{3.461842in}}%
\pgfpathlineto{\pgfqpoint{2.118000in}{3.461435in}}%
\pgfpathlineto{\pgfqpoint{2.133500in}{3.457184in}}%
\pgfpathlineto{\pgfqpoint{2.141250in}{3.453804in}}%
\pgfpathlineto{\pgfqpoint{2.149000in}{3.454281in}}%
\pgfpathlineto{\pgfqpoint{2.164500in}{3.451833in}}%
\pgfpathlineto{\pgfqpoint{2.180000in}{3.446791in}}%
\pgfpathlineto{\pgfqpoint{2.195500in}{3.446136in}}%
\pgfpathlineto{\pgfqpoint{2.218750in}{3.439983in}}%
\pgfpathlineto{\pgfqpoint{2.234250in}{3.439057in}}%
\pgfpathlineto{\pgfqpoint{2.257500in}{3.433400in}}%
\pgfpathlineto{\pgfqpoint{2.273000in}{3.432350in}}%
\pgfpathlineto{\pgfqpoint{2.296250in}{3.427046in}}%
\pgfpathlineto{\pgfqpoint{2.311750in}{3.425995in}}%
\pgfpathlineto{\pgfqpoint{2.335000in}{3.420913in}}%
\pgfpathlineto{\pgfqpoint{2.350500in}{3.419964in}}%
\pgfpathlineto{\pgfqpoint{2.373750in}{3.414982in}}%
\pgfpathlineto{\pgfqpoint{2.389250in}{3.414225in}}%
\pgfpathlineto{\pgfqpoint{2.420250in}{3.409351in}}%
\pgfpathlineto{\pgfqpoint{2.435750in}{3.407439in}}%
\pgfpathlineto{\pgfqpoint{2.451250in}{3.403621in}}%
\pgfpathlineto{\pgfqpoint{2.466750in}{3.403466in}}%
\pgfpathlineto{\pgfqpoint{2.482250in}{3.400679in}}%
\pgfpathlineto{\pgfqpoint{2.490000in}{3.398399in}}%
\pgfpathlineto{\pgfqpoint{2.505500in}{3.398359in}}%
\pgfpathlineto{\pgfqpoint{2.521000in}{3.396073in}}%
\pgfpathlineto{\pgfqpoint{2.528750in}{3.394077in}}%
\pgfpathlineto{\pgfqpoint{2.567500in}{3.389931in}}%
\pgfpathlineto{\pgfqpoint{2.575250in}{3.388181in}}%
\pgfpathlineto{\pgfqpoint{2.590750in}{3.388149in}}%
\pgfpathlineto{\pgfqpoint{2.606250in}{3.385904in}}%
\pgfpathlineto{\pgfqpoint{2.614000in}{3.384031in}}%
\pgfpathlineto{\pgfqpoint{2.645000in}{3.381943in}}%
\pgfpathlineto{\pgfqpoint{2.660500in}{3.378627in}}%
\pgfpathlineto{\pgfqpoint{2.676000in}{3.378705in}}%
\pgfpathlineto{\pgfqpoint{2.691500in}{3.376805in}}%
\pgfpathlineto{\pgfqpoint{2.707000in}{3.374228in}}%
\pgfpathlineto{\pgfqpoint{2.722500in}{3.373982in}}%
\pgfpathlineto{\pgfqpoint{2.738000in}{3.371905in}}%
\pgfpathlineto{\pgfqpoint{2.745750in}{3.370241in}}%
\pgfpathlineto{\pgfqpoint{2.784500in}{3.367259in}}%
\pgfpathlineto{\pgfqpoint{2.792250in}{3.365599in}}%
\pgfpathlineto{\pgfqpoint{2.815500in}{3.365014in}}%
\pgfpathlineto{\pgfqpoint{2.893000in}{3.357235in}}%
\pgfpathlineto{\pgfqpoint{2.916250in}{3.355931in}}%
\pgfpathlineto{\pgfqpoint{2.939500in}{3.353193in}}%
\pgfpathlineto{\pgfqpoint{2.962750in}{3.352105in}}%
\pgfpathlineto{\pgfqpoint{2.986000in}{3.349213in}}%
\pgfpathlineto{\pgfqpoint{3.009250in}{3.348423in}}%
\pgfpathlineto{\pgfqpoint{3.040250in}{3.345443in}}%
\pgfpathlineto{\pgfqpoint{3.063500in}{3.344111in}}%
\pgfpathlineto{\pgfqpoint{3.086750in}{3.341632in}}%
\pgfpathlineto{\pgfqpoint{3.110000in}{3.340774in}}%
\pgfpathlineto{\pgfqpoint{3.141000in}{3.337969in}}%
\pgfpathlineto{\pgfqpoint{3.164250in}{3.336833in}}%
\pgfpathlineto{\pgfqpoint{3.195250in}{3.334333in}}%
\pgfpathlineto{\pgfqpoint{3.218500in}{3.333055in}}%
\pgfpathlineto{\pgfqpoint{3.241750in}{3.330720in}}%
\pgfpathlineto{\pgfqpoint{3.265000in}{3.330112in}}%
\pgfpathlineto{\pgfqpoint{3.311500in}{3.327051in}}%
\pgfpathlineto{\pgfqpoint{3.334750in}{3.325198in}}%
\pgfpathlineto{\pgfqpoint{3.350250in}{3.323745in}}%
\pgfpathlineto{\pgfqpoint{3.373500in}{3.323288in}}%
\pgfpathlineto{\pgfqpoint{3.474250in}{3.317026in}}%
\pgfpathlineto{\pgfqpoint{3.505250in}{3.315284in}}%
\pgfpathlineto{\pgfqpoint{3.528500in}{3.313731in}}%
\pgfpathlineto{\pgfqpoint{3.559500in}{3.312516in}}%
\pgfpathlineto{\pgfqpoint{3.590500in}{3.310456in}}%
\pgfpathlineto{\pgfqpoint{3.621500in}{3.309227in}}%
\pgfpathlineto{\pgfqpoint{3.652500in}{3.307213in}}%
\pgfpathlineto{\pgfqpoint{3.683500in}{3.306093in}}%
\pgfpathlineto{\pgfqpoint{3.722250in}{3.303962in}}%
\pgfpathlineto{\pgfqpoint{3.753250in}{3.302553in}}%
\pgfpathlineto{\pgfqpoint{3.784250in}{3.300805in}}%
\pgfpathlineto{\pgfqpoint{3.815250in}{3.299733in}}%
\pgfpathlineto{\pgfqpoint{3.861750in}{3.297513in}}%
\pgfpathlineto{\pgfqpoint{3.900500in}{3.295381in}}%
\pgfpathlineto{\pgfqpoint{3.916000in}{3.294468in}}%
\pgfpathlineto{\pgfqpoint{3.954750in}{3.293484in}}%
\pgfpathlineto{\pgfqpoint{4.009000in}{3.291124in}}%
\pgfpathlineto{\pgfqpoint{4.047750in}{3.289080in}}%
\pgfpathlineto{\pgfqpoint{4.071000in}{3.288184in}}%
\pgfpathlineto{\pgfqpoint{4.109750in}{3.286971in}}%
\pgfpathlineto{\pgfqpoint{4.156250in}{3.284969in}}%
\pgfpathlineto{\pgfqpoint{4.202750in}{3.283052in}}%
\pgfpathlineto{\pgfqpoint{4.233750in}{3.281855in}}%
\pgfpathlineto{\pgfqpoint{4.280250in}{3.280305in}}%
\pgfpathlineto{\pgfqpoint{4.326750in}{3.278580in}}%
\pgfpathlineto{\pgfqpoint{4.373250in}{3.276849in}}%
\pgfpathlineto{\pgfqpoint{4.412000in}{3.275426in}}%
\pgfpathlineto{\pgfqpoint{4.466250in}{3.273634in}}%
\pgfpathlineto{\pgfqpoint{4.512750in}{3.272088in}}%
\pgfpathlineto{\pgfqpoint{4.567000in}{3.270254in}}%
\pgfpathlineto{\pgfqpoint{4.613500in}{3.268763in}}%
\pgfpathlineto{\pgfqpoint{4.675500in}{3.266759in}}%
\pgfpathlineto{\pgfqpoint{4.722000in}{3.265342in}}%
\pgfpathlineto{\pgfqpoint{4.791750in}{3.263200in}}%
\pgfpathlineto{\pgfqpoint{4.846000in}{3.261702in}}%
\pgfpathlineto{\pgfqpoint{4.923500in}{3.259293in}}%
\pgfpathlineto{\pgfqpoint{4.970000in}{3.258063in}}%
\pgfpathlineto{\pgfqpoint{5.055250in}{3.255781in}}%
\pgfpathlineto{\pgfqpoint{5.132750in}{3.253815in}}%
\pgfpathlineto{\pgfqpoint{5.241250in}{3.250650in}}%
\pgfpathlineto{\pgfqpoint{5.295500in}{3.249557in}}%
\pgfpathlineto{\pgfqpoint{5.419500in}{3.246508in}}%
\pgfpathlineto{\pgfqpoint{5.528000in}{3.244090in}}%
\pgfpathlineto{\pgfqpoint{5.721750in}{3.239510in}}%
\pgfpathlineto{\pgfqpoint{5.861250in}{3.236760in}}%
\pgfpathlineto{\pgfqpoint{7.225250in}{3.213359in}}%
\pgfpathlineto{\pgfqpoint{8.775250in}{3.194134in}}%
\pgfpathlineto{\pgfqpoint{8.891500in}{3.192766in}}%
\pgfpathlineto{\pgfqpoint{8.930250in}{3.192592in}}%
\pgfpathlineto{\pgfqpoint{9.000000in}{3.191879in}}%
\pgfpathlineto{\pgfqpoint{9.000000in}{3.191879in}}%
\pgfusepath{stroke}%
\end{pgfscope}%
\begin{pgfscope}%
\pgfpathrectangle{\pgfqpoint{1.250000in}{0.660000in}}{\pgfqpoint{7.750000in}{4.620000in}}%
\pgfusepath{clip}%
\pgfsetrectcap%
\pgfsetroundjoin%
\pgfsetlinewidth{3.011250pt}%
\definecolor{currentstroke}{rgb}{1.000000,0.498039,0.054902}%
\pgfsetstrokecolor{currentstroke}%
\pgfsetdash{}{0pt}%
\pgfpathmoveto{\pgfqpoint{1.265500in}{3.898545in}}%
\pgfpathlineto{\pgfqpoint{1.273250in}{3.779472in}}%
\pgfpathlineto{\pgfqpoint{1.281000in}{3.699938in}}%
\pgfpathlineto{\pgfqpoint{1.296500in}{3.604670in}}%
\pgfpathlineto{\pgfqpoint{1.304250in}{3.574120in}}%
\pgfpathlineto{\pgfqpoint{1.312000in}{3.550493in}}%
\pgfpathlineto{\pgfqpoint{1.319750in}{3.521475in}}%
\pgfpathlineto{\pgfqpoint{1.327500in}{3.502908in}}%
\pgfpathlineto{\pgfqpoint{1.335250in}{3.489080in}}%
\pgfpathlineto{\pgfqpoint{1.343000in}{3.470165in}}%
\pgfpathlineto{\pgfqpoint{1.350750in}{3.454523in}}%
\pgfpathlineto{\pgfqpoint{1.358500in}{3.446524in}}%
\pgfpathlineto{\pgfqpoint{1.366250in}{3.433937in}}%
\pgfpathlineto{\pgfqpoint{1.374000in}{3.418069in}}%
\pgfpathlineto{\pgfqpoint{1.381750in}{3.414103in}}%
\pgfpathlineto{\pgfqpoint{1.389500in}{3.405999in}}%
\pgfpathlineto{\pgfqpoint{1.397250in}{3.394513in}}%
\pgfpathlineto{\pgfqpoint{1.405000in}{3.387699in}}%
\pgfpathlineto{\pgfqpoint{1.412750in}{3.383027in}}%
\pgfpathlineto{\pgfqpoint{1.420500in}{3.375166in}}%
\pgfpathlineto{\pgfqpoint{1.428250in}{3.365085in}}%
\pgfpathlineto{\pgfqpoint{1.436000in}{3.363178in}}%
\pgfpathlineto{\pgfqpoint{1.443750in}{3.358271in}}%
\pgfpathlineto{\pgfqpoint{1.451500in}{3.350813in}}%
\pgfpathlineto{\pgfqpoint{1.459250in}{3.345342in}}%
\pgfpathlineto{\pgfqpoint{1.467000in}{3.342926in}}%
\pgfpathlineto{\pgfqpoint{1.474750in}{3.338043in}}%
\pgfpathlineto{\pgfqpoint{1.482500in}{3.331052in}}%
\pgfpathlineto{\pgfqpoint{1.498000in}{3.325903in}}%
\pgfpathlineto{\pgfqpoint{1.505750in}{3.321185in}}%
\pgfpathlineto{\pgfqpoint{1.513500in}{3.314715in}}%
\pgfpathlineto{\pgfqpoint{1.521250in}{3.314068in}}%
\pgfpathlineto{\pgfqpoint{1.529000in}{3.311386in}}%
\pgfpathlineto{\pgfqpoint{1.544500in}{3.302305in}}%
\pgfpathlineto{\pgfqpoint{1.552250in}{3.301467in}}%
\pgfpathlineto{\pgfqpoint{1.560000in}{3.298852in}}%
\pgfpathlineto{\pgfqpoint{1.575500in}{3.291290in}}%
\pgfpathlineto{\pgfqpoint{1.583250in}{3.290388in}}%
\pgfpathlineto{\pgfqpoint{1.591000in}{3.287913in}}%
\pgfpathlineto{\pgfqpoint{1.598750in}{3.284038in}}%
\pgfpathlineto{\pgfqpoint{1.606500in}{3.281437in}}%
\pgfpathlineto{\pgfqpoint{1.614250in}{3.280560in}}%
\pgfpathlineto{\pgfqpoint{1.622000in}{3.278271in}}%
\pgfpathlineto{\pgfqpoint{1.629750in}{3.274720in}}%
\pgfpathlineto{\pgfqpoint{1.637500in}{3.272556in}}%
\pgfpathlineto{\pgfqpoint{1.645250in}{3.271767in}}%
\pgfpathlineto{\pgfqpoint{1.653000in}{3.269695in}}%
\pgfpathlineto{\pgfqpoint{1.660750in}{3.266472in}}%
\pgfpathlineto{\pgfqpoint{1.668500in}{3.264490in}}%
\pgfpathlineto{\pgfqpoint{1.676250in}{3.263833in}}%
\pgfpathlineto{\pgfqpoint{1.691750in}{3.259108in}}%
\pgfpathlineto{\pgfqpoint{1.699500in}{3.257108in}}%
\pgfpathlineto{\pgfqpoint{1.715000in}{3.255038in}}%
\pgfpathlineto{\pgfqpoint{1.730500in}{3.250300in}}%
\pgfpathlineto{\pgfqpoint{1.746000in}{3.248683in}}%
\pgfpathlineto{\pgfqpoint{1.761500in}{3.243969in}}%
\pgfpathlineto{\pgfqpoint{1.777000in}{3.242833in}}%
\pgfpathlineto{\pgfqpoint{1.800250in}{3.238162in}}%
\pgfpathlineto{\pgfqpoint{1.815750in}{3.235817in}}%
\pgfpathlineto{\pgfqpoint{1.823500in}{3.233486in}}%
\pgfpathlineto{\pgfqpoint{1.846750in}{3.231045in}}%
\pgfpathlineto{\pgfqpoint{1.862250in}{3.227681in}}%
\pgfpathlineto{\pgfqpoint{1.877750in}{3.226545in}}%
\pgfpathlineto{\pgfqpoint{1.893250in}{3.222786in}}%
\pgfpathlineto{\pgfqpoint{1.908750in}{3.222258in}}%
\pgfpathlineto{\pgfqpoint{1.939750in}{3.218129in}}%
\pgfpathlineto{\pgfqpoint{1.955250in}{3.215664in}}%
\pgfpathlineto{\pgfqpoint{1.963000in}{3.214095in}}%
\pgfpathlineto{\pgfqpoint{1.978500in}{3.213511in}}%
\pgfpathlineto{\pgfqpoint{2.009500in}{3.209894in}}%
\pgfpathlineto{\pgfqpoint{2.025000in}{3.207734in}}%
\pgfpathlineto{\pgfqpoint{2.032750in}{3.206230in}}%
\pgfpathlineto{\pgfqpoint{2.048250in}{3.205825in}}%
\pgfpathlineto{\pgfqpoint{2.102500in}{3.199324in}}%
\pgfpathlineto{\pgfqpoint{2.125750in}{3.198234in}}%
\pgfpathlineto{\pgfqpoint{2.141250in}{3.195826in}}%
\pgfpathlineto{\pgfqpoint{2.156750in}{3.195529in}}%
\pgfpathlineto{\pgfqpoint{2.218750in}{3.189705in}}%
\pgfpathlineto{\pgfqpoint{2.234250in}{3.189155in}}%
\pgfpathlineto{\pgfqpoint{2.265250in}{3.186687in}}%
\pgfpathlineto{\pgfqpoint{2.319500in}{3.182633in}}%
\pgfpathlineto{\pgfqpoint{2.335000in}{3.181294in}}%
\pgfpathlineto{\pgfqpoint{2.358250in}{3.180032in}}%
\pgfpathlineto{\pgfqpoint{2.373750in}{3.178683in}}%
\pgfpathlineto{\pgfqpoint{2.397000in}{3.177571in}}%
\pgfpathlineto{\pgfqpoint{2.412500in}{3.176151in}}%
\pgfpathlineto{\pgfqpoint{2.435750in}{3.175233in}}%
\pgfpathlineto{\pgfqpoint{2.459000in}{3.173732in}}%
\pgfpathlineto{\pgfqpoint{2.482250in}{3.172247in}}%
\pgfpathlineto{\pgfqpoint{2.497750in}{3.171393in}}%
\pgfpathlineto{\pgfqpoint{2.521000in}{3.170205in}}%
\pgfpathlineto{\pgfqpoint{2.536500in}{3.169099in}}%
\pgfpathlineto{\pgfqpoint{2.559750in}{3.168234in}}%
\pgfpathlineto{\pgfqpoint{2.583000in}{3.166890in}}%
\pgfpathlineto{\pgfqpoint{2.606250in}{3.165694in}}%
\pgfpathlineto{\pgfqpoint{2.621750in}{3.164739in}}%
\pgfpathlineto{\pgfqpoint{2.645000in}{3.163933in}}%
\pgfpathlineto{\pgfqpoint{2.668250in}{3.162654in}}%
\pgfpathlineto{\pgfqpoint{2.691500in}{3.161646in}}%
\pgfpathlineto{\pgfqpoint{2.714750in}{3.160598in}}%
\pgfpathlineto{\pgfqpoint{2.776750in}{3.157962in}}%
\pgfpathlineto{\pgfqpoint{2.807750in}{3.156637in}}%
\pgfpathlineto{\pgfqpoint{2.877500in}{3.153577in}}%
\pgfpathlineto{\pgfqpoint{2.893000in}{3.152972in}}%
\pgfpathlineto{\pgfqpoint{2.924000in}{3.151822in}}%
\pgfpathlineto{\pgfqpoint{2.939500in}{3.151165in}}%
\pgfpathlineto{\pgfqpoint{2.970500in}{3.150155in}}%
\pgfpathlineto{\pgfqpoint{2.993750in}{3.149384in}}%
\pgfpathlineto{\pgfqpoint{3.063500in}{3.147027in}}%
\pgfpathlineto{\pgfqpoint{3.102250in}{3.145797in}}%
\pgfpathlineto{\pgfqpoint{3.365750in}{3.137830in}}%
\pgfpathlineto{\pgfqpoint{3.575000in}{3.132092in}}%
\pgfpathlineto{\pgfqpoint{3.714500in}{3.128954in}}%
\pgfpathlineto{\pgfqpoint{3.761000in}{3.127966in}}%
\pgfpathlineto{\pgfqpoint{3.799750in}{3.127340in}}%
\pgfpathlineto{\pgfqpoint{4.435250in}{3.115685in}}%
\pgfpathlineto{\pgfqpoint{5.109500in}{3.106169in}}%
\pgfpathlineto{\pgfqpoint{5.349750in}{3.103527in}}%
\pgfpathlineto{\pgfqpoint{6.117000in}{3.095860in}}%
\pgfpathlineto{\pgfqpoint{7.178750in}{3.087485in}}%
\pgfpathlineto{\pgfqpoint{7.566250in}{3.085048in}}%
\pgfpathlineto{\pgfqpoint{8.000250in}{3.082445in}}%
\pgfpathlineto{\pgfqpoint{8.511750in}{3.079546in}}%
\pgfpathlineto{\pgfqpoint{8.620250in}{3.079075in}}%
\pgfpathlineto{\pgfqpoint{8.883750in}{3.077767in}}%
\pgfpathlineto{\pgfqpoint{9.000000in}{3.077253in}}%
\pgfpathlineto{\pgfqpoint{9.000000in}{3.077253in}}%
\pgfusepath{stroke}%
\end{pgfscope}%
\begin{pgfscope}%
\pgfpathrectangle{\pgfqpoint{1.250000in}{0.660000in}}{\pgfqpoint{7.750000in}{4.620000in}}%
\pgfusepath{clip}%
\pgfsetbuttcap%
\pgfsetroundjoin%
\pgfsetlinewidth{1.505625pt}%
\definecolor{currentstroke}{rgb}{0.121569,0.466667,0.705882}%
\pgfsetstrokecolor{currentstroke}%
\pgfsetdash{{5.550000pt}{2.400000pt}}{0.000000pt}%
\pgfpathmoveto{\pgfqpoint{1.250000in}{3.740000in}}%
\pgfpathlineto{\pgfqpoint{9.000000in}{3.740000in}}%
\pgfusepath{stroke}%
\end{pgfscope}%
\begin{pgfscope}%
\pgfpathrectangle{\pgfqpoint{1.250000in}{0.660000in}}{\pgfqpoint{7.750000in}{4.620000in}}%
\pgfusepath{clip}%
\pgfsetbuttcap%
\pgfsetroundjoin%
\definecolor{currentfill}{rgb}{0.172549,0.627451,0.172549}%
\pgfsetfillcolor{currentfill}%
\pgfsetlinewidth{1.003750pt}%
\definecolor{currentstroke}{rgb}{0.172549,0.627451,0.172549}%
\pgfsetstrokecolor{currentstroke}%
\pgfsetdash{}{0pt}%
\pgfsys@defobject{currentmarker}{\pgfqpoint{-0.041667in}{-0.041667in}}{\pgfqpoint{0.041667in}{0.041667in}}{%
\pgfpathmoveto{\pgfqpoint{0.000000in}{-0.041667in}}%
\pgfpathcurveto{\pgfqpoint{0.011050in}{-0.041667in}}{\pgfqpoint{0.021649in}{-0.037276in}}{\pgfqpoint{0.029463in}{-0.029463in}}%
\pgfpathcurveto{\pgfqpoint{0.037276in}{-0.021649in}}{\pgfqpoint{0.041667in}{-0.011050in}}{\pgfqpoint{0.041667in}{0.000000in}}%
\pgfpathcurveto{\pgfqpoint{0.041667in}{0.011050in}}{\pgfqpoint{0.037276in}{0.021649in}}{\pgfqpoint{0.029463in}{0.029463in}}%
\pgfpathcurveto{\pgfqpoint{0.021649in}{0.037276in}}{\pgfqpoint{0.011050in}{0.041667in}}{\pgfqpoint{0.000000in}{0.041667in}}%
\pgfpathcurveto{\pgfqpoint{-0.011050in}{0.041667in}}{\pgfqpoint{-0.021649in}{0.037276in}}{\pgfqpoint{-0.029463in}{0.029463in}}%
\pgfpathcurveto{\pgfqpoint{-0.037276in}{0.021649in}}{\pgfqpoint{-0.041667in}{0.011050in}}{\pgfqpoint{-0.041667in}{0.000000in}}%
\pgfpathcurveto{\pgfqpoint{-0.041667in}{-0.011050in}}{\pgfqpoint{-0.037276in}{-0.021649in}}{\pgfqpoint{-0.029463in}{-0.029463in}}%
\pgfpathcurveto{\pgfqpoint{-0.021649in}{-0.037276in}}{\pgfqpoint{-0.011050in}{-0.041667in}}{\pgfqpoint{0.000000in}{-0.041667in}}%
\pgfpathlineto{\pgfqpoint{0.000000in}{-0.041667in}}%
\pgfpathclose%
\pgfusepath{stroke,fill}%
}%
\begin{pgfscope}%
\pgfsys@transformshift{1.939750in}{3.740000in}%
\pgfsys@useobject{currentmarker}{}%
\end{pgfscope}%
\end{pgfscope}%
\begin{pgfscope}%
\pgfsetrectcap%
\pgfsetmiterjoin%
\pgfsetlinewidth{0.803000pt}%
\definecolor{currentstroke}{rgb}{0.000000,0.000000,0.000000}%
\pgfsetstrokecolor{currentstroke}%
\pgfsetdash{}{0pt}%
\pgfpathmoveto{\pgfqpoint{1.250000in}{0.660000in}}%
\pgfpathlineto{\pgfqpoint{1.250000in}{5.280000in}}%
\pgfusepath{stroke}%
\end{pgfscope}%
\begin{pgfscope}%
\pgfsetrectcap%
\pgfsetmiterjoin%
\pgfsetlinewidth{0.803000pt}%
\definecolor{currentstroke}{rgb}{0.000000,0.000000,0.000000}%
\pgfsetstrokecolor{currentstroke}%
\pgfsetdash{}{0pt}%
\pgfpathmoveto{\pgfqpoint{9.000000in}{0.660000in}}%
\pgfpathlineto{\pgfqpoint{9.000000in}{5.280000in}}%
\pgfusepath{stroke}%
\end{pgfscope}%
\begin{pgfscope}%
\pgfsetrectcap%
\pgfsetmiterjoin%
\pgfsetlinewidth{0.803000pt}%
\definecolor{currentstroke}{rgb}{0.000000,0.000000,0.000000}%
\pgfsetstrokecolor{currentstroke}%
\pgfsetdash{}{0pt}%
\pgfpathmoveto{\pgfqpoint{1.250000in}{0.660000in}}%
\pgfpathlineto{\pgfqpoint{9.000000in}{0.660000in}}%
\pgfusepath{stroke}%
\end{pgfscope}%
\begin{pgfscope}%
\pgfsetrectcap%
\pgfsetmiterjoin%
\pgfsetlinewidth{0.803000pt}%
\definecolor{currentstroke}{rgb}{0.000000,0.000000,0.000000}%
\pgfsetstrokecolor{currentstroke}%
\pgfsetdash{}{0pt}%
\pgfpathmoveto{\pgfqpoint{1.250000in}{5.280000in}}%
\pgfpathlineto{\pgfqpoint{9.000000in}{5.280000in}}%
\pgfusepath{stroke}%
\end{pgfscope}%
\begin{pgfscope}%
\definecolor{textcolor}{rgb}{0.000000,0.000000,1.000000}%
\pgfsetstrokecolor{textcolor}%
\pgfsetfillcolor{textcolor}%
\pgftext[x=1.978500in,y=3.832400in,left,base]{\color{textcolor}\rmfamily\fontsize{16.000000}{19.200000}\selectfont \(\displaystyle N = 89\)}%
\end{pgfscope}%
\begin{pgfscope}%
\pgfsetbuttcap%
\pgfsetmiterjoin%
\definecolor{currentfill}{rgb}{1.000000,1.000000,1.000000}%
\pgfsetfillcolor{currentfill}%
\pgfsetfillopacity{0.800000}%
\pgfsetlinewidth{1.003750pt}%
\definecolor{currentstroke}{rgb}{0.800000,0.800000,0.800000}%
\pgfsetstrokecolor{currentstroke}%
\pgfsetstrokeopacity{0.800000}%
\pgfsetdash{}{0pt}%
\pgfpathmoveto{\pgfqpoint{5.879782in}{3.777732in}}%
\pgfpathlineto{\pgfqpoint{8.844444in}{3.777732in}}%
\pgfpathquadraticcurveto{\pgfqpoint{8.888889in}{3.777732in}}{\pgfqpoint{8.888889in}{3.822177in}}%
\pgfpathlineto{\pgfqpoint{8.888889in}{5.124444in}}%
\pgfpathquadraticcurveto{\pgfqpoint{8.888889in}{5.168889in}}{\pgfqpoint{8.844444in}{5.168889in}}%
\pgfpathlineto{\pgfqpoint{5.879782in}{5.168889in}}%
\pgfpathquadraticcurveto{\pgfqpoint{5.835338in}{5.168889in}}{\pgfqpoint{5.835338in}{5.124444in}}%
\pgfpathlineto{\pgfqpoint{5.835338in}{3.822177in}}%
\pgfpathquadraticcurveto{\pgfqpoint{5.835338in}{3.777732in}}{\pgfqpoint{5.879782in}{3.777732in}}%
\pgfpathlineto{\pgfqpoint{5.879782in}{3.777732in}}%
\pgfpathclose%
\pgfusepath{stroke,fill}%
\end{pgfscope}%
\begin{pgfscope}%
\pgfsetrectcap%
\pgfsetroundjoin%
\pgfsetlinewidth{3.011250pt}%
\definecolor{currentstroke}{rgb}{0.121569,0.466667,0.705882}%
\pgfsetstrokecolor{currentstroke}%
\pgfsetdash{}{0pt}%
\pgfpathmoveto{\pgfqpoint{5.924227in}{4.991111in}}%
\pgfpathlineto{\pgfqpoint{6.146449in}{4.991111in}}%
\pgfpathlineto{\pgfqpoint{6.368671in}{4.991111in}}%
\pgfusepath{stroke}%
\end{pgfscope}%
\begin{pgfscope}%
\definecolor{textcolor}{rgb}{0.000000,0.000000,0.000000}%
\pgfsetstrokecolor{textcolor}%
\pgfsetfillcolor{textcolor}%
\pgftext[x=6.546449in,y=4.913333in,left,base]{\color{textcolor}\rmfamily\fontsize{16.000000}{19.200000}\selectfont Saddle}%
\end{pgfscope}%
\begin{pgfscope}%
\pgfsetrectcap%
\pgfsetroundjoin%
\pgfsetlinewidth{3.011250pt}%
\definecolor{currentstroke}{rgb}{1.000000,0.498039,0.054902}%
\pgfsetstrokecolor{currentstroke}%
\pgfsetdash{}{0pt}%
\pgfpathmoveto{\pgfqpoint{5.924227in}{4.666651in}}%
\pgfpathlineto{\pgfqpoint{6.146449in}{4.666651in}}%
\pgfpathlineto{\pgfqpoint{6.368671in}{4.666651in}}%
\pgfusepath{stroke}%
\end{pgfscope}%
\begin{pgfscope}%
\definecolor{textcolor}{rgb}{0.000000,0.000000,0.000000}%
\pgfsetstrokecolor{textcolor}%
\pgfsetfillcolor{textcolor}%
\pgftext[x=6.546449in,y=4.588874in,left,base]{\color{textcolor}\rmfamily\fontsize{16.000000}{19.200000}\selectfont Bessel}%
\end{pgfscope}%
\begin{pgfscope}%
\pgfsetbuttcap%
\pgfsetroundjoin%
\pgfsetlinewidth{1.505625pt}%
\definecolor{currentstroke}{rgb}{0.121569,0.466667,0.705882}%
\pgfsetstrokecolor{currentstroke}%
\pgfsetdash{{5.550000pt}{2.400000pt}}{0.000000pt}%
\pgfpathmoveto{\pgfqpoint{5.924227in}{4.328858in}}%
\pgfpathlineto{\pgfqpoint{6.146449in}{4.328858in}}%
\pgfpathlineto{\pgfqpoint{6.368671in}{4.328858in}}%
\pgfusepath{stroke}%
\end{pgfscope}%
\begin{pgfscope}%
\definecolor{textcolor}{rgb}{0.000000,0.000000,0.000000}%
\pgfsetstrokecolor{textcolor}%
\pgfsetfillcolor{textcolor}%
\pgftext[x=6.546449in,y=4.251081in,left,base]{\color{textcolor}\rmfamily\fontsize{16.000000}{19.200000}\selectfont Classical limit \(\displaystyle F=2/3\)}%
\end{pgfscope}%
\begin{pgfscope}%
\pgfsetbuttcap%
\pgfsetmiterjoin%
\definecolor{currentfill}{rgb}{0.000000,0.000000,0.000000}%
\pgfsetfillcolor{currentfill}%
\pgfsetlinewidth{0.000000pt}%
\definecolor{currentstroke}{rgb}{0.000000,0.000000,0.000000}%
\pgfsetstrokecolor{currentstroke}%
\pgfsetstrokeopacity{0.000000}%
\pgfsetdash{}{0pt}%
\pgfpathmoveto{\pgfqpoint{5.924227in}{3.913303in}}%
\pgfpathlineto{\pgfqpoint{6.368671in}{3.913303in}}%
\pgfpathlineto{\pgfqpoint{6.368671in}{4.068859in}}%
\pgfpathlineto{\pgfqpoint{5.924227in}{4.068859in}}%
\pgfpathlineto{\pgfqpoint{5.924227in}{3.913303in}}%
\pgfpathclose%
\pgfusepath{fill}%
\end{pgfscope}%
\begin{pgfscope}%
\definecolor{textcolor}{rgb}{0.000000,0.000000,0.000000}%
\pgfsetstrokecolor{textcolor}%
\pgfsetfillcolor{textcolor}%
\pgftext[x=6.546449in,y=3.913303in,left,base]{\color{textcolor}\rmfamily\fontsize{16.000000}{19.200000}\selectfont Numeric}%
\end{pgfscope}%
\end{pgfpicture}%
\makeatother%
\endgroup%
}
    \caption*{Figure: The plots shows the maximum fidelity $F$ of quantum communication as a function of the chain length $N$ from 2 to 1000. First figure show  in a time interval [0, $3t_c$] while second [0, $3t_c$ + 4000/J]}
\end{figure}

The resulting transition amplitude is
\begin{equation*}
    g_N(t) = \sum_{m =1} ^{N} a_m ^ 2\cos^2 k_m/2 \exp \left[i \left\{Nk_m - 4Jt (1-\cos k_m)  \right\}  \right]
\end{equation*}
Further derivation reveal that the amplitude can be separated into three case based on critical time $t_c = N/4J$
\begin{equation*}
    g_N(t) =
    \begin{cases}
        \begin{aligned}
             & \sqrt{\frac{2}{\pi}} \left[ \frac{4Jt - i \sqrt{N^2 - (4Jt)^2}}{4Jt \left( N^2 - (4Jt)^2 \right)^{1/4}} \right]
            \exp\left[ i\left( \frac{\pi N}{2} - 4Jt \right) \right]                                                           \\
             & \qquad \times \exp\left[\sqrt{N^2 - (4Jt)^2} -N \operatorname{arcosh} \left(\frac{N}{4Jt}\right)  \right]
        \end{aligned}
         & t<t_c         \\\\
        \left(\frac{16}{N}\right)^{1/3} \exp \left[ i \left( \frac{\pi N}{2} - 4Jt \right) \right] \operatorname{Ai}\left[ \frac{4J(t_c - t) }{(2Jt_c)^{1/3}}   \right]
         & t \approx t_c \\\\
        \begin{aligned}
             & \sqrt{\frac{2 \cos^4 (k_1/2)}{\pi Jt |\cos k_1|}} \exp \left[ i\Phi(k_1) -i \frac{\pi }{4}\right]           \\
             & \qquad + \sqrt{\frac{2 \cos^4 (k_1/2)}{\pi Jt |\cos k_2|}} \exp \left[i\Phi(k_2) +i \frac{\pi }{4}  \right]
        \end{aligned}
         & t_c <t
    \end{cases}
\end{equation*}
The amplitude can also be expressed in terms of Bessel function, although it yields worse result
\begin{equation*}
    g_N(t)
    = e^{i(N\pi/2 -4Jt)}  \left[ J_{N}(4Jt) -i  J_N' (4Jt)\right]
\end{equation*}

\subsubsection{Derivation.}
With the $\mathbf{j}$ as the receiver site, we compute the transition amplitude from sender site to the receiver site
\begin{align*}
    f_r (t)
     & = \sum_{m=1 } ^N  \bra{\mathbf{r}} e ^{-iHt/\hbar } \ket{\tilde{m }} \braket{\tilde{m }  | \mathbf{s}} 
     = \sum_{m=1 } ^N  e ^{-iE_mt/\hbar }  \braket{\mathbf{r }  | \tilde{m}} \braket{\tilde{m }  | \mathbf{s}} \\
     & = \sum_{m=1 } ^N  a_m e ^{-iE_m t/\hbar } \braket{\mathbf{r}|\tilde{m }} \cos \left[ \frac{\pi }{2N }(m-1 ) (2s-1) \right]                                                                                        \\
     & = \sum_{m=1 } ^N  a_m^2  e ^{-iE_m t/\hbar }\cos \left[ \frac{\pi }{2N }(m-1 ) (2r-1) \right] \cos \left[ \frac{\pi }{2N }(m-1 ) (2s-1) \right]
\end{align*}
Due to the phase factor $e ^{-i E_m t/\hbar}$, this generally complex number.
However, note that we can write
\begin{equation*}
    \exp \left( -i E_m t/\hbar \right) =   \exp(-2iBt/\hbar ) \exp \left( -4Ji (1 - \cos \pi(m-1) t/N \hbar ) \right)
\end{equation*}
so that
\begin{equation*}
    f_r(t) = e ^{-2iBt/\hbar } g_N  = e ^{-2iBt/\hbar } | g_N | e ^{i\gamma(t)} = |g_N| e ^{i(\gamma(t)-2Bt)}
\end{equation*}
with $\hbar = 1$.
To maximize the fidelity, we impose that the phase is a multiple of $2\pi$
At time $t=t_0$ which is the time at which Bob observes the $\mathbf{r}$-th site, we have
\begin{equation*}
    B = \frac{\gamma_0- 2\pi k }{2 t_0}
\end{equation*}
In this condition, the fidelity is
\begin{equation*}
    F = \frac{1 }{2 }+ \frac{1 }{3 }|g_N| + \frac{1 }{6 }|g_N|^2
\end{equation*}
with
\begin{multline*}
    g_N = \sum_{m=1 } ^N  a_m^2 \exp \left[ -i4Jt (1 - \cos \pi(m-1) /N ) \right]  \times\\
    \cos \left[ \frac{\pi }{2N }(m-1 ) (2r-1) \right] \cos \left[ \frac{\pi }{2N }(m-1 ) (2s-1) \right]
\end{multline*}

\begin{figure}[t]
    \centering
    \resizebox{\textwidth}{!}{%% Creator: Matplotlib, PGF backend
%%
%% To include the figure in your LaTeX document, write
%%   \input{<filename>.pgf}
%%
%% Make sure the required packages are loaded in your preamble
%%   \usepackage{pgf}
%%
%% Also ensure that all the required font packages are loaded; for instance,
%% the lmodern package is sometimes necessary when using math font.
%%   \usepackage{lmodern}
%%
%% Figures using additional raster images can only be included by \input if
%% they are in the same directory as the main LaTeX file. For loading figures
%% from other directories you can use the `import` package
%%   \usepackage{import}
%%
%% and then include the figures with
%%   \import{<path to file>}{<filename>.pgf}
%%
%% Matplotlib used the following preamble
%%   
%%           \usepackage{amsmath}
%%           \usepackage{amssymb}
%%       
%%   \makeatletter\@ifpackageloaded{underscore}{}{\usepackage[strings]{underscore}}\makeatother
%%
\begingroup%
\makeatletter%
\begin{pgfpicture}%
\pgfpathrectangle{\pgfpointorigin}{\pgfqpoint{10.000000in}{8.000000in}}%
\pgfusepath{use as bounding box, clip}%
\begin{pgfscope}%
\pgfsetbuttcap%
\pgfsetmiterjoin%
\definecolor{currentfill}{rgb}{1.000000,1.000000,1.000000}%
\pgfsetfillcolor{currentfill}%
\pgfsetlinewidth{0.000000pt}%
\definecolor{currentstroke}{rgb}{1.000000,1.000000,1.000000}%
\pgfsetstrokecolor{currentstroke}%
\pgfsetdash{}{0pt}%
\pgfpathmoveto{\pgfqpoint{0.000000in}{0.000000in}}%
\pgfpathlineto{\pgfqpoint{10.000000in}{0.000000in}}%
\pgfpathlineto{\pgfqpoint{10.000000in}{8.000000in}}%
\pgfpathlineto{\pgfqpoint{0.000000in}{8.000000in}}%
\pgfpathlineto{\pgfqpoint{0.000000in}{0.000000in}}%
\pgfpathclose%
\pgfusepath{fill}%
\end{pgfscope}%
\begin{pgfscope}%
\pgfsetbuttcap%
\pgfsetmiterjoin%
\definecolor{currentfill}{rgb}{1.000000,1.000000,1.000000}%
\pgfsetfillcolor{currentfill}%
\pgfsetlinewidth{0.000000pt}%
\definecolor{currentstroke}{rgb}{0.000000,0.000000,0.000000}%
\pgfsetstrokecolor{currentstroke}%
\pgfsetstrokeopacity{0.000000}%
\pgfsetdash{}{0pt}%
\pgfpathmoveto{\pgfqpoint{0.294194in}{5.519538in}}%
\pgfpathlineto{\pgfqpoint{8.850354in}{5.519538in}}%
\pgfpathlineto{\pgfqpoint{8.850354in}{7.760000in}}%
\pgfpathlineto{\pgfqpoint{0.294194in}{7.760000in}}%
\pgfpathlineto{\pgfqpoint{0.294194in}{5.519538in}}%
\pgfpathclose%
\pgfusepath{fill}%
\end{pgfscope}%
\begin{pgfscope}%
\pgfpathrectangle{\pgfqpoint{0.294194in}{5.519538in}}{\pgfqpoint{8.556160in}{2.240462in}}%
\pgfusepath{clip}%
\pgfsetrectcap%
\pgfsetroundjoin%
\pgfsetlinewidth{1.505625pt}%
\definecolor{currentstroke}{rgb}{0.121569,0.466667,0.705882}%
\pgfsetstrokecolor{currentstroke}%
\pgfsetdash{}{0pt}%
\pgfpathmoveto{\pgfqpoint{0.294194in}{7.760000in}}%
\pgfpathlineto{\pgfqpoint{0.294537in}{7.740159in}}%
\pgfpathlineto{\pgfqpoint{0.295478in}{7.491691in}}%
\pgfpathlineto{\pgfqpoint{0.299927in}{5.519635in}}%
\pgfpathlineto{\pgfqpoint{0.300868in}{5.673513in}}%
\pgfpathlineto{\pgfqpoint{0.302751in}{6.633188in}}%
\pgfpathlineto{\pgfqpoint{0.305660in}{7.759752in}}%
\pgfpathlineto{\pgfqpoint{0.306088in}{7.723584in}}%
\pgfpathlineto{\pgfqpoint{0.307200in}{7.365508in}}%
\pgfpathlineto{\pgfqpoint{0.311392in}{5.519853in}}%
\pgfpathlineto{\pgfqpoint{0.311991in}{5.588031in}}%
\pgfpathlineto{\pgfqpoint{0.313360in}{6.133007in}}%
\pgfpathlineto{\pgfqpoint{0.317125in}{7.759746in}}%
\pgfpathlineto{\pgfqpoint{0.317382in}{7.745521in}}%
\pgfpathlineto{\pgfqpoint{0.318237in}{7.545710in}}%
\pgfpathlineto{\pgfqpoint{0.320548in}{6.286370in}}%
\pgfpathlineto{\pgfqpoint{0.322858in}{5.519653in}}%
\pgfpathlineto{\pgfqpoint{0.322943in}{5.521516in}}%
\pgfpathlineto{\pgfqpoint{0.323542in}{5.602002in}}%
\pgfpathlineto{\pgfqpoint{0.324997in}{6.213121in}}%
\pgfpathlineto{\pgfqpoint{0.328590in}{7.759954in}}%
\pgfpathlineto{\pgfqpoint{0.328762in}{7.755214in}}%
\pgfpathlineto{\pgfqpoint{0.329446in}{7.640932in}}%
\pgfpathlineto{\pgfqpoint{0.331157in}{6.832753in}}%
\pgfpathlineto{\pgfqpoint{0.334323in}{5.519825in}}%
\pgfpathlineto{\pgfqpoint{0.334409in}{5.519992in}}%
\pgfpathlineto{\pgfqpoint{0.334837in}{5.557052in}}%
\pgfpathlineto{\pgfqpoint{0.335949in}{5.912773in}}%
\pgfpathlineto{\pgfqpoint{0.340141in}{7.759917in}}%
\pgfpathlineto{\pgfqpoint{0.340740in}{7.700322in}}%
\pgfpathlineto{\pgfqpoint{0.342024in}{7.221490in}}%
\pgfpathlineto{\pgfqpoint{0.345960in}{5.519845in}}%
\pgfpathlineto{\pgfqpoint{0.346302in}{5.543044in}}%
\pgfpathlineto{\pgfqpoint{0.347329in}{5.829044in}}%
\pgfpathlineto{\pgfqpoint{0.351692in}{7.759736in}}%
\pgfpathlineto{\pgfqpoint{0.352548in}{7.649598in}}%
\pgfpathlineto{\pgfqpoint{0.354174in}{6.909864in}}%
\pgfpathlineto{\pgfqpoint{0.357511in}{5.519743in}}%
\pgfpathlineto{\pgfqpoint{0.357596in}{5.520435in}}%
\pgfpathlineto{\pgfqpoint{0.358110in}{5.574439in}}%
\pgfpathlineto{\pgfqpoint{0.359393in}{6.041593in}}%
\pgfpathlineto{\pgfqpoint{0.363329in}{7.759772in}}%
\pgfpathlineto{\pgfqpoint{0.363757in}{7.732430in}}%
\pgfpathlineto{\pgfqpoint{0.364783in}{7.439598in}}%
\pgfpathlineto{\pgfqpoint{0.369147in}{5.519863in}}%
\pgfpathlineto{\pgfqpoint{0.370003in}{5.629626in}}%
\pgfpathlineto{\pgfqpoint{0.371628in}{6.363828in}}%
\pgfpathlineto{\pgfqpoint{0.375051in}{7.759496in}}%
\pgfpathlineto{\pgfqpoint{0.375479in}{7.724534in}}%
\pgfpathlineto{\pgfqpoint{0.376591in}{7.379499in}}%
\pgfpathlineto{\pgfqpoint{0.380869in}{5.519838in}}%
\pgfpathlineto{\pgfqpoint{0.381554in}{5.596760in}}%
\pgfpathlineto{\pgfqpoint{0.383008in}{6.185568in}}%
\pgfpathlineto{\pgfqpoint{0.386687in}{7.759534in}}%
\pgfpathlineto{\pgfqpoint{0.386858in}{7.756469in}}%
\pgfpathlineto{\pgfqpoint{0.387543in}{7.651857in}}%
\pgfpathlineto{\pgfqpoint{0.389169in}{6.925288in}}%
\pgfpathlineto{\pgfqpoint{0.392591in}{5.519908in}}%
\pgfpathlineto{\pgfqpoint{0.392677in}{5.521228in}}%
\pgfpathlineto{\pgfqpoint{0.393190in}{5.577878in}}%
\pgfpathlineto{\pgfqpoint{0.394474in}{6.043546in}}%
\pgfpathlineto{\pgfqpoint{0.398495in}{7.759445in}}%
\pgfpathlineto{\pgfqpoint{0.398837in}{7.737535in}}%
\pgfpathlineto{\pgfqpoint{0.399864in}{7.461551in}}%
\pgfpathlineto{\pgfqpoint{0.404313in}{5.520266in}}%
\pgfpathlineto{\pgfqpoint{0.405254in}{5.646325in}}%
\pgfpathlineto{\pgfqpoint{0.406966in}{6.439624in}}%
\pgfpathlineto{\pgfqpoint{0.410217in}{7.759269in}}%
\pgfpathlineto{\pgfqpoint{0.410303in}{7.759125in}}%
\pgfpathlineto{\pgfqpoint{0.410730in}{7.723797in}}%
\pgfpathlineto{\pgfqpoint{0.411843in}{7.383196in}}%
\pgfpathlineto{\pgfqpoint{0.416206in}{5.520368in}}%
\pgfpathlineto{\pgfqpoint{0.416805in}{5.584170in}}%
\pgfpathlineto{\pgfqpoint{0.418174in}{6.098871in}}%
\pgfpathlineto{\pgfqpoint{0.422110in}{7.759294in}}%
\pgfpathlineto{\pgfqpoint{0.422367in}{7.747036in}}%
\pgfpathlineto{\pgfqpoint{0.423223in}{7.562502in}}%
\pgfpathlineto{\pgfqpoint{0.425447in}{6.406717in}}%
\pgfpathlineto{\pgfqpoint{0.428014in}{5.520217in}}%
\pgfpathlineto{\pgfqpoint{0.428271in}{5.530167in}}%
\pgfpathlineto{\pgfqpoint{0.429126in}{5.707046in}}%
\pgfpathlineto{\pgfqpoint{0.431265in}{6.803000in}}%
\pgfpathlineto{\pgfqpoint{0.434003in}{7.758986in}}%
\pgfpathlineto{\pgfqpoint{0.434431in}{7.724966in}}%
\pgfpathlineto{\pgfqpoint{0.435544in}{7.390771in}}%
\pgfpathlineto{\pgfqpoint{0.439907in}{5.520349in}}%
\pgfpathlineto{\pgfqpoint{0.440592in}{5.590825in}}%
\pgfpathlineto{\pgfqpoint{0.441961in}{6.110355in}}%
\pgfpathlineto{\pgfqpoint{0.445897in}{7.759086in}}%
\pgfpathlineto{\pgfqpoint{0.446153in}{7.747531in}}%
\pgfpathlineto{\pgfqpoint{0.447009in}{7.567236in}}%
\pgfpathlineto{\pgfqpoint{0.449148in}{6.474354in}}%
\pgfpathlineto{\pgfqpoint{0.451886in}{5.520625in}}%
\pgfpathlineto{\pgfqpoint{0.452314in}{5.552830in}}%
\pgfpathlineto{\pgfqpoint{0.453426in}{5.880344in}}%
\pgfpathlineto{\pgfqpoint{0.457875in}{7.758811in}}%
\pgfpathlineto{\pgfqpoint{0.458560in}{7.680526in}}%
\pgfpathlineto{\pgfqpoint{0.460014in}{7.108217in}}%
\pgfpathlineto{\pgfqpoint{0.463865in}{5.520727in}}%
\pgfpathlineto{\pgfqpoint{0.464036in}{5.526566in}}%
\pgfpathlineto{\pgfqpoint{0.464806in}{5.660874in}}%
\pgfpathlineto{\pgfqpoint{0.466603in}{6.500233in}}%
\pgfpathlineto{\pgfqpoint{0.469854in}{7.758829in}}%
\pgfpathlineto{\pgfqpoint{0.470196in}{7.740306in}}%
\pgfpathlineto{\pgfqpoint{0.471138in}{7.514310in}}%
\pgfpathlineto{\pgfqpoint{0.473790in}{6.115760in}}%
\pgfpathlineto{\pgfqpoint{0.475843in}{5.520870in}}%
\pgfpathlineto{\pgfqpoint{0.476015in}{5.524083in}}%
\pgfpathlineto{\pgfqpoint{0.476699in}{5.624358in}}%
\pgfpathlineto{\pgfqpoint{0.478325in}{6.319820in}}%
\pgfpathlineto{\pgfqpoint{0.481918in}{7.758609in}}%
\pgfpathlineto{\pgfqpoint{0.482261in}{7.738838in}}%
\pgfpathlineto{\pgfqpoint{0.483287in}{7.478934in}}%
\pgfpathlineto{\pgfqpoint{0.487908in}{5.521118in}}%
\pgfpathlineto{\pgfqpoint{0.488935in}{5.666724in}}%
\pgfpathlineto{\pgfqpoint{0.490817in}{6.553214in}}%
\pgfpathlineto{\pgfqpoint{0.493983in}{7.758473in}}%
\pgfpathlineto{\pgfqpoint{0.494239in}{7.749478in}}%
\pgfpathlineto{\pgfqpoint{0.495095in}{7.581287in}}%
\pgfpathlineto{\pgfqpoint{0.497149in}{6.568556in}}%
\pgfpathlineto{\pgfqpoint{0.500058in}{5.521138in}}%
\pgfpathlineto{\pgfqpoint{0.500400in}{5.538447in}}%
\pgfpathlineto{\pgfqpoint{0.501341in}{5.758171in}}%
\pgfpathlineto{\pgfqpoint{0.503908in}{7.095383in}}%
\pgfpathlineto{\pgfqpoint{0.506133in}{7.758303in}}%
\pgfpathlineto{\pgfqpoint{0.506218in}{7.757301in}}%
\pgfpathlineto{\pgfqpoint{0.506731in}{7.705836in}}%
\pgfpathlineto{\pgfqpoint{0.508015in}{7.271943in}}%
\pgfpathlineto{\pgfqpoint{0.512207in}{5.521374in}}%
\pgfpathlineto{\pgfqpoint{0.512635in}{5.546346in}}%
\pgfpathlineto{\pgfqpoint{0.513662in}{5.814631in}}%
\pgfpathlineto{\pgfqpoint{0.518282in}{7.757848in}}%
\pgfpathlineto{\pgfqpoint{0.519224in}{7.640462in}}%
\pgfpathlineto{\pgfqpoint{0.520935in}{6.892536in}}%
\pgfpathlineto{\pgfqpoint{0.524443in}{5.521546in}}%
\pgfpathlineto{\pgfqpoint{0.524785in}{5.539517in}}%
\pgfpathlineto{\pgfqpoint{0.525726in}{5.758288in}}%
\pgfpathlineto{\pgfqpoint{0.528293in}{7.084587in}}%
\pgfpathlineto{\pgfqpoint{0.530518in}{7.757647in}}%
\pgfpathlineto{\pgfqpoint{0.530689in}{7.755401in}}%
\pgfpathlineto{\pgfqpoint{0.531288in}{7.680715in}}%
\pgfpathlineto{\pgfqpoint{0.532742in}{7.127996in}}%
\pgfpathlineto{\pgfqpoint{0.536678in}{5.521792in}}%
\pgfpathlineto{\pgfqpoint{0.536849in}{5.525303in}}%
\pgfpathlineto{\pgfqpoint{0.537534in}{5.623433in}}%
\pgfpathlineto{\pgfqpoint{0.539160in}{6.298929in}}%
\pgfpathlineto{\pgfqpoint{0.542839in}{7.757667in}}%
\pgfpathlineto{\pgfqpoint{0.542924in}{7.756645in}}%
\pgfpathlineto{\pgfqpoint{0.543438in}{7.706121in}}%
\pgfpathlineto{\pgfqpoint{0.544721in}{7.280997in}}%
\pgfpathlineto{\pgfqpoint{0.548999in}{5.521988in}}%
\pgfpathlineto{\pgfqpoint{0.549427in}{5.548422in}}%
\pgfpathlineto{\pgfqpoint{0.550539in}{5.848717in}}%
\pgfpathlineto{\pgfqpoint{0.555160in}{7.757414in}}%
\pgfpathlineto{\pgfqpoint{0.556015in}{7.655798in}}%
\pgfpathlineto{\pgfqpoint{0.557641in}{6.985679in}}%
\pgfpathlineto{\pgfqpoint{0.561320in}{5.522387in}}%
\pgfpathlineto{\pgfqpoint{0.561406in}{5.522629in}}%
\pgfpathlineto{\pgfqpoint{0.561834in}{5.555288in}}%
\pgfpathlineto{\pgfqpoint{0.562946in}{5.867819in}}%
\pgfpathlineto{\pgfqpoint{0.567566in}{7.757135in}}%
\pgfpathlineto{\pgfqpoint{0.568336in}{7.669330in}}%
\pgfpathlineto{\pgfqpoint{0.569877in}{7.066430in}}%
\pgfpathlineto{\pgfqpoint{0.573727in}{5.522605in}}%
\pgfpathlineto{\pgfqpoint{0.573898in}{5.525376in}}%
\pgfpathlineto{\pgfqpoint{0.574583in}{5.618672in}}%
\pgfpathlineto{\pgfqpoint{0.576208in}{6.276349in}}%
\pgfpathlineto{\pgfqpoint{0.579973in}{7.756918in}}%
\pgfpathlineto{\pgfqpoint{0.580059in}{7.755935in}}%
\pgfpathlineto{\pgfqpoint{0.580572in}{7.706692in}}%
\pgfpathlineto{\pgfqpoint{0.581855in}{7.291074in}}%
\pgfpathlineto{\pgfqpoint{0.586219in}{5.522774in}}%
\pgfpathlineto{\pgfqpoint{0.586647in}{5.550005in}}%
\pgfpathlineto{\pgfqpoint{0.587759in}{5.846709in}}%
\pgfpathlineto{\pgfqpoint{0.592465in}{7.756619in}}%
\pgfpathlineto{\pgfqpoint{0.593321in}{7.651823in}}%
\pgfpathlineto{\pgfqpoint{0.594946in}{6.988416in}}%
\pgfpathlineto{\pgfqpoint{0.598711in}{5.523036in}}%
\pgfpathlineto{\pgfqpoint{0.599053in}{5.539696in}}%
\pgfpathlineto{\pgfqpoint{0.599995in}{5.747333in}}%
\pgfpathlineto{\pgfqpoint{0.602476in}{6.988591in}}%
\pgfpathlineto{\pgfqpoint{0.604957in}{7.756296in}}%
\pgfpathlineto{\pgfqpoint{0.605214in}{7.748584in}}%
\pgfpathlineto{\pgfqpoint{0.605984in}{7.617806in}}%
\pgfpathlineto{\pgfqpoint{0.607781in}{6.826960in}}%
\pgfpathlineto{\pgfqpoint{0.611203in}{5.523780in}}%
\pgfpathlineto{\pgfqpoint{0.611289in}{5.523444in}}%
\pgfpathlineto{\pgfqpoint{0.611374in}{5.525147in}}%
\pgfpathlineto{\pgfqpoint{0.611973in}{5.593429in}}%
\pgfpathlineto{\pgfqpoint{0.613428in}{6.116088in}}%
\pgfpathlineto{\pgfqpoint{0.617535in}{7.756001in}}%
\pgfpathlineto{\pgfqpoint{0.617792in}{7.748300in}}%
\pgfpathlineto{\pgfqpoint{0.618562in}{7.618435in}}%
\pgfpathlineto{\pgfqpoint{0.620358in}{6.832732in}}%
\pgfpathlineto{\pgfqpoint{0.623781in}{5.524557in}}%
\pgfpathlineto{\pgfqpoint{0.623866in}{5.523639in}}%
\pgfpathlineto{\pgfqpoint{0.624038in}{5.527867in}}%
\pgfpathlineto{\pgfqpoint{0.624722in}{5.624132in}}%
\pgfpathlineto{\pgfqpoint{0.626348in}{6.273054in}}%
\pgfpathlineto{\pgfqpoint{0.630198in}{7.755712in}}%
\pgfpathlineto{\pgfqpoint{0.630540in}{7.738228in}}%
\pgfpathlineto{\pgfqpoint{0.631482in}{7.531805in}}%
\pgfpathlineto{\pgfqpoint{0.633963in}{6.304851in}}%
\pgfpathlineto{\pgfqpoint{0.636530in}{5.523981in}}%
\pgfpathlineto{\pgfqpoint{0.636872in}{5.541222in}}%
\pgfpathlineto{\pgfqpoint{0.637813in}{5.746379in}}%
\pgfpathlineto{\pgfqpoint{0.640209in}{6.924173in}}%
\pgfpathlineto{\pgfqpoint{0.642861in}{7.755416in}}%
\pgfpathlineto{\pgfqpoint{0.643204in}{7.739625in}}%
\pgfpathlineto{\pgfqpoint{0.644145in}{7.538775in}}%
\pgfpathlineto{\pgfqpoint{0.646541in}{6.367382in}}%
\pgfpathlineto{\pgfqpoint{0.649193in}{5.524409in}}%
\pgfpathlineto{\pgfqpoint{0.649364in}{5.527013in}}%
\pgfpathlineto{\pgfqpoint{0.650049in}{5.615684in}}%
\pgfpathlineto{\pgfqpoint{0.651589in}{6.201721in}}%
\pgfpathlineto{\pgfqpoint{0.655610in}{7.755046in}}%
\pgfpathlineto{\pgfqpoint{0.655696in}{7.753710in}}%
\pgfpathlineto{\pgfqpoint{0.656295in}{7.689578in}}%
\pgfpathlineto{\pgfqpoint{0.657664in}{7.225120in}}%
\pgfpathlineto{\pgfqpoint{0.661942in}{5.524847in}}%
\pgfpathlineto{\pgfqpoint{0.662284in}{5.536964in}}%
\pgfpathlineto{\pgfqpoint{0.663225in}{5.726547in}}%
\pgfpathlineto{\pgfqpoint{0.665535in}{6.833718in}}%
\pgfpathlineto{\pgfqpoint{0.668359in}{7.754660in}}%
\pgfpathlineto{\pgfqpoint{0.668616in}{7.747315in}}%
\pgfpathlineto{\pgfqpoint{0.669386in}{7.622005in}}%
\pgfpathlineto{\pgfqpoint{0.671183in}{6.858323in}}%
\pgfpathlineto{\pgfqpoint{0.674691in}{5.526404in}}%
\pgfpathlineto{\pgfqpoint{0.674776in}{5.525037in}}%
\pgfpathlineto{\pgfqpoint{0.675033in}{5.532655in}}%
\pgfpathlineto{\pgfqpoint{0.675803in}{5.658304in}}%
\pgfpathlineto{\pgfqpoint{0.677600in}{6.420493in}}%
\pgfpathlineto{\pgfqpoint{0.681193in}{7.754264in}}%
\pgfpathlineto{\pgfqpoint{0.681364in}{7.751558in}}%
\pgfpathlineto{\pgfqpoint{0.682049in}{7.664158in}}%
\pgfpathlineto{\pgfqpoint{0.683589in}{7.088155in}}%
\pgfpathlineto{\pgfqpoint{0.687696in}{5.525553in}}%
\pgfpathlineto{\pgfqpoint{0.687782in}{5.527357in}}%
\pgfpathlineto{\pgfqpoint{0.688381in}{5.593518in}}%
\pgfpathlineto{\pgfqpoint{0.689835in}{6.094982in}}%
\pgfpathlineto{\pgfqpoint{0.694113in}{7.753972in}}%
\pgfpathlineto{\pgfqpoint{0.694370in}{7.745258in}}%
\pgfpathlineto{\pgfqpoint{0.695226in}{7.594781in}}%
\pgfpathlineto{\pgfqpoint{0.697194in}{6.721447in}}%
\pgfpathlineto{\pgfqpoint{0.700445in}{5.528147in}}%
\pgfpathlineto{\pgfqpoint{0.700616in}{5.525934in}}%
\pgfpathlineto{\pgfqpoint{0.700873in}{5.536998in}}%
\pgfpathlineto{\pgfqpoint{0.701728in}{5.694293in}}%
\pgfpathlineto{\pgfqpoint{0.703782in}{6.619898in}}%
\pgfpathlineto{\pgfqpoint{0.706948in}{7.751907in}}%
\pgfpathlineto{\pgfqpoint{0.707033in}{7.753477in}}%
\pgfpathlineto{\pgfqpoint{0.707375in}{7.740670in}}%
\pgfpathlineto{\pgfqpoint{0.708317in}{7.553971in}}%
\pgfpathlineto{\pgfqpoint{0.710541in}{6.513607in}}%
\pgfpathlineto{\pgfqpoint{0.713450in}{5.527604in}}%
\pgfpathlineto{\pgfqpoint{0.713536in}{5.526215in}}%
\pgfpathlineto{\pgfqpoint{0.713793in}{5.533459in}}%
\pgfpathlineto{\pgfqpoint{0.714563in}{5.655382in}}%
\pgfpathlineto{\pgfqpoint{0.716360in}{6.400368in}}%
\pgfpathlineto{\pgfqpoint{0.719953in}{7.751562in}}%
\pgfpathlineto{\pgfqpoint{0.720039in}{7.753080in}}%
\pgfpathlineto{\pgfqpoint{0.720381in}{7.740232in}}%
\pgfpathlineto{\pgfqpoint{0.721322in}{7.554759in}}%
\pgfpathlineto{\pgfqpoint{0.723547in}{6.520731in}}%
\pgfpathlineto{\pgfqpoint{0.726456in}{5.528811in}}%
\pgfpathlineto{\pgfqpoint{0.726627in}{5.526784in}}%
\pgfpathlineto{\pgfqpoint{0.726884in}{5.537870in}}%
\pgfpathlineto{\pgfqpoint{0.727739in}{5.693123in}}%
\pgfpathlineto{\pgfqpoint{0.729793in}{6.606872in}}%
\pgfpathlineto{\pgfqpoint{0.733044in}{7.751914in}}%
\pgfpathlineto{\pgfqpoint{0.733130in}{7.752738in}}%
\pgfpathlineto{\pgfqpoint{0.733301in}{7.748756in}}%
\pgfpathlineto{\pgfqpoint{0.733985in}{7.659076in}}%
\pgfpathlineto{\pgfqpoint{0.735611in}{7.050153in}}%
\pgfpathlineto{\pgfqpoint{0.739718in}{5.527242in}}%
\pgfpathlineto{\pgfqpoint{0.740146in}{5.555092in}}%
\pgfpathlineto{\pgfqpoint{0.741258in}{5.831472in}}%
\pgfpathlineto{\pgfqpoint{0.746221in}{7.752248in}}%
\pgfpathlineto{\pgfqpoint{0.747162in}{7.646396in}}%
\pgfpathlineto{\pgfqpoint{0.748873in}{6.982894in}}%
\pgfpathlineto{\pgfqpoint{0.752809in}{5.527481in}}%
\pgfpathlineto{\pgfqpoint{0.753066in}{5.534993in}}%
\pgfpathlineto{\pgfqpoint{0.753921in}{5.676929in}}%
\pgfpathlineto{\pgfqpoint{0.755889in}{6.519252in}}%
\pgfpathlineto{\pgfqpoint{0.759312in}{7.750476in}}%
\pgfpathlineto{\pgfqpoint{0.759397in}{7.751804in}}%
\pgfpathlineto{\pgfqpoint{0.759654in}{7.744732in}}%
\pgfpathlineto{\pgfqpoint{0.760424in}{7.626410in}}%
\pgfpathlineto{\pgfqpoint{0.762221in}{6.900656in}}%
\pgfpathlineto{\pgfqpoint{0.765986in}{5.528120in}}%
\pgfpathlineto{\pgfqpoint{0.766071in}{5.528186in}}%
\pgfpathlineto{\pgfqpoint{0.766499in}{5.555916in}}%
\pgfpathlineto{\pgfqpoint{0.767611in}{5.828378in}}%
\pgfpathlineto{\pgfqpoint{0.772660in}{7.751366in}}%
\pgfpathlineto{\pgfqpoint{0.773601in}{7.640394in}}%
\pgfpathlineto{\pgfqpoint{0.775312in}{6.978824in}}%
\pgfpathlineto{\pgfqpoint{0.779248in}{5.528599in}}%
\pgfpathlineto{\pgfqpoint{0.779333in}{5.528655in}}%
\pgfpathlineto{\pgfqpoint{0.779761in}{5.556085in}}%
\pgfpathlineto{\pgfqpoint{0.780874in}{5.826036in}}%
\pgfpathlineto{\pgfqpoint{0.785922in}{7.750846in}}%
\pgfpathlineto{\pgfqpoint{0.786863in}{7.647665in}}%
\pgfpathlineto{\pgfqpoint{0.788574in}{6.999570in}}%
\pgfpathlineto{\pgfqpoint{0.792596in}{5.528917in}}%
\pgfpathlineto{\pgfqpoint{0.792852in}{5.536096in}}%
\pgfpathlineto{\pgfqpoint{0.793622in}{5.652507in}}%
\pgfpathlineto{\pgfqpoint{0.795419in}{6.365362in}}%
\pgfpathlineto{\pgfqpoint{0.799184in}{7.748919in}}%
\pgfpathlineto{\pgfqpoint{0.799270in}{7.750323in}}%
\pgfpathlineto{\pgfqpoint{0.799612in}{7.738040in}}%
\pgfpathlineto{\pgfqpoint{0.800553in}{7.561972in}}%
\pgfpathlineto{\pgfqpoint{0.802778in}{6.568903in}}%
\pgfpathlineto{\pgfqpoint{0.805858in}{5.531609in}}%
\pgfpathlineto{\pgfqpoint{0.806029in}{5.529560in}}%
\pgfpathlineto{\pgfqpoint{0.806286in}{5.539847in}}%
\pgfpathlineto{\pgfqpoint{0.807141in}{5.686258in}}%
\pgfpathlineto{\pgfqpoint{0.809109in}{6.516067in}}%
\pgfpathlineto{\pgfqpoint{0.812532in}{7.746174in}}%
\pgfpathlineto{\pgfqpoint{0.812703in}{7.749878in}}%
\pgfpathlineto{\pgfqpoint{0.813045in}{7.736017in}}%
\pgfpathlineto{\pgfqpoint{0.813986in}{7.557258in}}%
\pgfpathlineto{\pgfqpoint{0.816211in}{6.566561in}}%
\pgfpathlineto{\pgfqpoint{0.819291in}{5.532452in}}%
\pgfpathlineto{\pgfqpoint{0.819462in}{5.530005in}}%
\pgfpathlineto{\pgfqpoint{0.819805in}{5.546269in}}%
\pgfpathlineto{\pgfqpoint{0.820831in}{5.756560in}}%
\pgfpathlineto{\pgfqpoint{0.823313in}{6.898850in}}%
\pgfpathlineto{\pgfqpoint{0.826051in}{7.747312in}}%
\pgfpathlineto{\pgfqpoint{0.826222in}{7.749148in}}%
\pgfpathlineto{\pgfqpoint{0.826478in}{7.738733in}}%
\pgfpathlineto{\pgfqpoint{0.827334in}{7.593510in}}%
\pgfpathlineto{\pgfqpoint{0.829302in}{6.772102in}}%
\pgfpathlineto{\pgfqpoint{0.832810in}{5.532517in}}%
\pgfpathlineto{\pgfqpoint{0.832981in}{5.530652in}}%
\pgfpathlineto{\pgfqpoint{0.833238in}{5.540957in}}%
\pgfpathlineto{\pgfqpoint{0.834093in}{5.685309in}}%
\pgfpathlineto{\pgfqpoint{0.836061in}{6.503290in}}%
\pgfpathlineto{\pgfqpoint{0.839569in}{7.746204in}}%
\pgfpathlineto{\pgfqpoint{0.839741in}{7.748737in}}%
\pgfpathlineto{\pgfqpoint{0.840083in}{7.732950in}}%
\pgfpathlineto{\pgfqpoint{0.841024in}{7.551949in}}%
\pgfpathlineto{\pgfqpoint{0.843334in}{6.525679in}}%
\pgfpathlineto{\pgfqpoint{0.846329in}{5.534856in}}%
\pgfpathlineto{\pgfqpoint{0.846500in}{5.531026in}}%
\pgfpathlineto{\pgfqpoint{0.846842in}{5.544126in}}%
\pgfpathlineto{\pgfqpoint{0.847783in}{5.717627in}}%
\pgfpathlineto{\pgfqpoint{0.850008in}{6.689367in}}%
\pgfpathlineto{\pgfqpoint{0.853174in}{7.745836in}}%
\pgfpathlineto{\pgfqpoint{0.853345in}{7.748144in}}%
\pgfpathlineto{\pgfqpoint{0.853687in}{7.732113in}}%
\pgfpathlineto{\pgfqpoint{0.854714in}{7.526285in}}%
\pgfpathlineto{\pgfqpoint{0.857195in}{6.402534in}}%
\pgfpathlineto{\pgfqpoint{0.860019in}{5.533317in}}%
\pgfpathlineto{\pgfqpoint{0.860104in}{5.531739in}}%
\pgfpathlineto{\pgfqpoint{0.860447in}{5.542564in}}%
\pgfpathlineto{\pgfqpoint{0.861388in}{5.708918in}}%
\pgfpathlineto{\pgfqpoint{0.863527in}{6.623770in}}%
\pgfpathlineto{\pgfqpoint{0.866778in}{7.742916in}}%
\pgfpathlineto{\pgfqpoint{0.866949in}{7.747549in}}%
\pgfpathlineto{\pgfqpoint{0.867377in}{7.729318in}}%
\pgfpathlineto{\pgfqpoint{0.868404in}{7.519659in}}%
\pgfpathlineto{\pgfqpoint{0.870971in}{6.356332in}}%
\pgfpathlineto{\pgfqpoint{0.873709in}{5.534048in}}%
\pgfpathlineto{\pgfqpoint{0.873794in}{5.532375in}}%
\pgfpathlineto{\pgfqpoint{0.874137in}{5.542652in}}%
\pgfpathlineto{\pgfqpoint{0.874992in}{5.683658in}}%
\pgfpathlineto{\pgfqpoint{0.876960in}{6.483395in}}%
\pgfpathlineto{\pgfqpoint{0.880554in}{7.744319in}}%
\pgfpathlineto{\pgfqpoint{0.880725in}{7.747029in}}%
\pgfpathlineto{\pgfqpoint{0.881067in}{7.732206in}}%
\pgfpathlineto{\pgfqpoint{0.882008in}{7.557632in}}%
\pgfpathlineto{\pgfqpoint{0.884233in}{6.598076in}}%
\pgfpathlineto{\pgfqpoint{0.887399in}{5.537268in}}%
\pgfpathlineto{\pgfqpoint{0.887570in}{5.532862in}}%
\pgfpathlineto{\pgfqpoint{0.887998in}{5.551214in}}%
\pgfpathlineto{\pgfqpoint{0.889025in}{5.758707in}}%
\pgfpathlineto{\pgfqpoint{0.891506in}{6.869035in}}%
\pgfpathlineto{\pgfqpoint{0.894329in}{7.743063in}}%
\pgfpathlineto{\pgfqpoint{0.894501in}{7.746449in}}%
\pgfpathlineto{\pgfqpoint{0.894843in}{7.733175in}}%
\pgfpathlineto{\pgfqpoint{0.895784in}{7.563881in}}%
\pgfpathlineto{\pgfqpoint{0.898009in}{6.615627in}}%
\pgfpathlineto{\pgfqpoint{0.901260in}{5.536431in}}%
\pgfpathlineto{\pgfqpoint{0.901431in}{5.533409in}}%
\pgfpathlineto{\pgfqpoint{0.901773in}{5.547308in}}%
\pgfpathlineto{\pgfqpoint{0.902715in}{5.717534in}}%
\pgfpathlineto{\pgfqpoint{0.904939in}{6.664141in}}%
\pgfpathlineto{\pgfqpoint{0.908191in}{7.742512in}}%
\pgfpathlineto{\pgfqpoint{0.908362in}{7.745819in}}%
\pgfpathlineto{\pgfqpoint{0.908704in}{7.732591in}}%
\pgfpathlineto{\pgfqpoint{0.909645in}{7.564741in}}%
\pgfpathlineto{\pgfqpoint{0.911784in}{6.666282in}}%
\pgfpathlineto{\pgfqpoint{0.915121in}{5.538345in}}%
\pgfpathlineto{\pgfqpoint{0.915292in}{5.534109in}}%
\pgfpathlineto{\pgfqpoint{0.915720in}{5.552291in}}%
\pgfpathlineto{\pgfqpoint{0.916747in}{5.756173in}}%
\pgfpathlineto{\pgfqpoint{0.919228in}{6.851268in}}%
\pgfpathlineto{\pgfqpoint{0.922137in}{7.742599in}}%
\pgfpathlineto{\pgfqpoint{0.922308in}{7.745128in}}%
\pgfpathlineto{\pgfqpoint{0.922651in}{7.730555in}}%
\pgfpathlineto{\pgfqpoint{0.923592in}{7.560598in}}%
\pgfpathlineto{\pgfqpoint{0.925816in}{6.622421in}}%
\pgfpathlineto{\pgfqpoint{0.929068in}{5.539615in}}%
\pgfpathlineto{\pgfqpoint{0.929239in}{5.534868in}}%
\pgfpathlineto{\pgfqpoint{0.929667in}{5.551481in}}%
\pgfpathlineto{\pgfqpoint{0.930693in}{5.750304in}}%
\pgfpathlineto{\pgfqpoint{0.933089in}{6.792849in}}%
\pgfpathlineto{\pgfqpoint{0.936084in}{7.740059in}}%
\pgfpathlineto{\pgfqpoint{0.936255in}{7.744411in}}%
\pgfpathlineto{\pgfqpoint{0.936683in}{7.726960in}}%
\pgfpathlineto{\pgfqpoint{0.937710in}{7.527109in}}%
\pgfpathlineto{\pgfqpoint{0.940191in}{6.444804in}}%
\pgfpathlineto{\pgfqpoint{0.943100in}{5.540132in}}%
\pgfpathlineto{\pgfqpoint{0.943271in}{5.535521in}}%
\pgfpathlineto{\pgfqpoint{0.943699in}{5.552178in}}%
\pgfpathlineto{\pgfqpoint{0.944726in}{5.749477in}}%
\pgfpathlineto{\pgfqpoint{0.947121in}{6.784679in}}%
\pgfpathlineto{\pgfqpoint{0.950116in}{7.737887in}}%
\pgfpathlineto{\pgfqpoint{0.950373in}{7.743760in}}%
\pgfpathlineto{\pgfqpoint{0.950801in}{7.721534in}}%
\pgfpathlineto{\pgfqpoint{0.951913in}{7.486490in}}%
\pgfpathlineto{\pgfqpoint{0.954737in}{6.227427in}}%
\pgfpathlineto{\pgfqpoint{0.957303in}{5.537244in}}%
\pgfpathlineto{\pgfqpoint{0.957389in}{5.536101in}}%
\pgfpathlineto{\pgfqpoint{0.957646in}{5.542243in}}%
\pgfpathlineto{\pgfqpoint{0.958416in}{5.645029in}}%
\pgfpathlineto{\pgfqpoint{0.960127in}{6.248320in}}%
\pgfpathlineto{\pgfqpoint{0.964405in}{7.742452in}}%
\pgfpathlineto{\pgfqpoint{0.964491in}{7.743117in}}%
\pgfpathlineto{\pgfqpoint{0.964662in}{7.739683in}}%
\pgfpathlineto{\pgfqpoint{0.965346in}{7.663415in}}%
\pgfpathlineto{\pgfqpoint{0.966886in}{7.174434in}}%
\pgfpathlineto{\pgfqpoint{0.971592in}{5.536797in}}%
\pgfpathlineto{\pgfqpoint{0.971763in}{5.540498in}}%
\pgfpathlineto{\pgfqpoint{0.972448in}{5.617472in}}%
\pgfpathlineto{\pgfqpoint{0.973988in}{6.106000in}}%
\pgfpathlineto{\pgfqpoint{0.978694in}{7.742401in}}%
\pgfpathlineto{\pgfqpoint{0.978865in}{7.739091in}}%
\pgfpathlineto{\pgfqpoint{0.979550in}{7.663968in}}%
\pgfpathlineto{\pgfqpoint{0.981090in}{7.180254in}}%
\pgfpathlineto{\pgfqpoint{0.985796in}{5.537561in}}%
\pgfpathlineto{\pgfqpoint{0.985967in}{5.539825in}}%
\pgfpathlineto{\pgfqpoint{0.986651in}{5.610553in}}%
\pgfpathlineto{\pgfqpoint{0.988191in}{6.085060in}}%
\pgfpathlineto{\pgfqpoint{0.992983in}{7.741639in}}%
\pgfpathlineto{\pgfqpoint{0.993154in}{7.737961in}}%
\pgfpathlineto{\pgfqpoint{0.993839in}{7.662070in}}%
\pgfpathlineto{\pgfqpoint{0.995379in}{7.180359in}}%
\pgfpathlineto{\pgfqpoint{1.000170in}{5.538391in}}%
\pgfpathlineto{\pgfqpoint{1.000256in}{5.539828in}}%
\pgfpathlineto{\pgfqpoint{1.000855in}{5.592700in}}%
\pgfpathlineto{\pgfqpoint{1.002224in}{5.967894in}}%
\pgfpathlineto{\pgfqpoint{1.007357in}{7.740746in}}%
\pgfpathlineto{\pgfqpoint{1.007785in}{7.718074in}}%
\pgfpathlineto{\pgfqpoint{1.008898in}{7.489324in}}%
\pgfpathlineto{\pgfqpoint{1.011721in}{6.261839in}}%
\pgfpathlineto{\pgfqpoint{1.014374in}{5.541290in}}%
\pgfpathlineto{\pgfqpoint{1.014545in}{5.539072in}}%
\pgfpathlineto{\pgfqpoint{1.014887in}{5.552952in}}%
\pgfpathlineto{\pgfqpoint{1.015828in}{5.712524in}}%
\pgfpathlineto{\pgfqpoint{1.017967in}{6.561878in}}%
\pgfpathlineto{\pgfqpoint{1.021475in}{7.732697in}}%
\pgfpathlineto{\pgfqpoint{1.021732in}{7.740120in}}%
\pgfpathlineto{\pgfqpoint{1.022160in}{7.722148in}}%
\pgfpathlineto{\pgfqpoint{1.023186in}{7.530860in}}%
\pgfpathlineto{\pgfqpoint{1.025582in}{6.534407in}}%
\pgfpathlineto{\pgfqpoint{1.028748in}{5.545522in}}%
\pgfpathlineto{\pgfqpoint{1.029005in}{5.539846in}}%
\pgfpathlineto{\pgfqpoint{1.029433in}{5.560551in}}%
\pgfpathlineto{\pgfqpoint{1.030545in}{5.781943in}}%
\pgfpathlineto{\pgfqpoint{1.033283in}{6.954916in}}%
\pgfpathlineto{\pgfqpoint{1.036106in}{7.737423in}}%
\pgfpathlineto{\pgfqpoint{1.036278in}{7.739216in}}%
\pgfpathlineto{\pgfqpoint{1.036534in}{7.730642in}}%
\pgfpathlineto{\pgfqpoint{1.037390in}{7.607073in}}%
\pgfpathlineto{\pgfqpoint{1.039272in}{6.924983in}}%
\pgfpathlineto{\pgfqpoint{1.043379in}{5.542672in}}%
\pgfpathlineto{\pgfqpoint{1.043550in}{5.540689in}}%
\pgfpathlineto{\pgfqpoint{1.043893in}{5.554632in}}%
\pgfpathlineto{\pgfqpoint{1.044919in}{5.734062in}}%
\pgfpathlineto{\pgfqpoint{1.047230in}{6.668750in}}%
\pgfpathlineto{\pgfqpoint{1.050566in}{7.732041in}}%
\pgfpathlineto{\pgfqpoint{1.050823in}{7.738522in}}%
\pgfpathlineto{\pgfqpoint{1.051251in}{7.719656in}}%
\pgfpathlineto{\pgfqpoint{1.052278in}{7.529632in}}%
\pgfpathlineto{\pgfqpoint{1.054673in}{6.546572in}}%
\pgfpathlineto{\pgfqpoint{1.057839in}{5.550325in}}%
\pgfpathlineto{\pgfqpoint{1.058096in}{5.541593in}}%
\pgfpathlineto{\pgfqpoint{1.058609in}{5.563941in}}%
\pgfpathlineto{\pgfqpoint{1.059722in}{5.785428in}}%
\pgfpathlineto{\pgfqpoint{1.062460in}{6.944729in}}%
\pgfpathlineto{\pgfqpoint{1.065283in}{7.734083in}}%
\pgfpathlineto{\pgfqpoint{1.065454in}{7.737655in}}%
\pgfpathlineto{\pgfqpoint{1.065882in}{7.720912in}}%
\pgfpathlineto{\pgfqpoint{1.066909in}{7.537175in}}%
\pgfpathlineto{\pgfqpoint{1.069305in}{6.566704in}}%
\pgfpathlineto{\pgfqpoint{1.072556in}{5.549698in}}%
\pgfpathlineto{\pgfqpoint{1.072813in}{5.542299in}}%
\pgfpathlineto{\pgfqpoint{1.073241in}{5.559140in}}%
\pgfpathlineto{\pgfqpoint{1.074267in}{5.742285in}}%
\pgfpathlineto{\pgfqpoint{1.076663in}{6.709071in}}%
\pgfpathlineto{\pgfqpoint{1.079914in}{7.728478in}}%
\pgfpathlineto{\pgfqpoint{1.080171in}{7.736723in}}%
\pgfpathlineto{\pgfqpoint{1.080684in}{7.714080in}}%
\pgfpathlineto{\pgfqpoint{1.081797in}{7.494908in}}%
\pgfpathlineto{\pgfqpoint{1.084535in}{6.347815in}}%
\pgfpathlineto{\pgfqpoint{1.087358in}{5.548919in}}%
\pgfpathlineto{\pgfqpoint{1.087615in}{5.543159in}}%
\pgfpathlineto{\pgfqpoint{1.088043in}{5.562385in}}%
\pgfpathlineto{\pgfqpoint{1.089155in}{5.772967in}}%
\pgfpathlineto{\pgfqpoint{1.091722in}{6.831911in}}%
\pgfpathlineto{\pgfqpoint{1.094802in}{7.731471in}}%
\pgfpathlineto{\pgfqpoint{1.095059in}{7.735781in}}%
\pgfpathlineto{\pgfqpoint{1.095401in}{7.721458in}}%
\pgfpathlineto{\pgfqpoint{1.096428in}{7.546307in}}%
\pgfpathlineto{\pgfqpoint{1.098738in}{6.636929in}}%
\pgfpathlineto{\pgfqpoint{1.102161in}{5.552305in}}%
\pgfpathlineto{\pgfqpoint{1.102417in}{5.544142in}}%
\pgfpathlineto{\pgfqpoint{1.102931in}{5.566272in}}%
\pgfpathlineto{\pgfqpoint{1.104043in}{5.781554in}}%
\pgfpathlineto{\pgfqpoint{1.106695in}{6.876647in}}%
\pgfpathlineto{\pgfqpoint{1.109690in}{7.730445in}}%
\pgfpathlineto{\pgfqpoint{1.109947in}{7.734922in}}%
\pgfpathlineto{\pgfqpoint{1.110289in}{7.721054in}}%
\pgfpathlineto{\pgfqpoint{1.111316in}{7.548689in}}%
\pgfpathlineto{\pgfqpoint{1.113626in}{6.648334in}}%
\pgfpathlineto{\pgfqpoint{1.117048in}{5.555843in}}%
\pgfpathlineto{\pgfqpoint{1.117391in}{5.544932in}}%
\pgfpathlineto{\pgfqpoint{1.117819in}{5.562973in}}%
\pgfpathlineto{\pgfqpoint{1.118845in}{5.743933in}}%
\pgfpathlineto{\pgfqpoint{1.121241in}{6.689918in}}%
\pgfpathlineto{\pgfqpoint{1.124578in}{7.725266in}}%
\pgfpathlineto{\pgfqpoint{1.124920in}{7.734012in}}%
\pgfpathlineto{\pgfqpoint{1.125348in}{7.713464in}}%
\pgfpathlineto{\pgfqpoint{1.126460in}{7.504335in}}%
\pgfpathlineto{\pgfqpoint{1.129113in}{6.424785in}}%
\pgfpathlineto{\pgfqpoint{1.132107in}{5.554043in}}%
\pgfpathlineto{\pgfqpoint{1.132364in}{5.545985in}}%
\pgfpathlineto{\pgfqpoint{1.132878in}{5.567425in}}%
\pgfpathlineto{\pgfqpoint{1.133990in}{5.777512in}}%
\pgfpathlineto{\pgfqpoint{1.136642in}{6.854454in}}%
\pgfpathlineto{\pgfqpoint{1.139637in}{7.724530in}}%
\pgfpathlineto{\pgfqpoint{1.139979in}{7.733064in}}%
\pgfpathlineto{\pgfqpoint{1.140407in}{7.712624in}}%
\pgfpathlineto{\pgfqpoint{1.141519in}{7.505577in}}%
\pgfpathlineto{\pgfqpoint{1.144086in}{6.473050in}}%
\pgfpathlineto{\pgfqpoint{1.147252in}{5.552595in}}%
\pgfpathlineto{\pgfqpoint{1.147509in}{5.546812in}}%
\pgfpathlineto{\pgfqpoint{1.147937in}{5.564671in}}%
\pgfpathlineto{\pgfqpoint{1.148963in}{5.742017in}}%
\pgfpathlineto{\pgfqpoint{1.151273in}{6.632922in}}%
\pgfpathlineto{\pgfqpoint{1.154782in}{7.723835in}}%
\pgfpathlineto{\pgfqpoint{1.155124in}{7.732084in}}%
\pgfpathlineto{\pgfqpoint{1.155552in}{7.711664in}}%
\pgfpathlineto{\pgfqpoint{1.156664in}{7.506485in}}%
\pgfpathlineto{\pgfqpoint{1.159231in}{6.482366in}}%
\pgfpathlineto{\pgfqpoint{1.162397in}{5.555477in}}%
\pgfpathlineto{\pgfqpoint{1.162653in}{5.547867in}}%
\pgfpathlineto{\pgfqpoint{1.163167in}{5.569306in}}%
\pgfpathlineto{\pgfqpoint{1.164279in}{5.775547in}}%
\pgfpathlineto{\pgfqpoint{1.166846in}{6.797273in}}%
\pgfpathlineto{\pgfqpoint{1.170012in}{7.723105in}}%
\pgfpathlineto{\pgfqpoint{1.170268in}{7.731115in}}%
\pgfpathlineto{\pgfqpoint{1.170782in}{7.710696in}}%
\pgfpathlineto{\pgfqpoint{1.171894in}{7.507418in}}%
\pgfpathlineto{\pgfqpoint{1.174461in}{6.491786in}}%
\pgfpathlineto{\pgfqpoint{1.177627in}{5.558669in}}%
\pgfpathlineto{\pgfqpoint{1.177969in}{5.548786in}}%
\pgfpathlineto{\pgfqpoint{1.178397in}{5.566614in}}%
\pgfpathlineto{\pgfqpoint{1.179509in}{5.762919in}}%
\pgfpathlineto{\pgfqpoint{1.181990in}{6.730564in}}%
\pgfpathlineto{\pgfqpoint{1.185327in}{7.722272in}}%
\pgfpathlineto{\pgfqpoint{1.185584in}{7.730122in}}%
\pgfpathlineto{\pgfqpoint{1.186097in}{7.709832in}}%
\pgfpathlineto{\pgfqpoint{1.187210in}{7.508725in}}%
\pgfpathlineto{\pgfqpoint{1.189777in}{6.501886in}}%
\pgfpathlineto{\pgfqpoint{1.193028in}{5.557180in}}%
\pgfpathlineto{\pgfqpoint{1.193285in}{5.549861in}}%
\pgfpathlineto{\pgfqpoint{1.193798in}{5.570987in}}%
\pgfpathlineto{\pgfqpoint{1.194910in}{5.772782in}}%
\pgfpathlineto{\pgfqpoint{1.197477in}{6.776876in}}%
\pgfpathlineto{\pgfqpoint{1.200729in}{7.721266in}}%
\pgfpathlineto{\pgfqpoint{1.200985in}{7.729087in}}%
\pgfpathlineto{\pgfqpoint{1.201499in}{7.709183in}}%
\pgfpathlineto{\pgfqpoint{1.202611in}{7.510752in}}%
\pgfpathlineto{\pgfqpoint{1.205178in}{6.513245in}}%
\pgfpathlineto{\pgfqpoint{1.208429in}{5.560672in}}%
\pgfpathlineto{\pgfqpoint{1.208771in}{5.550833in}}%
\pgfpathlineto{\pgfqpoint{1.209199in}{5.567973in}}%
\pgfpathlineto{\pgfqpoint{1.210226in}{5.737325in}}%
\pgfpathlineto{\pgfqpoint{1.212536in}{6.595187in}}%
\pgfpathlineto{\pgfqpoint{1.216130in}{7.714765in}}%
\pgfpathlineto{\pgfqpoint{1.216558in}{7.728051in}}%
\pgfpathlineto{\pgfqpoint{1.216985in}{7.708849in}}%
\pgfpathlineto{\pgfqpoint{1.218098in}{7.513828in}}%
\pgfpathlineto{\pgfqpoint{1.220579in}{6.563885in}}%
\pgfpathlineto{\pgfqpoint{1.223916in}{5.564820in}}%
\pgfpathlineto{\pgfqpoint{1.224258in}{5.552030in}}%
\pgfpathlineto{\pgfqpoint{1.224772in}{5.571588in}}%
\pgfpathlineto{\pgfqpoint{1.225884in}{5.766467in}}%
\pgfpathlineto{\pgfqpoint{1.228365in}{6.712968in}}%
\pgfpathlineto{\pgfqpoint{1.231702in}{7.713093in}}%
\pgfpathlineto{\pgfqpoint{1.232130in}{7.727045in}}%
\pgfpathlineto{\pgfqpoint{1.232558in}{7.708918in}}%
\pgfpathlineto{\pgfqpoint{1.233670in}{7.518266in}}%
\pgfpathlineto{\pgfqpoint{1.236151in}{6.579276in}}%
\pgfpathlineto{\pgfqpoint{1.239574in}{5.563923in}}%
\pgfpathlineto{\pgfqpoint{1.239916in}{5.552983in}}%
\pgfpathlineto{\pgfqpoint{1.240430in}{5.574816in}}%
\pgfpathlineto{\pgfqpoint{1.241542in}{5.772274in}}%
\pgfpathlineto{\pgfqpoint{1.244109in}{6.751758in}}%
\pgfpathlineto{\pgfqpoint{1.247446in}{7.716445in}}%
\pgfpathlineto{\pgfqpoint{1.247788in}{7.726008in}}%
\pgfpathlineto{\pgfqpoint{1.248216in}{7.709449in}}%
\pgfpathlineto{\pgfqpoint{1.249242in}{7.545372in}}%
\pgfpathlineto{\pgfqpoint{1.251553in}{6.708827in}}%
\pgfpathlineto{\pgfqpoint{1.255232in}{5.569205in}}%
\pgfpathlineto{\pgfqpoint{1.255660in}{5.554085in}}%
\pgfpathlineto{\pgfqpoint{1.256173in}{5.577450in}}%
\pgfpathlineto{\pgfqpoint{1.257371in}{5.798924in}}%
\pgfpathlineto{\pgfqpoint{1.260194in}{6.895242in}}%
\pgfpathlineto{\pgfqpoint{1.263275in}{7.718501in}}%
\pgfpathlineto{\pgfqpoint{1.263531in}{7.724871in}}%
\pgfpathlineto{\pgfqpoint{1.263959in}{7.710462in}}%
\pgfpathlineto{\pgfqpoint{1.264986in}{7.552766in}}%
\pgfpathlineto{\pgfqpoint{1.267211in}{6.766565in}}%
\pgfpathlineto{\pgfqpoint{1.271061in}{5.569340in}}%
\pgfpathlineto{\pgfqpoint{1.271489in}{5.555246in}}%
\pgfpathlineto{\pgfqpoint{1.272002in}{5.579299in}}%
\pgfpathlineto{\pgfqpoint{1.273200in}{5.800078in}}%
\pgfpathlineto{\pgfqpoint{1.276024in}{6.888250in}}%
\pgfpathlineto{\pgfqpoint{1.279104in}{7.715643in}}%
\pgfpathlineto{\pgfqpoint{1.279446in}{7.723661in}}%
\pgfpathlineto{\pgfqpoint{1.279874in}{7.705913in}}%
\pgfpathlineto{\pgfqpoint{1.280986in}{7.521597in}}%
\pgfpathlineto{\pgfqpoint{1.283467in}{6.609083in}}%
\pgfpathlineto{\pgfqpoint{1.286976in}{5.570263in}}%
\pgfpathlineto{\pgfqpoint{1.287403in}{5.556406in}}%
\pgfpathlineto{\pgfqpoint{1.287917in}{5.580205in}}%
\pgfpathlineto{\pgfqpoint{1.289115in}{5.798317in}}%
\pgfpathlineto{\pgfqpoint{1.291853in}{6.841049in}}%
\pgfpathlineto{\pgfqpoint{1.295018in}{7.711914in}}%
\pgfpathlineto{\pgfqpoint{1.295361in}{7.722558in}}%
\pgfpathlineto{\pgfqpoint{1.295874in}{7.701997in}}%
\pgfpathlineto{\pgfqpoint{1.296986in}{7.513705in}}%
\pgfpathlineto{\pgfqpoint{1.299468in}{6.603777in}}%
\pgfpathlineto{\pgfqpoint{1.302976in}{5.572050in}}%
\pgfpathlineto{\pgfqpoint{1.303404in}{5.557561in}}%
\pgfpathlineto{\pgfqpoint{1.303917in}{5.580068in}}%
\pgfpathlineto{\pgfqpoint{1.305115in}{5.793335in}}%
\pgfpathlineto{\pgfqpoint{1.307853in}{6.824520in}}%
\pgfpathlineto{\pgfqpoint{1.311019in}{7.706972in}}%
\pgfpathlineto{\pgfqpoint{1.311446in}{7.721376in}}%
\pgfpathlineto{\pgfqpoint{1.311960in}{7.699036in}}%
\pgfpathlineto{\pgfqpoint{1.313158in}{7.487185in}}%
\pgfpathlineto{\pgfqpoint{1.315896in}{6.460959in}}%
\pgfpathlineto{\pgfqpoint{1.319147in}{5.569258in}}%
\pgfpathlineto{\pgfqpoint{1.319489in}{5.558785in}}%
\pgfpathlineto{\pgfqpoint{1.320003in}{5.578871in}}%
\pgfpathlineto{\pgfqpoint{1.321115in}{5.763379in}}%
\pgfpathlineto{\pgfqpoint{1.323596in}{6.659305in}}%
\pgfpathlineto{\pgfqpoint{1.327190in}{7.706600in}}%
\pgfpathlineto{\pgfqpoint{1.327618in}{7.720132in}}%
\pgfpathlineto{\pgfqpoint{1.328131in}{7.697284in}}%
\pgfpathlineto{\pgfqpoint{1.329329in}{7.486489in}}%
\pgfpathlineto{\pgfqpoint{1.332067in}{6.468677in}}%
\pgfpathlineto{\pgfqpoint{1.335318in}{5.572845in}}%
\pgfpathlineto{\pgfqpoint{1.335746in}{5.560062in}}%
\pgfpathlineto{\pgfqpoint{1.336174in}{5.576707in}}%
\pgfpathlineto{\pgfqpoint{1.337286in}{5.752079in}}%
\pgfpathlineto{\pgfqpoint{1.339682in}{6.595858in}}%
\pgfpathlineto{\pgfqpoint{1.343447in}{7.705106in}}%
\pgfpathlineto{\pgfqpoint{1.343875in}{7.718901in}}%
\pgfpathlineto{\pgfqpoint{1.344388in}{7.696900in}}%
\pgfpathlineto{\pgfqpoint{1.345586in}{7.490014in}}%
\pgfpathlineto{\pgfqpoint{1.348238in}{6.518383in}}%
\pgfpathlineto{\pgfqpoint{1.351575in}{5.577875in}}%
\pgfpathlineto{\pgfqpoint{1.352003in}{5.561331in}}%
\pgfpathlineto{\pgfqpoint{1.352516in}{5.579773in}}%
\pgfpathlineto{\pgfqpoint{1.353629in}{5.757237in}}%
\pgfpathlineto{\pgfqpoint{1.356024in}{6.596781in}}%
\pgfpathlineto{\pgfqpoint{1.359789in}{7.702366in}}%
\pgfpathlineto{\pgfqpoint{1.360217in}{7.717625in}}%
\pgfpathlineto{\pgfqpoint{1.360730in}{7.697909in}}%
\pgfpathlineto{\pgfqpoint{1.361843in}{7.518854in}}%
\pgfpathlineto{\pgfqpoint{1.364324in}{6.645623in}}%
\pgfpathlineto{\pgfqpoint{1.368003in}{5.578332in}}%
\pgfpathlineto{\pgfqpoint{1.368431in}{5.562584in}}%
\pgfpathlineto{\pgfqpoint{1.368944in}{5.581449in}}%
\pgfpathlineto{\pgfqpoint{1.370057in}{5.757869in}}%
\pgfpathlineto{\pgfqpoint{1.372452in}{6.589766in}}%
\pgfpathlineto{\pgfqpoint{1.376217in}{7.698096in}}%
\pgfpathlineto{\pgfqpoint{1.376731in}{7.716272in}}%
\pgfpathlineto{\pgfqpoint{1.377244in}{7.693566in}}%
\pgfpathlineto{\pgfqpoint{1.378442in}{7.489427in}}%
\pgfpathlineto{\pgfqpoint{1.381094in}{6.534956in}}%
\pgfpathlineto{\pgfqpoint{1.384517in}{5.580280in}}%
\pgfpathlineto{\pgfqpoint{1.384945in}{5.563951in}}%
\pgfpathlineto{\pgfqpoint{1.385458in}{5.581589in}}%
\pgfpathlineto{\pgfqpoint{1.386570in}{5.753682in}}%
\pgfpathlineto{\pgfqpoint{1.388966in}{6.574279in}}%
\pgfpathlineto{\pgfqpoint{1.392816in}{7.698443in}}%
\pgfpathlineto{\pgfqpoint{1.393244in}{7.714899in}}%
\pgfpathlineto{\pgfqpoint{1.393757in}{7.697679in}}%
\pgfpathlineto{\pgfqpoint{1.394870in}{7.527359in}}%
\pgfpathlineto{\pgfqpoint{1.397266in}{6.712014in}}%
\pgfpathlineto{\pgfqpoint{1.401116in}{5.583908in}}%
\pgfpathlineto{\pgfqpoint{1.401629in}{5.565229in}}%
\pgfpathlineto{\pgfqpoint{1.402143in}{5.586569in}}%
\pgfpathlineto{\pgfqpoint{1.403340in}{5.784654in}}%
\pgfpathlineto{\pgfqpoint{1.405993in}{6.722367in}}%
\pgfpathlineto{\pgfqpoint{1.409501in}{7.697051in}}%
\pgfpathlineto{\pgfqpoint{1.409929in}{7.713520in}}%
\pgfpathlineto{\pgfqpoint{1.410442in}{7.696844in}}%
\pgfpathlineto{\pgfqpoint{1.411554in}{7.529489in}}%
\pgfpathlineto{\pgfqpoint{1.413950in}{6.724049in}}%
\pgfpathlineto{\pgfqpoint{1.417886in}{5.583098in}}%
\pgfpathlineto{\pgfqpoint{1.418314in}{5.566710in}}%
\pgfpathlineto{\pgfqpoint{1.418827in}{5.583220in}}%
\pgfpathlineto{\pgfqpoint{1.419940in}{5.749303in}}%
\pgfpathlineto{\pgfqpoint{1.422335in}{6.550010in}}%
\pgfpathlineto{\pgfqpoint{1.426271in}{7.693814in}}%
\pgfpathlineto{\pgfqpoint{1.426785in}{7.712246in}}%
\pgfpathlineto{\pgfqpoint{1.427298in}{7.691521in}}%
\pgfpathlineto{\pgfqpoint{1.428496in}{7.497899in}}%
\pgfpathlineto{\pgfqpoint{1.431063in}{6.609729in}}%
\pgfpathlineto{\pgfqpoint{1.434742in}{5.584241in}}%
\pgfpathlineto{\pgfqpoint{1.435170in}{5.568107in}}%
\pgfpathlineto{\pgfqpoint{1.435683in}{5.584395in}}%
\pgfpathlineto{\pgfqpoint{1.436795in}{5.748164in}}%
\pgfpathlineto{\pgfqpoint{1.439191in}{6.539631in}}%
\pgfpathlineto{\pgfqpoint{1.443213in}{7.694731in}}%
\pgfpathlineto{\pgfqpoint{1.443640in}{7.710723in}}%
\pgfpathlineto{\pgfqpoint{1.444154in}{7.694528in}}%
\pgfpathlineto{\pgfqpoint{1.445266in}{7.531879in}}%
\pgfpathlineto{\pgfqpoint{1.447662in}{6.745005in}}%
\pgfpathlineto{\pgfqpoint{1.451683in}{5.587392in}}%
\pgfpathlineto{\pgfqpoint{1.452197in}{5.569427in}}%
\pgfpathlineto{\pgfqpoint{1.452710in}{5.589760in}}%
\pgfpathlineto{\pgfqpoint{1.453908in}{5.779402in}}%
\pgfpathlineto{\pgfqpoint{1.456475in}{6.652842in}}%
\pgfpathlineto{\pgfqpoint{1.460154in}{7.687221in}}%
\pgfpathlineto{\pgfqpoint{1.460667in}{7.709280in}}%
\pgfpathlineto{\pgfqpoint{1.461181in}{7.693328in}}%
\pgfpathlineto{\pgfqpoint{1.462293in}{7.533032in}}%
\pgfpathlineto{\pgfqpoint{1.464603in}{6.788839in}}%
\pgfpathlineto{\pgfqpoint{1.468796in}{5.586692in}}%
\pgfpathlineto{\pgfqpoint{1.469224in}{5.570996in}}%
\pgfpathlineto{\pgfqpoint{1.469737in}{5.586761in}}%
\pgfpathlineto{\pgfqpoint{1.470849in}{5.745751in}}%
\pgfpathlineto{\pgfqpoint{1.473159in}{6.485262in}}%
\pgfpathlineto{\pgfqpoint{1.477352in}{7.690134in}}%
\pgfpathlineto{\pgfqpoint{1.477865in}{7.707917in}}%
\pgfpathlineto{\pgfqpoint{1.478379in}{7.688257in}}%
\pgfpathlineto{\pgfqpoint{1.479577in}{7.503184in}}%
\pgfpathlineto{\pgfqpoint{1.482144in}{6.645206in}}%
\pgfpathlineto{\pgfqpoint{1.485908in}{5.594500in}}%
\pgfpathlineto{\pgfqpoint{1.486507in}{5.572462in}}%
\pgfpathlineto{\pgfqpoint{1.487021in}{5.593808in}}%
\pgfpathlineto{\pgfqpoint{1.488218in}{5.781431in}}%
\pgfpathlineto{\pgfqpoint{1.490785in}{6.638420in}}%
\pgfpathlineto{\pgfqpoint{1.494550in}{7.684095in}}%
\pgfpathlineto{\pgfqpoint{1.495149in}{7.706343in}}%
\pgfpathlineto{\pgfqpoint{1.495662in}{7.685484in}}%
\pgfpathlineto{\pgfqpoint{1.496860in}{7.499938in}}%
\pgfpathlineto{\pgfqpoint{1.499427in}{6.648693in}}%
\pgfpathlineto{\pgfqpoint{1.503192in}{5.598838in}}%
\pgfpathlineto{\pgfqpoint{1.503791in}{5.573878in}}%
\pgfpathlineto{\pgfqpoint{1.504304in}{5.592107in}}%
\pgfpathlineto{\pgfqpoint{1.505416in}{5.752329in}}%
\pgfpathlineto{\pgfqpoint{1.507812in}{6.513253in}}%
\pgfpathlineto{\pgfqpoint{1.511919in}{7.681757in}}%
\pgfpathlineto{\pgfqpoint{1.512518in}{7.704862in}}%
\pgfpathlineto{\pgfqpoint{1.513032in}{7.685348in}}%
\pgfpathlineto{\pgfqpoint{1.514229in}{7.504750in}}%
\pgfpathlineto{\pgfqpoint{1.516796in}{6.665911in}}%
\pgfpathlineto{\pgfqpoint{1.520647in}{5.599223in}}%
\pgfpathlineto{\pgfqpoint{1.521246in}{5.575435in}}%
\pgfpathlineto{\pgfqpoint{1.521759in}{5.594059in}}%
\pgfpathlineto{\pgfqpoint{1.522871in}{5.753227in}}%
\pgfpathlineto{\pgfqpoint{1.525267in}{6.506287in}}%
\pgfpathlineto{\pgfqpoint{1.529460in}{7.683083in}}%
\pgfpathlineto{\pgfqpoint{1.529973in}{7.703246in}}%
\pgfpathlineto{\pgfqpoint{1.530486in}{7.687662in}}%
\pgfpathlineto{\pgfqpoint{1.531599in}{7.535518in}}%
\pgfpathlineto{\pgfqpoint{1.533909in}{6.826302in}}%
\pgfpathlineto{\pgfqpoint{1.538272in}{5.596185in}}%
\pgfpathlineto{\pgfqpoint{1.538786in}{5.577042in}}%
\pgfpathlineto{\pgfqpoint{1.539299in}{5.593366in}}%
\pgfpathlineto{\pgfqpoint{1.540412in}{5.746080in}}%
\pgfpathlineto{\pgfqpoint{1.542722in}{6.452813in}}%
\pgfpathlineto{\pgfqpoint{1.547085in}{7.681315in}}%
\pgfpathlineto{\pgfqpoint{1.547599in}{7.701615in}}%
\pgfpathlineto{\pgfqpoint{1.548112in}{7.686730in}}%
\pgfpathlineto{\pgfqpoint{1.549224in}{7.537834in}}%
\pgfpathlineto{\pgfqpoint{1.551535in}{6.838908in}}%
\pgfpathlineto{\pgfqpoint{1.555984in}{5.596435in}}%
\pgfpathlineto{\pgfqpoint{1.556497in}{5.578632in}}%
\pgfpathlineto{\pgfqpoint{1.557011in}{5.595736in}}%
\pgfpathlineto{\pgfqpoint{1.558123in}{5.748182in}}%
\pgfpathlineto{\pgfqpoint{1.560433in}{6.448255in}}%
\pgfpathlineto{\pgfqpoint{1.564882in}{7.682574in}}%
\pgfpathlineto{\pgfqpoint{1.565396in}{7.700086in}}%
\pgfpathlineto{\pgfqpoint{1.565909in}{7.682972in}}%
\pgfpathlineto{\pgfqpoint{1.567021in}{7.531420in}}%
\pgfpathlineto{\pgfqpoint{1.569332in}{6.835599in}}%
\pgfpathlineto{\pgfqpoint{1.573781in}{5.599750in}}%
\pgfpathlineto{\pgfqpoint{1.574294in}{5.580352in}}%
\pgfpathlineto{\pgfqpoint{1.574808in}{5.595303in}}%
\pgfpathlineto{\pgfqpoint{1.575920in}{5.741580in}}%
\pgfpathlineto{\pgfqpoint{1.578230in}{6.427758in}}%
\pgfpathlineto{\pgfqpoint{1.582765in}{7.680651in}}%
\pgfpathlineto{\pgfqpoint{1.583278in}{7.698413in}}%
\pgfpathlineto{\pgfqpoint{1.583792in}{7.682105in}}%
\pgfpathlineto{\pgfqpoint{1.584904in}{7.534002in}}%
\pgfpathlineto{\pgfqpoint{1.587214in}{6.848746in}}%
\pgfpathlineto{\pgfqpoint{1.591749in}{5.600322in}}%
\pgfpathlineto{\pgfqpoint{1.592262in}{5.581994in}}%
\pgfpathlineto{\pgfqpoint{1.592776in}{5.597458in}}%
\pgfpathlineto{\pgfqpoint{1.593888in}{5.742946in}}%
\pgfpathlineto{\pgfqpoint{1.596198in}{6.421792in}}%
\pgfpathlineto{\pgfqpoint{1.600733in}{7.675381in}}%
\pgfpathlineto{\pgfqpoint{1.601332in}{7.696696in}}%
\pgfpathlineto{\pgfqpoint{1.601845in}{7.678675in}}%
\pgfpathlineto{\pgfqpoint{1.603043in}{7.511587in}}%
\pgfpathlineto{\pgfqpoint{1.605525in}{6.755124in}}%
\pgfpathlineto{\pgfqpoint{1.609717in}{5.610910in}}%
\pgfpathlineto{\pgfqpoint{1.610402in}{5.583727in}}%
\pgfpathlineto{\pgfqpoint{1.610915in}{5.602079in}}%
\pgfpathlineto{\pgfqpoint{1.612113in}{5.768856in}}%
\pgfpathlineto{\pgfqpoint{1.614594in}{6.521462in}}%
\pgfpathlineto{\pgfqpoint{1.618872in}{7.672585in}}%
\pgfpathlineto{\pgfqpoint{1.619471in}{7.694985in}}%
\pgfpathlineto{\pgfqpoint{1.619985in}{7.678493in}}%
\pgfpathlineto{\pgfqpoint{1.621097in}{7.533654in}}%
\pgfpathlineto{\pgfqpoint{1.623407in}{6.865142in}}%
\pgfpathlineto{\pgfqpoint{1.628027in}{5.605877in}}%
\pgfpathlineto{\pgfqpoint{1.628626in}{5.585472in}}%
\pgfpathlineto{\pgfqpoint{1.629140in}{5.603376in}}%
\pgfpathlineto{\pgfqpoint{1.630338in}{5.767149in}}%
\pgfpathlineto{\pgfqpoint{1.632819in}{6.509334in}}%
\pgfpathlineto{\pgfqpoint{1.637097in}{7.665636in}}%
\pgfpathlineto{\pgfqpoint{1.637782in}{7.693201in}}%
\pgfpathlineto{\pgfqpoint{1.638295in}{7.676101in}}%
\pgfpathlineto{\pgfqpoint{1.639493in}{7.515075in}}%
\pgfpathlineto{\pgfqpoint{1.641889in}{6.809984in}}%
\pgfpathlineto{\pgfqpoint{1.646338in}{5.611408in}}%
\pgfpathlineto{\pgfqpoint{1.646937in}{5.587283in}}%
\pgfpathlineto{\pgfqpoint{1.647450in}{5.601399in}}%
\pgfpathlineto{\pgfqpoint{1.648562in}{5.738722in}}%
\pgfpathlineto{\pgfqpoint{1.650787in}{6.358551in}}%
\pgfpathlineto{\pgfqpoint{1.655664in}{7.673783in}}%
\pgfpathlineto{\pgfqpoint{1.656178in}{7.691364in}}%
\pgfpathlineto{\pgfqpoint{1.656691in}{7.677083in}}%
\pgfpathlineto{\pgfqpoint{1.657803in}{7.540327in}}%
\pgfpathlineto{\pgfqpoint{1.660028in}{6.924239in}}%
\pgfpathlineto{\pgfqpoint{1.664905in}{5.608593in}}%
\pgfpathlineto{\pgfqpoint{1.665504in}{5.589095in}}%
\pgfpathlineto{\pgfqpoint{1.666017in}{5.606588in}}%
\pgfpathlineto{\pgfqpoint{1.667215in}{5.765419in}}%
\pgfpathlineto{\pgfqpoint{1.669696in}{6.487852in}}%
\pgfpathlineto{\pgfqpoint{1.674146in}{7.665365in}}%
\pgfpathlineto{\pgfqpoint{1.674745in}{7.689470in}}%
\pgfpathlineto{\pgfqpoint{1.675258in}{7.676226in}}%
\pgfpathlineto{\pgfqpoint{1.676370in}{7.543372in}}%
\pgfpathlineto{\pgfqpoint{1.678595in}{6.938303in}}%
\pgfpathlineto{\pgfqpoint{1.683558in}{5.610128in}}%
\pgfpathlineto{\pgfqpoint{1.684156in}{5.590956in}}%
\pgfpathlineto{\pgfqpoint{1.684670in}{5.608139in}}%
\pgfpathlineto{\pgfqpoint{1.685868in}{5.764283in}}%
\pgfpathlineto{\pgfqpoint{1.688263in}{6.446922in}}%
\pgfpathlineto{\pgfqpoint{1.692884in}{7.664787in}}%
\pgfpathlineto{\pgfqpoint{1.693483in}{7.687637in}}%
\pgfpathlineto{\pgfqpoint{1.693996in}{7.673904in}}%
\pgfpathlineto{\pgfqpoint{1.695108in}{7.541847in}}%
\pgfpathlineto{\pgfqpoint{1.697333in}{6.943851in}}%
\pgfpathlineto{\pgfqpoint{1.702381in}{5.610178in}}%
\pgfpathlineto{\pgfqpoint{1.702895in}{5.592880in}}%
\pgfpathlineto{\pgfqpoint{1.703408in}{5.606083in}}%
\pgfpathlineto{\pgfqpoint{1.704520in}{5.736171in}}%
\pgfpathlineto{\pgfqpoint{1.706745in}{6.328600in}}%
\pgfpathlineto{\pgfqpoint{1.711879in}{7.671026in}}%
\pgfpathlineto{\pgfqpoint{1.712392in}{7.685760in}}%
\pgfpathlineto{\pgfqpoint{1.712905in}{7.670263in}}%
\pgfpathlineto{\pgfqpoint{1.714018in}{7.536409in}}%
\pgfpathlineto{\pgfqpoint{1.716328in}{6.913670in}}%
\pgfpathlineto{\pgfqpoint{1.721291in}{5.614483in}}%
\pgfpathlineto{\pgfqpoint{1.721890in}{5.594755in}}%
\pgfpathlineto{\pgfqpoint{1.722403in}{5.610291in}}%
\pgfpathlineto{\pgfqpoint{1.723515in}{5.743335in}}%
\pgfpathlineto{\pgfqpoint{1.725825in}{6.361844in}}%
\pgfpathlineto{\pgfqpoint{1.730788in}{7.661839in}}%
\pgfpathlineto{\pgfqpoint{1.731387in}{7.683766in}}%
\pgfpathlineto{\pgfqpoint{1.731900in}{7.670410in}}%
\pgfpathlineto{\pgfqpoint{1.733013in}{7.542677in}}%
\pgfpathlineto{\pgfqpoint{1.735237in}{6.962374in}}%
\pgfpathlineto{\pgfqpoint{1.740457in}{5.612275in}}%
\pgfpathlineto{\pgfqpoint{1.740970in}{5.596740in}}%
\pgfpathlineto{\pgfqpoint{1.741483in}{5.610630in}}%
\pgfpathlineto{\pgfqpoint{1.742596in}{5.738563in}}%
\pgfpathlineto{\pgfqpoint{1.744820in}{6.316112in}}%
\pgfpathlineto{\pgfqpoint{1.750040in}{7.664743in}}%
\pgfpathlineto{\pgfqpoint{1.750553in}{7.681707in}}%
\pgfpathlineto{\pgfqpoint{1.751066in}{7.669516in}}%
\pgfpathlineto{\pgfqpoint{1.752179in}{7.545926in}}%
\pgfpathlineto{\pgfqpoint{1.754403in}{6.977114in}}%
\pgfpathlineto{\pgfqpoint{1.759708in}{5.614564in}}%
\pgfpathlineto{\pgfqpoint{1.760221in}{5.598778in}}%
\pgfpathlineto{\pgfqpoint{1.760735in}{5.611883in}}%
\pgfpathlineto{\pgfqpoint{1.761847in}{5.736454in}}%
\pgfpathlineto{\pgfqpoint{1.764072in}{6.303550in}}%
\pgfpathlineto{\pgfqpoint{1.769377in}{7.662805in}}%
\pgfpathlineto{\pgfqpoint{1.769976in}{7.679679in}}%
\pgfpathlineto{\pgfqpoint{1.770489in}{7.663146in}}%
\pgfpathlineto{\pgfqpoint{1.771687in}{7.517392in}}%
\pgfpathlineto{\pgfqpoint{1.774083in}{6.877398in}}%
\pgfpathlineto{\pgfqpoint{1.778960in}{5.627623in}}%
\pgfpathlineto{\pgfqpoint{1.779644in}{5.600834in}}%
\pgfpathlineto{\pgfqpoint{1.780158in}{5.613791in}}%
\pgfpathlineto{\pgfqpoint{1.781270in}{5.736305in}}%
\pgfpathlineto{\pgfqpoint{1.783494in}{6.294645in}}%
\pgfpathlineto{\pgfqpoint{1.788885in}{7.661396in}}%
\pgfpathlineto{\pgfqpoint{1.789398in}{7.677602in}}%
\pgfpathlineto{\pgfqpoint{1.789912in}{7.665715in}}%
\pgfpathlineto{\pgfqpoint{1.791024in}{7.546262in}}%
\pgfpathlineto{\pgfqpoint{1.793249in}{6.995032in}}%
\pgfpathlineto{\pgfqpoint{1.798725in}{5.617514in}}%
\pgfpathlineto{\pgfqpoint{1.799238in}{5.602917in}}%
\pgfpathlineto{\pgfqpoint{1.799751in}{5.616150in}}%
\pgfpathlineto{\pgfqpoint{1.800864in}{5.737475in}}%
\pgfpathlineto{\pgfqpoint{1.803088in}{6.288216in}}%
\pgfpathlineto{\pgfqpoint{1.808564in}{7.660274in}}%
\pgfpathlineto{\pgfqpoint{1.809078in}{7.675523in}}%
\pgfpathlineto{\pgfqpoint{1.809591in}{7.663206in}}%
\pgfpathlineto{\pgfqpoint{1.810703in}{7.544620in}}%
\pgfpathlineto{\pgfqpoint{1.812928in}{7.000571in}}%
\pgfpathlineto{\pgfqpoint{1.818490in}{5.618634in}}%
\pgfpathlineto{\pgfqpoint{1.819003in}{5.605054in}}%
\pgfpathlineto{\pgfqpoint{1.819516in}{5.618776in}}%
\pgfpathlineto{\pgfqpoint{1.820629in}{5.739362in}}%
\pgfpathlineto{\pgfqpoint{1.822853in}{6.283110in}}%
\pgfpathlineto{\pgfqpoint{1.828415in}{7.659209in}}%
\pgfpathlineto{\pgfqpoint{1.828928in}{7.673388in}}%
\pgfpathlineto{\pgfqpoint{1.829442in}{7.660526in}}%
\pgfpathlineto{\pgfqpoint{1.830554in}{7.542562in}}%
\pgfpathlineto{\pgfqpoint{1.832778in}{7.005355in}}%
\pgfpathlineto{\pgfqpoint{1.838340in}{5.624488in}}%
\pgfpathlineto{\pgfqpoint{1.838939in}{5.607256in}}%
\pgfpathlineto{\pgfqpoint{1.839452in}{5.621485in}}%
\pgfpathlineto{\pgfqpoint{1.840565in}{5.741375in}}%
\pgfpathlineto{\pgfqpoint{1.842789in}{6.278201in}}%
\pgfpathlineto{\pgfqpoint{1.848351in}{7.653230in}}%
\pgfpathlineto{\pgfqpoint{1.848950in}{7.671185in}}%
\pgfpathlineto{\pgfqpoint{1.849463in}{7.657857in}}%
\pgfpathlineto{\pgfqpoint{1.850575in}{7.540667in}}%
\pgfpathlineto{\pgfqpoint{1.852800in}{7.010496in}}%
\pgfpathlineto{\pgfqpoint{1.858447in}{5.625699in}}%
\pgfpathlineto{\pgfqpoint{1.859046in}{5.609515in}}%
\pgfpathlineto{\pgfqpoint{1.859559in}{5.624080in}}%
\pgfpathlineto{\pgfqpoint{1.860757in}{5.756745in}}%
\pgfpathlineto{\pgfqpoint{1.863068in}{6.324050in}}%
\pgfpathlineto{\pgfqpoint{1.868458in}{7.646553in}}%
\pgfpathlineto{\pgfqpoint{1.869142in}{7.668923in}}%
\pgfpathlineto{\pgfqpoint{1.869656in}{7.655393in}}%
\pgfpathlineto{\pgfqpoint{1.870768in}{7.539506in}}%
\pgfpathlineto{\pgfqpoint{1.872993in}{7.017083in}}%
\pgfpathlineto{\pgfqpoint{1.878725in}{5.627360in}}%
\pgfpathlineto{\pgfqpoint{1.879324in}{5.611810in}}%
\pgfpathlineto{\pgfqpoint{1.879838in}{5.626357in}}%
\pgfpathlineto{\pgfqpoint{1.881036in}{5.757054in}}%
\pgfpathlineto{\pgfqpoint{1.883346in}{6.315511in}}%
\pgfpathlineto{\pgfqpoint{1.888822in}{7.644737in}}%
\pgfpathlineto{\pgfqpoint{1.889506in}{7.666621in}}%
\pgfpathlineto{\pgfqpoint{1.890020in}{7.653326in}}%
\pgfpathlineto{\pgfqpoint{1.891132in}{7.539630in}}%
\pgfpathlineto{\pgfqpoint{1.893357in}{7.026176in}}%
\pgfpathlineto{\pgfqpoint{1.899175in}{5.629650in}}%
\pgfpathlineto{\pgfqpoint{1.899774in}{5.614119in}}%
\pgfpathlineto{\pgfqpoint{1.900287in}{5.628126in}}%
\pgfpathlineto{\pgfqpoint{1.901485in}{5.755721in}}%
\pgfpathlineto{\pgfqpoint{1.903795in}{6.303878in}}%
\pgfpathlineto{\pgfqpoint{1.909357in}{7.642068in}}%
\pgfpathlineto{\pgfqpoint{1.910041in}{7.664278in}}%
\pgfpathlineto{\pgfqpoint{1.910555in}{7.651826in}}%
\pgfpathlineto{\pgfqpoint{1.911667in}{7.541553in}}%
\pgfpathlineto{\pgfqpoint{1.913892in}{7.038792in}}%
\pgfpathlineto{\pgfqpoint{1.919881in}{5.628391in}}%
\pgfpathlineto{\pgfqpoint{1.920394in}{5.616457in}}%
\pgfpathlineto{\pgfqpoint{1.920908in}{5.629238in}}%
\pgfpathlineto{\pgfqpoint{1.922020in}{5.739339in}}%
\pgfpathlineto{\pgfqpoint{1.924245in}{6.238895in}}%
\pgfpathlineto{\pgfqpoint{1.930234in}{7.648224in}}%
\pgfpathlineto{\pgfqpoint{1.930747in}{7.661841in}}%
\pgfpathlineto{\pgfqpoint{1.931261in}{7.650995in}}%
\pgfpathlineto{\pgfqpoint{1.932373in}{7.545710in}}%
\pgfpathlineto{\pgfqpoint{1.934512in}{7.079799in}}%
\pgfpathlineto{\pgfqpoint{1.940758in}{5.628484in}}%
\pgfpathlineto{\pgfqpoint{1.941186in}{5.618915in}}%
\pgfpathlineto{\pgfqpoint{1.941699in}{5.629631in}}%
\pgfpathlineto{\pgfqpoint{1.942812in}{5.733807in}}%
\pgfpathlineto{\pgfqpoint{1.944951in}{6.195381in}}%
\pgfpathlineto{\pgfqpoint{1.951282in}{7.651387in}}%
\pgfpathlineto{\pgfqpoint{1.951710in}{7.659431in}}%
\pgfpathlineto{\pgfqpoint{1.952224in}{7.647120in}}%
\pgfpathlineto{\pgfqpoint{1.953336in}{7.540505in}}%
\pgfpathlineto{\pgfqpoint{1.955561in}{7.054723in}}%
\pgfpathlineto{\pgfqpoint{1.961721in}{5.633443in}}%
\pgfpathlineto{\pgfqpoint{1.962234in}{5.621349in}}%
\pgfpathlineto{\pgfqpoint{1.962748in}{5.632958in}}%
\pgfpathlineto{\pgfqpoint{1.963860in}{5.737297in}}%
\pgfpathlineto{\pgfqpoint{1.965999in}{6.193507in}}%
\pgfpathlineto{\pgfqpoint{1.972416in}{7.649468in}}%
\pgfpathlineto{\pgfqpoint{1.972844in}{7.656921in}}%
\pgfpathlineto{\pgfqpoint{1.973358in}{7.644363in}}%
\pgfpathlineto{\pgfqpoint{1.974470in}{7.538909in}}%
\pgfpathlineto{\pgfqpoint{1.976695in}{7.060665in}}%
\pgfpathlineto{\pgfqpoint{1.982941in}{5.635856in}}%
\pgfpathlineto{\pgfqpoint{1.983454in}{5.623873in}}%
\pgfpathlineto{\pgfqpoint{1.983967in}{5.635094in}}%
\pgfpathlineto{\pgfqpoint{1.985080in}{5.736969in}}%
\pgfpathlineto{\pgfqpoint{1.987219in}{6.184174in}}%
\pgfpathlineto{\pgfqpoint{1.993721in}{7.646400in}}%
\pgfpathlineto{\pgfqpoint{1.994149in}{7.654387in}}%
\pgfpathlineto{\pgfqpoint{1.994663in}{7.642926in}}%
\pgfpathlineto{\pgfqpoint{1.995775in}{7.541386in}}%
\pgfpathlineto{\pgfqpoint{1.997914in}{7.097427in}}%
\pgfpathlineto{\pgfqpoint{2.004502in}{5.632864in}}%
\pgfpathlineto{\pgfqpoint{2.004930in}{5.626522in}}%
\pgfpathlineto{\pgfqpoint{2.005358in}{5.635953in}}%
\pgfpathlineto{\pgfqpoint{2.006470in}{5.732472in}}%
\pgfpathlineto{\pgfqpoint{2.008609in}{6.166575in}}%
\pgfpathlineto{\pgfqpoint{2.015283in}{7.645113in}}%
\pgfpathlineto{\pgfqpoint{2.015711in}{7.651740in}}%
\pgfpathlineto{\pgfqpoint{2.016139in}{7.642765in}}%
\pgfpathlineto{\pgfqpoint{2.017251in}{7.548158in}}%
\pgfpathlineto{\pgfqpoint{2.019390in}{7.119549in}}%
\pgfpathlineto{\pgfqpoint{2.026150in}{5.634920in}}%
\pgfpathlineto{\pgfqpoint{2.026492in}{5.629112in}}%
\pgfpathlineto{\pgfqpoint{2.027005in}{5.638919in}}%
\pgfpathlineto{\pgfqpoint{2.028118in}{5.734580in}}%
\pgfpathlineto{\pgfqpoint{2.030257in}{6.161926in}}%
\pgfpathlineto{\pgfqpoint{2.037102in}{7.644946in}}%
\pgfpathlineto{\pgfqpoint{2.037444in}{7.649101in}}%
\pgfpathlineto{\pgfqpoint{2.037872in}{7.640560in}}%
\pgfpathlineto{\pgfqpoint{2.038898in}{7.559107in}}%
\pgfpathlineto{\pgfqpoint{2.040952in}{7.172027in}}%
\pgfpathlineto{\pgfqpoint{2.048225in}{5.632610in}}%
\pgfpathlineto{\pgfqpoint{2.048396in}{5.631789in}}%
\pgfpathlineto{\pgfqpoint{2.048738in}{5.637393in}}%
\pgfpathlineto{\pgfqpoint{2.049679in}{5.701817in}}%
\pgfpathlineto{\pgfqpoint{2.051562in}{6.023566in}}%
\pgfpathlineto{\pgfqpoint{2.059348in}{7.646392in}}%
\pgfpathlineto{\pgfqpoint{2.059690in}{7.642319in}}%
\pgfpathlineto{\pgfqpoint{2.060631in}{7.582533in}}%
\pgfpathlineto{\pgfqpoint{2.062428in}{7.289456in}}%
\pgfpathlineto{\pgfqpoint{2.070471in}{5.634556in}}%
\pgfpathlineto{\pgfqpoint{2.070899in}{5.643242in}}%
\pgfpathlineto{\pgfqpoint{2.072011in}{5.733264in}}%
\pgfpathlineto{\pgfqpoint{2.074150in}{6.141502in}}%
\pgfpathlineto{\pgfqpoint{2.081252in}{7.640142in}}%
\pgfpathlineto{\pgfqpoint{2.081509in}{7.643586in}}%
\pgfpathlineto{\pgfqpoint{2.082022in}{7.634722in}}%
\pgfpathlineto{\pgfqpoint{2.083134in}{7.545073in}}%
\pgfpathlineto{\pgfqpoint{2.085273in}{7.140114in}}%
\pgfpathlineto{\pgfqpoint{2.092461in}{5.639958in}}%
\pgfpathlineto{\pgfqpoint{2.092717in}{5.637294in}}%
\pgfpathlineto{\pgfqpoint{2.093145in}{5.644392in}}%
\pgfpathlineto{\pgfqpoint{2.094172in}{5.719150in}}%
\pgfpathlineto{\pgfqpoint{2.096140in}{6.063322in}}%
\pgfpathlineto{\pgfqpoint{2.103926in}{7.640802in}}%
\pgfpathlineto{\pgfqpoint{2.104011in}{7.640675in}}%
\pgfpathlineto{\pgfqpoint{2.104525in}{7.627960in}}%
\pgfpathlineto{\pgfqpoint{2.105723in}{7.520881in}}%
\pgfpathlineto{\pgfqpoint{2.108033in}{7.059462in}}%
\pgfpathlineto{\pgfqpoint{2.114621in}{5.654710in}}%
\pgfpathlineto{\pgfqpoint{2.115220in}{5.640160in}}%
\pgfpathlineto{\pgfqpoint{2.115733in}{5.649642in}}%
\pgfpathlineto{\pgfqpoint{2.116846in}{5.738167in}}%
\pgfpathlineto{\pgfqpoint{2.118985in}{6.133322in}}%
\pgfpathlineto{\pgfqpoint{2.126343in}{7.635765in}}%
\pgfpathlineto{\pgfqpoint{2.126600in}{7.637926in}}%
\pgfpathlineto{\pgfqpoint{2.127028in}{7.630393in}}%
\pgfpathlineto{\pgfqpoint{2.128054in}{7.556615in}}%
\pgfpathlineto{\pgfqpoint{2.130022in}{7.221013in}}%
\pgfpathlineto{\pgfqpoint{2.137980in}{5.643067in}}%
\pgfpathlineto{\pgfqpoint{2.138065in}{5.643296in}}%
\pgfpathlineto{\pgfqpoint{2.138579in}{5.656199in}}%
\pgfpathlineto{\pgfqpoint{2.139777in}{5.761005in}}%
\pgfpathlineto{\pgfqpoint{2.142087in}{6.210079in}}%
\pgfpathlineto{\pgfqpoint{2.148846in}{7.621380in}}%
\pgfpathlineto{\pgfqpoint{2.149445in}{7.634990in}}%
\pgfpathlineto{\pgfqpoint{2.149958in}{7.625472in}}%
\pgfpathlineto{\pgfqpoint{2.151071in}{7.539247in}}%
\pgfpathlineto{\pgfqpoint{2.153210in}{7.155565in}}%
\pgfpathlineto{\pgfqpoint{2.160739in}{5.648126in}}%
\pgfpathlineto{\pgfqpoint{2.160996in}{5.646056in}}%
\pgfpathlineto{\pgfqpoint{2.161424in}{5.653346in}}%
\pgfpathlineto{\pgfqpoint{2.162451in}{5.724599in}}%
\pgfpathlineto{\pgfqpoint{2.164419in}{6.049315in}}%
\pgfpathlineto{\pgfqpoint{2.172547in}{7.631971in}}%
\pgfpathlineto{\pgfqpoint{2.172718in}{7.631129in}}%
\pgfpathlineto{\pgfqpoint{2.173403in}{7.606636in}}%
\pgfpathlineto{\pgfqpoint{2.174857in}{7.447601in}}%
\pgfpathlineto{\pgfqpoint{2.177681in}{6.826574in}}%
\pgfpathlineto{\pgfqpoint{2.183071in}{5.695556in}}%
\pgfpathlineto{\pgfqpoint{2.184183in}{5.649118in}}%
\pgfpathlineto{\pgfqpoint{2.184611in}{5.654728in}}%
\pgfpathlineto{\pgfqpoint{2.185638in}{5.720788in}}%
\pgfpathlineto{\pgfqpoint{2.187606in}{6.032814in}}%
\pgfpathlineto{\pgfqpoint{2.195906in}{7.628894in}}%
\pgfpathlineto{\pgfqpoint{2.196162in}{7.626948in}}%
\pgfpathlineto{\pgfqpoint{2.196932in}{7.593360in}}%
\pgfpathlineto{\pgfqpoint{2.198558in}{7.394172in}}%
\pgfpathlineto{\pgfqpoint{2.201809in}{6.639621in}}%
\pgfpathlineto{\pgfqpoint{2.206430in}{5.708364in}}%
\pgfpathlineto{\pgfqpoint{2.207713in}{5.652193in}}%
\pgfpathlineto{\pgfqpoint{2.208055in}{5.656522in}}%
\pgfpathlineto{\pgfqpoint{2.208997in}{5.710043in}}%
\pgfpathlineto{\pgfqpoint{2.210793in}{5.967050in}}%
\pgfpathlineto{\pgfqpoint{2.215072in}{7.010877in}}%
\pgfpathlineto{\pgfqpoint{2.218580in}{7.593808in}}%
\pgfpathlineto{\pgfqpoint{2.219521in}{7.625752in}}%
\pgfpathlineto{\pgfqpoint{2.220034in}{7.617504in}}%
\pgfpathlineto{\pgfqpoint{2.221146in}{7.538579in}}%
\pgfpathlineto{\pgfqpoint{2.223286in}{7.181754in}}%
\pgfpathlineto{\pgfqpoint{2.231328in}{5.655874in}}%
\pgfpathlineto{\pgfqpoint{2.231414in}{5.655408in}}%
\pgfpathlineto{\pgfqpoint{2.231756in}{5.658533in}}%
\pgfpathlineto{\pgfqpoint{2.232612in}{5.700894in}}%
\pgfpathlineto{\pgfqpoint{2.234323in}{5.922913in}}%
\pgfpathlineto{\pgfqpoint{2.237917in}{6.764707in}}%
\pgfpathlineto{\pgfqpoint{2.242109in}{7.564931in}}%
\pgfpathlineto{\pgfqpoint{2.243393in}{7.622532in}}%
\pgfpathlineto{\pgfqpoint{2.243820in}{7.617285in}}%
\pgfpathlineto{\pgfqpoint{2.244847in}{7.555264in}}%
\pgfpathlineto{\pgfqpoint{2.246815in}{7.261063in}}%
\pgfpathlineto{\pgfqpoint{2.255457in}{5.658604in}}%
\pgfpathlineto{\pgfqpoint{2.255799in}{5.662321in}}%
\pgfpathlineto{\pgfqpoint{2.256740in}{5.712140in}}%
\pgfpathlineto{\pgfqpoint{2.258537in}{5.955343in}}%
\pgfpathlineto{\pgfqpoint{2.262559in}{6.902581in}}%
\pgfpathlineto{\pgfqpoint{2.266409in}{7.575179in}}%
\pgfpathlineto{\pgfqpoint{2.267607in}{7.619252in}}%
\pgfpathlineto{\pgfqpoint{2.267949in}{7.614656in}}%
\pgfpathlineto{\pgfqpoint{2.268890in}{7.562992in}}%
\pgfpathlineto{\pgfqpoint{2.270773in}{7.302939in}}%
\pgfpathlineto{\pgfqpoint{2.275136in}{6.275023in}}%
\pgfpathlineto{\pgfqpoint{2.278730in}{5.695438in}}%
\pgfpathlineto{\pgfqpoint{2.279757in}{5.661914in}}%
\pgfpathlineto{\pgfqpoint{2.280185in}{5.667999in}}%
\pgfpathlineto{\pgfqpoint{2.281211in}{5.730088in}}%
\pgfpathlineto{\pgfqpoint{2.283179in}{6.017602in}}%
\pgfpathlineto{\pgfqpoint{2.291992in}{7.615942in}}%
\pgfpathlineto{\pgfqpoint{2.292420in}{7.610266in}}%
\pgfpathlineto{\pgfqpoint{2.293447in}{7.549748in}}%
\pgfpathlineto{\pgfqpoint{2.295415in}{7.266972in}}%
\pgfpathlineto{\pgfqpoint{2.304313in}{5.665296in}}%
\pgfpathlineto{\pgfqpoint{2.304741in}{5.670967in}}%
\pgfpathlineto{\pgfqpoint{2.305768in}{5.730855in}}%
\pgfpathlineto{\pgfqpoint{2.307736in}{6.010392in}}%
\pgfpathlineto{\pgfqpoint{2.316720in}{7.612519in}}%
\pgfpathlineto{\pgfqpoint{2.317148in}{7.606489in}}%
\pgfpathlineto{\pgfqpoint{2.318174in}{7.546384in}}%
\pgfpathlineto{\pgfqpoint{2.320142in}{7.268730in}}%
\pgfpathlineto{\pgfqpoint{2.329212in}{5.668806in}}%
\pgfpathlineto{\pgfqpoint{2.329640in}{5.675516in}}%
\pgfpathlineto{\pgfqpoint{2.330752in}{5.744461in}}%
\pgfpathlineto{\pgfqpoint{2.332805in}{6.045990in}}%
\pgfpathlineto{\pgfqpoint{2.341704in}{7.609030in}}%
\pgfpathlineto{\pgfqpoint{2.341875in}{7.608170in}}%
\pgfpathlineto{\pgfqpoint{2.342560in}{7.587106in}}%
\pgfpathlineto{\pgfqpoint{2.344014in}{7.452539in}}%
\pgfpathlineto{\pgfqpoint{2.346752in}{6.937292in}}%
\pgfpathlineto{\pgfqpoint{2.353169in}{5.711417in}}%
\pgfpathlineto{\pgfqpoint{2.354282in}{5.672323in}}%
\pgfpathlineto{\pgfqpoint{2.354709in}{5.676850in}}%
\pgfpathlineto{\pgfqpoint{2.355736in}{5.731629in}}%
\pgfpathlineto{\pgfqpoint{2.357619in}{5.978982in}}%
\pgfpathlineto{\pgfqpoint{2.361897in}{6.932257in}}%
\pgfpathlineto{\pgfqpoint{2.365747in}{7.559198in}}%
\pgfpathlineto{\pgfqpoint{2.367030in}{7.605441in}}%
\pgfpathlineto{\pgfqpoint{2.367373in}{7.601493in}}%
\pgfpathlineto{\pgfqpoint{2.368314in}{7.555552in}}%
\pgfpathlineto{\pgfqpoint{2.370111in}{7.335790in}}%
\pgfpathlineto{\pgfqpoint{2.373790in}{6.542694in}}%
\pgfpathlineto{\pgfqpoint{2.378239in}{5.743133in}}%
\pgfpathlineto{\pgfqpoint{2.379779in}{5.675897in}}%
\pgfpathlineto{\pgfqpoint{2.380036in}{5.677950in}}%
\pgfpathlineto{\pgfqpoint{2.380892in}{5.712164in}}%
\pgfpathlineto{\pgfqpoint{2.382603in}{5.899232in}}%
\pgfpathlineto{\pgfqpoint{2.385854in}{6.562289in}}%
\pgfpathlineto{\pgfqpoint{2.391073in}{7.534837in}}%
\pgfpathlineto{\pgfqpoint{2.392614in}{7.601809in}}%
\pgfpathlineto{\pgfqpoint{2.392870in}{7.599898in}}%
\pgfpathlineto{\pgfqpoint{2.393726in}{7.566535in}}%
\pgfpathlineto{\pgfqpoint{2.395437in}{7.382653in}}%
\pgfpathlineto{\pgfqpoint{2.398689in}{6.727614in}}%
\pgfpathlineto{\pgfqpoint{2.403993in}{5.745720in}}%
\pgfpathlineto{\pgfqpoint{2.405534in}{5.679584in}}%
\pgfpathlineto{\pgfqpoint{2.405790in}{5.681451in}}%
\pgfpathlineto{\pgfqpoint{2.406646in}{5.714289in}}%
\pgfpathlineto{\pgfqpoint{2.408357in}{5.895572in}}%
\pgfpathlineto{\pgfqpoint{2.411608in}{6.543139in}}%
\pgfpathlineto{\pgfqpoint{2.416999in}{7.533234in}}%
\pgfpathlineto{\pgfqpoint{2.418539in}{7.598077in}}%
\pgfpathlineto{\pgfqpoint{2.418796in}{7.596176in}}%
\pgfpathlineto{\pgfqpoint{2.419651in}{7.563603in}}%
\pgfpathlineto{\pgfqpoint{2.421363in}{7.384445in}}%
\pgfpathlineto{\pgfqpoint{2.424614in}{6.743897in}}%
\pgfpathlineto{\pgfqpoint{2.430004in}{5.753764in}}%
\pgfpathlineto{\pgfqpoint{2.431630in}{5.683362in}}%
\pgfpathlineto{\pgfqpoint{2.431887in}{5.685358in}}%
\pgfpathlineto{\pgfqpoint{2.432742in}{5.717863in}}%
\pgfpathlineto{\pgfqpoint{2.434454in}{5.895258in}}%
\pgfpathlineto{\pgfqpoint{2.437619in}{6.509783in}}%
\pgfpathlineto{\pgfqpoint{2.443181in}{7.525871in}}%
\pgfpathlineto{\pgfqpoint{2.444807in}{7.594246in}}%
\pgfpathlineto{\pgfqpoint{2.445063in}{7.592115in}}%
\pgfpathlineto{\pgfqpoint{2.445919in}{7.559539in}}%
\pgfpathlineto{\pgfqpoint{2.447630in}{7.383653in}}%
\pgfpathlineto{\pgfqpoint{2.450796in}{6.775357in}}%
\pgfpathlineto{\pgfqpoint{2.456443in}{5.753463in}}%
\pgfpathlineto{\pgfqpoint{2.458069in}{5.687256in}}%
\pgfpathlineto{\pgfqpoint{2.458326in}{5.689545in}}%
\pgfpathlineto{\pgfqpoint{2.459181in}{5.722270in}}%
\pgfpathlineto{\pgfqpoint{2.460892in}{5.896797in}}%
\pgfpathlineto{\pgfqpoint{2.464144in}{6.518029in}}%
\pgfpathlineto{\pgfqpoint{2.469705in}{7.519196in}}%
\pgfpathlineto{\pgfqpoint{2.471331in}{7.590336in}}%
\pgfpathlineto{\pgfqpoint{2.471588in}{7.589033in}}%
\pgfpathlineto{\pgfqpoint{2.472358in}{7.564522in}}%
\pgfpathlineto{\pgfqpoint{2.473898in}{7.426971in}}%
\pgfpathlineto{\pgfqpoint{2.476721in}{6.933492in}}%
\pgfpathlineto{\pgfqpoint{2.483395in}{5.740389in}}%
\pgfpathlineto{\pgfqpoint{2.484764in}{5.691166in}}%
\pgfpathlineto{\pgfqpoint{2.485107in}{5.693889in}}%
\pgfpathlineto{\pgfqpoint{2.486048in}{5.732253in}}%
\pgfpathlineto{\pgfqpoint{2.487845in}{5.922980in}}%
\pgfpathlineto{\pgfqpoint{2.491267in}{6.582659in}}%
\pgfpathlineto{\pgfqpoint{2.496572in}{7.512411in}}%
\pgfpathlineto{\pgfqpoint{2.498283in}{7.586373in}}%
\pgfpathlineto{\pgfqpoint{2.498540in}{7.584752in}}%
\pgfpathlineto{\pgfqpoint{2.499396in}{7.555324in}}%
\pgfpathlineto{\pgfqpoint{2.501021in}{7.402793in}}%
\pgfpathlineto{\pgfqpoint{2.504016in}{6.869305in}}%
\pgfpathlineto{\pgfqpoint{2.510348in}{5.753623in}}%
\pgfpathlineto{\pgfqpoint{2.511888in}{5.695194in}}%
\pgfpathlineto{\pgfqpoint{2.512144in}{5.696919in}}%
\pgfpathlineto{\pgfqpoint{2.513000in}{5.726326in}}%
\pgfpathlineto{\pgfqpoint{2.514626in}{5.877355in}}%
\pgfpathlineto{\pgfqpoint{2.517620in}{6.404896in}}%
\pgfpathlineto{\pgfqpoint{2.524038in}{7.525183in}}%
\pgfpathlineto{\pgfqpoint{2.525578in}{7.582290in}}%
\pgfpathlineto{\pgfqpoint{2.525834in}{7.580517in}}%
\pgfpathlineto{\pgfqpoint{2.526690in}{7.551312in}}%
\pgfpathlineto{\pgfqpoint{2.528401in}{7.391111in}}%
\pgfpathlineto{\pgfqpoint{2.531481in}{6.845108in}}%
\pgfpathlineto{\pgfqpoint{2.537728in}{5.761927in}}%
\pgfpathlineto{\pgfqpoint{2.539353in}{5.699323in}}%
\pgfpathlineto{\pgfqpoint{2.539610in}{5.701074in}}%
\pgfpathlineto{\pgfqpoint{2.540466in}{5.729850in}}%
\pgfpathlineto{\pgfqpoint{2.542177in}{5.887704in}}%
\pgfpathlineto{\pgfqpoint{2.545257in}{6.426756in}}%
\pgfpathlineto{\pgfqpoint{2.551589in}{7.515997in}}%
\pgfpathlineto{\pgfqpoint{2.553214in}{7.578119in}}%
\pgfpathlineto{\pgfqpoint{2.553471in}{7.576473in}}%
\pgfpathlineto{\pgfqpoint{2.554327in}{7.548395in}}%
\pgfpathlineto{\pgfqpoint{2.555952in}{7.404012in}}%
\pgfpathlineto{\pgfqpoint{2.558947in}{6.896720in}}%
\pgfpathlineto{\pgfqpoint{2.565621in}{5.759496in}}%
\pgfpathlineto{\pgfqpoint{2.567161in}{5.703542in}}%
\pgfpathlineto{\pgfqpoint{2.567418in}{5.704989in}}%
\pgfpathlineto{\pgfqpoint{2.568188in}{5.727826in}}%
\pgfpathlineto{\pgfqpoint{2.569813in}{5.863099in}}%
\pgfpathlineto{\pgfqpoint{2.572723in}{6.338555in}}%
\pgfpathlineto{\pgfqpoint{2.579739in}{7.523031in}}%
\pgfpathlineto{\pgfqpoint{2.581193in}{7.573826in}}%
\pgfpathlineto{\pgfqpoint{2.581536in}{7.571629in}}%
\pgfpathlineto{\pgfqpoint{2.582391in}{7.542579in}}%
\pgfpathlineto{\pgfqpoint{2.584102in}{7.388884in}}%
\pgfpathlineto{\pgfqpoint{2.587183in}{6.867569in}}%
\pgfpathlineto{\pgfqpoint{2.593685in}{5.772523in}}%
\pgfpathlineto{\pgfqpoint{2.595397in}{5.707867in}}%
\pgfpathlineto{\pgfqpoint{2.595568in}{5.708677in}}%
\pgfpathlineto{\pgfqpoint{2.596252in}{5.725209in}}%
\pgfpathlineto{\pgfqpoint{2.597707in}{5.828611in}}%
\pgfpathlineto{\pgfqpoint{2.600274in}{6.201622in}}%
\pgfpathlineto{\pgfqpoint{2.608744in}{7.551973in}}%
\pgfpathlineto{\pgfqpoint{2.609600in}{7.569463in}}%
\pgfpathlineto{\pgfqpoint{2.610113in}{7.564222in}}%
\pgfpathlineto{\pgfqpoint{2.611226in}{7.512915in}}%
\pgfpathlineto{\pgfqpoint{2.613279in}{7.286803in}}%
\pgfpathlineto{\pgfqpoint{2.617386in}{6.519914in}}%
\pgfpathlineto{\pgfqpoint{2.622092in}{5.789832in}}%
\pgfpathlineto{\pgfqpoint{2.623974in}{5.712280in}}%
\pgfpathlineto{\pgfqpoint{2.624146in}{5.712893in}}%
\pgfpathlineto{\pgfqpoint{2.624830in}{5.728221in}}%
\pgfpathlineto{\pgfqpoint{2.626285in}{5.827028in}}%
\pgfpathlineto{\pgfqpoint{2.628852in}{6.187549in}}%
\pgfpathlineto{\pgfqpoint{2.637664in}{7.552325in}}%
\pgfpathlineto{\pgfqpoint{2.638435in}{7.565010in}}%
\pgfpathlineto{\pgfqpoint{2.638948in}{7.559184in}}%
\pgfpathlineto{\pgfqpoint{2.640060in}{7.507931in}}%
\pgfpathlineto{\pgfqpoint{2.642114in}{7.286223in}}%
\pgfpathlineto{\pgfqpoint{2.646221in}{6.535509in}}%
\pgfpathlineto{\pgfqpoint{2.651012in}{5.798572in}}%
\pgfpathlineto{\pgfqpoint{2.652980in}{5.716803in}}%
\pgfpathlineto{\pgfqpoint{2.653066in}{5.716950in}}%
\pgfpathlineto{\pgfqpoint{2.653579in}{5.724385in}}%
\pgfpathlineto{\pgfqpoint{2.654777in}{5.784627in}}%
\pgfpathlineto{\pgfqpoint{2.656916in}{6.027222in}}%
\pgfpathlineto{\pgfqpoint{2.661536in}{6.880742in}}%
\pgfpathlineto{\pgfqpoint{2.665815in}{7.491289in}}%
\pgfpathlineto{\pgfqpoint{2.667611in}{7.560417in}}%
\pgfpathlineto{\pgfqpoint{2.667782in}{7.560006in}}%
\pgfpathlineto{\pgfqpoint{2.668467in}{7.546109in}}%
\pgfpathlineto{\pgfqpoint{2.669922in}{7.453399in}}%
\pgfpathlineto{\pgfqpoint{2.672488in}{7.110567in}}%
\pgfpathlineto{\pgfqpoint{2.681729in}{5.730776in}}%
\pgfpathlineto{\pgfqpoint{2.682414in}{5.721436in}}%
\pgfpathlineto{\pgfqpoint{2.682927in}{5.727110in}}%
\pgfpathlineto{\pgfqpoint{2.684039in}{5.776153in}}%
\pgfpathlineto{\pgfqpoint{2.686093in}{5.987994in}}%
\pgfpathlineto{\pgfqpoint{2.690114in}{6.695296in}}%
\pgfpathlineto{\pgfqpoint{2.695248in}{7.473763in}}%
\pgfpathlineto{\pgfqpoint{2.697302in}{7.555731in}}%
\pgfpathlineto{\pgfqpoint{2.697387in}{7.555480in}}%
\pgfpathlineto{\pgfqpoint{2.697986in}{7.545424in}}%
\pgfpathlineto{\pgfqpoint{2.699269in}{7.476051in}}%
\pgfpathlineto{\pgfqpoint{2.701580in}{7.205114in}}%
\pgfpathlineto{\pgfqpoint{2.712275in}{5.726186in}}%
\pgfpathlineto{\pgfqpoint{2.712361in}{5.726415in}}%
\pgfpathlineto{\pgfqpoint{2.712959in}{5.736181in}}%
\pgfpathlineto{\pgfqpoint{2.714243in}{5.804141in}}%
\pgfpathlineto{\pgfqpoint{2.716553in}{6.070420in}}%
\pgfpathlineto{\pgfqpoint{2.722542in}{7.134761in}}%
\pgfpathlineto{\pgfqpoint{2.726050in}{7.518562in}}%
\pgfpathlineto{\pgfqpoint{2.727334in}{7.550951in}}%
\pgfpathlineto{\pgfqpoint{2.727676in}{7.548727in}}%
\pgfpathlineto{\pgfqpoint{2.728617in}{7.519135in}}%
\pgfpathlineto{\pgfqpoint{2.730414in}{7.372104in}}%
\pgfpathlineto{\pgfqpoint{2.733580in}{6.890456in}}%
\pgfpathlineto{\pgfqpoint{2.740682in}{5.796827in}}%
\pgfpathlineto{\pgfqpoint{2.742564in}{5.731038in}}%
\pgfpathlineto{\pgfqpoint{2.742735in}{5.731760in}}%
\pgfpathlineto{\pgfqpoint{2.743505in}{5.748901in}}%
\pgfpathlineto{\pgfqpoint{2.745045in}{5.849095in}}%
\pgfpathlineto{\pgfqpoint{2.747698in}{6.195742in}}%
\pgfpathlineto{\pgfqpoint{2.756939in}{7.530185in}}%
\pgfpathlineto{\pgfqpoint{2.757880in}{7.546028in}}%
\pgfpathlineto{\pgfqpoint{2.758308in}{7.542167in}}%
\pgfpathlineto{\pgfqpoint{2.759334in}{7.504972in}}%
\pgfpathlineto{\pgfqpoint{2.761302in}{7.330717in}}%
\pgfpathlineto{\pgfqpoint{2.764810in}{6.777398in}}%
\pgfpathlineto{\pgfqpoint{2.771142in}{5.820179in}}%
\pgfpathlineto{\pgfqpoint{2.773281in}{5.735999in}}%
\pgfpathlineto{\pgfqpoint{2.773623in}{5.738120in}}%
\pgfpathlineto{\pgfqpoint{2.774564in}{5.766220in}}%
\pgfpathlineto{\pgfqpoint{2.776361in}{5.905944in}}%
\pgfpathlineto{\pgfqpoint{2.779527in}{6.366427in}}%
\pgfpathlineto{\pgfqpoint{2.786971in}{7.478362in}}%
\pgfpathlineto{\pgfqpoint{2.788853in}{7.541002in}}%
\pgfpathlineto{\pgfqpoint{2.789024in}{7.540341in}}%
\pgfpathlineto{\pgfqpoint{2.789709in}{7.527047in}}%
\pgfpathlineto{\pgfqpoint{2.791164in}{7.443793in}}%
\pgfpathlineto{\pgfqpoint{2.793730in}{7.138565in}}%
\pgfpathlineto{\pgfqpoint{2.803998in}{5.745611in}}%
\pgfpathlineto{\pgfqpoint{2.804511in}{5.741093in}}%
\pgfpathlineto{\pgfqpoint{2.805025in}{5.746016in}}%
\pgfpathlineto{\pgfqpoint{2.806137in}{5.788651in}}%
\pgfpathlineto{\pgfqpoint{2.808190in}{5.973894in}}%
\pgfpathlineto{\pgfqpoint{2.811955in}{6.563352in}}%
\pgfpathlineto{\pgfqpoint{2.818030in}{7.447460in}}%
\pgfpathlineto{\pgfqpoint{2.820255in}{7.535836in}}%
\pgfpathlineto{\pgfqpoint{2.820511in}{7.534979in}}%
\pgfpathlineto{\pgfqpoint{2.821282in}{7.518556in}}%
\pgfpathlineto{\pgfqpoint{2.822822in}{7.425401in}}%
\pgfpathlineto{\pgfqpoint{2.825474in}{7.103968in}}%
\pgfpathlineto{\pgfqpoint{2.835399in}{5.757351in}}%
\pgfpathlineto{\pgfqpoint{2.836169in}{5.746322in}}%
\pgfpathlineto{\pgfqpoint{2.836683in}{5.750340in}}%
\pgfpathlineto{\pgfqpoint{2.837795in}{5.789904in}}%
\pgfpathlineto{\pgfqpoint{2.839763in}{5.956478in}}%
\pgfpathlineto{\pgfqpoint{2.843357in}{6.491863in}}%
\pgfpathlineto{\pgfqpoint{2.849945in}{7.442945in}}%
\pgfpathlineto{\pgfqpoint{2.852170in}{7.530502in}}%
\pgfpathlineto{\pgfqpoint{2.852255in}{7.530569in}}%
\pgfpathlineto{\pgfqpoint{2.852426in}{7.529958in}}%
\pgfpathlineto{\pgfqpoint{2.853111in}{7.517618in}}%
\pgfpathlineto{\pgfqpoint{2.854565in}{7.440206in}}%
\pgfpathlineto{\pgfqpoint{2.857047in}{7.166980in}}%
\pgfpathlineto{\pgfqpoint{2.868170in}{5.752521in}}%
\pgfpathlineto{\pgfqpoint{2.868426in}{5.751665in}}%
\pgfpathlineto{\pgfqpoint{2.868854in}{5.755109in}}%
\pgfpathlineto{\pgfqpoint{2.869881in}{5.787960in}}%
\pgfpathlineto{\pgfqpoint{2.871763in}{5.933192in}}%
\pgfpathlineto{\pgfqpoint{2.875100in}{6.397746in}}%
\pgfpathlineto{\pgfqpoint{2.882544in}{7.450935in}}%
\pgfpathlineto{\pgfqpoint{2.884683in}{7.525164in}}%
\pgfpathlineto{\pgfqpoint{2.885026in}{7.523327in}}%
\pgfpathlineto{\pgfqpoint{2.885967in}{7.498698in}}%
\pgfpathlineto{\pgfqpoint{2.887764in}{7.375596in}}%
\pgfpathlineto{\pgfqpoint{2.890844in}{6.977322in}}%
\pgfpathlineto{\pgfqpoint{2.899400in}{5.803747in}}%
\pgfpathlineto{\pgfqpoint{2.901111in}{5.757129in}}%
\pgfpathlineto{\pgfqpoint{2.901368in}{5.758173in}}%
\pgfpathlineto{\pgfqpoint{2.902224in}{5.776817in}}%
\pgfpathlineto{\pgfqpoint{2.903849in}{5.874228in}}%
\pgfpathlineto{\pgfqpoint{2.906673in}{6.206387in}}%
\pgfpathlineto{\pgfqpoint{2.916598in}{7.501678in}}%
\pgfpathlineto{\pgfqpoint{2.917710in}{7.519592in}}%
\pgfpathlineto{\pgfqpoint{2.918138in}{7.516150in}}%
\pgfpathlineto{\pgfqpoint{2.919251in}{7.480602in}}%
\pgfpathlineto{\pgfqpoint{2.921218in}{7.329274in}}%
\pgfpathlineto{\pgfqpoint{2.924727in}{6.848725in}}%
\pgfpathlineto{\pgfqpoint{2.932085in}{5.841320in}}%
\pgfpathlineto{\pgfqpoint{2.934310in}{5.762780in}}%
\pgfpathlineto{\pgfqpoint{2.934566in}{5.763406in}}%
\pgfpathlineto{\pgfqpoint{2.935336in}{5.777421in}}%
\pgfpathlineto{\pgfqpoint{2.936876in}{5.858480in}}%
\pgfpathlineto{\pgfqpoint{2.939529in}{6.142276in}}%
\pgfpathlineto{\pgfqpoint{2.950652in}{7.509623in}}%
\pgfpathlineto{\pgfqpoint{2.951165in}{7.513929in}}%
\pgfpathlineto{\pgfqpoint{2.951679in}{7.510276in}}%
\pgfpathlineto{\pgfqpoint{2.952791in}{7.475355in}}%
\pgfpathlineto{\pgfqpoint{2.954759in}{7.328470in}}%
\pgfpathlineto{\pgfqpoint{2.958181in}{6.875553in}}%
\pgfpathlineto{\pgfqpoint{2.965882in}{5.843900in}}%
\pgfpathlineto{\pgfqpoint{2.968107in}{5.768530in}}%
\pgfpathlineto{\pgfqpoint{2.968363in}{5.769151in}}%
\pgfpathlineto{\pgfqpoint{2.969133in}{5.782680in}}%
\pgfpathlineto{\pgfqpoint{2.970674in}{5.860758in}}%
\pgfpathlineto{\pgfqpoint{2.973326in}{6.134576in}}%
\pgfpathlineto{\pgfqpoint{2.984877in}{7.505959in}}%
\pgfpathlineto{\pgfqpoint{2.985305in}{7.508082in}}%
\pgfpathlineto{\pgfqpoint{2.985733in}{7.504893in}}%
\pgfpathlineto{\pgfqpoint{2.986845in}{7.472011in}}%
\pgfpathlineto{\pgfqpoint{2.988813in}{7.331753in}}%
\pgfpathlineto{\pgfqpoint{2.992235in}{6.895061in}}%
\pgfpathlineto{\pgfqpoint{3.000278in}{5.844642in}}%
\pgfpathlineto{\pgfqpoint{3.002503in}{5.774400in}}%
\pgfpathlineto{\pgfqpoint{3.002845in}{5.775958in}}%
\pgfpathlineto{\pgfqpoint{3.003786in}{5.797311in}}%
\pgfpathlineto{\pgfqpoint{3.005583in}{5.904753in}}%
\pgfpathlineto{\pgfqpoint{3.008578in}{6.245303in}}%
\pgfpathlineto{\pgfqpoint{3.018417in}{7.471153in}}%
\pgfpathlineto{\pgfqpoint{3.019872in}{7.502111in}}%
\pgfpathlineto{\pgfqpoint{3.020214in}{7.500861in}}%
\pgfpathlineto{\pgfqpoint{3.021070in}{7.483524in}}%
\pgfpathlineto{\pgfqpoint{3.022781in}{7.390094in}}%
\pgfpathlineto{\pgfqpoint{3.025605in}{7.089805in}}%
\pgfpathlineto{\pgfqpoint{3.036471in}{5.793731in}}%
\pgfpathlineto{\pgfqpoint{3.037498in}{5.780470in}}%
\pgfpathlineto{\pgfqpoint{3.037926in}{5.783430in}}%
\pgfpathlineto{\pgfqpoint{3.039038in}{5.814272in}}%
\pgfpathlineto{\pgfqpoint{3.041006in}{5.946323in}}%
\pgfpathlineto{\pgfqpoint{3.044343in}{6.348291in}}%
\pgfpathlineto{\pgfqpoint{3.052985in}{7.433679in}}%
\pgfpathlineto{\pgfqpoint{3.055124in}{7.495993in}}%
\pgfpathlineto{\pgfqpoint{3.055209in}{7.495967in}}%
\pgfpathlineto{\pgfqpoint{3.055723in}{7.491702in}}%
\pgfpathlineto{\pgfqpoint{3.056921in}{7.454695in}}%
\pgfpathlineto{\pgfqpoint{3.059060in}{7.300501in}}%
\pgfpathlineto{\pgfqpoint{3.062739in}{6.837906in}}%
\pgfpathlineto{\pgfqpoint{3.070354in}{5.877212in}}%
\pgfpathlineto{\pgfqpoint{3.072750in}{5.787507in}}%
\pgfpathlineto{\pgfqpoint{3.073006in}{5.786654in}}%
\pgfpathlineto{\pgfqpoint{3.073434in}{5.789067in}}%
\pgfpathlineto{\pgfqpoint{3.074461in}{5.814254in}}%
\pgfpathlineto{\pgfqpoint{3.076343in}{5.928256in}}%
\pgfpathlineto{\pgfqpoint{3.079509in}{6.281576in}}%
\pgfpathlineto{\pgfqpoint{3.089178in}{7.446988in}}%
\pgfpathlineto{\pgfqpoint{3.090974in}{7.489706in}}%
\pgfpathlineto{\pgfqpoint{3.091145in}{7.489485in}}%
\pgfpathlineto{\pgfqpoint{3.091744in}{7.482807in}}%
\pgfpathlineto{\pgfqpoint{3.093113in}{7.433648in}}%
\pgfpathlineto{\pgfqpoint{3.095424in}{7.252440in}}%
\pgfpathlineto{\pgfqpoint{3.099616in}{6.707456in}}%
\pgfpathlineto{\pgfqpoint{3.106204in}{5.901401in}}%
\pgfpathlineto{\pgfqpoint{3.108771in}{5.795034in}}%
\pgfpathlineto{\pgfqpoint{3.109199in}{5.793028in}}%
\pgfpathlineto{\pgfqpoint{3.109627in}{5.795617in}}%
\pgfpathlineto{\pgfqpoint{3.110654in}{5.820416in}}%
\pgfpathlineto{\pgfqpoint{3.112536in}{5.930920in}}%
\pgfpathlineto{\pgfqpoint{3.115702in}{6.272573in}}%
\pgfpathlineto{\pgfqpoint{3.125713in}{7.444521in}}%
\pgfpathlineto{\pgfqpoint{3.127510in}{7.483272in}}%
\pgfpathlineto{\pgfqpoint{3.127681in}{7.482848in}}%
\pgfpathlineto{\pgfqpoint{3.128365in}{7.473977in}}%
\pgfpathlineto{\pgfqpoint{3.129820in}{7.417776in}}%
\pgfpathlineto{\pgfqpoint{3.132301in}{7.215649in}}%
\pgfpathlineto{\pgfqpoint{3.136921in}{6.611470in}}%
\pgfpathlineto{\pgfqpoint{3.142911in}{5.910471in}}%
\pgfpathlineto{\pgfqpoint{3.145563in}{5.801629in}}%
\pgfpathlineto{\pgfqpoint{3.145991in}{5.799550in}}%
\pgfpathlineto{\pgfqpoint{3.146504in}{5.802864in}}%
\pgfpathlineto{\pgfqpoint{3.147617in}{5.831567in}}%
\pgfpathlineto{\pgfqpoint{3.149585in}{5.950792in}}%
\pgfpathlineto{\pgfqpoint{3.152922in}{6.313596in}}%
\pgfpathlineto{\pgfqpoint{3.162505in}{7.424054in}}%
\pgfpathlineto{\pgfqpoint{3.164644in}{7.476659in}}%
\pgfpathlineto{\pgfqpoint{3.165071in}{7.474356in}}%
\pgfpathlineto{\pgfqpoint{3.166098in}{7.451398in}}%
\pgfpathlineto{\pgfqpoint{3.167981in}{7.348144in}}%
\pgfpathlineto{\pgfqpoint{3.171146in}{7.026224in}}%
\pgfpathlineto{\pgfqpoint{3.181756in}{5.838927in}}%
\pgfpathlineto{\pgfqpoint{3.183467in}{5.806253in}}%
\pgfpathlineto{\pgfqpoint{3.183639in}{5.806676in}}%
\pgfpathlineto{\pgfqpoint{3.184409in}{5.816896in}}%
\pgfpathlineto{\pgfqpoint{3.185949in}{5.877243in}}%
\pgfpathlineto{\pgfqpoint{3.188516in}{6.084262in}}%
\pgfpathlineto{\pgfqpoint{3.193393in}{6.704000in}}%
\pgfpathlineto{\pgfqpoint{3.199296in}{7.361378in}}%
\pgfpathlineto{\pgfqpoint{3.201949in}{7.467246in}}%
\pgfpathlineto{\pgfqpoint{3.202462in}{7.469866in}}%
\pgfpathlineto{\pgfqpoint{3.202976in}{7.466553in}}%
\pgfpathlineto{\pgfqpoint{3.204088in}{7.439224in}}%
\pgfpathlineto{\pgfqpoint{3.206141in}{7.320029in}}%
\pgfpathlineto{\pgfqpoint{3.209478in}{6.972736in}}%
\pgfpathlineto{\pgfqpoint{3.219404in}{5.865328in}}%
\pgfpathlineto{\pgfqpoint{3.221628in}{5.813136in}}%
\pgfpathlineto{\pgfqpoint{3.222056in}{5.815501in}}%
\pgfpathlineto{\pgfqpoint{3.223168in}{5.840346in}}%
\pgfpathlineto{\pgfqpoint{3.225136in}{5.947442in}}%
\pgfpathlineto{\pgfqpoint{3.228388in}{6.270032in}}%
\pgfpathlineto{\pgfqpoint{3.239083in}{7.426623in}}%
\pgfpathlineto{\pgfqpoint{3.240965in}{7.462895in}}%
\pgfpathlineto{\pgfqpoint{3.241051in}{7.462741in}}%
\pgfpathlineto{\pgfqpoint{3.241650in}{7.457256in}}%
\pgfpathlineto{\pgfqpoint{3.242933in}{7.419891in}}%
\pgfpathlineto{\pgfqpoint{3.245158in}{7.277808in}}%
\pgfpathlineto{\pgfqpoint{3.248923in}{6.867506in}}%
\pgfpathlineto{\pgfqpoint{3.257564in}{5.905563in}}%
\pgfpathlineto{\pgfqpoint{3.260131in}{5.821146in}}%
\pgfpathlineto{\pgfqpoint{3.260474in}{5.820190in}}%
\pgfpathlineto{\pgfqpoint{3.260901in}{5.822461in}}%
\pgfpathlineto{\pgfqpoint{3.262014in}{5.846245in}}%
\pgfpathlineto{\pgfqpoint{3.263982in}{5.948792in}}%
\pgfpathlineto{\pgfqpoint{3.267233in}{6.258736in}}%
\pgfpathlineto{\pgfqpoint{3.278356in}{7.423731in}}%
\pgfpathlineto{\pgfqpoint{3.280153in}{7.455760in}}%
\pgfpathlineto{\pgfqpoint{3.280324in}{7.455362in}}%
\pgfpathlineto{\pgfqpoint{3.281094in}{7.446135in}}%
\pgfpathlineto{\pgfqpoint{3.282634in}{7.391942in}}%
\pgfpathlineto{\pgfqpoint{3.285201in}{7.205435in}}%
\pgfpathlineto{\pgfqpoint{3.289822in}{6.668877in}}%
\pgfpathlineto{\pgfqpoint{3.296495in}{5.947952in}}%
\pgfpathlineto{\pgfqpoint{3.299319in}{5.832044in}}%
\pgfpathlineto{\pgfqpoint{3.300003in}{5.827412in}}%
\pgfpathlineto{\pgfqpoint{3.300517in}{5.830119in}}%
\pgfpathlineto{\pgfqpoint{3.301629in}{5.854006in}}%
\pgfpathlineto{\pgfqpoint{3.303597in}{5.953992in}}%
\pgfpathlineto{\pgfqpoint{3.306848in}{6.253815in}}%
\pgfpathlineto{\pgfqpoint{3.318314in}{7.418801in}}%
\pgfpathlineto{\pgfqpoint{3.320025in}{7.448428in}}%
\pgfpathlineto{\pgfqpoint{3.320282in}{7.447930in}}%
\pgfpathlineto{\pgfqpoint{3.321052in}{7.438689in}}%
\pgfpathlineto{\pgfqpoint{3.322592in}{7.386080in}}%
\pgfpathlineto{\pgfqpoint{3.325159in}{7.206337in}}%
\pgfpathlineto{\pgfqpoint{3.329694in}{6.698445in}}%
\pgfpathlineto{\pgfqpoint{3.336710in}{5.955642in}}%
\pgfpathlineto{\pgfqpoint{3.339533in}{5.840459in}}%
\pgfpathlineto{\pgfqpoint{3.340303in}{5.834840in}}%
\pgfpathlineto{\pgfqpoint{3.340817in}{5.837417in}}%
\pgfpathlineto{\pgfqpoint{3.341929in}{5.860211in}}%
\pgfpathlineto{\pgfqpoint{3.343897in}{5.955759in}}%
\pgfpathlineto{\pgfqpoint{3.347148in}{6.243333in}}%
\pgfpathlineto{\pgfqpoint{3.359042in}{7.414235in}}%
\pgfpathlineto{\pgfqpoint{3.360753in}{7.440908in}}%
\pgfpathlineto{\pgfqpoint{3.360924in}{7.440565in}}%
\pgfpathlineto{\pgfqpoint{3.361694in}{7.432246in}}%
\pgfpathlineto{\pgfqpoint{3.363234in}{7.382969in}}%
\pgfpathlineto{\pgfqpoint{3.365801in}{7.212367in}}%
\pgfpathlineto{\pgfqpoint{3.370250in}{6.734515in}}%
\pgfpathlineto{\pgfqpoint{3.377780in}{5.957095in}}%
\pgfpathlineto{\pgfqpoint{3.380603in}{5.847659in}}%
\pgfpathlineto{\pgfqpoint{3.381373in}{5.842465in}}%
\pgfpathlineto{\pgfqpoint{3.381887in}{5.845035in}}%
\pgfpathlineto{\pgfqpoint{3.382999in}{5.867029in}}%
\pgfpathlineto{\pgfqpoint{3.384967in}{5.958701in}}%
\pgfpathlineto{\pgfqpoint{3.388218in}{6.234850in}}%
\pgfpathlineto{\pgfqpoint{3.400625in}{7.411750in}}%
\pgfpathlineto{\pgfqpoint{3.402165in}{7.433189in}}%
\pgfpathlineto{\pgfqpoint{3.402422in}{7.432648in}}%
\pgfpathlineto{\pgfqpoint{3.403192in}{7.423971in}}%
\pgfpathlineto{\pgfqpoint{3.404732in}{7.375509in}}%
\pgfpathlineto{\pgfqpoint{3.407299in}{7.210321in}}%
\pgfpathlineto{\pgfqpoint{3.411748in}{6.748438in}}%
\pgfpathlineto{\pgfqpoint{3.419534in}{5.964518in}}%
\pgfpathlineto{\pgfqpoint{3.422443in}{5.855162in}}%
\pgfpathlineto{\pgfqpoint{3.423213in}{5.850288in}}%
\pgfpathlineto{\pgfqpoint{3.423727in}{5.852793in}}%
\pgfpathlineto{\pgfqpoint{3.424839in}{5.873887in}}%
\pgfpathlineto{\pgfqpoint{3.426807in}{5.961602in}}%
\pgfpathlineto{\pgfqpoint{3.429973in}{6.217880in}}%
\pgfpathlineto{\pgfqpoint{3.443064in}{7.409972in}}%
\pgfpathlineto{\pgfqpoint{3.444433in}{7.425260in}}%
\pgfpathlineto{\pgfqpoint{3.444775in}{7.424095in}}%
\pgfpathlineto{\pgfqpoint{3.445716in}{7.410638in}}%
\pgfpathlineto{\pgfqpoint{3.447513in}{7.344384in}}%
\pgfpathlineto{\pgfqpoint{3.450422in}{7.136977in}}%
\pgfpathlineto{\pgfqpoint{3.455984in}{6.538224in}}%
\pgfpathlineto{\pgfqpoint{3.462059in}{5.973406in}}%
\pgfpathlineto{\pgfqpoint{3.464968in}{5.864335in}}%
\pgfpathlineto{\pgfqpoint{3.465823in}{5.858310in}}%
\pgfpathlineto{\pgfqpoint{3.466337in}{5.860545in}}%
\pgfpathlineto{\pgfqpoint{3.467449in}{5.880331in}}%
\pgfpathlineto{\pgfqpoint{3.469417in}{5.963500in}}%
\pgfpathlineto{\pgfqpoint{3.472583in}{6.208156in}}%
\pgfpathlineto{\pgfqpoint{3.486273in}{7.406162in}}%
\pgfpathlineto{\pgfqpoint{3.487471in}{7.417128in}}%
\pgfpathlineto{\pgfqpoint{3.487813in}{7.415980in}}%
\pgfpathlineto{\pgfqpoint{3.488754in}{7.403052in}}%
\pgfpathlineto{\pgfqpoint{3.490551in}{7.339676in}}%
\pgfpathlineto{\pgfqpoint{3.493460in}{7.141130in}}%
\pgfpathlineto{\pgfqpoint{3.498851in}{6.582013in}}%
\pgfpathlineto{\pgfqpoint{3.505268in}{5.989765in}}%
\pgfpathlineto{\pgfqpoint{3.508262in}{5.874390in}}%
\pgfpathlineto{\pgfqpoint{3.509289in}{5.866552in}}%
\pgfpathlineto{\pgfqpoint{3.509717in}{5.868221in}}%
\pgfpathlineto{\pgfqpoint{3.510744in}{5.884004in}}%
\pgfpathlineto{\pgfqpoint{3.512626in}{5.954728in}}%
\pgfpathlineto{\pgfqpoint{3.515621in}{6.165266in}}%
\pgfpathlineto{\pgfqpoint{3.521525in}{6.778289in}}%
\pgfpathlineto{\pgfqpoint{3.527428in}{7.296148in}}%
\pgfpathlineto{\pgfqpoint{3.530423in}{7.403043in}}%
\pgfpathlineto{\pgfqpoint{3.531279in}{7.408797in}}%
\pgfpathlineto{\pgfqpoint{3.531792in}{7.406809in}}%
\pgfpathlineto{\pgfqpoint{3.532904in}{7.388594in}}%
\pgfpathlineto{\pgfqpoint{3.534872in}{7.311444in}}%
\pgfpathlineto{\pgfqpoint{3.538038in}{7.082949in}}%
\pgfpathlineto{\pgfqpoint{3.545054in}{6.358177in}}%
\pgfpathlineto{\pgfqpoint{3.550103in}{5.961864in}}%
\pgfpathlineto{\pgfqpoint{3.552841in}{5.878496in}}%
\pgfpathlineto{\pgfqpoint{3.553525in}{5.874991in}}%
\pgfpathlineto{\pgfqpoint{3.554038in}{5.877009in}}%
\pgfpathlineto{\pgfqpoint{3.555151in}{5.894952in}}%
\pgfpathlineto{\pgfqpoint{3.557119in}{5.970559in}}%
\pgfpathlineto{\pgfqpoint{3.560284in}{6.194336in}}%
\pgfpathlineto{\pgfqpoint{3.567044in}{6.882086in}}%
\pgfpathlineto{\pgfqpoint{3.572349in}{7.305437in}}%
\pgfpathlineto{\pgfqpoint{3.575172in}{7.395579in}}%
\pgfpathlineto{\pgfqpoint{3.575942in}{7.400239in}}%
\pgfpathlineto{\pgfqpoint{3.576456in}{7.398487in}}%
\pgfpathlineto{\pgfqpoint{3.577568in}{7.381439in}}%
\pgfpathlineto{\pgfqpoint{3.579536in}{7.308392in}}%
\pgfpathlineto{\pgfqpoint{3.582616in}{7.097830in}}%
\pgfpathlineto{\pgfqpoint{3.588777in}{6.483941in}}%
\pgfpathlineto{\pgfqpoint{3.594595in}{5.997877in}}%
\pgfpathlineto{\pgfqpoint{3.597590in}{5.891362in}}%
\pgfpathlineto{\pgfqpoint{3.598616in}{5.883653in}}%
\pgfpathlineto{\pgfqpoint{3.599044in}{5.884918in}}%
\pgfpathlineto{\pgfqpoint{3.600071in}{5.898647in}}%
\pgfpathlineto{\pgfqpoint{3.601953in}{5.961883in}}%
\pgfpathlineto{\pgfqpoint{3.604948in}{6.152799in}}%
\pgfpathlineto{\pgfqpoint{3.610424in}{6.678812in}}%
\pgfpathlineto{\pgfqpoint{3.617183in}{7.264294in}}%
\pgfpathlineto{\pgfqpoint{3.620349in}{7.382265in}}%
\pgfpathlineto{\pgfqpoint{3.621462in}{7.391465in}}%
\pgfpathlineto{\pgfqpoint{3.621889in}{7.390379in}}%
\pgfpathlineto{\pgfqpoint{3.622916in}{7.377333in}}%
\pgfpathlineto{\pgfqpoint{3.624799in}{7.316217in}}%
\pgfpathlineto{\pgfqpoint{3.627793in}{7.130438in}}%
\pgfpathlineto{\pgfqpoint{3.633098in}{6.632186in}}%
\pgfpathlineto{\pgfqpoint{3.640200in}{6.020487in}}%
\pgfpathlineto{\pgfqpoint{3.643366in}{5.902602in}}%
\pgfpathlineto{\pgfqpoint{3.644563in}{5.892535in}}%
\pgfpathlineto{\pgfqpoint{3.644991in}{5.893702in}}%
\pgfpathlineto{\pgfqpoint{3.646018in}{5.906689in}}%
\pgfpathlineto{\pgfqpoint{3.647900in}{5.966795in}}%
\pgfpathlineto{\pgfqpoint{3.650895in}{6.148940in}}%
\pgfpathlineto{\pgfqpoint{3.656200in}{6.638277in}}%
\pgfpathlineto{\pgfqpoint{3.663387in}{7.251358in}}%
\pgfpathlineto{\pgfqpoint{3.666639in}{7.372248in}}%
\pgfpathlineto{\pgfqpoint{3.667922in}{7.382468in}}%
\pgfpathlineto{\pgfqpoint{3.668264in}{7.381475in}}%
\pgfpathlineto{\pgfqpoint{3.669291in}{7.369143in}}%
\pgfpathlineto{\pgfqpoint{3.671088in}{7.314665in}}%
\pgfpathlineto{\pgfqpoint{3.673997in}{7.146418in}}%
\pgfpathlineto{\pgfqpoint{3.678960in}{6.705116in}}%
\pgfpathlineto{\pgfqpoint{3.686917in}{6.028744in}}%
\pgfpathlineto{\pgfqpoint{3.690168in}{5.911359in}}%
\pgfpathlineto{\pgfqpoint{3.691366in}{5.901639in}}%
\pgfpathlineto{\pgfqpoint{3.691794in}{5.902702in}}%
\pgfpathlineto{\pgfqpoint{3.692821in}{5.914953in}}%
\pgfpathlineto{\pgfqpoint{3.694703in}{5.972026in}}%
\pgfpathlineto{\pgfqpoint{3.697612in}{6.139514in}}%
\pgfpathlineto{\pgfqpoint{3.702660in}{6.583122in}}%
\pgfpathlineto{\pgfqpoint{3.710532in}{7.242968in}}%
\pgfpathlineto{\pgfqpoint{3.713783in}{7.361935in}}%
\pgfpathlineto{\pgfqpoint{3.715152in}{7.373262in}}%
\pgfpathlineto{\pgfqpoint{3.715495in}{7.372368in}}%
\pgfpathlineto{\pgfqpoint{3.716436in}{7.362250in}}%
\pgfpathlineto{\pgfqpoint{3.718233in}{7.312457in}}%
\pgfpathlineto{\pgfqpoint{3.721056in}{7.160580in}}%
\pgfpathlineto{\pgfqpoint{3.725762in}{6.765586in}}%
\pgfpathlineto{\pgfqpoint{3.734575in}{6.032938in}}%
\pgfpathlineto{\pgfqpoint{3.737827in}{5.920454in}}%
\pgfpathlineto{\pgfqpoint{3.739110in}{5.910971in}}%
\pgfpathlineto{\pgfqpoint{3.739452in}{5.911897in}}%
\pgfpathlineto{\pgfqpoint{3.740479in}{5.923370in}}%
\pgfpathlineto{\pgfqpoint{3.742276in}{5.974062in}}%
\pgfpathlineto{\pgfqpoint{3.745185in}{6.131046in}}%
\pgfpathlineto{\pgfqpoint{3.750062in}{6.538902in}}%
\pgfpathlineto{\pgfqpoint{3.758618in}{7.238595in}}%
\pgfpathlineto{\pgfqpoint{3.761870in}{7.352726in}}%
\pgfpathlineto{\pgfqpoint{3.763239in}{7.363834in}}%
\pgfpathlineto{\pgfqpoint{3.763581in}{7.363063in}}%
\pgfpathlineto{\pgfqpoint{3.764522in}{7.353646in}}%
\pgfpathlineto{\pgfqpoint{3.766319in}{7.306598in}}%
\pgfpathlineto{\pgfqpoint{3.769142in}{7.162156in}}%
\pgfpathlineto{\pgfqpoint{3.773763in}{6.791175in}}%
\pgfpathlineto{\pgfqpoint{3.783175in}{6.033650in}}%
\pgfpathlineto{\pgfqpoint{3.786340in}{5.929884in}}%
\pgfpathlineto{\pgfqpoint{3.787624in}{5.920506in}}%
\pgfpathlineto{\pgfqpoint{3.787966in}{5.921296in}}%
\pgfpathlineto{\pgfqpoint{3.788907in}{5.930590in}}%
\pgfpathlineto{\pgfqpoint{3.790704in}{5.976710in}}%
\pgfpathlineto{\pgfqpoint{3.793528in}{6.118102in}}%
\pgfpathlineto{\pgfqpoint{3.798148in}{6.481831in}}%
\pgfpathlineto{\pgfqpoint{3.807731in}{7.242012in}}%
\pgfpathlineto{\pgfqpoint{3.810897in}{7.344518in}}%
\pgfpathlineto{\pgfqpoint{3.812180in}{7.354179in}}%
\pgfpathlineto{\pgfqpoint{3.812608in}{7.353174in}}%
\pgfpathlineto{\pgfqpoint{3.813635in}{7.342172in}}%
\pgfpathlineto{\pgfqpoint{3.815517in}{7.291289in}}%
\pgfpathlineto{\pgfqpoint{3.818426in}{7.141679in}}%
\pgfpathlineto{\pgfqpoint{3.823303in}{6.754781in}}%
\pgfpathlineto{\pgfqpoint{3.832287in}{6.051984in}}%
\pgfpathlineto{\pgfqpoint{3.835624in}{5.940853in}}%
\pgfpathlineto{\pgfqpoint{3.836993in}{5.930269in}}%
\pgfpathlineto{\pgfqpoint{3.837336in}{5.930922in}}%
\pgfpathlineto{\pgfqpoint{3.838277in}{5.939509in}}%
\pgfpathlineto{\pgfqpoint{3.839988in}{5.980129in}}%
\pgfpathlineto{\pgfqpoint{3.842726in}{6.107122in}}%
\pgfpathlineto{\pgfqpoint{3.847175in}{6.436505in}}%
\pgfpathlineto{\pgfqpoint{3.857785in}{7.247942in}}%
\pgfpathlineto{\pgfqpoint{3.860865in}{7.337069in}}%
\pgfpathlineto{\pgfqpoint{3.861978in}{7.344298in}}%
\pgfpathlineto{\pgfqpoint{3.862405in}{7.343452in}}%
\pgfpathlineto{\pgfqpoint{3.863432in}{7.333231in}}%
\pgfpathlineto{\pgfqpoint{3.865229in}{7.288153in}}%
\pgfpathlineto{\pgfqpoint{3.868138in}{7.148086in}}%
\pgfpathlineto{\pgfqpoint{3.872844in}{6.793691in}}%
\pgfpathlineto{\pgfqpoint{3.882598in}{6.053221in}}%
\pgfpathlineto{\pgfqpoint{3.885850in}{5.950623in}}%
\pgfpathlineto{\pgfqpoint{3.887219in}{5.940255in}}%
\pgfpathlineto{\pgfqpoint{3.887561in}{5.940808in}}%
\pgfpathlineto{\pgfqpoint{3.888502in}{5.948807in}}%
\pgfpathlineto{\pgfqpoint{3.890213in}{5.987224in}}%
\pgfpathlineto{\pgfqpoint{3.892951in}{6.108070in}}%
\pgfpathlineto{\pgfqpoint{3.897315in}{6.416668in}}%
\pgfpathlineto{\pgfqpoint{3.908524in}{7.245441in}}%
\pgfpathlineto{\pgfqpoint{3.911604in}{7.328193in}}%
\pgfpathlineto{\pgfqpoint{3.912716in}{7.334204in}}%
\pgfpathlineto{\pgfqpoint{3.913144in}{7.333057in}}%
\pgfpathlineto{\pgfqpoint{3.914256in}{7.321123in}}%
\pgfpathlineto{\pgfqpoint{3.916224in}{7.269178in}}%
\pgfpathlineto{\pgfqpoint{3.919305in}{7.116498in}}%
\pgfpathlineto{\pgfqpoint{3.924438in}{6.728642in}}%
\pgfpathlineto{\pgfqpoint{3.933251in}{6.078298in}}%
\pgfpathlineto{\pgfqpoint{3.936674in}{5.964259in}}%
\pgfpathlineto{\pgfqpoint{3.938299in}{5.950452in}}%
\pgfpathlineto{\pgfqpoint{3.938556in}{5.950746in}}%
\pgfpathlineto{\pgfqpoint{3.939412in}{5.956600in}}%
\pgfpathlineto{\pgfqpoint{3.941037in}{5.988066in}}%
\pgfpathlineto{\pgfqpoint{3.943604in}{6.088728in}}%
\pgfpathlineto{\pgfqpoint{3.947626in}{6.347779in}}%
\pgfpathlineto{\pgfqpoint{3.960802in}{7.268052in}}%
\pgfpathlineto{\pgfqpoint{3.963540in}{7.321986in}}%
\pgfpathlineto{\pgfqpoint{3.964139in}{7.323908in}}%
\pgfpathlineto{\pgfqpoint{3.964653in}{7.322693in}}%
\pgfpathlineto{\pgfqpoint{3.965765in}{7.311038in}}%
\pgfpathlineto{\pgfqpoint{3.967733in}{7.260976in}}%
\pgfpathlineto{\pgfqpoint{3.970813in}{7.114290in}}%
\pgfpathlineto{\pgfqpoint{3.975861in}{6.747530in}}%
\pgfpathlineto{\pgfqpoint{3.985102in}{6.084761in}}%
\pgfpathlineto{\pgfqpoint{3.988525in}{5.974627in}}%
\pgfpathlineto{\pgfqpoint{3.990236in}{5.960862in}}%
\pgfpathlineto{\pgfqpoint{3.990407in}{5.961061in}}%
\pgfpathlineto{\pgfqpoint{3.991177in}{5.965503in}}%
\pgfpathlineto{\pgfqpoint{3.992717in}{5.991580in}}%
\pgfpathlineto{\pgfqpoint{3.995198in}{6.079302in}}%
\pgfpathlineto{\pgfqpoint{3.999049in}{6.308315in}}%
\pgfpathlineto{\pgfqpoint{4.013765in}{7.279345in}}%
\pgfpathlineto{\pgfqpoint{4.016161in}{7.313029in}}%
\pgfpathlineto{\pgfqpoint{4.016503in}{7.313384in}}%
\pgfpathlineto{\pgfqpoint{4.016931in}{7.312250in}}%
\pgfpathlineto{\pgfqpoint{4.018044in}{7.301133in}}%
\pgfpathlineto{\pgfqpoint{4.020012in}{7.253299in}}%
\pgfpathlineto{\pgfqpoint{4.023006in}{7.117708in}}%
\pgfpathlineto{\pgfqpoint{4.027969in}{6.773268in}}%
\pgfpathlineto{\pgfqpoint{4.037809in}{6.088250in}}%
\pgfpathlineto{\pgfqpoint{4.041231in}{5.984015in}}%
\pgfpathlineto{\pgfqpoint{4.042857in}{5.971464in}}%
\pgfpathlineto{\pgfqpoint{4.043113in}{5.971743in}}%
\pgfpathlineto{\pgfqpoint{4.043969in}{5.977123in}}%
\pgfpathlineto{\pgfqpoint{4.045595in}{6.005946in}}%
\pgfpathlineto{\pgfqpoint{4.048162in}{6.098233in}}%
\pgfpathlineto{\pgfqpoint{4.052183in}{6.337081in}}%
\pgfpathlineto{\pgfqpoint{4.066386in}{7.259130in}}%
\pgfpathlineto{\pgfqpoint{4.069039in}{7.301794in}}%
\pgfpathlineto{\pgfqpoint{4.069552in}{7.302678in}}%
\pgfpathlineto{\pgfqpoint{4.069980in}{7.301568in}}%
\pgfpathlineto{\pgfqpoint{4.071092in}{7.290867in}}%
\pgfpathlineto{\pgfqpoint{4.073060in}{7.244960in}}%
\pgfpathlineto{\pgfqpoint{4.076055in}{7.114807in}}%
\pgfpathlineto{\pgfqpoint{4.080932in}{6.789409in}}%
\pgfpathlineto{\pgfqpoint{4.091199in}{6.094921in}}%
\pgfpathlineto{\pgfqpoint{4.094622in}{5.994533in}}%
\pgfpathlineto{\pgfqpoint{4.096248in}{5.982264in}}%
\pgfpathlineto{\pgfqpoint{4.096504in}{5.982492in}}%
\pgfpathlineto{\pgfqpoint{4.097274in}{5.986718in}}%
\pgfpathlineto{\pgfqpoint{4.098815in}{6.010922in}}%
\pgfpathlineto{\pgfqpoint{4.101296in}{6.091900in}}%
\pgfpathlineto{\pgfqpoint{4.105146in}{6.303727in}}%
\pgfpathlineto{\pgfqpoint{4.120804in}{7.265759in}}%
\pgfpathlineto{\pgfqpoint{4.123114in}{7.291736in}}%
\pgfpathlineto{\pgfqpoint{4.123371in}{7.291742in}}%
\pgfpathlineto{\pgfqpoint{4.123713in}{7.290849in}}%
\pgfpathlineto{\pgfqpoint{4.124826in}{7.280869in}}%
\pgfpathlineto{\pgfqpoint{4.126793in}{7.237319in}}%
\pgfpathlineto{\pgfqpoint{4.129788in}{7.112968in}}%
\pgfpathlineto{\pgfqpoint{4.134580in}{6.806136in}}%
\pgfpathlineto{\pgfqpoint{4.145361in}{6.098994in}}%
\pgfpathlineto{\pgfqpoint{4.148783in}{6.004203in}}%
\pgfpathlineto{\pgfqpoint{4.150409in}{5.993240in}}%
\pgfpathlineto{\pgfqpoint{4.150665in}{5.993586in}}%
\pgfpathlineto{\pgfqpoint{4.151521in}{5.998829in}}%
\pgfpathlineto{\pgfqpoint{4.153147in}{6.025878in}}%
\pgfpathlineto{\pgfqpoint{4.155714in}{6.111658in}}%
\pgfpathlineto{\pgfqpoint{4.159735in}{6.333647in}}%
\pgfpathlineto{\pgfqpoint{4.174708in}{7.243892in}}%
\pgfpathlineto{\pgfqpoint{4.177275in}{7.280074in}}%
\pgfpathlineto{\pgfqpoint{4.177703in}{7.280748in}}%
\pgfpathlineto{\pgfqpoint{4.178217in}{7.279518in}}%
\pgfpathlineto{\pgfqpoint{4.179329in}{7.269259in}}%
\pgfpathlineto{\pgfqpoint{4.181297in}{7.226275in}}%
\pgfpathlineto{\pgfqpoint{4.184377in}{7.100827in}}%
\pgfpathlineto{\pgfqpoint{4.189254in}{6.795475in}}%
\pgfpathlineto{\pgfqpoint{4.199949in}{6.112246in}}%
\pgfpathlineto{\pgfqpoint{4.203457in}{6.015843in}}%
\pgfpathlineto{\pgfqpoint{4.205169in}{6.004385in}}%
\pgfpathlineto{\pgfqpoint{4.205340in}{6.004570in}}%
\pgfpathlineto{\pgfqpoint{4.206110in}{6.008397in}}%
\pgfpathlineto{\pgfqpoint{4.207650in}{6.030589in}}%
\pgfpathlineto{\pgfqpoint{4.210131in}{6.105190in}}%
\pgfpathlineto{\pgfqpoint{4.213896in}{6.296219in}}%
\pgfpathlineto{\pgfqpoint{4.222281in}{6.875189in}}%
\pgfpathlineto{\pgfqpoint{4.227928in}{7.178547in}}%
\pgfpathlineto{\pgfqpoint{4.231265in}{7.261004in}}%
\pgfpathlineto{\pgfqpoint{4.232720in}{7.269535in}}%
\pgfpathlineto{\pgfqpoint{4.233062in}{7.269048in}}%
\pgfpathlineto{\pgfqpoint{4.234003in}{7.262817in}}%
\pgfpathlineto{\pgfqpoint{4.235714in}{7.233404in}}%
\pgfpathlineto{\pgfqpoint{4.238452in}{7.140810in}}%
\pgfpathlineto{\pgfqpoint{4.242645in}{6.911082in}}%
\pgfpathlineto{\pgfqpoint{4.256848in}{6.066557in}}%
\pgfpathlineto{\pgfqpoint{4.259672in}{6.018036in}}%
\pgfpathlineto{\pgfqpoint{4.260442in}{6.015665in}}%
\pgfpathlineto{\pgfqpoint{4.260955in}{6.016711in}}%
\pgfpathlineto{\pgfqpoint{4.262068in}{6.026151in}}%
\pgfpathlineto{\pgfqpoint{4.264036in}{6.066350in}}%
\pgfpathlineto{\pgfqpoint{4.267030in}{6.180422in}}%
\pgfpathlineto{\pgfqpoint{4.271822in}{6.462759in}}%
\pgfpathlineto{\pgfqpoint{4.283287in}{7.162279in}}%
\pgfpathlineto{\pgfqpoint{4.286710in}{7.248330in}}%
\pgfpathlineto{\pgfqpoint{4.288335in}{7.258173in}}%
\pgfpathlineto{\pgfqpoint{4.288592in}{7.257838in}}%
\pgfpathlineto{\pgfqpoint{4.289448in}{7.253003in}}%
\pgfpathlineto{\pgfqpoint{4.291073in}{7.228265in}}%
\pgfpathlineto{\pgfqpoint{4.293640in}{7.149901in}}%
\pgfpathlineto{\pgfqpoint{4.297576in}{6.951512in}}%
\pgfpathlineto{\pgfqpoint{4.313576in}{6.054778in}}%
\pgfpathlineto{\pgfqpoint{4.315972in}{6.027419in}}%
\pgfpathlineto{\pgfqpoint{4.316314in}{6.027092in}}%
\pgfpathlineto{\pgfqpoint{4.316742in}{6.027951in}}%
\pgfpathlineto{\pgfqpoint{4.317854in}{6.036748in}}%
\pgfpathlineto{\pgfqpoint{4.319822in}{6.075032in}}%
\pgfpathlineto{\pgfqpoint{4.322817in}{6.184573in}}%
\pgfpathlineto{\pgfqpoint{4.327523in}{6.451819in}}%
\pgfpathlineto{\pgfqpoint{4.339416in}{7.156908in}}%
\pgfpathlineto{\pgfqpoint{4.342753in}{7.237139in}}%
\pgfpathlineto{\pgfqpoint{4.344379in}{7.246698in}}%
\pgfpathlineto{\pgfqpoint{4.344635in}{7.246382in}}%
\pgfpathlineto{\pgfqpoint{4.345491in}{7.241731in}}%
\pgfpathlineto{\pgfqpoint{4.347117in}{7.217857in}}%
\pgfpathlineto{\pgfqpoint{4.349684in}{7.142113in}}%
\pgfpathlineto{\pgfqpoint{4.353620in}{6.950001in}}%
\pgfpathlineto{\pgfqpoint{4.369962in}{6.062784in}}%
\pgfpathlineto{\pgfqpoint{4.372358in}{6.038705in}}%
\pgfpathlineto{\pgfqpoint{4.372700in}{6.038736in}}%
\pgfpathlineto{\pgfqpoint{4.372957in}{6.039332in}}%
\pgfpathlineto{\pgfqpoint{4.373983in}{6.046606in}}%
\pgfpathlineto{\pgfqpoint{4.375866in}{6.079878in}}%
\pgfpathlineto{\pgfqpoint{4.378775in}{6.178450in}}%
\pgfpathlineto{\pgfqpoint{4.383310in}{6.421579in}}%
\pgfpathlineto{\pgfqpoint{4.396315in}{7.163559in}}%
\pgfpathlineto{\pgfqpoint{4.399481in}{7.229004in}}%
\pgfpathlineto{\pgfqpoint{4.400764in}{7.235129in}}%
\pgfpathlineto{\pgfqpoint{4.401192in}{7.234487in}}%
\pgfpathlineto{\pgfqpoint{4.402219in}{7.227485in}}%
\pgfpathlineto{\pgfqpoint{4.404101in}{7.195001in}}%
\pgfpathlineto{\pgfqpoint{4.407011in}{7.098282in}}%
\pgfpathlineto{\pgfqpoint{4.411460in}{6.864221in}}%
\pgfpathlineto{\pgfqpoint{4.424808in}{6.115831in}}%
\pgfpathlineto{\pgfqpoint{4.427888in}{6.055573in}}%
\pgfpathlineto{\pgfqpoint{4.429086in}{6.050227in}}%
\pgfpathlineto{\pgfqpoint{4.429513in}{6.050831in}}%
\pgfpathlineto{\pgfqpoint{4.430540in}{6.057661in}}%
\pgfpathlineto{\pgfqpoint{4.432423in}{6.089555in}}%
\pgfpathlineto{\pgfqpoint{4.435332in}{6.184735in}}%
\pgfpathlineto{\pgfqpoint{4.439781in}{6.415444in}}%
\pgfpathlineto{\pgfqpoint{4.453214in}{7.159201in}}%
\pgfpathlineto{\pgfqpoint{4.456294in}{7.218330in}}%
\pgfpathlineto{\pgfqpoint{4.457492in}{7.223480in}}%
\pgfpathlineto{\pgfqpoint{4.457920in}{7.222841in}}%
\pgfpathlineto{\pgfqpoint{4.458947in}{7.215999in}}%
\pgfpathlineto{\pgfqpoint{4.460829in}{7.184356in}}%
\pgfpathlineto{\pgfqpoint{4.463738in}{7.090204in}}%
\pgfpathlineto{\pgfqpoint{4.468188in}{6.862227in}}%
\pgfpathlineto{\pgfqpoint{4.481621in}{6.126377in}}%
\pgfpathlineto{\pgfqpoint{4.484787in}{6.066571in}}%
\pgfpathlineto{\pgfqpoint{4.485985in}{6.061911in}}%
\pgfpathlineto{\pgfqpoint{4.486327in}{6.062434in}}%
\pgfpathlineto{\pgfqpoint{4.487354in}{6.068932in}}%
\pgfpathlineto{\pgfqpoint{4.489150in}{6.097798in}}%
\pgfpathlineto{\pgfqpoint{4.491974in}{6.185230in}}%
\pgfpathlineto{\pgfqpoint{4.496338in}{6.401579in}}%
\pgfpathlineto{\pgfqpoint{4.510456in}{7.157980in}}%
\pgfpathlineto{\pgfqpoint{4.513450in}{7.208521in}}%
\pgfpathlineto{\pgfqpoint{4.514477in}{7.211767in}}%
\pgfpathlineto{\pgfqpoint{4.514905in}{7.210958in}}%
\pgfpathlineto{\pgfqpoint{4.516017in}{7.202924in}}%
\pgfpathlineto{\pgfqpoint{4.517985in}{7.168170in}}%
\pgfpathlineto{\pgfqpoint{4.520980in}{7.068786in}}%
\pgfpathlineto{\pgfqpoint{4.525686in}{6.825486in}}%
\pgfpathlineto{\pgfqpoint{4.538092in}{6.152059in}}%
\pgfpathlineto{\pgfqpoint{4.541429in}{6.081282in}}%
\pgfpathlineto{\pgfqpoint{4.542969in}{6.073628in}}%
\pgfpathlineto{\pgfqpoint{4.543226in}{6.073934in}}%
\pgfpathlineto{\pgfqpoint{4.544167in}{6.078923in}}%
\pgfpathlineto{\pgfqpoint{4.545878in}{6.103333in}}%
\pgfpathlineto{\pgfqpoint{4.548531in}{6.178115in}}%
\pgfpathlineto{\pgfqpoint{4.552638in}{6.367374in}}%
\pgfpathlineto{\pgfqpoint{4.568381in}{7.168493in}}%
\pgfpathlineto{\pgfqpoint{4.570948in}{7.199225in}}%
\pgfpathlineto{\pgfqpoint{4.571462in}{7.200051in}}%
\pgfpathlineto{\pgfqpoint{4.571975in}{7.199086in}}%
\pgfpathlineto{\pgfqpoint{4.573087in}{7.190877in}}%
\pgfpathlineto{\pgfqpoint{4.575055in}{7.156290in}}%
\pgfpathlineto{\pgfqpoint{4.578050in}{7.058284in}}%
\pgfpathlineto{\pgfqpoint{4.582756in}{6.819434in}}%
\pgfpathlineto{\pgfqpoint{4.595077in}{6.163762in}}%
\pgfpathlineto{\pgfqpoint{4.598414in}{6.093379in}}%
\pgfpathlineto{\pgfqpoint{4.599954in}{6.085351in}}%
\pgfpathlineto{\pgfqpoint{4.600211in}{6.085562in}}%
\pgfpathlineto{\pgfqpoint{4.601066in}{6.089462in}}%
\pgfpathlineto{\pgfqpoint{4.602692in}{6.110259in}}%
\pgfpathlineto{\pgfqpoint{4.605259in}{6.177062in}}%
\pgfpathlineto{\pgfqpoint{4.609195in}{6.347840in}}%
\pgfpathlineto{\pgfqpoint{4.618692in}{6.898876in}}%
\pgfpathlineto{\pgfqpoint{4.623997in}{7.123009in}}%
\pgfpathlineto{\pgfqpoint{4.627163in}{7.182670in}}%
\pgfpathlineto{\pgfqpoint{4.628446in}{7.188324in}}%
\pgfpathlineto{\pgfqpoint{4.628874in}{7.187772in}}%
\pgfpathlineto{\pgfqpoint{4.629901in}{7.181482in}}%
\pgfpathlineto{\pgfqpoint{4.631698in}{7.153908in}}%
\pgfpathlineto{\pgfqpoint{4.634521in}{7.070734in}}%
\pgfpathlineto{\pgfqpoint{4.638885in}{6.865267in}}%
\pgfpathlineto{\pgfqpoint{4.653003in}{6.147747in}}%
\pgfpathlineto{\pgfqpoint{4.655997in}{6.100052in}}%
\pgfpathlineto{\pgfqpoint{4.656938in}{6.097068in}}%
\pgfpathlineto{\pgfqpoint{4.657452in}{6.097907in}}%
\pgfpathlineto{\pgfqpoint{4.658564in}{6.105677in}}%
\pgfpathlineto{\pgfqpoint{4.660532in}{6.138948in}}%
\pgfpathlineto{\pgfqpoint{4.663527in}{6.233766in}}%
\pgfpathlineto{\pgfqpoint{4.668233in}{6.465459in}}%
\pgfpathlineto{\pgfqpoint{4.680554in}{7.101578in}}%
\pgfpathlineto{\pgfqpoint{4.683891in}{7.169241in}}%
\pgfpathlineto{\pgfqpoint{4.685431in}{7.176619in}}%
\pgfpathlineto{\pgfqpoint{4.685687in}{7.176339in}}%
\pgfpathlineto{\pgfqpoint{4.686629in}{7.171622in}}%
\pgfpathlineto{\pgfqpoint{4.688340in}{7.148403in}}%
\pgfpathlineto{\pgfqpoint{4.690992in}{7.077136in}}%
\pgfpathlineto{\pgfqpoint{4.695099in}{6.896678in}}%
\pgfpathlineto{\pgfqpoint{4.710757in}{6.138401in}}%
\pgfpathlineto{\pgfqpoint{4.713324in}{6.109449in}}%
\pgfpathlineto{\pgfqpoint{4.713837in}{6.108755in}}%
\pgfpathlineto{\pgfqpoint{4.714351in}{6.109778in}}%
\pgfpathlineto{\pgfqpoint{4.715549in}{6.118805in}}%
\pgfpathlineto{\pgfqpoint{4.717602in}{6.155386in}}%
\pgfpathlineto{\pgfqpoint{4.720768in}{6.259462in}}%
\pgfpathlineto{\pgfqpoint{4.725731in}{6.508442in}}%
\pgfpathlineto{\pgfqpoint{4.736768in}{7.074804in}}%
\pgfpathlineto{\pgfqpoint{4.740276in}{7.153973in}}%
\pgfpathlineto{\pgfqpoint{4.742159in}{7.164969in}}%
\pgfpathlineto{\pgfqpoint{4.742244in}{7.164925in}}%
\pgfpathlineto{\pgfqpoint{4.742929in}{7.162873in}}%
\pgfpathlineto{\pgfqpoint{4.744298in}{7.149746in}}%
\pgfpathlineto{\pgfqpoint{4.746608in}{7.101260in}}%
\pgfpathlineto{\pgfqpoint{4.750116in}{6.971699in}}%
\pgfpathlineto{\pgfqpoint{4.756105in}{6.650806in}}%
\pgfpathlineto{\pgfqpoint{4.764490in}{6.228322in}}%
\pgfpathlineto{\pgfqpoint{4.768255in}{6.134896in}}%
\pgfpathlineto{\pgfqpoint{4.770394in}{6.120364in}}%
\pgfpathlineto{\pgfqpoint{4.770993in}{6.121555in}}%
\pgfpathlineto{\pgfqpoint{4.772191in}{6.130825in}}%
\pgfpathlineto{\pgfqpoint{4.774245in}{6.167546in}}%
\pgfpathlineto{\pgfqpoint{4.777410in}{6.271107in}}%
\pgfpathlineto{\pgfqpoint{4.782459in}{6.522195in}}%
\pgfpathlineto{\pgfqpoint{4.793154in}{7.062443in}}%
\pgfpathlineto{\pgfqpoint{4.796748in}{7.142877in}}%
\pgfpathlineto{\pgfqpoint{4.798544in}{7.153395in}}%
\pgfpathlineto{\pgfqpoint{4.798715in}{7.153322in}}%
\pgfpathlineto{\pgfqpoint{4.799486in}{7.150672in}}%
\pgfpathlineto{\pgfqpoint{4.800940in}{7.135407in}}%
\pgfpathlineto{\pgfqpoint{4.803336in}{7.082191in}}%
\pgfpathlineto{\pgfqpoint{4.806929in}{6.945303in}}%
\pgfpathlineto{\pgfqpoint{4.813518in}{6.590206in}}%
\pgfpathlineto{\pgfqpoint{4.820876in}{6.234493in}}%
\pgfpathlineto{\pgfqpoint{4.824555in}{6.145769in}}%
\pgfpathlineto{\pgfqpoint{4.826694in}{6.131902in}}%
\pgfpathlineto{\pgfqpoint{4.827293in}{6.133233in}}%
\pgfpathlineto{\pgfqpoint{4.828577in}{6.143742in}}%
\pgfpathlineto{\pgfqpoint{4.830716in}{6.183801in}}%
\pgfpathlineto{\pgfqpoint{4.833967in}{6.293336in}}%
\pgfpathlineto{\pgfqpoint{4.839272in}{6.560226in}}%
\pgfpathlineto{\pgfqpoint{4.849026in}{7.045544in}}%
\pgfpathlineto{\pgfqpoint{4.852620in}{7.129220in}}%
\pgfpathlineto{\pgfqpoint{4.854673in}{7.141909in}}%
\pgfpathlineto{\pgfqpoint{4.855272in}{7.140576in}}%
\pgfpathlineto{\pgfqpoint{4.856556in}{7.130087in}}%
\pgfpathlineto{\pgfqpoint{4.858695in}{7.090146in}}%
\pgfpathlineto{\pgfqpoint{4.861946in}{6.981006in}}%
\pgfpathlineto{\pgfqpoint{4.867337in}{6.710707in}}%
\pgfpathlineto{\pgfqpoint{4.876834in}{6.241170in}}%
\pgfpathlineto{\pgfqpoint{4.880513in}{6.155576in}}%
\pgfpathlineto{\pgfqpoint{4.882481in}{6.143316in}}%
\pgfpathlineto{\pgfqpoint{4.882567in}{6.143338in}}%
\pgfpathlineto{\pgfqpoint{4.883251in}{6.145183in}}%
\pgfpathlineto{\pgfqpoint{4.884620in}{6.157715in}}%
\pgfpathlineto{\pgfqpoint{4.886845in}{6.202373in}}%
\pgfpathlineto{\pgfqpoint{4.890267in}{6.323146in}}%
\pgfpathlineto{\pgfqpoint{4.896171in}{6.626984in}}%
\pgfpathlineto{\pgfqpoint{4.904471in}{7.029625in}}%
\pgfpathlineto{\pgfqpoint{4.908150in}{7.117074in}}%
\pgfpathlineto{\pgfqpoint{4.910203in}{7.130557in}}%
\pgfpathlineto{\pgfqpoint{4.910289in}{7.130542in}}%
\pgfpathlineto{\pgfqpoint{4.910973in}{7.128758in}}%
\pgfpathlineto{\pgfqpoint{4.912342in}{7.116365in}}%
\pgfpathlineto{\pgfqpoint{4.914567in}{7.071980in}}%
\pgfpathlineto{\pgfqpoint{4.917990in}{6.951762in}}%
\pgfpathlineto{\pgfqpoint{4.923893in}{6.649515in}}%
\pgfpathlineto{\pgfqpoint{4.932107in}{6.254133in}}%
\pgfpathlineto{\pgfqpoint{4.935787in}{6.167607in}}%
\pgfpathlineto{\pgfqpoint{4.937840in}{6.154610in}}%
\pgfpathlineto{\pgfqpoint{4.937926in}{6.154645in}}%
\pgfpathlineto{\pgfqpoint{4.938610in}{6.156583in}}%
\pgfpathlineto{\pgfqpoint{4.939979in}{6.169267in}}%
\pgfpathlineto{\pgfqpoint{4.942289in}{6.216344in}}%
\pgfpathlineto{\pgfqpoint{4.945797in}{6.342015in}}%
\pgfpathlineto{\pgfqpoint{4.952043in}{6.664312in}}%
\pgfpathlineto{\pgfqpoint{4.959658in}{7.023084in}}%
\pgfpathlineto{\pgfqpoint{4.963338in}{7.107524in}}%
\pgfpathlineto{\pgfqpoint{4.965306in}{7.119334in}}%
\pgfpathlineto{\pgfqpoint{4.965391in}{7.119295in}}%
\pgfpathlineto{\pgfqpoint{4.966076in}{7.117320in}}%
\pgfpathlineto{\pgfqpoint{4.967445in}{7.104572in}}%
\pgfpathlineto{\pgfqpoint{4.969755in}{7.057425in}}%
\pgfpathlineto{\pgfqpoint{4.973263in}{6.931826in}}%
\pgfpathlineto{\pgfqpoint{4.979594in}{6.606100in}}%
\pgfpathlineto{\pgfqpoint{4.987038in}{6.259469in}}%
\pgfpathlineto{\pgfqpoint{4.990632in}{6.177689in}}%
\pgfpathlineto{\pgfqpoint{4.992600in}{6.165755in}}%
\pgfpathlineto{\pgfqpoint{4.992686in}{6.165789in}}%
\pgfpathlineto{\pgfqpoint{4.993370in}{6.167719in}}%
\pgfpathlineto{\pgfqpoint{4.994739in}{6.180376in}}%
\pgfpathlineto{\pgfqpoint{4.997049in}{6.227361in}}%
\pgfpathlineto{\pgfqpoint{5.000557in}{6.352635in}}%
\pgfpathlineto{\pgfqpoint{5.006889in}{6.676976in}}%
\pgfpathlineto{\pgfqpoint{5.014247in}{7.016643in}}%
\pgfpathlineto{\pgfqpoint{5.017841in}{7.097100in}}%
\pgfpathlineto{\pgfqpoint{5.019723in}{7.108272in}}%
\pgfpathlineto{\pgfqpoint{5.019894in}{7.108182in}}%
\pgfpathlineto{\pgfqpoint{5.020664in}{7.105493in}}%
\pgfpathlineto{\pgfqpoint{5.022119in}{7.090310in}}%
\pgfpathlineto{\pgfqpoint{5.024515in}{7.037740in}}%
\pgfpathlineto{\pgfqpoint{5.028194in}{6.899682in}}%
\pgfpathlineto{\pgfqpoint{5.035809in}{6.504076in}}%
\pgfpathlineto{\pgfqpoint{5.041798in}{6.250729in}}%
\pgfpathlineto{\pgfqpoint{5.045135in}{6.184488in}}%
\pgfpathlineto{\pgfqpoint{5.046761in}{6.176745in}}%
\pgfpathlineto{\pgfqpoint{5.047018in}{6.177046in}}%
\pgfpathlineto{\pgfqpoint{5.047959in}{6.181702in}}%
\pgfpathlineto{\pgfqpoint{5.049670in}{6.204253in}}%
\pgfpathlineto{\pgfqpoint{5.052408in}{6.275761in}}%
\pgfpathlineto{\pgfqpoint{5.056601in}{6.453036in}}%
\pgfpathlineto{\pgfqpoint{5.069863in}{7.055960in}}%
\pgfpathlineto{\pgfqpoint{5.072686in}{7.095242in}}%
\pgfpathlineto{\pgfqpoint{5.073542in}{7.097375in}}%
\pgfpathlineto{\pgfqpoint{5.074055in}{7.096433in}}%
\pgfpathlineto{\pgfqpoint{5.075253in}{7.087781in}}%
\pgfpathlineto{\pgfqpoint{5.077307in}{7.052493in}}%
\pgfpathlineto{\pgfqpoint{5.080473in}{6.952369in}}%
\pgfpathlineto{\pgfqpoint{5.085692in}{6.703128in}}%
\pgfpathlineto{\pgfqpoint{5.094847in}{6.273289in}}%
\pgfpathlineto{\pgfqpoint{5.098355in}{6.197526in}}%
\pgfpathlineto{\pgfqpoint{5.100152in}{6.187539in}}%
\pgfpathlineto{\pgfqpoint{5.100323in}{6.187652in}}%
\pgfpathlineto{\pgfqpoint{5.101093in}{6.190454in}}%
\pgfpathlineto{\pgfqpoint{5.102633in}{6.207201in}}%
\pgfpathlineto{\pgfqpoint{5.105115in}{6.263921in}}%
\pgfpathlineto{\pgfqpoint{5.108879in}{6.409001in}}%
\pgfpathlineto{\pgfqpoint{5.124281in}{7.070071in}}%
\pgfpathlineto{\pgfqpoint{5.126505in}{7.086655in}}%
\pgfpathlineto{\pgfqpoint{5.127104in}{7.085775in}}%
\pgfpathlineto{\pgfqpoint{5.128216in}{7.078118in}}%
\pgfpathlineto{\pgfqpoint{5.130184in}{7.045875in}}%
\pgfpathlineto{\pgfqpoint{5.133265in}{6.951824in}}%
\pgfpathlineto{\pgfqpoint{5.138313in}{6.716032in}}%
\pgfpathlineto{\pgfqpoint{5.147810in}{6.275825in}}%
\pgfpathlineto{\pgfqpoint{5.151233in}{6.206180in}}%
\pgfpathlineto{\pgfqpoint{5.152858in}{6.198154in}}%
\pgfpathlineto{\pgfqpoint{5.153115in}{6.198422in}}%
\pgfpathlineto{\pgfqpoint{5.154056in}{6.202990in}}%
\pgfpathlineto{\pgfqpoint{5.155768in}{6.225494in}}%
\pgfpathlineto{\pgfqpoint{5.158506in}{6.297102in}}%
\pgfpathlineto{\pgfqpoint{5.162784in}{6.478208in}}%
\pgfpathlineto{\pgfqpoint{5.175019in}{7.029429in}}%
\pgfpathlineto{\pgfqpoint{5.178014in}{7.073775in}}%
\pgfpathlineto{\pgfqpoint{5.178869in}{7.076156in}}%
\pgfpathlineto{\pgfqpoint{5.179383in}{7.075340in}}%
\pgfpathlineto{\pgfqpoint{5.180495in}{7.067818in}}%
\pgfpathlineto{\pgfqpoint{5.182463in}{7.035701in}}%
\pgfpathlineto{\pgfqpoint{5.185543in}{6.941681in}}%
\pgfpathlineto{\pgfqpoint{5.190591in}{6.706476in}}%
\pgfpathlineto{\pgfqpoint{5.199747in}{6.285904in}}%
\pgfpathlineto{\pgfqpoint{5.203169in}{6.216385in}}%
\pgfpathlineto{\pgfqpoint{5.204709in}{6.208564in}}%
\pgfpathlineto{\pgfqpoint{5.205052in}{6.208887in}}%
\pgfpathlineto{\pgfqpoint{5.205993in}{6.213641in}}%
\pgfpathlineto{\pgfqpoint{5.207704in}{6.236566in}}%
\pgfpathlineto{\pgfqpoint{5.210442in}{6.308948in}}%
\pgfpathlineto{\pgfqpoint{5.214720in}{6.490766in}}%
\pgfpathlineto{\pgfqpoint{5.226442in}{7.016715in}}%
\pgfpathlineto{\pgfqpoint{5.229437in}{7.062946in}}%
\pgfpathlineto{\pgfqpoint{5.230378in}{7.065848in}}%
\pgfpathlineto{\pgfqpoint{5.230891in}{7.065027in}}%
\pgfpathlineto{\pgfqpoint{5.232004in}{7.057452in}}%
\pgfpathlineto{\pgfqpoint{5.233972in}{7.025118in}}%
\pgfpathlineto{\pgfqpoint{5.237052in}{6.930603in}}%
\pgfpathlineto{\pgfqpoint{5.242186in}{6.690981in}}%
\pgfpathlineto{\pgfqpoint{5.250913in}{6.294704in}}%
\pgfpathlineto{\pgfqpoint{5.254336in}{6.226011in}}%
\pgfpathlineto{\pgfqpoint{5.255876in}{6.218771in}}%
\pgfpathlineto{\pgfqpoint{5.256132in}{6.219053in}}%
\pgfpathlineto{\pgfqpoint{5.257074in}{6.223723in}}%
\pgfpathlineto{\pgfqpoint{5.258785in}{6.246609in}}%
\pgfpathlineto{\pgfqpoint{5.261523in}{6.319160in}}%
\pgfpathlineto{\pgfqpoint{5.265886in}{6.505362in}}%
\pgfpathlineto{\pgfqpoint{5.277095in}{7.005360in}}%
\pgfpathlineto{\pgfqpoint{5.280090in}{7.052606in}}%
\pgfpathlineto{\pgfqpoint{5.281117in}{7.055744in}}%
\pgfpathlineto{\pgfqpoint{5.281544in}{7.055032in}}%
\pgfpathlineto{\pgfqpoint{5.282657in}{7.047638in}}%
\pgfpathlineto{\pgfqpoint{5.284625in}{7.015454in}}%
\pgfpathlineto{\pgfqpoint{5.287705in}{6.920877in}}%
\pgfpathlineto{\pgfqpoint{5.292924in}{6.677168in}}%
\pgfpathlineto{\pgfqpoint{5.301224in}{6.304079in}}%
\pgfpathlineto{\pgfqpoint{5.304561in}{6.236502in}}%
\pgfpathlineto{\pgfqpoint{5.306101in}{6.228756in}}%
\pgfpathlineto{\pgfqpoint{5.306358in}{6.228967in}}%
\pgfpathlineto{\pgfqpoint{5.307213in}{6.232771in}}%
\pgfpathlineto{\pgfqpoint{5.308839in}{6.252948in}}%
\pgfpathlineto{\pgfqpoint{5.311406in}{6.317194in}}%
\pgfpathlineto{\pgfqpoint{5.315513in}{6.485474in}}%
\pgfpathlineto{\pgfqpoint{5.327406in}{7.005999in}}%
\pgfpathlineto{\pgfqpoint{5.330229in}{7.044302in}}%
\pgfpathlineto{\pgfqpoint{5.330914in}{7.045869in}}%
\pgfpathlineto{\pgfqpoint{5.331427in}{7.045026in}}%
\pgfpathlineto{\pgfqpoint{5.332540in}{7.037295in}}%
\pgfpathlineto{\pgfqpoint{5.334508in}{7.004363in}}%
\pgfpathlineto{\pgfqpoint{5.337588in}{6.908495in}}%
\pgfpathlineto{\pgfqpoint{5.342893in}{6.659679in}}%
\pgfpathlineto{\pgfqpoint{5.350764in}{6.310980in}}%
\pgfpathlineto{\pgfqpoint{5.354101in}{6.245188in}}%
\pgfpathlineto{\pgfqpoint{5.355556in}{6.238531in}}%
\pgfpathlineto{\pgfqpoint{5.355898in}{6.238990in}}%
\pgfpathlineto{\pgfqpoint{5.356925in}{6.244986in}}%
\pgfpathlineto{\pgfqpoint{5.358722in}{6.271801in}}%
\pgfpathlineto{\pgfqpoint{5.361545in}{6.352344in}}%
\pgfpathlineto{\pgfqpoint{5.366251in}{6.561944in}}%
\pgfpathlineto{\pgfqpoint{5.375492in}{6.973110in}}%
\pgfpathlineto{\pgfqpoint{5.378658in}{7.030991in}}%
\pgfpathlineto{\pgfqpoint{5.379941in}{7.036204in}}%
\pgfpathlineto{\pgfqpoint{5.380283in}{7.035751in}}%
\pgfpathlineto{\pgfqpoint{5.381310in}{7.029745in}}%
\pgfpathlineto{\pgfqpoint{5.383107in}{7.002828in}}%
\pgfpathlineto{\pgfqpoint{5.385931in}{6.921979in}}%
\pgfpathlineto{\pgfqpoint{5.390637in}{6.712088in}}%
\pgfpathlineto{\pgfqpoint{5.399706in}{6.310379in}}%
\pgfpathlineto{\pgfqpoint{5.402872in}{6.252991in}}%
\pgfpathlineto{\pgfqpoint{5.404070in}{6.248080in}}%
\pgfpathlineto{\pgfqpoint{5.404498in}{6.248644in}}%
\pgfpathlineto{\pgfqpoint{5.405524in}{6.254961in}}%
\pgfpathlineto{\pgfqpoint{5.407407in}{6.284308in}}%
\pgfpathlineto{\pgfqpoint{5.410316in}{6.370634in}}%
\pgfpathlineto{\pgfqpoint{5.415279in}{6.596812in}}%
\pgfpathlineto{\pgfqpoint{5.423493in}{6.959281in}}%
\pgfpathlineto{\pgfqpoint{5.426744in}{7.021105in}}%
\pgfpathlineto{\pgfqpoint{5.428027in}{7.026765in}}%
\pgfpathlineto{\pgfqpoint{5.428455in}{7.026199in}}%
\pgfpathlineto{\pgfqpoint{5.429482in}{7.019845in}}%
\pgfpathlineto{\pgfqpoint{5.431364in}{6.990335in}}%
\pgfpathlineto{\pgfqpoint{5.434273in}{6.903607in}}%
\pgfpathlineto{\pgfqpoint{5.439322in}{6.672786in}}%
\pgfpathlineto{\pgfqpoint{5.447279in}{6.323987in}}%
\pgfpathlineto{\pgfqpoint{5.450530in}{6.262721in}}%
\pgfpathlineto{\pgfqpoint{5.451814in}{6.257411in}}%
\pgfpathlineto{\pgfqpoint{5.452156in}{6.257871in}}%
\pgfpathlineto{\pgfqpoint{5.453183in}{6.263982in}}%
\pgfpathlineto{\pgfqpoint{5.454979in}{6.291355in}}%
\pgfpathlineto{\pgfqpoint{5.457889in}{6.376497in}}%
\pgfpathlineto{\pgfqpoint{5.462851in}{6.601595in}}%
\pgfpathlineto{\pgfqpoint{5.470894in}{6.953134in}}%
\pgfpathlineto{\pgfqpoint{5.474060in}{7.012233in}}%
\pgfpathlineto{\pgfqpoint{5.475343in}{7.017545in}}%
\pgfpathlineto{\pgfqpoint{5.475686in}{7.017073in}}%
\pgfpathlineto{\pgfqpoint{5.476712in}{7.010899in}}%
\pgfpathlineto{\pgfqpoint{5.478509in}{6.983321in}}%
\pgfpathlineto{\pgfqpoint{5.481418in}{6.897698in}}%
\pgfpathlineto{\pgfqpoint{5.486381in}{6.672209in}}%
\pgfpathlineto{\pgfqpoint{5.494253in}{6.329954in}}%
\pgfpathlineto{\pgfqpoint{5.497418in}{6.271466in}}%
\pgfpathlineto{\pgfqpoint{5.498616in}{6.266516in}}%
\pgfpathlineto{\pgfqpoint{5.499044in}{6.267124in}}%
\pgfpathlineto{\pgfqpoint{5.500071in}{6.273670in}}%
\pgfpathlineto{\pgfqpoint{5.501953in}{6.303835in}}%
\pgfpathlineto{\pgfqpoint{5.504948in}{6.395259in}}%
\pgfpathlineto{\pgfqpoint{5.510253in}{6.641138in}}%
\pgfpathlineto{\pgfqpoint{5.517355in}{6.944895in}}%
\pgfpathlineto{\pgfqpoint{5.520520in}{7.003616in}}%
\pgfpathlineto{\pgfqpoint{5.521718in}{7.008551in}}%
\pgfpathlineto{\pgfqpoint{5.522146in}{7.007923in}}%
\pgfpathlineto{\pgfqpoint{5.523173in}{7.001295in}}%
\pgfpathlineto{\pgfqpoint{5.525055in}{6.970875in}}%
\pgfpathlineto{\pgfqpoint{5.528050in}{6.878900in}}%
\pgfpathlineto{\pgfqpoint{5.533440in}{6.628419in}}%
\pgfpathlineto{\pgfqpoint{5.540371in}{6.335687in}}%
\pgfpathlineto{\pgfqpoint{5.543537in}{6.279326in}}%
\pgfpathlineto{\pgfqpoint{5.544649in}{6.275411in}}%
\pgfpathlineto{\pgfqpoint{5.545077in}{6.276187in}}%
\pgfpathlineto{\pgfqpoint{5.546189in}{6.284107in}}%
\pgfpathlineto{\pgfqpoint{5.548157in}{6.318362in}}%
\pgfpathlineto{\pgfqpoint{5.551237in}{6.417709in}}%
\pgfpathlineto{\pgfqpoint{5.557227in}{6.701486in}}%
\pgfpathlineto{\pgfqpoint{5.563216in}{6.943910in}}%
\pgfpathlineto{\pgfqpoint{5.566296in}{6.996461in}}%
\pgfpathlineto{\pgfqpoint{5.567323in}{6.999764in}}%
\pgfpathlineto{\pgfqpoint{5.567751in}{6.998971in}}%
\pgfpathlineto{\pgfqpoint{5.568863in}{6.990966in}}%
\pgfpathlineto{\pgfqpoint{5.570831in}{6.956444in}}%
\pgfpathlineto{\pgfqpoint{5.573911in}{6.856541in}}%
\pgfpathlineto{\pgfqpoint{5.580072in}{6.564411in}}%
\pgfpathlineto{\pgfqpoint{5.585890in}{6.334610in}}%
\pgfpathlineto{\pgfqpoint{5.588885in}{6.286578in}}%
\pgfpathlineto{\pgfqpoint{5.589740in}{6.284072in}}%
\pgfpathlineto{\pgfqpoint{5.590254in}{6.285036in}}%
\pgfpathlineto{\pgfqpoint{5.591366in}{6.293437in}}%
\pgfpathlineto{\pgfqpoint{5.593334in}{6.328747in}}%
\pgfpathlineto{\pgfqpoint{5.596500in}{6.433330in}}%
\pgfpathlineto{\pgfqpoint{5.603259in}{6.754948in}}%
\pgfpathlineto{\pgfqpoint{5.608479in}{6.948781in}}%
\pgfpathlineto{\pgfqpoint{5.611302in}{6.989596in}}%
\pgfpathlineto{\pgfqpoint{5.611987in}{6.991213in}}%
\pgfpathlineto{\pgfqpoint{5.612500in}{6.990251in}}%
\pgfpathlineto{\pgfqpoint{5.613612in}{6.981813in}}%
\pgfpathlineto{\pgfqpoint{5.615580in}{6.946311in}}%
\pgfpathlineto{\pgfqpoint{5.618746in}{6.841257in}}%
\pgfpathlineto{\pgfqpoint{5.625848in}{6.504316in}}%
\pgfpathlineto{\pgfqpoint{5.630725in}{6.330444in}}%
\pgfpathlineto{\pgfqpoint{5.633463in}{6.293625in}}%
\pgfpathlineto{\pgfqpoint{5.634062in}{6.292531in}}%
\pgfpathlineto{\pgfqpoint{5.634575in}{6.293626in}}%
\pgfpathlineto{\pgfqpoint{5.635773in}{6.303426in}}%
\pgfpathlineto{\pgfqpoint{5.637826in}{6.342897in}}%
\pgfpathlineto{\pgfqpoint{5.641078in}{6.455400in}}%
\pgfpathlineto{\pgfqpoint{5.654426in}{6.975593in}}%
\pgfpathlineto{\pgfqpoint{5.655880in}{6.982864in}}%
\pgfpathlineto{\pgfqpoint{5.656222in}{6.982377in}}%
\pgfpathlineto{\pgfqpoint{5.657249in}{6.975897in}}%
\pgfpathlineto{\pgfqpoint{5.659046in}{6.946942in}}%
\pgfpathlineto{\pgfqpoint{5.661955in}{6.857752in}}%
\pgfpathlineto{\pgfqpoint{5.667431in}{6.604558in}}%
\pgfpathlineto{\pgfqpoint{5.673677in}{6.351675in}}%
\pgfpathlineto{\pgfqpoint{5.676586in}{6.303654in}}%
\pgfpathlineto{\pgfqpoint{5.677527in}{6.300773in}}%
\pgfpathlineto{\pgfqpoint{5.677955in}{6.301575in}}%
\pgfpathlineto{\pgfqpoint{5.679068in}{6.309806in}}%
\pgfpathlineto{\pgfqpoint{5.681036in}{6.345357in}}%
\pgfpathlineto{\pgfqpoint{5.684201in}{6.451159in}}%
\pgfpathlineto{\pgfqpoint{5.697891in}{6.970983in}}%
\pgfpathlineto{\pgfqpoint{5.698918in}{6.974731in}}%
\pgfpathlineto{\pgfqpoint{5.699346in}{6.974034in}}%
\pgfpathlineto{\pgfqpoint{5.700373in}{6.966962in}}%
\pgfpathlineto{\pgfqpoint{5.702255in}{6.934819in}}%
\pgfpathlineto{\pgfqpoint{5.705250in}{6.838838in}}%
\pgfpathlineto{\pgfqpoint{5.711581in}{6.542233in}}%
\pgfpathlineto{\pgfqpoint{5.716801in}{6.348610in}}%
\pgfpathlineto{\pgfqpoint{5.719539in}{6.310090in}}%
\pgfpathlineto{\pgfqpoint{5.720138in}{6.308797in}}%
\pgfpathlineto{\pgfqpoint{5.720651in}{6.309777in}}%
\pgfpathlineto{\pgfqpoint{5.721763in}{6.318473in}}%
\pgfpathlineto{\pgfqpoint{5.723731in}{6.355061in}}%
\pgfpathlineto{\pgfqpoint{5.726897in}{6.462608in}}%
\pgfpathlineto{\pgfqpoint{5.739902in}{6.960992in}}%
\pgfpathlineto{\pgfqpoint{5.741186in}{6.966800in}}%
\pgfpathlineto{\pgfqpoint{5.741528in}{6.966303in}}%
\pgfpathlineto{\pgfqpoint{5.742555in}{6.959659in}}%
\pgfpathlineto{\pgfqpoint{5.744352in}{6.929974in}}%
\pgfpathlineto{\pgfqpoint{5.747261in}{6.838951in}}%
\pgfpathlineto{\pgfqpoint{5.753165in}{6.564999in}}%
\pgfpathlineto{\pgfqpoint{5.758555in}{6.359421in}}%
\pgfpathlineto{\pgfqpoint{5.761379in}{6.318012in}}%
\pgfpathlineto{\pgfqpoint{5.762063in}{6.316644in}}%
\pgfpathlineto{\pgfqpoint{5.762491in}{6.317553in}}%
\pgfpathlineto{\pgfqpoint{5.763603in}{6.326226in}}%
\pgfpathlineto{\pgfqpoint{5.765571in}{6.363042in}}%
\pgfpathlineto{\pgfqpoint{5.768737in}{6.471309in}}%
\pgfpathlineto{\pgfqpoint{5.781315in}{6.952667in}}%
\pgfpathlineto{\pgfqpoint{5.782684in}{6.959063in}}%
\pgfpathlineto{\pgfqpoint{5.783026in}{6.958480in}}%
\pgfpathlineto{\pgfqpoint{5.784053in}{6.951511in}}%
\pgfpathlineto{\pgfqpoint{5.785935in}{6.919055in}}%
\pgfpathlineto{\pgfqpoint{5.788930in}{6.821806in}}%
\pgfpathlineto{\pgfqpoint{5.802278in}{6.326618in}}%
\pgfpathlineto{\pgfqpoint{5.803133in}{6.324267in}}%
\pgfpathlineto{\pgfqpoint{5.803561in}{6.325154in}}%
\pgfpathlineto{\pgfqpoint{5.804673in}{6.333854in}}%
\pgfpathlineto{\pgfqpoint{5.806641in}{6.370981in}}%
\pgfpathlineto{\pgfqpoint{5.809807in}{6.480053in}}%
\pgfpathlineto{\pgfqpoint{5.821871in}{6.943568in}}%
\pgfpathlineto{\pgfqpoint{5.823326in}{6.951545in}}%
\pgfpathlineto{\pgfqpoint{5.823668in}{6.951103in}}%
\pgfpathlineto{\pgfqpoint{5.824609in}{6.945333in}}%
\pgfpathlineto{\pgfqpoint{5.826406in}{6.916316in}}%
\pgfpathlineto{\pgfqpoint{5.829315in}{6.825317in}}%
\pgfpathlineto{\pgfqpoint{5.835647in}{6.533423in}}%
\pgfpathlineto{\pgfqpoint{5.840524in}{6.362539in}}%
\pgfpathlineto{\pgfqpoint{5.843091in}{6.332036in}}%
\pgfpathlineto{\pgfqpoint{5.843433in}{6.331691in}}%
\pgfpathlineto{\pgfqpoint{5.843861in}{6.332516in}}%
\pgfpathlineto{\pgfqpoint{5.844973in}{6.341139in}}%
\pgfpathlineto{\pgfqpoint{5.846941in}{6.378405in}}%
\pgfpathlineto{\pgfqpoint{5.850107in}{6.488053in}}%
\pgfpathlineto{\pgfqpoint{5.861743in}{6.935230in}}%
\pgfpathlineto{\pgfqpoint{5.863284in}{6.944207in}}%
\pgfpathlineto{\pgfqpoint{5.863540in}{6.943942in}}%
\pgfpathlineto{\pgfqpoint{5.864396in}{6.939420in}}%
\pgfpathlineto{\pgfqpoint{5.866022in}{6.915771in}}%
\pgfpathlineto{\pgfqpoint{5.868674in}{6.839144in}}%
\pgfpathlineto{\pgfqpoint{5.873722in}{6.613356in}}%
\pgfpathlineto{\pgfqpoint{5.879626in}{6.381268in}}%
\pgfpathlineto{\pgfqpoint{5.882364in}{6.340384in}}%
\pgfpathlineto{\pgfqpoint{5.883048in}{6.338948in}}%
\pgfpathlineto{\pgfqpoint{5.883476in}{6.339886in}}%
\pgfpathlineto{\pgfqpoint{5.884589in}{6.348886in}}%
\pgfpathlineto{\pgfqpoint{5.886557in}{6.387030in}}%
\pgfpathlineto{\pgfqpoint{5.889808in}{6.501763in}}%
\pgfpathlineto{\pgfqpoint{5.900760in}{6.925381in}}%
\pgfpathlineto{\pgfqpoint{5.902471in}{6.937054in}}%
\pgfpathlineto{\pgfqpoint{5.902728in}{6.936854in}}%
\pgfpathlineto{\pgfqpoint{5.903583in}{6.932502in}}%
\pgfpathlineto{\pgfqpoint{5.905209in}{6.908977in}}%
\pgfpathlineto{\pgfqpoint{5.907862in}{6.832145in}}%
\pgfpathlineto{\pgfqpoint{5.912995in}{6.602499in}}%
\pgfpathlineto{\pgfqpoint{5.918642in}{6.385134in}}%
\pgfpathlineto{\pgfqpoint{5.921295in}{6.347162in}}%
\pgfpathlineto{\pgfqpoint{5.921894in}{6.345997in}}%
\pgfpathlineto{\pgfqpoint{5.922407in}{6.347230in}}%
\pgfpathlineto{\pgfqpoint{5.923605in}{6.358043in}}%
\pgfpathlineto{\pgfqpoint{5.925659in}{6.401101in}}%
\pgfpathlineto{\pgfqpoint{5.929081in}{6.528114in}}%
\pgfpathlineto{\pgfqpoint{5.938750in}{6.910182in}}%
\pgfpathlineto{\pgfqpoint{5.940974in}{6.930084in}}%
\pgfpathlineto{\pgfqpoint{5.941488in}{6.929201in}}%
\pgfpathlineto{\pgfqpoint{5.942600in}{6.920226in}}%
\pgfpathlineto{\pgfqpoint{5.944568in}{6.881743in}}%
\pgfpathlineto{\pgfqpoint{5.947819in}{6.766061in}}%
\pgfpathlineto{\pgfqpoint{5.958172in}{6.366114in}}%
\pgfpathlineto{\pgfqpoint{5.960055in}{6.352874in}}%
\pgfpathlineto{\pgfqpoint{5.960140in}{6.352934in}}%
\pgfpathlineto{\pgfqpoint{5.960825in}{6.355497in}}%
\pgfpathlineto{\pgfqpoint{5.962194in}{6.371560in}}%
\pgfpathlineto{\pgfqpoint{5.964589in}{6.432013in}}%
\pgfpathlineto{\pgfqpoint{5.968782in}{6.607284in}}%
\pgfpathlineto{\pgfqpoint{5.975798in}{6.887343in}}%
\pgfpathlineto{\pgfqpoint{5.978451in}{6.922709in}}%
\pgfpathlineto{\pgfqpoint{5.978878in}{6.923292in}}%
\pgfpathlineto{\pgfqpoint{5.979392in}{6.922068in}}%
\pgfpathlineto{\pgfqpoint{5.980590in}{6.911132in}}%
\pgfpathlineto{\pgfqpoint{5.982643in}{6.867481in}}%
\pgfpathlineto{\pgfqpoint{5.986151in}{6.735535in}}%
\pgfpathlineto{\pgfqpoint{5.995135in}{6.381585in}}%
\pgfpathlineto{\pgfqpoint{5.997445in}{6.359580in}}%
\pgfpathlineto{\pgfqpoint{5.997959in}{6.360431in}}%
\pgfpathlineto{\pgfqpoint{5.999071in}{6.369466in}}%
\pgfpathlineto{\pgfqpoint{6.001039in}{6.408418in}}%
\pgfpathlineto{\pgfqpoint{6.004290in}{6.525127in}}%
\pgfpathlineto{\pgfqpoint{6.014045in}{6.901576in}}%
\pgfpathlineto{\pgfqpoint{6.016013in}{6.916678in}}%
\pgfpathlineto{\pgfqpoint{6.016098in}{6.916633in}}%
\pgfpathlineto{\pgfqpoint{6.016783in}{6.914148in}}%
\pgfpathlineto{\pgfqpoint{6.018152in}{6.898047in}}%
\pgfpathlineto{\pgfqpoint{6.020462in}{6.839883in}}%
\pgfpathlineto{\pgfqpoint{6.024654in}{6.665357in}}%
\pgfpathlineto{\pgfqpoint{6.031414in}{6.399662in}}%
\pgfpathlineto{\pgfqpoint{6.033981in}{6.366574in}}%
\pgfpathlineto{\pgfqpoint{6.034408in}{6.366135in}}%
\pgfpathlineto{\pgfqpoint{6.034836in}{6.367179in}}%
\pgfpathlineto{\pgfqpoint{6.035949in}{6.376771in}}%
\pgfpathlineto{\pgfqpoint{6.037917in}{6.416864in}}%
\pgfpathlineto{\pgfqpoint{6.041253in}{6.538886in}}%
\pgfpathlineto{\pgfqpoint{6.050409in}{6.892709in}}%
\pgfpathlineto{\pgfqpoint{6.052462in}{6.910225in}}%
\pgfpathlineto{\pgfqpoint{6.052548in}{6.910217in}}%
\pgfpathlineto{\pgfqpoint{6.053147in}{6.908489in}}%
\pgfpathlineto{\pgfqpoint{6.054430in}{6.895072in}}%
\pgfpathlineto{\pgfqpoint{6.056655in}{6.842621in}}%
\pgfpathlineto{\pgfqpoint{6.060505in}{6.688106in}}%
\pgfpathlineto{\pgfqpoint{6.067778in}{6.402434in}}%
\pgfpathlineto{\pgfqpoint{6.070259in}{6.372725in}}%
\pgfpathlineto{\pgfqpoint{6.070601in}{6.372510in}}%
\pgfpathlineto{\pgfqpoint{6.071029in}{6.373591in}}%
\pgfpathlineto{\pgfqpoint{6.072141in}{6.383354in}}%
\pgfpathlineto{\pgfqpoint{6.074109in}{6.423950in}}%
\pgfpathlineto{\pgfqpoint{6.077446in}{6.546776in}}%
\pgfpathlineto{\pgfqpoint{6.086259in}{6.885859in}}%
\pgfpathlineto{\pgfqpoint{6.088398in}{6.903936in}}%
\pgfpathlineto{\pgfqpoint{6.088997in}{6.902297in}}%
\pgfpathlineto{\pgfqpoint{6.090281in}{6.888962in}}%
\pgfpathlineto{\pgfqpoint{6.092505in}{6.836371in}}%
\pgfpathlineto{\pgfqpoint{6.096441in}{6.677736in}}%
\pgfpathlineto{\pgfqpoint{6.103372in}{6.408312in}}%
\pgfpathlineto{\pgfqpoint{6.105853in}{6.378881in}}%
\pgfpathlineto{\pgfqpoint{6.106195in}{6.378747in}}%
\pgfpathlineto{\pgfqpoint{6.106538in}{6.379584in}}%
\pgfpathlineto{\pgfqpoint{6.107650in}{6.388945in}}%
\pgfpathlineto{\pgfqpoint{6.109618in}{6.429114in}}%
\pgfpathlineto{\pgfqpoint{6.112955in}{6.551658in}}%
\pgfpathlineto{\pgfqpoint{6.121511in}{6.879385in}}%
\pgfpathlineto{\pgfqpoint{6.123650in}{6.897806in}}%
\pgfpathlineto{\pgfqpoint{6.124163in}{6.896607in}}%
\pgfpathlineto{\pgfqpoint{6.125361in}{6.885365in}}%
\pgfpathlineto{\pgfqpoint{6.127415in}{6.840269in}}%
\pgfpathlineto{\pgfqpoint{6.131008in}{6.702160in}}%
\pgfpathlineto{\pgfqpoint{6.138624in}{6.408674in}}%
\pgfpathlineto{\pgfqpoint{6.140934in}{6.384780in}}%
\pgfpathlineto{\pgfqpoint{6.141276in}{6.384985in}}%
\pgfpathlineto{\pgfqpoint{6.141447in}{6.385456in}}%
\pgfpathlineto{\pgfqpoint{6.142474in}{6.393384in}}%
\pgfpathlineto{\pgfqpoint{6.144356in}{6.429619in}}%
\pgfpathlineto{\pgfqpoint{6.147522in}{6.541690in}}%
\pgfpathlineto{\pgfqpoint{6.156421in}{6.877199in}}%
\pgfpathlineto{\pgfqpoint{6.158303in}{6.891828in}}%
\pgfpathlineto{\pgfqpoint{6.158474in}{6.891692in}}%
\pgfpathlineto{\pgfqpoint{6.159244in}{6.888036in}}%
\pgfpathlineto{\pgfqpoint{6.160784in}{6.866161in}}%
\pgfpathlineto{\pgfqpoint{6.163351in}{6.791595in}}%
\pgfpathlineto{\pgfqpoint{6.169939in}{6.506928in}}%
\pgfpathlineto{\pgfqpoint{6.173704in}{6.402862in}}%
\pgfpathlineto{\pgfqpoint{6.175415in}{6.390643in}}%
\pgfpathlineto{\pgfqpoint{6.175672in}{6.390932in}}%
\pgfpathlineto{\pgfqpoint{6.176528in}{6.395905in}}%
\pgfpathlineto{\pgfqpoint{6.178153in}{6.421805in}}%
\pgfpathlineto{\pgfqpoint{6.180891in}{6.507297in}}%
\pgfpathlineto{\pgfqpoint{6.191587in}{6.883288in}}%
\pgfpathlineto{\pgfqpoint{6.192442in}{6.885984in}}%
\pgfpathlineto{\pgfqpoint{6.192870in}{6.885001in}}%
\pgfpathlineto{\pgfqpoint{6.193982in}{6.875237in}}%
\pgfpathlineto{\pgfqpoint{6.195950in}{6.833916in}}%
\pgfpathlineto{\pgfqpoint{6.199373in}{6.706036in}}%
\pgfpathlineto{\pgfqpoint{6.206988in}{6.417405in}}%
\pgfpathlineto{\pgfqpoint{6.209213in}{6.396403in}}%
\pgfpathlineto{\pgfqpoint{6.209726in}{6.397488in}}%
\pgfpathlineto{\pgfqpoint{6.210838in}{6.407484in}}%
\pgfpathlineto{\pgfqpoint{6.212806in}{6.449266in}}%
\pgfpathlineto{\pgfqpoint{6.216314in}{6.581441in}}%
\pgfpathlineto{\pgfqpoint{6.223673in}{6.858956in}}%
\pgfpathlineto{\pgfqpoint{6.225897in}{6.880307in}}%
\pgfpathlineto{\pgfqpoint{6.226411in}{6.879277in}}%
\pgfpathlineto{\pgfqpoint{6.227523in}{6.869366in}}%
\pgfpathlineto{\pgfqpoint{6.229491in}{6.827628in}}%
\pgfpathlineto{\pgfqpoint{6.232999in}{6.695491in}}%
\pgfpathlineto{\pgfqpoint{6.240272in}{6.422799in}}%
\pgfpathlineto{\pgfqpoint{6.242496in}{6.402021in}}%
\pgfpathlineto{\pgfqpoint{6.243010in}{6.403214in}}%
\pgfpathlineto{\pgfqpoint{6.244122in}{6.413503in}}%
\pgfpathlineto{\pgfqpoint{6.246176in}{6.458410in}}%
\pgfpathlineto{\pgfqpoint{6.249769in}{6.596369in}}%
\pgfpathlineto{\pgfqpoint{6.256614in}{6.852231in}}%
\pgfpathlineto{\pgfqpoint{6.258839in}{6.874736in}}%
\pgfpathlineto{\pgfqpoint{6.259352in}{6.873929in}}%
\pgfpathlineto{\pgfqpoint{6.260379in}{6.865525in}}%
\pgfpathlineto{\pgfqpoint{6.262261in}{6.827881in}}%
\pgfpathlineto{\pgfqpoint{6.265598in}{6.706454in}}%
\pgfpathlineto{\pgfqpoint{6.273213in}{6.424677in}}%
\pgfpathlineto{\pgfqpoint{6.275181in}{6.407506in}}%
\pgfpathlineto{\pgfqpoint{6.275352in}{6.407586in}}%
\pgfpathlineto{\pgfqpoint{6.276037in}{6.410447in}}%
\pgfpathlineto{\pgfqpoint{6.277491in}{6.429714in}}%
\pgfpathlineto{\pgfqpoint{6.279973in}{6.499246in}}%
\pgfpathlineto{\pgfqpoint{6.291353in}{6.869342in}}%
\pgfpathlineto{\pgfqpoint{6.291866in}{6.868330in}}%
\pgfpathlineto{\pgfqpoint{6.292978in}{6.858328in}}%
\pgfpathlineto{\pgfqpoint{6.294946in}{6.816109in}}%
\pgfpathlineto{\pgfqpoint{6.298454in}{6.683565in}}%
\pgfpathlineto{\pgfqpoint{6.305214in}{6.433990in}}%
\pgfpathlineto{\pgfqpoint{6.307438in}{6.412850in}}%
\pgfpathlineto{\pgfqpoint{6.307952in}{6.414057in}}%
\pgfpathlineto{\pgfqpoint{6.309064in}{6.424502in}}%
\pgfpathlineto{\pgfqpoint{6.311117in}{6.470001in}}%
\pgfpathlineto{\pgfqpoint{6.314797in}{6.612067in}}%
\pgfpathlineto{\pgfqpoint{6.321043in}{6.841127in}}%
\pgfpathlineto{\pgfqpoint{6.323267in}{6.864036in}}%
\pgfpathlineto{\pgfqpoint{6.323695in}{6.863523in}}%
\pgfpathlineto{\pgfqpoint{6.323781in}{6.863228in}}%
\pgfpathlineto{\pgfqpoint{6.324807in}{6.854714in}}%
\pgfpathlineto{\pgfqpoint{6.326690in}{6.816591in}}%
\pgfpathlineto{\pgfqpoint{6.330027in}{6.694667in}}%
\pgfpathlineto{\pgfqpoint{6.337128in}{6.435719in}}%
\pgfpathlineto{\pgfqpoint{6.339182in}{6.418077in}}%
\pgfpathlineto{\pgfqpoint{6.339267in}{6.418141in}}%
\pgfpathlineto{\pgfqpoint{6.339952in}{6.420975in}}%
\pgfpathlineto{\pgfqpoint{6.341321in}{6.438754in}}%
\pgfpathlineto{\pgfqpoint{6.343717in}{6.504348in}}%
\pgfpathlineto{\pgfqpoint{6.354840in}{6.858902in}}%
\pgfpathlineto{\pgfqpoint{6.355353in}{6.857619in}}%
\pgfpathlineto{\pgfqpoint{6.356551in}{6.845662in}}%
\pgfpathlineto{\pgfqpoint{6.358690in}{6.795589in}}%
\pgfpathlineto{\pgfqpoint{6.362712in}{6.635952in}}%
\pgfpathlineto{\pgfqpoint{6.368102in}{6.445645in}}%
\pgfpathlineto{\pgfqpoint{6.370327in}{6.423173in}}%
\pgfpathlineto{\pgfqpoint{6.370840in}{6.424152in}}%
\pgfpathlineto{\pgfqpoint{6.371952in}{6.434229in}}%
\pgfpathlineto{\pgfqpoint{6.373920in}{6.476916in}}%
\pgfpathlineto{\pgfqpoint{6.377599in}{6.616974in}}%
\pgfpathlineto{\pgfqpoint{6.383674in}{6.833886in}}%
\pgfpathlineto{\pgfqpoint{6.385813in}{6.853871in}}%
\pgfpathlineto{\pgfqpoint{6.386327in}{6.852661in}}%
\pgfpathlineto{\pgfqpoint{6.387439in}{6.842070in}}%
\pgfpathlineto{\pgfqpoint{6.389493in}{6.795987in}}%
\pgfpathlineto{\pgfqpoint{6.393343in}{6.646597in}}%
\pgfpathlineto{\pgfqpoint{6.398990in}{6.448295in}}%
\pgfpathlineto{\pgfqpoint{6.401129in}{6.428140in}}%
\pgfpathlineto{\pgfqpoint{6.401643in}{6.429326in}}%
\pgfpathlineto{\pgfqpoint{6.402755in}{6.439895in}}%
\pgfpathlineto{\pgfqpoint{6.404808in}{6.486003in}}%
\pgfpathlineto{\pgfqpoint{6.408659in}{6.635209in}}%
\pgfpathlineto{\pgfqpoint{6.414220in}{6.829157in}}%
\pgfpathlineto{\pgfqpoint{6.416359in}{6.848953in}}%
\pgfpathlineto{\pgfqpoint{6.416873in}{6.847658in}}%
\pgfpathlineto{\pgfqpoint{6.418071in}{6.835561in}}%
\pgfpathlineto{\pgfqpoint{6.420210in}{6.785020in}}%
\pgfpathlineto{\pgfqpoint{6.424402in}{6.618710in}}%
\pgfpathlineto{\pgfqpoint{6.429365in}{6.451954in}}%
\pgfpathlineto{\pgfqpoint{6.431418in}{6.432998in}}%
\pgfpathlineto{\pgfqpoint{6.431504in}{6.433021in}}%
\pgfpathlineto{\pgfqpoint{6.432103in}{6.435031in}}%
\pgfpathlineto{\pgfqpoint{6.433386in}{6.450009in}}%
\pgfpathlineto{\pgfqpoint{6.435696in}{6.509778in}}%
\pgfpathlineto{\pgfqpoint{6.441258in}{6.735117in}}%
\pgfpathlineto{\pgfqpoint{6.444937in}{6.834346in}}%
\pgfpathlineto{\pgfqpoint{6.446392in}{6.844153in}}%
\pgfpathlineto{\pgfqpoint{6.446734in}{6.843696in}}%
\pgfpathlineto{\pgfqpoint{6.447675in}{6.837017in}}%
\pgfpathlineto{\pgfqpoint{6.449472in}{6.803326in}}%
\pgfpathlineto{\pgfqpoint{6.452638in}{6.692687in}}%
\pgfpathlineto{\pgfqpoint{6.459483in}{6.452405in}}%
\pgfpathlineto{\pgfqpoint{6.461280in}{6.437741in}}%
\pgfpathlineto{\pgfqpoint{6.461451in}{6.437856in}}%
\pgfpathlineto{\pgfqpoint{6.462135in}{6.440963in}}%
\pgfpathlineto{\pgfqpoint{6.463590in}{6.461240in}}%
\pgfpathlineto{\pgfqpoint{6.466157in}{6.535938in}}%
\pgfpathlineto{\pgfqpoint{6.475397in}{6.837472in}}%
\pgfpathlineto{\pgfqpoint{6.476082in}{6.839465in}}%
\pgfpathlineto{\pgfqpoint{6.476595in}{6.838166in}}%
\pgfpathlineto{\pgfqpoint{6.477793in}{6.825956in}}%
\pgfpathlineto{\pgfqpoint{6.479932in}{6.775051in}}%
\pgfpathlineto{\pgfqpoint{6.484381in}{6.599267in}}%
\pgfpathlineto{\pgfqpoint{6.488916in}{6.457089in}}%
\pgfpathlineto{\pgfqpoint{6.490713in}{6.442374in}}%
\pgfpathlineto{\pgfqpoint{6.490884in}{6.442492in}}%
\pgfpathlineto{\pgfqpoint{6.491569in}{6.445626in}}%
\pgfpathlineto{\pgfqpoint{6.493023in}{6.466021in}}%
\pgfpathlineto{\pgfqpoint{6.495590in}{6.540908in}}%
\pgfpathlineto{\pgfqpoint{6.504574in}{6.832597in}}%
\pgfpathlineto{\pgfqpoint{6.505259in}{6.834885in}}%
\pgfpathlineto{\pgfqpoint{6.505772in}{6.833792in}}%
\pgfpathlineto{\pgfqpoint{6.506884in}{6.823266in}}%
\pgfpathlineto{\pgfqpoint{6.508938in}{6.776897in}}%
\pgfpathlineto{\pgfqpoint{6.512959in}{6.622141in}}%
\pgfpathlineto{\pgfqpoint{6.517836in}{6.463269in}}%
\pgfpathlineto{\pgfqpoint{6.519719in}{6.446902in}}%
\pgfpathlineto{\pgfqpoint{6.519890in}{6.447005in}}%
\pgfpathlineto{\pgfqpoint{6.520574in}{6.450091in}}%
\pgfpathlineto{\pgfqpoint{6.522029in}{6.470449in}}%
\pgfpathlineto{\pgfqpoint{6.524596in}{6.545295in}}%
\pgfpathlineto{\pgfqpoint{6.533323in}{6.827683in}}%
\pgfpathlineto{\pgfqpoint{6.534093in}{6.830414in}}%
\pgfpathlineto{\pgfqpoint{6.534607in}{6.829211in}}%
\pgfpathlineto{\pgfqpoint{6.535719in}{6.818417in}}%
\pgfpathlineto{\pgfqpoint{6.537772in}{6.771563in}}%
\pgfpathlineto{\pgfqpoint{6.541965in}{6.610125in}}%
\pgfpathlineto{\pgfqpoint{6.546585in}{6.465508in}}%
\pgfpathlineto{\pgfqpoint{6.548382in}{6.451327in}}%
\pgfpathlineto{\pgfqpoint{6.548553in}{6.451511in}}%
\pgfpathlineto{\pgfqpoint{6.549323in}{6.455661in}}%
\pgfpathlineto{\pgfqpoint{6.550863in}{6.479660in}}%
\pgfpathlineto{\pgfqpoint{6.553601in}{6.564632in}}%
\pgfpathlineto{\pgfqpoint{6.561302in}{6.819352in}}%
\pgfpathlineto{\pgfqpoint{6.562500in}{6.826041in}}%
\pgfpathlineto{\pgfqpoint{6.562928in}{6.825234in}}%
\pgfpathlineto{\pgfqpoint{6.563954in}{6.816482in}}%
\pgfpathlineto{\pgfqpoint{6.565837in}{6.777226in}}%
\pgfpathlineto{\pgfqpoint{6.569516in}{6.641947in}}%
\pgfpathlineto{\pgfqpoint{6.574735in}{6.471019in}}%
\pgfpathlineto{\pgfqpoint{6.576618in}{6.455655in}}%
\pgfpathlineto{\pgfqpoint{6.576703in}{6.455726in}}%
\pgfpathlineto{\pgfqpoint{6.577388in}{6.458722in}}%
\pgfpathlineto{\pgfqpoint{6.578842in}{6.478998in}}%
\pgfpathlineto{\pgfqpoint{6.581409in}{6.553687in}}%
\pgfpathlineto{\pgfqpoint{6.589709in}{6.818594in}}%
\pgfpathlineto{\pgfqpoint{6.590564in}{6.821766in}}%
\pgfpathlineto{\pgfqpoint{6.591078in}{6.820418in}}%
\pgfpathlineto{\pgfqpoint{6.592276in}{6.807947in}}%
\pgfpathlineto{\pgfqpoint{6.594415in}{6.756447in}}%
\pgfpathlineto{\pgfqpoint{6.599720in}{6.552789in}}%
\pgfpathlineto{\pgfqpoint{6.603228in}{6.466327in}}%
\pgfpathlineto{\pgfqpoint{6.604426in}{6.459870in}}%
\pgfpathlineto{\pgfqpoint{6.604853in}{6.460779in}}%
\pgfpathlineto{\pgfqpoint{6.605880in}{6.469805in}}%
\pgfpathlineto{\pgfqpoint{6.607848in}{6.511995in}}%
\pgfpathlineto{\pgfqpoint{6.611784in}{6.658306in}}%
\pgfpathlineto{\pgfqpoint{6.616490in}{6.804300in}}%
\pgfpathlineto{\pgfqpoint{6.618201in}{6.817593in}}%
\pgfpathlineto{\pgfqpoint{6.618458in}{6.817262in}}%
\pgfpathlineto{\pgfqpoint{6.619313in}{6.811764in}}%
\pgfpathlineto{\pgfqpoint{6.621025in}{6.781537in}}%
\pgfpathlineto{\pgfqpoint{6.624105in}{6.678351in}}%
\pgfpathlineto{\pgfqpoint{6.630265in}{6.476011in}}%
\pgfpathlineto{\pgfqpoint{6.631891in}{6.463998in}}%
\pgfpathlineto{\pgfqpoint{6.632148in}{6.464331in}}%
\pgfpathlineto{\pgfqpoint{6.633003in}{6.469842in}}%
\pgfpathlineto{\pgfqpoint{6.634715in}{6.500110in}}%
\pgfpathlineto{\pgfqpoint{6.637795in}{6.603192in}}%
\pgfpathlineto{\pgfqpoint{6.643870in}{6.801553in}}%
\pgfpathlineto{\pgfqpoint{6.645496in}{6.813509in}}%
\pgfpathlineto{\pgfqpoint{6.645752in}{6.813165in}}%
\pgfpathlineto{\pgfqpoint{6.646608in}{6.807608in}}%
\pgfpathlineto{\pgfqpoint{6.648319in}{6.777243in}}%
\pgfpathlineto{\pgfqpoint{6.651485in}{6.670839in}}%
\pgfpathlineto{\pgfqpoint{6.657389in}{6.479894in}}%
\pgfpathlineto{\pgfqpoint{6.659014in}{6.468037in}}%
\pgfpathlineto{\pgfqpoint{6.659271in}{6.468401in}}%
\pgfpathlineto{\pgfqpoint{6.660127in}{6.474026in}}%
\pgfpathlineto{\pgfqpoint{6.661838in}{6.504525in}}%
\pgfpathlineto{\pgfqpoint{6.665004in}{6.610893in}}%
\pgfpathlineto{\pgfqpoint{6.670822in}{6.797784in}}%
\pgfpathlineto{\pgfqpoint{6.672448in}{6.809514in}}%
\pgfpathlineto{\pgfqpoint{6.672704in}{6.809127in}}%
\pgfpathlineto{\pgfqpoint{6.673560in}{6.803421in}}%
\pgfpathlineto{\pgfqpoint{6.675271in}{6.772768in}}%
\pgfpathlineto{\pgfqpoint{6.678437in}{6.666428in}}%
\pgfpathlineto{\pgfqpoint{6.684170in}{6.483586in}}%
\pgfpathlineto{\pgfqpoint{6.685795in}{6.471990in}}%
\pgfpathlineto{\pgfqpoint{6.686052in}{6.472401in}}%
\pgfpathlineto{\pgfqpoint{6.686993in}{6.479137in}}%
\pgfpathlineto{\pgfqpoint{6.688790in}{6.513356in}}%
\pgfpathlineto{\pgfqpoint{6.692213in}{6.632142in}}%
\pgfpathlineto{\pgfqpoint{6.697346in}{6.792865in}}%
\pgfpathlineto{\pgfqpoint{6.698972in}{6.805611in}}%
\pgfpathlineto{\pgfqpoint{6.699229in}{6.805383in}}%
\pgfpathlineto{\pgfqpoint{6.699999in}{6.801016in}}%
\pgfpathlineto{\pgfqpoint{6.701539in}{6.776438in}}%
\pgfpathlineto{\pgfqpoint{6.704362in}{6.688407in}}%
\pgfpathlineto{\pgfqpoint{6.710865in}{6.483835in}}%
\pgfpathlineto{\pgfqpoint{6.712149in}{6.475845in}}%
\pgfpathlineto{\pgfqpoint{6.712576in}{6.476595in}}%
\pgfpathlineto{\pgfqpoint{6.713603in}{6.485314in}}%
\pgfpathlineto{\pgfqpoint{6.715486in}{6.524614in}}%
\pgfpathlineto{\pgfqpoint{6.719421in}{6.666115in}}%
\pgfpathlineto{\pgfqpoint{6.723785in}{6.791649in}}%
\pgfpathlineto{\pgfqpoint{6.725240in}{6.801796in}}%
\pgfpathlineto{\pgfqpoint{6.725582in}{6.801320in}}%
\pgfpathlineto{\pgfqpoint{6.726523in}{6.794393in}}%
\pgfpathlineto{\pgfqpoint{6.728320in}{6.759854in}}%
\pgfpathlineto{\pgfqpoint{6.731828in}{6.638331in}}%
\pgfpathlineto{\pgfqpoint{6.736705in}{6.490974in}}%
\pgfpathlineto{\pgfqpoint{6.738245in}{6.479619in}}%
\pgfpathlineto{\pgfqpoint{6.738587in}{6.480094in}}%
\pgfpathlineto{\pgfqpoint{6.739529in}{6.487022in}}%
\pgfpathlineto{\pgfqpoint{6.741325in}{6.521559in}}%
\pgfpathlineto{\pgfqpoint{6.744919in}{6.646053in}}%
\pgfpathlineto{\pgfqpoint{6.749625in}{6.786618in}}%
\pgfpathlineto{\pgfqpoint{6.751165in}{6.798061in}}%
\pgfpathlineto{\pgfqpoint{6.751507in}{6.797604in}}%
\pgfpathlineto{\pgfqpoint{6.752449in}{6.790723in}}%
\pgfpathlineto{\pgfqpoint{6.754245in}{6.756271in}}%
\pgfpathlineto{\pgfqpoint{6.757839in}{6.632215in}}%
\pgfpathlineto{\pgfqpoint{6.762545in}{6.493728in}}%
\pgfpathlineto{\pgfqpoint{6.764085in}{6.483318in}}%
\pgfpathlineto{\pgfqpoint{6.764342in}{6.483735in}}%
\pgfpathlineto{\pgfqpoint{6.765283in}{6.490513in}}%
\pgfpathlineto{\pgfqpoint{6.767080in}{6.524781in}}%
\pgfpathlineto{\pgfqpoint{6.770673in}{6.648308in}}%
\pgfpathlineto{\pgfqpoint{6.775294in}{6.783714in}}%
\pgfpathlineto{\pgfqpoint{6.776834in}{6.794410in}}%
\pgfpathlineto{\pgfqpoint{6.777091in}{6.794042in}}%
\pgfpathlineto{\pgfqpoint{6.777946in}{6.788371in}}%
\pgfpathlineto{\pgfqpoint{6.779657in}{6.757785in}}%
\pgfpathlineto{\pgfqpoint{6.782994in}{6.646837in}}%
\pgfpathlineto{\pgfqpoint{6.788043in}{6.496841in}}%
\pgfpathlineto{\pgfqpoint{6.789497in}{6.486922in}}%
\pgfpathlineto{\pgfqpoint{6.789839in}{6.487459in}}%
\pgfpathlineto{\pgfqpoint{6.790781in}{6.494563in}}%
\pgfpathlineto{\pgfqpoint{6.792577in}{6.529346in}}%
\pgfpathlineto{\pgfqpoint{6.796257in}{6.655731in}}%
\pgfpathlineto{\pgfqpoint{6.800706in}{6.781545in}}%
\pgfpathlineto{\pgfqpoint{6.802160in}{6.790831in}}%
\pgfpathlineto{\pgfqpoint{6.802417in}{6.790417in}}%
\pgfpathlineto{\pgfqpoint{6.803358in}{6.783644in}}%
\pgfpathlineto{\pgfqpoint{6.805155in}{6.749428in}}%
\pgfpathlineto{\pgfqpoint{6.808834in}{6.623923in}}%
\pgfpathlineto{\pgfqpoint{6.813283in}{6.499280in}}%
\pgfpathlineto{\pgfqpoint{6.814652in}{6.490456in}}%
\pgfpathlineto{\pgfqpoint{6.814995in}{6.490988in}}%
\pgfpathlineto{\pgfqpoint{6.815936in}{6.498079in}}%
\pgfpathlineto{\pgfqpoint{6.817733in}{6.532802in}}%
\pgfpathlineto{\pgfqpoint{6.821497in}{6.661429in}}%
\pgfpathlineto{\pgfqpoint{6.825776in}{6.778833in}}%
\pgfpathlineto{\pgfqpoint{6.827145in}{6.787342in}}%
\pgfpathlineto{\pgfqpoint{6.827487in}{6.786729in}}%
\pgfpathlineto{\pgfqpoint{6.828428in}{6.779424in}}%
\pgfpathlineto{\pgfqpoint{6.830225in}{6.744369in}}%
\pgfpathlineto{\pgfqpoint{6.834075in}{6.612849in}}%
\pgfpathlineto{\pgfqpoint{6.838182in}{6.502273in}}%
\pgfpathlineto{\pgfqpoint{6.839551in}{6.493919in}}%
\pgfpathlineto{\pgfqpoint{6.839893in}{6.494570in}}%
\pgfpathlineto{\pgfqpoint{6.840920in}{6.503052in}}%
\pgfpathlineto{\pgfqpoint{6.842803in}{6.541782in}}%
\pgfpathlineto{\pgfqpoint{6.847252in}{6.695416in}}%
\pgfpathlineto{\pgfqpoint{6.850760in}{6.778431in}}%
\pgfpathlineto{\pgfqpoint{6.851872in}{6.783918in}}%
\pgfpathlineto{\pgfqpoint{6.852300in}{6.782938in}}%
\pgfpathlineto{\pgfqpoint{6.853412in}{6.772472in}}%
\pgfpathlineto{\pgfqpoint{6.855466in}{6.726407in}}%
\pgfpathlineto{\pgfqpoint{6.864108in}{6.497300in}}%
\pgfpathlineto{\pgfqpoint{6.864279in}{6.497459in}}%
\pgfpathlineto{\pgfqpoint{6.865049in}{6.501569in}}%
\pgfpathlineto{\pgfqpoint{6.866589in}{6.525611in}}%
\pgfpathlineto{\pgfqpoint{6.869498in}{6.613957in}}%
\pgfpathlineto{\pgfqpoint{6.874974in}{6.772666in}}%
\pgfpathlineto{\pgfqpoint{6.876257in}{6.780572in}}%
\pgfpathlineto{\pgfqpoint{6.876685in}{6.779785in}}%
\pgfpathlineto{\pgfqpoint{6.877712in}{6.770974in}}%
\pgfpathlineto{\pgfqpoint{6.879680in}{6.729472in}}%
\pgfpathlineto{\pgfqpoint{6.888407in}{6.500613in}}%
\pgfpathlineto{\pgfqpoint{6.888835in}{6.501550in}}%
\pgfpathlineto{\pgfqpoint{6.889862in}{6.510706in}}%
\pgfpathlineto{\pgfqpoint{6.891830in}{6.552703in}}%
\pgfpathlineto{\pgfqpoint{6.900472in}{6.777289in}}%
\pgfpathlineto{\pgfqpoint{6.900728in}{6.776885in}}%
\pgfpathlineto{\pgfqpoint{6.901584in}{6.771108in}}%
\pgfpathlineto{\pgfqpoint{6.903295in}{6.740532in}}%
\pgfpathlineto{\pgfqpoint{6.906803in}{6.626899in}}%
\pgfpathlineto{\pgfqpoint{6.911081in}{6.512067in}}%
\pgfpathlineto{\pgfqpoint{6.912450in}{6.503862in}}%
\pgfpathlineto{\pgfqpoint{6.912793in}{6.504543in}}%
\pgfpathlineto{\pgfqpoint{6.913819in}{6.513090in}}%
\pgfpathlineto{\pgfqpoint{6.915787in}{6.554025in}}%
\pgfpathlineto{\pgfqpoint{6.924344in}{6.774086in}}%
\pgfpathlineto{\pgfqpoint{6.924686in}{6.773461in}}%
\pgfpathlineto{\pgfqpoint{6.925627in}{6.766145in}}%
\pgfpathlineto{\pgfqpoint{6.927509in}{6.729096in}}%
\pgfpathlineto{\pgfqpoint{6.932301in}{6.571228in}}%
\pgfpathlineto{\pgfqpoint{6.935467in}{6.509311in}}%
\pgfpathlineto{\pgfqpoint{6.936151in}{6.507029in}}%
\pgfpathlineto{\pgfqpoint{6.936665in}{6.508198in}}%
\pgfpathlineto{\pgfqpoint{6.937777in}{6.519036in}}%
\pgfpathlineto{\pgfqpoint{6.939916in}{6.567727in}}%
\pgfpathlineto{\pgfqpoint{6.947788in}{6.770831in}}%
\pgfpathlineto{\pgfqpoint{6.948044in}{6.770902in}}%
\pgfpathlineto{\pgfqpoint{6.948387in}{6.770035in}}%
\pgfpathlineto{\pgfqpoint{6.949413in}{6.760969in}}%
\pgfpathlineto{\pgfqpoint{6.951381in}{6.719398in}}%
\pgfpathlineto{\pgfqpoint{6.959681in}{6.510138in}}%
\pgfpathlineto{\pgfqpoint{6.959852in}{6.510306in}}%
\pgfpathlineto{\pgfqpoint{6.960622in}{6.514435in}}%
\pgfpathlineto{\pgfqpoint{6.962162in}{6.538336in}}%
\pgfpathlineto{\pgfqpoint{6.965242in}{6.630487in}}%
\pgfpathlineto{\pgfqpoint{6.970119in}{6.761246in}}%
\pgfpathlineto{\pgfqpoint{6.971317in}{6.767873in}}%
\pgfpathlineto{\pgfqpoint{6.971745in}{6.767006in}}%
\pgfpathlineto{\pgfqpoint{6.972772in}{6.758062in}}%
\pgfpathlineto{\pgfqpoint{6.974740in}{6.716811in}}%
\pgfpathlineto{\pgfqpoint{6.982868in}{6.513185in}}%
\pgfpathlineto{\pgfqpoint{6.983125in}{6.513404in}}%
\pgfpathlineto{\pgfqpoint{6.983895in}{6.517731in}}%
\pgfpathlineto{\pgfqpoint{6.985521in}{6.543817in}}%
\pgfpathlineto{\pgfqpoint{6.988772in}{6.642539in}}%
\pgfpathlineto{\pgfqpoint{6.993221in}{6.758053in}}%
\pgfpathlineto{\pgfqpoint{6.994419in}{6.764861in}}%
\pgfpathlineto{\pgfqpoint{6.994847in}{6.764068in}}%
\pgfpathlineto{\pgfqpoint{6.995874in}{6.755315in}}%
\pgfpathlineto{\pgfqpoint{6.997842in}{6.714502in}}%
\pgfpathlineto{\pgfqpoint{7.005885in}{6.516159in}}%
\pgfpathlineto{\pgfqpoint{7.006056in}{6.516267in}}%
\pgfpathlineto{\pgfqpoint{7.006740in}{6.519413in}}%
\pgfpathlineto{\pgfqpoint{7.008195in}{6.539841in}}%
\pgfpathlineto{\pgfqpoint{7.011018in}{6.619224in}}%
\pgfpathlineto{\pgfqpoint{7.016152in}{6.755713in}}%
\pgfpathlineto{\pgfqpoint{7.017350in}{6.761907in}}%
\pgfpathlineto{\pgfqpoint{7.017692in}{6.761238in}}%
\pgfpathlineto{\pgfqpoint{7.018719in}{6.752803in}}%
\pgfpathlineto{\pgfqpoint{7.020687in}{6.712630in}}%
\pgfpathlineto{\pgfqpoint{7.028644in}{6.519076in}}%
\pgfpathlineto{\pgfqpoint{7.028815in}{6.519187in}}%
\pgfpathlineto{\pgfqpoint{7.029500in}{6.522339in}}%
\pgfpathlineto{\pgfqpoint{7.030954in}{6.542717in}}%
\pgfpathlineto{\pgfqpoint{7.033778in}{6.621533in}}%
\pgfpathlineto{\pgfqpoint{7.038826in}{6.753208in}}%
\pgfpathlineto{\pgfqpoint{7.039938in}{6.759026in}}%
\pgfpathlineto{\pgfqpoint{7.040366in}{6.758225in}}%
\pgfpathlineto{\pgfqpoint{7.041393in}{6.749501in}}%
\pgfpathlineto{\pgfqpoint{7.043361in}{6.709043in}}%
\pgfpathlineto{\pgfqpoint{7.051147in}{6.521938in}}%
\pgfpathlineto{\pgfqpoint{7.051318in}{6.522018in}}%
\pgfpathlineto{\pgfqpoint{7.052003in}{6.525033in}}%
\pgfpathlineto{\pgfqpoint{7.053457in}{6.545084in}}%
\pgfpathlineto{\pgfqpoint{7.056281in}{6.623018in}}%
\pgfpathlineto{\pgfqpoint{7.061243in}{6.750500in}}%
\pgfpathlineto{\pgfqpoint{7.062356in}{6.756199in}}%
\pgfpathlineto{\pgfqpoint{7.062784in}{6.755363in}}%
\pgfpathlineto{\pgfqpoint{7.063810in}{6.746585in}}%
\pgfpathlineto{\pgfqpoint{7.065778in}{6.706217in}}%
\pgfpathlineto{\pgfqpoint{7.073479in}{6.524731in}}%
\pgfpathlineto{\pgfqpoint{7.073564in}{6.524766in}}%
\pgfpathlineto{\pgfqpoint{7.074163in}{6.526899in}}%
\pgfpathlineto{\pgfqpoint{7.075447in}{6.542172in}}%
\pgfpathlineto{\pgfqpoint{7.077928in}{6.604575in}}%
\pgfpathlineto{\pgfqpoint{7.083661in}{6.749888in}}%
\pgfpathlineto{\pgfqpoint{7.084516in}{6.753427in}}%
\pgfpathlineto{\pgfqpoint{7.085030in}{6.752325in}}%
\pgfpathlineto{\pgfqpoint{7.086142in}{6.741813in}}%
\pgfpathlineto{\pgfqpoint{7.088281in}{6.694892in}}%
\pgfpathlineto{\pgfqpoint{7.095212in}{6.527990in}}%
\pgfpathlineto{\pgfqpoint{7.095554in}{6.527472in}}%
\pgfpathlineto{\pgfqpoint{7.096067in}{6.528711in}}%
\pgfpathlineto{\pgfqpoint{7.097180in}{6.539488in}}%
\pgfpathlineto{\pgfqpoint{7.099319in}{6.586667in}}%
\pgfpathlineto{\pgfqpoint{7.106164in}{6.750196in}}%
\pgfpathlineto{\pgfqpoint{7.106506in}{6.750716in}}%
\pgfpathlineto{\pgfqpoint{7.107019in}{6.749486in}}%
\pgfpathlineto{\pgfqpoint{7.108132in}{6.738748in}}%
\pgfpathlineto{\pgfqpoint{7.110271in}{6.691747in}}%
\pgfpathlineto{\pgfqpoint{7.117030in}{6.530784in}}%
\pgfpathlineto{\pgfqpoint{7.117458in}{6.530172in}}%
\pgfpathlineto{\pgfqpoint{7.117886in}{6.531233in}}%
\pgfpathlineto{\pgfqpoint{7.118998in}{6.541627in}}%
\pgfpathlineto{\pgfqpoint{7.121137in}{6.588010in}}%
\pgfpathlineto{\pgfqpoint{7.127897in}{6.747500in}}%
\pgfpathlineto{\pgfqpoint{7.128324in}{6.748031in}}%
\pgfpathlineto{\pgfqpoint{7.128752in}{6.746894in}}%
\pgfpathlineto{\pgfqpoint{7.129865in}{6.736333in}}%
\pgfpathlineto{\pgfqpoint{7.132004in}{6.689838in}}%
\pgfpathlineto{\pgfqpoint{7.138678in}{6.533386in}}%
\pgfpathlineto{\pgfqpoint{7.139105in}{6.532807in}}%
\pgfpathlineto{\pgfqpoint{7.139533in}{6.533892in}}%
\pgfpathlineto{\pgfqpoint{7.140645in}{6.544303in}}%
\pgfpathlineto{\pgfqpoint{7.142785in}{6.590455in}}%
\pgfpathlineto{\pgfqpoint{7.149373in}{6.744697in}}%
\pgfpathlineto{\pgfqpoint{7.149801in}{6.745454in}}%
\pgfpathlineto{\pgfqpoint{7.150314in}{6.744172in}}%
\pgfpathlineto{\pgfqpoint{7.151512in}{6.732156in}}%
\pgfpathlineto{\pgfqpoint{7.153822in}{6.679677in}}%
\pgfpathlineto{\pgfqpoint{7.159811in}{6.537334in}}%
\pgfpathlineto{\pgfqpoint{7.160496in}{6.535367in}}%
\pgfpathlineto{\pgfqpoint{7.161009in}{6.536670in}}%
\pgfpathlineto{\pgfqpoint{7.162207in}{6.548703in}}%
\pgfpathlineto{\pgfqpoint{7.164517in}{6.601039in}}%
\pgfpathlineto{\pgfqpoint{7.170507in}{6.741258in}}%
\pgfpathlineto{\pgfqpoint{7.171106in}{6.742905in}}%
\pgfpathlineto{\pgfqpoint{7.171619in}{6.741743in}}%
\pgfpathlineto{\pgfqpoint{7.172731in}{6.731279in}}%
\pgfpathlineto{\pgfqpoint{7.174870in}{6.685460in}}%
\pgfpathlineto{\pgfqpoint{7.181288in}{6.538655in}}%
\pgfpathlineto{\pgfqpoint{7.181715in}{6.537890in}}%
\pgfpathlineto{\pgfqpoint{7.182229in}{6.539144in}}%
\pgfpathlineto{\pgfqpoint{7.183427in}{6.551015in}}%
\pgfpathlineto{\pgfqpoint{7.185737in}{6.602794in}}%
\pgfpathlineto{\pgfqpoint{7.191641in}{6.738798in}}%
\pgfpathlineto{\pgfqpoint{7.192240in}{6.740406in}}%
\pgfpathlineto{\pgfqpoint{7.192753in}{6.739226in}}%
\pgfpathlineto{\pgfqpoint{7.193865in}{6.728775in}}%
\pgfpathlineto{\pgfqpoint{7.196090in}{6.680930in}}%
\pgfpathlineto{\pgfqpoint{7.202250in}{6.541391in}}%
\pgfpathlineto{\pgfqpoint{7.202764in}{6.540367in}}%
\pgfpathlineto{\pgfqpoint{7.203277in}{6.541699in}}%
\pgfpathlineto{\pgfqpoint{7.204475in}{6.553682in}}%
\pgfpathlineto{\pgfqpoint{7.206785in}{6.605252in}}%
\pgfpathlineto{\pgfqpoint{7.212518in}{6.736055in}}%
\pgfpathlineto{\pgfqpoint{7.213202in}{6.737954in}}%
\pgfpathlineto{\pgfqpoint{7.213716in}{6.736640in}}%
\pgfpathlineto{\pgfqpoint{7.214914in}{6.724729in}}%
\pgfpathlineto{\pgfqpoint{7.217224in}{6.673432in}}%
\pgfpathlineto{\pgfqpoint{7.222957in}{6.544432in}}%
\pgfpathlineto{\pgfqpoint{7.223555in}{6.542786in}}%
\pgfpathlineto{\pgfqpoint{7.224069in}{6.543910in}}%
\pgfpathlineto{\pgfqpoint{7.225181in}{6.554165in}}%
\pgfpathlineto{\pgfqpoint{7.227406in}{6.601256in}}%
\pgfpathlineto{\pgfqpoint{7.233395in}{6.734377in}}%
\pgfpathlineto{\pgfqpoint{7.233909in}{6.735559in}}%
\pgfpathlineto{\pgfqpoint{7.234422in}{6.734408in}}%
\pgfpathlineto{\pgfqpoint{7.235534in}{6.724125in}}%
\pgfpathlineto{\pgfqpoint{7.237759in}{6.677160in}}%
\pgfpathlineto{\pgfqpoint{7.243663in}{6.546488in}}%
\pgfpathlineto{\pgfqpoint{7.244262in}{6.545172in}}%
\pgfpathlineto{\pgfqpoint{7.244775in}{6.546562in}}%
\pgfpathlineto{\pgfqpoint{7.245973in}{6.558545in}}%
\pgfpathlineto{\pgfqpoint{7.248283in}{6.609371in}}%
\pgfpathlineto{\pgfqpoint{7.253845in}{6.731522in}}%
\pgfpathlineto{\pgfqpoint{7.254529in}{6.733190in}}%
\pgfpathlineto{\pgfqpoint{7.254957in}{6.732140in}}%
\pgfpathlineto{\pgfqpoint{7.256069in}{6.722090in}}%
\pgfpathlineto{\pgfqpoint{7.258294in}{6.675784in}}%
\pgfpathlineto{\pgfqpoint{7.264198in}{6.548469in}}%
\pgfpathlineto{\pgfqpoint{7.264711in}{6.547494in}}%
\pgfpathlineto{\pgfqpoint{7.265224in}{6.548829in}}%
\pgfpathlineto{\pgfqpoint{7.266422in}{6.560621in}}%
\pgfpathlineto{\pgfqpoint{7.268732in}{6.610818in}}%
\pgfpathlineto{\pgfqpoint{7.274208in}{6.729206in}}%
\pgfpathlineto{\pgfqpoint{7.274893in}{6.730889in}}%
\pgfpathlineto{\pgfqpoint{7.275321in}{6.729864in}}%
\pgfpathlineto{\pgfqpoint{7.276433in}{6.719931in}}%
\pgfpathlineto{\pgfqpoint{7.278658in}{6.674130in}}%
\pgfpathlineto{\pgfqpoint{7.284390in}{6.551124in}}%
\pgfpathlineto{\pgfqpoint{7.284989in}{6.549775in}}%
\pgfpathlineto{\pgfqpoint{7.285503in}{6.551100in}}%
\pgfpathlineto{\pgfqpoint{7.286701in}{6.562804in}}%
\pgfpathlineto{\pgfqpoint{7.289096in}{6.614778in}}%
\pgfpathlineto{\pgfqpoint{7.294316in}{6.726495in}}%
\pgfpathlineto{\pgfqpoint{7.295000in}{6.728646in}}%
\pgfpathlineto{\pgfqpoint{7.295513in}{6.727597in}}%
\pgfpathlineto{\pgfqpoint{7.296626in}{6.717705in}}%
\pgfpathlineto{\pgfqpoint{7.298850in}{6.672313in}}%
\pgfpathlineto{\pgfqpoint{7.304498in}{6.553315in}}%
\pgfpathlineto{\pgfqpoint{7.305096in}{6.552012in}}%
\pgfpathlineto{\pgfqpoint{7.305610in}{6.553358in}}%
\pgfpathlineto{\pgfqpoint{7.306808in}{6.565033in}}%
\pgfpathlineto{\pgfqpoint{7.309203in}{6.616536in}}%
\pgfpathlineto{\pgfqpoint{7.314337in}{6.724353in}}%
\pgfpathlineto{\pgfqpoint{7.315022in}{6.726435in}}%
\pgfpathlineto{\pgfqpoint{7.315535in}{6.725355in}}%
\pgfpathlineto{\pgfqpoint{7.316647in}{6.715464in}}%
\pgfpathlineto{\pgfqpoint{7.318872in}{6.670435in}}%
\pgfpathlineto{\pgfqpoint{7.324434in}{6.555448in}}%
\pgfpathlineto{\pgfqpoint{7.325033in}{6.554206in}}%
\pgfpathlineto{\pgfqpoint{7.325460in}{6.555199in}}%
\pgfpathlineto{\pgfqpoint{7.326573in}{6.564902in}}%
\pgfpathlineto{\pgfqpoint{7.328797in}{6.609492in}}%
\pgfpathlineto{\pgfqpoint{7.334359in}{6.723135in}}%
\pgfpathlineto{\pgfqpoint{7.334872in}{6.724268in}}%
\pgfpathlineto{\pgfqpoint{7.335386in}{6.723154in}}%
\pgfpathlineto{\pgfqpoint{7.336498in}{6.713258in}}%
\pgfpathlineto{\pgfqpoint{7.338723in}{6.668595in}}%
\pgfpathlineto{\pgfqpoint{7.344199in}{6.557542in}}%
\pgfpathlineto{\pgfqpoint{7.344797in}{6.556356in}}%
\pgfpathlineto{\pgfqpoint{7.345225in}{6.557374in}}%
\pgfpathlineto{\pgfqpoint{7.346338in}{6.567074in}}%
\pgfpathlineto{\pgfqpoint{7.348562in}{6.611287in}}%
\pgfpathlineto{\pgfqpoint{7.354038in}{6.721052in}}%
\pgfpathlineto{\pgfqpoint{7.354552in}{6.722144in}}%
\pgfpathlineto{\pgfqpoint{7.355065in}{6.721006in}}%
\pgfpathlineto{\pgfqpoint{7.356263in}{6.709979in}}%
\pgfpathlineto{\pgfqpoint{7.358573in}{6.662584in}}%
\pgfpathlineto{\pgfqpoint{7.363707in}{6.560023in}}%
\pgfpathlineto{\pgfqpoint{7.364306in}{6.558442in}}%
\pgfpathlineto{\pgfqpoint{7.364819in}{6.559489in}}%
\pgfpathlineto{\pgfqpoint{7.365931in}{6.569143in}}%
\pgfpathlineto{\pgfqpoint{7.368156in}{6.612913in}}%
\pgfpathlineto{\pgfqpoint{7.373546in}{6.718986in}}%
\pgfpathlineto{\pgfqpoint{7.374060in}{6.720062in}}%
\pgfpathlineto{\pgfqpoint{7.374573in}{6.718929in}}%
\pgfpathlineto{\pgfqpoint{7.375771in}{6.707991in}}%
\pgfpathlineto{\pgfqpoint{7.378167in}{6.658994in}}%
\pgfpathlineto{\pgfqpoint{7.383044in}{6.562553in}}%
\pgfpathlineto{\pgfqpoint{7.383728in}{6.560503in}}%
\pgfpathlineto{\pgfqpoint{7.384242in}{6.561528in}}%
\pgfpathlineto{\pgfqpoint{7.385354in}{6.571066in}}%
\pgfpathlineto{\pgfqpoint{7.387579in}{6.614293in}}%
\pgfpathlineto{\pgfqpoint{7.392884in}{6.716921in}}%
\pgfpathlineto{\pgfqpoint{7.393397in}{6.718021in}}%
\pgfpathlineto{\pgfqpoint{7.393910in}{6.716933in}}%
\pgfpathlineto{\pgfqpoint{7.395023in}{6.707303in}}%
\pgfpathlineto{\pgfqpoint{7.397247in}{6.664153in}}%
\pgfpathlineto{\pgfqpoint{7.402467in}{6.563748in}}%
\pgfpathlineto{\pgfqpoint{7.403066in}{6.562534in}}%
\pgfpathlineto{\pgfqpoint{7.403579in}{6.563852in}}%
\pgfpathlineto{\pgfqpoint{7.404777in}{6.575066in}}%
\pgfpathlineto{\pgfqpoint{7.407173in}{6.623681in}}%
\pgfpathlineto{\pgfqpoint{7.411878in}{6.713972in}}%
\pgfpathlineto{\pgfqpoint{7.412648in}{6.716005in}}%
\pgfpathlineto{\pgfqpoint{7.413076in}{6.715029in}}%
\pgfpathlineto{\pgfqpoint{7.414189in}{6.705673in}}%
\pgfpathlineto{\pgfqpoint{7.416413in}{6.663281in}}%
\pgfpathlineto{\pgfqpoint{7.421547in}{6.565845in}}%
\pgfpathlineto{\pgfqpoint{7.422146in}{6.564504in}}%
\pgfpathlineto{\pgfqpoint{7.422659in}{6.565691in}}%
\pgfpathlineto{\pgfqpoint{7.423857in}{6.576534in}}%
\pgfpathlineto{\pgfqpoint{7.426253in}{6.624220in}}%
\pgfpathlineto{\pgfqpoint{7.430959in}{6.712302in}}%
\pgfpathlineto{\pgfqpoint{7.431643in}{6.714059in}}%
\pgfpathlineto{\pgfqpoint{7.432157in}{6.712875in}}%
\pgfpathlineto{\pgfqpoint{7.433355in}{6.702080in}}%
\pgfpathlineto{\pgfqpoint{7.435750in}{6.654698in}}%
\pgfpathlineto{\pgfqpoint{7.440456in}{6.567992in}}%
\pgfpathlineto{\pgfqpoint{7.441141in}{6.566459in}}%
\pgfpathlineto{\pgfqpoint{7.441569in}{6.567428in}}%
\pgfpathlineto{\pgfqpoint{7.442681in}{6.576657in}}%
\pgfpathlineto{\pgfqpoint{7.444906in}{6.618273in}}%
\pgfpathlineto{\pgfqpoint{7.449954in}{6.710956in}}%
\pgfpathlineto{\pgfqpoint{7.450553in}{6.712128in}}%
\pgfpathlineto{\pgfqpoint{7.451066in}{6.710833in}}%
\pgfpathlineto{\pgfqpoint{7.452264in}{6.699889in}}%
\pgfpathlineto{\pgfqpoint{7.454745in}{6.650799in}}%
\pgfpathlineto{\pgfqpoint{7.459194in}{6.570222in}}%
\pgfpathlineto{\pgfqpoint{7.459879in}{6.568343in}}%
\pgfpathlineto{\pgfqpoint{7.460392in}{6.569396in}}%
\pgfpathlineto{\pgfqpoint{7.461505in}{6.578686in}}%
\pgfpathlineto{\pgfqpoint{7.463815in}{6.621932in}}%
\pgfpathlineto{\pgfqpoint{7.468606in}{6.708727in}}%
\pgfpathlineto{\pgfqpoint{7.469291in}{6.710242in}}%
\pgfpathlineto{\pgfqpoint{7.469719in}{6.709293in}}%
\pgfpathlineto{\pgfqpoint{7.470831in}{6.700223in}}%
\pgfpathlineto{\pgfqpoint{7.473141in}{6.657487in}}%
\pgfpathlineto{\pgfqpoint{7.477933in}{6.571652in}}%
\pgfpathlineto{\pgfqpoint{7.478532in}{6.570206in}}%
\pgfpathlineto{\pgfqpoint{7.479045in}{6.571231in}}%
\pgfpathlineto{\pgfqpoint{7.480157in}{6.580382in}}%
\pgfpathlineto{\pgfqpoint{7.482467in}{6.622996in}}%
\pgfpathlineto{\pgfqpoint{7.487259in}{6.707214in}}%
\pgfpathlineto{\pgfqpoint{7.487858in}{6.708405in}}%
\pgfpathlineto{\pgfqpoint{7.488371in}{6.707175in}}%
\pgfpathlineto{\pgfqpoint{7.489569in}{6.696556in}}%
\pgfpathlineto{\pgfqpoint{7.492050in}{6.648932in}}%
\pgfpathlineto{\pgfqpoint{7.496414in}{6.573610in}}%
\pgfpathlineto{\pgfqpoint{7.497099in}{6.572042in}}%
\pgfpathlineto{\pgfqpoint{7.497612in}{6.573276in}}%
\pgfpathlineto{\pgfqpoint{7.498810in}{6.583855in}}%
\pgfpathlineto{\pgfqpoint{7.501291in}{6.631158in}}%
\pgfpathlineto{\pgfqpoint{7.505569in}{6.704795in}}%
\pgfpathlineto{\pgfqpoint{7.506254in}{6.706606in}}%
\pgfpathlineto{\pgfqpoint{7.506767in}{6.705569in}}%
\pgfpathlineto{\pgfqpoint{7.507965in}{6.695461in}}%
\pgfpathlineto{\pgfqpoint{7.510361in}{6.650898in}}%
\pgfpathlineto{\pgfqpoint{7.514810in}{6.575259in}}%
\pgfpathlineto{\pgfqpoint{7.515409in}{6.573823in}}%
\pgfpathlineto{\pgfqpoint{7.515922in}{6.574805in}}%
\pgfpathlineto{\pgfqpoint{7.517035in}{6.583703in}}%
\pgfpathlineto{\pgfqpoint{7.519345in}{6.625073in}}%
\pgfpathlineto{\pgfqpoint{7.523965in}{6.703510in}}%
\pgfpathlineto{\pgfqpoint{7.524564in}{6.704834in}}%
\pgfpathlineto{\pgfqpoint{7.525078in}{6.703769in}}%
\pgfpathlineto{\pgfqpoint{7.526275in}{6.693697in}}%
\pgfpathlineto{\pgfqpoint{7.528757in}{6.647803in}}%
\pgfpathlineto{\pgfqpoint{7.533035in}{6.577024in}}%
\pgfpathlineto{\pgfqpoint{7.533719in}{6.575595in}}%
\pgfpathlineto{\pgfqpoint{7.534147in}{6.576523in}}%
\pgfpathlineto{\pgfqpoint{7.535259in}{6.585264in}}%
\pgfpathlineto{\pgfqpoint{7.537570in}{6.625967in}}%
\pgfpathlineto{\pgfqpoint{7.542190in}{6.702052in}}%
\pgfpathlineto{\pgfqpoint{7.542789in}{6.703077in}}%
\pgfpathlineto{\pgfqpoint{7.543217in}{6.702136in}}%
\pgfpathlineto{\pgfqpoint{7.544329in}{6.693406in}}%
\pgfpathlineto{\pgfqpoint{7.546639in}{6.652954in}}%
\pgfpathlineto{\pgfqpoint{7.551174in}{6.578525in}}%
\pgfpathlineto{\pgfqpoint{7.551773in}{6.577298in}}%
\pgfpathlineto{\pgfqpoint{7.552286in}{6.578407in}}%
\pgfpathlineto{\pgfqpoint{7.553484in}{6.588424in}}%
\pgfpathlineto{\pgfqpoint{7.555966in}{6.633414in}}%
\pgfpathlineto{\pgfqpoint{7.560158in}{6.700119in}}%
\pgfpathlineto{\pgfqpoint{7.560757in}{6.701396in}}%
\pgfpathlineto{\pgfqpoint{7.561270in}{6.700343in}}%
\pgfpathlineto{\pgfqpoint{7.562468in}{6.690495in}}%
\pgfpathlineto{\pgfqpoint{7.564950in}{6.645978in}}%
\pgfpathlineto{\pgfqpoint{7.569057in}{6.580547in}}%
\pgfpathlineto{\pgfqpoint{7.569741in}{6.578984in}}%
\pgfpathlineto{\pgfqpoint{7.570255in}{6.580108in}}%
\pgfpathlineto{\pgfqpoint{7.571452in}{6.590059in}}%
\pgfpathlineto{\pgfqpoint{7.573934in}{6.634404in}}%
\pgfpathlineto{\pgfqpoint{7.578041in}{6.698417in}}%
\pgfpathlineto{\pgfqpoint{7.578640in}{6.699729in}}%
\pgfpathlineto{\pgfqpoint{7.579153in}{6.698735in}}%
\pgfpathlineto{\pgfqpoint{7.580351in}{6.689112in}}%
\pgfpathlineto{\pgfqpoint{7.582832in}{6.645413in}}%
\pgfpathlineto{\pgfqpoint{7.586939in}{6.581932in}}%
\pgfpathlineto{\pgfqpoint{7.587538in}{6.580630in}}%
\pgfpathlineto{\pgfqpoint{7.588052in}{6.581620in}}%
\pgfpathlineto{\pgfqpoint{7.589249in}{6.591180in}}%
\pgfpathlineto{\pgfqpoint{7.591731in}{6.634525in}}%
\pgfpathlineto{\pgfqpoint{7.595838in}{6.696939in}}%
\pgfpathlineto{\pgfqpoint{7.596437in}{6.698090in}}%
\pgfpathlineto{\pgfqpoint{7.596950in}{6.696985in}}%
\pgfpathlineto{\pgfqpoint{7.598148in}{6.687228in}}%
\pgfpathlineto{\pgfqpoint{7.600715in}{6.642199in}}%
\pgfpathlineto{\pgfqpoint{7.604651in}{6.583483in}}%
\pgfpathlineto{\pgfqpoint{7.605250in}{6.582248in}}%
\pgfpathlineto{\pgfqpoint{7.605763in}{6.583268in}}%
\pgfpathlineto{\pgfqpoint{7.606961in}{6.592791in}}%
\pgfpathlineto{\pgfqpoint{7.609442in}{6.635518in}}%
\pgfpathlineto{\pgfqpoint{7.613464in}{6.695312in}}%
\pgfpathlineto{\pgfqpoint{7.614063in}{6.696492in}}%
\pgfpathlineto{\pgfqpoint{7.614576in}{6.695439in}}%
\pgfpathlineto{\pgfqpoint{7.615774in}{6.685898in}}%
\pgfpathlineto{\pgfqpoint{7.618341in}{6.641718in}}%
\pgfpathlineto{\pgfqpoint{7.622191in}{6.585191in}}%
\pgfpathlineto{\pgfqpoint{7.622875in}{6.583849in}}%
\pgfpathlineto{\pgfqpoint{7.623303in}{6.584719in}}%
\pgfpathlineto{\pgfqpoint{7.624416in}{6.592892in}}%
\pgfpathlineto{\pgfqpoint{7.626811in}{6.632179in}}%
\pgfpathlineto{\pgfqpoint{7.631004in}{6.693877in}}%
\pgfpathlineto{\pgfqpoint{7.631603in}{6.694914in}}%
\pgfpathlineto{\pgfqpoint{7.632116in}{6.693767in}}%
\pgfpathlineto{\pgfqpoint{7.633314in}{6.684131in}}%
\pgfpathlineto{\pgfqpoint{7.635881in}{6.640476in}}%
\pgfpathlineto{\pgfqpoint{7.639731in}{6.586356in}}%
\pgfpathlineto{\pgfqpoint{7.640330in}{6.585405in}}%
\pgfpathlineto{\pgfqpoint{7.640758in}{6.586281in}}%
\pgfpathlineto{\pgfqpoint{7.641870in}{6.594382in}}%
\pgfpathlineto{\pgfqpoint{7.644266in}{6.633055in}}%
\pgfpathlineto{\pgfqpoint{7.648373in}{6.692286in}}%
\pgfpathlineto{\pgfqpoint{7.648972in}{6.693383in}}%
\pgfpathlineto{\pgfqpoint{7.649485in}{6.692317in}}%
\pgfpathlineto{\pgfqpoint{7.650683in}{6.682960in}}%
\pgfpathlineto{\pgfqpoint{7.653250in}{6.640220in}}%
\pgfpathlineto{\pgfqpoint{7.657015in}{6.588026in}}%
\pgfpathlineto{\pgfqpoint{7.657614in}{6.586910in}}%
\pgfpathlineto{\pgfqpoint{7.658127in}{6.587946in}}%
\pgfpathlineto{\pgfqpoint{7.659325in}{6.597186in}}%
\pgfpathlineto{\pgfqpoint{7.661892in}{6.639492in}}%
\pgfpathlineto{\pgfqpoint{7.665657in}{6.690863in}}%
\pgfpathlineto{\pgfqpoint{7.666256in}{6.691871in}}%
\pgfpathlineto{\pgfqpoint{7.666769in}{6.690759in}}%
\pgfpathlineto{\pgfqpoint{7.667967in}{6.681411in}}%
\pgfpathlineto{\pgfqpoint{7.670619in}{6.637693in}}%
\pgfpathlineto{\pgfqpoint{7.674213in}{6.589532in}}%
\pgfpathlineto{\pgfqpoint{7.674812in}{6.588399in}}%
\pgfpathlineto{\pgfqpoint{7.675325in}{6.589391in}}%
\pgfpathlineto{\pgfqpoint{7.676523in}{6.598428in}}%
\pgfpathlineto{\pgfqpoint{7.679090in}{6.639898in}}%
\pgfpathlineto{\pgfqpoint{7.682855in}{6.689582in}}%
\pgfpathlineto{\pgfqpoint{7.683368in}{6.690406in}}%
\pgfpathlineto{\pgfqpoint{7.683881in}{6.689430in}}%
\pgfpathlineto{\pgfqpoint{7.685079in}{6.680482in}}%
\pgfpathlineto{\pgfqpoint{7.687646in}{6.639418in}}%
\pgfpathlineto{\pgfqpoint{7.691325in}{6.590884in}}%
\pgfpathlineto{\pgfqpoint{7.691924in}{6.589865in}}%
\pgfpathlineto{\pgfqpoint{7.692438in}{6.590925in}}%
\pgfpathlineto{\pgfqpoint{7.693636in}{6.599998in}}%
\pgfpathlineto{\pgfqpoint{7.696288in}{6.642435in}}%
\pgfpathlineto{\pgfqpoint{7.699882in}{6.688138in}}%
\pgfpathlineto{\pgfqpoint{7.700395in}{6.688961in}}%
\pgfpathlineto{\pgfqpoint{7.700908in}{6.688012in}}%
\pgfpathlineto{\pgfqpoint{7.702106in}{6.679229in}}%
\pgfpathlineto{\pgfqpoint{7.704673in}{6.638963in}}%
\pgfpathlineto{\pgfqpoint{7.708352in}{6.592109in}}%
\pgfpathlineto{\pgfqpoint{7.708866in}{6.591289in}}%
\pgfpathlineto{\pgfqpoint{7.709379in}{6.592226in}}%
\pgfpathlineto{\pgfqpoint{7.710577in}{6.600929in}}%
\pgfpathlineto{\pgfqpoint{7.713229in}{6.642330in}}%
\pgfpathlineto{\pgfqpoint{7.716823in}{6.686818in}}%
\pgfpathlineto{\pgfqpoint{7.717336in}{6.687540in}}%
\pgfpathlineto{\pgfqpoint{7.717850in}{6.686518in}}%
\pgfpathlineto{\pgfqpoint{7.719048in}{6.677691in}}%
\pgfpathlineto{\pgfqpoint{7.721700in}{6.636507in}}%
\pgfpathlineto{\pgfqpoint{7.725208in}{6.593507in}}%
\pgfpathlineto{\pgfqpoint{7.725722in}{6.592690in}}%
\pgfpathlineto{\pgfqpoint{7.726235in}{6.593602in}}%
\pgfpathlineto{\pgfqpoint{7.727433in}{6.602142in}}%
\pgfpathlineto{\pgfqpoint{7.730085in}{6.642713in}}%
\pgfpathlineto{\pgfqpoint{7.733593in}{6.685341in}}%
\pgfpathlineto{\pgfqpoint{7.734107in}{6.686159in}}%
\pgfpathlineto{\pgfqpoint{7.734620in}{6.685261in}}%
\pgfpathlineto{\pgfqpoint{7.735818in}{6.676808in}}%
\pgfpathlineto{\pgfqpoint{7.738470in}{6.636657in}}%
\pgfpathlineto{\pgfqpoint{7.741978in}{6.594792in}}%
\pgfpathlineto{\pgfqpoint{7.742492in}{6.594066in}}%
\pgfpathlineto{\pgfqpoint{7.743005in}{6.595041in}}%
\pgfpathlineto{\pgfqpoint{7.744203in}{6.603601in}}%
\pgfpathlineto{\pgfqpoint{7.746855in}{6.643516in}}%
\pgfpathlineto{\pgfqpoint{7.750278in}{6.683974in}}%
\pgfpathlineto{\pgfqpoint{7.750877in}{6.684775in}}%
\pgfpathlineto{\pgfqpoint{7.751305in}{6.683941in}}%
\pgfpathlineto{\pgfqpoint{7.752503in}{6.675679in}}%
\pgfpathlineto{\pgfqpoint{7.755155in}{6.636383in}}%
\pgfpathlineto{\pgfqpoint{7.758663in}{6.595987in}}%
\pgfpathlineto{\pgfqpoint{7.759176in}{6.595429in}}%
\pgfpathlineto{\pgfqpoint{7.759604in}{6.596242in}}%
\pgfpathlineto{\pgfqpoint{7.760802in}{6.604396in}}%
\pgfpathlineto{\pgfqpoint{7.763455in}{6.643253in}}%
\pgfpathlineto{\pgfqpoint{7.766877in}{6.682710in}}%
\pgfpathlineto{\pgfqpoint{7.767390in}{6.683469in}}%
\pgfpathlineto{\pgfqpoint{7.767904in}{6.682570in}}%
\pgfpathlineto{\pgfqpoint{7.769102in}{6.674339in}}%
\pgfpathlineto{\pgfqpoint{7.771754in}{6.635747in}}%
\pgfpathlineto{\pgfqpoint{7.775177in}{6.597335in}}%
\pgfpathlineto{\pgfqpoint{7.775690in}{6.596737in}}%
\pgfpathlineto{\pgfqpoint{7.776118in}{6.597496in}}%
\pgfpathlineto{\pgfqpoint{7.777316in}{6.605405in}}%
\pgfpathlineto{\pgfqpoint{7.779968in}{6.643359in}}%
\pgfpathlineto{\pgfqpoint{7.783391in}{6.681532in}}%
\pgfpathlineto{\pgfqpoint{7.783904in}{6.682157in}}%
\pgfpathlineto{\pgfqpoint{7.784417in}{6.681153in}}%
\pgfpathlineto{\pgfqpoint{7.785615in}{6.672814in}}%
\pgfpathlineto{\pgfqpoint{7.788439in}{6.632083in}}%
\pgfpathlineto{\pgfqpoint{7.791605in}{6.598599in}}%
\pgfpathlineto{\pgfqpoint{7.792118in}{6.598029in}}%
\pgfpathlineto{\pgfqpoint{7.792546in}{6.598789in}}%
\pgfpathlineto{\pgfqpoint{7.793744in}{6.606592in}}%
\pgfpathlineto{\pgfqpoint{7.796396in}{6.643777in}}%
\pgfpathlineto{\pgfqpoint{7.799819in}{6.680417in}}%
\pgfpathlineto{\pgfqpoint{7.800332in}{6.680850in}}%
\pgfpathlineto{\pgfqpoint{7.800760in}{6.679988in}}%
\pgfpathlineto{\pgfqpoint{7.801958in}{6.671981in}}%
\pgfpathlineto{\pgfqpoint{7.804696in}{6.633616in}}%
\pgfpathlineto{\pgfqpoint{7.807947in}{6.599787in}}%
\pgfpathlineto{\pgfqpoint{7.808460in}{6.599304in}}%
\pgfpathlineto{\pgfqpoint{7.808888in}{6.600113in}}%
\pgfpathlineto{\pgfqpoint{7.810086in}{6.607933in}}%
\pgfpathlineto{\pgfqpoint{7.812824in}{6.645775in}}%
\pgfpathlineto{\pgfqpoint{7.816075in}{6.679154in}}%
\pgfpathlineto{\pgfqpoint{7.816589in}{6.679609in}}%
\pgfpathlineto{\pgfqpoint{7.817017in}{6.678786in}}%
\pgfpathlineto{\pgfqpoint{7.818215in}{6.670987in}}%
\pgfpathlineto{\pgfqpoint{7.820953in}{6.633521in}}%
\pgfpathlineto{\pgfqpoint{7.824204in}{6.600917in}}%
\pgfpathlineto{\pgfqpoint{7.824632in}{6.600518in}}%
\pgfpathlineto{\pgfqpoint{7.825145in}{6.601468in}}%
\pgfpathlineto{\pgfqpoint{7.826343in}{6.609405in}}%
\pgfpathlineto{\pgfqpoint{7.829167in}{6.647902in}}%
\pgfpathlineto{\pgfqpoint{7.832247in}{6.677957in}}%
\pgfpathlineto{\pgfqpoint{7.832760in}{6.678382in}}%
\pgfpathlineto{\pgfqpoint{7.833188in}{6.677557in}}%
\pgfpathlineto{\pgfqpoint{7.834386in}{6.669863in}}%
\pgfpathlineto{\pgfqpoint{7.837209in}{6.631945in}}%
\pgfpathlineto{\pgfqpoint{7.840375in}{6.602005in}}%
\pgfpathlineto{\pgfqpoint{7.840803in}{6.601756in}}%
\pgfpathlineto{\pgfqpoint{7.841231in}{6.602569in}}%
\pgfpathlineto{\pgfqpoint{7.842429in}{6.610179in}}%
\pgfpathlineto{\pgfqpoint{7.845252in}{6.647652in}}%
\pgfpathlineto{\pgfqpoint{7.848333in}{6.676817in}}%
\pgfpathlineto{\pgfqpoint{7.848760in}{6.677216in}}%
\pgfpathlineto{\pgfqpoint{7.849274in}{6.676308in}}%
\pgfpathlineto{\pgfqpoint{7.850472in}{6.668630in}}%
\pgfpathlineto{\pgfqpoint{7.853381in}{6.630269in}}%
\pgfpathlineto{\pgfqpoint{7.856375in}{6.603210in}}%
\pgfpathlineto{\pgfqpoint{7.856803in}{6.602931in}}%
\pgfpathlineto{\pgfqpoint{7.857231in}{6.603694in}}%
\pgfpathlineto{\pgfqpoint{7.858429in}{6.611067in}}%
\pgfpathlineto{\pgfqpoint{7.861252in}{6.647597in}}%
\pgfpathlineto{\pgfqpoint{7.864333in}{6.675724in}}%
\pgfpathlineto{\pgfqpoint{7.864761in}{6.676031in}}%
\pgfpathlineto{\pgfqpoint{7.865188in}{6.675308in}}%
\pgfpathlineto{\pgfqpoint{7.866386in}{6.668088in}}%
\pgfpathlineto{\pgfqpoint{7.869124in}{6.633247in}}%
\pgfpathlineto{\pgfqpoint{7.872290in}{6.604367in}}%
\pgfpathlineto{\pgfqpoint{7.872718in}{6.604089in}}%
\pgfpathlineto{\pgfqpoint{7.873146in}{6.604832in}}%
\pgfpathlineto{\pgfqpoint{7.874344in}{6.612045in}}%
\pgfpathlineto{\pgfqpoint{7.877167in}{6.647699in}}%
\pgfpathlineto{\pgfqpoint{7.880247in}{6.674665in}}%
\pgfpathlineto{\pgfqpoint{7.880675in}{6.674857in}}%
\pgfpathlineto{\pgfqpoint{7.881103in}{6.674039in}}%
\pgfpathlineto{\pgfqpoint{7.882301in}{6.666692in}}%
\pgfpathlineto{\pgfqpoint{7.885210in}{6.630166in}}%
\pgfpathlineto{\pgfqpoint{7.888119in}{6.605479in}}%
\pgfpathlineto{\pgfqpoint{7.888547in}{6.605227in}}%
\pgfpathlineto{\pgfqpoint{7.888975in}{6.605976in}}%
\pgfpathlineto{\pgfqpoint{7.890173in}{6.613092in}}%
\pgfpathlineto{\pgfqpoint{7.892996in}{6.647924in}}%
\pgfpathlineto{\pgfqpoint{7.895991in}{6.673499in}}%
\pgfpathlineto{\pgfqpoint{7.896419in}{6.673755in}}%
\pgfpathlineto{\pgfqpoint{7.896846in}{6.673020in}}%
\pgfpathlineto{\pgfqpoint{7.898044in}{6.665995in}}%
\pgfpathlineto{\pgfqpoint{7.900868in}{6.631604in}}%
\pgfpathlineto{\pgfqpoint{7.903863in}{6.606553in}}%
\pgfpathlineto{\pgfqpoint{7.904290in}{6.606348in}}%
\pgfpathlineto{\pgfqpoint{7.904718in}{6.607122in}}%
\pgfpathlineto{\pgfqpoint{7.905916in}{6.614191in}}%
\pgfpathlineto{\pgfqpoint{7.908825in}{6.649335in}}%
\pgfpathlineto{\pgfqpoint{7.911734in}{6.672508in}}%
\pgfpathlineto{\pgfqpoint{7.912162in}{6.672611in}}%
\pgfpathlineto{\pgfqpoint{7.912504in}{6.671995in}}%
\pgfpathlineto{\pgfqpoint{7.913702in}{6.665236in}}%
\pgfpathlineto{\pgfqpoint{7.916526in}{6.631791in}}%
\pgfpathlineto{\pgfqpoint{7.919520in}{6.607594in}}%
\pgfpathlineto{\pgfqpoint{7.919948in}{6.607451in}}%
\pgfpathlineto{\pgfqpoint{7.920376in}{6.608267in}}%
\pgfpathlineto{\pgfqpoint{7.921660in}{6.616072in}}%
\pgfpathlineto{\pgfqpoint{7.925339in}{6.659384in}}%
\pgfpathlineto{\pgfqpoint{7.927563in}{6.671612in}}%
\pgfpathlineto{\pgfqpoint{7.927649in}{6.671600in}}%
\pgfpathlineto{\pgfqpoint{7.928333in}{6.670145in}}%
\pgfpathlineto{\pgfqpoint{7.929788in}{6.659767in}}%
\pgfpathlineto{\pgfqpoint{7.935350in}{6.608454in}}%
\pgfpathlineto{\pgfqpoint{7.935863in}{6.609160in}}%
\pgfpathlineto{\pgfqpoint{7.937061in}{6.615774in}}%
\pgfpathlineto{\pgfqpoint{7.939970in}{6.649038in}}%
\pgfpathlineto{\pgfqpoint{7.942879in}{6.670464in}}%
\pgfpathlineto{\pgfqpoint{7.943307in}{6.670424in}}%
\pgfpathlineto{\pgfqpoint{7.943649in}{6.669725in}}%
\pgfpathlineto{\pgfqpoint{7.944933in}{6.662197in}}%
\pgfpathlineto{\pgfqpoint{7.950836in}{6.609499in}}%
\pgfpathlineto{\pgfqpoint{7.951778in}{6.611922in}}%
\pgfpathlineto{\pgfqpoint{7.953489in}{6.626028in}}%
\pgfpathlineto{\pgfqpoint{7.958537in}{6.669526in}}%
\pgfpathlineto{\pgfqpoint{7.959221in}{6.668175in}}%
\pgfpathlineto{\pgfqpoint{7.960676in}{6.658335in}}%
\pgfpathlineto{\pgfqpoint{7.966152in}{6.610509in}}%
\pgfpathlineto{\pgfqpoint{7.966665in}{6.611163in}}%
\pgfpathlineto{\pgfqpoint{7.967863in}{6.617438in}}%
\pgfpathlineto{\pgfqpoint{7.970858in}{6.649891in}}%
\pgfpathlineto{\pgfqpoint{7.973596in}{6.668432in}}%
\pgfpathlineto{\pgfqpoint{7.974024in}{6.668403in}}%
\pgfpathlineto{\pgfqpoint{7.974366in}{6.667743in}}%
\pgfpathlineto{\pgfqpoint{7.975649in}{6.660585in}}%
\pgfpathlineto{\pgfqpoint{7.981468in}{6.611508in}}%
\pgfpathlineto{\pgfqpoint{7.982409in}{6.613811in}}%
\pgfpathlineto{\pgfqpoint{7.984120in}{6.627231in}}%
\pgfpathlineto{\pgfqpoint{7.988997in}{6.667545in}}%
\pgfpathlineto{\pgfqpoint{7.989682in}{6.666532in}}%
\pgfpathlineto{\pgfqpoint{7.991051in}{6.658380in}}%
\pgfpathlineto{\pgfqpoint{7.996612in}{6.612474in}}%
\pgfpathlineto{\pgfqpoint{7.997297in}{6.613578in}}%
\pgfpathlineto{\pgfqpoint{7.998666in}{6.621803in}}%
\pgfpathlineto{\pgfqpoint{8.004227in}{6.666579in}}%
\pgfpathlineto{\pgfqpoint{8.004826in}{6.665594in}}%
\pgfpathlineto{\pgfqpoint{8.006195in}{6.657638in}}%
\pgfpathlineto{\pgfqpoint{8.011757in}{6.613429in}}%
\pgfpathlineto{\pgfqpoint{8.012356in}{6.614372in}}%
\pgfpathlineto{\pgfqpoint{8.013725in}{6.622172in}}%
\pgfpathlineto{\pgfqpoint{8.019286in}{6.665640in}}%
\pgfpathlineto{\pgfqpoint{8.019885in}{6.664681in}}%
\pgfpathlineto{\pgfqpoint{8.021254in}{6.656927in}}%
\pgfpathlineto{\pgfqpoint{8.026730in}{6.614352in}}%
\pgfpathlineto{\pgfqpoint{8.027415in}{6.615398in}}%
\pgfpathlineto{\pgfqpoint{8.028784in}{6.623209in}}%
\pgfpathlineto{\pgfqpoint{8.034260in}{6.664724in}}%
\pgfpathlineto{\pgfqpoint{8.034859in}{6.663796in}}%
\pgfpathlineto{\pgfqpoint{8.036228in}{6.656255in}}%
\pgfpathlineto{\pgfqpoint{8.041704in}{6.615262in}}%
\pgfpathlineto{\pgfqpoint{8.042303in}{6.616143in}}%
\pgfpathlineto{\pgfqpoint{8.043672in}{6.623519in}}%
\pgfpathlineto{\pgfqpoint{8.049148in}{6.663829in}}%
\pgfpathlineto{\pgfqpoint{8.049747in}{6.662942in}}%
\pgfpathlineto{\pgfqpoint{8.051116in}{6.655630in}}%
\pgfpathlineto{\pgfqpoint{8.056506in}{6.616145in}}%
\pgfpathlineto{\pgfqpoint{8.057190in}{6.617100in}}%
\pgfpathlineto{\pgfqpoint{8.058559in}{6.624444in}}%
\pgfpathlineto{\pgfqpoint{8.063950in}{6.662956in}}%
\pgfpathlineto{\pgfqpoint{8.064549in}{6.662120in}}%
\pgfpathlineto{\pgfqpoint{8.065918in}{6.655057in}}%
\pgfpathlineto{\pgfqpoint{8.071308in}{6.617010in}}%
\pgfpathlineto{\pgfqpoint{8.071907in}{6.617790in}}%
\pgfpathlineto{\pgfqpoint{8.073191in}{6.624060in}}%
\pgfpathlineto{\pgfqpoint{8.078667in}{6.662103in}}%
\pgfpathlineto{\pgfqpoint{8.079437in}{6.660848in}}%
\pgfpathlineto{\pgfqpoint{8.080977in}{6.651949in}}%
\pgfpathlineto{\pgfqpoint{8.086025in}{6.617857in}}%
\pgfpathlineto{\pgfqpoint{8.086282in}{6.618030in}}%
\pgfpathlineto{\pgfqpoint{8.087308in}{6.621278in}}%
\pgfpathlineto{\pgfqpoint{8.089533in}{6.638574in}}%
\pgfpathlineto{\pgfqpoint{8.093041in}{6.661132in}}%
\pgfpathlineto{\pgfqpoint{8.093555in}{6.661144in}}%
\pgfpathlineto{\pgfqpoint{8.093897in}{6.660575in}}%
\pgfpathlineto{\pgfqpoint{8.095180in}{6.654663in}}%
\pgfpathlineto{\pgfqpoint{8.100656in}{6.618685in}}%
\pgfpathlineto{\pgfqpoint{8.101426in}{6.619993in}}%
\pgfpathlineto{\pgfqpoint{8.102966in}{6.628677in}}%
\pgfpathlineto{\pgfqpoint{8.107929in}{6.660450in}}%
\pgfpathlineto{\pgfqpoint{8.108100in}{6.660362in}}%
\pgfpathlineto{\pgfqpoint{8.109041in}{6.657916in}}%
\pgfpathlineto{\pgfqpoint{8.111009in}{6.644344in}}%
\pgfpathlineto{\pgfqpoint{8.114945in}{6.619559in}}%
\pgfpathlineto{\pgfqpoint{8.115630in}{6.619934in}}%
\pgfpathlineto{\pgfqpoint{8.115801in}{6.620303in}}%
\pgfpathlineto{\pgfqpoint{8.117170in}{6.626797in}}%
\pgfpathlineto{\pgfqpoint{8.122389in}{6.659658in}}%
\pgfpathlineto{\pgfqpoint{8.122988in}{6.658979in}}%
\pgfpathlineto{\pgfqpoint{8.124271in}{6.653358in}}%
\pgfpathlineto{\pgfqpoint{8.129662in}{6.620276in}}%
\pgfpathlineto{\pgfqpoint{8.130346in}{6.621271in}}%
\pgfpathlineto{\pgfqpoint{8.131801in}{6.628488in}}%
\pgfpathlineto{\pgfqpoint{8.136849in}{6.658880in}}%
\pgfpathlineto{\pgfqpoint{8.137191in}{6.658620in}}%
\pgfpathlineto{\pgfqpoint{8.138304in}{6.654973in}}%
\pgfpathlineto{\pgfqpoint{8.140871in}{6.636076in}}%
\pgfpathlineto{\pgfqpoint{8.143780in}{6.621121in}}%
\pgfpathlineto{\pgfqpoint{8.144464in}{6.621427in}}%
\pgfpathlineto{\pgfqpoint{8.144635in}{6.621763in}}%
\pgfpathlineto{\pgfqpoint{8.146004in}{6.627797in}}%
\pgfpathlineto{\pgfqpoint{8.151224in}{6.658119in}}%
\pgfpathlineto{\pgfqpoint{8.151737in}{6.657548in}}%
\pgfpathlineto{\pgfqpoint{8.153020in}{6.652377in}}%
\pgfpathlineto{\pgfqpoint{8.158325in}{6.621784in}}%
\pgfpathlineto{\pgfqpoint{8.159095in}{6.622839in}}%
\pgfpathlineto{\pgfqpoint{8.160635in}{6.630357in}}%
\pgfpathlineto{\pgfqpoint{8.165513in}{6.657379in}}%
\pgfpathlineto{\pgfqpoint{8.165684in}{6.657288in}}%
\pgfpathlineto{\pgfqpoint{8.166625in}{6.655053in}}%
\pgfpathlineto{\pgfqpoint{8.168678in}{6.642370in}}%
\pgfpathlineto{\pgfqpoint{8.172358in}{6.622580in}}%
\pgfpathlineto{\pgfqpoint{8.173128in}{6.623061in}}%
\pgfpathlineto{\pgfqpoint{8.173213in}{6.623235in}}%
\pgfpathlineto{\pgfqpoint{8.174582in}{6.628967in}}%
\pgfpathlineto{\pgfqpoint{8.179716in}{6.656660in}}%
\pgfpathlineto{\pgfqpoint{8.180144in}{6.656266in}}%
\pgfpathlineto{\pgfqpoint{8.181342in}{6.652152in}}%
\pgfpathlineto{\pgfqpoint{8.186818in}{6.623237in}}%
\pgfpathlineto{\pgfqpoint{8.187673in}{6.624660in}}%
\pgfpathlineto{\pgfqpoint{8.189384in}{6.633432in}}%
\pgfpathlineto{\pgfqpoint{8.193748in}{6.655963in}}%
\pgfpathlineto{\pgfqpoint{8.194604in}{6.654908in}}%
\pgfpathlineto{\pgfqpoint{8.196144in}{6.647861in}}%
\pgfpathlineto{\pgfqpoint{8.200850in}{6.623917in}}%
\pgfpathlineto{\pgfqpoint{8.201021in}{6.623974in}}%
\pgfpathlineto{\pgfqpoint{8.201962in}{6.625898in}}%
\pgfpathlineto{\pgfqpoint{8.203930in}{6.636773in}}%
\pgfpathlineto{\pgfqpoint{8.207695in}{6.655243in}}%
\pgfpathlineto{\pgfqpoint{8.208550in}{6.654535in}}%
\pgfpathlineto{\pgfqpoint{8.210005in}{6.648650in}}%
\pgfpathlineto{\pgfqpoint{8.214882in}{6.624591in}}%
\pgfpathlineto{\pgfqpoint{8.215224in}{6.624792in}}%
\pgfpathlineto{\pgfqpoint{8.216337in}{6.627791in}}%
\pgfpathlineto{\pgfqpoint{8.218989in}{6.643859in}}%
\pgfpathlineto{\pgfqpoint{8.221727in}{6.654594in}}%
\pgfpathlineto{\pgfqpoint{8.222583in}{6.653821in}}%
\pgfpathlineto{\pgfqpoint{8.224037in}{6.648019in}}%
\pgfpathlineto{\pgfqpoint{8.228914in}{6.625259in}}%
\pgfpathlineto{\pgfqpoint{8.229171in}{6.625425in}}%
\pgfpathlineto{\pgfqpoint{8.230283in}{6.628273in}}%
\pgfpathlineto{\pgfqpoint{8.232936in}{6.643703in}}%
\pgfpathlineto{\pgfqpoint{8.235674in}{6.653952in}}%
\pgfpathlineto{\pgfqpoint{8.236529in}{6.653163in}}%
\pgfpathlineto{\pgfqpoint{8.238070in}{6.647045in}}%
\pgfpathlineto{\pgfqpoint{8.242775in}{6.625891in}}%
\pgfpathlineto{\pgfqpoint{8.242947in}{6.625960in}}%
\pgfpathlineto{\pgfqpoint{8.243973in}{6.628064in}}%
\pgfpathlineto{\pgfqpoint{8.246198in}{6.639764in}}%
\pgfpathlineto{\pgfqpoint{8.249449in}{6.653286in}}%
\pgfpathlineto{\pgfqpoint{8.250391in}{6.652562in}}%
\pgfpathlineto{\pgfqpoint{8.251931in}{6.646667in}}%
\pgfpathlineto{\pgfqpoint{8.256637in}{6.626520in}}%
\pgfpathlineto{\pgfqpoint{8.256808in}{6.626602in}}%
\pgfpathlineto{\pgfqpoint{8.257834in}{6.628716in}}%
\pgfpathlineto{\pgfqpoint{8.260145in}{6.640569in}}%
\pgfpathlineto{\pgfqpoint{8.263310in}{6.652703in}}%
\pgfpathlineto{\pgfqpoint{8.264252in}{6.651840in}}%
\pgfpathlineto{\pgfqpoint{8.265877in}{6.645516in}}%
\pgfpathlineto{\pgfqpoint{8.270412in}{6.627128in}}%
\pgfpathlineto{\pgfqpoint{8.271353in}{6.628527in}}%
\pgfpathlineto{\pgfqpoint{8.273236in}{6.636806in}}%
\pgfpathlineto{\pgfqpoint{8.277086in}{6.652119in}}%
\pgfpathlineto{\pgfqpoint{8.278027in}{6.651179in}}%
\pgfpathlineto{\pgfqpoint{8.279653in}{6.644959in}}%
\pgfpathlineto{\pgfqpoint{8.284102in}{6.627714in}}%
\pgfpathlineto{\pgfqpoint{8.285043in}{6.629044in}}%
\pgfpathlineto{\pgfqpoint{8.286926in}{6.636978in}}%
\pgfpathlineto{\pgfqpoint{8.290690in}{6.651522in}}%
\pgfpathlineto{\pgfqpoint{8.291632in}{6.650757in}}%
\pgfpathlineto{\pgfqpoint{8.293257in}{6.644954in}}%
\pgfpathlineto{\pgfqpoint{8.297707in}{6.628280in}}%
\pgfpathlineto{\pgfqpoint{8.297792in}{6.628303in}}%
\pgfpathlineto{\pgfqpoint{8.298819in}{6.629933in}}%
\pgfpathlineto{\pgfqpoint{8.300958in}{6.639258in}}%
\pgfpathlineto{\pgfqpoint{8.304295in}{6.650961in}}%
\pgfpathlineto{\pgfqpoint{8.305236in}{6.650218in}}%
\pgfpathlineto{\pgfqpoint{8.306862in}{6.644638in}}%
\pgfpathlineto{\pgfqpoint{8.311311in}{6.628842in}}%
\pgfpathlineto{\pgfqpoint{8.312338in}{6.630293in}}%
\pgfpathlineto{\pgfqpoint{8.314391in}{6.638676in}}%
\pgfpathlineto{\pgfqpoint{8.317814in}{6.650408in}}%
\pgfpathlineto{\pgfqpoint{8.318840in}{6.649576in}}%
\pgfpathlineto{\pgfqpoint{8.320552in}{6.643664in}}%
\pgfpathlineto{\pgfqpoint{8.324744in}{6.629371in}}%
\pgfpathlineto{\pgfqpoint{8.325771in}{6.630600in}}%
\pgfpathlineto{\pgfqpoint{8.327739in}{6.638042in}}%
\pgfpathlineto{\pgfqpoint{8.331333in}{6.649890in}}%
\pgfpathlineto{\pgfqpoint{8.332359in}{6.648997in}}%
\pgfpathlineto{\pgfqpoint{8.334156in}{6.642848in}}%
\pgfpathlineto{\pgfqpoint{8.338092in}{6.629899in}}%
\pgfpathlineto{\pgfqpoint{8.339119in}{6.630869in}}%
\pgfpathlineto{\pgfqpoint{8.340916in}{6.636991in}}%
\pgfpathlineto{\pgfqpoint{8.344851in}{6.649389in}}%
\pgfpathlineto{\pgfqpoint{8.345878in}{6.648314in}}%
\pgfpathlineto{\pgfqpoint{8.347761in}{6.641817in}}%
\pgfpathlineto{\pgfqpoint{8.351440in}{6.630413in}}%
\pgfpathlineto{\pgfqpoint{8.352467in}{6.631267in}}%
\pgfpathlineto{\pgfqpoint{8.354263in}{6.637035in}}%
\pgfpathlineto{\pgfqpoint{8.358199in}{6.648888in}}%
\pgfpathlineto{\pgfqpoint{8.359226in}{6.647872in}}%
\pgfpathlineto{\pgfqpoint{8.361194in}{6.641324in}}%
\pgfpathlineto{\pgfqpoint{8.364788in}{6.630900in}}%
\pgfpathlineto{\pgfqpoint{8.365900in}{6.631926in}}%
\pgfpathlineto{\pgfqpoint{8.367868in}{6.638363in}}%
\pgfpathlineto{\pgfqpoint{8.371376in}{6.648387in}}%
\pgfpathlineto{\pgfqpoint{8.372488in}{6.647491in}}%
\pgfpathlineto{\pgfqpoint{8.374371in}{6.641654in}}%
\pgfpathlineto{\pgfqpoint{8.378050in}{6.631377in}}%
\pgfpathlineto{\pgfqpoint{8.379162in}{6.632353in}}%
\pgfpathlineto{\pgfqpoint{8.381130in}{6.638483in}}%
\pgfpathlineto{\pgfqpoint{8.384638in}{6.647924in}}%
\pgfpathlineto{\pgfqpoint{8.385750in}{6.647015in}}%
\pgfpathlineto{\pgfqpoint{8.387718in}{6.641077in}}%
\pgfpathlineto{\pgfqpoint{8.391226in}{6.631844in}}%
\pgfpathlineto{\pgfqpoint{8.392339in}{6.632719in}}%
\pgfpathlineto{\pgfqpoint{8.394307in}{6.638496in}}%
\pgfpathlineto{\pgfqpoint{8.397815in}{6.647470in}}%
\pgfpathlineto{\pgfqpoint{8.398927in}{6.646603in}}%
\pgfpathlineto{\pgfqpoint{8.400895in}{6.640956in}}%
\pgfpathlineto{\pgfqpoint{8.404403in}{6.632286in}}%
\pgfpathlineto{\pgfqpoint{8.405515in}{6.633173in}}%
\pgfpathlineto{\pgfqpoint{8.407569in}{6.639017in}}%
\pgfpathlineto{\pgfqpoint{8.410991in}{6.647038in}}%
\pgfpathlineto{\pgfqpoint{8.412189in}{6.645958in}}%
\pgfpathlineto{\pgfqpoint{8.414414in}{6.639469in}}%
\pgfpathlineto{\pgfqpoint{8.417494in}{6.632718in}}%
\pgfpathlineto{\pgfqpoint{8.418692in}{6.633706in}}%
\pgfpathlineto{\pgfqpoint{8.420831in}{6.639680in}}%
\pgfpathlineto{\pgfqpoint{8.423997in}{6.646607in}}%
\pgfpathlineto{\pgfqpoint{8.425195in}{6.645682in}}%
\pgfpathlineto{\pgfqpoint{8.427334in}{6.639901in}}%
\pgfpathlineto{\pgfqpoint{8.430499in}{6.633139in}}%
\pgfpathlineto{\pgfqpoint{8.431697in}{6.634028in}}%
\pgfpathlineto{\pgfqpoint{8.433836in}{6.639642in}}%
\pgfpathlineto{\pgfqpoint{8.437002in}{6.646197in}}%
\pgfpathlineto{\pgfqpoint{8.438200in}{6.645319in}}%
\pgfpathlineto{\pgfqpoint{8.440425in}{6.639587in}}%
\pgfpathlineto{\pgfqpoint{8.443505in}{6.633538in}}%
\pgfpathlineto{\pgfqpoint{8.444788in}{6.634565in}}%
\pgfpathlineto{\pgfqpoint{8.447184in}{6.640788in}}%
\pgfpathlineto{\pgfqpoint{8.450008in}{6.645807in}}%
\pgfpathlineto{\pgfqpoint{8.451291in}{6.644747in}}%
\pgfpathlineto{\pgfqpoint{8.453858in}{6.638183in}}%
\pgfpathlineto{\pgfqpoint{8.456510in}{6.633923in}}%
\pgfpathlineto{\pgfqpoint{8.457794in}{6.635034in}}%
\pgfpathlineto{\pgfqpoint{8.460532in}{6.641905in}}%
\pgfpathlineto{\pgfqpoint{8.463013in}{6.645423in}}%
\pgfpathlineto{\pgfqpoint{8.464382in}{6.644096in}}%
\pgfpathlineto{\pgfqpoint{8.469773in}{6.634403in}}%
\pgfpathlineto{\pgfqpoint{8.470543in}{6.635168in}}%
\pgfpathlineto{\pgfqpoint{8.472938in}{6.640648in}}%
\pgfpathlineto{\pgfqpoint{8.475762in}{6.645057in}}%
\pgfpathlineto{\pgfqpoint{8.477131in}{6.643952in}}%
\pgfpathlineto{\pgfqpoint{8.480297in}{6.636719in}}%
\pgfpathlineto{\pgfqpoint{8.482521in}{6.634729in}}%
\pgfpathlineto{\pgfqpoint{8.483976in}{6.636503in}}%
\pgfpathlineto{\pgfqpoint{8.489024in}{6.644581in}}%
\pgfpathlineto{\pgfqpoint{8.489195in}{6.644473in}}%
\pgfpathlineto{\pgfqpoint{8.490907in}{6.641808in}}%
\pgfpathlineto{\pgfqpoint{8.495185in}{6.635040in}}%
\pgfpathlineto{\pgfqpoint{8.496639in}{6.636568in}}%
\pgfpathlineto{\pgfqpoint{8.501859in}{6.644192in}}%
\pgfpathlineto{\pgfqpoint{8.502115in}{6.644000in}}%
\pgfpathlineto{\pgfqpoint{8.504083in}{6.640703in}}%
\pgfpathlineto{\pgfqpoint{8.507677in}{6.635341in}}%
\pgfpathlineto{\pgfqpoint{8.509131in}{6.636473in}}%
\pgfpathlineto{\pgfqpoint{8.514607in}{6.643834in}}%
\pgfpathlineto{\pgfqpoint{8.515035in}{6.643471in}}%
\pgfpathlineto{\pgfqpoint{8.517431in}{6.639232in}}%
\pgfpathlineto{\pgfqpoint{8.520340in}{6.635662in}}%
\pgfpathlineto{\pgfqpoint{8.521880in}{6.636851in}}%
\pgfpathlineto{\pgfqpoint{8.527356in}{6.643455in}}%
\pgfpathlineto{\pgfqpoint{8.527527in}{6.643320in}}%
\pgfpathlineto{\pgfqpoint{8.529752in}{6.639840in}}%
\pgfpathlineto{\pgfqpoint{8.532918in}{6.635969in}}%
\pgfpathlineto{\pgfqpoint{8.534543in}{6.637159in}}%
\pgfpathlineto{\pgfqpoint{8.540019in}{6.643113in}}%
\pgfpathlineto{\pgfqpoint{8.540105in}{6.643047in}}%
\pgfpathlineto{\pgfqpoint{8.542330in}{6.639813in}}%
\pgfpathlineto{\pgfqpoint{8.545495in}{6.636264in}}%
\pgfpathlineto{\pgfqpoint{8.547121in}{6.637402in}}%
\pgfpathlineto{\pgfqpoint{8.552597in}{6.642811in}}%
\pgfpathlineto{\pgfqpoint{8.554907in}{6.639706in}}%
\pgfpathlineto{\pgfqpoint{8.558073in}{6.636549in}}%
\pgfpathlineto{\pgfqpoint{8.559784in}{6.637798in}}%
\pgfpathlineto{\pgfqpoint{8.564918in}{6.642655in}}%
\pgfpathlineto{\pgfqpoint{8.567057in}{6.640192in}}%
\pgfpathlineto{\pgfqpoint{8.570565in}{6.636822in}}%
\pgfpathlineto{\pgfqpoint{8.572362in}{6.638122in}}%
\pgfpathlineto{\pgfqpoint{8.577239in}{6.642459in}}%
\pgfpathlineto{\pgfqpoint{8.579378in}{6.640289in}}%
\pgfpathlineto{\pgfqpoint{8.582972in}{6.637077in}}%
\pgfpathlineto{\pgfqpoint{8.584854in}{6.638375in}}%
\pgfpathlineto{\pgfqpoint{8.589560in}{6.642243in}}%
\pgfpathlineto{\pgfqpoint{8.591785in}{6.640195in}}%
\pgfpathlineto{\pgfqpoint{8.595378in}{6.637326in}}%
\pgfpathlineto{\pgfqpoint{8.597432in}{6.638753in}}%
\pgfpathlineto{\pgfqpoint{8.601796in}{6.642043in}}%
\pgfpathlineto{\pgfqpoint{8.604020in}{6.640253in}}%
\pgfpathlineto{\pgfqpoint{8.607785in}{6.637571in}}%
\pgfpathlineto{\pgfqpoint{8.610010in}{6.639143in}}%
\pgfpathlineto{\pgfqpoint{8.613945in}{6.641853in}}%
\pgfpathlineto{\pgfqpoint{8.616256in}{6.640245in}}%
\pgfpathlineto{\pgfqpoint{8.620106in}{6.637798in}}%
\pgfpathlineto{\pgfqpoint{8.622502in}{6.639443in}}%
\pgfpathlineto{\pgfqpoint{8.626181in}{6.641643in}}%
\pgfpathlineto{\pgfqpoint{8.628577in}{6.640108in}}%
\pgfpathlineto{\pgfqpoint{8.632341in}{6.638006in}}%
\pgfpathlineto{\pgfqpoint{8.634908in}{6.639661in}}%
\pgfpathlineto{\pgfqpoint{8.638502in}{6.641421in}}%
\pgfpathlineto{\pgfqpoint{8.641240in}{6.639659in}}%
\pgfpathlineto{\pgfqpoint{8.644577in}{6.638210in}}%
\pgfpathlineto{\pgfqpoint{8.647400in}{6.639944in}}%
\pgfpathlineto{\pgfqpoint{8.650737in}{6.641218in}}%
\pgfpathlineto{\pgfqpoint{8.653817in}{6.639340in}}%
\pgfpathlineto{\pgfqpoint{8.656898in}{6.638432in}}%
\pgfpathlineto{\pgfqpoint{8.660320in}{6.640464in}}%
\pgfpathlineto{\pgfqpoint{8.663229in}{6.640942in}}%
\pgfpathlineto{\pgfqpoint{8.670844in}{6.639449in}}%
\pgfpathlineto{\pgfqpoint{8.675123in}{6.640838in}}%
\pgfpathlineto{\pgfqpoint{8.683422in}{6.639796in}}%
\pgfpathlineto{\pgfqpoint{8.687272in}{6.640663in}}%
\pgfpathlineto{\pgfqpoint{8.695230in}{6.639713in}}%
\pgfpathlineto{\pgfqpoint{8.699422in}{6.640492in}}%
\pgfpathlineto{\pgfqpoint{8.706952in}{6.639637in}}%
\pgfpathlineto{\pgfqpoint{8.711572in}{6.640327in}}%
\pgfpathlineto{\pgfqpoint{8.718674in}{6.639601in}}%
\pgfpathlineto{\pgfqpoint{8.723722in}{6.640171in}}%
\pgfpathlineto{\pgfqpoint{8.730481in}{6.639619in}}%
\pgfpathlineto{\pgfqpoint{8.735957in}{6.640010in}}%
\pgfpathlineto{\pgfqpoint{8.742289in}{6.639648in}}%
\pgfpathlineto{\pgfqpoint{8.748193in}{6.639873in}}%
\pgfpathlineto{\pgfqpoint{8.754353in}{6.639729in}}%
\pgfpathlineto{\pgfqpoint{8.760514in}{6.639756in}}%
\pgfpathlineto{\pgfqpoint{8.766845in}{6.639847in}}%
\pgfpathlineto{\pgfqpoint{8.773519in}{6.639617in}}%
\pgfpathlineto{\pgfqpoint{8.787038in}{6.639620in}}%
\pgfpathlineto{\pgfqpoint{8.845050in}{6.639768in}}%
\pgfpathlineto{\pgfqpoint{8.850354in}{6.639769in}}%
\pgfpathlineto{\pgfqpoint{8.850354in}{6.639769in}}%
\pgfusepath{stroke}%
\end{pgfscope}%
\begin{pgfscope}%
\pgfpathrectangle{\pgfqpoint{0.294194in}{5.519538in}}{\pgfqpoint{8.556160in}{2.240462in}}%
\pgfusepath{clip}%
\pgfsetrectcap%
\pgfsetroundjoin%
\pgfsetlinewidth{1.505625pt}%
\definecolor{currentstroke}{rgb}{1.000000,0.498039,0.054902}%
\pgfsetstrokecolor{currentstroke}%
\pgfsetdash{}{0pt}%
\pgfpathmoveto{\pgfqpoint{0.294194in}{6.639769in}}%
\pgfpathlineto{\pgfqpoint{0.297018in}{7.759850in}}%
\pgfpathlineto{\pgfqpoint{0.297103in}{7.759470in}}%
\pgfpathlineto{\pgfqpoint{0.297617in}{7.705493in}}%
\pgfpathlineto{\pgfqpoint{0.298900in}{7.226773in}}%
\pgfpathlineto{\pgfqpoint{0.302751in}{5.519560in}}%
\pgfpathlineto{\pgfqpoint{0.303093in}{5.538071in}}%
\pgfpathlineto{\pgfqpoint{0.304034in}{5.782309in}}%
\pgfpathlineto{\pgfqpoint{0.308483in}{7.759992in}}%
\pgfpathlineto{\pgfqpoint{0.309510in}{7.587509in}}%
\pgfpathlineto{\pgfqpoint{0.311478in}{6.561217in}}%
\pgfpathlineto{\pgfqpoint{0.314216in}{5.519554in}}%
\pgfpathlineto{\pgfqpoint{0.314558in}{5.538975in}}%
\pgfpathlineto{\pgfqpoint{0.315499in}{5.784193in}}%
\pgfpathlineto{\pgfqpoint{0.319949in}{7.759952in}}%
\pgfpathlineto{\pgfqpoint{0.320975in}{7.591983in}}%
\pgfpathlineto{\pgfqpoint{0.322943in}{6.573752in}}%
\pgfpathlineto{\pgfqpoint{0.325681in}{5.519741in}}%
\pgfpathlineto{\pgfqpoint{0.325767in}{5.520063in}}%
\pgfpathlineto{\pgfqpoint{0.326195in}{5.558072in}}%
\pgfpathlineto{\pgfqpoint{0.327307in}{5.917165in}}%
\pgfpathlineto{\pgfqpoint{0.331500in}{7.759848in}}%
\pgfpathlineto{\pgfqpoint{0.332099in}{7.695998in}}%
\pgfpathlineto{\pgfqpoint{0.333468in}{7.162823in}}%
\pgfpathlineto{\pgfqpoint{0.337232in}{5.519646in}}%
\pgfpathlineto{\pgfqpoint{0.337489in}{5.529260in}}%
\pgfpathlineto{\pgfqpoint{0.338345in}{5.713090in}}%
\pgfpathlineto{\pgfqpoint{0.340484in}{6.849361in}}%
\pgfpathlineto{\pgfqpoint{0.343051in}{7.759837in}}%
\pgfpathlineto{\pgfqpoint{0.343393in}{7.738203in}}%
\pgfpathlineto{\pgfqpoint{0.344420in}{7.455978in}}%
\pgfpathlineto{\pgfqpoint{0.348783in}{5.519893in}}%
\pgfpathlineto{\pgfqpoint{0.349724in}{5.651012in}}%
\pgfpathlineto{\pgfqpoint{0.351436in}{6.467113in}}%
\pgfpathlineto{\pgfqpoint{0.354601in}{7.759781in}}%
\pgfpathlineto{\pgfqpoint{0.354687in}{7.759207in}}%
\pgfpathlineto{\pgfqpoint{0.355200in}{7.705815in}}%
\pgfpathlineto{\pgfqpoint{0.356484in}{7.239311in}}%
\pgfpathlineto{\pgfqpoint{0.360420in}{5.519747in}}%
\pgfpathlineto{\pgfqpoint{0.360848in}{5.547246in}}%
\pgfpathlineto{\pgfqpoint{0.361874in}{5.840794in}}%
\pgfpathlineto{\pgfqpoint{0.366238in}{7.759736in}}%
\pgfpathlineto{\pgfqpoint{0.367094in}{7.648199in}}%
\pgfpathlineto{\pgfqpoint{0.368719in}{6.911020in}}%
\pgfpathlineto{\pgfqpoint{0.372056in}{5.519963in}}%
\pgfpathlineto{\pgfqpoint{0.372142in}{5.520187in}}%
\pgfpathlineto{\pgfqpoint{0.372570in}{5.556716in}}%
\pgfpathlineto{\pgfqpoint{0.373682in}{5.905641in}}%
\pgfpathlineto{\pgfqpoint{0.377960in}{7.759623in}}%
\pgfpathlineto{\pgfqpoint{0.378559in}{7.697160in}}%
\pgfpathlineto{\pgfqpoint{0.379928in}{7.176497in}}%
\pgfpathlineto{\pgfqpoint{0.383778in}{5.519853in}}%
\pgfpathlineto{\pgfqpoint{0.384035in}{5.529743in}}%
\pgfpathlineto{\pgfqpoint{0.384891in}{5.710314in}}%
\pgfpathlineto{\pgfqpoint{0.387030in}{6.826343in}}%
\pgfpathlineto{\pgfqpoint{0.389682in}{7.759493in}}%
\pgfpathlineto{\pgfqpoint{0.390110in}{7.725927in}}%
\pgfpathlineto{\pgfqpoint{0.391222in}{7.386160in}}%
\pgfpathlineto{\pgfqpoint{0.395500in}{5.520016in}}%
\pgfpathlineto{\pgfqpoint{0.396185in}{5.588998in}}%
\pgfpathlineto{\pgfqpoint{0.397554in}{6.116216in}}%
\pgfpathlineto{\pgfqpoint{0.401404in}{7.759563in}}%
\pgfpathlineto{\pgfqpoint{0.401661in}{7.749427in}}%
\pgfpathlineto{\pgfqpoint{0.402516in}{7.569642in}}%
\pgfpathlineto{\pgfqpoint{0.404655in}{6.460887in}}%
\pgfpathlineto{\pgfqpoint{0.407308in}{5.520025in}}%
\pgfpathlineto{\pgfqpoint{0.407650in}{5.539083in}}%
\pgfpathlineto{\pgfqpoint{0.408591in}{5.773185in}}%
\pgfpathlineto{\pgfqpoint{0.411500in}{7.328684in}}%
\pgfpathlineto{\pgfqpoint{0.413212in}{7.759464in}}%
\pgfpathlineto{\pgfqpoint{0.413468in}{7.748798in}}%
\pgfpathlineto{\pgfqpoint{0.414324in}{7.568406in}}%
\pgfpathlineto{\pgfqpoint{0.416463in}{6.463379in}}%
\pgfpathlineto{\pgfqpoint{0.419116in}{5.520140in}}%
\pgfpathlineto{\pgfqpoint{0.419372in}{5.529784in}}%
\pgfpathlineto{\pgfqpoint{0.420228in}{5.706494in}}%
\pgfpathlineto{\pgfqpoint{0.422367in}{6.805562in}}%
\pgfpathlineto{\pgfqpoint{0.425019in}{7.759195in}}%
\pgfpathlineto{\pgfqpoint{0.425105in}{7.758900in}}%
\pgfpathlineto{\pgfqpoint{0.425533in}{7.723145in}}%
\pgfpathlineto{\pgfqpoint{0.426645in}{7.383788in}}%
\pgfpathlineto{\pgfqpoint{0.431009in}{5.520308in}}%
\pgfpathlineto{\pgfqpoint{0.431693in}{5.597007in}}%
\pgfpathlineto{\pgfqpoint{0.433148in}{6.173340in}}%
\pgfpathlineto{\pgfqpoint{0.436913in}{7.759139in}}%
\pgfpathlineto{\pgfqpoint{0.437084in}{7.755870in}}%
\pgfpathlineto{\pgfqpoint{0.437768in}{7.653248in}}%
\pgfpathlineto{\pgfqpoint{0.439394in}{6.943490in}}%
\pgfpathlineto{\pgfqpoint{0.442902in}{5.520378in}}%
\pgfpathlineto{\pgfqpoint{0.442987in}{5.521676in}}%
\pgfpathlineto{\pgfqpoint{0.443501in}{5.576720in}}%
\pgfpathlineto{\pgfqpoint{0.444784in}{6.029673in}}%
\pgfpathlineto{\pgfqpoint{0.448891in}{7.758997in}}%
\pgfpathlineto{\pgfqpoint{0.449234in}{7.738245in}}%
\pgfpathlineto{\pgfqpoint{0.450260in}{7.471629in}}%
\pgfpathlineto{\pgfqpoint{0.454881in}{5.520690in}}%
\pgfpathlineto{\pgfqpoint{0.455822in}{5.663842in}}%
\pgfpathlineto{\pgfqpoint{0.457619in}{6.509365in}}%
\pgfpathlineto{\pgfqpoint{0.460784in}{7.758500in}}%
\pgfpathlineto{\pgfqpoint{0.460870in}{7.758801in}}%
\pgfpathlineto{\pgfqpoint{0.460956in}{7.756855in}}%
\pgfpathlineto{\pgfqpoint{0.461555in}{7.681210in}}%
\pgfpathlineto{\pgfqpoint{0.463009in}{7.110695in}}%
\pgfpathlineto{\pgfqpoint{0.466859in}{5.520711in}}%
\pgfpathlineto{\pgfqpoint{0.467031in}{5.526109in}}%
\pgfpathlineto{\pgfqpoint{0.467801in}{5.658331in}}%
\pgfpathlineto{\pgfqpoint{0.469597in}{6.493978in}}%
\pgfpathlineto{\pgfqpoint{0.472849in}{7.758787in}}%
\pgfpathlineto{\pgfqpoint{0.473105in}{7.749323in}}%
\pgfpathlineto{\pgfqpoint{0.473961in}{7.577849in}}%
\pgfpathlineto{\pgfqpoint{0.476015in}{6.553928in}}%
\pgfpathlineto{\pgfqpoint{0.478838in}{5.521109in}}%
\pgfpathlineto{\pgfqpoint{0.478924in}{5.521098in}}%
\pgfpathlineto{\pgfqpoint{0.479352in}{5.554209in}}%
\pgfpathlineto{\pgfqpoint{0.480464in}{5.879765in}}%
\pgfpathlineto{\pgfqpoint{0.484913in}{7.758623in}}%
\pgfpathlineto{\pgfqpoint{0.485598in}{7.689534in}}%
\pgfpathlineto{\pgfqpoint{0.486967in}{7.182114in}}%
\pgfpathlineto{\pgfqpoint{0.490988in}{5.521072in}}%
\pgfpathlineto{\pgfqpoint{0.491245in}{5.532674in}}%
\pgfpathlineto{\pgfqpoint{0.492100in}{5.709219in}}%
\pgfpathlineto{\pgfqpoint{0.494239in}{6.781483in}}%
\pgfpathlineto{\pgfqpoint{0.496977in}{7.758100in}}%
\pgfpathlineto{\pgfqpoint{0.497063in}{7.758237in}}%
\pgfpathlineto{\pgfqpoint{0.497491in}{7.726120in}}%
\pgfpathlineto{\pgfqpoint{0.498603in}{7.405445in}}%
\pgfpathlineto{\pgfqpoint{0.503138in}{5.521443in}}%
\pgfpathlineto{\pgfqpoint{0.503822in}{5.599048in}}%
\pgfpathlineto{\pgfqpoint{0.505277in}{6.159449in}}%
\pgfpathlineto{\pgfqpoint{0.509213in}{7.758083in}}%
\pgfpathlineto{\pgfqpoint{0.509384in}{7.751991in}}%
\pgfpathlineto{\pgfqpoint{0.510154in}{7.619624in}}%
\pgfpathlineto{\pgfqpoint{0.511951in}{6.796951in}}%
\pgfpathlineto{\pgfqpoint{0.515288in}{5.521416in}}%
\pgfpathlineto{\pgfqpoint{0.515630in}{5.540263in}}%
\pgfpathlineto{\pgfqpoint{0.516571in}{5.762198in}}%
\pgfpathlineto{\pgfqpoint{0.519224in}{7.139022in}}%
\pgfpathlineto{\pgfqpoint{0.521363in}{7.758016in}}%
\pgfpathlineto{\pgfqpoint{0.521534in}{7.754482in}}%
\pgfpathlineto{\pgfqpoint{0.522218in}{7.655439in}}%
\pgfpathlineto{\pgfqpoint{0.523844in}{6.974331in}}%
\pgfpathlineto{\pgfqpoint{0.527523in}{5.521703in}}%
\pgfpathlineto{\pgfqpoint{0.527951in}{5.551962in}}%
\pgfpathlineto{\pgfqpoint{0.529063in}{5.864021in}}%
\pgfpathlineto{\pgfqpoint{0.533598in}{7.757733in}}%
\pgfpathlineto{\pgfqpoint{0.534368in}{7.677816in}}%
\pgfpathlineto{\pgfqpoint{0.535823in}{7.121731in}}%
\pgfpathlineto{\pgfqpoint{0.539759in}{5.521821in}}%
\pgfpathlineto{\pgfqpoint{0.539930in}{5.525738in}}%
\pgfpathlineto{\pgfqpoint{0.540614in}{5.625270in}}%
\pgfpathlineto{\pgfqpoint{0.542240in}{6.302152in}}%
\pgfpathlineto{\pgfqpoint{0.545919in}{7.757609in}}%
\pgfpathlineto{\pgfqpoint{0.546005in}{7.756532in}}%
\pgfpathlineto{\pgfqpoint{0.546518in}{7.705767in}}%
\pgfpathlineto{\pgfqpoint{0.547801in}{7.280734in}}%
\pgfpathlineto{\pgfqpoint{0.552080in}{5.522049in}}%
\pgfpathlineto{\pgfqpoint{0.552507in}{5.547979in}}%
\pgfpathlineto{\pgfqpoint{0.553534in}{5.813151in}}%
\pgfpathlineto{\pgfqpoint{0.558240in}{7.757310in}}%
\pgfpathlineto{\pgfqpoint{0.559181in}{7.637264in}}%
\pgfpathlineto{\pgfqpoint{0.560893in}{6.898847in}}%
\pgfpathlineto{\pgfqpoint{0.564401in}{5.522635in}}%
\pgfpathlineto{\pgfqpoint{0.564486in}{5.522488in}}%
\pgfpathlineto{\pgfqpoint{0.564572in}{5.524444in}}%
\pgfpathlineto{\pgfqpoint{0.565171in}{5.596203in}}%
\pgfpathlineto{\pgfqpoint{0.566625in}{6.135883in}}%
\pgfpathlineto{\pgfqpoint{0.570647in}{7.757115in}}%
\pgfpathlineto{\pgfqpoint{0.570818in}{7.753117in}}%
\pgfpathlineto{\pgfqpoint{0.571502in}{7.654929in}}%
\pgfpathlineto{\pgfqpoint{0.573128in}{6.988556in}}%
\pgfpathlineto{\pgfqpoint{0.576893in}{5.522670in}}%
\pgfpathlineto{\pgfqpoint{0.577321in}{5.552121in}}%
\pgfpathlineto{\pgfqpoint{0.578433in}{5.855197in}}%
\pgfpathlineto{\pgfqpoint{0.583053in}{7.756662in}}%
\pgfpathlineto{\pgfqpoint{0.583909in}{7.662998in}}%
\pgfpathlineto{\pgfqpoint{0.585449in}{7.056162in}}%
\pgfpathlineto{\pgfqpoint{0.589299in}{5.522940in}}%
\pgfpathlineto{\pgfqpoint{0.589385in}{5.523281in}}%
\pgfpathlineto{\pgfqpoint{0.589898in}{5.568472in}}%
\pgfpathlineto{\pgfqpoint{0.591182in}{5.974107in}}%
\pgfpathlineto{\pgfqpoint{0.595545in}{7.756419in}}%
\pgfpathlineto{\pgfqpoint{0.596059in}{7.724282in}}%
\pgfpathlineto{\pgfqpoint{0.597171in}{7.418366in}}%
\pgfpathlineto{\pgfqpoint{0.601877in}{5.523282in}}%
\pgfpathlineto{\pgfqpoint{0.602647in}{5.612771in}}%
\pgfpathlineto{\pgfqpoint{0.604187in}{6.209052in}}%
\pgfpathlineto{\pgfqpoint{0.608123in}{7.756279in}}%
\pgfpathlineto{\pgfqpoint{0.608209in}{7.755126in}}%
\pgfpathlineto{\pgfqpoint{0.608722in}{7.705673in}}%
\pgfpathlineto{\pgfqpoint{0.610005in}{7.294841in}}%
\pgfpathlineto{\pgfqpoint{0.614369in}{5.523594in}}%
\pgfpathlineto{\pgfqpoint{0.614882in}{5.554817in}}%
\pgfpathlineto{\pgfqpoint{0.615995in}{5.856179in}}%
\pgfpathlineto{\pgfqpoint{0.620701in}{7.755974in}}%
\pgfpathlineto{\pgfqpoint{0.621471in}{7.676246in}}%
\pgfpathlineto{\pgfqpoint{0.622925in}{7.141736in}}%
\pgfpathlineto{\pgfqpoint{0.627032in}{5.523732in}}%
\pgfpathlineto{\pgfqpoint{0.627203in}{5.528290in}}%
\pgfpathlineto{\pgfqpoint{0.627888in}{5.625659in}}%
\pgfpathlineto{\pgfqpoint{0.629514in}{6.275511in}}%
\pgfpathlineto{\pgfqpoint{0.633364in}{7.755629in}}%
\pgfpathlineto{\pgfqpoint{0.633706in}{7.738115in}}%
\pgfpathlineto{\pgfqpoint{0.634647in}{7.531944in}}%
\pgfpathlineto{\pgfqpoint{0.637129in}{6.306542in}}%
\pgfpathlineto{\pgfqpoint{0.639696in}{5.524046in}}%
\pgfpathlineto{\pgfqpoint{0.640038in}{5.540713in}}%
\pgfpathlineto{\pgfqpoint{0.640979in}{5.744096in}}%
\pgfpathlineto{\pgfqpoint{0.643375in}{6.919037in}}%
\pgfpathlineto{\pgfqpoint{0.646027in}{7.755306in}}%
\pgfpathlineto{\pgfqpoint{0.646284in}{7.747320in}}%
\pgfpathlineto{\pgfqpoint{0.647139in}{7.595157in}}%
\pgfpathlineto{\pgfqpoint{0.649107in}{6.703136in}}%
\pgfpathlineto{\pgfqpoint{0.652359in}{5.524693in}}%
\pgfpathlineto{\pgfqpoint{0.652444in}{5.524563in}}%
\pgfpathlineto{\pgfqpoint{0.652530in}{5.526416in}}%
\pgfpathlineto{\pgfqpoint{0.653129in}{5.594201in}}%
\pgfpathlineto{\pgfqpoint{0.654583in}{6.106611in}}%
\pgfpathlineto{\pgfqpoint{0.658776in}{7.754974in}}%
\pgfpathlineto{\pgfqpoint{0.659033in}{7.746824in}}%
\pgfpathlineto{\pgfqpoint{0.659888in}{7.595208in}}%
\pgfpathlineto{\pgfqpoint{0.661856in}{6.708265in}}%
\pgfpathlineto{\pgfqpoint{0.665108in}{5.525415in}}%
\pgfpathlineto{\pgfqpoint{0.665193in}{5.524753in}}%
\pgfpathlineto{\pgfqpoint{0.665364in}{5.529328in}}%
\pgfpathlineto{\pgfqpoint{0.666049in}{5.624778in}}%
\pgfpathlineto{\pgfqpoint{0.667674in}{6.261293in}}%
\pgfpathlineto{\pgfqpoint{0.671610in}{7.754546in}}%
\pgfpathlineto{\pgfqpoint{0.672038in}{7.727294in}}%
\pgfpathlineto{\pgfqpoint{0.673150in}{7.442744in}}%
\pgfpathlineto{\pgfqpoint{0.678028in}{5.525161in}}%
\pgfpathlineto{\pgfqpoint{0.678883in}{5.626301in}}%
\pgfpathlineto{\pgfqpoint{0.680509in}{6.260610in}}%
\pgfpathlineto{\pgfqpoint{0.684445in}{7.754253in}}%
\pgfpathlineto{\pgfqpoint{0.684787in}{7.738113in}}%
\pgfpathlineto{\pgfqpoint{0.685728in}{7.540668in}}%
\pgfpathlineto{\pgfqpoint{0.688124in}{6.389620in}}%
\pgfpathlineto{\pgfqpoint{0.690862in}{5.525535in}}%
\pgfpathlineto{\pgfqpoint{0.691119in}{5.532716in}}%
\pgfpathlineto{\pgfqpoint{0.691889in}{5.656011in}}%
\pgfpathlineto{\pgfqpoint{0.693685in}{6.409915in}}%
\pgfpathlineto{\pgfqpoint{0.697279in}{7.753313in}}%
\pgfpathlineto{\pgfqpoint{0.697365in}{7.753821in}}%
\pgfpathlineto{\pgfqpoint{0.697536in}{7.749067in}}%
\pgfpathlineto{\pgfqpoint{0.698306in}{7.634611in}}%
\pgfpathlineto{\pgfqpoint{0.700017in}{6.940268in}}%
\pgfpathlineto{\pgfqpoint{0.703782in}{5.526037in}}%
\pgfpathlineto{\pgfqpoint{0.703953in}{5.528226in}}%
\pgfpathlineto{\pgfqpoint{0.704552in}{5.595493in}}%
\pgfpathlineto{\pgfqpoint{0.706006in}{6.095383in}}%
\pgfpathlineto{\pgfqpoint{0.710285in}{7.753413in}}%
\pgfpathlineto{\pgfqpoint{0.710541in}{7.746299in}}%
\pgfpathlineto{\pgfqpoint{0.711311in}{7.624525in}}%
\pgfpathlineto{\pgfqpoint{0.713108in}{6.878759in}}%
\pgfpathlineto{\pgfqpoint{0.716702in}{5.527740in}}%
\pgfpathlineto{\pgfqpoint{0.716787in}{5.526325in}}%
\pgfpathlineto{\pgfqpoint{0.717044in}{5.533466in}}%
\pgfpathlineto{\pgfqpoint{0.717814in}{5.654873in}}%
\pgfpathlineto{\pgfqpoint{0.719611in}{6.398120in}}%
\pgfpathlineto{\pgfqpoint{0.723290in}{7.752911in}}%
\pgfpathlineto{\pgfqpoint{0.723461in}{7.750639in}}%
\pgfpathlineto{\pgfqpoint{0.724060in}{7.683875in}}%
\pgfpathlineto{\pgfqpoint{0.725515in}{7.189274in}}%
\pgfpathlineto{\pgfqpoint{0.729878in}{5.526764in}}%
\pgfpathlineto{\pgfqpoint{0.730050in}{5.531588in}}%
\pgfpathlineto{\pgfqpoint{0.730820in}{5.644305in}}%
\pgfpathlineto{\pgfqpoint{0.732531in}{6.326734in}}%
\pgfpathlineto{\pgfqpoint{0.736381in}{7.752595in}}%
\pgfpathlineto{\pgfqpoint{0.736638in}{7.745308in}}%
\pgfpathlineto{\pgfqpoint{0.737408in}{7.624816in}}%
\pgfpathlineto{\pgfqpoint{0.739205in}{6.888881in}}%
\pgfpathlineto{\pgfqpoint{0.742884in}{5.527804in}}%
\pgfpathlineto{\pgfqpoint{0.742969in}{5.527145in}}%
\pgfpathlineto{\pgfqpoint{0.743141in}{5.531418in}}%
\pgfpathlineto{\pgfqpoint{0.743825in}{5.621721in}}%
\pgfpathlineto{\pgfqpoint{0.745451in}{6.229077in}}%
\pgfpathlineto{\pgfqpoint{0.749558in}{7.752091in}}%
\pgfpathlineto{\pgfqpoint{0.749986in}{7.725961in}}%
\pgfpathlineto{\pgfqpoint{0.751098in}{7.454885in}}%
\pgfpathlineto{\pgfqpoint{0.756146in}{5.527681in}}%
\pgfpathlineto{\pgfqpoint{0.757087in}{5.644209in}}%
\pgfpathlineto{\pgfqpoint{0.758798in}{6.317431in}}%
\pgfpathlineto{\pgfqpoint{0.762649in}{7.751113in}}%
\pgfpathlineto{\pgfqpoint{0.762734in}{7.751703in}}%
\pgfpathlineto{\pgfqpoint{0.762905in}{7.747368in}}%
\pgfpathlineto{\pgfqpoint{0.763590in}{7.657822in}}%
\pgfpathlineto{\pgfqpoint{0.765216in}{7.056947in}}%
\pgfpathlineto{\pgfqpoint{0.769323in}{5.528061in}}%
\pgfpathlineto{\pgfqpoint{0.769408in}{5.528700in}}%
\pgfpathlineto{\pgfqpoint{0.769922in}{5.570698in}}%
\pgfpathlineto{\pgfqpoint{0.771205in}{5.936891in}}%
\pgfpathlineto{\pgfqpoint{0.775997in}{7.751150in}}%
\pgfpathlineto{\pgfqpoint{0.776510in}{7.714765in}}%
\pgfpathlineto{\pgfqpoint{0.777708in}{7.395126in}}%
\pgfpathlineto{\pgfqpoint{0.782585in}{5.528619in}}%
\pgfpathlineto{\pgfqpoint{0.783355in}{5.595995in}}%
\pgfpathlineto{\pgfqpoint{0.784809in}{6.074074in}}%
\pgfpathlineto{\pgfqpoint{0.789259in}{7.750742in}}%
\pgfpathlineto{\pgfqpoint{0.789515in}{7.743596in}}%
\pgfpathlineto{\pgfqpoint{0.790285in}{7.627056in}}%
\pgfpathlineto{\pgfqpoint{0.792082in}{6.913001in}}%
\pgfpathlineto{\pgfqpoint{0.795847in}{5.530320in}}%
\pgfpathlineto{\pgfqpoint{0.795933in}{5.529054in}}%
\pgfpathlineto{\pgfqpoint{0.796189in}{5.536029in}}%
\pgfpathlineto{\pgfqpoint{0.796959in}{5.651634in}}%
\pgfpathlineto{\pgfqpoint{0.798756in}{6.362282in}}%
\pgfpathlineto{\pgfqpoint{0.802521in}{7.748499in}}%
\pgfpathlineto{\pgfqpoint{0.802606in}{7.750120in}}%
\pgfpathlineto{\pgfqpoint{0.802949in}{7.738745in}}%
\pgfpathlineto{\pgfqpoint{0.803890in}{7.565327in}}%
\pgfpathlineto{\pgfqpoint{0.806029in}{6.620913in}}%
\pgfpathlineto{\pgfqpoint{0.809195in}{5.532368in}}%
\pgfpathlineto{\pgfqpoint{0.809366in}{5.529569in}}%
\pgfpathlineto{\pgfqpoint{0.809708in}{5.545284in}}%
\pgfpathlineto{\pgfqpoint{0.810649in}{5.729084in}}%
\pgfpathlineto{\pgfqpoint{0.812960in}{6.770610in}}%
\pgfpathlineto{\pgfqpoint{0.815954in}{7.748145in}}%
\pgfpathlineto{\pgfqpoint{0.816040in}{7.749643in}}%
\pgfpathlineto{\pgfqpoint{0.816382in}{7.737946in}}%
\pgfpathlineto{\pgfqpoint{0.817323in}{7.565034in}}%
\pgfpathlineto{\pgfqpoint{0.819462in}{6.626243in}}%
\pgfpathlineto{\pgfqpoint{0.822628in}{5.533883in}}%
\pgfpathlineto{\pgfqpoint{0.822799in}{5.530062in}}%
\pgfpathlineto{\pgfqpoint{0.823141in}{5.543536in}}%
\pgfpathlineto{\pgfqpoint{0.824083in}{5.720314in}}%
\pgfpathlineto{\pgfqpoint{0.826307in}{6.704922in}}%
\pgfpathlineto{\pgfqpoint{0.829388in}{7.745690in}}%
\pgfpathlineto{\pgfqpoint{0.829559in}{7.749213in}}%
\pgfpathlineto{\pgfqpoint{0.829901in}{7.735247in}}%
\pgfpathlineto{\pgfqpoint{0.830842in}{7.557878in}}%
\pgfpathlineto{\pgfqpoint{0.833067in}{6.575278in}}%
\pgfpathlineto{\pgfqpoint{0.836147in}{5.534473in}}%
\pgfpathlineto{\pgfqpoint{0.836318in}{5.530608in}}%
\pgfpathlineto{\pgfqpoint{0.836660in}{5.543794in}}%
\pgfpathlineto{\pgfqpoint{0.837602in}{5.718509in}}%
\pgfpathlineto{\pgfqpoint{0.839826in}{6.695573in}}%
\pgfpathlineto{\pgfqpoint{0.842906in}{7.743690in}}%
\pgfpathlineto{\pgfqpoint{0.843078in}{7.748528in}}%
\pgfpathlineto{\pgfqpoint{0.843505in}{7.730288in}}%
\pgfpathlineto{\pgfqpoint{0.844532in}{7.517777in}}%
\pgfpathlineto{\pgfqpoint{0.847099in}{6.341446in}}%
\pgfpathlineto{\pgfqpoint{0.849837in}{5.531803in}}%
\pgfpathlineto{\pgfqpoint{0.849923in}{5.531173in}}%
\pgfpathlineto{\pgfqpoint{0.850094in}{5.535096in}}%
\pgfpathlineto{\pgfqpoint{0.850778in}{5.618730in}}%
\pgfpathlineto{\pgfqpoint{0.852404in}{6.186805in}}%
\pgfpathlineto{\pgfqpoint{0.856768in}{7.747898in}}%
\pgfpathlineto{\pgfqpoint{0.857195in}{7.722387in}}%
\pgfpathlineto{\pgfqpoint{0.858308in}{7.467351in}}%
\pgfpathlineto{\pgfqpoint{0.863527in}{5.531847in}}%
\pgfpathlineto{\pgfqpoint{0.864554in}{5.645520in}}%
\pgfpathlineto{\pgfqpoint{0.866351in}{6.324083in}}%
\pgfpathlineto{\pgfqpoint{0.870286in}{7.745848in}}%
\pgfpathlineto{\pgfqpoint{0.870372in}{7.747376in}}%
\pgfpathlineto{\pgfqpoint{0.870714in}{7.736481in}}%
\pgfpathlineto{\pgfqpoint{0.871655in}{7.570946in}}%
\pgfpathlineto{\pgfqpoint{0.873794in}{6.660893in}}%
\pgfpathlineto{\pgfqpoint{0.877046in}{5.538134in}}%
\pgfpathlineto{\pgfqpoint{0.877302in}{5.532448in}}%
\pgfpathlineto{\pgfqpoint{0.877645in}{5.548539in}}%
\pgfpathlineto{\pgfqpoint{0.878671in}{5.751856in}}%
\pgfpathlineto{\pgfqpoint{0.881153in}{6.862973in}}%
\pgfpathlineto{\pgfqpoint{0.883976in}{7.743430in}}%
\pgfpathlineto{\pgfqpoint{0.884147in}{7.746909in}}%
\pgfpathlineto{\pgfqpoint{0.884490in}{7.733670in}}%
\pgfpathlineto{\pgfqpoint{0.885431in}{7.563466in}}%
\pgfpathlineto{\pgfqpoint{0.887656in}{6.610137in}}%
\pgfpathlineto{\pgfqpoint{0.890821in}{5.538764in}}%
\pgfpathlineto{\pgfqpoint{0.891078in}{5.533045in}}%
\pgfpathlineto{\pgfqpoint{0.891420in}{5.548855in}}%
\pgfpathlineto{\pgfqpoint{0.892447in}{5.749874in}}%
\pgfpathlineto{\pgfqpoint{0.894928in}{6.852927in}}%
\pgfpathlineto{\pgfqpoint{0.897837in}{7.744573in}}%
\pgfpathlineto{\pgfqpoint{0.897923in}{7.746134in}}%
\pgfpathlineto{\pgfqpoint{0.898265in}{7.735711in}}%
\pgfpathlineto{\pgfqpoint{0.899206in}{7.574048in}}%
\pgfpathlineto{\pgfqpoint{0.901346in}{6.678826in}}%
\pgfpathlineto{\pgfqpoint{0.904682in}{5.538455in}}%
\pgfpathlineto{\pgfqpoint{0.904854in}{5.533713in}}%
\pgfpathlineto{\pgfqpoint{0.905281in}{5.550859in}}%
\pgfpathlineto{\pgfqpoint{0.906308in}{5.753688in}}%
\pgfpathlineto{\pgfqpoint{0.908789in}{6.852357in}}%
\pgfpathlineto{\pgfqpoint{0.911699in}{7.743497in}}%
\pgfpathlineto{\pgfqpoint{0.911870in}{7.745538in}}%
\pgfpathlineto{\pgfqpoint{0.912126in}{7.736226in}}%
\pgfpathlineto{\pgfqpoint{0.912982in}{7.601084in}}%
\pgfpathlineto{\pgfqpoint{0.914950in}{6.822005in}}%
\pgfpathlineto{\pgfqpoint{0.918629in}{5.537503in}}%
\pgfpathlineto{\pgfqpoint{0.918800in}{5.534204in}}%
\pgfpathlineto{\pgfqpoint{0.919143in}{5.547294in}}%
\pgfpathlineto{\pgfqpoint{0.920084in}{5.713796in}}%
\pgfpathlineto{\pgfqpoint{0.922223in}{6.606810in}}%
\pgfpathlineto{\pgfqpoint{0.925560in}{7.739530in}}%
\pgfpathlineto{\pgfqpoint{0.925816in}{7.744868in}}%
\pgfpathlineto{\pgfqpoint{0.926159in}{7.729145in}}%
\pgfpathlineto{\pgfqpoint{0.927185in}{7.532191in}}%
\pgfpathlineto{\pgfqpoint{0.929581in}{6.489943in}}%
\pgfpathlineto{\pgfqpoint{0.932576in}{5.539446in}}%
\pgfpathlineto{\pgfqpoint{0.932747in}{5.534978in}}%
\pgfpathlineto{\pgfqpoint{0.933175in}{5.552215in}}%
\pgfpathlineto{\pgfqpoint{0.934202in}{5.752002in}}%
\pgfpathlineto{\pgfqpoint{0.936683in}{6.835770in}}%
\pgfpathlineto{\pgfqpoint{0.939592in}{7.739831in}}%
\pgfpathlineto{\pgfqpoint{0.939763in}{7.744232in}}%
\pgfpathlineto{\pgfqpoint{0.940191in}{7.726975in}}%
\pgfpathlineto{\pgfqpoint{0.941218in}{7.527954in}}%
\pgfpathlineto{\pgfqpoint{0.943699in}{6.447944in}}%
\pgfpathlineto{\pgfqpoint{0.946608in}{5.540770in}}%
\pgfpathlineto{\pgfqpoint{0.946865in}{5.535706in}}%
\pgfpathlineto{\pgfqpoint{0.947207in}{5.551439in}}%
\pgfpathlineto{\pgfqpoint{0.948234in}{5.746125in}}%
\pgfpathlineto{\pgfqpoint{0.950630in}{6.777340in}}%
\pgfpathlineto{\pgfqpoint{0.953710in}{7.740629in}}%
\pgfpathlineto{\pgfqpoint{0.953881in}{7.743647in}}%
\pgfpathlineto{\pgfqpoint{0.954223in}{7.730502in}}%
\pgfpathlineto{\pgfqpoint{0.955164in}{7.567178in}}%
\pgfpathlineto{\pgfqpoint{0.957303in}{6.690926in}}%
\pgfpathlineto{\pgfqpoint{0.960726in}{5.541378in}}%
\pgfpathlineto{\pgfqpoint{0.960983in}{5.536407in}}%
\pgfpathlineto{\pgfqpoint{0.961325in}{5.552026in}}%
\pgfpathlineto{\pgfqpoint{0.962352in}{5.744873in}}%
\pgfpathlineto{\pgfqpoint{0.964747in}{6.768359in}}%
\pgfpathlineto{\pgfqpoint{0.967828in}{7.738584in}}%
\pgfpathlineto{\pgfqpoint{0.967999in}{7.742849in}}%
\pgfpathlineto{\pgfqpoint{0.968427in}{7.725844in}}%
\pgfpathlineto{\pgfqpoint{0.969453in}{7.530617in}}%
\pgfpathlineto{\pgfqpoint{0.971849in}{6.507791in}}%
\pgfpathlineto{\pgfqpoint{0.974929in}{5.541341in}}%
\pgfpathlineto{\pgfqpoint{0.975100in}{5.537056in}}%
\pgfpathlineto{\pgfqpoint{0.975528in}{5.553862in}}%
\pgfpathlineto{\pgfqpoint{0.976555in}{5.747842in}}%
\pgfpathlineto{\pgfqpoint{0.978951in}{6.766473in}}%
\pgfpathlineto{\pgfqpoint{0.982031in}{7.736991in}}%
\pgfpathlineto{\pgfqpoint{0.982288in}{7.742085in}}%
\pgfpathlineto{\pgfqpoint{0.982630in}{7.726981in}}%
\pgfpathlineto{\pgfqpoint{0.983657in}{7.537849in}}%
\pgfpathlineto{\pgfqpoint{0.986052in}{6.527000in}}%
\pgfpathlineto{\pgfqpoint{0.989133in}{5.544783in}}%
\pgfpathlineto{\pgfqpoint{0.989389in}{5.537687in}}%
\pgfpathlineto{\pgfqpoint{0.989817in}{5.556960in}}%
\pgfpathlineto{\pgfqpoint{0.990844in}{5.754738in}}%
\pgfpathlineto{\pgfqpoint{0.993325in}{6.812074in}}%
\pgfpathlineto{\pgfqpoint{0.996320in}{7.735971in}}%
\pgfpathlineto{\pgfqpoint{0.996577in}{7.741392in}}%
\pgfpathlineto{\pgfqpoint{0.997004in}{7.719508in}}%
\pgfpathlineto{\pgfqpoint{0.998117in}{7.491228in}}%
\pgfpathlineto{\pgfqpoint{1.000940in}{6.258884in}}%
\pgfpathlineto{\pgfqpoint{1.003593in}{5.540238in}}%
\pgfpathlineto{\pgfqpoint{1.003764in}{5.538610in}}%
\pgfpathlineto{\pgfqpoint{1.004020in}{5.547716in}}%
\pgfpathlineto{\pgfqpoint{1.004876in}{5.675336in}}%
\pgfpathlineto{\pgfqpoint{1.006844in}{6.413319in}}%
\pgfpathlineto{\pgfqpoint{1.010694in}{7.735516in}}%
\pgfpathlineto{\pgfqpoint{1.010951in}{7.740597in}}%
\pgfpathlineto{\pgfqpoint{1.011293in}{7.725948in}}%
\pgfpathlineto{\pgfqpoint{1.012320in}{7.541126in}}%
\pgfpathlineto{\pgfqpoint{1.014630in}{6.587626in}}%
\pgfpathlineto{\pgfqpoint{1.017882in}{5.545615in}}%
\pgfpathlineto{\pgfqpoint{1.018138in}{5.539217in}}%
\pgfpathlineto{\pgfqpoint{1.018566in}{5.558976in}}%
\pgfpathlineto{\pgfqpoint{1.019593in}{5.754589in}}%
\pgfpathlineto{\pgfqpoint{1.022074in}{6.797421in}}%
\pgfpathlineto{\pgfqpoint{1.025154in}{7.735524in}}%
\pgfpathlineto{\pgfqpoint{1.025326in}{7.739807in}}%
\pgfpathlineto{\pgfqpoint{1.025753in}{7.724027in}}%
\pgfpathlineto{\pgfqpoint{1.026780in}{7.537971in}}%
\pgfpathlineto{\pgfqpoint{1.029176in}{6.548337in}}%
\pgfpathlineto{\pgfqpoint{1.032342in}{5.547456in}}%
\pgfpathlineto{\pgfqpoint{1.032598in}{5.540015in}}%
\pgfpathlineto{\pgfqpoint{1.033026in}{5.557707in}}%
\pgfpathlineto{\pgfqpoint{1.034053in}{5.747167in}}%
\pgfpathlineto{\pgfqpoint{1.036449in}{6.737239in}}%
\pgfpathlineto{\pgfqpoint{1.039614in}{7.731930in}}%
\pgfpathlineto{\pgfqpoint{1.039871in}{7.739129in}}%
\pgfpathlineto{\pgfqpoint{1.040299in}{7.721202in}}%
\pgfpathlineto{\pgfqpoint{1.041326in}{7.532035in}}%
\pgfpathlineto{\pgfqpoint{1.043721in}{6.545292in}}%
\pgfpathlineto{\pgfqpoint{1.046887in}{5.548824in}}%
\pgfpathlineto{\pgfqpoint{1.047144in}{5.540867in}}%
\pgfpathlineto{\pgfqpoint{1.047657in}{5.565099in}}%
\pgfpathlineto{\pgfqpoint{1.048770in}{5.791699in}}%
\pgfpathlineto{\pgfqpoint{1.051593in}{6.998477in}}%
\pgfpathlineto{\pgfqpoint{1.054331in}{7.736214in}}%
\pgfpathlineto{\pgfqpoint{1.054502in}{7.738243in}}%
\pgfpathlineto{\pgfqpoint{1.054845in}{7.724543in}}%
\pgfpathlineto{\pgfqpoint{1.055786in}{7.569063in}}%
\pgfpathlineto{\pgfqpoint{1.057925in}{6.739353in}}%
\pgfpathlineto{\pgfqpoint{1.061518in}{5.549729in}}%
\pgfpathlineto{\pgfqpoint{1.061775in}{5.541711in}}%
\pgfpathlineto{\pgfqpoint{1.062289in}{5.565370in}}%
\pgfpathlineto{\pgfqpoint{1.063401in}{5.788928in}}%
\pgfpathlineto{\pgfqpoint{1.066139in}{6.948103in}}%
\pgfpathlineto{\pgfqpoint{1.068962in}{7.734017in}}%
\pgfpathlineto{\pgfqpoint{1.069134in}{7.737456in}}%
\pgfpathlineto{\pgfqpoint{1.069561in}{7.720454in}}%
\pgfpathlineto{\pgfqpoint{1.070588in}{7.536550in}}%
\pgfpathlineto{\pgfqpoint{1.072984in}{6.567435in}}%
\pgfpathlineto{\pgfqpoint{1.076235in}{5.550234in}}%
\pgfpathlineto{\pgfqpoint{1.076492in}{5.542538in}}%
\pgfpathlineto{\pgfqpoint{1.076920in}{5.558800in}}%
\pgfpathlineto{\pgfqpoint{1.077946in}{5.740262in}}%
\pgfpathlineto{\pgfqpoint{1.080257in}{6.663849in}}%
\pgfpathlineto{\pgfqpoint{1.083594in}{7.727332in}}%
\pgfpathlineto{\pgfqpoint{1.083936in}{7.736492in}}%
\pgfpathlineto{\pgfqpoint{1.084364in}{7.715429in}}%
\pgfpathlineto{\pgfqpoint{1.085476in}{7.500062in}}%
\pgfpathlineto{\pgfqpoint{1.088128in}{6.396397in}}%
\pgfpathlineto{\pgfqpoint{1.091038in}{5.550435in}}%
\pgfpathlineto{\pgfqpoint{1.091294in}{5.543371in}}%
\pgfpathlineto{\pgfqpoint{1.091722in}{5.560350in}}%
\pgfpathlineto{\pgfqpoint{1.092749in}{5.741784in}}%
\pgfpathlineto{\pgfqpoint{1.095145in}{6.699079in}}%
\pgfpathlineto{\pgfqpoint{1.098481in}{7.729618in}}%
\pgfpathlineto{\pgfqpoint{1.098738in}{7.735734in}}%
\pgfpathlineto{\pgfqpoint{1.099166in}{7.717348in}}%
\pgfpathlineto{\pgfqpoint{1.100193in}{7.533616in}}%
\pgfpathlineto{\pgfqpoint{1.102588in}{6.576963in}}%
\pgfpathlineto{\pgfqpoint{1.105925in}{5.550438in}}%
\pgfpathlineto{\pgfqpoint{1.106182in}{5.544249in}}%
\pgfpathlineto{\pgfqpoint{1.106610in}{5.562349in}}%
\pgfpathlineto{\pgfqpoint{1.107637in}{5.744645in}}%
\pgfpathlineto{\pgfqpoint{1.110032in}{6.696693in}}%
\pgfpathlineto{\pgfqpoint{1.113369in}{7.727535in}}%
\pgfpathlineto{\pgfqpoint{1.113626in}{7.734805in}}%
\pgfpathlineto{\pgfqpoint{1.114054in}{7.718669in}}%
\pgfpathlineto{\pgfqpoint{1.115081in}{7.541514in}}%
\pgfpathlineto{\pgfqpoint{1.117391in}{6.637390in}}%
\pgfpathlineto{\pgfqpoint{1.120813in}{5.554871in}}%
\pgfpathlineto{\pgfqpoint{1.121155in}{5.545215in}}%
\pgfpathlineto{\pgfqpoint{1.121583in}{5.564725in}}%
\pgfpathlineto{\pgfqpoint{1.122696in}{5.771860in}}%
\pgfpathlineto{\pgfqpoint{1.125262in}{6.812821in}}%
\pgfpathlineto{\pgfqpoint{1.128343in}{7.725533in}}%
\pgfpathlineto{\pgfqpoint{1.128599in}{7.733765in}}%
\pgfpathlineto{\pgfqpoint{1.129113in}{7.712557in}}%
\pgfpathlineto{\pgfqpoint{1.130225in}{7.502447in}}%
\pgfpathlineto{\pgfqpoint{1.132878in}{6.423763in}}%
\pgfpathlineto{\pgfqpoint{1.135872in}{5.554376in}}%
\pgfpathlineto{\pgfqpoint{1.136129in}{5.546235in}}%
\pgfpathlineto{\pgfqpoint{1.136642in}{5.567400in}}%
\pgfpathlineto{\pgfqpoint{1.137755in}{5.776464in}}%
\pgfpathlineto{\pgfqpoint{1.140407in}{6.850790in}}%
\pgfpathlineto{\pgfqpoint{1.143487in}{7.727977in}}%
\pgfpathlineto{\pgfqpoint{1.143744in}{7.732905in}}%
\pgfpathlineto{\pgfqpoint{1.144172in}{7.713547in}}%
\pgfpathlineto{\pgfqpoint{1.145284in}{7.509512in}}%
\pgfpathlineto{\pgfqpoint{1.147851in}{6.481397in}}%
\pgfpathlineto{\pgfqpoint{1.151017in}{5.553953in}}%
\pgfpathlineto{\pgfqpoint{1.151273in}{5.547067in}}%
\pgfpathlineto{\pgfqpoint{1.151701in}{5.563009in}}%
\pgfpathlineto{\pgfqpoint{1.152728in}{5.735712in}}%
\pgfpathlineto{\pgfqpoint{1.155038in}{6.619976in}}%
\pgfpathlineto{\pgfqpoint{1.158546in}{7.721610in}}%
\pgfpathlineto{\pgfqpoint{1.158889in}{7.732000in}}%
\pgfpathlineto{\pgfqpoint{1.159316in}{7.714339in}}%
\pgfpathlineto{\pgfqpoint{1.160343in}{7.538630in}}%
\pgfpathlineto{\pgfqpoint{1.162653in}{6.654023in}}%
\pgfpathlineto{\pgfqpoint{1.166161in}{5.558270in}}%
\pgfpathlineto{\pgfqpoint{1.166504in}{5.548031in}}%
\pgfpathlineto{\pgfqpoint{1.166931in}{5.565694in}}%
\pgfpathlineto{\pgfqpoint{1.167958in}{5.740605in}}%
\pgfpathlineto{\pgfqpoint{1.170268in}{6.621326in}}%
\pgfpathlineto{\pgfqpoint{1.173776in}{7.719557in}}%
\pgfpathlineto{\pgfqpoint{1.174119in}{7.731009in}}%
\pgfpathlineto{\pgfqpoint{1.174632in}{7.707841in}}%
\pgfpathlineto{\pgfqpoint{1.175744in}{7.499645in}}%
\pgfpathlineto{\pgfqpoint{1.178397in}{6.443410in}}%
\pgfpathlineto{\pgfqpoint{1.181477in}{5.557784in}}%
\pgfpathlineto{\pgfqpoint{1.181819in}{5.549119in}}%
\pgfpathlineto{\pgfqpoint{1.182247in}{5.568370in}}%
\pgfpathlineto{\pgfqpoint{1.183359in}{5.767595in}}%
\pgfpathlineto{\pgfqpoint{1.185926in}{6.774357in}}%
\pgfpathlineto{\pgfqpoint{1.189178in}{7.722465in}}%
\pgfpathlineto{\pgfqpoint{1.189434in}{7.729919in}}%
\pgfpathlineto{\pgfqpoint{1.189948in}{7.708953in}}%
\pgfpathlineto{\pgfqpoint{1.191060in}{7.506992in}}%
\pgfpathlineto{\pgfqpoint{1.193627in}{6.501048in}}%
\pgfpathlineto{\pgfqpoint{1.196878in}{5.557577in}}%
\pgfpathlineto{\pgfqpoint{1.197135in}{5.550136in}}%
\pgfpathlineto{\pgfqpoint{1.197648in}{5.570906in}}%
\pgfpathlineto{\pgfqpoint{1.198761in}{5.771524in}}%
\pgfpathlineto{\pgfqpoint{1.201328in}{6.772811in}}%
\pgfpathlineto{\pgfqpoint{1.204579in}{7.720248in}}%
\pgfpathlineto{\pgfqpoint{1.204921in}{7.728894in}}%
\pgfpathlineto{\pgfqpoint{1.205349in}{7.710176in}}%
\pgfpathlineto{\pgfqpoint{1.206461in}{7.514941in}}%
\pgfpathlineto{\pgfqpoint{1.208943in}{6.559842in}}%
\pgfpathlineto{\pgfqpoint{1.212280in}{5.562475in}}%
\pgfpathlineto{\pgfqpoint{1.212622in}{5.551107in}}%
\pgfpathlineto{\pgfqpoint{1.213135in}{5.573162in}}%
\pgfpathlineto{\pgfqpoint{1.214247in}{5.774396in}}%
\pgfpathlineto{\pgfqpoint{1.216814in}{6.769462in}}%
\pgfpathlineto{\pgfqpoint{1.220066in}{7.717673in}}%
\pgfpathlineto{\pgfqpoint{1.220408in}{7.727921in}}%
\pgfpathlineto{\pgfqpoint{1.220921in}{7.704438in}}%
\pgfpathlineto{\pgfqpoint{1.222119in}{7.477900in}}%
\pgfpathlineto{\pgfqpoint{1.224943in}{6.360694in}}%
\pgfpathlineto{\pgfqpoint{1.227937in}{5.558113in}}%
\pgfpathlineto{\pgfqpoint{1.228194in}{5.552155in}}%
\pgfpathlineto{\pgfqpoint{1.228622in}{5.567987in}}%
\pgfpathlineto{\pgfqpoint{1.229649in}{5.732441in}}%
\pgfpathlineto{\pgfqpoint{1.231959in}{6.577043in}}%
\pgfpathlineto{\pgfqpoint{1.235638in}{7.714595in}}%
\pgfpathlineto{\pgfqpoint{1.235980in}{7.726752in}}%
\pgfpathlineto{\pgfqpoint{1.236494in}{7.706616in}}%
\pgfpathlineto{\pgfqpoint{1.237606in}{7.512029in}}%
\pgfpathlineto{\pgfqpoint{1.240173in}{6.533633in}}%
\pgfpathlineto{\pgfqpoint{1.243510in}{5.563321in}}%
\pgfpathlineto{\pgfqpoint{1.243852in}{5.553245in}}%
\pgfpathlineto{\pgfqpoint{1.244365in}{5.576247in}}%
\pgfpathlineto{\pgfqpoint{1.245478in}{5.775528in}}%
\pgfpathlineto{\pgfqpoint{1.248045in}{6.755062in}}%
\pgfpathlineto{\pgfqpoint{1.251382in}{7.716326in}}%
\pgfpathlineto{\pgfqpoint{1.251724in}{7.725726in}}%
\pgfpathlineto{\pgfqpoint{1.252152in}{7.709056in}}%
\pgfpathlineto{\pgfqpoint{1.253178in}{7.545131in}}%
\pgfpathlineto{\pgfqpoint{1.255489in}{6.710290in}}%
\pgfpathlineto{\pgfqpoint{1.259168in}{5.570132in}}%
\pgfpathlineto{\pgfqpoint{1.259596in}{5.554360in}}%
\pgfpathlineto{\pgfqpoint{1.260109in}{5.576810in}}%
\pgfpathlineto{\pgfqpoint{1.261221in}{5.773103in}}%
\pgfpathlineto{\pgfqpoint{1.263788in}{6.742862in}}%
\pgfpathlineto{\pgfqpoint{1.267125in}{7.712280in}}%
\pgfpathlineto{\pgfqpoint{1.267467in}{7.724487in}}%
\pgfpathlineto{\pgfqpoint{1.267981in}{7.705407in}}%
\pgfpathlineto{\pgfqpoint{1.269093in}{7.516723in}}%
\pgfpathlineto{\pgfqpoint{1.271574in}{6.594694in}}%
\pgfpathlineto{\pgfqpoint{1.275082in}{5.566223in}}%
\pgfpathlineto{\pgfqpoint{1.275425in}{5.555516in}}%
\pgfpathlineto{\pgfqpoint{1.275938in}{5.576594in}}%
\pgfpathlineto{\pgfqpoint{1.277050in}{5.768280in}}%
\pgfpathlineto{\pgfqpoint{1.279617in}{6.726563in}}%
\pgfpathlineto{\pgfqpoint{1.283040in}{7.712830in}}%
\pgfpathlineto{\pgfqpoint{1.283382in}{7.723444in}}%
\pgfpathlineto{\pgfqpoint{1.283895in}{7.702472in}}%
\pgfpathlineto{\pgfqpoint{1.285008in}{7.511945in}}%
\pgfpathlineto{\pgfqpoint{1.287489in}{6.594714in}}%
\pgfpathlineto{\pgfqpoint{1.290997in}{5.568692in}}%
\pgfpathlineto{\pgfqpoint{1.291339in}{5.556781in}}%
\pgfpathlineto{\pgfqpoint{1.291853in}{5.575567in}}%
\pgfpathlineto{\pgfqpoint{1.292965in}{5.760817in}}%
\pgfpathlineto{\pgfqpoint{1.295446in}{6.669274in}}%
\pgfpathlineto{\pgfqpoint{1.298954in}{7.707090in}}%
\pgfpathlineto{\pgfqpoint{1.299382in}{7.722282in}}%
\pgfpathlineto{\pgfqpoint{1.299895in}{7.700482in}}%
\pgfpathlineto{\pgfqpoint{1.301008in}{7.510215in}}%
\pgfpathlineto{\pgfqpoint{1.303575in}{6.563766in}}%
\pgfpathlineto{\pgfqpoint{1.306997in}{5.572093in}}%
\pgfpathlineto{\pgfqpoint{1.307425in}{5.557866in}}%
\pgfpathlineto{\pgfqpoint{1.307938in}{5.580552in}}%
\pgfpathlineto{\pgfqpoint{1.309136in}{5.793663in}}%
\pgfpathlineto{\pgfqpoint{1.311874in}{6.822848in}}%
\pgfpathlineto{\pgfqpoint{1.315126in}{7.711453in}}%
\pgfpathlineto{\pgfqpoint{1.315468in}{7.721079in}}%
\pgfpathlineto{\pgfqpoint{1.315981in}{7.699607in}}%
\pgfpathlineto{\pgfqpoint{1.317094in}{7.511880in}}%
\pgfpathlineto{\pgfqpoint{1.319575in}{6.610915in}}%
\pgfpathlineto{\pgfqpoint{1.323083in}{5.576730in}}%
\pgfpathlineto{\pgfqpoint{1.323511in}{5.559190in}}%
\pgfpathlineto{\pgfqpoint{1.324024in}{5.577372in}}%
\pgfpathlineto{\pgfqpoint{1.325136in}{5.757628in}}%
\pgfpathlineto{\pgfqpoint{1.327618in}{6.647364in}}%
\pgfpathlineto{\pgfqpoint{1.331211in}{7.703946in}}%
\pgfpathlineto{\pgfqpoint{1.331639in}{7.719830in}}%
\pgfpathlineto{\pgfqpoint{1.332153in}{7.699929in}}%
\pgfpathlineto{\pgfqpoint{1.333265in}{7.517205in}}%
\pgfpathlineto{\pgfqpoint{1.335746in}{6.628215in}}%
\pgfpathlineto{\pgfqpoint{1.339340in}{5.576339in}}%
\pgfpathlineto{\pgfqpoint{1.339768in}{5.560342in}}%
\pgfpathlineto{\pgfqpoint{1.340281in}{5.579840in}}%
\pgfpathlineto{\pgfqpoint{1.341393in}{5.760818in}}%
\pgfpathlineto{\pgfqpoint{1.343875in}{6.644576in}}%
\pgfpathlineto{\pgfqpoint{1.347468in}{7.700550in}}%
\pgfpathlineto{\pgfqpoint{1.347982in}{7.718487in}}%
\pgfpathlineto{\pgfqpoint{1.348409in}{7.701410in}}%
\pgfpathlineto{\pgfqpoint{1.349522in}{7.526351in}}%
\pgfpathlineto{\pgfqpoint{1.351917in}{6.687795in}}%
\pgfpathlineto{\pgfqpoint{1.355682in}{5.577284in}}%
\pgfpathlineto{\pgfqpoint{1.356110in}{5.561605in}}%
\pgfpathlineto{\pgfqpoint{1.356623in}{5.580952in}}%
\pgfpathlineto{\pgfqpoint{1.357736in}{5.759743in}}%
\pgfpathlineto{\pgfqpoint{1.360217in}{6.634382in}}%
\pgfpathlineto{\pgfqpoint{1.363896in}{7.702017in}}%
\pgfpathlineto{\pgfqpoint{1.364324in}{7.717300in}}%
\pgfpathlineto{\pgfqpoint{1.364837in}{7.697745in}}%
\pgfpathlineto{\pgfqpoint{1.365950in}{7.519482in}}%
\pgfpathlineto{\pgfqpoint{1.368431in}{6.648806in}}%
\pgfpathlineto{\pgfqpoint{1.372110in}{5.579625in}}%
\pgfpathlineto{\pgfqpoint{1.372538in}{5.562975in}}%
\pgfpathlineto{\pgfqpoint{1.373051in}{5.580625in}}%
\pgfpathlineto{\pgfqpoint{1.374164in}{5.754153in}}%
\pgfpathlineto{\pgfqpoint{1.376559in}{6.581243in}}%
\pgfpathlineto{\pgfqpoint{1.380410in}{7.701897in}}%
\pgfpathlineto{\pgfqpoint{1.380838in}{7.716006in}}%
\pgfpathlineto{\pgfqpoint{1.381351in}{7.695576in}}%
\pgfpathlineto{\pgfqpoint{1.382463in}{7.517454in}}%
\pgfpathlineto{\pgfqpoint{1.384945in}{6.653715in}}%
\pgfpathlineto{\pgfqpoint{1.388624in}{5.583601in}}%
\pgfpathlineto{\pgfqpoint{1.389137in}{5.564206in}}%
\pgfpathlineto{\pgfqpoint{1.389650in}{5.585261in}}%
\pgfpathlineto{\pgfqpoint{1.390848in}{5.784326in}}%
\pgfpathlineto{\pgfqpoint{1.393501in}{6.728510in}}%
\pgfpathlineto{\pgfqpoint{1.397009in}{7.700268in}}%
\pgfpathlineto{\pgfqpoint{1.397437in}{7.714666in}}%
\pgfpathlineto{\pgfqpoint{1.397950in}{7.695111in}}%
\pgfpathlineto{\pgfqpoint{1.399062in}{7.520605in}}%
\pgfpathlineto{\pgfqpoint{1.401544in}{6.667647in}}%
\pgfpathlineto{\pgfqpoint{1.405308in}{5.582929in}}%
\pgfpathlineto{\pgfqpoint{1.405822in}{5.565636in}}%
\pgfpathlineto{\pgfqpoint{1.406335in}{5.588214in}}%
\pgfpathlineto{\pgfqpoint{1.407533in}{5.788394in}}%
\pgfpathlineto{\pgfqpoint{1.410185in}{6.726263in}}%
\pgfpathlineto{\pgfqpoint{1.413694in}{7.696990in}}%
\pgfpathlineto{\pgfqpoint{1.414121in}{7.713196in}}%
\pgfpathlineto{\pgfqpoint{1.414635in}{7.696335in}}%
\pgfpathlineto{\pgfqpoint{1.415747in}{7.529081in}}%
\pgfpathlineto{\pgfqpoint{1.418143in}{6.725396in}}%
\pgfpathlineto{\pgfqpoint{1.422079in}{5.584147in}}%
\pgfpathlineto{\pgfqpoint{1.422592in}{5.567007in}}%
\pgfpathlineto{\pgfqpoint{1.423105in}{5.589164in}}%
\pgfpathlineto{\pgfqpoint{1.424303in}{5.786349in}}%
\pgfpathlineto{\pgfqpoint{1.426956in}{6.713949in}}%
\pgfpathlineto{\pgfqpoint{1.430464in}{7.691658in}}%
\pgfpathlineto{\pgfqpoint{1.430977in}{7.711905in}}%
\pgfpathlineto{\pgfqpoint{1.431491in}{7.693137in}}%
\pgfpathlineto{\pgfqpoint{1.432603in}{7.523943in}}%
\pgfpathlineto{\pgfqpoint{1.434999in}{6.724622in}}%
\pgfpathlineto{\pgfqpoint{1.438934in}{5.587334in}}%
\pgfpathlineto{\pgfqpoint{1.439448in}{5.568335in}}%
\pgfpathlineto{\pgfqpoint{1.439961in}{5.588064in}}%
\pgfpathlineto{\pgfqpoint{1.441074in}{5.758242in}}%
\pgfpathlineto{\pgfqpoint{1.443469in}{6.555397in}}%
\pgfpathlineto{\pgfqpoint{1.447405in}{7.690583in}}%
\pgfpathlineto{\pgfqpoint{1.447919in}{7.710484in}}%
\pgfpathlineto{\pgfqpoint{1.448432in}{7.691942in}}%
\pgfpathlineto{\pgfqpoint{1.449544in}{7.525064in}}%
\pgfpathlineto{\pgfqpoint{1.451940in}{6.734955in}}%
\pgfpathlineto{\pgfqpoint{1.455961in}{5.586433in}}%
\pgfpathlineto{\pgfqpoint{1.456475in}{5.569853in}}%
\pgfpathlineto{\pgfqpoint{1.456988in}{5.591426in}}%
\pgfpathlineto{\pgfqpoint{1.458186in}{5.783181in}}%
\pgfpathlineto{\pgfqpoint{1.460753in}{6.656947in}}%
\pgfpathlineto{\pgfqpoint{1.464432in}{7.687192in}}%
\pgfpathlineto{\pgfqpoint{1.464945in}{7.708933in}}%
\pgfpathlineto{\pgfqpoint{1.465459in}{7.692805in}}%
\pgfpathlineto{\pgfqpoint{1.466571in}{7.532620in}}%
\pgfpathlineto{\pgfqpoint{1.468881in}{6.790136in}}%
\pgfpathlineto{\pgfqpoint{1.473074in}{5.587781in}}%
\pgfpathlineto{\pgfqpoint{1.473587in}{5.571310in}}%
\pgfpathlineto{\pgfqpoint{1.474101in}{5.592422in}}%
\pgfpathlineto{\pgfqpoint{1.475299in}{5.781110in}}%
\pgfpathlineto{\pgfqpoint{1.477865in}{6.644590in}}%
\pgfpathlineto{\pgfqpoint{1.481630in}{7.687920in}}%
\pgfpathlineto{\pgfqpoint{1.482144in}{7.707543in}}%
\pgfpathlineto{\pgfqpoint{1.482657in}{7.689863in}}%
\pgfpathlineto{\pgfqpoint{1.483769in}{7.528428in}}%
\pgfpathlineto{\pgfqpoint{1.486165in}{6.757817in}}%
\pgfpathlineto{\pgfqpoint{1.490272in}{5.591455in}}%
\pgfpathlineto{\pgfqpoint{1.490785in}{5.572735in}}%
\pgfpathlineto{\pgfqpoint{1.491299in}{5.591035in}}%
\pgfpathlineto{\pgfqpoint{1.492411in}{5.752786in}}%
\pgfpathlineto{\pgfqpoint{1.494807in}{6.520467in}}%
\pgfpathlineto{\pgfqpoint{1.498914in}{7.686031in}}%
\pgfpathlineto{\pgfqpoint{1.499427in}{7.706008in}}%
\pgfpathlineto{\pgfqpoint{1.499940in}{7.689248in}}%
\pgfpathlineto{\pgfqpoint{1.501053in}{7.531511in}}%
\pgfpathlineto{\pgfqpoint{1.503363in}{6.804535in}}%
\pgfpathlineto{\pgfqpoint{1.507641in}{5.591581in}}%
\pgfpathlineto{\pgfqpoint{1.508154in}{5.574275in}}%
\pgfpathlineto{\pgfqpoint{1.508668in}{5.593421in}}%
\pgfpathlineto{\pgfqpoint{1.509866in}{5.773750in}}%
\pgfpathlineto{\pgfqpoint{1.512433in}{6.614117in}}%
\pgfpathlineto{\pgfqpoint{1.516283in}{7.681346in}}%
\pgfpathlineto{\pgfqpoint{1.516882in}{7.704477in}}%
\pgfpathlineto{\pgfqpoint{1.517395in}{7.685138in}}%
\pgfpathlineto{\pgfqpoint{1.518593in}{7.505434in}}%
\pgfpathlineto{\pgfqpoint{1.521160in}{6.669354in}}%
\pgfpathlineto{\pgfqpoint{1.525010in}{5.600897in}}%
\pgfpathlineto{\pgfqpoint{1.525609in}{5.575828in}}%
\pgfpathlineto{\pgfqpoint{1.526123in}{5.593198in}}%
\pgfpathlineto{\pgfqpoint{1.527235in}{5.749395in}}%
\pgfpathlineto{\pgfqpoint{1.529631in}{6.497240in}}%
\pgfpathlineto{\pgfqpoint{1.533823in}{7.680298in}}%
\pgfpathlineto{\pgfqpoint{1.534422in}{7.702896in}}%
\pgfpathlineto{\pgfqpoint{1.534936in}{7.683709in}}%
\pgfpathlineto{\pgfqpoint{1.536133in}{7.506403in}}%
\pgfpathlineto{\pgfqpoint{1.538615in}{6.712501in}}%
\pgfpathlineto{\pgfqpoint{1.542636in}{5.600116in}}%
\pgfpathlineto{\pgfqpoint{1.543235in}{5.577434in}}%
\pgfpathlineto{\pgfqpoint{1.543748in}{5.596244in}}%
\pgfpathlineto{\pgfqpoint{1.544946in}{5.771711in}}%
\pgfpathlineto{\pgfqpoint{1.547428in}{6.560146in}}%
\pgfpathlineto{\pgfqpoint{1.551449in}{7.676007in}}%
\pgfpathlineto{\pgfqpoint{1.552048in}{7.701292in}}%
\pgfpathlineto{\pgfqpoint{1.552561in}{7.685020in}}%
\pgfpathlineto{\pgfqpoint{1.553674in}{7.533787in}}%
\pgfpathlineto{\pgfqpoint{1.555984in}{6.833601in}}%
\pgfpathlineto{\pgfqpoint{1.560433in}{5.596424in}}%
\pgfpathlineto{\pgfqpoint{1.560947in}{5.579040in}}%
\pgfpathlineto{\pgfqpoint{1.561460in}{5.596423in}}%
\pgfpathlineto{\pgfqpoint{1.562572in}{5.748972in}}%
\pgfpathlineto{\pgfqpoint{1.564882in}{6.447587in}}%
\pgfpathlineto{\pgfqpoint{1.569332in}{7.681475in}}%
\pgfpathlineto{\pgfqpoint{1.569845in}{7.699661in}}%
\pgfpathlineto{\pgfqpoint{1.570358in}{7.683358in}}%
\pgfpathlineto{\pgfqpoint{1.571471in}{7.533897in}}%
\pgfpathlineto{\pgfqpoint{1.573781in}{6.842215in}}%
\pgfpathlineto{\pgfqpoint{1.578230in}{5.602081in}}%
\pgfpathlineto{\pgfqpoint{1.578829in}{5.580738in}}%
\pgfpathlineto{\pgfqpoint{1.579342in}{5.599486in}}%
\pgfpathlineto{\pgfqpoint{1.580540in}{5.770698in}}%
\pgfpathlineto{\pgfqpoint{1.583022in}{6.540508in}}%
\pgfpathlineto{\pgfqpoint{1.587129in}{7.670512in}}%
\pgfpathlineto{\pgfqpoint{1.587813in}{7.697938in}}%
\pgfpathlineto{\pgfqpoint{1.588327in}{7.678961in}}%
\pgfpathlineto{\pgfqpoint{1.589524in}{7.508284in}}%
\pgfpathlineto{\pgfqpoint{1.592006in}{6.742567in}}%
\pgfpathlineto{\pgfqpoint{1.596198in}{5.605044in}}%
\pgfpathlineto{\pgfqpoint{1.596797in}{5.582390in}}%
\pgfpathlineto{\pgfqpoint{1.597311in}{5.599414in}}%
\pgfpathlineto{\pgfqpoint{1.598423in}{5.747591in}}%
\pgfpathlineto{\pgfqpoint{1.600733in}{6.428182in}}%
\pgfpathlineto{\pgfqpoint{1.605268in}{7.675617in}}%
\pgfpathlineto{\pgfqpoint{1.605867in}{7.696242in}}%
\pgfpathlineto{\pgfqpoint{1.606380in}{7.677779in}}%
\pgfpathlineto{\pgfqpoint{1.607578in}{7.510254in}}%
\pgfpathlineto{\pgfqpoint{1.610059in}{6.755054in}}%
\pgfpathlineto{\pgfqpoint{1.614252in}{5.611981in}}%
\pgfpathlineto{\pgfqpoint{1.614936in}{5.584136in}}%
\pgfpathlineto{\pgfqpoint{1.615450in}{5.601830in}}%
\pgfpathlineto{\pgfqpoint{1.616648in}{5.766681in}}%
\pgfpathlineto{\pgfqpoint{1.619129in}{6.515359in}}%
\pgfpathlineto{\pgfqpoint{1.623407in}{7.670248in}}%
\pgfpathlineto{\pgfqpoint{1.624006in}{7.694503in}}%
\pgfpathlineto{\pgfqpoint{1.624519in}{7.679754in}}%
\pgfpathlineto{\pgfqpoint{1.625632in}{7.538888in}}%
\pgfpathlineto{\pgfqpoint{1.627942in}{6.876962in}}%
\pgfpathlineto{\pgfqpoint{1.632648in}{5.603376in}}%
\pgfpathlineto{\pgfqpoint{1.633161in}{5.585900in}}%
\pgfpathlineto{\pgfqpoint{1.633675in}{5.600973in}}%
\pgfpathlineto{\pgfqpoint{1.634787in}{5.741587in}}%
\pgfpathlineto{\pgfqpoint{1.637097in}{6.400017in}}%
\pgfpathlineto{\pgfqpoint{1.641803in}{7.673555in}}%
\pgfpathlineto{\pgfqpoint{1.642402in}{7.692673in}}%
\pgfpathlineto{\pgfqpoint{1.642915in}{7.674114in}}%
\pgfpathlineto{\pgfqpoint{1.644113in}{7.510454in}}%
\pgfpathlineto{\pgfqpoint{1.646595in}{6.773666in}}%
\pgfpathlineto{\pgfqpoint{1.650958in}{5.611246in}}%
\pgfpathlineto{\pgfqpoint{1.651557in}{5.587712in}}%
\pgfpathlineto{\pgfqpoint{1.652071in}{5.602188in}}%
\pgfpathlineto{\pgfqpoint{1.653183in}{5.739790in}}%
\pgfpathlineto{\pgfqpoint{1.655493in}{6.388021in}}%
\pgfpathlineto{\pgfqpoint{1.660285in}{7.672674in}}%
\pgfpathlineto{\pgfqpoint{1.660798in}{7.690864in}}%
\pgfpathlineto{\pgfqpoint{1.661311in}{7.677329in}}%
\pgfpathlineto{\pgfqpoint{1.662424in}{7.542536in}}%
\pgfpathlineto{\pgfqpoint{1.664648in}{6.930453in}}%
\pgfpathlineto{\pgfqpoint{1.669611in}{5.605325in}}%
\pgfpathlineto{\pgfqpoint{1.670124in}{5.589511in}}%
\pgfpathlineto{\pgfqpoint{1.670638in}{5.605149in}}%
\pgfpathlineto{\pgfqpoint{1.671750in}{5.743306in}}%
\pgfpathlineto{\pgfqpoint{1.674060in}{6.385895in}}%
\pgfpathlineto{\pgfqpoint{1.678852in}{7.667791in}}%
\pgfpathlineto{\pgfqpoint{1.679451in}{7.689096in}}%
\pgfpathlineto{\pgfqpoint{1.679964in}{7.673595in}}%
\pgfpathlineto{\pgfqpoint{1.681076in}{7.536611in}}%
\pgfpathlineto{\pgfqpoint{1.683386in}{6.898704in}}%
\pgfpathlineto{\pgfqpoint{1.688263in}{5.609215in}}%
\pgfpathlineto{\pgfqpoint{1.688777in}{5.591468in}}%
\pgfpathlineto{\pgfqpoint{1.689290in}{5.604628in}}%
\pgfpathlineto{\pgfqpoint{1.690403in}{5.735968in}}%
\pgfpathlineto{\pgfqpoint{1.692627in}{6.334552in}}%
\pgfpathlineto{\pgfqpoint{1.697675in}{7.670274in}}%
\pgfpathlineto{\pgfqpoint{1.698189in}{7.687165in}}%
\pgfpathlineto{\pgfqpoint{1.698702in}{7.673419in}}%
\pgfpathlineto{\pgfqpoint{1.699814in}{7.541782in}}%
\pgfpathlineto{\pgfqpoint{1.702039in}{6.945839in}}%
\pgfpathlineto{\pgfqpoint{1.707087in}{5.611704in}}%
\pgfpathlineto{\pgfqpoint{1.707686in}{5.593357in}}%
\pgfpathlineto{\pgfqpoint{1.708200in}{5.610514in}}%
\pgfpathlineto{\pgfqpoint{1.709397in}{5.764090in}}%
\pgfpathlineto{\pgfqpoint{1.711793in}{6.435308in}}%
\pgfpathlineto{\pgfqpoint{1.716499in}{7.662716in}}%
\pgfpathlineto{\pgfqpoint{1.717098in}{7.685225in}}%
\pgfpathlineto{\pgfqpoint{1.717611in}{7.671931in}}%
\pgfpathlineto{\pgfqpoint{1.718724in}{7.542994in}}%
\pgfpathlineto{\pgfqpoint{1.720948in}{6.956592in}}%
\pgfpathlineto{\pgfqpoint{1.726082in}{5.612946in}}%
\pgfpathlineto{\pgfqpoint{1.726681in}{5.595330in}}%
\pgfpathlineto{\pgfqpoint{1.727194in}{5.612528in}}%
\pgfpathlineto{\pgfqpoint{1.728392in}{5.764194in}}%
\pgfpathlineto{\pgfqpoint{1.730788in}{6.426423in}}%
\pgfpathlineto{\pgfqpoint{1.735494in}{7.655903in}}%
\pgfpathlineto{\pgfqpoint{1.736178in}{7.683295in}}%
\pgfpathlineto{\pgfqpoint{1.736692in}{7.669390in}}%
\pgfpathlineto{\pgfqpoint{1.737804in}{7.540985in}}%
\pgfpathlineto{\pgfqpoint{1.740029in}{6.961312in}}%
\pgfpathlineto{\pgfqpoint{1.745248in}{5.613233in}}%
\pgfpathlineto{\pgfqpoint{1.745761in}{5.597260in}}%
\pgfpathlineto{\pgfqpoint{1.746275in}{5.610577in}}%
\pgfpathlineto{\pgfqpoint{1.747387in}{5.736903in}}%
\pgfpathlineto{\pgfqpoint{1.749612in}{6.310836in}}%
\pgfpathlineto{\pgfqpoint{1.754831in}{7.662520in}}%
\pgfpathlineto{\pgfqpoint{1.755430in}{7.681287in}}%
\pgfpathlineto{\pgfqpoint{1.755943in}{7.665944in}}%
\pgfpathlineto{\pgfqpoint{1.757056in}{7.536380in}}%
\pgfpathlineto{\pgfqpoint{1.759366in}{6.933553in}}%
\pgfpathlineto{\pgfqpoint{1.764500in}{5.617972in}}%
\pgfpathlineto{\pgfqpoint{1.765099in}{5.599268in}}%
\pgfpathlineto{\pgfqpoint{1.765612in}{5.614376in}}%
\pgfpathlineto{\pgfqpoint{1.766724in}{5.742583in}}%
\pgfpathlineto{\pgfqpoint{1.768949in}{6.312968in}}%
\pgfpathlineto{\pgfqpoint{1.774254in}{7.663364in}}%
\pgfpathlineto{\pgfqpoint{1.774767in}{7.679205in}}%
\pgfpathlineto{\pgfqpoint{1.775280in}{7.666555in}}%
\pgfpathlineto{\pgfqpoint{1.776393in}{7.544229in}}%
\pgfpathlineto{\pgfqpoint{1.778617in}{6.984635in}}%
\pgfpathlineto{\pgfqpoint{1.784008in}{5.616888in}}%
\pgfpathlineto{\pgfqpoint{1.784521in}{5.601364in}}%
\pgfpathlineto{\pgfqpoint{1.785035in}{5.614065in}}%
\pgfpathlineto{\pgfqpoint{1.786147in}{5.735619in}}%
\pgfpathlineto{\pgfqpoint{1.788372in}{6.291172in}}%
\pgfpathlineto{\pgfqpoint{1.793762in}{7.659504in}}%
\pgfpathlineto{\pgfqpoint{1.794361in}{7.677133in}}%
\pgfpathlineto{\pgfqpoint{1.794874in}{7.661964in}}%
\pgfpathlineto{\pgfqpoint{1.796072in}{7.521638in}}%
\pgfpathlineto{\pgfqpoint{1.798468in}{6.896710in}}%
\pgfpathlineto{\pgfqpoint{1.803516in}{5.626057in}}%
\pgfpathlineto{\pgfqpoint{1.804201in}{5.603483in}}%
\pgfpathlineto{\pgfqpoint{1.804714in}{5.618839in}}%
\pgfpathlineto{\pgfqpoint{1.805912in}{5.758591in}}%
\pgfpathlineto{\pgfqpoint{1.808308in}{6.379445in}}%
\pgfpathlineto{\pgfqpoint{1.813356in}{7.650127in}}%
\pgfpathlineto{\pgfqpoint{1.814040in}{7.675022in}}%
\pgfpathlineto{\pgfqpoint{1.814554in}{7.661716in}}%
\pgfpathlineto{\pgfqpoint{1.815666in}{7.541546in}}%
\pgfpathlineto{\pgfqpoint{1.817891in}{6.996860in}}%
\pgfpathlineto{\pgfqpoint{1.823452in}{5.619195in}}%
\pgfpathlineto{\pgfqpoint{1.823966in}{5.605597in}}%
\pgfpathlineto{\pgfqpoint{1.824479in}{5.619170in}}%
\pgfpathlineto{\pgfqpoint{1.825591in}{5.739017in}}%
\pgfpathlineto{\pgfqpoint{1.827816in}{6.280272in}}%
\pgfpathlineto{\pgfqpoint{1.833377in}{7.657468in}}%
\pgfpathlineto{\pgfqpoint{1.833891in}{7.672783in}}%
\pgfpathlineto{\pgfqpoint{1.834404in}{7.661188in}}%
\pgfpathlineto{\pgfqpoint{1.835516in}{7.546249in}}%
\pgfpathlineto{\pgfqpoint{1.837741in}{7.014662in}}%
\pgfpathlineto{\pgfqpoint{1.843388in}{5.622642in}}%
\pgfpathlineto{\pgfqpoint{1.843902in}{5.607837in}}%
\pgfpathlineto{\pgfqpoint{1.844415in}{5.619682in}}%
\pgfpathlineto{\pgfqpoint{1.845527in}{5.734268in}}%
\pgfpathlineto{\pgfqpoint{1.847752in}{6.262376in}}%
\pgfpathlineto{\pgfqpoint{1.853485in}{7.658450in}}%
\pgfpathlineto{\pgfqpoint{1.853998in}{7.670604in}}%
\pgfpathlineto{\pgfqpoint{1.854511in}{7.656372in}}%
\pgfpathlineto{\pgfqpoint{1.855709in}{7.523946in}}%
\pgfpathlineto{\pgfqpoint{1.858019in}{6.955346in}}%
\pgfpathlineto{\pgfqpoint{1.863410in}{5.631580in}}%
\pgfpathlineto{\pgfqpoint{1.864094in}{5.610072in}}%
\pgfpathlineto{\pgfqpoint{1.864608in}{5.624401in}}%
\pgfpathlineto{\pgfqpoint{1.865806in}{5.756068in}}%
\pgfpathlineto{\pgfqpoint{1.868116in}{6.320525in}}%
\pgfpathlineto{\pgfqpoint{1.873506in}{7.644375in}}%
\pgfpathlineto{\pgfqpoint{1.874191in}{7.668347in}}%
\pgfpathlineto{\pgfqpoint{1.874704in}{7.656171in}}%
\pgfpathlineto{\pgfqpoint{1.875816in}{7.543485in}}%
\pgfpathlineto{\pgfqpoint{1.878041in}{7.026957in}}%
\pgfpathlineto{\pgfqpoint{1.883859in}{5.625881in}}%
\pgfpathlineto{\pgfqpoint{1.884373in}{5.612342in}}%
\pgfpathlineto{\pgfqpoint{1.884886in}{5.624415in}}%
\pgfpathlineto{\pgfqpoint{1.885998in}{5.736038in}}%
\pgfpathlineto{\pgfqpoint{1.888223in}{6.248111in}}%
\pgfpathlineto{\pgfqpoint{1.894127in}{7.654596in}}%
\pgfpathlineto{\pgfqpoint{1.894640in}{7.666012in}}%
\pgfpathlineto{\pgfqpoint{1.895153in}{7.652075in}}%
\pgfpathlineto{\pgfqpoint{1.896351in}{7.524154in}}%
\pgfpathlineto{\pgfqpoint{1.898662in}{6.973992in}}%
\pgfpathlineto{\pgfqpoint{1.904223in}{5.635687in}}%
\pgfpathlineto{\pgfqpoint{1.904908in}{5.614687in}}%
\pgfpathlineto{\pgfqpoint{1.905421in}{5.628198in}}%
\pgfpathlineto{\pgfqpoint{1.906619in}{5.754227in}}%
\pgfpathlineto{\pgfqpoint{1.908929in}{6.298791in}}%
\pgfpathlineto{\pgfqpoint{1.914576in}{7.644889in}}%
\pgfpathlineto{\pgfqpoint{1.915175in}{7.663632in}}%
\pgfpathlineto{\pgfqpoint{1.915688in}{7.652796in}}%
\pgfpathlineto{\pgfqpoint{1.916801in}{7.546261in}}%
\pgfpathlineto{\pgfqpoint{1.918940in}{7.074415in}}%
\pgfpathlineto{\pgfqpoint{1.925100in}{5.627573in}}%
\pgfpathlineto{\pgfqpoint{1.925614in}{5.617124in}}%
\pgfpathlineto{\pgfqpoint{1.926042in}{5.627205in}}%
\pgfpathlineto{\pgfqpoint{1.927154in}{5.731274in}}%
\pgfpathlineto{\pgfqpoint{1.929293in}{6.196905in}}%
\pgfpathlineto{\pgfqpoint{1.935539in}{7.651875in}}%
\pgfpathlineto{\pgfqpoint{1.935967in}{7.661248in}}%
\pgfpathlineto{\pgfqpoint{1.936480in}{7.650182in}}%
\pgfpathlineto{\pgfqpoint{1.937592in}{7.544865in}}%
\pgfpathlineto{\pgfqpoint{1.939732in}{7.080312in}}%
\pgfpathlineto{\pgfqpoint{1.946063in}{5.626446in}}%
\pgfpathlineto{\pgfqpoint{1.946491in}{5.619554in}}%
\pgfpathlineto{\pgfqpoint{1.946919in}{5.629393in}}%
\pgfpathlineto{\pgfqpoint{1.948031in}{5.731265in}}%
\pgfpathlineto{\pgfqpoint{1.950170in}{6.188235in}}%
\pgfpathlineto{\pgfqpoint{1.956587in}{7.652352in}}%
\pgfpathlineto{\pgfqpoint{1.956930in}{7.658756in}}%
\pgfpathlineto{\pgfqpoint{1.957443in}{7.648494in}}%
\pgfpathlineto{\pgfqpoint{1.958555in}{7.546507in}}%
\pgfpathlineto{\pgfqpoint{1.960694in}{7.092177in}}%
\pgfpathlineto{\pgfqpoint{1.967112in}{5.629684in}}%
\pgfpathlineto{\pgfqpoint{1.967539in}{5.621981in}}%
\pgfpathlineto{\pgfqpoint{1.968053in}{5.634353in}}%
\pgfpathlineto{\pgfqpoint{1.969165in}{5.739842in}}%
\pgfpathlineto{\pgfqpoint{1.971390in}{6.219608in}}%
\pgfpathlineto{\pgfqpoint{1.977636in}{7.645143in}}%
\pgfpathlineto{\pgfqpoint{1.978149in}{7.656298in}}%
\pgfpathlineto{\pgfqpoint{1.978662in}{7.644125in}}%
\pgfpathlineto{\pgfqpoint{1.979775in}{7.539880in}}%
\pgfpathlineto{\pgfqpoint{1.981999in}{7.064825in}}%
\pgfpathlineto{\pgfqpoint{1.988331in}{5.634126in}}%
\pgfpathlineto{\pgfqpoint{1.988759in}{5.624582in}}%
\pgfpathlineto{\pgfqpoint{1.989272in}{5.634290in}}%
\pgfpathlineto{\pgfqpoint{1.990384in}{5.732635in}}%
\pgfpathlineto{\pgfqpoint{1.992524in}{6.173488in}}%
\pgfpathlineto{\pgfqpoint{1.999112in}{7.646787in}}%
\pgfpathlineto{\pgfqpoint{1.999540in}{7.653709in}}%
\pgfpathlineto{\pgfqpoint{1.999967in}{7.644769in}}%
\pgfpathlineto{\pgfqpoint{2.001080in}{7.549076in}}%
\pgfpathlineto{\pgfqpoint{2.003219in}{7.114759in}}%
\pgfpathlineto{\pgfqpoint{2.009893in}{5.633426in}}%
\pgfpathlineto{\pgfqpoint{2.010235in}{5.627173in}}%
\pgfpathlineto{\pgfqpoint{2.010748in}{5.636621in}}%
\pgfpathlineto{\pgfqpoint{2.011861in}{5.732774in}}%
\pgfpathlineto{\pgfqpoint{2.014000in}{6.164971in}}%
\pgfpathlineto{\pgfqpoint{2.020759in}{7.646379in}}%
\pgfpathlineto{\pgfqpoint{2.021101in}{7.651123in}}%
\pgfpathlineto{\pgfqpoint{2.021529in}{7.643086in}}%
\pgfpathlineto{\pgfqpoint{2.022556in}{7.561774in}}%
\pgfpathlineto{\pgfqpoint{2.024609in}{7.171079in}}%
\pgfpathlineto{\pgfqpoint{2.031797in}{5.630905in}}%
\pgfpathlineto{\pgfqpoint{2.031968in}{5.629741in}}%
\pgfpathlineto{\pgfqpoint{2.032310in}{5.634782in}}%
\pgfpathlineto{\pgfqpoint{2.033251in}{5.698532in}}%
\pgfpathlineto{\pgfqpoint{2.035134in}{6.022483in}}%
\pgfpathlineto{\pgfqpoint{2.042920in}{7.648410in}}%
\pgfpathlineto{\pgfqpoint{2.043177in}{7.644968in}}%
\pgfpathlineto{\pgfqpoint{2.044032in}{7.594495in}}%
\pgfpathlineto{\pgfqpoint{2.045743in}{7.328652in}}%
\pgfpathlineto{\pgfqpoint{2.050022in}{6.181101in}}%
\pgfpathlineto{\pgfqpoint{2.053187in}{5.651174in}}%
\pgfpathlineto{\pgfqpoint{2.053872in}{5.632440in}}%
\pgfpathlineto{\pgfqpoint{2.054385in}{5.643570in}}%
\pgfpathlineto{\pgfqpoint{2.055498in}{5.739921in}}%
\pgfpathlineto{\pgfqpoint{2.057637in}{6.161370in}}%
\pgfpathlineto{\pgfqpoint{2.064567in}{7.641344in}}%
\pgfpathlineto{\pgfqpoint{2.064909in}{7.645724in}}%
\pgfpathlineto{\pgfqpoint{2.065337in}{7.637844in}}%
\pgfpathlineto{\pgfqpoint{2.066364in}{7.559605in}}%
\pgfpathlineto{\pgfqpoint{2.068417in}{7.183859in}}%
\pgfpathlineto{\pgfqpoint{2.075861in}{5.635790in}}%
\pgfpathlineto{\pgfqpoint{2.075947in}{5.635245in}}%
\pgfpathlineto{\pgfqpoint{2.076289in}{5.638935in}}%
\pgfpathlineto{\pgfqpoint{2.077145in}{5.688776in}}%
\pgfpathlineto{\pgfqpoint{2.078856in}{5.947999in}}%
\pgfpathlineto{\pgfqpoint{2.082963in}{7.029006in}}%
\pgfpathlineto{\pgfqpoint{2.086300in}{7.615912in}}%
\pgfpathlineto{\pgfqpoint{2.087156in}{7.642917in}}%
\pgfpathlineto{\pgfqpoint{2.087583in}{7.634710in}}%
\pgfpathlineto{\pgfqpoint{2.088696in}{7.547030in}}%
\pgfpathlineto{\pgfqpoint{2.090835in}{7.146299in}}%
\pgfpathlineto{\pgfqpoint{2.098022in}{5.641666in}}%
\pgfpathlineto{\pgfqpoint{2.098364in}{5.638055in}}%
\pgfpathlineto{\pgfqpoint{2.098792in}{5.646440in}}%
\pgfpathlineto{\pgfqpoint{2.099904in}{5.733803in}}%
\pgfpathlineto{\pgfqpoint{2.102044in}{6.131342in}}%
\pgfpathlineto{\pgfqpoint{2.109316in}{7.637294in}}%
\pgfpathlineto{\pgfqpoint{2.109573in}{7.640102in}}%
\pgfpathlineto{\pgfqpoint{2.110001in}{7.633444in}}%
\pgfpathlineto{\pgfqpoint{2.111028in}{7.560703in}}%
\pgfpathlineto{\pgfqpoint{2.112995in}{7.223155in}}%
\pgfpathlineto{\pgfqpoint{2.120867in}{5.640919in}}%
\pgfpathlineto{\pgfqpoint{2.121038in}{5.641584in}}%
\pgfpathlineto{\pgfqpoint{2.121637in}{5.661495in}}%
\pgfpathlineto{\pgfqpoint{2.123006in}{5.805510in}}%
\pgfpathlineto{\pgfqpoint{2.125659in}{6.387413in}}%
\pgfpathlineto{\pgfqpoint{2.131306in}{7.601699in}}%
\pgfpathlineto{\pgfqpoint{2.132247in}{7.637173in}}%
\pgfpathlineto{\pgfqpoint{2.132760in}{7.628385in}}%
\pgfpathlineto{\pgfqpoint{2.133873in}{7.542489in}}%
\pgfpathlineto{\pgfqpoint{2.136012in}{7.155245in}}%
\pgfpathlineto{\pgfqpoint{2.143456in}{5.646258in}}%
\pgfpathlineto{\pgfqpoint{2.143712in}{5.643804in}}%
\pgfpathlineto{\pgfqpoint{2.144140in}{5.650654in}}%
\pgfpathlineto{\pgfqpoint{2.145167in}{5.721864in}}%
\pgfpathlineto{\pgfqpoint{2.147135in}{6.050014in}}%
\pgfpathlineto{\pgfqpoint{2.155178in}{7.634214in}}%
\pgfpathlineto{\pgfqpoint{2.155349in}{7.633502in}}%
\pgfpathlineto{\pgfqpoint{2.156033in}{7.609138in}}%
\pgfpathlineto{\pgfqpoint{2.157488in}{7.448387in}}%
\pgfpathlineto{\pgfqpoint{2.160312in}{6.819221in}}%
\pgfpathlineto{\pgfqpoint{2.165616in}{5.694455in}}%
\pgfpathlineto{\pgfqpoint{2.166729in}{5.646841in}}%
\pgfpathlineto{\pgfqpoint{2.167157in}{5.652432in}}%
\pgfpathlineto{\pgfqpoint{2.168183in}{5.719419in}}%
\pgfpathlineto{\pgfqpoint{2.170151in}{6.036371in}}%
\pgfpathlineto{\pgfqpoint{2.178365in}{7.631211in}}%
\pgfpathlineto{\pgfqpoint{2.178622in}{7.629140in}}%
\pgfpathlineto{\pgfqpoint{2.179392in}{7.594664in}}%
\pgfpathlineto{\pgfqpoint{2.181018in}{7.391410in}}%
\pgfpathlineto{\pgfqpoint{2.184269in}{6.625639in}}%
\pgfpathlineto{\pgfqpoint{2.188804in}{5.705772in}}%
\pgfpathlineto{\pgfqpoint{2.190087in}{5.649876in}}%
\pgfpathlineto{\pgfqpoint{2.190430in}{5.654646in}}%
\pgfpathlineto{\pgfqpoint{2.191371in}{5.710132in}}%
\pgfpathlineto{\pgfqpoint{2.193253in}{5.990197in}}%
\pgfpathlineto{\pgfqpoint{2.201809in}{7.628131in}}%
\pgfpathlineto{\pgfqpoint{2.202408in}{7.615443in}}%
\pgfpathlineto{\pgfqpoint{2.203606in}{7.516556in}}%
\pgfpathlineto{\pgfqpoint{2.205916in}{7.092474in}}%
\pgfpathlineto{\pgfqpoint{2.213104in}{5.662069in}}%
\pgfpathlineto{\pgfqpoint{2.213617in}{5.652976in}}%
\pgfpathlineto{\pgfqpoint{2.214130in}{5.662161in}}%
\pgfpathlineto{\pgfqpoint{2.215243in}{5.743419in}}%
\pgfpathlineto{\pgfqpoint{2.217382in}{6.105059in}}%
\pgfpathlineto{\pgfqpoint{2.225339in}{7.624200in}}%
\pgfpathlineto{\pgfqpoint{2.225510in}{7.624961in}}%
\pgfpathlineto{\pgfqpoint{2.225852in}{7.620462in}}%
\pgfpathlineto{\pgfqpoint{2.226794in}{7.567262in}}%
\pgfpathlineto{\pgfqpoint{2.228676in}{7.297343in}}%
\pgfpathlineto{\pgfqpoint{2.233296in}{6.181111in}}%
\pgfpathlineto{\pgfqpoint{2.236548in}{5.682372in}}%
\pgfpathlineto{\pgfqpoint{2.237403in}{5.656188in}}%
\pgfpathlineto{\pgfqpoint{2.237917in}{5.664198in}}%
\pgfpathlineto{\pgfqpoint{2.239029in}{5.741483in}}%
\pgfpathlineto{\pgfqpoint{2.241082in}{6.073716in}}%
\pgfpathlineto{\pgfqpoint{2.249468in}{7.621719in}}%
\pgfpathlineto{\pgfqpoint{2.249895in}{7.614702in}}%
\pgfpathlineto{\pgfqpoint{2.251008in}{7.540380in}}%
\pgfpathlineto{\pgfqpoint{2.253061in}{7.214973in}}%
\pgfpathlineto{\pgfqpoint{2.261532in}{5.659439in}}%
\pgfpathlineto{\pgfqpoint{2.261617in}{5.659803in}}%
\pgfpathlineto{\pgfqpoint{2.262216in}{5.675830in}}%
\pgfpathlineto{\pgfqpoint{2.263500in}{5.787039in}}%
\pgfpathlineto{\pgfqpoint{2.265896in}{6.227181in}}%
\pgfpathlineto{\pgfqpoint{2.272912in}{7.600608in}}%
\pgfpathlineto{\pgfqpoint{2.273682in}{7.618433in}}%
\pgfpathlineto{\pgfqpoint{2.274195in}{7.608907in}}%
\pgfpathlineto{\pgfqpoint{2.275393in}{7.521597in}}%
\pgfpathlineto{\pgfqpoint{2.277618in}{7.150610in}}%
\pgfpathlineto{\pgfqpoint{2.285575in}{5.665463in}}%
\pgfpathlineto{\pgfqpoint{2.285832in}{5.662788in}}%
\pgfpathlineto{\pgfqpoint{2.286345in}{5.670132in}}%
\pgfpathlineto{\pgfqpoint{2.287457in}{5.743005in}}%
\pgfpathlineto{\pgfqpoint{2.289511in}{6.058891in}}%
\pgfpathlineto{\pgfqpoint{2.298153in}{7.615095in}}%
\pgfpathlineto{\pgfqpoint{2.298238in}{7.614839in}}%
\pgfpathlineto{\pgfqpoint{2.298752in}{7.603578in}}%
\pgfpathlineto{\pgfqpoint{2.299949in}{7.514093in}}%
\pgfpathlineto{\pgfqpoint{2.302174in}{7.145621in}}%
\pgfpathlineto{\pgfqpoint{2.310131in}{5.670460in}}%
\pgfpathlineto{\pgfqpoint{2.310474in}{5.666178in}}%
\pgfpathlineto{\pgfqpoint{2.310987in}{5.673489in}}%
\pgfpathlineto{\pgfqpoint{2.312099in}{5.744802in}}%
\pgfpathlineto{\pgfqpoint{2.314153in}{6.053499in}}%
\pgfpathlineto{\pgfqpoint{2.322966in}{7.611616in}}%
\pgfpathlineto{\pgfqpoint{2.323051in}{7.611196in}}%
\pgfpathlineto{\pgfqpoint{2.323650in}{7.595645in}}%
\pgfpathlineto{\pgfqpoint{2.325019in}{7.479992in}}%
\pgfpathlineto{\pgfqpoint{2.327501in}{7.036012in}}%
\pgfpathlineto{\pgfqpoint{2.334602in}{5.690371in}}%
\pgfpathlineto{\pgfqpoint{2.335458in}{5.669653in}}%
\pgfpathlineto{\pgfqpoint{2.335886in}{5.675978in}}%
\pgfpathlineto{\pgfqpoint{2.336912in}{5.735855in}}%
\pgfpathlineto{\pgfqpoint{2.338880in}{6.009820in}}%
\pgfpathlineto{\pgfqpoint{2.348036in}{7.608112in}}%
\pgfpathlineto{\pgfqpoint{2.348463in}{7.601780in}}%
\pgfpathlineto{\pgfqpoint{2.349576in}{7.534827in}}%
\pgfpathlineto{\pgfqpoint{2.351629in}{7.239900in}}%
\pgfpathlineto{\pgfqpoint{2.360613in}{5.673224in}}%
\pgfpathlineto{\pgfqpoint{2.360870in}{5.674584in}}%
\pgfpathlineto{\pgfqpoint{2.361640in}{5.701934in}}%
\pgfpathlineto{\pgfqpoint{2.363180in}{5.856185in}}%
\pgfpathlineto{\pgfqpoint{2.366089in}{6.422574in}}%
\pgfpathlineto{\pgfqpoint{2.371993in}{7.549011in}}%
\pgfpathlineto{\pgfqpoint{2.373362in}{7.604547in}}%
\pgfpathlineto{\pgfqpoint{2.373704in}{7.601442in}}%
\pgfpathlineto{\pgfqpoint{2.374645in}{7.558011in}}%
\pgfpathlineto{\pgfqpoint{2.376442in}{7.343340in}}%
\pgfpathlineto{\pgfqpoint{2.380036in}{6.577422in}}%
\pgfpathlineto{\pgfqpoint{2.384656in}{5.743071in}}%
\pgfpathlineto{\pgfqpoint{2.386196in}{5.676816in}}%
\pgfpathlineto{\pgfqpoint{2.386453in}{5.678943in}}%
\pgfpathlineto{\pgfqpoint{2.387309in}{5.713205in}}%
\pgfpathlineto{\pgfqpoint{2.389020in}{5.899523in}}%
\pgfpathlineto{\pgfqpoint{2.392271in}{6.559392in}}%
\pgfpathlineto{\pgfqpoint{2.397491in}{7.531493in}}%
\pgfpathlineto{\pgfqpoint{2.399031in}{7.600856in}}%
\pgfpathlineto{\pgfqpoint{2.399373in}{7.598144in}}%
\pgfpathlineto{\pgfqpoint{2.400229in}{7.562527in}}%
\pgfpathlineto{\pgfqpoint{2.401940in}{7.375460in}}%
\pgfpathlineto{\pgfqpoint{2.405277in}{6.699922in}}%
\pgfpathlineto{\pgfqpoint{2.410411in}{5.752515in}}%
\pgfpathlineto{\pgfqpoint{2.412036in}{5.680524in}}%
\pgfpathlineto{\pgfqpoint{2.412293in}{5.682554in}}%
\pgfpathlineto{\pgfqpoint{2.413149in}{5.715734in}}%
\pgfpathlineto{\pgfqpoint{2.414860in}{5.896803in}}%
\pgfpathlineto{\pgfqpoint{2.418111in}{6.541703in}}%
\pgfpathlineto{\pgfqpoint{2.423502in}{7.530388in}}%
\pgfpathlineto{\pgfqpoint{2.425042in}{7.597117in}}%
\pgfpathlineto{\pgfqpoint{2.425384in}{7.594326in}}%
\pgfpathlineto{\pgfqpoint{2.426325in}{7.553664in}}%
\pgfpathlineto{\pgfqpoint{2.428122in}{7.350972in}}%
\pgfpathlineto{\pgfqpoint{2.431630in}{6.637268in}}%
\pgfpathlineto{\pgfqpoint{2.436593in}{5.752825in}}%
\pgfpathlineto{\pgfqpoint{2.438218in}{5.684338in}}%
\pgfpathlineto{\pgfqpoint{2.438475in}{5.686546in}}%
\pgfpathlineto{\pgfqpoint{2.439331in}{5.719563in}}%
\pgfpathlineto{\pgfqpoint{2.441042in}{5.897054in}}%
\pgfpathlineto{\pgfqpoint{2.444293in}{6.528595in}}%
\pgfpathlineto{\pgfqpoint{2.449769in}{7.523261in}}%
\pgfpathlineto{\pgfqpoint{2.451395in}{7.593296in}}%
\pgfpathlineto{\pgfqpoint{2.451652in}{7.591529in}}%
\pgfpathlineto{\pgfqpoint{2.452507in}{7.560332in}}%
\pgfpathlineto{\pgfqpoint{2.454133in}{7.399561in}}%
\pgfpathlineto{\pgfqpoint{2.457213in}{6.823096in}}%
\pgfpathlineto{\pgfqpoint{2.463117in}{5.752485in}}%
\pgfpathlineto{\pgfqpoint{2.464657in}{5.688213in}}%
\pgfpathlineto{\pgfqpoint{2.464999in}{5.690796in}}%
\pgfpathlineto{\pgfqpoint{2.465855in}{5.724082in}}%
\pgfpathlineto{\pgfqpoint{2.467566in}{5.898806in}}%
\pgfpathlineto{\pgfqpoint{2.470818in}{6.517810in}}%
\pgfpathlineto{\pgfqpoint{2.476379in}{7.516632in}}%
\pgfpathlineto{\pgfqpoint{2.478090in}{7.589351in}}%
\pgfpathlineto{\pgfqpoint{2.478262in}{7.588328in}}%
\pgfpathlineto{\pgfqpoint{2.479032in}{7.564979in}}%
\pgfpathlineto{\pgfqpoint{2.480572in}{7.430227in}}%
\pgfpathlineto{\pgfqpoint{2.483395in}{6.942037in}}%
\pgfpathlineto{\pgfqpoint{2.490155in}{5.739813in}}%
\pgfpathlineto{\pgfqpoint{2.491524in}{5.692159in}}%
\pgfpathlineto{\pgfqpoint{2.491866in}{5.695166in}}%
\pgfpathlineto{\pgfqpoint{2.492807in}{5.734065in}}%
\pgfpathlineto{\pgfqpoint{2.494604in}{5.924807in}}%
\pgfpathlineto{\pgfqpoint{2.498027in}{6.581900in}}%
\pgfpathlineto{\pgfqpoint{2.503331in}{7.509685in}}%
\pgfpathlineto{\pgfqpoint{2.505043in}{7.585338in}}%
\pgfpathlineto{\pgfqpoint{2.505299in}{7.584075in}}%
\pgfpathlineto{\pgfqpoint{2.506069in}{7.560444in}}%
\pgfpathlineto{\pgfqpoint{2.507610in}{7.427796in}}%
\pgfpathlineto{\pgfqpoint{2.510433in}{6.949949in}}%
\pgfpathlineto{\pgfqpoint{2.517364in}{5.741140in}}%
\pgfpathlineto{\pgfqpoint{2.518733in}{5.696226in}}%
\pgfpathlineto{\pgfqpoint{2.519075in}{5.699477in}}%
\pgfpathlineto{\pgfqpoint{2.520016in}{5.738117in}}%
\pgfpathlineto{\pgfqpoint{2.521813in}{5.924897in}}%
\pgfpathlineto{\pgfqpoint{2.525235in}{6.567878in}}%
\pgfpathlineto{\pgfqpoint{2.530711in}{7.508795in}}%
\pgfpathlineto{\pgfqpoint{2.532423in}{7.581250in}}%
\pgfpathlineto{\pgfqpoint{2.532679in}{7.579881in}}%
\pgfpathlineto{\pgfqpoint{2.533449in}{7.556538in}}%
\pgfpathlineto{\pgfqpoint{2.534990in}{7.427010in}}%
\pgfpathlineto{\pgfqpoint{2.537813in}{6.960696in}}%
\pgfpathlineto{\pgfqpoint{2.544915in}{5.743890in}}%
\pgfpathlineto{\pgfqpoint{2.546284in}{5.700376in}}%
\pgfpathlineto{\pgfqpoint{2.546626in}{5.703537in}}%
\pgfpathlineto{\pgfqpoint{2.547567in}{5.741030in}}%
\pgfpathlineto{\pgfqpoint{2.549364in}{5.922385in}}%
\pgfpathlineto{\pgfqpoint{2.552787in}{6.549815in}}%
\pgfpathlineto{\pgfqpoint{2.558348in}{7.498692in}}%
\pgfpathlineto{\pgfqpoint{2.560145in}{7.577037in}}%
\pgfpathlineto{\pgfqpoint{2.560316in}{7.576605in}}%
\pgfpathlineto{\pgfqpoint{2.560915in}{7.564183in}}%
\pgfpathlineto{\pgfqpoint{2.562284in}{7.473827in}}%
\pgfpathlineto{\pgfqpoint{2.564680in}{7.136854in}}%
\pgfpathlineto{\pgfqpoint{2.573749in}{5.709057in}}%
\pgfpathlineto{\pgfqpoint{2.574177in}{5.704613in}}%
\pgfpathlineto{\pgfqpoint{2.574691in}{5.710555in}}%
\pgfpathlineto{\pgfqpoint{2.575803in}{5.764972in}}%
\pgfpathlineto{\pgfqpoint{2.577856in}{6.001437in}}%
\pgfpathlineto{\pgfqpoint{2.582134in}{6.826797in}}%
\pgfpathlineto{\pgfqpoint{2.586584in}{7.506577in}}%
\pgfpathlineto{\pgfqpoint{2.588295in}{7.572767in}}%
\pgfpathlineto{\pgfqpoint{2.588552in}{7.571194in}}%
\pgfpathlineto{\pgfqpoint{2.589407in}{7.544236in}}%
\pgfpathlineto{\pgfqpoint{2.591033in}{7.405357in}}%
\pgfpathlineto{\pgfqpoint{2.594028in}{6.914923in}}%
\pgfpathlineto{\pgfqpoint{2.600958in}{5.761665in}}%
\pgfpathlineto{\pgfqpoint{2.602498in}{5.708956in}}%
\pgfpathlineto{\pgfqpoint{2.602755in}{5.710527in}}%
\pgfpathlineto{\pgfqpoint{2.603611in}{5.737130in}}%
\pgfpathlineto{\pgfqpoint{2.605236in}{5.873986in}}%
\pgfpathlineto{\pgfqpoint{2.608231in}{6.357882in}}%
\pgfpathlineto{\pgfqpoint{2.615247in}{7.515761in}}%
\pgfpathlineto{\pgfqpoint{2.616787in}{7.568373in}}%
\pgfpathlineto{\pgfqpoint{2.617044in}{7.566954in}}%
\pgfpathlineto{\pgfqpoint{2.617900in}{7.541195in}}%
\pgfpathlineto{\pgfqpoint{2.619525in}{7.407208in}}%
\pgfpathlineto{\pgfqpoint{2.622434in}{6.947263in}}%
\pgfpathlineto{\pgfqpoint{2.629707in}{5.761256in}}%
\pgfpathlineto{\pgfqpoint{2.631162in}{5.713420in}}%
\pgfpathlineto{\pgfqpoint{2.631504in}{5.715537in}}%
\pgfpathlineto{\pgfqpoint{2.632360in}{5.743091in}}%
\pgfpathlineto{\pgfqpoint{2.634071in}{5.888754in}}%
\pgfpathlineto{\pgfqpoint{2.637151in}{6.385955in}}%
\pgfpathlineto{\pgfqpoint{2.643996in}{7.502488in}}%
\pgfpathlineto{\pgfqpoint{2.645707in}{7.563872in}}%
\pgfpathlineto{\pgfqpoint{2.645878in}{7.563137in}}%
\pgfpathlineto{\pgfqpoint{2.646563in}{7.547644in}}%
\pgfpathlineto{\pgfqpoint{2.648018in}{7.450128in}}%
\pgfpathlineto{\pgfqpoint{2.650584in}{7.096132in}}%
\pgfpathlineto{\pgfqpoint{2.659483in}{5.731994in}}%
\pgfpathlineto{\pgfqpoint{2.660339in}{5.717985in}}%
\pgfpathlineto{\pgfqpoint{2.660766in}{5.722587in}}%
\pgfpathlineto{\pgfqpoint{2.661879in}{5.770294in}}%
\pgfpathlineto{\pgfqpoint{2.663932in}{5.983182in}}%
\pgfpathlineto{\pgfqpoint{2.667868in}{6.683900in}}%
\pgfpathlineto{\pgfqpoint{2.673002in}{7.476659in}}%
\pgfpathlineto{\pgfqpoint{2.674970in}{7.559243in}}%
\pgfpathlineto{\pgfqpoint{2.675141in}{7.558904in}}%
\pgfpathlineto{\pgfqpoint{2.675740in}{7.548142in}}%
\pgfpathlineto{\pgfqpoint{2.677109in}{7.468994in}}%
\pgfpathlineto{\pgfqpoint{2.679505in}{7.170396in}}%
\pgfpathlineto{\pgfqpoint{2.689601in}{5.723623in}}%
\pgfpathlineto{\pgfqpoint{2.689858in}{5.722631in}}%
\pgfpathlineto{\pgfqpoint{2.690200in}{5.725503in}}%
\pgfpathlineto{\pgfqpoint{2.691141in}{5.757877in}}%
\pgfpathlineto{\pgfqpoint{2.692938in}{5.913722in}}%
\pgfpathlineto{\pgfqpoint{2.696189in}{6.431967in}}%
\pgfpathlineto{\pgfqpoint{2.702863in}{7.483956in}}%
\pgfpathlineto{\pgfqpoint{2.704745in}{7.554560in}}%
\pgfpathlineto{\pgfqpoint{2.704917in}{7.553985in}}%
\pgfpathlineto{\pgfqpoint{2.705601in}{7.539945in}}%
\pgfpathlineto{\pgfqpoint{2.707056in}{7.449581in}}%
\pgfpathlineto{\pgfqpoint{2.709623in}{7.117780in}}%
\pgfpathlineto{\pgfqpoint{2.719120in}{5.736005in}}%
\pgfpathlineto{\pgfqpoint{2.719804in}{5.727387in}}%
\pgfpathlineto{\pgfqpoint{2.720318in}{5.733070in}}%
\pgfpathlineto{\pgfqpoint{2.721430in}{5.780589in}}%
\pgfpathlineto{\pgfqpoint{2.723484in}{5.984796in}}%
\pgfpathlineto{\pgfqpoint{2.727419in}{6.654174in}}%
\pgfpathlineto{\pgfqpoint{2.732810in}{7.464022in}}%
\pgfpathlineto{\pgfqpoint{2.734949in}{7.549724in}}%
\pgfpathlineto{\pgfqpoint{2.735377in}{7.545781in}}%
\pgfpathlineto{\pgfqpoint{2.736404in}{7.507648in}}%
\pgfpathlineto{\pgfqpoint{2.738371in}{7.329003in}}%
\pgfpathlineto{\pgfqpoint{2.741965in}{6.747914in}}%
\pgfpathlineto{\pgfqpoint{2.748040in}{5.818668in}}%
\pgfpathlineto{\pgfqpoint{2.750179in}{5.732259in}}%
\pgfpathlineto{\pgfqpoint{2.750521in}{5.734441in}}%
\pgfpathlineto{\pgfqpoint{2.751463in}{5.763309in}}%
\pgfpathlineto{\pgfqpoint{2.753259in}{5.906710in}}%
\pgfpathlineto{\pgfqpoint{2.756425in}{6.377801in}}%
\pgfpathlineto{\pgfqpoint{2.763698in}{7.480764in}}%
\pgfpathlineto{\pgfqpoint{2.765580in}{7.544782in}}%
\pgfpathlineto{\pgfqpoint{2.765751in}{7.544074in}}%
\pgfpathlineto{\pgfqpoint{2.766522in}{7.527357in}}%
\pgfpathlineto{\pgfqpoint{2.768062in}{7.429692in}}%
\pgfpathlineto{\pgfqpoint{2.770714in}{7.091186in}}%
\pgfpathlineto{\pgfqpoint{2.780126in}{5.752767in}}%
\pgfpathlineto{\pgfqpoint{2.781067in}{5.737270in}}%
\pgfpathlineto{\pgfqpoint{2.781495in}{5.740996in}}%
\pgfpathlineto{\pgfqpoint{2.782522in}{5.777134in}}%
\pgfpathlineto{\pgfqpoint{2.784404in}{5.936957in}}%
\pgfpathlineto{\pgfqpoint{2.787827in}{6.457546in}}%
\pgfpathlineto{\pgfqpoint{2.794500in}{7.457187in}}%
\pgfpathlineto{\pgfqpoint{2.796640in}{7.539730in}}%
\pgfpathlineto{\pgfqpoint{2.796982in}{7.537761in}}%
\pgfpathlineto{\pgfqpoint{2.797923in}{7.510668in}}%
\pgfpathlineto{\pgfqpoint{2.799720in}{7.375014in}}%
\pgfpathlineto{\pgfqpoint{2.802886in}{6.925582in}}%
\pgfpathlineto{\pgfqpoint{2.810586in}{5.798941in}}%
\pgfpathlineto{\pgfqpoint{2.812383in}{5.742383in}}%
\pgfpathlineto{\pgfqpoint{2.812554in}{5.742928in}}%
\pgfpathlineto{\pgfqpoint{2.813239in}{5.755473in}}%
\pgfpathlineto{\pgfqpoint{2.814693in}{5.835732in}}%
\pgfpathlineto{\pgfqpoint{2.817175in}{6.119839in}}%
\pgfpathlineto{\pgfqpoint{2.827955in}{7.533248in}}%
\pgfpathlineto{\pgfqpoint{2.828212in}{7.534543in}}%
\pgfpathlineto{\pgfqpoint{2.828640in}{7.531596in}}%
\pgfpathlineto{\pgfqpoint{2.829667in}{7.498740in}}%
\pgfpathlineto{\pgfqpoint{2.831549in}{7.349379in}}%
\pgfpathlineto{\pgfqpoint{2.834886in}{6.868954in}}%
\pgfpathlineto{\pgfqpoint{2.842159in}{5.818220in}}%
\pgfpathlineto{\pgfqpoint{2.844212in}{5.747633in}}%
\pgfpathlineto{\pgfqpoint{2.844298in}{5.747789in}}%
\pgfpathlineto{\pgfqpoint{2.844897in}{5.755886in}}%
\pgfpathlineto{\pgfqpoint{2.846180in}{5.813735in}}%
\pgfpathlineto{\pgfqpoint{2.848490in}{6.043692in}}%
\pgfpathlineto{\pgfqpoint{2.853111in}{6.790299in}}%
\pgfpathlineto{\pgfqpoint{2.857988in}{7.440721in}}%
\pgfpathlineto{\pgfqpoint{2.860212in}{7.529104in}}%
\pgfpathlineto{\pgfqpoint{2.860298in}{7.529235in}}%
\pgfpathlineto{\pgfqpoint{2.860555in}{7.528153in}}%
\pgfpathlineto{\pgfqpoint{2.861410in}{7.508647in}}%
\pgfpathlineto{\pgfqpoint{2.863036in}{7.406672in}}%
\pgfpathlineto{\pgfqpoint{2.865860in}{7.060113in}}%
\pgfpathlineto{\pgfqpoint{2.875357in}{5.775533in}}%
\pgfpathlineto{\pgfqpoint{2.876555in}{5.753020in}}%
\pgfpathlineto{\pgfqpoint{2.876983in}{5.756415in}}%
\pgfpathlineto{\pgfqpoint{2.878009in}{5.788918in}}%
\pgfpathlineto{\pgfqpoint{2.879892in}{5.932763in}}%
\pgfpathlineto{\pgfqpoint{2.883229in}{6.393537in}}%
\pgfpathlineto{\pgfqpoint{2.890758in}{7.451552in}}%
\pgfpathlineto{\pgfqpoint{2.892897in}{7.523789in}}%
\pgfpathlineto{\pgfqpoint{2.893240in}{7.521757in}}%
\pgfpathlineto{\pgfqpoint{2.894181in}{7.496786in}}%
\pgfpathlineto{\pgfqpoint{2.895978in}{7.373813in}}%
\pgfpathlineto{\pgfqpoint{2.899058in}{6.977668in}}%
\pgfpathlineto{\pgfqpoint{2.907614in}{5.807241in}}%
\pgfpathlineto{\pgfqpoint{2.909411in}{5.758541in}}%
\pgfpathlineto{\pgfqpoint{2.909582in}{5.759208in}}%
\pgfpathlineto{\pgfqpoint{2.910352in}{5.773656in}}%
\pgfpathlineto{\pgfqpoint{2.911892in}{5.857112in}}%
\pgfpathlineto{\pgfqpoint{2.914545in}{6.148654in}}%
\pgfpathlineto{\pgfqpoint{2.925411in}{7.512884in}}%
\pgfpathlineto{\pgfqpoint{2.926010in}{7.518209in}}%
\pgfpathlineto{\pgfqpoint{2.926523in}{7.513895in}}%
\pgfpathlineto{\pgfqpoint{2.927636in}{7.476767in}}%
\pgfpathlineto{\pgfqpoint{2.929689in}{7.314863in}}%
\pgfpathlineto{\pgfqpoint{2.933283in}{6.816421in}}%
\pgfpathlineto{\pgfqpoint{2.940384in}{5.848663in}}%
\pgfpathlineto{\pgfqpoint{2.942695in}{5.764224in}}%
\pgfpathlineto{\pgfqpoint{2.942780in}{5.764178in}}%
\pgfpathlineto{\pgfqpoint{2.942951in}{5.764755in}}%
\pgfpathlineto{\pgfqpoint{2.943721in}{5.778366in}}%
\pgfpathlineto{\pgfqpoint{2.945262in}{5.858140in}}%
\pgfpathlineto{\pgfqpoint{2.947914in}{6.138703in}}%
\pgfpathlineto{\pgfqpoint{2.959208in}{7.509543in}}%
\pgfpathlineto{\pgfqpoint{2.959636in}{7.512489in}}%
\pgfpathlineto{\pgfqpoint{2.960149in}{7.508799in}}%
\pgfpathlineto{\pgfqpoint{2.961262in}{7.474065in}}%
\pgfpathlineto{\pgfqpoint{2.963230in}{7.328331in}}%
\pgfpathlineto{\pgfqpoint{2.966652in}{6.878922in}}%
\pgfpathlineto{\pgfqpoint{2.974438in}{5.843339in}}%
\pgfpathlineto{\pgfqpoint{2.976663in}{5.769970in}}%
\pgfpathlineto{\pgfqpoint{2.976920in}{5.770734in}}%
\pgfpathlineto{\pgfqpoint{2.977690in}{5.784573in}}%
\pgfpathlineto{\pgfqpoint{2.979230in}{5.862732in}}%
\pgfpathlineto{\pgfqpoint{2.981882in}{6.135195in}}%
\pgfpathlineto{\pgfqpoint{2.993433in}{7.503891in}}%
\pgfpathlineto{\pgfqpoint{2.993861in}{7.506628in}}%
\pgfpathlineto{\pgfqpoint{2.994374in}{7.502972in}}%
\pgfpathlineto{\pgfqpoint{2.995487in}{7.469359in}}%
\pgfpathlineto{\pgfqpoint{2.997455in}{7.328771in}}%
\pgfpathlineto{\pgfqpoint{3.000877in}{6.893951in}}%
\pgfpathlineto{\pgfqpoint{3.008920in}{5.847843in}}%
\pgfpathlineto{\pgfqpoint{3.011145in}{5.775932in}}%
\pgfpathlineto{\pgfqpoint{3.011401in}{5.776498in}}%
\pgfpathlineto{\pgfqpoint{3.012171in}{5.789292in}}%
\pgfpathlineto{\pgfqpoint{3.013711in}{5.863457in}}%
\pgfpathlineto{\pgfqpoint{3.016278in}{6.114035in}}%
\pgfpathlineto{\pgfqpoint{3.022353in}{7.005025in}}%
\pgfpathlineto{\pgfqpoint{3.026546in}{7.439228in}}%
\pgfpathlineto{\pgfqpoint{3.028685in}{7.500612in}}%
\pgfpathlineto{\pgfqpoint{3.029113in}{7.497870in}}%
\pgfpathlineto{\pgfqpoint{3.030139in}{7.470877in}}%
\pgfpathlineto{\pgfqpoint{3.032022in}{7.350186in}}%
\pgfpathlineto{\pgfqpoint{3.035273in}{6.967333in}}%
\pgfpathlineto{\pgfqpoint{3.044343in}{5.832666in}}%
\pgfpathlineto{\pgfqpoint{3.046311in}{5.781993in}}%
\pgfpathlineto{\pgfqpoint{3.046396in}{5.782145in}}%
\pgfpathlineto{\pgfqpoint{3.046995in}{5.788741in}}%
\pgfpathlineto{\pgfqpoint{3.048279in}{5.834946in}}%
\pgfpathlineto{\pgfqpoint{3.050503in}{6.010614in}}%
\pgfpathlineto{\pgfqpoint{3.054525in}{6.543563in}}%
\pgfpathlineto{\pgfqpoint{3.061284in}{7.394250in}}%
\pgfpathlineto{\pgfqpoint{3.063766in}{7.493297in}}%
\pgfpathlineto{\pgfqpoint{3.064108in}{7.494421in}}%
\pgfpathlineto{\pgfqpoint{3.064536in}{7.491466in}}%
\pgfpathlineto{\pgfqpoint{3.065648in}{7.461347in}}%
\pgfpathlineto{\pgfqpoint{3.067616in}{7.332867in}}%
\pgfpathlineto{\pgfqpoint{3.070953in}{6.941178in}}%
\pgfpathlineto{\pgfqpoint{3.079851in}{5.846785in}}%
\pgfpathlineto{\pgfqpoint{3.081990in}{5.788230in}}%
\pgfpathlineto{\pgfqpoint{3.082333in}{5.789772in}}%
\pgfpathlineto{\pgfqpoint{3.083274in}{5.809567in}}%
\pgfpathlineto{\pgfqpoint{3.085071in}{5.908272in}}%
\pgfpathlineto{\pgfqpoint{3.088065in}{6.222318in}}%
\pgfpathlineto{\pgfqpoint{3.098761in}{7.466410in}}%
\pgfpathlineto{\pgfqpoint{3.100044in}{7.488116in}}%
\pgfpathlineto{\pgfqpoint{3.100472in}{7.486087in}}%
\pgfpathlineto{\pgfqpoint{3.101499in}{7.462412in}}%
\pgfpathlineto{\pgfqpoint{3.103381in}{7.353104in}}%
\pgfpathlineto{\pgfqpoint{3.106547in}{7.011023in}}%
\pgfpathlineto{\pgfqpoint{3.116558in}{5.832672in}}%
\pgfpathlineto{\pgfqpoint{3.118354in}{5.794657in}}%
\pgfpathlineto{\pgfqpoint{3.118525in}{5.795197in}}%
\pgfpathlineto{\pgfqpoint{3.119296in}{5.806594in}}%
\pgfpathlineto{\pgfqpoint{3.120836in}{5.872351in}}%
\pgfpathlineto{\pgfqpoint{3.123403in}{6.095630in}}%
\pgfpathlineto{\pgfqpoint{3.128536in}{6.787923in}}%
\pgfpathlineto{\pgfqpoint{3.133841in}{7.382942in}}%
\pgfpathlineto{\pgfqpoint{3.136408in}{7.480482in}}%
\pgfpathlineto{\pgfqpoint{3.136750in}{7.481623in}}%
\pgfpathlineto{\pgfqpoint{3.137178in}{7.479046in}}%
\pgfpathlineto{\pgfqpoint{3.138205in}{7.454863in}}%
\pgfpathlineto{\pgfqpoint{3.140087in}{7.347485in}}%
\pgfpathlineto{\pgfqpoint{3.143253in}{7.015117in}}%
\pgfpathlineto{\pgfqpoint{3.153521in}{5.837940in}}%
\pgfpathlineto{\pgfqpoint{3.155317in}{5.801221in}}%
\pgfpathlineto{\pgfqpoint{3.155488in}{5.801709in}}%
\pgfpathlineto{\pgfqpoint{3.156259in}{5.812495in}}%
\pgfpathlineto{\pgfqpoint{3.157799in}{5.875261in}}%
\pgfpathlineto{\pgfqpoint{3.160366in}{6.089355in}}%
\pgfpathlineto{\pgfqpoint{3.165328in}{6.736624in}}%
\pgfpathlineto{\pgfqpoint{3.170975in}{7.369445in}}%
\pgfpathlineto{\pgfqpoint{3.173628in}{7.473226in}}%
\pgfpathlineto{\pgfqpoint{3.174056in}{7.474963in}}%
\pgfpathlineto{\pgfqpoint{3.174483in}{7.472442in}}%
\pgfpathlineto{\pgfqpoint{3.175596in}{7.446133in}}%
\pgfpathlineto{\pgfqpoint{3.177564in}{7.333024in}}%
\pgfpathlineto{\pgfqpoint{3.180815in}{6.994094in}}%
\pgfpathlineto{\pgfqpoint{3.190997in}{5.849650in}}%
\pgfpathlineto{\pgfqpoint{3.192879in}{5.807966in}}%
\pgfpathlineto{\pgfqpoint{3.193050in}{5.808150in}}%
\pgfpathlineto{\pgfqpoint{3.193649in}{5.814038in}}%
\pgfpathlineto{\pgfqpoint{3.195018in}{5.857637in}}%
\pgfpathlineto{\pgfqpoint{3.197329in}{6.019287in}}%
\pgfpathlineto{\pgfqpoint{3.201350in}{6.491425in}}%
\pgfpathlineto{\pgfqpoint{3.208879in}{7.362519in}}%
\pgfpathlineto{\pgfqpoint{3.211532in}{7.465816in}}%
\pgfpathlineto{\pgfqpoint{3.211960in}{7.468142in}}%
\pgfpathlineto{\pgfqpoint{3.212473in}{7.465554in}}%
\pgfpathlineto{\pgfqpoint{3.213585in}{7.439992in}}%
\pgfpathlineto{\pgfqpoint{3.215553in}{7.331010in}}%
\pgfpathlineto{\pgfqpoint{3.218805in}{7.004356in}}%
\pgfpathlineto{\pgfqpoint{3.229329in}{5.854198in}}%
\pgfpathlineto{\pgfqpoint{3.231211in}{5.814874in}}%
\pgfpathlineto{\pgfqpoint{3.231382in}{5.815101in}}%
\pgfpathlineto{\pgfqpoint{3.232067in}{5.822375in}}%
\pgfpathlineto{\pgfqpoint{3.233521in}{5.871034in}}%
\pgfpathlineto{\pgfqpoint{3.235917in}{6.041624in}}%
\pgfpathlineto{\pgfqpoint{3.240110in}{6.530447in}}%
\pgfpathlineto{\pgfqpoint{3.247383in}{7.349340in}}%
\pgfpathlineto{\pgfqpoint{3.250121in}{7.458075in}}%
\pgfpathlineto{\pgfqpoint{3.250634in}{7.461138in}}%
\pgfpathlineto{\pgfqpoint{3.251147in}{7.458592in}}%
\pgfpathlineto{\pgfqpoint{3.252260in}{7.433992in}}%
\pgfpathlineto{\pgfqpoint{3.254228in}{7.329414in}}%
\pgfpathlineto{\pgfqpoint{3.257479in}{7.015249in}}%
\pgfpathlineto{\pgfqpoint{3.268431in}{5.856496in}}%
\pgfpathlineto{\pgfqpoint{3.270228in}{5.821976in}}%
\pgfpathlineto{\pgfqpoint{3.270399in}{5.822174in}}%
\pgfpathlineto{\pgfqpoint{3.271083in}{5.829049in}}%
\pgfpathlineto{\pgfqpoint{3.272452in}{5.871526in}}%
\pgfpathlineto{\pgfqpoint{3.274848in}{6.031107in}}%
\pgfpathlineto{\pgfqpoint{3.278955in}{6.487759in}}%
\pgfpathlineto{\pgfqpoint{3.286827in}{7.351034in}}%
\pgfpathlineto{\pgfqpoint{3.289479in}{7.450770in}}%
\pgfpathlineto{\pgfqpoint{3.290078in}{7.453937in}}%
\pgfpathlineto{\pgfqpoint{3.290592in}{7.450846in}}%
\pgfpathlineto{\pgfqpoint{3.291789in}{7.422999in}}%
\pgfpathlineto{\pgfqpoint{3.293843in}{7.311419in}}%
\pgfpathlineto{\pgfqpoint{3.297265in}{6.979858in}}%
\pgfpathlineto{\pgfqpoint{3.307790in}{5.877540in}}%
\pgfpathlineto{\pgfqpoint{3.310014in}{5.829252in}}%
\pgfpathlineto{\pgfqpoint{3.310442in}{5.831166in}}%
\pgfpathlineto{\pgfqpoint{3.311469in}{5.850501in}}%
\pgfpathlineto{\pgfqpoint{3.313351in}{5.937962in}}%
\pgfpathlineto{\pgfqpoint{3.316431in}{6.205189in}}%
\pgfpathlineto{\pgfqpoint{3.328924in}{7.431920in}}%
\pgfpathlineto{\pgfqpoint{3.330121in}{7.446564in}}%
\pgfpathlineto{\pgfqpoint{3.330549in}{7.445049in}}%
\pgfpathlineto{\pgfqpoint{3.331576in}{7.426999in}}%
\pgfpathlineto{\pgfqpoint{3.333458in}{7.342947in}}%
\pgfpathlineto{\pgfqpoint{3.336453in}{7.092009in}}%
\pgfpathlineto{\pgfqpoint{3.349544in}{5.845232in}}%
\pgfpathlineto{\pgfqpoint{3.350485in}{5.836727in}}%
\pgfpathlineto{\pgfqpoint{3.350999in}{5.839168in}}%
\pgfpathlineto{\pgfqpoint{3.352111in}{5.861474in}}%
\pgfpathlineto{\pgfqpoint{3.354079in}{5.955573in}}%
\pgfpathlineto{\pgfqpoint{3.357330in}{6.239683in}}%
\pgfpathlineto{\pgfqpoint{3.369395in}{7.414605in}}%
\pgfpathlineto{\pgfqpoint{3.371020in}{7.439000in}}%
\pgfpathlineto{\pgfqpoint{3.371277in}{7.438388in}}%
\pgfpathlineto{\pgfqpoint{3.372133in}{7.427559in}}%
\pgfpathlineto{\pgfqpoint{3.373758in}{7.370595in}}%
\pgfpathlineto{\pgfqpoint{3.376411in}{7.185166in}}%
\pgfpathlineto{\pgfqpoint{3.381202in}{6.659656in}}%
\pgfpathlineto{\pgfqpoint{3.388047in}{5.965095in}}%
\pgfpathlineto{\pgfqpoint{3.390956in}{5.849979in}}%
\pgfpathlineto{\pgfqpoint{3.391726in}{5.844401in}}%
\pgfpathlineto{\pgfqpoint{3.392240in}{5.846644in}}%
\pgfpathlineto{\pgfqpoint{3.393352in}{5.867747in}}%
\pgfpathlineto{\pgfqpoint{3.395320in}{5.957325in}}%
\pgfpathlineto{\pgfqpoint{3.398486in}{6.220512in}}%
\pgfpathlineto{\pgfqpoint{3.411235in}{7.413354in}}%
\pgfpathlineto{\pgfqpoint{3.412689in}{7.431220in}}%
\pgfpathlineto{\pgfqpoint{3.413031in}{7.429999in}}%
\pgfpathlineto{\pgfqpoint{3.413973in}{7.416015in}}%
\pgfpathlineto{\pgfqpoint{3.415769in}{7.347322in}}%
\pgfpathlineto{\pgfqpoint{3.418679in}{7.132800in}}%
\pgfpathlineto{\pgfqpoint{3.424326in}{6.508387in}}%
\pgfpathlineto{\pgfqpoint{3.430144in}{5.963503in}}%
\pgfpathlineto{\pgfqpoint{3.433053in}{5.856739in}}%
\pgfpathlineto{\pgfqpoint{3.433738in}{5.852282in}}%
\pgfpathlineto{\pgfqpoint{3.434251in}{5.854246in}}%
\pgfpathlineto{\pgfqpoint{3.435363in}{5.873999in}}%
\pgfpathlineto{\pgfqpoint{3.437331in}{5.958896in}}%
\pgfpathlineto{\pgfqpoint{3.440497in}{6.210172in}}%
\pgfpathlineto{\pgfqpoint{3.453845in}{7.410284in}}%
\pgfpathlineto{\pgfqpoint{3.455128in}{7.423239in}}%
\pgfpathlineto{\pgfqpoint{3.455470in}{7.422011in}}%
\pgfpathlineto{\pgfqpoint{3.456412in}{7.408503in}}%
\pgfpathlineto{\pgfqpoint{3.458208in}{7.342645in}}%
\pgfpathlineto{\pgfqpoint{3.461118in}{7.137037in}}%
\pgfpathlineto{\pgfqpoint{3.466594in}{6.552856in}}%
\pgfpathlineto{\pgfqpoint{3.472754in}{5.979329in}}%
\pgfpathlineto{\pgfqpoint{3.475749in}{5.866332in}}%
\pgfpathlineto{\pgfqpoint{3.476604in}{5.860349in}}%
\pgfpathlineto{\pgfqpoint{3.477118in}{5.862539in}}%
\pgfpathlineto{\pgfqpoint{3.478230in}{5.882051in}}%
\pgfpathlineto{\pgfqpoint{3.480198in}{5.964182in}}%
\pgfpathlineto{\pgfqpoint{3.483364in}{6.206077in}}%
\pgfpathlineto{\pgfqpoint{3.497139in}{7.403878in}}%
\pgfpathlineto{\pgfqpoint{3.498337in}{7.415070in}}%
\pgfpathlineto{\pgfqpoint{3.498765in}{7.413487in}}%
\pgfpathlineto{\pgfqpoint{3.499792in}{7.397765in}}%
\pgfpathlineto{\pgfqpoint{3.501674in}{7.326634in}}%
\pgfpathlineto{\pgfqpoint{3.504669in}{7.114280in}}%
\pgfpathlineto{\pgfqpoint{3.510573in}{6.496278in}}%
\pgfpathlineto{\pgfqpoint{3.516391in}{5.982497in}}%
\pgfpathlineto{\pgfqpoint{3.519386in}{5.874428in}}%
\pgfpathlineto{\pgfqpoint{3.520241in}{5.868636in}}%
\pgfpathlineto{\pgfqpoint{3.520755in}{5.870668in}}%
\pgfpathlineto{\pgfqpoint{3.521867in}{5.889145in}}%
\pgfpathlineto{\pgfqpoint{3.523835in}{5.967287in}}%
\pgfpathlineto{\pgfqpoint{3.527001in}{6.198436in}}%
\pgfpathlineto{\pgfqpoint{3.534188in}{6.946101in}}%
\pgfpathlineto{\pgfqpoint{3.539065in}{7.324065in}}%
\pgfpathlineto{\pgfqpoint{3.541803in}{7.404155in}}%
\pgfpathlineto{\pgfqpoint{3.542402in}{7.406675in}}%
\pgfpathlineto{\pgfqpoint{3.542915in}{7.404466in}}%
\pgfpathlineto{\pgfqpoint{3.544028in}{7.385949in}}%
\pgfpathlineto{\pgfqpoint{3.545996in}{7.308845in}}%
\pgfpathlineto{\pgfqpoint{3.549161in}{7.081794in}}%
\pgfpathlineto{\pgfqpoint{3.556177in}{6.362343in}}%
\pgfpathlineto{\pgfqpoint{3.561226in}{5.966636in}}%
\pgfpathlineto{\pgfqpoint{3.564049in}{5.880479in}}%
\pgfpathlineto{\pgfqpoint{3.564734in}{5.877138in}}%
\pgfpathlineto{\pgfqpoint{3.565247in}{5.879223in}}%
\pgfpathlineto{\pgfqpoint{3.566359in}{5.897146in}}%
\pgfpathlineto{\pgfqpoint{3.568327in}{5.972181in}}%
\pgfpathlineto{\pgfqpoint{3.571493in}{6.193870in}}%
\pgfpathlineto{\pgfqpoint{3.578167in}{6.867596in}}%
\pgfpathlineto{\pgfqpoint{3.583557in}{7.299600in}}%
\pgfpathlineto{\pgfqpoint{3.586467in}{7.393371in}}%
\pgfpathlineto{\pgfqpoint{3.587237in}{7.398063in}}%
\pgfpathlineto{\pgfqpoint{3.587750in}{7.396392in}}%
\pgfpathlineto{\pgfqpoint{3.588862in}{7.379677in}}%
\pgfpathlineto{\pgfqpoint{3.590830in}{7.307711in}}%
\pgfpathlineto{\pgfqpoint{3.593910in}{7.099787in}}%
\pgfpathlineto{\pgfqpoint{3.599985in}{6.500095in}}%
\pgfpathlineto{\pgfqpoint{3.605975in}{6.000274in}}%
\pgfpathlineto{\pgfqpoint{3.609055in}{5.892653in}}%
\pgfpathlineto{\pgfqpoint{3.609996in}{5.885862in}}%
\pgfpathlineto{\pgfqpoint{3.610510in}{5.887465in}}%
\pgfpathlineto{\pgfqpoint{3.611622in}{5.903718in}}%
\pgfpathlineto{\pgfqpoint{3.613590in}{5.973876in}}%
\pgfpathlineto{\pgfqpoint{3.616670in}{6.177005in}}%
\pgfpathlineto{\pgfqpoint{3.622659in}{6.757508in}}%
\pgfpathlineto{\pgfqpoint{3.628820in}{7.270384in}}%
\pgfpathlineto{\pgfqpoint{3.631900in}{7.380711in}}%
\pgfpathlineto{\pgfqpoint{3.633013in}{7.389245in}}%
\pgfpathlineto{\pgfqpoint{3.633440in}{7.387959in}}%
\pgfpathlineto{\pgfqpoint{3.634467in}{7.374565in}}%
\pgfpathlineto{\pgfqpoint{3.636349in}{7.313297in}}%
\pgfpathlineto{\pgfqpoint{3.639344in}{7.128475in}}%
\pgfpathlineto{\pgfqpoint{3.644735in}{6.625383in}}%
\pgfpathlineto{\pgfqpoint{3.651751in}{6.025241in}}%
\pgfpathlineto{\pgfqpoint{3.654917in}{5.905934in}}%
\pgfpathlineto{\pgfqpoint{3.656200in}{5.894788in}}%
\pgfpathlineto{\pgfqpoint{3.656542in}{5.895579in}}%
\pgfpathlineto{\pgfqpoint{3.657483in}{5.905900in}}%
\pgfpathlineto{\pgfqpoint{3.659195in}{5.954586in}}%
\pgfpathlineto{\pgfqpoint{3.661933in}{6.105805in}}%
\pgfpathlineto{\pgfqpoint{3.666553in}{6.507033in}}%
\pgfpathlineto{\pgfqpoint{3.675366in}{7.264254in}}%
\pgfpathlineto{\pgfqpoint{3.678446in}{7.371320in}}%
\pgfpathlineto{\pgfqpoint{3.679644in}{7.380190in}}%
\pgfpathlineto{\pgfqpoint{3.679986in}{7.379244in}}%
\pgfpathlineto{\pgfqpoint{3.680928in}{7.368697in}}%
\pgfpathlineto{\pgfqpoint{3.682724in}{7.316972in}}%
\pgfpathlineto{\pgfqpoint{3.685548in}{7.159584in}}%
\pgfpathlineto{\pgfqpoint{3.690339in}{6.744040in}}%
\pgfpathlineto{\pgfqpoint{3.698810in}{6.026517in}}%
\pgfpathlineto{\pgfqpoint{3.701976in}{5.914026in}}%
\pgfpathlineto{\pgfqpoint{3.703259in}{5.903952in}}%
\pgfpathlineto{\pgfqpoint{3.703602in}{5.904849in}}%
\pgfpathlineto{\pgfqpoint{3.704543in}{5.915070in}}%
\pgfpathlineto{\pgfqpoint{3.706340in}{5.965431in}}%
\pgfpathlineto{\pgfqpoint{3.709163in}{6.119031in}}%
\pgfpathlineto{\pgfqpoint{3.713955in}{6.526161in}}%
\pgfpathlineto{\pgfqpoint{3.722596in}{7.248586in}}%
\pgfpathlineto{\pgfqpoint{3.725848in}{7.361637in}}%
\pgfpathlineto{\pgfqpoint{3.727046in}{7.370930in}}%
\pgfpathlineto{\pgfqpoint{3.727473in}{7.369877in}}%
\pgfpathlineto{\pgfqpoint{3.728500in}{7.358000in}}%
\pgfpathlineto{\pgfqpoint{3.730383in}{7.302839in}}%
\pgfpathlineto{\pgfqpoint{3.733292in}{7.140905in}}%
\pgfpathlineto{\pgfqpoint{3.738340in}{6.710056in}}%
\pgfpathlineto{\pgfqpoint{3.746468in}{6.042353in}}%
\pgfpathlineto{\pgfqpoint{3.749720in}{5.925286in}}%
\pgfpathlineto{\pgfqpoint{3.751089in}{5.913328in}}%
\pgfpathlineto{\pgfqpoint{3.751431in}{5.913929in}}%
\pgfpathlineto{\pgfqpoint{3.752372in}{5.922974in}}%
\pgfpathlineto{\pgfqpoint{3.754083in}{5.966618in}}%
\pgfpathlineto{\pgfqpoint{3.756821in}{6.103601in}}%
\pgfpathlineto{\pgfqpoint{3.761271in}{6.456990in}}%
\pgfpathlineto{\pgfqpoint{3.771110in}{7.254414in}}%
\pgfpathlineto{\pgfqpoint{3.774276in}{7.354078in}}%
\pgfpathlineto{\pgfqpoint{3.775388in}{7.361442in}}%
\pgfpathlineto{\pgfqpoint{3.775816in}{7.360329in}}%
\pgfpathlineto{\pgfqpoint{3.776843in}{7.348752in}}%
\pgfpathlineto{\pgfqpoint{3.778725in}{7.295717in}}%
\pgfpathlineto{\pgfqpoint{3.781635in}{7.140476in}}%
\pgfpathlineto{\pgfqpoint{3.786597in}{6.733644in}}%
\pgfpathlineto{\pgfqpoint{3.795153in}{6.049418in}}%
\pgfpathlineto{\pgfqpoint{3.798490in}{5.934018in}}%
\pgfpathlineto{\pgfqpoint{3.799859in}{5.922924in}}%
\pgfpathlineto{\pgfqpoint{3.800202in}{5.923572in}}%
\pgfpathlineto{\pgfqpoint{3.801143in}{5.932392in}}%
\pgfpathlineto{\pgfqpoint{3.802854in}{5.974347in}}%
\pgfpathlineto{\pgfqpoint{3.805592in}{6.105613in}}%
\pgfpathlineto{\pgfqpoint{3.810041in}{6.445169in}}%
\pgfpathlineto{\pgfqpoint{3.820309in}{7.250473in}}%
\pgfpathlineto{\pgfqpoint{3.823389in}{7.343824in}}%
\pgfpathlineto{\pgfqpoint{3.824587in}{7.351729in}}%
\pgfpathlineto{\pgfqpoint{3.825015in}{7.350585in}}%
\pgfpathlineto{\pgfqpoint{3.826041in}{7.339354in}}%
\pgfpathlineto{\pgfqpoint{3.827924in}{7.288442in}}%
\pgfpathlineto{\pgfqpoint{3.830833in}{7.139700in}}%
\pgfpathlineto{\pgfqpoint{3.835710in}{6.755890in}}%
\pgfpathlineto{\pgfqpoint{3.844694in}{6.056808in}}%
\pgfpathlineto{\pgfqpoint{3.848031in}{5.944382in}}%
\pgfpathlineto{\pgfqpoint{3.849486in}{5.932745in}}%
\pgfpathlineto{\pgfqpoint{3.849828in}{5.933429in}}%
\pgfpathlineto{\pgfqpoint{3.850769in}{5.942020in}}%
\pgfpathlineto{\pgfqpoint{3.852480in}{5.982348in}}%
\pgfpathlineto{\pgfqpoint{3.855218in}{6.108153in}}%
\pgfpathlineto{\pgfqpoint{3.859667in}{6.434420in}}%
\pgfpathlineto{\pgfqpoint{3.870363in}{7.246132in}}%
\pgfpathlineto{\pgfqpoint{3.873443in}{7.334530in}}%
\pgfpathlineto{\pgfqpoint{3.874641in}{7.341793in}}%
\pgfpathlineto{\pgfqpoint{3.874983in}{7.341001in}}%
\pgfpathlineto{\pgfqpoint{3.876010in}{7.331009in}}%
\pgfpathlineto{\pgfqpoint{3.877807in}{7.286638in}}%
\pgfpathlineto{\pgfqpoint{3.880630in}{7.153518in}}%
\pgfpathlineto{\pgfqpoint{3.885251in}{6.812249in}}%
\pgfpathlineto{\pgfqpoint{3.895347in}{6.051796in}}%
\pgfpathlineto{\pgfqpoint{3.898598in}{5.952210in}}%
\pgfpathlineto{\pgfqpoint{3.899967in}{5.942789in}}%
\pgfpathlineto{\pgfqpoint{3.900310in}{5.943543in}}%
\pgfpathlineto{\pgfqpoint{3.901251in}{5.952014in}}%
\pgfpathlineto{\pgfqpoint{3.903048in}{5.993702in}}%
\pgfpathlineto{\pgfqpoint{3.905871in}{6.121573in}}%
\pgfpathlineto{\pgfqpoint{3.910406in}{6.446716in}}%
\pgfpathlineto{\pgfqpoint{3.921016in}{7.230943in}}%
\pgfpathlineto{\pgfqpoint{3.924182in}{7.323051in}}%
\pgfpathlineto{\pgfqpoint{3.925465in}{7.331657in}}%
\pgfpathlineto{\pgfqpoint{3.925807in}{7.331066in}}%
\pgfpathlineto{\pgfqpoint{3.926748in}{7.323190in}}%
\pgfpathlineto{\pgfqpoint{3.928460in}{7.285821in}}%
\pgfpathlineto{\pgfqpoint{3.931198in}{7.168618in}}%
\pgfpathlineto{\pgfqpoint{3.935561in}{6.868972in}}%
\pgfpathlineto{\pgfqpoint{3.947112in}{6.036759in}}%
\pgfpathlineto{\pgfqpoint{3.950193in}{5.958309in}}%
\pgfpathlineto{\pgfqpoint{3.951219in}{5.953033in}}%
\pgfpathlineto{\pgfqpoint{3.951647in}{5.953989in}}%
\pgfpathlineto{\pgfqpoint{3.952674in}{5.963828in}}%
\pgfpathlineto{\pgfqpoint{3.954556in}{6.008885in}}%
\pgfpathlineto{\pgfqpoint{3.957465in}{6.141464in}}%
\pgfpathlineto{\pgfqpoint{3.962171in}{6.474874in}}%
\pgfpathlineto{\pgfqpoint{3.972439in}{7.213638in}}%
\pgfpathlineto{\pgfqpoint{3.975776in}{7.311863in}}%
\pgfpathlineto{\pgfqpoint{3.977145in}{7.321302in}}%
\pgfpathlineto{\pgfqpoint{3.977487in}{7.320761in}}%
\pgfpathlineto{\pgfqpoint{3.978428in}{7.313305in}}%
\pgfpathlineto{\pgfqpoint{3.980139in}{7.277720in}}%
\pgfpathlineto{\pgfqpoint{3.982877in}{7.165775in}}%
\pgfpathlineto{\pgfqpoint{3.987156in}{6.884747in}}%
\pgfpathlineto{\pgfqpoint{3.999305in}{6.039430in}}%
\pgfpathlineto{\pgfqpoint{4.002300in}{5.968309in}}%
\pgfpathlineto{\pgfqpoint{4.003327in}{5.963494in}}%
\pgfpathlineto{\pgfqpoint{4.003755in}{5.964502in}}%
\pgfpathlineto{\pgfqpoint{4.004781in}{5.974128in}}%
\pgfpathlineto{\pgfqpoint{4.006664in}{6.017594in}}%
\pgfpathlineto{\pgfqpoint{4.009573in}{6.145047in}}%
\pgfpathlineto{\pgfqpoint{4.014279in}{6.466167in}}%
\pgfpathlineto{\pgfqpoint{4.024889in}{7.206679in}}%
\pgfpathlineto{\pgfqpoint{4.028226in}{7.301391in}}%
\pgfpathlineto{\pgfqpoint{4.029680in}{7.310734in}}%
\pgfpathlineto{\pgfqpoint{4.029937in}{7.310303in}}%
\pgfpathlineto{\pgfqpoint{4.030878in}{7.303396in}}%
\pgfpathlineto{\pgfqpoint{4.032589in}{7.269754in}}%
\pgfpathlineto{\pgfqpoint{4.035327in}{7.163124in}}%
\pgfpathlineto{\pgfqpoint{4.039605in}{6.893680in}}%
\pgfpathlineto{\pgfqpoint{4.052269in}{6.043032in}}%
\pgfpathlineto{\pgfqpoint{4.055263in}{5.977817in}}%
\pgfpathlineto{\pgfqpoint{4.056204in}{5.974157in}}%
\pgfpathlineto{\pgfqpoint{4.056632in}{5.975208in}}%
\pgfpathlineto{\pgfqpoint{4.057745in}{5.985844in}}%
\pgfpathlineto{\pgfqpoint{4.059713in}{6.031936in}}%
\pgfpathlineto{\pgfqpoint{4.062707in}{6.163042in}}%
\pgfpathlineto{\pgfqpoint{4.067584in}{6.491107in}}%
\pgfpathlineto{\pgfqpoint{4.077852in}{7.189236in}}%
\pgfpathlineto{\pgfqpoint{4.081274in}{7.288635in}}%
\pgfpathlineto{\pgfqpoint{4.082900in}{7.299987in}}%
\pgfpathlineto{\pgfqpoint{4.083157in}{7.299590in}}%
\pgfpathlineto{\pgfqpoint{4.084012in}{7.293962in}}%
\pgfpathlineto{\pgfqpoint{4.085638in}{7.265279in}}%
\pgfpathlineto{\pgfqpoint{4.088290in}{7.170854in}}%
\pgfpathlineto{\pgfqpoint{4.092312in}{6.935724in}}%
\pgfpathlineto{\pgfqpoint{4.106515in}{6.030712in}}%
\pgfpathlineto{\pgfqpoint{4.109168in}{5.986443in}}%
\pgfpathlineto{\pgfqpoint{4.109767in}{5.984991in}}%
\pgfpathlineto{\pgfqpoint{4.110280in}{5.986283in}}%
\pgfpathlineto{\pgfqpoint{4.111392in}{5.997072in}}%
\pgfpathlineto{\pgfqpoint{4.113360in}{6.042268in}}%
\pgfpathlineto{\pgfqpoint{4.116440in}{6.173992in}}%
\pgfpathlineto{\pgfqpoint{4.121403in}{6.499526in}}%
\pgfpathlineto{\pgfqpoint{4.131585in}{7.174344in}}%
\pgfpathlineto{\pgfqpoint{4.135093in}{7.276528in}}%
\pgfpathlineto{\pgfqpoint{4.136804in}{7.289061in}}%
\pgfpathlineto{\pgfqpoint{4.136975in}{7.288916in}}%
\pgfpathlineto{\pgfqpoint{4.137745in}{7.285115in}}%
\pgfpathlineto{\pgfqpoint{4.139286in}{7.262228in}}%
\pgfpathlineto{\pgfqpoint{4.141767in}{7.184545in}}%
\pgfpathlineto{\pgfqpoint{4.145532in}{6.985205in}}%
\pgfpathlineto{\pgfqpoint{4.161874in}{6.015353in}}%
\pgfpathlineto{\pgfqpoint{4.164013in}{5.996006in}}%
\pgfpathlineto{\pgfqpoint{4.164527in}{5.997150in}}%
\pgfpathlineto{\pgfqpoint{4.165639in}{6.007298in}}%
\pgfpathlineto{\pgfqpoint{4.167607in}{6.050344in}}%
\pgfpathlineto{\pgfqpoint{4.170601in}{6.172142in}}%
\pgfpathlineto{\pgfqpoint{4.175393in}{6.471873in}}%
\pgfpathlineto{\pgfqpoint{4.186345in}{7.174631in}}%
\pgfpathlineto{\pgfqpoint{4.189767in}{7.267099in}}%
\pgfpathlineto{\pgfqpoint{4.191393in}{7.277961in}}%
\pgfpathlineto{\pgfqpoint{4.191650in}{7.277660in}}%
\pgfpathlineto{\pgfqpoint{4.192505in}{7.272687in}}%
\pgfpathlineto{\pgfqpoint{4.194131in}{7.246645in}}%
\pgfpathlineto{\pgfqpoint{4.196698in}{7.163640in}}%
\pgfpathlineto{\pgfqpoint{4.200634in}{6.953470in}}%
\pgfpathlineto{\pgfqpoint{4.216121in}{6.038678in}}%
\pgfpathlineto{\pgfqpoint{4.218602in}{6.007555in}}%
\pgfpathlineto{\pgfqpoint{4.218944in}{6.007193in}}%
\pgfpathlineto{\pgfqpoint{4.219372in}{6.008090in}}%
\pgfpathlineto{\pgfqpoint{4.220484in}{6.017415in}}%
\pgfpathlineto{\pgfqpoint{4.222452in}{6.058091in}}%
\pgfpathlineto{\pgfqpoint{4.225447in}{6.174413in}}%
\pgfpathlineto{\pgfqpoint{4.230238in}{6.463012in}}%
\pgfpathlineto{\pgfqpoint{4.241618in}{7.170682in}}%
\pgfpathlineto{\pgfqpoint{4.244955in}{7.256327in}}%
\pgfpathlineto{\pgfqpoint{4.246581in}{7.266707in}}%
\pgfpathlineto{\pgfqpoint{4.246838in}{7.266405in}}%
\pgfpathlineto{\pgfqpoint{4.247693in}{7.261578in}}%
\pgfpathlineto{\pgfqpoint{4.249319in}{7.236428in}}%
\pgfpathlineto{\pgfqpoint{4.251886in}{7.156342in}}%
\pgfpathlineto{\pgfqpoint{4.255822in}{6.953326in}}%
\pgfpathlineto{\pgfqpoint{4.271651in}{6.047211in}}%
\pgfpathlineto{\pgfqpoint{4.274132in}{6.018705in}}%
\pgfpathlineto{\pgfqpoint{4.274474in}{6.018565in}}%
\pgfpathlineto{\pgfqpoint{4.274817in}{6.019350in}}%
\pgfpathlineto{\pgfqpoint{4.275929in}{6.028265in}}%
\pgfpathlineto{\pgfqpoint{4.277897in}{6.067357in}}%
\pgfpathlineto{\pgfqpoint{4.280891in}{6.179443in}}%
\pgfpathlineto{\pgfqpoint{4.285597in}{6.452815in}}%
\pgfpathlineto{\pgfqpoint{4.297405in}{7.165594in}}%
\pgfpathlineto{\pgfqpoint{4.300742in}{7.246292in}}%
\pgfpathlineto{\pgfqpoint{4.302282in}{7.255319in}}%
\pgfpathlineto{\pgfqpoint{4.302539in}{7.255036in}}%
\pgfpathlineto{\pgfqpoint{4.303394in}{7.250405in}}%
\pgfpathlineto{\pgfqpoint{4.305020in}{7.226183in}}%
\pgfpathlineto{\pgfqpoint{4.307587in}{7.148919in}}%
\pgfpathlineto{\pgfqpoint{4.311523in}{6.952643in}}%
\pgfpathlineto{\pgfqpoint{4.327694in}{6.055584in}}%
\pgfpathlineto{\pgfqpoint{4.330090in}{6.030114in}}%
\pgfpathlineto{\pgfqpoint{4.330432in}{6.030030in}}%
\pgfpathlineto{\pgfqpoint{4.330774in}{6.030840in}}%
\pgfpathlineto{\pgfqpoint{4.331887in}{6.039619in}}%
\pgfpathlineto{\pgfqpoint{4.333855in}{6.077679in}}%
\pgfpathlineto{\pgfqpoint{4.336849in}{6.186457in}}%
\pgfpathlineto{\pgfqpoint{4.341555in}{6.451797in}}%
\pgfpathlineto{\pgfqpoint{4.353534in}{7.156245in}}%
\pgfpathlineto{\pgfqpoint{4.356871in}{7.234834in}}%
\pgfpathlineto{\pgfqpoint{4.358411in}{7.243818in}}%
\pgfpathlineto{\pgfqpoint{4.358668in}{7.243586in}}%
\pgfpathlineto{\pgfqpoint{4.359523in}{7.239246in}}%
\pgfpathlineto{\pgfqpoint{4.361149in}{7.216070in}}%
\pgfpathlineto{\pgfqpoint{4.363716in}{7.141659in}}%
\pgfpathlineto{\pgfqpoint{4.367652in}{6.951855in}}%
\pgfpathlineto{\pgfqpoint{4.384165in}{6.063941in}}%
\pgfpathlineto{\pgfqpoint{4.386476in}{6.041602in}}%
\pgfpathlineto{\pgfqpoint{4.386818in}{6.041631in}}%
\pgfpathlineto{\pgfqpoint{4.387074in}{6.042220in}}%
\pgfpathlineto{\pgfqpoint{4.388101in}{6.049431in}}%
\pgfpathlineto{\pgfqpoint{4.389984in}{6.082437in}}%
\pgfpathlineto{\pgfqpoint{4.392893in}{6.180252in}}%
\pgfpathlineto{\pgfqpoint{4.397342in}{6.416359in}}%
\pgfpathlineto{\pgfqpoint{4.410604in}{7.165247in}}%
\pgfpathlineto{\pgfqpoint{4.413770in}{7.227320in}}%
\pgfpathlineto{\pgfqpoint{4.414968in}{7.232218in}}%
\pgfpathlineto{\pgfqpoint{4.415396in}{7.231436in}}%
\pgfpathlineto{\pgfqpoint{4.416508in}{7.223189in}}%
\pgfpathlineto{\pgfqpoint{4.418476in}{7.187079in}}%
\pgfpathlineto{\pgfqpoint{4.421471in}{7.083443in}}%
\pgfpathlineto{\pgfqpoint{4.426177in}{6.829595in}}%
\pgfpathlineto{\pgfqpoint{4.438412in}{6.136779in}}%
\pgfpathlineto{\pgfqpoint{4.441749in}{6.061693in}}%
\pgfpathlineto{\pgfqpoint{4.443289in}{6.053140in}}%
\pgfpathlineto{\pgfqpoint{4.443546in}{6.053368in}}%
\pgfpathlineto{\pgfqpoint{4.444401in}{6.057539in}}%
\pgfpathlineto{\pgfqpoint{4.446027in}{6.079744in}}%
\pgfpathlineto{\pgfqpoint{4.448594in}{6.151007in}}%
\pgfpathlineto{\pgfqpoint{4.452530in}{6.332970in}}%
\pgfpathlineto{\pgfqpoint{4.469300in}{7.200048in}}%
\pgfpathlineto{\pgfqpoint{4.471610in}{7.220521in}}%
\pgfpathlineto{\pgfqpoint{4.472038in}{7.220183in}}%
\pgfpathlineto{\pgfqpoint{4.472209in}{7.219685in}}%
\pgfpathlineto{\pgfqpoint{4.473321in}{7.211414in}}%
\pgfpathlineto{\pgfqpoint{4.475289in}{7.175861in}}%
\pgfpathlineto{\pgfqpoint{4.478284in}{7.074430in}}%
\pgfpathlineto{\pgfqpoint{4.482990in}{6.826524in}}%
\pgfpathlineto{\pgfqpoint{4.495311in}{6.145723in}}%
\pgfpathlineto{\pgfqpoint{4.498648in}{6.072970in}}%
\pgfpathlineto{\pgfqpoint{4.500188in}{6.064831in}}%
\pgfpathlineto{\pgfqpoint{4.500445in}{6.065085in}}%
\pgfpathlineto{\pgfqpoint{4.501300in}{6.069254in}}%
\pgfpathlineto{\pgfqpoint{4.502926in}{6.091083in}}%
\pgfpathlineto{\pgfqpoint{4.505493in}{6.160811in}}%
\pgfpathlineto{\pgfqpoint{4.509429in}{6.338567in}}%
\pgfpathlineto{\pgfqpoint{4.526285in}{7.189303in}}%
\pgfpathlineto{\pgfqpoint{4.528595in}{7.208830in}}%
\pgfpathlineto{\pgfqpoint{4.529194in}{7.207897in}}%
\pgfpathlineto{\pgfqpoint{4.530306in}{7.199614in}}%
\pgfpathlineto{\pgfqpoint{4.532274in}{7.164560in}}%
\pgfpathlineto{\pgfqpoint{4.535269in}{7.065079in}}%
\pgfpathlineto{\pgfqpoint{4.539975in}{6.822477in}}%
\pgfpathlineto{\pgfqpoint{4.552296in}{6.156293in}}%
\pgfpathlineto{\pgfqpoint{4.555633in}{6.084740in}}%
\pgfpathlineto{\pgfqpoint{4.557173in}{6.076553in}}%
\pgfpathlineto{\pgfqpoint{4.557429in}{6.076762in}}%
\pgfpathlineto{\pgfqpoint{4.558285in}{6.080705in}}%
\pgfpathlineto{\pgfqpoint{4.559911in}{6.101790in}}%
\pgfpathlineto{\pgfqpoint{4.562478in}{6.169578in}}%
\pgfpathlineto{\pgfqpoint{4.566413in}{6.342930in}}%
\pgfpathlineto{\pgfqpoint{4.575825in}{6.897846in}}%
\pgfpathlineto{\pgfqpoint{4.581130in}{7.127854in}}%
\pgfpathlineto{\pgfqpoint{4.584381in}{7.191235in}}%
\pgfpathlineto{\pgfqpoint{4.585750in}{7.197112in}}%
\pgfpathlineto{\pgfqpoint{4.586093in}{7.196603in}}%
\pgfpathlineto{\pgfqpoint{4.587119in}{7.190337in}}%
\pgfpathlineto{\pgfqpoint{4.588916in}{7.162562in}}%
\pgfpathlineto{\pgfqpoint{4.591740in}{7.078469in}}%
\pgfpathlineto{\pgfqpoint{4.596104in}{6.870332in}}%
\pgfpathlineto{\pgfqpoint{4.610307in}{6.138700in}}%
\pgfpathlineto{\pgfqpoint{4.613302in}{6.091077in}}%
\pgfpathlineto{\pgfqpoint{4.614243in}{6.088284in}}%
\pgfpathlineto{\pgfqpoint{4.614671in}{6.088976in}}%
\pgfpathlineto{\pgfqpoint{4.615783in}{6.096487in}}%
\pgfpathlineto{\pgfqpoint{4.617665in}{6.127632in}}%
\pgfpathlineto{\pgfqpoint{4.620574in}{6.218067in}}%
\pgfpathlineto{\pgfqpoint{4.625109in}{6.439374in}}%
\pgfpathlineto{\pgfqpoint{4.638200in}{7.119187in}}%
\pgfpathlineto{\pgfqpoint{4.641366in}{7.179404in}}%
\pgfpathlineto{\pgfqpoint{4.642735in}{7.185395in}}%
\pgfpathlineto{\pgfqpoint{4.643077in}{7.184951in}}%
\pgfpathlineto{\pgfqpoint{4.644104in}{7.178969in}}%
\pgfpathlineto{\pgfqpoint{4.645901in}{7.151990in}}%
\pgfpathlineto{\pgfqpoint{4.648724in}{7.069851in}}%
\pgfpathlineto{\pgfqpoint{4.653003in}{6.870688in}}%
\pgfpathlineto{\pgfqpoint{4.667377in}{6.146835in}}%
\pgfpathlineto{\pgfqpoint{4.670286in}{6.102595in}}%
\pgfpathlineto{\pgfqpoint{4.671227in}{6.100001in}}%
\pgfpathlineto{\pgfqpoint{4.671655in}{6.100749in}}%
\pgfpathlineto{\pgfqpoint{4.672768in}{6.108311in}}%
\pgfpathlineto{\pgfqpoint{4.674735in}{6.141148in}}%
\pgfpathlineto{\pgfqpoint{4.677730in}{6.235180in}}%
\pgfpathlineto{\pgfqpoint{4.682436in}{6.465495in}}%
\pgfpathlineto{\pgfqpoint{4.694757in}{7.098895in}}%
\pgfpathlineto{\pgfqpoint{4.698094in}{7.166336in}}%
\pgfpathlineto{\pgfqpoint{4.699634in}{7.173701in}}%
\pgfpathlineto{\pgfqpoint{4.699891in}{7.173424in}}%
\pgfpathlineto{\pgfqpoint{4.700832in}{7.168733in}}%
\pgfpathlineto{\pgfqpoint{4.702543in}{7.145612in}}%
\pgfpathlineto{\pgfqpoint{4.705196in}{7.074626in}}%
\pgfpathlineto{\pgfqpoint{4.709303in}{6.894876in}}%
\pgfpathlineto{\pgfqpoint{4.724875in}{6.142344in}}%
\pgfpathlineto{\pgfqpoint{4.727527in}{6.112277in}}%
\pgfpathlineto{\pgfqpoint{4.728041in}{6.111678in}}%
\pgfpathlineto{\pgfqpoint{4.728469in}{6.112485in}}%
\pgfpathlineto{\pgfqpoint{4.729581in}{6.120121in}}%
\pgfpathlineto{\pgfqpoint{4.731549in}{6.152804in}}%
\pgfpathlineto{\pgfqpoint{4.734544in}{6.245906in}}%
\pgfpathlineto{\pgfqpoint{4.739250in}{6.473152in}}%
\pgfpathlineto{\pgfqpoint{4.751399in}{7.087504in}}%
\pgfpathlineto{\pgfqpoint{4.754736in}{7.154560in}}%
\pgfpathlineto{\pgfqpoint{4.756276in}{7.162070in}}%
\pgfpathlineto{\pgfqpoint{4.756533in}{7.161837in}}%
\pgfpathlineto{\pgfqpoint{4.757389in}{7.157992in}}%
\pgfpathlineto{\pgfqpoint{4.759014in}{7.137854in}}%
\pgfpathlineto{\pgfqpoint{4.761581in}{7.073527in}}%
\pgfpathlineto{\pgfqpoint{4.765517in}{6.909672in}}%
\pgfpathlineto{\pgfqpoint{4.782031in}{6.142477in}}%
\pgfpathlineto{\pgfqpoint{4.784341in}{6.123322in}}%
\pgfpathlineto{\pgfqpoint{4.784683in}{6.123382in}}%
\pgfpathlineto{\pgfqpoint{4.785025in}{6.124193in}}%
\pgfpathlineto{\pgfqpoint{4.786223in}{6.132924in}}%
\pgfpathlineto{\pgfqpoint{4.788277in}{6.168702in}}%
\pgfpathlineto{\pgfqpoint{4.791357in}{6.267433in}}%
\pgfpathlineto{\pgfqpoint{4.796320in}{6.510581in}}%
\pgfpathlineto{\pgfqpoint{4.807357in}{7.064573in}}%
\pgfpathlineto{\pgfqpoint{4.810865in}{7.140744in}}%
\pgfpathlineto{\pgfqpoint{4.812662in}{7.150510in}}%
\pgfpathlineto{\pgfqpoint{4.812833in}{7.150367in}}%
\pgfpathlineto{\pgfqpoint{4.813689in}{7.146847in}}%
\pgfpathlineto{\pgfqpoint{4.815315in}{7.127439in}}%
\pgfpathlineto{\pgfqpoint{4.817881in}{7.064529in}}%
\pgfpathlineto{\pgfqpoint{4.821817in}{6.903366in}}%
\pgfpathlineto{\pgfqpoint{4.838160in}{6.154349in}}%
\pgfpathlineto{\pgfqpoint{4.840470in}{6.134862in}}%
\pgfpathlineto{\pgfqpoint{4.840812in}{6.134847in}}%
\pgfpathlineto{\pgfqpoint{4.841154in}{6.135578in}}%
\pgfpathlineto{\pgfqpoint{4.842267in}{6.143081in}}%
\pgfpathlineto{\pgfqpoint{4.844235in}{6.175154in}}%
\pgfpathlineto{\pgfqpoint{4.847229in}{6.266398in}}%
\pgfpathlineto{\pgfqpoint{4.851935in}{6.488374in}}%
\pgfpathlineto{\pgfqpoint{4.863743in}{7.066928in}}%
\pgfpathlineto{\pgfqpoint{4.867080in}{7.132051in}}%
\pgfpathlineto{\pgfqpoint{4.868620in}{7.139057in}}%
\pgfpathlineto{\pgfqpoint{4.868877in}{7.138764in}}%
\pgfpathlineto{\pgfqpoint{4.869818in}{7.134117in}}%
\pgfpathlineto{\pgfqpoint{4.871529in}{7.111504in}}%
\pgfpathlineto{\pgfqpoint{4.874267in}{7.039542in}}%
\pgfpathlineto{\pgfqpoint{4.878460in}{6.859741in}}%
\pgfpathlineto{\pgfqpoint{4.892920in}{6.183332in}}%
\pgfpathlineto{\pgfqpoint{4.895658in}{6.147747in}}%
\pgfpathlineto{\pgfqpoint{4.896428in}{6.146164in}}%
\pgfpathlineto{\pgfqpoint{4.896856in}{6.146907in}}%
\pgfpathlineto{\pgfqpoint{4.897968in}{6.154239in}}%
\pgfpathlineto{\pgfqpoint{4.899936in}{6.185910in}}%
\pgfpathlineto{\pgfqpoint{4.902931in}{6.276283in}}%
\pgfpathlineto{\pgfqpoint{4.907636in}{6.496074in}}%
\pgfpathlineto{\pgfqpoint{4.919187in}{7.055243in}}%
\pgfpathlineto{\pgfqpoint{4.922524in}{7.120527in}}%
\pgfpathlineto{\pgfqpoint{4.924064in}{7.127737in}}%
\pgfpathlineto{\pgfqpoint{4.924321in}{7.127483in}}%
\pgfpathlineto{\pgfqpoint{4.925262in}{7.122998in}}%
\pgfpathlineto{\pgfqpoint{4.926974in}{7.100733in}}%
\pgfpathlineto{\pgfqpoint{4.929626in}{7.032322in}}%
\pgfpathlineto{\pgfqpoint{4.933733in}{6.859788in}}%
\pgfpathlineto{\pgfqpoint{4.948279in}{6.191145in}}%
\pgfpathlineto{\pgfqpoint{4.951017in}{6.158363in}}%
\pgfpathlineto{\pgfqpoint{4.951616in}{6.157418in}}%
\pgfpathlineto{\pgfqpoint{4.952129in}{6.158409in}}%
\pgfpathlineto{\pgfqpoint{4.953327in}{6.167163in}}%
\pgfpathlineto{\pgfqpoint{4.955380in}{6.202604in}}%
\pgfpathlineto{\pgfqpoint{4.958546in}{6.303142in}}%
\pgfpathlineto{\pgfqpoint{4.963680in}{6.550646in}}%
\pgfpathlineto{\pgfqpoint{4.973605in}{7.029236in}}%
\pgfpathlineto{\pgfqpoint{4.977113in}{7.105971in}}%
\pgfpathlineto{\pgfqpoint{4.978996in}{7.116548in}}%
\pgfpathlineto{\pgfqpoint{4.979081in}{7.116499in}}%
\pgfpathlineto{\pgfqpoint{4.979766in}{7.114447in}}%
\pgfpathlineto{\pgfqpoint{4.981135in}{7.101547in}}%
\pgfpathlineto{\pgfqpoint{4.983445in}{7.054170in}}%
\pgfpathlineto{\pgfqpoint{4.986953in}{6.928357in}}%
\pgfpathlineto{\pgfqpoint{4.993284in}{6.603083in}}%
\pgfpathlineto{\pgfqpoint{5.000643in}{6.261843in}}%
\pgfpathlineto{\pgfqpoint{5.004236in}{6.180311in}}%
\pgfpathlineto{\pgfqpoint{5.006204in}{6.168518in}}%
\pgfpathlineto{\pgfqpoint{5.006290in}{6.168558in}}%
\pgfpathlineto{\pgfqpoint{5.006974in}{6.170538in}}%
\pgfpathlineto{\pgfqpoint{5.008343in}{6.183294in}}%
\pgfpathlineto{\pgfqpoint{5.010654in}{6.230432in}}%
\pgfpathlineto{\pgfqpoint{5.014162in}{6.355842in}}%
\pgfpathlineto{\pgfqpoint{5.020579in}{6.684341in}}%
\pgfpathlineto{\pgfqpoint{5.027766in}{7.014165in}}%
\pgfpathlineto{\pgfqpoint{5.031360in}{7.094457in}}%
\pgfpathlineto{\pgfqpoint{5.033242in}{7.105532in}}%
\pgfpathlineto{\pgfqpoint{5.033413in}{7.105433in}}%
\pgfpathlineto{\pgfqpoint{5.034183in}{7.102702in}}%
\pgfpathlineto{\pgfqpoint{5.035723in}{7.086133in}}%
\pgfpathlineto{\pgfqpoint{5.038205in}{7.029744in}}%
\pgfpathlineto{\pgfqpoint{5.041970in}{6.884861in}}%
\pgfpathlineto{\pgfqpoint{5.057970in}{6.194107in}}%
\pgfpathlineto{\pgfqpoint{5.060109in}{6.179454in}}%
\pgfpathlineto{\pgfqpoint{5.060708in}{6.180514in}}%
\pgfpathlineto{\pgfqpoint{5.061906in}{6.189402in}}%
\pgfpathlineto{\pgfqpoint{5.063959in}{6.225065in}}%
\pgfpathlineto{\pgfqpoint{5.067125in}{6.325680in}}%
\pgfpathlineto{\pgfqpoint{5.072344in}{6.575583in}}%
\pgfpathlineto{\pgfqpoint{5.081585in}{7.009858in}}%
\pgfpathlineto{\pgfqpoint{5.085093in}{7.085010in}}%
\pgfpathlineto{\pgfqpoint{5.086890in}{7.094677in}}%
\pgfpathlineto{\pgfqpoint{5.087061in}{7.094534in}}%
\pgfpathlineto{\pgfqpoint{5.087917in}{7.091045in}}%
\pgfpathlineto{\pgfqpoint{5.089542in}{7.071821in}}%
\pgfpathlineto{\pgfqpoint{5.092109in}{7.009697in}}%
\pgfpathlineto{\pgfqpoint{5.096045in}{6.852028in}}%
\pgfpathlineto{\pgfqpoint{5.110505in}{6.216131in}}%
\pgfpathlineto{\pgfqpoint{5.113072in}{6.190504in}}%
\pgfpathlineto{\pgfqpoint{5.113500in}{6.190257in}}%
\pgfpathlineto{\pgfqpoint{5.113928in}{6.191171in}}%
\pgfpathlineto{\pgfqpoint{5.115125in}{6.199877in}}%
\pgfpathlineto{\pgfqpoint{5.117179in}{6.235304in}}%
\pgfpathlineto{\pgfqpoint{5.120345in}{6.335643in}}%
\pgfpathlineto{\pgfqpoint{5.125650in}{6.588953in}}%
\pgfpathlineto{\pgfqpoint{5.134463in}{6.998864in}}%
\pgfpathlineto{\pgfqpoint{5.137971in}{7.074298in}}%
\pgfpathlineto{\pgfqpoint{5.139767in}{7.084016in}}%
\pgfpathlineto{\pgfqpoint{5.139939in}{7.083874in}}%
\pgfpathlineto{\pgfqpoint{5.140794in}{7.080371in}}%
\pgfpathlineto{\pgfqpoint{5.142420in}{7.061068in}}%
\pgfpathlineto{\pgfqpoint{5.144987in}{6.998729in}}%
\pgfpathlineto{\pgfqpoint{5.149008in}{6.836865in}}%
\pgfpathlineto{\pgfqpoint{5.162784in}{6.230911in}}%
\pgfpathlineto{\pgfqpoint{5.165436in}{6.201385in}}%
\pgfpathlineto{\pgfqpoint{5.165950in}{6.200790in}}%
\pgfpathlineto{\pgfqpoint{5.166377in}{6.201578in}}%
\pgfpathlineto{\pgfqpoint{5.167490in}{6.209069in}}%
\pgfpathlineto{\pgfqpoint{5.169458in}{6.241103in}}%
\pgfpathlineto{\pgfqpoint{5.172538in}{6.334973in}}%
\pgfpathlineto{\pgfqpoint{5.177586in}{6.570173in}}%
\pgfpathlineto{\pgfqpoint{5.186827in}{6.995825in}}%
\pgfpathlineto{\pgfqpoint{5.190249in}{7.065580in}}%
\pgfpathlineto{\pgfqpoint{5.191875in}{7.073554in}}%
\pgfpathlineto{\pgfqpoint{5.192132in}{7.073269in}}%
\pgfpathlineto{\pgfqpoint{5.193073in}{7.068626in}}%
\pgfpathlineto{\pgfqpoint{5.194784in}{7.045923in}}%
\pgfpathlineto{\pgfqpoint{5.197522in}{6.973905in}}%
\pgfpathlineto{\pgfqpoint{5.201800in}{6.792448in}}%
\pgfpathlineto{\pgfqpoint{5.213693in}{6.258906in}}%
\pgfpathlineto{\pgfqpoint{5.216688in}{6.213718in}}%
\pgfpathlineto{\pgfqpoint{5.217629in}{6.211139in}}%
\pgfpathlineto{\pgfqpoint{5.218057in}{6.211848in}}%
\pgfpathlineto{\pgfqpoint{5.219169in}{6.219172in}}%
\pgfpathlineto{\pgfqpoint{5.221137in}{6.251048in}}%
\pgfpathlineto{\pgfqpoint{5.224132in}{6.341654in}}%
\pgfpathlineto{\pgfqpoint{5.229095in}{6.570741in}}%
\pgfpathlineto{\pgfqpoint{5.238250in}{6.988546in}}%
\pgfpathlineto{\pgfqpoint{5.241587in}{7.055577in}}%
\pgfpathlineto{\pgfqpoint{5.243127in}{7.063303in}}%
\pgfpathlineto{\pgfqpoint{5.243384in}{7.063106in}}%
\pgfpathlineto{\pgfqpoint{5.244239in}{7.059384in}}%
\pgfpathlineto{\pgfqpoint{5.245865in}{7.039502in}}%
\pgfpathlineto{\pgfqpoint{5.248432in}{6.975981in}}%
\pgfpathlineto{\pgfqpoint{5.252453in}{6.812758in}}%
\pgfpathlineto{\pgfqpoint{5.265202in}{6.255645in}}%
\pgfpathlineto{\pgfqpoint{5.267940in}{6.222230in}}%
\pgfpathlineto{\pgfqpoint{5.268539in}{6.221295in}}%
\pgfpathlineto{\pgfqpoint{5.269052in}{6.222343in}}%
\pgfpathlineto{\pgfqpoint{5.270250in}{6.231396in}}%
\pgfpathlineto{\pgfqpoint{5.272304in}{6.267783in}}%
\pgfpathlineto{\pgfqpoint{5.275469in}{6.369892in}}%
\pgfpathlineto{\pgfqpoint{5.281031in}{6.635940in}}%
\pgfpathlineto{\pgfqpoint{5.288646in}{6.975411in}}%
\pgfpathlineto{\pgfqpoint{5.292069in}{7.045493in}}%
\pgfpathlineto{\pgfqpoint{5.293609in}{7.053259in}}%
\pgfpathlineto{\pgfqpoint{5.293865in}{7.053055in}}%
\pgfpathlineto{\pgfqpoint{5.294721in}{7.049282in}}%
\pgfpathlineto{\pgfqpoint{5.296347in}{7.029193in}}%
\pgfpathlineto{\pgfqpoint{5.298914in}{6.965137in}}%
\pgfpathlineto{\pgfqpoint{5.303021in}{6.797106in}}%
\pgfpathlineto{\pgfqpoint{5.315085in}{6.269646in}}%
\pgfpathlineto{\pgfqpoint{5.317908in}{6.232519in}}%
\pgfpathlineto{\pgfqpoint{5.318593in}{6.231229in}}%
\pgfpathlineto{\pgfqpoint{5.319106in}{6.232274in}}%
\pgfpathlineto{\pgfqpoint{5.320304in}{6.241382in}}%
\pgfpathlineto{\pgfqpoint{5.322358in}{6.278046in}}%
\pgfpathlineto{\pgfqpoint{5.325523in}{6.380798in}}%
\pgfpathlineto{\pgfqpoint{5.331256in}{6.655538in}}%
\pgfpathlineto{\pgfqpoint{5.338443in}{6.969796in}}%
\pgfpathlineto{\pgfqpoint{5.341780in}{7.036355in}}%
\pgfpathlineto{\pgfqpoint{5.343235in}{7.043432in}}%
\pgfpathlineto{\pgfqpoint{5.343577in}{7.043078in}}%
\pgfpathlineto{\pgfqpoint{5.344518in}{7.038142in}}%
\pgfpathlineto{\pgfqpoint{5.346230in}{7.014547in}}%
\pgfpathlineto{\pgfqpoint{5.348968in}{6.940506in}}%
\pgfpathlineto{\pgfqpoint{5.353417in}{6.748441in}}%
\pgfpathlineto{\pgfqpoint{5.363599in}{6.296484in}}%
\pgfpathlineto{\pgfqpoint{5.366679in}{6.244693in}}%
\pgfpathlineto{\pgfqpoint{5.367791in}{6.240943in}}%
\pgfpathlineto{\pgfqpoint{5.368219in}{6.241679in}}%
\pgfpathlineto{\pgfqpoint{5.369331in}{6.249232in}}%
\pgfpathlineto{\pgfqpoint{5.371299in}{6.282011in}}%
\pgfpathlineto{\pgfqpoint{5.374380in}{6.377937in}}%
\pgfpathlineto{\pgfqpoint{5.379770in}{6.630925in}}%
\pgfpathlineto{\pgfqpoint{5.387300in}{6.961787in}}%
\pgfpathlineto{\pgfqpoint{5.390637in}{7.027414in}}%
\pgfpathlineto{\pgfqpoint{5.392006in}{7.033824in}}%
\pgfpathlineto{\pgfqpoint{5.392348in}{7.033483in}}%
\pgfpathlineto{\pgfqpoint{5.393289in}{7.028534in}}%
\pgfpathlineto{\pgfqpoint{5.395000in}{7.004761in}}%
\pgfpathlineto{\pgfqpoint{5.397738in}{6.930144in}}%
\pgfpathlineto{\pgfqpoint{5.402187in}{6.737388in}}%
\pgfpathlineto{\pgfqpoint{5.411942in}{6.306682in}}%
\pgfpathlineto{\pgfqpoint{5.415022in}{6.254261in}}%
\pgfpathlineto{\pgfqpoint{5.416134in}{6.250438in}}%
\pgfpathlineto{\pgfqpoint{5.416562in}{6.251171in}}%
\pgfpathlineto{\pgfqpoint{5.417674in}{6.258782in}}%
\pgfpathlineto{\pgfqpoint{5.419642in}{6.291868in}}%
\pgfpathlineto{\pgfqpoint{5.422722in}{6.388568in}}%
\pgfpathlineto{\pgfqpoint{5.428198in}{6.646157in}}%
\pgfpathlineto{\pgfqpoint{5.435386in}{6.956633in}}%
\pgfpathlineto{\pgfqpoint{5.438637in}{7.018722in}}%
\pgfpathlineto{\pgfqpoint{5.439920in}{7.024439in}}%
\pgfpathlineto{\pgfqpoint{5.440348in}{7.023885in}}%
\pgfpathlineto{\pgfqpoint{5.441375in}{7.017544in}}%
\pgfpathlineto{\pgfqpoint{5.443257in}{6.988005in}}%
\pgfpathlineto{\pgfqpoint{5.446167in}{6.901143in}}%
\pgfpathlineto{\pgfqpoint{5.451215in}{6.670233in}}%
\pgfpathlineto{\pgfqpoint{5.459086in}{6.325968in}}%
\pgfpathlineto{\pgfqpoint{5.462338in}{6.264887in}}%
\pgfpathlineto{\pgfqpoint{5.463621in}{6.259713in}}%
\pgfpathlineto{\pgfqpoint{5.463964in}{6.260216in}}%
\pgfpathlineto{\pgfqpoint{5.464990in}{6.266467in}}%
\pgfpathlineto{\pgfqpoint{5.466787in}{6.294126in}}%
\pgfpathlineto{\pgfqpoint{5.469696in}{6.379759in}}%
\pgfpathlineto{\pgfqpoint{5.474659in}{6.605267in}}%
\pgfpathlineto{\pgfqpoint{5.482616in}{6.951802in}}%
\pgfpathlineto{\pgfqpoint{5.485782in}{7.010279in}}%
\pgfpathlineto{\pgfqpoint{5.486980in}{7.015277in}}%
\pgfpathlineto{\pgfqpoint{5.487408in}{7.014695in}}%
\pgfpathlineto{\pgfqpoint{5.488434in}{7.008225in}}%
\pgfpathlineto{\pgfqpoint{5.490317in}{6.978248in}}%
\pgfpathlineto{\pgfqpoint{5.493311in}{6.887179in}}%
\pgfpathlineto{\pgfqpoint{5.498531in}{6.645843in}}%
\pgfpathlineto{\pgfqpoint{5.505804in}{6.333063in}}%
\pgfpathlineto{\pgfqpoint{5.508969in}{6.273923in}}%
\pgfpathlineto{\pgfqpoint{5.510253in}{6.268769in}}%
\pgfpathlineto{\pgfqpoint{5.510595in}{6.269301in}}%
\pgfpathlineto{\pgfqpoint{5.511622in}{6.275701in}}%
\pgfpathlineto{\pgfqpoint{5.513504in}{6.305665in}}%
\pgfpathlineto{\pgfqpoint{5.516413in}{6.393682in}}%
\pgfpathlineto{\pgfqpoint{5.521633in}{6.634010in}}%
\pgfpathlineto{\pgfqpoint{5.528905in}{6.944825in}}%
\pgfpathlineto{\pgfqpoint{5.532071in}{7.002046in}}%
\pgfpathlineto{\pgfqpoint{5.533184in}{7.006336in}}%
\pgfpathlineto{\pgfqpoint{5.533611in}{7.005713in}}%
\pgfpathlineto{\pgfqpoint{5.534638in}{6.999082in}}%
\pgfpathlineto{\pgfqpoint{5.536521in}{6.968598in}}%
\pgfpathlineto{\pgfqpoint{5.539515in}{6.876432in}}%
\pgfpathlineto{\pgfqpoint{5.544906in}{6.625878in}}%
\pgfpathlineto{\pgfqpoint{5.551751in}{6.337550in}}%
\pgfpathlineto{\pgfqpoint{5.554831in}{6.282037in}}%
\pgfpathlineto{\pgfqpoint{5.555943in}{6.277592in}}%
\pgfpathlineto{\pgfqpoint{5.556371in}{6.278170in}}%
\pgfpathlineto{\pgfqpoint{5.557398in}{6.284729in}}%
\pgfpathlineto{\pgfqpoint{5.559280in}{6.315200in}}%
\pgfpathlineto{\pgfqpoint{5.562275in}{6.407569in}}%
\pgfpathlineto{\pgfqpoint{5.567751in}{6.662409in}}%
\pgfpathlineto{\pgfqpoint{5.574339in}{6.937962in}}%
\pgfpathlineto{\pgfqpoint{5.577419in}{6.993328in}}%
\pgfpathlineto{\pgfqpoint{5.578532in}{6.997613in}}%
\pgfpathlineto{\pgfqpoint{5.578959in}{6.996958in}}%
\pgfpathlineto{\pgfqpoint{5.579986in}{6.990181in}}%
\pgfpathlineto{\pgfqpoint{5.581869in}{6.959214in}}%
\pgfpathlineto{\pgfqpoint{5.584863in}{6.865974in}}%
\pgfpathlineto{\pgfqpoint{5.590425in}{6.606440in}}%
\pgfpathlineto{\pgfqpoint{5.596842in}{6.342133in}}%
\pgfpathlineto{\pgfqpoint{5.599922in}{6.289435in}}%
\pgfpathlineto{\pgfqpoint{5.600863in}{6.286206in}}%
\pgfpathlineto{\pgfqpoint{5.601377in}{6.287074in}}%
\pgfpathlineto{\pgfqpoint{5.602489in}{6.295291in}}%
\pgfpathlineto{\pgfqpoint{5.604457in}{6.330360in}}%
\pgfpathlineto{\pgfqpoint{5.607623in}{6.434747in}}%
\pgfpathlineto{\pgfqpoint{5.614382in}{6.755954in}}%
\pgfpathlineto{\pgfqpoint{5.619516in}{6.946334in}}%
\pgfpathlineto{\pgfqpoint{5.622340in}{6.987443in}}%
\pgfpathlineto{\pgfqpoint{5.623024in}{6.989106in}}%
\pgfpathlineto{\pgfqpoint{5.623538in}{6.988173in}}%
\pgfpathlineto{\pgfqpoint{5.624650in}{6.979773in}}%
\pgfpathlineto{\pgfqpoint{5.626618in}{6.944270in}}%
\pgfpathlineto{\pgfqpoint{5.629784in}{6.839101in}}%
\pgfpathlineto{\pgfqpoint{5.637056in}{6.494786in}}%
\pgfpathlineto{\pgfqpoint{5.641848in}{6.328660in}}%
\pgfpathlineto{\pgfqpoint{5.644500in}{6.295364in}}%
\pgfpathlineto{\pgfqpoint{5.645014in}{6.294614in}}%
\pgfpathlineto{\pgfqpoint{5.645527in}{6.295747in}}%
\pgfpathlineto{\pgfqpoint{5.646725in}{6.305660in}}%
\pgfpathlineto{\pgfqpoint{5.648778in}{6.345381in}}%
\pgfpathlineto{\pgfqpoint{5.652030in}{6.458300in}}%
\pgfpathlineto{\pgfqpoint{5.665292in}{6.973782in}}%
\pgfpathlineto{\pgfqpoint{5.666747in}{6.980803in}}%
\pgfpathlineto{\pgfqpoint{5.667003in}{6.980466in}}%
\pgfpathlineto{\pgfqpoint{5.667944in}{6.975190in}}%
\pgfpathlineto{\pgfqpoint{5.669656in}{6.949704in}}%
\pgfpathlineto{\pgfqpoint{5.672479in}{6.867240in}}%
\pgfpathlineto{\pgfqpoint{5.677527in}{6.638794in}}%
\pgfpathlineto{\pgfqpoint{5.684372in}{6.354995in}}%
\pgfpathlineto{\pgfqpoint{5.687367in}{6.305423in}}%
\pgfpathlineto{\pgfqpoint{5.688223in}{6.302792in}}%
\pgfpathlineto{\pgfqpoint{5.688736in}{6.303757in}}%
\pgfpathlineto{\pgfqpoint{5.689848in}{6.312355in}}%
\pgfpathlineto{\pgfqpoint{5.691816in}{6.348580in}}%
\pgfpathlineto{\pgfqpoint{5.694982in}{6.455340in}}%
\pgfpathlineto{\pgfqpoint{5.708415in}{6.968084in}}%
\pgfpathlineto{\pgfqpoint{5.709528in}{6.972729in}}%
\pgfpathlineto{\pgfqpoint{5.709956in}{6.972117in}}%
\pgfpathlineto{\pgfqpoint{5.710982in}{6.965227in}}%
\pgfpathlineto{\pgfqpoint{5.712865in}{6.933341in}}%
\pgfpathlineto{\pgfqpoint{5.715859in}{6.837576in}}%
\pgfpathlineto{\pgfqpoint{5.722191in}{6.541335in}}%
\pgfpathlineto{\pgfqpoint{5.727325in}{6.351042in}}%
\pgfpathlineto{\pgfqpoint{5.730063in}{6.312132in}}%
\pgfpathlineto{\pgfqpoint{5.730662in}{6.310773in}}%
\pgfpathlineto{\pgfqpoint{5.731175in}{6.311704in}}%
\pgfpathlineto{\pgfqpoint{5.732287in}{6.320317in}}%
\pgfpathlineto{\pgfqpoint{5.734255in}{6.356834in}}%
\pgfpathlineto{\pgfqpoint{5.737421in}{6.464404in}}%
\pgfpathlineto{\pgfqpoint{5.750341in}{6.959012in}}%
\pgfpathlineto{\pgfqpoint{5.751625in}{6.964849in}}%
\pgfpathlineto{\pgfqpoint{5.751967in}{6.964354in}}%
\pgfpathlineto{\pgfqpoint{5.752908in}{6.958544in}}%
\pgfpathlineto{\pgfqpoint{5.754705in}{6.929846in}}%
\pgfpathlineto{\pgfqpoint{5.757614in}{6.840012in}}%
\pgfpathlineto{\pgfqpoint{5.763347in}{6.575072in}}%
\pgfpathlineto{\pgfqpoint{5.768908in}{6.361219in}}%
\pgfpathlineto{\pgfqpoint{5.771732in}{6.319894in}}%
\pgfpathlineto{\pgfqpoint{5.772331in}{6.318550in}}%
\pgfpathlineto{\pgfqpoint{5.772844in}{6.319522in}}%
\pgfpathlineto{\pgfqpoint{5.773956in}{6.328313in}}%
\pgfpathlineto{\pgfqpoint{5.775924in}{6.365388in}}%
\pgfpathlineto{\pgfqpoint{5.779090in}{6.474081in}}%
\pgfpathlineto{\pgfqpoint{5.791497in}{6.950103in}}%
\pgfpathlineto{\pgfqpoint{5.792866in}{6.957174in}}%
\pgfpathlineto{\pgfqpoint{5.793208in}{6.956754in}}%
\pgfpathlineto{\pgfqpoint{5.794149in}{6.951090in}}%
\pgfpathlineto{\pgfqpoint{5.795860in}{6.924306in}}%
\pgfpathlineto{\pgfqpoint{5.798684in}{6.838788in}}%
\pgfpathlineto{\pgfqpoint{5.804245in}{6.584089in}}%
\pgfpathlineto{\pgfqpoint{5.809893in}{6.367027in}}%
\pgfpathlineto{\pgfqpoint{5.812631in}{6.327449in}}%
\pgfpathlineto{\pgfqpoint{5.813229in}{6.326129in}}%
\pgfpathlineto{\pgfqpoint{5.813743in}{6.327151in}}%
\pgfpathlineto{\pgfqpoint{5.814855in}{6.336136in}}%
\pgfpathlineto{\pgfqpoint{5.816823in}{6.373797in}}%
\pgfpathlineto{\pgfqpoint{5.820074in}{6.487242in}}%
\pgfpathlineto{\pgfqpoint{5.831797in}{6.940180in}}%
\pgfpathlineto{\pgfqpoint{5.833422in}{6.949687in}}%
\pgfpathlineto{\pgfqpoint{5.833679in}{6.949361in}}%
\pgfpathlineto{\pgfqpoint{5.834620in}{6.943901in}}%
\pgfpathlineto{\pgfqpoint{5.836331in}{6.917263in}}%
\pgfpathlineto{\pgfqpoint{5.839155in}{6.831533in}}%
\pgfpathlineto{\pgfqpoint{5.844802in}{6.573275in}}%
\pgfpathlineto{\pgfqpoint{5.850192in}{6.371036in}}%
\pgfpathlineto{\pgfqpoint{5.852845in}{6.334522in}}%
\pgfpathlineto{\pgfqpoint{5.853358in}{6.333514in}}%
\pgfpathlineto{\pgfqpoint{5.853872in}{6.334520in}}%
\pgfpathlineto{\pgfqpoint{5.854984in}{6.343560in}}%
\pgfpathlineto{\pgfqpoint{5.856952in}{6.381574in}}%
\pgfpathlineto{\pgfqpoint{5.860203in}{6.495857in}}%
\pgfpathlineto{\pgfqpoint{5.871412in}{6.930850in}}%
\pgfpathlineto{\pgfqpoint{5.873123in}{6.942400in}}%
\pgfpathlineto{\pgfqpoint{5.873380in}{6.942200in}}%
\pgfpathlineto{\pgfqpoint{5.874236in}{6.937882in}}%
\pgfpathlineto{\pgfqpoint{5.875861in}{6.914557in}}%
\pgfpathlineto{\pgfqpoint{5.878514in}{6.838278in}}%
\pgfpathlineto{\pgfqpoint{5.883562in}{6.612818in}}%
\pgfpathlineto{\pgfqpoint{5.889466in}{6.381784in}}%
\pgfpathlineto{\pgfqpoint{5.892204in}{6.341913in}}%
\pgfpathlineto{\pgfqpoint{5.892803in}{6.340717in}}%
\pgfpathlineto{\pgfqpoint{5.893316in}{6.341902in}}%
\pgfpathlineto{\pgfqpoint{5.894514in}{6.352531in}}%
\pgfpathlineto{\pgfqpoint{5.896567in}{6.395097in}}%
\pgfpathlineto{\pgfqpoint{5.899990in}{6.521290in}}%
\pgfpathlineto{\pgfqpoint{5.910001in}{6.916760in}}%
\pgfpathlineto{\pgfqpoint{5.912225in}{6.935294in}}%
\pgfpathlineto{\pgfqpoint{5.912824in}{6.933748in}}%
\pgfpathlineto{\pgfqpoint{5.914108in}{6.921127in}}%
\pgfpathlineto{\pgfqpoint{5.916247in}{6.873527in}}%
\pgfpathlineto{\pgfqpoint{5.919840in}{6.735192in}}%
\pgfpathlineto{\pgfqpoint{5.928910in}{6.372650in}}%
\pgfpathlineto{\pgfqpoint{5.931306in}{6.347804in}}%
\pgfpathlineto{\pgfqpoint{5.931648in}{6.347878in}}%
\pgfpathlineto{\pgfqpoint{5.931905in}{6.348536in}}%
\pgfpathlineto{\pgfqpoint{5.933017in}{6.357309in}}%
\pgfpathlineto{\pgfqpoint{5.934985in}{6.395402in}}%
\pgfpathlineto{\pgfqpoint{5.938236in}{6.510533in}}%
\pgfpathlineto{\pgfqpoint{5.948760in}{6.916080in}}%
\pgfpathlineto{\pgfqpoint{5.950557in}{6.928373in}}%
\pgfpathlineto{\pgfqpoint{5.950728in}{6.928223in}}%
\pgfpathlineto{\pgfqpoint{5.951498in}{6.924696in}}%
\pgfpathlineto{\pgfqpoint{5.953039in}{6.903927in}}%
\pgfpathlineto{\pgfqpoint{5.955605in}{6.832590in}}%
\pgfpathlineto{\pgfqpoint{5.960397in}{6.622367in}}%
\pgfpathlineto{\pgfqpoint{5.966301in}{6.392879in}}%
\pgfpathlineto{\pgfqpoint{5.968953in}{6.355511in}}%
\pgfpathlineto{\pgfqpoint{5.969467in}{6.354560in}}%
\pgfpathlineto{\pgfqpoint{5.969980in}{6.355701in}}%
\pgfpathlineto{\pgfqpoint{5.971178in}{6.366422in}}%
\pgfpathlineto{\pgfqpoint{5.973231in}{6.409673in}}%
\pgfpathlineto{\pgfqpoint{5.976654in}{6.537286in}}%
\pgfpathlineto{\pgfqpoint{5.985895in}{6.900911in}}%
\pgfpathlineto{\pgfqpoint{5.988119in}{6.921610in}}%
\pgfpathlineto{\pgfqpoint{5.988633in}{6.920832in}}%
\pgfpathlineto{\pgfqpoint{5.989659in}{6.913002in}}%
\pgfpathlineto{\pgfqpoint{5.991542in}{6.877822in}}%
\pgfpathlineto{\pgfqpoint{5.994707in}{6.768109in}}%
\pgfpathlineto{\pgfqpoint{6.005146in}{6.371655in}}%
\pgfpathlineto{\pgfqpoint{6.006772in}{6.361219in}}%
\pgfpathlineto{\pgfqpoint{6.007028in}{6.361503in}}%
\pgfpathlineto{\pgfqpoint{6.007884in}{6.366258in}}%
\pgfpathlineto{\pgfqpoint{6.009510in}{6.391005in}}%
\pgfpathlineto{\pgfqpoint{6.012248in}{6.473649in}}%
\pgfpathlineto{\pgfqpoint{6.018751in}{6.764381in}}%
\pgfpathlineto{\pgfqpoint{6.022943in}{6.895153in}}%
\pgfpathlineto{\pgfqpoint{6.025168in}{6.915052in}}%
\pgfpathlineto{\pgfqpoint{6.025681in}{6.914019in}}%
\pgfpathlineto{\pgfqpoint{6.026793in}{6.904537in}}%
\pgfpathlineto{\pgfqpoint{6.028761in}{6.864681in}}%
\pgfpathlineto{\pgfqpoint{6.032098in}{6.742996in}}%
\pgfpathlineto{\pgfqpoint{6.041339in}{6.385244in}}%
\pgfpathlineto{\pgfqpoint{6.043393in}{6.367716in}}%
\pgfpathlineto{\pgfqpoint{6.043478in}{6.367722in}}%
\pgfpathlineto{\pgfqpoint{6.044077in}{6.369428in}}%
\pgfpathlineto{\pgfqpoint{6.045360in}{6.382771in}}%
\pgfpathlineto{\pgfqpoint{6.047585in}{6.435042in}}%
\pgfpathlineto{\pgfqpoint{6.051435in}{6.589361in}}%
\pgfpathlineto{\pgfqpoint{6.058794in}{6.878832in}}%
\pgfpathlineto{\pgfqpoint{6.061275in}{6.908397in}}%
\pgfpathlineto{\pgfqpoint{6.061617in}{6.908603in}}%
\pgfpathlineto{\pgfqpoint{6.061960in}{6.907850in}}%
\pgfpathlineto{\pgfqpoint{6.062986in}{6.899892in}}%
\pgfpathlineto{\pgfqpoint{6.064869in}{6.864061in}}%
\pgfpathlineto{\pgfqpoint{6.068034in}{6.752972in}}%
\pgfpathlineto{\pgfqpoint{6.077703in}{6.386569in}}%
\pgfpathlineto{\pgfqpoint{6.079500in}{6.374049in}}%
\pgfpathlineto{\pgfqpoint{6.079671in}{6.374232in}}%
\pgfpathlineto{\pgfqpoint{6.080441in}{6.378024in}}%
\pgfpathlineto{\pgfqpoint{6.081981in}{6.399841in}}%
\pgfpathlineto{\pgfqpoint{6.084548in}{6.473735in}}%
\pgfpathlineto{\pgfqpoint{6.090024in}{6.715068in}}%
\pgfpathlineto{\pgfqpoint{6.094730in}{6.876827in}}%
\pgfpathlineto{\pgfqpoint{6.097126in}{6.902341in}}%
\pgfpathlineto{\pgfqpoint{6.097468in}{6.902173in}}%
\pgfpathlineto{\pgfqpoint{6.097725in}{6.901412in}}%
\pgfpathlineto{\pgfqpoint{6.098837in}{6.891886in}}%
\pgfpathlineto{\pgfqpoint{6.100805in}{6.851507in}}%
\pgfpathlineto{\pgfqpoint{6.104142in}{6.728774in}}%
\pgfpathlineto{\pgfqpoint{6.112698in}{6.399528in}}%
\pgfpathlineto{\pgfqpoint{6.114837in}{6.380215in}}%
\pgfpathlineto{\pgfqpoint{6.114923in}{6.380224in}}%
\pgfpathlineto{\pgfqpoint{6.115522in}{6.381988in}}%
\pgfpathlineto{\pgfqpoint{6.116805in}{6.395660in}}%
\pgfpathlineto{\pgfqpoint{6.119030in}{6.448951in}}%
\pgfpathlineto{\pgfqpoint{6.122966in}{6.608229in}}%
\pgfpathlineto{\pgfqpoint{6.129639in}{6.866252in}}%
\pgfpathlineto{\pgfqpoint{6.132121in}{6.896083in}}%
\pgfpathlineto{\pgfqpoint{6.132463in}{6.896244in}}%
\pgfpathlineto{\pgfqpoint{6.132805in}{6.895427in}}%
\pgfpathlineto{\pgfqpoint{6.133918in}{6.886076in}}%
\pgfpathlineto{\pgfqpoint{6.135886in}{6.845766in}}%
\pgfpathlineto{\pgfqpoint{6.139222in}{6.722937in}}%
\pgfpathlineto{\pgfqpoint{6.147522in}{6.405562in}}%
\pgfpathlineto{\pgfqpoint{6.149661in}{6.386230in}}%
\pgfpathlineto{\pgfqpoint{6.149747in}{6.386246in}}%
\pgfpathlineto{\pgfqpoint{6.150346in}{6.388078in}}%
\pgfpathlineto{\pgfqpoint{6.151629in}{6.401991in}}%
\pgfpathlineto{\pgfqpoint{6.153854in}{6.455888in}}%
\pgfpathlineto{\pgfqpoint{6.157875in}{6.619559in}}%
\pgfpathlineto{\pgfqpoint{6.164207in}{6.861012in}}%
\pgfpathlineto{\pgfqpoint{6.166688in}{6.890218in}}%
\pgfpathlineto{\pgfqpoint{6.166945in}{6.890334in}}%
\pgfpathlineto{\pgfqpoint{6.167287in}{6.889624in}}%
\pgfpathlineto{\pgfqpoint{6.168314in}{6.881611in}}%
\pgfpathlineto{\pgfqpoint{6.170196in}{6.845093in}}%
\pgfpathlineto{\pgfqpoint{6.173447in}{6.728900in}}%
\pgfpathlineto{\pgfqpoint{6.182004in}{6.407441in}}%
\pgfpathlineto{\pgfqpoint{6.183886in}{6.392100in}}%
\pgfpathlineto{\pgfqpoint{6.184057in}{6.392181in}}%
\pgfpathlineto{\pgfqpoint{6.184742in}{6.394985in}}%
\pgfpathlineto{\pgfqpoint{6.186196in}{6.413832in}}%
\pgfpathlineto{\pgfqpoint{6.188678in}{6.482210in}}%
\pgfpathlineto{\pgfqpoint{6.193897in}{6.707083in}}%
\pgfpathlineto{\pgfqpoint{6.198517in}{6.862291in}}%
\pgfpathlineto{\pgfqpoint{6.200742in}{6.884536in}}%
\pgfpathlineto{\pgfqpoint{6.201170in}{6.884044in}}%
\pgfpathlineto{\pgfqpoint{6.201255in}{6.883759in}}%
\pgfpathlineto{\pgfqpoint{6.202282in}{6.875520in}}%
\pgfpathlineto{\pgfqpoint{6.204164in}{6.838443in}}%
\pgfpathlineto{\pgfqpoint{6.207416in}{6.721473in}}%
\pgfpathlineto{\pgfqpoint{6.215630in}{6.414052in}}%
\pgfpathlineto{\pgfqpoint{6.217598in}{6.397820in}}%
\pgfpathlineto{\pgfqpoint{6.217683in}{6.397859in}}%
\pgfpathlineto{\pgfqpoint{6.218282in}{6.399881in}}%
\pgfpathlineto{\pgfqpoint{6.219566in}{6.414378in}}%
\pgfpathlineto{\pgfqpoint{6.221790in}{6.469584in}}%
\pgfpathlineto{\pgfqpoint{6.225983in}{6.641726in}}%
\pgfpathlineto{\pgfqpoint{6.231715in}{6.852462in}}%
\pgfpathlineto{\pgfqpoint{6.234111in}{6.878860in}}%
\pgfpathlineto{\pgfqpoint{6.234453in}{6.878672in}}%
\pgfpathlineto{\pgfqpoint{6.234625in}{6.878200in}}%
\pgfpathlineto{\pgfqpoint{6.235651in}{6.870132in}}%
\pgfpathlineto{\pgfqpoint{6.237534in}{6.833167in}}%
\pgfpathlineto{\pgfqpoint{6.240785in}{6.716228in}}%
\pgfpathlineto{\pgfqpoint{6.248742in}{6.420019in}}%
\pgfpathlineto{\pgfqpoint{6.250710in}{6.403403in}}%
\pgfpathlineto{\pgfqpoint{6.250796in}{6.403431in}}%
\pgfpathlineto{\pgfqpoint{6.251395in}{6.405397in}}%
\pgfpathlineto{\pgfqpoint{6.252678in}{6.419863in}}%
\pgfpathlineto{\pgfqpoint{6.254903in}{6.475218in}}%
\pgfpathlineto{\pgfqpoint{6.259181in}{6.650950in}}%
\pgfpathlineto{\pgfqpoint{6.264657in}{6.848384in}}%
\pgfpathlineto{\pgfqpoint{6.266967in}{6.873340in}}%
\pgfpathlineto{\pgfqpoint{6.267309in}{6.873165in}}%
\pgfpathlineto{\pgfqpoint{6.267481in}{6.872697in}}%
\pgfpathlineto{\pgfqpoint{6.268507in}{6.864604in}}%
\pgfpathlineto{\pgfqpoint{6.270390in}{6.827437in}}%
\pgfpathlineto{\pgfqpoint{6.273641in}{6.710214in}}%
\pgfpathlineto{\pgfqpoint{6.281342in}{6.425600in}}%
\pgfpathlineto{\pgfqpoint{6.283310in}{6.408851in}}%
\pgfpathlineto{\pgfqpoint{6.283395in}{6.408879in}}%
\pgfpathlineto{\pgfqpoint{6.283994in}{6.410863in}}%
\pgfpathlineto{\pgfqpoint{6.285278in}{6.425448in}}%
\pgfpathlineto{\pgfqpoint{6.287502in}{6.481161in}}%
\pgfpathlineto{\pgfqpoint{6.291866in}{6.660492in}}%
\pgfpathlineto{\pgfqpoint{6.297085in}{6.844554in}}%
\pgfpathlineto{\pgfqpoint{6.299395in}{6.868010in}}%
\pgfpathlineto{\pgfqpoint{6.299909in}{6.866972in}}%
\pgfpathlineto{\pgfqpoint{6.301021in}{6.856898in}}%
\pgfpathlineto{\pgfqpoint{6.302989in}{6.814528in}}%
\pgfpathlineto{\pgfqpoint{6.306583in}{6.678097in}}%
\pgfpathlineto{\pgfqpoint{6.313171in}{6.435675in}}%
\pgfpathlineto{\pgfqpoint{6.315396in}{6.414167in}}%
\pgfpathlineto{\pgfqpoint{6.315909in}{6.415298in}}%
\pgfpathlineto{\pgfqpoint{6.317021in}{6.425597in}}%
\pgfpathlineto{\pgfqpoint{6.319075in}{6.470898in}}%
\pgfpathlineto{\pgfqpoint{6.322754in}{6.612727in}}%
\pgfpathlineto{\pgfqpoint{6.329000in}{6.840648in}}%
\pgfpathlineto{\pgfqpoint{6.331225in}{6.862756in}}%
\pgfpathlineto{\pgfqpoint{6.331738in}{6.861746in}}%
\pgfpathlineto{\pgfqpoint{6.332850in}{6.851671in}}%
\pgfpathlineto{\pgfqpoint{6.334818in}{6.809153in}}%
\pgfpathlineto{\pgfqpoint{6.338412in}{6.672725in}}%
\pgfpathlineto{\pgfqpoint{6.344829in}{6.439754in}}%
\pgfpathlineto{\pgfqpoint{6.346968in}{6.419358in}}%
\pgfpathlineto{\pgfqpoint{6.347054in}{6.419372in}}%
\pgfpathlineto{\pgfqpoint{6.347653in}{6.421282in}}%
\pgfpathlineto{\pgfqpoint{6.348936in}{6.435871in}}%
\pgfpathlineto{\pgfqpoint{6.351161in}{6.491901in}}%
\pgfpathlineto{\pgfqpoint{6.355695in}{6.677600in}}%
\pgfpathlineto{\pgfqpoint{6.360487in}{6.837974in}}%
\pgfpathlineto{\pgfqpoint{6.362626in}{6.857634in}}%
\pgfpathlineto{\pgfqpoint{6.363139in}{6.856375in}}%
\pgfpathlineto{\pgfqpoint{6.364252in}{6.845718in}}%
\pgfpathlineto{\pgfqpoint{6.366305in}{6.799625in}}%
\pgfpathlineto{\pgfqpoint{6.370156in}{6.649946in}}%
\pgfpathlineto{\pgfqpoint{6.375888in}{6.445914in}}%
\pgfpathlineto{\pgfqpoint{6.378113in}{6.424420in}}%
\pgfpathlineto{\pgfqpoint{6.378626in}{6.425643in}}%
\pgfpathlineto{\pgfqpoint{6.379739in}{6.436250in}}%
\pgfpathlineto{\pgfqpoint{6.381792in}{6.482327in}}%
\pgfpathlineto{\pgfqpoint{6.385642in}{6.631808in}}%
\pgfpathlineto{\pgfqpoint{6.391289in}{6.831490in}}%
\pgfpathlineto{\pgfqpoint{6.393514in}{6.852626in}}%
\pgfpathlineto{\pgfqpoint{6.394027in}{6.851298in}}%
\pgfpathlineto{\pgfqpoint{6.395225in}{6.839166in}}%
\pgfpathlineto{\pgfqpoint{6.397364in}{6.788659in}}%
\pgfpathlineto{\pgfqpoint{6.401557in}{6.621883in}}%
\pgfpathlineto{\pgfqpoint{6.406605in}{6.449604in}}%
\pgfpathlineto{\pgfqpoint{6.408744in}{6.429365in}}%
\pgfpathlineto{\pgfqpoint{6.409258in}{6.430540in}}%
\pgfpathlineto{\pgfqpoint{6.410370in}{6.441098in}}%
\pgfpathlineto{\pgfqpoint{6.412423in}{6.487217in}}%
\pgfpathlineto{\pgfqpoint{6.416274in}{6.636324in}}%
\pgfpathlineto{\pgfqpoint{6.421750in}{6.827321in}}%
\pgfpathlineto{\pgfqpoint{6.423889in}{6.847744in}}%
\pgfpathlineto{\pgfqpoint{6.424402in}{6.846596in}}%
\pgfpathlineto{\pgfqpoint{6.425514in}{6.836071in}}%
\pgfpathlineto{\pgfqpoint{6.427568in}{6.789949in}}%
\pgfpathlineto{\pgfqpoint{6.431418in}{6.641067in}}%
\pgfpathlineto{\pgfqpoint{6.436809in}{6.454318in}}%
\pgfpathlineto{\pgfqpoint{6.438948in}{6.434194in}}%
\pgfpathlineto{\pgfqpoint{6.439461in}{6.435434in}}%
\pgfpathlineto{\pgfqpoint{6.440573in}{6.446177in}}%
\pgfpathlineto{\pgfqpoint{6.442627in}{6.492700in}}%
\pgfpathlineto{\pgfqpoint{6.446563in}{6.645352in}}%
\pgfpathlineto{\pgfqpoint{6.451782in}{6.823592in}}%
\pgfpathlineto{\pgfqpoint{6.453836in}{6.842969in}}%
\pgfpathlineto{\pgfqpoint{6.453921in}{6.842960in}}%
\pgfpathlineto{\pgfqpoint{6.454520in}{6.841043in}}%
\pgfpathlineto{\pgfqpoint{6.455804in}{6.826214in}}%
\pgfpathlineto{\pgfqpoint{6.458114in}{6.766597in}}%
\pgfpathlineto{\pgfqpoint{6.468723in}{6.438914in}}%
\pgfpathlineto{\pgfqpoint{6.468895in}{6.439101in}}%
\pgfpathlineto{\pgfqpoint{6.469665in}{6.443219in}}%
\pgfpathlineto{\pgfqpoint{6.471205in}{6.466988in}}%
\pgfpathlineto{\pgfqpoint{6.473943in}{6.551776in}}%
\pgfpathlineto{\pgfqpoint{6.482328in}{6.833051in}}%
\pgfpathlineto{\pgfqpoint{6.483440in}{6.838307in}}%
\pgfpathlineto{\pgfqpoint{6.483868in}{6.837335in}}%
\pgfpathlineto{\pgfqpoint{6.484980in}{6.827108in}}%
\pgfpathlineto{\pgfqpoint{6.486948in}{6.783800in}}%
\pgfpathlineto{\pgfqpoint{6.490799in}{6.637301in}}%
\pgfpathlineto{\pgfqpoint{6.496018in}{6.461473in}}%
\pgfpathlineto{\pgfqpoint{6.497986in}{6.443520in}}%
\pgfpathlineto{\pgfqpoint{6.498157in}{6.443607in}}%
\pgfpathlineto{\pgfqpoint{6.498841in}{6.446622in}}%
\pgfpathlineto{\pgfqpoint{6.500296in}{6.466795in}}%
\pgfpathlineto{\pgfqpoint{6.502863in}{6.541391in}}%
\pgfpathlineto{\pgfqpoint{6.511847in}{6.831667in}}%
\pgfpathlineto{\pgfqpoint{6.512531in}{6.833759in}}%
\pgfpathlineto{\pgfqpoint{6.513045in}{6.832516in}}%
\pgfpathlineto{\pgfqpoint{6.514157in}{6.821663in}}%
\pgfpathlineto{\pgfqpoint{6.516211in}{6.774760in}}%
\pgfpathlineto{\pgfqpoint{6.520318in}{6.616212in}}%
\pgfpathlineto{\pgfqpoint{6.525024in}{6.464456in}}%
\pgfpathlineto{\pgfqpoint{6.526906in}{6.448018in}}%
\pgfpathlineto{\pgfqpoint{6.527077in}{6.448116in}}%
\pgfpathlineto{\pgfqpoint{6.527762in}{6.451186in}}%
\pgfpathlineto{\pgfqpoint{6.529216in}{6.471525in}}%
\pgfpathlineto{\pgfqpoint{6.531783in}{6.546346in}}%
\pgfpathlineto{\pgfqpoint{6.540510in}{6.826898in}}%
\pgfpathlineto{\pgfqpoint{6.541280in}{6.829302in}}%
\pgfpathlineto{\pgfqpoint{6.541708in}{6.828282in}}%
\pgfpathlineto{\pgfqpoint{6.542821in}{6.817863in}}%
\pgfpathlineto{\pgfqpoint{6.544874in}{6.771585in}}%
\pgfpathlineto{\pgfqpoint{6.548981in}{6.614298in}}%
\pgfpathlineto{\pgfqpoint{6.553601in}{6.467830in}}%
\pgfpathlineto{\pgfqpoint{6.555484in}{6.452421in}}%
\pgfpathlineto{\pgfqpoint{6.555655in}{6.452621in}}%
\pgfpathlineto{\pgfqpoint{6.556425in}{6.456849in}}%
\pgfpathlineto{\pgfqpoint{6.557965in}{6.481004in}}%
\pgfpathlineto{\pgfqpoint{6.560703in}{6.566149in}}%
\pgfpathlineto{\pgfqpoint{6.568404in}{6.818902in}}%
\pgfpathlineto{\pgfqpoint{6.569602in}{6.824953in}}%
\pgfpathlineto{\pgfqpoint{6.569944in}{6.824256in}}%
\pgfpathlineto{\pgfqpoint{6.570971in}{6.815743in}}%
\pgfpathlineto{\pgfqpoint{6.572853in}{6.776860in}}%
\pgfpathlineto{\pgfqpoint{6.576447in}{6.645584in}}%
\pgfpathlineto{\pgfqpoint{6.581751in}{6.471752in}}%
\pgfpathlineto{\pgfqpoint{6.583548in}{6.456711in}}%
\pgfpathlineto{\pgfqpoint{6.583719in}{6.456819in}}%
\pgfpathlineto{\pgfqpoint{6.584404in}{6.459951in}}%
\pgfpathlineto{\pgfqpoint{6.585858in}{6.480507in}}%
\pgfpathlineto{\pgfqpoint{6.588425in}{6.555547in}}%
\pgfpathlineto{\pgfqpoint{6.596639in}{6.817403in}}%
\pgfpathlineto{\pgfqpoint{6.597495in}{6.820716in}}%
\pgfpathlineto{\pgfqpoint{6.598008in}{6.819450in}}%
\pgfpathlineto{\pgfqpoint{6.599121in}{6.808456in}}%
\pgfpathlineto{\pgfqpoint{6.601174in}{6.761199in}}%
\pgfpathlineto{\pgfqpoint{6.605623in}{6.591488in}}%
\pgfpathlineto{\pgfqpoint{6.609730in}{6.472415in}}%
\pgfpathlineto{\pgfqpoint{6.611356in}{6.460919in}}%
\pgfpathlineto{\pgfqpoint{6.611613in}{6.461330in}}%
\pgfpathlineto{\pgfqpoint{6.612554in}{6.468029in}}%
\pgfpathlineto{\pgfqpoint{6.614351in}{6.502126in}}%
\pgfpathlineto{\pgfqpoint{6.617688in}{6.618689in}}%
\pgfpathlineto{\pgfqpoint{6.623335in}{6.803090in}}%
\pgfpathlineto{\pgfqpoint{6.625046in}{6.816564in}}%
\pgfpathlineto{\pgfqpoint{6.625303in}{6.816259in}}%
\pgfpathlineto{\pgfqpoint{6.626158in}{6.810843in}}%
\pgfpathlineto{\pgfqpoint{6.627784in}{6.782823in}}%
\pgfpathlineto{\pgfqpoint{6.630779in}{6.684530in}}%
\pgfpathlineto{\pgfqpoint{6.637196in}{6.475339in}}%
\pgfpathlineto{\pgfqpoint{6.638736in}{6.465025in}}%
\pgfpathlineto{\pgfqpoint{6.638993in}{6.465439in}}%
\pgfpathlineto{\pgfqpoint{6.639934in}{6.472167in}}%
\pgfpathlineto{\pgfqpoint{6.641731in}{6.506342in}}%
\pgfpathlineto{\pgfqpoint{6.645068in}{6.622554in}}%
\pgfpathlineto{\pgfqpoint{6.650544in}{6.799140in}}%
\pgfpathlineto{\pgfqpoint{6.652255in}{6.812504in}}%
\pgfpathlineto{\pgfqpoint{6.652512in}{6.812176in}}%
\pgfpathlineto{\pgfqpoint{6.653367in}{6.806671in}}%
\pgfpathlineto{\pgfqpoint{6.655079in}{6.776395in}}%
\pgfpathlineto{\pgfqpoint{6.658244in}{6.670170in}}%
\pgfpathlineto{\pgfqpoint{6.664148in}{6.480365in}}%
\pgfpathlineto{\pgfqpoint{6.665688in}{6.469033in}}%
\pgfpathlineto{\pgfqpoint{6.666031in}{6.469498in}}%
\pgfpathlineto{\pgfqpoint{6.666972in}{6.476374in}}%
\pgfpathlineto{\pgfqpoint{6.668768in}{6.510820in}}%
\pgfpathlineto{\pgfqpoint{6.672191in}{6.630221in}}%
\pgfpathlineto{\pgfqpoint{6.677410in}{6.795426in}}%
\pgfpathlineto{\pgfqpoint{6.679122in}{6.808532in}}%
\pgfpathlineto{\pgfqpoint{6.679293in}{6.808351in}}%
\pgfpathlineto{\pgfqpoint{6.680063in}{6.804166in}}%
\pgfpathlineto{\pgfqpoint{6.681603in}{6.779935in}}%
\pgfpathlineto{\pgfqpoint{6.684426in}{6.692135in}}%
\pgfpathlineto{\pgfqpoint{6.691100in}{6.480690in}}%
\pgfpathlineto{\pgfqpoint{6.692384in}{6.472955in}}%
\pgfpathlineto{\pgfqpoint{6.692812in}{6.473787in}}%
\pgfpathlineto{\pgfqpoint{6.693838in}{6.482689in}}%
\pgfpathlineto{\pgfqpoint{6.695806in}{6.524708in}}%
\pgfpathlineto{\pgfqpoint{6.699999in}{6.677940in}}%
\pgfpathlineto{\pgfqpoint{6.704106in}{6.794272in}}%
\pgfpathlineto{\pgfqpoint{6.705646in}{6.804644in}}%
\pgfpathlineto{\pgfqpoint{6.705903in}{6.804227in}}%
\pgfpathlineto{\pgfqpoint{6.706844in}{6.797463in}}%
\pgfpathlineto{\pgfqpoint{6.708641in}{6.763193in}}%
\pgfpathlineto{\pgfqpoint{6.712149in}{6.641431in}}%
\pgfpathlineto{\pgfqpoint{6.717111in}{6.488936in}}%
\pgfpathlineto{\pgfqpoint{6.718737in}{6.476793in}}%
\pgfpathlineto{\pgfqpoint{6.718994in}{6.477121in}}%
\pgfpathlineto{\pgfqpoint{6.719849in}{6.482647in}}%
\pgfpathlineto{\pgfqpoint{6.721561in}{6.512993in}}%
\pgfpathlineto{\pgfqpoint{6.724726in}{6.618297in}}%
\pgfpathlineto{\pgfqpoint{6.730202in}{6.789361in}}%
\pgfpathlineto{\pgfqpoint{6.731742in}{6.800852in}}%
\pgfpathlineto{\pgfqpoint{6.732085in}{6.800409in}}%
\pgfpathlineto{\pgfqpoint{6.733026in}{6.793569in}}%
\pgfpathlineto{\pgfqpoint{6.734823in}{6.759177in}}%
\pgfpathlineto{\pgfqpoint{6.738331in}{6.637949in}}%
\pgfpathlineto{\pgfqpoint{6.743208in}{6.491482in}}%
\pgfpathlineto{\pgfqpoint{6.744748in}{6.480547in}}%
\pgfpathlineto{\pgfqpoint{6.745005in}{6.480873in}}%
\pgfpathlineto{\pgfqpoint{6.745860in}{6.486398in}}%
\pgfpathlineto{\pgfqpoint{6.747572in}{6.516740in}}%
\pgfpathlineto{\pgfqpoint{6.750823in}{6.624866in}}%
\pgfpathlineto{\pgfqpoint{6.756128in}{6.786677in}}%
\pgfpathlineto{\pgfqpoint{6.757668in}{6.797137in}}%
\pgfpathlineto{\pgfqpoint{6.757925in}{6.796729in}}%
\pgfpathlineto{\pgfqpoint{6.758780in}{6.790934in}}%
\pgfpathlineto{\pgfqpoint{6.760491in}{6.760113in}}%
\pgfpathlineto{\pgfqpoint{6.763828in}{6.648506in}}%
\pgfpathlineto{\pgfqpoint{6.768877in}{6.495517in}}%
\pgfpathlineto{\pgfqpoint{6.770417in}{6.484221in}}%
\pgfpathlineto{\pgfqpoint{6.770759in}{6.484713in}}%
\pgfpathlineto{\pgfqpoint{6.771700in}{6.491693in}}%
\pgfpathlineto{\pgfqpoint{6.773497in}{6.526293in}}%
\pgfpathlineto{\pgfqpoint{6.777091in}{6.649905in}}%
\pgfpathlineto{\pgfqpoint{6.781711in}{6.783657in}}%
\pgfpathlineto{\pgfqpoint{6.783166in}{6.793512in}}%
\pgfpathlineto{\pgfqpoint{6.783508in}{6.792960in}}%
\pgfpathlineto{\pgfqpoint{6.784449in}{6.785816in}}%
\pgfpathlineto{\pgfqpoint{6.786246in}{6.750958in}}%
\pgfpathlineto{\pgfqpoint{6.789925in}{6.624299in}}%
\pgfpathlineto{\pgfqpoint{6.794374in}{6.497542in}}%
\pgfpathlineto{\pgfqpoint{6.795829in}{6.487813in}}%
\pgfpathlineto{\pgfqpoint{6.796171in}{6.488396in}}%
\pgfpathlineto{\pgfqpoint{6.797112in}{6.495624in}}%
\pgfpathlineto{\pgfqpoint{6.798909in}{6.530604in}}%
\pgfpathlineto{\pgfqpoint{6.802674in}{6.660127in}}%
\pgfpathlineto{\pgfqpoint{6.806952in}{6.780218in}}%
\pgfpathlineto{\pgfqpoint{6.808406in}{6.789958in}}%
\pgfpathlineto{\pgfqpoint{6.808749in}{6.789377in}}%
\pgfpathlineto{\pgfqpoint{6.809690in}{6.782155in}}%
\pgfpathlineto{\pgfqpoint{6.811487in}{6.747204in}}%
\pgfpathlineto{\pgfqpoint{6.815251in}{6.618168in}}%
\pgfpathlineto{\pgfqpoint{6.819530in}{6.500110in}}%
\pgfpathlineto{\pgfqpoint{6.820899in}{6.491328in}}%
\pgfpathlineto{\pgfqpoint{6.821241in}{6.491870in}}%
\pgfpathlineto{\pgfqpoint{6.822182in}{6.498989in}}%
\pgfpathlineto{\pgfqpoint{6.823979in}{6.533748in}}%
\pgfpathlineto{\pgfqpoint{6.827744in}{6.662162in}}%
\pgfpathlineto{\pgfqpoint{6.831936in}{6.777391in}}%
\pgfpathlineto{\pgfqpoint{6.833305in}{6.786478in}}%
\pgfpathlineto{\pgfqpoint{6.833647in}{6.786013in}}%
\pgfpathlineto{\pgfqpoint{6.834589in}{6.779105in}}%
\pgfpathlineto{\pgfqpoint{6.836385in}{6.744716in}}%
\pgfpathlineto{\pgfqpoint{6.840150in}{6.617061in}}%
\pgfpathlineto{\pgfqpoint{6.844343in}{6.503245in}}%
\pgfpathlineto{\pgfqpoint{6.845712in}{6.494770in}}%
\pgfpathlineto{\pgfqpoint{6.846054in}{6.495390in}}%
\pgfpathlineto{\pgfqpoint{6.846995in}{6.502714in}}%
\pgfpathlineto{\pgfqpoint{6.848792in}{6.537762in}}%
\pgfpathlineto{\pgfqpoint{6.852728in}{6.671568in}}%
\pgfpathlineto{\pgfqpoint{6.856749in}{6.776029in}}%
\pgfpathlineto{\pgfqpoint{6.857947in}{6.783074in}}%
\pgfpathlineto{\pgfqpoint{6.858375in}{6.782334in}}%
\pgfpathlineto{\pgfqpoint{6.859402in}{6.773628in}}%
\pgfpathlineto{\pgfqpoint{6.861370in}{6.732226in}}%
\pgfpathlineto{\pgfqpoint{6.870183in}{6.498134in}}%
\pgfpathlineto{\pgfqpoint{6.870696in}{6.499320in}}%
\pgfpathlineto{\pgfqpoint{6.871808in}{6.510231in}}%
\pgfpathlineto{\pgfqpoint{6.873947in}{6.559426in}}%
\pgfpathlineto{\pgfqpoint{6.882332in}{6.779748in}}%
\pgfpathlineto{\pgfqpoint{6.882846in}{6.778548in}}%
\pgfpathlineto{\pgfqpoint{6.883958in}{6.767614in}}%
\pgfpathlineto{\pgfqpoint{6.886097in}{6.718447in}}%
\pgfpathlineto{\pgfqpoint{6.894311in}{6.501483in}}%
\pgfpathlineto{\pgfqpoint{6.894568in}{6.501533in}}%
\pgfpathlineto{\pgfqpoint{6.894825in}{6.502202in}}%
\pgfpathlineto{\pgfqpoint{6.895851in}{6.510969in}}%
\pgfpathlineto{\pgfqpoint{6.897819in}{6.552330in}}%
\pgfpathlineto{\pgfqpoint{6.906461in}{6.776483in}}%
\pgfpathlineto{\pgfqpoint{6.906803in}{6.775831in}}%
\pgfpathlineto{\pgfqpoint{6.907830in}{6.767368in}}%
\pgfpathlineto{\pgfqpoint{6.909798in}{6.726538in}}%
\pgfpathlineto{\pgfqpoint{6.918354in}{6.504656in}}%
\pgfpathlineto{\pgfqpoint{6.918782in}{6.505439in}}%
\pgfpathlineto{\pgfqpoint{6.919809in}{6.514224in}}%
\pgfpathlineto{\pgfqpoint{6.921777in}{6.555506in}}%
\pgfpathlineto{\pgfqpoint{6.930247in}{6.773297in}}%
\pgfpathlineto{\pgfqpoint{6.930590in}{6.772758in}}%
\pgfpathlineto{\pgfqpoint{6.931531in}{6.765676in}}%
\pgfpathlineto{\pgfqpoint{6.933328in}{6.731274in}}%
\pgfpathlineto{\pgfqpoint{6.937520in}{6.593472in}}%
\pgfpathlineto{\pgfqpoint{6.941028in}{6.512806in}}%
\pgfpathlineto{\pgfqpoint{6.942055in}{6.507812in}}%
\pgfpathlineto{\pgfqpoint{6.942483in}{6.508637in}}%
\pgfpathlineto{\pgfqpoint{6.943510in}{6.517505in}}%
\pgfpathlineto{\pgfqpoint{6.945477in}{6.558800in}}%
\pgfpathlineto{\pgfqpoint{6.953777in}{6.770168in}}%
\pgfpathlineto{\pgfqpoint{6.954119in}{6.769728in}}%
\pgfpathlineto{\pgfqpoint{6.955060in}{6.762926in}}%
\pgfpathlineto{\pgfqpoint{6.956857in}{6.729068in}}%
\pgfpathlineto{\pgfqpoint{6.961050in}{6.592998in}}%
\pgfpathlineto{\pgfqpoint{6.964558in}{6.515000in}}%
\pgfpathlineto{\pgfqpoint{6.965499in}{6.510903in}}%
\pgfpathlineto{\pgfqpoint{6.966012in}{6.512151in}}%
\pgfpathlineto{\pgfqpoint{6.967125in}{6.523128in}}%
\pgfpathlineto{\pgfqpoint{6.969264in}{6.571815in}}%
\pgfpathlineto{\pgfqpoint{6.976879in}{6.766848in}}%
\pgfpathlineto{\pgfqpoint{6.977136in}{6.767113in}}%
\pgfpathlineto{\pgfqpoint{6.977563in}{6.766186in}}%
\pgfpathlineto{\pgfqpoint{6.978590in}{6.757107in}}%
\pgfpathlineto{\pgfqpoint{6.980558in}{6.715684in}}%
\pgfpathlineto{\pgfqpoint{6.988687in}{6.513930in}}%
\pgfpathlineto{\pgfqpoint{6.988858in}{6.514081in}}%
\pgfpathlineto{\pgfqpoint{6.989628in}{6.518124in}}%
\pgfpathlineto{\pgfqpoint{6.991168in}{6.541794in}}%
\pgfpathlineto{\pgfqpoint{6.994248in}{6.632932in}}%
\pgfpathlineto{\pgfqpoint{6.998954in}{6.757215in}}%
\pgfpathlineto{\pgfqpoint{7.000152in}{6.764117in}}%
\pgfpathlineto{\pgfqpoint{7.000580in}{6.763360in}}%
\pgfpathlineto{\pgfqpoint{7.001606in}{6.754698in}}%
\pgfpathlineto{\pgfqpoint{7.003574in}{6.714060in}}%
\pgfpathlineto{\pgfqpoint{7.011617in}{6.516893in}}%
\pgfpathlineto{\pgfqpoint{7.011788in}{6.517034in}}%
\pgfpathlineto{\pgfqpoint{7.012558in}{6.521022in}}%
\pgfpathlineto{\pgfqpoint{7.014099in}{6.544526in}}%
\pgfpathlineto{\pgfqpoint{7.017179in}{6.634850in}}%
\pgfpathlineto{\pgfqpoint{7.021799in}{6.754617in}}%
\pgfpathlineto{\pgfqpoint{7.022997in}{6.761186in}}%
\pgfpathlineto{\pgfqpoint{7.023425in}{6.760319in}}%
\pgfpathlineto{\pgfqpoint{7.024452in}{6.751423in}}%
\pgfpathlineto{\pgfqpoint{7.026420in}{6.710581in}}%
\pgfpathlineto{\pgfqpoint{7.034291in}{6.519797in}}%
\pgfpathlineto{\pgfqpoint{7.034462in}{6.519900in}}%
\pgfpathlineto{\pgfqpoint{7.035147in}{6.523015in}}%
\pgfpathlineto{\pgfqpoint{7.036602in}{6.543307in}}%
\pgfpathlineto{\pgfqpoint{7.039425in}{6.621900in}}%
\pgfpathlineto{\pgfqpoint{7.044473in}{6.752742in}}%
\pgfpathlineto{\pgfqpoint{7.045586in}{6.758314in}}%
\pgfpathlineto{\pgfqpoint{7.046013in}{6.757419in}}%
\pgfpathlineto{\pgfqpoint{7.047040in}{6.748487in}}%
\pgfpathlineto{\pgfqpoint{7.049008in}{6.707752in}}%
\pgfpathlineto{\pgfqpoint{7.056794in}{6.522639in}}%
\pgfpathlineto{\pgfqpoint{7.056880in}{6.522686in}}%
\pgfpathlineto{\pgfqpoint{7.057479in}{6.524906in}}%
\pgfpathlineto{\pgfqpoint{7.058848in}{6.541910in}}%
\pgfpathlineto{\pgfqpoint{7.061415in}{6.608719in}}%
\pgfpathlineto{\pgfqpoint{7.066976in}{6.751370in}}%
\pgfpathlineto{\pgfqpoint{7.067917in}{6.755501in}}%
\pgfpathlineto{\pgfqpoint{7.068431in}{6.754319in}}%
\pgfpathlineto{\pgfqpoint{7.069543in}{6.743614in}}%
\pgfpathlineto{\pgfqpoint{7.071682in}{6.696240in}}%
\pgfpathlineto{\pgfqpoint{7.078698in}{6.525885in}}%
\pgfpathlineto{\pgfqpoint{7.079040in}{6.525425in}}%
\pgfpathlineto{\pgfqpoint{7.079468in}{6.526366in}}%
\pgfpathlineto{\pgfqpoint{7.080581in}{6.536529in}}%
\pgfpathlineto{\pgfqpoint{7.082720in}{6.583002in}}%
\pgfpathlineto{\pgfqpoint{7.089736in}{6.752290in}}%
\pgfpathlineto{\pgfqpoint{7.090078in}{6.752742in}}%
\pgfpathlineto{\pgfqpoint{7.090506in}{6.751794in}}%
\pgfpathlineto{\pgfqpoint{7.091618in}{6.741632in}}%
\pgfpathlineto{\pgfqpoint{7.093757in}{6.695276in}}%
\pgfpathlineto{\pgfqpoint{7.100773in}{6.528453in}}%
\pgfpathlineto{\pgfqpoint{7.101116in}{6.528181in}}%
\pgfpathlineto{\pgfqpoint{7.101543in}{6.529351in}}%
\pgfpathlineto{\pgfqpoint{7.102656in}{6.540044in}}%
\pgfpathlineto{\pgfqpoint{7.104795in}{6.587040in}}%
\pgfpathlineto{\pgfqpoint{7.111554in}{6.749303in}}%
\pgfpathlineto{\pgfqpoint{7.111982in}{6.750044in}}%
\pgfpathlineto{\pgfqpoint{7.112495in}{6.748724in}}%
\pgfpathlineto{\pgfqpoint{7.113693in}{6.736535in}}%
\pgfpathlineto{\pgfqpoint{7.116003in}{6.683264in}}%
\pgfpathlineto{\pgfqpoint{7.122250in}{6.532357in}}%
\pgfpathlineto{\pgfqpoint{7.122848in}{6.530821in}}%
\pgfpathlineto{\pgfqpoint{7.123362in}{6.532109in}}%
\pgfpathlineto{\pgfqpoint{7.124560in}{6.544201in}}%
\pgfpathlineto{\pgfqpoint{7.126870in}{6.597187in}}%
\pgfpathlineto{\pgfqpoint{7.133030in}{6.745658in}}%
\pgfpathlineto{\pgfqpoint{7.133629in}{6.747401in}}%
\pgfpathlineto{\pgfqpoint{7.134143in}{6.746298in}}%
\pgfpathlineto{\pgfqpoint{7.135255in}{6.735875in}}%
\pgfpathlineto{\pgfqpoint{7.137394in}{6.689640in}}%
\pgfpathlineto{\pgfqpoint{7.144068in}{6.533992in}}%
\pgfpathlineto{\pgfqpoint{7.144496in}{6.533464in}}%
\pgfpathlineto{\pgfqpoint{7.144924in}{6.534599in}}%
\pgfpathlineto{\pgfqpoint{7.146036in}{6.545123in}}%
\pgfpathlineto{\pgfqpoint{7.148175in}{6.591366in}}%
\pgfpathlineto{\pgfqpoint{7.154763in}{6.744193in}}%
\pgfpathlineto{\pgfqpoint{7.155191in}{6.744798in}}%
\pgfpathlineto{\pgfqpoint{7.155619in}{6.743745in}}%
\pgfpathlineto{\pgfqpoint{7.156731in}{6.733447in}}%
\pgfpathlineto{\pgfqpoint{7.158870in}{6.687655in}}%
\pgfpathlineto{\pgfqpoint{7.165373in}{6.536816in}}%
\pgfpathlineto{\pgfqpoint{7.165801in}{6.536000in}}%
\pgfpathlineto{\pgfqpoint{7.166314in}{6.537204in}}%
\pgfpathlineto{\pgfqpoint{7.167426in}{6.547764in}}%
\pgfpathlineto{\pgfqpoint{7.169651in}{6.596146in}}%
\pgfpathlineto{\pgfqpoint{7.175897in}{6.741082in}}%
\pgfpathlineto{\pgfqpoint{7.176411in}{6.742276in}}%
\pgfpathlineto{\pgfqpoint{7.176924in}{6.741096in}}%
\pgfpathlineto{\pgfqpoint{7.178036in}{6.730609in}}%
\pgfpathlineto{\pgfqpoint{7.180261in}{6.682482in}}%
\pgfpathlineto{\pgfqpoint{7.186507in}{6.539501in}}%
\pgfpathlineto{\pgfqpoint{7.187020in}{6.538521in}}%
\pgfpathlineto{\pgfqpoint{7.187534in}{6.539906in}}%
\pgfpathlineto{\pgfqpoint{7.188732in}{6.552051in}}%
\pgfpathlineto{\pgfqpoint{7.191042in}{6.604101in}}%
\pgfpathlineto{\pgfqpoint{7.196860in}{6.738008in}}%
\pgfpathlineto{\pgfqpoint{7.197544in}{6.739777in}}%
\pgfpathlineto{\pgfqpoint{7.197972in}{6.738753in}}%
\pgfpathlineto{\pgfqpoint{7.199085in}{6.728620in}}%
\pgfpathlineto{\pgfqpoint{7.201224in}{6.683618in}}%
\pgfpathlineto{\pgfqpoint{7.207555in}{6.541693in}}%
\pgfpathlineto{\pgfqpoint{7.207983in}{6.540976in}}%
\pgfpathlineto{\pgfqpoint{7.208496in}{6.542271in}}%
\pgfpathlineto{\pgfqpoint{7.209694in}{6.554157in}}%
\pgfpathlineto{\pgfqpoint{7.212005in}{6.605513in}}%
\pgfpathlineto{\pgfqpoint{7.217737in}{6.735547in}}%
\pgfpathlineto{\pgfqpoint{7.218422in}{6.737344in}}%
\pgfpathlineto{\pgfqpoint{7.218935in}{6.735957in}}%
\pgfpathlineto{\pgfqpoint{7.220133in}{6.723902in}}%
\pgfpathlineto{\pgfqpoint{7.222443in}{6.672512in}}%
\pgfpathlineto{\pgfqpoint{7.228090in}{6.545281in}}%
\pgfpathlineto{\pgfqpoint{7.228775in}{6.543388in}}%
\pgfpathlineto{\pgfqpoint{7.229288in}{6.544693in}}%
\pgfpathlineto{\pgfqpoint{7.230486in}{6.556539in}}%
\pgfpathlineto{\pgfqpoint{7.232796in}{6.607475in}}%
\pgfpathlineto{\pgfqpoint{7.238443in}{6.733247in}}%
\pgfpathlineto{\pgfqpoint{7.239128in}{6.734951in}}%
\pgfpathlineto{\pgfqpoint{7.239556in}{6.733913in}}%
\pgfpathlineto{\pgfqpoint{7.240668in}{6.723850in}}%
\pgfpathlineto{\pgfqpoint{7.242807in}{6.679542in}}%
\pgfpathlineto{\pgfqpoint{7.248882in}{6.546685in}}%
\pgfpathlineto{\pgfqpoint{7.249395in}{6.545758in}}%
\pgfpathlineto{\pgfqpoint{7.249823in}{6.546761in}}%
\pgfpathlineto{\pgfqpoint{7.250935in}{6.556708in}}%
\pgfpathlineto{\pgfqpoint{7.253075in}{6.600704in}}%
\pgfpathlineto{\pgfqpoint{7.259149in}{6.731839in}}%
\pgfpathlineto{\pgfqpoint{7.259577in}{6.732628in}}%
\pgfpathlineto{\pgfqpoint{7.260091in}{6.731453in}}%
\pgfpathlineto{\pgfqpoint{7.261203in}{6.721195in}}%
\pgfpathlineto{\pgfqpoint{7.263428in}{6.674687in}}%
\pgfpathlineto{\pgfqpoint{7.269246in}{6.549206in}}%
\pgfpathlineto{\pgfqpoint{7.269759in}{6.548062in}}%
\pgfpathlineto{\pgfqpoint{7.270273in}{6.549222in}}%
\pgfpathlineto{\pgfqpoint{7.271385in}{6.559422in}}%
\pgfpathlineto{\pgfqpoint{7.273610in}{6.605676in}}%
\pgfpathlineto{\pgfqpoint{7.279342in}{6.728999in}}%
\pgfpathlineto{\pgfqpoint{7.279941in}{6.730328in}}%
\pgfpathlineto{\pgfqpoint{7.280455in}{6.728979in}}%
\pgfpathlineto{\pgfqpoint{7.281652in}{6.717209in}}%
\pgfpathlineto{\pgfqpoint{7.284048in}{6.665042in}}%
\pgfpathlineto{\pgfqpoint{7.289267in}{6.552596in}}%
\pgfpathlineto{\pgfqpoint{7.290037in}{6.550340in}}%
\pgfpathlineto{\pgfqpoint{7.290465in}{6.551309in}}%
\pgfpathlineto{\pgfqpoint{7.291578in}{6.561060in}}%
\pgfpathlineto{\pgfqpoint{7.293802in}{6.606347in}}%
\pgfpathlineto{\pgfqpoint{7.299535in}{6.726984in}}%
\pgfpathlineto{\pgfqpoint{7.300048in}{6.728090in}}%
\pgfpathlineto{\pgfqpoint{7.300562in}{6.726917in}}%
\pgfpathlineto{\pgfqpoint{7.301760in}{6.715599in}}%
\pgfpathlineto{\pgfqpoint{7.304070in}{6.666698in}}%
\pgfpathlineto{\pgfqpoint{7.309375in}{6.554498in}}%
\pgfpathlineto{\pgfqpoint{7.310059in}{6.552554in}}%
\pgfpathlineto{\pgfqpoint{7.310572in}{6.553743in}}%
\pgfpathlineto{\pgfqpoint{7.311770in}{6.565057in}}%
\pgfpathlineto{\pgfqpoint{7.314081in}{6.613750in}}%
\pgfpathlineto{\pgfqpoint{7.319385in}{6.724218in}}%
\pgfpathlineto{\pgfqpoint{7.320070in}{6.725874in}}%
\pgfpathlineto{\pgfqpoint{7.320498in}{6.724869in}}%
\pgfpathlineto{\pgfqpoint{7.321610in}{6.715120in}}%
\pgfpathlineto{\pgfqpoint{7.323835in}{6.670375in}}%
\pgfpathlineto{\pgfqpoint{7.329396in}{6.555950in}}%
\pgfpathlineto{\pgfqpoint{7.329995in}{6.554750in}}%
\pgfpathlineto{\pgfqpoint{7.330423in}{6.555770in}}%
\pgfpathlineto{\pgfqpoint{7.331535in}{6.565522in}}%
\pgfpathlineto{\pgfqpoint{7.333760in}{6.610087in}}%
\pgfpathlineto{\pgfqpoint{7.339236in}{6.722341in}}%
\pgfpathlineto{\pgfqpoint{7.339835in}{6.723730in}}%
\pgfpathlineto{\pgfqpoint{7.340348in}{6.722492in}}%
\pgfpathlineto{\pgfqpoint{7.341546in}{6.711185in}}%
\pgfpathlineto{\pgfqpoint{7.343942in}{6.660926in}}%
\pgfpathlineto{\pgfqpoint{7.348990in}{6.558677in}}%
\pgfpathlineto{\pgfqpoint{7.349675in}{6.556871in}}%
\pgfpathlineto{\pgfqpoint{7.350188in}{6.558121in}}%
\pgfpathlineto{\pgfqpoint{7.351386in}{6.569412in}}%
\pgfpathlineto{\pgfqpoint{7.353782in}{6.619423in}}%
\pgfpathlineto{\pgfqpoint{7.358744in}{6.719589in}}%
\pgfpathlineto{\pgfqpoint{7.359429in}{6.721619in}}%
\pgfpathlineto{\pgfqpoint{7.359942in}{6.720549in}}%
\pgfpathlineto{\pgfqpoint{7.361054in}{6.710831in}}%
\pgfpathlineto{\pgfqpoint{7.363279in}{6.666879in}}%
\pgfpathlineto{\pgfqpoint{7.368669in}{6.560106in}}%
\pgfpathlineto{\pgfqpoint{7.369183in}{6.558960in}}%
\pgfpathlineto{\pgfqpoint{7.369696in}{6.560028in}}%
\pgfpathlineto{\pgfqpoint{7.370808in}{6.569708in}}%
\pgfpathlineto{\pgfqpoint{7.373033in}{6.613415in}}%
\pgfpathlineto{\pgfqpoint{7.378338in}{6.718206in}}%
\pgfpathlineto{\pgfqpoint{7.378937in}{6.719543in}}%
\pgfpathlineto{\pgfqpoint{7.379450in}{6.718303in}}%
\pgfpathlineto{\pgfqpoint{7.380648in}{6.707152in}}%
\pgfpathlineto{\pgfqpoint{7.383044in}{6.658013in}}%
\pgfpathlineto{\pgfqpoint{7.387921in}{6.562800in}}%
\pgfpathlineto{\pgfqpoint{7.388605in}{6.561014in}}%
\pgfpathlineto{\pgfqpoint{7.389119in}{6.562232in}}%
\pgfpathlineto{\pgfqpoint{7.390317in}{6.573292in}}%
\pgfpathlineto{\pgfqpoint{7.392712in}{6.622084in}}%
\pgfpathlineto{\pgfqpoint{7.397504in}{6.715458in}}%
\pgfpathlineto{\pgfqpoint{7.398274in}{6.717500in}}%
\pgfpathlineto{\pgfqpoint{7.398702in}{6.716514in}}%
\pgfpathlineto{\pgfqpoint{7.399814in}{6.707080in}}%
\pgfpathlineto{\pgfqpoint{7.402039in}{6.664286in}}%
\pgfpathlineto{\pgfqpoint{7.407258in}{6.564245in}}%
\pgfpathlineto{\pgfqpoint{7.407857in}{6.563032in}}%
\pgfpathlineto{\pgfqpoint{7.408370in}{6.564346in}}%
\pgfpathlineto{\pgfqpoint{7.409568in}{6.575530in}}%
\pgfpathlineto{\pgfqpoint{7.411964in}{6.623986in}}%
\pgfpathlineto{\pgfqpoint{7.416670in}{6.713592in}}%
\pgfpathlineto{\pgfqpoint{7.417354in}{6.715528in}}%
\pgfpathlineto{\pgfqpoint{7.417868in}{6.714459in}}%
\pgfpathlineto{\pgfqpoint{7.418980in}{6.704963in}}%
\pgfpathlineto{\pgfqpoint{7.421290in}{6.660450in}}%
\pgfpathlineto{\pgfqpoint{7.426253in}{6.566559in}}%
\pgfpathlineto{\pgfqpoint{7.426937in}{6.565007in}}%
\pgfpathlineto{\pgfqpoint{7.427365in}{6.565979in}}%
\pgfpathlineto{\pgfqpoint{7.428478in}{6.575270in}}%
\pgfpathlineto{\pgfqpoint{7.430702in}{6.617274in}}%
\pgfpathlineto{\pgfqpoint{7.435836in}{6.712502in}}%
\pgfpathlineto{\pgfqpoint{7.436349in}{6.713578in}}%
\pgfpathlineto{\pgfqpoint{7.436863in}{6.712513in}}%
\pgfpathlineto{\pgfqpoint{7.437975in}{6.703102in}}%
\pgfpathlineto{\pgfqpoint{7.440285in}{6.659139in}}%
\pgfpathlineto{\pgfqpoint{7.445162in}{6.568513in}}%
\pgfpathlineto{\pgfqpoint{7.445847in}{6.566933in}}%
\pgfpathlineto{\pgfqpoint{7.446275in}{6.567868in}}%
\pgfpathlineto{\pgfqpoint{7.447387in}{6.576996in}}%
\pgfpathlineto{\pgfqpoint{7.449611in}{6.618376in}}%
\pgfpathlineto{\pgfqpoint{7.454660in}{6.710526in}}%
\pgfpathlineto{\pgfqpoint{7.455259in}{6.711649in}}%
\pgfpathlineto{\pgfqpoint{7.455686in}{6.710687in}}%
\pgfpathlineto{\pgfqpoint{7.456799in}{6.701528in}}%
\pgfpathlineto{\pgfqpoint{7.459023in}{6.660314in}}%
\pgfpathlineto{\pgfqpoint{7.463986in}{6.570118in}}%
\pgfpathlineto{\pgfqpoint{7.464585in}{6.568815in}}%
\pgfpathlineto{\pgfqpoint{7.465098in}{6.569980in}}%
\pgfpathlineto{\pgfqpoint{7.466296in}{6.580575in}}%
\pgfpathlineto{\pgfqpoint{7.468692in}{6.626849in}}%
\pgfpathlineto{\pgfqpoint{7.473227in}{6.708068in}}%
\pgfpathlineto{\pgfqpoint{7.473911in}{6.709792in}}%
\pgfpathlineto{\pgfqpoint{7.474425in}{6.708641in}}%
\pgfpathlineto{\pgfqpoint{7.475622in}{6.698124in}}%
\pgfpathlineto{\pgfqpoint{7.478104in}{6.650208in}}%
\pgfpathlineto{\pgfqpoint{7.482553in}{6.572210in}}%
\pgfpathlineto{\pgfqpoint{7.483238in}{6.570680in}}%
\pgfpathlineto{\pgfqpoint{7.483665in}{6.571605in}}%
\pgfpathlineto{\pgfqpoint{7.484778in}{6.580557in}}%
\pgfpathlineto{\pgfqpoint{7.487088in}{6.622801in}}%
\pgfpathlineto{\pgfqpoint{7.491879in}{6.706748in}}%
\pgfpathlineto{\pgfqpoint{7.492478in}{6.707951in}}%
\pgfpathlineto{\pgfqpoint{7.492992in}{6.706737in}}%
\pgfpathlineto{\pgfqpoint{7.494190in}{6.696177in}}%
\pgfpathlineto{\pgfqpoint{7.496671in}{6.648757in}}%
\pgfpathlineto{\pgfqpoint{7.501035in}{6.573985in}}%
\pgfpathlineto{\pgfqpoint{7.501719in}{6.572500in}}%
\pgfpathlineto{\pgfqpoint{7.502147in}{6.573431in}}%
\pgfpathlineto{\pgfqpoint{7.503259in}{6.582320in}}%
\pgfpathlineto{\pgfqpoint{7.505569in}{6.624026in}}%
\pgfpathlineto{\pgfqpoint{7.510275in}{6.704922in}}%
\pgfpathlineto{\pgfqpoint{7.510874in}{6.706155in}}%
\pgfpathlineto{\pgfqpoint{7.511388in}{6.704993in}}%
\pgfpathlineto{\pgfqpoint{7.512585in}{6.694639in}}%
\pgfpathlineto{\pgfqpoint{7.515067in}{6.648008in}}%
\pgfpathlineto{\pgfqpoint{7.519345in}{6.575869in}}%
\pgfpathlineto{\pgfqpoint{7.520029in}{6.574270in}}%
\pgfpathlineto{\pgfqpoint{7.520543in}{6.575446in}}%
\pgfpathlineto{\pgfqpoint{7.521741in}{6.585780in}}%
\pgfpathlineto{\pgfqpoint{7.524222in}{6.632110in}}%
\pgfpathlineto{\pgfqpoint{7.528500in}{6.702991in}}%
\pgfpathlineto{\pgfqpoint{7.529099in}{6.704397in}}%
\pgfpathlineto{\pgfqpoint{7.529612in}{6.703408in}}%
\pgfpathlineto{\pgfqpoint{7.530725in}{6.694561in}}%
\pgfpathlineto{\pgfqpoint{7.533035in}{6.653606in}}%
\pgfpathlineto{\pgfqpoint{7.537655in}{6.577086in}}%
\pgfpathlineto{\pgfqpoint{7.538254in}{6.576027in}}%
\pgfpathlineto{\pgfqpoint{7.538682in}{6.576948in}}%
\pgfpathlineto{\pgfqpoint{7.539794in}{6.585650in}}%
\pgfpathlineto{\pgfqpoint{7.542104in}{6.626186in}}%
\pgfpathlineto{\pgfqpoint{7.546639in}{6.701330in}}%
\pgfpathlineto{\pgfqpoint{7.547238in}{6.702669in}}%
\pgfpathlineto{\pgfqpoint{7.547752in}{6.701650in}}%
\pgfpathlineto{\pgfqpoint{7.548949in}{6.691800in}}%
\pgfpathlineto{\pgfqpoint{7.551431in}{6.646855in}}%
\pgfpathlineto{\pgfqpoint{7.555623in}{6.579185in}}%
\pgfpathlineto{\pgfqpoint{7.556308in}{6.577733in}}%
\pgfpathlineto{\pgfqpoint{7.556736in}{6.578621in}}%
\pgfpathlineto{\pgfqpoint{7.557848in}{6.587158in}}%
\pgfpathlineto{\pgfqpoint{7.560158in}{6.627016in}}%
\pgfpathlineto{\pgfqpoint{7.564693in}{6.699907in}}%
\pgfpathlineto{\pgfqpoint{7.565292in}{6.700961in}}%
\pgfpathlineto{\pgfqpoint{7.565720in}{6.700066in}}%
\pgfpathlineto{\pgfqpoint{7.566832in}{6.691554in}}%
\pgfpathlineto{\pgfqpoint{7.569142in}{6.651967in}}%
\pgfpathlineto{\pgfqpoint{7.573591in}{6.580660in}}%
\pgfpathlineto{\pgfqpoint{7.574190in}{6.579396in}}%
\pgfpathlineto{\pgfqpoint{7.574704in}{6.580440in}}%
\pgfpathlineto{\pgfqpoint{7.575902in}{6.590193in}}%
\pgfpathlineto{\pgfqpoint{7.578383in}{6.634178in}}%
\pgfpathlineto{\pgfqpoint{7.582490in}{6.697992in}}%
\pgfpathlineto{\pgfqpoint{7.583089in}{6.699318in}}%
\pgfpathlineto{\pgfqpoint{7.583602in}{6.698341in}}%
\pgfpathlineto{\pgfqpoint{7.584800in}{6.688782in}}%
\pgfpathlineto{\pgfqpoint{7.587281in}{6.645296in}}%
\pgfpathlineto{\pgfqpoint{7.591388in}{6.582283in}}%
\pgfpathlineto{\pgfqpoint{7.591987in}{6.581039in}}%
\pgfpathlineto{\pgfqpoint{7.592501in}{6.582071in}}%
\pgfpathlineto{\pgfqpoint{7.593699in}{6.591698in}}%
\pgfpathlineto{\pgfqpoint{7.596180in}{6.634972in}}%
\pgfpathlineto{\pgfqpoint{7.600201in}{6.696294in}}%
\pgfpathlineto{\pgfqpoint{7.600886in}{6.697676in}}%
\pgfpathlineto{\pgfqpoint{7.601314in}{6.696803in}}%
\pgfpathlineto{\pgfqpoint{7.602426in}{6.688511in}}%
\pgfpathlineto{\pgfqpoint{7.604736in}{6.650161in}}%
\pgfpathlineto{\pgfqpoint{7.609100in}{6.583699in}}%
\pgfpathlineto{\pgfqpoint{7.609699in}{6.582660in}}%
\pgfpathlineto{\pgfqpoint{7.610212in}{6.583840in}}%
\pgfpathlineto{\pgfqpoint{7.611410in}{6.593679in}}%
\pgfpathlineto{\pgfqpoint{7.613977in}{6.638330in}}%
\pgfpathlineto{\pgfqpoint{7.617827in}{6.694813in}}%
\pgfpathlineto{\pgfqpoint{7.618426in}{6.696099in}}%
\pgfpathlineto{\pgfqpoint{7.618940in}{6.695144in}}%
\pgfpathlineto{\pgfqpoint{7.620137in}{6.685840in}}%
\pgfpathlineto{\pgfqpoint{7.622619in}{6.643780in}}%
\pgfpathlineto{\pgfqpoint{7.626640in}{6.585266in}}%
\pgfpathlineto{\pgfqpoint{7.627239in}{6.584236in}}%
\pgfpathlineto{\pgfqpoint{7.627753in}{6.585395in}}%
\pgfpathlineto{\pgfqpoint{7.628950in}{6.595084in}}%
\pgfpathlineto{\pgfqpoint{7.631517in}{6.638951in}}%
\pgfpathlineto{\pgfqpoint{7.635368in}{6.693515in}}%
\pgfpathlineto{\pgfqpoint{7.635967in}{6.694524in}}%
\pgfpathlineto{\pgfqpoint{7.636394in}{6.693685in}}%
\pgfpathlineto{\pgfqpoint{7.637507in}{6.685650in}}%
\pgfpathlineto{\pgfqpoint{7.639902in}{6.646924in}}%
\pgfpathlineto{\pgfqpoint{7.644095in}{6.586661in}}%
\pgfpathlineto{\pgfqpoint{7.644608in}{6.585769in}}%
\pgfpathlineto{\pgfqpoint{7.645122in}{6.586736in}}%
\pgfpathlineto{\pgfqpoint{7.646320in}{6.595910in}}%
\pgfpathlineto{\pgfqpoint{7.648886in}{6.638618in}}%
\pgfpathlineto{\pgfqpoint{7.652737in}{6.692046in}}%
\pgfpathlineto{\pgfqpoint{7.653336in}{6.692989in}}%
\pgfpathlineto{\pgfqpoint{7.653764in}{6.692125in}}%
\pgfpathlineto{\pgfqpoint{7.654876in}{6.684118in}}%
\pgfpathlineto{\pgfqpoint{7.657272in}{6.645954in}}%
\pgfpathlineto{\pgfqpoint{7.661379in}{6.588206in}}%
\pgfpathlineto{\pgfqpoint{7.661977in}{6.587306in}}%
\pgfpathlineto{\pgfqpoint{7.662405in}{6.588191in}}%
\pgfpathlineto{\pgfqpoint{7.663603in}{6.597133in}}%
\pgfpathlineto{\pgfqpoint{7.666170in}{6.638976in}}%
\pgfpathlineto{\pgfqpoint{7.669935in}{6.690420in}}%
\pgfpathlineto{\pgfqpoint{7.670534in}{6.691506in}}%
\pgfpathlineto{\pgfqpoint{7.671047in}{6.690467in}}%
\pgfpathlineto{\pgfqpoint{7.672245in}{6.681303in}}%
\pgfpathlineto{\pgfqpoint{7.674812in}{6.639546in}}%
\pgfpathlineto{\pgfqpoint{7.678577in}{6.589601in}}%
\pgfpathlineto{\pgfqpoint{7.679090in}{6.588766in}}%
\pgfpathlineto{\pgfqpoint{7.679603in}{6.589737in}}%
\pgfpathlineto{\pgfqpoint{7.680801in}{6.598702in}}%
\pgfpathlineto{\pgfqpoint{7.683368in}{6.639942in}}%
\pgfpathlineto{\pgfqpoint{7.687047in}{6.688948in}}%
\pgfpathlineto{\pgfqpoint{7.687646in}{6.690040in}}%
\pgfpathlineto{\pgfqpoint{7.688160in}{6.689035in}}%
\pgfpathlineto{\pgfqpoint{7.689357in}{6.680050in}}%
\pgfpathlineto{\pgfqpoint{7.692010in}{6.637533in}}%
\pgfpathlineto{\pgfqpoint{7.695604in}{6.591153in}}%
\pgfpathlineto{\pgfqpoint{7.696202in}{6.590236in}}%
\pgfpathlineto{\pgfqpoint{7.696630in}{6.591062in}}%
\pgfpathlineto{\pgfqpoint{7.697743in}{6.598751in}}%
\pgfpathlineto{\pgfqpoint{7.700138in}{6.635208in}}%
\pgfpathlineto{\pgfqpoint{7.704074in}{6.687622in}}%
\pgfpathlineto{\pgfqpoint{7.704673in}{6.688594in}}%
\pgfpathlineto{\pgfqpoint{7.705187in}{6.687517in}}%
\pgfpathlineto{\pgfqpoint{7.706384in}{6.678490in}}%
\pgfpathlineto{\pgfqpoint{7.709037in}{6.636606in}}%
\pgfpathlineto{\pgfqpoint{7.712545in}{6.592570in}}%
\pgfpathlineto{\pgfqpoint{7.713144in}{6.591657in}}%
\pgfpathlineto{\pgfqpoint{7.713572in}{6.592464in}}%
\pgfpathlineto{\pgfqpoint{7.714684in}{6.600013in}}%
\pgfpathlineto{\pgfqpoint{7.717080in}{6.635770in}}%
\pgfpathlineto{\pgfqpoint{7.721016in}{6.686413in}}%
\pgfpathlineto{\pgfqpoint{7.721529in}{6.687195in}}%
\pgfpathlineto{\pgfqpoint{7.722042in}{6.686239in}}%
\pgfpathlineto{\pgfqpoint{7.723240in}{6.677581in}}%
\pgfpathlineto{\pgfqpoint{7.725893in}{6.636728in}}%
\pgfpathlineto{\pgfqpoint{7.729401in}{6.593864in}}%
\pgfpathlineto{\pgfqpoint{7.729914in}{6.593035in}}%
\pgfpathlineto{\pgfqpoint{7.730427in}{6.593929in}}%
\pgfpathlineto{\pgfqpoint{7.731625in}{6.602401in}}%
\pgfpathlineto{\pgfqpoint{7.734278in}{6.642737in}}%
\pgfpathlineto{\pgfqpoint{7.737786in}{6.685032in}}%
\pgfpathlineto{\pgfqpoint{7.738299in}{6.685816in}}%
\pgfpathlineto{\pgfqpoint{7.738813in}{6.684890in}}%
\pgfpathlineto{\pgfqpoint{7.740010in}{6.676407in}}%
\pgfpathlineto{\pgfqpoint{7.742663in}{6.636396in}}%
\pgfpathlineto{\pgfqpoint{7.746171in}{6.595059in}}%
\pgfpathlineto{\pgfqpoint{7.746684in}{6.594412in}}%
\pgfpathlineto{\pgfqpoint{7.747198in}{6.595459in}}%
\pgfpathlineto{\pgfqpoint{7.748396in}{6.604142in}}%
\pgfpathlineto{\pgfqpoint{7.751134in}{6.645446in}}%
\pgfpathlineto{\pgfqpoint{7.754471in}{6.683757in}}%
\pgfpathlineto{\pgfqpoint{7.754984in}{6.684460in}}%
\pgfpathlineto{\pgfqpoint{7.755497in}{6.683484in}}%
\pgfpathlineto{\pgfqpoint{7.756695in}{6.675004in}}%
\pgfpathlineto{\pgfqpoint{7.759433in}{6.634246in}}%
\pgfpathlineto{\pgfqpoint{7.762770in}{6.596416in}}%
\pgfpathlineto{\pgfqpoint{7.763283in}{6.595745in}}%
\pgfpathlineto{\pgfqpoint{7.763797in}{6.596740in}}%
\pgfpathlineto{\pgfqpoint{7.764995in}{6.605202in}}%
\pgfpathlineto{\pgfqpoint{7.767733in}{6.645586in}}%
\pgfpathlineto{\pgfqpoint{7.771070in}{6.682568in}}%
\pgfpathlineto{\pgfqpoint{7.771583in}{6.683120in}}%
\pgfpathlineto{\pgfqpoint{7.772011in}{6.682317in}}%
\pgfpathlineto{\pgfqpoint{7.773209in}{6.674270in}}%
\pgfpathlineto{\pgfqpoint{7.775861in}{6.636019in}}%
\pgfpathlineto{\pgfqpoint{7.779284in}{6.597681in}}%
\pgfpathlineto{\pgfqpoint{7.779797in}{6.597059in}}%
\pgfpathlineto{\pgfqpoint{7.780310in}{6.598073in}}%
\pgfpathlineto{\pgfqpoint{7.781508in}{6.606463in}}%
\pgfpathlineto{\pgfqpoint{7.784332in}{6.647427in}}%
\pgfpathlineto{\pgfqpoint{7.787583in}{6.681440in}}%
\pgfpathlineto{\pgfqpoint{7.788011in}{6.681834in}}%
\pgfpathlineto{\pgfqpoint{7.788439in}{6.681100in}}%
\pgfpathlineto{\pgfqpoint{7.789551in}{6.674149in}}%
\pgfpathlineto{\pgfqpoint{7.792032in}{6.640127in}}%
\pgfpathlineto{\pgfqpoint{7.795712in}{6.598867in}}%
\pgfpathlineto{\pgfqpoint{7.796225in}{6.598356in}}%
\pgfpathlineto{\pgfqpoint{7.796653in}{6.599159in}}%
\pgfpathlineto{\pgfqpoint{7.797851in}{6.607044in}}%
\pgfpathlineto{\pgfqpoint{7.800503in}{6.644135in}}%
\pgfpathlineto{\pgfqpoint{7.803840in}{6.679964in}}%
\pgfpathlineto{\pgfqpoint{7.804353in}{6.680566in}}%
\pgfpathlineto{\pgfqpoint{7.804867in}{6.679574in}}%
\pgfpathlineto{\pgfqpoint{7.806150in}{6.670538in}}%
\pgfpathlineto{\pgfqpoint{7.809230in}{6.626615in}}%
\pgfpathlineto{\pgfqpoint{7.812140in}{6.599827in}}%
\pgfpathlineto{\pgfqpoint{7.812567in}{6.599647in}}%
\pgfpathlineto{\pgfqpoint{7.812910in}{6.600293in}}%
\pgfpathlineto{\pgfqpoint{7.814022in}{6.607018in}}%
\pgfpathlineto{\pgfqpoint{7.816503in}{6.639927in}}%
\pgfpathlineto{\pgfqpoint{7.820097in}{6.678776in}}%
\pgfpathlineto{\pgfqpoint{7.820610in}{6.679317in}}%
\pgfpathlineto{\pgfqpoint{7.821038in}{6.678572in}}%
\pgfpathlineto{\pgfqpoint{7.822236in}{6.670997in}}%
\pgfpathlineto{\pgfqpoint{7.824888in}{6.635195in}}%
\pgfpathlineto{\pgfqpoint{7.828225in}{6.601257in}}%
\pgfpathlineto{\pgfqpoint{7.828739in}{6.600860in}}%
\pgfpathlineto{\pgfqpoint{7.829167in}{6.601714in}}%
\pgfpathlineto{\pgfqpoint{7.830364in}{6.609510in}}%
\pgfpathlineto{\pgfqpoint{7.833188in}{6.647701in}}%
\pgfpathlineto{\pgfqpoint{7.836354in}{6.677827in}}%
\pgfpathlineto{\pgfqpoint{7.836782in}{6.678083in}}%
\pgfpathlineto{\pgfqpoint{7.837209in}{6.677272in}}%
\pgfpathlineto{\pgfqpoint{7.838407in}{6.669640in}}%
\pgfpathlineto{\pgfqpoint{7.841231in}{6.631960in}}%
\pgfpathlineto{\pgfqpoint{7.844311in}{6.602459in}}%
\pgfpathlineto{\pgfqpoint{7.844824in}{6.602061in}}%
\pgfpathlineto{\pgfqpoint{7.845252in}{6.602890in}}%
\pgfpathlineto{\pgfqpoint{7.846450in}{6.610516in}}%
\pgfpathlineto{\pgfqpoint{7.849274in}{6.647813in}}%
\pgfpathlineto{\pgfqpoint{7.852354in}{6.676567in}}%
\pgfpathlineto{\pgfqpoint{7.852782in}{6.676914in}}%
\pgfpathlineto{\pgfqpoint{7.853210in}{6.676213in}}%
\pgfpathlineto{\pgfqpoint{7.854407in}{6.668970in}}%
\pgfpathlineto{\pgfqpoint{7.857145in}{6.633578in}}%
\pgfpathlineto{\pgfqpoint{7.860311in}{6.603606in}}%
\pgfpathlineto{\pgfqpoint{7.860825in}{6.603245in}}%
\pgfpathlineto{\pgfqpoint{7.861252in}{6.604084in}}%
\pgfpathlineto{\pgfqpoint{7.862450in}{6.611620in}}%
\pgfpathlineto{\pgfqpoint{7.865359in}{6.649278in}}%
\pgfpathlineto{\pgfqpoint{7.868354in}{6.675518in}}%
\pgfpathlineto{\pgfqpoint{7.868782in}{6.675715in}}%
\pgfpathlineto{\pgfqpoint{7.869210in}{6.674887in}}%
\pgfpathlineto{\pgfqpoint{7.870408in}{6.667432in}}%
\pgfpathlineto{\pgfqpoint{7.873317in}{6.630232in}}%
\pgfpathlineto{\pgfqpoint{7.876311in}{6.604555in}}%
\pgfpathlineto{\pgfqpoint{7.876739in}{6.604415in}}%
\pgfpathlineto{\pgfqpoint{7.877082in}{6.605034in}}%
\pgfpathlineto{\pgfqpoint{7.878279in}{6.612042in}}%
\pgfpathlineto{\pgfqpoint{7.881017in}{6.646178in}}%
\pgfpathlineto{\pgfqpoint{7.884183in}{6.674357in}}%
\pgfpathlineto{\pgfqpoint{7.884611in}{6.674586in}}%
\pgfpathlineto{\pgfqpoint{7.885039in}{6.673809in}}%
\pgfpathlineto{\pgfqpoint{7.886237in}{6.666596in}}%
\pgfpathlineto{\pgfqpoint{7.889146in}{6.630365in}}%
\pgfpathlineto{\pgfqpoint{7.892141in}{6.605633in}}%
\pgfpathlineto{\pgfqpoint{7.892568in}{6.605572in}}%
\pgfpathlineto{\pgfqpoint{7.892911in}{6.606239in}}%
\pgfpathlineto{\pgfqpoint{7.894108in}{6.613293in}}%
\pgfpathlineto{\pgfqpoint{7.896932in}{6.647891in}}%
\pgfpathlineto{\pgfqpoint{7.899927in}{6.673237in}}%
\pgfpathlineto{\pgfqpoint{7.900355in}{6.673474in}}%
\pgfpathlineto{\pgfqpoint{7.900782in}{6.672726in}}%
\pgfpathlineto{\pgfqpoint{7.901980in}{6.665695in}}%
\pgfpathlineto{\pgfqpoint{7.904889in}{6.630383in}}%
\pgfpathlineto{\pgfqpoint{7.907798in}{6.606795in}}%
\pgfpathlineto{\pgfqpoint{7.908226in}{6.606635in}}%
\pgfpathlineto{\pgfqpoint{7.908654in}{6.607448in}}%
\pgfpathlineto{\pgfqpoint{7.909937in}{6.615343in}}%
\pgfpathlineto{\pgfqpoint{7.913446in}{6.657714in}}%
\pgfpathlineto{\pgfqpoint{7.915841in}{6.672407in}}%
\pgfpathlineto{\pgfqpoint{7.916526in}{6.671381in}}%
\pgfpathlineto{\pgfqpoint{7.917809in}{6.663245in}}%
\pgfpathlineto{\pgfqpoint{7.923713in}{6.607661in}}%
\pgfpathlineto{\pgfqpoint{7.924654in}{6.610052in}}%
\pgfpathlineto{\pgfqpoint{7.926365in}{6.624573in}}%
\pgfpathlineto{\pgfqpoint{7.931499in}{6.671344in}}%
\pgfpathlineto{\pgfqpoint{7.932184in}{6.670009in}}%
\pgfpathlineto{\pgfqpoint{7.933638in}{6.659884in}}%
\pgfpathlineto{\pgfqpoint{7.939200in}{6.608716in}}%
\pgfpathlineto{\pgfqpoint{7.939799in}{6.609601in}}%
\pgfpathlineto{\pgfqpoint{7.941082in}{6.617257in}}%
\pgfpathlineto{\pgfqpoint{7.946986in}{6.670296in}}%
\pgfpathlineto{\pgfqpoint{7.947927in}{6.667841in}}%
\pgfpathlineto{\pgfqpoint{7.949638in}{6.653629in}}%
\pgfpathlineto{\pgfqpoint{7.954687in}{6.609756in}}%
\pgfpathlineto{\pgfqpoint{7.955371in}{6.611089in}}%
\pgfpathlineto{\pgfqpoint{7.956826in}{6.620944in}}%
\pgfpathlineto{\pgfqpoint{7.962302in}{6.669280in}}%
\pgfpathlineto{\pgfqpoint{7.962815in}{6.668669in}}%
\pgfpathlineto{\pgfqpoint{7.964013in}{6.662454in}}%
\pgfpathlineto{\pgfqpoint{7.966922in}{6.630810in}}%
\pgfpathlineto{\pgfqpoint{7.969746in}{6.610889in}}%
\pgfpathlineto{\pgfqpoint{7.970173in}{6.610856in}}%
\pgfpathlineto{\pgfqpoint{7.970516in}{6.611471in}}%
\pgfpathlineto{\pgfqpoint{7.971714in}{6.617833in}}%
\pgfpathlineto{\pgfqpoint{7.974794in}{6.651074in}}%
\pgfpathlineto{\pgfqpoint{7.977446in}{6.668215in}}%
\pgfpathlineto{\pgfqpoint{7.977960in}{6.667990in}}%
\pgfpathlineto{\pgfqpoint{7.978216in}{6.667403in}}%
\pgfpathlineto{\pgfqpoint{7.979500in}{6.660101in}}%
\pgfpathlineto{\pgfqpoint{7.985232in}{6.611748in}}%
\pgfpathlineto{\pgfqpoint{7.986088in}{6.613502in}}%
\pgfpathlineto{\pgfqpoint{7.987628in}{6.624410in}}%
\pgfpathlineto{\pgfqpoint{7.992848in}{6.667303in}}%
\pgfpathlineto{\pgfqpoint{7.993104in}{6.667114in}}%
\pgfpathlineto{\pgfqpoint{7.994045in}{6.663837in}}%
\pgfpathlineto{\pgfqpoint{7.996013in}{6.646600in}}%
\pgfpathlineto{\pgfqpoint{8.000120in}{6.612886in}}%
\pgfpathlineto{\pgfqpoint{8.000634in}{6.612850in}}%
\pgfpathlineto{\pgfqpoint{8.000976in}{6.613505in}}%
\pgfpathlineto{\pgfqpoint{8.002259in}{6.620439in}}%
\pgfpathlineto{\pgfqpoint{8.007992in}{6.666344in}}%
\pgfpathlineto{\pgfqpoint{8.008848in}{6.664485in}}%
\pgfpathlineto{\pgfqpoint{8.010473in}{6.652936in}}%
\pgfpathlineto{\pgfqpoint{8.015522in}{6.613663in}}%
\pgfpathlineto{\pgfqpoint{8.015607in}{6.613699in}}%
\pgfpathlineto{\pgfqpoint{8.016377in}{6.615498in}}%
\pgfpathlineto{\pgfqpoint{8.018003in}{6.626898in}}%
\pgfpathlineto{\pgfqpoint{8.022966in}{6.665416in}}%
\pgfpathlineto{\pgfqpoint{8.023137in}{6.665362in}}%
\pgfpathlineto{\pgfqpoint{8.023992in}{6.663154in}}%
\pgfpathlineto{\pgfqpoint{8.025704in}{6.650520in}}%
\pgfpathlineto{\pgfqpoint{8.030409in}{6.614588in}}%
\pgfpathlineto{\pgfqpoint{8.031094in}{6.615441in}}%
\pgfpathlineto{\pgfqpoint{8.032377in}{6.622228in}}%
\pgfpathlineto{\pgfqpoint{8.037939in}{6.664504in}}%
\pgfpathlineto{\pgfqpoint{8.038709in}{6.663210in}}%
\pgfpathlineto{\pgfqpoint{8.040164in}{6.654429in}}%
\pgfpathlineto{\pgfqpoint{8.045383in}{6.615482in}}%
\pgfpathlineto{\pgfqpoint{8.045811in}{6.615874in}}%
\pgfpathlineto{\pgfqpoint{8.046923in}{6.620423in}}%
\pgfpathlineto{\pgfqpoint{8.049490in}{6.643572in}}%
\pgfpathlineto{\pgfqpoint{8.052570in}{6.663489in}}%
\pgfpathlineto{\pgfqpoint{8.053083in}{6.663455in}}%
\pgfpathlineto{\pgfqpoint{8.053426in}{6.662811in}}%
\pgfpathlineto{\pgfqpoint{8.054709in}{6.656334in}}%
\pgfpathlineto{\pgfqpoint{8.060271in}{6.616368in}}%
\pgfpathlineto{\pgfqpoint{8.061041in}{6.617775in}}%
\pgfpathlineto{\pgfqpoint{8.062581in}{6.627158in}}%
\pgfpathlineto{\pgfqpoint{8.067629in}{6.662742in}}%
\pgfpathlineto{\pgfqpoint{8.067886in}{6.662579in}}%
\pgfpathlineto{\pgfqpoint{8.068827in}{6.659707in}}%
\pgfpathlineto{\pgfqpoint{8.070880in}{6.643872in}}%
\pgfpathlineto{\pgfqpoint{8.074731in}{6.617342in}}%
\pgfpathlineto{\pgfqpoint{8.075244in}{6.617375in}}%
\pgfpathlineto{\pgfqpoint{8.075586in}{6.617993in}}%
\pgfpathlineto{\pgfqpoint{8.076870in}{6.624211in}}%
\pgfpathlineto{\pgfqpoint{8.082346in}{6.661892in}}%
\pgfpathlineto{\pgfqpoint{8.083116in}{6.660620in}}%
\pgfpathlineto{\pgfqpoint{8.084656in}{6.651747in}}%
\pgfpathlineto{\pgfqpoint{8.089704in}{6.618069in}}%
\pgfpathlineto{\pgfqpoint{8.089875in}{6.618166in}}%
\pgfpathlineto{\pgfqpoint{8.090817in}{6.620736in}}%
\pgfpathlineto{\pgfqpoint{8.092784in}{6.634875in}}%
\pgfpathlineto{\pgfqpoint{8.096720in}{6.660954in}}%
\pgfpathlineto{\pgfqpoint{8.097319in}{6.660803in}}%
\pgfpathlineto{\pgfqpoint{8.097576in}{6.660309in}}%
\pgfpathlineto{\pgfqpoint{8.098945in}{6.653742in}}%
\pgfpathlineto{\pgfqpoint{8.104250in}{6.618882in}}%
\pgfpathlineto{\pgfqpoint{8.104849in}{6.619587in}}%
\pgfpathlineto{\pgfqpoint{8.106132in}{6.625420in}}%
\pgfpathlineto{\pgfqpoint{8.111523in}{6.660253in}}%
\pgfpathlineto{\pgfqpoint{8.112293in}{6.659105in}}%
\pgfpathlineto{\pgfqpoint{8.113833in}{6.650815in}}%
\pgfpathlineto{\pgfqpoint{8.118795in}{6.619684in}}%
\pgfpathlineto{\pgfqpoint{8.119052in}{6.619837in}}%
\pgfpathlineto{\pgfqpoint{8.120079in}{6.622849in}}%
\pgfpathlineto{\pgfqpoint{8.122304in}{6.638947in}}%
\pgfpathlineto{\pgfqpoint{8.125726in}{6.659325in}}%
\pgfpathlineto{\pgfqpoint{8.126325in}{6.659264in}}%
\pgfpathlineto{\pgfqpoint{8.126582in}{6.658832in}}%
\pgfpathlineto{\pgfqpoint{8.127865in}{6.653334in}}%
\pgfpathlineto{\pgfqpoint{8.133256in}{6.620467in}}%
\pgfpathlineto{\pgfqpoint{8.133940in}{6.621429in}}%
\pgfpathlineto{\pgfqpoint{8.135395in}{6.628546in}}%
\pgfpathlineto{\pgfqpoint{8.140443in}{6.658689in}}%
\pgfpathlineto{\pgfqpoint{8.140785in}{6.658432in}}%
\pgfpathlineto{\pgfqpoint{8.141897in}{6.654819in}}%
\pgfpathlineto{\pgfqpoint{8.144464in}{6.636084in}}%
\pgfpathlineto{\pgfqpoint{8.147373in}{6.621303in}}%
\pgfpathlineto{\pgfqpoint{8.148058in}{6.621626in}}%
\pgfpathlineto{\pgfqpoint{8.148229in}{6.621964in}}%
\pgfpathlineto{\pgfqpoint{8.149598in}{6.627982in}}%
\pgfpathlineto{\pgfqpoint{8.154817in}{6.657929in}}%
\pgfpathlineto{\pgfqpoint{8.155331in}{6.657332in}}%
\pgfpathlineto{\pgfqpoint{8.156614in}{6.652138in}}%
\pgfpathlineto{\pgfqpoint{8.161919in}{6.621971in}}%
\pgfpathlineto{\pgfqpoint{8.162603in}{6.622863in}}%
\pgfpathlineto{\pgfqpoint{8.164058in}{6.629525in}}%
\pgfpathlineto{\pgfqpoint{8.169021in}{6.657205in}}%
\pgfpathlineto{\pgfqpoint{8.169363in}{6.656985in}}%
\pgfpathlineto{\pgfqpoint{8.170475in}{6.653652in}}%
\pgfpathlineto{\pgfqpoint{8.173042in}{6.636207in}}%
\pgfpathlineto{\pgfqpoint{8.175951in}{6.622733in}}%
\pgfpathlineto{\pgfqpoint{8.176807in}{6.623507in}}%
\pgfpathlineto{\pgfqpoint{8.178261in}{6.629873in}}%
\pgfpathlineto{\pgfqpoint{8.183224in}{6.656489in}}%
\pgfpathlineto{\pgfqpoint{8.183566in}{6.656260in}}%
\pgfpathlineto{\pgfqpoint{8.184679in}{6.652988in}}%
\pgfpathlineto{\pgfqpoint{8.187331in}{6.635504in}}%
\pgfpathlineto{\pgfqpoint{8.190155in}{6.623417in}}%
\pgfpathlineto{\pgfqpoint{8.191010in}{6.624301in}}%
\pgfpathlineto{\pgfqpoint{8.192550in}{6.631160in}}%
\pgfpathlineto{\pgfqpoint{8.197342in}{6.655792in}}%
\pgfpathlineto{\pgfqpoint{8.197598in}{6.655662in}}%
\pgfpathlineto{\pgfqpoint{8.198625in}{6.653122in}}%
\pgfpathlineto{\pgfqpoint{8.200935in}{6.639129in}}%
\pgfpathlineto{\pgfqpoint{8.204101in}{6.624179in}}%
\pgfpathlineto{\pgfqpoint{8.204871in}{6.624520in}}%
\pgfpathlineto{\pgfqpoint{8.205042in}{6.624844in}}%
\pgfpathlineto{\pgfqpoint{8.206497in}{6.630788in}}%
\pgfpathlineto{\pgfqpoint{8.211374in}{6.655114in}}%
\pgfpathlineto{\pgfqpoint{8.211716in}{6.654919in}}%
\pgfpathlineto{\pgfqpoint{8.212829in}{6.651914in}}%
\pgfpathlineto{\pgfqpoint{8.215481in}{6.635710in}}%
\pgfpathlineto{\pgfqpoint{8.218219in}{6.624786in}}%
\pgfpathlineto{\pgfqpoint{8.219075in}{6.625527in}}%
\pgfpathlineto{\pgfqpoint{8.220529in}{6.631328in}}%
\pgfpathlineto{\pgfqpoint{8.225406in}{6.654447in}}%
\pgfpathlineto{\pgfqpoint{8.225663in}{6.654298in}}%
\pgfpathlineto{\pgfqpoint{8.226775in}{6.651496in}}%
\pgfpathlineto{\pgfqpoint{8.229342in}{6.636513in}}%
\pgfpathlineto{\pgfqpoint{8.232166in}{6.625437in}}%
\pgfpathlineto{\pgfqpoint{8.233021in}{6.626154in}}%
\pgfpathlineto{\pgfqpoint{8.234476in}{6.631754in}}%
\pgfpathlineto{\pgfqpoint{8.239353in}{6.653795in}}%
\pgfpathlineto{\pgfqpoint{8.239610in}{6.653626in}}%
\pgfpathlineto{\pgfqpoint{8.240722in}{6.650821in}}%
\pgfpathlineto{\pgfqpoint{8.243460in}{6.635301in}}%
\pgfpathlineto{\pgfqpoint{8.246112in}{6.626051in}}%
\pgfpathlineto{\pgfqpoint{8.246968in}{6.626905in}}%
\pgfpathlineto{\pgfqpoint{8.248508in}{6.632974in}}%
\pgfpathlineto{\pgfqpoint{8.253129in}{6.653182in}}%
\pgfpathlineto{\pgfqpoint{8.253300in}{6.653126in}}%
\pgfpathlineto{\pgfqpoint{8.254326in}{6.651144in}}%
\pgfpathlineto{\pgfqpoint{8.256465in}{6.640364in}}%
\pgfpathlineto{\pgfqpoint{8.259802in}{6.626712in}}%
\pgfpathlineto{\pgfqpoint{8.260744in}{6.627418in}}%
\pgfpathlineto{\pgfqpoint{8.262284in}{6.633150in}}%
\pgfpathlineto{\pgfqpoint{8.266990in}{6.652556in}}%
\pgfpathlineto{\pgfqpoint{8.267075in}{6.652520in}}%
\pgfpathlineto{\pgfqpoint{8.268102in}{6.650619in}}%
\pgfpathlineto{\pgfqpoint{8.270241in}{6.640264in}}%
\pgfpathlineto{\pgfqpoint{8.273578in}{6.627297in}}%
\pgfpathlineto{\pgfqpoint{8.274519in}{6.628049in}}%
\pgfpathlineto{\pgfqpoint{8.276059in}{6.633653in}}%
\pgfpathlineto{\pgfqpoint{8.280680in}{6.651970in}}%
\pgfpathlineto{\pgfqpoint{8.280765in}{6.651941in}}%
\pgfpathlineto{\pgfqpoint{8.281792in}{6.650168in}}%
\pgfpathlineto{\pgfqpoint{8.283931in}{6.640284in}}%
\pgfpathlineto{\pgfqpoint{8.287268in}{6.627877in}}%
\pgfpathlineto{\pgfqpoint{8.288209in}{6.628621in}}%
\pgfpathlineto{\pgfqpoint{8.289835in}{6.634447in}}%
\pgfpathlineto{\pgfqpoint{8.294284in}{6.651402in}}%
\pgfpathlineto{\pgfqpoint{8.294370in}{6.651383in}}%
\pgfpathlineto{\pgfqpoint{8.295396in}{6.649785in}}%
\pgfpathlineto{\pgfqpoint{8.297450in}{6.640867in}}%
\pgfpathlineto{\pgfqpoint{8.300872in}{6.628448in}}%
\pgfpathlineto{\pgfqpoint{8.301814in}{6.629136in}}%
\pgfpathlineto{\pgfqpoint{8.303439in}{6.634694in}}%
\pgfpathlineto{\pgfqpoint{8.307888in}{6.650840in}}%
\pgfpathlineto{\pgfqpoint{8.307974in}{6.650818in}}%
\pgfpathlineto{\pgfqpoint{8.309001in}{6.649239in}}%
\pgfpathlineto{\pgfqpoint{8.311140in}{6.640209in}}%
\pgfpathlineto{\pgfqpoint{8.314477in}{6.628982in}}%
\pgfpathlineto{\pgfqpoint{8.315504in}{6.629920in}}%
\pgfpathlineto{\pgfqpoint{8.317300in}{6.636408in}}%
\pgfpathlineto{\pgfqpoint{8.321322in}{6.650301in}}%
\pgfpathlineto{\pgfqpoint{8.322349in}{6.649175in}}%
\pgfpathlineto{\pgfqpoint{8.324231in}{6.642188in}}%
\pgfpathlineto{\pgfqpoint{8.327996in}{6.629511in}}%
\pgfpathlineto{\pgfqpoint{8.329022in}{6.630472in}}%
\pgfpathlineto{\pgfqpoint{8.330819in}{6.636755in}}%
\pgfpathlineto{\pgfqpoint{8.334755in}{6.649771in}}%
\pgfpathlineto{\pgfqpoint{8.335782in}{6.648750in}}%
\pgfpathlineto{\pgfqpoint{8.337664in}{6.642130in}}%
\pgfpathlineto{\pgfqpoint{8.341429in}{6.630029in}}%
\pgfpathlineto{\pgfqpoint{8.342456in}{6.630958in}}%
\pgfpathlineto{\pgfqpoint{8.344253in}{6.636977in}}%
\pgfpathlineto{\pgfqpoint{8.348188in}{6.649262in}}%
\pgfpathlineto{\pgfqpoint{8.349215in}{6.648214in}}%
\pgfpathlineto{\pgfqpoint{8.351098in}{6.641807in}}%
\pgfpathlineto{\pgfqpoint{8.354777in}{6.630537in}}%
\pgfpathlineto{\pgfqpoint{8.355803in}{6.631383in}}%
\pgfpathlineto{\pgfqpoint{8.357600in}{6.637085in}}%
\pgfpathlineto{\pgfqpoint{8.361536in}{6.648766in}}%
\pgfpathlineto{\pgfqpoint{8.362563in}{6.647745in}}%
\pgfpathlineto{\pgfqpoint{8.364531in}{6.641256in}}%
\pgfpathlineto{\pgfqpoint{8.368124in}{6.631019in}}%
\pgfpathlineto{\pgfqpoint{8.369237in}{6.632065in}}%
\pgfpathlineto{\pgfqpoint{8.371205in}{6.638458in}}%
\pgfpathlineto{\pgfqpoint{8.374713in}{6.648274in}}%
\pgfpathlineto{\pgfqpoint{8.375825in}{6.647341in}}%
\pgfpathlineto{\pgfqpoint{8.377793in}{6.641191in}}%
\pgfpathlineto{\pgfqpoint{8.381387in}{6.631490in}}%
\pgfpathlineto{\pgfqpoint{8.382499in}{6.632514in}}%
\pgfpathlineto{\pgfqpoint{8.384552in}{6.638960in}}%
\pgfpathlineto{\pgfqpoint{8.387975in}{6.647815in}}%
\pgfpathlineto{\pgfqpoint{8.389087in}{6.646843in}}%
\pgfpathlineto{\pgfqpoint{8.391141in}{6.640584in}}%
\pgfpathlineto{\pgfqpoint{8.394563in}{6.631950in}}%
\pgfpathlineto{\pgfqpoint{8.395676in}{6.632899in}}%
\pgfpathlineto{\pgfqpoint{8.397729in}{6.639003in}}%
\pgfpathlineto{\pgfqpoint{8.401152in}{6.647366in}}%
\pgfpathlineto{\pgfqpoint{8.402264in}{6.646413in}}%
\pgfpathlineto{\pgfqpoint{8.404317in}{6.640437in}}%
\pgfpathlineto{\pgfqpoint{8.407654in}{6.632399in}}%
\pgfpathlineto{\pgfqpoint{8.408852in}{6.633372in}}%
\pgfpathlineto{\pgfqpoint{8.410991in}{6.639543in}}%
\pgfpathlineto{\pgfqpoint{8.414243in}{6.646929in}}%
\pgfpathlineto{\pgfqpoint{8.415440in}{6.645898in}}%
\pgfpathlineto{\pgfqpoint{8.417665in}{6.639527in}}%
\pgfpathlineto{\pgfqpoint{8.420745in}{6.632825in}}%
\pgfpathlineto{\pgfqpoint{8.421943in}{6.633778in}}%
\pgfpathlineto{\pgfqpoint{8.424082in}{6.639653in}}%
\pgfpathlineto{\pgfqpoint{8.427248in}{6.646502in}}%
\pgfpathlineto{\pgfqpoint{8.428446in}{6.645599in}}%
\pgfpathlineto{\pgfqpoint{8.430585in}{6.639904in}}%
\pgfpathlineto{\pgfqpoint{8.433751in}{6.633241in}}%
\pgfpathlineto{\pgfqpoint{8.434949in}{6.634121in}}%
\pgfpathlineto{\pgfqpoint{8.437088in}{6.639660in}}%
\pgfpathlineto{\pgfqpoint{8.440254in}{6.646098in}}%
\pgfpathlineto{\pgfqpoint{8.441537in}{6.645079in}}%
\pgfpathlineto{\pgfqpoint{8.443933in}{6.638782in}}%
\pgfpathlineto{\pgfqpoint{8.446756in}{6.633635in}}%
\pgfpathlineto{\pgfqpoint{8.448040in}{6.634676in}}%
\pgfpathlineto{\pgfqpoint{8.450521in}{6.641073in}}%
\pgfpathlineto{\pgfqpoint{8.453259in}{6.645711in}}%
\pgfpathlineto{\pgfqpoint{8.454543in}{6.644628in}}%
\pgfpathlineto{\pgfqpoint{8.457195in}{6.637905in}}%
\pgfpathlineto{\pgfqpoint{8.459762in}{6.634019in}}%
\pgfpathlineto{\pgfqpoint{8.461045in}{6.635162in}}%
\pgfpathlineto{\pgfqpoint{8.463869in}{6.642169in}}%
\pgfpathlineto{\pgfqpoint{8.466350in}{6.645309in}}%
\pgfpathlineto{\pgfqpoint{8.467719in}{6.643802in}}%
\pgfpathlineto{\pgfqpoint{8.473024in}{6.634516in}}%
\pgfpathlineto{\pgfqpoint{8.473537in}{6.634968in}}%
\pgfpathlineto{\pgfqpoint{8.475591in}{6.639207in}}%
\pgfpathlineto{\pgfqpoint{8.478928in}{6.644966in}}%
\pgfpathlineto{\pgfqpoint{8.480297in}{6.643934in}}%
\pgfpathlineto{\pgfqpoint{8.483206in}{6.637375in}}%
\pgfpathlineto{\pgfqpoint{8.485602in}{6.634777in}}%
\pgfpathlineto{\pgfqpoint{8.487056in}{6.636351in}}%
\pgfpathlineto{\pgfqpoint{8.492190in}{6.644508in}}%
\pgfpathlineto{\pgfqpoint{8.492532in}{6.644271in}}%
\pgfpathlineto{\pgfqpoint{8.494415in}{6.641057in}}%
\pgfpathlineto{\pgfqpoint{8.498179in}{6.635093in}}%
\pgfpathlineto{\pgfqpoint{8.499634in}{6.636323in}}%
\pgfpathlineto{\pgfqpoint{8.505024in}{6.644117in}}%
\pgfpathlineto{\pgfqpoint{8.505538in}{6.643677in}}%
\pgfpathlineto{\pgfqpoint{8.507848in}{6.639391in}}%
\pgfpathlineto{\pgfqpoint{8.510843in}{6.635423in}}%
\pgfpathlineto{\pgfqpoint{8.512297in}{6.636532in}}%
\pgfpathlineto{\pgfqpoint{8.517859in}{6.643697in}}%
\pgfpathlineto{\pgfqpoint{8.518201in}{6.643397in}}%
\pgfpathlineto{\pgfqpoint{8.520597in}{6.639228in}}%
\pgfpathlineto{\pgfqpoint{8.523506in}{6.635741in}}%
\pgfpathlineto{\pgfqpoint{8.525046in}{6.636923in}}%
\pgfpathlineto{\pgfqpoint{8.530522in}{6.643371in}}%
\pgfpathlineto{\pgfqpoint{8.530693in}{6.643237in}}%
\pgfpathlineto{\pgfqpoint{8.532918in}{6.639801in}}%
\pgfpathlineto{\pgfqpoint{8.536084in}{6.636044in}}%
\pgfpathlineto{\pgfqpoint{8.537709in}{6.637243in}}%
\pgfpathlineto{\pgfqpoint{8.543185in}{6.643024in}}%
\pgfpathlineto{\pgfqpoint{8.545410in}{6.639906in}}%
\pgfpathlineto{\pgfqpoint{8.548661in}{6.636337in}}%
\pgfpathlineto{\pgfqpoint{8.550372in}{6.637609in}}%
\pgfpathlineto{\pgfqpoint{8.555592in}{6.642832in}}%
\pgfpathlineto{\pgfqpoint{8.557731in}{6.640190in}}%
\pgfpathlineto{\pgfqpoint{8.561239in}{6.636623in}}%
\pgfpathlineto{\pgfqpoint{8.563036in}{6.638013in}}%
\pgfpathlineto{\pgfqpoint{8.567913in}{6.642651in}}%
\pgfpathlineto{\pgfqpoint{8.569966in}{6.640477in}}%
\pgfpathlineto{\pgfqpoint{8.573731in}{6.636895in}}%
\pgfpathlineto{\pgfqpoint{8.575613in}{6.638340in}}%
\pgfpathlineto{\pgfqpoint{8.580234in}{6.642440in}}%
\pgfpathlineto{\pgfqpoint{8.582287in}{6.640529in}}%
\pgfpathlineto{\pgfqpoint{8.586138in}{6.637149in}}%
\pgfpathlineto{\pgfqpoint{8.588106in}{6.638593in}}%
\pgfpathlineto{\pgfqpoint{8.592555in}{6.642215in}}%
\pgfpathlineto{\pgfqpoint{8.594694in}{6.640401in}}%
\pgfpathlineto{\pgfqpoint{8.598544in}{6.637398in}}%
\pgfpathlineto{\pgfqpoint{8.600683in}{6.638970in}}%
\pgfpathlineto{\pgfqpoint{8.604790in}{6.642007in}}%
\pgfpathlineto{\pgfqpoint{8.607015in}{6.640328in}}%
\pgfpathlineto{\pgfqpoint{8.610865in}{6.637628in}}%
\pgfpathlineto{\pgfqpoint{8.613090in}{6.639173in}}%
\pgfpathlineto{\pgfqpoint{8.617026in}{6.641796in}}%
\pgfpathlineto{\pgfqpoint{8.619336in}{6.640213in}}%
\pgfpathlineto{\pgfqpoint{8.623186in}{6.637855in}}%
\pgfpathlineto{\pgfqpoint{8.625582in}{6.639477in}}%
\pgfpathlineto{\pgfqpoint{8.629261in}{6.641587in}}%
\pgfpathlineto{\pgfqpoint{8.631742in}{6.639991in}}%
\pgfpathlineto{\pgfqpoint{8.635422in}{6.638062in}}%
\pgfpathlineto{\pgfqpoint{8.638074in}{6.639772in}}%
\pgfpathlineto{\pgfqpoint{8.641497in}{6.641383in}}%
\pgfpathlineto{\pgfqpoint{8.644234in}{6.639697in}}%
\pgfpathlineto{\pgfqpoint{8.647657in}{6.638265in}}%
\pgfpathlineto{\pgfqpoint{8.650566in}{6.640043in}}%
\pgfpathlineto{\pgfqpoint{8.653817in}{6.641162in}}%
\pgfpathlineto{\pgfqpoint{8.657069in}{6.639199in}}%
\pgfpathlineto{\pgfqpoint{8.660064in}{6.638511in}}%
\pgfpathlineto{\pgfqpoint{8.668192in}{6.639867in}}%
\pgfpathlineto{\pgfqpoint{8.672042in}{6.638646in}}%
\pgfpathlineto{\pgfqpoint{8.680599in}{6.639640in}}%
\pgfpathlineto{\pgfqpoint{8.684278in}{6.638843in}}%
\pgfpathlineto{\pgfqpoint{8.692149in}{6.639872in}}%
\pgfpathlineto{\pgfqpoint{8.696342in}{6.638994in}}%
\pgfpathlineto{\pgfqpoint{8.704128in}{6.639841in}}%
\pgfpathlineto{\pgfqpoint{8.708492in}{6.639160in}}%
\pgfpathlineto{\pgfqpoint{8.715850in}{6.639893in}}%
\pgfpathlineto{\pgfqpoint{8.720642in}{6.639319in}}%
\pgfpathlineto{\pgfqpoint{8.727572in}{6.639912in}}%
\pgfpathlineto{\pgfqpoint{8.732877in}{6.639485in}}%
\pgfpathlineto{\pgfqpoint{8.739294in}{6.639908in}}%
\pgfpathlineto{\pgfqpoint{8.745113in}{6.639629in}}%
\pgfpathlineto{\pgfqpoint{8.751273in}{6.639840in}}%
\pgfpathlineto{\pgfqpoint{8.757434in}{6.639757in}}%
\pgfpathlineto{\pgfqpoint{8.763680in}{6.639721in}}%
\pgfpathlineto{\pgfqpoint{8.770097in}{6.639881in}}%
\pgfpathlineto{\pgfqpoint{8.777370in}{6.639578in}}%
\pgfpathlineto{\pgfqpoint{8.813819in}{6.639727in}}%
\pgfpathlineto{\pgfqpoint{8.850354in}{6.639769in}}%
\pgfpathlineto{\pgfqpoint{8.850354in}{6.639769in}}%
\pgfusepath{stroke}%
\end{pgfscope}%
\begin{pgfscope}%
\pgfpathrectangle{\pgfqpoint{0.294194in}{5.519538in}}{\pgfqpoint{8.556160in}{2.240462in}}%
\pgfusepath{clip}%
\pgfsetbuttcap%
\pgfsetroundjoin%
\pgfsetlinewidth{1.505625pt}%
\definecolor{currentstroke}{rgb}{0.000000,0.000000,0.000000}%
\pgfsetstrokecolor{currentstroke}%
\pgfsetdash{{5.550000pt}{2.400000pt}}{0.000000pt}%
\pgfusepath{stroke}%
\end{pgfscope}%
\begin{pgfscope}%
\pgfsetrectcap%
\pgfsetmiterjoin%
\pgfsetlinewidth{0.803000pt}%
\definecolor{currentstroke}{rgb}{0.000000,0.000000,0.000000}%
\pgfsetstrokecolor{currentstroke}%
\pgfsetdash{}{0pt}%
\pgfpathmoveto{\pgfqpoint{0.294194in}{5.519538in}}%
\pgfpathlineto{\pgfqpoint{0.294194in}{7.760000in}}%
\pgfusepath{stroke}%
\end{pgfscope}%
\begin{pgfscope}%
\pgfsetrectcap%
\pgfsetmiterjoin%
\pgfsetlinewidth{0.803000pt}%
\definecolor{currentstroke}{rgb}{0.000000,0.000000,0.000000}%
\pgfsetstrokecolor{currentstroke}%
\pgfsetdash{}{0pt}%
\pgfpathmoveto{\pgfqpoint{8.850354in}{5.519538in}}%
\pgfpathlineto{\pgfqpoint{8.850354in}{7.760000in}}%
\pgfusepath{stroke}%
\end{pgfscope}%
\begin{pgfscope}%
\pgfsetrectcap%
\pgfsetmiterjoin%
\pgfsetlinewidth{0.803000pt}%
\definecolor{currentstroke}{rgb}{0.000000,0.000000,0.000000}%
\pgfsetstrokecolor{currentstroke}%
\pgfsetdash{}{0pt}%
\pgfpathmoveto{\pgfqpoint{0.294194in}{5.519538in}}%
\pgfpathlineto{\pgfqpoint{8.850354in}{5.519538in}}%
\pgfusepath{stroke}%
\end{pgfscope}%
\begin{pgfscope}%
\pgfsetrectcap%
\pgfsetmiterjoin%
\pgfsetlinewidth{0.803000pt}%
\definecolor{currentstroke}{rgb}{0.000000,0.000000,0.000000}%
\pgfsetstrokecolor{currentstroke}%
\pgfsetdash{}{0pt}%
\pgfpathmoveto{\pgfqpoint{0.294194in}{7.760000in}}%
\pgfpathlineto{\pgfqpoint{8.850354in}{7.760000in}}%
\pgfusepath{stroke}%
\end{pgfscope}%
\begin{pgfscope}%
\definecolor{textcolor}{rgb}{0.000000,0.000000,0.000000}%
\pgfsetstrokecolor{textcolor}%
\pgfsetfillcolor{textcolor}%
\pgftext[x=9.278162in,y=6.639769in,left,base]{\color{textcolor}\rmfamily\fontsize{16.000000}{19.200000}\selectfont \(\displaystyle t<t_c\)}%
\end{pgfscope}%
\begin{pgfscope}%
\pgfsetbuttcap%
\pgfsetmiterjoin%
\definecolor{currentfill}{rgb}{1.000000,1.000000,1.000000}%
\pgfsetfillcolor{currentfill}%
\pgfsetlinewidth{0.000000pt}%
\definecolor{currentstroke}{rgb}{0.000000,0.000000,0.000000}%
\pgfsetstrokecolor{currentstroke}%
\pgfsetstrokeopacity{0.000000}%
\pgfsetdash{}{0pt}%
\pgfpathmoveto{\pgfqpoint{0.294194in}{3.039076in}}%
\pgfpathlineto{\pgfqpoint{8.850354in}{3.039076in}}%
\pgfpathlineto{\pgfqpoint{8.850354in}{5.279538in}}%
\pgfpathlineto{\pgfqpoint{0.294194in}{5.279538in}}%
\pgfpathlineto{\pgfqpoint{0.294194in}{3.039076in}}%
\pgfpathclose%
\pgfusepath{fill}%
\end{pgfscope}%
\begin{pgfscope}%
\pgfpathrectangle{\pgfqpoint{0.294194in}{3.039076in}}{\pgfqpoint{8.556160in}{2.240462in}}%
\pgfusepath{clip}%
\pgfsetrectcap%
\pgfsetroundjoin%
\pgfsetlinewidth{1.505625pt}%
\definecolor{currentstroke}{rgb}{0.121569,0.466667,0.705882}%
\pgfsetstrokecolor{currentstroke}%
\pgfsetdash{}{0pt}%
\pgfpathmoveto{\pgfqpoint{0.294194in}{5.279538in}}%
\pgfpathlineto{\pgfqpoint{0.294537in}{5.259698in}}%
\pgfpathlineto{\pgfqpoint{0.295478in}{5.011253in}}%
\pgfpathlineto{\pgfqpoint{0.299927in}{3.039164in}}%
\pgfpathlineto{\pgfqpoint{0.300868in}{3.192541in}}%
\pgfpathlineto{\pgfqpoint{0.302751in}{4.151068in}}%
\pgfpathlineto{\pgfqpoint{0.305660in}{5.279348in}}%
\pgfpathlineto{\pgfqpoint{0.306088in}{5.243928in}}%
\pgfpathlineto{\pgfqpoint{0.307200in}{4.887961in}}%
\pgfpathlineto{\pgfqpoint{0.311392in}{3.039255in}}%
\pgfpathlineto{\pgfqpoint{0.311991in}{3.105121in}}%
\pgfpathlineto{\pgfqpoint{0.313360in}{3.645137in}}%
\pgfpathlineto{\pgfqpoint{0.317125in}{5.279465in}}%
\pgfpathlineto{\pgfqpoint{0.317382in}{5.266944in}}%
\pgfpathlineto{\pgfqpoint{0.318237in}{5.072888in}}%
\pgfpathlineto{\pgfqpoint{0.320462in}{3.871058in}}%
\pgfpathlineto{\pgfqpoint{0.322858in}{3.039118in}}%
\pgfpathlineto{\pgfqpoint{0.322943in}{3.040107in}}%
\pgfpathlineto{\pgfqpoint{0.323457in}{3.096678in}}%
\pgfpathlineto{\pgfqpoint{0.324740in}{3.574078in}}%
\pgfpathlineto{\pgfqpoint{0.328676in}{5.279238in}}%
\pgfpathlineto{\pgfqpoint{0.329018in}{5.255477in}}%
\pgfpathlineto{\pgfqpoint{0.330045in}{4.966924in}}%
\pgfpathlineto{\pgfqpoint{0.334409in}{3.039159in}}%
\pgfpathlineto{\pgfqpoint{0.335264in}{3.154585in}}%
\pgfpathlineto{\pgfqpoint{0.336890in}{3.902622in}}%
\pgfpathlineto{\pgfqpoint{0.340227in}{5.279438in}}%
\pgfpathlineto{\pgfqpoint{0.340313in}{5.277930in}}%
\pgfpathlineto{\pgfqpoint{0.340911in}{5.201002in}}%
\pgfpathlineto{\pgfqpoint{0.342366in}{4.602875in}}%
\pgfpathlineto{\pgfqpoint{0.346045in}{3.039248in}}%
\pgfpathlineto{\pgfqpoint{0.346216in}{3.045190in}}%
\pgfpathlineto{\pgfqpoint{0.346986in}{3.187021in}}%
\pgfpathlineto{\pgfqpoint{0.348869in}{4.120215in}}%
\pgfpathlineto{\pgfqpoint{0.351863in}{5.279361in}}%
\pgfpathlineto{\pgfqpoint{0.352206in}{5.258371in}}%
\pgfpathlineto{\pgfqpoint{0.353232in}{4.980800in}}%
\pgfpathlineto{\pgfqpoint{0.357682in}{3.039229in}}%
\pgfpathlineto{\pgfqpoint{0.358623in}{3.180363in}}%
\pgfpathlineto{\pgfqpoint{0.360420in}{4.050279in}}%
\pgfpathlineto{\pgfqpoint{0.363500in}{5.279247in}}%
\pgfpathlineto{\pgfqpoint{0.363671in}{5.275966in}}%
\pgfpathlineto{\pgfqpoint{0.364356in}{5.170196in}}%
\pgfpathlineto{\pgfqpoint{0.365981in}{4.440320in}}%
\pgfpathlineto{\pgfqpoint{0.369404in}{3.039342in}}%
\pgfpathlineto{\pgfqpoint{0.369746in}{3.060085in}}%
\pgfpathlineto{\pgfqpoint{0.370773in}{3.334250in}}%
\pgfpathlineto{\pgfqpoint{0.375222in}{5.279081in}}%
\pgfpathlineto{\pgfqpoint{0.376163in}{5.150889in}}%
\pgfpathlineto{\pgfqpoint{0.377874in}{4.352577in}}%
\pgfpathlineto{\pgfqpoint{0.381126in}{3.039413in}}%
\pgfpathlineto{\pgfqpoint{0.381382in}{3.048421in}}%
\pgfpathlineto{\pgfqpoint{0.382153in}{3.197434in}}%
\pgfpathlineto{\pgfqpoint{0.384035in}{4.125555in}}%
\pgfpathlineto{\pgfqpoint{0.387030in}{5.279186in}}%
\pgfpathlineto{\pgfqpoint{0.387286in}{5.269811in}}%
\pgfpathlineto{\pgfqpoint{0.388142in}{5.093150in}}%
\pgfpathlineto{\pgfqpoint{0.390281in}{3.990909in}}%
\pgfpathlineto{\pgfqpoint{0.392933in}{3.039515in}}%
\pgfpathlineto{\pgfqpoint{0.393190in}{3.048184in}}%
\pgfpathlineto{\pgfqpoint{0.393960in}{3.195162in}}%
\pgfpathlineto{\pgfqpoint{0.395843in}{4.115940in}}%
\pgfpathlineto{\pgfqpoint{0.398837in}{5.278838in}}%
\pgfpathlineto{\pgfqpoint{0.398923in}{5.278836in}}%
\pgfpathlineto{\pgfqpoint{0.399351in}{5.244522in}}%
\pgfpathlineto{\pgfqpoint{0.400463in}{4.908340in}}%
\pgfpathlineto{\pgfqpoint{0.404827in}{3.039548in}}%
\pgfpathlineto{\pgfqpoint{0.405511in}{3.113624in}}%
\pgfpathlineto{\pgfqpoint{0.406966in}{3.685559in}}%
\pgfpathlineto{\pgfqpoint{0.410730in}{5.278768in}}%
\pgfpathlineto{\pgfqpoint{0.410987in}{5.271770in}}%
\pgfpathlineto{\pgfqpoint{0.411757in}{5.131144in}}%
\pgfpathlineto{\pgfqpoint{0.413640in}{4.225769in}}%
\pgfpathlineto{\pgfqpoint{0.416720in}{3.039689in}}%
\pgfpathlineto{\pgfqpoint{0.416976in}{3.048516in}}%
\pgfpathlineto{\pgfqpoint{0.417747in}{3.193774in}}%
\pgfpathlineto{\pgfqpoint{0.419629in}{4.102758in}}%
\pgfpathlineto{\pgfqpoint{0.422709in}{5.278906in}}%
\pgfpathlineto{\pgfqpoint{0.422966in}{5.269350in}}%
\pgfpathlineto{\pgfqpoint{0.423821in}{5.096002in}}%
\pgfpathlineto{\pgfqpoint{0.425875in}{4.063548in}}%
\pgfpathlineto{\pgfqpoint{0.428699in}{3.039779in}}%
\pgfpathlineto{\pgfqpoint{0.428955in}{3.048973in}}%
\pgfpathlineto{\pgfqpoint{0.429811in}{3.220549in}}%
\pgfpathlineto{\pgfqpoint{0.431864in}{4.248435in}}%
\pgfpathlineto{\pgfqpoint{0.434688in}{5.278661in}}%
\pgfpathlineto{\pgfqpoint{0.434859in}{5.275727in}}%
\pgfpathlineto{\pgfqpoint{0.435544in}{5.176174in}}%
\pgfpathlineto{\pgfqpoint{0.437169in}{4.480119in}}%
\pgfpathlineto{\pgfqpoint{0.440763in}{3.039976in}}%
\pgfpathlineto{\pgfqpoint{0.441191in}{3.070802in}}%
\pgfpathlineto{\pgfqpoint{0.442303in}{3.391534in}}%
\pgfpathlineto{\pgfqpoint{0.446752in}{5.278588in}}%
\pgfpathlineto{\pgfqpoint{0.447522in}{5.195030in}}%
\pgfpathlineto{\pgfqpoint{0.448977in}{4.621838in}}%
\pgfpathlineto{\pgfqpoint{0.452827in}{3.040028in}}%
\pgfpathlineto{\pgfqpoint{0.452998in}{3.044621in}}%
\pgfpathlineto{\pgfqpoint{0.453683in}{3.149330in}}%
\pgfpathlineto{\pgfqpoint{0.455308in}{3.847303in}}%
\pgfpathlineto{\pgfqpoint{0.458902in}{5.278482in}}%
\pgfpathlineto{\pgfqpoint{0.459244in}{5.259475in}}%
\pgfpathlineto{\pgfqpoint{0.460186in}{5.035454in}}%
\pgfpathlineto{\pgfqpoint{0.462838in}{3.650513in}}%
\pgfpathlineto{\pgfqpoint{0.464977in}{3.040196in}}%
\pgfpathlineto{\pgfqpoint{0.465063in}{3.041580in}}%
\pgfpathlineto{\pgfqpoint{0.465662in}{3.111617in}}%
\pgfpathlineto{\pgfqpoint{0.467116in}{3.661400in}}%
\pgfpathlineto{\pgfqpoint{0.471052in}{5.278352in}}%
\pgfpathlineto{\pgfqpoint{0.471223in}{5.274328in}}%
\pgfpathlineto{\pgfqpoint{0.471908in}{5.173029in}}%
\pgfpathlineto{\pgfqpoint{0.473533in}{4.486526in}}%
\pgfpathlineto{\pgfqpoint{0.477127in}{3.040560in}}%
\pgfpathlineto{\pgfqpoint{0.477212in}{3.040665in}}%
\pgfpathlineto{\pgfqpoint{0.477640in}{3.073458in}}%
\pgfpathlineto{\pgfqpoint{0.478753in}{3.391970in}}%
\pgfpathlineto{\pgfqpoint{0.483287in}{5.278189in}}%
\pgfpathlineto{\pgfqpoint{0.484057in}{5.192202in}}%
\pgfpathlineto{\pgfqpoint{0.485598in}{4.581888in}}%
\pgfpathlineto{\pgfqpoint{0.489448in}{3.040597in}}%
\pgfpathlineto{\pgfqpoint{0.489533in}{3.042260in}}%
\pgfpathlineto{\pgfqpoint{0.490132in}{3.113047in}}%
\pgfpathlineto{\pgfqpoint{0.491587in}{3.657185in}}%
\pgfpathlineto{\pgfqpoint{0.495608in}{5.277853in}}%
\pgfpathlineto{\pgfqpoint{0.495780in}{5.271975in}}%
\pgfpathlineto{\pgfqpoint{0.496550in}{5.142882in}}%
\pgfpathlineto{\pgfqpoint{0.498346in}{4.335767in}}%
\pgfpathlineto{\pgfqpoint{0.501683in}{3.041208in}}%
\pgfpathlineto{\pgfqpoint{0.501769in}{3.040803in}}%
\pgfpathlineto{\pgfqpoint{0.501854in}{3.042518in}}%
\pgfpathlineto{\pgfqpoint{0.502453in}{3.113079in}}%
\pgfpathlineto{\pgfqpoint{0.503908in}{3.653297in}}%
\pgfpathlineto{\pgfqpoint{0.507929in}{5.277810in}}%
\pgfpathlineto{\pgfqpoint{0.508100in}{5.273337in}}%
\pgfpathlineto{\pgfqpoint{0.508785in}{5.172764in}}%
\pgfpathlineto{\pgfqpoint{0.510411in}{4.499792in}}%
\pgfpathlineto{\pgfqpoint{0.514090in}{3.041041in}}%
\pgfpathlineto{\pgfqpoint{0.514175in}{3.041324in}}%
\pgfpathlineto{\pgfqpoint{0.514603in}{3.074056in}}%
\pgfpathlineto{\pgfqpoint{0.515716in}{3.385792in}}%
\pgfpathlineto{\pgfqpoint{0.520336in}{5.277607in}}%
\pgfpathlineto{\pgfqpoint{0.521106in}{5.193560in}}%
\pgfpathlineto{\pgfqpoint{0.522646in}{4.598116in}}%
\pgfpathlineto{\pgfqpoint{0.526582in}{3.041176in}}%
\pgfpathlineto{\pgfqpoint{0.526668in}{3.042722in}}%
\pgfpathlineto{\pgfqpoint{0.527266in}{3.110974in}}%
\pgfpathlineto{\pgfqpoint{0.528721in}{3.639947in}}%
\pgfpathlineto{\pgfqpoint{0.532828in}{5.277310in}}%
\pgfpathlineto{\pgfqpoint{0.532999in}{5.272034in}}%
\pgfpathlineto{\pgfqpoint{0.533769in}{5.148534in}}%
\pgfpathlineto{\pgfqpoint{0.535566in}{4.363478in}}%
\pgfpathlineto{\pgfqpoint{0.538989in}{3.042123in}}%
\pgfpathlineto{\pgfqpoint{0.539074in}{3.041355in}}%
\pgfpathlineto{\pgfqpoint{0.539245in}{3.045987in}}%
\pgfpathlineto{\pgfqpoint{0.539930in}{3.145123in}}%
\pgfpathlineto{\pgfqpoint{0.541555in}{3.805405in}}%
\pgfpathlineto{\pgfqpoint{0.545320in}{5.277114in}}%
\pgfpathlineto{\pgfqpoint{0.545406in}{5.276516in}}%
\pgfpathlineto{\pgfqpoint{0.545919in}{5.230300in}}%
\pgfpathlineto{\pgfqpoint{0.547203in}{4.825788in}}%
\pgfpathlineto{\pgfqpoint{0.551652in}{3.041652in}}%
\pgfpathlineto{\pgfqpoint{0.552080in}{3.069739in}}%
\pgfpathlineto{\pgfqpoint{0.553192in}{3.364320in}}%
\pgfpathlineto{\pgfqpoint{0.557898in}{5.276721in}}%
\pgfpathlineto{\pgfqpoint{0.558753in}{5.185596in}}%
\pgfpathlineto{\pgfqpoint{0.560294in}{4.592492in}}%
\pgfpathlineto{\pgfqpoint{0.564229in}{3.041939in}}%
\pgfpathlineto{\pgfqpoint{0.564401in}{3.044628in}}%
\pgfpathlineto{\pgfqpoint{0.565085in}{3.134630in}}%
\pgfpathlineto{\pgfqpoint{0.566625in}{3.727617in}}%
\pgfpathlineto{\pgfqpoint{0.570561in}{5.276481in}}%
\pgfpathlineto{\pgfqpoint{0.570732in}{5.274150in}}%
\pgfpathlineto{\pgfqpoint{0.571331in}{5.203693in}}%
\pgfpathlineto{\pgfqpoint{0.572786in}{4.683095in}}%
\pgfpathlineto{\pgfqpoint{0.576978in}{3.042184in}}%
\pgfpathlineto{\pgfqpoint{0.577149in}{3.047380in}}%
\pgfpathlineto{\pgfqpoint{0.577919in}{3.167006in}}%
\pgfpathlineto{\pgfqpoint{0.579716in}{3.929648in}}%
\pgfpathlineto{\pgfqpoint{0.583310in}{5.276357in}}%
\pgfpathlineto{\pgfqpoint{0.583567in}{5.268433in}}%
\pgfpathlineto{\pgfqpoint{0.584422in}{5.117100in}}%
\pgfpathlineto{\pgfqpoint{0.586390in}{4.228587in}}%
\pgfpathlineto{\pgfqpoint{0.589642in}{3.043051in}}%
\pgfpathlineto{\pgfqpoint{0.589727in}{3.042398in}}%
\pgfpathlineto{\pgfqpoint{0.589898in}{3.047007in}}%
\pgfpathlineto{\pgfqpoint{0.590583in}{3.142772in}}%
\pgfpathlineto{\pgfqpoint{0.592208in}{3.780807in}}%
\pgfpathlineto{\pgfqpoint{0.596144in}{5.276006in}}%
\pgfpathlineto{\pgfqpoint{0.596572in}{5.248574in}}%
\pgfpathlineto{\pgfqpoint{0.597684in}{4.963324in}}%
\pgfpathlineto{\pgfqpoint{0.602561in}{3.042724in}}%
\pgfpathlineto{\pgfqpoint{0.603417in}{3.143585in}}%
\pgfpathlineto{\pgfqpoint{0.605043in}{3.777744in}}%
\pgfpathlineto{\pgfqpoint{0.608979in}{5.275801in}}%
\pgfpathlineto{\pgfqpoint{0.609321in}{5.260355in}}%
\pgfpathlineto{\pgfqpoint{0.610262in}{5.064767in}}%
\pgfpathlineto{\pgfqpoint{0.612572in}{3.960951in}}%
\pgfpathlineto{\pgfqpoint{0.615396in}{3.043174in}}%
\pgfpathlineto{\pgfqpoint{0.615481in}{3.043247in}}%
\pgfpathlineto{\pgfqpoint{0.615909in}{3.072428in}}%
\pgfpathlineto{\pgfqpoint{0.617022in}{3.358289in}}%
\pgfpathlineto{\pgfqpoint{0.621899in}{5.275474in}}%
\pgfpathlineto{\pgfqpoint{0.622754in}{5.183155in}}%
\pgfpathlineto{\pgfqpoint{0.624380in}{4.566957in}}%
\pgfpathlineto{\pgfqpoint{0.628401in}{3.043288in}}%
\pgfpathlineto{\pgfqpoint{0.628487in}{3.044276in}}%
\pgfpathlineto{\pgfqpoint{0.629000in}{3.089946in}}%
\pgfpathlineto{\pgfqpoint{0.630284in}{3.474522in}}%
\pgfpathlineto{\pgfqpoint{0.634904in}{5.275158in}}%
\pgfpathlineto{\pgfqpoint{0.635417in}{5.241731in}}%
\pgfpathlineto{\pgfqpoint{0.636615in}{4.918491in}}%
\pgfpathlineto{\pgfqpoint{0.641407in}{3.043721in}}%
\pgfpathlineto{\pgfqpoint{0.642177in}{3.113904in}}%
\pgfpathlineto{\pgfqpoint{0.643631in}{3.609645in}}%
\pgfpathlineto{\pgfqpoint{0.647995in}{5.274783in}}%
\pgfpathlineto{\pgfqpoint{0.648252in}{5.265255in}}%
\pgfpathlineto{\pgfqpoint{0.649107in}{5.115449in}}%
\pgfpathlineto{\pgfqpoint{0.651075in}{4.256330in}}%
\pgfpathlineto{\pgfqpoint{0.654412in}{3.046063in}}%
\pgfpathlineto{\pgfqpoint{0.654498in}{3.044202in}}%
\pgfpathlineto{\pgfqpoint{0.654840in}{3.055402in}}%
\pgfpathlineto{\pgfqpoint{0.655696in}{3.209761in}}%
\pgfpathlineto{\pgfqpoint{0.657749in}{4.117243in}}%
\pgfpathlineto{\pgfqpoint{0.661001in}{5.272643in}}%
\pgfpathlineto{\pgfqpoint{0.661086in}{5.274314in}}%
\pgfpathlineto{\pgfqpoint{0.661428in}{5.262463in}}%
\pgfpathlineto{\pgfqpoint{0.662370in}{5.082559in}}%
\pgfpathlineto{\pgfqpoint{0.664594in}{4.065544in}}%
\pgfpathlineto{\pgfqpoint{0.667589in}{3.046459in}}%
\pgfpathlineto{\pgfqpoint{0.667760in}{3.044551in}}%
\pgfpathlineto{\pgfqpoint{0.668017in}{3.055507in}}%
\pgfpathlineto{\pgfqpoint{0.668872in}{3.207814in}}%
\pgfpathlineto{\pgfqpoint{0.670926in}{4.106904in}}%
\pgfpathlineto{\pgfqpoint{0.674177in}{5.270956in}}%
\pgfpathlineto{\pgfqpoint{0.674348in}{5.274070in}}%
\pgfpathlineto{\pgfqpoint{0.674691in}{5.258323in}}%
\pgfpathlineto{\pgfqpoint{0.675632in}{5.070169in}}%
\pgfpathlineto{\pgfqpoint{0.677942in}{4.006145in}}%
\pgfpathlineto{\pgfqpoint{0.680851in}{3.046495in}}%
\pgfpathlineto{\pgfqpoint{0.680937in}{3.044859in}}%
\pgfpathlineto{\pgfqpoint{0.681279in}{3.056530in}}%
\pgfpathlineto{\pgfqpoint{0.682220in}{3.233507in}}%
\pgfpathlineto{\pgfqpoint{0.684445in}{4.237815in}}%
\pgfpathlineto{\pgfqpoint{0.687525in}{5.272475in}}%
\pgfpathlineto{\pgfqpoint{0.687611in}{5.273680in}}%
\pgfpathlineto{\pgfqpoint{0.687867in}{5.266421in}}%
\pgfpathlineto{\pgfqpoint{0.688637in}{5.149130in}}%
\pgfpathlineto{\pgfqpoint{0.690434in}{4.431784in}}%
\pgfpathlineto{\pgfqpoint{0.694199in}{3.046304in}}%
\pgfpathlineto{\pgfqpoint{0.694284in}{3.045130in}}%
\pgfpathlineto{\pgfqpoint{0.694541in}{3.052416in}}%
\pgfpathlineto{\pgfqpoint{0.695397in}{3.190578in}}%
\pgfpathlineto{\pgfqpoint{0.697279in}{3.970917in}}%
\pgfpathlineto{\pgfqpoint{0.700873in}{5.271643in}}%
\pgfpathlineto{\pgfqpoint{0.700958in}{5.273185in}}%
\pgfpathlineto{\pgfqpoint{0.701301in}{5.261454in}}%
\pgfpathlineto{\pgfqpoint{0.702242in}{5.086796in}}%
\pgfpathlineto{\pgfqpoint{0.704466in}{4.094676in}}%
\pgfpathlineto{\pgfqpoint{0.707547in}{3.048385in}}%
\pgfpathlineto{\pgfqpoint{0.707718in}{3.045558in}}%
\pgfpathlineto{\pgfqpoint{0.708060in}{3.061242in}}%
\pgfpathlineto{\pgfqpoint{0.709001in}{3.245138in}}%
\pgfpathlineto{\pgfqpoint{0.711311in}{4.288166in}}%
\pgfpathlineto{\pgfqpoint{0.714306in}{5.271015in}}%
\pgfpathlineto{\pgfqpoint{0.714392in}{5.272702in}}%
\pgfpathlineto{\pgfqpoint{0.714734in}{5.261759in}}%
\pgfpathlineto{\pgfqpoint{0.715675in}{5.090761in}}%
\pgfpathlineto{\pgfqpoint{0.717814in}{4.153487in}}%
\pgfpathlineto{\pgfqpoint{0.720980in}{3.050759in}}%
\pgfpathlineto{\pgfqpoint{0.721151in}{3.046047in}}%
\pgfpathlineto{\pgfqpoint{0.721579in}{3.065013in}}%
\pgfpathlineto{\pgfqpoint{0.722606in}{3.281359in}}%
\pgfpathlineto{\pgfqpoint{0.725172in}{4.470267in}}%
\pgfpathlineto{\pgfqpoint{0.727910in}{5.272268in}}%
\pgfpathlineto{\pgfqpoint{0.727996in}{5.272190in}}%
\pgfpathlineto{\pgfqpoint{0.728424in}{5.245658in}}%
\pgfpathlineto{\pgfqpoint{0.729536in}{4.985022in}}%
\pgfpathlineto{\pgfqpoint{0.734755in}{3.046422in}}%
\pgfpathlineto{\pgfqpoint{0.735697in}{3.153629in}}%
\pgfpathlineto{\pgfqpoint{0.737408in}{3.790109in}}%
\pgfpathlineto{\pgfqpoint{0.741515in}{5.271872in}}%
\pgfpathlineto{\pgfqpoint{0.741686in}{5.269740in}}%
\pgfpathlineto{\pgfqpoint{0.742371in}{5.193030in}}%
\pgfpathlineto{\pgfqpoint{0.743911in}{4.675846in}}%
\pgfpathlineto{\pgfqpoint{0.748360in}{3.046881in}}%
\pgfpathlineto{\pgfqpoint{0.748531in}{3.049632in}}%
\pgfpathlineto{\pgfqpoint{0.749215in}{3.128331in}}%
\pgfpathlineto{\pgfqpoint{0.750756in}{3.646822in}}%
\pgfpathlineto{\pgfqpoint{0.755205in}{5.271478in}}%
\pgfpathlineto{\pgfqpoint{0.755376in}{5.268926in}}%
\pgfpathlineto{\pgfqpoint{0.756060in}{5.191413in}}%
\pgfpathlineto{\pgfqpoint{0.757601in}{4.676948in}}%
\pgfpathlineto{\pgfqpoint{0.762135in}{3.047497in}}%
\pgfpathlineto{\pgfqpoint{0.762221in}{3.049117in}}%
\pgfpathlineto{\pgfqpoint{0.762820in}{3.107412in}}%
\pgfpathlineto{\pgfqpoint{0.764274in}{3.552864in}}%
\pgfpathlineto{\pgfqpoint{0.768980in}{5.271043in}}%
\pgfpathlineto{\pgfqpoint{0.769323in}{5.258114in}}%
\pgfpathlineto{\pgfqpoint{0.770264in}{5.088636in}}%
\pgfpathlineto{\pgfqpoint{0.772403in}{4.178321in}}%
\pgfpathlineto{\pgfqpoint{0.775740in}{3.050705in}}%
\pgfpathlineto{\pgfqpoint{0.775911in}{3.047831in}}%
\pgfpathlineto{\pgfqpoint{0.776253in}{3.062163in}}%
\pgfpathlineto{\pgfqpoint{0.777194in}{3.234398in}}%
\pgfpathlineto{\pgfqpoint{0.779419in}{4.187737in}}%
\pgfpathlineto{\pgfqpoint{0.782670in}{5.267602in}}%
\pgfpathlineto{\pgfqpoint{0.782842in}{5.270535in}}%
\pgfpathlineto{\pgfqpoint{0.783184in}{5.256449in}}%
\pgfpathlineto{\pgfqpoint{0.784125in}{5.085662in}}%
\pgfpathlineto{\pgfqpoint{0.786350in}{4.137284in}}%
\pgfpathlineto{\pgfqpoint{0.789601in}{3.052154in}}%
\pgfpathlineto{\pgfqpoint{0.789772in}{3.048343in}}%
\pgfpathlineto{\pgfqpoint{0.790200in}{3.067713in}}%
\pgfpathlineto{\pgfqpoint{0.791227in}{3.274910in}}%
\pgfpathlineto{\pgfqpoint{0.793708in}{4.376916in}}%
\pgfpathlineto{\pgfqpoint{0.796617in}{5.267728in}}%
\pgfpathlineto{\pgfqpoint{0.796788in}{5.269941in}}%
\pgfpathlineto{\pgfqpoint{0.797130in}{5.254673in}}%
\pgfpathlineto{\pgfqpoint{0.798157in}{5.058095in}}%
\pgfpathlineto{\pgfqpoint{0.800553in}{4.011703in}}%
\pgfpathlineto{\pgfqpoint{0.803548in}{3.053737in}}%
\pgfpathlineto{\pgfqpoint{0.803719in}{3.049014in}}%
\pgfpathlineto{\pgfqpoint{0.804147in}{3.065746in}}%
\pgfpathlineto{\pgfqpoint{0.805173in}{3.265159in}}%
\pgfpathlineto{\pgfqpoint{0.807569in}{4.310627in}}%
\pgfpathlineto{\pgfqpoint{0.810564in}{5.264528in}}%
\pgfpathlineto{\pgfqpoint{0.810735in}{5.269310in}}%
\pgfpathlineto{\pgfqpoint{0.811163in}{5.252909in}}%
\pgfpathlineto{\pgfqpoint{0.812189in}{5.055227in}}%
\pgfpathlineto{\pgfqpoint{0.814585in}{4.015032in}}%
\pgfpathlineto{\pgfqpoint{0.817580in}{3.055526in}}%
\pgfpathlineto{\pgfqpoint{0.817837in}{3.049443in}}%
\pgfpathlineto{\pgfqpoint{0.818264in}{3.071476in}}%
\pgfpathlineto{\pgfqpoint{0.819377in}{3.306924in}}%
\pgfpathlineto{\pgfqpoint{0.822200in}{4.571502in}}%
\pgfpathlineto{\pgfqpoint{0.824767in}{5.267582in}}%
\pgfpathlineto{\pgfqpoint{0.824853in}{5.268866in}}%
\pgfpathlineto{\pgfqpoint{0.825195in}{5.257998in}}%
\pgfpathlineto{\pgfqpoint{0.826136in}{5.100398in}}%
\pgfpathlineto{\pgfqpoint{0.828275in}{4.229679in}}%
\pgfpathlineto{\pgfqpoint{0.831698in}{3.057628in}}%
\pgfpathlineto{\pgfqpoint{0.831954in}{3.049969in}}%
\pgfpathlineto{\pgfqpoint{0.832382in}{3.068966in}}%
\pgfpathlineto{\pgfqpoint{0.833409in}{3.269328in}}%
\pgfpathlineto{\pgfqpoint{0.835890in}{4.342654in}}%
\pgfpathlineto{\pgfqpoint{0.838885in}{5.264632in}}%
\pgfpathlineto{\pgfqpoint{0.839056in}{5.268342in}}%
\pgfpathlineto{\pgfqpoint{0.839484in}{5.250004in}}%
\pgfpathlineto{\pgfqpoint{0.840511in}{5.052092in}}%
\pgfpathlineto{\pgfqpoint{0.842992in}{3.984870in}}%
\pgfpathlineto{\pgfqpoint{0.845987in}{3.055477in}}%
\pgfpathlineto{\pgfqpoint{0.846243in}{3.050732in}}%
\pgfpathlineto{\pgfqpoint{0.846586in}{3.066350in}}%
\pgfpathlineto{\pgfqpoint{0.847612in}{3.257251in}}%
\pgfpathlineto{\pgfqpoint{0.850008in}{4.272521in}}%
\pgfpathlineto{\pgfqpoint{0.853088in}{5.260236in}}%
\pgfpathlineto{\pgfqpoint{0.853345in}{5.267765in}}%
\pgfpathlineto{\pgfqpoint{0.853773in}{5.249178in}}%
\pgfpathlineto{\pgfqpoint{0.854800in}{5.052727in}}%
\pgfpathlineto{\pgfqpoint{0.857281in}{3.994508in}}%
\pgfpathlineto{\pgfqpoint{0.860276in}{3.058033in}}%
\pgfpathlineto{\pgfqpoint{0.860532in}{3.051143in}}%
\pgfpathlineto{\pgfqpoint{0.860960in}{3.070581in}}%
\pgfpathlineto{\pgfqpoint{0.861987in}{3.267893in}}%
\pgfpathlineto{\pgfqpoint{0.864468in}{4.323130in}}%
\pgfpathlineto{\pgfqpoint{0.867548in}{5.263674in}}%
\pgfpathlineto{\pgfqpoint{0.867720in}{5.267151in}}%
\pgfpathlineto{\pgfqpoint{0.868147in}{5.248967in}}%
\pgfpathlineto{\pgfqpoint{0.869174in}{5.055410in}}%
\pgfpathlineto{\pgfqpoint{0.871570in}{4.048585in}}%
\pgfpathlineto{\pgfqpoint{0.874736in}{3.056740in}}%
\pgfpathlineto{\pgfqpoint{0.874992in}{3.051915in}}%
\pgfpathlineto{\pgfqpoint{0.875335in}{3.066842in}}%
\pgfpathlineto{\pgfqpoint{0.876361in}{3.252038in}}%
\pgfpathlineto{\pgfqpoint{0.878671in}{4.205669in}}%
\pgfpathlineto{\pgfqpoint{0.881923in}{5.258506in}}%
\pgfpathlineto{\pgfqpoint{0.882180in}{5.266485in}}%
\pgfpathlineto{\pgfqpoint{0.882607in}{5.249483in}}%
\pgfpathlineto{\pgfqpoint{0.883634in}{5.060536in}}%
\pgfpathlineto{\pgfqpoint{0.886030in}{4.065798in}}%
\pgfpathlineto{\pgfqpoint{0.889196in}{3.060289in}}%
\pgfpathlineto{\pgfqpoint{0.889452in}{3.052444in}}%
\pgfpathlineto{\pgfqpoint{0.889880in}{3.069462in}}%
\pgfpathlineto{\pgfqpoint{0.890907in}{3.257440in}}%
\pgfpathlineto{\pgfqpoint{0.893303in}{4.247446in}}%
\pgfpathlineto{\pgfqpoint{0.896468in}{5.256714in}}%
\pgfpathlineto{\pgfqpoint{0.896811in}{5.265699in}}%
\pgfpathlineto{\pgfqpoint{0.897239in}{5.243340in}}%
\pgfpathlineto{\pgfqpoint{0.898351in}{5.019653in}}%
\pgfpathlineto{\pgfqpoint{0.901089in}{3.848529in}}%
\pgfpathlineto{\pgfqpoint{0.903912in}{3.056199in}}%
\pgfpathlineto{\pgfqpoint{0.904084in}{3.053064in}}%
\pgfpathlineto{\pgfqpoint{0.904426in}{3.064610in}}%
\pgfpathlineto{\pgfqpoint{0.905367in}{3.214888in}}%
\pgfpathlineto{\pgfqpoint{0.907506in}{4.040445in}}%
\pgfpathlineto{\pgfqpoint{0.911185in}{5.259028in}}%
\pgfpathlineto{\pgfqpoint{0.911442in}{5.265207in}}%
\pgfpathlineto{\pgfqpoint{0.911870in}{5.246049in}}%
\pgfpathlineto{\pgfqpoint{0.912896in}{5.056371in}}%
\pgfpathlineto{\pgfqpoint{0.915292in}{4.075872in}}%
\pgfpathlineto{\pgfqpoint{0.918544in}{3.060464in}}%
\pgfpathlineto{\pgfqpoint{0.918800in}{3.053742in}}%
\pgfpathlineto{\pgfqpoint{0.919228in}{3.071792in}}%
\pgfpathlineto{\pgfqpoint{0.920255in}{3.258026in}}%
\pgfpathlineto{\pgfqpoint{0.922651in}{4.231184in}}%
\pgfpathlineto{\pgfqpoint{0.925902in}{5.255833in}}%
\pgfpathlineto{\pgfqpoint{0.926159in}{5.264372in}}%
\pgfpathlineto{\pgfqpoint{0.926672in}{5.242258in}}%
\pgfpathlineto{\pgfqpoint{0.927784in}{5.023842in}}%
\pgfpathlineto{\pgfqpoint{0.930437in}{3.912719in}}%
\pgfpathlineto{\pgfqpoint{0.933346in}{3.061760in}}%
\pgfpathlineto{\pgfqpoint{0.933603in}{3.054467in}}%
\pgfpathlineto{\pgfqpoint{0.934030in}{3.071151in}}%
\pgfpathlineto{\pgfqpoint{0.935057in}{3.252364in}}%
\pgfpathlineto{\pgfqpoint{0.937367in}{4.172764in}}%
\pgfpathlineto{\pgfqpoint{0.940790in}{5.256428in}}%
\pgfpathlineto{\pgfqpoint{0.941047in}{5.263782in}}%
\pgfpathlineto{\pgfqpoint{0.941474in}{5.247407in}}%
\pgfpathlineto{\pgfqpoint{0.942501in}{5.067882in}}%
\pgfpathlineto{\pgfqpoint{0.944811in}{4.152968in}}%
\pgfpathlineto{\pgfqpoint{0.948234in}{3.064079in}}%
\pgfpathlineto{\pgfqpoint{0.948576in}{3.055315in}}%
\pgfpathlineto{\pgfqpoint{0.949004in}{3.076311in}}%
\pgfpathlineto{\pgfqpoint{0.950116in}{3.288830in}}%
\pgfpathlineto{\pgfqpoint{0.952769in}{4.382386in}}%
\pgfpathlineto{\pgfqpoint{0.955763in}{5.255954in}}%
\pgfpathlineto{\pgfqpoint{0.956020in}{5.263059in}}%
\pgfpathlineto{\pgfqpoint{0.956448in}{5.246685in}}%
\pgfpathlineto{\pgfqpoint{0.957475in}{5.069170in}}%
\pgfpathlineto{\pgfqpoint{0.959785in}{4.163804in}}%
\pgfpathlineto{\pgfqpoint{0.963207in}{3.067813in}}%
\pgfpathlineto{\pgfqpoint{0.963549in}{3.055918in}}%
\pgfpathlineto{\pgfqpoint{0.964063in}{3.080044in}}%
\pgfpathlineto{\pgfqpoint{0.965175in}{3.296403in}}%
\pgfpathlineto{\pgfqpoint{0.967828in}{4.384812in}}%
\pgfpathlineto{\pgfqpoint{0.970822in}{5.254374in}}%
\pgfpathlineto{\pgfqpoint{0.971079in}{5.262225in}}%
\pgfpathlineto{\pgfqpoint{0.971592in}{5.240423in}}%
\pgfpathlineto{\pgfqpoint{0.972705in}{5.029808in}}%
\pgfpathlineto{\pgfqpoint{0.975357in}{3.950623in}}%
\pgfpathlineto{\pgfqpoint{0.978437in}{3.062705in}}%
\pgfpathlineto{\pgfqpoint{0.978694in}{3.056675in}}%
\pgfpathlineto{\pgfqpoint{0.979122in}{3.074209in}}%
\pgfpathlineto{\pgfqpoint{0.980149in}{3.251285in}}%
\pgfpathlineto{\pgfqpoint{0.982459in}{4.145081in}}%
\pgfpathlineto{\pgfqpoint{0.985967in}{5.251405in}}%
\pgfpathlineto{\pgfqpoint{0.986309in}{5.261534in}}%
\pgfpathlineto{\pgfqpoint{0.986737in}{5.243402in}}%
\pgfpathlineto{\pgfqpoint{0.987764in}{5.066007in}}%
\pgfpathlineto{\pgfqpoint{0.990074in}{4.175564in}}%
\pgfpathlineto{\pgfqpoint{0.993582in}{3.068685in}}%
\pgfpathlineto{\pgfqpoint{0.993924in}{3.057464in}}%
\pgfpathlineto{\pgfqpoint{0.994437in}{3.081354in}}%
\pgfpathlineto{\pgfqpoint{0.995550in}{3.292471in}}%
\pgfpathlineto{\pgfqpoint{0.998202in}{4.359293in}}%
\pgfpathlineto{\pgfqpoint{1.001282in}{5.251985in}}%
\pgfpathlineto{\pgfqpoint{1.001625in}{5.260645in}}%
\pgfpathlineto{\pgfqpoint{1.002053in}{5.241149in}}%
\pgfpathlineto{\pgfqpoint{1.003165in}{5.040141in}}%
\pgfpathlineto{\pgfqpoint{1.005732in}{4.024960in}}%
\pgfpathlineto{\pgfqpoint{1.008983in}{3.066351in}}%
\pgfpathlineto{\pgfqpoint{1.009240in}{3.058423in}}%
\pgfpathlineto{\pgfqpoint{1.009753in}{3.078674in}}%
\pgfpathlineto{\pgfqpoint{1.010865in}{3.280249in}}%
\pgfpathlineto{\pgfqpoint{1.013432in}{4.291777in}}%
\pgfpathlineto{\pgfqpoint{1.016684in}{5.250955in}}%
\pgfpathlineto{\pgfqpoint{1.017026in}{5.259858in}}%
\pgfpathlineto{\pgfqpoint{1.017454in}{5.241129in}}%
\pgfpathlineto{\pgfqpoint{1.018566in}{5.044218in}}%
\pgfpathlineto{\pgfqpoint{1.021047in}{4.079393in}}%
\pgfpathlineto{\pgfqpoint{1.024384in}{3.070797in}}%
\pgfpathlineto{\pgfqpoint{1.024727in}{3.059136in}}%
\pgfpathlineto{\pgfqpoint{1.025240in}{3.081130in}}%
\pgfpathlineto{\pgfqpoint{1.026352in}{3.283745in}}%
\pgfpathlineto{\pgfqpoint{1.028919in}{4.287856in}}%
\pgfpathlineto{\pgfqpoint{1.032171in}{5.248153in}}%
\pgfpathlineto{\pgfqpoint{1.032513in}{5.259077in}}%
\pgfpathlineto{\pgfqpoint{1.033026in}{5.236289in}}%
\pgfpathlineto{\pgfqpoint{1.034138in}{5.033310in}}%
\pgfpathlineto{\pgfqpoint{1.036705in}{4.033065in}}%
\pgfpathlineto{\pgfqpoint{1.039957in}{3.071997in}}%
\pgfpathlineto{\pgfqpoint{1.040299in}{3.060026in}}%
\pgfpathlineto{\pgfqpoint{1.040812in}{3.080939in}}%
\pgfpathlineto{\pgfqpoint{1.041925in}{3.279032in}}%
\pgfpathlineto{\pgfqpoint{1.044492in}{4.270341in}}%
\pgfpathlineto{\pgfqpoint{1.047828in}{5.248560in}}%
\pgfpathlineto{\pgfqpoint{1.048171in}{5.258201in}}%
\pgfpathlineto{\pgfqpoint{1.048599in}{5.241321in}}%
\pgfpathlineto{\pgfqpoint{1.049625in}{5.074694in}}%
\pgfpathlineto{\pgfqpoint{1.051935in}{4.226497in}}%
\pgfpathlineto{\pgfqpoint{1.055615in}{3.075471in}}%
\pgfpathlineto{\pgfqpoint{1.056042in}{3.060876in}}%
\pgfpathlineto{\pgfqpoint{1.056556in}{3.085412in}}%
\pgfpathlineto{\pgfqpoint{1.057754in}{3.311560in}}%
\pgfpathlineto{\pgfqpoint{1.060577in}{4.422826in}}%
\pgfpathlineto{\pgfqpoint{1.063572in}{5.246896in}}%
\pgfpathlineto{\pgfqpoint{1.063914in}{5.257315in}}%
\pgfpathlineto{\pgfqpoint{1.064428in}{5.235007in}}%
\pgfpathlineto{\pgfqpoint{1.065540in}{5.037746in}}%
\pgfpathlineto{\pgfqpoint{1.068107in}{4.060013in}}%
\pgfpathlineto{\pgfqpoint{1.071444in}{3.075523in}}%
\pgfpathlineto{\pgfqpoint{1.071872in}{3.061817in}}%
\pgfpathlineto{\pgfqpoint{1.072299in}{3.079475in}}%
\pgfpathlineto{\pgfqpoint{1.073412in}{3.265885in}}%
\pgfpathlineto{\pgfqpoint{1.075893in}{4.191324in}}%
\pgfpathlineto{\pgfqpoint{1.079401in}{5.242811in}}%
\pgfpathlineto{\pgfqpoint{1.079829in}{5.256330in}}%
\pgfpathlineto{\pgfqpoint{1.080257in}{5.238742in}}%
\pgfpathlineto{\pgfqpoint{1.081369in}{5.053634in}}%
\pgfpathlineto{\pgfqpoint{1.083850in}{4.133712in}}%
\pgfpathlineto{\pgfqpoint{1.087358in}{3.078272in}}%
\pgfpathlineto{\pgfqpoint{1.087786in}{3.062687in}}%
\pgfpathlineto{\pgfqpoint{1.088300in}{3.084688in}}%
\pgfpathlineto{\pgfqpoint{1.089412in}{3.277768in}}%
\pgfpathlineto{\pgfqpoint{1.091979in}{4.238461in}}%
\pgfpathlineto{\pgfqpoint{1.095401in}{5.241650in}}%
\pgfpathlineto{\pgfqpoint{1.095829in}{5.255412in}}%
\pgfpathlineto{\pgfqpoint{1.096342in}{5.231569in}}%
\pgfpathlineto{\pgfqpoint{1.097540in}{5.013570in}}%
\pgfpathlineto{\pgfqpoint{1.100278in}{3.968343in}}%
\pgfpathlineto{\pgfqpoint{1.103444in}{3.077891in}}%
\pgfpathlineto{\pgfqpoint{1.103872in}{3.063661in}}%
\pgfpathlineto{\pgfqpoint{1.104385in}{3.086606in}}%
\pgfpathlineto{\pgfqpoint{1.105583in}{3.301381in}}%
\pgfpathlineto{\pgfqpoint{1.108321in}{4.339212in}}%
\pgfpathlineto{\pgfqpoint{1.111573in}{5.243250in}}%
\pgfpathlineto{\pgfqpoint{1.111915in}{5.254418in}}%
\pgfpathlineto{\pgfqpoint{1.112428in}{5.235072in}}%
\pgfpathlineto{\pgfqpoint{1.113540in}{5.050821in}}%
\pgfpathlineto{\pgfqpoint{1.116022in}{4.148519in}}%
\pgfpathlineto{\pgfqpoint{1.119615in}{3.080604in}}%
\pgfpathlineto{\pgfqpoint{1.120043in}{3.064649in}}%
\pgfpathlineto{\pgfqpoint{1.120557in}{3.084859in}}%
\pgfpathlineto{\pgfqpoint{1.121669in}{3.269633in}}%
\pgfpathlineto{\pgfqpoint{1.124150in}{4.168296in}}%
\pgfpathlineto{\pgfqpoint{1.127744in}{5.236132in}}%
\pgfpathlineto{\pgfqpoint{1.128172in}{5.253378in}}%
\pgfpathlineto{\pgfqpoint{1.128685in}{5.235053in}}%
\pgfpathlineto{\pgfqpoint{1.129797in}{5.055213in}}%
\pgfpathlineto{\pgfqpoint{1.132193in}{4.201985in}}%
\pgfpathlineto{\pgfqpoint{1.135958in}{3.080594in}}%
\pgfpathlineto{\pgfqpoint{1.136386in}{3.065645in}}%
\pgfpathlineto{\pgfqpoint{1.136899in}{3.086387in}}%
\pgfpathlineto{\pgfqpoint{1.138011in}{3.269839in}}%
\pgfpathlineto{\pgfqpoint{1.140493in}{4.158680in}}%
\pgfpathlineto{\pgfqpoint{1.144172in}{5.237471in}}%
\pgfpathlineto{\pgfqpoint{1.144600in}{5.252450in}}%
\pgfpathlineto{\pgfqpoint{1.145113in}{5.232081in}}%
\pgfpathlineto{\pgfqpoint{1.146225in}{5.050555in}}%
\pgfpathlineto{\pgfqpoint{1.148707in}{4.168014in}}%
\pgfpathlineto{\pgfqpoint{1.152386in}{3.084134in}}%
\pgfpathlineto{\pgfqpoint{1.152899in}{3.066826in}}%
\pgfpathlineto{\pgfqpoint{1.153327in}{3.084078in}}%
\pgfpathlineto{\pgfqpoint{1.154439in}{3.258200in}}%
\pgfpathlineto{\pgfqpoint{1.156835in}{4.092950in}}%
\pgfpathlineto{\pgfqpoint{1.160685in}{5.235144in}}%
\pgfpathlineto{\pgfqpoint{1.161113in}{5.251335in}}%
\pgfpathlineto{\pgfqpoint{1.161627in}{5.233089in}}%
\pgfpathlineto{\pgfqpoint{1.162739in}{5.058162in}}%
\pgfpathlineto{\pgfqpoint{1.165135in}{4.226602in}}%
\pgfpathlineto{\pgfqpoint{1.168985in}{3.085345in}}%
\pgfpathlineto{\pgfqpoint{1.169498in}{3.067855in}}%
\pgfpathlineto{\pgfqpoint{1.169926in}{3.084397in}}%
\pgfpathlineto{\pgfqpoint{1.171038in}{3.254560in}}%
\pgfpathlineto{\pgfqpoint{1.173434in}{4.076679in}}%
\pgfpathlineto{\pgfqpoint{1.177370in}{5.234972in}}%
\pgfpathlineto{\pgfqpoint{1.177798in}{5.250286in}}%
\pgfpathlineto{\pgfqpoint{1.178311in}{5.231658in}}%
\pgfpathlineto{\pgfqpoint{1.179424in}{5.058341in}}%
\pgfpathlineto{\pgfqpoint{1.181819in}{4.236766in}}%
\pgfpathlineto{\pgfqpoint{1.185755in}{3.084481in}}%
\pgfpathlineto{\pgfqpoint{1.186183in}{3.068902in}}%
\pgfpathlineto{\pgfqpoint{1.186696in}{3.086878in}}%
\pgfpathlineto{\pgfqpoint{1.187809in}{3.257708in}}%
\pgfpathlineto{\pgfqpoint{1.190204in}{4.072368in}}%
\pgfpathlineto{\pgfqpoint{1.194140in}{5.230687in}}%
\pgfpathlineto{\pgfqpoint{1.194654in}{5.249181in}}%
\pgfpathlineto{\pgfqpoint{1.195167in}{5.228054in}}%
\pgfpathlineto{\pgfqpoint{1.196365in}{5.031804in}}%
\pgfpathlineto{\pgfqpoint{1.198932in}{4.132078in}}%
\pgfpathlineto{\pgfqpoint{1.202611in}{3.087855in}}%
\pgfpathlineto{\pgfqpoint{1.203124in}{3.070020in}}%
\pgfpathlineto{\pgfqpoint{1.203638in}{3.091441in}}%
\pgfpathlineto{\pgfqpoint{1.204836in}{3.286986in}}%
\pgfpathlineto{\pgfqpoint{1.207402in}{4.181530in}}%
\pgfpathlineto{\pgfqpoint{1.211082in}{5.228064in}}%
\pgfpathlineto{\pgfqpoint{1.211595in}{5.248057in}}%
\pgfpathlineto{\pgfqpoint{1.212108in}{5.229152in}}%
\pgfpathlineto{\pgfqpoint{1.213221in}{5.060019in}}%
\pgfpathlineto{\pgfqpoint{1.215616in}{4.259656in}}%
\pgfpathlineto{\pgfqpoint{1.219638in}{3.089658in}}%
\pgfpathlineto{\pgfqpoint{1.220151in}{3.071144in}}%
\pgfpathlineto{\pgfqpoint{1.220665in}{3.091164in}}%
\pgfpathlineto{\pgfqpoint{1.221863in}{3.281070in}}%
\pgfpathlineto{\pgfqpoint{1.224429in}{4.160927in}}%
\pgfpathlineto{\pgfqpoint{1.228194in}{5.226998in}}%
\pgfpathlineto{\pgfqpoint{1.228708in}{5.246878in}}%
\pgfpathlineto{\pgfqpoint{1.229221in}{5.228586in}}%
\pgfpathlineto{\pgfqpoint{1.230333in}{5.063037in}}%
\pgfpathlineto{\pgfqpoint{1.232729in}{4.275103in}}%
\pgfpathlineto{\pgfqpoint{1.236836in}{3.090104in}}%
\pgfpathlineto{\pgfqpoint{1.237349in}{3.072344in}}%
\pgfpathlineto{\pgfqpoint{1.237863in}{3.092393in}}%
\pgfpathlineto{\pgfqpoint{1.239061in}{3.279706in}}%
\pgfpathlineto{\pgfqpoint{1.241627in}{4.147844in}}%
\pgfpathlineto{\pgfqpoint{1.245392in}{5.220463in}}%
\pgfpathlineto{\pgfqpoint{1.245991in}{5.245680in}}%
\pgfpathlineto{\pgfqpoint{1.246504in}{5.226746in}}%
\pgfpathlineto{\pgfqpoint{1.247617in}{5.062269in}}%
\pgfpathlineto{\pgfqpoint{1.250013in}{4.283896in}}%
\pgfpathlineto{\pgfqpoint{1.254120in}{3.095812in}}%
\pgfpathlineto{\pgfqpoint{1.254718in}{3.073641in}}%
\pgfpathlineto{\pgfqpoint{1.255232in}{3.094796in}}%
\pgfpathlineto{\pgfqpoint{1.256430in}{3.281820in}}%
\pgfpathlineto{\pgfqpoint{1.258997in}{4.140575in}}%
\pgfpathlineto{\pgfqpoint{1.262847in}{5.221878in}}%
\pgfpathlineto{\pgfqpoint{1.263446in}{5.244381in}}%
\pgfpathlineto{\pgfqpoint{1.263959in}{5.223900in}}%
\pgfpathlineto{\pgfqpoint{1.265157in}{5.039636in}}%
\pgfpathlineto{\pgfqpoint{1.267724in}{4.188335in}}%
\pgfpathlineto{\pgfqpoint{1.271574in}{3.100980in}}%
\pgfpathlineto{\pgfqpoint{1.272173in}{3.074823in}}%
\pgfpathlineto{\pgfqpoint{1.272687in}{3.091782in}}%
\pgfpathlineto{\pgfqpoint{1.273799in}{3.248769in}}%
\pgfpathlineto{\pgfqpoint{1.276109in}{3.973010in}}%
\pgfpathlineto{\pgfqpoint{1.280473in}{5.223751in}}%
\pgfpathlineto{\pgfqpoint{1.280986in}{5.243145in}}%
\pgfpathlineto{\pgfqpoint{1.281500in}{5.226528in}}%
\pgfpathlineto{\pgfqpoint{1.282612in}{5.071403in}}%
\pgfpathlineto{\pgfqpoint{1.284922in}{4.353410in}}%
\pgfpathlineto{\pgfqpoint{1.289286in}{3.098613in}}%
\pgfpathlineto{\pgfqpoint{1.289885in}{3.076123in}}%
\pgfpathlineto{\pgfqpoint{1.290398in}{3.095450in}}%
\pgfpathlineto{\pgfqpoint{1.291596in}{3.273286in}}%
\pgfpathlineto{\pgfqpoint{1.294077in}{4.071306in}}%
\pgfpathlineto{\pgfqpoint{1.298184in}{5.219530in}}%
\pgfpathlineto{\pgfqpoint{1.298783in}{5.241836in}}%
\pgfpathlineto{\pgfqpoint{1.299297in}{5.222739in}}%
\pgfpathlineto{\pgfqpoint{1.300494in}{5.046713in}}%
\pgfpathlineto{\pgfqpoint{1.302976in}{4.255085in}}%
\pgfpathlineto{\pgfqpoint{1.307083in}{3.102936in}}%
\pgfpathlineto{\pgfqpoint{1.307682in}{3.077446in}}%
\pgfpathlineto{\pgfqpoint{1.308195in}{3.093423in}}%
\pgfpathlineto{\pgfqpoint{1.309307in}{3.243745in}}%
\pgfpathlineto{\pgfqpoint{1.311618in}{3.943814in}}%
\pgfpathlineto{\pgfqpoint{1.316152in}{5.221677in}}%
\pgfpathlineto{\pgfqpoint{1.316666in}{5.240492in}}%
\pgfpathlineto{\pgfqpoint{1.317179in}{5.224729in}}%
\pgfpathlineto{\pgfqpoint{1.318291in}{5.076008in}}%
\pgfpathlineto{\pgfqpoint{1.320602in}{4.381952in}}%
\pgfpathlineto{\pgfqpoint{1.325136in}{3.100643in}}%
\pgfpathlineto{\pgfqpoint{1.325735in}{3.078794in}}%
\pgfpathlineto{\pgfqpoint{1.326249in}{3.097123in}}%
\pgfpathlineto{\pgfqpoint{1.327447in}{3.267554in}}%
\pgfpathlineto{\pgfqpoint{1.329928in}{4.039668in}}%
\pgfpathlineto{\pgfqpoint{1.334120in}{5.210404in}}%
\pgfpathlineto{\pgfqpoint{1.334805in}{5.239134in}}%
\pgfpathlineto{\pgfqpoint{1.335318in}{5.221212in}}%
\pgfpathlineto{\pgfqpoint{1.336516in}{5.052970in}}%
\pgfpathlineto{\pgfqpoint{1.338998in}{4.287777in}}%
\pgfpathlineto{\pgfqpoint{1.343276in}{3.105539in}}%
\pgfpathlineto{\pgfqpoint{1.343875in}{3.080237in}}%
\pgfpathlineto{\pgfqpoint{1.344388in}{3.094832in}}%
\pgfpathlineto{\pgfqpoint{1.345500in}{3.237689in}}%
\pgfpathlineto{\pgfqpoint{1.347810in}{3.912252in}}%
\pgfpathlineto{\pgfqpoint{1.352516in}{5.218384in}}%
\pgfpathlineto{\pgfqpoint{1.353030in}{5.237628in}}%
\pgfpathlineto{\pgfqpoint{1.353543in}{5.223717in}}%
\pgfpathlineto{\pgfqpoint{1.354655in}{5.083412in}}%
\pgfpathlineto{\pgfqpoint{1.356880in}{4.446622in}}%
\pgfpathlineto{\pgfqpoint{1.361757in}{3.098659in}}%
\pgfpathlineto{\pgfqpoint{1.362271in}{3.081606in}}%
\pgfpathlineto{\pgfqpoint{1.362784in}{3.097355in}}%
\pgfpathlineto{\pgfqpoint{1.363896in}{3.240219in}}%
\pgfpathlineto{\pgfqpoint{1.366206in}{3.906682in}}%
\pgfpathlineto{\pgfqpoint{1.370912in}{5.212278in}}%
\pgfpathlineto{\pgfqpoint{1.371511in}{5.236243in}}%
\pgfpathlineto{\pgfqpoint{1.372025in}{5.221651in}}%
\pgfpathlineto{\pgfqpoint{1.373137in}{5.082290in}}%
\pgfpathlineto{\pgfqpoint{1.375362in}{4.454428in}}%
\pgfpathlineto{\pgfqpoint{1.380324in}{3.099334in}}%
\pgfpathlineto{\pgfqpoint{1.380838in}{3.083071in}}%
\pgfpathlineto{\pgfqpoint{1.381351in}{3.098904in}}%
\pgfpathlineto{\pgfqpoint{1.382463in}{3.239617in}}%
\pgfpathlineto{\pgfqpoint{1.384773in}{3.895206in}}%
\pgfpathlineto{\pgfqpoint{1.389565in}{5.211063in}}%
\pgfpathlineto{\pgfqpoint{1.390164in}{5.234742in}}%
\pgfpathlineto{\pgfqpoint{1.390677in}{5.220669in}}%
\pgfpathlineto{\pgfqpoint{1.391790in}{5.084680in}}%
\pgfpathlineto{\pgfqpoint{1.394014in}{4.468866in}}%
\pgfpathlineto{\pgfqpoint{1.399062in}{3.101401in}}%
\pgfpathlineto{\pgfqpoint{1.399576in}{3.084600in}}%
\pgfpathlineto{\pgfqpoint{1.400089in}{3.099191in}}%
\pgfpathlineto{\pgfqpoint{1.401201in}{3.235073in}}%
\pgfpathlineto{\pgfqpoint{1.403426in}{3.846873in}}%
\pgfpathlineto{\pgfqpoint{1.408474in}{5.214142in}}%
\pgfpathlineto{\pgfqpoint{1.409073in}{5.233213in}}%
\pgfpathlineto{\pgfqpoint{1.409587in}{5.215948in}}%
\pgfpathlineto{\pgfqpoint{1.410784in}{5.059572in}}%
\pgfpathlineto{\pgfqpoint{1.413180in}{4.373962in}}%
\pgfpathlineto{\pgfqpoint{1.417886in}{3.111359in}}%
\pgfpathlineto{\pgfqpoint{1.418571in}{3.086171in}}%
\pgfpathlineto{\pgfqpoint{1.419084in}{3.103053in}}%
\pgfpathlineto{\pgfqpoint{1.420282in}{3.257301in}}%
\pgfpathlineto{\pgfqpoint{1.422678in}{3.936065in}}%
\pgfpathlineto{\pgfqpoint{1.427384in}{5.202502in}}%
\pgfpathlineto{\pgfqpoint{1.428068in}{5.231645in}}%
\pgfpathlineto{\pgfqpoint{1.428581in}{5.218145in}}%
\pgfpathlineto{\pgfqpoint{1.429694in}{5.087943in}}%
\pgfpathlineto{\pgfqpoint{1.431918in}{4.495173in}}%
\pgfpathlineto{\pgfqpoint{1.437138in}{3.105688in}}%
\pgfpathlineto{\pgfqpoint{1.437737in}{3.087804in}}%
\pgfpathlineto{\pgfqpoint{1.438250in}{3.104950in}}%
\pgfpathlineto{\pgfqpoint{1.439448in}{3.257179in}}%
\pgfpathlineto{\pgfqpoint{1.441844in}{3.924869in}}%
\pgfpathlineto{\pgfqpoint{1.446635in}{5.200956in}}%
\pgfpathlineto{\pgfqpoint{1.447320in}{5.230004in}}%
\pgfpathlineto{\pgfqpoint{1.447833in}{5.217247in}}%
\pgfpathlineto{\pgfqpoint{1.448945in}{5.090857in}}%
\pgfpathlineto{\pgfqpoint{1.451170in}{4.510909in}}%
\pgfpathlineto{\pgfqpoint{1.456475in}{3.108674in}}%
\pgfpathlineto{\pgfqpoint{1.457074in}{3.089367in}}%
\pgfpathlineto{\pgfqpoint{1.457587in}{3.104540in}}%
\pgfpathlineto{\pgfqpoint{1.458699in}{3.234704in}}%
\pgfpathlineto{\pgfqpoint{1.460924in}{3.815759in}}%
\pgfpathlineto{\pgfqpoint{1.466229in}{5.208749in}}%
\pgfpathlineto{\pgfqpoint{1.466828in}{5.228418in}}%
\pgfpathlineto{\pgfqpoint{1.467341in}{5.213929in}}%
\pgfpathlineto{\pgfqpoint{1.468454in}{5.086302in}}%
\pgfpathlineto{\pgfqpoint{1.470678in}{4.512499in}}%
\pgfpathlineto{\pgfqpoint{1.476069in}{3.109098in}}%
\pgfpathlineto{\pgfqpoint{1.476668in}{3.091069in}}%
\pgfpathlineto{\pgfqpoint{1.477181in}{3.106588in}}%
\pgfpathlineto{\pgfqpoint{1.478379in}{3.250127in}}%
\pgfpathlineto{\pgfqpoint{1.480775in}{3.889809in}}%
\pgfpathlineto{\pgfqpoint{1.485823in}{5.200744in}}%
\pgfpathlineto{\pgfqpoint{1.486507in}{5.226706in}}%
\pgfpathlineto{\pgfqpoint{1.487021in}{5.213244in}}%
\pgfpathlineto{\pgfqpoint{1.488133in}{5.089991in}}%
\pgfpathlineto{\pgfqpoint{1.490357in}{4.529838in}}%
\pgfpathlineto{\pgfqpoint{1.495919in}{3.107674in}}%
\pgfpathlineto{\pgfqpoint{1.496432in}{3.092785in}}%
\pgfpathlineto{\pgfqpoint{1.496946in}{3.105807in}}%
\pgfpathlineto{\pgfqpoint{1.498058in}{3.227020in}}%
\pgfpathlineto{\pgfqpoint{1.500283in}{3.780513in}}%
\pgfpathlineto{\pgfqpoint{1.505930in}{5.211488in}}%
\pgfpathlineto{\pgfqpoint{1.506443in}{5.224962in}}%
\pgfpathlineto{\pgfqpoint{1.506957in}{5.210870in}}%
\pgfpathlineto{\pgfqpoint{1.508069in}{5.088608in}}%
\pgfpathlineto{\pgfqpoint{1.510294in}{4.537593in}}%
\pgfpathlineto{\pgfqpoint{1.515941in}{3.109770in}}%
\pgfpathlineto{\pgfqpoint{1.516454in}{3.094589in}}%
\pgfpathlineto{\pgfqpoint{1.516967in}{3.106631in}}%
\pgfpathlineto{\pgfqpoint{1.518080in}{3.223570in}}%
\pgfpathlineto{\pgfqpoint{1.520304in}{3.763405in}}%
\pgfpathlineto{\pgfqpoint{1.526037in}{5.207619in}}%
\pgfpathlineto{\pgfqpoint{1.526550in}{5.223089in}}%
\pgfpathlineto{\pgfqpoint{1.527064in}{5.211678in}}%
\pgfpathlineto{\pgfqpoint{1.528176in}{5.097167in}}%
\pgfpathlineto{\pgfqpoint{1.530315in}{4.590473in}}%
\pgfpathlineto{\pgfqpoint{1.536305in}{3.106591in}}%
\pgfpathlineto{\pgfqpoint{1.536732in}{3.096392in}}%
\pgfpathlineto{\pgfqpoint{1.537246in}{3.108489in}}%
\pgfpathlineto{\pgfqpoint{1.538358in}{3.223275in}}%
\pgfpathlineto{\pgfqpoint{1.540583in}{3.752316in}}%
\pgfpathlineto{\pgfqpoint{1.546486in}{5.209103in}}%
\pgfpathlineto{\pgfqpoint{1.547000in}{5.221302in}}%
\pgfpathlineto{\pgfqpoint{1.547513in}{5.207304in}}%
\pgfpathlineto{\pgfqpoint{1.548626in}{5.089755in}}%
\pgfpathlineto{\pgfqpoint{1.550850in}{4.560790in}}%
\pgfpathlineto{\pgfqpoint{1.556754in}{3.111380in}}%
\pgfpathlineto{\pgfqpoint{1.557267in}{3.098240in}}%
\pgfpathlineto{\pgfqpoint{1.557781in}{3.110959in}}%
\pgfpathlineto{\pgfqpoint{1.558893in}{3.224764in}}%
\pgfpathlineto{\pgfqpoint{1.561118in}{3.744625in}}%
\pgfpathlineto{\pgfqpoint{1.567107in}{5.206471in}}%
\pgfpathlineto{\pgfqpoint{1.567620in}{5.219418in}}%
\pgfpathlineto{\pgfqpoint{1.568134in}{5.206846in}}%
\pgfpathlineto{\pgfqpoint{1.569246in}{5.094464in}}%
\pgfpathlineto{\pgfqpoint{1.571471in}{4.580502in}}%
\pgfpathlineto{\pgfqpoint{1.577546in}{3.111896in}}%
\pgfpathlineto{\pgfqpoint{1.578059in}{3.100180in}}%
\pgfpathlineto{\pgfqpoint{1.578572in}{3.113643in}}%
\pgfpathlineto{\pgfqpoint{1.579685in}{3.226730in}}%
\pgfpathlineto{\pgfqpoint{1.581909in}{3.737779in}}%
\pgfpathlineto{\pgfqpoint{1.587984in}{5.203759in}}%
\pgfpathlineto{\pgfqpoint{1.588498in}{5.217435in}}%
\pgfpathlineto{\pgfqpoint{1.589011in}{5.206268in}}%
\pgfpathlineto{\pgfqpoint{1.590123in}{5.099011in}}%
\pgfpathlineto{\pgfqpoint{1.592262in}{4.624495in}}%
\pgfpathlineto{\pgfqpoint{1.598680in}{3.109330in}}%
\pgfpathlineto{\pgfqpoint{1.599107in}{3.102200in}}%
\pgfpathlineto{\pgfqpoint{1.599535in}{3.112097in}}%
\pgfpathlineto{\pgfqpoint{1.600648in}{3.215490in}}%
\pgfpathlineto{\pgfqpoint{1.602787in}{3.680906in}}%
\pgfpathlineto{\pgfqpoint{1.609289in}{5.208120in}}%
\pgfpathlineto{\pgfqpoint{1.609717in}{5.215424in}}%
\pgfpathlineto{\pgfqpoint{1.610145in}{5.205935in}}%
\pgfpathlineto{\pgfqpoint{1.611257in}{5.104610in}}%
\pgfpathlineto{\pgfqpoint{1.613396in}{4.645812in}}%
\pgfpathlineto{\pgfqpoint{1.619985in}{3.111069in}}%
\pgfpathlineto{\pgfqpoint{1.620412in}{3.104243in}}%
\pgfpathlineto{\pgfqpoint{1.620840in}{3.113976in}}%
\pgfpathlineto{\pgfqpoint{1.621953in}{3.214852in}}%
\pgfpathlineto{\pgfqpoint{1.624092in}{3.669309in}}%
\pgfpathlineto{\pgfqpoint{1.630765in}{5.207451in}}%
\pgfpathlineto{\pgfqpoint{1.631108in}{5.213374in}}%
\pgfpathlineto{\pgfqpoint{1.631621in}{5.202664in}}%
\pgfpathlineto{\pgfqpoint{1.632733in}{5.100803in}}%
\pgfpathlineto{\pgfqpoint{1.634872in}{4.648651in}}%
\pgfpathlineto{\pgfqpoint{1.641546in}{3.113706in}}%
\pgfpathlineto{\pgfqpoint{1.641974in}{3.106283in}}%
\pgfpathlineto{\pgfqpoint{1.642488in}{3.118617in}}%
\pgfpathlineto{\pgfqpoint{1.643600in}{3.222714in}}%
\pgfpathlineto{\pgfqpoint{1.645824in}{3.697547in}}%
\pgfpathlineto{\pgfqpoint{1.652413in}{5.202716in}}%
\pgfpathlineto{\pgfqpoint{1.652841in}{5.211263in}}%
\pgfpathlineto{\pgfqpoint{1.653354in}{5.200580in}}%
\pgfpathlineto{\pgfqpoint{1.654466in}{5.100984in}}%
\pgfpathlineto{\pgfqpoint{1.656605in}{4.659439in}}%
\pgfpathlineto{\pgfqpoint{1.663450in}{3.114605in}}%
\pgfpathlineto{\pgfqpoint{1.663793in}{3.108485in}}%
\pgfpathlineto{\pgfqpoint{1.664306in}{3.118069in}}%
\pgfpathlineto{\pgfqpoint{1.665418in}{3.214294in}}%
\pgfpathlineto{\pgfqpoint{1.667557in}{3.647317in}}%
\pgfpathlineto{\pgfqpoint{1.674573in}{5.205146in}}%
\pgfpathlineto{\pgfqpoint{1.674916in}{5.209037in}}%
\pgfpathlineto{\pgfqpoint{1.675344in}{5.200033in}}%
\pgfpathlineto{\pgfqpoint{1.676456in}{5.106208in}}%
\pgfpathlineto{\pgfqpoint{1.678595in}{4.680341in}}%
\pgfpathlineto{\pgfqpoint{1.685697in}{3.114587in}}%
\pgfpathlineto{\pgfqpoint{1.686039in}{3.110685in}}%
\pgfpathlineto{\pgfqpoint{1.686467in}{3.119468in}}%
\pgfpathlineto{\pgfqpoint{1.687579in}{3.211706in}}%
\pgfpathlineto{\pgfqpoint{1.689718in}{3.631577in}}%
\pgfpathlineto{\pgfqpoint{1.696905in}{5.203081in}}%
\pgfpathlineto{\pgfqpoint{1.697248in}{5.206784in}}%
\pgfpathlineto{\pgfqpoint{1.697675in}{5.197959in}}%
\pgfpathlineto{\pgfqpoint{1.698788in}{5.106654in}}%
\pgfpathlineto{\pgfqpoint{1.700927in}{4.691832in}}%
\pgfpathlineto{\pgfqpoint{1.708200in}{3.116361in}}%
\pgfpathlineto{\pgfqpoint{1.708456in}{3.112968in}}%
\pgfpathlineto{\pgfqpoint{1.708970in}{3.122082in}}%
\pgfpathlineto{\pgfqpoint{1.710082in}{3.212953in}}%
\pgfpathlineto{\pgfqpoint{1.712221in}{3.623455in}}%
\pgfpathlineto{\pgfqpoint{1.719579in}{5.201488in}}%
\pgfpathlineto{\pgfqpoint{1.719836in}{5.204518in}}%
\pgfpathlineto{\pgfqpoint{1.720264in}{5.197965in}}%
\pgfpathlineto{\pgfqpoint{1.721291in}{5.124128in}}%
\pgfpathlineto{\pgfqpoint{1.723259in}{4.779198in}}%
\pgfpathlineto{\pgfqpoint{1.731301in}{3.115261in}}%
\pgfpathlineto{\pgfqpoint{1.731473in}{3.116373in}}%
\pgfpathlineto{\pgfqpoint{1.732157in}{3.143536in}}%
\pgfpathlineto{\pgfqpoint{1.733612in}{3.316029in}}%
\pgfpathlineto{\pgfqpoint{1.736435in}{3.982440in}}%
\pgfpathlineto{\pgfqpoint{1.741654in}{5.149007in}}%
\pgfpathlineto{\pgfqpoint{1.742852in}{5.202143in}}%
\pgfpathlineto{\pgfqpoint{1.743280in}{5.194562in}}%
\pgfpathlineto{\pgfqpoint{1.744307in}{5.120141in}}%
\pgfpathlineto{\pgfqpoint{1.746275in}{4.780698in}}%
\pgfpathlineto{\pgfqpoint{1.754489in}{3.117725in}}%
\pgfpathlineto{\pgfqpoint{1.754574in}{3.118235in}}%
\pgfpathlineto{\pgfqpoint{1.755173in}{3.137229in}}%
\pgfpathlineto{\pgfqpoint{1.756542in}{3.278175in}}%
\pgfpathlineto{\pgfqpoint{1.759109in}{3.833197in}}%
\pgfpathlineto{\pgfqpoint{1.765184in}{5.164602in}}%
\pgfpathlineto{\pgfqpoint{1.766125in}{5.199664in}}%
\pgfpathlineto{\pgfqpoint{1.766639in}{5.191129in}}%
\pgfpathlineto{\pgfqpoint{1.767751in}{5.106864in}}%
\pgfpathlineto{\pgfqpoint{1.769804in}{4.744281in}}%
\pgfpathlineto{\pgfqpoint{1.777933in}{3.120177in}}%
\pgfpathlineto{\pgfqpoint{1.778190in}{3.122211in}}%
\pgfpathlineto{\pgfqpoint{1.778960in}{3.156991in}}%
\pgfpathlineto{\pgfqpoint{1.780585in}{3.363062in}}%
\pgfpathlineto{\pgfqpoint{1.783751in}{4.122318in}}%
\pgfpathlineto{\pgfqpoint{1.788543in}{5.137250in}}%
\pgfpathlineto{\pgfqpoint{1.789826in}{5.197174in}}%
\pgfpathlineto{\pgfqpoint{1.790254in}{5.190999in}}%
\pgfpathlineto{\pgfqpoint{1.791281in}{5.123401in}}%
\pgfpathlineto{\pgfqpoint{1.793249in}{4.807276in}}%
\pgfpathlineto{\pgfqpoint{1.801805in}{3.122731in}}%
\pgfpathlineto{\pgfqpoint{1.802147in}{3.126537in}}%
\pgfpathlineto{\pgfqpoint{1.803088in}{3.179161in}}%
\pgfpathlineto{\pgfqpoint{1.804885in}{3.437195in}}%
\pgfpathlineto{\pgfqpoint{1.808907in}{4.441530in}}%
\pgfpathlineto{\pgfqpoint{1.812757in}{5.149962in}}%
\pgfpathlineto{\pgfqpoint{1.813869in}{5.194559in}}%
\pgfpathlineto{\pgfqpoint{1.814297in}{5.188925in}}%
\pgfpathlineto{\pgfqpoint{1.815324in}{5.124348in}}%
\pgfpathlineto{\pgfqpoint{1.817292in}{4.819311in}}%
\pgfpathlineto{\pgfqpoint{1.826105in}{3.125417in}}%
\pgfpathlineto{\pgfqpoint{1.826447in}{3.130328in}}%
\pgfpathlineto{\pgfqpoint{1.827388in}{3.184486in}}%
\pgfpathlineto{\pgfqpoint{1.829270in}{3.456102in}}%
\pgfpathlineto{\pgfqpoint{1.833549in}{4.511197in}}%
\pgfpathlineto{\pgfqpoint{1.837228in}{5.150845in}}%
\pgfpathlineto{\pgfqpoint{1.838340in}{5.191898in}}%
\pgfpathlineto{\pgfqpoint{1.838768in}{5.185676in}}%
\pgfpathlineto{\pgfqpoint{1.839795in}{5.121483in}}%
\pgfpathlineto{\pgfqpoint{1.841763in}{4.823357in}}%
\pgfpathlineto{\pgfqpoint{1.850747in}{3.128131in}}%
\pgfpathlineto{\pgfqpoint{1.851174in}{3.135185in}}%
\pgfpathlineto{\pgfqpoint{1.852287in}{3.208863in}}%
\pgfpathlineto{\pgfqpoint{1.854340in}{3.531831in}}%
\pgfpathlineto{\pgfqpoint{1.863153in}{5.189097in}}%
\pgfpathlineto{\pgfqpoint{1.863410in}{5.187527in}}%
\pgfpathlineto{\pgfqpoint{1.864180in}{5.157406in}}%
\pgfpathlineto{\pgfqpoint{1.865720in}{4.988683in}}%
\pgfpathlineto{\pgfqpoint{1.868629in}{4.373053in}}%
\pgfpathlineto{\pgfqpoint{1.874447in}{3.186045in}}%
\pgfpathlineto{\pgfqpoint{1.875731in}{3.130944in}}%
\pgfpathlineto{\pgfqpoint{1.876159in}{3.135689in}}%
\pgfpathlineto{\pgfqpoint{1.877185in}{3.193835in}}%
\pgfpathlineto{\pgfqpoint{1.879068in}{3.456800in}}%
\pgfpathlineto{\pgfqpoint{1.883346in}{4.470114in}}%
\pgfpathlineto{\pgfqpoint{1.887196in}{5.136683in}}%
\pgfpathlineto{\pgfqpoint{1.888480in}{5.186272in}}%
\pgfpathlineto{\pgfqpoint{1.888822in}{5.182241in}}%
\pgfpathlineto{\pgfqpoint{1.889763in}{5.133966in}}%
\pgfpathlineto{\pgfqpoint{1.891560in}{4.901808in}}%
\pgfpathlineto{\pgfqpoint{1.895239in}{4.061740in}}%
\pgfpathlineto{\pgfqpoint{1.899774in}{3.201129in}}%
\pgfpathlineto{\pgfqpoint{1.901228in}{3.133832in}}%
\pgfpathlineto{\pgfqpoint{1.901571in}{3.136657in}}%
\pgfpathlineto{\pgfqpoint{1.902426in}{3.174782in}}%
\pgfpathlineto{\pgfqpoint{1.904138in}{3.375683in}}%
\pgfpathlineto{\pgfqpoint{1.907474in}{4.100535in}}%
\pgfpathlineto{\pgfqpoint{1.912608in}{5.109872in}}%
\pgfpathlineto{\pgfqpoint{1.914148in}{5.183307in}}%
\pgfpathlineto{\pgfqpoint{1.914491in}{5.180467in}}%
\pgfpathlineto{\pgfqpoint{1.915346in}{5.142896in}}%
\pgfpathlineto{\pgfqpoint{1.917057in}{4.945362in}}%
\pgfpathlineto{\pgfqpoint{1.920309in}{4.251952in}}%
\pgfpathlineto{\pgfqpoint{1.925614in}{3.211359in}}%
\pgfpathlineto{\pgfqpoint{1.927239in}{3.136801in}}%
\pgfpathlineto{\pgfqpoint{1.927496in}{3.139120in}}%
\pgfpathlineto{\pgfqpoint{1.928352in}{3.174606in}}%
\pgfpathlineto{\pgfqpoint{1.930063in}{3.366046in}}%
\pgfpathlineto{\pgfqpoint{1.933314in}{4.046535in}}%
\pgfpathlineto{\pgfqpoint{1.938790in}{5.109354in}}%
\pgfpathlineto{\pgfqpoint{1.940330in}{5.180263in}}%
\pgfpathlineto{\pgfqpoint{1.940673in}{5.177596in}}%
\pgfpathlineto{\pgfqpoint{1.941528in}{5.141624in}}%
\pgfpathlineto{\pgfqpoint{1.943240in}{4.951683in}}%
\pgfpathlineto{\pgfqpoint{1.946491in}{4.279624in}}%
\pgfpathlineto{\pgfqpoint{1.952053in}{3.210494in}}%
\pgfpathlineto{\pgfqpoint{1.953678in}{3.139910in}}%
\pgfpathlineto{\pgfqpoint{1.953935in}{3.142313in}}%
\pgfpathlineto{\pgfqpoint{1.954791in}{3.176987in}}%
\pgfpathlineto{\pgfqpoint{1.956502in}{3.362244in}}%
\pgfpathlineto{\pgfqpoint{1.959668in}{4.002381in}}%
\pgfpathlineto{\pgfqpoint{1.965400in}{5.103405in}}%
\pgfpathlineto{\pgfqpoint{1.967026in}{5.177154in}}%
\pgfpathlineto{\pgfqpoint{1.967283in}{5.175539in}}%
\pgfpathlineto{\pgfqpoint{1.968053in}{5.148933in}}%
\pgfpathlineto{\pgfqpoint{1.969678in}{4.990747in}}%
\pgfpathlineto{\pgfqpoint{1.972588in}{4.439943in}}%
\pgfpathlineto{\pgfqpoint{1.979090in}{3.201883in}}%
\pgfpathlineto{\pgfqpoint{1.980545in}{3.143078in}}%
\pgfpathlineto{\pgfqpoint{1.980887in}{3.145811in}}%
\pgfpathlineto{\pgfqpoint{1.981828in}{3.185771in}}%
\pgfpathlineto{\pgfqpoint{1.983625in}{3.385649in}}%
\pgfpathlineto{\pgfqpoint{1.987048in}{4.080757in}}%
\pgfpathlineto{\pgfqpoint{1.992438in}{5.089792in}}%
\pgfpathlineto{\pgfqpoint{1.994235in}{5.173914in}}%
\pgfpathlineto{\pgfqpoint{1.994406in}{5.173079in}}%
\pgfpathlineto{\pgfqpoint{1.995090in}{5.154252in}}%
\pgfpathlineto{\pgfqpoint{1.996545in}{5.034885in}}%
\pgfpathlineto{\pgfqpoint{1.999197in}{4.587653in}}%
\pgfpathlineto{\pgfqpoint{2.006984in}{3.174374in}}%
\pgfpathlineto{\pgfqpoint{2.008010in}{3.146367in}}%
\pgfpathlineto{\pgfqpoint{2.008438in}{3.150833in}}%
\pgfpathlineto{\pgfqpoint{2.009465in}{3.199868in}}%
\pgfpathlineto{\pgfqpoint{2.011347in}{3.419679in}}%
\pgfpathlineto{\pgfqpoint{2.015026in}{4.171067in}}%
\pgfpathlineto{\pgfqpoint{2.020160in}{5.089407in}}%
\pgfpathlineto{\pgfqpoint{2.021957in}{5.170559in}}%
\pgfpathlineto{\pgfqpoint{2.022128in}{5.169813in}}%
\pgfpathlineto{\pgfqpoint{2.022813in}{5.151995in}}%
\pgfpathlineto{\pgfqpoint{2.024267in}{5.037997in}}%
\pgfpathlineto{\pgfqpoint{2.026834in}{4.625253in}}%
\pgfpathlineto{\pgfqpoint{2.035134in}{3.169848in}}%
\pgfpathlineto{\pgfqpoint{2.036075in}{3.149821in}}%
\pgfpathlineto{\pgfqpoint{2.036503in}{3.155263in}}%
\pgfpathlineto{\pgfqpoint{2.037615in}{3.211354in}}%
\pgfpathlineto{\pgfqpoint{2.039668in}{3.460369in}}%
\pgfpathlineto{\pgfqpoint{2.043861in}{4.321869in}}%
\pgfpathlineto{\pgfqpoint{2.048396in}{5.084014in}}%
\pgfpathlineto{\pgfqpoint{2.050278in}{5.167072in}}%
\pgfpathlineto{\pgfqpoint{2.050449in}{5.166168in}}%
\pgfpathlineto{\pgfqpoint{2.051219in}{5.144563in}}%
\pgfpathlineto{\pgfqpoint{2.052760in}{5.018589in}}%
\pgfpathlineto{\pgfqpoint{2.055498in}{4.571445in}}%
\pgfpathlineto{\pgfqpoint{2.063455in}{3.185876in}}%
\pgfpathlineto{\pgfqpoint{2.064653in}{3.153327in}}%
\pgfpathlineto{\pgfqpoint{2.065081in}{3.158182in}}%
\pgfpathlineto{\pgfqpoint{2.066107in}{3.204855in}}%
\pgfpathlineto{\pgfqpoint{2.068075in}{3.422204in}}%
\pgfpathlineto{\pgfqpoint{2.071840in}{4.151829in}}%
\pgfpathlineto{\pgfqpoint{2.077145in}{5.070938in}}%
\pgfpathlineto{\pgfqpoint{2.079113in}{5.163447in}}%
\pgfpathlineto{\pgfqpoint{2.079284in}{5.163101in}}%
\pgfpathlineto{\pgfqpoint{2.079883in}{5.151212in}}%
\pgfpathlineto{\pgfqpoint{2.081252in}{5.063241in}}%
\pgfpathlineto{\pgfqpoint{2.083648in}{4.731501in}}%
\pgfpathlineto{\pgfqpoint{2.093573in}{3.158606in}}%
\pgfpathlineto{\pgfqpoint{2.093830in}{3.156977in}}%
\pgfpathlineto{\pgfqpoint{2.094257in}{3.160899in}}%
\pgfpathlineto{\pgfqpoint{2.095284in}{3.203775in}}%
\pgfpathlineto{\pgfqpoint{2.097166in}{3.396859in}}%
\pgfpathlineto{\pgfqpoint{2.100675in}{4.038046in}}%
\pgfpathlineto{\pgfqpoint{2.106664in}{5.070006in}}%
\pgfpathlineto{\pgfqpoint{2.108717in}{5.159744in}}%
\pgfpathlineto{\pgfqpoint{2.108803in}{5.159482in}}%
\pgfpathlineto{\pgfqpoint{2.109402in}{5.148600in}}%
\pgfpathlineto{\pgfqpoint{2.110685in}{5.073136in}}%
\pgfpathlineto{\pgfqpoint{2.112995in}{4.778132in}}%
\pgfpathlineto{\pgfqpoint{2.123691in}{3.160778in}}%
\pgfpathlineto{\pgfqpoint{2.123776in}{3.160867in}}%
\pgfpathlineto{\pgfqpoint{2.124290in}{3.168031in}}%
\pgfpathlineto{\pgfqpoint{2.125488in}{3.228187in}}%
\pgfpathlineto{\pgfqpoint{2.127627in}{3.473277in}}%
\pgfpathlineto{\pgfqpoint{2.131905in}{4.286477in}}%
\pgfpathlineto{\pgfqpoint{2.136782in}{5.062584in}}%
\pgfpathlineto{\pgfqpoint{2.138921in}{5.155885in}}%
\pgfpathlineto{\pgfqpoint{2.139349in}{5.151685in}}%
\pgfpathlineto{\pgfqpoint{2.140375in}{5.110540in}}%
\pgfpathlineto{\pgfqpoint{2.142258in}{4.928525in}}%
\pgfpathlineto{\pgfqpoint{2.145680in}{4.338912in}}%
\pgfpathlineto{\pgfqpoint{2.152183in}{3.253156in}}%
\pgfpathlineto{\pgfqpoint{2.154237in}{3.164736in}}%
\pgfpathlineto{\pgfqpoint{2.154322in}{3.164760in}}%
\pgfpathlineto{\pgfqpoint{2.154750in}{3.169382in}}%
\pgfpathlineto{\pgfqpoint{2.155862in}{3.216069in}}%
\pgfpathlineto{\pgfqpoint{2.157830in}{3.412688in}}%
\pgfpathlineto{\pgfqpoint{2.161509in}{4.054964in}}%
\pgfpathlineto{\pgfqpoint{2.167584in}{5.052681in}}%
\pgfpathlineto{\pgfqpoint{2.169809in}{5.151871in}}%
\pgfpathlineto{\pgfqpoint{2.170066in}{5.150752in}}%
\pgfpathlineto{\pgfqpoint{2.170836in}{5.131668in}}%
\pgfpathlineto{\pgfqpoint{2.172376in}{5.025164in}}%
\pgfpathlineto{\pgfqpoint{2.175114in}{4.645775in}}%
\pgfpathlineto{\pgfqpoint{2.184611in}{3.186828in}}%
\pgfpathlineto{\pgfqpoint{2.185552in}{3.168810in}}%
\pgfpathlineto{\pgfqpoint{2.186066in}{3.173496in}}%
\pgfpathlineto{\pgfqpoint{2.187178in}{3.218389in}}%
\pgfpathlineto{\pgfqpoint{2.189232in}{3.417264in}}%
\pgfpathlineto{\pgfqpoint{2.192911in}{4.039926in}}%
\pgfpathlineto{\pgfqpoint{2.199157in}{5.043771in}}%
\pgfpathlineto{\pgfqpoint{2.201467in}{5.147660in}}%
\pgfpathlineto{\pgfqpoint{2.201638in}{5.147431in}}%
\pgfpathlineto{\pgfqpoint{2.201724in}{5.146901in}}%
\pgfpathlineto{\pgfqpoint{2.202494in}{5.129712in}}%
\pgfpathlineto{\pgfqpoint{2.204034in}{5.030365in}}%
\pgfpathlineto{\pgfqpoint{2.206686in}{4.685121in}}%
\pgfpathlineto{\pgfqpoint{2.216954in}{3.182331in}}%
\pgfpathlineto{\pgfqpoint{2.217638in}{3.173066in}}%
\pgfpathlineto{\pgfqpoint{2.218152in}{3.177440in}}%
\pgfpathlineto{\pgfqpoint{2.219264in}{3.219820in}}%
\pgfpathlineto{\pgfqpoint{2.221232in}{3.397888in}}%
\pgfpathlineto{\pgfqpoint{2.224740in}{3.956522in}}%
\pgfpathlineto{\pgfqpoint{2.231671in}{5.045923in}}%
\pgfpathlineto{\pgfqpoint{2.233981in}{5.143342in}}%
\pgfpathlineto{\pgfqpoint{2.234237in}{5.142523in}}%
\pgfpathlineto{\pgfqpoint{2.235008in}{5.125966in}}%
\pgfpathlineto{\pgfqpoint{2.236548in}{5.031360in}}%
\pgfpathlineto{\pgfqpoint{2.239200in}{4.702423in}}%
\pgfpathlineto{\pgfqpoint{2.249981in}{3.184015in}}%
\pgfpathlineto{\pgfqpoint{2.250580in}{3.177484in}}%
\pgfpathlineto{\pgfqpoint{2.251093in}{3.181816in}}%
\pgfpathlineto{\pgfqpoint{2.252206in}{3.222277in}}%
\pgfpathlineto{\pgfqpoint{2.254174in}{3.391506in}}%
\pgfpathlineto{\pgfqpoint{2.257682in}{3.925357in}}%
\pgfpathlineto{\pgfqpoint{2.265040in}{5.046595in}}%
\pgfpathlineto{\pgfqpoint{2.267350in}{5.138821in}}%
\pgfpathlineto{\pgfqpoint{2.267607in}{5.138079in}}%
\pgfpathlineto{\pgfqpoint{2.268377in}{5.122553in}}%
\pgfpathlineto{\pgfqpoint{2.269917in}{5.033400in}}%
\pgfpathlineto{\pgfqpoint{2.272569in}{4.721740in}}%
\pgfpathlineto{\pgfqpoint{2.283949in}{3.185804in}}%
\pgfpathlineto{\pgfqpoint{2.284463in}{3.182115in}}%
\pgfpathlineto{\pgfqpoint{2.284890in}{3.185636in}}%
\pgfpathlineto{\pgfqpoint{2.286003in}{3.222546in}}%
\pgfpathlineto{\pgfqpoint{2.287971in}{3.380387in}}%
\pgfpathlineto{\pgfqpoint{2.291393in}{3.871795in}}%
\pgfpathlineto{\pgfqpoint{2.299436in}{5.053691in}}%
\pgfpathlineto{\pgfqpoint{2.301661in}{5.134119in}}%
\pgfpathlineto{\pgfqpoint{2.301917in}{5.133388in}}%
\pgfpathlineto{\pgfqpoint{2.302687in}{5.118684in}}%
\pgfpathlineto{\pgfqpoint{2.304228in}{5.034546in}}%
\pgfpathlineto{\pgfqpoint{2.306880in}{4.739588in}}%
\pgfpathlineto{\pgfqpoint{2.318944in}{3.188113in}}%
\pgfpathlineto{\pgfqpoint{2.319201in}{3.186911in}}%
\pgfpathlineto{\pgfqpoint{2.319714in}{3.190591in}}%
\pgfpathlineto{\pgfqpoint{2.320827in}{3.226093in}}%
\pgfpathlineto{\pgfqpoint{2.322795in}{3.376006in}}%
\pgfpathlineto{\pgfqpoint{2.326217in}{3.843687in}}%
\pgfpathlineto{\pgfqpoint{2.334688in}{5.049738in}}%
\pgfpathlineto{\pgfqpoint{2.336912in}{5.129083in}}%
\pgfpathlineto{\pgfqpoint{2.336998in}{5.129224in}}%
\pgfpathlineto{\pgfqpoint{2.337255in}{5.128340in}}%
\pgfpathlineto{\pgfqpoint{2.338025in}{5.113954in}}%
\pgfpathlineto{\pgfqpoint{2.339565in}{5.033825in}}%
\pgfpathlineto{\pgfqpoint{2.342217in}{4.754169in}}%
\pgfpathlineto{\pgfqpoint{2.348121in}{3.811780in}}%
\pgfpathlineto{\pgfqpoint{2.352656in}{3.274775in}}%
\pgfpathlineto{\pgfqpoint{2.354966in}{3.192077in}}%
\pgfpathlineto{\pgfqpoint{2.355052in}{3.191932in}}%
\pgfpathlineto{\pgfqpoint{2.355394in}{3.193461in}}%
\pgfpathlineto{\pgfqpoint{2.356335in}{3.214968in}}%
\pgfpathlineto{\pgfqpoint{2.358046in}{3.316708in}}%
\pgfpathlineto{\pgfqpoint{2.360955in}{3.647804in}}%
\pgfpathlineto{\pgfqpoint{2.372335in}{5.108116in}}%
\pgfpathlineto{\pgfqpoint{2.373448in}{5.124065in}}%
\pgfpathlineto{\pgfqpoint{2.373875in}{5.121031in}}%
\pgfpathlineto{\pgfqpoint{2.374988in}{5.089521in}}%
\pgfpathlineto{\pgfqpoint{2.376956in}{4.954404in}}%
\pgfpathlineto{\pgfqpoint{2.380293in}{4.538885in}}%
\pgfpathlineto{\pgfqpoint{2.389876in}{3.259648in}}%
\pgfpathlineto{\pgfqpoint{2.392015in}{3.197194in}}%
\pgfpathlineto{\pgfqpoint{2.392357in}{3.198539in}}%
\pgfpathlineto{\pgfqpoint{2.393298in}{3.218392in}}%
\pgfpathlineto{\pgfqpoint{2.395009in}{3.313112in}}%
\pgfpathlineto{\pgfqpoint{2.397918in}{3.623717in}}%
\pgfpathlineto{\pgfqpoint{2.410154in}{5.109881in}}%
\pgfpathlineto{\pgfqpoint{2.411010in}{5.118695in}}%
\pgfpathlineto{\pgfqpoint{2.411437in}{5.115979in}}%
\pgfpathlineto{\pgfqpoint{2.412464in}{5.090248in}}%
\pgfpathlineto{\pgfqpoint{2.414346in}{4.975655in}}%
\pgfpathlineto{\pgfqpoint{2.417512in}{4.618782in}}%
\pgfpathlineto{\pgfqpoint{2.428635in}{3.234683in}}%
\pgfpathlineto{\pgfqpoint{2.430261in}{3.202702in}}%
\pgfpathlineto{\pgfqpoint{2.430518in}{3.203678in}}%
\pgfpathlineto{\pgfqpoint{2.431373in}{3.218777in}}%
\pgfpathlineto{\pgfqpoint{2.432999in}{3.296296in}}%
\pgfpathlineto{\pgfqpoint{2.435737in}{3.555049in}}%
\pgfpathlineto{\pgfqpoint{2.441213in}{4.344239in}}%
\pgfpathlineto{\pgfqpoint{2.446604in}{4.989961in}}%
\pgfpathlineto{\pgfqpoint{2.449342in}{5.110186in}}%
\pgfpathlineto{\pgfqpoint{2.449855in}{5.113056in}}%
\pgfpathlineto{\pgfqpoint{2.450368in}{5.109579in}}%
\pgfpathlineto{\pgfqpoint{2.451481in}{5.080488in}}%
\pgfpathlineto{\pgfqpoint{2.453449in}{4.960165in}}%
\pgfpathlineto{\pgfqpoint{2.456700in}{4.601962in}}%
\pgfpathlineto{\pgfqpoint{2.467909in}{3.248523in}}%
\pgfpathlineto{\pgfqpoint{2.469791in}{3.208484in}}%
\pgfpathlineto{\pgfqpoint{2.469962in}{3.208897in}}%
\pgfpathlineto{\pgfqpoint{2.470732in}{3.219125in}}%
\pgfpathlineto{\pgfqpoint{2.472272in}{3.279811in}}%
\pgfpathlineto{\pgfqpoint{2.474839in}{3.489419in}}%
\pgfpathlineto{\pgfqpoint{2.479374in}{4.084159in}}%
\pgfpathlineto{\pgfqpoint{2.486476in}{4.965604in}}%
\pgfpathlineto{\pgfqpoint{2.489299in}{5.100340in}}%
\pgfpathlineto{\pgfqpoint{2.490069in}{5.107130in}}%
\pgfpathlineto{\pgfqpoint{2.490583in}{5.104325in}}%
\pgfpathlineto{\pgfqpoint{2.491695in}{5.078303in}}%
\pgfpathlineto{\pgfqpoint{2.493663in}{4.968292in}}%
\pgfpathlineto{\pgfqpoint{2.496914in}{4.636027in}}%
\pgfpathlineto{\pgfqpoint{2.509150in}{3.242596in}}%
\pgfpathlineto{\pgfqpoint{2.510775in}{3.214566in}}%
\pgfpathlineto{\pgfqpoint{2.511032in}{3.215287in}}%
\pgfpathlineto{\pgfqpoint{2.511888in}{3.227813in}}%
\pgfpathlineto{\pgfqpoint{2.513513in}{3.293490in}}%
\pgfpathlineto{\pgfqpoint{2.516166in}{3.507136in}}%
\pgfpathlineto{\pgfqpoint{2.520872in}{4.102661in}}%
\pgfpathlineto{\pgfqpoint{2.528059in}{4.955712in}}%
\pgfpathlineto{\pgfqpoint{2.530968in}{5.092614in}}%
\pgfpathlineto{\pgfqpoint{2.531909in}{5.100865in}}%
\pgfpathlineto{\pgfqpoint{2.532337in}{5.098624in}}%
\pgfpathlineto{\pgfqpoint{2.533449in}{5.075407in}}%
\pgfpathlineto{\pgfqpoint{2.535417in}{4.975212in}}%
\pgfpathlineto{\pgfqpoint{2.538583in}{4.678353in}}%
\pgfpathlineto{\pgfqpoint{2.552188in}{3.235038in}}%
\pgfpathlineto{\pgfqpoint{2.553385in}{3.220983in}}%
\pgfpathlineto{\pgfqpoint{2.553813in}{3.222771in}}%
\pgfpathlineto{\pgfqpoint{2.554840in}{3.241593in}}%
\pgfpathlineto{\pgfqpoint{2.556722in}{3.327594in}}%
\pgfpathlineto{\pgfqpoint{2.559717in}{3.584821in}}%
\pgfpathlineto{\pgfqpoint{2.565706in}{4.343756in}}%
\pgfpathlineto{\pgfqpoint{2.571439in}{4.954004in}}%
\pgfpathlineto{\pgfqpoint{2.574434in}{5.086404in}}%
\pgfpathlineto{\pgfqpoint{2.575375in}{5.094281in}}%
\pgfpathlineto{\pgfqpoint{2.575803in}{5.092381in}}%
\pgfpathlineto{\pgfqpoint{2.576830in}{5.073935in}}%
\pgfpathlineto{\pgfqpoint{2.578712in}{4.990881in}}%
\pgfpathlineto{\pgfqpoint{2.581707in}{4.743136in}}%
\pgfpathlineto{\pgfqpoint{2.587439in}{4.039285in}}%
\pgfpathlineto{\pgfqpoint{2.593685in}{3.374825in}}%
\pgfpathlineto{\pgfqpoint{2.596766in}{3.237110in}}%
\pgfpathlineto{\pgfqpoint{2.597792in}{3.227769in}}%
\pgfpathlineto{\pgfqpoint{2.598220in}{3.229434in}}%
\pgfpathlineto{\pgfqpoint{2.599247in}{3.246685in}}%
\pgfpathlineto{\pgfqpoint{2.601129in}{3.325395in}}%
\pgfpathlineto{\pgfqpoint{2.604124in}{3.561860in}}%
\pgfpathlineto{\pgfqpoint{2.609686in}{4.220992in}}%
\pgfpathlineto{\pgfqpoint{2.616359in}{4.929112in}}%
\pgfpathlineto{\pgfqpoint{2.619525in}{5.075090in}}%
\pgfpathlineto{\pgfqpoint{2.620723in}{5.087302in}}%
\pgfpathlineto{\pgfqpoint{2.621151in}{5.085747in}}%
\pgfpathlineto{\pgfqpoint{2.622178in}{5.069387in}}%
\pgfpathlineto{\pgfqpoint{2.624060in}{4.994469in}}%
\pgfpathlineto{\pgfqpoint{2.627055in}{4.768568in}}%
\pgfpathlineto{\pgfqpoint{2.632360in}{4.163867in}}%
\pgfpathlineto{\pgfqpoint{2.639632in}{3.398627in}}%
\pgfpathlineto{\pgfqpoint{2.642884in}{3.248686in}}%
\pgfpathlineto{\pgfqpoint{2.644167in}{3.234968in}}%
\pgfpathlineto{\pgfqpoint{2.644595in}{3.236316in}}%
\pgfpathlineto{\pgfqpoint{2.645622in}{3.251569in}}%
\pgfpathlineto{\pgfqpoint{2.647504in}{3.322338in}}%
\pgfpathlineto{\pgfqpoint{2.650499in}{3.537209in}}%
\pgfpathlineto{\pgfqpoint{2.655633in}{4.098646in}}%
\pgfpathlineto{\pgfqpoint{2.663504in}{4.914110in}}%
\pgfpathlineto{\pgfqpoint{2.666841in}{5.065516in}}%
\pgfpathlineto{\pgfqpoint{2.668210in}{5.079882in}}%
\pgfpathlineto{\pgfqpoint{2.668553in}{5.078977in}}%
\pgfpathlineto{\pgfqpoint{2.669494in}{5.067253in}}%
\pgfpathlineto{\pgfqpoint{2.671205in}{5.012008in}}%
\pgfpathlineto{\pgfqpoint{2.673943in}{4.840097in}}%
\pgfpathlineto{\pgfqpoint{2.678478in}{4.388659in}}%
\pgfpathlineto{\pgfqpoint{2.688403in}{3.382176in}}%
\pgfpathlineto{\pgfqpoint{2.691569in}{3.254157in}}%
\pgfpathlineto{\pgfqpoint{2.692852in}{3.242627in}}%
\pgfpathlineto{\pgfqpoint{2.693195in}{3.243600in}}%
\pgfpathlineto{\pgfqpoint{2.694136in}{3.255022in}}%
\pgfpathlineto{\pgfqpoint{2.695933in}{3.311642in}}%
\pgfpathlineto{\pgfqpoint{2.698756in}{3.485162in}}%
\pgfpathlineto{\pgfqpoint{2.703376in}{3.932895in}}%
\pgfpathlineto{\pgfqpoint{2.713473in}{4.928932in}}%
\pgfpathlineto{\pgfqpoint{2.716724in}{5.059296in}}%
\pgfpathlineto{\pgfqpoint{2.718093in}{5.071970in}}%
\pgfpathlineto{\pgfqpoint{2.718435in}{5.071112in}}%
\pgfpathlineto{\pgfqpoint{2.719377in}{5.060482in}}%
\pgfpathlineto{\pgfqpoint{2.721088in}{5.010723in}}%
\pgfpathlineto{\pgfqpoint{2.723826in}{4.855630in}}%
\pgfpathlineto{\pgfqpoint{2.728190in}{4.461122in}}%
\pgfpathlineto{\pgfqpoint{2.739740in}{3.367350in}}%
\pgfpathlineto{\pgfqpoint{2.742821in}{3.260053in}}%
\pgfpathlineto{\pgfqpoint{2.744019in}{3.250804in}}%
\pgfpathlineto{\pgfqpoint{2.744446in}{3.252010in}}%
\pgfpathlineto{\pgfqpoint{2.745473in}{3.264540in}}%
\pgfpathlineto{\pgfqpoint{2.747356in}{3.321966in}}%
\pgfpathlineto{\pgfqpoint{2.750265in}{3.490939in}}%
\pgfpathlineto{\pgfqpoint{2.754971in}{3.916942in}}%
\pgfpathlineto{\pgfqpoint{2.765751in}{4.918716in}}%
\pgfpathlineto{\pgfqpoint{2.769088in}{5.048141in}}%
\pgfpathlineto{\pgfqpoint{2.770714in}{5.063492in}}%
\pgfpathlineto{\pgfqpoint{2.770971in}{5.062970in}}%
\pgfpathlineto{\pgfqpoint{2.771826in}{5.055446in}}%
\pgfpathlineto{\pgfqpoint{2.773452in}{5.017030in}}%
\pgfpathlineto{\pgfqpoint{2.776105in}{4.890650in}}%
\pgfpathlineto{\pgfqpoint{2.780212in}{4.568270in}}%
\pgfpathlineto{\pgfqpoint{2.794672in}{3.327022in}}%
\pgfpathlineto{\pgfqpoint{2.797410in}{3.262466in}}%
\pgfpathlineto{\pgfqpoint{2.798094in}{3.259588in}}%
\pgfpathlineto{\pgfqpoint{2.798607in}{3.260948in}}%
\pgfpathlineto{\pgfqpoint{2.799720in}{3.274178in}}%
\pgfpathlineto{\pgfqpoint{2.801688in}{3.331175in}}%
\pgfpathlineto{\pgfqpoint{2.804682in}{3.493252in}}%
\pgfpathlineto{\pgfqpoint{2.809474in}{3.894255in}}%
\pgfpathlineto{\pgfqpoint{2.821196in}{4.912120in}}%
\pgfpathlineto{\pgfqpoint{2.824704in}{5.038888in}}%
\pgfpathlineto{\pgfqpoint{2.826415in}{5.054385in}}%
\pgfpathlineto{\pgfqpoint{2.826586in}{5.054212in}}%
\pgfpathlineto{\pgfqpoint{2.827356in}{5.049567in}}%
\pgfpathlineto{\pgfqpoint{2.828897in}{5.021516in}}%
\pgfpathlineto{\pgfqpoint{2.831378in}{4.926146in}}%
\pgfpathlineto{\pgfqpoint{2.835143in}{4.679922in}}%
\pgfpathlineto{\pgfqpoint{2.842244in}{4.031174in}}%
\pgfpathlineto{\pgfqpoint{2.849432in}{3.453949in}}%
\pgfpathlineto{\pgfqpoint{2.853282in}{3.295977in}}%
\pgfpathlineto{\pgfqpoint{2.855506in}{3.269113in}}%
\pgfpathlineto{\pgfqpoint{2.855763in}{3.269182in}}%
\pgfpathlineto{\pgfqpoint{2.856020in}{3.269908in}}%
\pgfpathlineto{\pgfqpoint{2.857047in}{3.279344in}}%
\pgfpathlineto{\pgfqpoint{2.858929in}{3.323268in}}%
\pgfpathlineto{\pgfqpoint{2.861838in}{3.454279in}}%
\pgfpathlineto{\pgfqpoint{2.866202in}{3.766421in}}%
\pgfpathlineto{\pgfqpoint{2.881603in}{4.963716in}}%
\pgfpathlineto{\pgfqpoint{2.884683in}{5.038717in}}%
\pgfpathlineto{\pgfqpoint{2.885796in}{5.044515in}}%
\pgfpathlineto{\pgfqpoint{2.886223in}{5.043695in}}%
\pgfpathlineto{\pgfqpoint{2.887250in}{5.034849in}}%
\pgfpathlineto{\pgfqpoint{2.889133in}{4.993891in}}%
\pgfpathlineto{\pgfqpoint{2.892042in}{4.871695in}}%
\pgfpathlineto{\pgfqpoint{2.896405in}{4.578931in}}%
\pgfpathlineto{\pgfqpoint{2.913176in}{3.344152in}}%
\pgfpathlineto{\pgfqpoint{2.916170in}{3.282911in}}%
\pgfpathlineto{\pgfqpoint{2.917111in}{3.279362in}}%
\pgfpathlineto{\pgfqpoint{2.917539in}{3.280256in}}%
\pgfpathlineto{\pgfqpoint{2.918652in}{3.289870in}}%
\pgfpathlineto{\pgfqpoint{2.920534in}{3.329661in}}%
\pgfpathlineto{\pgfqpoint{2.923443in}{3.445613in}}%
\pgfpathlineto{\pgfqpoint{2.927807in}{3.721750in}}%
\pgfpathlineto{\pgfqpoint{2.945775in}{4.977969in}}%
\pgfpathlineto{\pgfqpoint{2.948684in}{5.030800in}}%
\pgfpathlineto{\pgfqpoint{2.949540in}{5.033751in}}%
\pgfpathlineto{\pgfqpoint{2.950053in}{5.032753in}}%
\pgfpathlineto{\pgfqpoint{2.951165in}{5.023510in}}%
\pgfpathlineto{\pgfqpoint{2.953133in}{4.983892in}}%
\pgfpathlineto{\pgfqpoint{2.956128in}{4.870180in}}%
\pgfpathlineto{\pgfqpoint{2.960577in}{4.602800in}}%
\pgfpathlineto{\pgfqpoint{2.979486in}{3.343614in}}%
\pgfpathlineto{\pgfqpoint{2.982396in}{3.293784in}}%
\pgfpathlineto{\pgfqpoint{2.983337in}{3.290638in}}%
\pgfpathlineto{\pgfqpoint{2.983850in}{3.291614in}}%
\pgfpathlineto{\pgfqpoint{2.984962in}{3.300211in}}%
\pgfpathlineto{\pgfqpoint{2.986930in}{3.336740in}}%
\pgfpathlineto{\pgfqpoint{2.989925in}{3.441484in}}%
\pgfpathlineto{\pgfqpoint{2.994289in}{3.683407in}}%
\pgfpathlineto{\pgfqpoint{3.002588in}{4.298030in}}%
\pgfpathlineto{\pgfqpoint{3.010460in}{4.812922in}}%
\pgfpathlineto{\pgfqpoint{3.014909in}{4.978188in}}%
\pgfpathlineto{\pgfqpoint{3.017733in}{5.019625in}}%
\pgfpathlineto{\pgfqpoint{3.018589in}{5.021901in}}%
\pgfpathlineto{\pgfqpoint{3.019102in}{5.020961in}}%
\pgfpathlineto{\pgfqpoint{3.020300in}{5.012081in}}%
\pgfpathlineto{\pgfqpoint{3.022353in}{4.975506in}}%
\pgfpathlineto{\pgfqpoint{3.025434in}{4.873110in}}%
\pgfpathlineto{\pgfqpoint{3.029883in}{4.640304in}}%
\pgfpathlineto{\pgfqpoint{3.038011in}{4.071555in}}%
\pgfpathlineto{\pgfqpoint{3.046567in}{3.528361in}}%
\pgfpathlineto{\pgfqpoint{3.051188in}{3.356238in}}%
\pgfpathlineto{\pgfqpoint{3.054268in}{3.307122in}}%
\pgfpathlineto{\pgfqpoint{3.055466in}{3.303124in}}%
\pgfpathlineto{\pgfqpoint{3.055894in}{3.303760in}}%
\pgfpathlineto{\pgfqpoint{3.057006in}{3.310468in}}%
\pgfpathlineto{\pgfqpoint{3.058888in}{3.338196in}}%
\pgfpathlineto{\pgfqpoint{3.061798in}{3.419583in}}%
\pgfpathlineto{\pgfqpoint{3.065990in}{3.608818in}}%
\pgfpathlineto{\pgfqpoint{3.072664in}{4.027981in}}%
\pgfpathlineto{\pgfqpoint{3.084728in}{4.781454in}}%
\pgfpathlineto{\pgfqpoint{3.089520in}{4.952528in}}%
\pgfpathlineto{\pgfqpoint{3.092686in}{5.003189in}}%
\pgfpathlineto{\pgfqpoint{3.094140in}{5.008708in}}%
\pgfpathlineto{\pgfqpoint{3.094482in}{5.008370in}}%
\pgfpathlineto{\pgfqpoint{3.095509in}{5.003630in}}%
\pgfpathlineto{\pgfqpoint{3.097306in}{4.982060in}}%
\pgfpathlineto{\pgfqpoint{3.100044in}{4.917977in}}%
\pgfpathlineto{\pgfqpoint{3.103980in}{4.766525in}}%
\pgfpathlineto{\pgfqpoint{3.109884in}{4.440967in}}%
\pgfpathlineto{\pgfqpoint{3.126397in}{3.484642in}}%
\pgfpathlineto{\pgfqpoint{3.130846in}{3.356013in}}%
\pgfpathlineto{\pgfqpoint{3.133841in}{3.319800in}}%
\pgfpathlineto{\pgfqpoint{3.134953in}{3.317125in}}%
\pgfpathlineto{\pgfqpoint{3.135381in}{3.317655in}}%
\pgfpathlineto{\pgfqpoint{3.136494in}{3.323065in}}%
\pgfpathlineto{\pgfqpoint{3.138462in}{3.346696in}}%
\pgfpathlineto{\pgfqpoint{3.141371in}{3.413142in}}%
\pgfpathlineto{\pgfqpoint{3.145478in}{3.564165in}}%
\pgfpathlineto{\pgfqpoint{3.151638in}{3.884295in}}%
\pgfpathlineto{\pgfqpoint{3.168751in}{4.816784in}}%
\pgfpathlineto{\pgfqpoint{3.173457in}{4.949420in}}%
\pgfpathlineto{\pgfqpoint{3.176622in}{4.989367in}}%
\pgfpathlineto{\pgfqpoint{3.178163in}{4.993774in}}%
\pgfpathlineto{\pgfqpoint{3.178419in}{4.993548in}}%
\pgfpathlineto{\pgfqpoint{3.179446in}{4.989915in}}%
\pgfpathlineto{\pgfqpoint{3.181243in}{4.973160in}}%
\pgfpathlineto{\pgfqpoint{3.183981in}{4.923031in}}%
\pgfpathlineto{\pgfqpoint{3.187831in}{4.806427in}}%
\pgfpathlineto{\pgfqpoint{3.193393in}{4.561173in}}%
\pgfpathlineto{\pgfqpoint{3.205628in}{3.887530in}}%
\pgfpathlineto{\pgfqpoint{3.213414in}{3.531982in}}%
\pgfpathlineto{\pgfqpoint{3.218462in}{3.389146in}}%
\pgfpathlineto{\pgfqpoint{3.221970in}{3.340853in}}%
\pgfpathlineto{\pgfqpoint{3.224110in}{3.333115in}}%
\pgfpathlineto{\pgfqpoint{3.224965in}{3.334650in}}%
\pgfpathlineto{\pgfqpoint{3.226505in}{3.344017in}}%
\pgfpathlineto{\pgfqpoint{3.228901in}{3.375096in}}%
\pgfpathlineto{\pgfqpoint{3.232324in}{3.452417in}}%
\pgfpathlineto{\pgfqpoint{3.237115in}{3.617452in}}%
\pgfpathlineto{\pgfqpoint{3.244302in}{3.952600in}}%
\pgfpathlineto{\pgfqpoint{3.260901in}{4.746259in}}%
\pgfpathlineto{\pgfqpoint{3.266463in}{4.904787in}}%
\pgfpathlineto{\pgfqpoint{3.270399in}{4.963593in}}%
\pgfpathlineto{\pgfqpoint{3.272966in}{4.976386in}}%
\pgfpathlineto{\pgfqpoint{3.273650in}{4.976354in}}%
\pgfpathlineto{\pgfqpoint{3.273907in}{4.975969in}}%
\pgfpathlineto{\pgfqpoint{3.275190in}{4.971015in}}%
\pgfpathlineto{\pgfqpoint{3.277329in}{4.951687in}}%
\pgfpathlineto{\pgfqpoint{3.280410in}{4.900477in}}%
\pgfpathlineto{\pgfqpoint{3.284688in}{4.787607in}}%
\pgfpathlineto{\pgfqpoint{3.290677in}{4.563701in}}%
\pgfpathlineto{\pgfqpoint{3.302057in}{4.033605in}}%
\pgfpathlineto{\pgfqpoint{3.312068in}{3.613189in}}%
\pgfpathlineto{\pgfqpoint{3.318143in}{3.442038in}}%
\pgfpathlineto{\pgfqpoint{3.322506in}{3.372632in}}%
\pgfpathlineto{\pgfqpoint{3.325501in}{3.353154in}}%
\pgfpathlineto{\pgfqpoint{3.326699in}{3.351882in}}%
\pgfpathlineto{\pgfqpoint{3.327041in}{3.352202in}}%
\pgfpathlineto{\pgfqpoint{3.328325in}{3.356085in}}%
\pgfpathlineto{\pgfqpoint{3.330378in}{3.371006in}}%
\pgfpathlineto{\pgfqpoint{3.333373in}{3.411348in}}%
\pgfpathlineto{\pgfqpoint{3.337565in}{3.502148in}}%
\pgfpathlineto{\pgfqpoint{3.343298in}{3.681001in}}%
\pgfpathlineto{\pgfqpoint{3.352197in}{4.039567in}}%
\pgfpathlineto{\pgfqpoint{3.367940in}{4.671534in}}%
\pgfpathlineto{\pgfqpoint{3.374528in}{4.848901in}}%
\pgfpathlineto{\pgfqpoint{3.379320in}{4.926269in}}%
\pgfpathlineto{\pgfqpoint{3.382742in}{4.952295in}}%
\pgfpathlineto{\pgfqpoint{3.384710in}{4.955990in}}%
\pgfpathlineto{\pgfqpoint{3.384796in}{4.955964in}}%
\pgfpathlineto{\pgfqpoint{3.385908in}{4.954224in}}%
\pgfpathlineto{\pgfqpoint{3.387705in}{4.945945in}}%
\pgfpathlineto{\pgfqpoint{3.390443in}{4.920642in}}%
\pgfpathlineto{\pgfqpoint{3.394208in}{4.862106in}}%
\pgfpathlineto{\pgfqpoint{3.399256in}{4.744996in}}%
\pgfpathlineto{\pgfqpoint{3.406272in}{4.524487in}}%
\pgfpathlineto{\pgfqpoint{3.420818in}{3.975067in}}%
\pgfpathlineto{\pgfqpoint{3.431085in}{3.634423in}}%
\pgfpathlineto{\pgfqpoint{3.437845in}{3.477213in}}%
\pgfpathlineto{\pgfqpoint{3.442893in}{3.405214in}}%
\pgfpathlineto{\pgfqpoint{3.446572in}{3.379293in}}%
\pgfpathlineto{\pgfqpoint{3.448882in}{3.374670in}}%
\pgfpathlineto{\pgfqpoint{3.449994in}{3.375637in}}%
\pgfpathlineto{\pgfqpoint{3.451706in}{3.381132in}}%
\pgfpathlineto{\pgfqpoint{3.454273in}{3.398309in}}%
\pgfpathlineto{\pgfqpoint{3.457866in}{3.439649in}}%
\pgfpathlineto{\pgfqpoint{3.462658in}{3.523633in}}%
\pgfpathlineto{\pgfqpoint{3.469075in}{3.679738in}}%
\pgfpathlineto{\pgfqpoint{3.478487in}{3.970061in}}%
\pgfpathlineto{\pgfqpoint{3.499877in}{4.643029in}}%
\pgfpathlineto{\pgfqpoint{3.507321in}{4.804195in}}%
\pgfpathlineto{\pgfqpoint{3.513054in}{4.885476in}}%
\pgfpathlineto{\pgfqpoint{3.517418in}{4.919941in}}%
\pgfpathlineto{\pgfqpoint{3.520412in}{4.929539in}}%
\pgfpathlineto{\pgfqpoint{3.522038in}{4.929971in}}%
\pgfpathlineto{\pgfqpoint{3.522124in}{4.929901in}}%
\pgfpathlineto{\pgfqpoint{3.523664in}{4.927075in}}%
\pgfpathlineto{\pgfqpoint{3.525974in}{4.917329in}}%
\pgfpathlineto{\pgfqpoint{3.529225in}{4.892743in}}%
\pgfpathlineto{\pgfqpoint{3.533589in}{4.840885in}}%
\pgfpathlineto{\pgfqpoint{3.539322in}{4.743316in}}%
\pgfpathlineto{\pgfqpoint{3.546937in}{4.571947in}}%
\pgfpathlineto{\pgfqpoint{3.559001in}{4.242312in}}%
\pgfpathlineto{\pgfqpoint{3.577311in}{3.749847in}}%
\pgfpathlineto{\pgfqpoint{3.586039in}{3.573106in}}%
\pgfpathlineto{\pgfqpoint{3.592798in}{3.476715in}}%
\pgfpathlineto{\pgfqpoint{3.598103in}{3.428864in}}%
\pgfpathlineto{\pgfqpoint{3.602124in}{3.409500in}}%
\pgfpathlineto{\pgfqpoint{3.604862in}{3.404684in}}%
\pgfpathlineto{\pgfqpoint{3.606488in}{3.405001in}}%
\pgfpathlineto{\pgfqpoint{3.608371in}{3.408283in}}%
\pgfpathlineto{\pgfqpoint{3.611023in}{3.418118in}}%
\pgfpathlineto{\pgfqpoint{3.614702in}{3.441517in}}%
\pgfpathlineto{\pgfqpoint{3.619494in}{3.487994in}}%
\pgfpathlineto{\pgfqpoint{3.625654in}{3.571635in}}%
\pgfpathlineto{\pgfqpoint{3.633697in}{3.714087in}}%
\pgfpathlineto{\pgfqpoint{3.645334in}{3.964365in}}%
\pgfpathlineto{\pgfqpoint{3.672457in}{4.558556in}}%
\pgfpathlineto{\pgfqpoint{3.682040in}{4.715588in}}%
\pgfpathlineto{\pgfqpoint{3.689740in}{4.808281in}}%
\pgfpathlineto{\pgfqpoint{3.695987in}{4.859208in}}%
\pgfpathlineto{\pgfqpoint{3.700949in}{4.883662in}}%
\pgfpathlineto{\pgfqpoint{3.704714in}{4.892753in}}%
\pgfpathlineto{\pgfqpoint{3.707281in}{4.894334in}}%
\pgfpathlineto{\pgfqpoint{3.709249in}{4.893053in}}%
\pgfpathlineto{\pgfqpoint{3.711816in}{4.888194in}}%
\pgfpathlineto{\pgfqpoint{3.715238in}{4.876254in}}%
\pgfpathlineto{\pgfqpoint{3.719687in}{4.851807in}}%
\pgfpathlineto{\pgfqpoint{3.725334in}{4.807268in}}%
\pgfpathlineto{\pgfqpoint{3.732436in}{4.732200in}}%
\pgfpathlineto{\pgfqpoint{3.741591in}{4.609796in}}%
\pgfpathlineto{\pgfqpoint{3.754426in}{4.405181in}}%
\pgfpathlineto{\pgfqpoint{3.791389in}{3.799090in}}%
\pgfpathlineto{\pgfqpoint{3.802683in}{3.657879in}}%
\pgfpathlineto{\pgfqpoint{3.812009in}{3.567271in}}%
\pgfpathlineto{\pgfqpoint{3.819881in}{3.510714in}}%
\pgfpathlineto{\pgfqpoint{3.826555in}{3.477415in}}%
\pgfpathlineto{\pgfqpoint{3.832116in}{3.459884in}}%
\pgfpathlineto{\pgfqpoint{3.836480in}{3.452504in}}%
\pgfpathlineto{\pgfqpoint{3.839817in}{3.450544in}}%
\pgfpathlineto{\pgfqpoint{3.842641in}{3.451314in}}%
\pgfpathlineto{\pgfqpoint{3.845806in}{3.454750in}}%
\pgfpathlineto{\pgfqpoint{3.849742in}{3.462686in}}%
\pgfpathlineto{\pgfqpoint{3.854705in}{3.478207in}}%
\pgfpathlineto{\pgfqpoint{3.860865in}{3.505448in}}%
\pgfpathlineto{\pgfqpoint{3.868395in}{3.549558in}}%
\pgfpathlineto{\pgfqpoint{3.877721in}{3.618326in}}%
\pgfpathlineto{\pgfqpoint{3.889614in}{3.723800in}}%
\pgfpathlineto{\pgfqpoint{3.906556in}{3.896022in}}%
\pgfpathlineto{\pgfqpoint{3.951733in}{4.363960in}}%
\pgfpathlineto{\pgfqpoint{3.967476in}{4.499587in}}%
\pgfpathlineto{\pgfqpoint{3.981081in}{4.598394in}}%
\pgfpathlineto{\pgfqpoint{3.993230in}{4.671082in}}%
\pgfpathlineto{\pgfqpoint{4.004182in}{4.723860in}}%
\pgfpathlineto{\pgfqpoint{4.014193in}{4.761755in}}%
\pgfpathlineto{\pgfqpoint{4.023263in}{4.787904in}}%
\pgfpathlineto{\pgfqpoint{4.031477in}{4.805231in}}%
\pgfpathlineto{\pgfqpoint{4.038835in}{4.815929in}}%
\pgfpathlineto{\pgfqpoint{4.045509in}{4.821940in}}%
\pgfpathlineto{\pgfqpoint{4.051584in}{4.824554in}}%
\pgfpathlineto{\pgfqpoint{4.057402in}{4.824674in}}%
\pgfpathlineto{\pgfqpoint{4.063392in}{4.822528in}}%
\pgfpathlineto{\pgfqpoint{4.069980in}{4.817711in}}%
\pgfpathlineto{\pgfqpoint{4.077595in}{4.809228in}}%
\pgfpathlineto{\pgfqpoint{4.086494in}{4.795796in}}%
\pgfpathlineto{\pgfqpoint{4.097103in}{4.775526in}}%
\pgfpathlineto{\pgfqpoint{4.110023in}{4.745766in}}%
\pgfpathlineto{\pgfqpoint{4.126451in}{4.701959in}}%
\pgfpathlineto{\pgfqpoint{4.149981in}{4.632231in}}%
\pgfpathlineto{\pgfqpoint{4.214752in}{4.437427in}}%
\pgfpathlineto{\pgfqpoint{4.238795in}{4.373711in}}%
\pgfpathlineto{\pgfqpoint{4.260442in}{4.322470in}}%
\pgfpathlineto{\pgfqpoint{4.280635in}{4.280229in}}%
\pgfpathlineto{\pgfqpoint{4.299972in}{4.244811in}}%
\pgfpathlineto{\pgfqpoint{4.318625in}{4.215170in}}%
\pgfpathlineto{\pgfqpoint{4.336849in}{4.190280in}}%
\pgfpathlineto{\pgfqpoint{4.354817in}{4.169420in}}%
\pgfpathlineto{\pgfqpoint{4.372700in}{4.152005in}}%
\pgfpathlineto{\pgfqpoint{4.390668in}{4.137573in}}%
\pgfpathlineto{\pgfqpoint{4.408893in}{4.125763in}}%
\pgfpathlineto{\pgfqpoint{4.427631in}{4.116258in}}%
\pgfpathlineto{\pgfqpoint{4.447225in}{4.108816in}}%
\pgfpathlineto{\pgfqpoint{4.468188in}{4.103256in}}%
\pgfpathlineto{\pgfqpoint{4.491204in}{4.099488in}}%
\pgfpathlineto{\pgfqpoint{4.517899in}{4.097441in}}%
\pgfpathlineto{\pgfqpoint{4.552809in}{4.097154in}}%
\pgfpathlineto{\pgfqpoint{4.635633in}{4.097270in}}%
\pgfpathlineto{\pgfqpoint{4.663612in}{4.094545in}}%
\pgfpathlineto{\pgfqpoint{4.687912in}{4.089978in}}%
\pgfpathlineto{\pgfqpoint{4.710244in}{4.083531in}}%
\pgfpathlineto{\pgfqpoint{4.731378in}{4.075111in}}%
\pgfpathlineto{\pgfqpoint{4.751742in}{4.064598in}}%
\pgfpathlineto{\pgfqpoint{4.771763in}{4.051745in}}%
\pgfpathlineto{\pgfqpoint{4.791785in}{4.036214in}}%
\pgfpathlineto{\pgfqpoint{4.811978in}{4.017696in}}%
\pgfpathlineto{\pgfqpoint{4.832684in}{3.995647in}}%
\pgfpathlineto{\pgfqpoint{4.854331in}{3.969286in}}%
\pgfpathlineto{\pgfqpoint{4.877604in}{3.937338in}}%
\pgfpathlineto{\pgfqpoint{4.903872in}{3.897344in}}%
\pgfpathlineto{\pgfqpoint{4.939466in}{3.838669in}}%
\pgfpathlineto{\pgfqpoint{4.983445in}{3.766832in}}%
\pgfpathlineto{\pgfqpoint{5.003124in}{3.738978in}}%
\pgfpathlineto{\pgfqpoint{5.018183in}{3.721298in}}%
\pgfpathlineto{\pgfqpoint{5.030590in}{3.709975in}}%
\pgfpathlineto{\pgfqpoint{5.041199in}{3.703174in}}%
\pgfpathlineto{\pgfqpoint{5.050440in}{3.699791in}}%
\pgfpathlineto{\pgfqpoint{5.058825in}{3.699042in}}%
\pgfpathlineto{\pgfqpoint{5.066697in}{3.700566in}}%
\pgfpathlineto{\pgfqpoint{5.074483in}{3.704388in}}%
\pgfpathlineto{\pgfqpoint{5.082441in}{3.710864in}}%
\pgfpathlineto{\pgfqpoint{5.090740in}{3.720588in}}%
\pgfpathlineto{\pgfqpoint{5.099639in}{3.734603in}}%
\pgfpathlineto{\pgfqpoint{5.109136in}{3.753896in}}%
\pgfpathlineto{\pgfqpoint{5.119404in}{3.780029in}}%
\pgfpathlineto{\pgfqpoint{5.130527in}{3.814732in}}%
\pgfpathlineto{\pgfqpoint{5.142591in}{3.859974in}}%
\pgfpathlineto{\pgfqpoint{5.155939in}{3.919057in}}%
\pgfpathlineto{\pgfqpoint{5.171083in}{3.996716in}}%
\pgfpathlineto{\pgfqpoint{5.189479in}{4.103480in}}%
\pgfpathlineto{\pgfqpoint{5.249630in}{4.463075in}}%
\pgfpathlineto{\pgfqpoint{5.260924in}{4.512161in}}%
\pgfpathlineto{\pgfqpoint{5.269993in}{4.542447in}}%
\pgfpathlineto{\pgfqpoint{5.277352in}{4.559955in}}%
\pgfpathlineto{\pgfqpoint{5.283341in}{4.569003in}}%
\pgfpathlineto{\pgfqpoint{5.288133in}{4.572636in}}%
\pgfpathlineto{\pgfqpoint{5.292154in}{4.573090in}}%
\pgfpathlineto{\pgfqpoint{5.296004in}{4.571234in}}%
\pgfpathlineto{\pgfqpoint{5.300283in}{4.566476in}}%
\pgfpathlineto{\pgfqpoint{5.305245in}{4.557335in}}%
\pgfpathlineto{\pgfqpoint{5.311063in}{4.541602in}}%
\pgfpathlineto{\pgfqpoint{5.317908in}{4.516165in}}%
\pgfpathlineto{\pgfqpoint{5.325866in}{4.477384in}}%
\pgfpathlineto{\pgfqpoint{5.335192in}{4.420066in}}%
\pgfpathlineto{\pgfqpoint{5.346401in}{4.336368in}}%
\pgfpathlineto{\pgfqpoint{5.361374in}{4.206223in}}%
\pgfpathlineto{\pgfqpoint{5.394572in}{3.912808in}}%
\pgfpathlineto{\pgfqpoint{5.404754in}{3.844575in}}%
\pgfpathlineto{\pgfqpoint{5.412626in}{3.805498in}}%
\pgfpathlineto{\pgfqpoint{5.418872in}{3.784665in}}%
\pgfpathlineto{\pgfqpoint{5.423664in}{3.775434in}}%
\pgfpathlineto{\pgfqpoint{5.427257in}{3.772583in}}%
\pgfpathlineto{\pgfqpoint{5.430081in}{3.772869in}}%
\pgfpathlineto{\pgfqpoint{5.433076in}{3.775646in}}%
\pgfpathlineto{\pgfqpoint{5.436669in}{3.782376in}}%
\pgfpathlineto{\pgfqpoint{5.441118in}{3.795854in}}%
\pgfpathlineto{\pgfqpoint{5.446594in}{3.820164in}}%
\pgfpathlineto{\pgfqpoint{5.453268in}{3.860857in}}%
\pgfpathlineto{\pgfqpoint{5.461397in}{3.925373in}}%
\pgfpathlineto{\pgfqpoint{5.471750in}{4.026864in}}%
\pgfpathlineto{\pgfqpoint{5.490060in}{4.233263in}}%
\pgfpathlineto{\pgfqpoint{5.504948in}{4.389513in}}%
\pgfpathlineto{\pgfqpoint{5.513932in}{4.462062in}}%
\pgfpathlineto{\pgfqpoint{5.520691in}{4.500350in}}%
\pgfpathlineto{\pgfqpoint{5.525911in}{4.518558in}}%
\pgfpathlineto{\pgfqpoint{5.529761in}{4.525085in}}%
\pgfpathlineto{\pgfqpoint{5.532414in}{4.526026in}}%
\pgfpathlineto{\pgfqpoint{5.534724in}{4.524436in}}%
\pgfpathlineto{\pgfqpoint{5.537547in}{4.519422in}}%
\pgfpathlineto{\pgfqpoint{5.541226in}{4.507823in}}%
\pgfpathlineto{\pgfqpoint{5.545932in}{4.484769in}}%
\pgfpathlineto{\pgfqpoint{5.551751in}{4.444112in}}%
\pgfpathlineto{\pgfqpoint{5.559109in}{4.375486in}}%
\pgfpathlineto{\pgfqpoint{5.568949in}{4.260607in}}%
\pgfpathlineto{\pgfqpoint{5.598896in}{3.895885in}}%
\pgfpathlineto{\pgfqpoint{5.605826in}{3.843593in}}%
\pgfpathlineto{\pgfqpoint{5.611045in}{3.819148in}}%
\pgfpathlineto{\pgfqpoint{5.614810in}{3.810391in}}%
\pgfpathlineto{\pgfqpoint{5.617377in}{3.808893in}}%
\pgfpathlineto{\pgfqpoint{5.619345in}{3.810250in}}%
\pgfpathlineto{\pgfqpoint{5.621826in}{3.815077in}}%
\pgfpathlineto{\pgfqpoint{5.625163in}{3.827028in}}%
\pgfpathlineto{\pgfqpoint{5.629527in}{3.851887in}}%
\pgfpathlineto{\pgfqpoint{5.635088in}{3.897803in}}%
\pgfpathlineto{\pgfqpoint{5.642276in}{3.977504in}}%
\pgfpathlineto{\pgfqpoint{5.652800in}{4.122296in}}%
\pgfpathlineto{\pgfqpoint{5.671538in}{4.380373in}}%
\pgfpathlineto{\pgfqpoint{5.678811in}{4.449980in}}%
\pgfpathlineto{\pgfqpoint{5.684116in}{4.481941in}}%
\pgfpathlineto{\pgfqpoint{5.687966in}{4.493684in}}%
\pgfpathlineto{\pgfqpoint{5.690447in}{4.495841in}}%
\pgfpathlineto{\pgfqpoint{5.692244in}{4.494699in}}%
\pgfpathlineto{\pgfqpoint{5.694554in}{4.489875in}}%
\pgfpathlineto{\pgfqpoint{5.697720in}{4.477177in}}%
\pgfpathlineto{\pgfqpoint{5.701913in}{4.449847in}}%
\pgfpathlineto{\pgfqpoint{5.707303in}{4.398448in}}%
\pgfpathlineto{\pgfqpoint{5.714490in}{4.306418in}}%
\pgfpathlineto{\pgfqpoint{5.727154in}{4.109379in}}%
\pgfpathlineto{\pgfqpoint{5.738876in}{3.941040in}}%
\pgfpathlineto{\pgfqpoint{5.745550in}{3.874173in}}%
\pgfpathlineto{\pgfqpoint{5.750341in}{3.845360in}}%
\pgfpathlineto{\pgfqpoint{5.753764in}{3.835972in}}%
\pgfpathlineto{\pgfqpoint{5.755817in}{3.835040in}}%
\pgfpathlineto{\pgfqpoint{5.757528in}{3.836998in}}%
\pgfpathlineto{\pgfqpoint{5.759924in}{3.843916in}}%
\pgfpathlineto{\pgfqpoint{5.763261in}{3.861546in}}%
\pgfpathlineto{\pgfqpoint{5.767625in}{3.897957in}}%
\pgfpathlineto{\pgfqpoint{5.773443in}{3.967247in}}%
\pgfpathlineto{\pgfqpoint{5.781914in}{4.098402in}}%
\pgfpathlineto{\pgfqpoint{5.799967in}{4.383201in}}%
\pgfpathlineto{\pgfqpoint{5.805957in}{4.442578in}}%
\pgfpathlineto{\pgfqpoint{5.810235in}{4.466324in}}%
\pgfpathlineto{\pgfqpoint{5.813144in}{4.472555in}}%
\pgfpathlineto{\pgfqpoint{5.814770in}{4.472409in}}%
\pgfpathlineto{\pgfqpoint{5.816481in}{4.469420in}}%
\pgfpathlineto{\pgfqpoint{5.819048in}{4.459515in}}%
\pgfpathlineto{\pgfqpoint{5.822556in}{4.435742in}}%
\pgfpathlineto{\pgfqpoint{5.827262in}{4.386693in}}%
\pgfpathlineto{\pgfqpoint{5.833679in}{4.293902in}}%
\pgfpathlineto{\pgfqpoint{5.847626in}{4.048544in}}%
\pgfpathlineto{\pgfqpoint{5.855925in}{3.926780in}}%
\pgfpathlineto{\pgfqpoint{5.861316in}{3.876455in}}%
\pgfpathlineto{\pgfqpoint{5.865080in}{3.858759in}}%
\pgfpathlineto{\pgfqpoint{5.867391in}{3.855642in}}%
\pgfpathlineto{\pgfqpoint{5.868760in}{3.856644in}}%
\pgfpathlineto{\pgfqpoint{5.870642in}{3.861491in}}%
\pgfpathlineto{\pgfqpoint{5.873380in}{3.875620in}}%
\pgfpathlineto{\pgfqpoint{5.877145in}{3.908135in}}%
\pgfpathlineto{\pgfqpoint{5.882278in}{3.974058in}}%
\pgfpathlineto{\pgfqpoint{5.889979in}{4.105965in}}%
\pgfpathlineto{\pgfqpoint{5.905466in}{4.374507in}}%
\pgfpathlineto{\pgfqpoint{5.910856in}{4.430632in}}%
\pgfpathlineto{\pgfqpoint{5.914621in}{4.450471in}}%
\pgfpathlineto{\pgfqpoint{5.916931in}{4.453960in}}%
\pgfpathlineto{\pgfqpoint{5.918215in}{4.452967in}}%
\pgfpathlineto{\pgfqpoint{5.920011in}{4.448051in}}%
\pgfpathlineto{\pgfqpoint{5.922664in}{4.433389in}}%
\pgfpathlineto{\pgfqpoint{5.926429in}{4.398183in}}%
\pgfpathlineto{\pgfqpoint{5.931562in}{4.326550in}}%
\pgfpathlineto{\pgfqpoint{5.939691in}{4.176873in}}%
\pgfpathlineto{\pgfqpoint{5.952012in}{3.955621in}}%
\pgfpathlineto{\pgfqpoint{5.957317in}{3.896666in}}%
\pgfpathlineto{\pgfqpoint{5.960996in}{3.876284in}}%
\pgfpathlineto{\pgfqpoint{5.963220in}{3.872982in}}%
\pgfpathlineto{\pgfqpoint{5.964504in}{3.874252in}}%
\pgfpathlineto{\pgfqpoint{5.966386in}{3.880319in}}%
\pgfpathlineto{\pgfqpoint{5.969124in}{3.897886in}}%
\pgfpathlineto{\pgfqpoint{5.972975in}{3.938995in}}%
\pgfpathlineto{\pgfqpoint{5.978365in}{4.023361in}}%
\pgfpathlineto{\pgfqpoint{5.989231in}{4.238788in}}%
\pgfpathlineto{\pgfqpoint{5.997189in}{4.373356in}}%
\pgfpathlineto{\pgfqpoint{6.001980in}{4.422009in}}%
\pgfpathlineto{\pgfqpoint{6.005232in}{4.436571in}}%
\pgfpathlineto{\pgfqpoint{6.006857in}{4.437809in}}%
\pgfpathlineto{\pgfqpoint{6.007028in}{4.437701in}}%
\pgfpathlineto{\pgfqpoint{6.008397in}{4.435203in}}%
\pgfpathlineto{\pgfqpoint{6.010537in}{4.425525in}}%
\pgfpathlineto{\pgfqpoint{6.013617in}{4.399659in}}%
\pgfpathlineto{\pgfqpoint{6.017980in}{4.341396in}}%
\pgfpathlineto{\pgfqpoint{6.024483in}{4.220400in}}%
\pgfpathlineto{\pgfqpoint{6.038772in}{3.949104in}}%
\pgfpathlineto{\pgfqpoint{6.043393in}{3.901947in}}%
\pgfpathlineto{\pgfqpoint{6.046473in}{3.888842in}}%
\pgfpathlineto{\pgfqpoint{6.047842in}{3.888078in}}%
\pgfpathlineto{\pgfqpoint{6.048098in}{3.888287in}}%
\pgfpathlineto{\pgfqpoint{6.049467in}{3.891274in}}%
\pgfpathlineto{\pgfqpoint{6.051607in}{3.902199in}}%
\pgfpathlineto{\pgfqpoint{6.054772in}{3.931766in}}%
\pgfpathlineto{\pgfqpoint{6.059222in}{3.997027in}}%
\pgfpathlineto{\pgfqpoint{6.066409in}{4.140924in}}%
\pgfpathlineto{\pgfqpoint{6.077446in}{4.355636in}}%
\pgfpathlineto{\pgfqpoint{6.082152in}{4.408127in}}%
\pgfpathlineto{\pgfqpoint{6.085233in}{4.422723in}}%
\pgfpathlineto{\pgfqpoint{6.086602in}{4.423761in}}%
\pgfpathlineto{\pgfqpoint{6.086858in}{4.423577in}}%
\pgfpathlineto{\pgfqpoint{6.088142in}{4.420856in}}%
\pgfpathlineto{\pgfqpoint{6.090195in}{4.410340in}}%
\pgfpathlineto{\pgfqpoint{6.093275in}{4.380956in}}%
\pgfpathlineto{\pgfqpoint{6.097639in}{4.314915in}}%
\pgfpathlineto{\pgfqpoint{6.104826in}{4.166310in}}%
\pgfpathlineto{\pgfqpoint{6.114923in}{3.966712in}}%
\pgfpathlineto{\pgfqpoint{6.119458in}{3.915663in}}%
\pgfpathlineto{\pgfqpoint{6.122452in}{3.901973in}}%
\pgfpathlineto{\pgfqpoint{6.123650in}{3.901323in}}%
\pgfpathlineto{\pgfqpoint{6.123907in}{3.901547in}}%
\pgfpathlineto{\pgfqpoint{6.125276in}{3.904911in}}%
\pgfpathlineto{\pgfqpoint{6.127415in}{3.917375in}}%
\pgfpathlineto{\pgfqpoint{6.130581in}{3.951087in}}%
\pgfpathlineto{\pgfqpoint{6.135201in}{4.027912in}}%
\pgfpathlineto{\pgfqpoint{6.144442in}{4.229890in}}%
\pgfpathlineto{\pgfqpoint{6.151543in}{4.360410in}}%
\pgfpathlineto{\pgfqpoint{6.155650in}{4.401816in}}%
\pgfpathlineto{\pgfqpoint{6.158217in}{4.411069in}}%
\pgfpathlineto{\pgfqpoint{6.159244in}{4.410968in}}%
\pgfpathlineto{\pgfqpoint{6.159415in}{4.410738in}}%
\pgfpathlineto{\pgfqpoint{6.160870in}{4.406335in}}%
\pgfpathlineto{\pgfqpoint{6.163180in}{4.390503in}}%
\pgfpathlineto{\pgfqpoint{6.166517in}{4.349822in}}%
\pgfpathlineto{\pgfqpoint{6.171480in}{4.258195in}}%
\pgfpathlineto{\pgfqpoint{6.187480in}{3.941011in}}%
\pgfpathlineto{\pgfqpoint{6.190902in}{3.916166in}}%
\pgfpathlineto{\pgfqpoint{6.192699in}{3.913133in}}%
\pgfpathlineto{\pgfqpoint{6.192870in}{3.913215in}}%
\pgfpathlineto{\pgfqpoint{6.193982in}{3.915317in}}%
\pgfpathlineto{\pgfqpoint{6.195865in}{3.925021in}}%
\pgfpathlineto{\pgfqpoint{6.198688in}{3.953426in}}%
\pgfpathlineto{\pgfqpoint{6.202795in}{4.020378in}}%
\pgfpathlineto{\pgfqpoint{6.210239in}{4.186429in}}%
\pgfpathlineto{\pgfqpoint{6.218197in}{4.346380in}}%
\pgfpathlineto{\pgfqpoint{6.222304in}{4.390745in}}%
\pgfpathlineto{\pgfqpoint{6.224785in}{4.399836in}}%
\pgfpathlineto{\pgfqpoint{6.225726in}{4.399565in}}%
\pgfpathlineto{\pgfqpoint{6.225897in}{4.399294in}}%
\pgfpathlineto{\pgfqpoint{6.227352in}{4.394253in}}%
\pgfpathlineto{\pgfqpoint{6.229662in}{4.376407in}}%
\pgfpathlineto{\pgfqpoint{6.233084in}{4.329675in}}%
\pgfpathlineto{\pgfqpoint{6.238389in}{4.221498in}}%
\pgfpathlineto{\pgfqpoint{6.250625in}{3.964725in}}%
\pgfpathlineto{\pgfqpoint{6.254304in}{3.929767in}}%
\pgfpathlineto{\pgfqpoint{6.256529in}{3.923916in}}%
\pgfpathlineto{\pgfqpoint{6.257470in}{3.925075in}}%
\pgfpathlineto{\pgfqpoint{6.259010in}{3.931622in}}%
\pgfpathlineto{\pgfqpoint{6.261491in}{3.953906in}}%
\pgfpathlineto{\pgfqpoint{6.265170in}{4.010641in}}%
\pgfpathlineto{\pgfqpoint{6.271245in}{4.145194in}}%
\pgfpathlineto{\pgfqpoint{6.280401in}{4.339710in}}%
\pgfpathlineto{\pgfqpoint{6.284336in}{4.382104in}}%
\pgfpathlineto{\pgfqpoint{6.286647in}{4.389628in}}%
\pgfpathlineto{\pgfqpoint{6.287588in}{4.388787in}}%
\pgfpathlineto{\pgfqpoint{6.289042in}{4.383035in}}%
\pgfpathlineto{\pgfqpoint{6.291353in}{4.363125in}}%
\pgfpathlineto{\pgfqpoint{6.294861in}{4.310217in}}%
\pgfpathlineto{\pgfqpoint{6.300593in}{4.183606in}}%
\pgfpathlineto{\pgfqpoint{6.310005in}{3.980801in}}%
\pgfpathlineto{\pgfqpoint{6.313770in}{3.940695in}}%
\pgfpathlineto{\pgfqpoint{6.315994in}{3.933817in}}%
\pgfpathlineto{\pgfqpoint{6.316936in}{3.934885in}}%
\pgfpathlineto{\pgfqpoint{6.318476in}{3.941726in}}%
\pgfpathlineto{\pgfqpoint{6.320872in}{3.964478in}}%
\pgfpathlineto{\pgfqpoint{6.324465in}{4.022900in}}%
\pgfpathlineto{\pgfqpoint{6.330882in}{4.171501in}}%
\pgfpathlineto{\pgfqpoint{6.338583in}{4.334710in}}%
\pgfpathlineto{\pgfqpoint{6.342262in}{4.373818in}}%
\pgfpathlineto{\pgfqpoint{6.344401in}{4.380142in}}%
\pgfpathlineto{\pgfqpoint{6.345257in}{4.379105in}}%
\pgfpathlineto{\pgfqpoint{6.346711in}{4.372674in}}%
\pgfpathlineto{\pgfqpoint{6.349107in}{4.349738in}}%
\pgfpathlineto{\pgfqpoint{6.352701in}{4.290119in}}%
\pgfpathlineto{\pgfqpoint{6.359289in}{4.134835in}}%
\pgfpathlineto{\pgfqpoint{6.366476in}{3.984203in}}%
\pgfpathlineto{\pgfqpoint{6.370070in}{3.947799in}}%
\pgfpathlineto{\pgfqpoint{6.371952in}{3.942959in}}%
\pgfpathlineto{\pgfqpoint{6.372123in}{3.943028in}}%
\pgfpathlineto{\pgfqpoint{6.373150in}{3.945226in}}%
\pgfpathlineto{\pgfqpoint{6.374861in}{3.955604in}}%
\pgfpathlineto{\pgfqpoint{6.377599in}{3.988629in}}%
\pgfpathlineto{\pgfqpoint{6.381792in}{4.070875in}}%
\pgfpathlineto{\pgfqpoint{6.395140in}{4.353399in}}%
\pgfpathlineto{\pgfqpoint{6.397963in}{4.370694in}}%
\pgfpathlineto{\pgfqpoint{6.398819in}{4.371277in}}%
\pgfpathlineto{\pgfqpoint{6.399247in}{4.370740in}}%
\pgfpathlineto{\pgfqpoint{6.400530in}{4.365837in}}%
\pgfpathlineto{\pgfqpoint{6.402669in}{4.347003in}}%
\pgfpathlineto{\pgfqpoint{6.405921in}{4.295592in}}%
\pgfpathlineto{\pgfqpoint{6.411568in}{4.163940in}}%
\pgfpathlineto{\pgfqpoint{6.419440in}{3.991581in}}%
\pgfpathlineto{\pgfqpoint{6.422948in}{3.955703in}}%
\pgfpathlineto{\pgfqpoint{6.424659in}{3.951452in}}%
\pgfpathlineto{\pgfqpoint{6.424830in}{3.951528in}}%
\pgfpathlineto{\pgfqpoint{6.425771in}{3.953576in}}%
\pgfpathlineto{\pgfqpoint{6.427482in}{3.964277in}}%
\pgfpathlineto{\pgfqpoint{6.430135in}{3.997541in}}%
\pgfpathlineto{\pgfqpoint{6.434327in}{4.082918in}}%
\pgfpathlineto{\pgfqpoint{6.446306in}{4.342637in}}%
\pgfpathlineto{\pgfqpoint{6.449130in}{4.362167in}}%
\pgfpathlineto{\pgfqpoint{6.450071in}{4.363093in}}%
\pgfpathlineto{\pgfqpoint{6.450499in}{4.362569in}}%
\pgfpathlineto{\pgfqpoint{6.451782in}{4.357479in}}%
\pgfpathlineto{\pgfqpoint{6.453921in}{4.337633in}}%
\pgfpathlineto{\pgfqpoint{6.457258in}{4.281743in}}%
\pgfpathlineto{\pgfqpoint{6.463504in}{4.130080in}}%
\pgfpathlineto{\pgfqpoint{6.470007in}{3.992943in}}%
\pgfpathlineto{\pgfqpoint{6.473258in}{3.962229in}}%
\pgfpathlineto{\pgfqpoint{6.474627in}{3.959382in}}%
\pgfpathlineto{\pgfqpoint{6.474970in}{3.959642in}}%
\pgfpathlineto{\pgfqpoint{6.476082in}{3.963168in}}%
\pgfpathlineto{\pgfqpoint{6.477964in}{3.978259in}}%
\pgfpathlineto{\pgfqpoint{6.480959in}{4.023673in}}%
\pgfpathlineto{\pgfqpoint{6.486007in}{4.140847in}}%
\pgfpathlineto{\pgfqpoint{6.494136in}{4.321861in}}%
\pgfpathlineto{\pgfqpoint{6.497387in}{4.352798in}}%
\pgfpathlineto{\pgfqpoint{6.498670in}{4.355459in}}%
\pgfpathlineto{\pgfqpoint{6.499013in}{4.355219in}}%
\pgfpathlineto{\pgfqpoint{6.500039in}{4.352097in}}%
\pgfpathlineto{\pgfqpoint{6.501922in}{4.337241in}}%
\pgfpathlineto{\pgfqpoint{6.504831in}{4.293189in}}%
\pgfpathlineto{\pgfqpoint{6.509793in}{4.177656in}}%
\pgfpathlineto{\pgfqpoint{6.517751in}{3.999665in}}%
\pgfpathlineto{\pgfqpoint{6.520917in}{3.969437in}}%
\pgfpathlineto{\pgfqpoint{6.522200in}{3.966820in}}%
\pgfpathlineto{\pgfqpoint{6.522542in}{3.967099in}}%
\pgfpathlineto{\pgfqpoint{6.523655in}{3.970843in}}%
\pgfpathlineto{\pgfqpoint{6.525537in}{3.986804in}}%
\pgfpathlineto{\pgfqpoint{6.528532in}{4.034528in}}%
\pgfpathlineto{\pgfqpoint{6.533922in}{4.164222in}}%
\pgfpathlineto{\pgfqpoint{6.540853in}{4.316457in}}%
\pgfpathlineto{\pgfqpoint{6.544018in}{4.346092in}}%
\pgfpathlineto{\pgfqpoint{6.545216in}{4.348230in}}%
\pgfpathlineto{\pgfqpoint{6.545559in}{4.347890in}}%
\pgfpathlineto{\pgfqpoint{6.546671in}{4.343876in}}%
\pgfpathlineto{\pgfqpoint{6.548639in}{4.326194in}}%
\pgfpathlineto{\pgfqpoint{6.551719in}{4.274553in}}%
\pgfpathlineto{\pgfqpoint{6.557708in}{4.126596in}}%
\pgfpathlineto{\pgfqpoint{6.563698in}{4.001066in}}%
\pgfpathlineto{\pgfqpoint{6.566692in}{3.975287in}}%
\pgfpathlineto{\pgfqpoint{6.567634in}{3.973817in}}%
\pgfpathlineto{\pgfqpoint{6.568061in}{3.974233in}}%
\pgfpathlineto{\pgfqpoint{6.569174in}{3.978470in}}%
\pgfpathlineto{\pgfqpoint{6.571142in}{3.996782in}}%
\pgfpathlineto{\pgfqpoint{6.574222in}{4.049757in}}%
\pgfpathlineto{\pgfqpoint{6.580639in}{4.209949in}}%
\pgfpathlineto{\pgfqpoint{6.586030in}{4.318032in}}%
\pgfpathlineto{\pgfqpoint{6.588853in}{4.340436in}}%
\pgfpathlineto{\pgfqpoint{6.589709in}{4.341421in}}%
\pgfpathlineto{\pgfqpoint{6.590137in}{4.340873in}}%
\pgfpathlineto{\pgfqpoint{6.591334in}{4.335672in}}%
\pgfpathlineto{\pgfqpoint{6.593388in}{4.314665in}}%
\pgfpathlineto{\pgfqpoint{6.596639in}{4.254754in}}%
\pgfpathlineto{\pgfqpoint{6.609730in}{3.984239in}}%
\pgfpathlineto{\pgfqpoint{6.611185in}{3.980428in}}%
\pgfpathlineto{\pgfqpoint{6.611527in}{3.980720in}}%
\pgfpathlineto{\pgfqpoint{6.612554in}{3.984311in}}%
\pgfpathlineto{\pgfqpoint{6.614436in}{4.001161in}}%
\pgfpathlineto{\pgfqpoint{6.617431in}{4.052074in}}%
\pgfpathlineto{\pgfqpoint{6.623592in}{4.206237in}}%
\pgfpathlineto{\pgfqpoint{6.628896in}{4.313081in}}%
\pgfpathlineto{\pgfqpoint{6.631634in}{4.334241in}}%
\pgfpathlineto{\pgfqpoint{6.632404in}{4.334979in}}%
\pgfpathlineto{\pgfqpoint{6.632832in}{4.334374in}}%
\pgfpathlineto{\pgfqpoint{6.634030in}{4.328854in}}%
\pgfpathlineto{\pgfqpoint{6.636084in}{4.306823in}}%
\pgfpathlineto{\pgfqpoint{6.639421in}{4.242696in}}%
\pgfpathlineto{\pgfqpoint{6.651228in}{3.992864in}}%
\pgfpathlineto{\pgfqpoint{6.653025in}{3.986686in}}%
\pgfpathlineto{\pgfqpoint{6.653196in}{3.986775in}}%
\pgfpathlineto{\pgfqpoint{6.654052in}{3.988991in}}%
\pgfpathlineto{\pgfqpoint{6.655677in}{4.001177in}}%
\pgfpathlineto{\pgfqpoint{6.658330in}{4.041438in}}%
\pgfpathlineto{\pgfqpoint{6.663036in}{4.154084in}}%
\pgfpathlineto{\pgfqpoint{6.669795in}{4.303999in}}%
\pgfpathlineto{\pgfqpoint{6.672619in}{4.327916in}}%
\pgfpathlineto{\pgfqpoint{6.673389in}{4.328912in}}%
\pgfpathlineto{\pgfqpoint{6.673902in}{4.328219in}}%
\pgfpathlineto{\pgfqpoint{6.675100in}{4.322413in}}%
\pgfpathlineto{\pgfqpoint{6.677154in}{4.299450in}}%
\pgfpathlineto{\pgfqpoint{6.680576in}{4.231249in}}%
\pgfpathlineto{\pgfqpoint{6.691186in}{4.001802in}}%
\pgfpathlineto{\pgfqpoint{6.693325in}{3.992628in}}%
\pgfpathlineto{\pgfqpoint{6.694095in}{3.993997in}}%
\pgfpathlineto{\pgfqpoint{6.695550in}{4.003273in}}%
\pgfpathlineto{\pgfqpoint{6.697945in}{4.036224in}}%
\pgfpathlineto{\pgfqpoint{6.702052in}{4.129801in}}%
\pgfpathlineto{\pgfqpoint{6.709753in}{4.302370in}}%
\pgfpathlineto{\pgfqpoint{6.712405in}{4.322595in}}%
\pgfpathlineto{\pgfqpoint{6.713004in}{4.323106in}}%
\pgfpathlineto{\pgfqpoint{6.713432in}{4.322533in}}%
\pgfpathlineto{\pgfqpoint{6.714630in}{4.316801in}}%
\pgfpathlineto{\pgfqpoint{6.716683in}{4.293507in}}%
\pgfpathlineto{\pgfqpoint{6.720106in}{4.224076in}}%
\pgfpathlineto{\pgfqpoint{6.730031in}{4.008272in}}%
\pgfpathlineto{\pgfqpoint{6.732170in}{3.998280in}}%
\pgfpathlineto{\pgfqpoint{6.732256in}{3.998289in}}%
\pgfpathlineto{\pgfqpoint{6.733026in}{3.999806in}}%
\pgfpathlineto{\pgfqpoint{6.734480in}{4.009583in}}%
\pgfpathlineto{\pgfqpoint{6.736876in}{4.043806in}}%
\pgfpathlineto{\pgfqpoint{6.741154in}{4.143826in}}%
\pgfpathlineto{\pgfqpoint{6.748085in}{4.297330in}}%
\pgfpathlineto{\pgfqpoint{6.750737in}{4.317234in}}%
\pgfpathlineto{\pgfqpoint{6.751251in}{4.317586in}}%
\pgfpathlineto{\pgfqpoint{6.751679in}{4.316990in}}%
\pgfpathlineto{\pgfqpoint{6.752876in}{4.311063in}}%
\pgfpathlineto{\pgfqpoint{6.754930in}{4.287046in}}%
\pgfpathlineto{\pgfqpoint{6.758352in}{4.216058in}}%
\pgfpathlineto{\pgfqpoint{6.767593in}{4.014500in}}%
\pgfpathlineto{\pgfqpoint{6.769818in}{4.003670in}}%
\pgfpathlineto{\pgfqpoint{6.770502in}{4.004761in}}%
\pgfpathlineto{\pgfqpoint{6.771786in}{4.012396in}}%
\pgfpathlineto{\pgfqpoint{6.774010in}{4.041725in}}%
\pgfpathlineto{\pgfqpoint{6.777861in}{4.128049in}}%
\pgfpathlineto{\pgfqpoint{6.785390in}{4.295168in}}%
\pgfpathlineto{\pgfqpoint{6.787871in}{4.312147in}}%
\pgfpathlineto{\pgfqpoint{6.788385in}{4.312232in}}%
\pgfpathlineto{\pgfqpoint{6.788727in}{4.311623in}}%
\pgfpathlineto{\pgfqpoint{6.789925in}{4.305341in}}%
\pgfpathlineto{\pgfqpoint{6.791978in}{4.280389in}}%
\pgfpathlineto{\pgfqpoint{6.795486in}{4.205602in}}%
\pgfpathlineto{\pgfqpoint{6.803872in}{4.021589in}}%
\pgfpathlineto{\pgfqpoint{6.806182in}{4.008825in}}%
\pgfpathlineto{\pgfqpoint{6.806866in}{4.009716in}}%
\pgfpathlineto{\pgfqpoint{6.808150in}{4.017142in}}%
\pgfpathlineto{\pgfqpoint{6.810374in}{4.046565in}}%
\pgfpathlineto{\pgfqpoint{6.814225in}{4.133530in}}%
\pgfpathlineto{\pgfqpoint{6.821326in}{4.290339in}}%
\pgfpathlineto{\pgfqpoint{6.823808in}{4.307175in}}%
\pgfpathlineto{\pgfqpoint{6.824321in}{4.307134in}}%
\pgfpathlineto{\pgfqpoint{6.824578in}{4.306652in}}%
\pgfpathlineto{\pgfqpoint{6.825776in}{4.300377in}}%
\pgfpathlineto{\pgfqpoint{6.827829in}{4.275084in}}%
\pgfpathlineto{\pgfqpoint{6.831423in}{4.197231in}}%
\pgfpathlineto{\pgfqpoint{6.839209in}{4.026971in}}%
\pgfpathlineto{\pgfqpoint{6.841519in}{4.013753in}}%
\pgfpathlineto{\pgfqpoint{6.842204in}{4.014627in}}%
\pgfpathlineto{\pgfqpoint{6.843487in}{4.022171in}}%
\pgfpathlineto{\pgfqpoint{6.845712in}{4.052154in}}%
\pgfpathlineto{\pgfqpoint{6.849648in}{4.142297in}}%
\pgfpathlineto{\pgfqpoint{6.856236in}{4.286019in}}%
\pgfpathlineto{\pgfqpoint{6.858632in}{4.302339in}}%
\pgfpathlineto{\pgfqpoint{6.859059in}{4.302417in}}%
\pgfpathlineto{\pgfqpoint{6.859402in}{4.301848in}}%
\pgfpathlineto{\pgfqpoint{6.860600in}{4.295489in}}%
\pgfpathlineto{\pgfqpoint{6.862653in}{4.269744in}}%
\pgfpathlineto{\pgfqpoint{6.866247in}{4.190986in}}%
\pgfpathlineto{\pgfqpoint{6.873519in}{4.032345in}}%
\pgfpathlineto{\pgfqpoint{6.875830in}{4.018483in}}%
\pgfpathlineto{\pgfqpoint{6.876514in}{4.019274in}}%
\pgfpathlineto{\pgfqpoint{6.877712in}{4.026060in}}%
\pgfpathlineto{\pgfqpoint{6.879851in}{4.054081in}}%
\pgfpathlineto{\pgfqpoint{6.883616in}{4.138959in}}%
\pgfpathlineto{\pgfqpoint{6.890290in}{4.283312in}}%
\pgfpathlineto{\pgfqpoint{6.892600in}{4.297818in}}%
\pgfpathlineto{\pgfqpoint{6.893113in}{4.297548in}}%
\pgfpathlineto{\pgfqpoint{6.893284in}{4.297171in}}%
\pgfpathlineto{\pgfqpoint{6.894482in}{4.290572in}}%
\pgfpathlineto{\pgfqpoint{6.896536in}{4.264168in}}%
\pgfpathlineto{\pgfqpoint{6.900215in}{4.182229in}}%
\pgfpathlineto{\pgfqpoint{6.906974in}{4.036597in}}%
\pgfpathlineto{\pgfqpoint{6.909285in}{4.022988in}}%
\pgfpathlineto{\pgfqpoint{6.909969in}{4.023962in}}%
\pgfpathlineto{\pgfqpoint{6.911253in}{4.031951in}}%
\pgfpathlineto{\pgfqpoint{6.913477in}{4.063260in}}%
\pgfpathlineto{\pgfqpoint{6.917670in}{4.161662in}}%
\pgfpathlineto{\pgfqpoint{6.923317in}{4.279353in}}%
\pgfpathlineto{\pgfqpoint{6.925627in}{4.293422in}}%
\pgfpathlineto{\pgfqpoint{6.926311in}{4.292539in}}%
\pgfpathlineto{\pgfqpoint{6.927509in}{4.285430in}}%
\pgfpathlineto{\pgfqpoint{6.929648in}{4.256475in}}%
\pgfpathlineto{\pgfqpoint{6.933584in}{4.166030in}}%
\pgfpathlineto{\pgfqpoint{6.939488in}{4.041424in}}%
\pgfpathlineto{\pgfqpoint{6.941798in}{4.027339in}}%
\pgfpathlineto{\pgfqpoint{6.942483in}{4.028264in}}%
\pgfpathlineto{\pgfqpoint{6.943766in}{4.036281in}}%
\pgfpathlineto{\pgfqpoint{6.945991in}{4.067909in}}%
\pgfpathlineto{\pgfqpoint{6.950269in}{4.168744in}}%
\pgfpathlineto{\pgfqpoint{6.955574in}{4.276423in}}%
\pgfpathlineto{\pgfqpoint{6.957798in}{4.289178in}}%
\pgfpathlineto{\pgfqpoint{6.958483in}{4.288096in}}%
\pgfpathlineto{\pgfqpoint{6.959766in}{4.279740in}}%
\pgfpathlineto{\pgfqpoint{6.961991in}{4.247486in}}%
\pgfpathlineto{\pgfqpoint{6.966440in}{4.141902in}}%
\pgfpathlineto{\pgfqpoint{6.971403in}{4.043605in}}%
\pgfpathlineto{\pgfqpoint{6.973628in}{4.031505in}}%
\pgfpathlineto{\pgfqpoint{6.974312in}{4.032838in}}%
\pgfpathlineto{\pgfqpoint{6.975681in}{4.042574in}}%
\pgfpathlineto{\pgfqpoint{6.978077in}{4.079853in}}%
\pgfpathlineto{\pgfqpoint{6.983895in}{4.219838in}}%
\pgfpathlineto{\pgfqpoint{6.987660in}{4.279297in}}%
\pgfpathlineto{\pgfqpoint{6.989200in}{4.285079in}}%
\pgfpathlineto{\pgfqpoint{6.989457in}{4.284856in}}%
\pgfpathlineto{\pgfqpoint{6.990398in}{4.281133in}}%
\pgfpathlineto{\pgfqpoint{6.992109in}{4.263204in}}%
\pgfpathlineto{\pgfqpoint{6.995104in}{4.204105in}}%
\pgfpathlineto{\pgfqpoint{7.002975in}{4.041893in}}%
\pgfpathlineto{\pgfqpoint{7.004601in}{4.035530in}}%
\pgfpathlineto{\pgfqpoint{7.004858in}{4.035778in}}%
\pgfpathlineto{\pgfqpoint{7.005799in}{4.039611in}}%
\pgfpathlineto{\pgfqpoint{7.007596in}{4.059060in}}%
\pgfpathlineto{\pgfqpoint{7.010762in}{4.123473in}}%
\pgfpathlineto{\pgfqpoint{7.017949in}{4.272712in}}%
\pgfpathlineto{\pgfqpoint{7.019746in}{4.281139in}}%
\pgfpathlineto{\pgfqpoint{7.019917in}{4.281065in}}%
\pgfpathlineto{\pgfqpoint{7.020687in}{4.278833in}}%
\pgfpathlineto{\pgfqpoint{7.022227in}{4.265338in}}%
\pgfpathlineto{\pgfqpoint{7.024879in}{4.218209in}}%
\pgfpathlineto{\pgfqpoint{7.033864in}{4.041612in}}%
\pgfpathlineto{\pgfqpoint{7.034805in}{4.039395in}}%
\pgfpathlineto{\pgfqpoint{7.035233in}{4.039935in}}%
\pgfpathlineto{\pgfqpoint{7.036345in}{4.045822in}}%
\pgfpathlineto{\pgfqpoint{7.038313in}{4.070966in}}%
\pgfpathlineto{\pgfqpoint{7.042077in}{4.154826in}}%
\pgfpathlineto{\pgfqpoint{7.047639in}{4.266800in}}%
\pgfpathlineto{\pgfqpoint{7.049693in}{4.277331in}}%
\pgfpathlineto{\pgfqpoint{7.049778in}{4.277288in}}%
\pgfpathlineto{\pgfqpoint{7.050548in}{4.275155in}}%
\pgfpathlineto{\pgfqpoint{7.052088in}{4.261735in}}%
\pgfpathlineto{\pgfqpoint{7.054741in}{4.214547in}}%
\pgfpathlineto{\pgfqpoint{7.063383in}{4.045578in}}%
\pgfpathlineto{\pgfqpoint{7.064409in}{4.043139in}}%
\pgfpathlineto{\pgfqpoint{7.064837in}{4.043789in}}%
\pgfpathlineto{\pgfqpoint{7.065949in}{4.050008in}}%
\pgfpathlineto{\pgfqpoint{7.068003in}{4.077305in}}%
\pgfpathlineto{\pgfqpoint{7.072024in}{4.168564in}}%
\pgfpathlineto{\pgfqpoint{7.076987in}{4.264148in}}%
\pgfpathlineto{\pgfqpoint{7.078869in}{4.273671in}}%
\pgfpathlineto{\pgfqpoint{7.079040in}{4.273600in}}%
\pgfpathlineto{\pgfqpoint{7.079811in}{4.271335in}}%
\pgfpathlineto{\pgfqpoint{7.081351in}{4.257570in}}%
\pgfpathlineto{\pgfqpoint{7.084089in}{4.207921in}}%
\pgfpathlineto{\pgfqpoint{7.092132in}{4.050181in}}%
\pgfpathlineto{\pgfqpoint{7.093329in}{4.046740in}}%
\pgfpathlineto{\pgfqpoint{7.093672in}{4.047183in}}%
\pgfpathlineto{\pgfqpoint{7.094698in}{4.052279in}}%
\pgfpathlineto{\pgfqpoint{7.096581in}{4.075313in}}%
\pgfpathlineto{\pgfqpoint{7.100174in}{4.153252in}}%
\pgfpathlineto{\pgfqpoint{7.105650in}{4.260975in}}%
\pgfpathlineto{\pgfqpoint{7.107533in}{4.270128in}}%
\pgfpathlineto{\pgfqpoint{7.107704in}{4.270013in}}%
\pgfpathlineto{\pgfqpoint{7.108560in}{4.267057in}}%
\pgfpathlineto{\pgfqpoint{7.110185in}{4.250974in}}%
\pgfpathlineto{\pgfqpoint{7.113094in}{4.195097in}}%
\pgfpathlineto{\pgfqpoint{7.120110in}{4.056087in}}%
\pgfpathlineto{\pgfqpoint{7.121565in}{4.050207in}}%
\pgfpathlineto{\pgfqpoint{7.121907in}{4.050496in}}%
\pgfpathlineto{\pgfqpoint{7.122848in}{4.054575in}}%
\pgfpathlineto{\pgfqpoint{7.124645in}{4.074934in}}%
\pgfpathlineto{\pgfqpoint{7.127982in}{4.144440in}}%
\pgfpathlineto{\pgfqpoint{7.133800in}{4.258621in}}%
\pgfpathlineto{\pgfqpoint{7.135512in}{4.266713in}}%
\pgfpathlineto{\pgfqpoint{7.135768in}{4.266551in}}%
\pgfpathlineto{\pgfqpoint{7.136624in}{4.263408in}}%
\pgfpathlineto{\pgfqpoint{7.138250in}{4.246921in}}%
\pgfpathlineto{\pgfqpoint{7.141244in}{4.188606in}}%
\pgfpathlineto{\pgfqpoint{7.147833in}{4.059710in}}%
\pgfpathlineto{\pgfqpoint{7.149373in}{4.053572in}}%
\pgfpathlineto{\pgfqpoint{7.149630in}{4.053814in}}%
\pgfpathlineto{\pgfqpoint{7.150571in}{4.057792in}}%
\pgfpathlineto{\pgfqpoint{7.152368in}{4.078061in}}%
\pgfpathlineto{\pgfqpoint{7.155704in}{4.147296in}}%
\pgfpathlineto{\pgfqpoint{7.161266in}{4.255173in}}%
\pgfpathlineto{\pgfqpoint{7.162977in}{4.263408in}}%
\pgfpathlineto{\pgfqpoint{7.163234in}{4.263257in}}%
\pgfpathlineto{\pgfqpoint{7.164090in}{4.260130in}}%
\pgfpathlineto{\pgfqpoint{7.165715in}{4.243604in}}%
\pgfpathlineto{\pgfqpoint{7.168710in}{4.185328in}}%
\pgfpathlineto{\pgfqpoint{7.175042in}{4.062988in}}%
\pgfpathlineto{\pgfqpoint{7.176582in}{4.056821in}}%
\pgfpathlineto{\pgfqpoint{7.176838in}{4.057068in}}%
\pgfpathlineto{\pgfqpoint{7.177780in}{4.061085in}}%
\pgfpathlineto{\pgfqpoint{7.179576in}{4.081485in}}%
\pgfpathlineto{\pgfqpoint{7.182999in}{4.152639in}}%
\pgfpathlineto{\pgfqpoint{7.188304in}{4.252763in}}%
\pgfpathlineto{\pgfqpoint{7.189929in}{4.260213in}}%
\pgfpathlineto{\pgfqpoint{7.190186in}{4.260053in}}%
\pgfpathlineto{\pgfqpoint{7.191042in}{4.256878in}}%
\pgfpathlineto{\pgfqpoint{7.192667in}{4.240215in}}%
\pgfpathlineto{\pgfqpoint{7.195748in}{4.179936in}}%
\pgfpathlineto{\pgfqpoint{7.201737in}{4.066144in}}%
\pgfpathlineto{\pgfqpoint{7.203277in}{4.059965in}}%
\pgfpathlineto{\pgfqpoint{7.203534in}{4.060218in}}%
\pgfpathlineto{\pgfqpoint{7.204475in}{4.064276in}}%
\pgfpathlineto{\pgfqpoint{7.206272in}{4.084795in}}%
\pgfpathlineto{\pgfqpoint{7.209694in}{4.155736in}}%
\pgfpathlineto{\pgfqpoint{7.214743in}{4.249579in}}%
\pgfpathlineto{\pgfqpoint{7.216368in}{4.257120in}}%
\pgfpathlineto{\pgfqpoint{7.216625in}{4.256967in}}%
\pgfpathlineto{\pgfqpoint{7.217481in}{4.253801in}}%
\pgfpathlineto{\pgfqpoint{7.219106in}{4.237112in}}%
\pgfpathlineto{\pgfqpoint{7.222186in}{4.176998in}}%
\pgfpathlineto{\pgfqpoint{7.227919in}{4.069387in}}%
\pgfpathlineto{\pgfqpoint{7.229459in}{4.063004in}}%
\pgfpathlineto{\pgfqpoint{7.229716in}{4.063228in}}%
\pgfpathlineto{\pgfqpoint{7.230657in}{4.067200in}}%
\pgfpathlineto{\pgfqpoint{7.232454in}{4.087609in}}%
\pgfpathlineto{\pgfqpoint{7.235962in}{4.160061in}}%
\pgfpathlineto{\pgfqpoint{7.240839in}{4.247687in}}%
\pgfpathlineto{\pgfqpoint{7.242379in}{4.254126in}}%
\pgfpathlineto{\pgfqpoint{7.242636in}{4.253908in}}%
\pgfpathlineto{\pgfqpoint{7.243577in}{4.249954in}}%
\pgfpathlineto{\pgfqpoint{7.245374in}{4.229562in}}%
\pgfpathlineto{\pgfqpoint{7.248882in}{4.157348in}}%
\pgfpathlineto{\pgfqpoint{7.253673in}{4.072206in}}%
\pgfpathlineto{\pgfqpoint{7.255214in}{4.065954in}}%
\pgfpathlineto{\pgfqpoint{7.255470in}{4.066206in}}%
\pgfpathlineto{\pgfqpoint{7.256411in}{4.070290in}}%
\pgfpathlineto{\pgfqpoint{7.258208in}{4.090908in}}%
\pgfpathlineto{\pgfqpoint{7.261802in}{4.165032in}}%
\pgfpathlineto{\pgfqpoint{7.266422in}{4.245360in}}%
\pgfpathlineto{\pgfqpoint{7.267877in}{4.251227in}}%
\pgfpathlineto{\pgfqpoint{7.268219in}{4.250878in}}%
\pgfpathlineto{\pgfqpoint{7.269246in}{4.245895in}}%
\pgfpathlineto{\pgfqpoint{7.271128in}{4.222701in}}%
\pgfpathlineto{\pgfqpoint{7.275235in}{4.135957in}}%
\pgfpathlineto{\pgfqpoint{7.279171in}{4.073475in}}%
\pgfpathlineto{\pgfqpoint{7.280455in}{4.068803in}}%
\pgfpathlineto{\pgfqpoint{7.280882in}{4.069312in}}%
\pgfpathlineto{\pgfqpoint{7.281909in}{4.074713in}}%
\pgfpathlineto{\pgfqpoint{7.283877in}{4.100002in}}%
\pgfpathlineto{\pgfqpoint{7.289011in}{4.207887in}}%
\pgfpathlineto{\pgfqpoint{7.292177in}{4.246749in}}%
\pgfpathlineto{\pgfqpoint{7.292947in}{4.248417in}}%
\pgfpathlineto{\pgfqpoint{7.293460in}{4.247658in}}%
\pgfpathlineto{\pgfqpoint{7.294658in}{4.240185in}}%
\pgfpathlineto{\pgfqpoint{7.296797in}{4.209430in}}%
\pgfpathlineto{\pgfqpoint{7.305268in}{4.071580in}}%
\pgfpathlineto{\pgfqpoint{7.305867in}{4.072408in}}%
\pgfpathlineto{\pgfqpoint{7.307064in}{4.080057in}}%
\pgfpathlineto{\pgfqpoint{7.309289in}{4.112637in}}%
\pgfpathlineto{\pgfqpoint{7.317332in}{4.245483in}}%
\pgfpathlineto{\pgfqpoint{7.317760in}{4.245621in}}%
\pgfpathlineto{\pgfqpoint{7.318102in}{4.244981in}}%
\pgfpathlineto{\pgfqpoint{7.319214in}{4.238394in}}%
\pgfpathlineto{\pgfqpoint{7.321353in}{4.208625in}}%
\pgfpathlineto{\pgfqpoint{7.329738in}{4.074260in}}%
\pgfpathlineto{\pgfqpoint{7.330337in}{4.075134in}}%
\pgfpathlineto{\pgfqpoint{7.331535in}{4.082873in}}%
\pgfpathlineto{\pgfqpoint{7.333760in}{4.115505in}}%
\pgfpathlineto{\pgfqpoint{7.341546in}{4.242761in}}%
\pgfpathlineto{\pgfqpoint{7.341974in}{4.243013in}}%
\pgfpathlineto{\pgfqpoint{7.342402in}{4.242223in}}%
\pgfpathlineto{\pgfqpoint{7.343600in}{4.234596in}}%
\pgfpathlineto{\pgfqpoint{7.345824in}{4.202169in}}%
\pgfpathlineto{\pgfqpoint{7.353525in}{4.077159in}}%
\pgfpathlineto{\pgfqpoint{7.353953in}{4.076888in}}%
\pgfpathlineto{\pgfqpoint{7.354380in}{4.077660in}}%
\pgfpathlineto{\pgfqpoint{7.355578in}{4.085237in}}%
\pgfpathlineto{\pgfqpoint{7.357803in}{4.117543in}}%
\pgfpathlineto{\pgfqpoint{7.365418in}{4.240183in}}%
\pgfpathlineto{\pgfqpoint{7.365846in}{4.240453in}}%
\pgfpathlineto{\pgfqpoint{7.366274in}{4.239680in}}%
\pgfpathlineto{\pgfqpoint{7.367472in}{4.232104in}}%
\pgfpathlineto{\pgfqpoint{7.369696in}{4.199853in}}%
\pgfpathlineto{\pgfqpoint{7.377140in}{4.079850in}}%
\pgfpathlineto{\pgfqpoint{7.377653in}{4.079419in}}%
\pgfpathlineto{\pgfqpoint{7.378081in}{4.080207in}}%
\pgfpathlineto{\pgfqpoint{7.379279in}{4.087820in}}%
\pgfpathlineto{\pgfqpoint{7.381504in}{4.120062in}}%
\pgfpathlineto{\pgfqpoint{7.388862in}{4.237552in}}%
\pgfpathlineto{\pgfqpoint{7.389376in}{4.237953in}}%
\pgfpathlineto{\pgfqpoint{7.389803in}{4.237140in}}%
\pgfpathlineto{\pgfqpoint{7.391001in}{4.229466in}}%
\pgfpathlineto{\pgfqpoint{7.393226in}{4.197203in}}%
\pgfpathlineto{\pgfqpoint{7.400499in}{4.082251in}}%
\pgfpathlineto{\pgfqpoint{7.401012in}{4.081888in}}%
\pgfpathlineto{\pgfqpoint{7.401440in}{4.082731in}}%
\pgfpathlineto{\pgfqpoint{7.402638in}{4.090480in}}%
\pgfpathlineto{\pgfqpoint{7.404862in}{4.122776in}}%
\pgfpathlineto{\pgfqpoint{7.411964in}{4.234992in}}%
\pgfpathlineto{\pgfqpoint{7.412477in}{4.235567in}}%
\pgfpathlineto{\pgfqpoint{7.412991in}{4.234642in}}%
\pgfpathlineto{\pgfqpoint{7.414189in}{4.226817in}}%
\pgfpathlineto{\pgfqpoint{7.416499in}{4.192900in}}%
\pgfpathlineto{\pgfqpoint{7.423344in}{4.085008in}}%
\pgfpathlineto{\pgfqpoint{7.423943in}{4.084235in}}%
\pgfpathlineto{\pgfqpoint{7.424456in}{4.085194in}}%
\pgfpathlineto{\pgfqpoint{7.425654in}{4.093087in}}%
\pgfpathlineto{\pgfqpoint{7.427964in}{4.127011in}}%
\pgfpathlineto{\pgfqpoint{7.434638in}{4.232199in}}%
\pgfpathlineto{\pgfqpoint{7.435323in}{4.233210in}}%
\pgfpathlineto{\pgfqpoint{7.435750in}{4.232493in}}%
\pgfpathlineto{\pgfqpoint{7.436948in}{4.225101in}}%
\pgfpathlineto{\pgfqpoint{7.439173in}{4.193524in}}%
\pgfpathlineto{\pgfqpoint{7.446103in}{4.087049in}}%
\pgfpathlineto{\pgfqpoint{7.446617in}{4.086555in}}%
\pgfpathlineto{\pgfqpoint{7.447045in}{4.087284in}}%
\pgfpathlineto{\pgfqpoint{7.448242in}{4.094696in}}%
\pgfpathlineto{\pgfqpoint{7.450467in}{4.126217in}}%
\pgfpathlineto{\pgfqpoint{7.457312in}{4.230432in}}%
\pgfpathlineto{\pgfqpoint{7.457825in}{4.230927in}}%
\pgfpathlineto{\pgfqpoint{7.458253in}{4.230201in}}%
\pgfpathlineto{\pgfqpoint{7.459451in}{4.222808in}}%
\pgfpathlineto{\pgfqpoint{7.461676in}{4.191400in}}%
\pgfpathlineto{\pgfqpoint{7.468350in}{4.089548in}}%
\pgfpathlineto{\pgfqpoint{7.468949in}{4.088799in}}%
\pgfpathlineto{\pgfqpoint{7.469376in}{4.089505in}}%
\pgfpathlineto{\pgfqpoint{7.470489in}{4.096057in}}%
\pgfpathlineto{\pgfqpoint{7.472628in}{4.125037in}}%
\pgfpathlineto{\pgfqpoint{7.479558in}{4.228353in}}%
\pgfpathlineto{\pgfqpoint{7.479986in}{4.228721in}}%
\pgfpathlineto{\pgfqpoint{7.480500in}{4.227801in}}%
\pgfpathlineto{\pgfqpoint{7.481697in}{4.220054in}}%
\pgfpathlineto{\pgfqpoint{7.484008in}{4.186820in}}%
\pgfpathlineto{\pgfqpoint{7.490254in}{4.092154in}}%
\pgfpathlineto{\pgfqpoint{7.490938in}{4.090970in}}%
\pgfpathlineto{\pgfqpoint{7.491452in}{4.091811in}}%
\pgfpathlineto{\pgfqpoint{7.492649in}{4.099369in}}%
\pgfpathlineto{\pgfqpoint{7.494960in}{4.132236in}}%
\pgfpathlineto{\pgfqpoint{7.501206in}{4.225558in}}%
\pgfpathlineto{\pgfqpoint{7.501890in}{4.226559in}}%
\pgfpathlineto{\pgfqpoint{7.502318in}{4.225850in}}%
\pgfpathlineto{\pgfqpoint{7.503516in}{4.218556in}}%
\pgfpathlineto{\pgfqpoint{7.505740in}{4.187670in}}%
\pgfpathlineto{\pgfqpoint{7.512158in}{4.093745in}}%
\pgfpathlineto{\pgfqpoint{7.512671in}{4.093094in}}%
\pgfpathlineto{\pgfqpoint{7.513184in}{4.093918in}}%
\pgfpathlineto{\pgfqpoint{7.514382in}{4.101408in}}%
\pgfpathlineto{\pgfqpoint{7.516692in}{4.133955in}}%
\pgfpathlineto{\pgfqpoint{7.522767in}{4.223362in}}%
\pgfpathlineto{\pgfqpoint{7.523452in}{4.224476in}}%
\pgfpathlineto{\pgfqpoint{7.523965in}{4.223597in}}%
\pgfpathlineto{\pgfqpoint{7.525163in}{4.216003in}}%
\pgfpathlineto{\pgfqpoint{7.527473in}{4.183439in}}%
\pgfpathlineto{\pgfqpoint{7.533463in}{4.096258in}}%
\pgfpathlineto{\pgfqpoint{7.534147in}{4.095165in}}%
\pgfpathlineto{\pgfqpoint{7.534661in}{4.096056in}}%
\pgfpathlineto{\pgfqpoint{7.535858in}{4.103655in}}%
\pgfpathlineto{\pgfqpoint{7.538169in}{4.136101in}}%
\pgfpathlineto{\pgfqpoint{7.544072in}{4.221306in}}%
\pgfpathlineto{\pgfqpoint{7.544757in}{4.222441in}}%
\pgfpathlineto{\pgfqpoint{7.545270in}{4.221587in}}%
\pgfpathlineto{\pgfqpoint{7.546468in}{4.214087in}}%
\pgfpathlineto{\pgfqpoint{7.548778in}{4.181899in}}%
\pgfpathlineto{\pgfqpoint{7.554597in}{4.098409in}}%
\pgfpathlineto{\pgfqpoint{7.555281in}{4.097166in}}%
\pgfpathlineto{\pgfqpoint{7.555794in}{4.097933in}}%
\pgfpathlineto{\pgfqpoint{7.556992in}{4.105223in}}%
\pgfpathlineto{\pgfqpoint{7.559303in}{4.137005in}}%
\pgfpathlineto{\pgfqpoint{7.565121in}{4.219359in}}%
\pgfpathlineto{\pgfqpoint{7.565805in}{4.220458in}}%
\pgfpathlineto{\pgfqpoint{7.566319in}{4.219589in}}%
\pgfpathlineto{\pgfqpoint{7.567517in}{4.212097in}}%
\pgfpathlineto{\pgfqpoint{7.569827in}{4.180187in}}%
\pgfpathlineto{\pgfqpoint{7.575559in}{4.100164in}}%
\pgfpathlineto{\pgfqpoint{7.576244in}{4.099132in}}%
\pgfpathlineto{\pgfqpoint{7.576757in}{4.100046in}}%
\pgfpathlineto{\pgfqpoint{7.577955in}{4.107613in}}%
\pgfpathlineto{\pgfqpoint{7.580351in}{4.140952in}}%
\pgfpathlineto{\pgfqpoint{7.585827in}{4.217174in}}%
\pgfpathlineto{\pgfqpoint{7.586597in}{4.218527in}}%
\pgfpathlineto{\pgfqpoint{7.587025in}{4.217878in}}%
\pgfpathlineto{\pgfqpoint{7.588137in}{4.211631in}}%
\pgfpathlineto{\pgfqpoint{7.590362in}{4.182767in}}%
\pgfpathlineto{\pgfqpoint{7.596351in}{4.101639in}}%
\pgfpathlineto{\pgfqpoint{7.596864in}{4.101024in}}%
\pgfpathlineto{\pgfqpoint{7.597378in}{4.101846in}}%
\pgfpathlineto{\pgfqpoint{7.598576in}{4.109165in}}%
\pgfpathlineto{\pgfqpoint{7.600886in}{4.140429in}}%
\pgfpathlineto{\pgfqpoint{7.606447in}{4.215649in}}%
\pgfpathlineto{\pgfqpoint{7.607132in}{4.216648in}}%
\pgfpathlineto{\pgfqpoint{7.607560in}{4.215982in}}%
\pgfpathlineto{\pgfqpoint{7.608672in}{4.209733in}}%
\pgfpathlineto{\pgfqpoint{7.610897in}{4.181120in}}%
\pgfpathlineto{\pgfqpoint{7.616715in}{4.103584in}}%
\pgfpathlineto{\pgfqpoint{7.617314in}{4.102889in}}%
\pgfpathlineto{\pgfqpoint{7.617742in}{4.103580in}}%
\pgfpathlineto{\pgfqpoint{7.618940in}{4.110605in}}%
\pgfpathlineto{\pgfqpoint{7.621250in}{4.141202in}}%
\pgfpathlineto{\pgfqpoint{7.626726in}{4.213827in}}%
\pgfpathlineto{\pgfqpoint{7.627410in}{4.214825in}}%
\pgfpathlineto{\pgfqpoint{7.627924in}{4.213919in}}%
\pgfpathlineto{\pgfqpoint{7.629122in}{4.206493in}}%
\pgfpathlineto{\pgfqpoint{7.631517in}{4.174132in}}%
\pgfpathlineto{\pgfqpoint{7.636651in}{4.106127in}}%
\pgfpathlineto{\pgfqpoint{7.637421in}{4.104672in}}%
\pgfpathlineto{\pgfqpoint{7.637934in}{4.105465in}}%
\pgfpathlineto{\pgfqpoint{7.639132in}{4.112619in}}%
\pgfpathlineto{\pgfqpoint{7.641528in}{4.144470in}}%
\pgfpathlineto{\pgfqpoint{7.646662in}{4.211685in}}%
\pgfpathlineto{\pgfqpoint{7.647432in}{4.213056in}}%
\pgfpathlineto{\pgfqpoint{7.647945in}{4.212215in}}%
\pgfpathlineto{\pgfqpoint{7.649143in}{4.204986in}}%
\pgfpathlineto{\pgfqpoint{7.651539in}{4.173210in}}%
\pgfpathlineto{\pgfqpoint{7.656587in}{4.107828in}}%
\pgfpathlineto{\pgfqpoint{7.657357in}{4.106426in}}%
\pgfpathlineto{\pgfqpoint{7.657870in}{4.107238in}}%
\pgfpathlineto{\pgfqpoint{7.659068in}{4.114376in}}%
\pgfpathlineto{\pgfqpoint{7.661464in}{4.145865in}}%
\pgfpathlineto{\pgfqpoint{7.666512in}{4.210111in}}%
\pgfpathlineto{\pgfqpoint{7.667282in}{4.211313in}}%
\pgfpathlineto{\pgfqpoint{7.667710in}{4.210630in}}%
\pgfpathlineto{\pgfqpoint{7.668908in}{4.203753in}}%
\pgfpathlineto{\pgfqpoint{7.671218in}{4.174126in}}%
\pgfpathlineto{\pgfqpoint{7.676352in}{4.109282in}}%
\pgfpathlineto{\pgfqpoint{7.677036in}{4.108128in}}%
\pgfpathlineto{\pgfqpoint{7.677550in}{4.108877in}}%
\pgfpathlineto{\pgfqpoint{7.678748in}{4.115819in}}%
\pgfpathlineto{\pgfqpoint{7.681143in}{4.146714in}}%
\pgfpathlineto{\pgfqpoint{7.686106in}{4.208458in}}%
\pgfpathlineto{\pgfqpoint{7.686876in}{4.209629in}}%
\pgfpathlineto{\pgfqpoint{7.687304in}{4.208941in}}%
\pgfpathlineto{\pgfqpoint{7.688502in}{4.202113in}}%
\pgfpathlineto{\pgfqpoint{7.690898in}{4.171537in}}%
\pgfpathlineto{\pgfqpoint{7.695860in}{4.110834in}}%
\pgfpathlineto{\pgfqpoint{7.696545in}{4.109793in}}%
\pgfpathlineto{\pgfqpoint{7.697058in}{4.110609in}}%
\pgfpathlineto{\pgfqpoint{7.698256in}{4.117636in}}%
\pgfpathlineto{\pgfqpoint{7.700652in}{4.148271in}}%
\pgfpathlineto{\pgfqpoint{7.705443in}{4.206708in}}%
\pgfpathlineto{\pgfqpoint{7.706213in}{4.208005in}}%
\pgfpathlineto{\pgfqpoint{7.706727in}{4.207170in}}%
\pgfpathlineto{\pgfqpoint{7.707925in}{4.200133in}}%
\pgfpathlineto{\pgfqpoint{7.710406in}{4.168347in}}%
\pgfpathlineto{\pgfqpoint{7.715026in}{4.112790in}}%
\pgfpathlineto{\pgfqpoint{7.715796in}{4.111406in}}%
\pgfpathlineto{\pgfqpoint{7.716310in}{4.112173in}}%
\pgfpathlineto{\pgfqpoint{7.717508in}{4.119031in}}%
\pgfpathlineto{\pgfqpoint{7.719903in}{4.149088in}}%
\pgfpathlineto{\pgfqpoint{7.724609in}{4.205142in}}%
\pgfpathlineto{\pgfqpoint{7.725379in}{4.206409in}}%
\pgfpathlineto{\pgfqpoint{7.725893in}{4.205574in}}%
\pgfpathlineto{\pgfqpoint{7.727091in}{4.198603in}}%
\pgfpathlineto{\pgfqpoint{7.729572in}{4.167301in}}%
\pgfpathlineto{\pgfqpoint{7.734107in}{4.114271in}}%
\pgfpathlineto{\pgfqpoint{7.734877in}{4.112984in}}%
\pgfpathlineto{\pgfqpoint{7.735390in}{4.113797in}}%
\pgfpathlineto{\pgfqpoint{7.736588in}{4.120686in}}%
\pgfpathlineto{\pgfqpoint{7.739069in}{4.151684in}}%
\pgfpathlineto{\pgfqpoint{7.743604in}{4.203732in}}%
\pgfpathlineto{\pgfqpoint{7.744374in}{4.204842in}}%
\pgfpathlineto{\pgfqpoint{7.744802in}{4.204167in}}%
\pgfpathlineto{\pgfqpoint{7.746000in}{4.197557in}}%
\pgfpathlineto{\pgfqpoint{7.748396in}{4.168369in}}%
\pgfpathlineto{\pgfqpoint{7.753016in}{4.115611in}}%
\pgfpathlineto{\pgfqpoint{7.753786in}{4.114537in}}%
\pgfpathlineto{\pgfqpoint{7.754214in}{4.115224in}}%
\pgfpathlineto{\pgfqpoint{7.755412in}{4.121831in}}%
\pgfpathlineto{\pgfqpoint{7.757807in}{4.150821in}}%
\pgfpathlineto{\pgfqpoint{7.762428in}{4.202437in}}%
\pgfpathlineto{\pgfqpoint{7.763112in}{4.203340in}}%
\pgfpathlineto{\pgfqpoint{7.763540in}{4.202719in}}%
\pgfpathlineto{\pgfqpoint{7.764738in}{4.196318in}}%
\pgfpathlineto{\pgfqpoint{7.767134in}{4.167765in}}%
\pgfpathlineto{\pgfqpoint{7.771669in}{4.117122in}}%
\pgfpathlineto{\pgfqpoint{7.772439in}{4.116024in}}%
\pgfpathlineto{\pgfqpoint{7.772866in}{4.116680in}}%
\pgfpathlineto{\pgfqpoint{7.774064in}{4.123140in}}%
\pgfpathlineto{\pgfqpoint{7.776460in}{4.151565in}}%
\pgfpathlineto{\pgfqpoint{7.780995in}{4.200977in}}%
\pgfpathlineto{\pgfqpoint{7.781679in}{4.201868in}}%
\pgfpathlineto{\pgfqpoint{7.782193in}{4.201026in}}%
\pgfpathlineto{\pgfqpoint{7.783476in}{4.193531in}}%
\pgfpathlineto{\pgfqpoint{7.786129in}{4.160769in}}%
\pgfpathlineto{\pgfqpoint{7.790150in}{4.118530in}}%
\pgfpathlineto{\pgfqpoint{7.790920in}{4.117481in}}%
\pgfpathlineto{\pgfqpoint{7.791348in}{4.118146in}}%
\pgfpathlineto{\pgfqpoint{7.792546in}{4.124561in}}%
\pgfpathlineto{\pgfqpoint{7.795027in}{4.153735in}}%
\pgfpathlineto{\pgfqpoint{7.799305in}{4.199342in}}%
\pgfpathlineto{\pgfqpoint{7.800075in}{4.200427in}}%
\pgfpathlineto{\pgfqpoint{7.800503in}{4.199790in}}%
\pgfpathlineto{\pgfqpoint{7.801701in}{4.193484in}}%
\pgfpathlineto{\pgfqpoint{7.804182in}{4.164652in}}%
\pgfpathlineto{\pgfqpoint{7.808460in}{4.119859in}}%
\pgfpathlineto{\pgfqpoint{7.809145in}{4.118873in}}%
\pgfpathlineto{\pgfqpoint{7.809658in}{4.119610in}}%
\pgfpathlineto{\pgfqpoint{7.810856in}{4.126048in}}%
\pgfpathlineto{\pgfqpoint{7.813337in}{4.154809in}}%
\pgfpathlineto{\pgfqpoint{7.817530in}{4.198050in}}%
\pgfpathlineto{\pgfqpoint{7.818215in}{4.199049in}}%
\pgfpathlineto{\pgfqpoint{7.818728in}{4.198333in}}%
\pgfpathlineto{\pgfqpoint{7.819926in}{4.191979in}}%
\pgfpathlineto{\pgfqpoint{7.822407in}{4.163534in}}%
\pgfpathlineto{\pgfqpoint{7.826600in}{4.121129in}}%
\pgfpathlineto{\pgfqpoint{7.827284in}{4.120259in}}%
\pgfpathlineto{\pgfqpoint{7.827798in}{4.121062in}}%
\pgfpathlineto{\pgfqpoint{7.829081in}{4.128254in}}%
\pgfpathlineto{\pgfqpoint{7.831819in}{4.160570in}}%
\pgfpathlineto{\pgfqpoint{7.835584in}{4.196806in}}%
\pgfpathlineto{\pgfqpoint{7.836268in}{4.197681in}}%
\pgfpathlineto{\pgfqpoint{7.836782in}{4.196894in}}%
\pgfpathlineto{\pgfqpoint{7.838065in}{4.189781in}}%
\pgfpathlineto{\pgfqpoint{7.840803in}{4.157822in}}%
\pgfpathlineto{\pgfqpoint{7.844568in}{4.122359in}}%
\pgfpathlineto{\pgfqpoint{7.845252in}{4.121616in}}%
\pgfpathlineto{\pgfqpoint{7.845680in}{4.122260in}}%
\pgfpathlineto{\pgfqpoint{7.846878in}{4.128400in}}%
\pgfpathlineto{\pgfqpoint{7.849359in}{4.155935in}}%
\pgfpathlineto{\pgfqpoint{7.853466in}{4.195594in}}%
\pgfpathlineto{\pgfqpoint{7.854151in}{4.196343in}}%
\pgfpathlineto{\pgfqpoint{7.854579in}{4.195710in}}%
\pgfpathlineto{\pgfqpoint{7.855776in}{4.189640in}}%
\pgfpathlineto{\pgfqpoint{7.858258in}{4.162407in}}%
\pgfpathlineto{\pgfqpoint{7.862279in}{4.123792in}}%
\pgfpathlineto{\pgfqpoint{7.862964in}{4.122903in}}%
\pgfpathlineto{\pgfqpoint{7.863477in}{4.123645in}}%
\pgfpathlineto{\pgfqpoint{7.864761in}{4.130520in}}%
\pgfpathlineto{\pgfqpoint{7.867499in}{4.161407in}}%
\pgfpathlineto{\pgfqpoint{7.871178in}{4.194402in}}%
\pgfpathlineto{\pgfqpoint{7.871777in}{4.195073in}}%
\pgfpathlineto{\pgfqpoint{7.872290in}{4.194349in}}%
\pgfpathlineto{\pgfqpoint{7.873573in}{4.187562in}}%
\pgfpathlineto{\pgfqpoint{7.876311in}{4.157042in}}%
\pgfpathlineto{\pgfqpoint{7.879905in}{4.124966in}}%
\pgfpathlineto{\pgfqpoint{7.880590in}{4.124185in}}%
\pgfpathlineto{\pgfqpoint{7.881017in}{4.124769in}}%
\pgfpathlineto{\pgfqpoint{7.882215in}{4.130600in}}%
\pgfpathlineto{\pgfqpoint{7.884697in}{4.156889in}}%
\pgfpathlineto{\pgfqpoint{7.888632in}{4.193006in}}%
\pgfpathlineto{\pgfqpoint{7.889317in}{4.193809in}}%
\pgfpathlineto{\pgfqpoint{7.889830in}{4.193038in}}%
\pgfpathlineto{\pgfqpoint{7.891114in}{4.186238in}}%
\pgfpathlineto{\pgfqpoint{7.893937in}{4.155231in}}%
\pgfpathlineto{\pgfqpoint{7.897445in}{4.125936in}}%
\pgfpathlineto{\pgfqpoint{7.898044in}{4.125433in}}%
\pgfpathlineto{\pgfqpoint{7.898472in}{4.126047in}}%
\pgfpathlineto{\pgfqpoint{7.899670in}{4.131875in}}%
\pgfpathlineto{\pgfqpoint{7.902237in}{4.158739in}}%
\pgfpathlineto{\pgfqpoint{7.906002in}{4.191837in}}%
\pgfpathlineto{\pgfqpoint{7.906686in}{4.192581in}}%
\pgfpathlineto{\pgfqpoint{7.907114in}{4.192001in}}%
\pgfpathlineto{\pgfqpoint{7.908312in}{4.186302in}}%
\pgfpathlineto{\pgfqpoint{7.910879in}{4.159818in}}%
\pgfpathlineto{\pgfqpoint{7.914643in}{4.127310in}}%
\pgfpathlineto{\pgfqpoint{7.915328in}{4.126644in}}%
\pgfpathlineto{\pgfqpoint{7.915756in}{4.127261in}}%
\pgfpathlineto{\pgfqpoint{7.916954in}{4.133019in}}%
\pgfpathlineto{\pgfqpoint{7.919606in}{4.160327in}}%
\pgfpathlineto{\pgfqpoint{7.923285in}{4.190862in}}%
\pgfpathlineto{\pgfqpoint{7.923884in}{4.191390in}}%
\pgfpathlineto{\pgfqpoint{7.924398in}{4.190614in}}%
\pgfpathlineto{\pgfqpoint{7.925681in}{4.183992in}}%
\pgfpathlineto{\pgfqpoint{7.928590in}{4.153388in}}%
\pgfpathlineto{\pgfqpoint{7.931841in}{4.128300in}}%
\pgfpathlineto{\pgfqpoint{7.932440in}{4.127814in}}%
\pgfpathlineto{\pgfqpoint{7.932868in}{4.128404in}}%
\pgfpathlineto{\pgfqpoint{7.934066in}{4.134009in}}%
\pgfpathlineto{\pgfqpoint{7.936719in}{4.160634in}}%
\pgfpathlineto{\pgfqpoint{7.940312in}{4.189672in}}%
\pgfpathlineto{\pgfqpoint{7.940911in}{4.190238in}}%
\pgfpathlineto{\pgfqpoint{7.941424in}{4.189517in}}%
\pgfpathlineto{\pgfqpoint{7.942708in}{4.183117in}}%
\pgfpathlineto{\pgfqpoint{7.945617in}{4.153344in}}%
\pgfpathlineto{\pgfqpoint{7.948868in}{4.129326in}}%
\pgfpathlineto{\pgfqpoint{7.949467in}{4.128990in}}%
\pgfpathlineto{\pgfqpoint{7.949895in}{4.129668in}}%
\pgfpathlineto{\pgfqpoint{7.951179in}{4.136021in}}%
\pgfpathlineto{\pgfqpoint{7.954173in}{4.166363in}}%
\pgfpathlineto{\pgfqpoint{7.957339in}{4.188801in}}%
\pgfpathlineto{\pgfqpoint{7.957938in}{4.189033in}}%
\pgfpathlineto{\pgfqpoint{7.958280in}{4.188499in}}%
\pgfpathlineto{\pgfqpoint{7.959478in}{4.183013in}}%
\pgfpathlineto{\pgfqpoint{7.962131in}{4.157306in}}%
\pgfpathlineto{\pgfqpoint{7.965639in}{4.130533in}}%
\pgfpathlineto{\pgfqpoint{7.966238in}{4.130072in}}%
\pgfpathlineto{\pgfqpoint{7.966665in}{4.130643in}}%
\pgfpathlineto{\pgfqpoint{7.967863in}{4.136032in}}%
\pgfpathlineto{\pgfqpoint{7.970516in}{4.161380in}}%
\pgfpathlineto{\pgfqpoint{7.974024in}{4.187592in}}%
\pgfpathlineto{\pgfqpoint{7.974623in}{4.187983in}}%
\pgfpathlineto{\pgfqpoint{7.975051in}{4.187371in}}%
\pgfpathlineto{\pgfqpoint{7.976334in}{4.181340in}}%
\pgfpathlineto{\pgfqpoint{7.979243in}{4.153162in}}%
\pgfpathlineto{\pgfqpoint{7.982409in}{4.131448in}}%
\pgfpathlineto{\pgfqpoint{7.983008in}{4.131212in}}%
\pgfpathlineto{\pgfqpoint{7.983350in}{4.131724in}}%
\pgfpathlineto{\pgfqpoint{7.984548in}{4.137016in}}%
\pgfpathlineto{\pgfqpoint{7.987286in}{4.162633in}}%
\pgfpathlineto{\pgfqpoint{7.990623in}{4.186511in}}%
\pgfpathlineto{\pgfqpoint{7.991222in}{4.186920in}}%
\pgfpathlineto{\pgfqpoint{7.991650in}{4.186340in}}%
\pgfpathlineto{\pgfqpoint{7.992933in}{4.180489in}}%
\pgfpathlineto{\pgfqpoint{7.995928in}{4.152238in}}%
\pgfpathlineto{\pgfqpoint{7.999008in}{4.132412in}}%
\pgfpathlineto{\pgfqpoint{7.999607in}{4.132319in}}%
\pgfpathlineto{\pgfqpoint{7.999949in}{4.132898in}}%
\pgfpathlineto{\pgfqpoint{8.001233in}{4.138891in}}%
\pgfpathlineto{\pgfqpoint{8.004398in}{4.168559in}}%
\pgfpathlineto{\pgfqpoint{8.007308in}{4.185781in}}%
\pgfpathlineto{\pgfqpoint{8.007907in}{4.185685in}}%
\pgfpathlineto{\pgfqpoint{8.008163in}{4.185218in}}%
\pgfpathlineto{\pgfqpoint{8.009447in}{4.179314in}}%
\pgfpathlineto{\pgfqpoint{8.012612in}{4.150081in}}%
\pgfpathlineto{\pgfqpoint{8.015436in}{4.133404in}}%
\pgfpathlineto{\pgfqpoint{8.016035in}{4.133369in}}%
\pgfpathlineto{\pgfqpoint{8.016377in}{4.133967in}}%
\pgfpathlineto{\pgfqpoint{8.017661in}{4.139929in}}%
\pgfpathlineto{\pgfqpoint{8.020998in}{4.170370in}}%
\pgfpathlineto{\pgfqpoint{8.023736in}{4.184845in}}%
\pgfpathlineto{\pgfqpoint{8.024420in}{4.184388in}}%
\pgfpathlineto{\pgfqpoint{8.024506in}{4.184206in}}%
\pgfpathlineto{\pgfqpoint{8.025789in}{4.178412in}}%
\pgfpathlineto{\pgfqpoint{8.029040in}{4.149235in}}%
\pgfpathlineto{\pgfqpoint{8.031778in}{4.134323in}}%
\pgfpathlineto{\pgfqpoint{8.032377in}{4.134453in}}%
\pgfpathlineto{\pgfqpoint{8.032634in}{4.134919in}}%
\pgfpathlineto{\pgfqpoint{8.033917in}{4.140687in}}%
\pgfpathlineto{\pgfqpoint{8.037254in}{4.170197in}}%
\pgfpathlineto{\pgfqpoint{8.039907in}{4.183837in}}%
\pgfpathlineto{\pgfqpoint{8.040591in}{4.183474in}}%
\pgfpathlineto{\pgfqpoint{8.040762in}{4.183113in}}%
\pgfpathlineto{\pgfqpoint{8.042131in}{4.176668in}}%
\pgfpathlineto{\pgfqpoint{8.048206in}{4.135190in}}%
\pgfpathlineto{\pgfqpoint{8.049062in}{4.136670in}}%
\pgfpathlineto{\pgfqpoint{8.050688in}{4.146004in}}%
\pgfpathlineto{\pgfqpoint{8.056249in}{4.182943in}}%
\pgfpathlineto{\pgfqpoint{8.056420in}{4.182847in}}%
\pgfpathlineto{\pgfqpoint{8.057362in}{4.180462in}}%
\pgfpathlineto{\pgfqpoint{8.059244in}{4.167885in}}%
\pgfpathlineto{\pgfqpoint{8.063950in}{4.136220in}}%
\pgfpathlineto{\pgfqpoint{8.064634in}{4.136518in}}%
\pgfpathlineto{\pgfqpoint{8.064806in}{4.136853in}}%
\pgfpathlineto{\pgfqpoint{8.066175in}{4.142971in}}%
\pgfpathlineto{\pgfqpoint{8.072164in}{4.182016in}}%
\pgfpathlineto{\pgfqpoint{8.073020in}{4.180625in}}%
\pgfpathlineto{\pgfqpoint{8.074645in}{4.171701in}}%
\pgfpathlineto{\pgfqpoint{8.080121in}{4.137061in}}%
\pgfpathlineto{\pgfqpoint{8.080292in}{4.137151in}}%
\pgfpathlineto{\pgfqpoint{8.081234in}{4.139436in}}%
\pgfpathlineto{\pgfqpoint{8.083201in}{4.152204in}}%
\pgfpathlineto{\pgfqpoint{8.087736in}{4.181028in}}%
\pgfpathlineto{\pgfqpoint{8.088421in}{4.180733in}}%
\pgfpathlineto{\pgfqpoint{8.088592in}{4.180408in}}%
\pgfpathlineto{\pgfqpoint{8.089961in}{4.174511in}}%
\pgfpathlineto{\pgfqpoint{8.095865in}{4.137953in}}%
\pgfpathlineto{\pgfqpoint{8.096635in}{4.139073in}}%
\pgfpathlineto{\pgfqpoint{8.098175in}{4.146723in}}%
\pgfpathlineto{\pgfqpoint{8.103736in}{4.180216in}}%
\pgfpathlineto{\pgfqpoint{8.103993in}{4.180036in}}%
\pgfpathlineto{\pgfqpoint{8.105020in}{4.177173in}}%
\pgfpathlineto{\pgfqpoint{8.107245in}{4.162098in}}%
\pgfpathlineto{\pgfqpoint{8.111095in}{4.139046in}}%
\pgfpathlineto{\pgfqpoint{8.111779in}{4.138986in}}%
\pgfpathlineto{\pgfqpoint{8.112122in}{4.139536in}}%
\pgfpathlineto{\pgfqpoint{8.113491in}{4.145272in}}%
\pgfpathlineto{\pgfqpoint{8.119309in}{4.179360in}}%
\pgfpathlineto{\pgfqpoint{8.119993in}{4.178446in}}%
\pgfpathlineto{\pgfqpoint{8.121448in}{4.171908in}}%
\pgfpathlineto{\pgfqpoint{8.127009in}{4.139668in}}%
\pgfpathlineto{\pgfqpoint{8.127523in}{4.140126in}}%
\pgfpathlineto{\pgfqpoint{8.128806in}{4.144763in}}%
\pgfpathlineto{\pgfqpoint{8.132229in}{4.169478in}}%
\pgfpathlineto{\pgfqpoint{8.134710in}{4.178532in}}%
\pgfpathlineto{\pgfqpoint{8.135566in}{4.177374in}}%
\pgfpathlineto{\pgfqpoint{8.137191in}{4.169533in}}%
\pgfpathlineto{\pgfqpoint{8.142411in}{4.140491in}}%
\pgfpathlineto{\pgfqpoint{8.142582in}{4.140542in}}%
\pgfpathlineto{\pgfqpoint{8.143523in}{4.142437in}}%
\pgfpathlineto{\pgfqpoint{8.145405in}{4.152956in}}%
\pgfpathlineto{\pgfqpoint{8.149855in}{4.177653in}}%
\pgfpathlineto{\pgfqpoint{8.150710in}{4.177077in}}%
\pgfpathlineto{\pgfqpoint{8.150796in}{4.176896in}}%
\pgfpathlineto{\pgfqpoint{8.152250in}{4.170805in}}%
\pgfpathlineto{\pgfqpoint{8.157726in}{4.141292in}}%
\pgfpathlineto{\pgfqpoint{8.158240in}{4.141735in}}%
\pgfpathlineto{\pgfqpoint{8.159523in}{4.146133in}}%
\pgfpathlineto{\pgfqpoint{8.165341in}{4.176930in}}%
\pgfpathlineto{\pgfqpoint{8.166283in}{4.175519in}}%
\pgfpathlineto{\pgfqpoint{8.167994in}{4.167160in}}%
\pgfpathlineto{\pgfqpoint{8.172871in}{4.142068in}}%
\pgfpathlineto{\pgfqpoint{8.173727in}{4.143102in}}%
\pgfpathlineto{\pgfqpoint{8.175352in}{4.150302in}}%
\pgfpathlineto{\pgfqpoint{8.180486in}{4.176163in}}%
\pgfpathlineto{\pgfqpoint{8.180657in}{4.176098in}}%
\pgfpathlineto{\pgfqpoint{8.181684in}{4.173960in}}%
\pgfpathlineto{\pgfqpoint{8.183737in}{4.162681in}}%
\pgfpathlineto{\pgfqpoint{8.187673in}{4.142948in}}%
\pgfpathlineto{\pgfqpoint{8.188443in}{4.143148in}}%
\pgfpathlineto{\pgfqpoint{8.188614in}{4.143421in}}%
\pgfpathlineto{\pgfqpoint{8.189983in}{4.148322in}}%
\pgfpathlineto{\pgfqpoint{8.195545in}{4.175413in}}%
\pgfpathlineto{\pgfqpoint{8.196144in}{4.174778in}}%
\pgfpathlineto{\pgfqpoint{8.197598in}{4.169420in}}%
\pgfpathlineto{\pgfqpoint{8.202989in}{4.143560in}}%
\pgfpathlineto{\pgfqpoint{8.203502in}{4.143993in}}%
\pgfpathlineto{\pgfqpoint{8.204786in}{4.148057in}}%
\pgfpathlineto{\pgfqpoint{8.210433in}{4.174696in}}%
\pgfpathlineto{\pgfqpoint{8.211288in}{4.173604in}}%
\pgfpathlineto{\pgfqpoint{8.212914in}{4.166775in}}%
\pgfpathlineto{\pgfqpoint{8.217877in}{4.144274in}}%
\pgfpathlineto{\pgfqpoint{8.217962in}{4.144299in}}%
\pgfpathlineto{\pgfqpoint{8.218904in}{4.145848in}}%
\pgfpathlineto{\pgfqpoint{8.220786in}{4.154671in}}%
\pgfpathlineto{\pgfqpoint{8.225064in}{4.173970in}}%
\pgfpathlineto{\pgfqpoint{8.226005in}{4.173163in}}%
\pgfpathlineto{\pgfqpoint{8.227545in}{4.167343in}}%
\pgfpathlineto{\pgfqpoint{8.232679in}{4.144972in}}%
\pgfpathlineto{\pgfqpoint{8.232850in}{4.145055in}}%
\pgfpathlineto{\pgfqpoint{8.233963in}{4.147380in}}%
\pgfpathlineto{\pgfqpoint{8.236273in}{4.159297in}}%
\pgfpathlineto{\pgfqpoint{8.239695in}{4.173248in}}%
\pgfpathlineto{\pgfqpoint{8.240722in}{4.172513in}}%
\pgfpathlineto{\pgfqpoint{8.242262in}{4.166898in}}%
\pgfpathlineto{\pgfqpoint{8.247310in}{4.145635in}}%
\pgfpathlineto{\pgfqpoint{8.247481in}{4.145699in}}%
\pgfpathlineto{\pgfqpoint{8.248508in}{4.147582in}}%
\pgfpathlineto{\pgfqpoint{8.250647in}{4.157734in}}%
\pgfpathlineto{\pgfqpoint{8.254326in}{4.172597in}}%
\pgfpathlineto{\pgfqpoint{8.255353in}{4.171856in}}%
\pgfpathlineto{\pgfqpoint{8.256979in}{4.165986in}}%
\pgfpathlineto{\pgfqpoint{8.261941in}{4.146299in}}%
\pgfpathlineto{\pgfqpoint{8.262027in}{4.146333in}}%
\pgfpathlineto{\pgfqpoint{8.263054in}{4.148093in}}%
\pgfpathlineto{\pgfqpoint{8.265193in}{4.157797in}}%
\pgfpathlineto{\pgfqpoint{8.268872in}{4.171970in}}%
\pgfpathlineto{\pgfqpoint{8.269899in}{4.171198in}}%
\pgfpathlineto{\pgfqpoint{8.271524in}{4.165476in}}%
\pgfpathlineto{\pgfqpoint{8.276401in}{4.146925in}}%
\pgfpathlineto{\pgfqpoint{8.277428in}{4.148394in}}%
\pgfpathlineto{\pgfqpoint{8.279396in}{4.156559in}}%
\pgfpathlineto{\pgfqpoint{8.283332in}{4.171365in}}%
\pgfpathlineto{\pgfqpoint{8.284359in}{4.170543in}}%
\pgfpathlineto{\pgfqpoint{8.286070in}{4.164544in}}%
\pgfpathlineto{\pgfqpoint{8.290690in}{4.147522in}}%
\pgfpathlineto{\pgfqpoint{8.291717in}{4.148736in}}%
\pgfpathlineto{\pgfqpoint{8.293600in}{4.155921in}}%
\pgfpathlineto{\pgfqpoint{8.297707in}{4.170778in}}%
\pgfpathlineto{\pgfqpoint{8.298733in}{4.169894in}}%
\pgfpathlineto{\pgfqpoint{8.300445in}{4.164018in}}%
\pgfpathlineto{\pgfqpoint{8.304979in}{4.148115in}}%
\pgfpathlineto{\pgfqpoint{8.306006in}{4.149287in}}%
\pgfpathlineto{\pgfqpoint{8.307888in}{4.156190in}}%
\pgfpathlineto{\pgfqpoint{8.311910in}{4.170185in}}%
\pgfpathlineto{\pgfqpoint{8.312937in}{4.169415in}}%
\pgfpathlineto{\pgfqpoint{8.314648in}{4.163870in}}%
\pgfpathlineto{\pgfqpoint{8.319183in}{4.148688in}}%
\pgfpathlineto{\pgfqpoint{8.320209in}{4.149831in}}%
\pgfpathlineto{\pgfqpoint{8.322092in}{4.156468in}}%
\pgfpathlineto{\pgfqpoint{8.326113in}{4.169636in}}%
\pgfpathlineto{\pgfqpoint{8.327140in}{4.168788in}}%
\pgfpathlineto{\pgfqpoint{8.328937in}{4.162983in}}%
\pgfpathlineto{\pgfqpoint{8.333215in}{4.149244in}}%
\pgfpathlineto{\pgfqpoint{8.334242in}{4.150194in}}%
\pgfpathlineto{\pgfqpoint{8.336124in}{4.156379in}}%
\pgfpathlineto{\pgfqpoint{8.340146in}{4.169084in}}%
\pgfpathlineto{\pgfqpoint{8.341258in}{4.168176in}}%
\pgfpathlineto{\pgfqpoint{8.343140in}{4.162136in}}%
\pgfpathlineto{\pgfqpoint{8.347162in}{4.149792in}}%
\pgfpathlineto{\pgfqpoint{8.348274in}{4.150718in}}%
\pgfpathlineto{\pgfqpoint{8.350156in}{4.156664in}}%
\pgfpathlineto{\pgfqpoint{8.354178in}{4.168568in}}%
\pgfpathlineto{\pgfqpoint{8.355290in}{4.167582in}}%
\pgfpathlineto{\pgfqpoint{8.357258in}{4.161340in}}%
\pgfpathlineto{\pgfqpoint{8.361108in}{4.150307in}}%
\pgfpathlineto{\pgfqpoint{8.362221in}{4.151223in}}%
\pgfpathlineto{\pgfqpoint{8.364103in}{4.156934in}}%
\pgfpathlineto{\pgfqpoint{8.368039in}{4.168054in}}%
\pgfpathlineto{\pgfqpoint{8.369151in}{4.167165in}}%
\pgfpathlineto{\pgfqpoint{8.371034in}{4.161593in}}%
\pgfpathlineto{\pgfqpoint{8.374969in}{4.150806in}}%
\pgfpathlineto{\pgfqpoint{8.376082in}{4.151708in}}%
\pgfpathlineto{\pgfqpoint{8.378050in}{4.157503in}}%
\pgfpathlineto{\pgfqpoint{8.381814in}{4.167555in}}%
\pgfpathlineto{\pgfqpoint{8.382927in}{4.166760in}}%
\pgfpathlineto{\pgfqpoint{8.384809in}{4.161508in}}%
\pgfpathlineto{\pgfqpoint{8.388745in}{4.151289in}}%
\pgfpathlineto{\pgfqpoint{8.389943in}{4.152320in}}%
\pgfpathlineto{\pgfqpoint{8.392082in}{4.158634in}}%
\pgfpathlineto{\pgfqpoint{8.395504in}{4.167072in}}%
\pgfpathlineto{\pgfqpoint{8.396702in}{4.166237in}}%
\pgfpathlineto{\pgfqpoint{8.398670in}{4.160846in}}%
\pgfpathlineto{\pgfqpoint{8.402435in}{4.151755in}}%
\pgfpathlineto{\pgfqpoint{8.403633in}{4.152754in}}%
\pgfpathlineto{\pgfqpoint{8.405772in}{4.158768in}}%
\pgfpathlineto{\pgfqpoint{8.409194in}{4.166625in}}%
\pgfpathlineto{\pgfqpoint{8.410392in}{4.165741in}}%
\pgfpathlineto{\pgfqpoint{8.412446in}{4.160253in}}%
\pgfpathlineto{\pgfqpoint{8.415954in}{4.152219in}}%
\pgfpathlineto{\pgfqpoint{8.417152in}{4.153026in}}%
\pgfpathlineto{\pgfqpoint{8.419205in}{4.158319in}}%
\pgfpathlineto{\pgfqpoint{8.422799in}{4.166188in}}%
\pgfpathlineto{\pgfqpoint{8.424082in}{4.165134in}}%
\pgfpathlineto{\pgfqpoint{8.426392in}{4.158925in}}%
\pgfpathlineto{\pgfqpoint{8.429558in}{4.152640in}}%
\pgfpathlineto{\pgfqpoint{8.430842in}{4.153682in}}%
\pgfpathlineto{\pgfqpoint{8.433237in}{4.159994in}}%
\pgfpathlineto{\pgfqpoint{8.436232in}{4.165755in}}%
\pgfpathlineto{\pgfqpoint{8.437516in}{4.164840in}}%
\pgfpathlineto{\pgfqpoint{8.439740in}{4.159290in}}%
\pgfpathlineto{\pgfqpoint{8.442992in}{4.153060in}}%
\pgfpathlineto{\pgfqpoint{8.444275in}{4.154030in}}%
\pgfpathlineto{\pgfqpoint{8.446671in}{4.159980in}}%
\pgfpathlineto{\pgfqpoint{8.449665in}{4.165349in}}%
\pgfpathlineto{\pgfqpoint{8.450949in}{4.164436in}}%
\pgfpathlineto{\pgfqpoint{8.453345in}{4.158682in}}%
\pgfpathlineto{\pgfqpoint{8.456339in}{4.153464in}}%
\pgfpathlineto{\pgfqpoint{8.457708in}{4.154481in}}%
\pgfpathlineto{\pgfqpoint{8.460275in}{4.160622in}}%
\pgfpathlineto{\pgfqpoint{8.463099in}{4.164960in}}%
\pgfpathlineto{\pgfqpoint{8.464468in}{4.163803in}}%
\pgfpathlineto{\pgfqpoint{8.467462in}{4.156715in}}%
\pgfpathlineto{\pgfqpoint{8.469858in}{4.153857in}}%
\pgfpathlineto{\pgfqpoint{8.471227in}{4.155169in}}%
\pgfpathlineto{\pgfqpoint{8.476874in}{4.164412in}}%
\pgfpathlineto{\pgfqpoint{8.477473in}{4.163838in}}%
\pgfpathlineto{\pgfqpoint{8.479783in}{4.159021in}}%
\pgfpathlineto{\pgfqpoint{8.482864in}{4.154222in}}%
\pgfpathlineto{\pgfqpoint{8.484318in}{4.155272in}}%
\pgfpathlineto{\pgfqpoint{8.487484in}{4.162051in}}%
\pgfpathlineto{\pgfqpoint{8.489794in}{4.164174in}}%
\pgfpathlineto{\pgfqpoint{8.491334in}{4.162495in}}%
\pgfpathlineto{\pgfqpoint{8.496639in}{4.154741in}}%
\pgfpathlineto{\pgfqpoint{8.496810in}{4.154858in}}%
\pgfpathlineto{\pgfqpoint{8.498693in}{4.157788in}}%
\pgfpathlineto{\pgfqpoint{8.502714in}{4.163862in}}%
\pgfpathlineto{\pgfqpoint{8.504254in}{4.162552in}}%
\pgfpathlineto{\pgfqpoint{8.509816in}{4.155137in}}%
\pgfpathlineto{\pgfqpoint{8.510073in}{4.155337in}}%
\pgfpathlineto{\pgfqpoint{8.512212in}{4.158835in}}%
\pgfpathlineto{\pgfqpoint{8.515634in}{4.163530in}}%
\pgfpathlineto{\pgfqpoint{8.517174in}{4.162463in}}%
\pgfpathlineto{\pgfqpoint{8.522907in}{4.155509in}}%
\pgfpathlineto{\pgfqpoint{8.523249in}{4.155798in}}%
\pgfpathlineto{\pgfqpoint{8.525730in}{4.159897in}}%
\pgfpathlineto{\pgfqpoint{8.528725in}{4.163206in}}%
\pgfpathlineto{\pgfqpoint{8.530351in}{4.161935in}}%
\pgfpathlineto{\pgfqpoint{8.535912in}{4.155855in}}%
\pgfpathlineto{\pgfqpoint{8.538051in}{4.158794in}}%
\pgfpathlineto{\pgfqpoint{8.541560in}{4.162903in}}%
\pgfpathlineto{\pgfqpoint{8.543271in}{4.161687in}}%
\pgfpathlineto{\pgfqpoint{8.548747in}{4.156115in}}%
\pgfpathlineto{\pgfqpoint{8.550886in}{4.158813in}}%
\pgfpathlineto{\pgfqpoint{8.554479in}{4.162606in}}%
\pgfpathlineto{\pgfqpoint{8.556276in}{4.161261in}}%
\pgfpathlineto{\pgfqpoint{8.561410in}{4.156310in}}%
\pgfpathlineto{\pgfqpoint{8.563549in}{4.158676in}}%
\pgfpathlineto{\pgfqpoint{8.567228in}{4.162329in}}%
\pgfpathlineto{\pgfqpoint{8.569025in}{4.161100in}}%
\pgfpathlineto{\pgfqpoint{8.574159in}{4.156579in}}%
\pgfpathlineto{\pgfqpoint{8.576383in}{4.158896in}}%
\pgfpathlineto{\pgfqpoint{8.579977in}{4.162059in}}%
\pgfpathlineto{\pgfqpoint{8.581945in}{4.160678in}}%
\pgfpathlineto{\pgfqpoint{8.586651in}{4.156759in}}%
\pgfpathlineto{\pgfqpoint{8.588790in}{4.158627in}}%
\pgfpathlineto{\pgfqpoint{8.592726in}{4.161792in}}%
\pgfpathlineto{\pgfqpoint{8.594779in}{4.160317in}}%
\pgfpathlineto{\pgfqpoint{8.599143in}{4.156968in}}%
\pgfpathlineto{\pgfqpoint{8.601368in}{4.158681in}}%
\pgfpathlineto{\pgfqpoint{8.605304in}{4.161554in}}%
\pgfpathlineto{\pgfqpoint{8.607528in}{4.160018in}}%
\pgfpathlineto{\pgfqpoint{8.611635in}{4.157186in}}%
\pgfpathlineto{\pgfqpoint{8.613945in}{4.158799in}}%
\pgfpathlineto{\pgfqpoint{8.617881in}{4.161318in}}%
\pgfpathlineto{\pgfqpoint{8.620277in}{4.159696in}}%
\pgfpathlineto{\pgfqpoint{8.624127in}{4.157406in}}%
\pgfpathlineto{\pgfqpoint{8.626609in}{4.159043in}}%
\pgfpathlineto{\pgfqpoint{8.630373in}{4.161099in}}%
\pgfpathlineto{\pgfqpoint{8.632940in}{4.159444in}}%
\pgfpathlineto{\pgfqpoint{8.636534in}{4.157608in}}%
\pgfpathlineto{\pgfqpoint{8.639186in}{4.159225in}}%
\pgfpathlineto{\pgfqpoint{8.642780in}{4.160898in}}%
\pgfpathlineto{\pgfqpoint{8.645518in}{4.159259in}}%
\pgfpathlineto{\pgfqpoint{8.649026in}{4.157824in}}%
\pgfpathlineto{\pgfqpoint{8.652021in}{4.159610in}}%
\pgfpathlineto{\pgfqpoint{8.655272in}{4.160681in}}%
\pgfpathlineto{\pgfqpoint{8.658609in}{4.158725in}}%
\pgfpathlineto{\pgfqpoint{8.661689in}{4.158092in}}%
\pgfpathlineto{\pgfqpoint{8.669647in}{4.159536in}}%
\pgfpathlineto{\pgfqpoint{8.673754in}{4.158204in}}%
\pgfpathlineto{\pgfqpoint{8.682481in}{4.159151in}}%
\pgfpathlineto{\pgfqpoint{8.686160in}{4.158404in}}%
\pgfpathlineto{\pgfqpoint{8.694117in}{4.159407in}}%
\pgfpathlineto{\pgfqpoint{8.698396in}{4.158561in}}%
\pgfpathlineto{\pgfqpoint{8.706182in}{4.159396in}}%
\pgfpathlineto{\pgfqpoint{8.710717in}{4.158736in}}%
\pgfpathlineto{\pgfqpoint{8.717904in}{4.159483in}}%
\pgfpathlineto{\pgfqpoint{8.723038in}{4.158905in}}%
\pgfpathlineto{\pgfqpoint{8.729797in}{4.159472in}}%
\pgfpathlineto{\pgfqpoint{8.735359in}{4.159061in}}%
\pgfpathlineto{\pgfqpoint{8.741690in}{4.159445in}}%
\pgfpathlineto{\pgfqpoint{8.747680in}{4.159195in}}%
\pgfpathlineto{\pgfqpoint{8.753926in}{4.159346in}}%
\pgfpathlineto{\pgfqpoint{8.760086in}{4.159314in}}%
\pgfpathlineto{\pgfqpoint{8.766418in}{4.159239in}}%
\pgfpathlineto{\pgfqpoint{8.773092in}{4.159449in}}%
\pgfpathlineto{\pgfqpoint{8.785840in}{4.159456in}}%
\pgfpathlineto{\pgfqpoint{8.832044in}{4.159319in}}%
\pgfpathlineto{\pgfqpoint{8.850354in}{4.159307in}}%
\pgfpathlineto{\pgfqpoint{8.850354in}{4.159307in}}%
\pgfusepath{stroke}%
\end{pgfscope}%
\begin{pgfscope}%
\pgfpathrectangle{\pgfqpoint{0.294194in}{3.039076in}}{\pgfqpoint{8.556160in}{2.240462in}}%
\pgfusepath{clip}%
\pgfsetrectcap%
\pgfsetroundjoin%
\pgfsetlinewidth{1.505625pt}%
\definecolor{currentstroke}{rgb}{1.000000,0.498039,0.054902}%
\pgfsetstrokecolor{currentstroke}%
\pgfsetdash{}{0pt}%
\pgfpathmoveto{\pgfqpoint{0.294194in}{4.159307in}}%
\pgfpathlineto{\pgfqpoint{0.297018in}{5.279385in}}%
\pgfpathlineto{\pgfqpoint{0.297103in}{5.279014in}}%
\pgfpathlineto{\pgfqpoint{0.297617in}{5.225112in}}%
\pgfpathlineto{\pgfqpoint{0.298900in}{4.746738in}}%
\pgfpathlineto{\pgfqpoint{0.302751in}{3.039109in}}%
\pgfpathlineto{\pgfqpoint{0.303093in}{3.057286in}}%
\pgfpathlineto{\pgfqpoint{0.304034in}{3.300437in}}%
\pgfpathlineto{\pgfqpoint{0.308483in}{5.279517in}}%
\pgfpathlineto{\pgfqpoint{0.309510in}{5.109869in}}%
\pgfpathlineto{\pgfqpoint{0.311478in}{4.087507in}}%
\pgfpathlineto{\pgfqpoint{0.314216in}{3.039137in}}%
\pgfpathlineto{\pgfqpoint{0.314302in}{3.039890in}}%
\pgfpathlineto{\pgfqpoint{0.314815in}{3.095398in}}%
\pgfpathlineto{\pgfqpoint{0.316098in}{3.572788in}}%
\pgfpathlineto{\pgfqpoint{0.319949in}{5.279293in}}%
\pgfpathlineto{\pgfqpoint{0.320376in}{5.254060in}}%
\pgfpathlineto{\pgfqpoint{0.321403in}{4.960759in}}%
\pgfpathlineto{\pgfqpoint{0.325767in}{3.039163in}}%
\pgfpathlineto{\pgfqpoint{0.326623in}{3.163153in}}%
\pgfpathlineto{\pgfqpoint{0.328334in}{3.975556in}}%
\pgfpathlineto{\pgfqpoint{0.331500in}{5.279362in}}%
\pgfpathlineto{\pgfqpoint{0.331585in}{5.278925in}}%
\pgfpathlineto{\pgfqpoint{0.332099in}{5.225928in}}%
\pgfpathlineto{\pgfqpoint{0.333382in}{4.757611in}}%
\pgfpathlineto{\pgfqpoint{0.337318in}{3.039147in}}%
\pgfpathlineto{\pgfqpoint{0.337660in}{3.058511in}}%
\pgfpathlineto{\pgfqpoint{0.338601in}{3.299837in}}%
\pgfpathlineto{\pgfqpoint{0.343136in}{5.279398in}}%
\pgfpathlineto{\pgfqpoint{0.344163in}{5.105581in}}%
\pgfpathlineto{\pgfqpoint{0.346216in}{4.043416in}}%
\pgfpathlineto{\pgfqpoint{0.348954in}{3.039261in}}%
\pgfpathlineto{\pgfqpoint{0.349297in}{3.060630in}}%
\pgfpathlineto{\pgfqpoint{0.350323in}{3.339666in}}%
\pgfpathlineto{\pgfqpoint{0.354773in}{5.279377in}}%
\pgfpathlineto{\pgfqpoint{0.355628in}{5.159472in}}%
\pgfpathlineto{\pgfqpoint{0.357339in}{4.363016in}}%
\pgfpathlineto{\pgfqpoint{0.360591in}{3.039260in}}%
\pgfpathlineto{\pgfqpoint{0.360848in}{3.049044in}}%
\pgfpathlineto{\pgfqpoint{0.361703in}{3.229863in}}%
\pgfpathlineto{\pgfqpoint{0.363842in}{4.348433in}}%
\pgfpathlineto{\pgfqpoint{0.366495in}{5.279078in}}%
\pgfpathlineto{\pgfqpoint{0.366922in}{5.244239in}}%
\pgfpathlineto{\pgfqpoint{0.368035in}{4.900904in}}%
\pgfpathlineto{\pgfqpoint{0.372313in}{3.039315in}}%
\pgfpathlineto{\pgfqpoint{0.372997in}{3.111969in}}%
\pgfpathlineto{\pgfqpoint{0.374366in}{3.645473in}}%
\pgfpathlineto{\pgfqpoint{0.378217in}{5.279224in}}%
\pgfpathlineto{\pgfqpoint{0.378473in}{5.267311in}}%
\pgfpathlineto{\pgfqpoint{0.379329in}{5.081839in}}%
\pgfpathlineto{\pgfqpoint{0.381554in}{3.916545in}}%
\pgfpathlineto{\pgfqpoint{0.384120in}{3.039515in}}%
\pgfpathlineto{\pgfqpoint{0.384548in}{3.072355in}}%
\pgfpathlineto{\pgfqpoint{0.385661in}{3.407997in}}%
\pgfpathlineto{\pgfqpoint{0.390024in}{5.279073in}}%
\pgfpathlineto{\pgfqpoint{0.390709in}{5.200031in}}%
\pgfpathlineto{\pgfqpoint{0.392163in}{4.615990in}}%
\pgfpathlineto{\pgfqpoint{0.395928in}{3.039491in}}%
\pgfpathlineto{\pgfqpoint{0.396099in}{3.044771in}}%
\pgfpathlineto{\pgfqpoint{0.396784in}{3.155977in}}%
\pgfpathlineto{\pgfqpoint{0.398495in}{3.932557in}}%
\pgfpathlineto{\pgfqpoint{0.401832in}{5.279062in}}%
\pgfpathlineto{\pgfqpoint{0.402089in}{5.270017in}}%
\pgfpathlineto{\pgfqpoint{0.402944in}{5.095995in}}%
\pgfpathlineto{\pgfqpoint{0.404998in}{4.053573in}}%
\pgfpathlineto{\pgfqpoint{0.407821in}{3.039686in}}%
\pgfpathlineto{\pgfqpoint{0.408249in}{3.071819in}}%
\pgfpathlineto{\pgfqpoint{0.409361in}{3.401513in}}%
\pgfpathlineto{\pgfqpoint{0.413725in}{5.278884in}}%
\pgfpathlineto{\pgfqpoint{0.414495in}{5.194438in}}%
\pgfpathlineto{\pgfqpoint{0.415950in}{4.609880in}}%
\pgfpathlineto{\pgfqpoint{0.419714in}{3.039688in}}%
\pgfpathlineto{\pgfqpoint{0.419886in}{3.043651in}}%
\pgfpathlineto{\pgfqpoint{0.420570in}{3.148182in}}%
\pgfpathlineto{\pgfqpoint{0.422196in}{3.857061in}}%
\pgfpathlineto{\pgfqpoint{0.425704in}{5.278876in}}%
\pgfpathlineto{\pgfqpoint{0.425789in}{5.277953in}}%
\pgfpathlineto{\pgfqpoint{0.426303in}{5.225564in}}%
\pgfpathlineto{\pgfqpoint{0.427586in}{4.780809in}}%
\pgfpathlineto{\pgfqpoint{0.431693in}{3.039843in}}%
\pgfpathlineto{\pgfqpoint{0.432121in}{3.065312in}}%
\pgfpathlineto{\pgfqpoint{0.433148in}{3.340699in}}%
\pgfpathlineto{\pgfqpoint{0.437768in}{5.278492in}}%
\pgfpathlineto{\pgfqpoint{0.438624in}{5.158273in}}%
\pgfpathlineto{\pgfqpoint{0.440335in}{4.393046in}}%
\pgfpathlineto{\pgfqpoint{0.443758in}{3.039922in}}%
\pgfpathlineto{\pgfqpoint{0.444100in}{3.057947in}}%
\pgfpathlineto{\pgfqpoint{0.445041in}{3.281536in}}%
\pgfpathlineto{\pgfqpoint{0.447693in}{4.674135in}}%
\pgfpathlineto{\pgfqpoint{0.449832in}{5.278434in}}%
\pgfpathlineto{\pgfqpoint{0.449918in}{5.276413in}}%
\pgfpathlineto{\pgfqpoint{0.450517in}{5.201312in}}%
\pgfpathlineto{\pgfqpoint{0.451972in}{4.638106in}}%
\pgfpathlineto{\pgfqpoint{0.455822in}{3.040215in}}%
\pgfpathlineto{\pgfqpoint{0.456079in}{3.047649in}}%
\pgfpathlineto{\pgfqpoint{0.456849in}{3.185438in}}%
\pgfpathlineto{\pgfqpoint{0.458731in}{4.068803in}}%
\pgfpathlineto{\pgfqpoint{0.461897in}{5.278360in}}%
\pgfpathlineto{\pgfqpoint{0.462068in}{5.275410in}}%
\pgfpathlineto{\pgfqpoint{0.462752in}{5.177636in}}%
\pgfpathlineto{\pgfqpoint{0.464378in}{4.493604in}}%
\pgfpathlineto{\pgfqpoint{0.467972in}{3.040436in}}%
\pgfpathlineto{\pgfqpoint{0.468057in}{3.040548in}}%
\pgfpathlineto{\pgfqpoint{0.468485in}{3.073609in}}%
\pgfpathlineto{\pgfqpoint{0.469597in}{3.394377in}}%
\pgfpathlineto{\pgfqpoint{0.474132in}{5.278244in}}%
\pgfpathlineto{\pgfqpoint{0.474817in}{5.205216in}}%
\pgfpathlineto{\pgfqpoint{0.476271in}{4.655627in}}%
\pgfpathlineto{\pgfqpoint{0.480207in}{3.040427in}}%
\pgfpathlineto{\pgfqpoint{0.480464in}{3.048800in}}%
\pgfpathlineto{\pgfqpoint{0.481234in}{3.187069in}}%
\pgfpathlineto{\pgfqpoint{0.483116in}{4.061228in}}%
\pgfpathlineto{\pgfqpoint{0.486282in}{5.277387in}}%
\pgfpathlineto{\pgfqpoint{0.486368in}{5.278117in}}%
\pgfpathlineto{\pgfqpoint{0.486539in}{5.273142in}}%
\pgfpathlineto{\pgfqpoint{0.487223in}{5.169208in}}%
\pgfpathlineto{\pgfqpoint{0.488849in}{4.483715in}}%
\pgfpathlineto{\pgfqpoint{0.492528in}{3.040692in}}%
\pgfpathlineto{\pgfqpoint{0.492956in}{3.070993in}}%
\pgfpathlineto{\pgfqpoint{0.494068in}{3.381154in}}%
\pgfpathlineto{\pgfqpoint{0.498689in}{5.277817in}}%
\pgfpathlineto{\pgfqpoint{0.499459in}{5.186033in}}%
\pgfpathlineto{\pgfqpoint{0.500999in}{4.572621in}}%
\pgfpathlineto{\pgfqpoint{0.504849in}{3.040796in}}%
\pgfpathlineto{\pgfqpoint{0.504935in}{3.042291in}}%
\pgfpathlineto{\pgfqpoint{0.505534in}{3.111208in}}%
\pgfpathlineto{\pgfqpoint{0.506988in}{3.647565in}}%
\pgfpathlineto{\pgfqpoint{0.511010in}{5.277746in}}%
\pgfpathlineto{\pgfqpoint{0.511266in}{5.269108in}}%
\pgfpathlineto{\pgfqpoint{0.512122in}{5.108027in}}%
\pgfpathlineto{\pgfqpoint{0.514175in}{4.128578in}}%
\pgfpathlineto{\pgfqpoint{0.517170in}{3.041389in}}%
\pgfpathlineto{\pgfqpoint{0.517256in}{3.041085in}}%
\pgfpathlineto{\pgfqpoint{0.517341in}{3.042874in}}%
\pgfpathlineto{\pgfqpoint{0.517940in}{3.113220in}}%
\pgfpathlineto{\pgfqpoint{0.519395in}{3.648649in}}%
\pgfpathlineto{\pgfqpoint{0.523416in}{5.277434in}}%
\pgfpathlineto{\pgfqpoint{0.523673in}{5.270195in}}%
\pgfpathlineto{\pgfqpoint{0.524443in}{5.138963in}}%
\pgfpathlineto{\pgfqpoint{0.526240in}{4.339600in}}%
\pgfpathlineto{\pgfqpoint{0.529662in}{3.041217in}}%
\pgfpathlineto{\pgfqpoint{0.529919in}{3.049188in}}%
\pgfpathlineto{\pgfqpoint{0.530689in}{3.181952in}}%
\pgfpathlineto{\pgfqpoint{0.532571in}{4.028652in}}%
\pgfpathlineto{\pgfqpoint{0.535908in}{5.277255in}}%
\pgfpathlineto{\pgfqpoint{0.536079in}{5.274267in}}%
\pgfpathlineto{\pgfqpoint{0.536764in}{5.181261in}}%
\pgfpathlineto{\pgfqpoint{0.538390in}{4.529802in}}%
\pgfpathlineto{\pgfqpoint{0.542240in}{3.041636in}}%
\pgfpathlineto{\pgfqpoint{0.542668in}{3.072084in}}%
\pgfpathlineto{\pgfqpoint{0.543780in}{3.373679in}}%
\pgfpathlineto{\pgfqpoint{0.548486in}{5.277095in}}%
\pgfpathlineto{\pgfqpoint{0.549256in}{5.194485in}}%
\pgfpathlineto{\pgfqpoint{0.550796in}{4.610255in}}%
\pgfpathlineto{\pgfqpoint{0.554818in}{3.041836in}}%
\pgfpathlineto{\pgfqpoint{0.554903in}{3.043746in}}%
\pgfpathlineto{\pgfqpoint{0.555502in}{3.113180in}}%
\pgfpathlineto{\pgfqpoint{0.556957in}{3.636345in}}%
\pgfpathlineto{\pgfqpoint{0.561064in}{5.276719in}}%
\pgfpathlineto{\pgfqpoint{0.561320in}{5.269822in}}%
\pgfpathlineto{\pgfqpoint{0.562090in}{5.142849in}}%
\pgfpathlineto{\pgfqpoint{0.563887in}{4.363754in}}%
\pgfpathlineto{\pgfqpoint{0.567395in}{3.042017in}}%
\pgfpathlineto{\pgfqpoint{0.567566in}{3.044622in}}%
\pgfpathlineto{\pgfqpoint{0.568251in}{3.134092in}}%
\pgfpathlineto{\pgfqpoint{0.569791in}{3.725249in}}%
\pgfpathlineto{\pgfqpoint{0.573812in}{5.276372in}}%
\pgfpathlineto{\pgfqpoint{0.573898in}{5.274440in}}%
\pgfpathlineto{\pgfqpoint{0.574497in}{5.205729in}}%
\pgfpathlineto{\pgfqpoint{0.575952in}{4.689219in}}%
\pgfpathlineto{\pgfqpoint{0.580144in}{3.042169in}}%
\pgfpathlineto{\pgfqpoint{0.580315in}{3.046513in}}%
\pgfpathlineto{\pgfqpoint{0.581000in}{3.141873in}}%
\pgfpathlineto{\pgfqpoint{0.582625in}{3.782421in}}%
\pgfpathlineto{\pgfqpoint{0.586476in}{5.276093in}}%
\pgfpathlineto{\pgfqpoint{0.586561in}{5.276041in}}%
\pgfpathlineto{\pgfqpoint{0.586989in}{5.246263in}}%
\pgfpathlineto{\pgfqpoint{0.588101in}{4.954000in}}%
\pgfpathlineto{\pgfqpoint{0.592893in}{3.042497in}}%
\pgfpathlineto{\pgfqpoint{0.593749in}{3.134860in}}%
\pgfpathlineto{\pgfqpoint{0.595289in}{3.717957in}}%
\pgfpathlineto{\pgfqpoint{0.599310in}{5.275989in}}%
\pgfpathlineto{\pgfqpoint{0.599396in}{5.275393in}}%
\pgfpathlineto{\pgfqpoint{0.599909in}{5.231047in}}%
\pgfpathlineto{\pgfqpoint{0.601192in}{4.842662in}}%
\pgfpathlineto{\pgfqpoint{0.605727in}{3.042858in}}%
\pgfpathlineto{\pgfqpoint{0.606241in}{3.073623in}}%
\pgfpathlineto{\pgfqpoint{0.607353in}{3.364111in}}%
\pgfpathlineto{\pgfqpoint{0.612230in}{5.275642in}}%
\pgfpathlineto{\pgfqpoint{0.613086in}{5.174833in}}%
\pgfpathlineto{\pgfqpoint{0.614711in}{4.544129in}}%
\pgfpathlineto{\pgfqpoint{0.618647in}{3.043115in}}%
\pgfpathlineto{\pgfqpoint{0.618733in}{3.043565in}}%
\pgfpathlineto{\pgfqpoint{0.619246in}{3.086392in}}%
\pgfpathlineto{\pgfqpoint{0.620530in}{3.467278in}}%
\pgfpathlineto{\pgfqpoint{0.625150in}{5.275407in}}%
\pgfpathlineto{\pgfqpoint{0.625663in}{5.241433in}}%
\pgfpathlineto{\pgfqpoint{0.626861in}{4.915119in}}%
\pgfpathlineto{\pgfqpoint{0.631653in}{3.043370in}}%
\pgfpathlineto{\pgfqpoint{0.632337in}{3.103857in}}%
\pgfpathlineto{\pgfqpoint{0.633706in}{3.548095in}}%
\pgfpathlineto{\pgfqpoint{0.638155in}{5.275052in}}%
\pgfpathlineto{\pgfqpoint{0.638498in}{5.261155in}}%
\pgfpathlineto{\pgfqpoint{0.639439in}{5.073174in}}%
\pgfpathlineto{\pgfqpoint{0.641749in}{3.991286in}}%
\pgfpathlineto{\pgfqpoint{0.644573in}{3.045934in}}%
\pgfpathlineto{\pgfqpoint{0.644744in}{3.043909in}}%
\pgfpathlineto{\pgfqpoint{0.645000in}{3.054968in}}%
\pgfpathlineto{\pgfqpoint{0.645856in}{3.209884in}}%
\pgfpathlineto{\pgfqpoint{0.647910in}{4.122413in}}%
\pgfpathlineto{\pgfqpoint{0.651161in}{5.273547in}}%
\pgfpathlineto{\pgfqpoint{0.651246in}{5.274717in}}%
\pgfpathlineto{\pgfqpoint{0.651503in}{5.267004in}}%
\pgfpathlineto{\pgfqpoint{0.652359in}{5.123226in}}%
\pgfpathlineto{\pgfqpoint{0.654327in}{4.273101in}}%
\pgfpathlineto{\pgfqpoint{0.657664in}{3.047498in}}%
\pgfpathlineto{\pgfqpoint{0.657835in}{3.044067in}}%
\pgfpathlineto{\pgfqpoint{0.658177in}{3.059497in}}%
\pgfpathlineto{\pgfqpoint{0.659118in}{3.248901in}}%
\pgfpathlineto{\pgfqpoint{0.661428in}{4.322400in}}%
\pgfpathlineto{\pgfqpoint{0.664338in}{5.273567in}}%
\pgfpathlineto{\pgfqpoint{0.664423in}{5.274367in}}%
\pgfpathlineto{\pgfqpoint{0.664594in}{5.270417in}}%
\pgfpathlineto{\pgfqpoint{0.665279in}{5.181924in}}%
\pgfpathlineto{\pgfqpoint{0.666904in}{4.580424in}}%
\pgfpathlineto{\pgfqpoint{0.671011in}{3.044451in}}%
\pgfpathlineto{\pgfqpoint{0.671097in}{3.045084in}}%
\pgfpathlineto{\pgfqpoint{0.671610in}{3.087233in}}%
\pgfpathlineto{\pgfqpoint{0.672894in}{3.455031in}}%
\pgfpathlineto{\pgfqpoint{0.677685in}{5.273846in}}%
\pgfpathlineto{\pgfqpoint{0.678199in}{5.236701in}}%
\pgfpathlineto{\pgfqpoint{0.679397in}{4.914639in}}%
\pgfpathlineto{\pgfqpoint{0.684274in}{3.044870in}}%
\pgfpathlineto{\pgfqpoint{0.685044in}{3.113642in}}%
\pgfpathlineto{\pgfqpoint{0.686498in}{3.594737in}}%
\pgfpathlineto{\pgfqpoint{0.690947in}{5.273590in}}%
\pgfpathlineto{\pgfqpoint{0.691204in}{5.266168in}}%
\pgfpathlineto{\pgfqpoint{0.692060in}{5.127248in}}%
\pgfpathlineto{\pgfqpoint{0.694028in}{4.300248in}}%
\pgfpathlineto{\pgfqpoint{0.697536in}{3.046566in}}%
\pgfpathlineto{\pgfqpoint{0.697621in}{3.045257in}}%
\pgfpathlineto{\pgfqpoint{0.697878in}{3.052109in}}%
\pgfpathlineto{\pgfqpoint{0.698648in}{3.167410in}}%
\pgfpathlineto{\pgfqpoint{0.700445in}{3.878061in}}%
\pgfpathlineto{\pgfqpoint{0.704295in}{5.272929in}}%
\pgfpathlineto{\pgfqpoint{0.704381in}{5.273017in}}%
\pgfpathlineto{\pgfqpoint{0.704809in}{5.246760in}}%
\pgfpathlineto{\pgfqpoint{0.705921in}{4.982795in}}%
\pgfpathlineto{\pgfqpoint{0.711055in}{3.045633in}}%
\pgfpathlineto{\pgfqpoint{0.712081in}{3.169111in}}%
\pgfpathlineto{\pgfqpoint{0.713878in}{3.875789in}}%
\pgfpathlineto{\pgfqpoint{0.717729in}{5.272113in}}%
\pgfpathlineto{\pgfqpoint{0.717814in}{5.272761in}}%
\pgfpathlineto{\pgfqpoint{0.717985in}{5.268763in}}%
\pgfpathlineto{\pgfqpoint{0.718670in}{5.183377in}}%
\pgfpathlineto{\pgfqpoint{0.720295in}{4.604428in}}%
\pgfpathlineto{\pgfqpoint{0.724574in}{3.046070in}}%
\pgfpathlineto{\pgfqpoint{0.724659in}{3.047179in}}%
\pgfpathlineto{\pgfqpoint{0.725258in}{3.103572in}}%
\pgfpathlineto{\pgfqpoint{0.726627in}{3.518829in}}%
\pgfpathlineto{\pgfqpoint{0.731333in}{5.272330in}}%
\pgfpathlineto{\pgfqpoint{0.731761in}{5.251458in}}%
\pgfpathlineto{\pgfqpoint{0.732788in}{5.032472in}}%
\pgfpathlineto{\pgfqpoint{0.735440in}{3.805024in}}%
\pgfpathlineto{\pgfqpoint{0.738092in}{3.046776in}}%
\pgfpathlineto{\pgfqpoint{0.738178in}{3.046666in}}%
\pgfpathlineto{\pgfqpoint{0.738264in}{3.048290in}}%
\pgfpathlineto{\pgfqpoint{0.738862in}{3.107631in}}%
\pgfpathlineto{\pgfqpoint{0.740317in}{3.561487in}}%
\pgfpathlineto{\pgfqpoint{0.744937in}{5.271825in}}%
\pgfpathlineto{\pgfqpoint{0.745280in}{5.259863in}}%
\pgfpathlineto{\pgfqpoint{0.746221in}{5.089970in}}%
\pgfpathlineto{\pgfqpoint{0.748360in}{4.168089in}}%
\pgfpathlineto{\pgfqpoint{0.751611in}{3.051093in}}%
\pgfpathlineto{\pgfqpoint{0.751782in}{3.046997in}}%
\pgfpathlineto{\pgfqpoint{0.752210in}{3.066669in}}%
\pgfpathlineto{\pgfqpoint{0.753237in}{3.280041in}}%
\pgfpathlineto{\pgfqpoint{0.755804in}{4.449850in}}%
\pgfpathlineto{\pgfqpoint{0.758542in}{5.269818in}}%
\pgfpathlineto{\pgfqpoint{0.758627in}{5.271297in}}%
\pgfpathlineto{\pgfqpoint{0.758970in}{5.260210in}}%
\pgfpathlineto{\pgfqpoint{0.759911in}{5.094206in}}%
\pgfpathlineto{\pgfqpoint{0.762050in}{4.183077in}}%
\pgfpathlineto{\pgfqpoint{0.765387in}{3.050130in}}%
\pgfpathlineto{\pgfqpoint{0.765558in}{3.047484in}}%
\pgfpathlineto{\pgfqpoint{0.765900in}{3.062457in}}%
\pgfpathlineto{\pgfqpoint{0.766841in}{3.237560in}}%
\pgfpathlineto{\pgfqpoint{0.769066in}{4.199073in}}%
\pgfpathlineto{\pgfqpoint{0.772232in}{5.265931in}}%
\pgfpathlineto{\pgfqpoint{0.772403in}{5.270763in}}%
\pgfpathlineto{\pgfqpoint{0.772831in}{5.253477in}}%
\pgfpathlineto{\pgfqpoint{0.773857in}{5.048324in}}%
\pgfpathlineto{\pgfqpoint{0.776339in}{3.938916in}}%
\pgfpathlineto{\pgfqpoint{0.779248in}{3.049436in}}%
\pgfpathlineto{\pgfqpoint{0.779333in}{3.048036in}}%
\pgfpathlineto{\pgfqpoint{0.779676in}{3.059121in}}%
\pgfpathlineto{\pgfqpoint{0.780617in}{3.222644in}}%
\pgfpathlineto{\pgfqpoint{0.782756in}{4.121338in}}%
\pgfpathlineto{\pgfqpoint{0.786093in}{5.265295in}}%
\pgfpathlineto{\pgfqpoint{0.786264in}{5.270224in}}%
\pgfpathlineto{\pgfqpoint{0.786692in}{5.253545in}}%
\pgfpathlineto{\pgfqpoint{0.787719in}{5.051733in}}%
\pgfpathlineto{\pgfqpoint{0.790200in}{3.952738in}}%
\pgfpathlineto{\pgfqpoint{0.793109in}{3.051372in}}%
\pgfpathlineto{\pgfqpoint{0.793280in}{3.048466in}}%
\pgfpathlineto{\pgfqpoint{0.793622in}{3.062416in}}%
\pgfpathlineto{\pgfqpoint{0.794564in}{3.231622in}}%
\pgfpathlineto{\pgfqpoint{0.796788in}{4.173195in}}%
\pgfpathlineto{\pgfqpoint{0.800040in}{5.264662in}}%
\pgfpathlineto{\pgfqpoint{0.800296in}{5.269714in}}%
\pgfpathlineto{\pgfqpoint{0.800639in}{5.253549in}}%
\pgfpathlineto{\pgfqpoint{0.801665in}{5.054978in}}%
\pgfpathlineto{\pgfqpoint{0.804061in}{4.008500in}}%
\pgfpathlineto{\pgfqpoint{0.807056in}{3.053784in}}%
\pgfpathlineto{\pgfqpoint{0.807227in}{3.049135in}}%
\pgfpathlineto{\pgfqpoint{0.807655in}{3.065960in}}%
\pgfpathlineto{\pgfqpoint{0.808681in}{3.265075in}}%
\pgfpathlineto{\pgfqpoint{0.811077in}{4.308449in}}%
\pgfpathlineto{\pgfqpoint{0.814072in}{5.263958in}}%
\pgfpathlineto{\pgfqpoint{0.814329in}{5.269216in}}%
\pgfpathlineto{\pgfqpoint{0.814671in}{5.253617in}}%
\pgfpathlineto{\pgfqpoint{0.815698in}{5.058501in}}%
\pgfpathlineto{\pgfqpoint{0.818093in}{4.022591in}}%
\pgfpathlineto{\pgfqpoint{0.821174in}{3.052740in}}%
\pgfpathlineto{\pgfqpoint{0.821345in}{3.049535in}}%
\pgfpathlineto{\pgfqpoint{0.821687in}{3.062387in}}%
\pgfpathlineto{\pgfqpoint{0.822628in}{3.225489in}}%
\pgfpathlineto{\pgfqpoint{0.824767in}{4.104368in}}%
\pgfpathlineto{\pgfqpoint{0.828190in}{5.263106in}}%
\pgfpathlineto{\pgfqpoint{0.828446in}{5.268711in}}%
\pgfpathlineto{\pgfqpoint{0.828874in}{5.246200in}}%
\pgfpathlineto{\pgfqpoint{0.829986in}{5.011344in}}%
\pgfpathlineto{\pgfqpoint{0.832810in}{3.753366in}}%
\pgfpathlineto{\pgfqpoint{0.835462in}{3.050395in}}%
\pgfpathlineto{\pgfqpoint{0.835548in}{3.050226in}}%
\pgfpathlineto{\pgfqpoint{0.835634in}{3.051643in}}%
\pgfpathlineto{\pgfqpoint{0.836233in}{3.105502in}}%
\pgfpathlineto{\pgfqpoint{0.837602in}{3.489176in}}%
\pgfpathlineto{\pgfqpoint{0.842650in}{5.268191in}}%
\pgfpathlineto{\pgfqpoint{0.843078in}{5.247044in}}%
\pgfpathlineto{\pgfqpoint{0.844190in}{5.017656in}}%
\pgfpathlineto{\pgfqpoint{0.846928in}{3.811063in}}%
\pgfpathlineto{\pgfqpoint{0.849666in}{3.051963in}}%
\pgfpathlineto{\pgfqpoint{0.849751in}{3.050749in}}%
\pgfpathlineto{\pgfqpoint{0.850094in}{3.061530in}}%
\pgfpathlineto{\pgfqpoint{0.851035in}{3.216019in}}%
\pgfpathlineto{\pgfqpoint{0.853174in}{4.071546in}}%
\pgfpathlineto{\pgfqpoint{0.856682in}{5.260573in}}%
\pgfpathlineto{\pgfqpoint{0.856939in}{5.267624in}}%
\pgfpathlineto{\pgfqpoint{0.857366in}{5.248346in}}%
\pgfpathlineto{\pgfqpoint{0.858393in}{5.050884in}}%
\pgfpathlineto{\pgfqpoint{0.860875in}{3.993596in}}%
\pgfpathlineto{\pgfqpoint{0.863869in}{3.058344in}}%
\pgfpathlineto{\pgfqpoint{0.864126in}{3.051296in}}%
\pgfpathlineto{\pgfqpoint{0.864554in}{3.070369in}}%
\pgfpathlineto{\pgfqpoint{0.865580in}{3.266379in}}%
\pgfpathlineto{\pgfqpoint{0.868062in}{4.318442in}}%
\pgfpathlineto{\pgfqpoint{0.871142in}{5.262929in}}%
\pgfpathlineto{\pgfqpoint{0.871313in}{5.266936in}}%
\pgfpathlineto{\pgfqpoint{0.871741in}{5.250162in}}%
\pgfpathlineto{\pgfqpoint{0.872768in}{5.060208in}}%
\pgfpathlineto{\pgfqpoint{0.875163in}{4.058857in}}%
\pgfpathlineto{\pgfqpoint{0.878329in}{3.058183in}}%
\pgfpathlineto{\pgfqpoint{0.878586in}{3.051934in}}%
\pgfpathlineto{\pgfqpoint{0.879014in}{3.071914in}}%
\pgfpathlineto{\pgfqpoint{0.880126in}{3.292866in}}%
\pgfpathlineto{\pgfqpoint{0.882778in}{4.432884in}}%
\pgfpathlineto{\pgfqpoint{0.885688in}{5.264163in}}%
\pgfpathlineto{\pgfqpoint{0.885859in}{5.266310in}}%
\pgfpathlineto{\pgfqpoint{0.886201in}{5.252482in}}%
\pgfpathlineto{\pgfqpoint{0.887142in}{5.094285in}}%
\pgfpathlineto{\pgfqpoint{0.889281in}{4.250664in}}%
\pgfpathlineto{\pgfqpoint{0.892875in}{3.058592in}}%
\pgfpathlineto{\pgfqpoint{0.893132in}{3.052586in}}%
\pgfpathlineto{\pgfqpoint{0.893559in}{3.072552in}}%
\pgfpathlineto{\pgfqpoint{0.894672in}{3.291189in}}%
\pgfpathlineto{\pgfqpoint{0.897324in}{4.421136in}}%
\pgfpathlineto{\pgfqpoint{0.900233in}{5.261970in}}%
\pgfpathlineto{\pgfqpoint{0.900404in}{5.265675in}}%
\pgfpathlineto{\pgfqpoint{0.900832in}{5.248881in}}%
\pgfpathlineto{\pgfqpoint{0.901859in}{5.062904in}}%
\pgfpathlineto{\pgfqpoint{0.904255in}{4.080592in}}%
\pgfpathlineto{\pgfqpoint{0.907506in}{3.059623in}}%
\pgfpathlineto{\pgfqpoint{0.907763in}{3.053233in}}%
\pgfpathlineto{\pgfqpoint{0.908191in}{3.072143in}}%
\pgfpathlineto{\pgfqpoint{0.909217in}{3.261777in}}%
\pgfpathlineto{\pgfqpoint{0.911613in}{4.244111in}}%
\pgfpathlineto{\pgfqpoint{0.914864in}{5.258751in}}%
\pgfpathlineto{\pgfqpoint{0.915121in}{5.265040in}}%
\pgfpathlineto{\pgfqpoint{0.915549in}{5.246170in}}%
\pgfpathlineto{\pgfqpoint{0.916576in}{5.057627in}}%
\pgfpathlineto{\pgfqpoint{0.918971in}{4.080181in}}%
\pgfpathlineto{\pgfqpoint{0.922223in}{3.061418in}}%
\pgfpathlineto{\pgfqpoint{0.922479in}{3.053946in}}%
\pgfpathlineto{\pgfqpoint{0.922907in}{3.070645in}}%
\pgfpathlineto{\pgfqpoint{0.923934in}{3.253398in}}%
\pgfpathlineto{\pgfqpoint{0.926244in}{4.180922in}}%
\pgfpathlineto{\pgfqpoint{0.929581in}{5.253983in}}%
\pgfpathlineto{\pgfqpoint{0.929923in}{5.264315in}}%
\pgfpathlineto{\pgfqpoint{0.930351in}{5.244685in}}%
\pgfpathlineto{\pgfqpoint{0.931464in}{5.032542in}}%
\pgfpathlineto{\pgfqpoint{0.934030in}{3.968000in}}%
\pgfpathlineto{\pgfqpoint{0.937111in}{3.059709in}}%
\pgfpathlineto{\pgfqpoint{0.937367in}{3.054704in}}%
\pgfpathlineto{\pgfqpoint{0.937795in}{3.075081in}}%
\pgfpathlineto{\pgfqpoint{0.938907in}{3.287845in}}%
\pgfpathlineto{\pgfqpoint{0.941560in}{4.387927in}}%
\pgfpathlineto{\pgfqpoint{0.944555in}{5.257883in}}%
\pgfpathlineto{\pgfqpoint{0.944811in}{5.263613in}}%
\pgfpathlineto{\pgfqpoint{0.945239in}{5.244647in}}%
\pgfpathlineto{\pgfqpoint{0.946351in}{5.036351in}}%
\pgfpathlineto{\pgfqpoint{0.948918in}{3.983376in}}%
\pgfpathlineto{\pgfqpoint{0.951999in}{3.063160in}}%
\pgfpathlineto{\pgfqpoint{0.952255in}{3.055421in}}%
\pgfpathlineto{\pgfqpoint{0.952769in}{3.078146in}}%
\pgfpathlineto{\pgfqpoint{0.953881in}{3.293565in}}%
\pgfpathlineto{\pgfqpoint{0.956533in}{4.387500in}}%
\pgfpathlineto{\pgfqpoint{0.959528in}{5.256085in}}%
\pgfpathlineto{\pgfqpoint{0.959785in}{5.262886in}}%
\pgfpathlineto{\pgfqpoint{0.960213in}{5.246111in}}%
\pgfpathlineto{\pgfqpoint{0.961239in}{5.068208in}}%
\pgfpathlineto{\pgfqpoint{0.963549in}{4.164261in}}%
\pgfpathlineto{\pgfqpoint{0.966972in}{3.068480in}}%
\pgfpathlineto{\pgfqpoint{0.967314in}{3.056129in}}%
\pgfpathlineto{\pgfqpoint{0.967828in}{3.079416in}}%
\pgfpathlineto{\pgfqpoint{0.968940in}{3.293538in}}%
\pgfpathlineto{\pgfqpoint{0.971592in}{4.377909in}}%
\pgfpathlineto{\pgfqpoint{0.974587in}{5.253011in}}%
\pgfpathlineto{\pgfqpoint{0.974929in}{5.262039in}}%
\pgfpathlineto{\pgfqpoint{0.975357in}{5.242186in}}%
\pgfpathlineto{\pgfqpoint{0.976469in}{5.036353in}}%
\pgfpathlineto{\pgfqpoint{0.979036in}{4.002091in}}%
\pgfpathlineto{\pgfqpoint{0.982202in}{3.064642in}}%
\pgfpathlineto{\pgfqpoint{0.982459in}{3.056958in}}%
\pgfpathlineto{\pgfqpoint{0.982972in}{3.078676in}}%
\pgfpathlineto{\pgfqpoint{0.984084in}{3.287327in}}%
\pgfpathlineto{\pgfqpoint{0.986737in}{4.358369in}}%
\pgfpathlineto{\pgfqpoint{0.989817in}{5.253389in}}%
\pgfpathlineto{\pgfqpoint{0.990074in}{5.261234in}}%
\pgfpathlineto{\pgfqpoint{0.990587in}{5.240118in}}%
\pgfpathlineto{\pgfqpoint{0.991700in}{5.033846in}}%
\pgfpathlineto{\pgfqpoint{0.994266in}{4.007484in}}%
\pgfpathlineto{\pgfqpoint{0.997432in}{3.067405in}}%
\pgfpathlineto{\pgfqpoint{0.997774in}{3.057691in}}%
\pgfpathlineto{\pgfqpoint{0.998202in}{3.075990in}}%
\pgfpathlineto{\pgfqpoint{0.999315in}{3.274727in}}%
\pgfpathlineto{\pgfqpoint{1.001796in}{4.251475in}}%
\pgfpathlineto{\pgfqpoint{1.005133in}{5.252324in}}%
\pgfpathlineto{\pgfqpoint{1.005389in}{5.260380in}}%
\pgfpathlineto{\pgfqpoint{1.005903in}{5.240244in}}%
\pgfpathlineto{\pgfqpoint{1.007015in}{5.038295in}}%
\pgfpathlineto{\pgfqpoint{1.009582in}{4.024317in}}%
\pgfpathlineto{\pgfqpoint{1.012833in}{3.066802in}}%
\pgfpathlineto{\pgfqpoint{1.013176in}{3.058624in}}%
\pgfpathlineto{\pgfqpoint{1.013603in}{3.078371in}}%
\pgfpathlineto{\pgfqpoint{1.014716in}{3.278266in}}%
\pgfpathlineto{\pgfqpoint{1.017283in}{4.286093in}}%
\pgfpathlineto{\pgfqpoint{1.020534in}{5.249647in}}%
\pgfpathlineto{\pgfqpoint{1.020876in}{5.259712in}}%
\pgfpathlineto{\pgfqpoint{1.021304in}{5.242548in}}%
\pgfpathlineto{\pgfqpoint{1.022331in}{5.071793in}}%
\pgfpathlineto{\pgfqpoint{1.024641in}{4.205628in}}%
\pgfpathlineto{\pgfqpoint{1.028235in}{3.073248in}}%
\pgfpathlineto{\pgfqpoint{1.028662in}{3.059420in}}%
\pgfpathlineto{\pgfqpoint{1.029090in}{3.078369in}}%
\pgfpathlineto{\pgfqpoint{1.030203in}{3.274101in}}%
\pgfpathlineto{\pgfqpoint{1.032684in}{4.231646in}}%
\pgfpathlineto{\pgfqpoint{1.036021in}{5.244763in}}%
\pgfpathlineto{\pgfqpoint{1.036449in}{5.258789in}}%
\pgfpathlineto{\pgfqpoint{1.036876in}{5.240294in}}%
\pgfpathlineto{\pgfqpoint{1.037989in}{5.046779in}}%
\pgfpathlineto{\pgfqpoint{1.040470in}{4.095528in}}%
\pgfpathlineto{\pgfqpoint{1.043893in}{3.070762in}}%
\pgfpathlineto{\pgfqpoint{1.044235in}{3.060184in}}%
\pgfpathlineto{\pgfqpoint{1.044748in}{3.083025in}}%
\pgfpathlineto{\pgfqpoint{1.045861in}{3.284320in}}%
\pgfpathlineto{\pgfqpoint{1.048427in}{4.276613in}}%
\pgfpathlineto{\pgfqpoint{1.051764in}{5.248800in}}%
\pgfpathlineto{\pgfqpoint{1.052107in}{5.257961in}}%
\pgfpathlineto{\pgfqpoint{1.052534in}{5.240597in}}%
\pgfpathlineto{\pgfqpoint{1.053647in}{5.052050in}}%
\pgfpathlineto{\pgfqpoint{1.056128in}{4.114051in}}%
\pgfpathlineto{\pgfqpoint{1.059551in}{3.076188in}}%
\pgfpathlineto{\pgfqpoint{1.059978in}{3.061081in}}%
\pgfpathlineto{\pgfqpoint{1.060492in}{3.084836in}}%
\pgfpathlineto{\pgfqpoint{1.061690in}{3.308676in}}%
\pgfpathlineto{\pgfqpoint{1.064513in}{4.415787in}}%
\pgfpathlineto{\pgfqpoint{1.067593in}{5.250113in}}%
\pgfpathlineto{\pgfqpoint{1.067850in}{5.257022in}}%
\pgfpathlineto{\pgfqpoint{1.068363in}{5.236825in}}%
\pgfpathlineto{\pgfqpoint{1.069476in}{5.044327in}}%
\pgfpathlineto{\pgfqpoint{1.071957in}{4.110566in}}%
\pgfpathlineto{\pgfqpoint{1.075380in}{3.078443in}}%
\pgfpathlineto{\pgfqpoint{1.075807in}{3.061996in}}%
\pgfpathlineto{\pgfqpoint{1.076321in}{3.083471in}}%
\pgfpathlineto{\pgfqpoint{1.077433in}{3.277302in}}%
\pgfpathlineto{\pgfqpoint{1.080000in}{4.245196in}}%
\pgfpathlineto{\pgfqpoint{1.083422in}{5.244683in}}%
\pgfpathlineto{\pgfqpoint{1.083765in}{5.256106in}}%
\pgfpathlineto{\pgfqpoint{1.084278in}{5.236068in}}%
\pgfpathlineto{\pgfqpoint{1.085390in}{5.046271in}}%
\pgfpathlineto{\pgfqpoint{1.087872in}{4.123558in}}%
\pgfpathlineto{\pgfqpoint{1.091380in}{3.077325in}}%
\pgfpathlineto{\pgfqpoint{1.091808in}{3.062940in}}%
\pgfpathlineto{\pgfqpoint{1.092321in}{3.086207in}}%
\pgfpathlineto{\pgfqpoint{1.093519in}{3.303663in}}%
\pgfpathlineto{\pgfqpoint{1.096257in}{4.350735in}}%
\pgfpathlineto{\pgfqpoint{1.099423in}{5.241460in}}%
\pgfpathlineto{\pgfqpoint{1.099850in}{5.255172in}}%
\pgfpathlineto{\pgfqpoint{1.100364in}{5.231438in}}%
\pgfpathlineto{\pgfqpoint{1.101562in}{5.014336in}}%
\pgfpathlineto{\pgfqpoint{1.104300in}{3.972183in}}%
\pgfpathlineto{\pgfqpoint{1.107551in}{3.073731in}}%
\pgfpathlineto{\pgfqpoint{1.107893in}{3.063890in}}%
\pgfpathlineto{\pgfqpoint{1.108407in}{3.085372in}}%
\pgfpathlineto{\pgfqpoint{1.109519in}{3.274452in}}%
\pgfpathlineto{\pgfqpoint{1.112000in}{4.183927in}}%
\pgfpathlineto{\pgfqpoint{1.115594in}{5.241001in}}%
\pgfpathlineto{\pgfqpoint{1.116022in}{5.254177in}}%
\pgfpathlineto{\pgfqpoint{1.116450in}{5.237399in}}%
\pgfpathlineto{\pgfqpoint{1.117562in}{5.059326in}}%
\pgfpathlineto{\pgfqpoint{1.119958in}{4.201937in}}%
\pgfpathlineto{\pgfqpoint{1.123722in}{3.078271in}}%
\pgfpathlineto{\pgfqpoint{1.124150in}{3.064920in}}%
\pgfpathlineto{\pgfqpoint{1.124664in}{3.088076in}}%
\pgfpathlineto{\pgfqpoint{1.125861in}{3.299950in}}%
\pgfpathlineto{\pgfqpoint{1.128599in}{4.323605in}}%
\pgfpathlineto{\pgfqpoint{1.131851in}{5.237402in}}%
\pgfpathlineto{\pgfqpoint{1.132279in}{5.253206in}}%
\pgfpathlineto{\pgfqpoint{1.132792in}{5.233322in}}%
\pgfpathlineto{\pgfqpoint{1.133904in}{5.050976in}}%
\pgfpathlineto{\pgfqpoint{1.136386in}{4.161219in}}%
\pgfpathlineto{\pgfqpoint{1.140065in}{3.080576in}}%
\pgfpathlineto{\pgfqpoint{1.140493in}{3.065901in}}%
\pgfpathlineto{\pgfqpoint{1.141006in}{3.086803in}}%
\pgfpathlineto{\pgfqpoint{1.142118in}{3.269978in}}%
\pgfpathlineto{\pgfqpoint{1.144600in}{4.156413in}}%
\pgfpathlineto{\pgfqpoint{1.148279in}{5.236302in}}%
\pgfpathlineto{\pgfqpoint{1.148707in}{5.252159in}}%
\pgfpathlineto{\pgfqpoint{1.149220in}{5.233010in}}%
\pgfpathlineto{\pgfqpoint{1.150332in}{5.054501in}}%
\pgfpathlineto{\pgfqpoint{1.152728in}{4.212873in}}%
\pgfpathlineto{\pgfqpoint{1.156493in}{3.086695in}}%
\pgfpathlineto{\pgfqpoint{1.157006in}{3.066959in}}%
\pgfpathlineto{\pgfqpoint{1.157520in}{3.088484in}}%
\pgfpathlineto{\pgfqpoint{1.158717in}{3.291574in}}%
\pgfpathlineto{\pgfqpoint{1.161370in}{4.253233in}}%
\pgfpathlineto{\pgfqpoint{1.164792in}{5.231324in}}%
\pgfpathlineto{\pgfqpoint{1.165306in}{5.251126in}}%
\pgfpathlineto{\pgfqpoint{1.165819in}{5.230031in}}%
\pgfpathlineto{\pgfqpoint{1.167017in}{5.029188in}}%
\pgfpathlineto{\pgfqpoint{1.169669in}{4.074257in}}%
\pgfpathlineto{\pgfqpoint{1.173177in}{3.084157in}}%
\pgfpathlineto{\pgfqpoint{1.173605in}{3.068094in}}%
\pgfpathlineto{\pgfqpoint{1.174119in}{3.085993in}}%
\pgfpathlineto{\pgfqpoint{1.175231in}{3.258445in}}%
\pgfpathlineto{\pgfqpoint{1.177627in}{4.081097in}}%
\pgfpathlineto{\pgfqpoint{1.181477in}{5.228499in}}%
\pgfpathlineto{\pgfqpoint{1.181990in}{5.250015in}}%
\pgfpathlineto{\pgfqpoint{1.182504in}{5.231362in}}%
\pgfpathlineto{\pgfqpoint{1.183616in}{5.058590in}}%
\pgfpathlineto{\pgfqpoint{1.186012in}{4.239678in}}%
\pgfpathlineto{\pgfqpoint{1.189948in}{3.085852in}}%
\pgfpathlineto{\pgfqpoint{1.190376in}{3.069267in}}%
\pgfpathlineto{\pgfqpoint{1.190889in}{3.085872in}}%
\pgfpathlineto{\pgfqpoint{1.192001in}{3.253364in}}%
\pgfpathlineto{\pgfqpoint{1.194397in}{4.062154in}}%
\pgfpathlineto{\pgfqpoint{1.198333in}{5.227809in}}%
\pgfpathlineto{\pgfqpoint{1.198846in}{5.248897in}}%
\pgfpathlineto{\pgfqpoint{1.199360in}{5.230543in}}%
\pgfpathlineto{\pgfqpoint{1.200472in}{5.060738in}}%
\pgfpathlineto{\pgfqpoint{1.202868in}{4.253417in}}%
\pgfpathlineto{\pgfqpoint{1.206804in}{3.092144in}}%
\pgfpathlineto{\pgfqpoint{1.207317in}{3.070332in}}%
\pgfpathlineto{\pgfqpoint{1.207830in}{3.087597in}}%
\pgfpathlineto{\pgfqpoint{1.208943in}{3.254044in}}%
\pgfpathlineto{\pgfqpoint{1.211338in}{4.053372in}}%
\pgfpathlineto{\pgfqpoint{1.215360in}{5.228872in}}%
\pgfpathlineto{\pgfqpoint{1.215873in}{5.247767in}}%
\pgfpathlineto{\pgfqpoint{1.216387in}{5.227947in}}%
\pgfpathlineto{\pgfqpoint{1.217499in}{5.057576in}}%
\pgfpathlineto{\pgfqpoint{1.219895in}{4.257854in}}%
\pgfpathlineto{\pgfqpoint{1.223916in}{3.090277in}}%
\pgfpathlineto{\pgfqpoint{1.224429in}{3.071433in}}%
\pgfpathlineto{\pgfqpoint{1.224943in}{3.090942in}}%
\pgfpathlineto{\pgfqpoint{1.226055in}{3.259511in}}%
\pgfpathlineto{\pgfqpoint{1.228451in}{4.053027in}}%
\pgfpathlineto{\pgfqpoint{1.232472in}{5.224883in}}%
\pgfpathlineto{\pgfqpoint{1.232986in}{5.246497in}}%
\pgfpathlineto{\pgfqpoint{1.233499in}{5.230122in}}%
\pgfpathlineto{\pgfqpoint{1.234611in}{5.069004in}}%
\pgfpathlineto{\pgfqpoint{1.236922in}{4.321904in}}%
\pgfpathlineto{\pgfqpoint{1.241114in}{3.093596in}}%
\pgfpathlineto{\pgfqpoint{1.241627in}{3.072687in}}%
\pgfpathlineto{\pgfqpoint{1.242141in}{3.089408in}}%
\pgfpathlineto{\pgfqpoint{1.243253in}{3.250039in}}%
\pgfpathlineto{\pgfqpoint{1.245563in}{3.992588in}}%
\pgfpathlineto{\pgfqpoint{1.249841in}{5.228475in}}%
\pgfpathlineto{\pgfqpoint{1.250355in}{5.245331in}}%
\pgfpathlineto{\pgfqpoint{1.250868in}{5.224922in}}%
\pgfpathlineto{\pgfqpoint{1.252066in}{5.038826in}}%
\pgfpathlineto{\pgfqpoint{1.254633in}{4.178767in}}%
\pgfpathlineto{\pgfqpoint{1.258483in}{3.095879in}}%
\pgfpathlineto{\pgfqpoint{1.259082in}{3.073961in}}%
\pgfpathlineto{\pgfqpoint{1.259596in}{3.095137in}}%
\pgfpathlineto{\pgfqpoint{1.260793in}{3.281546in}}%
\pgfpathlineto{\pgfqpoint{1.263360in}{4.137383in}}%
\pgfpathlineto{\pgfqpoint{1.267211in}{5.220201in}}%
\pgfpathlineto{\pgfqpoint{1.267810in}{5.244119in}}%
\pgfpathlineto{\pgfqpoint{1.268323in}{5.225046in}}%
\pgfpathlineto{\pgfqpoint{1.269521in}{5.044451in}}%
\pgfpathlineto{\pgfqpoint{1.272002in}{4.232447in}}%
\pgfpathlineto{\pgfqpoint{1.276024in}{3.097384in}}%
\pgfpathlineto{\pgfqpoint{1.276622in}{3.075181in}}%
\pgfpathlineto{\pgfqpoint{1.277136in}{3.095336in}}%
\pgfpathlineto{\pgfqpoint{1.278334in}{3.276922in}}%
\pgfpathlineto{\pgfqpoint{1.280901in}{4.118518in}}%
\pgfpathlineto{\pgfqpoint{1.284836in}{5.218982in}}%
\pgfpathlineto{\pgfqpoint{1.285435in}{5.242846in}}%
\pgfpathlineto{\pgfqpoint{1.285949in}{5.224503in}}%
\pgfpathlineto{\pgfqpoint{1.287061in}{5.066496in}}%
\pgfpathlineto{\pgfqpoint{1.289457in}{4.314640in}}%
\pgfpathlineto{\pgfqpoint{1.293735in}{3.098439in}}%
\pgfpathlineto{\pgfqpoint{1.294334in}{3.076467in}}%
\pgfpathlineto{\pgfqpoint{1.294847in}{3.096043in}}%
\pgfpathlineto{\pgfqpoint{1.296045in}{3.273754in}}%
\pgfpathlineto{\pgfqpoint{1.298526in}{4.069430in}}%
\pgfpathlineto{\pgfqpoint{1.302633in}{5.218052in}}%
\pgfpathlineto{\pgfqpoint{1.303232in}{5.241533in}}%
\pgfpathlineto{\pgfqpoint{1.303746in}{5.223637in}}%
\pgfpathlineto{\pgfqpoint{1.304858in}{5.068868in}}%
\pgfpathlineto{\pgfqpoint{1.307168in}{4.361338in}}%
\pgfpathlineto{\pgfqpoint{1.311618in}{3.099391in}}%
\pgfpathlineto{\pgfqpoint{1.312216in}{3.077795in}}%
\pgfpathlineto{\pgfqpoint{1.312730in}{3.096918in}}%
\pgfpathlineto{\pgfqpoint{1.313928in}{3.271027in}}%
\pgfpathlineto{\pgfqpoint{1.316409in}{4.053903in}}%
\pgfpathlineto{\pgfqpoint{1.320602in}{5.217052in}}%
\pgfpathlineto{\pgfqpoint{1.321201in}{5.240184in}}%
\pgfpathlineto{\pgfqpoint{1.321714in}{5.222763in}}%
\pgfpathlineto{\pgfqpoint{1.322826in}{5.071285in}}%
\pgfpathlineto{\pgfqpoint{1.325136in}{4.375829in}}%
\pgfpathlineto{\pgfqpoint{1.329671in}{3.100582in}}%
\pgfpathlineto{\pgfqpoint{1.330270in}{3.079150in}}%
\pgfpathlineto{\pgfqpoint{1.330784in}{3.097644in}}%
\pgfpathlineto{\pgfqpoint{1.331981in}{3.267771in}}%
\pgfpathlineto{\pgfqpoint{1.334463in}{4.037320in}}%
\pgfpathlineto{\pgfqpoint{1.338741in}{5.215632in}}%
\pgfpathlineto{\pgfqpoint{1.339340in}{5.238798in}}%
\pgfpathlineto{\pgfqpoint{1.339853in}{5.222178in}}%
\pgfpathlineto{\pgfqpoint{1.340965in}{5.074644in}}%
\pgfpathlineto{\pgfqpoint{1.343276in}{4.392160in}}%
\pgfpathlineto{\pgfqpoint{1.347896in}{3.102348in}}%
\pgfpathlineto{\pgfqpoint{1.348495in}{3.080530in}}%
\pgfpathlineto{\pgfqpoint{1.349008in}{3.097923in}}%
\pgfpathlineto{\pgfqpoint{1.350121in}{3.245849in}}%
\pgfpathlineto{\pgfqpoint{1.352431in}{3.924787in}}%
\pgfpathlineto{\pgfqpoint{1.357051in}{5.213434in}}%
\pgfpathlineto{\pgfqpoint{1.357650in}{5.237351in}}%
\pgfpathlineto{\pgfqpoint{1.358164in}{5.222144in}}%
\pgfpathlineto{\pgfqpoint{1.359276in}{5.079793in}}%
\pgfpathlineto{\pgfqpoint{1.361586in}{4.411957in}}%
\pgfpathlineto{\pgfqpoint{1.366292in}{3.105056in}}%
\pgfpathlineto{\pgfqpoint{1.366891in}{3.081973in}}%
\pgfpathlineto{\pgfqpoint{1.367404in}{3.097514in}}%
\pgfpathlineto{\pgfqpoint{1.368517in}{3.239377in}}%
\pgfpathlineto{\pgfqpoint{1.370827in}{3.902534in}}%
\pgfpathlineto{\pgfqpoint{1.375618in}{5.216391in}}%
\pgfpathlineto{\pgfqpoint{1.376217in}{5.235843in}}%
\pgfpathlineto{\pgfqpoint{1.376731in}{5.217572in}}%
\pgfpathlineto{\pgfqpoint{1.377928in}{5.054468in}}%
\pgfpathlineto{\pgfqpoint{1.380410in}{4.315308in}}%
\pgfpathlineto{\pgfqpoint{1.384859in}{3.109147in}}%
\pgfpathlineto{\pgfqpoint{1.385543in}{3.083509in}}%
\pgfpathlineto{\pgfqpoint{1.386057in}{3.101485in}}%
\pgfpathlineto{\pgfqpoint{1.387255in}{3.262648in}}%
\pgfpathlineto{\pgfqpoint{1.389650in}{3.964639in}}%
\pgfpathlineto{\pgfqpoint{1.394185in}{5.204961in}}%
\pgfpathlineto{\pgfqpoint{1.394870in}{5.234396in}}%
\pgfpathlineto{\pgfqpoint{1.395383in}{5.219685in}}%
\pgfpathlineto{\pgfqpoint{1.396495in}{5.082981in}}%
\pgfpathlineto{\pgfqpoint{1.398720in}{4.468187in}}%
\pgfpathlineto{\pgfqpoint{1.403768in}{3.102518in}}%
\pgfpathlineto{\pgfqpoint{1.404282in}{3.085022in}}%
\pgfpathlineto{\pgfqpoint{1.404795in}{3.098742in}}%
\pgfpathlineto{\pgfqpoint{1.405907in}{3.232277in}}%
\pgfpathlineto{\pgfqpoint{1.408132in}{3.839107in}}%
\pgfpathlineto{\pgfqpoint{1.413266in}{5.216892in}}%
\pgfpathlineto{\pgfqpoint{1.413779in}{5.232869in}}%
\pgfpathlineto{\pgfqpoint{1.414292in}{5.217981in}}%
\pgfpathlineto{\pgfqpoint{1.415405in}{5.083227in}}%
\pgfpathlineto{\pgfqpoint{1.417629in}{4.478701in}}%
\pgfpathlineto{\pgfqpoint{1.422763in}{3.104097in}}%
\pgfpathlineto{\pgfqpoint{1.423277in}{3.086608in}}%
\pgfpathlineto{\pgfqpoint{1.423790in}{3.099634in}}%
\pgfpathlineto{\pgfqpoint{1.424902in}{3.229448in}}%
\pgfpathlineto{\pgfqpoint{1.427127in}{3.823653in}}%
\pgfpathlineto{\pgfqpoint{1.432346in}{5.214277in}}%
\pgfpathlineto{\pgfqpoint{1.432860in}{5.231233in}}%
\pgfpathlineto{\pgfqpoint{1.433373in}{5.218024in}}%
\pgfpathlineto{\pgfqpoint{1.434485in}{5.088989in}}%
\pgfpathlineto{\pgfqpoint{1.436710in}{4.499699in}}%
\pgfpathlineto{\pgfqpoint{1.441929in}{3.107926in}}%
\pgfpathlineto{\pgfqpoint{1.442528in}{3.088133in}}%
\pgfpathlineto{\pgfqpoint{1.443041in}{3.103456in}}%
\pgfpathlineto{\pgfqpoint{1.444154in}{3.235642in}}%
\pgfpathlineto{\pgfqpoint{1.446378in}{3.825183in}}%
\pgfpathlineto{\pgfqpoint{1.451598in}{5.209184in}}%
\pgfpathlineto{\pgfqpoint{1.452197in}{5.229663in}}%
\pgfpathlineto{\pgfqpoint{1.452710in}{5.215306in}}%
\pgfpathlineto{\pgfqpoint{1.453822in}{5.086229in}}%
\pgfpathlineto{\pgfqpoint{1.456047in}{4.504667in}}%
\pgfpathlineto{\pgfqpoint{1.461352in}{3.108826in}}%
\pgfpathlineto{\pgfqpoint{1.461951in}{3.089785in}}%
\pgfpathlineto{\pgfqpoint{1.462464in}{3.105000in}}%
\pgfpathlineto{\pgfqpoint{1.463576in}{3.234675in}}%
\pgfpathlineto{\pgfqpoint{1.465801in}{3.813141in}}%
\pgfpathlineto{\pgfqpoint{1.471192in}{5.212191in}}%
\pgfpathlineto{\pgfqpoint{1.471705in}{5.227963in}}%
\pgfpathlineto{\pgfqpoint{1.472218in}{5.214960in}}%
\pgfpathlineto{\pgfqpoint{1.473331in}{5.090934in}}%
\pgfpathlineto{\pgfqpoint{1.475555in}{4.523891in}}%
\pgfpathlineto{\pgfqpoint{1.481031in}{3.107440in}}%
\pgfpathlineto{\pgfqpoint{1.481545in}{3.091521in}}%
\pgfpathlineto{\pgfqpoint{1.482058in}{3.104031in}}%
\pgfpathlineto{\pgfqpoint{1.483170in}{3.225909in}}%
\pgfpathlineto{\pgfqpoint{1.485395in}{3.786156in}}%
\pgfpathlineto{\pgfqpoint{1.490956in}{5.211825in}}%
\pgfpathlineto{\pgfqpoint{1.491470in}{5.226278in}}%
\pgfpathlineto{\pgfqpoint{1.491983in}{5.212648in}}%
\pgfpathlineto{\pgfqpoint{1.493095in}{5.089622in}}%
\pgfpathlineto{\pgfqpoint{1.495320in}{4.531725in}}%
\pgfpathlineto{\pgfqpoint{1.500882in}{3.109392in}}%
\pgfpathlineto{\pgfqpoint{1.501395in}{3.093298in}}%
\pgfpathlineto{\pgfqpoint{1.501908in}{3.104942in}}%
\pgfpathlineto{\pgfqpoint{1.503021in}{3.222770in}}%
\pgfpathlineto{\pgfqpoint{1.505245in}{3.769691in}}%
\pgfpathlineto{\pgfqpoint{1.510892in}{5.208183in}}%
\pgfpathlineto{\pgfqpoint{1.511491in}{5.224443in}}%
\pgfpathlineto{\pgfqpoint{1.512005in}{5.208714in}}%
\pgfpathlineto{\pgfqpoint{1.513203in}{5.069206in}}%
\pgfpathlineto{\pgfqpoint{1.515598in}{4.449838in}}%
\pgfpathlineto{\pgfqpoint{1.520818in}{3.120914in}}%
\pgfpathlineto{\pgfqpoint{1.521502in}{3.095031in}}%
\pgfpathlineto{\pgfqpoint{1.522016in}{3.107149in}}%
\pgfpathlineto{\pgfqpoint{1.523128in}{3.223677in}}%
\pgfpathlineto{\pgfqpoint{1.525353in}{3.760991in}}%
\pgfpathlineto{\pgfqpoint{1.531171in}{5.210352in}}%
\pgfpathlineto{\pgfqpoint{1.531684in}{5.222689in}}%
\pgfpathlineto{\pgfqpoint{1.532198in}{5.208318in}}%
\pgfpathlineto{\pgfqpoint{1.533395in}{5.074325in}}%
\pgfpathlineto{\pgfqpoint{1.535706in}{4.497443in}}%
\pgfpathlineto{\pgfqpoint{1.541182in}{3.119883in}}%
\pgfpathlineto{\pgfqpoint{1.541866in}{3.096831in}}%
\pgfpathlineto{\pgfqpoint{1.542379in}{3.110281in}}%
\pgfpathlineto{\pgfqpoint{1.543492in}{3.227272in}}%
\pgfpathlineto{\pgfqpoint{1.545716in}{3.757416in}}%
\pgfpathlineto{\pgfqpoint{1.551535in}{5.203995in}}%
\pgfpathlineto{\pgfqpoint{1.552134in}{5.220837in}}%
\pgfpathlineto{\pgfqpoint{1.552647in}{5.207086in}}%
\pgfpathlineto{\pgfqpoint{1.553759in}{5.090604in}}%
\pgfpathlineto{\pgfqpoint{1.555984in}{4.565066in}}%
\pgfpathlineto{\pgfqpoint{1.561888in}{3.113570in}}%
\pgfpathlineto{\pgfqpoint{1.562401in}{3.098789in}}%
\pgfpathlineto{\pgfqpoint{1.562914in}{3.109696in}}%
\pgfpathlineto{\pgfqpoint{1.564027in}{3.219223in}}%
\pgfpathlineto{\pgfqpoint{1.566166in}{3.706138in}}%
\pgfpathlineto{\pgfqpoint{1.572412in}{5.210742in}}%
\pgfpathlineto{\pgfqpoint{1.572840in}{5.218925in}}%
\pgfpathlineto{\pgfqpoint{1.573353in}{5.205504in}}%
\pgfpathlineto{\pgfqpoint{1.574465in}{5.091944in}}%
\pgfpathlineto{\pgfqpoint{1.576690in}{4.578282in}}%
\pgfpathlineto{\pgfqpoint{1.582765in}{3.112957in}}%
\pgfpathlineto{\pgfqpoint{1.583278in}{3.100655in}}%
\pgfpathlineto{\pgfqpoint{1.583792in}{3.113364in}}%
\pgfpathlineto{\pgfqpoint{1.584904in}{3.224347in}}%
\pgfpathlineto{\pgfqpoint{1.587129in}{3.730482in}}%
\pgfpathlineto{\pgfqpoint{1.593289in}{5.205240in}}%
\pgfpathlineto{\pgfqpoint{1.593803in}{5.216963in}}%
\pgfpathlineto{\pgfqpoint{1.594316in}{5.204012in}}%
\pgfpathlineto{\pgfqpoint{1.595428in}{5.093651in}}%
\pgfpathlineto{\pgfqpoint{1.597653in}{4.592299in}}%
\pgfpathlineto{\pgfqpoint{1.603813in}{3.117076in}}%
\pgfpathlineto{\pgfqpoint{1.604412in}{3.102721in}}%
\pgfpathlineto{\pgfqpoint{1.604926in}{3.116774in}}%
\pgfpathlineto{\pgfqpoint{1.606124in}{3.241213in}}%
\pgfpathlineto{\pgfqpoint{1.608434in}{3.776582in}}%
\pgfpathlineto{\pgfqpoint{1.614337in}{5.193554in}}%
\pgfpathlineto{\pgfqpoint{1.615022in}{5.214950in}}%
\pgfpathlineto{\pgfqpoint{1.615535in}{5.203013in}}%
\pgfpathlineto{\pgfqpoint{1.616648in}{5.096962in}}%
\pgfpathlineto{\pgfqpoint{1.618787in}{4.633372in}}%
\pgfpathlineto{\pgfqpoint{1.625289in}{3.113969in}}%
\pgfpathlineto{\pgfqpoint{1.625717in}{3.104728in}}%
\pgfpathlineto{\pgfqpoint{1.626231in}{3.115340in}}%
\pgfpathlineto{\pgfqpoint{1.627343in}{3.217556in}}%
\pgfpathlineto{\pgfqpoint{1.629482in}{3.672057in}}%
\pgfpathlineto{\pgfqpoint{1.636156in}{5.206505in}}%
\pgfpathlineto{\pgfqpoint{1.636498in}{5.212812in}}%
\pgfpathlineto{\pgfqpoint{1.637012in}{5.202817in}}%
\pgfpathlineto{\pgfqpoint{1.638124in}{5.102983in}}%
\pgfpathlineto{\pgfqpoint{1.640263in}{4.655556in}}%
\pgfpathlineto{\pgfqpoint{1.647022in}{3.112971in}}%
\pgfpathlineto{\pgfqpoint{1.647365in}{3.106852in}}%
\pgfpathlineto{\pgfqpoint{1.647878in}{3.116852in}}%
\pgfpathlineto{\pgfqpoint{1.648990in}{3.215588in}}%
\pgfpathlineto{\pgfqpoint{1.651129in}{3.657754in}}%
\pgfpathlineto{\pgfqpoint{1.657974in}{5.205149in}}%
\pgfpathlineto{\pgfqpoint{1.658317in}{5.210721in}}%
\pgfpathlineto{\pgfqpoint{1.658830in}{5.200179in}}%
\pgfpathlineto{\pgfqpoint{1.659942in}{5.101428in}}%
\pgfpathlineto{\pgfqpoint{1.662081in}{4.662933in}}%
\pgfpathlineto{\pgfqpoint{1.669012in}{3.113695in}}%
\pgfpathlineto{\pgfqpoint{1.669354in}{3.108974in}}%
\pgfpathlineto{\pgfqpoint{1.669782in}{3.117044in}}%
\pgfpathlineto{\pgfqpoint{1.670809in}{3.198474in}}%
\pgfpathlineto{\pgfqpoint{1.672777in}{3.569397in}}%
\pgfpathlineto{\pgfqpoint{1.680392in}{5.208466in}}%
\pgfpathlineto{\pgfqpoint{1.680477in}{5.208453in}}%
\pgfpathlineto{\pgfqpoint{1.680905in}{5.199215in}}%
\pgfpathlineto{\pgfqpoint{1.682017in}{5.105327in}}%
\pgfpathlineto{\pgfqpoint{1.684156in}{4.681215in}}%
\pgfpathlineto{\pgfqpoint{1.691258in}{3.115931in}}%
\pgfpathlineto{\pgfqpoint{1.691600in}{3.111201in}}%
\pgfpathlineto{\pgfqpoint{1.692114in}{3.122180in}}%
\pgfpathlineto{\pgfqpoint{1.693226in}{3.218543in}}%
\pgfpathlineto{\pgfqpoint{1.695365in}{3.642549in}}%
\pgfpathlineto{\pgfqpoint{1.702552in}{5.202910in}}%
\pgfpathlineto{\pgfqpoint{1.702809in}{5.206234in}}%
\pgfpathlineto{\pgfqpoint{1.703322in}{5.196860in}}%
\pgfpathlineto{\pgfqpoint{1.704435in}{5.104904in}}%
\pgfpathlineto{\pgfqpoint{1.706574in}{4.690976in}}%
\pgfpathlineto{\pgfqpoint{1.713847in}{3.117520in}}%
\pgfpathlineto{\pgfqpoint{1.714189in}{3.113520in}}%
\pgfpathlineto{\pgfqpoint{1.714617in}{3.121670in}}%
\pgfpathlineto{\pgfqpoint{1.715643in}{3.199653in}}%
\pgfpathlineto{\pgfqpoint{1.717611in}{3.552534in}}%
\pgfpathlineto{\pgfqpoint{1.725569in}{5.203946in}}%
\pgfpathlineto{\pgfqpoint{1.725654in}{5.203664in}}%
\pgfpathlineto{\pgfqpoint{1.726168in}{5.189917in}}%
\pgfpathlineto{\pgfqpoint{1.727366in}{5.079772in}}%
\pgfpathlineto{\pgfqpoint{1.729676in}{4.609512in}}%
\pgfpathlineto{\pgfqpoint{1.736435in}{3.130891in}}%
\pgfpathlineto{\pgfqpoint{1.737034in}{3.115882in}}%
\pgfpathlineto{\pgfqpoint{1.737547in}{3.125087in}}%
\pgfpathlineto{\pgfqpoint{1.738660in}{3.213434in}}%
\pgfpathlineto{\pgfqpoint{1.740799in}{3.611030in}}%
\pgfpathlineto{\pgfqpoint{1.748414in}{5.199717in}}%
\pgfpathlineto{\pgfqpoint{1.748671in}{5.201495in}}%
\pgfpathlineto{\pgfqpoint{1.749013in}{5.196062in}}%
\pgfpathlineto{\pgfqpoint{1.749954in}{5.135828in}}%
\pgfpathlineto{\pgfqpoint{1.751836in}{4.835168in}}%
\pgfpathlineto{\pgfqpoint{1.760307in}{3.118297in}}%
\pgfpathlineto{\pgfqpoint{1.760820in}{3.128647in}}%
\pgfpathlineto{\pgfqpoint{1.761933in}{3.217214in}}%
\pgfpathlineto{\pgfqpoint{1.764072in}{3.607805in}}%
\pgfpathlineto{\pgfqpoint{1.771858in}{5.197802in}}%
\pgfpathlineto{\pgfqpoint{1.772029in}{5.199072in}}%
\pgfpathlineto{\pgfqpoint{1.772457in}{5.192805in}}%
\pgfpathlineto{\pgfqpoint{1.773484in}{5.123662in}}%
\pgfpathlineto{\pgfqpoint{1.775452in}{4.800358in}}%
\pgfpathlineto{\pgfqpoint{1.783922in}{3.120848in}}%
\pgfpathlineto{\pgfqpoint{1.784179in}{3.123940in}}%
\pgfpathlineto{\pgfqpoint{1.785035in}{3.168372in}}%
\pgfpathlineto{\pgfqpoint{1.786746in}{3.403344in}}%
\pgfpathlineto{\pgfqpoint{1.790339in}{4.295201in}}%
\pgfpathlineto{\pgfqpoint{1.794532in}{5.138139in}}%
\pgfpathlineto{\pgfqpoint{1.795815in}{5.196537in}}%
\pgfpathlineto{\pgfqpoint{1.796158in}{5.192403in}}%
\pgfpathlineto{\pgfqpoint{1.797099in}{5.138537in}}%
\pgfpathlineto{\pgfqpoint{1.798896in}{4.877130in}}%
\pgfpathlineto{\pgfqpoint{1.803003in}{3.844330in}}%
\pgfpathlineto{\pgfqpoint{1.806682in}{3.170407in}}%
\pgfpathlineto{\pgfqpoint{1.807880in}{3.123422in}}%
\pgfpathlineto{\pgfqpoint{1.808222in}{3.128383in}}%
\pgfpathlineto{\pgfqpoint{1.809163in}{3.183745in}}%
\pgfpathlineto{\pgfqpoint{1.811046in}{3.461726in}}%
\pgfpathlineto{\pgfqpoint{1.815495in}{4.576053in}}%
\pgfpathlineto{\pgfqpoint{1.819003in}{5.162019in}}%
\pgfpathlineto{\pgfqpoint{1.819944in}{5.193899in}}%
\pgfpathlineto{\pgfqpoint{1.820457in}{5.185646in}}%
\pgfpathlineto{\pgfqpoint{1.821570in}{5.106764in}}%
\pgfpathlineto{\pgfqpoint{1.823623in}{4.767523in}}%
\pgfpathlineto{\pgfqpoint{1.832180in}{3.126041in}}%
\pgfpathlineto{\pgfqpoint{1.832265in}{3.126235in}}%
\pgfpathlineto{\pgfqpoint{1.832778in}{3.137789in}}%
\pgfpathlineto{\pgfqpoint{1.833976in}{3.232236in}}%
\pgfpathlineto{\pgfqpoint{1.836201in}{3.623517in}}%
\pgfpathlineto{\pgfqpoint{1.844158in}{5.187087in}}%
\pgfpathlineto{\pgfqpoint{1.844501in}{5.191212in}}%
\pgfpathlineto{\pgfqpoint{1.845014in}{5.182806in}}%
\pgfpathlineto{\pgfqpoint{1.846126in}{5.105710in}}%
\pgfpathlineto{\pgfqpoint{1.848180in}{4.775948in}}%
\pgfpathlineto{\pgfqpoint{1.856907in}{3.128807in}}%
\pgfpathlineto{\pgfqpoint{1.857078in}{3.129447in}}%
\pgfpathlineto{\pgfqpoint{1.857763in}{3.151048in}}%
\pgfpathlineto{\pgfqpoint{1.859217in}{3.293983in}}%
\pgfpathlineto{\pgfqpoint{1.861870in}{3.824719in}}%
\pgfpathlineto{\pgfqpoint{1.868372in}{5.150344in}}%
\pgfpathlineto{\pgfqpoint{1.869485in}{5.188417in}}%
\pgfpathlineto{\pgfqpoint{1.869913in}{5.182015in}}%
\pgfpathlineto{\pgfqpoint{1.870939in}{5.119575in}}%
\pgfpathlineto{\pgfqpoint{1.872907in}{4.832074in}}%
\pgfpathlineto{\pgfqpoint{1.877955in}{3.628529in}}%
\pgfpathlineto{\pgfqpoint{1.881207in}{3.156764in}}%
\pgfpathlineto{\pgfqpoint{1.882062in}{3.131660in}}%
\pgfpathlineto{\pgfqpoint{1.882576in}{3.138664in}}%
\pgfpathlineto{\pgfqpoint{1.883688in}{3.209700in}}%
\pgfpathlineto{\pgfqpoint{1.885742in}{3.520193in}}%
\pgfpathlineto{\pgfqpoint{1.894811in}{5.185511in}}%
\pgfpathlineto{\pgfqpoint{1.895068in}{5.184077in}}%
\pgfpathlineto{\pgfqpoint{1.895838in}{5.155512in}}%
\pgfpathlineto{\pgfqpoint{1.897378in}{4.994560in}}%
\pgfpathlineto{\pgfqpoint{1.900287in}{4.402412in}}%
\pgfpathlineto{\pgfqpoint{1.906362in}{3.189039in}}%
\pgfpathlineto{\pgfqpoint{1.907731in}{3.134546in}}%
\pgfpathlineto{\pgfqpoint{1.908073in}{3.138563in}}%
\pgfpathlineto{\pgfqpoint{1.909015in}{3.185756in}}%
\pgfpathlineto{\pgfqpoint{1.910811in}{3.412154in}}%
\pgfpathlineto{\pgfqpoint{1.914491in}{4.235477in}}%
\pgfpathlineto{\pgfqpoint{1.919111in}{5.109913in}}%
\pgfpathlineto{\pgfqpoint{1.920651in}{5.182558in}}%
\pgfpathlineto{\pgfqpoint{1.920993in}{5.179726in}}%
\pgfpathlineto{\pgfqpoint{1.921849in}{5.142470in}}%
\pgfpathlineto{\pgfqpoint{1.923560in}{4.946683in}}%
\pgfpathlineto{\pgfqpoint{1.926812in}{4.258436in}}%
\pgfpathlineto{\pgfqpoint{1.932202in}{3.208931in}}%
\pgfpathlineto{\pgfqpoint{1.933742in}{3.137577in}}%
\pgfpathlineto{\pgfqpoint{1.934084in}{3.140329in}}%
\pgfpathlineto{\pgfqpoint{1.934940in}{3.176802in}}%
\pgfpathlineto{\pgfqpoint{1.936651in}{3.368821in}}%
\pgfpathlineto{\pgfqpoint{1.939903in}{4.046423in}}%
\pgfpathlineto{\pgfqpoint{1.945379in}{5.106300in}}%
\pgfpathlineto{\pgfqpoint{1.947004in}{5.179517in}}%
\pgfpathlineto{\pgfqpoint{1.947261in}{5.177400in}}%
\pgfpathlineto{\pgfqpoint{1.948117in}{5.143393in}}%
\pgfpathlineto{\pgfqpoint{1.949828in}{4.958152in}}%
\pgfpathlineto{\pgfqpoint{1.952994in}{4.314678in}}%
\pgfpathlineto{\pgfqpoint{1.958726in}{3.211458in}}%
\pgfpathlineto{\pgfqpoint{1.960352in}{3.140680in}}%
\pgfpathlineto{\pgfqpoint{1.960609in}{3.142913in}}%
\pgfpathlineto{\pgfqpoint{1.961464in}{3.176760in}}%
\pgfpathlineto{\pgfqpoint{1.963176in}{3.359356in}}%
\pgfpathlineto{\pgfqpoint{1.966341in}{3.993418in}}%
\pgfpathlineto{\pgfqpoint{1.972160in}{5.104376in}}%
\pgfpathlineto{\pgfqpoint{1.973785in}{5.176359in}}%
\pgfpathlineto{\pgfqpoint{1.974042in}{5.174596in}}%
\pgfpathlineto{\pgfqpoint{1.974898in}{5.142831in}}%
\pgfpathlineto{\pgfqpoint{1.976523in}{4.978442in}}%
\pgfpathlineto{\pgfqpoint{1.979518in}{4.404280in}}%
\pgfpathlineto{\pgfqpoint{1.985850in}{3.206357in}}%
\pgfpathlineto{\pgfqpoint{1.987390in}{3.143879in}}%
\pgfpathlineto{\pgfqpoint{1.987646in}{3.145756in}}%
\pgfpathlineto{\pgfqpoint{1.988502in}{3.177360in}}%
\pgfpathlineto{\pgfqpoint{1.990213in}{3.351226in}}%
\pgfpathlineto{\pgfqpoint{1.993294in}{3.943261in}}%
\pgfpathlineto{\pgfqpoint{1.999454in}{5.101858in}}%
\pgfpathlineto{\pgfqpoint{2.001080in}{5.173065in}}%
\pgfpathlineto{\pgfqpoint{2.001336in}{5.171725in}}%
\pgfpathlineto{\pgfqpoint{2.002107in}{5.147050in}}%
\pgfpathlineto{\pgfqpoint{2.003647in}{5.008750in}}%
\pgfpathlineto{\pgfqpoint{2.006470in}{4.509254in}}%
\pgfpathlineto{\pgfqpoint{2.013657in}{3.191421in}}%
\pgfpathlineto{\pgfqpoint{2.014941in}{3.147231in}}%
\pgfpathlineto{\pgfqpoint{2.015369in}{3.151239in}}%
\pgfpathlineto{\pgfqpoint{2.016395in}{3.198788in}}%
\pgfpathlineto{\pgfqpoint{2.018278in}{3.414810in}}%
\pgfpathlineto{\pgfqpoint{2.021957in}{4.158790in}}%
\pgfpathlineto{\pgfqpoint{2.027091in}{5.081361in}}%
\pgfpathlineto{\pgfqpoint{2.028973in}{5.169701in}}%
\pgfpathlineto{\pgfqpoint{2.029144in}{5.168996in}}%
\pgfpathlineto{\pgfqpoint{2.029829in}{5.151504in}}%
\pgfpathlineto{\pgfqpoint{2.031283in}{5.038991in}}%
\pgfpathlineto{\pgfqpoint{2.033850in}{4.630601in}}%
\pgfpathlineto{\pgfqpoint{2.042235in}{3.170098in}}%
\pgfpathlineto{\pgfqpoint{2.043091in}{3.150679in}}%
\pgfpathlineto{\pgfqpoint{2.043604in}{3.156285in}}%
\pgfpathlineto{\pgfqpoint{2.044717in}{3.212238in}}%
\pgfpathlineto{\pgfqpoint{2.046770in}{3.459393in}}%
\pgfpathlineto{\pgfqpoint{2.050963in}{4.314766in}}%
\pgfpathlineto{\pgfqpoint{2.055583in}{5.085795in}}%
\pgfpathlineto{\pgfqpoint{2.057380in}{5.166154in}}%
\pgfpathlineto{\pgfqpoint{2.057551in}{5.165800in}}%
\pgfpathlineto{\pgfqpoint{2.058150in}{5.153498in}}%
\pgfpathlineto{\pgfqpoint{2.059519in}{5.062475in}}%
\pgfpathlineto{\pgfqpoint{2.061915in}{4.720052in}}%
\pgfpathlineto{\pgfqpoint{2.071498in}{3.157502in}}%
\pgfpathlineto{\pgfqpoint{2.071840in}{3.154239in}}%
\pgfpathlineto{\pgfqpoint{2.072353in}{3.159660in}}%
\pgfpathlineto{\pgfqpoint{2.073466in}{3.213262in}}%
\pgfpathlineto{\pgfqpoint{2.075519in}{3.450164in}}%
\pgfpathlineto{\pgfqpoint{2.079541in}{4.241918in}}%
\pgfpathlineto{\pgfqpoint{2.084503in}{5.074243in}}%
\pgfpathlineto{\pgfqpoint{2.086471in}{5.162567in}}%
\pgfpathlineto{\pgfqpoint{2.086557in}{5.162426in}}%
\pgfpathlineto{\pgfqpoint{2.087070in}{5.154532in}}%
\pgfpathlineto{\pgfqpoint{2.088268in}{5.090011in}}%
\pgfpathlineto{\pgfqpoint{2.090407in}{4.829592in}}%
\pgfpathlineto{\pgfqpoint{2.094942in}{3.927439in}}%
\pgfpathlineto{\pgfqpoint{2.099391in}{3.234093in}}%
\pgfpathlineto{\pgfqpoint{2.101273in}{3.157928in}}%
\pgfpathlineto{\pgfqpoint{2.101445in}{3.158799in}}%
\pgfpathlineto{\pgfqpoint{2.102215in}{3.178875in}}%
\pgfpathlineto{\pgfqpoint{2.103755in}{3.295466in}}%
\pgfpathlineto{\pgfqpoint{2.106493in}{3.711949in}}%
\pgfpathlineto{\pgfqpoint{2.115049in}{5.130994in}}%
\pgfpathlineto{\pgfqpoint{2.116161in}{5.158800in}}%
\pgfpathlineto{\pgfqpoint{2.116589in}{5.155129in}}%
\pgfpathlineto{\pgfqpoint{2.117616in}{5.114040in}}%
\pgfpathlineto{\pgfqpoint{2.119498in}{4.928003in}}%
\pgfpathlineto{\pgfqpoint{2.122921in}{4.323049in}}%
\pgfpathlineto{\pgfqpoint{2.129252in}{3.248595in}}%
\pgfpathlineto{\pgfqpoint{2.131306in}{3.161752in}}%
\pgfpathlineto{\pgfqpoint{2.131391in}{3.161990in}}%
\pgfpathlineto{\pgfqpoint{2.131990in}{3.172375in}}%
\pgfpathlineto{\pgfqpoint{2.133274in}{3.244914in}}%
\pgfpathlineto{\pgfqpoint{2.135584in}{3.529640in}}%
\pgfpathlineto{\pgfqpoint{2.141060in}{4.586165in}}%
\pgfpathlineto{\pgfqpoint{2.144910in}{5.098472in}}%
\pgfpathlineto{\pgfqpoint{2.146536in}{5.154887in}}%
\pgfpathlineto{\pgfqpoint{2.146793in}{5.153820in}}%
\pgfpathlineto{\pgfqpoint{2.147563in}{5.134265in}}%
\pgfpathlineto{\pgfqpoint{2.149103in}{5.024158in}}%
\pgfpathlineto{\pgfqpoint{2.151841in}{4.632077in}}%
\pgfpathlineto{\pgfqpoint{2.160996in}{3.187197in}}%
\pgfpathlineto{\pgfqpoint{2.162023in}{3.165718in}}%
\pgfpathlineto{\pgfqpoint{2.162451in}{3.169350in}}%
\pgfpathlineto{\pgfqpoint{2.163477in}{3.207985in}}%
\pgfpathlineto{\pgfqpoint{2.165360in}{3.381812in}}%
\pgfpathlineto{\pgfqpoint{2.168782in}{3.951659in}}%
\pgfpathlineto{\pgfqpoint{2.175542in}{5.062039in}}%
\pgfpathlineto{\pgfqpoint{2.177681in}{5.150851in}}%
\pgfpathlineto{\pgfqpoint{2.178023in}{5.148493in}}%
\pgfpathlineto{\pgfqpoint{2.178964in}{5.118380in}}%
\pgfpathlineto{\pgfqpoint{2.180761in}{4.969495in}}%
\pgfpathlineto{\pgfqpoint{2.183927in}{4.478098in}}%
\pgfpathlineto{\pgfqpoint{2.191627in}{3.238048in}}%
\pgfpathlineto{\pgfqpoint{2.193510in}{3.169851in}}%
\pgfpathlineto{\pgfqpoint{2.193681in}{3.170340in}}%
\pgfpathlineto{\pgfqpoint{2.194365in}{3.183505in}}%
\pgfpathlineto{\pgfqpoint{2.195820in}{3.269461in}}%
\pgfpathlineto{\pgfqpoint{2.198301in}{3.575735in}}%
\pgfpathlineto{\pgfqpoint{2.209339in}{5.145685in}}%
\pgfpathlineto{\pgfqpoint{2.209596in}{5.146621in}}%
\pgfpathlineto{\pgfqpoint{2.210023in}{5.142712in}}%
\pgfpathlineto{\pgfqpoint{2.211050in}{5.105744in}}%
\pgfpathlineto{\pgfqpoint{2.212932in}{4.942731in}}%
\pgfpathlineto{\pgfqpoint{2.216269in}{4.422802in}}%
\pgfpathlineto{\pgfqpoint{2.223713in}{3.253486in}}%
\pgfpathlineto{\pgfqpoint{2.225852in}{3.174169in}}%
\pgfpathlineto{\pgfqpoint{2.226280in}{3.178070in}}%
\pgfpathlineto{\pgfqpoint{2.227392in}{3.218951in}}%
\pgfpathlineto{\pgfqpoint{2.229360in}{3.393199in}}%
\pgfpathlineto{\pgfqpoint{2.232868in}{3.943990in}}%
\pgfpathlineto{\pgfqpoint{2.239970in}{5.049400in}}%
\pgfpathlineto{\pgfqpoint{2.242195in}{5.142131in}}%
\pgfpathlineto{\pgfqpoint{2.242280in}{5.142257in}}%
\pgfpathlineto{\pgfqpoint{2.242537in}{5.141085in}}%
\pgfpathlineto{\pgfqpoint{2.243393in}{5.120471in}}%
\pgfpathlineto{\pgfqpoint{2.245018in}{5.013034in}}%
\pgfpathlineto{\pgfqpoint{2.247842in}{4.646891in}}%
\pgfpathlineto{\pgfqpoint{2.257853in}{3.199768in}}%
\pgfpathlineto{\pgfqpoint{2.258965in}{3.178615in}}%
\pgfpathlineto{\pgfqpoint{2.259393in}{3.181755in}}%
\pgfpathlineto{\pgfqpoint{2.260420in}{3.214642in}}%
\pgfpathlineto{\pgfqpoint{2.262302in}{3.362948in}}%
\pgfpathlineto{\pgfqpoint{2.265553in}{3.830085in}}%
\pgfpathlineto{\pgfqpoint{2.273853in}{5.070013in}}%
\pgfpathlineto{\pgfqpoint{2.275906in}{5.137681in}}%
\pgfpathlineto{\pgfqpoint{2.275992in}{5.137482in}}%
\pgfpathlineto{\pgfqpoint{2.276591in}{5.129285in}}%
\pgfpathlineto{\pgfqpoint{2.277874in}{5.072315in}}%
\pgfpathlineto{\pgfqpoint{2.280099in}{4.857047in}}%
\pgfpathlineto{\pgfqpoint{2.284292in}{4.183985in}}%
\pgfpathlineto{\pgfqpoint{2.290281in}{3.303175in}}%
\pgfpathlineto{\pgfqpoint{2.292762in}{3.184446in}}%
\pgfpathlineto{\pgfqpoint{2.293019in}{3.183278in}}%
\pgfpathlineto{\pgfqpoint{2.293447in}{3.186056in}}%
\pgfpathlineto{\pgfqpoint{2.294473in}{3.216601in}}%
\pgfpathlineto{\pgfqpoint{2.296356in}{3.355778in}}%
\pgfpathlineto{\pgfqpoint{2.299607in}{3.798606in}}%
\pgfpathlineto{\pgfqpoint{2.308420in}{5.070299in}}%
\pgfpathlineto{\pgfqpoint{2.310388in}{5.132906in}}%
\pgfpathlineto{\pgfqpoint{2.310559in}{5.132673in}}%
\pgfpathlineto{\pgfqpoint{2.311158in}{5.124659in}}%
\pgfpathlineto{\pgfqpoint{2.312442in}{5.070434in}}%
\pgfpathlineto{\pgfqpoint{2.314752in}{4.856135in}}%
\pgfpathlineto{\pgfqpoint{2.318944in}{4.208589in}}%
\pgfpathlineto{\pgfqpoint{2.325276in}{3.305512in}}%
\pgfpathlineto{\pgfqpoint{2.327757in}{3.189760in}}%
\pgfpathlineto{\pgfqpoint{2.328100in}{3.188142in}}%
\pgfpathlineto{\pgfqpoint{2.328527in}{3.191107in}}%
\pgfpathlineto{\pgfqpoint{2.329554in}{3.220632in}}%
\pgfpathlineto{\pgfqpoint{2.331436in}{3.352948in}}%
\pgfpathlineto{\pgfqpoint{2.334602in}{3.760593in}}%
\pgfpathlineto{\pgfqpoint{2.344100in}{5.075152in}}%
\pgfpathlineto{\pgfqpoint{2.345982in}{5.127956in}}%
\pgfpathlineto{\pgfqpoint{2.346153in}{5.127640in}}%
\pgfpathlineto{\pgfqpoint{2.346838in}{5.117815in}}%
\pgfpathlineto{\pgfqpoint{2.348292in}{5.052409in}}%
\pgfpathlineto{\pgfqpoint{2.350774in}{4.814569in}}%
\pgfpathlineto{\pgfqpoint{2.355479in}{4.091423in}}%
\pgfpathlineto{\pgfqpoint{2.361212in}{3.317482in}}%
\pgfpathlineto{\pgfqpoint{2.363779in}{3.195747in}}%
\pgfpathlineto{\pgfqpoint{2.364207in}{3.193220in}}%
\pgfpathlineto{\pgfqpoint{2.364720in}{3.197033in}}%
\pgfpathlineto{\pgfqpoint{2.365833in}{3.230639in}}%
\pgfpathlineto{\pgfqpoint{2.367800in}{3.370521in}}%
\pgfpathlineto{\pgfqpoint{2.371137in}{3.795107in}}%
\pgfpathlineto{\pgfqpoint{2.380464in}{5.055794in}}%
\pgfpathlineto{\pgfqpoint{2.382688in}{5.122766in}}%
\pgfpathlineto{\pgfqpoint{2.383031in}{5.121138in}}%
\pgfpathlineto{\pgfqpoint{2.383972in}{5.100232in}}%
\pgfpathlineto{\pgfqpoint{2.385769in}{4.995881in}}%
\pgfpathlineto{\pgfqpoint{2.388763in}{4.662143in}}%
\pgfpathlineto{\pgfqpoint{2.400314in}{3.216165in}}%
\pgfpathlineto{\pgfqpoint{2.401512in}{3.198562in}}%
\pgfpathlineto{\pgfqpoint{2.401854in}{3.200489in}}%
\pgfpathlineto{\pgfqpoint{2.402796in}{3.221638in}}%
\pgfpathlineto{\pgfqpoint{2.404592in}{3.324172in}}%
\pgfpathlineto{\pgfqpoint{2.407587in}{3.649661in}}%
\pgfpathlineto{\pgfqpoint{2.419395in}{5.099826in}}%
\pgfpathlineto{\pgfqpoint{2.420593in}{5.117313in}}%
\pgfpathlineto{\pgfqpoint{2.421020in}{5.114707in}}%
\pgfpathlineto{\pgfqpoint{2.422047in}{5.089583in}}%
\pgfpathlineto{\pgfqpoint{2.423929in}{4.977255in}}%
\pgfpathlineto{\pgfqpoint{2.427010in}{4.637976in}}%
\pgfpathlineto{\pgfqpoint{2.438475in}{3.232243in}}%
\pgfpathlineto{\pgfqpoint{2.440015in}{3.204122in}}%
\pgfpathlineto{\pgfqpoint{2.440272in}{3.205082in}}%
\pgfpathlineto{\pgfqpoint{2.441128in}{3.219910in}}%
\pgfpathlineto{\pgfqpoint{2.442753in}{3.296027in}}%
\pgfpathlineto{\pgfqpoint{2.445491in}{3.550346in}}%
\pgfpathlineto{\pgfqpoint{2.450882in}{4.316953in}}%
\pgfpathlineto{\pgfqpoint{2.456443in}{4.984375in}}%
\pgfpathlineto{\pgfqpoint{2.459181in}{5.107696in}}%
\pgfpathlineto{\pgfqpoint{2.459780in}{5.111601in}}%
\pgfpathlineto{\pgfqpoint{2.460294in}{5.108205in}}%
\pgfpathlineto{\pgfqpoint{2.461406in}{5.079697in}}%
\pgfpathlineto{\pgfqpoint{2.463374in}{4.961657in}}%
\pgfpathlineto{\pgfqpoint{2.466625in}{4.609560in}}%
\pgfpathlineto{\pgfqpoint{2.478005in}{3.249686in}}%
\pgfpathlineto{\pgfqpoint{2.479887in}{3.209972in}}%
\pgfpathlineto{\pgfqpoint{2.480058in}{3.210335in}}%
\pgfpathlineto{\pgfqpoint{2.480743in}{3.218421in}}%
\pgfpathlineto{\pgfqpoint{2.482197in}{3.270220in}}%
\pgfpathlineto{\pgfqpoint{2.484593in}{3.450295in}}%
\pgfpathlineto{\pgfqpoint{2.488700in}{3.959195in}}%
\pgfpathlineto{\pgfqpoint{2.497000in}{4.983476in}}%
\pgfpathlineto{\pgfqpoint{2.499738in}{5.101143in}}%
\pgfpathlineto{\pgfqpoint{2.500422in}{5.105578in}}%
\pgfpathlineto{\pgfqpoint{2.500850in}{5.103162in}}%
\pgfpathlineto{\pgfqpoint{2.501962in}{5.078363in}}%
\pgfpathlineto{\pgfqpoint{2.503930in}{4.971670in}}%
\pgfpathlineto{\pgfqpoint{2.507096in}{4.657267in}}%
\pgfpathlineto{\pgfqpoint{2.519845in}{3.237270in}}%
\pgfpathlineto{\pgfqpoint{2.521300in}{3.216148in}}%
\pgfpathlineto{\pgfqpoint{2.521556in}{3.217006in}}%
\pgfpathlineto{\pgfqpoint{2.522412in}{3.229775in}}%
\pgfpathlineto{\pgfqpoint{2.524038in}{3.295017in}}%
\pgfpathlineto{\pgfqpoint{2.526690in}{3.505746in}}%
\pgfpathlineto{\pgfqpoint{2.531396in}{4.092899in}}%
\pgfpathlineto{\pgfqpoint{2.538669in}{4.950123in}}%
\pgfpathlineto{\pgfqpoint{2.541663in}{5.091000in}}%
\pgfpathlineto{\pgfqpoint{2.542605in}{5.099258in}}%
\pgfpathlineto{\pgfqpoint{2.543032in}{5.097148in}}%
\pgfpathlineto{\pgfqpoint{2.544059in}{5.077238in}}%
\pgfpathlineto{\pgfqpoint{2.545942in}{4.988199in}}%
\pgfpathlineto{\pgfqpoint{2.549022in}{4.714871in}}%
\pgfpathlineto{\pgfqpoint{2.555525in}{3.875282in}}%
\pgfpathlineto{\pgfqpoint{2.560744in}{3.342405in}}%
\pgfpathlineto{\pgfqpoint{2.563567in}{3.227946in}}%
\pgfpathlineto{\pgfqpoint{2.564337in}{3.222650in}}%
\pgfpathlineto{\pgfqpoint{2.564851in}{3.225428in}}%
\pgfpathlineto{\pgfqpoint{2.565963in}{3.248616in}}%
\pgfpathlineto{\pgfqpoint{2.567931in}{3.344933in}}%
\pgfpathlineto{\pgfqpoint{2.571097in}{3.627509in}}%
\pgfpathlineto{\pgfqpoint{2.578883in}{4.605520in}}%
\pgfpathlineto{\pgfqpoint{2.583503in}{5.010861in}}%
\pgfpathlineto{\pgfqpoint{2.586070in}{5.090814in}}%
\pgfpathlineto{\pgfqpoint{2.586498in}{5.092577in}}%
\pgfpathlineto{\pgfqpoint{2.587012in}{5.090274in}}%
\pgfpathlineto{\pgfqpoint{2.588124in}{5.068874in}}%
\pgfpathlineto{\pgfqpoint{2.590092in}{4.978035in}}%
\pgfpathlineto{\pgfqpoint{2.593258in}{4.708623in}}%
\pgfpathlineto{\pgfqpoint{2.599846in}{3.898823in}}%
\pgfpathlineto{\pgfqpoint{2.605322in}{3.357672in}}%
\pgfpathlineto{\pgfqpoint{2.608317in}{3.236210in}}%
\pgfpathlineto{\pgfqpoint{2.609172in}{3.229531in}}%
\pgfpathlineto{\pgfqpoint{2.609686in}{3.231652in}}%
\pgfpathlineto{\pgfqpoint{2.610798in}{3.251898in}}%
\pgfpathlineto{\pgfqpoint{2.612766in}{3.338330in}}%
\pgfpathlineto{\pgfqpoint{2.615846in}{3.587345in}}%
\pgfpathlineto{\pgfqpoint{2.621835in}{4.297310in}}%
\pgfpathlineto{\pgfqpoint{2.628081in}{4.935093in}}%
\pgfpathlineto{\pgfqpoint{2.631247in}{5.074711in}}%
\pgfpathlineto{\pgfqpoint{2.632360in}{5.085483in}}%
\pgfpathlineto{\pgfqpoint{2.632787in}{5.084159in}}%
\pgfpathlineto{\pgfqpoint{2.633814in}{5.068655in}}%
\pgfpathlineto{\pgfqpoint{2.635697in}{4.996342in}}%
\pgfpathlineto{\pgfqpoint{2.638691in}{4.776657in}}%
\pgfpathlineto{\pgfqpoint{2.643911in}{4.194227in}}%
\pgfpathlineto{\pgfqpoint{2.651440in}{3.405536in}}%
\pgfpathlineto{\pgfqpoint{2.654777in}{3.251168in}}%
\pgfpathlineto{\pgfqpoint{2.656146in}{3.236837in}}%
\pgfpathlineto{\pgfqpoint{2.656488in}{3.237870in}}%
\pgfpathlineto{\pgfqpoint{2.657429in}{3.250191in}}%
\pgfpathlineto{\pgfqpoint{2.659226in}{3.311385in}}%
\pgfpathlineto{\pgfqpoint{2.662050in}{3.498469in}}%
\pgfpathlineto{\pgfqpoint{2.666756in}{3.986063in}}%
\pgfpathlineto{\pgfqpoint{2.675825in}{4.922296in}}%
\pgfpathlineto{\pgfqpoint{2.679077in}{5.064025in}}%
\pgfpathlineto{\pgfqpoint{2.680446in}{5.077952in}}%
\pgfpathlineto{\pgfqpoint{2.680788in}{5.077057in}}%
\pgfpathlineto{\pgfqpoint{2.681729in}{5.065602in}}%
\pgfpathlineto{\pgfqpoint{2.683440in}{5.011731in}}%
\pgfpathlineto{\pgfqpoint{2.686178in}{4.844038in}}%
\pgfpathlineto{\pgfqpoint{2.690628in}{4.411831in}}%
\pgfpathlineto{\pgfqpoint{2.700981in}{3.378468in}}%
\pgfpathlineto{\pgfqpoint{2.704147in}{3.255258in}}%
\pgfpathlineto{\pgfqpoint{2.705344in}{3.244622in}}%
\pgfpathlineto{\pgfqpoint{2.705772in}{3.245743in}}%
\pgfpathlineto{\pgfqpoint{2.706799in}{3.258948in}}%
\pgfpathlineto{\pgfqpoint{2.708596in}{3.316891in}}%
\pgfpathlineto{\pgfqpoint{2.711505in}{3.496412in}}%
\pgfpathlineto{\pgfqpoint{2.716211in}{3.949925in}}%
\pgfpathlineto{\pgfqpoint{2.726136in}{4.917906in}}%
\pgfpathlineto{\pgfqpoint{2.729473in}{5.055133in}}%
\pgfpathlineto{\pgfqpoint{2.730928in}{5.069901in}}%
\pgfpathlineto{\pgfqpoint{2.731270in}{5.069271in}}%
\pgfpathlineto{\pgfqpoint{2.732125in}{5.060870in}}%
\pgfpathlineto{\pgfqpoint{2.733751in}{5.018533in}}%
\pgfpathlineto{\pgfqpoint{2.736404in}{4.880169in}}%
\pgfpathlineto{\pgfqpoint{2.740511in}{4.530857in}}%
\pgfpathlineto{\pgfqpoint{2.753516in}{3.339707in}}%
\pgfpathlineto{\pgfqpoint{2.756425in}{3.257234in}}%
\pgfpathlineto{\pgfqpoint{2.757281in}{3.252945in}}%
\pgfpathlineto{\pgfqpoint{2.757709in}{3.254255in}}%
\pgfpathlineto{\pgfqpoint{2.758735in}{3.266744in}}%
\pgfpathlineto{\pgfqpoint{2.760618in}{3.323064in}}%
\pgfpathlineto{\pgfqpoint{2.763527in}{3.488040in}}%
\pgfpathlineto{\pgfqpoint{2.768233in}{3.904350in}}%
\pgfpathlineto{\pgfqpoint{2.779356in}{4.920023in}}%
\pgfpathlineto{\pgfqpoint{2.782693in}{5.046108in}}%
\pgfpathlineto{\pgfqpoint{2.784319in}{5.061283in}}%
\pgfpathlineto{\pgfqpoint{2.784575in}{5.060825in}}%
\pgfpathlineto{\pgfqpoint{2.785431in}{5.053694in}}%
\pgfpathlineto{\pgfqpoint{2.787057in}{5.016759in}}%
\pgfpathlineto{\pgfqpoint{2.789623in}{4.899663in}}%
\pgfpathlineto{\pgfqpoint{2.793645in}{4.596894in}}%
\pgfpathlineto{\pgfqpoint{2.809046in}{3.314995in}}%
\pgfpathlineto{\pgfqpoint{2.811699in}{3.263029in}}%
\pgfpathlineto{\pgfqpoint{2.812212in}{3.261904in}}%
\pgfpathlineto{\pgfqpoint{2.812640in}{3.263194in}}%
\pgfpathlineto{\pgfqpoint{2.813752in}{3.275970in}}%
\pgfpathlineto{\pgfqpoint{2.815720in}{3.331093in}}%
\pgfpathlineto{\pgfqpoint{2.818715in}{3.488099in}}%
\pgfpathlineto{\pgfqpoint{2.823421in}{3.870199in}}%
\pgfpathlineto{\pgfqpoint{2.835742in}{4.917752in}}%
\pgfpathlineto{\pgfqpoint{2.839164in}{5.036643in}}%
\pgfpathlineto{\pgfqpoint{2.840875in}{5.051992in}}%
\pgfpathlineto{\pgfqpoint{2.841132in}{5.051686in}}%
\pgfpathlineto{\pgfqpoint{2.841988in}{5.045764in}}%
\pgfpathlineto{\pgfqpoint{2.843613in}{5.014024in}}%
\pgfpathlineto{\pgfqpoint{2.846180in}{4.912054in}}%
\pgfpathlineto{\pgfqpoint{2.850030in}{4.657992in}}%
\pgfpathlineto{\pgfqpoint{2.857731in}{3.962778in}}%
\pgfpathlineto{\pgfqpoint{2.864491in}{3.444932in}}%
\pgfpathlineto{\pgfqpoint{2.868255in}{3.297378in}}%
\pgfpathlineto{\pgfqpoint{2.870480in}{3.271598in}}%
\pgfpathlineto{\pgfqpoint{2.870822in}{3.271846in}}%
\pgfpathlineto{\pgfqpoint{2.870993in}{3.272392in}}%
\pgfpathlineto{\pgfqpoint{2.872020in}{3.281545in}}%
\pgfpathlineto{\pgfqpoint{2.873902in}{3.324002in}}%
\pgfpathlineto{\pgfqpoint{2.876812in}{3.450620in}}%
\pgfpathlineto{\pgfqpoint{2.881175in}{3.753104in}}%
\pgfpathlineto{\pgfqpoint{2.897261in}{4.969894in}}%
\pgfpathlineto{\pgfqpoint{2.900256in}{5.037056in}}%
\pgfpathlineto{\pgfqpoint{2.901282in}{5.041913in}}%
\pgfpathlineto{\pgfqpoint{2.901710in}{5.041166in}}%
\pgfpathlineto{\pgfqpoint{2.902737in}{5.032753in}}%
\pgfpathlineto{\pgfqpoint{2.904534in}{4.995947in}}%
\pgfpathlineto{\pgfqpoint{2.907357in}{4.884704in}}%
\pgfpathlineto{\pgfqpoint{2.911550in}{4.619133in}}%
\pgfpathlineto{\pgfqpoint{2.921304in}{3.806026in}}%
\pgfpathlineto{\pgfqpoint{2.927293in}{3.420700in}}%
\pgfpathlineto{\pgfqpoint{2.930973in}{3.301704in}}%
\pgfpathlineto{\pgfqpoint{2.933197in}{3.282085in}}%
\pgfpathlineto{\pgfqpoint{2.933796in}{3.283742in}}%
\pgfpathlineto{\pgfqpoint{2.935080in}{3.297143in}}%
\pgfpathlineto{\pgfqpoint{2.937219in}{3.348546in}}%
\pgfpathlineto{\pgfqpoint{2.940384in}{3.486038in}}%
\pgfpathlineto{\pgfqpoint{2.945176in}{3.803889in}}%
\pgfpathlineto{\pgfqpoint{2.960919in}{4.924536in}}%
\pgfpathlineto{\pgfqpoint{2.964428in}{5.018140in}}%
\pgfpathlineto{\pgfqpoint{2.966310in}{5.030893in}}%
\pgfpathlineto{\pgfqpoint{2.966395in}{5.030837in}}%
\pgfpathlineto{\pgfqpoint{2.967080in}{5.028412in}}%
\pgfpathlineto{\pgfqpoint{2.968449in}{5.013063in}}%
\pgfpathlineto{\pgfqpoint{2.970759in}{4.956360in}}%
\pgfpathlineto{\pgfqpoint{2.974182in}{4.807474in}}%
\pgfpathlineto{\pgfqpoint{2.979315in}{4.472692in}}%
\pgfpathlineto{\pgfqpoint{2.994460in}{3.427639in}}%
\pgfpathlineto{\pgfqpoint{2.998310in}{3.314206in}}%
\pgfpathlineto{\pgfqpoint{3.000620in}{3.293698in}}%
\pgfpathlineto{\pgfqpoint{3.000963in}{3.293782in}}%
\pgfpathlineto{\pgfqpoint{3.001219in}{3.294374in}}%
\pgfpathlineto{\pgfqpoint{3.002332in}{3.302151in}}%
\pgfpathlineto{\pgfqpoint{3.004214in}{3.334283in}}%
\pgfpathlineto{\pgfqpoint{3.007123in}{3.428325in}}%
\pgfpathlineto{\pgfqpoint{3.011401in}{3.650397in}}%
\pgfpathlineto{\pgfqpoint{3.018760in}{4.171903in}}%
\pgfpathlineto{\pgfqpoint{3.028086in}{4.790593in}}%
\pgfpathlineto{\pgfqpoint{3.032706in}{4.967102in}}%
\pgfpathlineto{\pgfqpoint{3.035701in}{5.015091in}}%
\pgfpathlineto{\pgfqpoint{3.036813in}{5.018740in}}%
\pgfpathlineto{\pgfqpoint{3.037241in}{5.018086in}}%
\pgfpathlineto{\pgfqpoint{3.038353in}{5.011062in}}%
\pgfpathlineto{\pgfqpoint{3.040236in}{4.981946in}}%
\pgfpathlineto{\pgfqpoint{3.043145in}{4.896476in}}%
\pgfpathlineto{\pgfqpoint{3.047338in}{4.698198in}}%
\pgfpathlineto{\pgfqpoint{3.054183in}{4.250169in}}%
\pgfpathlineto{\pgfqpoint{3.065306in}{3.537840in}}%
\pgfpathlineto{\pgfqpoint{3.070097in}{3.361562in}}%
\pgfpathlineto{\pgfqpoint{3.073263in}{3.311007in}}%
\pgfpathlineto{\pgfqpoint{3.074546in}{3.306463in}}%
\pgfpathlineto{\pgfqpoint{3.074889in}{3.306817in}}%
\pgfpathlineto{\pgfqpoint{3.075915in}{3.311815in}}%
\pgfpathlineto{\pgfqpoint{3.077712in}{3.334586in}}%
\pgfpathlineto{\pgfqpoint{3.080450in}{3.402213in}}%
\pgfpathlineto{\pgfqpoint{3.084386in}{3.561702in}}%
\pgfpathlineto{\pgfqpoint{3.090375in}{3.907980in}}%
\pgfpathlineto{\pgfqpoint{3.105606in}{4.822616in}}%
\pgfpathlineto{\pgfqpoint{3.110140in}{4.963108in}}%
\pgfpathlineto{\pgfqpoint{3.113135in}{5.002106in}}%
\pgfpathlineto{\pgfqpoint{3.114247in}{5.005159in}}%
\pgfpathlineto{\pgfqpoint{3.114675in}{5.004677in}}%
\pgfpathlineto{\pgfqpoint{3.115787in}{4.999143in}}%
\pgfpathlineto{\pgfqpoint{3.117670in}{4.975884in}}%
\pgfpathlineto{\pgfqpoint{3.120579in}{4.907050in}}%
\pgfpathlineto{\pgfqpoint{3.124686in}{4.749244in}}%
\pgfpathlineto{\pgfqpoint{3.130846in}{4.414687in}}%
\pgfpathlineto{\pgfqpoint{3.146932in}{3.505931in}}%
\pgfpathlineto{\pgfqpoint{3.151638in}{3.366435in}}%
\pgfpathlineto{\pgfqpoint{3.154804in}{3.325061in}}%
\pgfpathlineto{\pgfqpoint{3.156173in}{3.320907in}}%
\pgfpathlineto{\pgfqpoint{3.156515in}{3.321170in}}%
\pgfpathlineto{\pgfqpoint{3.157542in}{3.325071in}}%
\pgfpathlineto{\pgfqpoint{3.159339in}{3.342994in}}%
\pgfpathlineto{\pgfqpoint{3.162077in}{3.396524in}}%
\pgfpathlineto{\pgfqpoint{3.166013in}{3.524060in}}%
\pgfpathlineto{\pgfqpoint{3.171660in}{3.789587in}}%
\pgfpathlineto{\pgfqpoint{3.193136in}{4.880313in}}%
\pgfpathlineto{\pgfqpoint{3.197329in}{4.967861in}}%
\pgfpathlineto{\pgfqpoint{3.200067in}{4.988988in}}%
\pgfpathlineto{\pgfqpoint{3.200837in}{4.989680in}}%
\pgfpathlineto{\pgfqpoint{3.201264in}{4.989070in}}%
\pgfpathlineto{\pgfqpoint{3.202462in}{4.983602in}}%
\pgfpathlineto{\pgfqpoint{3.204516in}{4.961533in}}%
\pgfpathlineto{\pgfqpoint{3.207510in}{4.901781in}}%
\pgfpathlineto{\pgfqpoint{3.211703in}{4.768577in}}%
\pgfpathlineto{\pgfqpoint{3.217778in}{4.496073in}}%
\pgfpathlineto{\pgfqpoint{3.239083in}{3.475192in}}%
\pgfpathlineto{\pgfqpoint{3.243703in}{3.371605in}}%
\pgfpathlineto{\pgfqpoint{3.246869in}{3.340539in}}%
\pgfpathlineto{\pgfqpoint{3.248324in}{3.337511in}}%
\pgfpathlineto{\pgfqpoint{3.248580in}{3.337713in}}%
\pgfpathlineto{\pgfqpoint{3.249607in}{3.340715in}}%
\pgfpathlineto{\pgfqpoint{3.251404in}{3.354340in}}%
\pgfpathlineto{\pgfqpoint{3.254142in}{3.394986in}}%
\pgfpathlineto{\pgfqpoint{3.257992in}{3.489901in}}%
\pgfpathlineto{\pgfqpoint{3.263383in}{3.685189in}}%
\pgfpathlineto{\pgfqpoint{3.272110in}{4.096921in}}%
\pgfpathlineto{\pgfqpoint{3.285116in}{4.694673in}}%
\pgfpathlineto{\pgfqpoint{3.291191in}{4.877832in}}%
\pgfpathlineto{\pgfqpoint{3.295469in}{4.950318in}}%
\pgfpathlineto{\pgfqpoint{3.298378in}{4.970421in}}%
\pgfpathlineto{\pgfqpoint{3.299490in}{4.971751in}}%
\pgfpathlineto{\pgfqpoint{3.299832in}{4.971454in}}%
\pgfpathlineto{\pgfqpoint{3.301030in}{4.967814in}}%
\pgfpathlineto{\pgfqpoint{3.302998in}{4.953151in}}%
\pgfpathlineto{\pgfqpoint{3.305907in}{4.912331in}}%
\pgfpathlineto{\pgfqpoint{3.309929in}{4.821020in}}%
\pgfpathlineto{\pgfqpoint{3.315490in}{4.638419in}}%
\pgfpathlineto{\pgfqpoint{3.324218in}{4.267260in}}%
\pgfpathlineto{\pgfqpoint{3.339105in}{3.639333in}}%
\pgfpathlineto{\pgfqpoint{3.345437in}{3.461233in}}%
\pgfpathlineto{\pgfqpoint{3.350057in}{3.384467in}}%
\pgfpathlineto{\pgfqpoint{3.353309in}{3.359986in}}%
\pgfpathlineto{\pgfqpoint{3.355020in}{3.357108in}}%
\pgfpathlineto{\pgfqpoint{3.355191in}{3.357199in}}%
\pgfpathlineto{\pgfqpoint{3.356304in}{3.359462in}}%
\pgfpathlineto{\pgfqpoint{3.358186in}{3.369824in}}%
\pgfpathlineto{\pgfqpoint{3.361009in}{3.400366in}}%
\pgfpathlineto{\pgfqpoint{3.364860in}{3.469336in}}%
\pgfpathlineto{\pgfqpoint{3.370079in}{3.607072in}}%
\pgfpathlineto{\pgfqpoint{3.377523in}{3.869229in}}%
\pgfpathlineto{\pgfqpoint{3.401395in}{4.754754in}}%
\pgfpathlineto{\pgfqpoint{3.407299in}{4.880357in}}%
\pgfpathlineto{\pgfqpoint{3.411662in}{4.933222in}}%
\pgfpathlineto{\pgfqpoint{3.414743in}{4.948919in}}%
\pgfpathlineto{\pgfqpoint{3.416112in}{4.950105in}}%
\pgfpathlineto{\pgfqpoint{3.416368in}{4.949933in}}%
\pgfpathlineto{\pgfqpoint{3.417652in}{4.947210in}}%
\pgfpathlineto{\pgfqpoint{3.419705in}{4.936481in}}%
\pgfpathlineto{\pgfqpoint{3.422700in}{4.907162in}}%
\pgfpathlineto{\pgfqpoint{3.426807in}{4.842131in}}%
\pgfpathlineto{\pgfqpoint{3.432283in}{4.716031in}}%
\pgfpathlineto{\pgfqpoint{3.439898in}{4.483624in}}%
\pgfpathlineto{\pgfqpoint{3.468561in}{3.554146in}}%
\pgfpathlineto{\pgfqpoint{3.474636in}{3.445606in}}%
\pgfpathlineto{\pgfqpoint{3.479257in}{3.397769in}}%
\pgfpathlineto{\pgfqpoint{3.482508in}{3.382898in}}%
\pgfpathlineto{\pgfqpoint{3.484305in}{3.381371in}}%
\pgfpathlineto{\pgfqpoint{3.484391in}{3.381416in}}%
\pgfpathlineto{\pgfqpoint{3.485760in}{3.383593in}}%
\pgfpathlineto{\pgfqpoint{3.487813in}{3.391925in}}%
\pgfpathlineto{\pgfqpoint{3.490808in}{3.414697in}}%
\pgfpathlineto{\pgfqpoint{3.494915in}{3.465372in}}%
\pgfpathlineto{\pgfqpoint{3.500305in}{3.562668in}}%
\pgfpathlineto{\pgfqpoint{3.507492in}{3.736642in}}%
\pgfpathlineto{\pgfqpoint{3.518872in}{4.074553in}}%
\pgfpathlineto{\pgfqpoint{3.536413in}{4.587815in}}%
\pgfpathlineto{\pgfqpoint{3.544712in}{4.767223in}}%
\pgfpathlineto{\pgfqpoint{3.551044in}{4.860812in}}%
\pgfpathlineto{\pgfqpoint{3.555921in}{4.904309in}}%
\pgfpathlineto{\pgfqpoint{3.559514in}{4.919878in}}%
\pgfpathlineto{\pgfqpoint{3.561825in}{4.922493in}}%
\pgfpathlineto{\pgfqpoint{3.563279in}{4.921201in}}%
\pgfpathlineto{\pgfqpoint{3.565333in}{4.915566in}}%
\pgfpathlineto{\pgfqpoint{3.568242in}{4.900116in}}%
\pgfpathlineto{\pgfqpoint{3.572178in}{4.865859in}}%
\pgfpathlineto{\pgfqpoint{3.577311in}{4.799833in}}%
\pgfpathlineto{\pgfqpoint{3.583985in}{4.682732in}}%
\pgfpathlineto{\pgfqpoint{3.593140in}{4.478868in}}%
\pgfpathlineto{\pgfqpoint{3.629676in}{3.619212in}}%
\pgfpathlineto{\pgfqpoint{3.637291in}{3.509717in}}%
\pgfpathlineto{\pgfqpoint{3.643366in}{3.451274in}}%
\pgfpathlineto{\pgfqpoint{3.648072in}{3.424510in}}%
\pgfpathlineto{\pgfqpoint{3.651580in}{3.415151in}}%
\pgfpathlineto{\pgfqpoint{3.653890in}{3.413882in}}%
\pgfpathlineto{\pgfqpoint{3.655687in}{3.415543in}}%
\pgfpathlineto{\pgfqpoint{3.658168in}{3.421576in}}%
\pgfpathlineto{\pgfqpoint{3.661505in}{3.436347in}}%
\pgfpathlineto{\pgfqpoint{3.665954in}{3.467373in}}%
\pgfpathlineto{\pgfqpoint{3.671601in}{3.523963in}}%
\pgfpathlineto{\pgfqpoint{3.678788in}{3.620255in}}%
\pgfpathlineto{\pgfqpoint{3.688286in}{3.780022in}}%
\pgfpathlineto{\pgfqpoint{3.703602in}{4.081283in}}%
\pgfpathlineto{\pgfqpoint{3.724650in}{4.487062in}}%
\pgfpathlineto{\pgfqpoint{3.735687in}{4.658029in}}%
\pgfpathlineto{\pgfqpoint{3.744500in}{4.762729in}}%
\pgfpathlineto{\pgfqpoint{3.751773in}{4.825085in}}%
\pgfpathlineto{\pgfqpoint{3.757763in}{4.859384in}}%
\pgfpathlineto{\pgfqpoint{3.762554in}{4.875671in}}%
\pgfpathlineto{\pgfqpoint{3.766233in}{4.881535in}}%
\pgfpathlineto{\pgfqpoint{3.768886in}{4.882254in}}%
\pgfpathlineto{\pgfqpoint{3.771281in}{4.880429in}}%
\pgfpathlineto{\pgfqpoint{3.774276in}{4.874925in}}%
\pgfpathlineto{\pgfqpoint{3.778212in}{4.862428in}}%
\pgfpathlineto{\pgfqpoint{3.783260in}{4.838091in}}%
\pgfpathlineto{\pgfqpoint{3.789506in}{4.796058in}}%
\pgfpathlineto{\pgfqpoint{3.797292in}{4.727381in}}%
\pgfpathlineto{\pgfqpoint{3.807218in}{4.618467in}}%
\pgfpathlineto{\pgfqpoint{3.820908in}{4.441146in}}%
\pgfpathlineto{\pgfqpoint{3.866855in}{3.826096in}}%
\pgfpathlineto{\pgfqpoint{3.879347in}{3.697626in}}%
\pgfpathlineto{\pgfqpoint{3.889957in}{3.609999in}}%
\pgfpathlineto{\pgfqpoint{3.899197in}{3.550864in}}%
\pgfpathlineto{\pgfqpoint{3.907326in}{3.512265in}}%
\pgfpathlineto{\pgfqpoint{3.914342in}{3.488907in}}%
\pgfpathlineto{\pgfqpoint{3.920246in}{3.476194in}}%
\pgfpathlineto{\pgfqpoint{3.925208in}{3.470240in}}%
\pgfpathlineto{\pgfqpoint{3.929315in}{3.468458in}}%
\pgfpathlineto{\pgfqpoint{3.933080in}{3.469231in}}%
\pgfpathlineto{\pgfqpoint{3.937102in}{3.472501in}}%
\pgfpathlineto{\pgfqpoint{3.941893in}{3.479538in}}%
\pgfpathlineto{\pgfqpoint{3.947711in}{3.492403in}}%
\pgfpathlineto{\pgfqpoint{3.954727in}{3.513714in}}%
\pgfpathlineto{\pgfqpoint{3.963284in}{3.547355in}}%
\pgfpathlineto{\pgfqpoint{3.973722in}{3.598076in}}%
\pgfpathlineto{\pgfqpoint{3.986813in}{3.673511in}}%
\pgfpathlineto{\pgfqpoint{4.004610in}{3.790224in}}%
\pgfpathlineto{\pgfqpoint{4.075542in}{4.269321in}}%
\pgfpathlineto{\pgfqpoint{4.094194in}{4.373206in}}%
\pgfpathlineto{\pgfqpoint{4.111136in}{4.455432in}}%
\pgfpathlineto{\pgfqpoint{4.126965in}{4.521633in}}%
\pgfpathlineto{\pgfqpoint{4.142024in}{4.575361in}}%
\pgfpathlineto{\pgfqpoint{4.156484in}{4.618946in}}%
\pgfpathlineto{\pgfqpoint{4.170430in}{4.654124in}}%
\pgfpathlineto{\pgfqpoint{4.183949in}{4.682378in}}%
\pgfpathlineto{\pgfqpoint{4.197211in}{4.705077in}}%
\pgfpathlineto{\pgfqpoint{4.210217in}{4.723052in}}%
\pgfpathlineto{\pgfqpoint{4.223137in}{4.737215in}}%
\pgfpathlineto{\pgfqpoint{4.236142in}{4.748242in}}%
\pgfpathlineto{\pgfqpoint{4.249490in}{4.756667in}}%
\pgfpathlineto{\pgfqpoint{4.263351in}{4.762796in}}%
\pgfpathlineto{\pgfqpoint{4.278153in}{4.766903in}}%
\pgfpathlineto{\pgfqpoint{4.294496in}{4.769089in}}%
\pgfpathlineto{\pgfqpoint{4.313149in}{4.769273in}}%
\pgfpathlineto{\pgfqpoint{4.335480in}{4.767181in}}%
\pgfpathlineto{\pgfqpoint{4.364743in}{4.762062in}}%
\pgfpathlineto{\pgfqpoint{4.414198in}{4.750790in}}%
\pgfpathlineto{\pgfqpoint{4.515161in}{4.727824in}}%
\pgfpathlineto{\pgfqpoint{4.715977in}{4.684144in}}%
\pgfpathlineto{\pgfqpoint{4.756961in}{4.672226in}}%
\pgfpathlineto{\pgfqpoint{4.788020in}{4.660955in}}%
\pgfpathlineto{\pgfqpoint{4.813603in}{4.649354in}}%
\pgfpathlineto{\pgfqpoint{4.835764in}{4.636894in}}%
\pgfpathlineto{\pgfqpoint{4.855700in}{4.623148in}}%
\pgfpathlineto{\pgfqpoint{4.874096in}{4.607766in}}%
\pgfpathlineto{\pgfqpoint{4.891551in}{4.590248in}}%
\pgfpathlineto{\pgfqpoint{4.908321in}{4.570225in}}%
\pgfpathlineto{\pgfqpoint{4.924663in}{4.547210in}}%
\pgfpathlineto{\pgfqpoint{4.940920in}{4.520397in}}%
\pgfpathlineto{\pgfqpoint{4.957177in}{4.489198in}}%
\pgfpathlineto{\pgfqpoint{4.973605in}{4.452764in}}%
\pgfpathlineto{\pgfqpoint{4.990461in}{4.409855in}}%
\pgfpathlineto{\pgfqpoint{5.008001in}{4.358964in}}%
\pgfpathlineto{\pgfqpoint{5.026568in}{4.298080in}}%
\pgfpathlineto{\pgfqpoint{5.046932in}{4.223411in}}%
\pgfpathlineto{\pgfqpoint{5.071061in}{4.125992in}}%
\pgfpathlineto{\pgfqpoint{5.141136in}{3.836877in}}%
\pgfpathlineto{\pgfqpoint{5.155169in}{3.791716in}}%
\pgfpathlineto{\pgfqpoint{5.166377in}{3.762636in}}%
\pgfpathlineto{\pgfqpoint{5.175618in}{3.744418in}}%
\pgfpathlineto{\pgfqpoint{5.183319in}{3.733816in}}%
\pgfpathlineto{\pgfqpoint{5.189736in}{3.728497in}}%
\pgfpathlineto{\pgfqpoint{5.195212in}{3.726678in}}%
\pgfpathlineto{\pgfqpoint{5.200174in}{3.727312in}}%
\pgfpathlineto{\pgfqpoint{5.205223in}{3.730283in}}%
\pgfpathlineto{\pgfqpoint{5.210699in}{3.736257in}}%
\pgfpathlineto{\pgfqpoint{5.216859in}{3.746511in}}%
\pgfpathlineto{\pgfqpoint{5.223790in}{3.762633in}}%
\pgfpathlineto{\pgfqpoint{5.231661in}{3.786931in}}%
\pgfpathlineto{\pgfqpoint{5.240560in}{3.822062in}}%
\pgfpathlineto{\pgfqpoint{5.250656in}{3.871503in}}%
\pgfpathlineto{\pgfqpoint{5.262293in}{3.940201in}}%
\pgfpathlineto{\pgfqpoint{5.276411in}{4.037726in}}%
\pgfpathlineto{\pgfqpoint{5.297459in}{4.200893in}}%
\pgfpathlineto{\pgfqpoint{5.322272in}{4.389270in}}%
\pgfpathlineto{\pgfqpoint{5.334593in}{4.466245in}}%
\pgfpathlineto{\pgfqpoint{5.343919in}{4.511590in}}%
\pgfpathlineto{\pgfqpoint{5.351363in}{4.537754in}}%
\pgfpathlineto{\pgfqpoint{5.357267in}{4.551289in}}%
\pgfpathlineto{\pgfqpoint{5.361888in}{4.557063in}}%
\pgfpathlineto{\pgfqpoint{5.365567in}{4.558486in}}%
\pgfpathlineto{\pgfqpoint{5.368733in}{4.557389in}}%
\pgfpathlineto{\pgfqpoint{5.372241in}{4.553618in}}%
\pgfpathlineto{\pgfqpoint{5.376433in}{4.545540in}}%
\pgfpathlineto{\pgfqpoint{5.381567in}{4.530324in}}%
\pgfpathlineto{\pgfqpoint{5.387642in}{4.504814in}}%
\pgfpathlineto{\pgfqpoint{5.394915in}{4.463943in}}%
\pgfpathlineto{\pgfqpoint{5.403642in}{4.401219in}}%
\pgfpathlineto{\pgfqpoint{5.414508in}{4.305731in}}%
\pgfpathlineto{\pgfqpoint{5.431193in}{4.136413in}}%
\pgfpathlineto{\pgfqpoint{5.451300in}{3.937637in}}%
\pgfpathlineto{\pgfqpoint{5.461311in}{3.859798in}}%
\pgfpathlineto{\pgfqpoint{5.468841in}{3.816989in}}%
\pgfpathlineto{\pgfqpoint{5.474744in}{3.794913in}}%
\pgfpathlineto{\pgfqpoint{5.479194in}{3.785666in}}%
\pgfpathlineto{\pgfqpoint{5.482445in}{3.783148in}}%
\pgfpathlineto{\pgfqpoint{5.485012in}{3.783759in}}%
\pgfpathlineto{\pgfqpoint{5.487750in}{3.786974in}}%
\pgfpathlineto{\pgfqpoint{5.491258in}{3.794984in}}%
\pgfpathlineto{\pgfqpoint{5.495707in}{3.811404in}}%
\pgfpathlineto{\pgfqpoint{5.501183in}{3.840992in}}%
\pgfpathlineto{\pgfqpoint{5.507943in}{3.890906in}}%
\pgfpathlineto{\pgfqpoint{5.516413in}{3.971574in}}%
\pgfpathlineto{\pgfqpoint{5.528221in}{4.107736in}}%
\pgfpathlineto{\pgfqpoint{5.552521in}{4.391915in}}%
\pgfpathlineto{\pgfqpoint{5.560820in}{4.461565in}}%
\pgfpathlineto{\pgfqpoint{5.566981in}{4.496651in}}%
\pgfpathlineto{\pgfqpoint{5.571687in}{4.512313in}}%
\pgfpathlineto{\pgfqpoint{5.575024in}{4.517123in}}%
\pgfpathlineto{\pgfqpoint{5.577334in}{4.517289in}}%
\pgfpathlineto{\pgfqpoint{5.579558in}{4.514966in}}%
\pgfpathlineto{\pgfqpoint{5.582468in}{4.508240in}}%
\pgfpathlineto{\pgfqpoint{5.586318in}{4.492952in}}%
\pgfpathlineto{\pgfqpoint{5.591195in}{4.463449in}}%
\pgfpathlineto{\pgfqpoint{5.597355in}{4.411156in}}%
\pgfpathlineto{\pgfqpoint{5.605313in}{4.322724in}}%
\pgfpathlineto{\pgfqpoint{5.617634in}{4.157096in}}%
\pgfpathlineto{\pgfqpoint{5.633976in}{3.943290in}}%
\pgfpathlineto{\pgfqpoint{5.641762in}{3.869200in}}%
\pgfpathlineto{\pgfqpoint{5.647495in}{3.833640in}}%
\pgfpathlineto{\pgfqpoint{5.651688in}{3.819563in}}%
\pgfpathlineto{\pgfqpoint{5.654511in}{3.816158in}}%
\pgfpathlineto{\pgfqpoint{5.656394in}{3.816678in}}%
\pgfpathlineto{\pgfqpoint{5.658533in}{3.820002in}}%
\pgfpathlineto{\pgfqpoint{5.661442in}{3.829181in}}%
\pgfpathlineto{\pgfqpoint{5.665378in}{3.849994in}}%
\pgfpathlineto{\pgfqpoint{5.670426in}{3.890093in}}%
\pgfpathlineto{\pgfqpoint{5.676929in}{3.961288in}}%
\pgfpathlineto{\pgfqpoint{5.686169in}{4.089895in}}%
\pgfpathlineto{\pgfqpoint{5.707389in}{4.392545in}}%
\pgfpathlineto{\pgfqpoint{5.713977in}{4.453366in}}%
\pgfpathlineto{\pgfqpoint{5.718769in}{4.479905in}}%
\pgfpathlineto{\pgfqpoint{5.722191in}{4.488587in}}%
\pgfpathlineto{\pgfqpoint{5.724330in}{4.489459in}}%
\pgfpathlineto{\pgfqpoint{5.726127in}{4.487443in}}%
\pgfpathlineto{\pgfqpoint{5.728608in}{4.480531in}}%
\pgfpathlineto{\pgfqpoint{5.731945in}{4.463789in}}%
\pgfpathlineto{\pgfqpoint{5.736394in}{4.428776in}}%
\pgfpathlineto{\pgfqpoint{5.742213in}{4.363478in}}%
\pgfpathlineto{\pgfqpoint{5.750427in}{4.243206in}}%
\pgfpathlineto{\pgfqpoint{5.771646in}{3.921120in}}%
\pgfpathlineto{\pgfqpoint{5.777550in}{3.867439in}}%
\pgfpathlineto{\pgfqpoint{5.781828in}{3.845993in}}%
\pgfpathlineto{\pgfqpoint{5.784652in}{3.840701in}}%
\pgfpathlineto{\pgfqpoint{5.786277in}{3.840960in}}%
\pgfpathlineto{\pgfqpoint{5.788074in}{3.844080in}}%
\pgfpathlineto{\pgfqpoint{5.790641in}{3.853674in}}%
\pgfpathlineto{\pgfqpoint{5.794235in}{3.876998in}}%
\pgfpathlineto{\pgfqpoint{5.798941in}{3.923762in}}%
\pgfpathlineto{\pgfqpoint{5.805358in}{4.012170in}}%
\pgfpathlineto{\pgfqpoint{5.816738in}{4.206508in}}%
\pgfpathlineto{\pgfqpoint{5.827433in}{4.373759in}}%
\pgfpathlineto{\pgfqpoint{5.833422in}{4.436428in}}%
\pgfpathlineto{\pgfqpoint{5.837700in}{4.461305in}}%
\pgfpathlineto{\pgfqpoint{5.840524in}{4.467534in}}%
\pgfpathlineto{\pgfqpoint{5.842150in}{4.467303in}}%
\pgfpathlineto{\pgfqpoint{5.843861in}{4.464023in}}%
\pgfpathlineto{\pgfqpoint{5.846428in}{4.453302in}}%
\pgfpathlineto{\pgfqpoint{5.849936in}{4.427745in}}%
\pgfpathlineto{\pgfqpoint{5.854642in}{4.375358in}}%
\pgfpathlineto{\pgfqpoint{5.861230in}{4.274343in}}%
\pgfpathlineto{\pgfqpoint{5.883733in}{3.906474in}}%
\pgfpathlineto{\pgfqpoint{5.888353in}{3.871260in}}%
\pgfpathlineto{\pgfqpoint{5.891519in}{3.861015in}}%
\pgfpathlineto{\pgfqpoint{5.893230in}{3.860441in}}%
\pgfpathlineto{\pgfqpoint{5.894771in}{3.862941in}}%
\pgfpathlineto{\pgfqpoint{5.897081in}{3.872021in}}%
\pgfpathlineto{\pgfqpoint{5.900332in}{3.895332in}}%
\pgfpathlineto{\pgfqpoint{5.904781in}{3.945569in}}%
\pgfpathlineto{\pgfqpoint{5.911027in}{4.044603in}}%
\pgfpathlineto{\pgfqpoint{5.932504in}{4.408652in}}%
\pgfpathlineto{\pgfqpoint{5.936867in}{4.441182in}}%
\pgfpathlineto{\pgfqpoint{5.939691in}{4.449336in}}%
\pgfpathlineto{\pgfqpoint{5.941231in}{4.449305in}}%
\pgfpathlineto{\pgfqpoint{5.942771in}{4.446085in}}%
\pgfpathlineto{\pgfqpoint{5.945167in}{4.434795in}}%
\pgfpathlineto{\pgfqpoint{5.948504in}{4.406815in}}%
\pgfpathlineto{\pgfqpoint{5.953124in}{4.346994in}}%
\pgfpathlineto{\pgfqpoint{5.959969in}{4.225505in}}%
\pgfpathlineto{\pgfqpoint{5.975456in}{3.943515in}}%
\pgfpathlineto{\pgfqpoint{5.980333in}{3.894045in}}%
\pgfpathlineto{\pgfqpoint{5.983670in}{3.878594in}}%
\pgfpathlineto{\pgfqpoint{5.985467in}{3.876982in}}%
\pgfpathlineto{\pgfqpoint{5.985552in}{3.877025in}}%
\pgfpathlineto{\pgfqpoint{5.986921in}{3.879181in}}%
\pgfpathlineto{\pgfqpoint{5.988975in}{3.887595in}}%
\pgfpathlineto{\pgfqpoint{5.991969in}{3.910675in}}%
\pgfpathlineto{\pgfqpoint{5.996162in}{3.962626in}}%
\pgfpathlineto{\pgfqpoint{6.002237in}{4.069191in}}%
\pgfpathlineto{\pgfqpoint{6.019435in}{4.386039in}}%
\pgfpathlineto{\pgfqpoint{6.023799in}{4.424408in}}%
\pgfpathlineto{\pgfqpoint{6.026622in}{4.433843in}}%
\pgfpathlineto{\pgfqpoint{6.027906in}{4.433904in}}%
\pgfpathlineto{\pgfqpoint{6.027991in}{4.433813in}}%
\pgfpathlineto{\pgfqpoint{6.029446in}{4.430459in}}%
\pgfpathlineto{\pgfqpoint{6.031756in}{4.418181in}}%
\pgfpathlineto{\pgfqpoint{6.035007in}{4.387113in}}%
\pgfpathlineto{\pgfqpoint{6.039628in}{4.318825in}}%
\pgfpathlineto{\pgfqpoint{6.047157in}{4.168626in}}%
\pgfpathlineto{\pgfqpoint{6.057853in}{3.963270in}}%
\pgfpathlineto{\pgfqpoint{6.062644in}{3.909024in}}%
\pgfpathlineto{\pgfqpoint{6.065810in}{3.893020in}}%
\pgfpathlineto{\pgfqpoint{6.067436in}{3.891499in}}%
\pgfpathlineto{\pgfqpoint{6.067607in}{3.891608in}}%
\pgfpathlineto{\pgfqpoint{6.068890in}{3.894062in}}%
\pgfpathlineto{\pgfqpoint{6.070944in}{3.903939in}}%
\pgfpathlineto{\pgfqpoint{6.073938in}{3.930959in}}%
\pgfpathlineto{\pgfqpoint{6.078216in}{3.992548in}}%
\pgfpathlineto{\pgfqpoint{6.084976in}{4.127466in}}%
\pgfpathlineto{\pgfqpoint{6.096527in}{4.355469in}}%
\pgfpathlineto{\pgfqpoint{6.101147in}{4.406386in}}%
\pgfpathlineto{\pgfqpoint{6.104142in}{4.419788in}}%
\pgfpathlineto{\pgfqpoint{6.105340in}{4.420485in}}%
\pgfpathlineto{\pgfqpoint{6.105682in}{4.420187in}}%
\pgfpathlineto{\pgfqpoint{6.107051in}{4.416792in}}%
\pgfpathlineto{\pgfqpoint{6.109276in}{4.403854in}}%
\pgfpathlineto{\pgfqpoint{6.112527in}{4.369314in}}%
\pgfpathlineto{\pgfqpoint{6.117147in}{4.293237in}}%
\pgfpathlineto{\pgfqpoint{6.126217in}{4.097153in}}%
\pgfpathlineto{\pgfqpoint{6.133661in}{3.959151in}}%
\pgfpathlineto{\pgfqpoint{6.137939in}{3.914923in}}%
\pgfpathlineto{\pgfqpoint{6.140677in}{3.904587in}}%
\pgfpathlineto{\pgfqpoint{6.141704in}{3.904566in}}%
\pgfpathlineto{\pgfqpoint{6.141875in}{3.904769in}}%
\pgfpathlineto{\pgfqpoint{6.143244in}{3.908513in}}%
\pgfpathlineto{\pgfqpoint{6.145469in}{3.922502in}}%
\pgfpathlineto{\pgfqpoint{6.148720in}{3.959487in}}%
\pgfpathlineto{\pgfqpoint{6.153511in}{4.043484in}}%
\pgfpathlineto{\pgfqpoint{6.171565in}{4.389460in}}%
\pgfpathlineto{\pgfqpoint{6.174645in}{4.406936in}}%
\pgfpathlineto{\pgfqpoint{6.176014in}{4.408312in}}%
\pgfpathlineto{\pgfqpoint{6.176271in}{4.408124in}}%
\pgfpathlineto{\pgfqpoint{6.177554in}{4.405068in}}%
\pgfpathlineto{\pgfqpoint{6.179608in}{4.392950in}}%
\pgfpathlineto{\pgfqpoint{6.182688in}{4.359012in}}%
\pgfpathlineto{\pgfqpoint{6.187223in}{4.280273in}}%
\pgfpathlineto{\pgfqpoint{6.205533in}{3.929424in}}%
\pgfpathlineto{\pgfqpoint{6.208357in}{3.916388in}}%
\pgfpathlineto{\pgfqpoint{6.209298in}{3.916005in}}%
\pgfpathlineto{\pgfqpoint{6.209640in}{3.916363in}}%
\pgfpathlineto{\pgfqpoint{6.211009in}{3.920444in}}%
\pgfpathlineto{\pgfqpoint{6.213234in}{3.935935in}}%
\pgfpathlineto{\pgfqpoint{6.216485in}{3.976885in}}%
\pgfpathlineto{\pgfqpoint{6.221362in}{4.070468in}}%
\pgfpathlineto{\pgfqpoint{6.235908in}{4.367609in}}%
\pgfpathlineto{\pgfqpoint{6.239331in}{4.394161in}}%
\pgfpathlineto{\pgfqpoint{6.241042in}{4.397300in}}%
\pgfpathlineto{\pgfqpoint{6.241213in}{4.397230in}}%
\pgfpathlineto{\pgfqpoint{6.242325in}{4.395073in}}%
\pgfpathlineto{\pgfqpoint{6.244122in}{4.385416in}}%
\pgfpathlineto{\pgfqpoint{6.246860in}{4.356728in}}%
\pgfpathlineto{\pgfqpoint{6.250967in}{4.286667in}}%
\pgfpathlineto{\pgfqpoint{6.259181in}{4.097255in}}%
\pgfpathlineto{\pgfqpoint{6.265769in}{3.969927in}}%
\pgfpathlineto{\pgfqpoint{6.269534in}{3.932662in}}%
\pgfpathlineto{\pgfqpoint{6.271759in}{3.926468in}}%
\pgfpathlineto{\pgfqpoint{6.272700in}{3.927568in}}%
\pgfpathlineto{\pgfqpoint{6.274240in}{3.934133in}}%
\pgfpathlineto{\pgfqpoint{6.276721in}{3.956723in}}%
\pgfpathlineto{\pgfqpoint{6.280401in}{4.014383in}}%
\pgfpathlineto{\pgfqpoint{6.286647in}{4.154755in}}%
\pgfpathlineto{\pgfqpoint{6.295288in}{4.338194in}}%
\pgfpathlineto{\pgfqpoint{6.299139in}{4.379774in}}%
\pgfpathlineto{\pgfqpoint{6.301449in}{4.387191in}}%
\pgfpathlineto{\pgfqpoint{6.302390in}{4.386217in}}%
\pgfpathlineto{\pgfqpoint{6.303845in}{4.380163in}}%
\pgfpathlineto{\pgfqpoint{6.306240in}{4.358547in}}%
\pgfpathlineto{\pgfqpoint{6.309834in}{4.302021in}}%
\pgfpathlineto{\pgfqpoint{6.315909in}{4.164092in}}%
\pgfpathlineto{\pgfqpoint{6.324380in}{3.982772in}}%
\pgfpathlineto{\pgfqpoint{6.328144in}{3.942665in}}%
\pgfpathlineto{\pgfqpoint{6.330369in}{3.936168in}}%
\pgfpathlineto{\pgfqpoint{6.331310in}{3.937484in}}%
\pgfpathlineto{\pgfqpoint{6.332850in}{3.944832in}}%
\pgfpathlineto{\pgfqpoint{6.335332in}{3.969682in}}%
\pgfpathlineto{\pgfqpoint{6.339096in}{4.033904in}}%
\pgfpathlineto{\pgfqpoint{6.346369in}{4.206520in}}%
\pgfpathlineto{\pgfqpoint{6.352957in}{4.339252in}}%
\pgfpathlineto{\pgfqpoint{6.356466in}{4.373361in}}%
\pgfpathlineto{\pgfqpoint{6.358348in}{4.377866in}}%
\pgfpathlineto{\pgfqpoint{6.358433in}{4.377832in}}%
\pgfpathlineto{\pgfqpoint{6.359375in}{4.376076in}}%
\pgfpathlineto{\pgfqpoint{6.361000in}{4.367160in}}%
\pgfpathlineto{\pgfqpoint{6.363567in}{4.338662in}}%
\pgfpathlineto{\pgfqpoint{6.367503in}{4.266213in}}%
\pgfpathlineto{\pgfqpoint{6.382819in}{3.954564in}}%
\pgfpathlineto{\pgfqpoint{6.385215in}{3.945163in}}%
\pgfpathlineto{\pgfqpoint{6.386070in}{3.945896in}}%
\pgfpathlineto{\pgfqpoint{6.387439in}{3.951555in}}%
\pgfpathlineto{\pgfqpoint{6.389664in}{3.972133in}}%
\pgfpathlineto{\pgfqpoint{6.393086in}{4.028077in}}%
\pgfpathlineto{\pgfqpoint{6.399247in}{4.173869in}}%
\pgfpathlineto{\pgfqpoint{6.406520in}{4.329515in}}%
\pgfpathlineto{\pgfqpoint{6.410028in}{4.364971in}}%
\pgfpathlineto{\pgfqpoint{6.411739in}{4.369232in}}%
\pgfpathlineto{\pgfqpoint{6.411910in}{4.369166in}}%
\pgfpathlineto{\pgfqpoint{6.412851in}{4.367196in}}%
\pgfpathlineto{\pgfqpoint{6.414562in}{4.356750in}}%
\pgfpathlineto{\pgfqpoint{6.417215in}{4.324112in}}%
\pgfpathlineto{\pgfqpoint{6.421407in}{4.239909in}}%
\pgfpathlineto{\pgfqpoint{6.433814in}{3.972821in}}%
\pgfpathlineto{\pgfqpoint{6.436638in}{3.954253in}}%
\pgfpathlineto{\pgfqpoint{6.437493in}{3.953518in}}%
\pgfpathlineto{\pgfqpoint{6.437921in}{3.954021in}}%
\pgfpathlineto{\pgfqpoint{6.439204in}{3.958997in}}%
\pgfpathlineto{\pgfqpoint{6.441343in}{3.978486in}}%
\pgfpathlineto{\pgfqpoint{6.444595in}{4.031827in}}%
\pgfpathlineto{\pgfqpoint{6.450499in}{4.173242in}}%
\pgfpathlineto{\pgfqpoint{6.457515in}{4.324676in}}%
\pgfpathlineto{\pgfqpoint{6.460852in}{4.357677in}}%
\pgfpathlineto{\pgfqpoint{6.462392in}{4.361162in}}%
\pgfpathlineto{\pgfqpoint{6.462649in}{4.360989in}}%
\pgfpathlineto{\pgfqpoint{6.463675in}{4.358146in}}%
\pgfpathlineto{\pgfqpoint{6.465472in}{4.345041in}}%
\pgfpathlineto{\pgfqpoint{6.468296in}{4.305218in}}%
\pgfpathlineto{\pgfqpoint{6.472916in}{4.202850in}}%
\pgfpathlineto{\pgfqpoint{6.482413in}{3.991095in}}%
\pgfpathlineto{\pgfqpoint{6.485579in}{3.963303in}}%
\pgfpathlineto{\pgfqpoint{6.486777in}{3.961297in}}%
\pgfpathlineto{\pgfqpoint{6.487119in}{3.961612in}}%
\pgfpathlineto{\pgfqpoint{6.488232in}{3.965355in}}%
\pgfpathlineto{\pgfqpoint{6.490200in}{3.981905in}}%
\pgfpathlineto{\pgfqpoint{6.493280in}{4.030630in}}%
\pgfpathlineto{\pgfqpoint{6.498670in}{4.159013in}}%
\pgfpathlineto{\pgfqpoint{6.505943in}{4.319387in}}%
\pgfpathlineto{\pgfqpoint{6.509195in}{4.350863in}}%
\pgfpathlineto{\pgfqpoint{6.510564in}{4.353597in}}%
\pgfpathlineto{\pgfqpoint{6.510906in}{4.353266in}}%
\pgfpathlineto{\pgfqpoint{6.512018in}{4.349393in}}%
\pgfpathlineto{\pgfqpoint{6.513986in}{4.332348in}}%
\pgfpathlineto{\pgfqpoint{6.517066in}{4.282377in}}%
\pgfpathlineto{\pgfqpoint{6.522713in}{4.145525in}}%
\pgfpathlineto{\pgfqpoint{6.529387in}{4.000625in}}%
\pgfpathlineto{\pgfqpoint{6.532553in}{3.970903in}}%
\pgfpathlineto{\pgfqpoint{6.533751in}{3.968609in}}%
\pgfpathlineto{\pgfqpoint{6.534179in}{3.969028in}}%
\pgfpathlineto{\pgfqpoint{6.535377in}{3.973655in}}%
\pgfpathlineto{\pgfqpoint{6.537430in}{3.993046in}}%
\pgfpathlineto{\pgfqpoint{6.540596in}{4.047716in}}%
\pgfpathlineto{\pgfqpoint{6.547184in}{4.211597in}}%
\pgfpathlineto{\pgfqpoint{6.552660in}{4.321414in}}%
\pgfpathlineto{\pgfqpoint{6.555569in}{4.345223in}}%
\pgfpathlineto{\pgfqpoint{6.556511in}{4.346494in}}%
\pgfpathlineto{\pgfqpoint{6.556938in}{4.346002in}}%
\pgfpathlineto{\pgfqpoint{6.558136in}{4.341082in}}%
\pgfpathlineto{\pgfqpoint{6.560190in}{4.320936in}}%
\pgfpathlineto{\pgfqpoint{6.563441in}{4.262910in}}%
\pgfpathlineto{\pgfqpoint{6.577473in}{3.978056in}}%
\pgfpathlineto{\pgfqpoint{6.578671in}{3.975502in}}%
\pgfpathlineto{\pgfqpoint{6.579099in}{3.975891in}}%
\pgfpathlineto{\pgfqpoint{6.580211in}{3.980097in}}%
\pgfpathlineto{\pgfqpoint{6.582179in}{3.998482in}}%
\pgfpathlineto{\pgfqpoint{6.585260in}{4.051809in}}%
\pgfpathlineto{\pgfqpoint{6.591848in}{4.216691in}}%
\pgfpathlineto{\pgfqpoint{6.597067in}{4.318760in}}%
\pgfpathlineto{\pgfqpoint{6.599805in}{4.339097in}}%
\pgfpathlineto{\pgfqpoint{6.600575in}{4.339760in}}%
\pgfpathlineto{\pgfqpoint{6.601003in}{4.339146in}}%
\pgfpathlineto{\pgfqpoint{6.602201in}{4.333722in}}%
\pgfpathlineto{\pgfqpoint{6.604254in}{4.312227in}}%
\pgfpathlineto{\pgfqpoint{6.607591in}{4.249518in}}%
\pgfpathlineto{\pgfqpoint{6.620083in}{3.987437in}}%
\pgfpathlineto{\pgfqpoint{6.621795in}{3.982023in}}%
\pgfpathlineto{\pgfqpoint{6.621966in}{3.982110in}}%
\pgfpathlineto{\pgfqpoint{6.622907in}{3.984636in}}%
\pgfpathlineto{\pgfqpoint{6.624618in}{3.997916in}}%
\pgfpathlineto{\pgfqpoint{6.627356in}{4.040246in}}%
\pgfpathlineto{\pgfqpoint{6.632233in}{4.157569in}}%
\pgfpathlineto{\pgfqpoint{6.638993in}{4.307104in}}%
\pgfpathlineto{\pgfqpoint{6.641902in}{4.332282in}}%
\pgfpathlineto{\pgfqpoint{6.642758in}{4.333451in}}%
\pgfpathlineto{\pgfqpoint{6.643185in}{4.332936in}}%
\pgfpathlineto{\pgfqpoint{6.644383in}{4.327627in}}%
\pgfpathlineto{\pgfqpoint{6.646437in}{4.305809in}}%
\pgfpathlineto{\pgfqpoint{6.649774in}{4.241694in}}%
\pgfpathlineto{\pgfqpoint{6.661410in}{3.994673in}}%
\pgfpathlineto{\pgfqpoint{6.663207in}{3.988197in}}%
\pgfpathlineto{\pgfqpoint{6.663378in}{3.988264in}}%
\pgfpathlineto{\pgfqpoint{6.664234in}{3.990387in}}%
\pgfpathlineto{\pgfqpoint{6.665859in}{4.002479in}}%
\pgfpathlineto{\pgfqpoint{6.668512in}{4.042791in}}%
\pgfpathlineto{\pgfqpoint{6.673218in}{4.155709in}}%
\pgfpathlineto{\pgfqpoint{6.679892in}{4.303228in}}%
\pgfpathlineto{\pgfqpoint{6.682715in}{4.326616in}}%
\pgfpathlineto{\pgfqpoint{6.683485in}{4.327407in}}%
\pgfpathlineto{\pgfqpoint{6.683913in}{4.326780in}}%
\pgfpathlineto{\pgfqpoint{6.685111in}{4.321010in}}%
\pgfpathlineto{\pgfqpoint{6.687164in}{4.297988in}}%
\pgfpathlineto{\pgfqpoint{6.690587in}{4.229498in}}%
\pgfpathlineto{\pgfqpoint{6.701026in}{4.003428in}}%
\pgfpathlineto{\pgfqpoint{6.703165in}{3.994066in}}%
\pgfpathlineto{\pgfqpoint{6.703935in}{3.995407in}}%
\pgfpathlineto{\pgfqpoint{6.705304in}{4.003908in}}%
\pgfpathlineto{\pgfqpoint{6.707614in}{4.034772in}}%
\pgfpathlineto{\pgfqpoint{6.711550in}{4.122841in}}%
\pgfpathlineto{\pgfqpoint{6.719593in}{4.302792in}}%
\pgfpathlineto{\pgfqpoint{6.722159in}{4.321329in}}%
\pgfpathlineto{\pgfqpoint{6.722758in}{4.321659in}}%
\pgfpathlineto{\pgfqpoint{6.723186in}{4.320949in}}%
\pgfpathlineto{\pgfqpoint{6.724470in}{4.314136in}}%
\pgfpathlineto{\pgfqpoint{6.726609in}{4.288046in}}%
\pgfpathlineto{\pgfqpoint{6.730202in}{4.211780in}}%
\pgfpathlineto{\pgfqpoint{6.739186in}{4.013369in}}%
\pgfpathlineto{\pgfqpoint{6.741582in}{3.999685in}}%
\pgfpathlineto{\pgfqpoint{6.742352in}{4.000561in}}%
\pgfpathlineto{\pgfqpoint{6.743636in}{4.007733in}}%
\pgfpathlineto{\pgfqpoint{6.745860in}{4.036011in}}%
\pgfpathlineto{\pgfqpoint{6.749625in}{4.118571in}}%
\pgfpathlineto{\pgfqpoint{6.757753in}{4.299616in}}%
\pgfpathlineto{\pgfqpoint{6.760235in}{4.316095in}}%
\pgfpathlineto{\pgfqpoint{6.760748in}{4.316147in}}%
\pgfpathlineto{\pgfqpoint{6.761090in}{4.315531in}}%
\pgfpathlineto{\pgfqpoint{6.762288in}{4.309314in}}%
\pgfpathlineto{\pgfqpoint{6.764427in}{4.283373in}}%
\pgfpathlineto{\pgfqpoint{6.768021in}{4.206552in}}%
\pgfpathlineto{\pgfqpoint{6.776577in}{4.018263in}}%
\pgfpathlineto{\pgfqpoint{6.778887in}{4.005018in}}%
\pgfpathlineto{\pgfqpoint{6.779572in}{4.005665in}}%
\pgfpathlineto{\pgfqpoint{6.779657in}{4.005894in}}%
\pgfpathlineto{\pgfqpoint{6.780941in}{4.013240in}}%
\pgfpathlineto{\pgfqpoint{6.783166in}{4.042231in}}%
\pgfpathlineto{\pgfqpoint{6.787016in}{4.128368in}}%
\pgfpathlineto{\pgfqpoint{6.794545in}{4.294758in}}%
\pgfpathlineto{\pgfqpoint{6.797027in}{4.310961in}}%
\pgfpathlineto{\pgfqpoint{6.797540in}{4.310854in}}%
\pgfpathlineto{\pgfqpoint{6.797797in}{4.310348in}}%
\pgfpathlineto{\pgfqpoint{6.798995in}{4.304042in}}%
\pgfpathlineto{\pgfqpoint{6.801048in}{4.278964in}}%
\pgfpathlineto{\pgfqpoint{6.804556in}{4.203897in}}%
\pgfpathlineto{\pgfqpoint{6.812770in}{4.023329in}}%
\pgfpathlineto{\pgfqpoint{6.815080in}{4.010090in}}%
\pgfpathlineto{\pgfqpoint{6.815765in}{4.010867in}}%
\pgfpathlineto{\pgfqpoint{6.816963in}{4.017410in}}%
\pgfpathlineto{\pgfqpoint{6.819102in}{4.044455in}}%
\pgfpathlineto{\pgfqpoint{6.822781in}{4.125482in}}%
\pgfpathlineto{\pgfqpoint{6.830310in}{4.291157in}}%
\pgfpathlineto{\pgfqpoint{6.832706in}{4.306035in}}%
\pgfpathlineto{\pgfqpoint{6.833305in}{4.305601in}}%
\pgfpathlineto{\pgfqpoint{6.833391in}{4.305401in}}%
\pgfpathlineto{\pgfqpoint{6.834589in}{4.299038in}}%
\pgfpathlineto{\pgfqpoint{6.836642in}{4.273526in}}%
\pgfpathlineto{\pgfqpoint{6.840236in}{4.195322in}}%
\pgfpathlineto{\pgfqpoint{6.847851in}{4.028637in}}%
\pgfpathlineto{\pgfqpoint{6.850161in}{4.014966in}}%
\pgfpathlineto{\pgfqpoint{6.850845in}{4.015731in}}%
\pgfpathlineto{\pgfqpoint{6.852043in}{4.022382in}}%
\pgfpathlineto{\pgfqpoint{6.854182in}{4.049948in}}%
\pgfpathlineto{\pgfqpoint{6.857947in}{4.134137in}}%
\pgfpathlineto{\pgfqpoint{6.864963in}{4.286809in}}%
\pgfpathlineto{\pgfqpoint{6.867273in}{4.301253in}}%
\pgfpathlineto{\pgfqpoint{6.867787in}{4.301025in}}%
\pgfpathlineto{\pgfqpoint{6.868043in}{4.300436in}}%
\pgfpathlineto{\pgfqpoint{6.869327in}{4.292802in}}%
\pgfpathlineto{\pgfqpoint{6.871552in}{4.262421in}}%
\pgfpathlineto{\pgfqpoint{6.875573in}{4.169455in}}%
\pgfpathlineto{\pgfqpoint{6.881819in}{4.035167in}}%
\pgfpathlineto{\pgfqpoint{6.884215in}{4.019664in}}%
\pgfpathlineto{\pgfqpoint{6.884728in}{4.019918in}}%
\pgfpathlineto{\pgfqpoint{6.884899in}{4.020288in}}%
\pgfpathlineto{\pgfqpoint{6.886097in}{4.026817in}}%
\pgfpathlineto{\pgfqpoint{6.888151in}{4.053047in}}%
\pgfpathlineto{\pgfqpoint{6.891830in}{4.134742in}}%
\pgfpathlineto{\pgfqpoint{6.898675in}{4.282833in}}%
\pgfpathlineto{\pgfqpoint{6.900985in}{4.296729in}}%
\pgfpathlineto{\pgfqpoint{6.901670in}{4.295868in}}%
\pgfpathlineto{\pgfqpoint{6.902867in}{4.288877in}}%
\pgfpathlineto{\pgfqpoint{6.905006in}{4.260312in}}%
\pgfpathlineto{\pgfqpoint{6.908857in}{4.172534in}}%
\pgfpathlineto{\pgfqpoint{6.915103in}{4.038740in}}%
\pgfpathlineto{\pgfqpoint{6.917413in}{4.024131in}}%
\pgfpathlineto{\pgfqpoint{6.918012in}{4.024609in}}%
\pgfpathlineto{\pgfqpoint{6.918097in}{4.024823in}}%
\pgfpathlineto{\pgfqpoint{6.919295in}{4.031581in}}%
\pgfpathlineto{\pgfqpoint{6.921434in}{4.059921in}}%
\pgfpathlineto{\pgfqpoint{6.925285in}{4.147578in}}%
\pgfpathlineto{\pgfqpoint{6.931445in}{4.278460in}}%
\pgfpathlineto{\pgfqpoint{6.933755in}{4.292349in}}%
\pgfpathlineto{\pgfqpoint{6.934440in}{4.291387in}}%
\pgfpathlineto{\pgfqpoint{6.935723in}{4.283333in}}%
\pgfpathlineto{\pgfqpoint{6.937948in}{4.251714in}}%
\pgfpathlineto{\pgfqpoint{6.942226in}{4.150864in}}%
\pgfpathlineto{\pgfqpoint{6.947531in}{4.042093in}}%
\pgfpathlineto{\pgfqpoint{6.949841in}{4.028388in}}%
\pgfpathlineto{\pgfqpoint{6.950526in}{4.029453in}}%
\pgfpathlineto{\pgfqpoint{6.951809in}{4.037750in}}%
\pgfpathlineto{\pgfqpoint{6.954034in}{4.069857in}}%
\pgfpathlineto{\pgfqpoint{6.958397in}{4.173206in}}%
\pgfpathlineto{\pgfqpoint{6.963446in}{4.274964in}}%
\pgfpathlineto{\pgfqpoint{6.965670in}{4.288131in}}%
\pgfpathlineto{\pgfqpoint{6.966355in}{4.287158in}}%
\pgfpathlineto{\pgfqpoint{6.967638in}{4.278975in}}%
\pgfpathlineto{\pgfqpoint{6.969863in}{4.246914in}}%
\pgfpathlineto{\pgfqpoint{6.974312in}{4.141530in}}%
\pgfpathlineto{\pgfqpoint{6.979275in}{4.043987in}}%
\pgfpathlineto{\pgfqpoint{6.981414in}{4.032522in}}%
\pgfpathlineto{\pgfqpoint{6.982098in}{4.033782in}}%
\pgfpathlineto{\pgfqpoint{6.983382in}{4.042538in}}%
\pgfpathlineto{\pgfqpoint{6.985692in}{4.077203in}}%
\pgfpathlineto{\pgfqpoint{6.990654in}{4.196238in}}%
\pgfpathlineto{\pgfqpoint{6.995018in}{4.275223in}}%
\pgfpathlineto{\pgfqpoint{6.996901in}{4.284082in}}%
\pgfpathlineto{\pgfqpoint{6.997072in}{4.283983in}}%
\pgfpathlineto{\pgfqpoint{6.997927in}{4.281213in}}%
\pgfpathlineto{\pgfqpoint{6.999553in}{4.265880in}}%
\pgfpathlineto{\pgfqpoint{7.002377in}{4.213271in}}%
\pgfpathlineto{\pgfqpoint{7.010933in}{4.040486in}}%
\pgfpathlineto{\pgfqpoint{7.012216in}{4.036509in}}%
\pgfpathlineto{\pgfqpoint{7.012558in}{4.036901in}}%
\pgfpathlineto{\pgfqpoint{7.013585in}{4.041725in}}%
\pgfpathlineto{\pgfqpoint{7.015468in}{4.063946in}}%
\pgfpathlineto{\pgfqpoint{7.018890in}{4.137052in}}%
\pgfpathlineto{\pgfqpoint{7.025307in}{4.269985in}}%
\pgfpathlineto{\pgfqpoint{7.027275in}{4.280175in}}%
\pgfpathlineto{\pgfqpoint{7.027361in}{4.280161in}}%
\pgfpathlineto{\pgfqpoint{7.028045in}{4.278657in}}%
\pgfpathlineto{\pgfqpoint{7.029414in}{4.268412in}}%
\pgfpathlineto{\pgfqpoint{7.031810in}{4.230047in}}%
\pgfpathlineto{\pgfqpoint{7.042077in}{4.040403in}}%
\pgfpathlineto{\pgfqpoint{7.042591in}{4.040682in}}%
\pgfpathlineto{\pgfqpoint{7.042762in}{4.041086in}}%
\pgfpathlineto{\pgfqpoint{7.043960in}{4.048189in}}%
\pgfpathlineto{\pgfqpoint{7.046099in}{4.077925in}}%
\pgfpathlineto{\pgfqpoint{7.050463in}{4.179194in}}%
\pgfpathlineto{\pgfqpoint{7.055083in}{4.266602in}}%
\pgfpathlineto{\pgfqpoint{7.057051in}{4.276404in}}%
\pgfpathlineto{\pgfqpoint{7.057136in}{4.276367in}}%
\pgfpathlineto{\pgfqpoint{7.057907in}{4.274271in}}%
\pgfpathlineto{\pgfqpoint{7.059361in}{4.261945in}}%
\pgfpathlineto{\pgfqpoint{7.061928in}{4.217485in}}%
\pgfpathlineto{\pgfqpoint{7.070912in}{4.045419in}}%
\pgfpathlineto{\pgfqpoint{7.071682in}{4.044046in}}%
\pgfpathlineto{\pgfqpoint{7.072195in}{4.044901in}}%
\pgfpathlineto{\pgfqpoint{7.073393in}{4.052316in}}%
\pgfpathlineto{\pgfqpoint{7.075532in}{4.082668in}}%
\pgfpathlineto{\pgfqpoint{7.080153in}{4.189982in}}%
\pgfpathlineto{\pgfqpoint{7.084345in}{4.264639in}}%
\pgfpathlineto{\pgfqpoint{7.086142in}{4.272767in}}%
\pgfpathlineto{\pgfqpoint{7.086313in}{4.272639in}}%
\pgfpathlineto{\pgfqpoint{7.087169in}{4.269636in}}%
\pgfpathlineto{\pgfqpoint{7.088795in}{4.253540in}}%
\pgfpathlineto{\pgfqpoint{7.091704in}{4.197663in}}%
\pgfpathlineto{\pgfqpoint{7.098977in}{4.053058in}}%
\pgfpathlineto{\pgfqpoint{7.100431in}{4.047612in}}%
\pgfpathlineto{\pgfqpoint{7.100773in}{4.047993in}}%
\pgfpathlineto{\pgfqpoint{7.101800in}{4.052920in}}%
\pgfpathlineto{\pgfqpoint{7.103682in}{4.075707in}}%
\pgfpathlineto{\pgfqpoint{7.107276in}{4.153363in}}%
\pgfpathlineto{\pgfqpoint{7.112667in}{4.259639in}}%
\pgfpathlineto{\pgfqpoint{7.114549in}{4.269266in}}%
\pgfpathlineto{\pgfqpoint{7.114720in}{4.269193in}}%
\pgfpathlineto{\pgfqpoint{7.115490in}{4.266892in}}%
\pgfpathlineto{\pgfqpoint{7.117030in}{4.252955in}}%
\pgfpathlineto{\pgfqpoint{7.119768in}{4.203012in}}%
\pgfpathlineto{\pgfqpoint{7.127383in}{4.054912in}}%
\pgfpathlineto{\pgfqpoint{7.128581in}{4.051057in}}%
\pgfpathlineto{\pgfqpoint{7.129009in}{4.051583in}}%
\pgfpathlineto{\pgfqpoint{7.130121in}{4.057591in}}%
\pgfpathlineto{\pgfqpoint{7.132089in}{4.083311in}}%
\pgfpathlineto{\pgfqpoint{7.136111in}{4.172959in}}%
\pgfpathlineto{\pgfqpoint{7.140731in}{4.257870in}}%
\pgfpathlineto{\pgfqpoint{7.142442in}{4.265877in}}%
\pgfpathlineto{\pgfqpoint{7.142699in}{4.265699in}}%
\pgfpathlineto{\pgfqpoint{7.143555in}{4.262498in}}%
\pgfpathlineto{\pgfqpoint{7.145180in}{4.245893in}}%
\pgfpathlineto{\pgfqpoint{7.148175in}{4.187459in}}%
\pgfpathlineto{\pgfqpoint{7.154678in}{4.060597in}}%
\pgfpathlineto{\pgfqpoint{7.156218in}{4.054393in}}%
\pgfpathlineto{\pgfqpoint{7.156475in}{4.054626in}}%
\pgfpathlineto{\pgfqpoint{7.157416in}{4.058579in}}%
\pgfpathlineto{\pgfqpoint{7.159213in}{4.078823in}}%
\pgfpathlineto{\pgfqpoint{7.162549in}{4.147979in}}%
\pgfpathlineto{\pgfqpoint{7.168111in}{4.254904in}}%
\pgfpathlineto{\pgfqpoint{7.169822in}{4.262597in}}%
\pgfpathlineto{\pgfqpoint{7.169993in}{4.262480in}}%
\pgfpathlineto{\pgfqpoint{7.170849in}{4.259467in}}%
\pgfpathlineto{\pgfqpoint{7.172475in}{4.243125in}}%
\pgfpathlineto{\pgfqpoint{7.175469in}{4.185084in}}%
\pgfpathlineto{\pgfqpoint{7.181801in}{4.063478in}}%
\pgfpathlineto{\pgfqpoint{7.183256in}{4.057611in}}%
\pgfpathlineto{\pgfqpoint{7.183598in}{4.057934in}}%
\pgfpathlineto{\pgfqpoint{7.184539in}{4.062159in}}%
\pgfpathlineto{\pgfqpoint{7.186336in}{4.082924in}}%
\pgfpathlineto{\pgfqpoint{7.189844in}{4.156374in}}%
\pgfpathlineto{\pgfqpoint{7.194978in}{4.252118in}}%
\pgfpathlineto{\pgfqpoint{7.196603in}{4.259431in}}%
\pgfpathlineto{\pgfqpoint{7.196860in}{4.259247in}}%
\pgfpathlineto{\pgfqpoint{7.197716in}{4.255988in}}%
\pgfpathlineto{\pgfqpoint{7.199427in}{4.237932in}}%
\pgfpathlineto{\pgfqpoint{7.202593in}{4.174783in}}%
\pgfpathlineto{\pgfqpoint{7.208240in}{4.067707in}}%
\pgfpathlineto{\pgfqpoint{7.209865in}{4.060733in}}%
\pgfpathlineto{\pgfqpoint{7.210122in}{4.060976in}}%
\pgfpathlineto{\pgfqpoint{7.211063in}{4.065002in}}%
\pgfpathlineto{\pgfqpoint{7.212860in}{4.085476in}}%
\pgfpathlineto{\pgfqpoint{7.216368in}{4.158294in}}%
\pgfpathlineto{\pgfqpoint{7.221331in}{4.249282in}}%
\pgfpathlineto{\pgfqpoint{7.222957in}{4.256361in}}%
\pgfpathlineto{\pgfqpoint{7.223213in}{4.256132in}}%
\pgfpathlineto{\pgfqpoint{7.224154in}{4.252148in}}%
\pgfpathlineto{\pgfqpoint{7.225951in}{4.231726in}}%
\pgfpathlineto{\pgfqpoint{7.229459in}{4.159156in}}%
\pgfpathlineto{\pgfqpoint{7.234336in}{4.070648in}}%
\pgfpathlineto{\pgfqpoint{7.235962in}{4.063755in}}%
\pgfpathlineto{\pgfqpoint{7.236219in}{4.064017in}}%
\pgfpathlineto{\pgfqpoint{7.237160in}{4.068128in}}%
\pgfpathlineto{\pgfqpoint{7.238957in}{4.088777in}}%
\pgfpathlineto{\pgfqpoint{7.242550in}{4.163306in}}%
\pgfpathlineto{\pgfqpoint{7.247256in}{4.246927in}}%
\pgfpathlineto{\pgfqpoint{7.248796in}{4.253393in}}%
\pgfpathlineto{\pgfqpoint{7.249053in}{4.253178in}}%
\pgfpathlineto{\pgfqpoint{7.249994in}{4.249233in}}%
\pgfpathlineto{\pgfqpoint{7.251791in}{4.228850in}}%
\pgfpathlineto{\pgfqpoint{7.255299in}{4.156760in}}%
\pgfpathlineto{\pgfqpoint{7.260005in}{4.073201in}}%
\pgfpathlineto{\pgfqpoint{7.261545in}{4.066671in}}%
\pgfpathlineto{\pgfqpoint{7.261802in}{4.066877in}}%
\pgfpathlineto{\pgfqpoint{7.262743in}{4.070797in}}%
\pgfpathlineto{\pgfqpoint{7.264540in}{4.091145in}}%
\pgfpathlineto{\pgfqpoint{7.268134in}{4.164913in}}%
\pgfpathlineto{\pgfqpoint{7.272754in}{4.244826in}}%
\pgfpathlineto{\pgfqpoint{7.274208in}{4.250514in}}%
\pgfpathlineto{\pgfqpoint{7.274551in}{4.250120in}}%
\pgfpathlineto{\pgfqpoint{7.275577in}{4.245004in}}%
\pgfpathlineto{\pgfqpoint{7.277460in}{4.221604in}}%
\pgfpathlineto{\pgfqpoint{7.281738in}{4.131250in}}%
\pgfpathlineto{\pgfqpoint{7.285503in}{4.073645in}}%
\pgfpathlineto{\pgfqpoint{7.286701in}{4.069504in}}%
\pgfpathlineto{\pgfqpoint{7.287128in}{4.069990in}}%
\pgfpathlineto{\pgfqpoint{7.288155in}{4.075341in}}%
\pgfpathlineto{\pgfqpoint{7.290123in}{4.100547in}}%
\pgfpathlineto{\pgfqpoint{7.295342in}{4.209603in}}%
\pgfpathlineto{\pgfqpoint{7.298423in}{4.246240in}}%
\pgfpathlineto{\pgfqpoint{7.299193in}{4.247721in}}%
\pgfpathlineto{\pgfqpoint{7.299706in}{4.246837in}}%
\pgfpathlineto{\pgfqpoint{7.300904in}{4.239082in}}%
\pgfpathlineto{\pgfqpoint{7.303129in}{4.206332in}}%
\pgfpathlineto{\pgfqpoint{7.311171in}{4.072513in}}%
\pgfpathlineto{\pgfqpoint{7.311599in}{4.072293in}}%
\pgfpathlineto{\pgfqpoint{7.312027in}{4.073114in}}%
\pgfpathlineto{\pgfqpoint{7.313225in}{4.080820in}}%
\pgfpathlineto{\pgfqpoint{7.315450in}{4.113466in}}%
\pgfpathlineto{\pgfqpoint{7.323407in}{4.244771in}}%
\pgfpathlineto{\pgfqpoint{7.323835in}{4.244975in}}%
\pgfpathlineto{\pgfqpoint{7.324262in}{4.244137in}}%
\pgfpathlineto{\pgfqpoint{7.325460in}{4.236383in}}%
\pgfpathlineto{\pgfqpoint{7.327685in}{4.203710in}}%
\pgfpathlineto{\pgfqpoint{7.335557in}{4.075133in}}%
\pgfpathlineto{\pgfqpoint{7.335985in}{4.074977in}}%
\pgfpathlineto{\pgfqpoint{7.336327in}{4.075603in}}%
\pgfpathlineto{\pgfqpoint{7.337439in}{4.082149in}}%
\pgfpathlineto{\pgfqpoint{7.339578in}{4.111806in}}%
\pgfpathlineto{\pgfqpoint{7.347792in}{4.242383in}}%
\pgfpathlineto{\pgfqpoint{7.348391in}{4.241621in}}%
\pgfpathlineto{\pgfqpoint{7.349589in}{4.234096in}}%
\pgfpathlineto{\pgfqpoint{7.351814in}{4.201837in}}%
\pgfpathlineto{\pgfqpoint{7.359514in}{4.077743in}}%
\pgfpathlineto{\pgfqpoint{7.359942in}{4.077548in}}%
\pgfpathlineto{\pgfqpoint{7.360370in}{4.078397in}}%
\pgfpathlineto{\pgfqpoint{7.361568in}{4.086176in}}%
\pgfpathlineto{\pgfqpoint{7.363792in}{4.118739in}}%
\pgfpathlineto{\pgfqpoint{7.371236in}{4.239355in}}%
\pgfpathlineto{\pgfqpoint{7.371750in}{4.239828in}}%
\pgfpathlineto{\pgfqpoint{7.372177in}{4.239075in}}%
\pgfpathlineto{\pgfqpoint{7.373375in}{4.231554in}}%
\pgfpathlineto{\pgfqpoint{7.375600in}{4.199410in}}%
\pgfpathlineto{\pgfqpoint{7.383044in}{4.080389in}}%
\pgfpathlineto{\pgfqpoint{7.383472in}{4.080015in}}%
\pgfpathlineto{\pgfqpoint{7.383985in}{4.080943in}}%
\pgfpathlineto{\pgfqpoint{7.385183in}{4.088792in}}%
\pgfpathlineto{\pgfqpoint{7.387408in}{4.121329in}}%
\pgfpathlineto{\pgfqpoint{7.394680in}{4.236931in}}%
\pgfpathlineto{\pgfqpoint{7.395194in}{4.237345in}}%
\pgfpathlineto{\pgfqpoint{7.395622in}{4.236544in}}%
\pgfpathlineto{\pgfqpoint{7.396819in}{4.228902in}}%
\pgfpathlineto{\pgfqpoint{7.399044in}{4.196718in}}%
\pgfpathlineto{\pgfqpoint{7.406231in}{4.082961in}}%
\pgfpathlineto{\pgfqpoint{7.406745in}{4.082461in}}%
\pgfpathlineto{\pgfqpoint{7.407173in}{4.083190in}}%
\pgfpathlineto{\pgfqpoint{7.408370in}{4.090632in}}%
\pgfpathlineto{\pgfqpoint{7.410595in}{4.122470in}}%
\pgfpathlineto{\pgfqpoint{7.417782in}{4.234590in}}%
\pgfpathlineto{\pgfqpoint{7.418210in}{4.234975in}}%
\pgfpathlineto{\pgfqpoint{7.418723in}{4.234063in}}%
\pgfpathlineto{\pgfqpoint{7.419921in}{4.226271in}}%
\pgfpathlineto{\pgfqpoint{7.422231in}{4.192443in}}%
\pgfpathlineto{\pgfqpoint{7.428991in}{4.085743in}}%
\pgfpathlineto{\pgfqpoint{7.429590in}{4.084806in}}%
\pgfpathlineto{\pgfqpoint{7.430103in}{4.085624in}}%
\pgfpathlineto{\pgfqpoint{7.431301in}{4.093201in}}%
\pgfpathlineto{\pgfqpoint{7.433526in}{4.125071in}}%
\pgfpathlineto{\pgfqpoint{7.440456in}{4.232106in}}%
\pgfpathlineto{\pgfqpoint{7.440970in}{4.232637in}}%
\pgfpathlineto{\pgfqpoint{7.441397in}{4.231938in}}%
\pgfpathlineto{\pgfqpoint{7.442510in}{4.225378in}}%
\pgfpathlineto{\pgfqpoint{7.444649in}{4.196194in}}%
\pgfpathlineto{\pgfqpoint{7.451836in}{4.087371in}}%
\pgfpathlineto{\pgfqpoint{7.452264in}{4.087142in}}%
\pgfpathlineto{\pgfqpoint{7.452692in}{4.087949in}}%
\pgfpathlineto{\pgfqpoint{7.453890in}{4.095563in}}%
\pgfpathlineto{\pgfqpoint{7.456114in}{4.127313in}}%
\pgfpathlineto{\pgfqpoint{7.462788in}{4.229605in}}%
\pgfpathlineto{\pgfqpoint{7.463387in}{4.230373in}}%
\pgfpathlineto{\pgfqpoint{7.463900in}{4.229419in}}%
\pgfpathlineto{\pgfqpoint{7.465098in}{4.221574in}}%
\pgfpathlineto{\pgfqpoint{7.467408in}{4.188051in}}%
\pgfpathlineto{\pgfqpoint{7.473826in}{4.090287in}}%
\pgfpathlineto{\pgfqpoint{7.474425in}{4.089335in}}%
\pgfpathlineto{\pgfqpoint{7.474938in}{4.090129in}}%
\pgfpathlineto{\pgfqpoint{7.476136in}{4.097606in}}%
\pgfpathlineto{\pgfqpoint{7.478360in}{4.128987in}}%
\pgfpathlineto{\pgfqpoint{7.484949in}{4.227553in}}%
\pgfpathlineto{\pgfqpoint{7.485462in}{4.228183in}}%
\pgfpathlineto{\pgfqpoint{7.485976in}{4.227330in}}%
\pgfpathlineto{\pgfqpoint{7.487173in}{4.219737in}}%
\pgfpathlineto{\pgfqpoint{7.489484in}{4.186766in}}%
\pgfpathlineto{\pgfqpoint{7.495730in}{4.092638in}}%
\pgfpathlineto{\pgfqpoint{7.496414in}{4.091510in}}%
\pgfpathlineto{\pgfqpoint{7.496928in}{4.092391in}}%
\pgfpathlineto{\pgfqpoint{7.498125in}{4.100030in}}%
\pgfpathlineto{\pgfqpoint{7.500436in}{4.132951in}}%
\pgfpathlineto{\pgfqpoint{7.506596in}{4.224910in}}%
\pgfpathlineto{\pgfqpoint{7.507281in}{4.226042in}}%
\pgfpathlineto{\pgfqpoint{7.507794in}{4.225169in}}%
\pgfpathlineto{\pgfqpoint{7.508992in}{4.217563in}}%
\pgfpathlineto{\pgfqpoint{7.511302in}{4.184803in}}%
\pgfpathlineto{\pgfqpoint{7.517377in}{4.094804in}}%
\pgfpathlineto{\pgfqpoint{7.518061in}{4.093616in}}%
\pgfpathlineto{\pgfqpoint{7.518575in}{4.094442in}}%
\pgfpathlineto{\pgfqpoint{7.519773in}{4.101929in}}%
\pgfpathlineto{\pgfqpoint{7.522083in}{4.134412in}}%
\pgfpathlineto{\pgfqpoint{7.528158in}{4.222973in}}%
\pgfpathlineto{\pgfqpoint{7.528842in}{4.223947in}}%
\pgfpathlineto{\pgfqpoint{7.529270in}{4.223230in}}%
\pgfpathlineto{\pgfqpoint{7.530468in}{4.215957in}}%
\pgfpathlineto{\pgfqpoint{7.532778in}{4.183878in}}%
\pgfpathlineto{\pgfqpoint{7.538768in}{4.096815in}}%
\pgfpathlineto{\pgfqpoint{7.539452in}{4.095669in}}%
\pgfpathlineto{\pgfqpoint{7.539965in}{4.096517in}}%
\pgfpathlineto{\pgfqpoint{7.541163in}{4.104013in}}%
\pgfpathlineto{\pgfqpoint{7.543473in}{4.136261in}}%
\pgfpathlineto{\pgfqpoint{7.549377in}{4.220876in}}%
\pgfpathlineto{\pgfqpoint{7.550062in}{4.221935in}}%
\pgfpathlineto{\pgfqpoint{7.550575in}{4.221027in}}%
\pgfpathlineto{\pgfqpoint{7.551773in}{4.213418in}}%
\pgfpathlineto{\pgfqpoint{7.554083in}{4.181155in}}%
\pgfpathlineto{\pgfqpoint{7.559901in}{4.098712in}}%
\pgfpathlineto{\pgfqpoint{7.560586in}{4.097671in}}%
\pgfpathlineto{\pgfqpoint{7.561099in}{4.098587in}}%
\pgfpathlineto{\pgfqpoint{7.562297in}{4.106192in}}%
\pgfpathlineto{\pgfqpoint{7.564693in}{4.139804in}}%
\pgfpathlineto{\pgfqpoint{7.570255in}{4.218558in}}%
\pgfpathlineto{\pgfqpoint{7.571025in}{4.219971in}}%
\pgfpathlineto{\pgfqpoint{7.571538in}{4.219102in}}%
\pgfpathlineto{\pgfqpoint{7.572736in}{4.211623in}}%
\pgfpathlineto{\pgfqpoint{7.575046in}{4.179798in}}%
\pgfpathlineto{\pgfqpoint{7.580693in}{4.100828in}}%
\pgfpathlineto{\pgfqpoint{7.581378in}{4.099600in}}%
\pgfpathlineto{\pgfqpoint{7.581891in}{4.100364in}}%
\pgfpathlineto{\pgfqpoint{7.583089in}{4.107591in}}%
\pgfpathlineto{\pgfqpoint{7.585399in}{4.138953in}}%
\pgfpathlineto{\pgfqpoint{7.591046in}{4.216940in}}%
\pgfpathlineto{\pgfqpoint{7.591731in}{4.218059in}}%
\pgfpathlineto{\pgfqpoint{7.592244in}{4.217221in}}%
\pgfpathlineto{\pgfqpoint{7.593442in}{4.209854in}}%
\pgfpathlineto{\pgfqpoint{7.595752in}{4.178456in}}%
\pgfpathlineto{\pgfqpoint{7.601314in}{4.102591in}}%
\pgfpathlineto{\pgfqpoint{7.601998in}{4.101493in}}%
\pgfpathlineto{\pgfqpoint{7.602512in}{4.102342in}}%
\pgfpathlineto{\pgfqpoint{7.603709in}{4.109708in}}%
\pgfpathlineto{\pgfqpoint{7.606105in}{4.142400in}}%
\pgfpathlineto{\pgfqpoint{7.611410in}{4.214716in}}%
\pgfpathlineto{\pgfqpoint{7.612180in}{4.216200in}}%
\pgfpathlineto{\pgfqpoint{7.612694in}{4.215407in}}%
\pgfpathlineto{\pgfqpoint{7.613891in}{4.208188in}}%
\pgfpathlineto{\pgfqpoint{7.616202in}{4.177271in}}%
\pgfpathlineto{\pgfqpoint{7.621678in}{4.104340in}}%
\pgfpathlineto{\pgfqpoint{7.622362in}{4.103338in}}%
\pgfpathlineto{\pgfqpoint{7.622875in}{4.104243in}}%
\pgfpathlineto{\pgfqpoint{7.624073in}{4.111682in}}%
\pgfpathlineto{\pgfqpoint{7.626469in}{4.144143in}}%
\pgfpathlineto{\pgfqpoint{7.631688in}{4.213139in}}%
\pgfpathlineto{\pgfqpoint{7.632458in}{4.214370in}}%
\pgfpathlineto{\pgfqpoint{7.632886in}{4.213678in}}%
\pgfpathlineto{\pgfqpoint{7.634084in}{4.206690in}}%
\pgfpathlineto{\pgfqpoint{7.636394in}{4.176367in}}%
\pgfpathlineto{\pgfqpoint{7.641785in}{4.106099in}}%
\pgfpathlineto{\pgfqpoint{7.642469in}{4.105128in}}%
\pgfpathlineto{\pgfqpoint{7.642897in}{4.105792in}}%
\pgfpathlineto{\pgfqpoint{7.644095in}{4.112685in}}%
\pgfpathlineto{\pgfqpoint{7.646405in}{4.142733in}}%
\pgfpathlineto{\pgfqpoint{7.651710in}{4.211518in}}%
\pgfpathlineto{\pgfqpoint{7.652395in}{4.212622in}}%
\pgfpathlineto{\pgfqpoint{7.652908in}{4.211815in}}%
\pgfpathlineto{\pgfqpoint{7.654106in}{4.204674in}}%
\pgfpathlineto{\pgfqpoint{7.656501in}{4.173093in}}%
\pgfpathlineto{\pgfqpoint{7.661550in}{4.108192in}}%
\pgfpathlineto{\pgfqpoint{7.662320in}{4.106861in}}%
\pgfpathlineto{\pgfqpoint{7.662833in}{4.107716in}}%
\pgfpathlineto{\pgfqpoint{7.664031in}{4.114934in}}%
\pgfpathlineto{\pgfqpoint{7.666427in}{4.146433in}}%
\pgfpathlineto{\pgfqpoint{7.671389in}{4.209539in}}%
\pgfpathlineto{\pgfqpoint{7.672159in}{4.210903in}}%
\pgfpathlineto{\pgfqpoint{7.672673in}{4.210080in}}%
\pgfpathlineto{\pgfqpoint{7.673871in}{4.202960in}}%
\pgfpathlineto{\pgfqpoint{7.676266in}{4.171758in}}%
\pgfpathlineto{\pgfqpoint{7.681229in}{4.109744in}}%
\pgfpathlineto{\pgfqpoint{7.681999in}{4.108569in}}%
\pgfpathlineto{\pgfqpoint{7.682427in}{4.109257in}}%
\pgfpathlineto{\pgfqpoint{7.683625in}{4.116103in}}%
\pgfpathlineto{\pgfqpoint{7.686021in}{4.146789in}}%
\pgfpathlineto{\pgfqpoint{7.690983in}{4.208096in}}%
\pgfpathlineto{\pgfqpoint{7.691668in}{4.209237in}}%
\pgfpathlineto{\pgfqpoint{7.692181in}{4.208491in}}%
\pgfpathlineto{\pgfqpoint{7.693379in}{4.201602in}}%
\pgfpathlineto{\pgfqpoint{7.695775in}{4.171042in}}%
\pgfpathlineto{\pgfqpoint{7.700652in}{4.111391in}}%
\pgfpathlineto{\pgfqpoint{7.701422in}{4.110216in}}%
\pgfpathlineto{\pgfqpoint{7.701850in}{4.110890in}}%
\pgfpathlineto{\pgfqpoint{7.703047in}{4.117638in}}%
\pgfpathlineto{\pgfqpoint{7.705443in}{4.147843in}}%
\pgfpathlineto{\pgfqpoint{7.710320in}{4.206542in}}%
\pgfpathlineto{\pgfqpoint{7.711005in}{4.207607in}}%
\pgfpathlineto{\pgfqpoint{7.711518in}{4.206823in}}%
\pgfpathlineto{\pgfqpoint{7.712716in}{4.199913in}}%
\pgfpathlineto{\pgfqpoint{7.715112in}{4.169699in}}%
\pgfpathlineto{\pgfqpoint{7.719903in}{4.112857in}}%
\pgfpathlineto{\pgfqpoint{7.720588in}{4.111805in}}%
\pgfpathlineto{\pgfqpoint{7.721101in}{4.112590in}}%
\pgfpathlineto{\pgfqpoint{7.722299in}{4.119468in}}%
\pgfpathlineto{\pgfqpoint{7.724695in}{4.149456in}}%
\pgfpathlineto{\pgfqpoint{7.729401in}{4.204869in}}%
\pgfpathlineto{\pgfqpoint{7.730171in}{4.206004in}}%
\pgfpathlineto{\pgfqpoint{7.730599in}{4.205331in}}%
\pgfpathlineto{\pgfqpoint{7.731796in}{4.198677in}}%
\pgfpathlineto{\pgfqpoint{7.734192in}{4.169148in}}%
\pgfpathlineto{\pgfqpoint{7.738898in}{4.114452in}}%
\pgfpathlineto{\pgfqpoint{7.739583in}{4.113365in}}%
\pgfpathlineto{\pgfqpoint{7.740096in}{4.114105in}}%
\pgfpathlineto{\pgfqpoint{7.741294in}{4.120819in}}%
\pgfpathlineto{\pgfqpoint{7.743690in}{4.150228in}}%
\pgfpathlineto{\pgfqpoint{7.748310in}{4.203346in}}%
\pgfpathlineto{\pgfqpoint{7.749080in}{4.204461in}}%
\pgfpathlineto{\pgfqpoint{7.749508in}{4.203794in}}%
\pgfpathlineto{\pgfqpoint{7.750706in}{4.197219in}}%
\pgfpathlineto{\pgfqpoint{7.753102in}{4.168170in}}%
\pgfpathlineto{\pgfqpoint{7.757722in}{4.115912in}}%
\pgfpathlineto{\pgfqpoint{7.758406in}{4.114890in}}%
\pgfpathlineto{\pgfqpoint{7.758920in}{4.115658in}}%
\pgfpathlineto{\pgfqpoint{7.760118in}{4.122365in}}%
\pgfpathlineto{\pgfqpoint{7.762599in}{4.152625in}}%
\pgfpathlineto{\pgfqpoint{7.767048in}{4.201947in}}%
\pgfpathlineto{\pgfqpoint{7.767733in}{4.202979in}}%
\pgfpathlineto{\pgfqpoint{7.768246in}{4.202229in}}%
\pgfpathlineto{\pgfqpoint{7.769444in}{4.195596in}}%
\pgfpathlineto{\pgfqpoint{7.771925in}{4.165634in}}%
\pgfpathlineto{\pgfqpoint{7.776289in}{4.117536in}}%
\pgfpathlineto{\pgfqpoint{7.777059in}{4.116384in}}%
\pgfpathlineto{\pgfqpoint{7.777487in}{4.117005in}}%
\pgfpathlineto{\pgfqpoint{7.778685in}{4.123359in}}%
\pgfpathlineto{\pgfqpoint{7.781080in}{4.151560in}}%
\pgfpathlineto{\pgfqpoint{7.785615in}{4.200642in}}%
\pgfpathlineto{\pgfqpoint{7.786300in}{4.201502in}}%
\pgfpathlineto{\pgfqpoint{7.786728in}{4.200876in}}%
\pgfpathlineto{\pgfqpoint{7.787925in}{4.194544in}}%
\pgfpathlineto{\pgfqpoint{7.790321in}{4.166568in}}%
\pgfpathlineto{\pgfqpoint{7.794770in}{4.118785in}}%
\pgfpathlineto{\pgfqpoint{7.795455in}{4.117817in}}%
\pgfpathlineto{\pgfqpoint{7.795968in}{4.118584in}}%
\pgfpathlineto{\pgfqpoint{7.797166in}{4.125143in}}%
\pgfpathlineto{\pgfqpoint{7.799647in}{4.154369in}}%
\pgfpathlineto{\pgfqpoint{7.803926in}{4.199163in}}%
\pgfpathlineto{\pgfqpoint{7.804610in}{4.200087in}}%
\pgfpathlineto{\pgfqpoint{7.805123in}{4.199297in}}%
\pgfpathlineto{\pgfqpoint{7.806407in}{4.192025in}}%
\pgfpathlineto{\pgfqpoint{7.809059in}{4.160133in}}%
\pgfpathlineto{\pgfqpoint{7.812995in}{4.120230in}}%
\pgfpathlineto{\pgfqpoint{7.813680in}{4.119219in}}%
\pgfpathlineto{\pgfqpoint{7.814193in}{4.119932in}}%
\pgfpathlineto{\pgfqpoint{7.815391in}{4.126299in}}%
\pgfpathlineto{\pgfqpoint{7.817872in}{4.154870in}}%
\pgfpathlineto{\pgfqpoint{7.822065in}{4.197751in}}%
\pgfpathlineto{\pgfqpoint{7.822749in}{4.198703in}}%
\pgfpathlineto{\pgfqpoint{7.823263in}{4.197957in}}%
\pgfpathlineto{\pgfqpoint{7.824461in}{4.191558in}}%
\pgfpathlineto{\pgfqpoint{7.827027in}{4.162013in}}%
\pgfpathlineto{\pgfqpoint{7.831049in}{4.121614in}}%
\pgfpathlineto{\pgfqpoint{7.831819in}{4.120613in}}%
\pgfpathlineto{\pgfqpoint{7.832247in}{4.121265in}}%
\pgfpathlineto{\pgfqpoint{7.833445in}{4.127480in}}%
\pgfpathlineto{\pgfqpoint{7.835926in}{4.155436in}}%
\pgfpathlineto{\pgfqpoint{7.840033in}{4.196392in}}%
\pgfpathlineto{\pgfqpoint{7.840717in}{4.197354in}}%
\pgfpathlineto{\pgfqpoint{7.841231in}{4.196638in}}%
\pgfpathlineto{\pgfqpoint{7.842429in}{4.190382in}}%
\pgfpathlineto{\pgfqpoint{7.844996in}{4.161474in}}%
\pgfpathlineto{\pgfqpoint{7.849017in}{4.122705in}}%
\pgfpathlineto{\pgfqpoint{7.849702in}{4.121942in}}%
\pgfpathlineto{\pgfqpoint{7.850129in}{4.122570in}}%
\pgfpathlineto{\pgfqpoint{7.851327in}{4.128648in}}%
\pgfpathlineto{\pgfqpoint{7.853809in}{4.156003in}}%
\pgfpathlineto{\pgfqpoint{7.857916in}{4.195309in}}%
\pgfpathlineto{\pgfqpoint{7.858600in}{4.196010in}}%
\pgfpathlineto{\pgfqpoint{7.859028in}{4.195353in}}%
\pgfpathlineto{\pgfqpoint{7.860226in}{4.189237in}}%
\pgfpathlineto{\pgfqpoint{7.862793in}{4.160965in}}%
\pgfpathlineto{\pgfqpoint{7.866728in}{4.124013in}}%
\pgfpathlineto{\pgfqpoint{7.867413in}{4.123236in}}%
\pgfpathlineto{\pgfqpoint{7.867841in}{4.123837in}}%
\pgfpathlineto{\pgfqpoint{7.869039in}{4.129769in}}%
\pgfpathlineto{\pgfqpoint{7.871520in}{4.156510in}}%
\pgfpathlineto{\pgfqpoint{7.875541in}{4.194012in}}%
\pgfpathlineto{\pgfqpoint{7.876226in}{4.194735in}}%
\pgfpathlineto{\pgfqpoint{7.876654in}{4.194110in}}%
\pgfpathlineto{\pgfqpoint{7.877852in}{4.188155in}}%
\pgfpathlineto{\pgfqpoint{7.880418in}{4.160541in}}%
\pgfpathlineto{\pgfqpoint{7.884269in}{4.125302in}}%
\pgfpathlineto{\pgfqpoint{7.884953in}{4.124494in}}%
\pgfpathlineto{\pgfqpoint{7.885467in}{4.125268in}}%
\pgfpathlineto{\pgfqpoint{7.886750in}{4.132098in}}%
\pgfpathlineto{\pgfqpoint{7.889574in}{4.163283in}}%
\pgfpathlineto{\pgfqpoint{7.893082in}{4.192938in}}%
\pgfpathlineto{\pgfqpoint{7.893681in}{4.193497in}}%
\pgfpathlineto{\pgfqpoint{7.894194in}{4.192706in}}%
\pgfpathlineto{\pgfqpoint{7.895477in}{4.185884in}}%
\pgfpathlineto{\pgfqpoint{7.898301in}{4.155012in}}%
\pgfpathlineto{\pgfqpoint{7.901724in}{4.126364in}}%
\pgfpathlineto{\pgfqpoint{7.902408in}{4.125755in}}%
\pgfpathlineto{\pgfqpoint{7.902836in}{4.126422in}}%
\pgfpathlineto{\pgfqpoint{7.904119in}{4.133003in}}%
\pgfpathlineto{\pgfqpoint{7.906943in}{4.163359in}}%
\pgfpathlineto{\pgfqpoint{7.910365in}{4.191662in}}%
\pgfpathlineto{\pgfqpoint{7.911050in}{4.192254in}}%
\pgfpathlineto{\pgfqpoint{7.911478in}{4.191585in}}%
\pgfpathlineto{\pgfqpoint{7.912761in}{4.185046in}}%
\pgfpathlineto{\pgfqpoint{7.915585in}{4.155032in}}%
\pgfpathlineto{\pgfqpoint{7.919007in}{4.127445in}}%
\pgfpathlineto{\pgfqpoint{7.919606in}{4.126932in}}%
\pgfpathlineto{\pgfqpoint{7.920034in}{4.127517in}}%
\pgfpathlineto{\pgfqpoint{7.921232in}{4.133168in}}%
\pgfpathlineto{\pgfqpoint{7.923799in}{4.159230in}}%
\pgfpathlineto{\pgfqpoint{7.927563in}{4.190577in}}%
\pgfpathlineto{\pgfqpoint{7.928162in}{4.191097in}}%
\pgfpathlineto{\pgfqpoint{7.928676in}{4.190318in}}%
\pgfpathlineto{\pgfqpoint{7.929959in}{4.183717in}}%
\pgfpathlineto{\pgfqpoint{7.932868in}{4.153292in}}%
\pgfpathlineto{\pgfqpoint{7.936120in}{4.128543in}}%
\pgfpathlineto{\pgfqpoint{7.936719in}{4.128113in}}%
\pgfpathlineto{\pgfqpoint{7.937146in}{4.128740in}}%
\pgfpathlineto{\pgfqpoint{7.938430in}{4.135022in}}%
\pgfpathlineto{\pgfqpoint{7.941253in}{4.163898in}}%
\pgfpathlineto{\pgfqpoint{7.944590in}{4.189475in}}%
\pgfpathlineto{\pgfqpoint{7.945189in}{4.189937in}}%
\pgfpathlineto{\pgfqpoint{7.945617in}{4.189343in}}%
\pgfpathlineto{\pgfqpoint{7.946815in}{4.183792in}}%
\pgfpathlineto{\pgfqpoint{7.949467in}{4.157618in}}%
\pgfpathlineto{\pgfqpoint{7.953061in}{4.129656in}}%
\pgfpathlineto{\pgfqpoint{7.953660in}{4.129253in}}%
\pgfpathlineto{\pgfqpoint{7.954088in}{4.129881in}}%
\pgfpathlineto{\pgfqpoint{7.955371in}{4.136072in}}%
\pgfpathlineto{\pgfqpoint{7.958280in}{4.165190in}}%
\pgfpathlineto{\pgfqpoint{7.961532in}{4.188511in}}%
\pgfpathlineto{\pgfqpoint{7.962131in}{4.188764in}}%
\pgfpathlineto{\pgfqpoint{7.962473in}{4.188244in}}%
\pgfpathlineto{\pgfqpoint{7.963671in}{4.182827in}}%
\pgfpathlineto{\pgfqpoint{7.966323in}{4.157319in}}%
\pgfpathlineto{\pgfqpoint{7.969831in}{4.130787in}}%
\pgfpathlineto{\pgfqpoint{7.970430in}{4.130350in}}%
\pgfpathlineto{\pgfqpoint{7.970858in}{4.130934in}}%
\pgfpathlineto{\pgfqpoint{7.972056in}{4.136334in}}%
\pgfpathlineto{\pgfqpoint{7.974708in}{4.161558in}}%
\pgfpathlineto{\pgfqpoint{7.978216in}{4.187373in}}%
\pgfpathlineto{\pgfqpoint{7.978815in}{4.187697in}}%
\pgfpathlineto{\pgfqpoint{7.979243in}{4.187042in}}%
\pgfpathlineto{\pgfqpoint{7.980527in}{4.180919in}}%
\pgfpathlineto{\pgfqpoint{7.983521in}{4.151956in}}%
\pgfpathlineto{\pgfqpoint{7.986601in}{4.131647in}}%
\pgfpathlineto{\pgfqpoint{7.987115in}{4.131448in}}%
\pgfpathlineto{\pgfqpoint{7.987543in}{4.132087in}}%
\pgfpathlineto{\pgfqpoint{7.988826in}{4.138118in}}%
\pgfpathlineto{\pgfqpoint{7.991906in}{4.167529in}}%
\pgfpathlineto{\pgfqpoint{7.994901in}{4.186482in}}%
\pgfpathlineto{\pgfqpoint{7.995414in}{4.186618in}}%
\pgfpathlineto{\pgfqpoint{7.995842in}{4.185937in}}%
\pgfpathlineto{\pgfqpoint{7.997126in}{4.179842in}}%
\pgfpathlineto{\pgfqpoint{8.000291in}{4.149920in}}%
\pgfpathlineto{\pgfqpoint{8.003201in}{4.132578in}}%
\pgfpathlineto{\pgfqpoint{8.003800in}{4.132670in}}%
\pgfpathlineto{\pgfqpoint{8.004056in}{4.133138in}}%
\pgfpathlineto{\pgfqpoint{8.005340in}{4.139068in}}%
\pgfpathlineto{\pgfqpoint{8.008420in}{4.167687in}}%
\pgfpathlineto{\pgfqpoint{8.011329in}{4.185440in}}%
\pgfpathlineto{\pgfqpoint{8.011928in}{4.185519in}}%
\pgfpathlineto{\pgfqpoint{8.012270in}{4.184942in}}%
\pgfpathlineto{\pgfqpoint{8.013554in}{4.179026in}}%
\pgfpathlineto{\pgfqpoint{8.016719in}{4.149977in}}%
\pgfpathlineto{\pgfqpoint{8.019543in}{4.133624in}}%
\pgfpathlineto{\pgfqpoint{8.020142in}{4.133651in}}%
\pgfpathlineto{\pgfqpoint{8.020399in}{4.134082in}}%
\pgfpathlineto{\pgfqpoint{8.021682in}{4.139758in}}%
\pgfpathlineto{\pgfqpoint{8.024762in}{4.167487in}}%
\pgfpathlineto{\pgfqpoint{8.027671in}{4.184485in}}%
\pgfpathlineto{\pgfqpoint{8.028270in}{4.184472in}}%
\pgfpathlineto{\pgfqpoint{8.028613in}{4.183859in}}%
\pgfpathlineto{\pgfqpoint{8.029982in}{4.177335in}}%
\pgfpathlineto{\pgfqpoint{8.036142in}{4.134469in}}%
\pgfpathlineto{\pgfqpoint{8.036998in}{4.136063in}}%
\pgfpathlineto{\pgfqpoint{8.038623in}{4.145705in}}%
\pgfpathlineto{\pgfqpoint{8.044185in}{4.183665in}}%
\pgfpathlineto{\pgfqpoint{8.044356in}{4.183580in}}%
\pgfpathlineto{\pgfqpoint{8.045297in}{4.181218in}}%
\pgfpathlineto{\pgfqpoint{8.047180in}{4.168483in}}%
\pgfpathlineto{\pgfqpoint{8.051971in}{4.135501in}}%
\pgfpathlineto{\pgfqpoint{8.052656in}{4.135848in}}%
\pgfpathlineto{\pgfqpoint{8.052827in}{4.136200in}}%
\pgfpathlineto{\pgfqpoint{8.054196in}{4.142511in}}%
\pgfpathlineto{\pgfqpoint{8.060185in}{4.182716in}}%
\pgfpathlineto{\pgfqpoint{8.061041in}{4.181388in}}%
\pgfpathlineto{\pgfqpoint{8.062581in}{4.173062in}}%
\pgfpathlineto{\pgfqpoint{8.068228in}{4.136375in}}%
\pgfpathlineto{\pgfqpoint{8.068485in}{4.136561in}}%
\pgfpathlineto{\pgfqpoint{8.069511in}{4.139577in}}%
\pgfpathlineto{\pgfqpoint{8.071651in}{4.154886in}}%
\pgfpathlineto{\pgfqpoint{8.075758in}{4.181586in}}%
\pgfpathlineto{\pgfqpoint{8.076356in}{4.181670in}}%
\pgfpathlineto{\pgfqpoint{8.076699in}{4.181152in}}%
\pgfpathlineto{\pgfqpoint{8.077982in}{4.175813in}}%
\pgfpathlineto{\pgfqpoint{8.081576in}{4.146847in}}%
\pgfpathlineto{\pgfqpoint{8.083972in}{4.137282in}}%
\pgfpathlineto{\pgfqpoint{8.084742in}{4.138208in}}%
\pgfpathlineto{\pgfqpoint{8.086196in}{4.145085in}}%
\pgfpathlineto{\pgfqpoint{8.091929in}{4.180884in}}%
\pgfpathlineto{\pgfqpoint{8.092442in}{4.180365in}}%
\pgfpathlineto{\pgfqpoint{8.093726in}{4.175358in}}%
\pgfpathlineto{\pgfqpoint{8.097063in}{4.149274in}}%
\pgfpathlineto{\pgfqpoint{8.099629in}{4.138191in}}%
\pgfpathlineto{\pgfqpoint{8.100485in}{4.139099in}}%
\pgfpathlineto{\pgfqpoint{8.101940in}{4.145834in}}%
\pgfpathlineto{\pgfqpoint{8.107587in}{4.180009in}}%
\pgfpathlineto{\pgfqpoint{8.108100in}{4.179529in}}%
\pgfpathlineto{\pgfqpoint{8.109384in}{4.174714in}}%
\pgfpathlineto{\pgfqpoint{8.112721in}{4.149514in}}%
\pgfpathlineto{\pgfqpoint{8.115287in}{4.139042in}}%
\pgfpathlineto{\pgfqpoint{8.116143in}{4.140049in}}%
\pgfpathlineto{\pgfqpoint{8.117683in}{4.147325in}}%
\pgfpathlineto{\pgfqpoint{8.123159in}{4.179154in}}%
\pgfpathlineto{\pgfqpoint{8.123501in}{4.178919in}}%
\pgfpathlineto{\pgfqpoint{8.124614in}{4.175629in}}%
\pgfpathlineto{\pgfqpoint{8.127009in}{4.159282in}}%
\pgfpathlineto{\pgfqpoint{8.130518in}{4.140056in}}%
\pgfpathlineto{\pgfqpoint{8.131202in}{4.140071in}}%
\pgfpathlineto{\pgfqpoint{8.131459in}{4.140462in}}%
\pgfpathlineto{\pgfqpoint{8.132828in}{4.145802in}}%
\pgfpathlineto{\pgfqpoint{8.138646in}{4.178315in}}%
\pgfpathlineto{\pgfqpoint{8.139330in}{4.177391in}}%
\pgfpathlineto{\pgfqpoint{8.140871in}{4.170484in}}%
\pgfpathlineto{\pgfqpoint{8.146261in}{4.140694in}}%
\pgfpathlineto{\pgfqpoint{8.146603in}{4.140903in}}%
\pgfpathlineto{\pgfqpoint{8.147716in}{4.144006in}}%
\pgfpathlineto{\pgfqpoint{8.150111in}{4.159537in}}%
\pgfpathlineto{\pgfqpoint{8.153534in}{4.177331in}}%
\pgfpathlineto{\pgfqpoint{8.154218in}{4.177351in}}%
\pgfpathlineto{\pgfqpoint{8.154561in}{4.176824in}}%
\pgfpathlineto{\pgfqpoint{8.155930in}{4.171467in}}%
\pgfpathlineto{\pgfqpoint{8.161577in}{4.141494in}}%
\pgfpathlineto{\pgfqpoint{8.162176in}{4.142162in}}%
\pgfpathlineto{\pgfqpoint{8.163545in}{4.147415in}}%
\pgfpathlineto{\pgfqpoint{8.169192in}{4.176726in}}%
\pgfpathlineto{\pgfqpoint{8.169791in}{4.176032in}}%
\pgfpathlineto{\pgfqpoint{8.171245in}{4.170317in}}%
\pgfpathlineto{\pgfqpoint{8.176721in}{4.142263in}}%
\pgfpathlineto{\pgfqpoint{8.177235in}{4.142730in}}%
\pgfpathlineto{\pgfqpoint{8.178518in}{4.147049in}}%
\pgfpathlineto{\pgfqpoint{8.184251in}{4.175975in}}%
\pgfpathlineto{\pgfqpoint{8.185106in}{4.174821in}}%
\pgfpathlineto{\pgfqpoint{8.186732in}{4.167582in}}%
\pgfpathlineto{\pgfqpoint{8.191780in}{4.143014in}}%
\pgfpathlineto{\pgfqpoint{8.191866in}{4.143038in}}%
\pgfpathlineto{\pgfqpoint{8.192807in}{4.144663in}}%
\pgfpathlineto{\pgfqpoint{8.194689in}{4.154013in}}%
\pgfpathlineto{\pgfqpoint{8.199053in}{4.175205in}}%
\pgfpathlineto{\pgfqpoint{8.199994in}{4.174378in}}%
\pgfpathlineto{\pgfqpoint{8.201534in}{4.168237in}}%
\pgfpathlineto{\pgfqpoint{8.206754in}{4.143747in}}%
\pgfpathlineto{\pgfqpoint{8.207010in}{4.143898in}}%
\pgfpathlineto{\pgfqpoint{8.208123in}{4.146566in}}%
\pgfpathlineto{\pgfqpoint{8.210518in}{4.160029in}}%
\pgfpathlineto{\pgfqpoint{8.213855in}{4.174437in}}%
\pgfpathlineto{\pgfqpoint{8.214797in}{4.173874in}}%
\pgfpathlineto{\pgfqpoint{8.216251in}{4.168680in}}%
\pgfpathlineto{\pgfqpoint{8.221556in}{4.144446in}}%
\pgfpathlineto{\pgfqpoint{8.221984in}{4.144739in}}%
\pgfpathlineto{\pgfqpoint{8.223182in}{4.148010in}}%
\pgfpathlineto{\pgfqpoint{8.226262in}{4.165552in}}%
\pgfpathlineto{\pgfqpoint{8.228829in}{4.173819in}}%
\pgfpathlineto{\pgfqpoint{8.229770in}{4.172835in}}%
\pgfpathlineto{\pgfqpoint{8.231396in}{4.166377in}}%
\pgfpathlineto{\pgfqpoint{8.236358in}{4.145143in}}%
\pgfpathlineto{\pgfqpoint{8.237299in}{4.146513in}}%
\pgfpathlineto{\pgfqpoint{8.239182in}{4.154801in}}%
\pgfpathlineto{\pgfqpoint{8.243460in}{4.173128in}}%
\pgfpathlineto{\pgfqpoint{8.244401in}{4.172322in}}%
\pgfpathlineto{\pgfqpoint{8.246027in}{4.166280in}}%
\pgfpathlineto{\pgfqpoint{8.250989in}{4.145805in}}%
\pgfpathlineto{\pgfqpoint{8.251075in}{4.145833in}}%
\pgfpathlineto{\pgfqpoint{8.252102in}{4.147559in}}%
\pgfpathlineto{\pgfqpoint{8.254155in}{4.156965in}}%
\pgfpathlineto{\pgfqpoint{8.257920in}{4.172408in}}%
\pgfpathlineto{\pgfqpoint{8.258947in}{4.171788in}}%
\pgfpathlineto{\pgfqpoint{8.260487in}{4.166520in}}%
\pgfpathlineto{\pgfqpoint{8.265535in}{4.146450in}}%
\pgfpathlineto{\pgfqpoint{8.265706in}{4.146522in}}%
\pgfpathlineto{\pgfqpoint{8.266819in}{4.148645in}}%
\pgfpathlineto{\pgfqpoint{8.269214in}{4.160016in}}%
\pgfpathlineto{\pgfqpoint{8.272466in}{4.171803in}}%
\pgfpathlineto{\pgfqpoint{8.273492in}{4.171090in}}%
\pgfpathlineto{\pgfqpoint{8.275118in}{4.165484in}}%
\pgfpathlineto{\pgfqpoint{8.279995in}{4.147077in}}%
\pgfpathlineto{\pgfqpoint{8.280081in}{4.147105in}}%
\pgfpathlineto{\pgfqpoint{8.281107in}{4.148732in}}%
\pgfpathlineto{\pgfqpoint{8.283246in}{4.157906in}}%
\pgfpathlineto{\pgfqpoint{8.286840in}{4.171174in}}%
\pgfpathlineto{\pgfqpoint{8.287867in}{4.170553in}}%
\pgfpathlineto{\pgfqpoint{8.289493in}{4.165237in}}%
\pgfpathlineto{\pgfqpoint{8.294370in}{4.147687in}}%
\pgfpathlineto{\pgfqpoint{8.294455in}{4.147718in}}%
\pgfpathlineto{\pgfqpoint{8.295567in}{4.149555in}}%
\pgfpathlineto{\pgfqpoint{8.297878in}{4.159472in}}%
\pgfpathlineto{\pgfqpoint{8.301215in}{4.170611in}}%
\pgfpathlineto{\pgfqpoint{8.302241in}{4.169869in}}%
\pgfpathlineto{\pgfqpoint{8.303867in}{4.164592in}}%
\pgfpathlineto{\pgfqpoint{8.308659in}{4.148281in}}%
\pgfpathlineto{\pgfqpoint{8.309771in}{4.149912in}}%
\pgfpathlineto{\pgfqpoint{8.311995in}{4.158870in}}%
\pgfpathlineto{\pgfqpoint{8.315418in}{4.170028in}}%
\pgfpathlineto{\pgfqpoint{8.316445in}{4.169346in}}%
\pgfpathlineto{\pgfqpoint{8.318070in}{4.164318in}}%
\pgfpathlineto{\pgfqpoint{8.322776in}{4.148836in}}%
\pgfpathlineto{\pgfqpoint{8.323889in}{4.150282in}}%
\pgfpathlineto{\pgfqpoint{8.326028in}{4.158345in}}%
\pgfpathlineto{\pgfqpoint{8.329536in}{4.169466in}}%
\pgfpathlineto{\pgfqpoint{8.330648in}{4.168687in}}%
\pgfpathlineto{\pgfqpoint{8.332359in}{4.163340in}}%
\pgfpathlineto{\pgfqpoint{8.336809in}{4.149379in}}%
\pgfpathlineto{\pgfqpoint{8.337921in}{4.150659in}}%
\pgfpathlineto{\pgfqpoint{8.339974in}{4.157887in}}%
\pgfpathlineto{\pgfqpoint{8.343568in}{4.168924in}}%
\pgfpathlineto{\pgfqpoint{8.344680in}{4.168193in}}%
\pgfpathlineto{\pgfqpoint{8.346477in}{4.162750in}}%
\pgfpathlineto{\pgfqpoint{8.350755in}{4.149911in}}%
\pgfpathlineto{\pgfqpoint{8.351868in}{4.151036in}}%
\pgfpathlineto{\pgfqpoint{8.353921in}{4.157838in}}%
\pgfpathlineto{\pgfqpoint{8.357600in}{4.168428in}}%
\pgfpathlineto{\pgfqpoint{8.358713in}{4.167568in}}%
\pgfpathlineto{\pgfqpoint{8.360595in}{4.161861in}}%
\pgfpathlineto{\pgfqpoint{8.364616in}{4.150429in}}%
\pgfpathlineto{\pgfqpoint{8.365729in}{4.151409in}}%
\pgfpathlineto{\pgfqpoint{8.367697in}{4.157465in}}%
\pgfpathlineto{\pgfqpoint{8.371461in}{4.167921in}}%
\pgfpathlineto{\pgfqpoint{8.372574in}{4.167112in}}%
\pgfpathlineto{\pgfqpoint{8.374456in}{4.161681in}}%
\pgfpathlineto{\pgfqpoint{8.378392in}{4.150933in}}%
\pgfpathlineto{\pgfqpoint{8.379504in}{4.151776in}}%
\pgfpathlineto{\pgfqpoint{8.381387in}{4.157134in}}%
\pgfpathlineto{\pgfqpoint{8.385323in}{4.167446in}}%
\pgfpathlineto{\pgfqpoint{8.386520in}{4.166382in}}%
\pgfpathlineto{\pgfqpoint{8.388659in}{4.159969in}}%
\pgfpathlineto{\pgfqpoint{8.392082in}{4.151425in}}%
\pgfpathlineto{\pgfqpoint{8.393280in}{4.152271in}}%
\pgfpathlineto{\pgfqpoint{8.395248in}{4.157731in}}%
\pgfpathlineto{\pgfqpoint{8.399012in}{4.166973in}}%
\pgfpathlineto{\pgfqpoint{8.400210in}{4.165980in}}%
\pgfpathlineto{\pgfqpoint{8.402349in}{4.159905in}}%
\pgfpathlineto{\pgfqpoint{8.405772in}{4.151881in}}%
\pgfpathlineto{\pgfqpoint{8.406970in}{4.152739in}}%
\pgfpathlineto{\pgfqpoint{8.409023in}{4.158265in}}%
\pgfpathlineto{\pgfqpoint{8.412617in}{4.166516in}}%
\pgfpathlineto{\pgfqpoint{8.413815in}{4.165600in}}%
\pgfpathlineto{\pgfqpoint{8.415954in}{4.159862in}}%
\pgfpathlineto{\pgfqpoint{8.419376in}{4.152322in}}%
\pgfpathlineto{\pgfqpoint{8.420574in}{4.153179in}}%
\pgfpathlineto{\pgfqpoint{8.422713in}{4.158732in}}%
\pgfpathlineto{\pgfqpoint{8.426136in}{4.166075in}}%
\pgfpathlineto{\pgfqpoint{8.427419in}{4.165107in}}%
\pgfpathlineto{\pgfqpoint{8.429644in}{4.159308in}}%
\pgfpathlineto{\pgfqpoint{8.432895in}{4.152750in}}%
\pgfpathlineto{\pgfqpoint{8.434179in}{4.153724in}}%
\pgfpathlineto{\pgfqpoint{8.436489in}{4.159648in}}%
\pgfpathlineto{\pgfqpoint{8.439655in}{4.165659in}}%
\pgfpathlineto{\pgfqpoint{8.440938in}{4.164647in}}%
\pgfpathlineto{\pgfqpoint{8.443334in}{4.158590in}}%
\pgfpathlineto{\pgfqpoint{8.446329in}{4.153163in}}%
\pgfpathlineto{\pgfqpoint{8.447612in}{4.154101in}}%
\pgfpathlineto{\pgfqpoint{8.450008in}{4.159949in}}%
\pgfpathlineto{\pgfqpoint{8.453002in}{4.165249in}}%
\pgfpathlineto{\pgfqpoint{8.454286in}{4.164352in}}%
\pgfpathlineto{\pgfqpoint{8.456682in}{4.158685in}}%
\pgfpathlineto{\pgfqpoint{8.459676in}{4.153562in}}%
\pgfpathlineto{\pgfqpoint{8.461045in}{4.154576in}}%
\pgfpathlineto{\pgfqpoint{8.463612in}{4.160631in}}%
\pgfpathlineto{\pgfqpoint{8.466436in}{4.164863in}}%
\pgfpathlineto{\pgfqpoint{8.467805in}{4.163697in}}%
\pgfpathlineto{\pgfqpoint{8.470885in}{4.156519in}}%
\pgfpathlineto{\pgfqpoint{8.473281in}{4.153977in}}%
\pgfpathlineto{\pgfqpoint{8.474735in}{4.155587in}}%
\pgfpathlineto{\pgfqpoint{8.480126in}{4.164352in}}%
\pgfpathlineto{\pgfqpoint{8.480468in}{4.164095in}}%
\pgfpathlineto{\pgfqpoint{8.482350in}{4.160781in}}%
\pgfpathlineto{\pgfqpoint{8.486286in}{4.154311in}}%
\pgfpathlineto{\pgfqpoint{8.487741in}{4.155542in}}%
\pgfpathlineto{\pgfqpoint{8.493388in}{4.163941in}}%
\pgfpathlineto{\pgfqpoint{8.493901in}{4.163481in}}%
\pgfpathlineto{\pgfqpoint{8.496211in}{4.159076in}}%
\pgfpathlineto{\pgfqpoint{8.499377in}{4.154663in}}%
\pgfpathlineto{\pgfqpoint{8.500832in}{4.155770in}}%
\pgfpathlineto{\pgfqpoint{8.506565in}{4.163551in}}%
\pgfpathlineto{\pgfqpoint{8.507078in}{4.163078in}}%
\pgfpathlineto{\pgfqpoint{8.509645in}{4.158339in}}%
\pgfpathlineto{\pgfqpoint{8.512468in}{4.155003in}}%
\pgfpathlineto{\pgfqpoint{8.514008in}{4.156208in}}%
\pgfpathlineto{\pgfqpoint{8.519656in}{4.163186in}}%
\pgfpathlineto{\pgfqpoint{8.519912in}{4.162977in}}%
\pgfpathlineto{\pgfqpoint{8.522222in}{4.159231in}}%
\pgfpathlineto{\pgfqpoint{8.525388in}{4.155325in}}%
\pgfpathlineto{\pgfqpoint{8.527014in}{4.156482in}}%
\pgfpathlineto{\pgfqpoint{8.532661in}{4.162848in}}%
\pgfpathlineto{\pgfqpoint{8.532832in}{4.162716in}}%
\pgfpathlineto{\pgfqpoint{8.535142in}{4.159259in}}%
\pgfpathlineto{\pgfqpoint{8.538394in}{4.155638in}}%
\pgfpathlineto{\pgfqpoint{8.540019in}{4.156830in}}%
\pgfpathlineto{\pgfqpoint{8.545581in}{4.162539in}}%
\pgfpathlineto{\pgfqpoint{8.547806in}{4.159601in}}%
\pgfpathlineto{\pgfqpoint{8.551228in}{4.155933in}}%
\pgfpathlineto{\pgfqpoint{8.552939in}{4.157132in}}%
\pgfpathlineto{\pgfqpoint{8.558330in}{4.162309in}}%
\pgfpathlineto{\pgfqpoint{8.560554in}{4.159652in}}%
\pgfpathlineto{\pgfqpoint{8.564062in}{4.156219in}}%
\pgfpathlineto{\pgfqpoint{8.565859in}{4.157493in}}%
\pgfpathlineto{\pgfqpoint{8.570907in}{4.162133in}}%
\pgfpathlineto{\pgfqpoint{8.573047in}{4.159960in}}%
\pgfpathlineto{\pgfqpoint{8.576811in}{4.156491in}}%
\pgfpathlineto{\pgfqpoint{8.578694in}{4.157802in}}%
\pgfpathlineto{\pgfqpoint{8.583571in}{4.161888in}}%
\pgfpathlineto{\pgfqpoint{8.585795in}{4.159796in}}%
\pgfpathlineto{\pgfqpoint{8.589475in}{4.156747in}}%
\pgfpathlineto{\pgfqpoint{8.591528in}{4.158155in}}%
\pgfpathlineto{\pgfqpoint{8.596063in}{4.161689in}}%
\pgfpathlineto{\pgfqpoint{8.598287in}{4.159878in}}%
\pgfpathlineto{\pgfqpoint{8.602138in}{4.156998in}}%
\pgfpathlineto{\pgfqpoint{8.604277in}{4.158450in}}%
\pgfpathlineto{\pgfqpoint{8.608555in}{4.161474in}}%
\pgfpathlineto{\pgfqpoint{8.610865in}{4.159781in}}%
\pgfpathlineto{\pgfqpoint{8.614715in}{4.157234in}}%
\pgfpathlineto{\pgfqpoint{8.617026in}{4.158772in}}%
\pgfpathlineto{\pgfqpoint{8.621047in}{4.161255in}}%
\pgfpathlineto{\pgfqpoint{8.623528in}{4.159551in}}%
\pgfpathlineto{\pgfqpoint{8.627208in}{4.157453in}}%
\pgfpathlineto{\pgfqpoint{8.629689in}{4.159027in}}%
\pgfpathlineto{\pgfqpoint{8.633454in}{4.161053in}}%
\pgfpathlineto{\pgfqpoint{8.636021in}{4.159455in}}%
\pgfpathlineto{\pgfqpoint{8.639700in}{4.157670in}}%
\pgfpathlineto{\pgfqpoint{8.642438in}{4.159362in}}%
\pgfpathlineto{\pgfqpoint{8.645946in}{4.160835in}}%
\pgfpathlineto{\pgfqpoint{8.648855in}{4.159062in}}%
\pgfpathlineto{\pgfqpoint{8.652192in}{4.157888in}}%
\pgfpathlineto{\pgfqpoint{8.655358in}{4.159786in}}%
\pgfpathlineto{\pgfqpoint{8.658438in}{4.160615in}}%
\pgfpathlineto{\pgfqpoint{8.662117in}{4.158488in}}%
\pgfpathlineto{\pgfqpoint{8.664941in}{4.158192in}}%
\pgfpathlineto{\pgfqpoint{8.672299in}{4.159765in}}%
\pgfpathlineto{\pgfqpoint{8.677005in}{4.158292in}}%
\pgfpathlineto{\pgfqpoint{8.684962in}{4.159452in}}%
\pgfpathlineto{\pgfqpoint{8.689155in}{4.158428in}}%
\pgfpathlineto{\pgfqpoint{8.697283in}{4.159338in}}%
\pgfpathlineto{\pgfqpoint{8.701476in}{4.158605in}}%
\pgfpathlineto{\pgfqpoint{8.709091in}{4.159431in}}%
\pgfpathlineto{\pgfqpoint{8.713797in}{4.158779in}}%
\pgfpathlineto{\pgfqpoint{8.720898in}{4.159475in}}%
\pgfpathlineto{\pgfqpoint{8.726118in}{4.158945in}}%
\pgfpathlineto{\pgfqpoint{8.732706in}{4.159483in}}%
\pgfpathlineto{\pgfqpoint{8.738439in}{4.159097in}}%
\pgfpathlineto{\pgfqpoint{8.744770in}{4.159413in}}%
\pgfpathlineto{\pgfqpoint{8.750760in}{4.159225in}}%
\pgfpathlineto{\pgfqpoint{8.757006in}{4.159322in}}%
\pgfpathlineto{\pgfqpoint{8.763252in}{4.159347in}}%
\pgfpathlineto{\pgfqpoint{8.769755in}{4.159196in}}%
\pgfpathlineto{\pgfqpoint{8.776942in}{4.159495in}}%
\pgfpathlineto{\pgfqpoint{8.803723in}{4.159316in}}%
\pgfpathlineto{\pgfqpoint{8.829221in}{4.159291in}}%
\pgfpathlineto{\pgfqpoint{8.850354in}{4.159307in}}%
\pgfpathlineto{\pgfqpoint{8.850354in}{4.159307in}}%
\pgfusepath{stroke}%
\end{pgfscope}%
\begin{pgfscope}%
\pgfpathrectangle{\pgfqpoint{0.294194in}{3.039076in}}{\pgfqpoint{8.556160in}{2.240462in}}%
\pgfusepath{clip}%
\pgfsetbuttcap%
\pgfsetroundjoin%
\pgfsetlinewidth{1.505625pt}%
\definecolor{currentstroke}{rgb}{0.000000,0.000000,0.000000}%
\pgfsetstrokecolor{currentstroke}%
\pgfsetdash{{5.550000pt}{2.400000pt}}{0.000000pt}%
\pgfpathmoveto{\pgfqpoint{4.572274in}{3.039076in}}%
\pgfpathlineto{\pgfqpoint{4.572274in}{5.279538in}}%
\pgfusepath{stroke}%
\end{pgfscope}%
\begin{pgfscope}%
\pgfsetrectcap%
\pgfsetmiterjoin%
\pgfsetlinewidth{0.803000pt}%
\definecolor{currentstroke}{rgb}{0.000000,0.000000,0.000000}%
\pgfsetstrokecolor{currentstroke}%
\pgfsetdash{}{0pt}%
\pgfpathmoveto{\pgfqpoint{0.294194in}{3.039076in}}%
\pgfpathlineto{\pgfqpoint{0.294194in}{5.279538in}}%
\pgfusepath{stroke}%
\end{pgfscope}%
\begin{pgfscope}%
\pgfsetrectcap%
\pgfsetmiterjoin%
\pgfsetlinewidth{0.803000pt}%
\definecolor{currentstroke}{rgb}{0.000000,0.000000,0.000000}%
\pgfsetstrokecolor{currentstroke}%
\pgfsetdash{}{0pt}%
\pgfpathmoveto{\pgfqpoint{8.850354in}{3.039076in}}%
\pgfpathlineto{\pgfqpoint{8.850354in}{5.279538in}}%
\pgfusepath{stroke}%
\end{pgfscope}%
\begin{pgfscope}%
\pgfsetrectcap%
\pgfsetmiterjoin%
\pgfsetlinewidth{0.803000pt}%
\definecolor{currentstroke}{rgb}{0.000000,0.000000,0.000000}%
\pgfsetstrokecolor{currentstroke}%
\pgfsetdash{}{0pt}%
\pgfpathmoveto{\pgfqpoint{0.294194in}{3.039076in}}%
\pgfpathlineto{\pgfqpoint{8.850354in}{3.039076in}}%
\pgfusepath{stroke}%
\end{pgfscope}%
\begin{pgfscope}%
\pgfsetrectcap%
\pgfsetmiterjoin%
\pgfsetlinewidth{0.803000pt}%
\definecolor{currentstroke}{rgb}{0.000000,0.000000,0.000000}%
\pgfsetstrokecolor{currentstroke}%
\pgfsetdash{}{0pt}%
\pgfpathmoveto{\pgfqpoint{0.294194in}{5.279538in}}%
\pgfpathlineto{\pgfqpoint{8.850354in}{5.279538in}}%
\pgfusepath{stroke}%
\end{pgfscope}%
\begin{pgfscope}%
\definecolor{textcolor}{rgb}{0.000000,0.000000,0.000000}%
\pgfsetstrokecolor{textcolor}%
\pgfsetfillcolor{textcolor}%
\pgftext[x=9.278162in,y=4.159307in,left,base]{\color{textcolor}\rmfamily\fontsize{16.000000}{19.200000}\selectfont \(\displaystyle t=t_c\)}%
\end{pgfscope}%
\begin{pgfscope}%
\pgfsetbuttcap%
\pgfsetmiterjoin%
\definecolor{currentfill}{rgb}{1.000000,1.000000,1.000000}%
\pgfsetfillcolor{currentfill}%
\pgfsetlinewidth{0.000000pt}%
\definecolor{currentstroke}{rgb}{0.000000,0.000000,0.000000}%
\pgfsetstrokecolor{currentstroke}%
\pgfsetstrokeopacity{0.000000}%
\pgfsetdash{}{0pt}%
\pgfpathmoveto{\pgfqpoint{0.294194in}{0.558614in}}%
\pgfpathlineto{\pgfqpoint{8.850354in}{0.558614in}}%
\pgfpathlineto{\pgfqpoint{8.850354in}{2.799076in}}%
\pgfpathlineto{\pgfqpoint{0.294194in}{2.799076in}}%
\pgfpathlineto{\pgfqpoint{0.294194in}{0.558614in}}%
\pgfpathclose%
\pgfusepath{fill}%
\end{pgfscope}%
\begin{pgfscope}%
\pgfsetbuttcap%
\pgfsetroundjoin%
\definecolor{currentfill}{rgb}{0.000000,0.000000,0.000000}%
\pgfsetfillcolor{currentfill}%
\pgfsetlinewidth{0.803000pt}%
\definecolor{currentstroke}{rgb}{0.000000,0.000000,0.000000}%
\pgfsetstrokecolor{currentstroke}%
\pgfsetdash{}{0pt}%
\pgfsys@defobject{currentmarker}{\pgfqpoint{0.000000in}{-0.048611in}}{\pgfqpoint{0.000000in}{0.000000in}}{%
\pgfpathmoveto{\pgfqpoint{0.000000in}{0.000000in}}%
\pgfpathlineto{\pgfqpoint{0.000000in}{-0.048611in}}%
\pgfusepath{stroke,fill}%
}%
\begin{pgfscope}%
\pgfsys@transformshift{0.294194in}{0.558614in}%
\pgfsys@useobject{currentmarker}{}%
\end{pgfscope}%
\end{pgfscope}%
\begin{pgfscope}%
\definecolor{textcolor}{rgb}{0.000000,0.000000,0.000000}%
\pgfsetstrokecolor{textcolor}%
\pgfsetfillcolor{textcolor}%
\pgftext[x=0.294194in,y=0.461392in,,top]{\color{textcolor}\rmfamily\fontsize{16.000000}{19.200000}\selectfont \(\displaystyle 0\)}%
\end{pgfscope}%
\begin{pgfscope}%
\pgfsetbuttcap%
\pgfsetroundjoin%
\definecolor{currentfill}{rgb}{0.000000,0.000000,0.000000}%
\pgfsetfillcolor{currentfill}%
\pgfsetlinewidth{0.803000pt}%
\definecolor{currentstroke}{rgb}{0.000000,0.000000,0.000000}%
\pgfsetstrokecolor{currentstroke}%
\pgfsetdash{}{0pt}%
\pgfsys@defobject{currentmarker}{\pgfqpoint{0.000000in}{-0.048611in}}{\pgfqpoint{0.000000in}{0.000000in}}{%
\pgfpathmoveto{\pgfqpoint{0.000000in}{0.000000in}}%
\pgfpathlineto{\pgfqpoint{0.000000in}{-0.048611in}}%
\pgfusepath{stroke,fill}%
}%
\begin{pgfscope}%
\pgfsys@transformshift{2.433234in}{0.558614in}%
\pgfsys@useobject{currentmarker}{}%
\end{pgfscope}%
\end{pgfscope}%
\begin{pgfscope}%
\definecolor{textcolor}{rgb}{0.000000,0.000000,0.000000}%
\pgfsetstrokecolor{textcolor}%
\pgfsetfillcolor{textcolor}%
\pgftext[x=2.433234in,y=0.461392in,,top]{\color{textcolor}\rmfamily\fontsize{16.000000}{19.200000}\selectfont \(\displaystyle {\pi}/{4}\)}%
\end{pgfscope}%
\begin{pgfscope}%
\pgfsetbuttcap%
\pgfsetroundjoin%
\definecolor{currentfill}{rgb}{0.000000,0.000000,0.000000}%
\pgfsetfillcolor{currentfill}%
\pgfsetlinewidth{0.803000pt}%
\definecolor{currentstroke}{rgb}{0.000000,0.000000,0.000000}%
\pgfsetstrokecolor{currentstroke}%
\pgfsetdash{}{0pt}%
\pgfsys@defobject{currentmarker}{\pgfqpoint{0.000000in}{-0.048611in}}{\pgfqpoint{0.000000in}{0.000000in}}{%
\pgfpathmoveto{\pgfqpoint{0.000000in}{0.000000in}}%
\pgfpathlineto{\pgfqpoint{0.000000in}{-0.048611in}}%
\pgfusepath{stroke,fill}%
}%
\begin{pgfscope}%
\pgfsys@transformshift{4.572274in}{0.558614in}%
\pgfsys@useobject{currentmarker}{}%
\end{pgfscope}%
\end{pgfscope}%
\begin{pgfscope}%
\definecolor{textcolor}{rgb}{0.000000,0.000000,0.000000}%
\pgfsetstrokecolor{textcolor}%
\pgfsetfillcolor{textcolor}%
\pgftext[x=4.572274in,y=0.461392in,,top]{\color{textcolor}\rmfamily\fontsize{16.000000}{19.200000}\selectfont \(\displaystyle {\pi}/{2}\)}%
\end{pgfscope}%
\begin{pgfscope}%
\pgfsetbuttcap%
\pgfsetroundjoin%
\definecolor{currentfill}{rgb}{0.000000,0.000000,0.000000}%
\pgfsetfillcolor{currentfill}%
\pgfsetlinewidth{0.803000pt}%
\definecolor{currentstroke}{rgb}{0.000000,0.000000,0.000000}%
\pgfsetstrokecolor{currentstroke}%
\pgfsetdash{}{0pt}%
\pgfsys@defobject{currentmarker}{\pgfqpoint{0.000000in}{-0.048611in}}{\pgfqpoint{0.000000in}{0.000000in}}{%
\pgfpathmoveto{\pgfqpoint{0.000000in}{0.000000in}}%
\pgfpathlineto{\pgfqpoint{0.000000in}{-0.048611in}}%
\pgfusepath{stroke,fill}%
}%
\begin{pgfscope}%
\pgfsys@transformshift{6.711314in}{0.558614in}%
\pgfsys@useobject{currentmarker}{}%
\end{pgfscope}%
\end{pgfscope}%
\begin{pgfscope}%
\definecolor{textcolor}{rgb}{0.000000,0.000000,0.000000}%
\pgfsetstrokecolor{textcolor}%
\pgfsetfillcolor{textcolor}%
\pgftext[x=6.711314in,y=0.461392in,,top]{\color{textcolor}\rmfamily\fontsize{16.000000}{19.200000}\selectfont \(\displaystyle {3\pi}/{4}\)}%
\end{pgfscope}%
\begin{pgfscope}%
\pgfsetbuttcap%
\pgfsetroundjoin%
\definecolor{currentfill}{rgb}{0.000000,0.000000,0.000000}%
\pgfsetfillcolor{currentfill}%
\pgfsetlinewidth{0.803000pt}%
\definecolor{currentstroke}{rgb}{0.000000,0.000000,0.000000}%
\pgfsetstrokecolor{currentstroke}%
\pgfsetdash{}{0pt}%
\pgfsys@defobject{currentmarker}{\pgfqpoint{0.000000in}{-0.048611in}}{\pgfqpoint{0.000000in}{0.000000in}}{%
\pgfpathmoveto{\pgfqpoint{0.000000in}{0.000000in}}%
\pgfpathlineto{\pgfqpoint{0.000000in}{-0.048611in}}%
\pgfusepath{stroke,fill}%
}%
\begin{pgfscope}%
\pgfsys@transformshift{8.850354in}{0.558614in}%
\pgfsys@useobject{currentmarker}{}%
\end{pgfscope}%
\end{pgfscope}%
\begin{pgfscope}%
\definecolor{textcolor}{rgb}{0.000000,0.000000,0.000000}%
\pgfsetstrokecolor{textcolor}%
\pgfsetfillcolor{textcolor}%
\pgftext[x=8.850354in,y=0.461392in,,top]{\color{textcolor}\rmfamily\fontsize{16.000000}{19.200000}\selectfont \(\displaystyle \pi\)}%
\end{pgfscope}%
\begin{pgfscope}%
\pgfpathrectangle{\pgfqpoint{0.294194in}{0.558614in}}{\pgfqpoint{8.556160in}{2.240462in}}%
\pgfusepath{clip}%
\pgfsetrectcap%
\pgfsetroundjoin%
\pgfsetlinewidth{1.505625pt}%
\definecolor{currentstroke}{rgb}{0.121569,0.466667,0.705882}%
\pgfsetstrokecolor{currentstroke}%
\pgfsetdash{}{0pt}%
\pgfpathmoveto{\pgfqpoint{0.294194in}{2.799076in}}%
\pgfpathlineto{\pgfqpoint{0.294537in}{2.779236in}}%
\pgfpathlineto{\pgfqpoint{0.295478in}{2.530816in}}%
\pgfpathlineto{\pgfqpoint{0.299927in}{0.558693in}}%
\pgfpathlineto{\pgfqpoint{0.300868in}{0.711570in}}%
\pgfpathlineto{\pgfqpoint{0.302751in}{1.668947in}}%
\pgfpathlineto{\pgfqpoint{0.305660in}{2.798936in}}%
\pgfpathlineto{\pgfqpoint{0.306088in}{2.764263in}}%
\pgfpathlineto{\pgfqpoint{0.307200in}{2.410405in}}%
\pgfpathlineto{\pgfqpoint{0.311392in}{0.558697in}}%
\pgfpathlineto{\pgfqpoint{0.311991in}{0.622253in}}%
\pgfpathlineto{\pgfqpoint{0.313360in}{1.157297in}}%
\pgfpathlineto{\pgfqpoint{0.317125in}{2.799056in}}%
\pgfpathlineto{\pgfqpoint{0.317382in}{2.788237in}}%
\pgfpathlineto{\pgfqpoint{0.318237in}{2.599942in}}%
\pgfpathlineto{\pgfqpoint{0.320462in}{1.405728in}}%
\pgfpathlineto{\pgfqpoint{0.322858in}{0.558892in}}%
\pgfpathlineto{\pgfqpoint{0.322943in}{0.559011in}}%
\pgfpathlineto{\pgfqpoint{0.323371in}{0.595767in}}%
\pgfpathlineto{\pgfqpoint{0.324483in}{0.950321in}}%
\pgfpathlineto{\pgfqpoint{0.328676in}{2.799027in}}%
\pgfpathlineto{\pgfqpoint{0.329275in}{2.741417in}}%
\pgfpathlineto{\pgfqpoint{0.330558in}{2.267157in}}%
\pgfpathlineto{\pgfqpoint{0.334494in}{0.558706in}}%
\pgfpathlineto{\pgfqpoint{0.334837in}{0.579369in}}%
\pgfpathlineto{\pgfqpoint{0.335778in}{0.823169in}}%
\pgfpathlineto{\pgfqpoint{0.340313in}{2.798948in}}%
\pgfpathlineto{\pgfqpoint{0.341339in}{2.626043in}}%
\pgfpathlineto{\pgfqpoint{0.343307in}{1.619298in}}%
\pgfpathlineto{\pgfqpoint{0.346131in}{0.558721in}}%
\pgfpathlineto{\pgfqpoint{0.346473in}{0.578255in}}%
\pgfpathlineto{\pgfqpoint{0.347414in}{0.817360in}}%
\pgfpathlineto{\pgfqpoint{0.351949in}{2.798874in}}%
\pgfpathlineto{\pgfqpoint{0.353061in}{2.612570in}}%
\pgfpathlineto{\pgfqpoint{0.355200in}{1.498857in}}%
\pgfpathlineto{\pgfqpoint{0.357853in}{0.558839in}}%
\pgfpathlineto{\pgfqpoint{0.358195in}{0.579895in}}%
\pgfpathlineto{\pgfqpoint{0.359222in}{0.854696in}}%
\pgfpathlineto{\pgfqpoint{0.363671in}{2.798675in}}%
\pgfpathlineto{\pgfqpoint{0.364612in}{2.670964in}}%
\pgfpathlineto{\pgfqpoint{0.366324in}{1.874174in}}%
\pgfpathlineto{\pgfqpoint{0.369575in}{0.558932in}}%
\pgfpathlineto{\pgfqpoint{0.369660in}{0.559403in}}%
\pgfpathlineto{\pgfqpoint{0.370174in}{0.610633in}}%
\pgfpathlineto{\pgfqpoint{0.371457in}{1.062746in}}%
\pgfpathlineto{\pgfqpoint{0.375479in}{2.798687in}}%
\pgfpathlineto{\pgfqpoint{0.375907in}{2.773893in}}%
\pgfpathlineto{\pgfqpoint{0.376933in}{2.493413in}}%
\pgfpathlineto{\pgfqpoint{0.381468in}{0.559101in}}%
\pgfpathlineto{\pgfqpoint{0.382324in}{0.681313in}}%
\pgfpathlineto{\pgfqpoint{0.384035in}{1.463655in}}%
\pgfpathlineto{\pgfqpoint{0.387372in}{2.798748in}}%
\pgfpathlineto{\pgfqpoint{0.387714in}{2.780580in}}%
\pgfpathlineto{\pgfqpoint{0.388655in}{2.551565in}}%
\pgfpathlineto{\pgfqpoint{0.391393in}{1.092745in}}%
\pgfpathlineto{\pgfqpoint{0.393361in}{0.559147in}}%
\pgfpathlineto{\pgfqpoint{0.393532in}{0.565404in}}%
\pgfpathlineto{\pgfqpoint{0.394302in}{0.702922in}}%
\pgfpathlineto{\pgfqpoint{0.396099in}{1.552993in}}%
\pgfpathlineto{\pgfqpoint{0.399265in}{2.798451in}}%
\pgfpathlineto{\pgfqpoint{0.399436in}{2.795870in}}%
\pgfpathlineto{\pgfqpoint{0.400035in}{2.716633in}}%
\pgfpathlineto{\pgfqpoint{0.401490in}{2.137674in}}%
\pgfpathlineto{\pgfqpoint{0.405254in}{0.559287in}}%
\pgfpathlineto{\pgfqpoint{0.405426in}{0.561857in}}%
\pgfpathlineto{\pgfqpoint{0.406024in}{0.640670in}}%
\pgfpathlineto{\pgfqpoint{0.407479in}{1.216829in}}%
\pgfpathlineto{\pgfqpoint{0.411329in}{2.798399in}}%
\pgfpathlineto{\pgfqpoint{0.411500in}{2.792242in}}%
\pgfpathlineto{\pgfqpoint{0.412271in}{2.656942in}}%
\pgfpathlineto{\pgfqpoint{0.414067in}{1.818063in}}%
\pgfpathlineto{\pgfqpoint{0.417319in}{0.559187in}}%
\pgfpathlineto{\pgfqpoint{0.417661in}{0.576821in}}%
\pgfpathlineto{\pgfqpoint{0.418602in}{0.799694in}}%
\pgfpathlineto{\pgfqpoint{0.421255in}{2.192809in}}%
\pgfpathlineto{\pgfqpoint{0.423394in}{2.798270in}}%
\pgfpathlineto{\pgfqpoint{0.423479in}{2.796280in}}%
\pgfpathlineto{\pgfqpoint{0.424078in}{2.721356in}}%
\pgfpathlineto{\pgfqpoint{0.425533in}{2.158277in}}%
\pgfpathlineto{\pgfqpoint{0.429383in}{0.559470in}}%
\pgfpathlineto{\pgfqpoint{0.429640in}{0.566799in}}%
\pgfpathlineto{\pgfqpoint{0.430410in}{0.704251in}}%
\pgfpathlineto{\pgfqpoint{0.432292in}{1.586942in}}%
\pgfpathlineto{\pgfqpoint{0.435458in}{2.798168in}}%
\pgfpathlineto{\pgfqpoint{0.435629in}{2.795441in}}%
\pgfpathlineto{\pgfqpoint{0.436314in}{2.698634in}}%
\pgfpathlineto{\pgfqpoint{0.437939in}{2.016716in}}%
\pgfpathlineto{\pgfqpoint{0.441618in}{0.559637in}}%
\pgfpathlineto{\pgfqpoint{0.442046in}{0.591347in}}%
\pgfpathlineto{\pgfqpoint{0.443159in}{0.908577in}}%
\pgfpathlineto{\pgfqpoint{0.447693in}{2.798185in}}%
\pgfpathlineto{\pgfqpoint{0.448463in}{2.711143in}}%
\pgfpathlineto{\pgfqpoint{0.450004in}{2.097679in}}%
\pgfpathlineto{\pgfqpoint{0.453854in}{0.559778in}}%
\pgfpathlineto{\pgfqpoint{0.453939in}{0.561783in}}%
\pgfpathlineto{\pgfqpoint{0.454538in}{0.635091in}}%
\pgfpathlineto{\pgfqpoint{0.455993in}{1.184882in}}%
\pgfpathlineto{\pgfqpoint{0.459929in}{2.797887in}}%
\pgfpathlineto{\pgfqpoint{0.460186in}{2.790700in}}%
\pgfpathlineto{\pgfqpoint{0.460956in}{2.656908in}}%
\pgfpathlineto{\pgfqpoint{0.462838in}{1.793923in}}%
\pgfpathlineto{\pgfqpoint{0.466089in}{0.559898in}}%
\pgfpathlineto{\pgfqpoint{0.466260in}{0.562455in}}%
\pgfpathlineto{\pgfqpoint{0.466859in}{0.637350in}}%
\pgfpathlineto{\pgfqpoint{0.468314in}{1.185811in}}%
\pgfpathlineto{\pgfqpoint{0.472335in}{2.797734in}}%
\pgfpathlineto{\pgfqpoint{0.472507in}{2.791956in}}%
\pgfpathlineto{\pgfqpoint{0.473277in}{2.664261in}}%
\pgfpathlineto{\pgfqpoint{0.475073in}{1.863688in}}%
\pgfpathlineto{\pgfqpoint{0.478496in}{0.559902in}}%
\pgfpathlineto{\pgfqpoint{0.478753in}{0.568853in}}%
\pgfpathlineto{\pgfqpoint{0.479608in}{0.730499in}}%
\pgfpathlineto{\pgfqpoint{0.481662in}{1.708837in}}%
\pgfpathlineto{\pgfqpoint{0.484656in}{2.797140in}}%
\pgfpathlineto{\pgfqpoint{0.484742in}{2.797634in}}%
\pgfpathlineto{\pgfqpoint{0.484913in}{2.792372in}}%
\pgfpathlineto{\pgfqpoint{0.485683in}{2.668088in}}%
\pgfpathlineto{\pgfqpoint{0.487480in}{1.877940in}}%
\pgfpathlineto{\pgfqpoint{0.490902in}{0.560468in}}%
\pgfpathlineto{\pgfqpoint{0.490988in}{0.560228in}}%
\pgfpathlineto{\pgfqpoint{0.491074in}{0.562060in}}%
\pgfpathlineto{\pgfqpoint{0.491673in}{0.632111in}}%
\pgfpathlineto{\pgfqpoint{0.493127in}{1.163426in}}%
\pgfpathlineto{\pgfqpoint{0.497234in}{2.797437in}}%
\pgfpathlineto{\pgfqpoint{0.497405in}{2.792143in}}%
\pgfpathlineto{\pgfqpoint{0.498175in}{2.668927in}}%
\pgfpathlineto{\pgfqpoint{0.499972in}{1.886000in}}%
\pgfpathlineto{\pgfqpoint{0.503480in}{0.560268in}}%
\pgfpathlineto{\pgfqpoint{0.503822in}{0.576333in}}%
\pgfpathlineto{\pgfqpoint{0.504764in}{0.781715in}}%
\pgfpathlineto{\pgfqpoint{0.507159in}{1.972641in}}%
\pgfpathlineto{\pgfqpoint{0.509726in}{2.796985in}}%
\pgfpathlineto{\pgfqpoint{0.509812in}{2.797139in}}%
\pgfpathlineto{\pgfqpoint{0.509897in}{2.795260in}}%
\pgfpathlineto{\pgfqpoint{0.510496in}{2.725923in}}%
\pgfpathlineto{\pgfqpoint{0.511951in}{2.202244in}}%
\pgfpathlineto{\pgfqpoint{0.516058in}{0.560592in}}%
\pgfpathlineto{\pgfqpoint{0.516314in}{0.567597in}}%
\pgfpathlineto{\pgfqpoint{0.517085in}{0.694988in}}%
\pgfpathlineto{\pgfqpoint{0.518881in}{1.475153in}}%
\pgfpathlineto{\pgfqpoint{0.522389in}{2.796995in}}%
\pgfpathlineto{\pgfqpoint{0.522561in}{2.794340in}}%
\pgfpathlineto{\pgfqpoint{0.523245in}{2.704674in}}%
\pgfpathlineto{\pgfqpoint{0.524785in}{2.113197in}}%
\pgfpathlineto{\pgfqpoint{0.528807in}{0.560881in}}%
\pgfpathlineto{\pgfqpoint{0.528892in}{0.562750in}}%
\pgfpathlineto{\pgfqpoint{0.529491in}{0.630977in}}%
\pgfpathlineto{\pgfqpoint{0.530946in}{1.146263in}}%
\pgfpathlineto{\pgfqpoint{0.535138in}{2.796885in}}%
\pgfpathlineto{\pgfqpoint{0.535395in}{2.788086in}}%
\pgfpathlineto{\pgfqpoint{0.536251in}{2.633935in}}%
\pgfpathlineto{\pgfqpoint{0.538218in}{1.741636in}}%
\pgfpathlineto{\pgfqpoint{0.541470in}{0.561418in}}%
\pgfpathlineto{\pgfqpoint{0.541555in}{0.561006in}}%
\pgfpathlineto{\pgfqpoint{0.541641in}{0.562564in}}%
\pgfpathlineto{\pgfqpoint{0.542240in}{0.627978in}}%
\pgfpathlineto{\pgfqpoint{0.543694in}{1.133813in}}%
\pgfpathlineto{\pgfqpoint{0.547973in}{2.796516in}}%
\pgfpathlineto{\pgfqpoint{0.548144in}{2.791172in}}%
\pgfpathlineto{\pgfqpoint{0.548914in}{2.672553in}}%
\pgfpathlineto{\pgfqpoint{0.550711in}{1.919319in}}%
\pgfpathlineto{\pgfqpoint{0.554304in}{0.561788in}}%
\pgfpathlineto{\pgfqpoint{0.554390in}{0.561208in}}%
\pgfpathlineto{\pgfqpoint{0.554561in}{0.565883in}}%
\pgfpathlineto{\pgfqpoint{0.555331in}{0.681022in}}%
\pgfpathlineto{\pgfqpoint{0.557042in}{1.381148in}}%
\pgfpathlineto{\pgfqpoint{0.560807in}{2.796355in}}%
\pgfpathlineto{\pgfqpoint{0.561064in}{2.788832in}}%
\pgfpathlineto{\pgfqpoint{0.561834in}{2.664550in}}%
\pgfpathlineto{\pgfqpoint{0.563631in}{1.909024in}}%
\pgfpathlineto{\pgfqpoint{0.567224in}{0.562118in}}%
\pgfpathlineto{\pgfqpoint{0.567310in}{0.561450in}}%
\pgfpathlineto{\pgfqpoint{0.567481in}{0.565875in}}%
\pgfpathlineto{\pgfqpoint{0.568165in}{0.658938in}}%
\pgfpathlineto{\pgfqpoint{0.569791in}{1.282275in}}%
\pgfpathlineto{\pgfqpoint{0.573812in}{2.796062in}}%
\pgfpathlineto{\pgfqpoint{0.574155in}{2.778843in}}%
\pgfpathlineto{\pgfqpoint{0.575096in}{2.581278in}}%
\pgfpathlineto{\pgfqpoint{0.577492in}{1.439959in}}%
\pgfpathlineto{\pgfqpoint{0.580230in}{0.562429in}}%
\pgfpathlineto{\pgfqpoint{0.580315in}{0.561717in}}%
\pgfpathlineto{\pgfqpoint{0.580486in}{0.565976in}}%
\pgfpathlineto{\pgfqpoint{0.581171in}{0.657417in}}%
\pgfpathlineto{\pgfqpoint{0.582797in}{1.272985in}}%
\pgfpathlineto{\pgfqpoint{0.586818in}{2.795781in}}%
\pgfpathlineto{\pgfqpoint{0.586904in}{2.795335in}}%
\pgfpathlineto{\pgfqpoint{0.587417in}{2.753396in}}%
\pgfpathlineto{\pgfqpoint{0.588700in}{2.380090in}}%
\pgfpathlineto{\pgfqpoint{0.593406in}{0.562000in}}%
\pgfpathlineto{\pgfqpoint{0.593920in}{0.596578in}}%
\pgfpathlineto{\pgfqpoint{0.595118in}{0.918640in}}%
\pgfpathlineto{\pgfqpoint{0.599995in}{2.795511in}}%
\pgfpathlineto{\pgfqpoint{0.600679in}{2.733623in}}%
\pgfpathlineto{\pgfqpoint{0.602048in}{2.294613in}}%
\pgfpathlineto{\pgfqpoint{0.606583in}{0.562297in}}%
\pgfpathlineto{\pgfqpoint{0.606925in}{0.577424in}}%
\pgfpathlineto{\pgfqpoint{0.607866in}{0.764815in}}%
\pgfpathlineto{\pgfqpoint{0.610177in}{1.831559in}}%
\pgfpathlineto{\pgfqpoint{0.613086in}{2.793556in}}%
\pgfpathlineto{\pgfqpoint{0.613171in}{2.795123in}}%
\pgfpathlineto{\pgfqpoint{0.613513in}{2.783088in}}%
\pgfpathlineto{\pgfqpoint{0.614455in}{2.604533in}}%
\pgfpathlineto{\pgfqpoint{0.616679in}{1.595791in}}%
\pgfpathlineto{\pgfqpoint{0.619674in}{0.566347in}}%
\pgfpathlineto{\pgfqpoint{0.619845in}{0.562614in}}%
\pgfpathlineto{\pgfqpoint{0.620187in}{0.576950in}}%
\pgfpathlineto{\pgfqpoint{0.621129in}{0.760377in}}%
\pgfpathlineto{\pgfqpoint{0.623353in}{1.770437in}}%
\pgfpathlineto{\pgfqpoint{0.626348in}{2.791435in}}%
\pgfpathlineto{\pgfqpoint{0.626519in}{2.794910in}}%
\pgfpathlineto{\pgfqpoint{0.626861in}{2.780210in}}%
\pgfpathlineto{\pgfqpoint{0.627802in}{2.596841in}}%
\pgfpathlineto{\pgfqpoint{0.630027in}{1.590656in}}%
\pgfpathlineto{\pgfqpoint{0.633022in}{0.567150in}}%
\pgfpathlineto{\pgfqpoint{0.633193in}{0.562972in}}%
\pgfpathlineto{\pgfqpoint{0.633621in}{0.583860in}}%
\pgfpathlineto{\pgfqpoint{0.634647in}{0.807592in}}%
\pgfpathlineto{\pgfqpoint{0.637300in}{2.056738in}}%
\pgfpathlineto{\pgfqpoint{0.639867in}{2.794208in}}%
\pgfpathlineto{\pgfqpoint{0.639952in}{2.794452in}}%
\pgfpathlineto{\pgfqpoint{0.640038in}{2.792914in}}%
\pgfpathlineto{\pgfqpoint{0.640637in}{2.732851in}}%
\pgfpathlineto{\pgfqpoint{0.642006in}{2.306674in}}%
\pgfpathlineto{\pgfqpoint{0.646626in}{0.563427in}}%
\pgfpathlineto{\pgfqpoint{0.647054in}{0.582110in}}%
\pgfpathlineto{\pgfqpoint{0.648081in}{0.798662in}}%
\pgfpathlineto{\pgfqpoint{0.650648in}{1.991254in}}%
\pgfpathlineto{\pgfqpoint{0.653300in}{2.792485in}}%
\pgfpathlineto{\pgfqpoint{0.653386in}{2.794069in}}%
\pgfpathlineto{\pgfqpoint{0.653728in}{2.782861in}}%
\pgfpathlineto{\pgfqpoint{0.654669in}{2.612320in}}%
\pgfpathlineto{\pgfqpoint{0.656808in}{1.679619in}}%
\pgfpathlineto{\pgfqpoint{0.660059in}{0.566185in}}%
\pgfpathlineto{\pgfqpoint{0.660231in}{0.563743in}}%
\pgfpathlineto{\pgfqpoint{0.660573in}{0.579746in}}%
\pgfpathlineto{\pgfqpoint{0.661600in}{0.787286in}}%
\pgfpathlineto{\pgfqpoint{0.664081in}{1.920926in}}%
\pgfpathlineto{\pgfqpoint{0.666904in}{2.792420in}}%
\pgfpathlineto{\pgfqpoint{0.666990in}{2.793769in}}%
\pgfpathlineto{\pgfqpoint{0.667332in}{2.781875in}}%
\pgfpathlineto{\pgfqpoint{0.668273in}{2.611551in}}%
\pgfpathlineto{\pgfqpoint{0.670412in}{1.686592in}}%
\pgfpathlineto{\pgfqpoint{0.673664in}{0.567990in}}%
\pgfpathlineto{\pgfqpoint{0.673835in}{0.564065in}}%
\pgfpathlineto{\pgfqpoint{0.674263in}{0.584261in}}%
\pgfpathlineto{\pgfqpoint{0.675290in}{0.799319in}}%
\pgfpathlineto{\pgfqpoint{0.677856in}{1.973129in}}%
\pgfpathlineto{\pgfqpoint{0.680594in}{2.792021in}}%
\pgfpathlineto{\pgfqpoint{0.680680in}{2.793372in}}%
\pgfpathlineto{\pgfqpoint{0.681022in}{2.781739in}}%
\pgfpathlineto{\pgfqpoint{0.681963in}{2.614050in}}%
\pgfpathlineto{\pgfqpoint{0.684102in}{1.699609in}}%
\pgfpathlineto{\pgfqpoint{0.687439in}{0.566986in}}%
\pgfpathlineto{\pgfqpoint{0.687611in}{0.564508in}}%
\pgfpathlineto{\pgfqpoint{0.687953in}{0.579836in}}%
\pgfpathlineto{\pgfqpoint{0.688894in}{0.755966in}}%
\pgfpathlineto{\pgfqpoint{0.691119in}{1.719477in}}%
\pgfpathlineto{\pgfqpoint{0.694284in}{2.787936in}}%
\pgfpathlineto{\pgfqpoint{0.694456in}{2.792856in}}%
\pgfpathlineto{\pgfqpoint{0.694883in}{2.775793in}}%
\pgfpathlineto{\pgfqpoint{0.695910in}{2.571172in}}%
\pgfpathlineto{\pgfqpoint{0.698391in}{1.461825in}}%
\pgfpathlineto{\pgfqpoint{0.701301in}{0.566722in}}%
\pgfpathlineto{\pgfqpoint{0.701386in}{0.565067in}}%
\pgfpathlineto{\pgfqpoint{0.701728in}{0.575109in}}%
\pgfpathlineto{\pgfqpoint{0.702584in}{0.713538in}}%
\pgfpathlineto{\pgfqpoint{0.704552in}{1.502281in}}%
\pgfpathlineto{\pgfqpoint{0.708231in}{2.789869in}}%
\pgfpathlineto{\pgfqpoint{0.708402in}{2.792585in}}%
\pgfpathlineto{\pgfqpoint{0.708744in}{2.778182in}}%
\pgfpathlineto{\pgfqpoint{0.709686in}{2.607354in}}%
\pgfpathlineto{\pgfqpoint{0.711910in}{1.661652in}}%
\pgfpathlineto{\pgfqpoint{0.715162in}{0.570228in}}%
\pgfpathlineto{\pgfqpoint{0.715333in}{0.565465in}}%
\pgfpathlineto{\pgfqpoint{0.715761in}{0.582260in}}%
\pgfpathlineto{\pgfqpoint{0.716787in}{0.782690in}}%
\pgfpathlineto{\pgfqpoint{0.719269in}{1.874978in}}%
\pgfpathlineto{\pgfqpoint{0.722178in}{2.787530in}}%
\pgfpathlineto{\pgfqpoint{0.722349in}{2.792047in}}%
\pgfpathlineto{\pgfqpoint{0.722777in}{2.774858in}}%
\pgfpathlineto{\pgfqpoint{0.723803in}{2.574765in}}%
\pgfpathlineto{\pgfqpoint{0.726285in}{1.487235in}}%
\pgfpathlineto{\pgfqpoint{0.729279in}{0.567900in}}%
\pgfpathlineto{\pgfqpoint{0.729451in}{0.565861in}}%
\pgfpathlineto{\pgfqpoint{0.729707in}{0.574923in}}%
\pgfpathlineto{\pgfqpoint{0.730563in}{0.707167in}}%
\pgfpathlineto{\pgfqpoint{0.732531in}{1.473691in}}%
\pgfpathlineto{\pgfqpoint{0.736296in}{2.787869in}}%
\pgfpathlineto{\pgfqpoint{0.736467in}{2.791678in}}%
\pgfpathlineto{\pgfqpoint{0.736895in}{2.773169in}}%
\pgfpathlineto{\pgfqpoint{0.737921in}{2.572583in}}%
\pgfpathlineto{\pgfqpoint{0.740403in}{1.493521in}}%
\pgfpathlineto{\pgfqpoint{0.743397in}{0.569599in}}%
\pgfpathlineto{\pgfqpoint{0.743568in}{0.566224in}}%
\pgfpathlineto{\pgfqpoint{0.743911in}{0.578560in}}%
\pgfpathlineto{\pgfqpoint{0.744852in}{0.739180in}}%
\pgfpathlineto{\pgfqpoint{0.746991in}{1.610930in}}%
\pgfpathlineto{\pgfqpoint{0.750499in}{2.787222in}}%
\pgfpathlineto{\pgfqpoint{0.750670in}{2.791176in}}%
\pgfpathlineto{\pgfqpoint{0.751098in}{2.773465in}}%
\pgfpathlineto{\pgfqpoint{0.752125in}{2.577032in}}%
\pgfpathlineto{\pgfqpoint{0.754520in}{1.552378in}}%
\pgfpathlineto{\pgfqpoint{0.757601in}{0.572706in}}%
\pgfpathlineto{\pgfqpoint{0.757857in}{0.566754in}}%
\pgfpathlineto{\pgfqpoint{0.758285in}{0.588103in}}%
\pgfpathlineto{\pgfqpoint{0.759397in}{0.817058in}}%
\pgfpathlineto{\pgfqpoint{0.762135in}{2.021115in}}%
\pgfpathlineto{\pgfqpoint{0.764873in}{2.788717in}}%
\pgfpathlineto{\pgfqpoint{0.765045in}{2.790608in}}%
\pgfpathlineto{\pgfqpoint{0.765301in}{2.781796in}}%
\pgfpathlineto{\pgfqpoint{0.766157in}{2.654265in}}%
\pgfpathlineto{\pgfqpoint{0.768039in}{1.952189in}}%
\pgfpathlineto{\pgfqpoint{0.772061in}{0.569662in}}%
\pgfpathlineto{\pgfqpoint{0.772232in}{0.567257in}}%
\pgfpathlineto{\pgfqpoint{0.772574in}{0.580923in}}%
\pgfpathlineto{\pgfqpoint{0.773515in}{0.741014in}}%
\pgfpathlineto{\pgfqpoint{0.775654in}{1.596571in}}%
\pgfpathlineto{\pgfqpoint{0.779162in}{2.781960in}}%
\pgfpathlineto{\pgfqpoint{0.779419in}{2.790147in}}%
\pgfpathlineto{\pgfqpoint{0.779932in}{2.765332in}}%
\pgfpathlineto{\pgfqpoint{0.781045in}{2.533126in}}%
\pgfpathlineto{\pgfqpoint{0.783868in}{1.302512in}}%
\pgfpathlineto{\pgfqpoint{0.786606in}{0.568715in}}%
\pgfpathlineto{\pgfqpoint{0.786692in}{0.567760in}}%
\pgfpathlineto{\pgfqpoint{0.786949in}{0.573993in}}%
\pgfpathlineto{\pgfqpoint{0.787804in}{0.690994in}}%
\pgfpathlineto{\pgfqpoint{0.789687in}{1.368266in}}%
\pgfpathlineto{\pgfqpoint{0.793794in}{2.785642in}}%
\pgfpathlineto{\pgfqpoint{0.793965in}{2.789597in}}%
\pgfpathlineto{\pgfqpoint{0.794392in}{2.773200in}}%
\pgfpathlineto{\pgfqpoint{0.795419in}{2.586875in}}%
\pgfpathlineto{\pgfqpoint{0.797815in}{1.598917in}}%
\pgfpathlineto{\pgfqpoint{0.801066in}{0.574334in}}%
\pgfpathlineto{\pgfqpoint{0.801323in}{0.568314in}}%
\pgfpathlineto{\pgfqpoint{0.801751in}{0.588052in}}%
\pgfpathlineto{\pgfqpoint{0.802863in}{0.805039in}}%
\pgfpathlineto{\pgfqpoint{0.805516in}{1.930682in}}%
\pgfpathlineto{\pgfqpoint{0.808425in}{2.783615in}}%
\pgfpathlineto{\pgfqpoint{0.808681in}{2.789058in}}%
\pgfpathlineto{\pgfqpoint{0.809109in}{2.768611in}}%
\pgfpathlineto{\pgfqpoint{0.810222in}{2.551321in}}%
\pgfpathlineto{\pgfqpoint{0.812874in}{1.430357in}}%
\pgfpathlineto{\pgfqpoint{0.815783in}{0.575295in}}%
\pgfpathlineto{\pgfqpoint{0.816040in}{0.568866in}}%
\pgfpathlineto{\pgfqpoint{0.816468in}{0.587426in}}%
\pgfpathlineto{\pgfqpoint{0.817494in}{0.774913in}}%
\pgfpathlineto{\pgfqpoint{0.819890in}{1.750626in}}%
\pgfpathlineto{\pgfqpoint{0.823141in}{2.779352in}}%
\pgfpathlineto{\pgfqpoint{0.823484in}{2.788390in}}%
\pgfpathlineto{\pgfqpoint{0.823912in}{2.767050in}}%
\pgfpathlineto{\pgfqpoint{0.825024in}{2.550385in}}%
\pgfpathlineto{\pgfqpoint{0.827676in}{1.439855in}}%
\pgfpathlineto{\pgfqpoint{0.830671in}{0.573998in}}%
\pgfpathlineto{\pgfqpoint{0.830928in}{0.569627in}}%
\pgfpathlineto{\pgfqpoint{0.831270in}{0.583947in}}%
\pgfpathlineto{\pgfqpoint{0.832297in}{0.759621in}}%
\pgfpathlineto{\pgfqpoint{0.834607in}{1.673705in}}%
\pgfpathlineto{\pgfqpoint{0.838029in}{2.777442in}}%
\pgfpathlineto{\pgfqpoint{0.838372in}{2.787913in}}%
\pgfpathlineto{\pgfqpoint{0.838799in}{2.768918in}}%
\pgfpathlineto{\pgfqpoint{0.839912in}{2.560512in}}%
\pgfpathlineto{\pgfqpoint{0.842479in}{1.506233in}}%
\pgfpathlineto{\pgfqpoint{0.845559in}{0.578874in}}%
\pgfpathlineto{\pgfqpoint{0.845901in}{0.570223in}}%
\pgfpathlineto{\pgfqpoint{0.846329in}{0.591205in}}%
\pgfpathlineto{\pgfqpoint{0.847441in}{0.802935in}}%
\pgfpathlineto{\pgfqpoint{0.850094in}{1.894211in}}%
\pgfpathlineto{\pgfqpoint{0.853088in}{2.778230in}}%
\pgfpathlineto{\pgfqpoint{0.853431in}{2.787198in}}%
\pgfpathlineto{\pgfqpoint{0.853858in}{2.766890in}}%
\pgfpathlineto{\pgfqpoint{0.854971in}{2.558114in}}%
\pgfpathlineto{\pgfqpoint{0.857538in}{1.513302in}}%
\pgfpathlineto{\pgfqpoint{0.860703in}{0.577164in}}%
\pgfpathlineto{\pgfqpoint{0.860960in}{0.570701in}}%
\pgfpathlineto{\pgfqpoint{0.861388in}{0.587710in}}%
\pgfpathlineto{\pgfqpoint{0.862415in}{0.764565in}}%
\pgfpathlineto{\pgfqpoint{0.864725in}{1.662339in}}%
\pgfpathlineto{\pgfqpoint{0.868233in}{2.776270in}}%
\pgfpathlineto{\pgfqpoint{0.868575in}{2.786657in}}%
\pgfpathlineto{\pgfqpoint{0.869003in}{2.768677in}}%
\pgfpathlineto{\pgfqpoint{0.870030in}{2.590875in}}%
\pgfpathlineto{\pgfqpoint{0.872340in}{1.696473in}}%
\pgfpathlineto{\pgfqpoint{0.875848in}{0.582891in}}%
\pgfpathlineto{\pgfqpoint{0.876190in}{0.571363in}}%
\pgfpathlineto{\pgfqpoint{0.876704in}{0.594961in}}%
\pgfpathlineto{\pgfqpoint{0.877816in}{0.806154in}}%
\pgfpathlineto{\pgfqpoint{0.880468in}{1.876664in}}%
\pgfpathlineto{\pgfqpoint{0.883549in}{2.776572in}}%
\pgfpathlineto{\pgfqpoint{0.883891in}{2.785951in}}%
\pgfpathlineto{\pgfqpoint{0.884319in}{2.767272in}}%
\pgfpathlineto{\pgfqpoint{0.885431in}{2.567813in}}%
\pgfpathlineto{\pgfqpoint{0.887912in}{1.589961in}}%
\pgfpathlineto{\pgfqpoint{0.891249in}{0.581483in}}%
\pgfpathlineto{\pgfqpoint{0.891591in}{0.572074in}}%
\pgfpathlineto{\pgfqpoint{0.892019in}{0.590436in}}%
\pgfpathlineto{\pgfqpoint{0.893132in}{0.787776in}}%
\pgfpathlineto{\pgfqpoint{0.895613in}{1.758721in}}%
\pgfpathlineto{\pgfqpoint{0.898950in}{2.773699in}}%
\pgfpathlineto{\pgfqpoint{0.899292in}{2.785286in}}%
\pgfpathlineto{\pgfqpoint{0.899805in}{2.762892in}}%
\pgfpathlineto{\pgfqpoint{0.900918in}{2.558364in}}%
\pgfpathlineto{\pgfqpoint{0.903485in}{1.547264in}}%
\pgfpathlineto{\pgfqpoint{0.906736in}{0.583364in}}%
\pgfpathlineto{\pgfqpoint{0.907078in}{0.572737in}}%
\pgfpathlineto{\pgfqpoint{0.907592in}{0.596194in}}%
\pgfpathlineto{\pgfqpoint{0.908704in}{0.801375in}}%
\pgfpathlineto{\pgfqpoint{0.911271in}{1.808077in}}%
\pgfpathlineto{\pgfqpoint{0.914522in}{2.772697in}}%
\pgfpathlineto{\pgfqpoint{0.914864in}{2.784544in}}%
\pgfpathlineto{\pgfqpoint{0.915378in}{2.763282in}}%
\pgfpathlineto{\pgfqpoint{0.916490in}{2.563853in}}%
\pgfpathlineto{\pgfqpoint{0.919057in}{1.567720in}}%
\pgfpathlineto{\pgfqpoint{0.922394in}{0.583527in}}%
\pgfpathlineto{\pgfqpoint{0.922736in}{0.573466in}}%
\pgfpathlineto{\pgfqpoint{0.923164in}{0.589899in}}%
\pgfpathlineto{\pgfqpoint{0.924191in}{0.755867in}}%
\pgfpathlineto{\pgfqpoint{0.926501in}{1.605385in}}%
\pgfpathlineto{\pgfqpoint{0.930266in}{2.773363in}}%
\pgfpathlineto{\pgfqpoint{0.930608in}{2.783852in}}%
\pgfpathlineto{\pgfqpoint{0.931121in}{2.761300in}}%
\pgfpathlineto{\pgfqpoint{0.932234in}{2.562200in}}%
\pgfpathlineto{\pgfqpoint{0.934800in}{1.576535in}}%
\pgfpathlineto{\pgfqpoint{0.938137in}{0.587525in}}%
\pgfpathlineto{\pgfqpoint{0.938565in}{0.574324in}}%
\pgfpathlineto{\pgfqpoint{0.938993in}{0.592723in}}%
\pgfpathlineto{\pgfqpoint{0.940105in}{0.781960in}}%
\pgfpathlineto{\pgfqpoint{0.942587in}{1.715325in}}%
\pgfpathlineto{\pgfqpoint{0.946095in}{2.770058in}}%
\pgfpathlineto{\pgfqpoint{0.946523in}{2.782967in}}%
\pgfpathlineto{\pgfqpoint{0.946950in}{2.764587in}}%
\pgfpathlineto{\pgfqpoint{0.948063in}{2.576757in}}%
\pgfpathlineto{\pgfqpoint{0.950544in}{1.649814in}}%
\pgfpathlineto{\pgfqpoint{0.954052in}{0.590368in}}%
\pgfpathlineto{\pgfqpoint{0.954480in}{0.574995in}}%
\pgfpathlineto{\pgfqpoint{0.954993in}{0.597424in}}%
\pgfpathlineto{\pgfqpoint{0.956106in}{0.791982in}}%
\pgfpathlineto{\pgfqpoint{0.958672in}{1.757765in}}%
\pgfpathlineto{\pgfqpoint{0.962095in}{2.767862in}}%
\pgfpathlineto{\pgfqpoint{0.962523in}{2.782283in}}%
\pgfpathlineto{\pgfqpoint{0.963036in}{2.759122in}}%
\pgfpathlineto{\pgfqpoint{0.964234in}{2.542113in}}%
\pgfpathlineto{\pgfqpoint{0.966972in}{1.494504in}}%
\pgfpathlineto{\pgfqpoint{0.970223in}{0.586401in}}%
\pgfpathlineto{\pgfqpoint{0.970566in}{0.575818in}}%
\pgfpathlineto{\pgfqpoint{0.971079in}{0.596378in}}%
\pgfpathlineto{\pgfqpoint{0.972191in}{0.784374in}}%
\pgfpathlineto{\pgfqpoint{0.974673in}{1.696524in}}%
\pgfpathlineto{\pgfqpoint{0.978266in}{2.766647in}}%
\pgfpathlineto{\pgfqpoint{0.978694in}{2.781478in}}%
\pgfpathlineto{\pgfqpoint{0.979207in}{2.759619in}}%
\pgfpathlineto{\pgfqpoint{0.980320in}{2.570482in}}%
\pgfpathlineto{\pgfqpoint{0.982801in}{1.662142in}}%
\pgfpathlineto{\pgfqpoint{0.986395in}{0.592884in}}%
\pgfpathlineto{\pgfqpoint{0.986822in}{0.576661in}}%
\pgfpathlineto{\pgfqpoint{0.987336in}{0.596449in}}%
\pgfpathlineto{\pgfqpoint{0.988448in}{0.780046in}}%
\pgfpathlineto{\pgfqpoint{0.990929in}{1.677467in}}%
\pgfpathlineto{\pgfqpoint{0.994609in}{2.766125in}}%
\pgfpathlineto{\pgfqpoint{0.995036in}{2.780632in}}%
\pgfpathlineto{\pgfqpoint{0.995550in}{2.759196in}}%
\pgfpathlineto{\pgfqpoint{0.996662in}{2.573761in}}%
\pgfpathlineto{\pgfqpoint{0.999143in}{1.679411in}}%
\pgfpathlineto{\pgfqpoint{1.002823in}{0.593147in}}%
\pgfpathlineto{\pgfqpoint{1.003250in}{0.577510in}}%
\pgfpathlineto{\pgfqpoint{1.003764in}{0.597190in}}%
\pgfpathlineto{\pgfqpoint{1.004876in}{0.777706in}}%
\pgfpathlineto{\pgfqpoint{1.007357in}{1.661664in}}%
\pgfpathlineto{\pgfqpoint{1.011037in}{2.759741in}}%
\pgfpathlineto{\pgfqpoint{1.011550in}{2.779753in}}%
\pgfpathlineto{\pgfqpoint{1.012063in}{2.758287in}}%
\pgfpathlineto{\pgfqpoint{1.013261in}{2.554522in}}%
\pgfpathlineto{\pgfqpoint{1.015914in}{1.587857in}}%
\pgfpathlineto{\pgfqpoint{1.019422in}{0.593232in}}%
\pgfpathlineto{\pgfqpoint{1.019850in}{0.578385in}}%
\pgfpathlineto{\pgfqpoint{1.020363in}{0.598202in}}%
\pgfpathlineto{\pgfqpoint{1.021475in}{0.776128in}}%
\pgfpathlineto{\pgfqpoint{1.023871in}{1.611705in}}%
\pgfpathlineto{\pgfqpoint{1.027721in}{2.759779in}}%
\pgfpathlineto{\pgfqpoint{1.028235in}{2.778837in}}%
\pgfpathlineto{\pgfqpoint{1.028748in}{2.757301in}}%
\pgfpathlineto{\pgfqpoint{1.029946in}{2.556603in}}%
\pgfpathlineto{\pgfqpoint{1.032598in}{1.603912in}}%
\pgfpathlineto{\pgfqpoint{1.036106in}{0.599658in}}%
\pgfpathlineto{\pgfqpoint{1.036620in}{0.579298in}}%
\pgfpathlineto{\pgfqpoint{1.037133in}{0.599094in}}%
\pgfpathlineto{\pgfqpoint{1.038245in}{0.774118in}}%
\pgfpathlineto{\pgfqpoint{1.040641in}{1.596618in}}%
\pgfpathlineto{\pgfqpoint{1.044577in}{2.759437in}}%
\pgfpathlineto{\pgfqpoint{1.045090in}{2.777898in}}%
\pgfpathlineto{\pgfqpoint{1.045604in}{2.756645in}}%
\pgfpathlineto{\pgfqpoint{1.046802in}{2.559764in}}%
\pgfpathlineto{\pgfqpoint{1.049369in}{1.656677in}}%
\pgfpathlineto{\pgfqpoint{1.053048in}{0.600296in}}%
\pgfpathlineto{\pgfqpoint{1.053561in}{0.580248in}}%
\pgfpathlineto{\pgfqpoint{1.054075in}{0.599477in}}%
\pgfpathlineto{\pgfqpoint{1.055187in}{0.770511in}}%
\pgfpathlineto{\pgfqpoint{1.057583in}{1.578582in}}%
\pgfpathlineto{\pgfqpoint{1.061604in}{2.758349in}}%
\pgfpathlineto{\pgfqpoint{1.062117in}{2.776950in}}%
\pgfpathlineto{\pgfqpoint{1.062631in}{2.756714in}}%
\pgfpathlineto{\pgfqpoint{1.063829in}{2.565212in}}%
\pgfpathlineto{\pgfqpoint{1.066396in}{1.678358in}}%
\pgfpathlineto{\pgfqpoint{1.070160in}{0.601911in}}%
\pgfpathlineto{\pgfqpoint{1.070674in}{0.581253in}}%
\pgfpathlineto{\pgfqpoint{1.071187in}{0.598996in}}%
\pgfpathlineto{\pgfqpoint{1.072299in}{0.764201in}}%
\pgfpathlineto{\pgfqpoint{1.074695in}{1.555538in}}%
\pgfpathlineto{\pgfqpoint{1.078802in}{2.756111in}}%
\pgfpathlineto{\pgfqpoint{1.079315in}{2.775953in}}%
\pgfpathlineto{\pgfqpoint{1.079829in}{2.757833in}}%
\pgfpathlineto{\pgfqpoint{1.080941in}{2.593300in}}%
\pgfpathlineto{\pgfqpoint{1.083337in}{1.807758in}}%
\pgfpathlineto{\pgfqpoint{1.087444in}{0.604971in}}%
\pgfpathlineto{\pgfqpoint{1.088043in}{0.582320in}}%
\pgfpathlineto{\pgfqpoint{1.088556in}{0.603534in}}%
\pgfpathlineto{\pgfqpoint{1.089754in}{0.792307in}}%
\pgfpathlineto{\pgfqpoint{1.092321in}{1.659995in}}%
\pgfpathlineto{\pgfqpoint{1.096171in}{2.752190in}}%
\pgfpathlineto{\pgfqpoint{1.096770in}{2.774869in}}%
\pgfpathlineto{\pgfqpoint{1.097284in}{2.754150in}}%
\pgfpathlineto{\pgfqpoint{1.098481in}{2.568038in}}%
\pgfpathlineto{\pgfqpoint{1.101048in}{1.708626in}}%
\pgfpathlineto{\pgfqpoint{1.104899in}{0.610135in}}%
\pgfpathlineto{\pgfqpoint{1.105498in}{0.583320in}}%
\pgfpathlineto{\pgfqpoint{1.106011in}{0.600019in}}%
\pgfpathlineto{\pgfqpoint{1.107123in}{0.757420in}}%
\pgfpathlineto{\pgfqpoint{1.109433in}{1.486388in}}%
\pgfpathlineto{\pgfqpoint{1.113797in}{2.752781in}}%
\pgfpathlineto{\pgfqpoint{1.114311in}{2.773756in}}%
\pgfpathlineto{\pgfqpoint{1.114824in}{2.758518in}}%
\pgfpathlineto{\pgfqpoint{1.115936in}{2.605484in}}%
\pgfpathlineto{\pgfqpoint{1.118246in}{1.886661in}}%
\pgfpathlineto{\pgfqpoint{1.122696in}{0.603990in}}%
\pgfpathlineto{\pgfqpoint{1.123209in}{0.584405in}}%
\pgfpathlineto{\pgfqpoint{1.123722in}{0.600599in}}%
\pgfpathlineto{\pgfqpoint{1.124835in}{0.754136in}}%
\pgfpathlineto{\pgfqpoint{1.127145in}{1.468717in}}%
\pgfpathlineto{\pgfqpoint{1.131594in}{2.750734in}}%
\pgfpathlineto{\pgfqpoint{1.132193in}{2.772728in}}%
\pgfpathlineto{\pgfqpoint{1.132706in}{2.753310in}}%
\pgfpathlineto{\pgfqpoint{1.133904in}{2.576382in}}%
\pgfpathlineto{\pgfqpoint{1.136386in}{1.781641in}}%
\pgfpathlineto{\pgfqpoint{1.140493in}{0.614472in}}%
\pgfpathlineto{\pgfqpoint{1.141177in}{0.585514in}}%
\pgfpathlineto{\pgfqpoint{1.141690in}{0.604484in}}%
\pgfpathlineto{\pgfqpoint{1.142888in}{0.778849in}}%
\pgfpathlineto{\pgfqpoint{1.145370in}{1.565395in}}%
\pgfpathlineto{\pgfqpoint{1.149562in}{2.745417in}}%
\pgfpathlineto{\pgfqpoint{1.150161in}{2.771559in}}%
\pgfpathlineto{\pgfqpoint{1.150675in}{2.756640in}}%
\pgfpathlineto{\pgfqpoint{1.151787in}{2.609945in}}%
\pgfpathlineto{\pgfqpoint{1.154097in}{1.918490in}}%
\pgfpathlineto{\pgfqpoint{1.158717in}{0.607914in}}%
\pgfpathlineto{\pgfqpoint{1.159316in}{0.586674in}}%
\pgfpathlineto{\pgfqpoint{1.159830in}{0.605331in}}%
\pgfpathlineto{\pgfqpoint{1.161028in}{0.775830in}}%
\pgfpathlineto{\pgfqpoint{1.163509in}{1.547267in}}%
\pgfpathlineto{\pgfqpoint{1.167787in}{2.743058in}}%
\pgfpathlineto{\pgfqpoint{1.168472in}{2.770414in}}%
\pgfpathlineto{\pgfqpoint{1.168985in}{2.751756in}}%
\pgfpathlineto{\pgfqpoint{1.170183in}{2.582852in}}%
\pgfpathlineto{\pgfqpoint{1.172664in}{1.818635in}}%
\pgfpathlineto{\pgfqpoint{1.177028in}{0.612275in}}%
\pgfpathlineto{\pgfqpoint{1.177627in}{0.587863in}}%
\pgfpathlineto{\pgfqpoint{1.178140in}{0.602863in}}%
\pgfpathlineto{\pgfqpoint{1.179252in}{0.745461in}}%
\pgfpathlineto{\pgfqpoint{1.181563in}{1.416977in}}%
\pgfpathlineto{\pgfqpoint{1.186354in}{2.749556in}}%
\pgfpathlineto{\pgfqpoint{1.186953in}{2.769175in}}%
\pgfpathlineto{\pgfqpoint{1.187466in}{2.750532in}}%
\pgfpathlineto{\pgfqpoint{1.188664in}{2.584849in}}%
\pgfpathlineto{\pgfqpoint{1.191146in}{1.835193in}}%
\pgfpathlineto{\pgfqpoint{1.195595in}{0.614699in}}%
\pgfpathlineto{\pgfqpoint{1.196279in}{0.589142in}}%
\pgfpathlineto{\pgfqpoint{1.196793in}{0.607642in}}%
\pgfpathlineto{\pgfqpoint{1.197991in}{0.771425in}}%
\pgfpathlineto{\pgfqpoint{1.200472in}{1.513419in}}%
\pgfpathlineto{\pgfqpoint{1.204921in}{2.737969in}}%
\pgfpathlineto{\pgfqpoint{1.205606in}{2.767977in}}%
\pgfpathlineto{\pgfqpoint{1.206119in}{2.753325in}}%
\pgfpathlineto{\pgfqpoint{1.207231in}{2.615678in}}%
\pgfpathlineto{\pgfqpoint{1.209456in}{1.995196in}}%
\pgfpathlineto{\pgfqpoint{1.214504in}{0.609551in}}%
\pgfpathlineto{\pgfqpoint{1.215103in}{0.590387in}}%
\pgfpathlineto{\pgfqpoint{1.215616in}{0.608023in}}%
\pgfpathlineto{\pgfqpoint{1.216814in}{0.766761in}}%
\pgfpathlineto{\pgfqpoint{1.219210in}{1.461601in}}%
\pgfpathlineto{\pgfqpoint{1.223916in}{2.740859in}}%
\pgfpathlineto{\pgfqpoint{1.224601in}{2.766683in}}%
\pgfpathlineto{\pgfqpoint{1.225114in}{2.749886in}}%
\pgfpathlineto{\pgfqpoint{1.226312in}{2.594551in}}%
\pgfpathlineto{\pgfqpoint{1.228708in}{1.909014in}}%
\pgfpathlineto{\pgfqpoint{1.233499in}{0.616193in}}%
\pgfpathlineto{\pgfqpoint{1.234184in}{0.591709in}}%
\pgfpathlineto{\pgfqpoint{1.234697in}{0.609014in}}%
\pgfpathlineto{\pgfqpoint{1.235895in}{0.763862in}}%
\pgfpathlineto{\pgfqpoint{1.238291in}{1.443545in}}%
\pgfpathlineto{\pgfqpoint{1.243082in}{2.737166in}}%
\pgfpathlineto{\pgfqpoint{1.243766in}{2.765355in}}%
\pgfpathlineto{\pgfqpoint{1.244280in}{2.751334in}}%
\pgfpathlineto{\pgfqpoint{1.245392in}{2.620612in}}%
\pgfpathlineto{\pgfqpoint{1.247617in}{2.028084in}}%
\pgfpathlineto{\pgfqpoint{1.252922in}{0.610909in}}%
\pgfpathlineto{\pgfqpoint{1.253521in}{0.593073in}}%
\pgfpathlineto{\pgfqpoint{1.254034in}{0.609993in}}%
\pgfpathlineto{\pgfqpoint{1.255232in}{0.760839in}}%
\pgfpathlineto{\pgfqpoint{1.257628in}{1.425108in}}%
\pgfpathlineto{\pgfqpoint{1.262590in}{2.739630in}}%
\pgfpathlineto{\pgfqpoint{1.263275in}{2.763965in}}%
\pgfpathlineto{\pgfqpoint{1.263788in}{2.748041in}}%
\pgfpathlineto{\pgfqpoint{1.264986in}{2.600940in}}%
\pgfpathlineto{\pgfqpoint{1.267382in}{1.946576in}}%
\pgfpathlineto{\pgfqpoint{1.272430in}{0.618117in}}%
\pgfpathlineto{\pgfqpoint{1.273114in}{0.594444in}}%
\pgfpathlineto{\pgfqpoint{1.273628in}{0.610369in}}%
\pgfpathlineto{\pgfqpoint{1.274826in}{0.755892in}}%
\pgfpathlineto{\pgfqpoint{1.277221in}{1.402978in}}%
\pgfpathlineto{\pgfqpoint{1.282270in}{2.734348in}}%
\pgfpathlineto{\pgfqpoint{1.282954in}{2.762483in}}%
\pgfpathlineto{\pgfqpoint{1.283467in}{2.750408in}}%
\pgfpathlineto{\pgfqpoint{1.284580in}{2.629280in}}%
\pgfpathlineto{\pgfqpoint{1.286804in}{2.068696in}}%
\pgfpathlineto{\pgfqpoint{1.292452in}{0.610110in}}%
\pgfpathlineto{\pgfqpoint{1.292965in}{0.595859in}}%
\pgfpathlineto{\pgfqpoint{1.293478in}{0.609651in}}%
\pgfpathlineto{\pgfqpoint{1.294591in}{0.732887in}}%
\pgfpathlineto{\pgfqpoint{1.296815in}{1.291555in}}%
\pgfpathlineto{\pgfqpoint{1.302462in}{2.744945in}}%
\pgfpathlineto{\pgfqpoint{1.302976in}{2.760997in}}%
\pgfpathlineto{\pgfqpoint{1.303489in}{2.749430in}}%
\pgfpathlineto{\pgfqpoint{1.304601in}{2.632142in}}%
\pgfpathlineto{\pgfqpoint{1.306740in}{2.112913in}}%
\pgfpathlineto{\pgfqpoint{1.312644in}{0.610026in}}%
\pgfpathlineto{\pgfqpoint{1.313158in}{0.597376in}}%
\pgfpathlineto{\pgfqpoint{1.313671in}{0.611922in}}%
\pgfpathlineto{\pgfqpoint{1.314869in}{0.748131in}}%
\pgfpathlineto{\pgfqpoint{1.317179in}{1.334871in}}%
\pgfpathlineto{\pgfqpoint{1.322655in}{2.735947in}}%
\pgfpathlineto{\pgfqpoint{1.323340in}{2.759560in}}%
\pgfpathlineto{\pgfqpoint{1.323853in}{2.746057in}}%
\pgfpathlineto{\pgfqpoint{1.324965in}{2.627598in}}%
\pgfpathlineto{\pgfqpoint{1.327190in}{2.089857in}}%
\pgfpathlineto{\pgfqpoint{1.333094in}{0.612343in}}%
\pgfpathlineto{\pgfqpoint{1.333607in}{0.598914in}}%
\pgfpathlineto{\pgfqpoint{1.334120in}{0.611840in}}%
\pgfpathlineto{\pgfqpoint{1.335233in}{0.727724in}}%
\pgfpathlineto{\pgfqpoint{1.337457in}{1.257083in}}%
\pgfpathlineto{\pgfqpoint{1.343447in}{2.744753in}}%
\pgfpathlineto{\pgfqpoint{1.343960in}{2.757979in}}%
\pgfpathlineto{\pgfqpoint{1.344474in}{2.745271in}}%
\pgfpathlineto{\pgfqpoint{1.345586in}{2.631222in}}%
\pgfpathlineto{\pgfqpoint{1.347810in}{2.109246in}}%
\pgfpathlineto{\pgfqpoint{1.353885in}{0.613287in}}%
\pgfpathlineto{\pgfqpoint{1.354399in}{0.600520in}}%
\pgfpathlineto{\pgfqpoint{1.354912in}{0.613270in}}%
\pgfpathlineto{\pgfqpoint{1.356024in}{0.726015in}}%
\pgfpathlineto{\pgfqpoint{1.358249in}{1.241339in}}%
\pgfpathlineto{\pgfqpoint{1.364410in}{2.744195in}}%
\pgfpathlineto{\pgfqpoint{1.364923in}{2.756342in}}%
\pgfpathlineto{\pgfqpoint{1.365436in}{2.743392in}}%
\pgfpathlineto{\pgfqpoint{1.366549in}{2.631620in}}%
\pgfpathlineto{\pgfqpoint{1.368773in}{2.122494in}}%
\pgfpathlineto{\pgfqpoint{1.375019in}{0.613664in}}%
\pgfpathlineto{\pgfqpoint{1.375533in}{0.602200in}}%
\pgfpathlineto{\pgfqpoint{1.376046in}{0.615414in}}%
\pgfpathlineto{\pgfqpoint{1.377158in}{0.726344in}}%
\pgfpathlineto{\pgfqpoint{1.379383in}{1.229452in}}%
\pgfpathlineto{\pgfqpoint{1.385715in}{2.743803in}}%
\pgfpathlineto{\pgfqpoint{1.386228in}{2.754616in}}%
\pgfpathlineto{\pgfqpoint{1.386741in}{2.741169in}}%
\pgfpathlineto{\pgfqpoint{1.387939in}{2.618320in}}%
\pgfpathlineto{\pgfqpoint{1.390249in}{2.084913in}}%
\pgfpathlineto{\pgfqpoint{1.396324in}{0.622955in}}%
\pgfpathlineto{\pgfqpoint{1.397009in}{0.603964in}}%
\pgfpathlineto{\pgfqpoint{1.397522in}{0.617524in}}%
\pgfpathlineto{\pgfqpoint{1.398720in}{0.739066in}}%
\pgfpathlineto{\pgfqpoint{1.401030in}{1.265561in}}%
\pgfpathlineto{\pgfqpoint{1.407191in}{2.734359in}}%
\pgfpathlineto{\pgfqpoint{1.407875in}{2.752833in}}%
\pgfpathlineto{\pgfqpoint{1.408389in}{2.739368in}}%
\pgfpathlineto{\pgfqpoint{1.409587in}{2.619588in}}%
\pgfpathlineto{\pgfqpoint{1.411897in}{2.100630in}}%
\pgfpathlineto{\pgfqpoint{1.418143in}{0.624095in}}%
\pgfpathlineto{\pgfqpoint{1.418827in}{0.605745in}}%
\pgfpathlineto{\pgfqpoint{1.419341in}{0.618822in}}%
\pgfpathlineto{\pgfqpoint{1.420539in}{0.736197in}}%
\pgfpathlineto{\pgfqpoint{1.422849in}{1.246728in}}%
\pgfpathlineto{\pgfqpoint{1.429180in}{2.732334in}}%
\pgfpathlineto{\pgfqpoint{1.429865in}{2.751061in}}%
\pgfpathlineto{\pgfqpoint{1.430378in}{2.738747in}}%
\pgfpathlineto{\pgfqpoint{1.431576in}{2.624601in}}%
\pgfpathlineto{\pgfqpoint{1.433801in}{2.147076in}}%
\pgfpathlineto{\pgfqpoint{1.440475in}{0.618638in}}%
\pgfpathlineto{\pgfqpoint{1.440988in}{0.607540in}}%
\pgfpathlineto{\pgfqpoint{1.441501in}{0.618637in}}%
\pgfpathlineto{\pgfqpoint{1.442614in}{0.716964in}}%
\pgfpathlineto{\pgfqpoint{1.444753in}{1.150268in}}%
\pgfpathlineto{\pgfqpoint{1.451940in}{2.745649in}}%
\pgfpathlineto{\pgfqpoint{1.452197in}{2.749147in}}%
\pgfpathlineto{\pgfqpoint{1.452710in}{2.739795in}}%
\pgfpathlineto{\pgfqpoint{1.453822in}{2.646445in}}%
\pgfpathlineto{\pgfqpoint{1.455961in}{2.224960in}}%
\pgfpathlineto{\pgfqpoint{1.463320in}{0.612194in}}%
\pgfpathlineto{\pgfqpoint{1.463576in}{0.609457in}}%
\pgfpathlineto{\pgfqpoint{1.464004in}{0.616778in}}%
\pgfpathlineto{\pgfqpoint{1.465031in}{0.693813in}}%
\pgfpathlineto{\pgfqpoint{1.466999in}{1.049114in}}%
\pgfpathlineto{\pgfqpoint{1.475042in}{2.747244in}}%
\pgfpathlineto{\pgfqpoint{1.475127in}{2.746871in}}%
\pgfpathlineto{\pgfqpoint{1.475726in}{2.727993in}}%
\pgfpathlineto{\pgfqpoint{1.477010in}{2.595088in}}%
\pgfpathlineto{\pgfqpoint{1.479491in}{2.050002in}}%
\pgfpathlineto{\pgfqpoint{1.485737in}{0.639563in}}%
\pgfpathlineto{\pgfqpoint{1.486593in}{0.611445in}}%
\pgfpathlineto{\pgfqpoint{1.487106in}{0.621948in}}%
\pgfpathlineto{\pgfqpoint{1.488218in}{0.713601in}}%
\pgfpathlineto{\pgfqpoint{1.490357in}{1.118797in}}%
\pgfpathlineto{\pgfqpoint{1.497973in}{2.742196in}}%
\pgfpathlineto{\pgfqpoint{1.498229in}{2.745201in}}%
\pgfpathlineto{\pgfqpoint{1.498743in}{2.736090in}}%
\pgfpathlineto{\pgfqpoint{1.499855in}{2.648678in}}%
\pgfpathlineto{\pgfqpoint{1.501994in}{2.254362in}}%
\pgfpathlineto{\pgfqpoint{1.509866in}{0.614557in}}%
\pgfpathlineto{\pgfqpoint{1.510037in}{0.613508in}}%
\pgfpathlineto{\pgfqpoint{1.510379in}{0.617999in}}%
\pgfpathlineto{\pgfqpoint{1.511320in}{0.675021in}}%
\pgfpathlineto{\pgfqpoint{1.513117in}{0.950079in}}%
\pgfpathlineto{\pgfqpoint{1.517310in}{2.048374in}}%
\pgfpathlineto{\pgfqpoint{1.520903in}{2.704747in}}%
\pgfpathlineto{\pgfqpoint{1.521930in}{2.743122in}}%
\pgfpathlineto{\pgfqpoint{1.522358in}{2.736331in}}%
\pgfpathlineto{\pgfqpoint{1.523385in}{2.666154in}}%
\pgfpathlineto{\pgfqpoint{1.525353in}{2.341473in}}%
\pgfpathlineto{\pgfqpoint{1.533909in}{0.615687in}}%
\pgfpathlineto{\pgfqpoint{1.534251in}{0.619101in}}%
\pgfpathlineto{\pgfqpoint{1.535107in}{0.664108in}}%
\pgfpathlineto{\pgfqpoint{1.536818in}{0.899129in}}%
\pgfpathlineto{\pgfqpoint{1.540326in}{1.769808in}}%
\pgfpathlineto{\pgfqpoint{1.544690in}{2.673785in}}%
\pgfpathlineto{\pgfqpoint{1.546059in}{2.740927in}}%
\pgfpathlineto{\pgfqpoint{1.546401in}{2.737243in}}%
\pgfpathlineto{\pgfqpoint{1.547342in}{2.685110in}}%
\pgfpathlineto{\pgfqpoint{1.549139in}{2.428248in}}%
\pgfpathlineto{\pgfqpoint{1.552989in}{1.463844in}}%
\pgfpathlineto{\pgfqpoint{1.557011in}{0.678108in}}%
\pgfpathlineto{\pgfqpoint{1.558380in}{0.617935in}}%
\pgfpathlineto{\pgfqpoint{1.558722in}{0.622945in}}%
\pgfpathlineto{\pgfqpoint{1.559749in}{0.685619in}}%
\pgfpathlineto{\pgfqpoint{1.561631in}{0.969148in}}%
\pgfpathlineto{\pgfqpoint{1.566166in}{2.108651in}}%
\pgfpathlineto{\pgfqpoint{1.569674in}{2.701344in}}%
\pgfpathlineto{\pgfqpoint{1.570701in}{2.738642in}}%
\pgfpathlineto{\pgfqpoint{1.571214in}{2.730747in}}%
\pgfpathlineto{\pgfqpoint{1.572326in}{2.653930in}}%
\pgfpathlineto{\pgfqpoint{1.574380in}{2.321540in}}%
\pgfpathlineto{\pgfqpoint{1.583278in}{0.620212in}}%
\pgfpathlineto{\pgfqpoint{1.583449in}{0.621466in}}%
\pgfpathlineto{\pgfqpoint{1.584220in}{0.650815in}}%
\pgfpathlineto{\pgfqpoint{1.585760in}{0.820160in}}%
\pgfpathlineto{\pgfqpoint{1.588669in}{1.444795in}}%
\pgfpathlineto{\pgfqpoint{1.594573in}{2.678846in}}%
\pgfpathlineto{\pgfqpoint{1.595942in}{2.736283in}}%
\pgfpathlineto{\pgfqpoint{1.596284in}{2.731904in}}%
\pgfpathlineto{\pgfqpoint{1.597225in}{2.681503in}}%
\pgfpathlineto{\pgfqpoint{1.599022in}{2.441072in}}%
\pgfpathlineto{\pgfqpoint{1.602787in}{1.552648in}}%
\pgfpathlineto{\pgfqpoint{1.607236in}{0.690885in}}%
\pgfpathlineto{\pgfqpoint{1.608690in}{0.622621in}}%
\pgfpathlineto{\pgfqpoint{1.609033in}{0.625739in}}%
\pgfpathlineto{\pgfqpoint{1.609974in}{0.671862in}}%
\pgfpathlineto{\pgfqpoint{1.611771in}{0.901811in}}%
\pgfpathlineto{\pgfqpoint{1.615364in}{1.727588in}}%
\pgfpathlineto{\pgfqpoint{1.620070in}{2.656599in}}%
\pgfpathlineto{\pgfqpoint{1.621696in}{2.733821in}}%
\pgfpathlineto{\pgfqpoint{1.621953in}{2.731317in}}%
\pgfpathlineto{\pgfqpoint{1.622808in}{2.694045in}}%
\pgfpathlineto{\pgfqpoint{1.624519in}{2.494133in}}%
\pgfpathlineto{\pgfqpoint{1.627771in}{1.786555in}}%
\pgfpathlineto{\pgfqpoint{1.633161in}{0.701295in}}%
\pgfpathlineto{\pgfqpoint{1.634787in}{0.625123in}}%
\pgfpathlineto{\pgfqpoint{1.635044in}{0.627437in}}%
\pgfpathlineto{\pgfqpoint{1.635899in}{0.663383in}}%
\pgfpathlineto{\pgfqpoint{1.637610in}{0.857883in}}%
\pgfpathlineto{\pgfqpoint{1.640862in}{1.551250in}}%
\pgfpathlineto{\pgfqpoint{1.646423in}{2.657602in}}%
\pgfpathlineto{\pgfqpoint{1.648049in}{2.731263in}}%
\pgfpathlineto{\pgfqpoint{1.648306in}{2.728903in}}%
\pgfpathlineto{\pgfqpoint{1.649161in}{2.693503in}}%
\pgfpathlineto{\pgfqpoint{1.650873in}{2.503022in}}%
\pgfpathlineto{\pgfqpoint{1.654038in}{1.843452in}}%
\pgfpathlineto{\pgfqpoint{1.659771in}{0.705789in}}%
\pgfpathlineto{\pgfqpoint{1.661397in}{0.627768in}}%
\pgfpathlineto{\pgfqpoint{1.661654in}{0.629068in}}%
\pgfpathlineto{\pgfqpoint{1.662424in}{0.655329in}}%
\pgfpathlineto{\pgfqpoint{1.663964in}{0.803973in}}%
\pgfpathlineto{\pgfqpoint{1.666787in}{1.339531in}}%
\pgfpathlineto{\pgfqpoint{1.673632in}{2.675665in}}%
\pgfpathlineto{\pgfqpoint{1.675001in}{2.728607in}}%
\pgfpathlineto{\pgfqpoint{1.675344in}{2.725699in}}%
\pgfpathlineto{\pgfqpoint{1.676285in}{2.684563in}}%
\pgfpathlineto{\pgfqpoint{1.678082in}{2.479830in}}%
\pgfpathlineto{\pgfqpoint{1.681504in}{1.768027in}}%
\pgfpathlineto{\pgfqpoint{1.686980in}{0.716327in}}%
\pgfpathlineto{\pgfqpoint{1.688777in}{0.630454in}}%
\pgfpathlineto{\pgfqpoint{1.688948in}{0.631268in}}%
\pgfpathlineto{\pgfqpoint{1.689632in}{0.650253in}}%
\pgfpathlineto{\pgfqpoint{1.691087in}{0.771238in}}%
\pgfpathlineto{\pgfqpoint{1.693739in}{1.225814in}}%
\pgfpathlineto{\pgfqpoint{1.701697in}{2.699038in}}%
\pgfpathlineto{\pgfqpoint{1.702724in}{2.725820in}}%
\pgfpathlineto{\pgfqpoint{1.703151in}{2.720678in}}%
\pgfpathlineto{\pgfqpoint{1.704178in}{2.669682in}}%
\pgfpathlineto{\pgfqpoint{1.706146in}{2.431446in}}%
\pgfpathlineto{\pgfqpoint{1.709996in}{1.620684in}}%
\pgfpathlineto{\pgfqpoint{1.714959in}{0.721379in}}%
\pgfpathlineto{\pgfqpoint{1.716841in}{0.633312in}}%
\pgfpathlineto{\pgfqpoint{1.717012in}{0.634180in}}%
\pgfpathlineto{\pgfqpoint{1.717697in}{0.652532in}}%
\pgfpathlineto{\pgfqpoint{1.719152in}{0.767905in}}%
\pgfpathlineto{\pgfqpoint{1.721718in}{1.184444in}}%
\pgfpathlineto{\pgfqpoint{1.730275in}{2.704404in}}%
\pgfpathlineto{\pgfqpoint{1.731130in}{2.722906in}}%
\pgfpathlineto{\pgfqpoint{1.731644in}{2.716602in}}%
\pgfpathlineto{\pgfqpoint{1.732756in}{2.658821in}}%
\pgfpathlineto{\pgfqpoint{1.734809in}{2.407250in}}%
\pgfpathlineto{\pgfqpoint{1.738916in}{1.556065in}}%
\pgfpathlineto{\pgfqpoint{1.743708in}{0.726053in}}%
\pgfpathlineto{\pgfqpoint{1.745676in}{0.636307in}}%
\pgfpathlineto{\pgfqpoint{1.745761in}{0.636589in}}%
\pgfpathlineto{\pgfqpoint{1.746360in}{0.648402in}}%
\pgfpathlineto{\pgfqpoint{1.747644in}{0.730372in}}%
\pgfpathlineto{\pgfqpoint{1.749954in}{1.049811in}}%
\pgfpathlineto{\pgfqpoint{1.760307in}{2.719763in}}%
\pgfpathlineto{\pgfqpoint{1.760393in}{2.719841in}}%
\pgfpathlineto{\pgfqpoint{1.760564in}{2.718971in}}%
\pgfpathlineto{\pgfqpoint{1.761334in}{2.698211in}}%
\pgfpathlineto{\pgfqpoint{1.762874in}{2.577020in}}%
\pgfpathlineto{\pgfqpoint{1.765612in}{2.143754in}}%
\pgfpathlineto{\pgfqpoint{1.774168in}{0.668572in}}%
\pgfpathlineto{\pgfqpoint{1.775280in}{0.639432in}}%
\pgfpathlineto{\pgfqpoint{1.775708in}{0.643110in}}%
\pgfpathlineto{\pgfqpoint{1.776735in}{0.685375in}}%
\pgfpathlineto{\pgfqpoint{1.778617in}{0.877530in}}%
\pgfpathlineto{\pgfqpoint{1.782040in}{1.503462in}}%
\pgfpathlineto{\pgfqpoint{1.788372in}{2.621754in}}%
\pgfpathlineto{\pgfqpoint{1.790425in}{2.716616in}}%
\pgfpathlineto{\pgfqpoint{1.790511in}{2.716605in}}%
\pgfpathlineto{\pgfqpoint{1.790938in}{2.711735in}}%
\pgfpathlineto{\pgfqpoint{1.792051in}{2.662033in}}%
\pgfpathlineto{\pgfqpoint{1.794019in}{2.452562in}}%
\pgfpathlineto{\pgfqpoint{1.797698in}{1.771237in}}%
\pgfpathlineto{\pgfqpoint{1.803687in}{0.741592in}}%
\pgfpathlineto{\pgfqpoint{1.805826in}{0.642724in}}%
\pgfpathlineto{\pgfqpoint{1.805912in}{0.642757in}}%
\pgfpathlineto{\pgfqpoint{1.806340in}{0.647580in}}%
\pgfpathlineto{\pgfqpoint{1.807452in}{0.695969in}}%
\pgfpathlineto{\pgfqpoint{1.809420in}{0.899442in}}%
\pgfpathlineto{\pgfqpoint{1.813099in}{1.564155in}}%
\pgfpathlineto{\pgfqpoint{1.819260in}{2.611621in}}%
\pgfpathlineto{\pgfqpoint{1.821484in}{2.713282in}}%
\pgfpathlineto{\pgfqpoint{1.821741in}{2.712101in}}%
\pgfpathlineto{\pgfqpoint{1.822511in}{2.692416in}}%
\pgfpathlineto{\pgfqpoint{1.824051in}{2.582913in}}%
\pgfpathlineto{\pgfqpoint{1.826789in}{2.192948in}}%
\pgfpathlineto{\pgfqpoint{1.836458in}{0.663639in}}%
\pgfpathlineto{\pgfqpoint{1.837399in}{0.646150in}}%
\pgfpathlineto{\pgfqpoint{1.837912in}{0.651397in}}%
\pgfpathlineto{\pgfqpoint{1.839025in}{0.698118in}}%
\pgfpathlineto{\pgfqpoint{1.841078in}{0.902369in}}%
\pgfpathlineto{\pgfqpoint{1.844757in}{1.540533in}}%
\pgfpathlineto{\pgfqpoint{1.851174in}{2.602188in}}%
\pgfpathlineto{\pgfqpoint{1.853485in}{2.709630in}}%
\pgfpathlineto{\pgfqpoint{1.853570in}{2.709746in}}%
\pgfpathlineto{\pgfqpoint{1.853827in}{2.708416in}}%
\pgfpathlineto{\pgfqpoint{1.854682in}{2.685885in}}%
\pgfpathlineto{\pgfqpoint{1.856308in}{2.569281in}}%
\pgfpathlineto{\pgfqpoint{1.859132in}{2.174133in}}%
\pgfpathlineto{\pgfqpoint{1.868800in}{0.676693in}}%
\pgfpathlineto{\pgfqpoint{1.869998in}{0.649787in}}%
\pgfpathlineto{\pgfqpoint{1.870426in}{0.652974in}}%
\pgfpathlineto{\pgfqpoint{1.871453in}{0.687883in}}%
\pgfpathlineto{\pgfqpoint{1.873335in}{0.846363in}}%
\pgfpathlineto{\pgfqpoint{1.876586in}{1.345084in}}%
\pgfpathlineto{\pgfqpoint{1.884715in}{2.633008in}}%
\pgfpathlineto{\pgfqpoint{1.886768in}{2.706010in}}%
\pgfpathlineto{\pgfqpoint{1.886854in}{2.705831in}}%
\pgfpathlineto{\pgfqpoint{1.887453in}{2.697318in}}%
\pgfpathlineto{\pgfqpoint{1.888736in}{2.637061in}}%
\pgfpathlineto{\pgfqpoint{1.890961in}{2.408595in}}%
\pgfpathlineto{\pgfqpoint{1.895153in}{1.695634in}}%
\pgfpathlineto{\pgfqpoint{1.901143in}{0.772664in}}%
\pgfpathlineto{\pgfqpoint{1.903539in}{0.654837in}}%
\pgfpathlineto{\pgfqpoint{1.903795in}{0.653623in}}%
\pgfpathlineto{\pgfqpoint{1.904223in}{0.656597in}}%
\pgfpathlineto{\pgfqpoint{1.905250in}{0.688977in}}%
\pgfpathlineto{\pgfqpoint{1.907132in}{0.836154in}}%
\pgfpathlineto{\pgfqpoint{1.910384in}{1.303384in}}%
\pgfpathlineto{\pgfqpoint{1.919197in}{2.637721in}}%
\pgfpathlineto{\pgfqpoint{1.921164in}{2.702061in}}%
\pgfpathlineto{\pgfqpoint{1.921250in}{2.702001in}}%
\pgfpathlineto{\pgfqpoint{1.921763in}{2.696609in}}%
\pgfpathlineto{\pgfqpoint{1.922961in}{2.650935in}}%
\pgfpathlineto{\pgfqpoint{1.925100in}{2.462058in}}%
\pgfpathlineto{\pgfqpoint{1.928780in}{1.899846in}}%
\pgfpathlineto{\pgfqpoint{1.936223in}{0.769542in}}%
\pgfpathlineto{\pgfqpoint{1.938705in}{0.658355in}}%
\pgfpathlineto{\pgfqpoint{1.938876in}{0.657703in}}%
\pgfpathlineto{\pgfqpoint{1.939304in}{0.660096in}}%
\pgfpathlineto{\pgfqpoint{1.940330in}{0.689106in}}%
\pgfpathlineto{\pgfqpoint{1.942213in}{0.823593in}}%
\pgfpathlineto{\pgfqpoint{1.945379in}{1.243346in}}%
\pgfpathlineto{\pgfqpoint{1.955133in}{2.645604in}}%
\pgfpathlineto{\pgfqpoint{1.957015in}{2.697883in}}%
\pgfpathlineto{\pgfqpoint{1.957186in}{2.697378in}}%
\pgfpathlineto{\pgfqpoint{1.957871in}{2.686577in}}%
\pgfpathlineto{\pgfqpoint{1.959325in}{2.618000in}}%
\pgfpathlineto{\pgfqpoint{1.961807in}{2.371716in}}%
\pgfpathlineto{\pgfqpoint{1.966427in}{1.638830in}}%
\pgfpathlineto{\pgfqpoint{1.972416in}{0.794548in}}%
\pgfpathlineto{\pgfqpoint{1.975069in}{0.664431in}}%
\pgfpathlineto{\pgfqpoint{1.975497in}{0.661998in}}%
\pgfpathlineto{\pgfqpoint{1.976010in}{0.666024in}}%
\pgfpathlineto{\pgfqpoint{1.977122in}{0.700445in}}%
\pgfpathlineto{\pgfqpoint{1.979090in}{0.842928in}}%
\pgfpathlineto{\pgfqpoint{1.982427in}{1.275817in}}%
\pgfpathlineto{\pgfqpoint{1.992181in}{2.628306in}}%
\pgfpathlineto{\pgfqpoint{1.994406in}{2.693417in}}%
\pgfpathlineto{\pgfqpoint{1.994834in}{2.690473in}}%
\pgfpathlineto{\pgfqpoint{1.995946in}{2.659560in}}%
\pgfpathlineto{\pgfqpoint{1.997914in}{2.526552in}}%
\pgfpathlineto{\pgfqpoint{2.001165in}{2.127361in}}%
\pgfpathlineto{\pgfqpoint{2.011690in}{0.718780in}}%
\pgfpathlineto{\pgfqpoint{2.013657in}{0.666606in}}%
\pgfpathlineto{\pgfqpoint{2.013743in}{0.666620in}}%
\pgfpathlineto{\pgfqpoint{2.014256in}{0.670711in}}%
\pgfpathlineto{\pgfqpoint{2.015454in}{0.706696in}}%
\pgfpathlineto{\pgfqpoint{2.017508in}{0.849831in}}%
\pgfpathlineto{\pgfqpoint{2.020930in}{1.271895in}}%
\pgfpathlineto{\pgfqpoint{2.031112in}{2.619308in}}%
\pgfpathlineto{\pgfqpoint{2.033422in}{2.688637in}}%
\pgfpathlineto{\pgfqpoint{2.033594in}{2.688545in}}%
\pgfpathlineto{\pgfqpoint{2.033765in}{2.687726in}}%
\pgfpathlineto{\pgfqpoint{2.034620in}{2.672790in}}%
\pgfpathlineto{\pgfqpoint{2.036246in}{2.595948in}}%
\pgfpathlineto{\pgfqpoint{2.038984in}{2.338562in}}%
\pgfpathlineto{\pgfqpoint{2.044118in}{1.593819in}}%
\pgfpathlineto{\pgfqpoint{2.050193in}{0.819685in}}%
\pgfpathlineto{\pgfqpoint{2.053016in}{0.678322in}}%
\pgfpathlineto{\pgfqpoint{2.053786in}{0.671504in}}%
\pgfpathlineto{\pgfqpoint{2.054300in}{0.674708in}}%
\pgfpathlineto{\pgfqpoint{2.055412in}{0.702711in}}%
\pgfpathlineto{\pgfqpoint{2.057380in}{0.819701in}}%
\pgfpathlineto{\pgfqpoint{2.060631in}{1.171325in}}%
\pgfpathlineto{\pgfqpoint{2.072867in}{2.650100in}}%
\pgfpathlineto{\pgfqpoint{2.074578in}{2.683588in}}%
\pgfpathlineto{\pgfqpoint{2.074835in}{2.683006in}}%
\pgfpathlineto{\pgfqpoint{2.075605in}{2.672497in}}%
\pgfpathlineto{\pgfqpoint{2.077145in}{2.612875in}}%
\pgfpathlineto{\pgfqpoint{2.079712in}{2.408888in}}%
\pgfpathlineto{\pgfqpoint{2.084075in}{1.849312in}}%
\pgfpathlineto{\pgfqpoint{2.092204in}{0.825665in}}%
\pgfpathlineto{\pgfqpoint{2.095113in}{0.685201in}}%
\pgfpathlineto{\pgfqpoint{2.095969in}{0.676802in}}%
\pgfpathlineto{\pgfqpoint{2.096482in}{0.679135in}}%
\pgfpathlineto{\pgfqpoint{2.097594in}{0.702999in}}%
\pgfpathlineto{\pgfqpoint{2.099562in}{0.805778in}}%
\pgfpathlineto{\pgfqpoint{2.102728in}{1.110597in}}%
\pgfpathlineto{\pgfqpoint{2.116931in}{2.666062in}}%
\pgfpathlineto{\pgfqpoint{2.118044in}{2.678107in}}%
\pgfpathlineto{\pgfqpoint{2.118471in}{2.676243in}}%
\pgfpathlineto{\pgfqpoint{2.119498in}{2.657149in}}%
\pgfpathlineto{\pgfqpoint{2.121381in}{2.570316in}}%
\pgfpathlineto{\pgfqpoint{2.124375in}{2.310412in}}%
\pgfpathlineto{\pgfqpoint{2.130022in}{1.580386in}}%
\pgfpathlineto{\pgfqpoint{2.136440in}{0.849783in}}%
\pgfpathlineto{\pgfqpoint{2.139605in}{0.694436in}}%
\pgfpathlineto{\pgfqpoint{2.140718in}{0.682503in}}%
\pgfpathlineto{\pgfqpoint{2.141146in}{0.684000in}}%
\pgfpathlineto{\pgfqpoint{2.142172in}{0.701303in}}%
\pgfpathlineto{\pgfqpoint{2.144055in}{0.781732in}}%
\pgfpathlineto{\pgfqpoint{2.147049in}{1.025110in}}%
\pgfpathlineto{\pgfqpoint{2.152440in}{1.687481in}}%
\pgfpathlineto{\pgfqpoint{2.159541in}{2.490525in}}%
\pgfpathlineto{\pgfqpoint{2.162793in}{2.655409in}}%
\pgfpathlineto{\pgfqpoint{2.164162in}{2.672148in}}%
\pgfpathlineto{\pgfqpoint{2.164504in}{2.671283in}}%
\pgfpathlineto{\pgfqpoint{2.165445in}{2.658540in}}%
\pgfpathlineto{\pgfqpoint{2.167157in}{2.597353in}}%
\pgfpathlineto{\pgfqpoint{2.169895in}{2.406285in}}%
\pgfpathlineto{\pgfqpoint{2.174429in}{1.906610in}}%
\pgfpathlineto{\pgfqpoint{2.183927in}{0.848744in}}%
\pgfpathlineto{\pgfqpoint{2.187178in}{0.701500in}}%
\pgfpathlineto{\pgfqpoint{2.188462in}{0.688720in}}%
\pgfpathlineto{\pgfqpoint{2.188804in}{0.689772in}}%
\pgfpathlineto{\pgfqpoint{2.189745in}{0.702292in}}%
\pgfpathlineto{\pgfqpoint{2.191542in}{0.764453in}}%
\pgfpathlineto{\pgfqpoint{2.194365in}{0.954698in}}%
\pgfpathlineto{\pgfqpoint{2.198986in}{1.443565in}}%
\pgfpathlineto{\pgfqpoint{2.208911in}{2.506627in}}%
\pgfpathlineto{\pgfqpoint{2.212162in}{2.650558in}}%
\pgfpathlineto{\pgfqpoint{2.213617in}{2.665635in}}%
\pgfpathlineto{\pgfqpoint{2.213959in}{2.664625in}}%
\pgfpathlineto{\pgfqpoint{2.214900in}{2.652939in}}%
\pgfpathlineto{\pgfqpoint{2.216697in}{2.595203in}}%
\pgfpathlineto{\pgfqpoint{2.219521in}{2.418266in}}%
\pgfpathlineto{\pgfqpoint{2.224056in}{1.968983in}}%
\pgfpathlineto{\pgfqpoint{2.235008in}{0.844801in}}%
\pgfpathlineto{\pgfqpoint{2.238259in}{0.710277in}}%
\pgfpathlineto{\pgfqpoint{2.239713in}{0.695553in}}%
\pgfpathlineto{\pgfqpoint{2.240056in}{0.696282in}}%
\pgfpathlineto{\pgfqpoint{2.240997in}{0.706484in}}%
\pgfpathlineto{\pgfqpoint{2.242708in}{0.755201in}}%
\pgfpathlineto{\pgfqpoint{2.245446in}{0.908286in}}%
\pgfpathlineto{\pgfqpoint{2.249724in}{1.293034in}}%
\pgfpathlineto{\pgfqpoint{2.262815in}{2.551704in}}%
\pgfpathlineto{\pgfqpoint{2.265896in}{2.651148in}}%
\pgfpathlineto{\pgfqpoint{2.267008in}{2.658440in}}%
\pgfpathlineto{\pgfqpoint{2.267436in}{2.657138in}}%
\pgfpathlineto{\pgfqpoint{2.268462in}{2.644766in}}%
\pgfpathlineto{\pgfqpoint{2.270345in}{2.589007in}}%
\pgfpathlineto{\pgfqpoint{2.273254in}{2.425394in}}%
\pgfpathlineto{\pgfqpoint{2.277874in}{2.017514in}}%
\pgfpathlineto{\pgfqpoint{2.290538in}{0.839104in}}%
\pgfpathlineto{\pgfqpoint{2.293875in}{0.717703in}}%
\pgfpathlineto{\pgfqpoint{2.295500in}{0.703154in}}%
\pgfpathlineto{\pgfqpoint{2.295757in}{0.703589in}}%
\pgfpathlineto{\pgfqpoint{2.296613in}{0.710397in}}%
\pgfpathlineto{\pgfqpoint{2.298238in}{0.745685in}}%
\pgfpathlineto{\pgfqpoint{2.300805in}{0.857781in}}%
\pgfpathlineto{\pgfqpoint{2.304741in}{1.142621in}}%
\pgfpathlineto{\pgfqpoint{2.312955in}{1.949484in}}%
\pgfpathlineto{\pgfqpoint{2.319372in}{2.469106in}}%
\pgfpathlineto{\pgfqpoint{2.323137in}{2.624187in}}%
\pgfpathlineto{\pgfqpoint{2.325362in}{2.650376in}}%
\pgfpathlineto{\pgfqpoint{2.325789in}{2.649648in}}%
\pgfpathlineto{\pgfqpoint{2.326816in}{2.640349in}}%
\pgfpathlineto{\pgfqpoint{2.328613in}{2.598938in}}%
\pgfpathlineto{\pgfqpoint{2.331436in}{2.473347in}}%
\pgfpathlineto{\pgfqpoint{2.335715in}{2.167072in}}%
\pgfpathlineto{\pgfqpoint{2.353340in}{0.771598in}}%
\pgfpathlineto{\pgfqpoint{2.356164in}{0.714792in}}%
\pgfpathlineto{\pgfqpoint{2.357020in}{0.711743in}}%
\pgfpathlineto{\pgfqpoint{2.357447in}{0.712701in}}%
\pgfpathlineto{\pgfqpoint{2.358560in}{0.722886in}}%
\pgfpathlineto{\pgfqpoint{2.360528in}{0.767532in}}%
\pgfpathlineto{\pgfqpoint{2.363522in}{0.896310in}}%
\pgfpathlineto{\pgfqpoint{2.367972in}{1.198505in}}%
\pgfpathlineto{\pgfqpoint{2.386282in}{2.570996in}}%
\pgfpathlineto{\pgfqpoint{2.389362in}{2.636082in}}%
\pgfpathlineto{\pgfqpoint{2.390475in}{2.641185in}}%
\pgfpathlineto{\pgfqpoint{2.390902in}{2.640531in}}%
\pgfpathlineto{\pgfqpoint{2.391929in}{2.633066in}}%
\pgfpathlineto{\pgfqpoint{2.393726in}{2.600336in}}%
\pgfpathlineto{\pgfqpoint{2.396549in}{2.501049in}}%
\pgfpathlineto{\pgfqpoint{2.400656in}{2.266956in}}%
\pgfpathlineto{\pgfqpoint{2.407673in}{1.715454in}}%
\pgfpathlineto{\pgfqpoint{2.417512in}{0.979573in}}%
\pgfpathlineto{\pgfqpoint{2.422133in}{0.781438in}}%
\pgfpathlineto{\pgfqpoint{2.425213in}{0.725815in}}%
\pgfpathlineto{\pgfqpoint{2.426325in}{0.721665in}}%
\pgfpathlineto{\pgfqpoint{2.426753in}{0.722329in}}%
\pgfpathlineto{\pgfqpoint{2.427780in}{0.729008in}}%
\pgfpathlineto{\pgfqpoint{2.429662in}{0.759583in}}%
\pgfpathlineto{\pgfqpoint{2.432486in}{0.847666in}}%
\pgfpathlineto{\pgfqpoint{2.436593in}{1.054574in}}%
\pgfpathlineto{\pgfqpoint{2.443181in}{1.517597in}}%
\pgfpathlineto{\pgfqpoint{2.455416in}{2.373893in}}%
\pgfpathlineto{\pgfqpoint{2.460208in}{2.565333in}}%
\pgfpathlineto{\pgfqpoint{2.463459in}{2.623938in}}%
\pgfpathlineto{\pgfqpoint{2.464999in}{2.630374in}}%
\pgfpathlineto{\pgfqpoint{2.465256in}{2.630109in}}%
\pgfpathlineto{\pgfqpoint{2.466197in}{2.625887in}}%
\pgfpathlineto{\pgfqpoint{2.467909in}{2.605302in}}%
\pgfpathlineto{\pgfqpoint{2.470561in}{2.541758in}}%
\pgfpathlineto{\pgfqpoint{2.474411in}{2.387746in}}%
\pgfpathlineto{\pgfqpoint{2.480058in}{2.060175in}}%
\pgfpathlineto{\pgfqpoint{2.498882in}{0.894213in}}%
\pgfpathlineto{\pgfqpoint{2.503246in}{0.769723in}}%
\pgfpathlineto{\pgfqpoint{2.506155in}{0.735873in}}%
\pgfpathlineto{\pgfqpoint{2.507182in}{0.733567in}}%
\pgfpathlineto{\pgfqpoint{2.507610in}{0.734086in}}%
\pgfpathlineto{\pgfqpoint{2.508722in}{0.739482in}}%
\pgfpathlineto{\pgfqpoint{2.510690in}{0.763103in}}%
\pgfpathlineto{\pgfqpoint{2.513599in}{0.829501in}}%
\pgfpathlineto{\pgfqpoint{2.517706in}{0.980535in}}%
\pgfpathlineto{\pgfqpoint{2.523695in}{1.292863in}}%
\pgfpathlineto{\pgfqpoint{2.544401in}{2.445459in}}%
\pgfpathlineto{\pgfqpoint{2.549107in}{2.571969in}}%
\pgfpathlineto{\pgfqpoint{2.552359in}{2.611992in}}%
\pgfpathlineto{\pgfqpoint{2.554070in}{2.617065in}}%
\pgfpathlineto{\pgfqpoint{2.554241in}{2.616967in}}%
\pgfpathlineto{\pgfqpoint{2.555182in}{2.614480in}}%
\pgfpathlineto{\pgfqpoint{2.556894in}{2.601599in}}%
\pgfpathlineto{\pgfqpoint{2.559546in}{2.560985in}}%
\pgfpathlineto{\pgfqpoint{2.563311in}{2.463338in}}%
\pgfpathlineto{\pgfqpoint{2.568530in}{2.262291in}}%
\pgfpathlineto{\pgfqpoint{2.576658in}{1.850055in}}%
\pgfpathlineto{\pgfqpoint{2.591974in}{1.072695in}}%
\pgfpathlineto{\pgfqpoint{2.598220in}{0.866028in}}%
\pgfpathlineto{\pgfqpoint{2.602755in}{0.778883in}}%
\pgfpathlineto{\pgfqpoint{2.605921in}{0.751831in}}%
\pgfpathlineto{\pgfqpoint{2.607461in}{0.748798in}}%
\pgfpathlineto{\pgfqpoint{2.607718in}{0.748932in}}%
\pgfpathlineto{\pgfqpoint{2.608744in}{0.751278in}}%
\pgfpathlineto{\pgfqpoint{2.610541in}{0.762265in}}%
\pgfpathlineto{\pgfqpoint{2.613279in}{0.795318in}}%
\pgfpathlineto{\pgfqpoint{2.617129in}{0.872917in}}%
\pgfpathlineto{\pgfqpoint{2.622349in}{1.028491in}}%
\pgfpathlineto{\pgfqpoint{2.629878in}{1.327572in}}%
\pgfpathlineto{\pgfqpoint{2.654092in}{2.339516in}}%
\pgfpathlineto{\pgfqpoint{2.660424in}{2.497446in}}%
\pgfpathlineto{\pgfqpoint{2.665216in}{2.569272in}}%
\pgfpathlineto{\pgfqpoint{2.668724in}{2.594703in}}%
\pgfpathlineto{\pgfqpoint{2.670948in}{2.598989in}}%
\pgfpathlineto{\pgfqpoint{2.672061in}{2.597743in}}%
\pgfpathlineto{\pgfqpoint{2.673772in}{2.591502in}}%
\pgfpathlineto{\pgfqpoint{2.676424in}{2.571766in}}%
\pgfpathlineto{\pgfqpoint{2.680103in}{2.525218in}}%
\pgfpathlineto{\pgfqpoint{2.685066in}{2.430614in}}%
\pgfpathlineto{\pgfqpoint{2.691740in}{2.256074in}}%
\pgfpathlineto{\pgfqpoint{2.701751in}{1.928195in}}%
\pgfpathlineto{\pgfqpoint{2.724083in}{1.185127in}}%
\pgfpathlineto{\pgfqpoint{2.732639in}{0.978923in}}%
\pgfpathlineto{\pgfqpoint{2.739398in}{0.864433in}}%
\pgfpathlineto{\pgfqpoint{2.744874in}{0.805144in}}%
\pgfpathlineto{\pgfqpoint{2.749152in}{0.779653in}}%
\pgfpathlineto{\pgfqpoint{2.752147in}{0.772387in}}%
\pgfpathlineto{\pgfqpoint{2.753944in}{0.772074in}}%
\pgfpathlineto{\pgfqpoint{2.755655in}{0.774514in}}%
\pgfpathlineto{\pgfqpoint{2.758136in}{0.782644in}}%
\pgfpathlineto{\pgfqpoint{2.761644in}{0.802966in}}%
\pgfpathlineto{\pgfqpoint{2.766265in}{0.844295in}}%
\pgfpathlineto{\pgfqpoint{2.772340in}{0.920720in}}%
\pgfpathlineto{\pgfqpoint{2.780383in}{1.052818in}}%
\pgfpathlineto{\pgfqpoint{2.792105in}{1.286558in}}%
\pgfpathlineto{\pgfqpoint{2.826672in}{1.996666in}}%
\pgfpathlineto{\pgfqpoint{2.838394in}{2.182695in}}%
\pgfpathlineto{\pgfqpoint{2.848661in}{2.312172in}}%
\pgfpathlineto{\pgfqpoint{2.857902in}{2.402131in}}%
\pgfpathlineto{\pgfqpoint{2.866287in}{2.463461in}}%
\pgfpathlineto{\pgfqpoint{2.873988in}{2.504522in}}%
\pgfpathlineto{\pgfqpoint{2.881004in}{2.530793in}}%
\pgfpathlineto{\pgfqpoint{2.887421in}{2.546814in}}%
\pgfpathlineto{\pgfqpoint{2.893154in}{2.555618in}}%
\pgfpathlineto{\pgfqpoint{2.898288in}{2.559772in}}%
\pgfpathlineto{\pgfqpoint{2.903165in}{2.560975in}}%
\pgfpathlineto{\pgfqpoint{2.908213in}{2.559895in}}%
\pgfpathlineto{\pgfqpoint{2.914202in}{2.556183in}}%
\pgfpathlineto{\pgfqpoint{2.922074in}{2.548439in}}%
\pgfpathlineto{\pgfqpoint{2.935251in}{2.531801in}}%
\pgfpathlineto{\pgfqpoint{2.953219in}{2.509890in}}%
\pgfpathlineto{\pgfqpoint{2.963315in}{2.501047in}}%
\pgfpathlineto{\pgfqpoint{2.971786in}{2.496456in}}%
\pgfpathlineto{\pgfqpoint{2.979486in}{2.494747in}}%
\pgfpathlineto{\pgfqpoint{2.987016in}{2.495361in}}%
\pgfpathlineto{\pgfqpoint{2.994973in}{2.498314in}}%
\pgfpathlineto{\pgfqpoint{3.004128in}{2.504212in}}%
\pgfpathlineto{\pgfqpoint{3.016792in}{2.515312in}}%
\pgfpathlineto{\pgfqpoint{3.035787in}{2.531780in}}%
\pgfpathlineto{\pgfqpoint{3.043487in}{2.535505in}}%
\pgfpathlineto{\pgfqpoint{3.049477in}{2.536054in}}%
\pgfpathlineto{\pgfqpoint{3.054696in}{2.534308in}}%
\pgfpathlineto{\pgfqpoint{3.059744in}{2.530201in}}%
\pgfpathlineto{\pgfqpoint{3.065049in}{2.522854in}}%
\pgfpathlineto{\pgfqpoint{3.070782in}{2.510846in}}%
\pgfpathlineto{\pgfqpoint{3.077028in}{2.492214in}}%
\pgfpathlineto{\pgfqpoint{3.083787in}{2.464604in}}%
\pgfpathlineto{\pgfqpoint{3.091145in}{2.424615in}}%
\pgfpathlineto{\pgfqpoint{3.099103in}{2.368433in}}%
\pgfpathlineto{\pgfqpoint{3.107830in}{2.290025in}}%
\pgfpathlineto{\pgfqpoint{3.117413in}{2.182675in}}%
\pgfpathlineto{\pgfqpoint{3.128194in}{2.035421in}}%
\pgfpathlineto{\pgfqpoint{3.140943in}{1.828773in}}%
\pgfpathlineto{\pgfqpoint{3.159852in}{1.480650in}}%
\pgfpathlineto{\pgfqpoint{3.179874in}{1.123590in}}%
\pgfpathlineto{\pgfqpoint{3.189885in}{0.982511in}}%
\pgfpathlineto{\pgfqpoint{3.197329in}{0.906128in}}%
\pgfpathlineto{\pgfqpoint{3.202976in}{0.868574in}}%
\pgfpathlineto{\pgfqpoint{3.207083in}{0.853805in}}%
\pgfpathlineto{\pgfqpoint{3.209906in}{0.850237in}}%
\pgfpathlineto{\pgfqpoint{3.211703in}{0.850874in}}%
\pgfpathlineto{\pgfqpoint{3.213842in}{0.854656in}}%
\pgfpathlineto{\pgfqpoint{3.216751in}{0.865182in}}%
\pgfpathlineto{\pgfqpoint{3.220516in}{0.888185in}}%
\pgfpathlineto{\pgfqpoint{3.225222in}{0.931981in}}%
\pgfpathlineto{\pgfqpoint{3.231040in}{1.009026in}}%
\pgfpathlineto{\pgfqpoint{3.238056in}{1.133959in}}%
\pgfpathlineto{\pgfqpoint{3.246869in}{1.334106in}}%
\pgfpathlineto{\pgfqpoint{3.259618in}{1.680935in}}%
\pgfpathlineto{\pgfqpoint{3.279126in}{2.207114in}}%
\pgfpathlineto{\pgfqpoint{3.287255in}{2.368196in}}%
\pgfpathlineto{\pgfqpoint{3.293244in}{2.447784in}}%
\pgfpathlineto{\pgfqpoint{3.297608in}{2.480817in}}%
\pgfpathlineto{\pgfqpoint{3.300602in}{2.490238in}}%
\pgfpathlineto{\pgfqpoint{3.302228in}{2.490657in}}%
\pgfpathlineto{\pgfqpoint{3.302314in}{2.490586in}}%
\pgfpathlineto{\pgfqpoint{3.303854in}{2.487711in}}%
\pgfpathlineto{\pgfqpoint{3.306164in}{2.477682in}}%
\pgfpathlineto{\pgfqpoint{3.309330in}{2.452753in}}%
\pgfpathlineto{\pgfqpoint{3.313522in}{2.400037in}}%
\pgfpathlineto{\pgfqpoint{3.318827in}{2.302514in}}%
\pgfpathlineto{\pgfqpoint{3.325587in}{2.133705in}}%
\pgfpathlineto{\pgfqpoint{3.334999in}{1.836716in}}%
\pgfpathlineto{\pgfqpoint{3.357758in}{1.097322in}}%
\pgfpathlineto{\pgfqpoint{3.364346in}{0.960541in}}%
\pgfpathlineto{\pgfqpoint{3.369138in}{0.901142in}}%
\pgfpathlineto{\pgfqpoint{3.372475in}{0.882316in}}%
\pgfpathlineto{\pgfqpoint{3.374272in}{0.880236in}}%
\pgfpathlineto{\pgfqpoint{3.374357in}{0.880280in}}%
\pgfpathlineto{\pgfqpoint{3.375555in}{0.882268in}}%
\pgfpathlineto{\pgfqpoint{3.377437in}{0.890588in}}%
\pgfpathlineto{\pgfqpoint{3.380261in}{0.914961in}}%
\pgfpathlineto{\pgfqpoint{3.384026in}{0.969272in}}%
\pgfpathlineto{\pgfqpoint{3.388988in}{1.076824in}}%
\pgfpathlineto{\pgfqpoint{3.395491in}{1.271188in}}%
\pgfpathlineto{\pgfqpoint{3.405416in}{1.644564in}}%
\pgfpathlineto{\pgfqpoint{3.420647in}{2.209304in}}%
\pgfpathlineto{\pgfqpoint{3.427149in}{2.373026in}}%
\pgfpathlineto{\pgfqpoint{3.431855in}{2.442960in}}%
\pgfpathlineto{\pgfqpoint{3.435021in}{2.463735in}}%
\pgfpathlineto{\pgfqpoint{3.436561in}{2.465812in}}%
\pgfpathlineto{\pgfqpoint{3.436818in}{2.465639in}}%
\pgfpathlineto{\pgfqpoint{3.438016in}{2.462866in}}%
\pgfpathlineto{\pgfqpoint{3.439984in}{2.451269in}}%
\pgfpathlineto{\pgfqpoint{3.442893in}{2.418212in}}%
\pgfpathlineto{\pgfqpoint{3.446829in}{2.344283in}}%
\pgfpathlineto{\pgfqpoint{3.452048in}{2.199160in}}%
\pgfpathlineto{\pgfqpoint{3.459321in}{1.927184in}}%
\pgfpathlineto{\pgfqpoint{3.481567in}{1.049868in}}%
\pgfpathlineto{\pgfqpoint{3.486701in}{0.943758in}}%
\pgfpathlineto{\pgfqpoint{3.490294in}{0.907753in}}%
\pgfpathlineto{\pgfqpoint{3.492348in}{0.902407in}}%
\pgfpathlineto{\pgfqpoint{3.492433in}{0.902430in}}%
\pgfpathlineto{\pgfqpoint{3.493375in}{0.903980in}}%
\pgfpathlineto{\pgfqpoint{3.495000in}{0.912283in}}%
\pgfpathlineto{\pgfqpoint{3.497482in}{0.938608in}}%
\pgfpathlineto{\pgfqpoint{3.500990in}{1.003159in}}%
\pgfpathlineto{\pgfqpoint{3.505696in}{1.135916in}}%
\pgfpathlineto{\pgfqpoint{3.512284in}{1.392614in}}%
\pgfpathlineto{\pgfqpoint{3.534872in}{2.331744in}}%
\pgfpathlineto{\pgfqpoint{3.539407in}{2.418590in}}%
\pgfpathlineto{\pgfqpoint{3.542487in}{2.443604in}}%
\pgfpathlineto{\pgfqpoint{3.543685in}{2.445506in}}%
\pgfpathlineto{\pgfqpoint{3.544028in}{2.445234in}}%
\pgfpathlineto{\pgfqpoint{3.545140in}{2.441840in}}%
\pgfpathlineto{\pgfqpoint{3.547022in}{2.427371in}}%
\pgfpathlineto{\pgfqpoint{3.549846in}{2.385426in}}%
\pgfpathlineto{\pgfqpoint{3.553782in}{2.288591in}}%
\pgfpathlineto{\pgfqpoint{3.559172in}{2.093552in}}%
\pgfpathlineto{\pgfqpoint{3.567814in}{1.686052in}}%
\pgfpathlineto{\pgfqpoint{3.579707in}{1.145492in}}%
\pgfpathlineto{\pgfqpoint{3.585183in}{0.987887in}}%
\pgfpathlineto{\pgfqpoint{3.588948in}{0.932067in}}%
\pgfpathlineto{\pgfqpoint{3.591258in}{0.921405in}}%
\pgfpathlineto{\pgfqpoint{3.592028in}{0.921957in}}%
\pgfpathlineto{\pgfqpoint{3.592114in}{0.922145in}}%
\pgfpathlineto{\pgfqpoint{3.593397in}{0.928033in}}%
\pgfpathlineto{\pgfqpoint{3.595536in}{0.950504in}}%
\pgfpathlineto{\pgfqpoint{3.598616in}{1.009882in}}%
\pgfpathlineto{\pgfqpoint{3.602895in}{1.140902in}}%
\pgfpathlineto{\pgfqpoint{3.608969in}{1.403616in}}%
\pgfpathlineto{\pgfqpoint{3.628820in}{2.317110in}}%
\pgfpathlineto{\pgfqpoint{3.633098in}{2.405108in}}%
\pgfpathlineto{\pgfqpoint{3.635836in}{2.426681in}}%
\pgfpathlineto{\pgfqpoint{3.636606in}{2.427626in}}%
\pgfpathlineto{\pgfqpoint{3.637120in}{2.426995in}}%
\pgfpathlineto{\pgfqpoint{3.638317in}{2.421594in}}%
\pgfpathlineto{\pgfqpoint{3.640371in}{2.399626in}}%
\pgfpathlineto{\pgfqpoint{3.643366in}{2.339645in}}%
\pgfpathlineto{\pgfqpoint{3.647558in}{2.204767in}}%
\pgfpathlineto{\pgfqpoint{3.653633in}{1.927516in}}%
\pgfpathlineto{\pgfqpoint{3.671601in}{1.061446in}}%
\pgfpathlineto{\pgfqpoint{3.675879in}{0.964378in}}%
\pgfpathlineto{\pgfqpoint{3.678703in}{0.939484in}}%
\pgfpathlineto{\pgfqpoint{3.679473in}{0.938388in}}%
\pgfpathlineto{\pgfqpoint{3.679986in}{0.939027in}}%
\pgfpathlineto{\pgfqpoint{3.681184in}{0.944782in}}%
\pgfpathlineto{\pgfqpoint{3.683238in}{0.968428in}}%
\pgfpathlineto{\pgfqpoint{3.686232in}{1.033092in}}%
\pgfpathlineto{\pgfqpoint{3.690425in}{1.178006in}}%
\pgfpathlineto{\pgfqpoint{3.696671in}{1.481878in}}%
\pgfpathlineto{\pgfqpoint{3.712158in}{2.262930in}}%
\pgfpathlineto{\pgfqpoint{3.716607in}{2.378155in}}%
\pgfpathlineto{\pgfqpoint{3.719516in}{2.409223in}}%
\pgfpathlineto{\pgfqpoint{3.720543in}{2.411285in}}%
\pgfpathlineto{\pgfqpoint{3.720971in}{2.410755in}}%
\pgfpathlineto{\pgfqpoint{3.722083in}{2.405553in}}%
\pgfpathlineto{\pgfqpoint{3.724051in}{2.382916in}}%
\pgfpathlineto{\pgfqpoint{3.726960in}{2.319040in}}%
\pgfpathlineto{\pgfqpoint{3.731067in}{2.173093in}}%
\pgfpathlineto{\pgfqpoint{3.737228in}{1.863450in}}%
\pgfpathlineto{\pgfqpoint{3.751859in}{1.104307in}}%
\pgfpathlineto{\pgfqpoint{3.756308in}{0.985793in}}%
\pgfpathlineto{\pgfqpoint{3.759132in}{0.955835in}}%
\pgfpathlineto{\pgfqpoint{3.760073in}{0.954165in}}%
\pgfpathlineto{\pgfqpoint{3.760501in}{0.954797in}}%
\pgfpathlineto{\pgfqpoint{3.761698in}{0.961188in}}%
\pgfpathlineto{\pgfqpoint{3.763666in}{0.986339in}}%
\pgfpathlineto{\pgfqpoint{3.766661in}{1.058217in}}%
\pgfpathlineto{\pgfqpoint{3.770854in}{1.219377in}}%
\pgfpathlineto{\pgfqpoint{3.777442in}{1.571720in}}%
\pgfpathlineto{\pgfqpoint{3.789592in}{2.221213in}}%
\pgfpathlineto{\pgfqpoint{3.794212in}{2.357717in}}%
\pgfpathlineto{\pgfqpoint{3.797207in}{2.393834in}}%
\pgfpathlineto{\pgfqpoint{3.798148in}{2.396029in}}%
\pgfpathlineto{\pgfqpoint{3.798661in}{2.395355in}}%
\pgfpathlineto{\pgfqpoint{3.799774in}{2.389362in}}%
\pgfpathlineto{\pgfqpoint{3.801742in}{2.363719in}}%
\pgfpathlineto{\pgfqpoint{3.804651in}{2.291991in}}%
\pgfpathlineto{\pgfqpoint{3.808843in}{2.125864in}}%
\pgfpathlineto{\pgfqpoint{3.815517in}{1.757791in}}%
\pgfpathlineto{\pgfqpoint{3.826812in}{1.143746in}}%
\pgfpathlineto{\pgfqpoint{3.831346in}{1.006799in}}%
\pgfpathlineto{\pgfqpoint{3.834255in}{0.971128in}}%
\pgfpathlineto{\pgfqpoint{3.835197in}{0.968986in}}%
\pgfpathlineto{\pgfqpoint{3.835624in}{0.969553in}}%
\pgfpathlineto{\pgfqpoint{3.836737in}{0.975528in}}%
\pgfpathlineto{\pgfqpoint{3.838619in}{1.000310in}}%
\pgfpathlineto{\pgfqpoint{3.841528in}{1.073424in}}%
\pgfpathlineto{\pgfqpoint{3.845721in}{1.244267in}}%
\pgfpathlineto{\pgfqpoint{3.852480in}{1.627229in}}%
\pgfpathlineto{\pgfqpoint{3.863004in}{2.207125in}}%
\pgfpathlineto{\pgfqpoint{3.867454in}{2.344087in}}%
\pgfpathlineto{\pgfqpoint{3.870363in}{2.379776in}}%
\pgfpathlineto{\pgfqpoint{3.871218in}{2.381552in}}%
\pgfpathlineto{\pgfqpoint{3.871646in}{2.380934in}}%
\pgfpathlineto{\pgfqpoint{3.872759in}{2.374632in}}%
\pgfpathlineto{\pgfqpoint{3.874726in}{2.347057in}}%
\pgfpathlineto{\pgfqpoint{3.877636in}{2.269550in}}%
\pgfpathlineto{\pgfqpoint{3.881828in}{2.090764in}}%
\pgfpathlineto{\pgfqpoint{3.888930in}{1.675178in}}%
\pgfpathlineto{\pgfqpoint{3.898342in}{1.155036in}}%
\pgfpathlineto{\pgfqpoint{3.902791in}{1.017741in}}%
\pgfpathlineto{\pgfqpoint{3.905615in}{0.984444in}}%
\pgfpathlineto{\pgfqpoint{3.906385in}{0.983152in}}%
\pgfpathlineto{\pgfqpoint{3.906812in}{0.983894in}}%
\pgfpathlineto{\pgfqpoint{3.907925in}{0.990693in}}%
\pgfpathlineto{\pgfqpoint{3.909893in}{1.019747in}}%
\pgfpathlineto{\pgfqpoint{3.912802in}{1.100644in}}%
\pgfpathlineto{\pgfqpoint{3.917080in}{1.290345in}}%
\pgfpathlineto{\pgfqpoint{3.924695in}{1.748611in}}%
\pgfpathlineto{\pgfqpoint{3.933080in}{2.206826in}}%
\pgfpathlineto{\pgfqpoint{3.937358in}{2.336358in}}%
\pgfpathlineto{\pgfqpoint{3.940096in}{2.366842in}}%
\pgfpathlineto{\pgfqpoint{3.940695in}{2.367688in}}%
\pgfpathlineto{\pgfqpoint{3.941209in}{2.366732in}}%
\pgfpathlineto{\pgfqpoint{3.942406in}{2.358479in}}%
\pgfpathlineto{\pgfqpoint{3.944460in}{2.325016in}}%
\pgfpathlineto{\pgfqpoint{3.947540in}{2.231672in}}%
\pgfpathlineto{\pgfqpoint{3.952075in}{2.015598in}}%
\pgfpathlineto{\pgfqpoint{3.970728in}{1.034574in}}%
\pgfpathlineto{\pgfqpoint{3.973551in}{0.998341in}}%
\pgfpathlineto{\pgfqpoint{3.974321in}{0.996741in}}%
\pgfpathlineto{\pgfqpoint{3.974835in}{0.997671in}}%
\pgfpathlineto{\pgfqpoint{3.976032in}{1.006040in}}%
\pgfpathlineto{\pgfqpoint{3.978086in}{1.040261in}}%
\pgfpathlineto{\pgfqpoint{3.981166in}{1.135886in}}%
\pgfpathlineto{\pgfqpoint{3.985701in}{1.356754in}}%
\pgfpathlineto{\pgfqpoint{4.003498in}{2.309902in}}%
\pgfpathlineto{\pgfqpoint{4.006407in}{2.351842in}}%
\pgfpathlineto{\pgfqpoint{4.007348in}{2.354311in}}%
\pgfpathlineto{\pgfqpoint{4.007776in}{2.353612in}}%
\pgfpathlineto{\pgfqpoint{4.008888in}{2.346475in}}%
\pgfpathlineto{\pgfqpoint{4.010856in}{2.315293in}}%
\pgfpathlineto{\pgfqpoint{4.013765in}{2.228109in}}%
\pgfpathlineto{\pgfqpoint{4.018129in}{2.020112in}}%
\pgfpathlineto{\pgfqpoint{4.028226in}{1.389633in}}%
\pgfpathlineto{\pgfqpoint{4.034129in}{1.108225in}}%
\pgfpathlineto{\pgfqpoint{4.037723in}{1.023090in}}%
\pgfpathlineto{\pgfqpoint{4.039776in}{1.009914in}}%
\pgfpathlineto{\pgfqpoint{4.039862in}{1.009945in}}%
\pgfpathlineto{\pgfqpoint{4.040547in}{1.011875in}}%
\pgfpathlineto{\pgfqpoint{4.041916in}{1.024639in}}%
\pgfpathlineto{\pgfqpoint{4.044140in}{1.070098in}}%
\pgfpathlineto{\pgfqpoint{4.047477in}{1.191003in}}%
\pgfpathlineto{\pgfqpoint{4.052611in}{1.471480in}}%
\pgfpathlineto{\pgfqpoint{4.066044in}{2.239496in}}%
\pgfpathlineto{\pgfqpoint{4.069638in}{2.327276in}}%
\pgfpathlineto{\pgfqpoint{4.071691in}{2.341319in}}%
\pgfpathlineto{\pgfqpoint{4.071777in}{2.341312in}}%
\pgfpathlineto{\pgfqpoint{4.072376in}{2.339933in}}%
\pgfpathlineto{\pgfqpoint{4.073659in}{2.329161in}}%
\pgfpathlineto{\pgfqpoint{4.075798in}{2.288024in}}%
\pgfpathlineto{\pgfqpoint{4.078964in}{2.177874in}}%
\pgfpathlineto{\pgfqpoint{4.083841in}{1.917867in}}%
\pgfpathlineto{\pgfqpoint{4.098130in}{1.106779in}}%
\pgfpathlineto{\pgfqpoint{4.101553in}{1.032026in}}%
\pgfpathlineto{\pgfqpoint{4.103264in}{1.022749in}}%
\pgfpathlineto{\pgfqpoint{4.103435in}{1.022884in}}%
\pgfpathlineto{\pgfqpoint{4.104205in}{1.025883in}}%
\pgfpathlineto{\pgfqpoint{4.105745in}{1.043545in}}%
\pgfpathlineto{\pgfqpoint{4.108226in}{1.103607in}}%
\pgfpathlineto{\pgfqpoint{4.111906in}{1.256528in}}%
\pgfpathlineto{\pgfqpoint{4.117895in}{1.615265in}}%
\pgfpathlineto{\pgfqpoint{4.127649in}{2.184561in}}%
\pgfpathlineto{\pgfqpoint{4.131671in}{2.305005in}}%
\pgfpathlineto{\pgfqpoint{4.134066in}{2.328398in}}%
\pgfpathlineto{\pgfqpoint{4.134409in}{2.328618in}}%
\pgfpathlineto{\pgfqpoint{4.134836in}{2.327787in}}%
\pgfpathlineto{\pgfqpoint{4.135949in}{2.319896in}}%
\pgfpathlineto{\pgfqpoint{4.137917in}{2.285996in}}%
\pgfpathlineto{\pgfqpoint{4.140911in}{2.188617in}}%
\pgfpathlineto{\pgfqpoint{4.145446in}{1.955901in}}%
\pgfpathlineto{\pgfqpoint{4.160933in}{1.092572in}}%
\pgfpathlineto{\pgfqpoint{4.164013in}{1.039091in}}%
\pgfpathlineto{\pgfqpoint{4.165125in}{1.035309in}}%
\pgfpathlineto{\pgfqpoint{4.165553in}{1.036093in}}%
\pgfpathlineto{\pgfqpoint{4.166666in}{1.043935in}}%
\pgfpathlineto{\pgfqpoint{4.168634in}{1.078003in}}%
\pgfpathlineto{\pgfqpoint{4.171628in}{1.176158in}}%
\pgfpathlineto{\pgfqpoint{4.176163in}{1.410609in}}%
\pgfpathlineto{\pgfqpoint{4.191308in}{2.257394in}}%
\pgfpathlineto{\pgfqpoint{4.194388in}{2.312123in}}%
\pgfpathlineto{\pgfqpoint{4.195500in}{2.316201in}}%
\pgfpathlineto{\pgfqpoint{4.195928in}{2.315507in}}%
\pgfpathlineto{\pgfqpoint{4.197040in}{2.307833in}}%
\pgfpathlineto{\pgfqpoint{4.198923in}{2.275851in}}%
\pgfpathlineto{\pgfqpoint{4.201832in}{2.182366in}}%
\pgfpathlineto{\pgfqpoint{4.206195in}{1.960322in}}%
\pgfpathlineto{\pgfqpoint{4.221768in}{1.097604in}}%
\pgfpathlineto{\pgfqpoint{4.224677in}{1.050519in}}%
\pgfpathlineto{\pgfqpoint{4.225618in}{1.047629in}}%
\pgfpathlineto{\pgfqpoint{4.226046in}{1.048344in}}%
\pgfpathlineto{\pgfqpoint{4.227158in}{1.056126in}}%
\pgfpathlineto{\pgfqpoint{4.229041in}{1.088448in}}%
\pgfpathlineto{\pgfqpoint{4.231950in}{1.182752in}}%
\pgfpathlineto{\pgfqpoint{4.236399in}{1.411293in}}%
\pgfpathlineto{\pgfqpoint{4.251458in}{2.249458in}}%
\pgfpathlineto{\pgfqpoint{4.254453in}{2.300576in}}%
\pgfpathlineto{\pgfqpoint{4.255479in}{2.303965in}}%
\pgfpathlineto{\pgfqpoint{4.255907in}{2.303206in}}%
\pgfpathlineto{\pgfqpoint{4.257020in}{2.295272in}}%
\pgfpathlineto{\pgfqpoint{4.258987in}{2.260522in}}%
\pgfpathlineto{\pgfqpoint{4.261982in}{2.160340in}}%
\pgfpathlineto{\pgfqpoint{4.266517in}{1.922219in}}%
\pgfpathlineto{\pgfqpoint{4.280635in}{1.126990in}}%
\pgfpathlineto{\pgfqpoint{4.283801in}{1.065259in}}%
\pgfpathlineto{\pgfqpoint{4.285084in}{1.059776in}}%
\pgfpathlineto{\pgfqpoint{4.285426in}{1.060264in}}%
\pgfpathlineto{\pgfqpoint{4.286453in}{1.066649in}}%
\pgfpathlineto{\pgfqpoint{4.288250in}{1.095320in}}%
\pgfpathlineto{\pgfqpoint{4.291073in}{1.182731in}}%
\pgfpathlineto{\pgfqpoint{4.295352in}{1.396210in}}%
\pgfpathlineto{\pgfqpoint{4.310924in}{2.248489in}}%
\pgfpathlineto{\pgfqpoint{4.313747in}{2.290136in}}%
\pgfpathlineto{\pgfqpoint{4.314518in}{2.291888in}}%
\pgfpathlineto{\pgfqpoint{4.314945in}{2.291057in}}%
\pgfpathlineto{\pgfqpoint{4.316058in}{2.282885in}}%
\pgfpathlineto{\pgfqpoint{4.318026in}{2.247561in}}%
\pgfpathlineto{\pgfqpoint{4.321020in}{2.146296in}}%
\pgfpathlineto{\pgfqpoint{4.325641in}{1.901729in}}%
\pgfpathlineto{\pgfqpoint{4.339074in}{1.144991in}}%
\pgfpathlineto{\pgfqpoint{4.342325in}{1.078197in}}%
\pgfpathlineto{\pgfqpoint{4.343694in}{1.071776in}}%
\pgfpathlineto{\pgfqpoint{4.344037in}{1.072234in}}%
\pgfpathlineto{\pgfqpoint{4.344978in}{1.077749in}}%
\pgfpathlineto{\pgfqpoint{4.346689in}{1.103538in}}%
\pgfpathlineto{\pgfqpoint{4.349427in}{1.184793in}}%
\pgfpathlineto{\pgfqpoint{4.353534in}{1.383495in}}%
\pgfpathlineto{\pgfqpoint{4.369705in}{2.247836in}}%
\pgfpathlineto{\pgfqpoint{4.372272in}{2.279202in}}%
\pgfpathlineto{\pgfqpoint{4.372786in}{2.279949in}}%
\pgfpathlineto{\pgfqpoint{4.373299in}{2.278836in}}%
\pgfpathlineto{\pgfqpoint{4.374497in}{2.269035in}}%
\pgfpathlineto{\pgfqpoint{4.376550in}{2.229309in}}%
\pgfpathlineto{\pgfqpoint{4.379631in}{2.119841in}}%
\pgfpathlineto{\pgfqpoint{4.384508in}{1.853962in}}%
\pgfpathlineto{\pgfqpoint{4.396486in}{1.175515in}}%
\pgfpathlineto{\pgfqpoint{4.399909in}{1.094409in}}%
\pgfpathlineto{\pgfqpoint{4.401706in}{1.083675in}}%
\pgfpathlineto{\pgfqpoint{4.401877in}{1.083837in}}%
\pgfpathlineto{\pgfqpoint{4.402732in}{1.087742in}}%
\pgfpathlineto{\pgfqpoint{4.404358in}{1.109209in}}%
\pgfpathlineto{\pgfqpoint{4.406925in}{1.178839in}}%
\pgfpathlineto{\pgfqpoint{4.410775in}{1.353332in}}%
\pgfpathlineto{\pgfqpoint{4.418647in}{1.841024in}}%
\pgfpathlineto{\pgfqpoint{4.424979in}{2.165419in}}%
\pgfpathlineto{\pgfqpoint{4.428572in}{2.255591in}}%
\pgfpathlineto{\pgfqpoint{4.430455in}{2.268109in}}%
\pgfpathlineto{\pgfqpoint{4.430626in}{2.268015in}}%
\pgfpathlineto{\pgfqpoint{4.431396in}{2.265046in}}%
\pgfpathlineto{\pgfqpoint{4.432850in}{2.248167in}}%
\pgfpathlineto{\pgfqpoint{4.435246in}{2.189474in}}%
\pgfpathlineto{\pgfqpoint{4.438840in}{2.038260in}}%
\pgfpathlineto{\pgfqpoint{4.445171in}{1.658671in}}%
\pgfpathlineto{\pgfqpoint{4.453129in}{1.218891in}}%
\pgfpathlineto{\pgfqpoint{4.456893in}{1.113572in}}%
\pgfpathlineto{\pgfqpoint{4.459118in}{1.095489in}}%
\pgfpathlineto{\pgfqpoint{4.459717in}{1.096520in}}%
\pgfpathlineto{\pgfqpoint{4.460915in}{1.106086in}}%
\pgfpathlineto{\pgfqpoint{4.462968in}{1.145216in}}%
\pgfpathlineto{\pgfqpoint{4.466049in}{1.253272in}}%
\pgfpathlineto{\pgfqpoint{4.470926in}{1.515328in}}%
\pgfpathlineto{\pgfqpoint{4.482648in}{2.166024in}}%
\pgfpathlineto{\pgfqpoint{4.486070in}{2.245862in}}%
\pgfpathlineto{\pgfqpoint{4.487781in}{2.256325in}}%
\pgfpathlineto{\pgfqpoint{4.488038in}{2.256139in}}%
\pgfpathlineto{\pgfqpoint{4.488894in}{2.252209in}}%
\pgfpathlineto{\pgfqpoint{4.490519in}{2.230877in}}%
\pgfpathlineto{\pgfqpoint{4.493086in}{2.161977in}}%
\pgfpathlineto{\pgfqpoint{4.496937in}{1.989878in}}%
\pgfpathlineto{\pgfqpoint{4.505236in}{1.486657in}}%
\pgfpathlineto{\pgfqpoint{4.511226in}{1.196922in}}%
\pgfpathlineto{\pgfqpoint{4.514648in}{1.117632in}}%
\pgfpathlineto{\pgfqpoint{4.516359in}{1.107241in}}%
\pgfpathlineto{\pgfqpoint{4.516616in}{1.107426in}}%
\pgfpathlineto{\pgfqpoint{4.517472in}{1.111327in}}%
\pgfpathlineto{\pgfqpoint{4.519097in}{1.132507in}}%
\pgfpathlineto{\pgfqpoint{4.521664in}{1.200908in}}%
\pgfpathlineto{\pgfqpoint{4.525515in}{1.371699in}}%
\pgfpathlineto{\pgfqpoint{4.533900in}{1.875524in}}%
\pgfpathlineto{\pgfqpoint{4.539803in}{2.156875in}}%
\pgfpathlineto{\pgfqpoint{4.543226in}{2.234707in}}%
\pgfpathlineto{\pgfqpoint{4.544937in}{2.244582in}}%
\pgfpathlineto{\pgfqpoint{4.545108in}{2.244467in}}%
\pgfpathlineto{\pgfqpoint{4.545878in}{2.241467in}}%
\pgfpathlineto{\pgfqpoint{4.547419in}{2.223398in}}%
\pgfpathlineto{\pgfqpoint{4.549900in}{2.161878in}}%
\pgfpathlineto{\pgfqpoint{4.553579in}{2.006999in}}%
\pgfpathlineto{\pgfqpoint{4.560510in}{1.599061in}}%
\pgfpathlineto{\pgfqpoint{4.567611in}{1.230911in}}%
\pgfpathlineto{\pgfqpoint{4.571290in}{1.134718in}}%
\pgfpathlineto{\pgfqpoint{4.573429in}{1.118975in}}%
\pgfpathlineto{\pgfqpoint{4.574028in}{1.120119in}}%
\pgfpathlineto{\pgfqpoint{4.575226in}{1.129679in}}%
\pgfpathlineto{\pgfqpoint{4.577280in}{1.168068in}}%
\pgfpathlineto{\pgfqpoint{4.580446in}{1.276959in}}%
\pgfpathlineto{\pgfqpoint{4.585408in}{1.537322in}}%
\pgfpathlineto{\pgfqpoint{4.596617in}{2.139173in}}%
\pgfpathlineto{\pgfqpoint{4.600125in}{2.221431in}}%
\pgfpathlineto{\pgfqpoint{4.602007in}{2.232846in}}%
\pgfpathlineto{\pgfqpoint{4.602093in}{2.232801in}}%
\pgfpathlineto{\pgfqpoint{4.602777in}{2.230671in}}%
\pgfpathlineto{\pgfqpoint{4.604146in}{2.217049in}}%
\pgfpathlineto{\pgfqpoint{4.606457in}{2.166731in}}%
\pgfpathlineto{\pgfqpoint{4.609879in}{2.036192in}}%
\pgfpathlineto{\pgfqpoint{4.615697in}{1.713587in}}%
\pgfpathlineto{\pgfqpoint{4.624510in}{1.246198in}}%
\pgfpathlineto{\pgfqpoint{4.628275in}{1.147274in}}%
\pgfpathlineto{\pgfqpoint{4.630500in}{1.130710in}}%
\pgfpathlineto{\pgfqpoint{4.631099in}{1.131839in}}%
\pgfpathlineto{\pgfqpoint{4.632296in}{1.141197in}}%
\pgfpathlineto{\pgfqpoint{4.634350in}{1.178713in}}%
\pgfpathlineto{\pgfqpoint{4.637516in}{1.285059in}}%
\pgfpathlineto{\pgfqpoint{4.642478in}{1.539299in}}%
\pgfpathlineto{\pgfqpoint{4.653687in}{2.128042in}}%
\pgfpathlineto{\pgfqpoint{4.657195in}{2.209318in}}%
\pgfpathlineto{\pgfqpoint{4.659078in}{2.221103in}}%
\pgfpathlineto{\pgfqpoint{4.659249in}{2.221029in}}%
\pgfpathlineto{\pgfqpoint{4.660019in}{2.218328in}}%
\pgfpathlineto{\pgfqpoint{4.661473in}{2.202754in}}%
\pgfpathlineto{\pgfqpoint{4.663869in}{2.148433in}}%
\pgfpathlineto{\pgfqpoint{4.667463in}{2.008466in}}%
\pgfpathlineto{\pgfqpoint{4.673880in}{1.652813in}}%
\pgfpathlineto{\pgfqpoint{4.681666in}{1.257010in}}%
\pgfpathlineto{\pgfqpoint{4.685431in}{1.159502in}}%
\pgfpathlineto{\pgfqpoint{4.687655in}{1.142493in}}%
\pgfpathlineto{\pgfqpoint{4.688254in}{1.143336in}}%
\pgfpathlineto{\pgfqpoint{4.689367in}{1.151004in}}%
\pgfpathlineto{\pgfqpoint{4.691335in}{1.183569in}}%
\pgfpathlineto{\pgfqpoint{4.694329in}{1.276126in}}%
\pgfpathlineto{\pgfqpoint{4.699035in}{1.502229in}}%
\pgfpathlineto{\pgfqpoint{4.711613in}{2.136779in}}%
\pgfpathlineto{\pgfqpoint{4.714950in}{2.202185in}}%
\pgfpathlineto{\pgfqpoint{4.716490in}{2.209315in}}%
\pgfpathlineto{\pgfqpoint{4.716747in}{2.209046in}}%
\pgfpathlineto{\pgfqpoint{4.717688in}{2.204501in}}%
\pgfpathlineto{\pgfqpoint{4.719399in}{2.182114in}}%
\pgfpathlineto{\pgfqpoint{4.722051in}{2.113375in}}%
\pgfpathlineto{\pgfqpoint{4.726073in}{1.943352in}}%
\pgfpathlineto{\pgfqpoint{4.742415in}{1.179361in}}%
\pgfpathlineto{\pgfqpoint{4.744897in}{1.154711in}}%
\pgfpathlineto{\pgfqpoint{4.745324in}{1.154297in}}%
\pgfpathlineto{\pgfqpoint{4.745838in}{1.155300in}}%
\pgfpathlineto{\pgfqpoint{4.747036in}{1.163972in}}%
\pgfpathlineto{\pgfqpoint{4.749089in}{1.198955in}}%
\pgfpathlineto{\pgfqpoint{4.752255in}{1.298409in}}%
\pgfpathlineto{\pgfqpoint{4.757218in}{1.537068in}}%
\pgfpathlineto{\pgfqpoint{4.768769in}{2.110094in}}%
\pgfpathlineto{\pgfqpoint{4.772362in}{2.187232in}}%
\pgfpathlineto{\pgfqpoint{4.774245in}{2.197461in}}%
\pgfpathlineto{\pgfqpoint{4.774330in}{2.197415in}}%
\pgfpathlineto{\pgfqpoint{4.775100in}{2.194995in}}%
\pgfpathlineto{\pgfqpoint{4.776555in}{2.180676in}}%
\pgfpathlineto{\pgfqpoint{4.778951in}{2.130376in}}%
\pgfpathlineto{\pgfqpoint{4.782544in}{2.000218in}}%
\pgfpathlineto{\pgfqpoint{4.788790in}{1.676566in}}%
\pgfpathlineto{\pgfqpoint{4.797004in}{1.279015in}}%
\pgfpathlineto{\pgfqpoint{4.800855in}{1.183474in}}%
\pgfpathlineto{\pgfqpoint{4.803165in}{1.166216in}}%
\pgfpathlineto{\pgfqpoint{4.803764in}{1.166908in}}%
\pgfpathlineto{\pgfqpoint{4.804876in}{1.173839in}}%
\pgfpathlineto{\pgfqpoint{4.806844in}{1.203689in}}%
\pgfpathlineto{\pgfqpoint{4.809839in}{1.289020in}}%
\pgfpathlineto{\pgfqpoint{4.814459in}{1.494338in}}%
\pgfpathlineto{\pgfqpoint{4.827807in}{2.122573in}}%
\pgfpathlineto{\pgfqpoint{4.831058in}{2.179453in}}%
\pgfpathlineto{\pgfqpoint{4.832513in}{2.185520in}}%
\pgfpathlineto{\pgfqpoint{4.832855in}{2.185161in}}%
\pgfpathlineto{\pgfqpoint{4.833796in}{2.180671in}}%
\pgfpathlineto{\pgfqpoint{4.835507in}{2.159543in}}%
\pgfpathlineto{\pgfqpoint{4.838245in}{2.092885in}}%
\pgfpathlineto{\pgfqpoint{4.842352in}{1.929829in}}%
\pgfpathlineto{\pgfqpoint{4.858866in}{1.204492in}}%
\pgfpathlineto{\pgfqpoint{4.861433in}{1.178871in}}%
\pgfpathlineto{\pgfqpoint{4.862032in}{1.178224in}}%
\pgfpathlineto{\pgfqpoint{4.862459in}{1.179006in}}%
\pgfpathlineto{\pgfqpoint{4.863657in}{1.186688in}}%
\pgfpathlineto{\pgfqpoint{4.865711in}{1.218244in}}%
\pgfpathlineto{\pgfqpoint{4.868791in}{1.305636in}}%
\pgfpathlineto{\pgfqpoint{4.873583in}{1.514744in}}%
\pgfpathlineto{\pgfqpoint{4.886417in}{2.102237in}}%
\pgfpathlineto{\pgfqpoint{4.889839in}{2.165356in}}%
\pgfpathlineto{\pgfqpoint{4.891551in}{2.173444in}}%
\pgfpathlineto{\pgfqpoint{4.891807in}{2.173266in}}%
\pgfpathlineto{\pgfqpoint{4.892663in}{2.170053in}}%
\pgfpathlineto{\pgfqpoint{4.894289in}{2.152977in}}%
\pgfpathlineto{\pgfqpoint{4.896856in}{2.098109in}}%
\pgfpathlineto{\pgfqpoint{4.900706in}{1.960814in}}%
\pgfpathlineto{\pgfqpoint{4.907979in}{1.602435in}}%
\pgfpathlineto{\pgfqpoint{4.915166in}{1.291844in}}%
\pgfpathlineto{\pgfqpoint{4.919016in}{1.205783in}}%
\pgfpathlineto{\pgfqpoint{4.921326in}{1.190367in}}%
\pgfpathlineto{\pgfqpoint{4.922011in}{1.191272in}}%
\pgfpathlineto{\pgfqpoint{4.923209in}{1.198865in}}%
\pgfpathlineto{\pgfqpoint{4.925262in}{1.229244in}}%
\pgfpathlineto{\pgfqpoint{4.928343in}{1.312669in}}%
\pgfpathlineto{\pgfqpoint{4.933134in}{1.512012in}}%
\pgfpathlineto{\pgfqpoint{4.946396in}{2.093490in}}%
\pgfpathlineto{\pgfqpoint{4.949819in}{2.153376in}}%
\pgfpathlineto{\pgfqpoint{4.951616in}{2.161197in}}%
\pgfpathlineto{\pgfqpoint{4.951787in}{2.161068in}}%
\pgfpathlineto{\pgfqpoint{4.952642in}{2.158151in}}%
\pgfpathlineto{\pgfqpoint{4.954268in}{2.142277in}}%
\pgfpathlineto{\pgfqpoint{4.956835in}{2.090887in}}%
\pgfpathlineto{\pgfqpoint{4.960685in}{1.961613in}}%
\pgfpathlineto{\pgfqpoint{4.967701in}{1.633052in}}%
\pgfpathlineto{\pgfqpoint{4.975402in}{1.307973in}}%
\pgfpathlineto{\pgfqpoint{4.979338in}{1.220255in}}%
\pgfpathlineto{\pgfqpoint{4.981734in}{1.202843in}}%
\pgfpathlineto{\pgfqpoint{4.982161in}{1.202752in}}%
\pgfpathlineto{\pgfqpoint{4.982504in}{1.203341in}}%
\pgfpathlineto{\pgfqpoint{4.983616in}{1.209296in}}%
\pgfpathlineto{\pgfqpoint{4.985584in}{1.234671in}}%
\pgfpathlineto{\pgfqpoint{4.988579in}{1.307147in}}%
\pgfpathlineto{\pgfqpoint{4.993113in}{1.479151in}}%
\pgfpathlineto{\pgfqpoint{5.008343in}{2.102289in}}%
\pgfpathlineto{\pgfqpoint{5.011509in}{2.145010in}}%
\pgfpathlineto{\pgfqpoint{5.012793in}{2.148758in}}%
\pgfpathlineto{\pgfqpoint{5.013135in}{2.148411in}}%
\pgfpathlineto{\pgfqpoint{5.014162in}{2.143977in}}%
\pgfpathlineto{\pgfqpoint{5.015959in}{2.124165in}}%
\pgfpathlineto{\pgfqpoint{5.018782in}{2.063810in}}%
\pgfpathlineto{\pgfqpoint{5.022975in}{1.919165in}}%
\pgfpathlineto{\pgfqpoint{5.040772in}{1.237735in}}%
\pgfpathlineto{\pgfqpoint{5.043424in}{1.215742in}}%
\pgfpathlineto{\pgfqpoint{5.044023in}{1.215285in}}%
\pgfpathlineto{\pgfqpoint{5.044451in}{1.215981in}}%
\pgfpathlineto{\pgfqpoint{5.045649in}{1.222443in}}%
\pgfpathlineto{\pgfqpoint{5.047702in}{1.248654in}}%
\pgfpathlineto{\pgfqpoint{5.050782in}{1.321187in}}%
\pgfpathlineto{\pgfqpoint{5.055488in}{1.492955in}}%
\pgfpathlineto{\pgfqpoint{5.070547in}{2.083579in}}%
\pgfpathlineto{\pgfqpoint{5.073884in}{2.130723in}}%
\pgfpathlineto{\pgfqpoint{5.075424in}{2.136065in}}%
\pgfpathlineto{\pgfqpoint{5.075681in}{2.135924in}}%
\pgfpathlineto{\pgfqpoint{5.076622in}{2.132895in}}%
\pgfpathlineto{\pgfqpoint{5.078334in}{2.117386in}}%
\pgfpathlineto{\pgfqpoint{5.080986in}{2.069033in}}%
\pgfpathlineto{\pgfqpoint{5.084922in}{1.950218in}}%
\pgfpathlineto{\pgfqpoint{5.091852in}{1.659390in}}%
\pgfpathlineto{\pgfqpoint{5.100238in}{1.334860in}}%
\pgfpathlineto{\pgfqpoint{5.104430in}{1.247232in}}%
\pgfpathlineto{\pgfqpoint{5.106997in}{1.228485in}}%
\pgfpathlineto{\pgfqpoint{5.107596in}{1.228147in}}%
\pgfpathlineto{\pgfqpoint{5.108024in}{1.228845in}}%
\pgfpathlineto{\pgfqpoint{5.109222in}{1.234940in}}%
\pgfpathlineto{\pgfqpoint{5.111275in}{1.259285in}}%
\pgfpathlineto{\pgfqpoint{5.114355in}{1.326394in}}%
\pgfpathlineto{\pgfqpoint{5.118976in}{1.482390in}}%
\pgfpathlineto{\pgfqpoint{5.135318in}{2.080152in}}%
\pgfpathlineto{\pgfqpoint{5.138570in}{2.119266in}}%
\pgfpathlineto{\pgfqpoint{5.139939in}{2.123055in}}%
\pgfpathlineto{\pgfqpoint{5.140281in}{2.122807in}}%
\pgfpathlineto{\pgfqpoint{5.141308in}{2.119206in}}%
\pgfpathlineto{\pgfqpoint{5.143104in}{2.102734in}}%
\pgfpathlineto{\pgfqpoint{5.145842in}{2.054000in}}%
\pgfpathlineto{\pgfqpoint{5.149949in}{1.935040in}}%
\pgfpathlineto{\pgfqpoint{5.157308in}{1.642652in}}%
\pgfpathlineto{\pgfqpoint{5.165522in}{1.345960in}}%
\pgfpathlineto{\pgfqpoint{5.169800in}{1.261038in}}%
\pgfpathlineto{\pgfqpoint{5.172452in}{1.241861in}}%
\pgfpathlineto{\pgfqpoint{5.173137in}{1.241317in}}%
\pgfpathlineto{\pgfqpoint{5.173565in}{1.241899in}}%
\pgfpathlineto{\pgfqpoint{5.174762in}{1.247284in}}%
\pgfpathlineto{\pgfqpoint{5.176816in}{1.269128in}}%
\pgfpathlineto{\pgfqpoint{5.179896in}{1.329794in}}%
\pgfpathlineto{\pgfqpoint{5.184431in}{1.469017in}}%
\pgfpathlineto{\pgfqpoint{5.202485in}{2.077402in}}%
\pgfpathlineto{\pgfqpoint{5.205565in}{2.107312in}}%
\pgfpathlineto{\pgfqpoint{5.206763in}{2.109649in}}%
\pgfpathlineto{\pgfqpoint{5.207105in}{2.109350in}}%
\pgfpathlineto{\pgfqpoint{5.208217in}{2.105419in}}%
\pgfpathlineto{\pgfqpoint{5.210100in}{2.088627in}}%
\pgfpathlineto{\pgfqpoint{5.213009in}{2.038952in}}%
\pgfpathlineto{\pgfqpoint{5.217287in}{1.921391in}}%
\pgfpathlineto{\pgfqpoint{5.225073in}{1.631564in}}%
\pgfpathlineto{\pgfqpoint{5.233202in}{1.358939in}}%
\pgfpathlineto{\pgfqpoint{5.237565in}{1.276345in}}%
\pgfpathlineto{\pgfqpoint{5.240389in}{1.255769in}}%
\pgfpathlineto{\pgfqpoint{5.241244in}{1.254945in}}%
\pgfpathlineto{\pgfqpoint{5.241672in}{1.255482in}}%
\pgfpathlineto{\pgfqpoint{5.242870in}{1.260335in}}%
\pgfpathlineto{\pgfqpoint{5.244924in}{1.279913in}}%
\pgfpathlineto{\pgfqpoint{5.248004in}{1.334291in}}%
\pgfpathlineto{\pgfqpoint{5.252539in}{1.459792in}}%
\pgfpathlineto{\pgfqpoint{5.262122in}{1.806901in}}%
\pgfpathlineto{\pgfqpoint{5.268967in}{2.011268in}}%
\pgfpathlineto{\pgfqpoint{5.273159in}{2.079973in}}%
\pgfpathlineto{\pgfqpoint{5.275812in}{2.095410in}}%
\pgfpathlineto{\pgfqpoint{5.276496in}{2.095703in}}%
\pgfpathlineto{\pgfqpoint{5.276924in}{2.095115in}}%
\pgfpathlineto{\pgfqpoint{5.278207in}{2.089813in}}%
\pgfpathlineto{\pgfqpoint{5.280347in}{2.069430in}}%
\pgfpathlineto{\pgfqpoint{5.283512in}{2.014508in}}%
\pgfpathlineto{\pgfqpoint{5.288133in}{1.890489in}}%
\pgfpathlineto{\pgfqpoint{5.298400in}{1.533973in}}%
\pgfpathlineto{\pgfqpoint{5.305074in}{1.347566in}}%
\pgfpathlineto{\pgfqpoint{5.309181in}{1.284413in}}%
\pgfpathlineto{\pgfqpoint{5.311833in}{1.269514in}}%
\pgfpathlineto{\pgfqpoint{5.312604in}{1.269160in}}%
\pgfpathlineto{\pgfqpoint{5.312946in}{1.269578in}}%
\pgfpathlineto{\pgfqpoint{5.314229in}{1.274280in}}%
\pgfpathlineto{\pgfqpoint{5.316368in}{1.292904in}}%
\pgfpathlineto{\pgfqpoint{5.319534in}{1.343683in}}%
\pgfpathlineto{\pgfqpoint{5.324154in}{1.459415in}}%
\pgfpathlineto{\pgfqpoint{5.332967in}{1.751041in}}%
\pgfpathlineto{\pgfqpoint{5.340925in}{1.982321in}}%
\pgfpathlineto{\pgfqpoint{5.345460in}{2.058850in}}%
\pgfpathlineto{\pgfqpoint{5.348454in}{2.079828in}}%
\pgfpathlineto{\pgfqpoint{5.349567in}{2.081264in}}%
\pgfpathlineto{\pgfqpoint{5.349994in}{2.080893in}}%
\pgfpathlineto{\pgfqpoint{5.351192in}{2.077141in}}%
\pgfpathlineto{\pgfqpoint{5.353246in}{2.061568in}}%
\pgfpathlineto{\pgfqpoint{5.356240in}{2.019255in}}%
\pgfpathlineto{\pgfqpoint{5.360604in}{1.921748in}}%
\pgfpathlineto{\pgfqpoint{5.367877in}{1.699902in}}%
\pgfpathlineto{\pgfqpoint{5.378316in}{1.393980in}}%
\pgfpathlineto{\pgfqpoint{5.383107in}{1.311355in}}%
\pgfpathlineto{\pgfqpoint{5.386358in}{1.286390in}}%
\pgfpathlineto{\pgfqpoint{5.387813in}{1.283958in}}%
\pgfpathlineto{\pgfqpoint{5.388155in}{1.284177in}}%
\pgfpathlineto{\pgfqpoint{5.389353in}{1.287311in}}%
\pgfpathlineto{\pgfqpoint{5.391321in}{1.300324in}}%
\pgfpathlineto{\pgfqpoint{5.394230in}{1.336629in}}%
\pgfpathlineto{\pgfqpoint{5.398423in}{1.420420in}}%
\pgfpathlineto{\pgfqpoint{5.405011in}{1.603873in}}%
\pgfpathlineto{\pgfqpoint{5.417931in}{1.965438in}}%
\pgfpathlineto{\pgfqpoint{5.422722in}{2.041047in}}%
\pgfpathlineto{\pgfqpoint{5.425974in}{2.063773in}}%
\pgfpathlineto{\pgfqpoint{5.427514in}{2.065940in}}%
\pgfpathlineto{\pgfqpoint{5.427771in}{2.065760in}}%
\pgfpathlineto{\pgfqpoint{5.428969in}{2.062886in}}%
\pgfpathlineto{\pgfqpoint{5.430936in}{2.050995in}}%
\pgfpathlineto{\pgfqpoint{5.433846in}{2.017824in}}%
\pgfpathlineto{\pgfqpoint{5.438038in}{1.941036in}}%
\pgfpathlineto{\pgfqpoint{5.444455in}{1.776110in}}%
\pgfpathlineto{\pgfqpoint{5.459001in}{1.392859in}}%
\pgfpathlineto{\pgfqpoint{5.463792in}{1.323334in}}%
\pgfpathlineto{\pgfqpoint{5.467044in}{1.302071in}}%
\pgfpathlineto{\pgfqpoint{5.468669in}{1.299774in}}%
\pgfpathlineto{\pgfqpoint{5.468841in}{1.299858in}}%
\pgfpathlineto{\pgfqpoint{5.469953in}{1.301905in}}%
\pgfpathlineto{\pgfqpoint{5.471835in}{1.311236in}}%
\pgfpathlineto{\pgfqpoint{5.474659in}{1.338535in}}%
\pgfpathlineto{\pgfqpoint{5.478680in}{1.402296in}}%
\pgfpathlineto{\pgfqpoint{5.484584in}{1.536479in}}%
\pgfpathlineto{\pgfqpoint{5.503066in}{1.982795in}}%
\pgfpathlineto{\pgfqpoint{5.507515in}{2.034137in}}%
\pgfpathlineto{\pgfqpoint{5.510510in}{2.048526in}}%
\pgfpathlineto{\pgfqpoint{5.511879in}{2.049493in}}%
\pgfpathlineto{\pgfqpoint{5.512135in}{2.049281in}}%
\pgfpathlineto{\pgfqpoint{5.513504in}{2.046074in}}%
\pgfpathlineto{\pgfqpoint{5.515643in}{2.034166in}}%
\pgfpathlineto{\pgfqpoint{5.518809in}{2.001858in}}%
\pgfpathlineto{\pgfqpoint{5.523258in}{1.930302in}}%
\pgfpathlineto{\pgfqpoint{5.529932in}{1.781708in}}%
\pgfpathlineto{\pgfqpoint{5.546189in}{1.408166in}}%
\pgfpathlineto{\pgfqpoint{5.551323in}{1.342046in}}%
\pgfpathlineto{\pgfqpoint{5.554831in}{1.320254in}}%
\pgfpathlineto{\pgfqpoint{5.556970in}{1.316858in}}%
\pgfpathlineto{\pgfqpoint{5.558168in}{1.318247in}}%
\pgfpathlineto{\pgfqpoint{5.559965in}{1.324700in}}%
\pgfpathlineto{\pgfqpoint{5.562703in}{1.344302in}}%
\pgfpathlineto{\pgfqpoint{5.566553in}{1.390297in}}%
\pgfpathlineto{\pgfqpoint{5.571943in}{1.484582in}}%
\pgfpathlineto{\pgfqpoint{5.581526in}{1.699598in}}%
\pgfpathlineto{\pgfqpoint{5.591965in}{1.917047in}}%
\pgfpathlineto{\pgfqpoint{5.597698in}{1.994872in}}%
\pgfpathlineto{\pgfqpoint{5.601719in}{2.024315in}}%
\pgfpathlineto{\pgfqpoint{5.604372in}{2.031399in}}%
\pgfpathlineto{\pgfqpoint{5.605741in}{2.031167in}}%
\pgfpathlineto{\pgfqpoint{5.607281in}{2.027769in}}%
\pgfpathlineto{\pgfqpoint{5.609676in}{2.016027in}}%
\pgfpathlineto{\pgfqpoint{5.613099in}{1.986327in}}%
\pgfpathlineto{\pgfqpoint{5.617805in}{1.923566in}}%
\pgfpathlineto{\pgfqpoint{5.624650in}{1.798778in}}%
\pgfpathlineto{\pgfqpoint{5.644586in}{1.418229in}}%
\pgfpathlineto{\pgfqpoint{5.649976in}{1.361200in}}%
\pgfpathlineto{\pgfqpoint{5.653827in}{1.340237in}}%
\pgfpathlineto{\pgfqpoint{5.656308in}{1.335924in}}%
\pgfpathlineto{\pgfqpoint{5.657677in}{1.336641in}}%
\pgfpathlineto{\pgfqpoint{5.659474in}{1.340881in}}%
\pgfpathlineto{\pgfqpoint{5.662126in}{1.353800in}}%
\pgfpathlineto{\pgfqpoint{5.665891in}{1.384951in}}%
\pgfpathlineto{\pgfqpoint{5.671025in}{1.448536in}}%
\pgfpathlineto{\pgfqpoint{5.678554in}{1.573426in}}%
\pgfpathlineto{\pgfqpoint{5.698490in}{1.916444in}}%
\pgfpathlineto{\pgfqpoint{5.704394in}{1.977656in}}%
\pgfpathlineto{\pgfqpoint{5.708758in}{2.003289in}}%
\pgfpathlineto{\pgfqpoint{5.711838in}{2.010687in}}%
\pgfpathlineto{\pgfqpoint{5.713635in}{2.010882in}}%
\pgfpathlineto{\pgfqpoint{5.715346in}{2.008279in}}%
\pgfpathlineto{\pgfqpoint{5.717827in}{1.999784in}}%
\pgfpathlineto{\pgfqpoint{5.721335in}{1.978690in}}%
\pgfpathlineto{\pgfqpoint{5.726041in}{1.935285in}}%
\pgfpathlineto{\pgfqpoint{5.732544in}{1.852377in}}%
\pgfpathlineto{\pgfqpoint{5.743838in}{1.674206in}}%
\pgfpathlineto{\pgfqpoint{5.756502in}{1.485649in}}%
\pgfpathlineto{\pgfqpoint{5.763603in}{1.410470in}}%
\pgfpathlineto{\pgfqpoint{5.768908in}{1.375049in}}%
\pgfpathlineto{\pgfqpoint{5.772844in}{1.361370in}}%
\pgfpathlineto{\pgfqpoint{5.775496in}{1.358327in}}%
\pgfpathlineto{\pgfqpoint{5.777293in}{1.359067in}}%
\pgfpathlineto{\pgfqpoint{5.779432in}{1.362842in}}%
\pgfpathlineto{\pgfqpoint{5.782427in}{1.373241in}}%
\pgfpathlineto{\pgfqpoint{5.786534in}{1.396612in}}%
\pgfpathlineto{\pgfqpoint{5.792010in}{1.442245in}}%
\pgfpathlineto{\pgfqpoint{5.799540in}{1.525888in}}%
\pgfpathlineto{\pgfqpoint{5.814684in}{1.726272in}}%
\pgfpathlineto{\pgfqpoint{5.826321in}{1.864712in}}%
\pgfpathlineto{\pgfqpoint{5.833936in}{1.931137in}}%
\pgfpathlineto{\pgfqpoint{5.839839in}{1.965321in}}%
\pgfpathlineto{\pgfqpoint{5.844460in}{1.980653in}}%
\pgfpathlineto{\pgfqpoint{5.847797in}{1.985404in}}%
\pgfpathlineto{\pgfqpoint{5.850192in}{1.985590in}}%
\pgfpathlineto{\pgfqpoint{5.852503in}{1.983276in}}%
\pgfpathlineto{\pgfqpoint{5.855583in}{1.976514in}}%
\pgfpathlineto{\pgfqpoint{5.859690in}{1.961323in}}%
\pgfpathlineto{\pgfqpoint{5.864995in}{1.932265in}}%
\pgfpathlineto{\pgfqpoint{5.872011in}{1.880240in}}%
\pgfpathlineto{\pgfqpoint{5.882107in}{1.786741in}}%
\pgfpathlineto{\pgfqpoint{5.910514in}{1.514959in}}%
\pgfpathlineto{\pgfqpoint{5.919413in}{1.454536in}}%
\pgfpathlineto{\pgfqpoint{5.926600in}{1.419627in}}%
\pgfpathlineto{\pgfqpoint{5.932504in}{1.400786in}}%
\pgfpathlineto{\pgfqpoint{5.937295in}{1.392002in}}%
\pgfpathlineto{\pgfqpoint{5.941060in}{1.389076in}}%
\pgfpathlineto{\pgfqpoint{5.944226in}{1.389217in}}%
\pgfpathlineto{\pgfqpoint{5.947477in}{1.391724in}}%
\pgfpathlineto{\pgfqpoint{5.951498in}{1.397947in}}%
\pgfpathlineto{\pgfqpoint{5.956547in}{1.410250in}}%
\pgfpathlineto{\pgfqpoint{5.962964in}{1.432224in}}%
\pgfpathlineto{\pgfqpoint{5.971349in}{1.469584in}}%
\pgfpathlineto{\pgfqpoint{5.983157in}{1.533474in}}%
\pgfpathlineto{\pgfqpoint{6.025424in}{1.771223in}}%
\pgfpathlineto{\pgfqpoint{6.037574in}{1.822824in}}%
\pgfpathlineto{\pgfqpoint{6.048441in}{1.859750in}}%
\pgfpathlineto{\pgfqpoint{6.058366in}{1.886033in}}%
\pgfpathlineto{\pgfqpoint{6.067692in}{1.904763in}}%
\pgfpathlineto{\pgfqpoint{6.076591in}{1.917861in}}%
\pgfpathlineto{\pgfqpoint{6.085233in}{1.926757in}}%
\pgfpathlineto{\pgfqpoint{6.093874in}{1.932521in}}%
\pgfpathlineto{\pgfqpoint{6.102858in}{1.935862in}}%
\pgfpathlineto{\pgfqpoint{6.112955in}{1.937189in}}%
\pgfpathlineto{\pgfqpoint{6.125447in}{1.936459in}}%
\pgfpathlineto{\pgfqpoint{6.145554in}{1.932703in}}%
\pgfpathlineto{\pgfqpoint{6.180464in}{1.926390in}}%
\pgfpathlineto{\pgfqpoint{6.216229in}{1.920367in}}%
\pgfpathlineto{\pgfqpoint{6.228379in}{1.915726in}}%
\pgfpathlineto{\pgfqpoint{6.238560in}{1.909455in}}%
\pgfpathlineto{\pgfqpoint{6.247887in}{1.901054in}}%
\pgfpathlineto{\pgfqpoint{6.256956in}{1.889786in}}%
\pgfpathlineto{\pgfqpoint{6.266026in}{1.874839in}}%
\pgfpathlineto{\pgfqpoint{6.275352in}{1.855040in}}%
\pgfpathlineto{\pgfqpoint{6.285278in}{1.828496in}}%
\pgfpathlineto{\pgfqpoint{6.295973in}{1.793230in}}%
\pgfpathlineto{\pgfqpoint{6.308037in}{1.745357in}}%
\pgfpathlineto{\pgfqpoint{6.323096in}{1.675742in}}%
\pgfpathlineto{\pgfqpoint{6.356209in}{1.519536in}}%
\pgfpathlineto{\pgfqpoint{6.365193in}{1.489128in}}%
\pgfpathlineto{\pgfqpoint{6.371952in}{1.473355in}}%
\pgfpathlineto{\pgfqpoint{6.377086in}{1.466346in}}%
\pgfpathlineto{\pgfqpoint{6.381108in}{1.464216in}}%
\pgfpathlineto{\pgfqpoint{6.384444in}{1.464850in}}%
\pgfpathlineto{\pgfqpoint{6.387953in}{1.467972in}}%
\pgfpathlineto{\pgfqpoint{6.392060in}{1.474925in}}%
\pgfpathlineto{\pgfqpoint{6.397022in}{1.488152in}}%
\pgfpathlineto{\pgfqpoint{6.403012in}{1.511085in}}%
\pgfpathlineto{\pgfqpoint{6.410284in}{1.548581in}}%
\pgfpathlineto{\pgfqpoint{6.419525in}{1.609028in}}%
\pgfpathlineto{\pgfqpoint{6.436723in}{1.740253in}}%
\pgfpathlineto{\pgfqpoint{6.448445in}{1.820566in}}%
\pgfpathlineto{\pgfqpoint{6.455718in}{1.856527in}}%
\pgfpathlineto{\pgfqpoint{6.461023in}{1.873010in}}%
\pgfpathlineto{\pgfqpoint{6.464959in}{1.879030in}}%
\pgfpathlineto{\pgfqpoint{6.467697in}{1.879859in}}%
\pgfpathlineto{\pgfqpoint{6.470178in}{1.878146in}}%
\pgfpathlineto{\pgfqpoint{6.473173in}{1.872911in}}%
\pgfpathlineto{\pgfqpoint{6.477023in}{1.861101in}}%
\pgfpathlineto{\pgfqpoint{6.481900in}{1.838203in}}%
\pgfpathlineto{\pgfqpoint{6.488146in}{1.797161in}}%
\pgfpathlineto{\pgfqpoint{6.496788in}{1.723899in}}%
\pgfpathlineto{\pgfqpoint{6.518435in}{1.534192in}}%
\pgfpathlineto{\pgfqpoint{6.524339in}{1.503151in}}%
\pgfpathlineto{\pgfqpoint{6.528617in}{1.490835in}}%
\pgfpathlineto{\pgfqpoint{6.531526in}{1.487924in}}%
\pgfpathlineto{\pgfqpoint{6.533580in}{1.488652in}}%
\pgfpathlineto{\pgfqpoint{6.535890in}{1.492262in}}%
\pgfpathlineto{\pgfqpoint{6.539056in}{1.501987in}}%
\pgfpathlineto{\pgfqpoint{6.543248in}{1.523066in}}%
\pgfpathlineto{\pgfqpoint{6.548724in}{1.563228in}}%
\pgfpathlineto{\pgfqpoint{6.556596in}{1.639461in}}%
\pgfpathlineto{\pgfqpoint{6.574393in}{1.816145in}}%
\pgfpathlineto{\pgfqpoint{6.579869in}{1.848163in}}%
\pgfpathlineto{\pgfqpoint{6.583719in}{1.859555in}}%
\pgfpathlineto{\pgfqpoint{6.586201in}{1.861503in}}%
\pgfpathlineto{\pgfqpoint{6.587998in}{1.860183in}}%
\pgfpathlineto{\pgfqpoint{6.590308in}{1.855099in}}%
\pgfpathlineto{\pgfqpoint{6.593474in}{1.842058in}}%
\pgfpathlineto{\pgfqpoint{6.597752in}{1.813991in}}%
\pgfpathlineto{\pgfqpoint{6.603655in}{1.758896in}}%
\pgfpathlineto{\pgfqpoint{6.616832in}{1.607195in}}%
\pgfpathlineto{\pgfqpoint{6.624276in}{1.537962in}}%
\pgfpathlineto{\pgfqpoint{6.629068in}{1.511504in}}%
\pgfpathlineto{\pgfqpoint{6.632233in}{1.503958in}}%
\pgfpathlineto{\pgfqpoint{6.634116in}{1.503513in}}%
\pgfpathlineto{\pgfqpoint{6.635827in}{1.505775in}}%
\pgfpathlineto{\pgfqpoint{6.638308in}{1.513537in}}%
\pgfpathlineto{\pgfqpoint{6.641731in}{1.532617in}}%
\pgfpathlineto{\pgfqpoint{6.646437in}{1.572908in}}%
\pgfpathlineto{\pgfqpoint{6.653710in}{1.657304in}}%
\pgfpathlineto{\pgfqpoint{6.666458in}{1.804456in}}%
\pgfpathlineto{\pgfqpoint{6.671335in}{1.837148in}}%
\pgfpathlineto{\pgfqpoint{6.674587in}{1.847021in}}%
\pgfpathlineto{\pgfqpoint{6.676555in}{1.847908in}}%
\pgfpathlineto{\pgfqpoint{6.678180in}{1.845680in}}%
\pgfpathlineto{\pgfqpoint{6.680491in}{1.837940in}}%
\pgfpathlineto{\pgfqpoint{6.683827in}{1.817701in}}%
\pgfpathlineto{\pgfqpoint{6.688533in}{1.773340in}}%
\pgfpathlineto{\pgfqpoint{6.696405in}{1.673619in}}%
\pgfpathlineto{\pgfqpoint{6.706159in}{1.557407in}}%
\pgfpathlineto{\pgfqpoint{6.710865in}{1.524792in}}%
\pgfpathlineto{\pgfqpoint{6.713945in}{1.515904in}}%
\pgfpathlineto{\pgfqpoint{6.715657in}{1.515621in}}%
\pgfpathlineto{\pgfqpoint{6.717282in}{1.518483in}}%
\pgfpathlineto{\pgfqpoint{6.719678in}{1.528181in}}%
\pgfpathlineto{\pgfqpoint{6.723101in}{1.552673in}}%
\pgfpathlineto{\pgfqpoint{6.728063in}{1.606512in}}%
\pgfpathlineto{\pgfqpoint{6.746459in}{1.824146in}}%
\pgfpathlineto{\pgfqpoint{6.749711in}{1.836100in}}%
\pgfpathlineto{\pgfqpoint{6.751507in}{1.836994in}}%
\pgfpathlineto{\pgfqpoint{6.752962in}{1.834669in}}%
\pgfpathlineto{\pgfqpoint{6.755187in}{1.825905in}}%
\pgfpathlineto{\pgfqpoint{6.758438in}{1.802432in}}%
\pgfpathlineto{\pgfqpoint{6.763144in}{1.749929in}}%
\pgfpathlineto{\pgfqpoint{6.780684in}{1.536573in}}%
\pgfpathlineto{\pgfqpoint{6.783679in}{1.526062in}}%
\pgfpathlineto{\pgfqpoint{6.785219in}{1.525517in}}%
\pgfpathlineto{\pgfqpoint{6.785305in}{1.525586in}}%
\pgfpathlineto{\pgfqpoint{6.786759in}{1.528338in}}%
\pgfpathlineto{\pgfqpoint{6.788984in}{1.538247in}}%
\pgfpathlineto{\pgfqpoint{6.792235in}{1.564256in}}%
\pgfpathlineto{\pgfqpoint{6.797198in}{1.624764in}}%
\pgfpathlineto{\pgfqpoint{6.811401in}{1.808834in}}%
\pgfpathlineto{\pgfqpoint{6.814824in}{1.825898in}}%
\pgfpathlineto{\pgfqpoint{6.816706in}{1.827849in}}%
\pgfpathlineto{\pgfqpoint{6.816792in}{1.827808in}}%
\pgfpathlineto{\pgfqpoint{6.818075in}{1.825837in}}%
\pgfpathlineto{\pgfqpoint{6.820128in}{1.817468in}}%
\pgfpathlineto{\pgfqpoint{6.823209in}{1.793634in}}%
\pgfpathlineto{\pgfqpoint{6.827829in}{1.737469in}}%
\pgfpathlineto{\pgfqpoint{6.842289in}{1.548452in}}%
\pgfpathlineto{\pgfqpoint{6.845369in}{1.535035in}}%
\pgfpathlineto{\pgfqpoint{6.846824in}{1.534045in}}%
\pgfpathlineto{\pgfqpoint{6.846995in}{1.534160in}}%
\pgfpathlineto{\pgfqpoint{6.848364in}{1.536829in}}%
\pgfpathlineto{\pgfqpoint{6.850503in}{1.547106in}}%
\pgfpathlineto{\pgfqpoint{6.853669in}{1.574882in}}%
\pgfpathlineto{\pgfqpoint{6.858632in}{1.641088in}}%
\pgfpathlineto{\pgfqpoint{6.870011in}{1.797245in}}%
\pgfpathlineto{\pgfqpoint{6.873434in}{1.817259in}}%
\pgfpathlineto{\pgfqpoint{6.875316in}{1.819771in}}%
\pgfpathlineto{\pgfqpoint{6.875402in}{1.819736in}}%
\pgfpathlineto{\pgfqpoint{6.876600in}{1.817878in}}%
\pgfpathlineto{\pgfqpoint{6.878482in}{1.809882in}}%
\pgfpathlineto{\pgfqpoint{6.881391in}{1.786129in}}%
\pgfpathlineto{\pgfqpoint{6.885840in}{1.728689in}}%
\pgfpathlineto{\pgfqpoint{6.898418in}{1.556870in}}%
\pgfpathlineto{\pgfqpoint{6.901498in}{1.542514in}}%
\pgfpathlineto{\pgfqpoint{6.902782in}{1.541682in}}%
\pgfpathlineto{\pgfqpoint{6.903039in}{1.541889in}}%
\pgfpathlineto{\pgfqpoint{6.904408in}{1.545087in}}%
\pgfpathlineto{\pgfqpoint{6.906547in}{1.556959in}}%
\pgfpathlineto{\pgfqpoint{6.909798in}{1.589411in}}%
\pgfpathlineto{\pgfqpoint{6.915359in}{1.671670in}}%
\pgfpathlineto{\pgfqpoint{6.923573in}{1.787126in}}%
\pgfpathlineto{\pgfqpoint{6.927082in}{1.809915in}}%
\pgfpathlineto{\pgfqpoint{6.928878in}{1.812548in}}%
\pgfpathlineto{\pgfqpoint{6.929049in}{1.812464in}}%
\pgfpathlineto{\pgfqpoint{6.930247in}{1.810242in}}%
\pgfpathlineto{\pgfqpoint{6.932215in}{1.800513in}}%
\pgfpathlineto{\pgfqpoint{6.935210in}{1.772462in}}%
\pgfpathlineto{\pgfqpoint{6.940173in}{1.700857in}}%
\pgfpathlineto{\pgfqpoint{6.949328in}{1.570164in}}%
\pgfpathlineto{\pgfqpoint{6.952579in}{1.550426in}}%
\pgfpathlineto{\pgfqpoint{6.954119in}{1.548493in}}%
\pgfpathlineto{\pgfqpoint{6.954376in}{1.548653in}}%
\pgfpathlineto{\pgfqpoint{6.955574in}{1.551224in}}%
\pgfpathlineto{\pgfqpoint{6.957542in}{1.561803in}}%
\pgfpathlineto{\pgfqpoint{6.960622in}{1.592655in}}%
\pgfpathlineto{\pgfqpoint{6.965927in}{1.672921in}}%
\pgfpathlineto{\pgfqpoint{6.973713in}{1.783880in}}%
\pgfpathlineto{\pgfqpoint{6.976964in}{1.804141in}}%
\pgfpathlineto{\pgfqpoint{6.978419in}{1.806015in}}%
\pgfpathlineto{\pgfqpoint{6.978676in}{1.805864in}}%
\pgfpathlineto{\pgfqpoint{6.979874in}{1.803255in}}%
\pgfpathlineto{\pgfqpoint{6.981842in}{1.792327in}}%
\pgfpathlineto{\pgfqpoint{6.984922in}{1.760398in}}%
\pgfpathlineto{\pgfqpoint{6.990398in}{1.675396in}}%
\pgfpathlineto{\pgfqpoint{6.997414in}{1.575974in}}%
\pgfpathlineto{\pgfqpoint{7.000580in}{1.556390in}}%
\pgfpathlineto{\pgfqpoint{7.001949in}{1.554734in}}%
\pgfpathlineto{\pgfqpoint{7.002205in}{1.554901in}}%
\pgfpathlineto{\pgfqpoint{7.003403in}{1.557664in}}%
\pgfpathlineto{\pgfqpoint{7.005371in}{1.569110in}}%
\pgfpathlineto{\pgfqpoint{7.008451in}{1.602291in}}%
\pgfpathlineto{\pgfqpoint{7.014270in}{1.694834in}}%
\pgfpathlineto{\pgfqpoint{7.020516in}{1.781131in}}%
\pgfpathlineto{\pgfqpoint{7.023596in}{1.798950in}}%
\pgfpathlineto{\pgfqpoint{7.024794in}{1.799996in}}%
\pgfpathlineto{\pgfqpoint{7.025136in}{1.799668in}}%
\pgfpathlineto{\pgfqpoint{7.026420in}{1.795972in}}%
\pgfpathlineto{\pgfqpoint{7.028559in}{1.781495in}}%
\pgfpathlineto{\pgfqpoint{7.031981in}{1.739998in}}%
\pgfpathlineto{\pgfqpoint{7.044987in}{1.563826in}}%
\pgfpathlineto{\pgfqpoint{7.046783in}{1.560461in}}%
\pgfpathlineto{\pgfqpoint{7.046955in}{1.560554in}}%
\pgfpathlineto{\pgfqpoint{7.048067in}{1.562914in}}%
\pgfpathlineto{\pgfqpoint{7.049949in}{1.573654in}}%
\pgfpathlineto{\pgfqpoint{7.052944in}{1.606215in}}%
\pgfpathlineto{\pgfqpoint{7.058762in}{1.700297in}}%
\pgfpathlineto{\pgfqpoint{7.064495in}{1.778348in}}%
\pgfpathlineto{\pgfqpoint{7.067404in}{1.793925in}}%
\pgfpathlineto{\pgfqpoint{7.068345in}{1.794479in}}%
\pgfpathlineto{\pgfqpoint{7.068687in}{1.794121in}}%
\pgfpathlineto{\pgfqpoint{7.069971in}{1.790140in}}%
\pgfpathlineto{\pgfqpoint{7.072110in}{1.774657in}}%
\pgfpathlineto{\pgfqpoint{7.075532in}{1.730839in}}%
\pgfpathlineto{\pgfqpoint{7.086998in}{1.570873in}}%
\pgfpathlineto{\pgfqpoint{7.089051in}{1.565744in}}%
\pgfpathlineto{\pgfqpoint{7.089137in}{1.565769in}}%
\pgfpathlineto{\pgfqpoint{7.090078in}{1.567315in}}%
\pgfpathlineto{\pgfqpoint{7.091704in}{1.575355in}}%
\pgfpathlineto{\pgfqpoint{7.094356in}{1.601738in}}%
\pgfpathlineto{\pgfqpoint{7.099148in}{1.676871in}}%
\pgfpathlineto{\pgfqpoint{7.105822in}{1.773320in}}%
\pgfpathlineto{\pgfqpoint{7.108645in}{1.788830in}}%
\pgfpathlineto{\pgfqpoint{7.109501in}{1.789400in}}%
\pgfpathlineto{\pgfqpoint{7.109929in}{1.788942in}}%
\pgfpathlineto{\pgfqpoint{7.111212in}{1.784626in}}%
\pgfpathlineto{\pgfqpoint{7.113351in}{1.768119in}}%
\pgfpathlineto{\pgfqpoint{7.116859in}{1.720815in}}%
\pgfpathlineto{\pgfqpoint{7.126784in}{1.578373in}}%
\pgfpathlineto{\pgfqpoint{7.129095in}{1.570672in}}%
\pgfpathlineto{\pgfqpoint{7.129950in}{1.571548in}}%
\pgfpathlineto{\pgfqpoint{7.131405in}{1.577642in}}%
\pgfpathlineto{\pgfqpoint{7.133800in}{1.599380in}}%
\pgfpathlineto{\pgfqpoint{7.137907in}{1.661057in}}%
\pgfpathlineto{\pgfqpoint{7.145437in}{1.771439in}}%
\pgfpathlineto{\pgfqpoint{7.148089in}{1.784418in}}%
\pgfpathlineto{\pgfqpoint{7.148859in}{1.784517in}}%
\pgfpathlineto{\pgfqpoint{7.149202in}{1.784019in}}%
\pgfpathlineto{\pgfqpoint{7.150571in}{1.778735in}}%
\pgfpathlineto{\pgfqpoint{7.152881in}{1.758657in}}%
\pgfpathlineto{\pgfqpoint{7.156817in}{1.700607in}}%
\pgfpathlineto{\pgfqpoint{7.164517in}{1.587521in}}%
\pgfpathlineto{\pgfqpoint{7.167084in}{1.575464in}}%
\pgfpathlineto{\pgfqpoint{7.167854in}{1.575518in}}%
\pgfpathlineto{\pgfqpoint{7.168111in}{1.575920in}}%
\pgfpathlineto{\pgfqpoint{7.169480in}{1.581262in}}%
\pgfpathlineto{\pgfqpoint{7.171790in}{1.601665in}}%
\pgfpathlineto{\pgfqpoint{7.175726in}{1.660398in}}%
\pgfpathlineto{\pgfqpoint{7.183084in}{1.768162in}}%
\pgfpathlineto{\pgfqpoint{7.185651in}{1.780080in}}%
\pgfpathlineto{\pgfqpoint{7.186421in}{1.779901in}}%
\pgfpathlineto{\pgfqpoint{7.186592in}{1.779621in}}%
\pgfpathlineto{\pgfqpoint{7.187876in}{1.774781in}}%
\pgfpathlineto{\pgfqpoint{7.190101in}{1.755653in}}%
\pgfpathlineto{\pgfqpoint{7.193951in}{1.698819in}}%
\pgfpathlineto{\pgfqpoint{7.201395in}{1.590382in}}%
\pgfpathlineto{\pgfqpoint{7.203876in}{1.579663in}}%
\pgfpathlineto{\pgfqpoint{7.204732in}{1.580239in}}%
\pgfpathlineto{\pgfqpoint{7.204817in}{1.580419in}}%
\pgfpathlineto{\pgfqpoint{7.206186in}{1.586265in}}%
\pgfpathlineto{\pgfqpoint{7.208496in}{1.607849in}}%
\pgfpathlineto{\pgfqpoint{7.212689in}{1.672692in}}%
\pgfpathlineto{\pgfqpoint{7.219021in}{1.763966in}}%
\pgfpathlineto{\pgfqpoint{7.221588in}{1.775928in}}%
\pgfpathlineto{\pgfqpoint{7.222358in}{1.775611in}}%
\pgfpathlineto{\pgfqpoint{7.222529in}{1.775291in}}%
\pgfpathlineto{\pgfqpoint{7.223898in}{1.769523in}}%
\pgfpathlineto{\pgfqpoint{7.226208in}{1.747854in}}%
\pgfpathlineto{\pgfqpoint{7.230400in}{1.682715in}}%
\pgfpathlineto{\pgfqpoint{7.236561in}{1.594854in}}%
\pgfpathlineto{\pgfqpoint{7.239042in}{1.583737in}}%
\pgfpathlineto{\pgfqpoint{7.239898in}{1.584333in}}%
\pgfpathlineto{\pgfqpoint{7.241181in}{1.589508in}}%
\pgfpathlineto{\pgfqpoint{7.243406in}{1.609710in}}%
\pgfpathlineto{\pgfqpoint{7.247427in}{1.671252in}}%
\pgfpathlineto{\pgfqpoint{7.253759in}{1.761653in}}%
\pgfpathlineto{\pgfqpoint{7.256155in}{1.772021in}}%
\pgfpathlineto{\pgfqpoint{7.257010in}{1.771330in}}%
\pgfpathlineto{\pgfqpoint{7.258379in}{1.765406in}}%
\pgfpathlineto{\pgfqpoint{7.260690in}{1.743148in}}%
\pgfpathlineto{\pgfqpoint{7.265053in}{1.674266in}}%
\pgfpathlineto{\pgfqpoint{7.270615in}{1.597404in}}%
\pgfpathlineto{\pgfqpoint{7.273011in}{1.587529in}}%
\pgfpathlineto{\pgfqpoint{7.273866in}{1.588473in}}%
\pgfpathlineto{\pgfqpoint{7.275235in}{1.594861in}}%
\pgfpathlineto{\pgfqpoint{7.277631in}{1.619052in}}%
\pgfpathlineto{\pgfqpoint{7.282508in}{1.697689in}}%
\pgfpathlineto{\pgfqpoint{7.287299in}{1.759915in}}%
\pgfpathlineto{\pgfqpoint{7.289524in}{1.768326in}}%
\pgfpathlineto{\pgfqpoint{7.290380in}{1.767251in}}%
\pgfpathlineto{\pgfqpoint{7.291834in}{1.759982in}}%
\pgfpathlineto{\pgfqpoint{7.294316in}{1.733550in}}%
\pgfpathlineto{\pgfqpoint{7.300818in}{1.629066in}}%
\pgfpathlineto{\pgfqpoint{7.304412in}{1.594198in}}%
\pgfpathlineto{\pgfqpoint{7.305867in}{1.591176in}}%
\pgfpathlineto{\pgfqpoint{7.306123in}{1.591377in}}%
\pgfpathlineto{\pgfqpoint{7.307150in}{1.594376in}}%
\pgfpathlineto{\pgfqpoint{7.309032in}{1.608527in}}%
\pgfpathlineto{\pgfqpoint{7.312284in}{1.653438in}}%
\pgfpathlineto{\pgfqpoint{7.319728in}{1.758370in}}%
\pgfpathlineto{\pgfqpoint{7.321696in}{1.764781in}}%
\pgfpathlineto{\pgfqpoint{7.321781in}{1.764765in}}%
\pgfpathlineto{\pgfqpoint{7.322637in}{1.763234in}}%
\pgfpathlineto{\pgfqpoint{7.324177in}{1.754401in}}%
\pgfpathlineto{\pgfqpoint{7.326829in}{1.723486in}}%
\pgfpathlineto{\pgfqpoint{7.336583in}{1.595610in}}%
\pgfpathlineto{\pgfqpoint{7.337439in}{1.594631in}}%
\pgfpathlineto{\pgfqpoint{7.337867in}{1.595085in}}%
\pgfpathlineto{\pgfqpoint{7.339065in}{1.599663in}}%
\pgfpathlineto{\pgfqpoint{7.341204in}{1.619001in}}%
\pgfpathlineto{\pgfqpoint{7.345225in}{1.680743in}}%
\pgfpathlineto{\pgfqpoint{7.350616in}{1.753532in}}%
\pgfpathlineto{\pgfqpoint{7.352840in}{1.761400in}}%
\pgfpathlineto{\pgfqpoint{7.353696in}{1.759891in}}%
\pgfpathlineto{\pgfqpoint{7.355236in}{1.750963in}}%
\pgfpathlineto{\pgfqpoint{7.357889in}{1.719683in}}%
\pgfpathlineto{\pgfqpoint{7.367129in}{1.599060in}}%
\pgfpathlineto{\pgfqpoint{7.368070in}{1.597933in}}%
\pgfpathlineto{\pgfqpoint{7.368498in}{1.598450in}}%
\pgfpathlineto{\pgfqpoint{7.369696in}{1.603266in}}%
\pgfpathlineto{\pgfqpoint{7.371835in}{1.623165in}}%
\pgfpathlineto{\pgfqpoint{7.376113in}{1.689543in}}%
\pgfpathlineto{\pgfqpoint{7.380905in}{1.751265in}}%
\pgfpathlineto{\pgfqpoint{7.382958in}{1.758194in}}%
\pgfpathlineto{\pgfqpoint{7.383044in}{1.758162in}}%
\pgfpathlineto{\pgfqpoint{7.383900in}{1.756416in}}%
\pgfpathlineto{\pgfqpoint{7.385525in}{1.746221in}}%
\pgfpathlineto{\pgfqpoint{7.388349in}{1.710958in}}%
\pgfpathlineto{\pgfqpoint{7.396392in}{1.603784in}}%
\pgfpathlineto{\pgfqpoint{7.397675in}{1.601040in}}%
\pgfpathlineto{\pgfqpoint{7.398017in}{1.601302in}}%
\pgfpathlineto{\pgfqpoint{7.399130in}{1.605021in}}%
\pgfpathlineto{\pgfqpoint{7.401098in}{1.621596in}}%
\pgfpathlineto{\pgfqpoint{7.404862in}{1.677354in}}%
\pgfpathlineto{\pgfqpoint{7.410253in}{1.748643in}}%
\pgfpathlineto{\pgfqpoint{7.412221in}{1.755135in}}%
\pgfpathlineto{\pgfqpoint{7.412306in}{1.755103in}}%
\pgfpathlineto{\pgfqpoint{7.413162in}{1.753339in}}%
\pgfpathlineto{\pgfqpoint{7.414788in}{1.743005in}}%
\pgfpathlineto{\pgfqpoint{7.417611in}{1.707452in}}%
\pgfpathlineto{\pgfqpoint{7.425226in}{1.606923in}}%
\pgfpathlineto{\pgfqpoint{7.426595in}{1.604048in}}%
\pgfpathlineto{\pgfqpoint{7.426937in}{1.604391in}}%
\pgfpathlineto{\pgfqpoint{7.428050in}{1.608416in}}%
\pgfpathlineto{\pgfqpoint{7.430103in}{1.626587in}}%
\pgfpathlineto{\pgfqpoint{7.434210in}{1.688694in}}%
\pgfpathlineto{\pgfqpoint{7.438831in}{1.746375in}}%
\pgfpathlineto{\pgfqpoint{7.440627in}{1.752218in}}%
\pgfpathlineto{\pgfqpoint{7.440799in}{1.752163in}}%
\pgfpathlineto{\pgfqpoint{7.441654in}{1.750286in}}%
\pgfpathlineto{\pgfqpoint{7.443280in}{1.739662in}}%
\pgfpathlineto{\pgfqpoint{7.446189in}{1.702381in}}%
\pgfpathlineto{\pgfqpoint{7.453205in}{1.610333in}}%
\pgfpathlineto{\pgfqpoint{7.454660in}{1.606899in}}%
\pgfpathlineto{\pgfqpoint{7.454916in}{1.607100in}}%
\pgfpathlineto{\pgfqpoint{7.455943in}{1.610313in}}%
\pgfpathlineto{\pgfqpoint{7.457825in}{1.625537in}}%
\pgfpathlineto{\pgfqpoint{7.461419in}{1.677359in}}%
\pgfpathlineto{\pgfqpoint{7.466553in}{1.743519in}}%
\pgfpathlineto{\pgfqpoint{7.468350in}{1.749422in}}%
\pgfpathlineto{\pgfqpoint{7.468521in}{1.749366in}}%
\pgfpathlineto{\pgfqpoint{7.469376in}{1.747459in}}%
\pgfpathlineto{\pgfqpoint{7.471002in}{1.736708in}}%
\pgfpathlineto{\pgfqpoint{7.473911in}{1.699277in}}%
\pgfpathlineto{\pgfqpoint{7.480500in}{1.613490in}}%
\pgfpathlineto{\pgfqpoint{7.481954in}{1.609621in}}%
\pgfpathlineto{\pgfqpoint{7.482296in}{1.609853in}}%
\pgfpathlineto{\pgfqpoint{7.483323in}{1.613147in}}%
\pgfpathlineto{\pgfqpoint{7.485205in}{1.628577in}}%
\pgfpathlineto{\pgfqpoint{7.488970in}{1.683071in}}%
\pgfpathlineto{\pgfqpoint{7.493762in}{1.741909in}}%
\pgfpathlineto{\pgfqpoint{7.495387in}{1.746745in}}%
\pgfpathlineto{\pgfqpoint{7.495644in}{1.746609in}}%
\pgfpathlineto{\pgfqpoint{7.496585in}{1.744009in}}%
\pgfpathlineto{\pgfqpoint{7.498382in}{1.730442in}}%
\pgfpathlineto{\pgfqpoint{7.501805in}{1.682834in}}%
\pgfpathlineto{\pgfqpoint{7.507024in}{1.617140in}}%
\pgfpathlineto{\pgfqpoint{7.508650in}{1.612242in}}%
\pgfpathlineto{\pgfqpoint{7.508906in}{1.612372in}}%
\pgfpathlineto{\pgfqpoint{7.509847in}{1.614959in}}%
\pgfpathlineto{\pgfqpoint{7.511644in}{1.628532in}}%
\pgfpathlineto{\pgfqpoint{7.515067in}{1.676026in}}%
\pgfpathlineto{\pgfqpoint{7.520201in}{1.739637in}}%
\pgfpathlineto{\pgfqpoint{7.521826in}{1.744170in}}%
\pgfpathlineto{\pgfqpoint{7.522083in}{1.743977in}}%
\pgfpathlineto{\pgfqpoint{7.523110in}{1.740738in}}%
\pgfpathlineto{\pgfqpoint{7.524992in}{1.725316in}}%
\pgfpathlineto{\pgfqpoint{7.528842in}{1.670081in}}%
\pgfpathlineto{\pgfqpoint{7.533206in}{1.619063in}}%
\pgfpathlineto{\pgfqpoint{7.534746in}{1.614757in}}%
\pgfpathlineto{\pgfqpoint{7.535003in}{1.614909in}}%
\pgfpathlineto{\pgfqpoint{7.535944in}{1.617594in}}%
\pgfpathlineto{\pgfqpoint{7.537741in}{1.631370in}}%
\pgfpathlineto{\pgfqpoint{7.541249in}{1.680088in}}%
\pgfpathlineto{\pgfqpoint{7.546040in}{1.737503in}}%
\pgfpathlineto{\pgfqpoint{7.547580in}{1.741712in}}%
\pgfpathlineto{\pgfqpoint{7.547837in}{1.741541in}}%
\pgfpathlineto{\pgfqpoint{7.548864in}{1.738366in}}%
\pgfpathlineto{\pgfqpoint{7.550746in}{1.723016in}}%
\pgfpathlineto{\pgfqpoint{7.554682in}{1.667043in}}%
\pgfpathlineto{\pgfqpoint{7.558875in}{1.620538in}}%
\pgfpathlineto{\pgfqpoint{7.560244in}{1.617171in}}%
\pgfpathlineto{\pgfqpoint{7.560586in}{1.617443in}}%
\pgfpathlineto{\pgfqpoint{7.561613in}{1.620916in}}%
\pgfpathlineto{\pgfqpoint{7.563581in}{1.637709in}}%
\pgfpathlineto{\pgfqpoint{7.568201in}{1.703836in}}%
\pgfpathlineto{\pgfqpoint{7.571709in}{1.737312in}}%
\pgfpathlineto{\pgfqpoint{7.572821in}{1.739340in}}%
\pgfpathlineto{\pgfqpoint{7.573249in}{1.738859in}}%
\pgfpathlineto{\pgfqpoint{7.574447in}{1.733857in}}%
\pgfpathlineto{\pgfqpoint{7.576586in}{1.713146in}}%
\pgfpathlineto{\pgfqpoint{7.585057in}{1.619502in}}%
\pgfpathlineto{\pgfqpoint{7.585827in}{1.620381in}}%
\pgfpathlineto{\pgfqpoint{7.587196in}{1.627401in}}%
\pgfpathlineto{\pgfqpoint{7.589677in}{1.654790in}}%
\pgfpathlineto{\pgfqpoint{7.596608in}{1.735820in}}%
\pgfpathlineto{\pgfqpoint{7.597463in}{1.737074in}}%
\pgfpathlineto{\pgfqpoint{7.597977in}{1.736475in}}%
\pgfpathlineto{\pgfqpoint{7.599175in}{1.731224in}}%
\pgfpathlineto{\pgfqpoint{7.601399in}{1.709114in}}%
\pgfpathlineto{\pgfqpoint{7.609271in}{1.621881in}}%
\pgfpathlineto{\pgfqpoint{7.609870in}{1.621898in}}%
\pgfpathlineto{\pgfqpoint{7.610127in}{1.622328in}}%
\pgfpathlineto{\pgfqpoint{7.611325in}{1.627599in}}%
\pgfpathlineto{\pgfqpoint{7.613549in}{1.649708in}}%
\pgfpathlineto{\pgfqpoint{7.621250in}{1.734671in}}%
\pgfpathlineto{\pgfqpoint{7.621849in}{1.734759in}}%
\pgfpathlineto{\pgfqpoint{7.622191in}{1.734188in}}%
\pgfpathlineto{\pgfqpoint{7.623474in}{1.728152in}}%
\pgfpathlineto{\pgfqpoint{7.625870in}{1.703069in}}%
\pgfpathlineto{\pgfqpoint{7.632801in}{1.624784in}}%
\pgfpathlineto{\pgfqpoint{7.633571in}{1.623874in}}%
\pgfpathlineto{\pgfqpoint{7.634084in}{1.624538in}}%
\pgfpathlineto{\pgfqpoint{7.635368in}{1.630520in}}%
\pgfpathlineto{\pgfqpoint{7.637763in}{1.655478in}}%
\pgfpathlineto{\pgfqpoint{7.644608in}{1.731912in}}%
\pgfpathlineto{\pgfqpoint{7.645378in}{1.732767in}}%
\pgfpathlineto{\pgfqpoint{7.645806in}{1.732253in}}%
\pgfpathlineto{\pgfqpoint{7.647004in}{1.727149in}}%
\pgfpathlineto{\pgfqpoint{7.649229in}{1.705389in}}%
\pgfpathlineto{\pgfqpoint{7.656587in}{1.626194in}}%
\pgfpathlineto{\pgfqpoint{7.657186in}{1.626030in}}%
\pgfpathlineto{\pgfqpoint{7.657528in}{1.626557in}}%
\pgfpathlineto{\pgfqpoint{7.658812in}{1.632435in}}%
\pgfpathlineto{\pgfqpoint{7.661122in}{1.656023in}}%
\pgfpathlineto{\pgfqpoint{7.667881in}{1.729960in}}%
\pgfpathlineto{\pgfqpoint{7.668651in}{1.730726in}}%
\pgfpathlineto{\pgfqpoint{7.669079in}{1.730163in}}%
\pgfpathlineto{\pgfqpoint{7.670277in}{1.724933in}}%
\pgfpathlineto{\pgfqpoint{7.672502in}{1.703104in}}%
\pgfpathlineto{\pgfqpoint{7.679518in}{1.628338in}}%
\pgfpathlineto{\pgfqpoint{7.680117in}{1.627954in}}%
\pgfpathlineto{\pgfqpoint{7.680545in}{1.628526in}}%
\pgfpathlineto{\pgfqpoint{7.681742in}{1.633774in}}%
\pgfpathlineto{\pgfqpoint{7.684053in}{1.656642in}}%
\pgfpathlineto{\pgfqpoint{7.690812in}{1.728267in}}%
\pgfpathlineto{\pgfqpoint{7.691497in}{1.728754in}}%
\pgfpathlineto{\pgfqpoint{7.691924in}{1.728143in}}%
\pgfpathlineto{\pgfqpoint{7.693208in}{1.722219in}}%
\pgfpathlineto{\pgfqpoint{7.695604in}{1.697689in}}%
\pgfpathlineto{\pgfqpoint{7.701850in}{1.630900in}}%
\pgfpathlineto{\pgfqpoint{7.702705in}{1.629864in}}%
\pgfpathlineto{\pgfqpoint{7.703133in}{1.630401in}}%
\pgfpathlineto{\pgfqpoint{7.704331in}{1.635542in}}%
\pgfpathlineto{\pgfqpoint{7.706556in}{1.657062in}}%
\pgfpathlineto{\pgfqpoint{7.713229in}{1.726402in}}%
\pgfpathlineto{\pgfqpoint{7.713914in}{1.726864in}}%
\pgfpathlineto{\pgfqpoint{7.714342in}{1.726239in}}%
\pgfpathlineto{\pgfqpoint{7.715625in}{1.720297in}}%
\pgfpathlineto{\pgfqpoint{7.718021in}{1.695923in}}%
\pgfpathlineto{\pgfqpoint{7.724096in}{1.632654in}}%
\pgfpathlineto{\pgfqpoint{7.724866in}{1.631698in}}%
\pgfpathlineto{\pgfqpoint{7.725379in}{1.632322in}}%
\pgfpathlineto{\pgfqpoint{7.726663in}{1.638156in}}%
\pgfpathlineto{\pgfqpoint{7.729058in}{1.662299in}}%
\pgfpathlineto{\pgfqpoint{7.735048in}{1.724106in}}%
\pgfpathlineto{\pgfqpoint{7.735903in}{1.725063in}}%
\pgfpathlineto{\pgfqpoint{7.736331in}{1.724493in}}%
\pgfpathlineto{\pgfqpoint{7.737529in}{1.719292in}}%
\pgfpathlineto{\pgfqpoint{7.739839in}{1.696883in}}%
\pgfpathlineto{\pgfqpoint{7.746000in}{1.634258in}}%
\pgfpathlineto{\pgfqpoint{7.746770in}{1.633510in}}%
\pgfpathlineto{\pgfqpoint{7.747198in}{1.634071in}}%
\pgfpathlineto{\pgfqpoint{7.748396in}{1.639236in}}%
\pgfpathlineto{\pgfqpoint{7.750706in}{1.661506in}}%
\pgfpathlineto{\pgfqpoint{7.756781in}{1.722564in}}%
\pgfpathlineto{\pgfqpoint{7.757551in}{1.723310in}}%
\pgfpathlineto{\pgfqpoint{7.757979in}{1.722752in}}%
\pgfpathlineto{\pgfqpoint{7.759176in}{1.717604in}}%
\pgfpathlineto{\pgfqpoint{7.761487in}{1.695455in}}%
\pgfpathlineto{\pgfqpoint{7.767476in}{1.635974in}}%
\pgfpathlineto{\pgfqpoint{7.768246in}{1.635236in}}%
\pgfpathlineto{\pgfqpoint{7.768674in}{1.635797in}}%
\pgfpathlineto{\pgfqpoint{7.769872in}{1.640935in}}%
\pgfpathlineto{\pgfqpoint{7.772182in}{1.662973in}}%
\pgfpathlineto{\pgfqpoint{7.778086in}{1.720885in}}%
\pgfpathlineto{\pgfqpoint{7.778856in}{1.721613in}}%
\pgfpathlineto{\pgfqpoint{7.779284in}{1.721050in}}%
\pgfpathlineto{\pgfqpoint{7.780567in}{1.715364in}}%
\pgfpathlineto{\pgfqpoint{7.782963in}{1.691948in}}%
\pgfpathlineto{\pgfqpoint{7.788524in}{1.637842in}}%
\pgfpathlineto{\pgfqpoint{7.789380in}{1.636904in}}%
\pgfpathlineto{\pgfqpoint{7.789808in}{1.637467in}}%
\pgfpathlineto{\pgfqpoint{7.791091in}{1.643134in}}%
\pgfpathlineto{\pgfqpoint{7.793573in}{1.667437in}}%
\pgfpathlineto{\pgfqpoint{7.798877in}{1.718791in}}%
\pgfpathlineto{\pgfqpoint{7.799819in}{1.719976in}}%
\pgfpathlineto{\pgfqpoint{7.800246in}{1.719418in}}%
\pgfpathlineto{\pgfqpoint{7.801530in}{1.713785in}}%
\pgfpathlineto{\pgfqpoint{7.804011in}{1.689657in}}%
\pgfpathlineto{\pgfqpoint{7.809230in}{1.639713in}}%
\pgfpathlineto{\pgfqpoint{7.810172in}{1.638509in}}%
\pgfpathlineto{\pgfqpoint{7.810599in}{1.639053in}}%
\pgfpathlineto{\pgfqpoint{7.811797in}{1.644083in}}%
\pgfpathlineto{\pgfqpoint{7.814108in}{1.665545in}}%
\pgfpathlineto{\pgfqpoint{7.819669in}{1.717626in}}%
\pgfpathlineto{\pgfqpoint{7.820439in}{1.718403in}}%
\pgfpathlineto{\pgfqpoint{7.820867in}{1.717884in}}%
\pgfpathlineto{\pgfqpoint{7.822065in}{1.712936in}}%
\pgfpathlineto{\pgfqpoint{7.824375in}{1.691696in}}%
\pgfpathlineto{\pgfqpoint{7.829851in}{1.640887in}}%
\pgfpathlineto{\pgfqpoint{7.830621in}{1.640050in}}%
\pgfpathlineto{\pgfqpoint{7.831134in}{1.640709in}}%
\pgfpathlineto{\pgfqpoint{7.832418in}{1.646445in}}%
\pgfpathlineto{\pgfqpoint{7.834985in}{1.671265in}}%
\pgfpathlineto{\pgfqpoint{7.839862in}{1.715731in}}%
\pgfpathlineto{\pgfqpoint{7.840717in}{1.716891in}}%
\pgfpathlineto{\pgfqpoint{7.841231in}{1.716296in}}%
\pgfpathlineto{\pgfqpoint{7.842514in}{1.710731in}}%
\pgfpathlineto{\pgfqpoint{7.844996in}{1.687260in}}%
\pgfpathlineto{\pgfqpoint{7.849958in}{1.642581in}}%
\pgfpathlineto{\pgfqpoint{7.850814in}{1.641557in}}%
\pgfpathlineto{\pgfqpoint{7.851242in}{1.642049in}}%
\pgfpathlineto{\pgfqpoint{7.852440in}{1.646865in}}%
\pgfpathlineto{\pgfqpoint{7.854835in}{1.668524in}}%
\pgfpathlineto{\pgfqpoint{7.860055in}{1.714671in}}%
\pgfpathlineto{\pgfqpoint{7.860825in}{1.715388in}}%
\pgfpathlineto{\pgfqpoint{7.861252in}{1.714855in}}%
\pgfpathlineto{\pgfqpoint{7.862536in}{1.709420in}}%
\pgfpathlineto{\pgfqpoint{7.865017in}{1.686407in}}%
\pgfpathlineto{\pgfqpoint{7.869894in}{1.643910in}}%
\pgfpathlineto{\pgfqpoint{7.870750in}{1.643030in}}%
\pgfpathlineto{\pgfqpoint{7.871178in}{1.643582in}}%
\pgfpathlineto{\pgfqpoint{7.872461in}{1.649044in}}%
\pgfpathlineto{\pgfqpoint{7.875028in}{1.672884in}}%
\pgfpathlineto{\pgfqpoint{7.879648in}{1.712841in}}%
\pgfpathlineto{\pgfqpoint{7.880590in}{1.713953in}}%
\pgfpathlineto{\pgfqpoint{7.881017in}{1.713407in}}%
\pgfpathlineto{\pgfqpoint{7.882301in}{1.707987in}}%
\pgfpathlineto{\pgfqpoint{7.884868in}{1.684368in}}%
\pgfpathlineto{\pgfqpoint{7.889488in}{1.645359in}}%
\pgfpathlineto{\pgfqpoint{7.890344in}{1.644426in}}%
\pgfpathlineto{\pgfqpoint{7.890772in}{1.644940in}}%
\pgfpathlineto{\pgfqpoint{7.892055in}{1.650242in}}%
\pgfpathlineto{\pgfqpoint{7.894536in}{1.672623in}}%
\pgfpathlineto{\pgfqpoint{7.899242in}{1.711772in}}%
\pgfpathlineto{\pgfqpoint{7.900012in}{1.712593in}}%
\pgfpathlineto{\pgfqpoint{7.900526in}{1.711971in}}%
\pgfpathlineto{\pgfqpoint{7.901809in}{1.706504in}}%
\pgfpathlineto{\pgfqpoint{7.904376in}{1.683190in}}%
\pgfpathlineto{\pgfqpoint{7.908825in}{1.646731in}}%
\pgfpathlineto{\pgfqpoint{7.909681in}{1.645782in}}%
\pgfpathlineto{\pgfqpoint{7.910109in}{1.646275in}}%
\pgfpathlineto{\pgfqpoint{7.911306in}{1.650956in}}%
\pgfpathlineto{\pgfqpoint{7.913702in}{1.671574in}}%
\pgfpathlineto{\pgfqpoint{7.918494in}{1.710527in}}%
\pgfpathlineto{\pgfqpoint{7.919264in}{1.711242in}}%
\pgfpathlineto{\pgfqpoint{7.919692in}{1.710742in}}%
\pgfpathlineto{\pgfqpoint{7.920975in}{1.705563in}}%
\pgfpathlineto{\pgfqpoint{7.923542in}{1.682928in}}%
\pgfpathlineto{\pgfqpoint{7.927906in}{1.648055in}}%
\pgfpathlineto{\pgfqpoint{7.928761in}{1.647095in}}%
\pgfpathlineto{\pgfqpoint{7.929189in}{1.647568in}}%
\pgfpathlineto{\pgfqpoint{7.930387in}{1.652148in}}%
\pgfpathlineto{\pgfqpoint{7.932783in}{1.672334in}}%
\pgfpathlineto{\pgfqpoint{7.937489in}{1.709312in}}%
\pgfpathlineto{\pgfqpoint{7.938259in}{1.709933in}}%
\pgfpathlineto{\pgfqpoint{7.938686in}{1.709394in}}%
\pgfpathlineto{\pgfqpoint{7.939970in}{1.704166in}}%
\pgfpathlineto{\pgfqpoint{7.942622in}{1.680978in}}%
\pgfpathlineto{\pgfqpoint{7.946815in}{1.649146in}}%
\pgfpathlineto{\pgfqpoint{7.947671in}{1.648401in}}%
\pgfpathlineto{\pgfqpoint{7.948098in}{1.648968in}}%
\pgfpathlineto{\pgfqpoint{7.949382in}{1.654242in}}%
\pgfpathlineto{\pgfqpoint{7.952034in}{1.677274in}}%
\pgfpathlineto{\pgfqpoint{7.956141in}{1.707921in}}%
\pgfpathlineto{\pgfqpoint{7.956997in}{1.708670in}}%
\pgfpathlineto{\pgfqpoint{7.957425in}{1.708113in}}%
\pgfpathlineto{\pgfqpoint{7.958708in}{1.702896in}}%
\pgfpathlineto{\pgfqpoint{7.961361in}{1.680132in}}%
\pgfpathlineto{\pgfqpoint{7.965382in}{1.650445in}}%
\pgfpathlineto{\pgfqpoint{7.966238in}{1.649615in}}%
\pgfpathlineto{\pgfqpoint{7.966665in}{1.650125in}}%
\pgfpathlineto{\pgfqpoint{7.967949in}{1.655179in}}%
\pgfpathlineto{\pgfqpoint{7.970601in}{1.677566in}}%
\pgfpathlineto{\pgfqpoint{7.974623in}{1.706716in}}%
\pgfpathlineto{\pgfqpoint{7.975478in}{1.707460in}}%
\pgfpathlineto{\pgfqpoint{7.975906in}{1.706915in}}%
\pgfpathlineto{\pgfqpoint{7.977190in}{1.701798in}}%
\pgfpathlineto{\pgfqpoint{7.979928in}{1.678727in}}%
\pgfpathlineto{\pgfqpoint{7.983778in}{1.651562in}}%
\pgfpathlineto{\pgfqpoint{7.984634in}{1.650821in}}%
\pgfpathlineto{\pgfqpoint{7.985061in}{1.651360in}}%
\pgfpathlineto{\pgfqpoint{7.986345in}{1.656426in}}%
\pgfpathlineto{\pgfqpoint{7.989083in}{1.679220in}}%
\pgfpathlineto{\pgfqpoint{7.992933in}{1.705672in}}%
\pgfpathlineto{\pgfqpoint{7.993703in}{1.706303in}}%
\pgfpathlineto{\pgfqpoint{7.994131in}{1.705812in}}%
\pgfpathlineto{\pgfqpoint{7.995414in}{1.700911in}}%
\pgfpathlineto{\pgfqpoint{7.998067in}{1.679315in}}%
\pgfpathlineto{\pgfqpoint{8.002003in}{1.652537in}}%
\pgfpathlineto{\pgfqpoint{8.002773in}{1.651972in}}%
\pgfpathlineto{\pgfqpoint{8.003201in}{1.652491in}}%
\pgfpathlineto{\pgfqpoint{8.004484in}{1.657436in}}%
\pgfpathlineto{\pgfqpoint{8.007222in}{1.679653in}}%
\pgfpathlineto{\pgfqpoint{8.010987in}{1.704580in}}%
\pgfpathlineto{\pgfqpoint{8.011757in}{1.705160in}}%
\pgfpathlineto{\pgfqpoint{8.012185in}{1.704656in}}%
\pgfpathlineto{\pgfqpoint{8.013468in}{1.699787in}}%
\pgfpathlineto{\pgfqpoint{8.016206in}{1.677876in}}%
\pgfpathlineto{\pgfqpoint{8.019971in}{1.653572in}}%
\pgfpathlineto{\pgfqpoint{8.020741in}{1.653110in}}%
\pgfpathlineto{\pgfqpoint{8.021169in}{1.653671in}}%
\pgfpathlineto{\pgfqpoint{8.022538in}{1.659151in}}%
\pgfpathlineto{\pgfqpoint{8.025532in}{1.683503in}}%
\pgfpathlineto{\pgfqpoint{8.028869in}{1.703599in}}%
\pgfpathlineto{\pgfqpoint{8.029639in}{1.704028in}}%
\pgfpathlineto{\pgfqpoint{8.030067in}{1.703456in}}%
\pgfpathlineto{\pgfqpoint{8.031436in}{1.697987in}}%
\pgfpathlineto{\pgfqpoint{8.034516in}{1.673190in}}%
\pgfpathlineto{\pgfqpoint{8.037768in}{1.654518in}}%
\pgfpathlineto{\pgfqpoint{8.038452in}{1.654185in}}%
\pgfpathlineto{\pgfqpoint{8.038880in}{1.654721in}}%
\pgfpathlineto{\pgfqpoint{8.040249in}{1.660049in}}%
\pgfpathlineto{\pgfqpoint{8.043329in}{1.684437in}}%
\pgfpathlineto{\pgfqpoint{8.046581in}{1.702691in}}%
\pgfpathlineto{\pgfqpoint{8.047265in}{1.702961in}}%
\pgfpathlineto{\pgfqpoint{8.047693in}{1.702393in}}%
\pgfpathlineto{\pgfqpoint{8.049062in}{1.697016in}}%
\pgfpathlineto{\pgfqpoint{8.052228in}{1.672192in}}%
\pgfpathlineto{\pgfqpoint{8.055394in}{1.655413in}}%
\pgfpathlineto{\pgfqpoint{8.056078in}{1.655309in}}%
\pgfpathlineto{\pgfqpoint{8.056420in}{1.655796in}}%
\pgfpathlineto{\pgfqpoint{8.057789in}{1.661085in}}%
\pgfpathlineto{\pgfqpoint{8.060955in}{1.685528in}}%
\pgfpathlineto{\pgfqpoint{8.064035in}{1.701713in}}%
\pgfpathlineto{\pgfqpoint{8.064720in}{1.701908in}}%
\pgfpathlineto{\pgfqpoint{8.065062in}{1.701472in}}%
\pgfpathlineto{\pgfqpoint{8.066346in}{1.696863in}}%
\pgfpathlineto{\pgfqpoint{8.069255in}{1.675204in}}%
\pgfpathlineto{\pgfqpoint{8.072592in}{1.656594in}}%
\pgfpathlineto{\pgfqpoint{8.073362in}{1.656287in}}%
\pgfpathlineto{\pgfqpoint{8.073704in}{1.656722in}}%
\pgfpathlineto{\pgfqpoint{8.074987in}{1.661292in}}%
\pgfpathlineto{\pgfqpoint{8.077897in}{1.682649in}}%
\pgfpathlineto{\pgfqpoint{8.081234in}{1.700681in}}%
\pgfpathlineto{\pgfqpoint{8.082004in}{1.700867in}}%
\pgfpathlineto{\pgfqpoint{8.082346in}{1.700386in}}%
\pgfpathlineto{\pgfqpoint{8.083715in}{1.695238in}}%
\pgfpathlineto{\pgfqpoint{8.087052in}{1.670575in}}%
\pgfpathlineto{\pgfqpoint{8.089961in}{1.657290in}}%
\pgfpathlineto{\pgfqpoint{8.090731in}{1.657521in}}%
\pgfpathlineto{\pgfqpoint{8.090902in}{1.657808in}}%
\pgfpathlineto{\pgfqpoint{8.092271in}{1.662938in}}%
\pgfpathlineto{\pgfqpoint{8.095779in}{1.688454in}}%
\pgfpathlineto{\pgfqpoint{8.098517in}{1.699966in}}%
\pgfpathlineto{\pgfqpoint{8.099373in}{1.699462in}}%
\pgfpathlineto{\pgfqpoint{8.099458in}{1.699296in}}%
\pgfpathlineto{\pgfqpoint{8.100913in}{1.693553in}}%
\pgfpathlineto{\pgfqpoint{8.107245in}{1.658126in}}%
\pgfpathlineto{\pgfqpoint{8.108015in}{1.659082in}}%
\pgfpathlineto{\pgfqpoint{8.109555in}{1.665614in}}%
\pgfpathlineto{\pgfqpoint{8.115630in}{1.699103in}}%
\pgfpathlineto{\pgfqpoint{8.116143in}{1.698691in}}%
\pgfpathlineto{\pgfqpoint{8.117426in}{1.694566in}}%
\pgfpathlineto{\pgfqpoint{8.120336in}{1.674918in}}%
\pgfpathlineto{\pgfqpoint{8.123587in}{1.659244in}}%
\pgfpathlineto{\pgfqpoint{8.124357in}{1.659255in}}%
\pgfpathlineto{\pgfqpoint{8.124614in}{1.659626in}}%
\pgfpathlineto{\pgfqpoint{8.125983in}{1.664466in}}%
\pgfpathlineto{\pgfqpoint{8.129747in}{1.689894in}}%
\pgfpathlineto{\pgfqpoint{8.132229in}{1.698195in}}%
\pgfpathlineto{\pgfqpoint{8.133170in}{1.697132in}}%
\pgfpathlineto{\pgfqpoint{8.134796in}{1.690140in}}%
\pgfpathlineto{\pgfqpoint{8.140614in}{1.659935in}}%
\pgfpathlineto{\pgfqpoint{8.140871in}{1.660058in}}%
\pgfpathlineto{\pgfqpoint{8.141983in}{1.662608in}}%
\pgfpathlineto{\pgfqpoint{8.144208in}{1.675287in}}%
\pgfpathlineto{\pgfqpoint{8.148400in}{1.697106in}}%
\pgfpathlineto{\pgfqpoint{8.149170in}{1.697143in}}%
\pgfpathlineto{\pgfqpoint{8.149512in}{1.696649in}}%
\pgfpathlineto{\pgfqpoint{8.150967in}{1.691365in}}%
\pgfpathlineto{\pgfqpoint{8.157042in}{1.660795in}}%
\pgfpathlineto{\pgfqpoint{8.157726in}{1.661496in}}%
\pgfpathlineto{\pgfqpoint{8.159181in}{1.666767in}}%
\pgfpathlineto{\pgfqpoint{8.165256in}{1.696464in}}%
\pgfpathlineto{\pgfqpoint{8.165855in}{1.695832in}}%
\pgfpathlineto{\pgfqpoint{8.167309in}{1.690717in}}%
\pgfpathlineto{\pgfqpoint{8.173299in}{1.661627in}}%
\pgfpathlineto{\pgfqpoint{8.173983in}{1.662289in}}%
\pgfpathlineto{\pgfqpoint{8.175438in}{1.667376in}}%
\pgfpathlineto{\pgfqpoint{8.181427in}{1.695649in}}%
\pgfpathlineto{\pgfqpoint{8.182026in}{1.695060in}}%
\pgfpathlineto{\pgfqpoint{8.183481in}{1.690144in}}%
\pgfpathlineto{\pgfqpoint{8.189470in}{1.662442in}}%
\pgfpathlineto{\pgfqpoint{8.190069in}{1.663035in}}%
\pgfpathlineto{\pgfqpoint{8.191524in}{1.667902in}}%
\pgfpathlineto{\pgfqpoint{8.197427in}{1.694863in}}%
\pgfpathlineto{\pgfqpoint{8.198026in}{1.694336in}}%
\pgfpathlineto{\pgfqpoint{8.199481in}{1.689661in}}%
\pgfpathlineto{\pgfqpoint{8.205385in}{1.663214in}}%
\pgfpathlineto{\pgfqpoint{8.206069in}{1.663875in}}%
\pgfpathlineto{\pgfqpoint{8.207609in}{1.669135in}}%
\pgfpathlineto{\pgfqpoint{8.213342in}{1.694090in}}%
\pgfpathlineto{\pgfqpoint{8.213770in}{1.693775in}}%
\pgfpathlineto{\pgfqpoint{8.215053in}{1.690353in}}%
\pgfpathlineto{\pgfqpoint{8.218219in}{1.672997in}}%
\pgfpathlineto{\pgfqpoint{8.220957in}{1.664000in}}%
\pgfpathlineto{\pgfqpoint{8.221984in}{1.664813in}}%
\pgfpathlineto{\pgfqpoint{8.223609in}{1.670644in}}%
\pgfpathlineto{\pgfqpoint{8.229085in}{1.693347in}}%
\pgfpathlineto{\pgfqpoint{8.229257in}{1.693270in}}%
\pgfpathlineto{\pgfqpoint{8.230369in}{1.691168in}}%
\pgfpathlineto{\pgfqpoint{8.232594in}{1.680666in}}%
\pgfpathlineto{\pgfqpoint{8.236444in}{1.664830in}}%
\pgfpathlineto{\pgfqpoint{8.237385in}{1.665095in}}%
\pgfpathlineto{\pgfqpoint{8.237556in}{1.665359in}}%
\pgfpathlineto{\pgfqpoint{8.239096in}{1.670356in}}%
\pgfpathlineto{\pgfqpoint{8.244658in}{1.692634in}}%
\pgfpathlineto{\pgfqpoint{8.245000in}{1.692431in}}%
\pgfpathlineto{\pgfqpoint{8.246284in}{1.689448in}}%
\pgfpathlineto{\pgfqpoint{8.249278in}{1.674418in}}%
\pgfpathlineto{\pgfqpoint{8.252102in}{1.665442in}}%
\pgfpathlineto{\pgfqpoint{8.253129in}{1.666140in}}%
\pgfpathlineto{\pgfqpoint{8.254754in}{1.671507in}}%
\pgfpathlineto{\pgfqpoint{8.260145in}{1.691932in}}%
\pgfpathlineto{\pgfqpoint{8.260230in}{1.691899in}}%
\pgfpathlineto{\pgfqpoint{8.261343in}{1.690065in}}%
\pgfpathlineto{\pgfqpoint{8.263567in}{1.680456in}}%
\pgfpathlineto{\pgfqpoint{8.267332in}{1.666208in}}%
\pgfpathlineto{\pgfqpoint{8.268359in}{1.666564in}}%
\pgfpathlineto{\pgfqpoint{8.268444in}{1.666694in}}%
\pgfpathlineto{\pgfqpoint{8.269984in}{1.671342in}}%
\pgfpathlineto{\pgfqpoint{8.275460in}{1.691262in}}%
\pgfpathlineto{\pgfqpoint{8.275717in}{1.691117in}}%
\pgfpathlineto{\pgfqpoint{8.277000in}{1.688423in}}%
\pgfpathlineto{\pgfqpoint{8.279995in}{1.674640in}}%
\pgfpathlineto{\pgfqpoint{8.282733in}{1.666770in}}%
\pgfpathlineto{\pgfqpoint{8.283760in}{1.667451in}}%
\pgfpathlineto{\pgfqpoint{8.285386in}{1.672475in}}%
\pgfpathlineto{\pgfqpoint{8.290605in}{1.690625in}}%
\pgfpathlineto{\pgfqpoint{8.290690in}{1.690598in}}%
\pgfpathlineto{\pgfqpoint{8.291803in}{1.688938in}}%
\pgfpathlineto{\pgfqpoint{8.294027in}{1.680108in}}%
\pgfpathlineto{\pgfqpoint{8.297707in}{1.667473in}}%
\pgfpathlineto{\pgfqpoint{8.298819in}{1.667980in}}%
\pgfpathlineto{\pgfqpoint{8.300445in}{1.672692in}}%
\pgfpathlineto{\pgfqpoint{8.305664in}{1.689999in}}%
\pgfpathlineto{\pgfqpoint{8.305749in}{1.689971in}}%
\pgfpathlineto{\pgfqpoint{8.306862in}{1.688341in}}%
\pgfpathlineto{\pgfqpoint{8.309086in}{1.679835in}}%
\pgfpathlineto{\pgfqpoint{8.312766in}{1.668038in}}%
\pgfpathlineto{\pgfqpoint{8.313878in}{1.668692in}}%
\pgfpathlineto{\pgfqpoint{8.315589in}{1.673747in}}%
\pgfpathlineto{\pgfqpoint{8.320552in}{1.689404in}}%
\pgfpathlineto{\pgfqpoint{8.321664in}{1.688064in}}%
\pgfpathlineto{\pgfqpoint{8.323803in}{1.680531in}}%
\pgfpathlineto{\pgfqpoint{8.327653in}{1.668612in}}%
\pgfpathlineto{\pgfqpoint{8.328766in}{1.669309in}}%
\pgfpathlineto{\pgfqpoint{8.330477in}{1.674235in}}%
\pgfpathlineto{\pgfqpoint{8.335354in}{1.688823in}}%
\pgfpathlineto{\pgfqpoint{8.336466in}{1.687517in}}%
\pgfpathlineto{\pgfqpoint{8.338605in}{1.680284in}}%
\pgfpathlineto{\pgfqpoint{8.342370in}{1.669186in}}%
\pgfpathlineto{\pgfqpoint{8.343482in}{1.669826in}}%
\pgfpathlineto{\pgfqpoint{8.345194in}{1.674514in}}%
\pgfpathlineto{\pgfqpoint{8.349985in}{1.688270in}}%
\pgfpathlineto{\pgfqpoint{8.351183in}{1.686918in}}%
\pgfpathlineto{\pgfqpoint{8.353408in}{1.679574in}}%
\pgfpathlineto{\pgfqpoint{8.357001in}{1.669726in}}%
\pgfpathlineto{\pgfqpoint{8.358199in}{1.670510in}}%
\pgfpathlineto{\pgfqpoint{8.360082in}{1.675832in}}%
\pgfpathlineto{\pgfqpoint{8.364360in}{1.687724in}}%
\pgfpathlineto{\pgfqpoint{8.365558in}{1.686770in}}%
\pgfpathlineto{\pgfqpoint{8.367526in}{1.681063in}}%
\pgfpathlineto{\pgfqpoint{8.371547in}{1.670238in}}%
\pgfpathlineto{\pgfqpoint{8.372745in}{1.671085in}}%
\pgfpathlineto{\pgfqpoint{8.374713in}{1.676572in}}%
\pgfpathlineto{\pgfqpoint{8.378734in}{1.687191in}}%
\pgfpathlineto{\pgfqpoint{8.379932in}{1.686394in}}%
\pgfpathlineto{\pgfqpoint{8.381900in}{1.681069in}}%
\pgfpathlineto{\pgfqpoint{8.385921in}{1.670748in}}%
\pgfpathlineto{\pgfqpoint{8.387119in}{1.671547in}}%
\pgfpathlineto{\pgfqpoint{8.389087in}{1.676766in}}%
\pgfpathlineto{\pgfqpoint{8.393109in}{1.686704in}}%
\pgfpathlineto{\pgfqpoint{8.394307in}{1.685854in}}%
\pgfpathlineto{\pgfqpoint{8.396360in}{1.680405in}}%
\pgfpathlineto{\pgfqpoint{8.400210in}{1.671232in}}%
\pgfpathlineto{\pgfqpoint{8.401494in}{1.672165in}}%
\pgfpathlineto{\pgfqpoint{8.403633in}{1.677874in}}%
\pgfpathlineto{\pgfqpoint{8.407312in}{1.686222in}}%
\pgfpathlineto{\pgfqpoint{8.408595in}{1.685296in}}%
\pgfpathlineto{\pgfqpoint{8.410735in}{1.679718in}}%
\pgfpathlineto{\pgfqpoint{8.414328in}{1.671713in}}%
\pgfpathlineto{\pgfqpoint{8.415612in}{1.672529in}}%
\pgfpathlineto{\pgfqpoint{8.417751in}{1.677883in}}%
\pgfpathlineto{\pgfqpoint{8.421430in}{1.685763in}}%
\pgfpathlineto{\pgfqpoint{8.422713in}{1.684860in}}%
\pgfpathlineto{\pgfqpoint{8.424938in}{1.679290in}}%
\pgfpathlineto{\pgfqpoint{8.428446in}{1.672151in}}%
\pgfpathlineto{\pgfqpoint{8.429729in}{1.673037in}}%
\pgfpathlineto{\pgfqpoint{8.431954in}{1.678462in}}%
\pgfpathlineto{\pgfqpoint{8.435377in}{1.685309in}}%
\pgfpathlineto{\pgfqpoint{8.436660in}{1.684537in}}%
\pgfpathlineto{\pgfqpoint{8.438799in}{1.679586in}}%
\pgfpathlineto{\pgfqpoint{8.442393in}{1.672586in}}%
\pgfpathlineto{\pgfqpoint{8.443762in}{1.673551in}}%
\pgfpathlineto{\pgfqpoint{8.446157in}{1.679270in}}%
\pgfpathlineto{\pgfqpoint{8.449323in}{1.684893in}}%
\pgfpathlineto{\pgfqpoint{8.450692in}{1.683961in}}%
\pgfpathlineto{\pgfqpoint{8.453088in}{1.678411in}}%
\pgfpathlineto{\pgfqpoint{8.456254in}{1.673001in}}%
\pgfpathlineto{\pgfqpoint{8.457623in}{1.673941in}}%
\pgfpathlineto{\pgfqpoint{8.460104in}{1.679576in}}%
\pgfpathlineto{\pgfqpoint{8.463184in}{1.684492in}}%
\pgfpathlineto{\pgfqpoint{8.464553in}{1.683508in}}%
\pgfpathlineto{\pgfqpoint{8.467206in}{1.677567in}}%
\pgfpathlineto{\pgfqpoint{8.470029in}{1.673398in}}%
\pgfpathlineto{\pgfqpoint{8.471484in}{1.674455in}}%
\pgfpathlineto{\pgfqpoint{8.474393in}{1.680872in}}%
\pgfpathlineto{\pgfqpoint{8.477045in}{1.684094in}}%
\pgfpathlineto{\pgfqpoint{8.478500in}{1.682799in}}%
\pgfpathlineto{\pgfqpoint{8.484318in}{1.673955in}}%
\pgfpathlineto{\pgfqpoint{8.484917in}{1.674502in}}%
\pgfpathlineto{\pgfqpoint{8.487313in}{1.679182in}}%
\pgfpathlineto{\pgfqpoint{8.490479in}{1.683725in}}%
\pgfpathlineto{\pgfqpoint{8.491933in}{1.682735in}}%
\pgfpathlineto{\pgfqpoint{8.495014in}{1.676571in}}%
\pgfpathlineto{\pgfqpoint{8.497495in}{1.674166in}}%
\pgfpathlineto{\pgfqpoint{8.499035in}{1.675597in}}%
\pgfpathlineto{\pgfqpoint{8.504682in}{1.683162in}}%
\pgfpathlineto{\pgfqpoint{8.504939in}{1.682970in}}%
\pgfpathlineto{\pgfqpoint{8.507078in}{1.679555in}}%
\pgfpathlineto{\pgfqpoint{8.510757in}{1.674497in}}%
\pgfpathlineto{\pgfqpoint{8.512297in}{1.675586in}}%
\pgfpathlineto{\pgfqpoint{8.518287in}{1.682700in}}%
\pgfpathlineto{\pgfqpoint{8.518629in}{1.682393in}}%
\pgfpathlineto{\pgfqpoint{8.521281in}{1.677941in}}%
\pgfpathlineto{\pgfqpoint{8.524190in}{1.674835in}}%
\pgfpathlineto{\pgfqpoint{8.525816in}{1.676053in}}%
\pgfpathlineto{\pgfqpoint{8.531634in}{1.682372in}}%
\pgfpathlineto{\pgfqpoint{8.531720in}{1.682307in}}%
\pgfpathlineto{\pgfqpoint{8.534030in}{1.678975in}}%
\pgfpathlineto{\pgfqpoint{8.537367in}{1.675151in}}%
\pgfpathlineto{\pgfqpoint{8.539078in}{1.676277in}}%
\pgfpathlineto{\pgfqpoint{8.544896in}{1.682058in}}%
\pgfpathlineto{\pgfqpoint{8.547207in}{1.679014in}}%
\pgfpathlineto{\pgfqpoint{8.550629in}{1.675457in}}%
\pgfpathlineto{\pgfqpoint{8.552426in}{1.676700in}}%
\pgfpathlineto{\pgfqpoint{8.557902in}{1.681867in}}%
\pgfpathlineto{\pgfqpoint{8.560127in}{1.679302in}}%
\pgfpathlineto{\pgfqpoint{8.563720in}{1.675746in}}%
\pgfpathlineto{\pgfqpoint{8.565603in}{1.676991in}}%
\pgfpathlineto{\pgfqpoint{8.570822in}{1.681665in}}%
\pgfpathlineto{\pgfqpoint{8.573047in}{1.679421in}}%
\pgfpathlineto{\pgfqpoint{8.576811in}{1.676026in}}%
\pgfpathlineto{\pgfqpoint{8.578779in}{1.677355in}}%
\pgfpathlineto{\pgfqpoint{8.583656in}{1.681457in}}%
\pgfpathlineto{\pgfqpoint{8.585881in}{1.679515in}}%
\pgfpathlineto{\pgfqpoint{8.589817in}{1.676295in}}%
\pgfpathlineto{\pgfqpoint{8.591870in}{1.677683in}}%
\pgfpathlineto{\pgfqpoint{8.596491in}{1.681224in}}%
\pgfpathlineto{\pgfqpoint{8.598715in}{1.679491in}}%
\pgfpathlineto{\pgfqpoint{8.602737in}{1.676550in}}%
\pgfpathlineto{\pgfqpoint{8.604961in}{1.678063in}}%
\pgfpathlineto{\pgfqpoint{8.609239in}{1.681003in}}%
\pgfpathlineto{\pgfqpoint{8.611550in}{1.679380in}}%
\pgfpathlineto{\pgfqpoint{8.615571in}{1.676793in}}%
\pgfpathlineto{\pgfqpoint{8.617967in}{1.678391in}}%
\pgfpathlineto{\pgfqpoint{8.621903in}{1.680792in}}%
\pgfpathlineto{\pgfqpoint{8.624384in}{1.679206in}}%
\pgfpathlineto{\pgfqpoint{8.628234in}{1.677005in}}%
\pgfpathlineto{\pgfqpoint{8.630716in}{1.678511in}}%
\pgfpathlineto{\pgfqpoint{8.634652in}{1.680564in}}%
\pgfpathlineto{\pgfqpoint{8.637304in}{1.678919in}}%
\pgfpathlineto{\pgfqpoint{8.640983in}{1.677232in}}%
\pgfpathlineto{\pgfqpoint{8.643807in}{1.678951in}}%
\pgfpathlineto{\pgfqpoint{8.647315in}{1.680351in}}%
\pgfpathlineto{\pgfqpoint{8.650309in}{1.678567in}}%
\pgfpathlineto{\pgfqpoint{8.653646in}{1.677447in}}%
\pgfpathlineto{\pgfqpoint{8.656898in}{1.679340in}}%
\pgfpathlineto{\pgfqpoint{8.660064in}{1.680113in}}%
\pgfpathlineto{\pgfqpoint{8.668534in}{1.678788in}}%
\pgfpathlineto{\pgfqpoint{8.672385in}{1.679974in}}%
\pgfpathlineto{\pgfqpoint{8.681369in}{1.679053in}}%
\pgfpathlineto{\pgfqpoint{8.684962in}{1.679772in}}%
\pgfpathlineto{\pgfqpoint{8.693176in}{1.678795in}}%
\pgfpathlineto{\pgfqpoint{8.697369in}{1.679609in}}%
\pgfpathlineto{\pgfqpoint{8.705326in}{1.678774in}}%
\pgfpathlineto{\pgfqpoint{8.709775in}{1.679446in}}%
\pgfpathlineto{\pgfqpoint{8.717305in}{1.678726in}}%
\pgfpathlineto{\pgfqpoint{8.722267in}{1.679265in}}%
\pgfpathlineto{\pgfqpoint{8.729198in}{1.678692in}}%
\pgfpathlineto{\pgfqpoint{8.734674in}{1.679115in}}%
\pgfpathlineto{\pgfqpoint{8.741262in}{1.678732in}}%
\pgfpathlineto{\pgfqpoint{8.747166in}{1.678963in}}%
\pgfpathlineto{\pgfqpoint{8.753412in}{1.678792in}}%
\pgfpathlineto{\pgfqpoint{8.759658in}{1.678843in}}%
\pgfpathlineto{\pgfqpoint{8.766075in}{1.678915in}}%
\pgfpathlineto{\pgfqpoint{8.772749in}{1.678705in}}%
\pgfpathlineto{\pgfqpoint{8.785070in}{1.678714in}}%
\pgfpathlineto{\pgfqpoint{8.808686in}{1.678787in}}%
\pgfpathlineto{\pgfqpoint{8.850354in}{1.678845in}}%
\pgfpathlineto{\pgfqpoint{8.850354in}{1.678845in}}%
\pgfusepath{stroke}%
\end{pgfscope}%
\begin{pgfscope}%
\pgfpathrectangle{\pgfqpoint{0.294194in}{0.558614in}}{\pgfqpoint{8.556160in}{2.240462in}}%
\pgfusepath{clip}%
\pgfsetrectcap%
\pgfsetroundjoin%
\pgfsetlinewidth{1.505625pt}%
\definecolor{currentstroke}{rgb}{1.000000,0.498039,0.054902}%
\pgfsetstrokecolor{currentstroke}%
\pgfsetdash{}{0pt}%
\pgfpathmoveto{\pgfqpoint{0.294194in}{1.678845in}}%
\pgfpathlineto{\pgfqpoint{0.297018in}{2.798920in}}%
\pgfpathlineto{\pgfqpoint{0.297103in}{2.798558in}}%
\pgfpathlineto{\pgfqpoint{0.297617in}{2.744732in}}%
\pgfpathlineto{\pgfqpoint{0.298900in}{2.266704in}}%
\pgfpathlineto{\pgfqpoint{0.302751in}{0.558661in}}%
\pgfpathlineto{\pgfqpoint{0.303093in}{0.576503in}}%
\pgfpathlineto{\pgfqpoint{0.304034in}{0.818569in}}%
\pgfpathlineto{\pgfqpoint{0.308483in}{2.799023in}}%
\pgfpathlineto{\pgfqpoint{0.309510in}{2.632208in}}%
\pgfpathlineto{\pgfqpoint{0.311478in}{1.613800in}}%
\pgfpathlineto{\pgfqpoint{0.314216in}{0.558793in}}%
\pgfpathlineto{\pgfqpoint{0.314302in}{0.559120in}}%
\pgfpathlineto{\pgfqpoint{0.314729in}{0.597181in}}%
\pgfpathlineto{\pgfqpoint{0.315842in}{0.956529in}}%
\pgfpathlineto{\pgfqpoint{0.320034in}{2.798951in}}%
\pgfpathlineto{\pgfqpoint{0.320633in}{2.735130in}}%
\pgfpathlineto{\pgfqpoint{0.322002in}{2.202450in}}%
\pgfpathlineto{\pgfqpoint{0.325767in}{0.558717in}}%
\pgfpathlineto{\pgfqpoint{0.326024in}{0.567900in}}%
\pgfpathlineto{\pgfqpoint{0.326794in}{0.721916in}}%
\pgfpathlineto{\pgfqpoint{0.328762in}{1.727788in}}%
\pgfpathlineto{\pgfqpoint{0.331585in}{2.799017in}}%
\pgfpathlineto{\pgfqpoint{0.331927in}{2.779130in}}%
\pgfpathlineto{\pgfqpoint{0.332869in}{2.536768in}}%
\pgfpathlineto{\pgfqpoint{0.337403in}{0.558733in}}%
\pgfpathlineto{\pgfqpoint{0.338430in}{0.732130in}}%
\pgfpathlineto{\pgfqpoint{0.340398in}{1.740866in}}%
\pgfpathlineto{\pgfqpoint{0.343222in}{2.798957in}}%
\pgfpathlineto{\pgfqpoint{0.343564in}{2.778510in}}%
\pgfpathlineto{\pgfqpoint{0.344505in}{2.536670in}}%
\pgfpathlineto{\pgfqpoint{0.349040in}{0.558734in}}%
\pgfpathlineto{\pgfqpoint{0.350067in}{0.722939in}}%
\pgfpathlineto{\pgfqpoint{0.352035in}{1.717182in}}%
\pgfpathlineto{\pgfqpoint{0.354858in}{2.798665in}}%
\pgfpathlineto{\pgfqpoint{0.354944in}{2.798619in}}%
\pgfpathlineto{\pgfqpoint{0.355372in}{2.763249in}}%
\pgfpathlineto{\pgfqpoint{0.356484in}{2.418875in}}%
\pgfpathlineto{\pgfqpoint{0.360762in}{0.558787in}}%
\pgfpathlineto{\pgfqpoint{0.361446in}{0.631647in}}%
\pgfpathlineto{\pgfqpoint{0.362815in}{1.164919in}}%
\pgfpathlineto{\pgfqpoint{0.366666in}{2.798849in}}%
\pgfpathlineto{\pgfqpoint{0.366922in}{2.787276in}}%
\pgfpathlineto{\pgfqpoint{0.367778in}{2.603174in}}%
\pgfpathlineto{\pgfqpoint{0.369917in}{1.490821in}}%
\pgfpathlineto{\pgfqpoint{0.372570in}{0.558893in}}%
\pgfpathlineto{\pgfqpoint{0.372912in}{0.579241in}}%
\pgfpathlineto{\pgfqpoint{0.373939in}{0.849412in}}%
\pgfpathlineto{\pgfqpoint{0.378473in}{2.798804in}}%
\pgfpathlineto{\pgfqpoint{0.379415in}{2.660811in}}%
\pgfpathlineto{\pgfqpoint{0.381211in}{1.809666in}}%
\pgfpathlineto{\pgfqpoint{0.384377in}{0.559006in}}%
\pgfpathlineto{\pgfqpoint{0.384634in}{0.567442in}}%
\pgfpathlineto{\pgfqpoint{0.385404in}{0.712897in}}%
\pgfpathlineto{\pgfqpoint{0.387286in}{1.627824in}}%
\pgfpathlineto{\pgfqpoint{0.390367in}{2.798678in}}%
\pgfpathlineto{\pgfqpoint{0.390709in}{2.778596in}}%
\pgfpathlineto{\pgfqpoint{0.391736in}{2.512445in}}%
\pgfpathlineto{\pgfqpoint{0.396356in}{0.559289in}}%
\pgfpathlineto{\pgfqpoint{0.397297in}{0.705553in}}%
\pgfpathlineto{\pgfqpoint{0.399179in}{1.606951in}}%
\pgfpathlineto{\pgfqpoint{0.402260in}{2.798450in}}%
\pgfpathlineto{\pgfqpoint{0.402431in}{2.795766in}}%
\pgfpathlineto{\pgfqpoint{0.403030in}{2.716369in}}%
\pgfpathlineto{\pgfqpoint{0.404484in}{2.138159in}}%
\pgfpathlineto{\pgfqpoint{0.408335in}{0.559385in}}%
\pgfpathlineto{\pgfqpoint{0.408420in}{0.561625in}}%
\pgfpathlineto{\pgfqpoint{0.409019in}{0.639117in}}%
\pgfpathlineto{\pgfqpoint{0.410474in}{1.211915in}}%
\pgfpathlineto{\pgfqpoint{0.414324in}{2.798484in}}%
\pgfpathlineto{\pgfqpoint{0.414495in}{2.793103in}}%
\pgfpathlineto{\pgfqpoint{0.415265in}{2.661415in}}%
\pgfpathlineto{\pgfqpoint{0.417062in}{1.828737in}}%
\pgfpathlineto{\pgfqpoint{0.420313in}{0.559303in}}%
\pgfpathlineto{\pgfqpoint{0.420485in}{0.562488in}}%
\pgfpathlineto{\pgfqpoint{0.421169in}{0.662288in}}%
\pgfpathlineto{\pgfqpoint{0.422795in}{1.355058in}}%
\pgfpathlineto{\pgfqpoint{0.426388in}{2.798415in}}%
\pgfpathlineto{\pgfqpoint{0.426474in}{2.797244in}}%
\pgfpathlineto{\pgfqpoint{0.426987in}{2.744363in}}%
\pgfpathlineto{\pgfqpoint{0.428271in}{2.304958in}}%
\pgfpathlineto{\pgfqpoint{0.432463in}{0.559358in}}%
\pgfpathlineto{\pgfqpoint{0.432891in}{0.588206in}}%
\pgfpathlineto{\pgfqpoint{0.434003in}{0.900998in}}%
\pgfpathlineto{\pgfqpoint{0.438538in}{2.798287in}}%
\pgfpathlineto{\pgfqpoint{0.439308in}{2.710162in}}%
\pgfpathlineto{\pgfqpoint{0.440848in}{2.091588in}}%
\pgfpathlineto{\pgfqpoint{0.444613in}{0.559531in}}%
\pgfpathlineto{\pgfqpoint{0.444784in}{0.562823in}}%
\pgfpathlineto{\pgfqpoint{0.445469in}{0.661079in}}%
\pgfpathlineto{\pgfqpoint{0.447094in}{1.341598in}}%
\pgfpathlineto{\pgfqpoint{0.450774in}{2.798078in}}%
\pgfpathlineto{\pgfqpoint{0.451116in}{2.778651in}}%
\pgfpathlineto{\pgfqpoint{0.452143in}{2.525332in}}%
\pgfpathlineto{\pgfqpoint{0.456849in}{0.559826in}}%
\pgfpathlineto{\pgfqpoint{0.457875in}{0.699524in}}%
\pgfpathlineto{\pgfqpoint{0.459672in}{1.514307in}}%
\pgfpathlineto{\pgfqpoint{0.463009in}{2.797861in}}%
\pgfpathlineto{\pgfqpoint{0.463180in}{2.795141in}}%
\pgfpathlineto{\pgfqpoint{0.463865in}{2.700557in}}%
\pgfpathlineto{\pgfqpoint{0.465490in}{2.033306in}}%
\pgfpathlineto{\pgfqpoint{0.469170in}{0.560021in}}%
\pgfpathlineto{\pgfqpoint{0.469255in}{0.560049in}}%
\pgfpathlineto{\pgfqpoint{0.469683in}{0.591778in}}%
\pgfpathlineto{\pgfqpoint{0.470795in}{0.903101in}}%
\pgfpathlineto{\pgfqpoint{0.475416in}{2.797805in}}%
\pgfpathlineto{\pgfqpoint{0.476186in}{2.710480in}}%
\pgfpathlineto{\pgfqpoint{0.477726in}{2.107446in}}%
\pgfpathlineto{\pgfqpoint{0.481576in}{0.560090in}}%
\pgfpathlineto{\pgfqpoint{0.481747in}{0.562685in}}%
\pgfpathlineto{\pgfqpoint{0.482432in}{0.655312in}}%
\pgfpathlineto{\pgfqpoint{0.483972in}{1.266651in}}%
\pgfpathlineto{\pgfqpoint{0.487822in}{2.797638in}}%
\pgfpathlineto{\pgfqpoint{0.487908in}{2.796909in}}%
\pgfpathlineto{\pgfqpoint{0.488421in}{2.749241in}}%
\pgfpathlineto{\pgfqpoint{0.489705in}{2.337361in}}%
\pgfpathlineto{\pgfqpoint{0.494068in}{0.560142in}}%
\pgfpathlineto{\pgfqpoint{0.494582in}{0.595436in}}%
\pgfpathlineto{\pgfqpoint{0.495780in}{0.943277in}}%
\pgfpathlineto{\pgfqpoint{0.500314in}{2.797349in}}%
\pgfpathlineto{\pgfqpoint{0.500999in}{2.737870in}}%
\pgfpathlineto{\pgfqpoint{0.502368in}{2.270931in}}%
\pgfpathlineto{\pgfqpoint{0.506646in}{0.560368in}}%
\pgfpathlineto{\pgfqpoint{0.506988in}{0.578468in}}%
\pgfpathlineto{\pgfqpoint{0.507929in}{0.788470in}}%
\pgfpathlineto{\pgfqpoint{0.510411in}{2.027038in}}%
\pgfpathlineto{\pgfqpoint{0.512892in}{2.797082in}}%
\pgfpathlineto{\pgfqpoint{0.512978in}{2.796943in}}%
\pgfpathlineto{\pgfqpoint{0.513405in}{2.765978in}}%
\pgfpathlineto{\pgfqpoint{0.514518in}{2.465673in}}%
\pgfpathlineto{\pgfqpoint{0.519224in}{0.560624in}}%
\pgfpathlineto{\pgfqpoint{0.520079in}{0.653744in}}%
\pgfpathlineto{\pgfqpoint{0.521619in}{1.247485in}}%
\pgfpathlineto{\pgfqpoint{0.525555in}{2.796877in}}%
\pgfpathlineto{\pgfqpoint{0.525726in}{2.794565in}}%
\pgfpathlineto{\pgfqpoint{0.526325in}{2.724197in}}%
\pgfpathlineto{\pgfqpoint{0.527780in}{2.203930in}}%
\pgfpathlineto{\pgfqpoint{0.531972in}{0.560803in}}%
\pgfpathlineto{\pgfqpoint{0.532144in}{0.565725in}}%
\pgfpathlineto{\pgfqpoint{0.532914in}{0.684036in}}%
\pgfpathlineto{\pgfqpoint{0.534710in}{1.444110in}}%
\pgfpathlineto{\pgfqpoint{0.538304in}{2.796716in}}%
\pgfpathlineto{\pgfqpoint{0.538475in}{2.794087in}}%
\pgfpathlineto{\pgfqpoint{0.539160in}{2.705750in}}%
\pgfpathlineto{\pgfqpoint{0.540700in}{2.122279in}}%
\pgfpathlineto{\pgfqpoint{0.544721in}{0.561002in}}%
\pgfpathlineto{\pgfqpoint{0.544807in}{0.561705in}}%
\pgfpathlineto{\pgfqpoint{0.545320in}{0.606863in}}%
\pgfpathlineto{\pgfqpoint{0.546604in}{0.998114in}}%
\pgfpathlineto{\pgfqpoint{0.551138in}{2.796551in}}%
\pgfpathlineto{\pgfqpoint{0.551652in}{2.763758in}}%
\pgfpathlineto{\pgfqpoint{0.552764in}{2.468709in}}%
\pgfpathlineto{\pgfqpoint{0.557556in}{0.561446in}}%
\pgfpathlineto{\pgfqpoint{0.558411in}{0.648150in}}%
\pgfpathlineto{\pgfqpoint{0.559951in}{1.218577in}}%
\pgfpathlineto{\pgfqpoint{0.564058in}{2.796329in}}%
\pgfpathlineto{\pgfqpoint{0.564144in}{2.795354in}}%
\pgfpathlineto{\pgfqpoint{0.564657in}{2.749374in}}%
\pgfpathlineto{\pgfqpoint{0.565941in}{2.361626in}}%
\pgfpathlineto{\pgfqpoint{0.570561in}{0.561549in}}%
\pgfpathlineto{\pgfqpoint{0.571074in}{0.598532in}}%
\pgfpathlineto{\pgfqpoint{0.572272in}{0.930899in}}%
\pgfpathlineto{\pgfqpoint{0.577064in}{2.796010in}}%
\pgfpathlineto{\pgfqpoint{0.577748in}{2.732372in}}%
\pgfpathlineto{\pgfqpoint{0.579117in}{2.283420in}}%
\pgfpathlineto{\pgfqpoint{0.583567in}{0.561774in}}%
\pgfpathlineto{\pgfqpoint{0.583909in}{0.576566in}}%
\pgfpathlineto{\pgfqpoint{0.584850in}{0.766537in}}%
\pgfpathlineto{\pgfqpoint{0.587160in}{1.849501in}}%
\pgfpathlineto{\pgfqpoint{0.589984in}{2.793516in}}%
\pgfpathlineto{\pgfqpoint{0.590155in}{2.795584in}}%
\pgfpathlineto{\pgfqpoint{0.590412in}{2.784630in}}%
\pgfpathlineto{\pgfqpoint{0.591267in}{2.630388in}}%
\pgfpathlineto{\pgfqpoint{0.593321in}{1.720119in}}%
\pgfpathlineto{\pgfqpoint{0.596572in}{0.563682in}}%
\pgfpathlineto{\pgfqpoint{0.596658in}{0.562159in}}%
\pgfpathlineto{\pgfqpoint{0.597000in}{0.574684in}}%
\pgfpathlineto{\pgfqpoint{0.597941in}{0.756942in}}%
\pgfpathlineto{\pgfqpoint{0.600166in}{1.779069in}}%
\pgfpathlineto{\pgfqpoint{0.603160in}{2.793815in}}%
\pgfpathlineto{\pgfqpoint{0.603246in}{2.795367in}}%
\pgfpathlineto{\pgfqpoint{0.603588in}{2.783085in}}%
\pgfpathlineto{\pgfqpoint{0.604529in}{2.602427in}}%
\pgfpathlineto{\pgfqpoint{0.606754in}{1.585756in}}%
\pgfpathlineto{\pgfqpoint{0.609749in}{0.564778in}}%
\pgfpathlineto{\pgfqpoint{0.609920in}{0.562507in}}%
\pgfpathlineto{\pgfqpoint{0.610177in}{0.572873in}}%
\pgfpathlineto{\pgfqpoint{0.611032in}{0.722934in}}%
\pgfpathlineto{\pgfqpoint{0.613000in}{1.572433in}}%
\pgfpathlineto{\pgfqpoint{0.616423in}{2.793934in}}%
\pgfpathlineto{\pgfqpoint{0.616508in}{2.795135in}}%
\pgfpathlineto{\pgfqpoint{0.616765in}{2.787790in}}%
\pgfpathlineto{\pgfqpoint{0.617620in}{2.647996in}}%
\pgfpathlineto{\pgfqpoint{0.619588in}{1.814961in}}%
\pgfpathlineto{\pgfqpoint{0.623011in}{0.566180in}}%
\pgfpathlineto{\pgfqpoint{0.623182in}{0.562697in}}%
\pgfpathlineto{\pgfqpoint{0.623524in}{0.577456in}}%
\pgfpathlineto{\pgfqpoint{0.624465in}{0.761464in}}%
\pgfpathlineto{\pgfqpoint{0.626776in}{1.814897in}}%
\pgfpathlineto{\pgfqpoint{0.629770in}{2.793873in}}%
\pgfpathlineto{\pgfqpoint{0.629856in}{2.794827in}}%
\pgfpathlineto{\pgfqpoint{0.630113in}{2.786892in}}%
\pgfpathlineto{\pgfqpoint{0.630968in}{2.646815in}}%
\pgfpathlineto{\pgfqpoint{0.632936in}{1.820143in}}%
\pgfpathlineto{\pgfqpoint{0.636444in}{0.564704in}}%
\pgfpathlineto{\pgfqpoint{0.636530in}{0.563154in}}%
\pgfpathlineto{\pgfqpoint{0.636872in}{0.574815in}}%
\pgfpathlineto{\pgfqpoint{0.637813in}{0.748973in}}%
\pgfpathlineto{\pgfqpoint{0.640038in}{1.739539in}}%
\pgfpathlineto{\pgfqpoint{0.643118in}{2.791019in}}%
\pgfpathlineto{\pgfqpoint{0.643289in}{2.794481in}}%
\pgfpathlineto{\pgfqpoint{0.643631in}{2.780132in}}%
\pgfpathlineto{\pgfqpoint{0.644573in}{2.600107in}}%
\pgfpathlineto{\pgfqpoint{0.646797in}{1.607062in}}%
\pgfpathlineto{\pgfqpoint{0.649877in}{0.566428in}}%
\pgfpathlineto{\pgfqpoint{0.650049in}{0.563406in}}%
\pgfpathlineto{\pgfqpoint{0.650391in}{0.578480in}}%
\pgfpathlineto{\pgfqpoint{0.651332in}{0.759365in}}%
\pgfpathlineto{\pgfqpoint{0.653557in}{1.749474in}}%
\pgfpathlineto{\pgfqpoint{0.656637in}{2.790563in}}%
\pgfpathlineto{\pgfqpoint{0.656808in}{2.794119in}}%
\pgfpathlineto{\pgfqpoint{0.657150in}{2.780260in}}%
\pgfpathlineto{\pgfqpoint{0.658091in}{2.603446in}}%
\pgfpathlineto{\pgfqpoint{0.660316in}{1.621608in}}%
\pgfpathlineto{\pgfqpoint{0.663396in}{0.568990in}}%
\pgfpathlineto{\pgfqpoint{0.663567in}{0.563938in}}%
\pgfpathlineto{\pgfqpoint{0.663995in}{0.581660in}}%
\pgfpathlineto{\pgfqpoint{0.665022in}{0.793122in}}%
\pgfpathlineto{\pgfqpoint{0.667589in}{1.970596in}}%
\pgfpathlineto{\pgfqpoint{0.670327in}{2.792549in}}%
\pgfpathlineto{\pgfqpoint{0.670412in}{2.793714in}}%
\pgfpathlineto{\pgfqpoint{0.670669in}{2.786865in}}%
\pgfpathlineto{\pgfqpoint{0.671439in}{2.675349in}}%
\pgfpathlineto{\pgfqpoint{0.673236in}{1.987532in}}%
\pgfpathlineto{\pgfqpoint{0.677172in}{0.565341in}}%
\pgfpathlineto{\pgfqpoint{0.677257in}{0.564174in}}%
\pgfpathlineto{\pgfqpoint{0.677514in}{0.570942in}}%
\pgfpathlineto{\pgfqpoint{0.678284in}{0.681568in}}%
\pgfpathlineto{\pgfqpoint{0.680081in}{1.365015in}}%
\pgfpathlineto{\pgfqpoint{0.684017in}{2.791494in}}%
\pgfpathlineto{\pgfqpoint{0.684102in}{2.793152in}}%
\pgfpathlineto{\pgfqpoint{0.684445in}{2.782805in}}%
\pgfpathlineto{\pgfqpoint{0.685300in}{2.641549in}}%
\pgfpathlineto{\pgfqpoint{0.687268in}{1.840478in}}%
\pgfpathlineto{\pgfqpoint{0.690862in}{0.568133in}}%
\pgfpathlineto{\pgfqpoint{0.691033in}{0.564553in}}%
\pgfpathlineto{\pgfqpoint{0.691375in}{0.577611in}}%
\pgfpathlineto{\pgfqpoint{0.692316in}{0.747463in}}%
\pgfpathlineto{\pgfqpoint{0.694541in}{1.701963in}}%
\pgfpathlineto{\pgfqpoint{0.697792in}{2.789778in}}%
\pgfpathlineto{\pgfqpoint{0.697964in}{2.792929in}}%
\pgfpathlineto{\pgfqpoint{0.698306in}{2.779169in}}%
\pgfpathlineto{\pgfqpoint{0.699247in}{2.608506in}}%
\pgfpathlineto{\pgfqpoint{0.701472in}{1.657043in}}%
\pgfpathlineto{\pgfqpoint{0.704723in}{0.568700in}}%
\pgfpathlineto{\pgfqpoint{0.704894in}{0.564983in}}%
\pgfpathlineto{\pgfqpoint{0.705322in}{0.584708in}}%
\pgfpathlineto{\pgfqpoint{0.706349in}{0.793351in}}%
\pgfpathlineto{\pgfqpoint{0.708916in}{1.941617in}}%
\pgfpathlineto{\pgfqpoint{0.711739in}{2.790399in}}%
\pgfpathlineto{\pgfqpoint{0.711910in}{2.792368in}}%
\pgfpathlineto{\pgfqpoint{0.712167in}{2.782968in}}%
\pgfpathlineto{\pgfqpoint{0.713023in}{2.647693in}}%
\pgfpathlineto{\pgfqpoint{0.714991in}{1.868370in}}%
\pgfpathlineto{\pgfqpoint{0.718670in}{0.570041in}}%
\pgfpathlineto{\pgfqpoint{0.718841in}{0.565528in}}%
\pgfpathlineto{\pgfqpoint{0.719269in}{0.582833in}}%
\pgfpathlineto{\pgfqpoint{0.720295in}{0.783783in}}%
\pgfpathlineto{\pgfqpoint{0.722777in}{1.874333in}}%
\pgfpathlineto{\pgfqpoint{0.725686in}{2.787106in}}%
\pgfpathlineto{\pgfqpoint{0.725857in}{2.791876in}}%
\pgfpathlineto{\pgfqpoint{0.726285in}{2.775427in}}%
\pgfpathlineto{\pgfqpoint{0.727312in}{2.577560in}}%
\pgfpathlineto{\pgfqpoint{0.729707in}{1.536318in}}%
\pgfpathlineto{\pgfqpoint{0.732788in}{0.568682in}}%
\pgfpathlineto{\pgfqpoint{0.732959in}{0.565881in}}%
\pgfpathlineto{\pgfqpoint{0.733301in}{0.579585in}}%
\pgfpathlineto{\pgfqpoint{0.734242in}{0.745185in}}%
\pgfpathlineto{\pgfqpoint{0.736381in}{1.628609in}}%
\pgfpathlineto{\pgfqpoint{0.739804in}{2.786254in}}%
\pgfpathlineto{\pgfqpoint{0.740060in}{2.791450in}}%
\pgfpathlineto{\pgfqpoint{0.740403in}{2.776034in}}%
\pgfpathlineto{\pgfqpoint{0.741429in}{2.583141in}}%
\pgfpathlineto{\pgfqpoint{0.743825in}{1.555412in}}%
\pgfpathlineto{\pgfqpoint{0.746905in}{0.571766in}}%
\pgfpathlineto{\pgfqpoint{0.747162in}{0.566459in}}%
\pgfpathlineto{\pgfqpoint{0.747504in}{0.581552in}}%
\pgfpathlineto{\pgfqpoint{0.748531in}{0.772422in}}%
\pgfpathlineto{\pgfqpoint{0.750927in}{1.793773in}}%
\pgfpathlineto{\pgfqpoint{0.754007in}{2.784174in}}%
\pgfpathlineto{\pgfqpoint{0.754264in}{2.791103in}}%
\pgfpathlineto{\pgfqpoint{0.754691in}{2.771247in}}%
\pgfpathlineto{\pgfqpoint{0.755718in}{2.570713in}}%
\pgfpathlineto{\pgfqpoint{0.758200in}{1.502931in}}%
\pgfpathlineto{\pgfqpoint{0.761194in}{0.572300in}}%
\pgfpathlineto{\pgfqpoint{0.761451in}{0.566934in}}%
\pgfpathlineto{\pgfqpoint{0.761793in}{0.581597in}}%
\pgfpathlineto{\pgfqpoint{0.762820in}{0.769008in}}%
\pgfpathlineto{\pgfqpoint{0.765130in}{1.736609in}}%
\pgfpathlineto{\pgfqpoint{0.768381in}{2.785155in}}%
\pgfpathlineto{\pgfqpoint{0.768638in}{2.790506in}}%
\pgfpathlineto{\pgfqpoint{0.769066in}{2.768536in}}%
\pgfpathlineto{\pgfqpoint{0.770178in}{2.540169in}}%
\pgfpathlineto{\pgfqpoint{0.772916in}{1.343910in}}%
\pgfpathlineto{\pgfqpoint{0.775654in}{0.570393in}}%
\pgfpathlineto{\pgfqpoint{0.775825in}{0.567350in}}%
\pgfpathlineto{\pgfqpoint{0.776168in}{0.579670in}}%
\pgfpathlineto{\pgfqpoint{0.777109in}{0.735816in}}%
\pgfpathlineto{\pgfqpoint{0.779248in}{1.584622in}}%
\pgfpathlineto{\pgfqpoint{0.782842in}{2.784670in}}%
\pgfpathlineto{\pgfqpoint{0.783098in}{2.789991in}}%
\pgfpathlineto{\pgfqpoint{0.783526in}{2.768470in}}%
\pgfpathlineto{\pgfqpoint{0.784638in}{2.543883in}}%
\pgfpathlineto{\pgfqpoint{0.787376in}{1.360362in}}%
\pgfpathlineto{\pgfqpoint{0.790200in}{0.570012in}}%
\pgfpathlineto{\pgfqpoint{0.790371in}{0.567961in}}%
\pgfpathlineto{\pgfqpoint{0.790713in}{0.581956in}}%
\pgfpathlineto{\pgfqpoint{0.791654in}{0.740430in}}%
\pgfpathlineto{\pgfqpoint{0.793794in}{1.584233in}}%
\pgfpathlineto{\pgfqpoint{0.797387in}{2.782630in}}%
\pgfpathlineto{\pgfqpoint{0.797644in}{2.789532in}}%
\pgfpathlineto{\pgfqpoint{0.798072in}{2.771133in}}%
\pgfpathlineto{\pgfqpoint{0.799098in}{2.580998in}}%
\pgfpathlineto{\pgfqpoint{0.801494in}{1.591341in}}%
\pgfpathlineto{\pgfqpoint{0.804746in}{0.574020in}}%
\pgfpathlineto{\pgfqpoint{0.805002in}{0.568485in}}%
\pgfpathlineto{\pgfqpoint{0.805430in}{0.588904in}}%
\pgfpathlineto{\pgfqpoint{0.806542in}{0.806811in}}%
\pgfpathlineto{\pgfqpoint{0.809195in}{1.930752in}}%
\pgfpathlineto{\pgfqpoint{0.812104in}{2.783199in}}%
\pgfpathlineto{\pgfqpoint{0.812361in}{2.788940in}}%
\pgfpathlineto{\pgfqpoint{0.812788in}{2.769115in}}%
\pgfpathlineto{\pgfqpoint{0.813901in}{2.553969in}}%
\pgfpathlineto{\pgfqpoint{0.816553in}{1.437271in}}%
\pgfpathlineto{\pgfqpoint{0.819548in}{0.572588in}}%
\pgfpathlineto{\pgfqpoint{0.819719in}{0.569043in}}%
\pgfpathlineto{\pgfqpoint{0.820147in}{0.585697in}}%
\pgfpathlineto{\pgfqpoint{0.821174in}{0.768378in}}%
\pgfpathlineto{\pgfqpoint{0.823484in}{1.696615in}}%
\pgfpathlineto{\pgfqpoint{0.826906in}{2.781927in}}%
\pgfpathlineto{\pgfqpoint{0.827163in}{2.788385in}}%
\pgfpathlineto{\pgfqpoint{0.827591in}{2.770247in}}%
\pgfpathlineto{\pgfqpoint{0.828617in}{2.585491in}}%
\pgfpathlineto{\pgfqpoint{0.831013in}{1.619423in}}%
\pgfpathlineto{\pgfqpoint{0.834350in}{0.576615in}}%
\pgfpathlineto{\pgfqpoint{0.834607in}{0.569617in}}%
\pgfpathlineto{\pgfqpoint{0.835035in}{0.586608in}}%
\pgfpathlineto{\pgfqpoint{0.836061in}{0.767630in}}%
\pgfpathlineto{\pgfqpoint{0.838372in}{1.685360in}}%
\pgfpathlineto{\pgfqpoint{0.841794in}{2.778392in}}%
\pgfpathlineto{\pgfqpoint{0.842136in}{2.787684in}}%
\pgfpathlineto{\pgfqpoint{0.842564in}{2.767361in}}%
\pgfpathlineto{\pgfqpoint{0.843676in}{2.556500in}}%
\pgfpathlineto{\pgfqpoint{0.846243in}{1.502467in}}%
\pgfpathlineto{\pgfqpoint{0.849324in}{0.578931in}}%
\pgfpathlineto{\pgfqpoint{0.849666in}{0.570387in}}%
\pgfpathlineto{\pgfqpoint{0.850094in}{0.591362in}}%
\pgfpathlineto{\pgfqpoint{0.851206in}{0.802394in}}%
\pgfpathlineto{\pgfqpoint{0.853858in}{1.890534in}}%
\pgfpathlineto{\pgfqpoint{0.856939in}{2.781769in}}%
\pgfpathlineto{\pgfqpoint{0.857195in}{2.787117in}}%
\pgfpathlineto{\pgfqpoint{0.857623in}{2.768132in}}%
\pgfpathlineto{\pgfqpoint{0.858735in}{2.563167in}}%
\pgfpathlineto{\pgfqpoint{0.861302in}{1.524300in}}%
\pgfpathlineto{\pgfqpoint{0.864468in}{0.578949in}}%
\pgfpathlineto{\pgfqpoint{0.864725in}{0.570982in}}%
\pgfpathlineto{\pgfqpoint{0.865238in}{0.592367in}}%
\pgfpathlineto{\pgfqpoint{0.866351in}{0.801375in}}%
\pgfpathlineto{\pgfqpoint{0.868917in}{1.839012in}}%
\pgfpathlineto{\pgfqpoint{0.872083in}{2.778342in}}%
\pgfpathlineto{\pgfqpoint{0.872340in}{2.786363in}}%
\pgfpathlineto{\pgfqpoint{0.872853in}{2.765421in}}%
\pgfpathlineto{\pgfqpoint{0.873966in}{2.558723in}}%
\pgfpathlineto{\pgfqpoint{0.876532in}{1.528147in}}%
\pgfpathlineto{\pgfqpoint{0.879698in}{0.581712in}}%
\pgfpathlineto{\pgfqpoint{0.880040in}{0.571533in}}%
\pgfpathlineto{\pgfqpoint{0.880468in}{0.589355in}}%
\pgfpathlineto{\pgfqpoint{0.881495in}{0.765004in}}%
\pgfpathlineto{\pgfqpoint{0.883805in}{1.650348in}}%
\pgfpathlineto{\pgfqpoint{0.887399in}{2.776642in}}%
\pgfpathlineto{\pgfqpoint{0.887741in}{2.785765in}}%
\pgfpathlineto{\pgfqpoint{0.888169in}{2.766908in}}%
\pgfpathlineto{\pgfqpoint{0.889281in}{2.567709in}}%
\pgfpathlineto{\pgfqpoint{0.891763in}{1.592528in}}%
\pgfpathlineto{\pgfqpoint{0.895099in}{0.582444in}}%
\pgfpathlineto{\pgfqpoint{0.895442in}{0.572212in}}%
\pgfpathlineto{\pgfqpoint{0.895870in}{0.589409in}}%
\pgfpathlineto{\pgfqpoint{0.896896in}{0.761274in}}%
\pgfpathlineto{\pgfqpoint{0.899206in}{1.633403in}}%
\pgfpathlineto{\pgfqpoint{0.902800in}{2.771463in}}%
\pgfpathlineto{\pgfqpoint{0.903228in}{2.784997in}}%
\pgfpathlineto{\pgfqpoint{0.903656in}{2.765540in}}%
\pgfpathlineto{\pgfqpoint{0.904768in}{2.567646in}}%
\pgfpathlineto{\pgfqpoint{0.907249in}{1.603452in}}%
\pgfpathlineto{\pgfqpoint{0.910586in}{0.586727in}}%
\pgfpathlineto{\pgfqpoint{0.911014in}{0.573024in}}%
\pgfpathlineto{\pgfqpoint{0.911442in}{0.592005in}}%
\pgfpathlineto{\pgfqpoint{0.912554in}{0.787402in}}%
\pgfpathlineto{\pgfqpoint{0.915036in}{1.744245in}}%
\pgfpathlineto{\pgfqpoint{0.918458in}{2.773732in}}%
\pgfpathlineto{\pgfqpoint{0.918800in}{2.784408in}}%
\pgfpathlineto{\pgfqpoint{0.919314in}{2.761581in}}%
\pgfpathlineto{\pgfqpoint{0.920426in}{2.559788in}}%
\pgfpathlineto{\pgfqpoint{0.922993in}{1.563889in}}%
\pgfpathlineto{\pgfqpoint{0.926330in}{0.583652in}}%
\pgfpathlineto{\pgfqpoint{0.926672in}{0.573652in}}%
\pgfpathlineto{\pgfqpoint{0.927100in}{0.590022in}}%
\pgfpathlineto{\pgfqpoint{0.928127in}{0.755250in}}%
\pgfpathlineto{\pgfqpoint{0.930437in}{1.601628in}}%
\pgfpathlineto{\pgfqpoint{0.934202in}{2.772100in}}%
\pgfpathlineto{\pgfqpoint{0.934544in}{2.783620in}}%
\pgfpathlineto{\pgfqpoint{0.935057in}{2.762795in}}%
\pgfpathlineto{\pgfqpoint{0.936169in}{2.567819in}}%
\pgfpathlineto{\pgfqpoint{0.938651in}{1.626129in}}%
\pgfpathlineto{\pgfqpoint{0.942073in}{0.590257in}}%
\pgfpathlineto{\pgfqpoint{0.942501in}{0.574409in}}%
\pgfpathlineto{\pgfqpoint{0.943014in}{0.596880in}}%
\pgfpathlineto{\pgfqpoint{0.944127in}{0.793670in}}%
\pgfpathlineto{\pgfqpoint{0.946694in}{1.769394in}}%
\pgfpathlineto{\pgfqpoint{0.950116in}{2.771819in}}%
\pgfpathlineto{\pgfqpoint{0.950458in}{2.782864in}}%
\pgfpathlineto{\pgfqpoint{0.950972in}{2.762069in}}%
\pgfpathlineto{\pgfqpoint{0.952084in}{2.570023in}}%
\pgfpathlineto{\pgfqpoint{0.954565in}{1.641204in}}%
\pgfpathlineto{\pgfqpoint{0.958073in}{0.589752in}}%
\pgfpathlineto{\pgfqpoint{0.958501in}{0.575206in}}%
\pgfpathlineto{\pgfqpoint{0.959015in}{0.598421in}}%
\pgfpathlineto{\pgfqpoint{0.960213in}{0.816349in}}%
\pgfpathlineto{\pgfqpoint{0.962951in}{1.867506in}}%
\pgfpathlineto{\pgfqpoint{0.966202in}{2.772534in}}%
\pgfpathlineto{\pgfqpoint{0.966544in}{2.782094in}}%
\pgfpathlineto{\pgfqpoint{0.966972in}{2.766563in}}%
\pgfpathlineto{\pgfqpoint{0.967999in}{2.609203in}}%
\pgfpathlineto{\pgfqpoint{0.970223in}{1.831545in}}%
\pgfpathlineto{\pgfqpoint{0.974245in}{0.588319in}}%
\pgfpathlineto{\pgfqpoint{0.974673in}{0.576148in}}%
\pgfpathlineto{\pgfqpoint{0.975100in}{0.594187in}}%
\pgfpathlineto{\pgfqpoint{0.976213in}{0.776429in}}%
\pgfpathlineto{\pgfqpoint{0.978608in}{1.643502in}}%
\pgfpathlineto{\pgfqpoint{0.982288in}{2.762936in}}%
\pgfpathlineto{\pgfqpoint{0.982801in}{2.781142in}}%
\pgfpathlineto{\pgfqpoint{0.983229in}{2.763431in}}%
\pgfpathlineto{\pgfqpoint{0.984341in}{2.583334in}}%
\pgfpathlineto{\pgfqpoint{0.986737in}{1.723425in}}%
\pgfpathlineto{\pgfqpoint{0.990502in}{0.591851in}}%
\pgfpathlineto{\pgfqpoint{0.990929in}{0.576844in}}%
\pgfpathlineto{\pgfqpoint{0.991443in}{0.597885in}}%
\pgfpathlineto{\pgfqpoint{0.992555in}{0.783253in}}%
\pgfpathlineto{\pgfqpoint{0.995036in}{1.680054in}}%
\pgfpathlineto{\pgfqpoint{0.998716in}{2.765698in}}%
\pgfpathlineto{\pgfqpoint{0.999143in}{2.780417in}}%
\pgfpathlineto{\pgfqpoint{0.999657in}{2.759437in}}%
\pgfpathlineto{\pgfqpoint{1.000769in}{2.575623in}}%
\pgfpathlineto{\pgfqpoint{1.003250in}{1.685543in}}%
\pgfpathlineto{\pgfqpoint{1.006930in}{0.595094in}}%
\pgfpathlineto{\pgfqpoint{1.007357in}{0.577842in}}%
\pgfpathlineto{\pgfqpoint{1.007871in}{0.595383in}}%
\pgfpathlineto{\pgfqpoint{1.008983in}{0.770923in}}%
\pgfpathlineto{\pgfqpoint{1.011379in}{1.611198in}}%
\pgfpathlineto{\pgfqpoint{1.015229in}{2.762595in}}%
\pgfpathlineto{\pgfqpoint{1.015657in}{2.779428in}}%
\pgfpathlineto{\pgfqpoint{1.016170in}{2.761788in}}%
\pgfpathlineto{\pgfqpoint{1.017283in}{2.587487in}}%
\pgfpathlineto{\pgfqpoint{1.019678in}{1.753497in}}%
\pgfpathlineto{\pgfqpoint{1.023529in}{0.598314in}}%
\pgfpathlineto{\pgfqpoint{1.024042in}{0.578612in}}%
\pgfpathlineto{\pgfqpoint{1.024555in}{0.599725in}}%
\pgfpathlineto{\pgfqpoint{1.025753in}{0.800329in}}%
\pgfpathlineto{\pgfqpoint{1.028406in}{1.755696in}}%
\pgfpathlineto{\pgfqpoint{1.031914in}{2.759335in}}%
\pgfpathlineto{\pgfqpoint{1.032427in}{2.778615in}}%
\pgfpathlineto{\pgfqpoint{1.032941in}{2.757520in}}%
\pgfpathlineto{\pgfqpoint{1.034138in}{2.558564in}}%
\pgfpathlineto{\pgfqpoint{1.036705in}{1.645371in}}%
\pgfpathlineto{\pgfqpoint{1.040299in}{0.601888in}}%
\pgfpathlineto{\pgfqpoint{1.040812in}{0.579599in}}%
\pgfpathlineto{\pgfqpoint{1.041326in}{0.597246in}}%
\pgfpathlineto{\pgfqpoint{1.042438in}{0.767264in}}%
\pgfpathlineto{\pgfqpoint{1.044834in}{1.581604in}}%
\pgfpathlineto{\pgfqpoint{1.048770in}{2.755509in}}%
\pgfpathlineto{\pgfqpoint{1.049283in}{2.777614in}}%
\pgfpathlineto{\pgfqpoint{1.049796in}{2.760223in}}%
\pgfpathlineto{\pgfqpoint{1.050909in}{2.592146in}}%
\pgfpathlineto{\pgfqpoint{1.053304in}{1.785002in}}%
\pgfpathlineto{\pgfqpoint{1.057326in}{0.599376in}}%
\pgfpathlineto{\pgfqpoint{1.057839in}{0.580493in}}%
\pgfpathlineto{\pgfqpoint{1.058353in}{0.600667in}}%
\pgfpathlineto{\pgfqpoint{1.059551in}{0.792841in}}%
\pgfpathlineto{\pgfqpoint{1.062117in}{1.683091in}}%
\pgfpathlineto{\pgfqpoint{1.065882in}{2.757517in}}%
\pgfpathlineto{\pgfqpoint{1.066396in}{2.776713in}}%
\pgfpathlineto{\pgfqpoint{1.066909in}{2.757289in}}%
\pgfpathlineto{\pgfqpoint{1.068021in}{2.588009in}}%
\pgfpathlineto{\pgfqpoint{1.070417in}{1.789419in}}%
\pgfpathlineto{\pgfqpoint{1.074438in}{0.604652in}}%
\pgfpathlineto{\pgfqpoint{1.075037in}{0.581556in}}%
\pgfpathlineto{\pgfqpoint{1.075551in}{0.603100in}}%
\pgfpathlineto{\pgfqpoint{1.076749in}{0.794981in}}%
\pgfpathlineto{\pgfqpoint{1.079315in}{1.674071in}}%
\pgfpathlineto{\pgfqpoint{1.083080in}{2.751483in}}%
\pgfpathlineto{\pgfqpoint{1.083679in}{2.775687in}}%
\pgfpathlineto{\pgfqpoint{1.084193in}{2.755566in}}%
\pgfpathlineto{\pgfqpoint{1.085390in}{2.568324in}}%
\pgfpathlineto{\pgfqpoint{1.087957in}{1.699500in}}%
\pgfpathlineto{\pgfqpoint{1.091808in}{0.604746in}}%
\pgfpathlineto{\pgfqpoint{1.092407in}{0.582600in}}%
\pgfpathlineto{\pgfqpoint{1.092920in}{0.604011in}}%
\pgfpathlineto{\pgfqpoint{1.094118in}{0.792400in}}%
\pgfpathlineto{\pgfqpoint{1.096685in}{1.656951in}}%
\pgfpathlineto{\pgfqpoint{1.100535in}{2.750425in}}%
\pgfpathlineto{\pgfqpoint{1.101134in}{2.774664in}}%
\pgfpathlineto{\pgfqpoint{1.101647in}{2.755518in}}%
\pgfpathlineto{\pgfqpoint{1.102845in}{2.573569in}}%
\pgfpathlineto{\pgfqpoint{1.105326in}{1.754743in}}%
\pgfpathlineto{\pgfqpoint{1.109348in}{0.606916in}}%
\pgfpathlineto{\pgfqpoint{1.109947in}{0.583572in}}%
\pgfpathlineto{\pgfqpoint{1.110460in}{0.603013in}}%
\pgfpathlineto{\pgfqpoint{1.111658in}{0.783991in}}%
\pgfpathlineto{\pgfqpoint{1.114139in}{1.596468in}}%
\pgfpathlineto{\pgfqpoint{1.118161in}{2.746949in}}%
\pgfpathlineto{\pgfqpoint{1.118760in}{2.773553in}}%
\pgfpathlineto{\pgfqpoint{1.119273in}{2.757381in}}%
\pgfpathlineto{\pgfqpoint{1.120385in}{2.603179in}}%
\pgfpathlineto{\pgfqpoint{1.122696in}{1.885355in}}%
\pgfpathlineto{\pgfqpoint{1.127145in}{0.604971in}}%
\pgfpathlineto{\pgfqpoint{1.127658in}{0.584726in}}%
\pgfpathlineto{\pgfqpoint{1.128172in}{0.600042in}}%
\pgfpathlineto{\pgfqpoint{1.129284in}{0.751095in}}%
\pgfpathlineto{\pgfqpoint{1.131594in}{1.460191in}}%
\pgfpathlineto{\pgfqpoint{1.136129in}{2.754086in}}%
\pgfpathlineto{\pgfqpoint{1.136642in}{2.772482in}}%
\pgfpathlineto{\pgfqpoint{1.137156in}{2.755751in}}%
\pgfpathlineto{\pgfqpoint{1.138268in}{2.603267in}}%
\pgfpathlineto{\pgfqpoint{1.140578in}{1.897275in}}%
\pgfpathlineto{\pgfqpoint{1.145113in}{0.606067in}}%
\pgfpathlineto{\pgfqpoint{1.145626in}{0.585882in}}%
\pgfpathlineto{\pgfqpoint{1.146140in}{0.600389in}}%
\pgfpathlineto{\pgfqpoint{1.147252in}{0.746962in}}%
\pgfpathlineto{\pgfqpoint{1.149562in}{1.440711in}}%
\pgfpathlineto{\pgfqpoint{1.154183in}{2.751395in}}%
\pgfpathlineto{\pgfqpoint{1.154696in}{2.771240in}}%
\pgfpathlineto{\pgfqpoint{1.155209in}{2.756827in}}%
\pgfpathlineto{\pgfqpoint{1.156322in}{2.611866in}}%
\pgfpathlineto{\pgfqpoint{1.158632in}{1.925049in}}%
\pgfpathlineto{\pgfqpoint{1.163252in}{0.610637in}}%
\pgfpathlineto{\pgfqpoint{1.163851in}{0.586934in}}%
\pgfpathlineto{\pgfqpoint{1.164365in}{0.603244in}}%
\pgfpathlineto{\pgfqpoint{1.165477in}{0.750618in}}%
\pgfpathlineto{\pgfqpoint{1.167787in}{1.435579in}}%
\pgfpathlineto{\pgfqpoint{1.172493in}{2.751208in}}%
\pgfpathlineto{\pgfqpoint{1.173006in}{2.770095in}}%
\pgfpathlineto{\pgfqpoint{1.173520in}{2.755588in}}%
\pgfpathlineto{\pgfqpoint{1.174632in}{2.613279in}}%
\pgfpathlineto{\pgfqpoint{1.176857in}{1.970404in}}%
\pgfpathlineto{\pgfqpoint{1.181734in}{0.606586in}}%
\pgfpathlineto{\pgfqpoint{1.182247in}{0.588185in}}%
\pgfpathlineto{\pgfqpoint{1.182760in}{0.602747in}}%
\pgfpathlineto{\pgfqpoint{1.183873in}{0.743747in}}%
\pgfpathlineto{\pgfqpoint{1.186097in}{1.380277in}}%
\pgfpathlineto{\pgfqpoint{1.191060in}{2.752753in}}%
\pgfpathlineto{\pgfqpoint{1.191573in}{2.768945in}}%
\pgfpathlineto{\pgfqpoint{1.192087in}{2.752608in}}%
\pgfpathlineto{\pgfqpoint{1.193199in}{2.609426in}}%
\pgfpathlineto{\pgfqpoint{1.195509in}{1.944556in}}%
\pgfpathlineto{\pgfqpoint{1.200301in}{0.613204in}}%
\pgfpathlineto{\pgfqpoint{1.200900in}{0.589415in}}%
\pgfpathlineto{\pgfqpoint{1.201413in}{0.603782in}}%
\pgfpathlineto{\pgfqpoint{1.202525in}{0.741556in}}%
\pgfpathlineto{\pgfqpoint{1.204750in}{1.364561in}}%
\pgfpathlineto{\pgfqpoint{1.209798in}{2.749956in}}%
\pgfpathlineto{\pgfqpoint{1.210312in}{2.767631in}}%
\pgfpathlineto{\pgfqpoint{1.210825in}{2.753635in}}%
\pgfpathlineto{\pgfqpoint{1.211937in}{2.618028in}}%
\pgfpathlineto{\pgfqpoint{1.214162in}{2.002537in}}%
\pgfpathlineto{\pgfqpoint{1.219296in}{0.606804in}}%
\pgfpathlineto{\pgfqpoint{1.219809in}{0.590657in}}%
\pgfpathlineto{\pgfqpoint{1.220322in}{0.605751in}}%
\pgfpathlineto{\pgfqpoint{1.221435in}{0.742174in}}%
\pgfpathlineto{\pgfqpoint{1.223659in}{1.353998in}}%
\pgfpathlineto{\pgfqpoint{1.228793in}{2.747832in}}%
\pgfpathlineto{\pgfqpoint{1.229392in}{2.766306in}}%
\pgfpathlineto{\pgfqpoint{1.229905in}{2.748776in}}%
\pgfpathlineto{\pgfqpoint{1.231103in}{2.592648in}}%
\pgfpathlineto{\pgfqpoint{1.233499in}{1.908796in}}%
\pgfpathlineto{\pgfqpoint{1.238291in}{0.617706in}}%
\pgfpathlineto{\pgfqpoint{1.238975in}{0.591993in}}%
\pgfpathlineto{\pgfqpoint{1.239488in}{0.608125in}}%
\pgfpathlineto{\pgfqpoint{1.240686in}{0.759638in}}%
\pgfpathlineto{\pgfqpoint{1.243082in}{1.432657in}}%
\pgfpathlineto{\pgfqpoint{1.247959in}{2.739742in}}%
\pgfpathlineto{\pgfqpoint{1.248644in}{2.765016in}}%
\pgfpathlineto{\pgfqpoint{1.249157in}{2.749052in}}%
\pgfpathlineto{\pgfqpoint{1.250355in}{2.599490in}}%
\pgfpathlineto{\pgfqpoint{1.252751in}{1.934139in}}%
\pgfpathlineto{\pgfqpoint{1.257713in}{0.616787in}}%
\pgfpathlineto{\pgfqpoint{1.258398in}{0.593427in}}%
\pgfpathlineto{\pgfqpoint{1.258911in}{0.610329in}}%
\pgfpathlineto{\pgfqpoint{1.260109in}{0.760348in}}%
\pgfpathlineto{\pgfqpoint{1.262505in}{1.420954in}}%
\pgfpathlineto{\pgfqpoint{1.267467in}{2.737062in}}%
\pgfpathlineto{\pgfqpoint{1.268152in}{2.763635in}}%
\pgfpathlineto{\pgfqpoint{1.268665in}{2.749641in}}%
\pgfpathlineto{\pgfqpoint{1.269777in}{2.622491in}}%
\pgfpathlineto{\pgfqpoint{1.272002in}{2.046605in}}%
\pgfpathlineto{\pgfqpoint{1.277478in}{0.611440in}}%
\pgfpathlineto{\pgfqpoint{1.277991in}{0.594852in}}%
\pgfpathlineto{\pgfqpoint{1.278505in}{0.606944in}}%
\pgfpathlineto{\pgfqpoint{1.279617in}{0.728808in}}%
\pgfpathlineto{\pgfqpoint{1.281842in}{1.292788in}}%
\pgfpathlineto{\pgfqpoint{1.287403in}{2.744864in}}%
\pgfpathlineto{\pgfqpoint{1.288002in}{2.762151in}}%
\pgfpathlineto{\pgfqpoint{1.288516in}{2.746371in}}%
\pgfpathlineto{\pgfqpoint{1.289714in}{2.603532in}}%
\pgfpathlineto{\pgfqpoint{1.292109in}{1.967803in}}%
\pgfpathlineto{\pgfqpoint{1.297243in}{0.625171in}}%
\pgfpathlineto{\pgfqpoint{1.298013in}{0.596273in}}%
\pgfpathlineto{\pgfqpoint{1.298526in}{0.611752in}}%
\pgfpathlineto{\pgfqpoint{1.299724in}{0.752359in}}%
\pgfpathlineto{\pgfqpoint{1.302035in}{1.352286in}}%
\pgfpathlineto{\pgfqpoint{1.307425in}{2.739148in}}%
\pgfpathlineto{\pgfqpoint{1.308109in}{2.760626in}}%
\pgfpathlineto{\pgfqpoint{1.308623in}{2.744784in}}%
\pgfpathlineto{\pgfqpoint{1.309821in}{2.604957in}}%
\pgfpathlineto{\pgfqpoint{1.312131in}{2.011108in}}%
\pgfpathlineto{\pgfqpoint{1.317521in}{0.623341in}}%
\pgfpathlineto{\pgfqpoint{1.318206in}{0.597770in}}%
\pgfpathlineto{\pgfqpoint{1.318719in}{0.610048in}}%
\pgfpathlineto{\pgfqpoint{1.319832in}{0.726695in}}%
\pgfpathlineto{\pgfqpoint{1.322056in}{1.264326in}}%
\pgfpathlineto{\pgfqpoint{1.327960in}{2.746318in}}%
\pgfpathlineto{\pgfqpoint{1.328473in}{2.759167in}}%
\pgfpathlineto{\pgfqpoint{1.328987in}{2.745452in}}%
\pgfpathlineto{\pgfqpoint{1.330099in}{2.627256in}}%
\pgfpathlineto{\pgfqpoint{1.332324in}{2.092281in}}%
\pgfpathlineto{\pgfqpoint{1.338227in}{0.614378in}}%
\pgfpathlineto{\pgfqpoint{1.338741in}{0.599381in}}%
\pgfpathlineto{\pgfqpoint{1.339254in}{0.610528in}}%
\pgfpathlineto{\pgfqpoint{1.340367in}{0.722070in}}%
\pgfpathlineto{\pgfqpoint{1.342506in}{1.217515in}}%
\pgfpathlineto{\pgfqpoint{1.348752in}{2.748951in}}%
\pgfpathlineto{\pgfqpoint{1.349179in}{2.757574in}}%
\pgfpathlineto{\pgfqpoint{1.349693in}{2.744337in}}%
\pgfpathlineto{\pgfqpoint{1.350805in}{2.629898in}}%
\pgfpathlineto{\pgfqpoint{1.353030in}{2.109781in}}%
\pgfpathlineto{\pgfqpoint{1.359105in}{0.615019in}}%
\pgfpathlineto{\pgfqpoint{1.359618in}{0.600969in}}%
\pgfpathlineto{\pgfqpoint{1.360131in}{0.612227in}}%
\pgfpathlineto{\pgfqpoint{1.361244in}{0.721216in}}%
\pgfpathlineto{\pgfqpoint{1.363383in}{1.203991in}}%
\pgfpathlineto{\pgfqpoint{1.369800in}{2.748218in}}%
\pgfpathlineto{\pgfqpoint{1.370228in}{2.755901in}}%
\pgfpathlineto{\pgfqpoint{1.370656in}{2.746290in}}%
\pgfpathlineto{\pgfqpoint{1.371682in}{2.654276in}}%
\pgfpathlineto{\pgfqpoint{1.373736in}{2.219115in}}%
\pgfpathlineto{\pgfqpoint{1.380666in}{0.604114in}}%
\pgfpathlineto{\pgfqpoint{1.380838in}{0.602613in}}%
\pgfpathlineto{\pgfqpoint{1.381180in}{0.607777in}}%
\pgfpathlineto{\pgfqpoint{1.382121in}{0.677231in}}%
\pgfpathlineto{\pgfqpoint{1.384003in}{1.032850in}}%
\pgfpathlineto{\pgfqpoint{1.391533in}{2.754137in}}%
\pgfpathlineto{\pgfqpoint{1.391790in}{2.752128in}}%
\pgfpathlineto{\pgfqpoint{1.392560in}{2.710285in}}%
\pgfpathlineto{\pgfqpoint{1.394185in}{2.458167in}}%
\pgfpathlineto{\pgfqpoint{1.397693in}{1.456084in}}%
\pgfpathlineto{\pgfqpoint{1.401458in}{0.643469in}}%
\pgfpathlineto{\pgfqpoint{1.402399in}{0.604348in}}%
\pgfpathlineto{\pgfqpoint{1.402827in}{0.612772in}}%
\pgfpathlineto{\pgfqpoint{1.403854in}{0.698615in}}%
\pgfpathlineto{\pgfqpoint{1.405907in}{1.111598in}}%
\pgfpathlineto{\pgfqpoint{1.413180in}{2.751788in}}%
\pgfpathlineto{\pgfqpoint{1.413266in}{2.752391in}}%
\pgfpathlineto{\pgfqpoint{1.413608in}{2.748344in}}%
\pgfpathlineto{\pgfqpoint{1.414464in}{2.693574in}}%
\pgfpathlineto{\pgfqpoint{1.416175in}{2.409223in}}%
\pgfpathlineto{\pgfqpoint{1.420539in}{1.166238in}}%
\pgfpathlineto{\pgfqpoint{1.423619in}{0.626511in}}%
\pgfpathlineto{\pgfqpoint{1.424303in}{0.606147in}}%
\pgfpathlineto{\pgfqpoint{1.424817in}{0.617472in}}%
\pgfpathlineto{\pgfqpoint{1.425929in}{0.718324in}}%
\pgfpathlineto{\pgfqpoint{1.428068in}{1.162138in}}%
\pgfpathlineto{\pgfqpoint{1.435084in}{2.746304in}}%
\pgfpathlineto{\pgfqpoint{1.435426in}{2.750595in}}%
\pgfpathlineto{\pgfqpoint{1.435854in}{2.741951in}}%
\pgfpathlineto{\pgfqpoint{1.436881in}{2.659023in}}%
\pgfpathlineto{\pgfqpoint{1.438934in}{2.263324in}}%
\pgfpathlineto{\pgfqpoint{1.446550in}{0.608110in}}%
\pgfpathlineto{\pgfqpoint{1.446635in}{0.608070in}}%
\pgfpathlineto{\pgfqpoint{1.446721in}{0.608642in}}%
\pgfpathlineto{\pgfqpoint{1.447320in}{0.629692in}}%
\pgfpathlineto{\pgfqpoint{1.448689in}{0.785240in}}%
\pgfpathlineto{\pgfqpoint{1.451341in}{1.415635in}}%
\pgfpathlineto{\pgfqpoint{1.456903in}{2.708684in}}%
\pgfpathlineto{\pgfqpoint{1.457929in}{2.748667in}}%
\pgfpathlineto{\pgfqpoint{1.458357in}{2.739913in}}%
\pgfpathlineto{\pgfqpoint{1.459469in}{2.648583in}}%
\pgfpathlineto{\pgfqpoint{1.461609in}{2.232356in}}%
\pgfpathlineto{\pgfqpoint{1.469052in}{0.612244in}}%
\pgfpathlineto{\pgfqpoint{1.469309in}{0.609959in}}%
\pgfpathlineto{\pgfqpoint{1.469737in}{0.617918in}}%
\pgfpathlineto{\pgfqpoint{1.470764in}{0.695861in}}%
\pgfpathlineto{\pgfqpoint{1.472732in}{1.050554in}}%
\pgfpathlineto{\pgfqpoint{1.480775in}{2.746734in}}%
\pgfpathlineto{\pgfqpoint{1.480946in}{2.745855in}}%
\pgfpathlineto{\pgfqpoint{1.481630in}{2.719367in}}%
\pgfpathlineto{\pgfqpoint{1.483085in}{2.546866in}}%
\pgfpathlineto{\pgfqpoint{1.485908in}{1.874161in}}%
\pgfpathlineto{\pgfqpoint{1.491213in}{0.666629in}}%
\pgfpathlineto{\pgfqpoint{1.492411in}{0.611954in}}%
\pgfpathlineto{\pgfqpoint{1.492839in}{0.619170in}}%
\pgfpathlineto{\pgfqpoint{1.493866in}{0.693139in}}%
\pgfpathlineto{\pgfqpoint{1.495833in}{1.033716in}}%
\pgfpathlineto{\pgfqpoint{1.504133in}{2.744706in}}%
\pgfpathlineto{\pgfqpoint{1.504304in}{2.743545in}}%
\pgfpathlineto{\pgfqpoint{1.504989in}{2.716841in}}%
\pgfpathlineto{\pgfqpoint{1.506443in}{2.548446in}}%
\pgfpathlineto{\pgfqpoint{1.509267in}{1.894405in}}%
\pgfpathlineto{\pgfqpoint{1.514743in}{0.669156in}}%
\pgfpathlineto{\pgfqpoint{1.515941in}{0.614058in}}%
\pgfpathlineto{\pgfqpoint{1.516368in}{0.620053in}}%
\pgfpathlineto{\pgfqpoint{1.517395in}{0.688942in}}%
\pgfpathlineto{\pgfqpoint{1.519363in}{1.013603in}}%
\pgfpathlineto{\pgfqpoint{1.527919in}{2.742584in}}%
\pgfpathlineto{\pgfqpoint{1.528262in}{2.738328in}}%
\pgfpathlineto{\pgfqpoint{1.529203in}{2.683299in}}%
\pgfpathlineto{\pgfqpoint{1.531000in}{2.416596in}}%
\pgfpathlineto{\pgfqpoint{1.535021in}{1.384520in}}%
\pgfpathlineto{\pgfqpoint{1.538786in}{0.667769in}}%
\pgfpathlineto{\pgfqpoint{1.539984in}{0.616201in}}%
\pgfpathlineto{\pgfqpoint{1.540412in}{0.622410in}}%
\pgfpathlineto{\pgfqpoint{1.541438in}{0.689561in}}%
\pgfpathlineto{\pgfqpoint{1.543406in}{1.003495in}}%
\pgfpathlineto{\pgfqpoint{1.552219in}{2.740353in}}%
\pgfpathlineto{\pgfqpoint{1.552561in}{2.735619in}}%
\pgfpathlineto{\pgfqpoint{1.553503in}{2.681122in}}%
\pgfpathlineto{\pgfqpoint{1.555385in}{2.405349in}}%
\pgfpathlineto{\pgfqpoint{1.559578in}{1.349783in}}%
\pgfpathlineto{\pgfqpoint{1.563342in}{0.665193in}}%
\pgfpathlineto{\pgfqpoint{1.564540in}{0.618464in}}%
\pgfpathlineto{\pgfqpoint{1.564968in}{0.625364in}}%
\pgfpathlineto{\pgfqpoint{1.565995in}{0.691901in}}%
\pgfpathlineto{\pgfqpoint{1.567963in}{0.996887in}}%
\pgfpathlineto{\pgfqpoint{1.576947in}{2.738063in}}%
\pgfpathlineto{\pgfqpoint{1.577375in}{2.732743in}}%
\pgfpathlineto{\pgfqpoint{1.578401in}{2.670992in}}%
\pgfpathlineto{\pgfqpoint{1.580369in}{2.377611in}}%
\pgfpathlineto{\pgfqpoint{1.585246in}{1.174544in}}%
\pgfpathlineto{\pgfqpoint{1.588583in}{0.652106in}}%
\pgfpathlineto{\pgfqpoint{1.589524in}{0.620839in}}%
\pgfpathlineto{\pgfqpoint{1.590038in}{0.628027in}}%
\pgfpathlineto{\pgfqpoint{1.591150in}{0.701371in}}%
\pgfpathlineto{\pgfqpoint{1.593204in}{1.022067in}}%
\pgfpathlineto{\pgfqpoint{1.602273in}{2.735684in}}%
\pgfpathlineto{\pgfqpoint{1.602530in}{2.733980in}}%
\pgfpathlineto{\pgfqpoint{1.603300in}{2.703897in}}%
\pgfpathlineto{\pgfqpoint{1.604840in}{2.536997in}}%
\pgfpathlineto{\pgfqpoint{1.607749in}{1.926346in}}%
\pgfpathlineto{\pgfqpoint{1.613739in}{0.686528in}}%
\pgfpathlineto{\pgfqpoint{1.615193in}{0.623230in}}%
\pgfpathlineto{\pgfqpoint{1.615535in}{0.627308in}}%
\pgfpathlineto{\pgfqpoint{1.616477in}{0.675566in}}%
\pgfpathlineto{\pgfqpoint{1.618273in}{0.907375in}}%
\pgfpathlineto{\pgfqpoint{1.621867in}{1.729774in}}%
\pgfpathlineto{\pgfqpoint{1.626573in}{2.654240in}}%
\pgfpathlineto{\pgfqpoint{1.628199in}{2.733219in}}%
\pgfpathlineto{\pgfqpoint{1.628455in}{2.731178in}}%
\pgfpathlineto{\pgfqpoint{1.629311in}{2.695777in}}%
\pgfpathlineto{\pgfqpoint{1.631022in}{2.500729in}}%
\pgfpathlineto{\pgfqpoint{1.634188in}{1.823729in}}%
\pgfpathlineto{\pgfqpoint{1.639750in}{0.702827in}}%
\pgfpathlineto{\pgfqpoint{1.641375in}{0.625752in}}%
\pgfpathlineto{\pgfqpoint{1.641632in}{0.627741in}}%
\pgfpathlineto{\pgfqpoint{1.642488in}{0.662273in}}%
\pgfpathlineto{\pgfqpoint{1.644199in}{0.852716in}}%
\pgfpathlineto{\pgfqpoint{1.647365in}{1.516530in}}%
\pgfpathlineto{\pgfqpoint{1.653097in}{2.656620in}}%
\pgfpathlineto{\pgfqpoint{1.654723in}{2.730624in}}%
\pgfpathlineto{\pgfqpoint{1.654980in}{2.728498in}}%
\pgfpathlineto{\pgfqpoint{1.655835in}{2.694213in}}%
\pgfpathlineto{\pgfqpoint{1.657547in}{2.507251in}}%
\pgfpathlineto{\pgfqpoint{1.660712in}{1.855579in}}%
\pgfpathlineto{\pgfqpoint{1.666531in}{0.706453in}}%
\pgfpathlineto{\pgfqpoint{1.668242in}{0.628431in}}%
\pgfpathlineto{\pgfqpoint{1.668413in}{0.629594in}}%
\pgfpathlineto{\pgfqpoint{1.669183in}{0.655054in}}%
\pgfpathlineto{\pgfqpoint{1.670723in}{0.800979in}}%
\pgfpathlineto{\pgfqpoint{1.673547in}{1.329643in}}%
\pgfpathlineto{\pgfqpoint{1.680563in}{2.681126in}}%
\pgfpathlineto{\pgfqpoint{1.681846in}{2.727908in}}%
\pgfpathlineto{\pgfqpoint{1.682274in}{2.723571in}}%
\pgfpathlineto{\pgfqpoint{1.683301in}{2.672851in}}%
\pgfpathlineto{\pgfqpoint{1.685183in}{2.443301in}}%
\pgfpathlineto{\pgfqpoint{1.688948in}{1.637999in}}%
\pgfpathlineto{\pgfqpoint{1.693911in}{0.717161in}}%
\pgfpathlineto{\pgfqpoint{1.695707in}{0.631159in}}%
\pgfpathlineto{\pgfqpoint{1.695879in}{0.631837in}}%
\pgfpathlineto{\pgfqpoint{1.696563in}{0.650071in}}%
\pgfpathlineto{\pgfqpoint{1.698018in}{0.768470in}}%
\pgfpathlineto{\pgfqpoint{1.700584in}{1.198399in}}%
\pgfpathlineto{\pgfqpoint{1.708798in}{2.701998in}}%
\pgfpathlineto{\pgfqpoint{1.709740in}{2.725108in}}%
\pgfpathlineto{\pgfqpoint{1.710253in}{2.718418in}}%
\pgfpathlineto{\pgfqpoint{1.711365in}{2.657898in}}%
\pgfpathlineto{\pgfqpoint{1.713419in}{2.395624in}}%
\pgfpathlineto{\pgfqpoint{1.717697in}{1.479736in}}%
\pgfpathlineto{\pgfqpoint{1.722146in}{0.715818in}}%
\pgfpathlineto{\pgfqpoint{1.723943in}{0.634047in}}%
\pgfpathlineto{\pgfqpoint{1.724114in}{0.634647in}}%
\pgfpathlineto{\pgfqpoint{1.724799in}{0.651729in}}%
\pgfpathlineto{\pgfqpoint{1.726253in}{0.763481in}}%
\pgfpathlineto{\pgfqpoint{1.728820in}{1.172144in}}%
\pgfpathlineto{\pgfqpoint{1.737547in}{2.705672in}}%
\pgfpathlineto{\pgfqpoint{1.738403in}{2.722130in}}%
\pgfpathlineto{\pgfqpoint{1.738831in}{2.716955in}}%
\pgfpathlineto{\pgfqpoint{1.739943in}{2.662271in}}%
\pgfpathlineto{\pgfqpoint{1.741911in}{2.431190in}}%
\pgfpathlineto{\pgfqpoint{1.745761in}{1.654447in}}%
\pgfpathlineto{\pgfqpoint{1.750981in}{0.730988in}}%
\pgfpathlineto{\pgfqpoint{1.752949in}{0.637088in}}%
\pgfpathlineto{\pgfqpoint{1.753120in}{0.637501in}}%
\pgfpathlineto{\pgfqpoint{1.753719in}{0.649856in}}%
\pgfpathlineto{\pgfqpoint{1.755088in}{0.740168in}}%
\pgfpathlineto{\pgfqpoint{1.757483in}{1.079866in}}%
\pgfpathlineto{\pgfqpoint{1.767580in}{2.717855in}}%
\pgfpathlineto{\pgfqpoint{1.767837in}{2.719050in}}%
\pgfpathlineto{\pgfqpoint{1.768264in}{2.714308in}}%
\pgfpathlineto{\pgfqpoint{1.769291in}{2.669022in}}%
\pgfpathlineto{\pgfqpoint{1.771259in}{2.457910in}}%
\pgfpathlineto{\pgfqpoint{1.774853in}{1.777288in}}%
\pgfpathlineto{\pgfqpoint{1.780756in}{0.737192in}}%
\pgfpathlineto{\pgfqpoint{1.782896in}{0.640259in}}%
\pgfpathlineto{\pgfqpoint{1.783323in}{0.645091in}}%
\pgfpathlineto{\pgfqpoint{1.784436in}{0.695293in}}%
\pgfpathlineto{\pgfqpoint{1.786404in}{0.907524in}}%
\pgfpathlineto{\pgfqpoint{1.790083in}{1.596857in}}%
\pgfpathlineto{\pgfqpoint{1.795987in}{2.617760in}}%
\pgfpathlineto{\pgfqpoint{1.798126in}{2.715825in}}%
\pgfpathlineto{\pgfqpoint{1.798468in}{2.713397in}}%
\pgfpathlineto{\pgfqpoint{1.799409in}{2.680880in}}%
\pgfpathlineto{\pgfqpoint{1.801206in}{2.519095in}}%
\pgfpathlineto{\pgfqpoint{1.804372in}{1.987159in}}%
\pgfpathlineto{\pgfqpoint{1.811816in}{0.713091in}}%
\pgfpathlineto{\pgfqpoint{1.813612in}{0.643574in}}%
\pgfpathlineto{\pgfqpoint{1.813784in}{0.643925in}}%
\pgfpathlineto{\pgfqpoint{1.814383in}{0.654757in}}%
\pgfpathlineto{\pgfqpoint{1.815752in}{0.734310in}}%
\pgfpathlineto{\pgfqpoint{1.818147in}{1.036380in}}%
\pgfpathlineto{\pgfqpoint{1.824051in}{2.180994in}}%
\pgfpathlineto{\pgfqpoint{1.827816in}{2.661062in}}%
\pgfpathlineto{\pgfqpoint{1.829442in}{2.712407in}}%
\pgfpathlineto{\pgfqpoint{1.829698in}{2.710816in}}%
\pgfpathlineto{\pgfqpoint{1.830554in}{2.686441in}}%
\pgfpathlineto{\pgfqpoint{1.832180in}{2.562541in}}%
\pgfpathlineto{\pgfqpoint{1.835003in}{2.147045in}}%
\pgfpathlineto{\pgfqpoint{1.844158in}{0.679225in}}%
\pgfpathlineto{\pgfqpoint{1.845442in}{0.647041in}}%
\pgfpathlineto{\pgfqpoint{1.845784in}{0.649266in}}%
\pgfpathlineto{\pgfqpoint{1.846725in}{0.678720in}}%
\pgfpathlineto{\pgfqpoint{1.848522in}{0.825393in}}%
\pgfpathlineto{\pgfqpoint{1.851602in}{1.297562in}}%
\pgfpathlineto{\pgfqpoint{1.859987in}{2.650857in}}%
\pgfpathlineto{\pgfqpoint{1.861784in}{2.708812in}}%
\pgfpathlineto{\pgfqpoint{1.861955in}{2.708045in}}%
\pgfpathlineto{\pgfqpoint{1.862725in}{2.691043in}}%
\pgfpathlineto{\pgfqpoint{1.864265in}{2.592562in}}%
\pgfpathlineto{\pgfqpoint{1.866918in}{2.248988in}}%
\pgfpathlineto{\pgfqpoint{1.877699in}{0.658238in}}%
\pgfpathlineto{\pgfqpoint{1.878383in}{0.650752in}}%
\pgfpathlineto{\pgfqpoint{1.878897in}{0.656259in}}%
\pgfpathlineto{\pgfqpoint{1.880094in}{0.705632in}}%
\pgfpathlineto{\pgfqpoint{1.882234in}{0.911605in}}%
\pgfpathlineto{\pgfqpoint{1.886084in}{1.550847in}}%
\pgfpathlineto{\pgfqpoint{1.892672in}{2.592096in}}%
\pgfpathlineto{\pgfqpoint{1.895068in}{2.704539in}}%
\pgfpathlineto{\pgfqpoint{1.895239in}{2.705049in}}%
\pgfpathlineto{\pgfqpoint{1.895581in}{2.703010in}}%
\pgfpathlineto{\pgfqpoint{1.896522in}{2.676537in}}%
\pgfpathlineto{\pgfqpoint{1.898319in}{2.544763in}}%
\pgfpathlineto{\pgfqpoint{1.901400in}{2.115901in}}%
\pgfpathlineto{\pgfqpoint{1.910811in}{0.699499in}}%
\pgfpathlineto{\pgfqpoint{1.912437in}{0.654619in}}%
\pgfpathlineto{\pgfqpoint{1.912694in}{0.655579in}}%
\pgfpathlineto{\pgfqpoint{1.913464in}{0.671623in}}%
\pgfpathlineto{\pgfqpoint{1.915004in}{0.761218in}}%
\pgfpathlineto{\pgfqpoint{1.917656in}{1.072737in}}%
\pgfpathlineto{\pgfqpoint{1.929635in}{2.698856in}}%
\pgfpathlineto{\pgfqpoint{1.929977in}{2.701034in}}%
\pgfpathlineto{\pgfqpoint{1.930491in}{2.697252in}}%
\pgfpathlineto{\pgfqpoint{1.931603in}{2.660363in}}%
\pgfpathlineto{\pgfqpoint{1.933571in}{2.504338in}}%
\pgfpathlineto{\pgfqpoint{1.936994in}{2.017131in}}%
\pgfpathlineto{\pgfqpoint{1.945550in}{0.743491in}}%
\pgfpathlineto{\pgfqpoint{1.947860in}{0.658775in}}%
\pgfpathlineto{\pgfqpoint{1.948031in}{0.658961in}}%
\pgfpathlineto{\pgfqpoint{1.948117in}{0.659391in}}%
\pgfpathlineto{\pgfqpoint{1.948887in}{0.673352in}}%
\pgfpathlineto{\pgfqpoint{1.950427in}{0.754330in}}%
\pgfpathlineto{\pgfqpoint{1.952994in}{1.028932in}}%
\pgfpathlineto{\pgfqpoint{1.958299in}{1.903060in}}%
\pgfpathlineto{\pgfqpoint{1.963347in}{2.579603in}}%
\pgfpathlineto{\pgfqpoint{1.965828in}{2.694853in}}%
\pgfpathlineto{\pgfqpoint{1.966170in}{2.696792in}}%
\pgfpathlineto{\pgfqpoint{1.966684in}{2.693241in}}%
\pgfpathlineto{\pgfqpoint{1.967796in}{2.659236in}}%
\pgfpathlineto{\pgfqpoint{1.969764in}{2.515515in}}%
\pgfpathlineto{\pgfqpoint{1.973101in}{2.076299in}}%
\pgfpathlineto{\pgfqpoint{1.982684in}{0.729811in}}%
\pgfpathlineto{\pgfqpoint{1.984908in}{0.663134in}}%
\pgfpathlineto{\pgfqpoint{1.985336in}{0.666146in}}%
\pgfpathlineto{\pgfqpoint{1.986449in}{0.697789in}}%
\pgfpathlineto{\pgfqpoint{1.988417in}{0.833881in}}%
\pgfpathlineto{\pgfqpoint{1.991668in}{1.241399in}}%
\pgfpathlineto{\pgfqpoint{2.001935in}{2.636925in}}%
\pgfpathlineto{\pgfqpoint{2.003989in}{2.692276in}}%
\pgfpathlineto{\pgfqpoint{2.004074in}{2.692147in}}%
\pgfpathlineto{\pgfqpoint{2.004673in}{2.685762in}}%
\pgfpathlineto{\pgfqpoint{2.005957in}{2.640270in}}%
\pgfpathlineto{\pgfqpoint{2.008181in}{2.465604in}}%
\pgfpathlineto{\pgfqpoint{2.011946in}{1.960497in}}%
\pgfpathlineto{\pgfqpoint{2.020588in}{0.776466in}}%
\pgfpathlineto{\pgfqpoint{2.023155in}{0.669579in}}%
\pgfpathlineto{\pgfqpoint{2.023497in}{0.667799in}}%
\pgfpathlineto{\pgfqpoint{2.024011in}{0.670723in}}%
\pgfpathlineto{\pgfqpoint{2.025123in}{0.699865in}}%
\pgfpathlineto{\pgfqpoint{2.027091in}{0.824328in}}%
\pgfpathlineto{\pgfqpoint{2.030342in}{1.199193in}}%
\pgfpathlineto{\pgfqpoint{2.041722in}{2.646457in}}%
\pgfpathlineto{\pgfqpoint{2.043604in}{2.687432in}}%
\pgfpathlineto{\pgfqpoint{2.043690in}{2.687267in}}%
\pgfpathlineto{\pgfqpoint{2.044289in}{2.681166in}}%
\pgfpathlineto{\pgfqpoint{2.045572in}{2.639363in}}%
\pgfpathlineto{\pgfqpoint{2.047797in}{2.479915in}}%
\pgfpathlineto{\pgfqpoint{2.051476in}{2.027284in}}%
\pgfpathlineto{\pgfqpoint{2.061230in}{0.767341in}}%
\pgfpathlineto{\pgfqpoint{2.063797in}{0.673992in}}%
\pgfpathlineto{\pgfqpoint{2.064139in}{0.672791in}}%
\pgfpathlineto{\pgfqpoint{2.064567in}{0.675058in}}%
\pgfpathlineto{\pgfqpoint{2.065594in}{0.697454in}}%
\pgfpathlineto{\pgfqpoint{2.067476in}{0.798332in}}%
\pgfpathlineto{\pgfqpoint{2.070557in}{1.107202in}}%
\pgfpathlineto{\pgfqpoint{2.084247in}{2.671778in}}%
\pgfpathlineto{\pgfqpoint{2.085273in}{2.682244in}}%
\pgfpathlineto{\pgfqpoint{2.085701in}{2.679882in}}%
\pgfpathlineto{\pgfqpoint{2.086813in}{2.655394in}}%
\pgfpathlineto{\pgfqpoint{2.088781in}{2.549741in}}%
\pgfpathlineto{\pgfqpoint{2.091947in}{2.236817in}}%
\pgfpathlineto{\pgfqpoint{2.105723in}{0.693276in}}%
\pgfpathlineto{\pgfqpoint{2.106921in}{0.678175in}}%
\pgfpathlineto{\pgfqpoint{2.107348in}{0.679855in}}%
\pgfpathlineto{\pgfqpoint{2.108375in}{0.698979in}}%
\pgfpathlineto{\pgfqpoint{2.110257in}{0.787522in}}%
\pgfpathlineto{\pgfqpoint{2.113252in}{1.053736in}}%
\pgfpathlineto{\pgfqpoint{2.119070in}{1.822480in}}%
\pgfpathlineto{\pgfqpoint{2.125145in}{2.515233in}}%
\pgfpathlineto{\pgfqpoint{2.128226in}{2.665620in}}%
\pgfpathlineto{\pgfqpoint{2.129338in}{2.676652in}}%
\pgfpathlineto{\pgfqpoint{2.129766in}{2.674604in}}%
\pgfpathlineto{\pgfqpoint{2.130878in}{2.653055in}}%
\pgfpathlineto{\pgfqpoint{2.132846in}{2.559596in}}%
\pgfpathlineto{\pgfqpoint{2.135926in}{2.289407in}}%
\pgfpathlineto{\pgfqpoint{2.142001in}{1.509919in}}%
\pgfpathlineto{\pgfqpoint{2.148076in}{0.845065in}}%
\pgfpathlineto{\pgfqpoint{2.151242in}{0.695322in}}%
\pgfpathlineto{\pgfqpoint{2.152354in}{0.684002in}}%
\pgfpathlineto{\pgfqpoint{2.152782in}{0.685533in}}%
\pgfpathlineto{\pgfqpoint{2.153809in}{0.702456in}}%
\pgfpathlineto{\pgfqpoint{2.155691in}{0.780585in}}%
\pgfpathlineto{\pgfqpoint{2.158686in}{1.016854in}}%
\pgfpathlineto{\pgfqpoint{2.163991in}{1.652224in}}%
\pgfpathlineto{\pgfqpoint{2.171435in}{2.485670in}}%
\pgfpathlineto{\pgfqpoint{2.174772in}{2.653434in}}%
\pgfpathlineto{\pgfqpoint{2.176226in}{2.670579in}}%
\pgfpathlineto{\pgfqpoint{2.176568in}{2.669499in}}%
\pgfpathlineto{\pgfqpoint{2.177510in}{2.656536in}}%
\pgfpathlineto{\pgfqpoint{2.179306in}{2.592110in}}%
\pgfpathlineto{\pgfqpoint{2.182130in}{2.395144in}}%
\pgfpathlineto{\pgfqpoint{2.186836in}{1.880976in}}%
\pgfpathlineto{\pgfqpoint{2.196162in}{0.859600in}}%
\pgfpathlineto{\pgfqpoint{2.199499in}{0.705444in}}%
\pgfpathlineto{\pgfqpoint{2.200868in}{0.690370in}}%
\pgfpathlineto{\pgfqpoint{2.201210in}{0.691118in}}%
\pgfpathlineto{\pgfqpoint{2.202152in}{0.702449in}}%
\pgfpathlineto{\pgfqpoint{2.203863in}{0.757128in}}%
\pgfpathlineto{\pgfqpoint{2.206601in}{0.928830in}}%
\pgfpathlineto{\pgfqpoint{2.210965in}{1.365086in}}%
\pgfpathlineto{\pgfqpoint{2.222173in}{2.531099in}}%
\pgfpathlineto{\pgfqpoint{2.225339in}{2.653853in}}%
\pgfpathlineto{\pgfqpoint{2.226537in}{2.663912in}}%
\pgfpathlineto{\pgfqpoint{2.226965in}{2.662555in}}%
\pgfpathlineto{\pgfqpoint{2.227991in}{2.648741in}}%
\pgfpathlineto{\pgfqpoint{2.229874in}{2.585720in}}%
\pgfpathlineto{\pgfqpoint{2.232783in}{2.400888in}}%
\pgfpathlineto{\pgfqpoint{2.237489in}{1.936830in}}%
\pgfpathlineto{\pgfqpoint{2.248184in}{0.858704in}}%
\pgfpathlineto{\pgfqpoint{2.251607in}{0.714061in}}%
\pgfpathlineto{\pgfqpoint{2.253232in}{0.697377in}}%
\pgfpathlineto{\pgfqpoint{2.253489in}{0.697895in}}%
\pgfpathlineto{\pgfqpoint{2.254345in}{0.705806in}}%
\pgfpathlineto{\pgfqpoint{2.255970in}{0.746594in}}%
\pgfpathlineto{\pgfqpoint{2.258623in}{0.881080in}}%
\pgfpathlineto{\pgfqpoint{2.262644in}{1.216031in}}%
\pgfpathlineto{\pgfqpoint{2.277703in}{2.589889in}}%
\pgfpathlineto{\pgfqpoint{2.280441in}{2.654234in}}%
\pgfpathlineto{\pgfqpoint{2.281040in}{2.656519in}}%
\pgfpathlineto{\pgfqpoint{2.281554in}{2.655083in}}%
\pgfpathlineto{\pgfqpoint{2.282666in}{2.641306in}}%
\pgfpathlineto{\pgfqpoint{2.284634in}{2.582118in}}%
\pgfpathlineto{\pgfqpoint{2.287628in}{2.413920in}}%
\pgfpathlineto{\pgfqpoint{2.292334in}{2.004991in}}%
\pgfpathlineto{\pgfqpoint{2.304998in}{0.850991in}}%
\pgfpathlineto{\pgfqpoint{2.308506in}{0.721568in}}%
\pgfpathlineto{\pgfqpoint{2.310217in}{0.705201in}}%
\pgfpathlineto{\pgfqpoint{2.310474in}{0.705458in}}%
\pgfpathlineto{\pgfqpoint{2.311244in}{0.710458in}}%
\pgfpathlineto{\pgfqpoint{2.312784in}{0.739240in}}%
\pgfpathlineto{\pgfqpoint{2.315265in}{0.835828in}}%
\pgfpathlineto{\pgfqpoint{2.319030in}{1.084626in}}%
\pgfpathlineto{\pgfqpoint{2.325704in}{1.706892in}}%
\pgfpathlineto{\pgfqpoint{2.334003in}{2.421287in}}%
\pgfpathlineto{\pgfqpoint{2.338110in}{2.608882in}}%
\pgfpathlineto{\pgfqpoint{2.340592in}{2.647504in}}%
\pgfpathlineto{\pgfqpoint{2.341019in}{2.648198in}}%
\pgfpathlineto{\pgfqpoint{2.341447in}{2.647136in}}%
\pgfpathlineto{\pgfqpoint{2.342560in}{2.636208in}}%
\pgfpathlineto{\pgfqpoint{2.344528in}{2.588630in}}%
\pgfpathlineto{\pgfqpoint{2.347522in}{2.451926in}}%
\pgfpathlineto{\pgfqpoint{2.352057in}{2.125919in}}%
\pgfpathlineto{\pgfqpoint{2.368913in}{0.801950in}}%
\pgfpathlineto{\pgfqpoint{2.372079in}{0.722398in}}%
\pgfpathlineto{\pgfqpoint{2.373448in}{0.714079in}}%
\pgfpathlineto{\pgfqpoint{2.373790in}{0.714488in}}%
\pgfpathlineto{\pgfqpoint{2.374731in}{0.720713in}}%
\pgfpathlineto{\pgfqpoint{2.376442in}{0.750906in}}%
\pgfpathlineto{\pgfqpoint{2.379095in}{0.843466in}}%
\pgfpathlineto{\pgfqpoint{2.383031in}{1.069746in}}%
\pgfpathlineto{\pgfqpoint{2.389619in}{1.600666in}}%
\pgfpathlineto{\pgfqpoint{2.399715in}{2.386882in}}%
\pgfpathlineto{\pgfqpoint{2.404250in}{2.584313in}}%
\pgfpathlineto{\pgfqpoint{2.407159in}{2.635470in}}%
\pgfpathlineto{\pgfqpoint{2.408100in}{2.638660in}}%
\pgfpathlineto{\pgfqpoint{2.408614in}{2.637639in}}%
\pgfpathlineto{\pgfqpoint{2.409726in}{2.628788in}}%
\pgfpathlineto{\pgfqpoint{2.411694in}{2.591339in}}%
\pgfpathlineto{\pgfqpoint{2.414689in}{2.484128in}}%
\pgfpathlineto{\pgfqpoint{2.419052in}{2.236112in}}%
\pgfpathlineto{\pgfqpoint{2.426753in}{1.644913in}}%
\pgfpathlineto{\pgfqpoint{2.435994in}{0.987038in}}%
\pgfpathlineto{\pgfqpoint{2.440700in}{0.789432in}}%
\pgfpathlineto{\pgfqpoint{2.443866in}{0.730393in}}%
\pgfpathlineto{\pgfqpoint{2.445235in}{0.724420in}}%
\pgfpathlineto{\pgfqpoint{2.445577in}{0.724780in}}%
\pgfpathlineto{\pgfqpoint{2.446518in}{0.729568in}}%
\pgfpathlineto{\pgfqpoint{2.448229in}{0.752347in}}%
\pgfpathlineto{\pgfqpoint{2.450882in}{0.822064in}}%
\pgfpathlineto{\pgfqpoint{2.454732in}{0.989927in}}%
\pgfpathlineto{\pgfqpoint{2.460550in}{1.355196in}}%
\pgfpathlineto{\pgfqpoint{2.476893in}{2.433162in}}%
\pgfpathlineto{\pgfqpoint{2.481427in}{2.582648in}}%
\pgfpathlineto{\pgfqpoint{2.484422in}{2.624028in}}%
\pgfpathlineto{\pgfqpoint{2.485534in}{2.627341in}}%
\pgfpathlineto{\pgfqpoint{2.485962in}{2.626876in}}%
\pgfpathlineto{\pgfqpoint{2.486989in}{2.621849in}}%
\pgfpathlineto{\pgfqpoint{2.488871in}{2.598513in}}%
\pgfpathlineto{\pgfqpoint{2.491695in}{2.530722in}}%
\pgfpathlineto{\pgfqpoint{2.495716in}{2.373328in}}%
\pgfpathlineto{\pgfqpoint{2.501706in}{2.038598in}}%
\pgfpathlineto{\pgfqpoint{2.520529in}{0.921352in}}%
\pgfpathlineto{\pgfqpoint{2.525235in}{0.783238in}}%
\pgfpathlineto{\pgfqpoint{2.528487in}{0.741197in}}%
\pgfpathlineto{\pgfqpoint{2.529941in}{0.736979in}}%
\pgfpathlineto{\pgfqpoint{2.530284in}{0.737294in}}%
\pgfpathlineto{\pgfqpoint{2.531310in}{0.741207in}}%
\pgfpathlineto{\pgfqpoint{2.533107in}{0.758612in}}%
\pgfpathlineto{\pgfqpoint{2.535845in}{0.810021in}}%
\pgfpathlineto{\pgfqpoint{2.539781in}{0.932074in}}%
\pgfpathlineto{\pgfqpoint{2.545343in}{1.183131in}}%
\pgfpathlineto{\pgfqpoint{2.555610in}{1.771299in}}%
\pgfpathlineto{\pgfqpoint{2.565792in}{2.303845in}}%
\pgfpathlineto{\pgfqpoint{2.571696in}{2.508084in}}%
\pgfpathlineto{\pgfqpoint{2.575974in}{2.590178in}}%
\pgfpathlineto{\pgfqpoint{2.578883in}{2.612008in}}%
\pgfpathlineto{\pgfqpoint{2.579824in}{2.613136in}}%
\pgfpathlineto{\pgfqpoint{2.580252in}{2.612697in}}%
\pgfpathlineto{\pgfqpoint{2.581450in}{2.608327in}}%
\pgfpathlineto{\pgfqpoint{2.583503in}{2.590255in}}%
\pgfpathlineto{\pgfqpoint{2.586498in}{2.540914in}}%
\pgfpathlineto{\pgfqpoint{2.590691in}{2.430012in}}%
\pgfpathlineto{\pgfqpoint{2.596509in}{2.209786in}}%
\pgfpathlineto{\pgfqpoint{2.606177in}{1.743755in}}%
\pgfpathlineto{\pgfqpoint{2.619696in}{1.113329in}}%
\pgfpathlineto{\pgfqpoint{2.626456in}{0.895361in}}%
\pgfpathlineto{\pgfqpoint{2.631504in}{0.795626in}}%
\pgfpathlineto{\pgfqpoint{2.635183in}{0.759855in}}%
\pgfpathlineto{\pgfqpoint{2.637493in}{0.753472in}}%
\pgfpathlineto{\pgfqpoint{2.638520in}{0.754570in}}%
\pgfpathlineto{\pgfqpoint{2.640146in}{0.761172in}}%
\pgfpathlineto{\pgfqpoint{2.642627in}{0.782448in}}%
\pgfpathlineto{\pgfqpoint{2.646135in}{0.834431in}}%
\pgfpathlineto{\pgfqpoint{2.650927in}{0.942753in}}%
\pgfpathlineto{\pgfqpoint{2.657429in}{1.145600in}}%
\pgfpathlineto{\pgfqpoint{2.667954in}{1.554973in}}%
\pgfpathlineto{\pgfqpoint{2.684382in}{2.183607in}}%
\pgfpathlineto{\pgfqpoint{2.692253in}{2.403174in}}%
\pgfpathlineto{\pgfqpoint{2.698414in}{2.519023in}}%
\pgfpathlineto{\pgfqpoint{2.703205in}{2.572156in}}%
\pgfpathlineto{\pgfqpoint{2.706713in}{2.590338in}}%
\pgfpathlineto{\pgfqpoint{2.708852in}{2.592981in}}%
\pgfpathlineto{\pgfqpoint{2.710136in}{2.591566in}}%
\pgfpathlineto{\pgfqpoint{2.712104in}{2.585131in}}%
\pgfpathlineto{\pgfqpoint{2.715013in}{2.566476in}}%
\pgfpathlineto{\pgfqpoint{2.718949in}{2.524906in}}%
\pgfpathlineto{\pgfqpoint{2.724168in}{2.443749in}}%
\pgfpathlineto{\pgfqpoint{2.731099in}{2.298050in}}%
\pgfpathlineto{\pgfqpoint{2.741109in}{2.035260in}}%
\pgfpathlineto{\pgfqpoint{2.771998in}{1.194106in}}%
\pgfpathlineto{\pgfqpoint{2.781324in}{1.012601in}}%
\pgfpathlineto{\pgfqpoint{2.789110in}{0.901004in}}%
\pgfpathlineto{\pgfqpoint{2.795698in}{0.835756in}}%
\pgfpathlineto{\pgfqpoint{2.801174in}{0.801336in}}%
\pgfpathlineto{\pgfqpoint{2.805538in}{0.786133in}}%
\pgfpathlineto{\pgfqpoint{2.808789in}{0.781453in}}%
\pgfpathlineto{\pgfqpoint{2.811100in}{0.781395in}}%
\pgfpathlineto{\pgfqpoint{2.813410in}{0.783921in}}%
\pgfpathlineto{\pgfqpoint{2.816490in}{0.791095in}}%
\pgfpathlineto{\pgfqpoint{2.820597in}{0.806952in}}%
\pgfpathlineto{\pgfqpoint{2.825987in}{0.837537in}}%
\pgfpathlineto{\pgfqpoint{2.833004in}{0.891387in}}%
\pgfpathlineto{\pgfqpoint{2.842330in}{0.981735in}}%
\pgfpathlineto{\pgfqpoint{2.856448in}{1.142983in}}%
\pgfpathlineto{\pgfqpoint{2.891100in}{1.545185in}}%
\pgfpathlineto{\pgfqpoint{2.905218in}{1.679841in}}%
\pgfpathlineto{\pgfqpoint{2.917539in}{1.777021in}}%
\pgfpathlineto{\pgfqpoint{2.928577in}{1.847339in}}%
\pgfpathlineto{\pgfqpoint{2.938588in}{1.897616in}}%
\pgfpathlineto{\pgfqpoint{2.947657in}{1.932433in}}%
\pgfpathlineto{\pgfqpoint{2.955700in}{1.955073in}}%
\pgfpathlineto{\pgfqpoint{2.962716in}{1.968699in}}%
\pgfpathlineto{\pgfqpoint{2.968706in}{1.975924in}}%
\pgfpathlineto{\pgfqpoint{2.973839in}{1.978942in}}%
\pgfpathlineto{\pgfqpoint{2.978460in}{1.979177in}}%
\pgfpathlineto{\pgfqpoint{2.982995in}{1.977137in}}%
\pgfpathlineto{\pgfqpoint{2.987957in}{1.972335in}}%
\pgfpathlineto{\pgfqpoint{2.993690in}{1.963443in}}%
\pgfpathlineto{\pgfqpoint{3.000278in}{1.948776in}}%
\pgfpathlineto{\pgfqpoint{3.007722in}{1.926429in}}%
\pgfpathlineto{\pgfqpoint{3.016107in}{1.893821in}}%
\pgfpathlineto{\pgfqpoint{3.025434in}{1.848161in}}%
\pgfpathlineto{\pgfqpoint{3.035787in}{1.785772in}}%
\pgfpathlineto{\pgfqpoint{3.047252in}{1.702425in}}%
\pgfpathlineto{\pgfqpoint{3.060257in}{1.590663in}}%
\pgfpathlineto{\pgfqpoint{3.075915in}{1.435517in}}%
\pgfpathlineto{\pgfqpoint{3.123317in}{0.950292in}}%
\pgfpathlineto{\pgfqpoint{3.132044in}{0.890435in}}%
\pgfpathlineto{\pgfqpoint{3.138633in}{0.858964in}}%
\pgfpathlineto{\pgfqpoint{3.143681in}{0.844397in}}%
\pgfpathlineto{\pgfqpoint{3.147274in}{0.839670in}}%
\pgfpathlineto{\pgfqpoint{3.149841in}{0.839374in}}%
\pgfpathlineto{\pgfqpoint{3.152237in}{0.841527in}}%
\pgfpathlineto{\pgfqpoint{3.155232in}{0.847647in}}%
\pgfpathlineto{\pgfqpoint{3.158997in}{0.860958in}}%
\pgfpathlineto{\pgfqpoint{3.163617in}{0.886189in}}%
\pgfpathlineto{\pgfqpoint{3.169178in}{0.930010in}}%
\pgfpathlineto{\pgfqpoint{3.175681in}{1.000245in}}%
\pgfpathlineto{\pgfqpoint{3.183296in}{1.108241in}}%
\pgfpathlineto{\pgfqpoint{3.192366in}{1.270640in}}%
\pgfpathlineto{\pgfqpoint{3.203917in}{1.520679in}}%
\pgfpathlineto{\pgfqpoint{3.239340in}{2.318520in}}%
\pgfpathlineto{\pgfqpoint{3.246698in}{2.423178in}}%
\pgfpathlineto{\pgfqpoint{3.252260in}{2.474215in}}%
\pgfpathlineto{\pgfqpoint{3.256281in}{2.493951in}}%
\pgfpathlineto{\pgfqpoint{3.258933in}{2.498529in}}%
\pgfpathlineto{\pgfqpoint{3.260474in}{2.498003in}}%
\pgfpathlineto{\pgfqpoint{3.262270in}{2.494380in}}%
\pgfpathlineto{\pgfqpoint{3.264923in}{2.483031in}}%
\pgfpathlineto{\pgfqpoint{3.268431in}{2.456942in}}%
\pgfpathlineto{\pgfqpoint{3.272966in}{2.404595in}}%
\pgfpathlineto{\pgfqpoint{3.278613in}{2.311006in}}%
\pgfpathlineto{\pgfqpoint{3.285629in}{2.154782in}}%
\pgfpathlineto{\pgfqpoint{3.295041in}{1.890538in}}%
\pgfpathlineto{\pgfqpoint{3.323704in}{1.049211in}}%
\pgfpathlineto{\pgfqpoint{3.330036in}{0.938831in}}%
\pgfpathlineto{\pgfqpoint{3.334742in}{0.890298in}}%
\pgfpathlineto{\pgfqpoint{3.337993in}{0.875326in}}%
\pgfpathlineto{\pgfqpoint{3.339704in}{0.873824in}}%
\pgfpathlineto{\pgfqpoint{3.339876in}{0.873920in}}%
\pgfpathlineto{\pgfqpoint{3.341245in}{0.876305in}}%
\pgfpathlineto{\pgfqpoint{3.343298in}{0.885308in}}%
\pgfpathlineto{\pgfqpoint{3.346293in}{0.910116in}}%
\pgfpathlineto{\pgfqpoint{3.350229in}{0.963491in}}%
\pgfpathlineto{\pgfqpoint{3.355362in}{1.066791in}}%
\pgfpathlineto{\pgfqpoint{3.362036in}{1.251020in}}%
\pgfpathlineto{\pgfqpoint{3.371705in}{1.587749in}}%
\pgfpathlineto{\pgfqpoint{3.389844in}{2.223988in}}%
\pgfpathlineto{\pgfqpoint{3.396518in}{2.380633in}}%
\pgfpathlineto{\pgfqpoint{3.401309in}{2.447865in}}%
\pgfpathlineto{\pgfqpoint{3.404646in}{2.469138in}}%
\pgfpathlineto{\pgfqpoint{3.406358in}{2.471478in}}%
\pgfpathlineto{\pgfqpoint{3.406529in}{2.471388in}}%
\pgfpathlineto{\pgfqpoint{3.407727in}{2.469091in}}%
\pgfpathlineto{\pgfqpoint{3.409609in}{2.459593in}}%
\pgfpathlineto{\pgfqpoint{3.412433in}{2.431905in}}%
\pgfpathlineto{\pgfqpoint{3.416283in}{2.368795in}}%
\pgfpathlineto{\pgfqpoint{3.421331in}{2.244882in}}%
\pgfpathlineto{\pgfqpoint{3.428090in}{2.018406in}}%
\pgfpathlineto{\pgfqpoint{3.440583in}{1.506788in}}%
\pgfpathlineto{\pgfqpoint{3.450936in}{1.125057in}}%
\pgfpathlineto{\pgfqpoint{3.457011in}{0.974149in}}%
\pgfpathlineto{\pgfqpoint{3.461289in}{0.914210in}}%
\pgfpathlineto{\pgfqpoint{3.464198in}{0.898052in}}%
\pgfpathlineto{\pgfqpoint{3.465225in}{0.897297in}}%
\pgfpathlineto{\pgfqpoint{3.465567in}{0.897624in}}%
\pgfpathlineto{\pgfqpoint{3.466850in}{0.901443in}}%
\pgfpathlineto{\pgfqpoint{3.468904in}{0.916051in}}%
\pgfpathlineto{\pgfqpoint{3.471898in}{0.955862in}}%
\pgfpathlineto{\pgfqpoint{3.475920in}{1.042334in}}%
\pgfpathlineto{\pgfqpoint{3.481396in}{1.213904in}}%
\pgfpathlineto{\pgfqpoint{3.489439in}{1.546040in}}%
\pgfpathlineto{\pgfqpoint{3.505952in}{2.238625in}}%
\pgfpathlineto{\pgfqpoint{3.511599in}{2.384134in}}%
\pgfpathlineto{\pgfqpoint{3.515535in}{2.438185in}}%
\pgfpathlineto{\pgfqpoint{3.518102in}{2.450154in}}%
\pgfpathlineto{\pgfqpoint{3.518872in}{2.450045in}}%
\pgfpathlineto{\pgfqpoint{3.519129in}{2.449626in}}%
\pgfpathlineto{\pgfqpoint{3.520498in}{2.444164in}}%
\pgfpathlineto{\pgfqpoint{3.522723in}{2.423728in}}%
\pgfpathlineto{\pgfqpoint{3.525888in}{2.370523in}}%
\pgfpathlineto{\pgfqpoint{3.530166in}{2.256633in}}%
\pgfpathlineto{\pgfqpoint{3.536070in}{2.033559in}}%
\pgfpathlineto{\pgfqpoint{3.547193in}{1.507887in}}%
\pgfpathlineto{\pgfqpoint{3.556434in}{1.120476in}}%
\pgfpathlineto{\pgfqpoint{3.561825in}{0.976332in}}%
\pgfpathlineto{\pgfqpoint{3.565589in}{0.925568in}}%
\pgfpathlineto{\pgfqpoint{3.567814in}{0.916838in}}%
\pgfpathlineto{\pgfqpoint{3.568670in}{0.917803in}}%
\pgfpathlineto{\pgfqpoint{3.570039in}{0.924361in}}%
\pgfpathlineto{\pgfqpoint{3.572263in}{0.948100in}}%
\pgfpathlineto{\pgfqpoint{3.575429in}{1.009003in}}%
\pgfpathlineto{\pgfqpoint{3.579793in}{1.140833in}}%
\pgfpathlineto{\pgfqpoint{3.586039in}{1.405778in}}%
\pgfpathlineto{\pgfqpoint{3.606231in}{2.317079in}}%
\pgfpathlineto{\pgfqpoint{3.610595in}{2.407140in}}%
\pgfpathlineto{\pgfqpoint{3.613504in}{2.430794in}}%
\pgfpathlineto{\pgfqpoint{3.614360in}{2.431928in}}%
\pgfpathlineto{\pgfqpoint{3.614873in}{2.431323in}}%
\pgfpathlineto{\pgfqpoint{3.616071in}{2.426163in}}%
\pgfpathlineto{\pgfqpoint{3.618125in}{2.405170in}}%
\pgfpathlineto{\pgfqpoint{3.621119in}{2.347794in}}%
\pgfpathlineto{\pgfqpoint{3.625312in}{2.218434in}}%
\pgfpathlineto{\pgfqpoint{3.631301in}{1.955108in}}%
\pgfpathlineto{\pgfqpoint{3.650724in}{1.043548in}}%
\pgfpathlineto{\pgfqpoint{3.654917in}{0.956131in}}%
\pgfpathlineto{\pgfqpoint{3.657655in}{0.934941in}}%
\pgfpathlineto{\pgfqpoint{3.658339in}{0.934284in}}%
\pgfpathlineto{\pgfqpoint{3.658767in}{0.934824in}}%
\pgfpathlineto{\pgfqpoint{3.659965in}{0.940231in}}%
\pgfpathlineto{\pgfqpoint{3.661933in}{0.961490in}}%
\pgfpathlineto{\pgfqpoint{3.664842in}{1.020284in}}%
\pgfpathlineto{\pgfqpoint{3.668949in}{1.154107in}}%
\pgfpathlineto{\pgfqpoint{3.674938in}{1.432016in}}%
\pgfpathlineto{\pgfqpoint{3.692478in}{2.292149in}}%
\pgfpathlineto{\pgfqpoint{3.696757in}{2.390401in}}%
\pgfpathlineto{\pgfqpoint{3.699495in}{2.414294in}}%
\pgfpathlineto{\pgfqpoint{3.700265in}{2.415243in}}%
\pgfpathlineto{\pgfqpoint{3.700692in}{2.414664in}}%
\pgfpathlineto{\pgfqpoint{3.701890in}{2.408843in}}%
\pgfpathlineto{\pgfqpoint{3.703858in}{2.385941in}}%
\pgfpathlineto{\pgfqpoint{3.706853in}{2.320335in}}%
\pgfpathlineto{\pgfqpoint{3.711045in}{2.172239in}}%
\pgfpathlineto{\pgfqpoint{3.717292in}{1.861633in}}%
\pgfpathlineto{\pgfqpoint{3.732094in}{1.103016in}}%
\pgfpathlineto{\pgfqpoint{3.736543in}{0.984355in}}%
\pgfpathlineto{\pgfqpoint{3.739452in}{0.952399in}}%
\pgfpathlineto{\pgfqpoint{3.740479in}{0.950312in}}%
\pgfpathlineto{\pgfqpoint{3.740907in}{0.950877in}}%
\pgfpathlineto{\pgfqpoint{3.742019in}{0.956293in}}%
\pgfpathlineto{\pgfqpoint{3.743987in}{0.979734in}}%
\pgfpathlineto{\pgfqpoint{3.746896in}{1.045706in}}%
\pgfpathlineto{\pgfqpoint{3.751003in}{1.195991in}}%
\pgfpathlineto{\pgfqpoint{3.757249in}{1.517825in}}%
\pgfpathlineto{\pgfqpoint{3.771025in}{2.242858in}}%
\pgfpathlineto{\pgfqpoint{3.775474in}{2.365946in}}%
\pgfpathlineto{\pgfqpoint{3.778383in}{2.398078in}}%
\pgfpathlineto{\pgfqpoint{3.779239in}{2.399760in}}%
\pgfpathlineto{\pgfqpoint{3.779752in}{2.399053in}}%
\pgfpathlineto{\pgfqpoint{3.780950in}{2.392400in}}%
\pgfpathlineto{\pgfqpoint{3.782918in}{2.366415in}}%
\pgfpathlineto{\pgfqpoint{3.785913in}{2.292412in}}%
\pgfpathlineto{\pgfqpoint{3.790191in}{2.123040in}}%
\pgfpathlineto{\pgfqpoint{3.797036in}{1.747504in}}%
\pgfpathlineto{\pgfqpoint{3.808244in}{1.143912in}}%
\pgfpathlineto{\pgfqpoint{3.812865in}{1.004397in}}%
\pgfpathlineto{\pgfqpoint{3.815860in}{0.967545in}}%
\pgfpathlineto{\pgfqpoint{3.816801in}{0.965346in}}%
\pgfpathlineto{\pgfqpoint{3.817314in}{0.966064in}}%
\pgfpathlineto{\pgfqpoint{3.818512in}{0.973005in}}%
\pgfpathlineto{\pgfqpoint{3.820480in}{1.000227in}}%
\pgfpathlineto{\pgfqpoint{3.823475in}{1.077752in}}%
\pgfpathlineto{\pgfqpoint{3.827753in}{1.254565in}}%
\pgfpathlineto{\pgfqpoint{3.834854in}{1.657936in}}%
\pgfpathlineto{\pgfqpoint{3.844951in}{2.207348in}}%
\pgfpathlineto{\pgfqpoint{3.849486in}{2.346800in}}%
\pgfpathlineto{\pgfqpoint{3.852395in}{2.383000in}}%
\pgfpathlineto{\pgfqpoint{3.853336in}{2.385105in}}%
\pgfpathlineto{\pgfqpoint{3.853764in}{2.384488in}}%
\pgfpathlineto{\pgfqpoint{3.854876in}{2.378284in}}%
\pgfpathlineto{\pgfqpoint{3.856844in}{2.351206in}}%
\pgfpathlineto{\pgfqpoint{3.859753in}{2.275113in}}%
\pgfpathlineto{\pgfqpoint{3.863946in}{2.099337in}}%
\pgfpathlineto{\pgfqpoint{3.870962in}{1.693975in}}%
\pgfpathlineto{\pgfqpoint{3.880801in}{1.151936in}}%
\pgfpathlineto{\pgfqpoint{3.885251in}{1.015231in}}%
\pgfpathlineto{\pgfqpoint{3.888074in}{0.981269in}}%
\pgfpathlineto{\pgfqpoint{3.888844in}{0.979648in}}%
\pgfpathlineto{\pgfqpoint{3.889358in}{0.980411in}}%
\pgfpathlineto{\pgfqpoint{3.890470in}{0.987113in}}%
\pgfpathlineto{\pgfqpoint{3.892438in}{1.015697in}}%
\pgfpathlineto{\pgfqpoint{3.895347in}{1.095285in}}%
\pgfpathlineto{\pgfqpoint{3.899625in}{1.282190in}}%
\pgfpathlineto{\pgfqpoint{3.907069in}{1.725389in}}%
\pgfpathlineto{\pgfqpoint{3.915796in}{2.204352in}}%
\pgfpathlineto{\pgfqpoint{3.920160in}{2.337996in}}%
\pgfpathlineto{\pgfqpoint{3.922898in}{2.369857in}}%
\pgfpathlineto{\pgfqpoint{3.923583in}{2.371112in}}%
\pgfpathlineto{\pgfqpoint{3.924096in}{2.370273in}}%
\pgfpathlineto{\pgfqpoint{3.925294in}{2.362386in}}%
\pgfpathlineto{\pgfqpoint{3.927347in}{2.329838in}}%
\pgfpathlineto{\pgfqpoint{3.930428in}{2.238415in}}%
\pgfpathlineto{\pgfqpoint{3.934877in}{2.030425in}}%
\pgfpathlineto{\pgfqpoint{3.944289in}{1.452610in}}%
\pgfpathlineto{\pgfqpoint{3.950963in}{1.118355in}}%
\pgfpathlineto{\pgfqpoint{3.954898in}{1.014101in}}%
\pgfpathlineto{\pgfqpoint{3.957294in}{0.993558in}}%
\pgfpathlineto{\pgfqpoint{3.957636in}{0.993406in}}%
\pgfpathlineto{\pgfqpoint{3.958064in}{0.994201in}}%
\pgfpathlineto{\pgfqpoint{3.959177in}{1.001377in}}%
\pgfpathlineto{\pgfqpoint{3.961145in}{1.031910in}}%
\pgfpathlineto{\pgfqpoint{3.964139in}{1.119804in}}%
\pgfpathlineto{\pgfqpoint{3.968503in}{1.323547in}}%
\pgfpathlineto{\pgfqpoint{3.977401in}{1.874709in}}%
\pgfpathlineto{\pgfqpoint{3.984246in}{2.227774in}}%
\pgfpathlineto{\pgfqpoint{3.988182in}{2.335683in}}%
\pgfpathlineto{\pgfqpoint{3.990578in}{2.357374in}}%
\pgfpathlineto{\pgfqpoint{3.991006in}{2.357561in}}%
\pgfpathlineto{\pgfqpoint{3.991348in}{2.356901in}}%
\pgfpathlineto{\pgfqpoint{3.992460in}{2.349794in}}%
\pgfpathlineto{\pgfqpoint{3.994428in}{2.318893in}}%
\pgfpathlineto{\pgfqpoint{3.997337in}{2.232583in}}%
\pgfpathlineto{\pgfqpoint{4.001701in}{2.026518in}}%
\pgfpathlineto{\pgfqpoint{4.011027in}{1.444282in}}%
\pgfpathlineto{\pgfqpoint{4.017445in}{1.122123in}}%
\pgfpathlineto{\pgfqpoint{4.021209in}{1.024503in}}%
\pgfpathlineto{\pgfqpoint{4.023520in}{1.006673in}}%
\pgfpathlineto{\pgfqpoint{4.024119in}{1.007512in}}%
\pgfpathlineto{\pgfqpoint{4.025231in}{1.015049in}}%
\pgfpathlineto{\pgfqpoint{4.027199in}{1.047137in}}%
\pgfpathlineto{\pgfqpoint{4.030193in}{1.139322in}}%
\pgfpathlineto{\pgfqpoint{4.034643in}{1.356559in}}%
\pgfpathlineto{\pgfqpoint{4.052269in}{2.304843in}}%
\pgfpathlineto{\pgfqpoint{4.055092in}{2.342925in}}%
\pgfpathlineto{\pgfqpoint{4.055862in}{2.344528in}}%
\pgfpathlineto{\pgfqpoint{4.056290in}{2.343770in}}%
\pgfpathlineto{\pgfqpoint{4.057402in}{2.336305in}}%
\pgfpathlineto{\pgfqpoint{4.059370in}{2.303943in}}%
\pgfpathlineto{\pgfqpoint{4.062365in}{2.210478in}}%
\pgfpathlineto{\pgfqpoint{4.066814in}{1.990205in}}%
\pgfpathlineto{\pgfqpoint{4.083841in}{1.063616in}}%
\pgfpathlineto{\pgfqpoint{4.086750in}{1.021555in}}%
\pgfpathlineto{\pgfqpoint{4.087520in}{1.019566in}}%
\pgfpathlineto{\pgfqpoint{4.088034in}{1.020399in}}%
\pgfpathlineto{\pgfqpoint{4.089146in}{1.028114in}}%
\pgfpathlineto{\pgfqpoint{4.091114in}{1.061260in}}%
\pgfpathlineto{\pgfqpoint{4.094109in}{1.156600in}}%
\pgfpathlineto{\pgfqpoint{4.098558in}{1.380319in}}%
\pgfpathlineto{\pgfqpoint{4.114986in}{2.283208in}}%
\pgfpathlineto{\pgfqpoint{4.117895in}{2.328873in}}%
\pgfpathlineto{\pgfqpoint{4.118836in}{2.331782in}}%
\pgfpathlineto{\pgfqpoint{4.119350in}{2.330883in}}%
\pgfpathlineto{\pgfqpoint{4.120462in}{2.322928in}}%
\pgfpathlineto{\pgfqpoint{4.122430in}{2.289059in}}%
\pgfpathlineto{\pgfqpoint{4.125424in}{2.192033in}}%
\pgfpathlineto{\pgfqpoint{4.129874in}{1.965314in}}%
\pgfpathlineto{\pgfqpoint{4.145788in}{1.084317in}}%
\pgfpathlineto{\pgfqpoint{4.148783in}{1.035306in}}%
\pgfpathlineto{\pgfqpoint{4.149810in}{1.032198in}}%
\pgfpathlineto{\pgfqpoint{4.150238in}{1.033004in}}%
\pgfpathlineto{\pgfqpoint{4.151350in}{1.040869in}}%
\pgfpathlineto{\pgfqpoint{4.153318in}{1.074855in}}%
\pgfpathlineto{\pgfqpoint{4.156313in}{1.172636in}}%
\pgfpathlineto{\pgfqpoint{4.160847in}{1.406262in}}%
\pgfpathlineto{\pgfqpoint{4.176163in}{2.261345in}}%
\pgfpathlineto{\pgfqpoint{4.179243in}{2.315383in}}%
\pgfpathlineto{\pgfqpoint{4.180356in}{2.319284in}}%
\pgfpathlineto{\pgfqpoint{4.180783in}{2.318534in}}%
\pgfpathlineto{\pgfqpoint{4.181896in}{2.310744in}}%
\pgfpathlineto{\pgfqpoint{4.183864in}{2.276657in}}%
\pgfpathlineto{\pgfqpoint{4.186858in}{2.178235in}}%
\pgfpathlineto{\pgfqpoint{4.191393in}{1.943124in}}%
\pgfpathlineto{\pgfqpoint{4.206367in}{1.104546in}}%
\pgfpathlineto{\pgfqpoint{4.209447in}{1.048908in}}%
\pgfpathlineto{\pgfqpoint{4.210559in}{1.044564in}}%
\pgfpathlineto{\pgfqpoint{4.210987in}{1.045165in}}%
\pgfpathlineto{\pgfqpoint{4.212014in}{1.051752in}}%
\pgfpathlineto{\pgfqpoint{4.213896in}{1.082398in}}%
\pgfpathlineto{\pgfqpoint{4.216720in}{1.170790in}}%
\pgfpathlineto{\pgfqpoint{4.220998in}{1.384615in}}%
\pgfpathlineto{\pgfqpoint{4.237083in}{2.265758in}}%
\pgfpathlineto{\pgfqpoint{4.239821in}{2.305186in}}%
\pgfpathlineto{\pgfqpoint{4.240592in}{2.307005in}}%
\pgfpathlineto{\pgfqpoint{4.241105in}{2.305926in}}%
\pgfpathlineto{\pgfqpoint{4.242303in}{2.296304in}}%
\pgfpathlineto{\pgfqpoint{4.244356in}{2.257160in}}%
\pgfpathlineto{\pgfqpoint{4.247437in}{2.148868in}}%
\pgfpathlineto{\pgfqpoint{4.252228in}{1.888903in}}%
\pgfpathlineto{\pgfqpoint{4.265319in}{1.140710in}}%
\pgfpathlineto{\pgfqpoint{4.268656in}{1.065885in}}%
\pgfpathlineto{\pgfqpoint{4.270282in}{1.056751in}}%
\pgfpathlineto{\pgfqpoint{4.270538in}{1.056997in}}%
\pgfpathlineto{\pgfqpoint{4.271394in}{1.061143in}}%
\pgfpathlineto{\pgfqpoint{4.273020in}{1.082974in}}%
\pgfpathlineto{\pgfqpoint{4.275587in}{1.153014in}}%
\pgfpathlineto{\pgfqpoint{4.279437in}{1.328340in}}%
\pgfpathlineto{\pgfqpoint{4.286795in}{1.791177in}}%
\pgfpathlineto{\pgfqpoint{4.293811in}{2.173885in}}%
\pgfpathlineto{\pgfqpoint{4.297576in}{2.277930in}}%
\pgfpathlineto{\pgfqpoint{4.299801in}{2.294904in}}%
\pgfpathlineto{\pgfqpoint{4.300400in}{2.293548in}}%
\pgfpathlineto{\pgfqpoint{4.301598in}{2.283303in}}%
\pgfpathlineto{\pgfqpoint{4.303651in}{2.242911in}}%
\pgfpathlineto{\pgfqpoint{4.306817in}{2.128753in}}%
\pgfpathlineto{\pgfqpoint{4.311694in}{1.859635in}}%
\pgfpathlineto{\pgfqpoint{4.323844in}{1.163201in}}%
\pgfpathlineto{\pgfqpoint{4.327352in}{1.079372in}}%
\pgfpathlineto{\pgfqpoint{4.329063in}{1.068787in}}%
\pgfpathlineto{\pgfqpoint{4.329320in}{1.068976in}}%
\pgfpathlineto{\pgfqpoint{4.330175in}{1.072958in}}%
\pgfpathlineto{\pgfqpoint{4.331801in}{1.094575in}}%
\pgfpathlineto{\pgfqpoint{4.334368in}{1.164500in}}%
\pgfpathlineto{\pgfqpoint{4.338218in}{1.339842in}}%
\pgfpathlineto{\pgfqpoint{4.345833in}{1.816902in}}%
\pgfpathlineto{\pgfqpoint{4.352507in}{2.171155in}}%
\pgfpathlineto{\pgfqpoint{4.356186in}{2.268337in}}%
\pgfpathlineto{\pgfqpoint{4.358240in}{2.282933in}}%
\pgfpathlineto{\pgfqpoint{4.358325in}{2.282897in}}%
\pgfpathlineto{\pgfqpoint{4.359010in}{2.280753in}}%
\pgfpathlineto{\pgfqpoint{4.360379in}{2.266612in}}%
\pgfpathlineto{\pgfqpoint{4.362604in}{2.216484in}}%
\pgfpathlineto{\pgfqpoint{4.366026in}{2.080559in}}%
\pgfpathlineto{\pgfqpoint{4.371673in}{1.749764in}}%
\pgfpathlineto{\pgfqpoint{4.381171in}{1.204995in}}%
\pgfpathlineto{\pgfqpoint{4.384935in}{1.098898in}}%
\pgfpathlineto{\pgfqpoint{4.387160in}{1.080712in}}%
\pgfpathlineto{\pgfqpoint{4.387759in}{1.081759in}}%
\pgfpathlineto{\pgfqpoint{4.388957in}{1.091416in}}%
\pgfpathlineto{\pgfqpoint{4.391010in}{1.130891in}}%
\pgfpathlineto{\pgfqpoint{4.394091in}{1.239974in}}%
\pgfpathlineto{\pgfqpoint{4.398968in}{1.505169in}}%
\pgfpathlineto{\pgfqpoint{4.410946in}{2.180842in}}%
\pgfpathlineto{\pgfqpoint{4.414369in}{2.260868in}}%
\pgfpathlineto{\pgfqpoint{4.416080in}{2.271064in}}%
\pgfpathlineto{\pgfqpoint{4.416337in}{2.270819in}}%
\pgfpathlineto{\pgfqpoint{4.417192in}{2.266653in}}%
\pgfpathlineto{\pgfqpoint{4.418818in}{2.244718in}}%
\pgfpathlineto{\pgfqpoint{4.421385in}{2.174463in}}%
\pgfpathlineto{\pgfqpoint{4.425235in}{1.999431in}}%
\pgfpathlineto{\pgfqpoint{4.433364in}{1.497085in}}%
\pgfpathlineto{\pgfqpoint{4.439524in}{1.188592in}}%
\pgfpathlineto{\pgfqpoint{4.443032in}{1.103812in}}%
\pgfpathlineto{\pgfqpoint{4.444829in}{1.092532in}}%
\pgfpathlineto{\pgfqpoint{4.445000in}{1.092636in}}%
\pgfpathlineto{\pgfqpoint{4.445770in}{1.095648in}}%
\pgfpathlineto{\pgfqpoint{4.447310in}{1.114032in}}%
\pgfpathlineto{\pgfqpoint{4.449792in}{1.176878in}}%
\pgfpathlineto{\pgfqpoint{4.453471in}{1.335527in}}%
\pgfpathlineto{\pgfqpoint{4.460316in}{1.749912in}}%
\pgfpathlineto{\pgfqpoint{4.467589in}{2.141417in}}%
\pgfpathlineto{\pgfqpoint{4.471353in}{2.243284in}}%
\pgfpathlineto{\pgfqpoint{4.473493in}{2.259268in}}%
\pgfpathlineto{\pgfqpoint{4.474091in}{2.258034in}}%
\pgfpathlineto{\pgfqpoint{4.475289in}{2.248084in}}%
\pgfpathlineto{\pgfqpoint{4.477343in}{2.208389in}}%
\pgfpathlineto{\pgfqpoint{4.480509in}{2.095999in}}%
\pgfpathlineto{\pgfqpoint{4.485471in}{1.827224in}}%
\pgfpathlineto{\pgfqpoint{4.496766in}{1.200413in}}%
\pgfpathlineto{\pgfqpoint{4.500274in}{1.115857in}}%
\pgfpathlineto{\pgfqpoint{4.502070in}{1.104306in}}%
\pgfpathlineto{\pgfqpoint{4.502242in}{1.104371in}}%
\pgfpathlineto{\pgfqpoint{4.502926in}{1.106666in}}%
\pgfpathlineto{\pgfqpoint{4.504295in}{1.120933in}}%
\pgfpathlineto{\pgfqpoint{4.506605in}{1.173265in}}%
\pgfpathlineto{\pgfqpoint{4.510113in}{1.312847in}}%
\pgfpathlineto{\pgfqpoint{4.516103in}{1.659297in}}%
\pgfpathlineto{\pgfqpoint{4.524659in}{2.127828in}}%
\pgfpathlineto{\pgfqpoint{4.528424in}{2.230399in}}%
\pgfpathlineto{\pgfqpoint{4.530648in}{2.247516in}}%
\pgfpathlineto{\pgfqpoint{4.531247in}{2.246320in}}%
\pgfpathlineto{\pgfqpoint{4.532445in}{2.236550in}}%
\pgfpathlineto{\pgfqpoint{4.534499in}{2.197486in}}%
\pgfpathlineto{\pgfqpoint{4.537664in}{2.086836in}}%
\pgfpathlineto{\pgfqpoint{4.542627in}{1.822398in}}%
\pgfpathlineto{\pgfqpoint{4.553836in}{1.211139in}}%
\pgfpathlineto{\pgfqpoint{4.557344in}{1.127611in}}%
\pgfpathlineto{\pgfqpoint{4.559226in}{1.116046in}}%
\pgfpathlineto{\pgfqpoint{4.559312in}{1.116094in}}%
\pgfpathlineto{\pgfqpoint{4.559996in}{1.118273in}}%
\pgfpathlineto{\pgfqpoint{4.561365in}{1.132147in}}%
\pgfpathlineto{\pgfqpoint{4.563675in}{1.183345in}}%
\pgfpathlineto{\pgfqpoint{4.567098in}{1.316126in}}%
\pgfpathlineto{\pgfqpoint{4.572916in}{1.644202in}}%
\pgfpathlineto{\pgfqpoint{4.581729in}{2.119127in}}%
\pgfpathlineto{\pgfqpoint{4.585494in}{2.219287in}}%
\pgfpathlineto{\pgfqpoint{4.587718in}{2.235783in}}%
\pgfpathlineto{\pgfqpoint{4.588317in}{2.234531in}}%
\pgfpathlineto{\pgfqpoint{4.589515in}{2.224797in}}%
\pgfpathlineto{\pgfqpoint{4.591569in}{2.186245in}}%
\pgfpathlineto{\pgfqpoint{4.594735in}{2.077459in}}%
\pgfpathlineto{\pgfqpoint{4.599783in}{1.813008in}}%
\pgfpathlineto{\pgfqpoint{4.610820in}{1.223081in}}%
\pgfpathlineto{\pgfqpoint{4.614414in}{1.138869in}}%
\pgfpathlineto{\pgfqpoint{4.616211in}{1.127777in}}%
\pgfpathlineto{\pgfqpoint{4.616382in}{1.127842in}}%
\pgfpathlineto{\pgfqpoint{4.617152in}{1.130550in}}%
\pgfpathlineto{\pgfqpoint{4.618606in}{1.146357in}}%
\pgfpathlineto{\pgfqpoint{4.621002in}{1.201672in}}%
\pgfpathlineto{\pgfqpoint{4.624596in}{1.344377in}}%
\pgfpathlineto{\pgfqpoint{4.631013in}{1.706996in}}%
\pgfpathlineto{\pgfqpoint{4.638799in}{2.109482in}}%
\pgfpathlineto{\pgfqpoint{4.642564in}{2.207652in}}%
\pgfpathlineto{\pgfqpoint{4.644789in}{2.224043in}}%
\pgfpathlineto{\pgfqpoint{4.645388in}{2.222906in}}%
\pgfpathlineto{\pgfqpoint{4.646585in}{2.213581in}}%
\pgfpathlineto{\pgfqpoint{4.648639in}{2.176272in}}%
\pgfpathlineto{\pgfqpoint{4.651805in}{2.070586in}}%
\pgfpathlineto{\pgfqpoint{4.656767in}{1.817988in}}%
\pgfpathlineto{\pgfqpoint{4.668062in}{1.229810in}}%
\pgfpathlineto{\pgfqpoint{4.671570in}{1.150462in}}%
\pgfpathlineto{\pgfqpoint{4.673452in}{1.139534in}}%
\pgfpathlineto{\pgfqpoint{4.673538in}{1.139582in}}%
\pgfpathlineto{\pgfqpoint{4.674222in}{1.141676in}}%
\pgfpathlineto{\pgfqpoint{4.675591in}{1.154909in}}%
\pgfpathlineto{\pgfqpoint{4.677901in}{1.203632in}}%
\pgfpathlineto{\pgfqpoint{4.681409in}{1.333766in}}%
\pgfpathlineto{\pgfqpoint{4.687313in}{1.652126in}}%
\pgfpathlineto{\pgfqpoint{4.696040in}{2.099642in}}%
\pgfpathlineto{\pgfqpoint{4.699805in}{2.195775in}}%
\pgfpathlineto{\pgfqpoint{4.702030in}{2.212262in}}%
\pgfpathlineto{\pgfqpoint{4.702629in}{2.211322in}}%
\pgfpathlineto{\pgfqpoint{4.703827in}{2.202605in}}%
\pgfpathlineto{\pgfqpoint{4.705880in}{2.166973in}}%
\pgfpathlineto{\pgfqpoint{4.708960in}{2.068639in}}%
\pgfpathlineto{\pgfqpoint{4.713837in}{1.830296in}}%
\pgfpathlineto{\pgfqpoint{4.725645in}{1.234063in}}%
\pgfpathlineto{\pgfqpoint{4.729153in}{1.160294in}}%
\pgfpathlineto{\pgfqpoint{4.730864in}{1.151327in}}%
\pgfpathlineto{\pgfqpoint{4.731036in}{1.151439in}}%
\pgfpathlineto{\pgfqpoint{4.731806in}{1.154210in}}%
\pgfpathlineto{\pgfqpoint{4.733346in}{1.170779in}}%
\pgfpathlineto{\pgfqpoint{4.735827in}{1.227082in}}%
\pgfpathlineto{\pgfqpoint{4.739506in}{1.368820in}}%
\pgfpathlineto{\pgfqpoint{4.746351in}{1.738718in}}%
\pgfpathlineto{\pgfqpoint{4.753710in}{2.093099in}}%
\pgfpathlineto{\pgfqpoint{4.757474in}{2.185043in}}%
\pgfpathlineto{\pgfqpoint{4.759699in}{2.200438in}}%
\pgfpathlineto{\pgfqpoint{4.760298in}{2.199399in}}%
\pgfpathlineto{\pgfqpoint{4.761496in}{2.190739in}}%
\pgfpathlineto{\pgfqpoint{4.763549in}{2.155979in}}%
\pgfpathlineto{\pgfqpoint{4.766715in}{2.057329in}}%
\pgfpathlineto{\pgfqpoint{4.771678in}{1.820708in}}%
\pgfpathlineto{\pgfqpoint{4.783314in}{1.248770in}}%
\pgfpathlineto{\pgfqpoint{4.786822in}{1.173961in}}%
\pgfpathlineto{\pgfqpoint{4.788705in}{1.163200in}}%
\pgfpathlineto{\pgfqpoint{4.788876in}{1.163277in}}%
\pgfpathlineto{\pgfqpoint{4.789646in}{1.165804in}}%
\pgfpathlineto{\pgfqpoint{4.791100in}{1.180216in}}%
\pgfpathlineto{\pgfqpoint{4.793496in}{1.230363in}}%
\pgfpathlineto{\pgfqpoint{4.797090in}{1.359636in}}%
\pgfpathlineto{\pgfqpoint{4.803336in}{1.680591in}}%
\pgfpathlineto{\pgfqpoint{4.811635in}{2.078188in}}%
\pgfpathlineto{\pgfqpoint{4.815486in}{2.171929in}}%
\pgfpathlineto{\pgfqpoint{4.817796in}{2.188504in}}%
\pgfpathlineto{\pgfqpoint{4.818395in}{2.187691in}}%
\pgfpathlineto{\pgfqpoint{4.819593in}{2.179763in}}%
\pgfpathlineto{\pgfqpoint{4.821646in}{2.147084in}}%
\pgfpathlineto{\pgfqpoint{4.824726in}{2.056533in}}%
\pgfpathlineto{\pgfqpoint{4.829518in}{1.840215in}}%
\pgfpathlineto{\pgfqpoint{4.842010in}{1.250435in}}%
\pgfpathlineto{\pgfqpoint{4.845433in}{1.184011in}}%
\pgfpathlineto{\pgfqpoint{4.847229in}{1.175181in}}%
\pgfpathlineto{\pgfqpoint{4.847400in}{1.175305in}}%
\pgfpathlineto{\pgfqpoint{4.848256in}{1.178445in}}%
\pgfpathlineto{\pgfqpoint{4.849882in}{1.195840in}}%
\pgfpathlineto{\pgfqpoint{4.852449in}{1.252377in}}%
\pgfpathlineto{\pgfqpoint{4.856299in}{1.394376in}}%
\pgfpathlineto{\pgfqpoint{4.863743in}{1.773234in}}%
\pgfpathlineto{\pgfqpoint{4.870673in}{2.077670in}}%
\pgfpathlineto{\pgfqpoint{4.874438in}{2.162092in}}%
\pgfpathlineto{\pgfqpoint{4.876663in}{2.176471in}}%
\pgfpathlineto{\pgfqpoint{4.877262in}{2.175626in}}%
\pgfpathlineto{\pgfqpoint{4.878460in}{2.167947in}}%
\pgfpathlineto{\pgfqpoint{4.880513in}{2.136644in}}%
\pgfpathlineto{\pgfqpoint{4.883593in}{2.050173in}}%
\pgfpathlineto{\pgfqpoint{4.888385in}{1.843392in}}%
\pgfpathlineto{\pgfqpoint{4.901305in}{1.257895in}}%
\pgfpathlineto{\pgfqpoint{4.904727in}{1.195401in}}%
\pgfpathlineto{\pgfqpoint{4.906439in}{1.187290in}}%
\pgfpathlineto{\pgfqpoint{4.906695in}{1.187444in}}%
\pgfpathlineto{\pgfqpoint{4.907551in}{1.190543in}}%
\pgfpathlineto{\pgfqpoint{4.909177in}{1.207248in}}%
\pgfpathlineto{\pgfqpoint{4.911743in}{1.261155in}}%
\pgfpathlineto{\pgfqpoint{4.915594in}{1.396366in}}%
\pgfpathlineto{\pgfqpoint{4.922781in}{1.746173in}}%
\pgfpathlineto{\pgfqpoint{4.930139in}{2.062465in}}%
\pgfpathlineto{\pgfqpoint{4.933990in}{2.148327in}}%
\pgfpathlineto{\pgfqpoint{4.936300in}{2.164220in}}%
\pgfpathlineto{\pgfqpoint{4.936813in}{2.163942in}}%
\pgfpathlineto{\pgfqpoint{4.936984in}{2.163541in}}%
\pgfpathlineto{\pgfqpoint{4.938182in}{2.156433in}}%
\pgfpathlineto{\pgfqpoint{4.940236in}{2.127123in}}%
\pgfpathlineto{\pgfqpoint{4.943316in}{2.045741in}}%
\pgfpathlineto{\pgfqpoint{4.948107in}{1.849986in}}%
\pgfpathlineto{\pgfqpoint{4.961626in}{1.264873in}}%
\pgfpathlineto{\pgfqpoint{4.965049in}{1.206816in}}%
\pgfpathlineto{\pgfqpoint{4.966760in}{1.199576in}}%
\pgfpathlineto{\pgfqpoint{4.966931in}{1.199674in}}%
\pgfpathlineto{\pgfqpoint{4.967787in}{1.202401in}}%
\pgfpathlineto{\pgfqpoint{4.969413in}{1.217752in}}%
\pgfpathlineto{\pgfqpoint{4.971979in}{1.267924in}}%
\pgfpathlineto{\pgfqpoint{4.975830in}{1.394729in}}%
\pgfpathlineto{\pgfqpoint{4.982675in}{1.710326in}}%
\pgfpathlineto{\pgfqpoint{4.990632in}{2.045080in}}%
\pgfpathlineto{\pgfqpoint{4.994653in}{2.134326in}}%
\pgfpathlineto{\pgfqpoint{4.997049in}{2.151703in}}%
\pgfpathlineto{\pgfqpoint{4.997477in}{2.151840in}}%
\pgfpathlineto{\pgfqpoint{4.997905in}{2.151075in}}%
\pgfpathlineto{\pgfqpoint{4.999103in}{2.144158in}}%
\pgfpathlineto{\pgfqpoint{5.001156in}{2.116296in}}%
\pgfpathlineto{\pgfqpoint{5.004236in}{2.039476in}}%
\pgfpathlineto{\pgfqpoint{5.008942in}{1.858468in}}%
\pgfpathlineto{\pgfqpoint{5.023317in}{1.269425in}}%
\pgfpathlineto{\pgfqpoint{5.026654in}{1.218296in}}%
\pgfpathlineto{\pgfqpoint{5.028280in}{1.212077in}}%
\pgfpathlineto{\pgfqpoint{5.028536in}{1.212242in}}%
\pgfpathlineto{\pgfqpoint{5.029477in}{1.215521in}}%
\pgfpathlineto{\pgfqpoint{5.031189in}{1.232104in}}%
\pgfpathlineto{\pgfqpoint{5.033841in}{1.283588in}}%
\pgfpathlineto{\pgfqpoint{5.037777in}{1.409575in}}%
\pgfpathlineto{\pgfqpoint{5.045050in}{1.730702in}}%
\pgfpathlineto{\pgfqpoint{5.052750in}{2.035560in}}%
\pgfpathlineto{\pgfqpoint{5.056772in}{2.121017in}}%
\pgfpathlineto{\pgfqpoint{5.059253in}{2.138983in}}%
\pgfpathlineto{\pgfqpoint{5.059767in}{2.139207in}}%
\pgfpathlineto{\pgfqpoint{5.060194in}{2.138474in}}%
\pgfpathlineto{\pgfqpoint{5.061392in}{2.132001in}}%
\pgfpathlineto{\pgfqpoint{5.063446in}{2.106080in}}%
\pgfpathlineto{\pgfqpoint{5.066526in}{2.034653in}}%
\pgfpathlineto{\pgfqpoint{5.071232in}{1.865682in}}%
\pgfpathlineto{\pgfqpoint{5.086462in}{1.276731in}}%
\pgfpathlineto{\pgfqpoint{5.089799in}{1.230240in}}%
\pgfpathlineto{\pgfqpoint{5.091425in}{1.224853in}}%
\pgfpathlineto{\pgfqpoint{5.091681in}{1.225059in}}%
\pgfpathlineto{\pgfqpoint{5.092622in}{1.228274in}}%
\pgfpathlineto{\pgfqpoint{5.094334in}{1.243896in}}%
\pgfpathlineto{\pgfqpoint{5.096986in}{1.291866in}}%
\pgfpathlineto{\pgfqpoint{5.100922in}{1.409035in}}%
\pgfpathlineto{\pgfqpoint{5.107853in}{1.695394in}}%
\pgfpathlineto{\pgfqpoint{5.116323in}{2.019147in}}%
\pgfpathlineto{\pgfqpoint{5.120516in}{2.106378in}}%
\pgfpathlineto{\pgfqpoint{5.123168in}{2.125926in}}%
\pgfpathlineto{\pgfqpoint{5.123767in}{2.126291in}}%
\pgfpathlineto{\pgfqpoint{5.124195in}{2.125635in}}%
\pgfpathlineto{\pgfqpoint{5.125393in}{2.119749in}}%
\pgfpathlineto{\pgfqpoint{5.127446in}{2.096076in}}%
\pgfpathlineto{\pgfqpoint{5.130527in}{2.030625in}}%
\pgfpathlineto{\pgfqpoint{5.135147in}{1.878060in}}%
\pgfpathlineto{\pgfqpoint{5.151832in}{1.278838in}}%
\pgfpathlineto{\pgfqpoint{5.154998in}{1.241812in}}%
\pgfpathlineto{\pgfqpoint{5.156452in}{1.237952in}}%
\pgfpathlineto{\pgfqpoint{5.156709in}{1.238145in}}%
\pgfpathlineto{\pgfqpoint{5.157736in}{1.241530in}}%
\pgfpathlineto{\pgfqpoint{5.159532in}{1.257378in}}%
\pgfpathlineto{\pgfqpoint{5.162270in}{1.304623in}}%
\pgfpathlineto{\pgfqpoint{5.166292in}{1.417596in}}%
\pgfpathlineto{\pgfqpoint{5.173393in}{1.692933in}}%
\pgfpathlineto{\pgfqpoint{5.182121in}{2.006664in}}%
\pgfpathlineto{\pgfqpoint{5.186399in}{2.091781in}}%
\pgfpathlineto{\pgfqpoint{5.189137in}{2.112304in}}%
\pgfpathlineto{\pgfqpoint{5.189907in}{2.113010in}}%
\pgfpathlineto{\pgfqpoint{5.190335in}{2.112434in}}%
\pgfpathlineto{\pgfqpoint{5.191533in}{2.107169in}}%
\pgfpathlineto{\pgfqpoint{5.193586in}{2.085863in}}%
\pgfpathlineto{\pgfqpoint{5.196666in}{2.026717in}}%
\pgfpathlineto{\pgfqpoint{5.201201in}{1.890826in}}%
\pgfpathlineto{\pgfqpoint{5.219683in}{1.281464in}}%
\pgfpathlineto{\pgfqpoint{5.222763in}{1.253379in}}%
\pgfpathlineto{\pgfqpoint{5.223875in}{1.251463in}}%
\pgfpathlineto{\pgfqpoint{5.224303in}{1.251901in}}%
\pgfpathlineto{\pgfqpoint{5.225501in}{1.256580in}}%
\pgfpathlineto{\pgfqpoint{5.227469in}{1.275117in}}%
\pgfpathlineto{\pgfqpoint{5.230464in}{1.327679in}}%
\pgfpathlineto{\pgfqpoint{5.234913in}{1.451294in}}%
\pgfpathlineto{\pgfqpoint{5.243726in}{1.776282in}}%
\pgfpathlineto{\pgfqpoint{5.250999in}{2.005788in}}%
\pgfpathlineto{\pgfqpoint{5.255277in}{2.081023in}}%
\pgfpathlineto{\pgfqpoint{5.258015in}{2.098751in}}%
\pgfpathlineto{\pgfqpoint{5.258785in}{2.099231in}}%
\pgfpathlineto{\pgfqpoint{5.259213in}{2.098639in}}%
\pgfpathlineto{\pgfqpoint{5.260496in}{2.093208in}}%
\pgfpathlineto{\pgfqpoint{5.262635in}{2.072224in}}%
\pgfpathlineto{\pgfqpoint{5.265801in}{2.015619in}}%
\pgfpathlineto{\pgfqpoint{5.270507in}{1.885173in}}%
\pgfpathlineto{\pgfqpoint{5.289844in}{1.295114in}}%
\pgfpathlineto{\pgfqpoint{5.293010in}{1.267732in}}%
\pgfpathlineto{\pgfqpoint{5.294293in}{1.265487in}}%
\pgfpathlineto{\pgfqpoint{5.294635in}{1.265760in}}%
\pgfpathlineto{\pgfqpoint{5.295748in}{1.269173in}}%
\pgfpathlineto{\pgfqpoint{5.297630in}{1.283613in}}%
\pgfpathlineto{\pgfqpoint{5.300539in}{1.326291in}}%
\pgfpathlineto{\pgfqpoint{5.304732in}{1.425486in}}%
\pgfpathlineto{\pgfqpoint{5.311662in}{1.651621in}}%
\pgfpathlineto{\pgfqpoint{5.322015in}{1.977007in}}%
\pgfpathlineto{\pgfqpoint{5.326636in}{2.060037in}}%
\pgfpathlineto{\pgfqpoint{5.329716in}{2.083303in}}%
\pgfpathlineto{\pgfqpoint{5.330914in}{2.084944in}}%
\pgfpathlineto{\pgfqpoint{5.331256in}{2.084645in}}%
\pgfpathlineto{\pgfqpoint{5.332454in}{2.080925in}}%
\pgfpathlineto{\pgfqpoint{5.334422in}{2.065932in}}%
\pgfpathlineto{\pgfqpoint{5.337417in}{2.023033in}}%
\pgfpathlineto{\pgfqpoint{5.341695in}{1.925454in}}%
\pgfpathlineto{\pgfqpoint{5.348882in}{1.700976in}}%
\pgfpathlineto{\pgfqpoint{5.359235in}{1.390024in}}%
\pgfpathlineto{\pgfqpoint{5.364027in}{1.306267in}}%
\pgfpathlineto{\pgfqpoint{5.367192in}{1.282271in}}%
\pgfpathlineto{\pgfqpoint{5.368561in}{1.280173in}}%
\pgfpathlineto{\pgfqpoint{5.368904in}{1.280435in}}%
\pgfpathlineto{\pgfqpoint{5.370102in}{1.283823in}}%
\pgfpathlineto{\pgfqpoint{5.372069in}{1.297592in}}%
\pgfpathlineto{\pgfqpoint{5.375064in}{1.337144in}}%
\pgfpathlineto{\pgfqpoint{5.379342in}{1.427548in}}%
\pgfpathlineto{\pgfqpoint{5.386187in}{1.627147in}}%
\pgfpathlineto{\pgfqpoint{5.397909in}{1.963882in}}%
\pgfpathlineto{\pgfqpoint{5.402701in}{2.043350in}}%
\pgfpathlineto{\pgfqpoint{5.405952in}{2.067456in}}%
\pgfpathlineto{\pgfqpoint{5.407492in}{2.069866in}}%
\pgfpathlineto{\pgfqpoint{5.407749in}{2.069700in}}%
\pgfpathlineto{\pgfqpoint{5.408861in}{2.067115in}}%
\pgfpathlineto{\pgfqpoint{5.410744in}{2.055925in}}%
\pgfpathlineto{\pgfqpoint{5.413567in}{2.023751in}}%
\pgfpathlineto{\pgfqpoint{5.417674in}{1.947606in}}%
\pgfpathlineto{\pgfqpoint{5.423920in}{1.783549in}}%
\pgfpathlineto{\pgfqpoint{5.438637in}{1.385157in}}%
\pgfpathlineto{\pgfqpoint{5.443343in}{1.316999in}}%
\pgfpathlineto{\pgfqpoint{5.446509in}{1.297337in}}%
\pgfpathlineto{\pgfqpoint{5.447878in}{1.295725in}}%
\pgfpathlineto{\pgfqpoint{5.448220in}{1.295976in}}%
\pgfpathlineto{\pgfqpoint{5.449503in}{1.299230in}}%
\pgfpathlineto{\pgfqpoint{5.451557in}{1.311915in}}%
\pgfpathlineto{\pgfqpoint{5.454637in}{1.347361in}}%
\pgfpathlineto{\pgfqpoint{5.459001in}{1.427172in}}%
\pgfpathlineto{\pgfqpoint{5.465760in}{1.598379in}}%
\pgfpathlineto{\pgfqpoint{5.479621in}{1.953708in}}%
\pgfpathlineto{\pgfqpoint{5.484584in}{2.026829in}}%
\pgfpathlineto{\pgfqpoint{5.488007in}{2.050616in}}%
\pgfpathlineto{\pgfqpoint{5.489975in}{2.053771in}}%
\pgfpathlineto{\pgfqpoint{5.491087in}{2.052133in}}%
\pgfpathlineto{\pgfqpoint{5.492884in}{2.044334in}}%
\pgfpathlineto{\pgfqpoint{5.495622in}{2.020628in}}%
\pgfpathlineto{\pgfqpoint{5.499472in}{1.965208in}}%
\pgfpathlineto{\pgfqpoint{5.505034in}{1.848847in}}%
\pgfpathlineto{\pgfqpoint{5.526681in}{1.358013in}}%
\pgfpathlineto{\pgfqpoint{5.530788in}{1.321008in}}%
\pgfpathlineto{\pgfqpoint{5.533440in}{1.312616in}}%
\pgfpathlineto{\pgfqpoint{5.534638in}{1.312898in}}%
\pgfpathlineto{\pgfqpoint{5.534724in}{1.313014in}}%
\pgfpathlineto{\pgfqpoint{5.536264in}{1.317302in}}%
\pgfpathlineto{\pgfqpoint{5.538660in}{1.332036in}}%
\pgfpathlineto{\pgfqpoint{5.542082in}{1.369119in}}%
\pgfpathlineto{\pgfqpoint{5.546874in}{1.448385in}}%
\pgfpathlineto{\pgfqpoint{5.554318in}{1.614019in}}%
\pgfpathlineto{\pgfqpoint{5.568435in}{1.928450in}}%
\pgfpathlineto{\pgfqpoint{5.573911in}{2.003162in}}%
\pgfpathlineto{\pgfqpoint{5.577762in}{2.030510in}}%
\pgfpathlineto{\pgfqpoint{5.580243in}{2.036244in}}%
\pgfpathlineto{\pgfqpoint{5.581526in}{2.035513in}}%
\pgfpathlineto{\pgfqpoint{5.583238in}{2.030660in}}%
\pgfpathlineto{\pgfqpoint{5.585804in}{2.015290in}}%
\pgfpathlineto{\pgfqpoint{5.589398in}{1.978490in}}%
\pgfpathlineto{\pgfqpoint{5.594446in}{1.900972in}}%
\pgfpathlineto{\pgfqpoint{5.602318in}{1.740116in}}%
\pgfpathlineto{\pgfqpoint{5.616864in}{1.442753in}}%
\pgfpathlineto{\pgfqpoint{5.622682in}{1.367615in}}%
\pgfpathlineto{\pgfqpoint{5.626789in}{1.338594in}}%
\pgfpathlineto{\pgfqpoint{5.629612in}{1.331145in}}%
\pgfpathlineto{\pgfqpoint{5.631067in}{1.331336in}}%
\pgfpathlineto{\pgfqpoint{5.632693in}{1.334760in}}%
\pgfpathlineto{\pgfqpoint{5.635174in}{1.346366in}}%
\pgfpathlineto{\pgfqpoint{5.638682in}{1.375234in}}%
\pgfpathlineto{\pgfqpoint{5.643474in}{1.435494in}}%
\pgfpathlineto{\pgfqpoint{5.650404in}{1.554331in}}%
\pgfpathlineto{\pgfqpoint{5.671624in}{1.936907in}}%
\pgfpathlineto{\pgfqpoint{5.677100in}{1.991354in}}%
\pgfpathlineto{\pgfqpoint{5.681036in}{2.011905in}}%
\pgfpathlineto{\pgfqpoint{5.683688in}{2.016545in}}%
\pgfpathlineto{\pgfqpoint{5.685142in}{2.015932in}}%
\pgfpathlineto{\pgfqpoint{5.687025in}{2.011867in}}%
\pgfpathlineto{\pgfqpoint{5.689763in}{1.999543in}}%
\pgfpathlineto{\pgfqpoint{5.693528in}{1.970928in}}%
\pgfpathlineto{\pgfqpoint{5.698661in}{1.912821in}}%
\pgfpathlineto{\pgfqpoint{5.706020in}{1.801139in}}%
\pgfpathlineto{\pgfqpoint{5.728865in}{1.436724in}}%
\pgfpathlineto{\pgfqpoint{5.734769in}{1.382282in}}%
\pgfpathlineto{\pgfqpoint{5.739132in}{1.359460in}}%
\pgfpathlineto{\pgfqpoint{5.742213in}{1.352742in}}%
\pgfpathlineto{\pgfqpoint{5.744095in}{1.352478in}}%
\pgfpathlineto{\pgfqpoint{5.745892in}{1.354907in}}%
\pgfpathlineto{\pgfqpoint{5.748459in}{1.362810in}}%
\pgfpathlineto{\pgfqpoint{5.752052in}{1.382223in}}%
\pgfpathlineto{\pgfqpoint{5.756844in}{1.421837in}}%
\pgfpathlineto{\pgfqpoint{5.763347in}{1.496074in}}%
\pgfpathlineto{\pgfqpoint{5.773871in}{1.645923in}}%
\pgfpathlineto{\pgfqpoint{5.789272in}{1.860192in}}%
\pgfpathlineto{\pgfqpoint{5.796802in}{1.935415in}}%
\pgfpathlineto{\pgfqpoint{5.802534in}{1.972745in}}%
\pgfpathlineto{\pgfqpoint{5.806898in}{1.988414in}}%
\pgfpathlineto{\pgfqpoint{5.809978in}{1.992701in}}%
\pgfpathlineto{\pgfqpoint{5.812032in}{1.992473in}}%
\pgfpathlineto{\pgfqpoint{5.814256in}{1.989499in}}%
\pgfpathlineto{\pgfqpoint{5.817251in}{1.981159in}}%
\pgfpathlineto{\pgfqpoint{5.821272in}{1.962578in}}%
\pgfpathlineto{\pgfqpoint{5.826577in}{1.926453in}}%
\pgfpathlineto{\pgfqpoint{5.833679in}{1.861320in}}%
\pgfpathlineto{\pgfqpoint{5.844802in}{1.735663in}}%
\pgfpathlineto{\pgfqpoint{5.863882in}{1.521212in}}%
\pgfpathlineto{\pgfqpoint{5.872524in}{1.449031in}}%
\pgfpathlineto{\pgfqpoint{5.879284in}{1.409532in}}%
\pgfpathlineto{\pgfqpoint{5.884674in}{1.389742in}}%
\pgfpathlineto{\pgfqpoint{5.888867in}{1.381659in}}%
\pgfpathlineto{\pgfqpoint{5.891947in}{1.379737in}}%
\pgfpathlineto{\pgfqpoint{5.894514in}{1.380662in}}%
\pgfpathlineto{\pgfqpoint{5.897423in}{1.384393in}}%
\pgfpathlineto{\pgfqpoint{5.901188in}{1.393244in}}%
\pgfpathlineto{\pgfqpoint{5.906065in}{1.410940in}}%
\pgfpathlineto{\pgfqpoint{5.912396in}{1.443087in}}%
\pgfpathlineto{\pgfqpoint{5.920782in}{1.498238in}}%
\pgfpathlineto{\pgfqpoint{5.933701in}{1.600055in}}%
\pgfpathlineto{\pgfqpoint{5.957830in}{1.790448in}}%
\pgfpathlineto{\pgfqpoint{5.968782in}{1.858608in}}%
\pgfpathlineto{\pgfqpoint{5.977852in}{1.901927in}}%
\pgfpathlineto{\pgfqpoint{5.985552in}{1.928725in}}%
\pgfpathlineto{\pgfqpoint{5.992226in}{1.944598in}}%
\pgfpathlineto{\pgfqpoint{5.997959in}{1.953025in}}%
\pgfpathlineto{\pgfqpoint{6.002921in}{1.956676in}}%
\pgfpathlineto{\pgfqpoint{6.007285in}{1.957291in}}%
\pgfpathlineto{\pgfqpoint{6.011734in}{1.955617in}}%
\pgfpathlineto{\pgfqpoint{6.016783in}{1.951164in}}%
\pgfpathlineto{\pgfqpoint{6.022943in}{1.942495in}}%
\pgfpathlineto{\pgfqpoint{6.030644in}{1.927496in}}%
\pgfpathlineto{\pgfqpoint{6.040826in}{1.902334in}}%
\pgfpathlineto{\pgfqpoint{6.056312in}{1.857333in}}%
\pgfpathlineto{\pgfqpoint{6.088997in}{1.761568in}}%
\pgfpathlineto{\pgfqpoint{6.103628in}{1.726151in}}%
\pgfpathlineto{\pgfqpoint{6.116206in}{1.701243in}}%
\pgfpathlineto{\pgfqpoint{6.127500in}{1.683569in}}%
\pgfpathlineto{\pgfqpoint{6.137768in}{1.671420in}}%
\pgfpathlineto{\pgfqpoint{6.147265in}{1.663483in}}%
\pgfpathlineto{\pgfqpoint{6.156078in}{1.658915in}}%
\pgfpathlineto{\pgfqpoint{6.164463in}{1.657029in}}%
\pgfpathlineto{\pgfqpoint{6.172592in}{1.657457in}}%
\pgfpathlineto{\pgfqpoint{6.180891in}{1.660160in}}%
\pgfpathlineto{\pgfqpoint{6.189533in}{1.665382in}}%
\pgfpathlineto{\pgfqpoint{6.198774in}{1.673653in}}%
\pgfpathlineto{\pgfqpoint{6.208785in}{1.685702in}}%
\pgfpathlineto{\pgfqpoint{6.219822in}{1.702608in}}%
\pgfpathlineto{\pgfqpoint{6.232229in}{1.725877in}}%
\pgfpathlineto{\pgfqpoint{6.246774in}{1.758153in}}%
\pgfpathlineto{\pgfqpoint{6.268080in}{1.811602in}}%
\pgfpathlineto{\pgfqpoint{6.289898in}{1.864645in}}%
\pgfpathlineto{\pgfqpoint{6.300764in}{1.885509in}}%
\pgfpathlineto{\pgfqpoint{6.308807in}{1.896553in}}%
\pgfpathlineto{\pgfqpoint{6.315139in}{1.901793in}}%
\pgfpathlineto{\pgfqpoint{6.320358in}{1.903402in}}%
\pgfpathlineto{\pgfqpoint{6.324979in}{1.902532in}}%
\pgfpathlineto{\pgfqpoint{6.329599in}{1.899313in}}%
\pgfpathlineto{\pgfqpoint{6.334647in}{1.892920in}}%
\pgfpathlineto{\pgfqpoint{6.340380in}{1.881812in}}%
\pgfpathlineto{\pgfqpoint{6.346968in}{1.863791in}}%
\pgfpathlineto{\pgfqpoint{6.354498in}{1.836255in}}%
\pgfpathlineto{\pgfqpoint{6.363311in}{1.795052in}}%
\pgfpathlineto{\pgfqpoint{6.374348in}{1.731925in}}%
\pgfpathlineto{\pgfqpoint{6.411739in}{1.507732in}}%
\pgfpathlineto{\pgfqpoint{6.418413in}{1.484986in}}%
\pgfpathlineto{\pgfqpoint{6.423375in}{1.474970in}}%
\pgfpathlineto{\pgfqpoint{6.427055in}{1.471808in}}%
\pgfpathlineto{\pgfqpoint{6.429878in}{1.471999in}}%
\pgfpathlineto{\pgfqpoint{6.432787in}{1.474652in}}%
\pgfpathlineto{\pgfqpoint{6.436295in}{1.481229in}}%
\pgfpathlineto{\pgfqpoint{6.440659in}{1.494569in}}%
\pgfpathlineto{\pgfqpoint{6.446135in}{1.519168in}}%
\pgfpathlineto{\pgfqpoint{6.452980in}{1.561099in}}%
\pgfpathlineto{\pgfqpoint{6.462306in}{1.633613in}}%
\pgfpathlineto{\pgfqpoint{6.486264in}{1.826605in}}%
\pgfpathlineto{\pgfqpoint{6.492595in}{1.857867in}}%
\pgfpathlineto{\pgfqpoint{6.497130in}{1.870720in}}%
\pgfpathlineto{\pgfqpoint{6.500382in}{1.874406in}}%
\pgfpathlineto{\pgfqpoint{6.502606in}{1.874127in}}%
\pgfpathlineto{\pgfqpoint{6.505002in}{1.871228in}}%
\pgfpathlineto{\pgfqpoint{6.508082in}{1.863528in}}%
\pgfpathlineto{\pgfqpoint{6.512104in}{1.846883in}}%
\pgfpathlineto{\pgfqpoint{6.517323in}{1.814907in}}%
\pgfpathlineto{\pgfqpoint{6.524339in}{1.756662in}}%
\pgfpathlineto{\pgfqpoint{6.539655in}{1.604002in}}%
\pgfpathlineto{\pgfqpoint{6.548125in}{1.534019in}}%
\pgfpathlineto{\pgfqpoint{6.553601in}{1.504863in}}%
\pgfpathlineto{\pgfqpoint{6.557452in}{1.494264in}}%
\pgfpathlineto{\pgfqpoint{6.560019in}{1.492207in}}%
\pgfpathlineto{\pgfqpoint{6.561901in}{1.493338in}}%
\pgfpathlineto{\pgfqpoint{6.564297in}{1.498031in}}%
\pgfpathlineto{\pgfqpoint{6.567548in}{1.510154in}}%
\pgfpathlineto{\pgfqpoint{6.571912in}{1.536274in}}%
\pgfpathlineto{\pgfqpoint{6.577730in}{1.586114in}}%
\pgfpathlineto{\pgfqpoint{6.587484in}{1.692541in}}%
\pgfpathlineto{\pgfqpoint{6.598522in}{1.805580in}}%
\pgfpathlineto{\pgfqpoint{6.604083in}{1.842199in}}%
\pgfpathlineto{\pgfqpoint{6.607934in}{1.855321in}}%
\pgfpathlineto{\pgfqpoint{6.610415in}{1.857819in}}%
\pgfpathlineto{\pgfqpoint{6.612126in}{1.856718in}}%
\pgfpathlineto{\pgfqpoint{6.614265in}{1.852084in}}%
\pgfpathlineto{\pgfqpoint{6.617345in}{1.839169in}}%
\pgfpathlineto{\pgfqpoint{6.621538in}{1.810546in}}%
\pgfpathlineto{\pgfqpoint{6.627356in}{1.753534in}}%
\pgfpathlineto{\pgfqpoint{6.649004in}{1.524401in}}%
\pgfpathlineto{\pgfqpoint{6.652854in}{1.509246in}}%
\pgfpathlineto{\pgfqpoint{6.655250in}{1.506566in}}%
\pgfpathlineto{\pgfqpoint{6.656790in}{1.507685in}}%
\pgfpathlineto{\pgfqpoint{6.658843in}{1.512646in}}%
\pgfpathlineto{\pgfqpoint{6.661838in}{1.526810in}}%
\pgfpathlineto{\pgfqpoint{6.665945in}{1.558511in}}%
\pgfpathlineto{\pgfqpoint{6.671934in}{1.624481in}}%
\pgfpathlineto{\pgfqpoint{6.688619in}{1.817630in}}%
\pgfpathlineto{\pgfqpoint{6.692812in}{1.840076in}}%
\pgfpathlineto{\pgfqpoint{6.695464in}{1.845046in}}%
\pgfpathlineto{\pgfqpoint{6.696919in}{1.844556in}}%
\pgfpathlineto{\pgfqpoint{6.698715in}{1.840783in}}%
\pgfpathlineto{\pgfqpoint{6.701368in}{1.828943in}}%
\pgfpathlineto{\pgfqpoint{6.705133in}{1.800239in}}%
\pgfpathlineto{\pgfqpoint{6.710609in}{1.738732in}}%
\pgfpathlineto{\pgfqpoint{6.727293in}{1.539331in}}%
\pgfpathlineto{\pgfqpoint{6.731058in}{1.521019in}}%
\pgfpathlineto{\pgfqpoint{6.733368in}{1.517986in}}%
\pgfpathlineto{\pgfqpoint{6.734737in}{1.519260in}}%
\pgfpathlineto{\pgfqpoint{6.736705in}{1.525085in}}%
\pgfpathlineto{\pgfqpoint{6.739614in}{1.542014in}}%
\pgfpathlineto{\pgfqpoint{6.743721in}{1.580856in}}%
\pgfpathlineto{\pgfqpoint{6.750481in}{1.669898in}}%
\pgfpathlineto{\pgfqpoint{6.760491in}{1.797301in}}%
\pgfpathlineto{\pgfqpoint{6.764855in}{1.827603in}}%
\pgfpathlineto{\pgfqpoint{6.767593in}{1.834535in}}%
\pgfpathlineto{\pgfqpoint{6.768962in}{1.834230in}}%
\pgfpathlineto{\pgfqpoint{6.770588in}{1.830575in}}%
\pgfpathlineto{\pgfqpoint{6.773069in}{1.818265in}}%
\pgfpathlineto{\pgfqpoint{6.776663in}{1.787268in}}%
\pgfpathlineto{\pgfqpoint{6.782139in}{1.717695in}}%
\pgfpathlineto{\pgfqpoint{6.794631in}{1.554711in}}%
\pgfpathlineto{\pgfqpoint{6.798481in}{1.531533in}}%
\pgfpathlineto{\pgfqpoint{6.800791in}{1.527594in}}%
\pgfpathlineto{\pgfqpoint{6.801989in}{1.528654in}}%
\pgfpathlineto{\pgfqpoint{6.803786in}{1.534208in}}%
\pgfpathlineto{\pgfqpoint{6.806524in}{1.551451in}}%
\pgfpathlineto{\pgfqpoint{6.810545in}{1.593462in}}%
\pgfpathlineto{\pgfqpoint{6.817818in}{1.698355in}}%
\pgfpathlineto{\pgfqpoint{6.825348in}{1.795087in}}%
\pgfpathlineto{\pgfqpoint{6.829284in}{1.821262in}}%
\pgfpathlineto{\pgfqpoint{6.831594in}{1.825779in}}%
\pgfpathlineto{\pgfqpoint{6.832792in}{1.824767in}}%
\pgfpathlineto{\pgfqpoint{6.834503in}{1.819352in}}%
\pgfpathlineto{\pgfqpoint{6.837155in}{1.802120in}}%
\pgfpathlineto{\pgfqpoint{6.841091in}{1.759475in}}%
\pgfpathlineto{\pgfqpoint{6.848792in}{1.644966in}}%
\pgfpathlineto{\pgfqpoint{6.855380in}{1.562038in}}%
\pgfpathlineto{\pgfqpoint{6.859059in}{1.539133in}}%
\pgfpathlineto{\pgfqpoint{6.861198in}{1.535992in}}%
\pgfpathlineto{\pgfqpoint{6.862396in}{1.537665in}}%
\pgfpathlineto{\pgfqpoint{6.864279in}{1.545220in}}%
\pgfpathlineto{\pgfqpoint{6.867188in}{1.568007in}}%
\pgfpathlineto{\pgfqpoint{6.871552in}{1.622543in}}%
\pgfpathlineto{\pgfqpoint{6.885156in}{1.804933in}}%
\pgfpathlineto{\pgfqpoint{6.888065in}{1.817197in}}%
\pgfpathlineto{\pgfqpoint{6.889349in}{1.817779in}}%
\pgfpathlineto{\pgfqpoint{6.889520in}{1.817629in}}%
\pgfpathlineto{\pgfqpoint{6.890889in}{1.814497in}}%
\pgfpathlineto{\pgfqpoint{6.893028in}{1.802918in}}%
\pgfpathlineto{\pgfqpoint{6.896279in}{1.771239in}}%
\pgfpathlineto{\pgfqpoint{6.901755in}{1.691714in}}%
\pgfpathlineto{\pgfqpoint{6.910482in}{1.569186in}}%
\pgfpathlineto{\pgfqpoint{6.913991in}{1.546315in}}%
\pgfpathlineto{\pgfqpoint{6.915873in}{1.543399in}}%
\pgfpathlineto{\pgfqpoint{6.915958in}{1.543429in}}%
\pgfpathlineto{\pgfqpoint{6.917071in}{1.545118in}}%
\pgfpathlineto{\pgfqpoint{6.918868in}{1.552844in}}%
\pgfpathlineto{\pgfqpoint{6.921691in}{1.576603in}}%
\pgfpathlineto{\pgfqpoint{6.926140in}{1.636361in}}%
\pgfpathlineto{\pgfqpoint{6.937435in}{1.794748in}}%
\pgfpathlineto{\pgfqpoint{6.940429in}{1.809812in}}%
\pgfpathlineto{\pgfqpoint{6.941713in}{1.810811in}}%
\pgfpathlineto{\pgfqpoint{6.941969in}{1.810607in}}%
\pgfpathlineto{\pgfqpoint{6.943253in}{1.807570in}}%
\pgfpathlineto{\pgfqpoint{6.945392in}{1.795256in}}%
\pgfpathlineto{\pgfqpoint{6.948643in}{1.761121in}}%
\pgfpathlineto{\pgfqpoint{6.954547in}{1.670012in}}%
\pgfpathlineto{\pgfqpoint{6.961563in}{1.572614in}}%
\pgfpathlineto{\pgfqpoint{6.964900in}{1.551947in}}%
\pgfpathlineto{\pgfqpoint{6.966355in}{1.550096in}}%
\pgfpathlineto{\pgfqpoint{6.966611in}{1.550240in}}%
\pgfpathlineto{\pgfqpoint{6.967809in}{1.552776in}}%
\pgfpathlineto{\pgfqpoint{6.969777in}{1.563448in}}%
\pgfpathlineto{\pgfqpoint{6.972857in}{1.594735in}}%
\pgfpathlineto{\pgfqpoint{6.978248in}{1.677310in}}%
\pgfpathlineto{\pgfqpoint{6.985606in}{1.782252in}}%
\pgfpathlineto{\pgfqpoint{6.988858in}{1.802670in}}%
\pgfpathlineto{\pgfqpoint{6.990312in}{1.804458in}}%
\pgfpathlineto{\pgfqpoint{6.990569in}{1.804281in}}%
\pgfpathlineto{\pgfqpoint{6.991767in}{1.801512in}}%
\pgfpathlineto{\pgfqpoint{6.993735in}{1.790194in}}%
\pgfpathlineto{\pgfqpoint{6.996815in}{1.757465in}}%
\pgfpathlineto{\pgfqpoint{7.002548in}{1.667105in}}%
\pgfpathlineto{\pgfqpoint{7.009050in}{1.576318in}}%
\pgfpathlineto{\pgfqpoint{7.012216in}{1.557471in}}%
\pgfpathlineto{\pgfqpoint{7.013414in}{1.556209in}}%
\pgfpathlineto{\pgfqpoint{7.013756in}{1.556463in}}%
\pgfpathlineto{\pgfqpoint{7.014954in}{1.559497in}}%
\pgfpathlineto{\pgfqpoint{7.017008in}{1.572205in}}%
\pgfpathlineto{\pgfqpoint{7.020259in}{1.609202in}}%
\pgfpathlineto{\pgfqpoint{7.027532in}{1.726954in}}%
\pgfpathlineto{\pgfqpoint{7.032409in}{1.785873in}}%
\pgfpathlineto{\pgfqpoint{7.035147in}{1.798288in}}%
\pgfpathlineto{\pgfqpoint{7.036003in}{1.798521in}}%
\pgfpathlineto{\pgfqpoint{7.036345in}{1.798118in}}%
\pgfpathlineto{\pgfqpoint{7.037714in}{1.793691in}}%
\pgfpathlineto{\pgfqpoint{7.039938in}{1.777338in}}%
\pgfpathlineto{\pgfqpoint{7.043532in}{1.731198in}}%
\pgfpathlineto{\pgfqpoint{7.055169in}{1.568278in}}%
\pgfpathlineto{\pgfqpoint{7.057393in}{1.561805in}}%
\pgfpathlineto{\pgfqpoint{7.058334in}{1.562776in}}%
\pgfpathlineto{\pgfqpoint{7.059874in}{1.569095in}}%
\pgfpathlineto{\pgfqpoint{7.062356in}{1.590772in}}%
\pgfpathlineto{\pgfqpoint{7.066463in}{1.650056in}}%
\pgfpathlineto{\pgfqpoint{7.075276in}{1.779129in}}%
\pgfpathlineto{\pgfqpoint{7.078014in}{1.792774in}}%
\pgfpathlineto{\pgfqpoint{7.078869in}{1.793156in}}%
\pgfpathlineto{\pgfqpoint{7.079212in}{1.792778in}}%
\pgfpathlineto{\pgfqpoint{7.080495in}{1.788682in}}%
\pgfpathlineto{\pgfqpoint{7.082634in}{1.772890in}}%
\pgfpathlineto{\pgfqpoint{7.086142in}{1.727125in}}%
\pgfpathlineto{\pgfqpoint{7.097009in}{1.573435in}}%
\pgfpathlineto{\pgfqpoint{7.099233in}{1.567003in}}%
\pgfpathlineto{\pgfqpoint{7.100174in}{1.568228in}}%
\pgfpathlineto{\pgfqpoint{7.101715in}{1.575245in}}%
\pgfpathlineto{\pgfqpoint{7.104281in}{1.599618in}}%
\pgfpathlineto{\pgfqpoint{7.108731in}{1.667891in}}%
\pgfpathlineto{\pgfqpoint{7.115918in}{1.772947in}}%
\pgfpathlineto{\pgfqpoint{7.118656in}{1.787660in}}%
\pgfpathlineto{\pgfqpoint{7.119512in}{1.788161in}}%
\pgfpathlineto{\pgfqpoint{7.119939in}{1.787658in}}%
\pgfpathlineto{\pgfqpoint{7.121223in}{1.783170in}}%
\pgfpathlineto{\pgfqpoint{7.123447in}{1.765384in}}%
\pgfpathlineto{\pgfqpoint{7.127127in}{1.714150in}}%
\pgfpathlineto{\pgfqpoint{7.136282in}{1.581299in}}%
\pgfpathlineto{\pgfqpoint{7.138678in}{1.571913in}}%
\pgfpathlineto{\pgfqpoint{7.139619in}{1.572588in}}%
\pgfpathlineto{\pgfqpoint{7.140988in}{1.577971in}}%
\pgfpathlineto{\pgfqpoint{7.143298in}{1.598071in}}%
\pgfpathlineto{\pgfqpoint{7.147234in}{1.655943in}}%
\pgfpathlineto{\pgfqpoint{7.155106in}{1.771411in}}%
\pgfpathlineto{\pgfqpoint{7.157672in}{1.783350in}}%
\pgfpathlineto{\pgfqpoint{7.158442in}{1.783301in}}%
\pgfpathlineto{\pgfqpoint{7.158699in}{1.782905in}}%
\pgfpathlineto{\pgfqpoint{7.160068in}{1.777631in}}%
\pgfpathlineto{\pgfqpoint{7.162378in}{1.757448in}}%
\pgfpathlineto{\pgfqpoint{7.166314in}{1.699104in}}%
\pgfpathlineto{\pgfqpoint{7.173929in}{1.587901in}}%
\pgfpathlineto{\pgfqpoint{7.176496in}{1.576477in}}%
\pgfpathlineto{\pgfqpoint{7.177266in}{1.576771in}}%
\pgfpathlineto{\pgfqpoint{7.177437in}{1.577073in}}%
\pgfpathlineto{\pgfqpoint{7.178806in}{1.582551in}}%
\pgfpathlineto{\pgfqpoint{7.181116in}{1.603270in}}%
\pgfpathlineto{\pgfqpoint{7.185138in}{1.663938in}}%
\pgfpathlineto{\pgfqpoint{7.192240in}{1.767347in}}%
\pgfpathlineto{\pgfqpoint{7.194806in}{1.779046in}}%
\pgfpathlineto{\pgfqpoint{7.195577in}{1.778758in}}%
\pgfpathlineto{\pgfqpoint{7.195748in}{1.778452in}}%
\pgfpathlineto{\pgfqpoint{7.197117in}{1.772882in}}%
\pgfpathlineto{\pgfqpoint{7.199427in}{1.751807in}}%
\pgfpathlineto{\pgfqpoint{7.203534in}{1.689033in}}%
\pgfpathlineto{\pgfqpoint{7.210208in}{1.592417in}}%
\pgfpathlineto{\pgfqpoint{7.212689in}{1.580792in}}%
\pgfpathlineto{\pgfqpoint{7.213374in}{1.580846in}}%
\pgfpathlineto{\pgfqpoint{7.213630in}{1.581238in}}%
\pgfpathlineto{\pgfqpoint{7.214914in}{1.586186in}}%
\pgfpathlineto{\pgfqpoint{7.217138in}{1.605775in}}%
\pgfpathlineto{\pgfqpoint{7.221074in}{1.664947in}}%
\pgfpathlineto{\pgfqpoint{7.227919in}{1.764096in}}%
\pgfpathlineto{\pgfqpoint{7.230400in}{1.774962in}}%
\pgfpathlineto{\pgfqpoint{7.231256in}{1.774316in}}%
\pgfpathlineto{\pgfqpoint{7.232625in}{1.768579in}}%
\pgfpathlineto{\pgfqpoint{7.234935in}{1.746861in}}%
\pgfpathlineto{\pgfqpoint{7.239128in}{1.681585in}}%
\pgfpathlineto{\pgfqpoint{7.245203in}{1.595430in}}%
\pgfpathlineto{\pgfqpoint{7.247599in}{1.584737in}}%
\pgfpathlineto{\pgfqpoint{7.248369in}{1.585113in}}%
\pgfpathlineto{\pgfqpoint{7.248540in}{1.585452in}}%
\pgfpathlineto{\pgfqpoint{7.249909in}{1.591456in}}%
\pgfpathlineto{\pgfqpoint{7.252219in}{1.613740in}}%
\pgfpathlineto{\pgfqpoint{7.256583in}{1.682590in}}%
\pgfpathlineto{\pgfqpoint{7.262230in}{1.761000in}}%
\pgfpathlineto{\pgfqpoint{7.264625in}{1.771090in}}%
\pgfpathlineto{\pgfqpoint{7.265481in}{1.770260in}}%
\pgfpathlineto{\pgfqpoint{7.266850in}{1.764084in}}%
\pgfpathlineto{\pgfqpoint{7.269160in}{1.741383in}}%
\pgfpathlineto{\pgfqpoint{7.273695in}{1.669185in}}%
\pgfpathlineto{\pgfqpoint{7.278914in}{1.598191in}}%
\pgfpathlineto{\pgfqpoint{7.281225in}{1.588480in}}%
\pgfpathlineto{\pgfqpoint{7.282080in}{1.589268in}}%
\pgfpathlineto{\pgfqpoint{7.283449in}{1.595449in}}%
\pgfpathlineto{\pgfqpoint{7.285759in}{1.618315in}}%
\pgfpathlineto{\pgfqpoint{7.290380in}{1.692115in}}%
\pgfpathlineto{\pgfqpoint{7.295342in}{1.758309in}}%
\pgfpathlineto{\pgfqpoint{7.297653in}{1.767421in}}%
\pgfpathlineto{\pgfqpoint{7.298508in}{1.766341in}}%
\pgfpathlineto{\pgfqpoint{7.299963in}{1.759027in}}%
\pgfpathlineto{\pgfqpoint{7.302444in}{1.732457in}}%
\pgfpathlineto{\pgfqpoint{7.313738in}{1.592039in}}%
\pgfpathlineto{\pgfqpoint{7.314594in}{1.593227in}}%
\pgfpathlineto{\pgfqpoint{7.316048in}{1.600786in}}%
\pgfpathlineto{\pgfqpoint{7.318530in}{1.627837in}}%
\pgfpathlineto{\pgfqpoint{7.329311in}{1.763830in}}%
\pgfpathlineto{\pgfqpoint{7.329995in}{1.763575in}}%
\pgfpathlineto{\pgfqpoint{7.330166in}{1.763262in}}%
\pgfpathlineto{\pgfqpoint{7.331450in}{1.757782in}}%
\pgfpathlineto{\pgfqpoint{7.333674in}{1.736346in}}%
\pgfpathlineto{\pgfqpoint{7.338209in}{1.664726in}}%
\pgfpathlineto{\pgfqpoint{7.343001in}{1.602843in}}%
\pgfpathlineto{\pgfqpoint{7.345140in}{1.595458in}}%
\pgfpathlineto{\pgfqpoint{7.345995in}{1.596892in}}%
\pgfpathlineto{\pgfqpoint{7.347535in}{1.605662in}}%
\pgfpathlineto{\pgfqpoint{7.350188in}{1.636680in}}%
\pgfpathlineto{\pgfqpoint{7.359600in}{1.759532in}}%
\pgfpathlineto{\pgfqpoint{7.360455in}{1.760584in}}%
\pgfpathlineto{\pgfqpoint{7.360969in}{1.759986in}}%
\pgfpathlineto{\pgfqpoint{7.362252in}{1.754529in}}%
\pgfpathlineto{\pgfqpoint{7.364477in}{1.732913in}}%
\pgfpathlineto{\pgfqpoint{7.369183in}{1.658632in}}%
\pgfpathlineto{\pgfqpoint{7.373546in}{1.604996in}}%
\pgfpathlineto{\pgfqpoint{7.375514in}{1.598709in}}%
\pgfpathlineto{\pgfqpoint{7.375600in}{1.598743in}}%
\pgfpathlineto{\pgfqpoint{7.376456in}{1.600498in}}%
\pgfpathlineto{\pgfqpoint{7.378081in}{1.610678in}}%
\pgfpathlineto{\pgfqpoint{7.380905in}{1.645884in}}%
\pgfpathlineto{\pgfqpoint{7.389033in}{1.754617in}}%
\pgfpathlineto{\pgfqpoint{7.390317in}{1.757421in}}%
\pgfpathlineto{\pgfqpoint{7.390659in}{1.757181in}}%
\pgfpathlineto{\pgfqpoint{7.391686in}{1.753973in}}%
\pgfpathlineto{\pgfqpoint{7.393568in}{1.739008in}}%
\pgfpathlineto{\pgfqpoint{7.396991in}{1.690029in}}%
\pgfpathlineto{\pgfqpoint{7.402980in}{1.608498in}}%
\pgfpathlineto{\pgfqpoint{7.404948in}{1.601799in}}%
\pgfpathlineto{\pgfqpoint{7.405033in}{1.601820in}}%
\pgfpathlineto{\pgfqpoint{7.405889in}{1.603473in}}%
\pgfpathlineto{\pgfqpoint{7.407429in}{1.612840in}}%
\pgfpathlineto{\pgfqpoint{7.410167in}{1.646247in}}%
\pgfpathlineto{\pgfqpoint{7.418210in}{1.752047in}}%
\pgfpathlineto{\pgfqpoint{7.419408in}{1.754392in}}%
\pgfpathlineto{\pgfqpoint{7.419836in}{1.753971in}}%
\pgfpathlineto{\pgfqpoint{7.421034in}{1.749312in}}%
\pgfpathlineto{\pgfqpoint{7.423173in}{1.729386in}}%
\pgfpathlineto{\pgfqpoint{7.427707in}{1.659456in}}%
\pgfpathlineto{\pgfqpoint{7.431900in}{1.609871in}}%
\pgfpathlineto{\pgfqpoint{7.433611in}{1.604756in}}%
\pgfpathlineto{\pgfqpoint{7.433868in}{1.604904in}}%
\pgfpathlineto{\pgfqpoint{7.434809in}{1.607493in}}%
\pgfpathlineto{\pgfqpoint{7.436606in}{1.620866in}}%
\pgfpathlineto{\pgfqpoint{7.439857in}{1.665827in}}%
\pgfpathlineto{\pgfqpoint{7.445932in}{1.746297in}}%
\pgfpathlineto{\pgfqpoint{7.447644in}{1.751507in}}%
\pgfpathlineto{\pgfqpoint{7.447900in}{1.751367in}}%
\pgfpathlineto{\pgfqpoint{7.448841in}{1.748793in}}%
\pgfpathlineto{\pgfqpoint{7.450638in}{1.735399in}}%
\pgfpathlineto{\pgfqpoint{7.453890in}{1.690446in}}%
\pgfpathlineto{\pgfqpoint{7.459793in}{1.612813in}}%
\pgfpathlineto{\pgfqpoint{7.461505in}{1.607582in}}%
\pgfpathlineto{\pgfqpoint{7.461761in}{1.607724in}}%
\pgfpathlineto{\pgfqpoint{7.462703in}{1.610319in}}%
\pgfpathlineto{\pgfqpoint{7.464499in}{1.623793in}}%
\pgfpathlineto{\pgfqpoint{7.467836in}{1.670156in}}%
\pgfpathlineto{\pgfqpoint{7.473483in}{1.743521in}}%
\pgfpathlineto{\pgfqpoint{7.475195in}{1.748741in}}%
\pgfpathlineto{\pgfqpoint{7.475451in}{1.748591in}}%
\pgfpathlineto{\pgfqpoint{7.476393in}{1.745958in}}%
\pgfpathlineto{\pgfqpoint{7.478189in}{1.732382in}}%
\pgfpathlineto{\pgfqpoint{7.481526in}{1.685990in}}%
\pgfpathlineto{\pgfqpoint{7.487002in}{1.615501in}}%
\pgfpathlineto{\pgfqpoint{7.488714in}{1.610289in}}%
\pgfpathlineto{\pgfqpoint{7.488970in}{1.610444in}}%
\pgfpathlineto{\pgfqpoint{7.489911in}{1.613110in}}%
\pgfpathlineto{\pgfqpoint{7.491708in}{1.626774in}}%
\pgfpathlineto{\pgfqpoint{7.495131in}{1.674527in}}%
\pgfpathlineto{\pgfqpoint{7.500350in}{1.740847in}}%
\pgfpathlineto{\pgfqpoint{7.502061in}{1.746092in}}%
\pgfpathlineto{\pgfqpoint{7.502318in}{1.745937in}}%
\pgfpathlineto{\pgfqpoint{7.503259in}{1.743263in}}%
\pgfpathlineto{\pgfqpoint{7.505056in}{1.729557in}}%
\pgfpathlineto{\pgfqpoint{7.508478in}{1.681879in}}%
\pgfpathlineto{\pgfqpoint{7.513612in}{1.617712in}}%
\pgfpathlineto{\pgfqpoint{7.515238in}{1.612880in}}%
\pgfpathlineto{\pgfqpoint{7.515495in}{1.613023in}}%
\pgfpathlineto{\pgfqpoint{7.516436in}{1.615660in}}%
\pgfpathlineto{\pgfqpoint{7.518233in}{1.629328in}}%
\pgfpathlineto{\pgfqpoint{7.521655in}{1.676844in}}%
\pgfpathlineto{\pgfqpoint{7.526703in}{1.739023in}}%
\pgfpathlineto{\pgfqpoint{7.528329in}{1.743544in}}%
\pgfpathlineto{\pgfqpoint{7.528586in}{1.743348in}}%
\pgfpathlineto{\pgfqpoint{7.529612in}{1.740090in}}%
\pgfpathlineto{\pgfqpoint{7.531495in}{1.724629in}}%
\pgfpathlineto{\pgfqpoint{7.535431in}{1.668168in}}%
\pgfpathlineto{\pgfqpoint{7.539709in}{1.619257in}}%
\pgfpathlineto{\pgfqpoint{7.541163in}{1.615368in}}%
\pgfpathlineto{\pgfqpoint{7.541506in}{1.615617in}}%
\pgfpathlineto{\pgfqpoint{7.542532in}{1.619008in}}%
\pgfpathlineto{\pgfqpoint{7.544415in}{1.634692in}}%
\pgfpathlineto{\pgfqpoint{7.548522in}{1.693620in}}%
\pgfpathlineto{\pgfqpoint{7.552543in}{1.737688in}}%
\pgfpathlineto{\pgfqpoint{7.553912in}{1.741115in}}%
\pgfpathlineto{\pgfqpoint{7.554254in}{1.740859in}}%
\pgfpathlineto{\pgfqpoint{7.555281in}{1.737439in}}%
\pgfpathlineto{\pgfqpoint{7.557163in}{1.721699in}}%
\pgfpathlineto{\pgfqpoint{7.561356in}{1.661801in}}%
\pgfpathlineto{\pgfqpoint{7.565206in}{1.620899in}}%
\pgfpathlineto{\pgfqpoint{7.566575in}{1.617770in}}%
\pgfpathlineto{\pgfqpoint{7.566918in}{1.618104in}}%
\pgfpathlineto{\pgfqpoint{7.568030in}{1.622239in}}%
\pgfpathlineto{\pgfqpoint{7.570083in}{1.640864in}}%
\pgfpathlineto{\pgfqpoint{7.578982in}{1.738774in}}%
\pgfpathlineto{\pgfqpoint{7.579239in}{1.738625in}}%
\pgfpathlineto{\pgfqpoint{7.580180in}{1.735937in}}%
\pgfpathlineto{\pgfqpoint{7.581977in}{1.722122in}}%
\pgfpathlineto{\pgfqpoint{7.585656in}{1.671502in}}%
\pgfpathlineto{\pgfqpoint{7.589934in}{1.623595in}}%
\pgfpathlineto{\pgfqpoint{7.591388in}{1.620069in}}%
\pgfpathlineto{\pgfqpoint{7.591645in}{1.620286in}}%
\pgfpathlineto{\pgfqpoint{7.592672in}{1.623657in}}%
\pgfpathlineto{\pgfqpoint{7.594554in}{1.639324in}}%
\pgfpathlineto{\pgfqpoint{7.599089in}{1.702767in}}%
\pgfpathlineto{\pgfqpoint{7.602512in}{1.734627in}}%
\pgfpathlineto{\pgfqpoint{7.603538in}{1.736521in}}%
\pgfpathlineto{\pgfqpoint{7.603966in}{1.736115in}}%
\pgfpathlineto{\pgfqpoint{7.605078in}{1.731817in}}%
\pgfpathlineto{\pgfqpoint{7.607132in}{1.712980in}}%
\pgfpathlineto{\pgfqpoint{7.615603in}{1.622261in}}%
\pgfpathlineto{\pgfqpoint{7.615688in}{1.622279in}}%
\pgfpathlineto{\pgfqpoint{7.616544in}{1.624007in}}%
\pgfpathlineto{\pgfqpoint{7.618084in}{1.633845in}}%
\pgfpathlineto{\pgfqpoint{7.621079in}{1.671039in}}%
\pgfpathlineto{\pgfqpoint{7.626298in}{1.731332in}}%
\pgfpathlineto{\pgfqpoint{7.627581in}{1.734355in}}%
\pgfpathlineto{\pgfqpoint{7.627924in}{1.734094in}}%
\pgfpathlineto{\pgfqpoint{7.628950in}{1.730634in}}%
\pgfpathlineto{\pgfqpoint{7.630918in}{1.713902in}}%
\pgfpathlineto{\pgfqpoint{7.639389in}{1.624386in}}%
\pgfpathlineto{\pgfqpoint{7.639817in}{1.624720in}}%
\pgfpathlineto{\pgfqpoint{7.640929in}{1.628838in}}%
\pgfpathlineto{\pgfqpoint{7.642983in}{1.647341in}}%
\pgfpathlineto{\pgfqpoint{7.651197in}{1.732263in}}%
\pgfpathlineto{\pgfqpoint{7.651967in}{1.730934in}}%
\pgfpathlineto{\pgfqpoint{7.653421in}{1.722470in}}%
\pgfpathlineto{\pgfqpoint{7.656245in}{1.689305in}}%
\pgfpathlineto{\pgfqpoint{7.661635in}{1.628837in}}%
\pgfpathlineto{\pgfqpoint{7.662833in}{1.626445in}}%
\pgfpathlineto{\pgfqpoint{7.663175in}{1.626776in}}%
\pgfpathlineto{\pgfqpoint{7.664288in}{1.630911in}}%
\pgfpathlineto{\pgfqpoint{7.666341in}{1.649372in}}%
\pgfpathlineto{\pgfqpoint{7.674213in}{1.730256in}}%
\pgfpathlineto{\pgfqpoint{7.674983in}{1.729275in}}%
\pgfpathlineto{\pgfqpoint{7.676352in}{1.722094in}}%
\pgfpathlineto{\pgfqpoint{7.678919in}{1.693885in}}%
\pgfpathlineto{\pgfqpoint{7.684737in}{1.630093in}}%
\pgfpathlineto{\pgfqpoint{7.685764in}{1.628422in}}%
\pgfpathlineto{\pgfqpoint{7.686192in}{1.628922in}}%
\pgfpathlineto{\pgfqpoint{7.687390in}{1.633977in}}%
\pgfpathlineto{\pgfqpoint{7.689614in}{1.655479in}}%
\pgfpathlineto{\pgfqpoint{7.696545in}{1.727962in}}%
\pgfpathlineto{\pgfqpoint{7.697144in}{1.728272in}}%
\pgfpathlineto{\pgfqpoint{7.697571in}{1.727648in}}%
\pgfpathlineto{\pgfqpoint{7.698855in}{1.721694in}}%
\pgfpathlineto{\pgfqpoint{7.701251in}{1.697168in}}%
\pgfpathlineto{\pgfqpoint{7.707411in}{1.631449in}}%
\pgfpathlineto{\pgfqpoint{7.708267in}{1.630322in}}%
\pgfpathlineto{\pgfqpoint{7.708695in}{1.630813in}}%
\pgfpathlineto{\pgfqpoint{7.709892in}{1.635827in}}%
\pgfpathlineto{\pgfqpoint{7.712117in}{1.657141in}}%
\pgfpathlineto{\pgfqpoint{7.718791in}{1.725982in}}%
\pgfpathlineto{\pgfqpoint{7.719475in}{1.726389in}}%
\pgfpathlineto{\pgfqpoint{7.719903in}{1.725731in}}%
\pgfpathlineto{\pgfqpoint{7.721187in}{1.719702in}}%
\pgfpathlineto{\pgfqpoint{7.723582in}{1.695265in}}%
\pgfpathlineto{\pgfqpoint{7.729572in}{1.633153in}}%
\pgfpathlineto{\pgfqpoint{7.730427in}{1.632176in}}%
\pgfpathlineto{\pgfqpoint{7.730855in}{1.632737in}}%
\pgfpathlineto{\pgfqpoint{7.732053in}{1.637919in}}%
\pgfpathlineto{\pgfqpoint{7.734363in}{1.660348in}}%
\pgfpathlineto{\pgfqpoint{7.740609in}{1.723952in}}%
\pgfpathlineto{\pgfqpoint{7.741379in}{1.724602in}}%
\pgfpathlineto{\pgfqpoint{7.741807in}{1.723985in}}%
\pgfpathlineto{\pgfqpoint{7.743091in}{1.718096in}}%
\pgfpathlineto{\pgfqpoint{7.745486in}{1.694036in}}%
\pgfpathlineto{\pgfqpoint{7.751305in}{1.634947in}}%
\pgfpathlineto{\pgfqpoint{7.752160in}{1.633942in}}%
\pgfpathlineto{\pgfqpoint{7.752588in}{1.634484in}}%
\pgfpathlineto{\pgfqpoint{7.753786in}{1.639592in}}%
\pgfpathlineto{\pgfqpoint{7.756096in}{1.661739in}}%
\pgfpathlineto{\pgfqpoint{7.762171in}{1.722213in}}%
\pgfpathlineto{\pgfqpoint{7.762941in}{1.722863in}}%
\pgfpathlineto{\pgfqpoint{7.763369in}{1.722251in}}%
\pgfpathlineto{\pgfqpoint{7.764652in}{1.716406in}}%
\pgfpathlineto{\pgfqpoint{7.767134in}{1.691567in}}%
\pgfpathlineto{\pgfqpoint{7.772610in}{1.636889in}}%
\pgfpathlineto{\pgfqpoint{7.773551in}{1.635654in}}%
\pgfpathlineto{\pgfqpoint{7.773979in}{1.636198in}}%
\pgfpathlineto{\pgfqpoint{7.775177in}{1.641287in}}%
\pgfpathlineto{\pgfqpoint{7.777487in}{1.663210in}}%
\pgfpathlineto{\pgfqpoint{7.783391in}{1.720549in}}%
\pgfpathlineto{\pgfqpoint{7.784161in}{1.721179in}}%
\pgfpathlineto{\pgfqpoint{7.784588in}{1.720564in}}%
\pgfpathlineto{\pgfqpoint{7.785872in}{1.714742in}}%
\pgfpathlineto{\pgfqpoint{7.788353in}{1.690186in}}%
\pgfpathlineto{\pgfqpoint{7.793658in}{1.638525in}}%
\pgfpathlineto{\pgfqpoint{7.794599in}{1.637306in}}%
\pgfpathlineto{\pgfqpoint{7.795027in}{1.637850in}}%
\pgfpathlineto{\pgfqpoint{7.796225in}{1.642906in}}%
\pgfpathlineto{\pgfqpoint{7.798535in}{1.664571in}}%
\pgfpathlineto{\pgfqpoint{7.804268in}{1.718920in}}%
\pgfpathlineto{\pgfqpoint{7.805038in}{1.719559in}}%
\pgfpathlineto{\pgfqpoint{7.805466in}{1.718956in}}%
\pgfpathlineto{\pgfqpoint{7.806749in}{1.713205in}}%
\pgfpathlineto{\pgfqpoint{7.809230in}{1.689006in}}%
\pgfpathlineto{\pgfqpoint{7.814450in}{1.639906in}}%
\pgfpathlineto{\pgfqpoint{7.815305in}{1.638895in}}%
\pgfpathlineto{\pgfqpoint{7.815733in}{1.639412in}}%
\pgfpathlineto{\pgfqpoint{7.816931in}{1.644360in}}%
\pgfpathlineto{\pgfqpoint{7.819241in}{1.665660in}}%
\pgfpathlineto{\pgfqpoint{7.824803in}{1.717293in}}%
\pgfpathlineto{\pgfqpoint{7.825573in}{1.718006in}}%
\pgfpathlineto{\pgfqpoint{7.826001in}{1.717453in}}%
\pgfpathlineto{\pgfqpoint{7.827284in}{1.711882in}}%
\pgfpathlineto{\pgfqpoint{7.829765in}{1.688181in}}%
\pgfpathlineto{\pgfqpoint{7.834814in}{1.641565in}}%
\pgfpathlineto{\pgfqpoint{7.835669in}{1.640423in}}%
\pgfpathlineto{\pgfqpoint{7.836183in}{1.641032in}}%
\pgfpathlineto{\pgfqpoint{7.837466in}{1.646642in}}%
\pgfpathlineto{\pgfqpoint{7.839947in}{1.670241in}}%
\pgfpathlineto{\pgfqpoint{7.844910in}{1.715394in}}%
\pgfpathlineto{\pgfqpoint{7.845766in}{1.716514in}}%
\pgfpathlineto{\pgfqpoint{7.846279in}{1.715899in}}%
\pgfpathlineto{\pgfqpoint{7.847562in}{1.710299in}}%
\pgfpathlineto{\pgfqpoint{7.850044in}{1.686863in}}%
\pgfpathlineto{\pgfqpoint{7.854921in}{1.643051in}}%
\pgfpathlineto{\pgfqpoint{7.855776in}{1.641914in}}%
\pgfpathlineto{\pgfqpoint{7.856290in}{1.642512in}}%
\pgfpathlineto{\pgfqpoint{7.857573in}{1.648047in}}%
\pgfpathlineto{\pgfqpoint{7.860055in}{1.671254in}}%
\pgfpathlineto{\pgfqpoint{7.864932in}{1.714095in}}%
\pgfpathlineto{\pgfqpoint{7.865787in}{1.715027in}}%
\pgfpathlineto{\pgfqpoint{7.866215in}{1.714498in}}%
\pgfpathlineto{\pgfqpoint{7.867499in}{1.709087in}}%
\pgfpathlineto{\pgfqpoint{7.869980in}{1.686180in}}%
\pgfpathlineto{\pgfqpoint{7.874771in}{1.644402in}}%
\pgfpathlineto{\pgfqpoint{7.875627in}{1.643363in}}%
\pgfpathlineto{\pgfqpoint{7.876055in}{1.643834in}}%
\pgfpathlineto{\pgfqpoint{7.877253in}{1.648539in}}%
\pgfpathlineto{\pgfqpoint{7.879648in}{1.669709in}}%
\pgfpathlineto{\pgfqpoint{7.884697in}{1.712889in}}%
\pgfpathlineto{\pgfqpoint{7.885467in}{1.713611in}}%
\pgfpathlineto{\pgfqpoint{7.885894in}{1.713094in}}%
\pgfpathlineto{\pgfqpoint{7.887178in}{1.707769in}}%
\pgfpathlineto{\pgfqpoint{7.889659in}{1.685281in}}%
\pgfpathlineto{\pgfqpoint{7.894365in}{1.645660in}}%
\pgfpathlineto{\pgfqpoint{7.895221in}{1.644777in}}%
\pgfpathlineto{\pgfqpoint{7.895649in}{1.645313in}}%
\pgfpathlineto{\pgfqpoint{7.896932in}{1.650664in}}%
\pgfpathlineto{\pgfqpoint{7.899499in}{1.673941in}}%
\pgfpathlineto{\pgfqpoint{7.904034in}{1.711345in}}%
\pgfpathlineto{\pgfqpoint{7.904889in}{1.712235in}}%
\pgfpathlineto{\pgfqpoint{7.905317in}{1.711710in}}%
\pgfpathlineto{\pgfqpoint{7.906601in}{1.706416in}}%
\pgfpathlineto{\pgfqpoint{7.909167in}{1.683382in}}%
\pgfpathlineto{\pgfqpoint{7.913617in}{1.647073in}}%
\pgfpathlineto{\pgfqpoint{7.914472in}{1.646113in}}%
\pgfpathlineto{\pgfqpoint{7.914900in}{1.646598in}}%
\pgfpathlineto{\pgfqpoint{7.916098in}{1.651244in}}%
\pgfpathlineto{\pgfqpoint{7.918494in}{1.671744in}}%
\pgfpathlineto{\pgfqpoint{7.923285in}{1.710262in}}%
\pgfpathlineto{\pgfqpoint{7.924055in}{1.710902in}}%
\pgfpathlineto{\pgfqpoint{7.924483in}{1.710363in}}%
\pgfpathlineto{\pgfqpoint{7.925767in}{1.705091in}}%
\pgfpathlineto{\pgfqpoint{7.928333in}{1.682450in}}%
\pgfpathlineto{\pgfqpoint{7.932697in}{1.648229in}}%
\pgfpathlineto{\pgfqpoint{7.933553in}{1.647444in}}%
\pgfpathlineto{\pgfqpoint{7.933981in}{1.648002in}}%
\pgfpathlineto{\pgfqpoint{7.935264in}{1.653298in}}%
\pgfpathlineto{\pgfqpoint{7.937916in}{1.676660in}}%
\pgfpathlineto{\pgfqpoint{7.942109in}{1.708812in}}%
\pgfpathlineto{\pgfqpoint{7.942965in}{1.709614in}}%
\pgfpathlineto{\pgfqpoint{7.943392in}{1.709072in}}%
\pgfpathlineto{\pgfqpoint{7.944676in}{1.703851in}}%
\pgfpathlineto{\pgfqpoint{7.947328in}{1.680772in}}%
\pgfpathlineto{\pgfqpoint{7.951435in}{1.649580in}}%
\pgfpathlineto{\pgfqpoint{7.952291in}{1.648687in}}%
\pgfpathlineto{\pgfqpoint{7.952719in}{1.649176in}}%
\pgfpathlineto{\pgfqpoint{7.954002in}{1.654222in}}%
\pgfpathlineto{\pgfqpoint{7.956569in}{1.676050in}}%
\pgfpathlineto{\pgfqpoint{7.960762in}{1.707565in}}%
\pgfpathlineto{\pgfqpoint{7.961617in}{1.708376in}}%
\pgfpathlineto{\pgfqpoint{7.962045in}{1.707852in}}%
\pgfpathlineto{\pgfqpoint{7.963328in}{1.702746in}}%
\pgfpathlineto{\pgfqpoint{7.965981in}{1.680200in}}%
\pgfpathlineto{\pgfqpoint{7.970002in}{1.650727in}}%
\pgfpathlineto{\pgfqpoint{7.970858in}{1.649920in}}%
\pgfpathlineto{\pgfqpoint{7.971286in}{1.650438in}}%
\pgfpathlineto{\pgfqpoint{7.972569in}{1.655495in}}%
\pgfpathlineto{\pgfqpoint{7.975222in}{1.677781in}}%
\pgfpathlineto{\pgfqpoint{7.979243in}{1.706500in}}%
\pgfpathlineto{\pgfqpoint{7.980013in}{1.707187in}}%
\pgfpathlineto{\pgfqpoint{7.980441in}{1.706716in}}%
\pgfpathlineto{\pgfqpoint{7.981724in}{1.701817in}}%
\pgfpathlineto{\pgfqpoint{7.984377in}{1.679905in}}%
\pgfpathlineto{\pgfqpoint{7.988398in}{1.651714in}}%
\pgfpathlineto{\pgfqpoint{7.989168in}{1.651102in}}%
\pgfpathlineto{\pgfqpoint{7.989596in}{1.651607in}}%
\pgfpathlineto{\pgfqpoint{7.990880in}{1.656563in}}%
\pgfpathlineto{\pgfqpoint{7.993532in}{1.678322in}}%
\pgfpathlineto{\pgfqpoint{7.997468in}{1.705398in}}%
\pgfpathlineto{\pgfqpoint{7.998238in}{1.706013in}}%
\pgfpathlineto{\pgfqpoint{7.998666in}{1.705517in}}%
\pgfpathlineto{\pgfqpoint{7.999949in}{1.700620in}}%
\pgfpathlineto{\pgfqpoint{8.002602in}{1.679136in}}%
\pgfpathlineto{\pgfqpoint{8.006452in}{1.652927in}}%
\pgfpathlineto{\pgfqpoint{8.007308in}{1.652273in}}%
\pgfpathlineto{\pgfqpoint{8.007735in}{1.652834in}}%
\pgfpathlineto{\pgfqpoint{8.009019in}{1.657876in}}%
\pgfpathlineto{\pgfqpoint{8.011842in}{1.680856in}}%
\pgfpathlineto{\pgfqpoint{8.015522in}{1.704416in}}%
\pgfpathlineto{\pgfqpoint{8.016292in}{1.704847in}}%
\pgfpathlineto{\pgfqpoint{8.016719in}{1.704265in}}%
\pgfpathlineto{\pgfqpoint{8.018088in}{1.698705in}}%
\pgfpathlineto{\pgfqpoint{8.021169in}{1.673386in}}%
\pgfpathlineto{\pgfqpoint{8.024506in}{1.653700in}}%
\pgfpathlineto{\pgfqpoint{8.025190in}{1.653383in}}%
\pgfpathlineto{\pgfqpoint{8.025618in}{1.653941in}}%
\pgfpathlineto{\pgfqpoint{8.026987in}{1.659391in}}%
\pgfpathlineto{\pgfqpoint{8.030067in}{1.684325in}}%
\pgfpathlineto{\pgfqpoint{8.033319in}{1.703376in}}%
\pgfpathlineto{\pgfqpoint{8.034089in}{1.703741in}}%
\pgfpathlineto{\pgfqpoint{8.034431in}{1.703305in}}%
\pgfpathlineto{\pgfqpoint{8.035714in}{1.698569in}}%
\pgfpathlineto{\pgfqpoint{8.038452in}{1.677369in}}%
\pgfpathlineto{\pgfqpoint{8.042131in}{1.654837in}}%
\pgfpathlineto{\pgfqpoint{8.042902in}{1.654483in}}%
\pgfpathlineto{\pgfqpoint{8.043244in}{1.654918in}}%
\pgfpathlineto{\pgfqpoint{8.044527in}{1.659616in}}%
\pgfpathlineto{\pgfqpoint{8.047351in}{1.681282in}}%
\pgfpathlineto{\pgfqpoint{8.050859in}{1.702279in}}%
\pgfpathlineto{\pgfqpoint{8.051629in}{1.702707in}}%
\pgfpathlineto{\pgfqpoint{8.052057in}{1.702155in}}%
\pgfpathlineto{\pgfqpoint{8.053426in}{1.696845in}}%
\pgfpathlineto{\pgfqpoint{8.056592in}{1.672226in}}%
\pgfpathlineto{\pgfqpoint{8.059672in}{1.655761in}}%
\pgfpathlineto{\pgfqpoint{8.060356in}{1.655509in}}%
\pgfpathlineto{\pgfqpoint{8.060784in}{1.656077in}}%
\pgfpathlineto{\pgfqpoint{8.062153in}{1.661390in}}%
\pgfpathlineto{\pgfqpoint{8.065404in}{1.686367in}}%
\pgfpathlineto{\pgfqpoint{8.068399in}{1.701511in}}%
\pgfpathlineto{\pgfqpoint{8.069084in}{1.701625in}}%
\pgfpathlineto{\pgfqpoint{8.069426in}{1.701152in}}%
\pgfpathlineto{\pgfqpoint{8.070795in}{1.695972in}}%
\pgfpathlineto{\pgfqpoint{8.074046in}{1.671418in}}%
\pgfpathlineto{\pgfqpoint{8.077041in}{1.656644in}}%
\pgfpathlineto{\pgfqpoint{8.077725in}{1.656588in}}%
\pgfpathlineto{\pgfqpoint{8.078068in}{1.657083in}}%
\pgfpathlineto{\pgfqpoint{8.079437in}{1.662303in}}%
\pgfpathlineto{\pgfqpoint{8.082774in}{1.687183in}}%
\pgfpathlineto{\pgfqpoint{8.085683in}{1.700632in}}%
\pgfpathlineto{\pgfqpoint{8.086453in}{1.700426in}}%
\pgfpathlineto{\pgfqpoint{8.086624in}{1.700144in}}%
\pgfpathlineto{\pgfqpoint{8.087993in}{1.695025in}}%
\pgfpathlineto{\pgfqpoint{8.091415in}{1.669934in}}%
\pgfpathlineto{\pgfqpoint{8.094239in}{1.657511in}}%
\pgfpathlineto{\pgfqpoint{8.095009in}{1.657794in}}%
\pgfpathlineto{\pgfqpoint{8.095180in}{1.658090in}}%
\pgfpathlineto{\pgfqpoint{8.096549in}{1.663270in}}%
\pgfpathlineto{\pgfqpoint{8.100228in}{1.689775in}}%
\pgfpathlineto{\pgfqpoint{8.102795in}{1.699756in}}%
\pgfpathlineto{\pgfqpoint{8.103736in}{1.698967in}}%
\pgfpathlineto{\pgfqpoint{8.105191in}{1.693097in}}%
\pgfpathlineto{\pgfqpoint{8.111437in}{1.658357in}}%
\pgfpathlineto{\pgfqpoint{8.112122in}{1.659087in}}%
\pgfpathlineto{\pgfqpoint{8.113576in}{1.664725in}}%
\pgfpathlineto{\pgfqpoint{8.119822in}{1.698874in}}%
\pgfpathlineto{\pgfqpoint{8.120592in}{1.697974in}}%
\pgfpathlineto{\pgfqpoint{8.122132in}{1.691632in}}%
\pgfpathlineto{\pgfqpoint{8.128207in}{1.659275in}}%
\pgfpathlineto{\pgfqpoint{8.128635in}{1.659618in}}%
\pgfpathlineto{\pgfqpoint{8.129919in}{1.663542in}}%
\pgfpathlineto{\pgfqpoint{8.132742in}{1.682014in}}%
\pgfpathlineto{\pgfqpoint{8.135994in}{1.697714in}}%
\pgfpathlineto{\pgfqpoint{8.136764in}{1.697828in}}%
\pgfpathlineto{\pgfqpoint{8.137106in}{1.697357in}}%
\pgfpathlineto{\pgfqpoint{8.138475in}{1.692527in}}%
\pgfpathlineto{\pgfqpoint{8.144721in}{1.660151in}}%
\pgfpathlineto{\pgfqpoint{8.145577in}{1.661208in}}%
\pgfpathlineto{\pgfqpoint{8.147202in}{1.668073in}}%
\pgfpathlineto{\pgfqpoint{8.152935in}{1.697109in}}%
\pgfpathlineto{\pgfqpoint{8.153192in}{1.697001in}}%
\pgfpathlineto{\pgfqpoint{8.154304in}{1.694551in}}%
\pgfpathlineto{\pgfqpoint{8.156528in}{1.682218in}}%
\pgfpathlineto{\pgfqpoint{8.160721in}{1.661173in}}%
\pgfpathlineto{\pgfqpoint{8.161577in}{1.661329in}}%
\pgfpathlineto{\pgfqpoint{8.161748in}{1.661591in}}%
\pgfpathlineto{\pgfqpoint{8.163202in}{1.666630in}}%
\pgfpathlineto{\pgfqpoint{8.169277in}{1.696265in}}%
\pgfpathlineto{\pgfqpoint{8.169962in}{1.695545in}}%
\pgfpathlineto{\pgfqpoint{8.171416in}{1.690324in}}%
\pgfpathlineto{\pgfqpoint{8.177406in}{1.661842in}}%
\pgfpathlineto{\pgfqpoint{8.177919in}{1.662309in}}%
\pgfpathlineto{\pgfqpoint{8.179288in}{1.666625in}}%
\pgfpathlineto{\pgfqpoint{8.185449in}{1.695449in}}%
\pgfpathlineto{\pgfqpoint{8.186304in}{1.694344in}}%
\pgfpathlineto{\pgfqpoint{8.188015in}{1.687320in}}%
\pgfpathlineto{\pgfqpoint{8.193406in}{1.662629in}}%
\pgfpathlineto{\pgfqpoint{8.193491in}{1.662645in}}%
\pgfpathlineto{\pgfqpoint{8.194518in}{1.664228in}}%
\pgfpathlineto{\pgfqpoint{8.196401in}{1.672741in}}%
\pgfpathlineto{\pgfqpoint{8.201192in}{1.694640in}}%
\pgfpathlineto{\pgfqpoint{8.202133in}{1.693931in}}%
\pgfpathlineto{\pgfqpoint{8.203673in}{1.688512in}}%
\pgfpathlineto{\pgfqpoint{8.209406in}{1.663416in}}%
\pgfpathlineto{\pgfqpoint{8.209748in}{1.663651in}}%
\pgfpathlineto{\pgfqpoint{8.211032in}{1.666918in}}%
\pgfpathlineto{\pgfqpoint{8.213941in}{1.682862in}}%
\pgfpathlineto{\pgfqpoint{8.216936in}{1.693832in}}%
\pgfpathlineto{\pgfqpoint{8.217962in}{1.693216in}}%
\pgfpathlineto{\pgfqpoint{8.219502in}{1.688004in}}%
\pgfpathlineto{\pgfqpoint{8.225150in}{1.664157in}}%
\pgfpathlineto{\pgfqpoint{8.225492in}{1.664366in}}%
\pgfpathlineto{\pgfqpoint{8.226690in}{1.667155in}}%
\pgfpathlineto{\pgfqpoint{8.229342in}{1.680830in}}%
\pgfpathlineto{\pgfqpoint{8.232594in}{1.693076in}}%
\pgfpathlineto{\pgfqpoint{8.233620in}{1.692565in}}%
\pgfpathlineto{\pgfqpoint{8.235160in}{1.687626in}}%
\pgfpathlineto{\pgfqpoint{8.240808in}{1.664888in}}%
\pgfpathlineto{\pgfqpoint{8.241150in}{1.665118in}}%
\pgfpathlineto{\pgfqpoint{8.242433in}{1.668212in}}%
\pgfpathlineto{\pgfqpoint{8.245513in}{1.683891in}}%
\pgfpathlineto{\pgfqpoint{8.248251in}{1.692428in}}%
\pgfpathlineto{\pgfqpoint{8.249278in}{1.691690in}}%
\pgfpathlineto{\pgfqpoint{8.250904in}{1.686233in}}%
\pgfpathlineto{\pgfqpoint{8.256294in}{1.665587in}}%
\pgfpathlineto{\pgfqpoint{8.256380in}{1.665620in}}%
\pgfpathlineto{\pgfqpoint{8.257492in}{1.667466in}}%
\pgfpathlineto{\pgfqpoint{8.259717in}{1.677162in}}%
\pgfpathlineto{\pgfqpoint{8.263567in}{1.691695in}}%
\pgfpathlineto{\pgfqpoint{8.264679in}{1.691046in}}%
\pgfpathlineto{\pgfqpoint{8.266305in}{1.685806in}}%
\pgfpathlineto{\pgfqpoint{8.271610in}{1.666254in}}%
\pgfpathlineto{\pgfqpoint{8.271696in}{1.666281in}}%
\pgfpathlineto{\pgfqpoint{8.272808in}{1.668005in}}%
\pgfpathlineto{\pgfqpoint{8.274947in}{1.676835in}}%
\pgfpathlineto{\pgfqpoint{8.278797in}{1.691007in}}%
\pgfpathlineto{\pgfqpoint{8.279910in}{1.690483in}}%
\pgfpathlineto{\pgfqpoint{8.281535in}{1.685547in}}%
\pgfpathlineto{\pgfqpoint{8.286840in}{1.666910in}}%
\pgfpathlineto{\pgfqpoint{8.287011in}{1.666985in}}%
\pgfpathlineto{\pgfqpoint{8.288209in}{1.669118in}}%
\pgfpathlineto{\pgfqpoint{8.290776in}{1.680122in}}%
\pgfpathlineto{\pgfqpoint{8.294027in}{1.690422in}}%
\pgfpathlineto{\pgfqpoint{8.295140in}{1.689721in}}%
\pgfpathlineto{\pgfqpoint{8.296851in}{1.684391in}}%
\pgfpathlineto{\pgfqpoint{8.301899in}{1.667534in}}%
\pgfpathlineto{\pgfqpoint{8.303011in}{1.668968in}}%
\pgfpathlineto{\pgfqpoint{8.305150in}{1.676938in}}%
\pgfpathlineto{\pgfqpoint{8.309001in}{1.689787in}}%
\pgfpathlineto{\pgfqpoint{8.310113in}{1.689193in}}%
\pgfpathlineto{\pgfqpoint{8.311739in}{1.684506in}}%
\pgfpathlineto{\pgfqpoint{8.316873in}{1.668146in}}%
\pgfpathlineto{\pgfqpoint{8.316958in}{1.668175in}}%
\pgfpathlineto{\pgfqpoint{8.318156in}{1.669988in}}%
\pgfpathlineto{\pgfqpoint{8.320552in}{1.679183in}}%
\pgfpathlineto{\pgfqpoint{8.323889in}{1.689195in}}%
\pgfpathlineto{\pgfqpoint{8.325001in}{1.688619in}}%
\pgfpathlineto{\pgfqpoint{8.326712in}{1.683786in}}%
\pgfpathlineto{\pgfqpoint{8.331675in}{1.668726in}}%
\pgfpathlineto{\pgfqpoint{8.332787in}{1.670060in}}%
\pgfpathlineto{\pgfqpoint{8.334926in}{1.677389in}}%
\pgfpathlineto{\pgfqpoint{8.338691in}{1.688641in}}%
\pgfpathlineto{\pgfqpoint{8.339803in}{1.688010in}}%
\pgfpathlineto{\pgfqpoint{8.341515in}{1.683287in}}%
\pgfpathlineto{\pgfqpoint{8.346306in}{1.669280in}}%
\pgfpathlineto{\pgfqpoint{8.347418in}{1.670429in}}%
\pgfpathlineto{\pgfqpoint{8.349472in}{1.676955in}}%
\pgfpathlineto{\pgfqpoint{8.353408in}{1.688115in}}%
\pgfpathlineto{\pgfqpoint{8.354606in}{1.687226in}}%
\pgfpathlineto{\pgfqpoint{8.356488in}{1.681743in}}%
\pgfpathlineto{\pgfqpoint{8.360766in}{1.669828in}}%
\pgfpathlineto{\pgfqpoint{8.361964in}{1.670858in}}%
\pgfpathlineto{\pgfqpoint{8.364017in}{1.677020in}}%
\pgfpathlineto{\pgfqpoint{8.367868in}{1.687562in}}%
\pgfpathlineto{\pgfqpoint{8.369066in}{1.686833in}}%
\pgfpathlineto{\pgfqpoint{8.370948in}{1.681721in}}%
\pgfpathlineto{\pgfqpoint{8.375226in}{1.670352in}}%
\pgfpathlineto{\pgfqpoint{8.376424in}{1.671327in}}%
\pgfpathlineto{\pgfqpoint{8.378478in}{1.677195in}}%
\pgfpathlineto{\pgfqpoint{8.382328in}{1.687065in}}%
\pgfpathlineto{\pgfqpoint{8.383526in}{1.686274in}}%
\pgfpathlineto{\pgfqpoint{8.385494in}{1.681008in}}%
\pgfpathlineto{\pgfqpoint{8.389515in}{1.670868in}}%
\pgfpathlineto{\pgfqpoint{8.390713in}{1.671687in}}%
\pgfpathlineto{\pgfqpoint{8.392681in}{1.676872in}}%
\pgfpathlineto{\pgfqpoint{8.396702in}{1.686587in}}%
\pgfpathlineto{\pgfqpoint{8.397900in}{1.685695in}}%
\pgfpathlineto{\pgfqpoint{8.399954in}{1.680261in}}%
\pgfpathlineto{\pgfqpoint{8.403718in}{1.671362in}}%
\pgfpathlineto{\pgfqpoint{8.405002in}{1.672207in}}%
\pgfpathlineto{\pgfqpoint{8.407055in}{1.677496in}}%
\pgfpathlineto{\pgfqpoint{8.410820in}{1.686099in}}%
\pgfpathlineto{\pgfqpoint{8.412104in}{1.685234in}}%
\pgfpathlineto{\pgfqpoint{8.414243in}{1.679774in}}%
\pgfpathlineto{\pgfqpoint{8.417836in}{1.671831in}}%
\pgfpathlineto{\pgfqpoint{8.419120in}{1.672612in}}%
\pgfpathlineto{\pgfqpoint{8.421173in}{1.677609in}}%
\pgfpathlineto{\pgfqpoint{8.424938in}{1.685650in}}%
\pgfpathlineto{\pgfqpoint{8.426221in}{1.684761in}}%
\pgfpathlineto{\pgfqpoint{8.428446in}{1.679268in}}%
\pgfpathlineto{\pgfqpoint{8.431868in}{1.672277in}}%
\pgfpathlineto{\pgfqpoint{8.433152in}{1.673030in}}%
\pgfpathlineto{\pgfqpoint{8.435291in}{1.678022in}}%
\pgfpathlineto{\pgfqpoint{8.438885in}{1.685206in}}%
\pgfpathlineto{\pgfqpoint{8.440254in}{1.684280in}}%
\pgfpathlineto{\pgfqpoint{8.442564in}{1.678761in}}%
\pgfpathlineto{\pgfqpoint{8.445815in}{1.672700in}}%
\pgfpathlineto{\pgfqpoint{8.447184in}{1.673573in}}%
\pgfpathlineto{\pgfqpoint{8.449494in}{1.678914in}}%
\pgfpathlineto{\pgfqpoint{8.452746in}{1.684783in}}%
\pgfpathlineto{\pgfqpoint{8.454115in}{1.683921in}}%
\pgfpathlineto{\pgfqpoint{8.456510in}{1.678495in}}%
\pgfpathlineto{\pgfqpoint{8.459676in}{1.673105in}}%
\pgfpathlineto{\pgfqpoint{8.461045in}{1.673997in}}%
\pgfpathlineto{\pgfqpoint{8.463527in}{1.679521in}}%
\pgfpathlineto{\pgfqpoint{8.466607in}{1.684391in}}%
\pgfpathlineto{\pgfqpoint{8.467976in}{1.683437in}}%
\pgfpathlineto{\pgfqpoint{8.470628in}{1.677598in}}%
\pgfpathlineto{\pgfqpoint{8.473452in}{1.673495in}}%
\pgfpathlineto{\pgfqpoint{8.474906in}{1.674540in}}%
\pgfpathlineto{\pgfqpoint{8.477816in}{1.680856in}}%
\pgfpathlineto{\pgfqpoint{8.480468in}{1.683997in}}%
\pgfpathlineto{\pgfqpoint{8.481923in}{1.682701in}}%
\pgfpathlineto{\pgfqpoint{8.487741in}{1.674062in}}%
\pgfpathlineto{\pgfqpoint{8.488254in}{1.674516in}}%
\pgfpathlineto{\pgfqpoint{8.490564in}{1.678862in}}%
\pgfpathlineto{\pgfqpoint{8.493816in}{1.683630in}}%
\pgfpathlineto{\pgfqpoint{8.495270in}{1.682721in}}%
\pgfpathlineto{\pgfqpoint{8.498179in}{1.677024in}}%
\pgfpathlineto{\pgfqpoint{8.500746in}{1.674238in}}%
\pgfpathlineto{\pgfqpoint{8.502286in}{1.675481in}}%
\pgfpathlineto{\pgfqpoint{8.508105in}{1.683038in}}%
\pgfpathlineto{\pgfqpoint{8.508447in}{1.682755in}}%
\pgfpathlineto{\pgfqpoint{8.510757in}{1.678884in}}%
\pgfpathlineto{\pgfqpoint{8.514094in}{1.674582in}}%
\pgfpathlineto{\pgfqpoint{8.515720in}{1.675752in}}%
\pgfpathlineto{\pgfqpoint{8.521623in}{1.682623in}}%
\pgfpathlineto{\pgfqpoint{8.521795in}{1.682483in}}%
\pgfpathlineto{\pgfqpoint{8.524105in}{1.678911in}}%
\pgfpathlineto{\pgfqpoint{8.527442in}{1.674913in}}%
\pgfpathlineto{\pgfqpoint{8.529067in}{1.676016in}}%
\pgfpathlineto{\pgfqpoint{8.534971in}{1.682282in}}%
\pgfpathlineto{\pgfqpoint{8.535142in}{1.682146in}}%
\pgfpathlineto{\pgfqpoint{8.537624in}{1.678475in}}%
\pgfpathlineto{\pgfqpoint{8.540704in}{1.675229in}}%
\pgfpathlineto{\pgfqpoint{8.542415in}{1.676375in}}%
\pgfpathlineto{\pgfqpoint{8.548148in}{1.682017in}}%
\pgfpathlineto{\pgfqpoint{8.550458in}{1.679069in}}%
\pgfpathlineto{\pgfqpoint{8.553881in}{1.675530in}}%
\pgfpathlineto{\pgfqpoint{8.555677in}{1.676713in}}%
\pgfpathlineto{\pgfqpoint{8.561153in}{1.681810in}}%
\pgfpathlineto{\pgfqpoint{8.563378in}{1.679315in}}%
\pgfpathlineto{\pgfqpoint{8.567057in}{1.675821in}}%
\pgfpathlineto{\pgfqpoint{8.568940in}{1.677126in}}%
\pgfpathlineto{\pgfqpoint{8.574073in}{1.681599in}}%
\pgfpathlineto{\pgfqpoint{8.576298in}{1.679400in}}%
\pgfpathlineto{\pgfqpoint{8.580063in}{1.676094in}}%
\pgfpathlineto{\pgfqpoint{8.582031in}{1.677407in}}%
\pgfpathlineto{\pgfqpoint{8.586908in}{1.681387in}}%
\pgfpathlineto{\pgfqpoint{8.589132in}{1.679467in}}%
\pgfpathlineto{\pgfqpoint{8.593068in}{1.676362in}}%
\pgfpathlineto{\pgfqpoint{8.595207in}{1.677844in}}%
\pgfpathlineto{\pgfqpoint{8.599656in}{1.681175in}}%
\pgfpathlineto{\pgfqpoint{8.601881in}{1.679518in}}%
\pgfpathlineto{\pgfqpoint{8.605988in}{1.676618in}}%
\pgfpathlineto{\pgfqpoint{8.608298in}{1.678231in}}%
\pgfpathlineto{\pgfqpoint{8.612405in}{1.680950in}}%
\pgfpathlineto{\pgfqpoint{8.614801in}{1.679298in}}%
\pgfpathlineto{\pgfqpoint{8.618737in}{1.676846in}}%
\pgfpathlineto{\pgfqpoint{8.621133in}{1.678395in}}%
\pgfpathlineto{\pgfqpoint{8.625069in}{1.680739in}}%
\pgfpathlineto{\pgfqpoint{8.627550in}{1.679195in}}%
\pgfpathlineto{\pgfqpoint{8.631486in}{1.677073in}}%
\pgfpathlineto{\pgfqpoint{8.634138in}{1.678755in}}%
\pgfpathlineto{\pgfqpoint{8.637817in}{1.680511in}}%
\pgfpathlineto{\pgfqpoint{8.640555in}{1.678829in}}%
\pgfpathlineto{\pgfqpoint{8.644149in}{1.677285in}}%
\pgfpathlineto{\pgfqpoint{8.647058in}{1.679038in}}%
\pgfpathlineto{\pgfqpoint{8.650481in}{1.680298in}}%
\pgfpathlineto{\pgfqpoint{8.653646in}{1.678425in}}%
\pgfpathlineto{\pgfqpoint{8.656812in}{1.677501in}}%
\pgfpathlineto{\pgfqpoint{8.660320in}{1.679514in}}%
\pgfpathlineto{\pgfqpoint{8.663315in}{1.680031in}}%
\pgfpathlineto{\pgfqpoint{8.671358in}{1.678623in}}%
\pgfpathlineto{\pgfqpoint{8.675550in}{1.679919in}}%
\pgfpathlineto{\pgfqpoint{8.684278in}{1.678956in}}%
\pgfpathlineto{\pgfqpoint{8.688042in}{1.679736in}}%
\pgfpathlineto{\pgfqpoint{8.696256in}{1.678806in}}%
\pgfpathlineto{\pgfqpoint{8.700535in}{1.679552in}}%
\pgfpathlineto{\pgfqpoint{8.708150in}{1.678690in}}%
\pgfpathlineto{\pgfqpoint{8.712941in}{1.679390in}}%
\pgfpathlineto{\pgfqpoint{8.720214in}{1.678692in}}%
\pgfpathlineto{\pgfqpoint{8.725348in}{1.679232in}}%
\pgfpathlineto{\pgfqpoint{8.732193in}{1.678696in}}%
\pgfpathlineto{\pgfqpoint{8.737840in}{1.679065in}}%
\pgfpathlineto{\pgfqpoint{8.744257in}{1.678738in}}%
\pgfpathlineto{\pgfqpoint{8.750246in}{1.678938in}}%
\pgfpathlineto{\pgfqpoint{8.756578in}{1.678828in}}%
\pgfpathlineto{\pgfqpoint{8.762824in}{1.678813in}}%
\pgfpathlineto{\pgfqpoint{8.769327in}{1.678946in}}%
\pgfpathlineto{\pgfqpoint{8.776428in}{1.678666in}}%
\pgfpathlineto{\pgfqpoint{8.793969in}{1.678947in}}%
\pgfpathlineto{\pgfqpoint{8.829049in}{1.678861in}}%
\pgfpathlineto{\pgfqpoint{8.850354in}{1.678845in}}%
\pgfpathlineto{\pgfqpoint{8.850354in}{1.678845in}}%
\pgfusepath{stroke}%
\end{pgfscope}%
\begin{pgfscope}%
\pgfpathrectangle{\pgfqpoint{0.294194in}{0.558614in}}{\pgfqpoint{8.556160in}{2.240462in}}%
\pgfusepath{clip}%
\pgfsetbuttcap%
\pgfsetroundjoin%
\pgfsetlinewidth{1.505625pt}%
\definecolor{currentstroke}{rgb}{0.000000,0.000000,0.000000}%
\pgfsetstrokecolor{currentstroke}%
\pgfsetdash{{5.550000pt}{2.400000pt}}{0.000000pt}%
\pgfpathmoveto{\pgfqpoint{2.977154in}{0.558614in}}%
\pgfpathlineto{\pgfqpoint{2.977154in}{2.799076in}}%
\pgfusepath{stroke}%
\end{pgfscope}%
\begin{pgfscope}%
\pgfpathrectangle{\pgfqpoint{0.294194in}{0.558614in}}{\pgfqpoint{8.556160in}{2.240462in}}%
\pgfusepath{clip}%
\pgfsetbuttcap%
\pgfsetroundjoin%
\pgfsetlinewidth{1.505625pt}%
\definecolor{currentstroke}{rgb}{0.000000,0.000000,0.000000}%
\pgfsetstrokecolor{currentstroke}%
\pgfsetdash{{5.550000pt}{2.400000pt}}{0.000000pt}%
\pgfpathmoveto{\pgfqpoint{6.167395in}{0.558614in}}%
\pgfpathlineto{\pgfqpoint{6.167395in}{2.799076in}}%
\pgfusepath{stroke}%
\end{pgfscope}%
\begin{pgfscope}%
\pgfsetrectcap%
\pgfsetmiterjoin%
\pgfsetlinewidth{0.803000pt}%
\definecolor{currentstroke}{rgb}{0.000000,0.000000,0.000000}%
\pgfsetstrokecolor{currentstroke}%
\pgfsetdash{}{0pt}%
\pgfpathmoveto{\pgfqpoint{0.294194in}{0.558614in}}%
\pgfpathlineto{\pgfqpoint{0.294194in}{2.799076in}}%
\pgfusepath{stroke}%
\end{pgfscope}%
\begin{pgfscope}%
\pgfsetrectcap%
\pgfsetmiterjoin%
\pgfsetlinewidth{0.803000pt}%
\definecolor{currentstroke}{rgb}{0.000000,0.000000,0.000000}%
\pgfsetstrokecolor{currentstroke}%
\pgfsetdash{}{0pt}%
\pgfpathmoveto{\pgfqpoint{8.850354in}{0.558614in}}%
\pgfpathlineto{\pgfqpoint{8.850354in}{2.799076in}}%
\pgfusepath{stroke}%
\end{pgfscope}%
\begin{pgfscope}%
\pgfsetrectcap%
\pgfsetmiterjoin%
\pgfsetlinewidth{0.803000pt}%
\definecolor{currentstroke}{rgb}{0.000000,0.000000,0.000000}%
\pgfsetstrokecolor{currentstroke}%
\pgfsetdash{}{0pt}%
\pgfpathmoveto{\pgfqpoint{0.294194in}{0.558614in}}%
\pgfpathlineto{\pgfqpoint{8.850354in}{0.558614in}}%
\pgfusepath{stroke}%
\end{pgfscope}%
\begin{pgfscope}%
\pgfsetrectcap%
\pgfsetmiterjoin%
\pgfsetlinewidth{0.803000pt}%
\definecolor{currentstroke}{rgb}{0.000000,0.000000,0.000000}%
\pgfsetstrokecolor{currentstroke}%
\pgfsetdash{}{0pt}%
\pgfpathmoveto{\pgfqpoint{0.294194in}{2.799076in}}%
\pgfpathlineto{\pgfqpoint{8.850354in}{2.799076in}}%
\pgfusepath{stroke}%
\end{pgfscope}%
\begin{pgfscope}%
\definecolor{textcolor}{rgb}{0.000000,0.000000,0.000000}%
\pgfsetstrokecolor{textcolor}%
\pgfsetfillcolor{textcolor}%
\pgftext[x=9.278162in,y=1.678845in,left,base]{\color{textcolor}\rmfamily\fontsize{16.000000}{19.200000}\selectfont \(\displaystyle t_c<t\)}%
\end{pgfscope}%
\begin{pgfscope}%
\pgfsetbuttcap%
\pgfsetmiterjoin%
\definecolor{currentfill}{rgb}{1.000000,1.000000,1.000000}%
\pgfsetfillcolor{currentfill}%
\pgfsetfillopacity{0.800000}%
\pgfsetlinewidth{1.003750pt}%
\definecolor{currentstroke}{rgb}{0.800000,0.800000,0.800000}%
\pgfsetstrokecolor{currentstroke}%
\pgfsetstrokeopacity{0.800000}%
\pgfsetdash{}{0pt}%
\pgfpathmoveto{\pgfqpoint{8.518284in}{6.796374in}}%
\pgfpathlineto{\pgfqpoint{9.844444in}{6.796374in}}%
\pgfpathquadraticcurveto{\pgfqpoint{9.888889in}{6.796374in}}{\pgfqpoint{9.888889in}{6.840819in}}%
\pgfpathlineto{\pgfqpoint{9.888889in}{7.844444in}}%
\pgfpathquadraticcurveto{\pgfqpoint{9.888889in}{7.888889in}}{\pgfqpoint{9.844444in}{7.888889in}}%
\pgfpathlineto{\pgfqpoint{8.518284in}{7.888889in}}%
\pgfpathquadraticcurveto{\pgfqpoint{8.473839in}{7.888889in}}{\pgfqpoint{8.473839in}{7.844444in}}%
\pgfpathlineto{\pgfqpoint{8.473839in}{6.840819in}}%
\pgfpathquadraticcurveto{\pgfqpoint{8.473839in}{6.796374in}}{\pgfqpoint{8.518284in}{6.796374in}}%
\pgfpathlineto{\pgfqpoint{8.518284in}{6.796374in}}%
\pgfpathclose%
\pgfusepath{stroke,fill}%
\end{pgfscope}%
\begin{pgfscope}%
\pgfsetrectcap%
\pgfsetroundjoin%
\pgfsetlinewidth{1.505625pt}%
\definecolor{currentstroke}{rgb}{0.121569,0.466667,0.705882}%
\pgfsetstrokecolor{currentstroke}%
\pgfsetdash{}{0pt}%
\pgfpathmoveto{\pgfqpoint{8.562728in}{7.697994in}}%
\pgfpathlineto{\pgfqpoint{8.784950in}{7.697994in}}%
\pgfpathlineto{\pgfqpoint{9.007173in}{7.697994in}}%
\pgfusepath{stroke}%
\end{pgfscope}%
\begin{pgfscope}%
\definecolor{textcolor}{rgb}{0.000000,0.000000,0.000000}%
\pgfsetstrokecolor{textcolor}%
\pgfsetfillcolor{textcolor}%
\pgftext[x=9.184950in,y=7.620216in,left,base]{\color{textcolor}\rmfamily\fontsize{16.000000}{19.200000}\selectfont \(\displaystyle \Re f(k)\)}%
\end{pgfscope}%
\begin{pgfscope}%
\pgfsetrectcap%
\pgfsetroundjoin%
\pgfsetlinewidth{1.505625pt}%
\definecolor{currentstroke}{rgb}{1.000000,0.498039,0.054902}%
\pgfsetstrokecolor{currentstroke}%
\pgfsetdash{}{0pt}%
\pgfpathmoveto{\pgfqpoint{8.562728in}{7.347300in}}%
\pgfpathlineto{\pgfqpoint{8.784950in}{7.347300in}}%
\pgfpathlineto{\pgfqpoint{9.007173in}{7.347300in}}%
\pgfusepath{stroke}%
\end{pgfscope}%
\begin{pgfscope}%
\definecolor{textcolor}{rgb}{0.000000,0.000000,0.000000}%
\pgfsetstrokecolor{textcolor}%
\pgfsetfillcolor{textcolor}%
\pgftext[x=9.184950in,y=7.269522in,left,base]{\color{textcolor}\rmfamily\fontsize{16.000000}{19.200000}\selectfont \(\displaystyle \Im f(k)\)}%
\end{pgfscope}%
\begin{pgfscope}%
\pgfsetbuttcap%
\pgfsetroundjoin%
\pgfsetlinewidth{1.505625pt}%
\definecolor{currentstroke}{rgb}{0.000000,0.000000,0.000000}%
\pgfsetstrokecolor{currentstroke}%
\pgfsetdash{{5.550000pt}{2.400000pt}}{0.000000pt}%
\pgfpathmoveto{\pgfqpoint{8.562728in}{7.009723in}}%
\pgfpathlineto{\pgfqpoint{8.784950in}{7.009723in}}%
\pgfpathlineto{\pgfqpoint{9.007173in}{7.009723in}}%
\pgfusepath{stroke}%
\end{pgfscope}%
\begin{pgfscope}%
\definecolor{textcolor}{rgb}{0.000000,0.000000,0.000000}%
\pgfsetstrokecolor{textcolor}%
\pgfsetfillcolor{textcolor}%
\pgftext[x=9.184950in,y=6.931945in,left,base]{\color{textcolor}\rmfamily\fontsize{16.000000}{19.200000}\selectfont \(\displaystyle k_0\)}%
\end{pgfscope}%
\end{pgfpicture}%
\makeatother%
\endgroup%
}
    \caption*{Figure: Behavior of the integrand inside $g_N$ as a function of $k$. At $t < T_c$, there exist two saddle point, which to coalesce at $t = t_c$. It then moves into the complex plane at $t<t_c$.}
\end{figure}


\subsubsection{Derivation: Saddle point method.}
Using the condition $r=N$, $s=1$, further analytic consideration can be done with writing the amplitude back in terms of $k = \pi(m-1)/N$
\begin{align*}
    g_N
     & = \sum_{m } a_m ^ 2 \exp \left[ -i4Jt(1 - \cos k) \right] \cos (k(N-1/2)) \cos (k/2)                                     \\
     & = \sum_{m } a_m ^ 2 \exp \left[ -i4Jt(1 - \cos k) \right] \left[ \cos kN \cos k/2 + \sin kN \sin k/2 \right]  \cos (k/2) \\
     & = \sum_{m } a_m ^ 2\cos^2 k/2 \exp \left[ -i4Jt(1 - \cos k) \right](-1)^{m-1}                                            \\
     & = \sum_{m } a_m ^ 2\cos^2 k/2 \exp \left[i \left\{\pi(m-1) - 4Jt (1-\cos k)  \right\}  \right]                           \\
    g_N
     & = \sum_{m =1} ^{N} a_m ^ 2\cos^2 k_m/2 \exp \left[i \left\{Nk_m - 4Jt (1-\cos k_m)  \right\}  \right]
\end{align*}
where $kN = \pi (m-1)$, $a_1 = 1/\sqrt{N}$, $a_m = \sqrt{2/N}$ for $m>1$.
Considering that $a_m$ depend on the value of $m$, we can split the summation
\begin{equation*}
    g_N = \frac{1}{N} + \frac{2 }{N }\sum_{m=2 } ^{N} \cos^2 k/2 \exp \left[i \left\{Nk - 4Jt (1-\cos k)  \right\}  \right]
\end{equation*}

With $\Delta k = \pi /N$, we can approximate the sum as integral
\begin{equation*}
    g_N = \frac{1 }{N} +  \frac{2}{\pi}\int_{k_2} ^{k_N} dk\;  \cos^2 (k/2) \exp \left[i \left\{Nk - 4Jt (1-\cos k)  \right\}  \right]
\end{equation*}
with  $k_\text{2} = \pi/N$ and $k_N = \pi(N-1)/N$.
Under large $N$, the limit can be approximated as
\begin{equation*}
    g_N = \frac{2}{\pi} \int_{0} ^{\pi} dk\; \cos^2 (k/2) \exp \left[i \left\{Nk - 4Jt (1-\cos k)  \right\}  \right]
\end{equation*}
The  integral oscillates rapidly due to the phase $\Phi(k) $ of $e ^{i\Phi(k)}$.
Integral over infinitesimal contribution $A(k) e ^{i \Phi(k)} \;dk$ produces cancellation, except at the saddle point or stationary point $k_0$.

We determine this $k_0$ by
\begin{equation*}
    \frac{d\Phi(k )}{dk} = 0 \implies
    \sin k_0 = \frac{N }{4Jt} \implies
    k_0 = \arcsin \frac{N }{4Jt} = \arcsin \frac{t_c }{t}
\end{equation*}
From this value, we separate the integral based on the value of the arc sin argument: at $t<t_c$, $t \approx t_c$, and at $t_c<t$.
We will expand the phase around $k_0$
\begin{equation*}
    \Phi (k )  =  \Phi(k_0) +  \Phi'(k_0)(k-k_0) + \frac{1}{2}\Phi''(k_0)(k-k_0)^2 + \frac{1 }{6 }\Phi'''(k_0)(k-k_0)^3\dots
\end{equation*}
with $\Phi'(k_0) = 0$ and
\begin{align*}
    \Phi(k)    & =  Nk - 4Jt (1  - \cos k) \\
    \Phi'(k)   & =  N -4Jt\sin k           \\
    \Phi''(k)  & = - 4Jt \cos k            \\
    \Phi'''(k) & = 4Jt \sin k
\end{align*}

\subsubsection{Derivation: Complex saddle}
At $t<t_c$, we use pure saddle point approximation, where the saddle point is located at complex plane and the contour must be deformed accordingly.
This yields
\begin{multline*}
    g_N(t) = \sqrt{\frac{2}{\pi}} \left[ \frac{4Jt - i \sqrt{N^2 - (4Jt)^2}}{4Jt \left( N^2 - (4Jt)^2 \right)^{1/4}} \right]
    \exp\left[ i\left( \frac{\pi N}{2} - 4Jt \right) \right]\\
    \times \exp\left[ -N \operatorname{arcosh} \left(\frac{N}{4Jt}\right) + \sqrt{N^2 - (4Jt)^2} \right]
\end{multline*}

\subsubsection{Derivation: Airy regime}
At $t \approx t_c$, we want to use the Airy function.
Here, we not only expand the phase around $k_0 = \pi/2$, but around $t_c = N/4J$.
To do so, we use multivariable Taylor expansion
\begin{equation*}
    \Phi(k, t) \approx \sum_{m=0} \sum_{n=0} \frac{1}{m!n!} \frac{\partial^{m+n}\Phi}{\partial k^m \partial t^n}\left(k_0, t_c\right) q^m \tau^n
\end{equation*}
where $q=k-k_0 $ and $\tau = t-t_c$.
Inserting the phase $        \Phi(k)  =  Nk - 4Jt (1  - \cos k) $
\begin{align*}
    \Phi (k,t) & = \Phi(k_0,t_c) + q\partial_k\Phi  +\tau \partial_t \Phi  + \frac{1}{2} \big( q^2 \partial_k^2\Phi  + 2 q\tau \partial _k  \partial_t  \Phi  + \tau^2 \partial_t^2  \Phi\big) \\
               & \qquad + \frac{1}{6}\big(q^3\partial_k ^{3}\Phi  + 3q^2\tau \partial_k ^{2}\partial_t\Phi + 3 q\tau^2 \partial_k\partial_t^2\Phi + \tau^3 \partial_t^3\Phi \big)              \\
               & = \Phi(k_0,t_c) + q(N - 4Jt_c \sin k_0 ) + \tau \big( -4J(1-\cos k_0) \big)                                                                                                   \\
               & \qquad + \frac{1}{2} \big( q^2 (-4Jt_c \cos k_0) + 2q\tau (-4J \sin k_0) + \tau^2 (0) \big)                                                                                   \\
               & \qquad + \frac{1}{6} \big( q^3 (4Jt_c \sin k_0) + 3q^2\tau (-4J \cos k_0) + 3q\tau^2 (0) + \tau^3 (0) \big)                                                                   \\
               & = \left( \frac{\pi N}{2} - 4Jt_c \right) + q(0) - 4J\tau + \frac{1}{2} \big( q^2(0) - 8Jq\tau \big) + \frac{1}{6} \big( Nq^3 + 3q^2\tau(0) \big)                              \\
               & = \left( \frac{\pi N}{2} - 4Jt_c - 4J\tau \right) - 4J\tau q + \frac{N}{6}q^3                                                                                                 \\
               & = \left( \frac{\pi N}{2} - 4Jt \right) - 4J(t-t_c)q + \frac{N}{6}q^3
\end{align*}
Substituting this truncated  expansion back into the original integral equation yields
\begin{equation*}
    g_N(t) \approx \frac{2}{\pi} \int_{0}^{\pi} dk \; \cos^2\left( k_0/2\right) \exp \left[ i \left\{ \left( \frac{\pi N}{2} - 4Jt \right) - 4J(t-t_c)q + \frac{N}{6}q^3 \right\} \right]
\end{equation*}
Substituting $q$, factoring the integrand, and changing the integration limit
\begin{equation*}
    g_N(t) \approx \frac{1}{\pi} \exp \left( i\frac{\pi N}{2} -i 4Jt \right) \int_{-\infty}^{\infty} dq \;  \exp \left[ i \left\{\frac{N}{6}q^3 +\gamma_q q  \right\} \right]
\end{equation*}
with $\gamma_q = -4J(t-t_c)$ as the linear term
Introduce the transformation
\begin{equation*}
    \frac{1}{3}p^3 = \frac{N}{6}q^3 \quad dq = \left(\frac{2}{N}\right)^{1/3}dp
    \quad -4J(t-t_c)q = -4J(t-t_c)\left(\frac{2}{N}\right)^{1/3}p
\end{equation*}
into the integral
\begin{equation*}
    g_N(t) \approx \frac{1}{\pi} \left(\frac{2}{N}\right)^{1/3} \exp \left[ i \left( \frac{\pi N}{2} - 4Jt \right) \right] \int_{-\infty}^{\infty} dp \; \exp \left[ i \left( \frac{1}{3}p^3 + xp \right) \right]
\end{equation*}
with $\gamma_p = -4J(t-t_c)(2/N) ^{1/3}$.
Using the definition of Airy function
\begin{equation*}
    \operatorname{Ai} (x) = \frac{1 }{2\pi} \int_{-\infty }^{\infty } dp\; \exp \left( \frac{p ^3}{3 } + xp \right)
\end{equation*}
we write
\begin{equation*}
    g_N(t)= 2\left(\frac{2}{N}\right)^{1/3} \exp \left[ i \left( \frac{\pi N}{2} - 4Jt \right) \right] \operatorname{Ai}\left[ 4J(t_c - t) \left( \frac{2 }{N }\right) ^{1/3}   \right]
\end{equation*}

\subsubsection{Derivation: Gaussian regime}
At $t_c<t$, there are two possible saddle point
\begin{equation*}
    k_1 = \arcsin \frac{N }{4J}\qquad
    k_2 = \pi -k_1
\end{equation*}
Recall that horizontal line $|y(x)|<1$ touch the sine curve twice at the domain $[0,2\pi]$.
The phase $\Phi(k) =  Nk - 4Jt (1  - \cos k) $ expansion around $k_1$ is
\begin{align*}
    \Phi(k) & = \Phi(k_1) + \frac{1}{2}\Phi''(k_1) q_1^2                      \\
    \Phi(k) & =  \left[ N k_1 - 4Jt(1- \cos k_1) \right] - 2Jt \cos k_1 q_1^2
\end{align*}
and like wise for the $k_2$
\begin{equation*}
    \Phi(k)  =  \left[ N k_2 - 4Jt(1- \cos k_2) \right] - 2Jt \cos k_2 q_2^2
\end{equation*}
The integral splits after substituting this expansion
\begin{multline*}
    g_N = \frac{2}{\pi} \cos^2 (k_0/2) e ^{i\Phi(k_1)}\int_{-\infty} ^{\infty} dq_1\; \exp \left[ - i 2 Jt \cos k_1 q_1^2\right] \\
    + \frac{2}{\pi} \cos^2 (k_0/2) e ^{i\Phi(k_2)}\int_{-\infty} ^{\infty} dq_2\; \exp \left[ - i 2 Jt \cos k_2 q_2^2\right]
\end{multline*}
where we have also expanded the integration limit and using $q$ definition for both integral.
Evaluating the Fresnel integral
\begin{multline*}
    g_N = \frac{2}{\pi} \cos^2 (k_1/2) e ^{i\Phi(k_1)} \sqrt{\frac{\pi }{2Jt |\cos k_1|}} \exp \left( -i \frac{\pi }{4} \operatorname{sgn} (\cos k_1) \right)  \\
    + \frac{2}{\pi} \cos^2 (k_2/2) e ^{i\Phi(k_2)} \sqrt{\frac{\pi }{2Jt |\cos k_2|}} \exp \left( -i \frac{\pi }{4} \operatorname{sgn} (\cos k_2) \right)  \\
\end{multline*}
Now $k_1 \in (0,\pi/2)$ and $k_2 \in (\pi/2,\pi)$, so
\begin{multline*}
    g_N = \sqrt{\frac{2 \cos^4 (k_1/2)}{\pi Jt |\cos k_1|}} \exp \left[ i\Phi(k_1) -i \frac{\pi }{4}\right]  \\
    + \sqrt{\frac{2 \cos^4 (k_1/2)}{\pi Jt |\cos k_2|}} \exp \left[i\Phi(k_2) +i \frac{\pi }{4}  \right]
\end{multline*}

\subsubsection{Bessel function representation.}
Alternatively, starting from the integral
\begin{equation*}
    g_N = \frac{2 }{\pi} \int_{0} ^{\pi} dk\; \cos^2 (k/2) \exp \left[i \left\{Nk - 4Jt (1-\cos k)  \right\}  \right]
\end{equation*}
Using the Jacobi-Anger expansion
\begin{equation*}
    e^{i z \cos\theta} = \sum_{m=-\infty}^{\infty} i^{m} J_{m}(z)e^{i m \theta}
    = J_0(4Jt) + 2 \sum_{m=1}^{\infty} i^m J_m(4Jt) \cos(mk)
\end{equation*}
we can write
\begin{align*}
    g_N
     & = \frac{e^{-i4Jt}}{\pi}\int_{0}^{\pi} dk\;(1+\cos k) e^{iNk}\sum_{m=-\infty}^{\infty} i^{m} J_{m}(4Jt)e^{i m k}                   \\
     & = \frac{e^{-i4Jt}}{\pi} \sum_{m=-\infty}^{\infty} i^{m} J_{m}(4Jt) \int_{0}^{\pi} dk\; \big( e ^{i(N+m)k}                         \\
     & \qquad + \frac{1}{2}e ^{i(N+m+1)k} + \frac{1 }{2 } e ^{i(N+m-1)k}  \big)                                                          \\
     & = \frac{e^{-i4Jt}}{\pi} \left[ i ^{-N}\pi J_{-N}+ i ^{-N-1} \frac{\pi }{2 } J_{-N-1}+  i ^{-N+1} \frac{\pi }{2 } J_{-N+1} \right] \\
     & = e^{-i4Jt} i ^{-N} \left[ J_{-N} -  \frac{i }{2 } J_{-N+1} + \frac{i }{2 } J_{-N+1} \right]                                      \\
     & = e^{-i4Jt} i ^{N} \left[(-1)^{-N} J_{-N} -  \frac{i }{2 }(-1)^{-N} J_{-N+1} + \frac{i }{2 }(-1)^{-N} J_{-N+1} \right]            \\
     & = e^{-i4Jt} i ^{N} \left[ J_{N} -  \frac{i }{2 } J_{N-1} + \frac{i }{2 }J_{N-1} \right]                                           \\
    g_N(t)
     & = e^{i(N\pi/2 -4Jt)}  \left[ J_{N}(4Jt) -i  J_N' (4Jt)\right]
\end{align*}

\subsection{Optimization For XXX Hamiltonian: Dual Hole Method}
We start by considering the transition amplitude
\begin{equation*}
  f(t) = \sum_{j=1}^{N} e^{-i E_j t}\braket{\mathbf{r}|\tilde{m}_j }\braket{\tilde{m}_j |\mathbf{s} }
\end{equation*}
Defining the projection of eigenstates $\ket{\tilde{m}_j}$ to sender as $|\sigma_j| e ^{\phi_j}$ and to receiver as $|\rho_j| e ^{\psi_j}$, we have
\begin{equation*}
  f(t) = \sum_{j=1}^{N} e^{-i E_j t}\rho_j \sigma_j ^*
\end{equation*}
Then
\begin{align*}
  |f(t)|^2 
  &= \left(\sum_j e^{-iE_j t}\rho_j\sigma_j^*\right)\left(\sum_l e^{iE_l t}\rho_l^*\sigma_l\right) 
  =  \sum_{j,l} e^{-i(E_j-E_l)t}\rho_j\rho_l^*\sigma_j^*\sigma_l\\
  & = \sum_{j,l}|\rho_j||\rho_l||\sigma_j||\sigma_l|e^{-i(E_j-E_l)t} e^{i[(\psi_j-\psi_l)-(\phi_j-\phi_l)]}\\
  & = \sum_{j=1}^N |\sigma_j|^2 |\rho_j|^2 + 2\sum_{k<l} |\sigma_k||\sigma_l||\rho_k||\rho_l|\\
  &\qquad\times\cos\big([(E_l-E_k)t + (\psi_k-\psi_l) - (\phi_k-\phi_l)]\big) = f_m + f_t
\end{align*}
Since $|f(t)|^2\in [0,1]$ and $f_t$ oscillates rapidly, the allowed minimum value for $f_t$ is $-f_m$.
This implies that $2f_m$ is an accurate estimate for the maximum value of $|f(t)|^2$.
Hence, we investigate when
\begin{equation*}
    f_m = \sum_{j=1}^N |\sigma_j|^2 |\rho_j|^2
\end{equation*}
reaches its maximum value.
This search is bounded by the few normalization constraints.
First we write
\begin{align*}
    \braket{\mathbf{s} | \mathbf{s}} 
    & =  \sum_{j } \braket{\mathbf{s }  | \tilde{m}_j } \braket{\tilde{m}_j  | \mathbf{s}}
    = \sum_{j } |\sigma_j|^2 = 1\\
    \braket{\mathbf{r} | \mathbf{r}} 
    & =  \sum_{j } \braket{\mathbf{r }  | \tilde{m}_j } \braket{\tilde{m}_j  | \mathbf{r}}
    = \sum_{j } |\rho_j|^2 = 1 
\end{align*}
For fixed $j$, we can also expand in terms of other one-excitation state
\begin{equation*}
    \braket{\tilde{m}_j  | \tilde{m}_j } 
    = \sum_{ \mathbf{i} }  \braket{\tilde{m}_j  | \mathbf{i } } \braket{\mathbf{i }  | \tilde{m}_j }
    = \sum_{\mathbf{i} \neq (\mathbf{s},\mathbf{r})}  |\braket{\tilde{m}_j  | \mathbf{i }}|^2 + |\sigma_j|^2 + |\rho_j|^2 = 1
\end{equation*}
With the sum interior overlap as $|\gamma_j|^2$ (first term), we can write for fixed $j$
\begin{align*}
    |\sigma_j|^2|\rho_j|^2 = |\sigma_j|^2 \left( 1 - |\sigma_j|^2 -  |\gamma_j|^2\right) 
    = \left( 1 - |\rho_j|^2 -  |\gamma_j|^2\right) |\rho_j|^2 
\end{align*}
Setting the derivative with respect to $|\sigma_j|^2$ to zero 
\begin{align*}
    \frac{d }{d|\sigma_j|^2} |\sigma_j|^2|\rho_j|^2 & =  1 - 2|\sigma_j|^2 -|\gamma_j|^2 = 0\\
    \frac{d }{d|\rho_j|^2} |\sigma_j|^2|\rho_j|^2 & =  1 - 2|\rho_j|^2 -|\gamma_j|^2 = 0
\end{align*}
we find that 
\begin{equation*}
    |\rho_j|^2 = |\sigma_j|^2 = \frac{1 - |\gamma_j |^2}{2}
\end{equation*}
Suppose that the eigenstate is entirely localized in the boundary, which imply $|\gamma_j|^2=0$, we have then 
\begin{equation*}
   |\rho_j|^2 = |\sigma_j|^2 = \frac{1 }{2 } 
\end{equation*}
Explicitly, our eigenstates is in the form 
\begin{equation*}
    \ket{\tilde{m}_\pm} = \frac{1 }{\sqrt{2}} \left( \ket{\mathbf{s }} \pm \ket{\mathbf{r}} \right) 
\end{equation*}
In this case, we have unitary amplitude
\begin{align*}
    f(t) & =  
    e ^{-i E_+ t} \braket{\mathbf{r } |\tilde{m}_+ } \braket{\tilde{m}_+  | \mathbf{s}} +
    e ^{-i E_- t} \braket{\mathbf{r } |\tilde{m}_- } \braket{\tilde{m}_-  | \mathbf{s}}\\
    &= \frac{1 }{2 }e ^{-i E_+ t} - \frac{1 }{2} e ^{-i E_- t} = \frac{1}{2} e^{-i \bar{E} t}\left(e^{-i\Delta E t/2}-e^{i\Delta E t/2}\right)\\
    & = - i e ^{-i \bar{E} t}\sin \left( \Delta E t /2 \right) 
\end{align*}
with $\bar{E} = (E_++E_-)/2$ and $\Delta E = E_+ - E_-$.
Hence, the fidelity
\begin{equation*}
    F(t) = \frac{1 }{6 } \sin^2 \left( \Delta E t /2 \right) + \frac{1 }{3 }\sin \left( \Delta E t /2 \right) + \frac{1 }{2 }
\end{equation*}
will reach maximum value at $t_m = \pi/\Delta E$.
Dual hole method is then applied by setting the coupling constants $J_{2,i}$ and $J_{N-1,i}$ to zero.
Sender and receiver are still found at the two ends of the chain, but now their nearest neighboring sites are empty.

The eigenvalue is obtained by considering the matrix element which is restricted to $\left\{ \ket{1},\ket{N} \right\} $
\begin{equation*}
    H_\text{eff} = 
    \begin{bmatrix}
        2\sum_{k \neq 1}  J_{1k}  - 2B & -2 J_{1N}\\
        -2 J_{N1} & 2\sum_{k \neq N}  J_{Nk} - 2B
    \end{bmatrix} = 
    \begin{bmatrix}
        D_1 & A\\
        A & D_N
    \end{bmatrix}
\end{equation*}
The diagonal terms are equal, which can be seen from 
\begin{equation*}
    \sum_{k \neq 1}  J_{1k} = \sum_{k=2}^{N}\frac{C}{\bigl(a|k-1|\bigr)^{\nu}} = \sum_{m=1}^{N-1}\frac{1}{(am)^{\nu}}
\end{equation*}
and 
\begin{equation*}
    \sum_{k \neq N}  J_{Nk} = \sum_{k=1}^{N-1}\frac{C}{\bigl(a|k-N|\bigr)^{\nu}} = \sum_{m=1}^{N-1}\frac{1}{(am)^{\nu}}
\end{equation*}
The eigenvalue can be determined from characteristic equation
\begin{align*}
    (D_N-E)^2-A^2 & = 0 \\
    E_\pm = D \pm |A| & =  3 \sum_{k<l} J_{kl} - 4\sum_{k \neq 1}  J_{1k} \pm \frac{2C}{a^\nu(N-1)^\nu } 
\end{align*}
Shifting in the asymmetric spectrum, we have 
\begin{equation*}
    E_+ = 0 \qquad E_- = - \frac{4 C}{a^\nu(N-1)^{\nu}} \quad t = \frac{\pi }{4C }a^\nu(N-1 )^\nu
\end{equation*}



\begin{figure}
    \centering
    \resizebox{\textwidth}{!}{%% Creator: Matplotlib, PGF backend
%%
%% To include the figure in your LaTeX document, write
%%   \input{<filename>.pgf}
%%
%% Make sure the required packages are loaded in your preamble
%%   \usepackage{pgf}
%%
%% Also ensure that all the required font packages are loaded; for instance,
%% the lmodern package is sometimes necessary when using math font.
%%   \usepackage{lmodern}
%%
%% Figures using additional raster images can only be included by \input if
%% they are in the same directory as the main LaTeX file. For loading figures
%% from other directories you can use the `import` package
%%   \usepackage{import}
%%
%% and then include the figures with
%%   \import{<path to file>}{<filename>.pgf}
%%
%% Matplotlib used the following preamble
%%   
%%   \makeatletter\@ifpackageloaded{underscore}{}{\usepackage[strings]{underscore}}\makeatother
%%
\begingroup%
\makeatletter%
\begin{pgfpicture}%
\pgfpathrectangle{\pgfpointorigin}{\pgfqpoint{10.000000in}{6.000000in}}%
\pgfusepath{use as bounding box, clip}%
\begin{pgfscope}%
\pgfsetbuttcap%
\pgfsetmiterjoin%
\definecolor{currentfill}{rgb}{1.000000,1.000000,1.000000}%
\pgfsetfillcolor{currentfill}%
\pgfsetlinewidth{0.000000pt}%
\definecolor{currentstroke}{rgb}{1.000000,1.000000,1.000000}%
\pgfsetstrokecolor{currentstroke}%
\pgfsetdash{}{0pt}%
\pgfpathmoveto{\pgfqpoint{0.000000in}{0.000000in}}%
\pgfpathlineto{\pgfqpoint{10.000000in}{0.000000in}}%
\pgfpathlineto{\pgfqpoint{10.000000in}{6.000000in}}%
\pgfpathlineto{\pgfqpoint{0.000000in}{6.000000in}}%
\pgfpathlineto{\pgfqpoint{0.000000in}{0.000000in}}%
\pgfpathclose%
\pgfusepath{fill}%
\end{pgfscope}%
\begin{pgfscope}%
\pgfsetbuttcap%
\pgfsetmiterjoin%
\definecolor{currentfill}{rgb}{1.000000,1.000000,1.000000}%
\pgfsetfillcolor{currentfill}%
\pgfsetlinewidth{0.000000pt}%
\definecolor{currentstroke}{rgb}{0.000000,0.000000,0.000000}%
\pgfsetstrokecolor{currentstroke}%
\pgfsetstrokeopacity{0.000000}%
\pgfsetdash{}{0pt}%
\pgfpathmoveto{\pgfqpoint{1.250000in}{0.660000in}}%
\pgfpathlineto{\pgfqpoint{9.000000in}{0.660000in}}%
\pgfpathlineto{\pgfqpoint{9.000000in}{5.280000in}}%
\pgfpathlineto{\pgfqpoint{1.250000in}{5.280000in}}%
\pgfpathlineto{\pgfqpoint{1.250000in}{0.660000in}}%
\pgfpathclose%
\pgfusepath{fill}%
\end{pgfscope}%
\begin{pgfscope}%
\pgfsetbuttcap%
\pgfsetroundjoin%
\definecolor{currentfill}{rgb}{0.000000,0.000000,0.000000}%
\pgfsetfillcolor{currentfill}%
\pgfsetlinewidth{0.803000pt}%
\definecolor{currentstroke}{rgb}{0.000000,0.000000,0.000000}%
\pgfsetstrokecolor{currentstroke}%
\pgfsetdash{}{0pt}%
\pgfsys@defobject{currentmarker}{\pgfqpoint{0.000000in}{-0.048611in}}{\pgfqpoint{0.000000in}{0.000000in}}{%
\pgfpathmoveto{\pgfqpoint{0.000000in}{0.000000in}}%
\pgfpathlineto{\pgfqpoint{0.000000in}{-0.048611in}}%
\pgfusepath{stroke,fill}%
}%
\begin{pgfscope}%
\pgfsys@transformshift{1.250000in}{0.660000in}%
\pgfsys@useobject{currentmarker}{}%
\end{pgfscope}%
\end{pgfscope}%
\begin{pgfscope}%
\definecolor{textcolor}{rgb}{0.000000,0.000000,0.000000}%
\pgfsetstrokecolor{textcolor}%
\pgfsetfillcolor{textcolor}%
\pgftext[x=1.250000in,y=0.562778in,,top]{\color{textcolor}\rmfamily\fontsize{16.000000}{19.200000}\selectfont \(\displaystyle {0}\)}%
\end{pgfscope}%
\begin{pgfscope}%
\pgfsetbuttcap%
\pgfsetroundjoin%
\definecolor{currentfill}{rgb}{0.000000,0.000000,0.000000}%
\pgfsetfillcolor{currentfill}%
\pgfsetlinewidth{0.803000pt}%
\definecolor{currentstroke}{rgb}{0.000000,0.000000,0.000000}%
\pgfsetstrokecolor{currentstroke}%
\pgfsetdash{}{0pt}%
\pgfsys@defobject{currentmarker}{\pgfqpoint{0.000000in}{-0.048611in}}{\pgfqpoint{0.000000in}{0.000000in}}{%
\pgfpathmoveto{\pgfqpoint{0.000000in}{0.000000in}}%
\pgfpathlineto{\pgfqpoint{0.000000in}{-0.048611in}}%
\pgfusepath{stroke,fill}%
}%
\begin{pgfscope}%
\pgfsys@transformshift{2.800000in}{0.660000in}%
\pgfsys@useobject{currentmarker}{}%
\end{pgfscope}%
\end{pgfscope}%
\begin{pgfscope}%
\definecolor{textcolor}{rgb}{0.000000,0.000000,0.000000}%
\pgfsetstrokecolor{textcolor}%
\pgfsetfillcolor{textcolor}%
\pgftext[x=2.800000in,y=0.562778in,,top]{\color{textcolor}\rmfamily\fontsize{16.000000}{19.200000}\selectfont \(\displaystyle {20}\)}%
\end{pgfscope}%
\begin{pgfscope}%
\pgfsetbuttcap%
\pgfsetroundjoin%
\definecolor{currentfill}{rgb}{0.000000,0.000000,0.000000}%
\pgfsetfillcolor{currentfill}%
\pgfsetlinewidth{0.803000pt}%
\definecolor{currentstroke}{rgb}{0.000000,0.000000,0.000000}%
\pgfsetstrokecolor{currentstroke}%
\pgfsetdash{}{0pt}%
\pgfsys@defobject{currentmarker}{\pgfqpoint{0.000000in}{-0.048611in}}{\pgfqpoint{0.000000in}{0.000000in}}{%
\pgfpathmoveto{\pgfqpoint{0.000000in}{0.000000in}}%
\pgfpathlineto{\pgfqpoint{0.000000in}{-0.048611in}}%
\pgfusepath{stroke,fill}%
}%
\begin{pgfscope}%
\pgfsys@transformshift{4.350000in}{0.660000in}%
\pgfsys@useobject{currentmarker}{}%
\end{pgfscope}%
\end{pgfscope}%
\begin{pgfscope}%
\definecolor{textcolor}{rgb}{0.000000,0.000000,0.000000}%
\pgfsetstrokecolor{textcolor}%
\pgfsetfillcolor{textcolor}%
\pgftext[x=4.350000in,y=0.562778in,,top]{\color{textcolor}\rmfamily\fontsize{16.000000}{19.200000}\selectfont \(\displaystyle {40}\)}%
\end{pgfscope}%
\begin{pgfscope}%
\pgfsetbuttcap%
\pgfsetroundjoin%
\definecolor{currentfill}{rgb}{0.000000,0.000000,0.000000}%
\pgfsetfillcolor{currentfill}%
\pgfsetlinewidth{0.803000pt}%
\definecolor{currentstroke}{rgb}{0.000000,0.000000,0.000000}%
\pgfsetstrokecolor{currentstroke}%
\pgfsetdash{}{0pt}%
\pgfsys@defobject{currentmarker}{\pgfqpoint{0.000000in}{-0.048611in}}{\pgfqpoint{0.000000in}{0.000000in}}{%
\pgfpathmoveto{\pgfqpoint{0.000000in}{0.000000in}}%
\pgfpathlineto{\pgfqpoint{0.000000in}{-0.048611in}}%
\pgfusepath{stroke,fill}%
}%
\begin{pgfscope}%
\pgfsys@transformshift{5.900000in}{0.660000in}%
\pgfsys@useobject{currentmarker}{}%
\end{pgfscope}%
\end{pgfscope}%
\begin{pgfscope}%
\definecolor{textcolor}{rgb}{0.000000,0.000000,0.000000}%
\pgfsetstrokecolor{textcolor}%
\pgfsetfillcolor{textcolor}%
\pgftext[x=5.900000in,y=0.562778in,,top]{\color{textcolor}\rmfamily\fontsize{16.000000}{19.200000}\selectfont \(\displaystyle {60}\)}%
\end{pgfscope}%
\begin{pgfscope}%
\pgfsetbuttcap%
\pgfsetroundjoin%
\definecolor{currentfill}{rgb}{0.000000,0.000000,0.000000}%
\pgfsetfillcolor{currentfill}%
\pgfsetlinewidth{0.803000pt}%
\definecolor{currentstroke}{rgb}{0.000000,0.000000,0.000000}%
\pgfsetstrokecolor{currentstroke}%
\pgfsetdash{}{0pt}%
\pgfsys@defobject{currentmarker}{\pgfqpoint{0.000000in}{-0.048611in}}{\pgfqpoint{0.000000in}{0.000000in}}{%
\pgfpathmoveto{\pgfqpoint{0.000000in}{0.000000in}}%
\pgfpathlineto{\pgfqpoint{0.000000in}{-0.048611in}}%
\pgfusepath{stroke,fill}%
}%
\begin{pgfscope}%
\pgfsys@transformshift{7.450000in}{0.660000in}%
\pgfsys@useobject{currentmarker}{}%
\end{pgfscope}%
\end{pgfscope}%
\begin{pgfscope}%
\definecolor{textcolor}{rgb}{0.000000,0.000000,0.000000}%
\pgfsetstrokecolor{textcolor}%
\pgfsetfillcolor{textcolor}%
\pgftext[x=7.450000in,y=0.562778in,,top]{\color{textcolor}\rmfamily\fontsize{16.000000}{19.200000}\selectfont \(\displaystyle {80}\)}%
\end{pgfscope}%
\begin{pgfscope}%
\pgfsetbuttcap%
\pgfsetroundjoin%
\definecolor{currentfill}{rgb}{0.000000,0.000000,0.000000}%
\pgfsetfillcolor{currentfill}%
\pgfsetlinewidth{0.803000pt}%
\definecolor{currentstroke}{rgb}{0.000000,0.000000,0.000000}%
\pgfsetstrokecolor{currentstroke}%
\pgfsetdash{}{0pt}%
\pgfsys@defobject{currentmarker}{\pgfqpoint{0.000000in}{-0.048611in}}{\pgfqpoint{0.000000in}{0.000000in}}{%
\pgfpathmoveto{\pgfqpoint{0.000000in}{0.000000in}}%
\pgfpathlineto{\pgfqpoint{0.000000in}{-0.048611in}}%
\pgfusepath{stroke,fill}%
}%
\begin{pgfscope}%
\pgfsys@transformshift{9.000000in}{0.660000in}%
\pgfsys@useobject{currentmarker}{}%
\end{pgfscope}%
\end{pgfscope}%
\begin{pgfscope}%
\definecolor{textcolor}{rgb}{0.000000,0.000000,0.000000}%
\pgfsetstrokecolor{textcolor}%
\pgfsetfillcolor{textcolor}%
\pgftext[x=9.000000in,y=0.562778in,,top]{\color{textcolor}\rmfamily\fontsize{16.000000}{19.200000}\selectfont \(\displaystyle {100}\)}%
\end{pgfscope}%
\begin{pgfscope}%
\definecolor{textcolor}{rgb}{0.000000,0.000000,0.000000}%
\pgfsetstrokecolor{textcolor}%
\pgfsetfillcolor{textcolor}%
\pgftext[x=5.125000in,y=0.293874in,,top]{\color{textcolor}\rmfamily\fontsize{16.000000}{19.200000}\selectfont \(\displaystyle N\)}%
\end{pgfscope}%
\begin{pgfscope}%
\pgfsetbuttcap%
\pgfsetroundjoin%
\definecolor{currentfill}{rgb}{0.000000,0.000000,0.000000}%
\pgfsetfillcolor{currentfill}%
\pgfsetlinewidth{0.803000pt}%
\definecolor{currentstroke}{rgb}{0.000000,0.000000,0.000000}%
\pgfsetstrokecolor{currentstroke}%
\pgfsetdash{}{0pt}%
\pgfsys@defobject{currentmarker}{\pgfqpoint{-0.048611in}{0.000000in}}{\pgfqpoint{-0.000000in}{0.000000in}}{%
\pgfpathmoveto{\pgfqpoint{-0.000000in}{0.000000in}}%
\pgfpathlineto{\pgfqpoint{-0.048611in}{0.000000in}}%
\pgfusepath{stroke,fill}%
}%
\begin{pgfscope}%
\pgfsys@transformshift{1.250000in}{0.870000in}%
\pgfsys@useobject{currentmarker}{}%
\end{pgfscope}%
\end{pgfscope}%
\begin{pgfscope}%
\definecolor{textcolor}{rgb}{0.000000,0.000000,0.000000}%
\pgfsetstrokecolor{textcolor}%
\pgfsetfillcolor{textcolor}%
\pgftext[x=0.867364in, y=0.786667in, left, base]{\color{textcolor}\rmfamily\fontsize{16.000000}{19.200000}\selectfont \(\displaystyle {0.5}\)}%
\end{pgfscope}%
\begin{pgfscope}%
\pgfsetbuttcap%
\pgfsetroundjoin%
\definecolor{currentfill}{rgb}{0.000000,0.000000,0.000000}%
\pgfsetfillcolor{currentfill}%
\pgfsetlinewidth{0.803000pt}%
\definecolor{currentstroke}{rgb}{0.000000,0.000000,0.000000}%
\pgfsetstrokecolor{currentstroke}%
\pgfsetdash{}{0pt}%
\pgfsys@defobject{currentmarker}{\pgfqpoint{-0.048611in}{0.000000in}}{\pgfqpoint{-0.000000in}{0.000000in}}{%
\pgfpathmoveto{\pgfqpoint{-0.000000in}{0.000000in}}%
\pgfpathlineto{\pgfqpoint{-0.048611in}{0.000000in}}%
\pgfusepath{stroke,fill}%
}%
\begin{pgfscope}%
\pgfsys@transformshift{1.250000in}{1.710000in}%
\pgfsys@useobject{currentmarker}{}%
\end{pgfscope}%
\end{pgfscope}%
\begin{pgfscope}%
\definecolor{textcolor}{rgb}{0.000000,0.000000,0.000000}%
\pgfsetstrokecolor{textcolor}%
\pgfsetfillcolor{textcolor}%
\pgftext[x=0.867364in, y=1.626667in, left, base]{\color{textcolor}\rmfamily\fontsize{16.000000}{19.200000}\selectfont \(\displaystyle {0.6}\)}%
\end{pgfscope}%
\begin{pgfscope}%
\pgfsetbuttcap%
\pgfsetroundjoin%
\definecolor{currentfill}{rgb}{0.000000,0.000000,0.000000}%
\pgfsetfillcolor{currentfill}%
\pgfsetlinewidth{0.803000pt}%
\definecolor{currentstroke}{rgb}{0.000000,0.000000,0.000000}%
\pgfsetstrokecolor{currentstroke}%
\pgfsetdash{}{0pt}%
\pgfsys@defobject{currentmarker}{\pgfqpoint{-0.048611in}{0.000000in}}{\pgfqpoint{-0.000000in}{0.000000in}}{%
\pgfpathmoveto{\pgfqpoint{-0.000000in}{0.000000in}}%
\pgfpathlineto{\pgfqpoint{-0.048611in}{0.000000in}}%
\pgfusepath{stroke,fill}%
}%
\begin{pgfscope}%
\pgfsys@transformshift{1.250000in}{2.550000in}%
\pgfsys@useobject{currentmarker}{}%
\end{pgfscope}%
\end{pgfscope}%
\begin{pgfscope}%
\definecolor{textcolor}{rgb}{0.000000,0.000000,0.000000}%
\pgfsetstrokecolor{textcolor}%
\pgfsetfillcolor{textcolor}%
\pgftext[x=0.867364in, y=2.466667in, left, base]{\color{textcolor}\rmfamily\fontsize{16.000000}{19.200000}\selectfont \(\displaystyle {0.7}\)}%
\end{pgfscope}%
\begin{pgfscope}%
\pgfsetbuttcap%
\pgfsetroundjoin%
\definecolor{currentfill}{rgb}{0.000000,0.000000,0.000000}%
\pgfsetfillcolor{currentfill}%
\pgfsetlinewidth{0.803000pt}%
\definecolor{currentstroke}{rgb}{0.000000,0.000000,0.000000}%
\pgfsetstrokecolor{currentstroke}%
\pgfsetdash{}{0pt}%
\pgfsys@defobject{currentmarker}{\pgfqpoint{-0.048611in}{0.000000in}}{\pgfqpoint{-0.000000in}{0.000000in}}{%
\pgfpathmoveto{\pgfqpoint{-0.000000in}{0.000000in}}%
\pgfpathlineto{\pgfqpoint{-0.048611in}{0.000000in}}%
\pgfusepath{stroke,fill}%
}%
\begin{pgfscope}%
\pgfsys@transformshift{1.250000in}{3.390000in}%
\pgfsys@useobject{currentmarker}{}%
\end{pgfscope}%
\end{pgfscope}%
\begin{pgfscope}%
\definecolor{textcolor}{rgb}{0.000000,0.000000,0.000000}%
\pgfsetstrokecolor{textcolor}%
\pgfsetfillcolor{textcolor}%
\pgftext[x=0.867364in, y=3.306667in, left, base]{\color{textcolor}\rmfamily\fontsize{16.000000}{19.200000}\selectfont \(\displaystyle {0.8}\)}%
\end{pgfscope}%
\begin{pgfscope}%
\pgfsetbuttcap%
\pgfsetroundjoin%
\definecolor{currentfill}{rgb}{0.000000,0.000000,0.000000}%
\pgfsetfillcolor{currentfill}%
\pgfsetlinewidth{0.803000pt}%
\definecolor{currentstroke}{rgb}{0.000000,0.000000,0.000000}%
\pgfsetstrokecolor{currentstroke}%
\pgfsetdash{}{0pt}%
\pgfsys@defobject{currentmarker}{\pgfqpoint{-0.048611in}{0.000000in}}{\pgfqpoint{-0.000000in}{0.000000in}}{%
\pgfpathmoveto{\pgfqpoint{-0.000000in}{0.000000in}}%
\pgfpathlineto{\pgfqpoint{-0.048611in}{0.000000in}}%
\pgfusepath{stroke,fill}%
}%
\begin{pgfscope}%
\pgfsys@transformshift{1.250000in}{4.230000in}%
\pgfsys@useobject{currentmarker}{}%
\end{pgfscope}%
\end{pgfscope}%
\begin{pgfscope}%
\definecolor{textcolor}{rgb}{0.000000,0.000000,0.000000}%
\pgfsetstrokecolor{textcolor}%
\pgfsetfillcolor{textcolor}%
\pgftext[x=0.867364in, y=4.146667in, left, base]{\color{textcolor}\rmfamily\fontsize{16.000000}{19.200000}\selectfont \(\displaystyle {0.9}\)}%
\end{pgfscope}%
\begin{pgfscope}%
\pgfsetbuttcap%
\pgfsetroundjoin%
\definecolor{currentfill}{rgb}{0.000000,0.000000,0.000000}%
\pgfsetfillcolor{currentfill}%
\pgfsetlinewidth{0.803000pt}%
\definecolor{currentstroke}{rgb}{0.000000,0.000000,0.000000}%
\pgfsetstrokecolor{currentstroke}%
\pgfsetdash{}{0pt}%
\pgfsys@defobject{currentmarker}{\pgfqpoint{-0.048611in}{0.000000in}}{\pgfqpoint{-0.000000in}{0.000000in}}{%
\pgfpathmoveto{\pgfqpoint{-0.000000in}{0.000000in}}%
\pgfpathlineto{\pgfqpoint{-0.048611in}{0.000000in}}%
\pgfusepath{stroke,fill}%
}%
\begin{pgfscope}%
\pgfsys@transformshift{1.250000in}{5.070000in}%
\pgfsys@useobject{currentmarker}{}%
\end{pgfscope}%
\end{pgfscope}%
\begin{pgfscope}%
\definecolor{textcolor}{rgb}{0.000000,0.000000,0.000000}%
\pgfsetstrokecolor{textcolor}%
\pgfsetfillcolor{textcolor}%
\pgftext[x=0.867364in, y=4.986667in, left, base]{\color{textcolor}\rmfamily\fontsize{16.000000}{19.200000}\selectfont \(\displaystyle {1.0}\)}%
\end{pgfscope}%
\begin{pgfscope}%
\definecolor{textcolor}{rgb}{0.000000,0.000000,0.000000}%
\pgfsetstrokecolor{textcolor}%
\pgfsetfillcolor{textcolor}%
\pgftext[x=0.811809in,y=2.970000in,,bottom,rotate=90.000000]{\color{textcolor}\rmfamily\fontsize{16.000000}{19.200000}\selectfont \(\displaystyle F\)}%
\end{pgfscope}%
\begin{pgfscope}%
\pgfpathrectangle{\pgfqpoint{1.250000in}{0.660000in}}{\pgfqpoint{7.750000in}{4.620000in}}%
\pgfusepath{clip}%
\pgfsetrectcap%
\pgfsetroundjoin%
\pgfsetlinewidth{1.505625pt}%
\definecolor{currentstroke}{rgb}{0.121569,0.466667,0.705882}%
\pgfsetstrokecolor{currentstroke}%
\pgfsetdash{}{0pt}%
\pgfpathmoveto{\pgfqpoint{1.327500in}{5.070000in}}%
\pgfpathlineto{\pgfqpoint{1.405000in}{5.070000in}}%
\pgfpathlineto{\pgfqpoint{1.482500in}{4.495334in}}%
\pgfpathlineto{\pgfqpoint{1.560000in}{4.384150in}}%
\pgfpathlineto{\pgfqpoint{1.637500in}{4.641142in}}%
\pgfpathlineto{\pgfqpoint{1.715000in}{4.536356in}}%
\pgfpathlineto{\pgfqpoint{1.792500in}{4.999391in}}%
\pgfpathlineto{\pgfqpoint{1.870000in}{4.714679in}}%
\pgfpathlineto{\pgfqpoint{1.947500in}{4.957358in}}%
\pgfpathlineto{\pgfqpoint{2.025000in}{4.952211in}}%
\pgfpathlineto{\pgfqpoint{2.102500in}{4.802692in}}%
\pgfpathlineto{\pgfqpoint{2.180000in}{4.934627in}}%
\pgfpathlineto{\pgfqpoint{2.257500in}{4.997910in}}%
\pgfpathlineto{\pgfqpoint{2.335000in}{4.968111in}}%
\pgfpathlineto{\pgfqpoint{2.412500in}{4.881609in}}%
\pgfpathlineto{\pgfqpoint{2.490000in}{4.918921in}}%
\pgfpathlineto{\pgfqpoint{2.567500in}{4.808933in}}%
\pgfpathlineto{\pgfqpoint{2.645000in}{4.895462in}}%
\pgfpathlineto{\pgfqpoint{2.722500in}{4.913703in}}%
\pgfpathlineto{\pgfqpoint{2.800000in}{4.913883in}}%
\pgfpathlineto{\pgfqpoint{2.877500in}{4.684055in}}%
\pgfpathlineto{\pgfqpoint{2.955000in}{4.820131in}}%
\pgfpathlineto{\pgfqpoint{3.032500in}{4.433649in}}%
\pgfpathlineto{\pgfqpoint{3.110000in}{4.861506in}}%
\pgfpathlineto{\pgfqpoint{3.187500in}{4.977257in}}%
\pgfpathlineto{\pgfqpoint{3.265000in}{4.010425in}}%
\pgfpathlineto{\pgfqpoint{3.342500in}{4.777543in}}%
\pgfpathlineto{\pgfqpoint{3.420000in}{4.937447in}}%
\pgfpathlineto{\pgfqpoint{3.497500in}{4.858389in}}%
\pgfpathlineto{\pgfqpoint{3.575000in}{4.642638in}}%
\pgfpathlineto{\pgfqpoint{3.652500in}{4.730733in}}%
\pgfpathlineto{\pgfqpoint{3.730000in}{4.507285in}}%
\pgfpathlineto{\pgfqpoint{3.807500in}{4.605824in}}%
\pgfpathlineto{\pgfqpoint{3.885000in}{4.812733in}}%
\pgfpathlineto{\pgfqpoint{3.962500in}{4.802531in}}%
\pgfpathlineto{\pgfqpoint{4.040000in}{4.496048in}}%
\pgfpathlineto{\pgfqpoint{4.117500in}{4.552039in}}%
\pgfpathlineto{\pgfqpoint{4.195000in}{4.851376in}}%
\pgfpathlineto{\pgfqpoint{4.272500in}{4.636882in}}%
\pgfpathlineto{\pgfqpoint{4.350000in}{4.787352in}}%
\pgfpathlineto{\pgfqpoint{4.427500in}{4.726728in}}%
\pgfpathlineto{\pgfqpoint{4.505000in}{4.569415in}}%
\pgfpathlineto{\pgfqpoint{4.582500in}{4.692124in}}%
\pgfpathlineto{\pgfqpoint{4.660000in}{4.640270in}}%
\pgfpathlineto{\pgfqpoint{4.737500in}{4.585439in}}%
\pgfpathlineto{\pgfqpoint{4.815000in}{4.673639in}}%
\pgfpathlineto{\pgfqpoint{4.892500in}{4.507307in}}%
\pgfpathlineto{\pgfqpoint{4.970000in}{4.524629in}}%
\pgfpathlineto{\pgfqpoint{5.047500in}{4.797864in}}%
\pgfpathlineto{\pgfqpoint{5.125000in}{4.728121in}}%
\pgfpathlineto{\pgfqpoint{5.202500in}{4.791929in}}%
\pgfpathlineto{\pgfqpoint{5.280000in}{4.839548in}}%
\pgfpathlineto{\pgfqpoint{5.357500in}{4.856287in}}%
\pgfpathlineto{\pgfqpoint{5.435000in}{4.841653in}}%
\pgfpathlineto{\pgfqpoint{5.512500in}{4.805295in}}%
\pgfpathlineto{\pgfqpoint{5.590000in}{4.746139in}}%
\pgfpathlineto{\pgfqpoint{5.667500in}{4.825282in}}%
\pgfpathlineto{\pgfqpoint{5.745000in}{4.725330in}}%
\pgfpathlineto{\pgfqpoint{5.822500in}{4.437705in}}%
\pgfpathlineto{\pgfqpoint{5.900000in}{4.300627in}}%
\pgfpathlineto{\pgfqpoint{5.977500in}{4.517415in}}%
\pgfpathlineto{\pgfqpoint{6.055000in}{4.335095in}}%
\pgfpathlineto{\pgfqpoint{6.132500in}{4.374000in}}%
\pgfpathlineto{\pgfqpoint{6.210000in}{4.370041in}}%
\pgfpathlineto{\pgfqpoint{6.287500in}{4.262006in}}%
\pgfpathlineto{\pgfqpoint{6.365000in}{4.645858in}}%
\pgfpathlineto{\pgfqpoint{6.442500in}{4.425491in}}%
\pgfpathlineto{\pgfqpoint{6.520000in}{4.453111in}}%
\pgfpathlineto{\pgfqpoint{6.597500in}{4.574808in}}%
\pgfpathlineto{\pgfqpoint{6.675000in}{4.627087in}}%
\pgfpathlineto{\pgfqpoint{6.752500in}{4.126961in}}%
\pgfpathlineto{\pgfqpoint{6.830000in}{4.627621in}}%
\pgfpathlineto{\pgfqpoint{6.907500in}{4.502712in}}%
\pgfpathlineto{\pgfqpoint{6.985000in}{4.336020in}}%
\pgfpathlineto{\pgfqpoint{7.062500in}{4.521670in}}%
\pgfpathlineto{\pgfqpoint{7.140000in}{4.551736in}}%
\pgfpathlineto{\pgfqpoint{7.217500in}{4.694508in}}%
\pgfpathlineto{\pgfqpoint{7.295000in}{4.545164in}}%
\pgfpathlineto{\pgfqpoint{7.372500in}{4.545006in}}%
\pgfpathlineto{\pgfqpoint{7.450000in}{4.573549in}}%
\pgfpathlineto{\pgfqpoint{7.527500in}{4.440835in}}%
\pgfpathlineto{\pgfqpoint{7.605000in}{4.437643in}}%
\pgfpathlineto{\pgfqpoint{7.682500in}{4.325783in}}%
\pgfpathlineto{\pgfqpoint{7.760000in}{4.422493in}}%
\pgfpathlineto{\pgfqpoint{7.837500in}{4.506211in}}%
\pgfpathlineto{\pgfqpoint{7.915000in}{4.603280in}}%
\pgfpathlineto{\pgfqpoint{7.992500in}{4.633610in}}%
\pgfpathlineto{\pgfqpoint{8.070000in}{4.679332in}}%
\pgfpathlineto{\pgfqpoint{8.147500in}{4.710779in}}%
\pgfpathlineto{\pgfqpoint{8.225000in}{4.728041in}}%
\pgfpathlineto{\pgfqpoint{8.302500in}{4.736082in}}%
\pgfpathlineto{\pgfqpoint{8.380000in}{4.731090in}}%
\pgfpathlineto{\pgfqpoint{8.457500in}{4.713962in}}%
\pgfpathlineto{\pgfqpoint{8.535000in}{4.687296in}}%
\pgfpathlineto{\pgfqpoint{8.612500in}{4.651413in}}%
\pgfpathlineto{\pgfqpoint{8.690000in}{4.605472in}}%
\pgfpathlineto{\pgfqpoint{8.767500in}{4.550908in}}%
\pgfpathlineto{\pgfqpoint{8.845000in}{4.490988in}}%
\pgfpathlineto{\pgfqpoint{8.922500in}{4.423962in}}%
\pgfpathlineto{\pgfqpoint{9.000000in}{4.406296in}}%
\pgfpathlineto{\pgfqpoint{9.010000in}{4.397072in}}%
\pgfpathlineto{\pgfqpoint{9.010000in}{4.397072in}}%
\pgfusepath{stroke}%
\end{pgfscope}%
\begin{pgfscope}%
\pgfpathrectangle{\pgfqpoint{1.250000in}{0.660000in}}{\pgfqpoint{7.750000in}{4.620000in}}%
\pgfusepath{clip}%
\pgfsetrectcap%
\pgfsetroundjoin%
\pgfsetlinewidth{1.505625pt}%
\definecolor{currentstroke}{rgb}{0.682353,0.780392,0.909804}%
\pgfsetstrokecolor{currentstroke}%
\pgfsetdash{}{0pt}%
\pgfpathmoveto{\pgfqpoint{1.327500in}{5.070000in}}%
\pgfpathlineto{\pgfqpoint{1.482500in}{5.070000in}}%
\pgfpathlineto{\pgfqpoint{1.560000in}{4.985768in}}%
\pgfpathlineto{\pgfqpoint{1.637500in}{5.069999in}}%
\pgfpathlineto{\pgfqpoint{1.715000in}{5.052834in}}%
\pgfpathlineto{\pgfqpoint{1.792500in}{5.022366in}}%
\pgfpathlineto{\pgfqpoint{1.870000in}{5.026624in}}%
\pgfpathlineto{\pgfqpoint{1.947500in}{4.985798in}}%
\pgfpathlineto{\pgfqpoint{2.025000in}{5.020704in}}%
\pgfpathlineto{\pgfqpoint{2.102500in}{5.011594in}}%
\pgfpathlineto{\pgfqpoint{2.180000in}{5.059735in}}%
\pgfpathlineto{\pgfqpoint{2.257500in}{5.023927in}}%
\pgfpathlineto{\pgfqpoint{2.335000in}{5.003145in}}%
\pgfpathlineto{\pgfqpoint{2.412500in}{4.894842in}}%
\pgfpathlineto{\pgfqpoint{2.490000in}{4.990672in}}%
\pgfpathlineto{\pgfqpoint{2.567500in}{5.023505in}}%
\pgfpathlineto{\pgfqpoint{2.645000in}{5.013629in}}%
\pgfpathlineto{\pgfqpoint{2.722500in}{4.931994in}}%
\pgfpathlineto{\pgfqpoint{2.800000in}{5.017697in}}%
\pgfpathlineto{\pgfqpoint{2.877500in}{4.946169in}}%
\pgfpathlineto{\pgfqpoint{2.955000in}{4.970745in}}%
\pgfpathlineto{\pgfqpoint{3.032500in}{4.858460in}}%
\pgfpathlineto{\pgfqpoint{3.110000in}{4.958328in}}%
\pgfpathlineto{\pgfqpoint{3.187500in}{4.897330in}}%
\pgfpathlineto{\pgfqpoint{3.265000in}{4.849537in}}%
\pgfpathlineto{\pgfqpoint{3.342500in}{4.855746in}}%
\pgfpathlineto{\pgfqpoint{3.420000in}{4.798627in}}%
\pgfpathlineto{\pgfqpoint{3.497500in}{4.895126in}}%
\pgfpathlineto{\pgfqpoint{3.575000in}{4.868134in}}%
\pgfpathlineto{\pgfqpoint{3.652500in}{4.947851in}}%
\pgfpathlineto{\pgfqpoint{3.730000in}{4.595004in}}%
\pgfpathlineto{\pgfqpoint{3.807500in}{4.787300in}}%
\pgfpathlineto{\pgfqpoint{3.885000in}{4.697201in}}%
\pgfpathlineto{\pgfqpoint{3.962500in}{4.711369in}}%
\pgfpathlineto{\pgfqpoint{4.040000in}{4.723676in}}%
\pgfpathlineto{\pgfqpoint{4.117500in}{4.851313in}}%
\pgfpathlineto{\pgfqpoint{4.195000in}{4.655819in}}%
\pgfpathlineto{\pgfqpoint{4.272500in}{4.936523in}}%
\pgfpathlineto{\pgfqpoint{4.350000in}{4.796683in}}%
\pgfpathlineto{\pgfqpoint{4.427500in}{4.777525in}}%
\pgfpathlineto{\pgfqpoint{4.505000in}{4.739621in}}%
\pgfpathlineto{\pgfqpoint{4.582500in}{4.790155in}}%
\pgfpathlineto{\pgfqpoint{4.660000in}{4.513844in}}%
\pgfpathlineto{\pgfqpoint{4.737500in}{4.670733in}}%
\pgfpathlineto{\pgfqpoint{4.815000in}{4.583569in}}%
\pgfpathlineto{\pgfqpoint{4.892500in}{4.647257in}}%
\pgfpathlineto{\pgfqpoint{4.970000in}{4.477202in}}%
\pgfpathlineto{\pgfqpoint{5.047500in}{4.647095in}}%
\pgfpathlineto{\pgfqpoint{5.125000in}{4.637243in}}%
\pgfpathlineto{\pgfqpoint{5.202500in}{4.671980in}}%
\pgfpathlineto{\pgfqpoint{5.280000in}{4.522434in}}%
\pgfpathlineto{\pgfqpoint{5.357500in}{4.446870in}}%
\pgfpathlineto{\pgfqpoint{5.435000in}{4.530284in}}%
\pgfpathlineto{\pgfqpoint{5.512500in}{4.599985in}}%
\pgfpathlineto{\pgfqpoint{5.590000in}{4.555846in}}%
\pgfpathlineto{\pgfqpoint{5.667500in}{4.621750in}}%
\pgfpathlineto{\pgfqpoint{5.745000in}{4.596490in}}%
\pgfpathlineto{\pgfqpoint{5.822500in}{4.518241in}}%
\pgfpathlineto{\pgfqpoint{5.900000in}{4.408548in}}%
\pgfpathlineto{\pgfqpoint{5.977500in}{4.587564in}}%
\pgfpathlineto{\pgfqpoint{6.055000in}{4.435578in}}%
\pgfpathlineto{\pgfqpoint{6.132500in}{4.521746in}}%
\pgfpathlineto{\pgfqpoint{6.210000in}{4.427454in}}%
\pgfpathlineto{\pgfqpoint{6.287500in}{4.728690in}}%
\pgfpathlineto{\pgfqpoint{6.365000in}{4.590334in}}%
\pgfpathlineto{\pgfqpoint{6.442500in}{4.423776in}}%
\pgfpathlineto{\pgfqpoint{6.520000in}{4.436015in}}%
\pgfpathlineto{\pgfqpoint{6.597500in}{4.374270in}}%
\pgfpathlineto{\pgfqpoint{6.675000in}{4.457334in}}%
\pgfpathlineto{\pgfqpoint{6.752500in}{4.397499in}}%
\pgfpathlineto{\pgfqpoint{6.830000in}{4.417927in}}%
\pgfpathlineto{\pgfqpoint{6.907500in}{4.464825in}}%
\pgfpathlineto{\pgfqpoint{6.985000in}{4.459306in}}%
\pgfpathlineto{\pgfqpoint{7.062500in}{4.391438in}}%
\pgfpathlineto{\pgfqpoint{7.140000in}{4.439622in}}%
\pgfpathlineto{\pgfqpoint{7.217500in}{4.382886in}}%
\pgfpathlineto{\pgfqpoint{7.295000in}{4.379134in}}%
\pgfpathlineto{\pgfqpoint{7.372500in}{4.358716in}}%
\pgfpathlineto{\pgfqpoint{7.450000in}{4.419961in}}%
\pgfpathlineto{\pgfqpoint{7.527500in}{4.424660in}}%
\pgfpathlineto{\pgfqpoint{7.605000in}{4.499672in}}%
\pgfpathlineto{\pgfqpoint{7.682500in}{4.458696in}}%
\pgfpathlineto{\pgfqpoint{7.760000in}{4.409486in}}%
\pgfpathlineto{\pgfqpoint{7.837500in}{4.433080in}}%
\pgfpathlineto{\pgfqpoint{7.915000in}{4.392752in}}%
\pgfpathlineto{\pgfqpoint{7.992500in}{4.343780in}}%
\pgfpathlineto{\pgfqpoint{8.147500in}{4.381260in}}%
\pgfpathlineto{\pgfqpoint{8.225000in}{4.440587in}}%
\pgfpathlineto{\pgfqpoint{8.302500in}{4.339205in}}%
\pgfpathlineto{\pgfqpoint{8.380000in}{4.470263in}}%
\pgfpathlineto{\pgfqpoint{8.457500in}{4.385410in}}%
\pgfpathlineto{\pgfqpoint{8.535000in}{4.383250in}}%
\pgfpathlineto{\pgfqpoint{8.612500in}{4.424943in}}%
\pgfpathlineto{\pgfqpoint{8.690000in}{4.355928in}}%
\pgfpathlineto{\pgfqpoint{8.767500in}{4.378419in}}%
\pgfpathlineto{\pgfqpoint{8.845000in}{4.323071in}}%
\pgfpathlineto{\pgfqpoint{8.922500in}{4.402181in}}%
\pgfpathlineto{\pgfqpoint{9.000000in}{4.453005in}}%
\pgfpathlineto{\pgfqpoint{9.010000in}{4.433372in}}%
\pgfpathlineto{\pgfqpoint{9.010000in}{4.433372in}}%
\pgfusepath{stroke}%
\end{pgfscope}%
\begin{pgfscope}%
\pgfpathrectangle{\pgfqpoint{1.250000in}{0.660000in}}{\pgfqpoint{7.750000in}{4.620000in}}%
\pgfusepath{clip}%
\pgfsetrectcap%
\pgfsetroundjoin%
\pgfsetlinewidth{1.505625pt}%
\definecolor{currentstroke}{rgb}{1.000000,0.498039,0.054902}%
\pgfsetstrokecolor{currentstroke}%
\pgfsetdash{}{0pt}%
\pgfpathmoveto{\pgfqpoint{1.327500in}{5.070000in}}%
\pgfpathlineto{\pgfqpoint{1.482500in}{5.070000in}}%
\pgfpathlineto{\pgfqpoint{1.560000in}{5.060535in}}%
\pgfpathlineto{\pgfqpoint{1.637500in}{5.069977in}}%
\pgfpathlineto{\pgfqpoint{1.715000in}{5.062320in}}%
\pgfpathlineto{\pgfqpoint{1.792500in}{5.048246in}}%
\pgfpathlineto{\pgfqpoint{1.870000in}{5.063522in}}%
\pgfpathlineto{\pgfqpoint{1.947500in}{5.054168in}}%
\pgfpathlineto{\pgfqpoint{2.025000in}{5.047718in}}%
\pgfpathlineto{\pgfqpoint{2.102500in}{5.045019in}}%
\pgfpathlineto{\pgfqpoint{2.180000in}{5.018641in}}%
\pgfpathlineto{\pgfqpoint{2.257500in}{5.030410in}}%
\pgfpathlineto{\pgfqpoint{2.335000in}{5.000251in}}%
\pgfpathlineto{\pgfqpoint{2.412500in}{5.021511in}}%
\pgfpathlineto{\pgfqpoint{2.490000in}{4.965329in}}%
\pgfpathlineto{\pgfqpoint{2.567500in}{4.990261in}}%
\pgfpathlineto{\pgfqpoint{2.645000in}{4.998684in}}%
\pgfpathlineto{\pgfqpoint{2.722500in}{5.016287in}}%
\pgfpathlineto{\pgfqpoint{2.800000in}{4.986583in}}%
\pgfpathlineto{\pgfqpoint{2.877500in}{5.007045in}}%
\pgfpathlineto{\pgfqpoint{2.955000in}{5.008857in}}%
\pgfpathlineto{\pgfqpoint{3.032500in}{4.966690in}}%
\pgfpathlineto{\pgfqpoint{3.110000in}{4.947510in}}%
\pgfpathlineto{\pgfqpoint{3.187500in}{4.936948in}}%
\pgfpathlineto{\pgfqpoint{3.265000in}{4.930606in}}%
\pgfpathlineto{\pgfqpoint{3.342500in}{4.953072in}}%
\pgfpathlineto{\pgfqpoint{3.420000in}{4.931750in}}%
\pgfpathlineto{\pgfqpoint{3.497500in}{4.936614in}}%
\pgfpathlineto{\pgfqpoint{3.575000in}{4.927891in}}%
\pgfpathlineto{\pgfqpoint{3.652500in}{4.944542in}}%
\pgfpathlineto{\pgfqpoint{3.730000in}{4.827406in}}%
\pgfpathlineto{\pgfqpoint{3.807500in}{4.849392in}}%
\pgfpathlineto{\pgfqpoint{3.885000in}{4.951624in}}%
\pgfpathlineto{\pgfqpoint{3.962500in}{4.836948in}}%
\pgfpathlineto{\pgfqpoint{4.040000in}{4.859226in}}%
\pgfpathlineto{\pgfqpoint{4.117500in}{4.864436in}}%
\pgfpathlineto{\pgfqpoint{4.195000in}{4.808356in}}%
\pgfpathlineto{\pgfqpoint{4.272500in}{4.773115in}}%
\pgfpathlineto{\pgfqpoint{4.350000in}{4.800266in}}%
\pgfpathlineto{\pgfqpoint{4.427500in}{4.800867in}}%
\pgfpathlineto{\pgfqpoint{4.505000in}{4.807029in}}%
\pgfpathlineto{\pgfqpoint{4.582500in}{4.744115in}}%
\pgfpathlineto{\pgfqpoint{4.660000in}{4.758680in}}%
\pgfpathlineto{\pgfqpoint{4.737500in}{4.706181in}}%
\pgfpathlineto{\pgfqpoint{4.815000in}{4.817217in}}%
\pgfpathlineto{\pgfqpoint{4.892500in}{4.767322in}}%
\pgfpathlineto{\pgfqpoint{4.970000in}{4.772012in}}%
\pgfpathlineto{\pgfqpoint{5.047500in}{4.733095in}}%
\pgfpathlineto{\pgfqpoint{5.125000in}{4.715326in}}%
\pgfpathlineto{\pgfqpoint{5.202500in}{4.635407in}}%
\pgfpathlineto{\pgfqpoint{5.280000in}{4.665460in}}%
\pgfpathlineto{\pgfqpoint{5.357500in}{4.658211in}}%
\pgfpathlineto{\pgfqpoint{5.435000in}{4.790282in}}%
\pgfpathlineto{\pgfqpoint{5.512500in}{4.556119in}}%
\pgfpathlineto{\pgfqpoint{5.590000in}{4.485148in}}%
\pgfpathlineto{\pgfqpoint{5.667500in}{4.604758in}}%
\pgfpathlineto{\pgfqpoint{5.745000in}{4.567058in}}%
\pgfpathlineto{\pgfqpoint{5.822500in}{4.581036in}}%
\pgfpathlineto{\pgfqpoint{5.900000in}{4.484449in}}%
\pgfpathlineto{\pgfqpoint{5.977500in}{4.394209in}}%
\pgfpathlineto{\pgfqpoint{6.055000in}{4.728578in}}%
\pgfpathlineto{\pgfqpoint{6.132500in}{4.489850in}}%
\pgfpathlineto{\pgfqpoint{6.210000in}{4.640900in}}%
\pgfpathlineto{\pgfqpoint{6.287500in}{4.465152in}}%
\pgfpathlineto{\pgfqpoint{6.365000in}{4.472211in}}%
\pgfpathlineto{\pgfqpoint{6.442500in}{4.323200in}}%
\pgfpathlineto{\pgfqpoint{6.520000in}{4.517139in}}%
\pgfpathlineto{\pgfqpoint{6.597500in}{4.415420in}}%
\pgfpathlineto{\pgfqpoint{6.675000in}{4.510575in}}%
\pgfpathlineto{\pgfqpoint{6.752500in}{4.348814in}}%
\pgfpathlineto{\pgfqpoint{6.830000in}{4.550212in}}%
\pgfpathlineto{\pgfqpoint{6.907500in}{4.392906in}}%
\pgfpathlineto{\pgfqpoint{6.985000in}{4.270727in}}%
\pgfpathlineto{\pgfqpoint{7.062500in}{4.360211in}}%
\pgfpathlineto{\pgfqpoint{7.140000in}{4.514713in}}%
\pgfpathlineto{\pgfqpoint{7.217500in}{4.261735in}}%
\pgfpathlineto{\pgfqpoint{7.295000in}{4.231595in}}%
\pgfpathlineto{\pgfqpoint{7.372500in}{4.157081in}}%
\pgfpathlineto{\pgfqpoint{7.450000in}{4.237217in}}%
\pgfpathlineto{\pgfqpoint{7.527500in}{4.324024in}}%
\pgfpathlineto{\pgfqpoint{7.605000in}{4.281648in}}%
\pgfpathlineto{\pgfqpoint{7.682500in}{4.248799in}}%
\pgfpathlineto{\pgfqpoint{7.760000in}{4.435732in}}%
\pgfpathlineto{\pgfqpoint{7.837500in}{4.282012in}}%
\pgfpathlineto{\pgfqpoint{7.915000in}{4.280357in}}%
\pgfpathlineto{\pgfqpoint{7.992500in}{4.305369in}}%
\pgfpathlineto{\pgfqpoint{8.070000in}{4.160034in}}%
\pgfpathlineto{\pgfqpoint{8.147500in}{4.131150in}}%
\pgfpathlineto{\pgfqpoint{8.225000in}{4.122561in}}%
\pgfpathlineto{\pgfqpoint{8.380000in}{4.351452in}}%
\pgfpathlineto{\pgfqpoint{8.457500in}{4.298026in}}%
\pgfpathlineto{\pgfqpoint{8.535000in}{4.206534in}}%
\pgfpathlineto{\pgfqpoint{8.612500in}{4.248057in}}%
\pgfpathlineto{\pgfqpoint{8.690000in}{4.306155in}}%
\pgfpathlineto{\pgfqpoint{8.767500in}{4.273981in}}%
\pgfpathlineto{\pgfqpoint{8.845000in}{4.202885in}}%
\pgfpathlineto{\pgfqpoint{8.922500in}{4.033282in}}%
\pgfpathlineto{\pgfqpoint{9.000000in}{4.364372in}}%
\pgfpathlineto{\pgfqpoint{9.010000in}{4.323363in}}%
\pgfpathlineto{\pgfqpoint{9.010000in}{4.323363in}}%
\pgfusepath{stroke}%
\end{pgfscope}%
\begin{pgfscope}%
\pgfpathrectangle{\pgfqpoint{1.250000in}{0.660000in}}{\pgfqpoint{7.750000in}{4.620000in}}%
\pgfusepath{clip}%
\pgfsetrectcap%
\pgfsetroundjoin%
\pgfsetlinewidth{1.505625pt}%
\definecolor{currentstroke}{rgb}{1.000000,0.733333,0.470588}%
\pgfsetstrokecolor{currentstroke}%
\pgfsetdash{}{0pt}%
\pgfpathmoveto{\pgfqpoint{1.327500in}{5.070000in}}%
\pgfpathlineto{\pgfqpoint{1.560000in}{5.069863in}}%
\pgfpathlineto{\pgfqpoint{1.637500in}{5.064206in}}%
\pgfpathlineto{\pgfqpoint{1.715000in}{5.067510in}}%
\pgfpathlineto{\pgfqpoint{1.792500in}{5.066923in}}%
\pgfpathlineto{\pgfqpoint{1.870000in}{5.063300in}}%
\pgfpathlineto{\pgfqpoint{1.947500in}{5.064401in}}%
\pgfpathlineto{\pgfqpoint{2.025000in}{5.061957in}}%
\pgfpathlineto{\pgfqpoint{2.102500in}{5.057169in}}%
\pgfpathlineto{\pgfqpoint{2.180000in}{5.062528in}}%
\pgfpathlineto{\pgfqpoint{2.335000in}{5.062401in}}%
\pgfpathlineto{\pgfqpoint{2.412500in}{5.036085in}}%
\pgfpathlineto{\pgfqpoint{2.490000in}{5.049312in}}%
\pgfpathlineto{\pgfqpoint{2.567500in}{5.036394in}}%
\pgfpathlineto{\pgfqpoint{2.645000in}{5.046404in}}%
\pgfpathlineto{\pgfqpoint{2.722500in}{5.016584in}}%
\pgfpathlineto{\pgfqpoint{2.800000in}{5.039600in}}%
\pgfpathlineto{\pgfqpoint{2.877500in}{5.045483in}}%
\pgfpathlineto{\pgfqpoint{2.955000in}{5.053560in}}%
\pgfpathlineto{\pgfqpoint{3.032500in}{5.028073in}}%
\pgfpathlineto{\pgfqpoint{3.110000in}{5.036165in}}%
\pgfpathlineto{\pgfqpoint{3.187500in}{5.026727in}}%
\pgfpathlineto{\pgfqpoint{3.265000in}{5.044868in}}%
\pgfpathlineto{\pgfqpoint{3.342500in}{5.025067in}}%
\pgfpathlineto{\pgfqpoint{3.420000in}{5.029343in}}%
\pgfpathlineto{\pgfqpoint{3.497500in}{5.025249in}}%
\pgfpathlineto{\pgfqpoint{3.575000in}{5.013877in}}%
\pgfpathlineto{\pgfqpoint{3.652500in}{4.994175in}}%
\pgfpathlineto{\pgfqpoint{3.730000in}{5.020317in}}%
\pgfpathlineto{\pgfqpoint{3.807500in}{5.011835in}}%
\pgfpathlineto{\pgfqpoint{3.885000in}{4.985497in}}%
\pgfpathlineto{\pgfqpoint{3.962500in}{5.004395in}}%
\pgfpathlineto{\pgfqpoint{4.040000in}{5.028380in}}%
\pgfpathlineto{\pgfqpoint{4.117500in}{4.992472in}}%
\pgfpathlineto{\pgfqpoint{4.195000in}{4.991872in}}%
\pgfpathlineto{\pgfqpoint{4.272500in}{4.986076in}}%
\pgfpathlineto{\pgfqpoint{4.350000in}{5.017799in}}%
\pgfpathlineto{\pgfqpoint{4.427500in}{4.981573in}}%
\pgfpathlineto{\pgfqpoint{4.505000in}{4.965449in}}%
\pgfpathlineto{\pgfqpoint{4.582500in}{4.941687in}}%
\pgfpathlineto{\pgfqpoint{4.660000in}{5.019096in}}%
\pgfpathlineto{\pgfqpoint{4.737500in}{5.000281in}}%
\pgfpathlineto{\pgfqpoint{4.815000in}{4.917825in}}%
\pgfpathlineto{\pgfqpoint{4.892500in}{4.934148in}}%
\pgfpathlineto{\pgfqpoint{4.970000in}{4.910332in}}%
\pgfpathlineto{\pgfqpoint{5.047500in}{4.962528in}}%
\pgfpathlineto{\pgfqpoint{5.125000in}{4.943428in}}%
\pgfpathlineto{\pgfqpoint{5.202500in}{4.981752in}}%
\pgfpathlineto{\pgfqpoint{5.280000in}{4.951026in}}%
\pgfpathlineto{\pgfqpoint{5.357500in}{4.942313in}}%
\pgfpathlineto{\pgfqpoint{5.435000in}{4.866290in}}%
\pgfpathlineto{\pgfqpoint{5.512500in}{4.952091in}}%
\pgfpathlineto{\pgfqpoint{5.590000in}{4.896865in}}%
\pgfpathlineto{\pgfqpoint{5.667500in}{4.866163in}}%
\pgfpathlineto{\pgfqpoint{5.745000in}{4.870742in}}%
\pgfpathlineto{\pgfqpoint{5.822500in}{4.929975in}}%
\pgfpathlineto{\pgfqpoint{5.900000in}{4.908903in}}%
\pgfpathlineto{\pgfqpoint{5.977500in}{4.864087in}}%
\pgfpathlineto{\pgfqpoint{6.055000in}{4.900172in}}%
\pgfpathlineto{\pgfqpoint{6.132500in}{4.869685in}}%
\pgfpathlineto{\pgfqpoint{6.210000in}{4.793025in}}%
\pgfpathlineto{\pgfqpoint{6.287500in}{4.935747in}}%
\pgfpathlineto{\pgfqpoint{6.365000in}{4.820221in}}%
\pgfpathlineto{\pgfqpoint{6.442500in}{4.877370in}}%
\pgfpathlineto{\pgfqpoint{6.520000in}{4.885026in}}%
\pgfpathlineto{\pgfqpoint{6.597500in}{4.862544in}}%
\pgfpathlineto{\pgfqpoint{6.675000in}{4.905036in}}%
\pgfpathlineto{\pgfqpoint{6.752500in}{4.860417in}}%
\pgfpathlineto{\pgfqpoint{6.830000in}{4.846843in}}%
\pgfpathlineto{\pgfqpoint{6.907500in}{4.828058in}}%
\pgfpathlineto{\pgfqpoint{6.985000in}{4.814314in}}%
\pgfpathlineto{\pgfqpoint{7.062500in}{4.899748in}}%
\pgfpathlineto{\pgfqpoint{7.140000in}{4.855912in}}%
\pgfpathlineto{\pgfqpoint{7.217500in}{4.775582in}}%
\pgfpathlineto{\pgfqpoint{7.295000in}{4.832988in}}%
\pgfpathlineto{\pgfqpoint{7.372500in}{4.718363in}}%
\pgfpathlineto{\pgfqpoint{7.450000in}{4.734989in}}%
\pgfpathlineto{\pgfqpoint{7.527500in}{4.771742in}}%
\pgfpathlineto{\pgfqpoint{7.605000in}{4.728219in}}%
\pgfpathlineto{\pgfqpoint{7.682500in}{4.712625in}}%
\pgfpathlineto{\pgfqpoint{7.760000in}{4.754161in}}%
\pgfpathlineto{\pgfqpoint{7.837500in}{4.824662in}}%
\pgfpathlineto{\pgfqpoint{7.915000in}{4.687901in}}%
\pgfpathlineto{\pgfqpoint{7.992500in}{4.671005in}}%
\pgfpathlineto{\pgfqpoint{8.070000in}{4.716878in}}%
\pgfpathlineto{\pgfqpoint{8.147500in}{4.830220in}}%
\pgfpathlineto{\pgfqpoint{8.225000in}{4.671651in}}%
\pgfpathlineto{\pgfqpoint{8.302500in}{4.631018in}}%
\pgfpathlineto{\pgfqpoint{8.380000in}{4.745537in}}%
\pgfpathlineto{\pgfqpoint{8.457500in}{4.696850in}}%
\pgfpathlineto{\pgfqpoint{8.535000in}{4.729088in}}%
\pgfpathlineto{\pgfqpoint{8.612500in}{4.648300in}}%
\pgfpathlineto{\pgfqpoint{8.690000in}{4.720675in}}%
\pgfpathlineto{\pgfqpoint{8.767500in}{4.731687in}}%
\pgfpathlineto{\pgfqpoint{8.845000in}{4.679137in}}%
\pgfpathlineto{\pgfqpoint{8.922500in}{4.704948in}}%
\pgfpathlineto{\pgfqpoint{9.000000in}{4.817052in}}%
\pgfpathlineto{\pgfqpoint{9.010000in}{4.801245in}}%
\pgfpathlineto{\pgfqpoint{9.010000in}{4.801245in}}%
\pgfusepath{stroke}%
\end{pgfscope}%
\begin{pgfscope}%
\pgfpathrectangle{\pgfqpoint{1.250000in}{0.660000in}}{\pgfqpoint{7.750000in}{4.620000in}}%
\pgfusepath{clip}%
\pgfsetrectcap%
\pgfsetroundjoin%
\pgfsetlinewidth{1.505625pt}%
\definecolor{currentstroke}{rgb}{0.172549,0.627451,0.172549}%
\pgfsetstrokecolor{currentstroke}%
\pgfsetdash{}{0pt}%
\pgfpathmoveto{\pgfqpoint{1.327500in}{5.070000in}}%
\pgfpathlineto{\pgfqpoint{1.560000in}{5.069820in}}%
\pgfpathlineto{\pgfqpoint{1.637500in}{5.064816in}}%
\pgfpathlineto{\pgfqpoint{1.715000in}{5.065991in}}%
\pgfpathlineto{\pgfqpoint{1.792500in}{5.057176in}}%
\pgfpathlineto{\pgfqpoint{1.870000in}{5.067050in}}%
\pgfpathlineto{\pgfqpoint{1.947500in}{5.068730in}}%
\pgfpathlineto{\pgfqpoint{2.025000in}{5.068845in}}%
\pgfpathlineto{\pgfqpoint{2.102500in}{5.066984in}}%
\pgfpathlineto{\pgfqpoint{2.257500in}{5.069143in}}%
\pgfpathlineto{\pgfqpoint{2.335000in}{5.066368in}}%
\pgfpathlineto{\pgfqpoint{2.412500in}{5.069514in}}%
\pgfpathlineto{\pgfqpoint{2.567500in}{5.068138in}}%
\pgfpathlineto{\pgfqpoint{2.645000in}{5.069568in}}%
\pgfpathlineto{\pgfqpoint{2.722500in}{5.068120in}}%
\pgfpathlineto{\pgfqpoint{2.877500in}{5.068710in}}%
\pgfpathlineto{\pgfqpoint{2.955000in}{5.061881in}}%
\pgfpathlineto{\pgfqpoint{3.032500in}{5.068289in}}%
\pgfpathlineto{\pgfqpoint{3.110000in}{5.068885in}}%
\pgfpathlineto{\pgfqpoint{3.187500in}{5.067484in}}%
\pgfpathlineto{\pgfqpoint{3.265000in}{5.069152in}}%
\pgfpathlineto{\pgfqpoint{3.342500in}{5.068339in}}%
\pgfpathlineto{\pgfqpoint{3.497500in}{5.068579in}}%
\pgfpathlineto{\pgfqpoint{3.652500in}{5.068822in}}%
\pgfpathlineto{\pgfqpoint{3.730000in}{5.066845in}}%
\pgfpathlineto{\pgfqpoint{3.885000in}{5.067671in}}%
\pgfpathlineto{\pgfqpoint{3.962500in}{5.066746in}}%
\pgfpathlineto{\pgfqpoint{4.040000in}{5.068605in}}%
\pgfpathlineto{\pgfqpoint{4.195000in}{5.067876in}}%
\pgfpathlineto{\pgfqpoint{4.272500in}{5.069077in}}%
\pgfpathlineto{\pgfqpoint{4.350000in}{5.067845in}}%
\pgfpathlineto{\pgfqpoint{4.505000in}{5.068102in}}%
\pgfpathlineto{\pgfqpoint{4.582500in}{5.067801in}}%
\pgfpathlineto{\pgfqpoint{4.660000in}{5.064892in}}%
\pgfpathlineto{\pgfqpoint{4.737500in}{5.064991in}}%
\pgfpathlineto{\pgfqpoint{4.815000in}{5.066372in}}%
\pgfpathlineto{\pgfqpoint{4.892500in}{5.066253in}}%
\pgfpathlineto{\pgfqpoint{4.970000in}{5.067277in}}%
\pgfpathlineto{\pgfqpoint{5.047500in}{5.066614in}}%
\pgfpathlineto{\pgfqpoint{5.125000in}{5.068676in}}%
\pgfpathlineto{\pgfqpoint{5.202500in}{5.065977in}}%
\pgfpathlineto{\pgfqpoint{5.357500in}{5.066739in}}%
\pgfpathlineto{\pgfqpoint{5.435000in}{5.064092in}}%
\pgfpathlineto{\pgfqpoint{5.512500in}{5.066679in}}%
\pgfpathlineto{\pgfqpoint{5.667500in}{5.065622in}}%
\pgfpathlineto{\pgfqpoint{5.745000in}{5.067948in}}%
\pgfpathlineto{\pgfqpoint{5.822500in}{5.064840in}}%
\pgfpathlineto{\pgfqpoint{6.287500in}{5.066194in}}%
\pgfpathlineto{\pgfqpoint{6.365000in}{5.064820in}}%
\pgfpathlineto{\pgfqpoint{6.442500in}{5.067238in}}%
\pgfpathlineto{\pgfqpoint{6.520000in}{5.067788in}}%
\pgfpathlineto{\pgfqpoint{6.597500in}{5.065726in}}%
\pgfpathlineto{\pgfqpoint{6.675000in}{5.067918in}}%
\pgfpathlineto{\pgfqpoint{6.752500in}{5.063519in}}%
\pgfpathlineto{\pgfqpoint{6.830000in}{5.064791in}}%
\pgfpathlineto{\pgfqpoint{6.985000in}{5.063760in}}%
\pgfpathlineto{\pgfqpoint{7.062500in}{5.066959in}}%
\pgfpathlineto{\pgfqpoint{7.140000in}{5.064109in}}%
\pgfpathlineto{\pgfqpoint{7.217500in}{5.066249in}}%
\pgfpathlineto{\pgfqpoint{7.295000in}{5.064401in}}%
\pgfpathlineto{\pgfqpoint{7.372500in}{5.067311in}}%
\pgfpathlineto{\pgfqpoint{7.527500in}{5.064266in}}%
\pgfpathlineto{\pgfqpoint{7.605000in}{5.066979in}}%
\pgfpathlineto{\pgfqpoint{7.682500in}{5.060001in}}%
\pgfpathlineto{\pgfqpoint{7.760000in}{5.063687in}}%
\pgfpathlineto{\pgfqpoint{7.837500in}{5.062250in}}%
\pgfpathlineto{\pgfqpoint{7.992500in}{5.064259in}}%
\pgfpathlineto{\pgfqpoint{8.070000in}{5.063617in}}%
\pgfpathlineto{\pgfqpoint{8.147500in}{5.065665in}}%
\pgfpathlineto{\pgfqpoint{8.302500in}{5.062094in}}%
\pgfpathlineto{\pgfqpoint{8.380000in}{5.063857in}}%
\pgfpathlineto{\pgfqpoint{8.457500in}{5.060813in}}%
\pgfpathlineto{\pgfqpoint{8.535000in}{5.062524in}}%
\pgfpathlineto{\pgfqpoint{8.690000in}{5.059029in}}%
\pgfpathlineto{\pgfqpoint{8.767500in}{5.063281in}}%
\pgfpathlineto{\pgfqpoint{8.922500in}{5.064685in}}%
\pgfpathlineto{\pgfqpoint{9.000000in}{5.055379in}}%
\pgfpathlineto{\pgfqpoint{9.010000in}{5.056631in}}%
\pgfpathlineto{\pgfqpoint{9.010000in}{5.056631in}}%
\pgfusepath{stroke}%
\end{pgfscope}%
\begin{pgfscope}%
\pgfsetrectcap%
\pgfsetmiterjoin%
\pgfsetlinewidth{0.803000pt}%
\definecolor{currentstroke}{rgb}{0.000000,0.000000,0.000000}%
\pgfsetstrokecolor{currentstroke}%
\pgfsetdash{}{0pt}%
\pgfpathmoveto{\pgfqpoint{1.250000in}{0.660000in}}%
\pgfpathlineto{\pgfqpoint{1.250000in}{5.280000in}}%
\pgfusepath{stroke}%
\end{pgfscope}%
\begin{pgfscope}%
\pgfsetrectcap%
\pgfsetmiterjoin%
\pgfsetlinewidth{0.803000pt}%
\definecolor{currentstroke}{rgb}{0.000000,0.000000,0.000000}%
\pgfsetstrokecolor{currentstroke}%
\pgfsetdash{}{0pt}%
\pgfpathmoveto{\pgfqpoint{9.000000in}{0.660000in}}%
\pgfpathlineto{\pgfqpoint{9.000000in}{5.280000in}}%
\pgfusepath{stroke}%
\end{pgfscope}%
\begin{pgfscope}%
\pgfsetrectcap%
\pgfsetmiterjoin%
\pgfsetlinewidth{0.803000pt}%
\definecolor{currentstroke}{rgb}{0.000000,0.000000,0.000000}%
\pgfsetstrokecolor{currentstroke}%
\pgfsetdash{}{0pt}%
\pgfpathmoveto{\pgfqpoint{1.250000in}{0.660000in}}%
\pgfpathlineto{\pgfqpoint{9.000000in}{0.660000in}}%
\pgfusepath{stroke}%
\end{pgfscope}%
\begin{pgfscope}%
\pgfsetrectcap%
\pgfsetmiterjoin%
\pgfsetlinewidth{0.803000pt}%
\definecolor{currentstroke}{rgb}{0.000000,0.000000,0.000000}%
\pgfsetstrokecolor{currentstroke}%
\pgfsetdash{}{0pt}%
\pgfpathmoveto{\pgfqpoint{1.250000in}{5.280000in}}%
\pgfpathlineto{\pgfqpoint{9.000000in}{5.280000in}}%
\pgfusepath{stroke}%
\end{pgfscope}%
\begin{pgfscope}%
\pgfsetbuttcap%
\pgfsetmiterjoin%
\definecolor{currentfill}{rgb}{1.000000,1.000000,1.000000}%
\pgfsetfillcolor{currentfill}%
\pgfsetfillopacity{0.800000}%
\pgfsetlinewidth{1.003750pt}%
\definecolor{currentstroke}{rgb}{0.800000,0.800000,0.800000}%
\pgfsetstrokecolor{currentstroke}%
\pgfsetstrokeopacity{0.800000}%
\pgfsetdash{}{0pt}%
\pgfpathmoveto{\pgfqpoint{1.405556in}{0.771111in}}%
\pgfpathlineto{\pgfqpoint{2.893414in}{0.771111in}}%
\pgfpathquadraticcurveto{\pgfqpoint{2.937859in}{0.771111in}}{\pgfqpoint{2.937859in}{0.815556in}}%
\pgfpathlineto{\pgfqpoint{2.937859in}{2.415632in}}%
\pgfpathquadraticcurveto{\pgfqpoint{2.937859in}{2.460076in}}{\pgfqpoint{2.893414in}{2.460076in}}%
\pgfpathlineto{\pgfqpoint{1.405556in}{2.460076in}}%
\pgfpathquadraticcurveto{\pgfqpoint{1.361111in}{2.460076in}}{\pgfqpoint{1.361111in}{2.415632in}}%
\pgfpathlineto{\pgfqpoint{1.361111in}{0.815556in}}%
\pgfpathquadraticcurveto{\pgfqpoint{1.361111in}{0.771111in}}{\pgfqpoint{1.405556in}{0.771111in}}%
\pgfpathlineto{\pgfqpoint{1.405556in}{0.771111in}}%
\pgfpathclose%
\pgfusepath{stroke,fill}%
\end{pgfscope}%
\begin{pgfscope}%
\pgfsetrectcap%
\pgfsetroundjoin%
\pgfsetlinewidth{1.505625pt}%
\definecolor{currentstroke}{rgb}{0.121569,0.466667,0.705882}%
\pgfsetstrokecolor{currentstroke}%
\pgfsetdash{}{0pt}%
\pgfpathmoveto{\pgfqpoint{1.450000in}{2.282299in}}%
\pgfpathlineto{\pgfqpoint{1.672222in}{2.282299in}}%
\pgfpathlineto{\pgfqpoint{1.894444in}{2.282299in}}%
\pgfusepath{stroke}%
\end{pgfscope}%
\begin{pgfscope}%
\definecolor{textcolor}{rgb}{0.000000,0.000000,0.000000}%
\pgfsetstrokecolor{textcolor}%
\pgfsetfillcolor{textcolor}%
\pgftext[x=2.072222in,y=2.204521in,left,base]{\color{textcolor}\rmfamily\fontsize{16.000000}{19.200000}\selectfont nu = 1}%
\end{pgfscope}%
\begin{pgfscope}%
\pgfsetrectcap%
\pgfsetroundjoin%
\pgfsetlinewidth{1.505625pt}%
\definecolor{currentstroke}{rgb}{0.682353,0.780392,0.909804}%
\pgfsetstrokecolor{currentstroke}%
\pgfsetdash{}{0pt}%
\pgfpathmoveto{\pgfqpoint{1.450000in}{1.957839in}}%
\pgfpathlineto{\pgfqpoint{1.672222in}{1.957839in}}%
\pgfpathlineto{\pgfqpoint{1.894444in}{1.957839in}}%
\pgfusepath{stroke}%
\end{pgfscope}%
\begin{pgfscope}%
\definecolor{textcolor}{rgb}{0.000000,0.000000,0.000000}%
\pgfsetstrokecolor{textcolor}%
\pgfsetfillcolor{textcolor}%
\pgftext[x=2.072222in,y=1.880061in,left,base]{\color{textcolor}\rmfamily\fontsize{16.000000}{19.200000}\selectfont nu = 2}%
\end{pgfscope}%
\begin{pgfscope}%
\pgfsetrectcap%
\pgfsetroundjoin%
\pgfsetlinewidth{1.505625pt}%
\definecolor{currentstroke}{rgb}{1.000000,0.498039,0.054902}%
\pgfsetstrokecolor{currentstroke}%
\pgfsetdash{}{0pt}%
\pgfpathmoveto{\pgfqpoint{1.450000in}{1.633379in}}%
\pgfpathlineto{\pgfqpoint{1.672222in}{1.633379in}}%
\pgfpathlineto{\pgfqpoint{1.894444in}{1.633379in}}%
\pgfusepath{stroke}%
\end{pgfscope}%
\begin{pgfscope}%
\definecolor{textcolor}{rgb}{0.000000,0.000000,0.000000}%
\pgfsetstrokecolor{textcolor}%
\pgfsetfillcolor{textcolor}%
\pgftext[x=2.072222in,y=1.555601in,left,base]{\color{textcolor}\rmfamily\fontsize{16.000000}{19.200000}\selectfont nu = 3}%
\end{pgfscope}%
\begin{pgfscope}%
\pgfsetrectcap%
\pgfsetroundjoin%
\pgfsetlinewidth{1.505625pt}%
\definecolor{currentstroke}{rgb}{1.000000,0.733333,0.470588}%
\pgfsetstrokecolor{currentstroke}%
\pgfsetdash{}{0pt}%
\pgfpathmoveto{\pgfqpoint{1.450000in}{1.308919in}}%
\pgfpathlineto{\pgfqpoint{1.672222in}{1.308919in}}%
\pgfpathlineto{\pgfqpoint{1.894444in}{1.308919in}}%
\pgfusepath{stroke}%
\end{pgfscope}%
\begin{pgfscope}%
\definecolor{textcolor}{rgb}{0.000000,0.000000,0.000000}%
\pgfsetstrokecolor{textcolor}%
\pgfsetfillcolor{textcolor}%
\pgftext[x=2.072222in,y=1.231142in,left,base]{\color{textcolor}\rmfamily\fontsize{16.000000}{19.200000}\selectfont nu = 5}%
\end{pgfscope}%
\begin{pgfscope}%
\pgfsetrectcap%
\pgfsetroundjoin%
\pgfsetlinewidth{1.505625pt}%
\definecolor{currentstroke}{rgb}{0.172549,0.627451,0.172549}%
\pgfsetstrokecolor{currentstroke}%
\pgfsetdash{}{0pt}%
\pgfpathmoveto{\pgfqpoint{1.450000in}{0.984460in}}%
\pgfpathlineto{\pgfqpoint{1.672222in}{0.984460in}}%
\pgfpathlineto{\pgfqpoint{1.894444in}{0.984460in}}%
\pgfusepath{stroke}%
\end{pgfscope}%
\begin{pgfscope}%
\definecolor{textcolor}{rgb}{0.000000,0.000000,0.000000}%
\pgfsetstrokecolor{textcolor}%
\pgfsetfillcolor{textcolor}%
\pgftext[x=2.072222in,y=0.906682in,left,base]{\color{textcolor}\rmfamily\fontsize{16.000000}{19.200000}\selectfont nu = 10}%
\end{pgfscope}%
\end{pgfpicture}%
\makeatother%
\endgroup%
}
    \caption*{Figure: Maximum values of fidelity in dual-hole Hamiltonian as a function of the number of sites $N$ for various value of $\nu$.}
\end{figure}

\end{document}